\documentclass[11pt,twocolumn,twoside]{book}

\usepackage[margin=0.75in,top=1in,bottom=1in,headheight=14pt]{geometry}
\usepackage{titlesec}
\usepackage{amsmath}
\usepackage{multicol}
\usepackage{enumitem}
\usepackage{hyperref}
\usepackage{fancyhdr}
\usepackage{booktabs}
\usepackage{tabularx}
\usepackage{microtype}
\usepackage[utf8]{inputenc}
\usepackage[T1]{fontenc}
\usepackage{lmodern}
\usepackage{xcolor}
\usepackage{textcomp}
\usepackage{sectsty}
\usepackage{parskip}
\usepackage{makeidx}
\usepackage{graphicx}
\makeindex

\pagestyle{fancy}
\fancyhf{}
\fancyhead[LE]{\textit{\leftmark}}
\fancyhead[RO]{\textit{\leftmark}}
\fancyfoot[LE,RO]{\thepage}
\renewcommand{\headrulewidth}{0.4pt}
\renewcommand{\footrulewidth}{0.4pt}

\hypersetup{
    colorlinks=true,
    linkcolor=blue!60!black,
    urlcolor=blue!60!black,
    citecolor=blue!60!black,
}

\chapterfont{\centering\LARGE\bfseries\color{blue!70!black}}
\sectionfont{\Large\bfseries\color{blue!70!black}}
\subsectionfont{\normalsize\bfseries\color{blue!50!black}}

\newcommand{\medterm}[1]{\hypertarget{#1}{}{\raggedright\noindent\textbf{\large\textcolor{red!70!black}{#1}}\index{#1}\par}\nobreak\vspace{1pt}\noindent\ignorespaces}
\newcommand{\pronun}[1]{}
\newcommand{\abbrev}[1]{\textsc{#1}}
\newcommand{\1}{\ifmmode\mathbin{\%}\else\%\fi}

\raggedbottom
\sloppy
\setlength{\emergencystretch}{3em}

\setlength{\columnsep}{0.3in}
\setlength{\parindent}{0pt}
\setlength{\parskip}{2pt}

\title{\textbf{\Huge Medical Dictionary by Specialty}\\[0.5em]
\Large A Comprehensive Reference of Medical Terminology}
\author{Chaman Singh Verma}
\date{\today}

\begin{document}

\onecolumn
\frontmatter
\maketitle
\thispagestyle{fancy}

\chapter*{Preface}
\addcontentsline{toc}{chapter}{Preface}

\noindent Medical dictionaries have traditionally been organised alphabetically, an arrangement well suited for looking up an unfamiliar term when you already know its name. But medicine is not learned alphabetically, nor is it practised that way. A cardiologist thinking through a differential diagnosis reaches for cardiac terms, not a mixture of cardiology wedged between caesarean section and cataract.

\medskip

\noindent This dictionary is organised by medical specialty for precisely that reason. Each of the 43 chapters corresponds to a recognised clinical or preclinical discipline, with terms listed alphabetically within that chapter. The result is a reference that mirrors the structure of medical knowledge itself: renal terms in Nephrology, antiarrhythmics in Cardiology, autoimmune diseases in Rheumatology. Cross-references redirect the reader from variant names and alternative spellings to the main entry, so ``Pertussis'' still leads to ``Whooping Cough'' and ``Hepatolenticular Degeneration'' points to ``Wilson Disease.''

\medskip

\noindent The entries are drawn primarily from Black's Medical Dictionary (42nd Edition) and supplemented with definitions reflecting contemporary clinical practice. Each entry is written in a concise single-line format, with the term displayed in bold red and automatically indexed. The dictionary runs to over 5,500 entries across 960 pages, covering anatomy, physiology, pathology, pharmacology, clinical medicine, surgery, and the allied health sciences.

\medskip

\noindent This work is a living document. The LaTeX source is designed to be easily extended, corrected, and reorganised. Contributions, corrections, and suggestions are welcome.

\medskip

\noindent\textit{Chaman Singh Verma}


\chapter*{Medical Disclaimer}
\addcontentsline{toc}{chapter}{Medical Disclaimer}

\noindent \textbf{Educational and Reference Purpose Only}: This publication, \textit{Medical Dictionary by Specialty}, is prepared and compiled strictly for educational, academic, historical, and reference purposes. It is intended for medical students, healthcare professionals, educators, and researchers.

\medskip

\noindent \textbf{No Medical Advice}: The information contained herein does not constitute professional medical advice, clinical diagnosis, treatment protocols, or formal therapeutic recommendations. It should not be used as a substitute for the professional judgment of qualified physicians, surgeons, or certified healthcare providers.

\medskip

\noindent \textbf{Clinical Practice Guidelines}: Healthcare practitioners must always consult primary evidence-based medical literature, institutional protocols, manufacturer prescribing information, and national clinical guidelines (such as WHO, CDC, NICE, or national medical councils) when making diagnostic or therapeutic decisions for individual patients.

\medskip

\noindent \textbf{Limitation of Liability}: Neither the author nor the publishers accept any legal responsibility or liability for any errors, omissions, or clinical outcomes resulting from the use or application of any information presented in this work.

\newpage
\tableofcontents
\twocolumn
\mainmatter
\setcounter{tocdepth}{0}




\medterm{A \& E Medicine} Accident and emergency (A\&E) medicine provides rapid assessment, stabilization, diagnosis, and initial treatment for sudden illness and injury.

\medterm{A Fib} Alternate name; see \textit{Atrial fibrillation}.

\medterm{A-T} The adenine--thymine base pair in DNA. The two bases join by hydrogen bonds and help form the structure of DNA.

\medterm{Aagenaes Syndrome} A rare inherited disorder that causes reduced flow of bile from the liver beginning in infancy and long-lasting swelling from lymph-fluid buildup, usually in the legs. Liver problems may come and go with age, but some people develop vitamin deficiency, poor growth, or liver scarring.

\medterm{Aarskog Syndrome} Alternate name; see \textit{Aarskog-Scott Syndrome}.

\medterm{Aarskog-Scott Syndrome} A rare inherited developmental disorder, often related to FGD1 variants, characterized by distinctive facial features, short stature, skeletal differences, and genital anomalies.

\synonyms
Aarskog Syndrome

\medterm{Aase Syndrome} Alternate name; see \textit{Aase-Smith Syndrome II}.

\medterm{Aase-Smith Syndrome I} A rare condition present at birth that may cause excess fluid in the brain, a split in the roof of the mouth, joints that are fixed in bent positions, and unusual hand and facial features. Some people also have heart defects.

\medterm{Aase-Smith Syndrome II} A rare inherited disorder that causes reduced production of red blood cells and thumbs with three bones instead of the usual two.

\medterm{Abadie Sign} Reduced pain sensation when firm pressure is applied to the Achilles tendon, classically described in tabes dorsalis.

\medterm{Abandoned Lead} A wire from a pacemaker, defibrillator, or other implanted electrical device that is no longer connected or used but remains inside the body. A working lead carries electrical signals to the heart and senses its rhythm; its metal core conducts electricity and its silicone or polyurethane covering protects nearby tissue. These materials are chosen to be biocompatible, corrosion-resistant, strong, and flexible, but no implanted material is completely risk-free. A lead may be left in place when removal would be more dangerous because scar tissue can firmly attach it to a vein or the heart. Retained leads can later complicate infection treatment, blood-vessel access, or placement of another device.

\synonyms
Retained Cardiac Device Lead; Abandoned Pacemaker Lead; Abandoned ICD Lead

\medterm{Abasia} Inability to walk even though the legs may have enough strength to move in other situations. It may result from a problem with coordination, the nervous system, or the way movement is planned, rather than paralysis.

\medterm{Abate} To become less intense or to reduce in severity, frequency, or amount; symptoms, inflammation, pain, or fever may abate.

\medterm{Abatement} A reduction in the severity, intensity, frequency, or amount of a symptom, disease process, exposure, or other condition.

\medterm{ABCDE Assessment} A systematic emergency assessment of airway, breathing, circulation, disability, and exposure. Life-threatening
abnormalities are treated as they are identified, with repeated reassessment.

\medterm{ABCDE Rule For Melanoma} A screening mnemonic for suspicious pigmented lesions: asymmetry, border irregularity, color variation,
diameter greater than about 6 mm, and evolution. It does not diagnose melanoma; suspicious lesions require clinical
assessment.

\medterm{Abdomen} The body region between the diaphragm and pelvis. It contains most digestive organs, the kidneys, spleen, pancreas, and
major blood vessels and is described clinically by four quadrants or nine regions.

\medterm{Abdominal} Relating to the abdomen, the body region between the chest and pelvis that contains organs of the digestive, urinary, and reproductive systems.

\medterm{Abdominal Abscess} A localized collection of pus within the abdomen, usually caused by bacterial infection, inflammation, surgery, or
organ perforation. It may cause fever and abdominal pain and often requires antibiotics and drainage.

\medterm{Abdominal Adhesion} Alternate index label; see \textit{Abdominal Adhesions}.

\medterm{Abdominal Adhesions} Bands of scar-like tissue that form inside the abdomen and make two or more organs stick to each other or to the abdominal wall. They most often develop after surgery but may also follow infection, inflammation, endometriosis, or radiation; some are present from birth. They often cause no symptoms, but may cause long-lasting abdominal pain, infertility when they affect pelvic reproductive organs, or a partial or complete intestinal blockage. A blockage may cause pain, bloating, constipation, vomiting, or inability to pass gas.

\medterm{Abdominal Angina} Pain in the abdomen after eating, caused by reduced blood flow through the arteries supplying the intestines. It is usually due to hardening and narrowing of these arteries.

\medterm{Abdominal Aorta} The portion of the aorta below the diaphragm that extends to its bifurcation into the common iliac arteries. It supplies
the abdominal organs, pelvic organs, and lower limbs through its branches.

\medterm{Abdominal Aortic Aneurysm} A weakened, widened area of the abdominal aorta, usually defined as a diameter of at least 3 cm or more than 50 percent above the expected size. Many cause no symptoms, but rapid growth or pain can signal danger. Rupture may cause sudden severe abdominal or back pain, dizziness, fainting, or collapse and requires immediate emergency care.

\synonyms
AAA; Abdominal Aneurysm

\medterm{Abdominal Aortic Aneurysm Repair} Surgical or endovascular treatment of an abdominal aortic aneurysm to prevent rupture or treat a rupture.

\synonyms
AAA Repair

\medterm{Abdominal Aortic Aneurysm Screening} Ultrasound examination used to detect an abdominal aortic aneurysm before symptoms occur. Screening
recommendations depend on age, sex, smoking history, and clinical risk; follow-up intervals depend on aneurysm size.

\medterm{Abdominal Aortic Ultrasound} An ultrasound examination of the abdominal aorta used to look for an aneurysm, measure its size, or monitor it over time. It is painless and does not use radiation.

\medterm{Abdominal Binder} A wide supportive garment placed around the abdomen to provide gentle compression after surgery, injury, or selected abdominal-wall conditions.

\medterm{Abdominal Bloating} A feeling of fullness, pressure, or tightness in the abdomen. It may occur with or without visible enlargement and is commonly caused by intestinal gas, constipation, food intolerance, or functional gastrointestinal disorders. Persistent or severe bloating may require evaluation.

\medterm{Abdominal Bruit} An abnormal blowing or whooshing sound heard over the abdomen with a stethoscope,
usually caused by turbulent blood flow through a narrowed or enlarged artery. It may occur with
renal-artery stenosis or an abdominal aortic aneurysm. Further assessment depends on symptoms and
vascular risk factors.

\medterm{Abdominal Cavity} The cavity between the diaphragm and pelvis containing the stomach, liver, gallbladder,
spleen, pancreas, kidneys, and intestines.

\medterm{Abdominal Circumference} The distance around the abdomen. In pregnancy, fetal abdominal circumference is measured by ultrasound at a standardized level that includes the stomach, portal sinus, and spine; it is used with other measurements to estimate fetal size and assess growth.

\medterm{Abdominal Colic} Cramping abdominal pain caused by contractions or obstruction of a hollow organ,
such as the intestine, ureter, or bile duct. Pain often comes in waves and may be accompanied by
nausea, vomiting, sweating, or altered bowel or urinary function. Severe or persistent pain requires
prompt medical assessment.

\medterm{Abdominal Compartment Syndrome} A dangerous rise in pressure inside the abdomen, usually above 20 millimeters of mercury, that causes new problems with organ function. Causes include severe fluid buildup, an overly swollen or blocked intestine, internal bleeding, and some conditions after abdominal surgery.

\medterm{Abdominal CT} Computed tomography of the abdomen used to evaluate selected causes of abdominal pain, trauma,
bowel obstruction, inflammation, bleeding, and masses. The protocol and use of contrast
depend on the suspected condition and the patient's clinical status.

\medterm{Abdominal CT Scan} Alternate name; see \textit{Abdominal CT}.

\medterm{Abdominal Distension} Visible or measured enlargement of the abdomen, often with a feeling of fullness or tightness. Causes include gas, constipation, fluid, pregnancy, enlarged organs, bowel obstruction, or bleeding; severe pain, vomiting, fever, or inability to pass stool or gas requires urgent assessment.

\synonyms
Abdominal Swelling; Swollen Abdomen; Abdominal Distention

\medterm{Abdominal Distention} Alternate name; see \textit{Abdominal Distension}.

\medterm{Abdominal Exploration} Surgical inspection of the abdominal cavity to identify or treat a suspected source of pain, bleeding, infection, obstruction, or other disease.

\medterm{Abdominal Fat} Adipose tissue stored beneath the skin or around abdominal organs; excess visceral fat is metabolically active and associated with insulin resistance, fatty liver, and cardiovascular risk.

\medterm{Abdominal Fluid} Fluid in the space around the abdominal organs. A small amount may be normal, but excess fluid, called ascites, can cause swelling, discomfort, breathing difficulty, or infection and needs evaluation for its cause.

\synonyms
Peritoneal Fluid

\medterm{Abdominal Girth} The distance around the abdomen, measured at a specified level. It may be used to monitor abdominal distention, fetal growth, obesity, ascites, or changes in intra-abdominal volume.

\medterm{Abdominal Guarding} Tensing of the abdominal muscles when the abdomen is touched. It may occur when the lining around the abdominal organs is irritated or inflamed.

\medterm{Abdominal Hernia} Protrusion of an organ or tissue through a weakness or opening in the abdominal wall. Common types include incisional, umbilical, epigastric, and ventral hernias.

\medterm{Abdominal Hysterectomy} Surgical removal of the uterus through an incision in the abdominal wall; the cervix, ovaries, and fallopian tubes may or may not also be removed.

\medterm{Abdominal Mass} A palpable or imaging-detected abnormal collection, enlargement, or lesion within the
abdomen. Causes include distended bowel or bladder, organ enlargement, hernia,
inflammatory collection, benign tumour, and malignancy; assessment depends on location,
mobility, tenderness, systemic features, and appropriate imaging.

\medterm{Abdominal Migraine} A migraine-related disorder characterized mainly by recurrent episodes of moderate to severe midline abdominal pain, often with nausea or vomiting and little or no headache.

\medterm{Abdominal MRI Scan} A magnetic resonance imaging study that produces detailed images of the abdomen using magnetic fields and radiofrequency energy.

\medterm{Abdominal Muscles} The muscles forming the front and sides of the abdominal wall. They support the trunk, help with posture and breathing, and increase intra-abdominal pressure during coughing, lifting, and defecation.

\synonyms
Abdominal Wall Muscles

\medterm{Abdominal Neoplasm} An abnormal growth in an organ or tissue of the abdomen. It may be noncancerous or cancerous and can cause pain, swelling, bleeding, changes in bowel function, or no symptoms.

\synonyms
Abdominal Tumor

\medterm{Abdominal Pain} Pain or discomfort in the area between the chest and pelvis. It may be felt deep inside, come from the abdominal wall, or be felt in the abdomen even when the cause is elsewhere. Sudden severe pain, rigid muscles, fainting, persistent vomiting, fever, blood in stool, or pregnancy requires prompt medical assessment.

\synonyms
Stomach Pain

\medterm{Abdominal Point Tenderness} Pain felt in one specific area when the abdomen is gently pressed. It may occur with appendicitis, gallbladder inflammation, or irritation of the abdominal lining; it does not diagnose a condition by itself.

\synonyms
Focal Abdominal Tenderness; Localized Peritoneal Tenderness; Point Pain in Abdomen

\medterm{Abdominal Pregnancy} An ectopic pregnancy implanted in the peritoneal cavity rather than the uterus, fallopian tubes, or ovaries.
It is rare and can cause severe internal bleeding. Diagnosis requires specialist imaging and management.

\medterm{Abdominal Quadrant} One of four topographical anatomical divisions of the abdomen created by perpendicular vertical and horizontal imaginary planes intersecting at the umbilicus (right upper, left upper, right lower, and left lower quadrants), utilized in clinical triage to localize abdominal pain and organ pathology.

\medterm{Abdominal Radiography} A regular X-ray of the abdomen. It may be taken while a person is lying down, standing, or lying on their side.

\medterm{Abdominal Reflex} Contraction of the abdominal muscles after stimulation of the abdominal skin. An absent or asymmetric reflex may
indicate neurologic disease, although it can be absent normally in some people.

\medterm{Abdominal Regions} The nine named areas used to describe locations on the abdominal wall: the right and left hypochondriac, epigastric, right and left lumbar, umbilical, right and left iliac, and hypogastric regions.

\medterm{Abdominal Retractor} A surgical instrument used to hold back the abdominal wall or organs to improve access and visibility during an operation.

\synonyms
Balfour Retractor; Bookwalter Retractor

\medterm{Abdominal Rigidity} Involuntary tensing of the abdominal wall, usually caused by irritation of the peritoneum. It may occur with
peritonitis, perforation, severe infection, or intra-abdominal bleeding and requires clinical assessment.

\medterm{Abdominal Sounds} Sounds produced by movement of gas and fluid through the intestines; changes in frequency or character may support, but do not establish, a diagnosis.

\medterm{Abdominal Swelling} Alternate name; see \textit{Abdominal Distension}.

\medterm{Abdominal Tap} Alternate name; see \textit{Abdominocentesis}.

\medterm{Abdominal Thrusts} A first-aid maneuver for severe choking in a conscious adult or child older than 1 year. Give 5 back blows followed by 5 quick inward and upward pushes above the navel, repeating until the object comes out or the person becomes unconscious. Use chest thrusts for a pregnant or very large person; infants and unconscious people need different first aid.

\medterm{Abdominal Ultrasound} An ultrasound examination that uses sound waves to produce images of the abdomen and assess its structure and blood flow when applicable.

\medterm{Abdominal Wall} The layered muscular, fascial, and skin-covered boundary surrounding the abdominal cavity.
It supports and protects abdominal organs and helps with posture, breathing, coughing, and straining.
Weakness or defects may produce abdominal or incisional hernias; trauma, infection, and tumors may
also affect the wall.

\medterm{Abdominal Wall Fat Pad Biopsy} Removal of a small sample of abdominal fat to look for amyloid, usually with a special Congo-red stain. It is less invasive than biopsy of an affected organ, but a negative result does not rule out amyloidosis.

\medterm{Abdominal Wall Hernia} Alternate name; see \textit{Abdominal Hernia}.

\medterm{Abdominal Wall Surgery} Surgery to repair a defect or disease of the abdominal wall, often a hernia. It may be open or laparoscopic, and mesh may or may not be used depending on the defect, location, and patient factors.

\medterm{Abdominal X-Ray} Alternate name; see \textit{Abdominal Radiography}.

\medterm{Abdominocentesis} Inserting a needle through the abdominal wall into the space around the abdominal organs to remove or test extra fluid, called ascites. It is also called paracentesis.

\medterm{Abdominopelvic Cavity} The connected body space below the diaphragm containing abdominal and pelvic organs, including digestive, urinary, and reproductive structures, protected by surrounding muscles, bones, and lining tissues.

\medterm{Abdominoperineal Resection} An operation that removes the anus, rectum, and usually part of the sigmoid colon. Because the normal anal outlet is removed, the surgeon creates a permanent colostomy, bringing the colon through the abdominal wall to drain stool into a bag. It may be used for selected low rectal or anal cancers.

\synonyms
APR

\medterm{Abdominoplasty} A surgical procedure that removes excess abdominal skin and may tighten the abdominal wall. It is commonly called a tummy tuck.

\synonyms
Tummy Tuck

\medterm{Abducens Nerve} The sixth cranial nerve. It controls the muscle that moves the eye outward. Damage to this nerve can prevent outward eye movement and cause double vision.

\medterm{Abducens Nerve Palsy} Dysfunction of the sixth cranial nerve resulting in paralysis or paresis of the lateral rectus muscle. Clinically characterized by horizontal diplopia that worsens on gaze toward the affected side and convergent strabismus (esotropia) in primary gaze. Common etiologies include microvascular ischemia (diabetes mellitus, hypertension), elevated intracranial pressure (false localizing sign), trauma, and cavernous sinus lesions.

\synonyms
Sixth Nerve Palsy; CN VI Palsy

\medterm{Abducent Nerve} Alternate name; see \textit{Abducens Nerve}.

\medterm{Abduct} To move an arm, leg, finger, or other body part away from the center line of the body.

\medterm{Abduction} Movement of a limb or other body part away from the median plane of the body. In ocular motility, it is movement of an eye outward, away from the nose.

\medterm{Abductor Muscle} A muscle that moves a body part away from the midline or, for a digit, away from the axis of the hand or foot.

\medterm{Aberrant Conduction} Transient abnormal intraventricular conduction of a supraventricular electrical impulse through the ventricular conduction system, typically occurring when an impulse reaches bundle branch tissue during its relative refractory period. Most commonly presents as a wide QRS complex with right bundle branch block morphology following a sudden change in cycle length (Ashman phenomenon).

\synonyms
Aberrancy; Ventricular Aberration

\medterm{Aberration} A deviation from normal structure, function, or appearance; in optics, it is an imperfection that prevents light from forming a sharply focused image.

\medterm{Abetalipoproteinemia} A rare inherited disorder in which the body cannot make normal particles that carry fat and cholesterol in the blood. It causes poor absorption of dietary fat, deficiencies of vitamins carried with fat, abnormal red blood cells, and possible nerve or eye disease.

\medterm{Abfraction} Loss of tooth material near the gum line, possibly related to repeated mechanical stress or biting forces. It may cause a notch and tooth sensitivity; a dentist can assess the cause and need for treatment.

\medterm{Abiotrophy} Progressive loss of function as cells in a tissue or organ gradually degenerate, often because of an inherited or developmental problem. It may cause worsening weakness, loss of sensation, poor coordination, or other problems depending on the affected part of the body.

\medterm{Ablate} To remove or destroy tissue using surgery, heat, cold, chemicals, or another energy source. Ablation may treat abnormal heart rhythms, tumors, pain pathways, or other selected conditions.

\medterm{Ablation} A procedure that removes or destroys tissue or stops an unwanted body function using surgery, heat, cold, electricity, medicines, or other methods.

\medterm{Ablative Surgery} Surgery that destroys or removes abnormal tissue, often using excision, heat, cold, electrical energy, or other tissue-destructive methods.

\medterm{Ablepharon} A rare congenital condition in which one or both eyelids are absent or severely underdeveloped. It can leave the eye exposed and may require protective treatment and reconstructive surgery.

\medterm{Abnormal} Different from the usual healthy structure, function, or appearance. An abnormal finding may be harmless or may indicate a disease, depending on its type and context.

\medterm{Abnormal (Dysfunctional) Uterine Bleeding} Bleeding from the uterus that is abnormal in timing, regularity, duration, or amount and is not explained by pregnancy. Causes include structural lesions, ovulatory or hormonal disorders, medicines, bleeding disorders, and endometrial disease; evaluation depends on age and risk factors.

\medterm{Abnormal Color} A change in the usual color of skin, moist tissues, body fluids, or another structure. Its meaning depends on the color, location, duration, and associated symptoms.

\synonyms
Discoloration

\medterm{Abnormal Fissure} An atypical crack, groove, or cleft in a tissue or organ. It may be congenital or develop from injury, inflammation, scarring, pressure, or disease, and its importance depends on its location and depth.

\medterm{Abnormal Hemoglobins Testing} Blood testing that looks for unusual forms of hemoglobin, the oxygen-carrying protein in red blood cells. It helps diagnose inherited disorders such as sickle-cell disease and thalassemia; results may require specialized laboratory analysis.

\medterm{Abnormal Labor} Labor that is unusually slow, prolonged, arrested, or associated with abnormal contractions, fetal position, or fetal intolerance. Evaluation considers cervical change, contraction pattern, pelvic factors, fetal status, infection, and maternal condition before treatment or delivery decisions.

\medterm{Abnormal Tooth Shape} Abnormal tooth shape may result from developmental defects, genetic conditions, trauma, wear, or tumors and can affect function, appearance, and eruption.

\medterm{Abnormal Uterine Bleeding} Bleeding from the uterus that is unusually heavy, prolonged, irregular, or occurs between periods, after sex, or after menopause. Causes include pregnancy-related problems, hormonal changes, fibroids, infection, medicines, and cancer. Persistent or heavy bleeding may indicate an underlying condition.

\medterm{Abnormal Virilization} Development of masculinizing physical features in a person whose expected sex development does not include them, often from excess androgens.

\synonyms
Virilization

\medterm{Abnormality} A difference from the usual structure, function, development, or test result. An abnormality may be harmless, temporary, or a sign of disease; its importance depends on the finding, symptoms, and clinical evaluation.

\synonyms
Anomaly

\medterm{Abnormally Dark Or Light Skin} Skin that has become darker or lighter than a person's usual skin tone. Darkening is called
hyperpigmentation and lightening is called hypopigmentation. Causes include inflammation or injury, sun exposure,
medicines, infections, hormonal disorders, and genetic conditions.

\medterm{Abnormally Excessive Sweating} See \textit{Hyperhidrosis}.

\medterm{ABO Blood Group System} The major red-cell blood-group system based on whether A and/or B antigens are present on the red-cell surface and the corresponding antibodies are present in plasma. The four principal phenotypes are A, B, AB, and O; correct ABO matching is critical for safe transfusion.

\synonyms
ABO Blood Group
ABO Blood Groups
ABO Blood Types
ABO System

\medterm{ABO Blood Typing} A laboratory test that determines a person's ABO and Rh blood groups. It is performed before transfusion,
transplantation, and some obstetric procedures.

\medterm{ABO Incompatibility} An immune reaction that can occur when blood with an incompatible A, B, or O type is transfused, or when maternal and fetal blood types are incompatible. It can destroy red blood cells, cause a severe transfusion reaction, or affect a newborn.

\medterm{Abortion} The ending of a pregnancy, either spontaneously or intentionally. An intentional abortion uses medicines or a medical procedure; spontaneous abortion is also called miscarriage.

\medterm{Abortion Complications} Potential complications after spontaneous or induced abortion, including hemorrhage, retained tissue, infection, uterine injury, cervical trauma, thromboembolism, and emotional or psychological distress.

\medterm{Abortive} Incomplete, unsuccessful, or stopped before full development; an abortive infection produces limited disease without the usual complete clinical course.

\medterm{Abortive Polio} A mild, nonparalytic poliovirus infection causing nonspecific fever, sore throat, headache, nausea, or malaise without the neurologic findings of paralytic polio.

\medterm{Above The Knee Amputation} Amputation of the lower limb through the thigh above the knee joint, usually performed for severe vascular disease, infection, trauma, or cancer.

\medterm{ABR Test} Auditory brainstem response testing records electrical activity from the auditory nerve and brainstem through scalp electrodes to assess hearing and auditory-pathway function, especially in infants or people unable to provide reliable behavioral responses.

\medterm{Abrade} To wear away tissue by friction or erosion; an abrasion is the resulting scrape or superficial loss of skin.

\medterm{Abrasion} A superficial injury in which the skin is scraped or rubbed away, usually by friction. In dentistry, abrasion means mechanical wearing away of tooth structure by an external force.

\medterm{Abrin} A highly poisonous protein found in rosary pea seeds. It stops cells from making proteins, causing them to die. Swallowing chewed seeds or exposure to abrin powder can cause vomiting, diarrhea, dehydration, breathing problems, organ failure, and death. Suspected exposure requires urgent medical care; there is no routine antidote.

\medterm{Abruptio Placenta} Alternate name; see \textit{Abruptio Placentae}.

\medterm{Abruptio Placentae} Premature separation of the placenta from the uterine wall before delivery. It can cause vaginal bleeding,
abdominal pain, uterine tenderness, fetal distress, and maternal shock and requires urgent obstetric assessment.

\medterm{Abs} An informal name for the muscles at the front and sides of the abdomen. They help bend and rotate the trunk, support the spine, protect abdominal organs, and increase pressure for coughing, lifting, breathing, or bowel movements.

\medterm{Abscess} A pocket of pus that forms in tissue, an organ, or a body cavity, usually because of an infection. It may cause pain, swelling, warmth, and fever; treatment often requires drainage and sometimes antibiotics.

\medterm{Abscess, Bartholin's} A pocket of pus caused by infection in a Bartholin gland near the vaginal opening. It usually causes a painful, swollen lump and may need drainage; a fever or severe pain needs medical care.

\medterm{Abscess, Cutaneous} Alternate index label; see \textit{Boils}.

\medterm{Abscessed Tooth} A collection of pus caused by bacterial infection of the tooth or surrounding tissues, usually arising from
untreated dental decay, a cracked tooth, or periodontal disease. Pain, facial swelling, fever, and sensitivity may
occur. Dental treatment and drainage are usually required; swelling that affects swallowing or breathing is an
emergency.

\medterm{Abscission} The separation or removal of tissue, either naturally or by cutting. In surgery it may describe removing a growth, damaged tissue, or part of an organ; the exact meaning depends on the medical context.

\medterm{Absence Epilepsy} A generalized epilepsy characterized by brief episodes of impaired awareness, often with staring and subtle movements.

\synonyms
Childhood Absence Epilepsy; Pyknolepsy; Petit Mal Epilepsy

\medterm{Absence Of Sweating} Anhidrosis is diminished or absent sweating, which may be focal or generalized. It impairs thermoregulation and can result from autonomic dysfunction, peripheral neuropathy, skin disease, medications, or congenital conditions. Severe generalized anhidrosis can cause heat intolerance and exertional hyperthermia.

\medterm{Absence Of The Breast} Congenital absence of breast development, called amastia; it may affect one or both sides and can occur with other developmental anomalies.

\medterm{Absence Of The Nipple} Athelia: a rare congenital absence of one or both nipples, which may occur with hereditary disorders, chromosomal abnormalities, or developmental syndromes.

\medterm{Absence Seizure} A brief seizure causing sudden staring and reduced awareness, usually lasting only seconds with rapid recovery. The person may stop speaking or responding and may make small repetitive movements; an EEG helps confirm the diagnosis.

\synonyms
Petit Mal Seizure; Typical Absence Seizure

\medterm{Absence Seizure Cross-Reference} Alternate name; see \textit{Absence seizure}.

\medterm{Absent Breath Sounds} Lack of audible air movement over part or all of a lung. Causes include pneumothorax, pleural effusion, or
airway obstruction; sudden onset requires urgent assessment.

\medterm{Absent Eye} Anophthalmia, a congenital absence or severe underdevelopment of the eye globe and ocular tissue; it may affect one or both eyes and may occur with other birth defects.

\medterm{Absent Menstrual Periods} Alternate index label; see \textit{Amenorrhea}.

\medterm{Absent Menstrual Periods, Primary} Alternate index label; see \textit{Amenorrhea, Primary}.

\medterm{Absent Menstrual Periods, Secondary} Alternate index label; see \textit{Amenorrhea, Secondary}.

\medterm{Absent Pulmonary Valve} A birth defect in which the valve between the heart and the lung arteries is missing or severely underdeveloped. It can enlarge those arteries and cause breathing problems, poor blood flow to the lungs, or other heart defects.

\medterm{Absent Reflex} Failure of an expected automatic movement to occur when a tendon or other body area is stimulated. It may suggest a problem in a nerve, nerve root, spinal cord, or muscle, although reflexes can normally differ between people and sides.

\medterm{Absent Thirst} Reduced or absent desire to drink despite a potential need for fluids, which can occur with hypothalamic disease, aging, altered consciousness, medications, or severe illness.

\medterm{Absolute Eosinophil Count} A laboratory blood test measuring the exact number of eosinophils per microliter of blood to evaluate allergic disorders, asthma, parasitic infections, drug reactions, or hypereosinophilic syndromes.

\synonyms
AEC; Absolute Eosinophils; Blood Eosinophil Count

\medterm{Absolute Neutrophil Count} The number of neutrophils, a type of white blood cell that fights many bacterial and fungal infections, in a volume of blood. A low count increases infection risk, especially when it is severely reduced.

\synonyms
ANC

\medterm{Absorb} To take in a substance, such as through the skin or intestine, or to reduce the intensity of radiation or light by taking up its energy.

\medterm{Absorption} The process by which a drug enters the bloodstream from its site of administration, such
as the digestive tract, skin, or muscle.

\medterm{Abstinence} Deliberate avoidance of a substance or behavior, such as alcohol, drugs, tobacco, or sexual activity. Stopping some substances, especially alcohol, sedatives, or opioids, can cause dangerous withdrawal and may require medical supervision.

\medterm{Abulia} Marked reduction in motivation, initiative, and spontaneous activity, usually caused by dysfunction of frontal-subcortical brain circuits or related neurologic disease.

\medterm{Abusive Head Trauma} Injury to an infant or young child caused by inflicted blunt impact, violent shaking, or
both. Findings may include subdural hemorrhage, retinal hemorrhage, encephalopathy,
fractures, or other injuries; diagnosis requires careful multidisciplinary assessment
because no single finding is diagnostic.

\synonyms
Shaken Baby Syndrome; Non-Accidental Head Injury

\medterm{Abutment} A natural tooth, dental implant, or component that supports a dental prosthesis. Abutments may hold a bridge, denture, or implant crown in position and transfer chewing forces to the supporting tooth, implant, or bone.

\medterm{AC Joint} The acromioclavicular joint, a gliding joint between the acromion of the scapula and the clavicle at the top of the shoulder.

\medterm{Acalculia} Loss of the ability to do or understand numbers and calculations that a person could previously manage. It may result from injury or disease affecting areas of the brain involved in calculation.

\medterm{Acalculous Cholecystitis} Inflammation of the gallbladder without gallstones. It usually occurs suddenly in people who are critically ill, injured, recovering from surgery, unable to eat, or in shock. Reduced bile flow and poor blood supply can contribute, and severe cases may cause tissue death or a hole in the gallbladder. Fever and pain in the upper-right abdomen may occur; ultrasound is commonly used for evaluation.

\medterm{Acanthamoeba} A free-living amoeba found in soil and water that can cause corneal infection, skin or sinus infection, or rarely granulomatous amoebic encephalitis. Contact-lens exposure to contaminated water increases keratitis risk, while invasive disease is more likely with impaired immunity.

\medterm{Acanthamoeba Keratitis} A rare but potentially sight-threatening infection of the cornea caused by the free-living amoeba
Acanthamoeba. It is strongly associated with contaminated contact-lens water or poor lens hygiene and may cause
severe eye pain, redness, photophobia, blurred vision, and a ring-shaped corneal infiltrate. Prompt ophthalmic
assessment and prolonged antiamoebic treatment are required.

\medterm{Acanthocyte} An abnormal red blood cell with irregular, thorny or spike-like projections of varying length and width on its outer membrane. They are seen on a blood smear in advanced liver disease (such as cirrhosis), lipid metabolism disorders, abetalipoproteinemia, or neuroacanthocytosis.

\synonyms
Spur Cell

\medterm{Acanthocytosis} An abnormal hematologic condition characterized by the presence of circulating red blood cells with irregular, thorny or spiked projections (acanthocytes or spur cells). It is classically observed in abetalipoproteinemia, advanced liver cirrhosis (spur cell anemia), severe malnutrition, and neuroacanthocytosis syndromes.

\synonyms
Spur Cell Formation

\medterm{Acantholysis} Loss of attachment between cells in the outer layer of the skin, causing them to separate. This can produce fragile blisters or erosions and occurs in some autoimmune, inherited, infectious, or drug-related skin disorders.

\medterm{Acanthoma} A usually benign skin growth made of thickened squamous epidermal cells. Its appearance can resemble other lesions, so a changing or uncertain lesion may need examination or biopsy.

\medterm{Acanthosis Nigricans} Dark, thickened, velvety skin, usually in folds such as the neck, armpits, or groin. It is often associated with insulin resistance, obesity, type 2 diabetes, hormonal disorders, or medicines; rarely, sudden extensive disease can signal an underlying cancer. Treating the cause may improve the skin changes.

\medterm{Acapnia} An unusually low carbon-dioxide level in arterial blood, usually caused by breathing too quickly or deeply. It may cause light-headedness, tingling, or a change in the blood's acid level.

\medterm{Acardia} Absence or severe malformation of the fetal heart, usually occurring in the setting of a
rare twin perfusion sequence. The affected fetus cannot sustain independent circulation.

\medterm{Acariasis} Illness or skin infestation caused by mites or ticks. It may cause itching, a rash, burrows, skin nodules, or deeper tissue inflammation; some mites and ticks can also transmit infection.

\medterm{Acaricide} A pesticide used to kill mites and related arachnids such as ticks. It may be used on animals, plants, or environments; exposure can irritate skin and may cause poisoning if misused.

\medterm{Acatalasemia} A rare inherited condition in which catalase, an enzyme that breaks down hydrogen peroxide, is absent or greatly reduced. Many people have no symptoms, but some develop mouth ulcers, tissue damage, or increased sensitivity to oxidative stress.

\medterm{Acatalasia} A rare inherited condition in which catalase, an enzyme breaking down hydrogen peroxide, is absent or greatly reduced. Some people have mouth ulcers, infections, or tissue damage, while others have no symptoms.

\synonyms
Catalase Deficiency

\medterm{Accelerated Growth} An increase in a child's growth velocity above the expected rate for age and developmental stage, determined from serial height or length measurements. It may occur during normal puberty or catch-up growth; abnormal acceleration, especially with early pubertal signs or advanced bone age, may suggest precocious puberty or another endocrine disorder.

\medterm{Accelerated Idioventricular Rhythm} A heart rhythm that starts in the lower chambers and is usually slower than most dangerous ventricular rhythms. It has wide beats and may occur during recovery of blood flow to the heart. It is often temporary, but the person and the cause still need assessment.

\synonyms
AIVR

\medterm{Accelerometry} Measurement of movement or acceleration using a sensor. It can assess physical activity, walking, tremor, falls, sleep, or rehabilitation progress, depending on the device and the recording method.

\medterm{Acceptable Daily Intake} The estimated amount of a substance that a person can consume every day throughout life without a meaningful health risk, based on available evidence and safety margins. It is usually expressed per kilogram of body weight.

\synonyms
ADI

\medterm{Acceptance-Based Therapies} Psychotherapeutic approaches that help a person notice and accept thoughts, feelings, or bodily sensations while reducing unhelpful attempts to control or avoid them.

\synonyms
Acceptance-Based Psychotherapy; Acceptance-Focused Therapy

\medterm{Access Needle} A hollow needle used to enter a blood vessel, body cavity, organ, or fluid collection. It may be used to take a sample, drain fluid, give treatment, or help place a catheter.

\synonyms
Introducer Needle; Percutaneous Access Needle; Vascular Access Needle

\medterm{Accessory} Additional or supplementary; in anatomy, an accessory structure is an extra part or a structure that assists the usual organ or system.

\medterm{Accessory Digestive Organ} An organ that supports digestion but is not part of the alimentary canal, such as the liver, pancreas, gallbladder, or salivary glands.

\medterm{Accessory Dwelling Unit} A secondary, independent residential unit on the same property as a primary home; it is a housing and community-health term rather than a disease or body structure.

\medterm{Accessory Muscles} Muscles recruited to assist the primary muscles of respiration or another movement when the normal workload is increased; use of respiratory accessory muscles may indicate respiratory distress.

\medterm{Accessory Nerve} The eleventh cranial nerve (CN XI), a primarily motor nerve whose spinal component supplies the sternocleidomastoid and trapezius muscles. It enables head rotation and shoulder elevation; injury may cause weak head rotation, impaired shoulder shrug, and shoulder droop.

\synonyms
Cranial Nerve XI; Spinal Accessory Nerve

\medterm{Accessory Pathway} An abnormal electrical connection between the atria and ventricles that can bypass the atrioventricular node and cause pre-excitation or tachyarrhythmias, including atrioventricular reentrant tachycardia.

\medterm{Accessory Placenta} An extra placental lobe separate from the main placenta but connected to it by fetal blood vessels; it can increase the risk of retained placental tissue or vessel injury during delivery.

\medterm{Accident And Emergency Medicine} Accident and emergency medicine is the specialty concerned with immediate evaluation and treatment of acute illness, injury, poisoning, and life-threatening emergencies.

\medterm{Accidental Death} Death caused by an unintentional injury or event rather than disease, suicide, or homicide. The cause is determined through medical, scene, and sometimes legal investigation.

\synonyms
Unintentional Death

\medterm{Accidental Fall} An unintentional movement to the ground or a lower level. Falls can cause bruises, fractures, head injury, or internal bleeding, especially in older adults; loss of consciousness, severe pain, or confusion requires urgent care.

\synonyms
Fall Injury

\medterm{Accommodation} The ability of the eye to change the shape of its lens so that near objects can be focused clearly.

\medterm{Accommodation Reflex} The coordinated adjustment of the lens, ciliary muscle, and pupils that shifts the eye from distance to near focus; it includes lens thickening, convergence, and pupillary constriction.

\medterm{Accommodative Insufficiency} Inability of the eye to increase lens power adequately for near focusing, causing eyestrain, blurred near vision, headache, or difficulty reading.

\medterm{Accouchement} Accouchement is the process of childbirth, including labor, delivery of the infant, and expulsion of the placenta.

\medterm{Accretion} Gradual buildup or growth of material that makes an object or body tissue larger. Examples include crystals forming in fluid, extra bone developing after injury, or plaque collecting on teeth.

\medterm{ACE Blood Test} A blood test measuring angiotensin-converting enzyme. It may support evaluation of sarcoidosis and some other inflammatory or liver disorders, but an abnormal result is not specific and cannot diagnose a disease by itself.

\medterm{ACE Inhibitors} Medicines that reduce formation of angiotensin II, lowering blood pressure and pressure within the kidney's filtering units. They are used for hypertension, heart failure, and selected kidney or cardiovascular conditions; cough, high potassium, kidney-function changes, and rare facial or airway swelling are important risks, and they are generally avoided in pregnancy.

\medterm{Acebutolol} Acebutolol is a beta-adrenergic blocker with intrinsic sympathomimetic activity used in some settings for hypertension and cardiac rhythm disorders; it can slow heart rate and conduction.

\medterm{Acellular} Without cells. An acellular material may consist of fluid, minerals, or the supporting material around cells and may be natural, diseased, or used in a medical product.

\medterm{Acentric Chromosome} A chromosome fragment that lacks a centromere and therefore usually cannot attach normally to the spindle during cell division.

\medterm{Acephalostomia} An extremely rare and severe birth defect involving absence or major failure of development of the head and mouth. It is usually incompatible with survival and is identified during prenatal or newborn examination.

\medterm{Acetabular Labrum} A fibrocartilaginous rim attached to the margin of the acetabulum that deepens the hip
socket, helps maintain a fluid seal, and contributes to hip-joint stability.

\medterm{Acetabular Notch} The indentation in the inferior margin of the acetabulum, bridged by the transverse
ligament of the acetabulum and forming a passage for structures entering the hip joint.

\medterm{Acetabuloplasty} A surgical procedure that reshapes or augments the acetabular rim or socket to improve
coverage of the femoral head in selected structural hip disorders.

\medterm{Acetabulum} The cup-shaped socket of the hip formed by the fused ilium, ischium, and pubis. It articulates with the femoral head
and is deepened by the acetabular labrum.

\medterm{Acetaldehyde} A chemical produced when the body metabolizes alcohol. Accumulation may contribute to flushing, nausea, and vomiting.

\medterm{Acetaminophen} A widely used medication that relieves mild-to-moderate pain and reduces fever, but lacks significant anti-inflammatory effects. While safe at recommended therapeutic doses, excessive amounts can cause severe, life-threatening liver injury, especially when taken with alcohol or other products containing acetaminophen.

\synonyms
Paracetamol; APAP

\medterm{Acetaminophen And Codeine Overdose} A medical emergency caused by ingesting toxic amounts of a combination pain medicine. It produces dual life-threatening effects: codeine suppresses the central nervous system, causing extreme drowsiness, slowed or stopped breathing, and coma, while acetaminophen depletes protective liver antioxidants, leading to progressive liver failure.

\synonyms
Co-Codamol Overdose; Tylenol with Codeine Overdose

\medterm{Acetaminophen Overdose} Taking more acetaminophen than the body can safely process, causing serious liver damage. A person may have few symptoms at first, followed later by nausea, right-sided abdominal pain, jaundice, or liver failure. Get emergency help or contact a poison center immediately; the antidote N-acetylcysteine works best when started early but may still help when given later.

\synonyms
Acetaminophen Toxicity; Acetaminophen Poisoning; Paracetamol Toxicity; Paracetamol Overdose; APAP Overdose

\medterm{Acetaminophen Poisoning} Alternate name; see \textit{Acetaminophen Overdose}.

\medterm{Acetaminophen Toxicity} Alternate name; see \textit{Acetaminophen Overdose}.

\medterm{Acetate} A form of acetic acid and a small molecule used by cells in energy and building processes. The body can convert acetate into acetyl-CoA, which helps make or break down fats and other substances.

\medterm{Acetazolamide} A medicine that changes salt and fluid handling in the kidneys. It is used for glaucoma, prevention or treatment of altitude illness, and selected fluid, metabolic, or neurologic conditions.

\medterm{Acetone} A volatile ketone body produced during increased fat metabolism. It may contribute to the characteristic breath odor of ketosis and ketoacidosis.

\medterm{Acetone Body} A ketone produced when the body breaks down fat for energy. Acetone is one of three main ketones and can contribute to the distinctive breath odor of prolonged fasting or diabetic ketoacidosis.

\medterm{Acetone Breath} A distinctive sweet, fruity odor on the breath caused by the exhalation of volatile acetone, a ketone body produced during accelerated lipolysis and ketogenesis. It is a characteristic physical examination finding in diabetic ketoacidosis (DKA), as well as in severe starvation, prolonged fasting, and ketogenic metabolic states.
\synonyms{Fruity Breath; Ketotic Breath; Diabetic Breath Odor}

\medterm{Acetone Poisoning} Harm from swallowing or breathing a large amount of acetone, a solvent found in some products. It may cause irritation, dizziness, vomiting, drowsiness, low blood pressure, or breathing problems; significant exposure is a medical emergency requiring prompt treatment.

\medterm{Acetonuria} The presence of acetone or other ketone bodies in the urine, which may occur with fasting, vomiting, uncontrolled diabetes, or other states of increased fat metabolism.

\medterm{Acetyl-CoA Carboxylase} An enzyme that begins the process of making fatty acids by changing acetyl-CoA into malonyl-CoA. Its activity responds to the cell’s energy supply and to hormones, helping balance fat production with energy needs.

\medterm{Acetyl-Coenzyme A} A substance used in many energy and building processes inside cells. It helps release energy from food and provides material for making fatty acids, cholesterol, and ketones.

\medterm{Acetylation} A chemical change that adds a small acetyl group to a protein, medicine, or DNA-related molecule. This can alter how the molecule works, how long it lasts, how it is moved, or how strongly a gene is active.

\medterm{Acetylcholine} A major neurotransmitter synthesized from choline and acetyl-CoA by choline acetyltransferase. Operating throughout the central and peripheral nervous systems, it mediates neuromuscular transmission at motor endplates, autonomic ganglionic neurotransmission, parasympathetic postganglionic signaling, and central cognitive/memory processing.

\synonyms
ACh

\medterm{Acetylcholine Receptor Antibody} Acetylcholine receptor antibody testing detects antibodies associated with autoimmune myasthenia gravis and is interpreted with the neurologic examination and other studies.

\medterm{Acetylcholinesterase} A rapid-acting serine hydrolase enzyme localized at synaptic clefts and neuromuscular junctions that hydrolyzes the neurotransmitter acetylcholine into choline and acetate, terminating signal transmission and allowing tissue repolarization.

\synonyms
AChE; Acetylcholine Acetylhydrolase

\medterm{Acetylcholinesterase Inhibitors} A group of medicines that slow the breakdown of acetylcholine, a chemical messenger used by nerves. This increases acetylcholine activity and can help manage conditions such as Alzheimer disease and myasthenia gravis or reverse some medicine-related muscle paralysis.

\medterm{Acetylcysteine} A medicine that loosens thick mucus so it is easier to clear. It is also the main antidote for acetaminophen overdose because it restores glutathione, a substance the liver uses to neutralize the overdose’s toxic breakdown product.

\medterm{Acetylsalicylic Acid} Alternate index label; see \textit{Aspirin}.

\medterm{ACh (Acetylcholine)} Alternate name; see \textit{Acetylcholine}.

\medterm{Achalasia} A disorder in which the muscle ring between the esophagus and stomach does not relax properly, and the esophagus cannot move food normally into the stomach. It causes difficulty swallowing and regurgitation.

\synonyms
Cardiospasm; Esophageal Achalasia

\medterm{Achalasia Of The Cardia} A disorder in which the valve between the esophagus and stomach does not relax properly and the esophagus loses its normal squeezing movement. Food and liquid can collect in the esophagus, causing progressive difficulty swallowing, regurgitation, chest discomfort, and weight loss.

\medterm{AChE (Acetylcholinesterase)} Alternate name; see \textit{Acetylcholinesterase}.

\medterm{Achenbach's Syndrome} A benign condition of recurrent, painful, spontaneous hematoma formation in the pulp of
a finger or thumb, typically affecting middle-aged women. It presents with sudden onset
of throbbing pain and blue-black discoloration of the fingertip, resolving spontaneously
over one to two weeks without treatment.

\medterm{Achilles Tendinitis} Painful irritation or degeneration of the tendon joining the calf muscles to the heel. It usually follows repeated loading and can cause stiffness, swelling, tenderness, and pain during walking or exercise.

\synonyms
Achilles Tendinopathy; Calcaneal Tendinitis

\medterm{Achilles Tendinopathy} Pain and structural disorder of the Achilles tendon, usually caused by repeated loading. It causes posterior
heel pain and stiffness and is managed with activity modification, rehabilitation, and treatment of contributing factors.

\medterm{Achilles Tendon} The thickest and strongest tendon in the human body, connecting the calf muscles (gastrocnemius and soleus) to the heel bone (calcaneus). It transmits muscle force to the foot, enabling the ankle to point downward (plantarflexion) during walking, running, jumping, and standing on tiptoe.

\synonyms
Calcaneal Tendon; Tendo Calcaneus

\medterm{Achilles Tendon Repair} Surgery that reconnects a completely or severely torn Achilles tendon. The ankle is protected in a cast or brace afterward, followed by gradual exercises to restore movement and strength; recovery takes several months.

\medterm{Achilles Tendon Rupture} A partial or complete tear of the Achilles tendon, usually caused by sudden force. It causes posterior ankle
pain, weakness in plantar flexion, and difficulty walking or standing on tiptoe.

\medterm{Achilles Tendonitis} A commonly used term for pain and inflammation or degeneration of the Achilles tendon. It may affect the tendon
at its heel insertion or its mid-portion and is often related to repetitive loading.

\medterm{Achillobursitis} Inflammation of a bursa near the Achilles tendon, causing pain, tenderness, swelling, or friction-related discomfort at the back of the heel.

\medterm{Achlorhydria} A condition in which the stomach makes little or no hydrochloric acid. It may result from autoimmune gastritis, long-term acid-suppressing medicines, or loss of acid-producing cells and can impair absorption of iron and vitamin B12.

\medterm{Achondrogenesis} A group of severe inherited disorders in which cartilage and bone development is profoundly impaired. Affected fetuses have very short limbs and poorly formed bones, and most cases are fatal before or shortly after birth.

\medterm{Achondroplasia} An autosomal dominant bone-growth disorder caused by a disease-causing change in the \textit{FGFR3} gene. It causes short arms and legs, short stature, a larger head, and a prominent forehead. One altered gene copy usually causes the condition; two copies cause a severe, often fatal form.

\synonyms
Chondrodystrophy; Achondroplastic Dwarfism

\medterm{Achoo Syndrome} The photic sneeze reflex: repeated sneezing triggered by sudden exposure to bright light, especially after entering sunlight from darkness.

\medterm{Achromatopsia} An inherited disorder causing absent or greatly reduced color vision, often with light sensitivity and poor visual acuity.

\synonyms
Complete Color Blindness

\medterm{Achromia} Achromia is loss or marked reduction of pigment in skin, hair, eye, or another tissue; causes include genetic disorders, autoimmune depigmentation, injury, and chemical exposure.

\medterm{Aciclovir} Aciclovir, also called acyclovir, is an antiviral drug that inhibits herpesvirus DNA polymerase and is used for infections such as herpes simplex and varicella-zoster.

\medterm{Acid} A substance that increases acidity by releasing hydrogen ions or accepting electrons. Strong acids can damage tissue, while acids also take part in normal digestion, metabolism, and body-fluid regulation.

\medterm{Acid Base Balance} The body’s control of acids and bases in its fluids, especially the blood. The lungs remove carbon dioxide and the kidneys adjust acids and bicarbonate to keep blood pH within a narrow safe range.

\medterm{Acid Indigestion} Dyspepsia or upper-abdominal discomfort that may include burning, fullness, belching, or nausea; it can accompany reflux, gastritis, ulcer disease, or other conditions.

\medterm{Acid Loading Test} A specialized test of how well the kidneys remove acid from the body. After a carefully controlled acid dose, urine is measured for acidity and other substances to investigate suspected kidney-tubule problems.

\medterm{Acid Mucopolysaccharides} An older name for long sugar chains found in supporting tissues such as cartilage, skin, and joint fluid. Abnormal breakdown or storage of these substances causes a group of inherited storage disorders.

\medterm{Acid Phosphatase} An enzyme that works best in acidic conditions and is found in several tissues, including prostate and blood cells.

\synonyms
ACP

\medterm{Acid Reflux} Backward flow of acidic stomach contents into the esophagus, causing heartburn, sour regurgitation, cough, or throat irritation. Frequent reflux may be gastroesophageal reflux disease and can injure the esophagus.

\medterm{Acid Soldering Flux Poisoning} Injury from swallowing or inhaling acidic chemicals used in soldering. It may burn the mouth, throat, skin, or eyes and can cause vomiting, breathing difficulty, or dangerous changes in body chemistry; significant exposure requires urgent medical treatment.

\medterm{Acid Sphingomyelinase Deficiency} An inherited lysosomal storage disorder caused by deficient acid sphingomyelinase, usually because of disease-causing variants in the \textit{SMPD1} gene. Sphingomyelin accumulates in cells and may cause enlarged liver or spleen, lung disease, bone changes, or neurologic impairment; it is also called Niemann--Pick disease types A and B.

\medterm{Acid-Base Analysis} Evaluation of blood pH, bicarbonate, and carbon dioxide to identify metabolic or respiratory disturbances,
including those caused by poisoning.

\medterm{Acid-Base Disorder} A problem that makes the blood too acidic or too alkaline. It may result from changes in breathing, kidney function, metabolism, or severe illness. Blood tests help identify the cause and whether the body is compensating.

\medterm{Acid-Base, Fluid, and Electrolyte Disorders} Disorders that disturb blood acidity, body-water balance, or concentrations of minerals such as sodium, potassium, calcium, magnesium, and phosphate. They may arise from kidney, heart, endocrine, gastrointestinal, medication, or critical-illness problems and can affect the brain, muscles, and heart.

\medterm{Acid-Fast Bacilli} Rod-shaped bacteria that resist decolorization by acid-alcohol solutions during laboratory staining procedures due to the high lipid and mycolic acid content of their cell walls. The group classically includes \textit{Mycobacterium tuberculosis} and \textit{Mycobacterium leprae}.

\synonyms
AFB; Acid-Fast Bacteria

\medterm{Acid-Fast Bacillus Smear} Microscopic examination of a specimen, often sputum, for acid-fast organisms using a
special stain. A positive result supports mycobacterial infection but does not by itself
identify the species or establish drug susceptibility; culture or molecular testing is
also required.

\medterm{Acid-Fast Stain} A special dye test used to find mycobacteria and a few related organisms in samples such as sputum. A positive result does not identify the species or show which medicines will work.

\medterm{Acidemia} An abnormally low arterial blood pH (below 7.35), indicating an increased hydrogen ion concentration in circulating blood. It is the measurable laboratory state resulting from an underlying respiratory or metabolic acidosis.

\medterm{Acidophil} A cell or tissue part that absorbs acidic laboratory dyes strongly, especially the pink dye eosin. The pattern helps pathologists recognize certain cells in tissue samples, including some pituitary cells.

\medterm{Acidophilus} Lactobacillus acidophilus, a lactic-acid-producing bacterium found in the intestinal and genital microbiota and used in some fermented foods and probiotic products.

\medterm{Acidosis} A condition in which the body’s fluids become too acidic. It can result from breathing problems, kidney disease, severe infection, diabetes, poisoning, or other illness and may impair organ function when severe.

\medterm{Aciduria} Aciduria means an abnormally acidic urine or the presence of excess organic acids in urine; it may reflect diet, dehydration, medication, or an inherited metabolic disorder.

\medterm{Acinar Cell} A cell in a small cluster that makes and releases a fluid or substance. In the pancreas, acinar cells release digestive enzymes into small ducts.

\medterm{Acinetobacter Baumannii Resistance} Resistance means that Acinetobacter baumannii is not killed by one or more antibiotics. It can cause difficult hospital infections, especially in intensive-care units, so culture and susceptibility testing guide treatment.

\medterm{Acinus} A small working unit of a gland made of cells that release their product into a duct. Examples include the secretory units of the pancreas and salivary glands.

\medterm{ACL Injury} A sprain or tear of the strong ligament inside the knee that helps prevent the shinbone from sliding too far forward. It commonly follows twisting or sudden stopping and may cause a pop, swelling, pain, or a feeling that the knee gives way.

\medterm{ACL Reconstruction} Surgery that replaces a torn anterior cruciate ligament inside the knee with a tendon graft. It may restore stability after injury, especially when the knee gives way or the person plans activities requiring pivoting; rehabilitation is required afterward.

\medterm{Acne} A common inflammatory disorder of hair follicles and oil glands. Blocked follicles form blackheads or whiteheads, while inflammation can cause papules, pustules, nodules, scarring, and dark marks, especially on the face, chest, and back.

\synonyms
Acne Vulgaris; Pimples

\medterm{Acne Conglobata} A rare, severe nodulocystic form of acne with groups of inflammatory nodules, cysts, interconnected abscesses, and draining sinuses, usually on the trunk, buttocks, or limbs. It can be persistent and leave atrophic or hypertrophic scars, so early dermatologic treatment is important.

\medterm{Acne Fulminans} A sudden, severe ulcerating form of acne associated with fever, painful joints, malaise, and sometimes weight loss or abnormal blood counts. It requires urgent specialist care; isotretinoin can precipitate or worsen it if introduced too quickly.

\medterm{Acne Inversa} See \textit{Hidradenitis suppurativa}.

\medterm{Acne Keloid} A chronic inflammatory condition causing firm scar-like bumps around hair follicles, usually on the back of the neck.

\synonyms
Acne Keloidalis Nuchae

\medterm{Acne Keloidalis Nuchae (AKN)} Alternate name; see \textit{Acne Keloid}.

\medterm{Acne Rosacea} Acne rosacea is a chronic inflammatory facial skin disorder causing flushing, persistent redness, visible vessels, papules, pustules, or eye symptoms; it is distinct from acne vulgaris.

\medterm{Acne Vulgaris} The common acne disorder caused by follicular plugging, increased sebum production, Cutibacterium acnes activity, and inflammation. It produces comedones, papules, pustules, nodules, or cysts on the face, chest, back, and sometimes upper arms; deeper disease can scar.

\medterm{Acneiform Eruptions} Monomorphic follicular papules and pustules that resemble acne but usually lack comedones. They may be caused by medicines, occupational or chemical exposures, infections, or other disorders; timing, distribution, and trigger help distinguish them from acne vulgaris.

\medterm{Acoustic} Relating to sound, hearing, or how sound waves are produced, transmitted, measured, or received. Acoustic testing can assess hearing, rooms, medical devices, tissues, or signals used during diagnosis.

\medterm{Acoustic Nerve} The acoustic nerve, properly called the vestibulocochlear nerve, carries signals for hearing and balance from the inner ear to the brain. It is the eighth cranial nerve.

\medterm{Acoustic Neuroma} A benign, slow-growing neoplasm arising from Schwann cells of the vestibular division of the vestibulocochlear nerve (cranial nerve VIII) within the internal auditory canal or cerebellopontine angle. It classically presents with insidious, progressive unilateral sensorineural hearing loss, asymmetric high-pitched tinnitus, and dysequilibrium.

\synonyms
Vestibular Schwannoma; Acoustic Neurinoma; Acoustic Neurilemmoma

\medterm{Acoustic Reflex} An automatic tightening of a small middle-ear muscle in response to a loud sound. Testing it helps assess the middle ear, hearing pathways, and facial nerve.

\synonyms
Stapedial Reflex; Stapedius Reflex; Middle-Ear Muscle Reflex

\medterm{Acoustic Shock} Auditory and other symptoms after exposure to a sudden or intense sound, potentially including ear pain, tinnitus, hyperacusis, dizziness, or temporary hearing change.

\medterm{Acoustic Trauma} Acute damage to the auditory system from a brief, intense sound such as an explosion or
firearm. It may cause tinnitus, temporary or permanent sensorineural hearing loss, and
tympanic or ossicular injury; urgent audiologic assessment is indicated when hearing
loss is sudden.

\medterm{Acquired} Developed after birth rather than being present at birth or inherited; an acquired condition may result from infection, exposure, injury, aging, or another later event.

\medterm{Acquired C1-Inhibitor Deficiency Angioedema} A bradykinin-mediated disorder caused by acquired deficiency or dysfunction of C1 inhibitor, often associated with lymphoproliferative or autoimmune disease. Recurrent swelling of the skin, bowel wall, or upper airway usually occurs without hives or itching; low C1q and C1-inhibitor testing help distinguish it from hereditary or histamine-mediated angioedema.

\medterm{Acquired Cystic Kidney Disease} Development of multiple fluid-filled cysts in kidneys with long-standing chronic kidney disease, especially in people receiving dialysis. It can cause blood in the urine, pain, infection, or rarely kidney tumors and is monitored with clinical assessment and imaging.

\medterm{Acquired Feminization} Development of typically female secondary sexual characteristics after birth in a person not expected to develop them, often from hormonal change.

\medterm{Acquired Immunity} Immune protection that develops after infection, vaccination, or transfer of antibodies, involving adaptive immune responses directed toward specific antigens.

\medterm{Acquired Immunodeficiency Syndrome} See \textit{HIV/AIDS}.

\medterm{Acquired Lipodystrophy} A rare disorder characterized by selective, progressive loss of adipose tissue (lipoatrophy) occurring after birth without an inherited genetic defect. Frequently associated with autoimmune diseases, panniculitis, or antiretroviral therapy (e.g., Lawrence syndrome, Barraquer-Simons syndrome), it leads to ectopic lipid accumulation, severe insulin resistance, hypertriglyceridemia, and hepatic steatosis.

\synonyms
Acquired Generalized Lipodystrophy; AGL; Lawrence Syndrome; Acquired Partial Lipodystrophy

\medterm{Acquired Methemoglobinemia} An acquired hematologic disorder in which hemoglobin iron is oxidized from the ferrous ($\text{Fe}^{2+}$) to the ferric ($\text{Fe}^{3+}$) state by exogenous oxidizing agents (such as dapsone, benzocaine, lidocaine, nitrites, or sulfonamides), producing clinical cyanosis, chocolate-brown arterial blood, and cellular hypoxia refractory to supplemental oxygen.

\synonyms
Toxic Methemoglobinemia; Drug-Induced Methemoglobinemia; Chemically Induced Methemoglobinemia

\medterm{Acquired Nasolacrimal Duct Obstruction} Blockage of the tear-drainage system that develops after birth. It may cause excessive tearing, discharge, or repeated inflammation of the lacrimal sac and can result from aging, infection, injury, inflammation, scarring, or a mass.

\medterm{Acquired Platelet Function Defect} Impaired platelet adhesion or aggregation that develops after birth. Causes include antiplatelet medicines, kidney failure, hematologic disease, and other systemic disorders; platelet counts may remain normal, but mucocutaneous bleeding can occur.

\medterm{Acquired Tracheomalacia} Acquired tracheomalacia is weakening of the tracheal wall after injury, inflammation, compression, or prolonged ventilation, causing airway collapse and breathing symptoms.

\medterm{Acquisition} The process by which new information or a skill is learned and initially encoded, especially in discussions of memory and learning.

\synonyms
Learning Acquisition; Memory Acquisition

\medterm{Acquisition Duration} The time required to collect a measurement, image, or body signal. Longer acquisition may improve detail or accuracy but can increase motion problems, discomfort, radiation exposure, or equipment time.

\medterm{Acquisition Time} The time during which a medical imaging system collects data. The required time depends on the imaging method,
body region, protocol, and patient movement.

\medterm{Acral Lentiginous Melanoma} A subtype of melanoma arising on acral skin (palms, soles, or under nails), appearing as
a dark, irregularly pigmented patch or plaque. It is more common in people with darker
skin tones and often presents at a later stage due to less frequent surveillance of
these sites. Prognosis depends on thickness and stage at diagnosis.

\medterm{Acro-Osteolysis} Resorption or destruction of the terminal bones of the fingers or toes. It can occur with inherited disorders, autoimmune disease, severe vascular disease, trauma, or repeated occupational exposure and may cause shortened, painful, or deformed digits.

\medterm{Acrocentric Chromosome} A chromosome whose central joining point lies close to one end, giving it a very short section and a longer section. In humans, chromosomes 13, 14, 15, 21, and 22 have this shape.

\medterm{Acrochordon} Alternate name; see \textit{Skin Tag}.

\medterm{Acrocyanosis} Persistent bluish or purplish discoloration and coolness of the hands, feet, nose, or ears, usually without pain, caused by narrowing of small skin blood vessels. It is often worse in cold conditions and may be primary or due to another disorder.

\medterm{Acrodermatitis} Inflammation affecting skin on the hands, feet, or other body extremities. Causes vary and include infections, nutritional deficiency, inherited disorders, medicines, and immune reactions; appearance and treatment depend on cause.

\medterm{Acrodermatitis Chronica Atrophicans} A late cutaneous manifestation of Lyme borreliosis, usually caused by persistent \textit{Borrelia afzelii} infection in Europe. It begins with unilateral reddish-violet swelling on an extremity and can progress over months or years to thin, atrophic, fibrotic skin; diagnosis uses clinical findings, serology, and sometimes biopsy.

\medterm{Acrodermatitis Enteropathica} A zinc-deficiency disorder classically causing periorificial and acral dermatitis, diarrhea, and alopecia. The primary inherited form results from impaired intestinal zinc absorption; acquired deficiency may follow malnutrition, malabsorption, or other illness, and untreated disease can impair growth, healing, and immunity.

\medterm{Acrodynia} A painful condition with red, swollen, sweaty hands and feet, classically associated with mercury poisoning and historically called pink disease in children.

\medterm{Acrodysostosis} A rare inherited disorder of bone development causing short, broad fingers and toes, a small nose and middle face, and short stature. Some people have learning difficulties or reduced response to parathyroid and other hormones.

\medterm{Acrokeratosis Paraneoplastica} A rare skin condition causing thick, wart-like growths on the hands and feet, often around the nails, as a reaction to an internal cancer. It is most often linked with cancers of the digestive or respiratory tract and needs medical evaluation.

\medterm{Acromegaly} A disorder caused by excess growth hormone after height growth has ended, usually from a noncancerous pituitary adenoma. It enlarges the hands, feet, jaw, and soft tissues and may cause headaches, joint problems, diabetes, or vision changes.

\medterm{Acromioclavicular Joint} The small joint at the top of the shoulder where the collarbone meets the shoulder blade. It allows limited gliding as the arm moves and is held stable mainly by two groups of strong ligaments.

\medterm{Acromioclavicular Joint Injuries} Sprains or dislocations of the joint between the acromion and clavicle, usually caused by a fall or direct
shoulder impact. Severity ranges from ligament sprain to complete displacement and is assessed clinically and with imaging.

\synonyms
Acromioclavicular Joint Injury

\medterm{Acromion} A bony projection from the shoulder blade that forms the highest point of the shoulder. It joins the collarbone and forms a roof over the shoulder joint, with a tendon and lubricating sac passing beneath it.

\medterm{Acroparaesthesia} Acroparaesthesia is tingling, pins-and-needles, burning, or altered sensation in the hands or feet, commonly related to nerve compression, neuropathy, vascular disease, or metabolic causes.

\medterm{Acrophobia} An excessive, persistent fear of heights that causes marked anxiety, avoidance, or impairment beyond the ordinary caution associated with high places.

\medterm{Acropustulosis of Infancy} A recurrent itchy eruption of tiny vesicles and pustules, mainly on the palms and soles of infants and young children. Crops become less frequent and usually resolve by early childhood; scabies may precede or mimic it.

\medterm{Acrosin} An enzyme released from the cap-like structure at a sperm’s tip. It helps sperm pass through the egg’s surrounding layers during fertilization, although successful fertilization requires many additional steps.

\medterm{Acrosome} A cap-like structure on the head of a sperm cell containing enzymes that help it attach to and enter the egg during fertilization.

\medterm{Acrosome Reaction} The release of enzymes from the cap at the front of a sperm cell after it contacts an egg. These enzymes help the sperm pass through the egg’s outer coating; other steps are also needed for fertilization.

\medterm{ACTH Blood Test} A blood test that measures the level of ACTH, a hormone made by
the pituitary gland. It is usually checked with a cortisol test to help find problems
of the adrenal or pituitary glands.

\medterm{ACTH Stimulation Test} A test that shows whether the adrenal glands can make
cortisol. A synthetic form of ACTH is given, and cortisol levels are measured before
and after it is given. A small rise may suggest adrenal insufficiency.

\medterm{Actigraphy} Recording movement over days or weeks with a small wearable sensor, usually worn on the wrist. The pattern of movement helps estimate when a person is asleep or awake and how active they are, but it does not measure sleep directly and cannot by itself diagnose every sleep problem.

\synonyms
Actigraphic Monitoring

\medterm{Actin} A protein that helps cells keep their shape and move. In muscle, actin works with myosin to produce contraction; in other cells, it supports movement, internal transport, and changes in cell structure.

\synonyms
Actin Filament; Microfilament Protein

\medterm{Actin Myofilament} A thin muscle fiber made mainly of actin that interacts with myosin, allowing muscle cells to shorten and produce force during contraction and movement.

\synonyms
Thin Filament

\medterm{Actin-Binding Protein} A protein that attaches to actin filaments and helps build, organize, move, or anchor the cell's internal framework. Examples include myosin, gelsolin, and alpha-actinin.

\medterm{Acting Out} Expressing emotional conflict or distress through behavior rather than words.

\medterm{Actinic} Relating to ultraviolet radiation from sunlight or UV lamps; chronic actinic exposure can cause sunburn, photoaging, and actinic keratoses.

\medterm{Actinic Cheilitis} A precancerous sun-damage condition of the lips, usually the lower lip, caused by long-term ultraviolet exposure. It may cause persistent dryness, scaling, whitening, thickening, or blurring of the lip border and can progress to squamous cell carcinoma.

\medterm{Actinic Granuloma} A usually harmless inflammatory skin condition that occurs on sun-damaged skin. It causes ring-shaped patches with a thinner center and a raised red border, most often on the neck, arms, or backs of the hands.

\medterm{Actinic Keratosis} A rough, scaly, sandpaper-like lesion caused by cumulative ultraviolet damage, usually on sun-exposed skin such as the face, scalp, ears, neck, forearms, or hands. It is a keratinocytic precancer and may progress to cutaneous squamous-cell carcinoma, so changing, thick, tender, or bleeding lesions need assessment and treatment.

\synonyms
Solar Keratosis; AK; Senile Keratosis

\medterm{Actinic Keratosis Pathology} Atypical keratinocytes with disordered maturation in the lower epidermis, often with parakeratosis, alternating orthokeratosis, and solar elastosis. It is a premalignant squamous lesion that can progress to cutaneous squamous cell carcinoma.

\medterm{Actinic Purpura} Bruising caused by fragile, sun-damaged skin and connective tissue, usually on the backs of the hands and
forearms in older adults. The purple patches may appear after minor trauma and are generally harmless, although
medicines such as anticoagulants can worsen them. New, widespread, or unexplained bruising warrants assessment.

\medterm{Actinomycetoma} A slowly worsening bacterial infection of the skin and deeper tissue, usually affecting the foot after soil contamination. It causes swelling, firm lumps, draining channels, and grains in the discharge; untreated disease can damage bone.

\medterm{Actinomycin D} Actinomycin D, or dactinomycin, is an antineoplastic antibiotic that blocks DNA-dependent RNA synthesis and is used for selected cancers; myelosuppression and tissue injury are important risks.

\medterm{Actinomycosis} A long-lasting infection caused by a type of bacteria called Actinomyces that normally lives in the mouth, intestines, and female genital area. It starts when these bacteria get past a damaged surface, such as after a dental problem, dental work, abdominal surgery, or use of an intrauterine device. It most often affects the face and neck, chest, or belly and pelvis, where it forms hard, growing lumps and draining tunnels that contain small yellow grains called sulfur granules. It can spread to nearby tissue and look like cancer. Diagnosis uses cultures or tissue samples, and treatment is a long course of antibiotics, usually high-dose penicillin.

\medterm{Action Potential} A brief electrical signal that travels along a nerve or muscle cell. It begins when the cell’s electrical charge changes enough to trigger the signal and allows nerves to communicate and muscles to contract.

\medterm{Action Tremor} An involuntary rhythmic oscillation of a body part occurring during active voluntary muscle contraction. Subtypes include postural tremors (maintaining a position against gravity), kinetic tremors (moving through an arc), and intention tremors (worsening near the target).

\medterm{Activated Charcoal} A porous form of carbon that adsorbs some substances in the gastrointestinal tract and can reduce
their absorption after selected poisonings. It is not effective for every toxin and can
cause aspiration or other complications.

\medterm{Activated Partial Thromboplastin Time} A blood test that measures clotting through the intrinsic and common coagulation pathways.
It may be prolonged by unfractionated heparin, clotting-factor deficiencies, inhibitors,
or some liver and systemic disorders.

\medterm{Activated Protein C Resistance} An inherited thrombophilic disorder characterized by an impaired anticoagulant response of plasma to activated protein C (APC). Approximately 90 to 95 percent of cases are caused by the factor V Leiden mutation (G1691A), which alters the APC cleavage site on factor V, predisposing affected individuals to recurrent deep venous thrombosis and pulmonary embolism.

\synonyms
APC Resistance; Factor V Leiden Thrombophilia

\medterm{Active Caries} Tooth decay that is continuing to progress. The affected surface is often soft, rough, or chalky and requires dental
assessment and management.

\medterm{Active Management Of Third Stage} A set of measures used after birth to reduce postpartum hemorrhage, usually including a uterotonic
medicine, controlled cord traction when appropriate, and assessment of uterine tone and bleeding.

\medterm{Active Phase Of Labor} The part of the first stage of labor in which the cervix dilates progressively, generally beginning at about 6 cm
and continuing to full dilation. Progress varies between individuals.

\medterm{Active Site} The part of an enzyme where a target molecule attaches and a chemical reaction takes place. Its shape and chemical properties help the enzyme recognize particular molecules and change them into products.

\medterm{Active Surveillance} Close, structured monitoring of a cancer or precancerous condition without immediate
treatment. It includes scheduled examinations, laboratory tests, imaging, or biopsies,
with treatment started if evidence of progression develops.

\medterm{Active Transport} The movement of a substance across a cell membrane using cellular energy. It often moves the substance against a concentration or electrical gradient and helps cells control salts, nutrients, waste products, and other substances.

\medterm{Activities Of Daily Living} Routine tasks required for personal care and independent living, including basic activities and
more complex instrumental activities. Performance is assessed to describe functional
status and support care planning.

\medterm{Acupressure} A complementary treatment in which firm pressure is applied to selected points on the body, often with the fingers. It may provide short-term relief of symptoms such as nausea or pain, but evidence of benefit varies by condition.

\medterm{Acupuncture} A treatment originating in traditional Chinese medicine in which thin needles are inserted into the skin and sometimes moved or electrically stimulated. It may help some pain conditions or treatment-related nausea, but benefits vary.

\medterm{Acute} Describing symptoms or a condition that begins or worsens quickly and usually lasts a limited time. It is commonly contrasted with chronic, which develops or continues over a longer period.

\medterm{Acute Abdomen} Sudden or severe abdominal pain with clinical features suggesting a potentially serious intra-abdominal disorder that may require urgent investigation or surgery.

\medterm{Acute Acalculous Cholecystitis} Acute gallbladder inflammation without gallstones, usually occurring in critically ill, injured, postoperative, or fasting patients. Ischaemia, bile stasis, and infection may contribute; it can progress rapidly to necrosis or perforation and requires urgent assessment.

\medterm{Acute Adrenal Crisis} Acute adrenal crisis is life-threatening cortisol deficiency causing hypotension, vomiting, abdominal pain, electrolyte disturbance, hypoglycemia, or shock; urgent glucocorticoid and fluid treatment is required.

\medterm{Acute Agitation} A state of heightened motor activity, emotional tension, irritability, and impaired
behavioral control that may threaten the patient or others. Management begins with
de-escalation and assessment for delirium, intoxication, withdrawal, medical illness,
and psychosis.

\medterm{Acute Alcohol Intoxication} Neurobehavioral and physiologic impairment caused by recent ethanol exposure. Severity
varies with concentration, tolerance, co-ingestants, and comorbidity; dangerous effects
include respiratory depression, aspiration, hypoglycemia, hypothermia, trauma, and coma.

\medterm{Acute Angle-Closure Glaucoma} An ophthalmic emergency in which sudden closure of the anterior-chamber angle causes a rapid
rise in intraocular pressure. Severe eye pain, blurred vision, halos, headache, nausea,
vomiting, and a red eye may occur; urgent ophthalmic treatment is required to protect
vision.

\medterm{Acute Aortic Syndrome} A group of sudden, life-threatening problems in the wall of the aorta, the body’s largest artery. The wall may tear, bleed within itself, or develop an ulcer, causing severe chest or back pain and requiring emergency assessment.

\medterm{Acute Bacterial Prostatitis} A sudden bacterial infection of the prostate causing fever, chills, pain between the scrotum and anus or in the lower back, painful urination, and sometimes inability to pass urine. It requires prompt antibiotics and medical care; a forceful prostate examination is avoided because it can spread bacteria into the bloodstream.

\medterm{Acute Bacterial Rhinosinusitis} Bacterial infection of the paranasal sinuses suspected when symptoms persist without
improvement for at least 10 days, become substantially worse after initial improvement,
or are severe with high fever and purulent discharge. Treatment decisions depend on
severity, complications, resistance risk, and local guidance.

\medterm{Acute Bilirubin Encephalopathy} Brain dysfunction in a newborn caused by very high bilirubin levels. Poor feeding, sleepiness, weak muscle tone, irritability, a high-pitched cry, unusual movements, seizures, or coma may develop. It is an emergency requiring rapid treatment to prevent permanent brain injury.

\synonyms
Bilirubin Neurotoxicity; Kernicterus Spectrum Disorder

\medterm{Acute Blood Loss} Rapid loss of blood that reduces the amount circulating through the body. It can cause weakness, fast heartbeat, low blood pressure, confusion, organ injury, or shock and requires urgent control of bleeding and replacement of lost volume.

\medterm{Acute Bronchitis} Short-term inflammation of the large airways in the lungs, usually caused by a virus such as a cold or flu. It causes a cough that may bring up clear, yellow, or green mucus and may cause chest soreness, mild fever, tiredness, or wheezing. The cough often lasts 1 to 3 weeks. Because viruses cause more than 90\% of cases, antibiotics usually do not help; care focuses on rest, fluids, and symptom relief.

\synonyms
Chest Cold; Tracheobronchitis

\medterm{Acute Bursitis} Sudden, acute inflammation of a synovial bursa, typically triggered by direct trauma, acute crystal deposition (gout/pseudogout), bacterial infection (most commonly \textit{Staphylococcus aureus}), or unaccustomed repetitive strain. It presents with rapid-onset localized swelling, erythema, warmth, exquisite focal tenderness, and restricted active joint motion.

\synonyms
Acute Bursal Synovitis; Traumatic Bursitis; Septic Bursitis

\medterm{Acute Caries} Tooth decay that progresses rapidly, often producing a soft lesion or cavity. Risk is increased by factors such as frequent
sugar exposure, reduced saliva, or poor oral hygiene.

\medterm{Acute Cerebellar Ataxia} Sudden or rapidly developing incoordination caused by cerebellar dysfunction. In children it may follow an infection, but stroke, intoxication, inflammation, tumor, and other urgent causes must be considered.

\medterm{Acute Chest Syndrome} A potentially life-threatening pulmonary complication of sickle cell disease characterized by a new lung infiltrate plus fever or respiratory symptoms such as chest pain, cough, hypoxemia, or breathing difficulty.

\medterm{Acute Cholangitis} Alternate name; see \textit{Ascending Cholangitis}.

\medterm{Acute Cholecystitis} Acute inflammation of the gallbladder, most often caused by persistent cystic-duct
obstruction by a gallstone. It commonly causes sustained right-upper-quadrant pain,
fever, nausea, and a positive Murphy sign; ultrasound and timely antimicrobial and
surgical assessment are standard.

\medterm{Acute Colonic Pseudo-Obstruction} Severe widening of the large intestine without a physical blockage, usually in people who are seriously ill, injured, or recovering from surgery. It can cause pain, swelling, nausea, and vomiting. If it is not treated, the bowel may lose its blood supply or tear.

\synonyms
Ogilvie Syndrome; Colonic Pseudo-Obstruction

\medterm{Acute Compartment Syndrome} A rapid rise in pressure inside a closed muscle compartment that reduces blood flow and can permanently damage muscles and nerves. It usually follows a fracture, crush injury, or severe swelling and requires emergency assessment.

\synonyms
Compartment Syndrome; Intracompartmental Hypertension

\medterm{Acute Coronary Syndrome} A clinical spectrum of acute myocardial ischemia caused most often by coronary plaque disruption and thrombosis. It includes unstable angina, non-ST-elevation myocardial infarction (NSTEMI), and ST-elevation myocardial infarction (STEMI); ECG findings and serial cardiac troponin distinguish the major presentations and guide urgent treatment.

\medterm{Acute Coronary Syndromes} Another name for \textit{Acute Coronary Syndrome}.

\medterm{Acute Cough} A cough lasting less than three weeks. Common causes include viral respiratory infection, acute bronchitis, and pneumonia.

\medterm{Acute Cutaneous Lupus Erythematosus (ACLE)} A lupus-specific inflammatory skin eruption that commonly presents as a transient malar or generalized erythematous rash, often triggered by sunlight and associated with systemic lupus activity. The malar rash usually spares the nasolabial folds; evaluation should consider systemic symptoms and organ involvement.

\medterm{Acute Cystitis} Acute bacterial infection and inflammation confined to the lower urinary tract and bladder mucosa, typically caused by uropathogenic \textit{Escherichia coli} and presenting with acute dysuria, urinary frequency, urgency, suprapubic pain, and hematuria.

\synonyms
Acute Bacterial Cystitis; Bladder Infection; Acute UTI

\medterm{Acute Dacryocystitis} Acute infection and inflammation of the lacrimal sac, most commonly caused by
obstruction of the nasolacrimal duct leading to stasis and secondary bacterial
infection. It presents with painful swelling, redness, and tenderness at the medial
canthus with purulent discharge.

\medterm{Acute Decompensated Heart Failure} Sudden or rapidly worsening heart failure, causing breathlessness, tiredness, leg swelling, or fluid in the lungs. It may be triggered by infection, a heart attack, rhythm problems, uncontrolled blood pressure, or missed treatment and can require urgent hospital care.

\medterm{Acute Diarrhea} Passage of loose or watery stools, usually three or more times daily, lasting less than two weeks. Infection, food-related
illness, medicines, and other intestinal disorders can cause it. Oral rehydration is important; blood in stool,
severe dehydration, persistent vomiting, or severe pain requires prompt assessment.

\medterm{Acute Disease} An illness that begins suddenly or develops over a short period. It often lasts a limited time, although it can be mild, severe, or life-threatening and may sometimes lead to a longer-term condition.

\synonyms
Acute Illness

\medterm{Acute Disseminated Encephalomyelitis} A rare sudden inflammation of the brain and spinal cord, usually after an infection and less often after vaccination. It can cause confusion, weakness, poor coordination, vision problems, or seizures, especially in children.

\medterm{Acute Ear Infection} A sudden, rapid-onset bacterial or viral infection of the middle ear space behind the eardrum, characterized by earache, fever, hearing muffled sounds, and bulging of the tympanic membrane.

\synonyms
Acute Otitis Media; AOM; Middle Ear Infection

\medterm{Acute Endometritis} Acute inflammation of the endometrium, usually caused by ascending infection or retained tissue after pregnancy-related events or uterine instrumentation; fever, pelvic pain, and abnormal bleeding may occur.

\medterm{Acute Epididymitis} Sudden inflammation of the epididymis, usually caused by infection and sometimes involving the testis. It causes unilateral scrotal pain and swelling and must be distinguished urgently from testicular torsion.

\medterm{Acute Esophageal Necrosis} Ischemic necrosis of the esophageal mucosa, most commonly caused by low-flow states such
as shock, sepsis, or thromboembolism, presenting with hematemesis and chest pain.
Endoscopy shows circumferential black necrotic mucosa, typically in the distal
esophagus.

\medterm{Acute Exacerbation} A rapid worsening of a chronic pulmonary condition, with increased breathlessness, cough, oxygen need, or other symptoms; infection, aspiration, embolism, and other triggers may contribute.

\medterm{Acute Exacerbation Of Asthma} A progressive increase in dyspnea, cough, wheezing, or chest tightness with reduced
expiratory airflow compared with the patient’s baseline. Severe attacks may cause
exhaustion, silent chest, hypercapnia, respiratory failure, and need for urgent
treatment.

\medterm{Acute Exacerbation Of COPD} A sudden worsening of chronic obstructive pulmonary disease, with more breathlessness, cough, or sputum than usual. It is often triggered by a respiratory infection, air pollution, or another stress on the lungs.

\medterm{Acute Fatty Liver Of Pregnancy} A rare, potentially life-threatening liver disorder that usually develops late in pregnancy. Fat builds up in liver cells and may cause nausea, abdominal pain, jaundice, low blood sugar, abnormal bleeding, kidney failure, or other organ failure.

\medterm{Acute Febrile Illness} A sudden illness characterized by fever, often with headache, chills, muscle aches, gastrointestinal symptoms, or other signs of infection; the cause may be infectious or noninfectious.

\medterm{Acute Febrile Neutrophilic Dermatosis (Sweet Syndrome)} An inflammatory neutrophilic skin disease causing abrupt tender red or violaceous papules, plaques, or nodules with fever and neutrophilia. It may be triggered by infection, inflammatory disease, medicines, pregnancy, or an underlying malignancy; biopsy confirms the characteristic neutrophilic infiltrate.

\medterm{Acute Fever in Adults} A fever in an adult with recent onset, commonly caused by infection, inflammation, medication, heat illness, or other conditions. Severity, vital signs, immune status, source, and associated symptoms guide evaluation; confusion, breathing difficulty, stiff neck, shock, or severe localized pain requires urgent care.

\medterm{Acute Flaccid Myelitis} A rare serious illness, usually in children, causing sudden weakness and limpness of one or more arms or legs. Facial weakness, swallowing difficulty, or breathing failure may occur; immediate hospital assessment, imaging, and breathing monitoring are needed.

\medterm{Acute Gastritis} Sudden inflammation or irritation of the stomach lining. It may cause upper-abdominal pain, nausea, vomiting, or bleeding and can result from anti-inflammatory medicines, alcohol, severe illness, infection, or other irritants.

\synonyms
Acute Erosive Gastritis

\medterm{Acute Generalized Exanthematous Pustulosis} A sudden skin reaction, usually to a medicine, that causes many small pus-filled bumps on widespread red skin. Fever may occur. It often starts within days of exposure and improves after the triggering medicine is stopped.

\medterm{Acute Glomerulonephritis} Acute immune or infectious inflammation of the glomeruli producing a nephritic syndrome: hematuria with dysmorphic red cells or red-cell casts, variable proteinuria, edema, hypertension, and reduced kidney function. Causes include postinfectious disease, IgA nephropathy, vasculitis, and lupus; complement, serology, and sometimes kidney biopsy help identify the cause.

\medterm{Acute Granulocytic Leukemia} See \textit{Acute myelogenous leukemia}.

\medterm{Acute Hemorrhagic Edema of Infancy} A usually self-limited cutaneous small-vessel vasculitis of infants, characterized by fever, edema, and large bruise-like or targetoid purpuric lesions on the face and extremities. Internal-organ involvement is uncommon, but serious causes of purpura must be excluded; most cases resolve within one to three weeks.

\medterm{Acute Hemorrhagic Leukoencephalitis} A rare, fulminant inflammatory demyelinating disease of the brain, considered a severe form of acute disseminated encephalomyelitis. It can cause rapidly worsening headache, fever, confusion, seizures, weakness, coma, and brain swelling and requires emergency neurologic care.

\medterm{Acute Hepatitis} Sudden inflammation and injury of the liver, developing over days to months. It may cause tiredness, nausea, pain under the right ribs, dark urine, jaundice, or no symptoms. Causes include viruses, alcohol, medicines, toxins, and autoimmune disease.

\medterm{Acute HIV Syndrome} A flu-like illness that may occur about 2–6 weeks after HIV infection, when the virus is multiplying rapidly. Fever, sore throat, swollen glands, rash, muscle aches, headache, nausea, diarrhea, or mouth ulcers may occur, but some people have no symptoms.

\medterm{Acute Hypercapnic Respiratory Failure} A sudden inability to breathe out enough carbon dioxide, causing it to build up in the blood. Severe lung disease, asthma, weak breathing muscles, sedative medicines, or chest problems can cause confusion, sleepiness, or dangerous breathing failure.

\medterm{Acute Hypoxemic Respiratory Failure} A sudden inability of the lungs to put enough oxygen into the blood. Pneumonia, fluid in the lungs, severe inflammatory lung injury, collapsed lung tissue, or a blood clot in the lungs may cause severe breathlessness, bluish skin, confusion, or organ injury.

\medterm{Acute Idiopathic Polyneuritis} Guillain--Barré syndrome, an acute immune-mediated disorder of peripheral nerves that usually causes progressive symmetric weakness and reduced reflexes and may impair breathing.

\medterm{Acute Illness} A disease or health problem that begins suddenly or over a short period and generally lasts a limited time, although its severity may range from mild to life-threatening.

\medterm{Acute Inflammation} A rapid immune response to injury or infection that usually lasts hours to days. It can cause redness, warmth, swelling, pain, and reduced function as extra blood, fluid, and immune cells enter the affected area.

\medterm{Acute Inflammatory Demyelinating Polyneuropathy} See \textit{Guillain-Barre syndrome}.

\medterm{Acute Inflammatory Demyelinating Polyradiculoneuropathy} The most common form of Guillain–Barré syndrome. The immune system damages the protective covering of nerves, causing progressive weakness, tingling, and reduced reflexes that may spread to the breathing muscles.

\medterm{Acute Intermittent Porphyria} An inherited disorder of haem production caused by reduced hydroxymethylbilane synthase activity. Attacks may cause severe abdominal pain, nausea, constipation, dark urine, nerve symptoms, confusion, or seizures; medicines, hormonal changes, fasting, and illness can trigger them.

\medterm{Acute Interstitial Nephritis} Sudden inflammation of the tissue between the kidney’s filtering structures, most often caused by a reaction to a medicine. Infections and autoimmune disease can also cause it, leading to blood or protein in urine and a sudden decline in kidney function.

\medterm{Acute Ischemic Stroke} Sudden loss of brain, spinal-cord, or retinal function because a blood vessel is blocked and the tissue does not receive enough oxygen. Emergency imaging distinguishes it from bleeding and determines urgent treatment.

\medterm{Acute Kidney Failure} A sudden decline in kidney function that prevents adequate removal of waste, salt, and water. It may cause reduced urine, swelling, nausea, confusion, or dangerous electrolyte and acid changes.

\medterm{Acute Kidney Injury} An abrupt decline in kidney function, defined by a serum-creatinine rise of at least 0.3 mg/dL within 48 hours, a rise to at least 1.5 times baseline within 7 days, or urine output below 0.5 mL/kg/hour for 6 hours. Causes are commonly prerenal hypoperfusion, intrinsic kidney injury, or postrenal obstruction; complications include azotemia, fluid and electrolyte disorders, and acidosis.

\synonyms
AKI

\medterm{Acute Kidney Injury Staging} A standardized classification based on changes in serum creatinine and urine output.
KDIGO stages 1 through 3 quantify the severity of abrupt kidney-function decline and
support prognosis, monitoring, and treatment decisions.

\medterm{Acute Leukemia} A fast-growing blood cancer in which immature abnormal blood cells crowd the bone marrow and interfere with normal blood-cell production. It can cause anemia, infections, and bleeding and needs urgent treatment.

\medterm{Acute Limb Ischemia} Sudden reduction in arterial blood flow to a limb that threatens tissue viability,
usually caused by embolism, thrombosis, or graft occlusion. The six classic findings are
pain, pallor, pulselessness, paresthesia, paralysis, and poikilothermia; urgent vascular
assessment and revascularization are required.

\medterm{Acute Limb Ischemia Six Ps} Six warning findings of suddenly reduced artery blood flow to a limb: pain, paleness, absent pulse, tingling or numbness, weakness or paralysis, and a cold limb. This is an emergency requiring urgent vascular assessment.

\medterm{Acute Liver Failure} Rapid deterioration of liver function in a previously healthy individual without prior
liver disease. Causes include acetaminophen toxicity, viral hepatitis, drug reactions,
autoimmune hepatitis, and Wilson disease. Features include coagulopathy, encephalopathy,
and multiorgan failure. Presents with jaundice, confusion, and bleeding.

\synonyms
Fulminant Hepatic Failure; Acute Hepatic Necrosis

\medterm{Acute Low Back Pain} Lumbosacral pain lasting less than 4 to 6 weeks, most commonly caused by mechanical musculoligamentous strain, facet joint irritation, or acute lumbar disc herniation. Management emphasizes conservative therapy, staying active, and monitoring for neurological red flags.

\synonyms
Acute Lumbago; Mechanical Low Back Strain; Lumbosacral Strain

\medterm{Acute Lymphangitis} Acute inflammation of lymphatic vessels, typically resulting from streptococcal or
staphylococcal infection spreading from a distal wound. It presents with red, tender
streaks extending proximally from the site of infection, along with fever, malaise, and
regional lymphadenopathy. Treatment requires systemic antibiotics.

\medterm{Acute Lymphoblastic Leukemia} Acute lymphoblastic leukemia is a rapidly progressive cancer of immature lymphoid precursor cells that crowds the bone marrow and can involve blood, lymph nodes, or the central nervous system.

\medterm{Acute Lymphocytic Leukemia} A malignant proliferation of immature lymphoid precursors in the bone marrow and blood,
commonly presenting with cytopenias, fever, bone pain, or lymphadenopathy.

\synonyms
ALL; Acute Lymphoblastic Leukemia

\medterm{Acute Lymphoid Leukemia} See \textit{Acute lymphocytic leukemia}.

\medterm{Acute Megakaryoblastic Leukemia} A rare subtype of acute myeloid leukemia characterized by the proliferation of
megakaryoblasts in the bone marrow and peripheral blood.

\medterm{Acute Mesenteric Ischemia} A vascular surgical emergency caused by an abrupt reduction in intestinal blood flow (most commonly from superior mesenteric artery thromboembolism), classically presenting with severe abdominal pain out of proportion to physical examination findings, rapidly progressing to bowel necrosis and peritonitis.

\synonyms
Acute Mesenteric Infarction; AMI

\medterm{Acute Mitral Regurgitation} Sudden backward flow of blood from the left ventricle into the left atrium during systole, caused by acute leaflet, chordal, papillary-muscle, or ventricular injury. It can rapidly produce pulmonary edema, low cardiac output, or shock.

\medterm{Acute Mountain Sickness} A high-altitude illness caused by inadequate acclimatization, usually after ascent above about 2,500\,m (8,200\,ft). It is characterized by headache plus symptoms such as nausea, loss of appetite, dizziness, fatigue, or disturbed sleep. Worsening symptoms, confusion, inability to walk normally, or breathlessness at rest indicate progression to life-threatening high-altitude cerebral or pulmonary edema.

\medterm{Acute Myeloblastic Leukemia} See \textit{Acute Myeloid Leukemia}.

\medterm{Acute Myelogenous Leukemia} Alternate name; see \textit{Acute Myeloid Leukemia}.

\medterm{Acute Myeloid Leukemia} A clonal hematopoietic neoplasm characterized by the rapid proliferation and accumulation of abnormal, poorly differentiated myeloid progenitor cells (blasts) in the bone marrow and peripheral blood. This suppresses normal hematopoiesis, leading to progressive cytopenias manifested as fatigue from anemia, recurrent or opportunistic infections from neutropenia, and petechiae, ecchymoses, or mucosal bleeding from thrombocytopenia.

\synonyms
AML; Acute Myelogenous Leukemia; Acute Granulocytic Leukemia; Acute Non-Lymphocytic Leukemia

\medterm{Acute Myocardial Infarction} Sudden death of heart muscle caused by prolonged inadequate coronary blood flow, usually from an acutely obstructed coronary artery; it is a medical emergency.

\synonyms
AMI

\medterm{Acute Myocardial Infarction Pathology} Histopathologic evolution of myocardial infarction, from early ischemic injury and coagulative necrosis through inflammation, granulation tissue, and scar formation. The findings depend on the interval between infarction and examination.

\medterm{Acute Necrotizing Ulcerative Gingivitis} A painful, rapidly developing infection of the gums characterized by punched-out gingival papillae, ulceration, spontaneous bleeding, gray necrotic tissue, and foul breath. It is associated with anaerobic oral bacteria and risk factors such as poor oral hygiene, smoking, stress, malnutrition, or immunosuppression; urgent dental evaluation, cleaning, and treatment are needed.

\medterm{Acute Nephritic Syndrome} A clinical syndrome of glomerular inflammation marked by hematuria, variable proteinuria, edema, hypertension, and reduced kidney function.

\medterm{Acute Nonlymphocytic Leukemia} See \textit{Acute myelogenous leukemia}.

\medterm{Acute Obstructive Uropathy} Sudden impairment of urinary flow caused by obstruction of the urinary tract,
potentially producing hydronephrosis, pain, infection, and acute kidney injury.

\medterm{Acute Onset Pain} Pain that begins suddenly, often over seconds, minutes, or hours.

\synonyms
Sudden-Onset Pain

\medterm{Acute Otitis Media} Infection of the middle ear space with acute onset of symptoms and signs of middle ear
inflammation. Diagnosis requires middle ear effusion plus acute symptoms (ear pain,
fever, irritability) or signs (bulging tympanic membrane; erythema alone is not diagnostic). Most common
pathogen: Streptococcus pneumoniae.

\medterm{Acute Pain} Pain that begins suddenly and usually lasts for a limited period, often because of injury, illness, or a procedure.

\medterm{Acute Pancreatitis} Sudden inflammation of the pancreas, most often caused by gallstones or alcohol. It usually causes severe upper-abdominal pain that may reach the back, with nausea or vomiting; complications can include fluid loss, infection, organ failure, or pancreatic collections.

\medterm{Acute Pericarditis} Rapid-onset inflammation of the sac around the heart, often idiopathic or viral and sometimes related to autoimmune disease, heart injury, kidney disease, cancer, or medicines. It typically causes sharp chest pain that worsens with breathing or lying down and improves when sitting up and leaning forward; a friction rub, diffuse ST-segment elevation with PR-segment depression, or a pericardial effusion may occur.

\medterm{Acute Poisoning} Harmful effects that develop after a short-term exposure to a toxic substance by swallowing, inhaling, injection, skin contact, or another route. Effects depend on the substance, amount, route, and time of exposure.

\medterm{Acute Proctitis} Sudden inflammation of the rectal mucosa causing rectal pain, urgency, tenesmus, bleeding, mucus, or discharge. Infectious, inflammatory, ischemic, traumatic, and radiation-related causes are possible.

\medterm{Acute Promyelocytic Leukemia} A subtype of acute myeloid leukemia that can quickly cause severe bleeding or abnormal clotting. It is linked to a specific chromosome change and needs immediate specialist treatment.

\medterm{Acute Pulmonary Edema} Acute or rapidly worsening accumulation of fluid in the lungs; see \textit{Pulmonary Edema}.

\medterm{Acute Pulmonary Histoplasmosis} An acute lower respiratory tract infection caused by inhaling microconidia of the dimorphic fungus \textit{Histoplasma capsulatum}, commonly found in soil enriched with bird or bat guano. Manifestations range from mild flu-like symptoms to severe pneumonia with mediastinal lymphadenopathy.

\synonyms
Primary Pulmonary Histoplasmosis; Acute Histoplasmosis Pneumonia; Cave Disease

\medterm{Acute Pyelonephritis} A sudden bacterial infection of the kidney, usually spreading upward from the bladder. It can cause fever, chills, pain in the side or back, nausea, and painful or frequent urination. Prompt medical treatment is important because the infection can spread to the bloodstream and damage the kidney.

\medterm{Acute Radiation Sickness} Illness caused by receiving a high dose of penetrating radiation over a short time, potentially damaging bone marrow, intestines, skin, and other organs.

\medterm{Acute Radiation Syndrome} See \textit{Radiation sickness}.

\medterm{Acute Rehabilitation} Multidisciplinary rehabilitation provided during or soon after hospitalization for conditions such as stroke,
spinal-cord injury, traumatic brain injury, or major trauma. Treatment is tailored to functional needs.

\medterm{Acute Renal Artery Occlusion} Sudden obstruction of blood flow through one or both main renal arteries or their primary branches, most commonly caused by thromboembolism, in situ thrombosis, or trauma. It presents with severe flank or abdominal pain, hematuria, refractory hypertension, and rapid deterioration of kidney function.

\synonyms
Acute Renal Infarction; Renal Artery Thrombosis; Renal Artery Embolism

\medterm{Acute Renal Failure} See \textit{Acute kidney injury}.

\medterm{Acute Respiratory Disease} A general term for an illness affecting breathing that begins suddenly or over a short time. It may include infections, asthma attacks, or other lung problems and is not one specific diagnosis.

\synonyms
ARD

\medterm{Acute Respiratory Distress Syndrome} An acute, diffuse inflammatory lung injury causing hypoxemic respiratory failure and bilateral lung opacities, usually developing within one week of an illness or injury and not primarily explained by heart failure or fluid overload. Common triggers include sepsis, pneumonia, aspiration, trauma, pancreatitis, and major transfusion; treatment addresses the cause and supports breathing with oxygen, lung-protective ventilation, and, in severe cases, prone positioning.

\medterm{Acute Respiratory Failure} A sudden inability of the respiratory system to maintain adequate oxygenation, carbon-dioxide removal, or both. It may result from lung disease, airway obstruction, neuromuscular weakness, or impaired respiratory drive and requires urgent evaluation.

\medterm{Acute Retinal Necrosis} A rapidly progressive necrotizing viral retinitis, usually caused by herpes viruses, with vitreous inflammation, peripheral retinal necrosis, occlusive vasculitis, and risk of retinal detachment.

\medterm{Acute Retroviral Syndrome} The acute illness that may occur shortly after HIV infection, commonly involving fever, rash, sore throat, lymph-node enlargement, muscle aches, or headache.

\medterm{Acute Rhinitis} Sudden inflammation of the nasal lining, commonly caused by a viral infection or an irritant, producing congestion, nasal discharge, and sneezing.

\medterm{Acute Rhinosinusitis} See \textit{Acute sinusitis}.

\medterm{Acute Sinusitis} Acute inflammation of the mucosa of one or more paranasal sinuses, usually following a
viral upper-respiratory infection. It causes nasal obstruction, purulent or discoloured
discharge, facial pressure or pain, and reduced smell; persistent, severe, or worsening
symptoms may indicate acute bacterial rhinosinusitis or a complication.

\synonyms
Acute Rhinosinusitis; Paranasal Sinus Infection

\medterm{Acute Stress Disorder} A trauma-related mental-health condition beginning after exposure to actual or threatened death, serious injury, or sexual violence and lasting from three days to one month. Symptoms may include intrusive memories, avoidance, negative mood, dissociation, and increased arousal.

\medterm{Acute Symptomatic Hyponatremia} A recently developing low blood sodium concentration accompanied by symptoms, particularly neurologic symptoms such as headache, confusion, seizures, or reduced consciousness. Clinical significance depends on the sodium level, rate of change, underlying cause, and associated illness.

\medterm{Acute Transverse Myelitis} Inflammation across a segment of the spinal cord that can cause rapidly developing weakness, sensory changes, and bladder or bowel dysfunction.

\medterm{Acute Tubular Necrosis} Sudden injury and death of cells in the kidney’s tiny tubes, which normally help filter and balance the blood. It can follow prolonged low blood flow, severe illness, or toxic effects from medicines and other substances, causing acute kidney injury.

\synonyms
ATN

\medterm{Acute Urinary Retention} Sudden inability to pass urine despite a full bladder, often causing severe lower-abdominal pain and swelling. Causes include blockage, medicines, nerve problems, or impaired bladder contraction.

\medterm{Acute Viral Illness} A short-duration illness caused by a virus. Symptoms depend on the affected organ system and may
include fever, fatigue, cough, vomiting, diarrhea, or rash.

\medterm{Acute Vision Loss} Sudden reduction or loss of vision in one or both eyes caused by retinal vascular occlusion, retinal detachment, optic-nerve disease, stroke, acute glaucoma, inflammation, trauma, or other emergency. It requires immediate medical or ophthalmic assessment.

\medterm{Acute Weeping Dermatitis} A form of acute dermatitis with inflamed, often blistered skin that releases clear or yellowish fluid and may form crusts. It is a descriptive pattern that can occur with contact dermatitis, eczema, infection, or other inflammatory skin disorders.

\medterm{Acute-On-Chronic Liver Failure} Acute deterioration of pre-existing chronic liver disease associated with organ failure
and high short-term mortality. It is commonly precipitated by infection, alcohol-related
hepatitis, bleeding, or other acute insults and requires urgent multidisciplinary
management.

\medterm{Acute-Phase Reactants} A group of plasma proteins whose circulating concentrations increase or decrease by at least 25 percent during inflammatory states, tissue injury, or infection. Positive reactants include C-reactive protein, ferritin, fibrinogen, and hepcidin; negative reactants include albumin and transferrin.

\synonyms
Acute-Phase Proteins

\medterm{Acyclovir} An antiviral drug that inhibits herpesvirus DNA polymerase and is used for infections caused by herpes simplex virus and varicella-zoster virus.

\medterm{Ad Lib} A prescription or clinical instruction meaning as desired or as tolerated, from the Latin ad libitum.

\medterm{ADA Deficiency} A rare inherited disorder in which cells lack enough adenosine deaminase, an enzyme needed to break down certain cell chemicals. Toxic substances then build up and damage immune cells, often causing severe infections from early infancy.

\medterm{Adactyly} Absence of one or more fingers or toes from birth.

\medterm{Adam's Apple} The common name for the laryngeal prominence, the visible projection at the front of the
neck formed by the angle where the two laminae of the thyroid cartilage meet.

\medterm{Adamantine} Very hard or resembling diamond; in dentistry, relating to tooth enamel or enamel-like tissue. The term may describe an unusually hard structure or the appearance of a lesion containing enamel-forming tissue.

\medterm{Adamantinoma} A rare, low-grade malignant primary bone neoplasm characterized by a distinctive biphasic histology comprising epithelial cells embedded within an osteofibrous stroma. It shows a dramatic anatomical predilection for the anterior cortex of the mid-shaft (diaphysis) of the tibia (accounting for over 85\% to 90\% of cases), occasionally involving the adjacent fibula. Most commonly diagnosed in young adults aged 20 to 40 years, it presents as a slow-growing, painful pretibial swelling or pathological fracture. Plain radiographs characteristically show eccentric, multilocular, radiolucent cortical lesions with surrounding sclerosis producing a ``soap-bubble'' appearance. Because adamantinoma is resistant to chemotherapy and radiotherapy, standard treatment requires wide en bloc surgical resection with negative margins and bone reconstruction.

\synonyms
Adamantinoma of Long Bones; Tibial Adamantinoma; Malignant Angioblastoma of Bone

\medterm{Adams-Stokes Syndrome} A clinical syndrome characterized by sudden, transient episodes of syncope or loss of consciousness caused by abrupt bradyarrhythmias, high-grade or complete atrioventricular block, or ventricular asystole leading to acute cerebral hypoperfusion.

\synonyms
Stokes-Adams Syndrome; Adams-Stokes Attack

\medterm{Adaptation} Adaptation is the adjustment of a cell, organism, or person to a changing internal or external condition; in medicine it may involve physiologic compensation, behavior, or rehabilitation.

\medterm{Adaptive Design} A clinical-study design allowing planned changes based on accumulating data while preserving scientific validity. Changes may involve sample size, treatment groups, or enrollment and must follow a prespecified plan.

\medterm{Adaptive Equipment} Tools or devices modified to help a person perform daily activities safely despite disability, weakness, pain, or sensory loss. Examples include bathroom rails, reachers, adapted utensils, wheelchairs, and communication aids.

\medterm{Adaptive Immunity} The part of the immune system that learns to recognize specific germs or abnormal cells. B cells make targeted antibodies, T cells coordinate or destroy affected cells, and memory cells help produce a faster response to a later exposure.

\medterm{Addiction} A commonly used, non-diagnostic term for compulsive substance use or
recurrent behavior despite harm, with impaired control, craving, or difficulty stopping.
In clinical practice, the specific diagnosis is usually named as a substance-use disorder
or a recognized disorder due to addictive behavior; physical dependence alone does not
establish addiction.

\medterm{Addiction, Alcohol} Legacy label; see \textit{Alcohol Use Disorder}.

\medterm{Addiction, Gambling} Legacy label; see \textit{Gambling Disorder}.

\medterm{Addiction, Nicotine} Legacy label; see \textit{Nicotine Dependence}.

\medterm{Addiction, Sexual} Legacy label; see \textit{Compulsive Sexual Behavior Disorder}.

\medterm{Addison Anemia} An older name for pernicious anemia, a vitamin B12-deficiency megaloblastic anemia caused by inadequate intrinsic factor and impaired B12 absorption.

\synonyms
Addisonian Anemia; Addison's Anemia

\medterm{Addison Disease} Primary adrenal insufficiency caused by damage to the adrenal glands, so they produce too little cortisol and often too little aldosterone. It may cause fatigue, weight loss, low blood pressure, darkened skin, and adrenal crisis.

\synonyms
Addison's Disease
Chronic Adrenal Insufficiency
Primary Adrenal Insufficiency

\medterm{Addison's Disease} Destruction or dysfunction of the adrenal cortex causing deficiency of cortisol and
aldosterone, presenting with fatigue, hyperpigmentation, weight loss, and salt craving.
Lifelong hormone replacement is required.

\synonyms
Hypoadrenocorticism; Chronic Adrenocortical Insufficiency

\medterm{Addisonian Anemia} Alternate name; see \textit{Addison Anemia}.

\medterm{Addisonian Crisis} Life-threatening emergency resulting from acute cortisol deficiency in patients with
adrenal insufficiency. Precipitated by infection, trauma, surgery, or abrupt
discontinuation of chronic glucocorticoid therapy.

\medterm{Additive Genetic Effects} The contributions of individual alleles that combine approximately additively to influence a measurable trait, without requiring interaction between those alleles.

\medterm{Adduct} To move an arm, leg, finger, toe, or other body part toward the center line of the body.

\medterm{Adduction} Movement of a body part toward the body's midline or, for a finger or toe, toward the hand or foot's central axis. For example, bringing an arm back toward the side is shoulder adduction.

\medterm{Adductor Muscle} A muscle that moves a body part toward the body’s midline or, for a digit, toward the axis of the hand or foot.

\medterm{Adductor Muscles} A group of muscles on the inner thigh that pull the leg toward the body’s midline and help stabilize the hip. Strain of these muscles is a common cause of groin pain, especially during sports.

\medterm{Adenine} A purine nitrogenous nucleobase that forms complementary base pairs with thymine in DNA (via two hydrogen bonds) and with uracil in RNA. It also forms essential components of adenosine, cellular energy carriers (ATP, ADP, AMP), second messengers (cAMP), and key metabolic cofactors (NAD+, FAD).

\synonyms
6-Aminopurine; Ade

\medterm{Adenitis} Inflammation of a gland or lymph node, usually caused by infection or an immune disorder. It may cause swelling, tenderness, redness, or fever; treatment depends on the affected site and cause.

\medterm{Adenitis, Mesenteric} See \textit{Mesenteric lymphadenitis}.

\medterm{Adenocarcinoma} A cancer that begins in cells lining or forming glands. It can arise in organs such as the colon, breast, lung, pancreas, prostate, stomach, or uterus and may invade or spread.

\medterm{Adenocarcinoma In Situ Of Cervix} A noninvasive glandular precursor lesion of the cervix, usually associated with high-risk human papillomavirus; abnormal glandular cells remain confined to the epithelium but may progress to invasive adenocarcinoma.

\medterm{Adenocarcinoma Of The Colon} A malignant gland-forming epithelial tumor of the colon, commonly presenting with
altered bowel habits, occult bleeding, anemia, or an obstructing mass.

\medterm{Adenocarcinoma Of The Lung} The most common type of non-small-cell lung cancer. It begins in gland-forming cells, often toward the outer lung, and can occur in smokers or nonsmokers. Some tumors have gene changes, such as EGFR, ALK, or KRAS, that guide treatment choices.

\synonyms
Pulmonary Adenocarcinoma; Bronchial Adenocarcinoma

\medterm{Adenohypophysis} The front part of the pituitary gland, a small hormone-producing organ at the base of the brain. It releases hormones that regulate growth, thyroid activity, adrenal function, reproduction, milk production, and other body processes.

\medterm{Adenoid Basal Carcinoma} A rare, usually indolent cervical carcinoma composed of small basaloid nests, often associated with squamous intraepithelial lesions and requiring distinction from more aggressive tumors.

\medterm{Adenoid Cystic Carcinoma} A rare malignant cancer that arises in secretory glands, most commonly the major and minor salivary glands of the head and neck, as well as the breast, skin, and airways. It tends to grow slowly but is notable for tracking along nerve pathways (perineural invasion), causing pain or numbness, and recurring or spreading to the lungs years after initial treatment.

\synonyms
Cylindroma

\medterm{Adenoid Cystic Carcinoma Of The Vulva} A rare malignant tumor of vulvar Bartholin-type glands, often causing pain because of perineural invasion. It is usually slow-growing but may recur locally or metastasize after long intervals.

\medterm{Adenoid Hypertrophy} Enlargement of the adenoid tissue located in the nasopharynx behind the nasal cavity. Commonly occurring in children due to recurrent respiratory infections or allergies, it can obstruct nasal breathing, leading to persistent mouth breathing, loud snoring, obstructive sleep apnea, a nasal tone of voice, and recurrent middle-ear fluid accumulation (otitis media with effusion).

\synonyms
Enlarged Adenoids; Adenoidal Hyperplasia

\medterm{Adenoid Removal} Alternate name; see \textit{Adenoidectomy}.

\medterm{Adenoidectomy} Surgical removal of the adenoids (pharyngeal tonsil) from the back of the nasal passage. It is commonly performed in children to treat chronic nasal obstruction, pediatric obstructive sleep apnea, recurrent ear infections, or persistent rhinosinusitis that does not improve with medical therapy.

\synonyms
Adenoid Excision; Adenoid Removal

\medterm{Adenoiditis} Acute or chronic inflammation and infection of the adenoid tissue behind the nose. It causes persistent nasal congestion, purulent nasal discharge, mouth breathing, snoring, sore throat, and muffled speech.

\medterm{Adenoids} A collection of lymphatic tissue located in the upper airway at the back of the nasopharynx, forming part of Waldeyer's ring. In infants and young children, they help sample and trap inhaled germs to build immunity, but typically shrink in size during late childhood and adolescence.

\synonyms
Pharyngeal Tonsil; Nasopharyngeal Tonsil

\medterm{Adenolipoma} A benign tumor composed of glandular cells and mature fat cells. It may occur in the skin, breast, thyroid, parathyroid, or salivary glands.

\medterm{Adenolymphoma} A benign salivary-gland tumor, also called Warthin tumor, composed of oncocytic epithelium and lymphoid stroma and occurring most often in the parotid gland.

\medterm{Adenoma} A usually noncancerous gland-forming tumor arising from epithelial tissue. Some adenomas produce hormones or other secretions, and some—especially certain colon and rectal adenomas—contain precancerous changes and may later become cancerous.

\medterm{Adenoma Malignum} A rare, slow-looking but invasive cancer of the cervix that produces mucus. Persistent watery vaginal discharge or abnormal bleeding may occur. Diagnosis requires examination of cervical tissue by an experienced pathologist because the cells may look almost normal under the microscope.

\medterm{Adenoma-Carcinoma Sequence} The usual step-by-step pattern in which some colon polyps develop genetic changes, become increasingly abnormal, and eventually turn into colon cancer over several years. Removing suitable polyps can interrupt this progression.

\medterm{Adenomas} Plural of adenoma; see \textit{Adenoma}.

\medterm{Adenomatosis} A condition in which multiple gland-forming tumors called adenomas develop in one organ or tissue. The tumors are often benign, but their number and location may affect cancer risk.

\medterm{Adenomatous Polyp} A benign gland-forming polyp, most commonly in the colon, that may contain dysplasia and
can progress through the adenoma--carcinoma sequence. Tubular, villous, and
tubulovillous types are recognized.

\medterm{Adenomatous Polyposis} An autosomal dominant genetic disorder caused by germline mutations in the \textit{APC} tumor suppressor gene, characterized by the progressive development of hundreds to thousands of adenomatous polyps throughout the colon and rectum, carrying a nearly 100 percent lifetime risk of colorectal cancer if left untreated.

\synonyms
Familial Adenomatous Polyposis; FAP

\medterm{Adenomyoepithelioma Of The Breast} A rare biphasic breast tumor composed of proliferating epithelial and myoepithelial cells. Most are benign, but local recurrence and malignant transformation can occur, so complete excision is the usual treatment.

\medterm{Adenomyoma} A benign uterine lesion composed of endometrial glands and stroma within a proliferation of smooth muscle; it may cause abnormal uterine bleeding or mimic other uterine tumors.

\medterm{Adenomyosis} A condition in which tissue normally lining the uterus grows into its muscular wall. It can enlarge the uterus and cause heavy or painful periods, ongoing pelvic pain, or no symptoms; it is more common during the later reproductive years.

\medterm{Adenopathy} Disease or enlargement of a gland, most often used for lymph-node enlargement or abnormal lymph-node involvement.

\medterm{Adenosine} A purine nucleoside made of adenine linked to ribose that participates in energy transfer and cell signaling. Given rapidly by injection, it can briefly slow conduction through the atrioventricular node and is used for selected supraventricular tachycardias.

\medterm{Adenosine Deaminase Deficiency} An inherited purine-metabolism disorder in which toxic deoxyadenosine metabolites accumulate and damage lymphocytes, commonly causing severe combined immunodeficiency.

\medterm{Adenosine Diphosphate} A nucleotide formed when ATP loses one phosphate group. ADP can be converted back to ATP to support cellular energy use, and released ADP also activates nearby platelets and promotes platelet aggregation during clot formation.

\synonyms
ADP

\medterm{Adenosine Monophosphate} A small molecule made of adenine, ribose, and one phosphate group. It is formed when cells break down or recycle energy-carrying molecules and helps regulate energy use and cell signaling.

\synonyms
AMP; Adenylic Acid

\medterm{Adenosine Triphosphate} The main short-term energy source inside cells. Cells make it from nutrients and use it to power muscle contraction, movement of substances across cell membranes, chemical reactions, and many other essential processes; it is commonly called ATP.

\medterm{Adenosis} An increase in the number or size of glands or gland-like tissue. It is often noncancerous, but its meaning depends on the organ involved, such as the breast or lung, and the specific pattern.

\medterm{Adenovirus} A group of DNA viruses that commonly cause respiratory illness, including cold-like symptoms, sore throat, fever, bronchitis, croup, or pneumonia. Depending on the type, adenoviruses can also cause conjunctivitis, gastroenteritis, cystitis, or rarely neurologic disease. They spread through respiratory secretions, close contact, contaminated surfaces, and fecal--oral exposure. Most infections are mild, but severe disease is more likely in infants, people with heart or lung disease, and people with weakened immune systems.

\medterm{Adenovirus Infection} Illness caused by an adenovirus, commonly with sore throat, cough, fever, pink eye, diarrhea, or vomiting, especially in children. The virus spreads through droplets, close contact, contaminated surfaces, or infected stool. One pattern, pharyngoconjunctival fever, brings sore throat, red eyes, and fever together and spreads easily in daycare or swimming-pool settings. Most infections clear on their own, but severe breathing or widespread disease can occur in infants, older adults, people with heart or lung disease, and people with weak immune systems. Treatment is supportive, with fluids and fever relief; severe illness with breathing difficulty or dehydration may require hospital care.

\medterm{Adenoviruses} Plural of adenovirus; see \textit{Adenovirus}.

\medterm{Adenylate Cyclase} An enzyme that changes ATP into cyclic AMP, a chemical messenger inside cells. It is often activated by signals received at the cell surface and helps control many cell functions.

\synonyms
Adenylyl Cyclase

\medterm{Adequate Intake} A formal nutritional reference standard within the Dietary Reference Intakes (DRI) system. It represents the recommended average daily nutrient intake level based on observed or experimentally determined approximations in healthy populations, used when scientific evidence is insufficient to calculate an Estimated Average Requirement (EAR) and Recommended Dietary Allowance (RDA) (such as for vitamin K, potassium, or infant nutrition).

\synonyms
AI; Dietary Reference Intake: Adequate Intake

\medterm{Adherence} The extent to which a person's behavior corresponds with agreed healthcare recommendations, including taking medications, following diets, or executing lifestyle changes. Non-adherence may stem from adverse effects, regimen complexity, cognitive impairment, financial constraints, or lack of support.

\synonyms
Compliance; Treatment Adherence; Medication Adherence

\medterm{Adherens Junction} A connection that helps neighboring cells stick together. It links proteins on one cell to the internal support fibers of another, helping tissues keep their shape, strength, and protective barrier.

\medterm{Adhesiolysis} A procedure that cuts or removes bands of scar tissue that abnormally join organs or body surfaces. It may relieve pain or release a blocked intestine, but the bands can form again after surgery.

\medterm{Adhesion} A band of scar-like tissue that abnormally joins internal surfaces or organs, often after surgery, infection, inflammation, or injury, sometimes causing pain or blockage.

\medterm{Adhesion Molecules} Proteins on cells that help cells attach to one another or to the material supporting them. They help organize tissues, guide immune cells to inflamed areas, and can influence how cancer spreads.

\medterm{Adhesive Capsulitis} Alternate name; see \textit{Frozen Shoulder}.

\medterm{Adhesive Tape} A material used to secure dressings, devices, or tubing to the skin. It is available in paper, cloth, plastic, and other forms.

\medterm{Adie Tonic Pupil} A pupil that is unusually large and reacts slowly to light but constricts better when focusing on a near object. It usually affects one eye and results from injury to nerves controlling the pupil; the cause may be unknown or related to another disorder.

\medterm{Adipocere} A grayish-white, waxy, soap-like substance formed in decomposing bodies exposed to moisture and lack of oxygen (such as damp soil or submerged water). It is produced by the bacterial hydrolysis and hydrogenation of body fats into fatty acids (saponification). This process retards putrefaction and can preserve body contours, facial features, and traumatic injuries for months or years, assisting forensic identification.

\synonyms
Grave Wax; Saponification Wax; Mortuary Wax

\medterm{Adipocyte} A fat cell that stores energy and releases hormones and other signals affecting appetite, blood sugar, inflammation, and metabolism. White fat mainly stores energy, while brown and beige fat can produce heat; enlarged or excessive fat cells can contribute to insulin resistance.

\medterm{Adipocytes} Fat cells that store energy mainly as triglycerides and also release hormones and signaling molecules. White adipocytes primarily store energy, while brown and beige adipocytes can produce heat.

\synonyms
Fat Cell

\medterm{Adipogenesis} The process by which immature cells develop into fat cells. It helps form and maintain body-fat tissue, which stores energy and produces hormones; abnormal increases in the number or size of fat cells can contribute to obesity and metabolic disease.

\medterm{Adipokine} A chemical messenger released by body fat. Adipokines help control appetite, blood sugar, fat use, inflammation, blood pressure, and immune activity; abnormal levels may contribute to obesity, diabetes, and heart or blood-vessel disease.

\medterm{Adiponectin} A hormone produced mainly by fat cells that helps regulate fatty-acid use, glucose metabolism, insulin sensitivity, and inflammation. Levels are often lower with obesity and insulin resistance, but a blood level is not diagnostic by itself.

\medterm{Adipose} Relating to body fat or composed of adipose tissue, which stores energy and also functions as an endocrine and immune-active organ.

\medterm{Adipose Tissue} A specialized connective tissue composed mainly of adipocytes that store energy as
triglycerides. It provides thermal insulation, cushioning for organs, and endocrine
regulation through hormones and signaling molecules such as leptin and adiponectin.

\medterm{Adiposis Dolorosa} Adiposis dolorosa, or Dercum disease, is a rare disorder characterized by painful fatty deposits, often with obesity, fatigue, and other systemic symptoms.

\medterm{Adiposity} The amount and location of body fat. Too much body fat, especially around the abdomen, can increase the risk of diabetes, heart disease, high blood pressure, and joint or movement problems.

\medterm{Adipsia} Little or no feeling of thirst despite the body needing fluid. It can lead to dehydration and dangerously high blood sodium, especially when caused by illness affecting the brain, medicines, aging, or severe illness.

\medterm{Adjacent} Next to, close to, or sharing a boundary with another structure, tissue, organ, or lesion. The term helps describe location, spread of disease, or relationships seen on imaging or examination.

\medterm{Adjustment Disorder} Emotional or behavioral symptoms in response to an identifiable stressor, beginning within three months of the stressor. The symptoms cause significant distress or impairment and do not meet criteria for another mental disorder. They generally do not persist for more than six months after the stressor or its consequences have ended.

\medterm{Adjustment Disorder With Anxiety} An emotional or behavioral response to an identifiable stressor that causes distress or
impairment but does not meet criteria for another mental disorder.

\medterm{Adjustment Disorders} Alternate index label; see \textit{Adjustment Disorder}.

\medterm{Adjuvant} An added treatment or substance that enhances the effect of a primary treatment; in vaccines, an adjuvant strengthens the immune response to an antigen.

\medterm{Adjuvant Chemotherapy} Systemic antineoplastic drug therapy administered following definitive primary treatment, such as complete surgical tumor resection, designed to eradicate occult micrometastatic disease and reduce the risk of cancer recurrence.

\medterm{Adjuvant Therapy} Additional treatment given after the main treatment to reduce the chance of disease returning. Examples include chemotherapy, radiation, hormone treatment, or immune treatment after surgery for selected cancers.

\medterm{Adlerian Theory} A psychological theory developed by Alfred Adler, also called individual psychology. It emphasizes feelings of inferiority, goal-directed striving, a person’s style of life, and social interest—the capacity to belong to and contribute to the community. Early family experiences and relationships are considered important influences on development.

\synonyms
Individual Psychology

\medterm{ADME} The four processes that describe what the body does to a medicine: absorption into the bloodstream, distribution through tissues, chemical breakdown, and excretion from the body. These processes influence a medicine’s dose, timing, and effects.
\synonyms
Pharmacokinetics

\medterm{Administered Activity} The measured amount of radioactive material given to a patient for a scan or treatment. The amount is recorded with the radioactive substance, method, date, and time so the exposure and treatment can be checked safely.

\medterm{Administration Procedure} A method used to give a medicine, fluid, vaccine, nutrition, or other treatment, such as by mouth, injection, inhalation, skin application, or infusion. The method affects speed, dose, and risks.

\medterm{Adnexa} Structures next to an organ. In gynecology, the uterine adnexa are the ovaries, fallopian tubes, and supporting ligaments and tissues around the uterus.

\medterm{Adnexal Dysplasia} Abnormal cell development in tissue beside the uterus, such as an ovary or fallopian tube. The change may be related to development or may represent a precancerous process, so its importance depends on the organ and tissue findings.

\medterm{Adnexal Mass} A growth or enlargement involving structures beside the uterus, such as an ovary or fallopian tube.

\medterm{Adnexal Torsion} Twisting of an ovary, fallopian tube, or both around their supporting tissues. The twist can reduce or stop blood flow. It usually causes sudden one-sided pelvic pain, often with nausea or vomiting, and requires urgent medical assessment and usually surgery.

\medterm{Adnexal Tumors} Tumors arising in structures beside the uterus, especially the ovaries, fallopian tubes, or supporting tissues. They may be benign, borderline, or malignant and are evaluated with examination, imaging, laboratory tests, and sometimes surgery or biopsy.

\medterm{Adnexal Tumors And Masses} Abnormal growths or masses involving the ovary, fallopian tube, or nearby supporting tissues; they may be benign or malignant.

\medterm{Adnexitis} Inflammation of one or more uterine adnexa, usually the fallopian tubes and sometimes the ovaries. It commonly occurs as part of pelvic inflammatory disease and may cause lower abdominal pain, fever, abnormal discharge, or pain during sex; untreated infection can lead to an abscess, infertility, or ectopic pregnancy.

\synonyms
Salpingo-Oophoritis

\medterm{Adolescence} The developmental period between childhood and adulthood, involving physical maturation, puberty, and major psychological and social changes.

\medterm{Adolescent Development} The physical, cognitive, emotional, and social changes that occur as a person develops from childhood into adulthood, including puberty and the development of greater independence.

\medterm{Adolescent Health} The health and wellbeing of young people (10-19 years), addressing the physical,
psychological, and social transitions of puberty and adolescence.

\medterm{Adolescent Medicine} A medical specialty focused on the physical, mental, reproductive, and social health needs of adolescents and young adults, including prevention, risk behavior, chronic disease, and confidentiality-sensitive care.

\medterm{Adolescent Obesity} Alternate index label; see \textit{Pediatric Obesity}.

\medterm{Adolescent Schizophrenia} See \textit{Childhood schizophrenia}.

\medterm{Adoptive Cell Therapy} A personalized cellular immunotherapy in which a patient's own lymphocytes are expanded
or engineered ex vivo and reinfused to attack the tumor; examples include
tumor-infiltrating lymphocyte (TIL) therapy and chimeric antigen receptor (CAR) T-cell
therapy.

\medterm{Adoptive Immunotherapy} Treatment in which immune cells are collected, activated or genetically modified, and returned to the patient to attack disease, especially cancer. Fever, inflammation, and other immune reactions may occur.

\medterm{ADP-Ribosylation} A chemical change in which a small molecule is attached to a protein or other cell molecule. This can switch the molecule’s activity on or off and is used by normal cell processes and by some bacterial toxins.

\medterm{Adrenal Cancer} A rare cancer that begins in one of the adrenal glands above the kidneys. Some tumors start in the outer cortex and others in the inner medulla; some produce excess hormones, while others cause pain or other symptoms as they grow.

\medterm{Adrenal Carcinoma} A rare malignant tumor arising from the adrenal cortex. It may produce excess cortisol, aldosterone, androgens, or estrogens, or present as an incidental mass or pressure symptom; hormonal evaluation and imaging help guide diagnosis and staging.

\medterm{Adrenal Cortex} The outer part of each adrenal gland that produces steroid hormones, including cortisol, aldosterone, and adrenal androgens. These hormones help regulate stress responses, blood pressure, salt balance, and metabolism.

\medterm{Adrenal Cortex Physiology} The outer layer of each adrenal gland makes hormones that control salt and blood pressure, stress responses, metabolism, and some sexual development. Its activity is regulated by the kidneys, blood potassium, and the brain’s hormone-control system.

\medterm{Adrenal Cortical Carcinoma} Alternate name; see \textit{Adrenocortical Carcinoma}.

\medterm{Adrenal Crisis} A sudden, life-threatening shortage of adrenal hormones, especially cortisol. It may occur with severe adrenal disease or after abruptly stopping long-term steroid treatment, particularly during infection, injury, or surgery. Symptoms can include severe weakness, vomiting, abdominal pain, confusion, low blood pressure, and collapse; emergency treatment is required.

\medterm{Adrenal Failure} Inadequate production of adrenal hormones, especially cortisol and sometimes aldosterone; severe acute deficiency can cause adrenal crisis with hypotension and shock.

\medterm{Adrenal Gland} One of two hormone-producing glands above the kidneys. The outer cortex makes aldosterone, cortisol, and adrenal androgens; the inner medulla makes adrenaline and noradrenaline. Adrenal hormones help regulate blood pressure, metabolism, stress responses, and salt balance.

\medterm{Adrenal Gland Cancer} Alternate name for adrenal cancer, a malignant tumor arising in an adrenal gland.

\medterm{Adrenal Gland Removal} Surgical excision of an adrenal gland, performed for tumors, hyperplasia, or other adrenal disorders. Approaches include open, laparoscopic, or retroperitoneal; perioperative management requires assessment of hormone production and potential adrenal insufficiency.

\medterm{Adrenal Gland Tumor} A growth in one of the adrenal glands above the kidneys. It may be noncancerous or cancerous and may or may not produce hormones; a pheochromocytoma is one possible type.

\medterm{Adrenal Glands} Adrenal glands sit above the kidneys and produce cortisol, aldosterone, adrenal androgens, and catecholamines that regulate stress, blood pressure, salt balance, and metabolism.

\medterm{Adrenal Hemorrhage} Bleeding into an adrenal gland, often associated with severe infection, serious illness, a bleeding disorder, blood-thinning medicines, or injury. Bleeding into both glands can cause sudden loss of adrenal hormones, very low blood pressure, and shock.

\medterm{Adrenal Hyperplasia} Enlargement or overgrowth of adrenal tissue that may increase production of adrenal
hormones. Causes include congenital adrenal hyperplasia, chronic stimulation by ACTH, and some
adrenal tumors. Symptoms depend on the hormone affected and may include abnormal genital development,
early puberty, excess body hair, high blood pressure, or electrolyte disturbances.

\medterm{Adrenal Incidentaloma} An unexpected mass found in an adrenal gland during imaging for another reason. It may be benign or cancerous and may or may not produce hormones, so evaluation usually includes imaging review and selected hormone tests.

\medterm{Adrenal Insufficiency} A condition in which the body does not make enough cortisol and sometimes not enough aldosterone. Primary adrenal insufficiency results from damage or dysfunction of the adrenal glands; secondary and tertiary forms result from inadequate pituitary or hypothalamic signals, including suppression after prolonged steroid treatment. Features may include fatigue, low blood pressure, weight loss, salt imbalance, or adrenal crisis.

\medterm{Adrenal Mass} An abnormal growth or enlargement in or near an adrenal gland. It may be benign or cancerous and may or may not produce hormones; evaluation usually involves imaging and, when appropriate, hormone tests and specialist review.

\medterm{Adrenal Medulla} The inner part of each adrenal gland, which sits above a kidney. It releases adrenaline and noradrenaline during stress, helping increase heart rate, blood pressure, alertness, and blood flow to muscles.

\medterm{Adrenal Medulla Physiology} The inner part of each adrenal gland releases adrenaline and noradrenaline when the body’s stress-response nerves are activated. These hormones quickly increase alertness, heart activity, blood pressure, and blood flow to muscles.

\medterm{Adrenal Myelolipoma} A usually benign adrenal tumor composed of mature fat and blood-forming tissue. It is often
found incidentally on imaging and causes no symptoms, but larger lesions may produce mass
effect or bleeding. Management depends on size, symptoms, growth, and diagnostic
certainty.

\medterm{Adrenal Nodule} A localized lump in an adrenal gland. It may be a harmless nonfunctioning growth, a hormone-producing tumor, or rarely cancer; assessment depends on its size, imaging features, growth, and hormone production.

\medterm{Adrenal Tumors} Tumors arising in the adrenal cortex or medulla may be benign or malignant and may or may not produce
hormones. Evaluation uses clinical assessment, hormone testing when indicated, and
imaging to assess size and appearance; treatment depends on the diagnosis, hormone
secretion, symptoms, growth, and risk of malignancy.

\medterm{Adrenal Venous Sampling} A catheter-based test that measures hormones in blood from each adrenal gland. It helps determine whether one gland is producing too much aldosterone, which can guide treatment for difficult-to-control high blood pressure.

\medterm{Adrenalectomy} Surgery to remove one or both adrenal glands, usually for a tumor or severe hormone overproduction. Removing one gland may leave enough hormone production; removing both requires lifelong hormone replacement and an emergency steroid plan.

\medterm{Adrenaline} Also called epinephrine, a hormone and neurotransmitter released during stress. It increases heart rate and contractility, raises blood pressure in many vascular beds, relaxes airway muscle, and is the first-line emergency medicine for anaphylaxis.

\medterm{Adrenarche} An early stage of sexual development, typically occurring around age six to eight, marked by increased production and secretion of weak adrenal androgens (such as dehydroepiandrosterone) by the zona reticularis of the adrenal cortex. It contributes to early pubic hair growth, axillary odor, and oily skin.

\medterm{Adrenergic Agent} A medicine or substance that stimulates or blocks receptors responding to adrenaline (epinephrine) and noradrenaline (norepinephrine). These agents can change heart rate, blood pressure, airway size, and pupil size.

\medterm{Adrenergic Agonist} A medicine or substance that copies the effects of adrenaline or noradrenaline by activating specific cell receptors. Depending on the receptor, it may raise blood pressure, speed the heart, open the airways, or improve heart pumping.

\synonyms
Sympathomimetic

\medterm{Adrenergic Antagonists} Medicines or other substances that block adrenaline-like signaling at alpha or beta adrenergic receptors. They can lower blood pressure, slow the heart, reduce cardiac workload, or relax smooth muscle; effects and risks depend on the receptor targeted and the medicine used.

\synonyms
Adrenergic Blockers

\medterm{Adrenergic Bronchodilator Overdose} Excess exposure to a medicine that opens the airways by stimulating the body’s stress-response system. It may cause shaking, agitation, headache, fast heartbeat, low potassium, chest pain, or abnormal rhythms and needs prompt medical assessment.

\medterm{Adrenergic Receptor} A cell-surface receptor activated by adrenaline (epinephrine), noradrenaline (norepinephrine), or related medicines. Alpha and beta receptors help control heart rate, blood pressure, airway size, pupil size, and other body responses.

\medterm{Adrenergic Receptors} Cell-surface receptors activated by adrenaline and noradrenaline. Alpha and beta subtypes regulate cardiovascular, respiratory, metabolic, and other physiological responses.

\medterm{Adrenocortical Carcinoma} A rare cancer that begins in the outer layer, or cortex, of an adrenal gland, one of the glands above the kidneys. Some tumors are functioning and make too much of hormones such as cortisol, aldosterone, or adrenal androgens; others are nonfunctioning and make no excess hormones. Hormone-producing tumors may cause weight gain, high blood pressure, muscle weakness, unusual hair growth, or other changes in sexual development. A nonfunctioning tumor may cause no early symptoms, but a growing or spreading tumor may cause an abdominal lump, abdominal or back pain, or a feeling of fullness.

\medterm{Adrenocorticotropic Hormone} A hormone made by the pituitary gland that tells
the adrenal glands to make cortisol. Cortisol helps the body respond to stress and
helps control blood pressure, blood sugar, and inflammation.

\synonyms
Adrenocorticotrophic Hormone (ACTH); ACTH; Corticotropin

\medterm{Adrenogenital Syndrome} A group of disorders, usually caused by congenital adrenal hyperplasia, in which excess adrenal androgens cause atypical genital development or early puberty.

\medterm{Adrenoleukodystrophy} A group of inherited disorders, most often X-linked, in which very-long-chain fatty acids accumulate and damage myelin, the adrenal glands, or both. Childhood cerebral disease can cause behavior and learning changes, seizures, and loss of neurologic function; adrenal insufficiency may also occur.

\synonyms
ALD

\medterm{Adrenomedullin} A hormone-like substance made by many tissues that widens blood vessels and helps regulate blood pressure, circulation, fluid balance, kidney activity, and responses to injury.

\medterm{Adrenomyeloneuropathy} An adult-onset form of X-linked adrenoleukodystrophy that mainly affects the spinal cord and peripheral nerves. It may cause leg stiffness, weakness, walking difficulty, bladder or bowel problems, and sometimes adrenal insufficiency.

\medterm{Adson Test} A physical examination manoeuvre used to assess for thoracic outlet syndrome and compression of the subclavian artery by a cervical rib or scalene muscles. The examiner palpates the radial pulse while the patient abducts and extends the arm, extends the neck, and rotates the head toward the ipsilateral side with deep inspiration; diminution or obliteration of the radial pulse indicates a positive test.

\synonyms
Adson Manoeuvre; Adson's Test

\medterm{Adsorption} The sticking of molecules to the outside surface of a solid or liquid rather than entering its substance. Activated charcoal works partly by adsorbing some poisons in the digestive tract.

\medterm{Adult Day Services} Community-based daytime programs providing supervision, social activities, health-related support, or assistance with daily activities for adults who cannot safely remain alone throughout the day.

\synonyms
Adult Day Care; Adult Day Program

\medterm{Adult Hypertension} Blood pressure that stays too high over time. The exact level used for diagnosis depends on the guideline and the person's health. Untreated high blood pressure raises the risk of heart disease, stroke, and kidney disease.

\synonyms
Essential Hypertension; Primary Hypertension; Systemic Hypertension

\medterm{Adult Immunisation} The practice of administering vaccines to adults to protect against preventable infectious diseases, maintain waning immunity from childhood vaccinations, and reduce illness in high-risk groups. Essential adult vaccines include annual influenza, tetanus-diphtheria-pertussis (Td/Tdap) boosters every 10 years and during pregnancy, pneumococcal vaccines for adults aged 65 and older or those with chronic medical conditions, the recombinant herpes zoster vaccine for adults aged 50 and older to prevent shingles, and targeted vaccines for COVID-19, hepatitis B, HPV, and RSV.
\synonyms{Adult Immunization; Adult Vaccination; Preventive Adult Vaccines}

\medterm{Adult Informed Consent} The ethical and legal clinical communication process whereby a competent adult patient is provided with comprehensive, understandable information regarding the nature, indication, benefits, material risks, and alternatives of a proposed medical intervention or procedure, enabling an autonomous and voluntary decision.

\synonyms
Informed Medical Consent; Surgical Consent; Patient Autonomous Decision-Making

\medterm{Adult Nutritional Support in Illness} Dietary strategies providing increased caloric and protein intake during acute or chronic illness, convalescence, or hypermetabolic states to prevent involuntary weight loss and muscle wasting.

\synonyms
High-Calorie Diet in Illness; Nutritional Therapy in Recovery

\medterm{Adult Obstructive Sleep Apnea} Recurrent nocturnal collapse or partial obstruction of the pharyngeal airway in adults, producing repetitive episodes of hypopnea or apnea, transient arterial oxyhemoglobin desaturations, sympathetic surge, and fragmented sleep architecture leading to daytime somnolence and increased cardiovascular morbidity.

\synonyms
Adult OSA; Obstructive Sleep Apnea Syndrome; Upper Airway Sleep Apnea

\medterm{Adult Primary Brain Tumor} A neoplasm arising de novo within the brain parenchyma, cranial nerves, or meninges in adults, encompassing diffuse gliomas (astrocytoma, oligodendroglioma, glioblastoma), meningiomas, vestibular schwannomas, and primary central nervous system lymphomas.

\synonyms
Adult Primary Brain Tumor; Primary Intracranial Neoplasm; Adult Glioma

\medterm{Adult Spinal Deformity} Abnormal curvature or alignment of the spine in an adult, including scoliosis, kyphosis, or loss of
normal sagittal balance. It may cause back pain, nerve compression, difficulty walking,
or functional disability.

\medterm{Adult T-Cell Leukemia} A rare and aggressive peripheral T-cell neoplasm caused by infection with the human T-lymphotropic virus type 1 (HTLV-1). Characterized by generalized lymphadenopathy, hepatosplenomegaly, osteolytic bone lesions, severe hypercalcemia, skin infiltration, and atypical multi-lobed circulating lymphocytes known as flower cells.

\synonyms
Adult T-Cell Leukemia-Lymphoma; ATLL

\medterm{Adult-Onset Asthma} Adult-onset asthma is variable airway inflammation and narrowing that begins in adulthood, producing episodic wheezing, cough, chest tightness, or shortness of breath. Diagnosis uses symptoms and lung-function testing, and treatment reduces triggers and airway inflammation.

\synonyms
Late-Onset Asthma; Adult-Onset Reactive Airway Disease

\medterm{Adult-Onset Still Disease} A rare inflammatory disease that causes recurring high fevers, joint pain or swelling, a salmon-colored rash, and sore throat. It may also inflame the liver, heart, lungs, or other organs.

\medterm{Advance Directive} A legal document stating a person's wishes about future medical care and, in some cases, naming a health-care proxy to make decisions if the person cannot decide or communicate. It may address resuscitation, life support, artificial nutrition, and other treatments and is reviewed as circumstances change.

\synonyms
Living Will; Medical Power of Attorney; Healthcare Directive

\medterm{Advance Statements} Alternate name; see \textit{Living Will}.

\medterm{Advanced Cancer} Cancer that has grown beyond its original site. It may be locally advanced when it has
invaded nearby tissues or metastatic when it has spread to distant organs.

\medterm{Advanced Cardiovascular Life Support} A standardized set of advanced resuscitation measures for cardiac arrest and other
life-threatening cardiovascular emergencies. It includes high-quality CPR, defibrillation
when indicated, airway management, medicines, and treatment of reversible causes.

\synonyms
ACLS

\medterm{Advanced Glycation End-Products} Compounds formed when sugar attaches to proteins or fats without the usual enzyme-controlled process. They increase with aging and persistently high blood sugar and may promote inflammation and damage to blood vessels, nerves, kidneys, and other tissues.

\medterm{Advanced Prostate Cancer} Prostate cancer that has grown beyond the prostate or spread to distant parts of the body. It may involve nearby tissues, lymph nodes, bones, or other organs.

\medterm{Advanced Sleep-Wake Phase Disorder} A body-clock disorder in which a person becomes sleepy much earlier than desired and wakes unusually early. The person may sleep normally on this early schedule but cannot stay asleep until a more convenient time.

\medterm{Adventitia} The outer layer of connective tissue around a blood vessel or hollow organ. It supports and anchors the structure and contains small nerves and blood vessels that nourish the wall.

\synonyms
Tunica Adventitia

\medterm{Adventitious} Abnormal, unexpected, or arising from an unusual location or process; adventitious lung sounds include crackles, wheezes, and other added breath sounds.

\medterm{Adverse Effect} An unintended and harmful or bothersome effect of a medicine or other treatment given at normal doses or for an intended purpose.

\medterm{Adverse Event} Any new or worsening health problem that happens during or after a medicine, procedure, or other treatment, whether or not the treatment caused it. Clinicians record its timing, severity, and likely cause separately.

\medterm{Adverse Reaction} An unintended and harmful response to a medicine, treatment, or other exposure. It may be predictable from the medicine's effects, such as bleeding with an anticoagulant, or unpredictable, such as an allergic reaction. Severity ranges from mild nausea or rash to serious organ injury, breathing difficulty, or shock; suspected severe reactions need urgent medical attention.

\synonyms
Adverse Drug Reaction

\medterm{Aedes Aegypti} Aedes aegypti is a mosquito that can transmit dengue, yellow fever, chikungunya, Zika, and other viruses.

\medterm{Aegophony} An altered voice sound heard through a stethoscope when a person says “E” but it sounds like “A.” It can occur over lung tissue made denser by infection or fluid, but it does not identify the cause by itself.

\medterm{Aerobe} A microorganism that grows or produces energy using oxygen. Some aerobes cannot live without oxygen, while others can grow either with or without it. The distinction helps laboratories understand how an organism may behave in the body.

\medterm{Aerobic} Requiring or using oxygen. In microbiology, an aerobic organism grows with oxygen; in the body, aerobic energy production uses oxygen to release energy from nutrients for muscles and other cells.

\medterm{Aerobic Bacteria} Bacteria that need oxygen to grow. They may cause infection in tissues exposed to air, such as the lungs, skin, or surface of a wound; laboratory tests identify the specific species.

\medterm{Aerobic Bacterium} An aerobic bacterium grows using oxygen, although some species can tolerate or use alternative metabolic pathways; oxygen requirements help guide culture and infection assessment.

\medterm{Aerobic Exercise} Physical activity using large muscles continuously while the body produces energy with oxygen. Walking, cycling, swimming, and dancing can improve heart, lung, blood-vessel, muscle, mood, and endurance health.

\synonyms
Cardio Exercise

\medterm{Aerobiosis} Life or biological activity that occurs when oxygen is available. Many human cells use oxygen to release energy from food, and some microorganisms need oxygen to grow.

\medterm{Aerophagia} Excessive swallowing of air, which may cause frequent belching, abdominal bloating, discomfort, or increased intestinal gas. It can occur with rapid eating, chewing gum, anxiety, or swallowing disorders.

\medterm{Aerophagy} Repeated swallowing of too much air, which can cause frequent belching, abdominal fullness, bloating, or discomfort. It may occur with rapid eating, chewing gum, anxiety, or difficulty swallowing.

\synonyms
Aerophagia

\medterm{Aerophobia} An excessive fear of flying or air travel that causes marked anxiety, avoidance, or interference with daily life. It is a specific phobia and may improve with structured psychological treatment.

\medterm{Aerosol} Very small solid particles or liquid drops floating in a gas such as air. Aerosols can carry inhaled medicines, pollutants, or infectious particles and may remain in the air or settle on surfaces.

\medterm{Aerospace Medicine} The field of medicine concerned with maintaining the health, safety, and performance of people who travel or work in aircraft and spacecraft. It addresses effects such as changes in pressure, acceleration, low oxygen, radiation, isolation, and microgravity.

\medterm{Aerotitis} Ear pain or middle-ear injury caused by failure to equalize pressure during rapid changes in altitude, such as during aircraft ascent or descent.

\medterm{Aesthetic Dentistry} Dental care intended to improve the appearance of the teeth, gums, or smile. Procedures may include
tooth-colored restorations, veneers, whitening, or selected crowns.

\medterm{Aesthetic Medicine} Medical care intended to improve appearance or reduce visible signs of aging through procedures, devices, medicines, or skin care; appropriate assessment includes risks, contraindications, and realistic expectations.


\synonyms
Cosmetic Medicine; Aesthetic Dermatology

\medterm{Aetiology} The study of causation, or the specific originating cause or combination of pathogenic factors (such as genetic, environmental, infectious, or autoimmune mechanisms) that lead to the development of a disease or clinical disorder.

\synonyms
Etiology; Causation; Pathogenesis

\medterm{Affect} The observable expression of a person’s emotional state, including facial expression, voice, posture, and emotional responsiveness.

\medterm{Affective Disorder} Alternate name; see \textit{Mood Disorders}.

\synonyms
Mood Disorder

\medterm{Affective Flattening} Markedly reduced outward expression of emotion, including limited facial expression, voice modulation, and gestures.
It may occur in schizophrenia and other psychiatric or neurologic conditions.

\medterm{Afferent} Carrying signals toward a central organ or structure; afferent nerves transmit sensory information from the body toward the brain and spinal cord.

\medterm{Afferent Arteriole} A small branch of the renal arterial circulation that conveys blood directly into the glomerular capillary tuft of an individual nephron, playing a pivotal role in regulating glomerular filtration pressure and renal blood flow through autoregulatory and tubuloglomerular mechanisms.

\medterm{Afferent Loop Syndrome} A blockage or poor emptying in a section of small intestine made during some stomach operations. Digestive juices build up behind the blockage, causing upper-abdominal pain, nausea, vomiting, or poor absorption of nutrients.

\medterm{Afferent Lymphatic Vessel} A lymphatic vessel that carries lymph into a lymph node.

\medterm{Afferent Neuron} A sensory nerve cell carrying signals from receptors in the body toward the spinal cord and brain, where information is processed and appropriate responses begin.

\medterm{Afferent Pathway} A sensory nerve route carrying information from the skin, muscles, organs, or other tissues toward the spinal cord or brain, where the signals are processed and responses can begin.
\synonyms
Sensory Pathway

\medterm{Afferent Pupillary Defect} An asymmetry in pupillary light reactivity caused by a unilateral or asymmetric lesion in the anterior visual afferent pathway (optic nerve or retina). Identified during the swinging-flashlight test when swinging the light beam from the unaffected eye to the affected eye results in paradoxical bilateral pupillary dilation rather than constriction.

\synonyms
Marcus Gunn Pupil; Relative Afferent Pupillary Defect; RAPD

\medterm{Affinity} The strength of attraction or binding between two molecules, such as an antibody and antigen, receptor and ligand, or enzyme and substrate.

\medterm{Affinity Chromatography} A laboratory method for separating and purifying a substance by allowing it to stick to a partner designed to recognize it. Other substances are washed away, and the desired substance is then released.

\medterm{AFib} Alternate name; see \textit{Atrial Fibrillation}.

\medterm{Afibrinogenaemia} Alternate spelling; see \textit{Afibrinogenemia}.

\medterm{Afibrinogenemia} A rare congenital or severe acquired coagulation disorder characterized by the complete or near-complete absence of circulating fibrinogen (factor I, <0.1 g/L). Inherited in an autosomal recessive pattern (mutations in \textit{FGA}, \textit{FGB}, or \textit{FGG}), it manifests with umbilical cord bleeding in neonates, spontaneous mucosal and intracranial hemorrhages, hemarthroses, and prolonged bleeding times.

\synonyms
Congenital Afibrinogenemia; Fibrinogen Deficiency; Afibrinogenaemia

\medterm{Aflatoxin} A toxin made by certain molds that can contaminate foods such as grains, nuts, and spices. Long-term exposure can damage the liver and increase the risk of liver cancer.

\medterm{African Sleeping Sickness} An infection caused by \textit{Trypanosoma brucei} parasites and spread by infected tsetse flies in sub-Saharan Africa. It may cause fever and swollen lymph nodes, followed by sleep changes and nervous-system problems; untreated disease can be fatal.

\medterm{African Trypanosomiasis} An infection caused by \textit{Trypanosoma} parasites and spread by infected tsetse flies in sub-Saharan Africa. It may cause fever and swollen lymph nodes, followed by sleep changes and nervous-system problems; untreated disease can be fatal.

\synonyms
African Sleeping Sickness

\medterm{After-Cataract} Clouding of the thin membrane behind the artificial lens after cataract surgery. It can develop months or years later and cause blurred or hazy vision, glare around lights, or reduced sharpness of vision.

\synonyms
Posterior Capsule Opacification; Secondary Cataract

\medterm{Afterbirth} The placenta and the membranes that surrounded the baby, which leave the womb after childbirth. If any part remains inside, it can cause heavy bleeding or infection.

\medterm{Afterdepolarization} An abnormal electrical oscillation occurring during or immediately following myocardial action potential repolarization, capable of reaching threshold and triggering cardiac arrhythmias. Early afterdepolarizations occur during plateau or repolarization phases, while delayed afterdepolarizations occur after repolarization is complete.

\medterm{Afterload} The pressure or resistance the heart must overcome to push blood out of its ventricles. Higher blood pressure or narrowing of blood vessels increases the heart’s workload and can make pumping less efficient.

\medterm{Afterpains} Cramp-like pains in the lower abdomen during the first few days after childbirth, caused by the womb tightening and shrinking. They may become stronger during breastfeeding and help reduce bleeding.

\medterm{Aftershave Poisoning} Harm from swallowing aftershave, which may contain alcohol, fragrance, or other chemicals. It can cause vomiting, drowsiness, low blood sugar, breathing problems, or coma; significant ingestion is an emergency requiring prompt poison-center assessment and medical treatment.

\medterm{Agammaglobulinaemia} A serious immune deficiency in which the blood contains very few or no antibodies, usually because antibody-producing B cells do not develop normally. It causes repeated bacterial infections, often beginning in infancy or early childhood.

\medterm{Agammaglobulinemia} A serious immune deficiency in which the blood contains very few or no antibodies. It often results from abnormal B-cell development and causes repeated bacterial infections of the ears, sinuses, lungs, or bloodstream.

\medterm{Agar} A jelly-like substance made from red seaweed. Laboratories add it to nutrient mixtures to grow and identify bacteria and fungi on solid surfaces, such as culture plates. It is also used in some bulk laxatives.

\medterm{Agatston Score} A number calculated from a special CT scan that measures calcium in the heart’s arteries. A higher score usually means more coronary plaque and higher heart-disease risk, but it does not directly show blockage. Doctors interpret it with age, sex, blood pressure, cholesterol, diabetes, smoking, family history, and other health conditions.

\synonyms
Coronary Artery Calcium (CAC) Score; Agatston Calcium Score

\medterm{Age Spots} Flat, tan, brown, or black spots that develop on chronically sun-exposed areas of the skin, most commonly the face, backs of the hands, forearms, and shoulders. Caused by localized proliferation of pigment cells (melanocytes) from years of ultraviolet (UV) radiation exposure and aging, they are benign and harmless. Despite the colloquial name ``liver spots,'' they have no connection to liver function or liver disease.

\synonyms
Liver Spots; Solar Lentigines; Senile Lentigines; Sun Spots

\medterm{Age-Related Hearing Loss} Gradual hearing loss in both ears caused by aging. It often affects high-pitched sounds first and can make speech harder to understand.

\medterm{Age-Related Macular Degeneration, Dry} See \textit{Dry macular degeneration}.

\medterm{Age-Related Macular Degeneration, Wet} See \textit{Wet macular degeneration}.

\medterm{Age-Standardized Rate} A disease or death rate recalculated as if different populations had the same age distribution. This makes comparisons fairer when one population has more older or younger people than another.

\medterm{Ageing} Ageing is the natural, lifelong process of physical and biological change. Over time, body reserves may decrease and the risk of some diseases may increase, but the effects vary between people.

\synonyms
Aging

\medterm{Agenesis} Failure of an organ or body structure to form before birth, so it is absent. Agenesis may affect one organ, such as a kidney, or a specific body structure, such as the corpus callosum.

\medterm{Agenesis Of The Corpus Callosum} A birth condition in which the band of nerve fibres connecting the brain’s two halves is partly or completely absent. Effects vary widely and may include seizures, developmental delay, learning difficulties, or no noticeable symptoms.

\medterm{Ageusia} Complete loss of the ability to taste. It may result from infections, mouth or nose problems, medicines, chemotherapy, head injury, nerve damage, nutritional deficiency, or neurological disease. Loss of smell can also make food seem tasteless.

\medterm{Agger Nasi} The most anterior and superior ethmoidal air cell, located immediately anterior to the attachment of the middle nasal turbinate and serving as an important anatomical landmark in functional endoscopic sinus surgery.

\medterm{Agglutination} Visible clumping that occurs when antibodies or other binding substances join cells or particles together. Laboratory workers use this reaction to help identify blood groups, infections, and other substances in patient samples.

\medterm{Agglutination Test} A laboratory test that mixes a patient sample with a known antibody or antigen. Visible clumping suggests that the matching substance is present and may help identify blood groups or infections.

\medterm{Agglutinin} An antibody or other substance that causes cells or particles to stick together in visible clumps. Agglutinins are used in laboratory tests, including blood-group testing and some infection tests.

\medterm{Agglutinogen} An antigen on a cell or particle that can react with a matching antibody and cause visible clumping. Agglutinogens are especially important in identifying blood groups.

\medterm{Aggression} Words or actions intended to threaten, harm, or intimidate another person or damage property. Aggression may result from emotional distress, pain, mental illness, brain disorders, medicines, or substance use.

\medterm{Aggressive} Describing hostile or harmful behavior, or a disease that grows, damages nearby tissue, or spreads quickly. The meaning depends on whether it describes a person’s behavior or a medical condition.

\medterm{Aggressive Angiomyxoma} A rare, usually slow-growing soft-tissue tumor that most often develops in the vulva, perineum, or pelvis of women of reproductive age. It can grow into nearby tissue and often returns after treatment, but rarely spreads to distant organs.

\medterm{Aggressive Cancer} A cancer that is likely to grow, invade nearby tissue, or spread quickly. Doctors assess its aggressiveness using the cancer type, cell appearance, genetic features, stage, growth rate, and response to treatment.

\medterm{Agility} The ability to move, stop, or change direction quickly while staying balanced and controlled. Doctors and therapists may assess agility when evaluating movement, coordination, recovery from injury, or physical function.

\medterm{Agitation} An uncomfortable state of inner tension, restlessness, or increased activity. It may cause pacing, irritability, or inability to remain calm and can result from medical illness, pain, medicines, substance use, or mental health conditions.

\medterm{Agnathia} A rare birth condition in which the lower jaw is partly or completely absent. It may cause major facial abnormalities and serious problems with breathing, swallowing, and feeding.

\medterm{Agnosia} An inability to recognize familiar objects, people, sounds, or other sensations even though the related sense works normally. It usually results from damage to areas of the brain that interpret sensory information.

\medterm{Agonist} A drug or other ligand that binds to a receptor and activates it, producing a biological response. A full agonist can produce the maximum response in a given system, whereas a partial agonist produces a submaximal response even when occupying receptors.

\medterm{Agonist Muscle} The main muscle responsible for producing a particular movement. It is also called the prime mover.

\medterm{Agoraphobia} An anxiety disorder involving persistent, intense fear of places or situations where escape may be difficult, help may be unavailable, or panic-like symptoms may feel embarrassing. People may avoid crowds, public transportation, shops, or leaving home alone; the fear and avoidance cause significant distress or interfere with daily functioning.

\medterm{Agranulocytes} A group of white blood cells that do not have prominent granules when viewed under a microscope. The group mainly includes lymphocytes and monocytes, which help defend the body against infection.

\medterm{Agranulocytosis} A serious condition in which the blood has an extremely low number of granulocytes, especially neutrophils, white blood cells that fight infection. It may be caused by medicines, chemotherapy, infection, autoimmune disease, or bone-marrow problems and can cause fever, sore throat, and severe infections.

\medterm{Agraphesthesia} An inability to recognize a number or letter traced on the skin with the eyes closed, despite normal basic touch sensation. It may indicate damage to the opposite parietal lobe of the brain.

\medterm{Agraphia} Loss of a previously learned ability to write because of brain damage or disease. It may follow a stroke, head injury, brain tumour, or neurodegenerative disease and may occur with problems speaking or reading.

\medterm{Aicardi Syndrome} A rare disorder affecting mainly girls, usually involving abnormal or absent tissue connecting the brain’s two halves, distinctive retinal defects, and seizures beginning in infancy. Other brain and eye abnormalities may also occur.

\medterm{Aicardi-Goutieres Syndrome} A rare inherited disorder in which abnormal immune activity causes inflammation and damage mainly in the brain, skin, and immune system. It may begin in infancy with developmental regression, seizures, brain changes, and painful skin lesions.

\medterm{AIDS} The most advanced stage of HIV infection, when the immune system is severely weakened. AIDS is diagnosed when the CD4 count is below 200 cells/mm$^3$ or certain serious infections or cancers develop.

\synonyms
Acquired Immunodeficiency Syndrome; Stage 3 HIV Infection

\medterm{AIDS Dementia} A severe form of HIV-associated brain disease occurring mainly in advanced HIV infection. It may cause worsening memory and thinking, slowed movement, difficulty walking, and changes in behaviour.

\medterm{AIDS-Associated Nephropathy} Kidney disease caused by HIV, usually in untreated or poorly controlled infection. It often causes large amounts of protein in the urine and may rapidly worsen kidney function, leading to kidney failure.

\synonyms
HIV-Associated Nephropathy

\medterm{Ainhum} A rare disorder in which a painful fibrous band slowly tightens around a toe, usually the little toe. The band can damage bone and blood supply and may eventually cause loss of the affected part.

\synonyms
Dactylolysis Spontanea

\medterm{Air Abrasion} A dental technique that uses compressed air to spray fine abrasive particles onto a tooth. It can remove limited decay, stains, or old filling material and prepare the tooth for a restoration.

\medterm{Air Bag} An inflatable vehicle safety device that rapidly expands during a crash to cushion impact and reduce serious injury. It works with, but does not replace, a properly worn seat belt.

\medterm{Air Bronchogram} A branching pattern of air-filled airways seen within an area of lung that has lost its normal air, usually because the air sacs contain fluid, pus, blood, or cells. It is a sign on a chest X-ray or CT scan.

\medterm{Air Contaminants} Harmful substances in indoor or outdoor air, including smoke, dust, gases, chemicals, allergens, and germs. Health effects depend on the substance, amount breathed, duration of exposure, and a person’s health.

\medterm{Air Embolism} A blockage of blood flow caused by air or another gas bubble entering a blood vessel. It can occur during medical procedures, injury, or rapid pressure changes and may cause sudden breathing, heart, or brain problems.

\medterm{Air Passages} The connected tubes that carry air from the nose and mouth to the lungs. They include the throat, voice box, windpipe, and branching airways; swelling, mucus, or an object can narrow or block them.

\medterm{Air Polishing} A dental procedure that directs air, water, and a fine powder onto teeth to remove plaque and surface stains. Common powders include glycine, erythritol, and sodium bicarbonate.

\medterm{Air Pollution} Air contaminated with harmful gases, particles, smoke, chemicals, or germs. Breathing polluted air can irritate the eyes and lungs and increase the risk of asthma attacks, infections, heart disease, stroke, and other health problems.

\medterm{Air Quality} A measure of how clean or polluted the air is, based on the amounts of harmful substances it contains. Poor air quality can affect breathing and heart health, especially in people with existing medical conditions.

\medterm{Air Sac} A small, thin-walled space that contains air. In the human lung, an air sac is called an alveolus and is where oxygen enters the blood and carbon dioxide leaves it; birds also have larger air sacs that help breathing.

\synonyms
Air Space

\medterm{Air Trapping} A condition in which extra air remains in the lungs after breathing out because narrowed or damaged airways prevent complete exhalation. It commonly occurs in asthma, chronic obstructive pulmonary disease, and bronchiolitis.

\medterm{Air-Conditioner Lung} Lung inflammation caused by an immune reaction to mold, fungi, bacteria, or other particles released from contaminated air-conditioning or humidifying systems. Symptoms may include cough, shortness of breath, fever, and tiredness.

\medterm{Air-Sickness} Motion sickness caused by the movement of an aircraft. Mismatched signals from the eyes and inner ear can cause nausea, dizziness, sweating, headache, and sometimes vomiting.

\medterm{Airborne Disease} An infection that can spread when a person breathes in tiny particles containing germs that remain in the air. The risk depends on the germ, ventilation, distance, and length of exposure.

\medterm{Airborne Precautions} Measures used to prevent infections that remain suspended in the air from spreading. They include a special single-patient room with filtered exhaust air, limiting movement outside the room, and a properly fitted respirator such as an N95. They are used for tuberculosis, measles, and chickenpox.

\medterm{Airplane Ear} Ear pain, fullness, popping, or temporary hearing loss caused by unequal air pressure on the two sides of the eardrum during takeoff or landing. It is more likely with a blocked nose or ear infection.

\synonyms
Barotitis Media

\medterm{Airway} Any passage that carries air between the outside of the body and the lungs, including the nose, mouth, throat, voice box, windpipe, and branching tubes. Swelling, mucus, or an object can narrow or block an airway and restrict breathing.

\medterm{Airway Clearance Techniques} Methods used to loosen and remove mucus from the breathing tubes. They include controlled coughing, breathing exercises, drainage positions, devices that create resistance during exhalation, and vibrating vests. The choice depends on the illness and the person’s strength and ability.

\medterm{Airway Management} The methods and equipment used to keep a person’s airway open and support breathing when airflow is threatened. They may include positioning, suction, airway devices, or a breathing tube connected to a ventilator.

\medterm{Airway Obstruction} Partial or complete blockage of airflow through the breathing passages. Causes include an inhaled object, mucus, swelling, injury, tumour, or relaxed airway muscles. It may cause noisy breathing, wheezing, low oxygen, severe breathing difficulty, or respiratory arrest.

\medterm{Airway Resistance} The difficulty air encounters while moving through the breathing tubes. Narrowing from muscle tightening, swelling, mucus, or loss of support increases resistance and makes the breathing muscles work harder.

\medterm{Akathisia} A movement disorder causing an uncomfortable inner restlessness and a strong urge to move, often with pacing, rocking, or repeated shifting of the legs. It is commonly a side effect of antipsychotic or other medicines.

\medterm{Akinesia} An inability to initiate or perform voluntary movement, or a marked lack of spontaneous movement. It may occur in Parkinson disease and related disorders, severe depression, or other conditions affecting the brain’s movement-control systems.

\medterm{Akinetic} Describing a person or condition with little or no voluntary movement. It may result from a neurological disease, medicine effect, severe illness, or damage to the brain’s movement-control systems.

\medterm{Akinetic Mutism} A state in which a person is awake but makes little or no voluntary movement or speech. It usually results from severe damage or dysfunction in brain areas that control motivation, attention, and movement.

\medterm{Akinetopsia} A rare neuropsychological visual deficit characterized by an inability to perceive visual motion, resulting in moving objects appearing as a succession of static freeze-frames. It typically results from bilateral cortical lesions in the middle temporal visual area (area MT/V5).

\synonyms
Motion Blindness; Visual Motion Agnosia

\medterm{Alagille Syndrome} A genetic disorder in which too few or abnormal bile ducts can cause liver disease and jaundice. It may also affect the heart, kidneys, eyes, blood vessels, bones, and facial features. Most cases involve changes in the JAG1 gene.

\medterm{Alanine} An amino acid, or protein-building molecule, made by the body and obtained from food. Cells use it to make proteins and to help move nitrogen and produce glucose for energy.

\medterm{Alanine Aminotransferase} An enzyme found mainly inside liver cells. When liver cells are injured, alanine aminotransferase can leak into the blood, causing a raised test result. The result helps assess liver injury but does not identify its cause by itself.

\synonyms
ALT

\medterm{Alanine Transaminase Blood Test} A blood test that measures alanine aminotransferase, an enzyme released when liver cells are damaged. A high result may occur with hepatitis, medicine-related injury, fatty liver disease, or other conditions and must be interpreted with other information.

\medterm{ALARA Principle} A foundational radiation safety standard mandating that occupational and medical ionizing radiation exposures be kept as low as reasonably achievable, utilizing dose optimization, exposure time reduction, increased distance, and protective shielding.

\synonyms
As Low As Reasonably Achievable; ALARA; Radiation Safety Principle

\medterm{Alarmin} A signal released by injured or severely stressed cells that alerts the immune system. It helps start inflammation and attract protective cells to damaged tissue, even before the body identifies whether infection is present.

\medterm{Albendazole} An antiparasitic medicine that damages the internal structures parasites need to absorb nutrients and survive. It treats several worm infections, including some intestinal and tissue infections, when prescribed for the appropriate parasite.

\medterm{Albers-SchöNberg Disease} A usually inherited form of osteopetrosis in which bones become unusually dense because old bone is not removed normally. The bones can still break easily and may cause pain or pressure on nerves, especially in the skull.
\synonyms
Autosomal Dominant Osteopetrosis Type II (ADO II); Marble Bone Disease

\medterm{Albinism} A group of inherited genetic conditions in which the body produces little or no melanin, the pigment responsible for skin, hair, and eye coloration. Characterized by pale skin, white or light blond hair, and light-colored irises, it frequently causes visual impairments—including reduced visual sharpness, rapid involuntary eye movements (nystagmus), extreme sensitivity to bright light (photophobia), and abnormal optic nerve development. People with albinism have an elevated risk of severe sunburn and skin cancer, requiring strict sun protection.

\synonyms
Oculocutaneous Albinism; OCA; Congenital Achromia; Hypopigmentation Disorder

\medterm{Albino} An older noun for a person or animal with albinism, an inherited condition involving reduced melanin. For people, “person with albinism” is usually preferred because it describes the condition without defining the person by it.

\medterm{Albright Hereditary Osteodystrophy} A genetic condition that may cause short stature, a round face, obesity, dental problems, extra bone under the skin, and shortening of some finger bones. It is linked to changes in the \textit{GNAS} gene and may occur with hormone-resistance disorders.

\synonyms
AHO; Pseudohypoparathyroidism Type 1A Phenotype

\medterm{Albright's Hereditary Osteodystrophy} A group of inherited features that may include short stature, a round face, shortened fingers or toes, extra bone under the skin, and obesity. Some people also have resistance to parathyroid hormone or other hormones.

\medterm{Albumen} The white part of an egg, made mostly of water and proteins. Albumen is a food substance and is different from albumin, the main protein in blood plasma.

\medterm{Albumin} The main protein in blood plasma, made by the liver. It helps keep fluid inside blood vessels and carries substances such as hormones, fatty acids, and some medicines through the bloodstream.

\synonyms
Serum Albumin

\medterm{Albumin Clearance} A measure of how quickly the blood protein albumin leaves the bloodstream or another body space. It is used mainly in specialized kidney, liver, or research tests to study leakage, transport, or removal of the protein.

\medterm{Albumin-To-Creatinine Ratio} A urine test that estimates how much albumin, a blood protein, is leaking through the kidneys by comparing it with creatinine in the same sample. A persistent result above 30 mg/g may indicate kidney damage; results above 300 mg/g indicate severe leakage.

\medterm{Albuminocytologic Dissociation} A characteristic cerebrospinal fluid finding marked by an elevated protein concentration in the absence of a commensurate increase in white blood cell count. Classically seen in Guillain-Barr'{e} syndrome after the first week of illness, as well as in chronic inflammatory demyelinating polyneuropathy and spinal cord tumors.

\synonyms
Albumino-Cytological Dissociation; Protein-Cell Dissociation

\medterm{Albuminuria} An abnormally high amount of albumin in the urine, showing that the kidneys are allowing this blood protein to pass through. It may be temporary or persistent; ongoing albuminuria can indicate kidney damage and increased cardiovascular risk.

\medterm{ALCAPA Syndrome} A rare congenital cardiovascular anomaly in which the left coronary artery arises abnormally from the pulmonary artery instead of the aorta. Because the pulmonary artery carries desaturated blood at low pressure, myocardial perfusion is compromised, leading to retrograde coronary flow (coronary steal), infantile heart failure, myocardial infarction, and mitral regurgitation if uncorrected.

\synonyms
ALCAPA; Bland-White-Garland Syndrome; Anomalous Left Coronary Artery from the Pulmonary Artery

\medterm{Alcohol} Ethanol, a substance in beer, wine, and spirits that affects the central nervous system. Excessive consumption impairs motor coordination and cognitive judgment, produces acute intoxication, and can lead to physical dependence, metabolic toxicity, and chronic organ damage.

\medterm{Alcohol Dehydrogenase} An enzyme found mainly in the liver that begins breaking down alcohol. It changes alcohol into acetaldehyde, a harmful intermediate that is then processed into less toxic substances.

\synonyms
ADH; Ethanol Dehydrogenase

\medterm{Alcohol Intolerance} An unpleasant reaction after drinking alcohol, often caused by inherited difficulty breaking down alcohol or sensitivity to an ingredient in the drink. Symptoms may include flushing, headache, nausea, rapid heartbeat, or worsening asthma.

\medterm{Alcohol Intoxication} Temporary impairment caused by drinking alcohol. It can cause poor judgment, confusion, slurred speech, loss of coordination, vomiting, reduced consciousness, slowed breathing, or coma.

\synonyms
Alcoholic Intoxication

\medterm{Alcohol Poisoning} A life-threatening overdose of alcohol that severely depresses brain function and can cause vomiting, inability to wake, seizures, slow or irregular breathing, low body temperature, coma, or death.

\medterm{Alcohol Swab} A single-use pad moistened with alcohol, usually used to clean intact skin before an injection or other procedure. The skin is allowed to dry before the procedure.

\medterm{Alcohol Use Disorder} A medical condition in which repeated alcohol use causes difficulty controlling drinking, strong cravings, or problems at home, work, school, or in relationships. Symptoms may be mild, moderate, or severe.

\synonyms
Alcohol Abuse
Alcohol Abuse and Alcoholism
Alcohol Addiction
Alcohol Dependence
Alcohol Dependence (Alcohol Abuse and Alcoholism)
Alcohol Dependency
Alcoholism

\medterm{Alcohol Withdrawal} Symptoms that can occur after a person who drinks heavily for a long time suddenly stops or greatly reduces alcohol. Symptoms include shaking, sweating, anxiety, nausea, hallucinations, seizures, and delirium; severe withdrawal can be life-threatening.

\synonyms
Alcohol Withdrawal Syndrome

\medterm{Alcohol-Associated Liver Disease} Liver injury associated with alcohol use. It ranges from steatosis and inflammation to fibrosis, cirrhosis, liver failure, or liver cancer; severity depends on alcohol exposure, genetics, nutrition, and other health factors.

\medterm{Alcohol-Related Psychosis} Hallucinations, delusions, or other psychotic symptoms associated with alcohol use, intoxication, withdrawal, or a coexisting alcohol-related disorder. Other medical and psychiatric causes must be excluded.

\medterm{Alcoholic Cardiomyopathy} A disease of the heart muscle caused by long-term heavy alcohol use. The heart becomes enlarged and weak, reducing its ability to pump blood and potentially causing heart failure or abnormal heart rhythms.

\medterm{Alcoholic Foetor} A characteristic odor on the breath and body resulting from recent ethanol consumption and the pulmonary excretion of volatile alcohol and its primary metabolite, acetaldehyde. It serves as an immediate clinical sign of acute alcohol exposure or intoxication during physical assessment.
\synonyms{Alcoholic Fetor; Alcohol Breath; Fetor Alcoholicus}

\medterm{Alcoholic Hepatitis} An acute inflammation and injury of the liver caused by heavy alcohol use. It commonly causes jaundice, fatigue, fever, nausea, and abdominal pain and can become severe enough to cause liver failure.

\synonyms
Alcohol-Associated Hepatitis; Alcohol-Induced Hepatitis

\medterm{Alcoholic Ketoacidosis} A dangerous buildup of acids in the blood that can occur after prolonged heavy alcohol use, poor food intake, and repeated vomiting. It may cause abdominal pain, nausea, dehydration, rapid breathing, and a normal or low blood-sugar level.

\medterm{Alcoholic Liver Disease} Alternate name; see \textit{Alcohol-Associated Liver Disease}.

\medterm{Alcoholic Myopathy} Muscle damage caused by heavy alcohol use. It may cause muscle pain, tenderness, weakness, or wasting; severe acute injury can destroy muscle tissue and damage the kidneys.

\synonyms
Alcohol-Induced Myopathy; Ethanol-Induced Myopathy

\medterm{Alcoholic Neuropathy} Nerve damage associated with long-term heavy alcohol use and related nutritional deficiency. It commonly causes burning pain, numbness, tingling, or weakness, especially in the feet and lower legs.

\medterm{Alcuronium} A medicine that temporarily relaxes or paralyzes skeletal muscles during anesthesia. Because it can stop the muscles needed for breathing, it is given with controlled airway and breathing support.

\medterm{Aldolase Blood Test} A blood test that measures aldolase, an enzyme found mainly in muscle and liver. A high result may occur with muscle injury or disease, but it is not specific and is interpreted with symptoms and other tests.

\medterm{Aldosterone} A hormone made by the adrenal glands that helps the kidneys retain sodium and water and remove potassium. It helps control blood volume and blood pressure.

\medterm{Aldosterone Blood Test} A blood test that measures the hormone aldosterone. It helps investigate high or low blood pressure and abnormal potassium levels and is usually interpreted together with a renin test.

\medterm{Aldosterone Escape} A physiological phenomenon in primary hyperaldosteronism where persistent mineralocorticoid excess causes initial sodium and water retention, but subsequent increased atrial natriuretic peptide secretion and pressure natriuresis prevent progressive hypervolemia and severe edema.

\medterm{Aldosterone-to-Renin Ratio} A clinical screening blood test used in the evaluation of secondary hypertension to detect primary aldosteronism (Conn syndrome). A significantly elevated ratio of plasma aldosterone concentration to plasma renin activity prompts confirmatory testing.

\synonyms
ARR

\medterm{Aleukemic Leukemia} A presentation of acute or chronic leukemia in which malignant blasts or abnormal leukocytes are absent or present in extremely low numbers in the peripheral blood smear, despite extensive malignant infiltration of the bone marrow.

\medterm{Alexander Disease} A rare genetic disorder in which abnormal GFAP protein damages the brain’s supporting cells and white matter. It may begin in infancy, childhood, or adulthood and can cause enlarged head size, developmental problems, movement difficulty, seizures, or progressive weakness.

\synonyms
Fibrinoid Leukodystrophy; GFAP Leukodystrophy

\medterm{Alexia} Loss of the ability to understand written words despite adequate eyesight, usually because of brain damage from a stroke, injury, tumour, or degenerative disease. It may occur with other language or reading problems.

\medterm{Alexia Without Agraphia} A classical neurological disconnection syndrome in which a patient can write spontaneously and to dictation but is completely unable to read printed or handwritten text (including their own writing). Caused by an infarction of the left posterior cerebral artery affecting the left visual cortex and the splenium of the corpus callosum, isolating the intact right visual cortex from the language centers of the left hemisphere.

\synonyms
Pure Alexia; Dejerine Syndrome; Pure Word Blindness

\medterm{Alexithymia} Difficulty recognizing, understanding, or describing one’s own emotions. It is a trait or symptom rather than a diagnosis and may occur with autism, depression, trauma, stress, or other conditions.

\medterm{Alezzandrini Syndrome} An exceptionally rare acquired disorder of unilateral melanocyte destruction involving ipsilateral facial vitiligo and poliosis with retinal or visual abnormalities and sometimes hearing loss. The cause is uncertain; ophthalmic and audiologic follow-up is important because vision or hearing may progressively worsen.

\medterm{Alfacalcidol} A medicine related to vitamin D that the liver changes into an active form. It increases absorption of calcium and phosphate and is used for selected disorders of bones, kidneys, or mineral balance.

\medterm{Algia} A word ending that means pain. It appears in terms such as neuralgia for nerve pain, myalgia for muscle pain, and arthralgia for joint pain; it does not identify the cause.

\medterm{Alglucerase} An older enzyme-replacement medicine for Gaucher disease, a disorder in which fatty substances build up in cells. It replaces an enzyme the body lacks; newer enzyme preparations are used more often.

\medterm{Algor Mortis} The gradual cooling of the body after death toward the surrounding temperature. The amount of cooling can help estimate the time since death, but it is affected by the environment and other conditions.

\medterm{Aliasing Artifact} An ultrasound error in which blood flow moving faster than the selected display range appears to change direction or speed. It reflects a limitation of the image display, not necessarily abnormal blood flow.

\medterm{Alien Hand Syndrome} A neurological condition in which an upper limb performs autonomous, purposeful-appearing actions without conscious voluntary control, often accompanied by the patient's sensation that the limb is foreign or acting independently. It is associated with lesions of the corpus callosum, supplementary motor area, or medial frontal cortex.

\synonyms
Alien Limb Phenomenon; Anarchic Hand

\medterm{Alimentary} Relating to eating, digestion, or the movement of food through the body. In medical language, it usually refers to the digestive tract and the organs that help break down and absorb food.

\medterm{Alimentary Canal} The continuous muscular tube through which food passes from the mouth to the anus. It includes the throat, oesophagus, stomach, small intestine, and large intestine.

\medterm{Alkalemia} An abnormally elevated arterial blood pH (greater than 7.45), indicating a decreased hydrogen ion concentration in circulating blood. It represents the net systemic state caused by underlying metabolic or respiratory alkalosis.

\medterm{Alkali} A basic substance that can neutralize acids. Strong alkalis, such as some cleaning chemicals, can burn the skin, eyes, and digestive tract, while weaker alkalis have food, laboratory, or medical uses.

\medterm{Alkaline Diet} A diet based mainly on foods claimed to make the body more alkaline. The lungs and kidneys tightly control blood acidity, so food does not make the blood alkaline, although many such diets include nutritious foods.

\medterm{Alkaline Phosphatase} An enzyme found mainly in the liver, bile ducts, bones, intestine, and placenta. A high blood level may reflect liver or bile-duct disease, increased bone activity, pregnancy, or another condition and must be interpreted with other tests.

\medterm{Alkaline Phosphatase Test} A quantitative serum laboratory assay measuring alkaline phosphatase activity to evaluate suspected cholestatic liver disease, biliary obstruction, or increased bone turnover such as Paget disease or bone metastases.

\synonyms
ALP Blood Test; Serum ALP; Alkaline Phosphatase Level

\medterm{Alkaloid} A naturally occurring chemical containing nitrogen, often made by plants, that can strongly affect the body. Examples include morphine, quinine, nicotine, and atropine.

\synonyms
Plant Alkaloid; Basic Nitrogenous Organic Compound

\medterm{Alkalosis} A condition in which the blood becomes too alkaline. It may result from repeated vomiting, certain medicines, hormone disorders, or breathing too rapidly, and can cause muscle cramps, tingling, confusion, or abnormal heart rhythms.

\medterm{Alkaptonuria} A rare inherited metabolic disorder caused by reduced activity of homogentisate 1,2-dioxygenase, usually because of disease-causing variants in the \textit{HGD} gene. Homogentisic acid accumulates, darkens urine, and can deposit in cartilage, tendons, skin, and other tissues in a process called ochronosis. Over time it may cause arthritis and stiffness of the spine and large joints, urinary stones, or changes involving heart valves.

\medterm{Alkylating Agents} Drugs or chemicals that damage DNA by attaching chemical groups to it. Some are used to treat cancer by stopping cells from multiplying, but they can also injure healthy cells and cause side effects.

\medterm{Allantoin} A compound used in some skin products to soften, moisturize, and protect irritated or rough skin. Its effect and safety depend on the product and how it is used.

\medterm{Allantois} A temporary structure in the early embryo that helps with blood-vessel development and contributes to the umbilical cord and part of the urinary system.

\medterm{Allele} One of two or more versions of a gene. A person usually inherits one allele from each parent; different alleles can influence traits or the risk of inherited disease.

\medterm{Allelic} Relating to an allele, one of the different versions of a gene inherited from a parent.

\medterm{Allen Test} A bedside test that checks whether blood can reach the hand through both of its main arteries. It may be done before taking blood from or placing a catheter in the radial artery.

\medterm{Allergen} A usually harmless substance that triggers an allergic immune response in a sensitive person. Common allergens include foods, pollen, medicines, animal proteins, insect venom, and natural rubber.

\medterm{Allergen Immunotherapy} Treatment that gives gradually increasing amounts of a known allergen by injections or under the tongue. Over time, it can reduce allergic symptoms for selected allergies, but it can rarely cause a serious reaction.

\medterm{Allergic and Environmental Asthma} Asthma triggered or worsened by allergens or environmental exposures such as dust mites, molds, pollens, animals, smoke, air pollution, or workplace agents. Management combines exposure reduction, controller therapy, trigger planning, and evaluation for occupational or allergic disease when indicated.

\medterm{Allergic Bronchopulmonary Aspergillosis} An allergic reaction to the fungus Aspergillus in the airways, occurring mainly in people with asthma or cystic fibrosis. It can cause worsening breathing symptoms, mucus plugs, airway widening, and high levels of allergy-related cells or antibodies.

\medterm{Allergic Cascade} The chain of immune events that follows exposure to an allergen. Immune cells release substances such as histamine, which can cause itching, swelling, mucus, wheezing, hives, or anaphylaxis.

\medterm{Allergic Conjunctivitis} Alternate name; see \textit{Eye Allergies}.

\medterm{Allergic Contact Dermatitis} A delayed type-IV hypersensitivity reaction after skin contact with an allergen to which the immune system has been sensitized. It causes itching, redness, swelling, blisters, scaling, or chronic fissuring, often extending beyond the contact area; patch testing can help identify the trigger.

\medterm{Allergic Dermatitis} Skin inflammation caused by an immune reaction to an allergen. It commonly causes itching, redness, swelling, blisters, or scaling and may occur after contact with a triggering substance.

\medterm{Allergic Diseases} Conditions in which the immune system reacts abnormally to a usually harmless substance called an allergen, such as pollen, dust mites, mold, food proteins, medicines, insect venom, or animal proteins. Some reactions involve immunoglobulin E (IgE) antibodies, while others do not. Allergic diseases may affect the nose, eyes, skin, lungs, or digestive system and can cause sneezing, itching, hives, swelling, vomiting, wheezing, or anaphylaxis. Examples include allergic rhinitis, food allergy, drug allergy, allergic conjunctivitis, atopic dermatitis, urticaria, and allergic asthma. Food intolerance is different because it does not involve an immune reaction.

\medterm{Allergic Proctocolitis} A benign, non-IgE-mediated food protein-induced inflammatory disorder of the distal colon and rectum, seen predominantly in young infants exposed to cow's milk or soy protein in formula or breast milk, classically causing blood- and mucus-streaked stools in an otherwise thriving baby.

\synonyms
Food Protein-Induced Allergic Proctocolitis; FPIP

\medterm{Allergic Reaction} An immune response to a usually harmless substance. It may cause sneezing, itching, hives, swelling, vomiting, or breathing difficulty and can sometimes become life-threatening anaphylaxis.

\synonyms
Allergy

\synonyms
Allergic Reactions

\medterm{Allergic Rhinitis} Inflammation inside the nose caused by an allergic reaction to inhaled substances such as pollen, dust mites, mold, or animal proteins. It causes sneezing, itching, a blocked or runny nose, and sometimes itchy, watery eyes.

\synonyms
Hay Fever; Pollinosis; Allergic Rhinoconjunctivitis

\medterm{Allergic Urticaria} An acute, IgE-mediated type I hypersensitivity skin reaction triggered by specific external allergens such as foods, drugs, insect venoms, or latex. Cross-linking of allergen-specific IgE on the surface of dermal mast cells triggers rapid degranulation and release of histamine, prostaglandins, and leukotrienes, causing localized vasodilation and increased vascular permeability that manifests as intensely pruritic, transient, erythematous wheals that resolve within 24 hours.
\synonyms{Acute Allergic Urticaria; IgE-Mediated Urticaria; Allergic Hives}

\medterm{Allergy} An immune response to a usually harmless substance such as pollen, food, medicine, insect venom, or animal proteins. Symptoms range from sneezing, itching, and hives to swelling, breathing difficulty, or life-threatening anaphylaxis.

\synonyms
Allergic Hypersensitivity; Type I Hypersensitivity

\medterm{Allergy Shots} Injections used for allergen immunotherapy. They contain gradually increasing amounts of a specific allergen to reduce sensitivity over time and are given in a setting where a serious allergic reaction can be recognized and treated.

\medterm{Allergy Skin Test} An in vivo diagnostic procedure in which small amounts of suspected allergens are introduced into or onto the skin via prick, puncture, scratch, or intradermal injection to assess immediate IgE-mediated hypersensitivity through localized wheal-and-flare reactions.

\synonyms
Skin Prick Test; Epicutaneous Test; Allergy Skin Testing; Scratch Test

\medterm{Allergy Test} A skin or blood test used to look for sensitivity to a suspected allergen. A positive result shows sensitization but does not always prove that exposure causes symptoms, so the result must be considered with the person’s history.

\synonyms
Skin Prick Test; Specific IgE Test; Allergen Challenge

\medterm{Allis Clamp} A surgical instrument with small teeth used to hold firm tissue during an operation. It provides a strong grip but can injure delicate tissue if applied with too much pressure.

\medterm{Allocheiria} A neurological sign in which a person feels or identifies a touch on the opposite side of the body from where it was applied. It reflects impaired localization of sensation and may occur with disorders affecting the parietal lobe.

\synonyms
Allochiria

\medterm{Allodynia} Pain caused by something that normally should not hurt, such as light touch, clothing, or gentle pressure. It can occur with nerve injury, migraine, shingles, fibromyalgia, or other pain conditions.

\medterm{Allogeneic} Coming from another person of the same species who is genetically different, as in an allogeneic blood, tissue, stem-cell, or organ transplant.

\medterm{Allograft} An organ, tissue, or group of cells transplanted from one genetically different individual to another individual of the same species. The recipient’s immune system may reject it, so immunosuppressive medicines may be needed.

\medterm{Allosteric Regulation} Control of a protein or enzyme when another molecule attaches away from its main working site. The attachment changes the protein’s shape and can increase or decrease how well it works.

\medterm{Aloe} A plant used in some skin products, medicines, and supplements. Aloe gel applied to skin may soothe minor irritation, while products containing aloe latex taken by mouth can cause cramps, diarrhoea, and dangerous fluid or salt loss.

\medterm{Aloe Vera} A plant whose clear inner-leaf gel is used in some skin products and supplements. The gel may soothe minor irritation, while oral latex or whole-leaf products can cause diarrhoea, fluid or salt loss, medicine interactions, and toxicity.

\medterm{Alogia} Markedly reduced speech or speech that conveys little information. It may occur as a negative symptom of schizophrenia and can also be associated with depression, brain disease, or other conditions.

\medterm{Alongshan Virus} A virus spread by some ticks and reported mainly in parts of Asia and Europe. It may cause fever, headache, fatigue, nausea, or neurological symptoms, but its ability to cause human disease is still being studied.

\medterm{Alopecia} Hair loss from the scalp or another part of the body. It may be nonscarring, in which the hair follicles remain and hair may regrow, or scarring, in which inflammation or injury destroys the follicles and can cause permanent hair loss. Causes include inherited pattern hair loss, autoimmune disease, hormonal changes, illness, medicines, nutritional problems, infection, repeated pulling, or other injury. The amount, pattern, and speed of hair loss vary with the cause.

\medterm{Alopecia Areata} An autoimmune, usually nonscarring hair-loss disorder in which immune cells attack anagen hair follicles. It often causes sharply demarcated smooth patches with exclamation-mark hairs and may involve the scalp, beard, eyebrows, or nails; regrowth is possible but relapses and extensive forms can occur.

\synonyms
AA; Patchy Hair Loss; Autoimmune Alopecia

\medterm{Alopecia Mucinosa (Follicular Mucinosis)} A skin disorder in which mucin accumulates within hair follicles and sebaceous glands, causing itchy follicular papules or plaques with hair loss. It may be a primary self-limited condition or a cutaneous sign of lymphoma, especially mycosis fungoides, so persistent or widespread disease needs evaluation.

\medterm{Alopecia Totalis} Complete loss of hair from the scalp, usually caused by alopecia areata. Hair may regrow, remain absent, or fall out again, and other autoimmune conditions may occur.

\medterm{Alopecia Universalis} Complete or near-complete loss of hair from the scalp and body, usually caused by extensive alopecia areata. Hair may regrow or recur unpredictably.

\medterm{ALP Blood Test} Alternate index label; see \textit{Alkaline Phosphatase Test}.

\medterm{ALP Isoenzyme Test} A blood test that helps identify whether a high alkaline phosphatase level comes mainly from the liver, bones, intestine, or placenta. It can help distinguish liver or bile-duct disease from increased bone activity.

\medterm{Alpha Blockers} Medicines that relax certain blood vessels and smooth muscles by blocking alpha receptors. They may lower blood pressure or improve urine flow in an enlarged prostate, but can cause dizziness or fainting.

\medterm{Alpha Cells} Cells in the pancreas that make glucagon. Glucagon raises blood glucose when it falls too low or when the body needs energy between meals.

\synonyms
Pancreatic Alpha Cells; Glucagon-Secreting Cells

\medterm{Alpha Hydroxy Acids} A group of acids used in skin products to loosen the outer layer of skin and improve roughness, fine lines, or uneven colour. Common examples are glycolic acid and lactic acid; they can irritate skin and increase sun sensitivity.

\synonyms
AHAs; Fruit Acids

\medterm{Alpha Waves} A regular pattern of electrical activity in the brain, usually eight to thirteen cycles per second. It is strongest when a relaxed, awake person has closed eyes and decreases with attention or eye opening.

\synonyms
Alpha Rhythm; Posterior Dominant Rhythm

\medterm{Alpha-1 Adrenergic Antagonist} A medicine or other substance that blocks alpha-1 adrenergic receptors. This reduces selected effects of noradrenaline on vascular smooth muscle and other tissues; the precise effects depend on the receptor distribution and the substance used.

\synonyms
Alpha-1 Adrenergic Blocker

\medterm{Alpha-1 Antitrypsin Blood Test} A blood test that measures alpha-1 antitrypsin, a protein made by the liver that helps protect lung tissue. A low level may suggest inherited alpha-1 antitrypsin deficiency, which can affect the lungs and liver.

\medterm{Alpha-1 Antitrypsin Deficiency} An inherited condition in which low or abnormal alpha-1 antitrypsin increases the risk of emphysema, especially at a young age, and liver disease. Some people also develop inflammation of fatty tissue under the skin.

\synonyms
AATD; AAT Deficiency; Alpha-1 Deficiency

\medterm{Alpha-Agonists} Medicines that activate alpha receptors on cells. Depending on the medicine and tissue, they can narrow blood vessels, raise blood pressure, reduce eye fluid, or affect the nervous system.

\medterm{Alpha-Fetoprotein} A protein made mainly by a developing fetus. Measuring it in maternal blood can help screen for some fetal conditions, while measuring it in adults can help assess selected liver or germ-cell tumours. An abnormal result does not establish a diagnosis.

\medterm{Alpha-Gal Syndrome} An allergy to a sugar found in most mammals but not humans or other primates. It can cause delayed hives, stomach symptoms, or anaphylaxis after eating mammalian meat or using some animal-derived products.

\synonyms
Alpha-Gal Allergy

\medterm{Alpha-Globulin} One of several groups of proteins in the liquid part of blood. These proteins carry substances and help with inflammation, clotting, and immune defence; their levels may change with liver disease, inflammation, or other conditions.

\medterm{Alpha-Glucosidase Inhibitors} Diabetes medicines that slow the digestion and absorption of carbohydrates in the small intestine. They mainly reduce the rise in blood glucose after meals and may cause gas, bloating, or diarrhoea.

\medterm{Alpha-Ketoglutarate} A substance made during the breakdown of nutrients for energy. It also helps the body process amino acids and supports several chemical reactions inside cells.

\medterm{Alpha-Synuclein} A protein found mainly in nerve cells that helps regulate nerve-cell functions. Abnormal forms can collect into clumps in Parkinson disease, Lewy-body dementia, and related disorders, although the clumps are not the only cause of these diseases.

\medterm{Alpha-Thalassemia} An inherited blood disorder, usually autosomal recessive, in which the body makes too little alpha-globin, a part of hemoglobin that carries oxygen. The condition may cause no symptoms, mild anemia, or severe anemia before or shortly after birth, depending on the number of affected alpha-globin genes.

\synonyms
Alpha-Globin Deficiency; HbH Disease Carrier

\medterm{Alpha1-Antitrypsin (AAT) Deficiency} An inherited deficiency or dysfunction of alpha-1 antitrypsin that predisposes to early emphysema, especially with smoking, and sometimes liver disease from abnormal protein accumulation. Testing is considered in early or unexplained COPD, basilar emphysema, or unexplained liver disease.

\medterm{Alport Syndrome} An inherited disorder of collagen that can cause progressive kidney disease, hearing loss, and changes in the eyes. It may be caused by changes in the \textit{COL4A3}, \textit{COL4A4}, or \textit{COL4A5} genes and may follow X-linked, autosomal dominant, or autosomal recessive inheritance.

\medterm{Alprazolam} A benzodiazepine medicine used for anxiety and panic disorder. It calms brain activity but can cause sleepiness, impaired coordination, dependence, and dangerous slowing of breathing, especially with opioids or alcohol.

\medterm{Alstrom Syndrome} A rare inherited disorder caused by changes in the ALMS1 gene. It may cause progressive vision and hearing loss, obesity, diabetes or insulin resistance, heart disease, and kidney or liver problems.

\synonyms
Alström Syndrome

\medterm{Alternating Hemiplegia Of Childhood} A rare neurological disorder beginning in infancy or early childhood, causing recurrent episodes of temporary paralysis that may affect either side of the body. Episodes may also involve abnormal movements, eye-movement problems, breathing difficulty, seizures, and developmental impairment; symptoms typically disappear during sleep.

\medterm{Altitude Illness} Illness caused by rapid ascent to high altitude without enough time to adjust. It includes acute mountain sickness, high-altitude brain swelling, and high-altitude lung swelling; severe forms can be fatal.

\synonyms
Altitude Sickness

\medterm{Aluminum Phosphide} A highly poisonous pesticide that releases phosphine gas when it contacts moisture or stomach acid. Poisoning can rapidly damage the heart, lungs, digestive system, and other organs and requires emergency treatment.

\medterm{Aluminum Toxicity} Toxicity from excessive aluminum exposure or accumulation, especially in advanced kidney disease or with contaminated dialysis products. Chronic exposure can cause encephalopathy, speech or movement changes, microcytic anemia, and bone pain or osteomalacia; evaluation includes exposure review, kidney function, and specialist-directed aluminum testing.

\medterm{Alvarado Score} A clinical score used to estimate the likelihood of acute appendicitis from symptoms, examination findings, and blood-test results. It supports assessment but does not by itself confirm or exclude appendicitis.

\medterm{Alveolar} Relating to an alveolus, especially one of the tiny air sacs in the lungs where oxygen enters the blood and carbon dioxide leaves it.

\medterm{Alveolar Abnormalities} Changes in the lung air sacs that can interfere with oxygen and carbon-dioxide exchange. They may occur with infection, fluid buildup, emphysema, scarring, or a developmental condition.

\medterm{Alveolar Bone} The part of the jawbone that surrounds and supports the roots of the teeth. It can shrink after tooth loss and may need rebuilding before some dental implants.

\medterm{Alveolar Bone Loss} Loss of the bone that supports the teeth, most often caused by periodontitis. It can make teeth loose and may lead to tooth loss; dental treatment can slow further loss.

\medterm{Alveolar Dead Space} The volume of inhaled air that reaches alveoli that are adequately ventilated but lack sufficient pulmonary capillary blood perfusion, preventing gas exchange from taking place.

\medterm{Alveolar Echinococcosis} A rare, serious infection caused by the young form of the Echinococcus multilocularis tapeworm, usually after swallowing contaminated food or water. It creates slowly growing, tumor-like areas in the liver that can invade nearby tissue or spread.

\synonyms
Alveolar Hydatid Disease

\medterm{Alveolar Gas Equation} A formula used to estimate the oxygen level in the lung air sacs from the inhaled oxygen level and carbon-dioxide level. Comparing this estimate with arterial blood oxygen helps assess how well the lungs exchange gases.

\medterm{Alveolar Hypoventilation} Deficient ventilation of the pulmonary alveoli resulting in reduced carbon dioxide elimination, leading to arterial hypercapnia and respiratory acidosis.

\medterm{Alveolar Macrophage} An immune cell that lives in the lung air sacs and removes inhaled dust, microbes, and damaged cells. It also helps regulate inflammation in the lungs.
\synonyms
Dust Cell

\medterm{Alveolar Osteitis} Painful inflammation of a tooth socket after extraction, usually caused by loss or breakdown of the blood clot that normally protects the socket. It is also called dry socket.

\synonyms
Dry Socket; Fibrinolytic Alveolitis

\medterm{Alveolar Proteinosis} Alternate form; see \textit{Pulmonary Alveolar Proteinosis}.

\medterm{Alveolar Rhabdomyosarcoma} A cancer of immature muscle-forming cells, usually affecting children or adolescents. It may arise in the limbs, trunk, head, or neck and can spread early to lymph nodes, lungs, bones, or other tissues.

\medterm{Alveolar Ventilation} The amount of fresh air reaching the lung air sacs each minute and available for oxygen and carbon-dioxide exchange. It depends on breathing rate, breath size, and the amount of air remaining in the breathing passages.

\medterm{Alveolar Ventilation Equation} A relationship showing how breathing affects blood carbon-dioxide levels. When the amount of fresh air reaching the lung air sacs falls, carbon dioxide rises; when effective breathing increases, carbon dioxide falls, if production stays constant.

\medterm{Alveolar-Arterial Oxygen Gradient} The difference between oxygen pressure in the lung air sacs and oxygen pressure in arterial blood. It is calculated as PAO2 minus PaO2. A high value may indicate a problem with oxygen transfer, blood flow through the lungs, or lung tissue. A normal value with low blood oxygen may suggest slow breathing or low oxygen in the surrounding air.

\synonyms
A-A Gradient; A-a Gradient; A-a Oxygen Difference

\medterm{Alveolar-Capillary Barrier} The delicate biological membrane formed by the alveolar type I pneumocyte, capillary endothelial cell, and their intervening fused basement membranes, through which oxygen and carbon dioxide diffuse between alveolar air and capillary blood.

\synonyms
Alveolar-Capillary Membrane; Blood-Air Barrier

\medterm{Alveolitis} Inflammation of the lung air sacs, usually causing cough and shortness of breath, or inflammation of a tooth socket after extraction, called dry socket. The body site identifies the meaning.

\medterm{Alveolus} One of the microscopic air sacs at the ends of the lung airways. Oxygen passes from an alveolus into the blood and carbon dioxide passes from the blood into it.

\synonyms
Alveoli

\medterm{Alzheimer Disease} A progressive disease of the brain and the most common cause of dementia. It slowly damages nerve cells and the connections between them. Early symptoms often include problems with recent memory, finding words, or thinking through tasks. As the disease advances, it can affect language, judgment, orientation, behavior, and the ability to carry out everyday activities. Abnormal proteins called amyloid plaques and tau tangles build up in the brain and are linked to this damage. Alzheimer disease is not a normal part of aging, although the risk increases with age, and its symptoms and rate of progression vary.

\synonyms
Alzheimer's Disease; Alzheimer's
Alzheimer'S
Alzheimer'S Disease

\medterm{Amalgam} A dental filling material made by combining mercury with metals such as silver, tin, and copper. It is strong and durable but has a metallic appearance and may not be suitable for every person or tooth.

\medterm{Amastia} A birth condition in which one or both breasts do not develop. The nipple, chest muscle, or other nearby structures may also be absent or underdeveloped.

\medterm{Amaurosis Fugax} Sudden, temporary loss of vision in one eye, often described as a curtain or shade, caused by briefly reduced blood flow to the retina or optic nerve. It can warn of vascular disease or stroke risk.

\medterm{Ambidextrous} Able to use the left and right hands with similar skill. It describes a pattern of handedness, not a disease.

\medterm{Ambient} Surrounding or present in the immediate environment, such as ambient temperature, light, noise, or air.

\medterm{Ambiguous Genitalia} External genital anatomy that does not have typical male or female features. It usually reflects a difference in sex development, which may involve chromosomes, hormones, or the formation of reproductive organs.

\medterm{Amblyopia} Reduced vision in one eye (rarely both) resulting from abnormal brain visual processing during early childhood, occurring despite a structurally normal eye. When the brain receives conflicting, blurry, or misaligned visual signals, it suppresses input from the weaker eye and favors the dominant eye. Common causes include crossed or misaligned eyes (strabismic amblyopia), unequal refractive focusing power between the eyes (anisometropic amblyopia), or physical blockage of the visual axis (deprivation amblyopia from congenital cataract or severe ptosis). Because neural visual pathways develop rapidly in early life, treatment is time-sensitive and most effective during the critical period before age 7 to 8. Management involves correcting refractive errors with glasses, clearing anatomical obstructions, and patching the stronger eye (or using blurring atropine drops) to force the brain to strengthen neural pathways in the amblyopic eye.

\synonyms
Lazy Eye; Functional Amblyopia; Visual Suppression Amblyopia

\medterm{Ambu Bag} A self-inflating manual resuscitation bag used with a mask or airway device to push air into the lungs when a person is not breathing adequately.

\medterm{Ambulatory Blood Pressure Monitoring} Measurement of blood pressure at regular intervals, usually for twenty-four hours, with a portable cuff worn during normal activities and sleep. It can identify blood-pressure patterns that may not appear during a clinic visit.

\medterm{Amebiasis} An infection caused by the parasite Entamoeba histolytica, usually spread by swallowing contaminated food, water, or material from contaminated hands. It may cause no symptoms, diarrhea, abdominal pain, or bloody diarrhea. In some cases, the infection can spread from the intestine to the liver.

\synonyms
Entamoeba Histolytica Infection; Amoebic Dysentery

\medterm{Amebic Colitis} Inflammation of the large intestine caused by Entamoeba histolytica. It may cause abdominal pain, diarrhoea, bloody stools, or fever.

\synonyms
Amoebic Colitis

\medterm{Amebic Dysentery} An acute, invasive intestinal infection caused by the protozoan parasite Entamoeba histolytica, characterized by extensive ulceration of the colonic mucosa ("flask-shaped" ulcers). Patients typically present with severe lower abdominal cramping, tenesmus, and frequent passage of bloody, mucoid stools, with potential complications including fulminant necrotizing colitis, toxic megacolon, bowel perforation, and secondary amebic liver abscess.

\synonyms
Amoebic Dysentery; Intestinal Amebiasis; Amebic Colitis

\medterm{Amebic Liver Abscess} A collection of infected or damaged tissue in the liver caused by the parasite \textit{Entamoeba histolytica}. It often affects the right side of the liver and may cause fever, pain under the right ribs, and an enlarged or tender liver. Ultrasound or another scan and blood tests can support the diagnosis, but results must be interpreted with the person’s symptoms and history. Treatment usually includes a medicine that kills the parasite followed by a medicine that clears it from the intestine; drainage is needed in selected cases.

\synonyms
Amoebic Liver Abscess; Hepatic Amebiasis; Entamoeba Liver Abscess

\medterm{Ameboma} A chronic, localized inflammatory granulomatous pseudotumor of the bowel wall that develops in a minority of patients with untreated or inadequately treated Entamoeba histolytica infection, most commonly affecting the cecum or ascending colon. It presents as a palpable, tender right-lower-quadrant abdominal mass associated with intermittent bloody diarrhea, obstruction, or weight loss, frequently mimicking colonic adenocarcinoma on colonoscopy and cross-sectional imaging.

\synonyms
Amoeboma; Amebic Granuloma

\medterm{Amelanotic Melanoma} A melanoma that produces little or no dark pigment and may appear pink, red, or skin-coloured instead of black or brown. Its appearance can delay recognition; diagnosis requires examination of a tissue sample.

\medterm{Amelia} A rare birth condition in which one or more arms or legs are completely absent. It may occur alone or with other developmental differences.

\medterm{Amelioration} Improvement or reduction in the severity, frequency, or effect of a symptom, disease, impairment, or other harmful condition.

\medterm{Ameloblastoma} A usually slow-growing tumour that begins in cells involved in tooth development, most often in the lower jaw. It can grow into nearby bone, deform the jaw, and return after treatment.

\medterm{Amelogenesis} The process by which developing teeth form and harden enamel, the tough outer covering of each tooth. Problems with this process can make enamel thin, weak, discoloured, or easily damaged.

\medterm{Amelogenesis Imperfecta} A group of inherited conditions in which tooth enamel does not form normally. The enamel may be thin, soft, rough, discoloured, or easily worn on most or all teeth.

\medterm{Amenorrhea} Absence of menstrual periods. Primary amenorrhea means menstruation has not started by the expected age; secondary amenorrhea means periods stop after they have begun. It may be normal during pregnancy, breastfeeding, or menopause, or may result from hormonal, reproductive, medical, nutritional, or exercise-related conditions.

\medterm{American Trypanosomiasis} Alternate name; see \textit{Chagas disease}.

\medterm{Ametropia} A focusing problem in which the eye does not bring incoming light to a sharp focus on the retina. Nearsightedness, farsightedness, and astigmatism are common forms and can cause blurred vision.

\medterm{Amino Acid} A small molecule used to build proteins. Amino acids also support energy production and other chemical processes in the body.

\synonyms
Protein Building Block; Amino Monocarboxylic Acid

\medterm{Amino Acid Catabolism} The body’s breakdown of amino acids after proteins are digested. Their nitrogen part is removed and safely processed, while the remaining part is used to make energy, glucose, ketones, or other substances.

\medterm{Amino Acid Symbols} The standard one-letter and three-letter abbreviations used to identify the 20 amino acids most commonly used to build proteins. They make genetic, laboratory, and biochemical information shorter and easier to compare.

\medterm{Aminoaciduria} Abnormally large amounts of amino acids in the urine, usually because the kidney tubules cannot reabsorb them properly or because of an inherited metabolic disorder.

\medterm{Aminoglycoside Nephrotoxicity} Kidney injury caused by aminoglycoside antibiotics, which can build up in kidney cells. Risk increases with high doses, prolonged treatment, existing kidney disease, or other kidney-toxic medicines.

\medterm{Aminoglycoside Ototoxicity} Inner-ear damage caused by aminoglycoside antibiotics, which may lead to hearing loss, ringing in the ears, dizziness, or balance problems. Risk increases with high or prolonged exposure.

\medterm{Aminoglycosides} A group of antibiotics that kill certain bacteria by stopping them from making essential proteins. Examples include gentamicin, tobramycin, amikacin, and streptomycin; they can damage the kidneys and inner ear.

\medterm{Aminophylline Overdose} Poisoning caused by too much aminophylline, a medicine that opens narrowed airways. It can cause vomiting, tremor, rapid or irregular heartbeat, low blood pressure, seizures, or dangerous changes in blood chemistry.

\medterm{Aminotransferase} An enzyme that helps cells process amino acids, the building blocks of proteins. Blood levels of ALT and AST may rise when liver, muscle, or other cells are injured, but the result does not identify the cause alone.

\medterm{Amiodarone Pulmonary Toxicity} A serious, potentially life-threatening adverse reaction to chronic amiodarone therapy, manifesting as subacute chronic interstitial pneumonitis, organizing pneumonia, or acute respiratory distress syndrome. Pathologically marked by intra-alveolar foamy lipid-laden macrophages (phospholipidosis) and diffuse alveolar damage.

\synonyms
Amiodarone Lung; Amiodarone-Induced Pneumonitis

\medterm{Amiodarone-Induced Skin Discoloration} A blue-grey or slate-grey darkening of sun-exposed skin caused by long-term amiodarone treatment. It is more likely with higher doses and may fade partly after the medicine is stopped.

\medterm{Amiodarone-Induced Thyroid Dysfunction} An overactive or underactive thyroid caused by amiodarone. The medicine contains iodine and can either increase thyroid-hormone production or inflame and damage the thyroid.

\medterm{Amiodarone-Induced Vortex Keratopathy} Whorl-shaped deposits in the surface layer of the cornea caused by amiodarone. They usually cause no symptoms but may produce halos around lights and generally improve after the medicine is stopped.

\medterm{Amitriptyline} A tricyclic antidepressant that changes serotonin and noradrenaline signalling in the brain. It is also used for some nerve pain and migraine prevention; common effects include sleepiness, dry mouth, and constipation.

\medterm{Amitriptyline And Perphenazine Overdose} Poisoning caused by taking too much amitriptyline and perphenazine. It can cause severe sleepiness, confusion, seizures, low blood pressure, abnormal heart rhythms, coma, or breathing problems.

\medterm{Amitriptyline Hydrochloride Overdose} Poisoning caused by taking too much amitriptyline. It can cause severe sleepiness, confusion, seizures, low blood pressure, abnormal heart rhythms, coma, or breathing problems.

\medterm{Ammonia} A nitrogen-containing waste product made when the body processes proteins. The liver normally changes it into urea for removal; high blood levels can affect brain function, especially in severe liver disease.

\medterm{Ammonia Blood Test} A blood test that measures ammonia, a waste product normally cleared by the liver. A high result may occur with severe liver disease, inherited metabolic conditions, some medicines, or other illnesses and must be interpreted with symptoms.

\medterm{Ammonia Poisoning} Injury caused by breathing, swallowing, or getting concentrated ammonia on the skin or eyes. It can burn tissues and cause coughing, breathing difficulty, chest pain, or serious lung damage.

\medterm{Ammonia Toxicity} Irritant or corrosive injury after inhalation, ingestion, or skin or eye contact with ammonia. Concentrated exposure can cause burning, cough, bronchospasm, laryngeal injury, pulmonary edema, or chemical burns; treatment emphasizes immediate decontamination, airway and respiratory support, and specialist toxicology guidance.

\medterm{Ammonium Hydroxide Poisoning} Injury caused by exposure to ammonium hydroxide, a corrosive chemical found in some cleaners and industrial products. It can burn the skin, eyes, mouth, throat, lungs, and digestive tract.

\medterm{Amnesia} Loss or impairment of memory beyond ordinary forgetfulness. It may prevent a person from forming new memories, recalling past events, or doing both, and may result from brain injury, illness, medicines, or emotional trauma.

\medterm{Amnestic Syndrome} A pattern of memory impairment that significantly affects learning, storing, or recalling information. It may occur with brain injury, stroke, dementia, epilepsy, substance use, or other neurological conditions.

\medterm{Amniocentesis} A procedure in which a thin needle removes a small sample of fluid from around the developing baby. The sample can be tested for genetic, chromosomal, or infection-related conditions.

\medterm{Amnioinfusion} A procedure that adds sterile fluid to the sac around the developing baby, usually during labour, to reduce pressure on the umbilical cord or dilute thick meconium.

\medterm{Amnion} The thin inner membrane forming the fluid-filled sac around the developing baby. It holds the amniotic fluid and helps protect the fetus during pregnancy.

\medterm{Amnion Nodosum} Small deposits of fetal skin material or other debris attached to the inner membrane of the pregnancy sac. They are usually linked with too little amniotic fluid, often because the developing baby is making less urine.

\medterm{Amnioscope} An instrument formerly inserted through the cervix to examine the amniotic fluid and membranes. It is rarely used now because ultrasound is safer and more practical.

\medterm{Amniotic Band} A strand of amniotic membrane that can wrap around a developing body part, restricting growth or blood flow and causing a deformity or loss of tissue.

\synonyms
Amniotic Band Sequence

\medterm{Amniotic Band Sequence} A group of birth differences caused when strands of the amniotic membrane wrap around or attach to developing fetal parts. Effects may include constriction rings, missing digits, limb defects, or facial and body-wall abnormalities.

\medterm{Amniotic Band Syndrome} Another name for a group of birth differences caused by strands of amniotic membrane constricting or entangling developing fetal parts. The effects vary widely and may involve the limbs, face, or body wall.

\synonyms
Amniotic Constriction Band Sequence; ADAM Complex; Streeter Dysplasia

\medterm{Amniotic Fluid} The liquid surrounding the developing baby inside the amniotic sac. It cushions the baby, allows movement and growth, and helps maintain a stable environment during pregnancy. It is continually produced, swallowed, and replaced.

\medterm{Amniotic Fluid Embolism} A rare, life-threatening complication of pregnancy or childbirth in which amniotic material enters the mother’s bloodstream and triggers sudden breathing or heart problems, often with severe bleeding.

\medterm{Amniotic Fluid Index} An ultrasound estimate of the amount of fluid around the developing baby. It adds the deepest fluid measurements from four areas of the uterus and helps identify unusually low or high fluid levels.

\synonyms
AFI; Amniotic Fluid Volume Assessment

\medterm{Amniotic Sac} The thin, fluid-filled membrane around the developing baby in the uterus. It contains amniotic fluid, cushions the fetus, and usually breaks near or during labour when the waters break.

\medterm{Amok} An old, culture-related term for a sudden episode of uncontrolled, often violent behaviour followed by exhaustion or memory loss. Modern assessment considers medical, neurological, psychiatric, substance-related, and cultural factors.

\medterm{AMP-Activated Protein Kinase} A master cellular energy-sensing serine/threonine kinase activated during metabolic stress when intracellular AMP-to-ATP or ADP-to-ATP ratios rise (such as during exercise, hypoxia, or fasting). Once activated, it restores energy homeostasis by promoting ATP-generating catabolic processes (such as glucose uptake via GLUT4, fatty acid oxidation, glycolysis, and autophagy) while suppressing ATP-consuming anabolic pathways (such as lipid synthesis, protein synthesis via mTOR inhibition, and gluconeogenesis). It is a primary molecular target of the antidiabetic drug metformin.

\synonyms
AMPK; 5'-AMP-Activated Protein Kinase

\medterm{Amphetamine} A stimulant medicine that increases certain chemical signals in the brain. It is used for selected attention, activity, or sleep disorders; misuse can cause dependence, high blood pressure, abnormal heart rhythms, or psychosis.

\medterm{Amphetamine Substance Use} Use of amphetamine-type stimulants, which can be associated with insomnia, anxiety, psychosis, high blood pressure, arrhythmias, overheating, dependence, and overdose.

\medterm{Amphetamine Toxicity} A sympathomimetic toxidrome caused by amphetamine or related stimulants, with agitation, mydriasis, sweating, tachycardia, hypertension, hyperthermia, chest pain, seizures, or psychosis. Severe cases may cause stroke, rhabdomyolysis, or cardiac ischemia; benzodiazepines, cooling, and treatment of complications are central.

\medterm{Amphetamine-Related Psychiatric Disorders} Psychiatric syndromes associated with amphetamine exposure, including intoxication, withdrawal, substance-induced psychosis, mood symptoms, anxiety, or a stimulant use disorder.

\medterm{Amplified Musculoskeletal Pain Syndrome} A condition, often affecting children or teenagers, in which the nervous system increases normal pain signals. Pain may be severe or widespread and can occur with sensitivity, fatigue, swelling, or difficulty using the affected area.

\medterm{Ampulla} A widened part of a tube-shaped body structure. Examples include the ampulla where the bile and pancreatic ducts join the intestine and the ampulla of a fallopian tube.

\medterm{Ampulla of Vater} An anatomical dilation located within the medial wall of the descending duodenum, formed by the union of the common bile duct and the main pancreatic duct, draining into the duodenal lumen via the major duodenal papilla.

\synonyms
Hepatopancreatic Ampulla

\medterm{Ampullary Carcinoma} A malignant epithelial tumor arising in or around the ampulla of Vater, where the bile duct and pancreatic duct enter the duodenum.
It can block bile flow and cause jaundice, abnormal liver tests, pancreatitis, or gastrointestinal bleeding.

\synonyms
Ampullary Cancer

\medterm{Amputation} Surgical removal or traumatic loss of a limb or part of a limb, such as a finger, toe, arm, or leg. Common reasons include severe injury, serious infection, poor blood flow causing tissue death or severe pain, and cancer. The level of amputation depends on healthy tissue, blood supply, healing, expected function, and rehabilitation needs.

\synonyms
Surgical Amputation; Limb Removal

\medterm{Amputation Stump} The part of a limb left after amputation. It may develop wound problems, infection, poor circulation, nerve pain, phantom-limb pain, or skin irritation from a prosthesis.

\medterm{Amsler Grid} A small square chart used to check central vision. Wavy, missing, or distorted lines may indicate a problem in the macula and are evaluated with a complete eye examination.

\medterm{Amusia} Difficulty recognizing, understanding, or producing music despite adequate hearing. It may be present from birth or develop after brain injury, stroke, or another neurological condition.

\medterm{Amygdala} A group of nerve cells deep in each side of the brain that helps process emotions, especially fear, and contributes to learning, memory, and responses to important or threatening events.

\synonyms
Amygdaloid Nucleus

\medterm{Amylase} An enzyme that begins breaking down starch. It is made mainly by the salivary glands and pancreas and can be measured in blood or urine when pancreatic or salivary-gland disease is suspected.

\medterm{Amylase And Lipase} Blood tests for digestive enzymes used when pancreatic disease is suspected. Lipase is generally more specific for pancreatic injury, but neither test alone confirms acute pancreatitis.

\medterm{Amylase Test} A blood or urine test that measures the amount of amylase, an enzyme made mainly by the salivary glands and pancreas. It may help assess pancreatitis and other problems affecting the pancreas or salivary glands.

\synonyms
Serum Amylase; Blood Amylase Test; Amylase Blood Test

\medterm{Amylin Analogues} Medicines that imitate amylin, a hormone released by the pancreas with insulin. They slow stomach emptying, reduce the rise in blood glucose after meals, and reduce glucagon release.

\medterm{Amyloid} An abnormal protein material that forms insoluble deposits in tissues. These deposits can interfere with the function of organs such as the kidneys, heart, nerves, liver, or brain.
\synonyms
Amyloid Protein

\medterm{Amyloid Nephropathy} Renal involvement occurring in systemic amyloidosis, characterized by extracellular accumulation of insoluble amyloid fibrils in glomerular, vascular, and interstitial compartments, typically leading to heavy proteinuria, nephrotic syndrome, and progressive kidney failure.

\medterm{Amyloid Neuropathy} Nerve damage caused by amyloid deposits in or around peripheral nerves. It may cause numbness, tingling, burning pain, weakness, or problems with automatic body functions such as blood-pressure control.

\medterm{Amyloid PET} A brain scan using a small radioactive tracer that attaches to amyloid deposits. It can help assess selected cognitive disorders, but the result must be interpreted with symptoms and other tests.

\medterm{Amyloid Plaque} A deposit of amyloid-beta protein outside nerve cells in the brain. Plaques are associated with Alzheimer disease, but their presence alone does not explain every person’s memory or thinking problem.

\medterm{Amyloid Precursor Protein} A normal protein found in nerve-cell membranes. Enzymes can cut it into smaller pieces, including amyloid-beta; abnormal accumulation of amyloid-beta in the brain is associated with Alzheimer disease.

\synonyms
APP; Amyloid-Beta Precursor Protein

\medterm{Amyloid Purpura} Purple or bruise-like spots, often around the eyes, caused by amyloid deposits weakening small blood vessels. It may occur with systemic amyloidosis and can be accompanied by easy bruising elsewhere.

\medterm{Amyloidosis} A group of diseases in which proteins fold abnormally, clump together, and form deposits called amyloid in tissues or organs. Amyloidosis may be localized to one site or systemic, affecting several organs such as the kidneys, heart, liver, nerves, or digestive tract. Major forms include light-chain (AL), serum amyloid A (AA), and transthyretin (ATTR) amyloidosis; effects vary with the amyloid protein and organs involved.

\synonyms
Amyloid Disease; Amyloid Infiltration

\medterm{Amyotrophic Lateral Sclerosis} A progressive disease that damages and destroys motor neurons, the nerve cells in the brain and spinal cord that control voluntary muscle movement and breathing. As these cells are lost, muscles become weak, stiff, wasted, or twitchy, eventually causing paralysis. Symptoms may begin in an arm, leg, or the muscles used for speaking and swallowing, and may later cause difficulty walking, speaking, chewing, swallowing, or breathing. Sensation is usually preserved, but some people develop changes in thinking or behavior, sometimes with frontotemporal dementia. Most cases have no known cause; about one in ten are inherited. ALS is also called Lou Gehrig disease.

\synonyms
Lou Gehrig's Disease
ALS

\medterm{Anabolic Steroid Induced Hypogonadism} Reduced testicular hormone production and fertility caused by misuse of anabolic steroids. Extra steroid hormones suppress the body’s normal hormone signals and may lead to testicular shrinkage, low sex drive, erectile problems, infertility, or depression.

\synonyms
ASIH; Anabolic-Androgenic Steroid-Induced Hypogonadism

\medterm{Anabolic Steroid Use and Abuse} Use of anabolic-androgenic steroids without a medical indication or at excessive doses. It can suppress natural hormone production and affect fertility, mood, blood pressure, cholesterol, liver function, heart structure, and growth in adolescents.

\medterm{Anabolic Steroids} Synthetic medicines related to testosterone that promote muscle and bone growth. They have limited medical uses but are sometimes misused for appearance or performance; misuse can harm the heart, liver, fertility, mood, and hormone system.

\medterm{Anabolism} The set of body processes that use energy to build larger molecules from smaller ones, such as making proteins from amino acids or storing glucose as glycogen.

\medterm{Anaemia} A condition in which the blood has too few red blood cells or too little haemoglobin. This reduces the blood’s ability to carry oxygen and may cause tiredness, weakness, breathlessness, or dizziness.

\medterm{Anaerobe} A microorganism that can grow without oxygen. Some anaerobes are harmed or killed by oxygen, while others tolerate oxygen but do not use it to produce energy.

\medterm{Anaerobic} Able to live, grow, or function without oxygen. Anaerobic bacteria can cause infections in deep wounds, abscesses, the mouth, abdomen, or pelvis.

\medterm{Anaerobic Bacteria} Bacteria that can grow without oxygen. Obligate anaerobes cannot use oxygen and may be harmed by it, while facultative anaerobes can grow with or without oxygen. Anaerobic bacteria normally live in places such as the mouth, intestine, and vagina. They can cause infections in deep wounds, abscesses, the abdomen, pelvis, or other tissues with low oxygen.

\medterm{Anaerobic Glycolysis} The metabolic breakdown of glucose into lactate in the absence or insufficiency of cellular oxygen, producing a net yield of two molecules of adenosine triphosphate per glucose molecule and regenerating oxidized nicotinamide adenine dinucleotide (NAD+) to maintain energy generation.

\medterm{Anagen} The active growth phase of the hair cycle, when cells in the hair root make the hair shaft. Scalp hairs usually remain in this phase for several years before entering a resting or shedding phase.

\medterm{Anagen Effluvium} An abrupt, extensive shedding of hair during the active growth (anagen) phase, caused by acute toxic disruption of mitotic activity in the follicular matrix cells. It typically begins within days to weeks of exposure to cytotoxic chemotherapy, ionizing radiation, or acute heavy metal poisoning.

\medterm{Anal Abscess} A collection of pus near or within the anus, usually caused by an infected anal gland. It can cause throbbing pain, swelling, fever, and difficulty sitting.

\medterm{Anal Canal} The final part of the digestive tract, extending from the rectum to the anus. It contains muscles that control the passage of stool and has different nerve and blood supplies along its length.

\medterm{Anal Cancer} Cancer arising in the anal canal or perianal skin, most often squamous-cell carcinoma associated with persistent high-risk HPV infection. Bleeding, pain, a lump, itching, or altered bowel function may occur; examination and biopsy establish the diagnosis.

\medterm{Anal Carcinoma} Cancer arising in the anal canal or anal margin, most often squamous-cell cancer associated with persistent high-risk human papillomavirus infection. It may cause bleeding, pain, itching, or a lump.

\synonyms
Anal Cancer

\medterm{Anal Condyloma} A wart-like growth on or around the anus, usually caused by low-risk types of human papillomavirus. It is usually not cancerous, but multiple or recurring warts may occur with other abnormal anal cells.

\medterm{Anal Fissure} A small tear in the lining of the anus, usually caused by hard stools, constipation, or repeated diarrhea. It can cause sharp pain during bowel movements and bright-red blood on the toilet paper.

\medterm{Anal Fistula} An abnormal tunnel between the anal canal and nearby skin, often developing after an anal abscess. It may cause repeated discharge, pain, itching, or a small opening near the anus.

\medterm{Anal Intercourse} Sexual activity involving insertion into the anus. Because the anal lining can tear and does not provide natural lubrication, it can transmit infections and cause injury if precautions are not used.

\medterm{Anal Intraepithelial Neoplasia} Abnormal cells in the lining of the anus, usually associated with persistent high-risk human papillomavirus infection. Low-grade changes may resolve, while high-grade changes can progress to anal cancer.

\medterm{Anal Itching} Itching around the anus. Common causes include moisture, irritation, dermatitis, haemorrhoids, infection, and occasionally pinworms or another disease.

\synonyms
Pruritus Ani; Perianal Itching; Rectal Itching

\medterm{Anal Mucositis} Inflammation and open sores of the anal lining, often caused by radiation or chemotherapy. It can cause pain, burning, bleeding, discharge, and difficulty passing stool.

\medterm{Anal Necrosis} Death of tissue around the anus, anal canal, or nearby skin. It may result from severe infection, poor blood flow, injury, or radiation and can cause severe pain, skin discolouration, blisters, fever, or rapid worsening.

\medterm{Anal Pain} Pain around the anus or nearby skin. Common causes include a fissure, abscess, thrombosed haemorrhoid, muscle spasm, infection, or inflammation.

\synonyms
Proctalgia

\medterm{Anal Papilla} A small, usually normal projection of the anal lining near the junction between the upper and lower types of lining in the anal canal.

\medterm{Anal Sphincter} A muscular structure that controls passage of stool and gas from the anal canal. It includes the involuntary internal anal sphincter and the voluntary external anal sphincter.

\medterm{Analgesia} Relief or absence of pain without necessarily causing unconsciousness. It may result from medicines, a nerve block, or another treatment; touch, pressure, or movement may still be felt.

\medterm{Analgesics} Medicines that reduce or relieve pain. Non-opioid analgesics reduce pain and fever. Nonsteroidal anti-inflammatory drugs (NSAIDs) reduce pain, fever, and inflammation by lowering the production of substances involved in these responses. Opioid analgesics reduce the brain's and spinal cord's response to pain but may cause sleepiness, constipation, slowed breathing, tolerance, and dependence. Adjuvant analgesics are medicines mainly used for other conditions that can help with certain types of pain, such as nerve pain. The choice depends on the type and severity of pain, the person's health, and possible risks.

\synonyms
Painkillers; Pain-Relieving Medicines

\medterm{Analog} A substance, structure, or process that resembles another and has a similar function but is not identical in chemical makeup or structure.

\medterm{Analogous} Similar in function or purpose to another structure or process, even though its origin or physical structure is different.

\medterm{Analysis} A careful examination of information, samples, measurements, symptoms, or clinical problems to identify patterns, understand causes, or support a medical decision.

\medterm{Analysis Of Urine} Laboratory examination of urine for blood, cells, protein, glucose, germs, and other substances. Results can help assess urinary infection, kidney disease, diabetes, bleeding, dehydration, and other conditions.

\medterm{Analytic Sensitivity} The smallest amount of a substance that a laboratory method can reliably detect. It describes the test’s measurement ability and is different from clinical sensitivity, which describes how well a test detects disease.

\medterm{Anaphia} Complete loss of touch sensation in a body area. It may result from nerve injury, neurological disease, reduced blood flow, or anaesthetic treatment and can allow burns, wounds, or pressure injuries to go unnoticed.

\synonyms
Anaptic

\medterm{Anaphylactic Shock} A severe form of anaphylaxis in which a sudden allergic reaction causes dangerously low blood pressure and poor blood flow to organs. It may occur with airway swelling, breathing difficulty, collapse, or cardiac arrest.

\medterm{Anaphylactoid Reaction} An immediate, severe systemic hypersensitivity reaction that closely mimics IgE-mediated anaphylaxis in its clinical manifestations, but occurs via direct non-immune or non-IgE-mediated mast cell and basophil degranulation (such as that triggered by radiocontrast media or certain drugs).

\synonyms
Non-Allergic Anaphylaxis; Pseudoallergic Reaction

\medterm{Anaphylatoxin} A small peptide fragment (predominantly C3a, C4a, and C5a) generated during activation of the complement cascade. Anaphylatoxins induce mast cell degranulation, smooth muscle contraction, increased vascular permeability, and chemotaxis of phagocytic leukocytes.

\medterm{Anaphylaxis} A sudden, potentially life-threatening systemic hypersensitivity reaction that can involve the skin, airways, lungs, circulation, or gastrointestinal tract. It may cause hives or swelling, breathing difficulty, vomiting, dizziness, dangerously low blood pressure, or collapse; skin findings may be absent. Immediate intramuscular epinephrine and emergency medical care are required.

\medterm{Anaphylaxis During Anesthesia} A severe allergic reaction occurring during or soon after anaesthesia. It may cause sudden low blood pressure, airway narrowing, breathing difficulty, collapse, or skin flushing; triggers can include medicines, antiseptics, or latex.

\medterm{Anaplasia} Loss of the normal mature features and specialized function of cells. Anaplastic cells may vary greatly in size and shape, have abnormal nuclei, and divide rapidly or unusually. Anaplasia is most often seen in high-grade cancers and usually suggests more aggressive behavior. It is different from metaplasia, in which one mature cell type changes into another cell type.

\medterm{Anaplasmosis} A tick-borne infection caused by the bacterium \textit{Anaplasma phagocytophilum}. It may cause fever, headache, muscle aches, low platelet and white-cell counts, and raised liver enzymes. PCR or paired blood-antibody tests can help confirm it. Morulae seen in white blood cells can support the diagnosis but are not always present and do not confirm it alone.

\synonyms
Human Granulocytic Anaplasmosis; HGA

\medterm{Anaplastic} Describing cells or a tumor that has lost many features of its normal mature form. Anaplastic tumors often grow quickly, invade nearby tissue, or spread, although the outlook depends on the specific cancer and its treatment.

\medterm{Anaplastic Astrocytoma} An older term for a high-grade tumour arising from star-shaped supporting cells in the brain. Modern classification relies on its genetic features as well as its appearance; symptoms depend on its location.

\medterm{Anaplastic Ependymoma} A high-grade tumour arising from cells lining the brain’s fluid-filled spaces or the spinal canal. It can grow quickly and may obstruct the flow of cerebrospinal fluid.

\medterm{Anaplastic Large Cell Lymphoma} A rare type of non-Hodgkin lymphoma arising from abnormal T cells. The cells are usually large and carry the CD30 marker; the disease may affect the skin, lymph nodes, or other organs.

\synonyms
ALCL

\medterm{Anaplastic Thyroid Cancer} A rare, fast-growing thyroid cancer that may appear as a rapidly enlarging neck lump and cause difficulty swallowing, speaking, or breathing.

\medterm{Anaplastic Thyroid Carcinoma} An extremely aggressive, undifferentiated malignancy of the thyroid follicular epithelium characterized by rapid local invasion into adjacent cervical structures and early distant metastasis. Presenting symptoms include a rapidly enlarging hard neck mass, hoarseness, dysphagia, and dyspnea, carrying a very poor median survival despite multimodality therapy.

\synonyms
Undifferentiated Thyroid Carcinoma; ATC

\medterm{Anarithmetia} A primary neuropsychological impairment in basic arithmetic reasoning, computational ability, and understanding of numerical concepts, distinguishing it from secondary calculation deficits caused by spatial or linguistic impairments. It is typically associated with lesions of the dominant inferior parietal lobule (angular gyrus).

\medterm{Anasarca} Severe, widespread swelling caused by excess fluid under the skin and sometimes around organs. It can occur with serious heart, kidney, or liver disease or very low blood protein and may cause rapid weight gain or breathing difficulty.

\medterm{Anastomosis} A surgical, traumatic, or pathological connection created between two hollow organs, tubular structures, or blood vessels (such as an arterial anastomosis, arterio-venous fistula, or bowel resection re-connection), establishing communication and allowing continuous fluid or luminal flow.

\synonyms
Surgical Anastomosis; Vascular Anastomosis; Enteroanastomosis

\medterm{Anastomotic Leak} Leakage from a connection made during surgery, often where two sections of intestine or another hollow organ were joined. An incomplete seal can allow contents into nearby tissue and cause pain, fever, infection, or life-threatening illness.

\medterm{Anastrozole} A medicine that lowers the body’s production of oestrogen. It is used mainly for hormone-sensitive breast cancer, especially after menopause.

\medterm{Anatomic Orientation Terms} Standard words used to describe the position of one body structure relative to another, including above, below, front, back, near, far, inside, and outside.

\medterm{Anatomical Axis} An imaginary line around which a body part rotates or moves. Anatomical axes help describe movements such as rotation of the shoulder, hip, spine, or other joints.

\medterm{Anatomical Conduit} A natural or surgically created passage that carries air, fluid, tissue, or another structure between body regions, such as an airway, blood vessel, or surgical channel.

\medterm{Anatomical Orifice} A natural opening into or out of the body or through a body structure, such as the mouth, nostril, anus, urethra, or an opening of a blood vessel.

\synonyms
Body Opening

\medterm{Anatomical Planes} Imaginary flat surfaces used to divide the body into sections for describing anatomy and interpreting medical scans. Common planes are sagittal, coronal, and transverse.

\medterm{Anatomical Position} The standard reference posture for describing the body: standing upright, facing forward, arms at the sides, hands turned forward, and feet directed forward.

\medterm{Anatomical Snuffbox} A triangular hollow on the thumb side of the back of the wrist that appears when the thumb is extended. It overlies the scaphoid bone and radial artery; tenderness there after a fall may indicate a scaphoid fracture.

\synonyms
Anatomic Snuffbox

\medterm{Anatomist} A scientist or health professional who studies the structure of the human or animal body, including its organs, tissues, blood vessels, nerves, and their relationships.

\medterm{Anatomy} The study of body structure and how its parts relate to one another. It includes visible structures, tissues and cells seen with magnification, and how the body develops.

\medterm{ANCA-Associated Vasculitis} A group of rare diseases in which the immune system attacks small blood vessels. It can affect the kidneys, lungs, nerves, skin, eyes, and other organs. A blood test for antineutrophil cytoplasmic antibodies (ANCA) can help with diagnosis but cannot confirm the disease on its own.

\medterm{Ancillary Testing, Pathology} Supplementary diagnostic modalities—including immunohistochemistry, special histochemical stains, flow cytometry, in situ hybridization, and molecular genetic assays—performed in conjunction with standard light microscopy to achieve precise pathological classification.

\synonyms
Ancillary Diagnostic Testing; Diagnostic Immunohistochemistry; Special Pathology Stains

\medterm{Ancylostomiasis} A parasitic nematode infection caused by species of the genus \textit{Ancylostoma} (such as \textit{Ancylostoma duodenale}), acquired through transcutaneous penetration of infective larvae from contaminated soil, characteristically leading to chronic intestinal blood loss, iron deficiency anemia, and hypoproteinemia.

\synonyms
Hookworm Infection; Old World Hookworm Disease; Ancylostoma Duodenale Infection

\medterm{Andersen-Tawil Syndrome} A rare inherited disorder causing episodes of muscle weakness, abnormal heart rhythms, and distinctive physical features such as a small lower jaw, low-set ears, or widely spaced eyes.

\synonyms
ATS; Periodic Paralysis Type 3; Andersen Syndrome

\medterm{Anderson Disease} A rare inherited disorder in which abnormal glycogen builds up in cells because of changes in the \textit{GBE1} gene. It usually begins in infancy with poor growth and an enlarged liver and may progress to scarring, high blood pressure in the liver, and liver failure.

\synonyms
Glycogen Storage Disease Type IV (GSD IV); Amylopectinosis

\medterm{Androgen} A hormone such as testosterone that helps regulate sexual development, fertility, muscle, bone, and hair growth in people of all sexes.

\medterm{Androgen Deficiency} Too little production or effect of androgen hormones, especially testosterone. In males, it may cause reduced sexual desire, erectile difficulty, infertility, loss of muscle or bone, low energy, or reduced facial and body hair.

\medterm{Androgen Excess} Abnormally high levels or effects of androgen hormones. It may cause acne, increased facial or body hair, scalp hair loss, irregular periods, infertility, or a deeper voice, especially in females.
\synonyms
Hyperandrogenism

\medterm{Androgen Insensitivity Syndrome} An inherited condition in which cells cannot respond fully to testosterone and related hormones. A person with XY chromosomes may therefore have female-appearing external genitals, incomplete male development, undescended testes, and no uterus.

\medterm{Androgen Insufficiency} A level or effect of androgen hormones that is too low for normal body function. Symptoms may include reduced sexual desire, erectile difficulty, infertility, low energy, reduced muscle mass, or weaker bones.

\synonyms
Androgen Deficiency

\medterm{Androgen Receptor} A protein in cells that binds testosterone and related hormones. This allows the hormones to influence sexual development, reproduction, muscle, bones, hair, and some cancers.

\medterm{Androgen Suppression} Lowering the production of androgen hormones or blocking their action in the body. It is used mainly for some hormone-sensitive cancers, especially prostate cancer, and can cause hot flushes, sexual problems, bone loss, weight change, and metabolic effects.

\synonyms
Androgen Deprivation

\medterm{Androgenetic Alopecia} A common, inherited form of gradual hair loss caused by sensitivity of hair follicles to androgen hormones. It usually produces a receding hairline or crown thinning in males and diffuse thinning over the crown in females.

\medterm{Androgenic} Relating to androgen hormones, such as testosterone, or producing their effects. These effects may include facial or body hair, oily skin, a deeper voice, muscle growth, and development of male reproductive features.

\medterm{Android Pelvis} A relatively narrow, heart-shaped or funnel-shaped female pelvis with reduced space in the middle. During childbirth, this shape may make passage of the baby more difficult than some other pelvic shapes.

\medterm{Andropause} Alternate informal term; see \textit{Male Hypogonadism}.

\medterm{Anejaculation} Inability to release semen during orgasm or sexual activity. It may result from nerve or spinal disease, pelvic surgery, medicines, diabetes, hormonal problems, or psychological factors and may cause infertility.

\medterm{Anemia} A condition in which the blood cannot carry enough oxygen to the body’s tissues, usually because there are too few red blood cells or too little haemoglobin. Causes include blood loss, reduced production, and increased breakdown.

\synonyms
Anaemia; Low Red Blood Cell Count

\medterm{Anemia and Thrombocytopenia in Pregnancy} Low red-cell or platelet counts during pregnancy may reflect iron deficiency, dilution, nutritional disease, immune disorders, hypertensive disease, bleeding, infection, or marrow disease. Severity, cause, symptoms, and trends guide treatment and delivery planning.

\medterm{Anemia In Pregnancy} Anemia during pregnancy, commonly caused by iron deficiency, folate deficiency, vitamin B12 deficiency, blood loss, or the normal increase in the liquid part of blood. It can cause tiredness, weakness, breathlessness, or dizziness.

\synonyms
Gestational Anemia; Pregnancy-Associated Anemia

\medterm{Anemia Of Chronic Disease} Anemia associated with long-term inflammation, infection, cancer, kidney disease, or another chronic illness. Inflammation limits the use of stored iron and reduces red-cell production; mild cases may cause no symptoms.

\medterm{Anemia Of Inflammation} Anemia caused by long-lasting infection, inflammatory disease, cancer, or another condition that activates the immune system. Inflammation traps iron in storage and slows red-cell production.

\medterm{Anemia of Malignancy} A multifactorial normocytic or microcytic anemia developing in patients with neoplastic disease, resulting from cytokine-mediated blunting of erythropoiesis, impaired iron utilization, marrow infiltration, hemorrhage, or myelosuppressive antineoplastic therapies.

\synonyms
Cancer-Related Anemia; Tumor-Induced Anemia; Anemia of Cancer

\medterm{Anemic} Having anemia, meaning that the blood carries less oxygen than the body needs because it has too few red blood cells or too little haemoglobin.

\medterm{Anencephaly} A serious neural tube birth defect in which the upper part of the neural tube fails to close during the first weeks of pregnancy. As a result, much of the brain, especially the forebrain, and parts of the skull and scalp do not develop. It is usually detected before birth and is fatal; most affected babies are stillborn or die shortly after birth. There is no treatment that can correct the defect.

\medterm{Anergy} A state in which an immune cell does not respond to a substance that would normally activate it. The cell remains present but is functionally quiet; this can help prevent attacks against the body’s own tissues.

\medterm{Anesthesia} A medically controlled loss of feeling caused by medicines during a procedure. Anesthetic medicines mainly block nerve signals, especially pain signals, from reaching the brain. Local anesthesia numbs a small area while the person remains awake. Regional anesthesia blocks nerves supplying a larger part of the body and may be used while the person remains awake. General anesthesia also changes brain activity to cause controlled unconsciousness. Before modern anesthesia, surgery was extremely painful and had to be very quick; anesthesia made longer and more complex operations possible. Breathing, blood pressure, heart rate, and oxygen levels are monitored.

\synonyms
Anaesthesia

\medterm{Anesthesia Gas Analyzers} Devices that measure the amounts of oxygen, carbon dioxide, and inhaled anesthetic gases delivered to and breathed out by a patient during anesthesia. They help detect problems with breathing or gas delivery.

\medterm{Anesthesia Informatics} The use of electronic records, monitoring data, and computer systems to support anesthesia care. It can improve documentation, identify medicine interactions or dosing errors, and help evaluate outcomes.

\medterm{Anesthesia Machine} Equipment that supplies controlled mixtures of oxygen, air, and anesthetic gases, and supports breathing during anesthesia. It includes systems that remove carbon dioxide and safety devices that warn of delivery problems.

\medterm{Anesthesia Monitoring} Continuous assessment of vital functions during anesthesia, including oxygen levels, breathing, blood pressure, heart rate, temperature, and level of unconsciousness as appropriate for the procedure.

\medterm{Anesthesia Pharmacogenomics} The study of how inherited genetic differences affect responses to anesthetic and pain medicines. Examples include prolonged muscle paralysis from pseudocholinesterase deficiency and inherited susceptibility to malignant hyperthermia.

\medterm{Anesthesiologist} A doctor trained to provide anesthesia and manage pain, breathing, blood pressure, and other vital functions during procedures. Anesthesiologists also assess patients before surgery and may treat critical illness or severe pain.

\medterm{Anesthesiology} The medical specialty covering anesthesia, monitoring and support of vital functions during procedures, care before and after surgery, critical care, and treatment of acute or chronic pain.

\medterm{Anetoderma} A rare disorder caused by loss or weakening of elastic fibers in the dermis, producing soft, wrinkled, depressed, or pouch-like macules and plaques with a characteristic button-hole feel. It may be primary or follow infection, inflammation, autoimmune disease, or another skin lesion; biopsy and selected systemic evaluation may be needed.

\medterm{Aneuploidy} An abnormal number of whole chromosomes in a cell, caused by an error during cell division. A cell may have an extra or missing chromosome, as in trisomy 21, and this can affect development or pregnancy.

\medterm{Aneurysm} A weakened area of a blood-vessel wall that bulges or enlarges. Aneurysms can occur in arteries and rarely veins; they may enlarge, form clots, tear, or rupture and cause life-threatening bleeding.

\medterm{Aneurysm Clip} A small surgical device placed across the opening of a brain aneurysm to stop blood from entering it. It is usually placed during open brain surgery; rarely, it may move, block a nearby vessel, or allow bleeding to recur.

\medterm{Aneurysm Clipping} Brain surgery to place a small clip across the opening of a brain aneurysm, a weak bulge in a brain artery. The clip stops blood from entering the bulge and lowers the risk of bleeding, but the procedure has its own risks and cannot guarantee that a stroke or future bleeding will not occur.
\synonyms
Surgical Clipping; Cerebral Aneurysm Clipping; Microvascular Clipping

\medterm{Aneurysm Coils} Soft metal coils placed inside an aneurysm through a catheter. The coils promote clotting within the bulge and reduce blood flow into it, lowering the chance of rupture.

\medterm{Aneurysm Repair} A procedure that reinforces, closes, or replaces a weakened section of a blood vessel to prevent rupture or other complications. It may be performed through open surgery or through a catheter inside the vessel.

\medterm{Aneurysms} Weakened, bulging areas in blood-vessel walls. They may remain stable or enlarge, form clots, press on nearby structures, tear, or rupture and cause life-threatening bleeding.

\medterm{Angelman Syndrome} A genetic disorder affecting brain development and function, usually because the maternal copy of the \textit{UBE3A} gene is missing or inactive. It causes severe developmental delay, very limited speech, movement and balance problems, seizures, and a characteristic happy or excitable manner.

\medterm{Anger} An emotional response to frustration, unfairness, danger, or a perceived threat. It may range from irritation to intense rage; persistent or uncontrollable anger can affect relationships, safety, and health.

\medterm{Angiitis} Pathological inflammation and necrosis of blood-vessel walls, leading to lumen stenosis, ischemia, aneurysm formation, or hemorrhage in affected organ systems.

\synonyms
Vasculitis; Blood Vessel Inflammation; Arteritis

\medterm{Angina} Pressure, tightness, burning, or discomfort caused by temporary reduction of blood flow to the heart muscle, usually because of narrowed coronary arteries. It may spread to the arm, jaw, back, or neck and may occur with activity or at rest.

\medterm{Angina Bullosa Hemorrhagica} A usually benign, sudden blood-filled blister of the oral mucosa, most often the soft palate. It may follow minor trauma, inhaled or topical steroids, or occur with diabetes; blisters usually rupture and heal without treatment, but recurrent or atypical lesions need evaluation to exclude autoimmune blistering disease.

\medterm{Angina Pectoris} Chest discomfort caused by temporary reduced blood flow and oxygen to the heart muscle, usually because of narrowed coronary arteries. It may feel like pressure, squeezing, tightness, burning, or indigestion and may spread to the arm, shoulder, neck, jaw, back, or upper abdomen. Stable angina usually occurs with exercise or emotional stress and improves with rest. Unstable angina is new, worsening, prolonged, or occurs at rest and is associated with a higher risk of heart attack.

\medterm{Angina Symptoms} Symptoms of reduced blood flow to the heart may include pressure, squeezing, burning, or tightness in the chest, with discomfort spreading to the arm, jaw, back, or shoulder. Breathlessness, sweating, nausea, or unusual tiredness may also occur.

\synonyms
Ischemic Chest Pain; Anginal Chest Pain; Anginal Discomfort

\medterm{Angioblastoma} A usually slow-growing, highly vascular tumour of the central nervous system, most often found in the cerebellum or spinal cord. It may cause headache, balance problems, weakness, or other symptoms from pressure on nearby tissue.

\medterm{Angiocardiography} Imaging of the heart and major blood vessels, usually after contrast material is given. X-ray, CT, or MRI can show the heart chambers, valves, narrowed vessels, leaks, abnormal connections, or other structural problems.

\medterm{Angiodysplasia} A common acquired vascular malformation characterized by fragile, dilated, tortuous submucosal blood vessels in the gastrointestinal tract, most frequently found in the cecum and right colon of adults over age 60. It is a major cause of recurrent, painless lower gastrointestinal bleeding and unexplained iron-deficiency anemia. It is strongly associated with aortic valve stenosis (Heyde's syndrome, due to acquired von Willebrand factor deficiency) and chronic renal failure. Diagnosis is established on colonoscopy, where lesions appear as small, bright-red, spider-like mucosal patches; active or high-risk bleeding is treated endoscopically using argon plasma coagulation (APC) or electrocautery.
\synonyms{Colonic Angiodysplasia; Arteriovenous Malformation of the Colon; Vascular Ectasia of the Colon}

\medterm{Angiodysplasia Of The Colon} Alternate name; see \textit{Angiodysplasia}.

\medterm{Angioedema} Sudden, localized swelling of deeper skin or mucosal tissues, often affecting the lips, eyelids, tongue, throat, hands, feet, or bowel wall. Histamine-mediated episodes may occur with hives and itching, whereas bradykinin-mediated forms usually lack hives; tongue or throat swelling, voice change, or breathing difficulty is an emergency.

\medterm{Angioedema And Hives} A combination of raised, itchy skin patches called hives and deeper swelling called angioedema. Common triggers include foods, medicines, insect stings, infections, and physical exposures; breathing difficulty is an emergency.

\medterm{Angioedema, Hereditary} Alternate index label; see \textit{Hereditary Angioedema}.

\medterm{Angiofibroma} A usually noncancerous growth made of blood vessels and fibrous tissue. It commonly appears as a small facial papule in adolescents or, in a different form, as a vascular tumour in the back of the nose that may cause blockage or nosebleeds.

\medterm{Angiofollicular Lymph Node Hyperplasia} An uncommon disorder in which lymph nodes contain excess small blood vessels and abnormal immune-cell growth. It is also called Castleman disease and may affect one lymph-node region or many parts of the body.

\medterm{Angiogenesis} The formation of new blood vessels from existing vessels. It is normal during growth, menstruation, and wound healing, but abnormal angiogenesis can help cancers grow and can damage vision in some eye diseases.

\medterm{Angiogenesis Inhibition} Preventing or slowing the growth of new blood vessels. This can limit a tumor’s blood supply or reduce harmful vessel growth in some eye diseases. Treatment may increase the risk of high blood pressure, bleeding, or delayed wound healing.

\synonyms
Antiangiogenic Therapy

\medterm{Angiogenesis Inhibitor} A therapeutic agent that inhibits the formation of new blood vessels, predominantly by antagonizing vascular endothelial growth factor (VEGF) or its receptors, used in oncology to restrict tumor blood supply and in ophthalmology to treat neovascular macular degeneration.

\synonyms
Antiangiogenic Agent

\medterm{Angiogram} An image of blood vessels, usually made after contrast material is given. It can show narrowing, blockage, aneurysm, tearing, abnormal connections, or bleeding and may be produced by catheter X-ray, computed tomography, or magnetic resonance imaging.

\medterm{Angiography} Imaging of blood vessels to show their shape and blood flow. It may use X-rays with injected contrast material, computed tomography (CT), magnetic resonance imaging (MRI), or a catheter placed inside a blood vessel. Angiography can show narrowing, blockage, aneurysms, abnormal connections, bleeding, or other vascular abnormalities; catheter angiography may also allow treatment during the same procedure.

\medterm{Angioid Streak} Irregular dark lines caused by breaks in a tough layer beneath the retina. They may bleed or allow abnormal blood vessels to grow, reducing vision, and can be associated with some inherited disorders.

\medterm{Angioimmunoblastic T-Cell Lymphoma} An aggressive subtype of peripheral mature T-cell lymphoma characterized by generalized lymphadenopathy, constitutional symptoms (fever, night sweats, weight loss), hepatosplenomegaly, skin rashes, and polyclonal hypergammaglobulinemia.

\synonyms
AITL

\medterm{Angiokeratoma of the Scrotum} A benign vascular papule caused by superficial capillary dilation and overlying thickened skin, usually appearing as multiple dark-red, blue, or purple papules on the scrotum. Lesions may bleed after friction and can be mistaken for melanoma or other lesions; examination establishes the diagnosis.

\medterm{Angiolymphoid Hyperplasia With Eosinophilia} An uncommon benign vascular and inflammatory lesion, also called epithelioid hemangioma, that produces solitary or grouped red-brown papules or nodules, usually on the head and neck near the ears or scalp. Eosinophilia may occur; biopsy distinguishes it from Kimura disease and other vascular or neoplastic lesions.

\medterm{Angiomyofibroblastoma} A usually noncancerous, slow-growing tumour of blood-vessel and connective-tissue cells, most often occurring as a small, well-defined, painless mass in the vulva or vagina.

\medterm{Angiomyolipoma} A usually noncancerous kidney tumour made of blood vessels, muscle, and fat. It may occur alone or with tuberous sclerosis and can bleed, especially when large.

\medterm{Angioneurotic Edema, Hereditary} An older name for hereditary angioedema, an inherited disorder causing repeated swelling of the skin, bowel, or airway. Attacks usually occur without hives and throat swelling can block breathing.

\medterm{Angiopathy} Any disease or abnormal change in blood vessels. Causes include diabetes, inflammation, fatty deposits, injury, infection, and inherited disorders; effects depend on the vessels and organs involved.

\medterm{Angioplasty} A minimally invasive procedure that widens a narrowed or blocked blood vessel to improve blood flow. A thin catheter with a small balloon is guided to the narrowed area, often using X-ray imaging. The balloon is inflated briefly and then removed. A small mesh tube called a stent may be left in place to help keep the vessel open. Angioplasty may be performed on coronary arteries or other arteries and veins.

\medterm{Angioplasty Balloon} An inflatable cylindrical balloon catheter designed to dilate stenotic vascular segments, fracture atheromatous plaques, or deploy and seat endovascular stents.

\synonyms
Balloon Catheter; Dilatation Balloon; Angioplasty Catheter

\medterm{Angiosarcoma} A rare, fast-growing cancer that begins in cells lining blood or lymph vessels. It may occur in the skin, breast, liver, or bone and can spread to other organs.

\medterm{Angiosarcoma Of The Breast} A rare cancer arising from blood-vessel cells in the breast. It may cause an enlarging breast mass or bruise-like skin colour change and can develop after breast radiation or long-term lymphoedema.

\medterm{Angiosarcoma Of The Heart} A rare, aggressive cancer arising from blood-vessel cells in the heart, most often in the right upper chamber. It may obstruct blood flow, cause fluid around the heart or abnormal rhythms, and spread to other organs.

\medterm{Angiotensin} A group of hormones that help regulate blood pressure and the body’s salt and water balance. Angiotensin II narrows blood vessels and signals the kidneys and adrenal glands to retain salt and water.

\medterm{Angiotensin Converting Enzyme} An enzyme that changes angiotensin I into angiotensin II, which raises blood pressure by narrowing blood vessels. The enzyme also breaks down bradykinin, a substance that can widen vessels.

\medterm{Angiotensin I} A mostly inactive hormone made from angiotensinogen. Angiotensin-converting enzyme changes it into angiotensin II, which narrows blood vessels and promotes salt and water retention.

\medterm{Angiotensin II} A hormone that narrows blood vessels and causes the body to retain salt and water. These effects increase blood pressure and help maintain blood flow when circulating volume is low.

\medterm{Angiotensin Receptor Blockers} Medicines that prevent angiotensin II from activating its main receptor. They relax blood vessels and reduce salt and water retention, lowering blood pressure and reducing strain on the heart and kidneys.

\medterm{Angiotensin-Converting Enzyme 2} A cell-surface zinc metalloprotease that cleaves angiotensin II into the vasodilatory and cardioprotective peptide angiotensin-(1-7), acting as an essential counter-regulatory enzyme within the renin-angiotensin-aldosterone system; also serves as the host cell receptor for the SARS-CoV-2 spike protein.

\synonyms
ACE2

\medterm{Angle-Closure Glaucoma} A type of glaucoma in which the outer edge of the iris blocks the eye’s drainage angle. Fluid then builds up, causing eye pressure to rise and potentially damaging the optic nerve and vision. It may develop suddenly with severe eye pain, blurred vision, halos, headache, nausea, or vomiting, or develop slowly without symptoms.

\medterm{Angular Cheilitis} Inflammation, soreness, and cracking at one or both corners of the mouth. Moisture and saliva, yeast or bacterial infection, dentures, irritation, or iron and vitamin deficiencies can contribute.

\medterm{Anhedonia} Reduced or absent ability to feel pleasure or interest in activities that were previously enjoyable. It can occur with depression, schizophrenia, substance use, and other conditions.

\medterm{Anhidrosis} Absent or greatly reduced sweating. It may result from nerve damage, skin disease, dehydration, medicines, or an inherited condition. Without enough sweat, the body may overheat dangerously, especially during exercise or hot weather.

\medterm{Animal-Assisted Therapy} A planned healthcare intervention in which a trained animal is used with a qualified professional to support goals such as improving mood, communication, movement, or participation. Safety, hygiene, and allergies must be considered.

\medterm{Anion} A particle with a negative electrical charge. Important anions in body fluids include chloride, bicarbonate, and phosphate, which help maintain fluid balance, electrical activity, and the blood’s acid level.

\medterm{Anion Gap} A value calculated from several blood salts to estimate acids that were not directly measured. A high result can suggest excess acid in the blood, but it must be interpreted with albumin, other tests, symptoms, and the laboratory’s method.

\medterm{Anion Gap Metabolic Acidosis} A subtype of metabolic acidosis characterized by an elevated serum anion gap (exceeding 12 mEq/L), signifying the accumulation of unmeasured organic or inorganic fixed acids in plasma, such as ketoacids in diabetic ketoacidosis, lactic acid, or exogenous toxic metabolites.

\synonyms
AGMA; High Anion Gap Metabolic Acidosis

\medterm{Aniridia} A congenital, bilateral, complete or partial absence of the iris, typically caused by mutations in the PAX6 gene. It is accompanied by photophobia, nystagmus, reduced visual acuity, foveal hypoplasia, and an increased risk of early-onset glaucoma and cataracts, and may occur as part of the WAGR syndrome.

\medterm{Anisocoria} An unequal size of the pupils, the black centers of the eyes. In about 20\% of healthy people, a slight difference is natural and harmless. When abnormal, it is caused by eye injury, certain eye drops, or nerve disorders that prevent a pupil from widening or shrinking normally. Doctors test how the pupils react in bright light and darkness to identify the underlying cause.

\medterm{Anisocytosis} Greater-than-usual variation in the size of red blood cells. It is a laboratory finding rather than a disease and may occur with iron deficiency, vitamin deficiency, or other types of anemia.

\medterm{Anisometropia} A difference in the focusing power of the two eyes. It can cause unequal image size, blurred vision, eyestrain, or reduced development of vision in one eye during childhood.

\medterm{Ankle} The joint connecting the lower leg and foot. It is formed mainly by the tibia, fibula, and talus and allows the foot to move up and down while ligaments provide stability.

\medterm{Ankle Arthroplasty} Surgical replacement of the ankle joint with an artificial joint, usually for severe painful arthritis. It aims to reduce pain while preserving movement; suitability depends on bone, alignment, activity, and surrounding-joint condition.

\medterm{Ankle Arthroscopy} A procedure using a small camera and instruments inserted through small cuts around the ankle to diagnose or treat problems such as loose bone or cartilage, inflamed lining, or tissue impingement.

\medterm{Ankle Bone} Any of the bones forming the ankle joint, chiefly the lower ends of the tibia and fibula and the upper part of the talus.

\medterm{Ankle Dislocation} Complete displacement of the ankle bones from their normal joint position, usually after a severe injury. It often occurs with a fracture or ligament injury and can damage blood vessels or nerves, so it requires urgent reduction and imaging.

\synonyms
Dislocated Ankle; Talotibial Dislocation

\medterm{Ankle Fracture} A break in one or more bones forming the ankle, usually the tibia, fibula, or talus. It may cause pain, swelling, bruising, and inability to bear weight; displaced or unstable fractures need prompt assessment.

\medterm{Ankle Fractures} Breaks in one or more bones forming the ankle joint. They may also injure ligaments, blood vessels, or nerves; the seriousness depends on bone alignment, joint stability, and whether the skin is broken.

\medterm{Ankle Joint} The hinge-like joint between the lower ends of the tibia and fibula and the talus. It mainly moves the foot upward and downward and is supported by surrounding ligaments.

\medterm{Ankle Pain} Pain in or around the ankle, commonly caused by a sprain, fracture, tendon problem, arthritis, bursitis, or infection. Deformity, marked swelling, fever, numbness, or inability to bear weight requires prompt assessment.

\medterm{Ankle Replacement} Surgery that replaces the damaged surfaces of the ankle joint with an artificial joint, usually for severe arthritis causing persistent pain and disability. It aims to reduce pain and preserve ankle movement.

\medterm{Ankle Sprain} An injury in which one or more ankle ligaments are stretched or torn, most often on the outer side after the foot twists inward. Pain, swelling, bruising, and a feeling of instability may occur.

\medterm{Ankle Swelling} Enlargement around the ankle caused by injury, inflammation, infection, vein or lymph problems, or general fluid retention. Sudden one-sided swelling with pain, warmth, fever, or breathlessness needs urgent assessment.

\medterm{Ankle-Brachial Index} A comparison of blood pressure measured at the ankle with blood pressure measured at the arm. A low result may indicate narrowed leg arteries; a very high result may mean the arteries are too stiff to measure reliably.

\medterm{Ankle-Foot Orthosis} A brace that supports the ankle and foot and controls alignment or movement during standing and walking. It may improve walking or prevent foot drop caused by weakness, spasticity, instability, or deformity.

\medterm{Ankyloblepharon} Abnormal sticking or joining of the upper and lower eyelid edges. It may be present at birth or follow inflammation, burns, injury, or surgery and can partly block vision when severe.

\medterm{Ankyloglossia} Alternate name; see \textit{Tongue-Tie}.

\medterm{Ankylose} To become stiff or unable to move at a joint, sometimes because scar tissue or bone joins the joint surfaces.

\medterm{Ankylosing} Causing or involving progressive stiffness or loss of movement in a joint, sometimes because scar tissue or bone joins the joint surfaces.

\medterm{Ankylosing Spondylitis} A long-term inflammatory arthritis that mainly affects the sacroiliac joints, where the spine joins the pelvis, and the joints and ligaments of the spine. It usually causes lower-back or hip pain and stiffness that is worse after rest or inactivity and may improve with movement. Inflammation can reduce spinal movement and, in severe cases, cause the vertebrae to fuse, making the spine rigid. The hips, ribs, knees, or feet may also be affected. Some people develop inflammation of the eye (uveitis), psoriasis, or inflammatory bowel disease. The risk is higher in people with certain genetic factors, including HLA-B27, but having HLA-B27 does not by itself cause the disease.

\synonyms
Bekhterev Disease; Axial Spondyloarthritis; Marie-Strümpell Disease

\medterm{Ankylosis} Abnormal reduction or complete loss of movement in a joint. It may result from fibrous scar tissue, bone fusion, injury, infection, chronic inflammation, or surgery. Fibrous ankylosis involves scar tissue, while bony ankylosis involves fusion of neighboring bones. In dentistry, ankylosis may describe fusion of a tooth root to the surrounding jawbone.

\medterm{Ankyrin} A protein that connects selected cell-surface proteins to the cell’s internal support framework. It helps red blood cells keep their shape; inherited changes can make the cells fragile and cause anemia.

\medterm{Annihilation} The conversion of a particle and its opposite partner into energy when they collide. In PET scanning, a positron meets an electron and produces two gamma rays that travel in opposite directions to create the image.

\medterm{Annular Pancreas} A birth defect in which pancreatic tissue forms a ring around part of the duodenum, the first section of the small intestine. The ring may narrow the intestine and cause vomiting, abdominal swelling, or feeding problems.

\medterm{Annulus} A ring-shaped structure or band in the body. Examples include the tough outer ring of an intervertebral disc and the fibrous rings surrounding the heart valves.

\medterm{Annulus Fibrosus} The tough, outer fibrocartilaginous ring of an intervertebral disc composed of concentric lamellae of collagen fibers, which encircles and restrains the central gel-like nucleus pulposus and absorbs rotational and compressive mechanical loads.

\synonyms
Anulus Fibrosus

\medterm{ANOCA} Chest pain or other angina symptoms when the major heart arteries do not have a severe blockage. The symptoms may result from poor function of the heart’s small vessels or temporary spasm of a coronary artery.

\medterm{Anodontia} Complete absence of teeth from birth, caused by failure of tooth development. It is rare and may occur with inherited conditions affecting the skin, hair, nails, and teeth.

\medterm{Anogenital Tract} The connected region containing the anus, tissues around the anus, genital organs, and nearby skin and moist linings.

\medterm{Anomalous Anatomy} A difference from the usual structure, number, position, origin, or pathway of an organ, blood vessel, duct, or other body part. It may cause no symptoms or may affect health, diagnosis, or surgery.

\medterm{Anomalous Coronary Artery} A birth defect in which one of the heart’s arteries has an unusual starting point or path. Some arteries pass between the aorta and the pulmonary artery and may be squeezed during hard exercise. This can reduce blood flow to the heart and may cause chest pain, fainting, abnormal heart rhythms, or rarely sudden cardiac death.

\synonyms
AAOCA; Anomalous Aortic Origin of a Coronary Artery; Congenital Coronary Artery Anomaly; ACAOS

\medterm{Anomaly} A difference from the usual structure, development, function, or measurement of the body. It may be harmless, or it may cause symptoms or affect health.

\medterm{Anomaly Scan} A detailed ultrasound examination usually performed around 18–22 weeks of pregnancy to assess fetal growth and organs, the placenta, and amniotic fluid, and to screen for major structural birth defects.

\synonyms
Mid-Pregnancy Ultrasound; 20-Week Scan; Fetal Anomaly Scan

\medterm{Anomia} Difficulty finding or saying a familiar word, especially the name of an object or person, even though its meaning is known. It can occur after a stroke or with other brain or language disorders.

\medterm{Anophthalmia} A birth defect in which one or both eyes fail to develop. It may occur with other birth defects and can affect growth of the eye socket and face.

\medterm{Anorchia} Complete absence of testicular tissue, present from birth or caused by loss of testes that had formed. It causes very low or absent testosterone production and infertility unless appropriate specialist care is provided.

\medterm{Anorectal Abscess} A collection of pus in tissues around the anus or rectum, usually caused by an infected anal gland. It may cause severe pain, swelling, fever, or difficulty sitting and often needs prompt drainage.

\medterm{Anorectal Fistula} An abnormal tunnel between the anus or rectum and nearby skin or another organ. It often develops after an abscess or infection and may cause persistent drainage, pain, or repeated abscesses.

\synonyms
Anal Fistula

\medterm{Anorectal Infection} Infection of the anus, rectum, or nearby tissues. It may cause pain, swelling, discharge, fever, or difficulty passing stool and can lead to an abscess that requires drainage.

\medterm{Anorectal Malformation} A birth defect in which the anus or rectum does not form normally or does not connect correctly to the skin. The opening may be absent, misplaced, or abnormally connected to another organ.

\medterm{Anorectal Manometry} A test that measures pressure and sensation in the rectum and anal canal. A thin pressure-sensitive tube assesses how well the muscles relax and tighten and may help investigate constipation or loss of bowel control.

\medterm{Anorexia} Loss or reduction of appetite. It can result from illness, pain, medicines, emotional distress, or other causes and is different from anorexia nervosa, an eating disorder.

\medterm{Anorexia Nervosa} A serious eating disorder involving persistent restriction of food, intense fear of gaining weight, and a distorted view of body size or shape. It can cause severe weight loss, malnutrition, slow heart rate, low blood pressure, and organ damage.

\synonyms
Anorexia; Restrictive Eating Disorder

\medterm{Anorgasmia} Persistent or repeated difficulty reaching orgasm or inability to reach orgasm despite adequate sexual stimulation. It may result from medicines, illness, hormonal changes, pain, stress, trauma, or relationship factors.

\medterm{Anorgasmy} Another spelling of anorgasmia: persistent or repeated inability to reach orgasm despite adequate sexual stimulation. Possible causes include medicines, illness, hormonal changes, stress, trauma, or relationship factors.

\medterm{Anoscopy} Examination of the anal canal using a short, lighted tube called an anoscope. It helps investigate rectal bleeding, haemorrhoids, fissures, fistulas, pain, and other problems near the anus.

\medterm{Anoscopy And High-Resolution Anoscopy} Examination of the anal canal and lower rectum with a short viewing tube. High-resolution anoscopy adds magnification and special lighting to identify and assess abnormal tissue, including changes related to HPV.

\medterm{Anosmia} Complete loss of the sense of smell. It may follow a viral infection, blocked or inflamed nasal passages, head injury, medicines, toxic exposure, or a neurological or inherited disorder.

\medterm{Anosodiaphoria} A neurological state of apparent emotional indifference or unconcern toward severe physical disability, such as hemiplegia, despite intellectual awareness and verbal acknowledgement of the deficit. It is most commonly associated with non-dominant (right) hemisphere parietal lobe lesions.

\medterm{Anosognosia} Lack of awareness of, or inability to recognize, one’s own illness or neurological impairment. It can occur after certain brain injuries or strokes and is not simply deliberate denial.

\medterm{Anotia} A birth defect in which the outer ear is completely absent. The ear canal and middle-ear structures may also be abnormal, causing reduced hearing on the affected side.

\medterm{Anovulation} Failure to release an egg from an ovary during a menstrual cycle. It can cause irregular or absent periods, abnormal bleeding, or infertility and may result from hormonal disorders, low body weight, stress, excessive exercise, or polycystic ovary syndrome.

\medterm{Anoxia} Complete lack of oxygen reaching body tissues. It can quickly damage or kill cells, especially in the brain, heart, and kidneys, and may result from blocked blood flow, breathing failure, or severe poisoning.

\medterm{Anrep Effect} An intrinsic physiological autoregulation mechanism of the myocardium whereby an abrupt increase in ventricular afterload causes an initial stretch followed by a sustained, positive inotropic increase in ventricular contractility over the ensuing several minutes.

\medterm{Anserine Bursa} A synovial fluid-filled sac situated on the proximal medial aspect of the tibia, interposed between the tibial collateral ligament and the conjoined tendons of the sartorius, gracilis, and semitendinosus muscles; inflammation causes pes anserine bursitis.

\synonyms
Pes Anserinus Bursa

\medterm{Antacids} Medicines that neutralise stomach acid and provide short-term relief from heartburn, indigestion, and acid reflux. Different products contain different minerals and may interact with other medicines or alter bowel movements.

\medterm{Antagonist} A drug or other substance that reduces or blocks the effect of another substance. A receptor antagonist binds to a receptor without activating it and reduces the response to an agonist, often by preventing the agonist from binding or signaling.

\medterm{Antagonist Muscle} A muscle that produces the opposite movement to the main working muscle and helps control or reverse that movement.

\medterm{Antalgic Gait} An abnormal, compensatory walking pattern adopted to avoid or minimize pain in a weight-bearing lower limb or spinal structure. It is characterized by a significantly shortened stance phase on the affected, painful limb and a rapid, hurried transition to the contralateral unaffected limb, frequently accompanied by a compensatory trunk lean toward the affected side to reduce joint contact forces.
\synonyms{Pain-Relieving Gait; Limping Gait; Antalgic Walk}

\medterm{Antebrachium} The anatomical region of the upper limb situated between the elbow and the wrist, consisting of the radius and ulna surrounded by anterior and posterior muscular compartments.

\synonyms
Forearm

\medterm{Antecubital Fossa} The shallow hollow at the front of the elbow. It contains the biceps tendon, a major artery, and a major nerve and is commonly used to draw blood or measure blood pressure.

\medterm{Antegrade} Moving or occurring in the normal forward direction, such as blood flowing away from the heart or an electrical signal travelling through its usual pathway.

\medterm{Antemortem Injury} An injury that occurred before death. It may show bleeding into living tissue, clot formation, inflammation, swelling, or early healing, which can distinguish it from damage occurring after death.

\medterm{Antenatal} Relating to the period before birth or to care during pregnancy. Antenatal care monitors the pregnant person and fetus, screens for complications, and provides health advice and support.

\synonyms
Prenatal; Antepartum

\medterm{Antepartum Fetal Surveillance} Tests used before birth to check fetal well-being, especially in higher-risk pregnancies. They may include monitoring fetal movements, the fetal heart rate, ultrasound assessment, and blood-flow measurements.

\medterm{Antepartum Hemorrhage} Bleeding from the vagina during the second half of pregnancy and before birth. Causes include placenta previa, placental abruption, and problems with the cervix or vagina; any significant bleeding needs urgent assessment.

\medterm{Anterior} Relating to the front of the body or positioned nearer the front than another structure. The term is used to describe body parts, scan findings, injuries, and disease locations.

\medterm{Anterior Cerebral Artery Syndrome} A stroke syndrome resulting from occlusion or ischemia of the anterior cerebral artery, classically causing contralateral motor weakness and sensory loss predominantly affecting the lower extremity more than the upper extremity and face. Additional findings may include abulia, urinary incontinence, and frontal release signs due to medial frontal lobe involvement.

\synonyms
ACA Syndrome; ACA Infarction

\medterm{Anterior Chamber} The fluid-filled space between the cornea and iris at the front of the eye. It contains aqueous fluid, which nourishes the cornea and lens and helps maintain eye pressure; impaired drainage can contribute to glaucoma.

\medterm{Anterior Chamber Angle} The area where the cornea and iris meet at the front of the eye. It contains the main drainage tissue for aqueous fluid; blockage here can raise eye pressure and cause glaucoma.

\medterm{Anterior Commissure} A transverse bundle of myelinated nerve fibers crossing the midline anterior to the columns of the fornix and above the lamina terminalis, interconnecting the olfactory bulbs, amygdalae, and middle and inferior temporal gyri across the cerebral hemispheres.

\medterm{Anterior Cord Syndrome} A spinal-cord injury that weakens movement and reduces pain and temperature sensation below the injury, while vibration and awareness of limb position may be relatively preserved.

\medterm{Anterior Cruciate Ligament} A strong band of tissue inside the knee that connects the thigh bone to the shin bone. It helps prevent the shin bone from sliding too far forward and controls knee rotation.

\medterm{Anterior Cruciate Ligament Injury} Damage or tearing of the main stabilizing ligament inside the knee, often after sudden turning, stopping, landing, or overextension. It may cause a popping sensation, rapid swelling, pain, and a feeling that the knee gives way.

\medterm{Anterior Cruciate Ligament Tear} A partial or complete tear of the ligament inside the knee that limits forward sliding and rotation of the shin bone. It commonly causes a pop, rapid swelling, pain, and knee instability after a sudden twist or landing.

\medterm{Anterior Drawer Sign} A physical check for a torn knee ligament (anterior cruciate ligament). With the knee bent, the doctor gently pulls the shin forward; moving too far forward means the ligament is likely torn or stretched.

\medterm{Anterior Drawer Test} A clinical orthopedic physical examination maneuver used to assess the integrity of the anterior cruciate ligament (ACL) of the knee or the anterior talofibular ligament of the ankle. Excessive anterior translation of the tibia relative to the femur with a soft or mushy endpoint indicates an ACL tear.

\synonyms
Anterior Drawer Sign

\medterm{Anterior Interosseous Nerve} A pure motor branch of the median nerve arising in the proximal forearm that courses deeply along the anterior surface of the interosseous membrane to innervate the flexor pollicis longus, the lateral half of flexor digitorum profundus, and the pronator quadratus; damage produces the classic inability to make the "OK" sign.

\synonyms
AIN

\medterm{Anterior Interosseous Nerve Syndrome} A pure motor neuropathy of the forearm resulting from compression or neuritis of the anterior interosseous nerve (a motor branch of the median nerve). Causes weakness of the flexor pollicis longus, pronator quadratus, and the radial half of the flexor digitorum profundus, classically producing an inability to form the \"OK\" sign (pinching index finger and thumb pulp-to-pulp instead of tip-to-tip).

\synonyms
Kiloh-Nevin Syndrome; AIN Neuropathy

\medterm{Anterior Knee Pain} Pain at the front of or behind the kneecap, often worse with stairs, squatting, running, or prolonged sitting. Common causes include irritation around the kneecap, tendon problems, arthritis, and injury.

\medterm{Anterior Mediastinum} The front part of the lower mediastinum, located behind the sternum and in front of the heart sac. It contains fat, connective tissue, lymph nodes, and part of the thymus.

\medterm{Anterior Midline} The imaginary vertical line running down the center of the front of the body, separating the right and left sides.

\medterm{Anterior Myocardial Infarction} A heart attack affecting the front wall of the main pumping chamber, usually because the left anterior descending coronary artery is blocked. It can seriously weaken the heart and requires emergency treatment.

\synonyms
Anterior Heart Attack; Anterior MI

\medterm{Anterior Pituitary} The front part of the pituitary gland. It releases hormones that control the adrenal and thyroid glands, growth, milk production, egg or sperm production, and reproductive function.

\medterm{Anterior Talofibular Ligament} A short, flat ligament on the lateral aspect of the ankle that spans from the anterior margin of the lateral malleolus of the fibula to the neck of the talus. It is the primary stabilizer against anterior talar displacement and excessive inversion, and represents the most frequently torn ligament in lateral ankle sprains.

\synonyms
ATFL

\medterm{Anterior Tibial Artery} An artery that carries blood down the front of the leg and continues onto the top of the foot as the dorsalis pedis artery. Its pulse can usually be felt near the front of the ankle.

\medterm{Anterior Uveitis} Inflammation in the front part of the eye, usually involving the iris and sometimes nearby tissue. It can cause eye pain, redness, light sensitivity, blurred vision, and floating spots.

\synonyms
Iritis; Iridocyclitis

\medterm{Anterior Vaginal Prolapse} Bulging of the anterior vaginal wall into the vaginal vault, commonly resulting from herniation of the urinary bladder (cystocele) due to attenuation or tearing of the pubocervical fascia and pelvic floor musculature. Symptoms may include a palpable vaginal bulge, pelvic heaviness, urinary frequency, urgency, incomplete emptying, or stress incontinence.

\synonyms
Cystocele; Bladder Prolapse; Prolapsed Bladder; Anterior Wall Prolapse

\medterm{Anterior Vaginal Wall Repair} Surgery to strengthen the front wall of the vagina when the bladder bulges into it. The procedure restores support using the person’s own tissues; the approach and risks depend on the prolapse and overall health.

\medterm{Anterior/Posterior Drawer Test} A physical examination of knee stability. Pulling the shin bone forward tests the anterior cruciate ligament; pushing it backward tests the posterior cruciate ligament.

\medterm{Anterograde Amnesia} A selective impairment of memory in which an individual is unable to form, encode, or consolidate new long-term declarative (episodic and semantic) memories after the onset of a neurological insult, while remote memories formed prior to the event remain largely intact. It typically results from bilateral damage to the medial temporal lobes, hippocampi, mammillary bodies, or diencephalic structures following traumatic brain injury, stroke, herpes simplex encephalitis, or Korsakoff syndrome.
\synonyms{Post-Traumatic Amnesia; Anterograde Memory Loss; Forward-Acting Amnesia}

\medterm{Anterolisthesis} Forward slipping of one spinal bone relative to the bone below it. It may result from wear, a stress fracture, injury, or a problem present from birth and may cause back pain or nerve symptoms.

\medterm{Anteroposterior} Directed from the front of the body toward the back. In an anteroposterior X-ray view, the beam enters through the front and leaves through the back.

\medterm{Anthocyanin} A natural plant pigment that gives many fruits, vegetables, and flowers red, purple, or blue colours. It has antioxidant effects in laboratory studies and is found in foods such as berries, grapes, and red cabbage.

\medterm{Anthracosis} A benign, non-fibrotic form of pneumoconiosis characterized by the asymptomatic accumulation of inhaled carbon and soot pigment in the lung parenchyma, alveolar macrophages, and draining hilar lymph nodes, commonly seen in urban dwellers and cigarette smokers.

\synonyms
Black Lung Pigmentation; Carbon Pigmentation

\medterm{Anthrax} A serious infection caused by the bacterium \textit{Bacillus anthracis}. People can become infected through contact with infected animals or animal products, breathing in spores, or eating contaminated food; skin, lung, or intestinal disease may result.

\medterm{Anthrax Blood Test} A laboratory test that looks for the anthrax bacterium or the body’s immune response to it in blood or another sample. It helps confirm suspected anthrax, which can affect the skin, lungs, or digestive tract.

\medterm{Anthropometry} The measurement of body size and proportions, such as height, weight, waist size, limb length, and skinfold thickness. These measurements help assess growth, nutrition, and body composition.

\medterm{Anti-Anxiety Agents} Medicines used to reduce excessive anxiety or panic. Different medicines work in different ways and may cause sleepiness, dependence, or other side effects; the choice depends on the person and the condition.

\synonyms
Anxiolytics

\medterm{Anti-CCP} A blood-test antibody marker associated with rheumatoid arthritis. A positive result supports the diagnosis when symptoms fit, but it does not prove that rheumatoid arthritis is present or predict its course by itself.

\medterm{Anti-CCP Antibody} An antibody measured in blood that targets proteins changed by a process called citrullination. It can support a diagnosis of rheumatoid arthritis, but results must be interpreted with symptoms and other tests.

\medterm{Anti-Citrullinated Peptide Antibody} Alternate index label; see \textit{Anti-Cyclic Citrullinated Peptide Antibody}.

\medterm{Anti-Cyclic Citrullinated Peptide Antibody} An autoantibody directed against citrulline-containing peptide antigens, possessing high diagnostic specificity (greater than 95 percent) for rheumatoid arthritis and correlating with aggressive, erosive joint disease.

\synonyms
Anti-CCP Antibody; ACPA; Anti-Citrullinated Peptide Antibody; Anti-Cyclic Citrullinated Peptide; Cyclic Citrullinated Peptide Antibody

\medterm{Anti-D} An injection containing antibodies given to an RhD-negative pregnant person after certain pregnancy events, or around birth, to prevent the person’s immune system making antibodies against RhD-positive fetal blood. This helps protect a current or future pregnancy from fetal red-cell destruction.

\medterm{Anti-DNase B Blood Test} A blood test that measures the body’s antibodies against a protein made by group A streptococcal bacteria. It can support evidence of a recent strep infection when doctors investigate rheumatic fever, kidney inflammation, or another complication.

\medterm{Anti-Double-Stranded DNA Antibody} An antinuclear autoantibody specifically directed against the native double-helical structure of deoxyribonucleic acid, highly specific for systemic lupus erythematosus and closely correlating with lupus nephritis disease activity and flares.

\synonyms
Anti-dsDNA Antibody; dsDNA Antibody; Anti-Double-Stranded DNA; Double-Stranded DNA Antibody

\medterm{Anti-GBM Disease} An autoimmune disease in which antibodies attack supporting layers in the kidneys and sometimes the lungs. It can rapidly cause kidney failure, coughing up blood, or bleeding into the lungs and requires urgent specialist care.

\medterm{Anti-Glomerular Basement Membrane Blood Test} A blood test that detects antibodies attacking the kidney’s filtering membranes and sometimes the lungs’ supporting membranes. A positive result supports anti-GBM disease when symptoms and other test results are consistent.

\medterm{Anti-Glomerular Basement Membrane Disease} An autoimmune disease in which antibodies damage the kidney’s filtering membranes and sometimes the lungs. It may cause rapidly worsening kidney function, blood in the urine, coughing up blood, or lung bleeding.

\medterm{Anti-HMGCR Myopathy} A rare autoimmune muscle disease in which antibodies attack a muscle enzyme. It causes ongoing weakness and muscle injury, especially around the hips, thighs, shoulders, and upper arms, and may begin after treatment with a statin medicine.
\synonyms
Statin-Induced Immune-Mediated Necrotizing Myopathy; Anti-HMG-CoA Reductase Myopathy

\medterm{Anti-IgE Therapy} Treatment using an antibody that binds free immunoglobulin E, an allergy-related immune protein. This reduces activation of allergy cells and may lessen symptoms and attacks in selected allergic asthma and other allergic conditions.

\medterm{Anti-IL4 Receptor Therapy} Treatment that blocks a cell receptor used by interleukin-4 and interleukin-13, immune signals that promote allergic inflammation. It can reduce airway inflammation in selected asthma and other allergic diseases.

\medterm{Anti-IL5 Therapy} Treatment that blocks interleukin-5 or its receptor, reducing the production or survival of eosinophils, a type of white blood cell. It is used for selected eosinophilic diseases, including some severe asthma.

\medterm{Anti-Inflammatory} Pertaining to any drug, substance, or therapeutic intervention that reduces, counteracts, or suppresses inflammation, vascular permeability, tissue edema, and associated pain. Major therapeutic classes include nonsteroidal anti-inflammatory drugs and glucocorticoids.

\medterm{Anti-Insulin Antibody Test} A blood test that detects antibodies against insulin. These antibodies can bind insulin and release it unpredictably, causing low blood sugar; the test helps investigate unexplained hypoglycaemia and insulin autoimmune syndrome.

\medterm{Anti-MDA5 Dermatomyositis} A form of dermatomyositis associated with antibodies against a protein called MDA5. It often causes distinctive skin changes with little muscle weakness but carries a high risk of rapidly worsening inflammation and scarring in the lungs.

\medterm{Anti-Mullerian Hormone} A hormone made by developing ovarian follicles and by cells in the testes before birth. Its blood level can help estimate ovarian follicle number or assess some reproductive and developmental conditions, but it does not measure fertility by itself.

\medterm{Anti-NMDA Receptor Encephalitis} A serious immune attack on the brain involving a receptor important for communication between nerve cells. It may cause sudden behavior or psychiatric changes, memory loss, seizures, abnormal movements, reduced consciousness, and unstable heart rate or blood pressure. Urgent hospital treatment is needed.

\synonyms
Anti-N-Methyl-D-Aspartate Receptor Encephalitis; NMDA Receptor Encephalitis

\medterm{Anti-Reflux Surgery} An operation to reduce the backward flow of stomach contents into the oesophagus by strengthening or reshaping the valve between the oesophagus and stomach. It may be considered for selected people with severe reflux or a hiatus hernia.

\medterm{Anti-Reflux Surgery, Pediatric} Surgical reconstruction of the gastroesophageal junction—most commonly via laparoscopic Nissen fundoplication—in infants and children with medically refractory gastroesophageal reflux disease, recurrent aspiration, or severe failure to thrive.

\synonyms
Pediatric Fundoplication; Pediatric Anti-Reflux Surgery; Nissen Fundoplication in Children

\medterm{Anti-Rust Product Poisoning} Illness caused by swallowing, breathing, or skin exposure to a product used to remove or prevent rust. It may burn the mouth and digestive tract or damage the lungs and other organs; significant exposure requires urgent medical treatment.

\medterm{Anti-Smooth Muscle Antibody} An antibody that reacts with proteins in smooth muscle cells. A blood test may support autoimmune hepatitis, but low levels can also occur with other autoimmune diseases or infections and are not diagnostic alone.

\medterm{Anti-SRP Myopathy} A rare autoimmune muscle disease caused by antibodies against a protein involved in making new proteins. It causes rapidly progressive weakness, especially around the hips, shoulders, and neck, and may cause difficulty swallowing.

\synonyms
Anti-Signal Recognition Particle Myopathy; SRP-Associated Necrotizing Myopathy

\medterm{Anti-Thyroid Peroxidase Antibody} An autoantibody directed against thyroid peroxidase, the key enzyme required for thyroid hormone synthesis. Elevated titers are the hallmark laboratory finding in Hashimoto thyroiditis and are also present in Graves disease.

\synonyms
Anti-TPO Antibody; Thyroid Peroxidase Antibody

\medterm{Anti-VEGF Injections} Injections, usually given into the eye, that block vascular endothelial growth factor, a signal that promotes abnormal blood-vessel growth and leakage. They can improve or preserve vision in several retinal diseases.

\medterm{Antiandrogen} A medicine or substance that lowers the production or blocks the effects of androgen hormones such as testosterone. It is used for selected prostate conditions, androgen-related disorders, and some gender-affirming treatments.

\medterm{Antianginal} A medicine used to prevent or relieve angina, the chest discomfort caused by reduced blood flow to heart muscle. Different medicines lower the heart’s workload, widen blood vessels, or improve coronary blood flow.

\medterm{Antiarrhythmic} A medication or therapeutic intervention designed to prevent, suppress, or terminate abnormal cardiac rhythms and restore or maintain normal sinus rhythm. Pharmacological agents are conventionally categorized according to the Vaughan Williams classification into classes I through IV.

\medterm{Antibacterial} Able to kill bacteria or stop them from growing. The term may describe a medicine, disinfectant, surface treatment, or part of the body’s immune defence.

\medterm{Antibiogram} A report showing which antibiotics are likely to work against bacteria grown from patient samples in a laboratory. It helps clinicians choose treatment and monitor local patterns of antibiotic resistance.

\medterm{Antibiotic Misuse in Viral URIs} The inappropriate prescription or self-administration of antibacterial agents for acute viral upper respiratory tract infections (such as the common cold or influenza), which confers no therapeutic benefit and drives antimicrobial resistance and adverse drug events.

\synonyms
Antibiotic Stewardship in URI; Inappropriate Antibiotic Use in Colds; Antibiotic Overuse in Flu

\medterm{Antibiotic-Resistant Bacterial Infection} An infection caused by bacteria that are not adequately stopped by one or more antibiotics normally active against them. Resistance may be intrinsic or acquired and can limit effective treatment options.

\medterm{Antibiotics} Medicines that kill bacteria or stop them from multiplying. They treat susceptible bacterial infections but do not treat viral infections such as colds or influenza.

\medterm{Antibodies} Proteins made by immune cells that recognize and attach to specific foreign or abnormal substances. They can block germs or toxins, mark them for destruction, and activate other parts of the immune system.

\medterm{Antibody} A protein made by the immune system that attaches specifically to a foreign substance, germ, toxin, or abnormal cell and may block it or mark it for removal.

\medterm{Antibody Engineering} Designing or changing antibodies so they bind more accurately, last longer, cause fewer unwanted immune reactions, or work better as medicines. Researchers may alter their target, size, immune effects, or attachment to another medicine.

\medterm{Antibody Therapy} Treatment using specially prepared antibodies to block, activate, or target a particular molecule or cell. It can be used for cancer, autoimmune disease, infection, inflammation, and other conditions.

\synonyms
Immunotherapy

\medterm{Antibody Titer Blood Test} A blood test that estimates the amount of a specific antibody by repeatedly diluting the sample until the antibody can no longer be detected. Results may help show previous exposure, immunity, or a recent infection, depending on the antibody tested.

\medterm{Antibody-Dependent Cell-Mediated Cytotox.} An immune process in which antibodies coat an infected, abnormal, or cancer cell and attract killer immune cells. The killer cells then damage and destroy the marked cell.

\medterm{Antibody-Drug Conjugate} A treatment in which an antibody carrying a linked medicine binds to a target on a cell and delivers the medicine near or inside that cell. It is used for selected cancers and can still damage healthy tissues or cause medicine-specific toxicities.

\medterm{Anticholinergic} Blocking the action of acetylcholine, a chemical messenger used by nerves. These medicines can reduce secretions and muscle spasms but may cause dry mouth, blurred vision, constipation, difficulty urinating, fast heartbeat, or confusion.

\medterm{Anticholinergic Burden} The combined anticholinergic effect of one or more medicines. A high burden can cause dry mouth, constipation, blurred vision, difficulty urinating, confusion, or delirium, especially in older adults.

\medterm{Anticholinergic Toxicity} Excess blockade of muscarinic acetylcholine receptors, producing a hot, flushed, dry, tachycardic, mydriatic, delirious patient with urinary retention and reduced bowel sounds. Treatment is supportive with cooling and benzodiazepines; physostigmine is reserved for selected severe cases after toxicology assessment.

\medterm{Anticholinergic Toxidrome} A pattern of poisoning caused by excessive blockage of acetylcholine signals. It may cause dry skin and mouth, large pupils, fast heartbeat, inability to urinate, reduced bowel activity, overheating, agitation, or delirium.

\medterm{Anticoagulants and Thrombolytics in Pregnancy} Medicines that prevent clot formation or dissolve selected clots during pregnancy or after delivery. Choice and timing depend on maternal indication, placental transfer, bleeding risk, delivery plans, and whether neuraxial anesthesia may be used; specialist management is required.

\medterm{Anticoagulation} Treatment that reduces the formation or growth of blood clots by slowing the clotting process. It is used for conditions such as clots in veins, irregular heart rhythm, and selected heart or blood-vessel disorders.

\medterm{Anticodon} A group of three RNA building blocks on a transfer-RNA molecule. It pairs with a matching sequence on messenger RNA so the cell adds the correct amino acid while making a protein.

\medterm{Anticonvulsants} Medicines used to prevent or control seizures. Some are also used for nerve pain, migraine, or mood disorders; their effects, interactions, and side effects differ between medicines.

\medterm{Antideformity Position} A planned position for a limb or joint that helps prevent or limit shortening, stiffness, or deformity after injury, surgery, burns, nerve damage, or prolonged immobility.

\medterm{Antidepressants} Medicines used mainly to treat depression. Different types change the activity of chemical messengers in the brain and may also treat anxiety, nerve pain, obsessive thoughts, or other conditions.

\medterm{Antidiabetic} Relating to a medicine or treatment that lowers high blood sugar or helps control diabetes. Different treatments increase insulin action, reduce glucose production, or remove excess glucose from the body.

\medterm{Antidiuretic Hormone} A hormone made in the brain and released by the pituitary gland that helps the kidneys conserve water. It also narrows blood vessels when blood volume or blood pressure is low. It is also called vasopressin.

\medterm{Antidiuretic Hormone Blood Test} A blood test that measures vasopressin, the hormone that helps the kidneys conserve water. It may help investigate excessive thirst, very dilute urine, inappropriate water retention, or other disorders of water balance.

\medterm{Antidiuretic Hormone Physiology} Vasopressin is made in the hypothalamus and released from the posterior pituitary when the blood becomes too concentrated or blood volume falls. It makes kidney tubules retain water and can narrow blood vessels.

\medterm{Antidiuretic Hormone Resistance} Failure of the kidneys to respond properly to vasopressin. It causes production of large amounts of dilute urine and intense thirst and may result from medicines, mineral imbalance, kidney disease, or inherited changes.

\medterm{Antidote} A substance that counteracts or limits the harmful effects of a poison. The correct antidote depends on the substance involved; some antidotes block the poison, remove it, or restore a chemical process disrupted by it.

\medterm{Antiemetics} Medicines used to prevent or relieve nausea and vomiting. Different types act on chemical signals in the brain or digestive tract, and the choice depends on the cause, age, pregnancy, other medicines, and risk of side effects.

\medterm{Antiepileptic Drugs} Medicines used to prevent or control seizures by reducing excessive electrical activity in the brain. The best choice depends on the seizure type, age, other conditions, medicines, and possible side effects.

\medterm{Antifreeze Poisoning} Poisoning after swallowing or otherwise being exposed to antifreeze, especially products containing ethylene glycol. It may first cause intoxication and vomiting, then severe acid buildup, kidney failure, brain injury, or death; severe poisoning is an emergency requiring prompt medical treatment.

\medterm{Antifungal Prophylaxis} Use of an antifungal medicine to prevent a serious fungal infection in someone at high risk, such as a person with prolonged low white-cell counts, a stem-cell transplant, or severe immune suppression.

\medterm{Antifungal Resistance} The ability of a fungus to survive a medicine that would normally stop its growth or kill it. Resistance may be present naturally or develop after genetic changes or repeated exposure to antifungal medicines.

\medterm{Antigen} A substance or part of a substance that the immune system can recognize. It may come from a germ, toxin, transplanted tissue, or abnormal cell and can trigger antibodies or other immune responses.

\medterm{Antigen Detection Test} A test that looks for a protein or other marker from a specific germ in a patient sample. It is often rapid, but it may miss an infection when the amount of germ material is small.

\medterm{Antigen Presentation} The process by which an immune or other cell displays small pieces of proteins on its surface for T cells to inspect. This helps the immune system recognize infections, abnormal cells, transplanted tissue, and sometimes the body’s own tissues.

\medterm{Antigen Receptor} A protein on an immune cell that recognizes a particular substance or protein fragment. B-cell receptors can bind whole substances, while T-cell receptors recognize fragments displayed by other cells.

\medterm{Antigen-Presenting Cell} An immune cell that breaks proteins into small pieces and displays them on its surface so T cells can recognize them. Important antigen-presenting cells include dendritic cells, macrophages, and B cells.

\synonyms
APC

\medterm{Antigenic Drift} The gradual change in a virus’s surface proteins as it multiplies. These small changes can help viruses, especially influenza viruses, partly evade existing immunity and contribute to repeated seasonal outbreaks.

\medterm{Antigenic Shift} A sudden major change in a virus’s surface proteins, usually when two different influenza A viruses exchange genetic material inside the same cell. It can produce a virus to which many people have little immunity.

\medterm{Antigenic Variation} Changes in a germ’s surface markers that help it avoid recognition by the immune system. Variation may happen gradually through mutations or suddenly when genetic material is exchanged between different strains.

\medterm{Antigout} A medicine used to prevent or treat gout. Some lower uric-acid levels over time, while others reduce the inflammation and pain of an acute gout attack.

\medterm{Antihistamines} Medicines that block the effects of histamine. They can relieve itching, sneezing, runny nose, hives, or other allergy symptoms; some cause drowsiness, while others act mainly on stomach acid.

\medterm{Antihistamines For Allergies} Medicines that block histamine and relieve allergy symptoms such as itching, sneezing, runny nose, watery eyes, or hives. Some cause drowsiness and interactions or side effects vary between medicines.

\medterm{Antihypertensives In Pregnancy} Medicines used to lower high blood pressure during pregnancy. The choice depends on blood-pressure severity and pregnancy stage; some commonly used blood-pressure medicines are avoided because they can harm the fetus.

\medterm{Antileukotriene} A medicine that blocks leukotrienes, chemical signals that can cause airway narrowing and inflammation. It is used in selected asthma and allergy-related conditions, often as part of a broader treatment plan.

\synonyms
Leukotriene Modifier; Leukotriene Pathway Inhibitor

\medterm{Antimetabolite} A substance that resembles a chemical needed by cells and interferes with the cell’s normal chemical reactions or DNA production. Some are used to treat cancer or suppress an overactive immune system.

\medterm{Antimicrobial} A substance that kills or limits the growth of microscopic organisms, including bacteria, viruses, fungi, or parasites. Antibiotics are antimicrobials, but they act only against bacteria.

\medterm{Antimicrobial Adverse Effects} Unwanted effects caused by medicines that treat infections. Depending on the medicine, they may include allergy, diarrhoea, liver or kidney injury, hearing problems, tendon injury, or harmful effects on heart rhythm.

\medterm{Antimicrobial Agent} A substance that kills microorganisms or stops them from growing. Antimicrobial agents include antibacterial, antiviral, antifungal, and antiparasitic medicines.

\medterm{Antimicrobial Desensitization} A carefully supervised process of giving gradually increasing doses of an infection-fighting medicine to temporarily reduce an immediate allergy response. It is considered only when the medicine is essential and suitable alternatives are unavailable.

\medterm{Antimicrobial Dosing Optimization} Choosing the infection medicine, dose, timing, route, and duration so enough medicine reaches the infection without unnecessary harm. Kidney function, liver function, body size, infection severity, and resistance are considered.

\medterm{Antimicrobial Hypersensitivity Reactions} Immune reactions to an infection-fighting medicine, ranging from an itchy rash to life-threatening anaphylaxis. The timing and symptoms vary with the medicine and type of immune response.

\medterm{Antimicrobial Pharmacodynamics} How an infection-fighting medicine affects a microorganism, including how much medicine and how much exposure are needed to stop or kill it. This information helps select an effective dose and schedule.

\medterm{Antimicrobial Pharmacokinetics} How the body absorbs, spreads, changes, and removes an infection-fighting medicine. Kidney or liver disease, age, body size, and other medicines can change how much medicine is needed and how often it is given.

\medterm{Antimicrobial Resistance} The ability of a microorganism to survive or grow despite a medicine that would normally kill it or stop its growth. Resistance can affect bacteria, viruses, fungi, or parasites and can spread between organisms.

\synonyms
AMR; Drug Resistance

\medterm{Antimicrobial Resistance Mechanisms} Ways microorganisms avoid being killed by infection medicines. They may destroy or change the medicine, alter its target, prevent it entering, or actively pump it out of the cell.

\medterm{Antimicrobial Resistance Surveillance} Ongoing collection and analysis of resistance results from patient samples. It shows how resistance changes over time and helps hospitals and public-health services choose effective treatment and limit spread.

\medterm{Antimicrobial Sensitivity Test} A laboratory test that shows which infection-fighting medicines can stop or kill a microorganism grown from a patient sample. It helps guide treatment for a confirmed infection.

\medterm{Antimicrobial Stewardship} Coordinated efforts to use infection-fighting medicines only when needed and to choose the right medicine, dose, route, and duration. This improves outcomes and limits resistance and side effects.

\medterm{Antimicrobial Susceptibility Testing} Laboratory testing that estimates whether a microorganism is likely to respond to particular infection-fighting medicines. Methods include measuring growth around medicine-containing discs or in liquid; results help select treatment.

\medterm{Antimicrobial Therapy} Use of antibacterial, antiviral, antifungal, or antiparasitic medicines to prevent or treat an infection. The choice depends on the microorganism, infection site, severity, resistance, and the person’s health.

\medterm{Antimitochondrial Antibody} An antibody that reacts with structures inside mitochondria, the parts of cells that produce energy. It is strongly associated with primary biliary cholangitis, an autoimmune bile-duct disease, but the result must be interpreted with symptoms and liver tests.

\medterm{Antimony} A metallic chemical element whose compounds can be poisonous. Some antimony compounds have been used to treat leishmaniasis; exposure can affect the heart, liver, kidneys, or nervous system.

\medterm{Antineoplastic} Inhibiting or halting the proliferation, maturation, or spread of malignant neoplastic cells. The term designates chemotherapy agents, targeted molecular therapies, immunotherapies, and hormonal agents used in cancer therapy.

\synonyms
Anticancer Agent; Antitumor Agent

\medterm{Antinuclear Antibody} A diverse family of autoantibodies directed against nucleic acids, histones, ribonucleoproteins, and other nuclear and cytoplasmic components, evaluated via indirect immunofluorescence as a primary serological screen for systemic lupus erythematosus, systemic sclerosis, mixed connective tissue disease, and Sjogren syndrome.

\synonyms
ANA; Antinuclear Antibodies; Antinuclear Autoantibodies; FANA

\medterm{Antinuclear Antibody Panel} A comprehensive laboratory panel of specific autoantibody assays (such as anti-dsDNA, anti-Smith, anti-Ro/SSA, anti-La/SSB, and anti-RNP) utilized to characterize specific autoimmune rheumatic diseases following a positive antinuclear antibody screen.

\medterm{Antioxidant} A substance that helps neutralize or limit damage from reactive molecules, including free radicals and some reactive oxygen or nitrogen species. The body makes antioxidant enzymes and obtains other antioxidants from food; high-dose supplements are not automatically beneficial and may be harmful or interact with treatment.

\medterm{Antiparietal Cell Antibody Test} A blood test that detects antibodies against cells in the stomach lining that make acid and help absorb vitamin B12. A positive result may support autoimmune gastritis or pernicious anemia but does not prove either condition alone.

\medterm{Antiparkinson} Relating to medicines used to reduce symptoms of Parkinson disease. These medicines may increase or imitate dopamine activity or change other brain signals that influence movement.

\medterm{Antiphospholipid Antibodies} Autoantibodies that attach to proteins linked with cell membranes. Persistent antibodies may increase the risk of blood clots and pregnancy complications, but they can also appear temporarily after an infection or with some medicines.

\medterm{Antiphospholipid Antibody} An antibody directed against proteins that bind to cell-membrane fats. Persistent antibodies may be associated with blood clots or pregnancy complications, especially when clinical features of antiphospholipid syndrome are present.

\synonyms
APA; aPL Antibody

\medterm{Antiphospholipid Antibody Syndrome} Alternate name; see \textit{Antiphospholipid Syndrome}.

\medterm{Antiphospholipid Syndrome} A systemic autoimmune disorder characterized by recurrent arterial, venous, or microvascular thrombosis and obstetric complications (such as recurrent miscarriages or late fetal loss), occurring in the presence of persistent antiphospholipid antibodies.

\synonyms
APS; Antiphospholipid Antibody Syndrome; Hughes Syndrome; aPL Syndrome

\medterm{Antiphospholipid Syndrome and Pregnancy} An autoimmune clotting disorder associated with recurrent pregnancy loss, placental insufficiency, fetal growth restriction, preeclampsia, thrombosis, or pregnancy morbidity. Diagnosis requires clinical events plus persistent laboratory antiphospholipid antibodies; pregnancy care is specialist-led.

\medterm{Antiplatelet} A medicine that reduces platelet activation or aggregation and lowers the formation of arterial blood clots; see \textit{Antiplatelet Agents}.

\medterm{Antiplatelet Agents} Medicines that reduce the ability of platelets to stick together and form clots. They are used in selected people to lower the risk of clots in arteries, including after some heart or blood-vessel events.

\medterm{Antiplatelet Agents, P2Y12 Inhibitors} A subclass of oral antiplatelet medications—including clopidogrel, ticagrelor, and prasugrel—that selectively block the platelet P2Y12 adenosine diphosphate receptor, preventing ADP-mediated platelet activation and arterial thrombosis.

\synonyms
P2Y12 Inhibitors; P2Y12 Receptor Antagonists; ADP Receptor Antagonists

\medterm{Antiplatelet Therapy} Treatment with medicines that make platelets less likely to stick together. It is used in selected people to reduce clots in arteries, including after some heart attacks, strokes, or stent procedures.

\medterm{Antiprotozoal} A medicine or substance that kills protozoa or stops them growing. Protozoa are microscopic parasites that can cause infections such as malaria, giardiasis, or amoebiasis.

\medterm{Antipsychotic Medicine} A medicine used to reduce psychosis, including hallucinations, delusions, severe disorganization, or agitation. Antipsychotics differ in their effects and can cause movement problems, sleepiness, metabolic changes, or other adverse effects.

\medterm{Antipyretic} A substance or drug that lowers an abnormally elevated body temperature in fever by acting upon the hypothalamic thermoregulatory center to reset the temperature set point back toward normal. Common clinical examples include acetaminophen, ibuprofen, and aspirin.

\medterm{Antiretroviral} A medicine that acts against retroviruses, including HIV. Different antiretroviral medicines block different steps by which HIV enters cells, copies its genetic material, or makes new virus particles.

\medterm{Antiretroviral Therapy} Treatment of HIV infection with a combination of medicines that block several steps in HIV reproduction. It lowers the amount of virus, protects immune function, reduces illness, and greatly lowers the chance of sexual transmission when the virus remains undetectable.

\medterm{Antisepsis} Use of a chemical product on living skin or tissue to reduce harmful microorganisms and lower infection risk. It does not make tissue completely sterile and must be used in a way that avoids injury.

\medterm{Antiseptic} A substance used on skin or other living tissue to kill or slow harmful microorganisms. Antiseptics can help reduce infection risk but may damage tissue if used incorrectly.

\medterm{Antiserum} Blood serum containing antibodies against a specific germ, toxin, or other substance. It can provide temporary immune protection or help neutralize a toxin.

\medterm{Antisocial Personality} A long-term pattern of disregarding other people’s rights, safety, rules, and responsibilities, often including deceit, impulsive behaviour, aggression, or lack of remorse.

\synonyms
Antisocial Personality Disorder

\medterm{Antisocial Personality Disorder} A mental-health disorder involving a persistent pattern of violating other people’s rights and social rules, deceitfulness, impulsivity, aggression, irresponsibility, or lack of remorse. The pattern begins by adolescence and continues into adult life.

\synonyms
ASPD; Dissocial Personality Disorder; Sociopathy

\medterm{Antispasmodic Drugs} Medicines that relax involuntary muscle and reduce cramping or spasms, most often in the digestive tract. Some are also used for bladder or bile-duct spasms; side effects depend on the medicine.

\synonyms
Antispasmodics; Spasmolytics

\medterm{Antistreptolysin O Titer} A blood test that measures antibodies against a toxin made by group A streptococcal bacteria. A raised or rising result may show a recent strep infection when investigating rheumatic fever or kidney inflammation, but does not prove an active infection.

\medterm{Antisynthetase Syndrome} An autoimmune disease in which the immune system attacks proteins needed to make other proteins. It may cause muscle inflammation and weakness, arthritis, fever, colour changes in the fingers, rough cracked skin on the hands, or lung inflammation and scarring.

\medterm{Antithrombin} A natural blood protein that slows clotting by blocking thrombin and other clotting proteins, especially factor Xa. Too little antithrombin, whether inherited or acquired, increases the risk of clots in veins.

\medterm{Antithrombin Deficiency} A condition in which the blood has too little or too little active antithrombin, a natural anti-clotting protein. It may be inherited or acquired and increases the risk of clots, especially in the veins.

\medterm{Antithrombin III Blood Test} A blood test that measures the amount or activity of antithrombin, a protein that limits clotting. A low result may indicate inherited deficiency or may occur with liver disease, kidney protein loss, severe illness, or some treatments.

\medterm{Antithrombin III Deficiency} Reduced amount or activity of antithrombin, causing increased risk of blood clots in veins. It may be inherited or result from liver disease, kidney disease, severe illness, or other causes.

\medterm{Antithymocyte Globulin} A preparation of antibodies that removes or suppresses T cells, immune cells involved in transplant rejection and some autoimmune reactions. It is used in selected transplants and blood-cell disorders and can cause allergic reactions, low blood counts, and infections.

\synonyms
Rabbit Antithymocyte Globulin
rATG
Thymoglobulin

\medterm{Antithyroglobulin Antibody Test} A blood test that detects antibodies against thyroglobulin, a protein used to make thyroid hormones. Raised levels may support Hashimoto thyroiditis or other autoimmune thyroid disease but do not determine thyroid function by themselves.

\medterm{Antitoxin} An antibody or antibody preparation that binds to and neutralizes a toxin. It can provide temporary protection or treatment for selected poisonings and infections that produce toxins.

\medterm{Antitubercular Agents} Medicines used to treat tuberculosis, an infection caused by \textit{Mycobacterium tuberculosis}. Several medicines are used together for months to prevent resistance, with monitoring for liver injury, nerve problems, vision changes, and other side effects.

\medterm{Antitussive} An agent that suppresses or relieves coughing, either centrally by dampening cough centers within the medulla oblongata (such as dextromethorphan or codeine) or peripherally by desensitizing sensory stretch receptors in the respiratory tract.

\synonyms
Cough Suppressant

\medterm{Anton Syndrome} A rare neurological condition characterized by cortical blindness accompanied by visual anosognosia (complete denial of blindness) and visual confabulation. Patients genuinely behave and report as if they can see, despite profound vision loss, classically caused by bilateral occipital lobe infarctions.

\synonyms
Anton-Babinski Syndrome; Visual Anosognosia

\medterm{Antral Follicle Count} The number of small developing egg-containing sacs seen in both ovaries, usually measured by vaginal ultrasound early in the menstrual cycle. It helps estimate ovarian reserve and response to fertility medicines but does not by itself predict natural fertility.

\synonyms
AFC; Basal Antral Follicle Count

\medterm{Antral Web} A thin, circumferential mucosal diaphragm in the gastric antrum, usually located one to two centimetres proximal to the pylorus. It can be congenital or acquired and may cause gastric outlet obstruction, nausea, postprandial vomiting, and early satiety.

\synonyms
Gastric Antral Diaphragm; Antral Septum

\medterm{Antrochoanal Polyp} A usually noncancerous growth that starts in the lining of a cheek sinus and extends into the back of the nose. It is more common in children and young adults and may cause blockage or drainage on one side.

\medterm{Antrum} A cavity or chamber within an organ. The term commonly refers to the lower part of the stomach near its outlet or the air-filled sinus in the cheek.

\medterm{Anuria} Almost complete absence of urine production, often defined in an adult as less than about 100 millilitres in 24 hours. It can signal severe kidney failure, blockage, shock, or another emergency.

\medterm{Anus} The opening at the end of the digestive tract through which stool leaves the body. Two ring-shaped muscles help control when stool and gas are released.

\medterm{Anxiety} A feeling of worry, fear, or unease that can be a normal response to danger or stress. It becomes a disorder when it is excessive, persistent, difficult to control, or interferes with daily life.

\medterm{Anxiety Disorder} Alternate index label; see \textit{Anxiety Disorders}.

\medterm{Anxiety Disorders} A group of mental-health conditions involving excessive fear, worry, or anxiety that persists, is difficult to control, or causes significant distress or impairment. Major types include generalized anxiety disorder, panic disorder, social anxiety disorder, specific phobia, agoraphobia, separation anxiety disorder, and selective mutism. Physical symptoms may include a racing heart, sweating, muscle tension, poor sleep, or stomach upset.

\medterm{Aorta} The body’s largest artery, carrying oxygen-rich blood from the left side of the heart to the rest of the body. It passes through the chest and abdomen and divides into the arteries supplying the pelvis and legs.

\medterm{Aortal} Relating to the aorta, the large artery that carries blood from the heart to the body.

\medterm{Aortic} Relating to the aorta, the main artery carrying blood from the heart to the body.

\medterm{Aortic Aneurysm} A permanent pathological dilation of the aortic lumen exceeding normal diameter by at least 50 percent, most commonly involving the infrarenal abdominal aorta or ascending thoracic aorta. Risk factors include atherosclerosis, hypertension, smoking, and connective tissue disorders.

\medterm{Aortic Aneurysm Rupture} A break through the wall of an aortic aneurysm that causes severe internal bleeding. It may produce sudden chest, back, or abdominal pain, collapse, and shock and is a life-threatening emergency.

\medterm{Aortic Aneurysm Screening} Testing people at increased risk for an aortic aneurysm before symptoms appear. Abdominal ultrasound is commonly used; eligibility depends on age, sex, smoking history, and family history.

\medterm{Aortic Angiography} Imaging of the aorta after contrast material is given, usually with X-ray, computed tomography, or magnetic resonance imaging. It can show narrowing, blockage, aneurysm, tearing, or abnormal blood flow.

\medterm{Aortic Arch} The curved part of the aorta between its upward and downward sections. Three major arteries branch from it to supply the head, neck, and arms; disease here can affect blood flow or press on nearby structures.

\medterm{Aortic Arch Syndrome} Reduced blood flow caused by narrowing or blockage of the aortic arch or its major branches. It may cause unequal pulses or blood pressure, arm fatigue, dizziness, or neurological symptoms.

\medterm{Aortic Balloon Pump} A temporary heart-support device. A catheter with a balloon is placed in the main artery in the chest; the balloon inflates and deflates with each heartbeat to improve blood flow to the heart and reduce the heart's workload.

\synonyms
Intra-Aortic Balloon Pump; IABP

\medterm{Aortic Body} A small group of sensory cells near the aortic arch that detects changes in blood oxygen, carbon dioxide, and acidity. It sends signals to the brain to help adjust breathing and circulation.

\medterm{Aortic Coarctation} A narrowing of the aorta present from birth, usually near the part that connects to the fetal circulation. It can cause high blood pressure in the arms, weak leg pulses, poor growth, or heart failure in a newborn.

\medterm{Aortic Disease} Any disorder affecting the aorta, the body’s main artery, including aneurysm, dissection, narrowing, inflammation, or fatty deposits. Some cases cause no symptoms; sudden severe chest, back, or abdominal pain is an emergency.

\medterm{Aortic Dissection} A tear in the inner lining of the aorta that allows blood to split the layers of its wall. It can suddenly cause severe chest or back pain, stroke, organ damage, or death and needs emergency assessment.

\synonyms
Dissecting Aortic Aneurysm; Dissecting Aneurysm; Aortic Tear

\medterm{Aortic Embolism} Blockage of an artery by clot, fatty material, or other debris that breaks away from the aorta and travels downstream. It can suddenly reduce blood flow to a limb or organ.

\medterm{Aortic Graft Infection} Infection involving an artificial graft used to repair or replace part of the aorta. It can cause fever, bloodstream infection, bleeding, or a new abnormal connection with the intestine and requires urgent specialist care.

\medterm{Aortic Injury} Damage to the aorta, often caused by severe chest trauma or a medical procedure. A tear can cause internal bleeding, reduced blood flow, shock, and a life-threatening emergency.

\medterm{Aortic Insufficiency} Leakage through the aortic valve that allows blood to flow backward into the heart’s main pumping chamber. It can make the heart work harder and may eventually cause heart failure.

\medterm{Aortic Regurgitation} Leakage through an aortic valve that does not close completely, allowing blood to flow from the aorta back into the left ventricle during diastole. Acute regurgitation can cause sudden pulmonary edema or shock; chronic regurgitation may enlarge and weaken the left ventricle and cause breathlessness, chest discomfort, palpitations, or heart failure. Echocardiography assesses the cause and severity.

\medterm{Aortic Regurgitation Murmur} A blowing sound heard during the part of the heartbeat when the heart relaxes, caused by blood leaking backward through the aortic valve. It is often loudest along the left side of the breastbone.

\medterm{Aortic Rupture} A life-threatening break through the wall of the aorta, often following an aneurysm or dissection. It can cause sudden severe chest, back, or abdominal pain, internal bleeding, collapse, and shock.

\medterm{Aortic Sclerosis} Calcification, fibrosis, and thickening of the aortic valve leaflets without significant hemodynamic obstruction or restriction of valve opening. It is a common echocardiographic finding in older adults and is associated with increased cardiovascular risk.

\medterm{Aortic Stenosis} Narrowing of the aortic valve that obstructs blood leaving the left ventricle. Common causes include age-related calcification, a bicuspid valve, and rheumatic disease; severe disease can cause exertional breathlessness, chest pain, fainting, left-ventricular thickening, or heart failure. Echocardiography assesses severity.

\medterm{Aortic Stenosis Grading} Assessment of how severely the aortic valve is narrowed using ultrasound measurements of valve area and blood-flow speed or pressure difference. Results are generally described as mild, moderate, or severe and help guide monitoring and treatment.

\medterm{Aortic Valve} A valve between the heart’s left pumping chamber and the aorta. It opens to let blood leave the heart and closes to stop blood flowing backward.

\medterm{Aortic Valve Disease} Narrowing, leakage, or another abnormality of the valve between the heart’s left pumping chamber and the aorta. It may cause breathlessness, chest pain, tiredness, fainting, or no symptoms at first.

\medterm{Aortic Valve Insufficiency} Leakage of the aortic valve that allows blood to flow backward into the left pumping chamber between heartbeats. The extra workload can enlarge or weaken the heart over time.

\medterm{Aortic Valve Regurgitation} Backward leakage of blood through an aortic valve that does not close completely. The heart must pump harder, and severe or long-standing leakage can cause enlargement or heart failure.

\synonyms
Aortic Regurgitation; Aortic Insufficiency; Aortic Incompetence

\medterm{Aortic Valve Replacement} A cardiac surgical or catheter-directed interventional procedure in which a severely diseased or calcified aortic valve is excised or displaced and replaced with a mechanical prosthesis or a bioprosthetic tissue valve.

\synonyms
AVR; TAVR; TAVI

\medterm{Aortic Valve Stenosis} Alternate index label; see \textit{Aortic Stenosis}.

\medterm{Aortic Valve Surgery, Minimally Invasive} Surgical or transcatheter repair or replacement of the aortic valve performed without conventional full median sternotomy, utilizing mini-thoracotomy, upper hemisternotomy, or transcatheter arterial access.

\synonyms
Minimally Invasive Aortic Valve Replacement; MIAVR; TAVR; TAVI

\medterm{Aortic Valve Surgery, Open} Conventional surgical repair or replacement of a diseased aortic valve performed through a full median sternotomy with cardiopulmonary bypass and cardioplegic arrest.

\synonyms
Surgical Aortic Valve Replacement; SAVR; Open Aortic Valve Replacement

\medterm{Aortitis} Inflammation of the aorta, the body’s largest artery. It may result from an autoimmune disease or infection and can weaken the artery or cause narrowing, aneurysm, dissection, or leakage of the aortic valve.

\medterm{Aortocaval Compression Syndrome} Hemodynamic compromise occurring in late pregnancy when the gravid uterus compresses the inferior vena cava and abdominal aorta while the mother is in the supine position. Leads to decreased venous return, reduced cardiac output, hypotension, dizziness, and fetal distress, rapidly relieved by placing the patient in the left lateral decubitus position.

\synonyms
Supine Hypotensive Syndrome of Pregnancy

\medterm{Aortoenteric Fistula} A life-threatening direct communication between the aorta and the gastrointestinal tract, most frequently involving the third or fourth portion of the duodenum. It may occur primarily from an eroding abdominal aortic aneurysm or secondarily following prosthetic aortic reconstruction, classically presenting with minor herald gastrointestinal bleeding followed by catastrophic exsanguination.

\medterm{Aortography} Imaging of the aorta after contrast material is injected, often through a catheter. It can show aneurysms, tears, narrowing, blockages, or abnormal blood-vessel connections; computed tomography or magnetic resonance imaging is often used instead.

\medterm{Aortoiliac Occlusive Disease} Severe atherosclerotic stenosis or complete occlusion of the distal abdominal aorta and common iliac arteries. It is clinically characterized by the triad of buttock and thigh claudication, diminished or absent femoral pulses, and erectile dysfunction in men.

\synonyms
Leriche Syndrome

\medterm{Aortopulmonary Window} A rare birth defect in which an opening connects the aorta and pulmonary artery. Too much blood may flow to the lungs, causing breathing problems, poor growth, or heart failure in infancy.

\medterm{AP X-Ray} An X-ray view in which the beam enters through the front of the body and exits through the back. It is called an anteroposterior view.

\medterm{APACHE II Score} A hospital scoring system that combines vital signs, laboratory results, age, and serious long-term illnesses to estimate the severity of critical illness. A higher score generally indicates greater risk, but it does not predict an individual outcome with certainty.

\medterm{Apathy} Reduced motivation, interest, or emotional response. It can occur with depression, dementia, Parkinson disease, brain injury, other neurological illness, or medicines and is different from ordinary tiredness or sadness.

\medterm{APC Tumor Suppressor} A protein made from the APC gene that helps control cell growth and the removal of damaged cells. Harmful changes in this gene can allow excess intestinal polyps to form and increase colorectal-cancer risk.

\medterm{Apert Syndrome} A rare inherited condition in which some skull bones join too early. It causes an unusually shaped head and face and often joined fingers or toes; breathing, vision, hearing, or development may also be affected.

\medterm{Apex} The tip or pointed end of a body structure. Examples include the lower pointed end of the heart, formed mainly by the left ventricle, and the upper tip of a lung.

\medterm{Apex Beat} The outermost and lowest point on the precordium at which the cardiac impulse can be palpated or inspected during ventricular systole, typically located in the left fifth intercostal space along the midclavicular line.

\synonyms
Apical Impulse

\medterm{Apgar Score} A quick assessment of a newborn’s colour, heart rate, reflex response, muscle tone, and breathing, usually recorded one and five minutes after birth. It describes immediate adaptation and helps guide care but does not predict long-term health by itself.

\synonyms
Apgar Assessment; Newborn Apgar Score

\medterm{Aphagia} Complete inability to swallow. It may result from a blockage, neurological disease, severe pain, or loss of normal swallowing-muscle function and can lead to choking, dehydration, or aspiration.

\medterm{Aphakia} Absence of the eye’s natural lens, usually after cataract surgery, injury, or a birth difference. It causes major focusing problems and is usually corrected with an intraocular lens, contact lens, or strong glasses.

\medterm{Aphasia} A brain-related disorder that impairs speaking, understanding, reading, or writing language. It most often follows a stroke but can also result from head injury, brain tumour, infection, or degeneration.

\medterm{Apheresis} A procedure in which blood passes through a machine that removes a selected component, such as plasma, platelets, red cells, or abnormal cells, and returns the remaining blood to the person.

\medterm{Apheresis Donation} A blood donation in which a machine collects one component, such as platelets or plasma, and returns the other blood components to the donor. The process may allow more of the needed component to be collected than whole-blood donation.

\medterm{Aphonia} Complete loss of the ability to produce a normal voice, although whispering may remain possible. Causes include inflammation, vocal-cord paralysis, injury, growths, nerve disease, or psychological factors.

\medterm{Aphrasia} Loss or impairment of the ability to speak. The term is older and less precise than aphasia, which can affect speaking, understanding, reading, or writing.

\medterm{Aphthous Stomatitis} Repeated painful ulcers on the moist lining of the mouth, commonly called canker sores. They are usually shallow, heal without scarring, and are not contagious.

\medterm{Aphthous Ulcer} A small, painful ulcer on the moist lining of the mouth, commonly called a canker sore. It usually heals within one to two weeks and is not contagious.

\medterm{Apical} Relating to or located at the tip of a body structure, such as the tip of a lung, heart, or tooth.

\medterm{Apical Ballooning Syndrome} A temporary weakening and ballooning of the main pumping chamber of the heart, often after severe emotional or physical stress. It can resemble a heart attack but usually occurs without a blocked coronary artery.

\medterm{Apical Periodontitis} Inflammation or infection around the tip of a tooth root, usually caused by infected or dead pulp. It may cause pain when biting, swelling, or no symptoms and is assessed with dental examination and imaging.

\medterm{Apical Pulse} The heartbeat counted by listening over the tip of the heart, usually on the left side of the chest. It may be used when the pulse at the wrist is difficult to feel or when an accurate heart rate is needed.

\medterm{Apicoectomy} Dental surgery that removes the tip of a tooth root and nearby infected tissue, usually when infection persists after root-canal treatment. The end of the root is then sealed.

\medterm{Aplasia} Failure of an organ, tissue, or group of cells to develop, causing it to be absent or severely underdeveloped. It may be present from birth or result from inherited disease, injury, infection, or severe damage to the cells that produce it.

\medterm{Aplasia Cutis Congenita} A birth defect in which a localized area of skin, most often on the scalp, is absent or incompletely formed. It may heal as a scar or hairless patch, sometimes with an underlying skull or vascular abnormality; large, ulcerated, or atypical lesions need specialist evaluation.

\medterm{Aplastic} Describing a tissue that has failed to develop or produce its normal cells. The term is often used for bone-marrow failure, in which too few red cells, white cells, and platelets are made.

\medterm{Aplastic Anemia} A serious condition in which the bone marrow does not make enough red cells, white cells, or platelets. It can cause tiredness, infections, and bleeding and may result from immune attack, medicines, chemicals, infection, or inherited disease.

\medterm{Aplastic Crisis} A sudden, temporary stop in red-cell production that causes a rapid worsening of anemia. It most often follows parvovirus B19 infection in people whose red cells already have a shortened lifespan, such as those with sickle-cell disease.

\medterm{Apley Compression Test} A physical examination manoeuvre used to evaluate for meniscal tears of the knee joint. With the patient prone and the knee flexed to ninety degrees, the examiner applies downward axial force to the heel while rotating the tibia internally and externally; pain or joint line clicking indicates meniscal pathology.

\synonyms
Apley Grinding Test; Apley's Compression Test

\medterm{Apley Distraction Test} A clinical examination technique used to differentiate collateral ligament injuries from meniscal tears. With the patient prone and the knee flexed to ninety degrees, the examiner stabilizes the posterior thigh and applies upward traction to the lower leg with internal and external rotation; localized pain suggests ligamentous sprain rather than meniscal damage.

\synonyms
Apley's Distraction Test

\medterm{Apley Grind Test} A physical examination maneuver used in orthopedics to evaluate meniscus tears of the knee. With the patient prone and the knee flexed to 90 degrees, the examiner stabilizes the thigh, applies downward axial compressive force onto the heel while externally and internally rotating the tibia; pain or mechanical clicking indicates a torn meniscus.

\synonyms
Apley Compression Test; Apley Meniscal Test

\medterm{Apley Scratch Test} A rapid physical examination maneuver used in musculoskeletal medicine to assess the active range of motion of the shoulder joint complex. The patient is asked to reach behind the head and touch the superior medial angle of the opposite scapula (evaluating glenohumeral abduction and external rotation) and reach behind the back to touch the inferior angle of the opposite scapula (evaluating adduction and internal rotation). Asymmetry or restricted reach indicates rotator cuff pathology, adhesive capsulitis, or glenohumeral arthritis.
\synonyms{Apley Scratch Maneuver; Apley Shoulder Test}

\medterm{Apnea} A temporary pause in breathing. Prolonged or repeated episodes can lower oxygen, raise carbon dioxide, slow the heartbeat, cause loss of consciousness, or lead to death.

\medterm{Apnea Of Prematurity} Repeated pauses in breathing in a premature infant because breathing control is not fully developed. Episodes may lower oxygen or slow the heartbeat and require monitoring and, when necessary, breathing support.

\medterm{Apnea-Hypopnea Index} The average number of complete or partial breathing interruptions during each hour of sleep, measured in a sleep study. It helps describe sleep-apnea severity, but symptoms, oxygen levels, and other findings also matter.

\medterm{Apneic} Relating to a period when breathing has stopped temporarily. Apneic episodes may occur during sleep, illness, poisoning, seizures, or severe airway blockage and can quickly lower oxygen.

\medterm{Apneustic Breathing} An abnormal pattern with an unusually long pause after breathing in. It usually indicates serious injury to the brainstem and may occur with other signs of severe neurological damage.

\medterm{Apnoea} The temporary stopping of breathing; the spelling is commonly used in British English. It may result from airway blockage, reduced breathing signals from the brain, illness, medicines, or sleep-related breathing disorders.

\medterm{Apocrine Hidrocystoma} A benign cyst of an apocrine sweat gland, commonly presenting as a translucent or bluish, dome-shaped papule near the eyelid or elsewhere on the skin. It may enlarge with heat or sweating and is usually diagnosed clinically.

\medterm{Apoenzyme} The protein part of an enzyme before its required helper substance has attached. It is usually inactive or less active until the helper, called a cofactor, binds to it.

\medterm{Apolipoprotein B100} A protein on cholesterol-carrying particles such as low-density and very-low-density lipoproteins. Its blood level estimates the number of particles that can contribute to fatty buildup in artery walls.

\medterm{Apolipoprotein CII} A protein on fat-carrying particles in the blood that activates an enzyme needed to break down triglycerides. Too little or faulty apolipoprotein CII can cause triglycerides to rise to very high levels.

\medterm{Apolipoproteins} Proteins attached to fats in the blood that help transport and process cholesterol and triglycerides. Different apolipoproteins are found on different fat-carrying particles and have different effects.

\medterm{Aponeurosis} A broad, flat sheet of tough connective tissue that attaches a muscle to bone, another muscle, or other tissues and helps transmit the muscle’s force.

\medterm{Apophenia} Seeing meaningful patterns or connections in unrelated or random events. Mild experiences can occur normally, while marked or distressing experiences may occur with some psychiatric or neurological conditions.

\medterm{Apophyses} Bony projections that grow from or arise near a bone and usually provide attachment for muscles or ligaments. Unlike the main growth plate of a long bone, an apophysis does not normally increase the bone’s length.

\medterm{Apoptosis} A controlled process of cell death in which a cell dismantles itself and is removed with little inflammation. It helps shape developing organs and remove old, damaged, or abnormal cells; failure of this process can contribute to cancer.

\medterm{Apoptosis Extrinsic Pathway} A controlled cell-death process that begins when an outside signal attaches to a death receptor on a cell’s surface. Internal enzymes then take the cell apart in an orderly way, usually without causing the swelling and inflammation of an injury.

\medterm{Apoptosis Intrinsic Pathway} A controlled cell-death process that begins with severe stress inside a cell, such as DNA damage, lack of growth signals, or lack of oxygen. The cell’s energy-producing structures then release signals that activate enzymes to dismantle it.

\medterm{Apoptotic Body} A small membrane-bound piece of a cell that has begun controlled cell death. Immune or neighbouring cells can recognize and remove these pieces without causing much inflammation.

\medterm{Apparent Mineralocorticoid Excess} A rare inherited disorder in which cortisol acts too strongly on the kidneys because it is not changed into an inactive form normally. It can cause severe high blood pressure, low potassium, and unusually low levels of renin and aldosterone.

\medterm{Appendageal} Relating to structures that develop from the skin, including hair follicles, sweat glands, and oil-producing glands.

\medterm{Appendectomy} The surgical removal of the vermiform appendix, a small finger-shaped pouch attached to the cecum at the beginning of the large intestine. It is one of the most common emergency surgical operations worldwide, primarily performed to treat acute appendicitis (sudden inflammation and infection of the appendix) before it perforates or ruptures and causes widespread peritonitis (infection of the abdominal cavity) or an intra-abdominal abscess. An appendectomy can be performed using minimally invasive laparoscopic surgery (keyhole surgery through small puncture incisions using a miniature camera and instruments, enabling faster recovery and less scarring) or open surgery (through a single incision in the lower right abdomen, typically used if the appendix has ruptured or severe adhesions are present). Most patients recover rapidly and fully without long-term changes to digestion.

\synonyms
Appendicectomy; Laparoscopic Appendectomy; Appendix Removal; Open Appendectomy

\medterm{Appendiceal} Relating to the appendix, a narrow, finger-shaped pouch attached to the beginning of the large intestine.

\medterm{Appendiceal Adenocarcinoma} A cancer arising from gland-forming cells in the appendix. It may be found after appendicitis or unexpectedly and can spread through the abdominal lining; outlook depends on its type, size, and spread.

\medterm{Appendiceal Cancer} Cancer arising in the appendix. Types include neuroendocrine tumours, mucin-producing tumours, and adenocarcinoma; some can rupture and spread tumour cells or mucus through the abdomen.

\medterm{Appendiceal Rupture} A tear or hole in an inflamed appendix that allows infected material into the abdomen. It can cause peritonitis, an abscess, sepsis, or severe abdominal pain and requires urgent care.

\medterm{Appendiceal Tumors} Tumours arising in the appendix, including neuroendocrine tumours, mucin-producing tumours, and adenocarcinoma. Some cause no symptoms, while others cause appendicitis, pain, blockage, or abdominal swelling.

\medterm{Appendicitis} Inflammation of the appendix, often after its opening becomes blocked by hardened stool, mucus, or swollen tissue. It commonly causes pain that moves to the lower right abdomen, nausea, loss of appetite, and fever and can lead to rupture.

\synonyms
Epityphlitis; Acute Appendicitis

\medterm{Appendicular Skeleton} The bones of the arms and legs and the shoulder and pelvic girdles that attach them to the central skeleton.

\medterm{Appendix} A small, finger-shaped pouch attached to the beginning of the large intestine. It contains immune tissue and may help support healthy gut bacteria, although its exact function is not fully understood and it is not essential for digestion. Inflammation of the appendix is called appendicitis.

\medterm{Appendix Cancer} A malignant neoplasm arising within the appendix, including neuroendocrine tumors, goblet cell carcinoids, mucinous cystadenocarcinomas, and colonic-type adenocarcinomas.

\medterm{Appendix Testis} A small remnant of tissue near the upper part of the testis, left over from early development. It is usually harmless but can sometimes twist and cause pain in boys.

\medterm{Appetite} The psychological desire to eat food, distinct from physiological hunger, modulated by central nervous system pathways, neuroendocrine signaling, sensory stimuli, emotional states, and metabolic demands.

\medterm{Appetite Regulation} The control of hunger and fullness by the brain, digestive tract, hormones, emotions, habits, sleep, and surroundings. Signals such as ghrelin promote hunger, while leptin and other signals help promote fullness.

\medterm{Applied Behavior Analysis} A structured approach that observes behaviour and uses learning principles, practice, and reinforcement to build useful skills or reduce behaviour that interferes with daily life. Goals should be individualized and respectful of the person.

\medterm{Apposition} Placement of two structures next to or against each other. In wound care, good apposition means that the edges of a cut are accurately aligned so they can heal.

\medterm{Apprehension Test} A provocative orthopedic examination maneuver designed to detect joint instability, most commonly applied to the shoulder (glenohumeral joint) or knee (patellofemoral joint). In the anterior shoulder apprehension test, the examiner places the arm in 90 degrees of abduction and maximal external rotation; a positive test occurs when the patient exhibits visible fear, alarm, or resistance due to a feeling of impending anterior subluxation or dislocation. In the patellar apprehension test (Fairbank's sign), lateral displacement of the patella elicits anxiety and quadriceps contraction.
\synonyms{Anterior Apprehension Test; Crank Test; Fairbank's Apprehension Test}

\medterm{Appropriate For Gestational Age} Describing a fetus or newborn whose estimated size or weight is within the expected range for the number of weeks of pregnancy, commonly between the 10th and 90th percentiles.

\medterm{Apraxia} Difficulty performing a learned, purposeful movement even though strength, sensation, coordination, and understanding are sufficient. It results from a problem in the brain’s planning of movement rather than simple weakness.

\medterm{Apraxia of Speech} A neurogenic motor speech disorder characterized by an impaired ability to plan, sequence, and program the sensorimotor commands necessary for phonetically normal speech production in the absence of significant neuromuscular weakness or paralysis. Marked by inconsistent articulatory errors, effortful groping, and disrupted prosody.

\synonyms
Verbal Apraxia; Speech Apraxia

\medterm{Apraxic Gait} A higher-level neurological gait disorder characterized by the loss of the ability to coordinate voluntary stepping and leg movements in the absence of primary sensory loss, weakness, or cerebellar ataxia. Patients exhibit hesitation in initiating walking, a wide-based stance, short shuffling steps, and a tendency for the feet to appear glued to the floor (\"magnetic gait\"), commonly seen in normal pressure hydrocephalus, bilateral frontal lobe lesions, and advanced subcortical vascular encephalopathy.
\synonyms{Gait Apraxia; Frontal Gait Disorder; Magnetic Gait; Bruns Apraxia}

\medterm{Aprosodia} A neurological speech and language disorder characterized by an inability to express or comprehend the emotional tone, rhythm, pitch, and melody (prosody) of verbal communication. It is classically associated with lesions of the non-dominant (right) cerebral hemisphere in areas homologous to Broca and Wernicke expressive and receptive language regions.

\synonyms
Dysprosody; Affective Agnosia of Speech

\medterm{Aquagenic Keratoderma} Temporary painful whitening, thickening, and wrinkling of the palms or soles after contact with water. It may cause prickling or burning, is often associated with cystic fibrosis, and usually settles after the skin dries.

\medterm{Aquaporin} A channel in a cell membrane that lets water pass through. Aquaporin-2 in kidney tubules responds to vasopressin and helps the kidneys return water to the blood instead of losing it in urine.

\medterm{Aquatic Therapy} Exercise or rehabilitation performed in water. Water buoyancy reduces load on joints while water resistance strengthens muscles and supports balance, movement, and conditioning.

\medterm{Aqueous Flare} A hazy beam seen in the clear fluid at the front of the eye during slit-lamp examination. It means inflammation has allowed proteins to leak into the eye and may occur with uveitis, infection, injury, or surgery.

\medterm{Aqueous Humor} The clear fluid in the front part of the eye. It nourishes the cornea and lens, removes waste, and helps maintain pressure and the shape of the eye.

\medterm{Arachidonic Acid Pathway} A major cellular lipid metabolic cascade in which arachidonic acid, released from cell membrane phospholipids by phospholipase A2, is converted into bioactive eicosanoids (prostaglandins, prostacyclin, thromboxane A2, and leukotrienes) via cyclooxygenase and lipoxygenase enzymes.

\medterm{Arachnodactyly} Abnormally long, thin fingers or toes. It may occur alone or as a feature of inherited connective-tissue disorders such as Marfan syndrome.

\medterm{Arachnoid Cyst} A usually noncancerous, fluid-filled sac between layers of the membrane covering the brain or spinal cord. Many cause no symptoms, but a large cyst may cause headache, seizures, weakness, or pressure on nearby tissue.

\medterm{Arachnoid Cysts} Fluid-filled sacs between layers of the membrane covering the brain or spinal cord. Many are discovered by chance and cause no symptoms; larger cysts can press on nearby tissue and cause neurological problems.

\medterm{Arachnoid Granulations} Macroscopic tufted outgrowths of the arachnoid membrane that project through the dura mater into the dural venous sinuses, predominantly the superior sagittal sinus, serving as the primary structural pathway for the passive reabsorption of cerebrospinal fluid into venous blood.

\synonyms
Pacchionian Bodies; Arachnoid Villi

\medterm{Arachnoid Mater} The middle of three thin membranes covering the brain and spinal cord. The space beneath it contains the fluid that cushions the nervous system and carries important blood vessels.

\medterm{Arachnoiditis} Inflammation of the arachnoid, one of the membranes covering the brain and spinal cord, most often in the spine. Inflammation may cause thickening, scarring, and bands that make nerve roots stick together or interfere with the normal flow of cerebrospinal fluid. Possible causes include infection, bleeding around the spinal cord, spinal injury, chronic nerve compression, surgery or another spinal procedure, and chemical irritation. Symptoms may include severe or burning back and limb pain, numbness, tingling, weakness, abnormal reflexes, or problems with bladder or bowel function. The course may remain stable or worsen over time.

\medterm{Arbovirus} A virus spread by a biting insect or tick, such as a mosquito, tick, or sandfly. Arboviruses can cause fever, rash, joint pain, bleeding, or inflammation of the brain, depending on the virus.

\medterm{ARBs} Medicines that block the main receptor for angiotensin II, relaxing blood vessels and reducing salt and water retention. They lower blood pressure and can protect the heart and kidneys in selected conditions, but are usually not combined with ACE inhibitors.

\medterm{Arcanobacterium haemolyticum Infection} A bacterial infection that commonly causes pharyngitis in adolescents and young adults, with fever, sore throat, tonsillar exudate, cervical lymph-node enlargement, and sometimes a scarlatiniform or morbilliform rash. It can resemble group A streptococcal infection, so throat testing and clinical assessment guide treatment.

\medterm{Arciform} Having a curved or arch-like shape. The term may describe the appearance of a skin lesion, rash, blood vessel, or other body structure.

\medterm{Arcuate Uterus} A mild variation in the shape of the inside top of the uterus, present from birth, that creates a shallow curved indentation. It is usually considered a normal variant and generally has little or no effect on fertility or pregnancy.

\medterm{Arcus Senilis} A greyish-white, circular lipid deposition ring appearing in the peripheral corneal stroma, commonly seen in elderly individuals as a benign finding, or in younger patients (arcus juvenilis) associated with familial hyperlipidemia.

\synonyms
Corneal Arcus; Arcus Lipoides; Gerontoxon

\medterm{Area Under The Concentration-Time Curve} A measure of the total amount of a medicine in the blood over a period of time. It helps compare how medicines are absorbed and removed and can guide dosing for selected treatments.

\medterm{Arenaviruses} A family of viruses carried mainly by rodents. Some cause serious human diseases, including viral haemorrhagic fevers, after exposure to infected rodent urine, droppings, saliva, or contaminated dust.

\medterm{Areola} The darker circular area surrounding the nipple. It contains small glands that help lubricate and protect the skin, especially during breastfeeding.

\medterm{Argentine Hemorrhagic Fever} A serious viral illness caused by Junín virus, an arenavirus carried by rodents in parts of Argentina. It can cause fever, weakness, bleeding, and nervous-system problems and requires urgent medical care.

\medterm{Argininosuccinate Lyase Deficiency} An inherited urea-cycle disorder caused by deficient argininosuccinate lyase. Argininosuccinic acid and ammonia can accumulate, causing poor feeding, vomiting, liver problems, developmental impairment, or life-threatening hyperammonemia.

\medterm{Argininosuccinic Aciduria} A rare inherited disorder in which the body cannot safely remove waste nitrogen produced when protein is broken down. Ammonia and argininosuccinic acid build up and may cause vomiting, sleepiness, seizures, liver problems, or developmental impairment.

\medterm{Argon Plasma Coagulation} A procedure that uses electrically charged argon gas to deliver heat to the surface of tissue. It can stop bleeding or destroy selected abnormal tissue, often during an endoscopic examination.

\synonyms
APC; Argon Plasma Hemostasis

\medterm{Argyll Robertson Pupil} A small, irregular pupil that becomes smaller when focusing on a nearby object but reacts poorly or not at all to light. It is classically associated with late syphilis affecting the nervous system, but other causes are possible.

\medterm{Arias-Stella Reaction} A benign change in cells lining the uterus or cervix caused by pregnancy hormones or hormone medicines. Under a microscope it can resemble cancer, so the clinical setting and other findings are important.

\medterm{Arm} The part of the upper limb between the shoulder and elbow. In everyday language, “arm” may also mean the entire upper limb, including the forearm and hand.

\medterm{Arm CT Scan} A scan that uses X-rays and computer processing to create cross-sectional images of the arm. It can show fractures, bleeding, inflammation, masses, and other structural problems.

\medterm{Arm Fracture} A break in one or more bones of the upper limb, usually the upper arm bone or one of the forearm bones. It may cause pain, swelling, bruising, deformity, or inability to move the arm normally.

\medterm{Arm MRI Scan} A scan that uses a magnetic field and radio waves to create detailed images of the arm. It can show muscles, tendons, ligaments, nerves, bone marrow, blood vessels, and masses without X-ray radiation.

\medterm{Arm Pain} Pain in the shoulder, upper arm, elbow, forearm, or nearby tissues. Causes include muscle or tendon injury, nerve compression, arthritis, and blood-vessel disease; sudden severe pain with chest discomfort or breathlessness is an emergency.

\synonyms
Brachialgia; Upper-Limb Pain

\medterm{Armpit Lump} A lump or swelling in the armpit, commonly caused by an enlarged lymph node, skin infection, cyst, or noncancerous growth. A hard, fixed, enlarging, persistent, or unexplained lump needs medical assessment.

\medterm{Arnold-Chiari Malformation} A birth defect in which part of the lower brain extends through the opening at the base of the skull into the spinal canal. It may block the normal flow of brain fluid and cause headache, neck pain, balance problems, swallowing difficulty, or weakness.

\medterm{Aromatase} An enzyme that changes androgen hormones into oestrogen. It is found in tissues including the ovaries, fat, and breast and helps regulate hormone levels.

\medterm{Aromatase Excess Syndrome} A rare inherited condition in which too much aromatase causes excess oestrogen production. It may cause breast development in males and early puberty, rapid growth followed by short adult height, or other hormone-related changes.

\medterm{Arousal Index} The average number of brief awakenings or shifts to lighter sleep during each hour of a sleep study. A high result suggests repeatedly interrupted sleep, but its importance depends on symptoms and other sleep-study findings.

\medterm{Arrector Pili} A small involuntary muscle attached to a hair follicle. When it contracts, it raises the hair and produces goosebumps.

\medterm{Arrector Pili Muscle} The small involuntary muscle attached to each hair follicle. Nerve signals can make it contract, raising the hair and producing temporary goosebumps.

\synonyms
Piloerector Muscle; Arrectores Pilorum

\medterm{Arrested Caries} Tooth decay that has stopped progressing, often because the affected surface has become hard, smooth, and easier to clean. It may remain stable but still requires dental assessment and prevention of new decay.

\medterm{Arrhythmia} An abnormal heart rate or rhythm. The heartbeat may be too fast, too slow, irregular, or start from the wrong place; some arrhythmias are harmless, while others reduce blood flow or cause cardiac arrest.

\medterm{Arrhythmias} Abnormal heart rates or rhythms caused by changes in the heart’s electrical signals. They range from harmless extra beats to rhythms that cause fainting, poor circulation, or cardiac arrest.

\medterm{Arrhythmogenic Right Ventricular Cardiomyopathy} An inherited disease in which heart muscle, especially in the right pumping chamber, is replaced by scar tissue and sometimes fat. This can cause dangerous irregular rhythms, fainting, or sudden cardiac arrest, particularly during exercise.

\medterm{Arsenic} A naturally occurring element that is poisonous in many forms. Exposure can cause vomiting, diarrhoea, nerve problems, heart problems, skin changes, and increased cancer risk; some compounds have limited medical uses.

\medterm{Arsenic Poisoning} Illness caused by sudden or long-term exposure to arsenic. It may cause vomiting, diarrhoea, abdominal pain, heart problems, nerve damage, anemia, skin changes, or increased cancer risk and is an emergency requiring urgent medical treatment.

\synonyms
Arsenic Toxicity; Arsenicosis

\medterm{Arsenic Toxicity} Harm caused by arsenic exposure. It may affect the digestive system, heart, nerves, blood, skin, and nails; severe exposure can be fatal.

\medterm{Artemisinin} A compound originally obtained from the sweet-wormwood plant and used, with other medicines, to treat malaria. Combination treatment is important because malaria parasites can become resistant when artemisinin is used alone.

\medterm{Arteria Femoralis} The femoral artery, the main artery supplying the thigh and lower limb. It continues from the external iliac artery below the groin and can be felt or accessed near the groin.

\medterm{Arterial Blood Gas} A blood test, usually from an artery at the wrist, that measures acidity, oxygen, carbon dioxide, and bicarbonate. It helps assess breathing, oxygen delivery, and acid–base disorders.

\medterm{Arterial Blood Gas Analyzer} A machine that measures acidity, oxygen, carbon dioxide, and related values in arterial blood. The results help assess how well a person is breathing and whether the blood is too acidic or alkaline.

\medterm{Arterial Catheter} A thin tube placed in an artery, often at the wrist, to measure blood pressure continuously and take repeated blood samples. Possible complications include bleeding, infection, clotting, or reduced blood flow.

\medterm{Arterial Catheterization} Placement of a thin tube in an artery, usually at the wrist, for continuous blood-pressure monitoring or repeated arterial blood tests. It is used when close monitoring is needed and can cause bleeding, infection, clotting, or reduced circulation.

\medterm{Arterial Dissection} A pathological process in which a disruption or tear in the tunica intima allows circulating blood under pressure to enter and cleave the media, creating a false lumen that tracks longitudinally along the arterial wall and may lead to vessel occlusion or rupture.

\medterm{Arterial Duplex} An ultrasound test that shows the structure of arteries and measures blood flow through them. It can detect narrowing, blockage, aneurysm, or whether a vascular graft remains open.

\medterm{Arterial Embolism} Sudden blockage of an artery by material carried in the bloodstream, usually a blood clot. It can abruptly cut off blood flow to a limb, brain, kidney, intestine, or other organ and is an emergency.

\medterm{Arterial Embolization} A procedure that deliberately blocks a selected artery, usually to stop bleeding, reduce blood supply to a tumour, or treat an abnormal blood-vessel connection.

\medterm{Arterial Injury} Damage to an artery that may cause severe bleeding, a clot, narrowing, reduced blood flow, tissue loss, or life-threatening complications. It can result from trauma, surgery, or a medical procedure.

\medterm{Arterial Insufficiency} Reduced blood flow through an artery to a tissue or limb, usually because the artery is narrowed or blocked. It may cause pain with walking, pain at rest, cool skin, slow-healing wounds, or weak pulses.

\medterm{Arterial Line} A thin tube placed in an artery, usually at the wrist, for continuous blood-pressure monitoring and repeated arterial blood sampling. It is used in critical illness or major surgery and can cause bleeding, infection, clotting, or reduced circulation.

\medterm{Arterial Resistance} The opposition blood encounters as it flows through arteries, especially the small arteries that can narrow or widen. It helps control blood pressure and how much blood reaches each tissue.

\synonyms
Vascular Resistance; Systemic Vascular Resistance

\medterm{Arterial Spin Labeling} An MRI method that uses magnetically labeled blood as a natural tracer to measure blood flow, especially in the brain, without injecting a tracer.

\medterm{Arterial Stenosis} Narrowing of an artery that reduces blood flow to the tissues beyond it. Common causes include fatty plaque, inflammation, injury, or an inherited abnormality of the artery wall.

\medterm{Arterial Stick} Puncture of an artery with a needle, usually to obtain an arterial blood sample. It may cause temporary pain or bruising and, rarely, bleeding, arterial injury, or reduced blood flow.

\medterm{Arterial Stiffness} Reduced ability of arteries to expand and recoil with each heartbeat. It is associated with ageing, high blood pressure, diabetes, kidney disease, and fatty artery disease and is linked with higher cardiovascular risk.

\medterm{Arterial Switch Operation} Open-heart surgery usually performed in newborns whose two main arteries leave the wrong pumping chambers. The arteries and coronary arteries are reconnected so blood flows normally; lifelong heart follow-up is needed.

\synonyms
Jatene Procedure; ASO

\medterm{Arterial Tension} A measurement taken from arterial blood. Depending on the context, it may mean blood pressure inside an artery or the amount of oxygen or carbon dioxide dissolved in the blood.

\medterm{Arterial Thromboembolism} Blockage of an artery by a blood clot that formed elsewhere and travelled through the bloodstream. It can suddenly reduce blood flow to the brain, heart, limb, kidney, intestine, or another organ.

\medterm{Arterial Thrombosis} Formation of a blood clot inside an artery, often where a fatty plaque has damaged the vessel lining. The clot can block blood flow to the heart, brain, limb, intestine, or another organ.

\medterm{Arterial Ulcer} A slow-healing wound caused by poor blood flow through an artery, usually on a toe, heel, or other pressure point. It may be deep and sharply outlined, with a pale or dead-looking base and cool, shiny, hairless surrounding skin.

\medterm{Arteries} Blood vessels that carry blood away from the heart. Most carry oxygen-rich blood, but the pulmonary and umbilical arteries carry oxygen-poor blood; strong elastic and muscular walls help withstand and regulate pressure.

\synonyms
Arterial Vessels

\medterm{Arteriogram} An image of an artery, usually made after contrast material is injected. It can show narrowing, blockage, aneurysm, abnormal connections, or other changes in blood flow.

\medterm{Arteriole} A small branch of an artery that delivers blood to a network of tiny vessels called capillaries. Its muscular wall can narrow or widen to control blood flow and blood pressure.

\medterm{Arterioles} Small branches of arteries that control blood flow into capillaries by narrowing or widening their muscular walls. They are important in regulating blood pressure and the supply of blood to tissues.

\synonyms
Resistance Vessels; Small Arteries

\medterm{Arteriolosclerosis} A form of vascular disease affecting small arteries and arterioles, characterized by mural thickening, luminal narrowing, and loss of elasticity, classically manifesting as hyaline arteriolosclerosis (associated with aging and diabetes) or hyperplastic arteriolosclerosis (associated with severe hypertension).

\medterm{Arteriosclerosis} Thickening, hardening, and loss of elasticity in artery walls. It includes several processes, including age-related stiffening and plaque buildup, and can reduce blood flow to the heart, brain, legs, or other organs.

\medterm{Arteriosclerosis / Atherosclerosis} Arteriosclerosis means hardening or thickening of artery walls. Atherosclerosis is a type in which fatty, inflammatory plaque builds up inside the walls and may restrict blood flow or rupture.

\medterm{Arteriosclerotic Retinopathy} Damage or changes in the retina caused by thickened, narrowed, or hardened small blood vessels, most often after long-standing high blood pressure. Severe disease can reduce vision.

\medterm{Arteriovenous Fistula} An abnormal connection between an artery and a vein that bypasses the normal capillary network. It may be present from birth, follow injury or a procedure, or be created to provide access for haemodialysis.

\medterm{Arteriovenous Fistula Creation} Surgery to connect an artery to a nearby vein to create blood access for hemodialysis. Over time, the vein may become wider and stronger so it can be used for dialysis needles. It often has advantages over a graft or catheter, but its success, risks, and long-term usefulness vary from person to person.

\synonyms
AV Fistula Creation; Arteriovenous Fistula Surgery; Hemodialysis Fistula Placement; AVF Construction

\medterm{Arteriovenous Graft} An artificial tube connecting an artery and vein to provide access for haemodialysis, usually when a person’s veins cannot form a suitable fistula. It can clot or become infected and needs regular monitoring.

\medterm{Arteriovenous Malformation} An abnormal tangle of arteries and veins joined without the usual capillaries between them. It can bleed, cause seizures or other neurological symptoms, or put extra strain on the heart.

\medterm{Arteriovenous Malformation Resection} A complex neurosurgical operation performed to remove an arteriovenous malformation (AVM)—an abnormal, high-pressure tangle of blood vessels connecting arteries directly to veins in the brain without intervening capillaries. Because the thin-walled veins cannot tolerate direct arterial pressure, AVMs carry high risks of catastrophic brain bleeding (hemorrhagic stroke), seizures, and chronic headaches. Under high-magnification operating microscopes and continuous image guidance, a neurosurgeon enters through a craniotomy, painstakingly identifies and seals off every feeding artery supplying the malformation, coagulates and detaches the draining veins, and completely excises the vascular nidus from surrounding healthy brain tissue to eliminate the danger of hemorrhage permanently.
\synonyms
AVM Resection; Surgical Excision of Brain AVM; Cerebral AVM Removal

\medterm{Arteriovenous Nicking} A clinical fundoscopic finding in chronic hypertensive retinopathy characterized by the apparent tapering, deflection, or constriction of a retinal venule where it is crossed by a thickened, sclerotic retinal arteriole (Gunn's sign). Because the arteriole and venule share a common adventitial sheath at crossing points, progressive arteriolosclerosis compresses the thinner-walled venule, indicating sustained moderate-to-severe systemic vascular hypertension and microvascular end-organ damage.
\synonyms{AV Nicking; Gunn's Sign; Arteriovenous Crossing Changes}

\medterm{Arteritis} Inflammation of an artery wall. The inflammation can narrow, weaken, or block the artery, reducing blood flow and causing pain or organ damage in the affected area.

\medterm{Arteritis, Giant Cell} See \textit{Giant cell arteritis}.

\medterm{Arteritis, Takayasu's} See \textit{Takayasu's arteritis}.

\medterm{Artery} A blood vessel that carries blood away from the heart. Most arteries carry oxygen-rich blood, except the pulmonary arteries, which carry oxygen-poor blood to the lungs. Their elastic, muscular walls withstand and regulate pressure.

\medterm{Artery Occlusion} A blockage inside an artery, often caused by a clot, fatty plaque, or material that travelled from elsewhere. It can cut off blood flow; sudden pain, paleness, coldness, numbness, weakness, or absent pulse is an emergency.

\medterm{Arthralgia} Pain in one or more joints without clear signs of inflammation, such as swelling, warmth, or redness. Causes include overuse, injury, infection, autoimmune disease, and early joint wear.

\medterm{Arthritis} A group of conditions that affect joints and the tissues around them; it is not one single disease. Depending on the type, arthritis may involve inflammation, gradual breakdown of joint tissues, infection, or crystal deposits. It can cause joint pain, swelling, warmth, stiffness, and reduced movement. Some types can also affect the skin, eyes, heart, lungs, or other organs. Examples include osteoarthritis, rheumatoid arthritis, gout, and ankylosing spondylitis.

\synonyms
Joint Inflammation; Articular Inflammation

\medterm{Arthritis of the Wrist} Inflammation or degenerative disease affecting one or more wrist joints. Pain, swelling, stiffness, reduced grip, and limited motion may result from osteoarthritis, rheumatoid arthritis, crystal disease, injury, infection, or other inflammatory disorders.

\medterm{Arthritis, Ankylosing Spondylitis} See \textit{Ankylosing Spondylitis}.

\medterm{Arthritis, Basal Joint} See \textit{Thumb Arthritis}.

\medterm{Arthritis, Degenerative} See \textit{Osteoarthritis}.

\medterm{Arthritis, Gouty} See \textit{Gout}.

\medterm{Arthritis, Infectious} See \textit{Septic Arthritis}.

\medterm{Arthritis, Juvenile Idiopathic} See \textit{Juvenile Idiopathic Arthritis}.

\medterm{Arthritis, Lyme} See \textit{Lyme Disease}.

\medterm{Arthritis, Mixed Connective Tissue Disease} See \textit{Mixed Connective Tissue Disease}.

\medterm{Arthritis, Osteoarthritis} See \textit{Osteoarthritis}.

\medterm{Arthritis, Plant Thorn} See \textit{Plant Thorn Synovitis}.

\medterm{Arthritis, Pseudogout} See \textit{Pseudogout}.

\medterm{Arthritis, Psoriatic} See \textit{Psoriatic Arthritis}.

\medterm{Arthritis, Reactive} See \textit{Reactive Arthritis}.

\medterm{Arthritis, Reiter} See \textit{Reactive Arthritis}.

\medterm{Arthritis, Rheumatoid} See \textit{Rheumatoid Arthritis}.

\medterm{Arthritis, Sarcoid} See \textit{Sarcoidosis}.

\medterm{Arthritis, Scleroderma} See \textit{Systemic Sclerosis}.

\medterm{Arthritis, Septic} See \textit{Septic Arthritis}.

\medterm{Arthritis, Sjogren} See \textit{Sjogren Syndrome}.

\medterm{Arthritis, Still} See \textit{Still Disease}.

\medterm{Arthritis, Systemic Lupus Erythematosus} See \textit{Systemic Lupus Erythematosus}.

\medterm{Arthritis, Thumb} See \textit{Thumb Arthritis}.

\medterm{Arthrocentesis} A procedure in which a needle removes fluid from a joint. The fluid can be tested for infection, blood, or crystals and may be removed to reduce pain and pressure.

\medterm{Arthrodesis} Surgery that permanently joins the bones of a damaged joint so they heal as one solid unit. It can relieve severe pain and improve stability, but the fused joint no longer moves.

\medterm{Arthrogram} An image of a joint taken after contrast material is placed inside or around it. It can show tears or problems in cartilage, ligaments, the joint capsule, or other joint structures.

\medterm{Arthrography} Imaging of a joint after contrast material is injected into it. It can help examine the shoulder, hip, wrist, knee, or another joint when ordinary imaging does not show enough detail.

\medterm{Arthrogryposis} A group of conditions present from birth in which two or more joints are permanently stiff or limited in movement. It usually results from reduced movement before birth and may involve muscles, nerves, or connective tissue.

\medterm{Arthrogryposis Multiplex Congenita} A group of birth conditions causing stiffness and limited movement in several joints. The affected joints may be in the arms, legs, jaw, or spine, and muscle weakness or abnormal nerves and connective tissue may also be present.

\medterm{Arthrolysis} A procedure that releases scar tissue or other restrictions around a stiff joint to improve movement. It may be performed through small cuts with a camera or through open surgery.

\medterm{Arthropathy} Any disease or abnormal condition affecting a joint. It may result from inflammation, wear, injury, infection, bleeding, nerve damage, or crystals such as those that cause gout.

\synonyms
Joint Disease

\medterm{Arthroplasty} Surgery to repair, reshape, or replace a damaged joint to reduce pain and improve movement or alignment. It may replace part or all of the joint and is used for severe arthritis, fractures, or other joint damage.

\medterm{Arthropod Vector} An insect, tick, or related animal that carries and transmits an infectious organism between hosts. Examples include mosquitoes spreading malaria or ticks spreading Lyme disease.
\synonyms
Disease Vector

\medterm{Arthroscope} A thin instrument with a light and camera used to look inside a joint through a small cut. It can help diagnose damage and guide keyhole surgery to repair or remove abnormal tissue.

\medterm{Arthroscopy} Keyhole surgery in which a camera and small instruments are inserted through small cuts to examine or treat a joint. It may be used to repair cartilage or ligaments, remove loose tissue, or treat an inflamed joint lining.

\medterm{Arthrosis} Degenerative damage or wear of a joint, involving loss of cartilage and changes in the underlying bone. Osteoarthritis is the most common example and can cause pain, stiffness, and reduced movement.

\medterm{Arthus Reaction} A local immune reaction that causes redness, swelling, pain, and sometimes tissue damage several hours after exposure to a substance. It occurs when antibodies and the substance form deposits that activate inflammation in the affected tissue.

\medterm{Articular Cartilage} Smooth, tough tissue covering the ends of bones inside a joint. It allows bones to move with little friction and helps spread forces across the joint.

\medterm{Articular Cartilage Injuries} Damage to the smooth cartilage covering the ends of bones in a joint, sometimes extending into the underlying bone. It may cause pain, swelling, catching, or reduced movement and can follow injury, repeated stress, or joint disease.

\medterm{Articular Ligament} A strong band of connective tissue joining bones around a joint. It stabilizes the joint and limits excessive bending, sliding, or twisting.

\synonyms
Joint Ligament

\medterm{Articulation} A joint or connection between bones. In speech, articulation also means forming clear speech sounds using the tongue, lips, teeth, and other speech organs.

\medterm{Artifacts} Distortions or false features in a medical image that do not represent real body structures. They may result from movement, metal, air, the imaging beam, or limits of the scanner and image-processing method.

\medterm{Artificial Abortion} An older term for an induced abortion: ending a pregnancy with medicines, a procedure, or both, rather than through spontaneous miscarriage.

\medterm{Artificial Heart} A mechanical device that supports or replaces the heart’s pumping action. It may support one or both lower chambers temporarily or long term and requires management of power, infection, bleeding, and blood clots.

\medterm{Artificial Insemination} A fertility procedure in which prepared sperm is placed into the cervix or uterus to increase the chance of pregnancy. The sperm may come from a partner or a donor.

\medterm{Artificial Insemination By Donor} Fertility treatment in which sperm from a screened donor is placed into the uterus or cervix to attempt pregnancy. It may be used when partner sperm is unavailable or unlikely to result in pregnancy.

\medterm{Artificial Insemination By Husband} Fertility treatment in which prepared sperm from a woman’s husband or male partner is placed into the uterus or cervix to increase the chance of fertilization and pregnancy.

\medterm{Artificial Intelligence} Computer systems that learn from data to recognize patterns, make predictions, or support healthcare decisions. Medical AI can assist diagnosis, treatment, and research, but it must be tested, monitored for bias, and supervised by qualified professionals.

\synonyms
AI

\medterm{Artificial Kidney} A dialysis device that removes metabolic waste and excess fluid from blood and helps correct electrolyte disturbances when renal function is inadequate.

\medterm{Artificial Pacemaker} An implanted device that sends electrical signals to the heart when its natural rhythm is too slow or irregular. It helps maintain an adequate heart rate and may have one or more leads.

\medterm{Artificial Rupture Of Membranes} Deliberate breaking of the amniotic sac during labour or an examination. It may be used to start or strengthen labour, assess the amniotic fluid, or attach internal monitors; it carries risks such as infection and umbilical-cord prolapse.

\medterm{Artificial Saliva} A product that moistens the mouth and temporarily relieves discomfort from dry mouth. It may improve speaking, swallowing, and comfort but does not replace fluids, dental care, or treatment of the cause.

\synonyms
Saliva Substitute

\medterm{Artificial Skin} A biological or manufactured covering placed over a wound to protect it and support growth of new skin. It may reduce fluid loss and contamination, but some wounds later need a permanent skin graft.

\synonyms
Skin Substitute

\medterm{Artificial Sweeteners} Sweet-tasting substances that provide few or no calories and are used instead of sugar. They can reduce added sugar for some people but do not by themselves make a diet healthy; effects and safety depend on the product and amount.

\medterm{Artificial Urinary Sphincter} An implanted device that controls urine leakage by gently closing and releasing the urethra. It is used for selected severe stress incontinence and requires surgery and follow-up to operate and maintain the device.

\medterm{Artificial Ventilation} Support or replacement of breathing with a machine that moves air and oxygen into and out of the lungs when a person cannot breathe well enough alone. It may use a mask or a tube and requires close monitoring.

\medterm{Artificial Ventricle} A mechanical device that helps or replaces the pumping action of one or both lower heart chambers in severe heart failure. The exact device and treatment goal vary; it is not the same as every ventricular assist device or a total artificial heart.

\synonyms
Total Artificial Heart; TAH; Ventricular Assist Device

\medterm{Aryepiglottic Fold} A triangular mucosal fold enclosing ligamentous and muscular fibers, extending bilaterally from the lateral margins of the epiglottis to the arytenoid cartilages, forming the lateral anatomical boundaries of the laryngeal inlet.

\medterm{Arytenoid Cartilage} One of a pair of small cartilages at the back of the voice box. They anchor the vocal cords and help the laryngeal muscles open, close, and adjust the cords for breathing and speech.

\medterm{ASA Physical Status Classification System} A scale used before anesthesia to summarize a patient’s overall health and major medical conditions, from healthy to severely ill. It helps the care team communicate risk but does not by itself predict complications or outcome.

\medterm{Asbestos-Related Disorders} Diseases that may develop years after breathing asbestos fibres, including pleural plaques, lung scarring, lung cancer, and mesothelioma. Risk increases with greater exposure and smoking; symptoms may include breathlessness, cough, or chest pain.

\medterm{Asbestosis} Permanent scarring of the lungs caused by long-term inhalation of asbestos fibres. It can gradually cause breathlessness, dry cough, and reduced lung function and increases the risk of lung cancer and mesothelioma.

\medterm{Ascariasis} Infection with a roundworm acquired by swallowing eggs from contaminated soil, food, or water. It may cause no symptoms, abdominal pain, poor nutrition, intestinal blockage, or cough while larvae pass through the lungs.

\medterm{Ascending Aorta} The first section of the aorta, rising from the heart’s left pumping chamber. It gives rise to the coronary arteries and continues into the aortic arch; an aneurysm or dissection here can be life-threatening.

\medterm{Ascending Aortic Aneurysm} An enlarged, weakened area in the part of the aorta that rises from the heart. It may be associated with a bicuspid aortic valve, inherited connective-tissue disease, high blood pressure, or artery-wall disease and may rupture or dissect.

\medterm{Ascending Cholangitis} A life-threatening bacterial infection of the biliary tree caused by a combination of bile duct obstruction (most commonly a gallstone in the common bile duct or a tumor) and bacterial ascent from the duodenum into stagnant bile under pressure. It characteristically presents with Charcot's triad (fever with chills, right-upper-quadrant pain, and jaundice), which can rapidly progress to Reynolds' pentad (adding confusion and septic shock/hypotension in acute suppurative cases). Blood tests show leukocytosis and hyperbilirubinemia, and imaging (ultrasound, CT, or MRCP) confirms biliary duct dilation. Emergency management consists of resuscitation with intravenous fluids, broad-spectrum IV antibiotics, and urgent biliary decompression via endoscopic retrograde cholangiopancreatography (ERCP).
\synonyms{Acute Cholangitis; Acute Suppurative Cholangitis; Biliary Sepsis}

\medterm{Ascending Colon} The first part of the large intestine, running upward on the right side of the abdomen from the appendix region to the bend beneath the liver. It absorbs water and helps form stool.

\medterm{Ascites} Abnormal buildup of fluid inside the abdomen. It commonly results from advanced liver disease but may also be caused by cancer, heart failure, kidney disease, infection, or pancreatitis and can cause swelling, discomfort, or breathlessness.

\medterm{Ascitic Fluid} Fluid collected from the abdomen in a person with ascites. Laboratory testing can help identify liver disease, infection, cancer, pancreatitis, or other causes of the fluid buildup.

\synonyms
Peritoneal Fluid; Ascitic Effusion

\medterm{Ascorbic Acid} Vitamin C, a water-soluble vitamin needed to make collagen, heal wounds, absorb iron, protect cells, and support normal immune and connective-tissue function.

\medterm{Ascus} A Pap-test result meaning that some cells from the cervix look slightly unusual, but it is unclear whether the change is important. Follow-up depends on age, HPV testing, previous results, and clinical guidance; it is not a cancer diagnosis.

\medterm{Asepsis} The condition or practice of preventing contamination by disease-causing microorganisms, especially during medical or surgical procedures. It includes sterile technique, environmental controls, hand hygiene, and appropriate handling of instruments and materials.

\medterm{Aseptic} Free from disease-causing contamination, or performed using methods that prevent microorganisms from entering a wound, specimen, medicine, or body site. Aseptic does not necessarily mean completely sterile.

\medterm{Aseptic Loosening} Loss of firm attachment of an artificial joint or other implant without an active infection. Repeated stress or a reaction to tiny wear particles can cause pain, instability, and eventual implant failure.

\medterm{Aseptic Meningitis} Inflammation of the membranes around the brain and spinal cord without the usual evidence of bacterial meningitis. Viruses are common causes, but medicines, autoimmune disease, and other infections can also be responsible.

\medterm{Aseptic Necrosis} See \textit{Avascular necrosis (osteonecrosis)}.

\synonyms
Avascular Necrosis; Osteonecrosis; Ischemic Bone Necrosis

\medterm{Aseptic Technique} Practices used to prevent harmful microorganisms contaminating a wound, body site, medicine, or specimen. They include hand hygiene, clean or sterile equipment, surface cleaning, and limiting contact with exposed sterile items.

\medterm{Asherman Syndrome} Scar tissue inside the uterus that partly or completely joins the uterine walls. It most often follows a procedure inside the uterus or complications of pregnancy and may cause light or absent periods, infertility, or repeated miscarriage.

\medterm{Ashman Phenomenon} A temporary wide-looking heartbeat that occurs when a short beat follows a long pause, often during atrial fibrillation. It can resemble a dangerous rhythm on an ECG, so clinicians interpret it with the full tracing and the person’s symptoms.

\synonyms
Ashman's Phenomenon; Aberrant Intraventricular Conduction

\medterm{ASL} American Sign Language, a complete visual language used by many Deaf and hard-of-hearing people. Healthcare communication should respect each person’s language preference and provide qualified interpreters when needed.

\medterm{ASOT} An abbreviation for the antistreptolysin O test, a blood test measuring antibodies against a toxin made by group A streptococcal bacteria. It can help show a recent strep infection when investigating later complications.

\medterm{Asp} The standard three-letter abbreviation for aspartic acid, an amino acid used to build proteins and involved in energy production and other cell reactions.

\medterm{Aspartame} A low-calorie sweetener made from two amino-acid components. It is not suitable for people with phenylketonuria because it releases phenylalanine when digested.

\medterm{Aspartate Aminotransferase} An enzyme found mainly in the liver, heart, muscles, and red blood cells. When these tissues are injured, the enzyme can enter the blood and its level may rise.

\medterm{Aspartate Aminotransferase Blood Test} A blood test measuring the level of the enzyme aspartate aminotransferase. A high result can indicate injury to the liver, heart, muscles, or red blood cells and must be interpreted with other tests.

\medterm{Aspartic Acid} An amino acid that the body can make and uses to build proteins and carry out cell reactions. It is one of the components of aspartame but is not the same substance as the sweetener.

\medterm{Asperger Syndrome} A former diagnostic classification for a neurodevelopmental profile within the autism spectrum characterized by qualitative differences in social communication and social reciprocity, accompanied by restricted, repetitive patterns of behavior and intense focal interests, without significant delays in general language acquisition or cognitive development.

\synonyms
Asperger Disorder; Asperger's Syndrome; Asperger's Disorder; Autistic Psychopathy

\medterm{Aspergillosis} A group of infections and allergic conditions caused by Aspergillus molds. It may affect the airways, sinuses, or lungs and can include allergic disease, a fungal growth in an existing lung cavity, or invasive infection.

\synonyms
Aspergillus Infection; Aspergillus Disease

\medterm{Aspergillosis Precipitin} A qualitative or semi-quantitative serological test detecting circulating precipitating IgG antibodies against \textit{Aspergillus} antigens, useful in diagnosing chronic pulmonary aspergillosis and allergic bronchopulmonary aspergillosis.

\synonyms
Aspergillus Antibodies; Aspergillus IgG Precipitins

\medterm{Aspergillus} A genus of saprophytic molds whose airborne conidia are ubiquitous in the environment, capable of causing allergic bronchopulmonary aspergillosis, aspergilloma formation in pre-existing cavitary lung lesions, or severe invasive tissue infections in immunocompromised hosts.

\medterm{Aspergillus Infection} Alternate index label; see \textit{Aspergillosis}.

\medterm{Asphalt Cement Poisoning} Illness caused by swallowing, inhaling, or having significant skin contact with hot or solvent-containing asphalt cement. Burns, breathing problems, dizziness, or chemical injury can occur; significant exposure requires urgent medical treatment.

\medterm{Asphyxia} A life-threatening lack of oxygen caused by inability to breathe or by blocked blood oxygen delivery, such as from choking, drowning, strangulation, or severe airway blockage. It can quickly cause unconsciousness, brain injury, or death.

\medterm{Asphyxia Neonatorum} Failure of a newborn to start or maintain effective breathing at or soon after birth. It can reduce oxygen to the brain and other organs and requires immediate assessment and breathing support.

\medterm{Asphyxial Death Pathology} Autopsy findings associated with impaired oxygenation or ventilation, which may include congestion, petechiae, pulmonary edema, and visceral hypoxia. These findings are usually nonspecific and must be interpreted with scene, history, and toxicology.

\medterm{Asphyxiating Thoracic Dysplasia} A rare inherited disorder in which short ribs create a narrow, stiff chest that restricts breathing, especially in infancy. Kidney, liver, and limb abnormalities may also occur.

\synonyms
Jeune Syndrome

\medterm{Aspirate} Fluid, cellular material, tissue fragments, or particulate matter evacuated from a cavity, cyst, organ, or lesion via negative-pressure suction; or foreign liquid and solid material accidentally inhaled into the tracheobronchial tree.

\synonyms
Aspirated Material; Biopsy Aspirate

\medterm{Aspiration} Breathing food, liquid, saliva, vomit, blood, or another substance into the voice box or lungs instead of swallowing it into the oesophagus. It can cause choking, airway blockage, inflammation, or pneumonia.

\medterm{Aspiration Pneumonia} A lung infection that develops after inhaling saliva, food, liquid, or stomach contents. It is more likely with swallowing difficulty, reduced consciousness, stroke, or tube feeding and may cause cough, fever, breathlessness, or chest pain.

\medterm{Aspiration Pneumonitis} Chemical inflammation of the lungs after inhaling stomach contents, especially acidic fluid. It can cause coughing, wheezing, low oxygen, breathing difficulty, and fluid or collapse in parts of the lungs.

\medterm{Aspiration Pneumonitis and Pneumonia} Aspiration pneumonitis is an acute chemical lung injury after inhaling gastric contents; aspiration pneumonia is infection after inhaling contaminated material. Fever, cough, hypoxia, or infiltrates after an aspiration event require assessment, with treatment guided by severity and evidence of infection.

\medterm{Aspiration Risk} The likelihood that food, liquid, saliva, or stomach contents will enter the airway or lungs. Risk increases with swallowing difficulty, stroke, reduced alertness, nerve or muscle disease, reflux, and some dental or feeding problems.

\medterm{Aspirin} A medicine that can relieve pain, lower fever, reduce inflammation, and make platelets less likely to form blood clots. Low-dose aspirin is often used after a heart attack, stroke, or some heart procedures, but it is not suitable for everyone because it can cause serious bleeding. It should not be taken daily for prevention without medical advice.

\synonyms
Acetylsalicylic Acid; ASA

\medterm{Aspirin Overdose} Acute or chronic poisoning resulting from excessive salicylate ingestion. Characteristic features include tinnitus, nausea, vomiting, hyperventilation causing early respiratory alkalosis, followed by severe uncoupled oxidative phosphorylation, high anion gap metabolic acidosis, hyperthermia, confusion, seizures, and noncardiogenic pulmonary edema.

\synonyms
Salicylate Toxicity; Aspirin Poisoning; Salicylism

\medterm{Aspirin Poisoning} Alternate name; see \textit{Aspirin Overdose}.

\medterm{Aspirin-Exacerbated Respiratory Disease} A condition involving asthma, chronic nasal or sinus inflammation with polyps, and respiratory reactions to aspirin or some similar pain medicines. Exposure may trigger wheezing, nasal blockage, cough, or shortness of breath.

\medterm{Assay} A laboratory method used to detect or measure a substance, biological activity, genetic feature, or response in a sample.

\medterm{Assist-Control Ventilation} A ventilator setting that gives a full preset breath whenever the patient starts one and also gives breaths at a minimum rate if the patient does not breathe enough. Each breath is set by volume or pressure.

\medterm{Assisted Breech Delivery} A vaginal birth in which the baby’s buttocks or feet come first and a trained clinician helps deliver the body and head. It is considered only in carefully selected circumstances because breech birth can be dangerous for the baby or pregnant person.

\medterm{Assisted Circulation} Mechanical or procedural support that maintains blood flow when the heart cannot pump enough on its own. It may use a pump, oxygenator, or other device for a short time or longer-term support.
\synonyms
Mechanical Circulatory Support

\medterm{Assisted Delivery With Forceps} A vaginal birth in which a trained clinician places specially shaped instruments around the baby’s head to help complete delivery. It may be considered when labour is prolonged or the baby or pregnant person needs delivery sooner.

\medterm{Assisted Living} A residential setting where people receive help with daily activities such as bathing, dressing, meals, medicines, or transport while retaining some independence. It generally does not provide the continuous skilled nursing care of a hospital or nursing home.

\medterm{Assisted Reproductive Technology} Clinical, surgical, and laboratory procedures that involve the in vitro manipulation of human oocytes, spermatozoa, or embryos to achieve pregnancy, including in vitro fertilization (IVF), intracytoplasmic sperm injection (ICSI), gamete intrafallopian transfer (GIFT), and frozen embryo transfer.

\synonyms
ART; Assisted Reproduction; Assisted Conception

\medterm{Assistive Device} Equipment that helps a person move, communicate, perform daily activities, or remain independent. Examples include a cane, wheelchair, hearing aid, communication tool, brace, walker, or adapted household equipment.

\medterm{Assistive Device Training} Instruction in choosing, fitting, using, and maintaining equipment such as walkers, canes, crutches, braces, or wheelchairs, including safe walking, transfers, and stair use.

\medterm{Association} A relationship in which two findings, events, or conditions occur together more often than expected. An association alone does not prove that one causes the other.

\medterm{Association Cortex} Areas of the brain’s outer layer that combine information from different senses. They support memory, language, planning, judgment, learning, and other complex behavior rather than controlling one simple movement or sensation.

\medterm{Association Study} A research study examining whether a genetic variant, exposure, or other factor occurs more often with a particular disease or finding. Such a relationship does not prove that the factor causes the disease.

\medterm{Asteatotic Eczema} A dry-skin eczema pattern, also called eczema craquelé, in which fine fissures create a cracked or porcelain-like appearance, usually on the shins, arms, or trunk of older adults. Cold, dry weather, frequent washing, and diuretics can worsen it; emollients and treatment of inflammation help.

\medterm{Astereognosis} Inability to recognize a familiar object by touch alone even though basic touch sensation works. It usually results from damage to a brain area that combines sensory information, often after a stroke or other brain injury.

\medterm{Asterixis} Brief, repeated lapses in holding a posture that cause a flapping movement, often when the arms are stretched out and the wrists are bent back. It can occur with liver or kidney failure, breathing disorders, or some medicines.

\medterm{Asteroid Hyalosis} A usually harmless eye condition in which tiny calcium-and-fat particles float in the gel inside the eye and appear as bright spots. It often affects one eye and causes few symptoms, but dense particles can obscure the retina.

\medterm{Asthenia} Abnormal lack of physical strength or energy. It is a symptom rather than a disease and may occur with infection, chronic illness, malnutrition, depression, or medicine effects.

\medterm{Asthenic} Describing a person or condition marked by weakness, or sometimes a slender body build. In a clinical note, the surrounding context shows which meaning is intended; asthenia specifically means reduced strength or energy.

\medterm{Asthenopia} Eye strain or visual fatigue that may cause aching, burning, blurred vision, headache, or difficulty focusing. Refractive errors, dry eyes, poor lighting, and prolonged screen use can contribute.

\medterm{Asthenozoospermia} Abnormally reduced movement of sperm, which can make fertilization more difficult. It is assessed by semen analysis and may result from infection, a varicocele, toxins, illness, or an unknown cause.

\synonyms
Asthenospermia

\medterm{Asthma} A long-term inflammatory disease of the airways in the lungs. The airways are overly sensitive and can narrow when triggered because their lining swells, the surrounding muscles tighten, and extra mucus forms. This variable narrowing makes it harder for air to move, especially when breathing out. Symptoms may include wheezing, cough, chest tightness, and shortness of breath; they can come and go and may be triggered by allergens, infections, smoke, air pollution, exercise, certain medicines, or cold air.

\medterm{Asthma Attack} A sudden worsening of asthma in which inflamed, narrowed airways make breathing difficult. It may cause wheezing, cough, chest tightness, or breathlessness; trouble speaking, blue lips, severe exhaustion, or poor response to a reliever is an emergency.

\medterm{Asthma Controller Medications} Medicines used regularly to reduce airway inflammation and prevent asthma symptoms and attacks. Inhaled corticosteroids are the main controller medicine. Other medicines may be added when needed; long-acting bronchodilators should not be used alone for asthma.

\synonyms
Asthma Controllers; Long-Term Control Medications; Inhaled Corticosteroids for Asthma

\medterm{Asthma Exacerbation} A sudden or gradual worsening of asthma symptoms and airflow through the lungs. It may cause increasing wheeze, cough, chest tightness, or breathlessness and can become life-threatening.

\medterm{Asthma in Older Adults} Asthma in later life may be underdiagnosed because symptoms overlap with COPD, heart disease, deconditioning, and medication effects. Diagnosis uses history and objective airflow testing when possible; inhaler technique, comorbidities, cognition, and treatment adverse effects need special attention.

\medterm{Asthma in Pregnancy} Asthma may improve, worsen, or remain unchanged during pregnancy. Maintaining control with appropriate inhaled treatment is generally safer than allowing maternal hypoxia; worsening wheeze, breathlessness, or poor response to rescue medicine requires prompt care.

\medterm{Asthma Medications} Medicines for asthma include regular treatments that reduce airway inflammation and reliever treatments that quickly ease narrowed airways. The regimen depends on symptoms, lung function, attack risk, age, pregnancy, inhaler technique, and response.

\medterm{Asthma Quick-Relief Medications} Medicines used to ease asthma symptoms quickly by opening narrowed airways. They may include a short-acting bronchodilator or, for some people, an inhaled corticosteroid--formoterol inhaler. The correct reliever depends on the person’s treatment plan.

\synonyms
Asthma Relievers; Rescue Inhalers; Short-Acting Beta-Agonists; SABA

\medterm{Asthma, Adult-Onset} Alternate index label; see \textit{Adult-Onset Asthma}.

\medterm{Asthma, Exercise-Induced} See \textit{Exercise-induced asthma}.

\medterm{Asthma, Occupational} See \textit{Occupational asthma}.

\medterm{Asthma-COPD Overlap} A pattern in which a person has features of both asthma and chronic obstructive pulmonary disease, such as variable symptoms together with persistent airflow limitation. Diagnosis and treatment require individual assessment because the conditions and risks differ.

\medterm{Astigmatism} A focusing problem caused by an unevenly curved cornea or lens. Light is not focused at one clear point on the retina, causing blurred or distorted vision at near or far distances.

\medterm{Astringent} A substance that makes surface tissues contract and may reduce minor secretions or bleeding. It is used in some skin, mouth, or wound products but can irritate tissue and does not replace assessment of persistent bleeding or serious injury.

\medterm{Astrocyte} A star-shaped support cell in the brain and spinal cord. It helps nourish nerve cells, maintain the blood–brain barrier, regulate the chemical environment, and respond to injury.

\medterm{Astrocytoma} A tumour arising from astrocytes, the star-shaped support cells of the brain or spinal cord. Astrocytomas vary from slow-growing to aggressive, and symptoms depend on their location and size.

\medterm{Astrovirus Infection} An acute viral infection of the gastrointestinal tract caused by human astroviruses, a leading cause of non-bacterial gastroenteritis primarily in infants, young children, the elderly, and immunocompromised patients. It presents with watery diarrhea, mild abdominal pain, low-grade fever, and vomiting, typically resolving spontaneously within 2 to 4 days with oral rehydration therapy.

\synonyms
Astrovirus Gastroenteritis; Human Astrovirus Infection

\medterm{Asymmetrical} Unequal in shape, size, distribution, or appearance between corresponding sides. Asymmetry in a pigmented skin lesion can be a warning sign that needs assessment, especially with other changes.

\medterm{Asymptomatic} Having a disease, infection, or abnormal finding without noticeable symptoms. A person may still have detectable disease, develop complications, or transmit an infection; this differs from presymptomatic illness, in which symptoms appear later.

\medterm{Asymptomatic Adults} Adults who have no noticeable symptoms of the condition being assessed. Absence of symptoms does not exclude early or hidden disease, risk factors, or the ability to transmit an infection.

\medterm{Asymptomatic Bacteriuria} Bacteria detected in a properly collected urine sample when the person has no urinary symptoms. It does not always require treatment, but management is important in selected situations such as pregnancy or before some urinary-tract procedures.

\medterm{Asymptomatic Bacteriuria In Pregnancy} Bacteria growing in a pregnant person’s urine without burning, urgency, or other urinary symptoms. It is identified by urine culture and is important because untreated infection increases the risk of kidney infection and pregnancy complications.

\medterm{Asymptomatic Carrier} A person who carries an infectious organism without symptoms but may still transmit it to others. Whether transmission occurs depends on the organism, body site, and contact.

\medterm{Asymptomatic HIV Infection} HIV infection without noticeable symptoms at present. HIV can still multiply and gradually weaken the immune system, and the person may transmit it without effective treatment.

\medterm{Asymptomatic Infections} Infections in which the person has no noticeable symptoms. The organism may still be detectable, cause later complications, or spread to others; this differs from presymptomatic infection, in which symptoms develop later.

\synonyms
Subclinical Infections

\medterm{Asystole} Absence of effective electrical activity in the heart’s lower chambers and absence of useful pumping. It causes cardiac arrest and requires immediate resuscitation; a flat-looking ECG line must first be checked to ensure it is not an equipment or lead problem.

\medterm{Ataxia} Loss of smooth control and coordination of movement. It is a sign rather than one specific disease and may cause an unsteady walk, poor balance, clumsy movements, slurred speech, shaking, or abnormal eye movements. It can result from a problem in the cerebellum, spinal cord, sensory nerves, or balance organs and is different from muscle weakness alone. Causes include inherited disorders, stroke, multiple sclerosis, infection, vitamin deficiency, medicines, alcohol, or other toxins. The onset may be sudden or gradual, and the course depends on the cause.

\medterm{Ataxia Telangiectasia} A rare autosomal-recessive DNA-repair disorder causing progressive cerebellar ataxia, telangiectasias of the eyes or skin, immune deficiency, radiosensitivity, and increased risk of leukemia, lymphoma, or other cancers. Respiratory infections and elevated alpha-fetoprotein may occur.

\synonyms
A-T; Ataxia-Telangiectasia

\medterm{Ataxic Breathing} An irregular, unpredictable pattern in which the rate and depth of breaths vary randomly. It usually indicates severe injury to the lower brainstem and may precede breathing failure.

\medterm{Ataxic Gait} An abnormal, unsteady, and poorly coordinated walking pattern resulting from dysfunction in the motor coordination or sensory balance pathways. In cerebellar ataxia, the gait is characteristically wide-based, irregular, lurching, and clumsy, with difficulty maintaining a straight line and inability to perform tandem walking. In sensory ataxia (caused by dorsal column or peripheral proprioceptive loss), the gait is wide-based with high-stepping stamping of the feet, markedly worsening with eye closure (positive Romberg sign).
\synonyms{Cerebellar Gait; Sensory Ataxic Gait; Uncoordinated Gait}

\medterm{Atelectasis} Partial or complete collapse of a region of lung, reducing the amount of air it contains. It may result from blocked airways, pressure from outside the lung, shallow breathing, or scarring and can reduce oxygen exchange.

\synonyms
Lung Collapse

\medterm{Athelia} Absence of one or both nipples, present from birth or acquired after surgery, injury, or disease. It may occur alone or with other differences in breast or chest-wall development.

\medterm{Atherectomy} A minimally invasive endovascular catheter technique utilizing directional, rotational, orbital, or laser devices to physically excise, shave, vaporize, or pulverize calcified or atheromatous plaque from the intraluminal wall of diseased arteries.

\synonyms
Endovascular Atherectomy; Directional Atherectomy; Rotational Atherectomy; Orbital Atherectomy

\medterm{Atheroembolic Renal Disease} Kidney injury caused when cholesterol crystals break from an artery plaque and block small kidney vessels, often after a vascular procedure. It may cause worsening kidney function, mottled skin, painful toes, or raised eosinophils.

\medterm{Atheroma} A fatty, inflamed buildup within an artery wall that forms part of an atherosclerotic plaque. It may narrow the artery or rupture and trigger a blood clot.

\synonyms
Atheromatous Plaque; Arterial Plaque

\medterm{Atheromatous Plaque} An asymmetrical, fibrofatty lesion developing within the tunica intima of medium and large arteries, comprising a necrotic lipid-rich core of oxidized cholesterol crystals and apoptotic debris covered by a fibrous cap of collagen and vascular smooth muscle cells.

\synonyms
Atheroma; Atherosclerotic Plaque

\medterm{Atherosclerosis} A disease in which fatty, cholesterol-rich plaques build up inside artery walls. Plaques can narrow arteries or rupture and form clots, reducing blood flow to the heart, brain, legs, or other organs.

\medterm{Atherosclerosis Pathology} A chronic arterial intimal disease characterized by lipid-rich plaques with a fibrous cap, inflammation, smooth-muscle proliferation, calcification, and possible ulceration or thrombosis.

\medterm{Atherosclerotic} Relating to fatty plaque buildup and hardening inside artery walls. Atherosclerotic changes can narrow an artery or encourage a clot to form.

\medterm{Atherosclerotic Heart Disease} A disease in which cholesterol-containing plaque builds up inside the heart’s arteries, narrowing them and reducing blood flow. It may cause chest pain, heart attack, or other heart problems.

\medterm{Atherosclerotic Plaque} A buildup of cholesterol, fat, inflammatory cells, and scar-like tissue inside an artery wall. It can gradually narrow the artery or rupture and cause a blood clot.

\medterm{Atherothrombotic Stroke} An ischemic stroke caused by a clot forming on fatty plaque in a large artery supplying the brain or travelling from that plaque into a brain artery. It deprives brain tissue of oxygen.

\synonyms
Large-Artery Atherosclerotic Stroke; Atherothrombotic Ischemic Stroke

\medterm{Athetosis} Slow, involuntary, writhing movements that may affect the hands, fingers, feet, face, neck, or trunk. They result from abnormal function in brain areas that help control movement and may occur with inherited or acquired brain injury.

\medterm{Athlete Foot} Alternate name; see \textit{Athlete's Foot}.

\medterm{Athlete's Foot} A common superficial fungal infection (dermatophytosis) of the skin of the feet, most commonly caused by \textit{Trichophyton rubrum} or \textit{Trichophyton interdigitale}. Classically presents with pruritus, erythema, scaling, maceration, and fissuring, especially in the interdigital web spaces.

\synonyms
Tinea Pedis; Foot Ringworm; Foot Dermatophytosis

\medterm{Athlete's Heart} Normal physiological cardiac remodeling induced by intensive, prolonged physical endurance or strength training, manifested as physiological left ventricular hypertrophy, chamber dilation, and resting sinus bradycardia, without diastolic dysfunction.

\synonyms
Athletic Heart Syndrome; Athlete Heart Remodeling

\medterm{Atlanto-Occipital Dislocation} A severe injury in which the first neck bone separates from the base of the skull. It is extremely unstable and can damage the spinal cord or brainstem, so the neck must be immobilized and emergency care provided immediately.

\medterm{Atlanto-Occipital Joints} The paired joints between the upper surface of the atlas, the first cervical vertebra, and the occipital condyles at the base of the skull. They mainly allow nodding and small side-to-side movements while helping stabilize the head and upper neck.

\medterm{Atlantoaxial Instability} Excessive movement between the first and second cervical vertebrae, sometimes caused by trauma, rheumatoid arthritis, congenital conditions, or ligament weakness. It may cause neck pain, occipital headache, neurologic symptoms, or spinal-cord compression.

\medterm{Atlantoaxial Instability in Down Syndrome} Excessive movement between the first and second cervical vertebrae in a person with Down syndrome, often related to ligament laxity or altered bone development. Many people have no symptoms, but neck pain, weakness, gait change, or other neurologic signs require prompt assessment before activities that stress the neck.

\medterm{Atlantoaxial Joint} The complex synovial articulation between the first cervical vertebra (atlas, C1) and the second cervical vertebra (axis, C2). Comprising one median pivot joint (dens with anterior arch) and two lateral facet joints, it provides approximately 50\% of total cervical rotational range of motion. Instability may arise from trauma (odontoid fractures), rheumatoid arthritis, or Down syndrome.

\synonyms
Atlas-Axis Joint; C1-C2 Articulation

\medterm{Atlas} The first bone of the neck, supporting the skull and helping the head nod. It is ring-shaped and has openings through which important blood vessels pass.

\medterm{Atom} The smallest unit of a chemical element that retains the element’s properties. It contains a central nucleus surrounded by electrons.

\medterm{Atonic} Lacking normal muscle tone or ability to contract. For example, an atonic bladder does not squeeze effectively to empty urine.

\medterm{Atonic Seizure} A seizure causing sudden loss of muscle tone, often for only a few seconds. The person may drop their head, knees, or whole body and fall; injuries can occur during the fall.

\synonyms
Drop Attack; Akinetic Seizure

\medterm{Atopic} Having a tendency to develop allergic conditions such as eczema, allergic rhinitis, or asthma. This tendency may run in families and reflects increased sensitivity of the immune system.

\medterm{Atopic Dermatitis} A chronic, relapsing, intensely itchy inflammatory skin disease involving epidermal-barrier dysfunction and immune dysregulation. It causes dry, inflamed, scaly, or lichenified lesions whose distribution varies with age; scratching can produce infection, excoriations, and sleep disturbance.

\synonyms
Atopic Eczema; Eczema; Infantile Eczema; Allergic Dermatitis

\medterm{Atopic Eruption Of Pregnancy} An itchy eczema-like rash that begins or worsens during pregnancy, often on the face, neck, chest, or skin folds. It usually improves after delivery and is not harmful to the fetus.

\medterm{Atopic Rhinitis} Allergy-related inflammation inside the nose. Exposure to pollen, dust mites, mould, animals, or other triggers can cause sneezing, itching, a blocked or runny nose, and reduced smell.

\synonyms
Allergic Rhinitis

\medterm{Atopy} A genetic tendency to develop allergic conditions such as eczema, allergic rhinitis, asthma, or food allergy after exposure to common triggers.

\medterm{Atorvastatin} A medicine that lowers low-density lipoprotein cholesterol by reducing cholesterol production in the liver. It lowers the risk of heart attack and stroke in appropriate people; muscle symptoms and liver problems are uncommon but important side effects.

\medterm{Atovaquone–Proguanil} A combination of medicines used to prevent or treat some forms of malaria. The medicines interfere with energy production and folate use in malaria parasites; suitability depends on the type of malaria and the person’s health.

\medterm{ATP} Adenosine triphosphate, the main short-term energy-carrying molecule in cells. Cells release usable energy when ATP is broken down and remake it during metabolism.

\medterm{ATP Synthase} A protein machine in mitochondria and some other cell membranes that uses the movement of hydrogen ions to make ATP, the main molecule cells use for immediate energy.

\medterm{Atransferrinemia} A very rare inherited disorder in which the blood has little or no transferrin, the protein that carries iron. It can cause severe anemia, poor growth, and harmful iron buildup in organs.

\medterm{Atresia} A birth defect in which a normal opening or passage is absent, closed, or severely narrowed. Examples include intestinal, bile-duct, heart-valve, and ear-canal atresia; effects depend on the structure involved.

\medterm{Atresia Of Bile Ducts} See \textit{Biliary atresia}.

\medterm{Atria} The two upper chambers of the heart. The right atrium receives blood from the body and the left atrium receives blood from the lungs; each passes blood to the lower pumping chamber on its side.

\medterm{Atrial} Relating to either of the heart’s two upper chambers, which receive blood and pass it to the lower pumping chambers.

\medterm{Atrial Ablation} A procedure that uses heat, cold, or another form of energy to create small scars in heart tissue responsible for an abnormal atrial rhythm. The scars block unwanted electrical signals.

\synonyms
Cardiac Ablation

\medterm{Atrial Arrhythmia} An abnormal heart rhythm that begins in one of the heart’s upper chambers. It may cause palpitations, breathlessness, dizziness, chest discomfort, or fainting; examples include atrial fibrillation and atrial flutter.

\medterm{Atrial Fib} Alternate name; see \textit{Atrial Fibrillation}.

\medterm{Atrial Fibrillation} A supraventricular tachyarrhythmia characterized by rapid, disorganized, and chaotic electrical activation of the atria without effective mechanical atrial contraction. Electrocardiographically manifested by absent P waves and irregularly irregular QRS complexes, it predisposes to left atrial appendage thrombus formation, thromboembolic stroke, and heart failure. Management encompasses rate control, rhythm control, and systemic anticoagulation based on CHA2DS2-VASc risk scoring.

\synonyms
AF; AFib; Auricular Fibrillation

\medterm{Atrial Fibrillation Ablation} A catheter procedure that creates small scars, usually around the pulmonary veins, to block electrical signals that trigger atrial fibrillation. It may reduce symptoms, but recurrence and procedure-related risks remain possible.

\medterm{Atrial Fibrillation Medicines} Medicines used for atrial fibrillation may slow the heart rate, restore or maintain a regular rhythm, or reduce stroke risk by preventing clots. Choice depends on symptoms, heart structure, kidney function, bleeding risk, and other conditions.

\medterm{Atrial Flutter} A supraventricular arrhythmia characterized by a rapid, organized, macro-reentrant electrical circuit within the atria, typically circulating around the cavotricuspid isthmus of the right atrium. Manifested on ECG as rapid, regular sawtooth flutter waves (typically 250–350 bpm) with variable AV nodal conduction block (often 2:1), it presents with palpitations, fatigue, and elevated thromboembolic risk.

\synonyms
AFlutter; Cavotricuspid Isthmus-Dependent Flutter

\medterm{Atrial Myxoma} A usually noncancerous tumour arising inside the heart, most often in the left atrium. It may obstruct blood flow, cause breathlessness or fainting, or release fragments that block blood vessels.

\medterm{Atrial Natriuretic Peptide} A hormone released by the heart’s upper chambers when they stretch from extra blood. It tells the kidneys to remove more salt and water and helps relax blood vessels, lowering blood volume and pressure.

\medterm{Atrial Septal Defect} A hole present from birth in the wall between the heart’s upper chambers. A small hole may cause no problems, while a larger one can overwork the right side of the heart, raise lung pressure, or increase the risk of abnormal rhythms and stroke.

\medterm{Atrial Septum} The wall separating the heart’s right and left upper chambers. A hole in this wall is called an atrial septal defect.

\medterm{Atrial Standstill} Failure of the heart’s upper chambers to produce normal electrical signals or contractions. It can cause a very slow heartbeat, dizziness, fainting, poor circulation, or heart failure.

\medterm{Atrial Switch Operation} An older operation for transposition of the major heart arteries that redirects blood inside the atria using a tunnel-like passage. It restores circulation but leaves the right pumping chamber supporting the body and may cause later rhythm or heart-muscle problems.

\synonyms
Mustard Procedure; Senning Procedure

\medterm{Atrial Tachycardia} A fast heart rhythm that starts in an area of an upper chamber outside the normal pacemaker. It may cause palpitations, breathlessness, dizziness, chest discomfort, or fainting.

\medterm{Atrial Thrombosis} Formation of a blood clot inside one of the heart’s upper chambers, especially during atrial fibrillation. The clot can break loose and block blood flow to the brain or another organ.

\medterm{Atrioventricular} Relating to the connection between a heart’s upper chamber and the lower chamber on the same side. The term includes the mitral and tricuspid valves and the electrical pathway that coordinates the timing of their contractions.

\medterm{Atrioventricular Block} Delayed or interrupted electrical signals between the heart’s upper and lower chambers. It may cause a slow heartbeat, tiredness, dizziness, fainting, or no symptoms; severe block can require urgent treatment and pacing.

\medterm{Atrioventricular Bundle} A bundle of specialized electrical fibers that carries signals from the heart’s central relay point into the lower chambers. Its branches help both pumping chambers contract in a coordinated way.

\synonyms
Bundle of His

\medterm{Atrioventricular Canal Defect} A birth defect in which the wall between the heart chambers and the valves between upper and lower chambers do not form normally. It can allow abnormal blood flow and often causes heart failure or poor growth in infancy.

\synonyms
Atrioventricular Septal Defect (AVSD); Endocardial Cushion Defect

\medterm{Atrioventricular Dissociation} A cardiac rhythm abnormality in which the atria and ventricles depolarize and contract independently under the control of separate pacemakers, occurring when ventricular pacing is faster than atrial rate or during complete atrioventricular block.

\synonyms
AV Dissociation

\medterm{Atrioventricular Nodal Reentrant Tachycardia} A sudden, regular, fast heartbeat caused by an electrical loop near the heart’s central relay point. Episodes may cause pounding heartbeat, chest discomfort, breathlessness, dizziness, or fainting and may stop as suddenly as they begin.

\synonyms
AV Nodal Reentrant Tachycardia; AVNRT

\medterm{Atrioventricular Nodal Reentry Tachycardia} A rapid, regular heartbeat caused by a looping electrical signal near the atrioventricular node. It commonly begins and ends suddenly and may cause palpitations, dizziness, breathlessness, or fainting.

\medterm{Atrioventricular Node} A small area of specialized tissue near the junction of the heart’s upper and lower chambers that receives electrical signals and briefly delays them before passing them to the lower chambers. This delay allows the upper chambers to empty first.

\medterm{Atrioventricular Septal Defect} Alternate name; see \textit{Atrioventricular canal defect}.

\medterm{Atrioventricular Valves} The two heart valves between the upper and lower chambers: the tricuspid valve on the right and the mitral valve on the left. They open to let blood pass forward and close to prevent backflow.

\medterm{Atrium} One of the heart’s two upper chambers, which receive returning blood and pass it to the lower pumping chambers. The right atrium receives blood from the body; the left receives oxygen-rich blood from the lungs.

\medterm{Atrophic Glossitis} An inflammatory condition of the tongue characterized by the diffuse atrophy and loss of lingual filiform and fungiform papillae, resulting in a smooth, shiny, erythematous, and \"beefy red\" appearance often accompanied by lingual burning, pain, or altered taste. It is a key physical finding in nutritional deficiencies including vitamin B12 (Hunter's glossitis), folate, iron (Plummer-Vinson syndrome), niacin, and riboflavin, as well as systemic conditions such as celiac disease and severe protein-calorie malnutrition.
\synonyms{Hunter's Glossitis; Bald Tongue; Smooth Tongue; Moeller-Hunter Glossitis}

\medterm{Atrophic Vaginitis} Alternate name; see \textit{Vaginal Atrophy}.

\medterm{Atrophoderma Maculatum} A rare skin condition causing flat, slightly sunken patches in which the skin becomes thinner and changes colour. The patches usually have no open sore or active infection.

\medterm{Atrophy} Shrinkage and loss of function of a tissue, organ, or muscle because its cells become smaller or fewer. Causes include disuse, nerve damage, poor blood supply, malnutrition, hormone deficiency, chronic disease, and ageing; it may be partly reversible.

\medterm{Atrophy, Vaginal} Alternate index label; see \textit{Vaginal Atrophy}.

\medterm{Atropine} A medicine that blocks acetylcholine signals at certain nerve endings. It can increase heart rate, reduce secretions, widen the pupils, and counteract selected poisonings or dangerously slow heartbeats.

\medterm{Atropine In Organophosphate Poisoning} A medicine used in organophosphate poisoning to block excessive acetylcholine effects, such as wet secretions, narrowed airways, slow heartbeat, vomiting, and diarrhoea. Severe poisoning requires emergency treatment and monitoring.

\medterm{Attachment Disorder} See \textit{Reactive Attachment Disorder}.

\medterm{Attachment, Dental} A precision mechanical or magnetic interlocking device used to secure, stabilize, and retain a removable partial denture or overdenture to retained tooth roots, natural abutments, or dental implants.

\synonyms
Denture Attachment; Precision Attachment; Overdenture Attachment; Dental Attachment

\medterm{Attack Rate} The proportion of people at risk who develop a particular illness during a defined outbreak or exposure period. It is commonly used to describe foodborne outbreaks and other limited events.

\medterm{Attention} The ability to select information, focus on it, and maintain that focus while limiting distraction. Attention can be affected by sleep loss, stress, illness, medicines, or brain disorders.

\medterm{Attention Deficit Hyperactivity Disorder} Alternate index label; see \textit{Attention-Deficit/Hyperactivity Disorder}.

\medterm{Attention-Deficit/Hyperactivity Disorder} A neurodevelopmental disorder involving persistent inattention, hyperactivity, and impulsive behaviour that affects functioning or development. Presentations may be mainly inattentive, mainly hyperactive-impulsive, or combined.

\synonyms
ADHD; Hyperkinetic Disorder

\medterm{Attenuated} Weakened or reduced in force, severity, concentration, or biological activity. An attenuated vaccine contains a weakened form of a microorganism that is intended to stimulate immunity without usually causing disease.

\medterm{Attenuation Correction} A computer adjustment for the way body tissues weaken signals during medical imaging. It helps make PET, SPECT, and combined scan images more accurate by compensating for signal loss.

\medterm{Attributable Risk} The difference in disease risk between people exposed and not exposed to a possible cause. It estimates the extra risk associated with the exposure, but does not by itself prove causation.

\medterm{Atypia} An unusual appearance of cells or tissue seen under a microscope. It may result from irritation, inflammation, ageing, a precancerous change, or cancer and must be interpreted with the tissue pattern and clinical context.

\medterm{Atypia Morphology} An unusual appearance or arrangement of cells or tissue under a microscope. It may result from irritation or inflammation, or represent a precancerous or cancerous change; its meaning depends on the whole sample and clinical setting.

\synonyms
Cytologic Atypia

\medterm{Atypical} Different from the usual appearance, behaviour, structure, or course of a condition. Atypical does not by itself mean that a finding is harmless, precancerous, or cancerous.

\medterm{Atypical (Dysplastic) Nevus} A melanocytic mole with unusual clinical or microscopic features such as asymmetry, irregular border, varied color, or larger size. Most are benign, but multiple atypical nevi increase melanoma risk; a changing, bleeding, or distinctly different lesion should be examined and may require biopsy.

\medterm{Atypical Antipsychotics} A group of medicines used mainly for schizophrenia, bipolar disorder, and selected other conditions. They change several brain signals and can cause sleepiness, weight gain, movement problems, high blood sugar, or other metabolic changes.

\medterm{Atypical Endometrial Hyperplasia} An overgrowth of abnormal cells in the lining of the uterus. It is not cancer, but it can develop into uterine-lining cancer or occur beside cancer, so medical assessment and follow-up are important.

\medterm{Atypical Facial Pain} Persistent or recurring pain in the face that does not fit a clear nerve, dental, sinus, or other recognized cause. Evaluation looks for treatable causes before considering persistent facial-pain disorders.

\medterm{Atypical Features} A pattern sometimes used to describe depression in which mood can improve temporarily in response to positive events, with increased sleep, increased appetite or weight, heavy-feeling limbs, and marked sensitivity to rejection.

\medterm{Atypical Fibroxanthoma} A low-grade cutaneous sarcoma that usually arises as a rapidly growing red, pink, or flesh-coloured nodule on sun-damaged head or neck skin of an older adult. Diagnosis requires biopsy because it can mimic squamous-cell carcinoma or melanoma; complete excision or Mohs surgery is typical treatment.

\medterm{Atypical Genitalia} External genital anatomy that is not typically classified as male or female. It may result from differences in chromosomes, hormones, gonads, or genital development and requires respectful specialist assessment.

\medterm{Atypical Hemolytic Uremic Syndrome} A rare disease in which inherited or acquired abnormalities of the complement system cause uncontrolled immune activity and damage to small blood vessels, especially in the kidneys. It can destroy red blood cells, lower platelets, and cause sudden kidney injury, high blood pressure, or neurological problems. Unlike typical hemolytic uremic syndrome, it is not primarily caused by a toxin-producing diarrheal infection.

\synonyms
aHUS; Complement-Mediated HUS

\medterm{Atypical Hyperplasia Of The Breast} An overgrowth of unusual cells in a breast duct or lobule that is not cancer. It can increase the future risk of breast cancer and may require closer screening or preventive discussion.

\medterm{Atypical Lymphocyte} A lymphocyte, a type of white blood cell, with an unusual size or appearance under a microscope. It often reflects immune activity during a viral infection, but the finding alone does not identify a specific disease.

\medterm{Atypical Mole} A mole with unusual clinical or microscopic features, also called a dysplastic nevus. Most are noncancerous, but multiple atypical moles can increase melanoma risk; a changing, bleeding, or distinctly different mole needs examination.

\medterm{Atypical Nevi} Alternate plural name; see \textit{Dysplastic Nevi}.

\medterm{Atypical Pneumonia} A lung infection that often develops gradually and may cause more general symptoms than chest symptoms, such as headache, tiredness, sore throat, or muscle aches. Causes include certain bacteria and viruses.

\medterm{Atypical Spitz Nevus} An unusual mole made of pigment-producing cells that resembles a Spitz nevus but has concerning features. It can resemble melanoma under a microscope, so expert skin examination and tissue review may be needed.

\medterm{Atypical Teratoid Rhabdoid Tumor} A rare, fast-growing cancer of the brain or spinal cord that occurs mainly in infants and young children. It may cause vomiting, headache, sleepiness, balance problems, seizures, or developmental changes.

\synonyms
ATRT; Rhabdoid Brain Tumor

\medterm{Audiogram} A chart that shows a person’s ability to hear at different pitches or frequencies.

\medterm{Audiologist} A health professional who assesses hearing and fits hearing aids.

\medterm{Audiology} The healthcare field concerned with hearing, balance, and related communication problems. Audiologists test hearing and balance, provide hearing aids and other devices, offer training and support, and identify when medical or surgical care may be needed.

\medterm{Audiometry} Measurement of hearing sensitivity at different sound frequencies and intensities, usually using headphones or insert earphones.

\medterm{Auditory Acuity} The sensitivity or sharpness of hearing, commonly assessed by determining the softest sounds a person can detect at different frequencies.

\medterm{Auditory Area} A region of the cerebral cortex that processes sound information.

\synonyms
Auditory Cortex

\medterm{Auditory Fatigue} A temporary decrease in hearing sensitivity or listening ability after prolonged loud sound exposure; rest usually helps, but repeated exposure can cause lasting damage.

\medterm{Auditory Hallucinations} False sensory perceptions of sound occurring in the absence of any external acoustic stimulus. They may be non-verbal (such as clicks, ringing, or music) or verbal (such as single words, conversations, running commentary, or command voices). Verbal auditory hallucinations are a hallmark feature of psychotic disorders such as schizophrenia, but can also occur in severe mood disorders with psychotic features, delirium, temporal lobe epilepsy, sensory deprivation, and substance withdrawal.
\synonyms{Paracusia; Phonemes; Auditory Sensory Hallucinations}

\medterm{Auditory Information Processing Disorder} See \textit{Auditory processing disorder}.

\medterm{Auditory Nerve} The portion of the vestibulocochlear nerve that carries hearing signals from the inner ear to the brain.

\medterm{Auditory Neuropathy Spectrum Disorder} A form of sensorineural hearing impairment characterized by preserved cochlear outer hair cell function (demonstrated by normal otoacoustic emissions or cochlear microphonics) but abnormal, dyssynchronous neural transmission along the eighth cranial nerve.

\synonyms
ANSD; Auditory Neuropathy

\medterm{Auditory Ossicle} One of three small middle-ear bones—the malleus, incus, and stapes—that transmit and amplify vibrations from the eardrum to the inner ear. Disease, injury, or fixation of these bones can cause conductive hearing loss.

\synonyms
Ear Ossicle

\medterm{Auditory Pathway} A chain of nerve structures carrying sound signals from the inner ear through the brainstem to the auditory cortex for recognition and interpretation.

\medterm{Auditory Perception} The brain's process of detecting, recognizing, and interpreting sounds received through the auditory system.

\medterm{Auditory Processing Disorder} Difficulty recognizing or interpreting spoken sounds even though basic hearing sensitivity may be normal. Affected people may struggle to understand speech in noise, follow rapid instructions, or distinguish similar sounds; assessment requires specialized testing and management may include communication strategies and support.

\synonyms
APD

\medterm{Auditory Threshold} The quietest sound a person can hear at a particular pitch or frequency. It is measured during a hearing test and helps describe the degree and type of hearing loss in each ear.

\synonyms
Hearing Threshold

\medterm{Auer Rods} Needle-shaped structures sometimes seen inside immature blood cells under a microscope. They are made from fused cell granules and strongly suggest a cancer of immature myeloid blood cells, especially acute myeloid leukemia.

\medterm{Auerbach Plexus} An extensive intrinsic neural network of unmyelinated nerve fibers and postganglionic parasympathetic ganglia situated in the muscularis externa between the circular and longitudinal smooth muscle layers throughout the gastrointestinal tract, coordinating peristalsis and motor activity.

\synonyms
Myenteric Plexus

\medterm{Auger Electrons} Very low-energy electrons released by some radioactive atoms after a change inside the atom. They travel only a tiny distance and can damage nearby DNA, a property being studied for precisely targeted radiation treatment.

\medterm{Augmentation Mammaplasty} An operation to increase breast size or change breast shape, usually with implants or transferred fat. Possible risks include infection, bleeding, changes in sensation, scarring, and the need for further surgery.

\medterm{Augmentation Mammoplasty} An operation to increase breast size or change breast shape, usually with implants or transferred fat. Possible risks include infection, bleeding, changes in sensation, scarring, and the need for further surgery.
\synonyms
Breast Augmentation

\medterm{Augmentation Of Labor} Treatment used when labour has started but contractions are not strong or frequent enough for the cervix to continue opening. It may involve oxytocin or deliberately breaking the waters, with monitoring of parent and baby.

\medterm{Aura} A temporary change in vision, sensation, speech, movement, or awareness that may occur before or during a migraine or seizure. Symptoms usually develop gradually and then resolve, but a new or prolonged aura needs urgent assessment.

\medterm{Aural Fullness} A feeling of pressure, blockage, or fullness in the ear. It can result from poor pressure equalization, fluid, earwax, inner-ear disease, or jaw-joint problems and may occur with hearing loss, ringing, or dizziness.

\synonyms
Ear Pressure; Plugged-Ear Sensation; Blocked-Ear Sensation

\medterm{Aural Polyps} Soft, abnormal tissue growths in the ear canal or middle ear. They often result from long-lasting infection or inflammation, but persistent discharge, pain, or bleeding needs specialist examination to exclude a more serious cause.

\medterm{Auricle} The visible outer part of the ear attached to the side of the head. Its ridges collect and direct sound into the ear canal, helping the brain judge where sounds come from.

\medterm{Auricular Hematoma} A collection of blood between the cartilage and skin covering of the outer ear, usually after a blow or sports injury. Prompt drainage and compression help prevent permanent thickening and deformity of the ear.

\synonyms
Pinna Hematoma; Subperichondrial Auricular Hematoma

\medterm{Auscultation} Listening to sounds produced inside the body, usually with a stethoscope. Clinicians use it to assess the heart, lungs, abdomen, and blood vessels for findings such as murmurs, wheezing, crackles, or abnormal blood-flow sounds.

\medterm{Auscultatory Gap} A temporary disappearance of the sounds heard while a blood-pressure cuff is slowly deflated, followed by their return. If it is missed, the top blood-pressure number may be recorded falsely low.

\medterm{Auspitz Sign} A skin check for psoriasis where tiny pinpoint bleeding spots appear when the silvery scales of a rash are gently scraped away.

\medterm{Austin Flint Murmur} A low-pitched sound heard during the middle part of the heartbeat over the heart's apex. It is associated with severe leakage of the aortic valve and altered filling of the mitral valve.

\medterm{Autism} Alternate name; see \textit{Autism Spectrum Disorder}.

\medterm{Autism Spectrum Disorder} A lifelong neurodevelopmental condition involving persistent differences in social communication and interaction, along with restricted or repetitive behaviors, interests, or sensory patterns. Features begin during the developmental period, vary widely among individuals, and may affect communication, learning, relationships, daily activities, or the type and amount of support needed.

\synonyms
ASD; Autistic Spectrum Disorder; Pervasive Developmental Disorder

\medterm{Autism Spectrum Disorder, Childhood Disintegrative Disorder} Alternate index label; see \textit{Childhood Disintegrative Disorder}.

\medterm{Autistic Savant} An autistic person with an unusually strong, specific ability—such as music, art, calculation, or memory—compared with their other abilities. The term describes a pattern of skills and is not itself a medical diagnosis.

\medterm{Autoantibodies} Immune proteins that mistakenly recognize and attach to the body's own cells, tissues, or proteins. They may cause disease or serve as test clues; their presence alone does not prove that a person has an autoimmune disease.

\medterm{Autoantibody} An immune protein that mistakenly recognizes and binds a substance from the person's own body. It may damage tissues or interfere with their function, or it may be present only as a laboratory sign without causing disease.

\medterm{Autoantigen} A substance normally found in the body that is mistakenly recognized as foreign by the immune system. This mistaken response can lead to inflammation or damage to the person's own cells or tissues.

\medterm{Autoclave} A machine that sterilizes equipment by exposing it to pressurized steam at a controlled temperature and time. It destroys bacteria, viruses, fungi, and spores when the correct cycle is used.

\medterm{Autocrine Signaling} A mode of intercellular and intracellular communication in which a cell synthesizes and secretes a messenger molecule (such as a cytokine, hormone, or growth factor) that binds to surface receptors on the very same cell that produced it, initiating a feedback response.

\medterm{Autoerythrocyte Sensitivity} A rare condition causing painful bruises or raised purple patches, thought to involve an abnormal reaction to a person's own red-cell components. Other bleeding disorders, medicines, and causes of bruising must be excluded first.

\medterm{Autogenic Drainage} A controlled-breathing method that moves mucus from the smaller airways toward the throat so it can be coughed out. It uses different breathing depths and requires instruction and practice.

\medterm{Autograft} Tissue moved from one part of a person's body to another part of the same body. Because the tissue is genetically the person's own, rejection is unlikely, but the original site may have pain, scarring, or other complications.

\medterm{Autoimmune} Describing an immune response that mistakenly targets the body's own cells, tissues, or proteins. It can cause inflammation, tissue damage, or impaired organ function.

\medterm{Autoimmune Associated Epilepsy} See \textit{Autoimmune epilepsy}.

\medterm{Autoimmune Autonomic Ganglionopathy} A rare, severe autoimmune disorder of the autonomic nervous system characterized by widespread sympathetic and parasympathetic failure. Manifests clinically with severe orthostatic hypotension, anhidrosis, dry mouth and eyes, tonic pupils, gastrointestinal dysmotility (gastroparesis, intestinal pseudo-obstruction), and urinary retention; strongly associated with autoantibodies directed against the ganglionic nicotinic acetylcholine receptor ($\alpha_3\beta_4$ subtype).

\synonyms
AAG; Subacute Autonomic Neuropathy

\medterm{Autoimmune Disease} A disease in which the immune system mistakenly attacks the body's own cells, tissues, or organs. It may cause inflammation, pain, tissue damage, or loss of organ function.

\medterm{Autoimmune Diseases} Diseases caused by immune responses against the body's own tissues. They may affect one organ, such as the thyroid, or many parts of the body and can produce inflammation and lasting damage.

\medterm{Autoimmune Disorder} A condition in which the immune system mistakenly attacks the body's own cells, tissues, or organs. It may be limited to one organ or affect many organs, causing inflammation and impaired function.

\medterm{Autoimmune Encephalitis} Inflammation of the brain caused by an immune attack on nerve cells or their proteins. It may cause seizures, confusion, memory problems, unusual movements, behaviour changes, sleep problems, or reduced consciousness.

\medterm{Autoimmune Encephalopathy} See \textit{Autoimmune Encephalitis}.

\medterm{Autoimmune Endocrine Disease} An autoimmune disorder that damages a hormone-producing gland or hormone-producing cells. Examples include autoimmune thyroid disease, Addison disease, and type 1 diabetes; it may cause too little or, less often, abnormal hormone production.

\medterm{Autoimmune Enteropathy} A rare disease in which the immune system damages the lining of the small intestine, causing ongoing diarrhoea, poor nutrient absorption, weight loss, or growth problems. It may occur with other immune disorders.

\medterm{Autoimmune Epilepsy} Recurrent seizures caused by an immune attack on the brain. Seizures may begin suddenly or be difficult to control and can occur with memory problems, confusion, movement changes, or other signs of brain inflammation.

\medterm{Autoimmune Gastritis} Long-term inflammation that damages acid- and vitamin-B12-producing cells in the stomach. It can cause low stomach acid, iron deficiency, vitamin B12 deficiency, anaemia, digestive symptoms, and an increased risk of certain stomach growths.

\medterm{Autoimmune Hemolytic Anemia} Anaemia caused when immune proteins attach to and destroy red blood cells faster than the body can replace them. It may cause tiredness, pale or yellow skin, dark urine, rapid heartbeat, or shortness of breath.

\medterm{Autoimmune Hepatitis} Long-term liver inflammation caused by an immune attack on liver cells. It may cause tiredness, nausea, joint pain, or jaundice, while some people have no symptoms; untreated disease can lead to scarring and liver failure.

\medterm{Autoimmune Keratitis} Inflammation of the clear front surface of the eye caused by or associated with immune disease. It may cause eye pain, redness, light sensitivity, blurred vision, or a corneal scar and needs prompt eye assessment.

\medterm{Autoimmune Liver Disease Panel} A group of blood tests that looks for immune proteins linked with autoimmune liver diseases. Results are interpreted with liver tests, symptoms, examination, and sometimes a liver biopsy; the panel alone cannot establish a diagnosis.

\medterm{Autoimmune Lymphoproliferative Syndrome} An inherited immune disorder in which lymphocytes do not die off normally after an immune response, causing noncancerous enlargement of lymph nodes and spleen from early childhood. Autoimmune destruction of blood cells may cause anemia, low platelets, or low white-cell counts. It is often due to a defect in the Fas cell-death pathway; diagnosis uses blood counts, specialized immune-cell testing, and gene testing, and there is an increased risk of lymphoma.

\medterm{Autoimmune Pancreatitis} A rare form of pancreas inflammation caused by an abnormal immune response. It may cause painless jaundice, abdominal pain, weight loss, or a mass that resembles cancer and often improves with immune-suppressing treatment.

\medterm{Autoimmune Polyendocrine Syndromes} A group of inherited or acquired disorders in which immune attacks affect several hormone-producing glands. Depending on the type, features may include adrenal failure, low parathyroid hormone, thyroid disease, diabetes, chronic fungal infection, or skin changes.

\medterm{Autoimmune Polyendocrinopathy} A disorder in which immune attacks affect several hormone-producing glands. One inherited form can cause chronic fungal infections of skin or mouth, low parathyroid hormone, adrenal failure, and other autoimmune conditions.

\medterm{Autoimmune Polyglandular Syndrome} A group of disorders in which the immune system damages two or more hormone-producing glands. Possible effects include adrenal failure, thyroid disease, low parathyroid hormone, diabetes, or problems with reproductive hormones.

\medterm{Autoimmune Thyroid Disease and Pregnancy} Thyroid autoimmunity can cause hypothyroidism, hyperthyroidism, pregnancy loss, preterm birth, or fetal thyroid effects. Thyroid function and medicine doses are monitored closely before and during pregnancy and after delivery.

\medterm{Autoimmune Thyroiditis} Immune-mediated inflammation and damage of the thyroid gland. See \textit{Hashimoto Thyroiditis}.

\medterm{Autoimmunity} An immune response directed against the body's own cells, tissues, or proteins. It may cause no harm, or it may produce inflammation and disease when it damages tissues or interferes with organ function.

\medterm{Autoimmunity: Mechanisms} Autoimmunity can develop when the immune system fails to remove or control self-reactive cells, or when an infection triggers responses that also recognize body tissues. Genetic factors, immune regulation, and environmental triggers may all contribute.

\medterm{Autoinflammatory Diseases} Conditions in which the body's early, nonspecific immune response becomes overactive and causes repeated or ongoing inflammation. They may cause fever, rash, joint pain, abdominal pain, or organ damage, often without disease-specific autoantibodies.

\medterm{Autoinoculation} The spread of an infection from one part of a person's body to another, often by touching or scratching. This can transfer germs to broken skin, the eyes, the mouth, or another vulnerable site.

\medterm{Autologous} Taken from and returned to the same person. Examples include storing a person's own blood for later transfusion or collecting and returning their own stem cells during treatment.

\medterm{Autologous Blood Donation} Collection and storage of a person's own blood before an operation for possible transfusion during or after surgery. It avoids incompatibility with donor blood but may not be suitable or necessary for everyone.

\medterm{Autologous Breast Reconstruction} Surgical reconstruction of the breast contour utilizing the patient's own vascularized tissue flaps harvested from donor sites such as the abdomen (e.g., deep inferior epigastric perforator [DIEP] or transverse rectus abdominis myocutaneous [TRAM] flap), back (latissimus dorsi flap), or gluteal/thigh regions.

\synonyms
Autologous Breast Reconstruction; Flap Breast Reconstruction; Autogenous Tissue Reconstruction; DIEP Flap Reconstruction

\medterm{Autologous Fat Transplant} A procedure that removes fat from one area of a person's body and injects it into another area to restore volume or change shape. Some of the transferred fat may be reabsorbed, and repeat procedures may be needed.

\synonyms
Autologous Fat Grafting; Fat Transfer; Lipofilling

\medterm{Autologous Stem Cell Transplantation} A clinical procedure in which a patient's own hematopoietic stem cells are mobilized, harvested from peripheral blood or bone marrow, cryopreserved, and reinfused following high-dose myeloablative chemotherapy or radiation to rescue bone marrow hematopoiesis.

\synonyms
Autologous Bone Marrow Transplant; Auto-SCT

\medterm{Autolysis} The breakdown of cells or tissues by their own enzymes, usually after the cells die. This natural process contributes to tissue decomposition and can affect specimens examined after death.

\medterm{Automated External Defibrillator} A portable device that checks the heart rhythm during cardiac arrest and gives an electrical shock when an abnormal rhythm may respond to one. It provides spoken or visual instructions for using it safely.

\medterm{Automaticity} The physiological capacity of specialized cardiac pacemaker myocytes (principally in the sinoatrial node, atrioventricular junction, and His-Purkinje system) to undergo spontaneous diastolic phase 4 depolarization and generate action potentials without external neurohumoral stimulation.

\medterm{Autonomic Disorder} A problem affecting nerves that control automatic body functions such as blood pressure, heart rate, digestion, sweating, bladder emptying, sexual function, and pupil size. Causes include diabetes, nerve disease, infection, immune disease, and medicines; symptoms depend on the affected functions.

\medterm{Autonomic Dysreflexia} A sudden, potentially life-threatening rise in blood pressure in a person with a spinal-cord injury, usually at or above the upper chest. A problem such as a full bladder or bowel below the injury can trigger pounding headache, sweating, flushing, or heart-rate changes.

\medterm{Autonomic Ganglion} A small group of nerve-cell bodies outside the brain and spinal cord that passes signals between the autonomic nervous system and organs such as the heart, intestines, glands, and bladder.

\medterm{Autonomic Nervous System} The part of the nervous system that controls body functions that usually happen without conscious effort, including heart rate, blood pressure, digestion, sweating, pupil size, bladder activity, and sexual responses.

\medterm{Autonomic Neuropathy} Damage to nerves that control automatic functions such as blood pressure, heart rate, digestion, sweating, bladder emptying, and sexual function. Diabetes is common, but other diseases, toxins, and medicines can also cause it.

\synonyms
Dysautonomia; Autonomic Nerve Dysfunction

\medterm{Autonomic Neuropathy, Diabetic} Alternate index label; see \textit{Diabetic Autonomic Neuropathy}.

\medterm{Autonomy, Patient} A person's right and ability to make informed, voluntary decisions about medical care, including accepting or refusing tests or treatment. It requires understandable information and freedom from pressure.

\medterm{Autophagy} A normal process in which cells break down and recycle worn-out proteins and cell parts. It helps cells survive stress and remove damaged material, but abnormal autophagy may contribute to some cancers and other diseases.

\medterm{Autopsy} A medical examination of a body after death. It may include external examination, examination of internal organs, tissue testing, and sometimes toxicology to help determine the cause of death or study disease.

\synonyms
Post-Mortem Examination; Necropsy; Obduction

\medterm{Autopsy, Medicolegal} A specialized post-mortem examination performed under statutory or judicial authority by a forensic pathologist or medical examiner to determine the cause and manner of violent, sudden, suspicious, unnatural, or unattended deaths and preserve legal evidence.

\synonyms
Forensic Autopsy; Medico-Legal Autopsy; Coroner Post-Mortem

\medterm{Autoradiography} A laboratory method that uses a radioactive label and a photographic or digital detector to show where a substance is located in a tissue, cell, or biological sample.

\medterm{Autosomal Dominant} A pattern of inheritance in which one altered copy of a gene on a non-sex chromosome can cause a condition. A person with one altered copy may pass it to about half of their children.

\medterm{Autosomal Dominant Inheritance} Inheritance in which one altered copy of a gene on a non-sex chromosome can cause a condition. Each child of a parent with one altered copy commonly has a one-in-two chance of inheriting it.

\medterm{Autosomal Dominant PKD} An inherited disease in which many fluid-filled cysts gradually develop in the kidneys and may also occur in the liver. Cysts can cause high blood pressure, pain, blood in the urine, infections, and progressive kidney failure.

\medterm{Autosomal Dominant Polycystic Kidney Disease} An autosomal dominant inherited disease, usually caused by a change in the \textit{PKD1} or \textit{PKD2} gene, in which many cysts gradually enlarge the kidneys and can reduce kidney function. It may cause high blood pressure, pain, blood in the urine, kidney infections, liver cysts, or rarely brain-artery aneurysms.

\synonyms
ADPKD

\medterm{Autosomal Dominant Tubulointerstitial Kidney Disease} An inherited disease that gradually damages the kidney's small tubules and surrounding tissue. Kidney function declines, often without much protein or blood in the urine; some forms cause early gout, low blood pressure, or anaemia.

\medterm{Autosomal Inheritance} The passing of a genetic trait through one of the 22 pairs of chromosomes that are not sex chromosomes. Such traits can affect people of any sex and may follow dominant or recessive patterns.

\medterm{Autosomal Recessive} A pattern of inheritance in which a condition usually develops only when both copies of a gene are altered. Parents with one altered copy are often healthy carriers but can pass the condition to a child.

\medterm{Autosomal Recessive Disorders} Conditions that usually occur when a person inherits an altered copy of the same gene from each parent. People with one altered copy are often healthy carriers.

\medterm{Autosomal Recessive Inheritance} Inheritance in which a condition usually requires altered copies of a gene from both parents. When both parents are carriers, each pregnancy commonly has a one-in-four chance of producing an affected child.

\medterm{Autosomal Recessive PKD} An inherited disease that usually causes enlarged, poorly functioning kidneys in babies or children and scarring of the liver. Severe disease can interfere with lung development before birth; milder forms may appear later.

\medterm{Autosomal Recessive Polycystic Kidney Disease} An inherited disease that usually affects babies or children, causing enlarged kidneys with small cysts and progressive kidney problems. It can also cause liver scarring, high blood pressure, and breathing problems in severe early disease.

\synonyms
ARPKD

\medterm{Autosome} Any chromosome that is not a sex chromosome. Humans usually have 22 pairs of autosomes. Genes on these chromosomes can cause inherited conditions in people of any sex.

\synonyms
Non-Sex Chromosome

\medterm{Avascular Necrosis} Death of bone cellular components resulting from compromised subchondral microvascular blood supply, leading to bone trabecular death, structural collapse, and progressive secondary osteoarthritis. It most commonly affects the femoral head, scaphoid, talus, and humeral head; major etiologies include trauma/fractures, chronic corticosteroid therapy, excessive alcohol use, and sickle cell disease.

\synonyms
Osteonecrosis; Aseptic Necrosis; Ischemic Bone Necrosis; Femoral Head Avascular Necrosis; AVN

\medterm{Avastin} The brand name for bevacizumab, a monoclonal antibody that blocks vascular endothelial growth factor (VEGF). It is used for several cancers and, in ophthalmology, by injection into the eye for selected retinal conditions; risks include high blood pressure, bleeding, clotting, impaired wound healing, and rare gastrointestinal perforation.

\medterm{Avian Flu} Alternate name; see \textit{Avian Influenza}.

\medterm{Avian Influenza} An infection caused by avian-adapted influenza A virus strains (such as H5N1, H5N6, or H7N9) transmitted to humans primarily via direct contact with infected poultry or contaminated environments. Clinical manifestations range from conjunctivitis and upper respiratory symptoms to fulminant viral pneumonia, acute respiratory distress syndrome (ARDS), multi-organ failure, and death.

\synonyms
Bird Flu; Avian Flu; Highly Pathogenic Avian Influenza; HPAI

\medterm{AVN} Alternate name; see \textit{Avascular Necrosis}.

\medterm{Avoidance Behavior} Actions taken to stay away from a feared or distressing situation, thought, place, person, or activity. Avoidance may briefly reduce anxiety but can keep the fear going and interfere with treatment, work, relationships, or daily life.

\medterm{Avoidant Personality Disorder} A long-term pattern of avoiding social contact because of strong feelings of inadequacy and fear of criticism, embarrassment, or rejection. It can limit relationships and work; psychological treatment can help.

\medterm{Avoidant/Restrictive Food Intake Disorder} An eating disorder characterized by persistent failure to meet appropriate nutritional or energy needs, resulting in significant weight loss, nutritional deficiency, dependence on enteral feeding/supplements, or marked interference with psychosocial functioning. Unlike anorexia nervosa, the restriction is driven by sensory sensitivity, lack of interest in eating, or fear of aversive consequences (e.g., choking or vomiting), rather than body-image disturbance.

\synonyms
ARFID; Selective Eating Disorder

\medterm{Avolition} A marked reduction in motivation or ability to begin and complete purposeful activities. A person may neglect hygiene, household tasks, work, or social contact. It can occur in schizophrenia, severe depression, and other conditions.

\medterm{Avulsed Tooth} A tooth completely knocked out of its socket by an injury. It is a dental emergency; prompt treatment may allow the tooth to be replaced in its socket if it is handled by the crown and kept moist.

\medterm{Avulsion} An injury in which tissue is partly or completely torn away by force. It may involve skin, muscle, a nail, a fingertip, or a tooth and can cause severe bleeding, tissue loss, infection, and loss of function.

\medterm{Avulsion Fracture} A break in which a tendon or ligament pulls off a small piece of attached bone, usually during a sudden movement or injury. Treatment depends on the bone, fragment position, joint involvement, age, and associated damage.

\medterm{Avuncular} Relating to an uncle. In medical family-history or genetic writing, it describes the relationship between a person and their parent's brother or sister's husband.

\medterm{Awake Craniotomy} A specialized neurosurgical procedure in which the patient is deliberately kept awake, alert, and responsive during portions of open brain surgery. Performed primarily to remove tumors (such as gliomas) or epileptic seizure foci located within or adjacent to critical functional brain areas (the "eloquent cortex," which controls speech, language, motor function, and sensory perception), the operation begins under general anesthesia or conscious sedation while the skull bone flap is temporarily opened. The neurosurgical team then awakens the patient and applies gentle, low-voltage electrical stimulation to the exposed brain surface (cortical mapping) while a speech-language pathologist or neuropsychologist asks the patient to name pictures, read words, or move fingers and toes. If stimulating a particular brain spot temporarily stops speech or movement, surgeons know to spare that exact pathway. This real-time functional mapping maximizes tumor resection while safeguarding essential neurological abilities, reducing post-operative paralysis or loss of speech.
\synonyms
Awake Brain Surgery; Intraoperative Brain Mapping; Functional Craniotomy

\medterm{Awake Fiberoptic Intubation} A method of placing a breathing tube while the person remains awake and breathes independently. Numbing medicine and a flexible camera are used to guide the tube when a difficult airway is expected.

\medterm{Awareness} Conscious recognition of oneself, surroundings, thoughts, or sensations. Clinicians assess it along with alertness, attention, memory, and orientation when examining brain function.
\synonyms
Consciousness

\medterm{Awareness Under Anesthesia} Unintended awareness or memory of events during general anaesthesia. A person may hear sounds, feel pressure or pain, or recall part of the operation; although uncommon, it can cause significant distress and needs support.

\medterm{Axial Skeleton} The bones forming the central framework of the body: the skull, spine, ribs, breastbone, and associated small bones of the neck. It supports the body and protects the brain, spinal cord, heart, and lungs.

\medterm{Axial Spondyloarthritis} Long-term inflammation affecting the joints between the spine and pelvis and other parts of the spine. It commonly causes back pain and stiffness that improve with movement and may eventually limit spinal flexibility.

\medterm{Axial Tomography, Computerized} An imaging test that uses X-rays and computer processing to create detailed cross-sectional pictures of the body. It can show bones, organs, blood vessels, bleeding, tumours, and other abnormalities.

\medterm{Axilla} The hollow area beneath the shoulder, commonly called the armpit. It contains major blood vessels, nerves, lymph nodes, and tissues connecting the chest and neck with the arm.

\medterm{Axillary} Relating to the armpit or to structures within it, including the axillary blood vessels, nerves, and lymph nodes.

\synonyms
Armpit-Related; Pertaining to the Axilla

\medterm{Axillary Artery} The main artery passing through the armpit and supplying blood to the shoulder, chest wall, and arm. It continues as the brachial artery near the lower border of the armpit.

\medterm{Axillary Dissection} Surgery that removes selected lymph nodes and nearby tissue from the armpit, often for cancer assessment or treatment. It may cause numbness, shoulder stiffness, fluid collection, nerve injury, or swelling of the arm.

\synonyms
Axillary Node Dissection

\medterm{Axillary Lymph Node} A lymph node in the armpit that drains lymph from the arm, chest wall, and parts of the breast. Its size or involvement by cancer, infection, or inflammation may be assessed by examination or imaging.

\medterm{Axillary Lymphadenectomy} An operation to remove lymph nodes from the armpit, usually to assess or treat breast cancer or melanoma. Possible complications include pain, numbness, shoulder stiffness, a fluid collection, nerve injury, and arm swelling.

\medterm{Axillary Lymphadenopathy} An abnormal enlargement, induration, or matted consistency of the lymph nodes within the axilla. It receives lymphatic drainage from the ipsilateral breast, upper extremity, and chest wall. Major causes include metastatic breast carcinoma, lymphomas, melanoma, acute pyogenic infections of the upper limb, cat-scratch disease (Bartonella henselae), tuberculosis, and systemic autoimmune diseases.

\synonyms
Axillary Adenopathy; Enlarged Axillary Nodes

\medterm{Axillary Nerve} A nerve from the neck that passes around the shoulder and supplies muscles that lift the arm. It also provides feeling over the outer shoulder and may be injured by shoulder dislocation or upper-arm fracture.

\medterm{Axillary Nerve Dysfunction} Impairment of the axillary nerve, often from shoulder trauma, dislocation, compression, or surgery. It can cause weakness of shoulder abduction, deltoid wasting, and numbness over the lateral shoulder.

\medterm{Axillary Tail Of The Breast} The extension of breast tissue toward the armpit. It is also called the tail of Spence.

\medterm{Axillary Vein} A large vein in the armpit that drains blood from the arm and shoulder. It continues beneath the collarbone as the subclavian vein and returns blood toward the heart.

\medterm{Axis} The second bone in the neck. Its tooth-like projection acts as a pivot for the first neck bone, allowing the head and upper neck to rotate.

\medterm{Axis Vertebra} The second neck vertebra, whose tooth-like projection allows the first vertebra and head to rotate while protecting the spinal cord.

\synonyms
C2; Axis

\medterm{Axolemma} The thin outer membrane surrounding a nerve cell's axon. It maintains the electrical difference needed for nerve impulses and contains channels that allow those impulses to travel along the axon.

\medterm{Axon} The long extension of a nerve cell that carries electrical signals away from the cell body toward another nerve cell, muscle, or gland.

\medterm{Axon Fasciculation} The process during development in which growing nerve fibres group together to form organized nerve pathways. Chemical guidance signals and cell connections help them reach their correct targets.

\medterm{Axon Guidance} The process by which growing nerve fibres find and connect with their correct targets during development. Chemical signals and surrounding cells help direct their path.

\medterm{Axon Hillock} The cone-shaped part where a nerve cell's main body joins its axon, the long fiber that carries signals away. It combines incoming electrical signals and usually starts the nerve impulse when those signals are strong enough.

\medterm{Axon Terminal} The end of a nerve-cell axon where chemical messengers are released to communicate with another nerve cell, muscle, or gland.

\medterm{Axonal Transport} The movement of proteins, cell parts, and other materials along a nerve cell's axon between the cell body and the axon terminal. It supports nerve-cell maintenance and communication.

\medterm{Axoneme} The internal framework of tiny tubes inside a cilium or sperm tail. Together with motor proteins, it produces the bending movements that help move fluid over cells or propel sperm.

\medterm{Axonotmesis} A moderate degree of peripheral nerve injury in Seddon's classification characterized by disruption of the internal axon and its myelin sheath with subsequent Wallerian degeneration distal to the lesion site, but with anatomical preservation of the surrounding endoneurium, perineurium, and epineurium, allowing spontaneous axonal regeneration.

\medterm{Ayurveda} A traditional system of health care from India that uses ideas about bodily balance and may include diet, lifestyle practices, massage, and herbal products. Some products can interact with medicines or contain harmful contaminants.

\medterm{Ayurvedic Medicine} Health care based on traditional Indian theories of bodily balance, using individualized diet, lifestyle practices, physical treatments, and herbal products. Evidence and safety vary, and some products may interact with medicines or contain toxic ingredients.

\medterm{Azole Antifungals} A group of medicines that stop fungi making ergosterol, a substance needed for strong cell membranes. This weakens or kills the fungus. They can interact with other medicines and may affect the liver or heart rhythm.

\medterm{Azoospermia} The absence of sperm in semen on laboratory examination. It may result from a blockage that prevents sperm from leaving the reproductive tract or from reduced or absent sperm production in the testes.

\medterm{Azotaemia} Alternate spelling; see \textit{Azotemia}.

\medterm{Azotemia} A build-up of nitrogen-containing waste products in the blood, usually because the kidneys are filtering less effectively. It may occur with dehydration, reduced kidney blood flow, or kidney disease.

\medterm{Azygos Vein} A large vein running along the right side of the spine that drains blood from the chest wall into the superior vena cava. It provides an alternate route for blood when major veins are blocked.

\medterm{Azygous Vein} Alternate spelling; see \textit{Azygos Vein}.

\medterm{A1AT Deficiency} Abbreviation for alpha-1 antitrypsin deficiency, an inherited disorder that can cause early emphysema, liver disease, or both because protective alpha-1 antitrypsin levels are too low or abnormal.

\medterm{AAT} Abbreviation for alpha-1 antitrypsin, a protein made mainly in the liver that helps protect lung tissue from enzyme-related damage. Blood testing can measure its amount and type.

\medterm{AAT Deficiency} Alternate name for alpha-1 antitrypsin deficiency, an inherited condition associated with emphysema, liver disease, panniculitis, and rarely vasculitis.

\medterm{Abercrombie Syndrome} An older name for systemic amyloidosis, in which abnormal amyloid proteins accumulate in organs and impair their function.

\medterm{Abnormal Pap Test} A cervical screening result showing abnormal cells or evidence of human papillomavirus. It does not by itself mean cancer; follow-up depends on the cell changes, HPV result, age, and prior screening history.

\medterm{Abruption, Placental} Alternate name for placental abruption, premature separation of the placenta from the uterine wall before delivery. It can cause vaginal bleeding, abdominal pain, uterine tenderness, fetal distress, and maternal shock.

\medterm{Absence of Menstruation, Primary} Failure to have the first menstrual period by the expected age, called primary amenorrhea. Causes include normal variation, pregnancy, hormonal disorders, anatomic differences, chronic illness, undernutrition, and excessive exercise.

\medterm{Abuse, Child} Physical, sexual, or emotional harm to a child, or neglect of a child’s basic needs. Suspected abuse requires prompt attention to safety and reporting according to local law and professional responsibilities.

\medterm{ACC -- Adenoid Cystic Carcinoma} Abbreviation for adenoid cystic carcinoma, a malignant tumor of glandular tissue that may grow slowly but can invade nerves and recur or spread years after treatment.

\medterm{ACC -- Agenesis of Corpus Callosum} Abbreviation for agenesis of the corpus callosum, a congenital absence or incomplete development of the major nerve-fiber bridge connecting the brain’s two hemispheres.

\medterm{ACC -- Aplasia Cutis Congenita} Abbreviation for aplasia cutis congenita, a birth defect in which an area of skin, most often on the scalp, is absent or incompletely formed.

\medterm{ACH} Abbreviation commonly used for achondroplasia, a genetic skeletal disorder causing disproportionate short stature and characteristic changes in the skull, spine, and long bones.

\medterm{Achilles Tendon Problems} Disorders affecting the Achilles tendon, including tendinitis, tendinopathy, bursitis, partial tear, and rupture. Pain, stiffness, swelling, or a sudden loss of push-off strength warrants assessment.

\medterm{ADD/ADHD} Attention-deficit/hyperactivity disorder, a neurodevelopmental condition involving persistent inattention, hyperactivity, impulsivity, or combinations of these symptoms that impair functioning.

\medterm{Adenosine Deaminase Severe Combined Immunodeficiency (ADA-SCID)} A rare inherited immune deficiency caused by adenosine deaminase deficiency. Toxic metabolites prevent normal development of lymphocytes, leaving infants highly vulnerable to severe infections.

\medterm{ADHD in Children} Attention-deficit/hyperactivity disorder in childhood, characterized by developmentally inappropriate inattention, hyperactivity, or impulsivity across settings. Diagnosis requires persistent impairment and information from caregivers, teachers, and clinical assessment.

\medterm{Aganglionic Megacolon} Congenital absence of nerve cells in part of the large intestine, causing Hirschsprung disease. The affected segment cannot relax normally, leading to intestinal obstruction, abdominal swelling, and delayed passage of stool.

\medterm{Agenesis of Commissura Magna Cerebri} An older anatomic name for agenesis of the corpus callosum, in which the major commissural fiber tract connecting the cerebral hemispheres is absent or underdeveloped.

\medterm{Agnogenic Myeloid Metaplasia (AMM)} An older name for primary myelofibrosis, a chronic bone-marrow disorder in which abnormal scar tissue and blood-cell production outside the marrow cause anemia, enlarged spleen, constitutional symptoms, and abnormal blood counts.

\medterm{Agnosia, Primary Visual} An inability to recognize or identify objects by sight despite adequate vision. It results from dysfunction in visual association areas of the brain and may follow stroke, trauma, dementia, or other neurologic disease.

\medterm{Agnosis, Primary} An inability to recognize familiar objects, people, sounds, or other stimuli despite adequate sensation and understanding. The specific form depends on the affected sensory association system.

\medterm{Airsickness} Motion sickness caused by travel through the air. Conflicting signals from the eyes and inner ear can produce nausea, dizziness, sweating, headache, and vomiting.

\medterm{Albinismus} A genetic condition involving reduced or absent melanin production, causing lighter skin, hair, or eyes and often visual problems such as reduced acuity, nystagmus, and light sensitivity.

\medterm{ALD} Abbreviation commonly referring to adrenoleukodystrophy, an inherited disorder of fatty-acid metabolism that can damage the nervous system and adrenal glands.

\medterm{Aldosteronism With Normal Blood Pressure} A form of primary aldosteronism in which excess aldosterone is present without obvious hypertension. It may still cause low potassium, metabolic alkalosis, or cardiovascular and kidney risk and requires evaluation of the cause.

\medterm{Algodystrophy} An older term for complex regional pain syndrome, a persistent pain condition after injury or surgery that may cause swelling, color or temperature change, sweating abnormalities, and stiffness.

\medterm{Algoneurodystrophy} An older term for complex regional pain syndrome, involving disproportionate regional pain with sensory, vasomotor, sudomotor, and motor or trophic changes.

\medterm{Allergies to Insect Stings} Allergic reactions to venom from bees, wasps, hornets, or fire ants. Severe reactions can cause hives, throat swelling, wheezing, low blood pressure, or collapse and require immediate epinephrine and emergency care.

\medterm{Allergy to Natural Rubber (Latex)} An immune reaction to natural rubber latex that may cause contact dermatitis, hives, wheezing, or anaphylaxis. Avoidance and clearly documented latex precautions are important for people with significant reactions.

\medterm{AloBar Holoprosencephaly} A severe form of holoprosencephaly in which the developing forebrain fails to divide and the cerebral hemispheres remain largely fused. It is usually associated with profound developmental impairment and major facial or other congenital abnormalities.

\medterm{Alopecia Celsi} An older name for alopecia areata, an autoimmune disorder causing sharply defined patches of nonscarring hair loss.

\medterm{Alopecia Cicatrisata} An older term for scarring alopecia, in which inflammation or injury destroys hair follicles and replaces them with scar tissue, causing permanent hair loss in affected areas.

\medterm{Alopecia Circumscripta} An older name for alopecia areata, characterized by one or more circumscribed patches of nonscarring hair loss.

\medterm{Alpha-1} A shortened name commonly used for alpha-1 antitrypsin deficiency or alpha-1 antitrypsin testing; the intended meaning should be confirmed from context.

\medterm{Altitude Headache} Headache related to ascent to high altitude, often as part of acute mountain sickness. It may occur with nausea, dizziness, fatigue, or sleep disturbance and can signal more serious altitude illness when severe or accompanied by neurologic or breathing symptoms.

\medterm{Alzheimer’s} A common possessive form of Alzheimer disease, a progressive neurodegenerative disorder causing worsening memory, thinking, behavior, and ability to perform daily activities.

\medterm{Anaplastic Large-cell Lymphoma} A usually aggressive T-cell non-Hodgkin lymphoma characterized by large anaplastic cells and CD30 expression. It may involve lymph nodes, skin, bone, or other organs and is classified as systemic or cutaneous.

\medterm{Anxiety & Panic} Anxiety disorders involve excessive fear, worry, or physiological arousal; panic attacks are sudden episodes of intense fear with symptoms such as palpitations, breathlessness, chest discomfort, trembling, or fear of dying.

\medterm{APLS} Abbreviation for antiphospholipid syndrome, an autoimmune disorder associated with venous or arterial thrombosis and pregnancy complications, with persistent antiphospholipid antibodies.

\medterm{Apnea, Infantile} Repeated pauses in breathing in an infant. Causes vary with age and context and may include prematurity, infection, airway obstruction, reflux, seizures, metabolic disease, or unexplained infant events; color change, poor responsiveness, or prolonged pauses are emergencies.

\medterm{Appalachian Type Amyloidosis} A rare hereditary transthyretin amyloidosis variant associated with abnormal amyloid deposition, often causing neuropathy, autonomic symptoms, or heart involvement. Diagnosis requires genetic and tissue or specialized laboratory evaluation.

\medterm{APS} Abbreviation for antiphospholipid syndrome; it is also used for other conditions in different contexts, so the intended meaning should be confirmed.

\medterm{Arachnitis} An older spelling of arachnoiditis, chronic inflammation and scarring of the arachnoid membrane around spinal or cranial nerves that can cause persistent pain, sensory symptoms, weakness, or bladder and bowel dysfunction.

\medterm{Arachnoidal Fibroblastoma} A rare benign or low-grade tumor arising from fibroblastic cells of the arachnoid coverings. Symptoms depend on location and may result from compression of the brain, spinal cord, or nerves.

\medterm{Arbovirus A Chikungunya Type} An older classification term for chikungunya virus, a mosquito-borne alphavirus that causes fever and often severe inflammatory joint pain, rash, headache, and muscle pain.

\medterm{ARDS} Abbreviation for acute respiratory distress syndrome, a life-threatening inflammatory lung injury causing fluid leakage into the air sacs, severe oxygenation failure, and often the need for ventilatory support.

\medterm{Aregenerative Anemia} Anemia with an inappropriately low reticulocyte response, indicating inadequate red-cell production by the bone marrow. Causes include iron, vitamin B12, or folate deficiency, kidney disease, inflammation, endocrine disease, marrow failure, and some medicines.

\medterm{ARG Deficiency} Abbreviation for arginase deficiency, an inherited urea-cycle disorder in which arginine accumulates and can cause progressive spasticity, developmental problems, seizures, and elevated ammonia.

\medterm{Arteritis, Takayasu} Alternate name for Takayasu arteritis, a large-vessel vasculitis that can narrow or block the aorta and its major branches, causing pulse differences, limb claudication, blood-pressure differences, and organ ischemia.

\medterm{Arthritis (Osteoarthritis)} Osteoarthritis, a degenerative joint disorder involving cartilage loss, bone remodeling, pain, stiffness, and reduced function. It commonly affects the knees, hips, hands, spine, and feet.

\medterm{Arthritis Urethritica} An older name for reactive arthritis, inflammatory arthritis occurring after certain gastrointestinal or genitourinary infections and sometimes accompanied by urethritis, conjunctivitis, or skin findings.

\medterm{Arthritis, Juvenile Rheumatoid} An older name for juvenile idiopathic arthritis, chronic inflammatory arthritis beginning in childhood. It may affect one or many joints and can involve the eyes, skin, or other organs.

\medterm{AVM} Abbreviation for arteriovenous malformation, an abnormal direct connection between arteries and veins that bypasses capillaries. AVMs can bleed, cause seizures or neurologic deficits, or produce symptoms from high-flow effects.

\medterm{Axonal Neuropathy, Giant} A rare inherited neurodegenerative disorder in which giant axons develop in peripheral nerves and the central nervous system. It usually begins in childhood with progressive ataxia, weakness, sensory loss, and developmental or cognitive changes.



\medterm{B And T Cell Screen} A blood test that measures B cells and T cells, two major groups of immune cells. It may help evaluate suspected immune deficiency, unusual blood-cell counts, or some blood cancers, but results must be interpreted with other tests.

\medterm{B And T Cells} B cells make antibodies, while T cells coordinate immune responses or destroy infected and abnormal cells. Both are types of lymphocytes, a group of white blood cells that protect against infection and disease.
\medterm{B And T Lymphocytes} Two major types of immune cells. B lymphocytes can make antibodies and immune memory, while T lymphocytes help control immune responses or destroy infected and abnormal cells.

\medterm{B Cell} A type of white blood cell that develops in the bone marrow. After activation, it can become an antibody-producing cell or a memory cell that helps the immune system respond faster to a later infection.

\synonyms
B Lymphocyte; Bursa-Equivalent Lymphocyte

\medterm{B Hemophilia} Alternate name; see \textit{Hemophilia B}.

\medterm{B Lymphocyte} Alternate name; see \textit{B Cell}.

\medterm{B Virus} A herpes virus carried mainly by macaques. It can cause a rare but severe infection of the brain and spinal cord in people after a bite, scratch, or contact with infected animal tissue.

\medterm{B-Cell Activation} The process in which a B cell receives signals after recognizing a foreign substance, then multiplies and develops into antibody-producing cells or long-lasting memory cells.

\medterm{B-Cell Leukemia/Lymphoma Panel} A laboratory test that examines cells from blood, bone marrow, or tissue for abnormal groups of B cells. It uses cell features and marker proteins to help identify or classify some blood cancers.

\medterm{B-Cell Lymphoma} A group of cancers that begin in B lymphocytes. Different types grow at different speeds and may affect lymph nodes or organs; diagnosis and treatment depend on the exact type, extent, and person's health.

\medterm{B-Cell Receptor} A protein on the surface of a B cell that binds a particular foreign substance. This binding sends signals that can make the cell multiply, mature, and produce antibodies.

\medterm{B-Lymphocyte} Alternate spelling; see \textit{B Cell}.

\medterm{B-Lynch Suture} A surgical stitch placed around the uterus to control severe bleeding after childbirth when the uterus remains relaxed and medicines have not stopped the bleeding. It compresses the uterus while usually preserving the possibility of future pregnancy.

\medterm{B-Mode Ultrasound} The usual two-dimensional ultrasound image created from echoes of sound waves sent into the body. It shows the size, shape, and structure of organs and soft tissues without using radiation.

\medterm{B-Type Natriuretic Peptide} A hormone released mainly by the heart when its chambers are stretched or under pressure. A blood test for BNP or the related NT-proBNP can support assessment of heart failure, but age, kidney disease, rhythm, body size, and other illnesses affect the result.

\synonyms
BNP; Brain Natriuretic Peptide; NT-proBNP

\medterm{Ba} The chemical symbol for barium, a metallic element used in compounds such as barium sulfate. Barium sulfate can be swallowed or given as an enema to outline parts of the digestive tract on X-rays.

\medterm{Babcock Clamp} A surgical instrument with smooth, rounded jaws used to hold delicate tubular tissue gently. It may be used on the intestine, fallopian tubes, or other tissues that could be injured by a crushing clamp.

\medterm{Babesia} A group of parasites transmitted mainly by certain ticks. They infect red blood cells and can cause fever, tiredness, anaemia, low platelets, jaundice, or dark urine; severe illness is more likely without a spleen or with weak immunity.

\medterm{Babesia Divergens} A species of parasite that infects red blood cells and can cause babesiosis, especially in Europe. It may cause severe destruction of red cells, particularly in people without a spleen or with a weakened immune system.

\medterm{Babesia Microti} A parasite and the most common cause of babesiosis in North America. It is spread by black-legged ticks, infects red blood cells, and may cause no symptoms or severe anaemia, jaundice, and organ failure.

\medterm{Babesiosis} An infection caused by Babesia parasites, usually spread by ticks. The parasites enter red blood cells and may cause fever, tiredness, anaemia, low platelets, jaundice, or dark urine; severe disease can affect several organs.

\medterm{Babinski Reflex} A reflex tested by stroking the sole of the foot. In babies, the big toe normally moves upward and the other toes spread; in older children and adults, the toes usually curl downward. An upward response later in life may indicate a brain or spinal-cord problem.

\medterm{Babinski Response} Alternate name; see \textit{Babinski Reflex}.

\medterm{Babinski Sign} Alternate name; see \textit{Babinski Reflex}.

\medterm{Baby Acne} Small red or white bumps, usually on an infant's face, caused by temporary effects of hormones. It commonly improves without treatment; gentle washing is sufficient, and adult acne products are not used on infant skin unless a clinician advises.

\medterm{Baby Blues} A common short period of tearfulness, irritability, anxiety, or mood changes beginning in the first days after childbirth. It usually improves within two weeks; persistent, severe, or dangerous thoughts may indicate postnatal depression or an emergency.

\medterm{Baby Health} An infant's physical growth, development, feeding, sleep, immunizations, screening, and protection from illness or injury. Regular care also monitors hearing, vision, movement, communication, and social development.

\medterm{Baby Holding Techniques} Safe ways to support an infant's head, neck, back, and body while carrying, feeding, or moving the baby. Newborns need full head and neck support; never shake an infant or leave one unattended on a raised surface.

\medterm{Babyhood} The period of infancy, from birth through the first year of life. During this time, babies undergo rapid growth and development of movement, communication, feeding skills, immunity, and social responses.

\medterm{Bacillaceae} A family of rod-shaped bacteria that includes species able to form resistant spores. Some members are harmless or useful, while others can cause infections such as anthrax or food poisoning.

\medterm{Bacillary Angiomatosis} A bacterial infection, usually caused by Bartonella, that produces red or purple skin bumps in people with weakened immunity. It can also affect the liver, spleen, bones, or lymph nodes and may resemble other skin tumours.

\medterm{Bacillary Dysentery} An intestinal infection, commonly caused by Shigella bacteria, that produces fever, abdominal cramps, and frequent diarrhoea that may contain blood or mucus. Severe dehydration or complications require medical care.

\medterm{Bacille Calmette-Guérin} A weakened form of a mycobacterium used as a vaccine against tuberculosis in many countries. It is also placed into the bladder as an immune treatment for some early bladder cancers.

\medterm{Bacilluria} The presence of rod-shaped bacteria in urine, detected by microscopy or culture. It may represent contamination, harmless colonization, or a urinary infection; symptoms, the amount of bacteria, and the organism determine its significance.

\medterm{Bacillus} A rod-shaped bacterium. The word can describe bacterial shape or members of the Bacillus genus, many of which form resistant spores; some are harmless, while others cause disease.

\medterm{Bacillus Anthracis} A spore-forming bacterium that causes anthrax in people and animals. Infection may affect the skin, digestive tract, or lungs and can rapidly become life-threatening, requiring urgent specialist treatment and public-health action.

\medterm{Bacillus Cereus} A spore-forming bacterium that can cause food poisoning after contaminated food is eaten. It may cause sudden vomiting or watery diarrhoea and, rarely, serious infections in people with weakened immunity.

\medterm{Bacillus Cereus Food Poisoning} A foodborne illness caused by toxins produced by Bacillus cereus, manifesting in two distinct clinical forms: an emetic syndrome characterized by acute nausea and vomiting developing 1 to 5 hours after ingestion (classically linked to fried or reheated rice), and a diarrheal syndrome with abdominal cramps and watery diarrhea developing 8 to 16 hours after ingestion (associated with contaminated meats, vegetables, or sauces).

\medterm{Bacitracin} An antibiotic applied to the skin that prevents susceptible bacteria from building their protective cell wall. It is used for selected superficial infections, but it can cause skin irritation or allergic reactions.

\medterm{Bacitracin Overdose} Harm caused by too much exposure to bacitracin. Skin use usually causes local irritation, while swallowing a large amount or receiving excessive injected medicine can cause vomiting, allergic reactions, or kidney injury.

\medterm{Bacitracin Zinc Overdose} Harm caused by excessive exposure to bacitracin zinc. The effects depend on the amount and route; large exposures may cause vomiting, allergic reactions, or kidney injury and require prompt medical treatment.

\medterm{Back Injury} Damage to the muscles, ligaments, joints, discs, bones, or nerves of the back after trauma, lifting, or repeated strain. Pain and stiffness are common; weakness, numbness, loss of bladder or bowel control, or major-trauma pain requires urgent assessment.

\medterm{Back Pain} Pain or discomfort in any part of the back, from the neck to the lower back. It is often caused by muscles, joints, or discs, but may also arise from nerves or internal organs.

\synonyms
Dorsalgia; Backache; Spinal Pain

\medterm{Back Pressure} Pressure that builds up behind an obstruction or against the normal direction of flow. It may occur in an airway, blood vessel, duct, or hollow organ and can impair movement of air, blood, or fluid.

\medterm{Back-Blows} Firm blows delivered between the shoulder blades as first aid for severe choking in a responsive person. They are used with age-appropriate choking procedures and do not replace emergency help.

\medterm{Backache} Pain or discomfort in the back, from the neck to the buttocks. It is a symptom rather than a disease and may result from muscles, joints, discs, nerves, or pain referred from an internal organ.

\medterm{Backwash Ileitis} Mild, superficial inflammation of the terminal ileum that occurs in patients with extensive ulcerative colitis (pancolitis). It results from an open, incompetent ileocecal valve that allows inflamed colonic contents and bacteria to reflux backward into the small intestine. Unlike Crohn's disease, the inflammation is continuous with the cecum, is confined to the mucosa without strictures or granulomas, and resolves after surgical removal of the colon. Its presence is associated with severe colitis and an increased long-term risk of colorectal dysplasia.

\synonyms
Reflux Ileitis; Ulcerative Colitis Backwash Ileitis

\medterm{Baclofen} A medicine that reduces muscle spasm by lowering nerve signals in the spinal cord. It can cause sleepiness, weakness, dizziness, or confusion; stopping it suddenly after regular use can cause dangerous withdrawal symptoms.

\medterm{Bacteraemia} Alternate spelling; see \textit{Bacteremia}.

\medterm{Bacteremia} The presence of living bacteria in the bloodstream. It may be temporary, intermittent, or continuous. Some cases cause no illness, while others produce fever and can progress to sepsis or infection of distant organs.

\medterm{Bacteria} Microscopic single-celled organisms. Many live harmlessly or help the body, while others cause infections; bacteria can live alone or in communities and may develop resistance to antibiotics.

\medterm{Bacterial Adherence} The process by which bacteria attach to body cells or surfaces, often the first step in infection. Surface proteins and tiny hair-like structures help bacteria stay in place and resist removal.

\medterm{Bacterial Biofilm} A community of bacteria attached to a surface and enclosed in a protective material they produce. Biofilms can resist immune defenses and antibiotics and cause persistent infections on wounds, catheters, implants, or teeth.

\medterm{Bacterial Capsule} A protective outer coating around some bacteria, often made of sugars. It can help bacteria avoid being swallowed by immune cells, survive in the body, and cause disease.

\medterm{Bacterial Cell Wall} A strong outer layer that gives bacteria their shape and prevents them from bursting. Many bacterial cell walls contain peptidoglycan, a material made from linked sugars and short protein chains.

\medterm{Bacterial Classification} The grouping of bacteria according to their genetic features, appearance, chemistry, and evolutionary relationships. Classification helps identify organisms, understand relatedness, and guide diagnosis, infection control, and treatment.

\medterm{Bacterial Conjugation} The transfer of DNA from one bacterium to another through direct contact. The transferred DNA may include genes that help bacteria resist antibiotics or cause disease.

\medterm{Bacterial Conjunctivitis} A contagious bacterial infection of the thin surface covering the eye and inner eyelids. It commonly causes redness, irritation, gritty discomfort, and thick discharge that sticks the eyelids together. Eye pain, light sensitivity, reduced vision, or contact-lens use needs prompt assessment.

\medterm{Bacterial Dermatitis} Skin inflammation caused by bacterial infection. The skin may become red, warm, swollen, painful, or covered with pus or crusts. It can start in healthy skin or develop when bacteria infect an existing rash or wound.

\medterm{Bacterial Dormancy} A low-activity state that allows bacteria to survive conditions in which they cannot grow. Some form highly resistant spores; others remain as ordinary cells that temporarily tolerate stress and may become active again when conditions improve.

\medterm{Bacterial Endocarditis} A serious infection of the heart’s inner lining or valves. It may cause fever, chills, tiredness, a new or changing murmur, or weight loss. Bacteria can damage valves or break off and block blood vessels; diagnosis usually requires blood cultures and heart imaging.

\medterm{Bacterial Endocarditis Vegetations} Clumps of bacteria, clotting material, and platelets that form on a heart valve or the heart's inner lining during bacterial endocarditis. They can develop when bacteria in the blood attach to a damaged or abnormal heart surface.

\medterm{Bacterial Endospores} Highly resistant resting structures made by some bacteria, including Bacillus and Clostridium, when nutrients are scarce. They can survive drying, heat, chemicals, and long periods of time, then grow into active bacteria when conditions improve.

\medterm{Bacterial Flagella} Thread-like structures that help some bacteria move through liquid or across surfaces. A rotating base works like a motor and turns the attached filament; flagella can also influence how bacteria attach and spread.

\medterm{Bacterial Gastroenteritis} Acute inflammation of the stomach and intestines caused by pathogenic bacteria such as
Salmonella, Campylobacter, Shigella, or pathogenic Escherichia coli. It commonly causes
diarrhea, abdominal cramps, fever, and sometimes bloody stools or vomiting.

\medterm{Bacterial Genetics} The study of bacterial genes, DNA, inheritance, and genetic change. Bacteria usually have one main chromosome and may carry extra DNA rings called plasmids that can include antibiotic-resistance genes.

\medterm{Bacterial Growth Phases} The stages bacteria pass through when growing in a culture. They first adapt, then multiply rapidly, later reach a balance as nutrients fall and waste builds up, and finally decline as conditions become unsuitable.

\medterm{Bacterial Infection} An illness caused by single-celled microscopic organisms called bacteria entering, multiplying, and spreading in body tissues. While millions of harmless or beneficial bacteria normally live on the skin and in the gut (the microbiome), infections happen when disease-causing bacteria invade the body or when normal bacteria get into areas that are usually sterile, such as the lungs, kidneys, bladder, or bloodstream. Bacteria damage tissues directly, produce poisons called toxins, or trigger intense immune responses that cause swelling, redness, fever, and pus. Common examples include strep throat, urinary tract infections, bacterial pneumonia, skin boils, and food poisoning; severe infections can spread to the blood and cause life-threatening sepsis. Unlike viral infections, bacterial infections can be treated and cured with antibiotics, but taking antibiotics correctly is essential to prevent bacteria from becoming resistant to medicines.

\synonyms
Bacterial Disease; Bacterial Illness; Bacteriosis

\medterm{Bacterial Meningitis} A serious infection of the membranes around the brain and spinal cord. It can cause fever, severe headache, stiff neck, vomiting, confusion, seizures, or a rash and requires immediate hospital treatment.

\medterm{Bacterial Metabolism} The chemical reactions bacteria use to obtain energy, build cell material, grow, and survive. Different bacteria use oxygen, fermentation, light, or other substances, and some of these processes are targets for antibiotics.

\medterm{Bacterial Morphology} The size, shape, and arrangement of bacteria seen under a microscope. Common forms are spheres, rods, and spirals; bacteria may occur singly, in pairs, chains, clusters, or other patterns.

\medterm{Bacterial Overgrowth} Small-intestinal bacterial overgrowth is an excessive number of bacteria in the small intestine, where relatively few normally live. It can cause bloating, abdominal pain, diarrhoea, weight loss, or vitamin deficiencies and may result from poor intestinal movement or structural changes.

\synonyms
Small Intestinal Bacterial Overgrowth; SIBO

\medterm{Bacterial Pathogenesis} The ways bacteria enter the body, attach to cells, invade tissues, avoid immune defenses, and cause illness. Damage may result from direct invasion, toxins, inflammation, or disruption of normal body functions.

\medterm{Bacterial Peritonitis} A bacterial infection of the lining of the abdominal cavity. It may follow bowel perforation, surgery, another abdominal infection, or infection of abdominal fluid and can cause severe pain, fever, vomiting, and sepsis.

\medterm{Bacterial Pharyngitis} Bacterial infection and inflammation of the throat, often caused by group A Streptococcus. It may cause sudden sore throat, fever, painful swallowing, swollen neck glands, and sometimes white patches on the tonsils.

\medterm{Bacterial Pili} Tiny hair-like structures on the surface of some bacteria. They help bacteria attach to body cells or surfaces; some also help them move or take up DNA from their surroundings.

\medterm{Bacterial Pneumonia} A bacterial infection of the lungs that causes inflammation and may fill air spaces with fluid or pus. It can cause fever, cough, phlegm, chest pain, shortness of breath, and low oxygen levels.

\medterm{Bacterial Prostatitis} Acute or chronic bacterial infection of the prostate gland, typically caused by uropathogenic Gram-negative bacilli (such as \textit{Escherichia coli}), presenting with pelvic pain, dysuria, urinary frequency, fever, and tender prostate.

\synonyms
Acute Bacterial Prostatitis; Chronic Bacterial Prostatitis; Infectious Prostatitis

\medterm{Bacterial Protein} A protein produced by a bacterium that may form part of its structure, support its metabolism, help it attach to host cells, or contribute to disease.

\medterm{Bacterial Spore} A dormant, highly resistant form made by some bacteria when conditions are harsh. Spores can survive drying, heat, and many disinfectants and may later grow into active bacteria when conditions improve.

\medterm{Bacterial Toxemia} The presence or effects of harmful bacterial substances in the bloodstream. It can occur during severe infection and may cause widespread inflammation, low blood pressure, shock, or organ injury.

\medterm{Bacterial Toxin} A harmful substance made or released by bacteria. It may kill or injure cells, interfere with body functions, or trigger severe inflammation; toxins may be secreted or released when bacteria break apart.

\medterm{Bacterial Toxins} Harmful substances produced or released by bacteria. Some are secreted proteins that target particular cells, while others are parts of certain bacterial outer membranes that can trigger widespread inflammation when released.

\medterm{Bacterial Tracheitis} A bacterial infection and swelling of the windpipe. It can cause fever, a harsh cough, noisy or difficult breathing, and rapid narrowing of the airway, making urgent treatment necessary.

\medterm{Bacterial Transduction} The transfer of DNA from one bacterium to another by a virus that infects bacteria. The transferred DNA may give the receiving bacterium new traits, including toxin production or antibiotic resistance.

\medterm{Bacterial Transformation} The uptake of loose DNA from the environment by a bacterium able to absorb it. The new DNA may change the bacterium's features, including how it causes disease or responds to antibiotics.

\medterm{Bacterial Vaginosis} A common change in the vaginal bacteria in which protective lactobacilli decrease and other bacteria increase. It may cause a thin, grey or white discharge and a fish-like smell, but many people have no symptoms.

\medterm{Bacterial Virulence Factors} Features that help bacteria cause disease, such as attaching to cells, entering tissues, avoiding immune defenses, producing toxins, or obtaining nutrients. Examples include surface proteins, protective coatings, enzymes, and toxins.

\medterm{Bactericide} A substance that kills bacteria rather than only stopping them from multiplying. Whether it works depends on the organism, medicine, dose, infection site, and the person's immune defenses.

\medterm{Bacteriocin} A protein-like substance made by some bacteria that inhibits or kills other susceptible bacteria nearby. It gives the producing bacterium a competitive advantage, but its medical importance depends on the organism and setting.

\medterm{Bacteriology} The study of bacteria, including their structure, genes, growth, spread, effects on health, and role in disease. It includes laboratory methods used to identify bacteria and test antibiotic sensitivity.

\medterm{Bacteriophage} A virus that infects bacteria and uses bacterial cells to make more copies of itself. Some phages destroy the bacteria they infect and are being studied for selected antibiotic-resistant infections.

\medterm{Bacteriophage Lambda} A virus that infects the bacterium Escherichia coli. It may multiply and burst the bacterial cell or insert its genetic material into the cell’s DNA and remain inactive until later.

\medterm{Bacteriophage T4} A virus that infects Escherichia coli, makes many copies inside the bacterium, and usually breaks the cell open to release them.

\medterm{Bacteriophages} Viruses that infect bacteria and use them to make more viruses. They may destroy the bacterial cell or insert their genetic material into it and remain inactive for a time.

\medterm{Bacteriorhodopsin} A light-sensitive protein in some microorganisms. It uses light energy to move charged particles across the cell membrane, helping the organism produce energy.

\medterm{Bacteriostasis} The stopping or slowing of bacterial growth without necessarily killing the bacteria. The immune system can then remove the remaining organisms; this differs from an agent that directly kills bacteria.

\medterm{Bacteriostatic} Describes a substance that stops or slows bacteria from multiplying without directly killing them. The body's immune system can then help clear the remaining bacteria.

\medterm{Bacterium} A single bacterial cell. It usually lacks a nucleus, has genetic material in the cell interior, and may have an outer cell wall; some bacteria are helpful, while others cause infection.

\medterm{Bacteriuria} The presence of bacteria in urine detected by testing. It may occur with urinary symptoms as part of an infection or without symptoms; the result must be interpreted with the sample quality, amount of bacteria, pregnancy, and other medical factors.

\medterm{Bacteroidaceae} A family of bacteria that commonly live in the intestine and grow without oxygen. They are usually harmless in the gut but can cause abscesses or other infections if they enter wounds or body areas that should be free of bacteria.

\medterm{Bacteroides} A group of bacteria that normally live in the intestine and grow without oxygen. They are usually harmless there but can cause abscesses or bloodstream infection if bowel contents enter wounds or other sterile areas.

\medterm{Bacteroides Caccae} A bacterium that normally lives in the intestine without oxygen. It is usually harmless there but may be found in mixed abdominal or bloodstream infections when intestinal bacteria escape into other tissues.

\medterm{Bacteroides Dorei} A bacterium that normally lives in the human intestine and grows without oxygen. It is generally considered a harmless member of the gut microbiota.

\medterm{Bacteroides Fragilis} A bacterium that normally lives in the intestine without oxygen. If bowel contents leak into the abdomen, it is an important cause of abscesses, abdominal infection, bloodstream infection, and sepsis.

\medterm{Bacteroides Gingivalis} An older name for Porphyromonas gingivalis, a bacterium that grows without oxygen in dental plaque. It is associated with inflammation and destruction of the tissues supporting the teeth.

\medterm{Bacteroides Ovatus} An anaerobic Gram-negative bacterium that is part of the intestinal microbiota and can contribute to opportunistic infection if it enters damaged or normally sterile tissue.

\medterm{Bacteroides Thetaiotaomicron} An anaerobic Gram-negative bacterium that is an important member of the human gut microbiota and helps break down complex dietary carbohydrates; it can cause opportunistic infection outside the intestine.

\medterm{Bacteroides Ureolyticus} A bacterial name historically applied to an anaerobic, urease-producing organism now generally classified as Campylobacter ureolyticus, which can be associated with wound or bloodstream infection.

\medterm{Bacteroides Vulgatus} An anaerobic Gram-negative bacterium commonly found in the human intestine that can contribute to infection when the intestinal barrier is disrupted.

\medterm{Bad Breath} An unpleasant odor from the mouth, medically called halitosis. It most often results from bacteria, food debris, gum disease, tooth decay, dry mouth, or tonsil problems; less commonly, it may be related to nose, throat, digestive, or other medical conditions.

\synonyms
Halitosis; Oral Malodor; Fetor Oris

\medterm{BAER} Alternate name; see \textit{Brainstem Auditory Evoked Response}.

\medterm{Bag-Valve Mask} A hand-operated resuscitation device, also called a manual resuscitator, consisting of a self-inflating bag, one-way valves, a face mask or airway connector, and often an oxygen reservoir. It delivers positive-pressure ventilation and may be used with a face mask, supraglottic airway, or endotracheal tube.

\synonyms
BVM
Manual Resuscitator
Bag-Valve-Mask Device

\medterm{Bag-Valve-Mask Ventilation} A fundamental emergency airway and resuscitation procedure that delivers manual positive-pressure breathing support to a patient who has stopped breathing (apnea) or is breathing inadequately. A rescuer seals a cushioned anatomical mask tightly over the patient's nose and mouth using the specialized ``E-C clamp'' hand technique while maintaining an open airway, and rhythmically squeezes a self-inflating bag connected to high-flow oxygen. Each gentle squeeze forces oxygenated air past the vocal cords into the lungs, causing the chest to visibly rise. Bag-valve-mask ventilation is the essential first-line rescue bridge in emergency life support and CPR, sustaining vital oxygenation and carbon dioxide elimination until spontaneous breathing returns or an advanced artificial airway (such as an endotracheal tube) is secured.

\synonyms
BVM Ventilation; Bag-Mask Ventilation; Ambu Bag Ventilation; Manual Resuscitator Ventilation

\medterm{Bagassosis} An allergic lung inflammation caused by breathing dust from mouldy bagasse, the fibrous waste left after sugarcane juice is extracted. It may cause cough, fever, breathlessness, and fatigue after workplace exposure.

\medterm{Bags Under Eyes} Mild puffiness or loose skin beneath the eyes. It commonly develops with aging, inherited facial features, lack of sleep, allergies, or fluid retention and is usually harmless.

\synonyms
Puffy Eyes

\medterm{Bainbridge Reflex} An increase in heart rate that can occur when extra blood stretches the upper chambers of the heart. It helps the heart move increased venous blood forward.

\medterm{Bak Protein} A protein on the outer membrane of mitochondria, the cell’s energy-producing structures, that can help start controlled cell death. When activated, it makes openings that release cell-death signals; abnormal control may contribute to cancer.

\medterm{Baker Cyst} A fluid-filled swelling behind the knee caused by fluid from the knee joint collecting in a small sac. It often occurs with arthritis, a meniscus injury, or another condition causing knee inflammation.

\synonyms
Popliteal Cyst; Baker's Cyst

\medterm{Baker's Cyst} A fluid-filled lump behind the knee, usually caused by excess fluid from an irritated knee joint. It may cause tightness or pain and can occasionally rupture, producing calf swelling that can resemble a blood clot.

\medterm{Baker's Cyst} Alternate spelling; see \textit{Baker Cyst}.

\medterm{Baking Powder Overdose} Illness after swallowing a large amount of baking powder. Depending on its ingredients and amount, it may cause vomiting, abdominal swelling, excess sodium, disturbed body salts, or an abnormal blood-alkali level.

\medterm{Baking Soda Overdose} Harm caused by swallowing too much sodium bicarbonate. It can cause vomiting, abdominal swelling, excess sodium, disturbed body salts, abnormal blood acidity, seizures, or heart complications and requires prompt medical treatment.

\medterm{Balamuthia} A genus of free-living amoebas found in soil and water. Some species can rarely infect people, usually causing slowly progressive inflammation of the brain and sometimes skin disease.

\medterm{Balamuthia Mandrillaris} A free-living amoeba found in soil and water that can rarely infect humans. It may enter through broken skin or the nose and cause severe, gradually worsening infection of the brain or skin.

\medterm{Balamuthia Mandrillaris Infection} A rare, usually fatal infection caused by the free-living amoeba Balamuthia mandrillaris. It may
begin as a skin lesion and progress to granulomatous amoebic encephalitis with headache, confusion, seizures, or
focal neurologic deficits. Diagnosis is difficult and requires specialized testing; urgent infectious-disease and
neurologic evaluation is essential.

\medterm{Balance} The ability to keep the body steady while sitting, standing, or moving. It depends on information from the inner ear, eyes, muscles, joints, and brain; problems with any of these systems can increase the risk of falls.

\medterm{Balance Assessment} An evaluation of steadiness during sitting, standing, walking, and changing position. It may include movement tests and review of vision, medicines, strength, sensation, and dizziness to identify fall risk and guide support.

\medterm{Balance Disorders} Problems with steadiness or coordinated movement that may cause dizziness, unsteadiness, or falls. Causes include inner-ear disease, vision problems, brain or nerve disorders, muscle weakness, medicines, and blood-pressure changes.

\medterm{Balance Disorders And Vertigo} Conditions that cause unsteadiness, dizziness, or a false feeling of movement. Causes may involve the inner ear, brain, nerves, vision, muscles, medicines, or blood pressure and require assessment when severe or new.

\medterm{Balance Impairment} Reduced ability to remain steady while sitting, standing, or walking. It may cause swaying, stumbling, or falls and can result from inner-ear disease, vision problems, weakness, nerve disease, medicines, or brain disorders.

\medterm{Balance Problems} Difficulty maintaining body stability while standing or moving, often caused by disorders of the inner ear, nervous system, vision, or muscles.

\medterm{Balance Training} Exercises that improve steadiness and reduce falls by practising safe posture, weight shifting, walking, and responses to movement. A clinician or therapist can adapt exercises to a person's strength, dizziness, and medical conditions.

\medterm{Balanced Study} A clinical study design in which participant characteristics are evenly distributed between comparison groups at baseline, either by design or by randomization, to reduce confounding and allow fair comparison of interventions.

\medterm{Balanitis} Inflammation of the glans penis, often involving the prepuce (balanoposthitis). Causes
include poor hygiene, irritants, infection, and dermatologic conditions such as lichen
planus, psoriasis, and lichen sclerosus. It may cause erythema, swelling, discharge,
pruritus, and discomfort; treatment depends on the cause.

\medterm{Balanitis Xerotica Obliterans} A long-term inflammatory skin disorder affecting the foreskin and head of the penis. It may cause pale tight skin, itching, pain, painful erections, difficulty passing urine, and narrowing of the opening; it needs medical treatment and follow-up.

\synonyms
BZO; Penile Lichen Sclerosus

\medterm{Balanitis, Non-Infectious / Infectious} Alternate index label; see \textit{Balanitis}.

\medterm{Balanoposthitis} Inflammation of both the glans penis and the prepuce. It may cause pain, erythema,
swelling, discharge, or difficulty retracting the foreskin and can result from
infection, irritation, or inflammatory skin disease.

\medterm{Balantidiasis} Balantidiasis is intestinal infection caused by the protozoan Balantidium coli, producing diarrhea, abdominal pain, or colitis, especially after fecal exposure.

\medterm{Balantidium} A genus of large, single-celled parasites that can infect the large intestine. Infection may cause diarrhoea, abdominal pain, or inflammation of the colon, although many infected people have no symptoms.

\medterm{Balantidium Coli} A single-celled parasite that can infect the large intestine, usually after contact with food or water contaminated by animal or human stool. It may cause diarrhoea, abdominal pain, or colitis.

\medterm{Balbiani Body} A temporary cluster of cell parts and genetic messages near the nucleus of a developing egg cell. It helps organize materials that the early embryo needs after fertilization.

\medterm{Baldness} Commonly used term for substantial hair loss, especially androgenetic alopecia, also called pattern
hair loss. In this condition, genetically susceptible hair follicles gradually become smaller under the influence of
androgens.

\synonyms
Pattern Hair Loss

\medterm{Balint's Syndrome} A rare disorder caused by damage to both back portions of the brain. It can cause difficulty seeing more than one object at a time, reaching accurately for objects while looking at them, and directing visual attention.

\medterm{Ballance Sign} A physical examination finding characterized by fixed dullness to percussion in the left flank and shifting dullness in the right flank. It indicates a ruptured spleen with a localized left subcapsular or perisplenic hematoma and free blood in the peritoneal cavity.

\medterm{Ballard Score} A newborn examination used to estimate how many weeks of pregnancy had passed at birth. It assesses physical features and muscle tone or movement; the estimate is useful when pregnancy dates or early ultrasound information are uncertain.

\medterm{Ballism} A movement disorder causing sudden, involuntary, large, flinging movements of an arm or leg. It usually results from damage or dysfunction in deep movement-control areas of the brain.

\medterm{Balloon Angioplasty} A procedure in which a small balloon on a catheter is inflated inside a narrowed artery to improve blood flow. A stent may be placed to help keep the artery open.

\medterm{Balloon Aortic Valvuloplasty} A catheter procedure in which a balloon is inflated across a narrowed aortic valve to improve blood flow from the heart. The benefit may be temporary, and the valve can narrow again; some people later need valve replacement.

\medterm{Balloon Brachytherapy} A form of internal radiation treatment in which a balloon catheter is placed in a body space, often the cavity left after breast-cancer surgery, and radioactive material is temporarily put inside it to treat nearby tissue.

\medterm{Balloon Catheter} A thin flexible tube with an inflatable tip. It can widen a narrowed blood vessel or body passage, open a stent, stop bleeding by applying pressure, or help perform another procedure.

\medterm{Balloon Catheter, Vascular} An inflatable balloon mounted at the distal tip of a flexible vascular catheter used to dilate stenotic or occluded blood vessels during angioplasty, deliver and expand endovascular stents, or temporarily occlude vascular flow.

\synonyms
Angioplasty Balloon; Vascular Balloon Dilator; Balloon Catheter

\medterm{Balloon Dilation} A procedure that widens a narrowed blood vessel or body passage by inflating a small balloon at the narrowed area. The result and risks depend on the site and underlying cause.

\synonyms
Balloon Dilatation; Balloon Angioplasty When Performed in a Blood Vessel

\medterm{Balloon Occlusion} Temporary blockage of a blood vessel or other passage using an inflatable balloon on a catheter. It may control bleeding, test blood flow, or support another procedure; possible complications include reduced blood supply, clotting, and vessel injury.

\medterm{Balloon Tamponade} Temporary control of bleeding or blockage by inflating a balloon inside or against a body passage or blood vessel. It applies pressure to the bleeding area but can injure tissue and requires close monitoring.

\medterm{Balneology} The study of mineral waters, bathing, and spa treatments and their possible health effects. These treatments may ease some symptoms but do not replace proven care, and heat or water exposure can be unsafe for some people.

\medterm{Balneotherapy} Treatment using mineral or thermal water, baths, muds, or related spa methods. It may relieve symptoms in selected conditions, but evidence varies and heat, infection, falls, or heart strain can make it unsafe for some people.

\medterm{Balsam} A thick, resin-containing substance from certain plants, used in some fragrances, skin products, and traditional remedies. It can irritate the skin or cause allergic contact dermatitis in sensitive people.

\medterm{Band Cell} An immature neutrophil, a white blood cell that helps fight infection, whose nucleus has a curved band shape. Increased numbers may occur when the bone marrow responds to infection, inflammation, or severe stress.

\medterm{Band Keratopathy} A horizontal strip of calcium deposits in the clear front surface of the eye. It may result from long-term inflammation or high calcium levels and can cause irritation, blurred vision, or reduced sight.

\medterm{Bandage} Material placed over or around a body part to protect a wound, absorb fluid, apply pressure, or support an injured area. It may be made of gauze, adhesive material, elastic fabric, or a tubular sleeve.

\medterm{BANF} Alternate name; see \textit{Neurofibromatosis Type 2}.

\medterm{Bankart Lesion} A tear in the front-lower rim of the shoulder socket, usually after the upper arm has dislocated. It can make the shoulder unstable and increase the chance of repeated dislocation.

\medterm{Bankart Repair} An open or arthroscopic surgical procedure performed to reattach and reconstruct a torn anterior-inferior glenoid labrum and stretched joint capsule (Bankart lesion) to the glenoid rim to restore joint stability following recurrent anterior shoulder dislocations.

\medterm{Bannayan-Riley-Ruvalcaba Syndrome} A rare inherited condition that can cause many noncancerous growths, a large head size, intestinal polyps, extra freckles around the penis, and increased risk of certain cancers. It results from changes in the PTEN gene.

\medterm{Banti Syndrome} An older term for an enlarged spleen and low blood-cell counts caused by high pressure in the veins that carry blood through the liver. Modern evaluation looks for the liver, vein, or blood disorder responsible.

\synonyms
Banti Disease; Banti's Syndrome

\medterm{Bar Graph} A chart that represents values with rectangular bars whose lengths or heights correspond to the quantities being compared.

\medterm{Barakat Syndrome} An inherited disorder that can cause low parathyroid hormone, hearing loss, and kidney abnormalities. These problems may occur together because of genetic changes affecting the development of several organs.

\medterm{Barber Itch} An inflammatory infection of hair follicles in the beard area, often caused by bacteria or dermatophyte fungi and associated with itchy bumps or pustules.

\synonyms
Barber's Itch

\medterm{Barber Say Syndrome} A rare congenital disorder caused by changes in the TWIST2 gene, characterized by excessive facial hair growth (hypertrichosis), loose, atrophic skin, a wide mouth, distinctive facial features, and eyelid abnormalities such as ectropion.

\medterm{Barbiturate} A group of medicines that slow brain activity and can cause relaxation, sleep, anaesthesia, or seizure control. They can suppress breathing, cause dependence, and be fatal in overdose, so use requires careful medical supervision.

\medterm{Barbiturate Intoxication And Overdose} Poisoning caused by too much of a barbiturate medicine. It can cause extreme sleepiness, confusion, slow or stopped breathing, low blood pressure, coma, or death and requires emergency treatment.

\medterm{Barbiturates} Medicines that reduce activity in the brain and spinal cord. They may be used for anaesthesia or selected seizure disorders, but their risk of breathing suppression, dependence, and overdose means safer alternatives are often preferred.

\medterm{Bardet-Biedl Syndrome} A rare inherited disorder that may cause retinal degeneration, obesity, extra fingers or toes, kidney disease, and developmental or reproductive problems.

\synonyms
BBS

\medterm{Bare-Metal Stent} An expandable metal mesh scaffold without an antiproliferative drug coating, deployed via endovascular catheterization to maintain vascular patency in coronary or peripheral arteries following balloon angioplasty.

\synonyms
BMS; Uncoated Vascular Stent; Bare Stent

\medterm{Bariatric Embolization} A specialized catheter procedure that reduces blood flow to selected stomach arteries to alter hunger-related signaling and support weight loss. It remains limited or investigational in many settings and may cause bleeding, stomach injury, or other vascular complications.

\medterm{Bariatric Surgery} Operations that reduce stomach capacity or change digestion to support substantial weight loss and improve obesity-related disease. They require long-term changes in eating, supplements, follow-up, and monitoring for complications.

\medterm{Bariatrics} The medical field focused on preventing and treating obesity and related conditions. Care may include nutrition, physical activity, psychological support, medicines, and weight-loss surgery.

\medterm{Barium} A metallic element used medically mainly as barium sulfate, a substance that outlines the digestive tract during certain X-ray examinations. Soluble barium compounds are poisonous and are different from barium sulfate.

\medterm{Barium Enema} An X-ray examination of the large intestine after barium is put into the rectum. Sometimes air is added to show the bowel lining more clearly. It may be used to examine blockage, diverticula, or changes after surgery; other tests are now preferred for routine colon-cancer screening.

\medterm{Barium Meal} An X-ray examination in which swallowed barium coats the swallowing tube, stomach, and first part of the small intestine. It can show narrowing, ulcers, abnormal movement, or other structural changes.

\medterm{Barium Study} An X-ray examination that uses barium to outline parts of the digestive tract. Examples include a barium swallow, barium enema, and upper-GI series. It can show the shape and movement of the digestive tract.

\medterm{Barium Sulfate} A liquid or paste swallowed or placed in the rectum to outline the digestive tract during X-ray examinations. It is not absorbed normally, but breathing it into the lungs or leakage through a bowel hole can cause serious complications.

\medterm{Barium Swallow} An X-ray examination of the swallowing tube after a person drinks barium. It can show narrowing, abnormal shape, or problems with movement. A barium follow-through continues the examination into the small intestine.

\medterm{Barium-133} A radioactive form of barium used mainly to check and calibrate radiation detectors and imaging equipment. It is not ordinarily used as a diagnostic scan or treatment for patients.

\medterm{Barlow Maneuver} A clinical physical examination test used to detect developmental dysplasia of the hip in infants. The examiner adducts the flexed hip and applies gentle posterior pressure along the shaft of the femur; a palpable clunk or sensation of the femoral head slipping out of the acetabulum indicates hip instability.

\synonyms
Barlow Test; Barlow Dislocation Test

\medterm{Barlow Sign} A check for hip problems in newborns where the doctor gently presses the baby's bent thigh backward; feeling the hip slip out of place shows the joint is unstable or loose.

\medterm{Barlow Test} A physical examination maneuver for detecting hip instability in infants. With the hip flexed, gentle pressure is applied to see whether the femoral head can be displaced out of the socket, suggesting developmental dysplasia of the hip.

\medterm{Barlow's Disease} Barlow's disease is an older name for severe vitamin-C deficiency, or scurvy, causing bleeding gums, bruising, poor wound healing, and bone or joint symptoms.

\medterm{Barognosis} The ability to recognize and compare the weight of objects by holding them. Losing this ability despite normal basic touch sensation may indicate a problem in brain areas that interpret sensory information.

\medterm{Baroreceptor} A pressure sensor in the walls of major arteries, especially in the neck and chest. It detects blood-pressure changes and signals the brain to adjust heart rate and blood-vessel narrowing.

\medterm{Baroreceptor Reflex} A rapid automatic response that helps keep blood pressure stable. Pressure sensors in major arteries signal the brain to change heart rate and blood-vessel tightness when blood pressure rises or falls.

\medterm{Baroreceptors} Stretch-sensitive nerve endings in major arteries that detect changes in blood pressure. They send signals to the brain, which adjusts heart rate and blood-vessel narrowing to help maintain stable pressure.

\medterm{Baroreflex} A rapid reflex that helps regulate blood pressure by adjusting heart rate and blood-vessel tone in response to changes in arterial pressure.

\medterm{Barotitis Media} Alternate name; see \textit{Airplane Ear}.

\medterm{Barotrauma} Injury caused by a sudden or excessive difference in pressure, affecting the lungs, ears, sinuses, or other air-filled spaces. Lung barotrauma can cause a collapsed lung or air escaping into tissues.

\medterm{Barr Body} The tightly packed, inactive copy of an X chromosome seen inside the nucleus of many cells that contain more than one X chromosome. It helps balance gene activity between cells with different numbers of X chromosomes.

\medterm{Barraquer-Simons Disease} A rare acquired disorder in which fat gradually disappears from the face, arms, and upper body while remaining or increasing around the hips and legs. It may be associated with immune disease and kidney complications.

\medterm{Barrel Chest} A rounded, enlarged appearance of the chest caused by increased lung volume or changes in the chest wall.
It may occur with chronic obstructive pulmonary disease, severe asthma, or other long-standing
lung conditions.

\medterm{Barrett Esophagus} A change in the lower lining of the swallowing tube caused by long-term acid reflux. The changed lining can increase the risk of precancerous changes and cancer.

\medterm{Barrett Esophagus Surveillance} Planned repeat examinations of the lower swallowing tube, usually with an endoscope and tissue samples, to look for precancerous cell changes or early cancer. Timing depends on whether abnormal cell changes are present and their severity.

\medterm{Barrett's Esophagus} Alternate spelling; see \textit{Barrett Esophagus}.
Columnar Epithelium Lined Lower Esophagus

\medterm{Barrett's Esophagus} Alternate spelling; see \textit{Barrett Esophagus}.

\medterm{Barrett's Oesophagus} Alternate spelling; see \textit{Barrett Esophagus}.

\medterm{Barrier Creams} Skin products that form a protective layer against moisture, friction, irritants, and body fluids. They can help prevent or soothe irritation, but infected, broken, or worsening skin may need medical assessment.

\medterm{Barrier Nursing} Infection control and nursing protocols designed to prevent the transmission of pathogenic microorganisms between patients, healthcare personnel, and visitors through strict source isolation, rigorous hand hygiene, and personal protective equipment.

\synonyms
Isolation Nursing; Barrier Precautions; Transmission-Based Precautions

\medterm{Bart's Hemoglobinopathy} A severe inherited blood disorder caused when all four genes needed to make alpha-globin are missing or inactive. The fetus cannot make enough normal haemoglobin, causing severe anaemia, swelling, and often death before or soon after birth.

\medterm{Barth Syndrome} A rare inherited disorder of mitochondrial energy production caused by changes in the TAZ gene. It may cause heart-muscle disease, low muscle tone, weakness, poor growth, low neutrophil counts, and recurrent infections.

\medterm{Barthel Index} A scoring tool that measures how much help a person needs with daily activities such as eating, bathing, dressing, using the toilet, moving between surfaces, walking, and climbing stairs. It is used to track function and rehabilitation progress.

\medterm{Bartholin Abscess} A painful collection of pus caused by infection of a Bartholin gland near the vaginal opening. It can cause severe swelling and pain when sitting or walking and usually needs drainage; antibiotics may be needed in selected cases.

\synonyms
Infected Bartholin Cyst; Bartholin Gland Abscess

\medterm{Bartholin Cyst Or Abscess} A blocked gland near the vaginal opening can cause a painless fluid-filled cyst; infection can produce a painful collection of pus. A new lump in someone aged 40 or older needs specialist assessment.

\medterm{Bartholin Gland} One of two small glands beside the vaginal opening that release lubricating fluid. A blocked duct can cause a cyst, and infection can cause a painful abscess with swelling and difficulty sitting or walking.

\medterm{Bartholin Gland Cyst} A mucus-filled retention cyst resulting from obstruction of the main duct of a greater vestibular (Bartholin) gland located in the posterolateral labium majus. While typically painless when small, secondary infection leads to acute Bartholin abscess formation requiring drainage or marsupialization.

\synonyms
Bartholin Cyst; Bartholin Duct Cyst

\medterm{Bartholin Gland Hyperplasia} Enlargement caused by an increase in the cells of a Bartholin gland. It may appear as a lump near the vaginal opening and can resemble a cyst or tumour, so a persistent or unusual mass needs assessment.

\medterm{Bartholin Gland Mucinous Cyst} A mucus-filled lump near the vaginal opening caused by blockage of a Bartholin gland duct. It is often painless but can become infected and develop into an abscess.

\medterm{Bartholin's Cyst} Alternate spelling; see \textit{Bartholin Cyst}.

\synonyms
Bartholin Gland Cyst; Bartholin Cyst

\medterm{Bartholin's Glands} Alternate spelling; see \textit{Bartholin Gland}.

\medterm{Bartholinitis} Inflammation or infection of a Bartholin gland near the vaginal opening. It may cause pain, redness, swelling, or an abscess and can make sitting or walking uncomfortable.

\medterm{Bartonella} A group of bacteria usually spread by fleas, lice, ticks, or contact with animals. Different species can cause cat-scratch disease, trench fever, skin lesions, heart-valve infection, or bloodstream infection.

\medterm{Bartonella Elizabethae} A species of Bartonella bacteria that has rarely been linked to bloodstream infection and inflammation of the heart valves, especially in people with underlying health problems.

\medterm{Bartonella Henselae} The bacterium that most commonly causes cat-scratch disease. It is usually spread by a kitten scratch or bite and may cause a small skin bump, swollen nearby lymph nodes, fever, or, rarely, disease of other organs.

\medterm{Bartonella Infections} Infections caused by Bartonella bacteria. They may produce swollen lymph nodes, fever, skin lesions, anaemia, heart-valve infection, or inflammation of organs; symptoms and severity depend on the species and immune status.

\medterm{Bartonella Quintana} A Bartonella species spread mainly by body lice. It causes trench fever, with recurring fever, headache, bone pain, and weakness, and can sometimes cause bloodstream or heart-valve infection.

\medterm{Bartonellaceae} A family of bacteria that includes Bartonella species. Members can cause cat-scratch disease, trench fever, skin and blood-vessel lesions, bloodstream infection, or infection of the heart valves.

\medterm{Bartonellosis} Infection caused by a Bartonella bacterium. Depending on the species, it may cause swollen lymph nodes, recurring fever, skin lesions, anaemia, heart-valve infection, or disease affecting several organs.

\medterm{Bartter Disease} Alternate name; see \textit{Bartter Syndrome}.

\medterm{Bartter Syndrome} An inherited kidney disorder causing excessive loss of salt and potassium in urine. It may lead to dehydration, low blood pressure, muscle cramps or weakness, frequent urination, excessive thirst, and poor growth in children.

\medterm{Basal} Relating to the base of a structure or to a person's resting level of activity, such as resting metabolism or temperature measured after sleep.

\medterm{Basal Body} A small structure at the base of a cilium or flagellum, such as a cell hair or sperm tail. It anchors the structure to the cell and helps organize the parts that produce movement.

\medterm{Basal Cell} A cell in the deepest layer of an epithelial tissue, where it can divide and produce new cells. In the epidermis, basal cells help renew the skin and contain keratinocyte stem and progenitor cells.

\medterm{Basal Cell Carcinoma} The most common malignant skin tumor, arising from basal keratinocyte-lineage cells and
associated with cumulative ultraviolet exposure and Hedgehog-pathway abnormalities. It
commonly presents as a pearly papule or nodule with telangiectasia, ulceration, or a
rolled border.

\medterm{Basal Cell Skin Cancer} Alternate name; see \textit{Basal Cell Carcinoma}.

\medterm{Basal Ganglia} Deep brain structures that help control movement, habits, learning, and emotions. Problems in these structures can cause tremor, stiffness, slowed movement, or unwanted movements.

\medterm{Basal Ganglia Dysfunction} Abnormal function of deep brain structures that help control movement and learned habits. It may cause tremor, stiffness, slowed movement, twisting movements, or other involuntary movements, as in Parkinson disease or Huntington disease.

\medterm{Basal Joint Arthritis} Alternate name; see \textit{Thumb Arthritis}.

\medterm{Basal Lamina} A thin layer beneath many lining cells that anchors them to supporting tissue. It also helps control movement of substances between the lining and the tissue below.

\medterm{Basal Metabolic Rate} The amount of energy the body uses at complete rest to maintain essential functions such as breathing, circulation, temperature, and cell activity. It varies with body size, age, muscle mass, hormones, illness, and other factors.

\medterm{Basal Nuclei} Groups of nerve cells deep within the brain that help start, stop, and control movement. They also contribute to learning, habits, motivation, and other thinking processes.

\medterm{Basal Tear} A tear that is continuously present on the eye's surface. It keeps the cornea moist, smooth, and protected and helps wash away dust and microorganisms.

\medterm{Basal Temperature} Body temperature measured immediately after waking and before eating, exercise, or other activity. A small sustained rise often occurs after ovulation and can help track menstrual cycles.

\medterm{Base} A substance that can accept hydrogen ions or reduce acidity. In water, a base usually has a pH above 7; strong bases can burn tissue, while weaker bases are important in body fluids and chemical reactions.

\medterm{Base Excision Repair} A way cells repair small areas of DNA damage. Enzymes remove the faulty DNA building block, fill the gap with the correct one, and seal the DNA strand.

\medterm{Base Pair} Two matching building blocks joined inside DNA or RNA. In DNA, adenine pairs with thymine and cytosine pairs with guanine; in RNA, adenine usually pairs with uracil.

\medterm{Base Pairing} The matching of complementary DNA or RNA building blocks. This matching helps copy genetic information accurately and allows RNA molecules to fold into useful shapes or attach to other molecules.

\medterm{Base Sequence} The exact order of DNA or RNA building blocks in a molecule. It carries genetic information and can be examined to identify inherited changes, mutations, or infectious organisms.

\medterm{Basedow's Disease} Older name; see \textit{Graves disease}.

\medterm{Baseline} A starting measurement or description of a person's health before treatment, illness, or observation. Later findings are compared with it to identify improvement, worsening, or change.

\medterm{Baseline EKG} An electrocardiogram recorded before treatment, a procedure, or a change in health. It provides a starting record for comparison with later heart tracings.

\synonyms
Baseline ECG; Reference Electrocardiogram

\medterm{Basement Membrane} A thin supporting layer beneath many cells that line organs and body surfaces. It anchors the cells, helps control exchange with underlying tissue, and can act as a barrier to tumour spread.

\medterm{Basic First Aid} Immediate help given for illness or injury until professional care arrives. It may include making the area safe, calling emergency services, controlling bleeding, helping a choking person, starting resuscitation, and monitoring breathing.

\medterm{Basic Life Support} Emergency care for a person who is unresponsive and not breathing normally. It includes calling for help, chest compressions, rescue breaths when appropriate, and use of an automated defibrillator if advised.

\medterm{Basic Metabolic Panel} A routine blood test measuring eight vital substances: blood sugar (glucose), calcium, four key electrolytes (sodium, potassium, chloride, and bicarbonate), and two waste products filtered by the kidneys (BUN and creatinine). It provides a rapid overview of fluid and electrolyte balance, kidney function, blood sugar control, and blood acid-base status.
\synonyms
BMP; Chem-7; Chemistry Screen; Electrolyte and Kidney Panel

\medterm{Basic Reproduction Number} The average number of people one infected person would infect in a population where everyone is susceptible, under specified conditions. It helps describe spread but changes with behavior, immunity, treatment, and the environment.

\medterm{Basidiomycota} A major fungal group that reproduces by forming spores on club-shaped structures called basidia. It includes mushrooms, rusts, smuts, and some fungi that can cause human disease.

\medterm{Basilar} Relating to the base or lowest part of a body structure. For example, a basilar finding may be near the bottom of the skull, lung, or another organ.

\medterm{Basilar Artery} The artery formed by the union of the vertebral arteries that supplies the brainstem, cerebellum, and posterior brain.

\medterm{Basilar Artery Aneurysm} A weak, widened area in the artery supplying the brainstem and back of the brain. It may press on nearby nerves, cause double vision or weakness, or rupture and produce life-threatening bleeding around the brain.

\medterm{Basilar Membrane} A flexible membrane in the cochlea that supports the organ of Corti and vibrates in response to sound. Different parts respond best to different sound frequencies.

\medterm{Basilic Vein} A superficial vein of the upper limb that ascends along the medial side of the arm and joins the brachial veins to form the axillary vein. It is used for venous access in selected settings.

\medterm{Basiliximab} A monoclonal antibody that blocks the interleukin-2 receptor alpha chain on activated T cells. It is used as induction immunosuppression, especially in kidney transplantation, and increases susceptibility to infection.

\medterm{Basis Pontis} The basis pontis is the front portion of the pons containing descending motor fibers and pontine nuclei that relay signals to the cerebellum.

\medterm{Basket Cell} A basket cell is an inhibitory interneuron that forms synaptic contacts around the cell body or axon initial segment of another neuron.

\medterm{Basophil} A rare type of white blood cell involved in allergic and inflammatory reactions. It releases substances such as histamine and can increase in number with some allergies, infections, and blood disorders.

\medterm{Basophil Adenoma} A usually noncancerous tumour of hormone-producing cells in the pituitary gland. It may release too much hormone or press on nearby structures, causing headache, vision problems, or other hormone-related symptoms.

\medterm{Basophilia} An unusually high number of basophils in the blood. It may occur with allergies, inflammation, infection, an underactive thyroid, or disorders in which the bone marrow makes too many blood cells.

\medterm{Basophilic Focus} An area in a tissue sample that stains blue or purple with a laboratory dye. The finding describes appearance under the microscope and must be interpreted with the surrounding tissue and other tests.

\medterm{Basophilic Stippling} The abnormal presence of multiple fine or coarse, dark-blue granules scattered throughout erythrocytes on a Romanowsky-stained peripheral blood smear. The granules represent precipitated ribosomal ribonucleic acid (RNA) and are characteristically seen in lead poisoning, thalassemias, sideroblastic anemia, and myelodysplastic syndromes.

\synonyms
Punctate Basophilia

\medterm{Basosquamous Carcinoma} A skin cancer with features of both basal-cell and squamous-cell carcinoma. It usually develops on sun-damaged skin and may behave more aggressively than ordinary basal-cell cancer, so complete treatment and follow-up are important.

\medterm{Bassen-Kornzweig Syndrome} A rare inherited disorder in which the body cannot make certain particles that carry fat in the blood. It causes poor absorption of dietary fat, deficiencies of fat-soluble vitamins, unusual red cells, and progressive nerve or eye problems.

\medterm{Bateman Purpura} Alternate name; see \textit{Actinic Purpura}.

\medterm{Bath Salts Abuse And Addiction} Harmful use of synthetic stimulant drugs sold as “bath salts.” They can cause agitation, paranoia, hallucinations, seizures, dangerous overheating, heart problems, dependence, or death; severe symptoms require emergency care.

\medterm{Bather's Itch} An itchy rash caused by tiny parasite larvae released by infected snails in fresh or salt water. The larvae cannot mature in people, so the rash usually settles on its own but can be intensely uncomfortable.

\medterm{Bathing A Patient In Bed} Washing a person who cannot safely get out of bed. Care includes keeping water comfortably warm, protecting privacy and skin, preventing falls, checking for pressure injuries, and helping only as much as needed.

\medterm{Bathing An Infant} Washing a baby safely in shallow water while keeping one hand on the baby at all times. The water is comfortably warm, the room warm, and all supplies are ready before bathing; a baby is never left alone near water.

\medterm{Batrachotoxin} A very powerful poison found in some frogs and birds. It keeps nerve and heart sodium channels open, causing abnormal nerve activity, dangerous heart rhythms, paralysis, and potentially death.

\medterm{Batroxobin} An enzyme derived from snake venom that changes fibrinogen, a blood protein, into fibrin, the material that strengthens a blood clot. It has been studied or used in selected bleeding and clotting settings and can alter bleeding risk.

\medterm{Batson Plexus} A valveless network of epidural and paravertebral veins that connects the deep pelvic and thoracic veins to the internal vertebral venous system. Because of the lack of valves and low intravascular pressures, it provides an alternate route for the retrograde hematogenous spread of pelvic infections and metastases (especially prostate and breast carcinoma) to the spine and central nervous system.

\synonyms
Batson's Venous Plexus; Vertebral Venous Plexus

\medterm{Batten Disease} A group of inherited disorders in which waste material builds up in nerve cells. They usually begin in childhood and cause worsening seizures, vision loss, movement problems, learning difficulties, and loss of abilities.

\medterm{Battered Child Syndrome} An older term for repeated physical abuse of a child by a caregiver. Possible warning signs include injuries of different ages, unexplained bruises or fractures, head injury, inconsistent explanations, or delayed medical care; suspected abuse requires immediate safeguarding action.

\medterm{Battey Bacillus} An older name for the group of nontuberculous mycobacteria now called Mycobacterium avium complex. These organisms can cause lung disease, swollen lymph nodes, or widespread infection, especially when immunity is severely weakened.

\medterm{Battle Sign} Bruising behind the ear over the bone, which is a warning sign of a break in the base of the skull.

\medterm{Bauhin Valve} The valve between the last part of the small intestine and the first part of the large intestine. It helps regulate passage of intestinal contents and limits backward movement from the colon.

\medterm{Bauxite-Workers' Lung} A lung disorder caused by long-term workplace inhalation of bauxite dust. It may produce scarring or inflammation in the lungs and can cause cough, breathlessness, or reduced lung function.

\medterm{Bax Protein} A protein that helps start programmed cell death when a cell is damaged or no longer needed. Excessive or insufficient Bax activity can affect cell survival and has been studied in cancer and other diseases.

\medterm{Bayes Theorem} A rule for updating the likelihood of a diagnosis after new evidence. It combines the chance of the disease before the test with how accurately the test result supports or argues against that disease.

\medterm{Bazin Disease} A chronic, nodular, and often ulcerating form of lobular panniculitis predominantly affecting the posterior lower legs of young or middle-aged women. It represents a delayed-type hypersensitivity reaction associated with Mycobacterium tuberculosis or other systemic antigens.

\synonyms
Erythema Induratum of Bazin; Erythema Induratum

\medterm{Bazin's Disease} An older name for tender, inflamed lumps in the fatty tissue under the skin of the lower legs. It is sometimes linked with tuberculosis but can have other causes and may ulcerate as it heals.

\medterm{Bcgosis} An infection caused by the weakened mycobacteria used in BCG treatment or vaccination. It may remain local or spread to organs, especially in people with weakened immunity, and can cause fever, cough, or inflammation.

\medterm{Beach Chair Position} A semi-sitting position used during some operations, especially shoulder surgery. The head, neck, and pressure points must be supported, and blood pressure and brain blood flow require careful monitoring.

\medterm{Beacopp Regimen} A combination chemotherapy regimen used mainly to treat Hodgkin lymphoma; it includes bleomycin, etoposide, doxorubicin, cyclophosphamide, vincristine, procarbazine, and prednisone, with dosing and supportive care determined by the oncology team.

\medterm{Beals Syndrome} A rare inherited connective-tissue disorder present from birth. It may cause permanently bent joints, long slender fingers and toes, a curved spine, and ears with a crumpled appearance; severity varies.

\medterm{Beam Hardening Artifact} A false dark band or streak on a CT image caused when bone or metal removes some parts of the X-ray beam before it reaches the detector. It can obscure nearby body structures.

\medterm{Beat Number} The number assigned to a particular heartbeat or rhythmic event when beats are counted or analysed. It can help describe timing, rate, or changes in a recorded heart rhythm.

\medterm{Beats, Cardiac} Alternate index label; see \textit{Heartbeat}.

\medterm{Beau Lines} Transverse depressions or grooves across the nail plate caused by temporary interruption
of nail-matrix growth. They may follow severe illness, fever, trauma, malnutrition,
chemotherapy, or systemic inflammatory disease.

\medterm{Beau's Lines} Transverse grooves or depressions across a fingernail or toenail caused by temporary interruption of
nail-matrix growth. They may follow severe illness, injury, high fever, nutritional stress, or
local trauma; the timing of the groove can help estimate when the event occurred.

\medterm{Becaplermin} Becaplermin is a recombinant platelet-derived growth factor gel used in selected diabetic foot ulcers with appropriate wound care.

\medterm{Bechterew Disease} An older name for ankylosing spondylitis, a long-term inflammatory disease affecting the joints between the spine and pelvis and often the spine itself. It causes back pain and stiffness that may improve with movement.

\synonyms
Ankylosing Spondylitis; Marie-Strümpell Disease

\medterm{Beck Triad} Three findings that may occur when fluid or blood compresses the heart: low blood pressure, swollen neck veins, and quiet heart sounds. The triad may be absent, so suspected compression needs urgent examination and imaging.

\medterm{Becker Muscular Dystrophy} An inherited muscle disease caused by reduced or abnormal dystrophin, a protein that supports muscle fibres. It usually causes slowly progressive weakness beginning in childhood or adulthood and may also affect the heart.

\medterm{Becker Myotonia} An inherited muscle condition causing stiffness because muscles relax slowly after contraction. It often begins in childhood, may improve after repeated movement, and can cause muscle enlargement, difficulty starting movement, or brief weakness.

\synonyms
Autosomal Recessive Myotonia Congenita; Recessive Generalized Myotonia

\medterm{Becker Nevus} A usually harmless, irregularly shaped patch of darker skin that often becomes noticeable during adolescence and may develop increased hair growth. It commonly affects one shoulder or the upper trunk and is not usually cancerous.

\medterm{Becker's Muscular Dystrophy} An inherited muscle disease causing progressive weakness, usually beginning later and advancing more slowly than Duchenne muscular dystrophy. It can also weaken the heart and requires regular heart assessment.

\medterm{Becker’S Nevus} A usually harmless darkened patch of skin, often with extra hair, that commonly appears on one shoulder or the upper trunk during adolescence. It is not cancer, but changes or uncertainty require assessment.

\medterm{Beckwith-Wiedemann Syndrome} An overgrowth condition present from birth that may cause a large body or tongue, abdominal-wall defects, enlarged organs, low blood sugar in newborns, and increased risk of certain childhood cancers. Regular specialist monitoring is needed.

\medterm{Beclomethasone} Beclomethasone is a corticosteroid used by inhalation, nasal spray, or topical routes to reduce inflammation in asthma, rhinitis, or skin disease.

\medterm{Becquerel} The unit used to measure radioactivity. One becquerel means that one radioactive atom changes each second.

\medterm{Becquerel} The International System (SI) derived unit of radionuclide activity, defined as one nuclear decay or disintegration per second ($1\text{ Bq} = 1\text{ s}^{-1}$). Clinical diagnostic and therapeutic radioisotope doses are typically quantified in megabecquerels (MBq) or gigabecquerels (GBq).

\synonyms
Bq; SI Unit of Radioactivity

\medterm{Bed Rest} Staying in bed or remaining mostly reclined to limit activity. It is used only for selected medical reasons because prolonged inactivity can cause muscle loss, blood clots, constipation, pressure injuries, and reduced independence.

\medterm{Bed Rest During Pregnancy} Restricting activity or staying in bed during pregnancy for a specific medical reason. Routine bed rest is generally avoided because it has not prevented many pregnancy complications and can increase muscle loss and blood-clot risk.

\medterm{Bedbug} A small, flat, blood-feeding insect that hides in bedding, furniture, and cracks. Its bites can cause itchy skin bumps or allergic reactions, but bedbugs are not known to routinely spread human disease.

\medterm{Bedbug Bites} Itchy skin reactions caused by bedbugs feeding at night. Bites often appear as red bumps in clusters or lines on exposed skin, sometimes days after exposure; scratching can cause skin infection.

\medterm{Bedbug Bites} Alternate name; see \textit{Bed Bugs}.

\medterm{Bedbugs} Small blood-feeding insects that cause itchy bites and skin reactions and can spread through bedding, clothing, and furniture.

\medterm{Bednar Aphthae} Shallow sores on the back of a baby's palate caused by pressure or rubbing, often during feeding. They may make feeding painful and usually improve when the source of irritation is corrected.

\medterm{Bedpan} A container used by a person who cannot safely get out of bed to urinate or pass stool. It is positioned securely, privacy is maintained, and the skin is cleaned and dried afterward.

\medterm{Bedsore} An area of skin and deeper-tissue damage caused by prolonged pressure, usually over a bony area such as the heel, hip, or tailbone. Limited movement, moisture, poor nutrition, and reduced sensation increase risk.

\medterm{Bedsores} Localized ischemic injury to the skin and underlying subcutaneous tissue, muscle, or bone resulting from prolonged unrelieved pressure, shear forces, and friction over bony prominences in immobilized patients.

\synonyms
Pressure Ulcers; Decubitus Ulcers; Pressure Sores; Decubiti

\medterm{Bedtime Habits For Infants And Children} Regular bedtime routines that help children sleep safely and consistently. They may include a predictable wind-down period, suitable sleep conditions, and age-appropriate limits on screens; persistent sleep problems need assessment.

\medterm{Bedwetting} Alternate name; see \textit{Enuresis}.

\medterm{Bee And Wasp Sting} A sting from a bee or wasp that injects venom and usually causes immediate pain, redness, swelling, and itching. Trouble breathing, throat swelling, widespread hives, faintness, or vomiting may indicate anaphylaxis and requires emergency treatment.

\medterm{Bee Sting} A venom-injecting injury from a bee that usually causes local pain, redness, and swelling. A severe allergic reaction can cause widespread hives, breathing difficulty, faintness, or swelling of the mouth and throat.

\synonyms
Apis Envenomation; Honeybee Sting

\medterm{Bee Stings} Injuries caused by bees injecting venom into the skin. They may cause local pain and swelling, a large reaction extending beyond the sting, or life-threatening anaphylaxis in an allergic person.

\medterm{Beef or Pork Tapeworm Infection} Intestinal infection with Taenia saginata or Taenia solium acquired from undercooked beef or pork; T. solium eggs can also cause cysticercosis in tissues.
\medterm{Beer-Drinker's Heart} Alternate name; see \textit{Alcoholic Cardiomyopathy}.

\medterm{Beers Criteria} A regularly updated list of medicines that may cause extra harm in some older adults, such as confusion, falls, or bleeding. It supports medication review but does not replace individualized clinical judgment.

\medterm{Beeswax Poisoning} Illness after ingesting or inhaling a large amount of beeswax. Beeswax is generally of low toxicity, but large ingestions may cause gastrointestinal blockage or aspiration, especially in children.

\medterm{Behavior} The observable actions, responses, and activities of an organism, influenced by genetic, neurobiological, psychological, and environmental factors. In medicine, behavioral assessment helps diagnose psychiatric, developmental, and neurological conditions.

\medterm{Behavior Disorder} A persistent pattern of actions or emotional responses that causes significant distress, interferes with daily life, or harms relationships. Causes may be developmental, psychological, medical, or environmental.

\medterm{Behavior Test} A structured assessment of behavior, mood, attention, thinking, or social functioning. It may use observation, interviews, questionnaires, or standardized tasks to help identify developmental, psychiatric, or neurologic concerns.

\medterm{Behavior Therapy} A form of psychotherapy based on learning principles that uses reinforcement, extinction, and other behavioral techniques to modify maladaptive behaviors. It is effective for anxiety disorders, phobias, obsessive-compulsive disorder, and habit disorders.

\medterm{Behavioral Abnormality} A deviation from typical patterns of behavior that may indicate a psychiatric, neurodevelopmental, or neurologic condition. It is assessed by history, observation, standardized measures, and developmental norms to determine clinical significance.

\medterm{Behavioral Disorder} A persistent pattern of behavior that causes clinically significant distress or impairment in daily functioning, relationships, learning, or safety. The term is broad; interpretation considers age, development, mental health, medical conditions, environment, trauma, and substance exposure.

\medterm{Behavioral Genetics} The study of how genes and environmental factors influence behavior and mental traits. Most behaviors reflect many genes interacting with life experiences, not one gene alone.

\medterm{Behavioral Rehabilitation} Structured therapy that helps a person change behaviors, build coping skills, and regain daily function after illness, injury, addiction, or disability. It may include counseling, goal setting, and practice of useful routines.

\medterm{Behcet Disease} A long-term inflammatory illness that repeatedly causes painful mouth sores and genital sores. It may also inflame the eyes, skin, joints, blood vessels, brain, or digestive tract and can threaten vision or other organs.

\medterm{Behcet Syndrome} Another name for Behcet disease, an inflammatory illness that causes repeated mouth and genital sores and may affect the eyes, skin, joints, blood vessels, brain, or digestive tract.

\medterm{Behcet's Syndrome} Behcet syndrome is a relapsing inflammatory disorder that can cause recurrent mouth and genital ulcers, eye inflammation, skin lesions, and vascular, neurologic, or gastrointestinal disease. Treatment is tailored to the organs involved and may include immunosuppression.

\synonyms
Behçet Disease; Silk Road Disease; Adamantiades-Behçet Syndrome

\medterm{Behr Syndrome} A rare inherited neurological disorder causing progressive optic-nerve damage, poor vision, muscle stiffness, unsteady movement, and sometimes weakness or bladder problems.

\medterm{Behçet Disease} Standard spelling of Behcet disease, a relapsing inflammatory disorder that can affect the mouth, genitals, eyes, skin, joints, blood vessels, nervous system, or digestive tract.

\medterm{Behçet's Syndrome} Behçet's syndrome is a relapsing inflammatory disorder causing recurrent oral and genital ulcers, eye inflammation, skin lesions, and sometimes vascular, neurologic, or gastrointestinal disease.

\medterm{Beijerinck, Martinus W.} A Dutch microbiologist who helped establish the concept of viruses as infectious agents that reproduce only in living cells and contributed to the development of enrichment culture methods.

\medterm{Bejel} Bejel is endemic treponematosis, a chronic infection caused by Treponema pallidum pertenue or related treponemes and spread mainly by close nonsexual contact in endemic areas.

\medterm{Bekam} A Malay name for cupping, a practice in which cups create suction on the skin; wet cupping may also involve superficial skin incisions.

\medterm{Belching} The release of swallowed air or gas from the stomach through the mouth. Frequent
belching may occur with aerophagia, gastroesophageal reflux, functional dyspepsia, or
carbonated-drink consumption.

\medterm{Belinostat} A cancer medicine that blocks enzymes involved in turning genes on and off. It is used for some peripheral T-cell lymphomas that have returned or resisted earlier treatment.

\medterm{Bell Palsy} Sudden weakness or paralysis of one side of the face caused by dysfunction of the facial nerve. The eyelid may not close fully, and taste or hearing may change. Most cases improve, but urgent assessment is needed to distinguish it from stroke and to protect the eye.

\synonyms
Bell's Palsy

\medterm{Bell's Palsy} Acute peripheral facial-nerve weakness of uncertain cause, usually affecting one side of the face. Eye protection and timely clinical assessment are important; treatment is individualized.

\synonyms
Idiopathic Facial Palsy; Facial Nerve Palsy

\medterm{Bell's Phenomenon} Upward and outward movement of the eyes when a person closes the eyelids forcefully. This normal protective reflex helps shield the cornea and may become visible when facial-nerve weakness prevents complete eyelid closure.

\medterm{Bell-Clapper Deformity} A birth-related difference in how a testicle is attached inside the scrotum, allowing it to move and twist more freely than usual. It is often present on both sides and increases the risk of testicular torsion, a sudden emergency that cuts off blood flow.

\synonyms
Bell-Clapper Anomaly; High Tunica Vaginalis Insertion

\medterm{Belladonna} Atropa belladonna, the deadly nightshade plant, whose tropane alkaloids such as atropine and scopolamine can cause anticholinergic poisoning.

\medterm{Bells Palsy} Alternate name; see \textit{Bell's Palsy}.

\medterm{Belly} The abdomen, the region of the body between the chest and pelvis containing much of the digestive, urinary, and reproductive organs.

\medterm{Below-Knee Amputation} Surgery to remove the lower leg while preserving the knee joint. It may be needed for severe injury, infection, poor blood flow, or cancer; rehabilitation focuses on healing, strength, mobility, and a prosthesis when suitable.

\medterm{Bence Jones Protein} An abnormal antibody fragment made by some plasma-cell disorders. It may appear in blood or urine and help identify conditions such as multiple myeloma, but it can also damage the kidneys.

\medterm{Bence-Jones Protein} An abnormal antibody fragment made by a single group of plasma cells. Finding it in urine or blood may support evaluation for multiple myeloma or another plasma-cell disorder and possible kidney injury.

\medterm{Benchmarking} Comparing healthcare results, safety, costs, or processes with a standard or similar service. The comparison can identify areas for improvement but does not by itself prove that one service causes better outcomes.

\medterm{Bendamustine} A cancer medicine that damages tumour-cell DNA and is used for selected blood cancers. It can also damage healthy bone-marrow cells, causing low blood counts, infection, tiredness, nausea, or other side effects.

\medterm{Bendamustine Hydrochloride} The hydrochloride form of bendamustine, a chemotherapy medicine used for selected blood cancers. It damages DNA and may cause low blood counts, infection, infusion reactions, nausea, tiredness, or other side effects.

\medterm{Bendiocarb} A carbamate insecticide that can overstimulate nerves after significant exposure. Poisoning may cause sweating, salivation, vomiting, diarrhoea, small pupils, muscle twitching, breathing difficulty, or weakness and requires urgent medical treatment.

\medterm{Bendroflumethiazide} A diuretic medicine that helps the kidneys remove salt and water. It is used for high blood pressure or fluid retention and may lower potassium or sodium and increase uric acid or blood sugar.

\medterm{Bends} A form of decompression sickness caused by gas bubbles forming in the body when pressure falls too quickly, usually during or after diving. It commonly causes deep joint or muscle pain but can also injure the brain, spinal cord, lungs, or circulation.

\medterm{Benedikt Syndrome} A brainstem vascular stroke syndrome resulting from a paramedian lesion of the midbrain tegmentum involving the oculomotor nerve fascicles, red nucleus, and substantia nigra. It is clinically characterized by ipsilateral third cranial nerve palsy (ptosis, dilated pupil, downward-and-outward eye) and contralateral involuntary choreoathetosis, intention tremor, and ataxia.

\synonyms
Paramedian Midbrain Syndrome

\medterm{Benedikt’s Syndrome} A rare pattern of neurological problems caused by damage in the midbrain. It may cause a drooping eyelid or double vision on one side and tremor, poor coordination, or involuntary movements on the opposite side.

\medterm{Beneficence And Non-Maleficence} Two principles of medical ethics: beneficence means acting to benefit the patient, while non-maleficence means avoiding unnecessary harm. Good decisions balance expected benefits, risks, patient values, and available alternatives.

\medterm{Benefits And Risks Of Taking Amino Acid Supplement} Amino-acid supplements may help when a documented deficiency or specific medical need exists, but most people obtain enough from food. Unnecessary or excessive use can cause stomach symptoms, interactions, or metabolic, kidney, or liver problems.

\medterm{Benfotiamine} A fat-soluble form of thiamine, or vitamin B1, sold in some supplements and studied for diabetic nerve damage. Evidence of benefit is limited and it does not replace proven treatment for diabetes or neuropathy.

\medterm{Benfuracarb} A carbamate insecticide that can overstimulate nerves after significant exposure. It may cause sweating, salivation, vomiting, diarrhoea, small pupils, muscle twitching, breathing difficulty, or weakness and requires urgent medical treatment.

\medterm{Benign} Not cancerous. A benign growth does not invade nearby tissues or spread to other parts of the body, but it can still cause problems because of its size, location, or hormone production.

\synonyms
Nonmalignant; Noncancerous

\medterm{Benign Adrenal Tumors} Noncancerous growths in an adrenal gland. Some are found by chance and cause no problems, while others release excess hormones or become large enough to cause pressure symptoms.

\medterm{Benign Ameloblastoma} A usually slow-growing tumour arising from cells involved in tooth development, most often in the jaw. It is not cancer but can expand, damage bone, and displace teeth, so treatment usually requires specialist surgery.

\medterm{Benign Brain Tumor} A noncancerous growth in or around the brain. Although it does not spread like cancer, it can cause headache, seizures, vision or movement problems, or fluid blockage by pressing on nearby structures.

\medterm{Benign Chondroblastoma} A noncancerous cartilage-producing bone tumour that usually develops near the end of a long bone in an adolescent or young adult. It often causes pain, swelling, or reduced movement near a joint.

\medterm{Benign CNS Cyst} A noncancerous, fluid-filled sac in or around the brain or spinal cord. Many cause no symptoms, but a large or strategically located cyst can press on nerves or block the normal flow of brain fluid.

\medterm{Benign Ear Cyst Or Tumor} A noncancerous lump or fluid-filled sac in the outer, middle, or inner ear. It may cause no symptoms or lead to pain, discharge, hearing loss, dizziness, or pressure depending on its location and size.

\medterm{Benign Epithelioma} A noncancerous growth arising from cells that line the skin or an organ. It does not invade or spread, but removal or testing may be needed if it bleeds, grows, causes symptoms, or is difficult to distinguish from cancer.

\medterm{Benign Esophageal Stricture} A noncancerous narrowing of the esophageal lumen, most commonly resulting from chronic acid reflux injury (peptic stricture), corrosive ingestion, radiation, or prior surgery, causing progressive dysphagia.

\synonyms
Peptic Stricture; Esophageal Stenosis, Benign; Esophageal Stricture

\medterm{Benign Hemangiopericytoma} An older term for a usually slow-growing tumour with branching blood vessels, now often classified as a solitary fibrous tumour. Behaviour varies, so tissue examination and follow-up are important.

\medterm{Benign Intracranial Hypertension} Historical and misleading name; see \textit{Idiopathic Intracranial Hypertension}.

\medterm{Benign Liver Tumors} Noncancerous growths in the liver, including hemangiomas, focal nodular hyperplasia, and hepatocellular adenomas. Most are found by chance on ultrasound, CT, or MRI and do not require treatment. Some may cause pain, liver enlargement, bleeding, or rarely change into cancer.

\medterm{Benign Melanocytic Lesion} An older or imprecise term for a benign melanocytic growth, such as a melanocytic nevus or mole; unlike malignant melanoma, it is not cancerous.

\medterm{Benign Meningioma} A usually slow-growing, noncancerous tumour arising from the membranes around the brain or spinal cord. It may cause seizures, headache, weakness, or vision problems by pressing on nearby structures.

\medterm{Benign Mesenchymoma} A rare noncancerous tumour containing more than one type of tissue such as fat, cartilage, or bone. Its symptoms depend on the site and size, and diagnosis requires examination of tissue.

\medterm{Benign Migratory Glossitis} Alternate name; see \textit{Geographic Tongue}.

\medterm{Benign Mixed Tumor} A usually noncancerous salivary-gland tumour containing more than one tissue type. It often presents as a slow-growing, painless lump, but long-standing or recurrent growths require assessment because rare malignant change can occur.

\medterm{Benign Mullerian Papilloma} A rare noncancerous, wart-like growth in the vulva or vagina formed from Müllerian-type lining cells. It often appears as a small polyp and is diagnosed by examining a tissue sample.

\medterm{Benign Myoepithelioma} A noncancerous tumour of myoepithelial cells, often arising in a salivary gland or other glandular tissue. It usually grows slowly and may cause a painless lump; diagnosis depends on tissue examination.

\medterm{Benign Neoplasm} A noncancerous growth of cells that remains at its original site and does not spread to distant organs. It may still cause harm by pressing on nearby structures, blocking a passage, bleeding, or producing hormones.

\medterm{Benign Orgasmic Headache} A sudden severe headache during sexual activity or orgasm after bleeding, stroke, vessel disease, and other dangerous causes have been excluded. A first or unusually severe headache needs urgent medical assessment.

\synonyms
Primary Headache Associated with Sexual Activity; Sexual-Activity Headache

\medterm{Benign Paraganglioma} A usually slow-growing tumour of specialized nerve-related cells. It may release adrenaline-like hormones or press on nearby nerves and blood vessels; symptoms and treatment depend on its site and activity.

\medterm{Benign Paroxysmal Positional Vertigo} A common inner-ear disorder causing brief spinning attacks when the head changes position. Tiny calcium particles move into a balance canal and trigger dizziness, often lasting less than a minute.

\medterm{Benign Peripheral Nerve Tumor} A noncancerous tumour arising from or beside a nerve outside the brain and spinal cord. It may cause a lump, pain, tingling, numbness, weakness, or no symptoms at all.

\medterm{Benign Positional Vertigo} Alternate name; see \textit{Benign Paroxysmal Positional Vertigo}.

\medterm{Benign Prostatic Hyperplasia} Noncancerous enlargement of the prostate that can narrow the urine channel. It may cause a weak stream, difficulty starting, urgency, frequent night-time urination, dribbling, or incomplete emptying.

\medterm{Benign Prostatic Obstruction} Blockage at the bladder outlet caused by an enlarged prostate or increased muscle tone in the prostate. It may cause a weak or interrupted stream, straining, incomplete emptying, urinary retention, or recurrent infection.

\medterm{Benign Rolandic Epilepsy} A childhood seizure disorder, now often called self-limited epilepsy with centrotemporal spikes. Brief seizures commonly affect the face or mouth during sleep and may cause twitching, tingling, drooling, or temporary trouble speaking; it usually ends by adolescence.

\synonyms
Self-Limited Epilepsy with Centrotemporal Spikes; SeLECTS; Benign Childhood Epilepsy with Centrotemporal Spikes; Rolandic Epilepsy

\medterm{Benign Tumor} A noncancerous growth that remains localized and does not invade nearby tissue or spread to distant organs. It may still cause problems by pressing on structures, blocking a passage, bleeding, or making hormones.

\medterm{Benign Urethral Lesions} Noncancerous growths or structural changes in the urethra, such as polyps, caruncles, papillomas, or inflammatory lesions. They may cause bleeding, discharge, pain, narrowing, or difficulty passing urine, but some cause no symptoms.

\synonyms
Urethral Polyps; Urethral Caruncle; Benign Urethral Tumors

\medterm{Bennett Fracture} An intra-articular fracture-subluxation of the base of the first metacarpal bone. The smaller medial articular fragment remains attached to the trapezium by the anterior oblique ligament, while the main metacarpal shaft is displaced proximally, laterally, and dorsally by the pulling force of the abductor pollicis longus tendon.

\medterm{Benserazide} A medicine combined with levodopa in some Parkinson treatments. It blocks levodopa from being changed into dopamine outside the brain, allowing more to reach the brain and reducing some unwanted effects.

\medterm{Bent Penis} Curvature of the penis that may be present from birth or develop later. It can be painless or cause pain, difficulty with erection or intercourse, and sometimes problems passing urine.

\medterm{Bent Spine Syndrome} Severe forward bending of the trunk that improves when lying down. It results from weakness or abnormal contraction of back muscles and can occur with muscle, nerve, or movement disorders.

\synonyms
Camptocormia

\medterm{Bentall Procedure} Open-heart surgery that replaces a diseased aortic valve, the first part of the aorta, and the aortic root with an artificial valve and tube graft. The coronary arteries are reattached to the graft; lifelong monitoring and sometimes blood-thinning treatment are needed.

\synonyms
Bentall Operation; Composite Aortic Root Replacement

\medterm{Benzaldehyde} A chemical with an almond-like smell used in flavourings, fragrances, and manufacturing. Concentrated exposure can irritate the eyes, skin, or lungs; it is not a routine medical treatment.

\medterm{Benzazepines} A family of chemical structures made from joined benzene and seven-membered nitrogen-containing rings. Individual compounds may have different biological or medicinal effects, so the family name alone does not identify a treatment.

\medterm{Benzene} An industrial chemical and solvent found in some fuels and smoke. High short-term exposure can affect the brain and heart; repeated exposure can damage bone marrow and increase the risk of leukaemia.

\medterm{Benzene Poisoning} Illness caused by breathing, swallowing, or absorbing benzene. Short-term exposure may cause dizziness, drowsiness, or irregular heartbeat; repeated exposure can damage bone marrow and increase the risk of leukaemia.

\medterm{Benzhexol} Benzhexol, or trihexyphenidyl, is an anticholinergic drug used mainly to reduce tremor and rigidity in Parkinson disease or drug-induced extrapyramidal symptoms.

\medterm{Benzo[a]pyrene} A polycyclic aromatic hydrocarbon (PAH) formed during the incomplete combustion of organic matter, including tobacco smoke, vehicle exhaust, and grilled meats. Its active metabolite, benzo[a]pyrene-7,8-dihydrodiol-9,10-epoxide (BPDE), forms mutagenic DNA adducts that drive chemical carcinogenesis.

\synonyms
BaP; 3,4-Benzopyrene; Benzo(a)pyrene

\medterm{Benzoates} Salts or chemical forms of benzoic acid used mainly as food preservatives or ingredients in some products. Their effects and safety depend on the specific compound and amount; some people develop irritation or sensitivity.

\medterm{Benzocaine} A local anaesthetic applied to skin or moist body surfaces to reduce sensation. Rarely, especially after excessive use, it can cause methemoglobinaemia, a condition in which blood cannot carry oxygen normally.

\medterm{Benzodiazepine} A medicine that enhances the effect of the inhibitory chemical messenger GABA in the brain. It may reduce anxiety, induce sleep, relax muscles, or control seizures, but can cause sedation, impaired coordination, dependence, and dangerous breathing suppression when combined with opioids or other sedatives.

\medterm{Benzodiazepine Use Disorder} A pattern of benzodiazepine use that causes loss of control, harm, or problems at work, home, or in relationships. Regular use can cause dependence, and sudden stopping can produce life-threatening withdrawal.

\medterm{Benzodiazepine Withdrawal} Symptoms after reducing or stopping regular benzodiazepine use. Anxiety, insomnia, shaking, sweating, and sensitivity to light may occur; severe withdrawal can cause seizures, hallucinations, or delirium and requires supervised tapering.

\medterm{Benzopyran} A chemical ring structure formed by joining benzene and pyran rings. It occurs in some natural substances and medicines; the term describes chemistry rather than one specific disease or treatment.

\medterm{Benzopyrene} A group of chemicals produced by incomplete burning and found in tobacco smoke, vehicle exhaust, and heavily charred food. Some forms can damage DNA and increase cancer risk after the body changes them chemically.

\medterm{Benzopyrenes} Chemicals with a benzopyrene structure that form during incomplete combustion. Several occur in smoke and pollution, can damage DNA, and are linked with increased cancer risk.

\medterm{Benzoyl Peroxide} A topical medicine that kills acne-causing bacteria and helps remove blocked skin cells. It can cause dryness, peeling, irritation, or bleaching of hair and fabrics.

\medterm{Benzphetamine} A sympathomimetic amine and stimulant formerly used in some countries as a short-term adjunct for weight reduction; it can cause dependence, cardiovascular effects, and insomnia and is restricted or unavailable in many jurisdictions.

\medterm{Benztropine} An anticholinergic medicine used mainly for movement problems caused by some medicines and for acute muscle spasms called dystonia. It can cause dry mouth, blurred vision, constipation, urinary retention, confusion, or worsen glaucoma.

\medterm{Benzydamine} A topical anti-inflammatory pain-relieving medicine used for painful inflammation in the mouth, throat, or skin. It may cause temporary numbness, irritation, or an allergic reaction.

\medterm{Benzyl Benzoate} A topical chemical used to kill mites causing scabies and, in some settings, lice. It can irritate the skin and is applied as directed, with care to avoid the eyes, mouth, and broken skin.

\medterm{Benzylamines} Chemical compounds containing a benzyl group joined to an amino group. Some are used to make medicines or other chemicals, while their effects depend on the individual compound.

\medterm{Benzylpenicillin} A penicillin antibiotic used for serious infections caused by susceptible bacteria, including some streptococcal infections, meningitis, and syphilis. Allergic reactions can be severe.

\medterm{BEP Regimen} A chemotherapy combination of bleomycin, etoposide, and cisplatin, commonly used for some germ-cell cancers. It can cause low blood counts, infection, kidney damage, hearing changes, nerve symptoms, or lung injury and requires monitoring.

\medterm{Berberine} A plant compound sold in some supplements and studied for blood-sugar, cholesterol, and antimicrobial effects. Evidence varies, and it can cause stomach symptoms, low blood sugar, drug interactions, or harm during pregnancy.

\medterm{Bereavement} The emotional response to the death of someone important. Grief may include sadness, yearning, anger, numbness, sleep or appetite changes, and waves of distress; it varies widely and can sometimes lead to prolonged grief or depression.

\medterm{Bereavement Reaction} Emotional and physical responses after someone dies, such as sadness, disbelief, anger, poor sleep, or loss of appetite. These reactions are common and do not automatically indicate a mental disorder.

\medterm{Berger Disease} An older name for IgA nephropathy, a kidney disease in which an immune protein builds up in kidney filters. It may cause episodes of blood in the urine after infections, protein in the urine, high blood pressure, and kidney failure.

\medterm{Berger's Disease} Alternate name; see \textit{IgA Nephropathy}.

\medterm{Beriberi} Illness caused by severe vitamin B1 deficiency. It may damage nerves and cause weakness, numbness, or difficulty walking, or affect the heart and cause swelling, fast heartbeat, and heart failure; urgent vitamin replacement is needed.

\medterm{Bernard-Soulier Disease} A rare inherited bleeding disorder in which platelets are unusually large and do not stick properly to damaged blood vessels. It causes easy bruising, nosebleeds, heavy periods, and prolonged bleeding after injury or procedures.

\medterm{Bernard–Soulier Syndrome} A rare inherited bleeding disorder in which platelets are unusually large and cannot attach normally to injured blood vessels. It causes easy bruising, nosebleeds, heavy periods, and prolonged bleeding after injury or procedures.

\synonyms
Giant Platelet Disorder; Platelet-Type Bleeding Disorder 1

\medterm{Berry Aneurysm} A small, rounded balloon-like bulge at a branching point of an artery at the base of the brain. If it bursts, it can cause life-threatening bleeding around the brain.

\medterm{Berylliosis} A long-term lung disease caused by sensitivity to inhaled beryllium dust or fumes. It can cause cough, breathlessness, fatigue, and scarring or inflammation that resembles sarcoidosis.

\medterm{Beryllium} A light metal used in some industrial products. Breathing its dust or fumes can cause allergic lung inflammation, long-term lung disease, and increased cancer risk, especially after repeated workplace exposure.

\medterm{Beta 2-Microglobulin} A small protein released by most cells. Blood or urine levels may rise with inflammation, some blood cancers, or reduced kidney removal; the result is interpreted with symptoms and other tests.

\medterm{Beta Amyloid} Small protein fragments made from amyloid precursor protein. They can build up into plaques in the brain and are associated with Alzheimer disease, although plaques alone do not explain every symptom.

\synonyms
A$\beta$ Peptide; Amyloid-Beta Protein

\medterm{Beta Carotene} A yellow-orange plant pigment that the body can convert into vitamin A. It is found in foods such as carrots and leafy vegetables; excessive supplements can turn the skin yellow-orange.

\synonyms
Provitamin A Carotenoid; Carotenoid Provitamin

\medterm{Beta Cell} A cell in the pancreas that makes and releases insulin when blood glucose rises. Insulin helps body cells use glucose and keeps blood glucose from becoming too high.

\medterm{Beta Cells} Pancreatic cells that produce insulin and amylin. Insulin lowers blood glucose by helping cells take up glucose, while amylin slows the rise after meals; loss or failure of these cells contributes to diabetes.

\medterm{Beta Globulin} A group of blood proteins identified together during a protein blood test. They include proteins that carry iron or fats and support immunity; abnormal levels can occur with inflammation, liver disease, or immune disorders.

\medterm{Beta Glucans} Complex carbohydrates found in oats, barley, yeast, fungi, and some bacteria. Oat and barley forms can modestly lower cholesterol when eaten regularly; effects on immunity depend on the source and preparation.

\medterm{Beta Interferons} Immune-modifying medicines used to reduce relapses and new brain or spinal-cord lesions in some forms of multiple sclerosis. Common effects include flu-like symptoms and injection-site reactions; monitoring may detect liver or blood-count abnormalities.

\medterm{Beta Particle} A fast-moving electron or positively charged equivalent released by some radioactive atoms. It can damage nearby tissue and is therefore used carefully in selected radiation treatments and laboratory applications.

\medterm{Beta Radiation} Radiation made of fast electrons or their positively charged counterparts released by unstable atoms. It can injure living tissue but can also be directed at selected tissues for treatment or used in medical testing.

\medterm{Beta Thalassemia} An inherited blood disorder in which the body makes too little beta-globin, a part of haemoglobin. It can cause mild or severe anaemia, tiredness, jaundice, bone changes, or an enlarged spleen.

\medterm{Beta Tocopherol} One naturally occurring form of vitamin E. It acts as a fat-soluble antioxidant, but the body generally uses alpha-tocopherol more effectively; its specific clinical importance is still being studied.

\medterm{Beta-Adrenoceptor-Blocking Drugs} Medicines that block adrenaline-like signals at beta receptors. They can slow the heart, lower blood pressure, reduce tremor, and protect the heart in selected conditions, but may cause fatigue, low pulse, or breathing problems in susceptible people.

\medterm{Beta-Alanine} An amino acid that helps muscles make carnosine, which buffers acid during intense activity. Supplements may modestly improve some short, high-intensity exercise but commonly cause temporary tingling of the skin.

\medterm{Beta-Blockers} Medicines that block adrenaline-like effects, usually slowing the heartbeat and lowering blood pressure. They are used for several heart conditions, migraine prevention, tremor, glaucoma, and other problems; some can worsen asthma or mask low blood sugar.

\medterm{Beta-Blockers Overdose} Poisoning from taking too much of a beta-blocker. It can cause a very slow heartbeat, low blood pressure, wheezing, low blood sugar, confusion, seizures, or shock and requires emergency assessment.

\medterm{Beta-Carotene} A plant pigment that the body can convert into vitamin A. It occurs in foods such as carrots and leafy greens; high-dose supplements may increase lung-cancer risk in smokers and are not a substitute for a balanced diet.

\medterm{Beta-Carotene Blood Test} A blood test measuring beta-carotene, a plant pigment that can become vitamin A. A low result may suggest poor nutrition or fat absorption; a high result may cause harmless yellow-orange skin.

\medterm{Beta-Chemokines} Small immune signals that attract white blood cells to areas of infection or inflammation. They help coordinate defense, but excessive activity can contribute to long-lasting inflammation and tissue damage.

\medterm{Beta-Endorphin} A natural pain-relieving substance made in the brain and pituitary gland. It acts on opioid receptors and also influences stress responses, mood, and some hormone functions.

\medterm{Beta-Glucans} Complex carbohydrates found in foods such as oats, barley, mushrooms, and yeast. Their effects vary by source; some provide fibre and may modestly affect cholesterol, blood glucose, or immune activity.

\medterm{Beta-Haemolytic Streptococcus} A streptococcus that completely breaks down red blood cells around its colonies in a laboratory culture. Important group A and group B species can cause throat, skin, bloodstream, newborn, or other serious infections.

\medterm{Beta-Human Chorionic Gonadotrophin} The measurable part of the pregnancy hormone hCG. Blood or urine testing helps detect pregnancy and can help assess some pregnancy problems or selected tumours when interpreted with other findings.

\medterm{Beta-Lactam Antibiotics} A large group of antibiotics that stop susceptible bacteria building their protective cell wall. They include penicillins, cephalosporins, carbapenems, and related medicines; allergic reactions and bacterial resistance can limit their use.

\medterm{Beta-Lactamase} An enzyme made by some bacteria that breaks down certain penicillin-like antibiotics and makes them ineffective. Laboratory testing can identify this resistance, and some medicines block selected beta-lactamases.

\medterm{Beta-Mannosidosis} A rare inherited disorder in which cells cannot break down certain complex sugars. The stored material can cause developmental delay, hearing loss, repeated infections, seizures, or skeletal and skin problems.

\medterm{Beta-Oxidation} Alternate form; see \textit{Fatty Acid Oxidation}.

\medterm{Beta-Sheet} A common folded shape in proteins in which extended strands lie beside one another and are held by chemical bonds. Excessive beta-sheet folding can contribute to abnormal protein deposits in some diseases.

\medterm{Beta-Thalassemia} An inherited disorder in which the body makes too little or no beta-globin, a part of haemoglobin. Severity ranges from mild carrier status to severe anaemia requiring regular transfusions and monitoring for iron overload.

\medterm{Beta-Thalassemia Intermedia} A moderately severe inherited form of beta-thalassemia causing ongoing anaemia, enlarged spleen, jaundice, or bone changes. People may need occasional transfusions, but usually do not require regular lifelong transfusions.

\medterm{Beta-Thalassemia Minor} The carrier form of beta-thalassemia, caused by one altered beta-globin gene. It usually causes no symptoms or mild small-cell anaemia and is not treated with iron unless iron deficiency is also confirmed.

\synonyms
Beta-Thalassemia Trait; Cooley Trait

\medterm{Beta-Thalassemia Trait} The heterozygous carrier state for a beta-globin gene variant. It is usually asymptomatic or causes mild microcytosis with mild or no anemia; hemoglobin electrophoresis often shows increased HbA2. It is also called beta-thalassemia minor.

\synonyms
Beta-Thalassemia Minor; Thalassemia Minor; Cooley Trait

\medterm{Betacoronavirus} A group of coronaviruses that includes viruses causing respiratory illnesses in people and animals. Individual viruses differ in spread and severity; some cause mild colds, while others can cause severe pneumonia.

\medterm{Betaherpesvirinae} A subgroup of herpes viruses that includes cytomegalovirus and human herpesviruses 6 and 7. They remain in the body after infection and can reactivate, especially when immunity is weakened.

\medterm{Betaine} A natural substance involved in processing homocysteine and helping cells maintain their water balance. It is also sold as a supplement, but high doses can cause side effects or interact with medicines.

\medterm{Betamethasone} A strong corticosteroid medicine that reduces inflammation and immune activity. It may be applied to skin or given in other forms; prolonged or extensive use can thin the skin, worsen infections, or suppress the body's steroid production.

\medterm{Betamethasone Valerate} A corticosteroid applied to the skin to reduce inflammation and itching in conditions such as eczema or psoriasis. Excessive, prolonged, or covered application can thin the skin and increase absorption into the body.

\medterm{Betaxolol} A relatively beta-1-selective blocker used orally for some cardiovascular conditions and topically to lower intraocular pressure in glaucoma. It may cause bradycardia, hypotension, or bronchospasm.

\medterm{Bethanechol} Bethanechol is a muscarinic cholinergic agonist that stimulates bladder and gastrointestinal smooth muscle, used selectively for urinary retention or impaired motility.

\medterm{Bethesda System For Reporting Thyroid Cytopathology} A six-tier standardized system for reporting thyroid fine-needle aspiration results, from nondiagnostic and benign categories to follicular-pattern, suspicious, and malignant categories; each tier guides risk assessment and management.

\medterm{Bethlem Myopathy} A usually slowly progressive inherited muscle disease caused by changes in collagen VI. It causes weakness, often beginning in childhood, together with tight joints or shortened muscles, especially in the fingers, elbows, ankles, or spine; many people remain able to walk into adulthood.

\synonyms
Collagen VI-Related Myopathy (Mild Form); Benign Congenital Myopathy with Contractures

\medterm{Bevacizumab} A monoclonal antibody that blocks vascular endothelial growth factor and inhibits new blood-vessel formation. It is used for selected cancers and some eye diseases; important risks include hypertension, bleeding, thrombosis, proteinuria, impaired wound healing, and gastrointestinal perforation.

\medterm{Bexarotene} An oral retinoid X-receptor agonist used mainly for cutaneous T-cell lymphoma; it commonly causes hypertriglyceridaemia, hypothyroidism, and liver abnormalities and is teratogenic.

\medterm{Bezafibrate} Bezafibrate is a fibrate lipid-lowering drug that activates PPAR-alpha and lowers triglycerides; muscle and liver effects require monitoring.

\medterm{Bezoar} A compact mass of material that cannot be digested and becomes trapped in the stomach or intestine. It may contain hair, plant matter, medicines, or other substances and can cause pain, vomiting, blockage, or bleeding.

\medterm{Bezold-Jarisch Reflex} An inhibitory autonomic cardiovascular reflex characterized by the triad of profound bradycardia, peripheral vasodilation with severe hypotension, and transient apnea. It is triggered by the stimulation of non-myelinated vagal sensory C-fibers in the inferoposterior wall of the left ventricle by mechanical deformation, chemical agents, or acute ischemia during inferior myocardial infarction.

\medterm{Biacetyl} Biacetyl, or diacetyl, is a diketone compound with a buttery odor used in flavoring and industry; inhalational exposure can injure the small airways.

\medterm{Bias} A systematic deviation that distorts a measurement, study result, clinical judgment, or decision away from the truth or from an intended fair comparison.

\medterm{Bible Cyst} Alternate historical name; see \textit{Ganglion Cyst}.

\medterm{Bibliotherapy} The use of reading materials as a therapeutic intervention to support mental health, educate patients about conditions, or promote behavioral change. It may be used alone or alongside other treatments for depression, anxiety, or chronic disease self-management.

\medterm{Bicalutamide} A non-steroidal androgen-receptor antagonist used in prostate cancer, usually with androgen-deprivation therapy; adverse effects include hot flushes, breast enlargement, sexual dysfunction, and possible liver injury.

\medterm{Bicarbonate} A substance in the blood that helps keep the body’s acid level balanced. The lungs and kidneys regulate it, and a blood bicarbonate result can help identify acid–base or metabolic problems.

\medterm{Biceps} The biceps is a two-headed muscle, especially the biceps brachii of the upper arm, that flexes the elbow and supinates the forearm.

\medterm{Biceps Brachii Muscle} The two-part muscle on the front of the upper arm. It bends the elbow, turns the palm upward, and helps lift the arm. Injury or tendon rupture can cause weakness, pain, bruising, or a bulging “Popeye” appearance.

\medterm{Biceps Reflex} Bending of the elbow after the biceps tendon is gently tapped. This reflex tests the nerve supplying the biceps and the related spinal-cord pathways, mainly at the fifth and sixth neck levels.

\medterm{Bicipital} Bicipital means relating to a two-headed muscle, especially the biceps tendon, groove, or muscle of the upper arm.

\medterm{Bicipital Groove} The intertubercular groove of the humerus that transmits the tendon of the long head of
the biceps brachii, with the transverse humeral ligament and adjacent rotator-cuff
structures contributing to tendon stability.

\medterm{Bicipital Tuberosity} The radial tuberosity, a roughened prominence on the proximal radius where the biceps
brachii tendon inserts and produces flexion and supination of the forearm.

\medterm{Bickerstaff Brainstem Encephalitis} A rare autoimmune post-infectious central nervous system disorder closely related to the Miller Fisher variant of Guillain-Barré syndrome. It is characterized by the clinical triad of acute progressive ophthalmoplegia, severe ataxia, and central nervous system disturbance (altered consciousness or hyperreflexia), typically accompanied by serum anti-GQ1b IgG autoantibodies.

\medterm{Bicolor} Having two distinct colours. In medicine, the term may describe a skin lesion, eye, specimen, or imaging finding; its significance depends on the structure and clinical context.

\medterm{Biconcave} Curved inward on both sides. Normal red blood cells have this disc shape, which makes them flexible and provides a large surface area for carrying oxygen.

\medterm{Biconvex} Curved outward on both sides. The term may describe the shape of a lens, implant, anatomical structure, or finding on a scan.

\medterm{Bicornuate} Having two horn-like cavities or projections. A bicornuate uterus is a congenital difference in which the upper uterus forms two partly separated spaces.

\medterm{Bicornuate Uterus} A uterus with two upper cavities caused by incomplete joining during development before birth. Many people have no symptoms, but it can increase the risk of miscarriage, preterm birth, or abnormal baby position.

\medterm{Bicuspid} Having two points or flaps. The term describes some teeth, which have two main points, and some heart valves, which have two flaps instead of the usual three.

\medterm{Bicuspid Aortic Valve} A congenital aortic valve with two leaflets instead of the usual three. It can lead to
aortic stenosis, aortic regurgitation, enlargement of the aorta, or other congenital heart abnormalities.

\medterm{Bicuspid Valve} A heart valve with two leaflets rather than the usual three. The term most often refers
to congenital bicuspid aortic valve, which can cause aortic stenosis, regurgitation,
infective endocarditis, or ascending-aortic disease.

\medterm{Bid Protein} A protein that helps connect signals from outside a cell with programmed cell death inside it. When activated, it can help mitochondria release signals that lead to destruction of a damaged or unwanted cell.

\medterm{Bier Spots} Small, pale macules surrounded by an erythematous background on the extremities, typically elicited when the limbs are placed in a dependent position. They represent a benign physiological vascular anomaly caused by exaggerated local vasoconstriction in response to increased venous pressure and resolve upon limb elevation.

\medterm{Biermer's Anemia} Another name for pernicious anaemia, in which the body cannot absorb enough vitamin B12 because the stomach does not provide the needed absorption factor. It causes anaemia and can damage nerves if untreated.

\medterm{Bietti's Crystalline Dystrophy} A rare inherited eye disease in which tiny crystals and progressive damage develop in the retina and sometimes the cornea. It may cause difficulty seeing at night, loss of side vision, and gradual loss of central vision.

\medterm{Bifascicular Block} A cardiac intraventricular conduction defect in which two of the three main conduction fascicles below the bundle of His are blocked, characteristically presenting as complete right bundle branch block (RBBB) combined with either left anterior fascicular block (LAFB) or left posterior fascicular block (LPFB).

\medterm{Bifocal} Having two points of focus. Bifocal glasses contain one area for distance vision and another for near vision, allowing both prescriptions in one lens.

\medterm{Bifurcation} The point where one body structure, blood vessel, airway, nerve, or pathway divides into two branches.

\medterm{Big Toe Sign} Upward movement of the big toe, sometimes with spreading of the other toes, when the sole is stroked. In an adult, this may suggest a problem in pathways connecting the brain and spinal cord.

\medterm{Bigeminy} A heart-rhythm pattern in which each normal heartbeat is followed by an early extra beat. The extra beat often begins in the lower heart chambers and may cause palpitations or no symptoms.

\medterm{Biglycan} A protein-and-sugar component of the tissue framework around cells. It helps organize collagen and may influence inflammation, wound healing, and blood-vessel changes.

\medterm{Biguanide} A group of blood-glucose-lowering medicines that reduce glucose production by the liver and improve the body's response to insulin. Kidney function, side effects, and other health conditions guide their use.

\medterm{Biguanides} A group of medicines that lower blood glucose mainly by reducing glucose production in the liver and improving insulin sensitivity. They are used chiefly in type 2 diabetes.

\synonyms
Biguanide Drugs; Biguanide Antidiabetics

\medterm{Bilateral} Affecting both sides of the body or both members of a paired structure, such as both eyes, ears, knees, lungs, or kidneys.

\medterm{Bilateral Acoustic Neurofibromatosis} Alternate name; see \textit{Neurofibromatosis Type 2}.

\medterm{Bilateral Ankle Swelling} Alternate form; see \textit{Ankle Swelling}.

\medterm{Bilateral Gynecomastia} Enlargement of glandular breast tissue on both sides in a male. It may result from normal hormone changes, medicines, obesity, liver or kidney disease, or rarely cancer.

\medterm{Bilateral Hydronephrosis} Enlargement of the urine-collecting spaces in both kidneys because urine cannot drain normally. Causes include blockage, reflux, stones, or bladder problems; prolonged obstruction can damage kidney function.

\medterm{Bilateral Microphthalmos} A birth condition in which both eyes are unusually small. Vision may be reduced, and other eye or developmental differences may occur; assessment by an eye specialist is needed.

\medterm{Bilateral Oophorectomy} Surgery to remove both ovaries, sometimes performed for cancer, severe endometriosis, or high inherited cancer risk. Before natural menopause it causes immediate loss of ovarian hormones and infertility, so long-term effects need discussion.

\medterm{Bilateral Peritonsillar Abscesses} Collections of pus beside both tonsils, usually after a throat infection. They can cause severe sore throat, difficulty swallowing, trouble opening the mouth, muffled speech, and airway narrowing, requiring urgent hospital treatment.

\medterm{Bilateral Pneumonia} Infection or inflammation affecting both lungs. It may cause fever, cough, chest discomfort, breathlessness, or low oxygen; severity depends on the cause, extent, and person's health.

\medterm{Bilateral Tonic-Clonic Seizure} A seizure causing loss of awareness followed by stiffening and rhythmic jerking of both sides of the body. First aid aims to protect the person from injury without restraining them, and emergency help is needed if the seizure lasts five minutes or more.

\medterm{Bile} A fluid made by the liver and stored in the gallbladder that helps digest fats. It carries bile salts, cholesterol, pigments, and waste products into the intestine.

\medterm{Bile Acid Sequestrants For Cholesterol} Medicines that trap bile acids in the intestine so they leave in stool. The liver then uses cholesterol to make more bile acids, lowering LDL cholesterol; constipation and medicine interactions are common concerns.

\medterm{Bile Acids} Substances made by the liver and released into the intestine to help digest and absorb fats and vitamins A, D, E, and K. Abnormal blood levels can indicate liver or bile-flow problems.

\medterm{Bile Culture} A laboratory test that grows microorganisms from a bile sample to identify infection in the gallbladder or bile ducts and help select effective treatment.

\medterm{Bile Duct} A tube that carries bile from the liver and gallbladder into the first part of the small intestine. Blockage or inflammation can prevent bile flow and cause pain, jaundice, infection, or pancreatitis.

\medterm{Bile Duct Cancer} Cancer that begins in a bile duct. It can block bile flow and cause yellow skin, dark urine, pale stools, itching, abdominal pain, or weight loss; treatment depends on its location and spread.

\synonyms
Cholangiocarcinoma; Biliary Tract Cancer

\medterm{Bile Duct Carcinoma} Alternate name; see \textit{Bile Duct Cancer}.

\medterm{Bile Duct Obstruction} A blockage that prevents bile moving from the liver or gallbladder into the intestine. It may cause jaundice, dark urine, pale stools, itching, abdominal pain, infection, or pancreatitis.

\medterm{Bile Duct Stricture} Narrowing of a bile duct that slows or blocks bile flow. It may result from surgery, inflammation, stones, or cancer and can cause jaundice, pain, infection, or abnormal liver tests.

\medterm{Bile Ducts} Tubes that carry bile from the liver and gallbladder to the small intestine. A blockage or narrowing can cause jaundice, abdominal pain, infection, or inflammation of the pancreas.

\medterm{Bile Pigment} A coloured substance, mainly bilirubin, made when haemoglobin from old red blood cells is broken down. The liver processes it and excretes it in bile; build-up causes yellowing of the skin and eyes.

\medterm{Bile Reflux} Backward flow of bile and other intestinal contents into the stomach and sometimes the swallowing tube. It may cause upper-abdominal pain, nausea, bitter fluid, or inflammation and can occur after surgery or when normal stomach movement is impaired.

\synonyms
Duodenogastric Reflux; Biliary Gastritis

\medterm{Bile Salt Emulsification} The process in which bile salts break dietary fat into tiny droplets. This increases the surface area available to digestive enzymes and helps the intestine absorb fat and fat-soluble vitamins.

\medterm{Bile Salt Malabsorption} Failure to reabsorb enough bile salts in the lower small intestine. Excess bile salts enter the colon and draw in water, causing chronic watery diarrhoea, urgency, and abdominal discomfort.

\medterm{Bile Salts} Substances made from cholesterol in the liver and released in bile. They help break fats into small droplets so the intestine can digest and absorb fats and vitamins A, D, E, and K.

\medterm{Bilevel Positive Airway Pressure} Breathing support delivered through a mask that gives higher pressure when breathing in and lower pressure when breathing out. It can improve oxygen and carbon-dioxide levels in selected illnesses and requires monitoring.

\medterm{Bilharzia} An infection caused by Schistosoma parasites acquired when contaminated freshwater contacts the skin. It may affect the intestine, liver, urinary tract, or other organs and can cause long-term inflammation and scarring.

\medterm{Bilharzia} Alternate name; see \textit{Schistosomiasis}.

\medterm{Bili Lights} Special blue lights used to treat jaundice in newborns. The light changes bilirubin in the skin into forms the body can remove more easily, helping prevent bilirubin from damaging the brain.

\medterm{Biliary} Relating to bile, the gallbladder, or the tubes that carry bile from the liver to the intestine.

\medterm{Biliary Atresia} A serious disease in newborns in which the bile ducts outside the liver are blocked or absent. Bile builds up in the liver, causing jaundice and progressive liver damage unless treated promptly.

\synonyms
Congenital Biliary Atresia
Extrahepatic Biliary Atresia

\medterm{Biliary Calculi} Stones formed from hardened bile in the gallbladder or bile ducts. They may cause no symptoms or block bile flow, producing pain, jaundice, infection, or inflammation of the gallbladder or pancreas.

\medterm{Biliary Cirrhosis} Permanent liver scarring caused by long-term injury or blockage of the bile ducts. The term has also been used for primary biliary cholangitis, an immune disease affecting small bile ducts.

\medterm{Biliary Cirrhosis, Primary} Alternate name; see \textit{Primary Biliary Cholangitis}.

\medterm{Biliary Colic} Steady pain in the upper right abdomen or upper middle abdomen caused by a gallstone temporarily blocking the gallbladder outlet. It often follows a meal and lasts minutes to hours; fever or persistent pain suggests a complication.

\medterm{Biliary Decompression} A procedure that drains blocked bile ducts to restore bile flow. A doctor may place a stent or drainage tube through an endoscope or through the skin, depending on the cause and location of the blockage.

\medterm{Biliary Dyskinesia} Gallbladder pain without gallstones or another obvious structural problem, often associated with abnormal emptying. Diagnosis requires typical symptoms and exclusion of other causes; a scan may help but is not decisive alone.

\medterm{Biliary Fever} Fever caused by infection or inflammation in the gallbladder or bile ducts. Fever with upper-abdominal pain or jaundice may indicate a rapidly worsening infection and needs urgent medical care.

\medterm{Biliary Fistula} An abnormal passage allowing bile to leak from the gallbladder or bile ducts into another organ, the abdominal cavity, or the skin. It may cause pain, infection, bile drainage, or skin irritation.

\medterm{Biliary Sand} Fine crystals or thick sludge suspended in bile, usually containing cholesterol, calcium salts, or pigment. It may disappear, form gallstones, block a duct, or cause gallbladder inflammation.

\medterm{Biliary Scintigraphy} A specialized nuclear medicine scan that tracks the production and flow of bile from the liver into the gallbladder, bile ducts, and small intestine using an injected radioactive tracer (such as technetium-99m labeled iminodiacetic acid). It is the most accurate imaging test for diagnosing acute cholecystitis (indicated by failure of the gallbladder to fill due to cystic duct blockage). It is also used to detect post-surgical bile leaks, evaluate biliary dyskinesia by measuring gallbladder emptying after cholecystokinin administration, and differentiate biliary atresia from neonatal hepatitis in jaundiced infants.

\synonyms
Cholescintigraphy; HIDA Scan; Hepatobiliary Scintigraphy; Hepatobiliary Iminodiacetic Acid Scan

\medterm{Biliary Stent} A small tube placed inside a bile duct to keep it open and allow bile to drain when it is narrowed or blocked. Stents can move, clog, or cause infection or pancreatitis and require follow-up.

\medterm{Biliary System} The liver, gallbladder, and bile ducts that make, store, and carry bile into the small intestine to help digest fats.

\medterm{Biliary Tract} The network of ducts carrying bile from the liver and gallbladder to the small intestine. Blockage or inflammation can cause upper-abdominal pain, jaundice, infection, or pancreatitis.

\medterm{Biliary Tract Biopsy} The acquisition of tissue specimens or exfoliated cells from the extrahepatic or intrahepatic bile ducts or gallbladder for histopathologic examination, commonly performed via endoscopic retrograde cholangiopancreatography (ERCP), intraductal brush cytology, or image-guided percutaneous needle biopsy.

\synonyms
Bile Duct Biopsy; Gallbladder Biopsy; Biliary Cytology

\medterm{Biliary Tract Injuries} Damage to the gallbladder or bile ducts, often after abdominal surgery or trauma. Bile may leak or drainage may become blocked, causing abdominal pain, jaundice, infection, or later narrowing.

\medterm{Bilious} Containing or relating to bile. Bilious vomit is usually yellow or green because bile has entered the stomach contents.

\medterm{Bilious Vomiting} Vomiting green or yellow fluid containing bile. In a newborn or infant, clearly green vomit can indicate bowel blockage and requires immediate medical assessment.

\medterm{Bilirubin} A yellow substance made mainly when old red blood cells are broken down. The liver changes it and sends it into bile; too much bilirubin in the blood causes yellowing of the skin and eyes.

\medterm{Bilirubin Blood Test} A blood test measuring total and direct bilirubin. High results can occur when red blood cells break down too quickly, the liver cannot process bilirubin, or bile flow is blocked; newborn levels help assess risk of brain injury.

\medterm{Bilirubin Encephalopathy} Brain injury caused by very high levels of unconjugated bilirubin, especially in newborns. Early signs may include poor feeding, sleepiness, weak or stiff muscles, and seizures; permanent injury is called kernicterus.

\medterm{Bilirubin Metabolism} The process by which the body makes bilirubin from old red blood cells, carries it to the liver, changes it into a water-soluble form, and removes it in bile and stool.

\medterm{Bilirubin Test, Urine} A clinical urinalysis that detects conjugated (water-soluble) bilirubin in urine. A positive result indicates hepatobiliary disease, hepatocellular injury, or biliary tract obstruction; unconjugated bilirubin is bound to albumin and does not filter into urine.

\synonyms
Urine Bilirubin; Urinary Bilirubin Test

\medterm{Bilirubinuria} The presence of water-soluble bilirubin in urine. It may occur with liver disease or blocked bile flow and can make urine appear unusually dark.

\medterm{Biliverdin} A green pigment formed when haemoglobin is broken down. The body normally converts it into bilirubin, which the liver processes and removes in bile.

\medterm{Billroth I} A surgical procedure for partial gastrectomy in which the distal stomach (antrum and pylorus) is resected and the remaining proximal gastric remnant is anastomosed directly to the duodenum (gastroduodenostomy).

\synonyms
Billroth I Reconstruction; Gastroduodenostomy

\medterm{Billroth II} A surgical procedure for partial gastrectomy in which the distal stomach is resected, the duodenal stump is closed, and the remaining gastric remnant is anastomosed to a loop of the proximal jejunum (gastrojejunostomy), leaving a blind duodenal limb for biliopancreatic drainage.

\synonyms
Billroth II Reconstruction; Gastrojejunostomy

\medterm{Billroth's Operation} An older type of partial stomach-removal surgery. In Billroth I the remaining stomach is joined to the duodenum; in Billroth II it is joined to the jejunum, a later part of the small intestine.

\medterm{Bilophila} A group of bacteria that can live without oxygen in the human intestine. Most are harmless there, but Bilophila wadsworthia has been found in some appendicitis, abdominal, wound, and bloodstream infections.

\medterm{Bilophila Wadsworthia} A bacterium that tolerates bile and normally occurs in small numbers in the colon. If intestinal bacteria enter damaged or normally sterile tissue, it may contribute to appendicitis, abdominal infection, or wound infection.

\medterm{Bimanual} Performed using both hands. A bimanual pelvic examination uses one hand inside the vagina and the other on the abdomen to assess pelvic organs.

\medterm{Bimatoprost} Bimatoprost is a prostaglandin analogue used as eye drops to lower intraocular pressure and, in some formulations, promote eyelash growth.

\medterm{Bimedial} Relating to or located on the inner side of two structures. The exact structures involved must be specified because the term alone does not identify a disease or diagnosis.

\medterm{Bimellar} Made of two thin layers or plates. The term describes the structure of a membrane, tissue, or other object and has no disease meaning by itself.

\medterm{Binary Fission} A form of reproduction in which one cell divides into two genetically similar cells. It is common among bacteria and many single-celled organisms.

\medterm{Binder Excipient} An inactive ingredient added during tablet manufacture to help powder particles stick together. It has no intended treatment effect but affects the tablet's strength and breakdown.

\medterm{Binding Protein} A protein that attaches specifically to another molecule, such as a hormone, medicine, nutrient, DNA segment, or cell-surface signal. Binding can transport, store, activate, or block the other molecule.

\medterm{Binding Site} The particular area of a molecule where another molecule attaches. The interaction can change activity, transmit a signal, allow transport, or enable a laboratory test to detect the target.

\medterm{Binet-Simon Scale} An early intelligence test for children that compared performance with tasks expected at different ages. It influenced modern cognitive testing but is no longer the standard tool for clinical assessment.

\medterm{Binge Drinking} A pattern of consuming enough alcohol over a short period to produce a high blood alcohol concentration and increased risk of injury or acute toxicity.

\synonyms
Episodic Heavy Drinking; Heavy Episodic Drinking

\medterm{Binge Eating Disorder} Alternate spelling; see \textit{Binge-Eating Disorder}.

\medterm{Binge-Eating Disorder} A mental-health condition involving repeated episodes of eating large amounts of food with loss of control and distress, without regular behaviors intended to compensate. Weight alone does not establish the diagnosis.

\synonyms
BED; Compulsive Overeating Disorder

\medterm{Binocular Vision} The use of both eyes together to produce one image and judge depth. Misalignment or poor coordination between the eyes can cause double vision, eyestrain, or reduced depth perception.
\synonyms
Two-Eyed Vision; Binocular Visual Function

\medterm{Binocular Visual Loss} Reduced vision affecting both eyes. Causes may involve the eyes, optic nerves, brain, or visual pathways; sudden loss, weakness, speech trouble, or severe headache requires emergency assessment.

\medterm{Binswanger Disease} A form of small-vessel disease in the brain that damages deep white matter. It can cause gradually worsening thinking, walking difficulty, mood changes, and urinary problems.

\medterm{Binswanger's Dementia} A form of vascular cognitive impairment caused by diffuse small-vessel disease affecting the brain's deep white matter.

\synonyms
Subcortical Ischemic Vascular Disease; Subcortical Arteriosclerotic Encephalopathy

\medterm{Bio-Informatics} The use of computers and statistics to store, compare, and interpret biological information such as DNA sequences, genes, proteins, and medical data.

\medterm{Bioabsorbable Vascular Stent} A temporary endovascular scaffold constructed from biodegradable polymers (such as poly-L-lactic acid) or bioresorbable metallic alloys (such as magnesium) engineered to provide vessel patency during vascular remodeling before undergoing hydrolytic degradation and systemic resorption.

\synonyms
Bioabsorbable Vascular Scaffold; BVS; Bioresorbable Stent; BRS

\medterm{Bioaccumulation} The gradual build-up of a substance in an organism when absorption is faster than removal. Pollutants, medicines, and metals can accumulate in tissues and may cause harm when levels become high.

\medterm{Bioactive} Having an effect on living cells, tissues, or organisms, such as a bioactive drug, hormone, nutrient, or signaling molecule.

\medterm{Bioassay} A test that measures the effect or activity of a substance by observing its action on living cells, tissues, organisms, or a biological system. It can be used to estimate potency, toxicity, or biologic activity.

\medterm{Bioavailability} The amount of a medicine that reaches the bloodstream in an active form and how quickly it gets there. It depends on how the medicine is given, absorbed, and broken down; intravenous medicines have complete availability.

\medterm{Bioburden} The number and types of living microorganisms present on a surface, medical device, tissue, or medicine before sterilization. Testing it helps confirm that cleaning and sterilization methods are adequate.

\medterm{Biochanin A} A naturally occurring isoflavone found mainly in legumes and other plants. It is studied for biologic effects but is not established as a treatment for disease.

\synonyms
Biochanin a Isoflavone; 5,7-Dihydroxy-4'-Methoxyisoflavone

\medterm{Biochemical} Relating to the chemical substances and reactions that occur in living organisms, including metabolism and cellular signaling.

\medterm{Biochemical Assessment} Evaluation using laboratory measurements such as blood proteins, salts, glucose, fats, vitamins, and minerals. Results can help assess nutrition, metabolism, organ function, or disease but must be interpreted with the clinical picture.

\medterm{Biochemical Reaction} A chemical change occurring in a living cell or body, usually controlled by enzymes. It may build or break down molecules and supports energy production, growth, repair, and normal organ function.

\medterm{Biochemical Recurrence} A rise in a laboratory marker after treatment that suggests disease remains or has returned, even though symptoms or scans may not yet show it. The meaning depends on the disease and the test used.

\synonyms
Biochemical Relapse; Biochemical Failure

\medterm{Biochemistry} The study of the chemical substances and reactions occurring in living organisms. Clinical biochemistry uses laboratory measurements to assess health and disease.

\medterm{Biocompatibility} The ability of a material or device to perform its intended function without causing an unacceptable harmful reaction in body tissues.

\medterm{Biodegradable Implant} An implant designed to provide support temporarily and then gradually break down into substances the body can process or remove. Its safety depends on how quickly it breaks down and how tissues respond.

\medterm{Biodistribution} The way a medicine, tracer, or other substance spreads through the blood and different organs over time. It affects where treatment acts, how well an image is produced, and radiation or toxicity to healthy tissue.

\medterm{Bioengineering} The application of engineering principles to biology and medicine, including the design of medical devices, biomaterials, tissue-engineering systems, prostheses, and diagnostic technologies.

\medterm{Bioequivalence} A finding that two versions of a medicine enter the bloodstream at a similar rate and amount when given at the same dose. It supports the expectation of comparable effects and safety for approved uses.

\medterm{Bioethics} The study of moral questions in medicine, biology, research, and healthcare technology. It considers consent, privacy, fairness, benefit, harm, quality of life, and respect for patients and communities.

\medterm{Bioethics, Geriatric} The systematic exploration of moral and ethical considerations in the medical care of elderly patients, encompassing decisional capacity, informed consent, advanced healthcare directives, surrogate decision-making, palliative care transitions, and protection of vulnerable older adults.

\synonyms
Geriatric Bioethics; Ethics in Geriatric Medicine; Ethics of Aging

\medterm{Biofeedback} A technique that shows information about a body function, such as muscle tension, heart rate, or breathing, so a person can learn to control it. It may support treatment of selected pain, stress, or movement problems.

\medterm{Biofilm} A community of microorganisms attached to a surface and surrounded by a protective material they produce. Biofilms can form on teeth, wounds, catheters, and implants and may be difficult to remove.

\medterm{Biofilm-Associated Infections} Infections caused by microorganisms living in a protective layer on tissue, wounds, teeth, or medical devices. They can persist despite antibiotics and may require removal or replacement of an infected device.

\medterm{Biofluid} A fluid from a living body, such as blood, urine, saliva, spinal fluid, or joint fluid. Biofluids can be tested for cells, chemicals, microorganisms, medicines, or signs of disease.

\medterm{Biologic Therapy} Treatment using a medicine made from living cells or biological substances to target a specific immune or disease pathway. It is used for selected inflammatory, autoimmune, cancer, and other conditions.

\medterm{Biologic Therapy In Dermatology} Treatment with targeted biological agents, usually monoclonal antibodies or receptor
proteins, that inhibit selected immune pathways. Biologics are used for diseases such as
psoriasis, atopic dermatitis, hidradenitis suppurativa, and urticaria; infection risk,
vaccination status, screening, and monitoring must be considered.

\medterm{Biologic Variation} Normal change in a body measurement within the same person over time or between different people. It must be considered when deciding whether a laboratory change represents real disease or treatment-related change.

\medterm{Biological} Relating to living organisms, their parts, functions, processes, or products. In medicine, a biological product is made from or resembles a substance from a living source.

\medterm{Biological Availability} The amount and speed at which a medicine or other substance becomes available to produce an effect in the body. It depends on the route, absorption, breakdown, and removal of the substance.

\medterm{Biological Clock} The body's internal timing system that helps coordinate daily sleep, alertness, hormone release, body temperature, and other repeating functions in response to light and other environmental signals.

\medterm{Biological Half-Life} The time needed for the body to remove half of a medicine or other substance through metabolism and excretion. It differs from the physical half-life of a radioactive substance and helps guide dosing or safety decisions.

\medterm{Biological Model} A cell, tissue, organism, computational system, or other representation used to study biological mechanisms, disease, or treatment responses.

\medterm{Biological Process} A coordinated series of molecular, cellular, or physiological events in a living organism, such as metabolism, cell division, inflammation, or tissue repair.

\medterm{Biological Relative} A person who shares genetic ancestry with someone. Family medical history uses biological relatives to assess inherited risks, while relatives by marriage or legal relationship do not share that genetic information.

\medterm{Biological Therapy} Treatment using biological substances, cells, vaccines, antibodies, or similar products to change immune activity or target a specific disease pathway. It is used in selected cancers, immune disorders, and other conditions.

\medterm{Biological Transport} The movement of substances such as water, salts, nutrients, medicines, or cells within the body or across cell membranes. It may occur by diffusion, channels, pumps, filtration, or transport vesicles.

\medterm{Biological Variability} Natural differences in a body measurement between people or changes in the same person over time, separate from laboratory error. It can affect how test results are interpreted.

\synonyms
Biologic Variation; Physiologic Variability

\medterm{Biological Warfare} The deliberate use of microorganisms, toxins, or other biological agents to cause illness, death, or disruption in people, animals, or plants.

\medterm{Bioluminescence} Light produced by a living organism through a chemical reaction. It is used by some organisms for communication or defense and is used in laboratories to detect biological activity.

\medterm{Biomarker} A measurable feature that gives information about a normal body process, disease, likely outcome, or response to treatment. It may be measured in blood, tissue, genes, images, or other samples.

\medterm{Biomarker Testing} Laboratory testing of genes, proteins, metabolites, or other measurable features to help identify disease, estimate likely outcome, choose treatment, or monitor response.

\medterm{Biomaterial} A natural or manufactured material designed to contact or interact with the body. Examples include implant surfaces, artificial joints, wound dressings, prostheses, and materials that release medicines.

\medterm{Biomechanical Compliance} How easily a body structure changes size or shape when pressure or force is applied. Lung and blood-vessel compliance affect breathing and circulation, while joint compliance affects movement and stability.

\medterm{Biomechanics} The study of how forces and movement affect the body. It examines walking, posture, joints, muscles, and loads to understand injury, improve rehabilitation, or design supportive devices.

\medterm{Biomedical Cement} A medical material used to fill gaps, stabilize broken bone, anchor an implant, or repair tissue. Its use and strength depend on the procedure and the tissue being treated.

\medterm{Biomedical Device} An instrument, implant, software system, or other product intended to diagnose, monitor, prevent, or treat disease or injury. Its safety and usefulness depend on proper design, testing, maintenance, and correct use.

\medterm{Biomedical Engineering} The use of engineering, materials, computing, and design to create or improve medical devices, diagnostic tests, implants, prostheses, imaging systems, and other healthcare technologies.

\medterm{Biomedical Tape} Adhesive tape designed for medical use, such as securing dressings, tubes, sensors, or small devices. Different types vary in strength, flexibility, water resistance, and risk of skin irritation.

\medterm{Biomedicine} Medicine grounded in the biological sciences, including the study of disease mechanisms, diagnosis, prevention, and treatment through laboratory and clinical research.

\medterm{Biometry} Measurement of body structures or functions. In eye care, optical biometry measures the eye to help calculate the strength of an artificial lens before cataract surgery.

\medterm{Biomicroscopy} Examination of living tissue with a microscope. In eye care, a slit lamp provides magnified views of the cornea, front chamber, lens, and other eye structures.

\medterm{Biomolecule} A molecule made by or found in living organisms, such as a protein, fat, sugar, DNA, or RNA. Biomolecules provide structure, energy, information, and control of cell functions.

\medterm{Biomonitoring} Repeatedly measuring substances or body changes over time using blood, urine, breath, hair, or other samples. It can assess exposure to poisons or medicines, treatment levels, hormones, or disease activity.

\medterm{Biomphalaria} A group of freshwater snails that can carry the parasite Schistosoma mansoni. The parasite develops in the snails and can later infect people who contact contaminated freshwater.

\medterm{Bionics} The use of electronic or mechanical devices to replace or support body functions, such as hearing, movement, or sensation. Examples include cochlear implants, powered prostheses, and other assistive systems.

\medterm{Biophotonics} The use of light to study, diagnose, monitor, or treat living tissues and biological processes. Examples include optical imaging, fluorescence testing, laser surgery, and photodynamic therapy.

\medterm{Biophysical Profile} A pregnancy assessment combining fetal-heart-rate monitoring with ultrasound. It checks fetal movement, breathing movements, muscle tone, and amniotic-fluid volume to help assess fetal well-being.

\medterm{Biopsy} Removal of a small sample of cells or tissue so it can be examined under a microscope to diagnose disease, including cancer. Types include fine-needle aspiration (cells or fluid), core-needle (a tissue cylinder), incisional (part of an area), excisional (the whole area), punch (a circular piece of skin), and endoscopic biopsy (through a viewing tube). Possible complications include pain, bleeding, and infection; risks vary with the body site and procedure.

\synonyms
Biopsies

\medterm{Biopsy Lung} Removal of a small lung-tissue sample using a needle, bronchoscope, or operation. Examination can investigate infection, inflammation, scarring, or cancer; the method depends on the lesion and procedural risks.

\medterm{Biopsy Needle} A needle used to remove cells or tissue for laboratory examination. A fine needle usually collects loose cells, while a core needle removes a small cylinder that preserves the tissue's arrangement.

\medterm{Biopsy Site} The body location from which cells or tissue are removed for testing. It is chosen to sample the abnormal area while limiting injury to nearby organs, nerves, and blood vessels.

\medterm{Biopsy Specimen} Cells or tissue removed from the body for laboratory examination. Microscopy, special stains, immune tests, or genetic tests may be used to identify cancer, infection, inflammation, or another condition.

\medterm{Biosafety} Practices, equipment, and containment methods used to prevent accidental exposure to or release of infectious organisms, biological materials, or toxins.

\medterm{Biosafety Level} A classification describing the laboratory facilities, equipment, work practices, and protective measures needed to handle biological agents safely according to their risk.

\medterm{Biosecurity} Measures that prevent biological agents or toxins from being accidentally or deliberately released, stolen, or misused. It includes secure facilities, controlled access, staff training, risk assessment, and responsible handling of biological materials.

\medterm{Biosensor} A device that uses a biological material to recognize a substance and a detector to measure it. Medical examples include sensors for glucose, hormones, infectious organisms, or heart-related markers.

\medterm{Biosimilar} A biological medicine highly similar to an already approved biological medicine, with no meaningful differences in safety or effectiveness for approved uses. Because these medicines are complex, similarity is shown through comparative testing.

\medterm{Biospecimen} A sample of human tissue, cells, blood, urine, or another body fluid collected for testing or research. Correct collection, labelling, storage, and handling help keep the sample reliable.

\medterm{Biostatistics} The use of statistics to plan medical studies, analyse health data, measure uncertainty, and interpret evidence about disease, diagnosis, prevention, and treatment.

\medterm{Biosynthesis} The process by which living cells use enzymes and energy to build complex molecules from simpler substances, such as making proteins from amino acids.

\medterm{Biot Breathing} An irregular breathing pattern with groups of normal-sized breaths interrupted by unpredictable pauses in breathing. It may occur with serious damage affecting the lower brainstem and requires urgent assessment.

\medterm{Biot's Breathing} An abnormal pattern of irregular breaths interrupted by unpredictable pauses in breathing. It suggests disruption of the brain's breathing control and may occur with severe neurological illness.

\medterm{Biotechnology} The use of living cells, organisms, or biological molecules to develop products and processes such as vaccines, biological medicines, genetic tests, and gene therapies.

\medterm{Biotelemetry} Measuring and sending body information from a person or animal to a distant receiver. It can be used for heart-rhythm, temperature, breathing, glucose, or intensive-care monitoring.

\medterm{Bioterrorism} The deliberate release of infectious organisms, toxins, or other biological agents to cause illness, death, fear, or disruption. Response involves rapid detection, public-health coordination, protection, testing, isolation when needed, and treatment or prevention.

\medterm{Bioterrorism Agents} Infectious organisms or toxins that could be deliberately released to cause widespread illness or panic. Public-health authorities prepare for agents causing anthrax, smallpox, botulism, plague, tularemia, or severe viral haemorrhagic fever.

\medterm{Biotherapy} Treatment using living organisms, cells, genes, or biological substances to prevent or treat disease. Immunotherapy, some cell therapies, and selected biological medicines are forms of biotherapy.

\medterm{Biotin} A water-soluble B vitamin needed for enzymes involved in processing fats, sugars, and amino acids. Deficiency is uncommon but can cause a skin rash, hair loss, weakness, or neurological symptoms.

\medterm{Biotin Vitamin B7} Alternate name; see \textit{Biotin}.
\medterm{Biotinidase Deficiency} A rare inherited disorder in which the body cannot recycle biotin efficiently. Without treatment it can cause seizures, developmental problems, hearing or vision loss, skin rash, and weakness; lifelong biotin usually prevents symptoms.

\medterm{Biotinylation} The attachment of biotin to another molecule, usually a protein, so it can be found or separated in a laboratory test. The biotin acts as a tag that binds to special substances, making the target easier to detect or collect. Biotinylation can occur naturally or be done artificially.

\medterm{Biotransformation} The natural process by which the body changes drugs and other foreign substances, mainly through enzymes in the liver but also in other tissues. It often makes them less active and easier to remove. Phase I reactions modify the substance; phase II reactions may attach another molecule to help eliminate it. These steps are not always sequential, and the process can produce an active or harmful substance instead.

\medterm{Biotrauma} Inflammation caused by mechanical injury to the lungs during breathing-machine treatment. Inflammatory signals can spread through the body and contribute to organ injury; careful, lung-protective settings reduce the risk.

\medterm{BiPAP} Bilevel Positive Airway Pressure (BiPAP), a form of noninvasive breathing support delivered through a mask. It provides higher pressure during inhalation and lower pressure during exhalation. It may be used for conditions such as hypercapnic COPD exacerbation, sleep-related breathing disorders, and cardiogenic pulmonary edema.

\synonyms
Bilevel Positive Airway Pressure

\medterm{BiPAP Machine} A device that delivers two levels of breathing pressure through a mask—higher during inhalation and lower during exhalation. Settings are prescribed and adjusted according to the person's breathing problem and response.

\medterm{Biparental} Involving both biological parents. In genetics, biparental inheritance means that a child receives relevant genetic material from both the mother and the father.

\medterm{Biparietal Diameter} The width of a baby's head measured by ultrasound from one side to the other. It helps estimate pregnancy age and monitor growth, but is interpreted with other measurements because head shape can affect the result.

\medterm{Bipartite Ossification} Formation or persistence of two separate ossification centers in a bone. If they do not fuse, a bipartite bone, such as a bipartite patella, may result. This is usually an incidental developmental variant but can become painful after injury or repetitive stress.

\medterm{Biphenyl Compounds} Chemicals made of two linked benzene rings. Polychlorinated biphenyls (PCBs) are long-lasting industrial pollutants; exposure can harm the liver, immune system, and hormone system and may increase cancer risk.

\medterm{Bipolar} Alternate name; see \textit{Bipolar Disorder}.

\medterm{Bipolar Affective Disorder} Alternate name; see \textit{Bipolar Disorder}.

\medterm{Bipolar Disorder} A mental health condition involving episodes of mania, hypomania, and/or major depression. Bipolar I is defined by at least one manic episode; depressive episodes are common but not required. Bipolar II includes at least one hypomanic episode and at least one major depressive episode, without a history of mania. Episodes can affect mood, sleep, energy, judgment, and daily functioning.

\synonyms
Manic-Depressive Illness

\medterm{Bipolar I Disorder} A mood disorder defined by at least one manic episode lasting at least seven days or
requiring hospitalization. A manic episode is a distinct period of abnormally elevated,
expansive, or irritable mood and increased goal-directed activity or energy.

\medterm{Bipolar II Disorder} A mood disorder characterized by at least one hypomanic episode and at least one major depressive episode, with no history of a full manic episode. The combination of hypomania and depression can cause substantial distress or impairment even though hypomania is less severe than mania.

\medterm{Bipolar Neuron} A bipolar neuron has one dendrite and one axon extending from opposite sides of its cell body, as found in the retina and sensory pathways.

\medterm{Bipolaris} A genus of fungi found mainly in soil and plants. Some species can rarely cause eye, sinus, skin, or invasive infections, particularly in people with reduced immunity. Identification requires laboratory culture or molecular testing because treatment depends on the species and site of infection.

\medterm{Bird Fancier's Lung} A form of hypersensitivity pneumonitis caused by breathing proteins from bird feathers or droppings. It can cause cough and shortness of breath; repeated exposure may lead to permanent lung scarring.

\medterm{Bird Flu} An infectious zoonotic illness caused by avian influenza type A viruses (notably H5N1, H5N6, H7N9) that naturally circulate among wild waterfowl and domestic poultry. Transmission to humans occurs via direct contact with infected birds or contaminated environments, producing acute respiratory distress syndrome (ARDS), multi-organ dysfunction, and high case-fatality rates.

\synonyms
Avian Influenza; Avian Flu; Bird Flu; H5N1 Infection

\medterm{Birt-Hogg-Dubé Syndrome} An inherited condition caused by a change in the FLCN gene. It can cause small, usually harmless skin growths, cysts in the lungs with repeated collapsed lung, and an increased risk of kidney tumors. It is autosomal dominant, meaning one altered copy of the gene can cause the condition.

\medterm{Birth} The process by which a baby leaves the uterus, including labour, delivery of the baby, and delivery of the placenta.

\medterm{Birth Asphyxia} Too little oxygen or inadequate breathing before, during, or shortly after birth. It can injure the brain and other organs and requires immediate newborn assessment, breathing support, and treatment.

\medterm{Birth Control} Methods used to prevent pregnancy, including condoms, pills, implants, injections, intrauterine devices, fertility-awareness methods, and permanent procedures. Choice depends on effectiveness, health, side effects, future plans, access, and personal preference.

\medterm{Birth Control And Family Planning} Planning whether and when to have children using contraception, fertility care, pregnancy advice, and reproductive health services. Care respects personal choices and considers health conditions, safety, goals, and informed consent.

\medterm{Birth Control Methods} Birth-control methods include condoms, pills, implants, injections, intrauterine devices, patches, rings, fertility-awareness approaches, and permanent contraception. Effectiveness, sexually transmitted infection protection, side effects, reversibility, and medical eligibility vary by method.

\medterm{Birth Control Pill Overdose} Illness after taking too many contraceptive tablets. Most exposures cause little or no serious harm, but nausea, vomiting, drowsiness, or other symptoms may occur, and the product and amount taken help guide management.

\medterm{Birth Control Pills} Tablets taken on a schedule to prevent pregnancy, usually by stopping ovulation, thickening cervical mucus, or changing the uterine lining. They do not protect against sexually transmitted infections and work best when taken correctly.

\medterm{Birth Control, Extended-Release} Alternate name; see \textit{Long-Acting Reversible Contraception}.

\medterm{Birth Defect} A structural, functional, or metabolic abnormality present at birth resulting from chromosomal aberrations, monogenic mutations, prenatal teratogenic exposures (e.g., alcohol, retinoic acid, infections), or complex multifactorial interactions during embryogenesis and organogenesis.

\synonyms
Congenital Anomaly; Congenital Malformation; Birth Defects; Congenital Disorder

\medterm{Birth Defects} Structural or functional problems present at birth that may affect development or health. Causes include genetic or chromosomal changes, infections, medicines, nutritional deficiencies, and environmental exposures; sometimes no cause is found.

\medterm{Birth Injury} Physical harm to a newborn during labour or delivery. It may affect the skin, bones, nerves, muscles, or brain and can result from difficult delivery, the baby's position, or medical complications.

\medterm{Birth Order} The sequence in which siblings are born within a family. While birth order has been studied for psychological effects, its direct medical significance is limited to obstetric considerations such as risks associated with grand multiparity.

\medterm{Birth Pains} Birth pains are painful uterine contractions and related pressure during labor as the cervix dilates and the baby moves through the birth canal.

\medterm{Birth Rate} The number of live births in a population during a specified period, usually expressed per 1,000 people per year.

\medterm{Birth Spacing} The practice of planning the time between giving birth and becoming pregnant again. Health guidelines generally recommend waiting at least 18 to 24 months between pregnancies to allow the mother's body to recover and restore essential nutrients, reducing the risks of premature birth, low birth weight, and delivery complications.
\synonyms Interpregnancy interval; Child spacing; Pregnancy spacing.

\medterm{Birth Weight} A newborn's weight measured soon after birth. Weight below 2,500 grams is considered low and is linked with increased health risks; weight below 1,500 grams is very low and usually requires specialized newborn care.

\medterm{Birthmark} A localized mark on the skin or in blood vessels that is present at birth or appears soon afterward. Examples include pigmented moles, haemangiomas, and port-wine stains.

\medterm{Birthmarks} Marks on the skin present at birth or appearing early in life. Most are harmless, but rapid growth, colour or shape change, bleeding, ulceration, pain, or a mark affecting vision or breathing needs medical assessment.

\medterm{Birthmarks And Other Skin Pigmentation Problems} Birthmarks or changes in skin colour may result from blood vessels, pigment cells, inflammation, medicines, hormones, or inherited conditions. A rapidly changing, bleeding, ulcerated, painful, or distinctly different mark requires examination.

\medterm{Bis In Die} Latin for “twice daily,” traditionally abbreviated b.i.d. Clear prescriptions state when doses are to be taken rather than relying only on this abbreviation.

\medterm{Bisacodyl} A stimulant laxative that increases bowel movement and helps empty the bowel. It is used short term for constipation or before some procedures; cramps, diarrhoea, and dehydration can occur.

\medterm{Bisexual} Bisexual describes sexual or romantic attraction to more than one gender; it is an identity or orientation, not a disease or disorder.

\medterm{Bisexuality} Bisexuality is a sexual orientation involving attraction to more than one gender; it is not a disease or disorder.

\medterm{Bishop Score} A score used before inducing labour to describe how ready the cervix is for birth. It considers opening, thinning, softness, position, and the baby's head position; a higher score generally predicts an easier induction.

\medterm{Bismuth} A metal used in some medicines and industrial products. Certain bismuth compounds can soothe diarrhoea or stomach symptoms, but excessive or prolonged exposure can cause nervous-system or kidney toxicity.

\medterm{Bisoprolol} A beta-1-selective blocker used to treat hypertension, angina, and selected patients with heart failure. It lowers heart rate and blood pressure and may cause bradycardia, hypotension, or fatigue.

\medterm{Bispecific Antibody Therapy} Treatment using an engineered antibody that binds two different targets, such as a cancer-cell marker and an immune-cell marker, to bring cells together or block more than one disease pathway. It is used for selected cancers and can cause immune reactions or other treatment-specific effects.

\medterm{Bisphenol A} An industrial chemical used in some plastics and resins. It can act like a weak hormone, and possible health effects of exposure are studied; reducing unnecessary contact is especially reasonable for infants and children.

\medterm{Bisphenol A} An endocrine-disrupting industrial chemical used predominantly in the synthesis of polycarbonate plastics and epoxy resins lining food and beverage containers. It exhibits weak estrogenic and antiandrogenic activity and has been implicated in metabolic and reproductive disturbances.

\synonyms
BPA; 4,4'-(Propane-2,2-diyl)diphenol

\medterm{Bisphosphonate} A medicine that slows the removal of old bone. It is used mainly for osteoporosis, Paget disease, and some cancer-related bone problems. Rare but important risks include jaw-bone damage and unusual thigh-bone fractures.

\medterm{Bisphosphonates} Medicines that slow the breakdown of bone. They are used mainly to prevent fractures in osteoporosis and to treat Paget disease and some bone problems caused by cancer.

\synonyms
Bisphosphonate

\medterm{Bite, Chigger} Alternate index label; see \textit{Chiggers}.

\medterm{Bite, Flea} Alternate index label; see \textit{Flea Bites in Humans}.

\medterm{Bite, Snake} Alternate index label; see \textit{Snake Bite}.

\medterm{Bitemporal Hemianopia} Loss of side vision on the outer side of both eyes. It usually occurs when a growth or other problem presses on the place where the two optic nerves cross, often near the pituitary gland.

\medterm{Bitewing} A dental X-ray showing the crowns of upper and lower teeth and the bone around them in one image. It is especially useful for finding decay between teeth and early loss of supporting bone.

\medterm{Bitewing Radiograph} An X-ray taken inside the mouth that shows the crowns of back teeth and the nearby gum-supporting bone. It helps detect decay between teeth and bone loss from gum disease.

\medterm{Bitot's Spots} Distinctive, triangular, foamy, silvery-white or pearly-gray patches on the white part of the eye (bulbar conjunctiva), most commonly found on the outer (temporal) side near the cornea. They are a classic sign of vitamin A deficiency (Stage X1B of xerophthalmia) caused by the buildup of dried keratin and gas-forming bacteria on the dried eye surface. Bitot's spots commonly occur alongside night blindness and ocular dryness. Prompt treatment with high-dose vitamin A supplements usually resolves the spots within weeks and prevents progression to permanent corneal damage and blindness.
\synonyms{Bitot Spots; Conjunctival Foamy Patches; Keratinized Conjunctival Plaques}

\medterm{Bituminosis} A work-related breathing disorder caused by repeated exposure to bitumen or asphalt fumes and dust. It may irritate the airways and contribute to chronic cough, breathlessness, or other lung disease.

\medterm{Biuret Test} A laboratory test that detects protein by turning violet when protein reacts with a copper-containing solution. It is mainly used in teaching, research, and some chemical analyses rather than routine patient diagnosis.

\medterm{Bivalve} Bivalve means having two shells or valves; bivalve mollusks may be relevant to food allergy, poisoning, or parasitic exposure.

\medterm{Bizarre Parosteal Osteochondromatous Proliferation} A rare benign but locally recurrent bone-surface lesion, also called Nora Lesion, composed of cartilage, bone, and spindle cells; it most often affects the hands or feet.

\medterm{BK Virus} A common virus that usually remains inactive after childhood infection. It can reactivate when immunity is suppressed, especially after organ transplantation, and damage the transplanted kidney or urinary tract.

\medterm{BK Virus Nephropathy} Kidney inflammation and damage caused by reactivation of BK virus, usually in a kidney-transplant recipient. It can reduce transplant function and is managed mainly by carefully reducing immune-suppressing treatment and monitoring viral levels.

\medterm{Black Death} The major plague pandemic of the fourteenth century, caused mainly by the bacterium Yersinia pestis. It caused widespread severe illness and death and is distinct from the individual forms of plague seen today.

\medterm{Black Eye} Bruising and swelling around the eye, usually caused by blunt trauma. Vision may remain normal, but eye pain, double
vision, reduced vision, an irregular pupil, severe headache, vomiting, or loss of consciousness can indicate an eye
or orbital injury and requires urgent assessment.

\synonyms
Periorbital Hematoma; Shiner

\medterm{Black Fly} A small biting fly that can cause painful or itchy skin reactions. In some tropical regions, black flies spread the parasite that causes river blindness.

\medterm{Black Hair, Hairy Tongue} A harmless condition in which the tiny surface projections of the tongue become long and dark, trapping debris or pigment. Smoking, antibiotics, dry mouth, and poor oral care can contribute.

\medterm{Black Hairy Tongue} A usually harmless dark, furry appearance of the tongue caused by overgrown surface projections trapping debris or pigment. It may improve with tongue cleaning, good oral hygiene, and addressing smoking or dry mouth.

\medterm{Black Lung} Lung disease caused by long-term inhalation of coal dust, usually in miners. It can cause cough and breathlessness and may progress from small areas of dust-related damage to extensive permanent scarring.

\medterm{Black Lung Disease} A dust-related lung disease caused by prolonged coal-dust exposure. It may cause no symptoms at first but can lead to cough, breathlessness, reduced lung function, and severe permanent scarring.

\medterm{Black Nightshade Poisoning} Poisoning after eating black nightshade or its berries, which may contain toxic plant chemicals. It can cause nausea, vomiting, abdominal pain, confusion, enlarged pupils, seizures, or abnormal heart rhythms.

\medterm{Black Or Tarry Stools} Black or tarry stools may indicate upper gastrointestinal bleeding, but can also result from iron or bismuth. New or persistent findings require medical assessment.

\medterm{Black Stain} Dark discoloration on the tooth surface, often near the gumline, caused by external deposits or chromogenic bacteria. Dental assessment
determines the cause and appropriate removal.

\medterm{Black Triangle} A triangular gap between neighboring teeth near the gumline, usually caused by loss or shrinkage of the gum papilla after gum disease, injury, or dental treatment. It may trap food and can be treated cosmetically or orthodontically.

\medterm{Black Widow Spider} The black widow spider is a venomous spider whose bite can cause severe muscle pain, abdominal cramping, sweating, hypertension, and autonomic symptoms.

\medterm{Blackflesh} A nontechnical description of tissue that looks dark or black. Causes include bruising, severe loss of blood supply, dead tissue, infection, or staining. New, painful, cold, or spreading discoloration needs urgent medical assessment.

\medterm{Blackhead} An open blocked hair follicle filled with oil and dead skin cells. The material looks black because it reacts with air, not because the skin is dirty.

\medterm{Blackheads} Open comedones formed when a hair follicle becomes blocked with sebum and keratin and
the surface material oxidizes. They are a common non-inflammatory lesion of acne,
especially on the face, chest, and back.

\medterm{Blackout} A period for which a person cannot remember events, often after heavy alcohol use. Memory loss can also occur with drugs, seizures, head injury, fainting, or other causes of altered consciousness.

\synonyms
Amnestic Episode; Alcohol-Related Amnesia When Caused by Alcohol

\medterm{Blackwater Fever} A severe complication of malaria in which red blood cells break down inside blood vessels. It can cause dark urine, severe anaemia, jaundice, kidney failure, and shock and requires emergency treatment.

\medterm{Bladder} A muscular hollow organ that stores urine from the kidneys and contracts to release it through the urethra during urination.

\medterm{Bladder Augmentation} Surgery that enlarges the bladder, usually by adding a section of bowel or another tissue. It can increase urine storage when the bladder is too small or has high pressure.

\medterm{Bladder Biopsy} Removal of a small piece of bladder lining, usually during cystoscopy, for laboratory examination. It can investigate a visible growth, unexplained blood in the urine, inflammation, or suspected bladder cancer.

\medterm{Bladder Calculi} Alternate name; see \textit{Bladder Stones}.

\medterm{Bladder Cancer} A malignant tumor arising in the urinary bladder, most often from the urothelial lining. Blood in the urine is a common presenting
feature, although symptoms vary.

\synonyms
Urinary Bladder Neoplasm; Urothelial Carcinoma

\medterm{Bladder Carcinoma} Cancer that begins in the lining of the urinary bladder, most often in cells called urothelial cells. Blood in the urine is common; smoking and some workplace chemical exposures increase risk.

\medterm{Bladder Carcinoma} Alternate name; see \textit{Bladder Cancer}.

\medterm{Bladder Control Problems In Children} Difficulty controlling urine in a child who has passed the usual stage of toilet training. It may involve daytime leakage, bedwetting, urgency, frequent urination, or difficulty emptying and may be related to development, constipation, infection, nerve problems, or urinary-tract differences.

\medterm{Bladder Control, Loss of} Alternate index label; see \textit{Urinary Incontinence}.

\medterm{Bladder Diary} A record of drinks, urination times, urine amounts, urgency, leakage, and nighttime voiding over several days. It helps describe bladder habits and patterns of urinary symptoms.

\medterm{Bladder Diseases} Disorders of the bladder, including infection, stones, tumors, neurogenic dysfunction, inflammation, obstruction, and impaired emptying.
\medterm{Bladder Disorders} A broad spectrum of pathological conditions affecting the urinary bladder, encompassing infectious cystitis, interstitial cystitis (bladder pain syndrome), urothelial carcinoma, neurogenic bladder dysfunction, calculi, and bladder outlet obstruction.

\synonyms
Bladder Diseases; Urinary Bladder Diseases; Diseases of the Bladder

\medterm{Bladder Diverticula} Pockets that push outward from the bladder wall. They may be present from birth or develop when urine repeatedly faces resistance leaving the bladder, and can trap urine, leading to infection or stones.

\medterm{Bladder Diverticulum} A pouch extending from the bladder wall. It may cause no symptoms or trap urine and lead to infection, stones, or incomplete emptying; treatment addresses the cause and symptoms.

\medterm{Bladder Exstrophy} Bladder Exstrophy is a rare birth defect in which the bladder is open and exposed outside the body through a defect in the lower abdominal wall. It is part of the exstrophy–epispadias complex and may also affect the urethra, pelvic bones, pelvic floor, abdominal muscles, and genitals. It can cause difficulty storing urine, urinary leakage, urine-control problems, and risks to the kidneys.

\medterm{Bladder Exstrophy Repair} Reconstructive surgery for bladder exstrophy that places the bladder inside the body and repairs the lower abdomen and genitals. Treatment is often staged and aims to protect the kidneys and improve urine control.

\medterm{Bladder Fistula} An abnormal passage between the bladder and another organ or the skin, often the vagina, bowel, or uterus. It may cause constant urine leakage, repeated infection, gas or stool in urine, or unusual vaginal wetness.

\medterm{Bladder Health And Disorders} Healthy bladder function stores urine and releases it at an appropriate time. Common problems include infection, urgency, leakage, inability to empty, stones, inflammation, and nerve-related bladder dysfunction.

\medterm{Bladder Infection} An infection of the bladder, usually caused by bacteria. It may cause burning urination, frequent or urgent urination, lower-abdominal discomfort, cloudy or bloody urine, and sometimes fever.

\medterm{Bladder Infection} Alternate name; see \textit{Cystitis}.

\medterm{Bladder Inflammation} Alternate name; see \textit{Interstitial Cystitis}.

\medterm{Bladder Neck} The narrow lower part of the bladder where it joins the urethra. Its muscles help control the release of urine and can contribute to blockage if they scar or fail to relax.

\medterm{Bladder Neck Contracture} Scar-related narrowing of the bladder outlet, often after prostate surgery or other treatment. It can restrict urine flow and cause a weak stream, difficulty emptying, or urinary retention.

\medterm{Bladder Neck Obstruction} Narrowing or blockage where the bladder joins the urethra, making it difficult for urine to leave the bladder. It may cause a weak stream, straining, incomplete emptying, urinary retention, or recurrent infection.

\medterm{Bladder Neoplasm} An abnormal growth arising in the bladder. It may be noncancerous or cancerous; diagnosis requires examination, imaging, cystoscopy, and often a tissue sample to determine its type and extent.

\medterm{Bladder Outlet Obstruction} A blockage where urine leaves the bladder, caused by an enlarged prostate, scarred urethra, stone, tumour, or pelvic-organ prolapse. It may cause a weak stream, difficulty starting, retention, infection, or kidney damage.

\medterm{Bladder Pain} Pain or pressure felt in the lower abdomen or pelvis that may relate to infection, stones, inflammation, bladder pain syndrome, injury, or cancer. Blood in urine, fever, severe pain, or inability to urinate needs prompt assessment.

\medterm{Bladder Pain Syndrome} Long-lasting bladder or pelvic pain, pressure, or discomfort with urinary urgency or frequency when infection and other causes have been excluded. Symptoms often flare and settle and require individualized treatment.

\medterm{Bladder Perforation} A tear or hole in the bladder wall caused by trauma, pelvic surgery, instrumentation,
severe inflammation, or obstruction. Urine may leak into surrounding tissues, causing abdominal pain,
distention, blood in the urine, or infection. It requires urgent medical assessment and may need catheter
drainage or surgical repair.

\medterm{Bladder Problems} Conditions that interfere with storing or passing urine, such as leakage, urgency, repeated infection, stones, or inability to empty. Causes include infection, weak pelvic muscles, nerve disease, prostate enlargement, or blockage.

\medterm{Bladder Prolapse} Alternate name; see \textit{Anterior Vaginal Prolapse}.

\medterm{Bladder Retraining} A planned urination schedule that gradually increases time between bathroom visits to improve bladder control and reduce urgency.

\medterm{Bladder Spasm} An involuntary contraction of the bladder muscle that can cause sudden urgency, cramping,
urine leakage, or pain. It may occur with infection, catheter use, bladder irritation, neurological
conditions, or after surgery. Treatment addresses the cause and may include bladder training or medicines.

\medterm{Bladder Stones} Hard mineral deposits in the bladder that may cause pain, blood in urine, frequent urination, interrupted flow, or difficulty emptying. Some cause no symptoms and are found during tests for another problem.

\synonyms
Bladder Calculi

\medterm{Bladder Trauma} Injury to the bladder caused by a blow, pelvic fracture, penetrating wound, surgery, or a medical instrument. It may cause blood in the urine, lower-abdominal pain, difficulty urinating, or leakage of urine into nearby tissues.

\medterm{Bladder, Urinary} Alternate index label; see \textit{Urinary Bladder}.

\medterm{Blalock-Taussig Operation} Surgery that creates a connection between a major artery and a lung artery to increase blood flow to the lungs in selected blue-baby heart defects. Modern versions usually use a small artificial tube.

\medterm{Bland Diet} A temporary eating pattern using soft, mildly seasoned foods and avoiding items that worsen a person's symptoms. It may help some stomach or bowel problems but does not treat their underlying cause.

\medterm{Blaschko Lines} Invisible patterns formed by the movement of skin cells before birth. Some birthmarks and skin diseases follow these lines as stripes or swirls; they usually do not follow nerves or blood vessels.

\medterm{BLAST} An immature cell that normally develops into a blood cell. A high number of blasts in blood or bone marrow may indicate acute leukaemia or another serious marrow disorder.

\medterm{Blast Cell} An immature precursor cell, especially one in the bone marrow that will normally develop into a blood cell. Too many blasts can indicate acute leukaemia or another marrow disorder.

\medterm{Blast Crisis} An advanced, rapidly worsening phase of some blood cancers in which immature blast cells fill the bone marrow or blood and behave like acute leukaemia. It requires urgent specialist treatment.

\medterm{Blast Injury} Harm caused by an explosion. The pressure wave can injure lungs, ears, or organs; flying objects can cause wounds; being thrown can cause fractures; burns, smoke, toxins, and crushing can add further injuries.

\medterm{Blast Phase} An advanced phase of some blood cancers in which many immature blast cells appear in the blood or marrow. The disease becomes more aggressive and may resemble acute leukaemia.

\medterm{Blastic Plasmacytoid Dendritic-Cell Neoplasm} A rare, fast-growing blood cancer arising from immature immune cells. It often affects the skin, bone marrow, lymph nodes, or nervous system and requires urgent specialist treatment.

\medterm{Blastocyst} An early embryo formed about five to six days after fertilization. It is a hollow group of cells with an inner group that can become the fetus and an outer layer that helps form the placenta.

\medterm{Blastocyst Transfer} An assisted-reproduction procedure in which an embryo grown in the laboratory for about five or six days is placed inside the uterus. The decision to transfer one or more embryos considers success rates, patient factors, and the risk of multiple pregnancy.

\medterm{Blastocystis} A common, single-celled microscopic parasite that lives in the human digestive tract, spread by swallowing food or water contaminated with the organism. While many individuals carry it harmlessly without symptoms, it can cause intestinal illness (blastocystosis) leading to watery diarrhea, abdominal cramps, bloating, gas, nausea, and occasionally skin hives. Diagnosis is confirmed by detecting the parasite during a stool examination under a microscope or with stool DNA tests. When symptomatic and persistent, it is treated with antiparasitic medications such as metronidazole.

\synonyms
Blastocystis Hominis; Blastocystis spp.; Blastocystosis; Intestinal Blastocystis

\medterm{Blastocystis Hominis} Historical name; see \textit{Blastocystis}.

\medterm{Blastocystosis} Alternate name; see \textit{Blastocystis}.

\medterm{Blastoma} A cancer that develops from immature or embryonic cells. The specific name identifies the organ and tumour type, such as a kidney, nerve, liver, or eye blastoma.

\medterm{Blastomyces Dermatitidis} A fungus found in soil and damp organic material that can cause blastomycosis when its spores are inhaled. Infection usually affects the lungs but can spread to skin, bones, joints, or other organs.

\medterm{Blastomycosis} A fungal infection acquired by inhaling Blastomyces spores, usually from soil. It commonly causes pneumonia and may spread to the skin, bones, joints, or other organs; symptoms can develop gradually.

\medterm{Blasts} Immature cells in blood-forming tissue that normally develop into mature blood cells. An excess in blood or bone marrow may indicate acute leukaemia or another serious marrow disorder.

\medterm{Blautia Producta} A bacterium that normally lives in the human intestine and grows without oxygen. It is usually harmless but has occasionally been found in infections, especially in people with serious illness.

\medterm{Bleb} A small thin-walled pocket containing air or fluid. In the lung, a bleb can rupture and cause a collapsed lung; on the skin, the term may describe a blister-like lesion.

\medterm{Bleeder's Disease} Alternate name; see \textit{Hemophilia}.

\medterm{Bleeding} The escape of blood from blood vessels. It may be visible outside the body or occur inside a tissue, organ, or body cavity; severe, unexplained, or persistent bleeding needs medical attention.

\medterm{Bleeding Disorder} A condition causing unusually heavy, prolonged, or spontaneous bleeding because of problems with clotting proteins, platelets, blood vessels, or medicines. Examples include hemophilia, von Willebrand disease, platelet disorders, and some forms of disseminated intravascular coagulation.

\medterm{Bleeding Disorders} Alternate name; see \textit{Bleeding Disorder}.

\medterm{Bleeding During Pregnancy} Vaginal bleeding at any time during pregnancy. Causes range from minor cervical changes to miscarriage, placental problems, or labour; heavy bleeding, pain, dizziness, or reduced fetal movement requires urgent assessment.

\medterm{Bleeding Esophageal Varices} Severe bleeding from enlarged, fragile veins in the lower swallowing tube. These veins usually develop when liver disease raises pressure in veins entering the liver. Vomiting blood or passing black stools is an emergency.

\medterm{Bleeding Gums} Bleeding from the gum tissue, often caused by plaque-related gingivitis, periodontitis, trauma, or vigorous brushing.
It can also occur with blood disorders or medicines that affect clotting. Persistent bleeding, swelling,
loose teeth, or spontaneous bleeding requires dental or medical evaluation.

\medterm{Bleeding Into A Thyroid Nodule} Alternate name; see \textit{Hemorrhagic Thyroid Nodule}.

\medterm{Bleeding Tendency} An increased susceptibility to bleeding, which may result from platelet disorders, coagulation factor deficiencies, vascular abnormalities, or anticoagulant medications. Evaluation includes personal and family history, complete blood count, and coagulation studies.

\medterm{Bleeding Time} An older test measuring how long a small skin puncture takes to stop bleeding. It mainly reflected platelet function and small blood vessels and has largely been replaced by more reliable laboratory tests.

\medterm{Bleeding Varices} Enlarged fragile veins, usually in the swallowing tube or stomach, that have ruptured and are bleeding. Vomiting blood, black stools, weakness, or faintness indicates a life-threatening emergency.

\medterm{Blennorrhea} Thick discharge containing mucus and pus from a moist body surface, especially the urethra or eye. Causes include infection and inflammation, so the affected site and other symptoms guide testing and treatment.

\medterm{Bleomycin} A cancer medicine that damages tumour-cell DNA and is used for selected cancers. Its most important serious risk is lung injury; fever, skin changes, low blood counts, and allergic reactions can also occur.

\medterm{Bleomycin Sulfate} The sulphate form of bleomycin, a cancer medicine that damages DNA. It can cause potentially permanent lung injury, so breathing symptoms and lung function may need monitoring during treatment.

\medterm{Blepharitis} A common, usually chronic inflammation of the eyelid margins that can cause itchy or burning eyelids, redness, tearing, dry-eye symptoms, and crusting at the base of the eyelashes. Anterior blepharitis affects the skin, eyelash follicles, and lash bases; posterior blepharitis affects the meibomian oil glands. It is associated with bacterial overgrowth, dandruff, rosacea, allergies, or gland dysfunction and is not usually contagious.

\synonyms
Eyelid Margin Disease; Meibomian Gland Dysfunction

\medterm{Blepharoconjunctivitis} Inflammation of the eyelid edges and the surface covering the eye. It can cause crusting, itching, redness, tearing, and irritation due to infection, allergy, dry eye, or blocked eyelid glands.

\medterm{Blepharophimosis} Abnormally short horizontal openings between the eyelids. It may occur with drooping eyelids and other facial or developmental features and can interfere with vision.

\medterm{Blepharoplasty} Surgery to reshape or repair an eyelid, sometimes removing or repositioning excess skin, muscle, or fat. It may improve vision blocked by heavy lids, restore an injured lid, or be performed for appearance.

\medterm{Blepharoptosis} Drooping of the upper eyelid because of a weak muscle, damaged nerve, stretched tendon, birth difference, ageing, or a heavy eyelid. It can obstruct vision and may signal a neurological emergency if sudden.

\medterm{Blepharospasm} Involuntary tightening of the muscles around the eyes, causing repeated blinking or forceful eyelid closure. It may occur alone or with eye irritation, medicines, or a movement disorder.

\medterm{Blind} Having severe or complete loss of vision that substantially limits the ability to see. The cause and remaining level of vision are described because blindness varies in degree and may affect one or both eyes.

\medterm{Blind Loop Syndrome} Poor absorption caused by excess bacteria growing in a slow-moving or bypassed part of the small intestine. It may cause bloating, diarrhoea, weight loss, and vitamin deficiencies and can follow surgery or bowel disease.

\medterm{Blind Spot} An area within the field of vision where sight is absent. Everyone has a normal small blind spot where the optic nerve leaves the retina, but a new or enlarged blind spot may indicate eye or nerve disease.

\medterm{Blinded} In a research study, describing a participant, clinician, assessor, or analyst who does not know the assigned treatment or relevant result. This helps reduce expectations from influencing care or measurements.

\medterm{Blinding} Keeping treatment assignments or study information hidden from selected participants, clinicians, assessors, or analysts to reduce expectation and measurement bias. A study states exactly who is unaware and what is concealed.

\medterm{Blindness} Severe or complete loss of vision in one or both eyes caused by disease, injury, or a developmental condition affecting the eye, optic nerve, or brain. Sudden vision loss is an emergency.

\medterm{Blindness And Vision Loss} Reduced or absent ability to see caused by problems in the eye, optic nerve, brain, injury, or development. It may be sudden or gradual, partial or complete, and temporary or permanent.

\medterm{Blinking} Repeated brief closing and reopening of the eyelids that spreads tears, protects the cornea, and clears particles. It may increase with irritation, stress, tics, or vision problems and decrease during prolonged screen use.

\medterm{Blister} A raised pocket of fluid within or beneath the outer layer of skin. It can result from rubbing, burns, heat, infection, allergy, or a skin disease and commonly affects hands and feet.

\synonyms
Bulla; Vesicle When Small

\medterm{Bloating} A feeling of abdominal fullness, tightness, or visible swelling. Gas, constipation, food intolerance, or digestive disorders are common causes; persistent bloating with weight loss, vomiting, pain, or early fullness needs assessment.

\medterm{Block} An obstruction or interruption of flow, movement, electrical signals, or communication. The meaning depends on the structure involved, such as an airway, nerve, blood vessel, or heart.

\medterm{Blockage Of Upper Airway} Obstruction of airflow through the nose, mouth, throat, or voice box. It can rapidly cause suffocation; noisy breathing, inability to speak, blue skin, or severe distress requires immediate emergency help.

\medterm{Blocked Tear Duct} Partial or complete blockage of the passage that drains tears from the eye into the nose. It can cause constant watering, discharge, or repeated infection and is common in babies.

\medterm{Blocked Ureter} Alternate name; see \textit{Ureteral Obstruction}.

\medterm{Blocking Agent} A substance that prevents or reduces a biological process by blocking a receptor, channel, enzyme, or other target. Its effects and risks depend on the target and the tissue affected.

\synonyms
Blocker; Antagonist When It Acts by Receptor Antagonism

\medterm{Blocking Agents} Substances that reduce a body process by blocking a receptor, enzyme, ion channel, nerve signal, or other pathway. Their effects and risks vary according to the target and medicine.

\medterm{Blocking Antibody} An immune protein that prevents another molecule from attaching to its target or interferes with a biological interaction. It may be used as treatment, protect against disease, or contribute to an autoimmune disorder.

\medterm{Blood} A body fluid made of plasma, red blood cells, white blood cells, and platelets. It carries oxygen, nutrients, hormones, and waste and helps fight infection, stop bleeding, and regulate temperature.

\medterm{Blood Alcohol Concentration} The amount of alcohol in a person's blood, usually expressed as a percentage or mass per volume. Effects depend on the level, timing, tolerance, other substances, and health; one result does not fully predict impairment.

\medterm{Blood Alcohol Level} A measurement of alcohol concentration in blood. Its effects depend on the amount, timing, tolerance, other substances, and individual health, so the result must be interpreted with the person's condition.

\medterm{Blood Bank} A specialized medical facility and laboratory service that collects, tests, separates, and safely stores donated blood and its components (such as red blood cells, platelets, and plasma). It performs vital compatibility testing, blood typing, and antibody screening (crossmatching) before transfusions to prevent severe transfusion reactions and ensure patient safety.

\synonyms
Blood Banking; Transfusion Service; Immunohematology Laboratory

\medterm{Blood Banking} Alternate name; see \textit{Blood Bank}.

\medterm{Blood Biomarker} A measurable sign of a disease or condition that can be isolated or detected in a blood sample, such as a disease-associated antibody, hormone, protein, or other substance.

\medterm{Blood Blister} A raised pocket under the skin containing blood, usually caused by rubbing, pinching, pressure, or another injury that breaks small blood vessels.

\medterm{Blood Brain Barrier} A protective filter formed by specialized blood vessels and nearby brain cells. It controls what can pass from the bloodstream into the brain, limiting many medicines, chemicals, germs, and immune cells.

\medterm{Blood Cancer} A cancer affecting blood-forming tissue, blood cells, or the lymphatic system. Abnormal cells can crowd out normal cells, impair immunity, cause anaemia or bleeding, or spread to organs.

\medterm{Blood Cancer} Alternate name; see \textit{Leukemia}.

\medterm{Blood Cell} A cellular part of blood, including red blood cells that carry oxygen, white blood cells that help fight infection, and platelets that help stop bleeding. Most are made in the bone marrow.

\medterm{Blood Cholesterol} Cholesterol carried in the blood. The body needs it to build cell membranes and make hormones and bile acids. Too much LDL cholesterol can form fatty buildup in the arteries and increase the risk of heart disease and stroke. Blood cholesterol is usually checked as part of a lipid profile.

\medterm{Blood Circulation} The continuous movement of blood from the heart through the lungs and body and back again. It delivers oxygen and nutrients and carries away carbon dioxide and other waste.

\medterm{Blood Cleaner} A method used to reduce certain infectious organisms in donated blood or blood components while keeping them suitable for transfusion. It may use chemicals or light, but it does not remove every possible infection risk.

\medterm{Blood Clot} A mass of thickened blood formed during clotting. It is protective when it seals an injury, but a clot inside a blood vessel can block flow or travel to the lungs, brain, heart, or another organ.

\medterm{Blood Clot in Lung} Alternate name; see \textit{Pulmonary Embolism}.

\medterm{Blood Clots} Clumps of thickened blood. Clots can stop bleeding, but abnormal clots in veins or arteries may block blood flow and cause deep-vein thrombosis, pulmonary embolism, stroke, heart attack, or tissue damage.

\synonyms
Thrombi; Thrombosis; Blood Clot

\medterm{Blood Coagulation} The linked process that changes blood from a liquid to a firm clot after injury. Too little clotting causes bleeding, while excessive or misplaced clotting can cause dangerous vessel blockage.

\medterm{Blood Coagulation Disorder} A condition in which blood clots too little, too much, or at the wrong time. Bleeding disorders result from inadequate clot formation, whereas thrombophilic disorders cause excessive or misplaced clotting. Causes include genetic variants, medicines, liver disease, immune conditions, cancer, and other illness. Examples include hemophilia, von Willebrand disease, thrombocytopenia, thrombophilia, and disseminated intravascular coagulation.

\medterm{Blood Count} A blood test measuring the number and features of red cells, white cells, and platelets. It helps assess anaemia, infection, inflammation, bleeding risk, and disorders of the bone marrow.

\medterm{Blood Culture} A test in which blood is placed in special bottles to see whether bacteria or fungi grow. It helps investigate bloodstream infection, sepsis, or infection of the heart valves.

\medterm{Blood Culture Contamination} Accidental entry of skin or environmental organisms into a blood-culture bottle during collection or handling, rather than true infection in the blood. It can lead to unnecessary antibiotics, so results are interpreted with symptoms and other culture sets.

\medterm{Blood Differential Test} A blood test measuring the numbers and proportions of different white blood cells, including neutrophils, lymphocytes, monocytes, eosinophils, and basophils. It can help investigate infection, inflammation, allergy, leukaemia, or marrow disease.

\medterm{Blood Donation} Voluntary collection of whole blood or selected blood components from a screened donor for transfusion or manufacture of medical products. Eligibility checks, infection testing, sterile collection, and post-donation recovery protect donor and recipient.

\medterm{Blood Donation Before Surgery} Giving and storing a person's own blood before an operation for possible later transfusion. It is suitable only in selected situations and can cause anaemia or delays if the operation is postponed.

\medterm{Blood Donor} A person who gives blood or blood components for transfusion or medical use. Screening protects the donor and helps ensure that donated blood is safe for the recipient.

\medterm{Blood Draw} Collection of blood, usually with a needle placed in a vein, for testing, donation, or treatment. Temporary pain, bruising, dizziness, or rarely infection may occur.

\medterm{Blood Film} A thin layer of peripheral blood spread onto a glass microscope slide, stained with a Romanowsky stain (e.g., Wright-Giemsa), and examined under microscopy for quantitative and qualitative evaluation of erythrocyte morphology, leukocyte differentials, platelet adequacy, and identification of bloodborne parasites.

\synonyms
Peripheral Blood Smear; Blood Smear; Peripheral Smear

\medterm{Blood Flow} The movement of blood through vessels and organs. It depends mainly on the pressure difference and resistance in the vessels and delivers oxygen and nutrients while removing waste.

\medterm{Blood Gas Analysis} A test measuring oxygen, carbon dioxide, acidity, and related values in blood. It helps assess breathing, oxygen delivery, and acid-base balance in serious illness.

\medterm{Blood Gas, Venous} The clinical laboratory measurement of pH, partial pressure of carbon dioxide ($p\text{CO}_2$), bicarbonate ($\text{HCO}_3^-$), and base excess in a peripheral or central venous blood sample to assess systemic acid-base balance, ventilatory status, and tissue perfusion.

\synonyms
Venous Blood Gas; VBG; Central Venous Blood Gas

\medterm{Blood Gases} Measurements of oxygen, carbon dioxide, acidity, and related substances in blood. They help determine how well a person is breathing, exchanging gases, and maintaining acid-base balance.

\medterm{Blood Glucose} The amount of sugar called glucose in the blood. The body regulates it using food absorption, liver production, insulin, and other hormones; levels that are too low or high can cause illness.

\medterm{Blood Glucose Monitoring} The regular measurement of sugar (glucose) levels in the blood, primarily used to manage diabetes and detect dangerously high (hyperglycemia) or low (hypoglycemia) blood sugar. Testing is commonly performed using a handheld glucometer with a small drop of blood from a fingertip prick, or with a wearable continuous glucose monitor (CGM) that checks glucose in the fluid under the skin around the clock. It tracks how meals, physical activity, illness, and medications (such as insulin) affect glucose levels throughout the day.

\synonyms
Blood Sugar Monitoring; Self-Monitoring of Blood Glucose; SMBG; Continuous Glucose Monitoring; CGM

\medterm{Blood Group} A classification of blood based on specific inherited markers (antigens) on the surface of red blood cells. The two most vital systems are the ABO system (types A, B, AB, and O) and the Rhesus (Rh) system (positive or negative based on the RhD marker). Accurate blood grouping and crossmatching are essential before transfusions to prevent life-threatening immune reactions and during pregnancy to prevent Rh disease (hemolytic disease of the newborn).

\synonyms
Blood Type; Blood Groups; ABO and Rh Grouping; Blood Typing

\medterm{Blood Group, ABO} Alternate name; see \textit{ABO Blood Group System}.

\medterm{Blood Groups} Alternate name; see \textit{Blood Group}.

\medterm{Blood In Semen} Blood in the fluid released during ejaculation. It is often temporary and may follow infection, inflammation, injury, or a urinary-tract procedure. Persistent or recurrent blood, pain, fever, or urinary symptoms needs medical assessment.

\synonyms
Hematospermia; Hemospermia

\medterm{Blood in Stool} Alternate name; see \textit{Rectal Bleeding}.

\medterm{Blood In The Eye} Bleeding on the eye surface or inside the eye. A small surface bleed may be harmless, but bleeding inside the eye, eye pain, injury, or any vision change requires prompt eye assessment.

\medterm{Blood in Urine} The abnormal presence of red blood cells in the urine, detected either macroscopically as pink, red, or tea-colored urine (gross hematuria) or microscopically on urinalysis ($\ge 3$ erythrocytes per high-power field). Common etiologies include urinary tract infections, nephrolithiasis, glomerular nephritis, benign prostatic hyperplasia, and urothelial malignancies.

\synonyms
Hematuria; Gross Hematuria; Microscopic Hematuria; Blood in the Urine

\medterm{Blood Lancet} Alternate name; see \textit{Lancet}.

\medterm{Blood Lead Level Test} A quantitative diagnostic laboratory assay measuring venous or capillary whole-blood lead concentration in micrograms per deciliter ($\mu\text{g/dL}$). Utilized for pediatric screening, occupational exposure monitoring, and evaluating symptomatic plumbism.

\synonyms
BLL Assay; Venous Blood Lead Screen; Plumbism Blood Test
\medterm{Blood Loss} Loss of blood from the circulation because of injury, surgery, childbirth, internal bleeding, or disease. Severe loss can reduce oxygen delivery and cause shock, requiring urgent control of bleeding and replacement of blood or fluids.

\medterm{Blood Low Red Cell Count} Alternate name; see \textit{Anemia}.

\medterm{Blood PH} A measure of how acidic or alkaline the blood is. The lungs regulate carbon dioxide and the kidneys regulate acids and bicarbonate; normal arterial pH is about 7.35--7.45.

\medterm{Blood Plasma} The liquid part of blood, made mostly of water and containing proteins, salts, nutrients, hormones, gases, and waste products. It carries blood cells and dissolved substances throughout the body.

\medterm{Blood Poisoning} A non-technical, lay term historically used to describe a severe bloodstream bacterial infection (septicemia). In modern medicine, this condition is formally diagnosed and treated as \textit{Sepsis}---a life-threatening condition in which the body's overwhelming response to an infection causes organ dysfunction.

\synonyms
Septicemia; Sepsis; Bloodstream Infection; Bacteremia

\medterm{Blood Preservation} The controlled collection, processing, cooling, freezing, and storage of donated blood or its components so they remain safe and effective until transfusion.

\medterm{Blood Pressure} The force of circulating blood against artery walls, recorded as two numbers in millimetres of mercury. The systolic number is the pressure when the heart contracts; the diastolic number is the pressure between beats.

\medterm{Blood Pressure Chart: Reading By Age} A guide for interpreting blood-pressure readings in relation to age and clinical circumstances. Diagnosis usually requires repeated, correctly measured readings; pregnancy, illness, and treatment goals can change the appropriate range.

\medterm{Blood Pressure Cuff} The inflatable part of a blood-pressure measuring device that wraps around a limb and briefly compresses an artery. Cuffs come in different sizes and may be used with a manual sphygmomanometer or an automated monitor.

\medterm{Blood Pressure Measurement} Measurement of arterial pressure with an inflatable cuff or another approved device. The result records systolic pressure when the heart contracts and diastolic pressure when it relaxes; correct cuff size and positioning improve accuracy.

\medterm{Blood Pressure Monitors For Home} Devices that measure blood pressure outside a clinic, usually with an inflatable upper-arm cuff. A validated monitor, correctly sized cuff, proper positioning, and repeated readings give more reliable results.

\medterm{Blood Pressure, Low} A systolic blood pressure below $90\text{ mmHg}$ or diastolic blood pressure below $60\text{ mmHg}$, resulting in inadequate end-organ perfusion. Etiologies include hypovolemia, septic shock, cardiogenic impairment, adrenal insufficiency, orthostatic autonomic failure, and antihypertensive overdosage.

\synonyms
Low Blood Pressure; Hypotension; Orthostatic Hypotension

\medterm{Blood Product} A component made from donated blood for medical treatment, such as red blood cells, platelets, plasma, or cryoprecipitate. The product is selected according to the person's bleeding, anaemia, or clotting problem.

\medterm{Blood Products} Components prepared from donated blood for transfusion, including red blood cells, platelets, plasma, and cryoprecipitate. Each replaces a different blood function, and the choice depends on the person's symptoms, bleeding, and laboratory results.

\medterm{Blood Protein} A protein dissolved in blood plasma. These proteins help transport substances, maintain fluid balance, support immunity, and control clotting; abnormal amounts can occur with inflammation, liver or kidney disease, immune disorders, or some cancers.

\medterm{Blood Results} Measurements obtained from a blood sample, such as cell counts, glucose, salts, hormones, or organ-function markers. They must be interpreted using the laboratory's reference range, symptoms, medicines, and other findings.

\medterm{Blood Sedimentation} The settling of red blood cells in a tube of treated blood. The erythrocyte sedimentation rate is a nonspecific sign that may rise with inflammation, infection, anaemia, pregnancy, or other conditions.

\medterm{Blood Smear} A thin layer of blood spread on a glass slide, stained, and examined under a microscope. It shows the size, shape, and appearance of blood cells and can reveal abnormal cells, parasites, or cell fragments.

\medterm{Blood Spatter Analysis} The forensic examination of bloodstain size, shape, and distribution to help reconstruct events. It may distinguish drops caused by gravity, contact transfer, impact, or forceful expulsion, but conclusions have limits.

\medterm{Blood Substitutes} Fluids or artificial products intended to replace some functions of blood. Most restore circulating volume; some are designed to carry oxygen. They do not replace every function of blood and are used only in selected situations.

\medterm{Blood Sugar} The amount of glucose in the blood. Insulin, glucagon, the liver, food, activity, illness, and other hormones affect it. Persistently high or very low levels can cause symptoms and serious organ injury.

\medterm{Blood Sugar Test} A test measuring glucose in the blood. It may be done while fasting, at a random time, or after a glucose drink, and helps diagnose or monitor diabetes and other disorders of glucose control.

\medterm{Blood Test} Analysis of blood cells or substances dissolved in blood, including chemicals, proteins, hormones, antibodies, or genetic material. Results help investigate symptoms, diagnose disease, assess organ function, and monitor treatment.

\medterm{Blood Thinner} A common name for an anticoagulant medicine that lowers the blood's ability to form harmful clots. It does not actually make the blood thinner. These medicines help prevent or treat conditions such as deep vein thrombosis, pulmonary embolism, and stroke, but can increase bleeding risk.

\synonyms
Anticoagulant

\medterm{Blood Transfusion} The administration of whole blood or a specific blood component, such as red blood
cells, platelets, plasma, or cryoprecipitate, to treat blood loss, anaemia, thrombocytopenia,
or a clotting-factor deficiency. Red cells improve oxygen-carrying capacity; platelets,
plasma, and cryoprecipitate support haemostasis. ABO/Rh typing, antibody screening,
and compatibility testing help reduce transfusion reactions. Risks include febrile or
allergic reactions, acute haemolysis, circulatory overload, and rare transfusion-transmitted
infection; the decision depends on the patient's symptoms, laboratory results, bleeding,
and overall clinical context.

\medterm{Blood Typing} A laboratory test identifying markers on red blood cells, especially the ABO and Rh systems. It helps select compatible blood for transfusion and identifies some pregnancy-related risks from blood-group incompatibility.

\medterm{Blood Urea Nitrogen} The amount of nitrogen from urea in the blood. Urea is made when the body breaks down protein and is removed mainly by the kidneys; the result is interpreted with creatinine, hydration, and the clinical situation.

\medterm{Blood Urea Nitrogen Test} A blood test measuring urea nitrogen, a waste product made during protein breakdown. It helps assess kidney function and hydration, but abnormal results can also reflect diet, bleeding, medicines, or other illness.

\medterm{Blood Urea Nitrogen Test} A clinical serum chemistry test measuring the amount of nitrogen derived from urea, the principal end product of protein catabolism formed in the liver. Elevated serum concentrations (azotemia) occur in acute or chronic renal impairment, prerenal hypoperfusion, severe dehydration, and upper gastrointestinal hemorrhage.

\synonyms
BUN Test; Serum Blood Urea Nitrogen; Urea Nitrogen Test

\medterm{Blood Vessel} A tube that carries blood through the body. Arteries carry blood away from the heart, veins return it, and tiny capillaries allow oxygen, nutrients, and waste products to pass between blood and tissues.

\medterm{Blood Vessels} The arteries, arterioles, capillaries, venules, and veins that carry blood through the body. They deliver oxygen and nutrients, remove waste, and help control blood pressure, temperature, and fluid exchange.

\medterm{Blood Viscosity} How thick blood is and how much it resists flowing. It is influenced mainly by the number of red blood cells, but also by blood proteins and temperature. Excessively thick blood can slow circulation and increase clot risk.

\medterm{Blood Volume} The total amount of blood circulating in the body, including plasma and blood cells. Too little can cause low blood pressure and shock; too much can strain the heart and lungs.

\medterm{Blood-Borne Disease} An infection that can spread through infected blood or certain body fluids, for example through shared needles, needlestick injuries, sexual exposure, or contaminated medical equipment. Examples include HIV and hepatitis B and C.

\medterm{Blood-Borne Pathogens} Germs that can spread through infected blood or certain body fluids. Examples include HIV and hepatitis B and C viruses; prevention includes vaccination where available, gloves, safe sharps handling, and appropriate exposure testing.

\medterm{Blood-Brain Barrier} A protective filter formed by specialized cells lining brain blood vessels and their supporting tissues. It controls which medicines, chemicals, germs, and immune cells can move from the blood into brain tissue.

\medterm{Bloodborne Pathogens} Infectious agents spread through infected blood or certain body fluids. Important examples include HIV and hepatitis B and C viruses; exposure may require urgent medical assessment, testing, and preventive treatment.

\medterm{Bloom Syndrome} A rare inherited disorder caused by changes in both copies of the BLM gene. It causes short stature, a sun-sensitive facial rash, frequent infections, and a greatly increased risk of several cancers.

\medterm{Blount Disease} A growth disorder in which the upper part of the shinbone grows unevenly, causing worsening bowing of one or both legs. It can affect young children or adolescents and may cause an abnormal gait, unequal leg length, or knee damage.

\medterm{Blue Baby} An infant whose skin, lips, or nail beds look blue because the blood carries too little oxygen or circulation is inadequate. This can signal heart, lung, airway, blood, or other serious disease and needs urgent assessment.

\medterm{Blue Bloater} Outdated descriptive label; use the specific diagnosis, such as \textit{Chronic Bronchitis} or \textit{Chronic Obstructive Pulmonary Disease}.

\medterm{Blue Bottle} A common name for a bluebottle jellyfish. Its tentacles can inject venom that causes immediate stinging pain, raised marks, redness, and sometimes muscle cramps, breathing difficulty, or a severe allergic reaction.

\medterm{Blue Discoloration Of The Skin} A blue or gray tint of the skin, lips, or nail beds caused by low oxygen, reduced blood flow, or abnormal haemoglobin. Sudden or widespread discoloration, especially with breathing difficulty or confusion, is an emergency.

\medterm{Blue Nevi} Usually harmless moles formed by pigment-producing cells deeper in the skin, giving them a blue, blue-gray, or blue-black color. A new, enlarging, or changing lesion requires examination because melanoma can sometimes look similar.

\medterm{Blue Nevus} A usually harmless blue-gray or blue-black mole caused by pigment-producing cells located deeper than those in most moles. A new, changing, bleeding, or painful lesion requires examination.

\medterm{Blue Nightshade Poisoning} Illness caused by eating blue nightshade or a related poisonous plant containing glycoalkaloids. Symptoms may include vomiting, diarrhoea, abdominal pain, confusion, weakness, seizures, or abnormal heart rhythm; it is an emergency requiring urgent medical treatment.

\medterm{Blue NæVus} A usually harmless blue-gray or blue-black mole formed by pigment-producing cells deep in the skin. It requires assessment if it changes in size, shape, colour, or symptoms.

\medterm{Blue Sclera} A blue-gray appearance of the normally white outer layer of the eye, usually caused by thinning that allows deeper pigment to show through. It can occur in inherited connective-tissue or bone disorders and needs clinical assessment.

\medterm{Blue Sclerae} A distinct bluish or grayish discoloration of the ocular sclera caused by abnormal thinness or translucency of the scleral collagen fibers, which allows the pigment of the underlying uveal tract and choroid to be visible. It is a hallmark physical finding of osteogenesis imperfecta and can also occur in Ehlers-Danlos syndrome, Marfan syndrome, and severe iron deficiency anemia.

\medterm{Blue-Green Algae} Blue-green algae are cyanobacteria that can grow rapidly in water. Some release toxins that cause skin irritation, vomiting, liver injury, breathing problems, or neurologic illness when people or animals are exposed.

\medterm{Blumberg Sign} Sudden sharp pain when a doctor presses down gently on the belly and then quickly lets go, showing that the lining of the abdomen is swollen and irritated.

\medterm{Blunt Abdominal Trauma} An injury to the abdomen without an object piercing the skin, often caused by a crash, fall, or blow. It can damage the spleen, liver, kidneys, bowel, or pancreas, sometimes without early visible bruising.

\medterm{Blunt Force Trauma} Injury caused by a blow from a blunt object or by striking a blunt surface, without penetration by a sharp edge. It can produce bruises, grazes, torn skin, internal bleeding, fractures, or organ damage.

\medterm{Blunted Affect} Reduced outward expression of emotion. A person may show less facial movement, voice variation, or body language than expected, although their internal feelings may be stronger. It can occur in mental or neurologic disorders.

\medterm{Blurred Vision} Vision that appears out of focus, hazy, or unclear. Causes include an incorrect glasses prescription, dry eye, eye disease, medicines, migraine, or neurologic illness. Sudden blurring, pain, weakness, or vision loss needs urgent assessment.

\medterm{Blush} Temporary facial redness or warmth caused by increased blood flow near the skin surface. It may follow emotion, heat, exercise, fever, alcohol, medicines, hormonal changes, or some medical conditions.

\medterm{Blushing} Temporary redness or warmth of the face caused by increased blood flow in the skin. Common triggers include emotion, heat, exercise, alcohol, hormonal changes, medicines, fever, and some illnesses.

\medterm{Board-Like Rigidity} Severe involuntary tightening of the abdominal muscles, making the abdomen feel hard. It can indicate irritation of the abdominal lining from serious inflammation, infection, bleeding, or perforation and requires urgent assessment.

\medterm{Boas' Sign} Pain or unusual skin sensitivity below the right shoulder blade, traditionally associated with gallbladder disease. It is not specific enough to establish a diagnosis and must be interpreted with other findings.

\medterm{Bobbing, Ocular} Alternate index label; see \textit{Ocular Bobbing}.

\medterm{Bobble-Head Doll Syndrome} A rare childhood movement disorder causing repeated forward-and-back head movements. It is often linked to a fluid-filled cyst or other lesion near the third ventricle of the brain and requires neurologic evaluation and imaging.

\medterm{Bocavirus} A group of small viruses that can infect humans and animals. Human bocavirus most often affects the airways of young children, although infection may cause illness in other tissues or people with weakened immunity.

\medterm{Bocavirus Infection} A viral infection, usually of the airways in young children. It may cause a runny nose, cough, fever, wheezing, bronchiolitis, or pneumonia, often together with another infection. Treatment is usually supportive; breathing difficulty or dehydration needs prompt care.

\medterm{Bochdalek Hernia} A birth defect in which abdominal organs move through an opening in the back and side of the diaphragm into the chest. It can prevent normal lung development and cause severe breathing problems in a newborn.

\medterm{Bock's Disease} An older name for sarcoidosis, an inflammatory disease in which clusters of immune cells form in organs. It most often affects the lungs and lymph nodes but can also affect the skin, eyes, heart, or nervous system.

\medterm{BODE Index} A scoring tool for chronic obstructive pulmonary disease that combines body weight, breathing-test results, breathlessness, and walking distance. A higher score suggests more severe disease and a greater risk of hospital admission or death; it does not predict an individual outcome with certainty.

\synonyms
BODE Score; COPD Multidimensional Index

\medterm{Bodily Secretion} A substance made and released by cells or glands, such as sweat, saliva, mucus, tears, digestive juices, or hormones. Secretions protect tissues, aid digestion, regulate body functions, or help remove waste.

\medterm{Body} The physical organism, including its cells, tissues, organs, and interacting body systems. Health depends on the structure and function of these parts and how they respond to internal and external conditions.

\medterm{Body Aches} Widespread muscle or body pain that can occur with infection, fever, exercise, inflammation, medicines, or other illness. Severe pain with weakness, dark urine, breathing difficulty, chest pain, or high fever needs prompt assessment.

\medterm{Body Cavity} An enclosed space inside the body that contains and protects organs. Major cavities include the skull, spinal canal, chest, abdomen, and pelvis; their linings may produce small amounts of lubricating fluid.

\medterm{Body Cell Mass} The living, metabolically active cells of the body, including cells in muscles, organs, blood, and other tissues. It excludes stored fat and water outside cells and is sometimes estimated in nutrition assessment.

\medterm{Body Composition} The relative amounts of fat, muscle and other lean tissue, bone mineral, and water in the body. It provides more information about health and nutrition than body weight alone, although estimates vary by method.

\medterm{Body Dysmorphia} Alternate name; see \textit{Body Dysmorphic Disorder}.

\medterm{Body Dysmorphic Disorder} A mental-health disorder involving persistent distress about one or more perceived flaws in appearance that are minor or not noticeable to others. Repetitive checking, grooming, comparison, or avoidance interferes with daily life.

\synonyms
BDD

\medterm{Body Fat} Adipose tissue that stores energy, cushions organs, insulates the body, and releases hormones and other signals. Too much or too little body fat can affect metabolism, movement, fertility, immunity, and other aspects of health.

\medterm{Body Fluid} The water and dissolved substances (such as electrolytes, nutrients, proteins, and gases) that make up approximately 60\% of human body weight. Body fluids are organized into two major compartments: intracellular fluid (fluid inside cells, making up about two-thirds) and extracellular fluid (fluid outside cells, including interstitial fluid between tissues and blood plasma), along with specialized fluids such as cerebrospinal fluid, synovial joint fluid, lymph, and digestive secretions. They are essential for transporting nutrients and waste, regulating body temperature, and lubricating tissues and joints.

\synonyms
Body Fluids; Bodily Fluids; Total Body Water; Body Fluid Compartments

\medterm{Body Fluids} Alternate name; see \textit{Body Fluid}.

\medterm{Body Habitus} A person's general build, including body size, shape, muscle, fat distribution, and proportions. Clinicians may record it during an examination, but it does not by itself diagnose disease or predict health.

\medterm{Body Height} The vertical distance from the soles of the feet to the top of the head, measured while standing when possible. It helps assess growth and nutrition and is used in medication, kidney-function, and other calculations.

\medterm{Body Identification} The process of establishing who a deceased person is. It may use fingerprints, dental records, medical information, DNA, personal effects, or visual recognition, with the method chosen according to the circumstances.

\medterm{Body Image} A person's thoughts, feelings, and beliefs about their body and appearance. Distress or inaccurate beliefs about appearance can occur in body dysmorphic disorder, eating disorders, depression, or other mental-health conditions.

\medterm{Body Lice} Infestation with human body lice, which live mainly in clothing and move to the skin to feed. Bites cause intense itching and scratching; treatment includes washing or heat-treating clothing and bedding and prescribed insecticide when needed.

\synonyms
Pediculosis Corporis; Body Louse Infestation

\medterm{Body Mass Index} A screening measure calculated by dividing weight in kilograms by height in meters squared. It helps assess weight in relation to height, but it does not directly measure body fat or provide a complete assessment of health. Results are interpreted with age and other clinical factors; children and teens are assessed using age- and sex-specific percentiles.

\synonyms
BMI; Quetelet Index; BMI Calculation

\medterm{Body Odor} Smell produced when skin secretions, especially sweat, are broken down by skin microorganisms. A sudden or unusual change may be related to diet, medicines, infection, or a metabolic disorder.

\medterm{Body Odor and Sweating} Alternate index label; see \textit{Sweating and Body Odor}.

\medterm{Body Odour} Smell from the breakdown of skin secretions by microorganisms. The smell can change with sweating, hygiene, diet, medicines, infection, or some metabolic diseases.

\medterm{Body Piercing} Creation of an opening through skin or a mucous membrane to insert jewellery. Possible complications include pain, bleeding, infection, scarring, allergic reaction, and damage to teeth, nerves, or cartilage.

\medterm{Body Plethysmography} A breathing test performed inside a sealed cabin to measure lung volumes, including air that cannot be exhaled. It helps identify restricted lungs, over-expanded lungs, and some airway problems.

\medterm{Body Position} The arrangement of the body or a body part, such as standing, lying on the back, lying face down, or lying on the side. Position can affect breathing, circulation, examination results, and medical procedures.

\medterm{Body Region} A named area of the body used to describe the location of an organ, symptom, injury, or examination finding, such as the chest, abdomen, pelvis, or arm.

\medterm{Body Ringworm} A fungal infection of the skin on the body. It usually causes an expanding, ring-shaped, scaly rash with a more active edge and clearer centre; it spreads by contact and is treatable with antifungal medicine.

\synonyms
Tinea Corporis; Dermatophytosis of Glabrous Skin

\medterm{Body Space} A normal or potential compartment inside the body, such as the space around the lungs or abdominal organs. Air, blood, fluid, or infection can collect there and interfere with organ function.

\medterm{Body Substance} Material produced by or found in the body, including blood, saliva, urine, faeces, mucus, tissue, and other secretions. Some substances can carry infection and require appropriate handling precautions.

\medterm{Body Surface} The outside of the body, mainly the skin and accessible mucous membranes. Clinicians may estimate its area when assessing burns, calculating some medicine doses, or making physiological measurements.

\medterm{Body Surface Area} An estimate of the body's external area, usually calculated from height and weight. It is used in burn assessment, some medicine doses, and selected medical calculations; it is an estimate rather than a direct measurement.

\medterm{Body Temperature} The degree of heat in the body, determined by the balance between heat production and heat loss. It varies with age, activity, time of day, illness, and measurement site; fever and hypothermia require clinical interpretation.

\medterm{Body Temperature Norms} Expected temperature ranges vary with age, time of day, activity, illness, and whether temperature is measured by mouth, ear, armpit, skin, or rectum. The measurement method and symptoms are considered when interpreting fever or low temperature.

\medterm{Body Wall} The layers of skin, connective tissue, muscle, and sometimes bone that surround and protect a body compartment, such as the chest or abdomen. They also help movement, breathing, and support.

\medterm{Body Weight} The measured mass of a person, including fat, muscle, bone, water, and organs. It helps assess growth, nutrition, and fluid changes and is used with other information for some medical calculations.

\medterm{Boeck's Sarcoid} An older name for sarcoidosis, an inflammatory disease in which clusters of immune cells form in tissues. It commonly affects the lungs and nearby lymph nodes but may also involve the skin, eyes, heart, or nervous system.

\medterm{Boerhaave Syndrome} A sudden full-thickness tear of the food pipe, usually after forceful vomiting. Leakage into the chest can cause severe chest or upper-abdominal pain, infection, shock, and breathing problems; it is a surgical emergency.

\medterm{Boerhaave's Syndrome} A spontaneous full-thickness tear of the food pipe, most often following forceful vomiting. Contents can leak into the chest and cause severe infection, chest pain, breathing difficulty, and shock.

\medterm{Bogros' Space} A layer of tissue behind the groin region and in front of the abdominal lining. Surgeons use this natural plane when repairing some groin hernias.

\medterm{Bohn Nodules} Small, harmless cysts containing keratin that appear on a newborn's gums or the roof of the mouth. They usually disappear without treatment and do not interfere with feeding.

\medterm{Bohr Effect} The tendency of haemoglobin to release more oxygen when tissues become more acidic or contain more carbon dioxide. This helps active tissues receive oxygen where energy use and carbon-dioxide production are greatest.

\medterm{Bohring Syndrome} A rare genetic disorder associated with severe developmental delay, poor growth before birth, distinctive facial features, reduced movement, and frequent feeding or breathing difficulties.

\medterm{Boil} A painful, pus-filled skin abscess centered on an infected hair follicle and surrounding tissue, most commonly caused by \textit{Staphylococcus aureus} bacteria. It typically begins as a tender, red nodule that progressively enlarges, develops a central yellowish head of pus, and eventually ruptures to drain. Multiple or recurrent boils are termed furunculosis, whereas a deeper, connected cluster of several boils is classified as a \textit{Carbuncle}.

\synonyms
Furuncle; Furunculosis; Boils; Skin Abscess

\medterm{Boils} Alternate name; see \textit{Boil}.

\medterm{Boils and Carbuncles} Alternate name; see \textit{Boil} and \textit{Carbuncle}.

\medterm{Bolt Device} A bolt is a threaded fastener used to hold parts together. In medicine, the term may describe a device component or a specific product name, so its meaning depends on the clinical context.

\medterm{Bolus} A single measured amount of medicine or fluid given over a short time, usually by injection or infusion. In digestion, it means food that has been chewed and is ready to swallow.

\medterm{Bolus Gastrostomy Feeding} The administration of prescribed liquid enteral formula or hydration directly into the stomach through a gastrostomy tube over a short period using a syringe or gravity drip. Involves tube flushing before and after each feeding to maintain patency.

\synonyms
Bolus G-Tube Feeding; Bolus Enteral Nutrition; Syringe Feeding

\synonyms
G-Tube Bolus Feeding; Enteral Bolus Nutrition
\medterm{Bolus Infusion} Delivery of a measured amount of fluid or medicine into a vein over a short period. The dose and speed depend on the medicine, the person's condition, and the intended effect; errors can cause serious harm.

\medterm{Bombay Blood Group} A rare inherited blood type in which red blood cells lack the H marker needed to form A or B markers. It may look like type O but is usually compatible only with blood from another Bombay-type donor.

\medterm{Bombesin Receptor} A protein on or inside a cell that binds bombesin-like signalling substances and changes cell activity. These receptors help regulate processes such as hormone release, digestion, and cell growth and are studied in some cancers.

\medterm{Bone} A strong living tissue that forms the skeleton, supports the body, protects organs, enables movement with muscles, stores minerals, and contains marrow that makes blood cells. Bone is continually broken down and rebuilt.

\medterm{Bone-Anchored Hearing Implant} A hearing device fixed to the skull bone that transmits sound vibrations to the inner ear, bypassing the outer or middle ear. It may be used for selected conductive, mixed, or single-sided hearing loss.

\medterm{Bone And Mineral Disorders} Diseases affecting bone strength, bone formation, or the body's balance of calcium and phosphate. Examples include osteoporosis, rickets, osteomalacia, parathyroid disorders, and inherited bone diseases.

\medterm{Bone Callus} New tissue that forms around a broken bone during healing. It first consists of a soft blood clot and connective tissue, then cartilage and immature bone, which are gradually reshaped into stronger bone.

\medterm{Bone Cancer} Cancer that starts in bone, such as osteosarcoma or chondrosarcoma, or cancer that has spread there from another organ. Persistent focal pain, swelling, or a fracture after minor injury may occur; diagnosis uses imaging and tissue testing.

\synonyms
Bone Sarcoma; Primary Bone Malignancy

\medterm{Bone Cement} A material, usually polymethylmethacrylate, used to secure artificial joints or fill spaces in bone. It hardens chemically and may be used in joint replacement, vertebroplasty, kyphoplasty, or bone reconstruction. Possible complications include cement leakage, infection, loosening, and bone cement implantation syndrome.

\medterm{Bone Cement} A self-curing acrylic biomaterial composed of polymethylmethacrylate (PMMA) utilized in orthopedic and spine surgery to anchor prosthetic joint components, fill osseous voids, and stabilize vertebral compression fractures during vertebroplasty or kyphoplasty.

\synonyms
Polymethylmethacrylate; PMMA; Acrylic Bone Cement

\medterm{Bone Conduction} Hearing sound through vibrations carried by the skull to the inner ear, bypassing the outer and middle ear. Bone-conduction testing helps identify whether hearing loss comes from the outer or middle ear or from the inner ear and nerve.

\medterm{Bone Cyst} A fluid-filled or blood-filled space within bone. Many bone cysts are benign, but they can weaken the bone and cause pain, swelling, or fracture. Common types include simple bone cysts and aneurysmal bone cysts.

\medterm{Bone Densitometry} See \textit{Bone Density Screening}.

\medterm{Bone Density} The amount of mineral packed into a measured area of bone. It is usually estimated with a DXA scan; lower density is associated with greater fracture risk, but results must be interpreted with age and other risk factors.

\medterm{Bone Density Screening} A DXA scan that measures bone mineral at sites such as the hip and spine. It helps identify osteoporosis and estimate fracture risk, especially in people with age-related or other risk factors.

\medterm{Bone Development} The formation, growth, and shaping of bones from before birth through adulthood. Bones lengthen and widen as new tissue is made and old tissue is removed; hormones, nutrition, activity, and injury influence this process.

\medterm{Bone Disease} Any condition that weakens, deforms, inflames, infects, or otherwise damages bone. Examples include osteoporosis, fractures, infections, tumours, inherited skeletal disorders, and diseases caused by abnormal mineral or hormone levels.

\medterm{Bone Fracture} A partial or complete break in a bone caused by injury, repeated stress, or weakened bone. Pain, swelling, bruising, deformity, and loss of function may occur. An open fracture or loss of sensation or circulation needs emergency care.

\medterm{Bone Fractures} Alternate name; see \textit{Bone Fracture}.

\medterm{Bone Graft} Bone tissue or a substitute placed in a bone defect to support repair, join bones together, rebuild an area, or help a fracture heal. The material may come from the person, a donor, or a synthetic source.

\medterm{Bone Graft, Dental} An oral and maxillofacial surgical procedure involving the placement of autologous bone, allograft, xenograft, or synthetic osteoconductive matrix into the maxilla or mandible to reconstruct alveolar ridge deficiencies or augment the maxillary sinus prior to dental implant placement.

\synonyms
Dental Bone Graft; Alveolar Ridge Augmentation; Maxillofacial Bone Grafting

\medterm{Bone Health} The strength, structure, and ability of bones to grow, repair, and resist fractures. It is influenced by age, nutrition, hormones, physical activity, medicines, smoking, alcohol, and medical conditions.

\medterm{Bone Infection} Infection of bone, usually caused by bacteria. It may reach bone through the bloodstream, spread from nearby tissue, or enter through surgery or an open injury, causing pain, swelling, fever, or poor healing.

\medterm{Bone Lesion Biopsy} Removal of a small sample from an abnormal area of bone with a needle or operation for laboratory examination. Testing can distinguish infection, noncancerous growth, and cancer; the route must be planned with any later surgery in mind.

\medterm{Bone Marrow} Soft tissue inside many bones. Red marrow makes red cells, white cells, and platelets; yellow marrow stores fat and can produce more blood cells when the body needs them.

\medterm{Bone Marrow Aspiration} Removal of a small amount of liquid marrow, usually from the back of the hipbone, with a needle. It is examined to assess blood-cell development, abnormal cells, infection, or genetic changes.

\medterm{Bone Marrow Aspiration And Biopsy} Two ways of sampling marrow, usually from the back of the hipbone. Aspiration collects liquid marrow for cell and genetic testing; biopsy removes a small solid core to assess marrow structure, cell production, scarring, infection, or cancer.

\medterm{Bone Marrow Biopsy} A procedure that removes marrow, usually from the back of the hipbone, for examination. It may include liquid aspiration and a solid core to assess blood-cell production, abnormal cells, scarring, infection, or cancer.

\medterm{Bone Marrow Culture} A laboratory test in which marrow is placed in special conditions to grow and identify bacteria, fungi, or other organisms when an infection of the marrow or bloodstream is suspected.

\medterm{Bone Marrow Depression} Reduced production of blood cells by the marrow. It can cause anaemia, increased infection risk from too few white cells, bleeding from too few platelets, or a combination of these problems.

\medterm{Bone Marrow Disorder} Any disease that affects the marrow's production of red cells, white cells, or platelets. It may cause anaemia, frequent infections, bleeding, abnormal blood counts, or symptoms from abnormal marrow cells.

\medterm{Bone Marrow Donation} The harvesting of hematopoietic stem and progenitor cells (HSPCs) from an allogeneic or autologous donor via surgical iliac crest bone marrow aspiration or apheresis of granulocyte colony-stimulating factor (G-CSF)-mobilized peripheral blood stem cells for hematopoietic stem cell transplantation.

\synonyms
Hematopoietic Stem Cell Donation; Stem Cell Donation; Bone Marrow Harvest

\medterm{Bone Marrow Failure} A condition in which the marrow cannot make enough red cells, white cells, or platelets. Causes include inherited disorders, aplastic anaemia, medicines, toxins, infections, marrow cancers, and other diseases.

\synonyms
Marrow Failure
Bone Marrow Failure Syndrome

\medterm{Bone Marrow Hyperplasia} An increase in the number of marrow cells because blood-cell production has accelerated. It may be a response to ongoing red-cell destruction, blood loss, low oxygen, or another increased demand for blood cells.

\medterm{Bone Marrow Necrosis} Death of the blood-forming tissue inside bones. It can occur with severe sickle-cell disease, blood cancers, or infection and may cause severe bone pain, anaemia, low white-cell or platelet counts, and fever.

\medterm{Bone Marrow Transplantation} A treatment in which healthy blood-forming stem cells are infused to replace or restore bone-marrow cells damaged by disease or treatment. The cells may come from the patient or a donor. It is used for selected blood cancers, marrow failure, and some inherited disorders; risks include infection, graft failure, organ toxicity, and graft-versus-host disease after a donor transplant.

\synonyms
BMT; Bone Marrow Transplant

\medterm{Bone Mass} The amount of bone tissue in the skeleton or in a particular area. Greater bone mass generally provides more strength, while low bone mass can increase the risk of osteoporosis and fractures.

\synonyms
Bone Mineral Content; Skeletal Bone Mass

\medterm{Bone Matrix} The material surrounding bone cells that gives bone its shape and strength. It contains a flexible framework, mainly type I collagen, reinforced with mineral crystals made largely from calcium and phosphate.

\medterm{Bone Metastases} Areas of cancer that have spread to bone from another organ. They may cause bone pain, weakness or fracture, high blood calcium, or pressure on the spinal cord and nerves.

\medterm{Bone Metastasis} A cancer deposit that has spread to bone from another organ. It may cause pain, weakness, fracture, high blood calcium, or spinal-cord pressure; treatment depends on the original cancer, sites involved, and symptoms.

\medterm{Bone Mineral Density} The amount of mineral, mainly calcium, in a measured area of bone. A DXA scan estimates bone strength, helps diagnose osteoporosis, assesses fracture risk, and monitors treatment together with age and other risk factors.

\synonyms
BMD; Bone Density; Bone Mineral Content

\medterm{Bone Mineral Density Test} See \textit{Bone Density Screening}.

\medterm{Bone Necrosis} Death of bone tissue, usually because its blood supply is reduced or because of severe injury, infection, radiation, or prolonged corticosteroid use. It can cause pain, bone collapse, joint damage, and fracture.

\medterm{Bone Neoplasm} An abnormal growth of cells arising in bone. It may be noncancerous or cancerous; imaging and often a biopsy determine its type, behaviour, extent, and appropriate treatment.

\medterm{Bone Pain} Aching or other discomfort arising from bone. Causes include injury, fracture, infection, reduced blood supply, metabolic disease, overuse, or cancer. Persistent, severe, or unexplained pain requires medical assessment.

\medterm{Bone Pain Or Tenderness} Pain when a bone is used, pressed, or touched. It may result from injury, fracture, infection, inflammation, metabolic disease, or cancer; persistent or unexplained tenderness needs medical assessment.

\medterm{Bone Plate} A metal or other body-compatible plate fixed to bone with screws. It holds broken pieces in position, supports healing, or helps correct a deformity during orthopaedic surgery.

\medterm{Bone Regeneration} Formation of new bone to repair or replace bone that is damaged, missing, or diseased. It occurs naturally during fracture healing and may be supported by grafts, implants, or other treatments.

\medterm{Bone Remodeling} The lifelong replacement of old bone with new bone. Cells remove worn bone and other cells build replacement bone, helping maintain strength, shape, and the body's calcium balance.

\medterm{Bone Resorption} The removal and breakdown of bone tissue, which releases minerals into the blood. It is normal during bone renewal, but excessive resorption can weaken bones and contribute to osteoporosis.

\medterm{Bone Sarcoma} Alternate name; see \textit{Bone cancer}.

\medterm{Bone Saw} A surgical instrument used to cut bone during procedures such as amputation, joint replacement, fracture repair, or orthopaedic reconstruction.

\medterm{Bone Scan} A nuclear-medicine test that shows areas of increased or reduced bone activity after a small amount of tracer is given. It can help investigate fractures, infection, cancer spread, and some metabolic bone diseases, but abnormal uptake is not specific.

\medterm{Bone Scan} A nuclear medicine diagnostic imaging technique in which technetium-99m-labeled bisphosphonates (e.g., $^{99\text{m}}\text{Tc}$-MDP) are administered intravenously and imaged with a gamma camera to identify areas of osteoblastic activity, osteomyelitis, occult fractures, or skeletal metastases.

\synonyms
Bone Scintigraphy; Skeletal Scintigraphy; Radionuclide Bone Scan

\medterm{Bone Screws} Screws used during orthopaedic surgery to hold broken bone pieces together or attach plates, joint components, or other implants. Their size, material, and placement depend on the bone and the operation.

\medterm{Bone Spurs} Extra projections of bone that usually form near joints after long-term wear, arthritis, injury, or abnormal stress. They may cause no symptoms or may irritate nearby tissues and restrict movement.

\synonyms
Osteophytes

\medterm{Bone Transplantation} Transfer of bone tissue or a bone substitute to repair a defect, support fracture healing, or assist reconstruction of the spine, jaw, or a joint. The graft may come from the person, a donor, or a manufactured material.

\medterm{Bone Tumor} An abnormal mass that starts in bone. It may be noncancerous or cancerous; examination, scans, and often a biopsy determine its type, whether it is growing or spreading, and the best treatment.

\medterm{Bone Turnover Markers} Substances measured in blood or urine that reflect how quickly bone is being built or broken down. They may support assessment of some metabolic bone disorders or treatment response but do not replace bone-density testing.

\medterm{Bone Type} A way of grouping bones by shape. Long bones support movement, short bones provide stability, flat bones protect organs or provide broad muscle attachment, irregular bones have complex shapes, and sesamoid bones lie within tendons.

\medterm{Bone X-Ray} An X-ray picture of a bone or group of bones used to look for fractures, abnormal shape, arthritis, infection, bone loss, or some tumours.

\medterm{Bone-Marrow Transplantation} A treatment that replaces damaged or diseased blood-forming marrow with healthy stem cells from the patient or a donor. It is used for selected blood cancers, marrow failure, and some inherited disorders.

\medterm{Bones And Joints} Bones provide support and protection, while joints connect bones and allow movement. Joints may be fixed, slightly movable, or freely movable and are held together by tissues such as cartilage, ligaments, and joint fluid.

\medterm{Bones of the Skeleton} The body’s framework of bones and cartilage, providing support, protection, leverage for movement, mineral storage, and marrow-based blood-cell production.
\medterm{Bongkrek Poisoning} Severe poisoning caused by bongkrekic acid in contaminated fermented foods, especially coconut or corn products. It can cause vomiting, abdominal pain, weakness, liver failure, neurologic injury, and death; urgent treatment is required.

\medterm{Bonnevie-Ullrich Syndrome} An older name for the physical features sometimes associated with Turner syndrome, a condition involving an absent or altered X chromosome. Features may include short stature, a webbed neck, and reduced ovarian function.

\medterm{Bony Exostosis} A localized extra growth of bone. It may occur on the jaw or palate and is often harmless, but removal may be considered if it repeatedly causes injury, interferes with speech or dentures, or needs diagnosis.

\medterm{Bony Labyrinth} A system of channels inside the dense part of the skull bone that houses the inner ear. It contains the cochlea for hearing and the vestibule and semicircular canals for balance.

\medterm{Borborygmi} Rumbling or gurgling sounds made when gas and liquid move through the intestines. They are often normal, but major changes together with pain, vomiting, swelling, or inability to pass stool may signal bowel disease.

\medterm{Borborygmus} A single rumbling or gurgling intestinal sound caused by movement of gas and liquid through the bowel. Its significance depends on symptoms and the overall abdominal examination.

\medterm{Border Definition} The sharpness and shape of the edge of a skin lesion or abnormal growth. A poorly defined or irregular edge can be concerning, but the border alone cannot distinguish harmless from cancerous disease.

\medterm{Borderline} Describes a result or condition near the boundary between normal and abnormal or between diagnostic categories. It usually requires interpretation with symptoms, repeat testing, or other findings; the meaning depends on the clinical context.

\medterm{Borderline Personality Disorder} A long-term mental-health condition involving unstable emotions, self-image, relationships, and behavior. It may include intense fear of abandonment, impulsive actions, self-harm, rapidly changing moods, and difficulty managing feelings.

\synonyms
BPD; Emotionally Unstable Personality Disorder (EUPD)

\medterm{Bordetella} A group of small bacteria that can infect the airways of humans or animals. Bordetella pertussis causes whooping cough; other species usually affect animals but can rarely cause opportunistic human infection.

\medterm{Bordetella Avium} A bacterium that mainly causes respiratory disease in turkeys and other birds. Human infection is very rare and is usually reported in people with weakened immunity or serious underlying illness.

\medterm{Bordetella Bronchiseptica} A bacterium that commonly causes respiratory infection in dogs and other animals. It rarely infects humans, but may cause pneumonia or bloodstream infection, especially in people with weakened immunity.

\medterm{Bordetella Hinzii} A bacterium found mainly in poultry. It rarely causes respiratory or bloodstream infection in humans, usually in people with weakened immunity or significant underlying disease.

\medterm{Bordetella Parapertussis} A bacterium that can cause a whooping-cough-like respiratory infection. Illness is often milder than infection with Bordetella pertussis, but prolonged coughing can still occur and the infection can spread to others.

\medterm{Bordetella Pertussis} The bacterium that causes whooping cough, a highly contagious respiratory infection. It attaches to airway lining, produces toxins that disrupt airway function, and can cause repeated coughing, vomiting, breathing pauses, and serious illness in infants.

\medterm{Bordetella Petrii} An environmental bacterium that rarely causes human disease. When it is found in a specimen, clinicians decide whether it represents true infection by considering the body site, symptoms, and other test results.

\medterm{Boric Acid} A chemical used in some industrial products, insecticides, and limited topical preparations. Swallowing it or absorbing a large amount can poison the body; significant exposure requires prompt medical assessment.

\medterm{Boric Acid Poisoning} Illness caused by swallowing or absorbing too much boric acid. It may cause vomiting, diarrhoea, abdominal pain, skin changes, kidney injury, seizures, shock, or death and requires urgent poison or emergency care.

\medterm{Borna Disease} An uncommon viral infection of the nervous system caused by Borna disease virus. It is mainly recognized in animals; human infection is rare and may cause severe inflammation of the brain.

\medterm{Bornholm Disease} A contagious viral illness causing sudden attacks of severe chest or upper-abdominal muscle pain, often with fever, headache, or vomiting. It is usually caused by an enterovirus and generally resolves with supportive care.

\medterm{Boron} A chemical element found naturally in food, water, soil, and industrial materials. Small dietary amounts are generally tolerated, but excessive exposure can cause digestive symptoms, skin problems, reproductive effects, or developmental toxicity.

\medterm{Borrelia} A group of spiral-shaped bacteria spread mainly by ticks or lice. Different species cause Lyme disease, relapsing fever, or other infections, with symptoms and treatment depending on the species and exposure.

\medterm{Borrelia Burgdorferi} A spiral-shaped bacterium spread by infected blacklegged ticks and the main cause of Lyme disease in North America and parts of Europe. Infection may affect the skin, joints, nervous system, or heart.

\medterm{Borrelia Infection} Infection with a Borrelia bacterium. Depending on the species, it may cause Lyme disease or relapsing fever, with symptoms such as fever, headache, fatigue, rash, muscle pain, or joint and nerve problems.

\medterm{Borrelia Miyamotoi} A bacterium spread by some hard-bodied ticks that causes hard-tick relapsing fever. Illness commonly includes repeated episodes of fever, headache, fatigue, chills, and muscle or joint aches.

\medterm{Borreliosis} Infection caused by a Borrelia bacterium. Human forms include Lyme disease and relapsing fever; symptoms differ according to the species, stage of infection, route of spread, and organs involved.

\medterm{Bortezomib} A cancer medicine that blocks proteasomes, cell structures that remove unwanted proteins. This can kill susceptible cancer cells and is used mainly for multiple myeloma and some lymphomas; nerve damage, low blood counts, digestive symptoms, and infections can occur.

\medterm{Bosentan} A medicine that blocks endothelin, a substance that narrows blood vessels. It is used for pulmonary arterial hypertension, but can injure the liver and harm a developing fetus, so liver tests and pregnancy prevention are important.

\medterm{Bosniak Classification} A system for describing kidney cysts on contrast CT or MRI according to their walls, septations, mineral deposits, and solid parts. Categories help estimate cancer risk and guide monitoring, further imaging, or surgery.

\medterm{Bot Fly} A fly whose larvae can live in animal or human tissue, causing myiasis. In people, a larva may produce a painful boil-like skin lump with a small breathing opening; removal is performed by a clinician to avoid complications.

\medterm{Botfly} A fly whose larvae can infest skin or wounds in animals and occasionally people. The infestation is called myiasis and may cause a painful lump, discharge, itching, or tissue irritation.

\medterm{Bothriocephalus} An older name for a group of fish tapeworms. Related tapeworm infections may cause abdominal symptoms and, when prolonged, can reduce vitamin B12 absorption and lead to anaemia.

\medterm{Botryomycosis} A chronic bacterial infection that forms firm granules within lesions and can resemble a fungal infection. It most often affects the skin and tissue beneath it but may involve internal organs, especially in people with weakened immunity.

\medterm{Bottle Feeding} Feeding an infant expressed breast milk or infant formula from a bottle. Safe preparation, clean equipment, correct formula mixing, and responsive feeding help prevent infection, choking, overfeeding, and poor nutrition.

\medterm{Botulinum Toxin} A substance made by Clostridium botulinum that temporarily blocks nerve signals to muscles and some glands. Carefully measured injections treat selected muscle spasms, movement disorders, migraine, sweating, and cosmetic concerns; too much can cause weakness or breathing failure.

\medterm{Botulinum Toxin Injection, Laryngeal} Targeted electromyography- or laryngoscopy-guided injection of botulinum neurotoxin into specific intrinsic laryngeal muscles (such as the thyroarytenoid or posterior cricoarytenoid) to reduce involuntary spasms in focal laryngeal dystonias, predominantly adductor or abductor spasmodic dysphonia.

\synonyms
Laryngeal Botulinum Toxin Injection; Laryngeal Botox; Spasmodic Dysphonia Injection

\medterm{Botulism} A rare, life-threatening illness caused by botulinum toxin, which blocks nerve signals. It causes blurred or double vision, drooping eyelids, difficulty swallowing, progressive weakness, and sometimes breathing failure; urgent antitoxin and intensive care may be needed.

\medterm{Bouchard Nodes} Firm enlargements of the middle joints of the fingers, usually caused by osteoarthritis. They may be painless or may cause stiffness, pain, reduced movement, and finger deformity.

\medterm{Bouchard's Node} A hard enlargement of one middle finger joint, usually caused by osteoarthritis. It results from changes in the joint and surrounding bone and may cause pain, stiffness, or reduced movement.

\medterm{Bouchard's Nodes} Hard enlargements of the middle finger joints caused most often by osteoarthritis. They can limit movement or cause pain and stiffness, although some people have no symptoms.

\synonyms
Proximal Interphalangeal Osteophytes; PIP-Joint Nodes

\medterm{Bougie} A slender medical instrument used to examine or gradually widen a narrowed passage, such as the food pipe or urethra. It must be used carefully because it can cause pain, bleeding, perforation, or infection.

\medterm{Bougienage} Gradual widening or examination of a narrowed body passage with a bougie. It may relieve swallowing or urinary obstruction, but the cause of narrowing must be assessed and the procedure can cause bleeding or perforation.

\medterm{Bouin's Solution} A laboratory liquid containing picric acid, formaldehyde, and acetic acid that preserves tissue for microscopic examination. It is hazardous to handle and can interfere with some later genetic tests.

\medterm{Bounding Pulse} A hyperdynamic, high-amplitude arterial pulse characterized by a rapid, forceful upstroke and wide pulse pressure (Corrigan or water-hammer pulse), classically associated with aortic regurgitation, thyrotoxicosis, severe anemia, fever, patent ductus arteriosus, or septic shock.

\synonyms
Hyperdynamic Pulse; Water-Hammer Pulse; Corrigan Pulse

\medterm{Bourbon Virus} A rare tick-associated viral infection that can cause fever, tiredness, poor appetite, nausea, and low white-cell or platelet counts. No specific treatment or vaccine is established; care is supportive and other serious infections must be excluded.

\medterm{Boutonneuse Fever} A tick-borne infection caused usually by Rickettsia conorii. It commonly causes fever, headache, rash, and a dark sore at the bite site; prompt antibiotic treatment reduces the risk of serious illness.

\medterm{Boutonniere Deformity} A finger shape problem where the middle joint stays bent inward and the tip bends backward, often caused by an injury or rheumatoid arthritis.

\medterm{Bouveret Syndrome} A rare blockage near the stomach outlet caused by a gallstone that has entered the intestine through an abnormal connection from the gallbladder. It can cause persistent vomiting, upper-abdominal pain, bloating, and inability to eat.

\medterm{Bovine} Relating to cattle or to material obtained from cattle. In medicine it may describe an animal disease, a biological product, tissue used in treatment, or an infection that can pass between cattle and people.

\medterm{Bovine Anaplasmosis} A tick- or blood-spread infection of cattle caused by Anaplasma marginale. The bacterium infects red blood cells and destroys them, causing anaemia, weakness, fever, reduced milk production, or death.

\medterm{Bovine Papillomavirus} A group of papillomaviruses that infect cattle and can cause skin warts or other growths. Different strains vary in their effects; human disease is not the usual concern.

\medterm{Bovine Spongiform Encephalopathy} A fatal disease of cattle caused by misfolded proteins called prions. It progressively damages the brain and nervous system; people can rarely develop variant Creutzfeldt--Jakob disease after eating contaminated cattle tissue.

\medterm{Bow Legs} An angular deformational alignment of the lower extremity characterized by outward lateral bowing of the tibia and femur such that the knees remain abnormally separated when the medial malleoli of the ankles are approximated.

\synonyms
Genu Varum; Bowlegs; Bandy Legs

\medterm{Bowel} The small and large intestines. They receive partly digested food, absorb nutrients and water, move contents forward, and form and store stool before it leaves the body.

\medterm{Bowel Burn} Injury to the bowel lining caused by heat, chemicals, radiation, reduced blood flow, or another damaging exposure. It can cause pain, bleeding, perforation, infection, or narrowing and requires assessment according to the cause.

\medterm{Bowel Cancer} Cancer arising in the colon or rectum. It may cause a persistent change in bowel habit, blood in stool, anaemia, weight loss, abdominal pain, or no early symptoms; screening and timely assessment improve detection.

\medterm{Bowel Disorders} Pathological conditions affecting the small or large intestine, encompassing inflammatory bowel disease (Crohn disease, ulcerative colitis), infectious enterocolitis, ischemic bowel disease, bowel obstruction, and colorectal neoplasms.

\synonyms
Intestinal Diseases; Bowel Diseases; Intestinal Disorders

\medterm{Bowel Incontinence} Inability to control the passage of stool or liquid from the anus. It can result from diarrhoea, constipation, muscle or nerve injury, childbirth, bowel disease, or neurologic illness and may often be improved with treatment.

\synonyms
Fecal Incontinence

\medterm{Bowel Infarction} Death of part of the intestine because its blood supply is blocked or severely reduced. It can cause severe abdominal pain, bloody stool, vomiting, and abdominal tenderness and is a life-threatening emergency.

\medterm{Bowel Ischemia} Reduced blood flow to the intestine, causing injury from too little oxygen. It may cause sudden or worsening abdominal pain, vomiting, or bloody stool and can progress to bowel death if circulation is not restored.

\medterm{Bowel Obstruction} A blockage that prevents food, fluid, and gas from moving through the intestine. It can cause cramping pain, swelling, vomiting, and inability to pass stool or gas and may require urgent hospital treatment.

\medterm{Bowel Preparation} Emptying and cleansing the colon before a colonoscopy or other bowel procedure, usually with a special diet and a prescribed laxative solution. Following the instructions is important because poor preparation can hide lesions.

\synonyms
Bowel Prep; Colon Cleansing; Colonic Preparation

\medterm{Bowel Program} A planned routine using food, fluids, medicines, movement, and scheduled toilet visits to produce regular bowel movements and manage constipation or loss of bowel control.

\medterm{Bowel Retraining} A structured routine that uses regular toilet times, suitable food and fluids, physical activity, and helpful posture to establish predictable bowel movements and improve constipation or bowel-control problems.

\medterm{Bowel Sounds} Gurgling or clicking sounds made as the intestines move gas and liquid. Their presence and pattern can provide clues during an abdominal examination but cannot diagnose a disorder by themselves.

\medterm{Bowel Transit Time} The time food and other contents take to travel through the digestive tract. Testing it can help investigate unusually slow or fast movement when a bowel-movement or digestive disorder is suspected.

\medterm{Bowel Urgency} A sudden, hard-to-delay need to pass stool. It can occur with diarrhoea, infection, irritable bowel syndrome, inflammatory bowel disease, or disorders of the rectum and may lead to accidents.

\medterm{Bowen's Disease} An early form of skin cancer confined to the outermost skin layer. It usually appears as a slowly enlarging, persistent scaly red or brown patch and can invade deeper tissue if untreated.

\medterm{Bowenoid Papulosis} A condition linked to high-risk human papillomavirus that causes multiple reddish-brown or purple bumps on genital or nearby skin. The cells look cancerous under a microscope, but the lesions often behave less aggressively; treatment and follow-up reduce the risk of progression.

\medterm{Bowlegs} Alternate name; see \textit{Bow Legs}.

\medterm{Bowler's Thumb} Irritation or scarring of the digital nerve along the thumb caused by repeated pressure from a bowling ball. It may produce pain, numbness, tingling, or a tender lump and may improve with rest or equipment changes.

\medterm{Bowman Capsule} A cup-shaped structure at the beginning of each kidney filtering unit. It surrounds a knot of tiny blood vessels and collects the fluid filtered from blood before it passes through the rest of the kidney tubule.

\medterm{Bowman Membrane} A thin, strong layer near the front of the cornea that supports and protects the eye's clear surface. Significant injury can scar it because it does not normally grow back.

\medterm{Bowman's Capsule} The cup-shaped beginning of a kidney filtering unit. It surrounds a cluster of small blood vessels and receives the fluid filtered from blood.

\medterm{Bowman's Gland} A small gland high inside the nose that releases watery fluid onto the smell-sensing lining. The fluid helps dissolve scent molecules and clear the surface so odours can be detected.

\medterm{Boxed Warning} The strongest safety warning required in a medicine's prescribing information when there is a serious or life-threatening risk. It explains a major hazard but does not necessarily mean the medicine is unsuitable for every patient.

\medterm{Boxer Fracture} A transverse or comminuted fracture through the distal metaphysis or neck of the fifth metacarpal bone (and occasionally the fourth metacarpal), characteristically accompanied by volar angulation of the metacarpal head resulting from direct axial loading with a clenched fist.

\medterm{Boxer's Ear} A deformity of the outer ear caused by repeated injury or untreated blood collection around the ear cartilage. It can produce thickened, scarred tissue with a cauliflower-like appearance.

\medterm{Boxer's Fracture} A break near the knuckle end of the fifth metacarpal, the hand bone below the little finger. It commonly follows punching and causes pain, swelling, tenderness, or a changed knuckle shape.

\medterm{Braces and Retainers} Alternate name; see \textit{Dental Braces}.

\medterm{Brachial} Relating to the upper arm, including its bones, muscles, blood vessels, and nerves. The brachial artery supplies the arm, and the brachial plexus carries nerve signals to and from the arm and hand.

\medterm{Brachial Artery} The main artery of the upper arm. It carries blood from the armpit toward the elbow, where it divides into the radial and ulnar arteries, and it is commonly used to measure a pulse or blood pressure.

\medterm{Brachial Plexopathy} Damage or dysfunction of the nerve network supplying the shoulder, arm, and hand. It can cause pain, weakness, numbness, or paralysis after injury, pressure, inflammation, a tumour, or radiation.

\medterm{Brachial Plexus} A network of nerves from the lower neck and upper chest that supplies movement and sensation to the shoulder, arm, forearm, and hand. Its branches include the musculocutaneous, axillary, radial, median, and ulnar nerves.

\medterm{Brachial Plexus Cords} Three parts of the nerve network supplying the arm, named lateral, posterior, and medial because of their position around an artery in the armpit. They join to form nerves controlling movement and feeling in the upper limb.

\medterm{Brachial Plexus Injuries} Damage to the nerve network supplying the shoulder and upper limb, caused by stretching, pressure, or a penetrating injury. Symptoms range from pain or numbness to weakness or paralysis, depending on the nerves affected.

\medterm{Brachial Plexus Injury} Damage to the nerve network of the shoulder and arm, sometimes occurring during a difficult birth or after trauma. It may cause weakness, loss of movement, numbness, or paralysis; recovery depends on the severity and nerves involved.

\medterm{Brachial Plexus Injury In Newborns} Damage to the nerves supplying a newborn's shoulder, arm, and hand, usually after a difficult delivery. The affected arm may be weak or still, and recovery ranges from spontaneous improvement to requiring specialist treatment.

\medterm{Brachial Plexus Neuropathies} Disorders affecting the nerve network that supplies the shoulder, arm, and hand. They can cause pain, weakness, numbness, or loss of movement because of injury, pressure, inflammation, tumours, or radiation.

\medterm{Brachial Plexus Roots} The first portions of the nerve network that supplies the arm, arising from the lower neck and upper chest. Damage to different roots produces characteristic patterns of weakness, numbness, and pain.

\medterm{Brachial Plexus Trunks} Three sections of the brachial plexus—upper, middle, and lower—that carry nerve fibres toward the arm. Injury to the upper trunk can cause weakness of the shoulder and elbow, while lower-trunk injury can mainly affect the hand.

\medterm{Brachial Vein} Deep veins accompanying the brachial artery that drain the arm and join the basilic vein to form the axillary vein.

\medterm{Brachialgia} Pain that travels from the neck or shoulder down the arm. It often results from irritation of a neck nerve or the brachial plexus and may occur with numbness, tingling, or weakness.

\medterm{Brachialis Muscle} A strong muscle at the front of the upper arm, beneath the biceps. It is the main muscle that bends the elbow and works whether the palm faces up, down, or sideways.

\medterm{Brachiocephalic Trunk} The first large branch of the aorta in the chest. It carries blood toward the right side of the head, neck, and arm and divides into the right common carotid and right subclavian arteries.

\medterm{Brachiocephalic Veins} Two large veins formed by the joining of veins from the head, neck, and arm. They drain blood from these areas into the superior vena cava, which returns it to the heart.

\medterm{Brachioplasty} An operation to remove or reshape excess skin and tissue of the upper arm, sometimes called an arm lift. Possible complications include scars, altered sensation, bleeding, fluid collection, infection, and wound problems.

\medterm{Brachioradialis Reflex} A deep tendon stretch reflex elicited by briskly tapping the tendon of the brachioradialis muscle near the styloid process of the radius with a reflex hammer, producing forearm flexion and slight pronation. It is innervated by the radial nerve and tests the integrity of the C5 and C6 cervical spinal cord segments and nerve roots.

\synonyms
Supinator Reflex; Radial Reflex

\medterm{Brachium} The part of the upper limb between the shoulder and elbow. It contains the upper-arm bone, muscles that move the shoulder and elbow, and major nerves and blood vessels.

\medterm{Brachycephalic} Describes a skull that is relatively short from front to back and broad from side to side. It may be a normal family characteristic or part of a condition affecting skull growth.

\medterm{Brachycephaly} A broad, short skull shape. It may develop from pressure on an infant's head or from early joining of skull bones; assessment determines whether brain growth or treatment is affected.

\medterm{Brachydactyly} A birth condition in which one or more fingers or toes are shorter than usual because of altered development of their bones or surrounding tissues. It may occur alone or as part of an inherited syndrome.

\medterm{Brachytherapy} Internal radiation treatment in which a radioactive source is placed inside or close to a tumour. It delivers a concentrated dose to a limited area and may use a temporary applicator or a source left in place permanently.

\synonyms
Brachy Therapy
Internal Radiation Therapy

\medterm{Brachytherapy Seed} A tiny radioactive source placed in or near a tumour to deliver local radiation. Seeds may be temporary or permanent and are used most often for selected cancers, including prostate cancer.

\medterm{Bradycardia} A slower-than-usual heart rate, often below 60 beats per minute in adults. It can be normal during sleep or in trained athletes, but may result from heart-conduction disease or medicines and cause dizziness, fainting, tiredness, or breathlessness.

\synonyms
Slow Heartbeat

\medterm{Bradycardia Of Prematurity} A slow heart rate in a premature infant, often occurring with pauses in breathing and low oxygen. It commonly reflects immature breathing and heart control and is monitored and treated according to its frequency and severity.

\medterm{Bradycardia-Tachycardia Syndrome} Alternate name; see \textit{Sick Sinus Syndrome}.

\medterm{Bradykinesia} Slowness and reduced size of voluntary movements. It can make starting or repeating actions difficult and may occur with stiffness, tremor, reduced arm swing, a less expressive face, or a shuffling walk in Parkinson disease and related disorders.

\medterm{Bradykinesia and Athetosis} Hypokinetic and dyskinetic movement patterns characterized by generalized slowness of movement initiation (bradykinesia) or continuous slow, writhing, twisting involuntary movements (athetosis), frequently observed in Parkinson disease, secondary parkinsonism, or basal ganglia encephalopathies.

\synonyms
Slowness of Movement; Hypokinesia and Athetosis; Extrapyramidal Slowness
\medterm{Bradykinetic} Describes abnormally slow movement, especially the reduced speed and size of movement seen in Parkinson disease and related neurologic disorders.

\medterm{Bradykinin} A natural chemical messenger that widens some blood vessels and allows more fluid to leave them. It contributes to inflammation, swelling, pain, and rare reactions to certain blood-pressure medicines.

\medterm{Bradyphrenia} Unusually slow thinking or mental processing. It may cause delayed responses and difficulty completing tasks and can occur with neurologic disease, mental illness, medicines, depression, or other conditions.

\medterm{Bradypnea} Abnormally slow breathing, often fewer than 12 breaths per minute in an adult. It may result from sedative or opioid effects, brain disease, low body temperature, hormone disorders, or severe illness and can lead to respiratory failure.

\medterm{Bradypnoea} Abnormally slow breathing. Causes include medicines, neurologic disease, metabolic disturbance, low body temperature, or severe illness; breathing that is difficult, shallow, or associated with confusion requires urgent care.

\medterm{Braille System} A raised-dot writing and reading system used by people who are blind or have low vision. Patterns represent letters, numbers, punctuation, and symbols and can be read by touch.

\medterm{Brain} The main organ of the nervous system, protected inside the skull. It processes senses, controls movement and body functions, and supports thinking, memory, emotions, language, and consciousness.

\medterm{Brain Abscess} A pocket of pus and infection inside the brain. It can develop when infection spreads from the blood, nearby ear or sinus disease, surgery, or an injury and may cause headache, fever, seizures, or weakness.

\medterm{Brain Aneurysm} A weakened bulge in an artery supplying the brain. If it bursts, it can cause sudden severe headache, vomiting, neck stiffness, collapse, or other neurological problems.

\medterm{Brain Aneurysm Repair} Treatment to prevent a brain aneurysm from bleeding or bleeding again. A surgeon may place a clip or a specialist may seal the aneurysm from inside the artery with coils or another device.

\medterm{Brain Arteriovenous Malformation} A congenital vascular dysplasia consisting of a direct high-flow, low-resistance communication between cerebral arteries and veins through an abnormal vascular nidus without an intervening capillary bed, carrying a risk of intracranial hemorrhage, focal seizures, and neurological deficits.

\synonyms
Cerebral Arteriovenous Malformation; Brain AVM; Intracranial AVM

\medterm{Brain AVM} Alternate name; see \textit{Brain Arteriovenous Malformation}.

\medterm{Brain Cancer} Cancer affecting brain tissue or its coverings. Symptoms depend on the location and may include seizures, headache, weakness, speech or vision changes, personality change, or pressure-related vomiting; diagnosis usually requires brain imaging and sometimes tissue testing.

\synonyms
Malignant Brain Tumor; Primary Brain Neoplasm

\medterm{Brain Concussion} Alternate name; see \textit{Concussion}.

\medterm{Brain Death} Permanent loss of all brain and brainstem function, including consciousness, reflexes, and the ability to breathe independently. When formally confirmed according to law and medical standards, it means the person has died.

\medterm{Brain Disorder} Any illness or abnormality affecting the brain's structure or function, including stroke, infection, tumour, epilepsy, injury, degenerative disease, or a condition present from birth.

\medterm{Brain Edema} Swelling caused by extra fluid in brain tissue. It can increase pressure inside the skull after injury, stroke, infection, bleeding, tumour, or severe illness and may cause headache, vomiting, seizures, confusion, or reduced consciousness.

\medterm{Brain Embolism} Alternate name; see \textit{Embolic Stroke}.

\medterm{Brain Fog} An informal description of difficulty concentrating, thinking clearly, remembering, or finding words. It is a symptom rather than a diagnosis and may occur with poor sleep, stress, illness, medicines, depression, or other conditions.

\medterm{Brain Fornix} A curved bundle of nerve fibres deep in the brain that helps connect the hippocampus with other areas involved in memory. Damage may cause difficulty learning or remembering information.

\medterm{Brain Haemorrhage} Bleeding within or around the brain. It may result from high blood pressure, injury, an aneurysm or other blood-vessel problem, clotting disorders, or medicines and can cause sudden neurologic symptoms.

\medterm{Brain Health} The ability of the brain to support thinking, memory, learning, movement, emotions, communication, and consciousness. It is supported by controlling vascular risks, sleeping well, physical activity, balanced nutrition, social connection, and treatment of illness.

\medterm{Brain Hemorrhage} Bleeding within or around the brain, often caused by high blood pressure, an aneurysm, a blood-vessel malformation, injury, or blood-thinning treatment. Sudden severe headache, weakness, speech trouble, vomiting, seizure, or reduced consciousness requires emergency care.

\medterm{Brain Herniation} Dangerous displacement of brain tissue caused by severe swelling, bleeding, a tumour, or another increase in pressure inside the skull. It can compress the brainstem, impair breathing and consciousness, and requires emergency treatment.

\medterm{Brain Imaging} Tests that create pictures of brain structure or activity. CT and MRI show tissue, bleeding, swelling, tumours, and strokes; blood-vessel or nuclear scans answer more specific questions.

\synonyms
Neuroimaging; Cerebral Imaging

\medterm{Brain Injuries} Damage to the brain ranging from a concussion to severe injury, commonly caused by falls, crashes, sports impacts, or penetrating wounds. Effects depend on the area and severity and may include headache, confusion, memory loss, seizures, or unconsciousness.

\medterm{Brain Injury} Damage or impaired function of the brain caused by trauma, lack of oxygen, stroke, infection, toxins, or another disease. Effects vary with the location and severity and may affect movement, sensation, thinking, behaviour, or consciousness.

\medterm{Brain Injury} Alternate name; see \textit{Traumatic Brain Injury}.

\medterm{Brain Lesions} Areas of brain tissue that look or function abnormally because of a tumour, stroke, infection, inflammation, injury, loss of nerve insulation, or a blood-vessel disorder. Symptoms depend mainly on the lesion's location and size.

\medterm{Brain Lesions} Alternate name; see \textit{Brain Lesion}.

\medterm{Brain Malformations} Structural differences in the brain that develop before birth. They may result from genetic changes, disrupted fetal development, infection, or other prenatal influences and can affect movement, learning, seizures, or other neurologic functions.

\medterm{Brain Mapping} Methods used to identify where brain functions and structures are located. Imaging, electrical recordings, or direct stimulation may guide research or help surgeons avoid important areas during selected operations.

\medterm{Brain Metastases} Cancer deposits that have spread to the brain from another organ. They may cause headache, seizures, weakness, speech or vision changes, confusion, or no symptoms and are assessed with brain imaging.

\synonyms
Secondary Brain Cancer

\medterm{Brain Metastasis} A cancer deposit that has spread to the brain from another organ. Symptoms may include headache, seizures, weakness, speech or vision changes; treatment depends on the number, location, original cancer, and overall health.

\medterm{Brain Monitoring} Recording brain activity or pressure to assess neurologic function, seizures, sleep, or the depth of anaesthesia. The method used depends on the clinical question; a monitor result is interpreted with examination findings and other tests.

\medterm{Brain Natriuretic Peptide Test} A blood test measuring BNP or NT-proBNP, substances released when the heart is stretched. A raised result can support a diagnosis of heart failure, but it must be interpreted with symptoms, examination, age, kidney function, and other findings.

\medterm{Brain Neoplasm} An abnormal growth of cells in the brain or its coverings. It may be noncancerous or cancerous; imaging and sometimes biopsy determine its type, behaviour, extent, and treatment options.

\medterm{Brain PET Scan} A scan that uses a small amount of radioactive tracer to show chemical activity in the brain. It can help investigate selected tumours, memory disorders, seizures, or brain function when other scans do not answer the question.

\medterm{Brain Sand} Brain sand is an older name for corpora arenacea, tiny calcified deposits commonly found in the pineal gland and sometimes other brain regions.

\medterm{Brain SPECT} A nuclear-imaging scan that shows blood flow or selected brain activity using a radioactive tracer. It may help evaluate seizures, stroke-related blood-flow changes, dementia, or other specialist questions when interpreted with other findings.

\medterm{Brain Stem} The lower part of the brain connecting the cerebrum with the spinal cord and controlling essential functions such as breathing and consciousness.

\medterm{Brain Stem Glioma} A tumor arising from glial cells in the brainstem, the part of the brain controlling important functions such as breathing, swallowing, eye movement, and balance. Symptoms and treatment depend on its exact location and behavior.

\medterm{Brain Stimulation} Use of electrical, magnetic, or other energy to change brain activity. Methods include electroconvulsive therapy, transcranial magnetic stimulation, and deep-brain stimulation; their uses, benefits, and risks differ.

\medterm{Brain Surgery} An operation on the brain or nearby structures to remove or sample a tumor, relieve pressure, treat bleeding, repair an abnormal blood vessel, or control selected seizures. Risks depend on the area and procedure and can include bleeding, infection, or neurologic loss.

\medterm{Brain Tumor} An abnormal growth of cells in the brain or nearby structures. It may be noncancerous or cancerous, begin in the brain, or spread there from cancer elsewhere. It may cause seizures, headache, vomiting, weakness, changes in thinking, or no symptoms at first.

\medterm{Brain Tumor, Adult Primary} A neoplasm arising de novo within the brain parenchyma, cranial nerves, or meninges in adults, encompassing diffuse gliomas (astrocytoma, oligodendroglioma, glioblastoma), meningiomas, vestibular schwannomas, and primary central nervous system lymphomas.

\synonyms
Adult Primary Brain Tumor; Primary Intracranial Neoplasm; Adult Glioma

\medterm{Brain Tumor, Pediatric} A primary neoplasm originating within the parenchyma of the brain or surrounding meninges and cranial nerves during childhood or adolescence. Common histologic types include pilocytic astrocytoma, medulloblastoma, ependymoma, and diffuse intrinsic pontine glioma.

\synonyms
Pediatric Brain Tumor; Childhood Brain Tumor; Pediatric CNS Neoplasm

\medterm{Brain Tumour} An abnormal growth in or near the brain that may be noncancerous or cancerous. It can cause seizures, headache, vomiting, changes in thinking, weakness, or other problems by affecting nearby brain tissue.

\medterm{Brain Ventricles} Four connected spaces inside the brain that contain and circulate cerebrospinal fluid. This fluid cushions the brain and spinal cord, carries substances, and helps remove waste; blockage can cause pressure and hydrocephalus.

\medterm{Brain Wave} A pattern of electrical activity made by groups of nerve cells and recorded with an electroencephalogram. Patterns change with sleep, alertness, seizures, and some disorders of brain function.

\medterm{Brain Waves} Patterns of electrical activity produced by groups of nerve cells and recorded with an electroencephalogram. Their timing and shape can help assess sleep, seizures, brain function, and some disorders affecting consciousness.

\synonyms
Cerebral Electrical Activity; EEG Rhythms

\medterm{Brainstem} The lower part of the brain connecting the cerebrum with the spinal cord. It helps control breathing, heart rate, blood pressure, swallowing, eye movements, alertness, and pathways carrying signals between the brain and body.

\medterm{Brainstem Auditory Evoked Response} An electrodiagnostic test that records electrical potentials generated by the auditory nerve and brainstem pathways in response to acoustic click or tone stimuli. It is utilized to assess hearing acuity in newborns and uncooperative patients, and to evaluate auditory pathway lesions such as acoustic neuromas.

\synonyms
Brainstem Auditory Evoked Potential; BAEP; Brainstem Auditory Evoked Response; BAER; Auditory Brainstem Response; ABR

\medterm{Brainstem Tumor} A growth arising in the brainstem, which connects the brain with the spinal cord and controls important functions. Symptoms may include problems with eye movement, balance, swallowing, speech, strength, or breathing.

\medterm{Branchial} Relating to structures that develop from the paired tissue arches in an embryo's early head and neck. These arches contribute to parts of the face, throat, voice box, muscles, bones, blood vessels, and nerves.

\medterm{Branchial Arch} One of several paired structures in an embryo that develop into parts of the face, throat, voice box, muscles, bones, blood vessels, and nerves of the head and neck.

\medterm{Branchial Cleft} A groove between developing head-and-neck arches in an embryo. If part of it remains after birth, it may form a cyst, a narrow passage, or an opening, usually on the side of the neck.

\medterm{Branchial Cleft Anomaly} A birth-related remnant of developing head-and-neck tissue that may form a cyst, narrow passage, or skin opening. Repeated infection, swelling, or drainage from the side of the neck needs medical assessment.

\medterm{Branchial Cleft Cyst} A fluid-filled lump on the side of the neck caused by tissue left over from development before birth. It may enlarge or become painful after infection and occasionally causes drainage, swallowing difficulty, or pressure on the airway.

\medterm{Branchial Cleft Cyst} Alternate name; see \textit{Branchial Cyst}.

\medterm{Branchial Cyst} A benign, congenital lateral neck cyst that arises from the failure of obliteration of embryonic branchial cleft tissue, most commonly the second branchial cleft (accounting for over 90\% to 95\% of branchial anomalies). It typically presents in adolescents or young adults as a smooth, painless, fluctuant, non-transilluminating mass located along the anterior border of the upper third of the sternocleidomastoid muscle near the angle of the mandible, which often enlarges or becomes tender following an upper respiratory infection. Fine-needle aspiration yields characteristic turbid, straw-colored fluid containing cholesterol crystals. In adults over 40, lateral cystic neck masses require careful evaluation to rule out cystic nodal metastasis from occult oropharyngeal carcinoma. Definitive treatment is complete surgical excision of the cyst and its associated tract.

\synonyms
Branchial Cleft Cyst; Second Branchial Cleft Cyst; Lateral Cervical Cyst

\medterm{Branchio-Oto-Renal Syndrome} An inherited disorder that can cause neck cysts or openings, hearing loss, and kidney abnormalities. The combination and severity vary, and some affected people have only one or two of these features.

\medterm{Branchioma} A rare cyst or growth arising from tissue related to early head-and-neck development. It often presents as a lump in the neck and requires examination and imaging to establish its nature.

\medterm{Brass-Founder's Ague} An older name for metal-fume fever, a short-lived flu-like illness after breathing fumes from heated metals such as zinc or brass. Symptoms may include fever, chills, headache, muscle aches, cough, and nausea.

\medterm{Brassicasterol} A plant sterol found especially in rapeseed and related oils. It is mainly used to identify the composition of plant oils and has no routine role as a medical treatment.

\medterm{Brassidic Acid} A long-chain fatty acid found in some plant oils. It is mainly a chemistry term and is not an established medicine or routine clinical test.

\medterm{Braxton Hicks Contractions} Irregular tightening of the uterus during pregnancy, often becoming more noticeable later in pregnancy. They usually do not steadily strengthen or open the cervix and may ease with rest, fluids, or a change of activity.

\synonyms
False Labor; Practice Contractions

\medterm{BRCA} BRCA1 and BRCA2 are genes that help repair damaged DNA. An inherited harmful change in either gene increases the risk of breast, ovarian, fallopian-tube, primary peritoneal, prostate, and pancreatic cancers.

\medterm{BRCA1 And BRCA2} Genes that help repair breaks in DNA and limit abnormal cell growth. An inherited harmful variant raises the risk of several cancers and may affect screening, preventive surgery, and treatment choices for the person and relatives.

\medterm{BRCA1 And BRCA2 Gene Testing} A blood or saliva test looking for harmful inherited changes in BRCA1 or BRCA2. Results can guide cancer screening, preventive options, and treatment and may provide important information for biological relatives.

\medterm{BRCA1 Gene} A gene involved in repairing damaged DNA and limiting tumor formation. An inherited harmful variant increases the risk of breast, ovarian, fallopian-tube, peritoneal, prostate, and pancreatic cancers; genetic counseling helps interpret test results.

\medterm{BRCA2 Gene} A gene that helps repair serious DNA damage. An inherited harmful variant increases the risk of breast, ovarian, prostate, pancreatic, and some other cancers and may affect relatives; genetic counseling helps explain testing and options.

\medterm{Breakbone Fever} An older name for dengue fever, reflecting the severe muscle and joint pain that may occur with this mosquito-borne viral infection. Fever, headache, rash, nausea, or dangerous bleeding may also develop.

\medterm{Breakthrough Infection} An infection that occurs despite vaccination or another preventive measure intended to reduce its risk. Protection and illness severity vary with the germ, intervention, timing, and person's immune system.

\medterm{Breakthrough Pain} A short flare of pain that occurs despite otherwise controlled persistent pain. It may happen unexpectedly or after activity; treatment depends on the cause, pattern, regular pain medicines, and safety risks.

\medterm{Breast} One of the two organs on the front of the chest, made of milk-producing tissue, ducts, fat, connective tissue, blood vessels, nerves, and skin. It can be affected by benign conditions and cancer.

\medterm{Breast Abscess} A pocket of pus in breast tissue, often developing after mastitis during breastfeeding. It may cause a painful lump, redness, warmth, and fever; ultrasound can confirm it, and drainage plus antibiotics may be needed.

\medterm{Breast Anatomy} The breast contains skin, fat, milk-producing lobules, ducts, connective tissue, blood vessels, nerves, and lymph vessels that extend toward the armpit. Its appearance changes with age, hormones, pregnancy, and breastfeeding.

\medterm{Breast And Thoracic Cancer Surgery} Operations for cancers of the breast or chest, including selected cancers of the lung, pleura, food pipe, mediastinum, or chest wall. Surgery may remove a tumour and nearby lymph nodes, with further treatment based on the cancer type and stage.

\synonyms
Breast and Thoracic Oncologic Surgery

\medterm{Breast Augmentation} A procedure to increase breast size or change its shape, usually with an implant or transferred fat. Risks include infection, bleeding, scarring, altered sensation, implant problems, and the need for further surgery.

\synonyms
Breast Enlargement; Augmentation Mammoplasty

\medterm{Breast Augmentation Surgery} An operation that enlarges or reshapes the breasts using implants or transferred tissue. The surgeon discusses expected results, alternatives, risks, recovery, and possible later procedures before treatment.

\medterm{Breast Biopsy, Stereotactic} A minimally invasive percutaneous core-needle biopsy procedure utilizing digital mammographic triangulation from angled X-ray views to precisely target and sample nonpalpable microcalcifications or architectural distortions.

\synonyms
Stereotactic Breast Biopsy; Stereotactic Core Needle Biopsy; Mammographically Guided Breast Biopsy

\medterm{Breast Biopsy, Ultrasound-Guided} A percutaneous core-needle or vacuum-assisted tissue sampling procedure performed under continuous real-time sonographic guidance to biopsy solid masses, architectural abnormalities, or complex cystic breast lesions.

\synonyms
Ultrasound-Guided Breast Biopsy; Sonographically Guided Core Biopsy; Ultrasound Breast Biopsy

\medterm{Breast Bone} A common name for the breastbone, or sternum, a flat bone in the centre of the chest. It joins the ribs and helps protect the heart and lungs.

\medterm{Breast Calcifications} Small calcium deposits seen in breast tissue on a mammogram. Most are harmless, but their size, shape, and arrangement may require additional images or biopsy to exclude an early cancer.

\medterm{Breast Cancer} Cancer that begins in breast tissue. It may appear as a lump, skin or nipple change, or abnormal imaging, although early cancer may cause no symptoms; treatment depends on its type, stage, and biology.

\medterm{Breast Cancer And Lymphedema} Breast-cancer surgery or radiation can damage lymph drainage and cause lasting swelling, heaviness, tightness, or skin changes in the arm, breast, or chest. Early assessment and specialist treatment can reduce progression and complications.

\medterm{Breast Cancer Genetics} The study of inherited gene changes that raise breast-cancer risk. Important genes include BRCA1, BRCA2, PALB2, TP53, PTEN, and CHEK2; genetic counselling helps decide whether testing and extra screening are appropriate.

\medterm{Breast Cancer Prevention} Ways to lower, but not eliminate, breast-cancer risk include physical activity, limiting alcohol, avoiding tobacco, appropriate screening, and discussing inherited risk. Preventive medicine or surgery may be considered when risk is substantially high.

\medterm{Breast Cancer Recurrence} Return of breast cancer after treatment, either near the original site, in nearby lymph nodes, or in distant organs. New persistent symptoms or a lump require assessment, although many such symptoms have noncancerous causes.

\medterm{Breast Cancer Screening} Testing to find breast cancer before symptoms appear, usually with mammography. Ultrasound or MRI may be added for selected people according to age, breast density, inherited risk, and other factors.

\medterm{Breast Cancer Staging} Assessment of how far breast cancer has grown or spread, including the tumour's features, nearby lymph nodes, and distant organs. The stage helps estimate outlook and choose treatment.

\medterm{Breast Cancer, Inflammatory} Alternate name; see \textit{Inflammatory Breast Cancer}.

\medterm{Breast Cancer, Male} Alternate name; see \textit{Male Breast Cancer}.

\medterm{Breast Carcinoma} A cancer that begins in breast tissue, usually in the ducts or milk-producing lobules. Its treatment depends on the cell type, hormone and other markers, tumour size, lymph nodes, and spread.

\medterm{Breast Cyst} A benign, fluid-filled sac within breast tissue, most common in women between the ages of 35 and 50. It typically presents as a smooth, round, movable lump that can fluctuate in size and tenderness with the menstrual cycle. Ultrasound imaging distinguishes fluid-filled cysts from solid masses. Simple cysts are noncancerous and do not increase the risk of breast cancer; they generally require no intervention unless fluid is drained with a needle to relieve discomfort.

\medterm{Breast Cysts} Alternate name; see \textit{Breast Cyst}.

\medterm{Breast Cytology} Examination of individual cells from a breast lump, nipple discharge, or needle sample. It can suggest infection, a harmless change, or cancer, but a core biopsy is often needed to determine the full tissue pattern.

\medterm{Breast Density} The relative amount of milk-gland and connective tissue compared with fat on a mammogram. Dense breast tissue can hide small cancers and is linked with higher breast-cancer risk; screening plans depend on individual risk and local guidance.

\medterm{Breast Development} Growth and maturation of breast tissue, beginning mainly at puberty and changing with hormones, pregnancy, breastfeeding, ageing, and some medicines. Uneven development is common, but a new persistent change requires assessment.

\medterm{Breast Duct} A channel that carries milk from milk-producing areas toward the nipple. Blockage, inflammation, or abnormal cells in a duct can cause a lump, pain, or nipple discharge.

\medterm{Breast Engorgement} Overfilling and swelling of the breasts with milk and extra tissue fluid, usually a few days after birth. Both breasts may feel firm, warm, painful, and make the nipples harder for an infant to grasp.

\medterm{Breast Enlargement In Males} Enlargement of glandular breast tissue in a male, called gynecomastia. It may result from normal hormonal changes, medicines, obesity, liver or kidney disease, or other hormone disorders; a hard one-sided lump needs assessment.

\medterm{Breast Examination} Inspection and gentle examination of the breasts, nipples, and nearby lymph nodes to assess lumps, pain, skin or nipple changes, and other findings. It may be followed by imaging or biopsy when needed.

\medterm{Breast Feed} To feed an infant directly from the breast with human milk. Human milk provides nutrients and immune protection, but feeding may need support when the parent or infant has illness, difficulty latching, or medicine-related concerns.

\medterm{Breast Finding} Any noticeable breast change, such as a lump, pain, skin change, nipple inversion, or discharge. Its significance depends on the examination and imaging and sometimes requires a biopsy.

\medterm{Breast Health} Care and awareness aimed at recognising breast changes and detecting disease early. It includes awareness of what is normal, screening appropriate for age and risk, and prompt reporting of a new lump, discharge, or persistent skin change.

\medterm{Breast Imaging-Reporting And Data System} A standard system for describing mammogram, breast ultrasound, and breast MRI findings. Its numbered categories indicate whether results are normal, benign, probably benign, suspicious, or already known to be cancer and help guide follow-up.

\medterm{Breast Infection} Infection and inflammation of breast tissue, often occurring during breastfeeding. It can cause local pain, warmth, redness, swelling, and fever; worsening symptoms, a pus-filled lump, or feeling very unwell needs medical care because an abscess may develop.

\medterm{Breast Lift} Surgery that raises and reshapes sagging breast tissue by removing extra skin and, when needed, repositioning the nipple. It changes shape rather than necessarily increasing breast size and can leave scars or alter sensation.

\medterm{Breast Lobe} One of the larger sections of a breast's milk-producing gland. It contains smaller lobules and ducts that carry milk toward the nipple.

\medterm{Breast Lobule} A small milk-producing unit within the breast, made of tiny milk glands and ducts. Lobules enlarge and become more active during pregnancy and breastfeeding.

\medterm{Breast Lump} A swelling, mass, or thickened area in breast tissue. Most breast lumps (around 80\% to 90\%) are harmless and noncancerous, most often caused by fluid sacs (cysts), noncancerous growths (fibroadenomas), hormone changes, or past injury. Cancerous lumps are less common and usually feel hard, uneven, and fixed in place. Doctors identify the cause through physical examination, imaging such as ultrasound or mammograms, and a needle test (biopsy) if needed.

\medterm{Breast Lump Removal} An operation to remove a breast lump, usually so it can be examined or because it causes symptoms or concern. The need for surgery depends on examination, imaging, biopsy results, and the person's preferences.

\medterm{Breast Lumps} Alternate name; see \textit{Breast Lump}.

\medterm{Breast Milk} Milk made by the breasts after childbirth. It provides energy, nutrients, fluid, and immune substances that support an infant's growth and help protect against some infections.

\medterm{Breast Milk Expression and Storage} The manual or mechanical extraction of human breast milk and its subsequent storage following strict hygiene, refrigeration, and freezing protocols to preserve bioactive immunoglobulins, micronutrients, and enzymatic properties while preventing microbial contamination.

\synonyms
Pumping and Storing Breast Milk; Breast Milk Storage; Milk Expression

\medterm{Breast Milk Jaundice} Yellowing of the skin and eyes that begins after the first week in an otherwise well, feeding, and growing breastfed infant. It is usually temporary, but every newborn with jaundice needs assessment to exclude dangerous causes; feeding is generally continued unless clinicians advise otherwise.

\synonyms
Human Milk Jaundice; Late-Onset Neonatal Jaundice

\medterm{Breast MRI} An imaging test using a strong magnet and radio waves to create detailed pictures of breast tissue. Contrast is often used; MRI is added for selected screening, implant, or diagnostic questions rather than replacing mammography for everyone.

\medterm{Breast MRI Scan} A scan that uses magnetic fields and radio waves, sometimes with injected contrast, to show breast tissue in detail. It can help assess high cancer risk, an abnormality, implants, or the extent of known cancer.

\medterm{Breast Neoplasm} An abnormal growth of cells in breast tissue. It may be noncancerous or cancerous; imaging and often a biopsy determine its type, behaviour, extent, and treatment.

\medterm{Breast Pain} Pain or tenderness in one or both breasts. Common causes include hormonal changes, cysts, infection, injury, medicines, or chest-wall problems; persistent focal pain or pain with a lump needs assessment.

\synonyms
Mastalgia; Mastodynia

\medterm{Breast Rash} A change such as redness, scaling, itching, or irritation of breast skin. Causes include dermatitis, allergy, infection, or rarely inflammatory breast cancer; a persistent, one-sided, painful, or nipple-associated rash needs assessment.

\medterm{Breast Reconstruction} Breast Reconstruction is surgery to rebuild the shape and appearance of a breast, most often after mastectomy. It may use a breast implant, tissue transferred from another part of the body, or both. Reconstruction may begin during mastectomy (immediate reconstruction) or be performed later (delayed reconstruction). It may also include reconstruction of the nipple and areola and may require several procedures.

\medterm{Breast Reconstruction, Autologous} Surgical reconstruction of the breast contour utilizing the patient's own vascularized tissue flaps harvested from donor sites such as the abdomen (e.g., deep inferior epigastric perforator [DIEP] or transverse rectus abdominis myocutaneous [TRAM] flap), back (latissimus dorsi flap), or gluteal/thigh regions.

\synonyms
Autologous Breast Reconstruction; Flap Breast Reconstruction; Autogenous Tissue Reconstruction; DIEP Flap Reconstruction

\medterm{Breast Reconstruction, Implant-Based} Surgical reconstruction of the breast mound following mastectomy or trauma utilizing silicone gel-filled or sterile saline-filled prosthetic implants, often preceded by submuscular or prepectoral placement of a temporary tissue expander.

\synonyms
Prosthetic Breast Reconstruction; Implant Breast Reconstruction; Tissue Expander Breast Reconstruction

\medterm{Breast Reduction} Surgery that removes breast tissue, fat, and skin to reduce breast size and relieve symptoms such as back, neck, or shoulder pain. Possible effects include scars, altered sensation, breastfeeding difficulty, and uneven shape.

\medterm{Breast Screening} Testing people without breast symptoms to find breast cancer early, most commonly with mammography. The recommended starting age and interval depend on local guidance, personal risk, family history, and previous findings.

\medterm{Breast Self-Examination} Looking at and feeling one's own breasts to notice a new lump, skin or nipple change, or discharge. It may improve familiarity with normal changes but does not replace recommended screening or clinical assessment.

\medterm{Breast Skin And Nipple Changes} Changes such as dimpling, redness, scaling, inversion, crusting, or discharge may result from benign conditions, infection, hormonal changes, or cancer. New, persistent, one-sided, or bloody changes require assessment.

\medterm{Breast Surgeon} A surgeon who diagnoses and treats breast conditions, including benign disease, breast cancer, reconstruction, and operations that reduce cancer risk.

\medterm{Breast Surgery} An operation on the breast to remove or examine a lump, treat cancer, reduce size, reconstruct shape, or address another condition. The type of surgery depends on the diagnosis, extent of disease, and person's goals.

\medterm{Breast Ultrasound} Imaging that uses sound waves to examine breast tissue. It helps assess a lump, distinguish fluid from solid tissue, investigate a mammogram finding, guide a biopsy, and evaluate selected implants or dense breasts.

\medterm{Breast-Conserving Treatment} An oncological treatment approach for early-stage breast cancer that aims to treat the cancer effectively while preserving the natural appearance of the breast. It consists of breast-conserving surgery (a lumpectomy, wide local excision, or quadrantectomy to completely remove the tumor with clear microscopic tissue margins and assess sentinel lymph nodes) followed by adjuvant whole-breast radiation therapy to destroy any remaining microscopic cancer cells. Large randomized clinical trials demonstrate that breast-conserving therapy yields long-term survival rates equivalent to complete mastectomy. Contraindications include multicentric disease across multiple quadrants, extensive diffuse microcalcifications, prior chest wall radiation, or persistently positive surgical margins.

\synonyms
Breast-Conserving Therapy; BCT; Breast-Conserving Surgery; BCS; Lumpectomy with Radiotherapy

\medterm{Breastfeeding} Feeding an infant human milk from the breast. It provides nutrition and immune protection; many health organisations recommend exclusive breastfeeding for about six months, followed by complementary foods while breastfeeding continues as desired and feasible.

\medterm{Breastfeeding With Rheumatoid Arthritis} A person with rheumatoid arthritis may be able to breastfeed. Disease activity, fatigue, hand function, and medicine safety are reviewed because some treatments are compatible and others require avoidance or careful timing.

\medterm{Breastfeeding-Associated Nipple and Skin Changes} Dermatologic and mucosal complications of the nipple-areolar complex occurring during lactation, including painful erythema, skin fissures, milk blisters (blebs), contact dermatitis, Raynaud phenomenon of the nipple, and infectious mastitis or candidiasis.

\synonyms
Nipple Trauma in Breastfeeding; Sore Nipples; Nipple Fissures in Lactation

\medterm{Breath} Air moved into and out of the lungs. Its smell, depth, rate, pattern, and oxygen or carbon-dioxide content can provide clues about breathing, infection, metabolism, or other health problems.

\medterm{Breath Alcohol Test} A test measuring alcohol in exhaled air to estimate alcohol concentration in the blood. Results depend on timing, the device, sampling, and health factors and may be used in clinical or legal settings.

\medterm{Breath Analysis} Measurement of gases or other chemicals in exhaled air. Different tests can estimate alcohol, assess airway inflammation, detect infection, or study metabolism; results depend on the method and usually complement rather than replace standard testing.

\medterm{Breath Focus} A mindfulness or relaxation exercise that directs attention to the feeling of breathing. It may help some people reduce stress or emotional arousal but does not replace treatment for a medical or mental-health condition.

\synonyms
Breath-Focused Meditation; Breathing-Focused Mindfulness

\medterm{Breath Odor} An unpleasant smell from the mouth or exhaled air. Common causes include dental or gum disease, dry mouth, tobacco, food, medicines, and infection; persistent odor requires assessment by a dentist or clinician.

\medterm{Breath Sounds} Sounds made as air moves through the airways and lungs, heard with a stethoscope. Reduced, absent, or unusual sounds can provide clues about mucus, narrowing, fluid, collapse, or other lung problems.

\medterm{Breath Stacking} A technique in which several small breaths are taken without fully exhaling, increasing the amount of air in the lungs. It may improve cough strength in selected people with weak breathing muscles and is taught by a professional.

\medterm{Breath Test} Analysis of exhaled air to detect a substance made by the body or by microorganisms. Examples include tests for Helicobacter pylori and tests measuring hydrogen or methane after a sugar drink to assess digestion and absorption.

\medterm{Breath, Bad} Alternate name; see \textit{Bad Breath}.

\medterm{Breath-Holding Spell} A brief involuntary episode in a young child, usually after crying, pain, fear, or anger, in which breathing pauses and the child may change colour or briefly faint before recovering.

\medterm{Breath-Holding Spells} Brief involuntary episodes, usually between about six months and six years, triggered by upset, pain, or fear. A child may cry, exhale, change colour, stiffen or faint briefly, then recover; a first or atypical episode needs assessment.

\medterm{Breathalyzer} A device that estimates alcohol concentration in the blood by measuring alcohol in exhaled air. Results can be affected by timing, device accuracy, recent drinking, mouth alcohol, and sampling conditions.

\medterm{Breathing} Movement of air into and out of the lungs. It brings oxygen into the body and removes carbon dioxide through coordinated activity of the brain, nerves, muscles, and airways.

\medterm{Breathing Difficulty} Uncomfortable or difficult breathing, also called breathlessness. Causes include lung or heart disease, airway blockage, anaemia, infection, anxiety, and metabolic illness; sudden, severe, or worsening difficulty needs urgent assessment.

\medterm{Breathing Difficulty First Aid} Immediate emergency supportive care provided to an individual experiencing acute respiratory distress, including positioning the patient upright in a tripod or sitting posture, loosening restrictive clothing, assisting with prescribed bronchodilator or epinephrine autoinjector administration, and promptly contacting emergency medical services.

\synonyms
First Aid for Dyspnea; Emergency Respiratory Support; Respiratory Distress Management

\medterm{Breathing Difficulty, Recumbent} Alternate name; see \textit{Orthopnea}.

\medterm{Breathing Exercises} Deliberate changes in breathing pattern used to support relaxation, lung expansion, airway clearance, breath control, or rehabilitation. The safest technique depends on the person's condition and is taught when needed.

\medterm{Breathing, Slowed or Arrested} Pathologic depression of the respiratory rate (bradypnea) or complete cessation of ventilatory airflow (apnea), resulting from central nervous system depression, drug toxicity (such as opioid overdose), traumatic brain injury, upper airway obstruction, or terminal cardiopulmonary arrest requiring emergent airway management and resuscitation.

\synonyms
Bradypnea; Apnea; Respiratory Arrest; Hypopnea

\medterm{Breathlessness} Alternate name; see \textit{Dyspnea}.

\medterm{Breech Baby} A fetus positioned in the womb with the buttocks or feet pointing downward toward the birth canal instead of the head. Most babies turn into a head-down (vertex) position before late pregnancy, but those remaining breech near term may be delivered by cesarean section or planned vaginal breech birth, or turned via external cephalic version.
\synonyms Breech fetus; Breech presentation.

\medterm{Breech Birth} Birth in which the baby is positioned bottom or feet first rather than head first. The safest delivery plan depends on the exact position, pregnancy stage, baby and parent factors, and the team's experience.

\medterm{Breech Presentation} A fetal position in which the buttocks or feet point toward the birth canal. Types include frank, complete, and footling breech; risks include cord prolapse, difficulty delivering the head, and birth injury, so delivery needs careful planning.
\synonyms Breech position; Malpresentation.

\medterm{Bregma} The point on the skull where the main seam between the frontal and parietal bones meets the seam between the two parietal bones. In infants it lies beneath the soft spot at the front of the skull.

\medterm{Bremsstrahlung} Radiation produced when fast electrons slow down or change direction near matter. It is relevant to some imaging and radiation treatments and requires appropriate shielding and safety controls.

\medterm{Brenner Tumor} A rare ovarian growth made of cells resembling those lining the urinary tract. Most are small, one-sided, noncancerous, and found incidentally, but uncommon borderline or cancerous forms also occur.

\medterm{Brenner Tumour} A usually noncancerous ovarian growth whose cells resemble those lining the urinary tract. Rare borderline and cancerous forms occur and are distinguished by tissue examination.

\medterm{Breslau Bacillus} An older name for Salmonella Typhimurium, a bacterium that can cause foodborne gastroenteritis with diarrhoea, abdominal cramps, fever, and vomiting.

\medterm{Breslow Thickness} The depth, measured in millimetres, to which a melanoma has grown into the skin. It is an important part of staging and helps estimate the risk of spread and guide treatment.

\medterm{Bretylium Tosylate} An older medicine that can suppress certain dangerous abnormal heart rhythms, especially ventricular fibrillation. It is now rarely used because of limited availability, side effects, and newer treatments.

\medterm{Brevibacillus} A group of spore-forming bacteria found in soil, water, and food. Human infection is uncommon but can occur opportunistically, especially in people with weakened immunity or medical devices.

\medterm{Brevibacterium} A group of bacteria found on skin, in soil, and in food. They commonly contribute to body or cheese odour; rare species can cause opportunistic infection, particularly in people with serious illness or devices.

\medterm{Brevibacterium Casei} A species of Brevibacterium found in the environment and occasionally on people. It rarely causes bloodstream or device-related infection, usually in a person with major underlying illness or an indwelling device.

\medterm{Brevundimonas} A group of environmental water-associated bacteria that rarely cause human disease. Opportunistic infections, including bloodstream or wound infection, are more likely in people with weakened immunity or medical devices.

\medterm{Brevundimonas Diminuta} An environmental water-associated bacterium that rarely causes opportunistic human infection. Reported illnesses include bloodstream, wound, or device-related infection, mainly in people with serious underlying disease.

\medterm{Brevundimonas Vesicularis} An environmental bacterium that rarely causes human infection. When it is isolated from a clinical sample, symptoms, specimen quality, and other results determine whether it represents true disease or contamination.

\medterm{Bricklayer's Disease} An older occupational term for lung disease caused by breathing dust during masonry work, especially silica-containing dust. Long-term exposure can cause silicosis, which scars the lungs and impairs breathing.

\medterm{Bridge To Transplant} Temporary treatment or mechanical support that keeps a person with advanced heart failure alive and stable until a heart transplant can be performed. The device, timing, and risks depend on heart function and overall health.

\medterm{Bridge Work} A fixed dental replacement for one or more missing teeth. Artificial teeth are attached to prepared neighbouring teeth or implants, restoring some chewing function and appearance.

\medterm{Brief Psychotic Disorder} A sudden episode of delusions, hallucinations, severely disorganized speech, or markedly disorganized behavior lasting at least one day but less than one month, followed by a return toward usual functioning.

\synonyms
Brief Reactive Psychosis; Acute Transient Psychosis

\medterm{Brief Resolved Unexplained Event} A sudden, short, and resolved episode in an infant involving abnormal breathing, skin colour, muscle tone, or responsiveness, with no explanation after an initial assessment. Risk depends on the infant's age, history, examination, and whether it recurs.

\medterm{Brief Resolved Unexplained Event} A sudden, brief (under 1 minute), fully resolved episode in an infant under one year of age characterized by one or more of: cyanosis or pallor, absent, decreased, or irregular breathing, marked change in tone (hypertonia or hypotonia), and altered level of responsiveness, with no explanatory underlying cause identified following a thorough history and physical examination.

\synonyms
BRUE; Apparent Life-Threatening Event; ALTE

\medterm{Bright's Disease} An old term for several kidney diseases, especially conditions affecting the kidney's filters and causing protein in urine, swelling, or altered urine. Modern diagnosis identifies the specific kidney problem instead.

\medterm{Brill's Disease} Recurrence of epidemic typhus years after the first infection because dormant Rickettsia prowazekii becomes active again. It is usually milder than the original illness but can still be spread by body lice.

\medterm{Brill-Zinsser Disease} Recurrent epidemic typhus caused by reactivation of dormant Rickettsia prowazekii years after the original infection. The illness is usually milder but can be transmitted by body lice.

\medterm{Brimonidine Tartrate} A medicine that reduces certain nerve signals and blood-vessel activity. Eye drops lower pressure inside the eye, and skin preparations can reduce some facial redness; side effects and interactions may require medical attention.

\medterm{Brine Bath} Bathing in salty water, sometimes used for comfort or certain skin conditions. It does not treat every skin problem, and hot water or salt can worsen irritation, dry skin, wounds, or some medical conditions.

\medterm{Brittle Asthma} An uncommon term for severe asthma with sudden, unpredictable, or difficult-to-control attacks despite appropriate treatment. Specialist review is needed to confirm asthma, identify triggers or other diagnoses, and reduce risk.

\medterm{Brittle Bone Disease} An inherited disorder, usually osteogenesis imperfecta, in which bones break easily. It can also cause short stature, loose joints, blue-gray eye whites, hearing loss, dental problems, or other skeletal changes.

\medterm{Brittle Bone Disease} Alternate name; see \textit{Osteogenesis Imperfecta}.

\medterm{Brittle Diabetes} Alternate historical term; see \textit{Type 1 Diabetes}.

\medterm{Brivudine} An antiviral medicine used in some countries for shingles. It is not used with fluoropyrimidine cancer medicines or certain related treatments because the interaction can be fatal.

\medterm{Broad Ligament} A sheet of tissue extending from the sides of the uterus to the pelvic walls. It helps support the uterus, fallopian tubes, and ovaries and contains blood vessels, nerves, and connective tissue.

\medterm{Broad Nasal Bridge} An unusually wide upper part of the nose between the eyes. It may be a normal family feature or part of a condition present from birth or an inherited syndrome.

\medterm{Broad-Based Gait} An abnormal walking pattern in which the feet are placed abnormally far apart to increase the base of support and maintain balance. It is a cardinal clinical physical sign of midline cerebellar vermis lesions (cerebellar ataxia) or severe loss of lower limb proprioception (sensory ataxia).

\medterm{Broca Aphasia} A language disorder in which speech is slow, effortful, and difficult to produce, while understanding may be relatively preserved. It usually follows damage to the dominant frontal language area, often from a stroke.

\medterm{Broca's Aphasia} A nonfluent language disorder caused by damage to the dominant frontal language region. Speech is effortful and limited, while understanding may be better preserved than speaking or repeating words.

\medterm{Broca's Area} A region in the dominant frontal lobe that helps plan and produce spoken language. Damage can cause slow, effortful speech and difficulty forming words, while understanding may be relatively preserved.

\medterm{Brodie Abscess} A chronic, well-circumscribed subacute pyogenic bone abscess that represents a localized form of osteomyelitis. It characteristically develops in the metaphysis of long bones (most frequently the proximal or distal tibia) in children and young adults, presenting as a lytic bone lesion with a surrounding ring of sclerotic bone on radiography.

\medterm{Brodie's Abscess} A localized, slowly developing pocket of bone infection, often near the end of a long bone. It commonly causes persistent, sometimes night-time, bone pain with little fever and may require prolonged antibiotics or surgery.

\medterm{Brody Disease} A rare inherited muscle disorder causing painless stiffness and slow muscle relaxation after strenuous or repeated exercise. It often begins in childhood and may affect the face, eyelids, or hands; symptoms can resemble myotonia but arise from a different muscle problem.

\synonyms
Brody Myopathy; SERCA1-Related Muscle Disease

\medterm{Broken Ankle} A break in one or more bones forming the ankle joint. It usually causes pain, swelling, bruising, and difficulty bearing weight; deformity, an open wound, numbness, or poor circulation requires emergency care.

\medterm{Broken Arm} A break in one or more arm bones, usually causing pain, swelling, bruising, deformity, or loss of function. An open injury, severe deformity, numbness, or reduced circulation requires urgent care.

\medterm{Broken Blood Vessel in Eye} Alternate name; see \textit{Subconjunctival Hemorrhage}.

\medterm{Broken Bone} A partial or complete break in a bone caused by sudden injury, repeated stress, or weakened bone. Pain, swelling, bruising, deformity, or loss of function may occur and imaging usually confirms the diagnosis.

\medterm{Broken Bone} Alternate name; see \textit{Bone Fracture}.

\medterm{Broken Clavicle} Alternate name; see \textit{Broken Collarbone}.

\medterm{Broken Collarbone} A break in the clavicle, usually after a fall or direct blow. It causes pain, swelling, tenderness, and reduced shoulder movement; numbness, breathing difficulty, or a skin-threatening deformity requires urgent care.

\synonyms
Broken Clavicle

\medterm{Broken Finger} A break in a finger bone caused by a blow, crush, or twist. Pain, swelling, bruising, deformity, and reduced movement may occur; assessment checks alignment, tendon function, sensation, and circulation and guides splinting or surgery.

\synonyms
Finger Fracture

\medterm{Broken Foot} A break in one or more foot bones after a fall, blow, twist, or repeated stress. Pain, swelling, bruising, and difficulty walking are common; deformity, an open wound, numbness, or inability to walk needs urgent assessment.

\medterm{Broken Hand} A break in one or more hand bones, causing pain, swelling, bruising, deformity, or reduced movement. Assessment includes alignment, tendon function, sensation, and circulation because some breaks require reduction or surgery.

\medterm{Broken Heart Syndrome} Temporary weakening of the heart muscle, often after severe emotional or physical stress. It can resemble a heart attack with chest pain or breathlessness, but the coronary arteries usually have no blocking clot.

\medterm{Broken Hip} Alternate name; see \textit{Hip Fracture}.

\medterm{Broken Leg} A break in one or more leg bones, usually causing pain, swelling, bruising, deformity, or inability to bear weight. An open injury, numbness, or loss of circulation requires emergency care.

\synonyms
Leg Fracture

\medterm{Broken Little Finger} Alternate name; see \textit{Broken Finger}.
\medterm{Broken Nose} A break in one or more nasal bones that may cause pain, swelling, bruising, bleeding, deformity, or difficulty breathing through the nose. Persistent bleeding, clear fluid leakage, or severe deformity needs urgent assessment.

\medterm{Broken Or Dislocated Jaw} A break or displacement of the jaw that can cause pain, swelling, an altered bite, limited movement, loose teeth, or difficulty speaking, swallowing, or breathing. It needs urgent medical and dental assessment.

\medterm{Broken Or Knocked Out Tooth} A tooth damaged or completely forced out by injury. A knocked-out permanent tooth is a dental emergency; the tooth is handled by the crown, kept moist, and the root is not scrubbed.

\medterm{Broken Ribs} Breaks in one or more ribs, usually causing local chest pain that worsens with breathing, coughing, or movement. Severe breathing difficulty, increasing pain, coughing blood, or major injury requires urgent assessment.

\medterm{Broken Toe} A break in a toe bone, usually after a direct blow or crush. Pain, swelling, bruising, and tenderness are common; deformity, an open wound, numbness, or poor circulation needs prompt assessment.

\synonyms
Toe Fracture; Phalanx Fracture

\medterm{Broken Wrist} A break in one or more bones near the wrist, often causing pain, swelling, bruising, deformity, or reduced hand function. Numbness, pale or cold fingers, an open wound, or severe deformity needs urgent care.

\medterm{Bromelain} A mixture of enzymes from pineapple. It is sold as a supplement and studied for pain and inflammation, but evidence and product strength vary; it can cause allergy, stomach upset, or interactions with medicines.

\medterm{Bromhexine} A medicine that thins sticky mucus in the airways so it may be coughed up more easily. It can cause stomach upset or rash and rarely causes a serious skin reaction.

\medterm{Bromhidrosis} Unusually strong or unpleasant body odour produced when microorganisms break down sweat or skin secretions. It commonly affects the armpits or feet and may occur with excessive sweating or skin disease.

\medterm{Bromides} Salts or compounds containing bromide. Excessive or prolonged exposure can cause bromism, with symptoms affecting the nervous system, mood, digestion, or skin.

\medterm{Bromine} A reactive chemical used in industry and some flame-resistant materials. Bromine and many bromine compounds can irritate or burn tissue, while substantial exposure can cause poisoning affecting several organs.

\medterm{Bromism} Poisoning caused by prolonged or excessive exposure to bromide compounds. It may cause confusion, drowsiness, poor coordination, mood or behaviour changes, digestive symptoms, acne-like eruptions, or other neurologic problems.

\medterm{Bromobenzenes} Industrial chemicals made by attaching one or more bromine atoms to benzene. Exposure can irritate the eyes, skin, and airways or cause other toxic effects, depending on the specific compound and amount.

\medterm{Bromocriptine} A medicine that stimulates dopamine receptors and reduces prolactin release. It is used for high prolactin and selected movement or hormone disorders; nausea, dizziness, low blood pressure, hallucinations, and other effects may occur.

\medterm{Bromoderma} An inflammatory skin eruption linked to bromine or bromide exposure. It may appear as acne-like, pustular, or other irritated lesions and usually improves when the source is identified and removed.

\medterm{Bromodosis} Very unpleasant foot odour caused when sweat, moisture, and skin debris are broken down by bacteria or fungi. Washing, drying, breathable footwear, and treatment of any infection may help.

\medterm{Bromophos} An organophosphate insecticide. Exposure can overstimulate nerves and cause sweating, excess saliva, vomiting, breathing difficulty, weakness, seizures, or collapse and requires urgent poison or emergency care.

\medterm{Brompheniramine Overdose} Poisoning from taking too much brompheniramine, an antihistamine. It may cause severe sleepiness or agitation, confusion, dry mouth, rapid heartbeat, seizures, or dangerous heart rhythms; urgent help is needed.

\medterm{Bromthymol Blue} A chemical indicator that changes colour according to acidity and alkalinity. It is used in laboratories and is not a medicine or treatment.

\medterm{Bronchi} The two main airways carrying air from the windpipe into the lungs. Each divides into smaller branches; the right main bronchus is wider and more vertical, so inhaled objects are more likely to enter it.

\medterm{Bronchial} Relating to the bronchi, the large airways branching from the windpipe. Bronchial disease may cause cough, wheezing, mucus production, narrowing, or reduced airflow.

\medterm{Bronchial Artery} An artery from the body's general circulation that supplies the walls of the bronchi and supporting lung tissues. It differs from the pulmonary arteries, which carry blood to the lungs for oxygen exchange.

\medterm{Bronchial Artery Embolization} A catheter procedure that blocks bleeding arteries supplying the airways, usually to control severe or repeated coughing of blood. Imaging guides the catheter and the material used; careful technique is needed to avoid blocking arteries supplying the spinal cord.

\synonyms
BAE; Bronchial Embolization; Embolization of Bronchial Arteries

\medterm{Bronchial Asthma} A long-term condition in which sensitive, inflamed airways narrow intermittently. It can cause wheezing, cough, chest tightness, and breathlessness; symptoms often vary with triggers and improve with appropriate treatment.

\medterm{Bronchial Biopsy} Removal of a small tissue sample from an airway during bronchoscopy. Laboratory examination can investigate inflammation, infection, abnormal growths, or cancer; bleeding and breathing complications are possible.

\medterm{Bronchial Breath Sounds} Harsh, high-pitched, hollow breath sounds with a prolonged expiratory phase and an audible pause between inspiration and expiration. Normally auscultated over the trachea and main bronchi, their presence over peripheral lung fields indicates abnormal transmission through consolidated lung tissue, as in lobar pneumonia.

\medterm{Bronchial Carcinoid} A usually slow-growing neuroendocrine tumour arising in an airway. Typical carcinoids are generally less aggressive than atypical carcinoids, but either can block an airway, cause coughing or bleeding, or spread to other sites.

\medterm{Bronchial Fistula} An abnormal passage between an airway and another structure, such as the space around the lung, food pipe, or skin. It can cause persistent air leakage, coughing, infection, or breathing failure.

\medterm{Bronchial Hemorrhage} Bleeding from airway blood vessels, usually seen as coughing up blood. Heavy bleeding can block the airway or cause breathing failure and requires emergency treatment.

\medterm{Bronchial Obstruction} Partial or complete blockage of an airway by mucus, an inhaled object, a tumour, swelling, or pressure from outside. It can reduce airflow and lead to lung collapse, repeated infection, wheezing, or breathlessness.

\medterm{Bronchial Stenosis} Narrowing of an airway caused by scarring, inflammation, infection, injury, a tumour, or a previous procedure. It can cause localised wheezing, breathlessness, repeated infection, or poor mucus clearance and may require widening or surgery.

\medterm{Bronchial Stricture} Abnormal narrowing of an airway caused by scarring, inflammation, infection, a tumour, or previous instrumentation. It may produce localised wheezing, cough, repeated infection, or collapse of the lung beyond the narrowing.

\medterm{Bronchial Tree} The branching network of airways from the windpipe through the main, lobar, and smaller bronchi and bronchioles toward the air sacs. It conducts and distributes air throughout the lungs.

\medterm{Bronchial Tubes} The two main bronchi and their branches, which carry air from the windpipe into the lungs and divide into progressively smaller airways.

\synonyms
Bronchi; Bronchial Tree When Referring to the Branching Airway System

\medterm{Bronchial Vein} A vein that drains blood from the walls of the larger airways and nearby lung tissue. It empties partly into veins returning blood to the heart and partly into pulmonary veins.

\medterm{Bronchiectasis} Permanent widening and damage of the airways after repeated or severe inflammation and infection. It can cause daily cough, large amounts of mucus, breathlessness, repeated chest infections, and coughing up blood.

\medterm{Bronchiectasis, Acquired / Congenital} Alternate index label; see \textit{Bronchiectasis}.

\medterm{Bronchiole} A small airway without cartilage that branches from a bronchus and carries air toward the lung's air sacs. Its smooth muscle can narrow or widen the passage and regulate airflow.

\medterm{Bronchiolectasis} Permanent widening of the small airways, usually caused by chronic inflammation, infection, or blockage. It may occur with bronchiectasis and contribute to poor mucus clearance and repeated infection.

\medterm{Bronchioles} Small branching airways that carry air from the bronchi toward the air sacs and regulate airflow with their surrounding muscle.

\medterm{Bronchiolitis} Inflammation and blockage of the smallest airways, usually from a viral infection in infants and young children. It can cause cough, wheezing, fast breathing, chest retractions, poor feeding, and low oxygen.

\medterm{Bronchiolitis Obliterans} A rare disease in which inflammation and scarring permanently narrow or close the smallest airways. It can follow severe infection, inhaled injury, autoimmune disease, or transplantation and causes persistent cough and breathlessness.

\medterm{Bronchitis} Inflammation of the larger airways. Acute bronchitis is usually caused by a virus and causes cough, with or without mucus, while chronic bronchitis means mucus-producing cough for at least three months in two successive years.

\medterm{Bronchitis, Acute} Alternate name; see \textit{Acute Bronchitis}.

\medterm{Bronchitis, Chronic} Alternate name; see \textit{Chronic Bronchitis}.

\medterm{Bronchoalveolar Lavage} A bronchoscopy procedure in which sterile fluid is put into a small lung area and then collected. The sample can be tested for infection, inflammation, abnormal cells, or other lung disease.

\medterm{Bronchocele} Abnormal widening of a bronchus filled with mucus, usually because an airway is blocked or developed abnormally. It may cause cough, repeated infection, or no symptoms and is assessed with imaging.

\medterm{Bronchoconstriction} Narrowing of the airways because their muscle tightens, the lining swells, or mucus builds up. It increases resistance to airflow and commonly occurs in asthma, allergic reactions, or irritant exposure.

\medterm{Bronchoconstrictor Agents} Substances that narrow the airways by tightening airway muscle or causing inflammation. Examples include histamine, leukotrienes, irritant gases, and some medicines.

\medterm{Bronchodilator} A medicine that relaxes airway muscle and widens breathing tubes. It can relieve or prevent wheezing and breathlessness in asthma, chronic obstructive lung disease, and other conditions with narrowed airways.

\medterm{Bronchodilator Response} The change in airflow measured by spirometry after a bronchodilator is given. Improvement can support variable airway narrowing, but no response does not exclude asthma and a response does not prove one specific disease.

\medterm{Bronchodilator Reversibility} Improvement in measured airflow after an inhaled bronchodilator. It can support variable airway narrowing, but one result alone cannot confirm or exclude asthma or chronic obstructive pulmonary disease.

\medterm{Bronchoesophageal Fistula} An abnormal connection between an airway and the food pipe. It may be present from birth or develop from cancer, infection, injury, or treatment and can cause coughing during swallowing, repeated pneumonia, or breathing problems.

\medterm{Bronchogenic Carcinoma} A cancer that begins in the airways or nearby lung tissue. It may cause cough, coughing blood, chest pain, breathlessness, or no early symptoms; diagnosis and treatment depend on the cancer type and extent.

\medterm{Bronchogenic Cyst} A birth-related fluid-filled cyst that develops from early airway or food-pipe tissue, often near the windpipe or central chest. It may cause no symptoms or press on nearby structures and sometimes requires removal.

\medterm{Bronchomalacia} Weakness of an airway wall that allows it to collapse too much, especially during breathing out. It may cause noisy breathing, cough, breathlessness, repeated infections, or difficulty clearing mucus.

\medterm{Bronchophony} An examination finding in which spoken sounds are heard unusually clearly through a stethoscope over part of the lung. It can occur when lung tissue is solidified by infection or another process but is not diagnostic by itself.

\medterm{Bronchopleural Fistula} An abnormal connection between an airway and the space around the lung, often after lung surgery or severe infection. It can cause a persistent air leak, breathlessness, infection, or air under the skin.

\medterm{Bronchopneumonia} Infection causing scattered areas of inflammation and consolidation around the small airways in one or more lung lobes. It may cause fever, cough, mucus, chest pain, fast breathing, and low oxygen.

\medterm{Bronchopulmonary Dysplasia} A chronic lung disease in some premature infants whose developing lungs have been affected by immaturity, inflammation, oxygen, or breathing support. It may cause ongoing oxygen need, fast breathing, or wheezing.

\medterm{Bronchopulmonary Segments} Natural sections of a lung lobe, each supplied by its own airway branch and pulmonary artery branch. Their separation helps doctors describe disease and allows some operations to remove only an affected segment.

\medterm{Bronchoscope} A thin tube with a light and camera used to examine the windpipe and airways. Instruments passed through it can remove an object or collect tissue, fluid, or mucus for testing.

\medterm{Bronchoscopic Culture} Laboratory testing of fluid or tissue collected during bronchoscopy to look for bacteria, fungi, or mycobacteria when sputum or other ordinary samples are not sufficient.

\medterm{Bronchoscopy} Examination of the windpipe and airways with a flexible or rigid camera tube. It can investigate persistent cough, infection, bleeding, or a lung mass and can remove foreign material or collect samples; sedation or anaesthesia is used.

\medterm{Bronchospasm} Sudden tightening of airway muscles that narrows the breathing passages. It can cause wheezing, cough, chest tightness, and difficulty breathing and may be triggered by asthma, allergy, infection, irritants, or airway procedures.

\medterm{Bronchostenosis} Narrowing of a bronchus that restricts airflow. Causes include scarring, inflammation, infection, a tumour, an inhaled object, pressure from outside, or a previous procedure and may cause wheezing or repeated infection.

\medterm{Bronchus} One of the two main branches of the windpipe that carries air into a lung. Each divides into smaller branches; the right main bronchus is wider and more vertical, making it a more common site for inhaled objects to lodge.

\medterm{Bronchus Disease} Any condition affecting one or more bronchi, the air passages leading into the lungs. Infection, asthma, chronic inflammation, blockage, scarring, or a tumour may reduce airflow and cause cough, wheezing, or repeated infection.

\medterm{Bronchus Neoplasm} An abnormal growth arising in a bronchus. It may be noncancerous or cancerous; imaging, bronchoscopy, and tissue testing help determine its type, extent, and treatment.

\medterm{Bronze Diabetes} Historical name for hereditary hemochromatosis; see \textit{Hemochromatosis}.

\medterm{Brooke Disease} A rare inherited skin disorder involving abnormal thickening of the outer skin or hair follicles. Because the eponym has been used for more than one condition, the original source and clinical features are checked.

\medterm{Brooke-Spiegler Syndrome} An inherited disorder causing multiple usually noncancerous skin tumours, often on the head and neck. It is commonly linked to changes in the CYLD gene and may include cylindromas, spiradenomas, or trichoepitheliomas.

\medterm{Broselow Tape} A colour-coded tape that estimates a child's weight from body length during an emergency. It helps select medicine doses, fluid amounts, equipment sizes, and defibrillator settings, but clinical assessment and an actual weight are used when available.

\synonyms
Broselow Pediatric Emergency Tape; Pediatric Length-Based Resuscitation Tape

\medterm{Brow Presentation} A fetal position in which the forehead leads because the head is partly extended. The head presents a wide diameter, so persistent brow presentation at term may prevent vaginal birth and requires specialist delivery planning.

\medterm{Brown Fat} A type of body fat containing many energy-producing structures that generate heat. It is especially important for keeping newborns warm and can also be activated by cold in adults.

\medterm{Brown Lung Disease} An occupational breathing disorder caused by repeated exposure to cotton, flax, or hemp dust, traditionally called byssinosis. It can cause chest tightness, cough, wheezing, and work-related narrowing of the airways.

\medterm{Brown Recluse Spider} A venomous spider whose bite may initially be mild but can later cause painful skin damage, blistering, or tissue death. Rarely, its venom causes fever, red-cell destruction, kidney injury, or other systemic illness.

\medterm{Brown Recluse Spider Bite} Alternate name; see \textit{Spider Bites}.

\medterm{Brown Sequard Syndrome} A spinal-cord injury affecting mainly one side. It typically causes weakness and loss of vibration or position sense on the injured side and loss of pain and temperature sensation on the opposite side below the injury.

\medterm{Brown Stain} A brown deposit on teeth caused by pigments from food, drinks, tobacco, or bacteria. Professional cleaning may remove surface stains, but decay or enamel disease needs separate treatment.

\medterm{Brown Syndrome} Limited ability to lift the eye when it is turned inward, because the superior oblique tendon cannot glide freely through its pulley, the trochlea. The affected eye appears lower than the other in that gaze direction and double vision may occur. It can be present at birth or develop from inflammation, injury, sinus disease or surgery, or rheumatic disease such as rheumatoid arthritis. Some cases resolve on their own; persistent double vision is treated by an eye specialist.

\medterm{Brown Teeth} Brown discolouration of teeth caused by surface stains, decay, enamel defects, excess fluoride, medicines, injury, or other conditions. A dental examination identifies the cause and appropriate treatment.

\medterm{Brown-SéQuard Syndrome} A spinal-cord injury affecting mainly one side. It causes weakness and loss of position or vibration sense on the injured side and loss of pain and temperature sensation on the opposite side below the injury.

\medterm{Bruce Protocol} A graded treadmill test in which speed and slope increase at set stages. It assesses exercise capacity, symptoms, heart-rate and blood-pressure responses, and signs of reduced blood flow to the heart.

\medterm{Brucella} A group of bacteria that usually infect animals but can infect people through unpasteurised dairy products, inhalation, or contact with infected animals or tissues. Human infection causes brucellosis.

\medterm{Brucella Abortus} A Brucella species mainly associated with cattle. It can infect people exposed to infected animals, tissues, or unpasteurised dairy products and cause prolonged fever and other brucellosis symptoms.

\medterm{Brucella Canis} A Brucella species associated with dogs. It rarely infects people after contact with infected dogs or reproductive material and may cause prolonged fever or other brucellosis symptoms.

\medterm{Brucella Melitensis} A Brucella species mainly associated with goats and sheep and a frequent cause of severe human brucellosis, usually acquired through unpasteurised dairy products or infected animals.

\medterm{Brucella Suis} A Brucella species associated with pigs and some wildlife. It can infect people through animal contact or contaminated food and may cause prolonged illness or infection in bones and organs.

\medterm{Brucellaceae} A family of bacteria that includes Brucella and related animal-associated species. Some members cause brucellosis, an infection that can pass from animals to people.

\medterm{Brucellosis} An infection caused by Brucella bacteria, usually acquired from infected animals, animal tissues, or unpasteurised dairy products. It may cause prolonged fever, sweats, tiredness, headache, and muscle or joint pain.

\medterm{Bruch's Membrane} A thin supporting layer between the retina and the blood-rich tissue behind it. It helps nourish and maintain the retina and can become damaged in ageing, diabetes, and some retinal diseases.

\medterm{Bruck's Syndrome} A rare inherited disorder combining fragile bones with joint stiffness or fixed contractures present from birth. It results from changes affecting connective tissue and may cause repeated fractures and limited movement.

\medterm{Brudzinski Sign} Involuntary bending of the hips and knees when the neck is gently bent forward. It can occur with irritation of the tissues around the brain and spinal cord, including meningitis, but does not confirm the diagnosis.

\medterm{Brugada Phenocopy} A temporary ECG pattern resembling Brugada syndrome but caused by another problem, such as fever, abnormal blood salts, reduced heart blood flow, medicines, or poisoning. The pattern resolves when the cause is corrected and is distinguished from inherited disease.

\synonyms
Acquired Brugada Pattern; BrP

\medterm{Brugada Syndrome} An inherited heart-rhythm disorder that can cause dangerous ventricular arrhythmias and sudden cardiac death, often during rest, sleep, or fever. Diagnosis relies on a characteristic ECG pattern and clinical assessment.

\medterm{Brugia} A group of parasitic worms spread by mosquitoes that can cause lymphatic filariasis. The infection may damage lymph vessels and lead to recurrent fever, swelling, or long-term enlargement of limbs or genital tissues.

\medterm{Bruise} Bleeding under unbroken skin after small blood vessels are damaged, usually by a blow, pressure, or injury. It commonly changes colour as it heals; unexplained or frequent bruising needs medical assessment.

\medterm{Bruises} Areas of bleeding beneath the skin after small blood vessels are damaged. Frequent, large, spontaneous, painful, or unusually persistent bruises may result from medicines, platelet or clotting problems, liver disease, nutritional deficiency, or injury.

\medterm{Bruising} Purplish, blue, green, or yellow skin discolouration caused by blood leaking into tissue, usually after injury. Colour changes as the blood is broken down; widespread or unexplained bruising needs assessment.

\medterm{Bruit} An abnormal whooshing sound heard over a blood vessel, usually with a stethoscope, caused by turbulent blood flow. It may suggest narrowing but does not by itself establish the cause or severity.

\medterm{Brunner Gland} A gland in the first part of the small intestine that releases mucus and alkaline fluid. These secretions protect the intestinal lining and help neutralise acid arriving from the stomach.

\medterm{Brunner's Gland} A mucus- and bicarbonate-producing gland in the duodenum, the first part of the small intestine. Its alkaline secretion protects the lining from stomach acid and supports digestion.

\medterm{Brunner's Glands} Glands in the duodenum that release mucus and bicarbonate-rich fluid to protect the intestinal lining from stomach acid and help create a suitable environment for digestive enzymes.

\medterm{Brush Burns} Superficial friction injuries caused by scraping skin against a rough surface. They cause pain, redness, and grazes; cleaning and protecting the area usually allow healing, but deep, extensive, or infected injuries need medical care.

\medterm{Brushfield's Spots} Small gray or white spots arranged around the coloured part of the eye. They are common in people with Down syndrome but can also occur in healthy people and usually do not affect vision.

\medterm{Bruxism} Repeated clenching or grinding of the teeth, often during sleep. It can wear or fracture teeth and cause jaw, face, ear, or headache pain; dental review can assess damage and possible sleep, stress, or medicine-related causes.

\synonyms
Teeth Grinding; Tooth Clenching

\medterm{Bruxism, Sleep / Awake} Alternate index label; see \textit{Bruxism}.

\medterm{Bryant's Traction} An older treatment in which an infant's or young child's legs are gently elevated with traction to treat selected hip or thigh conditions. It is rarely used and requires careful specialist monitoring.

\medterm{Bubble Bath Soap Poisoning} Illness after swallowing or inhaling bubble-bath product. A small amount usually causes mouth or stomach irritation, but repeated vomiting, coughing, breathing difficulty, or drowsiness requires emergency assessment.

\medterm{Bubo} A tender, enlarged lymph node, usually in the groin, armpit, or neck, caused by infection or inflammation. Buboes are classically linked with plague but also occur in sexually transmitted and other infections.

\medterm{Bubonic Plague} A serious infection caused by Yersinia pestis, usually spread by infected flea bites. It causes fever and painful swollen lymph nodes and can spread to the blood or lungs without prompt antibiotic treatment.

\medterm{Buccal} Relating to the cheek or its inner lining. A buccal medicine is placed between the cheek and gum so it dissolves and is absorbed through the mouth.

\medterm{Buccal Gland} A small salivary gland in the cheek that releases mucus into the mouth. The secretion helps keep the cheek lining moist and supports chewing and swallowing.

\medterm{Buccal Mucosa} The non-keratinized stratified squamous mucous membrane lining the inside of the cheeks, bounded by the labial commissures anteriorly, the retromolar trigone posteriorly, and the alveolar sulci superiorly and inferiorly. It contains the opening of the parotid duct (Stensen duct) opposite the maxillary second molar and serves as a site for clinical examination of exanthems (such as Koplik spots in measles), oral lichen planus (Wickham striae), aphthous stomatitis, leukoplakia, and drug absorption.

\medterm{Buccal Smear} A sample of cells gently collected from the inside of the cheek for laboratory testing. It may be used for genetic analysis or, less commonly, to investigate infection or abnormal cells.

\medterm{Buccalmucosa} A spelling variant of buccal mucosa, the moist lining on the inside of the cheeks.

\medterm{Buccinator} A muscle in the cheek that presses the cheek against the teeth. It helps keep food between the teeth and assists sucking, blowing, chewing, and speaking.

\medterm{Buckle Fracture} Alternate name; see \textit{Torus Fracture}.

\medterm{Budd-Chiari Syndrome} Blockage of veins draining blood from the liver, usually from a blood clot or a condition that makes clotting more likely. It can cause abdominal pain, enlarged liver, fluid in the abdomen, jaundice, and liver failure.

\medterm{Budesonide} A corticosteroid that reduces inflammation. It is used in selected inhaled, nasal, digestive, or rectal treatments for conditions such as asthma, chronic obstructive lung disease, allergic symptoms, and inflammatory bowel disease; side effects depend on the route and dose.

\medterm{Budvicia} A group of water-associated bacteria related to Yersinia. They rarely cause human disease, but some isolates have been associated with urinary or other opportunistic infections.

\medterm{Budvicia Aquatica} A water-associated bacterium that has rarely been isolated from human urinary infections, including infections in people with urinary catheters. Its significance depends on symptoms and specimen quality.

\medterm{Buerger Disease} An inflammatory disease that blocks small and medium blood vessels, especially in the hands and feet. It is strongly linked to tobacco use and can cause pain, ulcers, poor circulation, and loss of tissue.

\medterm{Buerger Test} A bedside physical examination assessment used to evaluate the severity of peripheral arterial occlusive disease in the lower extremities. The patient's leg is elevated to 45 or 90 degrees until maximum pallor develops (Buerger's angle of circulatory sufficiency), and is then placed in a dependent sitting position; a prolonged delay in the return of pink skin color followed by a deep red or purple reactive erythema (rubor of dependency) indicates advanced arterial insufficiency.

\medterm{Buerger's Disease} A tobacco-associated inflammatory disease of small and medium arteries and veins, also called thromboangiitis obliterans. Reduced blood flow can cause pain, colour changes, ulcers, gangrene, and loss of fingers or toes.

\medterm{Buffalo Hump} A pad of fat at the back of the neck and upper back. It may occur with obesity, long-term corticosteroid excess, some HIV treatments, or Cushing syndrome and is assessed together with other findings.

\medterm{Buffer} A chemical solution that resists a change in acidity when small amounts of acid or base are added. Buffers help maintain stable conditions in blood, cells, medicines, and laboratory tests.

\medterm{Buffy Coat} The thin pale layer of white blood cells and platelets that forms between plasma and red blood cells when a blood sample is spun. It can be collected for laboratory testing.

\medterm{Bufonidae} The family of amphibians commonly called true toads. Some species produce bufotoxins in their skin glands; swallowing or absorbing these toxins can affect the heart and nervous system.

\medterm{Buformin} An older biguanide medicine for diabetes that has been withdrawn or restricted in many countries because it can cause dangerous lactic-acid buildup, especially when kidney function is poor.

\medterm{Bug Spray Poisoning} Illness after exposure to an insecticide. Effects depend on the chemical and may include skin or eye irritation, vomiting, breathing difficulty, unusual sweating, confusion, seizures, or other nervous-system problems; severe poisoning requires urgent medical treatment.

\medterm{Bulbar} Relating to the lower brainstem or to the nerves controlling swallowing, speaking, chewing, coughing, and movement of the tongue and soft palate. Bulbar weakness can allow food or saliva to enter the airway.

\medterm{Bulbar Conjunctiva} The part of the eye's clear moist covering that lies over the white of the eye. It joins the inner surface of the eyelids around the front of the eye and can become inflamed or injured.

\medterm{Bulbar Palsy} Weakness or paralysis of muscles controlled by lower cranial nerves. It can cause slurred speech, difficulty swallowing, weak coughing, poor handling of saliva, and sometimes breathing problems.

\medterm{Bulbar Paralysis} Weakness or paralysis of muscles used for speaking, swallowing, chewing, coughing, and moving the tongue or palate. It may result from brainstem, motor-neuron, or neuromuscular disease and can threaten breathing.

\medterm{Bulbocavernosus} A paired muscle in the pelvic floor that helps empty the urethra and contributes to erection and ejaculation. It surrounds structures at the base of the penis or the vaginal opening.

\medterm{Bulbospongiosus} A paired pelvic-floor muscle that helps empty the urethra and assists erection and the release of genital secretions during sexual activity.

\medterm{Bulbourethral Gland} One of two small glands below the prostate that release a slippery, alkaline fluid into the urethra during sexual arousal. The fluid helps lubricate the passage and reduce acidity.

\medterm{Bulbourethral Glands} Two small male reproductive glands, also called Cowper glands, that release mucus-like fluid into the urethra during sexual arousal. The fluid lubricates the urethra and may reduce acidity.

\medterm{Bulging} Projecting or swelling outward beyond the usual shape. A bulge may involve a disc, membrane, organ, skin area, or body wall, and its significance depends on the structure and cause.

\medterm{Bulging Eyes} Alternate name; see \textit{Exophthalmos}.
\medterm{Bulging Fontanelle} An outward protrusion or tense swelling of an infant's soft spot (fontanelle) when the baby is calm and upright, indicating elevated intracranial pressure from conditions such as meningitis, encephalitis, hydrocephalus, or intracranial hemorrhage.

\synonyms
Tense Fontanelle; Fontanelle Protrusion; Raised Intracranial Pressure in Infants
\medterm{Bulimia} Alternate name; see \textit{Bulimia Nervosa}.

\medterm{Bulimia Nervosa} An eating disorder involving recurrent binge-eating episodes, a sense of loss of control, and recurrent inappropriate compensatory behaviors such as vomiting, laxative misuse, fasting, or excessive exercise. Self-evaluation is unduly influenced by body shape or weight, and the pattern causes clinically significant distress or impairment.

\synonyms
Bulimia; Bulimic Disorder

\medterm{Bulla} A large fluid-filled blister on the skin or mucous membrane, or a large air-filled space in an organ such as the lung. The cause and treatment depend on its location.

\medterm{Bullae} Large fluid-filled blisters on or beneath the skin or a mucous membrane. They can result from friction, burns, infection, allergy, medicines, or autoimmune disease and may need care if widespread or infected.

\medterm{Bullectomy} An operation to remove one or more large air-filled spaces in damaged lung tissue. It may be considered when a giant bulla compresses healthier lung, causes severe breathlessness, or contributes to repeated pneumothorax.

\synonyms
Excision of Pulmonary Bullae; Resection of Giant Bulla

\medterm{Bullet Track} The path a firearm projectile takes through the body, from the entry wound through injured tissues to an exit wound or a retained bullet. The path may damage skin, organs, blood vessels, nerves, or bone.

\medterm{Bullet Wounds} Penetrating or perforating injuries caused by firearm projectiles. They may damage skin, organs, blood vessels, nerves, and bones and require urgent assessment even when the external wound looks small.

\medterm{Bullous} Describes a condition involving large, raised, fluid-filled blisters. Bullous disease can result from infection, burns, medicines, allergy, or an autoimmune reaction.

\medterm{Bullous Impetigo} A contagious superficial bacterial skin infection, usually caused by Staphylococcus, that produces fragile fluid-filled blisters and yellow or varnish-like crusts. It is especially common in infants and young children.

\medterm{Bullous Pemphigoid} An autoimmune blistering disease, usually in older adults, causing intensely itchy skin and firm blisters on normal or reddened skin. It results from antibodies attacking proteins that hold the skin layers together.

\synonyms
BP; Paraneoplastic Pemphigoid

\medterm{Bully} A person who repeatedly harms, threatens, or intimidates someone with less power. Bullying can cause injury, anxiety, depression, sleep problems, school avoidance, and other health effects.

\medterm{Bullying} Repeated physical, verbal, social, or online aggression involving a power imbalance. It can harm mental and physical health; safety planning, support, and intervention by responsible adults are important.

\medterm{Bumetanide} A diuretic that makes the kidneys remove extra salt and water in urine. It is used for swelling related to heart, liver, or kidney disease; excessive fluid loss can cause dehydration and abnormal blood salts.

\medterm{Bumps} Raised or palpable areas on or under the skin. Causes range from harmless cysts, blocked follicles, and insect bites to injury, allergy, infection, or tumours; a growing, painful, hard, or persistent bump requires assessment.

\medterm{BUN} Alternate name; see \textit{Blood Urea Nitrogen Test}.

\medterm{Bundle Branch Block} Delay or interruption of electrical signals through one of the heart's lower-chamber conduction branches. It may cause no symptoms or reflect heart disease and is diagnosed on an ECG.

\medterm{Bundle Of His} A short pathway of specialised heart tissue that carries electrical signals from the atria to the ventricles. It divides into right and left branches so the lower chambers contract in a coordinated way.

\medterm{Bundle-Branch Block} Delayed or interrupted electrical signalling through one of the heart's lower-chamber conduction branches. It can change the ECG and heartbeat pattern and may be harmless or related to underlying heart disease.

\medterm{Bungarotoxins} Poisons in the venom of krait snakes that disrupt nerve signals to muscles. A bite can cause drooping eyelids, weakness, progressive paralysis, and breathing failure and requires emergency treatment.

\medterm{Bunion} A change in the joint at the base of the big toe that pushes the toe toward the others and creates a bump on the inner foot. It may cause pain, redness, swelling, stiffness, and difficulty wearing shoes.

\medterm{Bunion Removal} An operation to correct a painful bunion by realigning bone and, when needed, balancing nearby soft tissues. It is considered when symptoms persist despite suitable shoes and other nonsurgical care.

\medterm{Bunionectomy} Surgery to correct a painful bunion by cutting or shifting bone and adjusting surrounding tissues. The operation is chosen according to the deformity, symptoms, activity, other health conditions, and response to nonsurgical care.

\medterm{Bunionette} A bony prominence at the outer side of the joint below the little toe. It may become painful, red, or swollen from pressure, tight shoes, or altered foot mechanics.

\synonyms
Tailor's Bunion; Bunion of the Fifth Toe

\medterm{Bunions} Deformities at the base of the big toe that push the toe toward the others and create a bump on the inner foot. They may cause pain, swelling, stiffness, and difficulty finding comfortable shoes.

\medterm{Bunioplasty} Surgery to correct a bunion by realigning the toe joint and, when needed, balancing nearby soft tissues. The operation is selected according to the deformity and symptoms.

\medterm{Bunyaviridae} An older virus-family grouping that included several viruses spread by insects, ticks, or rodents. Modern classification has changed; individual viruses differ in the illnesses they cause.

\medterm{Bunyaviruses} Alternate name; see \textit{Viral Hemorrhagic Fevers}.

\medterm{Buphthalmos} Enlargement of an infant's or young child's eye caused by high pressure inside the eye while its outer coat is still flexible. It is usually associated with congenital or infantile glaucoma and needs urgent eye care.

\medterm{Bupivacaine} A long-acting local anaesthetic that temporarily blocks nerve signals to prevent pain during local, nerve-block, or epidural procedures. Too much in the bloodstream can cause seizures, dangerous heart effects, or cardiac arrest.

\medterm{Bupivacaine Hydrochloride} The hydrochloride form of bupivacaine, a long-acting local anaesthetic used to block pain signals. Accidental injection into a blood vessel can cause severe neurologic or heart toxicity.

\medterm{Buprenorphine} A medicine that partly activates opioid receptors and is used to treat opioid-use disorder and some pain. It can cause dependence and overdose, especially with other sedatives, although respiratory suppression is limited compared with full opioid agonists.

\medterm{Bupropion} A medicine used for depression and, in some settings, to help stop smoking. It changes norepinephrine and dopamine signalling and may cause insomnia, dry mouth, or anxiety; it increases seizure risk and is unsuitable for some seizure or eating disorders.

\medterm{Burkholderia} A group of environmental bacteria found in soil and water. Some species cause melioidosis or glanders, while others cause difficult-to-treat lung, bloodstream, or wound infections, especially in people with chronic lung disease or weakened immunity.

\medterm{Burkholderia Cenocepacia} A bacterium in the Burkholderia cepacia complex that can cause severe, persistent lung infection, especially in people with cystic fibrosis or other chronic lung disease.

\medterm{Burkholderia Cepacia} An environmental bacterium in the Burkholderia cepacia complex. It can cause opportunistic lung, bloodstream, wound, or device-related infection, particularly in people with cystic fibrosis or weakened immunity.

\medterm{Burkholderia Gladioli} An environmental bacterium that can rarely cause lung, bloodstream, or wound infection, especially in people with cystic fibrosis or weakened immunity.

\medterm{Burkholderia Mallei} The bacterium that causes glanders, an infection mainly affecting horses and related animals. People can become infected through contact with infected animals or tissues.

\medterm{Burkholderia Pseudomallei} The bacterium that causes melioidosis. Infection follows contact with contaminated soil or water and may cause pneumonia, abscesses, bloodstream infection, or long-term illness.

\medterm{Burkholderiaceae} A family of mainly environmental bacteria that includes Burkholderia. Some members cause serious infections in people or animals, while others are harmless or used in environmental research.

\medterm{Burkitt Lymphoma} A very fast-growing cancer of B lymphocytes, often involving the abdomen, jaw, bone marrow, or nervous system. It requires prompt intensive treatment; the pattern differs between endemic, sporadic, and immunodeficiency-associated forms.

\synonyms
Burkitt's Tumor

\medterm{Burkitt's Lymphoma} A very fast-growing cancer of B lymphocytes. It may affect the abdomen, jaw, bone marrow, or brain and spinal cord and needs prompt intensive treatment with careful supportive care.

\medterm{Burkitt's Tumor} Alternate name; see \textit{Burkitt Lymphoma}.

\medterm{Burn} An injury to the skin or deeper tissue caused by heat, chemicals, electricity, radiation, friction, or extreme cold. Severity depends on depth, size, location, and associated injury; extensive or deep burns need urgent medical care.

\medterm{Burn Injury} Alternate name; see \textit{Burns}.

\medterm{Burn Surgery} Surgical care for serious burns, including assessing depth and area, removing dead tissue, covering wounds, performing skin grafts, and treating scars or movement-limiting contractures.

\medterm{Burn, Oral} Thermal, chemical, or electrical injury to the oral mucosal lining, palate, tongue, or lips, resulting in acute pain, erythema, epithelial desquamation, ulceration, and potential upper airway compromise.

\synonyms
Oral Burn; Oral Cavity Burn; Mouth Burn; Thermal Palatal Burn

\medterm{Burning Feet} Burning, tingling, or painful sensations in the feet, often caused by nerve damage from diabetes, nutritional deficiency, alcohol, medicines, or nerve compression. New, progressive, or weakness-associated symptoms need assessment.

\medterm{Burning Mouth Syndrome} A chronic oral pain disorder causing daily burning, scalding, or tingling discomfort in the mouth, often involving the tongue, lips, or palate, usually with a normal-appearing mucosa. Altered taste or dry mouth may occur. It may be primary, with no identifiable cause, or secondary to conditions or exposures such as nutritional deficiency, oral infection, medication effects, allergy, diabetes, or hormonal change; evaluation excludes treatable causes.

\synonyms
Scalded Mouth Syndrome
Stomatodynia

\medterm{Burning Thigh Pain} Alternate name; see \textit{Meralgia Paresthetica}.

\medterm{Burnout} A work-related state of emotional exhaustion, mental distance or cynicism, and reduced effectiveness caused by chronic workplace stress. It is not the same as depression, although both may occur together.

\medterm{Burns} Tissue injuries caused by heat, chemicals, electricity, radiation, friction, or extreme cold. Seriousness depends on depth, area, and location; burns affecting the face, airway, hands, genitals, major joints, or large areas need prompt care.

\synonyms
Thermal Injury; Burn Trauma

\medterm{Burns, First Aid} Alternate name; see \textit{Burns}.

\medterm{Burnt-Out Joint} An older description of a joint in which active inflammation has subsided but permanent damage, deformity, pain, or restricted movement remains.

\medterm{Burping} Release of swallowed air or gas from the stomach through the mouth. Occasional burping is normal; frequent, painful, or newly worsening burping may result from swallowed air, reflux, food intolerance, or another digestive disorder.

\medterm{Burr Cells} Erythrocytes that exhibit numerous (usually 10 to 30) short, blunt, evenly spaced, and symmetrical spicules across their entire surface membrane. They are typically observed in end-stage renal disease and uremia, severe hypomagnesemia, and pyruvate kinase deficiency.

\synonyms
Echinocytes; Crenated Red Cells

\medterm{Burr Hole Evacuation} An emergency or urgent neurosurgical procedure in which one or two small holes (approximately 10 to 14 mm in diameter) are drilled through the skull bone to drain an abnormal fluid or blood collection compressing the brain. Most frequently performed to treat chronic subdural hematomas (which commonly occur in older adults after minor, often forgotten head injuries), a neurosurgeon incises the scalp and dura mater, washes out the liquefied old blood and clots with warm sterile saline, and places a temporary flexible drain to prevent fluid reaccumulation. By rapidly decompressing the brain, burr hole drainage reverses life-threatening brain herniation, confusion, focal weakness, and lethargy with significantly less surgical trauma, blood loss, and recovery time than a formal craniotomy.

\synonyms
Burr Hole Drainage; Trephination; Chronic Subdural Hematoma Drainage; Burr Hole Craniostomy

\medterm{Burroughs} An inconsistent surname or eponym in medical sources rather than a precise diagnosis; the associated condition or person must be identified from context.

\medterm{Bursa} A small fluid-filled sac that reduces friction between moving tissues, such as a tendon and bone or skin and bone, especially near a joint. Inflammation of a bursa is called bursitis.

\synonyms
Bursae
Synovial Bursae

\medterm{Bursa Inflammation} Alternate name; see \textit{Bursitis}.

\medterm{Bursectomy} Surgical removal of a bursa, usually considered for persistent painful bursitis, recurrent infection, or a structural problem that has not improved with appropriate nonsurgical treatment.

\medterm{Bursitis} Inflammation of a bursa, the small fluid-filled sac that cushions movement between tendons, muscles, bones, and skin. It causes localized pain, tenderness, swelling, or reduced movement and may result from pressure, injury, infection, or arthritis.

\medterm{Bursitis Of The Heel} Inflammation of a bursa near the heel, causing localized tenderness, swelling, and pain when walking or when pressure is applied. It may follow repeated friction, injury, inflammatory arthritis, or infection.

\medterm{Bursitis of the Knee} Alternate name; see \textit{Knee Bursitis}.

\medterm{Bursitis, Calcific} Alternate name; see \textit{Calcific Bursitis}.

\medterm{Bursitis, Shoulder} Alternate name; see \textit{Shoulder Bursitis}.

\medterm{Bursopathy} Any disorder of a bursa, including inflammation, infection, bleeding, thickening, or degeneration. Affected bursae may cause localized pain, swelling, tenderness, and restricted movement around a joint.

\medterm{Buruli Ulcer} A skin infection caused by Mycobacterium ulcerans, usually affecting the arms or legs. It often begins as a painless lump or swelling and progresses to tissue destruction and a large ulcer; untreated disease can damage deeper tissue and cause permanent disability.

\medterm{Burundanga} A colloquial term for scopolamine (hyoscine hydrobromide), a tropane alkaloid and muscarinic receptor antagonist that induces central anticholinergic effects, anterograde amnesia, sedation, and severe toxicity.

\synonyms
Scopolamine; Hyoscine; Hyoscine Hydrobromide

\medterm{Buserelin} A synthetic gonadotropin-releasing hormone analogue. Continuous treatment first briefly increases and then suppresses pituitary reproductive hormones, lowering ovarian or testicular hormone production in selected fertility and hormone-sensitive conditions.

\medterm{Buspirone} A nonbenzodiazepine medicine used mainly for generalized anxiety disorder. It acts on serotonin signaling and usually takes days to weeks to help; it does not provide immediate relief from an acute panic attack.

\medterm{Busulfan} An alkylating chemotherapy medicine used mainly to prepare the body for blood-forming stem-cell transplantation and sometimes to treat chronic myeloid leukemia. It can suppress bone marrow and may cause seizures, infertility, lung injury, or liver damage.

\medterm{Butanol} A flammable industrial alcohol used as a solvent. Swallowing it or substantial exposure to its vapor can irritate the eyes, skin, and airways and may cause dizziness, vomiting, drowsiness, or other nervous-system effects.

\medterm{Butazolidin Overdose} Toxic effects from excessive phenylbutazone, an older nonsteroidal anti-inflammatory drug. Severe exposure may cause gastrointestinal bleeding, kidney injury, bone-marrow suppression, metabolic disturbances, or seizures.

\medterm{Butocarboxim} A carbamate insecticide that inhibits acetylcholinesterase. Significant exposure can cause cholinergic poisoning, with sweating, salivation, vomiting, diarrhea, wheezing, muscle weakness, seizures, or respiratory failure.

\medterm{Butter, Margarine, And Cooking Oils} Butter, margarine, and cooking oils are dietary fat sources. Their nutritional effects depend on portion size, fatty-acid composition, and the overall eating pattern.

\medterm{Buttiauxella} A genus of environmental gram-negative bacteria in the Enterobacterales; human infection is uncommon and usually opportunistic.

\medterm{Buttiauxella Agrestis} A species of environmental Gram-negative bacteria found in water, soil, and other natural settings. It is rarely reported from human clinical samples; when detected, the specimen source and the patient’s symptoms determine whether it represents infection or harmless contamination.

\medterm{Buttiauxella Ferragutiae} A species of environmental Gram-negative bacteria. It is uncommon in human disease, and a laboratory finding requires interpretation with the specimen source, repeat cultures, symptoms, and other evidence before deciding whether treatment is needed.

\medterm{Buttiauxella Gaviniae} A species of Buttiauxella bacteria found mainly in environmental samples. Human infection is uncommon; identification from a clinical specimen should be reviewed with the microbiology laboratory and interpreted according to the patient’s condition.

\medterm{Buttiauxella Noackiae} A species of environmental Gram-negative bacteria that is rarely associated with human infection. Its significance depends on where it was found, whether it was recovered repeatedly, and whether the patient has compatible symptoms or signs of infection.

\medterm{Buttock} The buttock is the rounded posterior region overlying the gluteal muscles, used for movement, sitting, examination, and selected injections.

\medterm{Buttocks} The paired fleshy regions overlying the gluteal muscles at the back of the pelvis, formed mainly by muscle and subcutaneous fat.

\medterm{Button Batteries} Button batteries are small batteries that can cause severe tissue injury if swallowed or inserted into the nose or ear. Suspected ingestion requires emergency evaluation.

\medterm{Button Magnets In The Nasal Cavity} Small magnets placed in or inhaled into the nose can attract across the nasal partition and cut off its blood supply. They may cause tissue death or a hole and require urgent medical removal.

\medterm{Buttonhole Deformity} Buttonhole deformity is a finger deformity with flexion of the proximal interphalangeal joint and hyperextension of the distal interphalangeal joint, often due to central slip injury or rheumatoid disease.

\medterm{Buttressing} Buttressing is reinforcement or thickening that supports a weakened structure; in bone it may describe cortical stress adaptation or surgical structural support.

\medterm{Butylamine} Butylamine is an irritating industrial amine whose vapor or liquid can injure the eyes, skin, and respiratory tract after exposure.

\medterm{Butyric Acid} Butyric acid is a short-chain fatty acid produced by gut bacteria during fiber fermentation and used by colon cells as an energy source; concentrated exposure irritates tissue.

\medterm{Bx} A common medical abbreviation for biopsy, the removal of cells or tissue for microscopic or laboratory examination.

\medterm{By-Pass} A route made around a blocked or narrowed passage. In medicine, a bypass may redirect blood, food, bile, urine, or another fluid through a natural or surgically created pathway.

\medterm{Bypass} Creation of an alternative route around a blocked or narrowed passage. A bypass can restore blood flow or redirect body contents, using a natural pathway, tube, vein, or artificial graft depending on the organ involved.

\medterm{Bypass Graft} A vein or artificial tube surgically connected to route blood around a blocked or weakened artery. Grafts are used in coronary, leg, and aortic operations; blood flow, infection, clotting, and long-term narrowing affect their durability.

\medterm{Bypass Surgery} An operation that creates a new route around a blocked or diseased passage. Examples include coronary and leg-artery bypasses and gastric bypass; possible complications include bleeding, infection, clotting, and narrowing of the connection.

\medterm{Byssinosis} An occupational lung disease caused by repeated inhalation of dust from cotton, flax, hemp, or other textile fibers. It can cause chest tightness, cough, and breathing difficulty that worsen during or after work exposure.


\medterm{C And T Cells} B cells make antibodies, while T cells coordinate immune responses or destroy infected and abnormal cells. Both are types of lymphocytes, a group of white blood cells that protect against infection and disease.
\medterm{C Fiber} A small, slow nerve fiber without an insulating covering that carries dull pain, warmth, and itch from the body to the spinal cord and brain.
\medterm{C-Peptide} A small protein fragment released when the pancreas makes insulin. Measuring it helps show how much insulin a person's body is making, rather than how much injected insulin is present. Kidney function and blood-glucose levels can affect the result.

\medterm{C-Reactive Protein} A protein made by the liver that rises when there is inflammation or infection. It can help monitor inflammatory diseases and possible infection, but it does not identify the cause by itself. A high-sensitivity CRP test may be used with other information to assess heart-disease risk.

\medterm{C-Reactive Protein Test} A blood test that measures C-reactive protein. The level may rise with inflammation or infection, but the test cannot identify the cause by itself and must be interpreted with symptoms and other findings.

\medterm{C-Section} A cesarean birth in which a baby is delivered through surgical openings in the abdomen and uterus. It may be planned or performed urgently when vaginal birth could endanger the mother or baby.
\medterm{C. Diff} Alternate index label; see \textit{C. difficile infection}.

\medterm{C. Difficile} Clostridioides difficile is a spore-forming bacterium that can multiply in the bowel, especially after antibiotics disrupt normal bacteria. Its toxins cause watery diarrhea and colon inflammation, which can become severe or recur.

\medterm{C. Difficile Colitis} Alternate index label; see \textit{Pseudomembranous colitis}.

\medterm{C. Difficile Infection} An infection caused by Clostridioides difficile, a bacterium that can grow excessively in the bowel after antibiotic use. It commonly causes watery diarrhea and inflammation of the colon and may lead to dehydration, severe colitis, or recurrence.

\synonyms
C. Diff

\medterm{C1Q Nephropathy} A kidney disorder in which the immune protein C1q is deposited in the glomeruli, the kidney's filtering units. It may cause protein or blood in the urine, swelling, high blood pressure, or reduced kidney function.

\medterm{C282Y Mutation} A change in the HFE gene that is the most common genetic cause of hereditary hemochromatosis, a condition in which the body stores too much iron. Having two copies of this mutation increases the risk, although not everyone with two copies develops the disease.

\medterm{CA 19-9} A substance measured in blood that may be increased in pancreatic and other digestive-system cancers. It can also rise with blocked bile ducts or inflammation, so it cannot screen for or diagnose cancer by itself.

\medterm{CA 27-29} A blood test that measures a protein sometimes released by breast-cancer cells. It may help monitor selected people with known breast cancer but cannot by itself screen for or diagnose cancer.

\medterm{CA-125} A blood protein that may rise in ovarian and some other cancers. It can also increase with menstruation, pregnancy, endometriosis, infection, liver disease, or other noncancerous conditions, so it cannot diagnose ovarian cancer alone.

\synonyms
CA125

\medterm{CA-125 Blood Test} See \textit{CA-125}.

\medterm{Cabazitaxel} A taxane chemotherapy medicine used with prednisone for metastatic castration-resistant prostate cancer after docetaxel. It can cause low white-cell counts, infection, diarrhea, and serious allergic reactions.

\medterm{Cabergoline} A dopamine D2-receptor agonist used to suppress prolactin secretion and treat hyperprolactinemia; it is also used in selected Parkinson disease settings and requires monitoring for adverse effects.

\medterm{CABG} Coronary artery bypass grafting, surgery that creates a new path around a narrowed or blocked heart artery to improve blood flow to the heart muscle.

\synonyms
Coronary artery bypass grafting; coronary bypass surgery.

\medterm{Cabot Ring} A threadlike or ring-shaped remnant seen inside a red blood cell on a blood smear. It may occur in severe anemia or when the bone marrow is producing abnormal red blood cells.

\medterm{Cachectic} Affected by cachexia, a severe wasting syndrome involving involuntary weight loss, muscle loss, weakness, and often reduced appetite in chronic disease.

\medterm{Cachetic} An alternate spelling of cachectic, describing severe unintentional loss of body weight and muscle caused by a serious long-term illness.

\medterm{Cachexia} A wasting syndrome caused by serious long-term illness, marked by ongoing unintentional loss of weight, muscle, and body fat. It may cause poor appetite, weakness, and inflammation and cannot be corrected by nutrition alone.

\medterm{Cacosmia} A distorted sense of smell in which ordinary odors seem unusually foul or unpleasant. It may occur with nose or sinus disease, infection, seizures, head injury, or another nervous-system disorder.

\medterm{CAD} Coronary artery disease, narrowing or blockage of the arteries supplying the heart, usually caused by fatty plaque. It can reduce oxygen to heart muscle and cause chest pain, heart attack, heart failure, or dangerous rhythms.

\medterm{CADASIL Syndrome} An inherited disorder of small arteries in the brain caused by changes in the NOTCH3 gene. It may cause migraine with aura, repeated small strokes, mood problems, and progressive memory or thinking difficulties.

\medterm{Cadaver} A dead human body used for anatomical study, surgical training, research, or identification. A donated cadaver may be preserved so tissues remain suitable for examination and teaching.
\medterm{Cadaveric Spasm} A rare form of muscle stiffening that occurs immediately at death in a specific group of voluntary muscles, without the usual period of relaxation. It may preserve the position held at the moment of death.

\medterm{Cadherin} A cell-surface protein that helps neighboring cells attach to one another and keeps tissues organized. Loss or alteration of cadherins can allow some cancer cells to invade nearby tissue and spread.

\medterm{Cadmium} A toxic metal found in tobacco smoke, some batteries and pigments, workplace dust, and contaminated food or soil. Long-term exposure can damage the kidneys, lungs, and bones and increases the risk of some cancers.

\medterm{Cadmium Chloride} A toxic compound containing cadmium, used mainly in industry and research. Inhalation or swallowing can injure the lungs, kidneys, and digestive tract; repeated exposure can damage bones and increase cancer risk.

\medterm{Caduceus} A winged staff with two entwined snakes, traditionally associated with Hermes. It is often used incorrectly as a medical symbol and differs from the single-snake Rod of Asclepius.

\medterm{Cadusafos} An organophosphate pesticide that blocks acetylcholinesterase, causing excess nerve signaling after exposure. Poisoning may produce sweating, salivation, vomiting, diarrhea, small pupils, muscle weakness, breathing difficulty, or seizures.

\medterm{Caecal} Relating to the caecum, the pouch at the beginning of the large intestine where the small intestine joins the colon and the appendix begins.

\medterm{Caecitis} Inflammation of the caecum, the first part of the large intestine. It may result from infection, inflammatory bowel disease, reduced blood flow, neutropenia, or spread of nearby inflammation and can cause right-sided abdominal pain and fever.

\medterm{Caecopexy} An operation that attaches the caecum, the first pouch of the large intestine, to nearby tissue to prevent or treat abnormal movement or twisting.

\medterm{Caecostomy} An operation that creates an opening from the caecum to the abdominal wall. It may temporarily release pressure, drain the bowel, or divert stool when the normal route is blocked or unsafe.

\medterm{Caecotomy} An operation that opens the caecum, usually to remove an object, access the bowel, or assist another abdominal procedure. The opening may be closed during the same operation.

\medterm{Caecum} The first pouch-like part of the large intestine, located where the ileum joins the colon and attached to the appendix.

\medterm{Caenorhabditis Elegans} A small, free-living roundworm widely used in laboratory research. Its simple nervous system, short life cycle, and well-studied genes help scientists investigate development, aging, nerve function, and disease mechanisms.

\medterm{Caesarean Birth} Delivery of a baby through surgical openings in the abdomen and uterus. It may be planned or urgent when problems such as fetal distress, labor complications, abnormal presentation, placenta previa, or multiple pregnancy make vaginal birth unsafe.

\medterm{Caesarean Section} A caesarean section is surgical delivery of a baby through incisions in the abdomen and uterus when vaginal birth is unsafe or not feasible.
\medterm{Cafe-Au-Lait Macules} Flat, evenly colored light-brown skin patches, named for their coffee-with-milk color. A few may be harmless; numerous or unusually large patches can be a sign of an inherited condition and should be assessed.

\medterm{Cafe-Au-Lait Spot} A flat, sharply outlined, evenly pigmented light-brown patch of skin. One or a few are often harmless, while many lesions or a characteristic pattern may indicate neurofibromatosis type 1 or another inherited disorder.

\medterm{Café Au Lait Spot} A flat, sharply demarcated, light-brown skin macule caused by increased melanin.
Multiple lesions may be a feature of neurofibromatosis type 1 or other genetic
disorders.

\medterm{Café Coronary} Sudden blockage of the airway by food, usually during a meal. Complete blockage prevents breathing and can rapidly cause collapse or death unless emergency choking first aid and medical treatment restore the airway.

\medterm{Café-Au-Lait Macules} Flat, light-brown skin patches with relatively uniform pigmentation. Multiple lesions may be associated with genetic disorders such as neurofibromatosis type 1.

\medterm{Caffeic Acid} A naturally occurring plant compound found in coffee, fruit, vegetables, and some herbs. It is studied in laboratory research for antioxidant and other biological effects but is not an established treatment for disease.

\medterm{Caffeine} A stimulant found in coffee, tea, chocolate, soft drinks, and energy drinks. It increases alertness and can raise heart rate, cause shakiness, or disturb sleep.

\medterm{Caffeine Intoxication} A toxic reaction to excessive caffeine, causing restlessness, anxiety, tremor, insomnia, nausea, vomiting, rapid heartbeat, or agitation. Severe poisoning can cause abnormal heart rhythms, seizures, or dangerous metabolic changes.

\medterm{Caffeine Overdose} A dangerous amount of caffeine can cause severe agitation, tremor, vomiting, rapid heartbeat, abnormal heart rhythms, low blood pressure, seizures, or disturbed body chemistry. Severe symptoms require urgent medical care.

\medterm{Caffeine Withdrawal} A group of symptoms after suddenly reducing or stopping regular caffeine use. Headache, tiredness, low mood, poor concentration, irritability, nausea, or flu-like feelings usually begin within a day and improve within several days.

\medterm{Caisson Disease} Decompression sickness caused by gas bubbles forming in tissues or blood after a rapid fall in surrounding pressure, such as after diving. It can cause joint pain, skin changes, breathing problems, brain or spinal-cord injury, or collapse.

\medterm{Cal (Calorie)} A unit of energy. In nutrition, a food Calorie is a kilocalorie, equal to 1,000 small calories, and describes the energy supplied by food or used by the body.

\medterm{Caladium Plant Poisoning} Poisoning from chewing or swallowing caladium, whose needle-like calcium oxalate crystals irritate the mouth and throat. Burning, pain, drooling, and swelling may occur; trouble breathing or swallowing is an emergency.

\medterm{Calamine} A skin preparation containing zinc compounds used to soothe itching and irritation from minor rashes, insect bites, and sunburn. It should not be applied to deep wounds, infected skin, or the eyes.
\medterm{Calbindin} A protein that binds calcium inside cells and helps control calcium levels and cell signaling, especially in nerve cells and intestinal cells. Laboratories may also use it as a marker to identify certain tissues.

\medterm{Calcaneal Spur} A bony growth from the heel bone, usually where the plantar fascia or Achilles tendon attaches. It may cause no symptoms; when heel pain occurs, the associated soft-tissue inflammation is often more important than the spur itself.

\medterm{Calcaneum} The calcaneum, or calcaneus, is the large bone forming the heel. It supports body weight, connects the ankle to the foot, and transmits force from the calf muscles through the Achilles tendon.
\medterm{Calcaneus} The largest bone in the ankle and foot, forming the heel. It supports body weight and connects with the ankle bone and cuboid; fractures commonly follow falls from height and may require imaging and specialist care.

\medterm{Calcifediol} The main circulating storage form of vitamin D, also called 25-hydroxyvitamin D. A blood test for calcifediol is used to assess vitamin-D status before the body converts it to the active hormone calcitriol.

\medterm{Calciferol} A fat-soluble form of vitamin D that helps the body absorb calcium and phosphate and maintain healthy bones, muscles, and immune function. The body changes it into active vitamin-D hormones.

\medterm{Calcific Bursitis} A condition in which calcium deposits form in or around a bursa, often near the shoulder, causing pain, tenderness, swelling, and restricted movement. Treatment depends on severity and may include rest, physical therapy, medicines, drainage, or deposit removal.

\synonyms
Bursitis Calcific (Calcific Bursitis)

\medterm{Calcific Tendinitis} A condition in which calcium crystals build up inside a tendon, most often a rotator-cuff tendon in the shoulder. It can cause sudden severe pain, stiffness, and restricted movement, sometimes followed by gradual improvement as deposits break down.

\medterm{Calcific Uremic Arteriolopathy} A rare, life-threatening disorder, usually in advanced kidney disease, in which calcium deposits and clotting block small blood vessels supplying the skin and fat. It causes extremely painful wounds, tissue death, and a high risk of infection.

\medterm{Calcification} Deposition of calcium salts in body tissue. It may be a normal part of bone formation, a response to injury or aging, or a sign of abnormal calcium or phosphate balance; its importance depends on the site and pattern.

\medterm{Calcified Granuloma} A small area of past inflammation that has healed and become hardened by calcium deposits. It is often found in the lung after an old infection and usually does not represent active disease.

\medterm{Calcineurin} An enzyme that helps activate T cells, a type of immune cell, after they receive signals. Medicines that block calcineurin reduce immune activity and are used to prevent transplant rejection and treat some immune diseases.

\medterm{Calcinos} An older or shortened term for calcinosis, in which calcium salts build up in the skin or tissues beneath it. Firm white or yellow lumps may form, sometimes ulcerate, and release chalky material.

\medterm{Calcinosis} Abnormal buildup of calcium salts in the skin or tissues beneath it. It may cause firm white or yellow lumps or plaques that are painful, limit movement, ulcerate, or release chalky material.

\medterm{Calciphylaxis} A rare, life-threatening disorder, usually associated with advanced kidney disease, in which small blood vessels supplying the skin become blocked by clotting and calcium deposits. It causes severe pain, skin wounds, tissue death, and serious infection risk.

\medterm{Calcitonin} A hormone made by specialized cells in the thyroid gland. It helps lower blood calcium by reducing calcium release from bone and increasing calcium loss through the kidneys, although it has a smaller role than parathyroid hormone in adults.

\medterm{Calcitonin Blood Test} A blood test that measures calcitonin, a hormone produced by thyroid C cells. It is used mainly to evaluate or monitor medullary thyroid carcinoma and may be affected by other thyroid, kidney, gastrointestinal, or medicine-related factors.

\medterm{Calcitriol} The active hormone made from vitamin D, mainly by the kidneys. It increases intestinal absorption of calcium and phosphate and helps regulate parathyroid hormone, bone formation, and blood-mineral levels.
\medterm{Calcitriol And Analogues} Medicines related to active vitamin D that reduce excessive parathyroid-hormone release and help manage mineral imbalance in some kidney disorders. They can raise calcium or phosphate, so blood levels require monitoring.

\medterm{Calcium} A mineral needed for strong bones and teeth, muscle movement, nerve signals, blood clotting, and normal cell function. Blood calcium exists in an active free form and forms attached to proteins or other substances.

\medterm{Calcium - Ionized} A blood test measuring ionized calcium, the biologically active fraction of circulating calcium. It is useful when total calcium may be misleading, such as with abnormal albumin, critical illness, or acid-base disturbances.

\medterm{Calcium - Urine} A urine test measuring calcium excretion over a specified collection period or in a spot specimen. It helps evaluate kidney stones, parathyroid disorders, bone disease, and abnormal calcium handling, and must be interpreted with kidney function and diet.


\medterm{Calcium Blood Test} A blood test that measures calcium, usually as total calcium. The result may be interpreted with albumin, a blood protein, or measured as ionized calcium when the active form is needed. It helps assess parathyroid, bone, and kidney disorders.

\medterm{Calcium Carbonate} A calcium supplement and antacid that neutralizes stomach acid. Excessive use can cause constipation, high blood calcium, or kidney problems and may interfere with medicine absorption.

\medterm{Calcium Carbonate Overdose} Overdose of calcium carbonate that may cause gastrointestinal symptoms, hypercalcemia, alkalosis, kidney injury, or kidney stones.

\medterm{Calcium Carbonate With Magnesium Overdose} Taking too much of these antacid ingredients can cause nausea, vomiting, constipation or diarrhea, weakness, confusion, abnormal calcium or magnesium levels, and kidney injury. Severe symptoms or reduced kidney function require urgent medical advice.
\medterm{Calcium Channel} A protein in the cell membrane that controls calcium entry into cells. Calcium channels help regulate muscle contraction, the release of substances, nerve signals, and heart rhythm.

\medterm{Calcium Chloride} A calcium salt given by injection for severe low blood calcium, dangerously high potassium, or certain toxic exposures. It can cause serious tissue injury if it leaks outside the vein and therefore requires careful administration.

\medterm{Calcium Citrate} A calcium supplement that is absorbed even when stomach acid is low. It may prevent or treat calcium deficiency but can cause constipation and, in some people, contribute to kidney stones.

\medterm{Calcium Deficiency} Too little calcium intake or availability for the body's needs. Long-term deficiency can weaken bones; severe deficiency can lower blood calcium and cause tingling, muscle cramps, spasms, or abnormal heart rhythms.

\medterm{Calcium Elevated (Hypercalcemia)} Alternate index label for Hypercalcemia.

\medterm{Calcium Gluconate} A calcium salt taken by mouth or given by injection to treat calcium deficiency and some emergencies, including dangerously high potassium or magnesium. Intravenous treatment requires medical monitoring because it can affect the heart and veins.

\medterm{Calcium Hydroxide Poisoning} Exposure to a strongly alkaline chemical that can burn the mouth, throat, esophagus, skin, or eyes. Do not induce vomiting; eye, breathing, or swallowing problems require immediate emergency assessment.

\medterm{Calcium Phosphate} A mineral that forms much of the hard structure of bones and teeth. Excess calcium phosphate can also deposit in soft tissues, especially after tissue damage or when blood calcium or phosphate levels are abnormal.

\medterm{Calcium Pyrophosphate Arthritis} Inflammation caused by calcium pyrophosphate crystals in cartilage or a joint. Sudden attacks, often called pseudogout, commonly affect the knee or wrist; chronic deposits can cause ongoing arthritis and joint damage.

\medterm{Calcium Pyrophosphate Deposition} Alternate index label; see \textit{Calcium Pyrophosphate Deposition Disease}.

\medterm{Calcium Pyrophosphate Deposition Disease} A disorder in which calcium pyrophosphate crystals collect in cartilage and joints. It may cause sudden painful swelling of one joint, usually the knee or wrist, or a chronic arthritis that resembles osteoarthritis or rheumatoid arthritis.

\synonyms
Calcium Pyrophosphate Deposition

\medterm{Calcium Sensing Receptor} A protein, especially in the parathyroid glands and kidneys, that detects the calcium level in blood. It helps control parathyroid hormone and determines how much calcium and phosphate the kidneys keep or remove.

\medterm{Calcium Signaling} A way cells communicate using changes in calcium concentration. These signals help control muscle contraction, secretion of substances, energy use, movement, and activity of genes.

\medterm{Calcium Stones} The most common type of kidney stone, made mainly of calcium oxalate or calcium phosphate. Low urine volume and certain mineral imbalances increase risk; stones may cause severe side pain, blood in the urine, blockage, or infection.

\medterm{Calculi} Hardened mineral concretions, such as gallstones, kidney stones, or salivary stones, formed when dissolved substances precipitate and aggregate within a body organ or duct.

\medterm{Calculus} A calculus is a hard deposit formed from mineral salts, such as a kidney stone, gallstone, or dental tartar.
\medterm{Calculus (Tartar)} Hardened dental plaque formed when minerals from saliva collect on teeth. Brushing cannot remove it; professional cleaning is needed because it traps bacteria and can cause gum inflammation and periodontal disease.
\synonyms
Tartar

\medterm{Calculus Removal} Professional dental treatment that removes hardened plaque from teeth and below the gumline using hand or ultrasonic instruments. It helps control gum inflammation and may include deep cleaning of affected tooth roots.

\medterm{Calculus, Renal} A kidney stone, a hard mass of mineral crystals formed in the urinary tract. It may pass without symptoms or cause severe wave-like pain, blood in the urine, blockage, vomiting, or urinary infection.

\medterm{Caldwell-Luc Operation} An operation that reaches the maxillary sinus through the upper gum and jawbone. It is used selectively to treat or examine certain sinus lesions, chronic disease, or foreign bodies when less invasive routes are unsuitable.

\medterm{Calf} The back part of the lower leg between the knee and ankle. It contains the gastrocnemius, soleus, blood vessels, nerves, and other tissues and helps point the foot downward and propel walking.

\medterm{Calf Pseudohypertrophy} An enlarged, firm appearance of the calf muscles caused mainly by replacement of damaged muscle with fat and fibrous tissue rather than increased muscle strength. It is seen especially in Duchenne and Becker muscular dystrophies and some related disorders.

\synonyms
Pseudohypertrophy of the Calves; Muscular Pseudohypertrophy

\medterm{Calibration} Checking and, when necessary, adjusting a medical or laboratory instrument against a known standard so its measurements are accurate and reliable. Regular calibration supports correct diagnosis, treatment, and patient safety.

\medterm{Caliciviridae} A family of viruses that includes noroviruses and sapoviruses, which commonly cause vomiting and diarrhea. Other members infect animals; disease depends on the virus, host, and route of exposure.

\medterm{Calipers} An instrument with two arms used to measure distance, thickness, or body dimensions. In health care, calipers may measure skinfold thickness, joint movement, wounds, or the size of body structures.

\medterm{Calla Lily Poisoning} Poisoning from chewing or swallowing calla lily, whose needle-like calcium oxalate crystals irritate the mouth and throat. Burning, pain, drooling, and swelling may occur; breathing or swallowing difficulty requires emergency care.

\medterm{Callosal} Relating to the corpus callosum, the broad bundle of nerve fibers connecting the left and right halves of the brain. Callosal abnormalities can affect communication between the two sides.
\medterm{Callosity} A callosity is a localized area of thickened skin caused by repeated pressure or friction.
\medterm{Callosotomy} Surgical division of the corpus callosum to reduce the spread of certain severe, treatment-resistant seizures between the cerebral hemispheres.

\medterm{Callus} A thick, hardened area of skin formed by repeated pressure or rubbing. Calluses commonly occur on hands or feet, usually cause no harm, and may become painful or crack when extensive.

\medterm{Calmodulin} A protein inside cells that senses calcium and passes its signals to other proteins. It helps control muscle contraction, movement of substances through cell membranes, chemical reactions, and activity of genes.

\medterm{Calnexin} A protein inside cells that helps newly made proteins fold into the correct shape. It checks proteins before they are sent to other parts of the cell or broken down.

\medterm{Calor} Heat or increased warmth of a body part, one of the classic local signs of inflammation.

\medterm{Calor, Dolor, Rubor, And Tumor} The traditional four signs of inflammation are heat, pain, redness, and swelling. Inflamed tissue may also lose some of its normal function.
\medterm{Caloric Restriction} Deliberately eating fewer calories without causing malnutrition. It may help with weight loss and some metabolic risks, but it should provide enough protein, vitamins, minerals, and other nutrients.
\medterm{Caloric Stimulation} A balance test in which warm or cool air or water is placed in the ear canal. The temperature briefly stimulates the inner-ear balance organs and helps compare the balance nerve on each side.

\medterm{Calorie} A unit used to describe the energy in food and the energy used by the body. On food labels, one Calorie means one kilocalorie, the energy needed to raise one kilogram of water by one degree Celsius.

\synonyms
a Calorie.




\medterm{Calories} Units used to express the energy supplied by food or expended by the body; food labels generally report kilocalories, commonly called calories.

\medterm{Calorimetry} Measurement of heat production or gas exchange to estimate how much energy the body uses. Indirect calorimetry uses oxygen consumption and carbon-dioxide production to help assess metabolism and nutritional needs.

\medterm{Calpain} A group of enzymes activated by calcium that cut selected proteins inside cells. They help regulate cell movement, signaling, programmed cell death, and platelet activity; abnormal calpain activity contributes to some muscle and nervous-system diseases.

\medterm{Calpainopathy} An inherited muscle disease caused by harmful changes in both copies of the CAPN3 gene. It slowly weakens hip, pelvic, shoulder, and trunk muscles, often causing shoulder-blade winging, back curvature, difficulty climbing, and walking problems while usually sparing facial and heart muscles.

\synonyms
Limb-Girdle Muscular Dystrophy R1 (LGMD R1); LGMD 2A; Calpain-3 Deficiency

\medterm{Calponin} A protein that binds to actin, a component of the cell's internal framework, and helps regulate contraction in smooth muscle. Laboratory staining for calponin can help identify muscle-related cells in tissue samples.

\medterm{Calprotectin} A protein released mainly by activated white blood cells. Measuring it in stool helps detect inflammation in the intestines, particularly in inflammatory bowel disease, and can help distinguish it from some noninflammatory conditions.
\medterm{Calreticulin} A protein inside cells that helps newly made proteins fold correctly and helps regulate calcium storage and immune signaling. Harmful changes in its gene are found in some blood cancers involving excessive production of blood cells.

\medterm{Calretinin} A calcium-binding protein found in certain nerve cells and cells lining body cavities. Pathologists may use calretinin staining to help distinguish mesothelioma from some other cancers in a tissue sample.

\medterm{Calvaria} The dome-shaped upper part of the skull that covers and protects the brain. It is formed mainly by the flat bones of the forehead, top and sides of the head, and back of the skull.

\medterm{Calvarium} The dome-like portion of the skull formed mainly by the frontal, parietal, and occipital bones. It encloses and protects the upper parts of the brain.

\medterm{Calyx} A cup-shaped body structure that supports tissue or collects fluid. In the kidney, renal calyces collect urine from the pyramids and carry it toward the renal pelvis.

\medterm{Cambridge Classification} A grading system used with endoscopic retrograde cholangiopancreatography to describe chronic pancreatitis. It rates abnormalities of the main pancreatic duct and its branches, from normal or uncertain findings to mild, moderate, or severe disease.

\medterm{Cameron Lesions} Long, shallow erosions or ulcers in the stomach where it is squeezed by the diaphragm, usually in people with a large hiatal hernia. They may cause chronic blood loss, iron-deficiency anemia, or visible gastrointestinal bleeding.

\medterm{Campesterol} A plant sterol chemically similar to cholesterol and found in many plant foods. Most people absorb only small amounts, but excessive absorption in the rare disorder sitosterolemia can contribute to abnormal sterol levels and vascular problems.

\medterm{Campho-Phenique Overdose} Poisoning after swallowing or excessive exposure to a topical product containing camphor and other ingredients. It may cause nausea, vomiting, drowsiness, agitation, or seizures and requires poison-control or emergency advice.

\medterm{Camphor} A strong-smelling substance used in some topical products and traditional remedies. Swallowing even a relatively small amount can cause nausea, agitation, drowsiness, or rapidly developing seizures and may be life-threatening.

\medterm{Camphor Overdose} Toxic exposure to too much camphor, causing nausea, vomiting, agitation, drowsiness, seizures, coma, or slowed breathing. Seizures may begin quickly; suspected ingestion requires immediate poison-control or emergency assessment.

\medterm{Campimetry} Measurement and mapping of the area a person can see while looking straight ahead. It helps detect and describe visual-field loss caused by disease of the retina, optic nerve, or brain pathways.

\medterm{Campomelic Dysplasia} A rare inherited disorder caused by changes affecting the SOX9 gene. It causes bowed long bones, poorly developed shoulder blades, characteristic facial features, and often serious problems with breathing or other organs.
\medterm{Camptocormia} Severe forward bending of the trunk that worsens while standing or walking and improves when lying down. It may result from muscle disease, Parkinson disease, dystonia, or another nerve or muscle disorder.

\medterm{Camptodactyly} A fixed bending deformity of one or more fingers, usually at the middle finger joint. It may be present from birth or develop later because of shortened tendons, muscles, connective tissue, or joint structures.

\medterm{Camptothecin} A plant-derived compound that blocks topoisomerase I, an enzyme needed to copy and repair DNA. It is mainly a research compound; related medicines are used in some cancer treatments and can damage rapidly dividing cells.

\medterm{Campylobacter} A group of bacteria that commonly cause foodborne intestinal infection. Illness may include fever, abdominal cramps, diarrhea, or bloody stools; some species can invade the bloodstream or cause later nerve or joint problems.

\medterm{Campylobacter Coli} A Campylobacter species that can cause foodborne intestinal infection, usually after contaminated meat, water, or other food. Symptoms include fever, abdominal cramps, diarrhea, and sometimes bloody stools.
\medterm{Campylobacter Concisus} A Campylobacter species found mainly in the mouth and digestive tract. Its role in disease is uncertain, but it has been associated with occasional intestinal or bloodstream infection, especially in people with weakened immunity.

\medterm{Campylobacter Curvus} A Campylobacter species found in the mouth and digestive tract. It rarely causes human disease but has been reported in infections outside the intestine, particularly when normal body defenses are impaired.

\medterm{Campylobacter Fetus} A Campylobacter species that can cause bloodstream infection and infections of blood vessels, joints, the brain, or other tissues. Serious disease is more likely in older adults and people with weakened immunity.

\medterm{Campylobacter Gracilis} A Campylobacter species found mainly in the mouth and upper airways. It is usually part of the normal microbial community but may contribute to gum disease or rarely cause infection elsewhere.

\medterm{Campylobacter Infection} A bacterial infection, usually caused by C. jejuni or C. coli, acquired from contaminated poultry, food, water, or unpasteurized milk. It causes fever, cramps, and diarrhea that may be bloody; rare complications include bloodstream infection, arthritis, or Guillain--Barré syndrome.

\medterm{Campylobacter Jejuni} A major cause of foodborne intestinal infection, usually acquired from undercooked poultry, contaminated water, or unpasteurized milk. It commonly causes fever, abdominal cramps, and diarrhea that may contain blood.
\medterm{Campylobacter Lari} A Campylobacter species associated with marine environments and poultry. It can rarely cause diarrhea or infection outside the intestine, especially in people with weakened immunity.

\medterm{Campylobacter Rectus} A bacterium found mainly in the mouth and linked with gum disease. It can rarely cause infections of the bloodstream, heart, or other tissues, usually when normal defenses are impaired.

\medterm{Campylobacter Serology Test} A blood test that detects antibodies produced after Campylobacter infection. It may support evaluation of selected complications after infection, but stool molecular testing or culture is generally used to diagnose active intestinal disease.

\medterm{Campylobacter Upsaliensis} A Campylobacter species carried by dogs, cats, and other animals that can occasionally cause human diarrhea or bloodstream infection. Risk is higher in people with weakened immunity.

\medterm{Campylobacteraceae} A family of curved or spiral bacteria that includes Campylobacter and related species. Some members cause intestinal infection, while others live harmlessly in animals or the environment and only rarely infect people.

\medterm{Campylobacteriosis} Infection caused by Campylobacter bacteria, usually producing fever, abdominal pain, and diarrhea that may be bloody. Dehydration is common; rare complications include bloodstream infection, reactive arthritis, and Guillain--Barré syndrome.



\medterm{Canagliflozin} A medicine that lowers blood glucose by causing the kidneys to remove more glucose in urine. It may also reduce some heart-failure and kidney risks, but can cause genital infections, dehydration, low blood pressure, or ketoacidosis.

\medterm{Canal Of Schlemm} A circular drainage channel at the edge of the cornea that carries aqueous fluid from the front of the eye into the bloodstream. Poor drainage can raise eye pressure and contribute to glaucoma.

\medterm{Canaliculus} A small anatomical channel. The tear canaliculi are tiny passages in the eyelids that carry tears from the eye surface into the lacrimal sac and then the nose.
\medterm{Canalith Repositioning Maneuver} A series of guided head and body movements that moves displaced inner-ear crystals out of a semicircular canal. It is commonly used to treat posterior-canal benign paroxysmal positional vertigo.

\medterm{Canalith Repositioning Maneuvers} Head and body movements that move loose crystals from an inner-ear semicircular canal back to the utricle, where they no longer trigger false movement signals. The maneuver is selected according to the affected canal and side.

\medterm{Canalization} Formation or restoration of a channel through tissue, a blood vessel, or an embryonic structure. In a blood clot, canalization means that small channels form through the clot and partly restore blood flow.
\medterm{Canarypox Virus} A poxvirus that mainly infects birds. It does not normally cause human disease but is used as a research and vaccine-delivery vector because it can enter human cells without usually multiplying in them.

\medterm{Canavan Disease} A rare inherited disorder in which deficiency of the ASPA enzyme damages the brain's white matter. It usually begins in infancy with poor muscle tone, delayed development, increasing head size, feeding problems, and later severe movement and communication impairment.

\medterm{Cancell} A shortened or misspelled form of cancellous, meaning the porous inner part of bone that contains marrow and helps spread mechanical forces.

\medterm{Cancellous Bone} Porous bone made of thin supporting plates or struts. It is found inside many bones, especially near joints and in the vertebrae, and contains marrow while helping distribute forces.

\medterm{Cancer} A disease in which abnormal cells grow without normal control and may invade nearby tissues or spread to distant organs. Cancers are classified by where they begin and the cell changes they contain; treatment depends on these features and the extent of disease.


\medterm{Cancer And Lymph Nodes} Cancer cells may travel from a primary tumor to nearby lymph nodes through lymphatic vessels. Examination, imaging, or biopsy of these nodes helps determine the cancer's extent and guide treatment.

\medterm{Cancer Care Team} The health professionals who work together to diagnose, treat, and support a person with cancer. Depending on the illness, the team may include medical and radiation oncologists, surgeons, pathologists, nurses, pharmacists, therapists, and palliative-care specialists.

\medterm{Cancer Causes} Cancer develops after changes in cells allow abnormal growth and survival. Risk can be increased by inherited variants, tobacco, alcohol, certain infections, ultraviolet or other radiation, harmful chemicals, chronic inflammation, and random changes during cell division.

\synonyms
Cancer Risk (Cancer Causes)

\medterm{Cancer Cell} A cell with changes that allow it to grow and survive without normal control. Cancer cells may lose their specialized function, resist cell death, invade nearby tissues, enter blood or lymph, and form tumors elsewhere.

\medterm{Cancer Cluster} A group or area with more cancer cases than expected during a specified time. Investigation compares the number and types of cancers with reliable background rates and examines possible shared exposures while considering chance variation.


\medterm{Cancer Genetics} The study of inherited and acquired genetic changes that affect cancer risk, behavior, or treatment. Testing may identify inherited risk syndromes or changes in a tumor that help guide treatment and family counseling.

\medterm{Cancer Grading} Assessment of how abnormal cancer cells and their tissue pattern look under a microscope. The grade gives information about likely growth behavior and is different from staging, which describes how far the cancer has spread.

\medterm{Cancer Immunotherapy} Cancer treatment that helps the immune system recognize or attack cancer cells. Methods include checkpoint inhibitors, engineered immune cells, immune-stimulating proteins, and some vaccines; use and side effects depend on the cancer and treatment.

\medterm{Cancer Mutations} Genetic changes in cancer cells that can alter growth, survival, or spread. Most develop during life and are not inherited. Some drive cancer development, while others have little effect; testing selected changes can help guide prognosis or targeted treatment.

\medterm{Cancer In Lung (Lung Cancer)} Alternate index label for Lung Cancer.

\medterm{Cancer Nutrition} Nutrition care for people with cancer, who may lose appetite or weight because of the disease or treatment. Individual plans may address calories, protein, fluids, swallowing, nausea, diarrhea, and cachexia while respecting the person's goals and treatment.

\medterm{Cancer Of Unknown Origin} Alternate index label; see \textit{Carcinoma of unknown primary}.

\medterm{Cancer Of Unknown Primary} Cancer that has spread to another site but whose original starting point cannot be identified after appropriate examination and testing. Tissue studies, scans, and selected molecular tests may help guide treatment.

\medterm{Cancer Pain} Pain caused by a tumor, its spread, or its treatment. It may be constant or intermittent and can arise from bone, nerves, organs, surgery, radiation, or medicines; treatment combines control of the cause with suitable pain relief and supportive care.

\synonyms
Pain Cancer (Cancer Pain)

\medterm{Cancer Prevention} Actions that lower the chance of developing cancer or find precancerous changes early. They include avoiding tobacco, limiting alcohol, vaccination against cancer-causing infections, protecting skin from ultraviolet light, reducing harmful exposures, and recommended screening.

\medterm{Cancer Registry} An organized record of cancer diagnoses, including cancer type, location, stage, treatment, and outcome. Registries help measure cancer rates, plan services, improve care, and support research while protecting personal information.

\medterm{Cancer Risk} The likelihood that a person will develop cancer. It is influenced by age, inherited variants, family history, tobacco, alcohol, certain infections, radiation, environmental exposures, body weight, and chance; having a risk factor does not guarantee cancer.

\medterm{Cancer Risk (Cancer Causes)} Alternate index label for Cancer Causes.

\medterm{Cancer Screening} Testing people who have no symptoms to find selected cancers or precancerous changes early, when treatment may work better. Each test has specific age, risk, interval, benefits, and possible harms.

\medterm{Cancer Staging} Describing how far a cancer has grown or spread. Assessment may include the primary tumor, nearby lymph nodes, distant metastases, scans, and tissue examination; stage helps guide treatment and estimate prognosis.

\medterm{Cancer Surveillance} Ongoing monitoring for cancer recurrence, progression, treatment effects, or new cancers after diagnosis or treatment. It may include examinations, imaging, laboratory tests, and symptom review according to the cancer and treatment.

\medterm{Cancer Survivor} A person living with cancer from diagnosis through treatment and afterward. The term includes people whose cancer is controlled or cured and those with ongoing disease; follow-up may address recurrence, late effects, emotional health, and quality of life.

\medterm{Cancer Susceptibility Genes} Genes in which inherited harmful variants increase the chance of developing one or more cancers. Examples include BRCA1, BRCA2, APC, and mismatch-repair genes; carrying a risk variant does not guarantee that cancer will develop.

\medterm{Cancer Treatment - Dealing With Pain} Pain during or after cancer treatment may come from the tumor, surgery, radiation, chemotherapy, or nerve injury. Care may combine treatment of the cause, pain medicines, physical rehabilitation, psychological support, and other supportive measures.

\medterm{Cancer Treatment - Early Menopause} Cancer treatment may stop menstrual periods earlier than expected, especially after chemotherapy, pelvic radiation, ovarian surgery, or hormone-blocking treatment. It can cause hot flashes, vaginal symptoms, infertility, and bone loss and requires individualized follow-up.


\medterm{Cancer Treatment: Dealing With Hot Flashes And Night Sweats} Cancer treatment or hormone changes can cause sudden heat, sweating, and disturbed sleep. Cooling measures, layered clothing, sleep habits, behavioral approaches, and selected medicines may help, depending on the cancer and other health risks.


\medterm{Cancer Treatments} Cancer treatment may include surgery, radiation, chemotherapy, immunotherapy, targeted medicines, hormone therapy, stem-cell transplantation, or supportive and palliative care. The choice depends on cancer type, location, stage, molecular features, overall health, and patient goals.

\medterm{Cancer, Carcinoid Tumors} Alternate index label; see \textit{Carcinoid tumors}.

\medterm{Cancer, Chronic Lymphocytic Leukemia} Alternate index label; see \textit{Chronic lymphocytic leukemia}.

\medterm{Cancer, Chronic Myelogenous Leukemia} Alternate index label; see \textit{Chronic myelogenous leukemia}.

\medterm{Cancer, Colon} Alternate index label; see \textit{Colon cancer}.

\medterm{Cancer, Endometrial} Alternate index label; see \textit{Endometrial cancer}.

\medterm{Cancer, Esophageal} Alternate index label; see \textit{Esophageal cancer}.

\medterm{Cancer, Eye Melanoma} Alternate index label; see \textit{Eye melanoma}.

\medterm{Cancer, Floor Of The Mouth} Alternate index label; see \textit{Floor of the mouth cancer}.

\medterm{Cancer, Gallbladder} Alternate index label; see \textit{Gallbladder cancer}.

\medterm{Cancer, Gastric} Alternate index label; see \textit{Stomach cancer}.

\medterm{Cancer, Hairy Cell Leukemia} Alternate index label; see \textit{Hairy cell leukemia}.

\medterm{Cancer, Hodgkin Lymphoma} Alternate index label; see \textit{Hodgkin lymphoma (Hodgkin disease)}.

\medterm{Cancer, Hurthle Cell} Alternate index label; see \textit{Hurthle cell cancer}.

\medterm{Cancer, Inflammatory Breast} Alternate index label; see \textit{Inflammatory breast cancer}.

\medterm{Cancer, Kidney} Alternate index label; see \textit{Kidney cancer}.

\medterm{Cancer, Liver} Alternate index label; see \textit{Liver cancer}.

\medterm{Cancer, Lung} Alternate index label; see \textit{Lung cancer}.

\medterm{Cancer, Merkel Cell} Alternate index label; see \textit{Merkel cell carcinoma}.

\medterm{Cancer, Mouth} Alternate index label; see \textit{Mouth cancer}.

\medterm{Cancer, Multiple Myeloma} Alternate index label; see \textit{Multiple myeloma}.

\medterm{Cancer, Nasopharyngeal} Alternate index label; see \textit{Nasopharyngeal carcinoma}.

\medterm{Cancer, Neuroblastoma} Alternate index label; see \textit{Neuroblastoma}.

\medterm{Cancer, Neuroendocrine} Alternate index label; see \textit{Neuroendocrine tumors}.

\medterm{Cancer, Non-Hodgkin Lymphoma} Alternate index label; see \textit{Non-Hodgkin lymphoma}.

\medterm{Cancer, Ovarian} Alternate index label; see \textit{Ovarian cancer}.

\medterm{Cancer, Pancreatic} Alternate index label; see \textit{Pancreatic cancer}.

\medterm{Cancer, Paraneoplastic Syndromes} Alternate index label; see \textit{Paraneoplastic syndromes of the nervous system}.

\medterm{Cancer, Prostate} Alternate index label; see \textit{Prostate cancer}.

\medterm{Cancer, Rectal} Alternate index label; see \textit{Rectal cancer}.

\medterm{Cancer, Retinoblastoma} Alternate index label; see \textit{Retinoblastoma}.

\medterm{Cancer, Skin} Alternate index label; see \textit{Skin cancer}.

\medterm{Cancer, Soft Tissue Sarcoma} Alternate index label; see \textit{Soft tissue sarcoma}.

\medterm{Cancer, Squamous Cell} Alternate index label; see \textit{Squamous cell carcinoma of the skin}.

\medterm{Cancer, Stomach} Alternate index label; see \textit{Stomach cancer}.

\medterm{Cancer, Testicular} Alternate index label; see \textit{Testicular cancer}.

\medterm{Cancer, Throat} Alternate index label; see \textit{Throat cancer}.

\medterm{Cancer, Thyroid} Alternate index label; see \textit{Thyroid cancer}.

\medterm{Cancer, Uterine} Alternate index label; see \textit{Endometrial cancer}.

\medterm{Cancer, Vagina} Alternate index label; see \textit{Vaginal cancer}.

\medterm{Cancer, Vulvar} Alternate index label; see \textit{Vulvar cancer}.

\medterm{Cancer, Wilms' Tumor} Alternate index label; see \textit{Wilms tumor}.

\medterm{Cancer-Associated Thrombosis} A blood clot in a vein or artery occurring with active cancer or its treatment. Cancer, surgery, immobility, central lines, and some treatments increase risk; anticoagulant treatment is chosen according to clot location, bleeding risk, and cancer status.

\medterm{Cancrum Oris} Noma, a rapidly destructive gangrenous infection of the mouth and face that occurs
mainly in severely malnourished or immunocompromised children. It begins with gingival
or mucosal necrosis and can extend through facial tissues.

\medterm{Candela} The international unit used to measure the brightness of a light source in a specified direction. It is a physical unit, not a measure of a medicine's light dose or the amount of light reaching tissue.

\medterm{Candida} A group of yeasts that may normally live on the skin, in the mouth, or in the digestive and genital tracts. Overgrowth can cause mouth, genital, skin, or, in vulnerable people, bloodstream and organ infections.

\medterm{Candida Albicans} The Candida species most often causing human infection. It can cause oral thrush, vaginal yeast infection, moist skin-fold rash, or serious bloodstream and organ infection in people with weakened defenses.

\medterm{Candida Albicans Infection} Infection caused by Candida albicans. It may be limited to the mouth, vagina, or skin, or may invade the bloodstream and organs, especially in people who are hospitalized or have weakened immunity.

\medterm{Candida Auris} A yeast that can live on skin without symptoms or cause bloodstream, wound, or ear infection. It often resists several antifungal medicines, spreads readily in health-care settings, and requires specialized testing and strict infection-control measures.

\medterm{Candida Auris Infection} Infection or symptom-free colonization caused by Candida auris. Invasive infection may affect the blood or organs, especially in hospitalized or immunocompromised people; treatment is guided by susceptibility testing because resistance is common.

\medterm{Candida Boidinii} A yeast found mainly in the environment and occasionally in people. Human infection is uncommon but may occur opportunistically, particularly in people with severe illness, invasive devices, or weakened immunity.

\medterm{Candida Dubliniensis} A yeast related to Candida albicans that can live in the mouth and cause oral thrush or, rarely, invasive infection. It is more likely to cause disease when immune defenses are impaired.

\medterm{Candida Esophagitis} Yeast infection of the esophagus, usually caused by Candida. It commonly causes painful or difficult swallowing and chest discomfort, especially in people with weakened immunity or after certain medicines.

\medterm{Candida Famata} A yeast found in the environment and occasionally on people. It rarely causes bloodstream or other invasive infection, mainly in hospitalized people with weakened immunity or indwelling medical devices.

\medterm{Candida Glabrata} A Candida species that can cause genital, urinary, bloodstream, or other invasive infection. It may resist some commonly used antifungal medicines, so laboratory susceptibility testing can help guide treatment.

\medterm{Candida Inconspicua} A yeast that is an uncommon cause of human infection. It may cause bloodstream or other invasive disease mainly in hospitalized people with weakened immunity, and laboratory identification may be needed for treatment selection.


\medterm{Candida Infection Of The Skin} A superficial yeast infection, often in warm, moist skin folds. It causes a red, itchy, sore rash that may have small separate spots around its edge and is more common with moisture, diabetes, obesity, or immune suppression.

\medterm{Candida Intermedia} A yeast found in the environment and occasionally in people. It rarely causes infection but can produce bloodstream or other invasive disease in hospitalized people with weakened immunity.

\medterm{Candida Orthopsilosis} A yeast related to Candida parapsilosis that can rarely cause bloodstream or device-associated infection, particularly in hospitalized people, newborns, or those with weakened immunity.

\medterm{Candida Parapsilosis} A Candida species that can cause bloodstream and catheter-related infection, especially in newborns, hospitalized people, and those receiving intravenous treatment. It may spread between patients through hands or contaminated equipment.
\medterm{Candida Sake} A yeast found mainly in the environment and food. Human infection is rare but may occur opportunistically, especially in people with severe illness or weakened immune defenses.

\medterm{Candida Tropicalis} A Candida species that can cause mouth or genital infection and may invade the bloodstream or organs. Serious infection is more likely in people with neutropenia, cancer, or other immune impairment.

\medterm{Candida Zeylanoides} A yeast found mainly in the environment. It is an uncommon cause of human disease but may rarely cause bloodstream or other infection in people with severe illness or weakened immunity.

\medterm{Candidal Intertrigo} A Candida yeast infection in warm, moist skin folds such as the armpits, groin, under the breasts, or beneath an abdominal fold. It causes a red, moist, itchy or burning rash, sometimes with small spots at the edge.

\medterm{Candidal Vaginitis} Inflammation of the vagina and vulva caused by overgrowth of Candida yeast. It usually causes intense itching, soreness, redness, pain with sex or urination, and a thick white discharge without a strong odor.

\medterm{Candidal Vulvovaginitis} Alternate index label; see \textit{Vulvovaginal Candidiasis}.

\medterm{Candidemia} Candida yeast infection of the bloodstream. It can spread to organs and is more common with central lines, surgery, intensive care, recent antibiotics, or weakened immunity; it requires prompt medical treatment and evaluation for spread.

\medterm{Candidiasis} Infection caused by Candida yeast. It may affect the mouth, esophagus, vagina, or moist skin folds, or may invade the blood and organs in seriously ill people. Symptoms and treatment depend on the site and whether the infection is invasive.

\medterm{Candidiasis, Oral} Alternate index label; see \textit{Oral thrush}.

\medterm{Candidiasis, Vaginal} Alternate index label; see \textit{Yeast infection (vaginal)}.

\medterm{Candles Poisoning} Illness after swallowing candle wax or exposure to candle products. Small amounts of plain wax usually cause mild stomach upset, but scented or decorative products may contain harmful ingredients; choking or breathing problems require urgent care.

\medterm{Cane} A walking aid used to improve balance and reduce some weight on a painful or weak leg. It should be fitted to the person's height and used on the side opposite the weaker or painful leg unless a clinician advises otherwise.

\medterm{Canertinib Dihydrochloride} An experimental anticancer compound that blocks several epidermal growth-factor receptors, proteins that can promote tumor-cell growth. It was studied in clinical research but is not a routine approved cancer treatment.

\medterm{Canfosfamide Hcl} An experimental anticancer prodrug designed to become active after processing by an enzyme found at increased levels in some tumors. It was studied in cancer research but is not a routine approved treatment.

\medterm{Canicola Fever} A form of leptospirosis caused by a strain of Leptospira bacteria commonly carried by dogs. It may cause fever, headache, muscle aches, red eyes, kidney injury, or liver involvement after contact with infected urine or contaminated water.

\medterm{Canine Teeth} The pointed teeth between the incisors and premolars. They help tear food and guide the jaws during biting and chewing.
\medterm{Canker} A nontechnical word often used for an aphthous mouth ulcer, a painful shallow sore on the moist lining of the mouth. The word can have other meanings, so the site and appearance should be clarified.
\medterm{Canker Sore} A painful, shallow ulcer on the moist lining of the mouth, usually round or oval with a red border. It is not contagious and commonly heals within one to two weeks; persistent, severe, or unusual sores need assessment.

\synonyms
Recurrent Oral Ulcers
Ulcer, Aphthous

\medterm{Canker Sores} Painful, shallow, noncontagious ulcers on the moist lining of the mouth. Minor injury, stress, nutritional deficiency, immune disease, medicines, or other triggers may contribute; frequent, severe, or nonhealing sores need medical or dental evaluation.

\synonyms
Sores Canker (Canker Sores)

\medterm{Cannabidiol} A cannabinoid from the cannabis plant that usually does not produce the “high” caused by THC. A purified form treats selected severe seizure disorders; other products vary in strength, purity, side effects, and medicine interactions.

\medterm{Cannabinoid} A substance that acts on cannabinoid receptors or related signaling systems. Cannabinoids may come from plants, the body, or laboratories; effects can include pain relief, appetite changes, sedation, impaired memory, or intoxication.

\medterm{Cannabinoid Hyperemesis Syndrome} Repeated episodes of severe nausea, vomiting, and abdominal pain associated with long-term frequent cannabis use. Hot showers may temporarily ease symptoms, but lasting improvement usually requires stopping cannabis.

\medterm{Cannabinoid Receptor} A cell-surface protein, mainly CB1 or CB2, that responds to cannabinoids made by the body or found in cannabis and medicines. It helps regulate pain, appetite, mood, memory, immune activity, and inflammation.

\medterm{Cannabinol} A cannabinoid formed as THC breaks down with age or exposure to air. It may have mild psychoactive effects, but its medical benefits and safety are less well established than those of some other cannabinoids.

\medterm{Cannabis} A plant containing cannabinoids such as THC and CBD. Use can alter perception, memory, coordination, mood, and heart rate; it may have selected medical uses but can also cause dependence, anxiety, psychosis, or other harms.
\medterm{Cannabis Use Disorder} A pattern of cannabis use that causes significant difficulty controlling use or problems at home, work, school, relationships, or health. Tolerance and withdrawal may occur, with irritability, anxiety, sleep problems, or reduced appetite after stopping.

\medterm{Cannot Intubate Cannot Oxygenate} A life-threatening airway emergency in which clinicians cannot place a breathing tube and cannot provide enough oxygen by mask or another airway device. It requires immediate emergency action, often including a surgical airway.

\medterm{Cannula} A thin hollow tube placed in a blood vessel, body cavity, or tissue to deliver medicines or fluids, remove material, measure pressure, provide oxygen, or allow access for another device.

\medterm{Cannula Hub} The connector at the outer end of a cannula that attaches it to tubing, a syringe, or another device. It may include a port or valve for giving fluids, taking samples, or flushing the cannula.

\medterm{Cannulation} Placement of a cannula into a blood vessel, body cavity, duct, or other passage. It provides access for giving fluids or medicines, removing material, monitoring, collecting samples, or performing a procedure.

\medterm{Canola Oil} A vegetable oil relatively low in saturated fat and containing mostly monounsaturated and polyunsaturated fats. Replacing solid fats with canola oil can support a heart-healthy eating pattern when total calories remain appropriate.

\medterm{Canrenone} A metabolite related to spironolactone that blocks aldosterone's action. It promotes loss of sodium and water while reducing potassium loss and has been used in some countries for conditions involving fluid retention or excess aldosterone.

\medterm{Cantharides} Preparations made from blister beetles that contain cantharidin, a powerful irritant. Swallowing or applying it can cause severe mouth, stomach, kidney, and urinary-tract injury and may be life-threatening.
\medterm{Canthaxanthin} A red-orange pigment used as a food color and formerly in tanning products. Excessive intake can cause yellow skin, liver injury, and crystal deposits in the retina that impair vision.

\medterm{Canthotomy} A surgical cut through the outer corner of the eyelids, usually performed urgently to relieve dangerous pressure around the eye and protect vision after severe orbital swelling or injury.

\medterm{Canthus} The corner where the upper and lower eyelids meet. The inner canthus is beside the nose, and the outer canthus is toward the temple.

\medterm{Cantlie Line} An imaginary line across the liver from the gallbladder to the large vein behind the liver. It follows the middle hepatic vein and marks the functional division between the right and left liver, including their blood supply and bile drainage.

\medterm{Cantuzumab Ravtansine} An experimental antibody-drug conjugate designed to attach to MUC1, a protein found on some cancer cells, and deliver a cell-killing drug. It has been studied in trials but is not a routine approved treatment.

\medterm{Capacity} A person's ability, at a particular time, to understand information about a medical decision, appreciate how it applies to them, weigh the options, and communicate a choice.

\medterm{Capacity And Competence} Capacity is the clinical ability to make a specific decision by understanding, using, and communicating relevant information. Competence is a legal determination; the two terms are related but not interchangeable.

\medterm{Capecitabine} An oral chemotherapy medicine that the body converts to fluorouracil, which interferes with cancer-cell DNA production. It is used for selected cancers and can cause diarrhea, mouth sores, hand-foot syndrome, low blood counts, or heart problems.

\medterm{Capgras Syndrome} A fixed false belief that someone familiar has been replaced by an identical impostor. It may occur with psychotic illness, dementia, brain injury, seizures, or other neurologic disorders and needs clinical assessment.

\medterm{Capillaria} A group of parasitic roundworms. Different species can cause intestinal infection, liver disease, or infection of other tissues; symptoms and treatment depend on the species and affected organ.

\medterm{Capillaria Philippinensis} A parasitic roundworm acquired by eating raw or undercooked freshwater fish. It infects the intestine and can cause persistent watery diarrhea, abdominal pain, weight loss, poor absorption of nutrients, and loss of blood proteins.

\medterm{Capillaries} The body's smallest blood vessels. Their thin walls allow oxygen, nutrients, hormones, fluid, and waste products to move between the blood and nearby tissues.
\medterm{Capillaroscopy} Examination of tiny blood vessels, usually in the skin beside a fingernail, with magnification. It can show widened or missing vessels, bleeding, and abnormal patterns that support evaluation of connective-tissue and small-vessel disease.

\medterm{Capillary} The smallest type of blood vessel, forming networks that connect arterioles with venules. Its thin wall allows exchange of oxygen, nutrients, fluid, and waste between blood and tissues.

\medterm{Capillary Bed} A network of capillaries within a tissue that connects small arteries to small veins. It provides a large surface for exchange of oxygen, carbon dioxide, nutrients, hormones, fluid, and waste.

\medterm{Capillary Electrophoresis} A laboratory method that separates charged molecules in a very thin tube using an electric field. It can analyze proteins, medicines, DNA, or other substances in a sample, depending on the test.
\medterm{Capillary Leak Syndrome} A condition in which fluid and proteins escape from small blood vessels into surrounding tissues. It can cause sudden swelling, low blood pressure, thickened blood, shock, and organ failure; severe infection, medicines, immune disease, or rare unexplained episodes may trigger it.

\medterm{Capillary Microscope} A microscope used to magnify tiny blood vessels, especially those beside the fingernails. It helps assess vessel shape, size, bleeding, and areas without vessels in some circulation and connective-tissue disorders.
\medterm{Capillary Nail Refill Test} A bedside check of circulation in which a nail bed is pressed until it turns pale and the time for color to return is observed. Temperature, lighting, age, skin color, shock, and vascular disease affect the result, so it is not used alone.

\medterm{Capillary Plasma} The liquid component of blood obtained from small vessels in skin or tissue, containing water, electrolytes, proteins, hormones, nutrients, and waste products.

\medterm{Capillary Refill} A bedside assessment of peripheral perfusion performed by briefly pressing a nail bed or
skin and measuring the time for normal color to return. Delayed refill may indicate poor
circulation, shock, hypothermia, or vascular disease.

\medterm{Capillary Refill Time} The time required for color to return after pressure briefly blanches a nail bed or area of skin. A delayed
result may suggest impaired peripheral perfusion, but interpretation depends on age,
temperature, technique, and the clinical setting.

\medterm{Capillary Sample} A small blood specimen obtained by puncturing the skin, usually a finger or heel, for bedside glucose, hemoglobin, newborn screening, or other point-of-care tests.

\medterm{Capillary Vessel} A microscopic vessel with a thin endothelial wall where blood exchanges oxygen, carbon dioxide, nutrients, fluid, and waste with surrounding tissue.

\medterm{Capitate} The largest wrist bone, located centrally in the distal row of carpal bones. It connects with several other wrist bones and the third metacarpal, helps transmit force through the wrist, and can rarely develop blood-supply loss after fracture.

\medterm{Capitation} A health-care payment arrangement in which a clinician or organization receives a fixed amount per enrolled person for specified services over a set period, whether or not that person uses the services.
\medterm{Caplan'S Syndrome} A lung disorder combining rheumatoid arthritis with pneumoconiosis caused by long-term inhalation of coal or mineral dust. It produces multiple lung nodules and may cause cough or breathing difficulty.
\medterm{Capnocytophaga} A group of bacteria found mainly in the mouths of dogs, cats, and people. Some species can cause serious human infection after animal bites or scratches, especially in people without a functioning spleen or with weakened immunity.

\medterm{Capnocytophaga Canimorsus} A bacterium commonly found in the mouths of dogs and cats. It can cause rapidly severe bloodstream infection, meningitis, or skin infection after a bite or scratch, especially in older adults or people without a functioning spleen or with weakened immunity.

\medterm{Capnocytophaga Gingivalis} A bacterium that can live in the human mouth and has been associated with gum disease. It rarely causes infection outside the mouth, mainly in people whose immune defenses are impaired.

\medterm{Capnocytophaga Infection} Infection caused by Capnocytophaga bacteria, often after a dog or cat bite, scratch, or contact with animal saliva. It may cause a wound infection, fever, bloodstream infection, or meningitis, with greatest risk in people with impaired immunity or no spleen.
\medterm{Capnocytophaga Ochracea} A bacterium found mainly in the human mouth and associated with dental plaque and gum disease. Infection outside the mouth is uncommon but can occur when immune defenses are severely weakened.

\medterm{Capnocytophaga Sputigena} A bacterium that may live in the mouth and contribute to gum disease. It can rarely cause infection of the bloodstream or other tissues, usually in people with serious illness or weakened immunity.

\medterm{Capnography} Continuous measurement and display of carbon dioxide in exhaled air. The waveform and end-tidal value help assess breathing, airway blockage, circulation, and correct placement of a breathing tube during anesthesia, sedation, or emergency care.

\medterm{Capreomycin} An injectable antibiotic formerly used for drug-resistant tuberculosis. It can damage the kidneys, hearing, or balance and is now used only in selected specialist-directed situations because safer options may be available.

\medterm{Capromab Pendetide} A radioactive antibody formerly used to image prostate-cancer tissue by binding a protein on prostate cells. Its clinical use is limited and depends on local approval and availability of more accurate imaging methods.

\medterm{Caps} A prescription abbreviation for capsules, oral medicines enclosed in a shell that dissolves in the digestive tract. Abbreviations can be misread, so complete dosage-form wording is safer on prescriptions.

\medterm{Capsaicin} The substance that makes chili peppers feel hot. Applied in controlled amounts to the skin, it can reduce some nerve and joint pains by temporarily lowering pain signaling; it commonly causes initial burning or redness.

\medterm{Capsid} The protein shell surrounding a virus's genetic material. It protects the genome and helps the virus attach to and enter host cells; its structure influences which cells can be infected and how the immune system responds.

\medterm{Capsid Proteins} Proteins that assemble around a virus's genetic material to form its protective shell. Some also help the virus attach to and enter cells and are recognized by the immune system.

\medterm{Capsular Contracture} Abnormal tightening and hardening of scar tissue around an implanted device, especially a breast implant. It may make the area firm, painful, distorted, or displaced and sometimes requires corrective surgery.

\medterm{Capsule} A protective covering around an organ, joint, or other structure, or an oral medicine enclosed in a dissolvable shell. The exact meaning depends on whether the term is used in anatomy or pharmacology.

\medterm{Capsule Endoscope} A swallowable capsule containing a tiny camera and transmitter. As it travels through the digestive tract, it sends images, especially of the small intestine, where conventional endoscopes may not reach easily.

\medterm{Capsule Endoscopy} A test in which a person swallows a small wireless camera that takes pictures of the digestive tract, especially the small intestine. The capsule usually passes naturally but can become stuck where the bowel is narrowed.

\medterm{Capsule Of Glisson} The layer of connective tissue surrounding the liver. It extends inward around the liver's blood vessels, bile ducts, and nerves and helps support and protect these structures.

\medterm{Capsule Of Tenon} A thin connective-tissue sheath surrounding the eyeball and separating it from the fat of the eye socket. It allows smooth eye movement and creates a space used for some eye injections and operations.

\medterm{Capsulectomy} Surgical removal of some or all of the scar-tissue capsule around an organ, implant, or other structure. It may be performed for painful tightening, infection, abnormal tissue, or revision of a reconstruction.

\medterm{Captopril} An angiotensin-converting enzyme inhibitor used for high blood pressure, heart failure, and some kidney problems. It can cause cough, high potassium, low blood pressure, or swelling.

\medterm{Captopril Renal Scan} A kidney imaging test using a radioactive tracer before and after captopril. It compares kidney blood flow and function and may help evaluate high blood pressure caused by narrowing of a kidney artery.
\medterm{Caput Medusae} Enlarged veins spreading out from the navel, caused by increased pressure in the portal vein system, usually from liver disease. They are a physical sign of abdominal collateral circulation rather than a disease by themselves.

\medterm{Caput Succedaneum} Soft swelling of a newborn's scalp caused by pressure during labor. It can cross the skull's suture lines and usually disappears within hours to a few days; rapidly enlarging swelling or pallor requires urgent assessment.

\medterm{CAR T-Cell} A T lymphocyte genetically changed to carry a receptor that recognizes a specific marker on cancer cells. These cells can be used as treatment for selected blood cancers but may cause severe immune reactions or nervous-system effects.

\medterm{Chimeric Antigen Receptor T-Cell Therapy} A cancer treatment in which a person's T cells are collected, genetically changed in a laboratory to recognize a cancer marker, multiplied, and returned to the body. Treatment can cause serious inflammation, infection risk, or neurologic effects.

\medterm{Carbachol} A medicine that mimics acetylcholine and stimulates muscarinic receptors. In selected eye procedures, it constricts the pupil and can lower pressure inside the eye; effects may include sweating, slow heartbeat, or breathing problems.

\medterm{Carbadox} An antimicrobial formerly used to promote growth and prevent disease in food-producing animals. Concerns about toxic breakdown products and possible cancer risk have led to restrictions or bans in many places.

\medterm{Carbamate} A chemical group found in some medicines and pesticides. Carbamate insecticides can temporarily block acetylcholinesterase, causing excess nerve activity with sweating, salivation, vomiting, weakness, wheezing, or seizures.

\medterm{Carbamazepine} A medicine used for focal seizures, trigeminal neuralgia, and selected bipolar disorder. Important risks include severe skin reactions, low blood counts, liver injury, low sodium, dizziness, and interactions with many other medicines.

\medterm{Carbamide} Another name for urea, a waste product made when the body breaks down protein and removed mainly by the kidneys. Urea is also used in some skin creams to soften and hydrate thick or dry skin.
\medterm{Carbapenem} A group of powerful beta-lactam antibiotics that kill bacteria by blocking construction of their cell wall. They are reserved mainly for serious infections because resistance and important side effects can limit treatment choices.

\medterm{Carbapenem Resistant Enterobacteriaceae (CRE Infection)} Alternate index label for CRE Infection.

\medterm{Carbapenem-Resistant Enterobacteriaceae Infection} Infection caused by Enterobacterales bacteria that resist carbapenem antibiotics. It may affect the urinary tract, lungs, blood, or wounds, can spread in health-care settings, and requires culture and susceptibility testing to select treatment.

\medterm{Carbapenem-Resistant Organisms} Bacteria that cannot be reliably treated with carbapenem antibiotics, often because they produce enzymes that destroy these medicines. They can cause difficult infections and spread in health-care settings, where strict infection control is important.

\medterm{Carbaryl} A carbamate insecticide that temporarily blocks acetylcholinesterase. Substantial exposure can cause sweating, excess saliva, vomiting, diarrhea, small pupils, muscle weakness, wheezing, seizures, or breathing failure.

\medterm{Carbazilquinone} An experimental quinone-containing compound studied for its ability to damage cancer cells, especially under low-oxygen conditions. It is a research substance and is not a routine clinical medicine.

\medterm{Carbenicillin} A penicillin antibiotic formerly used for selected serious infections caused by gram-negative bacteria. Resistance, allergic reactions, diarrhea, and kidney or electrolyte problems can limit its use.

\medterm{Carbidopa} A medicine given with levodopa for Parkinson disease. It blocks conversion of levodopa to dopamine outside the brain, allowing more levodopa to reach the brain and reducing nausea, vomiting, and low blood pressure.

\medterm{Carbimazole} An antithyroid medicine that the body converts to methimazole. It reduces production of thyroid hormones and can rarely cause dangerously low white-cell counts or liver injury.
\medterm{Carbofuran} A highly toxic carbamate pesticide that blocks acetylcholinesterase. Ingestion, inhalation, or skin exposure can rapidly cause sweating, vomiting, muscle weakness, seizures, breathing failure, and death.

\medterm{Carbogen} A controlled breathing mixture containing oxygen and a small amount of carbon dioxide. It has limited specialized or experimental uses to stimulate breathing or alter tissue oxygenation and is not for unsupervised use.

\medterm{Carbohydrate} A nutrient found in sugars, starches, fruits, grains, vegetables, and milk. The body breaks digestible carbohydrates into glucose for energy; fiber is a carbohydrate that is not fully digested.

\medterm{Carbohydrate Counting} A way to plan meals by estimating the grams of carbohydrate in food. People using insulin may match their insulin dose to the amount eaten; accurate reading of food labels and advice from a diabetes team are important.

\medterm{Carbohydrate Metabolism} The processes cells use to break down, store, and make sugars. These processes provide energy and help keep blood glucose stable; problems can occur in diabetes and inherited disorders of sugar storage.

\medterm{Carbohydrates} Sugars, starches, and fiber that form one of the main types of nutrients in food. Digestible carbohydrates are broken down into glucose for energy, while fiber is not fully digested and helps support bowel function and blood-sugar control.

\medterm{Carbolic Acid (Phenol)} Phenol, a chemical once used as a surgical antiseptic and still found in some industrial products. It kills microorganisms by damaging proteins and cell membranes but can severely burn skin, eyes, and internal organs.

\synonyms
Phenol

\medterm{Carbolic Acid Poisoning} Exposure to phenol, a corrosive and systemically toxic chemical that can cause burns, altered consciousness, metabolic acidosis, and organ injury.

\medterm{Carbon} A chemical element forming the backbone of organic molecules such as carbohydrates, fats, proteins, and nucleic acids.

\medterm{Carbon Dioxide} A gas produced when cells use energy and removed mainly by the lungs during exhalation. The body controls its blood level through breathing; excess can cause headache, confusion, drowsiness, or respiratory failure.
\medterm{Carbon Dioxide Absorber} A device in some anaesthesia breathing systems that removes carbon dioxide from exhaled air. It allows selected gases to be reused while helping prevent carbon-dioxide build-up during controlled breathing.

\medterm{Carbon Dioxide Transport} The movement of carbon dioxide from body tissues to the lungs. Blood carries it dissolved in plasma, attached to haemoglobin, and mainly as bicarbonate; the lungs remove it during exhalation.

\medterm{Carbon Ion Therapy} A form of particle radiation treatment that uses carbon ions to damage tumor DNA. It may deliver concentrated energy and has higher biological effectiveness than conventional X-rays for some tumors, but availability is limited.

\medterm{Carbon Monoxide} A colorless, odorless gas that binds hemoglobin and prevents blood from delivering enough oxygen. Poisoning may cause headache, dizziness, nausea, confusion, collapse, or death and requires immediate fresh air and emergency care.
\medterm{Carbon Monoxide Poisoning} Poisoning from breathing carbon monoxide, which prevents blood from carrying and delivering oxygen. It may cause headache, dizziness, nausea, confusion, fainting, seizures, or death; leave the area and obtain emergency care immediately.

\medterm{Carbon Tetrachloride} A toxic industrial solvent that can damage the liver and kidneys when swallowed or inhaled. Suspected exposure requires urgent contact with poison-control or emergency services.

\medterm{Carbonate} A chemical form of carbon dioxide found in minerals and body fluids. Carbonate salts are used in some antacids and supplements; their effects, absorption, and risks depend on the particular salt.

\medterm{Carbonic Anhydrase Inhibitors} Medicines that block carbonic anhydrase, an enzyme involved in fluid and acid balance. They reduce fluid production in the eye to lower eye pressure and have selected uses for swelling, seizures, or metabolic disorders.

\medterm{Carbophenothion} An organophosphate insecticide that blocks acetylcholinesterase. Significant exposure can cause sweating, excess saliva, vomiting, muscle weakness, confusion, seizures, breathing failure, and death.

\medterm{Carboplatin} A platinum chemotherapy medicine that damages cancer-cell DNA and prevents it from being copied. It can lower blood counts, especially platelets, and may cause nausea, anemia, kidney effects, nerve symptoms, or allergic reactions.

\medterm{Carboprost} A prostaglandin medicine that strongly contracts the uterus. It is used in selected cases of severe bleeding after childbirth and for some pregnancy-related procedures; it can cause diarrhea, vomiting, fever, and dangerous airway tightening in people with asthma.

\medterm{Carbosulfan} A carbamate insecticide that blocks acetylcholinesterase and overstimulates nerves. Exposure can cause sweating, salivation, vomiting, diarrhea, small pupils, muscle weakness, wheezing, seizures, or breathing failure.

\medterm{Carboxy-Lyases} Enzymes that break a carbon-carbon bond and release carbon dioxide. They participate in many biochemical reactions, including processes that make or break down amino acids and other molecules.

\medterm{Carboxyhaemoglobin} Hemoglobin bound to carbon monoxide. This reduces the blood's ability to carry and release oxygen; an increased level supports recent carbon-monoxide exposure but must be interpreted with symptoms and timing.
\medterm{Carboxyhemoglobin} A compound formed when carbon monoxide binds to hemoglobin. It prevents hemoglobin from carrying and releasing oxygen normally, causing headache, dizziness, confusion, collapse, or death in severe poisoning.

\medterm{Carboxyhemoglobinemia} An abnormally high level of carboxyhemoglobin in the blood after carbon-monoxide exposure. It reduces oxygen delivery and may cause headache, nausea, confusion, fainting, heart injury, brain injury, or death.

\medterm{Carboxyl Group} A chemical group containing carbon, oxygen, and hydrogen that gives many molecules acidic properties. It is found in amino acids, fatty acids, and numerous medicines and metabolic compounds.
\medterm{Carboxylesterase} An enzyme that breaks ester bonds in medicines, toxins, and natural fats. Differences in its activity can change how quickly some substances are broken down and can affect their effects, interactions, or toxicity.

\medterm{Carboxylic Acid} An organic compound containing a carboxyl group, which gives it acidic properties. Many important fatty acids, amino-acid products, and substances involved in metabolism are carboxylic acids.

\medterm{Carboxypeptidase H} An enzyme inside cells that helps convert precursor proteins into active peptide hormones and nerve-signaling molecules by removing amino acids from one end.

\medterm{Carbuncle} A group of connected boils, usually caused by Staphylococcus aureus, forming a painful swollen area with several pus-draining openings. Fever may occur, and risk is increased by diabetes or a weakened immune system.

\medterm{Carbuncles And Boils} Alternate index label; see \textit{Boils and carbuncles}.

\medterm{Carcinoembryonal Antigen} Another name for carcinoembryonic antigen, or CEA, a blood marker that may rise in colorectal and other cancers. It can also increase for noncancerous reasons and is not suitable for general cancer screening.

\medterm{Carcinoembryonic Antigen} A protein measured in blood that may be increased in colorectal and some other cancers. It can also rise with smoking, liver disease, bowel inflammation, or other noncancerous conditions, so it is mainly used to monitor selected known cancers rather than to screen for cancer.
\medterm{Carcinogen} A substance, form of radiation, or infection that can cause cancer by damaging cells or changing how they grow. Risk depends on the type and amount of exposure.
\medterm{Carcinogen Testing} Laboratory, animal, or population research used to determine whether a chemical, radiation exposure, infection, medicine, or other agent can cause or promote cancer.

\medterm{Carcinogenesis} The multistep process by which normal cells acquire changes that allow uncontrolled growth and become cancer. Changes may arise from inherited traits, random errors, infections, tobacco, radiation, chemicals, or other exposures and often accumulate over time.

\medterm{Carcinogenic} Capable of causing or promoting cancer by damaging genetic material or altering how cells grow, divide, repair themselves, or avoid cell death.
\medterm{Carcinoid} A usually slow-growing neuroendocrine tumor that can arise in the digestive tract, appendix, or lungs. Some release hormones or other substances that cause flushing, diarrhea, wheezing, or heart-valve disease.
\medterm{Carcinoid Heart Disease} Damage to the inner lining and valves of the heart caused by hormone-like substances released by a neuroendocrine tumor. It mainly affects the tricuspid and pulmonary valves and can lead to right-sided heart failure.

\medterm{Carcinoid Syndrome} A group of symptoms caused by hormone-like substances released by some neuroendocrine tumors, especially after spread to the liver. It may cause facial flushing, diarrhea, wheezing, rapid heartbeat, and eventually heart-valve disease.

\medterm{Carcinoid Tumor} A usually well-differentiated neuroendocrine tumor arising in hormone-producing cells, commonly in the appendix, small intestine, rectum, or lungs. Its behavior ranges from slow-growing to capable of spreading and causing hormone-related symptoms.

\medterm{Carcinoid Tumors} Neuroendocrine tumors that often grow slowly and may release hormones or other active substances. They commonly arise in the digestive tract or lungs and can sometimes invade or spread to distant organs.

\synonyms
Cancer, Carcinoid Tumors

\medterm{Carcinoma} A cancer that begins in epithelial cells, which cover body surfaces and line organs and glands. Common types include adenocarcinoma, squamous-cell carcinoma, basal-cell carcinoma, and urothelial carcinoma; behavior depends on site and cell features.

\medterm{Carcinoma En Cuirasse} An advanced form of breast cancer in which cancer cells spread through lymphatic vessels in the skin, causing widespread hardening and thickening that can make the chest wall resemble armor.

\medterm{Carcinoma Erysipeloides} Spread of an internal cancer through lymph vessels in the skin, producing a red, firm, sharply outlined area that can resemble the infection erysipelas. It most often occurs with cancers of the breast or digestive tract.

\medterm{Carcinoma In Situ} Abnormal cancer-like cells occupying the full thickness of an epithelial lining but not yet invading through its supporting boundary. It is an early, noninvasive stage that may progress if untreated.

\medterm{Carcinoma Of Unknown Primary} Cancer found outside its original site when the starting organ cannot be identified despite evaluation. Treatment uses the tumor’s appearance, genetic features, spread pattern, symptoms, and likely tissue of origin.

\synonyms
Cancer Of Unknown Origin
Occult Primary Cancer

\medterm{Carcinomatosis} Widespread spread of carcinoma to multiple surfaces or organs, often along the lining of the abdomen, chest, or other body cavities. It usually indicates advanced cancer and may cause fluid buildup, pain, bowel problems, or breathing difficulty.

\medterm{Carcinosarcoma} A malignant tumor containing both carcinomatous and sarcomatous elements, most commonly
of the uterus (malignant mixed Müllerian tumor), with an aggressive course.

\medterm{Cardia} The upper part of the stomach surrounding the opening where the esophagus enters. It helps receive swallowed food and can be affected by reflux, inflammation, hiatal hernia, or stomach cancer.

\medterm{Cardiac} Relating to the heart. The word may describe the heart muscle, chambers, valves, electrical system, blood supply, rhythm, or diseases such as coronary disease, cardiomyopathy, and heart failure; in “cardia,” it instead refers to the upper stomach.

\synonyms
Heart-related; cardiac-related.

\medterm{Cardiac Ablation} A procedure that uses heat, cold, or another form of energy to create a small scar in heart tissue causing an abnormal rhythm. The scar blocks faulty electrical signals and may restore a regular rhythm.

\medterm{Cardiac Ablation Procedures} Procedures that use heat, cold, or another energy source to treat heart-rhythm disorders by creating controlled scars that block abnormal electrical signals. Risks include bleeding, blood-vessel injury, recurrence, and damage to nearby heart structures.

\medterm{Cardiac Allograft} A donor heart transplanted into another person, usually for end-stage heart failure. Lifelong immune-suppressing treatment is needed to reduce rejection; complications include infection, rejection, coronary-vessel disease, cancer, and graft failure.

\medterm{Cardiac Allograft Vasculopathy} Progressive narrowing of the heart's blood vessels after transplantation caused by abnormal thickening of their inner walls. It may silently reduce blood flow and cause graft failure, rhythm problems, heart attack, or sudden death.

\medterm{Cardiac Amyloidosis} Buildup of abnormal amyloid protein in the heart, making the heart muscle stiff and impairing filling. It can cause breathlessness, swelling, rhythm problems, or heart failure; treatment depends on whether the protein is light-chain or transthyretin.

\medterm{Cardiac Amyloidosis Diagnosis} Evaluation combines symptoms, examination, blood and urine tests, heart imaging, and sometimes tissue biopsy. Specialized nuclear imaging may support transthyretin amyloidosis only after testing has excluded light-chain amyloidosis.

\medterm{Cardiac Amyloidosis Staging} A system for estimating the severity and outlook of transthyretin cardiac amyloidosis using markers of heart strain and injury, sometimes with other findings. Higher stages indicate more advanced disease; treatment and prognosis must be individualized.

\medterm{Cardiac Arrest} Sudden stopping of effective heart pumping, causing unresponsiveness and absent or abnormal breathing. It is an immediate emergency requiring cardiopulmonary resuscitation, defibrillation when appropriate, and advanced life support.

\medterm{Cardiac Arrest Survival} The chance of surviving cardiac arrest depends on its cause, where it occurs, the initial rhythm, how quickly CPR and defibrillation begin, and care afterward. Early recognition, effective CPR, rapid defibrillation, and post-arrest treatment improve outcomes.

\medterm{Cardiac Arrest, Sudden} Alternate index label; see \textit{Sudden cardiac arrest}.

\medterm{Cardiac Arrhythmias} Abnormal heart rates or rhythms caused by problems forming or conducting electrical impulses. They include rhythms that are too fast, too slow, or irregular and may cause palpitations, dizziness, fainting, or no symptoms.

\medterm{Cardiac Axis} The average direction in which the heart's lower chambers become electrically activated, estimated from an electrocardiogram. A leftward or rightward shift may reflect normal variation or heart, conduction, or lung disease.

\medterm{Cardiac Biomarkers} Substances measured in blood to detect or assess heart-muscle injury. Troponin is the preferred marker; results must be interpreted with symptoms, electrocardiography, imaging, and changes in repeated tests because many conditions can raise levels.

\medterm{Cardiac Cachexia} Severe unintentional loss of weight and muscle associated with advanced heart failure or other serious heart disease. It can cause weakness and reduced function and reflects inflammation, altered metabolism, poor appetite, and increased nutritional needs.

\medterm{Cardiac Care} Prevention, diagnosis, treatment, and follow-up for heart and circulation disorders. It may include risk reduction, medicines, monitoring, procedures, exercise rehabilitation, nutrition support, and care for symptoms and quality of life.

\medterm{Cardiac Catheter} A thin flexible tube inserted through a blood vessel and guided to the heart. It can measure pressures, take blood samples, inject contrast for imaging, or deliver treatment.

\medterm{Cardiac Catheterization} A procedure in which thin tubes are guided through a blood vessel to the heart to measure pressures, assess oxygen levels, inject contrast, image coronary arteries, or perform treatment. Possible risks include bleeding, rhythm problems, vessel injury, and kidney effects.


\medterm{Cardiac Conduction System} Specialized heart tissue that creates and carries electrical signals so the chambers contract in the correct order. It includes the sinoatrial node, atrioventricular node, bundle of His, bundle branches, and Purkinje fibers.

\medterm{Cardiac CT} Computed-tomography imaging of the heart and nearby blood vessels. Depending on the protocol, it can assess coronary calcium, coronary narrowing, heart structure, valves, and selected aspects of heart function.

\medterm{Cardiac CT Angiography} A noninvasive scan using contrast and computed tomography to show the coronary arteries and heart structures. It is useful for assessing selected people with chest pain and can show narrowing, plaque, and coronary calcium.

\medterm{Cardiac Cycle} The events in one complete heartbeat. During diastole the heart chambers relax and fill with blood; during systole the ventricles contract and pump blood into the lungs and the rest of the body.

\medterm{Cardiac Defibrillation} Delivery of an electrical shock to stop ventricular fibrillation or pulseless ventricular tachycardia during cardiac arrest. It must be combined with high-quality CPR and advanced resuscitation according to the rhythm and situation.

\medterm{Cardiac Device Infections} Infection involving a pacemaker, implantable defibrillator, its pocket, or a lead. It may cause redness, pain, drainage, fever, bloodstream infection, or infection of the heart lining and often requires antibiotics and removal of infected hardware.

\medterm{Cardiac Diet} An eating pattern emphasizing vegetables, fruits, whole grains, beans, fish, nuts, and unsaturated fats while limiting excess sodium, added sugar, saturated fat, and trans fat. It should be adapted for heart failure, kidney disease, diabetes, and nutrition needs.

\medterm{Cardiac Enzymes} An older term for blood tests used to detect heart-muscle injury, especially creatine kinase-MB. Troponin is now preferred; any result must be interpreted with symptoms, electrocardiography, imaging, and changes over time.

\medterm{Cardiac Event Monitors} Portable devices that record or send an electrocardiogram when symptoms occur or at planned times. They help detect intermittent abnormal rhythms that may be missed during a short clinic ECG.

\medterm{Cardiac Failure} Alternate index label; see \textit{Heart failure}.
\medterm{Cardiac Fibroma} A rare noncancerous tumor made of fibrous tissue in the heart, usually in a ventricle or septum. It may obstruct blood flow, interfere with electrical signals, or cause rhythm problems and heart failure.

\medterm{Cardiac Gated SPECT} A heart scan using a small radioactive tracer while images are synchronized with the heartbeat. It can assess blood flow to heart muscle, pumping strength, wall movement, and areas that are damaged or still viable.

\medterm{Cardiac Glycoside} A medicine or plant substance that can strengthen heart contraction and slow electrical conduction through the heart. Digoxin is the main clinical example; the safe dose is narrow and toxicity can cause dangerous rhythms.

\medterm{Cardiac Glycoside Overdose} Toxic exposure to a cardiac glycoside such as digoxin. It may cause nausea, vomiting, confusion, visual changes, slow or fast abnormal rhythms, high potassium, seizures, or cardiac arrest and requires urgent assessment.

\medterm{Cardiac Hemangioma} A rare noncancerous tumor made of blood-vessel tissue in or around the heart. Many cause no symptoms, but a large or strategically located tumor may obstruct flow, disturb rhythm, or cause fluid around the heart.

\medterm{Cardiac Imaging Tests} Tests that show the heart's structure, movement, blood flow, or tissue. They include echocardiography, CT, MRI, and nuclear scans; the appropriate test depends on symptoms, the suspected condition, kidney function, and the clinical question.

\medterm{Cardiac Index} The amount of blood the heart pumps each minute adjusted for body size. It helps assess pumping function, especially in seriously ill people, but a single value must be interpreted with blood pressure, symptoms, organ function, and the overall clinical situation.

\medterm{Cardiac Intravascular Ultrasound} An ultrasound probe on a catheter that creates cross-sectional images from inside a coronary artery. It helps measure plaque, vessel size, narrowing, tears, calcium, and stent expansion during selected procedures.

\medterm{Cardiac Ischemia} Alternate index label; see \textit{Myocardial ischemia}.

\medterm{Cardiac Massage} Manual compression of the heart, historically performed through the chest or during open surgery. Modern cardiac-arrest treatment uses external chest compressions and advanced resuscitation rather than routine direct heart massage.
\medterm{Cardiac Metastases} Cancer that has spread to the heart or its surrounding tissues from another organ. Lung and breast cancers, melanoma, lymphoma, and kidney cancer are common sources; spread may cause fluid around the heart, obstruction, rhythm problems, or heart failure.

\medterm{Cardiac MRI} Magnetic-resonance imaging that shows heart structure, pumping, blood flow, scar, swelling, infiltrative disease, and congenital abnormalities without x-rays. Contrast may be avoided or used cautiously in selected people with severe kidney disease.

\medterm{Cardiac Muscle} Involuntary muscle that forms the myocardium, the main muscular layer of the heart. Its cells are electrically connected so the chambers contract in a coordinated way to pump blood.

\medterm{Cardiac Myocyte} A specialized heart-muscle cell that contracts to pump blood. Electrical signals and movement of calcium inside the cell coordinate its contraction and relaxation.

\medterm{Cardiac Myocytes} The muscle cells that make up most of the heart and contract with each heartbeat. Damage to them can reduce pumping function and release proteins such as troponin into the blood.

\medterm{Cardiac Myxoma} The most common primary noncancerous tumor of the heart, usually arising in the left atrium. It can obstruct blood flow, cause breathlessness or constitutional symptoms, or release clots; surgical removal is usually recommended.

\medterm{Cardiac Output} The volume of blood pumped by the heart each minute. It depends mainly on heart rate and stroke volume, the amount ejected with each beat, and helps indicate whether circulation is adequate.

\medterm{Cardiac Output Measurement} Measurement of blood pumped by the heart each minute using methods such as echocardiography, a pulmonary-artery catheter, oxygen-consumption calculations, or noninvasive monitors. The method's accuracy depends on the patient and testing conditions.

\medterm{Cardiac Pacemaker} An implanted or temporary device that sends electrical impulses to maintain an adequate heartbeat when the heart's natural rhythm is too slow or unreliable. It may be placed inside the body or used externally during emergencies.
\medterm{Cardiac Pacemaker Types} Single-chamber pacemakers stimulate one heart chamber; dual-chamber devices coordinate the atrium and ventricle; biventricular or resynchronization devices stimulate both ventricles to improve timing in selected heart-failure patients.

\medterm{Cardiac Perforation} A hole or tear through the heart wall, sometimes caused by injury or a medical procedure. It can cause bleeding around the heart, cardiac tamponade, shock, and death and requires immediate treatment.

\medterm{Cardiac PET} A nuclear-imaging test using a positron-emitting tracer to assess heart-muscle blood flow, metabolism, and living or scarred tissue. Selected tracers can also help evaluate inflammation or infiltrative heart disease.

\medterm{Cardiac PET Imaging} Positron-emission tomography using a radioactive tracer to measure blood flow and metabolism in heart muscle. It can quantify blood flow and help distinguish living from scarred tissue, but availability, radiation exposure, and tracer choice affect its use.

\medterm{Cardiac Power Output} A measure combining blood pressure and cardiac output to estimate how effectively the heart is pumping. It is useful in severe heart failure or cardiogenic shock, but clinicians interpret it with examination findings, organ function, and other circulation measurements.

\synonyms
CPO; Cardiac Power Index (CPI)

\medterm{Cardiac Radiation Injury} Damage to the heart or its blood vessels that may appear months or years after radiation to the chest. It can cause inflammation or scarring around the heart, coronary disease, valve disease, heart-muscle weakness, or rhythm and conduction problems.

\medterm{Cardiac Radionuclide Imaging} Heart imaging performed after a small radioactive tracer is given, usually with SPECT or PET. It can assess blood flow, metabolism, living or scarred heart muscle, inflammation, and, when synchronized with the heartbeat, pumping function.

\medterm{Cardiac Rehab Outcomes} Cardiac rehabilitation can improve exercise ability, blood-pressure and cholesterol control, confidence, quality of life, and psychological well-being after selected heart conditions. Benefits vary with the program, health condition, and participation.

\medterm{Cardiac Rehabilitation} A medically supervised program after conditions such as heart attack, coronary procedures, bypass surgery, or heart failure. It combines gradual exercise, education, risk-factor reduction, medicine support, nutrition, and emotional care to improve recovery and reduce future risk.

\medterm{Cardiac Resynchronization Therapy} A pacing treatment that coordinates contraction of the heart's lower chambers. It is used in selected people with symptomatic heart failure and delayed electrical activation and can improve symptoms, pumping, and survival when appropriate.

\medterm{Cardiac Rhabdomyoma} A usually noncancerous tumor of heart muscle, most often found in children and associated with tuberous sclerosis. It may shrink on its own but can obstruct blood flow or disturb rhythm and pumping if large.

\medterm{Cardiac Sarcoidosis} Sarcoidosis affecting the heart, where clusters of inflammatory cells can form in heart muscle and the electrical system. It may cause heart block, dangerous ventricular rhythms, heart failure, or sudden cardiac arrest.

\medterm{Cardiac Sarcoma} A rare cancer arising from supporting or connective tissue of the heart. It may obstruct blood flow, disturb heart rhythm, impair pumping, or release clots that travel to other parts of the body.

\medterm{Cardiac Skeleton} A strong connective-tissue framework around the heart valves and parts of the wall between the chambers. It supports the valves, provides attachment for heart muscle, and helps electrically separate the upper and lower chambers.

\medterm{Cardiac SPECT} A nuclear scan that shows blood flow to the heart muscle and can help identify reduced flow, prior damage, or living muscle. When synchronized with the heartbeat, it can also assess chamber movement and pumping function.

\medterm{Cardiac Stress Test} A test that observes the heart while it works harder, usually during exercise or after a medicine increases its workload. It may include an ECG or imaging and can help detect reduced blood flow or abnormal rhythms.

\medterm{Cardiac Stress Testing} Assessment of the heart during exercise or medicine-induced stress to look for reduced blood flow, abnormal rhythms, or limited pumping reserve. The method is chosen according to the person's ability to exercise and the clinical question.

\medterm{Cardiac Syncope} Transient loss of consciousness caused by a cardiac disorder that reduces cerebral perfusion, usually through bradyarrhythmia, tachyarrhythmia, structural obstruction, or severe cardiac dysfunction. Syncope during exertion, while supine, or with palpitations or known heart disease requires urgent cardiovascular evaluation because it may indicate a risk of sudden cardiac death.

\medterm{Cardiac Tamponade} A medical emergency in which blood or fluid builds up under pressure around the heart and prevents the chambers from filling normally. It reduces blood flow to the body and can cause shock or cardiac arrest.

\medterm{Cardiac Tamponade Physiology} Pressure from blood or fluid around the heart prevents the chambers from relaxing and filling. This lowers cardiac output and may cause low blood pressure, a marked breathing-related change in pulse pressure, and shock.

\medterm{Cardiac Transplantation} Replacement of a failing heart with a donor heart in carefully selected people with end-stage heart failure. Lifelong immune-suppressing treatment and regular monitoring are needed for rejection, infection, coronary-vessel disease, and other complications.

\medterm{Cardiac Tumors} Tumors involving the heart may begin there or spread from another organ. They are rare and may obstruct blood flow, disturb rhythm, cause fluid around the heart, or release clots; treatment depends on type, location, size, and spread.

\medterm{Cardiac Tumors Management} Management depends on whether a tumor is noncancerous or cancerous, its size, location, symptoms, and spread. Observation may suit some harmless tumors; surgery, medicines, radiation, or chemotherapy may be needed for others.

\medterm{Cardiac Ventriculography} Imaging of the heart's lower chambers to assess their size, wall movement, and ability to pump blood. It may use x-rays with contrast or another imaging method and is now often replaced by echocardiography or MRI.

\medterm{Cardiac Volume} The amount of blood in a heart chamber at a particular time. The amount before contraction and after contraction helps assess chamber size, filling, and pumping function.

\medterm{Cardiac Wall} The three-layered wall of the heart: the inner endocardium, muscular myocardium, and outer epicardium. It supports the heart's structure, electrical activity, contraction, and interaction with surrounding tissues.
\medterm{Cardinal Body Part} A major body region or structure used in anatomical description or clinical examination, such as the head, neck, chest, abdomen, pelvis, limbs, or an important internal organ.

\medterm{Cardinal Cell Part} A main component of a cell, such as its membrane, cytoplasm, nucleus, or an organelle. Each part has a role in maintaining the cell, processing information, or performing specialized functions.

\medterm{Cardinal Ligament} A band of strong connective tissue extending from the cervix and upper vagina to the sidewall of the pelvis. It helps support the uterus and vagina and contains the uterine blood vessels, which cross over the ureter.

\synonyms
Mackenrodt's Ligament; Lateral Cervical Ligament; Transverse Cervical Ligament

\medterm{Cardio-Oncology} The medical field focused on preventing and treating heart and blood-vessel problems in people with cancer. It monitors effects of chemotherapy, targeted or immune treatment, and chest radiation while helping cancer therapy continue safely when possible.

\medterm{Cardiobacteriaceae} A family of bacteria that includes Cardiobacterium. Some members can enter the bloodstream and infect heart valves, causing infective endocarditis.

\medterm{Cardiobacterium} A slow-growing bacterium in the HACEK group that can enter the bloodstream and infect heart valves. It may cause prolonged fever, fatigue, weight loss, a new murmur, or embolic complications.

\medterm{Cardiobacterium Hominis} The Cardiobacterium species most often associated with human infection. It can cause slowly developing infective endocarditis with fever, fatigue, weight loss, a new heart murmur, valve damage, or clots traveling to other organs.

\medterm{Cardiofaciocutaneous Syndrome} A rare genetic disorder affecting the heart, facial appearance, skin, hair, growth, and development. It is usually caused by a change in a gene involved in cell-growth signals.
\synonyms
CFC Syndrome

\medterm{Cardiogenic Shock} A state of inadequate tissue perfusion caused by severe cardiac dysfunction. It may produce
low blood pressure, altered mental status, cool extremities, reduced urine output, metabolic acidosis, or other signs of
organ hypoperfusion.

\medterm{Cardiogenic Shock Management} Emergency management of shock caused by severe cardiac dysfunction. It includes
support of oxygenation and circulation, treatment of the underlying cause, and selected use of vasopressors, inotropes,
mechanical support, or revascularization.

\medterm{Cardiolipin} A fat-like molecule found mainly in the inner membrane of mitochondria, the cell's energy-producing structures. Antibodies against cardiolipin can be associated with antiphospholipid syndrome and abnormal blood clotting.

\medterm{Cardiologist} A doctor trained to diagnose and treat diseases of the heart and blood vessels, including coronary disease, rhythm disorders, valve disease, heart failure, and high blood pressure.

\medterm{Cardiology} The branch of medicine concerned with the heart and blood vessels. It covers prevention, diagnosis, and treatment of coronary disease, heart failure, rhythm and valve disorders, cardiomyopathy, congenital heart disease, and high blood pressure.

\medterm{Cardiomegaly} Enlargement of the heart seen on examination or imaging. It may result from enlarged chambers, thickened heart muscle, fluid around the heart, or a combination; causes include high blood pressure, valve disease, cardiomyopathy, and congenital disease.

\medterm{Cardiomyopathy} Disease of the heart muscle that can impair contraction, relaxation, electrical rhythm, or valve function and may lead to heart failure or sudden death. Major patterns include dilated, hypertrophic, and restrictive cardiomyopathy.

\medterm{Cardiomyopathy (Restrictive)} Alternate index label for Restrictive.

\medterm{Cardiomyopathy Classification} A system for describing cardiomyopathy by its heart structure and function, other organs involved, genetic pattern, cause, and functional severity. This information helps communicate diagnosis, guide testing, and estimate prognosis.

\medterm{Cardiomyopathy, Dilated} Alternate index label; see \textit{Dilated cardiomyopathy}.

\medterm{Cardiomyopathy, Hypertrophic} Alternate index label; see \textit{Hypertrophic cardiomyopathy}.

\medterm{Cardiopathy} Any disease or disorder of the heart, including problems of the heart muscle, valves, coronary arteries, electrical rhythm, or structure present from birth.

\medterm{Cardioplegia} Temporary stopping and protection of the heart during open-heart surgery, usually by delivering a cold solution that changes electrical and metabolic activity. It allows the surgeon to operate while limiting heart-muscle injury.

\medterm{Cardiopulmonary} Relating to the heart and lungs or to their combined function. For example, cardiopulmonary exercise testing assesses how the heart, lungs, blood, and muscles respond during physical activity.

\synonyms
Cardiorespiratory.

\medterm{Cardiopulmonary Bypass} A heart-lung machine temporarily takes over circulation and oxygen exchange during some operations. It pumps and oxygenates blood while the heart is stopped or opened; possible complications include bleeding, inflammation, clots, stroke, and organ injury.

\medterm{Cardiopulmonary Bypass Machine} A machine that temporarily pumps and oxygenates blood during selected heart or major-vessel operations. It allows surgeons to work while the heart is stopped or the lungs are not carrying out their usual gas exchange.

\synonyms
Heart-lung machine; cardiopulmonary bypass pump.

\medterm{Cardiopulmonary Exercise Testing} A monitored treadmill or bicycle test that measures how the heart, lungs, blood, and muscles respond as exercise becomes harder. It helps assess unexplained breathlessness, exercise limitation, and severity of heart or lung disease.

\medterm{Cardiopulmonary Resuscitation} Emergency chest compressions and, when appropriate, rescue breaths, defibrillation, and advanced life support for cardiac arrest. Its purpose is to maintain blood flow and oxygen delivery until the heartbeat returns or definitive treatment is provided.

\medterm{Cardiorenal Syndrome} A condition in which heart and kidney dysfunction worsen one another. Sudden or long-term heart disease can injure the kidneys, and sudden or long-term kidney disease can strain the heart; treatment must address both organs.

\medterm{Cardiorespiratory} Relating to the heart, blood vessels, and respiratory system considered together, as in cardiorespiratory fitness, monitoring, or failure.
\medterm{Cardiorespiratory Endurance} The ability of the heart, lungs, blood, and muscles to deliver and use oxygen during sustained physical activity. It is influenced by fitness, heart and lung disease, anemia, muscle condition, and other factors.
\synonyms
Cardiorespiratory fitness; aerobic endurance.

\medterm{Cardioselective Beta-Blockers} Beta-blocker medicines that mainly block beta-1 receptors in the heart, slowing the heart rate and reducing its workload. They may still narrow airways at higher doses and are used for selected heart, blood-pressure, and rhythm conditions.

\medterm{Cardiospasm} An older term for failure of the lower esophageal muscle to relax normally, as occurs in achalasia. It can cause difficulty swallowing, regurgitation of food or liquid, chest discomfort, and weight loss.

\medterm{Cardiospasm (Achalasia)} Alternate index label for Achalasia.

\medterm{Cardiotocography} Continuous recording of a fetus's heart rate and the mother's uterine contractions during pregnancy or labor. The pattern helps assess fetal well-being, but results must be interpreted with the stage of labor and the overall clinical situation.

\medterm{Cardiotoxicity} Injury or impaired function of the heart caused by a medicine, toxin, radiation, or disease. It may produce rhythm problems, reduced heart-muscle strength, chest pain from poor blood flow, or heart failure.

\medterm{Cardiotoxicity From Chemotherapy} Heart injury caused by some chemotherapy medicines. It may weaken heart muscle, reduce pumping, or cause rhythm problems during or after treatment; risk depends on the medicine, dose, other treatments, and existing heart disease.

\medterm{Cardiotoxin} A substance that directly injures heart muscle or disrupts the heart's electrical activity or blood supply. It may cause abnormal rhythms, reduced pumping, heart failure, or heart-muscle injury.

\medterm{Cardiovascular} Relating to the heart and blood vessels, including circulation, blood pressure, and diseases affecting the heart or vascular system.

\medterm{Cardiovascular Agents} Medicines that act on the heart or blood vessels. They include treatments for high blood pressure, abnormal rhythms, angina, heart failure, blood clots, and abnormal cholesterol levels.

\medterm{Cardiovascular Changes In Pregnancy} During pregnancy, blood volume and cardiac output increase to support the uterus and developing baby. The rise in plasma is greater than the rise in red cells, causing normal dilutional anemia; heart rate and blood-vessel resistance also change.

\medterm{Cardiovascular Disease} Disease of the heart or blood vessels, including coronary artery disease, heart failure, abnormal rhythms, valve disease, stroke-related vascular disease, and disorders of arteries or veins.

\medterm{Cardiovascular Health} The condition of the heart and blood vessels and the habits that protect them. Regular activity, not using tobacco, healthy eating, adequate sleep, and control of blood pressure, cholesterol, diabetes, and weight can lower cardiovascular risk.

\medterm{Cardiovascular Infection} An infection involving the bloodstream, heart, heart valves, or blood vessels. Examples include infective endocarditis, myocarditis, infected vascular grafts, and infected pacemakers or other devices.

\medterm{Cardiovascular Injury} Damage to the heart or blood vessels caused by trauma, reduced blood flow, inflammation, toxins, surgery, or another disease. Effects may include bleeding, impaired circulation, heart-muscle injury, or abnormal rhythms.

\medterm{Cardiovascular Neoplasm} A noncancerous or cancerous tumor arising in the heart, blood vessels, or related tissues. Its symptoms and treatment depend on the tissue type, location, size, effect on circulation, and whether it has spread.

\medterm{Cardiovascular Pharmacology} The study of how medicines affect the heart, blood vessels, blood pressure, rhythm, blood clotting, and cholesterol or lipid metabolism, including their uses, interactions, benefits, and risks.

\medterm{Cardiovascular Risk Reduction} Actions that lower the chance of heart attack, stroke, heart failure, and vascular death. They include avoiding tobacco, controlling blood pressure, cholesterol, and diabetes, staying active, eating well, maintaining a healthy weight, and taking indicated medicines.

\medterm{Cardiovascular Screening} Assessment for cardiovascular risk or disease, often before symptoms develop. It may include blood-pressure measurement, cholesterol and diabetes testing, review of family history and tobacco use, and targeted tests based on age and risk.

\medterm{Cardiovascular Surgery} Operations on the heart or major blood vessels, including bypass surgery, valve repair or replacement, aortic surgery, and heart transplantation. Many open procedures require a heart-lung machine and specialized care before and after surgery.

\medterm{Cardiovascular System} The heart and network of arteries, arterioles, capillaries, venules, and veins that circulate blood. It delivers oxygen and nutrients, removes waste, regulates temperature, and helps maintain blood pressure.

\medterm{Cardioversion} Restoration of a more regular heart rhythm using a synchronized electrical shock or medicine. Before cardioversion for some atrial rhythms, clinicians assess clot risk and may give anticoagulant treatment to reduce the risk of stroke.

\medterm{Cardioversion (Synchronized)} Delivery of an electrical shock timed with the heart's electrical cycle to restore a more organized rhythm. It treats selected fast rhythms; sedation, clot prevention, monitoring, and energy settings depend on the rhythm and patient.

\synonyms
Synchronized

\medterm{Cardiovirus} A group of viruses in the picornavirus family. Some infect people or animals and may cause fever, muscle inflammation, heart inflammation, or other illness; effects depend on the species and the person's defenses.

\medterm{Carditis} Inflammation of the heart muscle, inner lining, valves, or surrounding sac. It may result from infection, an immune reaction, medicines, or another inflammatory disease and can cause chest pain, breathlessness, rhythm problems, or heart failure.

\medterm{Care Coordination} Organizing communication and health-care activities among clinicians, services, patients, and caregivers so care remains safe, continuous, and appropriate when people move between settings or receive several types of treatment.

\medterm{Caregiver} A family member, friend, or paid worker who helps another person with health care, personal activities, daily tasks, transportation, medicines, or emotional support.

\medterm{Caregiver Burden} The emotional, physical, social, or financial strain of providing ongoing care for another person. It may lead to exhaustion, low mood, sleep problems, isolation, or difficulty meeting one's own health needs.
\medterm{Caregiver Stress} Strain caused by caring for another person, often leading to tiredness, worry, irritability, sleep problems, or reduced well-being. Support, respite, counseling, and practical help can reduce its effects.




\medterm{Caries} A biofilm-mediated, multifactorial disease in which acids from bacterial metabolism of fermentable carbohydrates cause net demineralization of enamel, dentin, or cementum. Early lesions may be arrested or remineralized, but progressive disease forms cavities and can reach the pulp, causing pain or infection.

\medterm{Caries (Dental)} Tooth decay caused when bacteria in plaque use sugars to produce acids that remove minerals from enamel, dentin, or root surface. Repeated attacks can form cavities, pain, infection, and tooth loss; fluoride and good cleaning lower risk.

\synonyms
Dental

\medterm{Caries Risk Assessment} Evaluation of a person's chance of developing tooth decay by reviewing past and current cavities, plaque, diet, fluoride exposure, saliva, oral hygiene, tooth structure, medicines, and medical conditions. Results guide prevention and follow-up.

\medterm{Carina} The ridge at the bottom of the windpipe where it divides into the right and left main bronchi. It is a landmark during bronchoscopy and breathing-tube placement; inhaled objects more often enter the wider, more vertical right bronchus.


\medterm{Cariogenic} Likely to promote tooth decay, especially by supplying sugars that plaque bacteria convert into acids or by supporting conditions that remove minerals from teeth.

\medterm{Cariogram} A computer-based tool that combines factors such as diet, plaque bacteria, fluoride exposure, saliva, previous cavities, and social or clinical history to estimate tooth-decay risk and show where prevention may help.

\medterm{Carisoprodol} A prescription muscle relaxant used short term for acute painful muscle conditions. It can cause drowsiness, dizziness, impaired coordination, misuse, dependence, and dangerous breathing suppression, especially with alcohol or sedatives.

\medterm{Carminative} A substance traditionally used to ease intestinal gas or bloating. It may relax intestinal muscle or help gas pass, but persistent pain, vomiting, swelling, or altered bowel habits need medical assessment.

\synonyms
Antiflatulent; gas-relieving agent.

\medterm{Carmine} A red coloring made from cochineal insects and used in some foods, cosmetics, and medicines. Rarely, it causes allergy, wheezing, hives, or anaphylaxis in sensitive people.

\medterm{Carmustine} A chemotherapy medicine that damages DNA and is used for selected brain tumors, blood cancers, and other cancers. It can lower blood counts and cause nausea, infertility, liver problems, or delayed lung injury.

\medterm{Carney Complex} A rare inherited disorder, usually caused by a PRKAR1A gene change, that can produce multiple noncancerous skin, heart, and endocrine tumors. It may also cause excess hormones and requires lifelong, organ-specific monitoring.

\medterm{Carney Triad} A rare condition involving combinations of a stomach gastrointestinal stromal tumor, a benign cartilage tumor in the lung, and a paraganglioma outside the adrenal gland. It usually occurs sporadically and more often affects young women.

\medterm{Carnitine} A substance that transports long-chain fats into mitochondria, the cell structures that make energy, so the fats can be used as fuel. The body usually makes enough; severe deficiency can cause weakness, low blood sugar, or heart and muscle disease.

\medterm{Carnitine Palmitoyltransferase II Deficiency} A rare inherited disorder in which muscle cells cannot use long-chain fats normally for energy. Exercise, fasting, cold, or infection can trigger muscle pain, stiffness, muscle breakdown, dark urine, and kidney injury; symptoms may be absent between attacks.

\synonyms
CPT II Deficiency; CPT2 Deficiency

\medterm{Carnobacteriaceae} A family of bacteria that produce lactic acid and live mainly in food, animals, or the environment. Most do not cause human disease, although rare opportunistic infections have been reported.

\medterm{Carnosine} A small molecule made from beta-alanine and histidine, found mainly in muscle and nervous tissue. It helps buffer acidity; claims that supplements prevent disease or improve performance remain uncertain.

\medterm{Carnosinemia} A rare inherited disorder in which deficient carnosinase causes excess carnosine in blood and urine. Children may have developmental delay, low muscle tone, seizures, or other neurologic problems of varying severity.

\medterm{Caroli Disease} A rare inherited disorder in which the larger bile ducts inside the liver are widened in segments. It can cause repeated bile-duct infections, stones, abdominal pain, jaundice, and, in some people, liver scarring or kidney disease.

\medterm{Carotene} An orange-yellow plant pigment found in foods such as carrots and leafy vegetables. Some forms, especially beta-carotene, can be converted to vitamin A; excess intake may harmlessly turn the skin yellow-orange.

\medterm{Carotenoid} A family of fat-soluble yellow, orange, or red pigments found mainly in plants and some microorganisms. Some, such as beta-carotene, can form vitamin A; others, such as lutein, support eye tissues.

\medterm{Carotenoids} Plant pigments that give many fruits and vegetables their yellow, orange, or red color. Some can be converted to vitamin A, while others have roles in eye and cellular function.

\synonyms
Carotenoid pigments; provitamin-A pigments.

\medterm{Carotid} Relating to either of the two major arteries in the neck that carry blood from the heart toward the brain, face, and other head and neck structures.

\medterm{Carotid Artery} One of the major arteries supplying the head and neck. Each common carotid divides into an internal branch that supplies the brain and an external branch that supplies much of the face, scalp, and neck.

\medterm{Carotid Artery Disease} Narrowing or blockage of a carotid artery, usually from fatty plaque. It can reduce blood flow to the brain or release a clot and cause a transient ischemic attack or ischemic stroke.

\medterm{Carotid Artery Stenting} A procedure that places a small mesh tube inside a narrowed carotid artery to keep it open. It is used selectively to reduce stroke risk; clot release, stroke, bleeding, and re-narrowing are possible complications.

\medterm{Carotid Body} A small group of sensing cells near the division of the common carotid artery. It detects changes in blood oxygen, carbon dioxide, and acidity and sends signals that help adjust breathing.

\medterm{Carotid Bruit} A blowing sound heard over a carotid artery when blood flow is turbulent. It can occur with plaque but does not reliably show how narrow the artery is; ultrasound or other vascular imaging is needed when disease is suspected.

\medterm{Carotid Cavernous Fistula} An abnormal connection between the carotid artery and the cavernous sinus, a large venous space behind the eye. It may cause a bulging or red eye, pulsation, a whooshing sound, double vision, or vision loss and needs specialist treatment.

\medterm{Carotid Duplex Ultrasound} An ultrasound test that shows plaque and measures blood flow through the carotid arteries. It helps estimate narrowing without radiation and may guide decisions about medicines, surgery, or stenting.

\medterm{Carotid Endarterectomy} Surgery to open a carotid artery and remove fatty plaque from its inner wall. It can reduce stroke risk in selected people with significant narrowing, especially after symptoms, but carries risks such as stroke, heart attack, and nerve injury.

\medterm{Carotid Stenosis} Narrowing of a carotid artery, usually from fatty plaque. It may reduce blood flow to the brain or release a clot, causing a transient ischemic attack or stroke; treatment depends on symptoms, severity, and overall risk.

\medterm{Carpal Bones} The eight small bones forming the wrist: scaphoid, lunate, triquetrum, pisiform, trapezium, trapezoid, capitate, and hamate. They connect the forearm to the hand and work together to allow wrist movement and transmit force.

\synonyms
Carpals

\medterm{Carpal Joints} The joints between the wrist bones and between the wrist, forearm, and hand bones. Together they allow the wrist to bend, straighten, move side to side, and rotate with the forearm.

\medterm{Carpal Tunnel} A narrow passage on the palm side of the wrist formed by the wrist bones and a strong ligament. It contains the median nerve and finger-flexing tendons; pressure on the nerve causes carpal tunnel syndrome.

\medterm{Carpal Tunnel Biopsy} Removal of a small tissue sample from the carpal tunnel or its contents when an unusual infiltrative, inflammatory, infectious, or tumor-related process is suspected. It is not a routine test for ordinary carpal tunnel syndrome.

\medterm{Carpal Tunnel Release} Surgery that cuts the ligament forming the roof of the wrist tunnel to relieve pressure on the median nerve. It may improve numbness, tingling, pain, and weakness when symptoms do not respond to suitable nonsurgical care.

\medterm{Carpal Tunnel Syndrome} Pressure on the median nerve as it passes through the narrow wrist tunnel. It commonly causes numbness, tingling, burning, or pain in the thumb, index, middle, and part of the ring finger, sometimes with thumb weakness.

\synonyms
Median Nerve Compression (Carpal Tunnel Syndrome)

\medterm{Carpal Tunnel Syndrome Tests} Examination and nerve tests used to assess suspected median-nerve compression. Wrist-position or tapping tests may reproduce symptoms, while nerve-conduction studies and ultrasound can support diagnosis and assess severity when needed.

\medterm{Carpometacarpal Joint} A joint between a wrist bone and the base of a hand bone. The thumb joint is a mobile saddle-shaped joint that allows opposition and is commonly affected by osteoarthritis.

\medterm{Carpus} The group of eight small bones forming the wrist between the forearm and the metacarpals of the hand. The carpus provides stability while allowing controlled wrist movement.
\medterm{Carrageenan} A substance extracted from red seaweed and used to thicken or stabilize some foods and medicines. Different forms have different properties; evidence about effects on human inflammation depends on the type and amount consumed.

\medterm{Carrier} A person who has one altered copy of a gene for a recessive condition and usually has no symptoms. A carrier can pass the altered gene to a child; risk depends on the other parent's genes and the inheritance pattern.

\medterm{Carrier-Mediated Transport} Movement of a substance across a cell membrane with help from a specific carrier protein. It may use no direct energy when moving down a concentration gradient or use energy to move against that gradient.

\medterm{Carrier Excipient} An inactive ingredient that supports an active medicine in a formulation. It may help control the medicine's dose, stability, absorption, appearance, or delivery and can occasionally cause allergy or intolerance.

\medterm{Carrier Screening} Genetic testing for people without symptoms to find whether they carry a harmful gene variant that could be passed to a child. Results can guide reproductive choices and may identify the need to test a partner or relatives.

\medterm{Carrier State} A condition in which a person carries and may spread a microorganism without symptoms. Some carriers need testing, treatment, vaccination of contacts, or other measures to prevent transmission, depending on the infection.

\medterm{Carry-Over Effect} Persistence of an earlier treatment's effect into a later measurement or treatment period. It can distort results in crossover studies or make a person's current clinical state difficult to interpret.

\medterm{Carteolol} A nonselective beta-blocker used mainly as eye drops to lower pressure inside the eye in open-angle glaucoma or ocular hypertension. It can sometimes slow the heartbeat, lower blood pressure, or worsen breathing problems.

\medterm{Cartilage} A tough, flexible tissue found in joints, the ear, nose, rib cage, and spinal discs. It cushions joints, supports body structures, and resists pressure, but it has limited blood supply and heals slowly after injury.

\medterm{Cartilage Disease} A disorder that damages cartilage in joints or supporting structures. It can cause pain, stiffness, swelling, deformity, impaired growth, or loss of smooth movement and may result from injury, aging, inflammation, or inherited disease.

\medterm{Cartilaginous Joint} A joint in which bones are connected by cartilage and allow little or limited movement. Examples include joints between spinal bones and the joint joining the two front pelvic bones.

\medterm{Caruncles} Small fleshy anatomical projections. Examples include the lacrimal caruncle at the inner corner of the eye and the urethral caruncle near the opening of the urethra.
\medterm{Carvedilol} A beta-blocker that also relaxes blood vessels. It slows the heart and reduces its workload and is used for selected cases of high blood pressure, heart failure, and after a heart attack; it can lower blood pressure and heart rate.

\medterm{CAS} Alternate index label; see \textit{Childhood apraxia of speech}.

\medterm{Casanthranol} A stimulant laxative that increases bowel movement and secretion of water into the intestine. Prolonged or excessive use can cause diarrhea, dehydration, electrolyte abnormalities, and reliance on laxatives.

\medterm{Cascade} A series of linked reactions in which one activated step triggers the next. Examples include the blood-clotting, complement, and inflammatory cascades.

\medterm{Cascara} A plant-derived stimulant laxative that increases bowel movement. Prolonged or excessive use can cause diarrhea, dehydration, electrolyte imbalance, and reliance on laxatives; its use is restricted in some places.
\medterm{Case Definition} Standardized clinical, laboratory, and exposure criteria used to classify a person as a suspected, probable, or confirmed case during surveillance or an outbreak investigation. It may differ from the criteria used to diagnose one patient.

\medterm{Case Fatality Rate} The proportion of people diagnosed with a disease who die from it during a stated period. It describes severity in a particular group, but depends on accurate diagnosis, follow-up, treatment, and the population studied.

\medterm{Case Management} Coordinated planning, referral, and follow-up of health and social services to help a person receive appropriate care, avoid gaps or duplication, and meet clinical, functional, and personal goals.

\medterm{Case Study} A detailed examination of one patient, disease episode, or clinical situation. Case studies support education, research, and quality improvement by illustrating diagnostic reasoning, treatment decisions, outcomes, or unusual presentations.

\medterm{Case-Control Studies} Observational studies that compare people with a disease or outcome with similar people without it and look backward for differences in earlier exposures, treatments, behaviors, or risk factors.

\medterm{Case-Control Study} An observational research study comparing people who have a disease or outcome with similar people who do not. Researchers examine earlier exposures or characteristics to identify possible causes or risk factors.

\medterm{Case-Fatality Rate} The proportion of diagnosed people who die from a disease during a stated period, calculated as deaths from that disease divided by diagnosed cases. It estimates severity in the group studied and is not the same as population mortality.

\medterm{Casein} The main family of proteins in cow's milk and other mammalian milk. It supplies amino acids and minerals and forms curds during digestion and cheese-making; it can trigger reactions in people with cow's-milk protein allergy.

\medterm{Casein Kinase} A family of enzymes that attach phosphate groups to proteins and thereby change their activity. Casein kinases help regulate cell signaling, growth, DNA repair, development, and the body's internal clock.
\medterm{Caseous} Having a soft, crumbly, cheese-like appearance. The term often describes dead tissue in granulomas caused by tuberculosis or certain other long-lasting infections.
\medterm{Caseous Necrosis} A form of tissue death producing soft, crumbly, white or yellow material that resembles cheese. It is classically seen in granulomas caused by tuberculosis, though similar changes can occur in some other chronic infections.

\medterm{Cashew Allergy} An immune reaction to cashew proteins that may cause itching, hives, swelling, vomiting, wheezing, or life-threatening anaphylaxis. People with this allergy should avoid cashews and follow an emergency action plan.

\medterm{Caspase} A group of enzymes that cut specific proteins inside cells. Some initiate orderly programmed cell death, while others activate inflammation and immune defense against infection.

\medterm{Caspase Activation} Conversion of inactive caspase enzymes into their active form inside a cell. Depending on the caspase and triggering signal, this can start programmed cell death or activate inflammation and immune defense.

\medterm{Caspofungin} An intravenous antifungal medicine that blocks construction of an important fungal cell-wall component. It is used for serious Candida infections and selected Aspergillus infections and may cause liver-test abnormalities, infusion reactions, or low potassium.
\medterm{Caspofungin Acetate} The acetate salt form of caspofungin, an intravenous antifungal medicine that blocks fungal cell-wall construction. It is used for selected serious Candida and Aspergillus infections under medical supervision.

\medterm{Cassette} A holder that contains recording or imaging material. In traditional radiography, it held the film or image receptor during exposure; the word may also describe holders used in laboratory or monitoring equipment.

\medterm{Cast} A rigid or semi-rigid support placed around an injured limb or joint to limit movement while a bone, ligament, tendon, or other tissue heals. It must be monitored for swelling, pressure injury, numbness, or poor circulation.

\medterm{Castleman Disease} A group of uncommon disorders causing abnormal enlargement and activity of lymph nodes. Disease may affect one lymph-node area or the whole body and can cause fever, fatigue, anemia, weight loss, and organ problems.

\medterm{Castor Oil} A vegetable oil from the seeds of Ricinus communis. Taken by mouth, it acts as a stimulant laxative; excessive use can cause cramping, diarrhea, dehydration, and electrolyte disturbances.

\medterm{Castor Oil Overdose} Excessive ingestion of castor oil, causing intestinal irritation, cramping, nausea, vomiting, and diarrhea. Severe exposure can cause dehydration, low blood pressure, or fluid and electrolyte imbalance.

\medterm{Castration} Removal or functional suppression of the testes or ovaries, reducing production of sex hormones. It may be surgical or achieved with medicines and can affect fertility, sexual function, mood, and bone health.

\medterm{CAT Scan} An older name for computed tomography, or CT, an imaging test that uses many x-ray measurements and computer processing to create cross-sectional pictures of the body.

\medterm{Cat'S Eye Reflex} A bright white reflection from the pupil, called leukocoria, rather than the usual red reflection in photographs. It can signal cataract, retinal disease, or retinoblastoma and requires urgent eye assessment, especially in a child.
\medterm{Cat-Nap} A short period of sleep, usually during the day. Frequent unintended sleep may indicate inadequate sleep, a medicine effect, or a sleep disorder and should be assessed if it affects safety or daily life.
\medterm{Cat-Scratch Disease} An infection caused by Bartonella henselae, usually after a scratch or bite from an infected cat. It commonly causes a small skin bump, nearby swollen lymph nodes, fever, and fatigue; rarely it affects the eye, liver, spleen, brain, or heart.

\synonyms
Cat-Scratch Fever

\medterm{Catabolism} The set of chemical processes that break down complex molecules such as carbohydrates, fats, and proteins into simpler substances, releasing energy that cells can use. Examples include digestion-related breakdown, glycolysis, and fat oxidation.

\medterm{Catagen} The short transition phase of the hair cycle between active growth and rest. The hair follicle shrinks, hair growth stops, and the hair remains attached until the later shedding phase.

\synonyms
Catagen phase; regression phase of hair growth.

\medterm{Catalase} An enzyme that breaks down hydrogen peroxide into water and oxygen. It protects cells from oxidative damage and helps some bacteria survive inside immune cells.

\medterm{Catalepsy} A state of greatly reduced movement and responsiveness in which a person may remain rigid or hold an imposed posture. It can occur with catatonia, narcolepsy, neurologic disease, or some psychiatric conditions.

\medterm{Catalyst} A substance that speeds a chemical reaction without being used up. It works by providing an easier reaction pathway; enzymes are biological catalysts, and catalysts change reaction speed without changing the final balance of products and reactants.

\medterm{Catamenial Pneumothorax} Repeated collapse of a lung occurring around the start of menstruation, usually because of thoracic endometriosis. It often affects the right side and causes sudden chest pain or breathlessness; recurrent cases need specialist evaluation and may require surgery or hormone treatment.

\synonyms
Endometriosis-Related Pneumothorax; Catamenial Thorax

\medterm{Cataplasm} A poultice made from a soft, moist substance spread on cloth and applied to the body. It was historically used to warm or soothe a local problem and is not a standard treatment for serious disease.

\medterm{Cataplexy} A sudden, brief loss of muscle strength triggered by strong emotion such as laughter, surprise, or anger. Awareness usually remains clear; it is strongly associated with narcolepsy and may cause the head, jaw, knees, or whole body to buckle.

\medterm{Cataract} Clouding of the eye's natural lens, which reduces the passage of light to the retina. It can cause blurred or faded vision, glare, poor contrast, and difficulty seeing at night; lens-replacement surgery can restore sight when daily activities are affected.

\medterm{Cataract - Adult} Clouding of the eye lens that develops during adulthood and gradually reduces sharpness, contrast, color perception, or tolerance of glare. Updating glasses may help temporarily, but lens-replacement surgery is the definitive treatment when vision limits daily activities.

\medterm{Cataract Removal} Surgery that removes a cloudy natural lens from the eye and usually replaces it with an artificial intraocular lens. It is performed when the cataract interferes with vision, work, driving, reading, or other daily activities.

\medterm{Cataract Surgery} Removal of a clouded eye lens and replacement with an artificial lens to improve vision. Possible complications include infection, bleeding, retinal problems, increased eye pressure, or a need for later treatment, though serious complications are uncommon.

\medterm{Cataracts} Clouding of one or both natural eye lenses, causing blurred or faded vision, glare, reduced contrast, or difficulty seeing in dim light. They usually develop with aging but can also follow injury, disease, medicines, or be present at birth.


\medterm{Catastrophic Aging} A rapid loss of physical or mental function in an older person after an illness, operation, fall, or other stress. Reduced physiologic reserve makes recovery harder; early assessment, treatment, rehabilitation, and support may limit decline.

\medterm{Catastrophic Reaction} An unusually intense emotional or behavioral response to a situation that may seem minor to others. It can occur with dementia, brain injury, or other neurologic disease and may involve shouting, fear, crying, or agitation.

\synonyms
Catastrophic response; catastrophic emotional reaction.

\medterm{Catatonia} A serious syndrome involving marked changes in movement, speech, and response to the environment. Features may include immobility, mutism, staring, unusual postures, rigidity, repetitive movements, or copying speech or actions; causes include psychiatric and medical disorders.

\medterm{Catatonia, Fatal} A severe, life-threatening form of catatonia with extreme motor disturbance, fever, unstable blood pressure or heart rate, dehydration, and risk of organ failure. It requires emergency treatment.
\medterm{Catbite} A puncture wound caused by a cat's teeth. Cat bites can push bacteria deep into tissue and may infect tendons, joints, or bone, so prompt cleaning and medical assessment are often needed.
\medterm{Catch-Up Immunization} Vaccination of a person who is behind the recommended schedule using age- and risk-appropriate intervals. A delayed series usually does not need to be restarted, but minimum ages, intervals, and medical precautions must be followed.

\medterm{Catechin} A group of plant compounds found in tea, cocoa, some fruits, and other foods. They show antioxidant effects in laboratory studies, but catechin supplements are not proven treatments and may cause side effects or interact with medicines.

\medterm{Catechol} A chemical compound with two neighboring hydroxyl groups, used in industry and involved in making some biologically active substances. Contact or inhalation can irritate tissues and substantial exposure may be toxic.

\medterm{Catecholamine} A group of hormones and nerve-signaling chemicals that includes epinephrine, norepinephrine, and dopamine. They help control the body's stress response, heart rate, blood pressure, alertness, and metabolism.

\medterm{Catecholamine Blood Test} A blood test measuring epinephrine, norepinephrine, and dopamine. It may help evaluate unusual hormone release from a pheochromocytoma or related tumor, although metanephrine testing is often preferred for initial screening.

\medterm{Catecholamine Crisis} A sudden, life-threatening surge of stress hormones causing very high blood pressure, rapid or irregular heartbeat, sweating, headache, and possible injury to the heart, brain, kidneys, or other organs. It can occur with certain adrenal or nerve tumors.

\medterm{Catecholaminergic Polymorphic Ventricular Tachycardia} A rare inherited rhythm disorder in which exercise or strong emotion triggers dangerous fast rhythms from the heart's lower chambers. The heart may look normal and the resting ECG may be normal; fainting or cardiac arrest can occur, especially in children and young adults.

\synonyms
CPVT; Bidirectional Ventricular Tachycardia; Stress-Induced Polymorphic VT

\medterm{Catecholamines} Adrenaline, noradrenaline, and dopamine are catecholamines. These hormones and nerve-signaling chemicals help regulate stress responses, blood pressure, heart function, alertness, and metabolism.
\medterm{Catecholamines (Urine/Plasma)} Blood or urine tests measuring catecholamines or their breakdown products. They may help investigate pheochromocytoma or paraganglioma, but medicines, stress, posture, diet, and collection conditions can affect results.

\synonyms
Urine/Plasma

\medterm{Catecholamines - Urine} A urine test measuring catecholamines or their breakdown products over a specified period. It may help investigate pheochromocytoma, paraganglioma, or unexplained episodes of high blood pressure, sweating, headache, or palpitations.

\medterm{Category Five Hurricane} A Saffir--Simpson hurricane with sustained winds of at least 157 miles per hour (252 km/h), capable of catastrophic damage and prolonged loss of habitability.

\medterm{Category I Fetal Heart Tracing} A reassuring fetal heart-rate pattern with a baseline of 110--160 beats per minute, moderate variability, and no late or variable decelerations. It is usually consistent with adequate fetal oxygenation at that time.

\medterm{Category II Fetal Heart Tracing} A fetal heart-rate pattern that is not clearly normal or abnormal. It requires continued monitoring and assessment of the mother, fetus, labor, and changes over time because it may become reassuring or concerning.

\medterm{Category III Fetal Heart Tracing} An abnormal fetal heart-rate pattern that may indicate inadequate fetal oxygenation or acid buildup. It includes absent variability with recurrent late or variable decelerations or a true sinusoidal pattern and requires prompt clinical evaluation and action.

\medterm{Category Three Hurricane} A Saffir--Simpson hurricane with sustained winds of 111--129 miles per hour (178--208 km/h), capable of devastating damage.

\medterm{Category Two Hurricane} A Saffir--Simpson hurricane with sustained winds of 96--110 miles per hour (154--177 km/h), capable of extensive damage to structures and vegetation.

\medterm{Catenin} A family of proteins that helps neighboring cells attach and communicate. Beta-catenin also carries growth signals inside cells; abnormal activity can allow some cancers to grow or spread.
\medterm{Caterpillar Dermatitis} An itchy or painful skin reaction after contact with irritating or venomous caterpillar hairs or spines. Washing the area and removing visible hairs may help; eye exposure, breathing difficulty, or widespread symptoms need medical care.
\medterm{Caterpillars} The larval stage of butterflies and moths. Hairs or spines on some species can cause painful skin inflammation, eye injury, or allergic reactions after contact or inhalation.

\medterm{Catharsis} Catharsis is emotional release or expression; it is a historical psychological concept and not a specific medical treatment or diagnosis.
\medterm{Cathartic} A substance that promotes bowel emptying, usually by increasing stool water, intestinal movement, or both. Excessive or prolonged use can cause diarrhea, dehydration, electrolyte imbalance, or dependence on laxatives.
\synonyms
Purgative; laxative.

\medterm{Cathepsin} A family of enzymes that break down proteins inside cells, especially in lysosomes, the cell's recycling compartments. They help with immune activity, tissue remodeling, and bone turnover; abnormal activity may contribute to inflammation, bone loss, or cancer.
\medterm{Cathepsin D} An enzyme in lysosomes that breaks down proteins in an acidic environment. It also participates in cell death and immune processes; abnormal activity has been studied in cancer and nervous-system disease.
\medterm{Cathepsin H} An enzyme in lysosomes that breaks down proteins and helps activate some biologically active peptides. It contributes to normal cell maintenance and may be altered in certain diseases.

\medterm{Cathepsin K} An enzyme made by osteoclasts, the cells that remove old bone. It breaks down the protein framework of bone; excessive activity can contribute to bone loss and osteoporosis.

\medterm{Cathepsin L} An enzyme in lysosomes that breaks down proteins and helps process antigens for immune recognition. It also participates in tissue remodeling and can assist entry or spread of some infectious agents.

\medterm{Cathepsin Z} An enzyme found mainly in lysosomes that cuts proteins into smaller pieces. It helps cells process materials and may contribute to immune responses and inflammation.

\medterm{Cathepsins B} A group of protein-cutting enzymes, including cathepsin B, that help recycle proteins inside cells. Abnormal activity can contribute to inflammation, tissue remodeling, invasion of nearby tissue by tumors, and spread of cancer.

\medterm{Catheter} A thin flexible tube placed in a blood vessel, body cavity, or duct to deliver medicines or fluids, drain urine or other material, measure pressure, collect samples, or guide a procedure.

\medterm{Catheter, Foley} A urinary catheter that remains in the bladder and uses a small inflatable balloon to stay in place. It continuously drains urine into a collection bag but increases the risk of infection and blockage.

\medterm{Catheter-Associated UTI} A urinary infection associated with an indwelling catheter. Risk rises with prolonged use, and symptoms may include fever, lower abdominal or flank pain, or confusion; avoiding unnecessary catheters and removing them early lowers risk.

\medterm{Catheter-Related Infections} Infections caused by microorganisms entering along or contaminating an indwelling catheter. They may cause local redness, pain, drainage, fever, or bloodstream infection; treatment may require antibiotics and catheter removal or replacement.

\medterm{Catheter-Related UTI} A urinary tract infection associated with an indwelling urinary catheter, usually caused by bacteria entering or traveling along it. Diagnosis requires compatible symptoms and appropriately collected urine testing because bacteria may be present without infection.

\medterm{Catheterization} Placement of a catheter into a blood vessel, body cavity, duct, or other passage to drain, infuse, measure, collect samples, or perform a procedure. Risks include pain, bleeding, infection, blockage, and injury to nearby tissue.

\medterm{Catheterization (Urinary)} Insertion of a flexible tube through the urethra or through the lower abdominal wall into the bladder to drain urine. It may be intermittent, indwelling with a balloon, or suprapubic; infection and injury are possible risks.

\synonyms
Urinary

\medterm{Cathexis} The investment of emotional or psychological energy in a person, idea, object, or activity, a term used mainly in psychology and psychoanalytic theory.

\medterm{Cation} An electrically positively charged atom or molecule, such as sodium, potassium, calcium, or hydrogen. Cations help regulate fluid balance, nerve signals, muscle contraction, and the body's acid-base balance.

\synonyms
Positively charged ion.

\medterm{Cauda Equina} The bundle of lumbar, sacral, and coccygeal nerve roots that continues below the end of the spinal cord. These nerves control sensation and movement in the legs and help control the bladder, bowel, and sexual organs.

\synonyms
Cauda equina nerve roots; horse-tail nerve bundle.

\medterm{Cauda Equina Syndrome} A medical emergency caused by pressure on the bundle of nerves below the spinal cord. Severe low-back pain, leg weakness or numbness, numbness around the genitals or anus, or new bladder or bowel problems require immediate assessment.

\synonyms
Syndrome Cauda Equina (Cauda Equina Syndrome)

\medterm{Caudad} Toward the tail or lower end of the body; in human anatomy, generally toward the feet or away from the head.

\medterm{Caudal} Positioned toward the tail or lower part of the body; in human anatomy, usually toward the feet or away from the head.

\medterm{Caudal Dysplasia} A group of birth defects in which the lower spine and nearby structures do not develop fully. It may cause abnormalities of the legs, pelvis, bowel, bladder, or genital organs and is more common with maternal diabetes.
\medterm{Caudal Vein} A vein near the tail end of an animal's body. In human anatomy, the term is used mainly in embryology or when describing structures in comparative anatomy.

\medterm{Caudate Nucleus} A curved group of nerve cells deep in the brain that helps control movement, learning, memory, motivation, and behavior. Disease affecting it may contribute to movement or cognitive problems.

\medterm{Cauliflower Ear} A thick, irregular deformity of the outer ear caused when blood or fluid collects between its cartilage and covering, usually after blunt injury or repeated friction. Prompt drainage and compression can protect the cartilage and prevent permanent change.

\synonyms
Perichondrial Hematoma (Cauliflower Ear)

\medterm{Caulking Compound Poisoning} Poisoning or injury after contact with or swallowing a caulking product. Symptoms depend on its ingredients and may include eye, skin, or airway irritation, vomiting, drowsiness, or aspiration into the lungs.

\medterm{Caulobacter} A group of bacteria found mainly in freshwater and soil. Human infection is extremely rare but may occur opportunistically in people with severe illness or weakened immunity.

\medterm{Caulobacteraceae} A family of mainly freshwater and soil bacteria that includes Caulobacter. Most members are environmental and are not recognized as common causes of human disease.

\medterm{Causalgia} Severe persistent burning pain, sensitivity, and other sensory changes after injury to a peripheral nerve. It is now usually considered complex regional pain syndrome type II and may include swelling, color or temperature changes, and movement difficulty.

\synonyms
Complex regional pain syndrome type II; causalgic pain.

\medterm{Cause Of Death} The disease, injury, or toxic exposure that starts the chain of events leading directly to death. It differs from the immediate mechanism, such as respiratory failure, and from contributing conditions.


\medterm{Causes Symptoms Of Hormonal Imbalances Women} Hormone-related symptoms in women can result from thyroid, ovarian, adrenal, or pituitary disease, pregnancy, menopause, medicines, or weight changes. Possible symptoms include irregular bleeding, acne, hair changes, infertility, hot flashes, tiredness, and mood changes.

\medterm{Caustic} Able to chemically burn or destroy living tissue, as strong acids and alkalis can. Swallowing or contact with a caustic substance can injure the mouth, esophagus, stomach, skin, eyes, or airway and may be an emergency.
\medterm{Cauterization} Intentional destruction or sealing of tissue with heat, electricity, chemicals, or another energy source. It may control bleeding, remove abnormal tissue, or close small vessels and wounds.

\medterm{Cautery} Use of heat, electricity, chemicals, or another energy source to destroy or seal tissue, commonly to control bleeding or remove a small abnormal area.
\medterm{Cautery Device} A medical instrument that uses heat, electricity, or another energy source to seal bleeding vessels or destroy unwanted tissue. Use can cause burns or damage surrounding structures if poorly controlled.

\medterm{Caveola} A small pocket-like fold in a cell membrane that helps organize signals, move materials into or out of the cell, and respond to stretching or mechanical stress.

\medterm{Caveolins} Proteins that form or regulate caveolae, small pockets in cell membranes. They help organize cell signals, handle fats, move materials across membranes, and respond to mechanical stress.

\medterm{Cavernous Hemangioma} A cluster of enlarged, thin-walled blood vessels that may occur in the skin, brain, spinal cord, or other tissues. Symptoms depend on location and may result from bleeding, pressure, pain, or altered function.

\medterm{Cavernous Malformations} Clusters of unusually wide, fragile blood vessels, usually in the brain or spinal cord. They may cause no symptoms, but bleeding or pressure on nerve tissue can produce headaches, seizures, weakness, numbness, vision problems, or other neurologic changes.

\medterm{Cavernous Sinus} A large venous channel beside the pituitary gland that drains blood from the eye and face. It contains important eye-movement nerves and the internal carotid artery; infection or clotting there can be life-threatening.
\medterm{Cavernous Sinus Thrombosis} A dangerous blood clot, often infected, in the venous channel behind the eyes. It usually follows infection of the face, nose, sinuses, teeth, or orbit and may cause severe headache, fever, bulging eyes, double vision, or vision loss.

\medterm{Cavities} Holes or damaged areas in teeth caused by acids made by plaque bacteria from sugars and other carbohydrates. Decay can reach the inner pulp and cause pain or infection; prevention and treatment may include fluoride, cleaning, fillings, or root-canal care.

\synonyms
Dental Caries (Cavities)

\medterm{Cavities And Tooth Decay} Progressive loss of tooth structure caused by acids produced when plaque bacteria use sugars. Decay can move from enamel into dentin and the pulp, causing sensitivity, pain, infection, abscess, or tooth loss.
\medterm{Cavity} A hole or weakened area in a tooth caused by decay. Without treatment, it may reach the inner pulp and cause sensitivity, pain, infection, abscess, fracture, or tooth loss.

\synonyms
Dental cavity; carious lesion.

\medterm{Cavum} A cavity or hollow space in the body. The precise meaning depends on the anatomical phrase, such as cavum septi pellucidi, a fluid-containing space in the brain.
\medterm{CBC} See \textit{Complete Blood Count}.

\synonyms
Complete blood count; full blood count.

\medterm{CBC Blood Test} See \textit{Complete Blood Count}.

\medterm{CCP} Cyclic citrullinated peptide, a target used in a blood test for antibodies associated with rheumatoid arthritis. A positive anti-CCP result supports the diagnosis and may indicate greater risk of joint damage, but it is not diagnostic alone.

\medterm{CCTA} Alternate name; see \textit{Coronary Computed Tomography Angiography}.

\medterm{CCU} A hospital unit providing continuous heart monitoring and specialized treatment for people with serious acute cardiac problems, such as heart attack, dangerous rhythm disturbance, or severe heart failure.

\synonyms
Coronary care unit; cardiac care unit.

\medterm{Cd Antigens} Proteins or other markers on the surface of cells that help identify cell types, especially immune cells. Laboratory antibodies against these markers can classify cells, diagnose blood disorders, and guide selected treatments.

\medterm{CD4 Count} The number of CD4 T lymphocytes in a blood sample. It helps assess immune function and the risk of opportunistic infection in people with HIV; results are interpreted with viral load, symptoms, treatment, and trends over time.

\medterm{CD4-Positive T Cell} A type of white blood cell that coordinates immune responses. Different CD4 cells help activate other immune cells, support antibody production, defend mucosal surfaces, or limit excessive inflammation.

\medterm{CD8-Positive T Cell} A type of white blood cell that recognizes abnormal signals on infected or altered cells. Cytotoxic CD8 cells can destroy virus-infected, cancerous, or otherwise abnormal cells.

\medterm{Cea} An abbreviation for carcinoembryonic antigen, a blood marker used mainly to monitor selected cancers rather than to screen people without symptoms.
\medterm{CEA Assay} A laboratory test measuring carcinoembryonic antigen in a specimen, usually blood. It is mainly used to monitor selected known cancers and cannot diagnose or exclude cancer by itself.

\medterm{CEA Blood Test} A blood test measuring carcinoembryonic antigen, used mainly to monitor selected cancers such as colorectal cancer and assess treatment response or recurrence. Smoking, inflammation, and other noncancerous conditions can also raise it.

\medterm{Cecal Hemorrhage} Bleeding from the caecum, the first part of the large intestine. Causes include inflammation, infection, reduced blood flow, polyps, diverticular disease, or cancer and may produce blood in the stool or anemia.
\medterm{Cecal Infection} Infection or inflammation of the caecum caused by microorganisms, reduced blood flow, inflammatory bowel disease, or nearby disease. It may cause right-sided abdominal pain, fever, diarrhea, or tenderness.

\medterm{Cecal Neoplasm} A noncancerous or cancerous growth arising in the caecum. It may cause bleeding, anemia, abdominal pain, altered bowel habits, or blockage and is usually evaluated with colonoscopy and tissue sampling.

\medterm{Cecal Volvulus} Twisting of the caecum and nearby colon around the tissue that supports them, causing large-bowel blockage. It can produce sudden abdominal pain, swelling, vomiting, and inability to pass stool or gas and usually requires urgent treatment.

\medterm{Cecostomy} A surgically created opening from the caecum to the abdominal wall to release pressure, drain bowel contents, or divert stool. It may be temporary or permanent depending on the underlying problem.

\medterm{Cecum} The first pouch-like part of the large intestine, in the lower right abdomen where the small intestine joins the colon. The appendix arises from it; disease can include inflammation, tumors, bleeding, or twisting that blocks the bowel.

\medterm{Cedar Leaf Oil Poisoning} Toxic exposure to cedar-leaf oil, usually after swallowing or inhaling it. It may irritate the mouth, stomach, lungs, or nervous system and can cause vomiting, dizziness, seizures, or breathing problems; urgent poison-control advice is needed.
\medterm{Cedecea} A group of gram-negative bacteria found in the environment that rarely cause human disease. When infection occurs, it is usually opportunistic and affects people with severe illness, weakened immunity, or medical devices.

\medterm{Cedecea Davisae} A rare environmental bacterium that has caused bloodstream, respiratory, wound, and urinary infections. Disease is uncommon and usually occurs in people with serious underlying illness or impaired immunity.

\medterm{Cedecea Lapagei} A rare environmental bacterium that can occasionally cause bloodstream, respiratory, or other opportunistic infection, mainly in people with severe illness or weakened immune defenses.
\medterm{Cedecea Neteri} An uncommon gram-negative bacterium that rarely causes opportunistic human infection. A culture result must be interpreted with symptoms, the body site, and the quality of the specimen.

\medterm{Cediranib Maleate} The maleate salt of cediranib, an experimental medicine that blocks vascular endothelial growth factor receptors and may limit a tumor's blood-vessel growth. It has been studied in cancer trials but is not a routine treatment.

\medterm{Cefaclor} A second-generation cephalosporin antibiotic used for selected susceptible bacterial infections. It can cause allergic reactions, nausea, diarrhea, or antibiotic-associated colitis and should be used only when appropriate.

\medterm{Cefadroxil} An oral first-generation cephalosporin antibiotic used for selected susceptible skin, throat, urinary, and other infections. Possible effects include allergy, nausea, diarrhea, and antibiotic-associated colitis.

\medterm{Cefamandole} A second-generation cephalosporin antibiotic used in some settings for susceptible bacterial infections. It can cause allergy or diarrhea and may increase bleeding risk, especially in people with poor nutrition or liver disease.

\medterm{Cefatrizine} An older oral cephalosporin antibiotic used in some countries for susceptible bacterial infections. Its usefulness depends on local availability, bacterial resistance, allergy risk, and current treatment guidance.

\medterm{Cefazolin} An injectable first-generation cephalosporin antibiotic used for selected susceptible infections and commonly given before surgery to prevent infection. Allergy, diarrhea, and dose adjustment in kidney disease may be important.

\medterm{Cefdinir} An oral third-generation cephalosporin antibiotic used for selected susceptible respiratory, skin, and other infections. It may cause diarrhea, nausea, allergic reactions, or antibiotic-associated colitis and can interact with iron-containing products.
\medterm{Cefepime} An injectable fourth-generation cephalosporin used for serious infections caused by susceptible bacteria, including some hospital-acquired infections. Kidney-dose adjustment is important because accumulation can cause confusion, seizures, or other neurologic toxicity.

\medterm{Cefixime} An oral third-generation cephalosporin used for selected susceptible urinary, respiratory, ear, and other bacterial infections. It can cause diarrhea or allergy, and unnecessary use promotes antibiotic resistance.

\medterm{Cefmenoxime} An injectable third-generation cephalosporin used in some countries for susceptible bacterial infections. Possible effects include allergic reactions, diarrhea, and antibiotic-associated colitis.

\medterm{Cefmetazole} An injectable cephamycin antibiotic used in some countries for selected abdominal, pelvic, and other infections caused by susceptible bacteria. Allergy, diarrhea, and antibiotic-associated colitis are possible.

\medterm{Cefonicid} A long-acting second-generation cephalosporin antibiotic used for selected susceptible infections. Its use is limited in many places; dose selection depends on kidney function, local resistance, and the infection site.

\medterm{Cefoperazone} A third-generation cephalosporin active against many gram-negative bacteria and sometimes combined with a beta-lactamase inhibitor. It can cause allergy, diarrhea, bleeding, or impaired bile flow and requires appropriate monitoring.

\medterm{Cefotaxime} An injectable third-generation cephalosporin used for serious susceptible infections, including meningitis, pneumonia, bloodstream infection, and sepsis. Allergy, diarrhea, and antibiotic-associated colitis are possible.
\medterm{Cefotetan} An injectable cephamycin antibiotic used for selected serious infections and surgical prevention. It covers some anaerobic bacteria but can cause allergy, diarrhea, bleeding, or a reaction if alcohol is consumed around treatment.

\medterm{Cefotiam} A second-generation cephalosporin used in some countries for susceptible respiratory, urinary, skin, and other bacterial infections. Allergy, gastrointestinal effects, and antibiotic-associated colitis are possible.

\medterm{Cefoxitin} An injectable cephamycin antibiotic used for selected abdominal, pelvic, skin, and other infections, including some involving anaerobic bacteria. Allergy, diarrhea, and antibiotic-associated colitis can occur.

\medterm{Cefsulodin} An older cephalosporin with activity against Pseudomonas aeruginosa. It is used mainly in laboratory or historical contexts because resistance, limited availability, and newer treatments restrict routine clinical use.

\medterm{Ceftazidime} An injectable third-generation cephalosporin with activity against Pseudomonas aeruginosa and other susceptible gram-negative bacteria. It is used for selected serious infections and can cause allergy, diarrhea, or neurologic effects when it accumulates in kidney disease.
\medterm{Ceftazidime-Avibactam} A combination antibiotic in which avibactam blocks some bacterial enzymes that destroy ceftazidime. It is used for selected difficult gram-negative infections, including some resistant organisms; kidney-dose adjustment and susceptibility testing are important.

\medterm{Ceftibuten} An oral cephalosporin antibiotic used for selected susceptible respiratory, urinary, and other bacterial infections. It may cause diarrhea or allergy, and inappropriate use promotes resistance and antibiotic-associated colitis.

\medterm{Ceftizoxime} An injectable third-generation cephalosporin used for selected serious infections caused by susceptible bacteria, including some abdominal, pelvic, urinary, and anaerobic infections. Allergy, diarrhea, and antibiotic-associated colitis are possible.

\medterm{Ceftriaxone} An injectable third-generation cephalosporin used for many serious susceptible bacterial infections, including meningitis, pneumonia, bloodstream infection, and gonorrhea. Allergy, diarrhea, biliary problems, and antibiotic-associated colitis are possible.

\medterm{Cefuroxime} A second-generation cephalosporin available by mouth or injection for selected respiratory, ear, urinary, skin, and other bacterial infections. It can cause allergy, nausea, diarrhea, or antibiotic-associated colitis.

\medterm{Celiac Artery} A major branch of the abdominal aorta that supplies the upper digestive organs, including the stomach, liver, gallbladder, spleen, pancreas, and first part of the small intestine. It divides into the left gastric, splenic, and common hepatic arteries.

\medterm{Celiac Crisis} A rare, life-threatening severe flare of celiac disease causing profuse diarrhea or vomiting, dehydration, major electrolyte disturbances, and sometimes shock. It requires urgent hospital treatment and strict control of gluten exposure.

\medterm{Celiac Disease} An immune disorder triggered by gluten in wheat, barley, and rye in genetically susceptible people. It damages the small intestine and can cause diarrhea, bloating, weight loss, anemia, poor absorption, bone disease, or symptoms outside the intestine.

\synonyms
Nontropical Sprue

\medterm{Celiac Disease (Gluten Enteropathy)} Alternate index label for Gluten Enteropathy.

\medterm{Celiac Disease - Nutritional Considerations} Nutrition care for celiac disease requires lifelong avoidance of wheat, barley, rye, and cross-contaminated foods. Dietitians can help maintain adequate fiber and nutrients and correct deficiencies of iron, folate, vitamin B12, calcium, or vitamin D.


\medterm{Celiac Disease - Sprue} Another description of celiac disease, an immune reaction to gluten that damages the small intestine. It can cause diarrhea, bloating, weight loss, anemia, poor nutrient absorption, bone disease, and other symptoms.

\medterm{Celiac Plexus} A network of nerves in the upper abdomen that carries signals to and from the stomach, liver, pancreas, intestines, and other organs. It can be targeted with an injection for selected severe abdominal pain.

\medterm{Celiac Trunk} The first major branch of the abdominal aorta, supplying the upper digestive organs. It divides into the left gastric, splenic, and common hepatic arteries and is surrounded by nerves of the celiac plexus.

\medterm{Celiac Trunk Stenosis} Narrowing of the celiac trunk, caused by plaque or, in some people, compression by the diaphragm's median arcuate ligament. It may cause upper abdominal pain after meals and weight loss, although imaging findings do not always explain symptoms.

\medterm{Cell} The smallest living unit that can perform life processes. Human cells usually contain DNA in a nucleus and specialized internal structures that produce energy, make proteins, communicate, and maintain the body's tissues.

\medterm{Cell Adhesion} The attachment of cells to one another or to the supporting material around them. It helps maintain tissue structure, transmit signals, guide movement, and organize healing and development.

\medterm{Cell Aggregation} Clumping of cells through interactions between their surfaces or surrounding proteins. It is important in development, immune responses, wound healing, and formation of blood clots.

\medterm{Cell Body} The main part of a cell containing the nucleus and most internal structures. In a nerve cell, it excludes long extensions such as the axon and dendrites.

\medterm{Cell Count} The number of cells in a sample of blood, body fluid, tissue, or laboratory culture. It helps assess blood-cell levels, inflammation, infection, immune function, or abnormal cell growth.

\medterm{Cell Culture} Growth of cells outside the body in a controlled nutrient environment. It is used for research, medicine testing, diagnosis, production of biological products, and study of how cells respond to disease or treatment.

\medterm{Cell Cycle} The series of stages through which a cell grows, copies its DNA, checks for damage, and divides into two cells. Disruption of cycle control can contribute to developmental disorders and cancer.

\medterm{Cell Death} The loss of a cell's ability to remain alive. It may be a controlled process that removes old or damaged cells, or an uncontrolled result of injury, infection, toxins, lack of oxygen, or severe disease.

\medterm{Cell Dedifferentiation} Loss of some specialized features by a mature cell, making it more like an immature cell. It can occur during repair or development and may contribute to abnormal growth in some cancers.

\medterm{Cell Density} The number of cells in a given area or volume of tissue or fluid. It can change with growth, inflammation, healing, infection, or cancer and is measured in laboratory samples and tissue studies.

\medterm{Cell Differentiation} The process by which an unspecialized cell develops the structure and function of a specialized cell type, such as a nerve, muscle, blood, or skin cell.

\medterm{Cell Division} The process by which one cell produces new cells. Mitosis makes most body cells for growth and repair, while meiosis makes reproductive cells with half the usual number of chromosomes.

\medterm{Cell Fusion} Joining of two cells or their outer membranes to form one larger cell, sometimes with several nuclei. It occurs in normal development and repair and can also occur during infection or disease.

\medterm{Cell Growth} Increase in a cell's size and contents through production of proteins, membranes, energy stores, and other components. Abnormal growth signals can contribute to tissue overgrowth or cancer.

\medterm{Cell Maturation} Development of an immature cell into a mature cell with a specialized structure and function. It is essential for producing healthy blood, nerve, muscle, immune, and other body cells.

\medterm{Cell Migration} Movement of cells through tissue or across a surface in response to chemical, mechanical, or environmental signals. It is important in development, immune defense, wound healing, and cancer spread.

\medterm{Cell Motility} The ability of cells to move through or across a surface using their internal framework, attachment structures, or specialized projections such as cilia and flagella. It supports development, immunity, wound healing, and cancer spread.

\synonyms
Cell Locomotion

\medterm{Cell Movement} Directed or random movement of a cell through tissue or across a surface. It is important in development, immune responses, wound repair, inflammation, and the spread of cancer.

\medterm{Cell Nucleus} The membrane-covered compartment containing most of a cell's DNA. It helps control gene activity, growth, division, repair, and production of proteins needed for normal cell function.

\medterm{Cell Pellet} The concentrated layer of cells collected at the bottom of a tube after spinning a sample in a centrifuge. It may be removed for microscopic, genetic, chemical, or other laboratory testing.


\medterm{Cell Physiology} The study of how cells obtain energy, transport substances, communicate, respond to signals, repair themselves, and perform specialized functions.

\medterm{Cell Polarity} The unequal arrangement of structures and activities within a cell, giving it distinct sides or directions. It enables organized transport, signaling, movement, development, and tissue structure.

\medterm{Cell Proliferation} Increase in cell number through repeated division. Normal proliferation supports growth and repair, while excessive or uncontrolled proliferation can cause tissue overgrowth, benign tumors, or cancer.

\medterm{Cell Secretion} Release of substances made by cells, such as hormones, enzymes, nerve chemicals, mucus, or immune signals. Secretion supports communication, digestion, protection, movement, and regulation of body processes.

\medterm{Cell Senescence} A state in which a cell remains active but permanently stops dividing, often because of aging, stress, or DNA damage. Senescent cells can help prevent cancer but may contribute to inflammation when they accumulate.

\synonyms
Cellular senescence; replicative senescence.

\medterm{Cell Sorting} Separation of cells into groups according to features such as size, density, surface markers, or fluorescence. It is used in research, diagnosis, blood-cell analysis, and preparation of selected cells for treatment.

\medterm{Cell Structure} The arrangement of a cell's membrane, cytoplasm, nucleus, and internal structures. These parts work together to produce energy, make proteins, communicate, divide, and carry out specialized functions.

\medterm{Cell Surface} The outer membrane of a cell and the molecules attached to it. It controls movement of substances, attachment, recognition, communication, and interactions with immune cells and surrounding tissue.

\medterm{Cell Survival} The ability of a cell to remain alive and functional despite injury, lack of nutrients, toxins, infection, or other stress. It depends on the cell type, the severity and duration of the stress, and the cell's repair systems.

\medterm{Cell Therapy} Treatment that uses living cells to replace damaged cells, restore a function, or change an immune response. Uses range from established blood-forming stem-cell transplantation to experimental treatments; benefits and risks depend on the cells and disease.

\medterm{Cell Transplantation} Transfer of living cells into a person to replace missing cells, restore tissue function, or modify immunity. Blood-forming stem-cell transplantation is established for selected diseases; rejection, infection, graft-versus-host disease, and other risks may occur.

\medterm{Cell Wall} A rigid outer layer surrounding plant, fungal, and bacterial cells. It provides shape and protection; human and animal cells have a flexible membrane but no cell wall.

\medterm{Cell-Cell Junction} A specialized connection between neighboring cells that provides attachment, communication, or controlled passage of substances. Examples include tight junctions, desmosomes, and gap junctions.

\medterm{Cell-Cell Signaling} The ways cells send and receive messages through direct contact, electrical activity, or chemicals such as hormones and neurotransmitters. Signaling coordinates growth, movement, immune responses, repair, and other body functions.

\medterm{Cell-Free Fetal DNA Screening} A prenatal blood test that analyzes small DNA fragments from the placenta to estimate the chance of selected chromosome conditions. It is a screening, not a diagnostic, test; a high-risk or unclear result needs counseling and confirmatory testing.

\medterm{Cell-Free System} A laboratory preparation containing selected cell components but no intact living cells. It is used to study biochemical reactions, make proteins, test medicines, or examine biological processes under controlled conditions.

\medterm{Cell-Matrix Junction} A specialized connection anchoring a cell to the supporting material around it. It also transmits signals that influence cell shape, movement, survival, and tissue organization.

\medterm{Cell-Mediated Immune Response} Alternate index label; see \textit{Cell-mediated immunity}.

\medterm{Cell-Mediated Immunity} Immune defense carried out mainly by T lymphocytes and other cells rather than by antibodies. It is important for controlling infected cells, cancer cells, and organisms that live inside cells.

\synonyms
Cell-Mediated Immune Response

\medterm{Cellophane} A thin transparent film made from cellulose. In medicine it may be used in some dressings, packaging, or laboratory techniques; its clinical meaning depends on the context.
\medterm{Cells Stem (Stem Cells)} Alternate index label; see \textit{Stem Cell}.

\medterm{Cellular} Relating to cells or to processes occurring within or between cells.

\medterm{Cellular Adaptation} A reversible change in a cell's size, number, structure, or function in response to altered demands or mild injury. Adaptation may help a cell survive but can become harmful if the stress continues.

\medterm{Cellular Immunity} Immune defense carried out mainly by T cells, natural killer cells, macrophages, and other cells rather than antibodies alone. It helps destroy infected or abnormal cells and control inflammation.

\medterm{Cellular Metabolic Process} A chemical reaction or linked pathway in a cell that produces energy, builds cell components, or breaks down nutrients and waste.

\medterm{Cellular Respiration} The process cells use to release energy from food. Most cells use oxygen to convert nutrients into ATP, the cell's usable energy, while limited energy can be made without oxygen.

\medterm{Cellular Swelling} Early reversible cell injury in which failure of membrane pumps causes sodium and water to collect inside the cell. It commonly follows reduced blood flow, low oxygen, toxins, or other stresses.

\synonyms
hydropic change.

\medterm{Cellular Vesicle} A small membrane-bound compartment that transports, stores, or releases substances inside or outside a cell. Vesicles help move proteins, fats, hormones, and waste.

\medterm{Cellulite} A common harmless dimpled or uneven appearance of skin, usually on the thighs, hips, or buttocks, caused by the arrangement of fat, connective tissue, and skin. It is different from cellulitis, a bacterial infection.

\medterm{Cellulitis} A bacterial infection of the deeper skin and tissue beneath it, causing spreading redness, warmth, swelling, tenderness, and sometimes fever or chills. It commonly follows a break in the skin and can spread rapidly or enter the bloodstream.

\medterm{Cellulomonas} A group of bacteria found mainly in soil, plants, and other environmental materials. They break down cellulose; human infection is rare and usually occurs opportunistically in people with severe illness or weakened immunity.

\medterm{Cellulomonas Hominis} A rare environmental bacterium that has occasionally caused human infection, usually in people with weakened immunity or serious underlying illness. Its significance depends on symptoms, the body site, and laboratory identification.

\medterm{Cellulose} A complex carbohydrate forming the structure of plant cell walls. Humans cannot digest it, but it is dietary fiber that adds bulk to stool and can support regular bowel movements when enough fluid is consumed.

\medterm{Cellulosimicrobium} A group of environmental bacteria found mainly in soil. Human infection is rare but may affect people with serious illness, implanted devices, or weakened immune defenses and can involve the blood, heart, or other tissues.

\medterm{Cellulosome} A group of enzymes joined together by some microorganisms to break down cellulose in plant cell walls. It helps the organism digest plant material and is studied for biotechnology and biofuel production.

\medterm{Celosomy} A rare birth defect involving incomplete closure of the body wall, allowing abdominal or chest organs to protrude outside the body. The term is used for severe omphalocele-like malformations and requires specialist newborn care.

\medterm{Cement} A material used to fill spaces or attach an artificial joint or other prosthesis to bone. Dental cement bonds or protects dental work; the appropriate type depends on the procedure.
\medterm{Cemento-Enamel Junction} The boundary where tooth enamel covering the crown meets cementum covering the root. Its shape varies, and an exposed junction may be sensitive or more vulnerable to decay.
\medterm{Cementoma} A usually noncancerous jaw lesion containing cementum-like mineralized tissue. Its behavior and treatment depend on the specific subtype, location, size, symptoms, and relationship to nearby teeth.

\medterm{Cementum} A thin mineralized layer covering a tooth root. Fibers of the periodontal ligament attach to it and help anchor the tooth in its socket; it can be lost with gum disease.

\medterm{Cenesthesia} The general awareness of one’s body, including its condition, position, and internal sensations. Disturbance may occur in some neurologic or psychiatric illnesses and can alter how bodily sensations are experienced.
\medterm{Centenarian} A person aged 100 years or older. Centenarians are a diverse group; some remain independent, while others need substantial support, and age alone does not determine health or cognitive ability.

\medterm{Center} A central point, region, or facility organized around a function. In medicine, the term may describe an anatomical location, a treatment facility, or the middle of a measured structure.
\medterm{Centigram} A metric unit of mass equal to one hundredth of a gram, or 10 milligrams, used for small quantities of medicines or specimens.

\medterm{Centimeter} A metric unit of length equal to one hundredth of a meter, commonly used in clinical measurements such as wound size, tumor dimensions, and anatomical distances.

\medterm{Centimeter (Cm)} A unit of length equal to one hundredth of a meter, or ten millimeters, commonly used to measure body structures and objects.

\medterm{Centipede} A many-legged arthropod whose bite can inject venom and cause local pain, swelling, and redness; severe allergic reactions are uncommon but possible.

\medterm{Centipoise} A unit of dynamic viscosity equal to one millipascal-second, used to describe the flow resistance of blood, fluids, medicines, and other materials.

\medterm{Central} Located near the middle of the body or an organ, rather than near its edge. The word may describe a body structure, nervous-system area, blood line, symptom, or pattern of disease.

\medterm{Central (Brain) Fatigue} A sense of reduced mental energy, concentration, alertness, or motivation arising primarily from central nervous-system processes.

\synonyms
Central fatigue; mental fatigue.

\medterm{Central Apnea} A pause in breathing caused by absent or reduced respiratory effort from impaired brainstem signaling, rather than upper-airway obstruction.

\medterm{Central Auditory Processing Disorder} Difficulty interpreting sounds in the brain despite adequate hearing. People may struggle to understand speech in noise, follow rapid instructions, identify sounds, or distinguish similar words.

\medterm{Central Core Disease} An inherited muscle disorder, often caused by an RYR1 gene change, that may cause low muscle tone, delayed motor development, and weakness. Some affected people can develop life-threatening overheating with certain anesthetics.

\synonyms
CCD; Core Myopathy; RYR1-Related Central Core Disease

\medterm{Central Deafness} Severe hearing loss caused by damage to sound-processing pathways or areas of the brain, rather than the outer, middle, or inner ear. Hearing tests may show poor understanding despite detected sounds.

\medterm{Central Diabetes Insipidus} Excessive urination and intense thirst caused by too little antidiuretic hormone from the brain. The kidneys cannot conserve enough water, so dehydration and high blood sodium may develop.

\medterm{Central Hypothyroidism} A condition with too little thyroid hormone because the brain's pituitary gland or hypothalamus does not provide enough stimulation to the thyroid.

\medterm{Central Line} A catheter inserted into a large vein with its tip in or near the heart, used to give medications or nutrition, draw blood, or monitor central venous pressures.

\medterm{Central Line Infections - Hospitals} An infection related to a central venous catheter, including a central-line-associated bloodstream infection, in which microorganisms enter the bloodstream from or through the catheter.

\medterm{Central Line Placement} Insertion of a catheter into a large vein, commonly in the neck, chest, or groin, with its tip near the heart. It allows long-term medicines, nutrition, fluids, blood sampling, or monitoring.

\medterm{Central Line-Associated Bloodstream Infection} A bloodstream infection linked to a central venous catheter. It may cause fever, chills, or low blood pressure and is evaluated with blood cultures; treatment usually requires antibiotics and sometimes catheter removal.

\medterm{Central Nervous System} The brain and spinal cord. It receives and interprets information, controls movement and many automatic body functions, and supports thought, memory, emotion, and regulation of hormones.
\medterm{Central Nervous System Lymphoma} A lymphoma involving the brain, spinal cord, cerebrospinal fluid, or eye that can cause neurologic or visual symptoms.

\medterm{Central Nervous System Vascular Malformations} Abnormal groups or connections of blood vessels in or around the brain or spinal cord. Some cause no symptoms, while others can cause headaches, seizures, weakness, or bleeding; treatment depends on the type, size, location, and risk.

\medterm{Central Neurofibromatosis} Alternate index label; see \textit{NF2-related schwannomatosis (neurofibromatosis type 2)}.

\medterm{Central Pontine Myelinolysis} Brain injury caused most often by correcting very low blood sodium too quickly. Damage in the pons may cause weakness, trouble speaking or swallowing, confusion, reduced consciousness, or locked-in syndrome.

\medterm{Central Precocious Puberty} Central precocious puberty begins when the hypothalamus and pituitary activate the normal puberty pathway unusually early. It causes early secondary sexual characteristics and accelerated bone maturation; evaluation identifies causes and treatment may delay progression when appropriate.

\medterm{Central Retinal Artery Occlusion} A sudden blockage of the main artery supplying the retina, causing severe, usually painless loss of vision in one eye. It is an eye stroke requiring emergency assessment for related vascular disease.

\medterm{Central Serous Chorioretinopathy} A condition in which fluid collects beneath the retina, usually near the macula, causing blurred or distorted central vision. It may be linked to corticosteroids or stress and often needs eye-specialist follow-up.

\medterm{Central Serous Choroidopathy} A disorder in which fluid accumulates beneath the neurosensory retina, usually near the macula, causing blurred or distorted central vision. It is associated with corticosteroid exposure and stress-related factors; persistent or recurrent disease requires ophthalmic evaluation.

\medterm{Central Sleep Apnea} Repeated pauses in breathing during sleep because the brain temporarily stops sending signals to breathe, without physical airway blockage. It may occur with heart failure, high altitude, or opioid use.

\synonyms
Sleep Apnea, Central

\medterm{Central Venous Catheter} A thin tube ending in a large vein near the heart. It provides reliable access for prolonged or irritating treatments and monitoring but can cause infection, blood clots, bleeding, or injury.

\medterm{Central Venous Catheter - Dressing Change} Replacement of the sterile dressing covering a central venous catheter using aseptic technique to keep the insertion site clean, dry, and protected from contamination and infection.

\medterm{Central Venous Catheter - Flushing} Instillation of a prescribed sterile solution through a central venous catheter to maintain patency and reduce blockage, using the device-specific technique and volume.

\medterm{Central Venous Catheters - Ports} Implantable venous access ports are central venous catheters placed beneath the skin, accessed with a special needle to deliver medication, fluids, blood products, or obtain blood samples over prolonged treatment.

\medterm{Central Venous Line} A vascular catheter whose tip lies in a major central vein, commonly the superior vena cava, for reliable long-term or high-concentration intravenous access.

\medterm{Central Venous Line - Infants} A catheter whose tip lies in a large central vein of an infant, used for irritant medicines, nutrition, fluids, blood sampling, or monitoring; infection and thrombosis are key risks.

\medterm{Central Venous Pressure} Pressure measured in or near the right atrium, usually through a central venous catheter. It reflects
right-sided filling conditions but is affected by ventilation, heart function, vascular tone,
and measurement technique; it should not be used alone to guide fluid treatment.

\medterm{Central Vision} Detailed vision directed at the point of gaze, provided primarily by the macula and fovea; it is needed for reading and recognizing faces.

\medterm{Centriacinar Emphysema} A subtype of emphysema in which destruction begins in the respiratory bronchioles and
extends peripherally. It predominantly affects the upper lobes and is strongly
associated with cigarette smoking. It is the most common form of emphysema and often
coexists with chronic bronchitis.

\medterm{Centric-Fusion Translocation} A Robertsonian translocation in which the long arms of two acrocentric chromosomes fuse near their centromeres, leaving a combined chromosome and usually losing the short arms.

\medterm{Centrifugation} Spinning a specimen at high speed to separate components by density, such as cells from plasma or precipitated material from a liquid.

\medterm{Centrifuge} Spins samples to separate components by density. Used for blood processing, urine
analysis, tissue preparation. Speed: RPM or RCF.

\medterm{Centrilobular Emphysema} A type of permanent lung damage in which the small airways and nearby air sacs are destroyed, usually from long-term smoking or inhaled irritants. It often affects the upper lungs and can cause breathlessness, cough, wheezing, and reduced ability to exercise.

\synonyms
Centriacinar Emphysema; Proximal Acinar Emphysema

\medterm{Centriole} A cylindrical microtubule structure in the centrosome that helps organize the mitotic spindle and, in some cells, forms the basal body of a cilium or flagellum.

\medterm{Centromere} The constricted chromosome region where sister chromatids remain attached and the kinetochore forms for accurate chromosome segregation during cell division.

\medterm{Centronuclear Myopathy} A group of inherited muscle diseases in which muscle fibers develop abnormally. Symptoms range from severe weakness and breathing problems at birth to milder adult-onset weakness, often affecting the shoulders, hips, eyelids, eye movements, face, or feet.

\synonyms
CNM; Myotubular Myopathy (General)

\medterm{Centrosome} The main microtubule-organizing center of an animal cell, containing a pair of centrioles and coordinating the spindle during mitosis.

\medterm{Centrum} Centrum means a central part, especially the main body of a vertebra around which the neural arch develops.
\medterm{Cephacetrile} An older first-generation cephalosporin antibiotic that inhibits bacterial cell-wall synthesis; it is not commonly used in current practice.

\medterm{Cephalalgia} Headache or pain in the head, including pain arising from cranial or related structures.
It is a symptom with many possible primary and secondary causes.

\medterm{Cephalexin} An oral first-generation cephalosporin that inhibits bacterial cell-wall synthesis and treats selected skin, respiratory, urinary, and other susceptible infections.

\medterm{Cephalgia} Pain in the head, commonly called a headache. It may arise from a primary headache disorder or from infection, injury, medication, high blood pressure, or another illness.

\medterm{Cephalhaematoma} A collection of blood beneath the membrane covering a newborn's skull bone, usually caused by birth-related pressure or injury. It is limited by skull sutures and usually resolves without treatment.

\medterm{Cephalic} Relating to the head or toward the head; in obstetrics, a cephalic presentation means the fetus presents head first.

\medterm{Cephalic Vein} A superficial vein on the lateral side of the upper limb that drains the hand and forearm and joins the axillary vein near the shoulder.

\medterm{Cephalohematoma} A collection of blood beneath the membrane covering a newborn's skull bone, usually caused by birth-related pressure or injury. It does not cross skull sutures and generally clears gradually.

\medterm{Cephalometry} Measurement of the skull and facial structures, usually on standardized radiographs, used in orthodontic planning and assessment of craniofacial growth.

\medterm{Cephalopelvic Disproportion} A situation in which the fetal head cannot pass easily through the mother's pelvis because of their relative size or shape. It may prevent vaginal delivery, but cannot always be confirmed before labor.

\medterm{Cephaloridine} An older injectable first-generation cephalosporin antibiotic whose use declined because of nephrotoxicity and the availability of safer alternatives.

\medterm{Cephalosporin} A beta-lactam antibiotic that kills susceptible bacteria by blocking construction of their protective cell wall. Different generations act against different bacteria; allergy, diarrhea, and resistance are possible.

\medterm{Cephalosporins} Cephalosporins are beta-lactam antibiotics that inhibit bacterial cell-wall synthesis; generations differ in spectrum, and allergy, diarrhea, and resistance are important considerations.

\medterm{Cephalothin} An older injectable first-generation cephalosporin that inhibits bacterial cell-wall synthesis and was used for susceptible serious infections.

\medterm{Cephapirin} A first-generation cephalosporin antibiotic formerly used for susceptible bacterial infections; current availability and indications vary by country.

\medterm{Cephradine} A first-generation cephalosporin antibiotic related to cephalexin that inhibits bacterial cell-wall synthesis and treats selected susceptible infections.

\medterm{Ceramide} A type of fat that forms part of cell membranes and helps the outer skin retain water and block irritants. Ceramides also participate in signals controlling cell growth and survival.

\medterm{Cercaria} A cercaria is a free-swimming larval stage of a trematode parasite; contact with some species can cause swimmer’s itch or other infection.
\medterm{Cercarial Dermatitis} Alternate index label; see \textit{Swimmer's itch}.

\medterm{Cerclage} A procedure in which a strong stitch is placed around the cervix to help prevent it opening too early during pregnancy. It may be placed through the vagina or abdomen and requires obstetric follow-up.

\medterm{Cerebellar} Cerebellar means relating to the cerebellum, which coordinates balance, posture, eye movements, and skilled movement.
\medterm{Cerebellar Ataxia} Poor coordination caused by disease or injury affecting the cerebellum or its connections. It may cause an unsteady walk, clumsy movements, slurred speech, tremor, or abnormal eye movements.

\medterm{Cerebellar Cortex} The folded outer layer of the cerebellum containing Purkinje and granular neurons that coordinates balance, posture, timing, and precision of movement.

\medterm{Cerebellar Disease} A disorder of the cerebellum causing impaired coordination, gait ataxia, dysmetria, tremor, abnormal eye movements, or speech changes.

\medterm{Cerebellar Dysmetria} Inaccurate control of the distance or force of a movement because of cerebellar dysfunction. A person may reach past a target, stop short, or have difficulty judging movement size.

\medterm{Cerebellar Hemangioblastoma} A usually benign, highly vascular tumor of the cerebellum that may cause headache, imbalance, vomiting, or hydrocephalus and can occur with von Hippel-Lindau disease.

\medterm{Cerebellar Hypoplasia} Incomplete development of the cerebellum, usually beginning before birth. It may cause poor coordination, abnormal eye movements, seizures, delayed development, or problems with balance and movement.

\medterm{Cerebellar Neoplasm} A benign or malignant tumor involving the cerebellum; symptoms may include poor coordination, imbalance, abnormal eye movements, headache, vomiting, or hydrocephalus.

\medterm{Cerebellar Nuclei} Deep clusters of cerebellar neurons that receive processed signals from the cerebellar cortex and send coordinated motor output to the brainstem and thalamus.

\medterm{Cerebellar Tumor} A growth in or near the cerebellum that may cause poor coordination, unsteady walking, abnormal eye movements, vomiting, or pressure buildup.

\medterm{Cerebellar Vermis} The midline portion of the cerebellum that helps control trunk posture, gait, balance, and coordination of eye and head movements.

\medterm{Cerebellitis} Inflammation of the cerebellum, often after an infection and sometimes from an immune reaction. It may cause sudden poor coordination, unsteady walking, slurred speech, abnormal eye movements, headache, or reduced alertness.

\medterm{Cerebellopontine Angle} The angle-shaped region at the junction of the cerebellum, pons, and petrous temporal
bone. Lesions in this space may affect cranial nerves and commonly include vestibular
schwannoma.

\medterm{Cerebellopontine Angle Tumor} Lesion at the cerebellopontine angle, an anatomic space between the cerebellum and
pons, medial to the temporal lobe and lateral to the brainstem, containing CN V
(trigeminal), VII (facial), VIII (vestibulocochlear), and the anterior inferior
cerebellar artery (AICA).

\medterm{Cerebellum} Largest brain portion (approximately 10 percent of volume but over 50 percent of
neurons) in the posterior cranial fossa behind the brainstem. Two hemispheres and a
vermis, divided into vestibulocerebellum (balance, eye movements), spinocerebellum
(posture, gait, coordination), and cerebrocerebellum (motor planning, cognition).

\medterm{Cerebral} Relating to the cerebrum or, more broadly, to the brain.

\medterm{Cerebral Amyloid Angiopathy} Deposition of amyloid beta in the walls of small and medium cerebral and leptomeningeal
arteries. It commonly causes recurrent lobar intracerebral hemorrhage in older adults
and may also cause transient focal neurological episodes or cognitive decline.

\medterm{Cerebral Aneurysm} An abnormal focal enlargement of an artery in or around the brain. It may remain
asymptomatic or rupture and cause subarachnoid hemorrhage. Risk factors include hypertension, smoking, certain inherited
disorders, and family history.

\medterm{Cerebral Angiography} An X-ray test in which contrast dye is put into an artery to make the blood vessels of the brain visible. It can show aneurysms, abnormal artery-to-vein connections, narrowing, or blockage and is usually done with a thin tube called a catheter.

\medterm{Cerebral Aqueduct} A narrow channel through the midbrain connecting the third and fourth brain ventricles. It allows cerebrospinal fluid to flow, and blockage can enlarge the ventricles and increase pressure.

\medterm{Cerebral Arteriovenous Malformation} An abnormal direct connection between cerebral arteries and veins without a normal
capillary bed. It may cause intracranial hemorrhage, seizures, headache, or focal
neurological deficits and is evaluated with brain imaging and angiography.

\medterm{Cerebral Artery} An artery that carries blood to part of the brain. A blockage can cause an ischemic stroke, while a break in the vessel can cause bleeding in or around the brain; the effects depend on which artery is involved.

\medterm{Cerebral Atrophy} Loss or shrinkage of brain tissue caused by aging, neurodegenerative disease, injury, infection, stroke, or another condition. Effects vary from none to memory loss, weakness, seizures, or impaired thinking.
\medterm{Cerebral Autoregulation} The intrinsic ability of cerebral arterioles to maintain relatively stable cerebral
blood flow across a range of perfusion pressures. It may be impaired by severe brain
injury, ischemia, hypercapnia, or anesthetic and vasoactive drugs.

\medterm{Cerebral Calcification} Abnormal calcium deposition in brain tissue, with causes including inherited, metabolic, infectious, vascular, and age-related disorders.

\medterm{Cerebral Cavernous Malformation} A low-flow vascular malformation composed of dilated thin-walled vascular spaces in the brain or spinal cord. It may cause seizures, headaches, or hemorrhage, although some lesions remain asymptomatic.

\medterm{Cerebral Cavernous Malformations} Low-flow vascular malformations composed of dilated, thin-walled capillaries without
intervening brain parenchyma, often multiple and familial (autosomal dominant with
mutations in CCM1, CCM2, or CCM3 genes). They may present with seizures, focal
neurological deficits, or intracerebral hemorrhage.

\medterm{Cerebral Cortex} The outermost layer of the cerebrum, composed of gray matter (neuronal cell bodies)
folded into gyri (ridges) and sulci (grooves), greatly increasing surface area.

\medterm{Cerebral Crus} A group of nerve fibers in the middle part of the brain that carries movement commands down from the brain's outer layer. Damage can affect strength, movement, coordination, or control on the opposite side of the body.

\medterm{Cerebral Diplegia} Cerebral diplegia is a form of cerebral palsy with predominant weakness and spasticity of both legs, often greater than in the arms.
\medterm{Cerebral Edema} Swelling caused by excess fluid in or around brain tissue. It can raise pressure inside the skull, reduce blood flow, and cause headache, vomiting, confusion, seizures, coma, or life-threatening brain displacement.

\medterm{Cerebral Embolism} Alternate index label; see \textit{Embolic stroke (brain embolism)}.

\medterm{Cerebral Haemorrhage} Cerebral haemorrhage is bleeding within brain tissue or an adjacent intracranial space, causing pressure, tissue injury, and neurologic deficits.
\medterm{Cerebral Hemisphere} One of the two large halves of the cerebrum, each containing specialized but interconnected cortical regions.

\medterm{Cerebral Hemorrhage} Bleeding into brain tissue, usually from a damaged blood vessel. It can injure nearby brain cells, cause swelling and pressure, and produce sudden weakness, speech trouble, headache, reduced alertness, or death.

\medterm{Cerebral Hypoxia} Too little oxygen reaching the brain, often because of cardiac arrest, severe breathing failure, low blood pressure, or choking. Prolonged oxygen shortage can permanently injure brain cells.

\medterm{Cerebral Infarction} Permanent brain injury caused by blocked blood flow in a brain artery. The blockage may result from a clot formed in a narrowed artery or a clot traveling from the heart or another vessel.

\medterm{Cerebral Ischemia} Reduced blood flow and oxygen reaching brain tissue. It may cause temporary symptoms, such as a transient ischemic attack, or permanent injury if the interruption lasts long enough.
\medterm{Cerebral Lobes} Major regions of the cerebrum with different but overlapping functions: the frontal lobe supports planning and movement, the parietal sensation and spatial awareness, the temporal hearing and memory, and the occipital vision.

\medterm{Cerebral Malaria} A life-threatening form of malaria in which parasites cause brain dysfunction, often with confusion, seizures, abnormal behavior, or coma. It is usually caused by Plasmodium falciparum and requires urgent treatment.
\medterm{Cerebral Palsy} A group of permanent disorders of movement and posture development causing activity
limitation, attributed to non-progressive disturbances occurring in the developing fetal
or infant brain. It is the most common motor disability in childhood, affecting
approximately 1 in 500 live births.

\synonyms
Palsy, Cerebral
Palsy Cerebral (Cerebral Palsy)


\medterm{Cerebral Peduncle} One of two bundles of nerve fibers in the midbrain that carry movement signals from the cerebrum to the brainstem and spinal cord. Damage can cause weakness or loss of coordinated movement on the opposite side.

\medterm{Cerebral Perfusion Pressure} The pressure pushing blood through the brain, estimated from average arterial pressure minus pressure inside the skull. It falls when blood pressure drops or brain swelling raises skull pressure.

\medterm{Cerebral Salt Wasting} Excessive loss of sodium and water in urine associated with brain injury or disease. It can cause low blood sodium and dehydration and requires careful replacement guided by clinical and laboratory monitoring.

\medterm{Cerebral Thrombosis} Formation of a blood clot within a cerebral artery, reducing blood flow and potentially causing an ischemic stroke.

\medterm{Cerebral Tumor} A neoplasm arising within or involving the brain, producing symptoms according to its
location, growth rate, and effect on intracranial pressure.

\medterm{Cerebral Tumours} Cerebral tumours are abnormal growths in the cerebrum, which may be benign or malignant and cause seizures, headache, cognitive change, or focal deficits.
\medterm{Cerebral Vasospasm} Abnormal narrowing of cerebral arteries that can reduce brain perfusion and cause delayed cerebral ischemia, especially after aneurysmal subarachnoid hemorrhage.

\medterm{Cerebral Venous Sinus Thrombosis} A blood clot in a large vein or drainage channel of the brain. It can cause headache, seizures, visual problems, weakness, brain bleeding, or increased pressure inside the skull.

\medterm{Cerebral Ventricle} One of four connected spaces inside the brain that contain cerebrospinal fluid. The fluid cushions the brain and spinal cord, carries nutrients and waste, and flows through the ventricles before being absorbed around the brain.

\medterm{Cerebral Ventriculography} Imaging of the brain’s fluid-filled ventricles to assess their size, shape, connections, or blockage. Modern scans usually use computed tomography or magnetic resonance imaging rather than older invasive methods.

\medterm{Cerebral Ventriculomegaly} Enlargement of one or more fluid-filled spaces inside the brain. It may result from blocked fluid drainage, loss of brain tissue, or developmental variation; its importance depends on the cause and severity.
\medterm{Cerebral Visual Impairment} Reduced vision caused by damage to the brain's visual pathways or visual-processing areas rather than the eyes. It may cause poor visual attention, difficulty recognizing objects, or inconsistent use of vision.

\medterm{Cerebriform Nevus Sebaceous} A variant of nevus sebaceous presenting as a soft, cerebriform (brain-like),
pedunculated or sessile nodule, often found on the scalp or face. It is a benign adnexal
hamartoma composed of sebaceous glands, but like other nevus sebaceous lesions, it
carries a low risk of secondary neoplasms and may warrant monitoring or excision.

\medterm{Cerebritis} Inflammation and infection of brain tissue before a definite abscess forms, usually caused by bacteria, fungi, or parasites. It can cause headache, fever, seizures, or neurologic deficits and requires urgent imaging and antimicrobial treatment.

\medterm{Cerebroplacental Ratio} A calculation from fetal Doppler ultrasound comparing blood-flow patterns in a brain artery and the umbilical artery. A low ratio may suggest reduced placental function and fetal adaptation, but it must be interpreted with gestational age and other findings.

\medterm{Cerebrosides} Glycosphingolipids composed of ceramide plus a single sugar, abundant in myelin and nervous tissue and involved in membrane structure.

\medterm{Cerebrospinal Fluid} A clear, colorless fluid that fills the ventricles of the brain, the central canal of
the spinal cord, and the subarachnoid space around the brain and spinal cord. Produced
primarily by the choroid plexus of the lateral, third, and fourth ventricles (~500
mL\/day, total volume ~150 mL in adults).

\synonyms
Spinal Fluid

\medterm{Cerebrospinal Fluid Collection} A procedure, usually lumbar puncture, that collects cerebrospinal fluid for laboratory analysis; it can help evaluate infection, bleeding, inflammation, malignancy, and other disorders of the central nervous system.

\medterm{Cerebrospinal Fluid Culture} A laboratory test that attempts to grow bacteria, fungi, or other microorganisms from cerebrospinal fluid to investigate meningitis or another central nervous system infection. Results depend on prior antimicrobial treatment, specimen volume, transport, and the clinical context.

\medterm{Cerebrospinal Fluid Analysis} Laboratory examination of cerebrospinal fluid, usually obtained by lumbar puncture. It may
assess infection, bleeding, inflammation, demyelinating disease, or malignant cells.
Analysis can include opening pressure, appearance, cell counts, protein, glucose,
microbiology, and other tests selected for the clinical question.

\medterm{Cerebrospinal Fluid Cytology} Microscopic examination of cells in cerebrospinal fluid to evaluate suspected malignancy, infection, inflammation, or hemorrhage. It is especially useful for detecting leptomeningeal tumor involvement when interpreted with clinical and imaging findings.

\medterm{Cerebrovascular} Relating to the brain's blood vessels and blood flow. Cerebrovascular problems can reduce oxygen delivery and cause strokes, brief warning episodes, bleeding, or lasting problems with movement, speech, vision, or thinking.

\synonyms
Brain-vascular; cerebral vascular.

\medterm{Cerebrovascular Accident} A cerebrovascular accident is a stroke caused by interrupted brain blood flow from ischemia or bleeding, producing sudden neurologic dysfunction.
\medterm{Cerebrovascular Accident Prevention} Measures to reduce the risk of stroke, including controlling blood pressure, diabetes, and lipids, avoiding tobacco, treating atrial fibrillation when indicated, and using appropriate antithrombotic therapy.


\medterm{Cerebrovascular Accident (Stroke)} Sudden brain dysfunction caused by interrupted blood flow or bleeding in the brain. See Stroke.
\synonyms
Stroke

\medterm{Cerebrovascular Disorder} A disorder of brain blood vessels or cerebral blood flow, including ischemia, hemorrhage, aneurysm, malformation, and venous thrombosis.

\medterm{Cerebrum} The largest part of the brain, comprising two cerebral hemispheres connected by the
corpus callosum. Each hemisphere has four lobes: frontal (executive function, motor,
Broca's area), parietal (somatosensory, spatial awareness), temporal (auditory, memory,
Wernicke's area), occipital (visual processing).

\medterm{Cerium} A rare-earth metal whose compounds have industrial and experimental medical uses; exposure to some forms can irritate tissues or cause toxicity.

\medterm{Certification} A formal process confirming that a person, service, laboratory, or facility meets specified training, quality, safety, or professional standards.

\medterm{Certified Nurse-Midwife} An advanced-practice nurse trained and credentialed to provide reproductive, pregnancy, birth, postpartum, and newborn care, including management of normal pregnancy and referral of complications.

\medterm{Certolizumab Pegol} A biologic medicine that blocks tumor necrosis factor and treats selected inflammatory diseases, including rheumatoid arthritis and Crohn disease. It can increase the risk of serious infections.

\medterm{Ceruloplasmin} The main protein that carries copper in the blood. It is made mostly by the liver; unusually low or high levels can occur in Wilson disease, inflammation, liver disease, or other conditions.
\synonyms
Caeruloplasmin

\medterm{Ceruloplasmin Blood Test} A blood test measuring ceruloplasmin, the main copper-carrying protein in plasma; it is used with other copper studies to evaluate Wilson disease and disorders of copper metabolism.

\medterm{Cerumen} A substance that helps keep dirt out of the ear and lubricates the skin in the ear. More
commonly known as earwax.

\medterm{Cerumen Impaction} Earwax packed tightly enough to block the ear canal or prevent examination of the eardrum. It may cause blocked hearing, pressure, ringing, itching, pain, or dizziness.

\medterm{Ceruminous Adenocarcinoma} A malignant tumor of the ceruminous glands in the external ear canal that may cause pain, discharge, hearing loss, or a persistent ear-canal mass.

\medterm{Cervical} Relating to the neck or to the cervix, the lower opening of the uterus. The meaning depends on context, such as cervical spine, cervical pain, or cervical cancer.

\medterm{Cervical Adenocarcinoma} A malignant gland-forming tumor arising in the uterine cervix; it may be HPV-associated or HPV-independent and commonly presents with abnormal bleeding or discharge.

\medterm{Cervical Arthritis} Degenerative or inflammatory disease of the neck joints that can cause neck pain, stiffness, reduced motion, and sometimes nerve or spinal-cord compression.

\medterm{Cervical Atlas} The first bone of the neck, called C1. It is a ring-shaped vertebra that supports the skull and moves around the second neck vertebra when the head turns.

\medterm{Cervical Biopsy} Removal of a small piece of tissue from the cervix for laboratory examination. It is commonly performed after an abnormal screening result to look for precancerous or cancerous changes.
\medterm{Cervical Cancer} Cancer arising in the cervix, most often after persistent infection with a high-risk human papillomavirus type. Early disease may cause no symptoms; later signs include unusual bleeding, discharge, or pelvic pain.

\medterm{Cervical Cancer - Screening And Prevention} Clinical measures used to prevent or detect cervical cancer early, including human papillomavirus vaccination and screening with HPV testing, cervical cytology, or both according to age and risk.

\medterm{Cervical Cap} A flexible contraceptive barrier placed over the cervix to keep sperm from entering the uterus. It is used with spermicide and must remain in place for the recommended time after sex.
\medterm{Cervical Carcinoma} A malignant tumor of the uterine cervix, most commonly associated with persistent
high-risk human papillomavirus infection.

\medterm{Cervical Cerclage} Surgical placement of a non-absorbable suture around the cervix to provide mechanical
support and prevent preterm delivery in patients with cervical insufficiency.

\medterm{Cervical Change} A change in the uterine cervix or cervical spine, depending on context; clinical significance is determined by the specific structural, inflammatory, or pregnancy-related finding.

\medterm{Cervical Dilatation} Opening of the cervix during labor, measured from closed to about 10 centimeters. The cervix also becomes thinner and shorter, allowing the baby's head to pass into the birth canal.

\medterm{Cervical Dilation} The opening of the cervical canal during labor, progressing from closed to about 10 centimeters. It occurs with thinning and shortening of the cervix and is needed for vaginal birth.

\medterm{Cervical Disc Herniation} A neck disc bulges or tears and presses on a nearby nerve root or the spinal cord. It may cause neck pain, arm pain, tingling, numbness, or weakness.

\medterm{Cervical Dislocation} Abnormal displacement of one cervical vertebra relative to another, usually from trauma, that can destabilize the neck and injure the spinal cord.

\medterm{Cervical Dysplasia} Abnormal precancerous cells on the surface of the cervix, usually caused by persistent high-risk human papillomavirus infection. Mild changes often clear; more severe changes may progress to cancer if untreated.

\synonyms
cervical intraepithelial.
Cervical Intraepithelial Neoplasia Cervical (Cervical Dysplasia)
Dysplasia Cervical (Cervical Dysplasia)
Squamous Intraepithelial Lesion Cervical (Cervical Dysplasia)

\medterm{Cervical Dystocia} Failure of the uterine cervix to dilate normally during labor, which can delay or prevent vaginal delivery.

\medterm{Cervical Dystonia} Involuntary tightening of neck muscles that pulls the head into an abnormal position. It may cause neck pain, twisting, jerky movements, or tremor.

\synonyms
Dystonia, Cervical
Torticollis, Spasmodic

\medterm{Cervical Effacement} Thinning and shortening of the cervix during labor, measured from 0\% to 100\%. Complete effacement leaves the cervix very thin and usually occurs with opening of the cervix.

\medterm{Cervical Endometriosis} Endometrial-like tissue in or on the cervix that may cause bleeding after sex, pelvic pain, or an abnormal cervical lesion.

\medterm{Cervical Exam In Labor} A digital vaginal examination used to assess cervical dilation, effacement, consistency,
and position, together with the station and position of the fetal head. Serial examinations help assess labor progress.

\medterm{Cervical Ganglion} A cluster of nerve-cell bodies in the sympathetic chain of the neck that relays autonomic signals to the head, neck, and upper chest.

\medterm{Cervical Incompetence} The inability of the cervix to maintain pregnancy in the second trimester due to
painless cervical dilatation. Results in recurrent second trimester pregnancy loss or
preterm delivery. Risk factors include prior cervical trauma, cone biopsy, and
congenital anomalies.

\medterm{Cervical Insufficiency} Painless shortening or opening of the cervix during pregnancy before term, increasing risk of second-trimester loss or spontaneous preterm birth.

\medterm{Cervical Intraepithelial Neoplasia} Abnormal precancerous changes in the squamous cells of the cervix, usually associated with persistent high-risk human papillomavirus infection. CIN is graded by the severity of dysplasia.

\synonyms
CIN

\medterm{Cervical Intraepithelial Neoplasia Cervical (Cervical Dysplasia)} Alternate index label for Cervical Dysplasia.

\medterm{Cervical Length Screening} Transvaginal ultrasound measurement of the cervix between 18-24 weeks to assess risk of
spontaneous preterm birth. Normal cervix measures greater than 25 mm. Short cervix (less
than 25 mm) increases preterm delivery risk. Cervical funneling and dilatation are
additional concerning findings.

\medterm{Cervical Lymphadenopathy} Enlarged lymph nodes in the neck, commonly caused by infection but sometimes reflecting malignancy, autoimmune disease, or other systemic illness.

\medterm{Cervical MRI Scan} Magnetic resonance imaging of the cervical spine that uses a magnetic field and radio waves to produce detailed images of the neck vertebrae, discs, spinal cord, nerves, and surrounding soft tissues.

\medterm{Cervical Mucus} Secretions from the cervical glands that change throughout the menstrual cycle under
hormonal influence. Estrogen produces thin, clear, stretchy mucus favorable for sperm
transport around ovulation (spinnbarkeit). Progesterone produces thick, sticky, opaque
mucus that forms a plug in the cervical canal during pregnancy.

\medterm{Cervical Myelopathy} Dysfunction of the spinal cord caused by compression in the cervical spine, often producing hand clumsiness, gait disturbance, weakness, sensory changes, or abnormal reflexes.

\medterm{Cervical Osteoarthritis} Alternate index label; see \textit{Cervical spondylosis}.

\medterm{Cervical Pain} Alternate index label; see \textit{Neck pain}.

\medterm{Cervical Pain (Neck Pain)} Alternate index label for Neck Pain.

\medterm{Cervical Polyp} A usually benign polypoid overgrowth of endocervical or ectocervical tissue that may cause intermenstrual or postcoital bleeding or may be found incidentally during examination.

\medterm{Cervical Polyps} Usually noncancerous growths arising from the cervix or its canal. They may cause bleeding between periods, bleeding after sex, or discharge and are often removed for laboratory examination.

\medterm{Cervical Pregnancy} A rare ectopic pregnancy implanted in the cervical canal rather than the main uterine cavity. It may cause painless early-pregnancy bleeding and can lead to severe hemorrhage, requiring urgent specialist care.

\medterm{Cervical Radiculopathy} Irritation or compression of a nerve leaving the neck, often from a disc problem or arthritis. It can cause neck pain spreading into the arm, with tingling, numbness, or weakness.

\medterm{Cervical Rib} An extra rib growing from a neck vertebra, usually the lowest one. It often causes no symptoms but can narrow the space for nerves or blood vessels and produce arm pain, numbness, or swelling.

\medterm{Cervical Ripening Agents} Medicines or mechanical devices used to soften and prepare the cervix before labor
induction. Options include prostaglandins and cervical balloons; selection and dosing depend on the cervix, pregnancy,
contraindications, and local obstetric protocol.

\medterm{Cervical Screening (Pap Smear)} Screening for cervical cancer by collecting cells from the transformation zone of the
cervix. Liquid-based cytology (ThinPrep) or conventional Pap smear. HPV co-testing
(high-risk HPV 16, 18, others) improves sensitivity.

\synonyms
Pap Smear

\medterm{Cervical Spine} The seven bones in the neck, numbered C1 through C7. They support the head, protect the upper spinal cord, allow neck movement, and contain openings for important blood vessels.

\medterm{Cervical Spine CT Scan} Computed tomography of the cervical spine that uses x-rays and computer reconstruction to show the bones and alignment of the neck; it is particularly useful for detecting fractures and other bony abnormalities.

\medterm{Cervical Spine Trauma} An injury to the neck bones, discs, ligaments, spinal cord, or exiting nerves. It may cause pain, weakness, numbness, or paralysis; suspected serious injury requires immobilization and urgent assessment.

\medterm{Cervical Spondyloarthropathy} Wear or inflammation affecting the joints and discs of the neck. It may cause neck stiffness and pain, or press on nerves or the spinal cord, producing arm symptoms, weakness, or walking problems.

\medterm{Cervical Spondylosis} Age-related wear of the neck, including changes in discs, bones, joints, and ligaments. It may cause neck pain and stiffness and, if nerves or the spinal cord are compressed, limb symptoms.

\synonyms
Cervical Osteoarthritis
Osteoarthritis, Cervical
Spondylosis, Cervical

\medterm{Cervical Squamous Cell Carcinoma} A malignant squamous tumor of the cervix, most often caused by persistent high-risk human papillomavirus infection; it may present with abnormal bleeding, discharge, or an abnormal screening result.

\medterm{Cervical Stenosis} Narrowing or closure of the canal through the cervix, sometimes after surgery, radiation, infection, or menopause. It may cause painful or absent periods, trapped blood in the uterus, infertility, or no symptoms.

\medterm{Cervical Vertebrae} The seven vertebrae of the neck. Typical cervical vertebrae have small bodies, bifid
spinous processes, and transverse foramina transmitting the vertebral vessels; C1, C2,
and C7 are specialized.

\medterm{Cervical Vertebrae (C1-C7)} The seven vertebrae of the neck, numbered C1 through C7; they support the head, protect the cervical spinal cord, and allow neck movement.
\medterm{Cervicitis} Inflammation of the cervix, commonly caused by a sexually transmitted infection, changes in normal vaginal bacteria, chemical irritation, or allergy. It may cause discharge, bleeding after sex, or pelvic discomfort, but can be symptomless.

\medterm{Cervicogenic Headache} A headache caused by a disorder of the cervical spine or its soft tissues, with pain referred from the neck to the head.

\synonyms
Cervical headache; neck-related headache.

\medterm{Cervix} The cervix is the lower part of the uterus projecting into the vagina and producing mucus, allowing menstrual flow out and sperm passage in.

\medterm{Cervix Carcinoma} Cancer arising from the cervix, most often related to persistent high-risk human papillomavirus infection; screening and HPV vaccination reduce risk, and treatment depends on stage.

\medterm{Cervix Cryosurgery} Freezing abnormal surface cells of the cervix with a very cold probe. It may treat selected precancerous changes after appropriate testing; watery discharge and cramping are common temporarily.

\medterm{Cervix Disease} Any disorder affecting the lower part of the uterus, including infection, polyps, precancer, cancer, scarring, or pregnancy-related changes. Symptoms and treatment depend on the specific condition.

\medterm{Cervix Neoplasm} A benign, precancerous, or malignant growth of the cervix; evaluation may include pelvic examination, HPV testing, cytology, colposcopy, and biopsy.

\medterm{Cervix Uteri} The cervix uteri is the lower portion of the uterus that projects into the vagina and forms the passage between the uterus and birth canal.
\medterm{Cesarean Delivery} Surgical delivery of the fetus through abdominal and uterine incisions. Performed when
vaginal delivery is contraindicated or fails. Indications include malpresentation,
placenta previa, previous cesarean, fetal distress, and failed induction. Most common
uterine incision is low transverse (classical for very preterm or transverse lie).

\medterm{Cesarean Section} Delivery of a baby through surgical openings in the abdomen and uterus. It may be planned or urgent when vaginal birth is unsafe or unsuccessful, such as with fetal distress, abnormal position, placenta covering the cervix, or labor problems. Risks include bleeding, infection, blood clots, and complications in later pregnancies.

\medterm{Cesarean Section (C-Section)} Delivery of a baby through surgical cuts in the abdomen and uterus. It may be planned or emergency and can cause bleeding, infection, blood clots, or complications in future pregnancies.
\synonyms
C-section

\medterm{Cesium} A chemical element whose radioactive isotope cesium-137 can cause serious external or internal radiation injury after contamination or exposure.

\medterm{Cestoda} The tapeworm class of flatworms; human infection follows ingestion of eggs or larval tissue and can cause intestinal disease or tissue cysts.

\medterm{Cestode} A cestode is a tapeworm, a segmented parasitic flatworm that can infect humans or animals through contaminated food, water, or intermediate hosts.
\medterm{Cestodiasis} Infection or infestation caused by cestodes, commonly called tapeworms. Disease may
involve the intestine or tissues, depending on the species and life-cycle stage.

\synonyms
tapeworms.

\medterm{Cetirizine} A second-generation H1 antihistamine used for allergic rhinitis and urticaria; it is less sedating than older antihistamines but can still cause drowsiness.

\medterm{Cetoleic Acid} A long-chain monounsaturated fatty acid found mainly in marine oils and studied for effects on lipid metabolism and cellular membranes.

\medterm{Cetrimide} Cetrimide is a quaternary ammonium antiseptic and detergent used in some topical or laboratory preparations; concentrated exposure irritates tissue.
\medterm{Cetrorelix} A hormone-blocking medicine used during assisted reproduction to prevent premature ovulation. It temporarily blocks gonadotropin-releasing hormone receptors; injection-site reactions, headache, nausea, or allergy may occur.
\medterm{Cetuximab} A monoclonal antibody that blocks the epidermal growth factor receptor. It is used for selected colorectal, head-and-neck, and other cancers and commonly causes acne-like rash.

\medterm{Cevimeline} A muscarinic receptor agonist used to stimulate salivary secretion and relieve dry mouth, especially in selected patients with Sjögren syndrome.

\medterm{Cevimeline Hydrochloride} A medicine that stimulates muscarinic receptors to increase saliva. It is used mainly for dry mouth in Sjögren syndrome and can cause sweating, nausea, or blurred vision.

\medterm{CF} Alternate index label; see \textit{Cystic fibrosis}.

\medterm{CFS} Chronic fatigue syndrome, a condition involving substantial persistent fatigue with post-exertional worsening and associated symptoms that are not explained by another disorder.

\synonyms
Chronic fatigue syndrome; myalgic encephalomyelitis/chronic fatigue syndrome.

\medterm{CFTR Modulator Therapy} Medicines that improve the amount or function of defective CFTR protein in people with cystic fibrosis. The appropriate combination depends on the person's CFTR variants; treatment can improve breathing, nutrition, and salt-water balance.

\medterm{CHA2DS2-VASc Score} A score estimating stroke risk in a person with atrial fibrillation. It considers heart failure, high blood pressure, age, diabetes, previous stroke, blood-vessel disease, and sex; the result helps clinicians decide whether a blood thinner may offer more benefit than harm.

\synonyms
CHA2DS2-VASc; Atrial Fibrillation Stroke Risk Score

\medterm{Chaetomium} A group of molds found in soil and decaying plant material. They rarely infect people but may cause skin, eye, sinus, or widespread infection, mainly when immune defenses are severely weakened.

\medterm{Chafing} Skin irritation caused by repeated friction, moisture, heat, or tight clothing, producing redness, burning, tenderness, and sometimes erosions in opposing or rubbing areas.

\medterm{Chagas Cardiomyopathy} Heart damage caused by long-term Chagas disease. It can cause an irregular heartbeat, an enlarged heart, heart failure, or blood clots that may lead to a stroke.

\medterm{Chagas Disease} An infection caused by the parasite \textit{Trypanosoma cruzi}, usually spread by infected kissing bugs. It can also pass from mother to baby or through contaminated blood, organs, or food. Early illness may be mild or absent; years later, some people develop heart disease or enlargement of the oesophagus or colon.

\synonyms
Trypanosomiasis, American
American Trypanosomiasis

\medterm{Chagoma} A red, swollen lump where the parasite causing Chagas disease enters the body. It may be painful and accompanied by swollen nearby lymph nodes during the early stage of infection.

\medterm{Chain Of Custody} Meticulous paper trail and chronological documentation tracking the seizure, custody,
control, transfer, analysis, and disposition of physical and electronic evidence in
medico-legal investigations (e.g., blood samples for toxicology, bullets, sexual assault
kits, biological fluids).

\medterm{Chain-Termination Codon} A codon that signals the end of protein translation; the three standard stop codons are UAA, UAG, and UGA in messenger RNA.

\medterm{Chalazion} A usually painless, firm eyelid lump caused by blockage and inflammation of a meibomian oil gland; persistent or atypical lesions require ophthalmic assessment.

\medterm{Challenge Agent} A substance, organism, allergen, or physiologic stimulus deliberately administered in a controlled test to provoke and measure a response.

\medterm{Challenge Testing} A controlled exposure to a suspected trigger, such as a food or medication, to assess whether it produces a reproducible reaction.

\synonyms
Provocation testing; challenge test.

\medterm{Chamomile Extract} An extract of chamomile flowers used in traditional and complementary products; allergic reactions and medicine interactions are possible, while evidence varies by indication.

\medterm{Chancre} A usually painless ulcer at the site of inoculation of an infection, especially the
primary lesion of syphilis. It may be accompanied by regional lymph-node enlargement and
heals without treatment, although infection persists.

\medterm{Chancroid} A sexually transmitted infection caused by Haemophilus ducreyi, characterized by painful ragged genital ulcers and tender suppurative inguinal lymph nodes.

\medterm{Changes In The Newborn At Birth} The newborn transition from placental support to independent breathing and circulation, including lung aeration, closure of fetal shunts, temperature control, and adaptation of glucose and feeding.



\medterm{Channelopathy} A disorder caused by abnormal proteins that control movement of charged particles through cell membranes. It can affect nerves, muscles, or the heart and may cause seizures, weakness, or abnormal rhythms.

\medterm{Chaplain} A trained spiritual-care professional who supports patients and families with religious, existential, emotional, or ethical concerns in healthcare settings.

\medterm{Chapped Hands} Dry, cracked, inflamed skin of the hands caused by irritants, cold, low humidity, frequent washing, eczema, or allergy; fissures can become painful or infected.

\medterm{Chapped Lips} Dry, scaling, fissured lips caused by weather, licking, dehydration, irritants, medicines, or cheilitis; persistent lesions need assessment for infection or another disorder.

\medterm{Charcoal} Fine black powder produced by controlled pyrolysis of organic materials, processed to
yield a massive micro-porous surface area (500-1500 m2/g). See activated charcoal.

\medterm{Charcot Arthropathy} Progressive joint damage caused by loss of protective sensation, so repeated injuries go unnoticed. The joint may become warm, swollen, unstable, and deformed, especially in people with diabetic nerve damage.

\medterm{Charcot Foot} Progressive weakening and deformity of the foot or ankle caused by repeated unnoticed injury in a person with reduced sensation, commonly from diabetes. Early warmth and swelling require urgent off-loading to prevent collapse and ulcers.

\medterm{Charcot Joint} Severe destruction of a joint caused by loss of pain and position sensation. Repeated unnoticed injury can fragment bone, loosen or dislocate the joint, and produce major deformity.

\medterm{Charcot Triad} The combination of fever, pain in the upper-right abdomen, and jaundice, traditionally associated with acute infection of blocked bile ducts. This pattern is important but may be absent even in serious disease.

\medterm{Charcot Triad (Cholangitis)} Classic diagnostic triad of acute ascending cholangitis: (1) Fever, (2) Right Upper
Quadrant abdominal pain, and (3) Jaundice. Reynolds Pentad adds Hypotension and Altered
Mental Status. See Acute Cholangitis.

\synonyms
Cholangitis

\medterm{Charcot'S Joints} Charcot’s joints are severely damaged, unstable joints caused by loss of protective sensation, classically in diabetic neuropathy or other neurologic disease.
\medterm{Charcot'S Triad} Charcot’s triad is fever, jaundice, and right-upper-quadrant pain, a classic but incomplete clinical pattern of acute ascending cholangitis.
\medterm{Charcot-Marie-Tooth Disease} A group of inherited disorders that damage nerves supplying the feet, legs, hands, and arms. Symptoms usually include gradually worsening weakness, muscle wasting, numbness, foot deformity, and difficulty walking.

\synonyms
Neuropathy, Hereditary Motor And Sensory
Peroneal Muscular Atrophy

\medterm{Charcot-Wilbrand Syndrome} A rare neuropsychiatric syndrome after brain injury characterized by loss of visual imagery and visual dreaming, sometimes with difficulty recognizing or interpreting visual information.

\medterm{Chargaff Rule} The observation that double-stranded DNA contains approximately equal amounts of adenine and thymine and of guanine and cytosine, reflecting complementary base pairing.

\medterm{CHARGE Syndrome} A genetic condition, usually involving a new change in the CHD7 gene, that can affect the eyes, ears, heart, breathing passages, growth, development, and reproductive system. Features vary widely.

\medterm{Charley Horse} A sudden painful muscle cramp, commonly in the calf or thigh. It may follow exercise, dehydration, prolonged positioning, nerve problems, or certain medicines and usually lasts from seconds to minutes.

\medterm{Chasing The Dragon} Inhalation of vapor produced by heating heroin, a method of drug use that can cause dependence, overdose, toxic exposures, and neurological injury.

\medterm{Chaulmoogra Oil} Chaulmoogra oil is a historical plant oil once used to treat leprosy and some skin diseases; it has largely been replaced by effective antimicrobial therapy.
\medterm{CHD} Alternate index label; see \textit{Hip dysplasia}.

\medterm{Checkpoints} Regulatory signals that help prevent immune cells from attacking healthy tissues. Some cancers use checkpoint pathways to evade immunity; checkpoint-inhibitor medicines block selected signals to restore antitumor immune activity.

\medterm{Chediak-Higashi Syndrome} Autosomal recessive (LYST mutation, lysosomal trafficking regulator). Giant lysosomal
granules in neutrophils, melanocytes, neurons. Partial oculocutaneous albinism, silvery
hair, photophobia, nystagmus.

\medterm{Cheek} The facial region forming the side wall of the mouth, containing skin, muscle, fat, salivary structures, nerves, and blood vessels.

\medterm{Cheek Pouch} A distensible pocket inside the cheek found in some rodents and other animals, used to temporarily store and transport food.

\medterm{Cheilectropion} An eversion deformity of the lip, usually involving the lower lip. It may result from
scarring, congenital malformation, or loss of tissue support.

\medterm{Cheilitis} Inflammation of one or both lips, causing dryness, scaling, fissuring, erythema,
swelling, pain, or ulceration. Causes include irritant or allergic contact exposure,
infection, nutritional deficiency, chronic sun exposure, and inflammatory skin disease;
persistent lesions require evaluation for premalignant or malignant change.

\medterm{Cheilodynia} Pain of the lips, which may result from inflammation, infection, trauma, allergy, dryness, nerve pain, or another oral or dermatologic disorder.

\medterm{Cheilognathopalatoschisis} A congenital cleft involving the lip, jaw, and palate, caused by incomplete fusion of facial processes during embryonic development.

\medterm{Cheilognathoschisis} A congenital cleft of the lip and alveolar part of the upper jaw, with or without associated palatal involvement.

\medterm{Cheiloplasty} Plastic or reconstructive surgery of the lip, performed to repair congenital clefting,
trauma, tumors, scars, or other deformity.

\medterm{Cheiloschisis} Cleft lip, a congenital fissure of the upper lip caused by incomplete fusion of facial
processes during embryonic development. It may occur alone or with cleft palate.

\medterm{Cheilosis} A disorder of the lips characterized by dryness, scaling, fissuring, or inflammation.
Causes include irritation, nutritional deficiency, infection, and systemic disease.

\medterm{Cheiromegaly} Abnormal enlargement of one or both hands, classically associated with excess growth hormone but also possible with local overgrowth or other disorders.

\medterm{Cheiropompholyx} A type of hand eczema that causes intensely itchy, small fluid-filled blisters on the palms or sides of the fingers. The skin may later peel, crack, or become painful.

\medterm{Chelating Agent} A compound that binds a metal ion through multiple chemical sites, allowing the metal to be measured, mobilized, or removed in selected toxic exposures.

\medterm{Chelating Agents} Chelating agents bind specific metal ions to form complexes that can facilitate their removal or reduce toxicity; examples include agents for lead, iron, or copper poisoning.
\medterm{Chelation} The formation of a stable complex between a metal ion and a molecule with multiple
binding sites. In medicine, chelation can be used under supervision to bind and promote
elimination of selected toxic metals.

\medterm{Chelation Study} A test that uses a metal-binding substance and imaging or blood measurements to assess how the body handles a mineral or to locate abnormal deposits. The exact method depends on the clinical question.

\medterm{Chelation Therapy} Treatment using a medicine that binds a specific harmful or excess metal so the body can remove it in urine or stool. It is used only for selected poisonings or disorders because it can harm the kidneys or alter essential minerals.

\medterm{Chemexfoliation} Application of a chemical solution to remove controlled layers of skin, improving selected scars, pigmentation, wrinkles, or superficial lesions.

\medterm{Chemical} A substance defined by its molecular composition and properties; medical effects depend on dose, route, exposure duration, metabolism, and biological target.

\medterm{Chemical Ablation} Destruction of unwanted tissue by injecting or applying a chemical irritant or sclerosant, used in selected lesions, vessels, cysts, or other conditions.

\medterm{Chemical Burn} Damage to skin, eyes, or internal tissues caused by a corrosive chemical. Injury may worsen while the chemical remains present; contaminated clothing should be removed and the area flushed promptly with plenty of water.

\medterm{Chemical Burn Or Reaction} Tissue injury caused by contact with a corrosive or irritating chemical, commonly affecting the skin or eyes; severity depends on the substance, concentration, contact time, and tissue involved.

\medterm{Chemical Carcinogenesis} Development of cancer after exposure to chemicals that cause DNA damage, abnormal cell signaling, chronic tissue injury, or altered gene regulation.

\medterm{Chemical Cleavage} Breaking a chemical bond with a reagent, used in laboratory analysis, molecular biology, synthesis, or controlled degradation of a larger molecule.

\medterm{Chemical Cofactor} A nonprotein chemical helper required for an enzyme or biologic reaction, such as a metal ion, vitamin-derived coenzyme, or other small molecule.


\medterm{Chemical Dependency} Legacy, nonpreferred label; see \textit{Substance Use Disorder}.

\medterm{Chemical Eye Injury} Ocular damage caused by exposure to an acidic or alkaline substance. Immediate copious
irrigation, removal of particulate material, measurement of ocular-surface pH, and
urgent ophthalmic assessment are required; alkali injuries may penetrate deeply and
continue to worsen after exposure.

\medterm{Chemical Fractionation} Separation of a mixture into fractions according to chemical or physical properties such as solubility, charge, size, or binding affinity.

\medterm{Chemical Impurity} An unintended chemical present in a medicine, specimen, or material, potentially affecting safety, stability, measurement accuracy, or biologic activity.

\medterm{Chemical Injury} Tissue damage caused by contact with a corrosive or irritating chemical. Acids and alkalis can destroy tissue differently, and immediate thorough rinsing, removal of contaminated clothing, and medical assessment may be needed.

\medterm{Chemical Peel} Application of a chemical solution to the skin to produce controlled exfoliation and regeneration of superficial or deeper skin layers.

\synonyms
Chemexfoliation; chemical resurfacing.

\medterm{Chemical Peels} Treatments that apply a chemical solution to remove a controlled amount of skin. Superficial peels affect the outer layer; deeper peels reach the underlying dermis and have greater risks of scarring, pigment changes, infection, and prolonged healing.

\medterm{Chemical Pneumonitis} Lung inflammation and injury after inhaling or breathing in irritating chemicals, stomach contents, fuels, or other toxic substances. It may cause cough, shortness of breath, wheezing, low oxygen, and chest imaging changes.

\medterm{Chemistry} The study of substances, their composition, properties, reactions, and transformations, forming the basis of biochemistry, pharmacology, toxicology, and laboratory medicine.

\medterm{Chemo} Informal shorthand for chemotherapy, treatment using drugs that destroy cancer cells, inhibit their growth, or suppress abnormal immune cells.

\medterm{Chemo Brain} Problems with memory, concentration, word-finding, or mental processing that may occur during or after cancer treatment. Symptoms can affect work and daily activities and may have several contributing causes.

\synonyms
Post-Chemotherapy Cognitive Impairment

\medterm{Chemoattractant} A chemical signal that guides cells, especially immune cells, toward its source. Such signals help direct cells to areas of infection, inflammation, tissue repair, or normal development.
\medterm{Chemodectoma} A tumor arising from paraganglionic tissue, especially a carotid-body tumor. It is
usually slow-growing and may compress nearby cranial nerves or vessels.

\medterm{Chemoembolization} A procedure that combines delivery of an antineoplastic agent with embolization of the
blood supply to a tumor, most commonly for selected liver cancers.

\medterm{Chemokine} A small signaling protein that directs immune cells to sites of infection, inflammation, or tissue repair. Different chemokines attract different cell types by binding to matching receptors.

\medterm{Chemokine Receptor} A cell-surface receptor that binds chemokines and directs immune-cell movement; some also serve as entry cofactors for viruses such as HIV.

\medterm{Chemokinesis} Increased or altered movement of a cell in response to a chemical, without movement specifically toward or away from the chemical source. It differs from chemotaxis, which is directional.

\medterm{Chemonucleolysis} An intradiscal treatment intended to dissolve part of a herniated intervertebral disk, historically using an enzyme preparation.

\synonyms
Enzymatic disk dissolution; chemonucleolysis procedure.

\medterm{Chemoprevention} The use of drugs or natural agents to prevent, delay, or reverse cancer development,
e.g., tamoxifen or raloxifene in high-risk breast cancer, HPV and hepatitis B
vaccination, and aspirin in colorectal cancer.

\medterm{Chemoradiation} The concurrent combination of chemotherapy and radiation therapy, used to treat
locoregionally advanced cancers such as anal, cervical, head and neck, esophageal, and
limited-stage small cell lung cancer.

\medterm{Chemoradiotherapy} Cancer treatment that combines chemotherapy with radiation treatment. The medicines may make cancer cells more sensitive to radiation or treat cells elsewhere, but combined treatment can increase tiredness, nausea, skin changes, and other side effects.

\medterm{Chemoreceptor} A sensory receptor that responds to changes in blood or CSF chemistry, involved in
respiratory and cardiovascular regulation.

\medterm{Chemoreceptor Cell} A sensory cell that detects chemical changes such as oxygen, carbon dioxide, pH, taste molecules, or odorants and converts them into neural or physiologic signals.

\medterm{Chemoreceptor Trigger Zone} A brain area that detects vomiting-provoking chemicals in the blood or cerebrospinal fluid and activates brain circuits that produce nausea and vomiting.

\medterm{Chemoreceptors} Chemoreceptors detect chemical changes such as oxygen, carbon dioxide, pH, or odor molecules and trigger neural or physiologic responses.
\medterm{Chemoresistance} Reduced response of cancer cells or other target cells to a medicine that previously inhibited or destroyed them. It may be present from the start or develop during treatment.

\medterm{Chemosensitivity} How strongly cells or a tumor respond to a chemical treatment. A sensitive tumor may shrink or stop growing, but laboratory sensitivity tests do not always predict a person's treatment response.

\medterm{Chemosensitivity Assay} A laboratory test exposing cells to medicines to assess response; results are mainly investigational and may not predict treatment benefit.

\medterm{Chemosensitization} Making cancer cells or microorganisms more vulnerable to treatment by adding an agent that improves delivery, blocks resistance, or increases treatment-related damage.

\medterm{Chemosis} Swelling of the clear membrane covering the white part of the eye. It may result from allergy, infection, irritation, injury, surgery, or fluid retention and can make the eye look blistered.

\medterm{Chemosurgery} Chemosurgery uses a chemical agent to destroy selected abnormal tissue, historically including chemical treatment of skin lesions; modern use is specialized.
\medterm{Chemotactic Factors} Chemical signals that attract or repel cells by forming a concentration gradient. During infection or tissue injury, they help direct immune cells to the affected area.

\medterm{Chemotaxis} Directed movement of a cell or organism toward or away from a chemical stimulus. It is
important in immune-cell recruitment and inflammatory responses.

\medterm{Chemotherapy} Treatment using medicines to kill cancer cells or slow their growth, often by disrupting cell division or damaging DNA. It may be used alone or with surgery or radiation. Healthy fast-growing cells can also be affected, causing nausea, hair loss, anemia, or reduced immunity.


\medterm{Chemotherapy Extravasation} Leakage of an intravenous cancer medicine into tissue around the vein. Some medicines can cause severe pain, blistering, or tissue death; the infusion must stop immediately and the affected area needs urgent protocol-based care.

\medterm{Chemotherapy Regimen} A specific combination of anticancer drugs, dosing schedule, and treatment cycles designed for a particular cancer type and stage. Regimens are standardized by evidence-based protocols and may be modified based on response, toxicity, and patient factors.

\medterm{Chemotherapy Treatment} Chemotherapy uses medicines that damage or inhibit rapidly dividing cancer cells. It may be given before or after local treatment, for advanced disease, or for symptom control; effects and monitoring depend on the drugs, cancer, treatment goal, and organ function.

\medterm{Cherry Angioma} A common benign proliferation of small dermal blood vessels that appears as a bright-red
to purple papule, usually on the trunk or limbs. Lesions may bleed after trauma but
generally require no treatment unless diagnosis is uncertain or removal is desired.

\medterm{Cherry Eye} Prolapse of the gland of the third eyelid, producing a rounded red mass at the inner
corner of the eye. The term is used especially in veterinary ophthalmology.

\medterm{Cherry-Red Spot} A bright-red area at the center of the retina surrounded by pale retina. It occurs when the surrounding retina becomes swollen or loses blood flow and may indicate a serious metabolic, vascular, or storage disorder.

\medterm{Cherubism} A rare inherited disorder characterized by bilateral enlargement of the jaws from
expansile giant-cell lesions, usually beginning in childhood and often regressing after
puberty.

\medterm{Chest} The thorax, the region between the neck and abdomen containing the heart, lungs, major vessels, and other structures protected by the ribs and sternum.

\medterm{Chest CT} Computed tomography of the chest using x-rays and computer reconstruction to produce cross-sectional images of the lungs, heart, blood vessels, mediastinum, pleura, and chest wall.

\medterm{Chest MRI} Magnetic resonance imaging of the chest using a magnetic field and radio waves to evaluate soft tissues, blood vessels, heart structures, mediastinum, chest wall, and selected lung or pleural abnormalities.

\medterm{Chest Pain} Pain, pressure, burning, or discomfort in the chest. Causes include heart, lung, digestive, muscle, bone, or anxiety-related problems. New severe pain, breathlessness, sweating, faintness, or spreading pain requires emergency care.

\synonyms
Pain, Chest

\medterm{Chest Pain On The Left Side Above A Female Breast} Left-sided chest or breast-area pain can arise from muscle or rib injury, breast tissue, reflux, lung disease, anxiety, or heart disease. Pressure, exertional pain, shortness of breath, sweating, faintness, or pain spreading to the arm or jaw requires emergency assessment.

\medterm{Chest Physiotherapy} Techniques to mobilize and clear secretions from the lungs, including postural drainage,
percussion, vibration, and huffing. It is used in bronchiectasis, CF, and post-operative
patients. Positive expiratory pressure devices and oscillating PEP devices are also
used.


\medterm{Chest Tightness} A sensation of pressure, constriction, or heaviness in the chest that may arise from heart, lung, gastrointestinal, musculoskeletal, or anxiety-related causes.

\medterm{Chest Tube} A flexible tube inserted into the pleural space to remove air, blood, or fluid and help the lung re-expand.

\medterm{Chest Tube Insertion} Placement of a flexible tube through the chest wall into the space around a lung to remove air, blood, pus, or fluid and help the lung expand. It is connected to a drainage system and requires monitoring for bleeding, infection, or persistent air leak.

\medterm{Chest Tube Thoracostomy} Insertion of a tube into the pleural space to drain air, fluid, or blood. Indications:
pneumothorax (large, tension, or on mechanical ventilation), hemothorax, pleural
effusion (empyema, malignant, transudative), and postoperative (cardiac, thoracic
surgery).

\medterm{Chest Wall} The musculoskeletal boundary of the thoracic cavity, formed by the skin, fascia,
muscles, ribs, costal cartilages, sternum, thoracic vertebrae, and associated
neurovascular structures. It protects the thoracic organs and participates in breathing
through coordinated movement of the ribs and respiratory muscles.

\medterm{Chest Wall Deformity} An abnormal shape of the ribs, breastbone, spine, or surrounding tissues. A sunken or protruding breastbone is often mainly cosmetic, but severe deformity or scoliosis can restrict breathing or affect the heart.

\medterm{Chest Wall Pain} Alternate index label; see \textit{Costochondritis}.

\medterm{Chest X-Ray} An imaging test using a small amount of radiation to show the lungs, heart, ribs, and chest structures. It can detect pneumonia, collapsed lung, fluid, some fractures, and changes in heart size.

\medterm{Chewing} The coordinated use of teeth, jaw muscles, and joints to break food into smaller pieces for swallowing and digestion; pain or missing teeth can interfere.

\medterm{Cheyne-Stokes Respiration} A repeating pattern in which breathing gradually becomes deeper and faster, then shallower, followed by a pause. It may occur during sleep with heart failure, stroke, brain disease, or at high altitude.

\medterm{CHF} Heart failure in which the heart cannot pump or fill well enough for the body's needs. Reduced circulation can cause breathlessness, tiredness, leg swelling, and fluid buildup in the lungs or elsewhere.

\medterm{Chiari I Malformation} A condition in which the lower part of the cerebellum extends through the opening at the skull base. It may cause headaches worsened by straining, neck pain, balance problems, or no symptoms.

\medterm{Chiari Malformation} A structural abnormality in which part of the cerebellum extends downward through the opening at the base of the skull. Different types vary in severity and may cause headache, balance problems, swallowing difficulty, or no symptoms.

\medterm{Chiari-Frommel Syndrome} Continued milk production and absent menstrual periods after pregnancy, usually caused by persistently high prolactin levels. Evaluation looks for pituitary disorders, medicines, thyroid disease, and other causes.

\medterm{Chick Embryo} The developing chicken embryo, widely used as an experimental model for embryology, organ development, virology, toxicology, and tissue growth.

\medterm{Chicken Soup And Sickness} A supportive-care topic concerning the use of warm fluids and nutrition during illness; chicken soup may ease hydration or congestion but does not replace diagnosis or treatment.

\medterm{Chickenpox} A highly contagious infection caused by varicella-zoster virus. It causes an itchy rash that appears as spots, fluid-filled blisters, and scabs at different stages, with fever and tiredness; complications are more likely in adults and people with weakened immunity.

\medterm{Chickenpox (Varicella)} Alternate index label for Varicella.


\medterm{Chigger Bites} Itchy skin reactions caused by larval mites. Small red bumps or blisters often develop where clothing is tight and may itch for days. Washing the skin, avoiding scratching, and itch treatment usually provide relief.

\medterm{Chiggers} Larval harvest mites whose bite causes intensely itchy grouped red papules or welts; symptoms result from skin irritation, not from the mite burrowing into human tissue.

\medterm{Chiggers (Bites)} Alternate index label for Bites.

\medterm{Chignon} Temporary swelling of a baby's scalp where a vacuum device was applied during assisted delivery. It results from the suction and usually disappears within a day or two without treatment.

\medterm{Chikungunya} A mosquito-borne viral disease caused by chikungunya virus, usually producing abrupt fever and severe, often symmetric joint pain. Headache, muscle pain, joint swelling, nausea, fatigue, and rash may occur; joint symptoms can persist for months or years.

\synonyms
Chikungunya fever; chikungunya virus disease; CHIKV infection

\medterm{Chikungunya Virus} An alphavirus transmitted primarily by infected female Aedes aegypti and Aedes albopictus mosquitoes. Infection can cause abrupt fever, severe polyarthralgia, joint swelling, rash, and prolonged inflammatory joint symptoms.

\medterm{Chilblain} An inflammatory skin lesion caused by an abnormal response to cold and damp conditions.
It produces itchy, red or violaceous swollen areas, usually on fingers, toes, ears, or
nose.

\medterm{Chilblains} Itchy, painful, red or purple swellings caused by an abnormal skin response after exposure to cold and damp conditions. They usually improve with warmth but can blister, ulcerate, or become infected.
\synonyms
Pernio

\medterm{Chilblains (Pernio)} Itchy, painful, red or purplish inflammatory lesions caused by an abnormal response of small blood vessels to cold and subsequent rewarming.

\synonyms
Pernio

\medterm{Child Abuse} Physical, sexual, or emotional harm or neglect of a child by a caregiver. Recognized by
inconsistent histories, injury patterns inconsistent with mechanism, and behavioral
indicators. Types: physical abuse, sexual abuse, neglect, and emotional abuse. Mandatory
reporting requirements for healthcare providers.

\synonyms
Child Maltreatment


\medterm{Child Adoption} Child adoption is the legal process by which parental rights and caregiving responsibility are transferred to adoptive parents; medical care includes records, consent, and family history.
\medterm{Child Behavior} The observable actions, responses, and conduct of children, influenced by developmental stage, temperament, environment, and neurobiological factors. Behavioral assessment helps identify developmental disorders, emotional difficulties, and parenting needs.


\medterm{Child Development} Progressive physical, cognitive, language, emotional, and social changes from infancy through adolescence, assessed against developmental milestones and individual variation.

\medterm{Child Health} The physical, developmental, mental, and social well-being of children, including prevention, nutrition, immunization, safety, and treatment of illness.

\medterm{Child Health Services} Child health services provide preventive care, immunization, growth and developmental monitoring, screening, treatment, and support for children and families.

\medterm{Child Neglect} The failure of a caregiver to provide adequate food, shelter, clothing, supervision,
education, or medical care. The most common form of child maltreatment. Types: physical
neglect, educational neglect, emotional neglect, and medical neglect. May result in
failure to thrive, developmental delay, and chronic illness.

\medterm{Child Neglect And Emotional Abuse} Failure to provide a child’s basic needs or repeated behavior that harms emotional development, safety, or wellbeing; suspected abuse requires safeguarding assessment and appropriate reporting under local law.

\medterm{Child Physical Abuse} Intentional physical injury or harm inflicted on a child by a caregiver or another person, including hitting, burning, shaking, or fractures; suspected abuse requires immediate safeguarding assessment.

\medterm{Child Poisoning} Exposure of a child to a medicine, chemical, or other substance that may cause harm. Suspected poisoning requires prompt expert
assessment.

\medterm{Child Psychiatry} The medical specialty concerned with assessment and treatment of emotional, behavioral, developmental, and psychiatric disorders in children and adolescents.

\medterm{Child Psychology} The study of children’s cognitive, emotional, social, and behavioral development and the effects of relationships, learning, biology, and environment.



\medterm{Child-Pugh Classification} Clinical scoring system assessing prognosis of chronic liver disease by evaluating five
parameters: total bilirubin, serum albumin, international normalized ratio, ascites, and
hepatic encephalopathy. Each parameter is scored one to three points.

\medterm{Childbed Fever} Puerperal fever, usually infection of the reproductive tract after childbirth or miscarriage, often involving postpartum endometritis and requiring prompt evaluation and treatment.

\medterm{Childbirth} The process of giving birth, encompassing three stages: (1) First stage: from the onset
of regular uterine contractions (cervical effacement and dilation) to full dilation (10
cm). See Parturition (Childbirth).

\medterm{Childhood} The developmental period between infancy and adolescence, during which physical growth, brain development, learning, and social skills progress.

\medterm{Childhood Absence Epilepsy} An epilepsy syndrome beginning in childhood with frequent brief absence seizures,
characterized by sudden impairment of awareness and a typical generalized spike-and-wave
EEG pattern.

\medterm{Childhood Angiosarcoma} A rare malignant tumor arising from blood-vessel or lymphatic-endothelial cells in a child; diagnosis and treatment require specialist pediatric oncology care.

\medterm{Childhood Apraxia Of Speech} A speech-sound disorder in which a child has difficulty planning and coordinating the movements needed for clear speech.

\synonyms
CAS

\medterm{Childhood Asthma} Asthma beginning during childhood, involving variable airway narrowing and inflammation that cause wheezing, cough, chest tightness, or shortness of breath.

\medterm{Childhood Ependymoma} A tumor arising from ependymal cells lining the brain ventricles or spinal canal in children; symptoms depend on location and obstruction of cerebrospinal-fluid flow.

\medterm{Childhood Fibrosarcoma} A malignant tumor of fibroblastic connective tissue in children, usually presenting as a growing mass; treatment depends on location, extent, and tumor biology.

\medterm{Childhood Ganglioglioma} A usually slow-growing mixed neuronal and glial brain tumor in children or young adults, often presenting with long-standing focal seizures.

\medterm{Childhood Leiomyosarcoma} A rare malignant tumor of smooth muscle in a child; it may arise in soft tissue or an organ and requires biopsy, staging, and specialist treatment.

\medterm{Childhood Leukemia} Cancer of blood-forming cells in a child, most often acute lymphoblastic or acute myeloid leukemia; symptoms may include pallor, bruising, fever, bone pain, or recurrent infection.

\medterm{Childhood Liposarcoma} A rare malignant tumor derived from fat-forming tissue in a child; diagnosis requires imaging and biopsy, and treatment is individualized by a sarcoma team.

\medterm{Childhood Obesity} Alternate index label; see \textit{Pediatric Obesity}.

\medterm{Childhood Osteosarcoma} A malignant bone tumor that commonly develops near the ends of long bones during adolescence; pain, swelling, or a fracture may occur and treatment combines chemotherapy and surgery.

\medterm{Childhood Rhabdomyosarcoma} A malignant tumor showing skeletal-muscle differentiation, occurring mainly in children and adolescents; it may arise in the head and neck, urinary or reproductive tract, or limbs.

\medterm{Childhood Schizophrenia} A rare psychotic disorder beginning in childhood, involving disturbances in thinking, perception, behavior, and emotional functioning.

\synonyms
Juvenile Schizophrenia
Schizophrenia, Adolescent
Schizophrenia, Childhood
Schizophrenia, Juvenile

\medterm{Children} Humans in the developmental period between infancy and adolescence; medical assessment must account for age-specific anatomy, physiology, growth, communication, and consent.

\medterm{Children And Grief} The emotional and behavioral response of children to death or major loss, which varies by developmental stage and may include sadness, regression, anger, worry, or physical symptoms.







\medterm{Children'S Cancer Centers} Specialized services that coordinate diagnosis, treatment, rehabilitation, psychosocial support, survivorship care, and research for children and adolescents with cancer.



\synonyms
Childrens Health (Children's Health)


\medterm{Chill} A chill is a sensation of cold, often with shivering, that can accompany fever, infection, exposure, anxiety, or other illness.
\medterm{Chills} A sensation of cold with involuntary shivering, commonly accompanying fever or infection; chills without measured fever may also occur with anxiety, hypoglycaemia, or environmental cold.

\medterm{Chilomastix Mesnili} A usually harmless single-celled organism that may be found in stool. Its presence generally does not cause illness or require treatment, but it can indicate exposure to contaminated food or water.

\medterm{Chimera} An organism or tissue containing genetically distinct cell populations derived from different zygotes or embryos, or a laboratory construct combining genetic material from different sources.

\medterm{Chimerism} The presence in one individual of cell populations derived from genetically different
zygotes or sources. It may occur naturally, after transplantation, or through exchange
of cells between a pregnant person and fetus.

\medterm{Chin} The lower facial prominence formed mainly by the mandible and overlying soft tissue, contributing to facial contour, speech, chewing, and airway anatomy.

\medterm{Chin Augmentation} A cosmetic procedure that increases chin projection with an implant, injectable filler, or bone surgery to change facial contour; risks include infection, asymmetry, nerve injury, and implant problems.

\medterm{Chinese Restaurant Syndrome} An older, controversial term for short-lived symptoms such as headache, flushing, or palpitations attributed to eating restaurant food; evidence does not show that monosodium glutamate causes these symptoms in most people at usual dietary exposures.

\medterm{Chionablepsia} Temporary eye pain and blurred vision caused by ultraviolet light reflected from snow or ice, commonly called snow blindness. Symptoms may appear hours after exposure and usually improve as the cornea heals.

\synonyms
snow blindness.

\medterm{Chironomidae} The nonbiting midge family of insects; their larvae and body parts can act as environmental allergens and rarely trigger respiratory or skin reactions.

\medterm{Chiropody} Chiropody is care of the feet, including assessment and treatment of skin, nail, pressure, and biomechanical problems; it is also called podiatry.
\medterm{Chiropractic} A healthcare practice that mainly uses hands-on treatment for muscle, joint, and spine problems, especially spinal manipulation. Benefits and risks vary by the condition, technique, and practitioner's training.

\medterm{Chiropractic Care For Back Pain} Manual therapy, primarily spinal manipulation and mobilization, provided for acute and chronic low back pain when indicated. Evidence supports its use for acute low back pain in particular, though it is not appropriate for all causes of back pain.

\medterm{Chiropractor} A healthcare practitioner who provides manual or manipulative treatment, especially for musculoskeletal conditions, according to the practitioner's training and local regulation.

\synonyms
Chiropractic practitioner; chiropractic clinician.


\medterm{Chitin} A structural polysaccharide in fungal cell walls and arthropod exoskeletons.

\medterm{Chitosan} A material made by chemically modifying chitin from shellfish or fungi. It is used in some wound dressings and research products; claims that oral chitosan causes major weight loss are not well supported.

\medterm{Chlamydia} A group of bacteria that must live inside host cells to multiply. Different species cause sexually transmitted infections, eye disease, respiratory illness, or infections acquired from birds.

\medterm{Chlamydia Infections In Women} Infections in women caused by Chlamydia trachomatis, commonly involving the cervix or urethra and often producing no symptoms; untreated infection can cause pelvic inflammatory disease, infertility, or ectopic pregnancy.

\medterm{Chlamydia Pneumoniae Infection} A respiratory infection caused by Chlamydia pneumoniae, formerly called Chlamydophila pneumoniae. It commonly causes gradual-onset atypical pneumonia, bronchitis, or upper-respiratory symptoms.

\medterm{Chlamydia Psittaci} A bacterium that causes psittacosis, a zoonotic infection usually acquired by inhaling organisms from infected birds.

\medterm{Chlamydia Trachomatis} A bacterium that commonly causes sexually transmitted infection of the urogenital tract and can also cause conjunctivitis or pneumonia.

\medterm{Chlamydia Trachomatis Infection} A bacterial infection spread mainly through sexual contact. It often causes no symptoms but may cause discharge or painful urination and can lead to pelvic inflammatory disease, infertility, eye infection, or newborn infection.

\medterm{Chlamydiaceae} A family of bacteria that includes Chlamydia species. Some cause sexually transmitted infections, eye disease, or respiratory infection; diagnosis depends on the species, specimen, symptoms, and laboratory testing.

\medterm{Chlamydial Infections - Male} Infections in men caused by Chlamydia trachomatis, most commonly urethritis and sometimes epididymitis or proctitis; many infections are asymptomatic and diagnosis is usually by nucleic acid amplification testing.

\medterm{Chlamydial Pneumonia} A lung infection caused by Chlamydia bacteria, usually with gradually worsening cough, fever, tiredness, headache, or sore throat. Diagnosis and treatment depend on the species, severity, and other medical conditions.

\medterm{Chlamydiosis} Infection caused by bacteria of the genus Chlamydia. Depending on the species and site,
it may cause sexually transmitted infection, conjunctivitis, pneumonia, or other
disease.

\medterm{Chloasma} Patchy brown or gray-brown darkening of the skin, usually on the face. It is often related to pregnancy, hormones, and sunlight and is commonly called melasma.

\synonyms
melasma.

\medterm{Chloasma (Melasma)} Alternate index label for Melasma.

\medterm{Chloracne} An acne-like eruption caused by exposure to certain halogenated aromatic compounds, such
as dioxins and related industrial chemicals.

\medterm{Chloral Hydrate} A sedative-hypnotic compound with limited modern use because of toxicity, dependence, respiratory depression, and cardiac or gastrointestinal adverse effects.

\medterm{Chlorambucil} An oral chemotherapy medicine that damages DNA and is used mainly for some chronic leukemias and lymphomas. It can lower blood counts and increase infection risk.

\medterm{Chloramphenicol} A broad-spectrum antibiotic with activity against gram-positive and gram-negative
organisms, as well as Rickettsia. Chloramphenicol inhibits bacterial protein synthesis
by binding to the 50S ribosomal subunit. It is indicated for serious infections when
other agents are not appropriate, including bacterial meningitis and typhoid fever.

\medterm{Chloranemia} An older term for anemia in which red blood cells contain too little hemoglobin and appear unusually pale. Iron deficiency is a common cause; modern evaluation uses blood counts and iron studies.

\medterm{Chlordane} A persistent organochlorine insecticide formerly used for termite control; exposure can affect the nervous system and liver and has environmental persistence.

\medterm{Chlordecone} A persistent organochlorine pesticide that can cause neurologic, reproductive, and liver effects after significant exposure and remains in the environment for long periods.

\medterm{Chlordiazepoxide} A benzodiazepine medicine that slows brain activity and may be used short term for alcohol-withdrawal symptoms or severe anxiety. It can cause drowsiness, dependence, impaired coordination, and dangerous breathing suppression with alcohol or opioids.

\medterm{Chlordiazepoxide Overdose} Overdose of chlordiazepoxide, a benzodiazepine, that can cause drowsiness, impaired coordination, respiratory depression, and coma, especially with other sedatives.

\medterm{Chlorethoxyfos} An organophosphate insecticide that can block a nerve enzyme, causing sweating, vomiting, weakness, breathing problems, seizures, or other poisoning symptoms after exposure.

\medterm{Chlorfenvinphos} An organophosphate insecticide that can cause cholinergic toxicity by inhibiting acetylcholinesterase; use is restricted or prohibited in many places.

\medterm{Chlorhexidine} A broad-spectrum antiseptic used for skin, mucosal, and oral preparation; it can cause irritation, allergy, or rare severe anaphylaxis.

\medterm{Chloride} The principal negatively charged electrolyte in extracellular fluid, important for fluid balance, electrical neutrality, and gastric acid formation.

\medterm{Chloride - Urine Test} A test that measures chloride excretion in urine; interpreted with serum electrolytes and clinical findings to evaluate fluid balance, acid-base disorders, and causes of abnormal chloride loss or retention.

\medterm{Chloride Channel} A protein forming a passage for chloride ions across a cell membrane. Chloride channels help control electrical signals, fluid movement, and muscle or nerve function.

\medterm{Chloride In Diet} Chloride consumed in food, usually as part of salt. It helps maintain fluid balance, nerve signals, and stomach acid; excessive intake often accompanies excess sodium.

\medterm{Chloride Ion} The negatively charged form of chlorine and a major extracellular electrolyte that helps maintain fluid balance, osmotic pressure, acid--base status, and electrical neutrality.

\medterm{Chloride Test - Blood} A blood test measuring serum chloride, an electrolyte that helps regulate extracellular fluid volume and acid-base balance; results are interpreted with sodium, bicarbonate, kidney function, and the clinical context.

\medterm{Chlorinated Lime} Chlorinated lime, or bleaching powder, releases chlorine and is used as a disinfectant; concentrated exposure can burn skin, eyes, and airways.
\medterm{Chlorinated Lime Poisoning} Exposure to chlorinated lime, a hypochlorite-containing corrosive that can injure the skin, eyes, airways, or gastrointestinal tract.

\medterm{Chlorine} A reactive halogen element and disinfectant; inhalation of chlorine gas can injure the eyes and airways and cause chemical pneumonitis or pulmonary edema.

\medterm{Chlorine Poisoning} Inhalation or contact exposure to chlorine gas, which irritates the eyes and airways and can cause cough, bronchospasm, chemical pneumonitis, or pulmonary edema.

\medterm{Chlormadinone Acetate} A synthetic progestogen used in some hormonal medicines; effects and risks include changes in bleeding, mood, metabolism, and thrombotic risk depending on the formulation.

\medterm{Chlormerodrin} A historical radioactive mercury compound once used in diagnostic or research settings; it is not a routine modern medicine and mercury exposure is toxic.

\medterm{Chlormethine} Chlormethine, or mechlorethamine, is an alkylating anticancer drug used systemically or topically for selected cancers and cutaneous T-cell lymphoma.
\medterm{Chlorodeoxyadenosine} Another name for cladribine, a purine-analog medicine used for selected blood cancers and multiple sclerosis; prolonged lymphocyte suppression and infection risk are important considerations.

\medterm{Chloroform} A volatile solvent and environmental contaminant that can irritate the lungs and nervous system and damage the liver or kidneys after significant exposure.

\medterm{Chlorogenic Acid} A natural plant compound found in coffee, fruits, and vegetables. It has antioxidant effects in laboratory studies, but evidence is insufficient to use supplements to prevent or treat disease.

\medterm{Chloroleukemia} A rare form of myeloid leukemia associated with greenish tumor masses composed of
myeloid cells, also called granulocytic sarcoma or chloroma.

\synonyms
granulocytic sarcoma or chloroma.

\medterm{Chloroma} A greenish extramedullary mass of immature myeloid cells, usually associated with acute
myeloid leukemia. It is also called granulocytic sarcoma or myeloid sarcoma.

\synonyms
granulocytic sarcoma or myeloid sarcoma.

\medterm{Chlorophyll} The green substance in plants and algae that absorbs light for photosynthesis, the process that makes food from light, water, and carbon dioxide. It is not a human nutrient or a proven treatment for disease.

\medterm{Chlorophyll Poisoning} Ingestion of chlorophyll-containing products, which usually causes mild gastrointestinal effects; product ingredients and dose determine clinical risk.

\medterm{Chlorophyllin} A water-soluble, copper-containing derivative of chlorophyll used in some products and studied for effects such as odor control; clinical uses depend on evidence and context.

\synonyms
Chlorophyllin copper complex; sodium copper chlorophyllin.

\medterm{Chlorophyta} The group of green algae; some species produce toxins or contribute to harmful algal blooms, while others are used as food or in research.

\medterm{Chloropicrin} A pungent soil fumigant and irritant that can injure the eyes, skin, and respiratory tract after exposure.

\medterm{Chloroplast} A plant and algal cell organelle that performs photosynthesis and contains its own DNA; it is not found in human cells.

\medterm{Chloroplast DNA} The genetic material inside chloroplasts, inherited mainly through the maternal line in many plants and used in studies of plant evolution and disease.

\medterm{Chloroprene} An industrial chemical used to make chloroprene rubber; repeated or high exposure can irritate tissues and has carcinogenic concern.

\medterm{Chloropsia} A visual disturbance in which objects appear green or greenish. It may result from retinal or visual-pathway disease, medicines, toxins, migraine, or other conditions and needs assessment if persistent or sudden.

\medterm{Chloroquine} An antimalarial medicine used for susceptible Plasmodium species and selected autoimmune conditions; retinal toxicity, cardiotoxicity, and drug interactions require attention.

\medterm{Chlorothiazide} A thiazide diuretic used to treat high blood pressure and fluid retention. It increases urine production and may cause low sodium, low potassium, dehydration, or high blood sugar.

\medterm{Chlorotoxin} A peptide from scorpion venom that binds some tumor and ion-channel targets; it is mainly studied experimentally for brain-tumor imaging or treatment.

\medterm{Chlorotrianisene} A synthetic estrogen formerly used for menopausal symptoms; it is rarely used now because of estrogen-related risks and safer alternatives.

\medterm{Chlorpheniramine} A first-generation antihistamine used for allergy symptoms; drowsiness, dry mouth, blurred vision, urinary retention, and other anticholinergic effects may occur.

\medterm{Chlorphentermine} An older stimulant appetite suppressant related to amphetamines; it is not routinely used because of cardiovascular, psychiatric, and dependence risks.

\medterm{Chlorpromazine} A first-generation antipsychotic and phenothiazine used for schizophrenia, severe nausea, and selected other conditions; adverse effects include sedation, hypotension, movement disorders, and QT prolongation.

\medterm{Chlorpromazine Overdose} Overdose of chlorpromazine that may cause sedation, hypotension, anticholinergic effects, seizures, cardiac conduction abnormalities, or neuroleptic malignant syndrome.

\medterm{Chlorpropamide} A first-generation sulfonylurea that lowers blood glucose by stimulating pancreatic insulin release; prolonged hypoglycemia and hyponatremia are important risks.

\medterm{Chlorpyrifos} An organophosphate insecticide that inhibits acetylcholinesterase; significant exposure can cause cholinergic toxicity with secretions, bronchospasm, bradycardia, weakness, or seizures.

\medterm{Chlortalidone} Chlortalidone, also called chlorthalidone, is a long-acting thiazide-like diuretic used mainly for hypertension and fluid retention.
\medterm{Chlorthalidone} A long-acting thiazide-like diuretic used for hypertension and edema; it can cause hyponatremia, hypokalemia, hyperuricemia, and kidney-related electrolyte changes.

\medterm{Chlorzoxazone} A centrally acting muscle relaxant used for painful muscle spasm in some countries; drowsiness and rare serious liver injury are important risks.

\medterm{CHO Cell} A Chinese hamster ovary cell line widely used to produce recombinant proteins, antibodies, vaccines, and other biologic medicines.

\medterm{Choanal Atresia} A birth defect in which tissue blocks one or both passages from the nose to the throat. Both sides can cause severe breathing difficulty in a newborn, especially during feeding, and require urgent treatment.

\medterm{Chocolate Cyst} A chocolate cyst is an ovarian endometrioma filled with old blood from endometriosis, which can cause pelvic pain, infertility, or an adnexal mass.
\medterm{Choice Behavior} The process of selecting among alternatives, studied in psychology and neuroscience to understand decision-making, impulsivity, and reward processing. Abnormal choice behavior may be relevant to addiction, eating disorders, and psychiatric conditions.

\medterm{Choke} Airway blockage by food or another object that prevents normal airflow and can rapidly cause inability to speak, cyanosis, loss of consciousness, or death.

\medterm{Choking} Blockage of the airway by food or another object, making breathing difficult or impossible. A person who cannot speak, cough effectively, or breathe needs immediate first aid and emergency medical help.

\medterm{Choking - Adult Or Child Over 1 Year} First aid for severe airway obstruction in a person older than one year, using back blows and abdominal thrusts when appropriate and CPR if the person becomes unresponsive.

\medterm{Choking - Infant Under 1 Year} First aid for severe airway obstruction in an infant, using alternating back slaps and chest thrusts and avoiding blind finger sweeps; unresponsiveness requires infant CPR.

\medterm{Choking - Unconscious Adult Or Child Over 1 Year} Resuscitation for an unresponsive person with suspected airway obstruction, combining emergency activation, airway inspection, CPR, and removal of a visible object without blind sweeping.

\medterm{Cholagogue} A substance that encourages the gallbladder to release bile into the small intestine. Bile helps the body digest fats; medicines or foods described as cholagogues should be used carefully when gallstones or blocked bile ducts are possible.

\synonyms
Gallbladder stimulant; biliary evacuant.

\medterm{Cholangiitis} Inflammation of the tubes that carry bile from the liver and gallbladder to the intestine. Infection, a blockage, autoimmune disease, or another bile-duct disorder may cause it.

\medterm{Cholangiocarcinoma} Cholangiocarcinoma is cancer arising from bile-duct lining cells and may cause jaundice, itching, abdominal symptoms, or bile-flow obstruction.

\synonyms
Bile Duct Cancer (Cholangiocarcinoma)

\medterm{Cholangiocarcinoma (Bile Duct Cancer (Cholangiocarcinoma))} Alternate index label for Bile Duct Cancer (Cholangiocarcinoma).

\medterm{Cholangiocarcinoma (Bile Duct Cancer)} Cancer arising in the bile ducts, which may cause jaundice, itching, abdominal discomfort, weight loss, or impaired bile flow.

\medterm{Cholangiography} Imaging of the bile ducts after contrast material makes them visible. It can show gallstones, narrowing, leaks, or blockage and may be performed during surgery, through the skin, or during endoscopy.
\medterm{Cholangiohepatitis} Inflammation involving both the bile ducts and liver, often associated with ascending bacterial infection, biliary obstruction, parasites, or recurrent cholangitis. It may cause fever, right-upper-quadrant pain, jaundice, and cholestatic liver-test abnormalities.

\medterm{Cholangitis} Cholangitis is inflammation or infection of the bile ducts, often from obstruction by a stone or stricture, and can rapidly cause sepsis.

\medterm{Cholangitis, Primary Sclerosing} Alternate index label; see \textit{Primary sclerosing cholangitis}.

\medterm{Cholate} The ionized form of cholic acid, a primary bile acid that helps emulsify dietary fat and is conjugated and secreted by the liver.

\medterm{Cholecalciferol} Vitamin D3, produced in skin by ultraviolet light and obtained from food or supplements, then converted to active metabolites that regulate calcium and phosphate.

\medterm{Cholecyst} The gallbladder, a small sac beneath the liver that stores and concentrates bile. When food enters the intestine, it contracts and releases bile to help digest dietary fats.

\medterm{Cholecystectomy} Surgical removal of the gallbladder, usually by a laparoscopic approach when appropriate. It may be performed for symptomatic gallstones, acute cholecystitis, selected gallstone pancreatitis, or other gallbladder disease.

\medterm{Cholecystitis} Inflammation of the gallbladder, usually because a gallstone blocks the tube through which bile leaves it. It commonly causes steady upper-right abdominal pain, fever, nausea, or vomiting and may require urgent treatment.

\synonyms
Inflamed Gallbladder

\medterm{Cholecystography} Radiographic examination of the gallbladder, usually after administration of a contrast
substance, to assess its filling, concentration, and emptying.

\medterm{Cholecystokinin} A hormone released by the small intestine after food, especially fat and protein, enters it. It makes the gallbladder release bile, stimulates digestive enzymes, and helps create a feeling of fullness.

\medterm{Cholecystoses} Non-inflammatory disorders causing deposits, wall thickening, or other structural changes in the gallbladder, including cholesterolosis and adenomyomatosis. Many are incidental, but symptoms and imaging findings determine whether further evaluation is needed.

\medterm{Choledochal Cyst} An abnormal widening of one or more bile ducts, usually present from birth. It may cause abdominal pain, jaundice, pancreatitis, or infection and can increase later cancer risk, so specialist evaluation is needed.

\medterm{Choledochal Cysts} Choledochal cysts are congenital or acquired dilations of the bile ducts that can cause abdominal pain, jaundice, pancreatitis, cholangitis, or stones. Imaging establishes the anatomy, and surgical removal is often recommended because of recurrent complications and cancer risk.

\synonyms
Cysts Choledochal (Choledochal Cysts)

\medterm{Choledocholithiasis} A gallstone within the common bile duct, which may cause biliary pain, jaundice, cholangitis, or pancreatitis and is usually evaluated with liver tests and imaging.

\medterm{Cholelithiasis} The presence of stones in the gallbladder or biliary tract. Stones may contain cholesterol, pigment, or mixed material and may cause no symptoms or lead to biliary pain and complications.

\medterm{Cholelithiasis (Gallstones)} Alternate index label for Gallstones.

\medterm{Cholera} An acute intestinal infection caused by toxigenic Vibrio cholerae, producing profuse watery diarrhoea and potentially life-threatening dehydration; prompt oral or intravenous rehydration is central treatment.


\medterm{Choleresis} Choleresis is production and secretion of bile by liver cells.
\medterm{Choleretic} A substance that encourages the liver to make or release bile. Bile helps digest fats and remove some waste products; a choleretic should be used carefully when gallstones or blocked bile ducts may be present.

\synonyms
Bile-secretory agent; bile-flow stimulant.

\medterm{Cholestasis} Reduced production or flow of bile. It may cause itching, yellow skin or eyes, dark urine, pale stools, and abnormal liver tests; a blockage and liver-cell causes must be distinguished.

\medterm{Cholestasis Of Pregnancy} A liver disorder of pregnancy in which impaired bile flow causes itching and elevated bile acids.

\synonyms
Obstetric Cholestasis

\medterm{Cholesteatoma} A keratin-filled sac of abnormal squamous epithelium in the middle ear or mastoid that can erode bone and cause chronic discharge, hearing loss, or intracranial complications.

\medterm{Cholesterol} A lipid required for cell membranes and synthesis of steroid hormones and bile acids. Excess atherogenic cholesterol can contribute to arterial plaque.

\synonyms
Serum cholesterol; blood cholesterol.

\medterm{Blood Cholesterol} Cholesterol carried in the blood. The body needs it to build cell membranes and make hormones and bile acids. Too much LDL cholesterol can form fatty buildup in the arteries and increase the risk of heart disease and stroke. Blood cholesterol is usually checked as part of a lipid profile.

\medterm{Cholesterol And Lifestyle} Diet, physical activity, weight, smoking, and alcohol use can affect blood cholesterol and cardiovascular risk. Lifestyle changes complement medicines when cholesterol remains high.

\medterm{Cholesterol And Lipids} Fatty substances in the blood, including LDL cholesterol, HDL cholesterol, and
triglycerides, that play critical roles in cellular function and energy metabolism.
Abnormal lipid levels, particularly high LDL and triglycerides with low HDL, promote
atherosclerosis and increase the risk of heart attack and stroke.

\medterm{Cholesterol Crystal Embolization} Small pieces of cholesterol plaque breaking loose from a large artery and blocking smaller arteries. It may damage the kidneys, skin, eyes, brain, or intestines, often after a vascular procedure.

\medterm{Cholesterol Crystals} Needle-shaped, cleft-like spaces seen in histological sections representing dissolved
cholesterol esters, commonly found in atherosclerotic plaques, atheroembolic disease,
and cholesterol-rich effusions. In atheroembolic disease, cholesterol crystals embolize
from ruptured plaques to distal organs causing tissue ischemia.

\medterm{Cholesterol Emboli} Cholesterol crystals released from an atherosclerotic plaque that lodge in small arteries and cause
ischemia or inflammation in downstream tissues.

\medterm{Cholesterol Embolism} Blockage of small arteries by cholesterol crystals released from an artery wall, often after vascular procedures or disruption of atherosclerotic plaque.

\medterm{Cholesterol Embolization} A condition in which tiny pieces of fatty material from an artery wall break loose and travel into smaller blood vessels. It can damage the kidneys, skin, eyes, brain, or intestines, sometimes after a vascular procedure or blood-thinning treatment.

\medterm{Cholesterol Lowering (Cholesterol Management)} Alternate index label for Cholesterol Management.

\medterm{Cholesterol Management} Cholesterol management combines cardiovascular-risk assessment with dietary pattern, physical activity, tobacco avoidance, weight care, and medicines when indicated. Statins and other agents lower atherogenic lipoproteins, with intensity chosen according to prior cardiovascular disease and overall risk.

\synonyms
Cholesterol Lowering (Cholesterol Management)

\medterm{Cholesterol Synthesis} The process by which cells, especially liver and intestinal cells, make cholesterol from smaller molecules. Cholesterol is needed for cell membranes, hormones, and bile acids, but excess production contributes to high blood cholesterol.

\medterm{Cholesterol Testing And Results} Laboratory assessment of blood lipids, usually including total cholesterol, LDL cholesterol, HDL cholesterol, and triglycerides, used to estimate cardiovascular risk and guide prevention or treatment.

\medterm{Cholesterolemia} The presence and concentration of cholesterol in the blood. Abnormal elevation is a risk
factor for atherosclerotic cardiovascular disease.

\medterm{Cholesterolosis} Deposition of cholesterol esters in the gallbladder mucosa, producing yellow speckling or a "strawberry gallbladder." It is often incidental but may coexist with gallstones or chronic gallbladder inflammation.

\medterm{Cholestyramine} Cholestyramine is a bile-acid sequestrant used to lower LDL cholesterol and sometimes to treat bile-acid diarrhea or itching from cholestasis.
\medterm{Choline} Choline is an essential nutrient used to make cell membranes and acetylcholine and to support liver and nervous-system function.
\medterm{Cholinergic} Cholinergic means related to acetylcholine or receptors and pathways activated by it, including muscarinic and nicotinic signaling.
\medterm{Cholinergic Crisis} A life-threatening excess of acetylcholine, often after organophosphate or carbamate poisoning. It can cause sweating, salivation, vomiting, diarrhea, pinpoint pupils, wheezing, muscle weakness, seizures, and respiratory failure.

\medterm{Cholinergic Neuron} A nerve cell that releases acetylcholine as its principal neurotransmitter.

\synonyms
Acetylcholine-releasing neuron; cholinergic nerve cell.

\medterm{Cholinergic Toxidrome} A poisoning pattern caused by excessive acetylcholine activity, often producing salivation, tearing, urination, diarrhea, vomiting, bronchial secretions, slow heart rate, muscle fasciculations, weakness, or respiratory failure.

\medterm{Cholinesterase} Cholinesterase enzymes break down acetylcholine; inhibition increases cholinergic activity and is used therapeutically but can also cause poisoning.
\medterm{Cholinesterase - Blood} A blood test measuring acetylcholinesterase or butyrylcholinesterase activity, useful in selected pesticide exposures, inherited enzyme variants, and some liver-function assessments.

\medterm{Chondritis} Inflammation of cartilage, such as the cartilage of the ear, nose, ribs, or larynx. It
may result from infection, trauma, or immune-mediated disease.

\medterm{Chondroblastic Osteosarcoma} A variant of osteosarcoma in which the malignant cells produce a cartilaginous matrix.
It typically affects the long bones of adolescents and young adults, with a predilection
for the metaphyseal regions.

\medterm{Chondroblastoma} A chondroblastoma is an uncommon usually benign cartilage-forming bone tumor, often arising near the end of a long bone in adolescents or young adults.
\medterm{Chondrocalcinosis} Calcification of cartilage visible on X-ray. Associated with CPPD disease. Found in
knees, wrists, pubic symphysis. May be incidental or symptomatic.

\medterm{Chondrocyte} A mature cell that makes and maintains cartilage, the firm flexible tissue found in joints, airways, and other places. Chondrocytes help keep cartilage smooth and strong, but cartilage usually heals slowly because it has little direct blood supply.

\synonyms
Cartilage cell.

\medterm{Chondrodysplasia} A disorder of cartilage and bone development that causes abnormal skeletal growth. Depending on the type, it may produce short limbs, short stature, enlarged joints, or other bone abnormalities.

\medterm{Chondroitin Sulfate} A sulfated glycosaminoglycan found in cartilage and other connective tissues; it is also marketed in supplements for osteoarthritis, with variable evidence of benefit.

\medterm{Chondroma} Benign cartilage-forming tumor composed of mature hyaline cartilage, most often
intramedullary (enchondroma) or subperiosteal (juxtacortical / periosteal chondroma).

\medterm{Chondromalacia} Softening and degeneration of cartilage, especially articular cartilage. The term is
often used for changes affecting the patella and causing anterior knee pain.

\medterm{Chondromyxoid Fibroma} A rare benign cartilage-forming bone tumor, usually arising in the metaphysis of a long bone and presenting with localized pain or an incidental radiographic lesion.

\medterm{Chondrosarcoma} A cancer that forms cartilage-producing tissue, usually in a bone or nearby soft tissue. It often causes persistent pain or a growing lump, and treatment depends on its grade, site, and spread.

\medterm{Chondrosarcoma Of The Breast} An exceptionally rare malignant mesenchymal breast tumor producing cartilage. It is distinguished from metaplastic carcinoma and other sarcomas by histology and is managed with sarcoma-oriented complete excision.






\medterm{Chopart'S Joint} The transverse tarsal joint, consisting of the talonavicular and calcaneocuboid joints,
between the hindfoot and the midfoot.

\medterm{Chorangioma} A benign vascular tumor of the placenta, usually small and incidental but occasionally associated with fetal anemia, hydrops, growth restriction, or other complications when large.

\medterm{Chordae Tendineae} Strong, thin cords inside the heart that connect the mitral and tricuspid valves to small muscles. They help keep the valves from turning inside out when the heart contracts.

\medterm{Chordee} Chordee is abnormal downward or other curvature of the penis, commonly associated with hypospadias or scar tissue and sometimes requiring surgical correction.
\medterm{Chordoma} A rare cancer arising from leftover cells of early spinal development. It usually begins near the tailbone, base of the skull, or spine and can cause pain, weakness, or nerve problems as it grows.

\medterm{Chorea} Uncontrolled, brief, irregular movements that flow unpredictably from one body part to another. It may result from inherited disease, autoimmune illness, medicines, infection, or damage to movement-control areas of the brain.

\medterm{Chorea, Huntington'S} Alternate index label; see \textit{Huntington's disease}.

\medterm{Choreoathetosis} Choreoathetosis combines irregular, dance-like chorea with slow twisting athetotic movements and may result from basal-ganglia injury or genetic disease.
\medterm{Chorioamnionitis} Infection or inflammation of the fluid, membranes, or tissues around a fetus, usually caused by bacteria moving upward from the vagina. Maternal fever, uterine tenderness, foul fluid, or a fast fetal heartbeat may occur.

\medterm{Chorioamnionitis (Intraamniotic Infection)} Infection or inflammation of the amniotic fluid, membranes, placenta, or decidua,
usually due to ascending polymicrobial infection. Maternal fever, uterine tenderness,
fetal tachycardia, and maternal leukocytosis may occur; treatment requires
broad-spectrum antibiotics and delivery planning.

\synonyms
Intraamniotic Infection

\medterm{Choriocarcinoma} Choriocarcinoma is an aggressive trophoblastic cancer producing human chorionic gonadotropin, usually arising after pregnancy or a molar pregnancy.

\medterm{Chorion} The outer membrane surrounding an early embryo. Its finger-like projections help form the fetal part of the placenta, where oxygen, nutrients, and waste are exchanged between fetal and maternal blood supplies.

\medterm{Chorion Villus Sampling} Chorion villus sampling removes placental tissue for genetic or chromosomal testing, usually early in pregnancy; it carries a small procedure-related miscarriage risk.
\medterm{Chorionic Gonadotrophin} Chorionic gonadotrophin, especially hCG, is a placental hormone that maintains early pregnancy and is measured in pregnancy tests and selected tumor evaluations.
\medterm{Chorionic Villus Sampling} A prenatal test usually performed around 10 to 13 weeks by removing a small placental sample through the abdomen or cervix. It can identify many chromosome and inherited disorders but carries a small miscarriage risk.

\medterm{Chorioretinitis} Inflammation involving the choroid and retina, often caused by infection, autoimmune
disease, or systemic inflammatory illness. It can produce floaters and impaired vision.

\medterm{Chorioretinitis (Uveitis)} Alternate index label for Uveitis.

\medterm{Chorioretinitis Toxoplasma (Toxoplasmosis)} Alternate index label for Toxoplasmosis.

\medterm{Choroid} The vascular layer between retina and sclera that supplies the outer retina and absorbs stray light.

\medterm{Choroid Plexus Carcinoma} A cancer arising from the choroid plexus, tissue that produces cerebrospinal fluid, potentially causing pressure buildup, headaches, vomiting, or neurologic problems.

\medterm{Choroid Plexus Cyst} A fluid-filled cyst within the choroid plexus of a fetal cerebral ventricle. An isolated cyst
with otherwise reassuring prenatal assessment is usually benign, but its significance
depends on other ultrasound findings and screening results.

\medterm{Choroid Plexus Tumor} A papillary neoplasm of the choroid plexus, occurring most often in children, and
frequently presenting with hydrocephalus from excessive cerebrospinal fluid production.

\medterm{Choroidal Dystrophies} Inherited disorders causing progressive degeneration of the choroid, retinal pigment epithelium, and retina, leading to night blindness, visual-field loss, and eventual vision impairment.

\medterm{Choroidal Neovascularization} Growth of fragile abnormal blood vessels beneath or through the retina. Leakage or bleeding can damage central vision, especially in wet age-related macular degeneration and some inflammatory or inherited eye disorders.

\medterm{Choroideremia} Choroideremia is an X-linked inherited retinal degeneration causing progressive night blindness and loss of peripheral then central vision, mainly in males.
\medterm{Choroiditis} Choroiditis is inflammation of the choroid, the vascular layer behind the retina, often causing blurred vision, floaters, or visual-field disturbance.
\medterm{Christmas Disease} An inherited bleeding disorder, also called hemophilia B, caused by low or abnormal clotting factor IX. It can cause prolonged bleeding or bleeding into joints and muscles and is treated with replacement or other specialist therapy.
\medterm{Chromaffin} Chromaffin describes cells that stain with chromium salts because they contain catecholamine-rich granules, especially adrenal medullary cells and related neuroendocrine cells.
\medterm{Chromaffinoma} A tumor of cells that can produce adrenaline-like hormones, usually in the adrenal gland or nearby nerve tissue. Hormone release may cause episodes of high blood pressure, headache, sweating, and palpitations.

\medterm{Chromate} Chromate is a hexavalent chromium compound that can cause irritant or allergic dermatitis, nasal injury, respiratory disease, and cancer after occupational exposure.
\medterm{Chromatic} Chromatic means relating to color or wavelengths of light; chromatic vision depends on cone photoreceptors and neural processing.
\medterm{Chromatids} Chromatids are the two copies of a replicated chromosome joined at the centromere before separating during cell division.
\medterm{Chromatin} The combination of DNA and proteins that packages genetic material inside the cell nucleus. Its arrangement controls which genes are available for use and helps regulate DNA copying and repair.

\medterm{Chromatography} A laboratory separation method in which substances partition between a stationary phase and a moving phase, allowing identification or measurement of components in a mixture.

\medterm{Chromatopsia} A visual disturbance in which objects appear abnormally colored. It may result from
retinal disease, toxic exposure, medication effects, or other disorders of the visual
system.

\medterm{Chromhidrosis} A rare condition in which sweat is abnormally colored, often yellow, green, blue, or
black, because of pigment in the sweat glands or on the skin surface.

\medterm{Chromium} Chromium is a metallic element; trivalent chromium is a trace nutrient, while hexavalent compounds can be toxic, corrosive, mutagenic, and carcinogenic.
\medterm{Chromium - Blood Test} A blood test measuring chromium exposure or body burden, generally used when occupational or medical exposure to chromium is suspected; results must be interpreted with the exposure history and clinical findings.

\medterm{Chromium In Diet} Reference information on chromium as a trace element involved in carbohydrate and fat metabolism. Chromium deficiency is rare; most people obtain adequate chromium from normal dietary intake, and routine chromium supplementation is generally not recommended.

\medterm{Chromoblastomycosis} A chronic fungal infection of skin and subcutaneous tissue that produces slowly enlarging warty or
scaly lesions, usually after traumatic implantation of the fungus.

\medterm{Chromopertubation} A fertility test performed during laparoscopy in which colored fluid is passed through the uterus while the fallopian tubes are viewed. Fluid spilling from the tube ends suggests that they are open; no spill may indicate blockage or surrounding scar tissue.

\medterm{Chromosomal Abnormality} Change in chromosome number or structure. Numerical: aneuploidy (abnormal number, e.g.,
trisomy 21), polyploidy (extra sets). Structural: deletion, duplication, inversion,
translocation.

\medterm{Chromosomal Microarray} A genetic test that detects gains and losses of chromosome material, including copy-number variants and some regions of homozygosity; it is commonly used for developmental delay, intellectual disability, congenital anomalies, or autism-spectrum presentations.

\medterm{Chromosome} A chromosome is a DNA-protein structure carrying genetic information; humans normally have 46 chromosomes in most body cells.

\medterm{Chromosome Complement} The complete set of chromosomes in a cell, including their number and sex-chromosome composition, usually described by a karyotype.

\medterm{Chromosome Duplication} A structural chromosome abnormality in which a segment of genetic material is copied one or more additional times, potentially altering gene dosage and development.

\medterm{Chronic} Persisting or recurring over a long period, in contrast to acute; the appropriate duration depends on the disease or clinical context.

\medterm{Chronic Adrenal Insufficiency} Alternate index label; see \textit{Addison's disease}.

\medterm{Chronic Allograft Nephropathy} Gradual loss of function in a transplanted kidney over months or years, caused by immune injury, high blood pressure, recurrent disease, medicines, or other damage. Testing and biopsy help identify the cause.

\medterm{Chronic Atrophic Gastritis} Long-term inflammation that destroys glands in the stomach lining. It may reduce stomach acid and intrinsic factor, impair iron or vitamin B12 absorption, and increase the risk of anemia or stomach tumors.

\medterm{Chronic Autoimmune Hashimoto Thyroiditis} A long-term immune-system attack on the thyroid gland that gradually damages thyroid tissue and often causes too little thyroid hormone. It may produce tiredness, cold intolerance, weight change, or an enlarged thyroid.

\medterm{Chronic Autoimmune Thyroiditis} Long-term inflammation caused by the immune system attacking the thyroid, usually called Hashimoto thyroiditis. It commonly leads to low thyroid hormone and may be associated with thyroid antibodies.

\medterm{Chronic Bacterial Prostatitis} A recurring bacterial infection of the prostate that may cause repeated urinary infections, burning or urgent urination, pelvic discomfort, or painful ejaculation. Diagnosis uses appropriate urine tests, and treatment often lasts several weeks.

\synonyms
NIH Category II Prostatitis; Recurrent Bacterial Prostatitis

\medterm{Chronic Bronchitis} A cough producing mucus on most days for at least three months in each of two consecutive years, after other causes are excluded. It reflects long-term airway irritation and commonly occurs with chronic obstructive pulmonary disease.

\synonyms
Bronchitis Chronic (Chronic Bronchitis)

\medterm{Chronic Caries} Tooth decay that develops slowly over time. It may cause a firm dark area, a cavity, sensitivity, or no symptoms; without dental treatment, decay can reach the inner tooth, causing pain, infection, or tooth loss.

\medterm{Chronic Cholecystitis} Repeated or persistent inflammation of the gallbladder, usually related to gallstones and impaired emptying. It may cause episodic right-upper-abdominal pain, nausea, or food intolerance and can lead to scarring or a contracted gallbladder.

\medterm{Chronic Compartment Syndrome} Alternate index label; see \textit{Chronic exertional compartment syndrome}.

\medterm{Chronic Cough} A cough lasting longer than eight weeks. Common causes include asthma, nasal drainage, acid reflux, smoking, and some blood-pressure medicines; persistent cough requires assessment for its cause.

\synonyms
Cough Chronic (Chronic Cough)

\medterm{Chronic Daily Headaches} Headaches occurring on at least fifteen days each month for more than three months, caused by migraine, tension, medication overuse, or other conditions.

\medterm{Chronic Disease} A disease or condition that persists for a prolonged period, often months or years, and may require continuing monitoring or management.

\medterm{Chronic Endometritis} Persistent low-grade inflammation of the endometrium, often associated with infection, retained products, intrauterine devices, or reproductive failure; plasma cells support the diagnosis histologically.
delivery. Risk factors include prolonged labor, multiple vaginal examinations, premature
rupture of membranes, and cesarean section. Presents with fever, uterine tenderness,
foul-smelling lochia, and leukocytosis.

\medterm{Chronic Eosinophilic Leukaemia} A clonal myeloid neoplasm characterized by persistent eosinophilia, increased or
dysplastic eosinophils, and evidence of an underlying clonal process or organ injury.
Symptoms may include fatigue, fever, weight loss, rash, respiratory symptoms,
splenomegaly, or cardiac and neurological complications.

\medterm{Chronic Eosinophilic Pneumonia} An idiopathic inflammatory lung disease characterized by eosinophilic infiltration of
the pulmonary parenchyma, presenting with fever, cough, dyspnea, and peripheral airspace
opacities on imaging.

\medterm{Chronic Exertional Compartment Syndrome} Exercise-related increased pressure within a muscle compartment that causes predictable pain, tightness, or neurologic symptoms during activity.

\synonyms
Chronic Compartment Syndrome

\medterm{Chronic Fatigue Syndrome} A long-term illness involving substantial reduced ability to function, worsening after physical or mental activity, unrefreshing sleep, and often difficulty thinking or standing. Symptoms are not explained by another condition and need careful evaluation.

\synonyms
CFS (Chronic Fatigue Syndrome)
SEID (Chronic Fatigue Syndrome)

\medterm{Chronic Gastritis} Chronic, progressive inflammation of the gastric mucosa leading over time to mucosal atrophy, loss of specialized parietal and chief cells, and intestinal metaplasia. Major causes include persistent \textit{Helicobacter pylori} infection and autoimmune gastritis targeting parietal cells and intrinsic factor.

\medterm{Chronic Glomerulonephritis} Long-standing inflammation and scarring of the kidney glomeruli that can progressively reduce filtration and lead to chronic kidney disease or kidney failure.

\medterm{Chronic Granulomatous Disease} An inherited immune disorder in which certain white blood cells cannot effectively kill some bacteria and fungi. It causes repeated serious infections, swollen lymph nodes, pneumonia, skin infections, and inflammation called granulomas.

\medterm{Chronic Hives} Recurrent itchy, raised welts occurring repeatedly for more than six weeks. Individual welts usually fade within a day, while causes may include immune activity, medicines, infections, or no identifiable trigger.

\medterm{Chronic Hypertension In Pregnancy} Hypertension that predates pregnancy or is diagnosed early in pregnancy.
It increases the risk of preeclampsia, fetal growth problems, preterm birth, and placental complications and requires
ongoing maternal and fetal monitoring.

\medterm{Chronic Inflammation} Inflammation that continues for weeks, months, or years. It can occur with persistent infection, autoimmune disease, long-term irritation, or obesity and may cause ongoing tissue damage while the body attempts repair.

\medterm{Chronic Inflammatory Demyelinating Polyneuropathy} An immune disorder that damages the insulating covering of nerves outside the brain and spinal cord. It causes gradually worsening or recurring weakness, numbness, poor balance, and reduced reflexes in the limbs.

\medterm{Chronic Inflammatory Response Syndrome} A proposed multisystem illness involving persistent inflammatory symptoms after an exposure
such as certain water-damaged building environments. Reported symptoms are nonspecific and may include fatigue,
headache, cognitive difficulty, and respiratory complaints. Symptoms require a careful medical evaluation for
established infectious, allergic, toxic, autoimmune, and other causes.

\medterm{Chronic Intestinal Pseudo-Obstruction} Repeated symptoms of bowel blockage without a physical obstruction, caused by poor nerve or muscle function in the intestine. It may cause abdominal swelling, pain, vomiting, constipation, and inability to absorb nutrition.

\medterm{Chronic Kidney Disease} Long-term damage or reduced function of the kidneys lasting at least three months. Common causes include diabetes and high blood pressure; advanced disease can cause anemia, bone problems, fluid buildup, and kidney failure.

\synonyms
Chronic Kidney Failure
Chronic Renal Failure
Kidney Disease, Chronic
Kidney Failure, Chronic

\medterm{Chronic Kidney Disease Staging} Classification of kidney disease by cause, eGFR category, and albuminuria category. G1
through G5 represent eGFR from normal or high to kidney failure; staging is applied when
abnormalities persist for at least 3 months and helps estimate progression and
complications.

\medterm{Chronic Kidney Failure} Alternate index label; see \textit{Chronic kidney disease}.

\medterm{Chronic Lead Poisoning} Illness caused by long-term or repeated exposure to lead. It may cause abdominal pain, constipation, tiredness, anemia, nerve damage, high blood pressure, or kidney injury. In children, it can impair learning, attention, growth, or development. Treatment begins with stopping the exposure; some people need specialist-directed chelation.

\medterm{Chronic Leukemia} A blood cancer that develops gradually as abnormal white blood cells multiply and remain in the blood or bone marrow. Some types progress slowly but still require regular medical monitoring.

\medterm{Chronic Limb-Threatening Ischemia} Severe, long-lasting reduction of blood flow to a limb, usually from narrowed or blocked arteries. It causes pain at rest, ulcers that do not heal, or gangrene and requires urgent vascular assessment to protect the limb.

\medterm{Chronic Liver Disease} Persistent hepatic injury and structural change lasting at least 6 months, caused by
viral, metabolic, alcohol-related, autoimmune, cholestatic, vascular, or drug-related
disease. Progressive fibrosis can lead to cirrhosis, portal hypertension, liver failure,
and hepatocellular carcinoma.

\medterm{Chronic Lymphocytic Leukemia} The most common leukemia in Western countries, characterized by the clonal proliferation
of mature-appearing but functionally incompetent B lymphocytes that accumulate in blood,
bone marrow, and lymphoid tissues.

\synonyms
Cancer, Chronic Lymphocytic Leukemia
CLL
Leukemia, Chronic Lymphocytic

\medterm{Chronic Lymphocytic Thyroiditis} Alternate index label; see \textit{Hashimoto's disease}.

\medterm{Chronic Mesenteric Ischemia} Reduced blood flow to the intestines, usually from hardening and narrowing of abdominal arteries. It often causes cramping after meals, food avoidance, and weight loss and needs vascular evaluation.

\medterm{Chronic Migraine} Headache on at least 15 days each month for more than three months, with migraine features on at least eight days. Pain may be one-sided or throbbing and may include nausea or sensitivity to light or sound.

\medterm{Chronic Motor Or Vocal Tic Disorder} A neurodevelopmental disorder with persistent motor or vocal tics lasting more than one year; symptoms often fluctuate and may be treated with education, behavioral therapy, or medicines when impairment is significant.

\medterm{Chronic Mucocutaneous Candidiasis} Repeated or persistent yeast infections of the skin, nails, or moist linings of the mouth, throat, or vagina. It may result from an inherited immune problem and can occur with autoimmune or hormone disorders.

\medterm{Chronic Myelogenous Leukemia} A blood cancer in which the bone marrow makes too many abnormal white blood cells, usually because of the BCR::ABL1 gene change. It may progress over time but is often controlled with targeted medicines.

\medterm{Chronic Myeloid Leukemia} A blood cancer usually caused by the Philadelphia chromosome, which creates the BCR::ABL1 gene. The bone marrow produces too many abnormal white blood cells, and targeted treatment can often control the disease.

\medterm{Chronic Myelomonocytic Leukemia} A rare blood cancer in which the bone marrow makes too many monocytes, a type of white blood cell, along with abnormal blood-forming cells. It can cause anemia, infections, bleeding, fatigue, or an enlarged spleen.

\medterm{Chronic Neutrophilic Leukaemia} A rare clonal myeloproliferative neoplasm characterized by sustained mature neutrophilic
leukocytosis without a primary infectious or inflammatory cause. It may cause fatigue,
fever, weight loss, splenomegaly, and progressive constitutional symptoms.

\medterm{Chronic Obstructive Lung Disease} An older abbreviation for chronic obstructive lung disease, now usually called chronic obstructive pulmonary disease, characterized by persistent airflow limitation.
\medterm{Chronic Obstructive Pulmonary Disease} A long-term lung disease in which damaged airways or air sacs limit airflow. It causes cough, mucus, wheezing, and breathlessness, and is most often related to smoking or prolonged exposure to irritating dusts or gases.

\synonyms
COPD (Chronic Obstructive Pulmonary Disease)

\medterm{Chronic Pain} Pain that persists or recurs for longer than three months. It may have an identifiable ongoing cause or persist after healing, and is influenced by
complex interactions among biological, psychological, and social factors.


\medterm{Chronic Pain Management} Care aimed at reducing long-lasting pain and improving sleep, movement, and daily function. It may combine exercise, physical therapy, psychological support, nonopioid medicines, procedures, and carefully monitored opioids when appropriate.

\medterm{Chronic Pain Syndrome} Long-lasting pain accompanied by substantial effects on movement, sleep, mood, relationships, or daily activities. Treatment addresses pain and these related effects rather than relying only on pain medicines.
\medterm{Chronic Pancreatitis} Repeated or long-term inflammation that permanently scars the pancreas. It may cause upper-abdominal pain, poor digestion, weight loss, diabetes, or oily stools; alcohol, smoking, inherited disorders, and blocked ducts are possible causes.

\medterm{Chronic Paroxysmal Hemicrania} A primary headache disorder causing frequent, short-lasting, severe one-sided headaches with autonomic symptoms and a characteristic response to indomethacin.

\synonyms
CPH; chronic paroxysmal hemicrania syndrome.

\medterm{Chronic Pelvic Pain} Pain in the lower abdomen or pelvis that persists or returns for at least six months. Causes may involve reproductive, urinary, digestive, muscle, nerve, or pain-processing disorders and require individualized evaluation.
\synonyms
Pelvic Pain, Chronic

\medterm{Chronic Pelvic Pain Syndrome} Persistent or recurring pain in the pelvis, area between the genitals and anus, lower abdomen, back, or scrotum without a proven ongoing bacterial infection. Urinary symptoms, painful ejaculation, muscle tension, and sexual problems may also occur.

\synonyms
CPPS; NIH Category III Prostatitis; Chronic Nonbacterial Prostatitis

\medterm{Chronic Progressive External Ophthalmoplegia} A slowly worsening muscle disorder causing drooping eyelids and limited movement of both eyes. It often begins in adulthood and may occur alone or with weakness of other muscles; inherited mitochondrial disease is a common cause.

\synonyms
CPEO; Progressive External Ophthalmoplegia (PEO)

\medterm{Chronic Recurrent Multifocal Osteomyelitis} A rare autoinflammatory bone disorder, often beginning in children or adolescents, causing recurrent sterile inflammation and pain at multiple bone sites; it is also called CRMO.

\medterm{Chronic Pyelonephritis} Long-term inflammation and scarring of the kidneys, often after repeated infections or backward flow of urine. It can gradually reduce kidney function and may cause high blood pressure, poor growth, or kidney failure.

\medterm{Chronic Renal Failure} Alternate index label; see \textit{Chronic kidney disease}.

\medterm{Chronic Rhinitis} Persistent inflammation or irritation of the nasal lining causing congestion, runny nose, sneezing, itching, or
postnasal drainage. Allergic, nonallergic, medication-related, and structural causes are possible. Management depends
on the cause and may include trigger avoidance, saline irrigation, and clinician-directed nasal treatment.

\synonyms
Post-Nasal Drip (Chronic Rhinitis)

\medterm{Chronic Rhinosinusitis} Persistent inflammation of the nose and paranasal sinuses lasting at least 12 weeks, usually
with nasal obstruction or discharge and sometimes facial pressure or reduced smell. It
may occur with or without nasal polyps.

\medterm{Chronic Sinusitis} Persistent inflammation of the paranasal sinuses for at least 12 weeks, causing nasal
obstruction, discharge, facial pressure, and reduced smell.

\synonyms
Sinusitis, Chronic

\medterm{Chronic Subdural Hematoma} A collection of blood and breakdown products between the dura and arachnoid that develops over weeks, often after minor trauma in older adults, causing headache, confusion, weakness, or gait change.

\medterm{Chronic Suppurative Otitis Media} Persistent inflammation and infection of the middle ear with recurrent or continuous ear discharge,
usually through a perforated tympanic membrane. It can cause conductive hearing loss and
may lead to cholesteatoma or other complications.

\synonyms
CSOM; chronic otitis media with discharge; chronic suppurative middle-ear infection.

\medterm{Chronic Tension-Type Headache} A recurring headache occurring on at least 15 days per month for more than three months, usually causing
bilateral pressing or tightening pain rather than throbbing. It may be associated with scalp or neck tenderness
but typically lacks prominent nausea; a new or changing pattern requires clinical assessment.

\medterm{Chronic Thromboembolic Pulmonary Hypertension} Pulmonary hypertension caused by organized, persistent thromboembolic obstruction and
secondary pulmonary-vessel remodeling after pulmonary embolism. It is potentially
curable in selected patients by pulmonary endarterectomy or balloon pulmonary
angioplasty.

\medterm{Chronic Thyroiditis (Hashimoto Disease)} An autoimmune disease in which immune attack gradually damages the thyroid, often causing hypothyroidism, fatigue, cold intolerance, constipation, dry skin, or goiter.

\medterm{Chronic Toxicity} Adverse effects from repeated or prolonged exposure to low doses. Cumulative damage may
be irreversible. Examples: heavy metals (lead, mercury), asbestos, alcohol, solvent
exposure. Monitoring: biomarkers, organ function tests.

\medterm{Chronic Traumatic Encephalopathy} A progressive brain disorder associated with repeated head impacts and characterized by cognitive, mood, or behavioral changes.

\medterm{Chronic Venous Disease} A spectrum of venous outflow disorders caused by reflux, obstruction, or
calf-muscle-pump dysfunction. Findings range from edema and varicose veins to skin
changes, lipodermatosclerosis, and venous ulceration.

\medterm{Chronic Venous Insufficiency} Impaired venous return from the lower extremities due to venous valve incompetence or
obstruction. Leads to venous hypertension, edema, skin changes (lipodermatosclerosis,
hyperpigmentation), and venous stasis ulcers (medial malleolus). Risk factors: varicose
veins, DVT, obesity, prolonged standing.

\medterm{Chronic Venous Stasis} Persistent pooling of blood in the lower limbs caused by impaired venous return. It can produce
edema, skin discoloration, dermatitis, and venous ulcers.

\medterm{Chronic Vulvar Pain} Alternate index label; see \textit{Vulvodynia}.

\medterm{Chronicity} How long a disease, symptom, or process has lasted or how long it tends to continue. A condition may be described as acute when it is recent and short-lasting, or chronic when it persists over a long period.

\medterm{Chronobiology} Chronobiology studies biological timing and rhythms, including circadian sleep-wake cycles, seasonal rhythms, and their effects on health.
\medterm{Chronotropic} Chronotropic describes an influence on heart rate; a positive chronotropic effect increases rate, while a negative effect decreases it.
\medterm{Chrysarobin} Chrysarobin is a plant-derived anthracene compound formerly used topically for psoriasis and other skin disorders; it can irritate and stain skin.
\medterm{Churg Strauss Syndrome (Churg-Strauss Syndrome)} Alternate index label for Churg-Strauss Syndrome.

\medterm{Churg-Strauss Syndrome} A rare disease in which inflamed small and medium blood vessels are associated with asthma and too many eosinophils, a type of white blood cell. It can damage the lungs, nerves, skin, heart, or kidneys.

\synonyms
Churg Strauss Syndrome (Churg-Strauss Syndrome)

\medterm{Chvostek Sign} Twitching of facial muscles after tapping over the facial nerve near the ear. It can occur with low blood calcium but is not reliable enough by itself to diagnose that condition.

\medterm{Chvostek'S Sign} Chvostek’s sign is twitching of facial muscles after tapping near the facial nerve, sometimes associated with hypocalcemia but neither sensitive nor specific.
\medterm{Chylangioma} A usually noncancerous abnormal collection of lymphatic vessels, often in the neck or chest. It may form a soft swelling and can press on nearby structures; the modern term is often lymphangioma.

\medterm{Chyle} A milky fluid formed in intestinal lymph vessels after a meal. It contains lymph and absorbed dietary fat and normally travels through the thoracic duct into the bloodstream.

\medterm{Chylomicron} A particle made by the intestine after eating fat. It carries dietary fat and cholesterol through the lymph and blood, then leaves remnants for the liver to remove.

\medterm{Chylomicronemia Syndrome} Severe elevation of triglyceride-rich chylomicrons in blood, usually from genetic or acquired impairment of triglyceride clearance; it can cause abdominal pain, eruptive xanthomas, and recurrent acute pancreatitis.

\medterm{Chylothorax} A buildup of chyle in the space around a lung, usually after injury or blockage of the thoracic duct. It can cause breathlessness and requires treatment of the leak, often with drainage and dietary or surgical measures.

\medterm{Chyluria} Passage of lymphatic fluid into the urine, giving it a milky appearance. It usually
results from an abnormal communication between lymphatic vessels and the urinary tract.

\medterm{Chyme} A semifluid mixture of partially digested food and gastric or intestinal secretions that passes from the stomach into the small intestine.

\synonyms
Digestive chyme; gastric chyme.

\medterm{Chymotrypsin} Chymotrypsin is a pancreatic digestive enzyme that cleaves proteins in the small intestine; abnormal pancreatic release can contribute to tissue injury.
\medterm{Ci (Curie)} A unit of radioactivity equal to 3.7 billion nuclear disintegrations per second; the SI unit is the becquerel.

\medterm{CIC} Commonly, clean intermittent catheterization: periodic insertion of a catheter to empty the bladder, followed by removal; other expansions are possible by context.

\medterm{Cicatricial Alopecia} Hair loss caused by destruction and scarring of hair follicles, usually leading to permanent loss in affected areas.

\synonyms
Scarring alopecia; cicatricial hair loss.

\medterm{Cicatrization} Cicatrization is healing by formation of a scar after tissue injury, with replacement of damaged tissue by fibrous connective tissue.
\medterm{Ciclopia} A rare severe congenital malformation in which the eyes are fused or a single median eye
develops, usually as part of holoprosencephaly.

\medterm{Cidofovir} Cidofovir is an antiviral nucleotide analogue used for selected serious DNA-virus infections; kidney toxicity requires careful dosing and monitoring.
\medterm{CIDP} Chronic inflammatory demyelinating polyneuropathy is a chronic, immune-mediated,
peripheral neuropathy characterized by progressive or relapsing symmetrical weakness and
sensory loss in all four limbs, with areflexia or hyporeflexia, lasting for more than
eight weeks.

\medterm{Cigarette} A small cylinder of tobacco or other material wrapped in paper for smoking; smoke exposure delivers nicotine and many toxic and carcinogenic substances.

\medterm{Ciguatera} Food poisoning caused by eating reef fish containing ciguatoxins. It may cause
gastrointestinal illness followed by neurologic symptoms such as paresthesias and
temperature-sensation disturbance.

\medterm{Ciguatera Poisoning} Ciguatera poisoning follows eating reef fish containing ciguatoxins. Gastrointestinal symptoms may be followed by neurologic sensations such as temperature reversal, tingling, itching, weakness, or abnormal heart rate; treatment is supportive and prevention depends on avoiding high-risk fish.

\synonyms
Poisoning Ciguatera (Ciguatera Poisoning)
Toxin Ciguatera (Ciguatera Poisoning)

\medterm{Ciguatoxin} A heat-stable toxin produced by marine microorganisms and accumulated in some reef fish, causing ciguatera poisoning with gastrointestinal, neurologic, and temperature-sensation symptoms.

\medterm{Cilia} Tiny hair-like structures on cells that move fluid, mucus, or particles across a tissue surface. In the airways, they help clear inhaled material; some cilia also sense signals and help move an egg cell.
\medterm{Ciliary Body} Circumferential structure between iris and choroid, with ciliary muscle (accommodation
control via zonular fibers) and ciliary processes (aqueous humor production). Part of
the uveal tract (iris, ciliary body, choroid).

\medterm{Ciliary Dyskinesia} Abnormal movement of cilia that impairs clearance of mucus from the airways and can cause recurrent respiratory infections.

\medterm{Ciliary Muscle} A ring of smooth muscle in the eye that changes the lens shape for focusing at different distances and helps regulate aqueous humor outflow through ciliary-body structures.

\medterm{Ciliary Neuralgia} Pain arising from or felt along the ciliary nerves around the eye, often described as deep orbital or eye pain; the cause requires clinical evaluation.

\medterm{Ciliopathy} A disorder caused by abnormal cell cilia, the tiny structures that move material or sense signals. It may affect development, vision, kidney function, breathing, or several organs.

\medterm{Cimetidine} Cimetidine is an H2-receptor antagonist that reduces gastric acid secretion; it has numerous drug interactions and is used less often than newer agents.
\medterm{Cinacalcet} A calcimimetic drug that activates calcium-sensing receptors and lowers parathyroid hormone, used for selected forms of hyperparathyroidism and hypercalcemia.

\medterm{Cineangiocardiography} Radiographic recording of the heart and blood vessels during passage of contrast
material, used to visualize cardiac chambers, valves, and coronary or great vessels.

\medterm{Cinefluorography} Motion-picture recording of fluoroscopic images, used to document movement of organs or
contrast material during a diagnostic examination.

\medterm{Cineplasty} Surgical creation of a muscle-controlled stump or related reconstruction, historically
used to improve control of a prosthetic limb.

\medterm{Cingulotomy} Surgical interruption of selected fibers of the cingulum, sometimes performed as a
psychosurgical treatment for severe, treatment-resistant psychiatric illness or chronic
pain.

\medterm{Cinnarizine} Cinnarizine is an antihistamine with calcium-channel-blocking properties used in some countries for motion sickness and vestibular symptoms.
\medterm{Ciprofloxacin} Ciprofloxacin is a fluoroquinolone antibiotic used for selected bacterial infections; tendon injury, nerve effects, QT effects, and resistance limit indiscriminate use.
\medterm{Circadian Rhythm} An internal cycle of about 24 hours that helps regulate sleep, alertness, hormones, body temperature, and metabolism. Light and darkness reset the cycle each day through signals from the eyes to the brain.

\medterm{Circadian Rhythms} Approximately 24-hour biological cycles that regulate sleep and wakefulness, hormone release, body temperature, metabolism, and other physiological functions in response chiefly to light and darkness.
\medterm{Circle Of Willis} A ring of arteries at the base of the brain that links the main front and back blood supplies. It can provide an alternate route for blood if one artery becomes narrowed or blocked.

\medterm{Circulating Tumor DNA} Fragments of tumor-derived DNA detectable in blood, enabling molecular profiling,
treatment monitoring, and detection of minimal residual disease without tissue biopsy
(liquid biopsy).

\medterm{Circulation} Continuous movement of blood through the heart and blood vessels, or of lymph through lymphatic vessels, to transport oxygen, nutrients, hormones, and waste.

\medterm{Circulatory} Relating to the movement of blood or lymph and to the organs and vessels that carry them.

\medterm{Circulatory System} The heart and blood-vessel network that pumps and distributes blood through pulmonary and systemic circuits; broadly, the term may include the lymphatic system.

\medterm{Circumcision} Surgical removal of the foreskin covering the head of the penis. It may be performed for cultural, religious, preventive, or medical reasons; possible complications include bleeding, infection, pain, and injury, although serious problems are uncommon with proper care.
medical reasons such as severe phimosis, recurrent balanoposthitis, paraphimosis, or
lichen sclerosus (balanitis xerotica obliterans); elective infant circumcision for
religious, cultural, or preventive health reasons.

\medterm{Circumflex} Curving around; the circumflex coronary artery is a branch of the left coronary artery that travels around the heart's lateral and posterior surface.

\medterm{Circumflex Artery} A branch of the left coronary artery that travels in the left atrioventricular groove.
It supplies the lateral wall of the left ventricle, the left atrium, and in
left-dominant systems, the inferior wall and posteromedial papillary muscle.

\medterm{Circumvallate Papillae} Circumvallate papillae are large taste papillae arranged near the back of the tongue and surrounded by a trench containing taste buds and serous glands.
\medterm{Circumvallate Placenta} An extrachorial placental configuration in which the chorionic plate is smaller than the basal plate, producing a raised circumferential edge; it may be associated with bleeding, abruption, or growth restriction.

\medterm{Cirrhosis} Advanced scarring of the liver that permanently changes its structure and reduces blood flow and function. Causes include viral hepatitis, alcohol-related disease, fatty-liver disease, and autoimmune disorders; complications include jaundice, fluid buildup, bleeding, and liver failure.


\medterm{Cirrhosis Primary Biliary (Primary Biliary Cirrhosis)} Alternate index label for Primary Biliary Cirrhosis.

\medterm{Cisplatin} A platinum-based chemotherapy drug that forms DNA cross-links and is used to treat several cancers; important toxicities include kidney injury, hearing loss, neuropathy, and nausea.

\medterm{Citalopram} Citalopram is a selective serotonin-reuptake inhibitor used to treat depression and some anxiety disorders; nausea, sexual effects, and QT prolongation are important considerations.
\medterm{Citrate} Citrate is a salt or ion of citric acid that participates in metabolism and binds calcium; it is used in anticoagulant blood collection and some stone-prevention treatments.
\medterm{Citrate Synthase} An enzyme that begins an important energy-producing process inside cells. It joins acetyl-CoA with oxaloacetate to make citrate, the first step of the citric acid cycle, which helps release energy from food.

\medterm{Citric Acid Cycle} The citric acid cycle is a mitochondrial pathway that oxidizes acetyl-CoA and generates reducing equivalents for ATP production.
\medterm{Citric Acid Cycle (Krebs Cycle)} A central metabolic pathway occurring in the mitochondrial matrix that oxidizes
acetyl-CoA (derived from carbohydrates, lipids, and proteins) into CO2, generating 3
NADH, 1 FADH2, and 1 GTP (ATP) per turn. Rate-limiting enzyme: isocitrate dehydrogenase
(inhibited by ATP and NADH; activated by ADP).

\synonyms
Krebs Cycle

\medterm{Citric Acid Urine Test} A urine test measuring citrate excretion, used mainly in the evaluation of recurrent kidney stones because low urinary citrate increases the risk of calcium stone formation.

\medterm{Citrullinaemia} An inherited disorder of the urea cycle causing citrulline and ammonia buildup, which can lead to vomiting, confusion, seizures, coma, or death.
\medterm{Citrullinemia} An inherited disorder in which the body cannot properly remove nitrogen through the urea cycle. Citrulline and ammonia build up, causing vomiting, poor feeding, sleepiness, confusion, seizures, coma, or death without urgent treatment.
\medterm{CJD} Alternate index label; see \textit{Creutzfeldt-Jakob disease}.

\medterm{CK} Creatine kinase, an enzyme found in skeletal muscle, heart muscle, and brain whose blood level may rise after tissue injury.

\synonyms
Creatine phosphokinase; creatine kinase enzyme.

\medterm{CKD-Mineral Bone Disorder} A complication of chronic kidney disease in which abnormal calcium, phosphate, vitamin D, and parathyroid hormone levels weaken bones and may cause calcium deposits in blood vessels and other tissues.

\medterm{Cl (Chloride)} The chemical symbol for chloride, the negatively charged ion of chlorine and a major extracellular electrolyte.

\medterm{Cladribine} Cladribine is a purine analogue used for selected blood cancers and multiple sclerosis; lymphocyte depletion and infection risk require monitoring.
\medterm{Clairvoyance} Clairvoyance is an alleged ability to perceive information beyond ordinary sensory means; it is not an established medical diagnostic finding.
\medterm{Clamp} A surgical instrument used to grasp, compress, occlude, or hold tissue, blood vessels, tubing, or other structures.

\medterm{Clang Association} Clang association is speech in which word choice is driven by sound, rhyme, or pun rather than meaning, seen particularly in mania or psychosis.
\medterm{Clarithromycin} Clarithromycin is a macrolide antibiotic used for selected bacterial and mycobacterial infections; interactions and QT prolongation are important risks.
\medterm{Clark Level Of Invasion} An older system describing how deeply a melanoma has grown through the skin layers. Modern melanoma assessment relies more on tumor thickness, ulceration, spread, and other features.

\medterm{Class IA Antiarrhythmics} Heart-rhythm medicines that slow electrical conduction and lengthen the time heart cells need before they can fire again. They treat selected abnormal rhythms but can themselves cause dangerous rhythm disturbances.

\medterm{Class IB Antiarrhythmics} Heart-rhythm medicines that reduce abnormal electrical activity, especially in injured heart muscle. They are used mainly for selected fast rhythms arising from the lower chambers and require monitoring.

\medterm{Class IC Antiarrhythmics} Strong heart-rhythm medicines that slow electrical conduction without greatly changing recovery time. They can treat selected upper-chamber rhythms but may be dangerous in people with structural heart disease.

\medterm{Class III Antiarrhythmics} Heart-rhythm medicines that lengthen the recovery period between electrical beats. They treat selected abnormal rhythms but may slow the heart, affect organs, or prolong the electrical interval and cause dangerous rhythms.

\medterm{Classic Migraine} A migraine attack preceded or accompanied by an aura, such as temporary visual, sensory, or speech changes. Headache may follow, with nausea and light or sound sensitivity, although patterns vary.

\medterm{Claudication} Muscle pain or cramping brought on by exertion and relieved by rest, usually because of reduced arterial blood flow to a limb.

\synonyms
Vascular Claudication

\medterm{Claudication Distance} The distance a person can walk before exertional limb pain begins. It is used to describe symptoms and monitor change.

\medterm{Claustrophobia} An intense fear of enclosed or confined spaces that may cause anxiety, panic, or avoidance.
Symptoms can occur in places such as elevators, tunnels, or MRI scanners.

\medterm{Clavicle} The collarbone, a curved bone running from the breastbone to the shoulder blade. It supports the shoulder and commonly breaks after a fall or direct blow, causing pain, swelling, or visible deformity.

\medterm{Clavicle Fracture} A break in the collarbone that may occur during birth or after trauma. A newborn may show
asymmetric arm movement or tenderness; older children and adults may have pain, swelling, or deformity. Diagnosis is
usually made by examination and imaging.

\medterm{Clavicle Fracture (Intentional)} Intentional fracture of the fetal clavicle as a last resort to reduce the bisacromial
diameter in shoulder dystocia. The operator applies upward pressure on the midportion of
the anterior fetal clavicle with a finger, fracturing it. The clavicle heals without
long-term sequelae.

\synonyms
Intentional

\medterm{Clavus} A clavus is a localized painful thickening of skin, or corn, caused by repeated pressure or friction.
\medterm{Claw Foot} A toe deformity in which the toes bend upward at their bases and downward at their middle and end joints. It may result from muscle imbalance, nerve damage, arthritis, or unsuitable footwear.

\medterm{Claw Hand} A deformity in which the large finger joints bend backward while the smaller joints bend inward. It is commonly caused by ulnar-nerve damage and may impair grip and hand function.

\medterm{Claw Toe} A toe deformity characterized by extension at the metatarsophalangeal joint and flexion
at the proximal and distal interphalangeal joints.

\medterm{Clean Catch Urine Sample} A midstream urine specimen collected after cleaning the genital area to reduce contamination and improve interpretation of urinalysis or culture.

\medterm{Clean-Catch Specimen} A urine sample collected after cleansing the surrounding area and beginning urination, with the middle portion collected in a sterile container to reduce contamination.

\medterm{Cleaning Supplies And Equipment} Materials and devices used to remove soil and microorganisms from surfaces or equipment; selection and disinfection level depend on the item and infection risk.

\medterm{Cleaning To Prevent The Spread Of Germs} Cleaning and disinfection practices that remove or inactivate infectious organisms on hands, surfaces, equipment, and shared objects to reduce transmission.

\medterm{Clear Cell Carcinoma} Clear cell carcinoma is a cancer whose cells appear clear because of glycogen or lipid; it can arise in organs such as kidney, ovary, or endometrium.

\medterm{Clear Cell Carcinoma Of The Ovary} A type of ovarian cancer whose cells look clear under a microscope. It is often associated with endometriosis and may be found at an earlier stage, but can resist some standard treatments and increase the risk of blood clots; treatment depends on stage and overall health.

\medterm{Clear Cell Renal Cell Carcinoma} The predominant histologic subtype of renal cell carcinoma, accounting for approximately 75--80\% of all adult primary kidney cancers. Arises from the proximal convoluted tubule and is histologically characterized by malignant cells with clear lipid- and glycogen-rich cytoplasm; strongly associated with loss or inactivation of the \textit{VHL} tumor suppressor gene.

\synonyms
ccRCC; Conventional renal cell carcinoma
\medterm{Clear Liquid Diet} A limited diet consisting of transparent liquids such as water, clear broth, apple juice, and black coffee that leave minimal residue. It is used for short periods before procedures, during acute illness, or after surgery to maintain hydration while resting the gastrointestinal tract.

\medterm{Clear Margins} A histological finding indicating that no abnormal cells (e.g., cervical intraepithelial
neoplasia or cancer cells) are present at the edge of a tissue specimen removed during a
procedure such as LLETZ, cone biopsy, or surgical excision.

\medterm{Clearance} The volume of plasma from which a drug is completely removed per unit time, typically
expressed in milliliters per minute or liters per hour. Total clearance is the sum of
all clearance mechanisms, primarily hepatic and renal clearance.

\medterm{Cleft Lip} A birth defect in which the upper lip does not join together fully
before birth. It may occur alone or with a cleft palate.

\medterm{Cleft Lip And Cleft Palate} Birth defects in which the upper lip, the roof of
the mouth, or both do not join together fully before birth. They can cause problems
with feeding, speech, hearing, or teeth and may need treatment by several specialists.

\medterm{Cleft Lip And Palate Repair} Surgery to close a cleft in the lip, palate, or both and improve feeding, speech, hearing, breathing, or appearance. Repair is usually staged during childhood and requires long-term dental and specialist follow-up.


\medterm{Cleft Palate} A birth defect in which there is an opening in the roof of the
mouth. It may occur alone or with a cleft lip and can cause problems with feeding,
speech, hearing, or teeth.



\medterm{Cleidocranial Dysostosis} An inherited bone disorder, usually caused by a RUNX2 gene change, in which the collarbones may be small or absent and skull bones close late. Extra teeth, delayed tooth eruption, short stature, and joint problems are common.

\medterm{Cleidocranial Dysplasia} An inherited disorder of bone development affecting the collarbones and skull. It can cause delayed closure of skull joints, unusual teeth, short stature, and other skeletal differences.

\medterm{Clemastine} Clemastine is a first-generation antihistamine used for allergic symptoms; sedation, anticholinergic effects, and impaired driving are concerns.
\medterm{Clendamycin} Clendamycin is a misspelling of clindamycin, a lincosamide antibiotic that inhibits bacterial protein synthesis and can cause severe antibiotic-associated colitis.
\medterm{Click-Murmur Syndrome} Alternate index label; see \textit{Mitral valve prolapse}.

\medterm{Climacteric} The period of physiologic change around the end of reproductive capacity, including the
menopausal transition. It may involve changes in menstruation, vasomotor symptoms, and
other endocrine effects.

\medterm{Climate} Climate is the long-term pattern of weather and environmental conditions; climate affects infectious disease, heat illness, allergies, and population health.
\medterm{Clindamycin} Clindamycin is a lincosamide antibiotic used for selected infections; diarrhea and Clostridioides difficile colitis are important adverse effects.
\medterm{Clinic Visit} A planned outpatient appointment for consultation, examination, treatment, follow-up, or monitoring. It may include reviewing symptoms, medicines, test results, prevention needs, and decisions about further care.

\medterm{Clinical} Relating to direct observation, examination, diagnosis, or treatment of a patient, rather than solely to laboratory or basic research findings.

\medterm{Clinical Communication} Exchange of information between healthcare professionals, patients, and caregivers to
support diagnosis, treatment, safety, and informed decisions.

\medterm{Clinical Counseling} Guidance provided in healthcare to help patients make informed decisions about health,
prevention, treatment, and self-management.

\medterm{Clinical Depression} A depressive illness causing persistent sadness or loss of interest, together with changes in sleep, appetite, energy, thinking, or self-worth that significantly interfere with daily life.

\medterm{Clinical Depression (Depression)} Alternate index label for Depression.

\medterm{Clinical Deterioration} Worsening health that increases the risk of organ failure, disability, or death. Warning signs may include new confusion, breathing difficulty, abnormal vital signs, reduced urine, worsening pain, or poor circulation.

\medterm{Clinical Disease} Disease that produces recognizable signs or symptoms or other findings detectable in a patient, as opposed to a silent or subclinical state.


\medterm{Clinical Frailty Scale} A scale used by a healthcare professional to describe how fit, independent, or frail an older person is. It considers activity, daily function, and illness burden and supports care planning.

\medterm{Clinical Guideline} Carefully developed recommendations helping healthcare professionals make decisions about prevention, diagnosis, treatment, monitoring, and follow-up for specific conditions.

\medterm{Clinical Handover} Transfer of responsibility and relevant patient information from one healthcare
professional or team to another. Effective handover communicates diagnosis, current
status, risks, pending tasks, contingency plans, and accountability.

\medterm{Clinical Practice Guideline} Evidence-informed recommendations intended to help healthcare professionals and patients make decisions about a specific
clinical condition or situation.

\medterm{Clinical Presentation} The symptoms, signs, examination findings, and test results seen in a person with a health problem. The presentation can differ with the cause, severity, age, other illnesses, and the time since the problem began.

\medterm{Clinical Reasoning} The cognitive process of collecting, weighting, and integrating clinical information to
formulate diagnostic and management decisions. It includes problem representation,
differential diagnosis, probability assessment, and recognition of cognitive bias.

\medterm{Clinical Research} Research involving people or human health data to study disease, diagnosis, prevention, treatment, or healthcare delivery.

\medterm{Clinical Research Bias} A systematic error in how a health study is designed, conducted, analyzed, or reported that can make results misleading. Examples include choosing unrepresentative participants, measuring inaccurately, or failing to account for other influences.
\medterm{Clinical Trial} A planned study in which people receive or are observed for an intervention to evaluate its safety, effects, or outcomes. Trials may
be conducted in phases and for many types of health interventions.

\medterm{Clinical Trial Phases} Stages for testing a medical treatment: Phase I mainly examines safety and dose, Phase II studies early effectiveness and safety, Phase III compares it with usual care in larger groups, and Phase IV monitors longer-term use.

\medterm{Clinical Trials} Research studies involving human participants that evaluate the safety, effectiveness, risks, or best use of medical interventions or health strategies.

\medterm{Clinically Proven} A description meaning that research in people supports a treatment or claim. The quality, number, participants, outcomes, and limitations of the studies should be considered before judging how strong the evidence is.
\medterm{Clinicopathologic Correlation} Combining a person's symptoms, examination, scans, laboratory tests, and tissue findings to reach the most accurate diagnosis, especially when a small sample may not show the whole disease.

\medterm{Clinistix} Clinistix is a reagent-strip method historically used to detect glucose in urine through an enzymatic color reaction.
\medterm{Clinitest} Clinitest is a tablet-based copper-reduction test historically used to detect reducing substances in urine, including glucose; it is less specific than modern methods.
\medterm{Clinodactyly} A finger or toe that curves sideways from birth or during development, most often the little finger. It is usually mild but may interfere with movement or occur with a genetic condition.

\medterm{Clinoid Process} A small bony projection at the base of the skull near the eye socket and major nerves and blood vessels. It helps form the support around the pituitary region and optic canal.

\medterm{CLIP} A device used to hold tissue or structures together; surgical clips can close blood vessels, mark tissue, or occlude structures such as the vas deferens.

\medterm{Clip (Vascular)} A small device placed during a procedure to close or control a blood vessel. It may stop bleeding, prevent blood flow to a selected area, or secure a vessel permanently.
\synonyms
Vascular

\medterm{Clitoral} Relating to the clitoris, an erectile organ with a dense concentration of sensory nerves at the front of the vulva.

\medterm{Clitorectomy} Surgical removal of all or part of the clitoris, sometimes with removal of adjacent vulvar tissue; nonconsensual forms constitute female genital mutilation and can cause serious harm.

\medterm{Clitoridectomy} Surgical removal of part or all of the clitoris. In medicine it may be performed for
selected tumors or severe disease; nonmedical removal is a harmful practice with no
health benefit.

\medterm{Clitoris} A small, sensitive erectile organ located at the anterior junction of the labia minora,
beneath the clitoral hood (prepuce). The clitoris is the primary source of female sexual
pleasure and is homologous to the male penis.

\medterm{Clitoromegaly} Clitoromegaly is enlargement of the clitoris, which may be congenital or acquired from androgen excess, endocrine disease, medication, or other causes.
\medterm{CLL} Alternate index label; see \textit{Chronic lymphocytic leukemia}.

\medterm{Cloacal Exstrophy} The most severe congenital anomaly of the ventral abdominal wall, also called the OEIS
complex (Omphalocele, Exstrophy of the bladder, Imperforate anus, Spinal defects).
Incidence 1 in 200,000-400,000 births.

\synonyms
the OEIS.

\medterm{Cloacitis} Inflammation of a cloaca or cloacal region, particularly in disorders affecting the
common terminal outlet of the urinary, genital, and intestinal tracts.

\medterm{Clodronate} Clodronate is a bisphosphonate that reduces bone resorption and is used in selected bone and cancer-related conditions.
\medterm{Clofazimine} Clofazimine is a riminophenazine antimicrobial used mainly in multidrug treatment of leprosy and selected mycobacterial infections; skin discoloration is common.
\medterm{Clomifene} Clomifene, or clomiphene, is a selective estrogen-receptor modulator used to induce ovulation in selected infertility settings.
\medterm{Clomipramine} Clomipramine is a tricyclic antidepressant particularly effective for obsessive-compulsive disorder; anticholinergic, cardiac, and overdose toxicity are important.
\medterm{Clonal Evolution} The selection and expansion of tumor-cell subclones with biologic or treatment-related
advantages over time. It contributes to progression, metastasis, and acquired resistance
to anticancer therapy.

\medterm{Clonal Expansion} Multiplication of cells descended from one original cell. It may occur when immune cells respond to a trigger or when an abnormal cell population grows in cancer or another disorder.
\synonyms
Clonal proliferation; clone expansion.

\medterm{Clonazepam} Clonazepam is a benzodiazepine anticonvulsant and anxiolytic that enhances GABA activity; sedation, dependence, and respiratory depression can occur.
\medterm{Clonic Seizure} A seizure characterized by rhythmic, repetitive jerking caused by recurrent involuntary
muscle contractions. It may be focal or generalized, with or without impaired awareness,
and is classified using the clinical pattern, electroencephalography, and the underlying
cause.

\medterm{Clonidine} Clonidine stimulates central alpha-2 receptors to reduce sympathetic outflow and is used for hypertension and selected other conditions; abrupt withdrawal can cause rebound hypertension.
\medterm{Clonorchiasis} Infection with the Chinese liver fluke, Clonorchis sinensis. It affects the biliary
tract and may cause inflammation, obstruction, and long-term hepatobiliary
complications.

\medterm{Clonus} Repeated, rhythmic muscle jerks triggered by suddenly stretching a muscle, often at the ankle. It commonly indicates overactive reflexes from a problem affecting the brain or spinal cord.

\medterm{Clopidogrel} Clopidogrel is an antiplatelet drug that irreversibly blocks the platelet P2Y12 receptor, reducing arterial thrombosis risk after selected cardiovascular events or procedures.
\medterm{Closed Comedo} A small skin lesion formed when a follicular opening is blocked by sebum and keratin but
remains covered by the epidermis. It is commonly called a whitehead.

\synonyms
a whitehead.

\medterm{Closed Fracture} A fracture in which the broken bone does not communicate through a wound with the
external environment.

\medterm{Closed Reduction} Nonoperative realignment of a fractured or dislocated bone or joint by manipulation
without surgically opening the site.

\medterm{Closed Reduction Of A Fractured Bone} A nonsurgical procedure that realigns a broken bone by manipulating it through the skin without opening the fracture site, usually followed by immobilization and imaging to confirm alignment.


\medterm{Closed Suction Drain With Bulb} A flexible drain connected to a compressible bulb that creates gentle negative pressure to remove blood or fluid from a surgical or injured area; the bulb is emptied and output recorded as directed.

\medterm{Closed-Angle Glaucoma} Glaucoma caused by blockage of the eye's drainage angle, so fluid cannot leave and pressure rises. Sudden attacks can cause severe eye pain, blurred vision, halos, headache, nausea, and permanent sight loss.

\medterm{Closed-Loop Insulin Delivery} Automated insulin delivery systems that combine continuous glucose monitoring with an
insulin pump and control algorithm. The algorithm adjusts basal insulin delivery based
on sensor glucose values in real time. Hybrid closed-loop systems require user-initiated
meal boluses but automatically manage basal rates.

\medterm{Clostridioides Difficile} A spore-forming bacterium that can overgrow after antibiotics disturb normal gut bacteria. Its toxins inflame the colon and cause watery diarrhea, abdominal pain, fever, and sometimes life-threatening colitis.

\medterm{Clostridioides Difficile Infection} Diarrhea and colon inflammation caused by overgrowth of Clostridioides difficile, often after antibiotic use. It may be mild or severe and can recur; diagnosis requires appropriate stool testing in a person with symptoms.

\medterm{Clostridioides Difficile Pathogenesis} The mechanisms by which C. difficile causes disease. Antibiotic disruption of the normal
gut flora allows the organism to colonize; its toxins injure the colonic mucosa and
cause inflammation and diarrhea. Spores persist in the environment and facilitate
transmission.

\medterm{Clostridium} A genus of anaerobic, usually spore-forming bacteria that includes species capable of
causing diseases such as tetanus, botulism, gas gangrene, and antibiotic-associated
colitis.

\medterm{Clostridium Botulinum} A bacterium that produces botulinum toxin, which blocks nerve signals to muscles and causes weakness or paralysis. Food, wounds, or swallowed spores can cause botulism, a medical emergency affecting breathing and swallowing.

\medterm{Clostridium Difficile} Clostridioides difficile infection: a toxin-producing anaerobic gram-positive rod
causing antibiotic-associated diarrhea and pseudomembranous colitis, most commonly after
broad-spectrum antibiotics (clindamycin, cephalosporins, fluoroquinolones).

\medterm{Clostridium Difficile Colitis} Clostridioides difficile colitis is antibiotic-associated inflammation of the colon caused by toxin-producing bacteria. It causes watery diarrhea, abdominal pain, fever, and sometimes severe colitis; diagnosis and treatment depend on symptoms, testing, recurrence, and severity.

\synonyms
Difficile Clostridium (Clostridium Difficile Colitis)

\medterm{Clostridium Difficile Infection} An anaerobic, spore-forming, toxin-producing gram-positive bacillus causing antibiotic-associated diarrhea and pseudomembranous colitis. See Clostridioides difficile.
\medterm{Clostridium Perfringens} A bacterium that can cause rapidly destructive muscle infection with gas in tissues or food poisoning with sudden diarrhea and abdominal cramps. Severe wound infection requires emergency surgery and antibiotics.

\medterm{Clostridium Tetani} A soil bacterium whose toxin causes tetanus after entering a wound. It produces painful muscle stiffness and spasms, often beginning in the jaw, and can interfere with swallowing and breathing; vaccination prevents it.

\medterm{Closure} Closure is the joining or sealing of a wound, opening, vessel, duct, or anatomical structure, either naturally or by treatment.
\medterm{Clot} A clot is a semisolid mass of coagulated blood; clots may stop bleeding or obstruct a vessel and cause thrombosis or embolism.
\medterm{Clot Blood (Blood Clots)} Alternate index label for Blood Clots.

\medterm{Clot Buster} A medicine that breaks down a blood clot, also called a thrombolytic medicine. It may be used urgently for selected strokes, heart attacks, or lung clots, but it can cause dangerous bleeding and is not safe for everyone.

\medterm{Cloth Dye Poisoning} Exposure to cloth dyes that may cause irritation, gastrointestinal symptoms, allergic reactions, or systemic toxicity depending on the chemical composition.

\medterm{Clotrimazole} Clotrimazole is an azole antifungal used topically or orally in selected formulations for fungal infections of skin, mouth, or vagina.
\medterm{Clotting} The body's process of stopping bleeding by narrowing a damaged vessel, forming a platelet plug, and producing a protein mesh that strengthens the plug. Excessive or inappropriate clotting can block blood vessels.

\medterm{Cloudy Cornea} Loss of corneal transparency from edema, scarring, infection, dystrophy, trauma, or deposits, potentially causing blurred vision, pain, photophobia, or glare.

\medterm{Clove Oil} Clove oil contains eugenol and is used in some dental and topical products; concentrated ingestion or exposure can cause irritation, liver injury, or seizures.
\medterm{Clozapine} Clozapine is an atypical antipsychotic for treatment-resistant schizophrenia and selected suicidality risk; neutropenia, myocarditis, seizures, constipation, and metabolic effects require monitoring.
\medterm{Club Drugs List And Side Effects} Club drugs include substances such as MDMA, ketamine, GHB, and others used recreationally in social settings. Effects may include agitation, hallucinations, seizures, dangerous temperature elevation, respiratory depression, cardiovascular toxicity, impaired consent, and overdose.

\medterm{Clubbing} Enlargement and rounding of the fingertips or toes with more curved nails and loss of the normal nail angle. It may accompany long-term lung, heart, bowel, liver, or other disease and should be medically assessed.

\medterm{Clubbing Of The Fingers Or Toes} Bulbous enlargement of fingertip or toe soft tissue with increased nail curvature caused by chronic changes in the nail bed; it can accompany lung, heart, gastrointestinal, or other disease.

\medterm{Clubfoot} A birth defect in which a baby's foot is turned inward and
downward. The foot may be stiff and difficult to move into a normal position.

\medterm{Clubfoot (Talipes Equinovarus)} Alternate name; see \textit{Clubfoot}.

\medterm{Clubfoot Repair} Treatment to gradually straighten a baby's turned-in foot, usually with gentle stretching and casts, followed by a brace. Some children need a small tendon procedure or later surgery if the deformity returns or is severe.

\medterm{Clubhand} A birth defect in which the hand and wrist curve toward the thumb or little-finger side because forearm bones or other tissues developed abnormally. Treatment depends on function and may include therapy or surgery.

\medterm{Clumping} The gathering of cells, particles, or proteins into groups. In a laboratory sample, clumped red blood cells or platelets can affect test results; in the body, clumping can sometimes block vessels.
\medterm{Cluster} A group of similar events, findings, cells, or cases occurring together in space, time, or clinical pattern.

\medterm{Cluster Headache} Repeated attacks of extremely severe pain around or behind one eye, lasting about 15 minutes to three hours. The same side may tear, redden, swell, or feel congested, and the person often becomes restless.

\medterm{Cluster Headaches} Cluster headache causes repeated attacks of severe one-sided orbital or temporal pain with tearing, nasal symptoms, eyelid drooping, or restlessness. Attacks occur in bouts and require targeted acute and preventive treatment; a first sudden severe headache needs urgent evaluation.

\medterm{Clutton'S Joints} Painless swelling of both knees or other large joints caused by long-term inflammation in children with congenital syphilis. Joint movement is usually preserved, but the underlying infection requires treatment.

\medterm{Clysis} An older method of giving fluid or medicine through the rectum, sometimes used for hydration or bowel treatment. It is rarely used today because safer and more effective routes are available.
\medterm{Cm (Centimeter)} The abbreviation for centimeter, a unit of length equal to one hundredth of a meter or 10 millimeters.

\medterm{CMV} Cytomegalovirus, a common herpesvirus that may cause mild or asymptomatic infection but can cause serious disease in fetuses and people with
weakened immunity.

\synonyms
Cytomegalovirus (CMV)

\medterm{CMV - Gastroenteritis/Colitis} Inflammation and ulceration of the stomach or intestine caused by cytomegalovirus, usually in people with substantial immune suppression. It may cause abdominal pain, fever, diarrhea, bleeding, or perforation and requires tissue or other targeted evaluation.

\medterm{CMV Blood Test} A blood test that looks for cytomegalovirus antibodies, viral proteins, or viral genetic material. Antibodies usually show past or recent exposure, while viral material can help detect active infection, especially in people with weakened immunity.

\medterm{CMV Colitis} Inflammation of the colon caused by cytomegalovirus, usually in someone with severely weakened immunity. It may cause bloody diarrhea, abdominal pain, fever, or weight loss and is confirmed with appropriate testing, often including biopsy.

\medterm{CMV Pneumonia} Pneumonia caused by cytomegalovirus, usually in people with substantial immune suppression; fever, cough, shortness of breath, and low oxygen may occur and require targeted evaluation and antiviral treatment.

\medterm{CMV Retinitis} Infection and inflammation of the retina caused by cytomegalovirus, usually in people with severely weakened immunity. It may cause floaters, blurred or missing areas of vision, and blindness unless treated urgently.

\medterm{CNS Infection} Infection of the brain, spinal cord, meninges, or cerebrospinal-fluid spaces caused by bacteria, viruses, fungi, parasites, or other organisms; presentation and treatment depend on the pathogen and anatomic site.

\medterm{CNS Lymphoma} A lymphoma involving the central nervous system, most often diffuse large B-cell lymphoma. It may involve brain parenchyma, meninges, eyes, or cerebrospinal fluid and requires urgent hematologic and neurologic evaluation.

\medterm{CO Poisoning} Toxicity from inhalation of carbon monoxide, a colorless, odorless gas binding
hemoglobin ~200-300 times more avidly than oxygen, forming carboxyhemoglobin (COHb) and
reducing oxygen delivery. See carbon monoxide poisoning.

\medterm{CO2 Blood Test} A blood chemistry measurement that usually estimates serum total carbon dioxide, which largely reflects bicarbonate. It helps assess acid-base and electrolyte disorders but must be interpreted with sodium, chloride, kidney function, and, when needed, blood-gas results.

\medterm{Coagulate} To change from a liquid to a semisolid or solid state, especially the formation of a blood clot through coagulation.

\synonyms
Clot; clot formation.

\medterm{Coagulation} The body's process for stopping bleeding by forming a stable blood clot. Platelets first gather at an injured vessel, then blood proteins create a mesh of fibrin that strengthens the plug; too little clotting causes bleeding, while too much can cause harmful clots.
\medterm{Coagulation Cascade} The coagulation cascade is a series of enzymatic reactions leading to fibrin clot
formation, involving serine protease activation cascades. The intrinsic pathway is
initiated by contact activation when factor XII binds exposed collagen, sequentially
activating factors XI, IX, and the cofactor VIII.

\medterm{Coagulation Defects} Coagulation defects are abnormalities of clotting proteins, platelets, vessels, or regulation that cause excessive bleeding or inappropriate thrombosis.
\medterm{Coagulation Factor} A blood protein required for the coagulation cascade and formation of a stable fibrin clot; deficiencies or inhibitors can cause bleeding, while excessive activation can contribute to thrombosis.

\medterm{Coagulation Factor Deficiencies} Inherited or acquired disorders in which one or more clotting factors are reduced or
function abnormally, causing excessive bleeding. The clinical pattern and laboratory
findings depend on the affected factor.

\medterm{Coagulation Profile} Panel of tests assessing the coagulation cascade. PT (prothrombin time, 11-14 sec):
measures extrinsic and common pathways; elevated in vitamin K deficiency, warfarin,
liver disease, DIC. INR standardizes PT for warfarin monitoring.

\medterm{Coagulation Time} Coagulation time is the time required for a blood sample to form a clot under specified test conditions; modern tests usually measure particular pathways more precisely.
\medterm{Coagulative Necrosis} The most common form of necrosis, typically resulting from ischemic injury or infarction
in most solid organs except the brain. The fundamental process involves denaturation of
both structural proteins and enzymatic proteins, which prevents the immediate
liquefaction of the dead tissue.

\medterm{Coagulopathy} A disorder of hemostasis that impairs normal clot formation or regulation and
predisposes to abnormal bleeding or thrombosis. It may result from coagulation-factor
deficiency, platelet dysfunction, anticoagulant exposure, liver disease, consumptive
coagulopathy, or systemic illness.

\medterm{Coal Worker'S Pneumoconiosis} Pneumoconiosis from coal dust. Simple CWP shows upper lobe macules. Complicated CWP
(progressive massive fibrosis) involves coalescent macules. May be associated with
Caplan syndrome.

\medterm{Coarctation} A localized narrowing of a tubular anatomical structure, especially coarctation of the
aorta, which obstructs blood flow and produces characteristic pressure differences.

\medterm{Coarctation Of The Aorta} A congenital narrowing of the aorta, typically near the ductus arteriosus. Presents with
upper extremity hypertension, diminished femoral pulses, and a blood pressure gradient
between arms and legs. Associated with bicuspid aortic valve and other congenital heart
defects.

\medterm{Coarse Crackles} Lower-pitched, longer-duration crackles heard throughout inspiration, caused by fluid or
secretions in larger airways. They are heard in COPD, bronchiectasis, pneumonia, and
pulmonary edema. They are less specific than fine crackles.

\medterm{Coats Disease} An uncommon retinal disorder characterized by abnormal retinal blood vessels, leakage,
exudation, and sometimes retinal detachment, usually affecting one eye in childhood.

\medterm{Coats Disease (Coats' Disease)} Alternate index label for Coats' Disease.

\medterm{Cobalamin} Vitamin B12, a nutrient needed for DNA production, healthy red blood cells, nerve function, and normal processing of certain fats and proteins.

\medterm{Cobalamin Deficiency} Too little vitamin B12, which is needed for healthy blood cells and normal nerve function. It may cause anemia, fatigue, sore tongue, numbness, balance problems, memory changes, or irreversible nerve injury if untreated.
\medterm{Cobalt} Cobalt is a metal present in vitamin B12 and some industrial materials; excess exposure can injure the heart, thyroid, lungs, or skin.
\medterm{Cobalt Poisoning} Toxic exposure to cobalt that may cause gastrointestinal, neurologic, thyroid, cardiac, or lung effects after significant or prolonged exposure.

\medterm{Cobalt-57} A radioactive form of cobalt used as a test source for checking and calibrating gamma cameras. It is mainly used for equipment quality control and is not normally given to patients as a diagnostic or treatment medicine.

\medterm{Cobalt-Chromium} A strong, corrosion-resistant metal mixture used for cast dental restorations, partial dentures, crowns, and selected medical devices. Some people may develop sensitivity, and prolonged exposure to cobalt can be harmful.

\medterm{Cocaine} A stimulant drug that blocks monoamine reuptake and can also act as a local anesthetic; use can cause dependence, cardiovascular complications, seizures, stroke, and overdose.

\medterm{Cocaine Use Disorder} A problematic pattern of cocaine use leading to clinically significant impairment or distress.
It may involve craving, impaired control, continued use despite harm, tolerance, or withdrawal.

\synonyms
Cocaine And Crack Abuse

\medterm{Cocaine Cardiotoxicity} Cocaine blocks norepinephrine and dopamine reuptake, causing sympathomimetic effects:
tachycardia, hypertension, and coronary vasospasm. Acute myocardial infarction can occur
even with normal coronary arteries due to vasospasm, accelerated atherosclerosis, and
prothrombotic effects. Treatment: benzodiazepines for sympatholysis.

\medterm{Cocaine Intoxication} A toxic reaction to cocaine that may cause agitation, overheating, high blood pressure, seizures, abnormal rhythms, heart attack, stroke, or sudden death.

\medterm{Cocaine Toxicity} A toxic state caused by cocaine, often producing agitation, high blood pressure, rapid heart
rate, hyperthermia, seizures, coronary vasospasm, myocardial infarction, or stroke. Treatment is supportive and may
include sedation, temperature control, and management of cardiovascular or neurologic complications.

\medterm{Cocaine Withdrawal} Physical and emotional symptoms after stopping or reducing cocaine, including fatigue, depressed mood, increased appetite, sleep changes, irritability, and cravings. Severe depression or suicidal thoughts require urgent help.

\medterm{Cocarboxylase} Cocarboxylase is thiamine pyrophosphate, the active coenzyme form of vitamin B1 required for several oxidative carbohydrate and amino-acid reactions.

\medterm{Coccidioides Complement Fixation} A blood test that detects complement-fixing antibodies to Coccidioides; antibody levels can support diagnosis and help assess the severity or progression of coccidioidomycosis.

\medterm{Coccidioides Immitis} A dimorphic soil fungus and a principal cause of coccidioidomycosis, or Valley fever, acquired by inhaling airborne arthroconidia. Infection is usually pulmonary but can disseminate, especially in people with impaired cellular immunity.

\medterm{Coccidioides Precipitin Test} A serologic test that detects antibodies to Coccidioides, typically early immunoglobulin M antibodies, and supports evaluation of suspected coccidioidomycosis when interpreted with symptoms and other testing.

\medterm{Coccidioidomycosis} A fungal infection caused by Coccidioides fungi, usually acquired by breathing spores from disturbed soil. It most often affects the lungs and causes cough, fever, tiredness, or chest pain; rarely, it spreads to the skin, bones, joints, or brain.
\commonname{Valley fever}

\medterm{Coccidiosis} Infection by coccidian protozoa. In humans it may cause intestinal disease or, in
immunocompromised patients, infection of other organs.

\medterm{Coccus} A round or nearly round bacterium. These bacteria may occur alone, in pairs, chains, or clusters; their shape helps describe them but does not identify the species or illness by itself.
\medterm{Coccydynia} Pain in the region of the coccyx, often worsened by sitting. Causes include trauma,
childbirth, degenerative change, and less commonly infection or tumor.

\synonyms
Pain Tailbone (Coccydynia)

\medterm{Coccygectomy} Surgical removal of the coccyx, reserved for selected cases of severe persistent
coccygeal pain or a lesion that has not responded to conservative treatment.

\medterm{Coccyx} The small pointed bone at the bottom of the spine, commonly called the tailbone. It connects to the sacrum, supports pelvic-floor tissues, and may become painful after a fall, childbirth, or prolonged sitting.

\medterm{Cochlea} The spiral-shaped part of the inner ear that changes sound vibrations into nerve signals. Fluid movement inside it stimulates sensory hair cells, allowing the brain to recognize pitch and loudness.

\medterm{Cochlear Implant} An electronic device for severe hearing loss that bypasses damaged inner-ear hair cells and stimulates the hearing nerve. It can provide useful sound perception but does not restore normal hearing and requires training.

\medterm{Cockayne Syndrome} A rare inherited disorder of DNA repair characterized by progressive neurologic
dysfunction, growth failure, developmental impairment, photosensitivity, and a
prematurely aged appearance.

\medterm{Code Black} A hospital emergency alert whose meaning varies by institution, often indicating a severe security threat, bomb threat, or other critical incident; local policy determines the response.

\medterm{Code Blue} A hospital emergency signal indicating that a person may be in cardiopulmonary arrest
and requires immediate resuscitation response.

\medterm{Code, Hospital} A standardized hospital announcement used to summon a specialized response to an emergency; the meaning of each code depends on the hospital's policy.

\medterm{Codeine Overdose} Overdose of codeine, an opioid, that can cause respiratory depression, impaired consciousness, pinpoint pupils, and coma. Naloxone may be required in emergency care.

\medterm{Codman Triangle} A triangular area of new bone seen on an X-ray when a rapidly growing bone lesion lifts the outer covering of bone. It suggests aggressive bone disease but does not identify a specific diagnosis.

\medterm{Codon} A group of three genetic building blocks in messenger RNA that specifies one amino acid or signals the end of protein production. Several different codons can specify the same amino acid.

\medterm{Coenzyme} A non-protein organic molecule that binds to an enzyme (apoenzyme) and is required for
its catalytic activity. Coenzymes often act as transient carriers of specific chemical
groups or electrons.

\medterm{Coenzyme Q} A fat-soluble molecule in mitochondrial membranes that transfers electrons during cellular energy production. It also acts as an antioxidant, and the body makes it naturally, although levels may change with age or illness.

\medterm{Coenzyme Q10} CoQ10, a lipid-soluble compound in mitochondrial electron transport and an antioxidant; it is also sold as a dietary supplement, although benefits vary by condition.

\medterm{Cofactor} A nonprotein helper, such as a mineral ion or vitamin-derived molecule, that an enzyme needs to perform a chemical reaction. Without its cofactor, the enzyme may not work normally.

\medterm{Coffin-Lowry Syndrome} A rare inherited developmental disorder, usually linked to an RPS6KA3 gene change. It can cause intellectual disability, distinctive facial features, short stature, spinal curvature, tapered fingers, and sudden falls triggered by surprise.

\synonyms
CLS

\medterm{Cogan’S Syndrome} A rare immune disorder causing inflammation of the cornea and inner ear. It may produce eye pain, reduced vision, dizziness, ringing, or hearing loss and can sometimes inflame large blood vessels.

\medterm{Cognition} The mental processes involved in acquiring, storing, retrieving, and using information, including attention, memory, language, reasoning, and judgment.

\medterm{Cognitive} Relating to mental processes such as attention, memory, learning, language, reasoning, perception, judgment, problem solving, and decision-making.

\medterm{Cognitive Behavior Therapy} A talking treatment that helps a person notice unhelpful thoughts and behaviors and develop more useful ways to respond. It is used for problems such as anxiety, depression, pain, insomnia, and some eating disorders.

\medterm{Cognitive Behavioral Therapy} A structured talking therapy that helps a person identify unhelpful thoughts and behaviors and practice more useful responses. It is used for depression, anxiety, insomnia, pain, and other conditions.

\medterm{Cognitive Behavioral Therapy For Back Pain} A structured psychological treatment that addresses pain-related thoughts, behaviors, sleep, activity, and coping to reduce disability and improve function in chronic back pain.

\medterm{Cognitive Behavioral Therapy For Insomnia} Cognitive behavioral therapy for insomnia is a structured, evidence-based program that
targets the thoughts and behaviors perpetuating poor sleep. Components include sleep
restriction, stimulus control, cognitive restructuring, and sleep hygiene education.

\medterm{Cognitive Decline} A reduction over time in one or more cognitive abilities, such as memory, attention, language, reasoning, or visuospatial function.

\medterm{Cognitive Distortions} Systematic errors in thinking that perpetuate negative emotional states and maladaptive
behaviors. Common distortions include all-or-nothing thinking, catastrophizing, mind
reading, overgeneralization, emotional reasoning, and should statements. CBT focuses on
identifying and modifying these patterns.

\medterm{Cognitive Fitness} The active maintenance of mental sharpness and brain health through stimulating
activities, physical exercise, social engagement, and proper nutrition. Regular
challenges such as learning new skills, puzzles, reading, and social interaction support
neuroplasticity and cognitive reserve.

\medterm{Cognitive Function} The mental abilities used for attention, memory, language, learning, reasoning, judgment, planning, problem solving, and understanding information.

\synonyms
Cognition; cognitive ability.

\medterm{Cognitive Health And Memory} Cognitive health encompasses the ability to think clearly, learn, remember, and make
decisions, with memory being a core component that involves encoding, storing, and
retrieving information. Age-related changes in processing speed and recall are normal,
but persistent decline may signal mild cognitive impairment or dementia.

\medterm{Cognitive Impairment} A decline in thinking abilities such as memory, attention, language, judgment, or problem solving. It may be mild with independence preserved or severe enough to interfere with everyday activities.

\medterm{Cognitive Rehabilitation} Structured therapy to improve cognitive function or help a person compensate for difficulties
with memory, attention, language, or executive skills after illness or injury.

\medterm{Cognitive Reserve} The brain's ability to use alternative strategies or networks to maintain function despite injury or disease.

\medterm{Cognitive Screening} A brief assessment used to identify possible problems with memory, attention, language, or other cognitive functions. It does not by
itself diagnose dementia or another disorder.

\medterm{Cognitive Therapy} A psychological treatment that helps a person identify and change unhelpful thoughts and beliefs and develop more adaptive behaviors; cognitive behavioral therapy combines it with behavioral methods.

\medterm{Cohort} A group of people who share a defined characteristic or experience during a specified period and are studied over time or compared with another group.

\medterm{Cohort Study} A research study following groups with different exposures over time to compare new disease rates, outcomes, or risks.

\medterm{Coil (Vascular)} A catheter-delivered wire device used to block selected blood vessels, often to treat aneurysms or bleeding.

\synonyms
Vascular

\medterm{Coincidence Detection} Principle of PET imaging. Two 511 keV photons detected simultaneously (within
nanoseconds) by opposing detectors. Line of response defined. Time-of-flight PET
improves localization.

\medterm{Coitus} Sexual intercourse, traditionally meaning penile-vaginal penetration; the term does not include every form of sexual activity or intimate contact.

\medterm{Coitus Interruptus} The withdrawal method of contraception, in which the penis is removed from the vagina before ejaculation; it is less reliable than many other contraceptive methods.

\medterm{Coke} Slang for cocaine, a stimulant drug; the word can also mean a carbon fuel product, so context is essential.

\medterm{Cold (Common Cold)} Alternate index label for Common Cold.

\medterm{Cold Abscess} A pocket of pus that develops slowly with little redness, warmth, or pain. It may occur with tuberculosis or other long-term infections and still requires medical evaluation and treatment.

\medterm{Cold Agglutinin} An antibody that causes red blood cells to agglutinate at relatively low temperatures.
Pathologic cold agglutinins can cause autoimmune hemolytic anemia.

\medterm{Cold Agglutinin Disease} An autoimmune blood disorder in which antibodies become active in cooler parts of the body and cause red blood cells to stick together and break down. It may cause anemia, fatigue, jaundice, or cold-related color changes.

\medterm{Cold Agglutinin Syndrome} An immune-mediated disorder in which antibodies react with red blood cells at low temperatures,
potentially causing hemolytic anemia or circulation-related symptoms.

\medterm{Cold Agglutinins} IgM autoantibodies that bind red blood cell antigens at temperatures below 37 degrees
Celsius, causing agglutination and potentially complement-mediated hemolysis in cooler
peripheral circulation.

\medterm{Cold Allergy} Alternate index label; see \textit{Cold urticaria}.

\medterm{Cold Exposure} Alternate index label; see \textit{Hypothermia}.

\medterm{Cold Flu Allergy Treatments (Cold, Flu, Allergy)} Alternate index label for Cold, Flu, Allergy.

\medterm{Cold Hands} Hands that feel unusually cold, which may result from environmental exposure, reduced blood flow, Raynaud phenomenon, anemia, thyroid disease, or other medical conditions.

\medterm{Cold Injury} Harm caused by low temperature, including hypothermia, frostbite, chilblains, and nonfreezing tissue injury. Risk increases with prolonged exposure, wetness, wind, or poor circulation and may affect the whole body or a local area.

\medterm{Cold Intolerance} Unusual sensitivity or inability to tolerate cold temperatures, which may occur with hypothyroidism, anemia, low body weight, vascular disease, or other conditions.

\medterm{Cold Medicines And Children} Safety guidance for giving medicines to children with cold symptoms; many combination products provide little benefit in young children and can cause dosing errors or serious adverse effects.

\medterm{Cold Sore} A contagious infection of the lips or skin around the mouth, usually caused by herpes simplex virus type 1 and occasionally by HSV-2. It often begins with tingling or burning, followed by grouped fluid-filled blisters that crust and heal in about 1--2 weeks. The virus remains in the body and can reactivate, sometimes after illness, stress, fever, or sunlight exposure.


\medterm{Cold Urticaria} Raised, itchy patches or swelling that appear after the skin is exposed to cold. A large area of cold exposure, such as swimming in cold water, can rarely cause fainting or a life-threatening allergic reaction.

\synonyms
Cold Allergy

\medterm{Cold Wave Lotion Poisoning} Harmful exposure to chemicals in a permanent-wave hair product. Swallowing or contact may burn the mouth, skin, eyes, or airways and can cause vomiting or breathing problems; urgent poison-control advice is needed.
\medterm{Cold, Common} Alternate index label; see \textit{Common cold}.

\medterm{Cold, Flu, Allergy} Colds, influenza, and allergies can all cause nasal symptoms and cough but differ in cause, fever pattern, duration, and contagiousness. Viral illness usually improves with supportive care, while allergies respond to trigger reduction and appropriate antihistamine or nasal treatment.

\synonyms
Cold Flu Allergy Treatments (Cold, Flu, Allergy)

\medterm{Colds And Flu - Antibiotics} An explanation that antibiotics treat susceptible bacterial infections, not uncomplicated viral colds or influenza; unnecessary use causes adverse effects and antimicrobial resistance.



\medterm{Colectomy} Surgery to remove part or all of the colon. It may treat cancer, severe inflammatory bowel disease, diverticular complications, blocked bowel, or inherited high cancer risk; the amount removed depends on the disease.

\medterm{Colibacillosis} Disease caused by pathogenic strains of Escherichia coli, which may produce intestinal
infection, urinary-tract infection, sepsis, or other illness.

\medterm{Colic} Severe wave-like pain caused by contraction or blockage of a hollow organ, such as the intestine, gallbladder, or urinary tract. The term also describes repeated crying in an otherwise healthy young infant.

\medterm{Colicky Pain} Intermittent, wave-like cramping pain caused by contraction, distension, or obstruction of a hollow organ; common examples include biliary, renal, and intestinal colic.


\medterm{Coliform} A group of bacteria used as an indicator of possible contamination in water or food. Finding them does not prove that disease-causing bacteria are present, because many coliforms are harmless environmental organisms.

\medterm{Colitis} Inflammation of the colon, which may cause diarrhea, abdominal pain, blood or mucus in stool, fever, or urgency. Causes include infection, inflammatory bowel disease, reduced blood flow, medicines, and radiation.

\medterm{Colitis Collagenous (Lymphocytic Colitis)} Alternate index label for Lymphocytic Colitis.


\medterm{Colitis Lymphocytic (Lymphocytic Colitis)} Alternate index label for Lymphocytic Colitis.

\medterm{Colitis Microscopic (Lymphocytic Colitis)} Alternate index label for Lymphocytic Colitis.

\medterm{Colitis Ulcerative (Ulcerative Colitis)} Alternate index label for Ulcerative Colitis.

\medterm{Colitis, Ischemic} Alternate index label; see \textit{Ischemic colitis}.

\medterm{Colitis, Microscopic} Alternate index label; see \textit{Microscopic colitis}.

\medterm{Colitis, Pseudomembranous} Alternate index label; see \textit{Pseudomembranous colitis}.

\medterm{Colitis, Ulcerative} Alternate index label; see \textit{Ulcerative colitis}.

\medterm{Collagen} A strong protein that forms much of the body's connective tissue. Different types support and strengthen skin, bone, cartilage, tendons, ligaments, blood vessels, and many internal structures.

\medterm{Collagen Disease} A systemic disorder involving connective tissue and often immune-mediated injury to
collagen-rich structures. The term has historically included diseases such as systemic
lupus erythematosus and systemic sclerosis.

\medterm{Collagen Implant} A medical material made from or based on collagen, a structural protein, used to support tissue repair or add soft-tissue volume. Sources may be animal-derived or manufactured. Its uses, durability, and risks depend on the product and the treated body area.

\medterm{Collagen Vascular Disease} An older term for systemic connective-tissue or autoimmune diseases such as lupus, systemic sclerosis, and inflammatory myopathies; modern diagnosis uses the specific disease name.

\medterm{Collagenase} An enzyme that breaks down collagen, a major structural protein in skin, bone, cartilage, and connective tissue. Natural collagenases help remodel tissue, while medicines or bacterial enzymes can be used in selected treatments or cause tissue damage.

\synonyms
Collagen-digesting enzyme; collagen-degrading enzyme.

\medterm{Collagenous Colitis} A type of microscopic colitis in which a thick layer of collagen develops beneath the colon lining. It usually causes persistent, watery, nonbloody diarrhea and may require colon biopsy for diagnosis.

\medterm{Collapsed Arch} Alternate index label; see \textit{Flatfeet}.

\medterm{Collapsed Lung} A common description of pneumothorax, in which air enters the pleural space and
partially or completely collapses the lung. It may cause sudden chest pain and shortness
of breath and can be life-threatening when tension develops. See also Pneumothorax.

\medterm{Collapsed Lung (Pneumothorax)} Air trapped between the lung and chest wall, causing part or all of the lung to collapse. It may cause sudden chest pain and breathlessness; pressure that shifts the heart and other structures is immediately life-threatening.

\medterm{Collateral} Providing an alternative or parallel route, such as collateral blood vessels that supply tissue around an obstruction; in other contexts it means accompanying or secondary.

\medterm{Collateral Circulation} Alternative vascular pathways that develop when coronary arteries become chronically
narrowed. Collaterals connect different coronary artery territories and can partially
compensate for stenosis. They develop gradually in response to ischemia and are more
prominent in patients with chronic coronary artery disease.


\medterm{Collateral Ligament Injury} MCL (medial) or LCL (lateral) knee ligament injury. MCL more common, often from valgus
stress. Treatment: braces, physical therapy, surgical repair if severe.

\medterm{Collecting Duct} Final segment of the nephron where fine-tuning of urine concentration and composition
occurs. Principal cells respond to antidiuretic hormone by inserting aquaporin-2 water
channels into the luminal membrane, enabling water reabsorption and urine concentration.

\medterm{College Students And The Flu} Influenza prevention and care for college students, emphasizing vaccination, hand and respiratory hygiene, staying home when ill, and prompt assessment for severe symptoms or high-risk conditions.

\medterm{Colles Fracture} A break near the wrist end of the larger forearm bone, with the broken hand-side fragment displaced backward. It usually follows a fall onto an outstretched hand and causes wrist pain, swelling, and deformity.


\medterm{Collimator} A device in an imaging camera that allows radiation traveling in selected directions to reach the detector. Its design affects the balance between image detail and detection sensitivity.

\medterm{Colloids} Fluids containing large particles that tend to remain in the bloodstream longer than simple salt solutions. Some are used for selected fluid replacement, but benefits and risks depend on the product and the patient's condition.

\medterm{Coloboma} A birth defect caused by incomplete formation of part of the eye, leaving a notch or gap in the iris, retina, choroid, or optic nerve. Vision problems vary with the structure and amount affected.

\medterm{Coloboma Of The Iris} A congenital keyhole- or notch-shaped defect in the iris caused by incomplete embryonic closure, sometimes associated with colobomas in the retina, optic nerve, or other organs.

\medterm{Colocolonic Intussusception} Telescoping of one part of the colon into the next section. It is uncommon in adults and may be caused by a polyp, tumor, or other lesion, producing intermittent pain, bloating, vomiting, or altered bowel movements.

\medterm{Cologne Poisoning} Ingestion of cologne containing alcohol and fragrances that may cause intoxication, hypoglycemia, respiratory depression, or gastrointestinal irritation.

\medterm{Cologuard} A stool-based colorectal-cancer screening test combining DNA-marker analysis with a fecal blood test; a positive result requires diagnostic colonoscopy.

\medterm{Colon} The main part of the large intestine between the small intestine and rectum. It absorbs water and salts and stores and moves stool through its ascending, transverse, descending, and sigmoid sections.

\medterm{Colon Cancer} Cancer arising in the lining of the colon, often developing gradually from a precancerous polyp. It may cause blood in stool, changed bowel habits, abdominal pain, or anemia, but early cancer may cause no symptoms.

\synonyms
Cancer, Colon
Colorectal Cancer (Colon Cancer)

\medterm{Colon Cancer Prevention} Colorectal-cancer risk can be reduced through recommended screening, physical activity, healthy weight, limiting alcohol, avoiding tobacco, and a diet rich in fiber-containing plant foods. Screening can detect precancerous polyps and early cancer before symptoms develop.

\medterm{Colon Cancer Screening} Testing intended to detect colorectal cancer or precancerous polyps before symptoms develop. Options include stool-based tests and direct visualization such as colonoscopy; the appropriate method and interval depend on age, risk, prior findings, and local guidelines.

\medterm{Colon Polyps} Growths arising from the inner lining of the colon; some types can develop into colorectal cancer.

\synonyms
Polyps, Colon
Polyps Colon (Colon Polyps)
Polyps Rectal (Colon Polyps)

\medterm{Colonic Diverticulitis (Diverticulosis)} Alternate index label for Diverticulosis.

\medterm{Colonic Inertia} Abnormally slow movement of stool through the colon, causing persistent constipation, bloating, abdominal discomfort, and infrequent bowel movements. It is diagnosed after other causes of constipation have been considered.

\medterm{Colonic Ischemia} Reduced blood flow to the colon causing tissue injury. It may produce sudden abdominal pain, bloody diarrhea, and tenderness and can result from low blood pressure, narrowed arteries, clots, medicines, or inflammation of blood vessels.

\medterm{Colonic Stenosis} Narrowing of part of the large intestine that makes it harder for stool and gas to pass. Causes include scar tissue, inflammation, cancer, or reduced blood supply; symptoms may include pain, bloating, constipation, vomiting, or bowel blockage.

\medterm{Colonization} Presence and multiplication of microorganisms on or in a body site without tissue
invasion or clinical disease. Colonization may precede infection but does not by itself
establish that symptoms are caused by the organism.

\medterm{Colonoscopy} An examination of the rectum and colon using a flexible camera. It can find inflammation, bleeding, tumors, and polyps and allows tissue sampling or removal of some polyps during the same procedure.


\medterm{Colony-Forming Unit} A laboratory estimate of the number of living microorganisms in a sample, based on the colonies that grow in culture. One colony may come from one organism or from a group of organisms.

\medterm{Color Blindness} Reduced or absent ability to distinguish certain colors under normal lighting. It is often
inherited, especially through X-linked variants, but may also result from eye disease, neurologic disease, aging, or
medicines.

\synonyms
Deficient Color Vision
Poor Color Vision

\medterm{Color Doppler} An ultrasound technique that uses colors to show the direction and relative speed of blood flow over a regular gray-scale image. It helps assess vessels, heart valves, and circulation.

\medterm{Color Vision Deficiency} Reduced ability to distinguish some colors, commonly red and green. It is often inherited because of differences in the eye’s color-sensing cells, but it can also develop from eye disease, optic-nerve or brain disorders, medicines, or toxins. The type and severity vary.

\medterm{Color Vision Test} An eye test that assesses the ability to distinguish colors or identify color patterns. It helps detect inherited color-vision deficiency and acquired abnormalities caused by eye, optic-nerve, brain, medicine, or toxic disorders.

\medterm{Colorado Tick Fever} An acute viral infection transmitted by ticks and characterized by fever, headache,
myalgia, and sometimes leukopenia or rash.

\medterm{Colorectal Adenocarcinoma} The most common cancer of the colon and rectum, arising from gland-forming cells in the bowel lining. It may cause bleeding, anemia, altered bowel habits, pain, or no early symptoms and can spread to other organs.

\medterm{Colorectal Adenoma} A usually noncancerous gland-forming growth in the colon or rectum that may gradually develop into colorectal cancer.

\medterm{Colorectal Cancer} A malignant tumor of the colon or rectum. It may cause bleeding, altered bowel habits, anemia, abdominal pain, or no early symptoms;
evaluation and treatment depend on the tumor's location and stage.

\medterm{Colorectal Cancer (Colon Cancer)} Alternate index label for Colon Cancer.


\medterm{Colorectal Cancer Genetics} The study of inherited variants that increase the risk of colorectal cancer.
Important inherited syndromes include Lynch syndrome, involving DNA mismatch-repair genes, and familial adenomatous
polyposis, involving the APC gene.

\medterm{Colorectal Cancer Screening} Public health strategy to detect colorectal cancer and precancerous polyps in
asymptomatic individuals. Current guidelines recommend screening beginning at age
forty-five.

\medterm{Colorectal Cancer Screening Tests} Tests used to detect colorectal cancer or precancerous polyps in people
without symptoms. They include stool-based tests and examinations that visualize the colon, such as colonoscopy, CT
colonography, and flexible sigmoidoscopy. The appropriate test and interval depend on age, risk, and local guidance.

\medterm{Colorectal Polyps} Growths arising from the lining of the colon or rectum. They may be non-neoplastic, adenomatous, or serrated; some can progress to colorectal cancer, so removal and microscopic examination guide surveillance.

\medterm{Colorectal Surgery} Surgical treatment of diseases of the colon, rectum, and anus. Procedures may include
segmental bowel resection, restorative reconstruction, stoma formation, or local
procedures, depending on the disease and its extent.

\medterm{Colostomy} A surgically created opening that brings part of the colon through the abdominal wall to divert stool into an external appliance.

\medterm{Colostrum} The thick, yellowish fluid produced by the mammary glands in late pregnancy and the
first few days postpartum. Rich in immunoglobulins (especially IgA), proteins,
leukocytes, and growth factors.

\medterm{Colpitis} Inflammation of the vagina caused by infection, altered vaginal bacteria or yeast, irritation, allergy, or hormonal changes. Symptoms may include discharge, itching, burning, pain, odor, or painful intercourse.
\medterm{Colporrhaphy} A surgical repair of vaginal wall defects performed to correct pelvic organ prolapse. Anterior colporrhaphy plicates and reinforces the attenuated pubocervical fascia to treat a cystocele (bladder prolapse), while posterior colporrhaphy plicates the rectovaginal fascia and perineal body to repair a rectocele (rectal prolapse) and restore normal vaginal caliber.

\medterm{Colposcope} A lighted magnifying instrument used to examine the cervix, vagina, or vulva closely and guide biopsy of suspicious areas.

\medterm{Colposcopy} Magnified visualization of the cervix, vagina, and vulva after application of acetic
acid. Indication: abnormal Pap smear (ASCUS with high-risk HPV, LSIL, HSIL, AGC). Acetic
acid causes acetowhite lesions in dysplastic epithelium. Lugol iodine (Schiller test):
normal glycogenated epithelium stains brown; abnormal epithelium does not.

\medterm{Colposcopy - Directed Biopsy} Removal of a small tissue sample from an unusual area of the cervix, vagina, or vulva during a magnified examination. Laboratory testing checks for infection, precancerous changes, or cancer.

\medterm{Columnar Epithelium Lined Lower Esophagus} Alternate index label; see \textit{Barrett's esophagus}.

\medterm{Coma} A state of profound unresponsiveness in which a person cannot be awakened and does not show purposeful
responses. It is a medical emergency with many possible causes.

\synonyms
Severe Brain Injury, Coma

\medterm{Coma Myxedema (Myxedema Coma)} Alternate index label for Myxedema Coma.

\medterm{Coma Scale} A standardized tool for describing level of consciousness using eye opening, verbal response, and motor
response.

\medterm{Combination Inhalers} Inhalers containing two or more medicines, such as an inhaled corticosteroid with a long-acting bronchodilator. They may simplify regular treatment for asthma or chronic obstructive pulmonary disease; the prescribed device and schedule must be followed.

\medterm{Combination Therapy} Use of two or more treatments together to address a disease or mechanism. The benefits and risks depend on the condition and
the specific combination.

\medterm{Combined Hormone Therapy} Hormone therapy containing both estrogen and a progestogen, generally used when endometrial protection is needed in a person with a uterus.

\synonyms
Estrogen-progestogen therapy; combined menopausal hormone therapy.

\medterm{Combined Hyperlipidemia} A blood-fat disorder with high cholesterol and triglycerides, often increasing harmful lipoproteins and raising the long-term risk of artery disease.

\synonyms
Mixed hyperlipidemia; combined hyperlipoproteinemia.

\medterm{Combined Spinal-Epidural Analgesia} A form of regional pain relief that gives medicine into the spinal fluid for rapid relief and places a catheter outside it for continued dosing. It is often used during labor or surgery and requires monitoring.

\medterm{Comedocarcinoma} A form of ductal carcinoma in situ of the breast characterized by central necrosis
within malignant ducts, which may produce a comedo-like appearance.

\medterm{Comedones} Plugs of sebum and keratin within hair follicles, appearing as open blackheads or closed whiteheads and representing a primary lesion of acne.

\medterm{Comfort Measures Only} A care plan focused on relief of pain, breathlessness, distress, and other symptoms rather than treatments intended to prolong life.

\medterm{Commando Procedure} Extensive heart surgery for severe infection destroying both the aortic and mitral valve areas and the tissue between them. The surgeon removes infected tissue, rebuilds the damaged area, and replaces both valves when necessary.

\synonyms
Aortomitral Fibrous Body Reconstruction; Intervalvular Fibrosa Reconstruction

\medterm{Commensal} Describing an organism that lives with another organism and benefits from the relationship without normally harming or helping the host; normal skin and intestinal flora may be commensal.

\medterm{Comminuted} Describing a fracture in which one bone breaks into several pieces. These fractures may be caused by high-energy injury and can be unstable or involve nearby tissues, so treatment depends on the bone, alignment, skin, nerves, and blood vessels.

\medterm{Comminuted Fracture} A broken bone divided into several pieces, usually by a high-energy injury. It may be unstable and damage nearby skin, muscles, nerves, or blood vessels, so treatment depends on the site and severity.

\medterm{Commissurotomy} Surgical incision or separation of a commissure, especially a fused cardiac valve
commissure, to relieve obstruction.

\medterm{Commode} A portable toilet or chair placed near a bed for a person who cannot easily reach a bathroom.

\medterm{Common Bile Duct} The tube formed when ducts from the liver and gallbladder join. It carries bile into the first part of the small intestine, where bile helps digest dietary fat.

\medterm{Common Cold} A short-term viral infection of the nose and throat causing runny or blocked nose, sneezing, sore throat, cough, and sometimes mild fever. Symptoms usually improve within one to two weeks.

\synonyms
Cold, Common
Cold (Common Cold)

\medterm{Common Cold In Babies} A viral infection of an infant's upper respiratory tract that commonly causes a runny nose, congestion, and cough.

\medterm{Common Cold Stages And Timeline Of Symptoms} The common cold often begins with sore throat, sneezing, or fatigue, followed by nasal congestion, runny nose, and cough. Symptoms usually peak within several days and improve over one to two weeks, although cough and fatigue may last longer.

\medterm{Common Early Symptoms Of Myasthenia Gravis} Early myasthenia gravis commonly causes fluctuating, fatigable weakness of the eyelids, eye movements, facial muscles, chewing, swallowing, speech, neck, or limbs. Symptoms worsen with use and improve with rest; breathing or swallowing difficulty is an emergency.

\medterm{Common Iliac Arteries} Terminal abdominal aortic branches at L4, each approximately 4 cm long, dividing into
internal and external iliac arteries. They supply the pelvis and lower extremities. The
right common iliac artery crosses anterior to the right common iliac vein.

\medterm{Common Medical Abbreviations And Terms} Medical abbreviations and terms summarize diagnoses, tests, medicines, anatomy, and treatment instructions. Because abbreviations can be ambiguous or unsafe, patients should ask clinicians or pharmacists to explain unfamiliar wording rather than infer its meaning.

\synonyms
Terms Medical Abbreviations (Common Medical Abbreviations And Terms)

\medterm{Common Migraine} A migraine attack without a preceding aura. It usually causes moderate or severe throbbing pain lasting four to 72 hours, often with nausea or sensitivity to light or sound and worsening with routine activity.

\synonyms
Migraine without aura; migraine without neurologic aura.

\medterm{Common Peroneal Nerve} A major nerve running from the back of the thigh around the outer knee into the lower leg and foot. It controls lifting and turning the foot and provides sensation to much of the lower leg and foot.

\medterm{Common Peroneal Nerve Dysfunction} Impairment of the common peroneal nerve, often from compression near the fibular head, trauma, surgery, or systemic nerve disease. It may cause foot drop, weakness of ankle eversion or dorsiflexion, and sensory loss over the lateral leg or dorsum of the foot.

\medterm{Common Peroneal Neuropathy} A peripheral nerve disorder causing weakness of ankle dorsiflexion and eversion, sensory
loss over the anterolateral leg and dorsum of the foot, and possible foot drop.

\medterm{Common Symptoms During Pregnancy} Frequent pregnancy-related symptoms such as nausea, fatigue, breast tenderness, heartburn, constipation, and urinary frequency; severe, sudden, or unusual symptoms may signal complications.

\medterm{Common Variable Immunodeficiency} An immune disorder in which the body makes too few effective antibodies. It causes repeated sinus, ear, or lung infections and may also cause digestive problems, autoimmune disease, or enlarged lymph nodes.

\medterm{Common Warts (Verruca Vulgaris)} Benign, rough-surfaced skin growths caused by infection with human papillomavirus, commonly affecting the hands and fingers.

\synonyms
Verruca Vulgaris

\medterm{Commotio Cordis} Sudden cardiac arrest caused by a blunt, nonpenetrating impact to the chest at a
vulnerable moment in the cardiac cycle, without structural heart injury.

\medterm{Communicable Disease} Any disease caused by bacteria, viruses, or other pathogens that is spread from person
to person.

\medterm{Communicating Hydrocephalus} A type of hydrocephalus in which there is free flow of cerebrospinal fluid between
the ventricles, but impaired absorption of CSF at the arachnoid granulations or
increased CSF production. The ventricles are uniformly dilated, including the fourth
ventricle.

\medterm{Communicating With Patients} The clinical exchange of information, concerns, preferences, and decisions between healthcare professionals and patients, using clear language, listening, empathy, and confirmation of understanding.

\medterm{Communicating With Someone With Aphasia} Communication strategies for a person with impaired language from brain injury or disease, including short sentences, extra time, visual cues, and supported conversation rather than assuming impaired intelligence.

\medterm{Communicating With Someone With Dysarthria} Communication strategies for a person whose speech muscles are weak or poorly coordinated, including reducing background noise, confirming messages, and using writing or augmentative aids when needed.

\medterm{Community Acquired Infection} An infection present or incubating at hospital admission, or developing within 48 hours
of admission without recent healthcare exposure. Common examples include pneumonia,
urinary tract infection, and meningitis caused by organisms such as Streptococcus
pneumoniae and Escherichia coli, generally susceptible to narrow-spectrum antibiotics.

\medterm{Community Health Worker} A trained member of a community who helps connect people with health services, education, prevention, and social support.

\medterm{Community Rehabilitation} Therapy and support delivered at home, in clinics, or in community settings to improve movement, communication, independence, participation, and daily functioning after illness or injury.
\medterm{Community Reintegration} The process of returning to everyday life and taking part in work, education, family, social, and leisure activities after illness or injury. Support may include rehabilitation, changes to the home or workplace, transport help, and practical or emotional support.

\medterm{Community-Acquired} Acquired outside a healthcare facility or before healthcare exposure. For example, community-acquired pneumonia begins outside a hospital and may involve different organisms and resistance patterns from hospital-acquired disease.
\medterm{Community-Acquired Pneumonia} Pneumonia acquired outside a hospital or long-term-care facility. It may cause cough, fever, chest pain, breathlessness, and low oxygen; diagnosis and treatment depend on examination, imaging, severity, likely organisms, comorbidities, and local antibiotic guidance. CURB-65 or PSI can help assess the need for hospital care but does not replace clinical judgment.

\medterm{Community-Acquired Pneumonia In Adults} Pneumonia acquired outside a hospital or healthcare facility; assessment considers severity, oxygenation, comorbidities, likely pathogens, and the need for outpatient or inpatient treatment.

\medterm{Comorbidity} Presence of two or more chronic conditions in same individual. Common in elderly:
hypertension, diabetes, arthritis, COPD, depression. Increases complexity of care,
polypharmacy risk, functional decline. Comprehensive assessment needed for each
condition and interactions.

\medterm{Computational Biomedicine} An interdisciplinary field that uses computers, mathematical models, simulations, and biomedical data to study health and disease. It may support research in areas such as genomics, drug development, medical imaging, disease prediction, diagnosis, and personalized treatment.

\synonyms
Computational Biomedical Science

\medterm{Compact Bone} Dense outer bone tissue arranged in strong cylindrical units called osteons. It provides most of the skeleton’s strength and protects inner marrow while supporting movement and weight-bearing.

\synonyms
Cortical bone; dense bone.

\medterm{Comparison} A reference in an imaging report to earlier studies used to assess whether a finding has changed.

\medterm{Compartment Pressure} The pressure inside a closed group of muscles, nerves, and blood vessels. If it rises too much, blood flow can fall and cause permanent tissue damage; measurements support but do not replace examination.

\medterm{Compartment Syndrome} Dangerous pressure buildup inside a confined muscle group that reduces blood flow and can permanently damage muscle and nerves. Severe pain, swelling, numbness, or weakness after injury requires emergency assessment and often surgery.

\synonyms
Syndrome Compartment (Compartment Syndrome)

\medterm{Competence} A legal determination that a person can make a particular decision or manage a matter. In clinical care, decision-making capacity is assessed for each specific choice.

\medterm{Competitive Antagonist} A medicine or substance that occupies the same receptor as another substance but does not activate it. By competing for that site, it reduces the other substance's effect; increasing the other substance may overcome the competition.

\medterm{Competitive Inhibition} Reduced enzyme activity caused when an inhibitor competes with the normal substance for the enzyme's active site. Because the binding is usually reversible, a higher concentration of the normal substance may lessen the inhibition.

\medterm{Complement} Complement is a system of blood proteins that enhances antibody and immune-cell defense by marking, inflaming, or directly damaging targets.

\medterm{Complement C1Q Deficiency} A rare inherited lack of a complement protein needed for immune defense and removal of antibody-bound material. It can cause repeated infections and autoimmune illness resembling lupus, including rash, joint, kidney, or blood problems.

\medterm{Complement C2/C4 Deficiency} Inherited or acquired lack of complement C2 or C4, proteins needed especially for the classical complement pathway. It can increase susceptibility to infections with encapsulated bacteria and autoimmune disease resembling lupus. Testing may show low CH50 with low C2 or C4; care includes vaccination, prompt infection treatment, and specialist management of autoimmune disease or recurrent infections.

\medterm{Complement C3} A major blood protein in the complement immune system. When activated, it helps mark microbes for removal and promotes inflammation; severe deficiency can cause repeated serious infections.

\medterm{Complement Cascade} A chain reaction involving many immune proteins in blood and tissues. It marks microbes for destruction, attracts immune cells, increases inflammation, and can directly break holes in some microbes.

\medterm{Complement Component 3 (C3)} A blood test measuring complement component C3, a major protein of the complement immune system; abnormal levels can support evaluation of immune-complex disease, infection, or complement deficiency.

\medterm{Complement Component 4} A blood test measuring complement component C4, part of the classical complement pathway; low levels may occur with immune-complex disease, hereditary angioedema, or complement consumption.

\medterm{Complement Deficiencies} Rare primary immunodeficiencies caused by genetic defects in complement pathway
proteins, increasing susceptibility to specific infections and autoimmune diseases.

\medterm{Complement Deficiency} Reduced or absent activity of one or more complement proteins, increasing susceptibility to some infections or autoimmune disease.

\medterm{Complement Fixation Test For C Burnetii} A serologic test detecting complement-fixing antibodies to Coxiella burnetii, the bacterium that causes Q fever; antibody phase patterns and titers help distinguish recent from persistent infection.

\medterm{Complement System} A group of blood proteins that supports immune defense by destroying microbes and promoting inflammation and phagocytosis.

\medterm{Complementary Sequence} A DNA or RNA sequence whose building blocks pair with those of another sequence. This matching allows genetic strands to bind and is used in copying, testing, and regulating genetic information.

\medterm{Complete Blood Count} A blood test that measures red blood cells, white blood cells, platelets, hemoglobin, and hematocrit. It may also measure red-cell size and hemoglobin content and list the different types of white blood cells. It helps assess anemia, infection, inflammation, bleeding, and some bone-marrow disorders.

\medterm{Complete Breech} A fetal position in which both hips and knees are bent, so the baby is sitting cross-legged with the buttocks or feet closest to the birth canal. Delivery planning depends on pregnancy, fetal, and maternal factors.

\synonyms
Flexed Breech; Complete Breech Presentation

\medterm{Complete Cystectomy} Surgery to remove the entire urinary bladder, usually for invasive bladder cancer. Because urine can no longer leave through the bladder, the surgeon creates a urinary diversion such as an ileal conduit, a continent reservoir, or a new bladder from bowel.

\medterm{Complete Denture} A removable artificial set of teeth replacing all teeth in the upper jaw, lower jaw, or both. It rests on the gums and may improve chewing, speech, and appearance but can require adjustment for comfort and stability.
\medterm{Complete Miscarriage} Pregnancy loss in which all pregnancy tissue has passed from the uterus, usually followed by easing bleeding and cramps. A clinician may use examination, hormone testing, or ultrasound to confirm that no tissue remains and to exclude an ectopic pregnancy; heavy bleeding, severe pain, dizziness, or fever needs urgent care.

\synonyms
Complete Spontaneous Abortion

\medterm{Complete Response} No detectable evidence of cancer after treatment on the tests used. It does not prove that every cancer cell has been eliminated, so continued follow-up is needed.

\medterm{Complex Carbohydrates} Carbohydrates made up of longer chains of sugars, chiefly starches and dietary fiber,
found in foods such as whole grains, legumes, vegetables, and potatoes. Fiber is not
fully digested and supports bowel function, satiety, and metabolic health, while starch
is digested into glucose at a rate that depends on food structure and processing.

\medterm{Complex Febrile Seizure} A seizure triggered by fever in a young child that lasts more than 15 minutes, affects only part of the body, or happens more than once during the same illness. It needs medical assessment to exclude brain infection or another cause. Most children recover fully, but the risk of another seizure is higher than with a simple febrile seizure.

\synonyms
Atypical Febrile Seizure; Complicated Febrile Seizure

\medterm{Complex Regional Pain Syndrome} Persistent severe pain and sensitivity in a limb after injury, surgery, or sometimes without an obvious trigger. Swelling, skin-color or temperature changes, sweating, stiffness, and weakness may also occur.

\medterm{Compliance} The extent to which a person follows a treatment recommendation; adherence is generally preferred because it emphasizes an agreed plan.

\medterm{Complicated Bereavement} Alternate index label; see \textit{Complicated grief}.

\medterm{Complicated Grief} Persistent, intense grief that substantially interferes with daily functioning and lasts longer than expected for the person's
cultural and social context.

\synonyms
Complicated Bereavement
Traumatic Grief

\medterm{Complicated Migraine} An older, nonspecific term for migraine accompanied by prolonged or prominent neurologic symptoms; the specific syndrome should be identified when possible.

\synonyms
Migraine with prolonged neurologic symptoms; complex migraine, historical usage.

\medterm{Complicated Urinary Tract Infection} A urinary infection associated with functional or structural abnormality, obstruction,
catheterization, renal impairment, male anatomy, pregnancy, or another factor increasing
treatment failure or complication risk. Management requires culture-directed therapy and
correction of the underlying problem when possible.

\medterm{Complication} Any adverse event arising during or after a surgical procedure, graded by the
Clavien-Dindo classification from minor deviation (Grade I) through organ failure (Grade
IV) to death (Grade V). Common types include surgical site infection, hemorrhage, wound
dehiscence, anastomotic leak, ileus, and venous thromboembolism.

\medterm{Complications (Vascular)} Problems involving blood vessels after disease, injury, or a procedure, such as bleeding, hematoma, thrombosis, embolism, vessel injury, or reduced blood flow to a limb. Symptoms and urgency depend on the vessel and tissue affected.

\synonyms
Vascular

\medterm{Composite Resin} A tooth-colored material used to repair cavities, chipped teeth, or other defects. It is bonded to the tooth and shaped to match it, but may wear or stain sooner than some harder restorations.

\medterm{Compound Fracture (Open Fracture)} A fracture associated with a wound that communicates between the fracture site and the external environment. It carries increased risks of infection and tissue damage.

\synonyms
Open Fracture

\medterm{Compound Heterozygote} A person with two different disease-causing variants in the same gene, one in each copy. Together, the variants may cause an inherited disorder even though neither copy has the same change.

\medterm{Compound Presentation} A birth position in which a hand, arm, foot, or another limb lies beside the baby's head or buttocks. The limb may move back during labor, but the position can obstruct delivery or increase injury risk.

\medterm{Compounding} Preparation, mixing, or altering of a medicine to meet an individual patient's needs when an appropriate approved product is unavailable.

\medterm{Compounding Pharmacy} A pharmacy that prepares customized formulations for an individual patient when a suitable commercially manufactured product is unavailable or unsuitable.

\synonyms
Compounding pharmacy practice; extemporaneous pharmacy.

\medterm{Comprehensive Geriatric Assessment} A multidimensional assessment of an older person's medical conditions, medicines, function, cognition, mood, and social
support, used to develop an individualized care plan.

\medterm{Comprehensive Metabolic Panel} A group of blood tests measuring electrolytes, glucose, kidney function, liver-related chemicals, proteins, and calcium. It helps assess general health and monitor or investigate metabolic, liver, and kidney disorders.

\medterm{Compress} To press together or apply pressure, as with a bandage or device, or to reduce the volume of something.

\medterm{Compressed Nerve} Alternate index label; see \textit{Pinched nerve}.

\medterm{Compression} Pressure applied to tissue or a structure, which may reduce swelling or bleeding but can also impair nerves, blood vessels, or organs when excessive.

\medterm{Compression Fracture} Collapse of a vertebra or other bone because of compression, often related to osteoporosis or trauma.

\medterm{Compression Fractures Of The Back} Collapse of one or more vertebral bodies caused by weakened bone, often from osteoporosis, or by substantial trauma. It can cause sudden back pain, height loss, kyphosis, and rarely nerve or spinal-cord compression.

\medterm{Compression Stockings} Elastic garments that apply graduated pressure to the legs, supporting venous return and reducing swelling or selected risks of venous pooling; correct fit and arterial circulation matter.

\medterm{Compression Therapy} Use of external pressure, usually from garments or bandages, to support venous or lymphatic flow and reduce swelling in selected
conditions.

\medterm{Compression Ultrasound} An ultrasound test in which gentle pressure is applied to a vein. A healthy vein usually collapses; failure to collapse may indicate a blood clot, especially in the deep veins of a limb.

\medterm{Compressive Atelectasis} Collapse of part of a lung because fluid, air, a tumor, or another process presses on it from outside. The lung may re-expand when the pressure is relieved.

\medterm{Compton Scattering} An interaction in which a photon transfers part of its energy to an electron and changes direction. It contributes to scattered
radiation in imaging.

\medterm{Compulsions} Repetitive behaviors or mental acts that the individual feels driven to perform in
response to an obsession or according to rigid rules. They are aimed at preventing or
reducing anxiety or preventing some dreaded event, but are not connected in a realistic
way with what they are designed to neutralize.

\medterm{Compulsive Gambling} Alternate, nonpreferred label; see \textit{Gambling Disorder}.

\medterm{Compulsive Obsessive Disorder (Obsessive Compulsive Disorder)} Alternate index label for Obsessive Compulsive Disorder.

\medterm{Compulsive Overeating} Alternate index label; see \textit{Binge-eating disorder}.

\medterm{Compulsive Sexual Behavior} Recurrent sexual thoughts or behaviors that are difficult to control and cause significant distress or impairment.

\synonyms
Sexual Obsession

\medterm{Compulsive Stealing} Alternate index label; see \textit{Kleptomania}.

\medterm{Computed Tomography} An imaging test that uses rotating X-rays and computer processing to create detailed cross-sectional pictures. It is fast and can show bones, bleeding, organs, blood vessels, tumors, infection, and internal injury.

\medterm{Computed Tomography Pulmonary Angiography} A CT scan performed after injected contrast highlights the lung arteries. It is commonly used to look for a pulmonary embolism and may also show strain on the heart or other chest abnormalities.

\synonyms
CTPA; CT Pulmonary Angiogram; Pulmonary CT Angiography

\medterm{Computed Tomography Scan} An X-ray scan processed by a computer to produce cross-sectional images of the body. It can rapidly detect bleeding, fractures, infection, blocked bowel, blood clots, tumors, and many other urgent conditions.

\medterm{Computerized Axial Tomography} Computed tomography, an imaging method that uses X-rays and computer processing to produce cross-sectional images of the body; it is also called a CT or CAT scan.

\medterm{Concato'S Disease} A historical term for primary polycythemia, now generally called polycythemia vera, a
myeloproliferative neoplasm causing increased red-cell mass and often elevated blood
viscosity.

\medterm{Concealed Conduction} A heart rhythm event in which an electrical impulse enters part of the heart's conduction system but stops before producing a visible heartbeat. The affected tissue may then delay or block the next impulse and alter the heartbeat pattern.

\medterm{Concealed Pregnancy} Pregnancy that is hidden or not recognized until delivery, often by denial (psychogenic)
or deliberate concealment; a medicolegal concept relevant to infanticide.

\medterm{Concentration} The amount of a substance in a given volume, such as the level of glucose or medicine in blood. It is reported in units such as milligrams per deciliter or millimoles per liter.

\medterm{Conception} The beginning of pregnancy, commonly referring to fertilization followed by implantation and establishment of a developing embryo in the uterus.

\medterm{Conception (Fertilization)} The joining of a sperm and egg, usually in a fallopian tube, to form the first cell of an embryo. The resulting cells divide as they travel toward the uterus, where implantation may occur several days later.

\synonyms
Fertilization

\medterm{Conchae} Curved bony structures covered by mucosa inside the nasal cavity that warm, humidify, and filter inhaled air.

\medterm{Conclusion} The part of a report that summarizes the main findings and their clinical significance or recommended next steps.

\medterm{Concomitant} Existing or occurring at the same time as another condition, treatment, or event. Concomitant medicines or illnesses may change diagnosis, treatment choices, expected effects, side effects, or monitoring needs.

\medterm{Concordance} Probability both twins have same trait. Monozygotic (identical) twins > dizygotic
(fraternal) for genetic traits. High concordance suggests genetic influence.

\medterm{Concrescence} Union of adjacent teeth by cementum only. Usually posterior teeth.

\medterm{Concretio Cordis} Abnormal adhesion or fusion of structures within or around the heart, including adhesion
of the pericardial surfaces or fusion of cardiac components.

\medterm{Concussion} A form of mild traumatic brain injury caused by a biomechanical force (direct blow to
the head, face, neck, or elsewhere on the body with impulsive force transmitted to the
head), resulting in a transient disturbance of brain function without structural damage
visible on conventional neuroimaging.





\medterm{Condition} A general term for a state of bodily health or, in medical usage, a specific deviation
from normal in structure or function that may or may not constitute a disease, including
complaints lacking a discrete pathologic basis and overt diagnoses.

\medterm{Conditioning} Learning in which repeated association or reinforcement changes a behavior or response; in medicine, conditioning can also mean preparing a patient or tissue for a treatment.

\medterm{Conditions Caused By Deletion Mutations} Deletion mutations remove DNA segments and can disrupt one gene or several neighboring genes. The clinical effect depends on size, location, inheritance, and whether the remaining copy compensates; genetic counseling and testing may clarify diagnosis and recurrence risk.

\medterm{Condom} A thin barrier sheath worn on the penis or inserted into the vagina that helps prevent pregnancy and reduces transmission of sexually transmitted infections when used correctly.

\medterm{Condoms - Male} A sheath placed over the penis during sexual activity to reduce exchange of semen and genital secretions and lower pregnancy and sexually transmitted infection risk when used correctly.

\medterm{Conduct Disorder} A childhood or adolescent disorder involving a persistent pattern of behavior that violates
the rights of others or major age-appropriate rules. Examples include aggression,
property destruction, deceitfulness, theft, and serious rule violations.

\medterm{Conduction Deafness} Hearing loss caused by impaired transmission of sound through the external or middle
ear, with the inner ear and neural auditory pathways relatively preserved.

\medterm{Conductive Deafness} Hearing loss resulting from a problem in the external auditory canal, tympanic membrane,
or middle ear that reduces conduction of sound to the inner ear.

\medterm{Conductive Hearing Loss} Hearing loss caused by a problem in the outer or middle ear that prevents sound from reaching the inner ear normally. Earwax, fluid, a damaged eardrum, bone changes, or ear-bone disease may be responsible.

\medterm{Condylar Process} The rounded articular projection of a bone, especially the mandibular process that
articulates with the temporal bone to form the temporomandibular joint.

\medterm{Condyle} A rounded projection at the end of a bone that forms part of a joint and helps guide movement against another bone.

\medterm{Condyloma} A wart-like growth, often in the genital or anal area and commonly caused by human papillomavirus; condyloma lata are lesions of secondary syphilis.

\medterm{Condyloma Acuminata} Soft, wart-like growths around the genitals or anus caused by certain types of human papillomavirus. They spread through close skin-to-skin sexual contact; treatment removes visible warts but may not remove the virus, and vaccination helps prevent many infections.

\medterm{Condyloma Acuminatum} A soft, wart-like growth around the genitals or anus caused by certain human papillomavirus types. It spreads through intimate skin contact; treatment removes visible warts, but the virus may remain and recur.

\medterm{Condyloma Latum} A broad, flat, moist, highly infectious lesion of secondary syphilis, usually around the genitals or anus. It differs from a genital wart and indicates syphilis requiring testing and antibiotic treatment.
\medterm{Condylomata Acuminata} Alternate index label; see \textit{Genital warts}.

\medterm{Condylomata Lata} Broad, flat, moist, grayish-white papules occurring in secondary syphilis, typically
found in the anogenital region and other intertriginous areas. They are highly
infectious and contain abundant Treponema pallidum organisms. Diagnosis is by darkfield
microscopy or serological testing for syphilis.

\medterm{Cone Biopsy} A surgical procedure that removes a cone-shaped portion of the cervix for diagnosis or treatment of high-grade cervical precancer and to determine whether invasive cancer is present.

\medterm{Cone-Beam Computed Tomography} A dental or facial scan that uses a cone-shaped X-ray beam to create three-dimensional images of teeth, jawbone, and nearby structures. It should be used when the extra detail will affect care.

\medterm{Cones} Light-sensitive retinal cells responsible for color vision and sharp detail, working best in daylight or other relatively bright conditions.

\synonyms
Cone photoreceptors; cone cells.

\medterm{Confabulation} Unintentional production of inaccurate memories or stories without an attempt to deceive. It may occur with memory disorders or injury affecting the front part of the brain.

\medterm{Confidentiality} The ethical and legal duty of healthcare professionals to protect patient information
from unauthorized disclosure. Founded on the principle of respect for patient autonomy
and trust, confidentiality is a core requirement of medical ethics (GMC guidelines) and
law (GDPR, Data Protection Act, Caldicott principles).

\medterm{Confirmatory Test} A test used to verify a preliminary or screening result, often with a different method
or greater specificity. Confirmation should be selected according to the condition,
specimen, timing, and accepted clinical guidelines.

\medterm{Conformal Radiation Therapy} Radiation treatment planned so that the high-dose region conforms closely to the tumor's three-dimensional shape while limiting exposure of nearby normal tissue.

\medterm{Confounding} A distortion that occurs when another factor is related to both an exposure and an outcome, making their relationship appear stronger, weaker, or different from reality. For example, age may affect both an exposure and disease risk. Study design and statistical adjustment can reduce confounding.

\medterm{Confusion} Impaired ability to think clearly, orient, attend, remember, or respond appropriately; sudden confusion is delirium and may indicate infection, medication effect, metabolic illness, stroke, or another emergency.

\medterm{Confusion (Delirium)} A sudden change in attention and awareness that fluctuates during the day. Infection, medicines, pain, dehydration, metabolic problems, constipation, or urinary retention may cause it and require prompt assessment.

\synonyms
Delirium

\medterm{Confusion Assessment Method} A bedside assessment for delirium. It checks for a sudden or fluctuating change, poor attention, disorganized thinking, and an altered level of alertness; the pattern helps identify delirium but does not replace finding its cause.

\medterm{Congenital} Present at birth, whether inherited, caused by an early developmental difference, or acquired before birth through infection, exposure, or another condition.

\medterm{Congenital Adrenal Hyperplasia} A group of inherited disorders in which the adrenal glands cannot make hormones normally. The most common form can cause salt loss, excess androgen effects, early puberty, or genital differences and may be life-threatening in newborns.

\medterm{Congenital Antithrombin III Deficiency} An inherited deficiency or dysfunction of antithrombin, a natural anticoagulant, causing increased risk of venous thromboembolism.

\medterm{Congenital Biliary Atresia} Alternate index label; see \textit{Biliary atresia}.

\medterm{Congenital Cataract} Clouding of the eye lens present at birth or developing soon afterward. It can interfere with visual development and may require early eye assessment and treatment to prevent lasting reduced vision.

\medterm{Congenital CMV} Cytomegalovirus infection acquired before birth. Some babies have small head size, hearing or vision loss, seizures, jaundice, an enlarged liver or spleen, or developmental problems; others appear well initially but develop later hearing loss.

\medterm{Congenital CNS Malformation} A structural abnormality of the brain, spinal cord, or related developmental structures present from birth, caused by disturbed neural development or other prenatal factors; severity ranges from incidental to life-threatening.

\medterm{Congenital Contractural Arachnodactyly} An inherited connective-tissue disorder causing unusually long, slender fingers and toes, tight joints present from birth, thin muscles, and sometimes a curved spine. Joint stiffness may improve with age.

\medterm{Congenital Cytomegalovirus} Cytomegalovirus infection present at birth after transmission during pregnancy. Some infants have hearing loss, vision problems, small head size, seizures, or developmental disability; others have no symptoms initially.

\medterm{Congenital Cytomegalovirus Infection} Cytomegalovirus infection acquired before birth. It may cause hearing loss,
developmental problems, microcephaly, or no symptoms.

\medterm{Congenital Diaphragmatic Hernia} A defect in the diaphragm allowing abdominal organs to herniate into the thoracic
cavity, causing pulmonary hypoplasia and respiratory distress at birth. Most commonly a
posterolateral (Bochdalek) defect on the left side. Presents with cyanosis, scaphoid
abdomen, and bowel sounds in the chest.

\medterm{Congenital Diaphragmatic Hernia Repair} Congenital diaphragmatic hernia repair is a procedure for congenital diaphragmatic hernia repair; indication, preparation, risks, technical details, and recovery depend on the patient and treatment goal.
\medterm{Congenital Disorder} A medical condition or structural abnormality present in an individual from birth, which
may be caused by genetic factors or environmental influences in the womb.

\medterm{Congenital Fibrinogen Deficiency} An inherited absence or abnormality of fibrinogen that can cause bleeding, thrombosis, or both depending on the specific molecular defect.

\medterm{Congenital Glaucoma} A rare form of glaucoma (1 in 10,000-20,000 births) caused by developmental anomalies of
the trabecular meshwork and anterior chamber angle, impairing aqueous humor outflow.
Most cases are sporadic; familial forms are autosomal recessive (CYP1B1 gene mutations).

\medterm{Congenital Heart Defect - Corrective Surgery} Surgery or catheter-based intervention to repair a structural heart abnormality present at birth, with the specific procedure depending on the defect, its severity, and the patient’s anatomy and condition.

\medterm{Congenital Heart Defects} Problems in the heart or major blood vessels
that are present at birth. They can be mild or serious and may affect how blood
flows through the heart and body. Examples include holes in the heart, abnormal
valves, and problems with the major blood vessels.

\medterm{Congenital Heart Defects In Children} Alternate index label; see \textit{Congenital Heart Defects}.

\medterm{Congenital Heart Disease} A heart or major blood-vessel problem present from birth. It may change blood flow, cause no symptoms, or lead to breathing difficulty, blue color, poor growth, heart failure, or abnormal rhythms. Treatment ranges from monitoring and medicines to catheter procedures or surgery.

\medterm{Congenital Heart Disease In Adults} A heart defect present from birth that continues into adulthood; see \textit{Congenital Heart Defects}.

\medterm{Congenital Heart Murmur} A heart murmur heard in a newborn or infant. It may be innocent or may indicate a
structural heart defect; assessment considers the murmur's characteristics and associated symptoms or examination
findings.

\medterm{Compensated Heart Failure} Heart failure in which symptoms and fluid buildup
are stable or controlled with treatment. The heart failure is still present and can
worsen, especially if a new illness or other problem develops.

\medterm{Congenital Hip Dislocation} Abnormal development or displacement of the hip joint present at birth or developing in
infancy, now generally called developmental dysplasia of the hip.

\medterm{Congenital Hyperinsulinism} An inherited or developmental disorder in which a baby's pancreas releases too much insulin. This can cause dangerously low blood sugar, poor feeding, seizures, brain injury, or death unless recognized and treated promptly.

\medterm{Congenital Hyperthyroidism} Excess thyroid hormone in a newborn, most often caused by thyroid-stimulating antibodies passing from a mother with Graves disease. It may cause a fast heartbeat, irritability, poor feeding, poor weight gain, or an enlarged thyroid.

\medterm{Congenital Hypothyroidism} Too little thyroid hormone from birth, usually because the thyroid is missing, misplaced, or underdeveloped. Newborn screening and early treatment are essential because untreated deficiency can impair growth and brain development.

\medterm{Congenital Malformation} A structural abnormality that develops before birth and is present at birth, caused by genetic, environmental, developmental, or unknown factors.

\medterm{Congenital Megacolon} A birth condition in which part of the intestine lacks the nerve cells needed to move stool. It causes blockage and enlargement of the bowel above it, with delayed first stool, swelling, vomiting, or severe constipation.

\medterm{Congenital Megacolon (Hirschsprung Disease)} Alternate index label for Hirschsprung Disease.

\medterm{Congenital Melanocytic Nevus} A dark mole present at birth or appearing soon afterward, caused by a collection of pigment-producing cells. Most remain harmless, but larger lesions have a higher melanoma risk and need specialist monitoring.

\medterm{Congenital Mitral Valve Anomalies} Structural abnormalities of the mitral valve that are present at birth and may impair blood flow.

\medterm{Congenital Muscular Dystrophy} A group of inherited muscle disorders causing weakness or low muscle tone at birth or during early infancy. Some types also affect the brain, eyes, breathing, heart, or joints; severity varies widely.

\synonyms
CMD; Congenital Dystrophy

\medterm{Congenital Myasthenic Syndromes} A group of inherited disorders that disrupt communication between nerves and muscles. They cause weakness that worsens with activity and may affect the eyes, swallowing, breathing, or limb movement.
\medterm{Congenital Nephrotic Syndrome} Severe kidney disease beginning before birth or during the first three months of life, causing large amounts of protein loss in urine, swelling, low blood protein, poor growth, infections, and increased clotting risk.

\medterm{Congenital Nevus} A melanocytic mole present at birth or developing soon afterward. It is composed of melanocytes in
the epidermis and dermis. Size, location, and changes over time determine the need for
clinical assessment and surveillance.

\medterm{Congenital Penile Curvature} Curvature of the penis present from birth, usually caused by unequal development of the erectile tissues. Mild curvature may need no treatment; significant curvature causing functional difficulty may require specialist assessment.

\medterm{Congenital Plagiocephaly} Alternate index label; see \textit{Craniosynostosis}.

\medterm{Congenital Platelet Function Defects} Inherited disorders in which platelets do not work normally even though their number may be normal. They can cause easy bruising, nosebleeds, heavy menstrual bleeding, or excessive bleeding after injury or surgery.

\medterm{Congenital Pneumonia} A lung infection present at birth, acquired before or during delivery. A newborn may have rapid or difficult breathing, temperature instability, poor feeding, or sepsis and needs urgent evaluation and treatment.

\medterm{Congenital Protein C Or S Deficiency} An inherited lack or abnormality of protein C or protein S, natural substances that limit clotting. It increases the risk of blood clots in veins, sometimes beginning in infancy or childhood.
\medterm{Congenital Rubella} Fetal infection with rubella virus acquired during pregnancy, potentially causing hearing loss, cataracts, congenital heart disease, developmental impairment, and other findings.

\medterm{Congenital Rubella Syndrome} Birth defects caused by rubella infection during pregnancy, especially early pregnancy. Affected babies may have hearing loss, cataracts, heart defects, small head size, developmental problems, enlarged organs, or a low platelet count.

\medterm{Congenital Scoliosis} A lateral curvature of the spine caused by congenital vertebral anomalies that arise
during embryogenesis (failure of formation or segmentation). Failure of formation:
hemivertebra (wedge-shaped vertebra causing progressive curvature). Failure of
segmentation: unilateral unsegmented bar (block of bone on one side).

\medterm{Congenital Syphilis} Treponema pallidum infection acquired transplacentally. Can cause stillbirth,
prematurity, hepatosplenomegaly, rash, snuffles, osteochondritis, and Hutchinson teeth.
Diagnosed by darkfield microscopy, VDRL/RPR in neonatal serum, or PCR.

\medterm{Congenital Toxoplasmosis} Toxoplasma gondii infection acquired in utero, potentially causing chorioretinitis, intracranial calcifications, hydrocephalus, seizures, or developmental impairment.

\medterm{Congenital Zika Syndrome} Fetal brain infection with Zika virus, causing microcephaly, brain abnormalities, eye
defects, hearing loss, and limb contractures. Transmission occurs via mosquito bite or
sexual contact. Diagnosis by Zika virus PCR in maternal serum or amniotic fluid. No
specific treatment. Prevention by avoiding endemic areas during pregnancy.

\medterm{Congestion} A passive vascular phenomenon of engorgement of a tissue or organ with blood from
impaired venous or lymphatic drainage. Chronic venous congestion produces stasis,
hypoxia, dilated vessels, edema, hemorrhage, and hemosiderin-laden macrophages,
exemplified by the lungs (heart failure cells) and liver (nutmeg liver).

\medterm{Congestive Heart Failure} A historical term for heart failure with clinical fluid congestion. Current terminology
is heart failure, which may occur with or without congestion and is classified by
left-ventricular ejection fraction, symptoms, and disease stage.

\medterm{Congestive Heart Failure Overview} Alternate index label; see \textit{Heart failure}.
\medterm{Conization} Surgical removal of a cone-shaped portion of the cervix for diagnosis or treatment of
cervical dysplasia or selected early cervical cancers.

\medterm{Conjoined Twins} Monozygotic, monochorionic, monoamniotic twins who are physically connected at birth,
resulting from incomplete fission of the fertilized ovum after day 13 of fertilization.
Incidence: 1 in 50,000-100,000 births.

\medterm{Conjugated Equine Estrogens} A mixture of estrogenic compounds historically derived from the urine of pregnant mares and used in some hormone-therapy formulations.

\synonyms
CEE; conjugated estrogens.

\medterm{Conjunctiva} A thin mucous membrane that lines the inner eyelids (palpebral conjunctiva) and covers the front surface of the sclera (bulbar conjunctiva). It protects and lubricates the eye and permits smooth movement of the eyelids.

\medterm{Conjunctival} Relating to the conjunctiva, the thin mucous membrane covering the front of the eye and lining the inner eyelids.

\medterm{Conjunctival Squamous-Cell Carcinoma} A malignant tumor of the conjunctival epithelium, typically arising from the limbal
area. It is associated with UV radiation exposure and can invade locally if not treated
promptly.

\medterm{Conjunctivitis} Inflammation of the conjunctiva, the thin, transparent mucous membrane lining the inner
surface of the eyelids and the anterior surface of the sclera. Conjunctivitis is the
most common cause of red eye and is classified by etiology.

\medterm{Conjunctivitis (Pink Eye)} Alternate index label for Pink Eye.

\medterm{Conjunctivitis Or Pink Eye} Inflammation of the conjunctiva causing redness, discharge, irritation, itching, or tearing. Causes include viral or bacterial infection, allergy, and chemical irritation; severe pain, photophobia, vision loss, or contact-lens use requires prompt eye assessment.

\medterm{Conjunctivoplasty} Plastic or reconstructive surgery of the conjunctiva, performed to repair defects,
scarring, adhesions, or other abnormalities of the ocular surface.

\medterm{Conn'S Syndrome} Primary hyperaldosteronism caused by autonomous secretion of aldosterone, leading to
hypertension and often low serum potassium.

\medterm{Connective Tissue} A diverse type of tissue that supports, binds, protects, and insulates organs and
structures throughout the body.

\medterm{Connective Tissue Disease} A disorder affecting tissues that support and connect organs, often because the immune system attacks the body’s own tissues. It may involve skin, joints, muscles, blood vessels, lungs, kidneys, or other organs. Examples include lupus, systemic sclerosis, rheumatoid arthritis, and Sjögren syndrome.

\medterm{Connective Tissue Disease-Related Interstitial Lung Disease} Interstitial lung disease associated with autoimmune connective-tissue disease such as systemic sclerosis, rheumatoid arthritis, polymyositis, dermatomyositis, or Sjogren syndrome; inflammation and fibrosis can impair gas exchange.

\medterm{Consanguinity} A biological relationship between people who share a recent common ancestor; it can increase the chance that both parents carry the same recessive disease variant.

\medterm{Conscious Sedation For Surgical Procedures} Use of medicines to make a person relaxed and sleepy during a procedure while allowing them to respond to voice or touch and usually breathe independently. Heart rate, breathing, oxygen, and alertness require continuous monitoring.
\medterm{Consolidation} A part of the lung whose air spaces have filled with fluid, pus, blood, cells, or another substance. It appears denser on imaging and may result from pneumonia, fluid overload, bleeding, or a tumor.

\medterm{Consolidation Therapy} Additional treatment given after an initial response or remission to destroy remaining disease and lower the chance of relapse. It may include medicines, radiation, surgery, or a transplant depending on the illness.
\medterm{Constipation} A condition involving fewer than three bowel movements per week, hard or lumpy stools, difficult or painful passage, straining, or a feeling that the bowel has not emptied completely. Causes include diet, medicines, dehydration, functional bowel disorders, and structural or systemic disease. Rectal bleeding, vomiting, constant abdominal pain, inability to pass gas, fever, or unexplained weight loss requires medical evaluation.

\medterm{Constipation (Laxatives For Constipation)} Alternate index label for Laxatives For Constipation.






\medterm{Constipation In Infants And Children} Infrequent, hard, painful, or difficult passage of stool in a child, caused by stool withholding, diet, dehydration, illness, medicines, or structural and neurologic disorders.

\medterm{Constitutional Delay In Puberty} A common pattern in which puberty starts later than average but eventually progresses normally. Growth and bone development may also be delayed, often with a family history of late puberty.

\medterm{Constraint-Induced Movement Therapy} Rehabilitation that limits use of a less-affected limb while providing structured practice with an affected limb to
improve functional use in selected people.

\medterm{Constrictive Pericarditis} Thickening and stiffening of the sac around the heart, preventing the heart from filling normally. It can cause breathlessness, leg or abdominal swelling, enlarged neck veins, and liver congestion; causes include infection, surgery, radiation, or unknown disease.

\medterm{Consultation-Liaison Psychiatry} The subspecialty of psychiatry focused on the interface between psychiatry and general
medical illness. Psychiatrists provide consultation on medical-surgical wards for
patients with psychiatric symptoms arising from or coexisting with medical conditions.


\medterm{Consumption} An older term for tuberculosis, especially wasting pulmonary tuberculosis; in general usage it also means the act or amount of using something.


\medterm{Contac Overdose} Overdose of a combination cold medicine that may cause acetaminophen toxicity, antihistamine or decongestant effects, sedation, seizures, or cardiac complications depending on ingredients.

\medterm{Contact Allergen Testing} Patch testing checks whether a substance causes allergic skin inflammation. Small amounts of suspected allergens are placed on the skin, usually the back, and the area is examined after one or more days for a delayed reaction.

\medterm{Contact Dermatitis} Red, itchy, burning, or blistered skin caused by contact with an irritating or allergy-triggering substance. Common triggers include soaps, chemicals, metals, fragrances, plants, and medicines.

\synonyms
Dermatitis, Contact

\medterm{Contact Precautions} Infection control measures used for patients colonised or infected with
multidrug-resistant organisms. Precautions include gown and glove use, patient placement
in single rooms, and dedicated equipment. Contact precautions are used for MRSA, VRE,
CRE, and other resistant organisms.

\medterm{Contact Tracing} Identifying and monitoring persons exposed to infectious disease. Critical for outbreak
control. Examples: COVID-19, Ebola, TB.




\medterm{Contagious Disease} An infectious disease that can be transmitted directly or indirectly from one person,
animal, or source to another.














\medterm{Contamination} Unwanted microorganisms or other material entering a specimen, surface, medicine, or normally sterile area without causing infection in the person. It can produce misleading test results and must be distinguished from true disease.

\medterm{Continuing Medical Education} Education after professional training that updates a healthcare worker's knowledge, skills, judgment, and ability to provide safe, effective patient care.

\medterm{Continuity Of Care} Consistent and coordinated care over time and across clinicians or settings, integrating
the patient’s history, current needs, treatment plan, and preferences. It supports
safety, adherence, prevention, and chronic-disease outcomes.

\medterm{Continuous Bladder Irrigation} Continuous flow of sterile fluid through a urinary catheter into and out of the bladder. It helps prevent blockage and wash out blood clots or debris, often after urinary-tract surgery.

\medterm{Continuous Combined Hormone Therapy} A hormone treatment regimen taking estrogen and progestogen continuously, often to manage menopausal symptoms while reducing regular withdrawal bleeding.

\synonyms
Continuous combined HRT; continuous estrogen-progestogen therapy.

\medterm{Continuous Electronic Fetal Monitoring} The continuous recording of fetal heart rate and uterine contractions using
external transducers (Doppler ultrasound and tocodynamometer) or internal devices (fetal
scalp electrode and intrauterine pressure catheter). Used intrapartum to detect fetal
hypoxia.

\medterm{Continuous Glucose Monitoring} A wearable sensor that measures glucose in fluid beneath the skin at regular intervals and shows current levels, trends, and alerts. It helps manage diabetes but may need confirmation with a blood-glucose test when readings do not fit symptoms.

\medterm{Continuous Kidney Replacement Therapy} A slow, continuous form of dialysis used mainly in intensive care for severe kidney failure. It gradually removes extra fluid and waste and is useful when a person is too unstable for conventional dialysis.

\medterm{Continuous Murmurs} Heart sounds heard throughout both contraction and relaxation because blood flows continuously through an abnormal opening or connection. A persistent murmur requires assessment to identify its cause and significance.
throughout the cardiac cycle.

\medterm{Continuous Positive Airway Pressure} Noninvasive respiratory support that delivers a constant positive pressure throughout
the breathing cycle to keep the upper airway or alveoli open.

\medterm{Continuous Subcutaneous Insulin Infusion} Delivery of rapid-acting insulin through a pump and a subcutaneous infusion set to provide basal insulin and mealtime doses.

\medterm{Contraception} Any method used to prevent pregnancy, including intrauterine devices, implants, pills, injections, condoms, fertility-awareness methods, and sterilization. Methods differ in effectiveness, duration, side effects, cost, and protection against sexually transmitted infections. The best choice depends on health, preferences, reproductive plans, and access to care.

\medterm{Contraception And Birth Control} Contraception encompasses methods used to prevent pregnancy, including hormonal options
like oral contraceptives, patches, rings, injections, and implants, as well as
non-hormonal methods like copper IUDs, barrier methods, diaphragms, and permanent
sterilization.

\medterm{Contraception For Adolescents} Methods to prevent unintended pregnancy. Options: combined oral contraceptives,
progestin-only pills, depot medroxyprogesterone acetate, contraceptive implant
(Nexplanon), intrauterine devices (IUDs), condoms, and abstinence.

\medterm{Contraceptive Device, Intrauterine} A small device placed inside the uterus to prevent pregnancy; copper IUDs impair sperm function, while hormonal IUDs release levonorgestrel and thicken cervical mucus.

\medterm{Contraceptive, Pill} An oral hormonal contraceptive, usually a combined estrogen--progestin pill or a progestin-only pill, that prevents ovulation and/or changes cervical mucus and the uterine lining.

\medterm{Contracted Pelvis} A pelvis whose size or shape is reduced or abnormal enough to interfere with passage of
the fetus during labor.

\medterm{Contractile Proteins} Proteins, especially actin and myosin, that interact inside muscle cells to generate force, shorten fibers, and produce body movement.

\synonyms
Muscle contractile proteins; actomyosin proteins.

\medterm{Contraction} Shortening or tightening of a muscle; uterine contractions during labor thin and dilate the cervix and help deliver the fetus.

\medterm{Contraction Stress Test} A pregnancy test that records the baby's heart rate while contractions occur, sometimes after gentle stimulation or medicine. Repeated slowing of the heart rate may suggest that the placenta cannot supply enough oxygen during labor.

\medterm{Contractions} Rhythmic tightening of the uterus during labor that thins and opens the cervix and helps move the baby through the birth canal. Before labor, Braxton Hicks contractions are often irregular and do not progressively open the cervix.

\medterm{Contractions Braxton-Hicks (Braxton Hicks Contractions)} Alternate index label for Braxton Hicks Contractions.

\medterm{Contracture} A lasting inability to fully move a joint because nearby muscles, tendons, ligaments, skin, or other tissues have shortened or stiffened. It may follow injury, burns, paralysis, inflammation, or prolonged immobility.

\medterm{Contracture Deformity} A fixed limitation of movement caused by shortening or tightening of muscle, tendon, ligament, joint capsule, skin, or other soft tissue. It may follow neurologic disease, prolonged immobility, injury, burns, or inflammation.

\medterm{Contracture Prevention} Measures such as positioning, stretching, exercise, splinting, orthoses, and serial casting that help preserve joint movement in people at risk of fixed shortening.

\medterm{Contraindication} A condition or circumstance in which a treatment, procedure, or medication should not be used because the potential harm outweighs the expected benefit.

\synonyms
Contraindications


\medterm{Contralateral} On or affecting the side opposite a stated body part or injury. For example, a stroke on the left side of the brain may cause weakness on the right side of the body.

\medterm{Contrast} A difference that makes structures distinguishable; in imaging, contrast material is administered to improve visualization of organs, vessels, or tissues.

\medterm{Contrast Agent} A substance given or placed in the body to make organs, tissues, blood vessels, or body spaces easier to see during imaging. Different agents are used for X-ray, CT, MRI, ultrasound, or digestive-tract studies.

\medterm{Contrast Echocardiography} Use of microbubble contrast agents to enhance endocardial border definition and assess
myocardial perfusion. Indicated for suboptimal non-contrast echo images and assessment
of myocardial viability. Microbubbles remain entirely within the intravascular space,
making them ideal for perfusion assessment.

\medterm{Contrast Medium} A substance administered or placed in the body to make tissues, organs, or body spaces more visible during imaging.

\synonyms
Contrast agent; radiographic contrast material.

\medterm{Contrast Radiology} Imaging of the gastrointestinal tract using contrast material, such as barium, to outline the esophagus, stomach, intestine, or colon on radiographs or fluoroscopy.

\medterm{Contrast Reaction} An unwanted response after an imaging contrast agent, ranging from nausea, itching, or a rash to breathing difficulty, low blood pressure, or a severe allergic-like reaction. Imaging staff should be told about previous reactions.
\medterm{Contrast-Associated Acute Kidney Injury} Acute worsening of kidney function occurring after administration of contrast, whether
or not contrast is causal. Risk assessment includes baseline kidney function,
hemodynamic status, diabetes, nephrotoxins, urgency, and the need to avoid unnecessary
contrast exposure.

\medterm{Contrast-Enhanced Ultrasound} An ultrasound examination using tiny gas-filled bubbles in the bloodstream to improve pictures of blood flow and help characterize selected liver, heart, blood-vessel, or other lesions.

\medterm{Contrast-Induced Nephropathy} A sudden decline in kidney function after iodinated contrast, although the contrast is not always the cause. Risk is higher with severe kidney disease, dehydration, low blood pressure, or other kidney stresses.

\medterm{Contrecoup} An injury occurring on the side of an organ opposite the site of impact, especially a
brain injury produced when the brain moves within the skull after trauma.

\medterm{Control (Research)} A comparison participant or group in a study that does not receive the experimental intervention, or receives standard care or a placebo, allowing the intervention's effects to be estimated.

\medterm{Control Group} The comparison group in a study, which may receive usual care, a placebo, or no intervention according to the study design.

\medterm{Controlled Cord Traction} A technique for active management of third-stage labor: gentle, steady traction on the
umbilical cord with one hand while the other hand applies suprapubic counter-pressure
(to prevent uterine inversion). Performed after signs of placental separation. The
Brandt-Andrews maneuver and the Finocchietti technique are CCT methods.

\medterm{Controlled Ovarian Hyperstimulation} Use of hormone medicines to make several ovarian follicles mature at the same time for fertility treatment. Ultrasound and blood tests guide dosing and help reduce the risk of ovarian hyperstimulation, a potentially serious complication.

\synonyms
COH; Ovarian Stimulation

\medterm{Controllers} Long-term asthma medicines taken regularly to reduce airway inflammation, prevent symptoms and attacks, and lower the need for quick-relief inhalers.

\synonyms
Controller medicines; maintenance asthma therapy.


\medterm{Contusion} A bruise caused by a blunt blow that damages small blood vessels without breaking the skin. It may cause pain, swelling, and discoloration; severe injuries can damage muscle, organs, or deeper tissues.

\medterm{Conus Medullaris Syndrome} A group of symptoms caused by damage to the lower end of the spinal cord. It may cause back pain, weakness or numbness in the legs and pelvic area, and loss of bladder, bowel, or sexual function.

\medterm{Convergence} The inward movement of both eyes so that they are directed toward the same nearby object.

\medterm{Convergence Insufficiency} Difficulty turning both eyes inward together to focus on something nearby. It can cause eyestrain, headaches, blurred or double vision, and reading difficulty, especially in children and young adults.

\medterm{Conversion Disorder} Alternate index label; see \textit{Functional neurologic disorder/conversion disorder}.

\medterm{Conversion Disorder (Functional Neurological Symptom Disorder)} A condition in which neurological symptoms, such as weakness, abnormal movements, sensory loss, or seizures, are not explained by a recognized structural neurological disease and cause significant impairment.

\synonyms
Functional Neurological Symptom Disorder

\medterm{Convulsion} A sudden episode of involuntary, often rhythmic muscle contractions that may occur with a seizure or other disturbance of brain function.

\synonyms
Convulsive episode; seizure-related motor activity.

\medterm{Cooking Without Salt} Preparing food with little or no added sodium to reduce dietary sodium intake, often useful in blood-pressure or fluid-retention management when combined with an individualized eating plan.

\medterm{Cooley'S Anemia} Alternate index label; see \textit{Thalassemia}.

\medterm{Coombs Test} A laboratory test that detects antibodies or complement attached to red blood cells, or
antibodies in serum that can react with red cells. It is used in evaluating immune
hemolysis and blood compatibility.

\medterm{Cooper'S Ligament} Fibrous connective-tissue septa extending from the deep surface of the breast toward the
skin and pectoral fascia. Tumor invasion or fibrosis may cause skin retraction.

\medterm{Coordination} The ability to control movements smoothly, accurately, and in the correct order. It depends on the brain, nerves, muscles, vision, and balance and may be tested with finger-to-nose, heel-to-shin, or rapid alternating movements.
movements.

\medterm{COPD} Chronic obstructive pulmonary disease is a progressive, largely irreversible airflow
limitation caused by chronic bronchitis, emphysema, or both, most commonly due to
cigarette smoking. It presents with chronic cough, sputum production, and exertional
dyspnea. Spirometry shows a reduced FEV1/FVC ratio.

\medterm{COPD - Control Drugs} Long-term medicines used to reduce COPD symptoms and exacerbations, such as long-acting bronchodilators and, for selected patients, inhaled corticosteroids.


\medterm{COPD - Quick-Relief Drugs} Fast-acting inhaled bronchodilators used for short-term relief of COPD breathlessness or bronchospasm; frequent need may indicate inadequate control or an exacerbation.


\medterm{COPD And Other Health Problems} The interaction between chronic obstructive pulmonary disease and conditions such as cardiovascular disease, osteoporosis, depression, diabetes, and lung cancer, which can complicate symptoms and care.

\medterm{COPD Exacerbation} Acute worsening of respiratory symptoms beyond normal day-to-day variation in a patient
with COPD that leads to additional therapy. Common triggers include viral or bacterial
infection, air pollution, and cardiopulmonary comorbidity.

\medterm{COPD Flare-Ups} Acute worsening of breathlessness, cough, or sputum beyond usual daily variation in chronic obstructive pulmonary disease, often triggered by infection, pollution, or other irritants.

\medterm{COPD Rehabilitation} A supervised program combining exercise, breathing strategies, education, nutrition, and support to help people with COPD breathe and function better, remain active, and reduce the effects of flare-ups.







\medterm{Copper} An essential trace element required for iron transport, connective-tissue cross-linking,
energy metabolism, neurotransmitter synthesis, and antioxidant enzymes. Deficiency can
cause anemia, neutropenia, impaired neurologic function, and bone abnormalities; excess
is toxic and accumulates in Wilson disease.

\medterm{Copper In Diet} Copper obtained from foods such as shellfish, nuts, seeds, beans, whole grains, and organ meats. The body needs small amounts, while excessive supplements or exposure can be harmful; intake should be tailored for deficiency or disease.

\medterm{Copper Poisoning} Toxic exposure to copper or copper salts that may cause vomiting, abdominal pain, hemolysis, liver injury, kidney injury, or shock.

\medterm{Coprolalia} Involuntary utterance of obscene, taboo, or socially inappropriate words, most often described as a complex vocal tic in Tourette syndrome.

\medterm{Copy Number Variant} A deletion or duplication that changes the number of copies of a DNA segment; its clinical significance depends on the region, size, genes involved, inheritance, and available evidence.

\medterm{Copy Number Variation} A deletion or duplication of a section of DNA that changes how many copies are present. Some variations are harmless, while others alter gene activity and cause developmental, neurologic, or other disorders.

\medterm{CoQ10} A fat-soluble substance made by the body that helps mitochondria produce energy and also acts as an antioxidant. It is sold as a supplement, but benefit varies by condition and it can interact with medicines.
\medterm{Cor} A Latin anatomical word meaning heart, appearing in medical terms that describe the heart, heart-related structures, or diseases affecting circulation.

\medterm{Cor Pulmonale} Enlargement and weakening of the right side of the heart caused by high pressure in the lung arteries, usually from long-term lung disease. It may cause breathlessness, leg swelling, enlarged neck veins, or liver congestion.

\medterm{Cor Triatriatum} A rare birth defect in which a membrane divides the left upper heart chamber into two spaces and may obstruct blood returning from the lungs. Symptoms range from none to breathlessness, poor growth, or heart failure.

\medterm{Coracoid Process} A hook-shaped bony projection from the shoulder blade that points forward. Several muscles and ligaments attach to it, and injury may occur with a direct blow or shoulder dislocation.

\medterm{Cord} A long, narrow structure such as the spinal cord, which contains the central nervous tissue, or the umbilical cord, which connects a fetus to the placenta.

\medterm{Cord Blood Testing} Laboratory testing of blood collected from the umbilical cord at birth, which may include blood type, complete blood count, blood gases, glucose, or tests for infection or inherited disease when clinically indicated.

\medterm{Cord Blood Stem Cells} Blood-forming stem cells collected from the umbilical cord and placenta after birth. They can make different types of blood and immune cells and may be used for stem-cell transplantation.

\medterm{Cord Presentation} A position in which the umbilical cord lies between the baby's presenting part and the cervix while the membranes are intact. If the membranes break, the cord may slip down and become compressed, threatening the baby's oxygen supply.

\medterm{Cordocentesis} A prenatal procedure using ultrasound guidance to take a small blood sample from the umbilical cord. It can assess fetal anemia, infection, or selected genetic conditions but carries risks including bleeding and pregnancy loss.

\medterm{Cordotomy} A procedure that cuts selected pain-carrying nerve fibers in the spinal cord to relieve severe pain that has not responded to other treatment. It can cause weakness, numbness, or other neurologic complications.

\medterm{Core Biopsy} Removal of one or more cylindrical tissue samples with a hollow cutting needle for microscopic examination. It preserves more tissue architecture than fine-needle aspiration and is used for breast, liver, prostate, and other lesions.

\medterm{Core Needle Biopsy} Removal of small cylinders of tissue with a hollow cutting needle for laboratory examination. Imaging may guide the needle, and the samples preserve tissue structure better than fluid or cell aspiration alone.

\medterm{Cori Cycle} A process in which muscles send lactate to the liver, where it is changed into glucose. The glucose returns to the blood and can be used again by muscles for energy.
\medterm{Cori Disease} A rare inherited disorder in which the body cannot break down stored glycogen normally. It can cause low blood sugar, an enlarged liver, muscle weakness, and sometimes heart disease.
\synonyms
Glycogen Storage Disease Type III (GSD III); Forbes Disease; Debrancher Deficiency

\medterm{Corn} A localized thickening of skin caused by repeated pressure or friction, often occurring on the toes or feet.

\synonyms
Clavus; heloma.

\medterm{Cornea} The clear, curved front surface of the eye. It protects the eye and bends light so images focus on the retina; scratches, infection, swelling, or scarring can cause pain and blurred vision.

\medterm{Corneal Abrasion} A scratch or loss of the outer layer of the cornea, often from a fingernail, foreign body, contact lens, makeup tool, or other injury. It can cause pain, tearing, redness, light sensitivity, and blurred vision.

\medterm{Corneal Cross-Linking} A treatment that uses riboflavin and ultraviolet light to strengthen the cornea. It can slow or stop worsening thinning and bulging, especially in keratoconus, but may cause pain during healing.

\medterm{Corneal Disease} Corneal disease affects the transparent front surface of the eye through infection, inflammation, trauma, dystrophy, degeneration, dryness, or abnormal shape. Pain, redness, light sensitivity, discharge, or reduced vision can signal urgent disease because scarring may impair sight.

\medterm{Corneal Edema} Swelling of the cornea from excess fluid, making it cloudy and causing blurred vision, glare, halos, or pain. Causes include damage to its inner cells, glaucoma, inflammation, injury, or eye surgery.

\medterm{Corneal Graft (Corneal Transplant)} Surgical replacement of part or all of a diseased or damaged cornea with donor corneal tissue to restore vision or relieve pain.

\synonyms
Corneal Transplant

\medterm{Corneal Implant} An implanted device or tissue substitute placed in or on the cornea to restore vision,
replace diseased tissue, or alter corneal shape.

\medterm{Corneal Inflammation (Keratitis)} Alternate index label for Keratitis.

\medterm{Corneal Injury} Damage to the clear front surface of the eye, including scratches, foreign bodies, chemical burns, infections, or cuts that may threaten vision.

\medterm{Corneal Transplant} Surgery that replaces part or all of a damaged or diseased cornea with healthy donor tissue. It may restore vision or relieve pain; risks include rejection, infection, increased eye pressure, and need for further treatment.


\medterm{Corneal Ulcer} A corneal ulcer is a defect in the corneal epithelium with associated stromal
inflammation and necrosis, typically caused by infectious or non-infectious mechanisms,
and it is a sight-threatening condition that requires prompt diagnosis and treatment.

\medterm{Corneal Ulcers And Infections} Infection or inflammation of the cornea that may produce an epithelial defect or ulcer, pain, redness, tearing, discharge, and light sensitivity. Bacterial, fungal, protozoal, and viral causes can threaten vision and require urgent ophthalmic assessment.

\medterm{Corneitis} Inflammation of the cornea, often producing pain, photophobia, tearing, and impaired
vision. It may be infectious, immune-mediated, traumatic, or related to exposure.

\medterm{Cornelia De Lange Syndrome} A genetic disorder caused by mutations in genes of the cohesin complex (NIPBL 60\%,
SMC1A, SMC3, RAD21, HDAC8).

\medterm{Corns} A corn is a localized thickening of skin caused by repeated pressure or friction, commonly on toes or other weight-bearing areas. Proper footwear, pressure relief, and treatment of the underlying deformity help; diabetes or poor circulation makes self-removal risky.

\medterm{Corns And Calluses} Corns are small, localized areas of thickened skin caused by repeated pressure or friction,
usually on the toes. Calluses are broader areas of thickened, hardened skin, usually on the
soles of the feet or palms of the hands. Both are usually harmless but may become painful.

\medterm{Coronal} Relating to a plane that divides the body into front and back portions, or to the crown or top of a structure.

\medterm{Coronal Caries} Tooth decay affecting the crown, the visible part of a tooth above the gum. It can begin in pits, grooves, or smooth surfaces when bacterial acids dissolve enamel and dentin.

\medterm{Coronal Plane} A vertical plane that divides the body or an organ into front (anterior) and back
(posterior) portions.

\medterm{Coronary} Relating to the arteries and blood flow that supply the heart muscle. Narrowing or blockage of these arteries can reduce oxygen to the heart and cause chest pain or a heart attack.

\synonyms
Cardiac blood-vessel-related; heart-coronary.

\medterm{Coronary Aneurysm} An abnormal widening or ballooning of a coronary artery. It may result from atherosclerosis, inflammation, infection, injury, or a congenital condition and can cause clotting, blockage, or rarely rupture.

\medterm{Coronary Angiography} An X-ray examination of the heart arteries after contrast is injected through a thin catheter. It shows narrowed or blocked areas and can guide treatment, but the procedure can cause bleeding, stroke, kidney injury, or other complications.

\medterm{Coronary Arteries} The blood vessels that arise from the aorta and supply oxygen to the heart muscle. The right and left arteries divide into branches, and blockage in one can cause angina or a heart attack.

\medterm{Coronary Artery} One of the blood vessels branching from the aorta that supplies oxygen-rich blood
directly to the heart muscle.

\medterm{Coronary Artery Aneurysm} A localized dilatation of a coronary artery exceeding the normal vessel diameter by
approximately 1.5 times. It may be caused by atherosclerosis, Kawasaki disease,
connective-tissue disorders, infection, trauma, or an iatrogenic procedure, and can lead
to thrombosis, distal embolization, myocardial ischemia, or rupture.

\medterm{Coronary Artery Anomaly} A congenital abnormality in the origin, course, structure, or termination of a coronary artery. Some anomalies are incidental, whereas others can impair myocardial perfusion or increase the risk of ischemia, arrhythmia, or sudden cardiac death.

\medterm{Coronary Artery Bypass} A surgical procedure that creates a new route for blood flow around a narrowed or
blocked coronary artery using a vascular graft.

\medterm{Coronary Artery Bypass Grafting} Surgery that uses a blood-vessel graft to create a route around severely narrowed or blocked heart arteries. It can improve blood flow and symptoms and may improve survival in selected people with extensive disease.

\medterm{Coronary Artery Bypass Surgery} An operation that uses a blood vessel from another part of the body to create a new route around narrowed or blocked heart arteries and improve blood flow to heart muscle.
\synonyms
CABG; coronary bypass surgery.

\medterm{Coronary Artery Disease} Narrowing or blockage of the arteries supplying the heart, usually from fatty plaque buildup. It may cause exertional chest pressure, heart attack, heart failure, abnormal rhythms, or sudden death.

\medterm{Coronary Artery Fistula} An abnormal connection between a coronary artery and a heart chamber or vessel that may cause ischemia, failure, arrhythmia, or endocarditis.

\medterm{Coronary Artery Spasm} Temporary tightening of a heart artery that reduces blood flow and may cause chest pain, often at rest. Smoking, stimulant drugs, stress, or other vessel problems can trigger it.

\medterm{Coronary Calcium Score} A CT measurement of calcium in the arteries supplying the heart. More calcium usually means more plaque and higher future cardiovascular risk, but a zero score does not exclude soft plaque or eliminate all risk.

\medterm{Coronary Calcium Scoring} A computed tomography scan measuring calcium-containing plaque in heart arteries. The score helps estimate future cardiovascular risk in selected people but does not show all plaque or replace clinical assessment.

\medterm{Coronary Care Unit} A hospital unit for people with serious or unstable heart conditions. Staff continuously monitor heart rhythm, blood pressure, oxygen, and other vital functions and provide urgent treatment for problems such as heart attack, dangerous rhythms, or heart failure.
\synonyms
CCU

\medterm{Coronary Circulation} The network of arteries and veins that supplies the heart muscle with oxygen and removes waste. Blood flow increases when the heart works harder and falls when an artery is narrowed or blocked.

\medterm{Coronary Computed Tomography Angiography} A noninvasive CT scan performed after injected contrast highlights the heart arteries. It can show plaque, narrowing, or blockage and is useful for selected people with suspected coronary artery disease.

\synonyms
CCTA; Coronary CT Angiography; Cardiac CTA

\medterm{Coronary Dominance} The pattern describing which heart artery gives rise to the vessel supplying the back part of the heart. Most people have right dominance; left or shared patterns are normal variations.

\medterm{Coronary Endarterectomy} Surgery removing hardened plaque from the inner lining of a heart artery. It may improve blood flow in selected complex blockages but is uncommon and carries risks of heart attack, bleeding, or vessel injury.

\synonyms
Coronary plaque endarterectomy; coronary endarterectomy procedure.

\medterm{Coronary Heart Disease} Disease caused by reduced blood flow through the coronary arteries, usually due to
atherosclerotic plaque. It may cause angina, myocardial infarction, heart failure, or
sudden cardiac death. Also called coronary artery disease.

\synonyms
coronary artery disease.

\medterm{Coronary Insufficiency} Inadequate blood flow through the coronary circulation to meet myocardial oxygen demand,
usually because of coronary artery disease.

\medterm{Coronary Microvascular Angina} Angina caused by impaired regulation or constriction of the small coronary vessels, often producing ischemic symptoms despite no flow-limiting epicardial coronary obstruction. Assessment may include evaluation for coronary microvascular dysfunction; management is individualized according to symptoms, cardiovascular risk, and the underlying mechanism.

\medterm{Coronary Microvascular Disease} Alternate index label; see \textit{Small vessel disease}.

\medterm{Coronary Microvascular Dysfunction} Abnormal function of the small heart vessels that limits blood flow even when the major arteries are not severely narrowed. It may cause chest pain, especially in women and people with diabetes or high blood pressure.

\medterm{Coronary Occlusion} Obstruction of a coronary artery, commonly by atherosclerotic plaque and thrombus, that
can cause myocardial ischemia or infarction.

\medterm{Coronary Sinus} A large vein on the back of the heart that collects deoxygenated blood from the heart
muscle and drains it into the right atrium.

\medterm{Coronary Slow Flow Phenomenon} Delayed movement of blood through the small heart vessels despite no major blockage on angiography. It may cause repeated chest pain at rest or with activity and can occasionally lead to abnormal rhythms or heart injury.

\synonyms
CSFP; Slow Coronary Flow; Primary Coronary Slow Flow

\medterm{Coronary Spasm} Temporary contraction of a coronary artery that reduces blood flow to the heart muscle and may cause angina or infarction.

\synonyms
Vasospastic angina mechanism; coronary artery spasm.

\medterm{Coronary Steal Phenomenon} Reduced blood flow to an already poorly supplied part of the heart when nearby healthy vessels widen during certain stress tests or treatments. This can worsen chest pain or reveal hidden coronary disease.

\synonyms
Coronary Steal Syndrome; Myocardial Steal

\medterm{Coronary Stent} A small mesh tube placed in a narrowed heart artery to help keep it open after angioplasty. Medicines are usually needed afterward to reduce the risk of clotting.
\medterm{Coronary Thrombosis} Formation of a blood clot in a coronary artery, potentially obstructing blood flow and causing
myocardial infarction.

\medterm{Coronavirus} A group of viruses that infect people and animals. Human coronaviruses can cause respiratory illness ranging from a mild cold to severe pneumonia, depending on the virus and the person.
\medterm{Coronavirus Disease 2019} A respiratory and systemic infectious disease caused by SARS-CoV-2. Illness ranges from
asymptomatic infection to pneumonia, respiratory failure, and multisystem complications.

\medterm{Corporeal} Relating to the physical body rather than thoughts, feelings, or the spirit. In medical writing, the word may describe a bodily symptom, physical examination, or treatment directed at the body.

\medterm{Corpse} The dead body of a human being; the subject of medicolegal examination. Forensic
assessment of a corpse includes identification, determination of death and time of death
(rigor mortis, livor mortis, algor mortis, vitreous potassium, gastric emptying),
external and internal examination, and collection of evidence.

\medterm{Corpus} Latin for body; in anatomy it refers to the main body of an organ or structure, such as the corpus callosum or corpus uteri.

\medterm{Corpus Callosum} Brain's largest white matter tract with approximately 200 million axons connecting left
and right hemispheres. Four parts: rostrum, genu, body, splenium. Facilitates
interhemispheric communication for sensory, motor, and cognitive integration.

\medterm{Corpus Cavernosum} One of two columns of erectile tissue in the penis that fill with blood during erection.

\synonyms
Cavernous body; penile corpus cavernosum.

\medterm{Corpus Luteum} The temporary endocrine structure formed from an ovarian follicle after ovulation; it produces progesterone during the luteal phase and early pregnancy.

\synonyms
Yellow body; luteal body.

\medterm{Corpus Luteum Cyst} A usually harmless ovarian cyst that forms after ovulation when the normal hormone-producing structure fills with fluid or blood. It often disappears within a few menstrual cycles but may cause pelvic pain, delayed menstruation, or bleeding if it ruptures.

\medterm{Corpus Spongiosum} A column of erectile tissue surrounding the penile urethra and expanding distally to form the glans.

\synonyms
Spongy body; corpus spongiosum penis.


\medterm{Corrected Age} The developmental age of a premature baby after allowing for the weeks born before 40 weeks of pregnancy. It helps interpret growth and milestones, usually during the first two years.

\medterm{Corrigan Pulse} A strong pulse that rises quickly and then collapses rapidly. It is classically associated with aortic-valve leakage but can occur with other high-flow states and requires clinical evaluation.

\medterm{Corrosive} A chemical that destroys tissue on contact, such as a strong acid or alkali. Exposure can burn skin, eyes, airways, or the digestive tract; prompt prolonged rinsing and urgent medical advice may be needed.

\medterm{Corrugator Muscle} A small facial muscle that pulls the eyebrows downward and inward, helping create frowning and other concentrated facial expressions.

\synonyms
Corrugator supercilii; glabellar muscle.

\medterm{Cortex} The outer layer of an organ or structure. Examples include the brain's cerebral cortex, the outer part of the adrenal gland, and the outer layer of a hair.

\synonyms
Cortical layer; outer cortex.

\medterm{Corti'S Ganglion} A group of nerve-cell bodies inside the cochlea that carries hearing signals from the inner ear to the brain.

\medterm{Cortical} Relating to the cortex, the outer layer of an organ such as the brain, kidney, or adrenal gland.

\medterm{Cortical Bone} Dense bone forming the outer shell of most bones and providing much of the skeleton's mechanical strength.

\synonyms
Compact bone; cortical tissue.

\medterm{Cortical Cataract} A cataract that begins in the outer part of the eye's lens and forms spoke-like cloudy areas extending inward. It can cause glare, halos, and blurred vision, especially in bright light.

\medterm{Corticobasal Degeneration} A rare, progressive brain disorder affecting regions involved in movement and higher brain functions. It usually affects one side more than the other and may cause stiffness and slow movement, difficulty carrying out learned movements (apraxia), involuntary jerks, abnormal muscle contractions (dystonia), reduced sensation, or an ``alien limb'' that feels or acts outside the person’s control.

\medterm{Corticosteroid} A steroid hormone from the adrenal cortex or a synthetic analogue; glucocorticoids reduce inflammation and immune activity, while mineralocorticoids regulate salt and water balance.

\medterm{Corticosteroids} Medicines similar to hormones made by the adrenal glands that reduce inflammation and suppress immune activity. They can help many conditions but may cause infection, high blood sugar, mood changes, bone loss, or adrenal suppression.

\medterm{Corticosteroids Mechanism} These medicines enter cells and bind hormone receptors that change gene activity. This reduces production of inflammatory signals and limits the activity and movement of several immune cells.

\medterm{Corticosteroids Overdose} Overdose of corticosteroids that may cause gastrointestinal symptoms, mood changes, hyperglycemia, fluid retention, infection risk, or adrenal suppression.

\medterm{Corticotropin-Releasing Factor} A hypothalamic hormone that stimulates the pituitary gland to release adrenocorticotropic hormone and helps coordinate the stress response.

\synonyms
Corticotropin-releasing hormone; CRH.

\medterm{Cortisol} A hormone made by the adrenal glands that helps the body respond to stress, maintain blood pressure and blood sugar, regulate metabolism, and control inflammation. Too much or too little can cause serious illness.

\medterm{Cortisol Blood Test} A blood test measuring cortisol, a hormone produced by the adrenal cortex. Results help evaluate adrenal insufficiency or cortisol excess, but vary substantially with time of day, stress, illness, medicines, and the assay used.

\medterm{Cortisol Urine Test} A test measuring cortisol excreted in a timed urine collection, commonly 24-hour urine free cortisol, used to evaluate suspected excess cortisol production such as Cushing syndrome.

\medterm{Cortisone} A corticosteroid that is converted in the body to cortisol and can be used to reduce inflammation or replace glucocorticoid activity.

\medterm{Coryza} Acute inflammation of the nasal mucous membranes with nasal discharge and obstruction,
commonly occurring in the common cold or other upper-respiratory infection.

\medterm{Cosmeceutical} A cosmetic product marketed as having biologically active, drug-like benefits; the term has no uniform formal regulatory category in many jurisdictions.

\medterm{Cosmetic Allergies} Cosmetic products can cause irritant dermatitis or allergic contact dermatitis, producing itching, redness, swelling, scaling, or blistering where applied. Stopping the suspected product and medical assessment, sometimes with patch testing, help identify the responsible ingredient.


\medterm{Cosmetic Dermatology} Dermatologic care intended to change appearance rather than treat disease, including procedures for wrinkles, scars, pigmentation, hair, or facial volume. Risks include infection, scarring, pigment changes, allergic reactions, and injury from improperly performed treatment.

\medterm{Cosmetic Ear Surgery} Surgery to change the size, shape, or position of the outer ear, often to correct prominent ears or a structural difference. Risks include bleeding, infection, scarring, altered sensation, and an unsatisfactory result.


\medterm{Costello Syndrome} A rare genetic condition, usually caused by an HRAS gene change, affecting growth, development, skin, heart, and muscles. It may cause feeding difficulty, short stature, distinctive facial features, heart disease, and increased risk of certain tumors.

\medterm{Costochondritis} Costochondritis is inflammation or irritation where ribs join the breastbone, causing localized chest pain often worse with movement or pressing the area.

\synonyms
Chest Wall Pain
Costosternal Chondrodynia
Costosternal Syndrome

\medterm{Costochondritis (Costochondritis And Tietze Syndrome)} Alternate index label for Costochondritis and Tietze Syndrome.

\medterm{Costochondritis And Tietze Syndrome} Costochondritis is inflammation or irritation at the junction of ribs and the breastbone, causing reproducible chest-wall pain. Tietze syndrome is a less common related condition with visible swelling at one or a few joints; heart and lung emergencies must still be excluded.

\synonyms
Costochondritis (Costochondritis And Tietze Syndrome)
Tietze Syndrome (Costochondritis And Tietze Syndrome)

\medterm{Costosternal Chondrodynia} Alternate index label; see \textit{Costochondritis}.

\medterm{Costosternal Syndrome} Alternate index label; see \textit{Costochondritis}.

\medterm{Cot Death} Sudden infant death syndrome, the sudden unexplained death of an apparently healthy
infant, usually during sleep, that remains unexplained after investigation.

\medterm{Cough} Cough is a protective reflex that clears the airways; its duration, triggers, sputum, associated symptoms, and examination findings help distinguish infectious, inflammatory, cardiac, and other causes.
foreign material from the tracheobronchial tree.

\medterm{Cough Augmentation} Techniques to improve cough effectiveness in patients with neuromuscular disease or
chest wall weakness. Methods include manually assisted cough (abdominal thrust timed
with cough), mechanical insufflation-exsufflation (CoughAssist device), and breath
stacking with a resuscitation bag.

\medterm{Cough Chronic (Chronic Cough)} Alternate index label for Chronic Cough.

\medterm{Cough Headaches} Head pain triggered by coughing, sneezing, straining, laughing, or another sudden pressure increase. New, severe, or persistent attacks need assessment because some are caused by a structural brain problem.
\medterm{Cough In Adults} A cough in an adult may be acute, subacute, or chronic and may result from infection, asthma, chronic obstructive pulmonary disease, reflux, upper-airway cough syndrome, medication, or other disease. Coughing blood, breathing difficulty, chest pain, low oxygen, or persistent unexplained cough requires medical assessment.

\medterm{Cough Suppressant} An antitussive medicine or measure that reduces the urge or reflex to cough; treatment should address the underlying cause and avoid suppressing a productive cough when clearance is needed.

\medterm{Coughing Up Blood} Expectoration of blood or blood-streaked sputum from the lower respiratory tract, called hemoptysis; heavy bleeding or breathing difficulty is an emergency.

\medterm{Could A Stiff Neck Be A Sign Of Something Serious} Neck stiffness commonly follows muscle strain, but stiffness with fever, severe headache, light sensitivity, confusion, rash, weakness, trauma, or neurologic change can signal meningitis, bleeding, spinal injury, or another emergency.


\medterm{Count Rate} The number of detected events each second during measurement. In medical imaging, very high rates can cause missed events because equipment needs recovery time, reducing accuracy unless corrected.

\medterm{Countercurrent Multiplication} A kidney process that creates a concentrated salt gradient deep in the kidney. This gradient allows the collecting ducts to reclaim more water when antidiuretic hormone is present.

\medterm{Counterpulsation} Timed inflation and deflation of a balloon or pressure cuffs with each heartbeat to improve blood flow to the heart or reduce the heart's workload. It is used only in selected serious cardiac conditions.

\medterm{Countertransference} The therapist's emotional reactions to the patient, which may be influenced by the
therapist's own unresolved conflicts or may be a response to the patient's transference.
Awareness of countertransference is essential for effective therapy and can provide
valuable clinical information about the patient's interpersonal style.

\medterm{Counting Carbohydrates} Estimating the grams of carbohydrate in foods and drinks to help plan meals, match insulin, and manage blood glucose, especially in diabetes.

\medterm{Coup-Contrecoup Injury} Brain injury in which the brain is damaged both beneath the point of impact and on the opposite side as it moves inside the skull. Falls and vehicle crashes are common causes.

\medterm{Courvoisier Sign} A painless, enlarged gallbladder in a person with jaundice, traditionally suggesting blockage of the bile duct by a tumor rather than a gallstone. It is a clue, not a diagnosis.

\medterm{Covered Stent (Vascular)} A vascular stent lined with a graft material that supports the vessel and excludes a leak or
aneurysm from the bloodstream.

\synonyms
Vascular

\medterm{COVID Toe} A chilblain-like skin lesion reported in association with COVID-19, usually affecting
the toes or fingers and appearing red, purple, or swollen.

\medterm{COVID-19} Coronavirus disease caused by infection with SARS-CoV-2. Illness ranges from asymptomatic or mild disease to pneumonia, respiratory failure, thrombosis, or death; common symptoms include fever, cough, shortness of breath, fatigue, muscle aches, and altered taste or smell. Some people develop persistent symptoms after the acute infection.

\medterm{COVID-19 Antibody Test} A blood test that detects antibodies produced in response to SARS-CoV-2 infection or vaccination; it does not reliably diagnose current infection or establish complete immunity.

\medterm{COVID-19 In Pregnancy} COVID-19 during pregnancy can be more severe than in nonpregnant adults and increases risks such as hospitalization, preterm birth, and stillbirth. Fetal infection before birth is uncommon but possible; vaccination and prompt care reduce risk.


\medterm{COVID-19 Symptoms} Symptoms of SARS-CoV-2 infection, which may include fever, cough, sore throat, congestion, fatigue, muscle aches, headache, loss of smell or taste, vomiting, or diarrhea.

\medterm{COVID-19 Test} A test for current SARS-CoV-2 infection using a respiratory sample. Molecular tests detect viral genetic material and antigen tests detect viral proteins; results depend on timing, sample quality, and the amount of virus present.


\medterm{COVID-19 Vaccines} Vaccines that prime immune responses against SARS-CoV-2 and reduce severe disease, hospitalization, and death; products, schedules, effectiveness, and recommendations vary by age and risk.

\medterm{COVID-19 Vaccines For Children Ages 6 Months And Older} COVID-19 vaccination guidance for children at least six months old, with product and dose schedules determined by age, prior doses, immune status, and current public-health recommendations.

\medterm{COVID-19 Virus Test} A test for current SARS-CoV-2 infection, usually using a nucleic acid amplification test or antigen test on a respiratory specimen; results must be interpreted with symptoms, timing, and test performance.



\medterm{Cow'S Milk - Infants} Guidance on cow’s milk in infancy: it should not replace breast milk or iron-fortified formula before 12 months because of low iron content and risks of renal solute load and intestinal blood loss.

\medterm{Cow'S Milk And Children} Nutritional guidance on cow’s milk for children, balancing protein, calcium, vitamin D, fat, iron, allergy, and total intake so milk does not displace a varied diet.

\medterm{Cowden Syndrome} An inherited condition, usually caused by a PTEN gene change, leading to multiple noncancerous growths and increased risks of breast, thyroid, uterine, kidney, and colorectal cancers. Regular specialist screening is important.

\medterm{Cowper'S Gland} One of a pair of small male accessory sex glands below the prostate. Its secretion
contributes mucus and fluid to the urethra before ejaculation.

\medterm{Cowpox} A zoonotic viral infection caused by cowpox virus or related orthopoxviruses, producing
localized pustular or ulcerative skin lesions and sometimes systemic illness.

\medterm{Cox-2} Cyclooxygenase-2, an inducible enzyme that produces inflammatory prostaglandins; selective COX-2 inhibitors reduce pain and inflammation but may have cardiovascular, renal, and gastrointestinal risks.

\medterm{COX-2 Inhibitors} Medicines that reduce the action of an enzyme involved in pain and inflammation. They may cause less stomach irritation than some older anti-inflammatory medicines, but can still cause bleeding, kidney problems, high blood pressure, or heart risks.

\synonyms
Cyclooxygenase-2 inhibitors; selective COX-2 inhibitors.

\medterm{Coxa Vara} A deformity in which the angle between the femoral neck and shaft is abnormally reduced,
potentially causing a limp, shortened limb, or altered hip mechanics.

\medterm{Coxalgia} Pain arising from the hip or nearby tissues, often felt in the groin and sometimes spreading to the thigh or knee. Causes vary with age and include infection, injury, growth problems, arthritis, cartilage damage, stress fractures, and reduced blood supply.

\medterm{Coxiella Burnetii} A bacterium that causes Q fever, usually spread by breathing dust contaminated by infected farm animals. It can cause fever, headache, hepatitis, or pneumonia and may persist as a serious infection of heart valves.

\medterm{Coxsackie (Coxsackie Virus)} Alternate index label for Coxsackie Virus.

\medterm{Coxsackie Virus Type A16} A virus that commonly causes hand, foot, and mouth disease, especially in children. Fever, painful mouth sores, and a rash or blisters on the hands and feet usually resolve without serious complications.

\medterm{Coxsackievirus} A group of enteroviruses that cause a range of human diseases. Group A viruses commonly
cause hand-foot-and-mouth disease and herpangina, whereas group B viruses are associated
with myocarditis, pericarditis, pleurodynia, and aseptic meningitis.

\medterm{CPAP} Continuous positive airway pressure delivered through a mask to splint open the upper airway during sleep. It is
commonly used for obstructive sleep apnea and can reduce apneas and related symptoms
when used consistently.

\medterm{CPAP Machine} A device that sends a steady stream of air through a mask to keep the upper airway open during sleep. It is commonly used for obstructive sleep apnea; benefit depends on correct mask fit and regular use.

\medterm{CpG} A DNA sequence in which cytosine is followed by guanine on the same strand, often occurring in regions whose methylation helps regulate gene expression.

\medterm{CPK Isoenzymes Test} A blood test that separates forms of creatine kinase to help identify whether increased enzyme levels come from skeletal muscle, heart muscle, or brain. Troponin testing is generally preferred for suspected heart attack.

\medterm{CPR} Cardiopulmonary resuscitation is an emergency procedure using chest compressions and, when appropriate, rescue breaths or an automated external defibrillator to support circulation and breathing during cardiac arrest.

\medterm{CPR - Infant} Cardiopulmonary resuscitation for an unresponsive infant who is not breathing normally, using emergency activation, infant chest compressions, rescue breaths when trained, and an AED according to local guidance.

\medterm{CPR - Young Child (Age 1 Year To Onset Of Puberty)} Cardiopulmonary resuscitation for a child from age one year until puberty who is unresponsive and not breathing normally, using emergency activation, compressions, breaths, and an AED as available.

\medterm{CPT Code For Foreign Body Removal From The Ear} Removal of an ear foreign body is coded according to the documented method, location, anesthesia, complications, and whether a separate evaluation was performed. Objects causing pain, bleeding, hearing loss, or suspected eardrum injury should be removed by a trained clinician.

\medterm{CPVT} Alternate name; see \textit{Catecholaminergic Polymorphic Ventricular Tachycardia}.

\medterm{Crabs} Alternate index label; see \textit{Pubic lice (crabs)}.

\medterm{Cracked Tooth Syndrome} A partial crack in a tooth that may cause sharp pain when biting or when pressure is released, and sensitivity to hot or cold. The crack can be difficult to locate and may worsen over time. Dental treatment depends on its depth and whether the inner pulp is involved.

\medterm{Crackles (Rales)} Discontinuous, non-musical adventitious lung sounds produced by the sudden explosive
opening of collapsed small airways and alveoli during inspiration or the bubbling of air
through fluid in the bronchioles.

\synonyms
Rales

\medterm{Cradle Cap} A common form of seborrheic dermatitis in infants, causing greasy or flaky scales on the
scalp with mild redness. It is usually harmless and improves with gentle scalp care;
persistent inflammation or signs of infection require assessment.

\medterm{Cramp, Muscle} Alternate index label; see \textit{Muscle cramp}.

\medterm{Cramping Leg Pain} Pain caused by involuntary, often painful muscle contraction in the leg; contributors include exertion, dehydration, electrolyte abnormalities, medications, nerve disease, and impaired circulation.

\medterm{Cramps But No Period} Pelvic cramps without menstruation can result from ovulation, pregnancy, implantation, constipation, urinary disease, infection, endometriosis, cysts, or other pelvic conditions. Severe one-sided pain, faintness, fever, vomiting, or possible pregnancy requires urgent assessment.

\medterm{Cramps Heat (Heat Cramps)} Alternate index label for Heat Cramps.

\medterm{Cramps Menstrual (Menstrual Cramps)} Alternate index label for Menstrual Cramps.

\medterm{Cramps, Menstrual} Alternate index label; see \textit{Menstrual cramps}.

\medterm{Cranial} Relating to the skull or toward the head, as opposed to caudal or toward the lower end of the body.

\medterm{Cranial Arteritis} Inflammation of arteries, usually in the scalp and head, that can cause a new headache, scalp tenderness, jaw pain while chewing, fever, or vision loss. Sudden visual symptoms require emergency treatment.

\medterm{Cranial Cavity} The hollow space inside the skull that contains and protects the brain, its coverings, blood vessels, and cerebrospinal fluid. A rise in pressure inside this space can compress the brain and become life-threatening.

\medterm{Cranial Fasciitis} A usually noncancerous, rapidly growing lump in the scalp or nearby soft tissue, most often in infants and young children. It can resemble cancer on examination or imaging and may require biopsy.

\medterm{Cranial Implant} An artificial material or device placed in or on the skull to repair a defect, protect the brain, drain fluid, stimulate nerves, or support treatment. Risks include infection, bleeding, and implant failure.

\medterm{Cranial Mononeuropathy III} Damage to the third nerve controlling several eye muscles and the eyelid. It can cause double vision, a drooping eyelid, an eye that points abnormally, or an enlarged pupil and needs evaluation.

\medterm{Cranial Mononeuropathy III - Diabetic Type} A third-nerve palsy associated with diabetes, often causing painful double vision and a drooping eyelid while the pupil remains normal. Other causes must be excluded, especially when the pupil is enlarged or symptoms are sudden.
\medterm{Cranial Mononeuropathy VI} Damage to the sixth cranial nerve, which controls the muscle that moves an eye outward. It can cause inward turning of the eye and horizontal double vision, especially when looking toward the affected side.

\medterm{Cranial Nerve} One of 12 pairs of nerves arising from the brain or brainstem. They control smell, vision, eye movement, facial sensation and movement, hearing, balance, swallowing, voice, and other head and neck functions.

\medterm{Cranial Nerve Palsy} Weakness or paralysis of one or more cranial nerves, causing deficits such as impaired eye
movement, facial weakness, hearing loss, swallowing difficulty, or tongue weakness. The
cause may be vascular, traumatic, compressive, inflammatory, infectious, or metabolic.

\medterm{Cranial Nerve X} The vagus nerve, the tenth cranial nerve, which carries parasympathetic and sensory fibers and contributes to swallowing, voice, and regulation of thoracic and abdominal organs.

\medterm{Cranial Nerves} The 12 pairs of nerves arising from the brain or brainstem that provide sensory, motor, and autonomic functions to the head, neck, and parts of the trunk.

\medterm{Cranial Sutures} Fibrous joints between skull bones that permit molding in infancy and accommodate brain growth before progressively fusing; premature fusion causes craniosynostosis.

\medterm{Craniectomy} Surgery removing part of the skull, sometimes leaving it off temporarily, to relieve dangerous brain pressure or access the brain. A later operation may replace the removed bone or use an implant.

\medterm{Craniocerebral Injury} An injury involving the scalp, skull, brain, or several of these structures, ranging from minor wounds to life-threatening traumatic brain injury.

\medterm{Craniofacial Abnormalities} Abnormalities of the skull, face, or related tissues that may be present at birth or develop later. Causes include genetic conditions, disrupted development, trauma, infection, and tumors.

\medterm{Craniofacial Dysostosis} A congenital disorder involving premature fusion of cranial sutures and characteristic
abnormalities of the face and skull, such as those seen in Crouzon syndrome.

\medterm{Craniopharyngioma} A usually slow-growing tumor near the pituitary gland and hypothalamus. It may cause headache, vision loss, excessive thirst or urination, hormone problems, growth changes, or increased pressure in the skull.

\medterm{Cranioplasty} Surgical repair or reconstruction of a skull defect using bone, an implant, or another
suitable material.

\medterm{Craniosynostosis} Craniosynostosis is premature fusion of one or more skull sutures, altering head shape and sometimes increasing pressure on the developing brain.

\synonyms
Congenital Plagiocephaly
Craniostenosis

\medterm{Craniosynostosis Repair} Surgery to separate prematurely fused skull bones and reshape the skull when needed. It aims to improve head shape and make room for brain growth; timing and technique depend on the child's age and anatomy.


\medterm{Craniotabes} Soft, thin, or pliable areas of an infant’s skull that indent under gentle pressure, commonly associated with normal newborn molding but also with rickets or other bone disease.

\medterm{Craniotomy} Surgery that temporarily removes a section of skull to reach the brain or its coverings. It may be used to remove a tumor, drain bleeding, repair a blood vessel, treat epilepsy, or relieve pressure.

\medterm{Cranium} The part of the skull that encloses and protects the brain, consisting of the cranial bones and their joints.

\medterm{Crapulence} An unpleasant state of sickness or discomfort caused by excessive eating or drinking, especially excessive alcohol use.

\medterm{Craving} A strong, often difficult-to-control desire for a substance, behavior, or experience, commonly seen in substance-use disorders.

\synonyms
Strong urge; compulsive desire.

\medterm{Crazy Paving Pattern} A lung-imaging pattern of hazy areas crossed by thick lines, resembling paving stones. It can occur with several conditions, including infection, bleeding, inflammation, or protein buildup, so it is not a diagnosis by itself.

\medterm{CRC} Colorectal cancer, malignant disease arising in the colon or rectum; screening and evaluation depend on age, risk, symptoms, and test results.

\medterm{CRE} Enterobacterales bacteria resistant to carbapenem antibiotics, important medicines for serious infections. They may colonize without illness or cause difficult-to-treat urinary, wound, abdominal, lung, or bloodstream infections.

\medterm{CRE Infection} Carbapenem-resistant Enterobacterales are bacteria resistant to important antibiotics and can cause urinary, bloodstream, abdominal, wound, or lung infections. Infection-control precautions, culture-based treatment, and specialist advice are important because therapeutic options may be limited.

\synonyms
Carbapenem Resistant Enterobacteriaceae (CRE Infection)

\medterm{Cream} A spreadable semisolid preparation applied to the skin, often containing a medicine. Creams are usually less greasy than ointments and may be useful on moist or broad areas, but the correct product depends on the skin condition.

\medterm{Creatine} A nitrogen-containing compound stored mainly in skeletal muscle that helps supply energy for short, high-intensity activity.

\medterm{Creatine Kinase} An enzyme found mainly in muscle and released into blood when muscle cells are damaged. Levels may rise after strenuous exercise, injury, seizures, muscle disease, or some medicines; other tests help identify the source.

\medterm{Creatine Phosphokinase Test} A blood test measuring creatine kinase, an enzyme released with muscle injury; elevated levels may occur in skeletal-muscle disease, strenuous exercise, rhabdomyolysis, or myocardial injury.

\medterm{Creating A Family Health History} Collecting a record of relatives’ diseases, ages at diagnosis, ancestry, and causes of death to identify inherited risk and guide screening, prevention, and genetic counseling.

\medterm{Creatinine} A breakdown product of muscle creatine phosphate produced at a constant rate, filtered
by the glomerulus and minimally reabsorbed; serum creatinine is the most used marker of
kidney function.

\medterm{Creatinine Blood Test} A blood test measuring serum creatinine, a waste product generated largely by skeletal muscle. It is used with age, body size, and other factors to estimate kidney filtration, but interpretation is affected by muscle mass, diet, medicines, and acute illness.

\medterm{Creatinine Clearance Test} A test estimating kidney filtration from creatinine in a timed urine collection and blood sample. It can be affected by incomplete urine collection and creatinine secretion, so estimated glomerular filtration rate is often preferred for routine assessment.

\medterm{Creatinine Test} A blood or urine test measuring creatinine to help assess kidney filtration and muscle-related creatinine production.

\synonyms
Creatinine measurement; renal function creatinine test.

\medterm{Creatinine Urine Test} A urine test measuring creatinine concentration or excretion, used to assess kidney handling, interpret other urine measurements, or judge the completeness of a timed collection. Results depend on muscle mass, diet, collection quality, and kidney function.

\medterm{Creeping Eruption} An intensely itchy winding skin track caused by larvae, usually hookworms from contaminated soil or sand, migrating within the superficial skin; it is also called cutaneous larva migrans.

\medterm{Crepitation (Crepitus)} A crackling, grating, or popping sensation or sound produced by movement of tissues, air, or bone. Causes include fractures, joint disease, subcutaneous air, and lung disease.

\synonyms
Crepitus

\medterm{Crepitus} A crackling, grating, or popping sensation or sound caused by movement of tissue, bone, joint surfaces, or air in soft tissue.

\synonyms
Crepitation; crackling sensation.

\medterm{Crescentic Glomerulonephritis} A rapidly worsening inflammation of the kidney's filtering units. It can cause blood or protein in urine, swelling, high blood pressure, and kidney failure and usually requires urgent specialist treatment.

\medterm{CREST Syndrome} A form of systemic sclerosis involving mainly the skin and certain organs. The name refers to calcium deposits, cold-sensitive fingers, swallowing difficulty, tight fingers, and small widened blood vessels; lung, heart, or digestive complications may occur.

\medterm{Creutzfeld-Jakob Disease} A rapidly progressive neurodegenerative disease caused by abnormal prion proteins, leading to
cognitive decline, ataxia, myoclonus, and death.

\medterm{Creutzfeldt Jakob Disease} Creutzfeldt-Jakob disease is a rare, fatal, transmissible spongiform
encephalopathy caused by the accumulation of misfolded prion protein (PrPSc) in the
brain, leading to rapid, progressive dementia and myoclonus.

\medterm{Creutzfeldt-Jakob Disease} A rare, rapidly progressive brain disease caused by abnormal prion proteins. It leads to worsening memory and thinking, poor coordination, involuntary jerks, and loss of function, and is ultimately fatal.
\synonyms
CJD
Mad Cow Disease
Variant CJD
VCJD

\medterm{Cri Du Chat Syndrome} A congenital genetic disorder caused by deletion of part of the short arm of chromosome
5, characterized in infancy by a high-pitched cry, developmental delay, and
characteristic facial features.

\medterm{Cri-Du-Chat Syndrome} A genetic disorder caused by a deletion of the short arm of chromosome 5 (5p-), with the
size of the deletion correlating with severity. Incidence 1 in 15,000-50,000.

\medterm{Crib Death} Alternate index label; see \textit{Sudden infant death syndrome}.

\medterm{Cricoid Cartilage} A complete ring of firm tissue at the lower part of the voice box, just above the windpipe. It supports the airway and connects with other structures involved in breathing and voice.

\medterm{Cricoid Pressure} Pressure applied to the cricoid cartilage during some emergency airway procedures. Its use is intended to reduce stomach contents entering the airway, but it can make ventilation or intubation more difficult and is clinician-dependent.

\medterm{Cricothyrotomy} An emergency operation creating an opening through the front of the neck into the airway when a person cannot be intubated or oxygenated by other means. It provides temporary access for breathing and has serious risks.

\medterm{Crigler-Najjar Syndrome} A rare inherited disorder in which the liver cannot properly process bilirubin. It causes persistent yellowing of the skin and eyes; severe forms can damage the brain unless treated urgently.

\medterm{Crime Scene Investigation} The systematic documentation, collection, preservation, and analysis of physical evidence from a suspected crime scene, often supporting forensic medical and legal conclusions.

\medterm{Criminal Abortion} Abortion performed illegally, outside permitted legal indications and procedures (e.g.,
unsafe methods, unqualified practitioners, or non-medical settings), carrying high
maternal risk.

\medterm{CRISPR-Cas9} A laboratory gene-editing tool that uses a guide molecule to direct an enzyme to a chosen DNA sequence. The enzyme cuts the DNA so the cell can remove, disable, or replace genetic material.

\medterm{CRISPR Therapy} A treatment that uses CRISPR-based genome-editing tools to make a targeted change to DNA in a patient’s cells. Cells may be edited outside the body and returned to the patient, or edited directly inside the body. It is used or studied for selected genetic and other diseases; possible concerns include unintended genetic changes, immune reactions, and effects that require long-term monitoring.

\medterm{Critical Care} Specialized care for people with life-threatening illness or injury who need close monitoring and support for failing organs.

\medterm{Critical Care Nursing} Specialized nursing care for critically ill patients requiring close monitoring, organ support, rapid assessment, complex therapies, and coordination with multidisciplinary intensive-care teams.
\medterm{Critical Findings} Findings requiring immediate communication. Examples: pneumothorax, brain herniation,
aortic dissection.

\medterm{Critical Illness Myopathy} A muscle weakness disorder that develops during severe illness, especially sepsis, organ failure, prolonged intensive-care treatment, or use of certain medicines. It can affect arms, legs, and breathing muscles, making movement and removal from a ventilator difficult. Rehabilitation supports recovery.

\synonyms
CIM; Acute Quadriplegic Myopathy; Critical Care Myopathy

\medterm{Critical Limb Ischemia} Severely reduced blood flow to an arm or leg causing pain even at rest, a wound that will not heal, or gangrene. It threatens the limb and requires urgent vascular assessment.


\medterm{Critical Results} Test findings suggesting a potentially life-threatening problem that require prompt communication, assessment, treatment, and documentation by the responsible healthcare team.

\medterm{Critical Value} A laboratory result indicating a potentially life-threatening condition that requires
prompt communication to a responsible clinician or care team. The threshold is test- and
laboratory-specific and should be defined by institutional policy.

\medterm{Crocodile Tears Syndrome} Gustatory lacrimation, in which eating or smelling food triggers excessive tearing, usually after facial-nerve injury or aberrant regeneration following Bell palsy.

\medterm{Crohn Disease} Chronic transmural inflammatory condition that can affect any segment of the
gastrointestinal tract from the oral cavity to the perianal region, though the terminal
ileum and ileocolonic region are most commonly involved.

\medterm{Crohn Disease (Crohn'S Disease)} Alternate index label for Crohn's Disease.



\medterm{Crohn Disease Cobblestoning} A bumpy bowel-lining appearance seen during endoscopy in some people with Crohn disease. Deep lengthwise ulcers surround raised areas of less-affected tissue, creating a pattern like cobblestones.

\medterm{Crohn'S Disease} A chronic inflammatory bowel disease that can cause patchy, transmural inflammation anywhere from the mouth to the anus, most often involving the terminal ileum or colon. Common features include abdominal cramping or pain, persistent diarrhoea, weight loss, fatigue, fever, anaemia, and extraintestinal inflammation of the joints, eyes, or skin.

\synonyms
Crohn Disease (Crohn's Disease)
Ileitis (Crohn's Disease)


\medterm{Crohn’S Disease Cause Dark Circles Around Eyes} Crohn disease does not directly cause dark circles in every patient, but anemia, fatigue, dehydration, poor sleep, nutritional deficiency, steroid effects, or atopic skin may contribute. Persistent pallor, fatigue, weight loss, bleeding, or active bowel symptoms deserve clinical review.

\medterm{Cross-Matching} Laboratory testing that checks whether donor blood is compatible with a recipient before transfusion. It helps prevent dangerous immune reactions by detecting antibodies that could attack transfused red cells.

\medterm{Cross-Section} A slice or plane made through a structure, or an image showing that slice, such as a cross-sectional CT image.

\medterm{Cross-Sectional Study} A study that measures possible causes and health outcomes at one point in time. It can show how common a condition is and whether factors are linked, but it usually cannot show which came first or prove cause and effect.

\medterm{Cross-Tolerance} Reduced response to one substance because the body has adapted to another substance with a similar effect. It can increase overdose risk when a person changes substances or combines them.

\medterm{Crossmatch} Laboratory test mixing donor and recipient blood to detect preformed antibodies that could cause transfusion or transplant reactions.

\medterm{Crotch} The region where the inner thighs meet the trunk, including the groin and nearby perineal structures.

\medterm{Croup} An acute viral infection causing inflammation and narrowing of the upper airway, typically producing a barking cough, hoarseness, and inspiratory stridor in young children.

\synonyms
Laryngotracheobronchitis (Croup)



\medterm{Crouzon Syndrome} An autosomal dominant craniosynostosis syndrome caused by mutations in the FGFR2 or
FGFR3 genes, leading to premature fusion of multiple cranial sutures. Features:
brachycephaly (short, broad head from bilateral coronal synostosis), proptosis (shallow
orbits), hypertelorism, midface hypoplasia, beaked nose, and Class III malocclusion.

\medterm{Crown} The visible part of a tooth above the gum line, or a dental restoration that covers and protects a damaged tooth.

\synonyms
Dental crown; anatomic tooth crown.

\medterm{Crown (Dental)} A custom-made cap placed over a weakened or damaged tooth to restore its shape, strength, and function. It may be made from metal, ceramic, or a combination; the tooth is reshaped first, and the crown is cemented in place.

\synonyms
Dental

\medterm{Crown Lengthening} A dental procedure that removes or reshapes gum tissue, and sometimes a small amount of bone, to expose more tooth. It may allow a damaged tooth to be restored or change a gummy smile.

\medterm{Crown-Rump Length} The ultrasound measurement from the top of an embryo's head to its bottom. It is the most accurate way to estimate pregnancy age during early pregnancy and helps monitor early growth.

\medterm{Crowning} The stage of childbirth when the baby's head remains visible at the vaginal opening between contractions because it has stretched the surrounding tissues. It occurs shortly before the head is born.

\medterm{CRP} C-reactive protein, a substance made by the liver that rises during inflammation or infection. A CRP test can help assess inflammation and response to treatment, but it does not identify the cause by itself.

\medterm{Crush Injury} Damage caused when a body part is squeezed or trapped under pressure. It can cause broken bones, muscle death, nerve and vessel injury, dangerous chemical changes in blood, shock, or kidney failure.

\medterm{Crush Syndrome} A life-threatening illness after prolonged muscle compression, often when someone is trapped. Damaged muscle releases substances that can cause shock, dangerous heart rhythms, and sudden kidney failure, especially after pressure is removed.

\medterm{Crushing Injury} Damage caused when a body part is compressed between heavy objects. It can produce fractures,
nerve and blood-vessel injury, tissue death, and rhabdomyolysis. Severe limb compression
or sudden release after prolonged entrapment requires emergency treatment.

\medterm{Crust} Dried serum, blood, pus, or other exudate on the skin surface, often covering an erosion, ulcer, vesicle, pustule, or inflamed lesion.

\medterm{Crusted Scabies} A severe form of scabies causing thick crusts full of very large numbers of mites. It is extremely contagious and occurs mainly in people with weakened immunity, older adults, or people with reduced sensation.

\medterm{Crutches} Walking aids that transfer some weight from the legs to the arms while a person stands or walks. Underarm and forearm types are adjusted to fit and require instruction to reduce falls, nerve pressure, and other injuries.
through loop.


\medterm{Crutches And Children - Sitting And Getting Up From A Chair} Instruction for a child using crutches to approach a chair, keep the injured limb protected as prescribed, reach for the armrests, and sit or stand safely without losing balance.

\medterm{Crutches And Children - Stairs} Instruction for safe stair use by a child with crutches, including handrail use, crutch placement, step sequencing, supervision, and adherence to weight-bearing restrictions.

\medterm{Crutches And Children - Standing And Walking} Instruction for fitting and using crutches in children, including posture, handgrip support, gait pattern, weight-bearing limits, turning, and fall prevention.

\medterm{Cry For Help} A verbal or behavioral expression of distress or need for assistance; in a mental-health setting it should be taken seriously and assessed for immediate safety concerns.

\medterm{Crying In Childhood} Vocal expression of distress or emotion in a child; persistent, inconsolable, unusual, or pain-associated crying may indicate illness, injury, abuse, or unmet developmental needs.

\medterm{Crying In Infancy} Crying is an infant’s normal communication of hunger, discomfort, fatigue, or need for contact, but sudden inconsolable crying with lethargy, fever, vomiting, or injury needs urgent assessment.

\medterm{Cryoablation} A treatment that uses controlled freezing to destroy selected abnormal tissue, such as a tumor or abnormal heart pathway. Imaging or specialized equipment guides the treatment, which may cause bleeding, injury, or recurrence.
\medterm{Cryobiopsy} A biopsy method that uses a freezing instrument to remove a larger tissue sample, often from the lung. It can improve diagnosis in selected cases but carries risks such as bleeding and a collapsed lung.

\medterm{Cryoglobulinemia} A condition in which abnormal blood proteins thicken or clump when cold. They can inflame small blood vessels and cause purplish skin spots, joint pain, nerve problems, kidney injury, or poor circulation.

\medterm{Cryoglobulinemic Vasculitis} Systemic vasculitis caused by cryoglobulin-containing immune complexes that deposit in
small and medium vessels, classically presenting with purpura, arthralgia, and
glomerulonephritis. Most cases are associated with hepatitis C virus infection.

\medterm{Cryoglobulins} Immunoglobulins that precipitate at low temperatures and dissolve on warming; circulating cryoglobulins can cause vasculitis, purpura, neuropathy, kidney disease, or hyperviscosity.

\medterm{Cryopexy} A procedure that uses controlled freezing to destroy or seal abnormal tissue, especially
retinal breaks or small retinal detachments.

\medterm{Cryoprecipitate} A blood-product concentrate containing mainly fibrinogen and some clotting proteins. It is given in selected serious bleeding disorders when a person's fibrinogen level is too low.

\medterm{Cryopreservation} Preserving cells, tissues, sperm, eggs, or embryos by cooling them to very low temperatures with protective chemicals. The material can later be thawed for treatment, research, or reproduction.

\medterm{Cryopyrin-Associated Autoinflammatory Syndromes} Alternate index label; see \textit{Cryopyrin-Associated Periodic Syndromes}; the group is also called CAPS.

\medterm{Cryopyrin-Associated Periodic Syndromes} A group of inherited inflammatory disorders causing repeated fever, rash, joint symptoms, eye inflammation, or hearing loss because immune activity is overactive.

\synonyms
Cryopyrin-Associated Autoinflammatory Syndromes

\medterm{Cryostat} An instrument that maintains very low temperatures and cuts thin frozen sections of tissue for rapid microscopic examination.

\medterm{Cryosurgery} Surgery that destroys abnormal tissue by controlled freezing, used in selected skin
lesions, tumors, and other conditions.

\medterm{Cryotherapy} Treatment using controlled cold to destroy or reduce abnormal tissue. It may treat selected skin lesions, tumors, eye problems, or prostate disease; pain, blistering, scarring, and damage to nearby tissue are possible.

\medterm{Cryotherapy For Prostate Cancer} A treatment that destroys prostate tissue by repeated freezing and thawing, used selectively for localized or recurrent prostate cancer.

\medterm{Cryotherapy For The Skin} Destruction of selected skin lesions by controlled freezing, usually with liquid nitrogen, for conditions such as warts, actinic keratoses, or some superficial tumors.

\medterm{Crypt} A small cavity, gland, or recessed space in body tissue, such as an intestinal crypt or tonsillar crypt.

\medterm{Cryptitis} Inflammation of a crypt or small anatomical recess, especially inflammation of the anal
glands and crypts.

\medterm{Crypto} Clinical shorthand that may refer to cryptococcosis or Cryptosporidium infection; the intended meaning depends on the clinical context.

\medterm{Cryptococcal Meningitis} Infection of the meninges by Cryptococcus, usually in people with impaired immunity, causing
headache, fever, altered mental status, or other neurologic symptoms.

\medterm{Cryptococcosis} Cryptococcosis is fungal infection caused by Cryptococcus, often affecting the lungs or brain in people with immune suppression.

\synonyms
Infection Cryptococcus (Cryptococcosis)

\medterm{Cryptococcus Neoformans} An encapsulated yeast that can cause pulmonary or central nervous system infection, especially in
people with impaired immunity.

\medterm{Cryptogenic Organizing Pneumonia} An inflammatory lung illness in which healing tissue blocks small airways and air sacs, without a known cause. It may cause cough, fever, and breathlessness and often improves with specialist-directed treatment.

\medterm{Cryptorchidism} Failure of one or both testicles to move into the scrotum before birth. An undescended testicle has higher risks of infertility, twisting, injury, and cancer and should be assessed during infancy.

\medterm{Cryptosporidiosis} An intestinal infection caused by Cryptosporidium parasites, spread through contaminated water, food, animals, or people. It usually causes watery diarrhea and cramps for days to weeks and can be severe or prolonged with weakened immunity.

\medterm{Cryptosporidium Enteritis} Diarrheal intestinal infection caused by Cryptosporidium parasites, transmitted through contaminated water, food, or contact; illness can be prolonged or severe in immunocompromised people.

\medterm{Crystal Arthropathy} Joint disease caused by crystal deposition. Gout (monosodium urate) and pseudogout
(calcium pyrophosphate). Diagnosis by joint aspiration and polarized microscopy.

\medterm{Crystal-Induced Synovitis} Acute or chronic joint inflammation caused by deposition of crystals such as monosodium urate, calcium pyrophosphate, or basic calcium phosphate in or around the joint.

\medterm{Crystalline Lens} The clear, flexible structure inside the eye that focuses light onto the retina. It changes shape for near vision and becomes cloudy in cataracts, causing blurred or reduced vision.

\medterm{Crystalloids} Aqueous solutions of mineral salts that distribute between intravascular and
interstitial compartments. Normal saline (0.9 percent NaCl) has 154 mEq/L each of sodium
and chloride, pH 5.4, and is associated with hyperchloremic metabolic acidosis at large
volumes.

\medterm{Crystalluria} The presence of crystals in urine, which may reflect precipitation of salts, medication
effects, metabolic disease, or a tendency to form urinary stones.

\medterm{CSF Analysis} Alternate index label; see \textit{Cerebrospinal fluid analysis}.

\medterm{CSF Cell Count} Laboratory measurement of red and white blood cells in cerebrospinal fluid, used to evaluate meningitis, encephalitis, subarachnoid hemorrhage, inflammation, malignancy, and other central nervous system disorders.

\medterm{CSF Coccidioides Complement Fixation Test} A cerebrospinal-fluid test for complement-fixing antibodies to Coccidioides, used with other findings to evaluate suspected coccidioidal meningitis.

\medterm{CSF Glucose Test} A laboratory measurement of glucose in cerebrospinal fluid, usually interpreted with a simultaneous blood-glucose value. Low cerebrospinal-fluid glucose may occur in bacterial, fungal, or tuberculous meningitis and other disorders, but results require the complete clinical and cerebrospinal-fluid profile.

\medterm{CSF Leak} Escape of cerebrospinal fluid through a defect in the skull base or spinal dura, causing positional headache, clear nasal or ear drainage, wound leakage, meningitis risk, or intracranial hypotension.

\medterm{CSF Leak (Cerebrospinal Fluid Leak)} Escape of cerebrospinal fluid through a tear in the membranes surrounding the brain or spinal cord, often causing positional headache.

\synonyms
Spinal Fluid Leak

\medterm{CSF Myelin Basic Protein} A test measuring a protein released when the insulating covering of nerve fibers is damaged. An increased level may support evaluation of multiple sclerosis or other nervous-system injury but is not specific by itself.

\medterm{CSF Oligoclonal Banding} An electrophoretic test for restricted immunoglobulin bands in cerebrospinal fluid, interpreted with serum bands to detect intrathecal antibody production such as in multiple sclerosis.

\medterm{CSF Rhinorrhea} Cerebrospinal fluid rhinorrhea is the leakage of cerebrospinal fluid from the
subarachnoid space through a defect in the skull base into the nasal cavity or paranasal
sinuses, resulting in a clear, watery nasal discharge that may be mistaken for allergic
rhinitis or sinusitis.

\medterm{CSF Smear} Microscopic examination of cerebrospinal fluid for cells, organisms, or abnormal material, supporting evaluation of meningitis, hemorrhage, malignancy, or inflammation.

\medterm{CSF Total Protein} A measurement of protein in cerebrospinal fluid, usually obtained by lumbar puncture. High or low levels may support evaluation of infection, bleeding, inflammation, nerve disorders, or other brain and spinal-cord disease.

\medterm{CSF-VDRL Test} A cerebrospinal-fluid test for reagin antibodies associated with Treponema pallidum infection; a reactive result is highly specific for neurosyphilis, but a nonreactive result does not exclude it.

\medterm{CT} Computed tomography, an imaging method that uses multiple X-ray measurements and computer reconstruction to produce cross-sectional images.

\synonyms
Computed tomography; CAT scan.

\medterm{CT Angiography} A CT scan performed after injected contrast highlights blood vessels. It can show narrowing, blockage, aneurysm, dissection, bleeding, or abnormal vessel connections in areas such as the brain, heart, chest, abdomen, or limbs.

\medterm{CT Angiography - Abdomen And Pelvis} Computed tomography angiography of the abdominal and pelvic blood vessels after intravenous contrast, used to assess aneurysm, stenosis, dissection, occlusion, bleeding, and vascular anatomy.

\medterm{CT Angiography - Arms And Legs} Computed tomography angiography of the peripheral arteries after intravenous contrast, used to evaluate arterial narrowing, blockage, aneurysm, injury, or other vascular abnormalities in the limbs.

\medterm{CT Angiography - Head And Neck} Computed tomography angiography of the blood vessels supplying the brain and neck after intravenous contrast, used to evaluate aneurysm, stenosis, dissection, occlusion, and vascular malformation.

\medterm{CT Angiography – Chest} Computed tomography angiography of the chest uses intravenous contrast and computed tomography to visualize the thoracic aorta, pulmonary arteries, or other chest vessels. It can evaluate embolism, aneurysm, dissection, bleeding, and vascular malformation.

\medterm{CT Chest} A CT scan of the chest that shows the lungs, heart, blood vessels, mediastinum, pleura, ribs, and chest wall. It can evaluate infection, tumors, injury, fluid, blood clots, and other abnormalities.
\medterm{CT Colonography} A CT examination of the colon after bowel cleansing and gentle inflation with air or carbon dioxide. It can screen for polyps or cancer when colonoscopy is incomplete or unsuitable, but abnormal findings usually require colonoscopy.

\medterm{CT Colonography (Virtual Colonoscopy)} A CT scan of the colon after bowel cleansing and gentle inflation with air or carbon dioxide. It can look for polyps or cancer when colonoscopy is incomplete or unsuitable, but an abnormal result usually requires colonoscopy for tissue removal or biopsy.

\synonyms
Virtual Colonoscopy

\medterm{CT Contrast} A substance administered during computed tomography to make selected tissues, vessels, or body spaces more visible.

\medterm{CT Coronary Angiography} A CT scan taken after injected contrast outlines the arteries supplying the heart. It can show plaque or narrowing and is a non-invasive option for selected people with possible coronary artery disease.

\medterm{CT Dose-Length Product} A measurement of the X-ray output from a CT examination that combines scan dose with scan length. It helps compare and optimize imaging protocols but does not directly predict an individual's radiation risk.

\medterm{CT Enterography} A CT scan of the small intestine after drinking contrast to expand the bowel and receiving injected contrast. It helps evaluate Crohn disease, inflammation, tumors, obstruction, and unexplained intestinal bleeding.

\medterm{CT Perfusion} A CT technique that estimates blood flow and related perfusion parameters in tissue, commonly to assess selected cases of acute stroke or tumors.

\medterm{CT Pulmonary Angiography} CT imaging performed after intravenous contrast to show the pulmonary arteries and evaluate suspected pulmonary embolism or other vascular abnormalities.

\medterm{CT Scan} Computed tomography, an imaging method that uses X-rays and computer processing to produce cross-sectional
images of the body.

\medterm{CT Scan (Computed Tomography)} An X-ray examination processed by a computer to create cross-sectional body images. Scans may be done without contrast or with contrast through a vein or mouth, depending on what needs to be assessed.

\synonyms
Computed Tomography

\medterm{Ctenocephalides FelFlea Bite (Flea Bites In Humans)} Alternate index label for Flea Bites in Humans.

\medterm{CTL} Cytotoxic T lymphocyte, an immune cell that recognizes and kills infected, abnormal, or target cells through antigen-specific mechanisms.

\medterm{CTPA} Alternate name; see \textit{Computed Tomography Pulmonary Angiography}.

\medterm{CTS} Commonly, carpal tunnel syndrome: compression of the median nerve at the wrist causing numbness, tingling, pain, or weakness in the hand.

\medterm{Cubital} Relating to the elbow or nearby forearm. Cubital structures include the joint, bones, muscles, nerves, and blood vessels, and problems there may cause pain, restricted movement, numbness, or weakness.

\medterm{Cubital Tunnel Syndrome} Pressure or stretching of the ulnar nerve at the elbow, causing tingling or numbness in the little and ring fingers, hand weakness, or pain. Symptoms often worsen when the elbow stays bent.

\medterm{Cubitus} The elbow or the forearm; the term is also used in names of structures and conditions around the elbow.

\medterm{Cubitus Valgus} An outward angle of the forearm at the elbow that is larger than usual when the arm is straight. It may be present from birth or develop after an elbow injury. In some people, it can put pressure on the ulnar nerve and cause tingling, numbness, or weakness in the ring and little fingers.

\medterm{Cubitus Varus} An abnormal inward angulation of the forearm relative to the upper arm when the elbow is
extended, often developing after a childhood elbow fracture.

\medterm{Cuboid} A cube-shaped bone on the outer side of the midfoot. It helps form the foot arch and connects several foot bones; injury or displacement can cause outer-foot pain, especially during walking.

\medterm{Cul-De-Sac} A blind-ending pouch or recess; in pelvic anatomy, the rectouterine pouch and rectovesical pouch are examples.

\medterm{Culdocentesis} A procedure that places a needle through the back of the vagina to sample or drain fluid from the lowest part of the pelvis. Ultrasound and less-invasive tests have replaced it in many situations.

\medterm{Cullen'S Sign} Bluish or violaceous periumbilical ecchymosis caused by blood tracking through the
fascial planes to the umbilicus. It is a sign of intra-abdominal or retroperitoneal
haemorrhage, classically severe haemorrhagic pancreatitis, and is associated with
serious underlying disease.

\medterm{Cultural Psychiatry} The study and clinical practice of how cultural factors influence mental health, illness
expression, diagnosis, and treatment. Cultural concepts of distress, idioms of distress,
and cultural syndromes must be considered. The DSM-5 includes a Cultural Formulation
Interview.

\medterm{Culture} A laboratory test in which a clinical specimen is placed in conditions that allow microorganisms to grow and be identified; susceptibility testing may help guide treatment.

\medterm{Culture - Colonic Tissue} A laboratory culture of a colon tissue specimen to detect bacteria, fungi, or other microorganisms in suspected infection; results are interpreted with histopathology and clinical findings.

\medterm{Culture - Duodenal Tissue} A laboratory culture of a duodenal tissue specimen to identify microorganisms causing infection or inflammation, interpreted with endoscopic findings and histopathology.

\medterm{Culture And Sensitivity} Microbiological culture of a clinical specimen (blood, urine, sputum, wound swab, CSF)
to identify pathogens and determine antimicrobial susceptibility. Culture media: blood
agar (most bacteria), MacConkey (Gram-negative), chocolate agar (fastidious: N.
gonorrhoeae, H. influenzae), and selective media.

\medterm{Culture Media} Nutrient materials prepared to grow microorganisms in a laboratory, with different formulas chosen for specific organisms, tests, or identification purposes.

\medterm{Culture Swab} A sterile swab used to collect material from skin, wounds, the throat, genital tract, or another site for laboratory culture. The test may identify bacteria or fungi and guide antibiotic choice; viral testing usually uses separate molecular or antigen tests.

\medterm{Culture-Negative Endocarditis} Infective endocarditis in which routine blood cultures do not identify an organism, often because of prior antibiotics, fastidious bacteria, fungi, or noninfectious mimics; diagnosis requires integrated clinical, imaging, and laboratory assessment.

\medterm{Cumulative Dose} The total amount of a substance or radiation received over time. Cumulative exposure matters for medicines, toxins, radiation, carcinogens, and pregnancy-related risks because effects may build gradually.

\medterm{Cumulative Trauma Disorder} Pain or injury that develops after repeated forceful, awkward, or prolonged movements. It may affect muscles, tendons, nerves, or other soft tissues in the hands, arms, neck, shoulders, or back.

\medterm{Cup-To-Disc Ratio} A comparison between the central hollow of the optic nerve and the whole visible optic nerve head. An enlarged or changing ratio may suggest glaucoma but must be interpreted with eye pressure and other findings.

\medterm{Cupping} Enlargement or deepening of the central depression in the optic disc, assessed during glaucoma evaluation and interpreted with eye pressure and other findings.

\synonyms
Optic-disc cupping; glaucomatous cupping.

\medterm{CURB-65 Score} A five-point tool estimating the severity of community-acquired pneumonia using confusion, blood urea, breathing rate, blood pressure, and age 65 or older. It helps guide hospital decisions but does not replace clinical judgment.

\medterm{Cure} Elimination or durable control of a disease so that it no longer causes active illness; the meaning and evidence for cure vary by condition.

\medterm{Curettage} Removal of tissue from a body cavity or surface with a curette, used for diagnosis,
treatment, or removal of abnormal tissue.

\medterm{Curettage Specimen} Tissue scraped or scooped from a body surface, cavity, or bone lesion for laboratory examination. Because the pieces are often small and fragmented, they may not show the full structure or depth of disease.

\medterm{Curette} A spoon-shaped or ring-shaped surgical instrument used to scrape, remove, or clean tissue from a cavity, bone, skin lesion, or tooth-supporting area. Its use and risks depend on the procedure and body site.

\medterm{Curling Ulcer} An acute ulcer in the stomach or duodenum that can develop after severe burns. Reduced blood flow and major physical stress weaken the digestive lining and may cause dangerous bleeding.

\medterm{Curling'S Ulcer} An acute peptic ulcer associated with severe burns, often occurring in the duodenum or
stomach during physiologic stress.

\medterm{Curvature Of The Penis} A bend of the erect or flaccid penis caused by congenital tissue asymmetry, scar tissue such as Peyronie disease, trauma, or other conditions; it may impair intercourse or cause pain.

\synonyms
Peyronie's Disease (Curvature Of The Penis)

\medterm{Curvularia Alcornii} A dark-pigmented environmental fungus that rarely causes infection. It may affect the skin, sinuses, or brain, mainly in people with weakened immunity, and requires specialist laboratory identification and treatment.

\medterm{Cushing Disease} Cushing syndrome caused by an adrenocorticotropic hormone-secreting pituitary tumor.
Excess cortisol can cause central weight gain, rounded facial features, skin changes, muscle weakness, high blood
pressure, high blood glucose, and menstrual or mood changes.

\medterm{Cushing Syndrome} Illness caused by prolonged exposure to too much cortisol or steroid medicine. It can cause central weight gain, easy bruising, muscle weakness, high blood pressure, diabetes, bone loss, infections, and mood or menstrual changes.

\medterm{Cushing Syndrome Complications} Untreated Cushing syndrome carries significant morbidity and mortality from
cardiovascular disease, infections, thromboembolic events, and metabolic derangements.
Cardiovascular complications include accelerated atherosclerosis, hypertension, and
dyslipidemia. Infection risk is increased from immunosuppression.

\medterm{Cushing Syndrome Due To Adrenal Tumor} Excess cortisol caused by an adrenal tumor, producing Cushing syndrome with features such as central weight gain, hypertension, diabetes, muscle weakness, and skin changes.

\medterm{Cushing Ulcer} An acute stomach or duodenal ulcer associated with severe brain disease, head injury, or neurosurgery. Excess nerve stimulation can increase acid production and cause abdominal pain or gastrointestinal bleeding.

\medterm{Cushing'S Disease} A pituitary tumor that makes too much adrenocorticotropic hormone (ACTH), causing the adrenal glands to produce too much cortisol. Features may include weight gain around the trunk, a round face, easy bruising, purple stretch marks, muscle weakness, high blood pressure, high blood glucose, and menstrual or mood changes.

\medterm{Cushing'S Syndrome} A disorder caused by prolonged exposure to too much cortisol or steroid medicine. It may cause weight gain around the trunk, a round face, fat at the back of the neck, thin easily bruised skin, purple stretch marks, muscle weakness, high blood pressure, and high blood glucose. Untreated disease can lead to infections, blood clots, bone loss, and heart or blood-vessel disease.

\synonyms
Hypercortisolism; cortisol excess syndrome.

\medterm{Cushingoid} Resembling the physical features of excess glucocorticoid exposure, such as a round face, central weight gain, thin skin, easy bruising, and muscle weakness.

\medterm{Cusp} A pointed projection, such as a projection of a tooth or one of the leaflets of a heart valve.

\medterm{Cusp Fracture} A break in one of the raised points of a tooth, often around a large filling or weakened tooth. It may cause pain when biting or sensitivity and may require a bonded repair, crown, or other dental treatment.
\medterm{Cut} A break or incision made in tissue, usually by a sharp object; the depth, contamination, and location determine the required care.

\medterm{Cutaneous} Relating to the skin. The skin forms a protective barrier, helps control temperature, senses touch and pain, and supports vitamin D production. Skin findings can provide clues to infections, immune diseases, hormone disorders, and skin cancers.

\medterm{Cutaneous Angiosarcoma} A fast-growing cancer of blood-vessel cells in the skin, usually on the scalp or face of older adults. It may look like a bruise, ulcerate, and spread.

\medterm{Cutaneous Anthrax} A skin infection caused by Bacillus anthracis, usually producing a painless black sore surrounded by swelling.

\medterm{Cutaneous B-Cell Lymphoma} A non-Hodgkin lymphoma mainly affecting B lymphocytes in the skin. It may cause persistent patches, plaques, or lumps and is treated according to subtype, extent, symptoms, and spread.

\medterm{Cutaneous Diphtheria} A skin infection caused by toxigenic or nontoxigenic Corynebacterium diphtheriae, producing chronic
ulcers or a gray membrane.

\medterm{Cutaneous Langerhans Cell Histiocytosis} An uncommon disorder in which abnormal immune cells build up in the skin. It may cause scaly or crusted spots, rash, or purple marks; testing checks whether other organs are involved.

\medterm{Cutaneous Larva Migrans} An intensely itchy, winding skin rash caused by hookworm larvae from contaminated soil or sand entering bare skin. The larvae move beneath the skin and usually die without reaching internal organs.

\medterm{Cutaneous Leiomyoma} A usually noncancerous smooth-muscle tumor in the skin, often causing firm, painful bumps. Multiple lesions may be associated with an inherited condition that also increases kidney-cancer risk and needs specialist assessment.

\medterm{Cutaneous Leishmaniasis} A skin infection caused by Leishmania parasites transmitted by sandflies, usually producing one
or more chronic ulcers.

\medterm{Cutaneous Lupus} A form of lupus erythematosus limited to the skin, manifesting as discoid lesions,
subacute cutaneous plaques, or acute malar erythema. It may occur in isolation or as a
feature of systemic lupus erythematosus.

\medterm{Cutaneous Lupus Erythematosus} An immune disorder causing inflamed skin patches or sores that may worsen with sunlight, scar, or change skin color. It may remain limited to skin or occur with systemic lupus and requires appropriate evaluation.

\medterm{Cutaneous Lymphoma} Cutaneous lymphomas are heterogeneous group of T-cell and B-cell malignancies primarily
involving the skin. Primary cutaneous B-cell lymphomas include marginal zone lymphoma,
follicle center lymphoma, and diffuse large B-cell lymphoma leg type. Primary cutaneous
T-cell lymphomas include mycosis fungoides and Sézary syndrome.

\medterm{Cutaneous Metastases} Secondary deposits of cancer in the skin, originating from a primary tumor elsewhere in
the body. They may present as nodules, plaques, or ulcers, and often indicate advanced
or recurrent disease.

\medterm{Cutaneous Metastasis} A deposit of malignant cells in the skin from a primary tumor elsewhere or, less often, another cutaneous site; it may present as a firm nodule, plaque, or inflammatory-appearing lesion.

\medterm{Cutaneous Polyarteritis Nodosa} A skin-limited necrotizing vasculitis of small- and medium-sized arteries, typically causing tender subcutaneous nodules, livedo reticularis, and ulcers on the lower legs without systemic organ involvement.

\medterm{Cutaneous Skin Tag} A small, soft, usually harmless growth of extra skin, often on a narrow stalk in the neck, armpit, groin, or another skin fold. It may catch on clothing or become irritated.
\medterm{Cutaneous Small-Vessel Vasculitis} Inflammation and destruction of small dermal blood vessels, commonly presenting as
palpable purpura on dependent areas. Causes include infections, medications, autoimmune
disease, and systemic vasculitis; skin biopsy with appropriate specimen handling is
central to diagnosis.

\medterm{Cutaneous T-Cell Lymphoma} A malignancy of skin-homing T lymphocytes. Mycosis fungoides is the most common subtype,
presenting as patches, plaques, and tumors in a clothing distribution. Sézary syndrome
is the leukemic variant with erythroderma, lymphadenopathy, and circulating malignant
cells.

\medterm{Cuticle} The thin outermost layer of a hair shaft formed by overlapping cells that protect the inner cortex.

\synonyms
Hair cuticle; cuticular layer.

\medterm{Cuticle Remover Poisoning} Exposure to cuticle-remover products, which may contain alkaline or acidic chemicals that can injure the skin, eyes, mouth, or gastrointestinal tract.

\medterm{Cutis Anserina} Gooseflesh or piloerection, in which contraction of tiny muscles around hair follicles makes the skin appear rough or dimpled.

\medterm{Cutis Laxa} A group of inherited or acquired connective-tissue disorders characterized by loose, inelastic skin caused by abnormal elastic fibers. Some forms also affect the lungs, blood vessels, heart, gastrointestinal tract, or skeleton.

\medterm{Cuts And Puncture Wounds} Breaks in the skin caused by sharp, blunt, or penetrating trauma; evaluation considers depth, contamination, retained foreign bodies, tendon or nerve injury, vascular damage, and tetanus protection.

\medterm{Cuts Wounds (Cuts, Scrapes And Puncture Wounds)} Alternate index label for Cuts, Scrapes and Puncture Wounds.

\medterm{Cuts, Scrapes And Puncture Wounds} Cuts, scrapes, and puncture wounds damage the skin to varying depths and can introduce bacteria or foreign material. Clean minor wounds with running water, control bleeding, cover them, and seek care for deep or contaminated wounds, bites, impaired circulation, infection, or uncertain tetanus protection.

\synonyms
Cuts Wounds (Cuts, Scrapes And Puncture Wounds)

\medterm{Cutting Balloon (Vascular)} A balloon catheter with tiny blades or scoring surfaces that make controlled cuts in hardened plaque while the balloon widens a narrowed blood vessel. It is used in selected vascular procedures.

\synonyms
Vascular

\medterm{Cutting For The Stone} An older term for surgically opening the urinary tract to remove a stone, historically called lithotomy.

\medterm{CVC} Central venous catheter, a catheter whose tip lies in a large central vein and provides intravenous access for medications, fluids, nutrition, or monitoring.

\medterm{Cyanide} A highly poisonous chemical that prevents cells from using oxygen. Exposure can rapidly cause headache, confusion, seizures, breathing failure, cardiovascular collapse, and death.

\medterm{Cyanide Poisoning} A medical emergency in which cyanide stops cells from using oxygen. It may follow smoke inhalation or industrial, accidental, or intentional exposure and can cause headache, confusion, seizures, collapse, and death within minutes.

\medterm{Cyanoacrylates Poisoning} Exposure to cyanoacrylate adhesives, commonly known as instant glue. Contact may bond tissues or irritate the skin, eyes, or airways; ingestion usually causes local adhesion or irritation.

\medterm{Cyanosis} A blue or gray-blue color of the skin, lips, or moist linings caused by too little oxygen in blood or abnormal blood flow. Sudden cyanosis, especially with breathing difficulty, requires emergency care.

\medterm{Cyanotic} Having a bluish or slate-gray discoloration of skin or mucous membranes, usually reflecting increased deoxygenated hemoglobin or abnormal blood flow.

\medterm{Cyanotic Heart Disease} A congenital or acquired heart disorder that reduces systemic oxygen saturation and causes cyanosis, often through right-to-left shunting or inadequate pulmonary blood flow.

\medterm{Cycle Time} Time for complete image acquisition. PET: minutes. SPECT: 20-30 minutes.

\medterm{Cyclic Citrullinated Peptide} A substance used in a blood test for antibodies associated with rheumatoid arthritis. A positive result supports the diagnosis but must be interpreted with symptoms, examination, and other tests.

\medterm{Cyclic Guanosine Monophosphate} A chemical messenger inside cells that helps regulate blood-vessel widening, smooth-muscle relaxation, vision, and other functions. Its level can be changed by some medicines.

\synonyms
cGMP; cyclic GMP.

\medterm{Cyclic Hormone Therapy} Hormone therapy in which estrogen is given continuously or for part of a cycle and progestogen is added for a defined number of days.

\synonyms
Sequential hormone therapy; cyclic estrogen-progestogen therapy.

\medterm{Cyclic Neutropenia} Repeated episodes of low neutrophils, the white blood cells that fight many infections. Episodes often recur at regular intervals and may cause fever, mouth ulcers, sore throat, or infections.

\medterm{Cyclic Vomiting Syndrome} Repeated, similar attacks of severe nausea and vomiting separated by periods of normal health. It is often related to migraine biology; dehydration and other causes must be considered during evaluation.

\medterm{Cyclooxygenase} An enzyme family that helps convert arachidonic acid into prostaglandins and related mediators involved in pain, inflammation, and platelet function.

\synonyms
COX enzyme; prostaglandin-endoperoxide synthase.

\medterm{Cyclooxygenase-2} An inducible cyclooxygenase enzyme that generates prostaglandins involved in inflammation, pain, and fever and is inhibited by COX-2-selective medicines.

\medterm{Cyclops} A term that may refer to cyclopia, a rare severe congenital malformation with a single central eye, or to a fibrous nodule causing loss of extension after knee surgery; context determines the meaning.


\medterm{Cyclospora Infection} An intestinal infection caused by Cyclospora parasites, usually acquired from contaminated food or water. It can cause prolonged or recurring watery diarrhea, cramps, nausea, tiredness, and weight loss.

\medterm{Cyclospora Infection (Cyclosporiasis)} Alternate index label for Cyclosporiasis.

\medterm{Cyclosporiasis} An intestinal infection caused by the parasite Cyclospora cayetanensis, usually acquired by consuming contaminated food or water. It commonly causes prolonged or relapsing watery diarrhea, abdominal cramps, nausea, fatigue, loss of appetite, and weight loss; diagnosis requires specialized stool testing, and trimethoprim-sulfamethoxazole is the usual effective treatment when appropriate.

\synonyms
Cyclospora Infection (Cyclosporiasis)

\medterm{Cyclosporiasis (Cyclospora)} An intestinal infection caused by Cyclospora parasites, typically producing prolonged watery diarrhea, cramps, nausea, and fatigue.

\medterm{Cyclothymia} A long-term mood disorder with repeated periods of hypomanic and depressive symptoms that do not meet criteria for full episodes.

\medterm{Cyclothymia (Cyclothymic Disorder)} A chronic mood disorder involving repeated periods of hypomanic and depressive symptoms that do not meet criteria for full episodes.

\medterm{Cyclothymic Disorder} A long-term mood disorder with repeated periods of elevated energy or activity and periods of depressive symptoms that are milder than full mania or major depression. Symptoms persist for years and can disrupt relationships or work.

\medterm{Cylindroma} A usually benign adnexal tumor composed of basaloid cells, often presenting as a firm pink or skin-coloured nodule on the scalp or face; multiple lesions may occur in Brooke-Spiegler syndrome.

\medterm{CYP2C19 Polymorphism} CYP2C19 is a cytochrome P450 enzyme involved in the metabolism of approximately 10
percent of drugs, including the antiplatelet agent clopidogrel, proton pump inhibitors,
some antidepressants, and the antiepileptic drug phenytoin.

\medterm{CYP2D6 Polymorphism} CYP2D6 is a cytochrome P450 enzyme responsible for the metabolism of approximately 25
percent of commonly prescribed drugs, including codeine, tamoxifen, many
antidepressants, and antipsychotics.

\medterm{CYP450 System} A family of enzymes, found mainly in the liver, that breaks down many medicines and foreign chemicals. Genetic differences and other medicines can speed up or slow down this process and alter drug effects.

\medterm{Cyproheptadine Overdose} Overdose of cyproheptadine, an antihistamine, that may cause sedation, agitation, anticholinergic effects, seizures, respiratory depression, or cardiac effects.

\medterm{Cyst} A closed sac containing fluid, air, cells, or semisolid material. Cysts can occur in the skin, organs, glands, or other tissues and may be harmless, infected, or related to another disease.

\medterm{Cyst Eyelid (Giardia Lamblia)} Alternate index label for Giardia Lamblia.

\medterm{Cyst, Bartholin'S} Alternate index label; see \textit{Bartholin's cyst}.

\medterm{Cyst, Ganglion} Alternate index label; see \textit{Ganglion cyst}.

\medterm{Cyst, Kidney} Alternate index label; see \textit{Kidney cysts}.

\medterm{Cyst, Ovarian} Alternate index label; see \textit{Ovarian cysts}.

\medterm{Cyst, Pancreatic} Alternate index label; see \textit{Pancreatic cysts}.

\medterm{Cyst, Pilonidal} Alternate index label; see \textit{Pilonidal cyst}.

\medterm{Cyst, Spermatic} Alternate index label; see \textit{Spermatocele}.

\medterm{Cystadenoma} A usually noncancerous gland-forming tumor containing fluid-filled spaces or cyst-like structures. It can occur in several organs, may enlarge or cause symptoms, and sometimes requires removal or monitoring.

\medterm{Cystic Acne} Severe acne with deep, painful lumps that can leave permanent scars. It often needs clinician-directed treatment; squeezing or picking the lesions increases scarring and infection risk.

\medterm{Cystic Echinococcosis} A parasitic infection in which slowly growing cysts form, most often in the liver or lungs, after swallowing Echinococcus eggs from infected dogs or contaminated food, water, or soil. Cysts may cause pain, cough, blockage, or serious allergic reactions if they rupture.

\medterm{Cystic Fibrosis} An inherited condition that makes mucus unusually thick and sticky, affecting the lungs and digestive system and sometimes the liver and reproductive organs. It can cause repeated lung infections, cough, poor digestion, poor growth, or infertility; treatment includes airway clearance, nutrition, enzymes, antibiotics, and selected targeted medicines.

\synonyms
CF

\medterm{Cystic Fibrosis - Nutrition} Nutritional care for cystic fibrosis, often requiring extra calories, salt, fat-soluble vitamins, and pancreatic enzymes when digestion is impaired. Growth, weight, stool changes, and blood tests guide the plan.
\medterm{Cystic Hygroma} A birth defect involving an abnormal collection of lymph vessels, usually forming a soft fluid-filled swelling in the neck. It may grow, become infected, or press on the airway and requires specialist assessment.

\medterm{Cystic Medial Necrosis} Weakening of the middle layer of the aorta caused by loss of supportive cells and elastic fibers. It can contribute to an aortic aneurysm or dissection and may occur with inherited connective-tissue disorders or aging.

\medterm{Cysticercosis} Cysticercosis is tissue infection by larval Taenia solium after swallowing eggs, potentially causing cysts in muscle, skin, eyes, or brain.

\medterm{Cystine And Triple Phosphate (Struvite) Crystals} Crystalline substances that may be found in urine sediment. Hexagonal cystine crystals
support cystinuria in the appropriate setting, while struvite crystals may occur in
infection-related alkaline urine; microscopy is interpreted with clinical and laboratory
findings.

\synonyms
struvite

\medterm{Cystine Stones} Kidney stones formed when inherited cystinuria causes excess cystine in urine. They often recur and may be difficult to dissolve; abundant fluids, medical treatment, and sometimes procedures help prevent or remove them.

\medterm{Cystinuria} An inherited disorder in which the kidneys release too much cystine into urine, allowing recurrent cystine kidney stones to form. It requires long-term fluid intake, monitoring, and sometimes medicines to prevent stones.
\medterm{Cystitis} See \textit{Bladder Infection}.

\synonyms
Bladder Infection

\medterm{Cystitis (Bladder Infection)} See \textit{Bladder Infection}.

\medterm{Cystitis - Acute} Acute inflammation of the urinary bladder, most often caused by bacterial infection and producing dysuria, urinary frequency, urgency, and suprapubic discomfort.

\medterm{Cystitis - Noninfectious} Inflammation of the urinary bladder without a proven bacterial infection. Causes include interstitial cystitis or bladder pain syndrome, medicines, radiation, chemical irritation, autoimmune disease, and injury; symptoms may include frequency, urgency, and painful urination.

\medterm{Cystitis, Interstitial} Alternate index label; see \textit{Interstitial cystitis}.

\medterm{Cystocele} Bulging of the bladder into the front wall of the vagina because supporting tissues have weakened. It may cause pressure, a vaginal bulge, urinary leakage, difficulty emptying the bladder, or no symptoms.

\medterm{Cystography} X-ray or other imaging of the urinary bladder after contrast is placed in it. It can show bladder injury, leaks, fistulas, structural abnormalities, or urine flowing backward toward the kidneys.

\medterm{Cystoid Macular Edema} Swelling at the center of the retina caused by fluid collecting in small cyst-like spaces. It can blur or distort central vision and may follow eye surgery, inflammation, diabetes, or retinal blood-vessel disease.
\medterm{Cystometric Study} A bladder test that measures pressure, sensation, and urine storage while the bladder fills and empties. It helps investigate urgency, leakage, difficulty emptying, weak bladder muscles, or nerve-related urinary problems.

\medterm{Cystometrogram} A recording of bladder pressure and sensation during filling and emptying. It is part of urodynamic testing used to investigate leakage, urgency, blockage, weak contractions, or nerve-related bladder problems.

\medterm{Cystometry} A bladder-function test measuring pressure and sensation while the bladder fills and empties. It helps assess urgency, leakage, obstruction, weak contractions, nerve problems, or other urinary difficulties.

\medterm{Cystoscope} A thin, lighted instrument used to examine the urethra and bladder and, when needed, perform procedures.

\medterm{Cystoscopy} Cystoscopy uses a camera passed through the urethra to inspect the bladder and urethra and may allow biopsy or treatment.

\medterm{Cystostomy} Surgical creation of an opening from the urinary bladder to the skin or another
structure to permit drainage of urine.

\medterm{Cystotome} A surgical instrument used to make an incision in the capsule of the lens during selected eye procedures.

\medterm{Cystourethrography} Radiographic imaging of the urinary bladder and urethra, usually during filling and
voiding after contrast has been introduced.

\medterm{Cysts} Closed sac-like structures containing fluid, air, semisolid material, or other contents; their causes range from benign developmental or obstructive processes to infection or neoplasia.

\medterm{Cysts Choledochal (Choledochal Cysts)} Alternate index label for Choledochal Cysts.

\medterm{Cysts Ovary (Ovarian Cysts)} Alternate index label for Ovarian Cysts.

\medterm{Cysts Pancreatic Inflammatory (Pancreatic Cysts)} Alternate index label for Pancreatic Cysts.

\medterm{Cysts True Pancreas (Pancreatic Cysts)} Alternate index label for Pancreatic Cysts.

\medterm{Cysts, Epidermoid} Alternate index label; see \textit{Epidermoid cysts}.

\medterm{Cysts, Sebaceous} Alternate index label; see \textit{Epidermoid cysts}.

\medterm{Cytochrome B5 Reductase} An enzyme that maintains hemoglobin iron in the ferrous state needed for oxygen binding; inherited deficiency can cause congenital methemoglobinemia, usually presenting with cyanosis and impaired oxygen delivery.

\medterm{Cytochrome C} A protein inside mitochondria that carries electrons during energy production. When released from damaged mitochondria, it can also help start the controlled death of a cell.

\medterm{Cytochrome P450} A family of enzymes, found mainly in liver cells, that chemically changes many medicines, toxins, and hormones so the body can remove them. Genetic differences and drug interactions can alter their activity.

\medterm{Cytogenetics} The study of chromosomes and their number, structure, and arrangement. Tests such as karyotyping, fluorescence in situ hybridization, and chromosome microarray can help identify inherited conditions, leukemia, other cancers, and some fetal abnormalities.

\medterm{Cytogenetics (Karyotype)} Visualization of chromosomes under light microscopy from dividing cells (bone marrow,
blood, amniotic fluid, tissue). Staining: G-banding (Giemsa, light and dark bands).
Resolution: 400-550 bands per haploid genome.

\synonyms
Karyotype

\medterm{Cytokine} A small protein messenger that allows cells to communicate during immunity, inflammation, infection, healing, and blood-cell production. Cytokines can protect tissues, but excessive or poorly controlled activity can cause disease.

\medterm{Cytokine Release Syndrome} A potentially serious whole-body inflammatory reaction caused by a sudden release of immune signaling proteins, sometimes after immune-based cancer treatment or infection. Fever, low blood pressure, breathing difficulty, and organ problems may occur.

\medterm{Cytokine Storm} A severe, dysregulated systemic inflammatory response involving excessive cytokine
release that can cause tissue injury, organ dysfunction, and shock.

\medterm{Cytokines} Small signaling proteins released by immune and other cells that regulate inflammation, immunity, blood-cell production, and tissue responses.

\synonyms
Immune signaling proteins; cytokine mediators.

\medterm{Cytologic Evaluation} Microscopic examination of individual cells or small cell groups from a fluid, brushing, scraping, or fine-needle sample to detect infection, inflammation, or malignancy.

\medterm{Cytology} The microscopic study of individual cells, distinct from histology in that tissue
architecture is not examined, used for screening and diagnosis of neoplastic and
infectious disease. Specimens are exfoliative (cervical smears, sputum) or aspirated
(fine needle aspiration, body fluids), interpreted with Papanicolaou or special stains.

\medterm{Cytology Exam Of Pleural Fluid} Microscopic examination of fluid around the lung for malignant cells and other cellular changes, performed with chemical and microbiologic testing when evaluating an effusion.

\medterm{Cytology Exam Of Urine} Microscopic examination of shed urinary-tract cells, used mainly to detect high-grade urothelial malignancy and selected abnormal cells; a negative result does not exclude cancer.

\medterm{Cytomegalovirus} A herpesvirus that causes significant disease in immunocompromised patients including
transplant recipients and HIV-infected individuals. CMV infects a wide range of cell
types and establishes latent infections in myeloid progenitor cells. Primary infection
is usually asymptomatic in immunocompetent individuals.

\medterm{Cytomegalovirus Infection} A widespread herpesvirus infection. Primary infection in immunocompetent individuals is
usually asymptomatic or causes mononucleosis-like syndrome. Congenital CMV: most common
congenital infection, causes sensorineural hearing loss, microcephaly, and developmental
delay.

\medterm{Cytopathology} The examination of individual cells under a microscope to help diagnose cancer, infection, inflammation, or other disease. Cells may come from fluid, a brushing, a smear, or a needle sample.

\medterm{Cytopenia} A lower-than-normal number of one or more blood-cell types: red blood cells, white blood cells, or platelets. The effects depend on which cells are low and how severe the reduction is.

\medterm{Cytoplasm} The material inside a cell between the cell membrane and nucleus, consisting of cytosol, organelles, and other cellular components.

\medterm{Cytoreductive Surgery} Surgery that removes as much visible cancer as safely possible to reduce the amount of disease. It may be combined with other treatment, and its benefits and risks depend on the cancer's site and spread.

\medterm{Cytoskeleton} A network of protein fibers inside a cell that maintains its shape, moves materials, enables movement, and helps the cell divide. It constantly changes in response to the cell's needs.

\medterm{Cytotoxic} Toxic to cells or capable of killing cells, as with some chemotherapy drugs, toxins, or cytotoxic immune cells.

\medterm{Cytotoxic Alopecia} Hair loss caused by exposure to a cytotoxic treatment or other toxic injury to actively growing hair follicles.

\synonyms
Drug-induced alopecia; chemotherapy-related hair loss.

\medterm{Cytotoxic T Cell} An immune cell that recognizes and kills infected, abnormal, or cancerous cells. It identifies small protein signals displayed on the target cell and can trigger that cell to die.

\medterm{Cytotoxicity} The capacity of an immune effector cell or immune mediator to injure or kill a target
cell. T lymphocytes and natural killer cells commonly use perforin-granzyme pathways or
death-receptor signaling to induce target-cell apoptosis.
\medterm{Common Carotid Artery} One of two major arteries carrying blood from the heart toward the head and neck. Each divides into an internal branch supplying the brain and an external branch supplying the face and scalp.

\medterm{Camelid} Relating to camels, llamas, or alpacas. Camelid antibodies and antibody fragments are used in research and some diagnostic or therapeutic applications; contact with camelids can also be relevant to certain infections or allergies.

\medterm{Cancer of the Ear} Cancer arising in the outer ear, ear canal, middle ear, or nearby tissues. Persistent sores, bleeding, pain, discharge, hearing loss, or a growing lump require examination; treatment depends on the site and tumor type.

\medterm{Cancer Treatment} Medical, surgical, radiation, systemic, or supportive care used to control cancer, remove it, slow its growth, prevent recurrence, or relieve symptoms. Treatment is selected according to cancer type, stage, biomarkers, overall health, and patient goals.

\medterm{Cantu Syndrome} A rare genetic condition, usually caused by an ABCC9 or KCNJ8 gene change. It may cause excessive hair growth, distinctive facial features, loose skin, enlarged bones, skeletal differences, and an enlarged or abnormal heart.

\medterm{Cardiometabolic} Relating to the interaction between heart and blood-vessel disease and metabolic conditions such as obesity, insulin resistance, diabetes, and abnormal blood lipids.

\medterm{Carotid Artery Stenosis} Narrowing of an artery in the neck that supplies the brain, usually from fatty plaque. It may reduce blood flow or release a clot and increase the risk of a transient ischemic attack or stroke.

\medterm{Carotid Occlusion} Complete blockage of an artery supplying the brain, usually from plaque or a blood clot. Some people have no symptoms because other vessels compensate, but sudden weakness, speech trouble, or vision loss is an emergency.

\medterm{Cat Allergy} An immune reaction to proteins in cat saliva, skin flakes, or other secretions. It can cause sneezing, itchy eyes, wheezing, hives, or asthma symptoms; reducing exposure and using appropriate allergy treatment may help.

\medterm{Cat Scratch Disease} An infection caused by \textit{Bartonella henselae}, usually after a scratch or bite from a kitten or cat. It commonly causes a small skin lesion and nearby swollen lymph nodes and is usually self-limited, though serious complications can occur.

\medterm{Catamenial Epilepsy} A pattern in which seizures become more frequent or severe at certain times in the menstrual cycle. Hormone changes may contribute, and a seizure diary can help guide neurologic treatment.

\medterm{Celebrex} A brand name for celecoxib, a selective cyclooxygenase-2 nonsteroidal anti-inflammatory drug used for pain and inflammation. It can cause stomach, kidney, cardiovascular, allergic, or bleeding complications and should be used as directed.

\medterm{Cell Signaling} The process by which cells detect information through receptors and transmit signals inside the cell to change functions such as growth, movement, metabolism, or immune activity.

\medterm{Central Pain Syndrome} Chronic pain caused by damage or dysfunction in the brain or spinal cord. Pain may be burning, aching, electric, or unusually sensitive to touch and may occur after stroke, spinal-cord injury, multiple sclerosis, or other lesions.

\medterm{Cerebrovascular Disease} A group of disorders affecting blood vessels and blood flow in the brain, including stroke, transient ischemic attack, aneurysm, and vascular malformations.

\medterm{Cervical Ectropion} A common condition in which cells normally lining the cervical canal extend onto the outer cervix. It may cause increased discharge or bleeding after sex, but it is usually benign and is not the same as cervical cancer.

\medterm{Char Syndrome} A rare genetic disorder, usually caused by variants in \textit{TFAP2B}, characterized by distinctive facial features, a heart defect such as patent ductus arteriosus, and abnormalities of the fifth fingers or toes.

\medterm{Charles Bonnet Syndrome} Visual hallucinations occurring in a person with significant vision loss who otherwise understands that the images are not real. The images may be simple patterns or detailed scenes; new hallucinations still require medical assessment.

\medterm{CHILD Syndrome} Congenital hemidysplasia with ichthyosiform nevus and limb defects, a rare disorder usually caused by an \textit{NSDHL} variant. It causes one-sided skin abnormalities with limb, skeletal, and sometimes neurologic or internal-organ differences.

\medterm{Chimeric Antigen Receptor} A synthetic receptor engineered to recognize a selected target on a cell surface. It is placed into immune cells, especially T cells, to direct them against cancer cells in some cellular immunotherapies.

\medterm{Chinese Herbal Medicine} The use of plant, mineral, or animal-derived preparations within traditional Chinese medicine. Products vary in composition and may interact with medicines or contain contaminants, so reliable sourcing and medical guidance are important.

\medterm{CHIP} Clonal hematopoiesis of indeterminate potential: an age-related finding in which blood-forming cells acquire a gene change and expand without meeting criteria for a blood cancer. It can increase the future risk of myeloid cancer and cardiovascular disease and may require monitoring.

\medterm{Choroidal Melanoma} A cancer arising from pigment-producing melanocytes in the choroid, the vascular layer beneath the retina. It may cause visual changes or be found incidentally and can spread, especially to the liver; treatment depends on size and location.

\medterm{Chromosome 1} The largest human chromosome, containing many genes involved in development, cell function, and disease. Extra, missing, or rearranged material can cause genetic disorders.

\medterm{Chromosome 13} A human chromosome containing genes important for development and cell function. Abnormal chromosome 13 number or structure can cause disorders such as trisomy 13.

\medterm{Chromosome 2} A large human chromosome containing many genes involved in growth, development, and normal body function. Changes in its number or structure can cause genetic disease.

\medterm{Chromosome 3} A human chromosome containing genes involved in development, immune function, and cell regulation. Deletions, duplications, or other changes may cause genetic disorders or cancer.

\medterm{Chromosome 4} A human chromosome containing many genes involved in development and cellular function. Structural or numerical abnormalities can cause congenital or inherited conditions.

\medterm{Chromosome 5} A human chromosome containing genes involved in development, immunity, and cell regulation. Loss of part of its short arm causes the genetic condition cri-du-chat syndrome.

\medterm{Chromosome X} One of the human sex chromosomes. Most females have two X chromosomes and most males have one X and one Y; genes on the X chromosome can cause X-linked disorders when altered.

\medterm{Chromosome Y} The sex chromosome typically present in most males and absent in most females. It contains the sex-determining \textit{SRY} gene and other genes important for testis development and sperm production.

\medterm{Chronic Idiopathic Constipation} Long-lasting difficult or infrequent bowel movements without an identified structural or metabolic cause. Symptoms may include hard stools, straining, incomplete emptying, or abdominal discomfort; treatment includes diet, fluids, activity, and selected medicines.

\medterm{Chronic Obstructive Pulmonary Disorder (COPD)} A chronic lung disease with persistent airflow limitation, usually involving emphysema, chronic bronchitis, or both. Smoking and other inhaled exposures are important causes; symptoms include breathlessness, cough, and sputum.

\medterm{Circulating Tumor Cell} A cancer cell that has detached from a tumor and entered the bloodstream. Testing may help study cancer spread or treatment response, but clinical use depends on the specific validated test and cancer.

\medterm{Climate Change} Long-term changes in temperature, precipitation, and weather patterns, largely driven by human-produced greenhouse gases. Health effects include heat illness, worsened air quality, infectious-disease changes, injuries, and effects on food, water, and mental health.

\medterm{Clinical Diagnostics} The use of history, examination, laboratory tests, imaging, pathology, and other methods to identify or evaluate a health condition and guide care.

\medterm{Clinical Imaging} Medical imaging used to visualize anatomy, physiology, or disease, including radiography, ultrasound, computed tomography, magnetic resonance imaging, and nuclear medicine.

\medterm{Club Foot} A congenital foot deformity in which the foot is turned inward and downward and the calf may be smaller. It is usually treated early with serial casting, bracing, and sometimes surgery.

\medterm{Coffee} A beverage containing caffeine and other compounds. Moderate intake is tolerated by many adults, but effects vary; excessive caffeine can cause insomnia, anxiety, palpitations, or gastrointestinal symptoms.

\medterm{Cognitive Behavioural Therapy} A structured psychological treatment that helps a person identify and change unhelpful thoughts and behaviors. It is used for conditions such as anxiety, depression, insomnia, pain, and some eating disorders.

\medterm{Colony-Stimulating Factors} Growth factors that stimulate bone marrow to produce specific blood cells. Medicines such as granulocyte colony-stimulating factor can prevent or treat some forms of chemotherapy-related neutropenia.

\medterm{Colorectal Polyp} A growth arising from the lining of the colon or rectum. Many are benign, but some adenomatous or serrated polyps can become cancerous over time and are removed during colonoscopy.

\medterm{Comet Assay} A laboratory technique that measures DNA strand breaks in individual cells by electrophoresis, producing a comet-like pattern. It is mainly used in research and toxicology rather than as a routine clinical test.

\medterm{Compression Bandage} A bandage that applies controlled pressure to a limb or body part. It may reduce swelling, support injured tissue, or help manage venous disease, but improper use can impair circulation or damage skin.

\medterm{Congenital Birth Defect} A structural or functional abnormality present before or at birth. Causes may include genetic changes, infections, nutritional deficiencies, medicines, environmental exposures, or unknown factors.

\medterm{Congenital Heart Defect} A structural abnormality of the heart or major vessels present at birth. It may alter blood flow and range from a small defect needing monitoring to a serious condition requiring urgent treatment or surgery.

\medterm{Conjugated Linoleic Acid} A group of related fats found mainly in meat and dairy products and sold as supplements. Claims about weight loss or muscle improvement are inconsistent, and supplements may adversely affect blood sugar or liver function.

\medterm{Contact Lenses} Thin lenses placed directly on the eye to correct vision or deliver selected treatment. Proper fitting, handwashing, cleaning, and replacement are essential because misuse can cause corneal infection, inflammation, or vision-threatening injury.

\medterm{Coronavirus Disease COVID-19} Illness caused by SARS-CoV-2 infection. It may cause no symptoms, a mild respiratory illness, pneumonia, or severe disease, and some people develop symptoms that persist after the initial infection.

\medterm{Cosmetic Surgery} Surgery intended primarily to change or improve appearance rather than treat disease or restore function. It still carries risks such as infection, bleeding, anesthesia complications, scarring, and unsatisfactory results.

\medterm{Counterfeit Drug} A product falsely labeled as an approved medicine or made to imitate one. It may contain the wrong ingredient, too little or too much active drug, contaminants, or no active drug and can cause treatment failure or poisoning.

\medterm{Cow's Milk Allergy} An immune reaction to proteins in cow's milk, most common in infants and young children. It may cause hives, vomiting, wheezing, eczema, diarrhea, or anaphylaxis and is different from lactose intolerance.

\medterm{Cracked Heels} Dry, thickened skin at the heel that develops fissures, sometimes deep enough to bleed or become infected. Moisturizers, gentle filing, and treatment of contributing skin disease can help; painful or infected cracks need assessment.

\medterm{Cramp} A sudden, involuntary, often painful muscle contraction. Cramps may result from exertion, dehydration, electrolyte changes, pregnancy, medicines, nerve disease, or reduced blood flow.

\medterm{Cranberries} Red fruits containing fiber and plant compounds. Cranberry products may modestly reduce recurrent urinary infections in some people but do not treat an established infection; they can interact with warfarin and other medicines.

\medterm{Craniofacial} Relating to the skull and face. Craniofacial conditions may involve bones, soft tissues, teeth, eyes, ears, or airway and may be congenital, traumatic, developmental, or acquired.

\medterm{Craniofacial Microsomia} A birth condition in which one side of the face develops less fully than the other. It may affect the jaw, ear, eye, facial muscles, and nearby tissues, with severity ranging from mild to complex.

\medterm{CRIA Syndrome} A rare inherited inflammatory disorder caused by a change in the RIPK1 gene. Repeated attacks of fever, swollen lymph nodes, muscle or joint pain, and inflammation often begin during childhood.

\medterm{CRISPR} A gene-editing method that uses a guide molecule and an enzyme to find and change a chosen DNA sequence. It is used in research and some treatments but can make unintended changes and requires careful monitoring.

\medterm{Cruciate Ligament} One of two strong bands crossing inside the knee that stabilize forward, backward, and twisting movement. The anterior cruciate ligament is commonly torn during sports or a sudden change in direction.

\medterm{CTLA-4} A protein on some immune cells that acts as a brake on T-cell activity. Blocking this brake can strengthen an immune attack against cancer but may also cause inflammation in healthy organs.

\medterm{Cupping Therapy} A complementary practice that places cups on the skin to create suction, sometimes after small skin cuts. It may leave bruises or marks and can cause burns, infection, or blood loss; evidence of benefit varies by condition.

\medterm{Curcumin} A yellow compound in turmeric that has anti-inflammatory effects in laboratory studies. Evidence that supplements treat disease is limited or variable, and concentrated products can interact with medicines or cause stomach or liver problems.

\medterm{Cutaneous Vasculitis} Inflammation of small blood vessels in the skin, causing red or purple spots, raised purpura, blisters, sores, or pain. Medicines, infections, autoimmune disease, or inflammation elsewhere in the body may be responsible.

\medterm{Cyclin-Dependent Kinase} An enzyme that helps control when a cell grows and divides after joining with a cyclin protein. Abnormal activity can promote cancer and is blocked by some targeted treatments.

\medterm{Cystinosis} A rare inherited disorder in which cystine builds up inside cells and damages organs. It often affects the kidneys and eyes and may also affect growth, thyroid, muscles, or other tissues; lifelong specialist care is needed.


\medterm{1p36 Deletion Syndrome} A genetic disorder caused by loss of a piece of chromosome 1. It may cause developmental and intellectual disability, small head size, low muscle tone, seizures, hearing loss, heart problems, and distinctive facial features; severity varies.

\synonyms
Monosomy 1p36; Chromosome 1p36 Deletion Syndrome

\medterm{22q11.2 Deletion Syndrome} A genetic condition caused by a missing piece of chromosome 22. It may cause heart defects, palate problems, reduced immune defenses, low calcium, distinctive facial features, learning or developmental difficulties, and increased risk of mental illness. The effects differ widely between people.

\synonyms
DiGeorge Syndrome Spectrum; Velocardiofacial Syndrome (VCFS); 22q11.2DS

\medterm{D And C} Dilation and curettage, a procedure that opens the cervix and removes tissue from the uterine lining for treatment or diagnosis.

\medterm{D And E} Dilation and evacuation, a procedure that gently opens the cervix and removes pregnancy tissue from the uterus using suction and instruments. It is performed by trained clinicians for selected pregnancy-related indications.

\medterm{D \& C} Dilation and curettage, a procedure in which the cervix is widened and tissue is removed from the uterine cavity for treatment or diagnosis.

\medterm{D2 Receptor} A protein on cells that responds to dopamine, a brain chemical involved in movement, motivation, and reward. Some medicines work partly by blocking or stimulating this receptor.

\medterm{Da Costa's Syndrome} An older term for recurring chest discomfort, awareness of heartbeat, tiredness, and breathlessness without a clear structural heart problem. Current evaluation should still exclude heart, lung, blood, and anxiety-related causes.


\medterm{Dacarbazine} An alkylating chemotherapy medicine used for Hodgkin lymphoma, melanoma, and selected other cancers; important toxicities include severe nausea, bone-marrow suppression, and liver injury.

\medterm{Daclizumab} A monoclonal antibody that blocks the interleukin-2 receptor on activated immune cells. It was used for selected immune conditions but was withdrawn or restricted in many settings because of serious liver and other risks.

\medterm{Dacomitinib} A targeted medicine that blocks several members of the epidermal-growth-factor receptor family and is used for selected EGFR-mutated lung cancers. Diarrhea, skin reactions, mouth sores, and lung or liver problems may occur.

\medterm{Dacrocytes (Teardrop Cells)} Teardrop-shaped red blood cells seen on a blood smear. They may occur when red blood cells pass through scarred or crowded bone marrow, and can be seen in myelofibrosis, bone-marrow cancers, and severe thalassemia.

\synonyms
Teardrop Cells

\medterm{Dacryoadenitis} Inflammation of a tear-producing gland, causing painful swelling near the outer upper eyelid. It may be caused by infection, inflammation, or a condition affecting the whole body.

\medterm{Dacryocystitis} Infection or inflammation of the tear sac near the inner corner of the eye, usually because the tear-drainage passage is blocked. It can cause pain, redness, swelling, and pus and may need prompt treatment.

\medterm{Dacryocystocele} A fluid-filled swelling near the inner corner of the eye caused by blockage and widening of the tear-drainage system. In a newborn it may interfere with breathing if it extends into the nose and needs prompt assessment.

\medterm{Dacryocystography} An imaging test in which contrast material is put into the tear-drainage passages to show narrowing, blockage, abnormal connections, or leakage. It may help plan treatment for persistent tearing or repeated tear-sac infection.

\medterm{Dacryocystorhinostomy} Surgery creating a new passage from the tear sac into the nose to bypass a blocked tear duct. It is used for persistent tearing or repeated tear-sac infections.

\medterm{Dacryocyte Analysis} Examination or counting of teardrop-shaped red blood cells on a blood smear; numerous dacryocytes may occur with marrow fibrosis, infiltration, severe anemia, or splenic disease.

\medterm{Dacryolithiasis} Formation of a stone or concretion within the lacrimal drainage system, which may cause excessive tearing, discharge, recurrent dacryocystitis, or obstruction.

\medterm{Dacryostenosis} Narrowing or blockage of the passages that drain tears from the eye into the nose. It can cause constant tearing, sticky discharge, or repeated infection of the tear sac.

\medterm{Dactinomycin} An anticancer medicine that binds DNA and interferes with gene activity and cell growth. It is used for selected cancers and can cause bone-marrow suppression, mouth sores, nausea, and tissue injury if it leaks outside a vein.

\medterm{Dactylitis} Swelling and inflammation of an entire finger or toe, producing a sausage-like appearance. It can occur with sickle-cell disease, psoriatic arthritis, infection, or other inflammatory conditions.

\medterm{Dactylomegaly} Abnormal enlargement of a finger or toe, which may occur as an isolated finding or as part of an endocrine, skeletal, vascular, or inflammatory disorder.

\medterm{Dactylosymphysis} A birth difference in which two neighboring fingers or toes are joined together by skin, bone, or both. The connection may limit movement or affect growth and hand or foot function.

\medterm{Daidzein} A naturally occurring isoflavone found in soy and other legumes that can be converted by intestinal bacteria into related compounds.

\medterm{Daily Bowel Care Program} A planned routine for preventing and treating constipation or fecal incontinence, often including timing, fluids, diet, activity, medication, and rectal measures when prescribed.

\medterm{Daily Value} A reference amount used on U.S. Nutrition Facts labels to show how much a nutrient in one serving contributes to a typical daily diet.

\medterm{Dakin Solution} A dilute sodium hypochlorite solution used as a topical antiseptic for selected contaminated or infected wounds. Concentration and exposure time must be controlled because stronger solutions can injure viable tissue.

\medterm{Dalteparin} An injectable blood thinner used to prevent or treat clots in veins and the lungs. It reduces clotting through the body’s natural anticoagulant system; bleeding is the main risk, and kidney disease may require dose adjustment.

\medterm{Dalteparin Sodium} The sodium salt of dalteparin, a low-molecular-weight heparin used to prevent or treat venous thromboembolism; bleeding is the principal serious adverse effect.

\medterm{Daltonism} Color-vision deficiency, especially difficulty distinguishing certain colors because of abnormal or absent retinal cone function.
\medterm{DALY} Disability-adjusted life year, a measure combining years of healthy life lost because of premature death and disability.

\medterm{Damage To The Amygdala} Damage to the amygdala can affect fear processing, emotional learning, memory, threat recognition, and social behavior. Causes include trauma, stroke, infection, tumor, epilepsy, or degenerative disease, and effects depend on the side and extent of injury.

\medterm{Damage-Control Resuscitation} Emergency treatment for severe injury that quickly controls bleeding, gives appropriate blood products, prevents dangerous cooling, and corrects clotting problems while avoiding harmful excess fluid.

\medterm{Damage-Control Surgery} A planned series of operations for a critically ill or severely injured person who cannot safely tolerate a long repair. The first operation controls bleeding or contamination; after stabilization, a later operation completes the repair.

\medterm{Damus-Kaye-Stansel Procedure} A heart operation for selected children with a single working ventricle and a blocked or narrowed route from the ventricle to the aorta. It connects the pulmonary artery to the aorta to create a wider route for blood leaving the heart.

\synonyms
DKS Procedure; Damus-Stansel-Kaye Operation

\medterm{Danazol} A synthetic androgen with antigonadotropic effects, used selectively for endometriosis, hereditary angioedema, or other conditions; androgenic, lipid, liver, and thrombotic adverse effects limit use.

\medterm{Dandruff} A common scalp condition causing small white or yellow flakes and sometimes itching. It is often related to oily skin, seborrheic dermatitis, or the scalp's reaction to the normal yeast \textit{Malassezia}; it is not usually contagious. Persistent severe flaking, redness, sores, or hair loss may indicate another scalp disorder.

\medterm{Dandruff And Scalp Conditions} Scalp flaking and itching may result from dandruff, seborrheic dermatitis, psoriasis, contact allergy, fungal infection, or head lice. Persistent redness, sores, pain, or hair loss needs medical assessment.

\medterm{Dandy Fever} An older name for dengue fever, a mosquito-borne viral infection causing sudden fever, headache, severe muscle or joint pain, nausea, or rash. Severe dengue can cause bleeding, shock, and organ failure.

\medterm{Dandy-Walker Malformation} A condition present at birth in which part of the cerebellum does not develop normally and one of the brain's fluid spaces becomes enlarged. It may cause fluid buildup in the brain, delayed development, balance problems, or seizures.

\medterm{Dandy-Walker Syndrome} A birth difference in the cerebellum and nearby fluid spaces of the brain. It may cause enlarged head size from fluid buildup, delayed development, balance problems, seizures, or no symptoms.

\medterm{Dane Particle} The complete infectious form of hepatitis B virus. It contains the virus's genetic material inside a protein shell and outer envelope, allowing it to enter liver cells and reproduce.
\medterm{Dangerous Abbreviations} Abbreviations that can be misread and cause medication or clinical errors; safer practice is to write the complete term when an abbreviation is ambiguous.

\medterm{Danon Disease} A rare inherited disorder, usually caused by an LAMP2 gene change, in which cells cannot properly recycle materials. It may cause heart-muscle disease, muscle weakness, intellectual disability, or eye problems.

\medterm{Dantrolene} A skeletal-muscle relaxant that reduces calcium release from the sarcoplasmic reticulum; it is the specific treatment for malignant hyperthermia and may cause weakness or liver injury.

\medterm{Dapsone} A sulfone antimicrobial used for leprosy, selected dermatitis, and other infections; haemolysis, especially with G6PD deficiency, methaemoglobinaemia, agranulocytosis, and severe skin reactions are important risks.

\medterm{Daptomycin} An antibiotic used for selected serious infections caused by gram-positive bacteria, including resistant strains. It can damage muscles and kidneys, so treatment requires appropriate testing and monitoring.

\medterm{Daptomycin Resistance} Reduced or absent response of bacteria to daptomycin, often because changes in the bacterial membrane prevent the antibiotic from working. It can develop during treatment and limits options for serious infections.

\medterm{Darbepoetin Alfa} A long-acting medicine that stimulates the bone marrow to make more red blood cells. It is used for selected anemia related to kidney disease or cancer treatment, but can raise blood pressure and the risk of blood clots.

\medterm{Darier's Disease} An inherited skin disorder causing firm, greasy, crusted bumps in areas such as the chest, back, scalp, or skin folds. Heat and sweating may worsen it, and nails or mouth lining can also be affected.

\medterm{Dark Adaptation} The gradual adjustment that allows the eyes to see better after moving from bright light into darkness. Poor dark adaptation may result from retinal disease, vitamin A deficiency, aging, or some medicines.

\medterm{Dartos} A thin muscle layer beneath the skin of the scrotum. It tightens or relaxes the scrotal skin in response to temperature and other signals, helping regulate the position of the testicles.

\medterm{Darunavir} An antiretroviral medicine that blocks HIV protease, an enzyme needed to make new viruses. It is used with other HIV medicines and has important interactions and possible liver or skin reactions.

\medterm{Darwin's Tubercle} A harmless small bump on the upper rim of the outer ear. It is a common normal variation in ear shape and usually needs no treatment.

\medterm{Dasatinib} A targeted cancer medicine that blocks abnormal tyrosine-kinase signals, especially those caused by BCR::ABL1. It treats selected blood cancers and may cause low blood counts, fluid around the lungs, bleeding, or abnormal heart rhythms.

\medterm{DASH Diet} An eating pattern designed to lower blood pressure by emphasizing vegetables, fruits, whole grains, beans, nuts, and low-fat dairy while limiting salt, saturated fat, and sugary foods. Individual needs may require adjustment.


\medterm{Dating Ultrasound} An ultrasound examination, usually performed in early pregnancy, used to estimate gestational age and calculate the expected date of delivery.
First-trimester crown--rump length is generally the most accurate ultrasound measurement for dating.

\synonyms
Pregnancy dating scan; gestational-age ultrasound.

\medterm{Daubenton's Plane} An anatomical reference plane extending from the lower margin of the orbit to the upper
margin of the external auditory opening, used in describing skull orientation.

\medterm{Daunorubicin} A chemotherapy medicine that damages cancer-cell DNA and is used mainly for selected blood cancers. It can lower blood counts and, especially after repeated doses, weaken or damage the heart.

\medterm{Daunorubicin Hydrochloride} The hydrochloride salt of daunorubicin, an anthracycline chemotherapy drug that interferes with DNA and is used mainly for acute leukemias. Important risks include bone-marrow suppression and heart damage.

\medterm{Dav Regimen} An abbreviation used for more than one cancer-treatment combination. In older leukemia reports, DAV meant daunorubicin, cytarabine, and etoposide; in newer reports, it may mean daunorubicin, cytarabine, and venetoclax. The source must identify the drugs and cancer before the regimen is used.

\medterm{David Procedure} An operation that replaces an enlarged or weakened aortic root while keeping the person's own aortic valve. The valve is supported inside a graft and the coronary arteries are reattached; bleeding, stroke, valve leakage, or recurrent aortic disease are possible.

\synonyms
David Operation; Valve-Sparing Root Replacement; Reimplantation of Aortic Valve

\medterm{Daw} An abbreviation meaning dispense as written, indicating that the pharmacist should dispense the prescribed product rather than automatically substitute a generic equivalent when permitted by law.

\medterm{Dawn Phenomenon} A rise in blood glucose during the early morning caused by normal hormone changes that help the body wake up. It may be more noticeable in people with diabetes whose insulin supply cannot match the rise.

\medterm{Day Hospital} A facility that provides multidisciplinary medical, psychiatric, rehabilitative, or diagnostic care during the day without overnight admission. It can support intensive treatment while allowing the patient to return home.

\medterm{Day-Blindness} Poor vision in bright light, also called hemeralopia. It can result from retinal disease, cone dysfunction, or other eye disorders.

\medterm{D\&C Instruments} Instruments used to dilate the cervix and remove tissue from the uterus during dilation and curettage.

\medterm{Dcc Gene} A gene that helps guide developing nerve fibers and may also help control cell growth. Harmful changes in this gene can cause some developmental nerve disorders and may contribute to certain cancers.

\medterm{DCIS} Alternate index label; see \textit{Ductal carcinoma in situ}.


\medterm{DDH} Alternate index label; see \textit{Hip dysplasia}.

\medterm{D-Dimer} A substance released into the blood when a blood clot breaks down. A high level can occur with a blood clot, but also with infection, pregnancy, recent surgery, cancer, or other conditions. A normal result can help rule out a clot in carefully selected patients whose initial risk is low or moderate.

\medterm{D-Dimer Test} A blood test measuring a substance released when blood clots are broken down. A normal result can help rule out a clot in carefully selected low-risk people; a high result has many possible causes and does not prove a clot.

\synonyms
D-dimer blood test

\medterm{De Lange Syndrome} A congenital developmental disorder, also called Cornelia de Lange syndrome, involving characteristic facial features, growth and developmental delay, limb abnormalities of variable severity, and intellectual disability.

\synonyms
Cornelia de Lange syndrome

\medterm{De Musset Sign} Repeated head nodding that matches the heartbeat. It can occur when the aortic valve leaks severely, but it is uncommon and cannot confirm the diagnosis without examination and heart tests.

\medterm{De Novo} Newly arising in an individual rather than inherited from either parent; the term is commonly used for genetic variants or disease presentations.

\medterm{De Novo Variant} A genetic change found in a person that was not found in either biological parent. It usually begins in an egg, sperm, or very early embryo and may cause disease; testing parents can help estimate the chance in future pregnancies.

\medterm{De Quervain Tendinitis} Alternate name; see \textit{De Quervain Tenosynovitis}.

\medterm{De Quervain Tenosynovitis} Painful swelling and narrowing around the tendons on the thumb side of the wrist. It causes pain with thumb movement, gripping, or lifting and may cause tenderness or swelling near the base of the thumb.

\synonyms
Tenosynovitis, De Quervain's

\medterm{De Quervain's Tendonitis} Alternate spelling; see \textit{De Quervain Tenosynovitis}.

\medterm{De Quervain's Tenosynovitis} Alternate spelling; see \textit{De Quervain Tenosynovitis}.

\medterm{De Winter ECG Pattern} An ECG pattern that can signal a sudden severe blockage in a major artery supplying the front of the heart, even without the usual ST elevation. It is a heart-attack emergency requiring immediate specialist assessment and treatment; exercise testing must not be used.

\synonyms
de Winter Pattern; de Winter T Waves; de Winter Sign

\medterm{Deacetylase Inhibitor} A medicine or research compound that blocks enzymes removing chemical groups called acetyl groups from proteins. Some alter gene activity and are used for selected cancers; effects and risks depend on the specific medicine.

\medterm{Dead Space} The part of a breath that does not exchange oxygen and carbon dioxide with the blood, including the nose, throat, and larger airways. It increases when air reaches lung areas with too little blood flow.

\medterm{Dead Space Ventilation} Breathing that moves air through areas of the lungs where little or no blood is available for gas exchange. It may increase with a pulmonary embolism, low blood flow, or emphysema and can make breathing less effective.

\medterm{Dead Time} The brief period after an event during which a detector or recording system cannot register another event, limiting the measured count rate.


\medterm{Deafferentation} Loss or interruption of nerve signals carrying sensation from a body area to the brain or spinal cord. It may cause numbness, unusual sensations, poor awareness of position, or persistent nerve pain.

\medterm{Deafness} Complete or severe loss of hearing in one or both ears. The problem may involve sound passing through the outer or middle ear, the inner ear or hearing nerve, or more than one of these areas.

\medterm{Deaner} A former brand name for deanol, also called dimethylaminoethanol, an older product once marketed for behavioral or learning problems. It is not a standard modern treatment, and evidence for medical benefit is limited.

\medterm{Death} The permanent end of the body's vital functions. It means that the heart, breathing, and brain functions have stopped and cannot be restored; medical and legal tests used to confirm death depend on the circumstances.

\medterm{Death Rate} The number of deaths in a defined population during a specified period, usually expressed per population size and used in public-health comparisons.

\medterm{Death Rattle} Noisy, rattling breathing caused by secretions accumulating in the throat or upper airways, often occurring near the end of life when swallowing and coughing weaken.

\medterm{Death Sudden Cardiac (Sudden Cardiac Death)} Alternate index label for Sudden Cardiac Death.

\medterm{Deauville Criteria} A five-point scale used with a PET scan to compare lymphoma activity with normal organs. The score helps clinicians judge whether treatment is working, but it must be interpreted with the type of lymphoma and other findings.

\medterm{Deaver Retractor} A hand-held surgical instrument with a broad, curved blade used to gently move organs or tissues aside. It is often used during chest or abdominal surgery to improve the surgeon's view.

\medterm{Debility} General physical weakness or reduced ability to function, often caused by illness, aging, poor nutrition, prolonged inactivity, or recovery from major stress. The underlying cause should be assessed when weakness persists.

\synonyms
Weakness

\medterm{Debride} To remove dead, damaged, infected, or contaminated tissue from a wound or body site to support healing and control infection.

\medterm{Debridement} Removal of dead, damaged, infected, or contaminated tissue from a wound. It may be done with instruments, irrigation, or other methods to reduce infection, reveal healthy tissue, and support healing.

\medterm{Debulk} To reduce the size or amount of a tumor or other abnormal tissue, usually by surgery or another treatment, without necessarily removing every abnormal cell.

\medterm{Debulking} Removal of as much of a tumor as possible when removing all of it is unsafe or impossible. The aim is to reduce symptoms or improve the effect of later treatment; benefit depends on the cancer and situation.

\medterm{Decalcification} Loss of calcium salts from bone, teeth, or other mineralized tissue. In teeth it may weaken enamel; in bone it can contribute to reduced strength.

\medterm{Decannulation} Removal of a tube, especially a tracheostomy tube, after a person can breathe safely, protect the airway, and clear secretions without it. The site is monitored for bleeding, narrowing, or breathing difficulty.

\medterm{Decanoic Acid} A ten-carbon saturated fatty acid, also called capric acid, found in some fats and studied for antimicrobial, metabolic, and antiseizure effects.

\medterm{Decapitation} Separation of the head from the body, seen in severe trauma (vehicle accidents, train
injuries, explosions), industrial accidents, hanging with judicial execution, or
homicide. Complete decapitation is almost always rapidly fatal due to exsanguination and
brain ischemia.

\medterm{Decay Correction} A calculation that adjusts a radioactive measurement for the gradual loss of activity over time. It allows imaging or laboratory results obtained at different times to be compared accurately.

\medterm{Deceased Donor} Individual who has died (usually from brain death or cardiac death) and donated organs, tissues, or both for transplantation.

\medterm{Decenoic Acid} An unsaturated ten-carbon fatty acid; its biologic effects depend on the position and geometry of its double bond and on tissue and dietary context.

\medterm{Decerebrate Posture} Abnormal straightening of the arms and legs caused by severe injury to the brainstem or nearby brain structures. It is a sign of serious neurologic injury and may occur after head injury, bleeding, or increased pressure in the skull.

\medterm{Decerebrate Posturing} Abnormal straightening of the arms and legs, sometimes with inward turning of the arms, caused by severe brain or brainstem injury. It indicates a medical emergency, but the posture alone cannot predict recovery.

\medterm{Decerebrate Rigidity} Severe, abnormal straightening of the arms and legs, sometimes with the arms turned inward. It usually indicates serious brain or brainstem injury or dangerously high pressure inside the skull and requires emergency assessment.

\medterm{Decerebrate State} A severe neurologic posture marked by extension and inward rotation of the arms and extension of the legs, usually indicating serious injury to the brainstem or structures above it.

\medterm{Decerebration} Severe brain dysfunction that causes abnormal straightening of the arms and legs. It usually results from major brain or brainstem injury or increased pressure inside the skull and is a medical emergency.

\medterm{Decibel} A unit used to describe sound intensity. Higher values mean louder sound; repeated or very loud exposure can damage hearing, depending on the level and duration.




\medterm{Deciding To Have Knee Or Hip Replacement} A decision based on pain, functional limitation, imaging, nonoperative treatment response, surgical risk, expectations, and shared discussion with an orthopedic clinician.

\medterm{Deciding To Quit Drinking Alcohol} The process of planning to stop or reduce alcohol use. Assessment should consider dependence, withdrawal risk, support, and treatment options; people with heavy use may need medical supervision.

\medterm{Decidua} The specialized lining of the uterus during pregnancy. It is shed after delivery and helps support the pregnancy and placenta.

\medterm{Deciduitis} Inflammation of the specialized uterine lining that forms during pregnancy. It may be associated with infection or other pregnancy complications and can require assessment when pain, fever, or bleeding occurs.

\medterm{Deciduoma} An old term for a mass-like thickening of the uterine lining caused by pregnancy hormones. It is not a true tumor and usually disappears when the hormonal stimulation ends.

\medterm{Deciduosis} Pregnancy-related tissue from the uterine lining growing in an unusually large area or outside the uterus, such as on the cervix or pelvic surfaces. It can look like a tumor during examination or surgery but is usually benign.

\medterm{Deciduous Teeth} The 20 primary teeth that erupt during childhood and are gradually replaced by permanent teeth. Caries, infection, trauma, or premature loss can affect nutrition, speech, space for permanent teeth, and oral development.

\medterm{Deciduous Tooth} A primary or baby tooth that is eventually shed and replaced by a permanent tooth. Early decay or loss can affect chewing, speech, and the space needed for permanent teeth.
\medterm{Decision Making} The process of comparing information, options, benefits, risks, and preferences to choose an action or plan, including shared decision-making in healthcare.

\medterm{Decision-Making Capacity} The ability to understand information about a medical decision, consider the possible results, compare choices, and communicate a choice. It applies to a specific decision, may change over time, and is different from legal competence.

\medterm{Decitabine} A cancer medicine that changes abnormal gene activity and interferes with growth of blood-forming cancer cells. It is used for selected blood cancers and may cause low blood counts, infection, bleeding, or fatigue.

\medterm{Declarative Memory} Memory for facts, information, and personal events that a person can consciously remember and describe.

\medterm{Decoction} A preparation made by simmering plant material in water to extract soluble substances. The strength and safety depend on the plant, preparation, dose, contamination, and interactions with medicines.

\medterm{Decompensated Cirrhosis} Advanced liver scarring that has begun to cause problems such as abdominal fluid buildup, internal bleeding, jaundice, or confusion from waste products. It requires specialist care and may require assessment for liver transplantation.

\medterm{Decomposition} The natural breakdown of a dead body after death. Cells and tissues first break down themselves, followed by the action of microorganisms; the speed and appearance depend on temperature, moisture, and other conditions.

\medterm{Decompression} Reduction of pressure on a tissue or body compartment, such as removing pressure from a spinal nerve, draining an obstructed organ, or treating pressure changes after diving.

\medterm{Decompression Sickness} Illness caused by nitrogen bubbles forming in the body when pressure falls too quickly, usually after diving. It may cause joint pain, skin changes, breathing problems, weakness, confusion, or paralysis and needs emergency care.

\medterm{Deconditioning} Loss of strength, fitness, balance, and function after inactivity, illness, or prolonged bed rest. It increases fall risk but can often improve gradually with safe movement, exercise, nutrition, and rehabilitation.

\medterm{Decongestant} A medicine that reduces swelling inside the nose and relieves blocked breathing. Some can raise blood pressure, cause fast heartbeat, or cause rebound congestion, so use should follow instructions and health precautions.

\medterm{Decontamination} Steps that remove or reduce a harmful substance from the skin, eyes, clothing, or digestive tract and prevent further absorption. The correct method depends on the substance; some exposures require washing, irrigation, or specialist treatment.

\medterm{Decorin} A small protein-sugar molecule found around cells. It helps organize collagen, influences tissue growth and repair, and may affect scar formation.

\medterm{Decorticate Posture} Abnormal bending of the arms with straightening of the legs, caused by severe injury to the brain. It is a sign of serious brain damage and reduced consciousness.

\medterm{Decorticate Posturing} Abnormal bending of the arms toward the chest with straightening of the legs, caused by severe brain injury. It signals impaired consciousness and requires emergency assessment; posture alone cannot predict outcome.

\medterm{Decreased Alertness} Decreased alertness is reduced awareness or responsiveness that may result from metabolic illness, intoxication, infection, seizure, stroke, trauma, or sleep-related causes.

\medterm{Decreased Breath Sounds} Softer-than-normal sounds heard when listening to the lungs, suggesting reduced air movement. Causes include fluid or air around the lung, partial lung collapse, blockage, or shallow breathing; sudden one-sided loss can be an emergency.

\medterm{Decreased Concentration} Reduced ability to sustain attention or mental focus; causes include sleep loss, stress, medications, mood disorders, substance use, and neurologic or medical illness.

\medterm{Decreased Dentin} Reduced thickness or formation of dentin, the mineralized tissue beneath tooth enamel; causes include developmental defects, wear, erosion, caries, or dental treatment.

\medterm{Decreased Libido} Reduced sexual desire that may be related to relationships, mood, stress, medications, hormones, chronic illness, pain, or other health factors.

\medterm{Decreased Tear Production} Alternate index label; see \textit{Dry eyes}.

\medterm{Decubitus} A lying or reclining position; the term also appears in decubitus ulcer, a pressure injury caused by prolonged pressure and shear.

\medterm{Decubitus Ulcer} Damage to skin and deeper tissue caused by prolonged pressure, sliding, or rubbing, usually over a bony area. It may begin with persistent redness and progress to an open wound, especially in people with limited movement.

\medterm{Deep} Located farther inside the body or a structure, away from its surface. For example, deep muscles lie beneath muscles closer to the skin.
\medterm{Deep Brain Stimulation} A treatment in which implanted electrodes deliver controlled electrical impulses to selected brain regions. It is used for conditions such as Parkinson disease, essential tremor, and some forms of dystonia.

\medterm{Deep Femoral Artery} The profunda femoris artery, a major branch of the femoral artery that supplies the deep muscles of the thigh and gives rise to perforating branches.

\medterm{Deep Infiltrating Endometriosis} Endometriosis that grows deeply into tissues around the uterus, bowel, bladder, or other pelvic structures. It may cause severe period pain, pain with sex or bowel movements, urinary symptoms, infertility, and sometimes requires complex specialist treatment.

\medterm{Deep Neck Space Infection} A bacterial infection in one or more of the tissue spaces deep in the neck. It may cause neck swelling, severe throat pain, fever, difficulty swallowing, or breathing problems and can spread to the chest or threaten the airway.

\synonyms
DIE; Deep Endometriosis

\medterm{Deep Sleep} The deepest stage of normal non-REM sleep, when a person is difficult to awaken and the body carries out important physical restoration. Too little may contribute to daytime sleepiness and poor function.
\medterm{Deep Tendon Reflex} An automatic muscle response caused by tapping a tendon with a reflex hammer. It helps assess nerves and the spinal cord. Common examples are the knee, ankle, biceps, and triceps reflexes.

\medterm{Deep Vein Thrombosis} A blood clot in a deep vein, usually in the leg, caused by reduced blood flow, increased clotting tendency, or injury to the vein wall. It may cause one-sided leg swelling, pain, warmth, or discoloration, although some clots cause no symptoms. Part of the clot can break off and travel to the lungs, causing a potentially life-threatening pulmonary embolism.

\synonyms
DVT (Deep Vein Thrombosis)

\medterm{Deep Vein Thrombosis Ultrasound} An ultrasound examination that checks whether deep veins compress normally and whether blood flows through them. A vein that does not compress may contain a clot; results are interpreted with symptoms and other tests.

\medterm{Deep Venous Reflux} Backward flow of blood through weakened valves in deep veins. It can cause long-term leg swelling, heaviness, skin changes, pain, and ulcers from increased pressure in the veins.

\medterm{Deep Venous Thrombosis} Formation of a blood clot in a deep vein, usually in the leg, which can cause swelling and travel to the lungs.

\medterm{De-Escalation} A set of verbal and nonverbal techniques used to reduce agitation and prevent violence
while preserving dignity and safety. It includes calm communication, appropriate
personal space, clear limits, environmental modification, and offering choices.

\medterm{Defaecation} The expulsion of stool from the rectum through the anus. Difficulty, pain, urgency, incontinence, bleeding, or a sustained change in bowel habit may indicate constipation, anorectal disease, neurologic dysfunction, or another condition.

\medterm{Defecation} The process of passing stool from the rectum through the anus. It involves signals from a full rectum, contraction of the bowel, and relaxation of the anal muscles, followed by voluntary control of timing.

\medterm{Defecation Syncope} Fainting during or immediately after defecation, usually a situational reflex involving straining, reduced venous return, and a temporary fall in blood pressure.

\medterm{Defecography} An imaging test that records the rectum and pelvic floor while a person passes a soft contrast material. It can show prolapse, a bulge, inward telescoping, or poor coordination causing difficult stool passage.

\medterm{Defective Virus} A virus missing genetic information needed to reproduce independently. It can multiply only with help from another virus or host system and may influence infection without always causing disease on its own.

\medterm{Defence Mechanisms} Mental processes that help a person manage conflict, anxiety, or distress, often without conscious awareness. They may be helpful or harmful depending on how rigidly they are used and how they affect daily life.

\medterm{Defense Mechanism} A psychological process, often unconscious, that helps a person manage conflict, anxiety, or distress; its usefulness depends on context and flexibility.

\medterm{Defense Mechanisms} Psychological processes, often outside conscious awareness, that help a person manage distress
or conflict. Examples include denial, projection, repression, humor, and sublimation.

\medterm{Defense Mechanisms (Forensic)} Injuries to the hands, arms, or other body areas that may occur when a person tries to protect against an assault. Their meaning must be assessed with the history and all other evidence.
\synonyms
Forensic defense injuries

\medterm{Defense Wounds} Injuries to the hands, wrists, arms, or other areas that may occur when a person tries to protect against an attack. Their appearance and location can support a forensic interpretation but cannot prove how an injury happened by themselves.

\medterm{Defense Wounds (Inscribed)} Cuts or other injuries on the hands or forearms that may occur when a person attempts to protect against an attack.
The distribution and morphology of defensive injuries may help
reconstruct an assault but are not independently proof of a particular event.

\synonyms
Defensive wounds

\medterm{Defenses} The body’s protective mechanisms against infection, injury, and other threats, including physical barriers and immune responses.

\synonyms
Body defenses

\medterm{Defensin} A small protective protein made by skin, lining tissues, and immune cells. It can damage microbes and is part of the body’s early defense against infection.

\medterm{Deferasirox} An oral iron-chelating medicine used to reduce chronic iron overload when clinically indicated; kidney, liver, hearing, and gastrointestinal toxicity require monitoring.

\medterm{Deferiprone} An oral iron chelator used for transfusional iron overload; neutropenia or agranulocytosis is an important adverse effect requiring blood-count monitoring.

\medterm{Deferoxamine} A parenteral iron chelator used for severe or chronic iron overload and sometimes aluminum toxicity; treatment requires monitoring for ocular, auditory, renal, and infectious complications.

\medterm{Defibrillation} Delivery of an electrical shock to stop ventricular fibrillation or pulseless ventricular tachycardia during cardiac arrest. The shock can restore an organized rhythm and must be combined with high-quality CPR and emergency care.

\medterm{Defibrillator} A device that delivers an electrical shock to treat certain life-threatening abnormal heart rhythms. External defibrillators are used during cardiac arrest; implanted devices can detect and treat dangerous rhythms automatically.

\medterm{Defibrillator Implant} An implanted device that continuously monitors heart rhythm and can use pacing or an electrical shock to stop dangerous fast rhythms. It may be recommended for selected people at high risk of sudden cardiac arrest.
\medterm{Deficiency Diseases} Illnesses caused by too little of a needed nutrient or other substance, such as iron, iodine, vitamin D, or vitamin B12. Symptoms depend on the missing substance.

\medterm{Deficiency, FAO} An inherited disorder in which cells cannot use certain fats normally for energy. Fasting or illness may cause low blood sugar, vomiting, muscle breakdown, weakness, heart problems, or sudden metabolic illness.

\medterm{Deficiency, Iron} Inadequate iron stores or availability, which can impair hemoglobin production and cause iron-deficiency anemia, fatigue, pallor, and reduced exercise tolerance.

\medterm{Deficiency Primary Immunodeficiency (Primary Immunodeficiency Disease PIDD)} Alternate index label for Primary Immunodeficiency Disease PIDD.

\medterm{Deficiency, UDP-Glucuronosyltransferase} Reduced activity of a liver enzyme that prepares bilirubin for removal from the body. Inherited deficiency can cause intermittent mild jaundice or, in severe forms, dangerous bilirubin buildup.

\medterm{Deficient Color Vision} Alternate index label; see \textit{Color blindness}.


\medterm{Deformation} An alteration in the shape or structure of a body part, tissue, or material caused by developmental forces, injury, disease, pressure, or mechanical stress.

\medterm{Deformity} An abnormal shape, alignment, or structure of a body part caused by development, injury, disease, pressure, or unusual growth. It may affect appearance, movement, comfort, or organ function.

\medterm{Degarelix} An injected hormone-blocking medicine that rapidly lowers testosterone. It is used for advanced prostate cancer; injection-site reactions, hot flushes, bone loss, and other hormone-related effects may occur.

\medterm{Degeneration} Progressive structural or functional deterioration of cells, tissues, or organs caused by aging, disease, injury, inherited factors, or other processes.

\medterm{Degenerative Arthritis} Alternate index label; see \textit{Osteoarthritis}.

\medterm{Degenerative Disc} Degenerative disc disease describes age- or wear-related changes in spinal discs that may cause neck or back pain, stiffness, or nerve compression. Imaging changes can occur without symptoms; treatment usually begins with activity, exercise-based rehabilitation, and pain management.

\synonyms
DDD (Degenerative Disc)

\medterm{Degenerative Disc Disease} Wear, drying, or cracking of the cushioning discs between spinal bones, often related to aging or injury. It may cause back or neck pain, stiffness, or nerve pressure; imaging changes can occur without symptoms.

\medterm{Degenerative Disease} A condition in which a body tissue or organ gradually becomes damaged or loses function. Causes include ageing, inherited problems, injury, and other disease; examples include osteoarthritis and some disorders of the brain or nerves.

\synonyms
Degenerative disorder

\medterm{Degenerative Disk Disease} Alternate spelling; see \textit{Degenerative Disc}.

\medterm{Degenerative Disorder} A disease in which progressive cellular or tissue deterioration causes loss of structure or function, as in neurodegenerative, joint, retinal, or muscle disorders.

\medterm{Degenerative Joint Disease} Gradual loss of joint cartilage with changes in nearby bone and tissues, commonly meaning osteoarthritis. It can cause pain, stiffness, swelling, reduced movement, and difficulty with daily activities.

\synonyms
Osteoarthritis

\medterm{Degenerative Spine Disease} Wear-related changes in the spinal discs, joints, or bones that can cause back or neck pain, stiffness, reduced movement, or pressure on nerves. Symptoms may include leg or arm pain, numbness, weakness, or walking difficulty.

\medterm{Degenerative Spondylolisthesis} Forward slipping of one spinal bone over another because joints and disks have worn down, sometimes causing back pain, leg pain, or nerve compression.

\medterm{Deglutition} Swallowing is the coordinated movement of food, liquid, or saliva from the mouth through the throat and esophagus into the stomach. It protects the airway and begins digestion.

\medterm{Deglutition (Swallowing)} The coordinated process of moving food, liquid, or saliva from the mouth through the throat and esophagus into the stomach. Problems can cause choking, coughing, or food entering the airway.

\synonyms
Swallowing

\medterm{Degranulation} Release of stored chemicals from cells such as mast cells, white blood cells, and platelets. The chemicals help respond to injury or infection but can also cause allergy, inflammation, blood-vessel changes, or clotting.

\medterm{Degrees Fahrenheit} A unit on the Fahrenheit temperature scale. In clinical care, body temperature may be reported in Fahrenheit; 98.6 degrees Fahrenheit is an approximate average normal oral temperature, not a universal threshold.

\medterm{Dehiscence} Partial or complete separation of a healing wound or surgical closure. It may involve only the skin or deeper tissues; increasing pain, drainage, or exposed tissue requires prompt medical assessment.

\medterm{DEHP} Di(2-ethylhexyl) phthalate, a chemical added to some plastics to make them flexible. Repeated or high exposure may affect reproduction, development, or hormone function, so medical exposure is limited when possible.

\medterm{Dehydration} Loss of more water and electrolytes than are replaced, often from vomiting, diarrhea, fever, sweating, or poor fluid intake. It can cause thirst, dry mouth, reduced urination, dizziness, fast heartbeat, abnormal electrolyte levels, or shock.

\medterm{Dehydration In Children} Loss of more body water than a child replaces, often from diarrhea, vomiting, fever, sweating, or poor intake. Warning signs include dry mouth, few tears, reduced urination, unusual sleepiness, dizziness, or rapid breathing.


\medterm{Dehydroascorbic Acid} The oxidized form of vitamin C. Cells can take it up and convert it back to active vitamin C, allowing it to participate in antioxidant and other chemical reactions.

\medterm{Dehydroepiandrosterone} A hormone made mainly by the adrenal glands that the body can convert into other sex hormones. Levels change with age and illness; supplements may cause acne, unwanted hair, mood changes, or medicine interactions.

\medterm{Deinstitutionalize} To help a person leave long-term institutional care and live in the community with suitable housing, healthcare, and support. The plan should protect safety, treatment needs, and independence.

\medterm{Deiodinase} An enzyme that removes iodine from thyroid hormones and changes how active they are. Different tissues use these enzymes to help control metabolism, growth, brain development, and responses to thyroid hormone.

\medterm{Deja Vu} A brief feeling that a new situation has happened before. Occasional episodes are common, but frequent episodes with loss of awareness, unusual movements, or confusion may occur with seizures and need assessment.

\medterm{Dejerine-Roussy Syndrome} Persistent burning, aching, or painful abnormal sensation on one side of the body after a stroke damages a deep brain sensory area. Normally harmless touch or temperature may become painful.

\medterm{Delavirdine} An older HIV medicine that blocks reverse transcriptase, an enzyme HIV needs to copy its genetic material. It was used with other HIV medicines, but resistance and drug interactions limit its use.

\medterm{Delayed Cord Clamping} Waiting, commonly 30–60 seconds or longer, before clamping the umbilical cord after birth. This allows extra blood to reach the newborn and may improve early circulation and iron stores; timing depends on both patients’ condition.

\medterm{Delayed Ejaculation} Persistent difficulty or inability to ejaculate despite adequate sexual stimulation, causing distress. Medicines, nerve or hormone problems, chronic illness, psychological factors, and relationship circumstances may contribute.

\medterm{Delayed Emergence} Failure to wake and respond as expected after anesthesia. Causes include lingering anesthetic medicines, low body temperature, abnormal blood sugar or salts, breathing problems, stroke, or other complications and require prompt assessment.

\medterm{Delayed Gastric Emptying} Alternate index label; see \textit{Gastroparesis}.

\medterm{Delayed Gastric Emptying (Gastroparesis)} Alternate index label for Gastroparesis.

\medterm{Delayed Growth} Slower-than-expected increase in a child’s height, weight, or head size. Nutrition, chronic disease, hormone problems, genetic conditions, and family growth pattern can contribute; growth should be assessed over time.

\medterm{Delayed Hypersensitivity} An immune reaction that develops hours or days after exposure rather than immediately. It can cause contact dermatitis, some medicine rashes, or a reaction to a tuberculosis skin test.

\medterm{Delayed Orgasm} Persistent difficulty or inability to reach orgasm despite adequate stimulation, causing distress. Medicines, illness, hormonal or nerve problems, anxiety, trauma, and relationship factors may contribute.

\medterm{Delayed Puberty} Puberty that has not begun or is not progressing at the expected age. Family pattern, poor nutrition, chronic illness, hormone disorders, or problems involving the ovaries, testes, brain, or pituitary may contribute.

\medterm{Delayed Puberty In Boys} Lack of testicular enlargement or other expected pubertal changes by the usual age. Causes include normal late family development, poor nutrition, chronic illness, low testosterone, or hypothalamus or pituitary disorders.

\medterm{Delayed Puberty In Girls} Lack of breast development or menstruation by the usual age, or puberty that stops progressing. Causes include normal late family development, low energy availability, chronic illness, ovarian failure, or hypothalamus or pituitary disorders.

\medterm{Delayed Sleep Phase} A body-clock pattern in which a person naturally falls asleep and wakes much later than desired. Sleep may be normal at the preferred times but interfere with school, work, or social activities.

\synonyms
delayed sleep-wake phase disorder.
Delayed Sleep-Wake Phase Sleep Disorder
DSPS
DSWPD

\medterm{Delayed Sleep Phase Syndrome} A persistent body-clock disorder in which sleep begins and ends much later than desired, despite the ability to sleep normally on that schedule. It causes difficulty meeting work, school, or social obligations.

\medterm{Delayed Sleep-Wake Phase Sleep Disorder} Alternate index label; see \textit{Delayed sleep phase}.

\medterm{Deletion} Loss of a section of DNA or part of a chromosome. The effects depend on the size, location, and genes involved; deletions may cause no problems or lead to inherited, developmental, or medical disorders.

\medterm{Deletion Mutation} A genetic change in which one or more DNA building blocks are missing. Its effect depends on the location and size of the deletion and whether it changes a gene or its control region.

\medterm{Delftia} A group of oxygen-using bacteria found mainly in soil and water. Human infection is uncommon but can occur as an opportunistic infection, especially in people with weakened immunity or medical devices.

\medterm{Delirium} A sudden change in attention, awareness, and thinking that fluctuates over hours or days. It usually results from illness, infection, medicines, substance use or withdrawal, dehydration, or another medical problem.

\medterm{Delirium Superimposed On Dementia} Delirium occurring in a person who already has dementia. The sudden, fluctuating worsening in attention and awareness differs from the person’s usual cognitive impairment and may be triggered by illness, medicines, dehydration, or hospitalization.

\medterm{Delirium In Anesthesia} Sudden confusion, poor attention, and changing awareness during recovery from anesthesia or surgery. Possible causes include medicines, infection, pain, low oxygen, dehydration, or abnormal body chemistry and require prompt assessment.

\medterm{Delirium Tremens} A life-threatening form of alcohol withdrawal causing severe confusion, agitation, hallucinations, shaking, sweating, fever, fast heartbeat, and high blood pressure. It commonly begins two to four days after heavy drinking stops and requires emergency treatment.

\medterm{Delivery} Childbirth, ending with birth of the baby and delivery of the placenta. It may occur through the vagina, sometimes with forceps or vacuum assistance, or through cesarean surgery.

\medterm{Delivery Of Posterior Arm} A maneuver for shoulder dystocia in which the clinician brings the baby’s back arm through the birth canal. This reduces the width of the shoulders and may help complete delivery.

\medterm{Delivery Presentations} The fetal body part positioned to enter the birth canal first. Common presentations are head-first, breech, and shoulder; the presentation affects delivery planning and risks.

\medterm{Delta Agent} An older name for hepatitis D virus, which can reproduce only when hepatitis B virus is present. Infection may cause severe acute hepatitis or worsen long-term liver disease.

\medterm{Delta Waves} Very slow electrical brain waves that are normal during deep non-REM sleep. Excessive delta activity while awake may indicate brain dysfunction and must be interpreted with symptoms and other tests.

\synonyms
Delta rhythm

\medterm{Delta-ALA Urine Test} A urine test measuring a substance used to make hemoglobin. Increased levels may occur with acute porphyria or lead poisoning and are interpreted with symptoms and other tests.

\medterm{Deltacoronavirus} A group of coronaviruses identified by their genetic features. Some infect birds or mammals, and their ability to infect people or cause disease depends on the specific virus and host.

\medterm{Deltaretrovirus} A group of viruses that includes human T-cell leukemia viruses. They copy their RNA into DNA inside infected cells; some can persist for years and cause blood-cell disorders or cancer.

\medterm{Delta-Tocopherol} One form of vitamin E found in foods and supplements. It has antioxidant activity, but alpha-tocopherol is the main form maintained in human blood and used to assess vitamin E status.

\medterm{Deltoid} The large, triangular muscle covering the shoulder. It helps lift the arm and stabilize the shoulder joint; injury or nerve damage can cause pain and weakness.

\medterm{Deltoid Muscle} A triangular shoulder muscle attached to the collarbone, shoulder blade, and upper arm. Its front, middle, and back fibers move the arm forward, outward, and backward; the axillary nerve supplies it.

\medterm{Delusion} A firmly held false belief that remains despite clear evidence to the contrary and is not explained by the person’s cultural background. It may involve persecution, grandiosity, illness, or other themes.

\medterm{Delusional Disorder} A mental disorder with one or more persistent delusions lasting at least one month, while overall functioning is relatively preserved. Hallucinations, if present, relate closely to the delusional belief.

\medterm{Demarcation} A clear boundary between two areas or conditions, such as the line separating healthy tissue from tissue that has died. It may guide diagnosis or treatment.

\medterm{Demeclocycline} A tetracycline antibiotic that blocks bacterial protein production. It was also used to reduce excess water retention in some patients with inappropriate antidiuretic hormone secretion; adverse effects include sun sensitivity and kidney or liver problems.

\medterm{Demecolcine} A colchicine-related substance that disrupts microtubules, structures needed for cell division. It has limited medical use and can be highly toxic, causing gastrointestinal, blood-cell, nerve, or organ injury.

\medterm{Dementia} A lasting decline in memory, thinking, language, judgment, or other mental abilities that interferes with independent daily life. Causes include brain diseases, blood-vessel disease, infections, medicines, and some treatable medical problems.

\synonyms
Major neurocognitive disorder

\medterm{Dementia And Driving} The effect of cognitive impairment on driving safety. Assessment considers memory, judgment, visuospatial skills, reaction time, medications, and local licensing requirements.

\medterm{Dementia - Behavior And Sleep Problems} Behavioral and sleep changes that can occur in dementia, including agitation, wandering, apathy, hallucinations, or altered sleep–wake timing; causes and reversible contributors should be assessed.


\medterm{Dementia Due To Metabolic Causes} Cognitive impairment caused or worsened by metabolic problems such as thyroid disease, vitamin deficiency, organ failure, electrolyte disturbance, or medication toxicity; some causes are treatable.

\medterm{Dementia, Frontotemporal} Alternate index label; see \textit{Frontotemporal dementia}.

\medterm{Dementia, Lewy Body} Alternate index label; see \textit{Lewy body dementia}.

\medterm{Dementia (Major Neurocognitive Disorder)} Alternate name; see \textit{Dementia}.

\medterm{Dementia Pugilistica} A progressive brain disorder associated with repeated head trauma, historically described in some professional boxers.
\medterm{Dementia, Vascular} Alternate index label; see \textit{Vascular dementia}.

\medterm{Demethylation} Removal of a small chemical group called a methyl group from a molecule such as DNA or a protein. In DNA, this change can alter which genes are active and affect cell development or disease.

\medterm{Demineralization} Loss of minerals from a tissue such as tooth enamel or bone. It can weaken the tissue, causing early tooth decay or contributing to fragile bones, depending on the cause.


\medterm{Demulcent} A substance that coats and soothes irritated moist body linings, temporarily reducing discomfort from dryness, inflammation, or minor injury.

\medterm{Demyelinating Disease} A disease in which the protective covering around nerve fibers is damaged. Nerve signals then slow or fail, causing problems such as weakness, numbness, vision changes, or poor coordination.

\medterm{Demyelinating Disease Of The Central Nervous System} A disorder in which the protective covering of nerves in the brain, spinal cord, or vision pathways is damaged. Causes include immune disease, infection, inherited conditions, toxins, and metabolic problems.

\medterm{Demyelination} Loss or damage of myelin, the protective coating that helps nerve fibers send signals quickly. It can disrupt movement, sensation, vision, balance, or thinking and may result from immune, infectious, inherited, or toxic disorders.

\medterm{Denaturant} A substance added to another product to change its taste, smell, or properties, often to make it unsuitable for drinking. Some denaturants are poisonous or irritating, so health effects depend on the product and amount swallowed or contacted.

\medterm{Dendrite} A branching extension of a nerve cell that receives signals from other nerve cells or sensory receptors and carries those signals toward the cell body.

\medterm{Dendrites} Branching extensions of nerve cells that receive signals from other neurons or sensory receptors and carry those messages toward the cell body.

\medterm{Dendritic Cell} An immune cell that captures foreign material and shows pieces of it to other immune cells. This helps start a targeted immune response against infections and abnormal cells.

\medterm{Dendritic Cells} Immune cells that collect material from germs or damaged cells and show identifying pieces to T cells. This helps the immune system start a targeted response and develop longer-lasting protection.

\medterm{Dendritic Spine} A tiny branch on a nerve cell that receives signals from other nerve cells. Its shape and number can change with learning, development, and some brain diseases.

\medterm{Dendritic Ulcer} A branching epithelial ulcer of the cornea, classically caused by herpes simplex virus. It may cause a red eye, pain, light sensitivity, tearing, or reduced vision and requires ophthalmic assessment.

\medterm{Denervate} To remove or interrupt the nerve supply to a tissue or organ, either deliberately during treatment or because of injury or disease. The change may reduce movement, sensation, pain, or automatic organ function.

\medterm{Denervation} Denervation is loss of nerve supply to a muscle or tissue, causing weakness, atrophy, altered sensation, or autonomic dysfunction.

\medterm{Dengue} A viral infection caused by a dengue virus and spread mainly by the bites of infected female \textit{Aedes} mosquitoes. Many infections cause no symptoms; symptomatic illness commonly causes sudden fever, severe headache, pain behind the eyes, muscle or joint pain, nausea, and rash. A small number of cases progress to severe dengue, which can cause leakage of fluid from blood vessels, fluid buildup with breathing difficulty, severe bleeding, serious organ impairment, shock, or death.

\medterm{Dengue Fever} Alternate name; see \textit{Dengue}.

\synonyms
Dengue Hemorrhagic Fever


\medterm{Dengue Hemorrhagic Fever} A severe form of dengue infection in which blood vessels leak and bleeding or shock may develop. Warning signs include severe abdominal pain, repeated vomiting, bleeding, extreme tiredness, cold skin, or reduced urine; urgent hospital care and careful fluids are needed.
\medterm{Dengue Shock Syndrome} Life-threatening circulatory shock caused by severe dengue with plasma leakage and reduced effective blood volume, often accompanied by bleeding or organ dysfunction.

\medterm{Dengue Virus} The virus that causes dengue fever and is spread mainly by infected Aedes mosquitoes. Four major types exist; protection after infection is strongest against the same type and does not reliably prevent later infection by other types.

\medterm{Denial} A psychological response in which a person cannot or will not accept a distressing fact, diagnosis, or situation. It may briefly reduce distress but can delay needed decisions or treatment when persistent.

\medterm{Denileukin Diftitox} A laboratory-made protein combining part of interleukin-2 with diphtheria toxin. It was used for selected cutaneous T-cell lymphomas by delivering the toxin to cells bearing the interleukin-2 receptor; toxicity limits use.

\medterm{Denominator} The number below the line in a fraction. In a medical rate or proportion, it represents the total population or group being compared with the number above it.

\medterm{Denosumab} A medicine that blocks a protein needed to activate bone-destroying cells. It is used for osteoporosis and some cancer-related bone problems; stopping it without follow-up can cause rapid bone loss and spine fractures.

\medterm{Denosumab For Osteoporosis} An injection given every six months that blocks a signal activating bone-destroying cells. It lowers fracture risk, including at the spine and hip, but calcium must be monitored and another treatment is needed when it is stopped.

\medterm{Dens} The tooth-shaped projection of the second neck vertebra. It fits into the first neck vertebra and acts as a pivot that allows the head to turn.

\medterm{Dens Evaginatus} An extra bump or cusp on a tooth, most often a premolar. It may contain an extension of the tooth’s nerve and can wear, break, or develop decay and pain.

\medterm{Dens Invaginatus} An inward folding of a tooth’s enamel and dentin that may reach the nerve. The pocket can trap bacteria and food, increasing the risk of early decay, nerve infection, and tooth loss.

\medterm{Densitometry} Measurement of how dense a tissue is, most often the mineral content of bone. A bone-density scan helps diagnose osteoporosis, estimate fracture risk, and monitor treatment.

\medterm{Density} The amount of mass in a given volume. In medicine, density measurements help describe bone mineral content, tissue composition, fluid collections, or findings on imaging.

\medterm{Dent Disease} An inherited disorder of the kidney’s filtering tubes, usually affecting boys. It causes loss of small proteins in urine, excess urinary calcium, kidney stones or calcium deposits, and sometimes progressive kidney failure.

\medterm{Dental Abscess} A pocket of pus caused by infection in a tooth or surrounding gum. It may cause severe toothache, swelling, fever, or difficulty opening the mouth; spreading facial or neck swelling needs urgent care.

\medterm{Dental Age} An estimate of a child’s developmental age based on tooth formation, hardening, and eruption. It can help assess growth but does not exactly determine the child’s actual age.

\medterm{Dental Alloy} A mixture of metals used to make dental fillings, crowns, bridges, dentures, orthodontic appliances, or other restorations. Its strength, appearance, corrosion resistance, and possible allergies depend on its composition.

\medterm{Dental Alveolus} The part of the upper or lower jawbone that forms a tooth socket and holds the tooth root. It can shrink after tooth loss or severe gum disease.

\medterm{Dental Amalgam} A durable filling material made by mixing liquid mercury with powdered metals, mainly silver, tin, and copper. It is strong but has a silver color and may be unsuitable in some situations.

\medterm{Dental Anesthesia} Use of local anesthetic, sedation, or occasionally general anesthesia to prevent pain or reduce anxiety during dental treatment. The method depends on the procedure, health, age, and patient preference.

\medterm{Dental Anxiety} Fear or distress about dental treatment that may cause avoidance, delayed care, or difficulty during appointments. Reassurance, communication, pain control, and behavioral or medicine-based support may help.

\medterm{Dental Avulsion} Complete displacement of a tooth from its socket, usually caused by an injury. A permanent tooth that is knocked out needs immediate dental care; hold it by the crown, keep it moist, and do not scrub the root.

\medterm{Dental Arch} The curved arrangement of teeth and the bone supporting them in the upper or lower jaw. Its shape affects tooth alignment, bite, speech, chewing, and space for dental restorations.

\medterm{Dental Articulator} A device that holds models of the upper and lower teeth and imitates jaw movements. It helps dental professionals design and adjust crowns, bridges, dentures, and other restorations so they fit the bite.

\medterm{Dental Braces} Fixed or removable appliances that apply gentle, controlled force to move teeth and improve their alignment or bite. Treatment requires cleaning and follow-up; possible problems include soreness, decay, gum disease, and root shortening.

\synonyms
Braces And Retainers (Dental Braces)

\medterm{Dental Braces (Orthodontics)} Fixed or removable appliances that gradually move teeth and may improve the bite or tooth alignment. Treatment length and the need for extraction or other appliances depend on the teeth, jaws, growth, and treatment goals.

\medterm{Dental Bridge} A fixed dental appliance that replaces one or more missing teeth by attaching artificial teeth to nearby teeth or implants. It can improve chewing and appearance but requires healthy supports and careful cleaning.

\medterm{Dental Calculus} Hardened plaque formed when minerals in saliva or gum fluid collect on teeth. It cannot be removed reliably by brushing alone and can contribute to gum inflammation and disease.

\medterm{Dental Caries} Tooth decay caused when bacteria use sugars and produce acids that remove minerals from enamel and dentin. Untreated decay can cause sensitivity, pain, infection, and tooth loss.
\medterm{Dental Caries (Cavities)} Alternate index label for Cavities.

\medterm{Dental Cavities} Permanent holes or weakened areas caused when plaque bacteria use sugars and make acids that dissolve a tooth. Decay can spread from enamel into the deeper tooth, causing sensitivity, pain, infection, or tooth loss.

\medterm{Dental Cement} A dental material used to attach crowns, bridges, or orthodontic bands, fill small spaces, or protect exposed tooth tissue. Its strength and use depend on the material and treatment.

\medterm{Dental Ceramic} A tooth-colored material used for crowns, veneers, bridges, fillings, and implant restorations. It provides good appearance and resistance to wear, but its strength and suitability vary by type and location.

\medterm{Dental Clasp} A part of a removable partial denture that rests against or around supporting teeth to help keep the denture in place. Poor fit or excessive pressure can irritate teeth or gums.

\medterm{Dental Crown} A custom-made covering placed over a damaged or weakened tooth to restore its shape, strength, chewing function, or appearance. It may be made from ceramic, metal, resin, or a combination of materials.

\medterm{Dental Crowns} Custom-made coverings placed over most or all of a damaged or weakened tooth. A crown restores shape, strength, chewing function, and appearance after decay, fracture, or root-canal treatment.

\medterm{Dental Deposit} Material that collects on a tooth, such as plaque, hardened tartar, food debris, or stain. Some deposits irritate the gums or support decay and may require professional cleaning.

\medterm{Dental Disease} A disorder of the teeth, gums, jaw, or supporting tissues, including decay, gum disease, infection, and tooth loss. Fluoride, daily cleaning, limited frequent sugar exposure, and dental visits reduce risk.

\medterm{Dental Dysplasia} Abnormal development of a tooth or its supporting tissue. It may change the tooth’s shape, color, strength, position, or eruption and can increase the risk of wear, decay, or infection.

\medterm{Dental Erosion} Gradual loss of tooth enamel caused by acid rather than bacteria. Acidic foods or drinks, reflux, vomiting, and some medicines may contribute, causing sensitivity or worn teeth; reducing acid exposure helps prevent progression.

\medterm{Dental Extraction} Removal of a tooth from its jaw socket because of severe decay, infection, gum disease, impaction, crowding, trauma, or treatment planning. Risks include bleeding, infection, nerve injury, fracture, and dry socket.
\medterm{Dental Drill} A powered dental handpiece with a rotating bur, driven by air or an electric motor, used to remove decayed tooth structure, shape teeth, prepare teeth for fillings or crowns, and perform other dental procedures. Different burs and speeds are selected for the material and procedure.

\synonyms
Dental Handpiece

\medterm{Dental Fillings} Materials placed in a tooth after decay or damage is removed to restore its shape and function. Resin, glass ionomer, metal amalgam, or other materials may be used depending on the tooth and clinical situation.

\medterm{Dental Fissure} A narrow natural groove or developmental defect in a tooth, especially on a chewing surface. Food and bacteria can collect there, increasing the risk of decay.

\medterm{Dental Fluorosis} Changes in tooth enamel caused by too much fluoride while the teeth are developing. It may appear as white lines or patches and, in more severe cases, brown discoloration or pits.

\medterm{Dental Floss} A thin thread used to remove plaque and food between teeth and just below the gumline. Daily cleaning between teeth helps prevent gum inflammation but does not replace brushing or dental care.

\medterm{Dental Hygiene} Daily care of the teeth, gums, tongue, and mouth, including brushing with fluoride toothpaste and cleaning between teeth. It removes plaque and helps prevent decay, gum disease, and bad breath.

\medterm{Dental Implant} A small biocompatible fixture placed in the jawbone to support a replacement tooth, bridge, or denture. Bone may grow around it for stability; risks include infection, poor healing, nerve injury, and implant failure.

\medterm{Dental Impression Material} A material placed in a tray to make an accurate mold of the teeth and mouth tissues. The mold helps make crowns, dentures, orthodontic appliances, restorations, or study models.

\medterm{Dental Injuries} Damage to teeth, gums, jawbone, or mouth tissues caused by a fall, blow, sports injury, accident, or biting force. Prompt assessment may save displaced or knocked-out teeth; severe bleeding or head injury is an emergency.

\synonyms
Dental Trauma (Dental Injuries)

\medterm{Dental Inlay} A custom-made filling fitted inside the grooves of a damaged tooth and bonded in place. It restores areas too large for a simple filling but does not cover the tooth’s cusps like an onlay or crown.

\medterm{Dental Material} A substance used in dental care, such as composite, amalgam, ceramic, metal, cement, bonding material, impression material, or implant material. Choice depends on the treatment, location, strength, appearance, and patient factors.

\medterm{Dental Occlusion} The way the upper and lower teeth meet when the mouth closes or the jaw moves. An abnormal bite may contribute to difficulty chewing, tooth wear, jaw pain, or muscle strain.

\medterm{Dental Overjet} The horizontal distance between the upper and lower front teeth when the teeth are closed. A large overjet may increase the risk of trauma to the upper front teeth and affect biting.

\medterm{Dental Pain (Toothache)} Pain from a tooth or nearby gum and jaw tissues. Common causes include decay, inflammation or infection inside the tooth, a crack, trauma, gum disease, or an impacted tooth.

\medterm{Dental Papilla} Soft embryonic tissue beneath a developing tooth. Its cells form the dentin and living pulp, which contains the tooth’s nerves and blood vessels.

\medterm{Dental Pellicle} A thin protective film of salivary proteins that forms on tooth enamel soon after cleaning. Bacteria can attach to it and build dental plaque.

\medterm{Dental Plaque} A sticky film of bacteria, saliva, and food material on teeth and along the gums. If not removed by brushing and cleaning between teeth, it can cause decay, gum inflammation, and gum disease.

\medterm{Dental Porcelain} A tooth-colored ceramic used for crowns, veneers, bridges, and other restorations. It provides a natural appearance but may chip or fracture under heavy force.

\medterm{Dental Procedure} An examination, preventive service, or treatment involving the teeth, gums, jaw, or mouth, such as cleaning, filling, extraction, root-canal treatment, or restoration. Benefits and risks depend on the procedure.

\medterm{Dental Prosthesis} An artificial replacement for one or more missing teeth or nearby mouth tissues, such as a denture, bridge, crown, or implant-supported restoration. It restores appearance, speech, and chewing.

\medterm{Dental Public Health} Efforts to prevent and control oral disease in communities. They include fluoride programs, oral-health education, screening, disease monitoring, access to dental care, and policies that reduce differences in oral health.

\medterm{Dental Pulp} The living soft tissue at the center of a tooth. It contains nerves, blood vessels, and connective tissue and can become inflamed or infected when decay or injury reaches it.

\medterm{Dental Pulp Necrosis} Death of the soft tissue inside a tooth, usually after untreated inflammation, deep decay, trauma, or loss of blood supply. The tooth may become discolored and develop infection around its root.

\medterm{Dental Pulpitis} Inflammation of the soft tissue inside a tooth, commonly caused by decay, trauma, a crack, or dental treatment. Brief pain may be reversible; spontaneous or lingering pain may require root-canal treatment or extraction.

\medterm{Dental Radiographs} X-ray images of the teeth and jaws used to find decay, bone loss, infection, impacted teeth, fractures, and developmental problems. Types include bitewing, periapical, occlusal, and panoramic views.

\medterm{Dental Radiography} Use of X-rays to examine teeth and jaws. Bitewing, periapical, panoramic, and three-dimensional scans answer different questions; exposure is kept as low as reasonably achievable while obtaining useful images.

\medterm{Dental Restorative Materials} Materials used to repair teeth, including composite resin, glass ionomer, amalgam, ceramic, and metal alloys. Selection depends on tooth location, size of damage, strength, appearance, moisture control, cost, and patient needs.

\medterm{Dental Sac} A layer of tissue surrounding a developing tooth. It forms the tissues that support the tooth, including the periodontal ligament, cementum, and part of the surrounding jawbone.

\medterm{Dental Screening} A check of the teeth, gums, and mouth to look for decay, gum disease, oral lesions, or risk factors before serious symptoms develop. Screening does not replace a complete dental examination.

\medterm{Dental Sealant} A thin protective coating placed in the grooves of teeth, especially molars. It blocks food and bacteria from collecting in deep pits and lowers the risk of decay.

\medterm{Dental Sealants} Thin resin or glass-ionomer coatings placed in the pits and grooves of teeth, especially molars. They reduce plaque and food retention and lower the risk of decay but do not replace brushing or fluoride.

\medterm{Dental Trauma (Dental Injuries)} Alternate index label for Dental Injuries.

\medterm{Dental Traumatology} The prevention, assessment, and treatment of injuries to teeth and surrounding mouth tissues. Injuries may involve cracks, fractures, displacement, root damage, or complete tooth loss and often need prompt care.

\medterm{Dental X-Rays} X-ray images of teeth and jaws used to detect decay, bone loss, infection, impacted teeth, fractures, developmental problems, and other disease not visible during an examination.

\medterm{Dentate Gyrus} A curved part of the hippocampus that helps process incoming information and form memories. It also produces some new nerve cells during adulthood.

\medterm{Dentatorubral-Pallidoluysian Atrophy} A rare inherited brain disorder caused by an extra-long repeated section of the ATN1 gene. Symptoms vary with age and may include seizures, poor coordination, involuntary movements, behavior changes, and worsening thinking ability.

\synonyms
DRPLA; Haw River Syndrome; Naito-Oyanagi Disease

\medterm{Dentifrice} A preparation, such as toothpaste or tooth powder, used to clean teeth. Fluoride dentifrice strengthens enamel and helps prevent tooth decay when used regularly.

\medterm{Dentigerous Cyst} A fluid-filled sac that forms around the crown of an unerupted tooth, usually in the jaw. It may cause swelling, move nearby teeth, or damage bone if it becomes large or infected.

\medterm{Dentin} The hard tissue beneath tooth enamel that forms most of the tooth. Tiny channels within it can transmit sensitivity, especially when dentin is exposed by decay, wear, gum recession, or injury.

\medterm{Dentin Sensitivity} Brief, sharp tooth pain caused by cold, heat, sweet foods, touch, or air reaching exposed dentin. Decay, worn enamel, gum recession, cracks, or aggressive brushing can expose it.

\medterm{Dentin Hypersensitivity} Alternate term; see \textit{Dentin Sensitivity}.

\medterm{Dentinogenesis} Formation of dentin by specialized cells during tooth development and repair. Abnormal dentin formation may produce weak, discolored, sensitive teeth that wear or fracture easily.

\medterm{Dentinogenesis Imperfecta} An inherited disorder in which dentin forms abnormally, causing teeth that may look blue-gray or amber and break, wear, or discolor easily. Some forms occur with fragile bones.

\medterm{Dentist} A licensed healthcare professional who examines, prevents, diagnoses, and treats diseases and injuries of the teeth, gums, jaw, and mouth.

\medterm{Dentistry} The healthcare field concerned with preventing, diagnosing, and treating disorders and injuries of the teeth, gums, jaw, and mouth, including restoration, replacement, and oral health education.

\medterm{Dentistry (Geriatric)} Dental care for older adults, who may have dry mouth, root decay, gum disease, tooth loss, denture problems, or oral cancer risk. Care includes medication review, prevention, examination, treatment, and adaptations for mobility or cognition.

\synonyms
Geriatric dentistry

\medterm{Dentition} The teeth present in the mouth, including their number, arrangement, development, and eruption. Children usually have primary teeth that are replaced by permanent teeth; disease, injury, or development can alter either set.

\medterm{Denture} A removable replacement for some or all missing teeth and nearby mouth tissues. Partial dentures replace selected teeth; complete dentures replace all teeth in one jaw. Correct fit, cleaning, and review help prevent sores and infection.

\medterm{Denture Base} The part of a denture that rests on the gums and mouth tissues and supports the artificial teeth. Its shape and fit distribute chewing pressure and help prevent soreness.

\medterm{Denture Problems} Difficulties caused by a removable denture, including poor fit, looseness, pain, sores, breakage, difficulty chewing or speaking, or infection. A dental review may correct the cause or determine whether replacement is needed.

\medterm{Dentures} Removable appliances that replace some or all missing teeth and nearby gum tissue. They can improve chewing, speech, and appearance but need proper fitting, daily cleaning, gradual adjustment, and regular dental review.

\medterm{Denturist} A dental professional authorized in some regions to make, fit, adjust, and repair removable dentures within the local scope of practice. Dentists may also assess oral disease and provide related treatment.

\medterm{Denver Developmental Screening Test} A screening questionnaire and task-based assessment for young children that checks personal-social, fine-motor, language, and gross-motor skills. It identifies possible developmental delay but does not diagnose its cause.

\medterm{Denys-Drash Syndrome} A rare genetic disorder caused by changes in the \textit{WT1} gene. It can cause severe kidney disease, differences in sex development, and a high risk of Wilms kidney tumor in childhood.

\medterm{Deodorant Poisoning} Harmful exposure to a deodorant product, usually by swallowing, inhaling, or getting it in the eyes or on the skin. Symptoms depend on the ingredients and may include irritation, vomiting, breathing difficulty, or chemical burns.

\medterm{Deoxycoformycin} An anticancer medicine that blocks an enzyme needed by some cells to make DNA. It was used for selected blood cancers and can suppress bone marrow, weaken immunity, and increase the risk of serious infections.

\medterm{Deoxycytidine} A DNA building block made of the base cytosine joined to the sugar deoxyribose. Cells use it to make and repair DNA.
\medterm{Deoxycytidine Kinase} An enzyme that chemically activates deoxycytidine and some similar medicines. The activated medicines can interfere with the DNA of cancer cells or viruses and stop them from multiplying.

\medterm{Deoxyglucose} A molecule resembling glucose that is used mainly in research and PET imaging. When labeled with a radioactive form of fluorine, it can show tissues using unusually large amounts of glucose, including many tumors and areas of inflammation.

\medterm{Deoxyguanosine} A DNA building block made of the base guanine joined to the sugar deoxyribose. Cells use it to make and repair DNA.
\medterm{Deoxyhemoglobin} Hemoglobin that is not carrying oxygen. A high amount in blood near the skin can make the lips, skin, or nail beds appear blue, a finding called cyanosis.

\medterm{Deoxyribonuclease} An enzyme that breaks DNA into smaller pieces. An inhaled laboratory-made form can thin thick mucus in some people with cystic fibrosis by breaking down DNA released from damaged cells.

\medterm{Deoxyribonucleic Acid} DNA, the molecule that stores hereditary information in cells. Its sequence helps instruct protein production and cell control; changes may be inherited or acquired during life and can affect health.

\synonyms
DNA

\medterm{Deoxyribonucleosides} Molecules made of one of the four DNA bases joined to the sugar deoxyribose. Cells use them as starting materials when making or repairing DNA.

\medterm{Deoxyribonucleotides} DNA building blocks made of a base, deoxyribose sugar, and phosphate groups. The four forms—dATP, dGTP, dCTP, and dTTP—are assembled into DNA strands.

\medterm{Dependence} Physical adaptation to repeated exposure to a medicine or substance, so that reducing or stopping it can cause withdrawal symptoms. Dependence can occur with or without addiction.

\medterm{Dependence, Nicotine} Alternate index label; see \textit{Nicotine dependence}.

\medterm{Dependent Personality Disorder} A long-term pattern of excessive need to be cared for, difficulty making decisions independently, fear of separation, and submissive behavior that causes distress or impairs relationships and daily life.

\medterm{Dependovirus} A group of small viruses that usually need another virus, such as an adenovirus or herpesvirus, to reproduce. They are generally not known to cause serious disease in healthy people.

\medterm{Depersonalization} A feeling of being detached from or observing one’s own body, thoughts, or emotions as if from outside. It can occur with anxiety, trauma, panic, substance use, or other mental health conditions.

\medterm{Depersonalization-Derealization Disorder} A disorder causing repeated or persistent feelings that one’s body, thoughts, emotions, or surroundings are unreal or distant, while the person knows these feelings are not literally true. Symptoms cause significant distress or impairment.

\medterm{Dephosphorylation} Removal of a phosphate group from a molecule. This can switch proteins and other cell components on or off and change how cells communicate, use energy, grow, or respond to signals.

\medterm{Depigmentation} Loss or reduction of normal pigment, especially melanin in the skin, hair, or mucous membranes. It may result from vitiligo, inflammation, injury, infection, chemical exposure, or treatment.

\medterm{Depilation} Removal of hair at or near the skin surface by shaving, cutting, waxing, or a chemical product. It differs from epilation, which removes hair together with its root.

\medterm{Depilatory Poisoning} Harmful exposure to a hair-removal product, usually by swallowing or contact with the eyes, skin, or airways. Ingredients may cause irritation, chemical burns, vomiting, or breathing problems.

\medterm{Depolarization} A change that makes the inside of a cell less electrically negative. In nerve and muscle cells, it can start or transmit an electrical signal that produces movement or sensation.

\medterm{Depot Implant} A small device placed in the body that releases medicine slowly over an extended period. It can provide steady treatment but may require a procedure for insertion and removal.

\medterm{Deprescribing} A planned, supervised reduction or stopping of a medicine when its harms, burden, or limited benefit outweigh expected benefit. The plan considers withdrawal, return of illness, safer alternatives, and patient goals.

\medterm{Depressed Fracture} A break in which bone is pushed inward, most often in the skull. It can injure the brain or other underlying tissue and requires urgent assessment, especially after head injury or with neurologic symptoms.

\medterm{Depression} A group of mental health disorders involving a persistent sad, empty, or irritable mood, or a loss of interest or pleasure in activities. It is more than the short-lived sadness or low mood that most people experience at times. Other symptoms may include changes in sleep, appetite, energy, concentration, movement, self-worth, or hope, as well as physical aches or digestive problems. Symptoms cause significant distress or interfere with daily life and may include thoughts of death or suicide. Major depression involves symptoms most of the day, nearly every day, for at least two weeks; persistent depressive disorder involves less severe symptoms lasting at least two years.

\synonyms
Clinical Depression (Depression)

\medterm{Depression (Geriatric)} Depression in an older adult may cause low mood, loss of interest, fatigue, sleep or appetite changes, slowed thinking, physical complaints, or memory problems. Medical illness, medicines, isolation, and bereavement may contribute.

\synonyms
Depression in older adults

\medterm{Depression In Children} Depression in a child may appear as persistent sadness or irritability, loss of interest, withdrawal, school decline, sleep or appetite changes, physical complaints, guilt, or thoughts of death. Evaluation includes development, family context, medical causes, and safety.

\medterm{Depression In Older Adults} Depression later in life may cause sadness, loss of interest, fatigue, sleep or appetite changes, slowed thinking, physical complaints, or memory problems. Assessment should also consider medical illness, medicines, grief, isolation, and suicide risk.

\medterm{Depression (Major Depressive Disorder)} Alternate index label; see \textit{Major depressive disorder}.



\medterm{Depression, Postpartum} Alternate index label; see \textit{Postpartum depression}.

\medterm{Depression Screening} Use of questions or a validated questionnaire to identify possible depression. A positive screen is not a diagnosis; a clinician must assess symptoms, duration, impairment, medical causes, treatment needs, and safety.

\medterm{Depression, Teen} Alternate index label; see \textit{Teen depression}.

\medterm{Depressive Disorder} A mood disorder involving persistent low mood, loss of interest, or both, with possible changes in sleep, appetite, energy, concentration, movement, or self-worth. Diagnosis depends on the pattern, duration, severity, and functional impact of symptoms.

\medterm{Depressive Illness} Broad, non-diagnostic term; see \textit{Depressive Disorder}.

\medterm{Depth Perception} The ability to judge how far away objects are and where they are in three-dimensional space. It depends mainly on using both eyes together and also on movement, size, shading, and perspective cues.

\medterm{Depurination} Loss of the DNA bases adenine or guanine, leaving a gap in the DNA strand. If the gap is not repaired correctly, it can cause a genetic mutation.

\medterm{DeQuervains Tenosynovitis (De Quervain's Tenosynovitis)} Alternate index label for De Quervain's Tenosynovitis.

\medterm{Derbyshire Neck} An older term for an enlarged thyroid gland, or goiter, historically associated with areas where iodine deficiency was common. The cause and significance require assessment of the thyroid and its function.

\medterm{Dercum's Disease} A rare disorder characterized by multiple painful fatty growths, often with obesity, fatigue, and widespread tenderness. The cause is uncertain, and diagnosis is based on symptoms and examination after other causes are excluded.

\medterm{Derealization} A feeling that the outside world is unreal, distant, dreamlike, or visually changed, while the person knows it is not literally unreal. It may occur with anxiety, panic, trauma, dissociation, or substance use.

\medterm{Dermabrasion} A procedure that uses a rapidly rotating instrument to remove outer skin layers. It may improve selected scars, wrinkles, sun damage, or lesions but can cause pain, pigment changes, infection, or scarring.

\medterm{Dermal} Relating to the dermis, the skin layer beneath the epidermis. The dermis contains connective tissue, blood vessels, nerves, hair follicles, and sweat glands that support and nourish the skin.

\medterm{Dermal Fillers} Injectable materials placed under or within the skin to restore volume or reduce wrinkles or scars. Common effects include swelling and bruising; infection, lumps, skin damage, or rare blockage of a blood vessel can be serious.

\medterm{Dermal Melanocytosis} A usually harmless blue-gray or slate-colored skin patch caused by pigment-producing cells located deep in the skin. It is often present at birth and commonly affects the lower back or buttocks.

\medterm{Dermal Papilla} A small structure at the base of a hair follicle that supplies signals and nutrients to developing hair cells. Its activity helps control hair growth and the hair-growth cycle.

\medterm{Dermal Structure} The dermis has a thin superficial layer and a deeper supportive layer. It contains collagen, elastic fibers, blood and lymph vessels, nerves, hair follicles, sweat glands, and oil glands.

\medterm{Dermatan Sulfate} A long chain of sugar molecules found in skin, connective tissue, blood vessels, and heart valves. It helps organize tissues and can influence wound healing and blood clotting.
\medterm{Dermatitis} Inflammation of the skin causing some combination of itching, redness, swelling, scaling, dryness, or blisters. It may result from irritation, allergy, infection, medicines, genetics, or another disease.

\medterm{Dermatitis, Atopic} Alternate index label; see \textit{Atopic dermatitis (eczema)}.

\medterm{Dermatitis, Cercarial} Alternate index label; see \textit{Swimmer's itch}.

\medterm{Dermatitis, Contact} Alternate index label; see \textit{Contact dermatitis}.

\medterm{Dermatitis Herpetiformis} A long-lasting, intensely itchy rash with groups of small bumps or blisters, commonly on the elbows, knees, buttocks, or scalp. It is strongly linked to celiac disease and usually improves with a strict gluten-free diet.

\medterm{Dermatitis, Scratch} Alternate index label; see \textit{Neurodermatitis}.

\medterm{Dermatitis, Seborrheic} Alternate index label; see \textit{Seborrheic dermatitis}.

\medterm{Dermatofibroma} A common harmless, firm skin nodule, often on the legs. It may be brown or pink and often dimples when pinched. It usually needs no treatment, but a changing or bleeding lesion should be examined.

\medterm{Dermatofibrosarcoma} A rare, slow-growing cancer of connective tissue in the skin. It usually begins as a firm plaque or nodule and can grow into nearby tissue, but it rarely spreads to distant organs.

\medterm{Dermatofibrosarcoma Protuberans} A rare cancer arising from connective tissue in the skin. It often appears as a slowly enlarging firm plaque or nodule, commonly on the trunk; it can recur locally but rarely spreads. Complete surgical removal is the usual treatment.

\synonyms
DFSP

\medterm{Dermatofibrosarcoma Protuberans Of The Breast} A rare, usually slow-growing skin cancer involving the breast skin or nearby superficial tissue. It can extend into surrounding tissue and recur if incompletely removed; treatment aims for complete surgical removal.

\medterm{Dermatoglyphics} The study of ridge patterns on the fingers, palms, toes, and soles. Patterns help with identification and may differ in some genetic or developmental conditions, but they cannot diagnose a condition by themselves.

\medterm{Dermatographia} A form of hives in which light scratching, rubbing, or pressure causes raised, red, itchy lines within minutes. It is usually harmless, though symptoms may be troublesome.

\medterm{Dermatographia (Dermatographism)} See \textit{Skin Writing}.

\synonyms
Skin Writing

\medterm{Dermatographic Urticaria} A form of hives in which scratching, rubbing, or pressure causes raised, itchy lines on the skin within minutes. The reaction results from release of histamine by skin immune cells.

\medterm{Dermatographism} Hives that appear as raised linear welts within minutes after the skin is stroked or scratched. They result from an exaggerated release of histamine and may occur with other allergic or hive disorders.

\medterm{Dermatologic} Relating to the skin, hair, nails, or medical care for these structures. Dermatologic problems include rashes, infections, wounds, hair loss, nail changes, acne, and skin cancers.

\medterm{Dermatologic Disorder} A disease affecting the skin, hair, nails, or skin glands. It may cause rash, scaling, pigment change, lumps, ulcers, pain, itching, hair loss, or abnormal nail growth.

\medterm{Dermatologist} A physician who specializes in diagnosing and treating disorders of the skin, hair, nails, and related mucous membranes, including infections, inflammatory diseases, tumors, and cosmetic concerns.

\medterm{Dermatology} The medical specialty concerned with preventing, diagnosing, and treating diseases and injuries of the skin, hair, nails, and related mucous membranes.

\medterm{Dermatome} Either a surgical tool that removes a thin layer of skin for a graft or a strip of skin supplied by one spinal nerve. The meaning depends on whether the term is being used in surgery or nerve anatomy.
\medterm{Dermatoscope} A handheld optical instrument that magnifies and illuminates skin lesions, often with polarized or nonpolarized light. It reveals structures below the skin surface and supports clinical assessment of pigmented and other lesions.

\medterm{Dermatomycosis} A fungal infection of the skin, hair, or nails. It may cause scaling, itching, discoloration, cracking, or hair loss; treatment depends on the site and may require a topical or oral antifungal medicine.

\medterm{Dermatomyositis} An autoimmune disease causing muscle inflammation and weakness with characteristic skin changes. It may affect swallowing or breathing muscles, and some adults have an increased risk of associated cancer.

\medterm{Dermatopathology} The study of skin disease by examining tissue samples with a microscope and related laboratory methods. It helps distinguish inflammatory, infectious, harmless, and cancerous skin disorders.

\medterm{Dermatophyte} A fungus that grows in keratin, the tough protein in skin, hair, and nails. It causes tinea infections such as athlete’s foot, ringworm, and fungal nail disease.

\medterm{Dermatophytosis} A superficial fungal infection of the skin, hair, or nails. It can cause ring-shaped rashes, scaling, itching, cracked skin, or thickened discolored nails.

\medterm{Dermatoplasty} Surgical repair or reconstruction of damaged skin, often using tissue rearrangement or a skin graft. It may restore coverage after injury, burns, removal of a lesion, or another operation.
\medterm{Dermatoses - Systemic} Skin changes caused by a disease affecting the whole body or several organs, such as autoimmune, hormone-related, infectious, nutritional, or cancerous disease. The skin finding may help identify the underlying illness.

\medterm{Dermatosis} Any disease or disorder of the skin, including inflammatory, infectious, allergic, inherited, or tumor-related conditions. It may cause rash, itching, scaling, pigment changes, lumps, ulcers, or pain.
\medterm{Dermis} The strong, flexible skin layer beneath the epidermis. It contains collagen and elastic fibers, blood vessels, nerves, hair follicles, and sweat and oil glands that support and nourish the skin.

\medterm{Dermoid} A usually benign cyst or tumor containing tissues such as skin, hair, oil glands, fat, or teeth. It may occur in the ovary or other sites and can cause pressure, rupture, infection, or twisting of an organ.

\medterm{Dermoid Cyst} A usually benign cyst lined by tissue resembling skin and sometimes containing hair, oil glands, fat, or other skin structures. It may be present from birth or develop during life.

\medterm{Dermoid Cyst Of The Ovary} A usually benign ovarian cyst that can contain skin, hair, fat, glands, or teeth. It may cause pain, pressure, rupture, or twisting of the ovary and is assessed with examination and imaging.

\medterm{Dermoscopy} Examination of a skin lesion with magnification and special lighting that reveals structures not visible to the unaided eye. It helps assess pigmented and other lesions but does not replace biopsy when cancer is suspected.

\medterm{Descending Colon} The left-sided part of the large intestine, extending downward from the transverse colon toward the sigmoid colon. It absorbs water and salts and helps compact stool before it moves toward the rectum.

\medterm{Descending Thoracic Aorta} The part of the body’s main artery that runs from the aortic arch through the chest to the diaphragm. It gives branches to the chest wall and organs and continues as the abdominal aorta.

\medterm{Desensitization} A treatment that gradually exposes a person to increasing or controlled amounts of a specific allergen or medicine. It can reduce immune reactions but must be supervised because serious reactions may occur.

\medterm{Desferrioxamine} A medicine that binds excess iron and helps remove it from the body. It is used for serious iron overload and selected aluminum toxicity; prolonged treatment can affect hearing, vision, kidneys, or growth.

\medterm{Desflurane} An inhaled medicine used to maintain general anesthesia. It acts quickly but may irritate the airway, lower blood pressure, cause nausea, or rarely trigger malignant hyperthermia in susceptible people.

\medterm{Desiccate} To dry out or remove moisture from a tissue, substance, surface, or specimen. Excessive drying can damage skin or wounds and can alter medicines or laboratory samples.

\medterm{Desiccated Thyroid} A thyroid-hormone medicine made from dried animal thyroid tissue. It contains mainly thyroxine and triiodothyronine; excess treatment can cause palpitations, bone loss, or other effects, so blood tests and monitoring are needed.

\synonyms
Desiccated thyroid extract; natural desiccated thyroid

\medterm{Desipramine} A tricyclic antidepressant that increases norepinephrine activity. It is used for selected mental health or pain conditions but may cause dry mouth, constipation, drowsiness, abnormal heart rhythms, or dangerous overdose.

\medterm{Desipramine Hydrochloride Overdose} Taking too much desipramine hydrochloride can cause severe drowsiness, confusion, seizures, low blood pressure, coma, and life-threatening heart-rhythm or conduction problems. It requires emergency assessment, even if symptoms are initially mild.

\medterm{Desmin} A structural protein that links parts of muscle cells and helps maintain their alignment. Changes in the desmin gene or abnormal desmin buildup can cause muscle weakness and heart-muscle disease.

\medterm{Desmoid Tumor} A locally aggressive growth of fibrous tissue that usually does not spread to distant organs but can invade nearby structures and recur. It may be associated with familial adenomatous polyposis; treatment depends on location and symptoms.

\medterm{Desmoid Tumors} Fibrous growths that usually do not spread to distant organs but can invade nearby tissues and recur. Some remain stable; observation, medicine, or surgery may be considered according to symptoms and location.

\medterm{Desmoid Tumour} A growth of fibrous connective tissue that usually does not spread elsewhere but can grow into nearby structures and return after treatment. Its management depends on size, site, growth, and symptoms.

\medterm{Desmoplasia} Formation of dense, scar-like tissue around a tumor, long-lasting inflammation, or an injury. Around some cancers it can make a lump feel firm or fixed, but its significance depends on the cells, location, and other findings.

\medterm{Desmoplastic} Describing the formation of dense, scar-like fibrous tissue, often around a tumor or in response to chronic inflammation or injury.

\medterm{Desmoplastic Fibroblastoma} A rare, usually harmless tumor made of collagen-producing cells. It is typically a slow-growing, painless lump in superficial or deep soft tissue and is treated according to its size, site, and symptoms.

\medterm{Desmoplastic Fibroma} A rare noncancerous but locally aggressive fibrous growth in bone. It can cause pain, swelling, or fracture and may return after treatment.

\medterm{Desmoplastic Infantile Ganglioglioma} A rare, usually slow-growing brain tumor of infancy containing both nerve-cell and support-cell components. It may cause seizures, increasing head size, vomiting, or other signs of pressure in the brain.

\medterm{Desmoplastic Melanoma} A rare form of skin melanoma with abundant scar-like tissue. It often develops on sun-damaged skin, may have little or no pigment, and can grow deeply or recur locally.

\medterm{Desmoplastic Reaction} Formation of dense, scar-like fibrous tissue in response to injury, chronic inflammation, or invasion by a cancer. It may make affected tissue firm and can complicate interpretation of imaging or biopsy.
\medterm{Desmoplastic Small Round Cell Tumors} Plural form; see \textit{Desmoplastic small round cell tumor}.

\synonyms
DSRCT

\medterm{Desmoplastic Small Round Cell Tumor} A rare and aggressive soft-tissue sarcoma that usually develops in the abdomen or pelvis of children, adolescents, or young adults, especially males. It often forms several masses on the lining of the abdomen and can spread within the abdomen or to other sites. Symptoms may include abdominal swelling or pain, nausea, vomiting, diarrhea, or constipation. Tumor cells typically have an \textit{EWSR1-WT1} gene fusion, which helps confirm the diagnosis.

\medterm{Desmoplastic Trichoepithelioma} A rare harmless tumor arising from a hair follicle, usually appearing as a small firm facial bump or ring-shaped lesion. It can resemble a skin cancer and may require biopsy for confirmation.

\medterm{Desmopressin} A medicine that mimics vasopressin, a hormone that reduces urine production and helps control bleeding. It is used for central diabetes insipidus, some bleeding disorders, and selected nighttime urination problems; excess fluid can dangerously lower blood sodium.

\medterm{Desmosome} A strong junction that fastens neighboring cells together, especially in the skin and heart muscle. It links internal support fibers and helps tissues resist stretching and friction.
\medterm{Desogestrel} A progestin used in some hormonal contraceptives. It prevents pregnancy mainly by suppressing ovulation and thickening cervical mucus; irregular bleeding, mood changes, breast tenderness, or other hormonal effects may occur.

\medterm{Desonide} A low-strength corticosteroid applied to the skin to reduce redness, itching, and inflammation from conditions such as eczema or dermatitis. Excessive or prolonged use can thin the skin or cause other local effects.

\medterm{Desoximetasone} A corticosteroid applied to the skin to reduce inflammation and itching in responsive rashes. Strength, body site, amount, and duration affect the risk of skin thinning and absorption into the body.

\medterm{Desoxycorticosterone} A steroid hormone with mineralocorticoid effects. It makes the kidneys retain sodium and water and lose potassium, thereby influencing blood volume and blood pressure.

\medterm{Desquamation} Shedding or peeling of the outer layer of skin or another body lining. It may be normal or result from dryness, irritation, infection, inflammation, burns, medicines, or another skin disorder.

\medterm{Desquamative Interstitial Pneumonia} A rare inflammatory lung disease strongly associated with smoking. Immune cells collect in the air sacs, causing cough and breathlessness; stopping smoking is essential, and additional treatment may be needed.

\medterm{Desulfovibrio} A group of bacteria that can live without oxygen and use sulfate during metabolism. They occur in the environment and intestine; invasive human infection is uncommon and must fit the patient’s symptoms and specimen findings.

\medterm{Desvenlafaxine Succinate} A form of desvenlafaxine, an antidepressant that increases serotonin and norepinephrine signaling. It is used for major depressive disorder and may cause nausea, insomnia, increased blood pressure, or withdrawal symptoms if stopped suddenly.

\medterm{Desynchronosis (Jet Lag)} Alternate index label for Jet Lag.

\medterm{Detached} Separated or no longer connected. In medicine it may describe a retina, tissue, device, or feeling of emotional distance; the clinical meaning depends on the surrounding term and findings.

\medterm{Detached Retina} Separation of the light-sensitive retina from the tissue beneath it. Flashes, new floaters, or a curtain-like shadow may occur; it is an emergency because delayed treatment can cause permanent vision loss.

\medterm{Detergent} A cleaning substance that helps remove oils and dirt from surfaces. Some detergents can irritate or burn the skin, eyes, airways, or digestive tract, especially when concentrated or swallowed.

\medterm{Detergent Poisoning} Harmful exposure to a detergent by swallowing, inhaling, or contact with the eyes or skin. Symptoms depend on the product and may include irritation, vomiting, chemical burns, coughing, or breathing difficulty.

\medterm{Deterministic Effects} Tissue injuries from radiation that usually occur only after exposure exceeds a threshold. Higher exposure generally causes more severe injury, such as skin damage, cataracts, or tissue death.

\medterm{Detoxification} Medical treatment that reduces the effects of a poison or helps the body remove it. Care may include supportive treatment, an antidote, activated charcoal, or dialysis; stomach washing is rarely appropriate.

\medterm{Detritic Synovitis} Inflammation of the lining of a joint caused by loose particles from damaged cartilage, bone, blood, crystals, or a worn implant. It can cause swelling, pain, stiffness, and reduced movement.

\medterm{Detrusor Instability} Involuntary contractions of the bladder muscle while the bladder is filling. They can cause sudden urinary urgency, frequent urination, nighttime urination, or urge leakage.

\medterm{Detrusor Muscle} The muscle forming the bladder wall. It relaxes while urine collects and contracts to empty the bladder; abnormal activity can cause urgency or leakage, while weak activity can make emptying difficult.

\medterm{Detumescence} Return of erectile tissue to its non-erect state after sexual stimulation or an erection ends. It occurs when blood leaves the erectile tissue and pressure falls.

\medterm{Deubiquitinating Enzymes} Enzymes that remove a small regulatory tag from other proteins. This can change whether a protein is broken down, where it works, and how cells control signals, repair DNA, and respond to infection.

\medterm{Deubiquitination} Removal of the ubiquitin tag from a protein. This process can prevent the protein from being broken down or can change its activity, location, or role in cell signaling.

\medterm{Deuteranopia} An inherited red-green color-vision deficiency caused by missing green-sensitive pigment in the retina. People may have difficulty distinguishing some shades of green, red, brown, and orange.

\medterm{Deuterium} A naturally occurring form of hydrogen whose nucleus contains one proton and one neutron. Small amounts occur in water; very high exposure to heavy water can disrupt cell function.

\medterm{Development} The progressive growth and change by which a person, organ, or ability acquires its structure and function. It includes physical, brain, language, emotional, social, and adaptive changes.

\medterm{Developmental Coordination Disorder} A childhood condition in which movement skills develop more slowly or less accurately than expected and interfere with dressing, writing, sports, school, or daily activities. It is not explained only by another medical or neurologic condition.

\medterm{Developmental Differences Of The Female Genital Tract} Differences in the formation or structure of the uterus, cervix, vagina, fallopian tubes, or external genitalia that arise before birth. Effects range from none to infertility, pain, menstrual problems, or pregnancy complications.
\medterm{Developmental Disabilities} Lifelong conditions that begin during childhood and affect movement, learning, thinking, language, behavior, or ability to manage daily activities. Support needs vary and may include medical, educational, communication, and social services.

\medterm{Developmental Disorder} A condition that begins during childhood and affects physical, intellectual, language, learning, behavioral, or social development. It may interfere with school, relationships, communication, or daily independence.

\medterm{Developmental Delay} A child’s acquisition of movement, language, thinking, social, or self-care skills is slower than expected for age. Delay may affect one or several areas and can result from genetic, neurologic, sensory, medical, or environmental causes.

\medterm{Developmental Disorders} Conditions beginning during childhood that affect thinking, language, movement, learning, behavior, or social functioning. Assessment considers developmental history, daily impact, coexisting conditions, strengths, and needed educational or clinical support.

\medterm{Developmental Dysplasia Of The Hip} Abnormal development of an infant’s or child’s hip, ranging from a shallow socket to partial or complete dislocation. Early detection and treatment help the hip develop normally and reduce later pain or arthritis.

\medterm{Developmental Expressive Language Disorder} A childhood disorder causing persistent difficulty expressing ideas with spoken words, sentences, or grammar. Hearing, development, autism, and other possible causes should be assessed before diagnosis.
\medterm{Developmental Milestone} A skill that most children acquire within a broad expected age range, such as smiling, sitting, walking, speaking, or interacting. A missed or lost milestone may require developmental assessment.

\medterm{Developmental Milestones} Skills children commonly acquire as they grow, including movement, language, thinking, social interaction, and self-care. Ages vary; delayed or lost skills should be discussed with a healthcare professional.

\medterm{Developmental Milestones Record} A written or electronic record of a child’s movement, language, thinking, social, and self-care skills. It helps follow progress, identify possible delay, and guide further assessment.

\medterm{Developmental Milestones Record - 12 Months} Skills commonly seen around one year include pulling to stand, moving while holding furniture, pointing or using gestures, responding to simple requests, and interactive communication. Children develop at different rates.

\medterm{Developmental Milestones Record - 18 Months} Skills commonly seen around 18 months include walking, using several words, following simple directions, pointing to show interest, and beginning pretend or imitative play. Individual variation is normal.

\medterm{Developmental Milestones Record - 2 Months} Skills commonly seen around two months include social smiling, watching a caregiver or moving object, reacting to loud sounds, making sounds, and holding the head up briefly. Individual development varies.

\medterm{Developmental Milestones Record - 2 Years} Skills commonly seen around two years include running, using two-word phrases, following simple two-step directions, pointing to body parts, and playing near other children. Individual development varies.

\medterm{Developmental Milestones Record - 3 Years} Skills commonly seen around three years include speaking in short sentences, climbing stairs with alternating feet, drawing simple shapes, using a fork, and joining brief interactive play. Individual development varies.

\medterm{Developmental Milestones Record - 4 Months} Skills commonly seen around four months include steady head control, reaching for objects, bringing hands to the mouth, laughing, smiling to gain attention, and responding to caregivers. Individual development varies.

\medterm{Developmental Milestones Record - 4 Years} Skills commonly seen around four years include speaking in understandable sentences, hopping on one foot, drawing a person with several parts, dressing with some help, and engaging in imaginative play.

\medterm{Developmental Milestones Record - 5 Years} Skills commonly seen around five years include hopping or skipping, telling a simple story, counting, drawing recognizable figures or letters, dressing with little help, and following game rules.

\medterm{Developmental Milestones Record - 6 Months} Skills commonly seen around six months include rolling, sitting with support, reaching for and transferring objects, babbling, responding to sounds, and recognizing familiar people. Individual development varies.

\medterm{Developmental Milestones Record - 9 Months} Skills commonly seen around nine months include sitting independently, moving purposefully, babbling, responding to a name, looking for a hidden object, and picking up small objects with finger and thumb.

\medterm{Developmental Process} The physical, brain, emotional, social, and functional changes through which a person or organ grows and acquires abilities. Development is influenced by genes, health, nutrition, relationships, learning, and environment.

\medterm{Developmental Psychology} The study of how thinking, emotions, relationships, behavior, and abilities change from infancy through old age. It examines the effects of genes, health, learning, experience, relationships, and environment.

\medterm{Developmental Reading Disorder} A learning disorder causing persistent difficulty learning accurate, fluent reading. Problems may include recognizing words, linking letters to sounds, spelling, or understanding written language despite appropriate teaching.

\medterm{Developmental Regression} Loss of skills a child previously acquired, such as speech, movement, social interaction, or self-care. Regression is a warning sign requiring prompt evaluation for neurologic, genetic, metabolic, infectious, or other causes.

\medterm{Developmental Stage} A period of growth characterized by typical physical, thinking, emotional, social, language, or self-care changes. Children vary, so stages describe broad patterns rather than exact deadlines.

\medterm{Developmental Surveillance} Ongoing monitoring of a child’s development through caregiver concerns, observation, history, and milestone review during routine care. It can identify concerns early but differs from formal screening with a standardized tool.

\medterm{Deviated Septum} A nasal partition displaced away from the center, often from normal development or injury. A marked deviation can block one side, causing difficulty breathing through the nose, congestion, nosebleeds, sinus problems, or sleep symptoms.

\medterm{Deviated Nasal Septum} Alternate name for deviated septum, in which the wall dividing the nasal passages is displaced from the midline.

\synonyms
Septum Deviated (Deviated Septum)

\medterm{Deviation} A difference from an expected position, direction, value, or pattern. In medicine, it may describe anatomy, movement, a laboratory result, or a measurement and must be interpreted in context.

\medterm{Device Embolisation} Movement of a medical device or broken piece of one through the bloodstream until it lodges in and blocks a blood vessel. It can reduce blood flow and may cause pain, organ injury, or an emergency requiring removal.

\medterm{Devices For Hearing Loss} Equipment that improves access to sound, including hearing aids, cochlear implants, bone-conduction devices, alerting systems, and communication accessories. Choice depends on the type and severity of hearing loss and personal needs.

\medterm{Dexamethasone} A strong, long-acting corticosteroid that reduces inflammation and immune activity. It is used for selected allergic, brain-swelling, cancer-related, endocrine, and other conditions; risks include infection, high blood sugar, mood changes, and adrenal suppression.

\medterm{Dexamethasone Suppression Test} A test of cortisol control. After taking a small dose of dexamethasone, blood or urine cortisol is measured; failure to suppress may suggest excess cortisol but can have other causes and needs clinical interpretation.

\medterm{Dexamphetamine} A stimulant that increases dopamine and norepinephrine activity in the brain. It is used for selected attention or sleep disorders and may cause reduced appetite, insomnia, fast heartbeat, high blood pressure, or dependence.

\medterm{Dexibuprofen} A nonsteroidal anti-inflammatory medicine used for short-term pain and inflammation. It can irritate the stomach and may cause ulcers, bleeding, kidney injury, fluid retention, or cardiovascular problems.

\medterm{Dexketoprofen} A nonsteroidal anti-inflammatory medicine used for short-term pain. It reduces substances that cause pain and inflammation but can cause stomach bleeding, kidney injury, fluid retention, heart problems, or allergic reactions.

\medterm{Dexlansoprazole} A proton-pump inhibitor that reduces stomach-acid production. It is used for acid reflux and related esophageal injury; prolonged use may contribute to infections, low magnesium or vitamin B12, fractures, or kidney inflammation.

\medterm{Dexmedetomidine} A sedative that activates alpha-2 receptors and is given by injection during monitored care or intensive care. It usually depresses breathing less than many sedatives but can slow the heart or lower blood pressure.

\medterm{Dexmethylphenidate} A stimulant used to treat attention-deficit/hyperactivity disorder. It increases dopamine and norepinephrine activity and may cause reduced appetite, insomnia, headache, increased heart rate or blood pressure, or misuse.

\medterm{Dexrazoxane} A medicine used in selected cancer treatments to reduce heart damage from anthracyclines or tissue injury after leakage of an anthracycline from a vein. It can suppress bone marrow and increase infection risk.

\medterm{Dexrazoxane Hydrochloride} A salt form of dexrazoxane used in selected settings to reduce anthracycline-related heart damage or tissue injury after medicine leakage from a vein. Its use requires specialist dosing and monitoring.
\medterm{Dextrin} A group of shorter carbohydrate chains formed when starch breaks down. Dextrins are used in foods, laboratory products, and medicines and are generally digested as carbohydrates.

\medterm{Dextroamphetamine} A stimulant that increases dopamine and norepinephrine activity in the brain. It is used for selected attention or sleep disorders and may cause insomnia, reduced appetite, fast heartbeat, high blood pressure, or dependence.

\medterm{Dextrocardia} A congenital condition in which the heart points or lies mainly on the right side of the chest. It may occur alone or with reversed organs and structural heart defects.

\medterm{Dextromethorphan} A cough suppressant that reduces activity in brain pathways controlling the cough reflex. Excessive amounts can cause drowsiness, agitation, hallucinations, seizures, breathing depression, or serotonin toxicity.

\medterm{Dextromethorphan Overdose} Taking too much dextromethorphan can cause drowsiness, confusion, agitation, hallucinations, fast heartbeat, seizures, breathing depression, or serotonin toxicity. It requires urgent poison or emergency assessment.

\medterm{Dextroposition Of The Heart} The heart is pushed toward the right side of the chest by nearby organs, the diaphragm, lung disease, or another space-occupying problem. The heart itself may be normally formed and is not necessarily reversed as in dextrocardia.

\medterm{Dextrose} D-glucose, a simple sugar used for energy. Oral or intravenous dextrose can treat or prevent low blood glucose, but excessive amounts can raise blood sugar and disturb body fluids.

\medterm{DFSP} Alternate index label; see \textit{Dermatofibrosarcoma protuberans}.

\medterm{DHA} Docosahexaenoic acid, an omega-3 fatty acid that is an important part of brain and retinal cell membranes. It is obtained mainly from oily fish, some algae, and supplements.

\medterm{DHEA} Dehydroepiandrosterone, a hormone precursor made mainly by the adrenal glands. The body can convert it into male and female sex hormones; levels vary with age, illness, and some hormone disorders.

\medterm{DHEA-Sulfate Test} A blood test measuring a hormone made mainly by the adrenal glands. It can help investigate early puberty, excess body hair, some hormone disorders, congenital adrenal hyperplasia, or an adrenal tumor; results vary with age and sex.

\medterm{Diabetes And Alcohol} Alcohol can cause delayed low blood glucose, interact with medicines, and make glucose control harder. Risk depends on amount, timing, food intake, diabetes type, liver health, and treatment.

\medterm{Diabetes And Exercise} Physical activity can improve glucose control, fitness, and heart health. People using insulin or medicines that lower glucose may need to check levels and adjust food or treatment with professional guidance to prevent low glucose.

\medterm{Diabetes And Eye Disease} Diabetes can damage retinal blood vessels, causing retinopathy and macular swelling that may reduce vision. Good glucose and blood-pressure control and regular dilated eye examinations help prevent or limit sight loss.

\medterm{Diabetes And Eye Problems} Diabetes may cause retinopathy, macular swelling, cataracts, glaucoma, and vision loss. Glucose and blood-pressure control, regular dilated eye examinations, and timely laser, injection, or surgical treatment reduce the risk of severe loss.

\medterm{Diabetes And Kidney Disease} Kidney damage caused by diabetes. It may begin with albumin leaking into urine and progress to reduced kidney function or kidney failure; glucose, blood-pressure, and kidney monitoring can slow progression.

\medterm{Diabetes And Metabolic Health} Diabetes is linked with insulin resistance, excess body fat, abnormal cholesterol, and high blood pressure. Together these problems increase the risk of heart disease and can worsen blood-glucose control.

\medterm{Diabetes And Nerve Damage} Long-term diabetes can injure nerves, causing burning pain, tingling, numbness, weakness, digestive or bladder problems, sexual difficulties, or loss of protective feeling in the feet.

\medterm{Diabetes Complications} Problems caused by diabetes, including low or high glucose emergencies and long-term damage to the eyes, kidneys, nerves, heart, blood vessels, teeth, and feet. Regular monitoring and risk-factor control reduce risk.

\medterm{Diabetes During Pregnancy (Gestational Diabetes)} Alternate index label for Gestational Diabetes.


\medterm{Diabetes Eye Exams} Eye examinations that look for diabetic retinopathy, macular swelling, cataracts, glaucoma, and other sight-threatening problems. Dilated examination or retinal photography may detect changes before vision is affected.

\medterm{Diabetes - Foot Ulcers} Open sores on the foot in a person with diabetes, often caused or worsened by reduced sensation, pressure, deformity, or poor circulation. They can become infected or reach bone, so prompt assessment and wound care are important.

\medterm{Diabetes, Gestational} Alternate index label; see \textit{Gestational diabetes}.

\medterm{Diabetes Insipidus} A disorder causing large amounts of very dilute urine and intense thirst because vasopressin is lacking or the kidneys do not respond to it. It is unrelated to diabetes mellitus and can cause dehydration.

\medterm{Diabetes - Insulin Therapy} Treatment with injected or infused insulin to replace missing insulin or improve glucose control. Doses are individualized and balanced with meals, activity, illness, and the risk of dangerously low glucose.


\medterm{Diabetes Management} Coordinated care using food choices, activity, medicines when needed, glucose monitoring, education, and screening for complications. The plan is adjusted for diabetes type, health conditions, life stage, and personal goals.

\medterm{Diabetes Medications} Medicines that lower blood glucose or replace insulin. They work in different ways, and the choice depends on the type of diabetes, kidney and heart health, other medicines, cost, and treatment goals.

\medterm{Diabetes Mellitus} A group of diseases in which blood glucose, also called blood sugar, stays too high. Insulin is a hormone made by the pancreas that helps glucose move from the blood into cells for energy. Diabetes develops when the body makes too little or no insulin, cannot use insulin properly, or both. The main types are type 1 diabetes, type 2 diabetes, and gestational diabetes, which develops during pregnancy; other forms can result from genetic conditions, pancreatic disease, or medicines. Persistently high blood glucose can damage the eyes, kidneys, nerves, heart, and blood vessels. Diabetes mellitus is different from diabetes insipidus, a disorder of water balance that is not caused by high blood glucose.

\synonyms
Diabetes

\medterm{Diabetes Mellitus Type 1} An autoimmune disorder in which the immune system destroys insulin-producing pancreatic cells, causing little or no insulin. It can begin at any age and requires insulin; untreated deficiency can cause ketoacidosis.

\medterm{Diabetes Mellitus Type 2} A disorder in which the body resists insulin and, over time, may not make enough of it. Risk is increased by family history, excess body weight, and inactivity; it can occur at any age and may have no early symptoms.


\medterm{Diabetes Pre (Prediabetes)} Alternate index label for Prediabetes.

\medterm{Diabetes - Preventing Heart Attack And Stroke} Reducing cardiovascular risk in diabetes through glucose, blood-pressure, and cholesterol control, not smoking, regular activity, healthy eating, weight management when appropriate, and indicated preventive medicines.




\medterm{Diabetes Prevention} The risk of type 2 diabetes can often be reduced through regular physical activity, nutritious meals, avoiding tobacco, adequate sleep, and weight loss when appropriate. People at increased risk may benefit from structured prevention programs and regular testing.

\medterm{Diabetes Symptoms In Men} Diabetes in men may cause thirst, frequent urination, fatigue, blurred vision, slow wound healing, recurrent infections, erectile dysfunction, or loss of sensation. Some people have no symptoms, so risk-based blood testing is important.

\medterm{Diabetes Symptoms In Women} Diabetes in women may cause thirst, frequent urination, fatigue, blurred vision, recurrent vaginal or urinary infections, slow wound healing, or pregnancy complications. Symptoms can be absent, and screening depends on age, risk factors, and pregnancy status.


\medterm{Diabetes Treatment} Diabetes treatment combines nutrition, physical activity, glucose monitoring, education, and medicines chosen for the type of diabetes and individual risks. Care also addresses blood pressure, lipids, kidneys, eyes, nerves, feet, vaccines, and cardiovascular health.

\medterm{Diabetes, Type 1} Alternate index label; see \textit{Type 1 diabetes}.


\medterm{Diabetes Type 1 (Type 1 Diabetes)} Alternate index label for Type 1 Diabetes.

\medterm{Diabetes, Type 2} Alternate index label; see \textit{Type 2 diabetes}.


\medterm{Diabetes Type 2 - Meal Planning} Meal planning for type 2 diabetes emphasizes suitable portions, fiber-rich foods, carbohydrate awareness, adequate protein, and unsaturated fats. The plan should fit medicines, kidney health, cultural foods, preferences, and glucose goals.

\medterm{Diabetes Type 2 (Type 2 Diabetes)} Alternate index label for Type 2 Diabetes.

\medterm{Diabetic Angiopathy} Damage to small or large blood vessels caused by diabetes. It can reduce circulation, injure the eyes, kidneys, heart, brain, or feet, and increase the risk of wounds, heart attack, or stroke.

\medterm{Diabetic Autonomic Neuropathy} Diabetes-related damage to nerves that control automatic body functions. It may cause dizziness on standing, fast heartbeat, sweating changes, delayed stomach emptying, bladder problems, sexual dysfunction, or reduced awareness of low glucose.

\medterm{Diabetic Cardiomyopathies} Heart-muscle disease associated with diabetes that is not fully explained by blocked heart arteries or high blood pressure. It can make the heart relax or pump poorly and may lead to heart failure.

\medterm{Diabetic Cardiomyopathy} Heart-muscle damage associated with diabetes, after other causes such as blocked arteries and high blood pressure are considered. It may cause poor relaxation, weak pumping, heart failure, or rhythm problems.

\medterm{Diabetic Coma} Profound unconsciousness caused by a severe diabetes-related problem, such as extremely low or high blood glucose, ketoacidosis, severe dehydration, or abnormal body chemistry. It is a medical emergency.

\medterm{Diabetic Complications} Acute or long-term problems caused by diabetes, including low or high glucose emergencies and damage to the eyes, kidneys, nerves, heart, blood vessels, teeth, and feet.

\medterm{Diabetic Diet} An individualized eating pattern for diabetes that balances carbohydrate portions, fiber, protein, and healthy fats with medicines, activity, culture, and health goals. It supports glucose control without requiring the same foods or schedule for everyone.

\synonyms
Diet Diabetes (Diabetic Diet)

\medterm{Diabetic Diet For Type 2 Diabetes} An eating pattern for type 2 diabetes emphasizing suitable portions, high-fiber foods, carbohydrate awareness, adequate protein, and unsaturated fats. Medicines, kidney function, culture, preferences, and glucose or weight goals guide the plan.

\medterm{Diabetic Foot} Foot problems associated with diabetes, including loss of sensation, deformity, ulcers, infection, poor circulation, and slow healing. Daily inspection, proper footwear, glucose control, and early care help prevent serious injury.

\medterm{Diabetic Foot Disease} Foot damage caused by a combination of diabetes-related nerve loss, poor circulation, deformity, and impaired healing or immunity. A small injury can progress to an ulcer, deep infection, or amputation without early care.

\medterm{Diabetic Foot Ulcer} An open wound on the foot caused or worsened by diabetes-related loss of sensation, pressure, deformity, poor circulation, or slow healing. It may be painless and can become deeply infected without prompt treatment.

\medterm{Diabetic Gastroparesis} A complication of diabetes in which damage to the stomach's nerves slows the movement of food from the stomach into the small intestine. It may cause early fullness, nausea, vomiting, bloating, abdominal pain, and difficult-to-control blood sugar levels.


\medterm{Diabetic Hyperglycemic Hyperosmolar Syndrome} A life-threatening diabetes emergency with extremely high blood glucose, severe dehydration, and concentrated blood, usually without major ketone buildup. It can cause confusion, seizures, coma, and shock and requires hospital treatment.

\medterm{Diabetic Hyperosmolar Hyperglycemic State} A severe diabetes emergency marked by extremely high glucose, severe dehydration, and increased blood concentration with little or no ketoacidosis. Infection, illness, or inadequate treatment may trigger it; urgent hospital care is required.

\medterm{Diabetic Hypoglycemia} Blood glucose that is too low in a person with diabetes, often because of insulin, certain medicines, missed food, alcohol, or unusual exercise. Shaking, sweating, confusion, seizures, or unconsciousness may occur.

\medterm{Diabetic Ketoacidosis} A life-threatening emergency caused by too little effective insulin, leading to high glucose, ketone buildup, and acidic blood. It can cause thirst, frequent urination, vomiting, abdominal pain, deep breathing, confusion, and coma.

\synonyms
DKA

\medterm{Diabetic Ketoacidosis Monitoring} Monitoring during treatment includes frequent glucose checks, blood salts, kidney function, ketones, acidity, fluid balance, urine output, and mental status. Close monitoring helps guide fluids, insulin, and electrolyte replacement safely.

\medterm{Diabetic Macular Edema} Swelling of the macula, the central part of the retina, caused by leaking blood vessels in diabetes. It can blur or distort central vision and may occur at any stage of diabetic retinopathy.

\medterm{Diabetic Mastopathy} A rare, noncancerous, firm breast thickening associated mainly with long-standing diabetes. It can feel or look like breast cancer, so examination and imaging—and sometimes a biopsy—may be needed to establish the diagnosis.

\medterm{Diabetic Microvascular Complications} Diabetes-related damage to small blood vessels, mainly affecting the eyes, kidneys, and nerves. Good glucose, blood-pressure, and cholesterol control, along with regular screening, lowers the risk of progression.
\medterm{Diabetic Nephropathy} Kidney damage caused by diabetes. It often begins with albumin leaking into urine and may progress to reduced filtering ability or kidney failure; regular urine and blood tests help detect and monitor it.

\synonyms
DKD

\medterm{Diabetic Nephropathy (Kidney Disease)} Kidney damage caused by diabetes, often beginning with albumin in the urine and potentially progressing to chronic kidney disease or kidney failure. Blood and urine monitoring can detect it early.

\medterm{Diabetic Neuropathy} Nerve damage caused by diabetes. It may cause burning pain, tingling, numbness, weakness, loss of protective feeling in the feet, or problems with digestion, blood pressure, sweating, bladder control, or sexual function.

\synonyms
Autonomic Neuropathy Diabetic (Diabetic Neuropathy)
Neuropathy, Diabetic
Neuropathy Autonomic Diabetic (Diabetic Neuropathy)
Neuropathy Diabetic (Diabetic Neuropathy)
Neuropathy Diabetic Peripheral (Diabetic Neuropathy)
Neuropathy Focal Diabetic (Diabetic Neuropathy)
Neuropathy Proximal Diabetic (Diabetic Neuropathy)
Peripheral Neuropathy Diabetic (Diabetic Neuropathy)
Proximal Neuropathy Diabetic (Diabetic Neuropathy)

\medterm{Diabetic Peripheral Neuropathy} Diabetes-related damage to nerves, usually beginning in the toes and feet and later affecting the hands. It can cause numbness, burning pain, weakness, and loss of protective sensation that increases injury risk.

\medterm{Diabetic Retinopathy} Damage to retinal blood vessels caused by diabetes. Early disease may have no symptoms; advanced disease can cause macular swelling, bleeding, retinal detachment, and vision loss. Regular dilated eye examinations are important.

\synonyms
Retinopathy, Diabetic

\medterm{Diabetogenic} Able to cause or promote diabetes or high blood glucose. The term may describe a disease, hormone, medicine, genetic change, or other factor that reduces insulin action or increases glucose production.

\medterm{Diabetologist} A physician specializing in diabetes prevention, diagnosis, and treatment, including glucose management and complications involving the heart, kidneys, nerves, eyes, and feet.

\medterm{Diabrosis} An old medical term for erosion or corrosion that eats through tissue and may cause an ulcer or perforation, especially in a blood-vessel or organ wall. Modern reports usually name the specific cause and injury.

\medterm{Diacetoxyscirpenol} A toxin produced by some Fusarium molds that can contaminate food or animal feed. Exposure may irritate the digestive tract, suppress blood-cell production, and impair immune function.


\medterm{Diacylglycerol} A fat-like molecule made from cell membranes that acts as an internal chemical signal. It helps control cell growth, movement, secretion, metabolism, and responses to hormones or other signals.

\medterm{Diagnosis} Identification of a disease or health condition using symptoms, medical history, examination, and appropriate tests. A clinician may compare several possible causes before reaching the most likely diagnosis.

\medterm{Diagnostic} Relating to identifying a disease, condition, or cause of symptoms. A diagnostic test provides information that must be interpreted with the person’s history, examination, and other findings rather than alone.

\medterm{Diagnostic Imaging} Use of images to detect, locate, or monitor disease or injury. Methods include X-ray, ultrasound, computed tomography, magnetic resonance imaging, and nuclear imaging; findings must be interpreted with the clinical situation.

\medterm{Diagnostic Laparoscopy} A minimally invasive procedure in which a camera is inserted through small abdominal incisions to inspect pelvic or abdominal organs. The clinician may take tissue samples or perform treatment during the same procedure.

\medterm{Diagnostic Microbiology} Laboratory testing used to identify microorganisms that may cause disease. Methods include microscopy, culture, antigen or antibody tests, and genetic tests; results must be interpreted with symptoms and the specimen site.

\medterm{Diagnostic Procedure} An examination or procedure used to identify, exclude, locate, or monitor a disease or injury. The purpose, preparation, risks, accuracy, and follow-up depend on the specific procedure and patient.

\medterm{Diagnostic Radiology} The medical specialty that uses imaging to diagnose and characterize disease or injury. It includes X-ray, ultrasound, computed tomography, magnetic resonance imaging, and selected nuclear-medicine studies.

\medterm{Diagnostic Service} A healthcare service that collects or interprets information to detect, identify, characterize, or monitor a disease or health condition, such as laboratory testing, imaging, or specialist assessment.

\medterm{Diagnostic Test} A test used to evaluate symptoms, signs, or an abnormal screening result and help confirm, exclude, or characterize a disease. Its value depends on accuracy, timing, the patient’s condition, and how the result will guide care.

\medterm{Diagnostic Trial} A clinical study that evaluates how accurately a test detects, rules out, or predicts a disease or important health outcome. Researchers may compare it with an established test or another accepted way of determining the diagnosis.

\medterm{Diagnostic Uncertainty} A situation in which available information does not establish one diagnosis with enough confidence. It is managed by explaining uncertainty, arranging follow-up, watching for warning signs, and ordering targeted further tests when appropriate.

\medterm{Diagnostics} The tests, examinations, and clinical reasoning used to detect, identify, classify, or monitor a health condition. Results must be interpreted in context because no test is perfect.

\medterm{Diakinesis} The final stage before the first division that produces eggs or sperm. Paired chromosomes become tightly packed and remain connected at points where they exchanged genetic material.

\medterm{Dialectical Behavior Therapy} A structured talking treatment that combines behavior-change skills with mindfulness and acceptance. It teaches ways to manage intense emotions, tolerate distress, improve relationships, and reduce harmful behaviors.

\medterm{Dialister Pneumosintes} A bacterium found in the mouth and respiratory tract that grows without oxygen. It may contribute to mixed infections, but a positive sample must be interpreted with symptoms and collection site.

\medterm{Dialysis} Treatment that replaces part of the kidneys’ work by removing waste, excess water, and some chemicals from the blood. The main types are hemodialysis and peritoneal dialysis.

\medterm{Dialysis Access} A route used to connect a person’s blood circulation to a hemodialysis machine. Options include an arteriovenous fistula, graft, or central venous catheter; each has different risks of infection, clotting, and inadequate flow.

\medterm{Dialysis Adequacy} How well dialysis removes waste and extra fluid and keeps body salts and acid levels within a safe range. Blood tests, treatment records, symptoms, nutrition, and weight help determine whether dialysis is sufficient.


\medterm{Dialysis Disequilibrium Syndrome} A rare brain reaction during or soon after starting or intensifying hemodialysis, caused by rapid changes in dissolved substances. Headache, nausea, confusion, seizures, or coma may occur; gradual treatment reduces risk.

\medterm{Dialysis - Hemodialysis} A form of dialysis in which blood travels through a filter outside the body to remove waste and excess fluid before returning to the patient. It requires reliable blood access and regular monitoring.

\medterm{Dialysis Machine} Equipment that moves blood or dialysis fluid across a filter to remove waste, excess water, and selected salts when the kidneys cannot do so adequately.

\medterm{Dialysis - Peritoneal} A form of dialysis in which cleansing fluid enters and leaves the abdomen through a catheter. The abdominal lining filters waste and water; cloudy fluid, pain, or fever may signal peritonitis and requires urgent care.

\medterm{Peritoneal Dialysis-Associated Peritonitis} Infection or inflammation of the abdominal lining during peritoneal dialysis, usually related to contamination of the dialysis catheter or fluid. It may cause cloudy drainage, abdominal pain, fever, or nausea.

\medterm{Dialysis Solutions} Specially prepared fluids used during hemodialysis or peritoneal dialysis. They allow waste and excess water to move out of the blood; their salt, glucose, and other contents are adjusted to the patient’s needs.

\medterm{Dialysis Therapy} Use of hemodialysis or peritoneal dialysis to replace part of kidney function by removing waste and excess fluid and correcting some salt and acid imbalances. Treatment is individualized according to kidney function, symptoms, health, and goals.

\medterm{Diamine} A molecule containing two amino groups. Diamines occur naturally in metabolism and are used in laboratory, pharmaceutical, and industrial chemistry; health effects depend on the specific compound and exposure.

\medterm{Diamond-Blackfan Anemia} A rare inherited disorder in which the bone marrow makes too few red blood cells, usually beginning in infancy. It may also cause short stature, thumb or facial differences, or heart defects.

\synonyms
DBA; Diamond-Blackfan syndrome; Congenital hypoplastic anemia


\medterm{Diamorphine} A strong opioid medicine used in selected settings for severe pain or during medical procedures. It can cause sleepiness, nausea, constipation, dependence, and dangerous breathing depression, especially in overdose or with other sedatives.

\medterm{Dianisidine} An industrial chemical used in some dyes and laboratory reagents. Significant or repeated exposure can irritate tissues and may damage blood or increase cancer risk; workplace protection is important.


\medterm{Diapedesis} Movement of blood cells, especially infection-fighting white cells, through the walls of tiny blood vessels into nearby tissue. This allows them to reach and respond to infection, injury, or inflammation.

\medterm{Diaper Rash} Inflammation of skin covered by a diaper, commonly caused by moisture, urine, stool, friction, or yeast infection. It causes redness and soreness; a persistent, severe, blistering, or spreading rash needs medical assessment.

\medterm{Diaphoresis} Excessive sweating. It may be a normal response to heat or exercise, or a sign of fever, pain, low blood glucose, anxiety, infection, heart problems, or another condition.

\medterm{Diaphoretic} Describing heavy sweating or something that causes sweating. Sudden cold sweating with chest pain, faintness, breathlessness, confusion, or weakness may signal a serious medical problem and needs urgent assessment.

\medterm{Diaphragm} The dome-shaped muscle separating the chest from the abdomen and the main muscle used for breathing. When it contracts, it flattens and draws air into the lungs; when it relaxes, air leaves.

\medterm{Diaphragm (Muscle)} The dome-shaped breathing muscle between the chest and abdomen. Contraction enlarges the chest cavity and draws air into the lungs; relaxation allows exhalation.

\medterm{Diaphragmatic Eventration} Abnormal thinning or weakness of part or all of the diaphragm, causing it to sit higher than usual. It may cause no symptoms or lead to breathlessness, especially if severe or present from birth.

\medterm{Diaphragmatic Hernia} Movement of abdominal organs through a hole or weak area in the diaphragm into the chest. It may be present from birth or follow injury and can impair breathing or cause bowel obstruction.

\medterm{Diaphragmatic Rupture} A tear in the diaphragm, usually caused by a severe blunt or penetrating injury to the chest or abdomen. Abdominal organs may move into the chest, causing breathing difficulty, obstruction, or strangulation.

\medterm{Diaphysis} The long central shaft of a bone such as the femur, humerus, or tibia. It contains a hollow marrow cavity and is surrounded by strong compact bone; fractures there may require prolonged healing.

\medterm{Diarrhea} Passing loose or watery stools more often than usual. Causes include infections, medicines, food intolerance, inflammation, malabsorption, or other illnesses. Diarrhea can lead to dehydration and, when prolonged, malabsorption.



\medterm{Diarrhea, Constipation, Gas, And Bloating} These digestive symptoms may result from diet, infection, medicines, food intolerance, constipation, or disorders such as irritable bowel syndrome. Persistent, severe, bloody, or unexplained symptoms require assessment.

\medterm{Diarrhea IBS (IBS-D Irritable Bowel Syndrome With Diarrhea)} Alternate index label for IBS-D Irritable Bowel Syndrome with Diarrhea.

\medterm{Diarrhea In Infants} Frequent loose or watery stools in a baby. Because infants can dehydrate quickly, poor feeding, fewer wet diapers, dry mouth, unusual sleepiness, fever, blood, or repeated vomiting requires prompt assessment.

\medterm{Diarrhea, Traveler's} Alternate index label; see \textit{Traveler's diarrhea}.

\medterm{Diarrhea Travelers (Traveler's Diarrhea)} Alternate index label for Traveler's Diarrhea.

\medterm{Diarrhoea} Alternate index label; see \textit{Diarrhea}.

\medterm{Diarthrosis} A freely movable joint in which the bones are separated by a space containing lubricating fluid. The knee, hip, shoulder, and elbow are examples.

\medterm{Diastase} An older name for enzymes that break starch into smaller sugars. Modern medical writing usually uses the term amylase for the related digestive enzyme activity.

\medterm{Diastasis} Separation or widening between structures that normally lie together. It commonly describes separation of the abdominal muscles after pregnancy and is not itself a true hernia, although both can occur together.

\medterm{Diastasis Recti} Widening and thinning of the tissue between the two front abdominal muscles, commonly during or after pregnancy or major weight change. It can cause a midline bulge when the abdomen is strained and is not a true hernia.

\medterm{Diastematomyelia} A birth defect in which the spinal cord is divided into two parts, often by bone or fibrous tissue. It may occur with abnormal vertebrae, a skin mark, weakness, altered sensation, or bladder problems.

\medterm{Diastole} The part of each heartbeat when the heart muscle relaxes and its chambers fill with blood. The lower blood-pressure number is the pressure in the arteries during this phase.

\medterm{Diastolic} Relating to diastole, when the heart relaxes and fills. Diastolic blood pressure is the lower number in a blood-pressure reading.

\medterm{Diastolic Blood Pressure} The pressure in the arteries while the heart relaxes between beats. It is the lower number in a blood-pressure reading and is interpreted with the upper number, symptoms, age, and other health conditions.

\medterm{Diastolic Dysfunction} Impaired relaxation or increased stiffness of a heart chamber, reducing its ability to fill normally. It can raise pressure in the lungs and cause breathlessness or heart failure, even when pumping strength appears preserved.

\medterm{Diastolic Heart Failure} An older term usually referring to heart failure caused by a stiff heart that does not fill easily, often with preserved pumping percentage. Diagnosis also requires symptoms, signs, and evidence of raised filling pressure.

\medterm{Diastolic Murmur} An extra or unusual heart sound heard while the heart relaxes between beats. It often reflects abnormal flow through a heart valve or nearby blood vessel and needs medical evaluation.

\medterm{Diastolic Murmurs} Extra heart sounds heard during the relaxation phase between beats. They commonly result from abnormal flow through a heart valve or major blood vessel and require clinical assessment.

\medterm{Diathermy} Use of high-frequency electrical or electromagnetic energy to produce heat in tissue. It may be used during surgery to cut tissue or stop bleeding, or in selected treatments to warm deeper tissues.

\medterm{Diathesis} A tendency to develop a particular disease, symptom, or body response. For example, a bleeding diathesis means a person is more likely to bleed than usual.

\medterm{Diatom Test} A forensic test that looks for microscopic algae in tissues after suspected drowning. It may support an investigation but cannot prove drowning by itself because algae can enter the body from other sources.

\medterm{Diatomaceous Earth} A powder made from the fossilized remains of microscopic algae. Inhalation of fine particles can irritate the lungs, and some forms containing crystalline silica can cause serious lung disease with repeated exposure.

\medterm{Diatrizoate Meglumine} An iodinated contrast substance used for selected X-ray or imaging studies. It improves visibility of body structures but may cause allergic-like reactions, kidney injury, aspiration problems, or thyroid effects.

\medterm{Diazepam} A benzodiazepine that calms brain activity and is used for selected anxiety, muscle-spasm, seizure, and alcohol-withdrawal conditions. It can cause drowsiness, falls, dependence, withdrawal, and dangerous breathing depression with opioids or alcohol.

\medterm{Diazepam Overdose} Taking too much diazepam can cause extreme sleepiness, confusion, poor coordination, low blood pressure, slowed breathing, coma, or death, especially when combined with alcohol, opioids, or other sedatives.

\medterm{Diazinon} An organophosphate insecticide that blocks acetylcholinesterase, causing excess nerve signals. Exposure can lead to sweating, salivation, vomiting, diarrhea, breathing difficulty, muscle weakness, seizures, or respiratory failure.

\medterm{Diazinon Poisoning} Poisoning by diazinon, causing excess cholinergic activity with sweating, salivation, vomiting, diarrhea, wheezing, slow heartbeat, muscle weakness, seizures, or breathing failure. It requires emergency treatment and decontamination.

\medterm{Diazonium Compounds} Chemicals containing a diazonium group, used in dyes, synthesis, and laboratory reactions. Some are unstable or toxic and may irritate tissues or release hazardous products.

\medterm{Diazoxide} A medicine that opens potassium channels and can reduce insulin release or relax blood vessels. It is used in selected cases of low blood sugar from excess insulin and severe high blood pressure; fluid retention and high glucose may occur.

\medterm{Dibekacin} An aminoglycoside antibiotic used in some regions for serious bacterial infections. It blocks bacterial protein production but can damage the kidneys, hearing, or balance, especially with prolonged or high-dose treatment.

\medterm{Dibenzepin} A tricyclic antidepressant used in some countries for depressive disorders. It can cause drowsiness, dry mouth, constipation, blurred vision, abnormal heart rhythms, and dangerous toxicity in overdose.

\medterm{Dibucaine} A long-acting local anesthetic applied to skin or mucous membranes to reduce pain or itching. Excessive use or absorption can cause numbness beyond the treatment area, seizures, heart problems, or other toxicity.

\medterm{Dicarboxylic Acid} An organic molecule containing two acid groups. These molecules take part in metabolism and may accumulate in some inherited or acquired metabolic disorders.

\medterm{Dichlorophen} A chemical formerly used in some disinfectants and medicines. Exposure can irritate the skin, eyes, or digestive tract; serious effects depend on the amount, route, and duration of exposure.

\medterm{Dichlorvos} An organophosphate insecticide that can block acetylcholinesterase. Poisoning may cause sweating, salivation, vomiting, diarrhea, muscle weakness, breathing difficulty, seizures, or respiratory failure.

\medterm{Dichorionic Diamniotic Twin Placenta} A twin pregnancy in which each fetus has its own outer membrane and amniotic sac, usually with a separate placenta. It generally has fewer shared-placenta complications, but routine maternal and fetal monitoring remains important.

\medterm{Dichorionic Twins} Twins that each have a separate outer membrane and amniotic sac. Most fraternal twins and some identical twins are dichorionic; risks still include premature birth, growth problems, and pregnancy complications.

\medterm{Diclofenac} A nonsteroidal anti-inflammatory medicine used for pain and inflammation. It can cause stomach irritation or bleeding, kidney or liver injury, fluid retention, and increased cardiovascular risk.

\medterm{Diclofenac Sodium} A salt form of diclofenac used for pain and inflammation. It can irritate or bleed the stomach and may affect the kidneys, blood pressure, liver, or heart.

\medterm{Diclofenac Sodium Overdose} Taking too much diclofenac sodium can cause nausea, vomiting, stomach bleeding, kidney injury, drowsiness, seizures, breathing problems, or abnormal body chemistry and requires urgent medical assessment.

\medterm{Dicloxacillin} A penicillin-type antibiotic that blocks bacterial cell-wall production. It is used for susceptible infections, especially some staphylococcal infections; allergy, diarrhea, and liver injury can occur.

\medterm{DICOM} An international format and communication standard for storing, sending, and viewing medical images and related information. It allows images from X-ray, CT, MRI, ultrasound, and other systems to be shared.

\medterm{Dicoumarol} An older oral anticoagulant that blocks vitamin K recycling and reduces production of several clotting factors. It has largely been replaced by medicines with more predictable effects and remains important mainly historically.

\medterm{Dicrotic} Describing a second small wave or notch in an artery's pressure pulse, usually produced when the aortic valve closes. It may be visible on a monitor or arterial pressure tracing.

\medterm{Dicrotophos} An organophosphate insecticide that can inhibit acetylcholinesterase. Poisoning may cause sweating, salivation, vomiting, diarrhea, muscle weakness, breathing difficulty, seizures, or respiratory failure.


\medterm{Dicyclomine} An anticholinergic medicine that relaxes intestinal muscle and may reduce cramping in selected bowel disorders. It can cause dry mouth, blurred vision, constipation, fast heartbeat, confusion, or difficulty urinating.

\medterm{Dieffenbachia Poisoning} Poisoning after chewing or swallowing parts of the Dieffenbachia plant. Needle-like calcium oxalate crystals can cause immediate burning, pain, drooling, and swelling of the mouth and throat; airway swelling requires emergency care.

\medterm{Diencephalon} A central part of the brain containing the thalamus and hypothalamus. It helps process sensory information and regulate hormones, temperature, appetite, sleep, movement, and other automatic functions.

\medterm{DIEP Flap} A reconstructive surgery using skin and fat from the lower abdomen, supplied by small blood vessels, to rebuild a breast or another body area while preserving most abdominal muscle.

\medterm{Diesel Oil} A petroleum fuel whose fumes or liquid can irritate the skin, eyes, and lungs. Swallowing or aspiration into the lungs can cause severe chemical pneumonia; suspected exposure requires prompt poison or medical advice.

\medterm{Diestrus} A phase of the reproductive cycle in many mammals between periods of sexual receptivity. It is characterized by activity and later regression of the corpus luteum.





\medterm{Diet} The foods and drinks a person regularly consumes. A health-supporting diet provides enough energy, protein, vitamins, minerals, fiber, and fluids and is adapted to health needs, culture, preferences, and available foods.

\medterm{Diet Diabetes (Diabetic Diet)} Alternate index label for Diabetic Diet.




\medterm{Diet Therapy} Use of an individualized eating plan to prevent, manage, or treat a health problem. It may address diabetes, kidney disease, food allergy, malnutrition, digestive disorders, or other needs and is tailored to the person's health and goals.

\medterm{Diet Ulcerative Colitis (Ulcerative Colitis Diet)} Alternate index label for Ulcerative Colitis Diet.



\medterm{Dietary Assessment} Systematic review of a person’s food, drink, nutrient, and supplement intake using methods such as a food diary, recall, or questionnaire. It helps identify deficiencies, excesses, patterns, and changes needed for health.

\medterm{Dietary Calcium} Calcium obtained from food or drinks, needed for bones, teeth, muscles, nerves, and blood clotting. Requirements vary with age and health; excess supplements may cause kidney stones or other problems.

\medterm{Dietary Carbohydrate} Carbohydrate in foods and drinks, including sugars, starches, and fiber. Digestible carbohydrate provides energy and affects blood glucose; fiber is only partly or not digested and supports bowel function.

\medterm{Dietary Fat} A nutrient made mainly of triglycerides and fatty acids that provides energy, supports cell membranes, and helps absorb vitamins A, D, E, and K. The type and amount of fat affect cardiovascular health.


\medterm{Dietary Fats} Fats obtained from food, including saturated, monounsaturated, polyunsaturated, and trans fats. They provide energy and support cell membranes and vitamin absorption; replacing saturated or trans fats with unsaturated fats is generally healthier.


\medterm{Dietary Fiber} Plant material that is not fully digested. It helps form stool, supports regular bowel movements, may improve cholesterol and glucose control, and can increase fullness.

\medterm{Dietary Fibre} Plant-derived carbohydrate that is not fully digested in the small intestine. Soluble and insoluble fiber support bowel function and may improve cholesterol, glucose control, and fullness; increase intake gradually with adequate fluids unless medically restricted.


\medterm{Dietary Iron} Iron from food or supplements, needed to make hemoglobin, muscle proteins, and enzymes. Too little can cause iron-deficiency anemia; too much can be toxic, especially in children or people with iron-storage disorders.

\medterm{Dietary Phospholipid} A food-derived fat containing phosphate that helps form cell membranes and transports fats in the blood. Its digestion and handling depend on food intake, absorption, and liver function.

\medterm{Dietary Phosphorus} Phosphorus from food, needed for bones, teeth, cell membranes, and energy use. People with reduced kidney function may need to limit highly absorbable sources to prevent high blood phosphorus.

\medterm{Dietary Protein} Protein from food that is digested into amino acids and small peptides. These support muscle and tissue repair, enzymes, hormones, immunity, and energy when needed.

\medterm{Dietary Reference Intakes} A set of nutrient-intake reference values used to plan and assess diets. They include the recommended allowance, adequate intake, estimated average requirement, and tolerable upper intake level, which vary with age, sex, and life stage.

\medterm{Dietary Sterol} A cholesterol-like substance found in foods, including cholesterol and plant sterols. Absorption and metabolism affect blood cholesterol; plant sterols can reduce intestinal cholesterol absorption in some people.

\medterm{Dietary Supplement} A product taken by mouth to provide vitamins, minerals, herbs, amino acids, or other dietary ingredients. Supplements may be useful for a documented need but can interact with medicines or cause toxicity.

\medterm{Dietary Supplementation} Deliberate addition of nutrients or other substances to the diet to correct a deficiency, meet increased needs, or support a health goal. The dose and need should be assessed because excess can be harmful.

\medterm{Dietary Supplements} Products containing vitamins, minerals, herbs, amino acids, or other substances intended to add to the diet. They do not replace balanced food and may interact with medicines or cause adverse effects.

\synonyms
Nutritional supplements

\medterm{Dietetics} The science and professional practice of using food and nutrition to promote health, prevent disease, and manage medical conditions through individualized dietary advice.

\medterm{Diethylcarbamazine} An antiparasitic medicine used mainly for lymphatic filariasis and some other filarial infections. Fever, rash, swelling, or inflammation may occur as the immune system reacts to dying parasites.

\medterm{Diethylhexylphthalate} Another name for DEHP, a chemical added to some plastics to make them flexible. Repeated or high exposure may affect reproduction, development, or hormone function.

\medterm{Diethylpropion} A stimulant appetite suppressant used short term in selected weight-management plans. It can cause insomnia, anxiety, fast heartbeat, high blood pressure, misuse, and dependence; medical supervision is required.
\medterm{Diethylstilbestrol} A synthetic estrogen formerly given to some pregnant people. Exposure before birth can increase the risk of reproductive-tract abnormalities, fertility problems, and certain rare cancers later in life.

\medterm{Dietitian} A nutrition professional who assesses food and nutrient needs and provides individualized advice for health, disease prevention, and medical conditions. Professional titles and licensing rules vary by region.

\medterm{Dieulafoy Lesion} An unusually large artery close to the surface of the stomach or another digestive organ that can suddenly cause severe bleeding, sometimes without a visible ulcer. It may produce vomiting of blood or black stools.

\medterm{Diff Dx} An abbreviation for differential diagnosis, the structured comparison of possible causes of a person’s symptoms, signs, or test results.

\medterm{Differences Of Sex Development} Variations present from birth in chromosomes, gonads, hormones, or internal or external genital anatomy. Medical evaluation explains the variation, addresses health needs, and supports informed decisions and well-being.

\medterm{Differential Blood Count} A blood test measuring the numbers or percentages of different white blood cells, including neutrophils, lymphocytes, monocytes, eosinophils, and basophils. Results help assess infection, inflammation, allergy, and blood disorders.

\medterm{Differential Diagnosis} The process of comparing possible diseases or conditions that could explain the same symptoms or findings. Clinicians consider serious causes first, use examination and targeted tests, and revise the possibilities as new information appears.

\medterm{Differentiated Cancer} A cancer whose cells still resemble the normal tissue from which they arose. The degree of resemblance, called differentiation, helps pathologists describe tumor grade and may provide information about behavior.

\medterm{Differentiation} The process by which cells develop specialized structures and functions. In cancer pathology, it describes how closely abnormal cells resemble the normal tissue from which they originated.

\medterm{Differentiation Antigens} Molecules on a cell’s surface that indicate its tissue type or stage of development. Laboratory tests for them help identify blood-cell disorders, classify cancers, and guide selected treatments.


\medterm{Difficile Clostridium (Clostridium Difficile Colitis)} Alternate index label for Clostridium Difficile Colitis.

\medterm{Difficult Airway} A situation in which a clinician may have trouble keeping a person breathing or placing a device into the airway. Causes include unusual anatomy, swelling, injury, illness, or limited neck movement and require careful planning.

\medterm{Difficult-to-Control Asthma} Asthma that remains poorly controlled despite appropriate treatment. Possible reasons include ongoing triggers, incorrect inhaler use, another illness, an incorrect diagnosis, severe airway inflammation, or a different response to medicines. Persistent symptoms need clinical review.

\medterm{Difficult Intubation} Difficulty placing a breathing tube in the windpipe despite appropriate technique and equipment. Limited mouth or neck movement, unusual anatomy, swelling, injury, obesity, or a previous difficult intubation may increase risk.

\medterm{Difficult Urination} Trouble starting or maintaining urine flow, a weak stream, pain, straining, or a feeling of incomplete emptying. Causes include infection, blockage, prostate or pelvic disorders, medicines, and nerve or bladder problems.

\medterm{Difficulty Speaking} Alternate index label; see \textit{Voice disorders}.

\medterm{Difficulty Swallowing} Trouble moving food or liquid safely from the mouth to the stomach. It may cause coughing, choking, pain, drooling, weight loss, or a feeling that food is stuck in the throat or chest.

\medterm{Difficulty Swallowing (Swallowing)} Alternate index label for Swallowing.

\medterm{Difficulty Walking} Trouble walking safely or normally because of pain, weakness, poor balance, numbness, joint disease, neurologic illness, or heart or lung limitations. Sudden onset requires urgent assessment.

\medterm{Diffuse} Spread broadly or throughout an area rather than limited to one small, clearly defined focus. The term may describe a disease, pain, swelling, imaging finding, or other process.

\medterm{Diffuse Astrocytoma} An older name for an infiltrating brain tumor made of astrocyte-like cells. Modern classification depends on molecular findings, including IDH status, which help determine diagnosis, outlook, and treatment.

\medterm{Diffuse Axonal Injury} Widespread damage to brain nerve fibers caused by rapid acceleration, deceleration, or rotation of the head. It can cause prolonged unconsciousness or severe brain dysfunction, and MRI may show the injury.

\medterm{Diffuse Esophageal Spasm} A disorder in which the esophagus contracts too early or unevenly, causing intermittent chest pain and difficulty swallowing liquids or solids. Testing helps distinguish it from blockage, reflux, and heart disease.

\medterm{Diffuse Idiopathic Skeletal Hyperostosis} A condition in which extra bone forms along spinal ligaments and where tendons attach to bone. It can cause stiffness, pain, swallowing difficulty, or fractures after minor injury.

\synonyms
DISH (Diffuse Idiopathic Skeletal Hyperostosis)

\medterm{Diffuse Large B-Cell Lymphoma} A fast-growing cancer of B lymphocytes, a type of white blood cell. It may cause painless swollen lymph nodes, fever, night sweats, or weight loss and usually requires prompt specialist treatment.

\medterm{Diffuse Midline Glioma} An aggressive brain or spinal-cord cancer arising in a midline structure, often in children. Symptoms depend on location and may include weakness, movement problems, vision changes, vomiting, or seizures.
\synonyms
DMG

\medterm{Diffuse Parenchymal Lung Diseases} Alternate index label; see \textit{Interstitial Lung Disease}.

\medterm{Diffuse Scleroderma} A form of systemic sclerosis in which skin thickening extends beyond the hands and face onto the trunk or limbs. It carries a higher risk of lung, kidney, heart, blood-vessel, and digestive complications.

\medterm{Diffusing Capacity} The ability of the lungs to move oxygen and other gases from the air sacs into the blood. It may be reduced by emphysema, interstitial lung disease, pulmonary vascular disease, or anemia.

\medterm{Diffusing Capacity Of The Lung} A measure of how efficiently gas crosses from the lung air sacs into nearby blood vessels. It is commonly measured as DLCO and helps evaluate lung tissue, blood vessels, and hemoglobin-related problems.

\medterm{Diffusion} Movement of molecules from an area where they are more concentrated to an area where they are less concentrated, caused by their random motion. It helps exchange gases and substances in the body.

\medterm{Diffusion Anisotropy} A difference in the direction water moves through tissue, especially along nerve fibers in the brain. MRI can use this pattern to study nerve-fiber structure and injury.

\medterm{Diffusion Capacity} A measure of how efficiently gases pass from lung air sacs into blood. It is commonly reported as DLCO and may be reduced in interstitial lung disease, emphysema, pulmonary vascular disease, or anemia.

\medterm{Diffusion Restriction} An area where water movement appears reduced on a special MRI sequence. It can occur with an acute stroke, abscess, densely packed tumor, seizure, or other problem and must be interpreted with the full examination.

\medterm{Diffusion-Weighted Imaging} An MRI technique that shows how water molecules move within tissue. Areas with restricted movement may occur in acute stroke, abscess, densely packed tumors, seizures, or other conditions.

\medterm{Diflunisal} A salicylate nonsteroidal anti-inflammatory medicine used for pain and inflammation. It can cause stomach irritation or bleeding, kidney injury, fluid retention, allergic reactions, or cardiovascular problems.

\medterm{Digastric Muscle} A two-part muscle beneath the jaw that helps open the mouth, lower the jaw, and raise the hyoid bone during swallowing and speech.

\medterm{DiGeorge Syndrome} A genetic disorder usually caused by loss of a small part of chromosome 22. It can affect the heart, immune system, calcium regulation, palate, hearing, learning, and development.

\synonyms
22q11.2 deletion syndrome

\medterm{DiGeorge Syndrome (22Q11.2 Deletion Syndrome)} Alternate name; see \textit{DiGeorge Syndrome}.

\medterm{Digestion} The process of breaking food into smaller substances that the body can absorb and use. It begins in the mouth and continues through the stomach and intestines with movement, acid, bile, and digestive enzymes.

\medterm{Digestive Diseases} Disorders affecting the digestive tract, liver, gallbladder, or pancreas. Examples include reflux, ulcers, infections, inflammatory bowel disease, blockage, malabsorption, and cancer.

\medterm{Digestive Disorder} A disease affecting the digestive tract or the process of digestion. It may cause abdominal pain, nausea, reflux, diarrhea, constipation, bleeding, vomiting, or poor nutrient absorption.

\medterm{Digestive Health} Healthy function of the digestive tract in breaking down food, absorbing nutrients and water, and removing waste. It is influenced by diet, fluids, medicines, movement, illness, and the gut lining and microbiome.

\synonyms
Gut Health

\medterm{Digestive System} The organs that process food, absorb nutrients and water, and remove solid waste. They include the mouth, esophagus, stomach, intestines, liver, gallbladder, pancreas, rectum, and anus.

\medterm{Digestive Tract} The continuous hollow passage through which food travels from the mouth to the anus. It includes the throat, esophagus, stomach, small intestine, large intestine, rectum, and anus.
\synonyms
Gastrointestinal tract; alimentary canal

\medterm{Digit} A finger or toe. Most fingers and toes contain three small bones, while the thumb and great toe usually contain two; muscles, tendons, nerves, and blood vessels allow movement and sensation.

\medterm{Digital Anorectal Exam - DARE} An examination in which a clinician uses a lubricated, gloved finger to feel the anal canal and rectum. It may assess tenderness, lumps, bleeding, muscle tone, and, when relevant, the prostate.

\medterm{Digital Clubbing} Enlargement and rounding of the fingertips with curved nails and loss of the normal angle at the nail base. It may be associated with lung, heart, digestive, liver, or other diseases and needs evaluation.

\medterm{Digital Fibrokeratoma} A rare, harmless, firm skin growth usually found on a finger or toe. It may resemble a wart; removal or microscopic examination may be needed when its appearance is unusual or it causes pain or interference.

\medterm{Digital Ischemia} Too little blood flow to a finger or toe, causing pain, pallor, coolness, numbness, or tissue death. It may result from arterial blockage, severe Raynaud phenomenon, inflammation, or an embolus and can be an emergency.

\medterm{Digital Mammography} Breast imaging that records X-ray pictures electronically. It helps detect masses, distortions, and calcium deposits and allows images to be enlarged, adjusted, stored, and compared.

\medterm{Digital Mucous Cyst} A usually harmless fluid-filled lump near the last joint of a finger, often linked to osteoarthritis. It may press on the nail, cause pain, or occasionally become infected.

\medterm{Digital Papillary Adenocarcinoma} A rare cancer of a sweat gland, usually on a finger, hand, or foot. It often appears as a slow-growing lump but can return locally or spread to other organs.

\medterm{Digital Rectal Exam} A physical examination using a lubricated, gloved finger to feel the anal canal and rectum and, when relevant, the prostate. It may detect tenderness, bleeding, lumps, hemorrhoids, or muscle problems.

\medterm{Digital Rectal Examination} An examination in which a clinician gently feels inside the rectum with a lubricated, gloved finger. It may check for pain, bleeding, lumps, hemorrhoids, prostate changes, or muscle problems after explanation and consent.

\medterm{Digital Subtraction Angiography} An X-ray study in which computer processing removes background structures from contrast images so blood vessels are clearer. It can diagnose narrowing or blockage and guide selected vascular procedures.

\medterm{Digitalis (Digoxin) Toxicity} Poisoning from digoxin or a related cardiac glycoside. Nausea, vomiting, confusion, visual changes, slow or fast abnormal rhythms, and dangerous potassium changes may occur; severe cases require emergency treatment.

\synonyms
Digoxin toxicity

\medterm{Digitalis Preparation} A preparation containing a cardiac glycoside such as digoxin or digitoxin. It can strengthen heart contraction and slow electrical conduction, but the dose range is narrow and toxicity can cause digestive, visual, or rhythm problems.

\medterm{Digitalis Toxicity} Harmful effects from digitalis or related cardiac glycosides, including nausea, vomiting, visual changes, confusion, slow or fast abnormal heart rhythms, and high potassium. Kidney disease and drug interactions increase risk.

\medterm{Digitate Warts (Verruca Digitata)} Finger-like or thread-like warts, usually on the face or neck, caused by human papillomavirus infection. They are generally harmless but can spread or recur.

\synonyms
Verruca Digitata

\medterm{Digitoxin} A cardiac glycoside that can strengthen heart contraction and slow electrical conduction. Its safe and toxic doses are close together, so abnormal rhythms, digestive symptoms, visual changes, and drug interactions are important concerns.

\medterm{Digoxin} A cardiac medicine used in selected people with heart failure or atrial fibrillation. It can strengthen contraction and slow the heart rate, but kidney function, electrolytes, other medicines, and blood levels affect toxicity risk.

\medterm{Digoxin Immune Fab} Antibody fragments given by injection that bind digoxin and related cardiac glycosides, reducing their effects. They are used for life-threatening abnormal rhythms, severe toxicity, or a potentially very large ingestion.

\medterm{Digoxin Test} A blood test measuring digoxin concentration to help guide dosing or assess possible toxicity. The result must be interpreted with timing, symptoms, kidney function, electrolytes, and interacting medicines; the number alone does not prove toxicity.

\medterm{Dihydrocodeine} An opioid medicine used for selected pain or cough conditions. It can cause constipation, nausea, drowsiness, dependence, dangerous breathing depression, and overdose, especially with alcohol or other sedatives.

\medterm{Dihydroergotamine} An ergot medicine used for selected acute migraine attacks. It narrows certain blood vessels and must be avoided in several vascular, heart, liver, kidney, and pregnancy-related conditions.

\medterm{Dihydrotestosterone} A potent androgen made from testosterone. It contributes to development of male external genitalia, facial and body hair, prostate growth, and other androgen effects.
\medterm{Dihydroxyacetone} A three-carbon sugar involved in metabolism. Its phosphate form is an intermediate used when cells break down glucose or make glucose.

\medterm{Dilantin Overdose} Taking too much phenytoin can cause involuntary eye movements, poor coordination, slurred speech, drowsiness, confusion, low blood pressure, abnormal heart rhythms, seizures, or coma.

\medterm{Dilatation} Widening or expansion of a hollow organ, blood vessel, duct, opening, or other body structure. It may be normal, caused by disease, or deliberately produced during a procedure.

\medterm{Dilatation And Curettage (D \& C)} A procedure that opens the cervix and removes tissue from inside the uterus using suction or an instrument. It may treat miscarriage or retained pregnancy tissue, investigate bleeding, or remove abnormal tissue; risks include bleeding, infection, and perforation.

\medterm{Dilate} To become wider or make a passage, blood vessel, opening, or body structure wider. Dilation may occur naturally, result from medicines or disease, or be created during a procedure.

\synonyms
Expand; enlarge

\medterm{Dilated Cardiomyopathy} A heart-muscle disease in which one or more chambers enlarge and pump weakly. Causes include inherited disease, infection, toxins, medicines, pregnancy, blocked arteries, and inflammation; it can lead to heart failure or abnormal rhythms.

\synonyms
Cardiomyopathy, Dilated

\medterm{Dilated Cardiomyopathy Prognosis} Outlook varies with heart-pumping strength, symptoms, cause, abnormal rhythms, other illnesses, and response to treatment. Medicines and devices can improve function for some people; severe disease may require advanced heart-failure care.

\medterm{Dilation} Widening or expansion of a hollow organ, blood vessel, duct, pupil, opening, or other body structure. It may be normal, caused by disease, or deliberately produced during a procedure.

\medterm{Dilation And Curettage} A procedure that opens the cervix and removes tissue from inside the uterus using suction or an instrument. It may treat miscarriage or retained pregnancy tissue, investigate abnormal bleeding, or remove abnormal uterine tissue.

\medterm{Dilation And Evacuation} A procedure that opens the cervix and removes pregnancy tissue from the uterus using suction and instruments. It is performed by trained clinicians for selected pregnancy-related indications, with risks including bleeding, infection, and uterine injury.

\medterm{Dilator} An instrument or device used to widen a body opening or passage gradually or temporarily. It may be used during examinations, surgery, childbirth care, dental treatment, or treatment of a narrowing.

\medterm{Diltiazem} A calcium-channel blocker used for high blood pressure, angina, and heart-rate control in some abnormal rhythms. It can slow the heart, lower blood pressure, cause constipation, or worsen heart failure in susceptible people.

\medterm{Dilute Russell Viper Venom Time} A blood-clotting test used mainly to detect lupus anticoagulant, an antibody associated with antiphospholipid syndrome. Anticoagulant medicines and other clotting problems can affect the result, so confirmatory testing is needed.

\medterm{Dilution} Lowering the concentration of a substance by adding solvent or another fluid. Accurate dilution is important when preparing medicines, laboratory reagents, standards, or specimens.

\medterm{Dimenhydrinate} An antihistamine used to prevent or treat motion sickness and nausea. It can cause drowsiness, dry mouth, blurred vision, constipation, confusion, or difficulty urinating.

\medterm{Dimenhydrinate Overdose} Taking too much dimenhydrinate can cause severe drowsiness or agitation, confusion, hallucinations, seizures, abnormal heart rhythms, breathing depression, or coma and requires urgent medical assessment.

\medterm{Dimercaprol} An injected medicine that binds certain toxic metals, including arsenic, mercury, and gold, so the body can remove them. It can cause high blood pressure, fast heartbeat, nausea, pain, and other adverse effects.

\medterm{Dimesna} A sulfur-containing compound studied as a protective or detoxifying medicine, including in some chemotherapy-related settings. Its clinical role depends on the specific treatment and available evidence.
\medterm{Dimethoate} An organophosphate insecticide that blocks acetylcholinesterase. Poisoning can cause sweating, excess saliva, vomiting, diarrhea, small pupils, muscle weakness, breathing problems, seizures, or coma.

\medterm{Dimethoxon} A poisonous breakdown product of some organophosphate pesticides. It can cause sweating, excess saliva, vomiting, diarrhea, small pupils, muscle twitching, weakness, breathing difficulty, or seizures and requires urgent poisoning advice.

\medterm{Dimethyl Fumarate} An oral medicine used for multiple sclerosis and selected skin conditions. It can cause flushing, digestive symptoms, low white-cell counts, liver injury, allergic reactions, or rare serious infections.

\medterm{Dimethyl Fumarate Allergy} An allergic or hypersensitivity reaction to dimethyl fumarate, a medicine used for some forms of multiple sclerosis. It may cause rash, itching, swelling, breathing difficulty, or other immune symptoms and requires medical assessment.

\medterm{Dimethyl Sulfoxide} An organic solvent used in laboratories and some specialized medical or investigational applications. Contact or ingestion can irritate tissues, and it may carry other substances through the skin.

\medterm{Dimethyl Sulphoxide} Another spelling of dimethyl sulfoxide, an organic solvent used in laboratories and selected medical or investigational applications. Exposure may cause irritation, a characteristic odor, or other adverse effects.


\medterm{Dimethylhydrazines} A group of highly toxic chemicals formerly used mainly in rocket fuel. Exposure can irritate or burn tissues and damage the nervous system, liver, blood, or other organs.

\medterm{Dimetilan} A poisonous organophosphate insecticide that disrupts nerve signals by blocking acetylcholinesterase. Exposure can cause sweating, vomiting, diarrhea, small pupils, weakness, breathing difficulty, seizures, or coma and needs urgent help.

\medterm{Diminished Ovarian Reserve} Fewer remaining eggs than expected for a person’s age, which may reduce response to fertility treatment and shorten the time available for conception. Tests estimate reserve but cannot precisely predict natural fertility.

\synonyms
Decreased Ovarian Reserve; DOR

\medterm{Dimocillin} A penicillin-type antibiotic used historically or in selected settings for susceptible bacterial infections. Choice depends on the organism, resistance testing, infection site, allergies, and local treatment guidance.

\medterm{Dimorphic} Existing in two distinct forms. Some fungi are dimorphic, growing as a mold in the environment and as a yeast in body tissues; this ability can contribute to infection.


\medterm{Dinophyceae} A group of aquatic microorganisms, including dinoflagellates. Some produce toxins during harmful algal blooms that can contaminate seafood or water and cause neurologic, digestive, or liver illness.

\medterm{Dinoprost} A prostaglandin medicine that stimulates uterine contraction and other smooth-muscle effects. It is used mainly in veterinary medicine and selected reproductive-health settings; pain, diarrhea, breathing difficulty, and blood-pressure changes may occur.

\medterm{Dinoprostone} A prostaglandin medicine used to soften the cervix and help start labor or manage selected pregnancy-related conditions. It can cause contractions that are too frequent, pain, fever, nausea, or fetal heart-rate changes.

\medterm{Dinoprostone (Cervidil, Prepidil)} A prostaglandin medicine used to soften the cervix and help start labor. It requires obstetric monitoring because it can cause overly frequent contractions, pain, fever, or changes in the fetal heart rate.

\synonyms
Cervidil, Prepidil

\medterm{Dinucleoside Phosphates} Molecules containing two nucleosides joined by phosphate groups. They occur in nucleotide metabolism and may take part in cellular signaling, energy transfer, or DNA and RNA-related reactions.

\medterm{Dinutuximab} An antibody-based cancer medicine that attaches to a marker found on many neuroblastoma cells and helps the immune system destroy them. It can cause severe nerve pain, allergic reactions, fever, and fluid leakage from blood vessels.

\medterm{Dioctyl Sodium Sulphosuccinate} A surfactant and stool-softening medicine also called docusate sodium. It helps water mix with stool and is used in selected constipation settings, although evidence of benefit is limited.

\medterm{Diopter} A unit of optical power equal to the reciprocal of focal length in meters. Eyeglass and contact-lens prescriptions use diopters to state the strength needed to focus light clearly.
\medterm{Dioxin} A group of long-lasting environmental chemicals produced by some industrial processes and burning. They can build up in body fat; long-term exposure may affect development, reproduction, immunity, hormones, and cancer risk.

\medterm{Dipeptide} A molecule made of two amino acids joined together. Dipeptides are produced during protein digestion and metabolism and can be absorbed or used in further chemical reactions.

\medterm{Dipeptidyl Peptidase-4 Inhibitors} Oral diabetes medicines that prolong the action of natural gut hormones. They help the body release insulin after meals and reduce glucagon, usually without causing weight gain or low blood glucose when used alone.


\medterm{Dipetalonema Infection} Infection with a Dipetalonema parasitic worm, usually acquired through an arthropod bite. It may affect skin or tissue beneath the skin; diagnosis and treatment depend on the species and clinical findings.

\medterm{Diphallia} A rare birth condition in which the penis is partly or completely duplicated. Abnormalities of the urinary tract, anus, spine, or other organs may also occur and require specialist assessment.

\medterm{Diphenhydramine Overdose} Taking too much diphenhydramine can cause severe drowsiness or agitation, confusion, hallucinations, seizures, abnormal heart rhythms, breathing depression, coma, or death and requires urgent medical assessment.

\medterm{Diphenoxylate} An opioid-related medicine that slows intestinal movement to reduce diarrhea, usually combined with atropine. Excessive use can cause drowsiness, constipation, dependence, breathing depression, or poisoning.

\medterm{Diphenylhydramine} An antihistamine with sedating and anticholinergic effects used in some allergy or sleep products. It can impair alertness and cause dry mouth, confusion, urinary retention, abnormal rhythms, or toxicity in overdose.

\medterm{Diphosphate} A chemical group containing two linked phosphate units. Diphosphate is involved in energy use, nucleotide formation, and other chemical reactions in cells.

\medterm{Diphosphonate} A chemical group found in compounds that can bind strongly to bone mineral. Some related medicines are used or studied for disorders of bone breakdown and calcium balance.

\medterm{Diphtheria} A serious infection caused by toxin-producing Corynebacterium bacteria. It can form a thick coating in the throat and cause airway blockage, heart inflammation, nerve damage, or death; vaccination prevents most cases.


\medterm{Diphtheria Toxin} A poisonous protein made by some Corynebacterium bacteria. It blocks protein production inside cells, causing local tissue injury and potentially severe heart, nerve, and airway complications.

\medterm{Diphtheroid Bacillus} A rod-shaped bacterium resembling Corynebacterium under the microscope. Many are normal skin bacteria or contaminants, but some can cause opportunistic infection in people with serious illness or medical devices.

\medterm{Diphyllobothriasis} An intestinal infection caused by a fish tapeworm acquired from raw or undercooked freshwater fish. It may cause abdominal symptoms and, rarely, vitamin B12 deficiency and anemia.

\medterm{Diphyllobothrium Latum} A fish tapeworm acquired by eating raw or undercooked freshwater fish. It lives in the intestine and may cause abdominal symptoms or compete for vitamin B12, occasionally leading to anemia.

\medterm{Diplacusis} Hearing one sound as two different pitches or qualities. It often reflects unequal hearing between the ears or inner-ear dysfunction and may follow infection, noise injury, or other ear disease.

\medterm{Diplegia} Severe weakness or paralysis affecting matching parts on both sides of the body. In spastic diplegia, the legs are affected more than the arms, often because of early brain injury.

\medterm{Diplococcus} A bacterium that appears as pairs of round cells under a microscope. The term describes shape, not one species; some diplococci cause infections in people.

\medterm{Diploid Cell} A cell containing two complete sets of chromosomes, one inherited from each parent. Most human body cells are diploid; eggs and sperm normally contain only one set.

\medterm{Diploidy} Having two complete sets of chromosomes, usually one from each parent. This is the normal chromosome state of most human body cells.

\medterm{Diplopia} Seeing two images of one object. It may result from eye misalignment, an eye-muscle or nerve problem, a brain disorder, or an eye-surface or lens problem; sudden onset needs urgent assessment.

\medterm{Diplopia (Double Vision)} Seeing two images of one object. It may occur with one eye or only when both eyes are open; causes include dry eye, cataract, eye misalignment, nerve disease, muscle disease, or brain disease.

\synonyms
Double Vision

\medterm{Dipole} A pair of equal and opposite electrical or magnetic charges separated by a distance. The concept is used in physiology, medical imaging, and physics to describe electrical or magnetic fields.

\medterm{Dipper} A person whose blood pressure falls during sleep, usually by about 10 percent or more compared with daytime readings. The pattern is measured with a portable 24-hour monitor and may help assess cardiovascular risk.
\medterm{Diprosopus} A rare birth defect in which part or all of the face is duplicated. It usually occurs with severe abnormalities of the brain, heart, or other organs and often causes major developmental problems.

\medterm{Dipsomania} Alternate historical term; see \textit{Alcohol Use Disorder}.

\medterm{Dipsosis} Excessive thirst or a persistent urge to drink. Causes include dehydration, high blood glucose, high blood concentration, medicines, dry mouth, or disorders affecting thirst regulation.

\medterm{Dipstick Device} A strip or small device containing chemical pads that change color when placed in a specimen, commonly urine. It provides a rapid screen for substances such as blood, protein, glucose, or infection markers.

\medterm{Diptera} The insect order containing flies and mosquitoes. Some species spread infections through bites or contamination, while others can cause infestations such as myiasis.

\medterm{Dipylidium Caninum} A tapeworm of dogs and cats that can occasionally infect people, especially children, after an infected flea is swallowed. Infection may cause mild digestive symptoms or visible worm segments in stool.

\medterm{Dipyridamole} A medicine that reduces platelet clumping and widens blood vessels. It is used in selected stroke-prevention plans and heart stress testing; headache, dizziness, low blood pressure, and bleeding may occur.

\medterm{Direct Action} An effect produced through an intervention’s immediate target, without requiring an intermediate pathway or mediator. The term may describe a medicine, hormone, nerve signal, or other biological process.

\medterm{Direct Antiglobulin Test} A blood test that detects antibodies or complement attached to a person’s red blood cells. A positive result may occur with autoimmune hemolysis, a transfusion reaction, or hemolysis affecting a newborn.

\medterm{Direct Coombs Test} Another name for the direct antiglobulin test, which detects antibodies or complement attached to red blood cells. It can help investigate autoimmune hemolysis, transfusion reactions, medicine-related hemolysis, or newborn hemolysis.

\medterm{Direct Hernia} An inguinal hernia in which abdominal tissue pushes through a weak area in the back wall of the groin canal. It may cause a groin lump or discomfort and can occasionally become trapped.


\medterm{Direct Oral Anticoagulants} Oral blood thinners that directly block a clotting protein, such as thrombin or factor Xa. They prevent and treat clots but can cause bleeding and require attention to kidney function, interactions, and procedures.

\medterm{Direct Transmission} Spread of an infectious organism through immediate contact or short-range respiratory droplets from an infected person or animal, without an intermediate object, food, water source, or vector.

\medterm{Direct Veneer} A tooth-colored composite layer shaped and hardened directly on a tooth to improve appearance or repair minor damage. It usually preserves more tooth than some alternatives but may stain, chip, or need replacement.


\synonyms
Direct composite veneer; composite dental veneer

\medterm{Directly Observed Therapy} A healthcare worker or trained supporter watches a person take each dose of medicine. It can support completion of complex treatment, especially for tuberculosis, and helps detect adverse effects or missed doses.

\medterm{Dirofilaria Repens} A mosquito-transmitted parasitic worm that rarely infects people. It usually causes a single lump beneath the skin or near the eye because humans are an accidental host; it does not usually mature or spread in the body.

\medterm{Dirt - Swallowing} Repeatedly eating soil or other nonfood substances, called geophagia and sometimes part of pica. It can expose a person to parasites, bacteria, toxic chemicals, or excess minerals and may cause constipation or blockage.

\medterm{Dirt-Eating} Repeatedly eating soil or other nonfood substances, a form of pica. It may be associated with iron deficiency, pregnancy, nutritional problems, or mental health conditions and can expose a person to infection or toxins.

\medterm{Dirty Necrosis} Dead tumor tissue mixed with broken cell material and many white blood cells. It is often seen in rapidly growing cancers and is interpreted by a pathologist with the rest of the tissue sample.

\medterm{Disability} A difficulty with body function, activity, or participation caused by a health condition and influenced by the person’s environment. Disability may be temporary or lifelong and varies in severity and support needs.

\medterm{Disability Evaluation} An assessment of how a health condition affects physical, mental, or daily functioning. History, examination, records, and appropriate tests may inform decisions about work, benefits, insurance, rehabilitation, or support.

\medterm{Disability-Adjusted Life Year} A population-health measure combining years lost from early death with years lived in less than full health. One DALY represents one lost year of healthy life under the method used.

\medterm{Disaccharidase} An intestinal enzyme that breaks a double sugar into simple sugars that can be absorbed. Lactase, for example, breaks lactose in milk into glucose and galactose.

\medterm{Disaccharide} A carbohydrate made of two linked simple sugars, such as lactose, sucrose, or maltose. It must be split by digestive enzymes before the intestine can absorb its component sugars.

\medterm{Disarticulation} Amputation performed through a joint rather than by cutting through the shaft of a bone. The level is chosen according to the injury, disease, circulation, function, and possibility of healing.

\medterm{Disc} A thin, round structure. In the spine, an intervertebral disc is a tough, flexible cushion between neighboring vertebrae that absorbs load and allows movement.

\medterm{Discectomy} Surgery to remove part of a damaged spinal disc that is pressing on a nerve. It may relieve leg or arm pain and weakness, but risks include infection, bleeding, nerve injury, and recurrent disc problems.


\medterm{Discitis} Inflammation or infection of a spinal disc, often involving the nearby vertebral bones. It may cause severe localized back or neck pain, stiffness, fever, or reduced movement and usually requires imaging and medical evaluation.


\medterm{Discography} An imaging procedure in which contrast is injected into a spinal disc to examine it and see whether the injection reproduces pain. Because it has limitations and risks, it is used selectively.

\medterm{Discoid} Coin-shaped or round. In medicine, it often describes a skin lesion or rash, such as the round plaques of discoid lupus; the associated condition must be identified separately.

\medterm{Discoid Lupus} A long-lasting form of skin lupus causing round, inflamed, scaly plaques. Lesions may scar, change skin color, or cause permanent hair loss on the scalp; some people later develop systemic lupus.

\medterm{Discoid Lupus Erythematosus} A chronic autoimmune skin disease causing persistent round, inflamed plaques that may scale, scar, or change pigment. It commonly affects the face, ears, and scalp and may cause permanent hair loss.
\medterm{Discoloration} An unusual change in the color of skin, mucous membranes, nails, or another tissue. Causes include bruising, low oxygen, jaundice, inflammation, infection, or altered pigment; sudden, painful, widespread, or unexplained change needs assessment.


\medterm{Discomfort} An unpleasant physical or emotional sensation that may be milder or less clearly located than pain. Its importance depends on the site, duration, associated symptoms, and underlying cause.

\medterm{Discordance} Lack of agreement between findings, measurements, genetic results, or the medical conditions of two related people. The significance depends on what is being compared and the clinical setting.

\medterm{Discordant} Not agreeing, matching, or corresponding. In medicine it may describe conflicting test results or different medical findings in related people, twins, or paired organs.

\medterm{Discordant Couple} A couple whose partners have different results or medical statuses relevant to their care, such as one partner having HIV and the other not. Counseling can reduce transmission risk and support reproductive decisions.


\medterm{Disease} An abnormal condition that affects the body or mind and changes normal structure or function. Diseases may be infectious, inflammatory, inherited, degenerative, metabolic, immune-related, or cancerous.


\medterm{Disease, Osgood-Schlatter} An overuse injury of the growing area where the knee tendon attaches to the shinbone. It causes pain and swelling below the knee, usually worse with running, jumping, or kneeling.

\medterm{Disease Progression} Change in a disease over time, often meaning worsening. It is assessed through symptoms, examination, laboratory results, imaging, function, complications, and response to treatment.

\medterm{Disease Reporting} Notification of specified diseases, infections, or health events to public-health authorities. Reports help detect outbreaks, protect contacts, arrange follow-up, and monitor health patterns; requirements vary by location.

\medterm{Disease Reservoir} A person, animal, environment, or material in which an infectious organism normally lives and from which it can spread to a susceptible host.

\medterm{Disease Response} The body’s biological or clinical reaction to a disease or its treatment. It is assessed using symptoms, signs, laboratory results, imaging, physical function, or survival.

\medterm{Disease, Subclinical} Disease that causes no obvious symptoms or signs but can be detected by laboratory tests, imaging, examination, or another evaluation.

\medterm{Disease Surveillance} Ongoing collection, analysis, interpretation, and reporting of health information to detect, monitor, and respond to disease patterns, outbreaks, and changes in population health.

\medterm{Disease Susceptibility} Increased likelihood of developing a disease because of genes, age, immune status, exposures, behavior, other illnesses, or environmental factors. Susceptibility raises risk but does not mean disease is certain.

\medterm{Disease Vector} A living carrier, commonly a mosquito, tick, flea, or other arthropod, that transfers an infectious organism between hosts. The vector may transmit infection without becoming ill itself.

\medterm{Disease-Free Survival} The length of time after treatment during which a person remains alive without detectable evidence of the disease. The exact definition and assessment method depend on the disease and study.

\medterm{Disfigurement} A visible change that substantially alters appearance, caused by injury, disease, surgery, a birth difference, or scarring. Its physical, emotional, and social effects vary between people.

\medterm{Disinfectant} A chemical or physical agent used on nonliving surfaces to destroy or inactivate many microorganisms. Disinfectants are not necessarily safe for skin, wounds, food, or internal use.

\medterm{Disinfection} Cleaning that destroys or inactivates most disease-causing microorganisms on nonliving surfaces. It may not destroy bacterial spores; the required level depends on the object, organism, and intended use.

\medterm{Disk} A flat, round structure. In the spine, an intervertebral disk lies between neighboring vertebrae, absorbs shock, and helps the back move.

\synonyms
Disc

\medterm{Disk Diffusion Susceptibility Test} A laboratory test of bacterial antibiotic sensitivity. Drug-containing disks are placed on growing bacteria, and the clear areas around them are measured and compared with standards to guide treatment.

\medterm{Disk Replacement - Lumbar Spine} Surgery replacing a damaged lower-back disk with an artificial disk in selected people. It aims to preserve movement and reduce pain, but risks include nerve injury, infection, persistent pain, and device problems.

\medterm{Diskectomy} Surgery to remove part or all of a damaged spinal disk pressing on a nerve or the spinal cord. It may relieve pain or weakness caused by disk herniation; risks include infection, bleeding, nerve injury, and recurrence.

\medterm{Diskitis} Inflammation or infection of a spinal disk and nearby vertebral bone. It commonly causes severe localized back pain, stiffness, fever, or reluctance to move and requires medical assessment and imaging.

\medterm{Dislocated Ankle (Ankle Dislocation)} Alternate index label for Ankle Dislocation.

\medterm{Dislocated Elbow} Loss of normal alignment of the elbow joint, usually after a fall or collision. It causes severe pain, deformity, and limited movement and needs urgent reduction with checks for fractures, nerve injury, and impaired circulation.

\medterm{Dislocated Hip} Displacement of the upper thighbone from its hip socket, usually after major trauma or around an artificial joint. It is an emergency because nearby nerves and the blood supply to the bone may be damaged.

\medterm{Dislocated Knee} Loss of alignment between the thighbone and shinbone, usually after severe injury. Major artery or nerve damage may be hidden, so urgent reduction and circulation checks are needed even if the knee returns to position.

\medterm{Dislocated Shoulder} Displacement of the upper arm bone from the shoulder socket. It causes pain, deformity, and inability to move normally and requires reduction with assessment of nerves, blood vessels, and fractures.


\medterm{Dislocation} Complete loss of normal contact between the surfaces of a joint, usually from injury or disease. It can cause severe pain and deformity and may occur with fractures or nerve and blood-vessel injury.

\medterm{Dislocation Of The Lens} Partial or complete movement of the eye’s natural lens away from its normal position because its supporting fibers are weak or torn. It may cause blurred or double vision, eye pain, or glaucoma.

\medterm{Disofenin Scan} A nuclear-imaging test that follows bile production and flow through the liver, gallbladder, and bile ducts. It can help detect blockage, leakage, inflammation, or selected causes of abdominal pain.

\medterm{Disopyramide} An antiarrhythmic medicine used selectively for certain abnormal heart rhythms. It can slow electrical conduction and weaken heart contraction and may cause dry mouth, urinary retention, low blood pressure, or new arrhythmias.

\medterm{Disorder} A disturbance of normal body or mind structure or function that causes symptoms, impairment, or important health risk. The term does not by itself indicate how severe the problem is.

\medterm{Disorder Antisocial Personality (Antisocial Personality Disorder)} Alternate index label for Antisocial Personality Disorder.

\medterm{Disorder Asperger (Asperger Syndrome)} Alternate index label for Asperger Syndrome.

\medterm{Disordered Thinking} A disturbance in the organization, clarity, logic, or flow of thoughts. It may occur with mental illness, brain disease, delirium, substance use, medicines, or other medical conditions.

\medterm{Disorders Sleep (Sleep)} Alternate index label for Sleep.



\medterm{Disorganized Schizophrenia} An older term for schizophrenia with markedly disorganized speech, behavior, or emotional expression. Modern classifications generally describe these features within schizophrenia-spectrum disorders rather than as a separate subtype.

\medterm{Disorientation} Difficulty knowing the time, place, situation, or one’s identity. It may occur with delirium, infection, abnormal body chemistry, intoxication or withdrawal, brain injury, or dementia; sudden onset needs urgent assessment.

\medterm{Dispensary} A place or service that stores, prepares, and supplies medicines and may provide basic medication advice. Safe dispensing includes checking the prescription, labeling the medicine, and explaining how to use it.

\medterm{Dispensed Amount} The quantity of medicine supplied to a patient. It is determined by the prescribed dose, treatment duration, formulation, refill rules, and applicable safety or legal requirements.

\medterm{Dispensing} The professional process of checking a prescription, selecting and preparing the medicine, labeling it, and supplying it with instructions for safe use.


\medterm{Displacement} Movement of an anatomical structure away from its normal position, as in a displaced fracture, joint injury, organ shift, or device migration. Effects depend on the structure involved and the amount of movement.

\medterm{Disruptive, Impulse-Control, And Conduct Disorders} A group of disorders involving persistent problems with emotional or behavioral self-control that violate the rights of others, major age-appropriate rules, or social norms. Examples include oppositional defiant disorder, conduct disorder, intermittent explosive disorder, kleptomania, and pyromania.

\medterm{Disruption} An interruption or breakdown of normal structure, function, continuity, or process. In medicine, it may describe damaged tissue, a broken barrier, disturbed development, or an interruption in the usual work of an organ or system.

\medterm{Dissect} To separate tissue layers or structures by cutting or careful blunt separation for surgery, examination, or anatomical study.

\medterm{Dissecting Aneurysm} Alternate index label; see \textit{Aortic dissection}.

\medterm{Dissection} Separation of tissue layers, either during surgery or because of disease. In an artery, a tear in the inner lining lets blood split the wall layers and can block blood flow or cause the artery to rupture.

\medterm{Dissection Aorta (Aortic Dissection)} Alternate index label for Aortic Dissection.

\medterm{Disseminate} To spread widely through the body, a tissue, or a population. Infection, cancer, inflammation, or another process may disseminate rather than remain in one limited area.

\medterm{Disseminated Coccidioidomycosis} A fungal infection in which Coccidioides spreads beyond the lungs to skin, bones, joints, brain coverings, or other organs. It can be severe and is more likely in people with impaired immunity.

\medterm{Disseminated Cutaneous Blastomycosis} A Blastomyces infection that spreads from the lungs or another site to multiple areas of skin. It may also involve bones or other organs and requires antifungal treatment.

\medterm{Disseminated Intravascular Coagulation} A life-threatening disorder in which clotting is activated throughout the bloodstream. Small clots can injure organs while using up platelets and clotting factors, causing serious bleeding at the same time.

\medterm{Disseminated Mycobacterium Avium–Intracellulare Infection} A widespread infection with Mycobacterium avium complex, usually in people with severe immune suppression. It may involve blood, bone marrow, lymph nodes, liver, spleen, or other organs.

\medterm{Disseminated Neuroblastoma} Neuroblastoma that has spread from its original site to distant organs, bone marrow, bones, liver, or lymph nodes. Treatment and outlook depend on age, spread, tumor biology, and response to specialist care.

\medterm{Disseminated Sclerosis} An older name for multiple sclerosis, an immune-mediated disease in which areas of damage develop in the brain, spinal cord, or optic nerves and cause symptoms that vary over time.

\medterm{Disseminated Tuberculosis} Tuberculosis that has spread through the bloodstream or lymphatic system to multiple organs. It may cause fever, weight loss, fatigue, organ damage, meningitis, or widespread small lesions and needs urgent treatment.

\medterm{Disseminated Varicella Infection} Chickenpox infection that spreads beyond the skin and may involve the lungs, brain, liver, or other organs. It is more likely and more severe in adults, pregnant people, and those with weakened immunity.

\medterm{Dissociation} A disruption in the usual connection between consciousness, memory, identity, emotions, perceptions, behavior, or sense of self. Mild forms may occur during daydreaming; severe or persistent forms can impair daily life.

\medterm{Dissociative Amnesia} Inability to recall important autobiographical information, often related to trauma or severe stress, beyond ordinary forgetfulness. Memory loss may involve a specific event, period, or broader part of a person’s life and is not better explained by ordinary forgetting, substances, a neurologic disorder, or another medical condition.

\medterm{Dissociative Disorder} Alternate index label for Dissociative Disorders.

\medterm{Dissociative Disorders} A group of mental disorders involving disruption in the usual integration of identity, memory, consciousness, perception, emotion, or behavior. Examples include dissociative amnesia, dissociative identity disorder, and depersonalization-derealization disorder. Temporary dissociative experiences can occur without constituting a disorder; a dissociative disorder causes significant distress or impairment.

\medterm{Dissociative Identity Disorder} A mental health disorder involving two or more distinct identity states and repeated gaps in memory for everyday events or personal information. It can cause changes in behavior, emotions, perception, and daily functioning.

\medterm{Distal} Farther from the point where a body part attaches or begins, or farther from the center of the body. For example, the hand is distal to the elbow.

\medterm{Distal Median Nerve Dysfunction} Abnormal function of the median nerve near the forearm, wrist, or hand, causing numbness, tingling, pain, or weakness in its area of supply. Compression at the wrist is a common cause.

\medterm{Distal Pancreatectomy} Surgery removing the body or tail of the pancreas, usually while leaving the head. It may be used for selected tumors, cysts, injury, or inflammation; diabetes and leakage from the pancreas are possible risks.

\medterm{Distal Radius Fractures} Breaks near the wrist end of the forearm bone, often after a fall. They may cause pain, swelling, deformity, or numbness; treatment depends on alignment, joint involvement, bone quality, and nerve or blood-vessel injury.

\medterm{Distal Renal Tubular Acidosis} A kidney disorder in which the kidneys cannot remove enough acid into the urine. Acid builds up in the blood and may cause low potassium, kidney stones, weak bones, muscle weakness, or poor growth in children.

\medterm{Distal Splenorenal Shunt} Surgery connecting the splenic vein to the left kidney vein to lower pressure in veins around the esophagus. It may reduce bleeding from varices while preserving some blood flow through the liver.

\medterm{Distal Tubule} A small kidney tube that adjusts the amount of salt, calcium, and water left in forming urine. Its activity helps regulate blood pressure and body chemistry and can be affected by some diuretic medicines.

\medterm{Distamycins} A group of antibiotic-like compounds that bind to DNA and have been studied for antimicrobial or anticancer effects. Their clinical use is limited and depends on the specific compound and evidence.

\medterm{Distant} Far from a reference point or original site. In cancer, distant disease means that cancer has spread to remote organs or lymph nodes rather than remaining near the primary tumor.


\medterm{Distended} Abnormally stretched, swollen, or enlarged by gas, fluid, blood, tissue, or pressure. The term may describe the abdomen, bladder, stomach, bowel, veins, or another body structure.

\medterm{Distension} Enlargement or stretching of a body part or hollow organ, often from gas, fluid, urine, stool, blood, pressure, or blockage. Abdominal distension with severe pain, vomiting, fever, or inability to pass stool or gas needs prompt assessment.

\medterm{Distention} Enlargement or stretching of an organ, cavity, or body part, such as abdominal distention from gas, fluid, swelling, or blockage.

\medterm{Distichia} Growth of an extra row of eyelashes, usually from openings of glands along the eyelid. Misguided lashes may rub the eye and cause irritation, tearing, corneal scratches, or ulcers.


\medterm{Distortion} An abnormal change in the shape, structure, appearance, or representation of tissue or an image. New breast, facial, or visual distortion may indicate disease and should be assessed with appropriate examination or testing.


\medterm{Distraction Osteogenesis} A surgical method for growing new bone by slowly separating two bone sections after a controlled cut. A device gradually widens the gap to lengthen or reshape the jaw, limbs, or other bones.

\medterm{Distress} Significant physical or emotional suffering, discomfort, or strain related to illness, injury, treatment, or difficult circumstances. Severe distress may affect sleep, functioning, decisions, or safety and may need support.

\medterm{Distributive Shock} Life-threatening shock caused by widespread blood-vessel dilation or abnormal blood-flow distribution. Effective circulation falls despite variable heart output; causes include sepsis, anaphylaxis, nerve injury, and some poisons.

\medterm{District Nurse} A community nurse who provides care in homes, clinics, or other local settings. Duties may include wound care, chronic-disease support, health education, medicine monitoring, prevention, and coordination with other services.

\medterm{Disulfiram-Like Reaction} An unpleasant reaction after alcohol is taken with disulfiram or another interacting substance. Flushing, headache, nausea, vomiting, palpitations, chest discomfort, or low blood pressure may occur.

\medterm{Disulfoton} A highly toxic organophosphate insecticide that inhibits acetylcholinesterase. Exposure can cause sweating, excess saliva, vomiting, diarrhea, muscle weakness, breathing failure, seizures, or coma.


\medterm{Dithiolethione} A sulfur-containing compound studied for activating cellular defense and detoxification pathways. It may influence inflammation and protection from chemical injury, but clinical uses depend on evidence for the specific compound.

\medterm{Ditiocarb} A term for dithiocarbamate compounds, some of which are used as pesticides or metal-binding chemicals. Toxic effects vary by compound and may involve the nervous system, liver, kidneys, or skin.

\medterm{Diuresis} Increased production of urine. It may result from excess fluid, high blood glucose, diabetes insipidus, medicines called diuretics, or other causes; excessive loss can lead to dehydration and electrolyte problems.

\medterm{Diuretic Drug} A medicine that makes the kidneys remove more sodium and water in urine. It is used for high blood pressure, heart failure, and swelling but can cause dehydration, low blood pressure, electrolyte changes, or kidney problems.

\medterm{Diurnal} Relating to daytime or occurring during a daily 24-hour cycle. The term may describe symptoms, hormone levels, body functions, or patterns of activity and sleep.

\medterm{Diverticula (Diverticulosis)} Alternate index label; see \textit{Diverticulosis}.

\medterm{Diverticular Disease} A condition involving small pouches in the wall of the colon. Diverticulosis means the pouches are present, while diverticular disease generally refers to symptoms or complications such as abdominal pain, altered bowel habits, bleeding, inflammation, infection, blockage, or perforation. Diverticulitis is inflammation or infection of one or more pouches and commonly causes lower-abdominal pain and fever.

\medterm{Diverticulitis} Inflammation of one or more small pouches in the wall of the colon. It commonly causes lower abdominal pain, fever, and changes in bowel habits and may lead to an abscess, perforation, or intestinal blockage.



\medterm{Diverticulosis} The presence of small pouches in the wall of the colon. It often causes no symptoms but can be associated with abdominal symptoms, bleeding, inflammation, or other complications.

\synonyms
Colonic diverticula

\medterm{Diverticulosis And Diverticulitis} Diverticulosis means small pouches have formed in the colon wall. Diverticulitis occurs when one or more pouches become inflamed or infected, causing abdominal pain, fever, and bowel changes.


\medterm{Diverticulum} A pouch-like projection from the wall of a hollow organ, especially the intestine. It may cause no symptoms or become inflamed, infected, blocked, bleed, or rarely perforate.

\medterm{Divisum Pancreas (Pancreas Divisum)} Alternate index label for Pancreas Divisum.

\medterm{Dix-Hallpike Maneuver} A bedside test for a common inner-ear cause of brief positional vertigo. The clinician moves the head and body into a specific position and watches for dizziness and characteristic eye movements.

\medterm{Dizygotic Twin} One of a pair of fraternal twins formed when two separate eggs are fertilized by two separate sperm. The twins share about half their genetic variation on average, like ordinary siblings.

\medterm{Dizygotic Twins} Twins formed when two separate eggs are fertilized by two separate sperm. They are genetically like ordinary siblings and may have separate or shared-appearing placentas; the pregnancy still requires appropriate monitoring.

\medterm{Dizziness} A sensation of spinning, lightheadedness, imbalance, faintness, or being unsteady. Causes include inner-ear, heart, blood-pressure, vision, medicine, blood, metabolic, or neurologic problems.

\medterm{Dizziness (Dizzy)} Alternate index label for Dizzy.


\medterm{Dizziness: When To See A Doctor?} Seek urgent care for sudden dizziness with weakness, numbness, trouble speaking, severe headache, chest pain, fainting, inability to walk, persistent vomiting, or new hearing loss. Recurrent or unexplained dizziness also needs medical evaluation.

\medterm{Déjà Vu} A feeling that a current situation has happened before, even though it is probably new. Occasional episodes are common; frequent episodes with impaired awareness or unusual movements may occur with seizures.

\medterm{DKD} Alternate index label; see \textit{Diabetic Nephropathy}.

\medterm{DLI} Donor lymphocyte infusion, in which immune cells from a stem-cell donor are given after transplantation to strengthen an immune attack against returning cancer. It can cause graft-versus-host disease and other complications.

\medterm{DM} An abbreviation most commonly meaning diabetes mellitus, a group of disorders with persistently high blood glucose. Because DM can have other meanings, the surrounding clinical context is important.

\medterm{DM2} An abbreviation for type 2 diabetes mellitus, in which the body resists insulin and may gradually produce too little of it. It can cause long-term damage to the eyes, kidneys, nerves, heart, and blood vessels.

\medterm{DMFT Index} A dental measure counting permanent teeth that are decayed, missing because of decay, or filled. It estimates a person's lifetime experience of tooth decay but does not show whether a current problem is active.

\medterm{DMG} Alternate index label; see \textit{Diffuse midline glioma}.

\medterm{DMSA Scan} A nuclear-imaging test using a small amount of labeled dimercaptosuccinic acid to show kidney tissue. It can identify differences in kidney function, scars, or injury, especially in children with repeated urinary infections.

\medterm{DNA} Deoxyribonucleic acid, the molecule that stores hereditary information in humans and most organisms. Its sequence helps control cell function and protein production; changes may be inherited or acquired during life.

\medterm{DNA Adduct} A chemical group attached to DNA after exposure to a reactive substance. If not repaired, it can alter DNA copying, cause mutations, and potentially contribute to cancer.

\medterm{DNA Alkylation} Attachment of an alkyl group to DNA. This can change base pairing, block copying or repair, cause mutations, or kill a cell; some anticancer medicines work partly this way.

\medterm{DNA Amplification} Production of many copies of a selected DNA sequence, naturally in cells or in the laboratory by methods such as polymerase chain reaction. It helps detect, identify, or study genes and microorganisms.

\medterm{DNA Binding} Attachment of a protein or other molecule to a DNA sequence. This can control gene activity, DNA copying, repair, chromosome organization, or the action of some medicines.

\medterm{DNA Damage} A chemical or structural change in DNA that can interfere with gene function or copying. Cells usually repair such damage, but persistent or incorrectly repaired damage can cause mutations, cell death, or cancer.
\medterm{DNA Demethylation} Removal or loss of small chemical groups called methyl groups from DNA. It can change how accessible a gene is and alter gene activity during development, cell reprogramming, or disease.

\medterm{DNA Double Helix} The twisted ladder-like structure of DNA, formed by two strands running in opposite directions. Sugar and phosphate form the outside rails, while paired bases join the strands inside.



\medterm{DNA Fragmentation} Breaking DNA strands into smaller pieces. It occurs normally during programmed cell death and can also result from injury, radiation, chemicals, disease, or laboratory processing.

\medterm{DNA Gyrase} A bacterial enzyme that twists and untwists DNA so it can be copied and separated. Fluoroquinolone antibiotics block this enzyme, stopping bacterial growth and survival.

\medterm{DNA Integration} Insertion of foreign or copied DNA into a cell’s genome. It occurs in some viral infections, evolution, and gene therapy and can alter nearby gene activity if the insertion affects a control region.

\medterm{DNA Ligases} Enzymes that join breaks in DNA strands during copying, repair, and genetic recombination. They help preserve the continuous structure of chromosomes.

\medterm{DNA Ligation} Joining DNA fragments by forming a chemical bond between them. Cells use it during DNA copying and repair, and laboratories use it in genetic testing and cloning.

\medterm{DNA Methylation} Addition of small chemical groups called methyl groups to DNA, usually at cytosine bases. It can change how strongly a gene is expressed without changing the DNA sequence.

\medterm{DNA Packaging} Arrangement of long DNA molecules around histone proteins and into compact structures called chromosomes. Packaging allows DNA to fit inside the nucleus and helps control access to genes.

\medterm{DNA Primase} An enzyme that makes short RNA starting pieces so other enzymes can begin copying DNA. It is required during cell division and DNA repair.

\medterm{DNA Primers} Short nucleic-acid sequences that provide a starting point for copying DNA. Cells make them during replication, and laboratories use designed primers to amplify selected genetic sequences.

\medterm{DNA Probe} A labeled single-stranded DNA sequence designed to bind a matching target. It can help detect a gene, mutation, chromosome region, or microorganism in a laboratory sample.

\medterm{DNA Profiling} Comparison of variable DNA markers from biological samples to assess identity or biological relationships. Results are statistical and require validated methods, suitable controls, and reliable sample handling.

\medterm{DNA Repair} Cellular processes that detect and correct damage or copying errors in DNA. Repair pathways address small base changes, bulky chemical damage, mismatched bases, and breaks in one or both DNA strands.

\medterm{DNA Repair Mechanisms} Cellular systems that correct DNA damage caused by copying errors, chemicals, radiation, or reactive oxygen. Failure of some pathways increases mutation risk and can contribute to inherited cancer syndromes.

\medterm{DNA Replication} The process by which DNA copies itself before cell division. Each new DNA molecule contains one original strand and one newly made strand, allowing genetic information to pass to daughter cells.

\medterm{DNA Sequencing} Determination of the order of DNA building blocks in a sample. It can examine a single gene, selected genes, all coding regions, or the whole genome to identify genetic changes.

\medterm{DNA Virus} A virus whose genetic material is DNA. DNA viruses vary widely in transmission, tissues affected, and disease; some cause persistent infection, while others cause short-term illness.

\medterm{D.O.} Doctor of Osteopathic Medicine, a physician degree whose graduates are licensed to practice medicine and surgery like MD graduates in the United States.

\medterm{Do Not Resuscitate} A medical order stating that cardiopulmonary resuscitation should not be attempted if a person’s breathing or heartbeat stops, according to the person’s informed wishes or legally authorized decision-maker.

\medterm{DOACs} Direct oral anticoagulants are oral blood thinners that directly block thrombin or factor Xa. They prevent and treat blood clots but can cause bleeding and require attention to kidney function, interactions, and planned procedures.

\medterm{DOB} An abbreviation for date of birth, the day, month, and year a person was born. It helps confirm identity and calculate age for records, screening, medicine doses, and treatment decisions.

\medterm{Dobutamine} An intravenous medicine that increases heart contraction and, sometimes, heart rate. It is used short term for selected cases of severe low-output heart failure or shock and requires close monitoring.

\medterm{Dobutamine Stress Echocardiography} An ultrasound heart test using dobutamine to increase heart workload when a person cannot exercise adequately. New areas of poor wall movement may suggest reduced blood flow, but results require clinical interpretation.

\medterm{Doc} Informal shorthand for doctor, a person qualified to practice medicine or another health profession. The exact role depends on the person’s training and professional license.

\medterm{Docetaxel} An anticancer medicine that interferes with microtubules, structures needed for cell division. It is used for several cancers and can cause low blood counts, infection, hair loss, nerve problems, or allergic reactions.

\medterm{Doconexent} A pharmaceutical name for docosahexaenoic acid, an omega-3 fatty acid found in oily fish, algae, and cell membranes. Its effects depend on the product, dose, and reason for use.

\medterm{Docosadienoic Acid} A 22-carbon fatty acid with two double bonds, present in small amounts in body and dietary lipids. It is studied mainly for possible effects on cell membranes and metabolism.

\medterm{Docosahexaenoic Acid} An omega-3 fatty acid that is an important structural part of brain and retinal cell membranes. It comes mainly from oily fish, algae, or supplements and is especially important during brain development.

\medterm{Docosapentaenoic Acid} A 22-carbon polyunsaturated fatty acid with five double bonds. It occurs as omega-3 and omega-6 forms in foods and tissues; its health effects depend on the form and amount.

\medterm{Docosatrienoic Acid} A 22-carbon fatty acid with three double bonds, found in small amounts in body and dietary lipids. It is studied mainly for its possible roles in cell membranes and metabolism.

\medterm{Doctor} A person who has completed appropriate professional training and is licensed or authorized to practice medicine. The title may also be used by other professionals with doctoral degrees, depending on local rules.



\medterm{Docusate} A stool-softening medicine that helps water mix with stool. It may be used when avoiding straining is important, but evidence for treating ordinary constipation is limited; diarrhea and abdominal cramps can occur.


\medterm{Dog Allergy} An immune reaction to proteins from a dog’s skin flakes, saliva, urine, or other secretions. It can cause sneezing, blocked or itchy nose, watery eyes, hives, cough, or asthma symptoms.

\medterm{Dog Bite} An injury caused by a dog’s teeth. Wash the wound promptly and seek assessment for tissue damage, infection, tetanus protection, and rabies risk, especially for deep bites or bites to the face or hand.

\synonyms
Dog bite injury

\medterm{Dog Bite Treatment} Care includes thorough wound cleaning, assessment of tendon, nerve, or bone injury, tetanus review, rabies-risk evaluation, and antibiotics when indicated. Deep bites, severe bleeding, infection, or bites to the hand or face need prompt care.

\synonyms
Bite Dog (Dog Bite Treatment)

\medterm{Dolasetron} A medicine that blocks serotonin 5-HT3 receptors and helps prevent nausea and vomiting in selected settings. It can cause headache, constipation, or dangerous heart-rhythm changes.

\medterm{Dolasetron Mesylate} A salt form of dolasetron, a medicine that blocks serotonin 5-HT3 receptors to prevent nausea and vomiting in selected settings. It can cause constipation, headache, or dangerous heart-rhythm changes.

\medterm{Dolichocephalic} Having a relatively long, narrow skull. It may be a normal family or positional variation, but marked or worsening skull narrowing in an infant may indicate early fusion of a skull suture and needs assessment.

\medterm{Dolichol} A long-chain fatty alcohol used inside cells to assemble sugar chains attached to proteins. This process is important for protein folding, transport, and function.

\medterm{Dolor} The Latin word for pain, one of the classic signs of inflammation. Pain may result from irritated nerves, swelling, injury, infection, or another cause and is interpreted with other findings.

\medterm{Domain} A defined region or area. In biology, a protein domain is a part with a particular structure or function; in genetics, a domain may be a genomic or regulatory region.

\medterm{Domestic Violence} A pattern of physical, sexual, emotional, psychological, or financial abuse or coercive control by an intimate partner or household member. It can cause injury, fear, isolation, and serious health effects; confidential support is available.

\medterm{Dominant} Describing a genetic variant whose effect can appear when only one copy is present. An affected parent with an autosomal dominant condition may pass the variant to each child with a one-in-two chance.

\medterm{Domperidone} A medicine that blocks dopamine receptors outside the brain and is used in some countries for nausea or slow stomach movement. It can prolong the heart rhythm and cause dangerous arrhythmias.

\medterm{Donath-Landsteiner Test} A blood test that detects an antibody associated with paroxysmal cold hemoglobinuria. The antibody binds red cells in the cold and destroys them when the blood warms.

\medterm{Donepezil} A medicine that increases acetylcholine activity in the brain and may temporarily improve symptoms of Alzheimer disease. It does not cure the disease and can cause nausea, diarrhea, slow heartbeat, fainting, or sleep problems.

\medterm{Donnan Equilibrium} The distribution of charged particles across a partly permeable membrane when large charged particles cannot cross. It helps explain cell volume, membrane electrical charge, and movement of salts and water.

\medterm{Donohue Syndrome} A rare inherited disorder in which the body cannot respond normally to insulin. It causes severe growth and feeding problems, very little body fat, unusual facial features, and serious metabolic complications beginning in infancy.

\medterm{Donor} A person who provides blood, cells, tissue, an organ, sperm, eggs, or another biological material for another person or for research. Screening, informed consent, compatibility, and traceability support safety.

\medterm{Donor Insemination} A fertility treatment in which sperm from a screened donor is placed in the reproductive tract to attempt pregnancy. The donor may be known or anonymous, and counseling and legal requirements vary by location.

\medterm{Donor Suitability Assessment} Evaluation of a potential donor’s medical history, examination, laboratory tests, infection risks, organ function, and compatibility to determine whether donation is safe and suitable for transplantation.

\medterm{Donor Testing} Laboratory and clinical evaluation of a potential donor for infections, cancer, organ function, blood type, and other factors that affect donation safety and recipient compatibility.
\medterm{Do-Not-Resuscitate Order} A medical order instructing clinicians not to perform cardiopulmonary resuscitation if a person’s breathing or heartbeat stops. It does not automatically refuse other treatments or comfort care.

\synonyms
DNR order

\medterm{Donovanosis} A sexually transmitted infection caused by Klebsiella granulomatis. It produces slowly enlarging, usually painless genital ulcers that may bleed easily and requires medical diagnosis and antibiotic treatment.

\synonyms
granuloma inguinale.

\medterm{Donovanosis (Granuloma Inguinale)} A sexually transmitted bacterial infection causing slowly progressive, usually painless genital ulcers that bleed easily. Diagnosis and antibiotic treatment require clinical evaluation, and sexual partners may need assessment.

\medterm{Dopa} A chemical precursor that the body can convert into dopamine. Levodopa, a form that reaches the brain, is used with another medicine to improve movement problems in Parkinson disease.

\medterm{Dopamine} A chemical messenger used by nerve cells and some hormone-producing cells. It helps control movement, motivation, reward, attention, learning, blood-vessel tone, and release of certain hormones.

\medterm{Dopamine Agonist} A medicine that stimulates dopamine receptors and mimics some effects of dopamine. It is used for selected movement, hormone, or other disorders and may cause nausea, sleepiness, low blood pressure, hallucinations, or impulse-control problems.
\medterm{Dopamine Agonists} Medicines that stimulate dopamine receptors and are used for selected Parkinson, restless-legs, or hormone disorders. Possible effects include nausea, sleepiness, low blood pressure, hallucinations, and impulse-control problems.

\medterm{Dopamine Antagonist} A medicine or compound that blocks dopamine receptors. It is used in selected psychiatric, nausea, or digestive treatments but can cause movement problems, drowsiness, hormone changes, or other adverse effects.

\medterm{Dopamine Receptor} A protein on or inside a cell that responds to dopamine. Different receptor types affect movement, mood, hormone release, blood vessels, and other functions and are targets of some medicines.

\medterm{Dopamine Transporter} A protein that clears dopamine from the space between nerve cells by moving it back into the releasing nerve cell. Stimulants and some imaging tracers affect or measure its activity.

\medterm{Dopaminergic Neuron} A nerve cell that makes and releases dopamine. These cells help control movement, motivation, reward, attention, and hormone regulation; loss of some contributes to Parkinson disease.

\medterm{Dopa-Responsive Dystonia} An inherited movement disorder causing sustained muscle contractions and abnormal postures, sometimes with Parkinson-like symptoms. It often improves markedly with a small dose of levodopa.

\medterm{Doppler Echocardiography} A heart ultrasound technique that measures the direction and speed of blood flow through the heart and major vessels. It helps assess valves, chambers, pressure, and abnormal flow.

\medterm{Doppler Ultrasound} An ultrasound technique that measures moving blood to assess its direction and speed. It can identify narrowed or blocked vessels, abnormal flow, blood clots, and circulation problems without ionizing radiation.

\medterm{Doppler Ultrasound Exam Of An Arm Or Leg} An ultrasound examination measuring blood flow through the arteries and veins of an arm or leg. It can help identify narrowing, blockage, a clot, aneurysm, or other circulation problem without ionizing radiation.

\medterm{Doppler Velocimetry} Ultrasound measurement of blood-flow speed and resistance in fetal or placental vessels. It helps assess circulation and fetal well-being when poor growth, placental problems, or other pregnancy risks are suspected.

\medterm{Dor Procedure} A heart operation that removes or excludes a scarred, bulging part of the left ventricle after a large heart attack. It reshapes the chamber in selected people with severe ischemic heart-muscle disease.

\synonyms
Endoventricular Circular Patch Plasty (EVCPP); Surgical Ventricular Restoration

\medterm{Doravirine} An HIV medicine that blocks reverse transcriptase, an enzyme HIV needs to copy its genetic material. It is used with other medicines for HIV-1 infection; interactions and resistance affect its choice.

\medterm{Dorea} A group of bacteria commonly found in the human intestine and able to grow without oxygen. Finding it usually reflects normal gut organisms unless it comes from a normally sterile site with compatible illness.

\medterm{Doripenem} An intravenous antibiotic that damages bacterial cell walls and is used for selected serious infections. It can cause allergy, diarrhea, seizures, or resistance and should be chosen according to the organism and patient.

\medterm{Dormant} Inactive or not multiplying. The term may describe a latent infection, quiet tumor cell, inactive disease process, or other biological state that can become active under certain conditions.

\medterm{Dorsal} Relating to the back or upper surface of a body part. In human anatomy it usually means toward the back, although in the hand or foot it refers to the opposite surface from the palm or sole.
\medterm{Dorsal Blocking Splint} A hand splint that limits wrist and finger extension to protect healing tendons, commonly after flexor-tendon repair. Position and duration are set by the treating surgical and therapy team.

\medterm{Dorsal Funiculus} A back column of the spinal cord containing nerve fibers that carry vibration, fine touch, and awareness of body position to the brain.

\medterm{Dorsal Kyphosis} Excessive forward curvature of the upper spine, causing a rounded upper back. It may result from posture, spinal wear, osteoporosis, or vertebral fractures and can cause pain, height loss, or breathing difficulty.
\medterm{Dorsal Root Ganglion} A swelling on the back root of a spinal nerve containing the cell bodies of sensory nerve cells. It relays information about touch, pain, temperature, and body position to the spinal cord.

\medterm{Dorsalis Pedis Artery} An artery running across the top of the foot as a continuation of the anterior tibial artery. Its pulse helps assess blood flow to the foot.

\medterm{Dorsiflexion} Upward movement of the foot at the ankle, bringing the toes toward the shin. It helps clear the foot during walking; weakness can cause foot drop and tripping.

\medterm{Dorsolateral} Describing a position toward the back and side of a body structure. For example, a dorsolateral area is behind and to one side rather than toward the front or middle.

\medterm{Dorsum} The back or upper surface of a body part, such as the back of the hand, top of the foot, or upper surface of the tongue.

\medterm{Dorzolamide} An eye medicine that lowers pressure inside the eye by reducing fluid production. It is used for glaucoma or raised eye pressure and may cause stinging, blurred vision, or a bitter taste.

\medterm{Dos And Donts During First Trimester Of Pregnancy} Early pregnancy care includes prenatal visits, folic acid, avoiding alcohol, tobacco, and non-prescribed drugs, reviewing medicines, following food-safety advice, and staying appropriately active. Severe pain, heavy bleeding, fainting, fever, or persistent vomiting needs urgent care.

\medterm{Dosage} The amount, frequency, route, and duration of a medicine given to achieve benefit while limiting harm. It may need adjustment for age, weight, kidney or liver function, interactions, and the condition treated.

\medterm{Dosage Forms} The physical form in which a medicine is supplied, such as a tablet, capsule, liquid, cream, inhaler, injection, patch, suppository, or implant. The form affects how and how quickly it is given.

\medterm{Dose} The amount of a medicine, radiation, or other substance given or received at one time. The appropriate dose depends on the substance, purpose, age, weight, organ function, interactions, and individual risk.

\medterm{Dose Rate} The amount of a medicine, radiation, or other exposure delivered during a specified time, such as milligrams per hour or grays per minute. It affects both benefit and toxicity.

\medterm{Dose-Limiting Toxicity} A treatment-related adverse effect severe enough to prevent increasing the dose or require reducing it. The term is used especially in early cancer trials to estimate a tolerable dose.

\medterm{Dose-Response Curve} A graph showing how the size of a biological effect changes as the dose of a medicine or substance increases. It helps describe potency, maximum effect, and the dose range that produces benefit or harm.

\medterm{Dose-Response Relationship} The relationship between the amount of a medicine or substance and the size of its effect. Increasing the dose may increase benefit, toxicity, or both, but the pattern depends on the substance and person.

\medterm{Dose Titration} Gradual adjustment of a medicine dose according to its effects, treatment goals, blood tests, or side effects. The dose may be increased, decreased, or kept the same as the person’s response is assessed.

\medterm{Drug-Food Interaction} A change in a medicine’s effect, absorption, or safety caused by food or drink. Food may increase or decrease absorption, while some foods can alter medicine metabolism or raise the risk of adverse effects.

\medterm{Pharmacodynamic Drug Interaction} An interaction in which two medicines have additive, opposing, or otherwise changed effects on the same body system, even when their blood concentrations are not altered.

\medterm{Pharmacokinetic Drug Interaction} An interaction in which one medicine changes another medicine’s absorption, distribution, metabolism, or elimination, altering its concentration or duration of action.

\medterm{Dosimeter} An instrument worn or placed near a person or object to measure or record exposure to ionizing radiation. It helps monitor occupational or treatment-related radiation.

\medterm{Dosimetry} Measurement and calculation of radiation dose received by a person, tissue, or object. In treatment, it helps plan and verify that the intended dose reaches the target while limiting exposure elsewhere.

\medterm{DOTATATE} A substance that binds somatostatin receptors found on many neuroendocrine tumors. When linked to a radioactive tracer, it is used in PET imaging to locate and stage tumors and assess possible targeted treatment.

\medterm{DOTS} A tuberculosis-control strategy involving reliable diagnosis, standardized combination treatment, medicine-supply support, monitoring, and sometimes directly observed doses. It aims to cure infection and limit drug resistance and transmission.

\medterm{Double Aortic Arch} A birth defect in which two aortic arches form a ring around the windpipe and esophagus. It may cause noisy breathing, cough, feeding difficulty, vomiting, or trouble swallowing.

\medterm{Double Blind} A study design in which participants and the researchers assessing outcomes do not know which treatment each participant receives during the blinded phase. It reduces expectation and assessment bias.

\synonyms
Double-masked

\medterm{Double Blind Trial} A clinical trial in which participants and the investigators measuring outcomes do not know the assigned treatment during the blinded phase. This helps reduce expectations from influencing results.

\medterm{Double Inlet Left Ventricle} A rare birth defect in which both upper heart chambers connect mainly to one lower chamber. It produces single-ventricle circulation and mixing of oxygen-rich and oxygen-poor blood and usually needs staged specialist treatment.

\medterm{Double Outlet Right Ventricle} A birth defect in which both major arteries arise mainly from the right lower heart chamber. An opening between the lower chambers allows blood from the left chamber to reach the arteries; anatomy and symptoms vary.

\medterm{Double Uterus} A birth difference in which the uterus develops as two separate cavities or two uteruses. Some people have no symptoms; others may have menstrual problems, infertility, miscarriage, or pregnancy complications.
\medterm{Double Vision} Seeing two images of one object. Causes include eye-surface or lens disease, eye misalignment, muscle or nerve problems, medicines, migraine, or brain disease; sudden onset with neurologic symptoms or eye pain is an emergency.

\medterm{Double-Balloon Enteroscopy} An endoscopic procedure using two balloons to move a flexible camera through much of the small intestine. It can locate bleeding or lesions, take biopsies, and perform selected treatments.

\medterm{Double-Blind Method} A research method in which treatment assignments are hidden from participants and relevant study personnel during the blinded phase to reduce expectation and assessment bias.

\medterm{Double-Blind Study} A clinical study in which participants and the investigators assessing outcomes do not know the assigned treatment group during the blinded phase.

\medterm{Double-Contrast Barium Enema} An X-ray examination in which barium coats the inside of the colon and air expands it. It can show polyps, masses, narrowing, or changes in the bowel lining.

\synonyms
Air-contrast barium enema

\medterm{Double-Outlet Right Ventricle} A birth defect in which both major arteries arise mainly from the right lower heart chamber. An opening between the lower chambers allows blood to reach the arteries; symptoms and treatment depend on the exact anatomy.
\medterm{Douche} A stream of water or solution directed into a body cavity, most often the vagina. Vaginal douching is not needed for hygiene and can irritate tissue or disturb normal bacteria and increase infection risk.

\medterm{Douglas, Pouch Of} A space in the pelvis between the uterus and rectum. In a person without a uterus, the comparable space lies between the bladder and rectum; fluid or blood can collect there.


\medterm{Dowager’S Hump} An exaggerated forward curve of the upper spine that creates a rounded back. It may result from osteoporosis and collapsed vertebrae; new pain, height loss, breathing difficulty, or worsening curvature needs assessment.

\synonyms
Kyphosis

\medterm{Down Syndrome} A genetic condition caused by an extra copy of chromosome 21. It may affect learning, muscle tone, growth, and appearance and is associated with heart, hearing, vision, thyroid, sleep, and digestive problems.

\synonyms
Trisomy 21 Syndrome
Down Syndrome Overview
Trisomy 21 (Down Syndrome Overview)
Down's Syndrome

\medterm{Downregulation} A decrease in the number, sensitivity, or activity of receptors, genes, or signaling pathways. It may occur after prolonged stimulation, medicine exposure, inflammation, or another regulatory change.

\medterm{Doxapram} A medicine that stimulates breathing and is used rarely in selected situations. It can cause agitation, high blood pressure, fast or abnormal heart rhythms, seizures, or other serious effects.

\medterm{Doxazosin} An alpha-1 blocker that relaxes blood vessels and prostate or bladder-neck muscle. It is used for high blood pressure or urinary symptoms from an enlarged prostate and can cause dizziness or fainting.

\medterm{Doxazosin Mesylate} A salt form of doxazosin, an alpha-1 blocker used for high blood pressure and urinary symptoms from an enlarged prostate. It can cause dizziness, low blood pressure, fainting, or falls, especially after the first dose.

\medterm{Doxepin} A tricyclic antidepressant also used in low doses for insomnia or skin itching. It can cause drowsiness, dry mouth, constipation, blurred vision, abnormal heart rhythms, and dangerous toxicity in overdose.

\medterm{Doxepin Overdose} Taking too much doxepin can cause severe drowsiness, confusion, seizures, low blood pressure, coma, and life-threatening heart-rhythm or conduction problems, especially with alcohol or other medicines.

\medterm{Doxercalciferol} A vitamin D-like medicine used in selected people with chronic kidney disease to reduce excess parathyroid hormone. Calcium and phosphate must be monitored because levels can become too high.

\medterm{Doxorubicin} An anthracycline cancer medicine that damages DNA and blocks cell division. It treats many cancers but can cause low blood counts, infection, nausea, hair loss, tissue injury after leakage, and dose-related heart damage.

\medterm{Doxorubicin Hydrochloride} A salt form of doxorubicin, an anthracycline cancer medicine used for many cancers. Important risks include low blood counts, infection, tissue injury after leakage, and dose-related heart damage.

\medterm{Doxycycline} A tetracycline antibiotic that blocks bacterial protein production. It treats many bacterial and tick-borne infections and can cause nausea, sun sensitivity, esophageal irritation, diarrhea, or tooth and bone effects in developing children.

\medterm{Doxycycline Hyclate} A salt form of doxycycline, a tetracycline antibiotic that blocks bacterial protein production. It is used for susceptible infections; nausea, sun sensitivity, esophageal irritation, diarrhea, and allergic reactions may occur.

\medterm{Doxylamine} A sedating antihistamine used in some sleep aids and pregnancy-related nausea treatments. It can cause drowsiness, dry mouth, blurred vision, constipation, urinary retention, confusion, or impaired driving.

\medterm{DPP-4 Inhibitor} An oral diabetes medicine that prolongs the action of hormones released after meals, increasing glucose-dependent insulin release and reducing glucagon. It is used mainly in type 2 diabetes and usually has a low risk of low glucose.

\medterm{DPP-4 Inhibitors} Oral medicines for type 2 diabetes that prolong hormones released after meals, helping the pancreas release insulin and reducing glucagon. They usually do not cause weight gain or low glucose when used alone but have modest glucose-lowering effects.

\medterm{DPT Immunization} Vaccination that protects against diphtheria, pertussis (whooping cough), and tetanus. Current childhood and booster vaccines use age-appropriate formulations and are given as a series according to the recommended schedule.
\medterm{DR Syndrome} An abbreviation that may refer to diabetic retinopathy syndrome or another context-specific disorder; the full clinical term should be specified.

\medterm{Dracunculiasis} Guinea-worm disease caused by a parasite acquired from unsafe drinking water containing infected tiny water fleas. About a year later, a painful blister usually forms on a leg as the adult worm emerges through the skin.

\medterm{Dracunculosis} Infection with the Guinea-worm parasite after drinking unsafe water containing infected water fleas. It causes a painful blister, usually on a leg, followed by emergence of a long adult worm through the skin.

\medterm{Dracunculus Medinensis} The parasitic worm that causes Guinea-worm disease. People become infected by drinking water containing tiny infected water fleas; roughly a year later, an adult female worm emerges through the skin.
\medterm{Drain} A tube or other device used to remove blood, pus, fluid, or air from a wound, body cavity, or organ. Drainage may prevent pressure, infection, or impaired healing after injury or surgery.

\medterm{Drain Cleaner Poisoning} Exposure to a caustic drain-cleaning product that can cause severe burns of the skin, eyes, airway, or gastrointestinal tract.

\medterm{Drain Opener Poisoning} Exposure to a caustic drain-cleaning product that can cause severe burns of the skin, eyes, airway, or gastrointestinal tract.

\medterm{Drain (Surgical)} A tube placed during or after surgery to remove blood, pus, or other fluid from an operative site. Possible complications include infection, blockage, leakage, bleeding, and movement of the tube.

\synonyms
Surgical drain

\medterm{Drainage} Removal of blood, pus, fluid, or air from a body site, wound, or abnormal collection. It may occur naturally or use a needle, catheter, tube, or surgical opening to relieve pressure and aid healing.

\medterm{Drainage Procedures} Treatments that remove pus, blood, or other collected fluid, sometimes with ultrasound or CT guidance. They may use a needle, catheter, tube, or surgery to treat abscesses and other fluid collections.

\medterm{Drainpipe Cleaners} Strong chemical products used to clear blocked pipes, commonly containing caustic alkalis or acids. Swallowing, inhaling, or splashing these products can cause severe burns to the mouth, throat, eyes, skin, or lungs and requires urgent poison-control or medical advice.


\medterm{Dream} A sequence of images, thoughts, sensations, or emotions experienced during sleep. Dreams are often most vivid during rapid-eye-movement sleep and may be influenced by stress, medicines, illness, or sleep disorders.

\medterm{Dreams} Experiences involving images, thoughts, sensations, or emotions during sleep, often vivid during rapid-eye-movement sleep. Dreams are normal, but distressing, violent, or disruptive dreams may occur with stress, medicines, or sleep disorders.

\medterm{Dreams (Geriatric)} Dreams in older adults may become more vivid or be affected by medicines, sleep disorders, depression, or changes in the sleep cycle. Repeated frightening dreams, acting them out, confusion, or sudden change requires clinical assessment.

\synonyms
Dreams in older adults

\medterm{Drepanocytosis} Sickle-cell disease or a related condition in which an abnormal form of hemoglobin makes red blood cells become rigid and sickle-shaped. This can cause anemia, painful blockages of blood flow, and damage to organs.

\medterm{DRESS Syndrome} A serious delayed allergic reaction to a medicine, usually beginning two to eight weeks after exposure. It may cause fever, widespread rash, facial swelling, abnormal blood counts, hepatitis, kidney injury, or other organ inflammation.

\medterm{Dressing} Material placed over a wound to protect it from contamination, absorb drainage, control minor bleeding, maintain moisture, or deliver local treatment. It should be changed as directed and monitored for infection.

\medterm{Dressler Syndrome} An immune-related inflammation of the sac around the heart that occurs weeks after a heart attack, heart surgery, or other heart injury. It can cause fever, sharp chest pain, and fluid around the heart or lungs.

\synonyms
Post-Cardiac Injury Syndrome
Postmyocardial Infarction Syndrome

\medterm{Dressler's Syndrome} Alternate spelling; see \textit{Dressler Syndrome}.

\medterm{Dressler’S Syndrome} Alternate spelling; see \textit{Dressler Syndrome}.

\medterm{DRI} Dietary reference values used to estimate the amounts of vitamins, minerals, protein, and other nutrients needed by healthy people. They help assess diets but are not individualized treatment targets for every illness.

\synonyms
Dietary Reference Intakes

\medterm{Dried Specimen} A biological sample collected and dried on a prepared card or other surface for transport and laboratory testing. Dried blood spots are commonly used for screening, but correct collection and handling are essential.


\medterm{Drift (Sit-To-Stand)} A tendency to lose balance or move away from a stable position while standing up from a chair. It may reflect weakness, poor coordination, dizziness, or impaired balance and should be interpreted with the full examination.

\synonyms
Sit-to-Stand

\medterm{Drinking Water} Water safe for human consumption. It should be free of disease-causing organisms and harmful chemicals; boiling, filtration, chlorination, or bottled water may be needed when the supply is contaminated or uncertain.

\medterm{Drip} A drop-by-drop flow of liquid, or fluid given slowly into a vein through tubing and a catheter. The infusion rate is controlled to deliver medicines, fluids, blood products, or nutrients safely.

\medterm{Drip (Intravenous Infusion)} Delivery of fluid, blood products, nutrients, or medicine directly into a vein through a catheter. The amount and speed are prescribed and monitored for reactions, fluid overload, infection, or vein irritation.

\synonyms
Intravenous Infusion

\medterm{Driving And Older Adults} Evaluation of whether an older person can drive safely, considering vision, attention, memory, hearing, movement, health conditions, and medicines. Age alone does not determine driving ability.

\medterm{Driving Evaluation} An assessment of driving ability after illness, injury, or a health change. It may examine vision, attention, memory, reaction time, strength, and vehicle control in a clinic and on the road; adaptive equipment may help.

\medterm{Driving Rehabilitation} Assessment and training to help a person resume driving safely or use adaptive vehicle controls after illness or injury. It may address vision, attention, reaction time, strength, vehicle operation, and road safety.

\medterm{Dronabinol} A medicine related to the active substance in cannabis, used for selected cases of chemotherapy-related nausea or severe appetite loss. It may cause dizziness, sleepiness, confusion, mood changes, or impaired coordination.

\medterm{Drool} Saliva that unintentionally escapes from the mouth. It may result from excess saliva, difficulty swallowing, poor lip closure, dental problems, or a neurologic condition affecting mouth control.

\medterm{Drooling} Unintentional loss of saliva from the mouth caused by excess saliva, difficulty swallowing, poor lip closure, or reduced control of the mouth. Persistent drooling may need assessment for dental, medication, or neurologic causes.

\medterm{Drop Foot} Alternate index label; see \textit{Foot drop}.

\medterm{Drop Hand Syndrome} Weakness or paralysis of the muscles that straighten the wrist and fingers, usually from injury or pressure on the radial nerve. The wrist hangs downward and gripping or releasing objects becomes difficult.

\medterm{Droperidol} A medicine used in selected clinical settings to treat severe nausea or agitation. It can prolong the heart’s electrical recovery time and rarely cause dangerous arrhythmias, so heart-risk factors and monitoring may be important.

\medterm{Droplet Precautions} Infection-control measures for infections spread by respiratory droplets released during coughing, sneezing, or talking. They may include a medical mask, hand hygiene, patient separation, and limiting transport according to the infection.

\medterm{Dropper Bottle} A bottle with a device that releases liquid one drop at a time. The medicine’s concentration, prescribed number of drops, route, and measuring device must be followed because drop size can vary.


\medterm{Drowning} Breathing impairment caused by submersion or immersion in liquid. It may be fatal or nonfatal. Lack of oxygen can injure the brain and other organs; complications include aspiration, lung injury, and respiratory failure.

\medterm{Drowning (Dry/Wet Signs)} Alternate index label; see \textit{Drowning}. The terms “dry drowning” and “wet drowning” are outdated and should not be used to classify drowning.

\medterm{Drowsiness} An increased tendency to fall asleep or reduced alertness. It may result from too little sleep, medicines, alcohol or other substances, metabolic illness, sleep disorders, infection, or neurologic disease; sudden severe drowsiness needs urgent assessment.

\medterm{Drug} A substance that changes body function or is used to prevent, diagnose, or treat disease. Drugs include prescription medicines, nonprescription medicines, biological products, and substances used recreationally or illegally.

\medterm{Drug Addiction} Legacy, nonpreferred label; see \textit{Substance Use Disorder}.

\medterm{Drug Caution Code} A label or code that warns about an important medicine risk, such as allergy, bleeding, interaction, overdose, or a need for special monitoring. Its meaning depends on the coding system and product information.
\medterm{Drug Delivery} The way a medicine is prepared and given so its active ingredient reaches the intended body site at an effective dose. Routes include swallowing, inhalation, skin application, injection, implantation, and targeted systems.

\medterm{Drug Desensitization} A supervised process that gives gradually increasing doses of a medicine to temporarily reduce an allergic reaction when that medicine is essential and no suitable alternative exists. The tolerance usually ends when doses are interrupted.

\medterm{Drug Discovery} The process of finding and developing substances that may prevent, diagnose, or treat disease. It includes laboratory research, safety testing, dose studies, clinical trials, manufacturing, and review by a regulatory authority.

\medterm{Drug Hypersensitivity} An immune reaction to a medicine that may cause rash, itching, swelling, breathing difficulty, fever, blood abnormalities, or organ inflammation. Severe forms include anaphylaxis and serious blistering or systemic skin reactions.

\medterm{Drug Implant} A solid medicine-containing device placed under the skin or in another tissue to release medicine slowly over time. It may provide local or whole-body treatment and can require removal when treatment ends.

\medterm{Drug Repurposing} Studying or using a medicine developed for one condition to treat a different condition. Earlier safety information may help, but the new use still needs evidence about effectiveness, dose, interactions, and risks.

\medterm{Drug Safety} The prevention, detection, and management of harmful effects, interactions, errors, and misuse involving medicines. It includes choosing the right medicine and dose, checking allergies and interactions, following instructions, monitoring effects, and reporting suspected problems.

\medterm{Drug Screening} Testing a sample such as urine, blood, saliva, or hair for selected drugs or their breakdown products. Initial results can be falsely positive or negative and may require confirmation with a more specific test.

\medterm{Drug Testing} Laboratory analysis of blood, urine, saliva, hair, or another sample to detect medicines, recreational drugs, their breakdown products, or poisons. The test’s purpose, detection limits, timing, and medicines taken affect interpretation.

\medterm{Drug-Induced Lupus} A lupus-like autoimmune reaction caused by certain medicines, often producing joint pain, muscle pain, fever, and inflammation around the lungs or heart. Symptoms and abnormal antibodies usually improve after the causative medicine is stopped under medical supervision.

\synonyms
Drug-Induced Lupus Erythematosus; DILE

\medterm{Drugs, Statin} A group of medicines that lowers LDL (“bad”) cholesterol by reducing cholesterol production in the liver. Statins lower the risk of heart attack and stroke; muscle symptoms and liver abnormalities can occur.

\medterm{Drugs That May Cause Erection Problems} Some medicines can contribute to difficulty getting or keeping an erection by affecting blood flow, hormones, nerves, or mood. A clinician should review the timing and all medicines before changing treatment.
\medterm{Drusen} Yellow-white deposits that form beneath the retina, often with aging. Small deposits may cause no symptoms; many or larger deposits can be associated with age-related macular degeneration and need eye examination.

\medterm{Dry Age-Related Macular Degeneration} Alternate index label; see \textit{Dry macular degeneration}.

\medterm{Dry AMD} Alternate index label; see \textit{Dry macular degeneration}.

\medterm{Dry Cell Battery Poisoning} Injury or poisoning caused by swallowing, leaking, or otherwise contacting a dry-cell battery. A swallowed button battery can rapidly burn the esophagus and is an emergency; do not induce vomiting and seek urgent medical help.
\medterm{Dry Cough} A cough that produces little or no mucus. Common causes include viral infection, asthma, acid reflux, allergies, smoking, certain medicines, and other airway or lung conditions; persistent cough needs assessment.

\medterm{Dry Cupping} A complementary treatment in which cups create suction on intact skin without making cuts. It may temporarily relieve some muscle pain, but evidence is limited; bruising, burns, skin injury, infection, and worsening of skin disease can occur.

\medterm{Dry Eye} A disorder in which tears are insufficient or evaporate too quickly, causing burning, grittiness, redness, watering, light sensitivity, or blurred vision. Long-lasting disease can damage the eye surface.

\synonyms
Dry-eye disease; keratoconjunctivitis sicca; Dry Eyes

\medterm{Dry Eye Disease} Alternate name; see \textit{Dry Eye}.

\medterm{Dry Eye Syndrome} Alternate name; see \textit{Dry Eye}.


\medterm{Dry Gangrene} Tissue death caused mainly by severely reduced arterial blood flow. The affected area becomes dry, shrunken, and dark, often with a clear boundary; infection can develop and urgent vascular assessment is needed.

\medterm{Dry Hair} Hair that feels rough, brittle, dull, or frizzy because its outer layer is damaged or lacks moisture. Heat, chemical treatments, weather, frequent washing, nutritional problems, and thyroid disease may contribute.

\medterm{Dry Ice} Solid carbon dioxide that changes directly into a cold gas. Contact can cause severe cold burns, and gas in an enclosed space can displace oxygen and cause unconsciousness or suffocation; never seal it in a container.

\medterm{Dry Macular Degeneration} The more common form of age-related macular degeneration, in which light-sensing cells in the macula gradually deteriorate without bleeding from new abnormal vessels. It can cause slow loss of detailed central vision.

\synonyms
Dry Age-Related Macular Degeneration
Dry AMD
Macular Degeneration, Dry

\medterm{Dry Mouth} A feeling or condition of having too little saliva. It can cause thirst, sticky mouth, difficulty speaking or swallowing, altered taste, tooth decay, mouth sores, and infections; medicines are a common cause.
\synonyms
xerostomia.

\medterm{Dry Mouth During Cancer Treatment} Reduced saliva during or after cancer treatment, especially head-and-neck radiation or certain medicines. It can cause mouth discomfort, difficulty chewing or swallowing, tooth decay, taste changes, and oral infection.

\medterm{Dry Powder Inhaler} A breath-powered device that delivers powdered medicine into the lungs. It does not require pressing and breathing at the same time, but the user must inhale strongly and correctly for the dose to reach the airways.


\medterm{Dry Skin} Skin that is rough, scaly, tight, or itchy because it has lost moisture or its protective barrier is damaged. Aging, cold or dry air, frequent washing, eczema, medicines, and some illnesses can contribute.


\medterm{Dry Socket} Painful inflammation after tooth extraction when the protective blood clot is lost or breaks down too early. Severe pain may begin a few days later and spread to the ear or jaw; dental treatment is needed.

\medterm{Alveolar Osteitis} Painful inflammation of a tooth socket after extraction, usually caused by loss or breakdown of the blood clot that normally protects the socket. It is also called dry socket.

\synonyms
alveolar osteitis.
Alveolar Osteitis

\medterm{Dry Socket (Alveolar Osteitis)} Alternate name; see \textit{Dry Socket}.

\medterm{Dry Socket Overview} Alternate name; see \textit{Dry Socket}.

\medterm{Dryness Vaginal (Vaginal Dryness And Vaginal Atrophy)} Alternate index label for Vaginal Dryness and Vaginal Atrophy.

\medterm{Dry-Powder Inhaler} A device that delivers powdered medicine to the lungs when a person inhales through it. Correct preparation, a strong breath, and proper technique are needed; the device should not usually be used with a spacer.

\medterm{DSM} The Diagnostic and Statistical Manual of Mental Disorders, a reference that describes recognized mental-health diagnoses and their criteria. Clinicians use it with history and examination; it does not replace individual clinical judgment.

\medterm{DSPS} Alternate index label; see \textit{Delayed sleep phase}.

\medterm{DSRCT} Alternate index label; see \textit{Desmoplastic small round cell tumor}.

\medterm{DSWPD} Alternate index label; see \textit{Delayed sleep phase}.

\medterm{DTaP} A vaccine for infants and young children that protects against diphtheria, tetanus, and pertussis (whooping cough). It is given as a series, with later booster vaccines used as children grow older.


\medterm{Dual Antiplatelet Therapy} Treatment with two medicines that reduce platelet clumping, usually after a heart attack or coronary stent. It lowers clot risk but increases bleeding risk; the duration depends on the procedure, heart condition, and bleeding risk.

\medterm{Dual Energy X-Ray Absorptiometry} See \textit{Bone Mineral Density Test}.

\synonyms
DXA; DEXA

\medterm{Dual-Energy CT} A CT scan that uses two X-ray energy levels to help distinguish materials in the body. It can improve evaluation of kidney stones, gout, blood vessels, bleeding, and some tissues while using the usual CT principles and risks.

\medterm{Dual-Energy X-Ray Absorptiometry} Alternate spelling; see \textit{Dual Energy X-Ray Absorptiometry}.

\medterm{Dual-Photon Absorptiometry} An older test that used two photon energy levels to estimate bone mineral content. It has largely been replaced by dual-energy X-ray absorptiometry for measuring bone density and fracture risk.
\medterm{Duane Retraction Syndrome} A birth-related eye-movement disorder caused by abnormal nerve control of an eye muscle. Turning the eye inward or outward may be limited, and the eyeball may pull backward and the eyelids narrow.

\medterm{Dubin-Johnson Syndrome} An inherited liver-transport disorder in which conjugated bilirubin is not moved normally into bile. It causes intermittent mild jaundice, often triggered by illness or stress, but usually does not cause serious liver damage.

\medterm{Dubowitz Syndrome} A rare developmental condition associated with poor growth, a small head, characteristic facial features, developmental or learning difficulties, and variable immune, skin, or behavioral problems. Features differ among affected people.

\medterm{Duchenne Muscular Dystrophy} An inherited muscle disease, usually affecting boys, caused by changes in the dystrophin gene. Progressive muscle weakness begins in childhood and may affect walking, breathing, the heart, and learning or behavior.

\medterm{Duck Embryo Vaccine} An older rabies vaccine made by growing inactivated rabies virus in duck embryos. It has largely been replaced by safer, more consistent cell-culture vaccines and is not commonly used today.

\medterm{Duct} A tube or channel that carries fluid within the body, such as bile, pancreatic juice, saliva, tears, or secretions from a gland. Blockage or injury can cause pain, swelling, infection, or impaired organ function.

\medterm{Ductal Carcinoma} A cancer that begins in cells lining a duct. The term most often describes breast cancer or pancreatic cancer; its behavior and treatment depend on the organ, spread, and tumor features.

\medterm{Ductal Carcinoma In Situ} Abnormal cancer-like cells confined to the milk ducts of the breast without invasion into nearby breast tissue. It has not spread elsewhere, but some cases can become invasive if untreated.

\synonyms
DCIS
Intraductal Carcinoma

\medterm{Ductal-Dependent Lesions} Serious congenital heart defects in which blood flow to the lungs or body depends on the fetal ductus arteriosus remaining open after birth. Affected newborns may rapidly become blue, breathless, pale, or shocked and need urgent specialist care.

\medterm{Ductless Glands} Glands that release hormones directly into the bloodstream rather than through a tube. Examples include the pituitary, thyroid, adrenal glands, and pancreatic islets; their hormones regulate growth, metabolism, stress, and reproduction.

\medterm{Ductus Arteriosus} A blood vessel in the fetus connecting the pulmonary artery to the aorta, allowing blood to bypass the lungs before birth. It normally closes soon after birth; failure to close can strain the heart and lungs.

\medterm{Ductus Venosus} A fetal blood vessel that directs oxygen-rich blood from the umbilical vein toward the heart while partly bypassing the liver. It normally closes after birth; abnormal fetal flow can indicate circulatory or placental problems.

\medterm{Duffy-Null Associated Neutrophil Count} A naturally lower measured neutrophil count associated with the Duffy-null blood-group pattern, common in some populations. It usually does not increase infection risk and should not automatically be treated as disease.

\medterm{Duke Criteria} A set of clinical, blood-culture, and heart-imaging findings used to classify suspected infective endocarditis. The criteria combine major and minor findings; the result supports diagnosis but must be interpreted with the patient’s history and examination.

\synonyms
Modified Duke Criteria; Duke Endocarditis Criteria

\medterm{Dukes Classification} An older system for describing how far colorectal cancer has spread through the bowel wall, into nearby lymph nodes, or to distant organs. Modern staging systems are more detailed and are generally preferred.

\medterm{Dull Pain} Aching, heavy, or pressure-like pain that is not sharp or stabbing. It can arise from muscles, organs, inflammation, or other tissues, and its significance depends on location, duration, associated symptoms, and examination findings.

\medterm{Duloxetine} A medicine that increases serotonin and norepinephrine signaling in the brain. It treats depression, anxiety, some nerve pain, fibromyalgia, and chronic muscle or bone pain; nausea, sleep changes, sweating, and withdrawal symptoms may occur.

\medterm{Duloxetine Hydrochloride} The salt form of duloxetine, a medicine used for depression, anxiety, some nerve pain, fibromyalgia, and chronic muscle or bone pain. It may cause nausea, sleep changes, sweating, blood-pressure changes, or withdrawal symptoms.

\medterm{Dumbbell Ossification} Bone formation or a bone lesion shaped like a dumbbell, often extending through an opening between the spinal canal and surrounding tissues. Imaging is needed to identify its cause and relationship to nerves or other structures.

\medterm{Dumdum Fever} An older name for visceral leishmaniasis, a parasitic infection transmitted by sandflies. It can cause prolonged fever, weight loss, enlarged spleen or liver, anemia, low blood counts, and life-threatening illness without treatment.

\medterm{Dumping Syndrome} Symptoms caused by food moving too quickly from the stomach into the small intestine, often after stomach surgery. Early symptoms include cramps, diarrhea, flushing, sweating, dizziness, and a fast heartbeat; later low blood sugar may occur.

\medterm{Dunbar Syndrome} Alternate index label; see \textit{Median arcuate ligament syndrome}.

\medterm{Duodenal} Relating to the duodenum, the first part of the small intestine immediately beyond the stomach. This segment receives bile and pancreatic juice and begins digestion and absorption of nutrients.
\medterm{Duodenal Adenocarcinoma} A cancer arising from gland-forming cells in the duodenum. It may cause abdominal pain, vomiting, bleeding, anemia, jaundice, or blockage; diagnosis usually requires imaging, endoscopy, and tissue sampling.

\medterm{Duodenal Atresia} A birth defect in which the duodenum is completely blocked, preventing food and digestive fluid from passing. Newborns usually develop green vomiting and feeding difficulty; it is often associated with other birth defects and requires surgery.

\medterm{Duodenal Bulb} The widened first part of the duodenum just beyond the stomach outlet. It receives partly digested food and is a common site for peptic ulcers; inflammation or scarring can narrow it and delay passage.

\medterm{Duodenal Disease} Any disorder affecting the duodenum, such as inflammation, ulcer, infection, bleeding, blockage, poor nutrient absorption, or a growth. Symptoms may include upper-abdominal pain, vomiting, anemia, or weight loss.

\medterm{Duodenal Diverticulum} A pouch that bulges outward from the duodenal wall, often near the openings of the bile and pancreatic ducts. It commonly causes no symptoms but can rarely lead to bleeding, inflammation, blockage, pancreatitis, or bile-duct infection.

\medterm{Duodenal Fistula} An abnormal connection between the duodenum and another organ, bowel segment, or skin. Digestive fluid may leak through it, causing infection, skin damage, dehydration, electrolyte loss, or poor nutrition.

\medterm{Duodenal Hemorrhage} Bleeding from the duodenum, often caused by a peptic ulcer, inflammation, abnormal blood vessel, injury, or cancer. It may cause black stools, vomiting blood, weakness, dizziness, anemia, or shock.

\medterm{Duodenal Infection} Infection affecting the duodenum, caused by bacteria, parasites, viruses, or spread from nearby digestive disease. Symptoms can include upper-abdominal pain, nausea, vomiting, diarrhea, fever, or bleeding.

\medterm{Duodenal Neoplasm} An abnormal growth in the duodenum that may be benign or cancerous. It can cause pain, bleeding, vomiting, blockage, anemia, or jaundice; evaluation depends on its size, location, appearance, and tissue type.

\medterm{Duodenal Obstruction} A partial or complete blockage that prevents food and digestive fluid from passing through the duodenum. It can cause repeated vomiting, abdominal swelling, pain, and dehydration and may need urgent imaging and treatment.

\medterm{Duodenal Papilla} A small raised opening in the duodenal wall where bile and pancreatic juice enter the intestine. Narrowing, blockage, or inflammation at this site can cause jaundice, pancreatitis, or digestive symptoms.

\medterm{Duodenal Perforation} A hole through the full thickness of the duodenal wall, usually caused by a severe ulcer, injury, or a medical procedure. Sudden severe abdominal pain, rigid muscles, and infection of the abdominal lining require emergency treatment.

\medterm{Duodenal Stenosis} Narrowing of the duodenum that partly blocks passage of food and digestive fluid. It may cause vomiting, abdominal swelling, pain, early fullness, poor feeding, or dehydration; treatment depends on the cause and severity.

\medterm{Duodenal Ulcer} An open sore in the duodenal lining, commonly related to Helicobacter pylori infection or anti-inflammatory medicines. It may cause upper-abdominal pain and can bleed, perforate the bowel, or obstruct food passage.

\medterm{Duodenitis} Inflammation of the duodenal lining, which may result from Helicobacter pylori, excess acid, anti-inflammatory medicines, infection, celiac disease, or inflammatory bowel disease. It can cause upper-abdominal pain, nausea, vomiting, or bleeding.

\medterm{Duodenoscope} A flexible viewing tube passed through the mouth and stomach to examine the duodenum and reach the openings of the bile and pancreatic ducts. It may be used during procedures to remove stones or place stents.

\medterm{Duodenostomy} A surgically created opening from the duodenum to the skin or another part of the intestine. It may be used in selected cases for feeding, drainage, or bypassing a blockage.

\medterm{Duodenum} The first part of the small intestine, between the stomach and jejunum. It receives bile and pancreatic juice, mixes them with stomach contents, and begins much of the digestion and absorption of nutrients.

\medterm{Duplex Doppler Ultrasound Scanning} An ultrasound test that shows the structure of blood vessels and measures the direction and speed of blood flow. It can detect narrowing, blockage, clots, aneurysms, or abnormal connections without radiation.
\synonyms
Duplex ultrasound

\medterm{Duplex Surveillance} Planned examinations and ultrasound checks of a blood vessel, graft, stent, aneurysm, or dialysis access. Monitoring can detect recurrent narrowing, clotting, leakage, or failure before serious symptoms develop.

\medterm{Duplex Ultrasound} An ultrasound examination that combines pictures of body structures with Doppler measurement of blood flow. It is commonly used to assess blood vessels and can identify narrowing, blockage, clots, or abnormal flow.

\medterm{Duplication} A genetic change in which a section of DNA or a chromosome is copied one or more extra times. The additional genetic material can alter how much of a gene is active and may cause a health condition, although some duplications have little or no effect.

\medterm{Dupuytren Contracture} Progressive thickening and shortening of tissue in the palm that pulls one or more fingers toward the palm and can limit hand function.

\synonyms
Dupuytrens Contracture (Dupuytren Contracture)

\medterm{Dupuytren, Guillaume} A French surgeon whose name is associated with Dupuytren contracture, progressive thickening and shortening of palmar fascia that bends the fingers toward the palm.

\medterm{Dupuytren's Contracture} A progressive fibrosing disorder of the palmar fascia that produces nodules and cords and may permanently flex one or more fingers toward the palm. Treatment depends on severity and functional impairment and may include observation, injection, needle fasciotomy, or surgery.

\medterm{Dupuytrens Contracture (Dupuytren Contracture)} Alternate index label for Dupuytren Contracture.

\medterm{Dupuytren's Disease} A usually progressive disorder in which the palmar fascia, a layer of fibrous tissue beneath the skin of the palm and fingers, becomes thick and forms nodules and cords. The cords can gradually prevent one or more fingers, especially the ring or little finger, from straightening and may limit hand function. It usually begins in middle age or later and is more common in men, although anyone can develop it. The condition begins in the palm and is not primarily a disease of the finger joints; the later bending deformity is called Dupuytren contracture.

\medterm{Dura} The tough outermost layer of protective tissue surrounding the brain and spinal cord. It helps contain cerebrospinal fluid, supports blood vessels, and can become inflamed, torn, thickened, or infected.

\medterm{Dura Mater} The tough outer layer of tissue surrounding and protecting the brain and spinal cord. It helps support blood vessels and contains the fluid-filled coverings around the central nervous system.

\medterm{Durable Power Of Attorney} A legal authorization allowing a designated person to act or make decisions for another person when specified conditions are
met. In healthcare, it may appoint a healthcare proxy.

\medterm{Dural Arteriovenous Fistulas} Abnormal links between arteries and veins in the protective covering around the brain or spinal cord. They can cause headaches, pulsating sounds, seizures, bleeding, or neurologic problems and may need vascular treatment.

\medterm{Dural Sac} The tough, tubular covering around the spinal cord, its surrounding fluid, and the spinal nerve roots inside the spinal canal. A tear can allow spinal-fluid leakage and cause positional headaches.

\medterm{Dural Venous Sinuses} Channels between layers of the brain’s tough covering that collect blood and cerebrospinal fluid and drain them toward the neck veins. A clot in these channels can cause headache, seizures, vision problems, or stroke-like symptoms.

\medterm{Duramorph} A brand name for preservative-free morphine used by injection, including in epidural or intrathecal analgesia; respiratory depression is a serious risk.

\medterm{Durapatite} Another name for hydroxyapatite, the calcium-phosphate mineral that gives bones and teeth much of their hardness. Its amount and arrangement contribute to bone strength and can change in metabolic bone disease.
\medterm{Durules} A formulation designed to release a medicine slowly over an extended period. Such tablets or pellets must be used exactly as directed because crushing or chewing may release too much medicine at once.

\medterm{Dust} Small solid particles suspended in air or settled on surfaces. Inhaled dust can irritate the nose and lungs; some types can cause allergy, asthma, infection, poisoning, or long-term lung disease.

\medterm{Dust Disease} Lung or other illness caused by repeated exposure to harmful dust, often at work. Effects depend on the dust and exposure and may include asthma, scarring of lung tissue, chronic cough, or cancer.

\medterm{Dust Mite Allergy} An immune reaction to proteins from tiny mites living in household dust. It can cause sneezing, blocked or itchy nose, watery eyes, cough, wheezing, eczema flares, or asthma symptoms.

\synonyms
Allergy, Dust Mite

\medterm{Dutasteride} A medicine that lowers the hormone dihydrotestosterone and gradually shrinks an enlarged prostate. It improves urinary symptoms and lowers retention risk but may reduce sexual function and lowers the blood-test level used to screen for prostate cancer.

\medterm{DVT} A blood clot in a deep vein, usually in the leg. It can cause swelling, pain, warmth, or redness and may travel to the lungs, causing a life-threatening pulmonary embolism.

\medterm{DVT Prevention} Steps to reduce the risk of a deep-vein blood clot, including early movement, leg exercises, compression devices, and preventive anticoagulant medicine when appropriate. The choice depends on clot risk and bleeding risk.

\medterm{Dwarfism} A medical condition causing unusually short adult height, often defined as about 4 feet 10 inches (147 cm) or less. Causes include skeletal disorders, growth-hormone problems, genetic conditions, and chronic disease; features vary by cause.

\medterm{Dwarfism (Achondroplasia)} Alternate index label for Achondroplasia.

\medterm{DWI FLAIR Mismatch} A pattern on brain MRI in which an area is abnormal on diffusion images but not yet clearly abnormal on FLAIR images. In selected people with an unknown stroke onset time, it may help identify potentially treatable recent stroke.
\medterm{Dx} A medical abbreviation for diagnosis or diagnostic assessment. Because abbreviations can be misunderstood, the full word should be used in patient instructions and other documents where clarity matters.
\medterm{DXA Scan (Dual-Energy X-Ray Absorptiometry)} Alternate name; see \textit{Dual Energy X-Ray Absorptiometry}.

\medterm{D-Xylose Absorption} An older test of how well the small intestine absorbs a simple sugar after it is swallowed. Low amounts in blood or urine may suggest intestinal lining disease, but kidney function and collection quality affect results.

\medterm{Dyad} A pair of related people, structures, cells, or measurements considered together. Medical examples include two linked findings or two parts of a cellular or body structure that function together.

\medterm{Dye} A colored substance used to stain tissue, identify cells, make structures visible, or provide contrast during a medical test. Some dyes can cause allergic reactions or affect the kidneys, especially when injected.

\medterm{Dye Lasers} Lasers that use a liquid organic dye to produce adjustable wavelengths of light. In medicine, selected wavelengths can treat blood-vessel or pigmented skin lesions; eye and skin protection is essential.

\medterm{Dye Remover Poisoning} Harm from swallowing, inhaling, or getting dye-removing chemicals on the skin or eyes. Depending on the ingredients, exposure can cause burns, breathing problems, vomiting, poisoning, or organ injury and needs poison-control advice.

\medterm{Dying} The physical process of approaching death, often involving increasing weakness, sleepiness, reduced appetite and drinking, changes in breathing, and less responsiveness. Emotional, social, cultural, and spiritual support may help the person and family.

\medterm{Dynamin Gtpase} A family of GTP-binding proteins that constrict membranes during vesicle budding, endocytosis, and other cell-membrane processes.

\medterm{Dynamometer} An instrument that measures force, commonly hand-grip or muscle strength. Clinicians use the result to compare sides, identify weakness, monitor recovery, and assess physical function during rehabilitation or neurologic examination.

\medterm{Dynein Atpase} The energy-using activity of dynein, a motor protein that moves materials inside cells along microtubules. Dynein also powers the movement of cilia and flagella, which help move fluid or cells.

\medterm{Dynorphin} A naturally occurring opioid peptide that mainly activates kappa-opioid receptors in the nervous system. It influences pain, stress, mood, movement, and hormone regulation.

\medterm{Dyphylline} A bronchodilator used in some countries to relax airway muscles in asthma or chronic obstructive lung disease. It may cause nausea, tremor, headache, nervousness, or a fast heartbeat.

\medterm{Dysaesthesia} An unpleasant abnormal sensation, such as burning, electric-shock pain, pins and needles, or excessive sensitivity to light touch. It commonly results from nerve injury, diabetes, medicines, or other nerve disease.

\medterm{Dysarthria} Slurred, weak, slow, or poorly coordinated speech caused by a problem in the brain, nerves, or muscles used for speaking. Sudden dysarthria can be a sign of stroke.

\medterm{Dysautonomia} A disorder of the automatic nervous system that controls blood pressure, heart rate, sweating, temperature, digestion, bladder function, and sexual function. Causes include nerve disease, diabetes, medicines, and inherited conditions.

\medterm{Dysbarism} Illness caused by a change in surrounding pressure, usually during diving, flying, or hyperbaric treatment. It includes decompression sickness and pressure injury to the ears, sinuses, lungs, or other tissues.

\medterm{Dysbiosis} An imbalance or change in the community of microorganisms living in the body, especially the intestine. Antibiotics, infection, diet, and illness may contribute; dysbiosis is associated with some digestive and inflammatory disorders, but it is not always a proven cause.

\synonyms
Gut Microbiome Imbalance

\medterm{Dyscalculia} A persistent difficulty understanding numbers, quantities, or arithmetic that is not explained by inadequate teaching or general intellectual disability. It may begin during development or follow brain injury and can affect school or daily tasks.

\medterm{Dyschezia} Difficult or painful passage of stool, often with excessive straining or a feeling of incomplete emptying. Causes include constipation, pelvic-floor coordination problems, anal disorders, and other bowel disease.

\medterm{Dyscrasia} An older broad term for an abnormal blood or bone-marrow condition. It may describe abnormal blood cells, clotting, or proteins in the blood, so the specific disorder should be named when possible.

\synonyms
Blood dyscrasia

\medterm{Dyscrasias} An older plural term for disorders affecting blood, bone marrow, or body-fluid composition. It may refer to abnormal blood cells, clotting, or abnormal proteins, so the specific disorder must be identified.

\medterm{Dysembryoplastic Neuroepithelial Tumor} A usually slow-growing, low-grade brain tumor that develops most often in the outer part of the brain in children or young adults. It commonly causes long-standing focal seizures and is often treatable with surgery.

\medterm{Dysentery} Severe intestinal infection causing frequent diarrhea with blood or mucus, abdominal cramps, fever, and an urgent feeling of needing to pass stool. Bacteria or parasites may cause it; dehydration and complications require medical care.

\medterm{Dysesthesia} An unpleasant abnormal sensation, such as burning, electric-shock pain, or painful sensitivity to light touch. It may result from injury or disease of nerves in the brain, spinal cord, or limbs.

\medterm{Dysferlinopathy} A group of inherited muscle diseases caused by changes in both copies of the DYSF gene, which makes a muscle-membrane repair protein. Weakness often begins in adolescence or early adulthood, affecting the hips, shoulders, or calves.

\synonyms
DYSF-Related Muscular Dystrophy; Dysferlin Deficiency

\medterm{Dysfibrinogenemia} An inherited or acquired condition in which fibrinogen, a blood-clotting protein, has abnormal structure or function. It may cause unusual bleeding, excessive clotting, or no symptoms and is diagnosed with clotting and specialized blood tests.

\medterm{Dysfunction} Abnormal, reduced, or poorly controlled function of an organ, tissue, body system, or process. It may be temporary or permanent and can result from disease, injury, medicines, inherited conditions, stress, or aging.

\medterm{Dysfunctional Voiding} Abnormal coordination between the bladder muscle and the muscles around the urethra during urination. It may cause a weak or interrupted stream, straining, incomplete emptying, urinary infections, wetting, or constipation, especially in children.

\medterm{Dysfunctional Uterine Bleeding} An older term for bleeding from the uterus that is unusually heavy, prolonged, frequent, or irregular without an identified structural cause. Current evaluation considers pregnancy, medicines, hormones, ovulation, and uterine disease.

\medterm{Dysgammaglobulinemia} An abnormal level or pattern of antibodies in the blood, with some antibody types low or absent and others normal or increased. It may cause repeated infections, autoimmune problems, or no symptoms and is identified by blood tests.

\medterm{Dysgerminoma} A cancer of the ovary arising from germ cells, the cells that can develop into eggs. It occurs most often in adolescents and young adults and may cause pelvic pain, swelling, or an abdominal mass.

\medterm{Dysgeusia} An altered sense of taste, including reduced, distorted, or unpleasant taste. Causes include viral infection, mouth or nose disease, nerve injury, nutritional deficiency, dry mouth, medicines, and other illnesses.

\medterm{Dysgraphia} A persistent difficulty with handwriting, spelling, or expressing ideas in writing despite appropriate learning opportunities. It may reflect problems with motor skills, language, attention, or brain function and can occur with other learning disorders.

\medterm{Dyshidrosis} A skin condition causing small, intensely itchy blisters on the palms, sides of the fingers, or soles. The skin may peel, crack, or become sore; triggers include sweating, irritation, allergy, and stress.

\synonyms
Eczema, Dyshidrotic

\medterm{Dyskeratosis Congenita} An inherited disorder affecting chromosome-end maintenance. It may cause abnormal skin coloring, nail changes, white patches in the mouth, bone-marrow failure, lung or liver disease, developmental problems, and increased cancer risk.

\medterm{Dyskinesia} Abnormal movement that is often involuntary, excessive, repetitive, twisting, or irregular. It may result from neurologic disease or medicines, including long-term treatment for movement disorders, and can interfere with walking, speaking, or daily activities.

\medterm{Dyslexia} A learning disorder involving persistent difficulty recognizing words accurately or fluently, decoding them, spelling, or linking letters with speech sounds. It is not caused by low intelligence or poor motivation and may continue into adulthood.

\synonyms
Reading Disability, Specific
Specific Reading Disability

\medterm{Dyslipidemia} An unhealthy level or pattern of fats in the blood, such as high LDL cholesterol, high triglycerides, low HDL cholesterol, or a combination. It can promote artery disease and, when triglycerides are very high, pancreatitis.

\medterm{Dyslipidemias} Conditions involving abnormal blood fats, including high LDL cholesterol, high triglycerides, low HDL cholesterol, or mixed changes. They can increase the risk of artery disease and pancreatitis when triglycerides are very high.

\medterm{Dysmenorrhea} Painful menstrual periods, usually causing cramping in the lower abdomen and sometimes back pain, nausea, diarrhea, headache, or dizziness. Severe, new, or worsening pain may indicate an underlying pelvic disorder.

\medterm{Dysmenorrhoea} Painful menstrual periods. Primary pain occurs without another pelvic disease, while secondary pain results from a condition such as endometriosis, fibroids, adenomyosis, pelvic infection, or an intrauterine device.

\medterm{Dysmetria} Inability to judge or control the distance, force, or speed of a voluntary movement. A person may overshoot or undershoot a target; it commonly reflects dysfunction of the cerebellum or its connecting pathways.

\medterm{Dysorexia} An abnormal appetite, including little desire to eat, excessive hunger, or disordered eating. It may result from illness, medicines, emotional distress, metabolic disease, or an eating disorder and can affect nutrition and weight.

\medterm{Dysostoses} Inherited or developmental conditions in which one or more bones form or grow abnormally. They may affect the skull, face, spine, ribs, arms, or legs and can alter body shape, movement, growth, hearing, breathing, or other functions.

\medterm{Dyspareunia} Recurrent or persistent genital or pelvic pain before, during, or after sexual activity, especially attempted vaginal penetration. Pain may be at the vaginal opening or deeper in the pelvis and can result from dryness, infection, skin disease, pelvic-floor muscle spasm, endometriosis, inflammation, injury, or emotional factors.

\medterm{Dyspepsia} Recurrent discomfort in the upper abdomen, such as burning, early fullness, bloating, nausea, or belching. Causes include ulcers, reflux, medicines, infection, and other illness; sometimes no structural cause is found.

\medterm{Dyspeptic} Describing indigestion symptoms such as upper-abdominal discomfort or burning, early fullness, bloating, belching, or nausea, especially after eating.
\medterm{Dysphagia} See \textit{Difficulty Swallowing}.

\synonyms
Difficulty Swallowing
Trouble Swallowing

\medterm{Dysphagia Diet} Food and drinks altered in texture or thickness to make swallowing safer after a swallowing assessment. The plan should reduce aspiration risk while providing enough calories, protein, fluids, and enjoyable foods.

\medterm{Dysphagia (Geriatric)} Difficulty swallowing in an older adult. It may cause coughing or choking during meals, food remaining in the mouth, weight loss, dehydration, or aspiration and can result from stroke, dementia, muscle disease, blockage, or medicines.

\synonyms
Swallowing difficulty in older adults

\medterm{Dysphagia Rehabilitation} Therapy to improve safe swallowing, nutrition, and hydration. It may include swallowing exercises, posture or pacing strategies, food-texture changes, and caregiver education based on the cause and assessment.

\medterm{Dysphasia} Difficulty using or understanding language because of a problem in the brain's language areas. It may affect speaking, understanding, reading, or writing; many clinicians use the clearer term aphasia.

\medterm{Dysphonia} An abnormal voice, such as hoarseness, breathiness, roughness, strain, weakness, or loss of voice. Causes include infection, vocal-cord swelling or growths, nerve problems, voice overuse, reflux, and cancer; persistent symptoms need evaluation.

\medterm{Dysphoria} A state of marked emotional distress, unease, or dissatisfaction. Gender dysphoria specifically means distress related to a difference between experienced gender and sex assigned at birth; it is not itself a mental illness.

\medterm{Dysplasia} Abnormal size, shape, or arrangement of cells in a tissue. It is not cancer, but some forms—especially when severe or persistent—can increase the chance of cancer and may need monitoring or treatment.

\medterm{Dysplasia Cervical (Cervical Dysplasia)} Alternate index label for Cervical Dysplasia.

\medterm{Dysplastic Nevi} Atypical moles with unusual size, shape, color, border, or microscopic features. Most do not become melanoma, but having many can increase melanoma risk; a changing, bleeding, or suspicious mole should be examined.

\medterm{Dysplastic Nevus} An atypical mole that may be larger, unevenly colored, or have an irregular border. Most do not become cancer, but having many atypical moles can increase melanoma risk; changing lesions should be examined.

\medterm{Dyspnea} The feeling of difficult, uncomfortable, or insufficient breathing. Causes include asthma, lung or heart disease, anemia, infection, muscle weakness, anxiety, and other conditions.

\medterm{Dyspnea On Exertion} Shortness of breath that occurs or worsens with physical activity. It may result from heart or lung disease, anemia, poor fitness, or other conditions; new, severe, or rapidly worsening symptoms need prompt assessment.

\medterm{Dyspnoea} The feeling of difficult, uncomfortable, or insufficient breathing. It may result from heart or lung disease, anemia, infection, poor fitness, anxiety, or other causes; sudden or severe symptoms require emergency care.

\medterm{Dyspraxia} Difficulty planning and coordinating purposeful movements. It may affect dressing, writing, balance, speaking, or sports; in children it may be part of developmental coordination disorder.

\medterm{Dyspraxia (Developmental Coordination Disorder)} A developmental difficulty planning and coordinating movements that interferes with daily activities, learning, or motor skills. It is not explained by muscle weakness, loss of sensation, or inadequate opportunity to learn.

\synonyms
Developmental Coordination Disorder

\medterm{Dysraphism} A birth defect caused by incomplete joining of developing body structures, especially the spinal cord or spinal bones. Effects range from a small skin mark to weakness, loss of sensation, or bladder and bowel problems.

\medterm{Dysreflexia (Autonomic Dysreflexia)} A dangerous sudden rise in automatic nervous-system activity, usually after spinal-cord injury above the upper chest. A full bladder, bowel problem, skin irritation, or other stimulus may cause severe high blood pressure, headache, sweating, and slow heartbeat.

\synonyms
Autonomic Dysreflexia

\medterm{Dysrhythmia} An abnormal heart rhythm that may be too fast, too slow, or irregular. Some cause no symptoms; others cause palpitations, dizziness, fainting, chest pain, stroke, or inadequate blood flow.


\medterm{Dyssomnia} A broad or older term for a sleep disorder involving difficulty sleeping, excessive sleepiness, or abnormal sleep timing. Modern diagnosis usually specifies the problem, such as insomnia, hypersomnia, or a circadian rhythm disorder.
\medterm{Dyssomnias} Sleep disorders that mainly cause difficulty falling asleep, staying asleep, waking at the desired time, or staying awake. Causes include sleep schedules, medical conditions, medicines, mood disorders, and abnormal sleep regulation.

\medterm{Dysthymia} Historical term; see \textit{Persistent Depressive Disorder}.

\medterm{Dysthymia (Persistent Depressive Disorder)} Alternate index label for Persistent Depressive Disorder.

\medterm{Dystocia} Difficult or unusually slow labor caused by ineffective contractions, the baby’s size or position, or the mother’s pelvis or soft tissues. It can increase risks for both parent and baby and may require assisted delivery or cesarean birth.

\medterm{Dystonia} A neurological movement disorder in which changes in the brain’s movement signals make muscles tighten or contract without the person intending them to. This can cause slow, repeated or twisting movements, tremor, or abnormal and sometimes painful postures. It may affect one body part, several connected areas, or most of the body, and symptoms may be triggered by movement, stress, or tiredness. Dystonia may be inherited, have no known cause, or result from brain disease or injury, infection, poisoning, or certain medicines. It can begin at any age and may remain stable or worsen over time.

\medterm{Dystonia, Cervical} Alternate index label; see \textit{Cervical dystonia}.

\medterm{Dystonia Disorder} A disorder causing repeated or sustained involuntary muscle contractions, which produce twisting movements, tremor, or abnormal postures. Symptoms may be limited to one area or involve the whole body.

\medterm{Dystonin} A protein made from the DST gene that helps attach the internal framework of cells to other cell structures. Harmful gene changes can cause inherited nerve disease, muscle weakness, and fragile or blistering skin.

\medterm{Dystrophia Unguium} Abnormal nail growth, shape, color, or texture. Causes include injury, fungal infection, inflammation, poor nutrition, medicines, reduced blood flow, and inherited disorders; treatment depends on the cause.

\medterm{Dystrophic Calcification} Buildup of calcium salts in damaged, dead, or chronically inflamed tissue even though the blood calcium level is normal. It can occur in scars, diseased valves, tumors, or injured muscles and may impair function.

\medterm{Dystrophin} A protein that strengthens the membrane of muscle cells during contraction. Too little or abnormal dystrophin damages muscle fibers and causes Duchenne or Becker muscular dystrophy.

\medterm{Dystrophin-Glycoprotein Complex} A group of proteins that connects the internal framework of muscle cells to their surrounding support tissue and protects the cell membrane during contraction. Changes in its components can cause several inherited muscular dystrophies.

\synonyms
DGC; Dystrophin-Associated Glycoprotein Complex (DAGC)

\medterm{Dystrophy} Abnormal development, maintenance, or degeneration of a tissue or organ, often caused by an inherited or metabolic condition. Effects depend on the tissue and may include muscle weakness, poor vision, fragile structures, or loss of function.

\medterm{Dysuria} Pain, burning, or discomfort during urination. Common causes include bladder or urethral infection, vaginal or penile inflammation, stones, irritation, and sexually transmitted infections; fever, flank pain, pregnancy, or blood in urine needs medical assessment.
\synonyms
Painful Urination

\medterm{Eagle Syndrome} A rare cause of throat, face, jaw, or neck pain related to an unusually long bony projection near the skull base or a hardened nearby ligament. Pain may worsen with swallowing or turning the head.

\medterm{Eales Disease} An uncommon inflammatory disorder of small retinal veins, usually affecting otherwise healthy young adults. Blocked vessels can cause abnormal new vessels and bleeding inside the eye, leading to floaters or reduced vision.

\medterm{Ear} The organ for hearing and balance. The outer ear collects sound, the middle ear transmits vibrations through the eardrum and tiny bones, and the inner ear converts sound and movement into signals sent to the brain.

\medterm{Earache} Pain felt in the ear or nearby area. It may arise inside the ear or be referred from the teeth, jaw, throat, or neck; identifying the source requires an examination when pain persists or is severe.

\synonyms
Otalgia

\medterm{Ear Barotrauma} Ear injury caused by unequal pressure between the middle ear and surrounding air or water, often during flying or diving. It may cause pain, fullness, blocked hearing, ringing, dizziness, bleeding, or a torn eardrum.

\medterm{Ear Canal} The passage between the outer ear and eardrum. Its skin and earwax protect the middle ear, but infection, injury, foreign objects, or wax buildup can cause pain, discharge, or reduced hearing.

\medterm{Ear Cartilage} Flexible connective tissue that forms much of the outer ear and gives it shape. Injury or infection can cause pain, swelling, deformity, or damage that permanently changes the ear.

\medterm{Ear Disease} Any disorder of the outer, middle, or inner ear. It may cause pain, discharge, hearing loss, ringing, dizziness, infection, or balance problems; diagnosis requires examination of the ear and sometimes hearing or imaging tests.

\medterm{Ear Diseases} Disorders affecting the outer ear, ear canal, eardrum, middle-ear bones, pressure-equalizing tube, cochlea, or balance organs. They include infection, injury, blockage, inflammation, hearing disorders, and growths.

\medterm{Ear, Diseases Of} Disorders of the outer, middle, or inner ear. Symptoms, examination findings, causes, and treatment depend on the affected structure and may include pain, discharge, hearing loss, ringing, dizziness, or imbalance.

\medterm{Ear Drainage Culture} A laboratory test of fluid draining from the ear to identify bacteria or fungi. Results help guide treatment but must be interpreted with the examination because normal organisms, contamination, or colonization can affect the result.

\medterm{Eardrum} A thin membrane separating the ear canal from the middle ear. Sound makes it vibrate, and the vibration moves the tiny middle-ear bones; infection, pressure, or injury can tear it or reduce hearing.

\medterm{Eardrum Repair} Treatment to close a hole or repair damage in the eardrum, often using a tissue graft or other surgical technique. It may improve hearing and prevent infection, but healing is not always complete.

\medterm{Eardrum, Ruptured} Alternate name; see \textit{Perforated Eardrum}.

\medterm{Ear Examination} Inspection of the outer ear and ear canal and viewing of the eardrum, usually with an otoscope. Hearing, balance, or additional tests may be added when symptoms suggest middle- or inner-ear disease.

\medterm{Ear Foreign Body} An object lodged in the ear canal, most often in a child. It may cause pain, reduced hearing, bleeding, discharge, or injury; removal should be done with direct visualization and suitable instruments.

\synonyms
Auditory Canal Foreign Body; Aural Foreign Body

\medterm{Ear Infection} Infection or inflammation affecting part of the ear, including the outer ear, ear canal, middle ear, or inner ear. It may cause ear pain, fever, discharge, reduced hearing, dizziness, or irritability in children.

\medterm{Ear Infection, Outer Ear} Alternate name; see \textit{Swimmer's Ear}.

\medterm{Earliest Signs Of Increased Intracranial Pressure} Early warning signs may include worsening headache, vomiting, blurred or double vision, drowsiness, slowed thinking, or behavior change. Causes include bleeding, swelling, infection, a growth, or excess brain fluid; progressive symptoms need urgent assessment.

\medterm{Earlobes} The soft lower parts of the outer ears, made mainly of skin and fatty or connective tissue without cartilage. They may be affected by piercing injury, infection, cysts, allergic reactions, or tears.

\medterm{Early Afterdepolarization} An extra abnormal electrical signal during the heart’s recovery phase after each beat. It is more likely when the QT interval is prolonged and can trigger dangerous abnormal rhythms.
\synonyms
EAD; Early After-Depolarization

\medterm{Early Childhood Tooth Decay} Early-childhood tooth decay is bacterial acid destruction of primary teeth, promoted by frequent sugars, poor brushing, enamel defects, and prolonged bottle or nursing exposure.

\medterm{Early Decelerations} A gradual temporary slowing of the fetal heart rate that begins and ends with a uterine contraction, usually from pressure on the baby’s head. When isolated, it is generally reassuring, but the complete tracing and clinical situation matter.

\medterm{Early Detection} Finding a disease at an early stage, before it causes major complications or spreads. It may result from screening or prompt evaluation of symptoms and can improve treatment options for some conditions.

\medterm{Early Diagnosis} Identifying a disease before it has progressed or caused major complications. Earlier diagnosis may provide more treatment choices, but its benefit depends on the disease and the accuracy and usefulness of available tests.

\medterm{Early Endosome} A small membrane-bound compartment inside a cell that receives material brought in from the cell surface. It sorts that material for return to the surface, delivery elsewhere in the cell, or breakdown.

\medterm{Early Gastric Cancer} Stomach cancer limited to the inner lining or the layer just beneath it, whether or not nearby lymph nodes are involved. It may be found early because of screening or testing for digestive symptoms.

\synonyms
EGC

\medterm{Early Infantile Epileptic Encephalopathy} A group of severe seizure disorders beginning in infancy that can disrupt brain development. Causes may be genetic, structural, metabolic, or unknown; frequent seizures and developmental concerns require specialist evaluation.

\medterm{Early Intervention} Support provided to infants and young children with developmental delays or disabilities. It may include physical, occupational, speech, behavioral, educational, and family services to improve development, function, and participation.

\medterm{Early Puberty} Alternate name; see \textit{Precocious Puberty}.

\medterm{Early Repolarization Pattern} An electrocardiogram pattern with a small rise or notch near the end of the QRS wave, often in the lower or side leads. It is often harmless, but its meaning depends on symptoms, family history, and the rest of the tracing.
\synonyms
Early Repolarization ECG Pattern; BER

\medterm{Early Repolarization Syndrome} A rare heart-rhythm condition associated with a characteristic ECG pattern and, in some people, an increased risk of dangerous ventricular arrhythmias. Risk depends on the pattern, symptoms, and personal or family history.
\synonyms
ERS; Malignant Early Repolarization

\medterm{Early Satiety} Feeling full sooner than expected while eating, sometimes after only a small amount of food. It may result from stomach or digestive disease and should be assessed if persistent or associated with weight loss.

\medterm{Early-Stage Prostate Cancer} Prostate cancer that remains confined to the prostate or nearby tissues and has not spread to distant organs. The term includes cancers with different grades and risks of progression.

\synonyms
Early Prostatic Carcinoma; Localized Prostate Cancer

\medterm{Ear Microbiome} The community of microorganisms living on or in the ear, especially the ear canal. They interact with skin and earwax and may influence infection risk; moisture, cleaning, hearing devices, inflammation, and illness can change the community.

\medterm{Ear Neoplasm} An abnormal growth in the outer, middle, or inner ear that may be benign or cancerous. It can cause a persistent lump, pain, discharge, hearing loss, dizziness, or facial weakness and may require imaging or biopsy.

\medterm{Ear, Nose, And Throat} The medical specialty caring for disorders of the ears, nose, throat, voice, head, and neck. Specialists diagnose and treat problems such as hearing loss, sinus disease, swallowing or voice disorders, and head and neck tumors.

\medterm{Ear Speculum} A small cone-shaped tip fitted to an otoscope to direct light into the external ear canal and permit examination of the canal and tympanic membrane. Disposable and reusable types are available in different sizes.

\medterm{Ear Tag} A small, skin-covered projection on or near the outer ear, usually present from birth. It is generally harmless, but examination may be appropriate if it causes problems or occurs with other birth findings.

\medterm{Ear Thermometer} A device that estimates body temperature by measuring heat in the ear canal near the eardrum. Correct positioning is important; earwax, a small ear canal, or an ear infection can affect the reading.

\medterm{Ear Tube Insertion} A procedure placing a small ventilation tube through the eardrum to allow air into the middle ear and fluid to drain. It may help repeated infections or persistent fluid but can cause drainage, blockage, scarring, or a persistent opening.

\medterm{Ear Wax} A mixture of gland secretions and shed skin that protects and cleans the ear canal. A buildup can cause reduced hearing, fullness, ringing, itching, or discomfort; cotton swabs and inserted objects can worsen blockage or injure the ear.

\synonyms
Cerumen; Earwax

\medterm{Earwax} A waxy mixture made in the outer ear canal that lubricates and protects the skin and traps dust and microorganisms. It usually clears naturally, but buildup can block hearing or cause discomfort.

\medterm{Earwax Blockage} A buildup of wax that partly or completely blocks the ear canal. It may cause reduced hearing, fullness, ringing, itching, pain, or cough; safe removal may require a clinician.

\medterm{Earwax Removal} Removal of excess or impacted earwax by a clinician after examining the ear. Instruments, gentle suction, or irrigation may be used when safe; inserting objects at home can cause injury or push wax deeper.

\medterm{Eating Disorder} A group of mental-health conditions involving harmful eating patterns that cause distress, medical risk, or impaired functioning. Some involve intense concern about food, weight, or body shape, but avoidant/restrictive food intake disorder (ARFID) does not require body-image concerns. Examples include anorexia nervosa, bulimia nervosa, binge-eating disorder, and ARFID.

\medterm{Eating Disorders} Alternate name; see \textit{Eating Disorder}.

\medterm{Eating Disorders, Anorexia} Alternate name; see \textit{Anorexia Nervosa}.

\medterm{Eating Disorders, Binge Eating} Alternate name; see \textit{Binge-Eating Disorder}.

\medterm{Eating Disorders, Bulimia} Alternate name; see \textit{Bulimia Nervosa}.

\medterm{Eaton-Lambert Syndrome} An uncommon condition in which the immune system interferes with nerve signals to muscles. It causes weakness, often in the hips and shoulders, reduced reflexes, dry mouth, and other automatic-body problems. It can be linked to small-cell lung cancer.
\synonyms
Lambert-Eaton Myasthenic Syndrome; LEMS

\medterm{EBCT} Electron-beam computed tomography, an older rapid CT method that was used especially to measure calcium in the coronary arteries. Newer CT technology is now more commonly used for heart and vessel imaging.

\synonyms
Electron-Beam Computed Tomography; Ultrafast CT

\medterm{Ebola} Alternate name; see \textit{Ebola virus disease}.

\medterm{Ebola Hemorrhagic Fever} Alternate name; see \textit{Ebola Virus Disease}.

\medterm{Ebola Virus} A virus that causes Ebola disease. It spreads mainly through direct contact with infected blood or body fluids, usually after symptoms begin; illness may include fever, vomiting, diarrhea, bleeding, shock, and organ failure.

\medterm{Ebolavirus} A group of viruses that includes those causing Ebola disease. Infection can produce severe fever, vomiting, diarrhea, bleeding, shock, and organ damage and spreads through contact with infected body fluids.

\medterm{Ebola Virus Disease} A severe, often fatal viral disease caused by Ebola viruses and affecting humans and other primates. After an incubation period of about 2--21 days, it may begin with fever, tiredness, muscle pain, headache, or sore throat and progress to vomiting, diarrhea, rash, bleeding, shock, or organ failure. It spreads mainly when blood or other body fluids from an infected person or animal, or contaminated objects such as clothing or bedding, enter another person through broken skin or the eyes, nose, or mouth. People generally cannot spread the disease before symptoms begin, but can spread it after becoming ill.

\medterm{Ebstein Anomaly} A birth defect in which the tricuspid valve is positioned too low in the right ventricle and does not close normally. It can cause leakage of blood, enlargement of the right heart, blue coloring, heart failure, or abnormal rhythms.

\medterm{Ebstein's Anomaly} A birth defect in which the tricuspid valve is displaced downward and leaks. Effects range from no symptoms to blue coloring, heart failure, enlargement of the right heart, or abnormal heart rhythms.

\medterm{EBUS-TBNA} A procedure that uses a flexible breathing-tube camera and ultrasound to locate lymph nodes or masses beside the airways, then obtains tissue with a needle. It helps diagnose and stage lung cancer, sarcoidosis, lymphoma, and tuberculosis.

\synonyms
Endobronchial Ultrasound Transbronchial Needle Aspiration; Convex Probe EBUS

\medterm{Ecbolics} Medicines that stimulate contractions of the uterus. They may be used to start or strengthen labor or control heavy bleeding after birth; dosing requires medical supervision because excessive contractions can harm the parent or baby.

\medterm{Eccentric} Positioned away from the center. In exercise, an eccentric muscle action occurs when a muscle produces force while lengthening, such as lowering a weight; it can improve strength but may cause soreness or injury.

\medterm{Eccentric Action} Muscle contraction that produces force while the muscle lengthens, such as the controlled lowering of an object. It helps control movement and build strength but may cause muscle soreness when new or intense.

\synonyms
Eccentric Contraction

\medterm{Eccentrocyte Count} The number or proportion of red blood cells whose hemoglobin is pushed toward one side of the cell. This smear finding may indicate oxidative injury or a hemolytic disorder but must be interpreted with other blood results.

\medterm{Ecchymosis} A flat purple, blue, or brown skin mark caused by blood leaking into tissues, commonly after an injury. Easy or unexplained bruising can also result from medicines, low platelets, clotting disorders, fragile vessels, or illness.

\medterm{Eccrine Gland} A small sweat gland found throughout the skin that releases watery sweat onto the surface. It helps cool the body and is especially numerous on the palms, soles, and forehead.

\medterm{Eccrine Glands} Sweat glands distributed through the skin that release watery sweat directly onto the surface. They help control body temperature; disorders can cause excessive sweating, reduced sweating, blockage, or tumors.

\medterm{Eccrine Porocarcinoma} A rare skin cancer arising from cells of sweat-gland ducts. It may appear as a slowly enlarging lump or patch and can recur or spread; biopsy and complete removal are often needed.

\medterm{Eccrine Poroma} A usually harmless growth arising from a sweat-gland duct, often appearing as a flesh-colored, pink, or reddish lump on a hand or foot. It may bleed or be tender and can be removed if needed.

\medterm{ECG} A recording of the heart’s electrical activity made with skin electrodes. It can show heart rate, rhythm, conduction, reduced blood flow, prior heart damage, or other cardiac problems.

\synonyms
Electrocardiogram; EKG

\medterm{ECG/EKG Machine} A device that records the heart’s electrical activity through electrodes placed on the skin. A standard 12-lead recording helps assess rhythm, conduction, heart strain, reduced blood flow, and prior heart injury.

\medterm{Echinacea} A group of plants used in some herbal products for colds or claimed immune support. Evidence of benefit is inconsistent; allergic reactions and interactions with medicines can occur, especially in people with plant allergies.

\medterm{Echinocandins} A group of antifungal medicines that stop fungi building their protective cell wall. They are given mainly by injection for serious Candida infections and selected Aspergillus infections; liver reactions and infusion effects can occur.

\medterm{Echinococcal Cysts} Fluid-filled cysts caused by the larval stage of Echinococcus tapeworms, most often in the liver or lungs. They may remain silent or cause pressure, pain, breathing symptoms, or a severe allergic reaction if they rupture.

\medterm{Echinococcosis} A parasitic infection caused by larval Echinococcus tapeworms. Cystic disease usually forms cysts in the liver or lungs; alveolar disease grows invasively in the liver and may spread, requiring specialist treatment.

\medterm{Echinococcus} A group of tapeworms that can cause disease when people swallow eggs from contaminated soil, food, water, or contact with infected animals. Larvae can form cysts or invasive lesions, especially in the liver or lungs.

\medterm{Echinococcus Granulosus} A tapeworm whose larvae cause cystic echinococcosis. Dogs commonly carry the adult worm, and people become infected by swallowing eggs from contaminated material; cysts most often develop in the liver or lungs.

\medterm{Echinococcus Multilocularis} A tapeworm species whose larvae cause alveolar echinococcosis. The infection mainly affects the liver, where it grows slowly and invasively and may spread to other organs.

\medterm{Echinococcus Vogeli} A tapeworm species that causes polycystic echinococcosis, a disease producing multiple parasitic cysts mainly in the liver and sometimes other abdominal organs.

\medterm{Echinocytes} Red blood cells with many small, evenly spaced projections around their edges. They may reflect kidney disease or an inherited blood disorder, but can also be created by improper storage or preparation of the blood sample.

\synonyms
Burr Cells

\medterm{Echo} A common abbreviation for echocardiography, an ultrasound test that examines the heart’s chambers, valves, pumping action, and blood flow. The full term is preferable in patient instructions.

\medterm{Echocardiogram} An ultrasound picture of the heart that shows its chambers, valves, pumping movement, and blood flow without ionizing radiation. It helps evaluate heart failure, valve disease, congenital defects, clots, and other heart problems.

\medterm{Echocardiography} Ultrasound imaging of the heart. It shows the chambers, valves, pumping action, and blood flow; the probe usually scans through the chest and may be passed into the esophagus for closer views.

\medterm{Echoencephalography} Ultrasound imaging of the brain, historically used mainly in infants to look for enlarged fluid spaces or displacement of brain structures. CT or MRI is now usually preferred for detailed brain imaging.

\medterm{Echogenic Bowel} An ultrasound finding in which a fetus’s bowel appears unusually bright, sometimes as bright as bone. It may be a normal variation or may be associated with infection, chromosome conditions, cystic fibrosis, bleeding, or other fetal problems.

\medterm{Echogenic Intracardiac Focus} A small bright spot, usually in a muscle inside the fetal heart, seen on ultrasound. When isolated and all other findings and screening are reassuring, it is commonly a harmless variation; its significance depends on the full evaluation.

\medterm{Echography} Medical ultrasonography, which uses high-frequency sound waves to create images of tissues, organs, blood flow, or a developing fetus.

\medterm{Echolalia} Repeating words or phrases spoken by someone else, either immediately or after a delay. It can occur in autism, developmental conditions, brain disorders, or some psychiatric illnesses and must be interpreted with other symptoms.

\medterm{Echophenomena} Automatic, non-purposeful imitation or repetition of actions, words, or sounds observed in others without explicit awareness or intent. Subtypes include echolalia (vocal repetition) and echopraxia (imitation of movements), commonly observed in Tourette syndrome, catatonia, and certain neurodevelopmental or neurodegenerative disorders.

\medterm{Echopraxia} Involuntary imitation or repetition of another person’s movements. It may occur in some neurologic, developmental, or psychiatric conditions and is interpreted with other symptoms.

\medterm{Echovirus} A group of enteroviruses spread mainly by the fecal-oral route and sometimes by respiratory secretions. Infection is often mild or silent but can cause rash, respiratory illness, meningitis, encephalitis, heart inflammation, or severe newborn infection.

\medterm{E-Cigarettes And E-Hookahs} Electronic devices that heat a liquid to produce an inhaled aerosol. Many contain nicotine, which can cause dependence; the aerosol may contain irritants and toxic substances and is not harmless water vapor.

\medterm{Eclampsia} A new seizure during pregnancy or after delivery related to preeclampsia or pregnancy-related high blood pressure. It is an emergency requiring protection of breathing, seizure treatment, blood-pressure control, and urgent obstetric care.

\medterm{Eclampsia in Pregnancy} Alternate name; see \textit{Eclampsia}.

\medterm{ECMO} Alternate name; see \textit{Extracorporeal Membrane Oxygenation}.

\medterm{ECOG Performance Status} A scale from 0 to 5 describing how much a person’s illness limits daily activity. It is used to estimate prognosis and decide whether someone may tolerate certain cancer treatments; 0 means fully active and 5 means death.

\medterm{E. Coli} A group of bacteria commonly found in the intestines of people and animals. Most are harmless, but some strains cause diarrhea, urinary infection, or serious illness such as bloodstream infection or kidney injury.

\medterm{E Coli Enteritis} Infection and inflammation of the intestine caused by disease-producing strains of Escherichia coli. It may cause watery or bloody diarrhea, cramps, vomiting, and dehydration; some strains can injure the kidneys.

\medterm{E Coli Infection} Infection caused by Escherichia coli bacteria. Depending on the strain and site, it may cause diarrhea, urinary infection, wound or abdominal infection, newborn meningitis, or bloodstream infection; some strains cause kidney injury.

\medterm{E. Coli O157:H7} A strain of Escherichia coli that produces Shiga toxin and can cause severe cramps and bloody diarrhea. It may injure the kidneys, especially in children; urgent medical advice is needed for bloody diarrhea or reduced urination.

\medterm{Econazole} A topical antifungal that damages the fungal cell membrane and treats ringworm, athlete’s foot, jock itch, and some yeast skin infections. It may cause local irritation, burning, itching, or allergic contact dermatitis.

\medterm{Economy Class Syndrome} Alternate informal name; see \textit{Deep Vein Thrombosis}.

\medterm{Ecstasy} A street name for MDMA, an illegal stimulant and hallucinogenic drug. It can cause overheating, dehydration or dangerously low sodium, fast heartbeat, seizures, confusion, liver injury, and death.

\medterm{Ectasia} Abnormal widening of a blood vessel, duct, tube, or organ. Its importance depends on the structure involved; it may cause pressure, poor flow, bleeding, blockage, or rupture.

\medterm{Ecthyma} A bacterial skin infection that extends deeper than a superficial sore and causes painful, punched-out ulcers covered by thick crusts. It commonly affects the legs and may need wound care and antibiotic treatment.

\medterm{Ecthyma Gangrenosum} A serious skin infection causing rapidly developing dark or black ulcers from tissue death, often due to Pseudomonas bloodstream infection. It occurs mainly in people with severely weakened immunity or very low white-cell counts and requires emergency treatment.

\medterm{Ectoderm} The outer layer of cells in an early embryo. It develops into the skin, hair, nails, nervous system, eyes, ears, and other sense-related structures.

\medterm{Ectodermal Dysplasia} A group of inherited disorders affecting tissues such as hair, teeth, nails, skin, and sweat glands. Features may include sparse hair, abnormal teeth, reduced sweating, overheating, dry skin, and repeated infections.

\medterm{Ectodermal Dysplasias} Inherited disorders affecting structures such as hair, teeth, nails, skin, and sweat glands. Depending on the type, they may cause abnormal teeth or hair, reduced sweating, overheating, dry skin, and repeated infections.

\medterm{Ectomesenchymoma} A rare cancer containing both nerve-related and connective-tissue parts, usually occurring in children. It may appear as a growing mass or cause pain; diagnosis requires tissue examination and specialist cancer treatment.

\medterm{Ectoparasite} A parasite that lives on the outer surface of a host, such as a louse, flea, tick, or mite. It may cause itching, bites, skin disease, anemia, or transmission of infections.

\medterm{Ectopia Cordis} A rare birth defect in which the heart lies partly or completely outside the chest because the chest wall did not form normally. It is usually life-threatening and requires immediate specialist care.

\medterm{Ectopia Lentis} Displacement of the eye’s natural lens from its normal position. Injury or inherited connective-tissue disorders may cause it, leading to blurred or double vision, unequal focusing, or increased eye pressure.

\medterm{Ectopic} Located or occurring in an abnormal position. For example, an ectopic pregnancy implants outside the main cavity of the uterus, most often in a fallopian tube.

\medterm{Ectopic ACTH Syndrome} Excess cortisol production caused by a tumor outside the pituitary making adrenocorticotropic hormone (ACTH). It can cause weight gain, muscle weakness, high blood pressure, diabetes, low potassium, infections, and other features of cortisol excess.

\medterm{Ectopic Beat} An extra or early heartbeat that starts outside the heart’s usual pacemaker. It may feel like a skipped beat, flutter, or thump and can occur in healthy people or with heart disease.

\medterm{Ectopic Cushing Syndrome} Excess cortisol caused by a tumor outside the pituitary producing ACTH or, rarely, a related hormone. It may cause high blood pressure, diabetes, low potassium, muscle weakness, infections, and rapid physical changes.

\medterm{Ectopic Heartbeat} A heartbeat that starts outside the heart’s usual pacemaker. Early beats may feel like a skipped beat, flutter, or sudden thump and can occur in healthy people or with structural heart disease.

\medterm{Ectopic Heartbeats} Extra or early beats arising outside the heart’s usual pacemaker. They may feel like a skipped beat, thump, or flutter and are often harmless, but frequent episodes, fainting, chest pain, or breathlessness need assessment.

\medterm{Ectopic Kidney} A kidney located somewhere other than its usual position in the upper back of the abdomen. It may work normally or be associated with abnormal drainage, urine reflux, stones, or infection.

\medterm{Ectopic Pacemaker} An excitable cardiac focus or group of pacemaker cells located outside the sinoatrial node (in the atria, atrioventricular junction, or ventricles) that prematurely initiates cardiac depolarizations and rhythm, frequently producing ectopic beats or tachycardia.

\synonyms
Ectopic Focus; Latent Pacemaker

\medterm{Ectopic Pregnancy} A pregnancy implanted outside the main cavity of the uterus, most often in a fallopian tube. It may cause missed periods, abdominal pain, or vaginal bleeding; rupture can cause life-threatening internal bleeding and requires emergency care.

\synonyms
Tubal Pregnancy; Extrauterine Gestation

\medterm{Ectopic Thyroid} Thyroid tissue located outside its usual position in the lower front of the neck, most often at the base of the tongue. It may be the only functioning thyroid tissue and can cause swallowing, voice, or breathing problems.

\medterm{Ectopic Thyroid Tissue} Thyroid tissue that developed outside the usual thyroid location because it did not move normally during development. It is often at the tongue base and may be the only working thyroid tissue, sometimes causing swallowing or airway symptoms.

\medterm{Ectopic Tissue} Tissue located in an unusual position, usually because it developed there, failed to move normally during growth, or became implanted elsewhere. Effects depend on the tissue and whether it causes pressure, bleeding, blockage, or abnormal function.

\medterm{Ectopic Ureter} A ureter that does not enter the bladder in its normal position. It may drain into the urethra, vagina, or another abnormal site and can cause urine leakage, infection, blockage, or kidney damage.

\medterm{Ectopion} Outward turning or eversion of a body structure, most often the lower eyelid or the cervix. The effect depends on the site and may include irritation, exposure, bleeding, or abnormal discharge.

\medterm{Ectrodactyly} A birth difference in which one or more central fingers or toes are missing, leaving a deep cleft in the hand or foot. It may occur alone or with other limb or genetic differences and can affect function.

\medterm{Ectromelia} A birth defect in which one or more limbs are absent or severely underdeveloped. In laboratory medicine, the same word names a mouse poxvirus infection; these are separate uses and should not be confused.

\medterm{Ectropion} Outward turning of an eyelid, usually the lower lid, so the eye may not close or drain tears normally. It can cause tearing, dryness, irritation, and corneal injury; aging, weakness, scarring, or paralysis may contribute.

\medterm{Eculizumab} An antibody medicine that blocks part of the immune system called complement. It treats selected blood, kidney, and nerve disorders but greatly increases the risk of meningococcal and other serious infections, so vaccination and monitoring are essential.

\medterm{Eczema} A group of skin conditions that cause itchy, dry, inflamed, red, scaly, cracked, or sometimes oozing skin. The skin barrier is weakened, so the skin loses moisture and becomes easier to irritate. Atopic dermatitis is the most common type. Symptoms may worsen because of irritants, dry weather, heat, sweating, stress, or skin infection. Eczema is not contagious.

\medterm{Eczema Herpeticum} A rapidly spreading herpes simplex infection in skin affected by eczema or another barrier disorder. It causes many similar painful blisters or punched-out sores, often with fever, and needs urgent antiviral treatment because the eyes and other organs may be affected.

\medterm{Eczematous Dermatitis} A pattern of inflammatory skin disease characterized by itching, redness, scaling, dryness, oozing, or vesicles. Causes include atopic, allergic contact, irritant contact, and other forms of dermatitis.

\medterm{ED} Erectile dysfunction, difficulty getting or keeping an erection sufficient for sexual activity. It may reflect blood-vessel, nerve, hormone, medicine, psychological, or relationship factors and can signal broader health problems.

\synonyms
Erectile Dysfunction; Impotence

\medterm{Edema} Swelling caused by extra fluid in body tissues. It may leave a dent when pressed or feel firm. Common causes include heart, kidney, or liver disease, poor vein drainage, blocked lymph flow, and some medicines.

\medterm{Edema, Pulmonary} Alternate index label; see \textit{Pulmonary Edema}.

\medterm{Edema Trunk} Abnormal accumulation of fluid in the soft tissues of the chest or abdomen, producing swelling from heart, kidney, liver, lymphatic, venous, inflammatory, or medication-related causes.

\medterm{Edentulism} Complete or partial absence of natural teeth. It can impair mastication, speech,
nutrition, facial support, and quality of life and may be managed with removable, fixed,
or implant-supported prostheses.

\medterm{Edentulous} Having some or all natural teeth missing. Tooth loss may result from decay, gum disease, injury, extraction, or a developmental condition and can affect chewing, speech, nutrition, jawbone support, and appearance; dentures or implants may restore function.

\medterm{Edentulousness} Complete or partial absence of natural teeth. It may result from dental disease, trauma, developmental abnormalities, or tooth extraction and can impair chewing, speech, and oral function.

\medterm{Edetic Acid} Another name for EDTA, a chemical that binds certain metals. It is used in blood-collection tubes, laboratories, and selected treatments for heavy-metal poisoning.

\medterm{EDTA} A chemical that binds metals. It is used in some blood-collection tubes and, in selected cases, as a supervised treatment for heavy-metal poisoning; unnecessary use can harm the kidneys and lower essential minerals.

\medterm{Edwardsiella} A group of bacteria associated with water and animals. Some species, especially E. tarda, can rarely cause diarrhea, wound infection, bloodstream infection, or other illness in people.

\medterm{Edwardsiella Tarda} A bacterium found in fish, reptiles, and water that can rarely infect people. It may cause diarrhea, wound infection, gallbladder or abdominal infection, or bloodstream infection.

\medterm{Edwards Syndrome} A chromosome condition caused by an extra chromosome 18 in some or all cells. It commonly causes poor growth, developmental impairment, feeding difficulty, heart defects, and other birth abnormalities; medical needs and survival vary widely.

\synonyms
Trisomy 18; Trisomy E

\medterm{EEG} A test that records the brain’s electrical activity through sensors placed on the scalp. It helps evaluate seizures, unexplained episodes, altered awareness, certain sleep disorders, and some brain illnesses.

\synonyms
Electroencephalogram

\medterm{EFAST Examination} A rapid bedside ultrasound used after serious injury to look for internal bleeding, fluid around the heart, and air or blood around the lungs. A normal scan does not exclude every internal injury, so repeat assessment or CT may be needed.

\medterm{Efavirenz} An antiretroviral medicine that blocks HIV reverse transcriptase and is used in combination treatment for HIV infection. It can cause vivid dreams, dizziness, mood or nervous-system effects, rash, and important medicine interactions.

\medterm{Effacement} Thinning and shortening of the cervix during labor, measured as a percentage from 0 to 100. Effacement helps the cervix open and usually occurs along with dilation as contractions progress.

\medterm{Effective Dose} An estimate that combines radiation exposure to different organs into one number. It helps compare the approximate risk of imaging procedures but cannot predict an individual person’s exact risk.

\medterm{Effective Half-Life} The time needed for the amount of a radioactive substance in the body or a tissue to fall by half through both natural decay and removal by the body. It helps estimate how long radiation exposure lasts.

\medterm{Effectiveness} How well a treatment, test, or public-health intervention achieves its intended result in usual real-world practice. It differs from efficacy, which describes performance under ideal, controlled study conditions.

\medterm{Effective Reproduction Number} The average number of people one infected person infects at a particular time, taking account of immunity, behavior, and control measures. Continued spread generally occurs when the number is above 1.

\synonyms
Rt

\medterm{Effector} A cell, molecule, organ, or mechanism that carries out a response after receiving a signal. For example, an immune cell may destroy an infected cell after activation.

\medterm{Efferent} Carrying signals away from a central structure. In the nervous system, efferent pathways carry movement or automatic-control signals from the brain or spinal cord to muscles, glands, or organs.

\medterm{Efferent Arteriole} The small blood vessel that carries blood away from the kidney’s filtering unit. It helps control pressure in the filter and leads to the small vessels around the kidney tubules.

\medterm{Efferent Ductules} A series of 10 to 20 delicate, ciliated convoluted tubules that penetrate the tunica albuginea of the testis to transport spermatozoa from the rete testis into the head (caput) of the epididymis, reabsorbing the majority of testicular luminal fluid.

\synonyms
Ductuli Efferentes

\medterm{Efferent Lymphatic Vessel} A lymphatic vessel that carries lymph out of a lymph node.

\medterm{Efferent Nerve} A nerve pathway carrying signals away from the brain or spinal cord to muscles, glands, or organs. Damage can cause weakness, impaired automatic function, or loss of control in the affected area.

\medterm{Efferent Neuron} A nerve cell that carries movement or automatic-control signals from the brain or spinal cord to muscles, glands, or organs. Damage can cause weakness or impaired control of body functions.

\medterm{Efferent Pathway} A nerve route carrying signals away from the brain or spinal cord to muscles, glands, or organs. It controls movement and automatic functions such as heart rate, digestion, sweating, and bladder activity.

\medterm{Effervescent Tablets} Tablets dissolved in water before use, releasing bubbles as they form a drinkable solution. Some contain substantial sodium or other ingredients, so people with heart, kidney, or blood-pressure problems should check the label.

\medterm{Efficacy} How well a treatment produces its intended result under ideal, carefully controlled study conditions. It differs from effectiveness, which describes how well the treatment works in ordinary clinical practice.

\medterm{Effort Angina} A classic presentation of stable ischemic heart disease characterized by substernal chest discomfort, tightness, or pressure predictably provoked by physical exertion, emotional stress, or cold exposure, and consistently relieved within minutes by rest or sublingual nitroglycerin.

\synonyms
Classic Angina; Exertional Angina

\medterm{Effort Syncope} Transient loss of consciousness occurring abruptly during or immediately following physical exertion, serving as a cardinal red-flag symptom of severe cardiac outflow obstruction (such as critical aortic stenosis or hypertrophic obstructive cardiomyopathy) or exertion-induced ventricular arrhythmia.

\medterm{Effusion} Abnormal buildup of fluid in a body space, such as around a joint, lung, or heart. It may cause swelling, pain, stiffness, or breathing difficulty; the cause may be injury, infection, inflammation, heart failure, or cancer.

\medterm{Effusion-Constriction Syndrome} A rare condition in which fluid around the heart and a stiff heart covering occur together. The fluid and stiffness can prevent the heart from filling normally and may continue to cause pressure problems even after fluid is drained.

\synonyms
Effusive-Constrictive Pericarditis

\medterm{Effusion, Pleural} Alternate name; see \textit{Pleural Effusion}.

\medterm{EGD} An examination of the esophagus, stomach, and duodenum using a flexible camera passed through the mouth. It can investigate pain, swallowing problems, bleeding, or anemia and can obtain samples or treat selected lesions.

\medterm{EGFR} A protein on the cell surface that receives growth signals. Changes that overactivate EGFR can help some cancers grow; testing the tumor may guide use of treatments that block this pathway.

\medterm{EGFR Inhibition} Cancer treatment that blocks growth signals sent through the epidermal growth factor receptor. It may slow selected colorectal, lung, or head-and-neck cancers when the tumor has suitable features; rash, diarrhea, and other effects may occur.

\medterm{Egg} A nutrient-rich food containing protein, fat, vitamins, and minerals, but also a common allergen. Raw or undercooked eggs can transmit Salmonella; thorough cooking and individual allergy or dietary needs guide safe use.

\medterm{Egg Allergy} An immune reaction to egg proteins that may cause hives, swelling, vomiting, stomach pain, wheezing, or anaphylaxis. It is more common in children; reactions vary in severity and require an individualized avoidance and emergency plan.

\synonyms
Egg Hypersensitivity; Hen's Egg Allergy

\medterm{Egg-Cell} A mature female reproductive cell, also called an oocyte, that can unite with sperm during fertilization. It contains half the usual genetic material and can develop into an embryo after fertilization.

\medterm{Eggerthella} A group of bacteria that can live in the human intestine without causing disease. Rarely, especially in people with serious illness, they can cause bloodstream infection, abscesses, or other opportunistic infections.

\medterm{Egg Shell} The hard calcium-containing outer covering of an egg. In medical imaging, “eggshell calcification” describes a thin rim of calcium around a lesion or lymph node.

\medterm{Eggshell Nail} A nail that becomes unusually thin, soft, flexible, and sometimes pale or concave. Possible causes include nutritional deficiency, chronic illness, aging, or repeated trauma; persistent changes should be medically assessed.

\medterm{Egophony} A change in voice sounds heard through a stethoscope in which a spoken “E” sounds like an “A.” It may occur when lung tissue becomes dense or filled with fluid, such as in pneumonia.

\medterm{Ehlers-Danlos Syndrome} A group of inherited connective-tissue disorders that can cause unusually flexible joints, stretchy or fragile skin, easy bruising, and poor wound healing. Some types also weaken blood vessels or internal organs and require specialist care.
\synonyms
EDS; Hypermobility Syndrome

\medterm{EHR} Electronic health record: a digital collection of a person’s health information, such as diagnoses, medicines, allergies, test results, and clinical notes, that can be shared among authorized healthcare professionals.

\medterm{Ehrlichia} A group of bacteria spread mainly by ticks that infect certain white blood cells and cause ehrlichiosis. Infection may produce fever, headache, muscle aches, low blood counts, and serious illness if untreated.

\medterm{Ehrlichia And Anaplasma} Tick-borne bacteria that infect white blood cells. They can cause fever, headache, muscle aches, weakness, low blood counts, and abnormal liver tests; severe disease is more likely in older or immunocompromised people.

\medterm{Ehrlichiosis} A tick-borne infection caused by Ehrlichia bacteria. It commonly causes fever, headache, muscle aches, weakness, low white-cell or platelet counts, and raised liver enzymes; untreated infection can damage the brain, lungs, or other organs.

\medterm{Ehrlichiosis And Anaplasmosis} Tick-borne bacterial infections that may cause fever, headache, muscle aches, weakness, low blood counts, and abnormal liver tests. Prompt medical assessment and appropriate antibiotics reduce the risk of severe complications.

\medterm{Eicosanoid} A signaling substance made from certain fats that helps regulate inflammation, blood-vessel width, platelet activity, pain, and muscle contraction. Prostaglandins, leukotrienes, and thromboxanes are examples.

\medterm{Eicosanoid Degradation} Breakdown and inactivation of eicosanoids, the body’s short-lived fatty-acid signals. This limits how long and how strongly they affect inflammation, blood vessels, platelets, airways, and other muscles.

\medterm{Eicosanoid Metabolism} The body’s production, modification, and breakdown of prostaglandins, leukotrienes, thromboxanes, and related fatty-acid signals that influence inflammation, blood vessels, platelets, pain, and airway function.

\medterm{Eicosanoid Receptor} A protein on or inside a cell that responds to prostaglandins, leukotrienes, thromboxanes, or related lipid signals. Activation can change inflammation, blood-vessel tone, platelet activity, pain, or smooth-muscle contraction.

\medterm{Eicosapentaenoic Acid} The omega-3 fat EPA, found in oily fish and some supplements. It forms part of cell membranes and can be changed into signaling substances involved in inflammation and blood clotting.

\medterm{Eicosatetraenoic Acid} A 20-carbon fatty acid with four double bonds. Arachidonic acid is the best-known form and is used by cells to make prostaglandins, leukotrienes, thromboxanes, and other signaling substances.

\medterm{Eighth Cranial Nerve} The vestibulocochlear nerve, which carries hearing and balance signals from the inner ear to the brain. Damage can cause hearing loss, ringing, dizziness, imbalance, or abnormal eye movements.

\synonyms
Vestibulocochlear Nerve; Auditory Nerve

\medterm{Eikenella} A group of bacteria found in the mouth. Some species can cause deep wound, bone, joint, or bloodstream infection, especially after human bites or injuries involving a person’s teeth.

\medterm{Eikenella Corrodens} A slow-growing bacterium normally found in the mouth that can cause infection after human bites or clenched-fist injuries. It may produce abscesses, bone or joint infection, or rarely bloodstream infection.

\medterm{Eisenmenger Complex} Alternate form; see \textit{Eisenmenger syndrome}.

\medterm{Eisenmenger's Syndrome} Alternate form; see \textit{Eisenmenger syndrome}.

\medterm{Eisenmenger syndrome} A form of advanced pulmonary vascular disease caused by a longstanding congenital heart defect that initially produces excessive blood flow to the lungs. Progressive elevation of pulmonary vascular resistance leads to right-to-left or bidirectional shunting, reduced blood oxygen, and cyanosis.

\medterm{Ejaculation} The release of semen through the penis, usually during sexual climax. It involves coordinated contractions of reproductive ducts, glands, pelvic muscles, and nerves; pain, blood, or major change in ejaculation needs assessment.

\medterm{Ejaculatory Duct} One of two small tubes formed where the vas deferens joins a seminal-vesicle duct. Each passes through the prostate and opens into the urethra, carrying semen during ejaculation.

\medterm{Ejaculatory Ducts} Two small tubes within the prostate that carry semen from the vas deferens and seminal vesicles into the urethra. Blockage or inflammation can contribute to painful ejaculation, infertility, or reduced semen volume.

\medterm{Ejaculatory Dysfunction} Difficulty ejaculating, such as delayed ejaculation, absent ejaculation, painful ejaculation, or ejaculation into the bladder instead of out through the penis. Causes include medicines, nerve or pelvic disorders, diabetes, prostate procedures, and psychological factors.

\medterm{Ejection Click} A sharp heart sound heard just after the heart begins to contract, caused by opening of an abnormal aortic or pulmonary valve. It may occur with a bicuspid valve or valve narrowing and is assessed with other findings.

\medterm{Ejection Fraction} The percentage of blood pumped out of a heart chamber with each beat. It is used to describe the heart’s pumping strength, but a normal value does not rule out heart failure.

\medterm{Ejection Murmur} A midsystolic crescendo-decrescendo heart murmur that begins shortly after the first heart sound (S1), peaks in midsystole, and ends before the second heart sound (S2), produced by turbulent blood flow across an obstructed or narrowed ventricular outflow tract or semilunar valve, as seen in aortic stenosis and pulmonic stenosis.
\synonyms
Systolic Ejection Murmur; Midsystolic Murmur

\medterm{Ekbom Syndrome} An ambiguous historical term that may refer to Willis--Ekbom disease, now usually called restless legs syndrome, or to delusional infestation. The specific modern diagnosis should be used instead.

\medterm{EKG} A recording of the heart’s electrical activity made through electrodes on the skin. It can show heart rate, rhythm, conduction problems, reduced blood flow, prior heart injury, or signs of heart enlargement.

\synonyms
Electrocardiogram; ECG

\medterm{Elastance} A physical and physiological measure of the stiffness and tendency of a hollow organ or vascular structure to recoil toward its original resting dimensions after being distended by pressure, defined mathematically as the change in pressure divided by the change in volume (the inverse of compliance).

\medterm{Elastic Cartilage} Flexible supporting tissue containing many elastic fibers. It helps structures such as the outer ear, epiglottis, and pressure-equalizing ear tube keep their shape while bending.

\medterm{Elastic Fiber} A thin connective-tissue fiber that allows skin, lungs, blood vessels, and other structures to stretch and recoil. Damage or loss of elastic fibers can reduce flexibility or weaken tissues.

\medterm{Elasticity} The ability of a tissue or material to return toward its original shape after stretching or compression. It contributes to the function of skin, lungs, blood vessels, muscles, and ligaments.

\medterm{Elastics} Elastic bands used in medical or dental care, such as orthodontic bands, compression supports, wound devices, or surgical ties. Their purpose and safe use depend on the specific treatment.

\medterm{Elastic Skin} Alternate name; see \textit{Ehlers-Danlos Syndrome}.

\medterm{Elastic Tissue} Connective tissue containing fibers that let structures stretch and return toward their original shape. It helps arteries pulse, lungs expand, skin recoil, and some ligaments support movement.

\medterm{Elastin} A flexible protein in connective tissue that lets skin, lungs, blood vessels, and other structures stretch and recoil. Aging, inflammation, and chronic sun exposure can damage elastin.

\medterm{Elastofibroma} A harmless, slowly growing mass made of fatty and elastic tissue, most often beneath the lower shoulder blade in older adults. It may cause swelling or movement-related discomfort but often needs no treatment.

\medterm{Elastography} An imaging method that estimates how stiff tissue is. It is commonly used to assess liver scarring and may help evaluate breast, thyroid, prostate, or other abnormal areas without tissue removal.

\medterm{Elastomer} A flexible rubber-like material that can stretch and return toward its original shape. It is used in some medical tubing, seals, gloves, devices, and implants; biocompatibility depends on the material.

\medterm{Elastosis} Abnormal buildup, breakdown, or alteration of elastic fibers in tissue. It is commonly associated with long-term sun exposure and can occur in some degenerative, inherited, or lung disorders.

\medterm{Elastosis Perforans Serpiginosa} A rare perforating dermatosis characterized by transepidermal elimination of abnormal elastic fibers, presenting clinically as small, keratotic, umbilicated papules arranged in a serpiginous or annular configuration on the neck, face, or upper extremities, often associated with systemic disorders such as Down syndrome, Ehlers-Danlos syndrome, and pseudoxanthoma elasticum.
\synonyms
EPS; Lutz Disease

\medterm{Elbow} The joint connecting the upper arm and forearm, formed by the humerus, radius, and ulna. It allows bending and straightening and works with nearby joints to rotate the forearm.

\medterm{Elbow Extensor} A muscle that straightens the elbow, mainly the triceps and anconeus. These muscles help push, reach, lift, throw, and control lowering of the arm.

\medterm{Elbow Flexors} Muscles that bend the elbow, mainly the biceps, brachialis, and brachioradialis. They help lift the forearm and turn the palm upward; injury, pain, nerve problems, or tendon damage can weaken movement.

\medterm{Elbow Joint} A joint complex between the humerus, radius, and ulna that allows the arm to bend and straighten and the forearm to rotate. Cartilage, ligaments, muscles, and fluid-filled coverings provide stability and movement.

\medterm{Elbow Pain} Pain in or around the elbow caused by tendon injury, arthritis, bursitis, nerve pressure, fracture, instability, infection, or referred pain from the neck or shoulder. Swelling, deformity, weakness, or injury needs assessment.

\synonyms
Elbow Arthralgia; Cubital Pain

\medterm{Elbow Replacement} Surgery that removes damaged parts of the elbow and replaces them with artificial components to relieve pain and improve movement. It may be used for severe arthritis or injury; infection, loosening, and nerve or bone problems can occur.

\medterm{Elder Abuse} An intentional act or failure to provide care that harms an older person. It includes physical, emotional, sexual, or financial abuse and neglect; warning signs require sensitive assessment and safeguarding according to local law.

\medterm{Elective} Planned in advance rather than performed immediately for an emergency. An elective procedure may still be medically necessary; its timing follows assessment of benefits, risks, alternatives, and the person’s preferences.

\medterm{Elective Abortion} A planned ending of a pregnancy using medicines or a procedure, usually before an immediate medical emergency develops. The method depends on pregnancy duration, health, personal choice, and local clinical and legal requirements.

\medterm{Elective Surgery} Surgery scheduled in advance because it is not needed immediately to prevent serious harm. It may still be important or medically necessary, and timing follows discussion of benefits, risks, alternatives, and preferences.

\medterm{Electrical Alternans} A beat-to-beat change in the size or direction of ECG complexes. It can occur when the heart moves within a large fluid collection around it and may indicate cardiac tamponade, an emergency.

\medterm{Electrical Cardioversion} A controlled electrical shock delivered at a specific point in the heartbeat to restore a normal rhythm, often for atrial fibrillation or flutter. Sedation, clot risk, heart stability, and follow-up treatment must be considered.

\medterm{Electrical Conductance} The ability of a material or tissue to carry electrical current. It depends on how easily charged particles move and is important in nerve, muscle, heart, and medical-device measurements.

\medterm{Electrical Injury} Damage caused by electrical current passing through the body. Severity depends on voltage, current, duration, pathway, and skin resistance; effects may include burns, heart-rhythm problems, muscle breakdown, nerve injury, falls, or breathing failure.

\medterm{Electrical Stimulation} Use of controlled electrical currents to activate nerves or muscles. It may help selected pain, weakness, movement, or rehabilitation problems; benefits and risks depend on the device, body site, and underlying condition.

\medterm{Electrical Synapse} A direct connection between cells through tiny channels that let charged particles and electrical signals pass rapidly from one cell to another. It helps coordinate some brain, heart, and other tissue activity.

\medterm{Electric Axis of the Heart} The net average direction and magnitude of the electrical electromotive forces traveling through the ventricular myocardium during ventricular depolarization, represented on the frontal plane electrocardiogram and normally oriented between -30 degrees and +90 degrees.

\synonyms
Mean QRS Axis

\medterm{Electric Impedance} The opposition a material or tissue offers to alternating electrical current, measured in ohms. Medical devices use impedance changes to assess body composition, monitor signals, or detect contact and function.

\medterm{Electric Organ} A specialized organ in some fish that produces electrical discharges for defense, hunting, navigation, or communication. It is mainly a biological term rather than a human medical structure.

\medterm{Electric Shock} Injury caused when electrical current passes through the body. It can cause burns, muscle damage, falls, abnormal heart rhythms, breathing failure, nerve injury, or internal damage even when skin burns are small.

\medterm{Electrocardiogram} A recording of the heart’s electrical signals from skin electrodes. It can help identify abnormal rhythms, conduction problems, reduced blood flow, heart attack, chamber strain, and some effects of medicines or electrolyte changes.

\synonyms
ECG; EKG

\medterm{Electrocardiographic Gating} An imaging synchronization technique used in cardiac computed tomography, magnetic resonance imaging, and nuclear cardiology that synchronizes image acquisition with specific phases of the patient's cardiac cycle (prospective or retrospective gating) to eliminate cardiac motion artifacts.

\synonyms
ECG Gating

\medterm{Electrocardiography} The process of recording the heart’s electrical activity through skin electrodes, usually as a 12-lead ECG. It assesses rate, rhythm, conduction, chamber strain, reduced blood flow, prior injury, and some medicine or electrolyte effects.

\medterm{Electrocauterization} A procedure using controlled electrical energy to heat tissue, destroy selected tissue, or seal bleeding vessels. It can cause burns, scarring, smoke exposure, or injury to nearby tissue and requires appropriate precautions.

\medterm{Electrocautery} Use of heat produced by electrical current to cut tissue, destroy unwanted tissue, or stop bleeding during a procedure. Possible complications include burns, scarring, smoke inhalation, and damage to nearby structures.

\medterm{Electrocoagulation} A surgical technique using electrical energy to heat tissue, seal blood vessels, or destroy selected tissue. It is used to control bleeding and may cause burns, scarring, or unintended injury.

\synonyms
Surgical Diathermy; Electrocautery

\medterm{Electroconvulsive Therapy} A treatment that causes a brief controlled seizure while the person is under general anesthesia. It can rapidly treat severe depression, catatonia, or mania; temporary confusion and memory problems are common risks.

\medterm{Electrocution} Serious injury or death caused by electrical current passing through the body. It may produce burns, abnormal heart rhythms, muscle or nerve damage, breathing failure, falls, or hidden internal injury and requires emergency evaluation.

\medterm{Electrode} A conductive contact that carries electrical signals or current to or from the body, a device, or a laboratory sample. Electrodes are used in tests such as ECG and EEG and in selected treatments.

\medterm{Electrodessication} A procedure using electrical energy to dry and destroy selected superficial tissue or control bleeding. It may treat some skin lesions, often with scraping, and can cause pain, scarring, infection, or pigment change.

\medterm{Electroencephalogram} A recording of brain electrical activity through scalp electrodes. It helps evaluate seizures, prolonged altered awareness, encephalopathy, selected sleep disorders, and some intensive-care conditions; a normal result does not exclude epilepsy.

\medterm{Electroencephalography} The process of recording brain electrical activity with electrodes on the scalp. It can help identify abnormal patterns associated with seizures, sleep disorders, or diffuse brain dysfunction.

\medterm{Electrogastrogram} A recording of the stomach’s electrical activity, usually through electrodes on the abdominal skin. It measures stomach rhythm and is used mainly in specialist assessment of suspected motility disorders.

\medterm{Electrolysis} Hair removal using electrical energy delivered through a fine probe into individual hair follicles. Repeated treatment may permanently reduce hair growth, but pain, irritation, pigment changes, scarring, or infection can occur.

\medterm{Electrolyte} A mineral that carries an electrical charge when dissolved in body fluids, such as sodium, potassium, calcium, magnesium, chloride, or phosphate. Electrolytes help control nerves, muscles, heart rhythm, fluid balance, and acid-base balance.

\medterm{Electrolyte Balance} Maintenance of appropriate levels of minerals such as sodium, potassium, calcium, magnesium, chloride, and phosphate in body fluids. The kidneys, hormones, intestines, and lungs regulate them; imbalance can affect nerves, muscles, heart rhythm, and hydration.

\medterm{Electrolyte Imbalance} An abnormal level of minerals such as sodium, potassium, calcium, magnesium, chloride, or phosphate in the blood. It may cause weakness, cramps, confusion, seizures, or abnormal heart rhythms; causes include fluid loss, kidney disease, medicines, and hormone disorders.

\medterm{Electrolytes} Electrically charged minerals in body fluids, including sodium, potassium, chloride, bicarbonate, calcium, magnesium, and phosphate. They support nerve conduction, muscle contraction, fluid balance, and acid-base regulation.

\medterm{Electromagnetic Navigation Bronchoscopy} A computer-guided form of bronchoscopy that helps a flexible camera reach small areas near the outer lung. It can guide a biopsy of a lesion not visible with ordinary bronchoscopy, but bleeding, infection, or a collapsed lung can occur.

\medterm{Electromechanical Dissociation} A cardiac arrest in which the monitor shows electrical activity but the heart does not pump enough blood to produce a pulse. It requires immediate CPR and treatment of the underlying cause.

\synonyms
EMD; Pulseless Electrical Activity; PEA

\medterm{Electromyography} A test that records the electrical activity of skeletal muscles, often together with nerve-conduction studies, to evaluate peripheral nerve, nerve-root, neuromuscular-junction, and muscle disorders.

\medterm{Electron-Beam Computed Tomography} A rapid CT technique that uses an electronically steered X-ray beam to image the heart and measure calcium deposits in coronary arteries. It has largely been replaced by newer CT methods.
\synonyms
EBCT

\medterm{Electronic Cigarette} A battery-powered device that heats liquid into an inhaled aerosol, often containing nicotine and other chemicals. It can cause dependence, airway irritation, lung injury, or poisoning from liquid nicotine and is not harmless.

\medterm{Electronic Fetal Monitoring} Continuous or intermittent recording of the fetal heart rate and uterine contractions during pregnancy or labor. It helps identify patterns suggesting that the baby may not be receiving enough oxygen, but results require clinical interpretation.

\medterm{Electron Microscope} A microscope that uses electrons rather than visible light to produce highly magnified images of cells, tissues, microorganisms, or materials.

\medterm{Electron Transport Chain} A series of protein systems inside mitochondria that transfer energy from nutrients to oxygen and use it to make ATP, the cell’s main energy source. Problems in this process can impair organs that need substantial energy, such as the brain and muscles.

\medterm{Electronystagmography} A test that records eye movements with sensors around the eyes, often while the balance system is stimulated. It helps evaluate involuntary eye movements and inner-ear or balance-nerve disorders; video testing is now often preferred.

\medterm{Electrophoresis} A laboratory method that separates charged molecules by how they move through a gel or liquid in an electric field. It can analyze proteins, hemoglobin, DNA, and other substances.

\medterm{Electrophysiologic Testing} Testing that records or deliberately stimulates the heart’s electrical system, usually with catheters, to investigate abnormal rhythms and guide treatment. It is invasive and carries risks such as bleeding, infection, or rhythm disturbance.

\medterm{Electrophysiology} The study of electrical activity in living cells and tissues. Clinical electrophysiology includes tests and treatments involving heart rhythm, nerves, muscles, and the connections between them.

\medterm{Electrophysiology Study} An invasive test in which thin electrode catheters are placed in the heart to map electrical pathways and trigger or measure abnormal rhythms. It can guide ablation or other treatment; bleeding, infection, and rhythm complications are possible.

\medterm{Electroretinography} A test that records the retina’s electrical response to flashes or patterned light. It helps evaluate inherited or acquired retinal disorders when examination and imaging do not fully explain visual loss.

\medterm{Electrosurgery} Use of electrical energy during a procedure to cut tissue, destroy tissue, or control bleeding. Risks include burns, scarring, smoke exposure, unintended tissue injury, and interference with implanted electronic devices.

\medterm{Electrosurgical Unit} A generator that delivers high-frequency electrical energy through an electrode to cut tissue or stop bleeding. Safe use requires attention to current pathways, burns, nearby structures, and possible interference with implanted electronic devices.

\medterm{Element} A pure substance made of atoms with the same number of protons. Some elements, such as iron and iodine, are essential nutrients; others can be toxic, and health effects depend on dose and exposure.

\medterm{Elephant Ear Poisoning} Irritation or poisoning after eating or handling elephant-ear or taro plants, which contain sharp calcium-oxalate crystals. It can cause intense burning and swelling of the mouth and throat and may obstruct breathing.

\medterm{Elephantiasis} Severe, long-lasting swelling and thickening of skin, usually in a limb or the scrotum, caused by blocked lymph drainage. Lymphatic filariasis is one cause; other causes include repeated infection, surgery, cancer, or injury.

\medterm{Elevated Blood Pressure} Blood pressure readings that are repeatedly above the healthy range. Persistent elevation increases the risk of heart disease, stroke, kidney disease, and blood-vessel problems, although a single high reading does not establish a diagnosis.

\medterm{Elevated Hemidiaphragm} An imaging finding in which one side of the diaphragm sits higher than usual. Causes include reduced lung volume, nerve or muscle weakness, abdominal disease, or a normal variation; significance depends on symptoms and other findings.

\medterm{Elevated Liver Enzymes} Higher-than-expected blood levels of enzymes released by stressed or injured liver cells. Causes include fatty liver, medicines, alcohol, viral hepatitis, blocked bile flow, and muscle injury; the pattern and degree guide further evaluation.

\synonyms
Transaminitis; Elevated Aminotransferases

\medterm{Elevation} Upward movement or position of a body part, such as raising the arm, shoulder, jaw, or foot. The term can also describe raising an injured limb to reduce swelling.

\medterm{Elevator} A surgical or dental instrument used to lift, separate, or loosen tissue, bone, or a tooth from nearby structures. The exact design and use depend on the procedure.

\medterm{Eligible} Meeting the requirements for a treatment, clinical trial, screening program, benefit, or healthcare service. Eligibility does not guarantee that an intervention is appropriate or that it will be available.

\medterm{Elimination} Removal of waste products or medicines from the body, mainly through the kidneys, liver, lungs, skin, or intestines. Reduced elimination can make a medicine remain longer in the body and increase its effects or toxicity.

\medterm{Elimination Diets} Short-term eating plans that remove suspected foods and later reintroduce them in a planned way to identify food-related symptoms. They should be supervised when many foods are excluded or nutrition could be compromised.
\synonyms
Exclusion Diet; Diagnostic Elimination Diet

\medterm{Elimination Disorders} A group of childhood-onset disorders involving inappropriate urination or defecation, including enuresis and encopresis. These disorders are diagnosed in a developmentally inappropriate context and are distinguished from incontinence caused by a known medical, neurologic, structural, infectious, or substance-related condition.

\medterm{Elimination Half-Life} The time required for the amount of a medicine in the blood to fall by half. It helps guide dosing intervals and predicts accumulation and clearance; kidney or liver disease, age, and interactions can change it.

\medterm{Elimination Patterns} The timing, frequency, amount, and appearance of urination and bowel movements. Changes can provide clues about hydration, kidney or bowel function, continence, infection, bleeding, or other illness.

\medterm{ELISA Blood Test} A blood test that uses an enzyme-linked reaction to detect or measure an antibody, antigen, hormone, protein, or other substance. Results depend on the test design, timing, sample, and likelihood of disease before testing.

\medterm{Elixir} A sweetened liquid preparation, historically sometimes containing alcohol, used to carry a medicine. The exact ingredients and strength determine its medical effects, interactions, and poisoning risk.

\medterm{Elizabethkingia} A group of environmental bacteria that can rarely cause severe infection, especially in newborns, older adults, or people with weakened immunity. Illness may involve the blood, lungs, brain, or other organs.

\medterm{Elliptocyte Count} The number or proportion of oval or elongated red blood cells seen on a blood smear. Increased elliptocytes may occur in an inherited red-cell disorder or some anemias and must be interpreted with other results.

\medterm{Elliptocytosis} A condition in which many red blood cells are oval or elongated rather than round. Inherited forms may cause red cells to break down early and produce anemia, while some people have no symptoms.

\synonyms
Ovalocytosis

\medterm{Ellis-Van Creveld Syndrome} An inherited skeletal condition causing short limbs, extra fingers or toes, a narrow chest, and abnormal teeth or nails. Congenital heart disease is common and may affect breathing, growth, or physical function.

\medterm{Ellsworth-Howard Test} A specialized test of how the kidneys respond to parathyroid hormone. It measures changes in urine after the hormone is given and may help distinguish pseudohypoparathyroidism from other causes of low parathyroid hormone effects.
\synonyms
PTH Responsiveness Test; Parathyroid Hormone Infusion Test

\medterm{Elongin} A protein complex that helps RNA polymerase II continue copying genetic information into RNA. It contributes to gene regulation and cell function and is mainly a molecular-biology term.

\medterm{Elotuzumab} An antibody medicine that helps the immune system recognize and attack myeloma cells. It is used with other medicines for selected multiple-myeloma patients and can cause infusion reactions, infection, fatigue, or low blood counts.

\medterm{Eltrombopag} An oral medicine that stimulates the thrombopoietin receptor in bone marrow and increases platelet production. It is used for selected low-platelet conditions and can affect the liver, clotting risk, or blood counts.

\medterm{Eltrombopag Olamine} The salt form of eltrombopag, an oral medicine that increases platelet production in selected blood disorders. Liver tests, blood counts, medicine interactions, and clotting risk may require monitoring.

\medterm{Emaciation} Extreme loss of body weight, muscle, and fat caused by severe malnutrition, chronic disease, cancer, infection, or other illness. It can weaken immunity and function and requires assessment of causes and nutritional needs.

\medterm{Embolectomy} Removal of a traveling blood clot or other blockage from a blood vessel with surgery or a catheter procedure to restore blood flow. It may be urgent when a limb or organ is threatened.

\synonyms
Arterial Embolectomy; Surgical Embolectomy

\medterm{Emboli} Particles or material carried through the bloodstream that can lodge in a smaller vessel and block blood flow. Emboli may be clots, fat, air, infected material, or rarely tumor fragments.

\medterm{Embolic Stroke} A stroke caused by a clot or other material traveling to and blocking an artery that supplies the brain. Emergency treatment may restore blood flow and limit permanent brain injury.

\synonyms
Brain Embolism
Cerebral Embolism

\medterm{Embolic Stroke of Undetermined Source} A non-lacunar ischemic stroke detected on neuroimaging without proximal arterial stenosis exceeding 50 percent, a major cardioembolic source (such as atrial fibrillation), or other established specific cause despite comprehensive diagnostic evaluation.

\synonyms
ESUS

\medterm{Embolism} Blockage of a blood vessel by material that traveled through the bloodstream from another site. The material is most often a clot but may be fat, air, infected material, or rarely tumor tissue.

\medterm{Embolism, Pulmonary} Alternate name; see \textit{Pulmonary Embolism}.

\medterm{Embolization} Deliberate blockage of a blood vessel or abnormal vessel connection using a catheter and materials such as coils, particles, or plugs. It can control bleeding, reduce tumor blood flow, or treat selected aneurysms and malformations.

\medterm{Embolization Therapy} An image-guided procedure that blocks selected blood vessels with coils, particles, plugs, or other materials. It can control bleeding, reduce blood supply to a tumor, treat aneurysms, or close abnormal vessel connections.

\medterm{Embolus} Material traveling through the bloodstream that may lodge in a smaller vessel and block it. It may be a blood clot, fat, air bubble, infected material, or rarely tumor tissue.

\medterm{Embryo} A developing human from fertilization through the end of the eighth week after fertilization. From the ninth week until birth, the developing human is called a fetus.

\medterm{Embryogenesis} The process by which an embryo develops after fertilization through cell division, movement, specialization, and formation of organs and body structures.

\medterm{Embryology} The study of development from fertilization through the embryonic and fetal stages, including how organs and body structures form, grow, and become specialized.

\medterm{Embryoma} A rare tumor made of primitive cells resembling early embryonic tissue. The term can refer to different tumor types, so diagnosis depends on the organ involved and microscopic examination.

\medterm{Embryonal Carcinoma} An aggressive germ-cell tumor composed of primitive malignant cells, occurring in the testis, ovary, or other sites and often treated with chemotherapy and surgery.

\medterm{Embryonal Rhabdomyosarcoma} A malignant skeletal-muscle-lineage tumor that commonly affects young children and may arise in the head and neck, urinary tract, or reproductive tract.

\medterm{Embryonal Tumor Of The Central Nervous System} A highly cellular malignant pediatric CNS tumor with primitive or embryonal morphology. Diagnosis depends on integrated histologic and molecular classification, and treatment usually combines surgery, radiation when appropriate, and chemotherapy.

\medterm{Embryonal Tumors} Cancers that arise from cells resembling developing embryonic tissue, often occurring in children and commonly affecting the nervous system.

\medterm{Embryo Neoplasm} A tumor arising from primitive cells resembling developing embryonic tissue, usually in a child. Diagnosis and treatment depend on the tissue of origin, microscopic and molecular features, location, and spread.

\medterm{Embryonic} Relating to an embryo, the developing human or animal during the early period after fertilization. In humans, the embryonic period ends at about eight weeks after fertilization, followed by the fetal period.

\medterm{Embryonic Cell} A cell of an early embryo capable of division and differentiation into specialized tissues; embryonic stem cells have broad developmental potential.

\medterm{Embryonic Induction} A developmental process in which one group of embryonic cells signals nearby cells to develop into a particular tissue or structure.

\medterm{Embryonic Mosaic} A developing embryo containing two or more groups of cells with different genetic makeups, produced by a change after fertilization. Effects depend on the cells involved and the proportion of each cell group.

\medterm{Embryonic Stage} The early developmental period after fertilization during which the basic body plan and major organs begin to form. In human pregnancy it lasts through about eight weeks after fertilization.

\medterm{Embryo Transfer} A fertility procedure in which one or more laboratory-grown embryos are placed into the uterus through a thin catheter. It may be done soon after fertilization or after freezing and thawing; implantation and pregnancy are not guaranteed.

\synonyms
ET; Intrauterine Embryo Transfer

\medterm{Emergence Agitation} Temporary confusion, restlessness, fear, or inconsolable crying while waking from anesthesia. It is more common in children and usually resolves with reassurance and monitoring, but pain, low oxygen, or another complication must be excluded.

\medterm{Emergency Airway Puncture} An emergency procedure that creates an airway through the neck when severe upper-airway blockage prevents breathing and a breathing tube cannot be placed. It is a temporary, high-risk lifesaving measure performed only by trained professionals.

\medterm{Emergency Care} Immediate assessment and treatment for a serious or time-sensitive illness or injury. Clinicians first protect breathing, circulation, brain function, and control of bleeding, pain, or infection, then arrange continued treatment or transfer.

\medterm{Emergency Code} A hospital alert that summons a designated response to an urgent event, such as cardiac arrest, fire, or security threat. The meaning of each code varies by institution and should be confirmed locally.

\medterm{Emergency Contraception} Birth control used after unprotected sex or contraceptive failure to reduce the chance of pregnancy. Options include an oral medicine or copper intrauterine device; effectiveness depends mainly on the method and how soon it is used.

\medterm{Emergency Department} The hospital service that evaluates and treats sudden illness and injury. Staff prioritize people by seriousness, stabilize life-threatening problems, diagnose causes, provide initial treatment, and arrange discharge, admission, or transfer.

\medterm{Emergency Medicine} The medical specialty focused on rapid assessment, stabilization, diagnosis, and initial treatment of sudden illness and injury at any age. It includes resuscitation, urgent procedures, poisoning care, and decisions about admission or transfer.

\medterm{Emergency Procedures} Planned actions for responding safely to urgent events such as cardiac arrest, airway blockage, contamination, equipment failure, or radiation exposure while protecting patients, staff, and the public.

\medterm{Emerging Infectious Disease} An infection that is newly recognized, has recently increased, or is spreading into a new population or region. Causes may include animal-to-human transmission, genetic change, travel, climate, or loss of effective control.

\medterm{Emery-Dreifuss Muscular Dystrophy} A rare inherited muscle disease causing early tightness around the elbows, ankles, and neck followed by slowly progressive weakness. Heart conduction problems, abnormal rhythms, and cardiomyopathy may occur, so regular heart monitoring is essential.

\medterm{Emesis} Forceful emptying of stomach contents through the mouth, commonly called vomiting. It may result from infection, medicines, poisoning, pregnancy, blockage, brain disease, or other illness and can cause dehydration.

\medterm{Emetic} A substance or medicine that causes vomiting. Emetics are generally not used for poisoning because vomiting can cause aspiration, worsen burns or injury, and delay safer treatment.

\medterm{Emetics} Substances that induce vomiting. They are generally avoided for poisoning because vomit can enter the lungs or cause additional injury, and vomiting does not reliably remove a dangerous substance.

\medterm{EMG} A test that records electrical activity in muscles, usually with nerve-conduction testing. It helps investigate weakness, numbness, nerve injury, trapped nerves, muscle disease, and disorders of the nerve-muscle connection.

\synonyms
Electromyography

\medterm{Eminentia} An anatomical term for a raised or prominent part of a body structure. It may provide attachment for a muscle or ligament, form part of a joint, or serve as a landmark during examination or surgery.

\medterm{Emission} The discharge or release of a substance, usually a fluid. In male reproduction, emission is the movement of semen into the urethra before ejaculation.

\medterm{Emmenagogue} A substance traditionally claimed to start or increase menstrual bleeding. Evidence and safety vary, and some products can cause poisoning, miscarriage, or interactions; missed or abnormal periods need medical evaluation rather than self-treatment.

\medterm{Emmetropia} The usual focusing state of the eye in which incoming parallel light focuses on the retina without extra effort. Nearsightedness, farsightedness, or astigmatism causes light to focus away from the retina.

\medterm{Emodin} A plant chemical studied for laxative, anti-inflammatory, and anticancer effects. Evidence and safety are not established for routine treatment, and concentrated products may cause diarrhea, dehydration, or medicine interactions.

\medterm{Emollient} A substance applied to skin to soften and moisturize it, reduce water loss, and strengthen the protective barrier. Emollients are commonly used for dry skin and eczema; greasy products can make floors or clothing slippery.

\medterm{Emotion} A feeling state involving subjective experience, brain and body responses, behavior, and interpretation of events. Emotions are normal, but persistent or extreme changes may occur with mental or physical illness.

\medterm{Emotional Adjustment} The process of adapting thoughts, feelings, and behavior to illness, loss, stress, or a major life change. Difficulty adjusting may cause persistent distress or impaired functioning and may benefit from support or treatment.

\medterm{Emotional Dependency} Reliance on another person for reassurance, self-worth, or emotional stability to a degree that limits independence or causes distress. It is not a diagnosis by itself and may reflect anxiety, low self-esteem, trauma, or relationship difficulties.

\medterm{Emotional Disorders} Mental-health conditions involving persistent or severe changes in mood or emotional responses that cause distress or interfere with daily life, relationships, school, or work. Examples include depressive, anxiety, and trauma-related disorders.

\medterm{Emotional Dysregulation Disorder} Alternate name; see \textit{Borderline Personality Disorder}.

\medterm{Emotional Eating} Eating in response to stress, boredom, sadness, anger, or other feelings rather than mainly physical hunger. Occasional emotional eating is common, but frequent episodes can affect health and may benefit from behavioral or psychological support.

\medterm{Empathy} The ability to understand another person’s feelings and point of view. It supports respectful communication and compassionate care but does not require agreeing with the person or experiencing the same feelings.

\medterm{Emphysema} A type of chronic obstructive pulmonary disease in which air sacs are permanently damaged and enlarged, reducing air exchange. Smoking is a common cause; breathlessness, cough, and reduced exercise capacity may develop gradually.

\medterm{Emphysematous Aortitis} A rare, life-threatening infection of the aortic wall that produces gas within the artery, usually caused by organisms such as \textit{Clostridium septicum}. It may cause severe back or abdominal pain, sepsis, bleeding, or an abnormal connection with the intestine.

\medterm{Emphysematous Cholecystitis} A severe gallbladder infection in which gas-forming bacteria produce gas in the gallbladder wall or cavity. It is more common in people with diabetes and can rapidly cause tissue death, perforation, or sepsis.

\medterm{Emphysematous Hepatitis} A rare, life-threatening gas-forming infection of liver tissue, usually caused by bacteria in people with diabetes or weakened immunity. It can cause fever, abdominal pain, sepsis, and liver failure.

\medterm{Emphysematous Prostatitis} A rare, severe prostate infection in which gas-forming bacteria produce gas within prostate tissue. It may cause fever, pelvic or perineal pain, difficulty urinating, urinary retention, and sepsis.

\medterm{Emphysematous Pyelonephritis} A severe kidney infection in which gas-forming bacteria destroy kidney tissue. It occurs mainly in people with diabetes and can cause fever, flank pain, shock, and sepsis; urgent drainage, antibiotics, and sometimes surgery are needed.

\medterm{Empirical} Based on observation, measurement, or experience rather than theory alone. Empirical treatment begins before the exact cause is confirmed and is revised when test results or the clinical course provide more information.

\medterm{Empirical Treatment} Treatment started on the basis of symptoms and likely causes before a diagnosis is confirmed. It should be reassessed when test results or the response to treatment becomes available.

\medterm{Empiric Antibiotic Therapy} Starting an antibiotic before the exact germ and its medicine sensitivity are known. The first choice is based on the symptoms, likely source of infection, local resistance patterns, and the person’s health; it may be changed when test results return.

\medterm{Empiric Risk} An estimate of risk based on observed data from a population or clinical setting. Its usefulness depends on how similar the studied population, exposure, and outcome are to the person or situation being assessed.

\medterm{Employee} A person who works for an employer under an agreement. In occupational health, assessment may consider workplace exposures, safety, physical or mental function, and ability to perform duties safely.

\medterm{Empty Calories} A nontechnical term for foods or drinks that provide energy but little protein, fiber, vitamins, or minerals, such as many sugary drinks and highly processed snacks. Frequent intake can displace more nutritious foods.

\medterm{Empty Can Test} A clinical examination manoeuvre used to evaluate the integrity of the supraspinatus muscle and tendon. The patient's arms are abducted ninety degrees in the scapular plane with full internal rotation (thumbs pointing downward), and the examiner applies downward pressure against active resistance; weakness or pain indicates supraspinatus tendinopathy or tear.

\synonyms
Jobe Test; Supraspinatus Test

\medterm{Empty Nose Syndrome} A poorly understood, controversial condition that can follow extensive surgery on the nasal turbinates, in which the nose feels blocked even though the airway is open. People may describe a dry, suffocating or breathless feeling, inability to sense airflow, crusting, headache, fatigue, and poor concentration, with little relief from decongestants. Findings of obstruction may be absent, so diagnosis is clinical; management is supportive and includes moisturizing, saline care, and consultation with clinicians experienced with the condition.

\medterm{Empty Sella Syndrome} A condition in which fluid presses the pituitary gland flat within its bony space. It is often found incidentally and causes no symptoms, but some people have headaches, vision problems, or reduced pituitary-hormone production.

\medterm{Empyema} A collection of pus in a normally sterile body space, most often infected fluid around a lung. It can cause fever, chest pain, cough, breathlessness, and sepsis and usually requires antibiotics plus drainage.

\medterm{Empyema Necessitans} Extension of infected pus from the space around a lung through the chest wall, forming a painful swelling. It is a serious complication requiring imaging, antibiotics, and drainage.

\medterm{Empyema Necessitatis} The direct extension and burrowing of an intrapleural empyema through the parietal pleura and chest wall tissues into the subcutaneous tissue and skin, creating an extrapleural abscess, classically caused by Mycobacterium tuberculosis, Actinomyces israelii, or neglected bacterial empyema.
\synonyms
Empyema Necessitatis; Pleurocutaneous Fistula

\medterm{EMR} Electronic medical record: a digital record of health information maintained within a particular healthcare organization or practice. It may include diagnoses, notes, medicines, allergies, test results, and treatment plans.

\medterm{Emtricitabine} An antiretroviral medicine that blocks HIV reverse transcriptase and is used with other medicines to treat or prevent HIV infection. Kidney function and hepatitis B status are important because stopping it can worsen hepatitis B.

\medterm{Emulsion} A mixture of two liquids that normally do not blend, such as oil and water, with one dispersed in the other. Emulsions are used in some medicines, nutrition products, cosmetics, and laboratory preparations.

\medterm{Enalapril} A medicine that blocks formation of angiotensin II, relaxing blood vessels and reducing strain on the heart. It treats high blood pressure and heart failure but can cause cough, high potassium, kidney problems, or rarely facial swelling.

\medterm{Enalaprilat} The active injectable form related to enalapril that blocks angiotensin II formation and lowers blood pressure. It is used in selected settings and can cause low blood pressure, high potassium, kidney problems, or facial swelling.

\medterm{Enamel} The hard, mineral-rich outer covering of a tooth crown. It protects the sensitive inner tooth but cannot repair itself after major damage; acids and bacteria can cause erosion and decay.

\medterm{Enamel Hypoplasia} Incomplete formation of tooth enamel while a tooth is developing, leaving it thin, grooved, pitted, or discolored. Affected teeth are more vulnerable to sensitivity, wear, and decay.

\medterm{Enamel Organ} The developing tissue of a tooth that produces enamel and helps shape the crown. It is present during tooth formation and later contributes to the protective covering around the developing tooth.

\medterm{Enamel Pearl} A small lump of enamel on a tooth root, often near where the roots divide. It can make cleaning difficult and interfere with gum attachment, increasing the risk of localized gum disease.

\medterm{Enanthem} A rash, spot, ulcer, or other lesion on a mucous membrane, especially inside the mouth. It may occur with infection, inflammation, an allergic or medicine reaction, or a systemic illness.

\medterm{Enanthema} A rash or lesion on a moist membrane lining the mouth, throat, genitals, or another body surface. Infection, inflammation, medicines, or systemic disease may cause it.

\medterm{Enarthrosis} A ball-and-socket joint that allows movement in many directions, including bending, straightening, rotation, and movement away from or toward the body. The hip and shoulder are examples.

\medterm{Encapsulated} Enclosed by a capsule or clearly defined layer of tissue. This may describe a cyst, growth, infection, or organ; being encapsulated does not by itself prove that a lesion is harmless or cancerous.

\medterm{Encapsulated Peritoneal Sclerosis} A rare serious complication of long-term peritoneal dialysis in which scar tissue forms around the intestines. It can block the bowel and cause repeated pain, vomiting, and poor nutrition.

\synonyms
Sclerosing Encapsulating Peritonitis; EPS

\medterm{Encapsulating Peritoneal Sclerosis} A rare condition in which thick scar tissue surrounds the intestines and can cause repeated blockage, pain, vomiting, and poor nutrition. It is most associated with long-term peritoneal dialysis but can have other causes.

\medterm{Encephalitis} Inflammation of brain tissue, often caused by an infection or an abnormal immune reaction. It may cause fever, headache, confusion, seizures, behavior change, weakness, or movement problems and requires urgent medical care.

\medterm{Encephalitis Lethargica} A rare inflammatory brain disorder described mainly during the early twentieth century, causing sleepiness, abnormal movements, eye-movement problems, and other neurologic symptoms. Its cause remains uncertain.

\medterm{Encephalitis Virus} A virus capable of infecting the brain and causing encephalitis. It may spread through mosquitoes, ticks, animals, respiratory secretions, food, or other routes depending on the virus.

\medterm{Encephalitozoon} A group of tiny spore-forming parasites that can infect the intestine, eyes, kidneys, or several organs. Serious disease occurs mainly in people with weakened immunity.

\medterm{Encephalocele} A birth defect in which the coverings of the brain and sometimes brain tissue protrude through an opening in the skull. It may cause a visible head swelling and can be associated with seizures, fluid buildup, developmental problems, or other brain abnormalities.

\medterm{Encephalomalacia} Softening or loss of brain tissue after injury, stroke, bleeding, infection, or reduced blood flow. The resulting permanent damage may cause seizures, weakness, movement problems, memory changes, or other neurologic effects.

\medterm{Encephalomyelitis} Inflammation affecting both the brain and spinal cord, caused by infection or an abnormal immune response. It may cause confusion, seizures, weakness, numbness, balance problems, or bladder changes and requires urgent neurologic assessment.

\medterm{Encephalopathy} A disorder that broadly affects brain function, causing confusion, drowsiness, personality change, difficulty thinking, or coma. Causes include infection, toxins, organ failure, medicines, metabolic problems, seizures, and reduced oxygen.

\medterm{Encephalotrigeminal Angiomatosis} A birth condition involving abnormal blood vessels in the skin of the face and the coverings of the brain. It may cause a port-wine stain, seizures, weakness, learning problems, or glaucoma, but symptoms vary.

\synonyms
Sturge-Weber Syndrome; SWS

\medterm{Enchondroma} A usually harmless cartilage growth inside a bone, often in the hands or long bones and found by chance. Pain, enlargement, fracture, or unusual imaging features may require evaluation for a more serious tumor.

\medterm{Enchondromatosis} A disorder in which multiple cartilage growths develop inside bones. It can cause uneven growth, limb deformity, pain, or fractures; a small number of lesions may become cancerous and require monitoring.

\medterm{Encoding} The process by which sensory or other information is changed into a form the brain can store in memory. Attention, sleep, stress, illness, and brain injury can affect encoding.

\medterm{Encopresis} Repeated passage of stool in inappropriate places by a child at an age or developmental level when bowel control is expected. It is often related to long-term constipation and retained stool but may have other medical, developmental, or emotional contributors.

\medterm{Endangiitis} Inflammation of the inner lining of a blood vessel. Swelling or scarring may narrow the vessel and reduce blood flow; the effects depend on the cause and the vessels involved.

\medterm{Endarterectomy} Surgery that removes fatty plaque and diseased inner lining from an artery to improve blood flow, most often a carotid artery in the neck. Possible complications include bleeding, stroke, heart attack, or nerve injury.

\medterm{Endarteritis} Inflammation of the inner lining of an artery that can narrow the vessel and reduce blood flow. Causes include infection, autoimmune disease, injury, and other vascular disorders.

\medterm{Endemic} Regular presence of a disease or infectious agent in a particular population or region at an expected baseline level. Endemic does not mean that every person is infected or that cases cannot increase.

\medterm{Endemic Diseases} Diseases that remain present at a predictable level in a particular population or region. Their frequency is shaped by local environment, immunity, behavior, exposures, healthcare access, and prevention measures.

\medterm{Endemic Fungi} Fungi naturally found in particular geographic regions whose spores can cause infection when inhaled. Disease may affect the lungs and sometimes spread to other organs, especially in people with weakened immunity.

\medterm{Endemic Goiter} Thyroid enlargement occurring commonly in a population, traditionally because of insufficient iodine intake. It may occur with normal, low, or high thyroid function depending on the cause.

\medterm{Endergonic} Describing a chemical reaction that requires an input of energy. In cells, energy from food or another reaction drives processes such as building proteins, DNA, or other complex molecules.

\medterm{Ending Pregnancy With Medicines} Ending a pregnancy with prescribed medicines that cause the uterus to empty. The appropriate method depends on pregnancy duration, health, symptoms, and local clinical requirements; follow-up confirms completion and checks for complications.

\medterm{Endobronchial Foreign Body} An object lodged in a bronchus, partly or completely blocking an airway. It may cause sudden coughing, one-sided wheezing, collapse of part of the lung, or pneumonia and often requires urgent bronchoscopy for removal.

\medterm{Endobronchial Tumor} A growth inside a bronchus that can narrow the airway and cause cough, wheezing, coughing blood, lung collapse, or repeated pneumonia. It may be benign or cancerous; diagnosis usually requires imaging and tissue sampling.

\medterm{Endobronchial Ultrasound} A procedure combining bronchoscopy with ultrasound to examine structures beside the airways and guide needle sampling of lymph nodes or masses. It is commonly used to diagnose and stage lung cancer and other chest diseases.

\medterm{Endobronchial Valve} A one-way valve placed through a bronchoscope into a selected airway to reduce air trapping in severe emphysema. It lets air and mucus leave but limits new air entering; pneumothorax and infection are possible complications.

\medterm{Endocannabinoid} A signaling molecule made naturally by the body that activates cannabinoid receptors. Endocannabinoids help regulate pain, appetite, mood, memory, movement, and immune activity.

\medterm{Endocardial Cushion Defect} A birth defect in which the central wall of the heart and the valves between its chambers do not form normally. It can allow abnormal blood flow and cause heart failure, high pressure in lung vessels, poor growth, or breathing difficulty.

\medterm{Endocardial Fibroelastosis} Thickening of the inner lining of the heart from excess fibrous and elastic tissue, usually in infants or children. It can stiffen the heart or weaken its pumping and may lead to heart failure.

\medterm{Endocarditis} Infection or inflammation of the heart’s inner lining, usually involving a valve. It can cause fever, fatigue, a new or changed murmur, stroke or other emboli, valve destruction, or heart failure.

\synonyms
Infective Endocarditis; Valvular Endocarditis

\medterm{Endocarditis Prophylaxis} Preventive antibiotics given before selected dental procedures to people at highest risk of serious infective endocarditis complications. They are not needed for everyone; eligibility depends on the heart condition and procedure.

\medterm{Endocarditis Surgical Indications} Reasons surgery may be needed for endocarditis include heart failure from valve damage, uncontrolled infection or abscess, infection of a prosthetic valve, fungal infection, or a high risk of emboli. Decisions require a specialist heart team.

\medterm{Endocardium} The thin inner lining of the heart chambers and valves. It provides a smooth surface for blood flow and continues into the inner lining of blood vessels; inflammation or infection can damage it.

\medterm{Endocervical} Relating to the inner canal or lining of the cervix, the lower part of the uterus that connects the uterine cavity with the vagina.

\medterm{Endocervical Culture} A laboratory test that attempts to grow bacteria or fungi from a sample collected from the inner cervix. It may help investigate cervicitis or selected infections, although molecular tests are preferred for many sexually transmitted infections.

\medterm{Endocervical Gram Stain} A microscope test that stains cells and organisms from the inner cervix. It may show inflammation or some bacteria but cannot by itself reliably diagnose every sexually transmitted infection.

\medterm{Endocrine} Relating to hormones released into the bloodstream by glands or specialized cells. Endocrine organs include the thyroid, pituitary, adrenal glands, pancreas, and reproductive glands.

\medterm{Endocrine Disease} A disorder of a hormone-producing gland or hormone action. It may affect growth, metabolism, reproduction, stress responses, blood pressure, or fluid and electrolyte balance.

\medterm{Endocrine Disorders} Diseases in which hormone-producing glands make too much or too little hormone, or the body does not respond to hormones properly. They may affect growth, metabolism, reproduction, mood, blood pressure, and fluid balance. Examples include thyrotoxicosis, hypothyroidism, hypogonadism, Cushing syndrome, and Addison's disease.

\medterm{Endocrine Disrupting Chemicals} Chemicals that can interfere with hormone production, transport, breakdown, or action. Some exposures are associated with reproductive, developmental, thyroid, metabolic, or cancer outcomes, but risk depends on the substance, dose, and timing.

\medterm{Endocrine Disruptor} A chemical that interferes with hormone systems. Possible effects include altered development, reproduction, thyroid function, or metabolism; health risk depends on the chemical, exposure level, duration, and timing.

\medterm{Endocrine Finding} A symptom, examination result, imaging result, or laboratory result that may suggest abnormal hormone production, hormone action, or an endocrine-gland disorder.

\medterm{Endocrine Function} The production, release, transport, and action of hormones that regulate metabolism, growth, reproduction, stress responses, blood pressure, and fluid or electrolyte balance.

\medterm{Endocrine Gland} An organ that releases hormones directly into the bloodstream. Its chemical messages regulate growth, metabolism, reproduction, stress responses, blood pressure, and other body functions.

\medterm{Endocrine Glands} Ductless organs that release hormones into the bloodstream to regulate distant tissues. They include the pituitary, thyroid, adrenal glands, pancreas, and reproductive glands.

\medterm{Endocrine Hypertension} High blood pressure caused by excess or abnormal hormone activity. Causes include primary aldosteronism, excess cortisol, or excess adrenaline-like hormones from an adrenal or related tumor.

\medterm{Endocrine Neoplasm} A benign or cancerous growth arising in a hormone-producing gland or cell. It may release too much hormone or cause symptoms by pressing on nearby structures; evaluation depends on its location and behavior.

\medterm{Endocrine Pancreas} The hormone-producing part of the pancreas, made mainly of pancreatic islets. Its cells release insulin, glucagon, somatostatin, and other hormones that help regulate blood glucose and digestion.

\medterm{Endocrine System} The network of glands and hormone-producing cells and tissues that release chemical messengers into the blood. It includes structures such as the hypothalamus, pituitary, thyroid, parathyroids, adrenal glands, pancreatic islets, and gonads; its hormones regulate growth, metabolism, reproduction, stress responses, sleep, blood pressure, fluid balance, and other aspects of homeostasis.

\medterm{Endocrine Testing} Blood, urine, saliva, stimulation, suppression, and sometimes imaging tests used to assess hormone production and action. Results depend on timing, medicines, age, sex, illness, and the test method.

\medterm{Endocrine Therapy} Treatment that lowers, blocks, or changes hormone signals to control hormone-sensitive cancer. It may be used for breast, prostate, or other cancers; effects depend on the medicine and can include hot flashes, sexual changes, or bone loss.

\medterm{Endocrine Tissues} Cells or groups of cells that produce and release hormones into the bloodstream or nearby tissues. They are found in glands and in organs such as the pancreas, ovaries, testes, kidneys, and intestine.

\medterm{Endocrinologist} A physician who diagnoses and treats disorders of hormones, endocrine glands, and metabolism, including diabetes, thyroid disease, adrenal disease, osteoporosis, and reproductive hormone problems.

\medterm{Endocrinologists} Physicians specializing in hormone, endocrine-gland, and metabolic disorders, including diabetes, thyroid disease, adrenal disease, osteoporosis, and reproductive hormone problems.

\medterm{Endocrinology} The medical specialty concerned with hormones, endocrine glands, and metabolism. It diagnoses and treats conditions affecting diabetes, thyroid function, growth, reproduction, bones, adrenal glands, and fluid or electrolyte balance.

\medterm{Endocrinopathy} A disorder of an endocrine gland or hormone system that causes too much, too little, or abnormal hormone action. Effects depend on the hormone and may involve growth, metabolism, reproduction, mood, or fluid balance.

\medterm{Endocytosis} The process by which a cell folds its outer membrane around material and brings it inside in a small sac. Cells use it to absorb nutrients, remove particles, and control surface signals.

\medterm{Endoderm} The innermost of the three early embryonic cell layers. It develops into the linings of the digestive and respiratory systems and contributes to organs such as the liver, pancreas, thyroid, and bladder.

\medterm{Endodontics} The dental specialty concerned with the soft tissue inside teeth, root canals, and tissues around tooth roots. Endodontists diagnose and treat tooth infection, inflammation, injury, and related pain.

\medterm{Endodontist} A dentist with advanced training in diagnosing and treating disease or injury of the tooth’s inner tissue, root canals, and tissues around tooth roots. They commonly perform complex root-canal treatment.

\medterm{End-Of-Life Care} Care for a person nearing death that focuses on comfort, symptom relief, dignity, communication, and personal goals. It also supports family and caregivers and may include hospice or palliative services.

\medterm{Endogenous} Originating within the body or produced by its own cells or tissues, rather than introduced from outside. Endogenous hormones, chemicals, infections, or substances may be contrasted with external sources.

\medterm{Endogenous Opiates} Opioid-like substances naturally made by the body, including endorphins, enkephalins, and dynorphins. They act on opioid receptors and influence pain, stress, mood, reward, movement, and other nervous-system functions.

\medterm{Endogenous Opioids} Opioid-like peptides made naturally by the body, including endorphins, enkephalins, and dynorphins. They activate opioid receptors to influence pain, stress, mood, reward, movement, and hormone regulation.

\synonyms
Endorphins and Enkephalins; Endogenous Opioid Peptides

\medterm{Endogenous Retrovirus} A piece of viral DNA inherited in the human genome from an ancient virus that infected a reproductive cell. Most such sequences are inactive, although some can influence gene control or normal development.

\medterm{Endoleak} Continued blood flow into the weakened part of an artery after endovascular aneurysm repair. Some types indicate an inadequate seal or device problem and may need urgent repair; others are monitored according to sac growth and risk.

\medterm{Endolimax} A group of amoebas that can live in the intestine. They are usually harmless colonizers rather than causes of disease, and a stool finding should be interpreted with symptoms and other test results.

\medterm{Endolimax Nana} An amoeba that can live harmlessly in the intestine and appear in stool tests. It usually does not cause illness; its presence may indicate exposure to contaminated food or water and should not alone explain digestive symptoms.

\medterm{Endolymph} Fluid inside the membranous part of the inner ear that helps sensory organs detect sound and head movement. Abnormal pressure or composition can contribute to dizziness, hearing changes, or ringing.

\medterm{Endolymphatic Hydrops} Pathological distension and increased volume of the endolymphatic fluid compartment of the membranous inner ear (scala media and vestibular apparatus) resulting from impaired fluid resorption in the endolymphatic sac, constituting the underlying pathophysiologic hallmark of Ménière disease.
\synonyms
Meniere Hydrops; Inner Ear Hydrops

\medterm{Endometrial} Relating to the endometrium, the inner lining of the uterus. This lining changes during the menstrual cycle and can support implantation of a pregnancy; abnormal growth may cause bleeding or affect fertility.

\medterm{Endometrial Ablation} A procedure that destroys much of the uterine lining to reduce heavy menstrual bleeding. It is not contraception, pregnancy afterward can be dangerous, and it is generally intended for people who do not plan future pregnancy.

\medterm{Endometrial Adenocarcinoma} A gland-forming cancer of the uterine lining. It often causes abnormal bleeding, especially after menopause; treatment depends on the cancer’s type, grade, stage, molecular features, and overall health.

\medterm{Endometrial Biopsy} Removal of a small sample of the uterine lining for laboratory examination, often to investigate abnormal bleeding or suspected overgrowth or cancer. Cramping and light bleeding are common; occasionally the sample is insufficient for diagnosis.

\medterm{Endometrial Cancer} Cancer arising from the inner lining of the uterus. Abnormal bleeding—especially after menopause—is a common warning sign; evaluation usually includes examination, imaging, and a sample of the uterine lining.

\synonyms
Uterine Endometrial Neoplasm; Endometrial Carcinoma

\medterm{Endometrial Carcinoma} Cancer arising from the uterine lining, often causing abnormal bleeding, particularly after menopause. Diagnosis requires examination of uterine-lining tissue, and treatment depends on type and spread.

\medterm{Endometrial Disorder} Any condition affecting the uterine lining, including inflammation, polyps, abnormal overgrowth, endometriosis-related changes, or cancer. It may cause heavy or irregular bleeding, pelvic pain, infertility, or no symptoms.

\medterm{Endometrial Hyperplasia} Abnormal overgrowth and thickening of the uterine lining, often from too much estrogen without enough progesterone. It may cause heavy, irregular, or postmenopausal bleeding; atypical changes can precede cancer and need evaluation.

\medterm{Endometrial Implants} Alternate name; see \textit{Endometriosis}.

\medterm{Endometrial Intraepithelial Neoplasia} A precancerous overgrowth of abnormal glands in the uterine lining. It can exist with or progress to endometrial cancer, so treatment may involve removal of the uterus or hormone treatment when preserving fertility is important.

\synonyms
EIN; Atypical Endometrial Hyperplasia

\medterm{Endometrial Neoplasm} An abnormal growth of the uterine lining that may be benign, precancerous, or cancerous. Symptoms, tissue examination, imaging, and—if cancer is present—stage determine treatment.

\medterm{Endometrial Polyp} A localized growth of the uterine lining that projects into the uterine cavity. It may cause irregular or heavy bleeding, infertility, or no symptoms; some polyps contain precancerous or cancerous cells.

\medterm{Endometrial Polyps} Localized growths of the uterine lining projecting into the uterine cavity. They may cause irregular or heavy bleeding, infertility, or no symptoms, and some require removal for tissue examination.

\medterm{Endometrial Receptivity Analysis} A test that examines activity of genes in a sample of the uterine lining to estimate when it may best support embryo implantation. Its ability to improve pregnancy rates is uncertain, so results need specialist interpretation.

\synonyms
ERA; Endometrial Receptivity Assay

\medterm{Endometrial Stromal Nodule} A rare, usually harmless, well-defined growth of connective-tissue cells in the uterus. It differs from cancer by not invading nearby tissue or blood vessels, but microscopic examination is needed for reliable classification.

\medterm{Endometrial Stromal Sarcoma} A rare cancer arising from connective-tissue cells of the uterine lining. Some forms grow slowly and others aggressively; abnormal bleeding or pelvic symptoms may occur, and treatment depends on grade, spread, hormone receptors, and molecular findings.

\synonyms
ESS

\medterm{Endometrioid Carcinoma Of The Vulva} A rare gland-forming cancer of the vulvar skin or nearby tissue that resembles cancer of the uterine lining. Diagnosis requires tissue examination and confirmation that it did not spread from the uterus or ovary.

\medterm{Endometrioid Ovarian Carcinoma} A type of ovarian cancer whose cells resemble glands of the uterine lining. It may be associated with endometriosis; symptoms, treatment, and outlook depend on stage, grade, and tumor features.

\medterm{Endometrioma} An ovarian cyst caused by endometriosis in which old blood collects. It may cause pelvic pain, painful periods, pain during sex, or infertility, although some people have no symptoms.

\medterm{Endometriosis} A chronic inflammatory condition in which tissue resembling the lining of the uterus grows outside the uterus, most often in the pelvis, ovaries, fallopian tubes, or tissue around the pelvic organs. It can also affect the bowel or bladder. This tissue responds to hormones and can cause repeated inflammation, bleeding, scar tissue, adhesions, or ovarian cysts. Symptoms may include severe period pain, ongoing pelvic pain, heavy or irregular bleeding, pain during sex or bowel movements, painful urination, bloating, nausea, tiredness, or infertility. Some people have no symptoms.

\synonyms
Endometrioma; Pelvic Endometriosis

\medterm{Endometritis} Inflammation or infection of the uterine lining, often after childbirth, miscarriage, abortion, or a procedure entering the uterus. It may cause fever, pelvic pain, tenderness, abnormal discharge, or bleeding and needs treatment.

\medterm{Endometrium} The inner lining of the uterus. It thickens each month under hormone control to prepare for pregnancy and sheds during menstruation; abnormal thickening, inflammation, polyps, or cancer can cause bleeding or fertility problems.

\medterm{Endomyocardial Biopsy} A procedure that removes tiny samples from the heart muscle through a catheter in a vein. The samples can help check for rejection after a heart transplant or investigate some heart-muscle diseases.

\synonyms
EMB

\medterm{Endomyocardial Fibrosis} Scarring and thickening of the inner heart lining and nearby muscle that makes the heart stiff and reduces filling. It can cause breathlessness, swelling, abnormal rhythms, or heart failure.

\medterm{Endomysium} A delicate layer of reticular connective tissue that individually invests each single skeletal muscle fiber, containing an intricate network of capillaries and nerve terminals essential for metabolic exchange and muscle contraction.

\medterm{Endoneurium} A delicate layer of loose connective tissue and collagen fibrils that directly surrounds each individual myelinated or unmyelinated axon and its accompanying Schwann cell sheath within a peripheral nerve fascicle.

\medterm{Endoparasite} A parasite that lives inside its host, such as a worm in the intestine or a microscopic parasite in blood or tissue. It may cause no symptoms or lead to inflammation, anemia, organ damage, or infection.

\medterm{Endopeptidase} An enzyme that cuts proteins or peptides at internal points rather than removing amino acids only from an end. Digestive enzymes such as pepsin and trypsin are examples.

\medterm{Endophthalmitis} A severe infection or inflammation inside the eye, usually after eye surgery, an injury, or spread through the bloodstream. It can cause severe pain, redness, swelling, and rapid vision loss and requires emergency eye treatment.

\medterm{Endoplasmic Reticulum} A network of tiny membrane-lined spaces inside cells. It helps make and transport proteins and fats, stores calcium, and processes substances needed for cell structure and function.

\medterm{Endoprosthesis} An artificial device placed inside the body to replace, support, or restore a body structure. Examples include joint replacements, vascular grafts or stents, and some devices used to keep an airway open.

\medterm{Endoribonuclease} An enzyme that cuts RNA at internal chemical bonds. It helps process RNA into usable forms or break down RNA that is damaged, unnecessary, or produced by viruses.

\medterm{Endorphin} A natural chemical made by the body that acts on opioid receptors in the brain and other tissues. It can reduce the feeling of pain and influence stress, mood, and other body responses.

\medterm{Endorphins} Natural opioid-like chemicals made by the body that reduce pain and influence mood. They are released during stress, exercise, injury, and other situations and act through opioid receptors in the nervous system.

\synonyms
Endogenous Opioid Peptides; Beta-Endorphin

\medterm{Endoscope} A flexible or rigid instrument with a light and camera used to view an internal body passage or cavity. It can sometimes collect tissue, stop bleeding, remove growths, or perform other treatments.

\medterm{Endoscopic Endonasal Surgery} A minimally invasive operation performed through the nostrils and nasal passages using an endoscope, commonly to reach selected skull-base, pituitary, sinus, or nearby lesions. Risks include bleeding, infection, cerebrospinal-fluid leak, and injury to nearby structures.

\medterm{Endoscopic Mucosal Resection} An endoscopic procedure that removes a superficial growth from the digestive tract after lifting it away from deeper tissue. The removed tissue is examined to determine whether it was completely removed and whether cancer is present.

\medterm{Endoscopic Retrograde Cholangiopancreatography} A procedure combining an endoscope and X-ray imaging to examine and treat the bile and pancreatic ducts. It can remove stones, open narrowed ducts, place stents, or obtain samples; pancreatitis, bleeding, and infection are possible risks.

\medterm{Endoscopic Submucosal Dissection} An endoscopic procedure that removes a larger or flatter superficial growth from the digestive tract in one piece by dissecting beneath it. It allows detailed tissue examination but carries risks of bleeding and perforation.

\medterm{Endoscopic Third Ventriculostomy} A brain procedure used to treat some types of hydrocephalus, in which fluid builds up because its normal path is blocked. A small camera is used to make an opening in one of the brain's fluid spaces so the fluid can drain. It may reduce the need for a shunt, but not everyone is a suitable candidate and the opening can close later.
\synonyms
ETV; Third Ventriculostomy; Endoscopic Ventriculostomy

\medterm{Endoscopic Thoracic Sympathectomy} A minimally invasive operation that interrupts selected nerves in the chest, usually for severe sweating of the hands. Possible complications include increased sweating elsewhere, drooping eyelid, collapsed lung, nerve injury, and heart or blood-pressure changes.

\medterm{Endoscopic Ultrasound} A procedure combining an endoscope with ultrasound to view the digestive-tract wall and nearby organs. It can guide needle sampling of lymph nodes, pancreas, or other masses and help assess the depth or spread of disease.

\medterm{Endoscopy} Examination of an internal body passage or cavity with a camera-equipped instrument. Depending on the site, it can directly inspect tissue, collect samples, remove growths, widen narrow areas, or stop bleeding.

\medterm{Endosome} A small membrane-bound compartment inside a cell that receives material taken in from the cell surface and sorts it for recycling, transport, or breakdown.

\medterm{Endosteum} A thin membrane lining the inner surfaces of bones, including the marrow cavity and small internal channels. It contains cells involved in bone formation, remodeling, and repair.

\medterm{Endostosis} Abnormal formation of bone on or within an existing bone. It may be a localized developmental finding or part of a broader bone disorder; the significance depends on its location, extent, and symptoms.

\medterm{Endothelial} Relating to the thin layer of cells lining blood and lymphatic vessels and some internal surfaces. These cells help control blood flow, vessel leakiness, inflammation, clotting, and new-vessel growth.

\medterm{Endothelial Cell} A cell lining a blood or lymphatic vessel. It regulates movement between blood and tissues and helps control vessel widening, inflammation, clotting, and formation of new blood vessels.

\medterm{Endothelial Nitric Oxide Synthase} A calcium- and calmodulin-dependent enzyme constitutively expressed in vascular endothelial cells that synthesizes nitric oxide from L-arginine, mediating basal vasodilation, inhibition of leukocyte adhesion, and suppression of platelet aggregation.

\synonyms
eNOS; NOS3

\medterm{Endothelin} A small protein made mainly by cells lining blood vessels that strongly narrows vessels and affects blood pressure, circulation, inflammation, and tissue growth. Excess signaling contributes to some cases of pulmonary arterial hypertension.

\medterm{Endothelin Receptor Antagonists} Medicines that block signals that narrow blood vessels, especially in the lungs. They are used for selected pulmonary arterial hypertension; liver injury, anemia, swelling, and serious fetal harm are important risks.

\medterm{Endothelium} The thin layer of specialized cells lining blood and lymphatic vessels. It controls exchange between blood and tissues and helps regulate vessel diameter, inflammation, clotting, and new-vessel growth.

\medterm{Endothelium-Derived Relaxing Factor} A substance released by vessel-lining cells that relaxes blood vessels. It is largely identified with nitric oxide, which also helps reduce platelet sticking and vessel-wall inflammation.

\synonyms
EDRF; Nitric Oxide

\medterm{Endotoxin} A component of the outer membrane of certain bacteria, especially Gram-negative bacteria. When released into the body, it can trigger fever, inflammation, low blood pressure, organ injury, and life-threatening sepsis.

\medterm{Endotracheal Intubation} Placement of a breathing tube through the mouth or nose into the windpipe to keep the airway open, deliver ventilator support, give anesthetic gases, or protect the lungs. Incorrect placement can cause severe harm and must be confirmed.

\medterm{Endotracheal Tube} A breathing tube passed through the mouth or nose into the windpipe to maintain an airway and support ventilation. An inflatable cuff may seal the airway; irritation, injury, infection, or incorrect placement can occur.

\medterm{Endourologist} A urologist specializing in minimally invasive procedures performed through the urinary tract or small access sites. Common work includes treating urinary stones, narrowing, blockage, and selected tumors.

\medterm{Endovaginal Ultrasound} Ultrasound imaging performed with a covered probe placed gently in the vagina to examine the uterus, uterine lining, ovaries, cervix, and early pregnancy. It uses sound waves rather than radiation and may cause temporary discomfort.

\medterm{Endovascular} Located inside a blood vessel or performed from within it, usually with a catheter passed through a small skin puncture. Endovascular procedures can treat aneurysms, clots, narrowing, bleeding, or abnormal vessels.

\medterm{Endovascular Aneurysm Repair} A catheter procedure that places a fabric-covered graft inside an aortic aneurysm to reinforce the artery and reduce pressure on its weakened wall. It usually recovers faster than open surgery but requires lifelong imaging for leakage or device problems.

\synonyms
EVAR; Endovascular Stent-Graft Repair; Endovascular Aortic Repair

\medterm{Endovascular Coiling} A minimally invasive, catheter-based procedure performed to treat brain aneurysms from within the blood vessels without opening the skull. Guided by real-time X-ray fluoroscopy, an interventional neurosurgeon or neuroradiologist threads a long, micro-thin catheter through an artery in the groin (femoral artery) or wrist (radial artery) all the way up into the blood vessels of the brain and directly inside the bulging aneurysm sac. Soft, ultra-fine platinum coils (each as thin as a human hair) are deployed through the catheter into the aneurysm, where they pack together and induce natural blood clot formation (thrombosis). The resulting clot seals the aneurysm from the inside, preventing blood from entering and drastically reducing the risk of aneurysm rupture and fatal intracranial bleeding.
\synonyms
Endovascular Aneurysm Coiling; Guglielmi Detachable Coiling; GDC; Aneurysm Embolization Coils

\medterm{Endovascular Embolization} A catheter procedure that deliberately blocks a selected blood vessel using coils, particles, plugs, or another material. It may control bleeding or treat an aneurysm, malformation, or tumor blood supply; non-target blockage is a possible risk.

\medterm{Endovascular Graft} A fabric-covered artificial tube placed inside a diseased blood vessel through a catheter. It can reinforce a weak artery or exclude an aneurysm; suitability and follow-up depend on anatomy, symptoms, and imaging.

\medterm{Endovascular Repair} A minimally invasive procedure that uses a catheter-delivered graft, balloon, stent, or other device to repair a blood vessel from inside. Bleeding, clotting, infection, leakage, or need for further treatment can occur.

\medterm{Endovenous Laser Ablation} A treatment for some varicose veins in which a thin laser fiber is placed inside the vein through a small needle opening. Heat closes the treated vein, and blood flows through healthier veins. It is usually done with local numbing medicine; bruising, pain, clots, nerve irritation, or return of the vein can occur.

\synonyms
EVLA; Endovenous Laser Therapy; EVLT; Endovenous Thermal Ablation

\medterm{End Point} A predefined outcome or measurement used to judge a clinical study, treatment, diagnostic test, or disease process. It may be a symptom, laboratory result, event, survival measure, or patient-reported outcome.

\medterm{End Stage} The advanced phase of a disease in which normal function is severely impaired. Meaning, criteria, expected course, and treatment options depend on the disease, such as advanced kidney, liver, lung, or heart failure.

\medterm{End-Stage Kidney Disease} Advanced chronic kidney disease in which the kidneys can no longer adequately remove waste or control fluid, electrolytes, and acid levels. Treatment options include dialysis, kidney transplantation, or supportive care based on health and preferences.

\medterm{End-Stage Renal Disease} The most advanced stage of chronic kidney disease, usually when kidney filtration is below 15 or dialysis is needed. Waste and fluid can accumulate, causing nausea, itching, confusion, swelling, anemia, or heart complications.

\synonyms
ESRD

\medterm{Endurance} The ability to continue physical or mental activity without becoming tired quickly. It depends on heart and lung function, muscle conditioning, energy, sleep, and overall health and can be assessed with activity or exercise tests.

\medterm{Enema} Introduction of liquid into the rectum through the anus to soften stool, stimulate bowel emptying, deliver selected treatment, or assist a test. Overuse can cause irritation, dehydration, or electrolyte problems.

\medterm{Energy} The capacity needed for body functions, movement, growth, and temperature control, obtained mainly from food. Persistent low energy may reflect poor sleep, inadequate nutrition, anemia, hormone disease, mood disorder, infection, or another illness.

\medterm{Energy Conservation} Strategies that reduce unnecessary effort and fatigue during daily activities, such as planning, pacing, prioritizing, resting, sitting for tasks, changing position, and using helpful equipment.

\medterm{Energy Expenditure} The energy the body uses for basic functions, movement, digestion, growth, pregnancy, and temperature control. It varies with body size, muscle mass, age, activity, illness, and other conditions.

\medterm{Energy Increased} A noticeable rise in energy or activity. It may be normal, but unusually high energy with little need for sleep, racing thoughts, impulsive behavior, or irritability can occur with stimulant use, thyroid disease, mania, or hypomania.

\medterm{Energy Intake} The amount of energy obtained from food and drink, usually measured in calories. Appropriate intake depends on age, body size, activity, illness, pregnancy, and nutritional goals.

\medterm{Energy Metabolism} The chemical processes that convert nutrients into usable energy for cells. They include breaking down carbohydrates, fats, and proteins and storing or using the resulting energy for body functions.

\medterm{Energy Resolution} The ability of a radiation detector to distinguish radiation particles with different energies. Better resolution allows more accurate identification and separation of radioactive substances during imaging or laboratory measurement.

\medterm{Energy Window} The range of radiation energies accepted by a detector or gamma camera to create an image. Choosing it appropriately improves image quality by accepting desired signals and reducing scattered radiation.

\medterm{Enfuvirtide} An injectable HIV medicine that blocks fusion of HIV with human immune cells. It is used with other antiretroviral medicines for resistant infection; injection-site reactions and allergic reactions can occur.

\medterm{Engagement} Descent of the widest part of the baby’s presenting head or body into the parent’s pelvis. It may occur before labor or during labor and indicates that the baby has entered the birth passage.

\medterm{Enhanced Elimination} Medical methods used to remove an absorbed poison more quickly, such as dialysis or changing urine chemistry. They are used only for selected toxins because benefits, timing, and risks depend on the substance and the person’s condition.

\medterm{Enhanced Pause} An unusually long pause in breathing or another measured body signal. Its significance depends on the test and setting and may require interpretation with symptoms and the complete recording.

\medterm{Enhanced Recovery After Surgery} A coordinated care plan before, during, and after surgery that supports recovery through education, appropriate nutrition and fasting, pain control, early movement, early eating, and prevention of complications.

\medterm{Enhancement Artifact} An imaging effect in which tissue behind a fluid-filled or low-density structure appears brighter than it really is, especially on ultrasound. It is a technical artifact, not necessarily a sign of increased blood flow or tissue density.

\medterm{Enhancer} A stretch of DNA that increases activity of a gene when regulatory proteins bind to it. It may be far from the gene and can work in either direction, helping control when and where the gene is active.

\medterm{Enkephalin} A natural opioid-like peptide made by the body that activates mainly delta and mu opioid receptors. It helps regulate pain, stress, mood, movement, and other nervous-system functions.

\synonyms
Leu-Enkephalin; Met-Enkephalin

\medterm{Enlarged Adenoids} Overgrowth of immune tissue behind the nose that can block nasal breathing and cause mouth breathing, snoring, repeated ear-fluid problems, or sleep-disordered breathing. Symptoms and age determine whether treatment is needed.

\medterm{Enlarged Fontanelle} A soft spot on an infant's skull that is significantly larger than normal for age, which may be associated with conditions such as congenital hypothyroidism, rickets, skeletal dysplasias, or increased intracranial pressure.

\synonyms
Large Fontanelle; Delayed Fontanelle Closure

\medterm{Enlarged Heart} An increase in heart size caused by conditions such as high blood pressure, valve disease, heart-muscle disease, anemia, or athletic adaptation. It may cause no symptoms or lead to breathlessness, swelling, rhythm problems, or heart failure.
\synonyms
Cardiomegaly

\medterm{Enlarged Liver} Alternate name; see \textit{Hepatomegaly}.

\synonyms
Hepatomegaly; Hepatic Enlargement

\medterm{Enlarged Prostate} Alternate name; see \textit{Benign Prostatic Hyperplasia}.

\medterm{Enlarged Spleen} Increase in spleen size caused by infection, liver or blood disease, inflammation, cancer, or congestion. It may cause fullness or pain and can increase rupture risk; evaluation considers symptoms, blood counts, and the underlying cause.

\synonyms
Splenomegaly; Splenic Enlargement

\medterm{Enophthalmos} Abnormal backward sinking of an eyeball within the eye socket. It may follow injury, loss of tissue around the eye, dehydration, aging, or a nerve disorder and can cause unequal eye appearance or double vision.

\medterm{Enoxaparin} A low-molecular-weight heparin that reduces clotting by enhancing antithrombin activity, especially against factor Xa. It prevents or treats blood clots but can cause bleeding, low platelets, or injection-site bruising.

\medterm{Enrichment} Adding nutrients to food after processing has removed some of them. Enriched flour, for example, commonly has certain B vitamins and iron added back to help prevent nutritional deficiencies.

\medterm{Entamoeba} A group of single-celled amoebas that live in the intestine or other tissues. Some species are harmless, while Entamoeba histolytica can cause diarrhea, intestinal ulcers, or liver abscesses.

\medterm{Entamoeba Coli} A usually harmless amoeba found in the human intestine and spread by fecal contamination. Its presence in stool indicates exposure to contaminated food or water but usually does not cause illness or require treatment.

\medterm{Entamoeba Hartmanni} A small, usually harmless intestinal amoeba spread through fecal contamination. It generally causes no disease; identification in stool mainly indicates possible exposure to contaminated food, water, or hands.

\medterm{Entamoeba Histolytica} A parasite spread when contaminated food, water, or hands are swallowed. It can cause diarrhea, abdominal pain, and bloody stools and may invade the intestine or spread to the liver, causing an abscess.

\medterm{Entamoeba Polecki} An intestinal amoeba that can infect people after fecal contamination, often through contact with animals or contaminated food and water. It is usually harmless but may occasionally be associated with diarrhea.

\medterm{Enteral} Relating to or delivered through the digestive tract. Enteral nutrition provides food and fluids by mouth or through a feeding tube into the stomach or intestine.

\medterm{Enteral Formula} A liquid nutrition preparation given by mouth or feeding tube. It may provide complete nutrition or selected nutrients and is chosen according to digestion, absorption, fluid needs, kidney or liver disease, and metabolic requirements.

\medterm{Enteral Nutrition} Delivery of food, fluids, and nutrients through the digestive tract, by mouth or feeding tube. It may be used when swallowing or eating is unsafe or inadequate; aspiration, diarrhea, blockage, and tube displacement are possible complications.

\medterm{Enteric} Relating to the intestine or digestive tract. The term can also describe organisms, infections, medicines, or coatings associated with the intestines.

\medterm{Enteric Bacteria} Bacteria that live in or affect the intestinal tract. Some are harmless normal residents, while others cause diarrhea or infections outside the intestine, such as urinary, wound, or bloodstream infection.

\medterm{Enteric-Coated} Covered with a protective layer that resists stomach acid and dissolves in the small intestine. This can protect the stomach or protect a medicine from stomach acid; such tablets should not usually be crushed.

\medterm{Enteric Duplication Cyst} A birth-related cyst attached to or within the digestive tract, most often the small intestine. It may cause pain, blockage, bleeding, infection, or no symptoms and can require surgery.

\medterm{Enteric Fever} A serious infection caused by \textit{Salmonella} Typhi or \textit{Salmonella} Paratyphi, usually spread through food or water contaminated with human waste. It can cause gradually rising fever, headache, stomach pain, weakness, diarrhea or constipation, and sometimes a rash. Untreated illness can cause bleeding or a hole in the intestine; testing and antibiotic treatment are needed.

\synonyms
Typhoid Fever; Paratyphoid Fever; Salmonella Enteric Infection

\medterm{Enteric Nervous System} The network of nerve cells within the digestive tract that helps control movement, secretion, blood flow, and communication with the brain and spinal cord.

\medterm{Enteritis} Inflammation of the small intestine caused by infection, inflammatory disease, radiation, reduced blood flow, medicines, or other injury. It may cause abdominal pain, diarrhea, vomiting, fever, bleeding, or dehydration.

\medterm{Enterobacter} A group of Gram-negative bacteria that may live in the intestine but can cause healthcare-associated urinary, bloodstream, wound, respiratory, or abdominal infection, especially in vulnerable people.

\medterm{Enterobacter Aerogenes} A former name for Klebsiella aerogenes, a Gram-negative bacterium that can cause urinary, lung, wound, abdominal, or bloodstream infection, especially in healthcare settings. It may resist several antibiotics.

\medterm{Enterobacter Cloacae} A Gram-negative bacterium that can cause healthcare-associated urinary, bloodstream, lung, wound, or abdominal infection. It may resist several antibiotics, so laboratory susceptibility testing guides treatment.

\medterm{Enterobacter Hormaechei} A bacterium in the Enterobacter cloacae group that can cause opportunistic healthcare-associated infection of blood, urine, lungs, wounds, or devices. Testing for antibiotic resistance is often important.

\medterm{Enterobacteriaceae} A family of Gram-negative bacteria that includes Escherichia, Klebsiella, Enterobacter, Salmonella, and Shigella. Some normally live in the intestine, while others cause intestinal, urinary, wound, or bloodstream infection.

\medterm{Enterobacteriaceae Carbapenem Resistance} Resistance in certain Gram-negative bacteria to carbapenem antibiotics, often caused by enzymes that destroy these medicines. Such infections can be difficult to treat and require laboratory testing, infection control, and specialist guidance.

\medterm{Enterobiasis} Infection with pinworms, small intestinal worms whose eggs spread easily between people. It commonly causes intense itching around the anus at night, disturbed sleep, and sometimes irritation in the genital area.

\medterm{Enterobius} A group of pinworms that live in the intestine. Enterobius vermicularis causes the common pinworm infection, which spreads when eggs are swallowed and often causes nighttime anal itching.

\medterm{Enterobius Vermicularis} A small intestinal worm commonly called the pinworm or threadworm. Adult females lay eggs around the anus, usually at night, causing itching and allowing infection to spread through contaminated hands or surfaces.

\synonyms
Pinworm; Threadworm; Enterobius

\medterm{Enterocele} A pelvic organ prolapse in which a pouch of abdominal lining and sometimes small intestine bulges into the space between the vagina and rectum. It may cause pelvic pressure, a vaginal bulge, back or rectal discomfort, or no symptoms.

\synonyms
Small Bowel Prolapse; Douglas Pouch Hernia

\medterm{Enterochromaffin Cells} Specialized enteroendocrine cells located throughout the mucosal epithelium of the gastrointestinal tract that synthesize, store, and release serotonin (5-hydroxytryptamine) and substance P to modulate gastrointestinal motility, secretion, and visceral pain perception.
\synonyms
EC Cells; Kultschitzky Cells

\medterm{Enterochromaffin-Like Cells} Specialized neuroendocrine cells situated within the oxyntic mucosa of the gastric fundus and body that synthesize and secrete histamine in response to gastrin and acetylcholine stimulation, directly stimulating parietal cell H+/K+-ATPase proton pumps.

\synonyms
ECL Cells

\medterm{Enteroclysis} An imaging test in which contrast material is delivered through a tube into the small intestine to look for inflammation, narrowing, tumors, or blockage. CT or MRI enterography is now more commonly used in many settings.

\medterm{Enterococcaceae} A family of bacteria that includes Enterococcus species. These bacteria may normally live in the intestine but can cause urinary, wound, bloodstream, or heart-valve infections when they enter protected body sites.

\medterm{Enterococcal Infection} An infection caused by Enterococcus bacteria, which normally live in the intestine but can cause urinary-tract, wound, abdominal, bloodstream, or heart-valve infections. Enterococcal infections are more common after hospitalization, surgery, antibiotic exposure, or use of medical devices. Some strains resist multiple antibiotics, so culture and susceptibility testing guide treatment.

\synonyms
Enterococcosis; Enterococcal Sepsis

\medterm{Enterococcus} A group of bacteria normally found in the intestine that can cause urinary, wound, abdominal, bloodstream, or heart-valve infection. Some strains resist several antibiotics and require laboratory susceptibility testing.

\medterm{Enterococcus Durans} A bacterium that may live in the bowel but can rarely cause infection in seriously ill or immunocompromised people. It may affect blood, urine, wounds, or other sites, and testing guides antibiotic treatment.

\medterm{Enterococcus Faecalis} A common intestinal bacterium that can cause urinary, wound, abdominal, bloodstream, or heart-valve infection. Disease is more likely after surgery, hospitalization, medical-device use, or when immunity is weakened.

\medterm{Enterococcus Faecium} An Enterococcus species that can cause healthcare-associated urinary, wound, abdominal, bloodstream, or heart-valve infection. Some strains resist vancomycin and other antibiotics, making laboratory testing essential.

\medterm{Enterocolitis} Inflammation of the small intestine and colon caused by infection, immune disease, reduced blood flow, toxins, or other injury. It may cause abdominal pain, diarrhea, fever, vomiting, bleeding, dehydration, or severe illness.

\medterm{Enterocyte} A cell lining the small intestine that digests and absorbs nutrients, water, and electrolytes. It also helps form the intestinal barrier and participates in immune defense.

\medterm{Enterocytozoon} A group of tiny spore-forming parasites that can cause persistent diarrhea and intestinal disease, especially in people with weakened immunity. Some species can also affect other organs.

\medterm{Enteroendocrine Cell} A specialized cell in the intestinal lining that senses food and releases hormones regulating digestion, appetite, gallbladder and pancreatic activity, and blood-glucose control.

\medterm{Enterohepatic Circulation} Recycling of a substance from the liver into bile and the intestine, then back to the liver through the bloodstream. This can prolong the effects of some medicines and natural body chemicals.

\medterm{Enteromonas Hominis} A usually harmless microorganism that may live in the intestine and appear in stool samples. It generally does not cause disease, so its presence should be interpreted with symptoms and other test results.

\medterm{Enteropathic} Relating to disease of the intestine or caused by intestinal disease. The term may describe arthritis, skin changes, protein loss, or other problems occurring with inflammatory bowel disease or impaired intestinal absorption.

\medterm{Enteropathic Arthritis} Inflammatory joint disease associated with inflammatory bowel disease such as Crohn disease or ulcerative colitis. It may affect the spine, joints between the spine and pelvis, or joints of the arms and legs.

\medterm{Enteropathy} Any disease or disorder affecting the intestine. Causes include infection, inflammation, gluten sensitivity, poor blood flow, toxins, medicines, inherited conditions, or impaired absorption.

\medterm{Enteropathy-Associated T-Cell Lymphoma} A rare, fast-growing cancer of immune cells in the intestine. It can be linked to long-standing celiac disease and may cause severe abdominal pain, weight loss, poor nutrient absorption, blockage, or a hole in the bowel.

\synonyms
EATL; Enteropathy-Type Intestinal T-Cell Lymphoma

\medterm{Enteroscopy} Examination of the small intestine with a camera-equipped instrument. It can identify disease, obtain tissue samples, treat bleeding, or remove selected growths beyond the reach of routine upper endoscopy.

\medterm{Enterospasm} An involuntary tightening or contraction of the intestine that may cause cramping, abdominal pain, bloating, or altered movement of intestinal contents.

\medterm{Enterostomy} A surgically created opening that brings part of the intestine to the abdominal surface so stool or intestinal contents can leave the body or a feeding solution can enter.

\medterm{Enterostomy Procedure} Surgery that creates an opening from the intestine to the abdominal wall, such as an ileostomy or colostomy. It may divert stool or support feeding; bleeding, infection, blockage, and skin problems can occur.

\medterm{Enterotoxin} A toxin made by some microorganisms that acts on the intestinal lining and can cause vomiting, cramps, watery diarrhea, dehydration, or other intestinal illness.

\medterm{Enterovesical Fistula} An abnormal connection between the intestine and urinary bladder. It can cause repeated urinary infections, air or stool in the urine, abdominal pain, or urine leaking into the intestine and may result from inflammation, cancer, or surgery.

\medterm{Enterovirulent E. Coli} A group of Escherichia coli strains that cause intestinal disease by making toxins, attaching to the bowel, or invading it. Illness may include watery or bloody diarrhea, cramps, vomiting, or dehydration.

\medterm{Enterovirus} A group of viruses spread mainly by respiratory secretions or contact with infected stool. Many infections are mild or cause no symptoms, but some cause hand-foot-and-mouth disease, rash, meningitis, heart inflammation, or severe illness in newborns.

\medterm{Enterovirus D68} A virus that usually causes a respiratory illness but can rarely cause severe pneumonia or sudden weakness of the limbs, especially in children. It spreads mainly through respiratory secretions and contaminated surfaces.

\medterm{Enterovirus Infection} Alternate name; see \textit{Enterovirus}.

\medterm{Enthesis} The place where a tendon, ligament, or joint capsule attaches to bone. Repeated stress or inflammation there can cause localized pain and stiffness and occurs in several inflammatory joint diseases.

\medterm{Enthesitis} Inflammation where a tendon, ligament, or joint capsule attaches to bone. It can cause localized pain, tenderness, and stiffness and is common in spondyloarthritis, including psoriatic and inflammatory-bowel-disease-related arthritis.

\medterm{Enthesitis And Enthesopathy} Enthesitis is inflammation where a tendon, ligament, or fascia attaches to bone. Enthesopathy includes inflammation, wear, injury, or other disease at that site and may cause localized pain and stiffness.

\synonyms
Enthesopathy; Insertion Tendinopathy

\medterm{Enthesitis-Related Arthritis} A type of arthritis that begins in childhood and causes inflammation where tendons or ligaments attach to bone. It often affects the legs, feet, or lower back and may also cause swelling in joints or pain where the heel tendon attaches.

\synonyms
ERA; Juvenile Spondyloarthritis

\medterm{Entorhinal Cortex} A region of the brain that links the hippocampus with other memory and navigation areas. It helps process spatial information and new memories and is affected early in many cases of Alzheimer disease.

\medterm{ENT Physician} An otolaryngologist, a physician who diagnoses and treats disorders of the ear, nose, throat, and related head and neck structures.

\medterm{Entrance Wound} The wound where a projectile enters the body. Its appearance may help estimate the direction and distance of the shot, but findings vary with the weapon, clothing, body site, and circumstances and require forensic examination.

\medterm{Entropion} Inward turning of an eyelid, usually the lower lid, so the eyelashes rub against the eye. It can cause irritation, tearing, pain, corneal scratches, infection, and reduced vision.

\medterm{Enucleation} The surgical removal of an entire organ or well-defined encapsulated tumor in one piece without cutting into surrounding healthy tissue. In ophthalmology, it specifically refers to the complete surgical removal of the eyeball (globe) and optic nerve stump while preserving the eyelids, orbital fat, and extraocular muscles. It is performed for intraocular malignancies (e.g., retinoblastoma, uveal melanoma), severe irreparable ocular trauma, endophthalmitis, or a blind, painfully disfigured eye. An orbital implant is placed and attached to the eye muscles to support a prosthetic eye (artificial eye). In general surgery, it describes the clean ``shelling out'' of an intact benign tumor (such as a uterine fibroid, insulinoma, or cyst).
\synonyms{Ocular Enucleation; Eye Enucleation; Surgical Enucleation}

\medterm{Enucleation Procedure} Alternate name; see \textit{Enucleation}.

\medterm{Enuresis} Repeated involuntary urination, usually during sleep in a child expected to remain dry. It may be primary or begin again after dryness; constipation, infection, diabetes, sleep problems, stress, or urinary disorders may contribute.

\medterm{Environment} The physical, chemical, biological, and social surroundings in which a person lives or works. Environmental factors such as air, water, housing, noise, climate, and pollutants can influence health.

\medterm{Environmental Exposure} Contact with a chemical, pollutant, organism, physical hazard, or other factor in the surroundings. Health effects depend on the substance, dose, route, duration, timing, and individual susceptibility.

\medterm{Environmental Health} The field concerned with how surroundings affect human health. It includes air and water quality, food safety, housing, workplace hazards, climate, radiation, pollution, and prevention of harmful exposures.

\medterm{Environmental Illness} Illness caused or worsened by an environmental exposure. Evaluation considers the timing, dose, setting, symptoms, other possible causes, and whether symptoms improve when exposure stops.

\medterm{Enzyme} A protein, or occasionally an RNA molecule, that speeds a specific chemical reaction without being used up. Enzymes are essential for digestion, energy production, DNA processes, and many other cell functions.

\medterm{Enzyme Activity} The rate at which an enzyme carries out a chemical reaction under specified conditions. It is affected by temperature, acidity, inhibitors, helper molecules, substrate amount, and enzyme concentration.

\medterm{Enzyme Deficiency} Reduced or absent activity of an enzyme, usually inherited or sometimes acquired. It can cause buildup of substances, lack of needed products, and characteristic metabolic or cellular disease.

\medterm{Enzyme Induction} Increased production or activity of an enzyme, often after exposure to a medicine or chemical. If liver drug-metabolizing enzymes are induced, some medicines may be broken down faster and become less effective.

\medterm{Enzyme Inhibition} Reduced activity of an enzyme, sometimes caused by a medicine or chemical. If a drug-metabolizing enzyme is inhibited, another medicine may accumulate and cause stronger effects or toxicity.

\medterm{Enzyme Kinetics} The study of how quickly enzyme reactions occur and how their rate changes with substrate amount, temperature, acidity, inhibitors, helper molecules, and enzyme concentration.

\medterm{Enzyme-Linked Immunosorbent Assay} A laboratory test that uses antibodies and an enzyme-generated color signal to detect or measure a target such as an antibody, antigen, hormone, or infection marker.

\medterm{Enzyme Markers} Enzymes measured in blood or another body fluid to help assess tissue injury or organ function. Abnormal levels can suggest liver, bile-duct, muscle, heart, bone, or pancreatic disease but are not a diagnosis by themselves.

\medterm{Enzyme Precursors} Inactive or weakly active proteins that become working enzymes after being cut or chemically changed. This allows the body to activate digestive and other enzymes at the correct place and time.

\medterm{Enzyme Therapy} Treatment that replaces, supplements, blocks, or uses an enzyme to correct a deficiency or alter a disease process. The purpose, dose, benefits, and risks depend on the specific enzyme and condition.

\medterm{Eosinophil} A type of white blood cell involved in allergic inflammation and defense against some parasites. Eosinophils normally make up a small fraction of blood cells; high levels can occur with allergy, infection, medicines, or other disease.

\medterm{Eosinophilia} An abnormally high number of eosinophils in the blood, commonly caused by allergy, parasites, medicines, autoimmune disease, or some blood disorders. Severe or persistent elevation can injure the heart, lungs, skin, nerves, or other organs.

\medterm{Eosinophilic Asthma} Asthma in which eosinophils, a type of white blood cell, contribute to airway inflammation. It may cause recurring wheezing, cough, chest tightness, or breathlessness and can be difficult to control in some people.

\medterm{Eosinophilic Bronchitis} A cause of long-lasting cough in which eosinophils inflame the airways without the variable narrowing typical of asthma. The cough may resemble asthma or reflux and needs assessment to identify the cause.

\medterm{Eosinophilic Esophagitis} A chronic immune-related disease in which eosinophils inflame the esophagus. It may cause difficulty swallowing, food becoming stuck, chest or upper-abdominal pain, reflux-like symptoms, and narrowing of the esophagus.

\medterm{Eosinophilic Fasciitis} A rare inflammatory disease causing painful swelling and hardening of the deeper skin and tissue around muscles, usually in the arms or legs. Blood eosinophils may be high, and joint stiffness or limited movement can develop.

\synonyms
Shulman Syndrome; Diffuse Fasciitis with Eosinophilia

\medterm{Eosinophilic Gastroenteritis} A rare immune-related disease in which eosinophils inflame the stomach or intestines. Symptoms may include abdominal pain, nausea, vomiting, diarrhea, bleeding, poor absorption, or weight loss.

\medterm{Eosinophilic Gastrointestinal Disorders} Conditions in which eosinophils, a type of white blood cell, build up in one or more parts of the digestive tract and cause inflammation. Symptoms may include pain, vomiting, swallowing difficulty, diarrhea, or poor growth; diagnosis usually requires endoscopy and biopsy.

\medterm{Eosinophilic Granuloma} A localized form of Langerhans-cell disease in which abnormal immune cells collect, often in bone or skin. Bone lesions can cause pain, swelling, or fractures and may require imaging and tissue examination.

\medterm{Eosinophilic Granulomatosis With Polyangiitis} A rare disease in which inflamed blood vessels occur with asthma and high levels of eosinophils, a type of white blood cell. It can damage the lungs, skin, nerves, heart, kidneys, or digestive organs.

\synonyms
Allergic Granulomatosis
Allergic Granulomatosis and Angiitis

\medterm{Eosinophils} White blood cells that help control allergic inflammation and defend against some parasites. High levels may occur with allergies, asthma, infection, medicines, autoimmune disease, or blood disorders; low levels usually have little significance.

\synonyms
Eosinophil Granulocytes

\medterm{Ependyma} A specialized layer of cells lining the brain’s fluid-filled ventricles and the spinal cord’s central canal. It helps separate nervous tissue from cerebrospinal fluid and supports movement of that fluid.

\medterm{Ependymal Cell} A supporting cell lining the brain ventricles and spinal cord central canal. It helps maintain the cerebrospinal-fluid boundary and, in some areas, helps move or produce the fluid.

\medterm{Ependymal Tumor} A growth arising from cells lining the brain ventricles or spinal canal. It may block cerebrospinal-fluid flow and cause headache, vomiting, balance problems, weakness, or hydrocephalus; treatment depends on its type and extent.

\medterm{Ependymoblastoma} A rare, aggressive brain tumor of childhood made from very immature cells. It may cause headache, vomiting, seizures, or changes in movement or development; tissue and genetic testing guide treatment and outlook.

\medterm{Ependymoma} A tumor arising from cells lining the brain ventricles or spinal canal. It may block fluid flow and cause headache, vomiting, seizures, balance problems, weakness, or hydrocephalus, especially in children.

\medterm{Ephedrine} A stimulant that activates the sympathetic nervous system and can widen airways and raise blood pressure and heart rate. It has limited medical uses and may cause palpitations, insomnia, anxiety, high blood pressure, or misuse.

\medterm{Ephedrinum} Another name for ephedrine, a stimulant that can increase heart rate, blood pressure, alertness, and airway diameter. Excessive use may cause dangerous heart-rhythm, blood-pressure, or nervous-system effects.

\medterm{Ephelides} Small, flat freckles caused by increased pigment production, often inherited. They become darker with sunlight and usually fade when exposure decreases; sun protection helps prevent further pigment change and skin damage.

\medterm{Ephrins} Proteins on the surface of cells that help guide cell position, nerve growth, blood-vessel development, and tissue organization. Abnormal ephrin signaling can contribute to some diseases and cancers.

\medterm{Epiblast} A layer of cells in the early embryo that develops into the three main cell layers of the body and ultimately forms most tissues and organs.

\medterm{Epicanthal Folds} Skin folds of the upper eyelid that partly cover the inner corner of the eye. They are normal in many people but can also occur with some chromosome or developmental conditions.

\medterm{Epicardium} The thin outer layer of the heart wall and the inner layer of its surrounding protective sac. It contains connective tissue, fat, blood vessels, and nerves and helps reduce friction as the heart beats.

\medterm{Epichlorohydrin} An industrial chemical used to make resins and other materials. Exposure can irritate or burn the skin, eyes, and airways; swallowing or substantial exposure can cause serious poisoning and requires urgent medical advice.

\medterm{Epicondyle} A bony projection beside a rounded joint surface where muscles and ligaments attach. Repeated strain or injury around an elbow epicondyle can cause localized pain and weakness.

\medterm{Epicondylitis} Pain and tenderness where a tendon attaches near an epicondyle, usually from repeated strain rather than infection. At the elbow it can weaken gripping and wrist movement and is commonly called tennis or golfer’s elbow.

\synonyms
Tennis Elbow; Golfer’s Elbow

\medterm{Epicondylitis, Lateral} Alternate name; see \textit{Tennis Elbow}.

\medterm{Epidemic} An increase in disease cases above the number normally expected in a population, place, or period. An epidemic may result from infection, environmental exposure, behavior, or another shared cause.

\medterm{Epidemic Keratoconjunctivitis} A very contagious viral infection of the eye. It can cause red eyes, pain, watering, swollen eyelids, light sensitivity, blurred vision, and swollen glands near the ear; symptoms may last for weeks.

\synonyms
EKC; Adenoviral Keratoconjunctivitis; Shipyard Eye

\medterm{Epidemiological Study} A study of health and disease in groups of people. It may measure how often a condition occurs, compare possible causes or exposures, or examine outcomes over time to help guide prevention and healthcare planning.

\synonyms
Epidemiologic Study

\medterm{Epidemiologic Method} A structured way to measure how often a health problem occurs and to study whether exposures, behaviors, or other factors are associated with it in a population. Results guide prevention and policy.

\medterm{Epidemiologic Studies} Research that examines the distribution, causes, risk factors, prevention, or outcomes of health events in populations.

\medterm{Epidemiologist} A scientist or public-health professional who studies the distribution, causes, risk factors, and prevention of health conditions in populations.

\medterm{Epidemiology} The study of how often diseases and other health events occur, who is affected, and which factors influence them. It uses population data to identify causes, guide prevention, investigate outbreaks, and evaluate healthcare.

\medterm{Epidermal} Relating to the epidermis, the outermost layer of skin. An epidermal problem affects this protective surface layer rather than the deeper skin tissues.

\medterm{Epidermal Acanthosis} Thickening of the epidermis caused by increased growth of its cells, seen in conditions such as psoriasis, chronic rubbing, and some inflammatory or genetic skin disorders.

\medterm{Epidermal Cells} Cells forming the outer layer of skin, chiefly keratinocytes, which create a protective barrier; melanocytes, immune cells, and touch-associated cells are also present.

\medterm{Epidermal Growth Factor} A signaling protein that binds the epidermal-growth-factor receptor and promotes cell growth, survival, and tissue repair; abnormal signaling can contribute to cancer.

\medterm{Epidermal Inclusion Cyst} A common, usually harmless lump beneath the skin filled with keratin, a protein found in skin. It often has a small central opening and may become swollen, painful, infected, or recur if its wall remains.

\medterm{Epidermal Layers} The epidermis has several layers of cells that make and protect the skin. From deepest to most superficial they are basal, spinous, granular, clear, and horny layers; the clear layer is mainly in the palms and soles.

\medterm{Epidermal Nevus} A usually benign, often present-from-birth or childhood skin lesion caused by an overgrowth of epidermal cells. It commonly appears as a localized, linear, or warty thickened patch and may occur with other developmental abnormalities.

\medterm{Epidermis} The outer layer of skin that forms a protective barrier against injury, germs, and water loss. It contains cells that make keratin and pigment, as well as immune and touch-sensing cells, and it is replaced continuously.

\medterm{Epidermodysplasia Verruciformis} A rare inherited disorder of skin immunity associated with susceptibility to certain human papillomavirus infections, widespread wart-like lesions, and increased risk of some skin cancers.

\medterm{Epidermoid Cyst} Alternate index label; see \textit{Epidermoid Cysts}.

\medterm{Epidermoid Cysts} Slow-growing, usually harmless sacs beneath the skin filled with keratin. They often appear as movable lumps and may become painful, red, swollen, or infected if they rupture.
\synonyms
Epidermal Inclusion Cysts; Infundibular Cysts

\medterm{Epidermolysis Bullosa} A group of inherited disorders in which the skin blisters or tears after minor rubbing or injury. Severity varies; some forms heal without scars, while others cause scarring, tight joints, poor nutrition, infection, or increased skin-cancer risk.

\medterm{Epidermolysis Bullosa Acquisita} An autoimmune blistering disease in which the body attacks a protein that helps hold skin layers together. The skin becomes fragile and blisters easily after minor injury; scars may form as it heals.

\medterm{Epidermophyton} A group of fungi that grow in keratin, a tough protein in skin, hair, and nails. They can cause ringworm infections, including itchy, scaly skin disease and fungal infection of the nails, but do not usually infect living hair shafts.

\medterm{Epidermophyton Floccosum} A dermatophyte fungus that commonly causes tinea of the body, groin, and feet and can infect nails.

\medterm{Epididymal Cyst} A usually benign, fluid-filled cyst arising in the epididymis. It may be painless and may require treatment only if it causes significant symptoms or complications.

\medterm{Epididymal Neoplasm} A benign or malignant growth arising in the epididymis; most palpable epididymal masses are benign, but persistent or enlarging lesions require evaluation.

\medterm{Epididymides} The paired coiled ducts on the posterior surface of the testes where sperm mature and are stored before entering the vasa deferentia.

\synonyms
Epididymal Ducts; Convoluted Seminiferous Ducts

\medterm{Epididymis} A tightly coiled tube along the back of each testicle where sperm mature, gain the ability to move, and are stored before passing into the vas deferens. It has a head, body, and tail.

\medterm{Epididymitis} Inflammation of the epididymis, usually caused by infection. It commonly causes gradually developing pain and swelling on one side of the scrotum. Causes include sexually transmitted infections, urinary bacteria, urine-flow problems, and less commonly injury or medicines.

\medterm{Epididymitis And Orchitis} Alternate name for epididymo-orchitis, inflammation of the epididymis and testicle.

\medterm{Epididymo-Orchitis} Inflammation of the epididymis and testicle, usually caused by an infection. It can cause one-sided scrotal pain, swelling, redness, painful urination, and fever. The cause may be a sexually transmitted infection or bacteria from the urinary tract, depending on the person and situation. Sudden severe pain needs emergency assessment to rule out testicular torsion; treatment depends on the cause.

\synonyms
Epididymo-orchitis; Testicular and Epididymal Inflammation; Orchiepididymitis

\medterm{Epidural} Pain relief or regional anesthesia given through a small catheter placed in the lower back, near the spinal nerves. It is commonly used during labor and some operations; possible effects include temporary low blood pressure, itching, headache, or incomplete relief.

\medterm{Epidural Abscess} A pocket of pus between the tough outer covering of the brain or spinal cord and nearby bone. It can compress the brain, spinal cord, or nerves, causing fever, severe pain, weakness, numbness, or loss of bladder control and requires emergency treatment.

\medterm{Epidural Analgesia} Pain relief produced by delivering numbing medicine, sometimes with an opioid, through a small catheter into the space around the spinal nerves. It may be used during labor or surgery while the person remains awake. Low blood pressure, itching, headache, infection, or incomplete relief can occur.

\medterm{Epidural Analgesia For Labor} Pain relief during labor provided through a small catheter placed in the lower back. Numbing medicine, sometimes with an opioid, is given gradually to reduce pain while allowing the laboring person to remain awake and participate in birth. Temporary low blood pressure, itching, headache, or incomplete relief can occur.

\medterm{Epidural Anesthesia} Regional anesthesia produced by injecting numbing medicine, sometimes with an opioid, into the space around the spinal nerves. It blocks pain from part of the body while the person stays awake; low blood pressure, headache, infection, or nerve injury can rarely occur.

\synonyms
Epidural Anaesthesia

\medterm{Epidural Blood Patch} A procedure that injects a small amount of the person’s own blood into the space around the spinal nerves to seal a leak after spinal anesthesia or an accidental puncture. It usually treats a severe positional headache; back pain, infection, or another puncture can occur.

\medterm{Epidural Hematoma} A collection of blood between the skull and the brain’s tough outer covering, usually after a head injury. It can rapidly compress the brain, causing headache, vomiting, confusion, weakness, loss of consciousness, or death and requires emergency treatment.

\medterm{Epidural Neoplasm} A tumor in the space outside the dura surrounding the spinal cord or brain; spinal epidural tumors can compress nerves or the spinal cord and cause pain, weakness, or sensory loss.

\medterm{Epidural Space} The space between the tough outer covering of the spinal canal and the canal wall. Medicines injected there can reduce pain during childbirth, surgery, or some pain treatments.

\medterm{Epidural Steroid Injection} A minimally invasive outpatient spine procedure used to relieve severe radiating back, neck, arm, or leg pain (such as sciatica or lumbar radiculopathy) caused by an inflamed or compressed spinal nerve. Under live X-ray guidance (fluoroscopy) and local anesthesia, a pain specialist or spine clinician carefully directs a thin needle into the epidural space surrounding the spinal nerves and dural sac, injects a tiny amount of contrast dye to verify accurate placement, and administers a powerful long-acting corticosteroid combined with a local anesthetic. The medication coats the irritated nerve roots, flushing away inflammatory chemicals and dramatically reducing swelling caused by herniated discs, spinal stenosis, degenerative disc disease, or bone spurs. Epidural steroid injections can be performed through translaminar, transforaminal, or caudal routes, offering substantial pain relief that allows patients to engage in active physical therapy and frequently avoid spinal surgery.

\synonyms
ESI; Epidural Injections for Back Pain; Translaminar Epidural Injection; Transforaminal Epidural Injection; Caudal Epidural Steroid Injection

\medterm{Epigastric} Relating to the upper middle area of the abdomen between the ribs. Pain there may arise from the stomach, esophagus, pancreas, gallbladder, intestine, or other organs.

\medterm{Epigastric Artery} An artery supplying the front abdominal wall. The superior and inferior epigastric arteries provide blood to abdominal muscles and skin and are important landmarks during some procedures.

\medterm{Epigastric Pain} Pain or discomfort in the upper middle abdomen, below the ribs and above the navel. Causes include indigestion, ulcer, reflux, gallbladder disease, pancreatitis, and other conditions; severe or persistent pain needs assessment.

\medterm{Epigastric Region} The upper middle section of the abdomen between the rib margins and above the navel. It overlies parts of the stomach, duodenum, pancreas, and liver; pain there can have several digestive or vascular causes.

\medterm{Epigastrium} The upper middle part of the abdomen, below the breastbone and between the ribs. It overlies parts of the stomach, liver, pancreas, and small intestine; pain there has many possible causes.

\medterm{Epigenetics} Changes in how genes are switched on or off without changing the DNA sequence itself. Chemical tags and DNA-packaging changes can be influenced by development, environment, and illness and may sometimes be reversible.

\medterm{Epigenome} The collection of chemical tags and DNA-packaging features that help control which genes are active in a cell without changing the DNA sequence.

\medterm{Epigenomics} The study of chemical and structural changes that control which genes are active across the genome without changing the DNA sequence. These changes can influence development, disease, and treatment response.

\medterm{Epiglottic Vallecula} A shallow mucosal depression situated in the glossoepiglottic space between the base of the tongue and the anterior surface of the epiglottis, bounded by median and lateral glossoepiglottic folds, serving as the anatomical seat for the tip of a curved Macintosh laryngoscope blade during endotracheal intubation.
\synonyms
Vallecula Epiglottica; Vallecular Fossa

\medterm{Epiglottic Vallecula} A shallow mucosal depression or pocket located at the base of the tongue immediately anterior to the epiglottis, bounded medially by the median glossoepiglottic fold and laterally by lateral glossoepiglottic folds; serves as an anatomical landmark for laryngoscope blade placement during intubation.

\synonyms
Vallecula

\medterm{Epiglottis} A leaf-shaped flap of flexible tissue at the base of the tongue that helps direct food and liquid into the esophagus during swallowing and away from the airway. Swelling can rapidly obstruct breathing.

\medterm{Epiglottitis} A potentially life-threatening infection or inflammation causing swelling around the epiglottis and upper airway. Fever, severe sore throat, drooling, trouble swallowing, noisy breathing, or sitting forward requires emergency care; the throat should not be examined forcefully.

\medterm{Epilepsia Partialis Continua} Repeated or nearly continuous twitching in one part of the body, often an arm, leg, or side of the face. Awareness may remain normal, but the condition can last for hours or longer and needs urgent medical assessment.

\synonyms
EPC; Kojewnikoff Syndrome

\medterm{Epilepsy} A chronic brain disorder that causes a lasting tendency to have recurrent unprovoked seizures. A seizure is a brief episode caused by excessive, abnormal electrical activity in the brain and may change movement, sensation, awareness, behavior, or consciousness. Seizures may begin in one area of the brain (focal seizures) or involve both sides from the start (generalized seizures). Causes include genetic conditions, brain injury or malformation, stroke, infection, metabolic or immune disease, or no known cause. A single seizure, especially one triggered by a temporary illness or abnormal body chemistry, does not always mean epilepsy.

\synonyms
Seizure Disorder

\medterm{Epilepsy, Frontal Lobe} Alternate name; see \textit{Frontal Lobe Epilepsy}.

\medterm{Epilepsy Surgery} Surgery or a procedure used for selected people whose seizures continue despite appropriate medicines and arise from an area that can be safely treated. Detailed testing identifies the seizure source and estimates benefits and risks.

\medterm{Epileptic Spasms} Brief seizures, most common in infants, in which the head, body, arms, or legs suddenly bend or stiffen. They often occur in clusters and need prompt assessment because they can affect development.

\synonyms
Infantile Spasms

\medterm{Epileptiform} Resembling a seizure or the electrical pattern associated with seizure activity. It may describe movements or an EEG finding but does not by itself prove that a person has epilepsy.

\medterm{Epimysium} A substantial sheath of dense irregular collagenous connective tissue that completely envelops the entire exterior surface of a skeletal muscle, continuous with tendons, muscular fascial planes, and the periosteum of bone.

\medterm{Epinephrin} Alternate spelling; see \textit{Epinephrine}.

\medterm{Epinephrine} A hormone and neurotransmitter released during stress that increases heart activity, opens the airways, and narrows many blood vessels. As a medicine, it treats anaphylaxis and selected cardiac or breathing emergencies.

\synonyms
Adrenaline

\medterm{Epineurium} The thick, outermost fibrous connective tissue sheath that encloses an entire peripheral nerve, surrounding multiple nerve fascicles, longitudinal blood vessels (vasa nervorum), and adipose tissue to provide tensile mechanical protection.

\medterm{Epiphora} Excessive tearing caused by blocked tear drainage or irritation that triggers extra tears. Causes include dry eye, allergy, infection, eyelid position problems, and blockage of the tear duct.

\medterm{Epiphyseal Plate} The growth plate, a layer of cartilage near the end of a growing bone that allows the bone to lengthen. It gradually closes after puberty; injury can cause pain, deformity, or unequal growth.

\synonyms
Growth Plate; Physis

\medterm{Epiphyseal Plate Disorders} Conditions affecting a growing bone’s growth plate, including fractures, inflammation, infection, abnormal blood supply, premature closure, and disorders in which an epiphysis shifts or separates. They may cause pain, swelling, limping, deformity, or unequal limb growth.

\medterm{Epiphysis} The expanded end of a long bone, usually covered by smooth cartilage where it forms a joint. It contains mostly spongy bone and helps determine the shape and movement of the joint.

\medterm{Epiphysitis} Inflammation or injury involving a growth plate or the area near the end of a growing bone. It may cause localized pain, swelling, or altered growth, and the cause depends on age and the affected bone.

\medterm{Epiploic Foramen} The narrow anatomical communication aperture connecting the greater peritoneal sac with the lesser sac (omental bursa), bounded anteriorly by the hepatoduodenal ligament (containing the portal triad), posteriorly by the inferior vena cava, superiorly by the caudate lobe of the liver, and inferiorly by the first part of the duodenum.

\synonyms
Foramen of Winslow

\medterm{Epipodophyllotoxin} A group of cancer medicines derived from a plant compound that damage DNA by blocking an enzyme needed for cell division. They treat selected cancers but may cause low blood counts, nausea, hair loss, infection, or rarely a later blood cancer.

\medterm{Epiretinal Membrane} A thin layer of scar-like tissue growing on the retina, usually over the macula. It may cause blurred or distorted central vision and sometimes needs surgery if visual impairment is significant.

\medterm{Epirubicin} An anthracycline cancer medicine that damages DNA and blocks cell division. Important risks include low blood counts, infection, nausea, hair loss, tissue injury after leakage, and dose-related heart damage.

\medterm{Epirubicin Hydrochloride} The salt form of epirubicin, an anthracycline cancer medicine that damages tumor-cell DNA. It is used for selected cancers but can cause low blood counts, tissue injury after leakage, and dose-related heart damage.

\medterm{Episcleritis} Usually short-lived inflammation of the thin tissue over the white of the eye, causing a localized red patch and mild tenderness or discomfort. Severe pain, light sensitivity, or vision loss suggests deeper eye disease.

\medterm{Episiotomy} A cut made in the tissue between the vagina and anus during vaginal birth to enlarge the opening. Routine use is not recommended; selected use may help urgent or assisted delivery but can cause pain, bleeding, infection, or deeper tearing.

\medterm{Episiotomy Repair} Closing an episiotomy or related tear with stitches after birth. The area is examined to identify deeper injury, and care includes pain relief, hygiene, monitoring for infection, and follow-up of healing and pelvic-floor function.

\medterm{Episodic Ataxia} An inherited condition causing repeated attacks of poor coordination and balance. Episodes may last seconds to hours and can include dizziness, slurred speech, muscle tightening, or abnormal eye movements, with improvement between attacks.

\medterm{Epispadias} A birth defect in which the urethral opening is on the upper surface of the penis or in an unusual position in a female. It may affect urine control and is sometimes associated with bladder-exstrophy conditions.

\medterm{Epistasis} A genetic interaction in which one gene changes or masks the effect of another gene. This can alter inherited traits and complicate how a disease risk or laboratory finding is predicted from genetic testing.

\medterm{Epistaxis} Medical term; see \textit{Nosebleeds}.

\medterm{Posterior Epistaxis} Nasal hemorrhage originating from arteries situated in the posterior nasal cavity or nasopharynx (most commonly branches of the sphenopalatine artery or Woodruff plexus), typically severe, profuse, draining down the oropharynx, and requiring posterior nasal packing or arterial embolization.

\medterm{Epithalamic} Relating to the epithalamus, a small region near the center of the brain that includes the pineal gland and structures involved in sleep timing, hormonal rhythms, and responses to reward or stress.

\medterm{Epithalamus} A dorsal posterior division of the diencephalon comprising the pineal gland (epiphysis cerebri), the habenular nuclei, the stria medullaris thalami, and the posterior commissure, involved in circadian melatonin rhythmicity and limbic-midbrain connections.

\medterm{Epithelial} Relating to cells that cover the skin, line body passages and organs, or form glands. These cells protect surfaces and may absorb substances, release fluids, or help detect changes.

\medterm{Epithelial Cells} Cells that cover the skin, line body passages and organs, and form many glands. They create a protective barrier and may absorb substances, release fluids, move material, or sense changes in their surroundings.

\synonyms
Epithelial Tissue Cells

\medterm{Epithelial Cyst} A usually benign, enclosed sac lined by surface-type cells and filled with fluid, mucus, or a skin protein called keratin. It may enlarge, become painful, or become infected and sometimes requires removal.

\medterm{Epithelial Tissue} Tissue made of closely packed cells that covers the skin, lines internal spaces and tubes, and forms glands. It provides protection and may absorb, secrete, transport, or exchange substances.

\medterm{Epithelioid Cell} A rounded cell that resembles a surface-lining cell. The term often describes an activated immune cell in a granuloma, but it can also describe the appearance of some tumors.

\medterm{Epithelioid Histiocyte} An activated macrophage that has undergone structural transformation under the influence of interferon-gamma in chronic granulomatous inflammation, acquiring abundant pale eosinophilic cytoplasm, elongated slipper-shaped nuclei, and indistinct cell borders resembling epithelial cells.

\medterm{Epithelioid Sarcoma} A rare cancer of connective or soft tissue, often beginning in the skin or deeper tissues of the hands, forearms, feet, or lower limbs. It may recur or spread and requires specialist treatment.

\medterm{Epithelioid Trophoblastic Tumor} A rare pregnancy-related cancer arising from cells that normally help form the placenta. It may appear months or years after pregnancy, cause abnormal bleeding, and requires specialist cancer treatment.

\medterm{Epithelioma} An older or broad term for a tumor arising from epithelial cells; the specific diagnosis should identify whether it is benign, precancerous, or malignant.

\medterm{Epitheliopathy} Disease or injury affecting epithelial cells, which cover body surfaces and line organs. The term is nonspecific and may describe conditions of the cornea, skin, gastrointestinal tract, or other epithelial tissues.

\medterm{Epithelium} A layer or sheet of closely packed cells covering the skin, lining body spaces and tubes, or forming glands. It protects surfaces and may absorb substances, release fluids, transport material, or allow exchange.

\medterm{Epitope} The small, specific part of an antigen that is recognized by an antibody or immune-cell receptor. Recognition can trigger an immune response; epitopes are important in infections, vaccines, allergies, and antibody-based treatments.

\medterm{Epitope Mapping} The process of identifying an epitope, the small part of an antigen recognized by an antibody or specific immune-cell receptor. It helps researchers understand immune responses and develop vaccines, diagnostic tests, and antibody-based treatments.

\medterm{Eplerenone} A medicine that blocks aldosterone, a hormone that promotes salt retention and heart remodeling. It treats high blood pressure and selected heart-failure patients; kidney function and blood potassium require monitoring.

\medterm{Epley Maneuver} A sequence of head and body movements that returns displaced inner-ear particles to their usual position. It can relieve brief spinning triggered by head movement in posterior-canal positional vertigo.

\medterm{Epoetin Alfa} A laboratory-made form of erythropoietin that stimulates the bone marrow to produce red blood cells. It treats selected anemias, but can raise blood pressure and the risk of blood clots if hemoglobin rises too high.

\medterm{Epoetin Beta} A laboratory-made form of erythropoietin used for selected anemias, including some caused by chronic kidney disease. Treatment is individualized and requires monitoring because it may cause high blood pressure or blood clots.

\medterm{Eponym} A medical name derived from a person’s name or another proper name, such as the name of a disease, sign, procedure, or body structure. Eponyms may be replaced by descriptive terms to improve clarity.

\medterm{Epoprostenol} A short-acting medicine that widens blood vessels in the lungs and reduces platelet clumping. Continuous intravenous infusion treats severe pulmonary arterial hypertension; interruption can cause rapid worsening.

\medterm{EPO Test} A blood test measuring erythropoietin, a kidney-made hormone that stimulates red-blood-cell production. It may help investigate anemia or an abnormally high red-cell count alongside blood counts and other tests.

\medterm{Epoxide Hydrolase} An enzyme that changes reactive chemicals called epoxides into less reactive substances. It helps process some medicines, toxins, and natural fats and can influence their effects and the risk of tissue injury.

\medterm{Epoxy Compounds} Chemicals containing a reactive group used in adhesives, coatings, plastics, and other materials. Uncured components can irritate the skin and eyes, cause allergic skin inflammation, or irritate the airways; workplace exposure may cause asthma.

\medterm{Epsom Salts} Magnesium sulfate crystals used in some baths and, occasionally, as a laxative. Drinking them can cause diarrhea or magnesium poisoning, especially in people with kidney disease, so medical advice is important.

\medterm{Epstein-Barr Virus} A common virus spread mainly through saliva. It can cause infectious mononucleosis and remains dormant in the body; it is also linked to some lymphomas, other cancers, and immune disorders.

\medterm{Epstein-Barr Virus Antibody Test} A blood test that measures antibodies produced against Epstein-Barr virus. The pattern of antibodies, symptoms, and timing can help distinguish recent infection from past infection, but results need clinical interpretation.

\medterm{Epstein Pearls} Small, white, harmless cysts containing keratin that appear on a newborn’s gums or roof of the mouth. They usually disappear within weeks without treatment and should not be squeezed.

\medterm{Eptifibatide} An injectable antiplatelet medicine used during selected urgent coronary treatments. It prevents platelets from clumping by blocking a platelet-binding receptor, but can cause serious bleeding and requires blood-count and kidney monitoring.

\medterm{Epulis} A localized growth arising from the gum or tissues around a tooth. Most are reactive and benign, but persistent, painful, bleeding, or enlarging growths need dental examination and sometimes biopsy.

\medterm{Equilibrium Receptors} Sensory cells in the inner ear that detect head movement and position. Signals from them, the eyes, and muscles help the brain maintain balance and coordinate posture and eye movements.

\medterm{Equilin} An estrogen hormone found mainly in pregnant mares and included in some mixtures of conjugated estrogens. Estrogen treatment can increase the risk of blood clots and other adverse effects and requires individualized medical advice.

\medterm{Equine Morbillivirus} A virus related to measles that has caused severe respiratory or neurologic disease in horses and can infect other mammals. Suspected cases require veterinary investigation and infection-control measures.

\medterm{Equinovarus Deformity} A deformity in which the foot points downward and turns inward. It may be present at birth or develop from muscle, nerve, joint, or bone problems, limiting walking and requiring cause-specific treatment.

\medterm{Equipment Safety} Measures that reduce harm from medical devices by using the correct equipment, checking it before use, cleaning it properly, maintaining it, and training users. Faulty equipment should be removed from service.

\medterm{Equivalence Trial} A clinical study designed to determine whether a new treatment and an established treatment have effects within a predefined clinically acceptable range of one another. The range must be set before results are analyzed.

\medterm{Eradication} Complete removal of an infection or disease from a person or population. For an individual infection, eradication may require a full treatment course and a follow-up test to confirm that the organism is gone.

\medterm{Erbb Receptors} A family of cell-surface proteins that receive growth signals and regulate cell division, survival, and specialization. Abnormal activity in some members can promote cancers and may guide targeted treatment.

\medterm{Erbium} A chemical element used in specialized lasers and materials. Erbium lasers can precisely remove or remodel selected skin and dental tissue; treatment may cause pain, burns, pigment changes, or scarring.

\medterm{Erb Palsy} A peripheral obstetric or traumatic traction injury to the upper trunk of the brachial plexus involving spinal nerve roots C5 and C6 (and sometimes C7), causing paralysis of the deltoid, supraspinatus, infraspinatus, and biceps brachii muscles, resulting in an adducted, internally rotated shoulder and extended, pronated forearm ('waiter's tip' posture).
\synonyms
Erb-Duchenne Palsy; Upper Trunk Brachial Plexus Palsy

\medterm{Erb's Palsy} Weakness or paralysis of the shoulder and arm caused by injury to the upper nerves of the brachial plexus, often during difficult birth. The arm may hang inward with the elbow straight and palm facing backward.

\medterm{ERCP} A procedure using a flexible camera and X-ray guidance to diagnose or treat blocked, narrowed, or leaking bile or pancreatic ducts. It can remove stones or place a stent, but may cause pancreatitis, bleeding, infection, or perforation.

\medterm{Erdheim-Chester Disease} A rare non-Langerhans cell histiocytic neoplasm driven by activating kinase mutations in the MAPK pathway (predominantly BRAF V600E), characterized pathologically by lipid-laden foamy histiocyte infiltration causing symmetric diaphyseal osteosclerosis of long bones, periaortic fibrosis ('coated aorta'), and perinephric infiltration ('hairy kidney').
\synonyms
ECD; Non-Langerhans Histiocytosis

\medterm{Erectile Disorder} Alternate name; see \textit{Erectile Dysfunction}.

\medterm{Erectile Dysfunction} Repeated difficulty getting or keeping an erection firm enough for satisfactory sexual activity. It may result from blood-vessel disease, diabetes, nerve or hormonal disorders, medicines, penile conditions, anxiety, depression, or mixed causes and can signal cardiovascular risk.

\medterm{Erection} Temporary enlargement and stiffening of the penis when blood fills erectile tissue and drainage through veins decreases. It can occur with sexual stimulation, during sleep, or spontaneously and normally subsides afterward.

\medterm{Erector Spinae} A group of muscles running along both sides of the spine. They straighten and stabilize the back and help bend the spine sideways during standing, walking, and lifting.

\synonyms
Sacrospinalis Muscle; Spinalis and Longissimus Complex

\medterm{Ergocalciferol} Vitamin D2, a form of vitamin D used to prevent or treat deficiency. It helps the body absorb calcium and maintain bones, but excessive doses can cause high blood calcium, nausea, confusion, or kidney injury.

\medterm{Ergolines} A family of compounds related to chemicals made by ergot fungi. Some affect serotonin, dopamine, or adrenaline receptors and are used or studied for migraine, hormone disorders, or neurologic conditions.

\medterm{Ergometrine} A medicine that contracts the uterus and narrows some blood vessels. It may help control severe bleeding after childbirth, but can dangerously raise blood pressure and is unsuitable for some people with hypertension or heart disease.

\medterm{Ergonomics} Designing work, tools, furniture, and environments to match the body’s abilities and limits. Good ergonomics can reduce strain, awkward posture, repetitive-injury risk, fatigue, and errors.

\medterm{Ergosterol} A fat-like substance that helps maintain fungal cell membranes. Many antifungal medicines damage fungi by binding ergosterol or blocking its production; human cells use cholesterol instead.

\medterm{Ergot} A fungus that can infect grains such as rye and produce toxic chemicals. Eating contaminated grain or taking excessive ergot-derived medicines can cause severe blood-vessel narrowing, neurologic symptoms, or uterine contractions.

\medterm{Ergotamine} A medicine that narrows certain blood vessels and can stop some migraine attacks. It has many interactions and should not be overused because it can cause dangerous loss of blood flow to the limbs or organs.

\medterm{Ergotism} Poisoning caused by ergot toxins or excessive ergot-derived medicines. Severe blood-vessel narrowing can cause burning pain, numbness, cold or discolored limbs, confusion, seizures, tissue death, or dangerous uterine contractions.

\medterm{Ergots} Chemicals produced by ergot fungi, some of which are used in medicines for migraine or bleeding. They can strongly narrow blood vessels and interact with other medicines, causing serious circulation problems or poisoning.

\medterm{Eribulin Mesylate} A cancer medicine that interferes with microtubules, structures needed for cell division. It is used for selected advanced breast cancers and liposarcomas; common risks include low blood counts, nerve damage, fatigue, and nausea.

\medterm{Erlotinib Hydrochloride} A cancer medicine that blocks growth signals from the epidermal growth-factor receptor. It treats selected cancers with suitable molecular changes and may cause rash, diarrhea, liver injury, or lung inflammation.

\medterm{Erosion} Loss of only the surface layer of skin or a moist lining, without the deeper tissue loss of an ulcer. Friction, inflammation, infection, injury, or chemicals may cause it and may produce pain or bleeding.

\medterm{Erosion Lesion} A shallow area where the surface layer of skin or a moist lining has been lost. It may result from rubbing, inflammation, infection, injury, or chemicals and usually heals when the cause is removed.

\medterm{Erosive Esophagitis} Inflammation of the swallowing tube with visible breaks in its lining, usually from repeated stomach-acid reflux. It can cause painful swallowing, heartburn, or bleeding and may lead to narrowing if untreated.

\medterm{Erosive Gastritis} Irritation and inflammation of the stomach lining with shallow areas of damage. Common causes include anti-inflammatory pain medicines, alcohol, severe illness, and stress on the body; it may cause pain, vomiting, or gastrointestinal bleeding.

\medterm{Erosive Lichen Planus} A long-lasting immune-mediated inflammation that causes painful raw areas, burning, bleeding, or soreness in the mouth or genital tissues. Scarring and narrowing may develop, so examination and treatment can help protect function.

\synonyms
Vulvovaginal Gingival Syndrome; Erosive Vulvovaginal Lichen Planus

\medterm{Erratic Results} Test results that change unexpectedly or do not agree with repeated measurements. Possible causes include sample problems, device error, normal biological variation, medicines, or a test that does not fit the clinical question.

\medterm{Ertapenem} A once-daily antibiotic for selected serious infections caused by susceptible bacteria. It does not reliably treat Pseudomonas or enterococcal infections and can cause diarrhea, allergic reactions, seizures, or antibiotic-resistant organisms.

\medterm{Erucic Acid} A fatty acid found in some plant oils, including mustard and older varieties of rapeseed oil. Modern edible oils are usually bred to contain little of it; unusually high exposure may harm the heart.

\medterm{Erucic Acids} Long-chain fatty acids found mainly in certain plant oils. The amount consumed depends on the oil, and unusually high exposure may be harmful; ordinary low-erucic edible oils provide much less.

\medterm{Eructation} Release of swallowed air or gas from the stomach through the mouth, commonly called belching. Frequent troublesome belching may be related to swallowed air, reflux, food intolerance, or other digestive problems.

\synonyms
Belching; Burping

\medterm{Eruption Cyst} A soft, fluid-filled swelling over a tooth that is about to come through the gum, often blue or purple. It usually disappears when the tooth erupts, but painful or infected swelling needs dental review.

\medterm{Eruptive Fever} Fever occurring with a widespread rash. Viral infections are common causes, but bacteria, medicines, allergic reactions, and immune disorders can also be responsible; urgent assessment is needed for severe illness or breathing difficulty.

\medterm{Eruptive Xanthomas} Sudden crops of small yellow-red bumps, often on the buttocks, elbows, knees, or trunk, caused by very high blood triglycerides. They signal a serious lipid disorder requiring prompt evaluation and treatment.

\medterm{Eruptive Xanthomatosis} A sudden outbreak of many small yellow-red skin bumps caused by fat collecting in skin immune cells. It is usually linked to very high triglycerides, sometimes from uncontrolled diabetes or an inherited lipid disorder.

\medterm{Eryptosis} A controlled process by which aging or damaged red blood cells shrink and are marked for removal from the bloodstream. It helps maintain blood-cell balance but may increase in some diseases or during oxidative stress.

\medterm{Erysipelas} A sudden bacterial infection of the upper skin and lymph vessels, usually causing a painful, raised, sharply bordered red patch with fever. Prompt antibiotic treatment is important to prevent spread or complications.

\medterm{Erysipeloid} A skin infection caused by bacteria acquired through cuts while handling fish, meat, poultry, or animals. It usually causes a slowly expanding painful purple-red patch on a finger or hand and needs antibiotic treatment.

\medterm{Erysipelothrix} A group of bacteria found in animals and animal products that can enter broken skin and cause a painful localized infection. Rarely, infection spreads to the bloodstream or heart valves.

\medterm{Erythema} Redness of the skin or a moist body lining caused by increased blood flow near the surface. It may result from inflammation, infection, heat, allergy, or another illness and often fades briefly when pressed.

\medterm{Erythema Ab Igne} A net-like red-brown skin discoloration caused by repeated exposure to moderate heat, such as heating pads, hot-water bottles, or heaters. Stopping the exposure may allow improvement; persistent lesions need examination.

\medterm{Erythema Annulare Centrifugum} A reactive gyrate erythema presenting with slowly expanding, polycyclic, annular erythematous plaques characterized by central clearing and a distinctive trailing rim of fine scale along the inner border of the advancing erythematous margin.
\synonyms
EAC; Centrifugal Annular Erythema

\medterm{Erythema Chronicum Migrans} An expanding skin rash that can appear days to weeks after a tick transmits Lyme disease bacteria. It may be evenly red or ring-shaped and does not always look like a classic bull’s-eye.

\medterm{Erythema Dose} The amount of ultraviolet or other radiation needed to produce visible skin redness. The minimal erythema dose is the smallest amount that causes a defined response and varies with skin type, medicines, and exposure conditions.

\medterm{Erythema Elevatum Diutinum} A rare chronic cutaneous leukocytoclastic vasculitis presenting as persistent, symmetrical, red-violaceous to yellowish-brown papules, nodules, and plaques over the extensor surfaces of joints, associated with IgA paraproteinemia, hematologic disorders, and chronic infections.
\synonyms
EED; Chronic Extensor Leukocytoclastic Vasculitis

\medterm{Erythema Gyratum Repens} A rapidly spreading, ringed skin rash with a wood-grain pattern. It is often associated with an internal cancer, especially of the lung, so new or persistent cases require prompt medical evaluation.

\medterm{Erythema Induratum} Recurrent tender lumps caused by inflammation of the fat beneath the skin, usually on the backs of the lower legs. Tuberculosis and other infections, inflammatory diseases, or medicines may be responsible.

\medterm{Erythema Infectiosum} A common childhood infection caused by parvovirus B19, also called fifth disease. Mild fever, tiredness, or cold symptoms may come first, followed by bright red cheeks and a lacy rash on the body and limbs.

\synonyms
Fifth Disease; Slapped Cheek Disease

\medterm{Erythema Marginatum} A temporary, non-itchy, ring-shaped or wavy red rash that usually appears on the trunk and upper limbs. It is a recognized sign of acute rheumatic fever and may come and go over days or weeks.

\medterm{Erythema Migrans} An expanding rash that develops days to weeks after a tick transmits Lyme disease bacteria. It may be uniformly red or ring-shaped and is usually diagnosed from its appearance and exposure history.

\medterm{Erythema Multiforme} A short-lived immune reaction, often triggered by herpes simplex infection or, less commonly, another infection or medicine. It causes symmetrical target-shaped spots, usually on the hands, feet, arms, or face; mouth involvement may occur.

\medterm{Erythema Nodosum} Inflammation of the fat beneath the skin causing tender red or purple lumps, usually on both shins. Infections, medicines, inflammatory bowel disease, sarcoidosis, and other conditions may trigger it.

\medterm{Erythema Toxicum} A common harmless newborn rash with red spots, small bumps, or tiny pus-filled areas. It often appears during the first days of life and disappears on its own without treatment.

\medterm{Erythema Toxicum Neonatorum} A harmless newborn rash that usually appears during the first few days as red blotches, bumps, or tiny pustules on the face, trunk, or limbs. It usually spares the palms and soles and clears within one to two weeks.

\synonyms
Erythema Neonatorum; Toxic Erythema of the Newborn

\medterm{Erythrasma} A superficial bacterial infection in warm, moist skin folds that causes well-defined reddish-brown patches, often in the groin, armpits, or between toes. It usually responds to prescribed antibacterial treatment and keeping folds dry.

\medterm{Erythroblast} An immature red-blood-cell precursor made in the bone marrow. It develops through several stages before becoming a mature red blood cell; abnormal numbers may indicate altered blood-cell production.

\medterm{Erythroblastosis Fetalis} A condition in which antibodies from the pregnant person destroy a fetus’s or newborn’s red blood cells. It can cause anemia, jaundice, widespread swelling, heart failure, or death and is often prevented by blood-group screening and preventive treatment.

\medterm{Erythroblastosis Foetalis} Destruction of fetal or newborn red blood cells by antibodies from the mother. It can cause anemia, jaundice, widespread fluid buildup, heart failure, or death. Pregnancy blood-group testing and preventive treatment can reduce the risk.

\medterm{Erythrocyte} A mature red blood cell that carries oxygen from the lungs to tissues and helps return carbon dioxide to the lungs. Its flexible shape allows it to pass through very small blood vessels.

\medterm{Erythrocyte Casts} Tube-shaped groups of red blood cells seen in urine. They form inside the kidney's small tubes and strongly suggest bleeding or inflammation in the kidney's filtering units, such as glomerulonephritis or vasculitis.
\synonyms
Red Blood Cell Casts; RBC Casts

\medterm{Erythrocyte Enzyme Deficiency} An inherited or acquired lack of a protein that red blood cells need for energy or protection. The cells may break down too early and cause anemia; examples include G6PD and pyruvate-kinase deficiency.

\medterm{Erythrocyte Membrane} The flexible outer layer of a red blood cell. It maintains the cell’s shape, allows it to bend through small vessels, and controls movement of water and substances into and out of the cell.

\medterm{Erythrocyte Sedimentation Rate} A blood test measuring how quickly red blood cells settle in a tube over one hour. A faster rate can occur with inflammation, but it does not identify the cause and is affected by anemia, age, pregnancy, and other conditions.

\medterm{Erythrocytosis} An abnormally high number or volume of red blood cells, reflected by high hemoglobin or hematocrit. It may result from dehydration, low oxygen, excess erythropoietin, or a bone-marrow disorder and can increase clot risk.

\medterm{Erythroderma} Widespread redness and scaling affecting most of the skin. Causes include psoriasis, eczema, medicines, infection, or lymphoma; severe cases can cause fluid loss, temperature problems, infection, and heart strain.

\medterm{Erythrodermic Psoriasis} A rare, potentially life-threatening form of psoriasis causing nearly all the skin to become intensely red, painful, and scaly. Fever, chills, dehydration, infection, and unstable blood pressure may occur and require urgent care.

\medterm{Erythrodysplasia} Abnormal development or maturation of red-blood-cell precursors in the bone marrow. It can produce ineffective blood formation and anemia and may occur with a marrow disorder or other illness.

\medterm{Erythroid Hyperplasia} An increase in immature red-blood-cell precursors in the bone marrow, usually because the body is responding to anemia, blood loss, red-cell destruction, or low oxygen. The underlying cause needs evaluation.

\medterm{Erythroleukoplakia} A mouth lesion containing both red and white areas. It has a substantial risk of advanced precancerous change or cancer, so a dentist or specialist should examine it and usually perform a biopsy.

\medterm{Erythromelalgia} Repeated attacks of burning pain, warmth, and redness in the feet, hands, or sometimes face, often triggered by heat or exercise and eased by cooling. It may be inherited or linked to blood or nerve disorders.

\medterm{Erythromycin} An antibiotic that blocks bacterial protein production. It treats selected infections and can be used when some other antibiotics are unsuitable; nausea, diarrhea, liver injury, and abnormal heart rhythm are possible.

\medterm{Erythromycin Estolate} An oral form of the antibiotic erythromycin. It blocks bacterial protein production and treats selected infections; stomach upset, liver inflammation, allergic reactions, and medicine interactions can occur.

\medterm{Erythromycin Ophthalmic Prophylaxis} Antibiotic ointment placed in a newborn’s eyes soon after birth to reduce the risk of severe eye infection from gonorrhea acquired during delivery. It does not replace testing and treatment of infections during pregnancy.

\medterm{Erythroplakia} A persistent, sharply defined red patch on a moist lining, especially inside the mouth. It may contain advanced precancerous change or cancer and requires prompt examination, usually with biopsy.

\medterm{Erythroplasia} A sharply bordered red patch on skin or a moist lining. Persistent lesions, especially in the mouth or genital area, may represent a serious precancerous change or cancer and need medical assessment.

\medterm{Erythroplasia Of Queyrat} Squamous-cell carcinoma in situ affecting the head of the penis or foreskin. It usually appears as a persistent, well-defined red or velvety patch and can progress to invasive cancer without treatment.

\medterm{Erythropoiesis} The process of making mature red blood cells, mainly in the bone marrow. It is stimulated by erythropoietin when oxygen is low and requires healthy marrow plus adequate iron, vitamin B12, and folate.

\medterm{Erythropoiesis-Stimulating Agents} Medicines that signal the bone marrow to make more red blood cells. They are used for selected anemias, including some caused by kidney disease or cancer treatment, but may increase blood pressure and clot risk.

\medterm{Erythropoietin} A hormone made mainly by the kidneys when body tissues receive too little oxygen. It signals the bone marrow to produce red blood cells; kidney disease can reduce this signal and contribute to anemia.

\medterm{Erythropoietin Deficiency} An insufficient amount of erythropoietin for the body's need, resulting in reduced stimulation of red-blood-cell production in the bone marrow. It commonly contributes to anemia in chronic kidney disease; reduced response to available erythropoietin is a separate phenomenon.

\medterm{Erythropoietin Receptor} A protein on immature red-blood-cell precursors that receives the erythropoietin signal. Activation helps these cells survive, multiply, and mature into red blood cells.

\medterm{Erythropoietin Test} A blood test measuring the kidney-made hormone that stimulates red-blood-cell production. It can help investigate selected anemias or high red-cell counts and is interpreted with the blood count, oxygen status, and kidney function.

\medterm{Erythrosine} A red dye used in some foods, medicines, and laboratory procedures. Exposure is usually small, but regulations vary and people with a known sensitivity should follow product and medical advice.

\medterm{ESBL} An enzyme made by some bacteria that destroys many penicillin and cephalosporin antibiotics. Infections caused by these bacteria can be harder to treat and require laboratory testing to select an effective medicine.

\medterm{ESBL-Producing Organisms} Bacteria that make enzymes capable of breaking down many penicillin and cephalosporin antibiotics. They commonly include certain intestinal bacteria such as E. coli and Klebsiella; treatment depends on susceptibility testing and infection site.

\medterm{Escape Reaction} A body response that restores an organ’s activity after it has been reduced by a medicine, hormone, nerve signal, or feedback system. The response may lessen the effect of continued treatment or stimulation.

\medterm{Eschar} A firm layer of dead tissue, often black or brown, over a burn, pressure injury, wound, infection, or area with poor blood supply. It can hide deeper damage and may need specialist assessment or removal.

\medterm{Escharotomy} A surgical cut through stiff dead tissue around a deep burn to release pressure. It may restore circulation to a limb or allow the chest to expand, and is performed urgently when swelling threatens function.

\medterm{Escherichia} A group of bacteria that includes E. coli and related species found in intestines or the environment. Some strains are harmless, while others cause diarrhea, urinary infection, wound infection, or bloodstream infection.

\medterm{Escherichia Albertii} An uncommon intestinal bacterium that can cause diarrheal illness. It may be mistaken for certain E. coli strains, so specialized laboratory identification may be needed when infection is suspected.

\medterm{Escherichia Blattae} A rarely reported Escherichia species associated mainly with insects and the environment. Its significance in a human sample is uncertain and must be judged from symptoms, the specimen source, and other test results.

\medterm{Escherichia Coli} A bacterium that normally lives in the intestine. Some strains are harmless, while others cause urinary, intestinal, wound, or bloodstream infections, including severe bloody diarrhea and hemolytic uremic syndrome.

\medterm{Escherichia Coli Infection} Infection caused by a harmful strain of E. coli. It may affect the intestine, urinary tract, wounds, or bloodstream; symptoms, contagiousness, and treatment depend on the strain and site.

\medterm{Escherichia Coli O157:H7 Infection} Infection with a toxin-producing E. coli strain, usually acquired from contaminated food, water, animals, or people. It commonly causes severe cramps and bloody diarrhea and can lead to kidney failure; urgent medical advice is important.

\medterm{Escherichia Fergusonii} An uncommon bacterium that can occasionally cause urinary, wound, intestinal, or bloodstream infection, especially in people with serious illness or weakened immunity. Laboratory identification and susceptibility testing guide treatment.

\medterm{Escherichia Hermannii} An uncommon Escherichia species that rarely causes opportunistic infection, including urinary or bloodstream infection. Its importance depends on the patient’s symptoms, immune status, and the specimen in which it is found.

\medterm{Escherichia Vulneris} An uncommon bacterium rarely linked to wound, urinary, or bloodstream infection, usually in people with other serious illnesses. A laboratory result requires clinical correlation to distinguish infection from contamination.

\medterm{Escitalopram} An antidepressant that increases available serotonin in the brain. It treats depression and generalized anxiety disorder; nausea, sleep changes, sexual problems, withdrawal symptoms, bleeding risk, and rarely serotonin toxicity may occur.

\medterm{ESKD} Severe permanent loss of kidney function, usually meaning kidney filtration below about 15 milliliters per minute or a need for dialysis or transplantation. It can cause fluid, electrolyte, anemia, and toxin problems and needs kidney-specialist care.

\medterm{Esomeprazole} A medicine that reduces stomach-acid production by blocking the acid pump in stomach cells. It treats reflux, ulcers, and acid-related disorders; prolonged or inappropriate use may cause infections, low magnesium, or vitamin deficiencies.

\medterm{Esomeprazole Magnesium} A magnesium salt of esomeprazole that reduces stomach-acid production by blocking the stomach’s acid pump. It treats reflux and ulcers; long-term use may cause diarrhea, infections, low magnesium, or nutrient deficiencies.

\medterm{Esophageal} Relating to the esophagus, the muscular tube that carries swallowed food and liquid from the throat to the stomach.

\medterm{Esophageal Achalasia} A disorder in which the lower esophageal valve does not relax properly and the esophagus cannot push food normally into the stomach. It causes difficulty swallowing, regurgitation, chest discomfort, and weight loss.

\medterm{Esophageal Adenocarcinoma} A cancer arising from gland-forming cells, usually in the lower esophagus and often linked to Barrett esophagus. Reflux, obesity, and smoking increase risk; progressive swallowing difficulty or weight loss requires evaluation and biopsy.

\medterm{Esophageal Atresia} A birth defect in which the esophagus ends in a blind pouch instead of connecting normally to the stomach. A connection to the windpipe may also occur, causing drooling, choking, coughing, and breathing problems during feeds.

\medterm{Esophageal Bleeding} Bleeding from the esophagus, which may result from enlarged veins, inflammation, ulcers, tears, injury, or cancer. Vomiting blood, passing black stools, dizziness, or fainting requires emergency medical care.

\medterm{Esophageal Cancer} A malignant tumor of the swallowing tube. The main types are squamous-cell cancer and adenocarcinoma; progressive difficulty swallowing, weight loss, pain, or bleeding may occur, and treatment depends on the stage.

\synonyms
Esophageal Carcinoma; Esophagus Malignancy

\medterm{Esophageal Candidiasis} A yeast infection of the esophagus, more common with weakened immunity, diabetes, or certain medicines. It often causes painful or difficult swallowing and is treated with antifungal medicine, sometimes after endoscopy confirms it.

\medterm{Esophageal Carcinoma} Cancer arising in the esophagus, most often squamous-cell carcinoma or adenocarcinoma. Progressive difficulty swallowing, weight loss, chest pain, or bleeding may occur; endoscopic biopsy and staging guide treatment.

\medterm{Esophageal Culture} Laboratory growth testing on a sample from the esophagus to identify bacteria or fungi when infection is suspected. Results depend on how the sample was collected and must be interpreted with symptoms and endoscopic findings.

\medterm{Esophageal Dilation} A procedure that gently stretches a narrowed area of the esophagus with a tube or balloon passed through an endoscope. It can improve swallowing caused by scarring, inflammation, injury, or other narrowing; tearing and bleeding are uncommon but serious risks.

\medterm{Esophageal Diseases} Disorders affecting the swallowing tube between the throat and stomach. They include reflux, inflammation, narrowing, movement problems, cancer, infection, and birth defects, often causing swallowing difficulty, pain, heartburn, or regurgitation.

\medterm{Esophageal Diverticulum} A pouch that bulges outward from the wall of the esophagus. It may cause difficulty swallowing, food coming back up, coughing, bad breath, or no symptoms; treatment depends on its size, location, and effects.

\medterm{Esophageal Duplication} A rare birth defect consisting of a cyst or extra segment beside the esophagus. It may press on the esophagus or airway, causing swallowing difficulty, breathing problems, infection, or no symptoms.

\medterm{Esophageal Fistula} An abnormal connection between the esophagus and another organ, often the airway. It can allow food or liquid into the lungs, causing coughing during swallowing, repeated pneumonia, poor nutrition, or severe illness.

\medterm{Esophageal Foreign Body} Food or another object stuck in the esophagus, causing pain, difficulty swallowing, drooling, or breathing problems. A battery, sharp object, complete blockage, or suspected tear requires emergency removal.

\medterm{Esophageal Infection} Infection of the esophagus caused by a yeast, virus, or bacterium, usually in people with weakened immunity or damaged lining. It may cause painful swallowing, difficulty swallowing, chest pain, or ulcers.

\medterm{Esophageal Intramural Pseudodiverticulosis} A rare condition in which small drainage ducts of glands within the esophageal wall become widened. It may cause repeated swallowing difficulty, food impaction, or narrowing and is diagnosed with imaging or endoscopy.

\medterm{Esophageal Leukoplakia} A persistent white patch on the esophageal lining caused by thickened surface cells. It is uncommon but may contain precancerous change, so evaluation and biopsy are often needed to determine its significance.

\medterm{Esophageal Manometry} A test that measures pressure and muscle coordination in the esophagus during swallowing. A thin pressure-sensitive tube passes through the nose into the stomach and helps diagnose achalasia, spasm, and other movement disorders.

\medterm{Esophageal Motility Disorder} A disorder of the muscle contractions or nerve control that moves food through the esophagus. It may cause difficulty swallowing, chest pain, regurgitation, or food sticking; testing may include esophageal manometry.

\medterm{Esophageal Necrosis} Death of esophageal lining, sometimes producing a black appearance. It can result from severe low blood flow, shock, acid injury, or caustic substances and may lead to bleeding or perforation.

\medterm{Esophageal Perforation} A hole through the full thickness of the esophageal wall, caused by forceful vomiting, an instrument, injury, or disease. It can release contents into the chest, causing severe infection and requiring emergency treatment.

\medterm{Esophageal pH Monitoring} A test that records how much acid enters the esophagus, usually for 24 hours, using a thin tube or a small wireless sensor. It helps determine whether acid reflux explains persistent symptoms when the diagnosis is uncertain.

\medterm{Esophageal Spasm} A movement disorder in which esophageal muscles contract too strongly or without normal coordination. It can cause intermittent chest pain and difficulty swallowing; pressure testing helps distinguish it from other disorders.

\medterm{Esophageal Spasms} Uncoordinated or forceful contractions of the esophagus that may cause sudden chest pain, food sticking, or difficulty swallowing. Symptoms can resemble heart pain and should be assessed when severe or new.

\medterm{Esophageal Squamous Cell Carcinoma} A cancer arising from the flat cells lining the esophagus, often in its upper or middle part. Smoking, heavy alcohol use, certain infections, and nutritional factors increase risk; diagnosis requires endoscopic biopsy.

\medterm{Esophageal Stenosis} Narrowing of the passage through the esophagus caused by scarring, inflammation, a tumor, radiation, or a birth defect. It commonly causes progressively worsening difficulty swallowing, especially solid foods.

\medterm{Esophageal Stent} A tube placed inside a narrowed or leaking esophagus to keep it open and allow swallowing. It is often used for cancer or injury; movement, blockage, bleeding, pain, or perforation can occur.

\medterm{Esophageal Stricture} A narrowed area of the esophagus, commonly caused by acid reflux, caustic injury, inflammation, radiation, or procedures. It usually causes gradually worsening difficulty swallowing solids and may require stretching during endoscopy.

\medterm{Esophageal Surgery} An operation on the esophagus or the area where it meets the stomach. It may treat cancer, severe narrowing, perforation, reflux, or movement disorders; the procedure and risks depend on the disease and overall health.

\medterm{Esophageal Ulceration} An open sore in the esophageal lining caused by reflux, infection, medicines, caustic injury, inflammation, or cancer. It may cause painful swallowing, chest pain, or bleeding and often requires endoscopic evaluation.

\medterm{Esophageal Varices} Enlarged veins in the esophageal wall, usually caused by high pressure in the portal vein from advanced liver disease. They may rupture without warning and cause life-threatening vomiting of blood.

\synonyms
Esophageal Dilated Veins; Portosystemic Varices

\medterm{Esophageal Varices Cross-Reference} Alternate name; see \textit{Esophageal varices}.

\medterm{Esophageal Web} A thin shelf of tissue projecting into the upper esophagus. It may cause difficulty swallowing solid food and can be associated with iron-deficiency anemia; persistent symptoms may require endoscopy.

\medterm{Esophagectomy} Surgery to remove part or all of the esophagus, usually for cancer or severe irreversible damage. The remaining digestive tract is reconstructed, commonly using the stomach or a segment of intestine.

\medterm{Esophagitis} Inflammation or injury of the esophageal lining. Common causes include acid reflux, infection, medicines that irritate the esophagus, eosinophilic or other immune-mediated disease, swallowed irritants, and radiation. Symptoms may include heartburn, painful or difficult swallowing, chest pain, food sticking, or bleeding; persistent inflammation can cause ulcers or narrowing.

\medterm{Esophagogastric Bypass} Surgery that creates a new route between the esophagus and stomach to restore passage of food, usually after esophageal removal or severe blockage. The method depends on anatomy, disease, and surgical goals.

\medterm{Esophagogastric Junction} The area where the esophagus joins the stomach. Its muscular valve and surrounding structures help food enter the stomach while limiting backward flow of stomach contents.

\medterm{Esophagogastric Tamponade} Temporary control of severe bleeding from esophageal or stomach veins by inflating a balloon device to apply pressure. It is a rescue measure used while definitive treatment is arranged and can obstruct breathing or cause aspiration.

\medterm{Esophagogastroduodenoscopy} An upper endoscopic diagnostic and therapeutic procedure in which a thin, flexible video endoscope is passed through the mouth to examine the lining of the esophagus, stomach, and duodenum.

\synonyms
EGD; Upper Endoscopy; Upper Gastrointestinal Endoscopy

\medterm{Esophagoscope} A flexible or rigid instrument used to look inside the esophagus. It may allow a clinician to take tissue samples, remove a foreign object, treat bleeding, or perform another procedure.

\medterm{Esophagoscopy} Examination of the esophagus with a camera-equipped instrument. It may investigate swallowing difficulty, bleeding, reflux damage, narrowing, foreign objects, infection, or suspected cancer and may allow treatment or biopsy.

\medterm{Esophagram} An X-ray examination performed while a person swallows contrast liquid. It shows the shape and movement of the esophagus and can reveal narrowing, leaks, pouches, blockage, or abnormal swallowing.

\medterm{Esophagus} A muscular tube that carries food and liquid from the throat to the stomach. Coordinated muscle contractions move swallowed material downward, while valves at each end help limit reflux and air entry.

\medterm{Esophagus Disease} A disorder of the esophagus causing swallowing difficulty, pain, reflux, regurgitation, bleeding, obstruction, inflammation, motility disturbance, or cancer.

\medterm{Esophagus Neoplasm} A benign, precancerous, or malignant growth in the esophagus; evaluation may require endoscopy, biopsy, imaging, and staging.

\medterm{Esotropia} A condition in which one eye turns inward toward the nose, causing the eyes to point in different directions. It may begin in infancy or result from focusing problems, illness, or injury; glasses, patching, or surgery may help depending on the cause.

\medterm{ESR} The erythrocyte sedimentation rate, a nonspecific blood test measuring how quickly red cells settle in a tube; it can reflect inflammation but cannot identify its cause.

\synonyms
SED Rate

\medterm{ESRD} Alternate index label; see \textit{End-Stage Renal Disease}.

\medterm{Essential} Necessary or fundamental; in a diagnosis such as essential hypertension, it indicates that no single secondary cause has been identified.

\medterm{Essential Fats} Fatty acids the body cannot make in sufficient amounts and must obtain from food. They support cell membranes, skin, brain function, and production of signaling molecules involved in inflammation.
\synonyms
Essential Fatty Acids

\medterm{Essential Fatty Acid} A fatty acid the body cannot make in sufficient amounts and therefore must obtain from food; the main human examples are linoleic acid and alpha-linolenic acid.

\medterm{Essential Hypertension} Long-term high blood pressure with no single identifiable underlying cause. Many genetic, kidney, blood-vessel, metabolic, and lifestyle factors contribute, and untreated disease raises the risk of stroke, heart disease, kidney damage, and vision problems.

\medterm{Essential Medicines} Medicines considered necessary to meet the most important health needs of a population. They are selected for proven benefit, safety, affordability, and practical availability and should be reliably accessible through healthcare services.

\medterm{Essential Mixed Cryoglobulinemia} An immune disorder in which abnormal proteins in the blood thicken or clump in the cold and inflame small blood vessels. It may cause purple spots, joint pain, weakness, nerve problems, or kidney disease.

\medterm{Essential Nutrient} A nutrient the body cannot make in sufficient amounts and must obtain from food or supplements. Essential nutrients include certain amino acids, fatty acids, vitamins, minerals, and water.

\medterm{Essential Thrombocythemia} A long-term bone-marrow disorder that produces too many platelets. It can cause abnormal clotting or bleeding and may cause headaches, vision changes, burning pain in fingers or toes, or no symptoms.

\medterm{Essential Tremor} A nervous-system disorder causing rhythmic shaking, usually of both hands during movement or when holding a position. It may also affect the head or voice, gradually worsen, and interfere with writing, eating, or other tasks.

\medterm{Essiac} A herbal mixture promoted as a treatment for cancer and other illnesses. Reliable evidence has not shown that it cures cancer, and its ingredients may interact with medicines or cause liver, kidney, or other harm.

\medterm{Estazolam} A sedative medicine sometimes used short-term for insomnia. It slows brain activity but can cause daytime drowsiness, falls, memory problems, slowed breathing, dependence, and withdrawal; alcohol and opioids increase its dangers.

\medterm{Ester} A chemical formed when an acid combines with an alcohol. Esters occur naturally in fats, medicines, and flavorings; their health effects depend on the specific compound, amount, route of exposure, and duration.

\medterm{Esterase} An enzyme that breaks ester chemical bonds. Esterases help process fats, medicines, and other substances; abnormal activity or inhibition can alter drug effects and may contribute to poisoning.

\medterm{Esterified Estrogens} A mixture of estrogen hormones used for selected menopausal symptoms or estrogen deficiency. Treatment can cause bleeding and increase risks of blood clots, stroke, and breast or uterine cancer, so benefits and risks must be individualized.

\medterm{Esthesia} The ability to sense touch, pain, temperature, movement, or body position. Reduced, increased, or altered sensation may result from nerve injury, brain disease, medicines, or changes in consciousness.

\medterm{Esthesiometer} An instrument used to measure sensitivity to touch, pressure, or other sensations. Eye clinicians may use one to assess corneal sensation when nerve injury or an eye-surface problem is suspected.

\medterm{Esthesioneuroblastoma} A rare cancer arising from nerve-related cells high in the nasal cavity. It may cause blocked or bleeding nose, reduced smell, facial pain, or vision symptoms and is diagnosed with imaging and biopsy.

\medterm{Estimated Average Glucose} An estimate of a person’s average blood-glucose level calculated from the hemoglobin A1c result. It is reported in familiar glucose units but does not show daily highs, lows, or sudden changes.

\medterm{Estimated Date Of Confinement} The approximate date when a pregnancy is expected to reach delivery, calculated from the last menstrual period or an early ultrasound. Only a small proportion of births occur on the estimated date.

\medterm{Estimated Due Date} The approximate date of delivery, usually calculated as 40 weeks from the first day of the last menstrual period and adjusted when needed using an early ultrasound. Most births occur before or after this date.

\medterm{Estimated Fetal Weight} An ultrasound-based estimate of a fetus’s weight using measurements of the head, abdomen, and thigh bone. It helps assess growth and plan care, but commonly differs from actual birth weight by about 10–15 percent.

\medterm{Estimated Glomerular Filtration Rate} An estimate of how well the kidneys filter blood, calculated mainly from blood creatinine and other factors. It helps assess kidney disease and adjust medicines but is less reliable during rapidly changing illness or unusual body size.

\medterm{Estradiol} The main active estrogen hormone during most reproductive years. It helps regulate ovulation and menstruation and supports the vagina, uterus, bones, brain, and other tissues; levels fall after menopause.

\synonyms
17$\beta$-Estradiol

\medterm{Estradiol Blood Test} A blood test measuring estradiol. It may help evaluate puberty, menstrual or fertility problems, ovarian or testicular function, menopause, or selected hormone disorders; results vary with age, sex, cycle timing, pregnancy, and medicines.

\medterm{Estramustine} A cancer medicine combining an estrogen-like component with a compound that interferes with cell division. It has been used for selected advanced prostate cancers; blood clots, fluid retention, nausea, and heart effects may occur.

\medterm{Estriol} An estrogen produced mainly by the placenta during pregnancy from substances made by the fetus and mother. Maternal blood estriol may be included in prenatal screening, but an abnormal result does not by itself diagnose a problem.

\medterm{Estrogen} A group of hormones involved in reproductive development, menstruation, fertility, bone strength, and many other body functions. The main forms and amounts vary with age, sex, pregnancy, and ovarian function.

\medterm{Estrogen Overdose} Taking or receiving more estrogen than prescribed. It may cause nausea, vomiting, breast tenderness, headache, or abnormal bleeding and can increase the risk of blood clots; urgent advice is needed after a large exposure.

\medterm{Estrogen Receptor} A cell protein that binds estrogen and changes cell activity. Testing for estrogen receptors helps classify some cancers, especially breast cancer, and can guide hormone-based treatment.

\medterm{Estrogen Replacement Therapy} Treatment supplying estrogen to relieve selected menopausal symptoms or replace deficient hormone. If the uterus is present, a progestogen is usually added to protect its lining; risks and benefits must be reviewed individually.

\medterm{Estrogen-Replacement Therapy} Treatment that supplies estrogen when natural levels are low, such as after menopause or ovarian failure. It may relieve hot flushes and protect bone, but can cause bleeding and increase some health risks.

\synonyms
Estrogen Therapy

\medterm{Estrogens} Hormones, including estradiol, that regulate reproductive development, menstruation, fertility, bones, and other tissues. Before menopause they are produced mainly by the ovaries; after menopause, levels are lower and come from several tissues.

\medterm{Estrogen Test} A blood, urine, or saliva test measuring one or more estrogen hormones. It may help evaluate puberty, menstrual problems, fertility, menopause, pregnancy, or hormone disorders; interpretation depends on the test type and timing.

\synonyms
Estrogen Levels Test
Estrogen Blood Test

\medterm{Estrone} A naturally occurring estrogen made partly by converting other hormones in body tissues. It is the main estrogen present after menopause, although estradiol is usually more active before menopause.

\medterm{Estrous Cycle} The recurring reproductive cycle of many female mammals other than humans. It includes hormonal and physical changes that prepare for ovulation, mating, and possible pregnancy.

\medterm{Eszopiclone} A sleep medicine that enhances inhibitory signaling in the brain. It can help insomnia but may cause drowsiness, impaired coordination, memory problems, dependence, unusual sleep behaviors, or a bitter taste.

\medterm{Etacrynic Acid} A diuretic that makes the kidneys remove more salt and water, reducing swelling and sometimes blood pressure. It can cause dehydration, abnormal electrolytes, kidney problems, or hearing damage, especially at high doses.

\medterm{Etanercept} A biologic medicine that blocks tumor-necrosis factor, an inflammation signal. It treats selected arthritis, psoriasis, and other inflammatory diseases but can increase serious infection risk, so tuberculosis and hepatitis screening may be needed.

\medterm{Ethacrynic Acid} A loop diuretic that increases kidney excretion of salt and water. It can reduce fluid overload but may cause dehydration, low electrolytes, kidney injury, or dose-related hearing damage.

\medterm{Ethambutol} An antibiotic used with other medicines to treat tuberculosis and selected infections caused by related bacteria. It can damage the optic nerve, causing blurred vision or trouble distinguishing colors, so vision testing is important.

\medterm{Ethambutol Toxicity} A dose- and duration-dependent adverse pharmacological effect of the first-line antimycobacterial agent ethambutol, characterized by retrobulbar or axial optic neuritis, presenting with bilateral painless reduction in visual acuity, central scotomas, and loss of red-green color discrimination.

\medterm{Ethane} A colorless, highly flammable gas used mainly to make industrial chemicals. It is not usually poisonous through direct action, but high concentrations can displace oxygen and cause suffocation or fire hazards.

\medterm{Ethanol} The alcohol in alcoholic drinks and some disinfectants, medicines, and fuels. Drinking too much can impair judgment, suppress breathing, lower blood sugar, cause dependence, and damage the liver and other organs.

\medterm{Ethanolamine} A chemical used in some industrial products, medicines, and cosmetic ingredients. Concentrated exposure can irritate or burn the skin, eyes, and airways; health effects depend on the product and amount of exposure.

\medterm{Ethanol Metabolism} The body’s breakdown of alcohol, mainly in the liver, into acetaldehyde and then acetate. The process affects blood-alcohol levels, drug interactions, intoxication, hangovers, and long-term alcohol-related organ damage.

\medterm{Ethanol Poisoning} Harm from drinking a dangerous amount of alcohol, causing confusion, poor coordination, vomiting, slow or irregular breathing, low blood sugar, seizures, coma, or death. An unconscious or poorly breathing person needs emergency care.

\medterm{Ether} A family of chemicals containing an oxygen atom between carbon groups. Some are volatile solvents; diethyl ether was once an anesthetic but is highly flammable, and inhalation can cause irritation, dizziness, or unconsciousness.

\medterm{Ethers} A group of chemicals containing an oxygen atom between two carbon groups. Some are used as solvents or fuels; exposure can cause irritation, dizziness, unconsciousness, or fire and explosion hazards.

\medterm{Ethics} Principles used to decide what is right, fair, respectful, and responsible in healthcare and research. Medical ethics includes informed choice, privacy, avoiding harm, helping patients, and treating people fairly.

\medterm{Ethinyl Estradiol} A synthetic estrogen used in many hormonal contraceptives and some other treatments. It can cause nausea or irregular bleeding and increases blood-clot risk, especially with smoking, migraine with aura, or certain medical conditions.

\medterm{Ethiodized Oil} An iodine-containing oil used as a contrast material in selected imaging or treatment procedures, including some lymphatic and liver procedures. It can affect thyroid function or rarely enter blood vessels and cause an embolism.

\medterm{Ethion} An organophosphate pesticide that can poison the nervous system by blocking acetylcholinesterase. Exposure may cause sweating, vomiting, pinpoint pupils, muscle twitching, weakness, breathing failure, or death and requires urgent treatment.

\medterm{Ethionamide} An antibiotic used with other medicines for drug-resistant or difficult-to-treat tuberculosis. It interferes with mycobacterial cell-wall production and may cause nausea, liver injury, nerve symptoms, or thyroid problems.

\medterm{Ethmocephaly} A severe birth defect in which the eyes are very close together and the nose is replaced or accompanied by a tube-like structure. It usually occurs with major brain-development abnormalities and is often life-limiting.

\medterm{Ethmoidal Air Cells} Multiple small air-filled spaces within the ethmoid bone between the nose and eye sockets. They are part of the sinus system and can become blocked or infected, causing sinus pressure and discharge.

\medterm{Ethmoidal Sinus} A group of small, thin-walled air spaces in the ethmoid bone between the eye sockets and above the nasal cavity. They are divided into front and back groups according to where they drain.

\synonyms
Ethmoid Air Cells; Paranasal Ethmoid Sinus

\medterm{Ethmoid Bone} A light, complex bone between the eye sockets that forms part of the nasal septum, nasal roof, and inner eye socket. It contains air cells and openings for smell nerves, so injury can affect smell or the eye.

\medterm{Ethmoiditis} Inflammation or infection of the ethmoid air spaces between the nose and eyes. It may cause fever, nasal blockage, discharge, facial pain, or swelling and can rarely spread to the eye socket or brain.

\medterm{Ethmoid Sinus} One of several small air spaces in the ethmoid bone between the nose and eye sockets. Inflammation or infection can cause nasal blockage, headache, or eye-area pain and may spread to nearby tissues.

\medterm{Ethmoid Sinusitis} Inflammation or infection of the air spaces between the nose and eyes. It may cause congestion, discharge, fever, headache, or pain around the eyes; eye swelling, vision changes, or severe illness requires urgent care.

\medterm{Ethnicity} A social and cultural identity shaped by shared ancestry, history, language, traditions, or community; it is not a biological diagnosis and should not be used as a substitute for individual risk assessment.

\medterm{Ethoglucid} An older sclerosing substance used historically to treat selected vascular lesions or hemorrhoids; modern use is limited by tissue-injury risks and alternatives.

\medterm{Ethosuximide} An anticonvulsant used mainly for absence seizures; common adverse effects include gastrointestinal symptoms, fatigue, and headache, with rare blood or skin reactions.

\medterm{Ethoxzolamide} A carbonic-anhydrase-inhibiting medicine that reduces formation of certain body fluids. It has been used for selected eye and metabolic conditions; possible effects include tingling, nausea, kidney stones, and electrolyte changes.

\medterm{Ethyl Chloride} A volatile liquid sprayed on skin to produce brief cooling and numbness before minor procedures, but misuse can cause cold injury or toxicity.

\medterm{Ethylene} A colorless, flammable gas used in chemical manufacturing and to ripen fruit. High concentrations can displace oxygen and cause headache, dizziness, drowsiness, unconsciousness, or suffocation. Remove a person from exposure only when it is safe and seek emergency help for breathing or consciousness problems.

\medterm{Ethylene Chlorohydrin} A toxic industrial chemical used in some manufacturing processes. Exposure can irritate the eyes, skin, and airways and may affect the nervous system, liver, or kidneys. It can be absorbed through skin, so contaminated clothing should be removed safely and urgent poison-control advice obtained.

\medterm{Ethylene Dibromide} A toxic chemical formerly used as a soil fumigant and fuel additive. It can irritate the eyes and lungs and may damage the liver, kidneys, or nervous system; long-term exposure is associated with cancer risk. Suspected exposure needs urgent poison-control or occupational-health assessment.

\medterm{Ethylene Glycol} A toxic alcohol found in some antifreeze and industrial products. Ingestion can cause severe metabolic acidosis, neurologic impairment, kidney injury, and death.

\medterm{Ethylene Glycol Blood Test} A blood test measuring ethylene glycol, a toxic alcohol found in some antifreeze products. Detection supports diagnosis of poisoning, but treatment must not wait for the result when the exposure and metabolic findings are strongly suggestive.

\medterm{Ethylene Glycol Poisoning} Poisoning after swallowing ethylene glycol, found in some antifreeze products. Early effects may resemble alcohol intoxication, followed by severe acid buildup, seizures, coma, and kidney failure; emergency antidote and sometimes dialysis are needed.

\medterm{Ethylene Oxide} A flammable gas used to sterilize heat-sensitive medical equipment and make other chemicals. It can irritate the skin and airways, cause nerve injury, and with repeated exposure increase cancer risk. Sterilized equipment must be adequately aerated before use, and workplace exposure requires strict controls.

\medterm{Ethyl Ether} A highly flammable liquid formerly used as an anesthetic and still used as a laboratory solvent. Breathing its vapor can cause irritation, dizziness, drowsiness, or loss of consciousness; swallowing it can injure the lungs if it is aspirated. Exposure requires fresh air, fire precautions, and medical assessment when symptoms are severe.

\medterm{Ethynodiol Diacetate} A synthetic progestogen formerly used in some oral contraceptives. It helps prevent pregnancy mainly by suppressing ovulation and changing cervical mucus and the uterine lining.

\medterm{Etidronic Acid} A medicine related to bisphosphonates that can slow the breakdown of bone by attaching to bone mineral. It has also been used to bind minerals in some nonmedical products; effects and safety depend on the use and dose.

\medterm{Etilefrine} A medicine that stimulates some of the body’s adrenaline-like receptors and can raise blood pressure. It increases blood-vessel tightening and heart activity, so it may cause fast heartbeat, headache, or other effects and requires medical supervision.

\medterm{Etiology} The cause or set of causes of a disease or health condition, including genetic, infectious, environmental, immune, traumatic, and behavioural factors.

\medterm{Etodolac} A nonsteroidal anti-inflammatory drug used for pain and inflammation; gastrointestinal bleeding, kidney injury, cardiovascular events, and hypersensitivity are important risks.

\medterm{Etomidate} A short-acting intravenous anesthetic used for induction, particularly when cardiovascular stability is desired; it can cause myoclonus, nausea, and transient adrenal suppression.

\medterm{Etonogestrel} A man-made form of the hormone progesterone used in a small contraceptive implant placed under the skin of the upper arm. It mainly prevents ovulation and also thickens cervical mucus; irregular bleeding and other hormone-related effects can occur.

\medterm{Etoposide} A chemotherapy medicine that interferes with DNA repair and cell division. It treats several cancers but can cause low blood counts, infection, hair loss, nausea, mouth sores, or a later risk of leukemia.

\medterm{Etoposide Phosphate} A water-soluble form of etoposide that changes into etoposide inside the body. It provides another method of giving the chemotherapy medicine, with similar cancer treatment effects and potential toxicities.

\medterm{Etoricoxib} An anti-inflammatory pain medicine that selectively blocks an enzyme involved in pain and swelling. It may help arthritis or acute pain but can raise blood pressure and increase heart, kidney, fluid-retention, and stomach risks.

\medterm{Eubacterium} A group of bacteria that usually live without oxygen in the mouth and intestine. Most are part of normal microbial communities, but some can rarely cause opportunistic infection in seriously ill people.

\medterm{Eubacterium Limosum} A bacterium that lives without oxygen in the intestine and helps break down nutrients. It is rarely associated with human infection, and its importance in a sample depends on symptoms and the collection site.

\medterm{Eubacterium Rectale} A common intestinal bacterium that helps ferment dietary fiber and produce short-chain fatty acids. These products support the colon and contribute to normal gut microbial activity.

\medterm{Eucalyptus Oil} A concentrated plant oil containing eucalyptol and other chemicals. Swallowing it can poison the nervous system and cause vomiting, drowsiness, seizures, or breathing problems; it should not be taken internally without medical advice.

\medterm{Eucalyptus Oil Overdose} Poisoning after swallowing too much eucalyptus oil. Burning, vomiting, drowsiness, confusion, seizures, and breathing problems may occur; do not induce vomiting and seek emergency poison-control advice.

\medterm{Euchromatin} A relatively open form of DNA packaging inside the cell nucleus. Its structure makes nearby genes more accessible for reading and activity than DNA in tightly packed regions.

\medterm{Eugenicist} A person who supports eugenics, a discredited belief that society should control reproduction to alter inherited traits in a population. Eugenic policies caused coercion, discrimination, and forced sterilization and have no place in ethical healthcare.

\medterm{Eugenics} A discredited ideology that seeks to change a population by encouraging some people to reproduce and preventing others from doing so. Its history includes racism, forced sterilization, institutional abuse, and genocide.

\medterm{Eugenol Oil Overdose} Poisoning after excessive exposure to clove-derived eugenol oil. It may irritate the mouth and stomach and cause liver injury, bleeding, seizures, low blood pressure, slowed breathing, or coma.

\medterm{Euglycemic Diabetic Ketoacidosis} A severe metabolic complication of diabetes mellitus characterized by profound anion gap metabolic acidosis and elevated serum ketones occurring with normal or only mildly elevated blood glucose levels (less than 250 mg/dL), frequently triggered by SGLT2 inhibitor therapy, restricted carbohydrate intake, or acute physiological stress.
\synonyms
EuDKA; Normal Glycemic Ketoacidosis

\medterm{Euphoria} An unusually intense feeling of happiness or well-being. It may occur normally, with mania, brain disease, or some medicines and drugs; sudden extreme euphoria with unsafe behavior needs clinical assessment.

\medterm{Eustachian Tube} A passage connecting the middle ear to the back of the nose. It opens during swallowing or yawning to equalize pressure, drain fluid, and help protect the middle ear from throat secretions.

\medterm{Eustachian Tube Dysfunction} Poor opening or closing of the passage between the middle ear and back of the nose. It can cause ear pressure, popping, muffled hearing, pain, or fluid after colds, allergies, or pressure changes.

\medterm{Eustachian-Tube Dysfunction} Failure of the passage between the middle ear and back of the nose to equalize pressure or clear fluid. It may cause fullness, popping, reduced hearing, or ear pain after infection, allergies, or pressure changes.

\medterm{Eustachian Tube Patency} The ability of the passage between the middle ear and back of the nose to remain open enough to equalize pressure and drain fluid. Poor patency can cause fullness, muffled hearing, or middle-ear fluid.

\medterm{Eustachian Tube Problems} Problems with the ear-to-nose passage can cause fullness, popping, pain, muffled hearing, or repeated middle-ear fluid or infection. Colds, allergies, inflammation, and anatomy may contribute; persistent hearing loss needs examination.

\synonyms
Eustachian Tube Dysfunction (ETD); Auditory Tube Obstruction

\medterm{Eustachian Valve} A crescentic embryological flap of tissue at the junction of the inferior vena cava and the inferior right atrium that in fetal life directs oxygenated umbilical venous blood across the foramen ovale into the left atrium; usually persists as a benign vestigial remnant in adults.

\synonyms
Valve of the Inferior Vena Cava

\medterm{Euthanasia} The deliberate ending of a person’s life by another person to relieve suffering. It is ethically and legally distinct from allowing a natural death, stopping unwanted treatment, or providing comfort-focused palliative care.

\medterm{Euthanasia And Assisted Dying} Terms describing legally regulated practices intended to end life or help a person end life to relieve suffering. Laws and safeguards differ widely; they are distinct from palliative care and withdrawal of unwanted treatment.

\medterm{Euthyroid} Having normal thyroid function, usually shown by thyroid-stimulating hormone and thyroid-hormone levels within the laboratory’s reference ranges. A person can be euthyroid even when another illness temporarily changes some thyroid test results.

\medterm{Euthyroid Goiter} An enlarged thyroid gland with normal thyroid-hormone production. It may result from iodine imbalance, autoimmune disease, nodules, medicines, or inherited factors and should be assessed if persistent, enlarging, or causing swallowing or breathing problems.

\medterm{Euthyroid Sick Syndrome} Temporary changes in thyroid blood tests during severe illness, starvation, or major surgery without a primary thyroid disorder. Treatment usually focuses on the underlying illness rather than thyroid hormone.

\medterm{Eutocia} Normal labor and birth progressing without major obstruction or complication. The term describes the course of delivery, not a separate disease or treatment.

\medterm{Euvolemic Hyponatremia} Low blood sodium with no obvious clinical signs of excess or deficient body fluid volume. Common mechanisms include inappropriate antidiuretic hormone activity, adrenal insufficiency, hypothyroidism, and excess water intake.

\medterm{Evacuation} Removal of people from danger or removal of material from a body cavity, organ, or wound. In clinical care it may refer to clearing a hematoma, bowel, uterus, or other collection; the context determines the meaning.

\medterm{Evacuation Of Uterus} Removal of pregnancy tissue or other retained material from the uterus using medicines, suction, or a surgical procedure. It may treat miscarriage, abortion, retained placenta, or bleeding and carries risks of infection, bleeding, or uterine injury.

\medterm{Evans Syndrome} An autoimmune disorder characterized by the simultaneous or sequential development of autoimmune hemolytic anemia (positive direct antiglobulin test) and immune thrombocytopenic purpura, occasionally accompanied by autoimmune neutropenia.
\synonyms
Evans's Syndrome; AIHA and ITP Syndrome

\medterm{Event-Free Survival} The length of time after treatment begins or a study starts before a defined event occurs, such as disease progression, relapse, a complication, or death. The exact event must be stated for comparisons to be meaningful.

\medterm{Event Monitor} A portable device that records the heart’s electrical rhythm when symptoms occur or at scheduled intervals over days or weeks. It can help detect intermittent abnormal rhythms that a brief ECG may miss.

\medterm{Everolimus} A medicine that suppresses a cell-growth pathway and immune activity. It is used for selected cancers, transplant care, and some genetic disorders; mouth sores, infection, high blood sugar, lung inflammation, and low blood counts may occur.

\medterm{Eversion} Turning the sole of the foot outward, away from the body’s midline. The movement uses several ankle and foot joints and can be limited by injury, arthritis, or nerve and muscle disorders.

\medterm{Evidence-Based Practice} Care that combines the best available research, professional clinical judgment, and the patient’s goals, values, and circumstances. It supports informed decisions while recognizing that evidence and individual needs may differ.

\medterm{Evisceration} Protrusion of internal organs through a surgical abdominal wound that has opened. This is a surgical emergency; the wound should be covered with sterile moist material, and urgent assessment is needed for repair and infection prevention.

\medterm{Evoked Potential} An electrical response recorded from the nervous system after a visual, sound, or touch stimulus. The test assesses whether signals travel normally along sensory pathways and may reveal nerve damage or slowed transmission.

\medterm{Evoked Potentials} Tests that record nervous-system responses to visual, sound, or touch stimuli. They can assess pathways for sight, hearing, or sensation and may help identify slowed signaling from disorders affecting nerves, the brain, or spinal cord.

\synonyms
VEP; BAER; SSEP; Evoked Responses

\medterm{Evolution} Gradual change in inherited traits of a population across generations. In clinical writing, the term can also describe how a disease, symptom, or treatment response changes over time.

\medterm{Evolutionarily Conserved Gene} A gene whose sequence or function remains similar across different species. This similarity suggests that the gene performs an important biological role, although it does not by itself prove a disease association.

\medterm{Ewart Sign} A physical-examination finding of a large fluid collection around the heart: dullness, reduced breath sounds, and bronchial breathing over the lower left back. Modern imaging is more reliable for confirming the cause.

\medterm{Ewing Sarcoma} A rare, aggressive cancer of bone or soft tissue, most often affecting children, teenagers, and young adults. It may cause persistent pain, swelling, fever, or a fracture and requires biopsy and specialist treatment.

\medterm{Ewing's Sarcoma} A malignant tumor of bone or soft tissue, most common in children and young adults. It may cause localized pain, swelling, or fever and is diagnosed with imaging, biopsy, and molecular testing; treatment usually combines chemotherapy with surgery and/or radiation.

\medterm{Ewing's Tumour} A rare aggressive cancer that usually begins in bone or nearby soft tissue in children or young adults. Persistent pain, swelling, fever, or an unexplained fracture may occur; diagnosis requires imaging and biopsy.

\medterm{Exacerbate} To make a disease, symptom, or medical problem worse. Infection, stress, exertion, allergens, medicines, or delayed treatment may exacerbate symptoms, depending on the condition.

\medterm{Exacerbation} A worsening or flare of a disease or its symptoms after a period of relative stability. Triggers may include infection, allergens, treatment interruption, exertion, or another physical or emotional stress.

\medterm{Exacerbation Of COPD} A sudden worsening of breathlessness, cough, or sputum beyond usual daily variation in chronic obstructive pulmonary disease. Infections and air pollution are common triggers; severe breathlessness, confusion, or blue lips requires urgent care.

\medterm{Exaggerated Placental Site} A noncancerous increase in placental cells at the place where the placenta attaches to the uterus. It may resemble a pregnancy-related cancer in a tissue sample, so an experienced pathologist must distinguish it from cancer.

\medterm{Examination} A systematic assessment of health using observation, touch, listening, measurements, and selected tests. It may be general or focused on a complaint and helps identify normal findings, illness, injury, and urgency.

\medterm{Examination Table} A padded, adjustable surface that supports a person during examination, testing, minor procedures, or treatment. It should be stable, cleanable, accessible, and used with appropriate safety measures.

\medterm{Exanthem} A widespread skin rash, often caused by an infection or medicine. Its appearance, timing, fever, mucosal involvement, and other symptoms help clinicians determine whether it may be contagious or serious.

\medterm{Exanthema} A widespread skin eruption associated with an illness affecting the body. Viral infections are common causes, but medicines, bacteria, and immune disorders may also cause it; blisters, skin pain, mouth sores, or severe illness need prompt assessment.

\medterm{Exanthema Subitum} A common viral illness in infants and toddlers, usually caused by human herpesvirus 6 or 7. It often causes three to five days of high fever followed by a pink rash as the fever settles.

\synonyms
Roseola Infantum; Sixth Disease

\medterm{Exanthematous Pustulosis} A sudden medicine-related eruption of many small, noninfectious pustules on red skin, often with fever. It usually improves after the responsible medicine is stopped under medical guidance, but severe reactions require urgent assessment.

\medterm{Exanthem Subitum} A common viral illness in infants and toddlers, usually caused by human herpesvirus 6 or 7. It often causes three to five days of high fever followed by a pink rash as the fever settles.
\synonyms
Roseola Infantum; Sixth Disease

\medterm{Excavation} A hollow, depression, or removal of tissue or material. In medicine it may describe a normal body feature, a disease-related defect, or the surgical creation or clearing of a space.

\medterm{Exceptional Release} Permission to use a donated organ that does not meet usual acceptance criteria when clinicians judge it suitable for a particular recipient. The decision requires careful risk assessment, documentation, and informed consent.

\medterm{Excessive Bruising} Bruises that occur often, are unusually large, appear without clear injury, or last unusually long. Causes include medicines, fragile skin, low platelets, clotting disorders, liver disease, or other illness.

\medterm{Excessive Daytime Sleepiness} An excessive tendency to feel drowsy or fall asleep during waking hours, despite an opportunity for adequate nighttime sleep. It may occur with insufficient sleep, sleep disorders, medicines, medical or neurologic conditions, or mood disorders.

\medterm{Excessive Menstrual Bleeding} Alternate name; see \textit{Heavy Menstrual Bleeding}.

\medterm{Excessive Scarring} Alternate name; see \textit{Keloid}.

\medterm{Excessive Sweating} Sweating more than needed to cool the body, either all over or in specific areas such as the hands, feet, or armpits. It may result from hyperhidrosis, medicines, infection, hormones, or another illness.

\medterm{Excessive Tearing} Tears overflowing onto the face because the eyes make too many tears or drainage is blocked. Irritation, allergy, infection, dry eye, eyelid problems, or a blocked tear duct may be responsible.

\medterm{Excessive Thirst} Polydipsia, an abnormally strong or persistent desire to drink, commonly associated with dehydration, diabetes mellitus, diabetes insipidus, medicines, or psychiatric conditions.

\medterm{Excessive Yawning} Excessive yawning is unusually frequent or difficult-to-control yawning. It may accompany sleep deprivation, excessive daytime sleepiness, medication effects, vasovagal reactions, or less commonly neurologic or cardiopulmonary disease; persistent unexplained episodes merit clinical evaluation.

\medterm{Excess Post-Exercise Oxygen Consumption} Increased oxygen use after exercise as the body returns to its usual metabolic and physiologic state.

\medterm{Exchange Transfusion} A procedure that repeatedly removes small amounts of a baby’s blood and replaces them with compatible donor blood. It can rapidly reduce dangerous bilirubin or antibody-related anemia when other treatment is insufficient.

\medterm{Excise} To remove tissue, a growth, or a lesion by cutting it out. The removed material may be examined in a laboratory to identify the cause and determine whether abnormal cells reach the edges.

\medterm{Excision} Surgical removal of part or all of a lesion, tissue, or organ. The specimen may be examined to establish a diagnosis, and removal may also treat disease when the affected area can be safely removed.

\medterm{Excisional Biopsy} Removal of an entire lump or abnormal area, sometimes with a small margin of nearby tissue, so a laboratory can examine it. It may provide a diagnosis and, for selected small lesions, complete treatment.

\medterm{Excision Biopsy} Alternate name; see \textit{Excisional Biopsy}.

\medterm{Excitation} Activation or increased responsiveness of a cell, nerve, muscle, molecule, or body system after a stimulus. In nerves and muscles, it can start electrical signals that produce movement or other functions.

\medterm{Excitation-Contraction Coupling} The chain of events that links an electrical signal in a muscle cell to contraction. The signal causes calcium release inside the cell, allowing muscle proteins to interact and generate force.

\medterm{Excitatory Neurotransmitter} A chemical messenger that makes a receiving nerve cell more likely to generate an electrical signal. Examples include glutamate and, in some settings, acetylcholine; excessive activity can injure nerve cells.

\medterm{Excitatory Synapse} A junction between nerve cells where chemical signaling makes the receiving cell more likely to produce an electrical impulse. Balanced excitatory and inhibitory signals are needed for normal brain and nerve function.

\medterm{Excitotoxicity} The pathological process of neuronal injury and cell death resulting from excessive or prolonged synaptic stimulation by excitatory neurotransmitters, primarily glutamate, triggering massive intracellular calcium influx, mitochondrial dysfunction, reactive oxygen species generation, and apoptotic protease activation.

\medterm{Excitotoxin} A substance that overstimulates nerve cells and can damage or kill them. Excessive signaling may allow too much calcium into cells, contributing to injury in some brain diseases or poisoning states.

\medterm{Exclamation Point Hair} A short broken hair that is thinner at its outer end and wider near the scalp. It is a characteristic finding of alopecia areata but can occasionally appear in other hair disorders.

\medterm{Exclusion Diet} A planned diet that removes a suspected food and later reintroduces it to see whether symptoms recur. It should be supervised when possible and must not replace emergency allergy evaluation when severe reactions or anaphylaxis are possible.

\medterm{Excoriate} To remove or injure the skin surface by scratching, rubbing, or picking. Repeated excoriation can cause open sores, infection, scarring, and difficulty healing.

\medterm{Excoriation} A superficial skin injury caused by scratching, rubbing, or picking. It may cause pain, bleeding, crusting, infection, or scarring and usually heals when the irritation or behavior stops.

\medterm{Excoriation Disorder} A mental-health disorder involving recurrent skin picking that causes skin lesions and clinically significant distress or impairment. Repeated unsuccessful efforts to reduce or stop the behavior are common.

\synonyms
Skin-Picking Disorder

\medterm{Excreta} Material expelled from the body, including stool, urine, sweat, vomit, and respiratory secretions. Some excreta carry infectious organisms, so handwashing, protective equipment, cleaning, and safe disposal help prevent spread.

\medterm{Excretion} Removal of waste, medicines, or other substances from the body, mainly through urine, bile, breath, stool, sweat, or breast milk. Kidney and liver function strongly affect how long many substances remain in the body.

\medterm{Executive Function} Mental skills used to plan, focus, remember instructions, control impulses, solve problems, make decisions, and switch between tasks. Problems may follow brain injury, dementia, developmental conditions, or mental illness.

\medterm{Executive Functions} Mental abilities that support planning, organization, attention, memory, decision-making, self-control, and flexible behavior. Difficulties can affect work, school, relationships, and self-care and may have neurologic, developmental, or psychiatric causes.

\synonyms
Executive Control

\medterm{Exemestane} A medicine that lowers estrogen production by blocking an enzyme used after menopause. It treats selected hormone-sensitive breast cancers and may cause hot flushes, joint pain, bone loss, fatigue, or abnormal cholesterol.

\medterm{Exenatide} An injectable medicine that imitates a gut hormone, increasing insulin release when glucose is high, reducing glucagon, and slowing stomach emptying. It treats type 2 diabetes but may cause nausea, dehydration, or pancreatitis.

\medterm{Exencephaly} A severe birth defect in which much of the developing brain lies outside the skull because the skull and protective coverings did not form normally. It usually progresses to anencephaly and is not compatible with long-term survival.

\medterm{Exenteration} Extensive surgery removing organs and tissues from a body compartment, such as the pelvis or eye socket. It is reserved for selected advanced cancers or destructive disease and can cause major functional and emotional effects.

\medterm{Exercise} Planned physical activity intended to improve or maintain strength, endurance, flexibility, balance, function, or health. Aerobic, resistance, and balance activities should be chosen and increased according to a person’s abilities and medical risks.

\synonyms
Physical Exercise

\medterm{Exercise And Activity For Weight Loss} Physical activity used with eating changes to increase energy use, preserve muscle, and support gradual weight loss. The safest plan depends on fitness, medical conditions, joint health, and personal ability.

\medterm{Exercise And Age} The way physical activity needs and responses may change with age; safe exercise should match fitness, health, and ability.

\medterm{Exercise And Immunity} The relationship between physical activity and immune function; regular moderate activity generally supports health, while excessive exertion may increase short-term illness risk.

\medterm{Exercise Headaches} Headaches occurring during or after strenuous physical activity. Primary exercise headaches are recurrent and usually not caused by an underlying disorder, but a first, abrupt, severe, or atypical exertional headache requires prompt medical assessment to exclude secondary causes such as subarachnoid hemorrhage, arterial dissection, or reversible cerebral vasoconstriction syndrome.

\medterm{Exercise-Induced Arrhythmias} Abnormal heart rhythms that begin or worsen during physical exertion. They may result from inherited electrical disorders, heart muscle disease, or reduced blood flow; fainting, chest pain, or palpitations during exercise needs urgent evaluation.

\medterm{Exercise-Induced Asthma} Airway narrowing during or soon after exercise, causing cough, wheezing, chest tightness, or breathlessness. It can occur with or without persistent asthma and is managed with an individualized prevention and treatment plan.

\medterm{Exercise-Induced Bronchoconstriction} Temporary narrowing of the airways during or after exercise, causing cough, wheezing, chest tightness, or shortness of breath. Cold, dry air can worsen it; diagnosis and preventive inhaler use should be guided clinically.

\medterm{Exercise Intolerance} An inability to do as much physical activity as expected because of unusual tiredness, breathlessness, chest discomfort, leg pain, weakness, or dizziness. Heart, lung, blood-vessel, muscle, nerve, blood, metabolic, and conditioning problems can cause it.

\medterm{Exercises To Help Prevent Falls} Activities that improve leg strength, balance, flexibility, walking, and confidence. A program should match ability and medical conditions; supervised practice may be safest for people with previous falls or severe weakness.

\medterm{Exercise Stress Test} A test that records heart rate, rhythm, blood pressure, and symptoms while a person walks on a treadmill or pedals a stationary bicycle. It assesses exercise capacity and may suggest reduced heart blood flow or an abnormal rhythm.

\medterm{Exercise Tolerance} The ability to perform physical activity before symptoms such as breathlessness, fatigue, chest discomfort, or pain limit activity.

\medterm{Exfoliation Syndrome} An age-related eye disorder in which abnormal material collects on the lens and other front-eye structures. It increases the risk of glaucoma and can make cataract surgery more complicated, so regular eye checks are important.

\medterm{Exfoliative Cytology} Examination of cells that naturally shed from a body surface or are gently collected. It may help investigate infection, inflammation, or cancer, but abnormal or uncertain results may require tissue biopsy.

\medterm{Exfoliative Dermatitis} Severe redness, scaling, and shedding affecting most of the skin. Medicines, psoriasis, eczema, lymphoma, and other diseases may cause it; fluid loss, infection, temperature problems, and heart strain make urgent medical assessment necessary.

\medterm{Exhalation} Breathing air out of the lungs. It is usually passive at rest as the lungs recoil, but muscles actively assist exhalation during exercise, coughing, or breathing difficulty.

\medterm{Exhaled Nitric Oxide} A breath test that estimates one type of airway inflammation. A raised result can support asthma assessment and suggest likely benefit from inhaled steroids, but smoking, infection, allergies, and current treatment can affect the result.

\medterm{Exhaustion} Extreme physical or mental tiredness that reduces the ability to function. It may follow exertion, poor sleep, illness, dehydration, anemia, hormone disorders, medicines, or stress; severe or persistent symptoms need evaluation.

\medterm{Exhaustion, Heat} Alternate name; see \textit{Heat Exhaustion}.

\medterm{Exhibitionism} Sexual arousal from exposing one’s genitals to an unsuspecting person. Repeated nonconsensual exposure can harm others and becomes a diagnosable disorder when urges or actions cause significant distress, impairment, or risk.

\medterm{Exhibitionistic Disorder} A mental-health disorder involving repeated sexual arousal from exposing the genitals to an unsuspecting person, with action on the urges or clinically significant distress or impairment.

\medterm{Exhumation} Lawful removal of a buried body or remains for identification, further testing, investigation, or repatriation. It requires legal authorization and respectful handling; testing may include toxicology, DNA analysis, or examination of injuries.

\medterm{Existential Distress} Emotional suffering related to meaning, identity, purpose, freedom, mortality, or loss of control. It can occur during serious illness and may improve with empathetic communication, counseling, spiritual care, and symptom relief.

\medterm{Exit Wound} The wound where a projectile leaves the body. Its size and shape vary with the projectile, tissue, clothing, and bone, so appearance alone cannot prove the direction or identify every injury; complete examination is required.

\medterm{Exocervix} The part of the cervix that projects into the vagina and is covered mainly by flat squamous cells. It is examined during pelvic examination and cervical screening.

\medterm{Exocrine} Relating to glands that release substances through ducts onto a body surface or into a cavity. Examples include sweat, salivary, tear, and digestive glands.

\medterm{Exocrine Gland} A gland that sends its secretion through a duct to a body surface or internal cavity. Sweat, salivary, tear, and digestive glands are examples; their secretions help protect, lubricate, or digest.

\medterm{Exocrine Glands} Glands that release substances through ducts onto a surface or into a body cavity. They include sweat, salivary, tear, and digestive glands and support protection, lubrication, digestion, and temperature control.

\medterm{Exocrine Pancreas} The part of the pancreas that releases digestive enzymes and bicarbonate into the small intestine. These secretions help break down fats, proteins, and carbohydrates and neutralize stomach acid.

\medterm{Exocrine Pancreatic Insufficiency} Failure of the pancreas to release enough digestive enzymes, causing poor digestion and absorption. It may lead to greasy stools, bloating, diarrhea, weight loss, and vitamin deficiency.

\medterm{Exocrine Pancreatic Insufficiency EPI} Alternate index label; see \textit{Exocrine Pancreatic Insufficiency}.

\synonyms
EPI; Pancreatic Exocrine Insufficiency

\medterm{Exocytosis} The process by which a cell releases substances outside itself. Small sacs inside the cell join its outer membrane and release hormones, nerve signals, digestive enzymes, or other proteins.

\medterm{Exogenous} Coming from outside the body rather than being produced within it. An exogenous substance may be a medicine, hormone, nutrient, toxin, infection, or other externally introduced material.

\medterm{Exogenous Cushing Syndrome} A condition caused by prolonged or excessive corticosteroid treatment, producing effects similar to too much natural cortisol. Features may include weight gain, easy bruising, high blood pressure, diabetes, weakness, and bone loss.

\medterm{Exomphalos} A birth defect in which abdominal organs protrude through the belly button into the base of the umbilical cord and are covered by a membrane. Newborns need protection of the sac and specialist surgical care.

\medterm{Exon} A segment of a gene that remains in the mature messenger RNA after unused sections are removed. It may contain instructions for making part of a protein or, in some genes, help regulate gene activity.

\medterm{Exopeptidase} An enzyme that removes amino acids from the end of a protein. Digestive exopeptidases help break dietary proteins into smaller units that the intestine can absorb.

\medterm{Exophiala} A group of darkly pigmented environmental fungi that can rarely cause skin, lung, eye, or widespread infection. Disease is more likely in people with weakened immunity or implanted medical devices.

\medterm{Exophthalmic Goitre} An enlarged thyroid associated with Graves disease, in which immune stimulation causes excess thyroid hormone. Graves disease may also cause eye dryness, pain, swelling, or bulging.

\medterm{Exophthalmos} The abnormal forward protrusion or displacement of one or both eyeballs from the orbit, most frequently caused by Graves disease (thyroid eye disease), orbital tumors, or inflammation.

\synonyms
Bulging Eyes; Proptosis; Ocular Proptosis

\medterm{Exosomes} Very small membrane-covered particles released by cells that carry proteins, fats, and genetic material to other cells. They help cells communicate and are being studied for diagnosis and drug delivery, but many medical uses remain experimental.

\medterm{Exostosis} A benign growth of bone. It may be harmless, but location can cause pain, pressure, restricted movement, or blockage; repeated cold-water exposure can produce bony growths in the ear canal.

\synonyms
Surfer's Ear; Osteoma of External Auditory Canal

\medterm{Exotoxins} Poisonous proteins released by some bacteria that damage particular cells or disrupt body functions. They cause diseases such as botulism, tetanus, diphtheria, or cholera and may require urgent treatment.

\medterm{Exotropia} A form of misaligned eyes in which one eye turns outward. It may be occasional or constant and can cause double vision, reduced depth perception, or loss of vision development in children; eye assessment guides treatment.

\medterm{Expectant Treatment} Careful observation with regular review instead of immediate medicine or surgery when immediate intervention is not necessary. It includes clear follow-up and warning signs because treatment may be needed if symptoms worsen or tests change.

\medterm{Expectorants} Medicines intended to loosen or increase airway mucus so it can be coughed up more easily. Benefits vary, and fluids, humidification, and treatment of the underlying illness may be more important than an expectorant.

\medterm{Expectoration} Coughing up and bringing mucus or phlegm out of the respiratory tract. Changes in amount, color, smell, blood, fever, or breathlessness may indicate infection or another lung problem.

\medterm{Experience} A person’s direct knowledge, perception, or emotional response to an event. In healthcare, patient experience includes how care is accessed, communicated, coordinated, and received.

\synonyms
Subjective Experience; Patient-Reported Experience

\medterm{Expiration} Movement of air out of the lungs, usually caused by relaxation of the breathing muscles and natural recoil of the lungs. During forceful breathing, abdominal and chest muscles actively push air out.

\medterm{Expiratory Grunting} A grunting sound made during breathing out, caused by partially closing the voice box. It can briefly help keep small air spaces open but, especially in a baby or child, often signals respiratory distress and needs assessment.

\medterm{Expiratory Reserve Volume} The maximum additional volume of air that can be actively and forcefully exhaled from the lungs at the end of a normal, quiet tidal expiration, normally measuring approximately 1.0 to 1.2 liters in healthy adults.

\synonyms
ERV

\medterm{Explant} To remove an implanted device, prosthesis, or tissue. The term can also mean tissue removed from an organism and maintained in laboratory culture for study.

\medterm{Explantation} Surgical removal of an implanted device or prosthesis. It may be needed because of infection, malfunction, migration, breakage, tissue injury, intolerance, or completion of temporary treatment.

\medterm{Exploratory Laparotomy} Open surgery to enter the abdomen and directly find or treat a serious problem such as bleeding, perforation, blocked bowel, or loss of blood flow when less-invasive tests or treatment are insufficient.

\medterm{Exploratory Operation} Surgery performed to inspect an area directly when scans, tests, or less-invasive procedures cannot establish the cause or extent of disease. Tissue may be sampled and treatment may be performed during the operation.

\medterm{Exposure} Contact with a physical, chemical, biological, or psychological agent. Health risk depends on what the agent is, how much contact occurred, how it entered the body, and how long exposure lasted.

\medterm{Exposure Route} The way a substance enters the body, such as by breathing, swallowing, skin contact, injection, or the eyes. The route affects how quickly and how much of the substance is absorbed and how toxic it may be.

\medterm{Expression} In genetics, the process by which a gene’s information is used to make RNA or protein. In clinical examination, expression can also mean the outward display of emotion or symptoms; altered gene activity can contribute to disease.

\medterm{Expression Hopping} A shift from one form of addictive behavior to another, such as replacing substance use with compulsive gambling. The underlying loss of control and harm may persist even when the specific behavior changes.

\medterm{Expressivity} The degree to which a genetic variant’s effects appear in a person. Variable expressivity means that people with the same variant may have different features or different severity of the associated condition.

\medterm{Expulsion, Stage Of} The second stage of labor, beginning when the cervix is fully open and ending when the baby is born. It includes pushing and may be shortened or assisted when the mother or baby is at risk.

\medterm{Exsanguination} Life-threatening loss of blood from severe internal or external bleeding. It can rapidly cause shock, organ failure, and death; immediate pressure or bleeding control, emergency care, and blood replacement may be needed.

\medterm{Exserohilum} A group of darkly pigmented molds that can rarely cause sinus, eye, skin, or invasive infection. Serious disease is more likely in people with weakened immunity and requires specialist testing and treatment.

\medterm{Extend} To straighten a joint or increase the angle between two bones. For example, extending the elbow moves the forearm away from the upper arm toward a straight position.

\synonyms
Extension

\medterm{Extension} Movement that straightens a joint or increases the angle between bones, such as straightening the knee or elbow. Excessive or forced extension can strain or tear muscles, tendons, ligaments, cartilage, or the joint.

\medterm{Extensively Drug-Resistant Tuberculosis} Tuberculosis caused by bacteria resistant to several key medicines, including isoniazid, rifampin, a fluoroquinolone, and at least one other important medicine. It is difficult to treat and requires specialist, susceptibility-guided therapy.

\synonyms
XDR-TB; XDR Tuberculosis

\medterm{Extensor} A muscle or other structure that straightens a joint or increases the angle between bones. Extensor muscles commonly act opposite flexor muscles and help position limbs and fingers.

\medterm{Extensor Digitorum} A forearm muscle whose tendons straighten the fingers and help extend the wrist. Overuse near its origin can contribute to outer-elbow pain commonly called tennis elbow.

\medterm{Extensor Plantar Response} An abnormal neurological reflex in which firm tactile stroking of the lateral plantar border of the foot produces slow extension (dorsiflexion) of the great toe accompanied by fanning of the other toes, indicating corticospinal (pyramidal) tract impairment.
\synonyms
Positive Babinski Sign; Babinski Reflex

\medterm{Extent} The size, severity, range, or spread of a disease, injury, symptom, or treatment effect. Describing extent helps determine prognosis, further testing, and treatment.

\medterm{External Auditory Canal} The passage from the outer ear to the eardrum. It is lined with skin, hair, and wax-producing glands that help protect the ear, but it can become blocked, irritated, or infected.

\medterm{External Beam Partial Breast Irradiation} An accelerated conformal radiation therapy technique delivered via external beam linac that targets exclusively the lumpectomy cavity margin following breast-conserving surgery, sparing the remainder of the ipsilateral breast tissue in selected early-stage breast cancer patients.

\synonyms
EB-APBI; Accelerated Partial Breast Irradiation; Targeted Breast Radiotherapy

\medterm{External Beam Radiation Therapy} Cancer treatment delivered from a machine outside the body to a carefully planned area. It can damage cancer cells while limiting nearby tissue exposure; fatigue, skin changes, and organ-specific effects may occur.

\medterm{External Capsule} A thin bundle of nerve fibers in the brain near the basal ganglia. It carries signals between the cerebral cortex and deeper brain regions involved in movement, attention, and other functions.

\medterm{External Carotid Artery} A major branch of the neck artery that supplies blood to much of the face, scalp, neck, and parts of the brain coverings through several branches.

\medterm{External Cephalic Version} A procedure in which a clinician applies gentle pressure through the maternal abdomen to turn a fetus from a breech or transverse presentation into a head-down (cephalic) position for vaginal delivery. It is typically performed near term with continuous ultrasound and fetal heart rate monitoring.
\synonyms ECV; Cephalic version.

\medterm{External Ear} The visible part of the ear and the ear canal leading to the eardrum. It collects sound, directs it inward, and helps protect the middle and inner ear.

\medterm{External Fixation} Stabilization of a broken bone or injured limb with pins or wires connected to a frame outside the body. It may be temporary or definitive, especially when swelling, wounds, or contamination prevent internal fixation.

\medterm{External Iliac Artery} A major artery carrying blood from the pelvis toward the leg. It becomes the femoral artery below the groin and gives branches to the lower abdominal wall and pelvic region.

\medterm{External Iliac Vein} A continuation of the femoral vein above the inguinal ligament that joins the internal iliac vein to form the common iliac vein.

\medterm{External Incontinence Devices} Products worn or placed outside the body to collect urine or stool or protect clothing from leakage. Examples include urinary sheaths, absorbent products, and fecal collection systems; skin care and correct fitting help prevent injury.

\medterm{External Jugular Vein} A visible vein running along the side of the neck that drains blood from the scalp and face toward the chest. Its fullness can provide clues about pressure in the right side of the heart.

\medterm{External Nose} The visible nose formed by bone, cartilage, skin, and soft tissue. It shapes the nasal passages, helps warm and filter inhaled air, supports smell, and contributes to facial structure.

\medterm{External Oblique Muscle} The broad, superficial muscle on each side of the abdomen. It compresses abdominal contents and helps bend and rotate the trunk; its lower border contributes to the groin ligament and canal.

\medterm{External Otitis} Inflammation or infection of the ear canal outside the eardrum, often called swimmer’s ear. Moisture, scratching, or skin disease may contribute, causing pain, itching, discharge, swelling, or reduced hearing.

\synonyms
Otitis Externa; Swimmer's Ear

\medterm{External Resorption} Loss of tooth root or other tooth structure caused by a process starting outside the tooth. Injury, pressure, inflammation, or replacement by bone may contribute; treatment depends on location and severity.

\medterm{External Ventricular Drain} A critical neurosurgical bedside procedure in which a flexible plastic catheter is surgically placed through a small burr hole in the skull directly into one of the fluid-filled cavities (lateral ventricles) of the brain. The catheter connects to an external closed drainage and pressure-transducer system. An external ventricular drain serves two vital functions: continuously monitoring intracranial pressure (the pressure inside the skull) and therapeutically draining excess cerebrospinal fluid or blood. It is essential in managing acute hydrocephalus, severe traumatic brain injury, subarachnoid hemorrhage, intraventricular hemorrhage, or intracranial hypertension, preventing brain herniation and death.

\synonyms
EVD; Ventriculostomy; External Ventricular Drainage; Ventricular Catheter

\medterm{Extinction} In learning, the fading of a learned response when reinforcement stops. In neurology, it can mean failing to notice one side of the body or space when both sides are stimulated, suggesting an attention disorder.

\medterm{Extirpation} Complete surgical removal of a tissue, organ, growth, or foreign body. The term describes the extent of removal rather than a particular disease or surgical technique.

\medterm{Extraamniotic} Located outside the amniotic sac surrounding a fetus. The term may describe a procedure, device, or fluid placed between the amnion and nearby tissues during pregnancy care.

\medterm{Extracapsular Cataract Surgery} Cataract surgery that removes the cloudy lens while leaving most of its surrounding capsule in place. An artificial lens is usually inserted into the remaining capsule to restore focusing ability.

\medterm{Extracellular} Located outside cells. Extracellular fluid, proteins, ions, and signals surround cells and help provide nutrients, remove waste, support tissues, and coordinate cell activity.

\medterm{Extracellular Exosome} A tiny membrane-covered vesicle released by a cell that can carry proteins, fats, and genetic material to other cells. Exosomes help cells communicate and are being studied as disease markers and treatment vehicles.

\medterm{Extracellular Fluid} Fluid outside cells, including blood plasma and the fluid between cells. It carries nutrients, hormones, and waste and is regulated mainly by the kidneys, circulation, and hormones.

\medterm{Extracellular Matrix} A network of proteins and other substances outside cells that supports tissues and helps cells attach, communicate, grow, and repair injury. Its composition differs among bone, skin, muscle, and other organs.

\medterm{Extracellular Pathogens} Microorganisms that live and multiply outside human cells. Examples include Streptococcus, Staphylococcus, and some Escherichia coli. They can damage tissues directly or through toxins, and the immune system usually fights them with antibodies and other defenses.

\medterm{Extracellular Region} The area outside cells. It contains fluid, minerals, proteins, and signals that support tissues and influence how nearby cells grow, move, and function.

\medterm{Extracellular Space} The fluid- and matrix-filled area outside cells. Nutrients, waste products, hormones, and signaling molecules move through it, and its composition helps control cell survival, movement, and function.

\medterm{Extracellular Vesicle} A membrane-bound particle released by a cell that transfers molecular cargo between cells and is studied in physiology, disease, biomarkers, and drug delivery.

\medterm{Extracorporeal} Occurring or performed outside the body; extracorporeal circulation and hemodialysis route blood through an external machine.

\medterm{Extracorporeal Circulation} Movement of a person's blood through equipment outside the body and back again. Heart-lung machines, dialysis, and some life-support systems use it to add oxygen, remove waste, or support circulation; bleeding, clotting, infection, and organ injury are possible.

\medterm{Extracorporeal Membrane Oxygenation} An advanced, temporary mechanical life-support technique that drains deoxygenated blood from the body, circulates it through an external artificial membrane lung to add oxygen and remove carbon dioxide, and pumps it back into the bloodstream. It is indicated for severe, life-threatening respiratory failure or cardiogenic shock refractory to standard mechanical ventilation and intensive care. It is utilized in two main configurations: veno-venous (VV-ECMO) for isolated respiratory support, and veno-arterial (VA-ECMO) for combined hemodynamic and respiratory support. Major risks include hemorrhage from systemic anticoagulation, thromboembolism, hemolysis, cannula-site infection, and limb ischemia.

\synonyms
ECMO; Extracorporeal Life Support; ECLS; VV-ECMO; VA-ECMO

\medterm{Extracorporeal Shock Wave Lithotripsy} A treatment that uses focused shock waves from outside the body to break kidney or ureter stones into smaller pieces that can pass in urine. Bruising, bleeding, pain, blockage, or incomplete stone removal can occur.

\medterm{Extracranial} Located or occurring outside the skull or cranial cavity. For example, extracranial arteries supply the scalp and face, while extracranial injury affects tissues outside the skull and brain.

\medterm{Extracranial Hematoma} A collection of blood outside the skull or cranial cavity, usually caused by bleeding into soft tissues after injury.

\medterm{Extract} A concentrated preparation made by separating active substances from a plant, herb, food, or other source. Strength, purity, interactions, and safety vary widely, especially in unregulated supplements.
\synonyms
Botanical Extract; Pharmacological Extract

\medterm{Extraction} Removal of a tooth, tissue, foreign object, or other material from the body using a procedure. The method and risks depend on what is removed, where it is located, and the person’s health.

\medterm{Extra-Dural} Located outside the dura mater, the tough protective membrane around the brain and spinal cord. The term may describe a space, collection of blood, injection, or other structure outside this membrane.

\medterm{Extra-Gastrointestinal Stromal Tumor} A tumor with features of a gastrointestinal stromal tumor that begins outside the digestive tract, often in abdominal supporting tissues. Diagnosis uses imaging, tissue examination, and molecular tests to guide cancer treatment.

\medterm{Extrahepatic Biliary Atresia} Alternate name; see \textit{Biliary Atresia}.

\medterm{Extrahepatic Cholestasis} Reduced or blocked bile flow in the ducts outside the liver. Gallstones, scarring, inflammation, or a tumor may cause yellow skin, dark urine, pale stools, itching, pain, or infection and may need urgent treatment.

\medterm{Extramammary Paget Disease} A rare skin cancer usually affecting the vulva, anus, scrotum, or nearby skin. It often appears as a persistent itchy, sore, red, or scaly patch and requires biopsy to confirm diagnosis and assess for related cancer.

\medterm{Extramammary Paget's Disease} A rare skin cancer involving gland-bearing skin outside the breast, often around the vulva, anus, or penis. It may appear as a persistent itchy, red, scaly, or oozing patch and requires biopsy.

\medterm{Extramedullary Hematopoiesis} The making of blood cells outside the bone marrow, usually in the liver or spleen. It can happen when the bone marrow is damaged, crowded, or unable to make enough blood cells.
\synonyms
EMH; Extramedullary Blood Formation

\medterm{Extranodal} Located outside the lymph nodes. An extranodal disease or lymphoma affects an organ or other tissue, such as the stomach, skin, lung, bone, or brain.

\medterm{Extranodal Marginal Zone Lymphoma} Alternate name; see \textit{MALT Lymphoma}.

\medterm{Extranodal Marginal Zone Lymphoma Of Mucosa-Associated Lymphoid Tissue} A usually slow-growing cancer of B lymphocytes arising in tissues such as the stomach, salivary glands, thyroid, lungs, or eye. Symptoms depend on the organ involved and treatment may include observation, infection treatment, or cancer therapy.

\synonyms
Extranodal Marginal Zone Lymphoma
MALT Lymphoma
Mucosa-Associated Lymphoid Tissue Lymphoma

\medterm{Extranodal Natural Killer T-Cell Lymphoma} An aggressive lymphoma arising from immune cells, often affecting the nose and upper throat. It is strongly associated with Epstein-Barr virus and may cause blockage, bleeding, swelling, fever, or tissue destruction.

\medterm{Extraocular} Outside the eyeball or involving the muscles, nerves, and tissues around it. Extraocular structures move the eye, protect it, support vision, and connect it with the brain and face.

\medterm{Extraocular Muscle} One of the muscles outside the eyeball that moves the eye in different directions. Weakness or nerve injury can cause double vision, eye misalignment, or difficulty looking in a particular direction.

\medterm{Extraocular Muscle Function Testing} Examination of eye movements to assess the muscles and nerves that move the eyes. The person follows targets in different directions; abnormal movement can indicate muscle, nerve, brain, or eye-socket disease.

\medterm{Extraocular Muscles} Six muscles outside each eyeball that move the eyes up, down, inward, outward, and in rotation. Weakness or nerve injury can cause double vision, eye misalignment, or difficulty looking in a particular direction.

\medterm{Extrapulmonary Manifestation Of Tuberculosis} Tuberculosis affecting tissues outside the lungs, such as lymph nodes, bones, kidneys, the brain coverings, abdomen, or reproductive organs. Symptoms vary by site and treatment requires multiple antibiotics.

\medterm{Extrapyramidal Disorder} A disorder of brain circuits that regulate movement, posture, and muscle tone. It may cause stiffness, slowness, tremor, abnormal movements, or restlessness and can result from disease or medicines.

\medterm{Extrapyramidal Symptoms} Movement problems caused by medicines or diseases affecting movement circuits. They include stiffness, tremor, slowed movement, restlessness, sudden spasms, or involuntary movements, commonly after dopamine-blocking medicines.

\medterm{Extrapyramidal System} Interconnected brain structures and pathways, including the basal ganglia, that help control automatic movement, posture, muscle tone, and coordination. Damage or medicine effects can produce abnormal movements.

\medterm{Extrapyramidal Tracts} Nerve pathways outside the main voluntary movement tract that influence posture, muscle tone, automatic movements, and coordination. Disturbance can cause stiffness, tremor, slowness, or involuntary movements.

\medterm{Extraskeletal Chondroma} A benign cartilage-forming lump arising in soft tissue rather than from bone, often in the hands or feet. It may calcify, cause pain or pressure, and sometimes needs removal to confirm the diagnosis.

\medterm{Extraskeletal Mesenchymal Chondrosarcoma} A rare aggressive cancer of soft tissue that contains malignant cartilage. It may arise outside bone and requires specialist biopsy, staging, and combined cancer treatment.

\medterm{Extraskeletal Osteosarcoma} A rare aggressive cancer that produces bone-like material in soft tissue rather than beginning inside bone. It usually affects adults and requires imaging, biopsy, staging, and specialist treatment.

\medterm{Extrasystole} Alternate name; see \textit{Premature Ventricular Contractions}.

\medterm{Extrauterine Pregnancy} A pregnancy implanted outside the normal uterine cavity, most often in a fallopian tube. It can rupture and cause life-threatening internal bleeding; one-sided pain, fainting, or bleeding during pregnancy requires urgent care.

\medterm{Extravasation} Leakage of blood, medicine, urine, or another body fluid from a vessel or duct into nearby tissues. It can cause swelling, pain, inflammation, infection, or tissue death; treatment depends on the substance, location, and severity.

\medterm{Extravascular Hemolysis} The early removal and breakdown of red blood cells mainly in the spleen and liver. It can cause anemia, yellowing of the skin or eyes, and an enlarged spleen.
\synonyms
Reticuloendothelial Hemolysis; Splenic Erythrocyte Clearance

\medterm{Extraventricular Drain} A temporary tube placed through the skull into a fluid-filled space in the brain. It can measure pressure and drain extra brain fluid or blood, especially when fluid buildup is dangerous.

\synonyms
EVD; Ventriculostomy Catheter

\medterm{Extremities} The arms and legs, including their bones, joints, muscles, nerves, blood vessels, and skin. Problems in an extremity may affect movement, sensation, circulation, or function.

\medterm{Extremity} A limb of the body. The upper extremity includes the shoulder, arm, forearm, wrist, and hand; the lower extremity includes the hip, thigh, leg, ankle, and foot.

\medterm{Extremity Angiography} Imaging of blood vessels in an arm or leg, usually after contrast is introduced. It can show narrowing, blockage, aneurysm, abnormal connections, or injury and may sometimes be combined with treatment.

\medterm{Extremity Hardware Removal} Surgical extraction of orthopedic implants such as plates, screws, intramedullary rods, or wires from an arm or leg. Typically indicated for persistent localized pain, implant infection, mechanical irritation, soft-tissue impingement, nonunion, or planned subsequent procedures.

\synonyms
Orthopedic Hardware Removal; Implant Removal; Plate and Screw Removal

\medterm{Extremity X-Ray} An X-ray image of an arm or leg used to look for fractures, joint alignment problems, arthritis, infection, tumors, foreign objects, or other bone abnormalities.

\medterm{Extrinsic} Coming from outside a structure, organ, or person. An extrinsic cause may be an external injury, medicine, infection, toxin, environmental exposure, or other outside influence.

\medterm{Extrinsic Allergic Alveolitis} An immune-mediated lung disease caused by repeated inhalation of organic dusts, molds, bird proteins, or other environmental substances. It can cause cough, breathlessness, fever, and, with ongoing exposure, permanent lung scarring; it is also called hypersensitivity pneumonitis.

\medterm{Extrinsic Factor} An influence from outside a person, tissue, or body process. Examples include medicines, infections, toxins, temperature, nutrition, stress, and environmental exposure; effects depend on the type and amount.

\medterm{Extrinsic Ureteral Obstruction} Blockage of a ureter caused by pressure from outside the ureter, such as a tumor, scar tissue, enlarged lymph nodes, blood-vessel abnormality, or nearby inflammation. It can reduce urine drainage and cause kidney swelling or loss of kidney function.

\synonyms
Extrinsic Ureteral Compression; Secondary Ureteral Obstruction

\medterm{Extubate} To remove a breathing tube after a person can breathe adequately, protect the airway, clear secretions, and maintain safe oxygen levels without the ventilator.

\medterm{Extubation} Removal of a breathing tube from the windpipe after a person can breathe independently, keep the airway open, clear secretions, and maintain adequate oxygen and carbon-dioxide levels.

\medterm{Extubation Criteria} Findings used to judge whether a person can safely have a breathing tube removed. They include adequate breathing and oxygenation, stable circulation, alertness, airway protection, effective cough, and recovery from muscle-paralyzing medicines.

\medterm{Exudate} Protein-rich fluid that leaks from inflamed or injured blood vessels into tissues or a body cavity. Its presence usually indicates inflammation, infection, or tissue injury rather than pressure alone.

\medterm{Exudative Pleural Effusion} A protein-rich collection of fluid between the lung and chest wall caused by local inflammation or injury. Pneumonia, cancer, tuberculosis, and blood clots are possible causes; testing the fluid helps guide treatment.

\medterm{Eye} The organ that detects light and forms vision. It includes the cornea, lens, iris, retina, optic nerve, and supporting tissues within the eye socket; disease can affect sight, movement, comfort, or appearance.

\medterm{Eye Allergies} An allergic reaction affecting the eyes and often the inner eyelids. It usually causes itching, redness, watering, burning, or swollen eyelids and is commonly called allergic conjunctivitis.

\synonyms
Eye Allergy
Ocular Allergy

\medterm{Eye And Orbit Ultrasound} An ultrasound scan of the eye and surrounding socket. It can assess the retina, optic nerve, eye muscles, masses, bleeding, or other structures when the view through the front of the eye is obscured.

\medterm{Eyeball} The roughly spherical organ of vision containing the cornea, sclera, iris, lens, retina, internal fluids, and other tissues. It focuses light onto the retina, which sends visual signals through the optic nerve.

\synonyms
Globe

\medterm{Eye Bank} An organization that recovers, evaluates, stores, and distributes donated corneas or other eye tissue for transplantation, research, or education. Donor screening and tissue testing help protect recipients.

\medterm{Eye, Black} Alternate name; see \textit{Black Eye}.

\medterm{Eye Bleed} Bleeding on, inside, or around the eye. A bright-red painless patch on the white of the eye is often superficial, but bleeding inside the eye, pain, injury, or vision change requires prompt assessment.

\medterm{Eyebrow} A strip of hair above the eye socket that helps divert sweat, water, and particles away from the eye and contributes to facial expression.

\medterm{Eye Cancer} A cancer that begins in the eyelid, surface of the eye, uvea, retina, or tissues around the eye. Warning signs can include a growing or changing spot, persistent redness, vision changes, eye pain, bulging, or a new pupil abnormality; diagnosis usually requires specialist examination and imaging or biopsy.

\medterm{Eye Color} The visible color of the iris, determined mainly by pigment and light scattering. A new, one-sided, uneven, or painful color change may indicate injury, inflammation, medicine effects, or another eye disorder.

\medterm{Eye Development} Formation and maturation of the eyes, optic nerves, and visual pathways before and after birth. Disruption can cause structural abnormalities, reduced vision, misalignment, or delayed development of visual skills.

\medterm{Eye Disease} Any disorder affecting the eye or visual pathways, including refractive errors, cataract, glaucoma, retinal disease, infection, inflammation, and tumors. Symptoms such as sudden vision loss, severe pain, flashes, or a curtain-like shadow require urgent assessment.

\medterm{Eye, Diseases Of} Diseases affecting the eye or visual pathways, including refractive problems, cataracts, glaucoma, retinal disease, infection, inflammation, injury, and tumors. Symptoms may involve vision, pain, redness, discharge, or eye movement.

\medterm{Eye Disorder} A disease or structural problem affecting the eye or visual pathways. It may cause changed vision, pain, redness, discharge, double vision, field loss, or other findings and may require routine or urgent assessment.

\medterm{Eye Disorders} Diseases or structural problems affecting the eye, optic nerve, or visual pathways. They may cause blurred or lost vision, pain, redness, discharge, double vision, or field loss and can require routine or emergency assessment.

\medterm{Eye Drops} Liquid medicines, lubricants, or diagnostic preparations placed on the eye surface. They may relieve dryness or treat disease; contaminated bottles, shared drops, or incorrect use can cause injury or infection.

\medterm{Eye Emergencies} Conditions needing immediate eye assessment include sudden vision loss, chemical exposure, penetrating injury, severe pain, new neurologic symptoms, or rapidly spreading infection. Do not rub an injured eye or press on it.

\medterm{Eye Floaters} Alternate name; see \textit{Floaters}.

\medterm{Eyeglass} A lens mounted in a frame to correct focusing problems such as nearsightedness, farsightedness, astigmatism, or presbyopia. The prescription is selected after measuring vision and the eyes.

\medterm{Eyeglasses and Sunglasses} Alternate name; see \textit{Eyeglasses, Sunglasses, and Magnifiers}.

\medterm{Eyeglasses, Sunglasses, And Magnifiers} Eyeglasses correct focusing errors, sunglasses reduce glare and ultraviolet exposure, and magnifiers enlarge nearby targets. Correct prescription, fit, lens protection, and the underlying eye condition determine the appropriate aid.

\synonyms
Corrective Eyewear; Optical Aids

\medterm{Eye Health} Maintaining vision and preventing avoidable eye disease through appropriate examinations, injury and ultraviolet protection, control of diabetes and blood pressure, safe contact-lens care, and prompt attention to new symptoms.

\medterm{Eye Hemorrhage} Bleeding on, inside, or around the eye. A small surface bleed may be painless, while bleeding in the retina, vitreous, or eye socket can impair vision or cause pressure and needs medical assessment.

\medterm{Eye Infection} Infection of the eyelid, conjunctiva, cornea, internal eye, or surrounding tissues caused by bacteria, viruses, fungi, or parasites. Pain, light sensitivity, reduced vision, or swelling around the eye needs prompt care.

\medterm{Eye Injury} Damage to the eye or nearby tissues from blunt or penetrating trauma, chemicals, foreign objects, radiation, or heat. Vision change, severe pain, bleeding, or an embedded object requires urgent assessment.

\medterm{Eye Irritation} Burning, itching, redness, or watering caused by dryness, allergy, smoke, chemicals, infection, or a foreign particle. Chemical exposure requires immediate prolonged rinsing with clean water and urgent medical advice.

\medterm{Eyelash} A hair growing from the eyelid edge that helps shield the eye from particles and triggers blinking when touched. Misrouted or infected lashes can irritate the eye surface.

\medterm{Eyelashes} Short curved hairs along the eyelid edges that help protect the eye from dust, moisture, and bright light. Inward-growing lashes can rub the cornea and require treatment.

\medterm{Eyelid} A movable fold of skin, muscle, and glands that protects the eye, spreads tears across its surface, and helps regulate light. Swelling, malposition, or weakness can impair protection and vision.

\medterm{Eyelid Bump} A localized eyelid swelling, often a painful infected gland called a stye or a firmer blocked oil gland called a chalazion. Warm compresses may help; persistent, recurrent, or vision-affecting bumps need examination.

\medterm{Eyelid Cancer} Cancer arising in the skin or other tissues of the eyelid, most commonly basal cell carcinoma. A persistent or changing lump, sore, bleeding area, loss of eyelashes, or eyelid distortion needs ophthalmic or dermatologic evaluation.

\medterm{Eyelid Coloboma} A birth defect causing a notch or missing section of the eyelid. It may leave the eye exposed, impair closure, and occur with other eye or facial abnormalities requiring specialist assessment.

\medterm{Eyelid Disorder} A disease or structural problem of the eyelid, such as infection, inflammation, a blocked gland, malposition, weakness, or tumor. It may cause pain, swelling, poor eye protection, irritation, or visual problems.

\medterm{Eyelid Drooping} Alternate name; see \textit{Ptosis}.

\medterm{Eyelid Inflammation} Alternate name; see \textit{Blepharitis}.

\medterm{Eyelid Lift} Alternate name; see \textit{Blepharoplasty}.

\medterm{Eyelid Neoplasm} A benign, precancerous, or cancerous growth of the eyelid. A persistent, enlarging, changing, bleeding, crusted, or ulcerated lesion needs examination by an eye or skin specialist.

\medterm{Eyelids} Movable folds of skin, muscle, and glands that protect the eyes, spread tears, and help control light. Their position and movement are important for keeping the eye surface moist and safe.

\medterm{Eyelid Swelling} Puffiness of one or both eyelids caused by allergy, infection, inflammation, injury, a blocked gland, or fluid accumulation. Fever, severe pain, impaired eye movement, bulging eye, or vision change may indicate an emergency.

\medterm{Eyelid Twitch} Alternate name; see \textit{Eye Twitching}.

\medterm{Eye Melanoma} A cancer arising from pigment-producing cells in the eye, most often in the uvea. It may cause blurred vision, flashes, a dark spot, or no symptoms; treatment depends on size, location, spread, and vision potential.

\synonyms
Uveal Melanoma; Ocular Melanoma; Choroidal Melanoma

\medterm{Eye Movement} Coordinated movement of the eyes that directs gaze and keeps images aligned. It depends on the eye muscles, their nerves, and brain control; abnormal movement may cause double vision or imbalance.

\medterm{Eye Muscle Repair} Treatment to correct a damaged or misaligned eye muscle or its attachment, usually by eye-muscle surgery. It may improve alignment or double vision but does not always restore normal movement and carries surgical risks.

\medterm{Eye Neoplasm} A benign or malignant growth arising in the eye or surrounding tissues. Symptoms and treatment depend on the tissue involved, size, location, effect on vision, and whether the growth has spread.

\medterm{Eye Pain} Pain in or around the eye caused by surface irritation, infection, injury, inflammation, glaucoma, or disease behind the eye. Severe pain with vision loss, light sensitivity, nausea, or trauma requires urgent assessment.

\medterm{Eye Prosthesis} An artificial eye or implant used to restore appearance after an eye is lost or severely damaged. It generally does not restore sight; fitting, cleaning, and periodic specialist care are needed.

\medterm{Eye Redness} Redness caused by enlarged or inflamed blood vessels on or inside the eye. Irritation and conjunctivitis are common causes, but pain, light sensitivity, discharge, injury, or vision change may indicate serious disease.

\medterm{Eyesight} Alternate name; see \textit{Vision}.

\medterm{Eyestrain} Tired, dry, sore, or uncomfortable eyes, sometimes with blurred vision or headache, after prolonged visual effort. Screen use, glare, poor lighting, dry eye, or an uncorrected focusing problem may contribute; rest often helps.

\medterm{Eye Trauma} Injury ranging from a surface scratch to a penetrating wound or ruptured eyeball. Rinse chemical exposure immediately, avoid pressure or rubbing, shield penetrating injuries, and obtain urgent care for vision change or severe pain.

\medterm{Eye Twitching} Repeated involuntary eyelid or eye-area spasms, commonly linked to fatigue, stress, caffeine, irritation, or eye strain. Persistent spasms, facial weakness, eye closure, or other neurologic symptoms need assessment.

\medterm{Ezetimibe} An oral medicine that reduces absorption of cholesterol from the intestine. It lowers blood cholesterol alone or with another medicine and may cause diarrhea, abdominal discomfort, or rarely liver or muscle problems.


\medterm{Fab Immunoglobulins} The target-binding part of an antibody. A laboratory-made Fab fragment can recognize and block a harmful substance or help remove it from the bloodstream in selected treatments.

\medterm{Fabry Disease} An inherited disorder in which low activity of an enzyme causes a fatty substance to accumulate in cells. It may cause burning pain in the hands and feet, skin spots, reduced sweating, kidney disease, heart disease, or stroke.

\medterm{Fabry'S Disease} An inherited enzyme disorder that causes a fatty substance to accumulate in cells. It may cause burning limb pain, skin findings, digestive symptoms, kidney or heart disease, and blood-vessel complications.

\medterm{Face} The front part of the head, including the eyes, nose, mouth, cheeks, and jaw. Its muscles produce facial expression, its nerves provide movement and sensation, and its structures support breathing, eating, speech, and vision.

\medterm{Face Pain} Pain in the face that may arise from teeth, sinuses, nerves, eyes, jaw joints, skin, blood vessels, or other tissues. Location, duration, triggers, swelling, fever, and neurologic signs help identify the cause.

\medterm{Face Powder Poisoning} Illness after swallowing or inhaling face powder or its ingredients. Coughing, choking, vomiting, breathing difficulty, drowsiness, or persistent symptoms require poison-control or emergency advice; do not induce vomiting.

\medterm{Face Presentation} A birth position in which the baby’s neck is extended so the face, rather than the top of the head, leads through the birth canal. The baby’s chin direction helps determine whether vaginal birth is safe.

\medterm{Facelift} Cosmetic surgery that removes or repositions selected facial and neck skin and deeper tissues to reduce sagging. It does not stop aging and may cause bleeding, infection, nerve injury, scarring, or temporary swelling and numbness.

\synonyms
Rhytidectomy

\medterm{Facet Joint} A small joint at the back of the spine connecting neighboring vertebrae and guiding movement. Arthritis, injury, or inflammation can cause localized neck or back pain and stiffness.

\medterm{Facet Joints} Small paired joints at the back of the spine that guide bending and rotation and help stabilize vertebrae. Wear, inflammation, or injury may cause localized neck or back pain.

\medterm{Facet Rhizotomy} A procedure that uses heat or another method to interrupt small nerves carrying pain signals from joints at the back of the spine. It may reduce selected chronic neck or back pain, but relief may be temporary.
\medterm{Facial Asymmetry} A difference in the size, position, movement, or shape of the two sides of the face. It may be present from birth or result from injury, nerve weakness, muscle disease, swelling, or a growth.

\medterm{Facial Cleft} A birth defect caused by incomplete joining of facial tissues during development. It may affect the lip, nose, mouth, or other facial structures and can interfere with feeding, speech, breathing, hearing, or appearance.

\medterm{Facial Edema} Swelling of the face caused by excess fluid in tissues. Allergy, infection, injury, kidney or heart disease, and blocked veins may cause it; sudden lip, tongue, or throat swelling with breathing difficulty is an emergency.

\medterm{Facial Expression} The visible movement or position of facial muscles that communicates emotion, pain, intention, or other information. Reduced or unequal expression can indicate facial-nerve or muscle dysfunction.

\medterm{Facial Hair, Excess In Women} Coarse dark hair on the face or other androgen-sensitive areas in a woman, called hirsutism. It may result from polycystic ovary syndrome, medicines, or rarely an adrenal or ovarian disorder; sudden rapid growth needs assessment.

\synonyms
Hirsutism

\medterm{Facial Muscle} A muscle that moves the skin of the face for expression and helps with speaking, eating, breathing through the mouth, or closing the eyes. Weakness may result from nerve, muscle, or brain disease.

\medterm{Facial Neoplasm} A benign or malignant growth arising in facial skin, soft tissue, bone, salivary glands, nerves, or other structures. Symptoms and treatment depend on its tissue of origin, size, location, and spread.

\medterm{Facial Nerve} The seventh cranial nerve, which controls facial expression and carries taste from part of the tongue. It also contributes to tear and saliva production and helps protect the ear from loud sounds.

\medterm{Facial Nerve Palsy} Weakness or paralysis of facial muscles caused by dysfunction of the seventh cranial nerve. It may affect eye closure, smiling, taste, tears, or saliva; sudden weakness with arm or speech symptoms may indicate a stroke.

\medterm{Facial Nerve Palsy Due To Birth Trauma} One-sided facial weakness in a newborn caused by pressure on the facial nerve during birth. It often improves as the nerve recovers, but eye protection, feeding, and other causes of weakness must be assessed.

\medterm{Facial Nerve Problems} Disorders affecting the nerve that controls facial movement and some taste, tear, saliva, and sound-protection functions. Causes include Bell palsy, stroke, infection, injury, or a tumor; sudden weakness needs emergency assessment.

\synonyms
Facial nerve disorder; facial nerve dysfunction

\medterm{Facial Nucleus} A group of nerve cells in the pons, part of the brainstem, whose fibers form the facial nerve and control facial-expression muscles. Injury can cause weakness of the face and related functions.

\medterm{Facial Pain} Pain in the face that may arise from dental or sinus disease, jaw-joint problems, nerve disorders, infection, injury, or other causes. The pattern and associated symptoms guide evaluation.

\medterm{Facial Palsy} Weakness or paralysis of facial muscles, usually on one side, caused by facial-nerve dysfunction. Bell palsy is one cause, but stroke, infection, injury, Lyme disease, and tumors are also possible; sudden onset with other neurologic signs is urgent.

\synonyms
Bell's palsy; facial paralysis

\medterm{Facial Paralysis} Partial or complete loss of facial-muscle movement caused by disease or injury of the facial nerve or its brain pathways. Sudden onset, especially with limb weakness or speech trouble, requires urgent stroke assessment.

\medterm{Facial Paresis} Weakness of muscles supplied by the facial nerve, causing reduced facial movement and sometimes poor eye closure, altered taste, or abnormal tearing. The cause may be nerve inflammation, injury, infection, stroke, or a tumor.

\medterm{Facial Spasm} Involuntary repeated tightening of facial muscles. It may affect one side around the eye or mouth and can result from nerve irritation, medicines, stress, or another neurologic disorder; persistent spasms need assessment.

\medterm{Facial Swelling} Enlargement of part or all of the face caused by infection, allergy, injury, dental or salivary disease, sinus disease, fluid retention, or a growth. Rapid swelling of the mouth or throat with breathing difficulty is an emergency.

\medterm{Facial Tics} Sudden, repeated, involuntary facial movements or sounds, such as blinking, grimacing, or nose wrinkling. They may increase with stress and occur in tic disorders; persistent or disruptive tics can be treated.

\medterm{Facial Trauma} Injury affecting facial skin, soft tissue, teeth, nerves, blood vessels, eyes, nose, or bones. Assessment must consider airway, bleeding, vision, brain injury, and structural damage.

\medterm{Facial Vein} A vein that drains blood from the superficial face toward the internal jugular vein. Connections with deeper veins mean facial infections can rarely spread toward the eye socket or brain.

\medterm{Facies} A characteristic facial appearance associated with a particular disease, inherited condition, or physiologic state. It can provide clues but is not diagnostic by itself and must be interpreted with other findings.
\medterm{Facilitated Diffusion} Movement of a substance across a cell membrane through a channel or carrier protein, from higher to lower concentration. The process does not directly use cellular energy and helps move substances that cannot cross the membrane unaided.

\medterm{Facioscapulohumeral Dystrophy} An inherited muscle disease causing gradually progressive weakness of the face, shoulder blades, upper arms, and sometimes legs or trunk. Onset and severity vary, and breathing, hearing, and eye complications may occur.

\medterm{Facioscapulohumeral Muscular Dystrophy} An inherited disorder causing progressive, often uneven weakness of facial, shoulder-blade, and upper-arm muscles. Symptoms may begin in adolescence or adulthood and can later involve the trunk, hips, or legs.

\medterm{Factitious Disorder} A mental-health disorder in which a person deliberately creates, exaggerates, or falsely reports illness to be treated as sick, without a clear outside reward. Care should be nonjudgmental and avoid unnecessary tests or procedures.

\synonyms
Munchausen Syndrome

\medterm{Factitious Disorders} Factitious disorder involves deliberately producing or falsifying illness behavior to assume the sick role, without an obvious external reward. It is a mental-health disorder requiring nonjudgmental assessment, protection from unnecessary procedures, and coordinated psychiatric and medical care.

\medterm{Factitious Hyperthyroidism} A state of excessive thyroid hormone in the body caused by ingestion of thyroid hormone or another external source rather than increased thyroid-gland production.

\medterm{Factor H Deficiency} A rare inherited or acquired problem in which a protein that keeps part of the immune system under control is missing or faulty. It can cause repeated infections or inflammation that damages small blood vessels, especially in the kidneys.

\medterm{Factor} A blood protein needed to form a clot. Different clotting factors work together to stop bleeding; too little or abnormal activity can cause a bleeding disorder, while excessive clotting can contribute to thrombosis.

\medterm{Factor I Deficiency} A rare condition in which a protein that helps control immune-system activity is missing or does not work properly. It may cause repeated bacterial infections or inflammation and injury to the kidneys.

\medterm{Factor II (Prothrombin) Assay} A blood test measuring factor II activity, which helps assess the clotting pathway and investigate abnormal bleeding, prolonged prothrombin time, or suspected prothrombin deficiency or inhibition.

\medterm{Factor IX} A vitamin-K-dependent blood-clotting protein that works with factor VIII to form clots. Inherited deficiency causes hemophilia B, leading to prolonged or spontaneous bleeding, especially into joints and muscles.
\medterm{Factor IX Assay} A blood test measuring factor IX activity, used to diagnose or monitor factor IX deficiency, the cause of hemophilia B, and to guide replacement treatment.

\medterm{Factor V} A blood-clotting protein that helps generate thrombin and fibrin, the mesh that strengthens a clot. Too little can cause bleeding; an inherited change affecting its regulation can increase venous-clot risk.

\medterm{Factor V Assay} A blood test measuring how well factor V supports clotting. It helps investigate abnormal bleeding or clotting-test results and diagnose inherited deficiency or an acquired reduction from liver disease or other illness.

\medterm{Factor V Deficiency} An inherited or acquired bleeding disorder caused by too little active factor V. It may cause easy bruising, nosebleeds, heavy menstrual bleeding, or excessive bleeding after injury, dental work, surgery, or childbirth.

\medterm{Factor V Leiden} An inherited change in factor V that makes it resist a natural clot-limiting protein. It increases the risk of clots in veins, but the risk varies with family history, pregnancy, hormones, surgery, and other factors.

\medterm{Factor V Leiden Mutation} A factor V gene variant that reduces the normal response to activated protein C, a clot-limiting substance. It causes inherited tendency toward venous thrombosis, although many carriers never develop a clot.

\medterm{Factor VII} A liver-made blood-clotting protein that starts clot formation when tissue is injured. Inherited or acquired deficiency can cause easy bruising, nosebleeds, or serious bleeding after injury or procedures.

\medterm{Factor VII Assay} A blood test measuring factor VII activity. It helps investigate a prolonged prothrombin time and identify inherited deficiency, liver disease, vitamin-K deficiency, or medicine-related reduction.

\medterm{Factor VII Deficiency} A bleeding disorder caused by too little or poorly functioning factor VII. Bleeding severity varies and may include bruising, nosebleeds, heavy periods, bleeding after surgery, or severe hemorrhage; prothrombin time is usually prolonged.

\medterm{Factor VIII} A blood-clotting protein that works with factor IX to form a stable clot. Inherited deficiency causes hemophilia A, with prolonged or spontaneous bleeding, often into joints, muscles, or other tissues.

\medterm{Factor VIII Assay} A blood test measuring factor VIII activity. It helps diagnose or monitor hemophilia A and investigate acquired reductions caused by inhibitors, liver disease, or other conditions.

\medterm{Factor X Assay} A blood test measuring factor X activity, a key clotting step shared by two clotting pathways. It helps investigate abnormal bleeding and diagnose or monitor inherited or acquired factor X deficiency.

\medterm{Factor X Deficiency} A rare inherited or acquired bleeding disorder caused by too little active factor X. Bleeding severity varies from bruising to serious hemorrhage, and both major clotting-time tests may be prolonged.

\medterm{Factor XI Deficiency} An inherited or acquired bleeding disorder caused by reduced factor XI activity. Bleeding is variable and is often noticed after surgery, dental procedures, childbirth, or trauma rather than occurring spontaneously.

\medterm{Factor Xii} A blood-clotting protein that starts one pathway measured by laboratory clotting tests. Factor XII deficiency can greatly prolong the activated partial-thromboplastin time but usually does not cause abnormal bleeding.

\medterm{Factor XII (Hageman Factor) Deficiency} A deficiency of factor XII that can greatly prolong one laboratory clotting test but usually does not cause bleeding. It is often discovered unexpectedly during testing for another reason.

\medterm{Factor XII Assay} A blood test measuring factor XII activity. It may help explain an isolated prolonged activated partial-thromboplastin time, especially when the person has no unusual bleeding.

\medterm{Factor Xiii} The clot-stabilizing protein that strengthens the fibrin mesh after a clot forms. Factor XIII deficiency can cause delayed bleeding, poor wound healing, or repeated pregnancy loss despite initially normal routine clotting tests.





\medterm{Facultative Anaerobe} A microorganism that can grow with or without oxygen. It uses oxygen when available but can switch to other energy-producing processes when oxygen is absent; many intestinal bacteria have this ability.
\medterm{Faecal Fistula} An abnormal connection between the intestine and another organ, the skin, or the body surface that allows stool or intestinal fluid to leak. It can cause infection, dehydration, skin injury, and poor nutrition.

\medterm{Faecalibacterium} A group of oxygen-sensitive bacteria commonly found in the intestine. They can produce butyrate, a substance that supports the colon lining; their presence is studied in relation to gut health and inflammation.

\medterm{Faeces} The solid waste passed from the intestine through the anus, also called stool. Its color, consistency, frequency, and contents can provide clues about digestion, bleeding, infection, inflammation, or other illness.

\medterm{Fahr Disease} A rare disorder involving abnormal calcium deposits in parts of the brain. It may cause movement problems, seizures, cognitive changes, or psychiatric symptoms; similar deposits can also result from other metabolic or genetic conditions.

\medterm{Failed Induction Of Labor} Failure of labor to progress after medicines or procedures intended to start it. The decision to continue induction or perform cesarean birth depends on cervical change, contractions, time, the mother’s condition, and fetal well-being.

\medterm{Failed Intubation} Inability to place a breathing tube safely in the windpipe. It is an airway emergency requiring prompt help, continued oxygen delivery, an alternative airway method, and a plan to avoid repeated traumatic attempts.

\medterm{Failure To Thrive} Poor weight gain or growth in a child. Causes may include feeding difficulty, inadequate intake, poor absorption, chronic disease, developmental problems, or social stress; assessment considers the child’s growth pattern, diet, development, and health.

\medterm{Fainting} A brief loss of consciousness and muscle tone caused by temporarily reduced blood flow to the brain, followed by spontaneous recovery. Reflex episodes are common, but fainting during exercise, with chest pain, or without warning needs urgent assessment.

\synonyms
syncope.

\medterm{Fainting (Syncope)} A temporary loss of consciousness and posture from reduced brain blood flow, followed by recovery. Causes include reflex reactions, standing-related blood-pressure drop, medicines, dehydration, and heart disease; warning features require urgent evaluation.

\medterm{Falciform Ligament} A fold of tissue attaching the liver to the front abdominal wall and diaphragm. It marks the division between the liver’s right and left lobes and contains the remnant of a fetal umbilical vein.

\medterm{Falciparum Malaria} Malaria caused by Plasmodium falciparum parasites transmitted by infected mosquitoes. It can worsen rapidly and cause brain involvement, severe anemia, kidney injury, breathing failure, shock, or death, so urgent treatment is required.

\medterm{Fall} An unplanned movement to the ground or another lower level. Falls may result from poor balance, weakness, vision problems, medicines, dizziness, low blood pressure, or environmental hazards and can cause serious injury.

\medterm{Fall Prevention} Steps to reduce falling, including strength and balance training, medication review, vision and hearing checks, safe footwear, adequate lighting, removal of hazards, and support with walking when needed.

\medterm{Fall Risk} The likelihood of falling based on factors such as previous falls, balance, walking ability, strength, vision, medicines, dizziness, blood-pressure changes, cognition, and environmental hazards.

\medterm{Fall Screening} A brief assessment of previous falls, balance, walking, strength, medicines, vision, dizziness, cognition, and surroundings to identify fall risk and guide prevention or further evaluation.

\medterm{Fallen Arches} Flattening of the foot’s inner arch, also called flat feet. It may be flexible or rigid and may cause no symptoms or lead to foot, ankle, knee, or back discomfort; treatment depends on symptoms and cause.

\synonyms
Pes planus; flat feet

\medterm{Fallopian Tube} One of two tubes between the ovaries and uterus. An egg usually travels through a tube, where fertilization may occur, before reaching the uterus; blockage or scarring can affect fertility.

\medterm{Fallopian Tube Cancer} A rare cancer beginning in a fallopian tube. It may cause pelvic pain, a mass, abnormal or watery discharge, or no early symptoms; diagnosis, staging, and treatment resemble those for ovarian cancer.

\medterm{Fallopian Tube Torsion} Twisting of a fallopian tube, sometimes around a cyst or fluid-filled tube, that can cut off blood flow. Sudden severe pelvic pain, nausea, or vomiting requires urgent surgical assessment.

\medterm{Fallopian Tubes} The two tubes connecting the ovaries and uterus. They transport eggs and are common sites of fertilization; blockage, inflammation, scarring, pregnancy, or tumors can affect them.

\medterm{Fallopian-Tube Endometriosis} Endometriosis involving or surrounding a fallopian tube. It may cause pelvic pain, adhesions, or infertility, or be found unexpectedly; treatment depends on symptoms, disease extent, and fertility goals.

\medterm{Fallot'S Tetralogy} A birth defect combining four heart abnormalities that reduce blood flow to the lungs and may cause bluish skin, breathing difficulty, poor growth, or fainting. Specialist care and usually surgery are required.

\medterm{Fallot, Tetralogy Of} A congenital heart defect consisting of four abnormalities: ventricular septal defect, right ventricular outflow tract obstruction, overriding aorta, and right ventricular hypertrophy. Causes cyanosis due to right-to-left shunting. Treatment includes surgical repair.

\synonyms
Tetralogy of Fallot

\medterm{Falls} Unplanned events in which a person comes to rest on the ground or another lower surface. Falls may cause fractures or head injury and often involve balance, weakness, vision, medicines, dizziness, or environmental hazards.

\medterm{False Labor} Contractions that may be uncomfortable but do not progressively open and thin the cervix. They are often irregular and may lessen with rest, hydration, or a change in activity, unlike established labor.

\medterm{False Negative} A test result that says a condition is absent when it is actually present. It can occur because testing was too early, the sample was inadequate, the condition level was low, or the test was not sensitive enough.

\medterm{False Passages} Abnormal channels made through tissue during instrumentation, surgery, or injury instead of following the normal body pathway. They can cause bleeding, perforation, infection, or failure to place a tube correctly.

\medterm{False Positive} A test result that indicates a condition is present when it is not. It may result from cross-reaction, contamination, technical error, or testing in a group where the condition is uncommon.
\medterm{False Ribs} The lower ribs that do not attach directly to the breastbone. Their cartilage joins the cartilage above or ends freely, allowing flexibility but leaving them vulnerable to injury in some areas.

\medterm{False-Negative Result} A negative result when the tested condition is actually present. Early testing, an inadequate sample, prior treatment, low levels of the target, technical error, or limited test sensitivity can cause it.

\medterm{False-Positive Result} A positive result when the tested condition is actually absent. Cross-reaction, contamination, technical error, limited specificity, or testing in a low-risk population can produce it.

\medterm{Faltering Weight} Slower-than-expected weight gain or a downward change in a child’s growth pattern. Feeding difficulty, inadequate intake, poor absorption, chronic disease, developmental problems, or social stress may contribute and require assessment.

\medterm{Famciclovir} An antiviral medicine that changes into penciclovir in the body and limits multiplication of herpes viruses. It can treat selected shingles, cold-sore, or genital-herpes episodes but does not remove the virus permanently.

\medterm{Familial} Occurring more often within a family because of shared genes, environment, behaviors, or a combination. Familial does not necessarily mean that one specific gene causes the condition.

\medterm{Familial Adenomatous Polyposis} An inherited condition causing numerous precancerous polyps in the colon and rectum. Without preventive treatment, colorectal cancer is very likely, so lifelong specialist screening and risk-reducing surgery are often needed.

\medterm{Familial Combined Hyperlipidemia} An inherited tendency to have high LDL cholesterol, high triglycerides, or both, often affecting several relatives. It increases early cardiovascular risk and is managed with lifestyle measures and lipid-lowering treatment when indicated.

\medterm{Familial Dysautonomia} An inherited nerve disorder affecting sensation and automatic functions such as blood pressure, heart rate, sweating, tears, digestion, and swallowing. Symptoms may include dizziness, temperature problems, repeated lung infections, and poor growth.

\medterm{Familial Dysbetalipoproteinemia} An inherited disorder in which the body clears certain cholesterol- and triglyceride-rich particles poorly. It can cause high blood fats, yellowish skin deposits, and premature artery disease, especially with other metabolic risk factors.

\medterm{Familial Hypercholesterolaemia} An inherited disorder causing very high LDL cholesterol from childhood. Without effective treatment, cholesterol can build up in artery walls and cause early heart attack, stroke, or disease of other arteries.

\medterm{Familial Hypercholesterolemia} An inherited disorder causing persistently high LDL cholesterol because the body removes it from blood poorly. Family history, cholesterol levels, physical findings, and sometimes genetic testing support the diagnosis.

\medterm{Familial Hypertriglyceridemia} An inherited tendency toward high blood triglycerides, often worsened by diabetes, obesity, alcohol, or certain medicines. Very high levels can cause abdominal pain, eruptive skin bumps, and acute pancreatitis.

\medterm{Familial Hypocalciuric Hypercalcemia} An inherited condition causing lifelong, usually mild high blood calcium with low urinary calcium. People often have no symptoms, and recognizing it prevents unnecessary parathyroid surgery; specialist testing may confirm the diagnosis.

\medterm{Familial Lipoprotein Lipase Deficiency} A rare inherited disorder in which triglycerides cannot be broken down normally. It causes extremely high triglycerides, recurrent pancreatitis, abdominal pain, and yellow-red skin bumps and requires specialist dietary and medical care.

\medterm{Familial Mediterranean Fever} An inherited inflammatory disorder causing repeated short attacks of fever with abdominal, chest, or joint pain. Without treatment, ongoing inflammation can damage the kidneys; preventive medicine can reduce attacks and complications.

\medterm{Familial Multiple Lipomatosis} An inherited tendency to develop many soft, slow-growing, usually painless fatty lumps under the skin, often on the trunk and limbs. They are generally benign, but painful, rapidly growing, or unusual lumps need assessment.

\medterm{Familial Multiple Trichoepithelioma} An inherited condition causing many benign tumors from hair follicles, usually small skin-colored facial bumps. They may increase over time and can affect appearance; treatment is chosen according to number, size, and location.

\medterm{Familial Paroxysmal Peritonitis} An inherited inflammatory disorder causing repeated short attacks of fever with severe abdominal, chest, or joint pain. Without treatment, ongoing inflammation can damage the kidneys; preventive treatment can reduce attacks and complications.

\synonyms
Familial Mediterranean fever

\medterm{Familial Testotoxicosis} An inherited condition in which boys produce excess testosterone without normal brain stimulation. It causes early sexual development, rapid growth, and advanced bone age but may reduce adult height without treatment.

\medterm{Family} People related by biology, law, or social bonds. Family history, shared genes, environment, and support relationships can affect health risks, decisions, caregiving, and recovery.

\medterm{Family Involvement} Participation of family members or chosen support people in healthcare planning, treatment, rehabilitation, or daily care, with the patient’s permission and attention to privacy, safety, preferences, and caregiver burden.

\medterm{Family Planning} Personal and healthcare decisions about whether and when to have children. It includes contraception, fertility care, preconception health, pregnancy planning, and support for informed reproductive choices.

\medterm{Family Study} A structured review of health information in relatives to understand how a disease or trait occurs in a family. It may include family history, medical records, examination, genetic counseling, or testing.

\medterm{Family Therapy} Counseling involving family members or other close relationships to improve communication, manage conflict, and support recovery. It may help with mental illness, addiction, chronic disease, caregiving stress, or family adjustment.


\medterm{Famophos} An organophosphate insecticide that can poison the nervous system by blocking acetylcholinesterase. Exposure may cause sweating, salivation, vomiting, pinpoint pupils, muscle weakness, breathing failure, or death and requires emergency care.

\medterm{Famotidine} A medicine that blocks histamine signals to reduce stomach-acid production. It treats reflux, heartburn, and ulcers; dose reduction may be needed in kidney disease, and headache, diarrhea, or confusion can occur.

\medterm{Fanconi Anemia} An inherited disorder affecting DNA repair and bone-marrow blood-cell production. It may cause low blood counts, infections, bleeding, short stature, birth differences, organ problems, and a high risk of leukemia or other cancers.

\medterm{Fanconi Syndrome} A kidney-tubule disorder causing excess loss of glucose, phosphate, bicarbonate, amino acids, and other substances in urine. It can lead to dehydration, abnormal blood acidity, weak bones, poor growth, or kidney disease.

\medterm{FAP} An inherited condition causing many precancerous polyps in the colon and rectum. Colorectal cancer is very likely without preventive treatment, so regular specialist surveillance and risk-reducing surgery are often needed.

\synonyms
Familial adenomatous polyposis

\medterm{Farad} The standard unit of electrical capacitance, equal to one coulomb per volt. It is used in medical equipment and electrophysiology to describe how much electrical charge a component can store.

\medterm{Farber Lipogranulomatosis} A rare inherited storage disorder caused by deficient acid ceramidase. Fatty material accumulates and may cause painful skin lumps, joint stiffness, hoarseness, bone changes, breathing problems, and neurologic disease.

\medterm{Farmer'S Lung} Lung inflammation caused by repeatedly inhaling moldy hay, grain, or other farm dust. It can cause fever, cough, chills, and breathlessness hours after exposure and may lead to permanent scarring if exposure continues.

\medterm{Farmer’S Lung} An immune lung disease caused by inhaling particles from moldy hay, grain, or other farm materials. Repeated exposure can cause cough, fever, and breathlessness and eventually permanent lung scarring; avoiding the source is central to care.

\medterm{Farsightedness} A focusing problem in which close objects look blurred because light would focus behind the retina. It may cause eyestrain or headaches and is usually corrected with glasses or contact lenses.

\synonyms
Hyperopia

\medterm{Fart} A colloquial term for passing intestinal gas through the anus. Gas comes from swallowed air and bacterial breakdown of food; frequent painful or newly changed gas may reflect diet, constipation, intolerance, or digestive disease.

\medterm{Fascia} Sheets or bands of connective tissue that surround, separate, support, and connect muscles, organs, nerves, and blood vessels. Injury or inflammation of fascia can cause pain, stiffness, or restricted movement.

\medterm{Fascia Adherens} A strong connection joining neighboring cells, especially heart-muscle cells. It links the internal contractile framework of one cell to the membrane of another so force can pass through the tissue.

\medterm{Fascial Loop} A loop or sling made from a person’s tissue or synthetic material to support a body structure. For example, a sling may support the urethra in selected treatment of stress urinary incontinence.

\medterm{Fasciculation} A brief, involuntary twitch visible under the skin from contraction of a small group of muscle fibers. It may be harmless, but widespread or persistent twitching with weakness needs neurologic evaluation.

\medterm{Fasciitis} Inflammation of the connective tissue around muscles and other structures. It may be localized, as in plantar fasciitis, or rapidly destructive, as in necrotizing infection requiring emergency surgery and antibiotics.

\medterm{Fasciola} A genus of parasitic flatworms, or liver flukes, that infect some animals and people. Infection is acquired by swallowing immature parasites on contaminated water plants or water and may affect the liver and bile ducts.

\medterm{Fasciola Hepatica} A liver fluke that infects people after contaminated aquatic plants or water are swallowed. Immature worms migrate through the liver, while adult worms live in bile ducts and can cause pain, fever, or obstruction.

\medterm{Fascioliasis} Infection with liver flukes after eating contaminated aquatic plants or drinking contaminated water. It may cause fever, right-upper-abdominal pain, enlarged liver, or high eosinophils and later cause bile-duct pain, infection, or blockage.

\medterm{Fasciolopsis} A genus of large intestinal flukes that can infect people who eat contaminated aquatic plants. Heavy infection may cause abdominal pain, diarrhea, swelling, poor nutrition, or intestinal obstruction.

\medterm{Fasciolopsis Buski} A large intestinal fluke acquired by eating aquatic plants carrying immature parasites. Heavy infection may cause abdominal pain, diarrhea, intestinal inflammation, malnutrition, or swelling.

\medterm{Fasciotomy} Emergency surgery that cuts tight fascia to release dangerous pressure in a muscle compartment. It may prevent permanent nerve and muscle damage after crush injury, fracture, burns, or loss and return of blood flow.


\medterm{Fast Heart Rate} A heart rate faster than expected for a person’s age and activity. Exercise, fever, pain, stress, dehydration, anemia, medicines, or an abnormal rhythm may cause it; chest pain or fainting is urgent.

\medterm{Fast Heart Rate - Tachycardia} A rapid heartbeat that may be a normal response to exercise, fever, pain, anxiety, dehydration, anemia, or medicines, or may reflect an abnormal rhythm. Chest pain, fainting, or severe breathlessness requires urgent care.

\medterm{Fast-Twitch Fiber} A muscle fiber capable of quick, powerful contraction but more rapid fatigue than a slow-twitch fiber. It supports sprinting, jumping, lifting, and other brief high-intensity activity.

\synonyms
Type II muscle fiber

\medterm{Fasting} Going without food, and sometimes without drinks, for a defined period. It may be required before blood tests, anesthesia, surgery, or certain procedures, or undertaken for personal or religious reasons.

\medterm{Fasting Blood Glucose} Blood sugar measured after no caloric intake, usually for at least eight hours. It helps screen for and diagnose diabetes or impaired glucose regulation, but results must be interpreted with the clinical context.

\medterm{Fasting Guidelines} Instructions about when to stop food and drinks before anesthesia, surgery, or a test. Timing depends on the procedure, age, health, and local policy; the treating team’s instructions must take priority.

\medterm{Fasting Lipid Profile} A blood test measuring cholesterol and triglycerides after a period without food. It helps assess cardiovascular risk and guide treatment; fasting is particularly useful when triglycerides need accurate measurement.

\synonyms
Lipid panel

\medterm{Fasting Plasma Glucose Test} A blood test measuring blood glucose after fasting, usually at least eight hours. It helps identify diabetes or prediabetes and is interpreted with repeat testing or other results when needed.

\medterm{Fat} Energy-rich substances that provide fuel, insulation, and components for cell membranes and hormones. Dietary fats also help absorb vitamins A, D, E, and K; types include saturated, unsaturated, and trans fats.

\synonyms
Lipid

\medterm{Fat Embolism} Entry of fat droplets into the bloodstream, usually after major trauma, a thigh or pelvic fracture, or orthopedic surgery. It may cause breathing difficulty, confusion, rash, or low oxygen and can become life-threatening.

\medterm{Fat Measurement} Estimation of body-fat amount or distribution using methods such as skinfolds, bioimpedance, DEXA, imaging, or laboratory analysis. Results vary by method and should be interpreted with overall health, strength, and metabolic findings.

\medterm{Fat Necrosis} Death of fat cells, often after pancreatitis, injury, or breast surgery. It can cause pain, inflammation, or a firm lump that resembles cancer on examination or imaging and may require testing to distinguish the cause.

\medterm{Fat Pad} A localized cushion of fatty tissue that protects, supports, insulates, or permits movement around an organ or joint. Inflammation or injury of a fat pad can cause localized pain and swelling.

\medterm{Fat-Soluble Toxins} Poisonous chemicals that dissolve in body fat and may remain for a long time. Some pesticides, industrial chemicals, polychlorinated biphenyls, and dioxins can accumulate after repeated exposure and harm several organs.

\medterm{Fat-Soluble Vitamin} A vitamin that dissolves in fat and can be stored in body tissues. Vitamins A, D, E, and K are fat-soluble; deficiency or excessive supplementation can cause illness.

\medterm{Fat-Soluble Vitamins} Vitamins A, D, E, and K, which dissolve in fat and can accumulate in the body. They support vision, bones, immunity, cell protection, and blood clotting; excessive supplements can be toxic.

\medterm{Fatal Familial Insomnia} A very rare inherited prion disease that progressively disrupts sleep, blood-pressure and temperature control, movement, and thinking. It causes severe neurologic decline and is ultimately fatal; specialist genetic and supportive care is needed.

\medterm{Fatigue} Persistent tiredness or reduced energy that limits physical or mental activity and may not improve fully with rest. Sleep problems, anemia, infection, hormone or organ disease, medicines, mood disorders, and neurologic conditions can cause it.

\medterm{Fatigue (Cancer)} Persistent tiredness or low energy related to cancer or its treatment. Anemia, pain, poor sleep, infection, medicines, emotional distress, and the cancer itself may contribute; evaluation can identify treatable causes.

\synonyms
Cancer-related fatigue

\medterm{Fatty Acid} A building block of fats that can provide energy, form cell membranes, and serve as a starting material for signaling chemicals. Fatty acids may be saturated, monounsaturated, or polyunsaturated.

\medterm{Fatty Acid Synthase} A group of linked enzymes that helps cells make fatty acids, especially palmitate. These fatty acids are used for energy storage, cell membranes, and other cellular functions.
\medterm{Fatty Acids} Components of fats used for energy, cell membranes, and chemical signals. Their effects on health depend on their type and amount.
\medterm{Fatty Degeneration} Abnormal buildup of fat inside cells, especially liver cells, caused by metabolic problems, alcohol, toxins, poor nutrition, obesity, or other injury. The change may reverse or progress to inflammation and scarring.

\medterm{Fatty Liver (Steatosis)} Excess fat stored in liver cells. Alcohol, obesity, diabetes, some medicines, and other causes may contribute; it may remain harmless or progress to liver inflammation, scarring, cirrhosis, or cancer.
\synonyms
Hepatic steatosis; fatty liver

\medterm{Fatty Liver} Alternate index label; see \textit{Fatty Liver (Steatosis)}.
\medterm{Fatty Liver Disease} Excess fat in the liver, often associated with obesity, diabetes, abnormal cholesterol, or alcohol use. Some people develop inflammation and scarring, so cause, metabolic health, and liver risk should be assessed.

\synonyms
Hepatic steatosis; steatotic liver disease

\medterm{Fatty Streak} An early artery-wall change in which cholesterol-rich immune cells collect beneath the lining. It can be an early stage of atherosclerosis, although some fatty streaks never progress to a harmful plaque.

\medterm{Favism} Rapid destruction of red blood cells after eating fava beans, usually in a person with glucose-6-phosphate dehydrogenase deficiency. It may cause jaundice, dark urine, weakness, breathlessness, or severe anemia and needs medical care.

\medterm{Favus} A long-lasting fungal infection of the scalp, usually causing yellow cup-shaped crusts around hairs. Without treatment it can destroy follicles and cause permanent scarring hair loss.

\medterm{Fc Receptor} A protein on an immune cell that binds the tail end of an antibody. This helps the cell recognize, ingest, or destroy the antibody’s target and can also influence inflammation.

\medterm{Fear} An emotional and body response to a perceived threat that prepares a person to protect themselves. Persistent or excessive fear can cause avoidance, panic, sleep problems, and impaired daily functioning.

\medterm{Fears} Feelings of threat or apprehension. Fears are normal when proportionate to danger, but persistent, excessive, or disabling fears may indicate an anxiety disorder and can often be treated.

\medterm{Febrile} Relating to fever, a temporary rise in body temperature usually caused by infection, inflammation, or another illness. Fever differs from heat-related hyperthermia; severe illness, dehydration, confusion, or fever with very low white cells is urgent.

\medterm{Febrile (Warm) And Cold Agglutinins} Antibodies that make red blood cells clump at warmer or colder temperatures. They can interfere with laboratory testing or cause autoimmune destruction of red cells and anemia.

\medterm{Febrile Convulsion} A seizure during fever in a young child without evidence of brain infection or another immediate cause. Most are brief and generalized, but a prolonged, repeated, or one-sided seizure requires emergency assessment.

\medterm{Febrile Convulsions} Seizures triggered by fever in young children, usually between six months and five years, without evidence of infection or another direct brain problem. Most resolve quickly, but first or prolonged seizures require medical assessment.

\medterm{Febrile Neutropenia} Fever in a person with a dangerously low neutrophil count, the white cells that fight infection. Infection may be present without obvious symptoms, so urgent examination, cultures, and prompt antibiotics are often needed.

\medterm{Febrile Seizure} A seizure during fever in a child, usually between six months and five years, without evidence of brain infection or a previous seizure without fever. Most are brief, but prolonged, repeated, or one-sided events require urgent assessment.

\synonyms
Seizure, Febrile

\medterm{Febrile Seizures} Seizures associated with fever in young children without evidence of brain infection or a previous unprovoked seizure. Most do not cause lasting harm, but the first seizure or any prolonged event needs medical evaluation.

\synonyms
Seizure Febrile (Febrile Seizures)
Seizure Fever Induced (Febrile Seizures)


\medterm{Febuxostat} A medicine that lowers uric-acid production for long-term prevention of gout attacks. It does not relieve sudden gout pain immediately and may require monitoring for liver problems and cardiovascular risk.

\medterm{FEC} Alternate index label; see \textit{FECD}.

\medterm{Fec Protocol} A chemotherapy regimen combining fluorouracil, epirubicin, and cyclophosphamide for selected cancers. It can cause low blood counts, infection, nausea, hair loss, mouth sores, infertility, and heart or bladder-related toxicity.

\medterm{Fecal Calprotectin} A stool test measuring a protein released during intestinal inflammation. A raised result can support evaluation for inflammatory bowel disease, but infection, medicines, age, and other conditions can also raise it.

\medterm{Fecal Culture} Laboratory testing that grows bacteria from a stool sample to investigate infectious diarrhea. It may identify a treatable cause and guide antibiotics, but some infections require toxin, molecular, or parasite testing instead.

\medterm{Fecal Elastase} A stool test estimating whether the pancreas releases enough digestive enzymes. A low result may suggest pancreatic insufficiency, causing poor digestion, greasy stools, weight loss, or vitamin deficiency.

\medterm{Fecal Fat} Measurement of excess fat in stool. High fecal fat suggests poor digestion or absorption from pancreatic disease, reduced bile, intestinal disease, or another condition and may explain greasy, difficult-to-flush stools.

\medterm{Fecal Immunochemical Test} A stool screening test that detects human blood, mainly from the lower intestine, using antibodies. A positive result does not diagnose cancer but usually requires colon evaluation, often colonoscopy.

\medterm{Fecal Impaction} A large, hard mass of stool trapped in the rectum or colon. It may cause pain, abdominal swelling, nausea, urinary symptoms, bowel blockage, or watery leakage around the mass and may require professional removal.

\medterm{Fecal Incontinence} Inability to control passage of stool or gas, ranging from soiling to complete loss of bowel control. Causes include diarrhea, constipation, sphincter injury, nerve disease, rectal prolapse, and cognitive impairment.

\synonyms
Incontinence, Bowel
Incontinence, Fecal

\medterm{Fecal Microbiota Transplant} Transfer of carefully screened and prepared donor stool or its microbial contents into the intestine to restore helpful microbes, mainly for selected recurrent C. difficile infection. Infection transmission remains a possible risk.

\medterm{Fecal Microbiota Transplantation} Transfer of screened donor stool or microbial material into a patient’s intestine to restore microbial balance. It is used mainly for selected recurrent C. difficile infections after standard therapy and requires regulated donor screening.

\medterm{Fecal Occult Blood Test} A stool test that detects blood too small to see. It may screen for colorectal cancer or investigate digestive-tract bleeding, but it does not identify the cause; a positive result usually requires further colon evaluation.

\medterm{Fecal Smear} Examination of a thin stool sample with a microscope or stain. Depending on the test, it may help identify parasites, inflammatory cells, blood, fat, or other abnormalities.

\medterm{Fecalith} A hard mass of dried stool, often in the appendix or colon. It may contribute to appendicitis, blockage, inflammation, or perforation.

\medterm{Fecaluria} Stool, gas, or a fecal smell in urine, usually indicating an abnormal connection between the intestine and urinary tract. It often accompanies recurrent urinary infection and requires medical evaluation.

\medterm{FECD} An abbreviation for Fuchs endothelial corneal dystrophy, a progressive loss of corneal cells that normally remove fluid. It can cause morning blur, glare, painful swelling, and reduced vision; advanced disease may need a corneal transplant.

\synonyms
Fuchs dystrophy

\medterm{Feces} Solid or semisolid waste passed from the bowel, containing water, undigested material, bacteria, shed cells, and metabolic waste. Changes in stool can provide clues about digestion, bleeding, infection, or inflammation.

\medterm{Fecolith} A hard mass of dried stool that can block the appendix or bowel and contribute to inflammation, appendicitis, or perforation.

\medterm{Fecundity} The biological ability to conceive and produce a live birth. It depends on age, ovulation, ovarian reserve, fallopian tubes, uterus, sperm health, and other factors; it differs from the chance of pregnancy in one cycle.

\medterm{Feedback} Information about a system’s output that changes later activity. Negative feedback stabilizes functions such as temperature, glucose, and hormones, while positive feedback strengthens processes such as labor contractions.

\medterm{Feedback Inhibition} A control process in which the end product of a chemical pathway slows an earlier step. This prevents unnecessary buildup of the product and helps cells use nutrients and energy efficiently.

\medterm{Feedback Mechanism} A regulatory process in which a system’s result influences what happens next. Negative feedback maintains stability, while positive feedback amplifies a response until a defined endpoint is reached.

\medterm{Feeding Jejunostomy} Placement of a feeding tube into the jejunum, part of the small intestine, when eating or stomach feeding is unsafe or impossible. It provides nutrition but can cause infection, blockage, leakage, or bowel injury.

\medterm{Feeding Methods} Ways of providing food or liquid, including breast or bottle feeding, oral feeding, tube feeding, and intravenous nutrition. The safest method depends on swallowing ability, age, illness, and nutritional needs.

\medterm{Feeding Patterns And Diet - Babies And Infants} Guidance on how babies receive breast milk, formula, and later complementary foods. Advice should match age, growth, development, swallowing safety, allergies, and medical needs.

\medterm{Feeding Patterns And Diet - Children 6 Months To 2 Years} Guidance on introducing varied foods and drinks while continuing appropriate milk feeding. Portion size, texture, allergy prevention, choking safety, growth, and iron intake are important.

\medterm{Feeding Poor} Taking too little food or fluid for age or health needs. In an infant it may signal illness, swallowing difficulty, breathing problems, pain, reflux, or developmental difficulty and requires assessment if persistent.

\medterm{Feeding Tube} A flexible tube used to deliver nutrition, fluids, or medicines to the stomach or intestine when eating is inadequate or unsafe. Risks include blockage, leakage, infection, aspiration, and injury.

\medterm{Feeding Tube - Infants} A tube delivering breast milk, formula, fluids, or medicines to an infant unable to feed safely or adequately by mouth. The route, feeding plan, and duration depend on swallowing, breathing, growth, and the underlying illness.

\medterm{Feeding Tube Insertion - Gastrostomy} Placement of a tube through the abdominal wall into the stomach to provide nutrition, fluids, or medicines. It may be temporary or long-term and can cause bleeding, infection, leakage, blockage, or injury.

\medterm{Feingold Diet} A restrictive diet that avoids certain food additives and some naturally occurring chemicals, promoted for behavior or attention problems. Evidence of benefit is limited, and restriction can cause nutritional problems without professional guidance.

\medterm{Felbamate} An antiseizure medicine reserved for selected severe or drug-resistant epilepsy. It can rarely cause life-threatening bone-marrow failure or liver failure, so specialist prescribing and regular blood and liver tests are essential.

\medterm{Felodipine} A calcium-channel blocker that relaxes artery walls and lowers blood pressure. It may cause ankle swelling, flushing, headache, fast heartbeat, or gum enlargement and should be taken as prescribed.

\medterm{Felon} A painful bacterial infection in the fleshy fingertip, often after a small puncture. Pus and pressure can damage tissue, bone, or tendons, so prompt assessment and sometimes drainage and antibiotics are needed.

\medterm{Felty Syndrome} A rare complication of long-standing rheumatoid arthritis involving a very low neutrophil count and often an enlarged spleen. It increases infection risk and requires specialist treatment of arthritis and neutropenia.

\medterm{Felty'S Syndrome} A rare complication of long-standing rheumatoid arthritis causing neutropenia and often an enlarged spleen. It raises the risk of serious infection and may require disease-modifying treatment, growth factors, or other specialist care.

\medterm{Female} Describing the sex category generally associated with ovaries, eggs, and female reproductive anatomy. Biological traits vary among individuals, and sex, gender identity, and gender expression are distinct concepts.

\medterm{Female Athlete Triad} A pattern of low energy availability, menstrual disturbance, and reduced bone strength in physically active females. Inadequate nutrition or excessive exercise can increase the risk of stress fractures and other health problems.

\medterm{Female Breast} A mammary gland and its skin, fat, connective tissue, ducts, and lobules on the front of the chest. Its structure changes with age, hormones, pregnancy, breastfeeding, and body composition.

\medterm{Female Circumcision} Removal or injury of part of the external female genitalia for nonmedical reasons. It has no health benefit and can cause severe pain, bleeding, infection, urinary and sexual problems, childbirth complications, and lasting psychological harm.

\medterm{Female Condom} A soft internal sheath placed in the vagina before sex that forms a barrier against semen and can reduce pregnancy and transmission of some sexually transmitted infections when used correctly.

\medterm{Female Condoms} Flexible internal barriers placed in the vagina before sex to help prevent pregnancy and reduce exposure to semen and some sexually transmitted infections. They should not be used at the same time as an external condom.

\medterm{Female Genital Mutilation} Cutting or removing external female genital tissue for nonmedical reasons. It has no health benefit and can cause pain, bleeding, infection, urinary and sexual problems, childbirth complications, and long-term psychological harm.

\medterm{Female Genitalia} The external and internal reproductive organs, including the vulva, vagina, cervix, uterus, fallopian tubes, and ovaries. These structures support urination, sexual function, menstruation, fertility, pregnancy, and birth.

\medterm{Female Gonad} An ovary, the reproductive gland that stores and releases eggs and produces hormones including estrogen and progesterone.

\medterm{Female Infertility} Difficulty achieving pregnancy because of problems with ovulation, ovarian reserve, fallopian tubes, uterus, hormones, or other factors. Evaluation considers both partners and may be recommended after a defined period of trying or sooner with risk factors.

\synonyms
Infertility, Female

\medterm{Female Orgasmic Disorder (Anorgasmia)} Persistent or recurrent difficulty achieving orgasm that causes personal distress. It may be lifelong or acquired and can relate to health conditions, medicines, mood, trauma, sexual pain, stimulation, or relationship factors.

\synonyms
Anorgasmia

\medterm{Female Pattern Baldness} Gradual scalp hair thinning, usually widest near the part or crown, influenced by genes and hormones. Iron deficiency, thyroid disease, medicines, and other causes should be considered when loss is new or rapid.

\medterm{Female Pelvis} The pelvic bones and soft tissues in a person with typical female anatomy. It is often broader with a wider outlet than typical male anatomy, but substantial individual variation exists.

\medterm{Female Phenotype} The observable anatomical, physiologic, and developmental traits associated with female sex. These traits vary with genes, hormones, development, age, and individual biology.

\medterm{Female Reproductive System} The ovaries, fallopian tubes, uterus, cervix, vagina, vulva, and related hormones. It supports egg release, menstruation, pregnancy, birth, and sexual function and can be affected by developmental, hormonal, infectious, structural, or cancerous disorders.

\medterm{Female Sexual Dysfunction} Persistent sexual difficulty causing distress, such as low desire, arousal or lubrication problems, difficulty with orgasm, or pain. Medical conditions, medicines, mood, trauma, relationships, and life changes may contribute.

\synonyms
Sexual Dysfunction, Female

\medterm{Female Sexual Interest/Arousal Disorder} Persistent or recurrent low sexual interest or arousal causing personal distress. Medical conditions, medicines, mood, trauma, pain, relationship factors, hormonal changes, and other social or psychological factors may contribute.

\medterm{Female Urethral Opening} The opening through which urine leaves the female urethra, located in the vulvar vestibule in front of the vaginal opening. Its position and appearance may vary normally.

\medterm{Feminization} Development of physical or hormonal traits commonly associated with female puberty, such as breast growth or reduced body hair. It may be normal, part of gender-affirming treatment, or caused by hormones, medicines, or disease.

\medterm{Femoral} Relating to the femur or thigh, including the thigh bone, artery, vein, nerve, or neck of the femur.

\medterm{Femoral Artery} A major artery carrying blood through the groin and thigh to the leg. Its pulse can be felt in the groin, and it is used for some catheter procedures; injury or blockage can threaten the limb.

\medterm{Femoral Hernia} A bulge of abdominal tissue through an opening below the groin ligament and beside the femoral vessels. It has a relatively high risk of trapping or losing blood supply to the contents and often needs surgical assessment.

\medterm{Femoral Hernia Repair} Surgery that returns trapped tissue to the abdomen and closes or reinforces the opening in the groin. Repair is often recommended because femoral hernias can become trapped or strangulated.

\medterm{Femoral Nerve} A large nerve from the lower back that supplies muscles that flex the hip and straighten the knee and provides sensation to the front thigh and inner leg. Injury can cause weakness, numbness, or falls.

\medterm{Femoral Nerve Dysfunction} Reduced function of the nerve supplying the front thigh and knee-straightening muscles. It can cause hip or knee weakness and altered sensation and may result from diabetes, compression, injury, surgery, or pelvic disease.

\medterm{Femoral Neuropathy} Damage to the femoral nerve causing weakness in hip flexion or knee extension and numbness or pain in the front thigh or inner leg. Causes include diabetes, injury, bleeding, compression, and pelvic disease.
\medterm{Femoral Pulse} The pulse felt over the femoral artery in the groin. Comparing it with pulses elsewhere helps assess blood flow to the leg and can provide clues about arterial blockage or shock.

\medterm{Femoral Shaft Fractures} Breaks in the middle portion of the thigh bone, often from high-energy trauma but sometimes from weakened bone. They can cause major blood loss, deformity, and inability to walk and usually require urgent stabilization and specialist treatment.

\medterm{Femoral Triangle} A triangular area at the upper front of the thigh bordered by the groin ligament and nearby muscles. It contains the femoral nerve, artery, vein, and lymph nodes and is an important examination and access site.

\medterm{Femoral Vein} The major deep vein of the thigh, carrying blood from the leg toward the pelvis and heart. A clot in it can cause deep-vein thrombosis and pulmonary embolism.

\medterm{Femtogram} A unit of mass equal to one quadrillionth of a gram, or $10^{-15}$ gram. It is used in laboratory measurements of extremely small quantities.

\medterm{Femtoliter} A unit of volume equal to one quadrillionth of a liter, or $10^{-15}$ liter. It is used to describe very small laboratory or cellular volumes.

\medterm{Femur} The thigh bone and the longest, strongest bone in the body. It joins the hip to the knee and includes a head, neck, shaft, and lower end; fractures can threaten mobility and blood supply.

\medterm{Femur Fracture} A break in the thigh bone caused by trauma or weakened bone. It may cause severe pain, deformity, blood loss, and inability to bear weight; treatment depends on location, alignment, and overall health.


\medterm{Femur Length} The ultrasound-measured length of a fetus’s thigh bone. It helps estimate pregnancy age and assess growth, but interpretation depends on other measurements, family size, pregnancy dating, and possible skeletal or chromosomal conditions.

\medterm{Fencer'S Pose} A posture in which turning the head makes the arm and leg on that side straighten while the opposite arm and leg bend. It is a normal reflex in young infants but may indicate a brain or nerve problem if it persists or appears unexpectedly.

\medterm{Fenfluramine} Fenfluramine is a serotonergic medicine used in some jurisdictions for selected severe epileptic encephalopathies; historical appetite-suppressant use was limited by valvular heart disease and pulmonary hypertension, so cardiac monitoring is required.

\medterm{Fenitrothion} A poisonous organophosphate insecticide. Exposure can cause sweating, excess saliva, vomiting, diarrhea, small pupils, muscle twitching, breathing difficulty, seizures, or weakness because it disrupts nerve signals; suspected poisoning needs urgent medical help.

\medterm{FeNO (Fractional Exhaled Nitric Oxide)} A breath test that estimates type 2, often eosinophilic, airway inflammation. The result is interpreted with symptoms, age, treatment, and other clinical findings.

\synonyms
Fractional Exhaled Nitric Oxide

\medterm{Fenobucarb} A carbamate insecticide that can cause nerve poisoning after significant absorption, producing sweating, vomiting, small pupils, muscle weakness, breathing difficulty, or seizures.

\medterm{Fenofibrate} A medicine that lowers blood triglycerides and can modestly change cholesterol levels by increasing the liver's processing of fats. It is used for selected lipid disorders; liver, kidney, muscle, and gallbladder risks require monitoring.

\medterm{Fenoldopam} A short-acting dopamine D1-receptor agonist used intravenously for severe hypertension in selected settings; it can cause hypotension, flushing, headache, and increased intraocular pressure.

\medterm{Fenoprofen} Fenoprofen is a nonsteroidal anti-inflammatory drug used for pain and inflammatory arthritis; gastrointestinal bleeding, renal injury, fluid retention, and cardiovascular events are important risks.

\medterm{Fenoprofen Calcium Overdose} Taking too much fenoprofen, an anti-inflammatory pain medicine, can cause nausea, vomiting, stomach bleeding, drowsiness, kidney injury, breathing or acid-base problems, seizures, or coma. Suspected overdose requires immediate poison-control or emergency advice; do not induce vomiting.

\medterm{Fenoterol} A short-acting beta-2 bronchodilator that relaxes airway muscle and relieves bronchospasm. It may cause tremor, nervousness, low potassium, or a fast or irregular heartbeat and is not approved everywhere.

\medterm{Fenoxycarb} An insect-growth regulator that imitates juvenile hormone and disrupts development of immature insects. It is used for pest control, and exposure should be minimized because products can irritate skin, eyes, or airways.

\medterm{Fentanyl Citrate} The citrate salt of fentanyl, a potent opioid analgesic used for severe pain and anesthesia; respiratory depression, sedation, dependence, and overdose are major risks.

\medterm{Fentanyl Derivatives} Strong synthetic opioid medicines related to fentanyl, including sufentanil, remifentanil, and alfentanil. They may relieve severe pain or support anesthesia but can rapidly slow breathing, cause overdose, and produce dependence, so close monitoring is essential.

\medterm{Fenthion} A poisonous organophosphate insecticide. Exposure can cause sweating, excess saliva, vomiting, diarrhea, small pupils, muscle twitching, breathing difficulty, seizures, or weakness because it disrupts nerve signals; suspected poisoning needs urgent medical help.

\medterm{Fermentation} A metabolic process in which microorganisms or cells convert organic substrates without using oxidative respiration, producing products such as acids, alcohols, or gases.

\medterm{Ferritin} A protein that stores iron inside cells and releases it when needed. Blood ferritin often reflects body iron stores, but inflammation, infection, liver disease, or cancer can raise it even when iron is low.

\medterm{Ferritin Blood Test} A blood test measuring stored iron indirectly. Low ferritin strongly supports iron deficiency; high ferritin may reflect iron excess, inflammation, infection, liver disease, or cancer and must be interpreted with other tests.

\medterm{Ferroptosis} A controlled form of cell death caused by iron-dependent damage to cell fats and failure of antioxidant protection. It differs from apoptosis and ordinary tissue necrosis and is being studied in cancer and other diseases.

\medterm{Ferrous Sulfate} An oral iron supplement used to prevent or treat iron-deficiency anemia. It may cause nausea, constipation, abdominal discomfort, or dark stools; accidental overdose can be life-threatening, especially in children.

\medterm{Fertile Period} The several days in a menstrual cycle when sex can result in pregnancy, including the days before and shortly after ovulation. Timing varies, and calendar estimates are not reliable contraception for everyone.

\medterm{Fertilisation} Union of a sperm and an egg to form a zygote, the first cell of an embryo. It usually occurs in a fallopian tube before the developing embryo moves toward the uterus.

\medterm{Fertility} The biological ability to conceive and produce offspring. It depends on ovulation, sperm production, reproductive anatomy, hormones, age, general health, and other factors affecting either partner.

\medterm{Fertility Domain} The part of health concerned with reproductive function, ability to conceive, fertility goals, infertility, preservation of fertility, and related medical, emotional, and social care.

\medterm{Fertility Drugs} Medicines used to stimulate ovulation, support reproductive treatment, or improve the chance of pregnancy. They can cause side effects and multiple pregnancy, so selection and monitoring require fertility-specialist care.

\medterm{Fertility Preservation} Steps taken to protect the possibility of future pregnancy before treatment or disease may damage reproductive function. Options include freezing eggs, embryos, sperm, or selected reproductive tissue, depending on age, timing, and goals.

\medterm{Fertilization} Fusion of sperm and egg genetic material to form a zygote, usually in a fallopian tube. This begins embryonic development, which may continue after the zygote reaches the uterus and implants.

\medterm{Ferumoxytol} An intravenous iron preparation used for selected adults with iron-deficiency anemia. It can cause allergic reactions, low blood pressure, or iron overload, so monitoring is required during and after infusion.

\medterm{FESS} Functional endoscopic sinus surgery uses instruments passed through the nostrils to open blocked sinus drainage pathways and remove selected diseased tissue. It may improve chronic sinus disease but can cause bleeding, infection, eye injury, or fluid leakage.

\medterm{Festinating Gait} Walking with increasingly short, quick steps and a forward-leaning posture, commonly associated with Parkinson disease or other movement disorders. It increases fall risk and may improve with specialist treatment and physical therapy.

\medterm{Fetal Acid-Base Status} An assessment of a fetus’s oxygen supply and blood acidity, often using umbilical-cord blood gases at birth. Results must be interpreted with the sampling site, timing, fetal condition, and complete clinical record.

\medterm{Fetal Alcohol Spectrum Disorder} Lifelong physical, learning, behavioral, or developmental effects caused by alcohol exposure before birth. Effects vary widely and may include growth problems, attention difficulties, poor coordination, or characteristic facial features.

\synonyms
Fetal Alcohol Spectrum Disorders


\medterm{Fetal Alcohol Syndrome} The most severe recognized form of fetal alcohol spectrum disorder, involving characteristic facial features, poor growth, and brain-development problems after prenatal alcohol exposure. No known safe amount of alcohol exists during pregnancy.

\synonyms
Alcohol Pregnancy (Fetal Alcohol Syndrome)

\medterm{Fetal Attitude} The position of a fetus’s head, arms, legs, and trunk in relation to one another. Normal flexion helps birth; unusual extension or flexion can change the presenting part and affect labor.

\medterm{Fetal Bradyarrhythmia} An unusually slow or irregular fetal heart rhythm. Causes include conduction block, extra beats, maternal or fetal disease, or medicines; persistent abnormalities require prompt specialist assessment.

\medterm{Fetal Bradycardia} A fetal heart-rate baseline below 110 beats per minute for at least ten minutes. Medicines, maternal illness, fetal heart disease, low oxygen, cord problems, or placental separation may cause it and persistent cases require urgent obstetric assessment.

\medterm{Fetal Circulation} Blood flow before birth in which the placenta supplies oxygen and nutrients and special connections direct much blood around the lungs, which are not yet used for breathing. These connections normally close after birth.

\medterm{Fetal Death} Death of a fetus before delivery. Medical and legal definitions vary by gestational age and location; care includes confirming the diagnosis, supporting the family, and discussing delivery and follow-up.

\medterm{Fetal Development} Growth and maturation from the end of the embryonic period until birth. Organs enlarge and mature, movements develop, and body systems gradually acquire the functions needed after delivery.

\medterm{Fetal Disease} A disease, structural abnormality, or developmental problem affecting a fetus before birth. Causes include genetic changes, infection, medicines, toxins, and poor placental function; ultrasound and prenatal tests may identify it.

\medterm{Fetal Distress} A nonspecific term for concerning fetal findings during labor, especially abnormal heart-rate patterns that may indicate low oxygen, cord compression, infection, or placental problems. Clinicians assess the whole pattern and may act urgently.

\medterm{Fetal Echocardiography} A specialized ultrasound examination of the fetal heart used to assess cardiac structure, rhythm, and function, including possible congenital heart disease.

\medterm{Fetal Growth Restriction} A fetus that is smaller than expected because it may not be reaching its growth potential, often because the placenta is not supplying enough oxygen or nutrients. Repeat ultrasound, fluid, blood-flow, and well-being checks guide care.

\medterm{Fetal Heart} The developing heart of a fetus, whose rate and rhythm can be assessed by ultrasound or electronic monitoring to evaluate fetal well-being.

\medterm{Fetal Heart Rate} The number of times a fetus’s heart beats each minute, usually about 110–160 during pregnancy and labor. It can be checked with a Doppler or electronic monitor; changes may reflect normal activity, medicines, or reduced oxygen.

\medterm{Fetal Hemoglobin} The main oxygen-carrying protein in a fetus’s red blood cells. It binds oxygen more strongly than adult hemoglobin, helping the fetus obtain oxygen from the placenta; its amount normally falls after birth.

\medterm{Fetal Hydrops} Severe fluid buildup in at least two parts of a fetus, such as under the skin, around the lungs or heart, or in the abdomen. Causes include anemia, heart disease, infection, chromosome conditions, and blood-group incompatibility and require specialist care.

\medterm{Fetal Hypoxia} Too little oxygen reaching a fetus before or during birth. Placental disease, maternal breathing or circulation problems, cord compression, or prolonged labor can contribute; abnormal heart-rate patterns require prompt obstetric assessment.

\medterm{Fetal Kick Counts (Fetal Movement Counting)} Observation and recording of fetal movements by a pregnant person. A noticeable reduction from the usual pattern should prompt timely obstetric advice.

\synonyms
Fetal Movement Counting

\medterm{Fetal Lie} The direction of a fetus’s long body axis compared with the pregnant person’s spine. The fetus may lie lengthwise, sideways, or diagonally; a sideways or diagonal position near birth can affect delivery planning.

\medterm{Fetal Macrosomia} A fetus that is larger than expected, often estimated at 4,000–4,500 grams or more at birth depending on the clinical definition. Diabetes and previous large babies increase risk; complications include difficult birth, injury, and cesarean delivery.

\medterm{Fetal Malpresentation} A fetal presentation other than the usual vertex presentation, including breech, transverse, or oblique lie, which may affect labor and delivery planning.

\medterm{Fetal Monitoring} Assessment of fetal heart rate and, sometimes, uterine contractions to evaluate fetal wellbeing during pregnancy or labor.

\medterm{Fetal Monitoring, Electronic} Continuous or intermittent recording of fetal heart rate and uterine activity using external sensors or an internal electrode when indicated; patterns are interpreted with gestational age and labor context.

\medterm{Fetal Presentation} The part of the baby positioned to enter the pregnant person’s pelvis first. The usual presentation is head first; breech, shoulder, face, brow, and mixed presentations can affect labor and delivery planning.

\medterm{Fetal Pyelectasis} Dilation of the fetal renal pelvis seen on prenatal imaging; it may be transient or associated with urinary tract obstruction or aneuploidy and requires context-specific follow-up.

\medterm{Fetal Scalp Electrode} A small electrode placed on the fetus’s scalp through the cervix to record the heart rate directly during labor. It may give a clearer signal when external monitoring is inadequate but can cause infection, bleeding, or scalp injury.

\medterm{Fetal Scalp PH Testing} A test that analyzes a small blood sample from the baby's scalp during labor to assess whether reduced oxygen may be causing acid buildup. It may clarify an abnormal fetal heart-rate pattern, but use is limited and depends on local expertise and clinical circumstances.

\medterm{Fetal Scalp Stimulation} A test during labor in which the baby’s scalp is gently touched during a vaginal examination. A brief rise in the baby’s heart rate may be reassuring, but the result must be interpreted with the full heart-rate recording and clinical situation.

\medterm{Fetal Stage} The fetal stage begins after the embryonic period and continues until birth, characterized mainly by growth and maturation of established organs.

\medterm{Fetal Tachycardia} A sustained fetal heart-rate baseline above 160 beats per minute for at least 10 minutes. Causes may include maternal fever, infection, medicines, fetal hypoxia, or fetal arrhythmia.

\medterm{Fetal Tissue} Fetal tissue is tissue from a developing fetus and may be relevant to pregnancy care, pathology, research, transplantation, or genetic testing under applicable ethical and legal rules.

\medterm{Fetal Ultrasound} Imaging during pregnancy using sound waves to assess the fetus, placenta, fluid, growth, movement, and anatomy. It can estimate gestational age, identify some structural problems, assess growth or blood flow, and guide selected procedures.

\medterm{Fetal Viability} The possibility that a fetus could survive outside the uterus with intensive medical care. It depends on pregnancy age, fetal condition, treatment available, and other factors and is not defined by one fixed week for every pregnancy.

\medterm{Fetal Weight} An estimate of a fetus’s mass, usually calculated from ultrasound measurements. It helps assess growth and identify possible growth restriction or unusually large size, but it is not exact.

\medterm{Fetal-Maternal Erythrocyte Distribution Blood Test} A test that estimates the amount of fetal red blood cells in maternal circulation, helping guide Rh immune globulin dosing after fetomaternal hemorrhage.

\medterm{Fetishism} A recurrent sexual interest in a nonliving object or a specific non-genital body part; it is a disorder only when it causes clinically significant distress or impairment or involves harm or lack of consent.

\medterm{Fetishistic Disorder} A mental-health disorder involving repeated, intense sexual arousal from a nonliving object or a specific non-genital body part. It is diagnosed only when it causes substantial distress or impairment or involves harm or lack of consent.

\medterm{Fetomaternal Hemorrhage} Passage of fetal blood into the pregnant person’s bloodstream through the placenta. Small amounts can occur normally, but heavy bleeding across the placenta after trauma, placental separation, or a procedure can cause fetal anemia and affect preventive treatment.

\medterm{Fetomaternal Transfusion} Movement of fetal blood cells across the placenta into the pregnant person's bloodstream. Small amounts are common, but heavy transfer after trauma, placental problems, or procedures can cause fetal anemia and may require urgent testing and treatment.

\medterm{Fetoscope} An instrument used to hear or assess the fetal heartbeat through the pregnant person's abdominal wall, especially later in pregnancy.

\medterm{Fetus} A developing human from the start of the ninth week after fertilization until birth. During this period, established organs grow and mature; medical descriptions may instead use pregnancy weeks counted from the last menstrual period.

\medterm{FEV1/FVC Ratio} The amount of air forcefully exhaled in the first second divided by the total amount exhaled. A lower-than-expected ratio suggests narrowed airways, but interpretation uses age, sex, height, reference values, and bronchodilator response.

\medterm{Fever} A rise in body temperature caused by the body’s response to infection, inflammation, certain medicines, heat illness, cancer, or autoimmune disease.

\medterm{Fever And Headache} Fever with headache may occur with viral illness, influenza, sinus disease, dehydration, meningitis, encephalitis, or other infection. Severe sudden headache, neck stiffness, confusion, rash, seizure, weakness, or reduced consciousness requires emergency assessment.

\medterm{Fever Blister} A small, painful blister on or around the lips caused by herpes simplex virus, usually type 1. The virus remains in nerve tissue and may reactivate with illness, stress, or sunlight. Blisters usually heal in about one to two weeks and can spread by close contact.

\synonyms
Cold sore; herpes labialis

\medterm{Fever Of Unknown Origin} Persistent or recurrent fever whose cause remains unexplained after an appropriate initial evaluation. Causes may include infection, inflammation, cancer, or medicines and require systematic reassessment.
\medterm{Fever, Acute, In Adults} A fever in an adult with recent onset, commonly caused by infection, inflammation, medication, heat illness, or other conditions. Severity, vital signs, immune status, source, and associated symptoms guide evaluation; confusion, breathing difficulty, stiff neck, shock, or severe localized pain requires urgent care.

\medterm{Fever, Chronic} A fever that continues or repeatedly returns over a long period. Causes include infection, inflammatory disease, cancer, and medicines; persistent unexplained fever needs medical evaluation based on its pattern and accompanying symptoms.

\medterm{Fever, Valley} Coccidioidomycosis: a fungal infection caused by Coccidioides species, endemic to arid regions of the southwestern US and parts of Mexico and Central/South America. Inhalation of spores causes respiratory illness ranging from asymptomatic to severe pneumonia. Dissemination can affect skin, bones, and meninges.

\synonyms
Valley fever; coccidioidomycosis

\medterm{Fevers, Viral Hemorrhagic} A group of serious viral illnesses that can damage blood vessels and affect several organs. Some cause bleeding, shock, or organ failure. They spread in different ways, including through mosquitoes, ticks, rodents, infected animals, or people; urgent isolation and specialist care may be needed.

\synonyms
Viral hemorrhagic fevers

\medterm{Fexofenadine} An antihistamine used to relieve sneezing, itching, runny nose, and watery eyes from allergies and itching from hives. It usually causes less drowsiness than older antihistamines, but headache or nausea can occur.

\medterm{Fiber} Dietary fiber is plant material that the small intestine cannot digest. It adds bulk to stool, supports regular bowel movements, and some types help lower blood cholesterol and slow rises in blood glucose.

\medterm{Fiber Optic} A technology that sends light through thin glass or plastic fibers. Flexible fiber-optic instruments are used in endoscopy to view body passages and guide some minimally invasive procedures.

\medterm{Fibrates} A group of medicines that mainly lower high triglyceride levels by changing how the liver and tissues handle fats. They can cause indigestion, gallstones, liver problems, or muscle injury, especially with some other lipid medicines.

\medterm{Fibre} A threadlike structure in the body or another material. In nutrition, fibre is plant material that is not digested in the small intestine and helps regulate bowel movements.

\medterm{Fibrillation} Rapid, disorganized activity of muscle fibers. In the heart, atrial fibrillation causes an irregular rhythm and can lead to blood clots; ventricular fibrillation stops effective pumping.

\medterm{Fibrin} An insoluble, threadlike protein that forms a mesh at an injured blood vessel. The mesh traps blood cells and strengthens the clot that helps stop bleeding.

\medterm{Fibrin Degradation Products Blood Test} A blood test that measures fragments released when blood clots are broken down. High levels may occur with widespread clotting, severe infection, liver disease, recent surgery, or pregnancy, so results require clinical interpretation.

\medterm{Fibrin Sealant} A surgical adhesive containing clotting proteins that form a fibrin clot when applied to tissue. It can help control bleeding or seal selected tissues, but it does not replace repair of major blood-vessel injuries.

\medterm{Fibrinogen} A blood protein made mainly by the liver that is changed into fibrin during clot formation. Low levels can cause bleeding; levels may rise with inflammation and fall in severe liver disease or widespread clotting.

\medterm{Fibrinogen Assay} A laboratory test that measures the amount or clotting activity of fibrinogen in blood. It helps investigate unexplained bleeding, abnormal clotting, liver disease, or suspected widespread clotting disorders.

\medterm{Fibrinogen Blood Test} A blood test measuring fibrinogen, a liver-made protein needed to form blood clots. Low levels may contribute to bleeding, while high levels commonly accompany inflammation, infection, pregnancy, or cardiovascular risk.

\medterm{Fibrinogen Decreased} A lower-than-normal blood fibrinogen level, which can weaken clot formation and cause bleeding. Causes include inherited deficiency, severe liver disease, major blood loss or transfusion, and widespread activation of clotting.

\medterm{Fibrinoid Necrosis} Severe damage to a blood-vessel wall in which blood proteins and immune material collect within the wall. It can occur in some inflammatory vessel diseases and is identified by examining affected tissue.

\medterm{Fibrinolysis} The natural process that dissolves a blood clot after healing begins. An enzyme called plasmin breaks down fibrin; excessive breakdown can cause bleeding, while insufficient breakdown can allow harmful clots to persist.

\medterm{Fibrinolysis - Primary Or Secondary} Primary fibrinolysis is early clot breakdown without preceding widespread clot formation. Secondary fibrinolysis follows clotting activation, as the body breaks down fibrin formed during a clotting disorder or injury.

\medterm{Fibrinolytic Agent} A medicine that dissolves a blood clot by helping the body convert plasminogen into plasmin, which breaks down fibrin. It may treat selected emergencies but can cause serious bleeding, including bleeding in the brain.

\medterm{Fibrinolytic Therapy} Treatment that uses clot-dissolving medicine to restore blood flow in selected emergencies, such as certain strokes, heart attacks, or pulmonary embolisms. The main risk is serious bleeding, so eligibility is assessed urgently.

\synonyms
Thrombolytic therapy

\medterm{Fibrinopeptide A Blood Test} A blood test that measures a protein fragment released when the body begins converting fibrinogen into fibrin. It can indicate clotting activity but is not specific enough to diagnose a condition by itself.

\medterm{Fibrinous Pericarditis} Inflammation of the sac surrounding the heart, with a sticky protein called fibrin coating its surfaces. It may cause sharp chest pain, a scratchy heart sound, and fluid around the heart.

\medterm{Fibroadenoma} A common noncancerous breast lump made of glandular and connective tissue. It is usually firm, smooth, movable, and painless, especially in adolescents and young adults; imaging helps distinguish it from other breast masses.

\medterm{Fibroadenoma Of The Breast} A noncancerous breast lump containing glandular and connective tissue. It is commonly smooth, firm, movable, and painless, but any new or changing breast lump should be examined and appropriately imaged.

\medterm{Fibroblast} A cell that makes collagen and other supporting materials in connective tissue. Fibroblasts help maintain skin, tendons, and organs and become active during wound healing and scar formation.

\medterm{Fibrocartilage} A tough type of cartilage that combines strength with limited flexibility. It helps absorb pressure and resist pulling in the discs between spinal bones, the pubic joint, and the knee menisci.

\medterm{Fibrocystic Breast Changes} Common noncancerous breast changes involving thickened tissue or fluid-filled cysts. Breast fullness, lumpiness, and tenderness often vary with hormonal changes during the menstrual cycle.

\medterm{Fibrocystic Breast Condition} A common noncancerous pattern of breast lumpiness, tenderness, or cysts that may change during the menstrual cycle. A persistent or changing lump, skin change, nipple inversion, or bloody discharge needs medical assessment.

\medterm{Fibrocystic Breast Disease} An older name for common noncancerous breast changes causing lumpiness, swelling, cysts, or tenderness. Symptoms may vary with the menstrual cycle, and the changes themselves do not increase breast-cancer risk.

\medterm{Fibrocystic Breasts} Common noncancerous changes in breast tissue that may cause lumpiness, thickening, cysts, swelling, tenderness, or discomfort. Symptoms may change during the menstrual cycle and usually improve after menopause. Fibrocystic breast changes do not increase the risk of breast cancer.

\medterm{Fibrodysplasia Ossificans Progressiva} A rare inherited disorder in which painful episodes gradually turn muscles, tendons, and other soft tissues into bone. This progressive extra bone can severely restrict movement and complicate medical procedures.

\synonyms
FOP; Stone Man Syndrome

\medterm{Fibroepithelial Polyp Of The Vulva} A usually noncancerous, soft tissue growth on the vulva or nearby vaginal tissue. It may enlarge during pregnancy or the reproductive years and can cause irritation, bleeding, or discomfort if large.

\medterm{Fibrofolliculoma} A small, usually noncancerous skin growth arising around a hair follicle, often on the face, neck, or upper trunk. Multiple lesions may occur with an inherited disorder that also raises kidney-tumor risk.

\medterm{Fibroid} A noncancerous growth of smooth muscle, most often in the uterus. It may cause heavy menstrual bleeding, pelvic pressure, pain, frequent urination, fertility problems, or no symptoms.

\synonyms
Uterine leiomyoma

\medterm{Fibroid Tumor} A noncancerous growth of the uterus made from smooth muscle and connective tissue. It may cause heavy or painful periods, pelvic pressure, frequent urination, infertility, or no symptoms.

\medterm{Fibroids} Noncancerous growths in the muscle of the uterus. They may cause heavy or painful menstrual bleeding, pelvic pressure, frequent urination, constipation, fertility problems, or no symptoms.

\medterm{Fibroids, Cysts, And Polyps} Three usually noncancerous types of growth: fibroids arise from uterine muscle, cysts are fluid-filled sacs often found in the ovaries, and polyps are tissue projections from the uterine lining or cervix.

\medterm{Fibroids, Uterine} Noncancerous smooth-muscle growths in the uterus. They may cause heavy bleeding, pelvic pain, pressure, frequent urination, or fertility problems; management depends on symptoms, size, location, age, and pregnancy plans.

\synonyms
Uterine fibroids; leiomyomas

\medterm{Fibrolipoma} A usually noncancerous, slow-growing lump made of mature fat and fibrous tissue. It is often painless and lies beneath the skin, but a new or enlarging lump should be assessed.

\medterm{Fibroma} A usually noncancerous tumor made of fibrous connective tissue. Cardiac fibromas most often occur in a child’s heart muscle or dividing wall and may obstruct blood flow or disturb the heartbeat.

\medterm{Fibroma Of Tendon Sheath} A usually noncancerous, slow-growing lump attached to the covering of a tendon, most often in a finger, hand, or wrist. It is firm and may recur after removal.

\medterm{Fibromatoses} A group of connective-tissue growth disorders that do not spread to distant organs but may grow into nearby structures and return after treatment. Effects depend on the affected body site.

\medterm{Fibromatosis} An overgrowth of connective-tissue cells that does not spread to distant organs but can invade nearby tissue and recur. Desmoid-type disease may cause pain, a firm mass, or impaired organ function.

\medterm{Fibromuscular Dysplasia} An abnormal structure of medium-sized artery walls, most often in the kidney or neck arteries. It can cause high blood pressure, headache, or stroke, or produce an artery bulge or tear.

\medterm{Fibromyalgi} Alternate index label; misspelling of \textit{Fibromyalgia}.

\medterm{Fibromyalgia} A long-term condition involving increased sensitivity to pain, with widespread aching, fatigue, unrefreshing sleep, and problems with concentration or memory. Diagnosis is based on symptoms and examination; no single laboratory test confirms it.


\medterm{Fibromyalgia Syndrome} A long-term pain condition causing widespread body pain, fatigue, poor-quality sleep, and difficulty thinking or concentrating. Doctors diagnose it from the symptom pattern and examination, while treatment combines activity, sleep support, and selected medicines.

\synonyms
Fibromyalgia

\medterm{Fibromyxosarcoma} A rare cancer of soft tissue containing fibrous and jelly-like material. It may present as a growing lump or pain; diagnosis requires tissue examination, and treatment depends on its site and spread.

\medterm{Fibronectin} A protein found in blood and connective tissue that helps cells attach, move, and organize during tissue repair. It supports wound healing and the structure of tissues around cells.

\medterm{Fibroproliferation} Excess growth of connective-tissue cells with increased production of scar tissue. It can thicken or stiffen organs, including the lungs, and may develop after injury, inflammation, or surgery.

\medterm{Fibrosarcoma} A rare cancer arising from cells that make fibrous connective tissue. It usually develops in deep soft tissue or bone and may cause a growing mass, pain, or reduced movement.

\medterm{FibroScan} A painless ultrasound-based test that estimates liver stiffness. Greater stiffness can suggest scarring, but inflammation, congestion, obesity, or fluid in the abdomen can affect accuracy and results may need other tests.

\medterm{Fibrosing Alopecia} Hair loss caused by inflammation and permanent scarring of hair follicles. The affected skin may look smooth or shiny, and regrowth is unlikely once follicles are destroyed; early treatment may limit progression.

\medterm{Fibrosis} Excess formation of scar tissue after ongoing injury or inflammation. Fibrosis can make an organ stiff and reduce its function, such as causing breathlessness in the lungs or impaired filtration in the kidneys.

\medterm{Fibrosis, Interstitial Pulmonary} Scarring in the supporting tissue between the lung air sacs, which can reduce oxygen transfer. It may cause progressive breathlessness and dry cough; causes include autoimmune disease, medicines, radiation, or no identified cause.

\synonyms
Interstitial lung disease

\medterm{Fibrosis, Pulmonary} Scarring of lung tissue that makes the lungs stiff and reduces oxygen exchange. It often causes gradually worsening breathlessness and dry cough; examination, lung tests, and scans help identify the cause and extent.

\synonyms
Pulmonary fibrosis

\medterm{Fibrositis} An outdated term formerly used for muscle and soft-tissue pain. It is generally avoided because true inflammation is usually not shown and a more specific diagnosis, such as fibromyalgia, may be appropriate.

\medterm{Fibrous Dysplasia} A noncancerous bone disorder in which normal bone is replaced by weaker abnormal tissue. It may affect one or several bones and cause pain, deformity, fractures, or pressure on nearby nerves.

\medterm{Fibrous Plaque} A thickened area in an artery wall containing scar-like tissue, cholesterol, and other material. It can narrow blood flow or rupture, allowing a blood clot to form and block the artery.

\synonyms
Atherosclerotic plaque

\medterm{Fibula} The slender bone on the outer side of the lower leg. It provides muscle attachment and forms the outer ankle. The nerve near its upper end can be injured by pressure or fracture.

\medterm{Fick Principle} A method for estimating blood flow from how much oxygen the body uses and the difference in oxygen content between blood entering and leaving the lungs. It can estimate cardiac output during specialized heart testing.

\synonyms
Fick Method; Fick Cardiac Output Determination

\medterm{Fiducial Markers} Small visible markers placed on or inside the body to identify a precise location during imaging, radiation treatment, or surgery. They help clinicians align equipment and target the intended tissue accurately.

\medterm{Field Fever} An older name for leptospirosis, a bacterial infection acquired through water or soil contaminated with urine from infected animals. It may cause fever, headache, muscle pain, jaundice, kidney injury, or meningitis.

\medterm{Field Of View} The area included in a medical image or examination. A wider field shows more of the body, while a narrower field provides greater detail about a smaller area.

\medterm{Fifth Disease} A usually mild infection caused by parvovirus B19. Children may develop fever, cold-like symptoms, bright-red cheeks, and a lacy rash. It can cause serious fetal anemia during pregnancy and prolonged anemia in some vulnerable people.

\synonyms
Erythema Infectiosum (Fifth Disease)

\medterm{FIGO Leiomyoma Classification} A system that describes where a uterine fibroid lies in relation to the uterine cavity and outer surface. The location category helps clinicians predict symptoms and choose treatment.

\medterm{Filamentous Fungus} A fungus that grows as branching threads called hyphae, which form a network called mycelium. Some filamentous fungi cause skin, nail, lung, or invasive infections, especially in people with weakened immunity.
\medterm{Filarial Elephantiases} Severe enlargement and thickening of body parts caused by long-term lymphatic obstruction from filarial worm infection. It most often affects the legs, arms, breasts, or genitals and develops gradually.

\medterm{Filarial Elephantiasis} Long-lasting swelling and thickening of skin and underlying tissues caused by lymphatic blockage from filarial worms. It commonly affects the legs or genitals and may require antiparasitic treatment, skin care, and swelling management.

\medterm{Filariasis} Infection with threadlike parasitic worms spread by mosquitoes, flies, or other insects. Depending on the species, worms may affect lymph vessels, skin, or eyes, causing swelling, inflammation, itching, or vision problems.

\medterm{Filler} A material injected or placed in tissue to restore volume, smooth wrinkles, change contours, or replace lost structure. Risks include pain, infection, lumps, and accidental blockage of a blood vessel.

\medterm{Filler Excipient} An inactive ingredient added to a medicine to provide bulk, improve stability, or help form a tablet, capsule, or liquid. It may rarely cause allergy or intolerance.

\medterm{Filopodium} A thin, finger-like projection from a cell surface. It senses nearby surroundings and helps cells attach, move, communicate, or form connections, including connections made by developing nerve cells.

\medterm{Filoviridae} A family of viruses that includes Ebola and Marburg viruses. They can cause severe fever, bleeding, organ failure, and death, and are mainly spread through direct contact with blood or other infected body fluids.

\medterm{Filoviruses} Viruses such as Ebola and Marburg that can cause severe hemorrhagic fever. Symptoms may include fever, headache, muscle pain, vomiting, diarrhea, bleeding, and organ failure; spread occurs through infected body fluids.

\medterm{Filter (Vascular)} A small device placed in a large vein, usually the inferior vena cava, to catch clots traveling toward the lungs. It is used selectively when anticoagulant treatment is unsuitable or ineffective.

\synonyms
Vena cava filter; inferior vena cava filter

\medterm{Filtered Back Projection} A computer method for creating an image from scan measurements, especially in computed tomography. It is fast, but images may contain more noise or artifacts than those made with newer reconstruction methods.

\medterm{Filtering Procedure} Eye surgery that creates a new drainage route for fluid to lower pressure inside the eye. It may treat glaucoma when medicines or other procedures do not control pressure adequately.

\medterm{Filtration} Separation of particles or dissolved substances from a fluid by passing it through a barrier. In the kidneys, tiny filters remove water and small waste products from blood to begin urine formation.

\medterm{Filtration Fraction} The percentage of plasma flowing through the kidneys that is filtered to begin urine formation. It helps describe kidney blood flow and filtering function and is calculated from kidney filtration compared with plasma flow.

\medterm{Filum Terminale} A thin strand of connective tissue that anchors the lower end of the spinal cord to the tailbone. Abnormal attachment can restrict cord movement and cause symptoms of tethered spinal cord.

\medterm{Fimbria} A fringe-like structure. In the fallopian tube, the fimbria are finger-like projections near the ovary that help guide a released egg into the tube; the term is also used for bacterial attachment structures.

\medterm{Fimbriae} Finger-like projections at the end of each fallopian tube that help collect an egg released from the ovary and guide it into the tube. Their movement and tiny surface hairs help direct the egg toward the uterus.

\medterm{Fimbriae Proteins} Bacterial surface proteins that form hair-like projections used to attach bacteria to human cells, medical devices, or other surfaces. Attachment can help bacteria form communities and begin infection.

\medterm{Fine Crackles} Brief, high-pitched, popping sounds heard mainly during breathing in. They may occur when small airways or air sacs open in conditions such as lung scarring, fluid congestion, or other lung disease.

\medterm{Fine Motor Control} The coordinated use of small muscles for accurate movements, especially of the hands and fingers. It supports tasks such as writing, buttoning clothes, using utensils, and handling small objects.

\medterm{Fine Motor Skills} Abilities that use small muscles for controlled movements, especially in the hands and fingers. Children develop skills such as grasping, feeding themselves, drawing, dressing, and writing at different individual rates.

\medterm{Fine Needle Aspiration} A procedure using a thin needle to remove cells or fluid from a lump or organ for laboratory examination. It usually causes little tissue injury and may be guided by touch or ultrasound.

\medterm{Fine Needle Aspiration Of The Thyroid} A procedure using a thin needle, often guided by ultrasound, to collect thyroid cells from a nodule. Laboratory examination helps determine whether the nodule is benign, malignant, or needs further testing.

\medterm{Finegoldia} A group of bacteria that grow without oxygen and normally live on skin or mucous membranes. They can occasionally cause wound, soft-tissue, bone, joint, or bloodstream infections.

\medterm{Finegoldia Magna} A bacterium normally found on skin and mucous membranes that grows without oxygen. It can cause infections after wounds or surgery, including abscesses and deeper soft-tissue infections.

\medterm{Finger} One of the hand’s digits, made of bones, joints, muscles, tendons, nerves, blood vessels, and skin. Fingers support grasping, touch, balance, and precise movements.

\medterm{Finger Dislocation} An injury in which the bones of a finger joint are forced out of their normal alignment. Pain, swelling, deformity, and limited movement are common; fractures, tendon damage, and impaired circulation must also be checked.

\medterm{Finger Flexors} Muscles and tendons that bend the fingers at their joints. Injury can cause pain or inability to bend a finger and may require prompt assessment, especially after a cut or forceful injury.

\medterm{Finger Injury} Damage to a finger’s bone, joint, tendon, ligament, nerve, blood vessel, nail, or skin, usually from trauma. Assessment checks alignment, movement, sensation, blood flow, and possible hidden fractures.

\medterm{Finger Joint} A connection between the bones of a finger that allows bending, straightening, and limited side-to-side movement. Finger joints can be affected by injury, arthritis, inflammation, infection, or dislocation.

\medterm{Finger Pain} Discomfort in one or more fingers caused by injury, arthritis, tendon problems, nerve compression, infection, inflammation, or reduced blood flow. Sudden severe pain, deformity, numbness, or discoloration needs prompt care.

\medterm{Finger Stick} A small fingertip puncture used to collect a drop of capillary blood. The sample can be tested immediately, commonly for blood glucose or other point-of-care measurements.

\medterm{Fingernail} The hard protective plate covering the end of a finger, made mainly of keratin and growing from the nail root. Changes in color, shape, or growth can sometimes signal injury or disease.

\medterm{Fingers That Change Color} Fingers that become white, blue, purple, or red because of changes in blood flow, temperature, or oxygen. Recurrent attacks may be Raynaud phenomenon; persistent changes, pain, numbness, or sores need medical assessment.

\medterm{Finkelstein Test} A wrist examination that gently stresses the tendons on the thumb side. Pain there may support de Quervain tenosynovitis, but the result is not specific and must be interpreted with symptoms and examination.

\medterm{Fire Ants} Ants whose stings inject venom and cause immediate burning pain, redness, and itchy pus-filled bumps. Multiple stings can cause severe swelling, vomiting, breathing difficulty, or a life-threatening allergic reaction.


\medterm{Firesetting Behavior} Deliberately starting fires. Repeated fire-setting may be linked to an impulse-control disorder, mental illness, substance use, developmental problems, or social stress and requires assessment of safety and risk to others.

\medterm{First Aid} Immediate basic care for injury or sudden illness before professional treatment arrives. It includes ensuring safety, calling emergency services, checking breathing, starting resuscitation when needed, and controlling severe bleeding.

\medterm{First Aid Kit} A group of supplies for treating minor injuries and giving temporary care during emergencies. Useful items include gloves, dressings, bandages, antiseptic, scissors, and emergency contact information.

\medterm{First Stage Of Labor} The time from regular labor contractions that cause cervical change until the cervix is fully open. It includes early and active phases, with progress varying according to the person, fetus, and clinical circumstances.

\medterm{First Trimester} The first twelve weeks of pregnancy, when the embryo’s organs begin forming and the placenta develops. Nausea, fatigue, and breast tenderness are common; prenatal care includes dating and early health assessment.

\medterm{First-Degree AV Block} A mild delay in the electrical signal traveling from the heart’s upper chambers to its lower chambers. Every signal still reaches the lower chambers; it is often symptomless but may reflect medicines or heart disease.

\medterm{First-Degree Burn} A superficial burn affecting only the outer skin layer. It causes redness, pain, and dryness without blisters, as with mild sunburn, and usually heals within several days without scarring.

\medterm{First-Degree Heart Block} A delay in the electrical signal between the heart’s upper and lower chambers, shown on an electrocardiogram by a prolonged but constant interval. It is often harmless, though medicines or heart disease may contribute.

\medterm{First-Episode Psychosis} The first recognized period of losing contact with reality, such as hearing or seeing things that are not present or holding fixed false beliefs. Urgent assessment looks for medical, medication, substance, and psychiatric causes.

\medterm{First-Pass Effect} The reduction of an orally swallowed medicine by the intestine and liver before it reaches the bloodstream. This can reduce the amount available to act and may make another route necessary.

\medterm{First-Pass Metabolism} Breakdown of a swallowed medicine by the intestinal wall and liver before it enters general circulation. Extensive breakdown lowers the active amount reaching the body and can influence dose or route of administration.

\medterm{Fish Oil} Oil containing omega-3 fatty acids, usually obtained from oily fish. Supplements may lower triglycerides at prescribed doses, but benefits vary and side effects or medicine interactions should be reviewed with a clinician.

\medterm{Fish Poisoning} Illness caused by toxins or germs in contaminated or improperly stored fish. Symptoms may include vomiting, diarrhea, flushing, tingling, weakness, or breathing difficulty; severe symptoms require urgent medical care.

\medterm{Fish Tapeworm Infection} An intestinal infection acquired by eating raw or undercooked fish containing tapeworm larvae. It may cause abdominal discomfort, diarrhea, weight loss, or vitamin B12 deficiency and is usually treatable with antiparasitic medicine.

\medterm{Fish Venom} Poison produced by certain fish, usually delivered through spines or sharp body parts. It can cause immediate severe pain, swelling, numbness, tissue injury, or rarely whole-body poisoning.

\medterm{Fish-Handler'S Disease} A skin infection caused by Erysipelothrix bacteria after handling infected fish, shellfish, meat, or animals. It usually causes a slowly expanding, painful purple-red lesion on the hand or fingers.

\medterm{Fisher'S Syndrome} A rare nerve disorder, usually after an infection, causing eye-muscle weakness, poor coordination, and reduced or absent reflexes. It is a variant of Guillain-Barré syndrome and needs prompt medical assessment.

\medterm{Fishhook Removal} Removal of a fishhook from skin or deeper tissue. The method depends on the hook’s location and depth; clinicians assess nerves, blood vessels, tendons, joints, tetanus protection, and infection risk.

\medterm{Fissure} A narrow crack, groove, or split in a body structure. An anal fissure is a painful tear in the anal lining, often causing sharp pain and bright-red bleeding during bowel movements.

\medterm{Fissured Tongue} A usually harmless condition in which grooves develop on the tongue’s surface. Food may collect in the grooves and cause irritation or odor; gentle tongue cleaning and routine oral care can help.

\medterm{Fistula} An abnormal passage connecting two organs or a body organ to the skin. It may develop because of infection, inflammation, surgery, injury, cancer, or abnormal development and can cause persistent drainage.

\medterm{Fistula (Pathology)} Alternate index label; see \textit{Fistula}.

\medterm{Fistula, Anal} An abnormal tunnel between the anal canal or rectum and nearby skin, often left after an abscess. It may cause persistent drainage, pain, skin irritation, or repeated infections and often needs surgical treatment.

\medterm{Fistula, Arteriovenous} Alternate word order; see \textit{Arteriovenous Fistula}.

\medterm{Fistula-In-Ano} An abnormal tunnel between the anal canal and skin near the anus. It commonly follows an anal abscess and may occur with Crohn disease, causing ongoing drainage, discomfort, or recurrent abscesses.

\medterm{Fit} A seizure or sudden episode of altered movement or awareness; this everyday term is also used to mean physically healthy. In medical writing, the specific symptom or diagnosis should be stated instead.

\medterm{Fits (Seizures)} Episodes caused by abnormal electrical activity in the brain. They may involve staring, unusual sensations, jerking, stiffening, loss of awareness, or unconsciousness; a first or prolonged seizure needs urgent assessment.

\synonyms
Seizures

\medterm{Fitz-Hugh-Curtis Syndrome} Inflammation and scarring around the liver caused by pelvic inflammatory disease, usually from gonorrhea or chlamydia. It can cause sharp right-upper-abdominal pain, fever, and pain that worsens with movement or breathing.

\medterm{Five-Year Survival Rate} The percentage of people alive five years after a diagnosis or treatment. It summarizes outcomes for groups, not an individual’s exact prognosis, which also depends on cancer type, stage, health, and treatment.

\medterm{Fixation} Holding a body part, tissue, or broken bone in the correct position while it heals. Methods include stitches, wires, screws, plates, rods, casts, splints, or external frames.

\medterm{Fixation, Internal} Surgery that stabilizes a broken or unstable bone with devices placed inside the body, such as plates, screws, pins, or rods. The aim is to restore alignment and allow healing and movement.

\medterm{Fixative} A chemical used to preserve cells or tissue for laboratory examination by preventing decay and maintaining structure. Common fixatives can irritate skin, eyes, and airways and must be handled safely.

\medterm{Fixed Joint} A joint that permits little or no movement because the bones are joined by bone or dense connective tissue. Some fixed joints are normal; others result from injury, disease, or surgical fusion.
\synonyms
Synarthrosis

\medterm{Fixed Macrophage} An immune cell that stays within a particular tissue instead of circulating in blood. It removes dead cells and microbes and helps coordinate inflammation and repair.

\medterm{Flaccid} Soft and lacking normal firmness or muscle tone. It may describe a limb, muscle, bladder, or penis; sudden flaccid weakness, especially on one side, can signal a neurologic emergency.

\medterm{Flagellin} The main protein in a bacterial flagellum, the whip-like structure that helps some bacteria move. Human immune cells can recognize flagellin and respond with inflammation.

\medterm{Flagellum} A long, whip-like structure that helps a cell or microorganism move. Human sperm and some bacteria have flagella, which differ in structure and function.
\medterm{Flail Chest} A serious chest injury in which several ribs break in more than one place, leaving part of the chest wall unstable. Severe pain and difficult breathing can result, often with lung bruising and need for urgent care.

\medterm{Flakka} A street name for illicit synthetic stimulant drugs, especially alpha-PVP. Overdose can cause extreme agitation, hallucinations, seizures, dangerously high temperature, muscle breakdown, irregular heartbeat, or death and requires emergency treatment.

\medterm{Flanged Needle} A needle with a widened rim or attached flange that helps stabilize it, limit how far it enters, or connect it to another device. The flange may reduce movement and tissue injury during sampling, injection, drainage, or infusion.

\medterm{Flank} The side of the body between the lower ribs and the hip, including part of the abdomen and back. Pain in this area may arise from the kidney, muscles, spine, bowel, or nearby organs.

\medterm{Flank Pain} Pain between the lower ribs and hip, often caused by kidney stones, kidney infection, muscle strain, spine problems, or nearby abdominal disease. Fever, vomiting, blood in urine, injury, or severe pain needs prompt evaluation.

\medterm{Flap Operation} Surgery that moves a piece of skin, tissue, or muscle, while preserving its blood supply, to cover a wound or reconstruct a damaged body part. Risks include infection, bleeding, and loss of the flap.

\medterm{Flare} A temporary increase in symptoms or inflammation after a disease has been stable or less active. Flares may be triggered by infection, stress, injury, or treatment changes; severe or unusual symptoms need assessment.

\synonyms
Exacerbation

\medterm{Flare-Up} A sudden or gradual worsening of a long-term disease after a period of relative stability. The cause may be infection, stress, missed treatment, or disease progression and should be assessed if severe or persistent.

\medterm{Flashes} Brief flickers or streaks of light seen without an outside light source. They may result from changes in the eye’s gel, migraine, or inflammation; sudden flashes with floaters or missing vision require urgent eye assessment.

\medterm{Flashing} Seeing brief flashes or streaks of light without an external source. New flashing, especially with floaters, a curtain-like shadow, eye pain, or reduced vision, can indicate a retinal emergency.

\medterm{Flat Affect} Very little outward expression of emotion, such as reduced facial movement, gestures, or voice variation. It is a clinical sign, not a diagnosis, and may occur with psychiatric, neurologic, or medication-related conditions.

\medterm{Flat Bones} Broad, relatively thin bones such as the skull, ribs, breastbone, and shoulder blades. They protect organs, provide surfaces for muscle attachment, and contain marrow that produces blood cells.

\medterm{Flat Feet} A lowered or absent arch along the inner side of the foot, causing more of the sole to touch the ground. It is often painless, but rigid or painful flat feet may reflect tendon, joint, or bone problems.

\medterm{Flatfoot} A foot with a reduced or collapsed inner arch. It may be flexible and painless or rigid and painful; causes include normal development, loose ligaments, tendon problems, injury, arthritis, or bone shape.

\medterm{Flatulence} Passage of gas through the anus, produced by swallowed air and bacterial digestion of food in the intestine. It may increase with diet, constipation, swallowing air, or digestive disorders; persistent troublesome symptoms need assessment.

\medterm{Flatus} Gas released through the anus after air is swallowed or intestinal bacteria break down food. Passing gas is normal; inability to pass gas with increasing abdominal swelling, vomiting, or pain may indicate bowel obstruction.

\synonyms
Intestinal gas

\medterm{Flavin Adenine Dinucleotide} A helper molecule made from vitamin B2 that allows enzymes to transfer electrons during energy production and other chemical reactions. It is commonly abbreviated FAD.

\synonyms
Flavin Adenine Dinucleotide

\medterm{Flaviviridae} A family of single-stranded RNA viruses that includes dengue, yellow fever, Zika, hepatitis C, and West Nile viruses. Different members spread through insects, blood, or other routes and cause different illnesses.

\medterm{Flavivirus} A group of RNA viruses that includes dengue, yellow fever, West Nile, Japanese encephalitis, and Zika viruses. Many spread through mosquitoes or ticks and can cause fever, inflammation, or nervous-system disease.

\medterm{Flaviviruses} Viruses in the flavivirus group, including dengue, Zika, yellow fever, and West Nile viruses. They are often spread by mosquitoes or ticks and may cause fever, bleeding, liver disease, birth defects, or brain inflammation.

\medterm{Flavobacteriaceae} A family of mainly gram-negative bacteria found in water, soil, and other environments. Most do not infect people, but some can cause opportunistic infections, especially in hospitalized or immunocompromised patients.

\medterm{Flavobacterium} A group of gram-negative bacteria commonly found in water and soil. A few species can cause opportunistic infections in people or animals, particularly when immunity is weakened.

\medterm{Flavoxate} A medicine used to relieve bladder spasms, urgency, frequency, or painful urination. It treats symptoms rather than infection or blockage and can cause dry mouth, blurred vision, dizziness, or constipation.

\medterm{Flaxseed Extract} A preparation made from flaxseed that may contain omega-3 fat, lignans, and soluble fiber. Effects vary by product; it may interact with medicines affecting clotting, blood sugar, or blood pressure.

\medterm{Flea} A small wingless insect that feeds on blood. Flea bites can cause itchy skin reactions, and fleas can transmit infections between animals or, less commonly, to people.

\medterm{Flea Bites} Itchy red bumps caused by fleas feeding on the skin, often clustered around the ankles or lower legs. Scratching can cause infection; treating pets and the environment helps prevent new bites.

\medterm{Flea Bites In Humans} Small itchy bumps, often around the ankles or lower legs, caused by flea feeding. Allergic reactions may produce larger areas of swelling; control requires treating infested animals, bedding, rooms, and symptoms.

\synonyms
Bites Flea (Flea Bites In Humans)
Ctenocephalides Flea Bite (Flea Bites In Humans)

\medterm{Fleas} Small wingless insects that feed on blood and infest animals, homes, or clothing. They can cause itchy bites, allergic skin disease, anemia in heavily infested animals, and transmission of some infections.

\medterm{Flecainide} A prescription antiarrhythmic medicine that slows electrical signals in the heart to prevent selected abnormal rhythms. It is not suitable for everyone and can worsen dangerous rhythms, especially with structural heart disease.

\medterm{Flesh-Eating Bug} An informal name for necrotizing fasciitis, a rapidly spreading infection of deep skin and soft tissue. Severe pain out of proportion to visible changes, swelling, fever, blisters, or skin discoloration requires emergency care.

\medterm{Flex} To bend a joint so the angle between connected bones becomes smaller. For example, bending the elbow brings the forearm toward the upper arm; pain, swelling, stiffness, or weakness can limit this movement.

\synonyms
Flexion

\medterm{Flexibility} The range of movement available at a joint or group of joints. It depends on bones, muscles, tendons, ligaments, joint health, nervous-system control, age, training, and pain.

\medterm{Flexible Nasolaryngoscopy} An examination using a thin flexible camera passed through the nose to view the nasal passages, throat, and voice box. It helps investigate hoarseness, swallowing difficulty, blockage, lesions, and vocal-cord movement.

\medterm{Flexible Sigmoidoscopy} An examination using a flexible camera to view the rectum and lower colon. It can investigate bleeding, diarrhea, or bowel symptoms and monitor known disease; preparation and biopsy may be needed.

\medterm{Flexion} Movement that bends a joint and decreases the angle between connected bones, such as bending the elbow, knee, hip, or fingers.

\medterm{Flexor} A muscle or tendon that bends a joint or body part. Flexors in the arm and hand bend the fingers and wrist; injury can cause pain, weakness, or loss of movement.

\medterm{Flexor Digitorum Profundus} A deep forearm muscle whose tendons bend the end joints of the fingers and also assist wrist movement. Its inner part is supplied by the ulnar nerve and its outer part by the median nerve.

\medterm{Flexor Digitorum Superficialis} A forearm muscle whose tendons bend the middle joints of the fingers and assist wrist movement. Damage to the muscle, tendon, or median nerve can cause pain and difficulty bending the fingers.

\medterm{Flexural} Relating to a skin fold or the inner surface of a bent joint, such as the elbow crease, back of the knee, armpit, groin, or fold beneath the breast.

\medterm{Flibanserin} A prescription medicine for persistent, distressing low sexual desire in some premenopausal women. It is taken regularly rather than as needed and can cause dizziness, sleepiness, low blood pressure, or fainting, especially with alcohol.

\medterm{Flight-Or-Fight Response} The body’s rapid response to perceived danger, releasing stress hormones that increase heart rate, breathing, alertness, and muscle readiness so a person can escape or confront a threat.
\synonyms
Fight-or-flight response; stress response

\medterm{Floater} A spot, thread, or cobweb-like shape that drifts across vision, usually caused by changes in the gel inside the eye. Sudden new floaters, flashes, or missing vision require urgent eye assessment.

\medterm{Floaters} Moving spots, threads, cobwebs, or rings in the field of vision caused by shadows from material in the eye’s vitreous gel. They are often age-related, but sudden onset can indicate bleeding or a retinal tear.

\medterm{Floating Kidney} Unusual downward movement of a kidney when a person stands, sometimes called nephroptosis. It is often symptomless but may cause position-related flank pain, urinary symptoms, or rarely impaired urine drainage.

\medterm{Floating Rib} One of the eleventh or twelfth ribs, attached to the spine but not connected to the breastbone. These ribs are less protected and can be injured by direct trauma.

\medterm{Floating Ribs} The eleventh and twelfth pairs of ribs, which attach to the spine but not the breastbone. They are relatively mobile and less protected; fractures can cause pain and may injure nearby organs.

\medterm{Flocculation} The joining of tiny particles suspended in a liquid into larger clumps. Laboratory tests may use this visible reaction to detect antibodies or other substances; it is also used in water treatment and pharmaceutical preparation.

\medterm{Floor Of The Mouth Cancer} Cancer arising in the tissues beneath the tongue, most often from surface cells. A persistent ulcer, lump, pain, bleeding, numbness, or difficulty swallowing or speaking needs prompt examination.

\synonyms
Cancer, Floor Of The Mouth

\medterm{Floppy Baby Syndrome} An informal description of an infant with unusually low muscle tone and a limp posture. Causes include brain, spinal cord, nerve, muscle, genetic, or metabolic disorders and require medical evaluation.

\medterm{Floppy Iris Syndrome} A cataract-surgery complication in which the iris becomes unusually loose, moves with fluid currents, and narrows the working space. It is associated with some prostate medicines and should be anticipated by the surgeon.

\medterm{Floppy Valve Syndrome} A condition in which the mitral valve’s leaflets bulge backward into the upper heart chamber when the heart contracts. It is often harmless but may cause palpitations, chest discomfort, or leakage through the valve.

\synonyms
Mitral valve prolapse

\medterm{Flora} The microorganisms living normally or temporarily in a body site or environment. Some protect against harmful organisms, but they can cause infection if they enter another body site or defenses are weakened.


\medterm{Flossing} Cleaning between teeth with dental floss or another interdental device to remove plaque and food that a toothbrush cannot reach. Gentle daily cleaning can help prevent gum disease and cavities.

\medterm{Flotation} A laboratory method that separates material by density. For example, a flotation test can concentrate parasite eggs from stool so they are easier to identify under a microscope.

\medterm{Flow Cytometry} A laboratory test that examines cells or particles one at a time as they pass through a laser beam. It can count cells and identify their features, helping diagnose blood cancers and immune disorders.

\medterm{Flow-Volume Loop} A graph showing how much air moves in and out of the lungs at different speeds during forceful breathing. Its shape can help identify airway narrowing, reduced lung size, or an upper-airway blockage.

\medterm{Flowmeter} A device that measures the rate at which a gas or liquid moves. In medicine, it may measure oxygen delivery, intravenous fluid flow, or breathing flow such as peak expiratory flow.

\medterm{Floxacillin} An antibacterial medicine used in some countries to treat infections caused by susceptible bacteria, especially staphylococci. It can cause allergic reactions, diarrhea, liver problems, or other antibiotic-related effects.

\medterm{Floxuridine} An anticancer medicine that interferes with DNA production, slowing or killing rapidly dividing cells. It is used only in selected treatment plans and may cause low blood counts, mouth sores, diarrhea, or other toxicity.

\medterm{Flu} An acute respiratory infection caused by influenza viruses. It commonly causes sudden fever, cough, sore throat, headache, muscle aches, and fatigue and can lead to pneumonia or worsen chronic diseases.


\synonyms
Influenza; influenza infection

\medterm{Flu, Swine} Influenza A infection mainly affecting pigs. Some strains can infect people, causing an illness like seasonal flu; severity varies, and human-to-human spread can occur when a strain adapts to spread efficiently.

\synonyms
H1N1 flu; swine influenza

\medterm{Flucloxacillin} An antibiotic used for selected bacterial infections, especially those caused by susceptible staphylococci. It can cause nausea, diarrhea, allergy, or liver injury and should be taken only as prescribed.

\medterm{Fluconazole} An antifungal medicine used for selected yeast and other fungal infections. It can interact with many medicines and may cause nausea, rash, or liver injury; treatment depends on the infection and patient.

\medterm{Flucytosine} An antifungal medicine converted inside susceptible fungi into substances that interfere with genetic-material production. It is used for selected serious infections, usually with another antifungal, and requires blood-count and liver monitoring.

\medterm{Fludarabine} A chemotherapy medicine that interferes with DNA production in abnormal blood cells. It treats selected blood cancers but can severely suppress immunity and blood counts, increasing the risk of serious infections.

\medterm{Fludarabine Phosphate} A form of fludarabine used in selected blood-cancer treatment plans. It interferes with DNA production and can cause prolonged low blood counts, immune suppression, infections, and neurologic or other serious effects.

\medterm{Fludeoxyglucose F-18} A radioactive form of glucose used in positron-emission tomography scans. Active tissues absorb more of it, helping doctors assess some cancers, brain disorders, and heart muscle; uptake is not specific for cancer.

\medterm{Fludrocortisone} A medicine that helps the body retain salt and water and raise blood pressure. It treats some adrenal-hormone deficiencies and selected orthostatic disorders; monitoring is needed for swelling, high blood pressure, and low potassium.

\medterm{Fluid Around The Heart} Extra fluid in the sac surrounding the heart. Small amounts may cause no symptoms, while rapid or large collections can restrict heart filling and cause breathlessness, low blood pressure, or life-threatening tamponade.

\synonyms
Pericardial effusion

\medterm{Fluid Balance} The relationship between water entering and leaving the body. The kidneys, hormones, thirst, breathing, sweat, and stool help regulate it; disruption can cause dehydration, swelling, or abnormal blood salts.

\medterm{Fluid Imbalance} Too much or too little body water or an abnormal level of salts such as sodium or potassium. Causes include vomiting, diarrhea, kidney or heart disease, hormone disorders, dehydration, or excess fluid treatment.

\medterm{Fluid Management} Planning and monitoring fluids, salts, and urine output to maintain circulation and organ function. Treatment may involve drinking, intravenous fluids, diuretics, or fluid restriction according to the underlying problem.

\medterm{Fluid Overload} Excess water and salt in the body, causing swelling, rapid weight gain, high blood pressure, or fluid in the lungs. It may result from heart, kidney, or liver disease or too much fluid treatment.

\medterm{Fluid Therapy} Treatment using fluids, often through a vein, to restore hydration, blood volume, or salt balance. The type and amount depend on the cause; excess fluid can worsen swelling, heart failure, or lung congestion.

\synonyms
Intravenous fluid therapy

\medterm{Fluke} A flat, leaf-shaped parasitic worm. Different species infect the liver, lungs, intestines, or blood vessels and are acquired through contaminated food, water, or contact with freshwater.

\medterm{Flumazenil} An antidote that can reverse sedation and slowed breathing caused by some benzodiazepine medicines. It can trigger seizures or sudden withdrawal and is not suitable for every overdose, especially mixed overdoses.

\medterm{Flunarizine} A medicine used in some countries to prevent migraine and treat selected balance disorders. It can cause sleepiness, weight gain, depression, or movement problems and should be prescribed with appropriate monitoring.

\medterm{Flunisolide} An inhaled or nasal corticosteroid used to reduce inflammation in asthma or allergic rhinitis. It may cause throat irritation, nosebleeds, or oral thrush; regular use does not relieve a sudden severe asthma attack.

\medterm{Fluocinolone} A corticosteroid applied to skin, scalp, or other affected areas to reduce inflammation and itching. Prolonged or extensive use can thin the skin or rarely cause effects throughout the body.

\medterm{Fluocinolone Acetonide} A corticosteroid medicine applied to the skin or scalp to reduce redness, swelling, and itching. Overuse can thin the skin, worsen infections, or rarely affect the body’s hormone system.

\medterm{Fluocinonide} A strong corticosteroid applied to selected inflammatory skin conditions to reduce itching, redness, and swelling. It should be used only as directed because prolonged or extensive use can damage skin or cause systemic effects.

\medterm{Fluocortolone} A corticosteroid used in selected skin or other inflammatory conditions. It reduces inflammation but can cause skin thinning, infection, or other steroid effects, especially when used for a long time or over large areas.

\medterm{Fluorescein} A bright fluorescent dye used in eye examinations and some imaging tests. In the eye, it highlights scratches, ulcers, and abnormal blood-vessel leakage; rare allergic reactions can occur after injection.

\medterm{Fluorescein Angiography} An eye test in which fluorescent dye is injected into a vein and photographs track its movement through retinal blood vessels. It can show leakage, blockage, or abnormal vessels in retinal disease.

\medterm{Fluorescein Eye Stain} A test in which dye is placed on the eye and viewed under blue light. It highlights scratches, ulcers, foreign-body injuries, and some dry-eye surface damage.

\medterm{Fluorescence} The emission of visible light by a substance after it absorbs light or other energy. Clinicians and researchers use it to identify cells, chemicals, microbes, or structures in laboratory and imaging tests.

\medterm{Fluorescence In Situ Hybridization} A laboratory test that uses glowing DNA probes to find specific chromosome regions or gene changes in cells. It can detect selected missing, extra, rearranged, or amplified genetic material.

\synonyms
Fluorescence In Situ Hybridization

\medterm{Fluorescence In-Situ Hybridization} A targeted genetic test using fluorescent DNA probes to locate selected chromosome sequences. It can identify some extra or missing chromosome material and gene rearrangements but does not examine the entire genome.

\medterm{Fluorescence Microscopy} A microscope method that uses fluorescent labels to make selected cells, molecules, or microbes visible. It helps examine tissue structure, cell proteins, genetic material, and infectious organisms.


\medterm{Fluorescent Dye} A substance that absorbs light of one color and emits light of another. It can label cells, tissues, molecules, or microbes so they can be detected during laboratory or imaging examinations.

\medterm{Fluoridation} Adding a controlled amount of fluoride to water or dental products to strengthen tooth enamel and reduce cavities. Excessive fluoride during tooth development can cause dental fluorosis, which appears as enamel discoloration or pits.

\medterm{Fluoride} A mineral that strengthens tooth enamel, helps repair early damage, and reduces acid production by cavity-causing bacteria. It is found in some water supplies, toothpaste, mouth rinses, and professional dental treatments.

\medterm{Fluoride In Diet} Fluoride obtained from drinking water and foods or drinks made with fluoridated water. The right amount helps protect teeth, while excessive intake, especially in children, can damage developing enamel.

\medterm{Fluoride Overdose} Excessive fluoride swallowed over a short time. It can cause nausea, vomiting, abdominal pain, low blood calcium, muscle spasms, abnormal heart rhythms, or seizures and requires urgent poison-control or medical advice.

\medterm{Fluoride Poisoning} Illness caused by swallowing too much fluoride. Early symptoms include nausea, vomiting, and abdominal pain; severe poisoning can lower blood calcium and cause muscle spasms, seizures, or dangerous heart rhythms.

\medterm{Fluoride Varnish} A concentrated fluoride coating painted onto teeth by a dental professional to prevent or slow cavities. It hardens quickly and is used at intervals based on age, tooth condition, and cavity risk.

\medterm{Fluorimetry} A laboratory method that measures light emitted by a fluorescent substance. It can detect or estimate very small amounts of chemicals, medicines, proteins, or biological material.

\medterm{Fluorine} A highly reactive chemical element found in compounds such as fluoride. Small controlled amounts of fluoride help protect teeth, while excessive exposure can poison the body or damage developing bones and teeth.

\medterm{Fluorine-18} A radioactive form of fluorine used to label tracers for positron-emission tomography scans. Its short half-life allows useful imaging while limiting how long radioactivity remains in the body.

\medterm{Fluorometholone} A corticosteroid eye medicine used for selected inflammatory conditions. Prolonged use can raise eye pressure, delay healing, worsen infection, or increase cataract risk, so it requires medical supervision.

\medterm{Fluorophotometry} A method that measures fluorescent light from the eye or another sample. It can help estimate substances or tissue properties, but results must be interpreted with the test method and clinical findings.

\medterm{Fluoropyrimidine} A class of anticancer medicines that resembles substances used to build DNA and RNA. They disrupt genetic-material production in rapidly dividing cells and can cause mouth sores, diarrhea, low blood counts, and hand-foot symptoms.

\medterm{Fluoroquinolone} An antibiotic that kills bacteria by blocking enzymes needed for bacterial DNA copying. It is reserved for selected infections because it can cause tendon injury, nerve symptoms, heart-rhythm changes, blood-sugar changes, or severe diarrhea.

\medterm{Fluoroquinolones} A group of antibiotics that kill bacteria by blocking enzymes required for bacterial DNA replication. They treat selected infections but are avoided when safer options work because of potentially serious tendon, nerve, heart, and blood-sugar effects.

\medterm{Fluoroscopy} An X-ray method that produces moving images inside the body. It helps guide procedures and examine swallowing, blood vessels, joints, and digestive or urinary passages; contrast dye may be used and radiation exposure occurs.

\medterm{Fluorosis} A change in developing tooth enamel caused by too much fluoride during childhood. It usually appears as white streaks or spots; more severe cases cause brown discoloration or pits.

\medterm{Fluorouracil} An anticancer medicine that interferes with DNA production. It is applied to selected precancerous skin lesions or given for some cancers; possible effects include mouth sores, diarrhea, low blood counts, and serious toxicity in susceptible people.

\medterm{Fluorouracil Side Effect} An unwanted effect of fluorouracil, such as mouth sores, diarrhea, nausea, low blood counts, hand-foot skin reactions, or heart problems. Severe symptoms require prompt medical review and treatment adjustment.

\medterm{Fluoxetine} An antidepressant that increases available serotonin in the brain. It treats depression and several other conditions, but may cause nausea, sleep or sexual problems, agitation, or rarely dangerous serotonin excess or suicidal thoughts.

\medterm{Fluoxymesterone} A synthetic androgen medicine used in selected hormone-related conditions. It can affect the liver, cholesterol, mood, fertility, prostate, and blood count and should be used only with medical supervision.

\medterm{Flupentixol} An antipsychotic medicine used in some countries for schizophrenia and other psychiatric conditions. It can reduce hallucinations and delusions but may cause movement problems, sleepiness, hormonal changes, or rare serious reactions.

\medterm{Fluphenazine} An antipsychotic medicine used to treat schizophrenia and related psychotic symptoms. It can cause stiffness, tremor, restlessness, sleepiness, hormonal changes, and rare but serious involuntary movements or fever.

\medterm{Fluphenazine Hydrochloride} The hydrochloride salt form of fluphenazine, an antipsychotic medicine for selected psychotic disorders. Monitoring is needed because it can cause movement disorders, sedation, hormonal effects, and other serious adverse reactions.

\medterm{Flurandrenolide} A corticosteroid applied to the skin to reduce inflammation, redness, and itching. Prolonged, extensive, or covered use can thin the skin, worsen infection, or allow steroid effects throughout the body.

\medterm{Flurazepam} A long-acting sedative medicine used for short-term treatment of insomnia in some settings. It can cause next-day drowsiness, falls, memory problems, dependence, and dangerous breathing suppression with alcohol or opioids.

\medterm{Flurbiprofen} A nonsteroidal anti-inflammatory medicine used for pain and swelling and, in some forms, eye inflammation. It can cause stomach bleeding, kidney problems, heart risks, or allergic reactions and should be used as directed.

\medterm{Flush} A sudden feeling of warmth or visible reddening caused by increased skin blood flow. In healthcare, flush can also mean washing a tube, body passage, or intravenous line with fluid.

\medterm{Flushable Reagent Stool Blood Test} Alternate index label; see \textit{Fecal occult blood test}.

\medterm{Flushing} Temporary redness and warmth of the skin caused by increased blood flow. Triggers include heat, exercise, emotion, alcohol, medicines, fever, or hormone disorders; flushing with wheezing, diarrhea, faintness, or palpitations needs evaluation.

\medterm{Flutamide} A medicine that blocks male-hormone effects and is used with other hormone treatment for prostate cancer. It can seriously injure the liver and may cause nausea, breast tenderness, or reduced blood counts.

\medterm{Flutter} A rapid, regular, repetitive movement or contraction. Atrial flutter is an abnormal heart rhythm caused by a fast electrical circuit in the upper chambers and can increase the risk of stroke.

\medterm{Flutter Valve} A handheld breathing device that creates gentle pressure and vibrations during exhalation. It can help loosen and move mucus in some people with bronchiectasis, cystic fibrosis, or chronic lung disease.

\medterm{Fluvastatin} A statin medicine that lowers low-density lipoprotein cholesterol by reducing cholesterol production in the liver. It lowers cardiovascular risk but can cause muscle symptoms, liver abnormalities, or rare muscle breakdown.

\medterm{Fluvoxamine} An antidepressant that increases serotonin signaling and is used mainly for obsessive-compulsive disorder and selected other conditions. It can cause nausea, sleep changes, sexual problems, medicine interactions, or rarely serotonin toxicity.

\medterm{Flux} An old, nonspecific term for diarrhea or an excessive discharge, especially from the bowel. Modern medical descriptions should identify the actual symptom, duration, severity, and suspected cause.

\medterm{FM} An abbreviation sometimes used for fibromyalgia, a long-term condition causing widespread pain, fatigue, poor sleep, and concentration problems. Abbreviations can be unclear, so the full term is preferred in patient records.

\synonyms
Fibromyalgia

\medterm{FMRI} Functional magnetic resonance imaging maps changes in brain blood flow and oxygen use while a person performs tasks or rests. It can help identify language or movement areas and does not use X-ray radiation.

\medterm{FMS} An abbreviation sometimes used for fibromyalgia syndrome, a long-term condition with widespread pain, fatigue, poor sleep, and cognitive symptoms. The full term is preferred because FMS can have other meanings.

\synonyms
Fibromyalgia

\medterm{FNA} An abbreviation for fine-needle aspiration, a procedure using a thin needle to collect cells or fluid from a lump or organ for laboratory examination. The full term is clearer in patient information.

\medterm{FND} Functional neurological disorder causes real weakness, abnormal movements, seizure-like episodes, sensory changes, or other nervous-system symptoms because nerve signaling is disrupted. Symptoms are not imagined; diagnosis is based on characteristic examination findings and may require tests to exclude other causes.

\synonyms
Conversion disorder; functional neurologic disorder

\medterm{Foam Cell} A cell containing many fat droplets, giving it a foamy appearance under a microscope. Foam cells, often immune cells that have taken up cholesterol, are common in early artery-wall plaque.

\medterm{Foam Cells} Cells filled with fat droplets, often immune cells that have taken up altered cholesterol. They commonly form early in artery-wall plaque and contribute to inflammation and atherosclerosis.

\medterm{Focal} Limited to one specific area rather than spread widely. A focal finding may involve an image, seizure, infection, injury, or neurologic problem and must be interpreted according to its location.

\medterm{Focal Adhesion} A specialized attachment site where a cell connects to its surrounding support material. It helps cells stick, sense physical forces, move, and send signals during growth and tissue repair.

\medterm{Focal Glomerulosclerosis} Scarring affecting some kidney filters or only parts of affected filters. It can cause protein to leak into urine, swelling, high blood pressure, and gradual loss of kidney function.
\medterm{Focal Neurologic Deficits} Loss of a particular function, such as strength, sensation, speech, vision, or coordination, because of a localized problem in the brain, spinal cord, or a nerve. Sudden onset may indicate stroke.

\medterm{Focal Neuropathy} Damage or irritation affecting one peripheral nerve or a small group of nerves. It may cause pain, tingling, numbness, weakness, or altered movement in the area supplied by those nerves.

\synonyms
Mononeuropathy

\medterm{Focal Nodular Hyperplasia} A usually harmless overgrowth of liver tissue related to an unusual local blood supply. It often causes no symptoms, is found on imaging, and rarely needs treatment unless its diagnosis is uncertain.

\medterm{Focal Sclerosis} Hardening or scarring limited to a particular area of tissue. In the kidney, it may describe scarring of some filters or parts of filters, which can impair filtration and cause protein in urine.

\medterm{Focal Segmental Glomerulosclerosis} A kidney disorder in which scars affect some filters or parts of them. It can cause heavy urine protein, swelling, high blood pressure, and kidney failure; causes include genetic changes, infections, medicines, and excess filtration.

\medterm{Focal Seizure} A seizure that begins in one area of the brain. It may cause unusual movement, sensation, behavior, emotion, or awareness and can spread to involve both sides of the brain.

\medterm{Focal-Onset Seizure} A seizure that begins in one area of the brain. Awareness may remain normal or become impaired, and symptoms depend on the affected area; the seizure may spread and cause whole-body stiffening and jerking.

\medterm{Focus Words} Words or short phrases used during relaxation or stress-management exercises to direct attention and support calm breathing. They may help coping but are not a treatment for an underlying medical disorder.

\medterm{FODMAP Diet} A short-term eating plan that reduces certain poorly absorbed carbohydrates that can cause gas, bloating, pain, or diarrhea. Foods are gradually reintroduced with professional guidance to identify personal triggers.

\medterm{FODMAPs} A group of short-chain carbohydrates that may be poorly absorbed in the small intestine and fermented by gut bacteria. In susceptible people they can cause gas, bloating, abdominal pain, or diarrhea.

\medterm{Foetal Acidosis} Excess acidity in fetal blood, usually caused by inadequate oxygen delivery or impaired gas exchange. Fetal heart-rate monitoring and, in selected situations, blood testing help assess its likelihood and severity.

\medterm{Foetal Alcohol Syndrome} A lifelong condition caused by alcohol exposure before birth, with characteristic facial features, poor growth, and brain-development problems affecting learning, behavior, or coordination. No amount of alcohol is known to be safe during pregnancy.

\medterm{Foetus} The developing human offspring from the end of the embryonic period, about nine weeks after fertilization, until birth. During this period, organs grow and mature and body functions become more coordinated.

\medterm{Folate} Vitamin B9, needed for DNA production, cell division, and healthy red blood cells. Adequate folate before and during early pregnancy lowers neural-tube-defect risk; high-dose supplements can hide vitamin B12 deficiency.

\medterm{Folate Antagonist} A medicine that blocks folate-related processes needed to make DNA. Depending on the drug, it may treat cancer, inflammatory disease, or bacterial infection, but can cause low blood counts, mouth sores, or liver problems.

\medterm{Folate Deficiency} Too little vitamin B9, caused by poor intake, increased need, malabsorption, alcohol use, or some medicines. It can cause large abnormal red blood cells, anemia, fatigue, sore tongue, and pregnancy complications.

\medterm{Folate Deficiency Anemia} Anemia caused by too little folate, or vitamin B9. Red blood cells become unusually large and immature, causing fatigue, weakness, pale skin, shortness of breath, or a sore tongue.

\synonyms
Vitamin B9 deficiency anemia

\medterm{Folate-Deficiency Anemia} Anemia caused by inadequate folate, poor absorption, increased requirements, alcohol use, or certain medicines. Blood tests help confirm the cause, and folate replacement treats the deficiency after vitamin B12 deficiency is considered.

\medterm{Foley Catheter} A flexible tube placed through the urethra into the bladder to drain urine. A small balloon holds it in place; risks include discomfort, blockage, urinary infection, and urethral injury.

\medterm{Folic Acid} A manufactured form of vitamin B9 needed for DNA production and healthy blood cells. Supplements prevent many fetal neural-tube defects when taken before conception and treat folate deficiency.

\medterm{Folic Acid - Test} A blood test measuring folate levels to investigate suspected deficiency or anemia with unusually large red blood cells. Results are interpreted with blood counts, vitamin B12 testing, diet, medicines, and symptoms.

\medterm{Folic Acid And Birth Defect Prevention} Taking folic acid before conception and during early pregnancy reduces the risk of neural-tube defects, which affect development of the fetal brain and spinal cord. Higher-risk people may need a prescribed dose.
\medterm{Folic Acid In Diet} Folate occurs naturally in leafy vegetables, beans, citrus fruits, and liver, while folic acid is added to some fortified foods. Intake supports blood formation and fetal development; needs increase during pregnancy.
\medterm{Folic Acid Supplementation} Taking vitamin B9 before conception and during early pregnancy to reduce neural-tube-defect risk. Dose depends on individual risk; supplementation also treats deficiency and supports DNA production and red-cell formation.

\medterm{Folinic Acid} An active form of folate used to treat or prevent folate-related toxicity and to support selected chemotherapy treatments. It can also treat certain folate deficiencies; its use depends on the clinical situation.

\medterm{Follicle} A small sac or organized tissue structure. A hair follicle produces a hair, while an ovarian follicle contains a developing egg; other follicles can occur in glands or lymphatic tissue.

\medterm{Follicle Stimulating Hormone} A hormone released by the pituitary gland. In females it helps ovarian follicles grow and produce estrogen; in males it supports sperm production. Levels vary with age, sex, and reproductive stage.

\medterm{Follicle-Stimulating Hormone} A pituitary hormone that supports egg-containing follicle development and estrogen production in females and sperm production in males.

\medterm{Follicle-Stimulating Hormone Blood Test} A blood test measuring a hormone involved in ovarian function and sperm production. It helps evaluate infertility, menstrual changes, puberty problems, menopause, testicular or ovarian failure, and pituitary disorders.

\medterm{Follicular} Relating to a follicle, especially a hair follicle or an ovarian follicle. A follicular skin problem is centered around a hair and may appear as a bump, pustule, plug, or inflamed area.

\medterm{Follicular Cyst} A usually harmless ovarian cyst that forms when a follicle does not release an egg or does not shrink normally. It often causes no symptoms and disappears within a few menstrual cycles without treatment.

\medterm{Follicular Fluid} The liquid inside an ovarian follicle surrounding a developing egg. It contains hormones and other substances that support follicle growth and help coordinate ovulation.

\medterm{Follicular Hyperkeratosis} Excess keratin blocking hair follicles and causing small rough bumps. It commonly occurs with keratosis pilaris and, less often, with nutritional deficiency or other skin conditions.

\medterm{Follicular Lymphoma} A usually slow-growing cancer of immune cells that often causes painless swelling of lymph nodes. It may involve bone marrow or other organs and can sometimes change into a faster-growing lymphoma.

\medterm{Follicular Thyroid Carcinoma} A thyroid cancer that begins in follicle-forming cells and can spread through the bloodstream, especially to bone or lungs. Tissue examination is needed to show invasion and distinguish it from a benign growth.

\medterm{Follicular Unit} A natural group of one or more hair follicles along with nearby skin glands, nerves, and tiny muscles. Hair-transplant procedures may move these units from a fuller area to a thinning area.

\medterm{Follicular Unit Transplantation} Hair-restoration surgery that moves natural groups of hair follicles from a donor area to thinning or bald areas. Results take months, and possible problems include scarring, infection, and continued hair loss.

\synonyms
FUT

\medterm{Folliculitis} Inflammation or infection around a hair follicle, causing small itchy, red, or tender bumps and sometimes pus. Causes include bacteria, yeast, friction, blocked skin, shaving, and ingrown hairs.

\synonyms
Barber's Itch


\medterm{Follow-Up Recommendations} A plan for repeat examination, testing, imaging, treatment, or specialist review after an initial finding. The plan depends on the diagnosis, symptoms, test results, personal risks, and accepted clinical guidance.

\medterm{Fomentation} Applying a warm, moist cloth or material to part of the body to soothe pain, relax muscles, or improve comfort. Excessive heat can burn skin and should be avoided.

\medterm{Fomepizole} An antidote for poisoning with methanol or ethylene glycol. It blocks formation of their toxic breakdown products and may be combined with intensive supportive care and dialysis.

\medterm{Fomite} An object or surface contaminated with an infectious organism that can transfer it to another person. Examples include towels, equipment, toys, or needles, although the importance of surface spread varies by organism.

\medterm{Fomites} Objects or surfaces that may carry infectious organisms and transfer them through contact. Examples include shared towels, medical equipment, toys, and other contaminated items; cleaning and hand hygiene reduce risk.

\medterm{Fondaparinux} An injectable anticoagulant that prevents formation or extension of blood clots by blocking activated factor X. It is used for selected clotting disorders and prevention after surgery; bleeding is its main risk.

\medterm{Fondaparinux Sodium} The sodium salt form of fondaparinux, an injectable medicine that reduces clotting by blocking activated factor X. It treats or prevents selected blood clots but can cause serious bleeding and requires dose adjustment in kidney disease.

\medterm{Fonofos} A highly poisonous organophosphate insecticide. Exposure can cause sweating, excess saliva, vomiting, small pupils, muscle twitching, breathing difficulty, seizures, or weakness and requires urgent poison-control or emergency care.

\medterm{Fonsecaea} A group of darkly pigmented fungi that can cause slow-growing skin and soft-tissue infections called chromoblastomycosis. Infection usually follows injury and may produce warty plaques or nodules, especially in tropical regions.

\medterm{Fontan Procedure} An operation for some children born with only one effective pumping chamber. It directs blood returning from the body to the lungs without a separate pumping chamber, improving circulation but requiring lifelong follow-up for rhythm, heart, liver, clotting, and protein-loss problems.

\synonyms
Fontan Operation; Fontan Circulation; Total Cavopulmonary Connection (TCPC)

\medterm{Fontanel} A soft gap between an infant’s skull bones that allows the brain and skull to grow and helps the skull mold during birth. Its size and firmness can provide clues about hydration or pressure inside the skull.

\medterm{Fontanel (Fontanelle)} A soft membranous space between an infant’s skull bones. It normally changes as the skull grows; a persistently bulging or markedly sunken soft spot can signal illness and needs assessment.

\medterm{Fontanelle} A soft membranous space between an infant’s skull bones. Fontanelles allow growth and molding during birth; the front one usually closes during the second year, while the back one closes earlier.

\medterm{Fontanelles - Bulging} A tense soft spot that remains outwardly bulging when an infant is calm and upright. It may indicate increased pressure inside the skull, infection, bleeding, or another serious problem requiring urgent assessment.

\medterm{Fontanelles - Enlarged} Soft spots on an infant’s skull that are larger than expected for age. Causes include normal variation, prematurity, thyroid disease, vitamin D deficiency, bone disorders, or increased pressure inside the skull.

\medterm{Fontanelles - Sunken} A noticeably depressed soft spot in an infant who is calm and upright. It commonly suggests dehydration, especially with poor feeding, fewer wet diapers, vomiting, or diarrhea, and needs medical assessment.

\medterm{Food Additives} Substances added to food to preserve it or change its flavor, color, texture, or stability. Most are safe in approved amounts, but some people may have intolerance or allergy to particular additives.

\medterm{Food Allergies} Immune reactions to particular food proteins. They may cause hives, swelling, vomiting, wheezing, or life-threatening anaphylaxis; strict avoidance and an emergency plan are important for people with severe reactions.
\synonyms
Food allergy

\medterm{Food Allergy} An immune reaction to a food protein. It may cause hives, swelling, vomiting, abdominal symptoms, wheezing, or anaphylaxis. Common triggers include milk, egg, peanut, tree nuts, fish, shellfish, soy, and wheat. Severe reactions require an emergency plan.

\synonyms
Allergy Food (Food Allergy)
Allergy, Food

\medterm{Food Allergy, Egg} An immune reaction to egg proteins, most common in children. It can cause hives, vomiting, breathing problems, or anaphylaxis; evaluation guides avoidance, emergency planning, and whether the allergy has resolved.

\synonyms
Egg allergy

\medterm{Food Allergy, Milk} An immune reaction to milk proteins that can cause eczema, hives, vomiting, diarrhea, wheezing, or anaphylaxis. It is common in young children, and many outgrow it; safe nutritional substitutes are important.

\synonyms
Milk allergy

\medterm{Food Consumption} The amount, type, and pattern of food and drink a person takes in. It is influenced by appetite, culture, access, habits, health, activity, medicines, and the body’s nutritional needs.

\medterm{Food Hygiene} Practices that prevent food contamination and foodborne illness, including washing hands, separating raw and cooked foods, cooking thoroughly, chilling promptly, and using safe water.

\medterm{Food Intolerance} A troublesome reaction to a food that does not involve the immune system. Lactose intolerance is common; symptoms may include gas, bloating, cramps, or diarrhea and depend on the food and amount.

\medterm{Food Jags} A child’s repeated preference for only a few foods or refusal of foods previously accepted. It is common during development but needs assessment if it limits nutrition, growth, social eating, or causes distress.

\medterm{Food Labeling} Information on packaged food showing ingredients, allergens, serving size, nutrients, calories, dates, and health claims. Reading labels helps people manage allergies, medical diets, and overall nutrition.

\medterm{Food Poisoning} Illness from eating food contaminated with germs, toxins, or harmful chemicals. Nausea, vomiting, diarrhea, abdominal cramps, and fever are common.

\medterm{Food Poisoning (Foodborne Illness)} Illness caused by eating or drinking material contaminated with germs, toxins, or harmful chemicals. Symptoms commonly affect the digestive system, but some infections can also damage the nervous system or other organs.

\medterm{Food Poisoning Prevention} Reduce risk by washing hands, separating raw foods, cooking thoroughly, refrigerating promptly, avoiding unsafe water, and following food recalls. High-risk people should avoid foods commonly linked to serious infection.

\medterm{Food Poisoning, Campylobacter} A foodborne infection caused by Campylobacter bacteria, often from undercooked poultry, unpasteurized milk, or contaminated water. It causes diarrhea, cramps, and fever and may rarely trigger bloodstream infection or nerve inflammation.

\medterm{Food Safety} Handling, storing, preparing, and labeling food in ways that prevent germs, toxins, harmful chemicals, and allergens from causing illness.

\medterm{Food Supplement} A product taken to add nutrients or other substances to the diet, such as vitamins, minerals, herbs, or protein. Supplements can interact with medicines and are not automatically safe or effective.

\medterm{Food-Borne Disease} Illness acquired from food or water contaminated with infectious organisms, toxins, or harmful chemicals. It may cause digestive symptoms or affect other organs depending on the contaminant.

\medterm{Food-Borne Illness} Illness caused by eating food contaminated with bacteria, viruses, parasites, toxins, or chemicals. Symptoms commonly include nausea, vomiting, diarrhea, cramps, and fever; safe handling and cooking help prevent it.

\synonyms
Food poisoning; foodborne illness

\medterm{Foodborne Disease} Illness acquired by eating or drinking material contaminated with germs, toxins, or harmful chemicals. Severity ranges from brief diarrhea to dehydration, organ damage, or life-threatening infection.

\medterm{Foodborne Illness} Illness resulting from contaminated food or water, including infections and poisoning from germs, toxins, or chemicals. Serious disease is more likely in pregnancy, infancy, older age, or weakened immunity.


\medterm{Foot} The part of the lower limb below the ankle, made of bones, joints, muscles, tendons, ligaments, nerves, and blood vessels. Its shape and arches support standing, balance, shock absorption, and walking.



\medterm{Foot Drop} Difficulty lifting the front of the foot because of weakness in the ankle and foot muscles. It may cause tripping or a high-stepping walk and can result from nerve injury, spinal disease, muscle disease, or stroke.

\synonyms
Drop Foot

\medterm{Foot Fracture} A break in one or more foot bones caused by sudden injury, repeated stress, or weakened bone. Pain, swelling, bruising, and difficulty bearing weight are common; treatment ranges from protection and immobilization to surgery.

\synonyms
Broken foot

\medterm{Footling Breech} A fetal position in which one or both feet point downward toward the birth canal. It increases the risk of umbilical-cord prolapse and difficulty delivering the baby, so delivery planning requires specialist obstetric assessment.

\synonyms
Incomplete Breech; Footling Breech Presentation

\medterm{Foot Pain} Pain in the foot caused by injury, overuse, arthritis, nerve problems, infection, or poor circulation. Persistent pain, inability to bear weight, numbness, discoloration, fever, or an open wound requires assessment.


\medterm{Foot, Leg, And Ankle Swelling} Enlargement of the lower limb from fluid, inflammation, injury, infection, vein disease, medicines, or heart, kidney, or liver problems. Sudden one-sided swelling, pain, redness, or breathlessness needs urgent assessment.

\medterm{Foramen} A natural opening or passage through a bone or other body structure, usually allowing a nerve, blood vessel, or other tissue to pass.

\medterm{Foramen Ovale} An opening between the upper chambers of the fetal heart that directs blood away from the unused lungs. It usually closes after birth; persistence is called a patent foramen ovale.

\medterm{Foraminotomy} Surgery that enlarges an opening where a spinal nerve leaves the backbone, relieving pressure caused by arthritis, a disc problem, or other narrowing. Risks include infection, bleeding, nerve injury, and persistent symptoms.

\medterm{Force Plates} Platforms that measure forces beneath the feet during standing, walking, or other movement. Clinicians and researchers use them to assess balance, gait, symmetry, and how weight is distributed.

\medterm{Forced Expiratory Volume} The amount of air forcefully exhaled during a measured time after taking a full breath. The first-second value helps assess narrowed or restricted airways and is interpreted with other lung-test results.

\medterm{Forced Oscillation Technique} A breathing test in which a device sends gentle pressure waves through the airways while a person breathes normally. It measures how easily air moves and can help assess lung disease when forceful breathing tests are difficult.

\medterm{Forced Vital Capacity} The total amount of air forcefully exhaled after a person breathes in fully. A low value may occur when the lungs are small or stiff or when airways prevent complete emptying.

\medterm{Forceps} A handheld instrument used to grasp, hold, or move tissue or medical materials. Different designs have smooth or toothed tips and are chosen according to the procedure.

\medterm{Forceps (Obstetric)} Paired curved instruments placed around the fetal head during selected assisted vaginal births. They can help complete delivery when the baby or pregnant person needs delivery sooner and conditions are suitable.

\synonyms
Obstetric forceps; delivery forceps

\medterm{Forceps (Surgical)} Handheld instruments used during surgery to hold, grasp, move, or manipulate tissue, sutures, dressings, or other materials. Their shape and pressure are selected to limit tissue injury.

\synonyms
Surgical grasping instrument

\medterm{Forceps (Vascular)} Surgical forceps designed to handle blood vessels, clots, tissue, or medical materials during vascular procedures. Appropriate use helps avoid crushing or tearing delicate vessels.

\synonyms
Vascular surgical forceps

\medterm{Forceps Delivery} An assisted vaginal birth in which forceps guide the baby’s head out of the birth canal. It may be used for prolonged labor, concerning fetal findings, or inability to push safely; tears and newborn injury are possible.

\medterm{Forceps-Assisted Delivery} A vaginal delivery assisted by placing forceps around the fetal head. It is considered only when the cervix is fully open, the head’s position is known, and vaginal birth is judged safe; maternal and newborn injuries can occur.

\medterm{Forcible Intercourse} Sexual intercourse without freely given consent, involving force, threats, coercion, or inability to consent. It is sexual violence; survivors may need immediate safety support, medical care, emergency contraception, infection prevention, and counseling.

\medterm{Fordyce'S Angiokeratomas} Small, usually harmless dark-red or purple spots caused by widened tiny blood vessels, commonly on the scrotum or vulva. They may bleed after injury; unusual or changing lesions should be examined.

\medterm{Forearm} The part of the arm between the elbow and wrist, containing the radius and ulna bones, muscles, tendons, nerves, and blood vessels.

\medterm{Forearm Circumference} The distance around the forearm at a specified point. Measurement can help monitor growth, muscle loss, swelling, nutrition, or changes during treatment and rehabilitation.

\medterm{Forehead} The part of the face above the eyebrows and in front of the frontal skull bone. It contains skin, muscles, nerves, and blood vessels and helps with facial expression and protection.

\medterm{Forehead Lift} Cosmetic surgery that raises drooping eyebrows or forehead skin and may reduce wrinkles. Risks include scarring, numbness, asymmetry, hairline changes, bleeding, infection, and nerve injury.

\synonyms
Brow lift

\medterm{Foreign Body} An object that is not normally part of the body and becomes lodged in tissue or a body passage, such as glass, metal, food, or a splinter. It may cause injury, blockage, infection, or inflammation.

\medterm{Foreign Body Airway Obstruction} Partial or complete blockage of the breathing passages by an inhaled object or material. Severe choking with inability to speak, cough, or breathe is life-threatening and requires immediate first aid and emergency help.

\medterm{Foreign Body Aspiration} Inhalation of an object or food into the airway, especially in young children. Choking, persistent cough, wheezing, noisy breathing, or reduced breath sounds may occur; complete blockage is an emergency.

\medterm{Foreign Body In The Nose} An object lodged in a nasal passage, often in a child. One-sided blockage, foul discharge, bleeding, pain, or difficulty breathing may occur; safe removal prevents tissue damage and infection.

\medterm{Foreign Body Reaction} Inflammation around material that remains in tissue, such as a splinter, stitch, implant, or device. It may cause a lump, redness, pain, drainage, or scar tissue and can resemble infection.

\medterm{Foreign Object - Inhaled} An object or material that enters the airway during breathing. Coughing, choking, wheezing, noisy breathing, or one-sided reduced breath sounds may occur; breathing difficulty requires emergency care.

\medterm{Foreign Object - Swallowed} An object that is accidentally swallowed and enters the digestive tract. Many small smooth objects pass naturally, but batteries, sharp objects, magnets, or objects causing pain or blockage require urgent medical assessment.

\medterm{Foreign-Body Aspiration} Entry of an object or food into the voice box, windpipe, or bronchial tubes, especially in young children. Sudden choking, cough, wheezing, or noisy breathing may occur; complete blockage is an immediate emergency.

\medterm{Forensic} Relating to the use of scientific or medical knowledge in legal investigations, court proceedings, identification, injury assessment, or determining the cause and manner of death.

\medterm{Forensic Anthropology} The examination of human remains, especially bones, for legal investigations. It can help determine whether remains are human, estimate biological characteristics, assess injury, and support identification.


\medterm{Forensic Evidence} Physical, biological, digital, or documentary material collected and analyzed during a legal investigation. Examples include blood, DNA, fibers, fingerprints, toxicology samples, photographs, and medical records.

\medterm{Forensic Examination} A structured medical or scientific assessment performed for legal purposes. It may document injuries, collect evidence after assault, assess impairment, examine a death, or analyze a scene or specimen.

\medterm{Forensic Genetics} The use of DNA testing to identify people, establish biological relationships, or compare biological evidence in legal investigations. Results require proper collection, chain of custody, and expert interpretation.

\medterm{Forensic Odontology} The use of dental knowledge in legal investigations. Comparing teeth, dental records, and radiographs can help identify remains and document dental injuries or other evidence.

\medterm{Forensic Pathology} The medical study of deaths and injuries for legal purposes. Forensic pathologists investigate sudden, unexpected, violent, or suspicious deaths and help determine the cause and manner of death.

\medterm{Forensic Psychiatry} The application of psychiatry to legal questions, such as ability to stand trial, mental state during an alleged offense, risk assessment, involuntary treatment, and decision-making capacity.

\medterm{Forensic Radiography} The use of X-rays, computed tomography, magnetic resonance imaging, or other imaging in legal investigations. It can identify people, document injuries, locate objects, assess trauma, and support death investigations.

\medterm{Forensic Toxicology} The testing of blood, urine, hair, tissue, or other specimens for medicines, alcohol, drugs, poisons, and their breakdown products to assess possible roles in illness, impairment, poisoning, or death.

\medterm{Forensic Toxicology Screen} An initial laboratory test of biological samples for medicines, alcohol, drugs, poisons, or their breakdown products. A positive screen usually needs confirmation and measurement before legal or medical conclusions are made.

\medterm{Foreskin} The movable fold of skin covering the head of the penis. It protects the sensitive surface and contains nerves; it may be removed by circumcision.

\medterm{Foreskin (Prepuce)} The fold of skin covering and protecting the head of the penis. It normally becomes retractable over time; forceful retraction can cause pain, tearing, or scarring.

\synonyms
Prepuce

\medterm{Formaldehyde} A strong-smelling chemical used in manufacturing and tissue preservation. It can irritate the eyes, skin, and airways, cause allergic reactions, and increase cancer risk with substantial or long-term exposure.

\medterm{Formalin} A water-based solution of formaldehyde used to preserve tissue for laboratory examination. It can irritate or burn the skin, eyes, and airways and must be handled with appropriate ventilation and protection.

\medterm{Formalin-Fixed Paraffin-Embedded Tissue} Tissue preserved in formalin, processed, and embedded in wax so thin sections can be examined. These blocks support microscopic, protein, and many genetic tests, although preservation can alter some results.

\medterm{Formate} A salt or ester of formic acid. In medicine, formate is important because it is a toxic breakdown product of methanol and can contribute to severe acidosis and vision loss after poisoning.

\medterm{Formed Elements Of Blood} The cells and cell fragments suspended in blood: red blood cells carry oxygen, white blood cells help defend against infection, and platelets help form clots and stop bleeding.

\medterm{Formerly Churg-Strauss Syndrome} An older name for eosinophilic granulomatosis with polyangiitis, an inflammatory blood-vessel disease associated with asthma and high eosinophil levels. It can damage the lungs, nerves, skin, heart, kidneys, or digestive organs.
\medterm{Formetanate} A carbamate insecticide that can poison the nervous system by disrupting acetylcholine breakdown. Exposure may cause sweating, vomiting, small pupils, muscle weakness, breathing difficulty, or seizures and requires urgent medical advice.

\medterm{Formication} The feeling that insects are crawling on or under the skin when none are present. It may occur with nerve disease, menopause, substance withdrawal, medicines, or psychiatric illness and should be medically evaluated if persistent.

\medterm{Formoterol Fumarate} A long-acting inhaled medicine that relaxes airway muscles and helps prevent asthma or chronic obstructive lung disease symptoms. It is not a rescue medicine for sudden severe breathing difficulty and is often combined with an inhaled corticosteroid.

\medterm{Formothion} An organophosphate insecticide that can poison the nervous system. Exposure may cause sweating, excess saliva, vomiting, muscle twitching, breathing difficulty, seizures, or weakness and requires urgent emergency or poison-control advice.

\medterm{Formparanate} A pesticide-related chemical whose health effects depend on the formulation, dose, and route of exposure. Suspected poisoning should be treated as urgent and assessed through poison-control or emergency services.

\medterm{Formula Feed} Prepared infant nutrition given by bottle, cup, or feeding tube when breast milk is unavailable, insufficient, or not chosen. The formula and preparation method should match the infant’s age and medical needs.

\medterm{Formula Feeding} Feeding an infant with commercially prepared formula instead of or in addition to breast milk. Safe preparation, clean equipment, correct dilution, and age-appropriate amounts help prevent infection, dehydration, and poor nutrition.

\medterm{Formula, Infant} Commercially prepared nutrition for infants, available in standard, soy-based, extensively hydrolyzed, amino-acid, or other specialized forms. Choice depends on age, tolerance, growth, and medical needs; preparation must follow instructions exactly.

\medterm{Formulae} The plural of formula. In healthcare, the word may mean mathematical equations, medicine formulations, or prepared nutritional products such as infant formula.

\medterm{Formulary} A list of medicines approved or preferred by a health system, insurer, or institution. It may show coverage, restrictions, alternatives, and requirements for approval, but does not replace a clinician’s prescribing judgment.

\medterm{Fornix} A recess or arch-shaped structure. In the eye, the conjunctival fornix is the fold where the eyelid lining joins the surface covering the eyeball; the term also describes parts of the brain and vagina.

\medterm{Forrester Classification} A system that groups people with severe heart failure or shock according to blood flow and congestion. “Warm” means adequate flow and “cold” reduced flow; “dry” means little congestion and “wet” significant congestion, helping guide urgent treatment.

\synonyms
Forrester Hemodynamic Subsets; Forrester Score

\medterm{Fortification} Adding vitamins or minerals to foods to improve nutritional intake and prevent deficiency. Examples include iodine in salt, vitamin D in some dairy products, and folic acid in many grain products.

\medterm{Fortified Wine} Wine with added distilled alcohol, giving it a higher alcohol content than ordinary wine. Excess use can cause intoxication, dependence, liver disease, injuries, and interactions with medicines.


\medterm{Fosaprepitant Dimeglumine} An intravenous medicine converted to aprepitant that blocks a brain signaling pathway involved in nausea and vomiting. It helps prevent chemotherapy-related sickness and can interact with other medicines.

\medterm{Foscarnet Sodium} An antiviral medicine used for selected serious or resistant cytomegalovirus and herpesvirus infections. It can damage the kidneys and disturb blood salts, so hydration and laboratory monitoring are important.

\medterm{Fosfomycin} An antibiotic that blocks an early step in bacterial cell-wall production. It is used for selected urinary or difficult-to-treat infections; diarrhea, allergic reactions, and resistance can occur.

\medterm{Fosinopril} An angiotensin-converting enzyme inhibitor used for high blood pressure and some forms of heart failure. It relaxes blood vessels and reduces fluid strain; cough, high potassium, kidney problems, or swelling can occur.

\medterm{Fossa} A shallow hollow or depression in a bone or other body structure. It may provide space for a joint, muscle, blood vessel, nerve, gland, or part of an organ.

\medterm{Fossula} A small pit or shallow depression in an anatomical structure. Examples include small depressions in a bone or on a tooth; the term describes shape and location rather than a disease.


\medterm{Fotemustine} A chemotherapy medicine that can cross into the brain and is used in some countries for selected advanced cancers. It can suppress bone marrow, increasing the risk of infection and bleeding.

\medterm{Found Dead} A circumstance in which a person is discovered deceased without the death having been observed. Medical or forensic examination may determine identity, timing, cause, and manner of death.

\medterm{Four-Dimensional Ultrasound} Three-dimensional ultrasound shown as movement over time, creating a live view. It may help assess fetal anatomy or other structures, but image quality and diagnostic value depend on the examination and equipment.

\medterm{Fourchette} The small fold of tissue at the back of the opening of the vulva, where the inner labial folds meet. It may stretch or tear during childbirth or sexual injury.

\medterm{Fourier Analysis} A mathematical method that separates a complex signal into simpler repeating frequencies. It is used in medical imaging, heart and brain signal analysis, and other measurements to identify patterns.

\medterm{Fournier Gangrene} A rapidly spreading, life-threatening infection that destroys tissue around the genitals, perineum, or anus. Severe pain, swelling, fever, skin discoloration, or feeling very ill requires immediate surgery and antibiotics.


\medterm{Fourth Finger} The ring finger, located between the middle finger and little finger. It contains bones, joints, tendons, nerves, blood vessels, and skin that support movement and grasping.

\medterm{Fourth Stage Of Labor} The first hours after the placenta is delivered, when the birthing person is closely observed for heavy bleeding, abnormal blood pressure, poor uterine contraction, pain, and other early complications.

\medterm{Fourth Venereal Disease} An old name for lymphogranuloma venereum, a sexually transmitted infection caused by particular Chlamydia strains. It may cause a small genital sore followed by painful groin swelling and, untreated, rectal complications.

\medterm{Fourth-Degree Tear} A severe childbirth tear extending from the vaginal area through the anal sphincter and into the rectal lining. It needs careful repair, antibiotics when appropriate, follow-up, and monitoring for bowel-control problems.

\medterm{Fovea} A tiny central depression in the retina responsible for the sharpest detailed and color vision. Damage to it can impair reading, face recognition, and central vision.

\medterm{Fovea Centralis} The central pit of the retina with the greatest concentration of cone cells, providing the sharpest central and color vision. It has no major blood vessels crossing directly over its center.

\medterm{Fowler Position} A semi-upright position with the head and upper body raised, commonly about 45 to 60 degrees. It may improve comfort and breathing or be used during examinations and procedures.

\medterm{Foxglove Poisoning} Poisoning from foxglove plants or medicines containing cardiac glycosides. It can cause nausea, vomiting, confusion, visual changes, dangerously slow or irregular heart rhythms, or cardiac arrest and requires emergency care.

\medterm{Fractional Excretion Of Sodium} The percentage of filtered sodium leaving the body in urine. It can help assess kidney function and distinguish reduced kidney blood flow from some tubular injuries, but diuretics and other illnesses can make it misleading.

\medterm{Fractional Excretion Of Urea (FEUrea)} The percentage of filtered urea excreted in urine. It can support assessment of kidney blood flow when diuretics affect sodium testing, but infection, bleeding, nutrition, and chronic kidney disease also influence the result.

\synonyms
FEUrea

\medterm{Fractional Flow Reserve} A pressure measurement made during coronary angiography to determine whether a narrowed heart artery significantly limits blood flow. The result helps guide decisions about stents or medical treatment.

\medterm{Fractional IPV} A smaller-than-usual dose of inactivated polio vaccine given into the skin rather than the muscle. When used in an approved schedule, it helps protect against polio; the number of doses and timing follow local immunization guidance.

\medterm{Fracture} A break or crack in a bone caused by injury, repeated stress, or weakened bone. Pain, swelling, bruising, deformity, or inability to use the part may occur; treatment restores alignment and supports healing.

\medterm{Fracture Displacement} Movement of broken bone pieces out of their normal alignment, including sideways shift, angling, shortening, or rotation. The need for reduction or surgery depends on the bone, joint involvement, age, and function.

\medterm{Fracture Fixation} Stabilization of a broken bone with a cast, splint, plate, screw, nail, wire, or external frame. The goal is to maintain alignment while the bone heals and, when possible, preserve movement.

\medterm{Fracture Liaison Service} A coordinated healthcare program for people with fragility fractures. It evaluates bone strength and underlying causes, arranges treatment and fall prevention, and follows patients to reduce future fractures.

\medterm{Fracture Of The Temporal Bone} A break in a skull bone containing parts of the ear and nearby nerves and blood vessels. It may cause hearing loss, dizziness, facial weakness, bleeding, or leakage of brain fluid and needs urgent assessment.

\medterm{Fracture Prevention} Measures that lower the chance of broken bones, including treating osteoporosis, improving strength and balance, preventing falls, correcting vitamin or mineral deficiency, and reviewing medicines that increase risk.

\medterm{Fracture, Arm} A break in the humerus, radius, or ulna caused by injury, a fall, or weakened bone. Pain, swelling, deformity, and loss of movement may occur; treatment ranges from immobilization to surgery.

\synonyms
Broken arm

\medterm{Fracture, Growth Plate} A break involving the cartilage growth area of a child’s or adolescent’s bone. It can disturb future growth if displaced or untreated, so prompt imaging and specialist management may be needed.

\synonyms
Growth plate fractures; physeal fractures

\medterm{Fracture, Hip} A break in the upper femur near the hip joint, often after a fall in older adults or major trauma in younger people. It causes pain and inability to walk and commonly needs urgent surgery.

\synonyms
Hip fracture

\medterm{Fracture, Leg} A break in the tibia, fibula, or both lower-leg bones, usually from trauma, a fall, or sports injury. Pain, swelling, deformity, and inability to bear weight may occur; treatment depends on stability and alignment.

\synonyms
Broken leg

\medterm{Fracture, Stress} A small crack caused by repeated loading or overuse rather than one major injury. It commonly affects the foot or lower leg and causes pain that worsens with activity; continued stress can complete the break.

\synonyms
Stress fractures; fatigue fractures

\medterm{Fracture-Fixation Device} A plate, screw, rod, pin, wire, or external frame used to hold broken bone pieces in alignment during healing. Choice depends on the fracture, bone quality, soft-tissue injury, and expected movement.

\medterm{Fractured Clavicle In The Newborn} A broken collarbone in a newborn, usually related to birth stress. The baby may move one arm less or cry when handled; most heal well with gentle support and follow-up.

\medterm{Fractured Elbow} A break involving one of the bones forming the elbow. It can cause pain, swelling, deformity, and limited movement and may injure nearby nerves or blood vessels.

\medterm{Fractured Ilium} A break in the broad upper part of the pelvic bone, usually from major trauma or weakened bone. It may cause severe pain and can accompany injury to pelvic organs or blood vessels.

\medterm{Fractured Nose} A break in the nasal bones or supporting cartilage, usually after a blow. Pain, swelling, bleeding, and deformity may occur; a collection of blood inside the septum requires urgent treatment.

\synonyms
Broken nose; nasal fracture

\medterm{Fractured Rib} A break in one or more ribs, usually from injury or forceful coughing. Sharp pain worsens with breathing or movement; lung collapse, bleeding, or injury to nearby organs must be considered after significant trauma.

\synonyms
Broken ribs; rib fracture

\medterm{Fragile X Syndrome} An inherited condition caused by a change in the FMR1 gene. It can affect learning, speech, behavior, social development, and physical features and is a common inherited cause of intellectual disability.

\medterm{FRAIL Scale} A brief questionnaire assessing fatigue, ability to climb stairs, walking ability, illness burden, and unintentional weight loss. It screens for frailty risk but does not replace a complete medical and functional assessment.

\medterm{Frailty} A state of reduced physical reserve that makes an older or seriously ill person more vulnerable to falls, infection, hospitalization, disability, and other stress. Signs include weakness, slow walking, exhaustion, inactivity, and unintentional weight loss.

\medterm{Framboesia} Another name for yaws, a chronic infection caused by Treponema pallidum pertenue. It spreads mainly through skin contact in tropical areas and can cause skin lesions, bone pain, and deformity if untreated.

\medterm{Frameshift Mutation} A genetic change caused by inserting or deleting DNA bases in numbers that are not multiples of three. This shifts how the genetic code is read and often produces a severely altered or nonworking protein.

\medterm{Framycetin} An aminoglycoside antibiotic applied to selected bacterial skin, ear, or eye infections. It can cause local allergy or irritation; prolonged or extensive use may allow absorption and antibiotic-related toxicity.

\medterm{Francisella} A group of small bacteria that includes Francisella tularensis, which causes tularemia. People can acquire infection through ticks, animals, contaminated water, food, or inhaled particles.

\medterm{Francisella Tularensis} A bacterium that causes tularemia, usually acquired through ticks, deer flies, infected animals, contaminated water or food, or inhalation. It can cause fever, an ulcer with swollen lymph nodes, pneumonia, or severe systemic illness.

\medterm{Francisellaceae} A family of small bacteria that includes Francisella species. Some cause tularemia or related infections in animals and people; disease type depends on the species and route of exposure.

\medterm{Frank Breech} A fetal position in which the buttocks are nearest the birth canal, the hips are bent, and the knees are straight with the feet near the head. Delivery planning depends on gestational age, fetal size, and clinical circumstances.

\synonyms
Extended Breech; Frank Breech Presentation

\medterm{Frank-Starling Mechanism} The heart normally contracts more strongly when its chambers fill with more blood, up to a limit. This helps match the amount pumped with the amount returning to the heart.

\medterm{Frasier Syndrome} A rare inherited disorder caused by changes in the WT1 gene. It can cause abnormal sexual development, progressive kidney disease, and increased risk of tumors in the gonads.

\medterm{Fraternal Twins} Twins formed when two separate eggs are fertilized by two separate sperm. They are genetically like ordinary siblings and may be the same or different sex.

\medterm{FRAX Tool} A calculator estimating a person’s ten-year probability of a major osteoporosis-related fracture or hip fracture. It uses factors such as age, sex, bone density, previous fracture, and other risks and supports treatment decisions.
\synonyms
FRAX fracture-risk calculator

\medterm{Freckle} A small, flat, tan or light-brown area of increased skin pigment that often darkens with sunlight. Freckles are usually harmless, but a changing, bleeding, irregular, or unusual spot needs examination.

\medterm{Freckles} Small, flat pigmented spots that often become darker after sun exposure and reflect inherited skin characteristics. They are usually harmless; changes in size, shape, color, symptoms, or bleeding require assessment.

\medterm{Freckles (Ephelides)} Small, flat, light-brown skin spots caused by increased pigment production. They commonly darken with sunlight and fade when exposure decreases and are usually harmless.

\synonyms
Ephelides

\medterm{Free Fragment} A piece of spinal-disc material that has broken away from the main disc and may press on a nearby nerve. It can cause back or leg pain, numbness, weakness, or bladder and bowel problems.

\synonyms
Sequestered disk fragment

\medterm{Free Radical} A highly reactive molecule with an unpaired electron. The body makes some during normal metabolism, while excessive amounts can damage cell membranes, proteins, and DNA as part of oxidative stress.

\medterm{Free Radical Injury} Cell damage caused by excessive reactive molecules that alter fats, proteins, or DNA. It can contribute to inflammation and tissue injury, although the importance varies with the disease and exposure.

\medterm{Free T4 Test} A blood test measuring unbound thyroxine, the thyroid hormone available to tissues. It is usually interpreted with thyroid-stimulating hormone to assess overactive or underactive thyroid function and treatment response.

\medterm{Free Thyroxine Index} A calculated estimate of the thyroid hormone available to tissues, using total thyroxine and a measure of hormone binding. It can help assess thyroid function when binding-protein levels affect total hormone results.

\medterm{Free Water Clearance} A kidney measurement showing whether urine contains more or less water than blood plasma. Positive values indicate excretion of dilute urine; negative values indicate water retention and concentrated urine.

\medterm{Freeze Drying} A preservation process that freezes material and removes ice under low pressure. It helps preserve medicines, vaccines, biological samples, and foods while limiting damage from heat and moisture.


\medterm{Freezing} Exposure to very low temperature that can freeze tissue, cause frostbite, or dangerously lower body temperature. The term also describes preserving food, medicines, or samples by storing them at low temperature.

\medterm{Fremitus} A vibration felt or heard through body tissues. During a chest examination, transmitted voice vibrations can change with fluid, air, or solid tissue in or around the lungs.

\medterm{French Disease} An old name for syphilis, a sexually transmitted infection caused by Treponema pallidum. Untreated infection can progress from sores to widespread disease and later damage the brain, nerves, heart, or other organs.

\medterm{Frenectomy} A procedure that removes or releases a frenulum, a fold of tissue that may restrict tongue movement, pull on the gums, separate teeth, or interfere with dentures. Bleeding, infection, and scarring are possible.

\medterm{Frenulum} A small fold or band of skin or mucous membrane that anchors and limits movement of a body part, such as beneath the tongue, inside the lips, or on the penis.

\medterm{Frenulum Breve} An abnormally short penile frenulum that may cause pain, tearing, or bleeding during erection or sexual activity. Treatment depends on symptoms and may include observation or a minor surgical procedure.

\medterm{Frenum} A narrow fold of tissue that attaches one body surface to another and limits movement. Examples include the folds beneath the tongue, inside the lips, and beneath the foreskin.

\medterm{Frenum (Frenulum)} A fold of skin or mucous membrane that anchors a body part and limits its movement. Examples include the tongue fold, lip folds, and penile frenulum.

\synonyms
Frenulum

\medterm{Frequency} Urinating more often than usual, sometimes with only small amounts passed. It differs from producing a large total urine volume and may result from infection, pregnancy, diabetes, medicines, or bladder disorders.

\medterm{Frequency Of Urination} Needing to urinate more often than usual. Causes include increased fluid intake, urinary infection, pregnancy, diabetes, diuretic medicines, bladder irritation, prostate enlargement, or other urinary disorders.

\medterm{Frequency Range} The lowest to highest sound or signal frequencies that a device can detect, produce, or amplify effectively. In hearing care, it helps describe which pitches a hearing aid can process.

\medterm{Frequency Response} How strongly a device or body system responds to different sound or signal frequencies. In a hearing aid, it describes the amount of amplification provided for low, middle, and high pitches.

\medterm{Frequency, Urinary} The need to urinate more often than usual, sometimes passing only small amounts. Infection, irritation, pregnancy, diuretics, diabetes, overactive bladder, or prostate disease may cause it.

\medterm{Frequent Bowel Movements} Passing stool more often than is usual for a person. Causes include diet, infection, medicines, inflammation, poor absorption, or functional bowel disorders; blood, fever, dehydration, or weight loss needs assessment.

\medterm{Frequent Or Urgent Urination} Urinating often or having sudden urges that are difficult to postpone. Infection, bladder irritation, diabetes, pregnancy, medicines, obstruction, or overactive bladder can cause these symptoms.

\medterm{Frequent Urination} Needing to urinate more often than usual. It may result from drinking more, infection, diabetes, pregnancy, medicines, prostate enlargement, overactive bladder, or another urinary condition.

\medterm{Fresh Specimen} A recently collected biological sample tested promptly or preserved correctly before changes affect its cells, organisms, chemistry, or diagnostic accuracy.




\medterm{Frey’S Syndrome} Sweating, warmth, or flushing over the cheek or temple while eating, caused by nerve regrowth after parotid-gland surgery or injury. Symptoms may be treated with topical medicines, injections, or other specialist care.

\medterm{Friable} Describes tissue that is unusually soft, fragile, and easily torn. Friable tissue may bleed readily when touched and can occur with inflammation, infection, cancer, or poor healing.

\medterm{Friction} Rubbing between surfaces. Skin friction can cause irritation, blisters, or sores; reducing moisture, pressure, and rubbing helps prevent injury.

\medterm{Friedman Curve} An older graph describing cervical opening over time during labor. Modern labor care recognizes wider normal variation and does not use this historical curve alone to diagnose slow labor or decide on intervention.

\medterm{Friedreich Ataxia} An inherited disorder that progressively damages nerves and coordination. It can cause unsteady walking, weakness, reduced sensation, spinal curvature, heart muscle disease, and sometimes diabetes.

\medterm{Friedreich'S Ataxia} An inherited condition that gradually affects balance, coordination, sensation, and strength. Some people also develop heart muscle disease, spinal curvature, or diabetes and need regular neurologic and cardiac care.
\medterm{Friedreichs Ataxia Syndrome} An inherited disorder that gradually damages nerves and coordination. It can cause unsteady walking, weakness, reduced sensation, spinal curvature, heart muscle disease, and diabetes and usually requires long-term specialist care.

\medterm{Frigidity} An outdated and stigmatizing term for sexual difficulties. Modern assessment describes the specific problem, such as reduced desire, arousal, orgasm, or pain, and considers physical, emotional, relationship, and medicine-related causes.
\medterm{Frontal} Relating to the forehead, frontal skull bone, or frontal lobe of the brain. The frontal lobe helps control voluntary movement, planning, attention, speech, behavior, and decision-making.

\medterm{Frontal Bone} The skull bone forming the forehead, roof of the eye sockets, and part of the nasal region. It contains the frontal sinuses and protects the front of the brain.

\medterm{Frontal Bossing} An unusually prominent forehead caused by forward bulging of the frontal skull bones. It may be familial or associated with bone disease, excess pressure in the skull, or skeletal disorders.

\medterm{Frontal Lobe} The front portion of the brain’s outer layer. It helps control movement, planning, attention, judgment, behavior, personality, and speech production in the language-dominant hemisphere.

\medterm{Frontal Lobe Epilepsy} Epilepsy caused by seizures beginning in the frontal lobe. Episodes may involve sudden movements, unusual behavior, impaired awareness, or events during sleep and can be difficult to distinguish from sleep disorders.

\synonyms
Frontal lobe seizures

\medterm{Frontal Lobe Seizures} Seizures beginning in the front part of the brain, often causing brief unusual movements, behaviors, sensations, speech changes, or impaired awareness. Symptoms depend on the exact brain area involved.

\synonyms
Epilepsy, Frontal Lobe

\medterm{Frontal Sinus} An air-filled cavity within the frontal skull bone above the eye. It connects with the nasal passages and can become inflamed or blocked during sinusitis.

\medterm{Frontal Sinusitis} Inflammation of the air-filled cavities above the eyes, usually from a viral infection or blocked drainage. It can cause forehead pain, congestion, and discharge; severe headache, eye swelling, fever, confusion, or neurologic symptoms require urgent care.

\medterm{Frontotemporal Dementia} A group of progressive brain disorders affecting the frontal and temporal lobes. Early changes often involve personality, behavior, judgment, language, or movement rather than memory alone.

\synonyms
Dementia, Frontotemporal

\medterm{Frontotemporal Lobar Degeneration} Progressive loss of nerve cells in the frontal and temporal brain regions. It can cause early changes in behavior, personality, language, judgment, or movement and may occur at a relatively young age.

\medterm{Frostbite} A freezing injury to the skin and deeper tissues that can cause numbness, pale or waxy skin, swelling, blisters, and tissue death.

\medterm{Frotteuristic Disorder} A mental disorder involving persistent sexual arousal from touching or rubbing against a person without consent, causing distress or impairment or involving acting on the urges. Nonconsensual touching is sexual assault.

\medterm{Frozen Section} A rapid tissue examination performed during surgery. A specimen is frozen, cut into thin slices, and examined to provide immediate information about a lesion or surgical margin, although permanent sections are often more accurate.

\medterm{Frozen Shoulder} A painful condition in which the shoulder capsule tightens, progressively limiting movement. It often improves slowly over months or years; diabetes, thyroid disease, and prolonged immobility increase risk.

\synonyms
Adhesive Capsulitis Of Shoulder (Frozen Shoulder)


\medterm{Fructosamine} A blood measurement of sugar attached to blood proteins, reflecting average blood glucose over roughly the previous two to three weeks. It can help when the longer-term hemoglobin A1c test is unreliable.

\medterm{Fructose} A natural sugar found in fruit, honey, some vegetables, and many sweetened foods. It is processed mainly by the liver; poor absorption can cause gas, bloating, cramps, or diarrhea after eating it.

\synonyms
Fruit sugar

\medterm{Fructose Intolerance} Difficulty digesting or processing fructose, depending on the condition. Poor intestinal absorption causes gas, bloating, and diarrhea, while inherited inability to process fructose can cause low blood sugar and liver injury.

\medterm{Fructose Malabsorption} Incomplete absorption of fructose in the small intestine. Unabsorbed sugar is fermented by bacteria and can cause gas, bloating, cramps, or diarrhea; reducing personal dietary triggers may help.

\medterm{Fructose Metabolism} The body’s processing of fructose, mainly in the liver, into substances used for energy or stored. Rare enzyme defects can cause harmless fructose in urine or serious low blood sugar and liver injury after intake.


\medterm{Fruity Breath} A sweet, acetone-like breath odor caused by ketones. In someone with diabetes, especially with high glucose, vomiting, abdominal pain, deep breathing, or confusion, it may signal life-threatening ketoacidosis.


\medterm{Frustration} An emotional response to blocked goals, unmet needs, or difficult circumstances. It may cause irritability, anger, distress, withdrawal, or behavior changes and becomes clinically important when persistent or impairing.

\medterm{Fryns Syndrome} A rare inherited condition present from birth, often involving a hole in the diaphragm, characteristic facial features, and abnormal fingers or toes. Brain, heart, kidney, or genital abnormalities may also occur.

\medterm{FSH} An abbreviation for follicle-stimulating hormone, which supports ovarian follicle development and sperm production. Blood levels help assess infertility, puberty, menopause, and disorders of the pituitary gland, ovaries, or testes.

\synonyms
Follicle-stimulating hormone

\medterm{Ft.} An abbreviation for foot or feet, a length of 12 inches or about 30.48 centimeters. In medical records, the surrounding context should distinguish it from other abbreviations.

\medterm{FTA-ABS Blood Test} A blood test detecting antibodies against Treponema pallidum, the bacterium that causes syphilis. It supports diagnosis but often stays positive after treatment and cannot by itself show whether infection is active.

\medterm{FTA-ABS Test} A syphilis blood test that detects antibodies against Treponema pallidum. Because antibodies may remain for life after infection, clinicians interpret the result with other tests, symptoms, treatment history, and exposure risk.

\medterm{Fuchs Dystrophy} A progressive eye disorder in which cells lining the inner cornea gradually fail to remove fluid. It can cause blurred morning vision, glare, corneal swelling, pain, and worsening sight.

\medterm{Fuchs Endothelial Corneal Dystrophy} A usually bilateral disorder in which the cornea’s inner pumping cells are lost and tiny deposits develop. Fluid then accumulates, causing morning blur, glare, swelling, pain, and reduced vision.

\synonyms
Fuchs dystrophy

\medterm{Fuchs' Dystrophy} A progressive disorder of the cornea’s inner cells that causes fluid buildup, swelling, blurred vision, glare, and sometimes painful blisters on the corneal surface.
\medterm{Fucose} A sugar found in some cell-surface proteins and fats. It helps cells recognize and attach to one another and participates in immune signaling and other biological processes.

\medterm{Fucosidosis} A rare inherited disorder in which cells cannot break down certain sugar-containing substances. Their buildup progressively affects development, movement, bones, growth, sweating, and organ function.

\medterm{Fucosyltransferase} An enzyme that attaches the sugar fucose to proteins or fats. This changes how cells recognize one another, develop, interact with microbes, and display some blood-group markers.

\medterm{Fuel Oil Poisoning} Illness after swallowing, inhaling, or having fuel oil on the skin. Inhalation can damage the lungs and cause coughing or breathing difficulty; vomiting should not be induced, and urgent poison-control advice is needed.

\medterm{Fugue} A rare dissociative state involving unexpected travel or wandering with loss of memory about personal identity or the past. It needs assessment for trauma, mental illness, substances, neurologic disease, and safety risks.

\medterm{Fukuyama Congenital Muscular Dystrophy} A rare inherited muscle disorder beginning in infancy, causing low muscle tone, weakness, contractures, and delayed development. Abnormalities of the brain and eyes, seizures, and heart disease may also occur.

\synonyms
FCMD; Fukuyama Type Congenital Muscular Dystrophy

\medterm{Fulguration} Destruction or sealing of tissue using electrical energy during a procedure. It may control bleeding or treat selected abnormal tissue, but nearby tissue can be burned or injured.

\medterm{Full Agonist} A medicine or chemical that binds to a cell receptor and produces the strongest response that receptor can generate. The actual effect still depends on dose, receptor number, tissue, and individual response.

\medterm{Full Blood Count} See \textit{Complete Blood Count}.

\synonyms
FBC

\medterm{Full Liquid Diet} A diet containing liquids and foods that become liquid at room temperature, such as milk, strained soups, yogurt, and custard. It may be used temporarily, but prolonged use can lack fiber and some nutrients.

\medterm{Full Sibling} A person who shares both biological parents with another person. This relationship is relevant to family medical history, genetic-risk assessment, and some tissue or organ-donor matching.

\medterm{Full-Thickness Skin Graft} A graft containing the outer skin layer and the entire dermis, moved to cover a wound or reconstruct a defect. It usually gives durable coverage but requires a suitable donor site and blood supply.

\medterm{Fulminant} Beginning suddenly and becoming severe very quickly. The term describes a dangerous pattern of illness, such as rapidly progressive liver failure, infection, or inflammation requiring urgent treatment.

\medterm{Fulminant Hepatitis} Rapid, severe liver injury causing loss of liver function, abnormal bleeding, and sometimes confusion or coma from toxin buildup. It is a medical emergency that may require intensive care and liver transplantation.

\medterm{Fulminant Myocarditis} Sudden, severe inflammation of the heart muscle causing poor pumping, dangerous rhythms, or shock. It requires intensive care and treatment of the cause and complications; temporary mechanical heart support may be needed.

\medterm{Fulvestrant} An injectable hormone treatment that blocks and helps destroy estrogen receptors in selected hormone-sensitive breast cancers. It can cause injection-site pain, hot flushes, joint symptoms, fatigue, and liver-test changes.
\medterm{Fumigation} Treating an enclosed area with gas, vapor, or smoke to kill pests or microorganisms. People and animals must leave beforehand and return only after the area is declared safe and adequately ventilated.

\medterm{Functional Asplenia} Poor or absent spleen function even though the spleen is present. It increases the risk of severe infections, especially from certain bacteria, so vaccination, prompt fever assessment, and sometimes preventive antibiotics are important.

\medterm{Functional Assessment} Evaluation of how well a person moves, communicates, thinks, and performs daily activities. It identifies support needs, tracks change, guides rehabilitation, and helps plan safe care.

\medterm{Functional Capacity} The level of physical, mental, communication, and self-care activity a person can safely perform under particular conditions. It depends on health, environment, assistance, and personal goals.

\medterm{Functional Disorder} A condition in which symptoms or impaired body function arise from altered processing or regulation without a structural abnormality sufficient to explain them. The symptoms are real and can cause substantial disability.

\medterm{Functional Dyspepsia} Ongoing upper-abdominal discomfort, burning, early fullness, or unpleasant fullness after meals without an ulcer or other structural cause explaining symptoms. It may involve altered stomach movement or gut sensitivity.

\synonyms
Nonulcer Dyspepsia
Nonulcer Stomach Pain

\medterm{Functional Electrical Stimulation} Electrical impulses applied to nerves or muscles to produce a useful movement, such as lifting the foot during walking. It may assist rehabilitation after stroke, spinal injury, or other nerve damage.

\medterm{Functional Gastrointestinal Disorders} Digestive symptoms caused by altered gut movement or sensitivity without an obvious structural disease. Symptoms may include pain, bloating, constipation, diarrhea, nausea, or early fullness.

\synonyms
Disorders of gut-brain interaction

\medterm{Functional Hypothalamic Amenorrhea} Missing menstrual periods because stress, inadequate energy intake, excessive exercise, or disordered eating suppresses reproductive hormones. It can lower estrogen, impair fertility, and weaken bones, but often improves when the underlying imbalance is corrected.

\synonyms
Hypothalamic Amenorrhea; FHA

\medterm{Functional Incontinence} Urine leakage because a person cannot reach, recognize, or use the toilet in time, even when bladder storage may be adequate. Mobility, memory, vision, clothing, communication, or environmental barriers may contribute.

\medterm{Functional Independence Measure} A structured tool that rates how much help a person needs with self-care, toileting, transfers, walking, communication, and thinking. It is used mainly in rehabilitation to track function and progress.

\medterm{Functional MRI} A magnetic-resonance scan that estimates brain activity from changes in blood flow and oxygen use. It can map language or movement areas before selected brain procedures and does not use X-ray radiation.

\medterm{Functional Murmur} An extra heart sound caused by increased blood flow through a structurally normal heart. It may occur with fever, anemia, pregnancy, or exercise and is assessed to exclude valve or heart disease.

\medterm{Functional Neurologic Disorder} A disorder in which altered nervous-system function causes genuine symptoms such as weakness, abnormal movements, seizures, or sensory changes without the structural damage expected for them.
\synonyms
FND; functional neurological disorder; functional neurological symptom disorder; conversion disorder

\medterm{Functional Neurologic Disorder/Conversion Disorder} Alternate index label; see \textit{Functional Neurologic Disorder}.

\medterm{Functional Neurological Disorder} Alternate index label; see \textit{Functional Neurologic Disorder}.

\medterm{Functional Neurological Symptom Disorder} Alternate index label; see \textit{Functional Neurologic Disorder}.

\medterm{Functional Neurological Symptom Disorder (Conversion Disorder)} Alternate index label; see \textit{Functional Neurologic Disorder}.

\medterm{Functional Residual Capacity} The amount of air left in the lungs after a normal relaxed breath out. It helps keep air sacs open and may be reduced by obesity or restrictive lung disease or increased by air trapping.

\medterm{Functional Screening} A brief check of a person’s ability to move safely, communicate, perform daily activities, and manage basic self-care. Abnormal results may lead to a fuller medical, occupational, or rehabilitation assessment.

\medterm{Functional Status} A person’s ability to manage daily activities such as bathing, dressing, eating, toileting, walking, shopping, cooking, medicines, finances, and transportation. Decline may signal illness or a need for support.

\medterm{Functional Urinary Incontinence} Urine leakage that happens because a person cannot reach, recognize, or use the toilet in time, even when the bladder may work normally. Problems with walking, memory, vision, clothing, communication, or the environment can contribute; treatment addresses these barriers and other causes of leakage.

\medterm{Fundal Height} The distance from the pubic bone to the top of the uterus, measured during pregnancy. It roughly follows weeks of pregnancy after about 20 weeks; differences may reflect dating error, growth problems, twins, or fluid changes.

\medterm{Fundoplication} Surgery that wraps the upper stomach around the lower esophagus to strengthen the valve and reduce acid reflux. It may also repair a hiatal hernia; swallowing difficulty, gas, or recurrent reflux can occur afterward.

\medterm{Fundoscopy} Examination of the back of the eye using an ophthalmoscope or camera. It views the retina, optic nerve, macula, and blood vessels to assess conditions such as diabetes, high blood pressure, or retinal disease.

\medterm{Fundus} The rounded or farthest part of a hollow organ. Examples include the upper part of the stomach, the top of the uterus, and the inside back surface of the eye.

\medterm{Fundus Oculi} The inside back surface of the eye visible through the pupil. It includes the retina, optic nerve, macula, and retinal blood vessels.

\medterm{Fundus Photography} Photographing the back of the eye to record the retina, optic nerve, macula, and blood vessels. Images help detect and monitor retinal disease and allow comparison over time.

\medterm{Fundus Uteri} The rounded upper part of the uterus above the openings of the fallopian tubes. Its position and size can be assessed during pelvic examination and pregnancy.

\medterm{Fundus, Retinal} The inside back surface of the eye viewed during fundoscopy, including the retina, optic nerve, macula, and retinal blood vessels.

\medterm{Fungal} Relating to fungi or caused by a fungus. Fungal infections may affect the skin, nails, mouth, lungs, bloodstream, or other organs.

\medterm{Fungal Arthritis} A fungal infection inside a joint, usually after surgery, injury, injection, bloodstream spread, or with weakened immunity. It often develops slowly with persistent pain, swelling, and reduced movement and may need prolonged treatment and surgery.

\medterm{Fungal Endophthalmitis} A fungal infection inside the eye, often after bloodstream infection, surgery, injury, or weakened immunity. It can rapidly damage vision and requires urgent specialist treatment with antifungal medicine and sometimes surgery.

\medterm{Fungal Infection} An infection caused by a fungus. It may affect the skin, nails, or mouth, or spread to the lungs, bloodstream, or other organs, especially when the immune system is weakened.

\medterm{Fungal Infection (Tinea/Dermatophytosis)} A superficial infection of skin, hair, or nails caused by dermatophyte fungi. Ringworm, athlete’s foot, and jock itch commonly cause itchy, scaly, or ring-shaped rashes in warm, moist areas.

\synonyms
Tinea/Dermatophytosis

\medterm{Fungal Infection, Nail} A fungal infection of a fingernail or toenail causing thickening, discoloration, brittleness, or separation from the nail bed. Diagnosis may require testing because injury and skin disease can look similar.

\synonyms
Onychomycosis; nail fungus

\medterm{Fungal Nail Infection} Infection of a fingernail or toenail by a fungus, causing thickening, discoloration, brittleness, debris beneath the nail, or separation from its bed. Treatment depends on the organism, severity, and medical history.

\medterm{Fungal Nails} Nails infected by fungi may become thick, yellow or brown, brittle, crumbly, or separated from the nail bed. Testing may be needed before treatment, which can be topical or oral.

\synonyms
Onychomycosis (Fungal Nails)

\medterm{Fungal Pathogenesis} The ways fungi cause disease, including attaching to cells, growing into tissue, avoiding immune defenses, and releasing damaging substances. The process varies with the fungus, body site, and person’s immunity.

\medterm{Fungal Resistance Mechanisms} Changes that allow fungi to survive antifungal medicines, such as altering a drug target, pumping medicine out of the cell, or changing cell-wall production. Resistance can make infections harder to treat.

\medterm{Fungal Skin Infections} Infections of the skin caused by dermatophytes or yeasts, including ringworm, athlete’s foot, and candidiasis. They often cause red, itchy, scaly, or ring-shaped areas and thrive in warm, moist conditions.

\medterm{Fungating} Describing a mass or wound that grows outward with an irregular, ulcerated, dead, or foul-smelling surface. It may occur with advanced cancer or severe chronic wounds and can bleed, leak, and become infected.

\medterm{Fungemia} The presence of living fungi in the bloodstream, most often Candida in hospitalized patients. It can spread to organs and requires urgent cultures, source evaluation, and appropriate antifungal treatment.

\medterm{Fungi} Organisms including yeasts, molds, and mushrooms. Some live harmlessly in the environment or body, while others cause infections of the skin, mouth, lungs, bloodstream, or other organs.

\medterm{Fungicide} A chemical used to kill or inhibit fungi. Some fungicides are used on crops or surfaces, while selected medicines act against fungi in people; toxicity depends on the substance and exposure.

\medterm{Fungoid} Resembling a fungus or mushroom in appearance. In medicine it may describe an outward-growing, irregular lesion or mass, but appearance alone does not establish its cause.

\medterm{Fungus} A type of organism including yeasts and molds, with distinctive cell walls and membranes. Fungi reproduce in various ways and can cause superficial or invasive infection, especially when immunity is weakened.

\medterm{Funiculus} A cord-like bundle of nerve fibers or other structures. In the spinal cord, a funiculus contains nerve pathways; in the spermatic cord, it is the bundle carrying the vas deferens, blood vessels, nerves, and lymph vessels.

\medterm{Funisitis} Inflammation of the umbilical cord, usually associated with infection or inflammation of the membranes around the fetus. It may increase newborn infection risk and requires assessment of the placenta, mother, and baby.

\medterm{Funnel Chest} A chest-wall shape present from birth in which the breastbone and nearby ribs curve inward. It is often mild, but severe deformity may affect exercise, breathing, body image, or heart function.

\synonyms
Pectus excavatum

\medterm{Funnel-Web Spider Bite} A potentially life-threatening bite from a funnel-web spider. Venom may cause pain, sweating, tingling, muscle spasms, breathing difficulty, or collapse; keep the person still, apply pressure immobilization, and seek emergency care.

\medterm{Furan} A chemical compound used in manufacturing and found in small amounts in some heated foods. High or repeated exposure may harm the liver and other organs, so occupational exposure should be controlled.

\medterm{Furniture Polish Poisoning} Illness after swallowing or inhaling furniture-polish chemicals. Some products can damage the lungs if aspirated and cause drowsiness or chemical injury; do not induce vomiting and contact poison-control services urgently.

\medterm{Furocoumarin} A plant chemical found in some citrus and other plants that can make skin unusually sensitive to ultraviolet light. Contact followed by sunlight may cause redness, blistering, or dark pigmentation.

\medterm{Fursultiamin} A fat-soluble form of thiamine, or vitamin B1, used in some countries as a vitamin preparation. Its clinical value and safety depend on the indication, dose, and other medicines.

\medterm{Furuncle} A painful, pus-filled infection of a hair follicle, usually caused by Staphylococcus aureus. It forms a red lump and may drain; larger, recurrent, facial, or fever-associated lesions need medical assessment.

\medterm{Furuncle And Carbuncle} A furuncle is a deep infected hair follicle, or boil. A carbuncle is a connected group of boils forming a larger, deeper infection; both may require drainage and sometimes antibiotics.

\medterm{Furuncular Myiasis} A fly-larva infestation forming a boil-like skin lump with a central opening through which the larva breathes. It may cause pain, movement sensations, or drainage and requires safe removal.

\medterm{Furunculosis} Repeated or multiple boils caused by infection of hair follicles, often with Staphylococcus aureus. Risk factors include close contact, skin disease, diabetes, and weakened immunity; treatment may require drainage, antibiotics, and prevention measures.

\medterm{Fusariosis} An opportunistic infection caused by Fusarium fungi, mainly affecting people with severely weakened immunity or prolonged low neutrophil counts. It can cause skin lesions, bloodstream infection, pneumonia, and disease in multiple organs.

\medterm{Fusarium} A group of environmental molds that can cause nail, eye, skin, sinus, lung, or widespread infections. Severe disease occurs mainly in people with weakened immunity, especially prolonged neutropenia.

\medterm{Fusarium Infection} An infection caused by Fusarium molds, which may affect the nails, eyes, skin, sinuses, lungs, or bloodstream. Invasive disease mainly affects people with severe immune suppression and requires urgent specialist antifungal care.
\medterm{Fusarium Solani} An environmental mold that can cause eye, nail, skin, or invasive infections, especially in people with weakened immunity. Its clinical importance depends on the specimen, body site, symptoms, and laboratory confirmation.

\medterm{Fused Placenta} A placenta in which areas that normally form separate lobes are joined. It is often an anatomical variation, but careful examination after delivery helps identify any retained placental tissue or abnormal blood vessels.

\medterm{Fused Teeth} Alternate index label for dental fusion.

\medterm{Fusidic Acid} An antibiotic that blocks bacterial protein production and is used for selected staphylococcal infections. Resistance can develop, so it is often combined with another antibiotic and used according to local guidance.

\medterm{Fusiform} Shaped like a spindle, wider in the middle and narrowing toward both ends. The term may describe a muscle, cell, swelling, or blood-vessel enlargement such as an aneurysm.

\medterm{Fusion} Joining of two structures that are normally separate. It may describe bones, joints, teeth, organs, genes, or tissues and can occur naturally, during development, after injury, or through surgery.

\synonyms
Dental fusion
Tooth fusion

\medterm{Fusion Gene} A gene formed when parts of two separate genes become joined. It may produce an abnormal protein that drives some cancers and can sometimes be identified by genetic testing.

\medterm{Fusion Of The Ear Bones} Abnormal joining or stiffening of the tiny middle-ear bones, which can reduce sound transmission and cause conductive hearing loss. Evaluation identifies the cause and whether surgery or hearing support may help.
\medterm{Fusion Protein} A protein made when parts of two different proteins become joined, usually because their genes have fused. Some fusion proteins drive cancer; others are deliberately engineered for research or treatment.

\medterm{Fusobacterium} A group of bacteria that grow without oxygen and commonly live in the mouth and intestine. Some species can cause dental, throat, abdominal, wound, bloodstream, or other invasive infections.

\medterm{Fusobacterium Infection} Infection caused by Fusobacterium bacteria, which may spread from the mouth or intestine to the throat, lungs, abdomen, skin, joints, or bloodstream. Symptoms and treatment depend on the site and severity.

\medterm{Fusobacterium Necrophorum} A bacterium that can cause severe throat infection and Lemierre syndrome, in which infection spreads to neck veins and may produce infected clots, lung complications, and bloodstream infection.

\medterm{Fusobacterium Nucleatum} A common mouth bacterium associated with gum disease. It can occasionally cause bloodstream or other invasive infection and has also been studied for possible links with intestinal and pregnancy-related disease.
\medterm{Fusospirochetal Gingivitis} A painful gum infection involving fusobacteria and spirochetes, causing bleeding, ulcers, swelling, and foul breath. It is associated with poor oral hygiene, stress, smoking, or illness and needs dental treatment.
\medterm{Fibular Artery} Also called the peroneal artery, it is a branch of the leg’s main artery that supplies the outer and back lower leg and contributes blood flow around the ankle.

\medterm{Fahr's Syndrome} A rare disorder involving calcium deposits in deep brain regions, sometimes causing movement problems, seizures, thinking changes, or psychiatric symptoms. Doctors look for inherited causes and treatable problems with calcium or parathyroid hormones.

\medterm{Fallen Bladder} A common term for bladder prolapse, in which weakened pelvic-support tissues allow the bladder to bulge into the front wall of the vagina. It may cause pressure, leakage, incomplete emptying, or no symptoms.

\medterm{Fatty Liver of Pregnancy (AFLP)} A rare, serious liver disorder usually occurring late in pregnancy. It can cause nausea, abdominal pain, jaundice, low blood sugar, abnormal bleeding, kidney injury, or liver failure and requires urgent delivery planning and specialist care.

\medterm{Fecal Transplant} A procedure placing carefully screened donor stool into the intestine to restore helpful gut microbes. It is mainly used for selected recurrent C. difficile infection and has small but important infection and procedure-related risks.

\medterm{Femoral Anteversion} An inward twist of the thigh bone that turns the knees or feet inward, often noticed in childhood. It commonly improves with growth; persistent pain, severe asymmetry, or walking problems need orthopedic assessment.

\medterm{Fertility Treatment} Medical or procedural care intended to help achieve pregnancy. Options include treating underlying disease, ovulation medicines, surgery, insemination, or in-vitro fertilization, chosen according to cause, age, health, and reproductive goals.

\medterm{Fetal Alcohol Spectrum Disorders (FASD)} A group of lifelong conditions caused by alcohol exposure before birth. Effects may include poor growth, learning or attention problems, behavior changes, coordination difficulties, or characteristic facial features; no amount of alcohol is known to be safe during pregnancy.

\medterm{Fetor Hepaticus} A musty or sweet breath odor that can occur with severe liver dysfunction when certain substances are not adequately cleared. It is a warning sign requiring evaluation for advanced liver disease and complications.

\medterm{Fish Odour Syndrome} A rare metabolic disorder in which trimethylamine builds up and causes a fish-like smell in sweat, breath, or urine. Symptoms may vary with diet; testing, dietary guidance, and emotional support can help.

\medterm{Focal Segmental Glomerular Sclerosis} A kidney disorder in which scars affect some filters or parts of them. It can cause protein in urine, swelling, high blood pressure, and declining kidney function and may result from genetic, immune, or other conditions.

\medterm{FODMAP} An abbreviation for fermentable short-chain carbohydrates that may be poorly absorbed and cause gas, bloating, pain, or diarrhea. A supervised reduction followed by gradual reintroduction may help some people with irritable bowel syndrome.

\medterm{Folliculitis Decalvans} A rare inflammatory scalp disorder causing repeated pustules, crusting, hair loss, and permanent scarring. Early dermatologic treatment aims to control inflammation and infection and limit further follicle destruction.

\medterm{Food Addiction} A nonformal term for intense cravings or loss of control around eating. It is not a universally accepted diagnosis; assessment should consider binge-eating disorder, mood, stress, medicines, sleep, and access to food.

\medterm{Food Contamination} The presence of harmful germs, toxins, chemicals, allergens, or foreign materials in food. Contamination can occur during production, storage, preparation, or serving and may cause infection, poisoning, allergic reactions, or injury.

\medterm{Forensics} The use of scientific methods to answer legal questions, including identification, toxicology, injury analysis, and interpretation of biological or physical evidence. Reliable findings require validated methods, documentation, and expert interpretation.

\medterm{Fowler's Syndrome} A rare cause of persistent difficulty emptying the bladder, usually in young women, related to abnormal relaxation of the urethral sphincter. Evaluation excludes blockage and nerve disease; treatment may require drainage or neuromodulation.

\medterm{Fox-Fordyce Disease} A rare chronic skin disorder caused by blocked and inflamed sweat glands. It produces intensely itchy small bumps, commonly in the armpits, pubic region, or around the nipples.

\medterm{Frontal Lobe Disorder} A condition affecting brain areas involved in planning, judgment, movement, attention, speech, and behavior. Causes include stroke, injury, tumor, dementia, infection, or other neurologic disease.

\medterm{Frontalis Suspension} An operation connecting the upper eyelid to the forehead muscle with a sling to lift a drooping eyelid. It is used when the eyelid-lifting muscle is weak, especially in congenital or childhood ptosis.

\medterm{G6PD Deficiency} An inherited condition in which red blood cells can break down after infections, fava beans, or certain medicines. Episodes may cause anemia, jaundice, fatigue, shortness of breath, or dark urine.

\medterm{Gabapentin} A medicine used for certain seizures and nerve-related pain. It can cause sleepiness, dizziness, swelling, or impaired coordination; the dose often needs adjustment in kidney disease and stopping suddenly can cause problems.

\medterm{Gaba Receptor} A brain-cell protein that responds to GABA, a chemical messenger that usually reduces nerve activity. GABA receptors help regulate sleep, anxiety, muscle control, and seizures and are targets of several medicines.

\medterm{GAD} An abbreviation for generalized anxiety disorder, involving excessive, difficult-to-control worry about several areas of life for at least six months, with symptoms such as restlessness, tension, fatigue, poor concentration, or sleep disturbance.

\synonyms
Generalized Anxiety Disorder

\medterm{Gadolinium} A metal used in a chemically bound form to improve some MRI images. Serious reactions are uncommon, but kidney function, pregnancy, prior reactions, and possible tissue retention are considered before use.

\medterm{Gadolinium-Based Contrast Agent} A substance injected during some MRI scans to make organs, blood vessels, or abnormal tissue easier to see. The agent and dose are selected according to kidney function, pregnancy, allergies, and the scan’s purpose.

\medterm{Gadolinium-Induced Nephrogenic Systemic Fibrosis} A rare illness causing thick, tight, hardened skin and sometimes damage to internal organs after certain MRI contrast exposure, mainly in severe kidney disease. Safer contrast agents have made it very uncommon.

\medterm{Gadopentetate Dimeglumine} An injectable gadolinium contrast agent used to improve selected MRI images. It can cause allergic reactions and is used cautiously or avoided in severe kidney disease because of rare serious complications.

\medterm{Gag Gene} A retroviral gene that encodes structural proteins forming the virus’s inner support, shell, and core. These proteins help assemble new virus particles inside infected cells.

\medterm{Gagging} Retching caused by stimulation of the throat or back of the tongue. It may occur with swallowing difficulty, reflux, infection, examination, aspiration, or distress.

\medterm{Gag Protein} A large retroviral protein that is cut into components forming the virus’s inner support, protective shell, and genetic-material core. These parts are required to assemble new virus particles.

\medterm{Gag Reflex} A protective response in which touching the back of the throat causes throat-muscle contraction and sometimes retching. It helps prevent food or objects from entering the airway.

\synonyms
Pharyngeal Reflex

\medterm{Gain} The amount by which a hearing aid or amplifier increases sound intensity, measured in decibels. Correct gain improves hearing while limiting excessive loudness, feedback, discomfort, and further hearing damage.

\medterm{Gain-Of-Function Mutation} A genetic change that makes a gene or its protein more active, gives it a new activity, or causes activity at the wrong time or place. Its effects range from harmless variation to inherited disease or cancer.

\medterm{Gaisbock Syndrome} A form of relative or stress polycythemia characterized by an elevated hematocrit resulting from decreased plasma volume rather than increased red cell mass, typically occurring in middle-aged, hypertensive, anxious, and overweight individuals.
\synonyms
Stress Polycythemia; Polycythemia Hypertonica

\medterm{Gait} A person’s pattern of walking, including speed, balance, step length, posture, arm movement, and turning. Gait reflects the combined function of the brain, nerves, muscles, joints, vision, and balance system.

\medterm{Gait Abnormality} A change from a person’s usual walking pattern caused by problems with muscles, joints, nerves, balance, vision, heart function, or general health. Assessment examines speed, stability, symmetry, strength, sensation, and pain.

\medterm{Gait Apraxia} Difficulty starting or carrying out walking even though basic muscle strength is adequate. It may occur with disease or injury affecting the front of the brain and can cause short steps, freezing, or instability.

\medterm{Gait Assessment} An examination of how a person walks, including speed, steps, balance, posture, turning, and use of aids. It can reveal problems with strength, sensation, coordination, joints, vision, or thinking.

\medterm{Gait Ataxia} Unsteady, poorly coordinated walking caused by problems in the cerebellum, sensation pathways, balance organs, or other parts of the nervous system.

\medterm{Gait Cycle} The complete sequence of walking movements from one heel strike to the next heel strike of the same foot. It includes the stance phase when the foot bears weight and the swing phase when it moves forward.

\medterm{Gait Disturbance} An abnormal walking pattern involving speed, balance, coordination, strength, or step movement. Causes include neurologic disease, joint or muscle problems, inner-ear disorders, pain, medicines, and systemic illness.

\medterm{Gaits} Patterns of walking. Different gaits may reflect pain, weakness, joint disease, balance problems, or neurologic conditions; assessment considers the person’s usual walking pattern and associated findings.

\medterm{Gait Speed} Walking speed measured over a set distance. It is a simple indicator of mobility and physical reserve, helps estimate fall or disability risk, and can track changes during illness or rehabilitation.

\medterm{Gait Training} Therapy that teaches or retrains safer walking after injury, illness, weakness, or balance problems. It may use exercises, parallel bars, walkers, canes, stairs, or other aids to improve independence.

\medterm{Galactocele} A milk-filled breast cyst caused by blockage of a milk duct, usually during breastfeeding or after weaning. It commonly feels like a painless lump; examination and ultrasound distinguish it from other breast masses.

\medterm{Galactocerebrosidase Deficiency} An inherited lack of an enzyme needed to break down certain fats in nerve tissue. Severe deficiency causes Krabbe disease, which progressively damages the brain and nerves.

\synonyms
Krabbe Disease; Galactosylceramidase Deficiency

\medterm{Galactogogue} A medicine, herb, food, or other substance claimed to increase breast-milk production. Evidence and safety vary, and effects may depend on feeding technique, milk removal, medical causes, and interactions.

\medterm{Galactokinase Deficiency} An autosomal recessive inborn error of galactose metabolism caused by mutations in the \textit{GALK1} gene, which encodes galactokinase. The resulting accumulation of galactose in tissues leads to reduction by aldose reductase into galactitol in the crystalline lens, causing the progressive development of bilateral infantile cataracts without the hepatic, renal, or cognitive damage seen in classic galactosemia.

\synonyms
GALK Deficiency; Type II Galactosemia

\medterm{Galactorrhea} The inappropriate or spontaneous secretion of milk or milk-like fluid from one or both breasts unrelated to normal pregnancy or breastfeeding. It is most commonly caused by hyperprolactinaemia secondary to pituitary prolactinomas, dopamine-antagonist medications (such as antipsychotics and antiemetics), primary hypothyroidism (via elevated thyrotropin-releasing hormone), chronic kidney disease, or local chest wall stimulation.

\synonyms
Galactorrhoea; Non-Puerperal Lactation; Hyperprolactinemic Discharge

\medterm{Galactorrhoea} Alternate name; see \textit{Galactorrhea}.

\medterm{Galactosaemia} A rare inherited disorder in which the body cannot properly process galactose, a sugar in milk. Untreated infants may develop feeding problems, jaundice, liver injury, cataracts, serious infection, or developmental problems.

\medterm{Galactose} A simple sugar that forms part of lactose in milk. The body normally processes it for energy, but inherited galactosemia can make milk and other galactose-containing foods dangerous for infants.

\medterm{Galactose-1-Phosphate Uridyltransferase Blood Test} A blood test measuring an enzyme needed to process galactose. It helps screen for or diagnose classic galactosemia, especially in newborns or infants with symptoms after milk feeding.

\medterm{Galactose Metabolism} The body’s process for changing galactose, part of milk sugar, into energy or other useful substances. Inherited enzyme problems can block this process and cause galactosemia.

\medterm{Galactosemia} An inherited disorder in which the body cannot properly process galactose. Classic disease can cause vomiting, jaundice, liver injury, cataracts, serious infection, and developmental problems after milk feeding.

\medterm{Galactosialidosis} A rare inherited storage disorder in which cells accumulate certain sugar-containing substances. It can cause developmental delay, muscle weakness, bone changes, kidney disease, and eye problems, with severity varying widely.

\medterm{Galactoside} A compound in which the sugar galactose is joined to another molecule, such as another sugar, a fat, or a protein. Galactosides occur in cells and in some foods.

\medterm{Galactosyltransferase} An enzyme that attaches galactose to proteins, fats, or other sugars while cells build complex molecules. These molecules support cell structure, recognition, signaling, and nerve insulation.

\medterm{Galant Reflex} A newborn reflex in which stroking one side of the back causes the infant’s trunk to curve toward that side. It normally fades during early infancy; marked asymmetry or persistence can occur with nervous-system problems.

\medterm{Galeazzi Fracture} A fracture of the middle to distal third of the radial shaft associated with dislocation or subluxation of the distal radioulnar joint (DRUJ), typically requiring anatomic open reduction and internal fixation to restore forearm pronation and supination.
\synonyms
Galeazzi Fracture-Dislocation; Piedmont Fracture

\medterm{Galeazzi Fracture-Dislocation} A forearm injury involving a break in the lower part of the radius and dislocation of the nearby joint between the radius and ulna. It can cause wrist pain, instability, and difficulty turning the forearm.

\synonyms
Galeazzi Fracture; Reverse Monteggia Fracture

\medterm{Galeazzi Sign} A finding in an infant lying with hips and knees bent, in which one knee is lower than the other. It may suggest hip dislocation, developmental hip dysplasia, or unequal thigh-bone length.

\synonyms
Allis Sign; Galeazzi Test

\medterm{Gallavardin Phenomenon} A finding in severe aortic-valve narrowing where the heart murmur sounds harsh near the upper chest but musical near the apex. It reflects different sound frequencies transmitted through the chest.

\medterm{Gall-Bladder} A small pear-shaped sac beneath the liver that stores and concentrates bile. Bile is released into the small intestine to help digest fats.

\medterm{Gallbladder Absence} Having no gallbladder because it did not develop or was surgically removed. Bile then flows directly from the liver into the intestine; most people can digest normally, though some develop digestive symptoms.

\medterm{Gallbladder Adenomyomatosis} A benign, non-inflammatory hyperplastic alteration of the gallbladder wall characterized by excessive epithelial proliferation with smooth muscle hypertrophy and deep, pouch-like invaginations of the luminal mucosa extending into the muscularis propria (termed Rokitansky-Aschoff sinuses). Often identified on ultrasound by intramural diverticula containing cholesterol crystals that produce characteristic V-shaped \"comet-tail\" reverberation artifacts.

\synonyms
Adenomyoma of Gallbladder; Diverticular Gallbladder Disease

\medterm{Gallbladder Cancer} A cancer arising in the gallbladder, sometimes associated with long-term inflammation or stones. It may cause abdominal pain, jaundice, weight loss, or blocked bile flow, but can also be found unexpectedly after surgery.

\synonyms
Gallbladder Carcinoma; Biliary Malignancy

\medterm{Gallbladder Carcinoma} A cancer of the gallbladder, most often arising from its lining cells. It may be found incidentally or cause right-upper-abdominal pain, jaundice, weight loss, or bile-duct blockage; treatment depends on its stage and spread.

\medterm{Gall Bladder Disease} Any disorder of the gallbladder, including stones, inflammation, infection, polyps, abnormal movement, or cancer. Symptoms may include right-upper-abdominal pain, nausea, vomiting, fever, or jaundice.

\medterm{Gallbladder Diseases} Conditions affecting the gallbladder, including stones, inflammation, infection, polyps, abnormal emptying, and cancer. Symptoms and treatment depend on the exact disorder and whether bile flow is blocked.

\medterm{Gallbladder Duplication} A rare birth variation in which two gallbladders develop, with separate or shared drainage ducts. It may cause no symptoms but can be associated with stones or inflammation and is important to identify before surgery.

\medterm{Gallbladder Hydrops} Severe, non-inflammatory distension and dilation of the gallbladder filled with clear, sterile, watery or mucoid transudate resulting from chronic, complete obstruction of the cystic duct (most often by an impacted gallstone in Hartmann's pouch) in the absence of acute infection.
\synonyms
Mucocele of Gallbladder; Hydrops Vesicae Felleae

\medterm{Gallbladder Infection} Inflammation or infection of the gallbladder, often after a stone blocks bile flow. Right-upper-abdominal pain, fever, nausea, vomiting, and tenderness may occur; treatment can include antibiotics, drainage, or surgery.

\medterm{Gallbladder Inflammation} Inflammation of the gallbladder, most often caused by a gallstone blocking its drainage. It can cause persistent right-upper-abdominal pain, fever, nausea, and vomiting.

\synonyms
Cholecystitis

\medterm{Gallbladder Mucocele} Enlargement of the gallbladder with mucus or watery fluid because its outlet is blocked, often by an impacted stone. It may cause pain or infection; drainage or removal may be used.

\medterm{Gallbladder Neck} The narrow, tapered S-shaped constriction of the gallbladder that connects the gallbladder body to the cystic duct, characterized internally by spiral mucosal folds (valves of Heister); an impacted gallstone here leads to acute cholecystitis or Mirizzi syndrome.

\synonyms
Collum Vesicae Biliaris

\medterm{Gallbladder Necrosis} Death of gallbladder tissue caused by severe inflammation or inadequate blood supply. It can lead to perforation, abdominal infection, and sepsis.

\medterm{Gallbladder Pain} Pain from temporary blockage of the gallbladder outlet, usually by a stone. It is often steady, severe, and felt under the right ribs after meals and may include nausea or vomiting.

\synonyms
Biliary Colic; Gallstone Pain

\medterm{Gallbladder Perforation} A hole through the gallbladder wall that allows bile or infected contents into the abdomen. It can cause severe inflammation, abscess, peritonitis, or sepsis.

\medterm{Gallbladder Polyp} A mucosal projection extending into the lumen of the gallbladder, categorized as pseudopolyps (cholesterol polyps, inflammatory polyps, adenomyomatosis) or true epithelial neoplasms (adenomas); polyps larger than 10 mm carry an increased malignancy risk warranting cholecystectomy.

\medterm{Gallbladder Radionuclide Scan} A nuclear medicine test that follows a tracer through the liver, gallbladder, and bile ducts. It helps assess suspected gallbladder inflammation, blockage, or abnormal bile flow.

\medterm{Gallbladder Volvulus} Twisting of the gallbladder around its supporting tissues and duct, cutting off blood flow. It causes sudden severe right-upper-abdominal pain, nausea, vomiting, and abdominal irritation.

\medterm{Gallium} A metallic element used in some radioactive tracers, medical compounds, and technology. Its medical effects depend on the isotope or compound and how it is administered and distributed in the body.

\medterm{Gallium-68} A radioactive isotope used to label tracers for PET scans. Linked to a suitable targeting molecule, it can help locate some neuroendocrine tumors or prostate-cancer lesions; it is used for diagnosis, not direct treatment.

\medterm{Gallium Nitrate} A gallium-containing compound formerly used in selected cancer settings to reduce high calcium levels. Its use is limited because it can damage the kidneys and has largely been replaced by other treatments.

\medterm{Gallium Scan} A nuclear medicine scan using radioactive gallium to locate some infections, inflammation, or tumors. It is used less often now because newer imaging tests may provide more specific or faster information.

\medterm{Gallop Rhythm} An extra heart sound creating a three- or four-beat rhythm. It may reflect abnormal heart filling, stiff heart muscle, or heart failure and is interpreted with symptoms and other examination findings.

\medterm{Gallstone Ileus} A blockage of the intestine caused by a large gallstone that has entered through an abnormal connection from the gallbladder. It can cause pain, vomiting, swelling, and inability to pass stool or gas.

\medterm{Gallstone Pancreatitis} Sudden inflammation of the pancreas caused by a gallstone blocking the shared opening of the bile and pancreatic ducts. It usually causes severe persistent upper-abdominal pain and may be associated with complications or remaining duct stones.

\medterm{Gallstones} Solid deposits that form in the gallbladder when bile contains too much cholesterol or bilirubin, too little bile salt, or remains in the gallbladder without emptying fully. They may cause no symptoms or block bile flow, producing severe upper-abdominal pain, inflammation, jaundice, infection, or pancreatitis.

\synonyms
Cholelithiasis; Biliary Calculi

\medterm{Galvanocautery} Use of heat generated by electric current to destroy selected tissue or control bleeding during a procedure. It can cause burns or damage nearby structures if not carefully controlled.

\medterm{Gambling Disorder} A persistent inability to control gambling despite financial, relationship, work, health, or emotional harm. The disorder can cause debt, secrecy, distress, and suicide risk and benefits from professional behavioral and mental-health support.

\synonyms
Pathological Gambling; Compulsive Gambling; Ludomania

\medterm{Gamete} A mature reproductive cell, such as a sperm or egg, containing one set of chromosomes. Two gametes combine during fertilization to form a cell with genetic material from both biological parents.

\medterm{Gamete Intrafallopian Transfer} An assisted-reproduction procedure in which eggs and sperm are placed into a fallopian tube so fertilization may occur inside the body. It is now used less often than in-vitro fertilization.

\medterm{Gametogenesis} The formation and maturation of sperm or eggs in the testes or ovaries. It includes cell division that reduces chromosome number so gametes can combine during fertilization.

\medterm{Gamma-Aminobutyric Acid} A chemical messenger that usually reduces nerve activity in the brain and spinal cord. It helps regulate sleep, anxiety, muscle activity, and seizures through two main receptor types.

\synonyms
GABA

\medterm{Gamma-Benzene Hexachloride} An insecticide, also called lindane, that can cause nervous-system toxicity and seizures after excessive exposure. Medical use is restricted in many places because safer treatments are available.

\medterm{Gamma Camera} A device used in some nuclear-medicine scans to detect radiation from a small tracer inside the body and create pictures showing an organ’s structure or activity.

\medterm{Gamma-Chemokines} Signaling proteins that attract or activate immune cells through specific receptors. They help coordinate inflammation, immune surveillance, and movement of defensive cells into affected tissue.

\medterm{Gamma Globulin} The part of blood plasma containing most immunoglobulins, or antibodies, which help defend against infections and can be measured or given as immune treatment.

\medterm{Gamma-Globulin} A group of blood proteins that includes antibodies. The term may refer to antibody-containing plasma fractions or immunoglobulin preparations used to provide temporary immune protection or treatment.

\medterm{Gamma-Glutamyl Transferase Blood Test} A blood test measuring an enzyme found mainly in the liver and bile ducts. It can support evaluation of liver or bile-duct disease and help explain a high alkaline-phosphatase result.

\medterm{Gammaherpesvirinae} A subgroup of herpesviruses that establish long-lasting infection in lymphoid or related cells. Some members, including Epstein-Barr virus, can cause human disease and contribute to certain cancers.

\medterm{Gamma Hydroxybutyrate} A central-nervous-system depressant that can cause drowsiness, confusion, slowed breathing, coma, or death in overdose. It is also misused recreationally and can cause dependence and withdrawal.

\medterm{Gamma-Interferon} A cytokine made by activated immune cells that strengthens macrophage activity and cell-mediated defense against certain infections. It also contributes to inflammation and immune regulation.

\medterm{Gamma Knife} A noninvasive radiation treatment that aims many focused beams at selected brain tumors or other lesions. It avoids an open operation but can cause swelling, bleeding, or injury to nearby brain tissue.

\medterm{Gamma Knife Radiosurgery} Gamma Knife is a focused radiation system that delivers multiple precisely aimed beams, mainly to selected brain tumors and vascular or other brain lesions.

\medterm{Gamma Rays} High-energy electromagnetic radiation produced by radioactive decay or nuclear reactions. It can damage cells and is used in controlled amounts for some cancer treatments and medical sterilization.

\medterm{Gamma-Ray Therapy} Treatment that uses controlled gamma radiation to damage the DNA of selected abnormal cells. It can treat some cancers, but nearby healthy tissue may also be affected and side effects depend on the treated area.

\medterm{Gamna-Gandy Bodies} Small, firm, yellow-brown or reddish-brown siderofibrotic nodules in the splenic red pulp containing deposits of hemosiderin, iron, and calcium encrusting connective tissue and elastic fibers, resulting from organized focal hemorrhages in chronic congestive splenomegaly and sickle cell disease.
\synonyms
Gamna-Gandy Nodules; Siderofibrotic Bodies

\medterm{Ganciclovir} An antiviral medicine that interferes with viral DNA production and is used mainly for serious cytomegalovirus infection. It can suppress bone marrow, affect the kidneys, and cause other serious toxicity.

\medterm{Ganglia} Groups of nerve-cell bodies outside the brain and spinal cord. They relay and organize signals between the central nervous system and the rest of the body.

\medterm{Ganglioglioma} A usually slow-growing tumor containing nerve cells and supporting brain cells. It often causes seizures and may occur in the brain or spinal cord; treatment commonly involves surgical removal.

\medterm{Ganglion} A cluster of nerve-cell bodies outside the brain and spinal cord. It acts as a relay point for signals between the central nervous system and body tissues.

\medterm{Ganglion Cell} A nerve cell located in a ganglion. Retinal ganglion cells collect visual signals and send them through the optic nerve to the brain.

\medterm{Ganglion Cyst} A noncancerous, fluid-filled lump near a joint or tendon, most often at the wrist or hand. It may be painless or cause pressure, pain, numbness, or limited movement.

\synonyms
Synovial Cyst; Bible Cyst

\medterm{Ganglionectomy} Surgical removal of a ganglion, a cluster of nerve-cell bodies, or occasionally a ganglion cyst. The purpose, benefits, and risks depend on the exact structure and body site being treated.

\medterm{Ganglioneuroblastoma} A tumor made of both mature and immature nerve-related cells, usually in the adrenal or tissue beside the spine. It mainly affects children and may cause a mass, pain, high blood pressure, or other symptoms.

\medterm{Ganglioneuroma} A usually noncancerous tumor arising from nerve tissue of the autonomic nervous system. It often causes no symptoms but may produce a mass, pain, or hormone-related effects depending on its location.

\medterm{Ganglioneuromatosis} A condition in which nerve-related ganglion tissue grows diffusely or in multiple areas. Effects depend on the organs involved and may include masses, pain, bowel symptoms, or hormone-related changes.

\medterm{Ganglionic Blockers} Medicines that interrupt nerve signals at autonomic ganglia, affecting blood pressure, heart rate, digestion, and other automatic functions. They are rarely used because of widespread side effects.

\medterm{Ganglioside} A fat-containing molecule in cell membranes, especially nerve cells. Gangliosides help cells communicate and recognize one another; abnormal accumulation occurs in several inherited storage disorders.

\medterm{Gangliosidosis} A group of inherited disorders in which gangliosides accumulate inside cells because a breakdown enzyme is missing. The buildup often damages the nervous system and may also affect the eyes, bones, organs, and development.

\medterm{Gangrene} Death of body tissue caused by severely reduced blood flow, infection, injury, or a combination of these factors. It may affect skin, muscle, or an internal organ. The area may become discolored, cold, swollen, painful, numb, or foul-smelling; fever, confusion, or feeling very ill may occur with spreading infection.

\medterm{Ganser Syndrome} A rare dissociative or psychiatric syndrome characterized by approximate answers, confusion, altered awareness, and sometimes hallucinations or neurologic-like symptoms. It may occur with severe stress, trauma, other mental disorders, or medical and substance-related conditions.

\medterm{Gap Junction} A tiny connection between neighboring cells that allows ions and small molecules to pass directly between them. Gap junctions help coordinate heartbeats, nerve activity, and tissue growth.

\medterm{Gardasil} A brand of vaccine that protects against infection with selected human papillomavirus (HPV) types. It helps prevent HPV-related cancers and genital warts and does not treat an existing HPV infection.

\medterm{Gardnerella} A group of bacteria that may live in the genital tract. Overgrowth of Gardnerella, especially G. vaginalis, can contribute to bacterial vaginosis with discharge, odor, and altered vaginal bacteria.

\medterm{Gardnerella Vaginalis} A bacterium associated with bacterial vaginosis, in which normal vaginal bacteria become imbalanced. It may be present without illness; when symptomatic, discharge, odor, irritation, or burning can occur.

\medterm{Gardner Syndrome} An inherited condition caused by an APC gene change, producing many colon polyps and increasing colorectal-cancer risk. It may also cause bone growths, dental abnormalities, skin cysts, and desmoid tumors.

\medterm{Gargle} To rinse the mouth and throat with liquid while making a bubbling sound. Gargling may soothe irritation or help deliver a prescribed treatment.

\medterm{Garré's Osteomyelitis} A chronic inflammatory reaction in the jawbone, often near an infected tooth, causing firm swelling. Dental infection, persistent swelling, or pain may be present.

\medterm{Gartner Duct} A remnant of an embryonic duct in the female reproductive tract. It can persist harmlessly or form a cyst along the side wall of the vagina or broad ligament.

\medterm{Gartner Duct Cyst} A usually harmless fluid-filled cyst arising from an embryonic duct remnant along the side wall of the vagina. It often causes no symptoms, but a large cyst may cause pressure, pain, or difficulty with intercourse.

\medterm{Gas} A state of matter without a fixed shape or volume. In medicine, gas may mean air or another gas in the digestive tract, blood, tissues, or a body cavity.

\medterm{Gas And Gas Pains} Fullness, bloating, cramps, or discomfort from swallowed air or gas made by intestinal bacteria. Foods, constipation, food intolerance, and digestive disorders can contribute.

\medterm{Gas Chromatography} A laboratory method that separates vaporized chemicals so they can be identified or measured. It is often combined with mass spectrometry to test medicines, poisons, pollutants, or biological samples.

\medterm{Gas Exchange} Movement of oxygen from inhaled air into blood and carbon dioxide from blood into exhaled air. Oxygen then moves to tissues while carbon dioxide returns to the lungs for removal.

\medterm{Gas –  Flatulence} Passage of gas from the digestive tract through the anus, caused by swallowed air and bacterial digestion of food. It is normal, but excess gas with bloating or pain may reflect diet, constipation, or digestive disease.

\synonyms
Flatulence; Flatus Expulsion

\medterm{Gas Gangrene} A life-threatening infection in which bacteria produce toxins and gas inside damaged tissue, rapidly destroying muscle. Severe pain, swelling, blisters, discoloration, or a crackling feeling may occur.

\medterm{Gaskin Maneuver} A maneuver for shoulder dystocia in which the birthing person moves onto hands and knees. The position can change pelvic dimensions and help release the baby’s shoulder.

\synonyms
All-Fours Position; Hands-and-Knees Posture

\medterm{Gasoline Poisoning} Harm from swallowing, inhaling, or getting gasoline on the skin. It can cause nausea, dizziness, drowsiness, seizures, abnormal heart rhythms, or chemical lung injury, especially if vomited and inhaled.

\synonyms
Petrol Poisoning; Hydrocarbon Toxicity

\medterm{Gas Poisoning} Illness caused by inhaling or otherwise being exposed to a harmful gas. Effects range from irritation and headache to oxygen deprivation, lung injury, confusion, or collapse.

\medterm{Gas Scavenger} A substance or device that removes unwanted gases from a fluid, container, or breathing system. In healthcare it may help prevent toxic exposure or maintain safe gas delivery.

\medterm{Gasserian Ganglion} A cluster of nerve cells near the base of the skull that carries sensation from the face and controls some chewing muscles. It is involved in trigeminal neuralgia and other facial-nerve disorders.

\medterm{Gastralgia} Pain in the upper abdomen or stomach area. Causes include indigestion, reflux, ulcers, gallbladder disease, pancreatic inflammation, infection, or other conditions; severe pain may occur with bleeding, vomiting, fever, or weight loss.

\medterm{Gas Transport} Movement of oxygen from the lungs to body tissues and carbon dioxide back to the lungs through the blood. Hemoglobin carries most oxygen and part of the carbon dioxide.

\medterm{Gastrectomy} Surgery removing part or all of the stomach, usually for cancer or severe disease. Food passage is reconstructed, and possible effects include early fullness, weight loss, nutritional deficiencies, reflux, or dumping symptoms.

\medterm{Gastric Acid} Strong acid made by the stomach lining that helps break down food, activate digestive enzymes, and kill many swallowed microbes. Excess can contribute to reflux or ulcers; too little can impair digestion and nutrient absorption.

\medterm{Gastric Adenocarcinoma} The most common stomach cancer, arising from cells lining the stomach. It may cause indigestion, pain, early fullness, vomiting, bleeding, or weight loss; tissue examination and staging identify its type and extent.

\medterm{Gastric Analysis} Examination of stomach contents for volume, acidity, or selected substances. It is used only in specific situations because endoscopy, imaging, and blood tests are more useful for many stomach disorders.

\medterm{Gastric Antral Vascular Ectasia} Enlarged fragile blood vessels in the lower stomach that can cause slow internal bleeding and iron-deficiency anemia. It may appear as red stripes or spots during endoscopy and is associated with some liver and autoimmune diseases.

\synonyms
GAVE; Watermelon Stomach

\medterm{Gastric Antrum} The distal, non-acid-secreting muscular portion of the stomach located between the incisura angularis and the pylorus, containing G cells that secrete gastrin and D cells that secrete somatostatin, serving as the site of mechanical food grinding and triturition.

\medterm{Gastric Balloon} A temporary balloon placed inside the stomach to create fullness and support weight loss. Nausea, vomiting, reflux, deflation, blockage, bleeding, or perforation can occur.

\medterm{Gastric Bezoar} A mass of material that cannot be digested and collects in the stomach. It may contain plant fiber, hair, or medicines and can cause fullness, pain, vomiting, blockage, or poor nutrition.

\medterm{Gastric Biopsy} Removal of a small sample of stomach lining, usually during endoscopy, for laboratory examination. It can help assess inflammation, infection, precancerous changes, or cancer.

\medterm{Gastric Bypass} An operation that makes the stomach smaller and connects it to a lower part of the small intestine. It limits food intake and nutrient absorption and can improve severe obesity and related diseases.

\medterm{Gastric Bypass Surgery} Surgery that creates a small stomach pouch and reroutes food to the small intestine. It produces substantial weight loss and may improve diabetes, high blood pressure, sleep apnea, and other obesity-related conditions.

\synonyms
Roux-en-Y Gastric Bypass; Bariatric Bypass

\medterm{Gastric Cancer} A cancer that begins in the stomach lining or other stomach tissue. It may cause indigestion, pain, early fullness, vomiting, weight loss, or bleeding, although early disease may cause few symptoms.

\medterm{Gastric Body} The main middle part of the stomach between the fundus and the antrum. It stores and mixes food and contains glands that release acid and digestive substances.

\medterm{Gastric Carcinoid Tumor} A neuroendocrine tumor arising in the stomach, often from cells that make digestive hormones. Some are linked to chronic gastritis and high gastrin levels; behavior ranges from slow-growing to aggressive.

\medterm{Gastric Cardia} The anatomical transition zone of the stomach immediately surrounding the gastroesophageal junction where the tubular esophagus empties into the stomach, lined by simple columnar mucous-secreting glands without parietal or chief cells.

\medterm{Gastric Culture} A laboratory test that grows bacteria from a stomach sample, usually collected during endoscopy. It can identify Helicobacter pylori and help select antibiotics, although other tests are more commonly used.

\synonyms
Helicobacter Pylori Culture; Gastric Biopsy Culture

\medterm{Gastric Dysplasia} Abnormal cells in the stomach lining that have some features of cancer but have not invaded deeper tissue. Low- or high-grade changes may increase the risk of stomach cancer.

\medterm{Gastric Emptying} The passage of food and liquid from the stomach into the small intestine through coordinated stomach contractions and opening of the pylorus. Abnormally slow or rapid emptying can cause digestive symptoms.

\medterm{Gastric Emptying Study} A test that measures how quickly food leaves the stomach. The patient eats a meal containing a small tracer, and a scanner records the remaining food over several hours to detect unusually slow or rapid emptying.

\medterm{Gastric Fistula} An abnormal passage between the stomach and another organ or the skin. Stomach contents may leak through it, causing drainage, infection, skin irritation, dehydration, or poor nutrition.

\medterm{Gastric Fundus} The upper, rounded part of the stomach, located above the main stomach body and near the diaphragm. It can expand to store food and contains glands that produce digestive secretions.

\medterm{Gastric Glands} Glands in the stomach lining that make acid, digestive enzymes, mucus, intrinsic factor, and hormones.

\medterm{Gastric Hemorrhage} Bleeding from the stomach, commonly caused by an ulcer, severe inflammation, a tumor, or blood-vessel damage. It may cause vomiting blood, black stools, dizziness, anemia, or shock.

\medterm{Gastric Juice} Fluid made by the stomach containing acid, digestive enzymes, mucus, water, salts, and intrinsic factor. It helps break down food, protect the stomach lining, and absorb vitamin B12 later in digestion.

\medterm{Gastric Lavage} A procedure that flushes the stomach through a tube passed through the mouth or nose. It is rarely used because it can cause aspiration or injury and is reserved for selected emergencies.

\medterm{Gastric Leiomyosarcoma} A rare cancer of smooth muscle in the stomach wall. It may cause pain, bleeding, vomiting, or a mass and is evaluated with imaging and biopsy; treatment usually involves surgery when possible.

\medterm{Gastric Lymphoma} A cancer of immune-system cells that starts in the stomach. Some types are linked to Helicobacter pylori infection and may improve after eradication treatment; others are treated with chemotherapy, immunotherapy, or surgery.

\medterm{Gastric Malignancy} Cancer that begins in the stomach, most often in the stomach lining. It may cause indigestion, early fullness, pain, vomiting, weight loss, anemia, or bleeding, though early disease may have no symptoms; endoscopy with biopsy confirms the type and scans assess spread.

\medterm{Gastric Motility} The coordinated muscle activity that mixes stomach contents and moves them toward the small intestine.

\medterm{Gastric Mucosa} The moist inner lining of the stomach. It contains glands that produce acid, digestive enzymes, mucus, and intrinsic factor and is protected by a mucus barrier; inflammation or injury can cause pain or bleeding.

\medterm{Gastric Muscularis} The muscle layers in the stomach wall that mix food and move it toward the pylorus. The stomach has an extra inner muscle layer that helps with grinding.

\medterm{Gastric Necrosis} Death of stomach tissue caused by severely reduced blood flow, corrosive injury, extreme stretching, infection, or another serious insult. It can lead to perforation, severe infection, and shock.

\medterm{Gastric Obstruction} Partial or complete blockage that prevents food and liquid from leaving the stomach normally. It may cause repeated vomiting, upper-abdominal pain, fullness, swelling, dehydration, and weight loss.

\medterm{Gastric Outlet Obstruction} Blockage where the stomach empties into the small intestine, usually at the pylorus or duodenum. Causes include ulcer scarring, tumors, inflammation, and infantile pyloric stenosis; vomiting can cause dehydration.

\medterm{Gastric Pneumatosis} Gas within the stomach wall. It may result from pressure-related injury and be relatively mild, or from severe infection, which can destroy tissue and cause sepsis.

\medterm{Gastric Polyps} Growths projecting from the stomach lining. Most are harmless, but some types are linked to chronic inflammation, inherited conditions, or long-term acid-suppressing treatment and may contain precancerous changes.

\medterm{Gastric Perforation} A hole through the stomach wall that allows stomach contents to leak into the abdomen. It can cause severe inflammation, infection, and shock.

\medterm{Gastric Rugae} Folds in the inner lining of the empty stomach that flatten as the stomach fills. They allow the stomach to expand and increase the lining's surface area.

\medterm{Gastric Secretion} The release of acid, enzymes, mucus, intrinsic factor, and hormones by the stomach lining.

\medterm{Gastric Serosa} The thin outer covering of the stomach, formed by the visceral peritoneum. It reduces friction as the stomach moves inside the abdomen.

\medterm{Gastric Stenosis} Abnormal narrowing of the stomach or its outlet that slows or blocks emptying. It may result from birth defects, ulcer scarring, inflammation, tumors, or thickening of the outlet muscle.

\medterm{Gastric Submucosa} The connective-tissue layer beneath the stomach lining. It contains blood vessels, lymph vessels, nerves, and supporting tissue.

\medterm{Gastric Suction} Removal of stomach contents through a tube passed through the nose or mouth. It may relieve severe distension or vomiting from obstruction and is used only when clinically indicated because complications can occur.

\synonyms
Nasogastric Aspiration; Stomach Pumping

\medterm{Gastric Surgery} An operation on the stomach to remove disease, repair damage, reduce its size, or bypass a blockage. The exact procedure and risks depend on the condition and may include bleeding, infection, leakage, or altered digestion.

\medterm{Gastric Tissue Biopsy And Culture} Removal of small stomach-lining samples during endoscopy for microscopic examination and, when needed, laboratory testing for infection. It can investigate inflammation, Helicobacter pylori, precancerous changes, or tumors.

\synonyms
Endoscopic Gastric Biopsy; Mucosal Culture

\medterm{Gastric Ulcer} An open sore in the stomach lining, commonly caused by Helicobacter pylori infection or anti-inflammatory medicines. It may cause burning upper-abdominal pain, bleeding, anemia, or perforation.

\medterm{Gastric Varices} Enlarged, fragile veins in the stomach caused by increased pressure in the portal circulation, often from advanced liver disease. Rupture can cause sudden, life-threatening vomiting of blood.

\medterm{Gastric Volvulus} Twisting of the stomach that can block food passage and cut off blood supply. It may cause sudden severe upper-abdominal pain, retching without vomiting, bloating, and inability to pass a tube.

\medterm{Gastric Wall} The layered wall of the stomach, made up of the mucosa, submucosa, muscle layers, and outer covering. It allows the stomach to store, mix, and move food.

\medterm{Gastrin} A hormone made mainly in the lower stomach and upper small intestine. It stimulates stomach acid production and helps maintain the stomach lining; levels rise after eating and may become excessive in certain diseases.

\medterm{Gastrin Blood Test} A blood test measuring the hormone gastrin. It helps investigate excessive stomach acid, recurrent ulcers, or loss of acid-producing cells. Acid-suppressing medicines and fasting can significantly affect the result.

\synonyms
Serum Gastrin Level; Fasting Gastrin Assay

\medterm{Gastrinoma} A tumor that makes too much gastrin, a hormone that causes the stomach to produce acid. It usually starts in the pancreas or duodenum and may cause Zollinger--Ellison syndrome, repeated peptic ulcers, diarrhea, and acid reflux. Some cases are linked to MEN1 and may spread to the liver or nearby lymph nodes.

\medterm{Gastrin-Releasing Peptide} A regulatory neuropeptide released by postganglionic parasympathetic vagal nerve endings in the gastric antrum that binds to receptors on G cells, stimulating gastrin release and subsequent gastric acid secretion.

\synonyms
GRP; Mammalian Bombesin

\medterm{Gastrin Secretion} Release of the hormone gastrin, mainly from cells in the lower stomach after food enters. Gastrin increases stomach acid and digestive-enzyme release and supports growth and function of the stomach lining.

\medterm{Gastritis} Inflammation or irritation of the stomach lining. Common causes include \textit{H. pylori} infection, anti-inflammatory medicines, alcohol, autoimmune disease, and severe illness. It may cause upper-abdominal pain, nausea, vomiting, or bleeding.

\synonyms
Type A Gastritis; Autoimmune Metaplastic Atrophic Gastritis

\medterm{Gastrocnemius} The large muscle at the back of the calf. It helps point the foot downward and bend the knee during standing, walking, running, and jumping; injury can cause calf pain, swelling, or weakness.

\medterm{Gastrocnemius Muscle} The large superficial calf muscle that joins the Achilles tendon. It points the foot downward and helps bend the knee and push the body forward during walking, running, and jumping.

\medterm{Gastroduodenal Artery} A major arterial branch arising from the common hepatic artery, descending posterior to the first part of the duodenum before dividing into the right gastro-omental and superior pancreaticoduodenal arteries; erosion by a posterior duodenal ulcer produces massive upper gastrointestinal hemorrhage.

\synonyms
GDA

\medterm{Gastroduodenal Crohn Disease} Crohn disease affecting the stomach or first part of the small intestine. It may cause upper-abdominal pain, nausea, vomiting, weight loss, ulcers, narrowing, or blockage; treatment varies with symptoms and disease severity.

\medterm{Gastroduodenostomy} An operation creating a new opening between the stomach and duodenum so food can pass directly into the small intestine. It may bypass a blockage or restore passage after part of the stomach is removed.

\medterm{Gastroenteritis} Infection or inflammation of the stomach and intestines, commonly caused by viruses and sometimes by bacteria, parasites, or toxins. It usually causes diarrhea, vomiting, cramps, and fever and can lead to dehydration.

\medterm{Gastroenteritis, Acute} Alternate name; see \textit{Stomach Flu}.

\medterm{Gastroenteritis Bacterial} Inflammation of the stomach or intestines caused by bacteria, usually after contaminated food or water. It can cause diarrhea, cramps, fever, vomiting, and dehydration; blood in stool or severe illness may occur.

\medterm{Gastroenteritis Viral} Alternate name; see \textit{Viral gastroenteritis}.

\medterm{Gastroenterologist} A physician specializing in disorders of the digestive tract, liver, gallbladder, bile ducts, and pancreas. They diagnose and treat digestive problems using medicines, endoscopy, procedures, lifestyle care, and surgery referrals.

\medterm{Gastroenterology} The medical specialty concerned with the digestive tract, liver, gallbladder, bile ducts, and pancreas. Gastroenterologists diagnose and treat these conditions with medicines, endoscopy, procedures, and referrals for surgery.

\medterm{Gastroenterostomy} Surgical creation of an opening between the stomach and a segment of the small intestine, usually the jejunum, to bypass or replace a blocked or diseased gastric outlet.

\medterm{Gastroepiploic Artery} The gastroepiploic arteries run along the greater curve of the stomach and supply blood to the stomach and nearby omentum.

\medterm{Gastroesophageal} Relating to both the stomach and esophagus, especially their joining point or the backward movement of stomach contents into the esophagus.

\medterm{Gastroesophageal Junction Adenocarcinoma} A gland-forming cancer arising where the esophagus joins the stomach. It may cause difficulty swallowing, reflux, weight loss, anemia, or bleeding; endoscopy, biopsy, and scans determine treatment and extent.

\medterm{Gastroesophageal Reflux} Backward flow of stomach contents into the esophagus. It may cause heartburn, sour fluid in the throat, cough, chest discomfort, or no symptoms; frequent reflux can injure the esophageal lining.

\medterm{Gastroesophageal Reflux Disease} A long-term condition in which stomach contents repeatedly flow back into the esophagus, causing symptoms such as heartburn, regurgitation, chest pain, cough, or difficulty swallowing and sometimes damaging the esophageal lining.

\synonyms
GERD; Acid Reflux Disease

\medterm{Gastrogastric Intussusception} A rare condition in which one part of the stomach folds into an adjoining part, sometimes because of a polyp or tumor. It may cause pain, vomiting, blockage, or bleeding.

\medterm{Gastrografin Study} An X-ray study using a water-soluble contrast liquid to show the digestive tract. It may be chosen when a leak or perforation is possible, because the contrast is safer outside the tract than barium.

\medterm{Gastrointestinal} Relating to the stomach and intestines, or to the complete digestive tract from the mouth to the anus. These organs digest food, absorb nutrients and water, move contents forward, and remove waste.

\medterm{Gastrointestinal Agents} Medicines or other substances that act on the digestive tract to relieve or treat problems such as nausea, vomiting, diarrhea, constipation, ulcers, reflux, or intestinal spasms. Their effects and risks depend on the specific agent.

\medterm{Gastrointestinal Amyloidosis} Buildup of abnormal protein deposits in the digestive tract. It can impair movement, absorption, or bleeding control and may occur with disease affecting other organs such as the heart, kidneys, or nerves.

\medterm{Gastrointestinal Bleeding} Bleeding anywhere in the digestive tract. It may be acute or chronic and may come from the upper digestive tract, small intestine, or lower digestive tract. It can cause vomiting blood, black stools, or red blood in the stool.

\medterm{Gastrointestinal Bleeding Scan} A nuclear-medicine scan that uses a small radioactive tracer in the blood to find active bleeding in the digestive tract. It can help localize bleeding when endoscopy does not identify the source.

\medterm{Gastrointestinal Cancer} Cancer arising in the digestive tract or related organs, including the esophagus, stomach, intestine, rectum, liver, pancreas, and biliary tract. Symptoms and treatment depend on the organ, tumor type, and stage.

\medterm{Gastrointestinal Decontamination} Methods used after some poisonings to reduce absorption, such as activated charcoal or whole-bowel irrigation. The choice depends on the substance, timing, symptoms, and risk of aspiration; stomach washing is rarely used.

\medterm{Gastrointestinal Disease} Any disorder affecting the digestive tract or related organs. Examples include reflux, ulcers, inflammatory bowel disease, infection, malabsorption, liver disease, and cancer.

\medterm{Gastrointestinal Disorder} A disease affecting any part of the digestive tract, from the esophagus to the anus. It may cause pain, nausea, vomiting, diarrhea, constipation, bleeding, blockage, or poor nutrient absorption, depending on the affected structure.

\medterm{Gastrointestinal Fistula} An abnormal connection between the digestive tract and another organ, the skin, or another section of bowel. Digestive fluid may leak through it, causing infection, dehydration, skin damage, or poor nutrition.

\medterm{Gastrointestinal Gas} Air or other gas within the stomach or intestines, usually from swallowed air or bacterial digestion of food. Excess gas may cause belching, bloating, cramps, or passing gas and can increase with constipation or food intolerance.

\medterm{Gastrointestinal Health} The normal structure and function of the digestive system, including digestion, absorption, bowel movement, and a balanced relationship with gut microorganisms. Problems may include bleeding, persistent pain, vomiting, or unexplained weight loss.

\medterm{Gastrointestinal Hemorrhage} Bleeding from any part of the digestive tract, from the esophagus to the anus. It may cause vomiting blood, black stools, red blood in stool, anemia, dizziness, or shock.

\medterm{Gastrointestinal Hormone} A chemical messenger made by cells in the digestive tract that helps regulate stomach acid, intestinal movement, digestion, appetite, gallbladder activity, or blood-glucose control.

\medterm{Gastrointestinal Motility Disorder} Abnormal movement of food, liquid, or stool through the digestive tract. It can cause difficulty swallowing, nausea, vomiting, bloating, abdominal pain, constipation, or diarrhea; causes may involve digestive muscles, nerves, hormones, or a blockage.

\medterm{Gastrointestinal Neoplasm} An abnormal growth arising in the digestive tract. It may be a harmless growth, a precancerous change, or cancer; symptoms and treatment depend on its tissue, location, and behavior.

\medterm{Gastrointestinal Obstruction} Partial or complete blockage of food, fluid, or stool through the digestive tract. It can cause cramping pain, vomiting, swelling, constipation, dehydration, loss of blood supply, perforation, or sepsis.

\medterm{Gastrointestinal Pain} Pain arising from the digestive tract or related organs. Causes include infection, ulcers, inflammation, blockage, gallstones, and diseases of the liver or pancreas. It may be severe, persistent, or worsening.

\medterm{Gastrointestinal Perforation} A hole through the wall of the digestive tract that allows its contents to leak into surrounding tissue. It can cause severe abdominal pain, peritonitis, sepsis, and shock.

\medterm{Gastrointestinal Stromal Tumor} A tumor that begins in supporting cells in the wall of the digestive tract, most often in the stomach or small intestine. Its behavior varies; treatment may include surgery or medicines designed to block specific tumor signals.

\medterm{Gastrointestinal System} The connected organs that move and digest food, absorb nutrients and water, and remove solid waste. It includes the mouth, esophagus, stomach, intestines, rectum, and anus, with help from the liver, gallbladder, and pancreas.

\medterm{Gastrointestinal Tract} The continuous digestive passage from the mouth through the esophagus, stomach, small and large intestines, rectum, and anus. It moves food, digests it, absorbs nutrients and water, and removes waste.

\medterm{Gastrointestinal Villi} Tiny finger-like projections lining the small intestine that greatly increase the area available to absorb nutrients and water. Each contains small blood vessels and a lymph vessel that carry absorbed substances away.

\medterm{Gastrojejunostomy} An operation creating a new connection between the stomach and jejunum, the middle part of the small intestine. It may bypass a blocked stomach outlet or replace a damaged section of digestive tract.

\medterm{Gastromalacia} Abnormal softening of the stomach wall. It is a rare descriptive finding that may result from severe inflammation, poor blood supply, or tissue injury and can increase the risk of weakness or perforation.

\medterm{Gastroparesis} A disorder in which the stomach empties too slowly, or sometimes stops emptying, even though there is no physical blockage. The stomach muscles or the nerves that control them do not move food normally. It may cause early fullness, nausea, vomiting, bloating, belching, upper-abdominal pain, heartburn, and poor appetite. Diabetes is the most common known cause; other causes include stomach surgery, nerve or muscle disorders, some autoimmune or viral illnesses, although the cause is often unknown. Some medicines can slow stomach emptying or cause similar symptoms without causing gastroparesis itself. Severe disease can lead to dehydration, poor nutrition, and difficult-to-control blood glucose.

\synonyms
Delayed Gastric Emptying; Gastric Hypomotility

\medterm{Gastropathy} Damage or disease of the stomach lining without the prominent inflammation seen in gastritis. Causes include medicines, alcohol, poor blood flow, bile reflux, and severe illness; symptoms may include pain, nausea, or bleeding.

\medterm{Gastroplasty} Surgery that changes the size or shape of the stomach, usually to support weight loss or repair a defect. Possible complications include bleeding, infection, leakage, reflux, and nutritional deficiencies.

\medterm{Gastroschisis} A birth defect in which the intestine, and sometimes other organs, protrudes through an opening beside the umbilicus without a protective covering. The exposed bowel may be swollen or damaged.

\medterm{Gastroschisis Repair} Surgical repair of gastroschisis, a congenital abdominal wall defect where intestines protrude through an opening near the umbilicus. The bowel is returned to the abdomen and the defect is closed, sometimes in stages using a temporary covering.

\synonyms
Abdominal Wall Defect Repair; Congenital Evisceration Closure

\medterm{Gastroscope} A thin, flexible instrument with a light and camera used to view the esophagus, stomach, and first part of the small intestine. It can also pass tools for biopsy or treatment.

\medterm{Gastroscopy} Examination of the esophagus, stomach, and duodenum using a flexible camera tube passed through the mouth. It can identify bleeding, ulcers, inflammation, narrowing, or tumors and allows biopsies or selected treatments.

\medterm{Gastrospenic Fistulization} An abnormal connection between the stomach and spleen, usually caused by a severe stomach ulcer or a splenic abscess. It can cause upper-abdominal pain, digestive-tract bleeding, severe infection, and sepsis.

\medterm{Gastrostomy} A surgically created opening through the abdominal wall into the stomach. A tube placed through it can provide food, fluids, or medicines or remove stomach contents when normal feeding is unsafe or impossible.

\medterm{Gastrostomy Tube} A tube passed through the abdominal wall into the stomach to provide food, fluids, or medicines when swallowing is unsafe or inadequate. Problems can include skin infection, leakage, blockage, bleeding, and movement of the tube.

\medterm{Gastrostomy Tube Insertion} A surgical, laparoscopic, or endoscopic procedure in which a feeding tube is placed directly through the abdominal wall into the stomach for long-term enteral nutrition.

\synonyms
G-Tube Placement; Percutaneous Endoscopic Gastrostomy; PEG Tube Insertion

\medterm{Gastrula} An early embryo stage formed after gastrulation, when cell layers reorganize into ectoderm, mesoderm, and endoderm that give rise to body tissues.

\medterm{Gastrulation} An early stage of embryo development in which cells move and reorganize into three layers that later form all major tissues and organs.

\medterm{Gated Blood Pool Scan} A radionuclide cardiac study, also called radionuclide ventriculography or MUGA, that uses ECG-gated images of labeled blood to measure ventricular ejection fraction and assess ventricular wall motion. It is used selectively when quantitative ventricular function is needed.

\medterm{Gatifloxacin} Gatifloxacin is a fluoroquinolone antibiotic that inhibits bacterial DNA gyrase and topoisomerase IV; systemic use is restricted in many settings because of dysglycaemia and other class toxicities, while ophthalmic use remains available in some regions.

\medterm{Gaucher Disease} An autosomal recessive lysosomal storage disorder caused by reduced activity of the enzyme glucocerebrosidase, usually because of disease-causing variants in the \textit{GBA1} gene. Glucosylceramide accumulates in macrophages and other cells, affecting the spleen, liver, bone marrow, bones, lungs, and sometimes the nervous system. It can cause enlargement of the spleen or liver, anemia, low platelets, easy bruising, bone pain or fractures, and neurologic problems. Type 1 is usually non-neuronopathic, while types 2 and 3 involve the nervous system.

\medterm{Gaucher's Disease} Alternate spelling; see \textit{Gaucher Disease}.

\medterm{Gaucher Splenomegaly} Alternate name; see \textit{Gaucher disease}.

\medterm{Gaultheria Oil} Oil of wintergreen, an essential oil containing concentrated methyl salicylate. It may be used in small amounts in topical products, but swallowing even a small amount can cause serious salicylate poisoning, especially in children.

\medterm{Gauss} A unit of magnetic flux density equal to $10^{-4}$ tesla, used in physics, imaging, and descriptions of magnetic fields.

\medterm{Gauze} A woven or nonwoven material used to clean, dress, pack, or absorb fluid from a wound. It may be sterile or nonsterile and comes in several forms.

\medterm{Gavage} Giving liquid food or medicine through a tube into the stomach. It is used when swallowing is unsafe or oral intake is inadequate, including for premature infants who cannot coordinate sucking, swallowing, and breathing.

\medterm{Gaze Palsy} An inability to move both eyes conjugately in a particular direction (horizontal, vertical, or upgaze), caused by lesions involving the cortical frontal eye fields, the paramedian pontine reticular formation (PPRF), or the midbrain rostral interstitial nucleus of the medial longitudinal fasciculus (riMLF).
\synonyms
Conjugate Gaze Palsy; Ocular Conjugate Palsy

\medterm{GBM} An abbreviation for glioblastoma, an aggressive cancer that starts in supporting cells of the brain. It can grow quickly and cause headache, seizures, weakness, speech or memory problems, or personality changes. Treatment may include surgery, radiation, and medicines.

\synonyms
Glioblastoma Multiforme; Grade IV Astrocytoma

\medterm{Gefitinib} An oral epidermal growth factor receptor tyrosine-kinase inhibitor used for selected advanced non-small-cell lung cancers with activating EGFR mutations.

\medterm{Gegenhalten} A velocity-independent, variable increase in resistance to passive limb movement in which the patient involuntarily opposes the examiner's manipulation with a force proportional to that applied by the examiner, commonly observed in diffuse frontal lobe disease, dementia, and metabolic encephalopathy.
\synonyms
Paratonia; Oppositional Paratonia

\medterm{Geiger Counter} An instrument that detects and measures ionizing radiation. It can identify radioactive contamination or exposure, but its readings do not by themselves show how much biological harm the radiation caused.

\medterm{Gel} A semisolid preparation in which liquid is held in a soft network. Gels may be applied to the skin, eyes, mouth, or other areas to soothe tissue, deliver medicine, or form a protective layer.

\medterm{Gelastic Seizure} A rare seizure that causes sudden, repeated bursts of laughter without feeling amused. It is often linked to a small growth called a hypothalamic hamartoma, but other brain conditions can also cause it.

\medterm{Gelatin} Gelatin is a protein made from collagen and used in foods, medicines, capsules, and some medical materials.

\medterm{Gel Electrophoresis} Gel electrophoresis separates DNA, RNA, or proteins by moving them through a gel in an electric field, usually according to size and charge.

\medterm{Gelfoam} An absorbable sponge made from purified animal gelatin that helps control bleeding during surgery. It provides a surface for clot formation and is gradually absorbed; swelling or infection can rarely occur.

\medterm{Gelineau Syndrome} A sleep disorder causing repeated, uncontrollable daytime sleep episodes. Some people also have sudden muscle weakness triggered by emotion, sleep paralysis, or vivid dream-like experiences while falling asleep or waking.

\medterm{Gemcitabine} A chemotherapy medicine that becomes active inside cells and interferes with DNA production, slowing or killing rapidly dividing cancer cells. It is used for several cancers and may cause low blood counts, nausea, rash, or organ toxicity.

\medterm{Gemcitabine Hydrochloride} The hydrochloride form of gemcitabine, a chemotherapy medicine that interferes with DNA production in dividing cells. It is used for several cancers; important risks include infection from low blood counts and treatment-related organ injury.

\medterm{Gemella} Gemella is a group of bacteria that normally may inhabit the mouth and upper airway but can rarely cause serious infections, including bloodstream infection.

\medterm{Gemfibrozil} A medicine that lowers very high blood triglycerides and may reduce the risk of pancreatitis. It changes how the liver processes fats but can interact with statins and increase the risk of serious muscle injury.

\medterm{Gemifloxacin} An antibiotic in the fluoroquinolone group used for selected bacterial respiratory infections. It can cause nausea and, rarely, tendon injury, nerve symptoms, abnormal heart rhythm, severe diarrhea, or dangerous allergic reactions.

\medterm{Gemination} An attempt by one developing tooth bud to divide, producing a large tooth with a split or partly divided crown. The total number of teeth is usually normal, and dental assessment guides treatment.

\medterm{Gemistocytic Astrocytoma} A former histologic variant of diffuse astrocytoma composed partly of gemistocytic astrocytes. The term is now used mainly descriptively because integrated molecular classification determines the modern tumor diagnosis and grade.

\medterm{Gemtuzumab Ozogamicin} An antibody-linked chemotherapy medicine used for selected acute myeloid leukemia. The antibody attaches to CD33-positive leukemia cells and delivers a toxic drug; liver injury, infusion reactions, and low blood counts are important risks.

\medterm{Gender} A social, cultural, and personal concept involving identity, roles, expression, and expectations. Gender is related to, but distinct from, sex characteristics and sexual orientation.

\medterm{Gender Dysphoria} Clinically significant distress or impairment associated with incongruence between a person’s experienced or expressed gender and sex assigned at birth. Having a gender identity different from sex assigned at birth does not by itself constitute gender dysphoria.

\medterm{Gender Identity} A person’s internal sense of being a woman, man, both, neither, or another gender, which may or may not correspond to sex assigned at birth.

\medterm{Gender Identity Disorder} Alternate historical term; see \textit{Gender Dysphoria}.

\medterm{Gender Incongruence} A difference between a person’s experienced or expressed gender and sex assigned at birth. Gender incongruence is a descriptive term and does not by itself imply distress, impairment, or a mental disorder.

\medterm{Gene} A section of DNA that contains instructions for making a functional product, usually a protein or RNA. Genes are arranged at specific locations on chromosomes and influence inherited traits and cell function.

\medterm{Gene Amplification} An increase in the number of copies of a gene or genomic region, which can raise gene expression and drive some inherited conditions or cancers.

\medterm{Gene Conversion} A nonreciprocal transfer of sequence information from one DNA region to a homologous region, potentially changing an allele during recombination or repair.

\medterm{Gene Deletion} Loss of part or all of a gene from DNA, potentially abolishing its product or altering dosage and causing disease when clinically significant.

\medterm{Gene Dosage} The number of functional copies of a gene and the amount of product they produce, influenced by deletions, duplications, chromosome number, and regulation.

\medterm{Gene Duplication} An extra copy of a gene created by replication or rearrangement, which may have no effect, alter dosage, or provide material for new gene function.

\medterm{Gene-Environment Interaction} A gene-environment interaction occurs when an environmental exposure changes disease risk differently according to a person's genotype, so inherited susceptibility and exposure jointly influence phenotype.

\medterm{Gene Expression} The process by which a cell uses information in a gene to make RNA or protein. Cells control when and how strongly genes are expressed, allowing different tissues to develop and perform specialized functions.

\medterm{Gene Frequency} The proportion of all gene copies in a population that carry a particular version, or allele, of that gene. It describes population genetics and does not by itself predict an individual’s disease risk.

\medterm{Gene Fusion} A hybrid gene formed when segments of two separate genes join, sometimes producing an abnormal protein that drives cancer or other disease.

\medterm{Gene Modification} An intentional or naturally occurring alteration of DNA sequence or gene regulation that changes cellular function or inherited traits.

\medterm{Gene-Modified} Describing cells or organisms whose genetic material has been deliberately altered using recombinant DNA or genome-editing methods.

\medterm{Gene Mutation} A permanent change in the DNA sequence of a gene, which may be harmless, beneficial, or pathogenic depending on its effect and inheritance.

\medterm{Gene Panel} A genetic test that examines a selected group of genes linked to a person’s symptoms or suspected condition. It can provide focused results but may miss changes in genes or DNA regions not included in the panel.

\medterm{General Anesthesia} A controlled, temporary state of unconsciousness produced by medicines so a person does not feel or remember a procedure. Breathing, blood pressure, heart rate, and other body functions are monitored throughout.

\medterm{General Anesthetic} A medicine that produces a controlled, reversible state of unconsciousness and reduced awareness of pain, often with muscle relaxation, so surgery or another procedure can be performed safely.

\medterm{Generalised} Widespread or involving multiple regions of the body rather than being confined to a single localized site.

\medterm{Generalized} Distributed widely or involving most or all of a body region, system, or population rather than being confined to one focus.

\medterm{Generalized Anxiety Disorder} A mental-health disorder involving excessive, difficult-to-control worry about several areas of life on most days for at least six months. Restlessness, tiredness, poor concentration, irritability, muscle tension, and sleep problems may occur.

\medterm{Generalized Anxiety Disorder In Children} Excessive, difficult-to-control worry in children about multiple topics (school performance, health, family safety) occurring more days than not for at least 6 months. Associated with restlessness, fatigue, difficulty concentrating, irritability, muscle tension, and sleep disturbance. Treatment includes cognitive behavioral therapy and possibly medication.

\synonyms
Pediatric GAD; Childhood Generalized Anxiety

\medterm{Generalized Granuloma Annulare} A widespread skin condition causing many small, skin-colored, red, or purple bumps over the trunk and limbs. It may be associated with diabetes and can last longer than localized granuloma annulare.

\medterm{Generalized Lipodystrophy} A disorder with near-total loss of adipose tissue, severe insulin resistance, diabetes, high triglycerides, fatty liver, and abnormal fat deposition in other tissues.

\medterm{Generalized Lymphadenopathy} Pathological enlargement of lymph nodes involving two or more non-contiguous anatomical lymph node chains. Unlike localized adenopathy, it reflects systemic, hematologic, infectious, or autoimmune disease. Major etiologies include lymphomas, leukemias (such as chronic lymphocytic leukemia), systemic viral infections (HIV, Epstein-Barr virus, cytomegalovirus), secondary syphilis, disseminated tuberculosis, systemic lupus erythematosus, sarcoidosis, and drug-induced hypersensitivity reactions (such as DRESS syndrome).

\synonyms
Diffuse Lymphadenopathy; Systemic Lymphadenopathy

\medterm{Generalized-Onset Seizure} A seizure that begins in networks on both sides of the brain. It may cause staring, brief muscle jerks, stiffening, loss of muscle tone, or whole-body stiffening and jerking, depending on the type.

\medterm{Generalized Seizure} A seizure arising from or rapidly involving both cerebral hemispheres, commonly producing impaired awareness, bilateral motor activity, or brief absence episodes.

\medterm{Generalized Tonic-Clonic Seizure} A seizure causing loss of consciousness, stiffening of the whole body, and then repeated jerking of the limbs. Confusion and sleepiness often follow; a seizure lasting five minutes or longer is prolonged.

\medterm{Generalized Weakness} Reduced strength affecting multiple body regions, caused by systemic illness, metabolic or electrolyte disturbance, medication, infection, muscle disease, nerve disease, neuromuscular-junction disease, or brain or spinal-cord disease. It may occur with speech, facial, breathing, or walking difficulty.

\medterm{General Paresis} A late form of syphilis affecting the brain and spinal cord. It can gradually cause memory loss, personality or mood changes, poor judgment, weakness, tremor, or difficulty walking. Blood and spinal-fluid tests help identify the infection.

\synonyms
General Paresis of the Insane; Paralytic Dementia; Neurosyphilis

\medterm{Generator} A device that produces a short-lived radioactive substance from a longer-lived one for nuclear-medicine tests. The product is separated, checked for purity and sterility, and then used to make a tracer.

\medterm{Gene Rearrangement} Reorganization of DNA segments within or between chromosomes that can alter gene structure, expression, immune-receptor diversity, or cancer biology.

\medterm{Generic} Describing a medicine by its active ingredient rather than a brand name; an approved generic drug must meet applicable standards for quality and therapeutic equivalence to its reference product.

\medterm{Genes} Sections of DNA that contain instructions for making proteins or functional RNA. They are arranged on chromosomes and influence inherited traits, cell activity, disease risk, and responses to some treatments.

\medterm{Gene Silencing} Reduction or stopping of a gene's activity through DNA changes, chromatin control, RNA interference, or other processes that limit protein production.

\medterm{Genesis} The origin, formation, or development of something, such as a disease, organ, symptom, or process. Understanding its genesis can help explain causes, risk factors, progression, prevention, and treatment.

\medterm{Gene Targeting} A molecular method that changes a chosen genomic locus through homologous recombination or engineered editing to study or alter gene function.

\medterm{Gene Therapy} Treatment that adds, removes, replaces, or edits genetic material in a patient’s cells to treat disease. It may correct a missing gene function or alter harmful activity; risks include immune reactions and unintended genetic effects.

\medterm{Genetic Anticipation} Earlier onset or increasing severity of a hereditary disorder in successive generations, often caused by expansion of unstable repeated DNA sequences.

\medterm{Genetic Carrier} A person with one pathogenic variant for a recessive disorder who is usually unaffected but can transmit the variant to children.

\medterm{Genetic Code} The rules cells use to read genetic instructions in groups of three DNA or RNA building blocks. Each group directs the cell to add a particular amino acid to a protein or to stop making the protein.

\medterm{Genetic Counseling} A professional discussion that explains how inherited factors may affect a person or family. It may cover risk, testing choices, results, reproductive options, emotional support, and implications for relatives.

\medterm{Genetic Counselling} A process in which a trained professional explains inherited conditions, test results, health implications, and reproductive choices to a person or family.

\medterm{Genetic Counselor} A trained health professional who explains inherited conditions, estimates genetic risk, discusses testing, and helps patients and families understand results and choices. They may coordinate testing of relatives when appropriate.

\medterm{Genetic Disease} Alternate name; see \textit{Genetic Disorder}.

\medterm{Genetic Disorder} A disease or disorder caused wholly or partly by a pathogenic change in genetic material. The change may involve one or more genes, chromosome number or structure, mitochondrial DNA, or gene regulation. It may be inherited through germline DNA, arise de novo, or occur only in certain tissues as a somatic change. Some conditions also result from interactions among multiple genetic variants and environmental factors; not every genetic change causes disease.

\medterm{Genetic Drift} Random changes in allele frequencies from one generation to the next, especially in small populations, independent of natural selection.

\medterm{Genetic Engineering} Deliberate changing, adding, removing, or transferring genetic material using laboratory methods. It supports research, diagnosis, treatment development, agriculture, and biotechnology but requires careful assessment of safety, ethics, and unintended effects.

\medterm{Genetic Heterogeneity} The situation in which similar features or the same disease can result from different genetic causes. Different genes or different changes within one gene may produce overlapping effects.

\medterm{Genetic Linkage} The tendency of nearby genes or DNA markers on the same chromosome to be inherited together because crossing over is less likely between them.

\medterm{Genetic Marker} A detectable DNA variation or feature used to track inheritance, identify a chromosome region, assess disease risk, or distinguish samples.

\medterm{Genetic Polymorphism} A common DNA variation with at least two alleles in a population, which may influence traits, drug response, or disease susceptibility.

\medterm{Genetic Profile} A set of genetic variants or DNA markers describing an individual, tumor, organism, or sample for diagnosis, risk assessment, identification, or research.

\medterm{Genetic Recombination} Exchange or rearrangement of DNA segments that creates new allele combinations during meiosis, repair, immune-cell development, or laboratory manipulation.

\medterm{Genetics} Genetics is the study of genes, heredity, genetic variation, and how DNA and gene regulation influence traits, disease risk, diagnosis, and treatment response.

\medterm{Genetic Screening} Testing people without symptoms to estimate the chance of a genetic condition. Examples include newborn screening, carrier screening before or during pregnancy, and prenatal screening; screening results are not always diagnostic.

\medterm{Genetic Sequencing} Analysis of the order of DNA bases to identify variants associated with genetic disorders. The method is selected according to the gene, variant, and clinical question.

\synonyms
DNA Sequencing; Next-Generation Sequencing (NGS)

\medterm{Genetic Sonogram} A detailed pregnancy ultrasound, usually performed around 18–22 weeks, that examines fetal anatomy and may look for findings associated with chromosome conditions. It cannot detect every genetic disorder.

\medterm{Genetic Susceptibility} Inherited tendency to develop a disease or respond to an exposure, which usually modifies risk rather than determining outcome alone.

\medterm{Genetic Technique} A laboratory method used to detect, amplify, sequence, edit, compare, or otherwise analyze DNA, RNA, chromosomes, or gene expression.

\medterm{Genetic Testing} Analysis of DNA or chromosomes to assess disease risk, explain symptoms, identify carriers, or guide treatment. Results may also provide information about inherited traits and conditions affecting relatives.

\medterm{Genetic Transcription} The first step in using a gene: an enzyme copies information from DNA into RNA. The RNA may then guide protein production or perform another function inside the cell.

\medterm{Genetic Variant Of Uncertain Significance} A DNA change whose effect on health is not yet known. It does not by itself show that a person has a disease or predict future disease and may later be classified as harmless or disease-causing as evidence improves.

\medterm{Gene Transfer} Movement of genetic material into a cell or organism, naturally or through laboratory vectors, potentially changing gene expression or correcting a defect.

\medterm{Geniculate Body} One of two paired relay structures in the brain. The lateral geniculate body processes visual signals, while the medial geniculate body processes hearing signals before sending them to the cerebral cortex.

\medterm{Geniculate Ganglion} An L-shaped sensory ganglion of the facial nerve (cranial nerve VII) located within the facial canal of the temporal bone, housing the primary pseudo-unipolar sensory neuronal cell bodies conveying special taste sensation from the anterior two-thirds of the tongue (via chorda tympani).

\medterm{Genital} Relating to the reproductive organs or external sexual structures. Genital symptoms may involve skin, mucosa, nerves, glands, or the urinary and reproductive systems.

\medterm{Genital Edema} Swelling of genital tissues from fluid accumulation caused by infection, inflammation, allergy, trauma, venous or lymphatic obstruction, or systemic disease.

\medterm{Genital Herpes} An infection of the genital or anal area caused by herpes simplex virus type 1 or type 2 (HSV-1 or HSV-2). It spreads mainly through sexual skin-to-skin contact and can be transmitted even when no sores are present. It may cause painful blisters or ulcers, burning urination, swollen glands, or no symptoms; outbreaks can recur. Infection during pregnancy can threaten the newborn.

\medterm{Genital Herpes in Women} Herpes simplex virus infection of the female genitalia, causing painful blisters or ulcers. It may have no symptoms. Antiviral treatment can reduce symptoms and transmission; infection during pregnancy can affect the newborn.

\synonyms
Female Genital Herpes; HSV-2 Infection in Women

\medterm{Genitalia} The external and internal reproductive organs. The term includes the vulva and vagina,
or the penis, scrotum, and associated structures, together with the gonads and
reproductive ducts when used broadly; anatomy varies with sex, age, and developmental
stage.

\medterm{Genital Injury} Damage to the external or internal genital tissues from childbirth, accidents, sexual contact, assault, medical procedures, or disease. It may cause pain, swelling, bleeding, or difficulty urinating.

\medterm{Genital Mycoplasmas} Bacteria that may live in or infect the genital or urinary tract. Some cause inflammation of the urethra or cervix, pelvic inflammatory disease, or few symptoms; testing helps distinguish infection from harmless colonization.

\medterm{Genital Prolapse} The dropping or bulging of one or more pelvic organs (such as the bladder, uterus, or rectum) into or through the vagina due to weakening of the pelvic floor muscles and supporting tissues. It commonly causes a feeling of a vaginal bulge or heaviness, lower back discomfort, and urinary or bowel difficulties; risks include vaginal childbirth, aging, menopause, and chronic straining.
\synonyms Pelvic organ prolapse; Uterovaginal prolapse; Urogenital prolapse.

\medterm{Genitals} The external and internal reproductive organs. Female genitals: external (vulva — mons
pubis, labia majora, labia minora, clitoris, vaginal opening, Bartholin glands) and
internal (vagina, cervix, uterus, fallopian tubes, ovaries).

\medterm{Genital System} The organs involved in reproduction and sexual function, including the gonads, reproductive ducts, accessory glands, and external genital structures. Anatomy differs between individuals and changes with development and age.

\medterm{Genital Ulcer Disease} A clinical syndrome characterized by single or multiple ulcerative breaches of the genital or perianal skin and mucosa, predominantly caused by sexually transmitted pathogens including herpes simplex virus (HSV-2), \textit{Treponema pallidum} (syphilis chancre), and \textit{Haemophilus ducreyi} (chancroid).

\synonyms
GUD

\medterm{Genital Wart} A wart in the genital or anal area caused by certain types of human papillomavirus. Warts may be flat or raised and can recur; treatment removes visible lesions but may not eliminate the virus immediately.

\medterm{Genital Warts} Warts of the genital or anal region caused most often by low-risk types of human papillomavirus. They spread through direct skin-to-skin sexual contact and may be visible or asymptomatic.

\synonyms
Condylomata Acuminata; Anogenital Warts

\medterm{Genitofemoral Nerve} A branch of the lumbar plexus (arising from L1 and L2 nerve roots) that pierces the anterior surface of the psoas major muscle and divides into a genital branch (innervating the cremaster muscle and scrotal/labial skin; mediating the cremasteric reflex) and a femoral branch (cutaneous sensation to the anterior upper thigh).

\medterm{Genito-Pelvic Pain/Penetration Disorder} Persistent or recurrent difficulty with vaginal penetration, including pain, fear or anxiety about pain, involuntary pelvic-floor tightening, or pain during attempted penetration, causing clinically significant distress. Causes may include pelvic-floor problems, infection, hormonal changes, injury, or emotional factors.

\medterm{Genitourinary} Relating to the urinary system and the reproductive organs, which share developmental and anatomical relationships.

\medterm{Genitourinary Fistula} Alternate name for urogenital fistula, an abnormal connection between the urinary and genital tracts.

\medterm{Genito-Urinary Medicine} The medical specialty concerned with urinary and reproductive health. It includes evaluation and treatment of urinary infections, sexual health problems, sexually transmitted infections, and some reproductive conditions.

\medterm{Genitourinary Neoplasm} An abnormal growth in the urinary or reproductive organs, such as the kidney, bladder, prostate, testis, ovary, or uterus. It may be benign or cancerous; examination, imaging, and sometimes biopsy identify its nature.

\medterm{Genitourinary Syndrome Of Menopause} A group of changes in the genital and urinary tissues after menopause, caused mainly by lower estrogen levels. It may include vaginal dryness, burning, pain during sex, urinary urgency, painful urination, or repeated urinary infections.

\medterm{Genitourinary Tract} The combined urinary and reproductive organs and their associated ducts.

\medterm{Genodermatoses} Inherited disorders that mainly affect the skin, hair, nails, sweat glands, or related organs. Examples include inherited scaling disorders, blistering diseases, and some neurocutaneous syndromes; severity varies widely.

\medterm{Genome} The complete genetic material of an organism, including genes and DNA that does not directly code proteins. It contains instructions influencing development, body function, inherited traits, disease risk, and treatment response.

\medterm{Genomic Imprinting} A pattern in which a gene’s activity depends on whether its copy came from the mother or father. Abnormal imprinting can cause conditions such as Prader–Willi or Angelman syndrome.

\medterm{Genomic Instability} An increased tendency for DNA changes, chromosome damage, and mutations to accumulate, contributing to cancer development, treatment resistance, and disease progression.

\medterm{Genomics} The study of an organism’s complete genetic material, including its DNA sequence, variation, activity, and interaction with disease and treatment. It supports research, diagnosis, risk assessment, and personalized care.

\medterm{Genotoxicity} The ability of a substance or physical agent to damage genetic material. Genotoxicity may cause mutations or chromosome damage, but a positive test does not by itself establish that an agent causes cancer in humans.

\medterm{Genotype} A person’s genetic makeup, meaning the versions of one or more genes they carry. The genotype helps influence observable traits and disease risk together with other genes and environmental factors.

\medterm{Genotype Determination} Genotype determination is the use of laboratory testing to identify an individual's genetic variants or alleles at one or more locations.

\medterm{Genotype-Phenotype Correlation} The relationship between a person’s genetic makeup and the observable features, severity, age of onset, or clinical course of a condition; correlations may be incomplete because other genetic and environmental factors contribute.

\medterm{Gentamicin} An aminoglycoside antibiotic that interferes with bacterial protein production. It can injure the kidneys or hearing and is used with dose and blood-level monitoring.

\medterm{Gentamicin Sulfate} The sulfate salt of gentamicin, an aminoglycoside antibiotic that interferes with bacterial protein production. It can injure the kidneys or hearing and is used with careful dosing and monitoring.

\medterm{Gentian} A plant of the Gentiana genus whose roots or other parts have been used in herbal preparations, traditionally as a bitter digestive; clinical evidence and product quality vary.

\medterm{Gentian Violet} A purple antiseptic and dye with antifungal and antibacterial activity, used in selected topical applications but capable of irritating tissue and staining skin.

\medterm{Genu} The knee or a knee-shaped part of an anatomical structure; the term also describes a bend, such as the genu of the corpus callosum.

\medterm{Genu Recurvatum} Hyperextension of the knee beyond the usual straight position, caused by laxity, deformity, muscle imbalance, neurologic disease, or injury.

\medterm{Genu Valgum} An alignment abnormality in which the knees angle inward and may touch while the ankles remain apart. It may be developmental or caused by bone, joint, or metabolic disease.

\synonyms
Knock-Knees; Tibial Valgus

\medterm{Genu Varum} Outward bowing of the legs, commonly called bowlegs. It is often normal in young children and improves with growth, but persistent, worsening, unequal, painful, or severe bowing may result from rickets, growth-plate disease, or another condition.

\synonyms
Bowlegs; Bandy Legs

\medterm{Geographic Atrophy} An advanced form of dry age-related macular degeneration in which areas of light-sensing and supporting cells in the center of the retina are permanently lost. It causes slowly worsening blind spots and loss of detailed central vision.

\medterm{Geographic Tongue} Benign inflammatory condition of the tongue characterized by irregular red patches with white or yellow borders, alternating with areas of depapillation. Often asymptomatic and of unknown cause. May cause burning or irritation with certain foods.

\synonyms
Benign Migratory Glossitis
Benign Migratory Glossitis

\medterm{Geophagia} Persistent eating of soil or other nonfood substances; it may cause poisoning, infection, intestinal obstruction, or iron-deficiency anemia.

\medterm{Geophagy} The deliberate ingestion of soil or clay; it may cause lead or parasite exposure, gastrointestinal obstruction, or iron-deficiency anemia and can occur with pica.

\medterm{Geotrichum} Geotrichum is a genus of fungi that may colonize humans and can rarely cause superficial or invasive infection, especially in people with weakened immunity.

\medterm{Geriatric Assessment} A broad review of an older person’s medical conditions, medicines, mobility, memory, mood, daily activities, nutrition, safety, home situation, and goals. It identifies treatable problems and guides an individualized care plan.

\medterm{Geriatric Care Manager} A professional who helps coordinate medical, social, housing, and support services for an older adult. They may assess needs, organize care, communicate with families, and help plan for changing abilities.

\synonyms
Aging-Life Care Manager; Eldercare Coordinator

\medterm{Geriatric Dentistry} Dental care adapted to the needs of older adults, including tooth decay, gum disease, dry mouth, dentures, swallowing concerns, medicine effects, and examination for oral cancer.

\medterm{Geriatric Depression} Depression in an older adult may cause low mood, loss of interest, fatigue, sleep or appetite changes, slowed thinking, physical complaints, or memory problems. Medical illness, medicines, isolation, and bereavement may contribute.

\synonyms
Late-Life Depression; Depression in the Elderly; Late-Onset Depressive Disorder

\medterm{Geriatric Depression Scale} A questionnaire used to screen older adults for symptoms of depression. It can identify people with possible depression, but a positive result does not by itself establish depression or explain its cause.

\medterm{Geriatric Dream Changes} Dreams in older adults may become more vivid or be affected by medicines, sleep disorders, depression, or changes in the sleep cycle. Repeated frightening dreams, acting them out, confusion, or sudden change may occur.

\synonyms
Geriatric Sleep Patterns; Dreaming in Older Adults

\medterm{Geriatric Dysphagia} Difficulty swallowing in an older adult. It may cause coughing or choking during meals, food remaining in the mouth, weight loss, dehydration, or aspiration and can result from stroke, dementia, muscle disease, blockage, or medicines.

\synonyms
Presbyphagia; Age-Related Swallowing Dysfunction

\medterm{Geriatric Hospitalization} Hospital care for an older adult, with attention to medical illness as well as mobility, medicines, nutrition, hearing, vision, and cognition. Preventing falls, delirium, pressure injuries, and loss of independence supports safer recovery.

\synonyms
Inpatient Geriatric Care; Acute Care for Elders (ACE)

\medterm{Geriatrician} A physician specializing in the health care of older adults, including complex illness, multiple conditions, medications, function, and goals of care.

\medterm{Geriatric Infection} Infection in an older adult may cause confusion, weakness, reduced appetite, or a fall without causing fever. Common sites include the urine, lungs, skin, and bloodstream; illness can worsen quickly.

\synonyms
Infection in the Elderly; Late-Life Infection

\medterm{Geriatric Informed Consent} The process of ensuring that an older patient understands a proposed intervention, its benefits, risks, and alternatives before choosing freely. If decision-making ability is impaired, an authorized decision-maker may help according to local law.

\synonyms
Geriatric Decision-Making Capacity; Surrogate Consent in Older Adults

\medterm{Geriatric Medical Ethics} Ethical decision-making in the care of older adults, including respect for the person’s wishes, decision-making capacity, likely benefits and harms, fairness, privacy, and appropriate involvement of family or caregivers.

\medterm{Geriatric Nursing} Nursing care focused on the health, function, medication safety, cognition, mobility, and goals of older adults.

\medterm{Geriatric Nutrition} Nutrition care for older adults that considers aging, illness, medicines, frailty, dental or swallowing problems, appetite, and access to food. Goals include maintaining strength, muscle, hydration, function, and quality of life.

\medterm{Geriatric Osteoarthritis} Osteoarthritis occurring in later life, commonly affecting the knees, hips, hands, or spine
and causing joint pain, stiffness, and reduced movement. Management may include exercise,
weight management, medicines, assistive measures, or joint replacement when appropriate.

\synonyms
Age-Related Osteoarthritis; Senile Osteoarthritis

\medterm{Geriatric Rehabilitation} Therapy that helps an older adult regain or maintain movement, strength, independence, and safety after illness, injury, surgery, or functional decline. The plan is adapted to health, abilities, goals, and support needs.

\medterm{Geriatrics} The medical specialty focused on health, function, independence, and quality of life in older adults. Care considers multiple illnesses, medicines, memory, mobility, nutrition, mood, safety, goals, and support.

\medterm{Geriatric Screening} Checking older adults who have no symptoms for diseases or risks that may benefit from early detection. Examples include falls, depression, osteoporosis, and selected cancers.

\synonyms
Geriatric Screening; Senior Health Assessment

\medterm{Geriatric Sleep Changes} Sleep patterns and sleep problems in older adults, including insomnia, sleep apnea, circadian changes,
and medication-related disturbances.

\synonyms
Geriatric Sleep Disorders; Elder Sleep Management

\medterm{Geriatric Syndrome} A common problem in older adults caused by several interacting factors rather than one disease. Examples include falls, frailty, delirium, incontinence, immobility, and pressure injuries; causes may be medical, social, or both.

\medterm{Geriatric Trauma} Injury in an older adult, often involving falls, fractures, head injury, or bleeding. Age-related frailty, osteoporosis, balance problems, and medicines that increase bleeding can make injuries more serious and recovery slower.

\medterm{Germ} Alternate informal term; see \textit{Microorganism}.

\medterm{Germanium} A chemical element sometimes included in unproven health supplements. Germanium compounds have no established medical benefit and can cause kidney damage, nerve damage, or other serious toxicity when taken internally.

\medterm{German Measles} Rubella: a contagious viral infection caused by rubella virus, characterized by fever, lymph-node enlargement, and a fine rash. Infection during pregnancy can harm the fetus, causing congenital rubella syndrome with cardiac, auditory, and developmental defects.

\synonyms
Rubella; Three-Day Measles

\medterm{Germ Cell} A reproductive cell or precursor that gives rise to sperm or eggs and carries genetic material to the next generation.

\medterm{Germ Cell Tumor} A tumor arising from cells that normally develop into sperm or eggs. It may occur in the testicles, ovaries, or other sites such as the chest or abdomen; some types are highly treatable with surgery, chemotherapy, or both.

\medterm{Germ Cell Tumors} Tumors arising from cells that normally develop into sperm or eggs. They can occur in the testicles, ovaries, or other body sites and may be benign or cancerous, depending on the type.

\medterm{Germinoma} A usually slow-growing germ-cell cancer that most often develops in the pineal or suprasellar area of the brain. Related tumors occur in the testicles or ovaries; treatment commonly uses radiation, chemotherapy, surgery, or combinations.

\medterm{Germ Layer} One of the three early layers of an embryo—ectoderm, mesoderm, or endoderm—from which all major tissues and organs develop. Each layer gives rise to characteristic body structures during development.

\medterm{Germline Mutation} A genetic change present in an egg or sperm and therefore usually present in nearly every cell of the person. Some changes can be passed to children and may increase the risk of inherited disorders or cancer.

\medterm{Germline Testing} Genetic testing of blood, saliva, or another sample to look for changes that may be inherited or passed to children. Results can affect the person’s care, relatives, cancer screening, and reproductive decisions.

\medterm{Germline Variant} A DNA change present in cells that can give rise to eggs or sperm and therefore may be passed to children. It is usually present in most cells of the person, unlike a change limited to a tumor or other tissue.

\medterm{Germ Theory Of Disease} The evidence-based principle that particular microorganisms can cause particular infectious diseases. It underlies infection control, vaccination, antimicrobial treatment, sterilization, and methods used to identify disease-causing organisms.

\medterm{Gerontologist} A specialist studying aging and its biological, psychological, social, or economic effects. Gerontologists may work in research, healthcare, policy, counseling, education, housing, or community services for older people.

\medterm{Gerstmann-Straussler-Scheinker Syndrome} A rare, autosomal dominant inherited prion disease caused by mutations in the PRNP gene (most commonly P102L), characterized pathologically by multicentric PrP-amyloid plaques throughout the cerebellum and cerebral cortex, presenting with slowly progressive cerebellar ataxia, spasticity, and late dementia.
\synonyms
GSS Syndrome; GSS Prion Disease

\medterm{Gerstmann Syndrome} A neurological cognitive syndrome caused by a lesion of the dominant angular gyrus in the inferior parietal lobule (typically left hemisphere), characterized by the classic tetrad of agraphia (dysgraphia), acalculia (dyscalculia), finger agnosia, and right-left disorientation.
\synonyms
Angular Gyrus Syndrome; Dominant Parietal Lobe Syndrome

\medterm{Gesell Developmental Schedules} Standardized developmental assessment schedules that describe expected motor, language, adaptive, and social milestones in children; they support screening but do not by themselves diagnose a disorder.

\medterm{Gestation} The period of pregnancy during which an embryo and fetus develop in the uterus. Pregnancy is usually dated as about 40 weeks from the first day of the last menstrual period, or about 38 weeks from fertilization.

\medterm{Gestational Age} The length of a pregnancy measured in weeks from the first day of the last menstrual period. It is usually about two weeks longer than the time since fertilization and helps estimate development, due date, and appropriate care.

\synonyms
Menstrual Age; Fetal Postmenstrual Age

\medterm{Gestational Diabetes} Carbohydrate intolerance of variable severity with onset or first recognition during pregnancy. It occurs when pancreatic insulin secretion cannot overcome pregnancy-induced insulin resistance driven by placental hormones, including human placental lactogen, progesterone, and cortisol. Maternal risks include preeclampsia and operative delivery; fetal and neonatal risks include macrosomia, shoulder dystocia, birth injury, neonatal hypoglycemia, hyperbilirubinemia, and respiratory distress. Screening is typically performed between 24 and 28 weeks of gestation using an oral glucose challenge or tolerance test. Management centers on medical nutrition therapy, physical activity, and blood glucose monitoring, with insulin or oral agents added when targets are not met. Glucose levels typically normalize after delivery, but affected individuals carry a significantly increased lifetime risk of developing type 2 diabetes mellitus.

\synonyms
GDM; Gestational Diabetes Mellitus

\medterm{Gestational Diabetes Mellitus} Alternate medical term; see \textit{Gestational Diabetes}.

\medterm{Gestational Hypertension} New-onset high blood pressure after 20 weeks of pregnancy without protein in the urine or other features of preeclampsia. It may later develop into preeclampsia.

\medterm{Gestational Pyelonephritis} A kidney infection occurring during pregnancy, usually causing fever, flank pain, nausea, or urinary symptoms. It can lead to sepsis, preterm labor, or breathing complications.

\synonyms
Pyelonephritis in Pregnancy; Gravid Kidney Infection

\medterm{Gestational Sac} The first fluid-filled structure usually seen inside the uterus on an early pregnancy ultrasound. Its size and appearance help estimate pregnancy age and assess early development, but an early scan may not yet show an embryo.

\medterm{Gestational Trophoblastic Disease} A group of uncommon conditions caused by abnormal growth of cells that normally help form the placenta. It includes molar pregnancy and some pregnancy-related cancers and may cause unusually high pregnancy hormone levels, bleeding, or a uterus larger than expected.

\medterm{Gestational Trophoblastic Neoplasia} Persistent or cancerous growth of placenta-related cells after a pregnancy, miscarriage, abortion, or molar pregnancy. It is often detected by pregnancy-hormone levels that remain high or rise and is usually very responsive to specialist cancer treatment.

\medterm{Gestation Period} The time from conception to birth or, clinically, from the last menstrual period to delivery; human pregnancy is usually about 40 weeks from the last menstrual period.

\medterm{GFR} Glomerular filtration rate, an estimate of how well the kidneys filter blood. It is usually calculated from blood creatinine and other information, and a persistently low value may indicate chronic kidney disease.

\synonyms
Glomerular Filtration Rate; Estimated GFR (eGFR)

\medterm{Ghon Complex} The combination of a primary tuberculous lung focus (Ghon focus) and caseous granulomatous lymphadenitis in the draining regional hilar or mediastinal lymph nodes, representing the hallmark of primary pulmonary tuberculosis that often undergoes fibrocalcification.
\synonyms
Primary Tuberculous Complex; Ranke Precursor Complex

\medterm{Ghon Focus} A primary granulomatous tuberculous pulmonary lesion consisting of a circumscribed area of caseous necrosis typically located subpleurally in the mid-to-lower lung zone, representing the initial site of Mycobacterium tuberculosis infection.
\synonyms
Primary Focus of Ghon; Subpleural Tuberculous Focus

\medterm{Ghrelin} A hormone produced mainly in the stomach that increases hunger and helps regulate energy balance. It also stimulates growth-hormone release and changes with eating, fasting, sleep, and body weight.

\medterm{Gianotti-Crosti Syndrome} A childhood skin condition characterized by a papular or papulovesicular rash on the face, buttocks, and extremities, often following viral infections (EBV, hepatitis B, CMV). Self-limited, lasting 2-8 weeks. Also called papular acrodermatitis of childhood.

\synonyms
Papular Acrodermatitis of Childhood; Infantile Lichenoid Acrodermatitis

\medterm{Giant Cell Arteritis} Inflammation of arteries, usually in people over 50, often affecting vessels near the temples and eyes. It can cause new headache, scalp tenderness, jaw pain while chewing, or vision changes and may permanently damage sight.

\medterm{Giant Cell Arteritis Emergency} A severe presentation of suspected giant cell arteritis with visual symptoms, jaw or tongue pain during chewing, or reduced blood flow to the nervous system. It carries a risk of permanent visual loss.

\medterm{Giant Cell Myocarditis} A rare, severe inflammation of the heart muscle caused by an abnormal immune reaction. It can weaken the heart and cause dangerous heart rhythms or heart failure.

\medterm{Giant Cells} Large cells formed by fusion of macrophages or abnormal cellular growth, seen in granulomatous inflammation, foreign-body reactions, and some tumors.

\medterm{Giant Cell Tumor} A usually noncancerous but locally destructive bone tumor containing many large abnormal cells. It commonly affects bones near major joints, may cause pain or swelling, and can return after treatment or rarely spread to the lungs.

\medterm{Giant Cell Tumor Of Bone} A locally aggressive osteoclast-rich bone tumor, usually arising near the end of a long bone in skeletally mature adults. It may cause pain, swelling, cortical destruction, and occasional pulmonary metastases.

\medterm{Giant Cell Tumor Of Soft Tissue} A usually low-grade soft-tissue neoplasm containing numerous osteoclast-like giant cells, most often arising in the extremities of adults and generally having a favorable outcome after excision.

\medterm{Giant Cell Tumor Of Tendon Sheath} A common benign localized tenosynovial tumor, usually presenting as a slow-growing firm nodule on a finger or hand; recurrence can occur after excision.

\medterm{Giant Congenital Nevus} A very large dark or brown skin mark present at birth. It may grow with the child and has a higher risk of melanoma than smaller birthmarks.

\synonyms
Giant Pigmented Nevus; Garment Nevus; Bathing Trunk Nevus

\medterm{Giant Coronary Aneurysm} An unusually large bulge in an artery supplying the heart. It can form a clot, reduce blood flow, or rarely rupture.

\medterm{Giant Hypertrophic Gastritis} Severe thickening of the stomach lining and its folds, which can cause stomach pain, nausea, swelling, and loss of blood proteins. It may increase stomach-cancer risk.

\medterm{Giant Lymph Node Hyperplasia} A rare condition in which lymph-node tissue grows excessively. It may affect one area or many parts of the body and can cause enlarged nodes, fever, fatigue, anemia, weight loss, or organ problems.

\synonyms
Castleman Disease; Angiofollicular Lymph Node Hyperplasia

\medterm{Giant Platelet Syndrome} An inherited platelet disorder caused by deficiency of glycoprotein Ib/IX/V complex, resulting in large platelets and impaired platelet adhesion. Presents with mucocutaneous bleeding, thrombocytopenia, and giant platelets on blood smear.

\synonyms
Bernard-Soulier Syndrome; Platelet Glycoprotein Ib Deficiency

\medterm{Giant Right Atrium} Marked enlargement of the heart’s right upper chamber, usually from long-standing valve disease, a congenital heart defect, or abnormal rhythm. It can disturb the heartbeat, compress nearby structures, or allow clots to form and is assessed with heart imaging.

\medterm{Giardia} A microscopic parasite that infects the small intestine and spreads through contaminated water, food, or hands. It can cause foul-smelling diarrhea, cramps, bloating, nausea, weight loss, or no symptoms.

\medterm{Giardia Infection} An intestinal infection caused by Giardia parasites, usually spread through contaminated water, food, or contact with an infected person. It can cause prolonged foul-smelling diarrhea, cramps, bloating, weight loss, and poor nutrient absorption.

\medterm{Giardia Lamblia} The parasite that causes giardiasis. It is swallowed in contaminated water, food, or material from an infected person and can cause foul-smelling diarrhea, abdominal cramps, gas, bloating, and weight loss.

\medterm{Giardiasis} An intestinal infection caused by the parasite Giardia duodenalis. It spreads when contaminated water, food, surfaces, or material from an infected person is swallowed. It may cause diarrhea, foul-smelling or greasy stools, gas, abdominal cramps, nausea, bloating, weight loss, and poor nutrient absorption.

\medterm{Gibbus Deformity} A sharp, angular, structural kyphosis of the spine produced by the anterior collapse or destruction of one or more adjacent vertebral bodies, classically seen in tuberculosis of the spine (Pott disease) or severe osteoporotic wedge fractures.
\synonyms
Pott Gibbus; Angular Kyphosis

\medterm{Gibraltar Fever} An old name for brucellosis, an infection caused by Brucella bacteria, usually acquired through unpasteurized dairy products or contact with infected animals.

\medterm{Giemsa Stain} A laboratory dye used to show blood cells and some bacteria, parasites, and other organisms under a microscope. It can help identify malaria parasites and other infections in blood or tissue samples.

\medterm{Gigabecquerel} A unit of radioactivity equal to one billion nuclear disintegrations per second. It measures the activity of a radioactive source, not directly the radiation dose absorbed by a person.

\medterm{Gigantism} Excessive growth caused by too much growth hormone during childhood, before the bones finish lengthening. It usually results from a pituitary growth and causes unusual height and enlargement of body tissues.

\medterm{Gigli Saw} A flexible wire saw used during some orthopedic, neurosurgical, or amputation procedures to cut bone. It is handled carefully to limit injury to nearby soft tissues.

\medterm{Gilbert's Syndrome} A common inherited condition in which the liver processes bilirubin slowly, causing occasional mild yellowing of the eyes, especially during illness, fasting, or stress.

\medterm{Gilbert Syndrome} A common inherited condition in which the liver processes bilirubin more slowly than usual. It can cause occasional mild yellowing of the eyes or skin, especially during fasting, illness, dehydration, or stress, and usually does not harm the liver.

\medterm{Gilles De La Tourette Syndrome} A neurological disorder causing repeated involuntary movements and sounds called motor and vocal tics. Symptoms begin in childhood and often change over time.

\medterm{Gingiva} The firm pink tissue, commonly called the gums, that surrounds and protects the teeth and the bone around their roots. Healthy gingiva usually does not bleed with gentle brushing; redness, swelling, or bleeding may indicate gum disease.

\synonyms
Gums; Gingival Tissue

\medterm{Gingiva Disease} A disorder of the gums, including plaque-induced gingivitis, periodontitis, infection, immune disease, recession, or tumor.

\medterm{Gingival Disease} Inflammation of the gums, usually caused by plaque buildup. The gums may become red, swollen, tender, or bleed during brushing. It usually does not damage the supporting bone and can often improve with daily cleaning and professional dental care.

\synonyms
Marginal Gingivitis; Periodontal Gum Inflammation

\medterm{Gingival Hemorrhage} Bleeding from gum tissue, commonly due to inflammation and plaque but also caused by trauma, periodontal disease, medicines, or bleeding disorders.

\medterm{Gingival Hyperplasia} An abnormal enlargement or overgrowth of gingival tissue around the teeth that can impair cleaning, cause bleeding, and increase periodontal risk. Causes include plaque-related inflammation and medicines such as phenytoin, cyclosporine, and some calcium-channel blockers; hereditary and systemic conditions can also contribute.

\medterm{Gingival Infections} Infections of the gums, usually involving bacteria in dental plaque. They can cause red, swollen, tender, or bleeding gums, bad breath, pus, and sometimes spreading facial swelling; dental assessment helps prevent tissue and bone loss.

\medterm{Gingival Margin} The edge of the gum tissue that surrounds and touches a tooth. It normally forms a shallow seal; inflammation or recession can make it bleed or expose the tooth root.

\medterm{Gingival Melanoma} A rare, aggressive cancer of pigment-producing cells in the gums. It may appear as a dark or non-pigmented lump, patch, or ulcer; biopsy can identify it.

\medterm{Gingival Neoplasm} An abnormal growth in the gums that may be reactive, benign, precancerous, or cancerous. It may cause persistent swelling, ulceration, bleeding, pain, or a lump; examination and sometimes biopsy identify its nature.

\medterm{Gingival Overgrowth} Gingival overgrowth is an abnormal increase in gum tissue, which may be caused by medicines, inflammation, or systemic disease.

\medterm{Gingival Pocket} The small space between the gum and tooth surface, measured during a dental examination. A persistently deep pocket usually indicates loss of attachment from periodontal disease and can collect plaque and bacteria.

\medterm{Gingival Recession} Pulling back of the gums so that part of a tooth root becomes exposed. Gum disease, forceful brushing, tooth position, thin gum tissue, or treatment can cause it and may lead to sensitivity or root decay.

\medterm{Gingivectomy} Surgical removal of gum tissue to treat overgrowth, reshape the gums, expose decay, or improve access. Possible complications include bleeding, sensitivity, infection, or gum recession.

\medterm{Gingivitis} Inflammation of the gingiva characterized by redness, swelling, and bleeding with brushing or probing, usually due to dental plaque. It is generally reversible with effective oral hygiene and professional cleaning; it can precede periodontitis, but does not inevitably progress to it.

\synonyms
ANUG; Trench Mouth; Vincent Stomatitis

\medterm{Gingivitis Transmission} Gingivitis itself is not directly contagious, but the bacteria that cause it (particularly Porphyromonas gingivalis) can be transmitted through saliva. Poor oral hygiene is the primary risk factor. Prevention focuses on personal oral care.

\medterm{Gingivoplasty} Surgical reshaping or contouring of the gums to improve their form, remove diseased tissue, or make cleaning easier. It may be combined with gingivectomy or other periodontal treatment.

\medterm{Gingivostomatitis} Common viral infection of the mouth and gums, most often caused by primary herpes simplex virus infection but also associated with other infections, trauma, or systemic disease. It may cause painful ulcers, swollen bleeding gums, fever, and difficulty eating or drinking.

\medterm{Ginkgo Biloba} An herbal product derived from Ginkgo biloba leaves and marketed for cognitive or circulatory benefits; evidence is variable and it may increase bleeding risk or interact with medicines.

\medterm{Ginseng} A group of plants whose roots are used in traditional medicine and dietary supplements; products vary in composition and may interact with medicines or affect blood glucose and clotting.

\medterm{Ginseng Preparation} A product made from ginseng root and sold as a supplement. Products vary in strength and purity; interactions and side effects can occur.

\medterm{Girdle Pain} Pain that wraps around the chest or abdomen in a band-like pattern. It may result from nerve irritation, spinal disease, shingles, muscle injury, or internal-organ disease.

\medterm{Gitelman Syndrome} A rare inherited kidney disorder that causes the body to lose too much salt and magnesium in urine. It may lead to muscle cramps, weakness, tiredness, thirst, or abnormal heart rhythms.

\medterm{Githism} An old term for poisoning by corn cockle, a weed whose seeds can contaminate grain and cause stomach upset, weakness, or other toxic effects.


\synonyms
Corn-Cockle Poisoning; Agrostemma Githago Intoxication

\medterm{Glabella} The smooth area of the forehead between the eyebrows and above the bridge of the nose.

\medterm{Gland} An organ or group of cells that produces and releases substances. Endocrine glands release hormones into blood, while exocrine glands release substances through ducts or onto a surface.

\medterm{Glanders} A serious infection caused by the bacterium \textit{Burkholderia mallei}, which mainly affects horses and can rarely infect people exposed to infected animals or tissues. It may cause skin sores, fever, pneumonia, or bloodstream infection.

\medterm{Glandular Fever} Alternate name; see \textit{Infectious Mononucleosis}.

\medterm{Glans Penis} The expanded, sensitive distal end of the penis containing the external urethral opening and covered by the foreskin when present.

\medterm{Glanzmann Thrombasthenia} An inherited platelet function disorder caused by deficiency or dysfunction of platelet glycoprotein IIb/IIIa, producing mucocutaneous bleeding despite a normal platelet count.

\medterm{Glasgow Coma Scale} A bedside score describing a person’s level of consciousness by assessing eye opening, speaking, and movement. Scores range from 3 to 15; lower scores indicate less response, but medicines and injuries can affect the result.

\medterm{Glass Eye} A glass eye is an ocular prosthesis made to replace the appearance of an eye that is missing or severely damaged.

\medterm{Glass Ionomer} A tooth-colored dental material that bonds to tooth structure and slowly releases fluoride. It is used for selected fillings, repairs, protective bases, and children’s teeth; durability depends on the location and amount of stress.

\medterm{Glassy Cell Carcinoma} A rare poorly differentiated adenosquamous carcinoma of the cervix with distinctive glassy cytoplasm, prominent nucleoli, and aggressive clinical behavior.

\medterm{Glaucoma} A group of eye diseases in which the optic nerve is damaged, often in association with high pressure inside the eye. Types include open-angle, angle-closure, normal-tension, secondary, and congenital glaucoma. It can cause permanent loss of side or central vision, while early disease may have no symptoms.

\medterm{Glaucoma Drainage Implant} A surgically implanted tube-and-plate device that diverts aqueous humor from the anterior chamber to a subconjunctival reservoir, lowering intraocular pressure in selected glaucomas.

\medterm{Glaze Poisoning} Toxic exposure to ceramic or pottery glaze, which may contain lead, cadmium, or other heavy metals. Symptoms depend on the metal involved and may include abdominal pain, neurological symptoms, and organ damage. Severe cases may be treated with chelation therapy.

\synonyms
Ceramic Glaze Poisoning; Lead or Cadmium Glaze Toxicity

\medterm{Gleason Grading System} A histological grading system for prostate cancer based on glandular pattern differentiation. Scores historically ranged from 2 to 10, but current reporting summarizes them as Grade Groups 1 to 5 (scores 6 to 10), with higher groups indicating more aggressive cancer. Used to guide treatment decisions and predict prognosis.

\synonyms
Gleason Score; Histological Prostate Grade

\medterm{Gleason Score} A grading system estimating prostate cancer aggressiveness by examining tissue patterns under a microscope and combining the two most common patterns.

\synonyms
Gleason Grade Group; Prostate Adenocarcinoma Score

\medterm{Gleason Tumor Grade} A histopathologic grade of prostate adenocarcinoma based on the architectural pattern of the tumor glands; it contributes to the Gleason score and Grade Group.

\medterm{Glenn Procedure} A staged heart operation for some children born with only one effective pumping chamber. It connects the large vein returning blood from the upper body directly to the lung arteries, allowing blood to pick up oxygen while reducing the workload on the heart.

\synonyms
Bidirectional Glenn (BDG) Procedure; Cavopulmonary Anastomosis

\medterm{Glenohumeral Joint} The ball-and-socket shoulder joint formed by the upper-arm bone and shoulder blade. It allows the arm to move in many directions but is relatively easy to dislocate because it sacrifices some stability for movement.

\medterm{Glenoid} The shallow socket on the shoulder blade that joins the head of the upper-arm bone to form the shoulder joint. Its shape and surrounding labrum help balance mobility with stability.
\synonyms
Glenoid Cavity; Scapular Glenoid Fossa

\medterm{Glenoid Cavity} The shallow socket on the shoulder blade that receives the rounded head of the upper-arm bone. A rim of tough tissue deepens the socket, but the joint’s wide range of movement also makes it prone to dislocation.

\medterm{Glenoid Labrum} A fibrocartilaginous rim attached to the peripheral margin of the glenoid fossa of the scapula, which deepens the shallow glenoid cavity to enhance glenohumeral joint stability and serves as the superior anchor for the long head of the biceps brachii tendon.

\synonyms
Glenoid Ligament

\medterm{Gliadin} A group of wheat storage proteins that, together with glutenin, form gluten; gliadin peptides trigger the immune response in celiac disease.

\medterm{Glial Cell} A non-neuronal cell that supports, nourishes, insulates, and regulates neurons; examples include astrocytes, oligodendrocytes, microglia, and Schwann cells.

\medterm{Glial Cells} Support cells of the brain and spinal cord that nourish and protect nerve cells, form insulation, maintain the chemical environment, and help repair tissue. Some brain tumors arise from glial cells.

\medterm{Glial Fibrillary Acidic Protein} A protein found mainly in support cells of the brain and nervous system. Laboratory tests can use it to help identify some brain tumors or detect injury to support cells after brain trauma.

\synonyms
GFAP

\medterm{Glibenclamide} An oral sulfonylurea medicine, also called glyburide, used for type 2 diabetes. It stimulates the pancreas to release insulin and can cause prolonged, sometimes severe low blood sugar and weight gain.

\medterm{Gliclazide} Gliclazide is a sulfonylurea for type 2 diabetes that stimulates pancreatic beta-cell insulin release; hypoglycaemia, weight gain, and dose accumulation in organ impairment are important considerations.

\medterm{GLIM Criteria} An international framework for diagnosing malnutrition. It considers body changes such as weight or muscle loss together with causes such as poor intake, poor absorption, or illness-related inflammation and grades severity according to the findings.

\medterm{Glimepiride} An oral sulfonylurea medicine for type 2 diabetes that increases insulin release from functioning pancreatic beta cells. Low blood sugar is the main serious risk, particularly with missed meals, kidney disease, or other glucose-lowering medicines.

\medterm{Glioblastoma} A fast-growing cancerous tumor of the central nervous system. It forms from glial tissue, which supports nerve cells, and usually develops in the brain of an adult; it can also occur in the spinal cord. The tumor can grow into nearby healthy tissue. Symptoms depend on its location and may include headache, seizures, weakness, changes in speech, memory, vision, balance, or behaviour. Also called glioblastoma multiforme (GBM).

\medterm{Glioblastoma Multiforme} An aggressive cancer arising from supporting cells of the brain or spinal cord. It may cause headache, seizures, personality or memory changes, weakness, or speech problems. Treatment often combines surgery, radiation, and medicines.

\medterm{Glioma} A tumor arising from supporting cells of the brain or spinal cord. Some gliomas grow slowly, while others are aggressive; symptoms and treatment depend on the tumor’s type, location, molecular features, and spread.

\medterm{Gliomatosis Cerebri} A diffuse infiltrative growth pattern of glioma involving multiple brain regions, rather than a separate tumor type in current classifications; it may cause progressive cognitive, motor, or seizure-related symptoms.

\medterm{Gliosarcoma} A rare aggressive brain cancer containing both glial and connective-tissue-like tumor features. Symptoms depend on location and may include headache, seizures, weakness, or cognitive changes; tissue examination identifies the tumor type.

\medterm{Gliosis} A reactive increase in glial cells after injury, infection, ischemia, or other damage to the central nervous system.

\medterm{Glipizide} Glipizide is a short-acting sulfonylurea that lowers blood glucose by stimulating insulin secretion from functioning pancreatic beta cells; hypoglycaemia is the principal adverse effect.

\medterm{Glisson's Capsule} A thin layer of connective tissue surrounding the liver and extending between its sections. It contains nerves and blood vessels and can cause pain when the liver becomes swollen or inflamed.

\medterm{Global Aphasia} A severe language disorder caused by extensive damage to the dominant cerebral hemisphere, impairing speech production, comprehension, reading, and writing.

\medterm{Global Developmental Delay} Noticeably delayed development in at least two areas, such as movement, speech, learning, social skills, or daily activities, in a child younger than five years. Causes vary and may include genetic, neurologic, hearing, vision, or environmental factors; assessment guides support and further testing.

\synonyms
GDD; Neurodevelopmental Delay

\medterm{Global Health} The study and practice of improving health across countries and populations. It includes problems that cross borders, such as outbreaks, unequal access to care, pollution, climate change, malnutrition, and differences in life expectancy.

\medterm{Global Longitudinal Strain} An echocardiogram measurement showing how much the heart muscle shortens lengthwise as it contracts. It can detect early weakening of the left ventricle before the usual pumping measurement changes and is interpreted with the complete heart scan and clinical findings.

\synonyms
GLS; 2D Speckle-Tracking GLS; Left Ventricular Longitudinal Strain

\medterm{Globin} A protein that joins with heme to form hemoglobin and related oxygen-carrying molecules. Different globin chains help determine the type of hemoglobin made at different stages of life.

\medterm{Globozoospermia} A rare sperm disorder in which sperm heads are round and lack a normal acrosome, causing impaired fertilization and male infertility.

\medterm{Globulin} A group of blood proteins that includes antibodies and proteins involved in transport, inflammation, and clotting. Blood tests may divide globulins into several groups; an abnormal level can occur with infection, immune disease, liver disease, kidney disease, or some cancers.

\medterm{Globus} A persistent or intermittent sensation of a lump or tightness in the throat without a true obstructing mass, often called globus pharyngeus.

\medterm{Globus Hystericus} A sensation of a lump or tightness in the throat without a structural obstruction, now called globus sensation; reflux, muscle tension, and anxiety may contribute.

\medterm{Globus Pallidus} The globus pallidus is a group of deep brain structures involved in regulating movement as part of the basal ganglia.

\medterm{Globus Sensation} The persistent or intermittent sensation of a lump or foreign body in the throat without
an anatomic lesion or dysphagia, also termed globus pharyngeus, typically noticed
between meals. There is no impedance to swallowing and no weight loss; the feeling
worsens with stress, anxiety, or laryngopharyngeal reflux.

\medterm{Glomangiomyoma} A benign glomus tumor showing both vascular and smooth-muscle differentiation, usually presenting as a small painful subcutaneous nodule.

\medterm{Glomangiosarcoma} A rare cancerous glomus tumor with abnormal cells and invasive growth. It may occur in the skin or deeper soft tissue.

\medterm{Glomerular Basement Membrane} The thick extracellular matrix layer of the renal filtration barrier interposed between glomerular endothelial cells and podocyte foot processes, composed of type IV collagen, laminin, and heparan sulfate proteoglycans that impart mechanical structural support and negative charge selectivity.

\synonyms
GBM

\medterm{Glomerular Capsule} A cup-shaped structure around a tiny group of blood vessels in the kidney. It collects the fluid filtered from blood and passes it into a kidney tubule, where useful water and substances are reclaimed and waste is removed.

\medterm{Glomerular Filtrate} The watery fluid filtered from blood into the kidney tubules by a glomerulus, a small group of kidney blood vessels. It normally contains water and small dissolved substances but very little protein or no blood cells.

\medterm{Glomerular Filtration} The first step in urine formation, in which pressure moves water and small dissolved substances from blood through kidney filters into the tubules. Blood cells and most proteins normally remain in the bloodstream.

\medterm{Glomerular Filtration Rate} An estimate of how much blood fluid the kidneys filter each minute. It is calculated mainly from blood creatinine and other factors; a persistently low result can indicate chronic kidney disease, but the result is interpreted with other findings.

\medterm{Glomerular Hyperfiltration} An abnormally elevated single-nephron or whole-kidney glomerular filtration rate exceeding physiological upper limits, representing an early hemodynamic adaptive response in early diabetic kidney disease, obesity, and after unilateral nephrectomy that eventually accelerates glomerulosclerosis.

\medterm{Glomeruli} Microscopic clusters of capillaries in the kidneys where blood filtration begins. Each glomerulus is surrounded by a capsule that collects filtered fluid and passes it into a kidney tubule.

\synonyms
Renal Glomeruli; Malpighian Corpuscles

\medterm{Glomerulomegaly} Abnormal enlargement of kidney glomeruli, seen with diabetes, obesity-related hyperfiltration, congenital syndromes, or other causes of glomerular stress.

\medterm{Glomerulonephritis} Inflammation or other immune-related injury of the kidney’s filtering units. It can cause blood or protein in urine, swelling, high blood pressure, and reduced kidney function and may follow an infection or occur with an autoimmune disease.

\synonyms
GN; Glomerular Nephritis

\medterm{Glomerulonephritis Acute} Sudden inflammation of the kidney glomeruli that can cause blood or protein in urine, edema, hypertension, and reduced kidney function.

\medterm{Glomerulopathy} Any disease affecting the kidney’s filtering units. It may cause blood or protein in urine, swelling, high blood pressure, reduced filtration, or kidney failure and may result from diabetes, infection, immune disease, inherited disorders, or other injury.

\medterm{Glomerulosclerosis} Scarring and hardening of the kidney’s filtering units. It can reduce filtration and cause protein in urine, swelling, high blood pressure, or kidney failure; causes include diabetes, inflammation, genetic disease, and other injury.

\medterm{Glomerulus} A tuft of tiny blood vessels inside a kidney capsule where water and small dissolved substances are filtered from blood to begin urine formation. Damage can allow blood or excess protein into urine and reduce kidney function.

\medterm{Glomus Jugulare Tumor} A usually noncancerous but locally invasive growth arising near the jugular vein at the skull base. It may cause pulsating noise in the ear, hearing loss, dizziness, or weakness of nearby nerves.

\synonyms
Jugular Paraganglioma; Tympanojugular Glomus

\medterm{Glomus Tumor} A usually harmless growth arising from cells that help control blood flow and temperature, most often beneath a fingernail or in a finger. It typically causes severe pinpoint pain and sensitivity to cold or pressure.

\medterm{Glomus Tympanicum} A benign vascular tumor arising from the glomus bodies on the promontory of the middle
ear, presenting with pulsatile tinnitus, hearing loss, and a reddish mass behind the
tympanic membrane.

\medterm{Glomus Tympanicum Tumor} A paraganglioma arising from the glomus bodies on the promontory of the middle ear. Benign tumor causing pulsatile tinnitus and conductive hearing loss. Diagnosis by otoscopy and imaging. Treatment includes surgery or observation.

\synonyms
Middle Ear Paraganglioma; Tympanic Glomus

\medterm{Glossal} Relating to the tongue, the muscular organ used for chewing, swallowing, tasting, and speaking.

\medterm{Glossalgia} Pain or burning of the tongue, with causes including trauma, infection, nutritional deficiency, neuropathy, medication, or burning-mouth syndrome.

\medterm{Glossectomy} Surgical removal of all or part of the tongue, performed for malignancies, severe trauma, or an unusually large tongue. Procedures range from partial to total removal. Speech and swallowing rehabilitation may follow.

\medterm{Glossinidae} The family of flies belonging to the genus Glossina, commonly known as tsetse flies. These insects are vectors of trypanosomes that cause African trypanosomiasis (sleeping sickness) in humans and nagana in animals.

\synonyms
Tsetse Flies; Glossina Vectors

\medterm{Glossitis} Inflammation of the tongue, which can cause pain, swelling, and changes in tongue appearance. Causes include nutritional deficiencies (iron, B12, folate), infections (bacterial, viral, fungal), allergic reactions, and autoimmune conditions. Treatment depends on the underlying cause.

\medterm{Glossodynia} Persistent burning, soreness, or tingling in the tongue, lips, palate, or other mouth tissues, often without visible changes. It may occur with dry mouth, altered taste, nutritional deficiency, medicines, nerve problems, or other illness.

\synonyms
Burning Mouth Syndrome; Orofacial Dysesthesia

\medterm{Glossopharyngeal Nerve} The ninth cranial nerve, carrying sensation and taste from the back of the tongue and throat, signals from blood-pressure and blood-chemistry sensors, and motor control for one throat muscle. It also stimulates saliva production.

\medterm{Glossopharyngeal Neuralgia} Recurrent, brief attacks of severe stabbing or electric-shock-like pain in the throat,
tonsillar region, base of the tongue, or ear, caused by irritation of the
glossopharyngeal nerve. Swallowing, coughing, or talking may trigger attacks.

\medterm{Glossoptosis} Backward displacement of the tongue that can narrow the upper airway, particularly in infants with craniofacial disorders.

\medterm{Glottis} The part of the voice box containing the vocal folds and the opening between them. It opens for breathing, closes to protect the airway during swallowing, and vibrates to produce the voice.

\medterm{GLP-1 Agonists} Medicines that mimic a natural gut hormone. They help lower blood glucose after meals, slow stomach emptying, and reduce appetite. They are used for selected people with diabetes or obesity and commonly cause nausea, vomiting, diarrhea, or constipation.

\medterm{GLP-1 Receptor Agonists} A class of medicines that activate the glucagon-like peptide-1 receptor. They increase glucose-dependent insulin release, reduce inappropriate glucagon release, slow gastric emptying, and increase fullness. They are used for selected people with type 2 diabetes or chronic weight management. Examples include semaglutide, liraglutide, and dulaglutide.

\synonyms
GLP-1 Agonists; Incretin Mimetics

\medterm{GLS} Alternate name; see \textit{Global Longitudinal Strain}.

\medterm{Glu} Commonly an abbreviation for glucose in laboratory or clinical notes; in biochemical notation, Glu can also denote the amino acid glutamic acid, so context matters.

\medterm{Glucagon} A hormone made by pancreatic alpha cells that raises blood glucose when it falls too low. It prompts the liver to release stored glucose and can be given for severe hypoglycemia.

\medterm{Glucagon Blood Test} Laboratory measurement of glucagon levels in the blood, used to evaluate suspected glucagonoma, hypoglycemia, or other pancreatic endocrine disorders. Levels are affected by fasting, meals, and medications.

\synonyms
Serum Glucagon Level; Plasma Glucagon Assay

\medterm{Glucagon-like Peptide-1} A hormone released mainly by the intestine after eating. It helps regulate blood glucose by enhancing glucose-dependent insulin secretion, reducing glucagon release, slowing stomach emptying, and influencing appetite.

\medterm{Glucagonoma} A rare pancreatic neuroendocrine tumor that secretes excess glucagon, causing diabetes mellitus (mild, new onset), dermatitis (necrolytic migratory erythema, characteristic painful, erythematous, blistering rash on groin, perineum, extremities), diarrhea, and deep vein thrombosis.

\medterm{Glucagon Signaling} The chain of cell signals started when glucagon binds its receptor. These signals increase cyclic AMP and activate proteins that help the liver release glucose and the body maintain energy balance.

\medterm{Glucan} A complex carbohydrate made of many linked glucose molecules. Different glucans occur in plants, fungi, bacteria, and foods and may act as dietary fiber or interact with the immune system.

\medterm{Glucarpidase} An enzyme that rapidly breaks down methotrexate in the bloodstream and is used for toxic methotrexate concentrations when renal clearance is inadequate.

\medterm{Glucocerebrosidase Deficiency} Gaucher Disease: an inherited lysosomal storage disorder caused by deficient glucocerebrosidase, leading to lipid accumulation in cells and possible enlargement of the spleen and liver, anemia, bone disease, or neurologic disease.

\synonyms
Gaucher Disease; Acid Beta-Glucosidase Deficiency

\medterm{Glucocorticoid} A corticosteroid hormone or medicine that regulates inflammation, immunity, metabolism, and stress responses; prolonged treatment has systemic adverse effects.

\medterm{Glucocorticoid-Induced Osteoporosis} Bone loss and increased fracture risk caused or accelerated by prolonged or high-dose glucocorticoid treatment; risk depends on dose, duration, baseline bone health, and other factors.

\medterm{Glucocorticoid Remediable Aldosteronism} An autosomal dominant cause of secondary endocrine hypertension caused by an unequal crossing-over mutation that fuses the ACTH-responsive promoter of the 11-beta-hydroxylase gene (\textit{CYP11B1}) to the aldosterone synthase coding sequence (\textit{CYP11B2}), rendering aldosterone secretion suppressible by dexamethasone.

\synonyms
GRA; Familial Hyperaldosteronism Type I

\medterm{Glucocorticoid-Remediable Aldosteronism} An autosomal dominant form of secondary hypertension caused by an unequal crossover genetic recombination fusing the promoter of the 11$\beta$-hydroxylase gene (\textit{CYP11B1}) to the coding region of the aldosterone synthase gene (\textit{CYP11B2}). This places aldosterone production under the ectopic regulatory control of adrenocorticotropic hormone (ACTH), resulting in hyperaldosteronism, hypokalemia, and early-onset hypertension that responds promptly to exogenous glucocorticoid suppression.

\synonyms
GRA; Familial Hyperaldosteronism Type I

\medterm{Glucocorticoid Replacement Therapy} Treatment of adrenal insufficiency with a glucocorticoid medicine to replace the cortisol the body lacks. Doses may vary with symptoms, illness, and physical stress.

\medterm{Glucocorticoids} Steroid hormones made by the adrenal glands, along with medicines that act in the same way. They reduce inflammation and immune activity but can cause infection, high blood glucose, bone loss, or adrenal suppression with prolonged use.

\medterm{Glucocorticoid-Sparing Therapy} Use of a non-glucocorticoid treatment strategy to control inflammation while reducing
cumulative corticosteroid exposure. It aims to limit osteoporosis, infection, diabetes,
hypertension, adrenal suppression, and other glucocorticoid complications.

\medterm{Glucometer} A handheld device that measures blood glucose, usually from a small capillary-blood sample obtained by finger prick.

\medterm{Gluconates} Salts or esters of gluconic acid used in some medicines and supplements. Examples include calcium gluconate and ferrous gluconate; their effects, dosing, and risks depend on the specific compound.

\medterm{Gluconeogenesis} A metabolic pathway that generates glucose from non-carbohydrate precursors such as lactate, glycerol, and glucogenic amino acids. It occurs mainly in the liver and, to a variable extent, in the kidneys and other tissues.

\medterm{Glucosamine} A naturally occurring compound used in cartilage formation and sometimes taken as a supplement for joint symptoms.

\medterm{Glucose} A simple sugar that is the main energy source for the body’s cells. It comes from digesting carbohydrates and can also be released from stored energy or made by the liver. Hormones, especially insulin, help keep blood glucose within a healthy range.

\medterm{Glucose-6-Phosphatase} An enzyme in the liver and kidneys that releases glucose from glucose-6-phosphate. Inherited deficiency causes glycogen-storage disease type I, which can produce low blood glucose, enlarged liver, and metabolic problems.

\medterm{Glucose-6-Phosphate Dehydrogenase Deficiency} An X-linked red-blood-cell enzyme disorder. Infections, fava beans, or certain medicines can cause oxidative damage and sudden destruction of red blood cells, leading to hemolytic anemia, jaundice, dark urine, fatigue, or shortness of breath.

\medterm{Glucose-6-Phosphate Dehydrogenase Test} Laboratory test to measure G6PD enzyme activity in red blood cells, used to diagnose G6PD deficiency. Important before prescribing certain medications (primaquine, dapsone) that can cause hemolytic anemia in deficient individuals.

\synonyms
G6PD Assay; Erythrocyte Enzyme Screen

\medterm{Glucose Challenge Test} A screening test for gestational diabetes mellitus, commonly using a nonfasting 50-g glucose drink followed by a one-hour plasma-glucose measurement.
An abnormal screening result is followed by a diagnostic oral glucose tolerance test; the threshold depends on the protocol used.

\synonyms
GCT; One-Hour Glucose Challenge; O'Sullivan Screening Test

\medterm{Glucose Intolerance} Impaired handling of glucose resulting in blood-glucose concentrations above normal but
below the diagnostic threshold for diabetes in a defined testing context. It often
reflects insulin resistance or inadequate insulin secretion and is associated with
increased risk of type 2 diabetes and cardiovascular disease.

\medterm{Glucose Management Indicator} An estimate of the hemoglobin A1c value calculated from average continuous-glucose-monitoring readings. It helps summarize recent glucose exposure but is not the same as a laboratory A1c and may differ from it.

\medterm{Glucose Screening Tests During Pregnancy} Laboratory tests to detect gestational diabetes mellitus, including the oral glucose tolerance test and glucose challenge test. They are usually performed between 24 and 28 weeks of gestation; abnormal results are followed by additional testing or management.

\synonyms
Gestational Diabetes Screening; Maternal Glucose Screen

\medterm{Glucose Sensor} A small sensor placed under the skin that measures glucose in the fluid between cells. Used in continuous glucose monitoring, it sends repeated readings to a receiver or compatible device and can show trends and alerts.

\medterm{Glucose Test} A laboratory or point-of-care test that measures glucose in blood or another sample. It is used to screen for, diagnose, or monitor diabetes and low or high blood sugar; interpretation depends on timing, fasting status, and clinical context.

\medterm{Glucose Tolerance Test} A test of how the body handles a measured glucose drink. Blood samples are taken before and at set times afterward to help diagnose diabetes or gestational diabetes; thresholds depend on the protocol.

\synonyms
Oral Glucose Tolerance Test (OGTT); 3-Hour Diagnostic OGTT

\medterm{Glucose Transport} Movement of glucose across cell membranes through specialized transport proteins. Insulin increases glucose entry into muscle and fat cells, while other transporters allow glucose to enter the brain, liver, kidneys, and intestines.

\medterm{Glucose Transporter} A membrane protein that moves glucose across a cell membrane; transporter defects or altered regulation can affect glucose metabolism.

\medterm{Glucose Urine Test} Laboratory test to detect glucose in the urine (glycosuria). Normally absent; presence indicates blood glucose exceeding renal threshold, suggesting diabetes mellitus or renal tubular dysfunction. May be detected by dipstick or laboratory analysis.

\synonyms
Urinary Glucose; Glycosuria Screen

\medterm{Glucosuria} The presence of glucose in urine, most commonly from blood glucose exceeding the renal reabsorption threshold, diabetes, pregnancy, or a proximal tubular reabsorption defect.

\medterm{Glucosylceramidase} A lysosomal enzyme that breaks down glucosylceramide, a fat-containing cell substance. Inherited deficiency causes Gaucher disease, allowing fatty material to accumulate in the spleen, liver, bone marrow, bones, or nervous system.

\medterm{Glucuronate} The anion or salt of glucuronic acid, a sugar-acid component of connective-tissue substances and a molecule used in hepatic conjugation reactions.

\medterm{Glucuronic Acid} A sugar acid used by the liver to conjugate bilirubin, drugs, hormones, and toxins, increasing their water solubility for excretion.

\medterm{Glucuronide} A compound formed when glucuronic acid is conjugated to a drug, hormone, bilirubin, or other molecule, usually increasing water solubility and facilitating excretion.

\medterm{Glucuronyl Transferase} UDP-glucuronosyltransferase: a liver enzyme involved in phase II drug metabolism, catalyzing glucuronidation of bilirubin, drugs, and toxins to make them water-soluble for excretion. Deficiency causes Gilbert syndrome or Crigler-Najjar syndrome.

\synonyms
UDP-Glucuronosyltransferase; UGT Enzyme

\medterm{Glue Ear} Persistent fluid in the middle ear without acute infection, often causing conductive hearing loss and a feeling of blockage. It is common in children.

\synonyms
Secretory Otitis Media; Otitis Media with Effusion (OME)

\medterm{GLUT4 Transporter} The major insulin- and exercise-regulated glucose transporter isoform expressed primarily in striated muscle (skeletal and cardiac) and adipose tissue, which translocates from intracellular vesicles to the plasma membrane upon insulin stimulation to facilitate cellular glucose uptake.

\medterm{Glutamate} The main excitatory chemical messenger in the brain and spinal cord. It helps nerve cells communicate, learn, and form memories; excessive glutamate activity can contribute to nerve-cell injury.

\medterm{Glutamate Decarboxylase} Glutamate decarboxylase is an enzyme that converts glutamate into the neurotransmitter GABA; antibodies against it are associated with some autoimmune disorders.

\medterm{Glutamate Receptor} A neuronal receptor activated by glutamate, the major excitatory neurotransmitter, including NMDA, AMPA, kainate, and metabotropic receptor classes.

\medterm{Glutamic Acid} An amino acid that acts as a metabolic intermediate and neurotransmitter precursor; its charged form, glutamate, mediates most fast excitatory signaling in the brain.

\medterm{Glutamic Acid Decarboxylase} An enzyme that converts glutamate to GABA; antibodies to it can occur in type 1 diabetes.

\synonyms
GAD; Glutamate Decarboxylase

\medterm{Glutamine} An amino acid used to build proteins, transport nitrogen, and provide energy for some cells, especially cells in the intestine and immune system. The body usually makes enough, but needs may change during severe illness.

\medterm{Glutaral} Glutaral is a chemical disinfectant and tissue fixative that can irritate the skin, eyes, and airways.

\medterm{Glutaric Acidemia Type 1} A rare inherited disorder in which the body cannot break down certain amino acids, allowing harmful substances to build up. Babies may have a large head and can develop sudden brain injury, abnormal movements, or stiffness during illness; newborn screening and dietary treatment can reduce harm.

\synonyms
Glutaric Aciduria Type 1 (GA-1); Glutaryl-CoA Dehydrogenase Deficiency

\medterm{Glutaric Aciduria} An inherited metabolic disorder in which the body cannot properly break down certain amino acids. Type I can cause brain injury, movement problems, and enlarged head size.

\medterm{Glutathione} A small antioxidant molecule made from glutamate, cysteine, and glycine. Reduced glutathione helps control oxidative stress and helps the liver neutralize reactive chemicals, including NAPQI produced from acetaminophen. N-acetylcysteine can restore glutathione after acetaminophen overdose.

\medterm{Glutathione Disulfide} The oxidized form of glutathione, produced after glutathione neutralizes reactive oxygen compounds. Cells normally convert it back to active glutathione to maintain antioxidant protection.

\medterm{Glutathione Peroxidase} A selenium-dependent antioxidant enzyme family that catalyzes the reduction of toxic hydrogen peroxide and organic lipid hydroperoxides to water and corresponding alcohols, utilizing reduced glutathione (GSH) as an essential electron donor to protect cell membranes from oxidative injury.

\synonyms
GPx

\medterm{Glutathione Transferase} A family of enzymes that attaches glutathione to selected medicines, toxins, and products of cell metabolism. This usually makes them less reactive and easier for the body to process or remove.

\medterm{Gluteal} Relating to the buttock or the gluteal muscles, which extend and rotate the hip and help stabilize the pelvis.

\medterm{Gluteal Gait} An abnormal walking gait caused by weakness or paralysis of the hip abductor muscles (gluteus medius and minimus) innervated by the superior gluteal nerve, manifested by a downward pelvic tilt toward the unsupported swing leg during single-leg stance, compensated by leaning the trunk toward the affected stance side.
\synonyms
Trendelenburg Gait; Gluteus Medius Gait

\medterm{Gluteal Injury} Damage to the muscles, tendons, nerves, or soft tissues of the buttock caused by trauma, overuse, injection, or falls. Pain, bruising, weakness, difficulty walking, severe swelling, numbness, inability to bear weight, or bowel or bladder changes may occur.

\medterm{Gluteal Muscle} One of the large muscles of the buttock, especially the gluteus maximus, medius, or minimus. These muscles extend, rotate, and move the hip and help stabilize the pelvis during standing and walking.

\medterm{Gluteal Tendinopathy} Tendon degeneration, microtrauma, and tearing affecting the insertion of the gluteus medius or gluteus minimus tendons onto the greater trochanter of the femur, representing the predominant cause of greater trochanteric pain syndrome in middle-aged and older adults.

\medterm{Gluten} A complex of storage proteins (prolamins and glutelins) found in wheat (gliadin), barley
(hordein), and rye (secalin). Gluten provides elasticity to dough. In susceptible
individuals, gluten triggers an autoimmune response (celiac disease) or non-autoimmune
sensitivity (non-celiac gluten sensitivity, NCGS).

\medterm{Gluten Enteropathy} Celiac disease, an immune reaction to gluten that damages the small intestine and can cause diarrhea, weight loss, anemia, fatigue, or poor nutrient absorption.

\synonyms
Celiac Disease; Celiac Sprue

\medterm{Gluten-Free Diet} A diet that avoids wheat, barley, rye, and foods contaminated with them. It is essential for celiac disease and may help selected people with non-celiac gluten sensitivity; poorly planned diets can lack fiber or nutrients.

\medterm{Gluten Intolerance} Alternate name; see \textit{Non-Celiac Gluten Sensitivity}.

\medterm{Gluten-Sensitive Enteropathy} Alternate name; see \textit{Celiac Disease}.

\medterm{Gluten Sensitivity} Alternate name; see \textit{Nonceliac gluten sensitivity}.

\medterm{Gluteus} Relating to the buttock muscles. The gluteus maximus, medius, and minimus move and stabilize the hip and pelvis during standing, walking, climbing, and other movements.

\medterm{Gluteus Maximus Muscle} Largest, most superficial gluteal muscle forming the buttock bulk, from posterior iliac
gluteal line, sacrum, coccyx, and sacrotuberous ligament to the iliotibial tract and
gluteal tuberosity. Innervated by inferior gluteal nerve (L5-S2). Primary hip extensor
and lateral rotator.

\medterm{Gluteus Medius Muscle} Thick fan-shaped muscle between maximus and minimus, from external ilium to the greater
trochanter. Innervated by superior gluteal nerve (L4-S1). Primary hip abductor essential
for pelvic stability during single-leg stance. Superior gluteal nerve injury causes
Trendelenburg sign with contralateral pelvic drop.

\medterm{Gluteus Minimus Muscle} The smallest and deepest gluteal muscle, extending from the outer ilium to the greater trochanter. It abducts and
medially rotates the hip and helps stabilize the pelvis during walking.

\medterm{Glyburide} An oral sulfonylurea that stimulates pancreatic insulin release for type 2 diabetes; prolonged or severe hypoglycemia is an important risk.

\medterm{Glycaemia} The amount of glucose in the blood. Insulin lowers blood glucose, while glucagon and other hormones raise it when needed. Abnormally high or low levels can cause symptoms and, if persistent, harm organs.

\medterm{Glycated Hemoglobin} A blood test showing the average blood glucose level over roughly the previous two to three months. It is used to diagnose diabetes and monitor treatment, but results may be misleading when red blood cell survival is abnormal.

\synonyms
Hemoglobin A1c; HbA1c; Glycohemoglobin

\medterm{Glycation} A chemical process in which sugar attaches to proteins, fats, or genetic material without the usual enzyme control. Excess glycation during prolonged high blood glucose can damage blood vessels, nerves, kidneys, eyes, and other tissues.

\medterm{Glycemia} The amount of glucose in the blood. It is affected by food, activity, illness, the liver, and hormones such as insulin and glucagon. Too little is hypoglycemia; too much is hyperglycemia.

\medterm{Glycemic Control} The level and pattern of blood-glucose management over time. It includes glucose values, symptoms, and the effects of food, activity, illness, medicines, pregnancy, and other circumstances.

\medterm{Glycemic Index} A ranking of carbohydrate-containing foods according to how quickly they raise blood glucose compared with a reference food. It does not account for the amount of carbohydrate eaten in a serving.

\medterm{Glycemic Index And Diabetes} The glycemic index ranks carbohydrate-containing foods by their effect on blood glucose levels. Low GI foods (\(\leq 55\)) cause gradual glucose rise; high GI foods (\(\geq 70\)) cause rapid spikes. Important for diabetes management and dietary planning.

\synonyms
GI; Glycemic Load Assessment

\medterm{Glycemic Load} A measure that combines a food’s glycemic index with the amount of available carbohydrate in a serving to estimate its effect on blood glucose.

\medterm{Glyceraldehyde} A small sugar with three carbon atoms that appears as the body breaks down glucose and other carbohydrates. It is quickly changed into other substances during the steps that release energy from food.

\medterm{Glyceride} A glyceride is a compound formed when glycerol is linked to one or more fatty acids, including many dietary fats.

\medterm{Glycerin} A colorless, odorless, sweet-tasting liquid also known as glycerol. It attracts water, dissolves some substances, and is used in medicines, skin products, foods, and selected treatments.

\medterm{Glycerol} A colorless, sweet-tasting alcohol that forms part of fats. It is used in medicines, foods, and skin products and can draw water into the bowel when used in some laxatives.
\synonyms
Glycerine; Trihydroxypropane

\medterm{Glycerol Kinase} Glycerol kinase is an enzyme that adds phosphate to glycerol, allowing it to enter pathways for energy production and fat metabolism.

\medterm{Glycerophosphates} Compounds containing glycerol and phosphate that help build cell membranes and participate in energy metabolism, chemical signaling, and other cellular processes.

\medterm{Glycerophospholipids} The major membrane phospholipids built from glycerol, fatty acids, phosphate, and a polar head group, including phosphatidylcholine and phosphatidylethanolamine.

\medterm{Glycerylphosphorylcholine} A choline-containing substance formed when certain membrane fats are broken down. It occurs in body fluids and cells and participates in membrane metabolism and the supply of choline.

\medterm{Glyceryl Trinitrate} A nitrate medicine, also called nitroglycerin, that releases nitric oxide and relaxes blood vessels. It can relieve or prevent angina, but may cause headache, dizziness, or low blood pressure and can interact dangerously with phosphodiesterase-5 inhibitors.

\medterm{Glycine} A small amino acid used to make proteins and heme. It also acts as an inhibitory chemical messenger, especially in the spinal cord and brainstem, helping regulate movement and reflexes.

\medterm{Glycine Receptor} A ligand-gated chloride channel activated by glycine that mediates inhibitory neurotransmission, especially in the spinal cord and brainstem.

\medterm{Glycocalyx} A sugar-rich coating on the outside of many cells, including cells lining blood vessels. It helps protect surfaces, control passage of substances, support cell attachment, and influence inflammation and clotting.

\medterm{Glycogen} The body’s stored form of glucose. It is kept mainly in the liver, which releases it between meals, and in muscles, which use it during activity.

\medterm{Glycogenesis} The process of converting glucose into glycogen for storage, mainly in the liver and muscles. It occurs after eating, when blood glucose and insulin levels rise.

\medterm{Glycogenolysis} The breakdown of stored glycogen. The liver uses it to release glucose into the blood between meals, while muscles use their glycogen locally for energy during exercise.

\medterm{Glycogen Storage Disease} Any of several inherited disorders in which the body cannot make or break down glycogen normally. Depending on the type, these disorders may cause low blood glucose, an enlarged liver, muscle weakness, exercise intolerance, or heart problems.

\medterm{Glycogen Storage Diseases} Inherited disorders that prevent cells from making or breaking down glycogen normally. They can cause low blood sugar, enlarged liver, muscle weakness, or heart disease.

\medterm{Glycohemoglobin} Hemoglobin with glucose chemically attached to it; glycated hemoglobin A1c reflects average blood glucose over roughly the preceding two to three months.

\medterm{Glycolipid} A lipid containing a carbohydrate group that contributes to cell membranes, recognition, signaling, and myelin structure.

\medterm{Glycolysis} A series of reactions in the cytosol that breaks one glucose molecule into two pyruvate molecules. It produces ATP and other molecules used in further energy production and can proceed when oxygen is limited.

\medterm{Glycopeptide} A molecule made of a short protein chain with one or more attached sugar groups. Glycopeptides occur in cell-surface proteins, antibodies, and some medicines, and the attached sugars can affect structure and function.

\medterm{Glycophorin} A sugar-coated protein in the red-blood-cell membrane that helps maintain the cell surface and carries markers used to distinguish blood-cell types.

\medterm{Glycoprotein} A protein with covalently attached carbohydrate chains, important in receptors, antibodies, enzymes, mucus, cell adhesion, and extracellular matrix.

\medterm{Glycopyrrolate} Glycopyrrolate is an anticholinergic medicine used to reduce secretions and, in some inhaled forms, to help control chronic obstructive lung disease; it may cause dry mouth, constipation, or urinary retention.

\medterm{Glycosaminoglycan} A long, unbranched carbohydrate chain found in the extracellular matrix and connective tissues. Abnormal breakdown or accumulation of GAGs occurs in mucopolysaccharidoses.

\medterm{Glycosaminoglycans} Glycosaminoglycans are long chains of sugars found in connective tissue and involved in its structure, hydration, and signaling.

\medterm{Glycoside} A compound in which a sugar is chemically linked to another molecule. Glycosides occur in plants, foods, and medicines; their effects depend on the attached molecule and how the body breaks the bond.

\medterm{Glycosphingolipid} A membrane lipid containing ceramide and one or more sugars, involved in cell recognition, signaling, and nervous-system structure.

\medterm{Glycosuria} The presence of glucose in the urine. It most often occurs when blood glucose is high, but it can also result from pregnancy, certain medicines, or a kidney-tubule problem that prevents normal glucose reabsorption.

\medterm{Glycosylated Hemoglobin} Hemoglobin with glucose attached; the A1C measurement reflects average blood sugar over roughly two to three months and helps diagnose or monitor diabetes.

\synonyms
Hemoglobin A1c; HbA1c; Glycohemoglobin

\medterm{Glycosylation} Attachment of sugar chains to proteins or fats inside cells. This process helps determine their shape, stability, location, and activity; abnormal glycosylation can contribute to inherited disease, cancer, and immune disorders.

\medterm{Glycosylphosphatidylinositol} A lipid-and-sugar anchor that attaches certain proteins to the outer surface of a cell membrane. Defects in making or attaching this anchor can cause rare blood, nerve, or developmental disorders.

\medterm{Glycosyltransferase} Glycosyltransferase is an enzyme that transfers a sugar to a protein, lipid, or other molecule during glycan synthesis.

\medterm{Glycyrrhetinic Acid} An active metabolite of licorice glycosides that inhibits 11-beta-hydroxysteroid dehydrogenase; excess exposure can cause hypertension, low potassium, and fluid retention.

\medterm{Glycyrrhizic Acid} The major saponin in licorice root that is converted to glycyrrhetinic acid and can produce mineralocorticoid-like toxicity when consumed excessively.

\medterm{Glyphosate} A broad-spectrum herbicide that inhibits plant EPSP synthase; concentrated exposure can irritate tissues and cause systemic poisoning, especially after ingestion.

\medterm{Gm} Abbreviation; see \textit{Gram}.

\medterm{GN} Abbreviation; see \textit{Glomerulonephritis}.

\medterm{Gnathic} Relating to the jaw, especially the upper jaw or lower jaw and their development, structure, or position.

\medterm{Gnathion} The lowest midline point on the outer surface of the lower jaw. It is used as a landmark in dental, orthodontic, facial, and skull measurements.

\medterm{Gnathitis} Inflammation of the jaw or jaw tissues. It may cause pain, swelling, tenderness, difficulty opening the mouth, or chewing problems and can result from infection, injury, or inflammatory disease.

\medterm{Gnathoplasty} Plastic or reconstructive surgery to reshape or repair the jaw. It may improve function, facial balance, or a structural defect; risks include infection, nerve injury, altered bite, and recurrence.

\medterm{Gnathostoma} A group of roundworms that can infect people who eat raw or undercooked freshwater fish, frogs, or other carriers. Larvae may migrate through skin, organs, or the nervous system and cause swelling, pain, or serious neurologic disease.

\medterm{Gnathostoma Spinigerum} A roundworm that can cause gnathostomiasis after people eat raw or undercooked freshwater animals. Migrating larvae may produce recurring painful skin swellings and, rarely, eye or brain disease.

\medterm{Gnathostomiasis} A parasitic infection caused by Gnathostoma larvae acquired from undercooked freshwater fish or other hosts, producing migratory skin swellings or deeper neurologic or ocular disease.

\medterm{Gnawing Pain} A dull, persistent pain often described as an ache or a feeling that something is eating away at the body. It can arise from muscles, nerves, internal organs, or inflammation. Its location, duration, triggers, and accompanying symptoms help identify the cause.

\medterm{GNE Myopathy} A rare inherited muscle disease that usually begins in early adulthood and slowly weakens the muscles of the lower legs, causing foot drop, then the hands and arms. The thigh muscles are often spared until late disease. There is no cure, but rehabilitation helps maintain function.

\synonyms
Nonaka Myopathy; Distal Myopathy with Rimmed Vacuoles (DMRV); Hereditary Inclusion Body Myopathy (HIBM)

\medterm{GnRH} Gonadotropin-releasing hormone, a hormone from the brain that signals the pituitary gland to release hormones controlling ovulation, sperm production, and sex-hormone production.

\synonyms
Gonadotropin-Releasing Hormone

\medterm{GnRH Agonist} A medicine that first increases and then suppresses signals from the pituitary gland to the reproductive organs. It lowers sex-hormone production and is used in selected fertility treatments, endometriosis, fibroids, and hormone-sensitive cancers.

\medterm{GnRH Antagonist} A medicine that directly blocks the brain-to-pituitary signal controlling reproductive hormones. It lowers these hormones quickly without the temporary hormone rise that can occur when a GnRH agonist is started.

\medterm{GnRH Antagonists} Medicines that block a brain hormone signal and quickly reduce pituitary release of reproductive hormones. They are used in some fertility treatments and hormone-sensitive cancers, with effects depending on the specific medicine.

\medterm{Goal Attainment Scaling} A way to measure progress toward personal healthcare goals. The goal is defined in advance with levels ranging from much less than expected to much better than expected, then scored after treatment.

\medterm{Goal Setting} Choosing clear, realistic outcomes and deciding how and when to work toward them. In healthcare, goals can reflect a person’s priorities, include a way to measure progress, and change as needs or circumstances change.

\medterm{Goals-Of-Care Discussion} A clinical conversation that aligns treatment recommendations with a patient's values,
priorities, prognosis, acceptable burdens, and preferred outcomes. It may include
decisions about hospitalization, resuscitation, artificial nutrition, and
comfort-focused care.

\medterm{Goblet Cell} A mucus-secreting epithelial cell, abundant in the respiratory and intestinal tracts, that protects surfaces with a mucin-rich layer.

\medterm{Goblet Cell Adenocarcinoma} A distinctive appendiceal adenocarcinoma composed of goblet-like mucinous cells, formerly called goblet cell carcinoid. It can spread within the peritoneum or to the ovaries and is managed according to grade, stage, and dissemination.

\medterm{Goblet Cells} Cells lining the intestines and airways that make and release mucus. Mucus moistens and protects these surfaces and helps trap particles, germs, and other substances for removal.

\medterm{Goeckerman Regimen} A dermatological therapy used for extensive or recalcitrant plaque psoriasis combining the daily topical application of crude coal tar with exposure to ultraviolet B (UVB) phototherapy.
\synonyms
Goeckerman Therapy; Coal Tar and UVB Regimen

\medterm{Goiter} Enlargement of the thyroid gland, either throughout the gland or as one or more lumps. It may occur with normal, overactive, or underactive thyroid function and can result from iodine deficiency, immune disease, or a tumor.

\synonyms
Thyromegaly; Struma; Thyroid Enlargement

\medterm{Goitre} Enlargement of the thyroid gland. It may occur with normal, low, or high thyroid hormone levels and can result from iodine deficiency, autoimmune disease, nodules, or inflammation.

\medterm{Goitrogenic Foods} Foods or food components that can interfere with thyroid hormone production or iodine use under some conditions; risk depends on intake, preparation, iodine status, and thyroid health.

\medterm{Gold} A metal used in selected dental restorations, implants, laboratory applications, and historical medicinal preparations; exposure can cause contact allergy or toxicity depending on the compound and dose.

\medterm{Gold Alloy} A mixture of gold and other metals used for some dental restorations. It is durable, fits closely, and is usually well tolerated by body tissues, but it costs more and has a distinctive appearance.

\medterm{Gold Alloys} Metal mixtures containing gold and other elements, used in selected dental restorations or medical devices; properties and biocompatibility depend on composition.

\medterm{Goldenhar Syndrome} A birth condition, also called oculo-auriculo-vertebral spectrum, in which one side of the face, ear, eye, or spine develops differently. Features may include facial asymmetry, a small or misshapen ear, hearing loss, eyelid or eye-surface growths, and spinal differences.

\medterm{Golden Hour} The golden hour is the early period after a serious injury or medical emergency when rapid assessment and treatment may improve the chance of survival and recovery.

\medterm{Goldmann Tonometry} A slit-lamp examination that estimates intraocular pressure by measuring the force needed to flatten a defined area of the cornea.

\medterm{Gold Therapy} Treatment with gold-containing medicines, formerly used mainly for rheumatoid arthritis. Its use is now limited because benefit is slow and serious effects can include skin or mouth inflammation, kidney injury, and reduced blood-cell counts.

\medterm{Golfer Elbow} An overuse injury of the tendons on the inner side of the elbow. Repeated gripping or wrist movement can cause pain and tenderness, usually worsened by lifting, bending the wrist, or making a fist.

\medterm{Golfer's Elbow} Painful irritation or degeneration of tendons attached to the inner side of the elbow. Repeated gripping or wrist movement can cause tenderness, reduced strength, and pain during lifting or bending.

\medterm{Golgi Apparatus} A structure inside cells that modifies, sorts, and packages proteins and fats made by the cell, sending them to other parts of the cell or releasing them outside it.

\medterm{Golgi Cisternae} Flattened membrane-bound sacs that form the Golgi apparatus and modify, sort, and package proteins and lipids for transport within or outside the cell.

\medterm{Golgi Tendon Organ} A sensory structure where a muscle joins its tendon. It detects tension in the muscle and tendon and helps the nervous system control muscle force and protect against excessive strain.

\medterm{Gomori's Stain} Gomori's stain is a group of laboratory stains used to highlight structures such as connective tissue, carbohydrates, fungi, or muscle components in tissue samples.

\medterm{Gomphosis} A fixed joint in which a tooth is held in its socket by a strong ligament.

\medterm{Gonad} A testis or ovary, an organ producing reproductive cells and sex hormones. Gonadal function influences puberty, fertility, menstrual cycles, sexual development, bone health, and other body processes.

\medterm{Gonadal Agenesis} Congenital absence or severe underdevelopment of a gonad, resulting in impaired sex-hormone production and infertility or differences in sexual development.

\medterm{Gonadal Dysgenesis} Incomplete or abnormal development of the ovaries or testes, often resulting from
chromosomal abnormalities, gene variants, or disordered sex development. It may cause
atypical genital development, delayed or absent puberty, infertility, and abnormal
gonadal hormone production.

\medterm{Gonadal Vein} A vein draining a testis or ovary. The right gonadal vein usually drains directly into the inferior vena cava, while the left usually drains into the left renal vein.

\medterm{Gonad Disorder} A disease affecting testes or ovaries and causing abnormal reproductive-cell production, hormone secretion, development, fertility, or sexual function.

\medterm{Gonadoblastoma} A rare tumor that develops in an abnormal or incompletely developed ovary or testicle. It may produce sex hormones and has a risk of becoming a malignant germ-cell tumor, especially when a Y-chromosome change is present.

\medterm{Gonadorelin} Synthetic gonadotropin-releasing hormone that stimulates pituitary release of luteinizing hormone and follicle-stimulating hormone when given intermittently.

\medterm{Gonadotrophin} A hormone that acts on the ovaries or testes. The pituitary hormones follicle-stimulating hormone and luteinizing hormone regulate egg development, ovulation, sperm production, and sex-hormone release.

\medterm{Gonadotrophin-Releasing Factors} Signals from the brain that control release of the pituitary hormones acting on the ovaries or testes. The main regulator is gonadotropin-releasing hormone, released in pulses from the hypothalamus.

\medterm{Gonadotrophins} Hormones that act on the ovaries or testes. Follicle-stimulating hormone helps develop eggs and sperm, while luteinizing hormone triggers ovulation and supports sex-hormone production.

\medterm{Gonadotrophs} Pituitary cells that release luteinizing hormone and follicle-stimulating hormone after stimulation by GnRH, regulating the ovaries, testes, fertility, and sex-hormone production.

\medterm{Gonadotropin} A hormone, chiefly luteinizing hormone or follicle-stimulating hormone, that acts on the ovaries or testes to regulate gamete production and sex-hormone secretion.

\medterm{Gonadotropin Hormone} A hormone that acts on the ovaries or testes. The pituitary gonadotropins follicle-stimulating hormone and luteinizing hormone regulate reproductive development, ovulation, sperm production, and sex-hormone production.

\medterm{Gonadotropin-Releasing Hormone} A hormone made in the brain that signals the pituitary gland to release two hormones, FSH and LH. These hormones control the ovaries and testes, including egg development, ovulation, sperm production, and sex-hormone production.

\synonyms
GnRH

\medterm{Gonadotropins} Hormones that act on the gonads. In humans, the pituitary gonadotropins are follicle-stimulating hormone and luteinizing hormone, which regulate gamete production and sex-hormone secretion.

\medterm{Gonads} The reproductive organs: testes in males and ovaries in females. They produce sperm or eggs and make sex hormones such as testosterone, estrogen, and progesterone. These hormones influence puberty, fertility, menstrual cycles, pregnancy, and bone health.

\medterm{Gondoic Acid} A 20-carbon monounsaturated fatty acid found in some seed oils and incorporated into dietary and tissue lipids.

\medterm{Goniometer} An instrument measuring joint angles and range of motion, used to assess stiffness, document impairment, and track progress during rehabilitation.

\medterm{Gonioscopy} An eye examination using a special contact lens to view the angle where the iris meets the cornea. It helps determine whether glaucoma is open-angle or angle-closure and guides treatment.

\medterm{Gonococcal} Relating to Neisseria gonorrhoeae, the bacterium that causes gonorrhea. Gonococcal infection commonly affects the genitals, rectum, or throat and may spread to joints or blood.

\medterm{Gonococcal Arthritis} A disseminated infection of joints caused by Neisseria gonorrhoeae, often with fever, migrating joint or tendon pain, skin pustules, and inflammation around tendons.

\medterm{Gonococcal Pharyngitis} Infection of the throat by Neisseria gonorrhoeae, usually acquired through sexual contact. It may
cause sore throat or no symptoms.

\medterm{Gonococcal Urethritis} Infection and inflammation of the urethra caused by Neisseria gonorrhoeae. It may cause burning urination or discharge, but some people have no symptoms; untreated infection can spread.

\medterm{Gonococcus} A spherical bacterium of the genus Neisseria. Neisseria gonorrhoeae, the gonococcus, causes the sexually transmitted infection gonorrhea.

\medterm{Gonorrhea} A sexually transmitted infection caused by the bacterium Neisseria gonorrhoeae. It may cause discharge, painful urination, pelvic or rectal symptoms, or no symptoms. Untreated infection can lead to infertility, joint infection, or infection in a newborn.

\medterm{Gonorrhoea} Alternate spelling; see \textit{Gonorrhea}.

\medterm{Goodell Sign} A clinical sign of early pregnancy characterized by marked softening of the uterine cervix (produced by increased pelvic vascularity and edema), identifiable upon digital pelvic bimanual examination around the sixth to eighth week of gestation.

\medterm{Goodpasture's Syndrome} An autoimmune disease in which antibodies attack the lungs and kidney filters, causing coughing blood, kidney inflammation, blood in urine, and sometimes rapid kidney failure.

\medterm{Goodpasture Syndrome} An autoimmune disease in which antibodies attack the lung air sacs and kidney filters. It can cause coughing blood, breathlessness, blood in the urine, and rapidly worsening kidney failure.

\medterm{Good Samaritan} A person who voluntarily helps someone in an emergency, especially by providing immediate aid before professional responders arrive.

\medterm{Good Samaritan Law} A law that may protect a person who provides reasonable emergency assistance from certain civil claims, subject to jurisdiction-specific conditions and exceptions such as gross negligence.

\medterm{Goose Bump} A small raised area around a hair follicle caused by contraction of tiny arrector pili muscles, commonly in response to cold, fear, or strong emotion.

\medterm{Gooseflesh} Temporary raised bumps around hair follicles caused by contraction of tiny skin muscles, often in response to cold, fear, or strong emotion. It is usually harmless.

\medterm{Goose-Skin} Goose skin is the temporary raising of hair follicles, usually caused by cold, fear, or strong emotion.

\medterm{Gorlin's Syndrome} An inherited disorder that increases the risk of multiple basal-cell skin cancers, jaw cysts, and abnormalities of the bones, eyes, and nervous system.

\medterm{Gorlin Syndrome} Alternate name; see \textit{Gorlin's Syndrome}.

\medterm{Goserelin} A gonadotropin-releasing hormone agonist that initially stimulates and then suppresses pituitary gonadotropins; it is used in hormone-sensitive prostate or breast cancer and selected gynecologic conditions.

\medterm{Gottron Papules} Raised red, purple, or scaly patches over the knuckles and other finger joints. They are a characteristic skin finding of dermatomyositis, an inflammatory disease that can also cause muscle weakness.

\medterm{Gottron Sign} Symmetrical, violaceous, macular or papular erythematous patches and scaling located over the extensor aspects of the joints (particularly the elbows, patellae, and medial malleoli), serving as a characteristic cutaneous hallmark of dermatomyositis.
\synonyms
Gottron's Sign; Dermatomyositis Extensor Erythema

\medterm{Gougerot-Carteaud Syndrome} A skin disorder causing brown, slightly raised spots that join into net-like patches, usually on the upper chest, back, or neck of young people.

\medterm{Gout} A type of inflammatory arthritis caused by a buildup of urate, a substance made when the body breaks down purines. When urate levels stay high, needle-shaped crystals can form in and around joints. These crystals cause sudden attacks, or flares, of severe pain, swelling, warmth, redness, and tenderness, often in the big toe but also in other joints. Repeated flares can damage joints. Long-standing gout may cause firm crystal deposits called tophi and uric acid kidney stones.

\synonyms
Gouty Arthritis; Monosodium Urate Arthropathy

\medterm{Gout Flare} A sudden attack of gout caused by inflammation around urate crystals in a joint. The joint becomes very painful, swollen, warm, and red, often beginning in the big toe. Infection can cause similar symptoms.

\medterm{Gout Nodulosis} A condition in which many firm urate deposits, called tophi, form under the skin or around joints because of long-standing gout. They may occur at the ears, fingers, toes, elbows, or Achilles tendons and can restrict movement or break through the skin.

\medterm{Gout Suppressants} Medicines that lower uric-acid production or increase its excretion to prevent recurrent gout attacks and urate crystal deposition.

\medterm{Gouty Tophus} A nodular, chalky subcutaneous or periarticular deposit of crystallized monosodium urate monohydrate crystals surrounded by chronic foreign-body granulomatous inflammation, occurring in chronic uncontrolled hyperuricemia on the helix of the ear, olecranon bursa, Achilles tendon, and distal joints.

\synonyms
Tophus; Tophi

\medterm{Gowers Sign} A maneuver in which a person uses the hands to push up from the floor or thighs because of proximal muscle weakness, classically seen in muscular dystrophy.

\medterm{Gower Syndrome} A form of reflex syncope, also called vasovagal or vasodepressor syncope, in which a trigger causes a fall in blood pressure and sometimes heart rate, briefly reducing brain perfusion.

\medterm{G-Protein Coupled Receptors} Proteins on the cell surface that receive signals from hormones, medicines, and other substances and pass those signals into the cell. They help control heart rate, vision, digestion, mood, and inflammation.

\medterm{Graafian Follicle} A mature, fluid-filled sac in an ovary that contains an egg ready for release. It enlarges during the menstrual cycle and normally opens at ovulation; after release, the remaining tissue forms the corpus luteum.

\medterm{GRACE Score} A risk score used in people with an acute coronary syndrome to estimate the chance of death or another heart attack during the hospital stay and over the next six months. It uses age, vital signs, kidney function, cardiac arrest, ECG changes, heart-failure signs, and heart-injury blood tests to guide treatment intensity.

\synonyms
Global Registry of Acute Coronary Events Score; GRACE Risk Model

\medterm{Gracilis Muscle} A long, slender strap-like muscle located on the medial aspect of the thigh originating from the inferior pubic ramus and body of the pubis, descending to insert into the pes anserinus on the proximal medial tibia; acts to adduct the thigh and flex and internally rotate the leg.

\medterm{Grade} A measure of how abnormal cells look under a microscope and, in many cancers, how quickly the growth is likely to behave. Lower grades usually look more like normal cells; higher grades usually look more abnormal. The grading system differs by disease.

\medterm{Graded Potentials} Small changes in an electrically active cell caused by a stimulus. Their size depends on the stimulus strength, and they can either increase or decrease the chance that the cell will send a nerve signal.

\medterm{Grade IV Astrocytoma} An older term for glioblastoma, an aggressive brain tumor that grows into nearby brain tissue. It may cause headache, seizures, weakness, speech or vision problems, or changes in thinking. Treatment often combines surgery, radiation, and medicines.

\synonyms
Glioblastoma; GBM

\medterm{Gradenigo Syndrome} A classic clinical triad consisting of suppurative otitis media (otorrhea), retro-orbital or temporal pain in the sensory distribution of the trigeminal nerve (cranial nerve V), and ipsilateral abducens nerve palsy (cranial nerve VI) caused by petrous apicitis.
\synonyms
Gradenigo's Syndrome; Petrous Apicitis Triad

\medterm{Graefe's Sign} A delayed downward movement of the upper eyelid during downward gaze, classically associated with thyroid eye disease but not specific to it.

\medterm{Graft} Tissue or a synthetic material transferred to another site to replace, repair, reinforce, or bypass a damaged body structure.

\medterm{Grafting} The surgical transfer of tissue or a synthetic substitute to repair, replace, or support a body structure. An autograft comes from the same person, whereas an allograft comes from another person of the same species.

\medterm{Graft Versus Host} Alternate name; see \textit{Graft Versus Host Disease}.

\medterm{Graft Versus Host Disease} A complication of a stem-cell or bone-marrow transplant in which immune cells from the donor attack the recipient’s tissues. It commonly affects the skin, liver, and digestive tract and may be short-term or long-lasting; prevention and treatment use immune-suppressing medicines.
\synonyms
GVHD; Transfusion-Associated or Allogeneic GVHD

\medterm{Graft-Versus-Host Disease} Alternate spelling; see \textit{Graft Versus Host Disease}.

\medterm{Graft-Versus-Host Reaction} An immune reaction in which donor immune cells attack the recipient's tissues after transplantation, potentially affecting skin, intestines, liver, lungs, or other organs.

\medterm{Gram} A metric unit used to measure mass, equal to one-thousandth of a kilogram. Grams are commonly used for medicines, laboratory samples, food, and body substances; one thousand milligrams make one gram.

\synonyms
Metric Gram; Unit G

\medterm{Gramicidin} Gramicidin is a topical antibiotic that forms channels in bacterial membranes; it is not generally used systemically because it can be toxic.

\medterm{Gram Mass Unit} A metric unit of mass equal to one-thousandth of a kilogram; medication doses and laboratory measurements are commonly expressed in grams or smaller metric units.

\medterm{Gram Metric Unit} A metric unit of mass equal to one-thousandth of a kilogram or approximately 15.432 grains.

\medterm{Gram-Negative} Describing bacteria with a thin peptidoglycan cell wall and an outer membrane containing lipopolysaccharide; they stain pink or red with the Gram stain.

\medterm{Gram-Negative Bacilli} Rod-shaped bacteria that appear pink or red with a Gram stain because they have a thin peptidoglycan layer and an outer membrane. They include harmless body organisms and bacteria that can cause urinary, respiratory, intestinal, bloodstream, wound, or other infections.

\medterm{Gram-Negative Bacteria} Bacteria that do not retain crystal violet during Gram staining and therefore appear pink or red.
They have a thin peptidoglycan layer and an outer membrane that may contain
lipopolysaccharide. Examples include Escherichia coli, Klebsiella, Pseudomonas, and
Neisseria.

\medterm{Gram-Negative Cocci And Coccobacilli} Gram-negative bacteria that are either round or short rod-shaped. Cocci are round, while coccobacilli are short rods that may look almost round. They include bacteria that can cause respiratory, genital, bloodstream, intestinal, and other infections.

\medterm{Gram-Negative Meningitis} Meningitis caused by gram-negative bacteria, including Escherichia coli, Klebsiella, and Neisseria species. It can progress rapidly to sepsis, neurologic injury, and death.

\medterm{Gram-Positive} Describing bacteria with a thick peptidoglycan cell wall that retains crystal violet during Gram staining and therefore appears purple.

\medterm{Gram-Positive Bacilli} Rod-shaped bacteria that retain the purple part of the Gram stain because of their thick peptidoglycan cell wall. Some are harmless colonizers, while others can cause wound, intestinal, respiratory, bloodstream, or nervous-system infections.

\medterm{Gram-Positive Bacteria} Bacteria that retain crystal violet in the Gram stain (appear purple/blue),
characterized by a thick peptidoglycan cell wall and no outer membrane.

\medterm{Gram-Positive Cocci} Round bacteria that appear purple with a Gram stain because they have a thick peptidoglycan cell wall. Important groups include staphylococci, streptococci, and enterococci, which may normally live on the body or cause skin, throat, lung, urinary, heart-valve, or bloodstream infections.

\medterm{Gram's Stain} Gram's stain is a laboratory method that classifies bacteria by how their cell walls retain the stain, helping guide identification and initial antibiotic choices.

\medterm{Gram Stain} A laboratory test that colors bacteria to group them as Gram-positive or Gram-negative. It helps identify possible causes of infection quickly and can guide initial treatment while more specific tests are performed.

\medterm{Gram Stain Of Skin Lesion} Microscopic examination of a skin lesion specimen stained with Gram stain to identify bacteria. Useful for diagnosing bacterial skin infections and guiding antibiotic therapy. Results available within hours.

\medterm{Gram Stain Of Tissue Biopsy} Microscopic examination of tissue biopsy specimen stained with Gram stain to identify bacteria. Used to diagnose tissue infections and guide antimicrobial therapy. May be performed alongside culture and histopathology.

\synonyms
Tissue Gram Stain; Histochemical Bacterial Stain

\medterm{Grand Mal} Alternate historical term; see \textit{Tonic-Clonic Seizure}.

\medterm{Grand Mal Seizure} Alternate name; see \textit{Tonic-clonic seizure}.

\medterm{Grand Multipara} A person who has given birth five or more times beyond the gestational threshold used to count parity; definitions can vary by obstetric system.

\medterm{Grand Rounds} An educational medical conference in which clinicians present a case, diagnosis, management, or clinical topic to a professional audience.

\medterm{Granisetron} Granisetron is a selective 5-HT3-receptor antagonist used to prevent chemotherapy-, radiotherapy-, or surgery-related nausea and vomiting; headache, constipation, and QT prolongation can occur.

\medterm{Granisetron Hydrochloride} The hydrochloride form of granisetron, a medicine that blocks serotonin 5-HT3 receptors to prevent nausea and vomiting from chemotherapy, radiation, or surgery. Headache, constipation, and heart-rhythm changes may occur.

\medterm{Granular Casts} Small grainy particles formed inside kidney tubules and seen in urine. They may occur when kidney cells are injured or break down and can be found in acute kidney injury or some chronic kidney diseases.

\medterm{Granular Cell Tumor Of The Breast} An uncommon tumor showing Schwann-cell differentiation that can occur in the breast and mimic carcinoma clinically or on imaging. Histologic examination distinguishes it from malignancy; most are benign.

\medterm{Granular Leukocyte} A granulocyte, a white blood cell containing cytoplasmic granules; the main types are neutrophils, eosinophils, and basophils.

\medterm{Granular Parakeratosis} A rare skin condition causing brown or reddish, sometimes itchy patches or papules, often in skin folds such as the armpits. It reflects abnormal keratinization and may be difficult to distinguish from other rashes.

\medterm{Granulation} The formation of new vascular connective tissue during wound healing, producing moist red or pink granulation tissue made of capillaries, fibroblasts, and extracellular matrix.

\medterm{Granulation Tissue} New tissue that grows in a healing wound and contains small blood vessels, connective tissue, and immune cells. Healthy granulation tissue is usually pink or red and moist.

\medterm{Granulicatella} Granulicatella is a group of bacteria that may live in the mouth and can rarely cause bloodstream infection or endocarditis.

\medterm{Granulicatella Adiacens} A mouth bacterium that can rarely cause serious infection, especially infection of heart valves or the bloodstream. It may be difficult to grow in culture, and treatment depends on susceptibility testing.

\medterm{Granulicatella Elegans} A bacterium normally found in the mouth that can rarely cause bloodstream infection or infection of a heart valve. It may grow poorly in routine culture, so repeated positive cultures and compatible symptoms are important when judging its significance.

\medterm{Granulocyte} A type of white blood cell containing visible granules. Neutrophils, eosinophils, and basophils are granulocytes that help fight infection, trigger inflammation, and respond to allergic or parasitic disease.

\medterm{Granulocyte Colony-Stimulating Factor} A hematopoietic glycoprotein growth factor that stimulates the bone marrow to produce and differentiate granulocyte precursors into mature neutrophils and mobilize them into circulating blood; recombinant forms (filgrastim, pegfilgrastim) are used to treat and prevent chemotherapy-induced neutropenia.

\synonyms
G-CSF

\medterm{Granulocyte Count} A granulocyte count measures the number of neutrophils, eosinophils, and basophils in the blood and helps assess infection, inflammation, and bone-marrow function.

\medterm{Granulocytic Sarcoma} A mass of immature myeloid blood-forming cells outside the bone marrow, also called myeloid sarcoma. It may occur with acute myeloid leukemia or precede it.

\medterm{Granulocytopenia} An abnormally low number of granulocytes, especially neutrophils, which can increase susceptibility to bacterial and fungal infection.

\medterm{Granuloma} A small area of organized inflammation formed when immune cells surround material they cannot easily remove, such as certain germs, foreign material, or substances released by the body.

\medterm{Granuloma Annulare} A benign, noninfectious inflammatory skin condition that causes raised red, pink, or skin-colored bumps arranged in rings, often with clearer skin in the center. It commonly affects the hands, feet, or elbows and may cause little or no discomfort. Widespread disease may be associated with diabetes or thyroid disease, but the cause is often unknown.

\medterm{Granuloma Inguinale} A chronic sexually transmitted infection causing painless, slowly enlarging ulcers in the genital or groin region. It is caused by bacteria and spreads through sexual contact; antibiotics are used in treatment.

\medterm{Granulomatosis} A disease in which long-lasting inflammation forms granulomas, small organized collections of immune cells. Depending on the cause, granulomas may affect the lungs, skin, lymph nodes, intestines, or other organs.

\medterm{Granulomatosis With Polyangiitis} An autoimmune disease that inflames and damages small and medium-sized blood vessels, often affecting the nose, sinuses, lungs, and kidneys. It may cause sinus problems, cough, coughing blood, or kidney injury.
\synonyms
GPA; Wegener Granulomatosis

\medterm{Granulomatous Amebic Encephalitis} A rare, severe brain infection caused by free-living amoebae, usually in people with weakened immunity. It develops over days or weeks and may cause headache, confusion, seizures, weakness, or coma.

\medterm{Granulomatous Colitis} Inflammation of the colon in which granulomas are found on tissue examination, often associated with Crohn disease but also possible with infections or other inflammatory conditions.

\medterm{Granulomatous Inflammation} A pattern of long-lasting inflammation in which immune cells form small clusters around material they cannot easily remove. It can occur with infections, foreign materials, autoimmune disease, or other inflammatory conditions.

\medterm{Granulomatous Mastitis} A chronic inflammatory breast disease characterized by noncaseating granulomas after other causes, especially infection and sarcoidosis, have been excluded. It may present as a painful mass, abscesses, or draining sinuses.

\medterm{Granulosa Cell} A granulosa cell is a cell surrounding an ovarian egg that supports its development and produces hormones such as estrogen and inhibin.

\medterm{Granulosa Cell Tumor} A tumor arising from hormone-producing cells in the ovary. It may release estrogen, causing irregular bleeding or early puberty, and can recur years after treatment even when it initially appears slow-growing.

\medterm{Granulosa Tumor} Alternate form; see \textit{Granulosa Cell Tumor}.

\medterm{Granzyme} A protease released by cytotoxic T lymphocytes and natural killer cells that enters target cells through perforin-associated pores and promotes programmed cell death.

\medterm{Grasp Reflex} An automatic closing of a baby’s fingers when the palm is touched. It is normal in early infancy and usually fades as voluntary hand control develops; persistence or asymmetry can occur with nervous-system problems.

\medterm{Grass Allergy} Hypersensitivity reaction to grass pollen, causing allergic rhinitis (hay fever) with symptoms including sneezing, nasal congestion, itchy eyes, and throat irritation. Seasonal, typically spring through fall. Treatment includes antihistamines, nasal corticosteroids, and allergen avoidance.

\synonyms
Grass Pollen Allergy; Gramineae Pollinosis

\medterm{Grass And Weed Killer Poisoning} Illness caused by swallowing, inhaling, or getting herbicide on the skin or eyes. Effects range from irritation, vomiting, and breathing problems to severe poisoning, depending on the product and amount.

\synonyms
Herbicide Intoxication; Weedicide Poisoning

\medterm{Graves Disease} An autoimmune disease that causes an overactive thyroid (hyperthyroidism). The immune system makes antibodies that act like thyroid-stimulating hormone and make the thyroid produce too much thyroid hormone. Symptoms may include weight loss, a fast or irregular heartbeat, sweating, heat intolerance, tremor, nervousness, trouble sleeping, frequent bowel movements, muscle weakness, or an enlarged thyroid (goiter). Some people develop Graves eye disease, which can cause dry or gritty eyes, puffiness, bulging eyes, light sensitivity, or double vision. Rarely, the skin over the shins becomes thickened and red.

\medterm{Graves Ophthalmopathy} An autoimmune inflammatory disorder of the retro-orbital tissues and extraocular muscles associated with Graves disease, driven by autoantibodies against the TSH receptor leading to fibroblastic proliferation, glycosaminoglycan deposition, proptosis, upper lid retraction, and diplopia.
\synonyms
Graves Orbitopathy; Thyroid Eye Disease

\medterm{Grave Wax} Alternate name; see \textit{Adipocere}.

\medterm{Gravida} The number of times a person has been pregnant, regardless of pregnancy outcome; the current pregnancy is included in the count.

\medterm{Gravidity} Gravidity is the number of times a person has been pregnant, including the current pregnancy and pregnancies that did not result in a live birth.

\medterm{Gray} A unit used to measure the amount of ionizing radiation absorbed by a material or body tissue. One gray equals one joule of radiation energy absorbed per kilogram. It is used in radiation treatment and safety assessment.

\synonyms
Absorbed Dose Unit Gy; Radiation Unit Gray

\medterm{Gray Baby Syndrome} A serious form of chloramphenicol toxicity in newborns, caused by immature drug metabolism and characterized by abdominal distension, vomiting, cyanosis or gray skin color, cardiovascular collapse, and potentially death.

\medterm{Gray Matter} Areas of the brain and spinal cord containing many nerve-cell bodies and connections. Gray matter processes information and helps control movement, sensation, memory, speech, and other functions.

\medterm{Greater Curvature} The long, rounded outer border of the stomach, extending from its upper opening to its outlet. It provides attachment for the greater omentum and receives blood from several stomach arteries.

\medterm{Greater Omentum} A fold of fatty peritoneal tissue that hangs from the stomach and partly covers the abdominal organs.

\medterm{Greater Trochanter} The large bony prominence on the outer upper thigh, where several hip muscles attach. It can be felt through the skin and is a common site of pain from tendon problems, bursitis, or injury.

\medterm{Greater Trochanteric Pain Syndrome} Pain over the lateral hip centered on the greater trochanter, commonly caused by tendinopathy or bursitis of the gluteal tendons. Presents with lateral hip pain worse with walking, climbing stairs, or lying on the affected side. Treatment includes physical therapy and activity modification.

\synonyms
GTPS; Trochanteric Bursitis; Gluteal Tendinopathy

\medterm{Greater Vestibular Glands} Paired glands near the vaginal opening that secrete mucus into the vestibule, contributing to lubrication. They are also called Bartholin glands.

\synonyms
Bartholin Glands; Vulvovaginal Glands

\medterm{Great Saphenous Vein} The longest superficial vein of the lower limb, ascending from the medial foot to drain into the femoral vein at the saphenofemoral junction.

\medterm{Transposition of the Great Vessels} A cyanotic congenital heart defect resulting from embryological failure of the aorticopulmonary septum to spiral, such that the aorta arises directly from the right ventricle and the pulmonary artery originates from the left ventricle, creating two parallel, independent circulations requiring a mixing communication for survival.

\synonyms
TGA; D-Transposition of the Great Arteries; d-TGA

\medterm{Green Blindness} A form of red-green color-vision deficiency in which distinguishing green from certain other colors is difficult, usually because of an inherited cone-photopigment difference.

\medterm{Green Nail Syndrome} Green discoloration of one or more nails, usually caused by colonization or infection with Pseudomonas aeruginosa, often in an onycholytic or chronically moist nail.

\medterm{Greenstick Fracture} An incomplete fracture in which a young bone bends and cracks rather than breaking all the way through.

\medterm{Green Stool} Stool with a green color that may result from diet, food coloring, rapid intestinal transit, medicines, or infection.

\medterm{Grenz Zone} A narrow, uninvolved zone of normal collagenous dermis located immediately beneath the epidermis, separating the epidermal basement membrane from a dense underlying dermal cellular infiltrate, classically seen in lepromatous leprosy, granuloma faciale, and cutaneous lymphomas.
\synonyms
Subepidermal Grenz Zone; Clear Dermal Zone

\medterm{Grey Ankle Phenomenon} A gray or dusky discoloration around the ankle associated with chronic venous insufficiency and hemosiderin deposition. It occurs with other signs of venous disease.

\medterm{Grey Matter} Brain and spinal-cord tissue containing many nerve-cell bodies, connections, and supporting cells. It processes information and helps control movement, sensation, thinking, memory, and other functions.

\medterm{Grey Turner Sign} Bruising on the sides of the belly between the ribs and the hips, which is a sign of internal bleeding from severe pancreas inflammation or an injured blood vessel.

\medterm{Grey Turner's Sign} Bruising or blue-purple discoloration of the skin at the sides of the lower back. It can occur when bleeding deep inside the abdomen spreads through the tissues, including in severe pancreatitis.

\medterm{Grid Implant} A flexible electrode array placed on or near the brain or another tissue surface to record electrical
activity or deliver stimulation.

\medterm{Grief} The natural emotional response to loss, particularly the death of a loved one, but also following other significant losses such as divorce, job loss, or declining health. It encompasses a range of feelings including sadness, anger, guilt, and numbness, often occurring in waves rather than linear stages.

\medterm{Grief and Bereavement} The natural emotional response to loss, encompassing the process of coping with the death of a loved one or other significant losses. Includes grief (emotional response), loss (the event), and bereavement (the period of mourning). May involve physical, emotional, social, and spiritual dimensions.

\synonyms
Bereavement; Mourning; Complicated Grief

\medterm{Grimontia Hollisae} A salt-loving bacterium from seawater and seafood, formerly called \textit{Vibrio hollisae}. Eating contaminated raw or undercooked seafood can cause watery diarrhea, abdominal cramps, nausea, and fever; severe infection is uncommon but more likely in people with serious underlying illness.

\medterm{Griping} Cramping or spasmodic abdominal pain, often caused by contractions, gas, inflammation, or obstruction of a hollow digestive organ.

\medterm{Grip Strength} The force generated when the hand closes around an object, commonly measured with a dynamometer to assess hand function, muscle strength, nutrition, or rehabilitation progress.

\medterm{Griscelli Syndrome} A rare inherited disorder causing unusually light or silvery hair and skin pigmentation. Some forms also cause neurologic problems or severe immune overactivation, which can lead to fever, enlarged organs, low blood counts, and life-threatening illness.

\medterm{Griseofulvin} Griseofulvin is an oral antifungal that deposits in newly formed keratin and disrupts dermatophyte mitosis, treating infections of skin, hair, and nails; treatment is prolonged and hepatic or drug-interaction risks apply.

\medterm{Groin} The region where the lower abdomen meets the upper thigh, containing the inguinal canal and related muscles, vessels, nerves, and lymph nodes.

\medterm{Groin Lump} A palpable mass in the inguinal region, which may represent enlarged lymph nodes, a hernia, cyst, tumor, or another condition.

\synonyms
Inguinal Mass; Inguinal Lymphadenopathy

\medterm{Groin Pain} Pain in the inguinal region, which may originate from musculoskeletal, genitourinary, gastrointestinal, or vascular causes. Common causes include muscle strain, hernia, enlarged lymph nodes, and hip disease.

\synonyms
Inguinal Pain; Pubalgia

\medterm{Grommet Tube} A miniature, spool-shaped synthetic ventilation tube surgically inserted through an incision in the tympanic membrane (myringotomy) to equalize middle ear pressure and facilitate fluid drainage in chronic otitis media with effusion.

\synonyms
Tympanostomy Tube; Ventilation Tube

\medterm{Groove} A long, narrow depression or channel on a body structure, such as a groove on a bone for a tendon or a groove between folds of the brain. Its exact function depends on its location.

\medterm{Gross Examination} Inspection of a tissue or other specimen with the unaided eye before it is examined under a microscope. The pathologist records its size, shape, color, and visible abnormalities to help guide diagnosis.

\medterm{Gross Hematuria} Visible blood in the urine, appearing pink, red, or brown. Causes include urinary tract infections, kidney stones, bladder or kidney cancer, trauma, and disease of the kidney filters.

\synonyms
Macroscopic Hematuria; Visible Hematuria

\medterm{Gross Motor Control} The ability to use large muscle groups for coordinated movements such as walking, running, jumping, and maintaining posture. Develops progressively from infancy through childhood and is assessed as part of developmental milestones.

\synonyms
Large Motor Coordination; Fundamental Motor Skills

\medterm{Gross Motor Skills} Movements that use the large muscles for posture, balance, walking, running, jumping, and climbing. Children develop these skills at different rates.

\medterm{Ground-Glass Opacity} A hazy area on a chest CT scan in which underlying lung vessels remain partly visible. It is a radiologic finding, not a diagnosis, and may result from infection, inflammation, fluid, bleeding, scarring, or sometimes a tumor; its meaning depends on the pattern and clinical context.

\medterm{Ground Itch} An itchy skin rash caused by hookworm larvae entering the skin, usually after bare-foot contact with contaminated soil. It produces small, red, intensely itchy bumps.

\medterm{Ground Substance} The gel-like, noncellular part of connective tissue surrounding cells and fibers. It contains water, proteoglycans, and other molecules that support tissues, allow nutrient movement, and influence inflammation and repair.

\medterm{Group A Strep} A bacterium that commonly causes sore throat and skin infections. It can occasionally cause scarlet fever, pneumonia, severe soft-tissue infection, bloodstream infection, or immune complications affecting the heart or kidneys.

\medterm{Group A Streptococcal Infections} Infections caused by Streptococcus pyogenes, including pharyngitis, scarlet fever, impetigo, cellulitis, pneumonia, necrotizing soft-tissue infection, and bloodstream infection. Immune-mediated complications can include rheumatic fever and poststreptococcal glomerulonephritis.

\medterm{Group A Streptococcus} Streptococcus pyogenes, a bacterium that can cause pharyngitis, skin infections, invasive disease, and immune-mediated complications.

\medterm{Group B Strep} A bacterium that can live harmlessly in the bowel or genital tract but can cause serious infection in newborns during or soon after birth. It can also cause infection in pregnant people and vulnerable adults. Pregnancy screening and care during labor can reduce newborn risk.
\synonyms
Group B Streptococcus; Streptococcus Agalactiae

\medterm{Group B Strep Disease} Illness caused by group B Streptococcus bacteria, which can cause serious infections in newborns and infections in pregnant people or adults with medical risks.

\medterm{Group B Streptococcal Disease} Infection caused by group B Streptococcus bacteria, which can be serious in newborns and vulnerable adults.

\medterm{Group B Streptococcal Infection In Pregnancy} Maternal colonization or infection with Streptococcus agalactiae, which can be transmitted to the newborn during labor and cause serious early-onset infection. Screening and intrapartum management reduce neonatal risk.

\medterm{Group B Streptococcal Septicemia Of The Newborn} Serious bloodstream infection caused by Group B Streptococcus in neonates, typically acquired during birth. It can cause fever, lethargy, poor feeding, and breathing difficulty; treatment includes antibiotics and supportive care.

\synonyms
Early-Onset GBS Sepsis; Neonatal Streptococcus Agalactiae Sepsis

\medterm{Group B Streptococcus} A bacterium that may live harmlessly in the bowel or genital tract but can cause serious infection in newborns during or soon after birth. It can also cause infection in pregnant people and vulnerable adults; pregnancy screening helps reduce newborn risk.

\medterm{Group B Streptococcus in Pregnancy} Asymptomatic colonization of the maternal rectovaginal tract with \textit{Streptococcus agalactiae}. Routine culture screening is performed between 36 and 37 weeks of gestation so that colonized individuals can receive intrapartum intravenous antibiotic prophylaxis to prevent neonatal early-onset sepsis, pneumonia, and meningitis.

\synonyms
Maternal GBS Colonization; Intrapartum GBS Prophylaxis; Maternal Group B Strep

\synonyms
Maternal GBS Colonization; Intrapartum GBS Prophylaxis

\medterm{Group Practice} A healthcare organization in which several clinicians share facilities, staff, records, or administrative services. Patients may receive coordinated care from different professionals within the same practice.

\medterm{Group Therapy} A form of psychotherapy where a small group of individuals with similar issues meet to discuss, interact, and develop coping strategies under the guidance of a trained therapist. Used for depression, anxiety, substance use disorders, grief, and chronic illness. Benefits include peer support and shared experiences.

\medterm{Grover Disease} A transient or recurrent acantholytic dermatosis causing itchy red papules or papulovesicles, usually on the trunk of middle-aged or older adults; heat, sweating, and occlusion may aggravate it.

\medterm{Growing Pains} Recurrent, bilateral leg pain in children ages 3-12, typically occurring in the late
afternoon or evening and resolving by morning. Pain is deep, aching in the calves,
shins, and thighs (not joints). No limping, fever, swelling, or functional impairment.

\medterm{Growth} The process of increasing in physical size and development, including physical, cognitive, and emotional maturation. In medicine, growth is monitored using growth charts and anthropometric measurements to assess health and development in children.

\medterm{Growth Plate} A layer of cartilage near the end of a growing bone where the bone lengthens. It gradually closes after puberty.

\medterm{Growth Chart} A graphical representation of growth (weight, height/length, head circumference, BMI)
plotted against age, standardized by population (WHO growth standards 0-5 years; CDC
growth charts 2-20 years; percentiles/z-scores).

\medterm{Growth Cone} A specialized, moving tip at the end of a growing nerve fiber. It senses chemical and physical signals and guides the axon toward its target during nervous-system development or repair.

\medterm{Growth Disorder} A condition in which a child or adult grows unusually slowly, quickly, or unevenly because of genetic, hormonal, nutritional, chronic, or other medical causes.

\medterm{Growth Factor} A protein that sends signals controlling cell growth, division, specialization, survival, or repair. Different growth factors act on different tissues and are involved in wound healing, blood-vessel formation, immunity, and cancer.

\medterm{Growth Failure} Growth that is slower than expected or weight gain that is inadequate for age. It can result from poor nutrition, chronic illness, hormone problems, genetic conditions, or difficulty absorbing nutrients.

\medterm{Growth Faltering} A pattern of inadequate weight gain or growth in a child. It is the preferred modern term for many cases previously called failure to thrive and may result from feeding difficulty, inadequate intake, poor absorption, chronic disease, or other causes.

\medterm{Growth Hormone} A hormone released by the anterior pituitary that promotes growth partly through insulin-like growth factor 1 and also affects protein, fat, and carbohydrate metabolism.

\medterm{Growth Hormone Deficiency} Insufficient secretion of growth hormone from the anterior pituitary, causing short
stature in children and metabolic effects in adults. In children: delayed growth, short
stature, delayed bone age, and increased fat mass. In adults: decreased muscle mass,
increased fat mass, reduced bone density, fatigue, and dyslipidemia.

\medterm{Growth Hormone Physiology} Growth hormone is released in pulses by the pituitary gland, especially during sleep. It acts directly on tissues and also stimulates the liver to produce insulin-like growth factor 1, which supports bone and tissue growth.

\medterm{Growth Hormone-Releasing Hormone} A 44-amino-acid hypothalamic peptide hormone synthesized by neurons of the arcuate nucleus that travels via the hypophyseal portal system to bind G-protein coupled receptors on anterior pituitary somatotropes, stimulating growth hormone synthesis and pulsatile secretion.

\synonyms
GHRH; Somatoliberin

\medterm{Growth Hormone Stimulation Test} A test for suspected growth-hormone deficiency. A medicine or other controlled stimulus is given, and blood samples are taken over time to see whether the pituitary gland releases enough growth hormone. Results depend on the protocol and laboratory.

\synonyms
GH Provocation Test; Insulin Tolerance Test for GH

\medterm{Growth Hormone Suppression Test} A test for excess growth hormone, usually performed after drinking a glucose solution. In healthy people, glucose lowers growth-hormone levels; failure to suppress them supports a diagnosis such as acromegaly.

\synonyms
GH Suppression Test; Glucose Loading GH Assay

\medterm{Growth Hormone Test} Laboratory measurement of growth hormone levels in the blood, used to evaluate GH deficiency or excess. Random GH levels have limited utility; stimulation or suppression tests are more informative.

\synonyms
Serum Growth Hormone; Somatotropin Assay

\medterm{Growth Medium} A nutrient preparation supporting microbial multiplication in vitro, supplying carbon,
nitrogen, minerals, and growth factors in liquid or solid form.

\medterm{Growth Plate Fracture} A break involving the cartilage growth area of a child’s or adolescent’s bone. It can disturb future growth if displaced or untreated; imaging and specialist management may be used.

\synonyms
Physeal Fractures; Salter-Harris Fractures

\medterm{Growth Plate Fractures} Injuries involving the cartilage growth plates at the ends of growing children's bones, with possible effects on future bone growth.

\medterm{Growth Retardation} Older wording for growth that is slower than expected for age or developmental stage. Possible causes include inadequate nutrition, chronic illness, hormone or genetic disorders, poor absorption, and normal family variation; growth is assessed over time rather than by one measurement.

\medterm{Growth Scan} An ultrasound examination used to assess fetal size, growth pattern, amniotic fluid, and sometimes
blood flow. Measurements are compared with gestational-age standards to identify possible growth restriction or excessive
growth.

\medterm{Gtp Binding} GTP binding is the attachment of guanosine triphosphate to a protein, often switching that protein into an active state for cell signaling.

\medterm{Gtt.} An abbreviation from the Latin \textit{guttae}, meaning drops. It refers to a liquid dose or the number of drops.

\medterm{Guaifenesin} An expectorant used to loosen respiratory mucus and make coughing more productive. It provides
symptomatic relief in some respiratory illnesses; dosing depends on the formulation and patient.

\medterm{Guanfacine} Guanfacine is a centrally acting alpha-2A adrenergic agonist used for hypertension and, in extended-release form, ADHD; sedation, hypotension, bradycardia, and rebound hypertension after abrupt withdrawal are important risks.

\medterm{Guanine} Guanine is a purine nitrogenous base in DNA and RNA; it pairs with cytosine and is present in nucleotides such as GMP and GTP that support energy transfer and signaling.

\medterm{Guanosinetriphosphatase} An enzyme that hydrolyzes guanosine triphosphate to guanosine diphosphate and phosphate, regulating signaling and other cellular processes.

\medterm{Guanosine Triphosphate} A guanine nucleotide that provides energy and acts as a molecular switch in protein synthesis, cytoskeletal regulation, and G-protein signaling.

\medterm{Guanylate Cyclase} Guanylate cyclase is an enzyme that converts GTP to cyclic GMP, a second messenger involved in nitric-oxide signaling, smooth-muscle relaxation, phototransduction, and intestinal fluid secretion.

\medterm{Guarana} A plant whose seeds contain caffeine and other stimulants. It is used in energy products and supplements; excessive intake may cause anxiety, insomnia, rapid heartbeat, high blood pressure, or interactions with medicines.

\medterm{Guardianship} A legal arrangement in which an appointed person makes decisions for someone unable to make them independently.

\synonyms
Legal Guardianship; Conservatorship

\medterm{Guarding} Involuntary tightening of the abdominal muscles when the abdomen is examined, often reflecting irritation or inflammation of the abdominal lining.

\medterm{Guevedoces} A term used in the Dominican Republic for individuals with 46,XY differences in sex development caused by 5-alpha-reductase type 2 deficiency; they may have female-appearing or ambiguous external genitalia at birth and increased virilization at puberty.

\medterm{Guide Device} A tool used to direct, position, or stabilize another instrument during a procedure, such as a guidewire, drill guide, or needle guide. It helps ensure accurate placement and reduces injury to surrounding tissue.

\medterm{Guided Imagery} Guided imagery is a relaxation technique that uses imagined sights, sounds, or other sensations to promote calm and sometimes help manage pain or stress.

\medterm{Guided Tissue Regeneration} A dental procedure placing a barrier membrane to encourage regrowth of bone, root covering, and supporting ligament around a tooth. It may repair selected periodontal defects after inflammation is controlled.

\medterm{Guillain-Barre Syndrome} A rare disorder in which the immune system attacks peripheral nerves, often after an infection. Weakness and tingling usually begin in the legs and may spread upward; severe cases affect breathing or swallowing.

\synonyms
GBS; Acute Inflammatory Demyelinating Polyneuropathy (AIDP)

\medterm{Guillain-Barré Syndrome} An uncommon immune-mediated disorder in which peripheral nerves are attacked, usually causing rapidly progressive weakness and reduced reflexes. It can impair breathing or swallowing and is monitored closely in severe cases.

\medterm{Guinea Worm} A parasitic infection caused by Dracunculus medinensis, acquired by drinking water containing infected copepods and characterized by emergence of a long worm through the skin.

\medterm{Gulf War Syndrome} A term used for a group of persistent symptoms reported by some Gulf War veterans, such as fatigue, pain, cognitive difficulties, gastrointestinal symptoms, and respiratory problems. It is also called Gulf War illness or chronic multisymptom illness; causes are complex and treatment focuses on individual symptoms and function.

\medterm{Gullet} The esophagus, a muscular tube carrying swallowed food and liquid from the throat to the stomach. It moves contents by coordinated squeezing and is protected from injury by its moist inner lining.

\synonyms
Esophagus; Food Pipe

\medterm{Gum} The soft tissue, also called gingiva, that surrounds and protects the teeth and covers the jawbone. Healthy gums are usually firm and do not bleed with gentle brushing.

\medterm{Gum Biopsy} Histological examination of gingival tissue obtained by biopsy, used to diagnose gingival diseases, infections, or malignancies. May involve excisional or incisional biopsy techniques.

\synonyms
Gingival Biopsy; Periodontal Tissue Sampling

\medterm{Gum Disease} Inflammation or infection affecting the gums and tissues supporting the teeth. Early disease may cause bleeding and swelling, while advanced disease can destroy bone, loosen teeth, and cause tooth loss.

\synonyms
Periodontal Disease; Periodontitis

\medterm{Gumma} A soft area of tissue damage caused by late untreated syphilis. It can develop in the skin, bones, liver, or other organs and may cause a lump, pain, or organ damage.

\synonyms
Syphilitic Gumma; Tertiary Syphilis Granuloma

\medterm{Gummata} More than one gumma: soft inflammatory masses caused by late untreated syphilis. They can affect skin, bone, liver, or other organs and may damage tissue.

\medterm{Gummy Smile} A smile showing more gum than usual. It may result from lip movement, tooth eruption, jaw shape, or gum position; treatment is optional and depends on the cause and the person’s goals.

\medterm{Gum Problems} Gum problems include inflammation, bleeding, recession, ulcers, infection, trauma, and growths. Causes include plaque, tobacco, dry mouth, medicines, nutritional or immune disease, and poorly fitting appliances.

\medterm{Gums} The firm pink tissue, also called gingiva, surrounding and supporting the teeth. Healthy gums usually do not bleed with gentle brushing; swelling or bleeding can indicate gum disease.

\synonyms
Gingiva; Gingival Mucosa

\medterm{Gunn Sign} An ophthalmologic sign of hypertensive retinopathy in which an arteriosclerotic, rigid retinal arteriole crosses and compresses an underlying retinal venule, producing visible tapering and apparent deflection of the venule on both sides of the crossing.
\synonyms
Gunn's Sign; Arteriovenous Nicking

\medterm{Gunshot Wound} An injury caused by a firearm projectile. Damage may include bleeding, tissue destruction, fractures, nerve injury, organ damage, and infection.


\synonyms
Firearm Injury; Ballistic Trauma; Bullet Wound

\medterm{Gunther Disease} A severe autosomal recessive disorder of heme biosynthesis caused by a deficiency of uroporphyrinogen III synthase, leading to massive accumulation of isomer I porphyrins in tissues, manifesting with extreme cutaneous photosensitivity, blistering, mutilating scar formation, erythrodontia, and hemolytic anemia.
\synonyms
Congenital Erythropoietic Porphyria; CEP

\medterm{Gustation} The sense of taste. Special cells in taste buds detect sweet, salty, sour, bitter, and savory flavors. Most taste buds are on the tongue, with some in the soft palate and throat. Signals travel through several head nerves to the brain. Taste can change with infections, medicines, smoking, aging, or damage to the mouth or nerves.

\medterm{Gustatory} Relating to taste or the sense of taste, including taste buds, taste nerves, and the brain pathways that interpret flavors.

\medterm{Gustatory Lacrimation} Excessive tearing triggered by eating or smelling food, usually after an injury to the facial nerve or recovery from facial paralysis. The nerve signals for saliva and tears become incorrectly linked.

\medterm{Gustatory Sweating} Sweating triggered by eating or thinking about food, usually over the cheek or temple after injury or surgery involving the parotid region; it is also called Frey syndrome.

\medterm{Gut} Alternate informal term; see \textit{Digestive Tract}.

\medterm{Gut-Brain Axis} The two-way communication between the digestive tract and brain through nerves, hormones, immune signals, and substances made by gut microorganisms. It influences digestion, appetite, stress responses, and mood, but its clinical importance is still being studied.

\medterm{Gut Epithelium} The layer of cells lining the intestines. It absorbs nutrients and water, forms a barrier against harmful substances, and interacts with immune cells and gut microorganisms.

\medterm{Guthrie Test} An older name for a newborn blood-spot screening test, originally used to detect phenylketonuria. Modern screening covers more conditions; an abnormal result is followed by confirmatory testing.

\medterm{Gut Microbiome} The complex community of trillions of microorganisms, including bacteria, viruses, and
fungi, residing in the gastrointestinal tract. The gut microbiome plays a crucial role
in digestion, immune function, vitamin synthesis, and protection against pathogens.

\medterm{Guttae} Small wart-like bumps on the inner surface of the cornea, the clear front window of the eye. Numerous or severe guttae can impair fluid removal, causing corneal swelling, blurred vision, glare, or eye pain.

\medterm{Guttate Psoriasis} A form of psoriasis causing many small, scaly, drop-shaped spots, often after a streptococcal throat infection. It commonly affects children or young adults and may clear, recur, or later become plaque psoriasis.

\medterm{Guttural} Relating to the throat. A guttural sound is produced mainly in the throat or back of the mouth rather than at the front of the tongue and lips.

\medterm{Guyon Canal Syndrome} An entrapment neuropathy of the ulnar nerve as it passes through Guyon canal (the fibro-osseous ulnar tunnel at the wrist between the pisiform and hook of the hamate), causing sensory loss over the little finger and medial half of the ring finger and motor weakness of the hypothenar and interosseous muscles.
\synonyms
Ulnar Tunnel Syndrome; Handlebar Palsy

\medterm{Gymnast Wrist} A chronic overuse repetitive stress injury of the distal radial growth plate (physis) occurring in skeletally immature gymnasts and athletes who subject the upper extremities to intense axial loading and hyperextension. Radiographs typically demonstrate distal radial physeal widening, irregular metaphyseal margins, and cystic changes.

\synonyms
Distal Radial Physeal Stress Syndrome; Gymnast's Epiphysitis

\medterm{Gynaecologist} A medical doctor who specializes in the health of the female reproductive system
(vagina, cervix, uterus, fallopian tubes, ovaries) and breasts.

\medterm{Gynaecology} The medical specialty concerned with the female reproductive system, including the ovaries, fallopian tubes, uterus, cervix, vagina, and vulva, as well as related hormonal health.

\medterm{Gynaecomastia} Alternate name; see \textit{Gynecomastia}.

\medterm{Gynatresia} Absence of or blockage in an opening of the female reproductive tract, present from birth or acquired later. If menstrual blood cannot leave the body, it may cause absent periods, pelvic pain, or enlargement of the uterus or vagina.

\medterm{Gynecic} Relating to the female reproductive system or to gynecology.

\medterm{Gynecoid Pelvis} A commonly described pelvic shape with a rounded inlet, relatively broad dimensions, and a spacious outlet; it is associated with typical female pelvic anatomy but individual variation is substantial.

\medterm{Gynecological Examination} A gynecological examination assesses the female reproductive organs and may include inspection, a pelvic examination, and screening tests when appropriate.

\medterm{Gynecologist} A gynecologist is a physician who provides care for the female reproductive system, including preventive care, diagnosis, and treatment.

\medterm{Gynecology} The medical specialty concerned with the health of the female reproductive system
(uterus, ovaries, fallopian tubes, cervix, vagina, and vulva).

\medterm{Gynecomastia} A benign enlargement of male breast tissue caused by proliferation of glandular ductal components and stroma secondary to an increased ratio of estrogen to androgen activity. It commonly occurs as a physiological finding in neonates, pubertal adolescents, and older men. Pathological causes include primary or secondary hypogonadism (such as Klinefelter syndrome), liver cirrhosis, chronic kidney disease, estrogen-producing tumors, and medications such as spironolactone, antiandrogens, and anabolic steroids.

\synonyms
Gynaecomastia; Male Breast Enlargement

\medterm{Gynecomastia, Male} Alternate name; see \textit{Gynecomastia}.

\medterm{Gyrate Atrophy} An inherited retinal disorder causing gradually worsening peripheral and night vision loss. It results from abnormal ornithine metabolism and may be managed with dietary treatment, supplements, and regular eye monitoring.

\medterm{Gyrus} A raised fold on the surface of the brain, separated from neighboring folds by grooves called sulci. These folds increase the cerebral cortex’s surface area within the skull and contain areas responsible for different functions.



\medterm{21-Hydroxylase Deficiency} An inherited condition in which the adrenal glands cannot make enough of certain hormones. It can cause salt loss, low blood pressure, and excess male-type hormones, and is the most common cause of congenital adrenal hyperplasia.

\medterm{Acute Pulmonary Histoplasmosis} An acute lower respiratory tract infection caused by inhaling microconidia of the dimorphic fungus \textit{Histoplasma capsulatum}, commonly found in soil enriched with bird or bat guano. Manifestations range from mild flu-like symptoms to severe pneumonia with mediastinal lymphadenopathy.

\synonyms
Primary Pulmonary Histoplasmosis; Acute Histoplasmosis Pneumonia; Cave Disease

\medterm{Adult Hypertension} A chronic cardiovascular condition defined by sustained elevation of systemic arterial blood pressure ($\ge 130/80\text{ mmHg}$ in adults). A major modifiable risk factor for myocardial infarction, heart failure, stroke, and chronic kidney disease.

\synonyms
Essential Hypertension; Primary Hypertension; Systemic Hypertension

\medterm{Biological Half-Life} The time needed for the amount of a medicine in the blood to fall by half. A medicine usually takes several half-lives to reach a stable level after regular dosing or to leave the body after it is stopped.

\synonyms
T1/2

\medterm{Cardiac Electrophysiology Study} An electrophysiology study in which thin tubes are placed through a blood vessel into the heart to record and stimulate its electrical activity. It helps find the cause of some abnormal rhythms or unexplained fainting.

\medterm{Combined Electrolyte Imbalances} Abnormal blood levels of calcium, phosphate, and potassium. This combination does not have one single cause; each result must be confirmed and interpreted with kidney function, medicines, symptoms, and the rest of the electrolyte panel.

\medterm{Congenital Hemolytic Jaundice} Jaundice present from or caused by an inherited condition that increases red-blood-cell destruction, producing excess unconjugated bilirubin; examples include hereditary spherocytosis and some enzymatic defects.

\medterm{Developmental History} A record of a child's developmental milestones and abilities, including motor, language, cognitive, social, and adaptive skills, used to identify possible delays or regression.

\medterm{Dorsocervical Fat Pad} A visible accumulation of fat at the back of the neck and upper shoulders, associated with obesity, corticosteroid excess, some medicines, or lipodystrophy; it differs from spinal kyphosis.

\medterm{Extremity Hardware Removal} Surgical extraction of orthopedic implants such as plates, screws, intramedullary rods, or wires from an arm or leg. Typically indicated for persistent localized pain, implant infection, mechanical irritation, soft-tissue impingement, nonunion, or planned subsequent procedures.

\synonyms
Orthopedic Hardware Removal; Implant Removal; Plate and Screw Removal

\medterm{Geriatric Hospitalization} Hospital care for an older adult, with attention to medical illness as well as mobility, medicines, nutrition, hearing, vision, and cognition. Preventing falls, delirium, pressure injuries, and loss of independence supports safer recovery.

\synonyms
Inpatient Geriatric Care; Acute Care for Elders (ACE)

\medterm{H And E Staining} A standard laboratory stain using hematoxylin and eosin to show cell nuclei, cytoplasm, and tissue structure under a microscope. It is widely used to examine biopsies and identify inflammation, infection, or tumors.

\medterm{H And H} Abbreviation for hemoglobin and hematocrit, two blood measurements commonly used to assess red-blood-cell levels and anemia.

\medterm{H And P} Abbreviation for a patient’s history and physical examination. The history covers symptoms, past health, medicines, and relevant background; the examination checks the body for signs of illness.

\medterm{H Influenzae Meningitis} A serious infection of the coverings of the brain and spinal cord caused by Haemophilus influenzae bacteria. It may cause fever, headache, stiff neck, confusion, seizures, or brain injury, especially in unvaccinated children.

\medterm{H1N1 Flu} Alternate name; see \textit{H1N1 Influenza}.

\medterm{H1N1 Influenza} Influenza A infection caused by an H1N1 virus. It spreads through respiratory droplets and commonly causes sudden fever, cough, sore throat, muscle aches, headache, and tiredness; some people develop pneumonia or severe breathing problems.

\medterm{H1N2} A subtype of influenza A virus that can infect humans and pigs. Human infections are uncommon and usually cause an influenza-like respiratory illness; surveillance and testing help distinguish it from other influenza viruses.

\medterm{H2 Blockers} Medicines that reduce stomach-acid production. They can relieve heartburn, indigestion, and some ulcers; possible adverse effects include headache, diarrhea, constipation, and interactions with other medicines.

\medterm{H2 Receptor Antagonists Overdose} Taking too much of an H2 blocker medicine. It may cause drowsiness, confusion, agitation, abnormal heartbeat, or rarely seizures; urgent poison-control or emergency advice is needed after a significant overdose.

\medterm{H5N1} A subtype of avian influenza A virus that mainly infects birds but can occasionally infect mammals and humans after close contact with infected animals or contaminated environments. Human illness can be severe and requires public-health evaluation.

\medterm{Habenula} A small brain structure helping process reward, unpleasant experiences, stress, motivation, and learned behavior. It connects several brain regions and may influence mood, sleep, pain, and responses to negative outcomes.

\medterm{Habit} A repeated behavior or pattern of action that may be ordinary, helpful, or harmful; in clinical assessment, its frequency, control, and effect on functioning determine whether it is a concern.

\medterm{Habituation} A gradual reduction in response to a repeated, harmless stimulus. For example, a person may stop noticing a constant sound; habituation differs from addiction or drug tolerance.

\medterm{Habitus} A person’s general body build and appearance as observed during a medical examination. Clinicians may describe body habitus when noting nutritional state, muscle mass, or features relevant to a condition.

\medterm{Haem} The iron-containing part of hemoglobin that binds oxygen inside red blood cells. Its structure allows blood to carry oxygen from the lungs to tissues and return carbon dioxide.

\medterm{Haemangioma} A usually harmless growth made of extra blood vessels. It is common in babies and may appear as a red or purple mark or lump.

\medterm{Haemarthrosis} Bleeding into a joint, causing swelling, pain, and reduced movement. It may follow an injury or occur with hemophilia, blood-thinning medicines, surgery, or certain joint disorders.

\medterm{Haematemesis} Vomiting blood or material that looks like coffee grounds, usually from bleeding in the esophagus, stomach, or upper small intestine. It can be life-threatening and requires emergency medical assessment.

\medterm{Haematinics} Nutrients or medicines needed to make healthy red blood cells, especially iron, vitamin B12, and folate. Deficiency can cause anemia, although treatment depends on the specific cause.

\medterm{Haematocele} A collection of blood in a body cavity or tissue, often in the scrotum after injury, surgery, or bleeding. It may cause swelling, pain, or pressure and sometimes needs imaging or drainage.

\medterm{Haematocolpos} Accumulation of menstrual blood in the vagina, usually because the vaginal opening is blocked or too narrow. It can cause absent periods, monthly pelvic pain, and a swollen vagina or uterus.

\medterm{Haematocrit} The proportion of blood volume occupied by red blood cells, expressed as a percentage or decimal fraction. The reference range varies with age, sex, altitude, pregnancy, and laboratory method.

\medterm{Haematogenous} Spread through the bloodstream. An infection, cancer, or other material that travels this way can reach organs or tissues far from its original site.

\medterm{Haematologist} A medical doctor who specializes in the diagnosis, treatment, and prevention of diseases
of the blood and blood-forming organs.

\medterm{Haematoma} A collection of blood outside a blood vessel, usually caused by injury, surgery, bleeding disorders, or blood-thinning medicines. A hematoma may cause swelling and pain and can compress nearby organs or nerves.

\medterm{Haematometra} A collection of menstrual blood in the uterus, usually because the cervix or vagina is blocked or absent. It may cause missed periods, recurring pelvic pain, and uterine enlargement.

\medterm{Haematomyelia} Bleeding into the spinal cord, usually after severe injury or because of an abnormal blood vessel or bleeding disorder. It can cause sudden back pain, weakness, numbness, paralysis, or loss of bladder control.

\medterm{Haematopericardium} Blood within the sac around the heart, usually after chest injury, heart surgery, or a tear in the heart or a blood vessel. It can compress the heart and cause life-threatening tamponade.

\medterm{Haematopoiesis} The process of making blood cells, mainly inside the bone marrow. It produces red cells that carry oxygen, white cells that fight infection, and platelets that help stop bleeding.

\medterm{Haematoporphyria} An older term for a disorder involving porphyrins, substances used to make hemoglobin. Porphyrin disorders may cause sensitivity to sunlight, abdominal or nerve symptoms, or red-brown urine, depending on the type.

\medterm{Haematosalpinx} Blood collected inside a fallopian tube, often from ectopic pregnancy, endometriosis, pelvic injury, or blockage. It may cause pelvic pain and can affect fertility or require urgent treatment.

\medterm{Haematuria} Blood in the urine. It may be visible or found only with a urine test and can have many causes, including infection, stones, medicines, or kidney disease.

\medterm{Haemochromatosis} A disorder in which the body absorbs and stores too much iron. Over time, iron can damage the liver, heart, pancreas, joints, skin, and reproductive organs; early treatment can prevent complications.

\medterm{Haemodialysis} A treatment for severe kidney failure in which a machine filters waste products and extra fluid from the blood. It may be needed temporarily or regularly when the kidneys cannot do enough of this work.

\medterm{Haemoflagellate} A parasite with a whip-like structure that lives in blood or tissues. Examples include the parasites that cause sleeping sickness and Chagas disease.

\medterm{Haemoglobin} The iron-containing protein inside red blood cells that carries oxygen from the lungs to tissues and helps return carbon dioxide to the lungs. A low level is anemia; an abnormal form can cause disease.

\medterm{Haemoglobinopathy} An inherited disorder affecting the amount or structure of hemoglobin, the oxygen-carrying protein in red blood cells. Examples include sickle-cell disease and thalassemia, which can cause anemia or other complications.

\medterm{Haemoglobinuria} Free hemoglobin in the urine, often released when red blood cells break apart in the bloodstream. It can make urine red or brown and may occur with infections, medicines, burns, or blood disorders.

\medterm{Haemoglobinuria} Alternate British spelling; see \textit{Hemoglobinuria}.

\medterm{Haemolysis} The breakdown of red blood cells, which releases hemoglobin into the blood. It can happen normally as old cells are removed or too quickly because of disease.

\medterm{Haemolytic Anaemia} Anemia caused by red blood cells being destroyed faster than the bone marrow can replace them. It may cause tiredness, jaundice, dark urine, rapid heartbeat, breathlessness, or an enlarged spleen.

\medterm{Haemolytic Disease Of The Newborn} Anemia and jaundice in a newborn caused by antibodies from the pregnant person crossing the placenta and destroying the baby’s red blood cells. Severe cases can cause brain injury without prompt treatment.

\medterm{Haemoperitoneum} Bleeding into the abdominal cavity, commonly after injury, a ruptured ectopic pregnancy, or a bleeding liver, spleen, artery, or other organ. It can cause shock and requires urgent assessment.

\medterm{Haemophilia} Alternate spelling; see \textit{Hemophilia}.

\medterm{Haemophilia A} Hemophilia A, an inherited bleeding disorder caused by deficiency or dysfunction of clotting factor VIII. It can cause prolonged bleeding, easy bruising, and bleeding into joints or muscles; treatment replaces or stimulates factor VIII activity.

\medterm{Haemophilus} A group of bacteria that may live harmlessly in the nose and throat or cause infection. Depending on the species, illness can affect the ears, sinuses, lungs, blood, joints, or brain coverings.

\medterm{Haemophilus Aegyptius} A bacterium associated with eye infections and, rarely, severe bloodstream disease. It can spread through respiratory secretions or direct contact, and laboratory identification helps guide treatment and public-health precautions.

\medterm{Haemophilus Haemolyticus} A bacterium that commonly lives harmlessly in the upper airway but can occasionally cause respiratory or invasive infection, especially in people with weakened immunity or serious underlying illness.

\medterm{Haemophilus Infections} Infections caused by Haemophilus bacteria, which may affect the ears, sinuses, lungs, bloodstream, joints, or brain coverings. Severity depends on the species, infection site, age, vaccination, and immune health.

\medterm{Haemophilus Influenzae} A type of bacteria that can cause ear infections, pneumonia, meningitis, and other infections. Type b, or Hib, can cause severe disease in children, but vaccination greatly reduces this risk.

\medterm{Haemophilus Parainfluenzae} A small gram-negative bacterium that normally colonizes the upper respiratory tract but can rarely cause respiratory, bloodstream, or other opportunistic infections.

\medterm{Haemoptysis} Coughing up blood or blood-stained mucus from the airways or lungs. Causes include infection, bronchiectasis, a blood clot, cancer, or irritation; coughing up more than a small streak, or having breathlessness or chest pain, needs urgent medical assessment.

\medterm{Haemorrhage} British spelling of Hemorrhage; see Hemorrhage.

\medterm{Haemorrhagic Fever} A group of severe infections caused by certain viruses that can damage blood vessels and disrupt clotting. Fever, weakness, bleeding, shock, and organ failure may occur; urgent isolation and specialist care may be needed.

\medterm{Haemorrhagic Fever With Renal Syndrome} A hantavirus infection usually acquired by inhaling dust contaminated with infected rodent urine or droppings. It causes fever, headache, bleeding, low blood pressure, and temporary kidney failure and may be severe.

\medterm{Haemorrhagic Stroke} A stroke caused by a burst blood vessel in the brain or in the space around it. The bleeding injures brain tissue and may cause sudden weakness, speech trouble, severe headache, vomiting, or loss of consciousness.

\medterm{Haemorrhoids} Alternate spelling; see \textit{Hemorrhoids}.

\medterm{Haemosiderin} A brownish iron-storage material that forms inside cells when iron accumulates, often after repeated bleeding, transfusions, or excess iron absorption. Large deposits can contribute to tissue damage.

\medterm{Haemosiderosis} Excess storage of iron-containing haemosiderin in tissues, often after repeated transfusions, bleeding into the lungs, or increased iron absorption. Effects depend on the organs involved and may include organ damage.

\medterm{Haemostasis} The body’s normal process for stopping bleeding. Blood vessels narrow, platelets form a temporary plug, and clotting proteins strengthen it with a stable clot.

\medterm{Haemostat} A substance, device, or instrument used to control bleeding. It may work by pressing on a vessel, sealing tissue, promoting clotting, or temporarily blocking blood flow during an operation.

\medterm{Haemothorax} Blood collecting between the lung and chest wall, usually after chest injury, surgery, or a medical procedure. It can compress the lung, cause breathing difficulty or shock, and may require drainage.

\medterm{Haff Disease} Sudden muscle breakdown and severe muscle pain after eating certain contaminated fish or seafood. It can release muscle proteins that damage the kidneys; urgent treatment focuses on fluids and monitoring for complications.

\medterm{Hafnia} A genus of gram-negative intestinal bacteria; Hafnia alvei is the species most often identified in human specimens.

\medterm{Hafnia Alvei} A Gram-negative bacterium found in the intestine and environment that rarely causes opportunistic infection, usually in people with serious illness or weakened immunity.

\medterm{Hageman Factor} An older name for factor XII, a blood protein that helps start one pathway of clotting. Factor XII deficiency can prolong a laboratory clotting test but usually does not cause abnormal bleeding.

\medterm{Hailey-Hailey Disease} An inherited skin disorder that causes repeated blisters, raw areas, and painful cracks, especially in skin folds such as the neck, armpits, and groin. Heat, sweating, friction, and infection can trigger flares.

\medterm{Hair} Keratinized filaments produced by follicles in the skin. Hair protects the scalp, contributes to temperature regulation and sensation, and grows through anagen, catagen, and telogen phases.

\medterm{Hair Bleach Poisoning} Illness after swallowing or inhaling hair-bleaching chemicals, which may contain peroxide or persulfates. They can irritate or burn the mouth, throat, skin, eyes, or lungs; contact poison-control services promptly.

\medterm{Hair Bulb} The rounded base of a hair follicle where new hair is produced. It contains growing cells and a small blood-supplied area that provides nutrients and helps determine hair growth and color.

\medterm{Hair Cells} Special cells in the inner ear that change sound vibrations and head movement into nerve signals for hearing and balance.

\medterm{Hair Disease} A disorder affecting hair growth, the hair shaft, or the scalp; causes include genetics, inflammation, infection, medicines, and nutritional problems.

\medterm{Hair Dye Poisoning} Harm caused by swallowing or inhaling hair-dye chemicals. Depending on the ingredients and amount, symptoms may include vomiting, swelling, breathing difficulty, allergic reaction, kidney injury, or muscle damage and require urgent advice.

\medterm{Hair Follicle} A small tube in the skin that produces a hair. Each follicle has a growth phase, a brief transition phase, and a resting phase.

\medterm{Hair Loss} Loss or thinning of hair on the scalp or body. It may be temporary or permanent, patchy or widespread, and caused by genetics, illness, medicines, autoimmune disease, infection, or scarring of the scalp.

\synonyms
Alopecia; Effluvium; Baldness

\medterm{Hair Loss} The partial or complete loss of hair from areas where it normally grows, most commonly
the scalp. Causes include genetic predisposition (androgenetic alopecia), autoimmune
conditions (alopecia areata), hormonal changes, medications, stress, and nutritional
deficiencies.

\synonyms
Alopecia; Baldness; Effluvium

\medterm{Hair Loss Classification} Hair loss is classified as scarring or nonscarring. Scarring hair loss destroys follicles and is usually permanent; nonscarring loss leaves follicles intact and may improve when its cause, such as illness, medicines, stress, or autoimmune disease, is treated.

\medterm{Hair Removal} Removal of unwanted hair by shaving, trimming, waxing, depilatory products, electrolysis, or laser treatment. Possible problems include irritation, burns, ingrown hairs, pigment changes, and infection.

\medterm{Hair Shaft} The part of a hair that extends above the skin. It is made of dead, hardened cells produced by the follicle and can be affected by cutting, heat, chemicals, friction, or damage to the hair root.

\medterm{Hair Spray Poisoning} Illness after swallowing or inhaling hair spray. Small accidental exposures may irritate the mouth or lungs, while larger exposures can cause vomiting, drowsiness, breathing problems, or chemical injury; contact poison-control services.

\medterm{Hair Straightener Poisoning} Harm after swallowing a hair-straightening product, which may contain corrosive chemicals. It can burn the mouth, throat, esophagus, or stomach and cause breathing or swallowing problems; do not induce vomiting and seek emergency advice.

\medterm{Hair Tonic Poisoning} Illness after swallowing or inhaling a hair-tonic product. Ingredients such as alcohols or solvents may cause vomiting, drowsiness, low blood sugar, or breathing problems; the product and amount determine emergency treatment.

\medterm{Hair Transplant} A procedure that moves hair follicles, usually from a donor area of the scalp, to areas affected by permanent hair loss; results depend on diagnosis, donor supply, technique, and postoperative care.

\medterm{Hair-Ball} A hair-ball, or trichobezoar, is a mass of swallowed hair in the stomach or intestine, usually associated with repeated hair eating.

\medterm{Hairy Cell Leukaemia} Hairy cell leukaemia is a rare slow-growing blood cancer in which abnormal B lymphocytes accumulate in the bone marrow, spleen, and blood.

\medterm{Hairy Cell Leukemia} A rare, usually slow-growing cancer of a type of white blood cell. Abnormal cells build up in the bone marrow, blood, and spleen, causing low blood counts, infections, tiredness, or an enlarged spleen.

\synonyms
HCL; Leukemic Reticuloendotheliosis

\medterm{Hairy Tongue} A usually harmless overgrowth and discoloration of the tongue's papillae, often appearing black, brown,
or white. It may be associated with smoking, poor oral hygiene, dry mouth, or antibiotics;
gentle tongue cleaning and removal of contributing factors usually help.

\medterm{Hajdu-Cheney Syndrome} A rare inherited bone disorder causing fragile bones, short stature, abnormal skull and facial development, and sometimes loose joints or dental problems. Severity varies, and specialist monitoring helps prevent complications.

\medterm{Halcinonide} A potent topical corticosteroid used to reduce inflammation and itching in steroid-responsive dermatitis and other inflammatory skin disorders; prolonged or extensive use can cause skin atrophy and systemic effects.

\medterm{Haldane Effect} The effect of oxygen binding to hemoglobin in the lungs, which helps release carbon dioxide for exhalation. In body tissues, hemoglobin without oxygen carries more carbon dioxide back toward the lungs.

\medterm{Half-Life} The time it takes for the amount of a medicine in the body, or its blood level, to decrease by half. It helps determine how often a medicine should be taken and how long it remains in the body.

\medterm{Halitosis} Persistent unpleasant breath odor, most often from bacteria on the tongue, teeth, or gums. Dry mouth, smoking, foods, sinus or throat disease, reflux, and some illnesses or medicines can also contribute.

\medterm{Hallermann's Syndrome} A rare congenital disorder characterized by distinctive craniofacial development, sparse hair, dental abnormalities, and sometimes eye or skeletal findings; it is also called Hallermann--Streiff syndrome.

\medterm{Hallervorden-Spatz Disease} Alternate historical term; see \textit{Neurodegeneration With Brain Iron Accumulation}.

\medterm{Hallucination} A sensory experience that feels real but occurs without a matching external stimulus. It may involve seeing, hearing, feeling, smelling, or tasting something that is not present.

\medterm{Hallucinations} Hallucinations are perceptions without an external sensory stimulus and can arise from psychiatric illness, neurologic disease, delirium, substances, medications, sensory loss, or sleep transitions.

\medterm{Hallucinogen} A psychoactive substance that can alter perception, thought, and mood and may produce hallucinations; effects and toxicity vary by substance, dose, co-ingestants, and mental state.

\medterm{Hallux} The great toe, or big toe. It has two bones and helps the foot stay stable and push off during walking. Common problems include hallux valgus (a bunion), hallux rigidus (a stiff, arthritic joint), turf toe (a sprain), and gout causing sudden pain and swelling at the joint where the toe meets the foot.

\medterm{Hallux Rigidus} Arthritis of the joint where the big toe meets the foot, causing pain, stiffness, swelling, and reduced upward movement. A bony bump may develop and make walking or wearing shoes difficult.

\medterm{Hallux Valgus} A deformity in which the big toe angles toward the other toes and the joint at its base becomes prominent on the inner side of the foot. It may cause pain, rubbing, swelling, and difficulty wearing shoes.

\medterm{Hallux Valgus} A bunion: a change in the alignment of the joint at the base of the big toe, with the toe pointing inward and a prominent area on the inner forefoot. It may become painful from pressure or inflammation.

\synonyms
Bunion; Hallux Abductovalgus

\medterm{Hallux Varus} Abnormal medial deviation of the great toe, often following surgery, trauma, neuromuscular imbalance, or congenital deformity; it can cause shoe difficulty and pressure-related pain.

\medterm{Halo Nevus} A usually harmless mole surrounded by a ring of lighter skin. The immune system causes the mole and nearby pigment cells to lose color; it often occurs in children or young adults and may slowly fade.

\medterm{Halogen} An element in group 17 of the periodic table, including fluorine, chlorine, bromine, iodine, and astatine; halogens and their compounds have medical, biologic, and toxicologic importance.

\medterm{Haloperidol} Haloperidol is a first-generation antipsychotic that blocks dopamine D2 receptors and is used for psychosis, agitation, delirium, and tics; extrapyramidal reactions, tardive dyskinesia, hyperprolactinaemia, and QT prolongation are important risks.

\medterm{Halothane} A volatile inhaled anesthetic formerly used to induce or maintain general anesthesia; it can cause hypotension, arrhythmias, and rare severe immune-mediated liver injury.

\medterm{Hamamelis} A botanical preparation from witch hazel used topically as an astringent for minor skin or mucosal irritation; evidence and product composition vary.

\medterm{Hamartoma} A usually harmless growth made of an abnormal arrangement of tissue types normally found in that part of the body. It may occur in organs such as the lung, skin, brain, or intestine and may cause problems if it grows or presses on nearby structures.

\medterm{Hamartomatous Polyposis} A condition in which multiple hamartomatous gastrointestinal polyps develop from disorganized but benign tissue overgrowth, sometimes as part of an inherited cancer-predisposition syndrome.

\medterm{Hamate} A small, wedge-shaped bone on the little-finger side of the wrist. Its hook provides attachment for ligaments and tendons; fractures there may cause wrist pain or pressure on a nearby nerve.

\medterm{Hamate Bone} The hamate is a wedge-shaped wrist bone on the little-finger side that has a hook serving as an attachment and passage point for tendons and nerves.

\medterm{Hamman Sign} A crunching or clicking noise heard with a stethoscope over the chest in time with the heartbeat, caused by air trapped in the center of the chest.

\medterm{Hamman's Sign} A crackling or crunching sound heard over the chest in time with the heartbeat, usually caused by air in the central chest around the heart and major vessels. It can occur with pneumomediastinum and needs clinical assessment.

\medterm{Hammer} The common name for the malleus, one of the three small bones of the middle ear.

\medterm{Hammer Toe} A deformity in which a lesser toe bends downward at its middle joint and may press against the shoe. It can cause pain, stiffness, corns, or calluses and may be flexible at first or become fixed.

\medterm{Hammer Toe Repair} Surgery to correct a fixed or painful flexion deformity of a lesser toe, using soft-tissue release, tendon balancing, bone correction, or temporary fixation as indicated.

\medterm{Hammertoe} A deformity in which a toe bends downward at its middle joint, giving it a claw-like
appearance. It may cause pain, stiffness, corns, or calluses, especially when wearing shoes.

\medterm{Hampton Hump} Wedge-shaped opacification with its base abutting the pleura on chest radiography,
representing pulmonary infarction secondary to Pulmonary Embolism. See Pulmonary
Embolism (PE).

\medterm{Hampton's Hump} Alternate spelling; see \textit{Hampton Hump}.

\medterm{Hamstring} One of three muscles forming the posterior compartment of the thigh: semitendinosus,
semimembranosus, and biceps femoris (long and short heads). Common origin: ischial
tuberosity (except short head of biceps femoris which originates from the linea aspera
of the femur).

\medterm{Hamstring Injury} A strain or tear of one or more muscles or tendons at the back of the thigh.

\medterm{Hamstring Muscles} The muscles at the back of the thigh. They bend the knee and help move the hip, especially during walking, running, and jumping. A sudden stretch or forceful contraction can strain or tear them.

\medterm{Hand} The part of the upper limb beyond the wrist, containing the palm, fingers, bones, joints, muscles, nerves, and blood vessels used for grasping and sensation.

\medterm{Hand Deformity} An abnormal shape, alignment, or loss of motion of the hand or fingers caused by congenital difference, trauma, arthritis, nerve dysfunction, tendon disease, or other conditions.

\medterm{Hand Foot And Mouth Disease} A contagious viral illness, most common in young children, that causes fever, painful mouth sores, and a rash or small blisters on the hands, feet, and sometimes buttocks. It spreads through close contact and contaminated secretions or stool.

\medterm{Hand, Foot, And Mouth Disease} A contagious enterovirus infection, usually in children, causing fever, painful mouth sores, and small spots or blisters on the hands, feet, and sometimes buttocks. It usually improves on its own, but dehydration can occur.

\medterm{Hand Hygiene} Cleaning the hands with soap and water or an alcohol-based hand rub to remove or kill germs. It is especially important before and after patient contact, after touching body fluids, and when hands are visibly dirty.

\medterm{Hand Lotion Poisoning} Illness after swallowing hand lotion. Most products cause mild stomach upset, but some contain alcohols or other chemicals that can cause drowsiness, low blood sugar, or breathing problems; contact poison-control services.

\medterm{Hand Microsurgery} Specialized hand surgery performed with magnification and fine instruments to repair or reconstruct small blood vessels, nerves, and other delicate tissues. It may be used for complex injuries, replantation, and restoration of blood flow to a threatened digit or hand.

\medterm{Hand of Benediction} An abnormal hand posture observed when a patient attempts to make a fist, resulting in extension of the index and middle fingers with flexion of the fourth and fifth fingers. It occurs due to proximal median nerve injury, which paralyses the flexor digitorum superficialis and the lateral half of the flexor digitorum profundus.

\synonyms
Benediction Sign; Preacher's Hand

\medterm{Hand Or Foot Spasms} Sudden, involuntary tightening of muscles in a hand or foot. Causes include overuse, dehydration, abnormal calcium or other salts, nerve disease, medicines, and reduced blood flow; repeated or painful spasms need assessment.

\medterm{Hand Pain} Pain or discomfort in the hand, fingers, or wrist caused by injury, arthritis, nerve compression, tendon problems, infection, or other conditions. Sudden severe pain, deformity, numbness, or loss of circulation needs prompt assessment.

\medterm{Hand Prolapse} A fetal hand or arm lying beside or below the part leading into the birth canal. It can interfere with labor or reduce blood flow to the limb; the position requires prompt obstetric assessment and should not be pushed back at home.

\medterm{Hand X-Ray} Radiographic imaging of the hand used to evaluate fractures, dislocations, arthritis, infection, tumors, bone age, and other abnormalities of the bones and joints.

\medterm{Hand-Arm Vibration Syndrome} An occupational disorder from prolonged exposure to vibrating tools, causing numbness, tingling, reduced grip, and episodic whitening of the fingers from vascular spasm.

\medterm{Hand-Foot Syndrome} A dose-limiting toxic cutaneous reaction triggered by certain systemic antineoplastic and targeted chemotherapy agents (notably capecitabine, 5-fluorouracil, pegylated liposomal doxorubicin, cytarabine, and multi-kinase inhibitors like sorafenib and sunitinib). It is characterized by progressive bilateral burning dysesthesia, tingling, painful symmetric erythema, swelling, and subsequent blistering or desquamation primarily localized to the friction-bearing pressure areas of the palms and soles.
\synonyms{Palmar-Plantar Erythrodysesthesia; PPE; Chemotherapy-Induced Acral Erythema; Burgdorf Reaction}

\medterm{Hand-Foot-And-Mouth Disease} A contagious viral illness causing sores in the mouth and a rash on the hands and feet, usually in children.

\medterm{Hand-Foot-And-Mouth Syndrome} A usually mild, contagious viral illness, commonly caused by enteroviruses, producing fever, mouth sores,
and small spots or blisters on the hands and feet. It mainly affects young children but can occur in adults. Hydration,
pain relief, and hygiene are important; severe dehydration, neurologic symptoms, or persistent fever requires care.

\medterm{Hand-Foot-Mouth Disease} A common enteroviral infection, usually in young children, causing fever, painful oral
ulcers, and a vesicular rash on the hands, feet, and sometimes buttocks. It is highly
contagious and usually self-limited.

\medterm{Handling Sharps And Needles} Handling sharps and needles means using, passing, disposing of, and storing sharp medical instruments safely to prevent needlestick injuries and infection transmission.

\medterm{Hand-SchüLler-Christian Disease} Hand-Schüller-Christian disease is an older name for a multisystem form of Langerhans cell histiocytosis, often involving bone, skin, and the pituitary gland.

\medterm{Handwashing} Handwashing is cleaning the hands with soap and water or an alcohol-based product to remove germs and reduce transmission of infection.

\medterm{Hanging And Strangulation Signs} Injuries or other findings caused by pressure around the neck that restricts breathing or blood flow. Findings may include neck marks, swelling, voice changes, trouble swallowing, breathing difficulty, confusion, or loss of consciousness; urgent medical assessment is required.

\medterm{Hang-Nail} A small torn flap of skin beside a fingernail or toenail, often caused by dryness or trauma and vulnerable to local infection if torn or contaminated.

\medterm{Hangover} A group of symptoms after alcohol use, such as headache, nausea, thirst, fatigue, dizziness, and sensitivity to light; symptoms reflect effects of alcohol and its metabolites, dehydration, sleep disruption, and inflammation.

\medterm{Hansen Disease} A long-lasting infection caused by Mycobacterium leprae that mainly affects the skin and peripheral nerves. Pale or reddish numb patches, weakness, and eye problems may occur; multidrug antibiotics cure infection and prevent disability.

\synonyms
Leprosy; Hanseniasis; Mycobacterium Leprae Infection

\medterm{Hantavirus} A group of viruses spread mainly by breathing dust contaminated with infected rodent urine, droppings, or nesting material. They can cause severe lung disease or a fever that damages the kidneys.

\medterm{Hantavirus Infection} A zoonotic viral infection acquired mainly through inhalation of contaminated rodent excreta. It may cause hemorrhagic fever with renal syndrome or hantavirus pulmonary syndrome, which can progress rapidly to respiratory failure.

\medterm{Hantavirus Pulmonary Syndrome} A severe lung infection acquired mainly by inhaling dust contaminated with infected rodent urine, droppings, or nesting material. Fever, muscle aches, headache, and stomach symptoms can rapidly progress to fluid-filled lungs and respiratory failure.

\synonyms
HPS; Sin Nombre Virus Infection

\medterm{Haploid} Having one complete set of chromosomes, as in a human sperm or egg. Haploid cells contain half the chromosome number of most body cells and restore the full number when they join during fertilization.

\medterm{Haploidentical Donor} A donor who shares one human leukocyte antigen haplotype with the recipient, commonly used for selected hematopoietic stem-cell or organ transplantation when a fully matched donor is unavailable.

\medterm{Haploidy} The state of having one complete set of chromosomes, as in human gametes; haploid cells contain 23 chromosomes.

\medterm{Haploinsufficiency} A genetic state in which one functional copy of a gene does not produce enough gene product for normal function, so loss or inactivation of the other copy can cause disease.

\medterm{Haplotype} A group of nearby gene versions inherited together on one chromosome. Haplotype information can help study inherited disease risk, identify tissue donors, and trace how genetic traits are passed through families.

\medterm{Hapten} A hapten is a small molecule that is not immunogenic by itself but can trigger an immune response after binding to a larger carrier protein, as in some drug allergies.

\medterm{Haptoglobin} A blood protein made mainly by the liver that binds free hemoglobin released when red blood cells break down. Low levels can support evaluation of hemolysis, but liver disease and inflammation also affect the result.

\medterm{Haptoglobin Analyses} Laboratory tests measuring haptoglobin, a plasma protein that binds free hemoglobin; low values can support evaluation of intravascular hemolysis but are also affected by liver disease and inflammation.

\medterm{Haptoglobin Blood Test} A blood test measuring haptoglobin, a protein that binds free hemoglobin; low levels may support intravascular hemolysis but can also reflect liver disease or inflammation.

\medterm{Haptoglobin Decreased} A low blood haptoglobin result, commonly caused by intravascular red-cell destruction, reduced hepatic production, or severe liver disease; it is interpreted with hemoglobin, bilirubin, reticulocytes, and a blood smear.

\medterm{Hard Palate} The bony, mucosa-covered front portion of the roof of the mouth that separates the oral and nasal cavities and provides a surface for chewing and speech.

\medterm{Harlequin Color Change} A temporary, harmless change in a newborn’s skin color in which one side becomes red while the other remains pale. It often appears when the baby lies on one side and usually fades within minutes.

\medterm{Harlequin Foetus} An older term for a baby born with harlequin ichthyosis, a severe inherited skin disorder causing thick, hard plates of skin separated by deep cracks. The condition can impair temperature control, movement, feeding, and protection from infection.

\medterm{Harlequin Ichthyosis} A rare severe inherited skin disorder in which newborns have thick plates of skin separated by deep cracks. It can impair temperature control, movement, feeding, breathing, and protection from infection and requires intensive neonatal care.

\medterm{Harm Reduction} Evidence-based strategies that decrease health, social, and legal harms associated with
substance use without requiring immediate abstinence. Examples include naloxone access,
sterile injecting equipment, drug-checking where available, safer-use education, and
linkage to treatment.

\medterm{Harm Reduction Therapy} Care that reduces illness, injury, infection, overdose, and other harms related to substance use or another risky behavior, even when stopping immediately is not possible. It can include safer-use support and treatment for dependence.

\synonyms
Harm Reduction; Risk Mitigation Strategy

\medterm{Harmless Murmur} An innocent or physiologic heart murmur produced by normal blood flow without structural heart disease; it is diagnosed only after appropriate clinical assessment.

\medterm{Harrison's Sulcus} Harrison's sulcus is a horizontal groove along the lower chest wall caused by inward pulling of the ribs, classically in long-standing childhood breathing difficulty.

\medterm{Hartmann Procedure} An operation that removes the diseased lower colon, brings the remaining colon through the abdominal wall as a colostomy, and closes the downstream bowel. It is often used when infection, blockage, perforation, or severe illness makes immediate reconnection unsafe.

\medterm{Hartnup Disease} A rare inherited problem absorbing and reusing certain amino acids. It may cause a sun-sensitive rash, poor coordination, or mood and thinking changes, although some people have no symptoms.

\medterm{Hartnup Disorder} An inherited disorder of neutral amino-acid transport that can cause pellagra-like skin changes, neurologic symptoms, or no symptoms.

\medterm{HAS-BLED Score} A score estimating bleeding risk in a person with atrial fibrillation who may need a blood thinner. It considers high blood pressure, kidney or liver disease, previous stroke or bleeding, unstable clotting results, age, medicines, and alcohol; a high score prompts closer monitoring and correction of changeable risks.

\synonyms
HAS-BLED; Anticoagulant Bleeding Risk Index

\medterm{Hashimoto Disease} An autoimmune disorder in which antibodies and lymphocytes progressively damage the thyroid, commonly causing hypothyroidism, goiter, fatigue, cold intolerance, constipation, and weight change.

\medterm{Hashimoto Thyroiditis} An autoimmune disorder in which immune-mediated inflammation gradually damages the thyroid gland and may cause
hypothyroidism. The gland may be enlarged or atrophic, and thyroid autoantibodies are
often present.

\medterm{Hashimoto's Disease} An autoimmune disorder in which the immune system gradually damages the thyroid, often causing low thyroid-hormone levels. It may produce tiredness, cold intolerance, weight change, constipation, or an enlarged thyroid.
\synonyms
Hashimoto Thyroiditis; Chronic Autoimmune Thyroiditis

\medterm{Hashimoto's Thyroiditis} A chronic autoimmune inflammation of the thyroid that commonly causes progressive hypothyroidism and may produce a firm, enlarged thyroid gland.

\medterm{Hashish} A concentrated form of cannabis made from resin of the cannabis plant. It contains cannabinoids that can impair memory, coordination, and judgment and may cause anxiety, dependence, or psychosis in susceptible people.

\medterm{Hashitoxicosis} A temporary period of excess thyroid hormone caused by release of stored hormone during Hashimoto thyroiditis.

\medterm{Hassall's Corpuscles} Hassall's corpuscles are rounded structures in the thymus made of specialized epithelial cells and involved in the development of immune tolerance.

\medterm{Haustra} The series of small pouches along the outside of the large intestine. They form as the bowel muscles contract and help the colon mix and move its contents.

\medterm{Haverhill Fever} An illness caused by Streptobacillus bacteria, usually after drinking contaminated milk or water. It may cause fever, vomiting, headache, muscle or joint pain, and a rash and requires antibiotic treatment.

\medterm{Hay Fever} A common name for seasonal allergic rhinitis, an allergic response to airborne outdoor plant pollens and mold spores that causes sneezing, a blocked or runny nose, an itchy palate, and watery, irritated eyes. See also \textit{Seasonal Allergies}.
\synonyms
Seasonal Allergies; Seasonal Allergic Rhinitis; Pollinosis

\medterm{Hazardous Materials} Substances that can cause injury, poisoning, fire, radiation exposure, or environmental harm; safe handling requires labeling, protective equipment, exposure control, and emergency procedures.

\medterm{Hazardous Substance} A chemical, biological agent, or other material that can harm health through swallowing, inhalation, skin contact, injection, radiation, or environmental exposure. Risk depends on its properties, amount, route, and duration of exposure.

\medterm{Hb} Hemoglobin, the iron-containing protein in red blood cells that carries oxygen and contributes to carbon-dioxide transport and blood's acid--base buffering.

\medterm{HbA1c} Abbreviation; see \textit{Glycated Hemoglobin}.
\synonyms
Glycated Hemoglobin; Hemoglobin A1c; A1c Test

\medterm{HCG} Human chorionic gonadotropin is a hormone produced mainly during pregnancy; its measurement is used in pregnancy testing and evaluation of selected tumors.

\medterm{HCG In Urine} A urine test detecting human chorionic gonadotropin, commonly used as a pregnancy test; accuracy depends on timing, urine concentration, and test performance.

\medterm{Hct} The percentage of blood volume occupied by red blood cells, measured as part of a complete blood count.

\medterm{HCV PCR} A polymerase chain reaction test that detects or measures hepatitis C virus RNA in blood, helping identify active infection and monitor treatment response.

\medterm{HCV-Associated Rheumatic Disease} Joint, muscle, blood-vessel, or autoimmune problems associated with hepatitis C virus infection. Possible manifestations include joint pain or inflammation, mixed cryoglobulinemia, vasculitis, fatigue, or muscle symptoms.

\medterm{HDL} High-density lipoprotein, a class of particles that transports cholesterol and other lipids in the blood. HDL
concentration is one component of cardiovascular risk assessment; it should not be interpreted
as a protective or harmful measure by itself.

\medterm{HDL Cholesterol} Cholesterol carried in high-density lipoprotein particles; HDL participates in transporting cholesterol from tissues toward the liver, but its concentration alone does not determine cardiovascular risk.

\medterm{HDL Test} A blood test measuring high-density lipoprotein cholesterol, which is used with the rest of the lipid profile to assess cardiovascular risk.

\medterm{HDR Syndrome} A genetic disorder characterized by hypoparathyroidism, sensorineural deafness, and renal dysplasia or other kidney abnormalities, usually caused by a pathogenic variant in \textit{GATA3}.

\medterm{Head} The body region containing the brain, eyes, ears, nose, mouth, skull, and face, connected to the trunk by the neck and supplied by major cranial nerves and vessels.

\medterm{Head And Face Reconstruction} Reconstructive surgery to restore the form or function of the skull, face, or associated soft tissues after trauma, cancer treatment, congenital difference, or other disease.

\medterm{Head And Neck Cancer} Cancers arising in the oral cavity, pharynx, larynx, sinonasal tract, and salivary
glands, most commonly squamous cell carcinoma. Risk factors include tobacco, alcohol,
and HPV (oropharyngeal). Treatment involves surgery, radiation, and chemotherapy, often
combined to preserve function.

\medterm{Head And Neck Oncologic Surgery} Surgical treatment of cancers of the oral cavity, throat, larynx, nose, sinuses, salivary glands, or neck lymph nodes. Procedures may remove the tumor, treat involved lymph nodes, and reconstruct affected tissues while preserving breathing, swallowing, speech, and appearance when possible.

\synonyms
Head and Neck Surgical Oncology; Maxillofacial Oncologic Surgery

\medterm{Head Circumference} The measurement around the largest part of a baby’s or child’s head, taken above the eyebrows and ears. It is plotted over time to assess brain growth; unusually small, large, or rapidly changing measurements may need evaluation.

\medterm{Head CT Scan} Computed tomography of the head using x-rays and computer reconstruction to evaluate the brain, skull, bleeding, stroke, tumors, trauma, and other abnormalities.

\medterm{Head Injury} Damage to the scalp, skull, or brain caused by a blow, fall, collision, or penetrating object. It may cause a wound, skull fracture, concussion, bleeding, or brain damage. Worsening headache, repeated vomiting, seizure, confusion, weakness, or loss of consciousness requires urgent assessment.

\synonyms
Traumatic Brain Injury (TBI); Craniocerebral Trauma

\medterm{Head Injury First Aid} Immediate emergency first-aid protocols following head trauma, including stabilizing the cervical spine, maintaining an open airway, controlling external cranial bleeding with gentle pressure, and continuously monitoring level of consciousness and vital signs while urgently seeking medical evaluation for warning signs such as loss of consciousness, repeated vomiting, or focal neurologic deficits.

\synonyms
Cranial Trauma First Aid; Head Trauma Immediate Care; Concussion First Aid

\medterm{Head Lice} Infestation of the scalp by the parasitic insect \textit{Pediculus humanus capitis}. It causes itching and visible nits or lice and spreads mainly through close head-to-head contact; treatment requires an effective product and careful follow-up.

\medterm{Head Movement} Motion of the head produced mainly by the cervical spine and neck muscles, including flexion, extension, rotation, and side bending; pain or restricted movement may indicate injury or disease.

\medterm{Head MRI} Magnetic resonance imaging of the head using a magnetic field and radio waves to produce detailed images of the brain, cranial nerves, blood vessels, and surrounding structures.

\medterm{Head Trauma} An injury to the scalp, skull, brain, or other tissues of the head. Concussion, bleeding, skull fracture, and brain injury can result; worsening headache, repeated vomiting, seizure, confusion, weakness, or loss of consciousness requires emergency care.

\medterm{Headache} Pain or discomfort in the head, scalp, or upper neck. Primary headaches include migraine and tension headache; secondary headaches result from another problem such as infection, injury, bleeding, very high blood pressure, or a brain disorder.

\medterm{Headache and Migraine} Alternate name; see \textit{Headache}.

\medterm{Headache Red Flags} Critical clinical warning signs in a patient presenting with headache that indicate potentially life-threatening intracranial pathology (such as subarachnoid hemorrhage, meningitis, intracranial tumor, or temporal arteritis). Warning signs include sudden thunderclap onset, fever with meningismus, new focal neurological deficits, onset after age 50, papilledema, or postural worsening.

\synonyms
Headache Warning Signs; Red Flag Headaches; Danger Signs in Headache

\medterm{Head-Lice} Head lice are small insects that live on the scalp and cause itching; they spread mainly through close head-to-head contact.

\medterm{Healing} The body’s process of repairing damaged tissue. It usually involves stopping bleeding, reducing inflammation, forming new tissue, and remodeling a scar; healing may be delayed by infection, poor blood flow, diabetes, or inadequate nutrition.

\medterm{Health Anxiety Disorder} Excessive worry about having or developing a serious illness, often despite reassurance or limited medical findings. It can cause repeated checking or health-care visits and is treated with supportive care and evidence-based psychological therapy.

\medterm{Health Behavior} An action or pattern of action that affects health, such as eating, physical activity, tobacco use, alcohol use, sexual practices, or taking medicines.

\medterm{Health Behavior Change} Modification of behaviors that affect health, such as tobacco use, diet, physical
activity, and adherence to treatment.

\medterm{Health Behavior Motivation} Factors that influence initiation and maintenance of health-related behaviors, including
adherence, physical activity, and smoking cessation.

\medterm{Health Care} The prevention, diagnosis, treatment, and ongoing support provided to maintain or improve health, delivered by professionals, institutions, public-health systems, and informal caregivers.

\medterm{Health Care Agents} People legally chosen to make healthcare decisions for someone who cannot make or communicate them. Their authority and responsibilities depend on local law and the document that appoints them.

\medterm{Health Care Facility} A place where health services are delivered, such as a clinic, hospital, laboratory, nursing facility, or rehabilitation center; its services and level of care determine appropriate referrals.

\medterm{Health Care Personnel} Professionals and support workers involved in patient care, including clinicians, nurses, therapists, technicians, aides, and public-health staff.

\medterm{Health Care Proxy} A person designated to make healthcare decisions for someone who cannot make or communicate those decisions. The appointment may be recorded in an advance directive or other legal document.

\medterm{Health Checkup} A preventive visit that reviews symptoms, medical history, medicines, lifestyle, vaccinations, examination findings, and age- or risk-appropriate screening. The content should be individualized rather than based on unnecessary tests.

\medterm{Health Disparities} Differences in health outcomes or access to care among groups, often associated with social,
economic, environmental, geographic, or structural factors.

\medterm{Health Disparity} A difference in health outcomes or access to healthcare between population groups. It may be linked to income, housing, education, discrimination, environment, language, disability, geography, or unequal treatment.

\medterm{Health Education} The process of giving people understandable information and skills to prevent illness, recognize symptoms, use treatment safely, and make informed health decisions.

\medterm{Health Equity} The fair opportunity for all people to attain their highest possible level of health, without avoidable differences caused by social or economic disadvantage.

\medterm{Health For All} A public-health goal emphasizing that all people should have access to needed health services and conditions that support health, without financial hardship.

\medterm{Health Inequity} A systematic, avoidable, and unjust difference in health or access to health resources
between social groups. It differs from health variation that is not preventable or not
linked to unfair social conditions.

\medterm{Health Information Technology} The use of computers and digital systems to collect, store, exchange, protect, and use health information. Examples include electronic records, telemedicine, patient portals, decision-support tools, and remote monitoring; safe systems protect privacy and support accurate care.

\medterm{Health Literacy} The ability to find, understand, judge, and use health information and services to make appropriate decisions. It includes communicating with clinicians, following instructions, reading medicine labels, and recognizing when to seek help; clear communication by healthcare staff also improves health literacy.

\medterm{Health Maintenance Organization} A health-insurance arrangement that usually uses a network of clinicians and requires a primary-care clinician to coordinate specialist care. Coverage rules, referrals, costs, and services vary by plan.

\medterm{Health Promotion} Actions that help people improve or protect their health and prevent illness. They may include vaccination, healthy eating, physical activity, sleep, screening, safer environments, social connection, and avoiding tobacco. Advice should reflect a person’s age, abilities, culture, risks, and preferences.

\medterm{Health Risk} The likelihood that a behavior, exposure, condition, or other factor will cause harm to health, considered together with the severity and timing of possible harm.

\medterm{Health Risks Of Alcohol Use} Alcohol use can increase the risk of liver disease, several cancers, high blood pressure, injuries, dependence, and problems affecting mental health and relationships.

\medterm{Health Surveillance} The ongoing collection and analysis of health data to identify patterns, monitor disease, and guide public-health action.

\medterm{Health Survey} A structured set of questions used to collect information about health, symptoms, behaviors, healthcare use, or population needs.

\medterm{Health Visitor} A community health professional, especially in the United Kingdom, who supports families with infants and young children. They provide health advice, developmental monitoring, prevention, and referral to other services.

\medterm{Health Workforce} The people and organizations that provide, support, manage, or study health services, including clinicians, public-health workers, and support staff.

\medterm{Healthcare Associated Infection} An infection acquired while receiving healthcare or as a result of a medical procedure. It may occur in hospitals, clinics, nursing facilities, or at home and can involve a wound, urinary tract, lungs, bloodstream, or other site.

\medterm{Healthcare Provider} A healthcare professional or organization that provides medical care, such as a
physician, nurse, therapist, or clinic.

\medterm{Healthcare Proxy} A person designated to make healthcare decisions for someone who cannot make or communicate those decisions. The appointment may
be recorded in an advance directive or other legal document.

\medterm{Health-Related QoL} A person’s physical, emotional, social, and daily functioning as affected by health and illness. It reflects the patient’s experience and is assessed separately from laboratory results or survival.

\medterm{Healthy Aging} Developing and maintaining the physical and mental abilities that support well-being and independence in later life. It is influenced by health conditions, activity, nutrition, relationships, environment, and access to care.

\medterm{Healthy Years} Years lived without substantial disease, disability, or loss of function. This population measure helps describe quality of life and health needs, but it does not predict an individual’s exact lifespan.

\medterm{Healthy Years Lost} The years lost because of early death or time lived with illness or disability. This population measure estimates the overall effect of disease; it does not predict exactly how long one person will live.

\medterm{Hearing} The ability to detect and understand sound. Sound vibrations are converted into nerve signals in the inner ear and interpreted by the brain. Hearing loss may result from the outer or middle ear, inner ear, hearing nerve, or brain pathways.

\medterm{Hearing Aid} A small electronic device worn in or behind the ear that makes sounds louder or clearer for someone with hearing loss. It contains a microphone, processor, and speaker and is fitted and adjusted according to hearing-test results.

\medterm{Hearing Aids And Devices} Electronic devices worn in or behind the ear that amplify sound for individuals with
hearing loss. Modern hearing aids include digital signal processing, directional
microphones, and connectivity to smartphones and other audio sources.

\medterm{Hearing Disorder} A condition that reduces or distorts hearing through disease of the outer or middle ear, cochlea, auditory nerve, or central pathways; audiometry and examination help distinguish conductive from sensorineural loss.

\medterm{Hearing Impairment} Partial or complete reduction in the ability to detect or understand sound. It may be
conductive, sensorineural, or mixed and can result from congenital disorders, infection,
noise, ototoxic injury, trauma, ageing, or disease of the ear or auditory pathways.

\medterm{Hearing Loss} Reduced ability to hear. Conductive loss results from a problem in the outer or middle ear; sensorineural loss results from damage to the inner ear or hearing nerve; mixed loss has both types. Causes include aging, noise, infection, injury, and some medicines.

\synonyms
Hearing Impairment; Hypoacusis

\medterm{Hearing Loss And Music} Difficulty hearing or distinguishing musical sounds because of reduced hearing, damage to inner-ear cells, or problems processing sound. Music may seem quieter, distorted, or lacking in certain pitches; hearing protection and individualized treatment may help.

\medterm{Hearing Loss: Sudden Deafness} A rapid sensorineural hearing loss, usually developing over hours to a few days, that may affect one or both ears. It is an otologic emergency; urgent hearing testing and evaluation are needed because early treatment may improve recovery.

\medterm{Hearing Screening} A quick test used to identify possible hearing loss before symptoms are obvious. Newborns are commonly screened before leaving hospital; an abnormal result requires repeat testing and a full hearing assessment rather than confirming permanent loss by itself.

\medterm{Hearing Test} An audiologic assessment that measures hearing thresholds and speech perception across frequencies to identify conductive, sensorineural, or mixed hearing loss.

\medterm{Heart} A four-chambered muscular pump (two atria, two ventricles) located in the mediastinum,
responsible for circulating blood throughout the body.

\medterm{Heart And Vascular Services} Heart and vascular services include prevention, diagnosis, and treatment of diseases affecting the heart and blood vessels.

\medterm{Heart Aneurysm} A localized bulging or outpouching of a weakened heart wall, often after myocardial infarction or from cardiomyopathy, that can impair contraction or form clots and arrhythmias.

\medterm{Heart Arrhythmia} An abnormal heart rate or rhythm caused by altered electrical activity in the heart. At rest, the adult heart usually beats regularly at about 60--100 times per minute; an arrhythmia may make it beat too fast, too slowly, or irregularly.

\synonyms
Cardiac Dysrhythmias; Irregular Heart Rhythm

\medterm{Heart Attack} A heart attack occurs when blood flow carrying oxygen to part of the heart muscle suddenly becomes blocked, usually by a blood clot in a narrowed heart artery. The lack of oxygen can injure or destroy heart muscle and may cause chest pressure, shortness of breath, sweating, nausea, or other symptoms.

\medterm{Heart Attack First Aid} Heart attack first aid means calling emergency services immediately, helping the person rest, and following dispatcher instructions; aspirin may be appropriate for some adults who can safely take it.

\medterm{Heart Attack Prevention} Heart-attack prevention targets blood pressure, atherogenic cholesterol, diabetes, tobacco exposure, diet, physical activity, weight, sleep, and adherence to prescribed medicines. Risk-based preventive treatment is more effective than relying on one supplement or symptom.

\medterm{Heart Block} A delay or interruption in the electrical signal traveling from the heart’s upper chambers to its lower chambers. It may cause a slow or irregular heartbeat, dizziness, fainting, or no symptoms.

\medterm{Heart Bypass Surgery} Coronary artery bypass grafting creates new routes around blocked coronary arteries using graft vessels to improve myocardial blood flow and relieve ischemic symptoms or reduce selected risks.

\medterm{Heart Catheterization} A procedure in which a catheter is passed through a blood vessel into the heart to measure pressures, assess coronary arteries, inject contrast, obtain samples, or perform treatment.

\medterm{Heart Contraction} The squeezing of heart muscle that pushes blood from the lower chambers to the lungs and the rest of the body. Electrical signals coordinate the heartbeat, and abnormal contraction can reduce blood flow or cause rhythm problems.

\medterm{Heart CT Scan} Computed tomography of the heart and nearby structures, sometimes with ECG synchronization or contrast, used to evaluate coronary arteries, cardiac anatomy, calcium, or selected masses.

\medterm{Heart Defect} An abnormality of the heart or major blood vessels. It may be present from birth or develop later and can affect blood flow, heart rhythm, valve function, or pumping ability.

\medterm{Heart Disease} A broad term for conditions affecting the heart or blood vessels, including coronary artery disease, heart failure, abnormal rhythms, valve disease, and congenital defects. Symptoms and treatment depend on the part affected and may include chest pain, breathlessness, palpitations, or swelling.

\medterm{Heart Disease And Depression} Depression can occur with heart disease and may worsen quality of life, treatment adherence, and recovery; both conditions should be assessed and treated.

\medterm{Heart Disease And Intimacy} Heart disease or its treatment may affect sexual activity, desire, and confidence; individualized medical advice can help address symptoms and safety.

\medterm{Heart Disease And Women} Heart disease in women may involve different risk factors, symptoms, and timing than in men; chest pressure, shortness of breath, nausea, fatigue, or jaw or back pain can be important symptoms.

\medterm{Heart Failure} A condition in which the heart cannot pump or fill with enough blood to meet the body's needs. It does not mean that the heart has stopped. In systolic heart failure, the heart does not squeeze strongly enough; in diastolic heart failure, the heart is stiff and does not relax or fill properly. Symptoms may include shortness of breath, tiredness, reduced ability to exercise, and swelling caused by fluid buildup.

\medterm{Heart Failure In Children} A condition in which a child's heart cannot pump or fill well enough to meet the body's needs. It may cause fast breathing, sweating or tiring during feeds, poor growth, swelling, or reduced activity.

\medterm{Heart Failure Stages} A framework describing progression from risk without heart disease, to heart changes without symptoms, to heart failure with symptoms, and finally advanced disease with persistent symptoms. Stages help guide prevention and treatment and are different from scales that describe how limited a person feels during activity.

\medterm{Heart Failure With Mildly Reduced Ejection Fraction} Heart failure in which the left ventricle pumps out a somewhat smaller-than-usual proportion of its blood, commonly about 41--49 percent. Symptoms may include breathlessness, tiredness, and fluid-related swelling.

\medterm{Heart Failure With Preserved Ejection Fraction} Heart failure in which the left ventricle still squeezes out a usual proportion of its blood but is too stiff to fill normally. Fluid then backs up, causing breathlessness, tiredness, and swelling.

\medterm{Heart Failure With Reduced Ejection Fraction} Heart failure in which the left ventricle cannot squeeze strongly enough to supply the body. It may cause breathlessness, tiredness, and swelling and can result from blocked arteries, heart-muscle disease, high blood pressure, or valve disease.

\medterm{Heart Injury} Damage to heart muscle or surrounding cardiac structures caused by trauma, ischemia, inflammation, toxins, or procedures; symptoms and tests depend on the injured tissue and severity.

\medterm{Heart MRI} Magnetic resonance imaging of the heart that evaluates cardiac structure, function, tissue characteristics, blood flow, and selected congenital, ischemic, inflammatory, or infiltrative disorders.

\medterm{Heart Murmur} An extra or unusual sound heard when blood flows through the heart or nearby blood vessels. Some murmurs are harmless, while others result from a valve problem, heart defect, or altered blood flow.

\medterm{Heart Murmurs} Heart murmurs are extra or unusual sounds caused by turbulent blood flow through the heart or nearby vessels; some are harmless and others reflect heart disease.

\medterm{Heart Muscle} The myocardium, the specialized cardiac muscle that contracts rhythmically to pump blood through the heart and circulation.

\medterm{Heart Neoplasm} An abnormal growth in the heart or its surrounding tissues. It may be benign or cancerous and can interfere with blood flow, heart rhythm, or pumping; evaluation uses imaging and sometimes biopsy.

\medterm{Heart Pacemaker} A heart pacemaker is a device that sends electrical impulses to maintain an appropriate heartbeat when the heart's natural rhythm is too slow or unreliable.

\medterm{Heart Palpitations} An awareness of the heartbeat as pounding, racing, fluttering, or irregular.
Palpitations may occur with normal rhythm or arrhythmia and can be triggered by stress,
exertion, fever, stimulants, endocrine disease, or structural heart disease.

\synonyms
Cardiac Palpitations; Ectopic Beats

\medterm{Heart PET Scan} Positron-emission tomography of the heart using a radiotracer to assess myocardial perfusion, metabolism, viability, inflammation, or selected cardiac masses.

\medterm{Heart Rate} The number of heartbeats per minute. Resting rates vary with age, fitness, activity, illness, medicines, and emotion; a persistently very fast, slow, or irregular rate may need assessment.

\medterm{Heart Rhythm} The rate and pattern of cardiac electrical activation and contraction, normally initiated by the sinoatrial node; abnormal rhythms include bradycardia, tachycardia, atrial fibrillation, and ventricular arrhythmias.

\medterm{Heart Rhythm Disorders} Arrhythmias, in which the heart beats too fast, too slowly, or irregularly because of abnormal electrical impulse formation or conduction.

\synonyms
Cardiac Arrhythmias; Dysrhythmias

\medterm{Heart Septum} The muscular or membranous wall separating the right and left sides of the heart, including the atrial and ventricular septa; defects can permit abnormal blood flow between chambers.

\medterm{Heart Sounds} The sounds made as the heart valves open and close during each heartbeat. The usual two sounds are called “lub-dub”; extra or unusual sounds may suggest a valve problem, heart failure, or another condition.

\medterm{Heart Surgery} Surgical procedures performed on the heart or great vessels to treat structural or
functional cardiac conditions. Common operations include coronary artery bypass grafting
(CABG), valve repair or replacement, and repair of congenital heart defects.

\medterm{Heart Tamponade} A life-threatening condition in which fluid or blood builds up around the heart and prevents it from filling normally. It can cause breathlessness, faintness, low blood pressure, and shock and requires emergency treatment.

\medterm{Heart Transplant} Replacement of a failing heart with a donor heart, followed by lifelong immunosuppression and surveillance for rejection, infection, coronary allograft disease, kidney injury, and medication complications.

\medterm{Heart Transplantation} Surgical replacement of a failing heart with a donor heart, usually for selected patients with advanced heart failure who meet specialist criteria.

\medterm{Heart Valve} A flap-like structure that directs one-way blood flow through the heart. The four valves are the tricuspid, pulmonary, mitral, and aortic
valves.

\medterm{Heart Valve Disease} Structural or functional abnormality of one or more cardiac valves causing stenosis, regurgitation, or both. Common causes include degeneration, congenital malformation, rheumatic disease, and infective endocarditis.

\medterm{Heart Valve Disorder} Alternate term; see \textit{Heart valve disease}.

\medterm{Heart Valve Prosthesis} An artificial valve used to replace or repair a diseased heart valve. Mechanical valves are durable but usually require long-term blood-thinning medicine. Tissue valves may avoid lifelong blood thinners but can wear out and sometimes need another procedure.

\medterm{Heart Valve Surgery} Repair or replacement of a diseased cardiac valve to correct severe narrowing or leakage and improve blood flow; choice of repair, mechanical valve, or tissue valve depends on anatomy and patient factors.

\medterm{Heart Ventricle} One of the two lower heart chambers that pumps blood to the lungs or systemic circulation; the right ventricle pumps to the pulmonary artery and the left to the aorta.

\medterm{Heart Ventricles} The two lower chambers of the heart. The right ventricle pumps blood to the lungs, and the left ventricle pumps blood to the body.

\medterm{Heartburn} A burning feeling behind the breastbone caused by stomach contents flowing back into the esophagus. It may be accompanied by a sour taste or regurgitation and is not usually caused by the heart. Frequent symptoms may indicate reflux disease.

\medterm{Heartland Virus} A tick-borne virus associated with fever, fatigue, headache, muscle aches, diarrhea, and low white-cell or platelet counts. It can resemble ehrlichiosis and may be severe, especially in older or medically vulnerable people.

\medterm{Heart-Lung Machine} A heart-lung machine temporarily circulates and oxygenates blood during some heart or major blood-vessel operations while the heart is stopped.

\medterm{Heart-Lung Machine} A device that temporarily takes over the pumping and gas-exchange functions of the heart and lungs during selected cardiac operations.

\synonyms
Cardiopulmonary Bypass Pump; Extracorporeal Circuit

\medterm{Heat} Thermal energy that raises tissue or body temperature; excessive heat exposure can cause heat cramps, heat exhaustion, heatstroke, or thermal injury.

\medterm{Heat Cramps} Painful, involuntary muscle spasms occurring during or after strenuous activity in heat, often related to fluid and
electrolyte loss. Resting in a cool place and replacing fluids and electrolytes may help. Confusion, fainting, or
ongoing symptoms may signal heat illness requiring urgent medical attention.

\medterm{Heat Emergencies} Heat emergencies range from heat cramps and exhaustion to heat stroke, in which dangerous core-temperature elevation causes central nervous system dysfunction and organ injury.

\medterm{Heat Exhaustion} A heat-related illness caused by losing too much water and salt through sweating. It may cause heavy sweating, weakness, dizziness, headache, nausea, or fainting. Confusion, seizures, or very hot skin suggest heat stroke and require emergency care.

\medterm{Heat Intolerance} An abnormal inability to tolerate warm environments or increased body heat, often
causing excessive sweating, flushing, fatigue, or dizziness. Common causes include
hyperthyroidism, menopause, fever, medications, autonomic disorders, and dehydration.

\medterm{Heat Pack} A device that applies warmth to the body surface to relieve some muscle pain or stiffness. Excessive heat or prolonged contact can cause burns.

\medterm{Heat Prostration} An older term for severe heat exhaustion, marked by weakness, fatigue, dizziness, headache, nausea, and sweating after heat exposure; confusion or loss of consciousness suggests heat stroke.

\medterm{Heat Rash} An itchy or prickly rash caused by blocked sweat ducts during hot, humid conditions. Small red bumps or clear vesicles
often appear in skin folds or areas covered by clothing. Cooling, keeping the skin dry, and avoiding occlusive products
usually help; infection or persistent symptoms warrants assessment.

\medterm{Heat Shock Protein} A stress-response protein that helps other proteins fold, avoid aggregation, and recover after heat, oxidative, infectious, or other cellular stress; some also influence immunity and cancer biology.

\medterm{Heat Stress} The strain placed on the body when it cannot adequately lose heat, often during hot or humid conditions, exertion, or inadequate hydration. It can progress from heat cramps or exhaustion to life-threatening heat stroke.

\medterm{Heat Stroke} A life-threatening heat-related illness with markedly elevated body temperature and altered mental status. It may cause confusion, seizures, loss of consciousness, and damage to multiple organs.

\medterm{Heavy Menstrual Bleeding} Menstrual bleeding that is unusually heavy, prolonged, or disruptive to daily life. Causes include fibroids, hormone changes, bleeding disorders, medicines, pregnancy-related problems, or cancer and may cause anemia.

\synonyms
Menorrhagia; Hypermenorrhea

\medterm{Heavy Metal Poisoning} Harm from excessive exposure to metals such as lead, mercury, arsenic, cadmium, or thallium. Effects may include abdominal pain, anemia, nerve injury, kidney damage, or confusion; treatment starts with removing exposure and may include specialist antidote therapy.

\medterm{Hebefrenia} Alternate historical term; see \textit{Schizophrenia}.

\medterm{Heberden Node} A firm bony bump at the finger joint nearest the fingertip, usually caused by osteoarthritis. It may be painful or tender while developing and can cause stiffness or reduced finger movement.

\medterm{Heberden Nodes} Firm bony enlargements at the finger joints nearest the fingertips, commonly caused by osteoarthritis. They may be painful or tender while forming, then remain as permanent bumps that limit movement in some people.

\synonyms
Heberden's Arthritic Nodes; DIP Joint Osteophytes

\medterm{Heberden's Node} A hard, bony enlargement of the distal interphalangeal joint caused by osteoarthritis, usually affecting the fingers and sometimes producing pain, stiffness, or reduced dexterity.

\medterm{Heberden's Nodes} Firm bony enlargements of the distal interphalangeal finger joints, usually caused by osteoarthritis.

\medterm{Hecht's Pneumonia} A rare, severe form of measles-related pneumonia, usually affecting people with weakened immunity. It can cause rapidly worsening breathing problems and requires urgent hospital treatment.

\medterm{Hedgehog Proteins} Signaling proteins that help guide the growth and organization of tissues before birth and during repair. Abnormal activity in this pathway can contribute to some birth differences and cancers.

\medterm{Heel} The heel is the back and underside of the foot, formed mainly by the calcaneus and a cushioning fat pad.

\medterm{Heel Pain} Pain beneath or around the heel, commonly caused by plantar fasciitis, Achilles
tendinopathy, bursitis, stress fracture, nerve compression, or inflammatory arthritis.

\synonyms
Plantar Fasciitis Discomfort; Calcaneal Spur Pain

\medterm{Heel Spur} A bony growth arising from the heel bone, often associated with plantar-fascia strain. The spur itself
may cause no symptoms; heel pain is more often related to surrounding soft-tissue inflammation.

\medterm{Heel-to-Shin Test} A neurological test of lower limb coordination in which the patient places the heel of one foot on the opposite knee and slides it smoothly down the crest of the tibia to the ankle. Inability to maintain contact or an irregular, jerky trajectory indicates cerebellar ataxia or sensory ataxia.

\medterm{Hegar's Dilators} Hegar's dilators are graduated metal instruments used to gently widen the cervix or other body openings during selected procedures.

\medterm{Hegar's Sign} Softening of the lower part of the uterus, once used as an early sign of pregnancy. It is not specific and is not sufficient to confirm pregnancy without testing.

\medterm{Height} The vertical distance from the base of a person to the top of the head, measured for growth assessment, nutritional evaluation, drug dosing, and calculation of body-mass index.

\medterm{Heimlich Maneuver} A common name for abdominal thrusts used to help a responsive person with severe choking. First-aid steps differ for adults, children, infants, and pregnant people; emergency services should be called for serious airway obstruction.

\medterm{Heimlich Maneuver On Self} A self-administered abdominal-thrust technique sometimes used for severe choking when a person cannot breathe, speak, or cough effectively; emergency services should be activated immediately.

\medterm{Heimlich Manoeuvre} An emergency procedure for relieving upper airway obstruction caused by a foreign body
(choking). The maneuver is performed on a conscious, choking adult who cannot speak,
cough effectively, or breathe.

\medterm{Heine-Medin Disease} An older name for polio, a viral infection that can inflame the spinal cord and brain. It may cause fever and weakness and sometimes permanent paralysis; vaccination prevents most cases.

\medterm{Heinz Bodies} Clumps of damaged hemoglobin inside red blood cells, usually formed after oxidative injury. They may occur in glucose-6-phosphate dehydrogenase deficiency or after exposure to certain medicines, chemicals, or infections.

\medterm{Hejama} An Arabic term for cupping, a traditional practice in which suction is applied to the skin; wet cupping also involves superficial skin incisions and carries bleeding and infection risks.

\medterm{Helcococcus} A genus of gram-positive cocci that can rarely cause opportunistic skin, wound, bone, or bloodstream infections.

\medterm{Helcococcus Kunzii} A Gram-positive round bacterium usually found on skin that rarely causes opportunistic infection. Reported infections may involve wounds, blood, or implanted devices in people with serious illness or weakened defenses.

\medterm{Helcococcus Pyogenes} A rarely encountered Gram-positive round bacterium that may cause opportunistic human infection. It has been found in wounds or blood, particularly in people with serious illness, chronic wounds, or weakened immunity.

\medterm{Helicase} An enzyme that uses cellular energy to separate the paired strands of DNA or RNA. This allows genetic material to be copied, repaired, or read.

\medterm{Helicobacter} A genus of spiral-shaped bacteria that colonize the stomach or intestine; some species cause ulcers, gastritis, or other disease.

\medterm{Helicobacter Cinaedi} A Helicobacter species that can cause bloodstream, skin, bone, joint, or other invasive infections, especially in people with weakened immunity.

\medterm{Helicobacter Fennelliae} An intestinal Helicobacter species that rarely causes bloodstream or other opportunistic infection. Disease is more likely in people with weakened defenses, and diagnosis requires specialized laboratory identification.

\medterm{Helicobacter Pylori} A gram-negative, spiral-shaped bacterium adapted to colonize the stomach. It can cause chronic
gastritis, peptic ulcer disease, and increase the risk of gastric cancer and gastric MALT
lymphoma.

\medterm{Helicobacter Pylori Infection} Infection of the stomach by H. pylori, a spiral-shaped bacterium associated with chronic
gastritis, peptic ulcer disease, gastric adenocarcinoma, and gastric MALT lymphoma. It
is diagnosed with noninvasive or endoscopic testing and treated with combination
eradication therapy.

\medterm{Helicobacteraceae} A family of spiral-shaped bacteria that includes Helicobacter and related genera, many of which inhabit the digestive tract.

\medterm{Heliotherapy} Treatment using controlled exposure to sunlight or artificial ultraviolet light. It has been used for selected skin disorders, but excessive exposure can cause burns, premature aging, eye injury, and skin cancer.

\medterm{Helium} A chemically inert gas used in respiratory mixtures, cryogenics, imaging, and laboratory applications; inhalation can displace oxygen and cause hypoxia despite its lack of direct toxicity.

\medterm{Helix} The curved outer rim of the external ear. Its shape is formed by cartilage covered with skin; swelling, tenderness, or a skin change on the helix may follow injury, infection, sun damage, or other skin disease.

\medterm{Heller Myotomy} A surgical procedure that cuts the muscles of the lower esophageal sphincter to relieve obstruction in achalasia, usually performed laparoscopically and often combined with an antireflux procedure.

\medterm{Heller's Operation} A surgical procedure that cuts the lower esophageal muscle to relieve achalasia, usually with an anti-reflux procedure.

\medterm{HELLP Syndrome} A dangerous pregnancy complication involving destruction of red blood cells, liver injury, and a low platelet count. It may cause upper-abdominal pain, nausea, headache, or high blood pressure and can occur with preeclampsia. Urgent obstetric assessment is required.

\medterm{Helminth} A parasitic worm, including roundworms, tapeworms, and flukes, that may live in or on a host.

\medterm{Helminthiasis} Infection or infestation caused by parasitic worms, including roundworms, tapeworms, and flukes. Symptoms vary by species and may include abdominal problems, anemia, cough, rash, malnutrition, or organ damage.

\synonyms
Worm Infestation; Helminthic Parasitosis

\medterm{Helminths} Parasitic worms that cause disease in humans. The major groups are nematodes
(roundworms), cestodes (tapeworms), and trematodes (flukes). Helminths have complex life
cycles often involving intermediate hosts. Infection occurs through ingestion of eggs or
larvae, penetration of skin, or insect vectors.

\medterm{Help Prevent Hospital Errors} Actions that reduce hospital errors include confirming identity, checking medicines, communicating clearly, asking questions, and reporting changes or concerns promptly.

\medterm{Helper T Cell} A type of white blood cell that coordinates immune responses by sending chemical signals that activate and guide other immune cells, including antibody-producing cells and cells that destroy infected cells.

\medterm{Helper T Cells} A type of white blood cell that helps organize the immune response. These cells release chemical messages that activate and guide other immune cells, including cells that make antibodies or destroy infected cells.

\synonyms
CD4+ T Lymphocytes; T Helper Cells

\medterm{Helper Virus} A virus that supplies functions needed for another virus or defective viral genome to replicate.

\medterm{HELPERR Mnemonic} A memory aid for managing shoulder dystocia during birth: call for help, consider an episiotomy, flex the birthing person’s legs, apply suprapubic pressure, perform internal maneuvers, remove the baby’s posterior arm, and roll to an all-fours position. It requires trained obstetric staff.

\medterm{Helping A Loved One With A Drinking Problem} Support for someone with alcohol use disorder includes calm nonjudgmental communication, boundaries, encouragement toward medical treatment, attention to withdrawal risk, and emergency action for overdose or dangerous withdrawal.

\medterm{Hemagglutination Test} A test in which antibodies or other agents cause red blood cells to clump; the reaction can help detect antibodies, antigens, or selected infectious agents.

\medterm{Hemagglutinin} A surface glycoprotein, especially on influenza viruses, that binds sialic-acid receptors and helps the virus attach to and enter host cells; it is a major target of protective antibodies.

\medterm{Hemangioblastoma} A usually slow-growing tumor made of blood-vessel-forming cells, most often found in the cerebellum or spinal cord. It may cause headache, balance problems, or weakness and can be associated with von Hippel–Lindau disease.

\medterm{Hemangioendothelioma} A rare tumor arising from cells that line blood vessels. Its behavior ranges from slow-growing to cancerous, depending on the subtype and location, so diagnosis requires tissue examination and specialist assessment.

\medterm{Hemangioma} A benign vascular tumor formed by the abnormal proliferation of blood vessel endothelial cells. Hemangiomas are broadly classified based on vessel caliber and clinical presentation into capillary hemangiomas (such as infantile hemangiomas that proliferate in infancy and spontaneously involute), cavernous hemangiomas (composed of larger, mature blood-filled vascular spaces that do not regress), and lobular capillary hemangiomas (pyogenic granulomas). While many are asymptomatic or undergo spontaneous involution, problematic lesions causing functional impairment, ulceration, or visual compromise (e.g., periocular capillary hemangiomas) are treated medically with beta-blockers (propranolol), laser therapy, or surgical excision.

\synonyms
Vascular Hemangioma; Benign Vascular Tumor

\medterm{Hemangiopericytoma} A rare tumor formerly defined by branching “staghorn” vessels and perivascular cells; many lesions classified this way are now recognized as solitary fibrous tumors.

\medterm{Hemangiosarcoma} A malignant tumor of vascular endothelial cells, an aggressive sarcoma that may arise in
the skin, liver, or bone.

\medterm{Hemarthrosis} Bleeding into a joint, most often the knee, causing rapid swelling, warmth, pain, and limited movement. Injury and bleeding disorders such as hemophilia are common causes; repeated episodes can damage cartilage and lead to long-term joint disease.

\medterm{Hematemesis} Vomiting blood or material that looks like coffee grounds, usually from bleeding in the esophagus, stomach, or first part of the small intestine. It can be life-threatening and requires emergency medical assessment.

\medterm{Hematidrosis} A very rare condition in which blood-tinged fluid appears through apparently intact skin or other surfaces, sometimes during extreme stress. The diagnosis requires careful examination to exclude injury, skin disease, bleeding disorders, or contamination.

\medterm{Hematocele} A collection of blood around a testicle, usually after scrotal injury or surgery and sometimes from a tumor. It can cause a heavy, firm, painful swelling; ultrasound is needed to check the testicle and blood flow, and large collections may need surgery.

\synonyms
Scrotal Hematocele; Tunica Vaginalis Blood Collection

\medterm{Hematochezia} Passing bright-red or maroon blood from the rectum. It often comes from the large intestine, rectum, or anus, but heavy bleeding higher in the digestive tract can look the same and requires urgent assessment.

\medterm{Hematocolpos} Accumulation of menstrual blood in the vagina, usually because an imperforate hymen or another obstructing anomaly prevents outflow; it may cause cyclic pelvic pain and a bulging vaginal membrane.

\medterm{Hematocrit} The proportion of blood volume occupied by packed red blood cells, expressed as a percentage. One of the most commonly ordered laboratory tests in clinical medicine.

\medterm{Hematogenous} Spread through the bloodstream, as when infection, tumor cells, or embolic material travels from one body site to another.

\medterm{Hematologic Agents} Medicines or biologic therapies that act on blood cells, coagulation, marrow, or immune pathways, including agents used for anemia, thrombosis, leukemia, and lymphoma.

\medterm{Hematologic Changes In Pregnancy} Pregnancy normally increases the fluid part of blood more than the red-cell mass, so blood counts may show mild dilutional anemia. White-cell counts also commonly rise, mainly because of more neutrophils. Results must be interpreted using pregnancy-specific ranges.

\medterm{Hematologic Malignancy} Alternate name; see \textit{Blood Cancer}.

\medterm{Hematologic Response} A measurable improvement or change in blood counts or marrow findings after treatment, such as recovery of neutrophils, platelets, or hemoglobin or reduction of malignant cells.

\medterm{Hematological Disease} A disorder of blood cells, bone marrow, lymphatic tissue, or coagulation, including anemia, leukemia, lymphoma, bleeding disorders, and clotting disorders.

\medterm{Hematologist} A physician who diagnoses and treats disorders of blood, bone marrow, lymph nodes, clotting, and hematologic cancers.

\medterm{Hematology} The branch of medicine concerned with the study, diagnosis, treatment, and prevention of
disorders of the blood and blood-forming organs (bone marrow, lymph nodes, spleen).

\medterm{Hematology-Oncology} The medical specialty concerned with blood disorders and cancers, including anemia, leukemia, lymphoma, myeloma, and clotting disorders.

\medterm{Hematoma} A collection of blood outside a blood vessel, usually after injury, surgery, or bleeding caused by illness or medicines. It may cause swelling and pain and can press on nearby nerves or organs.

\medterm{Hematometra} Blood trapped inside the uterus because the normal outlet is blocked or too narrow. It can cause absent periods, monthly pelvic pain, and an enlarged tender uterus; the blockage may be present from birth or develop after scarring, surgery, infection, or radiation.

\medterm{Hematometrocolpos} A buildup of menstrual blood in the uterus and vagina because the normal outlet is blocked, often by tissue present from birth. It may cause absent periods, monthly pelvic pain, and a lower-abdominal mass.

\medterm{Hematomyelia} Bleeding inside the spinal cord, usually after a serious injury or from an abnormal blood vessel or bleeding disorder. Sudden neck or back pain, weakness, numbness, or loss of bladder or bowel control is an emergency requiring immediate imaging and treatment.

\medterm{Hematophobia} An excessive or persistent fear of blood, blood-related sights, or medical procedures involving blood that causes distress or avoidance; some people have a vasovagal fainting response.

\medterm{Hematopoiesis} The process of making and maturing blood cells, mainly in the bone marrow. It produces red blood cells for oxygen transport, white blood cells for immune defense, and platelets for stopping bleeding.

\medterm{Hematopoietic Growth Factor} A cytokine or medicine that stimulates production, maturation, or survival of blood-forming cells in bone marrow, such as erythropoietin or granulocyte colony-stimulating factor.

\medterm{Hematopoietic Neoplasm} A cancer arising from blood-forming cells in the bone marrow or immune tissues. Leukemias, lymphomas, and plasma-cell cancers are examples; symptoms and treatment depend on the exact cell type and extent.

\medterm{Hematopoietic Stem Cell Transplantation} Infusion of blood-forming stem cells to replace damaged or diseased bone marrow. The cells may come from the patient or a donor. Treatment can cause serious infection, bleeding, infertility, or an immune reaction against donor and recipient tissues.

\medterm{Hematopoietic System} The bone marrow and related organs and cells that produce, mature, circulate, and remove blood cells, including red cells, white cells, and platelets.

\medterm{Hematoporphyrin Derivative} A porphyrin-based photosensitizer that accumulates in some tumors and can be activated by light in photodynamic treatment; it may also cause prolonged photosensitivity.

\medterm{Hematosalpinx} Accumulation of blood within a fallopian tube, often caused by ectopic pregnancy, endometriosis, pelvic surgery, or obstruction, and potentially producing pelvic pain or infertility.

\medterm{Hematosis} The exchange of gases in the lungs: oxygen moves from inhaled air into the blood, while carbon dioxide moves from the blood into the air to be breathed out.

\medterm{Hematospermia} Blood in the semen, often temporary and associated with infection, inflammation, trauma, or a recent urologic procedure. Persistent or recurrent bleeding, pain, fever, urinary symptoms, or bleeding in an older adult requires medical evaluation.

\medterm{Hematoxylin} A dye used in histology, usually with eosin, to stain cell nuclei blue-purple and help show tissue structure.

\medterm{Hematuria} Blood in the urine. It may be visible as pink, red, or brown urine or found only under a microscope. Causes include infection, stones, medicines, injury, prostate disease, and kidney or urinary-tract disorders.

\medterm{Heme} Heme is an iron-containing porphyrin group that binds oxygen in haemoglobin and myoglobin and forms the functional centre of several enzymes; abnormal heme metabolism contributes to porphyrias and jaundice.

\medterm{Hemeralopia} Hemeralopia is reduced vision in bright light, sometimes associated with retinal disorders or abnormal cone function.

\medterm{Hemianopia} Loss of vision in half of the area a person can see, affecting one or both eyes. It often results from a problem in the brain’s visual pathways, such as a stroke or tumor, and sudden onset is an emergency.

\medterm{Hemianopsia} Loss of vision in half of the visual field of one or both eyes, usually from a lesion of the optic tract, radiations, or occipital cortex; sudden onset may indicate stroke.

\medterm{Hemiballism} A movement disorder causing sudden, large-amplitude, flinging movements on one side of the body, usually from a lesion involving the subthalamic region.

\medterm{Hemiballismus} Hemiballismus is involuntary, rapid, flinging movement of one side of the body, usually caused by a lesion of the opposite subthalamic region.

\medterm{Hemicolectomy} Surgical removal of the right or left portion of the colon, sometimes with nearby lymph nodes and reconnection of the bowel.

\medterm{Hemidesmosome} A specialized junction that anchors basal epithelial cells to the underlying basement membrane by linking intracellular keratin filaments with extracellular matrix proteins.

\medterm{Hemidiaphragm} One of the two muscular halves of the diaphragm separating the chest from the abdomen and assisting breathing.

\medterm{Hemifacial Atrophy} Progressive wasting of skin, subcutaneous tissue, muscle, and sometimes bone on one side of the face, also called Parry--Romberg syndrome.

\medterm{Hemifacial Microsomia} A birth difference in which one side of the face is smaller or less developed than the other, often affecting the jaw, ear, cheek, and nearby soft tissues. Hearing, vision, and spinal differences may also occur.

\synonyms

\medterm{Hemifacial Spasm} Repeated involuntary twitching on one side of the face, usually beginning around the eye and sometimes spreading to the cheek or mouth. Irritation of the facial nerve is a common cause; persistent spasms need evaluation.

\medterm{Hemihypertrophy} Unequal overgrowth of one side or region of the body, producing asymmetry of limbs, trunk, or face; some childhood cases carry increased risk of embryonal tumors.

\medterm{Hemilaryngectomy} Hemilaryngectomy is surgery to remove part of the larynx, usually to treat a selected laryngeal cancer while preserving as much speech and swallowing function as possible.

\medterm{Hemimandibulectomy} Hemimandibulectomy is surgery to remove one side of the lower jaw, usually for a tumor or severe bone disease, sometimes followed by reconstruction.

\medterm{Hemimegalencephaly} A rare developmental brain malformation in which one cerebral hemisphere is abnormally enlarged and dysplastic, often causing infantile seizures, developmental impairment, and contralateral neurologic signs.

\medterm{Hemimelia} A congenital limb-reduction defect in which part or all of a limb segment is absent, such as fibular or radial hemimelia.

\medterm{Hemiparesis} Weakness on one side of the body, affecting the face, arm, leg, or a combination. It is less severe than paralysis and may result from stroke, brain injury, a tumor, multiple sclerosis, or another nervous-system disorder.

\medterm{Hemipelvectomy} Hemipelvectomy is major surgery that removes part or all of one side of the pelvis, with or without removal of the leg, usually for extensive cancer or severe trauma.

\medterm{Hemiplegia} Complete or near-complete paralysis of one side of the body, including the face, arm, or leg. It usually results from injury to movement pathways in the brain, commonly a stroke, and sudden onset is an emergency.

\medterm{Hemispatial Neglect} Hemispatial neglect is reduced awareness of one side of the body or space, usually after injury to the opposite cerebral hemisphere.

\medterm{Hemispherectomy} Hemispherectomy is surgery that removes or disconnects much of one cerebral hemisphere to treat severe, otherwise uncontrolled epilepsy in selected patients.

\medterm{Hemisternebra} A congenital or developmental partial segment of the sternum, caused by incomplete formation or fusion of sternal elements and sometimes associated with chest-wall anomalies.

\medterm{Hemithorax} One side of the thorax, including the lung, pleural space, ribs, and other structures on that side of the chest.

\medterm{Hemivertebrae} Congenital vertebrae formed from only part of the normal vertebral body, producing wedge-shaped bone and potentially causing congenital scoliosis or spinal-cord abnormalities.

\medterm{Hemizygous} Having only one copy of a gene or chromosome region rather than the usual pair, as occurs for many X-linked genes in people with one X chromosome.

\medterm{Hemochromatosis} A disorder in which too much iron builds up in the body and can damage the liver, pancreas, heart, joints, skin, and hormone-producing organs. It may be inherited or develop from another illness or repeated transfusions.

\medterm{Hemochromatosis Cardiomyopathy} Heart-muscle damage caused by excess iron buildup. The heart may become stiff or weak and cause breathlessness, swelling, abnormal rhythms, or heart failure; treatment aims to remove excess iron and support heart function.

\medterm{Hemocyanin} A copper-containing oxygen-transport protein in many mollusks and arthropods; unlike human hemoglobin, it circulates dissolved in their hemolymph.

\medterm{Hemocytometer} A ruled counting chamber used with a microscope to determine the concentration of cells or other particles in a measured volume of fluid.

\medterm{Hemodialysis} A treatment for kidney failure in which blood passes through a filter outside the body to remove waste and extra fluid before returning to the person. It usually requires repeated sessions and a blood-vessel access; low blood pressure, infection, clotting, and access problems can occur.

\medterm{Hemodialysis Access Care} Daily maintenance protocols and hygiene measures for an arteriovenous fistula, arteriovenous graft, or central venous catheter used for hemodialysis. Includes daily palpation for a vascular thrill, avoiding blood pressure cuffs or constrictive clothing on the access arm, and promptly reporting signs of infection or thrombosis.

\synonyms
Vascular Access Care; Fistula and Graft Maintenance; Dialysis Access Management

\medterm{Hemodialysis Access Procedures} Procedures that create, maintain, repair, or revise vascular access for hemodialysis, including fistula or graft creation, catheter placement, angioplasty, thrombectomy, and treatment of stenosis or infection.

\medterm{Hemodialysis Clearance} The volume of plasma from which hemodialysis removes a solute per unit time; it depends on blood and dialysate flow, membrane properties, treatment time, and the substance’s distribution.

\medterm{Hemodialysis Solution} The dialysate fluid circulated across a semipermeable membrane during hemodialysis to remove waste and help correct electrolyte, acid--base, and fluid abnormalities.

\medterm{Hemodynamics} The movement of blood through the heart and blood vessels, including pressure, flow, heart output, and resistance. Hemodynamic measurements help assess circulation and guide treatment in conditions such as shock, heart failure, and critical illness.

\medterm{Hemofiltration} A blood-cleaning treatment that removes extra water and dissolved waste through a filter. Replacement fluid is added to maintain the right blood volume; it is used in some forms of temporary or long-term kidney support.

\medterm{Hemoglobin} The iron-containing protein in red blood cells that carries oxygen from the lungs to tissues and helps carry carbon dioxide back to the lungs. Low or abnormal hemoglobin can cause anemia or other blood disorders.

\medterm{Hemoglobin A} The main type of hemoglobin in most adults. It is the oxygen-carrying protein in red blood cells and is made from two alpha and two beta protein parts.

\medterm{Hemoglobin A1c} A routine blood test measuring the percentage of red blood cell hemoglobin coated with sugar (glycated hemoglobin). Because red blood cells live for approximately 120 days (3 months), the test reflects average blood sugar levels over the preceding two to three months, unlike daily fingerstick glucose checks that fluctuate from hour to hour. It is the gold standard for diagnosing and monitoring diabetes and prediabetes (normal is below 5.7\%, prediabetes is 5.7\% to 6.4\%, and diabetes is 6.5\% or higher).
\synonyms
HbA1c; Glycated Hemoglobin; Glycosylated Hemoglobin

\medterm{Hemoglobin A1C Variability} Factors that can make hemoglobin A1c less reliable include altered red-cell survival, hemoglobin variants, anemia,
pregnancy, kidney disease, transfusion, and some treatments. An alternative measure of
glucose may be needed when A1c does not reflect recent glycemia.

\medterm{Hemoglobin C} A beta-globin variant caused by a hemoglobin gene mutation; homozygosity can cause mild hemolytic anemia and splenomegaly, while inheritance with hemoglobin S causes a clinically significant sickling disorder.

\medterm{Hemoglobin C Disease} An inherited hemoglobinopathy caused by a beta-globin variant that produces hemoglobin
C. It may cause mild hemolytic anemia, jaundice, splenomegaly, or, when combined with
hemoglobin S, clinically significant sickling disease.

\medterm{Hemoglobin Electrophoresis} A blood test that separates the different types of hemoglobin. It helps identify inherited conditions such as sickle-cell disease, sickle-cell trait, and thalassemia. Results are interpreted with age, pregnancy, blood counts, recent transfusion, and other tests.

\medterm{Hemoglobin H} An unusual form of hemoglobin made from four beta protein parts when the body makes too little alpha hemoglobin. It is unstable and can cause ongoing breakdown of red blood cells in alpha-thalassemia.

\medterm{Hemoglobin Increased} An increased hemoglobin concentration, which may reflect dehydration, chronic hypoxemia, erythrocytosis, or a myeloproliferative neoplasm and must be interpreted with hematocrit and clinical context.

\medterm{Hemoglobin M} A group of abnormal hemoglobins that stabilize iron in the ferric state and can cause congenital methemoglobinemia and cyanosis.

\medterm{Hemoglobin Normal Values} Reference ranges for blood hemoglobin concentration, which vary with age, sex, pregnancy, altitude, laboratory method, and other factors; results should be interpreted using the reporting laboratory's range.

\medterm{Hemoglobin S} An inherited form of hemoglobin that can make red blood cells become stiff and sickle-shaped when oxygen is low. It can cause anemia, painful blockages, and organ damage in sickle-cell disease.

\medterm{Hemoglobin SC Disease} An inherited form of sickle cell disease caused by having one hemoglobin S gene and one hemoglobin C gene. It can cause anemia, pain episodes, and organ complications.
\synonyms
HbSC; Sickle-Hemoglobin C Disease

\medterm{Hemoglobin Trait} The inherited presence of a hemoglobin variant in a heterozygous state, usually causing no symptoms but potentially affecting screening results or reproductive risk.

\medterm{Hemoglobinopathies: Thalassemias} Thalassemias are inherited disorders of globin chain synthesis imbalance:
alpha-thalassemia (HBA1/HBA2 deletions) and beta-thalassemia (HBB mutations).

\medterm{Hemoglobinopathy} An inherited disorder affecting the structure or production of hemoglobin. Examples include sickle-cell disease and thalassemia; effects range from no symptoms to anemia, painful blood-vessel blockages, and organ damage.

\medterm{Hemoglobinuria} Haemoglobinuria is free haemoglobin in urine, usually from intravascular haemolysis; urine may be dipstick-positive for blood with few or no red cells, and causes include transfusion reactions, malaria, and paroxysmal nocturnal haemoglobinuria.

\medterm{Hemoglobinuria Test} A laboratory test evaluating urine for hemoglobin or related findings, used to investigate intravascular hemolysis, hemoglobinuria, transfusion reactions, or other causes of dark urine.

\medterm{Hemolysin} A toxin or immune substance that damages the outer covering of red blood cells and makes them break open. Some bacteria produce hemolysins, and the resulting destruction of red cells is called hemolysis.

\medterm{Hemolysis} The early destruction of red blood cells, releasing hemoglobin into the blood. It can occur in the bloodstream or in the spleen and liver and may cause anemia, jaundice, dark urine, tiredness, or shortness of breath.

\medterm{Hemolytic} Relating to hemolysis, the premature destruction of red blood cells, which can release hemoglobin and produce anemia or jaundice.

\medterm{Hemolytic Anemia} Anemia caused by red blood cells being destroyed faster than the body can replace them. It may result from an inherited cell problem, an immune reaction, infection, medicine, toxin, or another illness and can cause jaundice and dark urine.

\medterm{Hemolytic Anemia Caused By Chemicals And Toxins} Anemia caused by accelerated destruction of red blood cells after exposure to a chemical or toxin; evaluation identifies hemolysis and removes or treats the responsible exposure.

\medterm{Hemolytic Anemias: Autoimmune} Anemias in which the immune system mistakenly destroys red blood cells. Warm-antibody and cold-antibody forms differ in triggers and temperature sensitivity; symptoms may include tiredness, jaundice, dark urine, and breathlessness.

\medterm{Hemolytic Anemias: Hereditary} Inherited conditions in which red blood cells are destroyed too quickly. Hereditary spherocytosis is one example; it can cause anemia, jaundice, an enlarged spleen, and gallstones. Other inherited red-cell or enzyme problems can cause similar anemia.

\medterm{Hemolytic Crisis} A sudden increase in red-cell destruction causing worsening anemia, jaundice, dark urine, pain, or other manifestations of hemolysis.

\medterm{Hemolytic Disease Of The Newborn} Anemia and jaundice caused when antibodies from the pregnant person cross the placenta and destroy the fetus’s or newborn’s red blood cells. Severe disease can cause widespread fluid buildup or brain injury without prompt treatment.

\medterm{Hemolytic Index} A laboratory-quality indicator estimating hemolysis in a blood specimen; substantial in-vitro hemolysis can interfere with measurement of several analytes.

\medterm{Hemolytic Jaundice} Jaundice resulting from excessive hemolysis (breakdown of red blood cells) causing
overproduction of unconjugated bilirubin that exceeds the liver's capacity for
conjugation and excretion. See Jaundice for complete overview.

\medterm{Hemolytic Transfusion Reaction} An immune or nonimmune complication in which transfused red blood cells are destroyed, potentially causing fever, pain, hemoglobinuria, kidney injury, shock, or disseminated intravascular coagulation.

\medterm{Hemolytic Uremic Syndrome} A serious illness in which tiny blood clots damage small blood vessels, causing red blood cells to break apart, platelets to fall, and the kidneys to become injured. It often follows severe diarrheal infection but can have other causes.

\medterm{Hemolytic-Uremic Syndrome} A serious illness in which small blood-vessel clots destroy red blood cells, lower the platelet count, and injure the kidneys. It often follows severe diarrhea caused by certain bacteria, but inherited immune problems and other triggers can also cause it.

\medterm{Hemoperfusion} A treatment in which blood passes through a material such as charcoal or resin that traps selected poisons or medicines. It is used only for certain severe overdoses because effectiveness depends on the substance and timing.

\medterm{Hemopericardium} Accumulation of blood in the pericardial sac around the heart. If pressure rises and prevents the heart from filling, cardiac tamponade can occur and requires emergency treatment.

\medterm{Hemoperitoneum} Blood collecting inside the abdominal cavity, usually from trauma, a ruptured blood vessel or organ, surgery, or a ruptured ectopic pregnancy. It can cause abdominal pain, distention, dizziness, or shock and may require emergency imaging and surgery or another procedure.

\medterm{Hemopexin} A plasma glycoprotein that binds free heme and transports it to the liver, helping limit oxidative injury during hemolysis.

\medterm{Hemophagocytic Lymphohistiocytosis} A rare, life-threatening illness in which the immune system remains excessively active and injures the body’s organs. It may be inherited or triggered by infection, cancer, or immune disease. Persistent fever, enlarged spleen, low blood counts, liver problems, bleeding, or confusion require urgent specialist care.

\medterm{Hemophagocytosis Of Sickle Cells} A finding in which immune cells engulf blood cells in the bone marrow or other tissues during severe inflammation, infection, or sickle-cell complications. It can worsen anemia and low platelet counts and requires specialist evaluation.

\medterm{Hemophilia} An inherited bleeding disorder in which the blood cannot clot properly due to a shortage of essential clotting proteins (most commonly Factor VIII in Hemophilia A or Factor IX in Hemophilia B). Because the condition is carried on the X chromosome, it almost exclusively affects males, while females are usually carriers. Beyond prolonged bleeding from minor cuts, its hallmark is spontaneous, painful bleeding inside deep muscles and large joints (such as the knees, elbows, and ankles), which can lead to chronic joint stiffness and deformity over time. Historically known as the ``royal disease'' and ``bleeder's disease,'' it is managed by replacing the missing clotting factor through intravenous infusions.

\synonyms
Haemophilia; Bleeder's Disease; The Royal Disease; Factor Deficiency Disorder; Classic Hemophilia

\medterm{Hemophilia A} An X-linked recessive deficiency of coagulation factor VIII, causing prolonged bleeding into joints, muscles, and soft tissues. Severity correlates with factor VIII activity levels.

\medterm{Hemophilia B} An X-linked recessive deficiency of coagulation factor IX (Christmas disease), causing prolonged bleeding similar to hemophilia A. Severity correlates with factor IX activity levels.

\medterm{Hemophilus} An older spelling or variant of Haemophilus, a group of small gram-negative bacteria that includes respiratory and genital pathogens.

\medterm{Hemophilus Ducreyi} A gram-negative bacterium that causes chancroid, a sexually transmitted infection with painful genital ulcers and tender lymph nodes.

\medterm{Hemophobia} An excessive or persistent fear of blood, also called hematophobia, that causes significant distress, avoidance, or sometimes a vasovagal faint.

\medterm{Hemopoietic Stem Cell} A blood-forming stem cell found mainly in the bone marrow. It can make copies of itself and develop into red blood cells, white blood cells, or platelets.

\synonyms
Hematopoietic Stem Cell (HSC); Blood-Forming Cell

\medterm{Hemopoietic Tissue} Tissue that produces blood cells, chiefly bone marrow, containing stem cells and supporting cells that generate erythrocytes, leukocytes, and platelets.

\medterm{Hemoprotein} A protein containing a heme group with an iron atom, such as hemoglobin, myoglobin, cytochromes, and many oxidoreductive enzymes.

\medterm{Hemoptysis} Coughing up blood or blood-stained mucus from the airways or lungs. Causes include infection, irritated or widened airways, a blood clot, or cancer. More than a few streaks, breathlessness, or chest pain requires urgent care.

\medterm{Hemorrhage} Abnormal escape of blood from a damaged blood vessel or, less commonly, from the heart. It may be external or internal, acute or chronic, and range from minor bleeding to life-threatening blood loss causing anemia, shock, or organ compression; the cause and site determine management.

\medterm{Hemorrhagic} Pertaining to or characterized by hemorrhage (extravasation of blood from vessels). Used in contexts such as hemorrhagic stroke, hemorrhagic shock, and hemorrhagic fever.

\medterm{Hemorrhagic Anemia} Anemia caused by losing blood faster than the body can replace it. Blood loss may be sudden, such as after an injury, or slow, such as from heavy menstrual or intestinal bleeding.

\medterm{Hemorrhagic Disorder} A condition that causes abnormal bleeding because of platelet number or function, clotting-factor deficiency, excessive fibrinolysis, or blood-vessel fragility.

\medterm{Hemorrhagic Shock} Circulatory failure caused by severe blood loss. It may produce rapid pulse, low blood pressure, cool
skin, reduced urine output, confusion, or collapse. Treatment requires urgent control of
bleeding, restoration of circulating volume, blood products when indicated, and correction
of associated injuries or clotting abnormalities.

\medterm{Hemorrhagic Stroke} A stroke caused by a blood vessel breaking and bleeding in or around the brain. It can cause sudden weakness, speech or vision problems, severe headache, vomiting, confusion, or loss of consciousness. The effects depend on the location and amount of bleeding.

\medterm{Hemorrhagic Thyroid Nodule} A sudden bleed inside a thyroid nodule (a lump in the thyroid gland), usually caused by the rupture of fragile blood vessels within an existing fluid-filled cyst or noncancerous growth. It produces an abrupt, tender swelling in the front of the neck that may radiate pain toward the ear or jaw, sometimes accompanied by throat tightness, hoarseness, or mild difficulty swallowing. Although the sudden enlargement can appear alarming, most hemorrhagic nodules are benign and gradually shrink as the trapped blood is reabsorbed over several weeks. It is diagnosed with a thyroid ultrasound, and needle aspiration can drain the fluid to relieve local pressure and confirm the diagnosis.

\synonyms
Intrathyroidal Hemorrhage; Hemorrhagic Thyroid Cyst; Bleeding Into A Thyroid Nodule; Thyroid Nodule Bleed

\medterm{Hemorrhoid} Singular form; see \textit{Hemorrhoids}.

\medterm{Hemorrhoid Surgery} An operation to remove or reduce symptomatic internal or external hemorrhoids when bleeding, prolapse, thrombosis, or pain persists despite appropriate conservative treatment.

\medterm{Hemorrhoidal Hemorrhage} Rectal bleeding from enlarged or injured hemorrhoidal vessels, usually bright red and associated with defecation; heavy, persistent, painful, or unexplained bleeding requires evaluation for other causes.

\medterm{Hemorrhoidectomy} An operation to remove enlarged hemorrhoids that continue to bleed, prolapse, clot, or cause severe symptoms despite other treatment. Recovery may include significant pain, bleeding, difficulty passing urine, or temporary bowel discomfort.

\medterm{Hemorrhoids} Swollen, engorged vascular cushions and surrounding connective tissue in the lower rectum and anal canal, commonly triggered by chronic straining, constipation, pregnancy, or heavy lifting. They are classified by their position relative to the pain-sensitive dentate line: internal hemorrhoids (above the dentate line) are typically painless and manifest with bright red rectal bleeding or mucosal prolapse during defecation; external hemorrhoids (below the dentate line) are covered by sensitive skin and can develop acute, exquisitely painful blood clots (thrombosed external hemorrhoids). Management ranges from dietary fiber, stool softeners, and sitz baths to office procedures (rubber band ligation) and surgical hemorrhoidectomy for advanced disease.
\synonyms{Piles; Hemorrhoidal Disease; Haemorrhoids}

\medterm{Hemosiderin} An iron-containing pigment left in cells after hemoglobin from red blood cells is broken down. It can build up after repeated bleeding, excessive red-cell destruction, or too much iron in the body.

\medterm{Hemosiderosis} Excess deposition of hemosiderin, an iron-storage pigment, in tissues. It may follow repeated bleeding or transfusions and can affect organs such as the lungs or liver.

\medterm{Hemosiderotic Synovitis} Synovial inflammation with hemosiderin deposition caused by recurrent bleeding into a joint, commonly associated with hemophilia, anticoagulation, trauma, or a vascular synovial lesion.

\medterm{Hemospermia} The presence of blood in semen, usually caused by inflammation or infection of the genitourinary tract and often self-limited, but persistent or recurrent cases require evaluation.

\medterm{Hemostasis} The body’s process for stopping bleeding after a blood vessel is injured. Blood vessels narrow, platelets form a plug, and clotting proteins strengthen it; later, the body removes the clot as healing occurs.

\medterm{Hemostasis And Coagulation} The linked processes that stop bleeding. Blood vessels narrow, platelets make a plug, and clotting proteins form a stronger clot. Problems with these processes can cause excessive bleeding or unwanted clots.

\medterm{Hemostasis Valve} A valve or sealing device used with an intravascular catheter or sheath to limit blood loss while allowing passage of wires or instruments.

\medterm{Hemostat} A surgical clamping instrument with jaws that compress a blood vessel or tissue to control bleeding during a procedure. Many hemostats have a ratcheted locking mechanism that maintains pressure without continuous hand force.

\synonyms
Hemostatic Forceps; Artery Clamp

\medterm{Hemostatic Agent} A medicine or material used to stop or reduce bleeding. It may form a physical seal, help the blood clot, slow clot breakdown, or replace a missing clotting protein.

\medterm{Hemothorax} Blood collecting in the space between the lung and chest wall, usually after injury or surgery. It can cause chest pain, breathlessness, low blood pressure, and reduced breath sounds and often requires urgent drainage and treatment of the bleeding source.

\medterm{Hemovac Drain} A closed-suction surgical drain with a compressible reservoir used to remove blood or other fluid from a wound or operative site.

\medterm{Hendra Virus} A rare zoonotic virus carried by fruit bats that can infect horses and occasionally humans through close contact with infected horses. Human disease may cause severe respiratory or neurologic illness and requires public-health management.

\medterm{Henipavirus} A genus of paramyxoviruses that includes Nipah and Hendra viruses, which can cause severe encephalitis and respiratory disease.

\medterm{Henna} A plant dye used on skin, hair, and textiles. Natural henna is usually orange-brown; so-called black henna may contain para-phenylenediamine, which can cause severe allergic skin reactions and permanent scarring.

\medterm{Henoch-SchöNlein Purpura} An older name for IgA vasculitis, a disease in which small blood vessels become inflamed. It can cause purple spots on the legs or buttocks, joint pain, abdominal pain, and kidney problems.

\medterm{Henoch-Schonlein Purpura} An inflammation of small blood vessels, now called IgA vasculitis, that mainly affects children. It can cause raised purple spots on the legs or buttocks, joint pain, abdominal pain, and kidney inflammation. Most cases improve, but kidney monitoring is important.

\synonyms
IgA Vasculitis; Anaphylactoid Purpura; Schönlein-Henoch Syndrome

\medterm{HEPA Filter} A high-efficiency air filter that removes most very small airborne particles. It can reduce dust, allergens, and some infectious particles when properly fitted and maintained, but does not remove every gas or chemical.

\synonyms
High-Efficiency Particulate Air Filter

\medterm{Hepar} Latin for liver; it is used in medical terminology and in names of substances or preparations relating to the liver.

\medterm{Heparan Sulfate} A sulfated glycosaminoglycan in cell surfaces and extracellular matrix that binds growth factors, regulates signaling, and contributes to filtration-barrier and connective-tissue structure.

\medterm{Heparin} An anticoagulant medicine that reduces blood clotting by enhancing antithrombin activity. It may be given by injection to prevent or treat blood clots.

\medterm{Heparin Binding} The attachment of a protein or other molecule to heparin, a highly charged substance related to blood clotting. Binding can alter the molecule’s activity, movement, or laboratory measurement.

\medterm{Heparin-Induced Thrombocytopenia} An immune reaction to heparin that lowers the platelet count while making dangerous blood clots more likely. It usually begins several days after heparin exposure. Suspected cases require stopping heparin and urgent medical treatment with a different anticoagulant.

\medterm{Hepatectomy} Surgery to remove part or, rarely, nearly all of the liver, usually to treat a tumor or severe localized liver disease. The remaining liver must be large and healthy enough to work. Bleeding, bile leakage, infection, and liver failure are important risks.

\medterm{Hepatic} Relating to the liver. The liver processes nutrients and medicines, makes bile and important blood proteins, stores energy, and helps remove harmful substances; hepatic disease can affect one or several of these functions.

\medterm{Hepatic Artery} The main artery supplying oxygen-rich blood to the liver. It divides into branches that deliver blood to the right and left parts of the organ.

\medterm{Hepatic Artery Occlusion} Partial or complete blockage of the hepatic artery, usually caused by a blood clot, embolus, inflammation, injury, or a procedure. It may reduce blood flow to the liver and cause liver infarction or ischemic hepatitis, especially when portal-vein blood flow is also reduced.

\medterm{Hepatic Biopsy} Alternate name; see \textit{Liver Biopsy}.

\medterm{Hepatic Coma} Severe loss of consciousness caused by advanced liver failure and buildup of toxic substances in the blood. It is a medical emergency, often preceded by confusion, sleepiness, personality change, or hand-flapping movements.

\medterm{Hepatic Cyst} A fluid-filled sac in the liver. Most simple liver cysts are harmless, cause no symptoms, and are found by chance on imaging. Large, infected, bleeding, parasitic, or complex cysts may cause pain, pressure, fever, or other complications.

\synonyms
Simple Liver Cyst; Biliary Cyst

\medterm{Hepatic Dose Adjustment} Changing a medicine or its dose because liver disease may slow its breakdown or removal. The correct adjustment depends on the medicine, liver function, other illnesses, and monitoring of benefit and toxicity.

\medterm{Hepatic Duct} The duct formed by the union of the right and left liver ducts, carrying bile from the liver toward the common bile duct and intestine.

\medterm{Hepatic Encephalopathy} Confusion or changes in alertness, behavior, movement, or consciousness caused by severe liver disease and buildup of substances that affect the brain. Symptoms can progress to coma and may be precipitated by factors such as infection, bleeding, constipation, or medicines.

\medterm{Hepatic Fissure} A natural groove or cleft on the liver's surface. It separates parts of the liver and may contain or provide a route for blood vessels, bile ducts, or supporting ligaments.

\medterm{Hepatic Flexure} The bend where the right-sided ascending colon turns into the transverse colon near the liver. It is a normal anatomical landmark that may be seen during imaging, surgery, or colonoscopy.

\medterm{Hepatic Granulomas} Small collections of immune cells that form in the liver in response to infections, medicines, sarcoidosis, or other inflammatory conditions. They often cause no symptoms, but the underlying condition may cause abnormal liver tests, bile-duct blockage, fibrosis, or portal hypertension.

\medterm{Hepatic Hemangioma} A common, usually harmless growth made from small blood vessels in the liver. Most cause no symptoms and are found by chance on imaging; larger ones may cause discomfort or rarely bleeding.

\medterm{Hepatic Hydrothorax} A buildup of fluid around a lung caused by severe liver disease and high pressure in the abdomen. It is usually on the right side and can cause breathlessness; treatment focuses on the liver disease and fluid control.

\medterm{Hepatic Ischemia} Liver injury caused by inadequate hepatic blood flow or oxygen delivery, often during
shock, severe hypoxemia, cardiac failure, or hepatic vascular occlusion. It can produce
a rapid, marked rise in aminotransferases and may progress to acute liver failure.

\medterm{Hepatic Lamina} A sheet of cells in a developing embryo that contributes to the liver and its bile ducts.

\medterm{Hepatic Necrosis} Death of liver cells caused by severe ischemia, toxins, viral infection, autoimmune injury, or other insults; extensive necrosis can produce acute liver failure.

\medterm{Hepatic Portal System} A network of veins that carries blood from the digestive organs, spleen, and pancreas to the liver before the blood returns to the heart. The liver processes nutrients, medicines, and harmful substances in this blood.

\medterm{Hepatic Steatosis} Excess accumulation of fat in liver cells, commonly called fatty liver; causes include metabolic disease, alcohol, medicines, and other conditions, and inflammation or fibrosis may develop in some people.

\medterm{Hepatic Surgery} Surgical treatment of liver disease, including partial liver resection, transplantation, and
procedures for tumors or biliary disorders. The operation depends on the diagnosis, tumor extent, liver function, and
overall health.

\medterm{Hepatic Synthetic Dysfunction} Reduced ability of the liver to make important proteins, including albumin and blood-clotting factors. It can cause swelling, easy bleeding, or abnormal clotting tests, although other problems such as vitamin K deficiency or blood-thinning medicines can also alter these tests.

\medterm{Hepatic Trauma} Liver injury caused by blunt or penetrating trauma, potentially producing internal bleeding and requiring urgent imaging, monitoring, or surgery based on severity.

\medterm{Hepatic Vein} One of the veins that carries blood away from the liver and into the large vein leading to the heart. The main hepatic veins drain different regions of the liver.

\medterm{Hepatic Vein Obstruction} Obstruction of hepatic venous outflow, usually from thrombosis or compression, causing abdominal pain, hepatomegaly, ascites, and potentially liver failure.

\medterm{Hepatitis} Hepatitis is inflammation of the liver caused by viruses, alcohol, medications, toxins, metabolic disease, or autoimmunity and may be acute or chronic.

\medterm{Hepatitis A} A short-term liver infection caused by hepatitis A virus. It spreads when virus from infected stool is swallowed, often through close contact or contaminated food or water. It may cause fever, tiredness, nausea, abdominal discomfort, or jaundice. Most people recover completely; chronic infection is not expected. Vaccination helps prevent it.

\medterm{Hepatitis A Virus} The virus that causes hepatitis A. It spreads when virus from infected stool is swallowed, commonly through contaminated food, water, or close contact. It causes a short-term liver infection that may produce fever, nausea, tiredness, dark urine, pale stools, or jaundice; vaccination helps prevent it.

\medterm{Hepatitis B} A liver infection caused by the hepatitis B virus (HBV). It spreads when infected blood, semen, or certain other body fluids enter another person’s body, including through sex, shared needles or equipment, or from an infected parent to a baby around the time of birth. Infection may be acute (short-term) or chronic (long-term). Many people have no symptoms. When symptoms occur, they may include tiredness, fever, loss of appetite, nausea, abdominal pain, joint pain, dark urine, pale stools, or yellow skin and eyes (jaundice). Chronic infection, especially when acquired in infancy or childhood, can cause liver damage, cirrhosis, liver failure, or liver cancer. Vaccination prevents hepatitis B.

\medterm{Hepatitis B Virus} The virus that causes hepatitis B. HBV spreads when infected blood, semen, or certain other body fluids enter another person’s body, including through sex, shared needles or equipment, or around childbirth. It can cause acute (short-term) or chronic (long-term) liver infection. Chronic infection can cause liver damage, cirrhosis, liver failure, or liver cancer. Vaccination prevents hepatitis B.

\medterm{Hepatitis C} A liver infection caused by the hepatitis C virus (HCV). It spreads mainly when infected blood enters another person’s body, especially through shared injection equipment; it can also spread through some other blood exposures and, less often, through sex or from an infected parent to a baby. Acute infection often causes no symptoms, and some people clear the virus without treatment. In many others it becomes chronic (long-term) and can cause liver damage, cirrhosis, liver failure, or liver cancer. Direct-acting antiviral medicines cure more than 95\% of infections in most people; no vaccine is available.

\medterm{Hepatitis C Virus} The virus that causes hepatitis C. HCV spreads mainly when infected blood enters another person’s body, especially through shared needles or other injection equipment. It can also spread through some medical exposures, from an infected parent to a baby, or less often through sex. Many people have no symptoms. The infection may be acute (short-term) or become chronic (long-term), which can cause liver damage, cirrhosis, liver failure, or liver cancer. Modern antiviral medicines cure more than 95\% of infections in most people, but no vaccine is available.

\medterm{Hepatitis D} A liver infection caused by the hepatitis D virus (HDV). HDV can infect only people who also have hepatitis B because it needs hepatitis B virus (HBV) to reproduce. Infection with both viruses at the same time is called coinfection; getting HDV after already having HBV is called superinfection. Superinfection is more likely to cause long-term liver disease, cirrhosis, liver failure, or death.

\medterm{Hepatitis D} Another name for hepatitis D, a liver infection caused by HDV. HDV needs HBV to reproduce. It may occur as coinfection with HBV or as superinfection in someone who already has HBV; superinfection is more likely to cause long-term liver disease, cirrhosis, liver failure, or death.

\medterm{Hepatitis D Virus} A virus that can multiply only when hepatitis B virus is also present. It spreads through infected blood or body fluids and can make hepatitis B more severe, increasing the risk of cirrhosis and liver cancer. Hepatitis B vaccination also prevents hepatitis D.

\medterm{Hepatitis E} A liver infection caused by the hepatitis E virus (HEV). It usually spreads when people swallow water or food contaminated with stool, and it can also spread through undercooked meat from infected animals. It usually causes a short-term illness, often with tiredness, nausea, vomiting, abdominal pain, dark urine, pale stools, or jaundice. Severe illness is more likely during pregnancy and in some people with weakened immune systems. Rarely, infection becomes chronic in people with weakened immunity. No vaccine is available in the United States.

\medterm{Hepatitis E Virus} The virus that causes hepatitis E. HEV usually spreads through water or food contaminated with stool and can also spread through undercooked meat from infected animals. It commonly causes short-term liver infection, but it may cause severe illness during pregnancy and can occasionally become chronic in people with weakened immune systems. No vaccine is available in the United States.

\medterm{Hepatitis Serology} Blood testing for antibodies, antigens, and other markers of hepatitis viruses. The results may show a current or past infection, immunity from vaccination, or whether infection has become chronic. Different viruses require different tests. For example, hepatitis B testing commonly includes hepatitis B surface antigen and antibodies to hepatitis B surface and core antigens; hepatitis C testing usually starts with an antibody test and uses an HCV RNA test to show current infection. Results depend on the virus, timing, vaccination history, and other laboratory findings, so a single result may not give the full answer.

\medterm{Hepatitis Virus} Any of the viruses that can infect and inflame the liver. The five main hepatitis viruses are hepatitis A, B, C, D, and E. They differ in how they spread and in whether they cause short-term or long-term infection. Hepatitis A usually causes a short-term illness; hepatitis B and C can become chronic and damage the liver; hepatitis D occurs only with hepatitis B; and hepatitis E usually causes short-term illness but can be severe during pregnancy or in people with weakened immunity. Vaccines prevent hepatitis A and B, and hepatitis B vaccination also prevents hepatitis D.

\medterm{Hepatitis Virus Panel} A group of blood tests, usually done on one blood sample, that looks for signs of hepatitis A, B, and C infection. The tests may detect viral antigens or antibodies and may show a current infection, a past infection, or immunity after vaccination. The usual panel does not automatically test for hepatitis D or E, which require separate tests when appropriate. A positive antibody result does not always mean that the virus is still present; additional tests may be needed to confirm current infection.

\medterm{Hepatitis Viruses} The five main viruses that infect and inflame the liver: hepatitis A, B, C, D, and E. Hepatitis A and E usually cause short-term infection, although hepatitis E can be severe during pregnancy and may become chronic in some people with weakened immunity. Hepatitis B and C can become chronic and cause cirrhosis or liver cancer. Hepatitis D occurs only in people who also have hepatitis B. Vaccines prevent hepatitis A and B, hepatitis B vaccination also prevents hepatitis D, and effective medicines can cure hepatitis C; no vaccine is available for hepatitis C or E in the United States.

\medterm{Hepatobiliary} Relating to the liver, gallbladder, and bile ducts. The hepatobiliary system makes bile in the liver, stores and concentrates it in the gallbladder, and carries it through the bile ducts to the small intestine, where it helps digest fats. It also includes the liver’s roles in processing nutrients, medicines, and waste products.

\medterm{Hepatobiliary Scintigraphy} A nuclear-medicine scan, often called a HIDA scan, that uses a small injected radioactive tracer to follow bile movement through the liver, bile ducts, gallbladder, and small intestine. A special camera records the tracer over time. The scan can show blocked bile flow, whether the gallbladder fills and empties normally, bile leaks, and findings that support acute gallbladder inflammation or other biliary disorders.

\synonyms
Biliary Scintigraphy; Cholescintigraphy; HIDA Scan; Hepatobiliary Iminodiacetic Acid Scan

\medterm{Hepatobiliary System} The liver, gallbladder, and bile ducts. The liver makes and secretes bile, the gallbladder stores and concentrates it, and the bile ducts carry it to the small intestine, where it helps digest fats and carry some waste products out of the liver. The liver also processes nutrients and many medicines and toxins.

\medterm{Hepatoblastoma} A rare cancer that begins in the liver and is the most common type of liver cancer in children. It usually affects children younger than 3 years. A child may have no symptoms until the tumor grows larger; possible findings include a lump or swelling in the abdomen, abdominal pain, poor appetite, nausea, vomiting, or unexplained weight loss. Evaluation may include imaging and blood tests, especially the tumor marker alpha-fetoprotein (AFP); a tissue sample may be used when needed to confirm the diagnosis.

\medterm{Hepatocellular Adenoma} A rare, usually noncancerous tumor made of mature liver cells. It occurs most often in women of reproductive age and is associated with estrogen exposure, although it can also occur with anabolic steroid use, certain inherited metabolic diseases, or without a known cause. It is often found by chance on imaging and may cause no symptoms. Larger tumors can bleed or rupture, and a small number can change into liver cancer, especially certain subtypes.

\medterm{Hepatocellular Carcinoma} A cancer that begins in liver cells called hepatocytes. It is the most common primary liver cancer in adults. It often develops in a liver affected by long-term disease, especially cirrhosis, chronic hepatitis B or C, fatty liver inflammation, or long-term heavy alcohol use, but it can occur without a known risk factor. Early disease may cause no symptoms. Later findings may include a lump or pain in the upper-right abdomen, abdominal swelling, jaundice, tiredness, poor appetite, or unexplained weight loss. Diagnosis may use imaging, liver blood tests, and the tumor marker alpha-fetoprotein (AFP); a tissue sample is sometimes needed.

\medterm{Hepatocellular Jaundice} Yellowing of the skin or eyes caused by damaged liver cells processing or releasing bilirubin poorly. Causes include hepatitis, alcohol, medicines, toxins, fatty liver disease, and other liver disorders.

\synonyms
Hepatocellular Jaundice; Hepatocellular Icterus

\medterm{Hepatocellular Jaundice} Jaundice caused by impaired liver cells. Injured hepatocytes cannot take up, change into a water-soluble form, or remove bilirubin normally, so bilirubin builds up in the blood and makes the skin and whites of the eyes yellow. Urine may become dark and stools may become pale; itching, tiredness, nausea, or abdominal discomfort may also occur. Causes include viral or autoimmune hepatitis, alcohol-related liver disease, medicines or toxins, cirrhosis, hemochromatosis, and Wilson disease.

\medterm{Hepatocerebral Degeneration} Usually referring to acquired hepatocerebral degeneration, a rare long-term brain disorder linked to advanced liver disease or a portosystemic shunt, an abnormal connection that lets blood bypass the liver. Harmful substances that are normally cleared by the liver, including excess manganese and other neurotoxins, can affect brain areas involved in movement and thinking. Symptoms may include tremor, slowed movement or parkinsonism, poor coordination, unsteady walking, speech changes, abnormal movements, mood or personality changes, and problems with thinking or memory. It is different from the short-lived episodes of hepatic encephalopathy, although the conditions can overlap.

\medterm{Hepatocyte} The main working cell of the liver. Hepatocytes process carbohydrates, fats, proteins, medicines, and other substances from the blood. They make bile, albumin, clotting factors, and other blood proteins; store energy as glycogen; change ammonia into urea; and help remove or break down harmful substances. Injury to hepatocytes can reduce these functions and raise liver enzymes in the blood.

\medterm{Hepatojugular Reflux} A bedside physical sign, also called abdominojugular reflux. Firm, steady pressure on the upper abdomen briefly increases the amount of blood returning to the heart. In a normal response, the jugular veins in the neck rise briefly and then fall while the pressure continues. A sustained rise in jugular venous pressure suggests elevated right-sided heart filling pressure or limited right-ventricle function, often in heart failure. It is a clinical finding, not a diagnosis by itself.

\medterm{Hepatolenticular Degeneration} Another name for Wilson disease, a rare inherited disorder caused by changes in the \textit{ATP7B} gene. The altered gene prevents the liver from removing excess copper into bile, so copper builds up to harmful levels in the liver, brain, eyes, kidneys, and other tissues. It may cause hepatitis, cirrhosis, or liver failure; tremor, stiffness, poor coordination, speech or mood changes; and green, gold, or brown Kayser--Fleischer rings around the cornea. The disorder is inherited in an autosomal recessive pattern.

\medterm{Hepatologist} A physician who specializes in diseases of the liver and its connected bile ducts and gallbladder. Hepatologists commonly diagnose and manage hepatitis, fatty liver disease, cirrhosis, jaundice, liver failure, liver tumors, and complications of chronic liver disease. Some also care for related diseases of the pancreas and people being assessed for liver transplantation. Hepatology is usually a subspecialty of gastroenterology.

\medterm{Hepatology} The branch of medicine concerned mainly with the liver and also with the gallbladder, bile ducts, and related pancreas disorders. It covers how these organs work, how their diseases are diagnosed and managed, and complications such as hepatitis, fatty liver disease, cirrhosis, jaundice, bile-flow disorders, liver cancer, and liver failure. Hepatology is usually a subspecialty of gastroenterology.

\medterm{Hepatoma} An older, nonspecific term for a liver tumor. In some contexts it is used to mean hepatocellular carcinoma, a cancer arising from liver cells, but the word alone does not show whether a tumor is benign or malignant. \textit{Hepatocellular carcinoma} is the preferred term when a malignant liver-cell tumor is meant.

\medterm{Hepatomegaly} An enlarged liver. It may be noticed when the liver is felt below the right ribs during an examination or seen on ultrasound, CT, or MRI. Hepatomegaly is a physical finding, not a disease itself. Causes include acute hepatitis, metabolic dysfunction-associated steatotic liver disease, alcohol-related liver disease, congestion from heart or blood-flow problems, bleeding into the liver, bile-duct blockage, storage or infiltrative diseases, and primary or metastatic cancer. It may cause no symptoms or may be associated with right-upper-abdominal discomfort.

\medterm{Hepatopulmonary Syndrome} A complication of liver disease or portal hypertension in which blood vessels inside the lungs become abnormally widened. Blood then passes through the lungs without receiving enough oxygen, causing low oxygen levels in the blood. Symptoms may include shortness of breath, bluish skin, or clubbing of the fingers. Breathing difficulty and low oxygen may become worse when sitting or standing and improve when lying down; this pattern is called platypnea and orthodeoxia. The condition is identified by liver disease or portal hypertension, impaired oxygenation, and evidence of widened vessels inside the lungs.

\medterm{Hepatorenal Syndrome} Serious kidney dysfunction that occurs mainly in advanced cirrhosis, usually with portal hypertension and ascites. Widening of blood vessels in the abdominal circulation lowers the amount of blood effectively reaching the kidneys. The kidneys respond by narrowing their blood vessels and retaining salt and water, so filtration falls even though the kidneys may not initially have a primary structural disease. Hepatorenal syndrome may cause acute kidney injury (HRS-AKI), which develops rapidly, or more persistent kidney dysfunction. Diagnosis requires excluding other important causes of kidney injury, such as shock, dehydration, medicines, infection, or intrinsic kidney disease.

\medterm{Hepatosplenomegaly} Enlargement of both the liver and spleen at the same time. It is a physical finding rather than a disease itself and may be found during an examination or on imaging. Causes include infections, chronic liver disease with portal hypertension, blood-cell destruction, blood cancers, inherited storage or metabolic disorders, and other infiltrative or inflammatory diseases. It may cause abdominal fullness or discomfort, or no symptoms; other findings depend on the underlying cause. Blood tests and imaging may help identify the cause.

\medterm{Hepatotoxicity} Toxic injury to the liver caused by a medicine, herbal product, dietary supplement, poison, chemical, or other exposure. The injury may mainly damage liver cells (hepatocellular injury), slow or block bile flow (cholestatic injury), or show both patterns. It may cause no symptoms but can lead to tiredness, nausea, poor appetite, abdominal pain, itching, dark urine, or jaundice, with abnormal liver blood tests. Most cases improve after the cause is removed, but severe injury can cause acute liver failure.

\medterm{Hepatovirus} A genus in the Picornaviridae family of small RNA viruses. Its best-known human member is hepatitis A virus (HAV), which causes hepatitis A and usually spreads by the fecal--oral route, when virus from stool reaches the mouth through contaminated food, water, hands, or close contact. Other hepatoviruses infect different animal hosts and are not established causes of human disease.

\medterm{Hepcidin} A peptide hormone made mainly by liver cells and the main regulator of iron balance. When iron stores are high or inflammation is present, hepcidin binds to ferroportin, the protein that exports iron from intestinal cells, macrophages, and liver cells. Ferroportin is then removed, reducing iron absorption from food and the release of recycled or stored iron into the blood. When iron is low or the bone marrow needs more iron to make red blood cells, hepcidin levels usually fall.

\medterm{Heptachlor} A persistent organochlorine insecticide formerly used mainly on crops and for termite control. It breaks down slowly into heptachlor epoxide, which can remain in soil and the food chain for a long time. Exposure can occur through contaminated soil, food, water, air, or contact with old treated materials. Studies, especially in animals, associate exposure with effects on the liver, nervous system, reproductive system, and development; evidence of health effects in humans is limited. Most uses have been canceled or severely restricted in the United States.

\medterm{HER2} Short for human epidermal growth factor receptor 2, a protein on the surface of cells that helps control normal cell growth. It is made from the \textit{ERBB2} gene. Some cancer cells have too much HER2 protein or extra copies of the \textit{ERBB2} gene; these cancers may grow more quickly and spread more easily. Testing HER2 protein levels or gene copies can help classify a cancer and identify whether HER2-targeted treatment may be useful.

\medterm{HER2-Positive} A cancer whose cells have increased expression or amplification of human epidermal growth factor receptor 2 (HER2). Some HER2-targeted medicines may be used according to the cancer type and clinical context.

\medterm{HER2-Positive Breast Cancer} Breast cancer cells with unusually high levels of the HER2 growth receptor or extra copies of its gene. This can make the cancer grow faster, but it also allows treatment with medicines that specifically block HER2.

\medterm{Herbal Remedies And Supplements For Weight Loss} Plant-based weight-loss products have variable evidence, uncertain composition, and potential stimulant, liver, cardiovascular, or drug-interaction risks; they should not replace evidence-based obesity treatment.

\medterm{Herbal Supplement} A concentrated plant-derived product marketed for health or symptom relief; composition and evidence vary, and contamination, drug interactions, liver injury, and pregnancy risks must be considered.

\medterm{Herbaspirillum} A genus of gram-negative bacteria found in soil and plants; human infection is rare but can occur opportunistically.

\medterm{Herbicide} A chemical used to kill or control unwanted plants; exposure can irritate skin, eyes, lungs, or the gastrointestinal tract, with toxicity determined by the specific compound and dose.

\medterm{Herd Immunity} Protection that occurs when enough people in a community are immune to an infection that it spreads less easily. The proportion needed varies widely by infection, and herd immunity does not protect everyone equally.

\medterm{Hereditary} Transmitted through genetic material from a parent or arising from an inherited variant; a hereditary trait may be dominant, recessive, X-linked, mitochondrial, or multifactorial.

\medterm{Hereditary Amyloidosis} A group of inherited disorders in which mutations cause misfolded proteins to form amyloid deposits in organs such as the nerves, heart, kidneys, or gastrointestinal tract.

\medterm{Hereditary Angioedema} An inherited, usually bradykinin-mediated disorder causing repeated episodes of swelling in the skin, hands, feet, stomach or bowel wall, or throat. Types 1 and 2 commonly result from deficiency or dysfunction of C1 inhibitor, often due to variants in the \textit{SERPING1} gene; other forms may have normal C1-inhibitor levels. Attacks may cause severe abdominal pain or swelling that narrows the upper airway and often occur without itchy hives.

\synonyms
HAE; Hereditary Angioneurotic Edema; HANE; C1-INH Deficiency Angioedema

\medterm{Hereditary Cancer} Cancer occurring because a person inherited a disease-causing genetic variant that raises cancer risk. Genetic counseling and appropriate risk-based screening may be recommended.

\medterm{Hereditary Coproporphyria} An inherited disorder of a chemical pathway used to make heme, the oxygen-carrying part of hemoglobin. Attacks may cause severe abdominal pain, vomiting, nerve or mental symptoms, and sometimes blistering skin after sunlight exposure.

\medterm{Hereditary Diffuse Gastric Cancer} An inherited condition, often caused by a harmful CDH1 gene change, that greatly increases the risk of a hard-to-detect stomach cancer and also raises the risk of lobular breast cancer. Genetic counseling, specialist screening, and discussion of risk-reducing stomach surgery may be recommended.

\synonyms
HDGC; Familial Gastric Cancer Syndrome

\medterm{Hereditary Elliptocytosis} An inherited red-cell membrane disorder in which erythrocytes are elliptical on the blood smear; most people are asymptomatic, but hemolytic anemia can occur in severe forms.

\medterm{Hereditary Fructose Intolerance} An inherited disorder in which the body cannot properly break down fructose, a sugar in fruit, honey, and table sugar. Eating it can cause vomiting, sweating, low blood glucose, abdominal pain, or liver injury.

\medterm{Hereditary Hemorrhagic Telangiectasia} An inherited disorder in which fragile, abnormal blood vessels form in the nose, skin, lungs, brain, or liver. Recurrent nosebleeds and anemia are common; bleeding or abnormal vessels in internal organs may be serious.

\synonyms
HHT; Osler-Weber-Rendu Syndrome

\medterm{Hereditary Hemorrhagic Telangiectasia} An inherited blood-vessel disorder causing fragile vessels that can bleed, especially in the nose and digestive tract. Abnormal vessels may also form in the lungs, brain, or liver.

\synonyms
HHT; Osler-Weber-Rendu Syndrome

\medterm{Hereditary Hyperbilirubinemia} An inherited disorder of bilirubin production, conjugation, or excretion that causes persistent or intermittent jaundice, such as Gilbert, Crigler--Najjar, Dubin--Johnson, or Rotor syndrome.

\medterm{Hereditary Leiomyomatosis And Renal Cell Cancer} An inherited condition that can cause painful skin muscle tumors, early or multiple uterine fibroids, and an aggressive type of kidney cancer. It is usually related to an FH gene change; genetic counseling and regular specialist skin, uterine, and kidney assessment are important.

\synonyms
HLRCC; Reed Syndrome

\medterm{Hereditary Multiple Exostoses} An inherited disorder in which many benign, cartilage-capped bone growths called osteochondromas form near the ends of growing bones. They usually develop during childhood or adolescence and stop forming when skeletal growth ends. They may cause pain, limited joint movement, unequal limb lengths, short stature, bone deformity, or pressure on nearby nerves and blood vessels. The condition is usually caused by a change in the \textit{EXT1} or \textit{EXT2} gene and follows an autosomal dominant inheritance pattern. Most growths remain benign, but a small number can change into a cancerous cartilage tumor. Also called hereditary multiple osteochondromas.

\synonyms
Hereditary Multiple Osteochondromas; Diaphyseal Aclasis

\medterm{Hereditary Ovalocytosis} An inherited red-cell membrane disorder in which many erythrocytes are oval or elliptical; it is often mild but can cause hemolytic anemia in some individuals.

\medterm{Hereditary Paraganglioma-Pheochromocytoma Syndrome} An inherited condition that increases the risk of tumors arising from nerve-related cells near blood vessels or in the adrenal glands. Some tumors release hormones that cause high blood pressure, headache, sweating, or palpitations; genetic counseling and lifelong imaging and hormone monitoring may be needed.

\synonyms
Hereditary Paraganglioma; SDH-Associated Pheochromocytoma

\medterm{Hereditary Persistence of Fetal Hemoglobin} An inherited condition in which production of fetal hemoglobin continues into adulthood. It is usually harmless and may lessen the severity of some beta-globin disorders such as sickle-cell disease.

\medterm{Hereditary Spastic Paraplegia} A group of inherited neurologic disorders causing progressive stiffness and weakness of the legs because of corticospinal-tract degeneration. Some forms also affect speech, vision, sensation, coordination, or other organs.

\medterm{Hereditary Spherocytic Anemia} An inherited red-cell membrane disorder, usually hereditary spherocytosis, in which fragile spherical erythrocytes are destroyed prematurely, causing hemolytic anemia, jaundice, and splenomegaly.

\medterm{Hereditary Spherocytosis} An inherited blood disorder in which red blood cells are round rather than disc-shaped and are broken down too quickly, mainly in the spleen. It may cause anemia, jaundice, gallstones, and an enlarged spleen.

\medterm{Hereditary Urea Cycle Abnormality} An inherited problem removing ammonia, a waste product, from the blood. Ammonia can build up and cause vomiting, sleepiness, confusion, seizures, coma, and brain injury, especially during illness or fasting.

\medterm{Heredity} The passing of genetic characteristics from parents to children through genes and chromosomes. Heredity influences appearance, body function, disease risk, medicine response, and susceptibility to some inherited conditions.

\medterm{Heredo-Ataxia} An inherited disorder causing gradually worsening problems with balance, walking, coordination, or speech. Different gene changes cause different forms, so diagnosis usually includes neurologic assessment and genetic testing.

\medterm{Heredopolyneuritis} An inherited disorder affecting many peripheral nerves. It may cause gradually worsening weakness, muscle wasting, numbness, reduced reflexes, or difficulty walking, often beginning in the feet and lower legs.

\medterm{Hering-Breuer Reflex} An automatic lung response that slows or stops breathing in when the lungs are stretched too much. It helps protect the lungs from overinflation, especially during large breaths.

\medterm{Hermanski-Pudlak Syndrome} Alternate spelling; see \textit{Hermansky-Pudlak Syndrome}.

\medterm{Hermansky-Pudlak Syndrome} A group of inherited disorders that can cause reduced skin, hair, and eye pigment; easy bleeding because platelets do not work normally; and, in some types, lung scarring, bowel inflammation, or low white-cell counts. Genetic and specialized blood testing help confirm the diagnosis.

\medterm{Hermaphroditism} An outdated and potentially stigmatizing term historically used for differences in sex development. Modern clinical language generally uses intersex or a specific difference of sex development, based on the person's anatomy, chromosomes, hormones, and lived identity.

\medterm{Hernia} A bulge formed when an organ or tissue pushes through a weak area in the muscle or connective tissue that normally holds it in place. Hernias commonly occur in the abdominal wall and may become trapped or lose their blood supply, causing severe pain.

\medterm{Hernia Mesh} A sheet of synthetic or biologic material used to strengthen the body wall during some hernia repairs. It may reduce recurrence, but complications can include infection, persistent pain, movement of the mesh, or another hernia.

\medterm{Hernia Repair} Surgery that returns the displaced tissue to its proper place and closes or strengthens the weak area in the body wall. Repair may use stitches, mesh, or a minimally invasive approach; the choice depends on the hernia and the person’s health.

\synonyms
Herniorrhaphy; Hernioplasty

\medterm{Herniated Disc} A spinal disk that bulges or tears so that its inner material irritates or presses on a nearby nerve. It can cause back or neck pain, pain spreading into an arm or leg, numbness, tingling, or weakness.

\medterm{Herniated Disk} Displacement of disk material through a tear in the outer ring of an intervertebral disc, most commonly in the lumbar spine (often at L4-L5 and L5-S1) or the cervical spine, which may press on a nearby nerve and cause pain, numbness, tingling, or weakness.

\medterm{Herniation} Protrusion of an organ, tissue, or structure through an abnormal opening or weak area, as in a herniated disk or abdominal-wall hernia.

\medterm{Herniorrhaphy} An operation to repair a hernia, in which tissue or an organ has pushed through a weak area in the body wall. The surgeon returns the tissue and closes or reinforces the weakness, sometimes with mesh; recurrence, pain, infection, or injury to nearby structures are possible.

\medterm{Heroin} Heroin is a highly addictive opioid converted to morphine in the body; overdose causes respiratory depression, coma, and death, while repeated use causes dependence and withdrawal.

\medterm{Heroin Dependence} Legacy, drug-specific label; see \textit{Opioid Use Disorder}.

\medterm{Heroin Overdose} A life-threatening reaction to too much heroin, causing extreme sleepiness, pinpoint pupils, slow or stopped breathing, blue lips, coma, or death. Call emergency services and give naloxone if available; repeat doses may be needed.

\medterm{Herpangina} A common childhood viral illness causing fever, sore throat, and small blisters or ulcers at the back of the mouth. It spreads through respiratory or stool-contaminated secretions and usually improves on its own.

\medterm{Herpes} Infection caused by a herpesvirus. The term commonly refers to herpes simplex infection, which causes oral or
genital lesions, but other herpesviruses cause distinct diseases.

\medterm{Herpes Encephalitis} Inflammation of the brain caused by herpes simplex virus infection.

\medterm{Herpes Gestationis} An older name for pemphigoid gestationis, an uncommon autoimmune blistering rash of pregnancy. It usually causes intensely itchy hives followed by blisters around the abdomen and may recur in later pregnancies.

\medterm{Herpes Labialis} Alternate name; see \textit{Cold Sore}.

\medterm{Herpes Simplex} Infection with herpes simplex virus type 1 (HSV-1) or type 2 (HSV-2). HSV-1 most often causes oral herpes or cold sores, while HSV-2 most often causes genital herpes, although either type can infect either site. Many infections cause no symptoms; when symptoms occur, painful blisters or ulcers may appear and recur. The virus remains inactive in nerve cells and can reactivate later. It spreads mainly through skin-to-skin or oral contact, including when no sore is visible, and rarely passes from a parent to a baby during delivery. Rare complications include eye infection, brain infection, or severe infection in a newborn or a person with weakened immunity.

\medterm{Herpes Simplex Encephalitis} A life-threatening brain infection caused most often by herpes simplex virus type 1. It can cause fever, headache, confusion, seizures, behavior changes, or loss of consciousness and needs urgent antiviral treatment.

\medterm{Herpes Simplex Infection} A viral infection caused by herpes simplex virus type 1 (HSV-1) or type 2 (HSV-2). HSV-1 usually causes oral herpes, while HSV-2 usually causes genital herpes, but either type can affect the mouth or genitals. Many infections are symptom-free; painful blisters or ulcers may develop and return because the virus remains inactive in nerve cells and can reactivate. HSV spreads mainly through oral or skin-to-skin contact, including contact with skin that looks normal. Rare complications include infection of the eye or brain and severe infection in newborns or people with weakened immunity.

\synonyms
HSV Infection; Cold Sores and Genital Herpes

\medterm{Herpes Simplex Keratitis} Infection or inflammation of the cornea caused by herpes simplex virus, usually producing eye pain, redness, tearing, light sensitivity, or blurred vision. Recurrent disease can scar the cornea and impair sight.

\medterm{Herpes Simplex Skin Infection} A herpes simplex infection of the skin or moist body linings that causes grouped blisters or shallow sores, often around the mouth or genitals. Widespread infection in someone with eczema can be severe.

\medterm{Herpes Simplex Type 1} HSV-1, a DNA virus commonly causing oral or perioral blisters; it can also cause genital infection and remains latent with possible recurrence.

\medterm{Herpes Simplex Virus} A double-stranded DNA virus belonging to the Herpesviridae family that causes oral and
genital herpes. HSV has two serotypes: HSV-1, which typically causes oral herpes, and
HSV-2, which typically causes genital herpes. The virus establishes latent infections in
sensory ganglia and can reactivate periodically.

\medterm{Herpes Simplex Virus Infection} Infection with herpes simplex virus type 1 or type 2, commonly causing recurrent oral,
genital, skin, eye, or central nervous system disease.

\medterm{Herpes Vegetans} A rare long-lasting herpes simplex infection that causes thick, wart-like or ulcerated growths, usually on genital or anal skin. It occurs mainly in people with weakened immunity and may be difficult to treat.

\medterm{Herpes Viral Culture Of Lesion} A laboratory culture of fluid or tissue from a lesion to detect herpes simplex virus or varicella-zoster virus; molecular testing is generally more sensitive and is often preferred.

\medterm{Herpes Zoster} A painful, usually one-sided blistering rash caused by reactivation of latent varicella-zoster virus in a sensory nerve
ganglion. It typically follows a band-like distribution and may cause persistent nerve
pain after the rash resolves.

\synonyms
Shingles; Zoster; Acute Posterior Ganglionitis

\medterm{Herpes Zoster} A painful, vesicular rash caused by reactivation of latent varicella-zoster virus (VZV,
human herpesvirus 3) in the dorsal root ganglia.

\synonyms
Shingles; Acute Zoster Ganglionitis

\medterm{Herpes Zoster Ophthalmicus} Shingles affecting the forehead, eyelid, or nose near the eye. It can inflame the eye and threaten vision, so a painful blistering rash in this area needs same-day medical assessment and treatment.

\medterm{Herpesviridae} A family of DNA viruses that can remain inactive in the body after infection and reactivate later. Human members include herpes simplex viruses, varicella-zoster virus, cytomegalovirus, Epstein--Barr virus, and human herpesviruses 6, 7, and 8. They can cause skin or mouth lesions, chickenpox, shingles, glandular fever, congenital infection, or disease in people with weakened immunity.

\medterm{Herpesviridae Infection} An infection caused by a herpesvirus; symptoms range from localized lesions to serious disease in newborns or people with weakened immunity.

\medterm{Herpetic Gingivostomatitis} The classic and most common manifestation of primary herpes simplex virus type 1 (HSV-1) infection in infants and young children aged 6 months to 5 years. It is characterized by high fever, irritability, painful submandibular and cervical lymphadenopathy, and widespread friable, inflamed gums with clusters of painful vesicles that rupture into shallow, coalescing ulcers across the buccal mucosa, tongue, and lips, often severely impairing oral intake.

\synonyms
Primary Herpetic Gingivostomatitis; Oral HSV-1 Stomatitis

\medterm{Herpetic Glossitis} Painful inflammation and sores of the tongue caused by herpes simplex virus. It may occur with blisters or ulcers elsewhere in the mouth and can be more severe or persistent in people with weakened immunity.

\medterm{Herpetic Keratitis} An infection of the cornea caused by herpes simplex virus (predominantly HSV-1), recognized globally as the leading infectious cause of corneal blindness. It classically presents with unilateral eye pain, photophobia, foreign-body sensation, tearing, and pathognomonic branching dendritic ulcers with terminal bulbs visualized under fluorescein dye examination; recurrent episodes predispose to deep stromal keratitis, corneal scarring, and permanent vision loss.

\synonyms
Dendritic Keratitis; Herpes Simplex Corneal Infection

\medterm{Herpetic Stomatitis} Inflammation and ulceration of the mouth caused by herpes simplex virus, most commonly primary HSV-1 infection, with painful vesicles or ulcers and sometimes fever or gingivitis.

\medterm{Herpetic Whitlow} A painful herpes simplex infection of a finger, causing redness, swelling, and grouped blisters. It can spread by contact and should not be cut open because this may spread infection.

\medterm{Herpetiform} Resembling herpes lesions, typically describing grouped vesicles or an arrangement of lesions; the term does not by itself establish herpesvirus infection.

\medterm{Hertz} A unit used to measure frequency, equal to one complete cycle or repetition per second. It is used for sound, electrical signals, brain waves, heart signals, and other repeating events.

\medterm{Herxheimer Reaction} A temporary worsening of fever, chills, headache, muscle aches, or other symptoms soon after antibiotics begin for some infections, especially syphilis. It usually settles within a day and is caused by the body’s response to dying bacteria.

\medterm{Hesitation Wounds} Multiple superficial, tentative parallel incised cuts surrounding a deeper primary fatal
cut wound, classically observed in suicides involving sharp instruments (wrists, neck).

\medterm{Hesselbach Triangle} A naturally weaker area of the lower front abdominal wall. Direct groin hernias can push through this area, which lies between a central abdominal muscle, blood vessels, and the groin ligament.

\medterm{Heterochromatin} DNA that is tightly packed inside a cell’s nucleus and is usually less active in making proteins. It helps organize chromosomes and control which genes are available for use.

\medterm{Heterochromia} A difference in iris color between the eyes or within one iris, which may be congenital or acquired from disease, trauma, medication, or inflammation.

\medterm{Heterocyclic Amine} A nitrogen-containing chemical compound formed mainly when meat or fish is cooked at high temperatures; several heterocyclic amines are mutagens and potential dietary carcinogens.

\medterm{Heterogeneity} Variation in composition, structure, or behavior within a tissue, sample, lesion, or population; in tumors, heterogeneity can affect biopsy interpretation and treatment response.

\medterm{Heterogeneous Echotexture} An ultrasound appearance in which different parts of a tissue look unlike one another. It is a finding, not a diagnosis, and may result from inflammation, scarring, fat, bleeding, a tumor, or normal variation.

\medterm{Heterologous} Coming from a different species, source, or type. In medicine, the term may describe a graft, tissue, antibody, hormone, or genetic material obtained from a different organism or cell type.

\medterm{Heterophil Antibody} A heterophil antibody reacts with antigens from unrelated animal species; testing for these antibodies can help diagnose some infections, including infectious mononucleosis.

\medterm{Heterophile Antibody} A heterophile antibody is an antibody that reacts with antigens from unrelated species and can interfere with some laboratory immunoassays.

\medterm{Heterophyes Heterophyes} A small intestinal fluke acquired by eating raw or undercooked fish; it can cause abdominal symptoms and occasionally tissue complications.

\medterm{Heteroplasia} Development of tissue in an abnormal location or of a type not normally found there. The effects depend on the tissue involved and may range from no symptoms to obstruction, pain, or impaired function.

\medterm{Heteroplasmy} The presence of more than one type of mitochondrial DNA in a cell or person, usually a mixture of normal and changed copies. The proportion of changed copies can influence whether disease develops and how severe it is.

\medterm{Heterosexual} Describing a person who experiences sexual or romantic attraction mainly toward people of a different sex or gender. Sexual orientation is distinct from gender identity and is not a disease or disorder.

\medterm{Heterotaxy Syndrome} A birth condition in which organs in the chest and abdomen are arranged differently from usual. Heart defects, abnormal spleen function, and problems with the intestines or other organs may occur.

\medterm{Heterotopic} Located in an abnormal or unusual place; a heterotopic pregnancy is simultaneous implantation outside the uterus and, usually, inside the uterus.

\medterm{Heterotopic Breast Tissue Of The Vulva} Ectopic mammary-like tissue in the vulva that may undergo benign, hormonal, inflammatory, or malignant changes similar to breast tissue; diagnosis is histologic.

\medterm{Heterotopic Pregnancy} A pregnancy in which one embryo implants inside the uterus while another implants outside it, usually in a fallopian tube. It is uncommon but more likely after fertility treatment; an intrauterine pregnancy does not rule out an ectopic pregnancy.

\medterm{Heterotransplant} A transplant of cells, tissue, or an organ between members of different species. Because the immune difference is substantial, rejection is a major concern and long-term use requires specialized treatment.

\medterm{Heterozygote} A heterozygote has two different alleles at a genetic locus, one inherited from each parent; phenotype depends on dominance, the variant's effect, and the surrounding genetic and environmental context.

\medterm{Heterozygous} Having two different alleles of a particular gene at a genetic locus. The observable effect depends on the gene and the relationship between the alleles.

\medterm{Hexachlorobenzene} A persistent organochlorine compound that can cause porphyria, liver injury, neurologic effects, and other toxicity after significant exposure.

\medterm{Hexachlorophane} An antibacterial chemical formerly used in some skin cleansers. Repeated or extensive exposure can damage the nervous system, especially in infants, so its medical use is highly restricted.

\medterm{Hexachlorophene} A topical antiseptic phenolic compound whose excessive or prolonged use can cause neurotoxicity and skin injury; many uses are restricted.

\medterm{Hexadactyly} The presence of six fingers or six toes on a hand or foot, also called polydactyly.

\medterm{Hexadecadienoic Acid} A long-chain fatty acid containing 16 carbon atoms and two double bonds, occurring in lipid metabolism and studied as a biochemical or signaling molecule.

\medterm{Hexadecatetraenoic Acid} A 16-carbon polyunsaturated fatty acid containing four double bonds, present in lipid pathways and investigated for effects on membranes and signaling.

\medterm{Hexamethonium} An older medicine that blocks nerve signals at autonomic ganglia and can lower blood pressure. It is rarely used because it also causes blurred vision, dry mouth, constipation, urinary retention, and other widespread effects.

\medterm{Hexamine} A chemical used in industry and, in some settings, as methenamine, a urinary antiseptic. It releases formaldehyde in acidic urine and is unsuitable for some kidney or liver conditions.

\medterm{Hexane} A solvent found in some industrial products and glue. Long-term inhalation can damage peripheral nerves, causing numbness, weakness, and difficulty walking; swallowing or heavy exposure can also injure the lungs and brain.

\medterm{Hexanoic Acid} A six-carbon saturated fatty acid, also called caproic acid, found in some fats and odors and used in chemical synthesis; concentrated exposure irritates skin and mucosa.

\medterm{Hexokinase} An enzyme that begins the use of glucose by adding a phosphate group to it inside cells. This traps glucose in the cell and allows it to enter energy-producing or storage pathways. Most tissues use hexokinase; the liver also uses a related enzyme called glucokinase.

\medterm{Hexose} A six-carbon monosaccharide, such as glucose, fructose, or galactose, that serves as an energy substrate or metabolic intermediate.

\medterm{Hexylresorcinol} A phenolic antiseptic and local anesthetic used in some throat lozenges and topical products; excessive use or allergy can cause adverse effects.

\medterm{Hg} The chemical symbol for mercury; in measurements such as mmHg, it refers to millimeters of mercury, a unit of pressure.

\medterm{Hiatal} Relating to a hiatus, an opening or passage in an anatomical structure; a hiatal hernia occurs when abdominal tissue passes through the diaphragm's esophageal hiatus.

\medterm{Hiatal Hernia} Protrusion of part of the stomach through the esophageal opening in the diaphragm into the
chest. It may be sliding or paraesophageal and can be associated with reflux or other symptoms.

\medterm{Hiatus} A natural opening or passage in a body structure. For example, the diaphragm has an opening through which the esophagus passes; tissue can sometimes bulge through such an opening.

\medterm{Hiatus Hernia} Another term for hiatal hernia: protrusion of part of the stomach through the esophageal hiatus
of the diaphragm into the chest. It may be sliding or paraesophageal.

\medterm{Hib} Haemophilus influenzae type b, a bacterium that can cause serious infections such as meningitis, pneumonia, bloodstream infection, and swelling of the throat. Routine vaccination greatly reduces these infections in children.

\medterm{Hibernoma} A rare benign tumor of brown fat, usually a slow-growing painless soft-tissue mass that may feel warm because of its high metabolic and vascular activity.

\medterm{Hiccough} Alternate spelling; see \textit{Hiccups}.

\medterm{Hiccough} Alternate spelling; see \textit{Hiccups}.

\medterm{Hiccup} Singular form; see \textit{Hiccups}.

\medterm{Hiccups} Involuntary, spasmodic contractions of the diaphragm followed by sudden closure of the glottis,
producing a characteristic sound. Persistent or intractable hiccups may reflect
irritation of the diaphragm or phrenic and vagal pathways, metabolic disease, medication
exposure, or central nervous system disease.

\medterm{Hidradenitis} A chronic inflammatory skin disorder causing repeated painful lumps, abscesses, draining tunnels, and scars, usually in the armpits, groin, buttocks, or under the breasts.

\medterm{Hidradenitis Suppurativa} A long-lasting skin disease causing repeated painful lumps, pockets of pus, draining tunnels, and scars, usually in the armpits, groin, buttocks, or under the breasts. It is not caused by poor hygiene and is not contagious.

\medterm{Hidradenocarcinoma} A rare malignant sweat-gland tumor that may arise de novo or from a hidradenoma, presenting as a firm nodule and capable of local recurrence and metastasis.

\medterm{Hidradenoma} A usually harmless tumor arising from sweat-gland cells, typically appearing as a single slow-growing skin lump. It may be flesh-colored, red, cyst-like, or ulcerated, so examination or removal may be needed to distinguish it from other tumors.

\medterm{Hidradenoma Papilliferum} A benign papillary tumor of apocrine or mammary-like glands, usually presenting as a solitary flesh-colored or red nodule in the vulvar or anogenital region.

\medterm{Hidrocystoma} A benign cyst of a sweat gland near the eyelid or skin, appearing as a small translucent or bluish papule that may enlarge with heat or sweating.

\medterm{Hidrosis} Sweating or perspiration; abnormal hidrosis may mean excessive sweating, reduced sweating, or sweating in an unusual pattern.

\medterm{High Altitude} Elevation high enough that reduced atmospheric pressure lowers available oxygen, commonly beginning around 2,500 meters and increasing the risk of altitude illness.

\medterm{High Altitude Illness} Health problems caused by traveling to high elevation before the body adapts to lower oxygen levels. Symptoms range from headache and nausea to dangerous brain or lung swelling requiring immediate descent.

\medterm{High Arch} An unusually high curve along the inside of the foot. It may run in families or result from a nerve or muscle disorder and can cause pain, calluses, ankle instability, or difficulty finding comfortable shoes.

\synonyms
Pes Cavus; High-Arched Foot; Cavus Foot

\medterm{High Blood Cholesterol Levels} Too much cholesterol in the blood, especially LDL cholesterol, which can build up in artery walls and increase the risk of heart attack and stroke. It usually causes no symptoms and is found with a blood test.

\medterm{High Blood Potassium} Elevated potassium concentration in the bloodstream; see \textit{Hyperkalemia}.

\medterm{High Blood Pressure} Persistently elevated arterial blood pressure. Diagnosis generally requires properly obtained readings on more than one occasion or other appropriate confirmation. It often causes no symptoms but increases the risk of cardiovascular, cerebrovascular, kidney, and eye disease.

\medterm{High Blood Pressure And Eye Disease} Long-standing high blood pressure can damage retinal blood vessels and the optic nerve, causing hypertensive retinopathy or vision loss.

\medterm{High Blood Pressure Medications} Medicines used to lower arterial blood pressure, including diuretics, ACE inhibitors, ARBs, calcium-channel blockers, beta blockers, and other agents selected according to comorbidities and risk.

\medterm{High Blood Protein} An abnormally increased concentration of proteins in the blood, which may result from dehydration, inflammation, infection, or a plasma-cell disorder.

\medterm{High Blood Sugar} Hyperglycemia, an abnormally elevated concentration of glucose in the blood; persistent hyperglycemia is commonly associated with diabetes but can have other causes.

\medterm{High Cholesterol} An abnormally high level of cholesterol in the blood, especially low-density lipoprotein cholesterol. It often causes no symptoms but can contribute to plaque buildup in the arteries and cardiovascular disease.

\medterm{High Hemoglobin Count} A hemoglobin level above the expected range, which may reflect dehydration, chronic low oxygen, smoking, medication effects, or a blood disorder.

\medterm{High Output Failure} High-output heart failure, in which cardiac output is elevated or normal but inadequate for the body's unusually high metabolic or circulatory demand. Causes include severe anemia, thyrotoxicosis, beriberi, and large arteriovenous fistulas; treatment targets the cause and may include heart-failure therapy.

\medterm{High Potassium Level} Hyperkalemia is an elevated blood potassium concentration that can cause dangerous cardiac conduction abnormalities; confirmation, medication review, kidney assessment, ECG, and urgent treatment depend on level and symptoms.

\medterm{High Protein Diet} An eating pattern providing more protein than a person's usual or estimated requirement; suitability depends on nutritional needs, kidney and liver function, and the overall diet.

\medterm{High Red Blood Cell Count} An increased number or concentration of red blood cells in the blood, which may thicken the blood and result from dehydration, low oxygen, smoking, medications, or a bone marrow disorder.

\synonyms
Erythrocytosis; Polycythemia

\medterm{High Uric Acid Level} An increased concentration of uric acid in the blood, which may contribute to gout or kidney stones and can result from reduced kidney clearance, diet, medicines, or metabolic disease.

\medterm{High White Blood Cell Count} An increased number of white blood cells, often associated with infection, inflammation, stress, medicines, or a blood or bone marrow disorder.

\medterm{High White Blood Cell Count In Pregnancy} A mild rise in white blood cells can be normal during pregnancy, labor, or stress. Marked or persistent elevation with fever, pain, urinary or respiratory symptoms, or abnormal differential counts may indicate infection, inflammation, medication effect, or another disorder.

\medterm{High-Altitude Pulmonary Edema} Life-threatening fluid buildup in the lungs after rapid travel to high altitude. It causes breathlessness, cough, weakness, and sometimes pink or frothy sputum; immediate descent and emergency treatment are needed.

\medterm{High-Density Lipoprotein} A particle that carries cholesterol through the blood. Often called “HDL cholesterol,” its level is one part of heart-risk assessment and does not by itself determine whether someone will have heart disease.

\medterm{High-Dependency Unit} A hospital area for people who need closer observation and more treatment than a general ward provides but do not need the full support of an intensive-care unit. Staff monitor vital signs frequently and provide specialized care.

\medterm{High-Fiber Foods} Foods rich in soluble or insoluble dietary fiber, including legumes, whole grains, vegetables, fruit, nuts, and seeds, support bowel regularity and may improve cholesterol and glucose control when increased with adequate fluid.

\medterm{High-Flow Nasal Cannula} A device that delivers warmed, moist oxygen through small tubes in the nose at a high flow rate. It can improve breathing and oxygen levels in people with serious lung illness and is also used after removal of a breathing tube.

\medterm{High-Fructose Corn Syrup} A liquid sweetener containing varying proportions of glucose and fructose, used in many processed foods and drinks. Like other added sugars, high intake increases total energy exposure and is associated with dental caries and cardiometabolic risk; it is not inherently unique in effect compared with equivalent added sugars.

\medterm{High-Grade Endometrial Stromal Sarcoma} An aggressive cancer arising from supporting tissue in the uterus. It can grow into the uterine wall, blood vessels, or distant organs and requires specialist treatment based on its extent and cell features.

\medterm{High-Grade Serous Ovarian Carcinoma} An aggressive cancer involving the ovary, fallopian tube, or nearby lining of the abdomen. It often spreads before causing clear symptoms and may cause bloating, abdominal pain, or changes in bowel or bladder habits.

\medterm{High-Grade Surface Osteosarcoma} An aggressive bone cancer that starts on the outer surface of a bone. It may cause pain or swelling and can spread to the lungs or other sites, so treatment usually requires specialist cancer care.

\medterm{High-Intensity Focused Ultrasound} A treatment that concentrates sound-wave energy to heat and destroy a selected area of tissue while limiting damage nearby. Uses include some tumors and other conditions, depending on the device and treatment site.

\medterm{High-Level Disinfection} A cleaning process that kills nearly all microorganisms but may not destroy large numbers of bacterial spores. It is used for some reusable medical equipment that touches moist body surfaces and cannot be heat-sterilized.

\medterm{Highly Pathogenic Avian Influenza} Highly pathogenic influenza A infection in birds that can cause severe disease and mortality in poultry; selected strains can rarely infect humans after animal exposure.

\medterm{High-Pitched Cry} An unusually shrill infant cry that can occur with pain, neurologic dysfunction, drug withdrawal, or serious illness and warrants assessment when persistent or accompanied by abnormal behavior or feeding.

\medterm{High-Resolution Computed Tomography} A CT scan that uses very thin images and special settings to show fine detail, especially in the lungs. It can help detect scarring, inflammation, widened airways, and small-airway disease, but it uses X-rays.

\synonyms
HRCT; High-Resolution Chest CT

\medterm{High-Sensitivity Cardiac Troponin} A blood test that detects very small amounts of proteins released when heart muscle is injured. A raised result can occur with a heart attack or many other illnesses, so clinicians interpret repeated results with symptoms, ECG findings, and other tests.

\synonyms
hs-cTn; High-Sensitivity Troponin Assay

\medterm{Hilar} Relating to a hilum, the gateway where vessels, nerves, ducts, or bronchi enter or leave an organ, especially the lung or kidney.

\medterm{Hilar Cholangiocarcinoma} A cancer at the place where the liver’s main bile ducts join. It can block bile flow and cause yellow skin or eyes, itching, dark urine, pale stools, weight loss, or repeated bile-duct infections.

\medterm{Hill's Sign} Hill's sign is an older physical finding described in some cases of aortic regurgitation; it is not sufficient by itself to diagnose the condition.

\medterm{Hill-Sachs Lesion} A dent or compression fracture in the back of the upper-arm bone caused when the shoulder dislocates forward, sometimes contributing to repeated dislocations.

\medterm{Hilum} A depression or opening where blood vessels, nerves, ducts, or other structures enter or leave an organ.

\medterm{Hilus} The hilus is an indentation or entry point on an organ where vessels, nerves, ducts, or airways enter or leave.

\medterm{Hindbrain} The hindbrain is the lower part of the developing brain, giving rise to the medulla, pons, and cerebellum and helping control vital functions and coordination.

\medterm{Hindgut} The hindgut is the embryonic part of the digestive tract that forms much of the distal colon, rectum, and upper anal canal.

\medterm{Hip} The joint and surrounding region where the thigh bone meets the pelvis, allowing movement and bearing body weight.

\medterm{Hip Arthroscopy} A minimally invasive procedure that uses a camera and instruments inserted through small incisions to diagnose and treat selected disorders inside or around the hip joint.

\medterm{Hip Bursitis} Hip bursitis is inflammation of a bursa near the hip, commonly causing pain over the outer hip that worsens with lying on that side, stairs, or prolonged walking. Activity changes, physical therapy, medicines, or injection may help after other causes are considered.

\synonyms
Trochanteric Bursitis; Greater Trochanteric Pain Syndrome

\medterm{Hip Circumference} The measurement around the widest part of the hips or buttocks, used in anthropometry and sometimes combined with waist circumference to assess body-fat distribution.

\medterm{Hip Dislocation} The ball at the top of the thighbone comes out of the hip socket. It is usually caused by major injury and causes severe pain and inability to move the leg; nearby nerves and the blood supply to the bone may also be injured.

\medterm{Hip Disorders} The hip joints are ball-and-socket joints connecting the femur to the pelvis, providing
stability and a wide range of motion for walking, running, and weight-bearing
activities. Common hip conditions include osteoarthritis, fractures, bursitis, and
labral tears.

\medterm{Hip Dysplasia} A problem present from birth or developing early in life in which the hip socket and upper thighbone do not fit together securely. The hip may be unstable or dislocate. Early examination and ultrasound or X-ray help guide treatment and protect hip development.

\synonyms
Developmental Dysplasia of the Hip (DDH); Congenital Hip Dislocation

\medterm{Hip Fracture} A break in the upper thighbone, commonly near the femoral neck or just below it. It often follows a fall in someone with fragile bones and causes hip or groin pain, inability to bear weight, and sometimes a shortened or outward-turned leg.

\medterm{Hip Fracture In Older Adults} A fracture of the upper thighbone in an older person, commonly after a fall or minor force when bone strength is reduced. It can cause pain, loss of mobility, delirium, blood clots, and loss of independence.

\synonyms
Femoral Neck Fracture; Geriatric Hip Fracture

\medterm{Hip Fracture Surgery} An operation to stabilize or replace a fractured proximal femur, using internal fixation, hemiarthroplasty, or total hip arthroplasty according to fracture pattern and patient factors.

\medterm{Hip Impingement} Abnormal contact between parts of the hip joint during movement, often due to cam or pincer morphology. It may cause groin pain, stiffness, and reduced motion and can contribute to labral or cartilage damage.

\medterm{Hip Injury} Damage involving the hip joint, surrounding muscles, tendons, ligaments, or bones; causes range from strain and bursitis to fracture and dislocation, with inability to bear weight requiring prompt assessment.

\medterm{Hip Joint} The ball-and-socket joint where the top of the thighbone fits into the pelvis. It supports body weight and allows the leg to move in many directions; strong ligaments and a cartilage rim help keep it stable.

\medterm{Hip Joint Injection} Injection of local anesthetic, corticosteroid, or another prescribed agent into or around the hip joint for diagnosis or temporary relief of pain and inflammation.

\medterm{Hip Joint Replacement} Arthroplasty replaces damaged parts of the hip with prosthetic components to relieve pain and improve mobility, commonly for advanced osteoarthritis, fracture, or other destructive joint disease.

\medterm{Hip Labral Tear} A tear in the ring of cartilage surrounding the hip socket. It may cause groin pain, clicking, catching, stiffness, or reduced movement and can result from injury, repeated stress, or abnormal hip shape.

\medterm{Hip Pain} Pain in or around the hip caused by joint disease, tendon or muscle injury, bursitis, fracture, nerve problems, or pain referred from the back.

\medterm{Hip Prosthesis} An artificial hip joint used to replace a damaged hip. It usually has a stem placed in the thighbone, a ball, and a socket; parts may be metal, ceramic, or durable plastic and may be fixed with or without bone cement.

\medterm{Hip Replacement} A surgical procedure in which the diseased hip joint is replaced with a prosthetic
implant (arthroplasty). Total hip replacement (THR) replaces both the femoral head and
the acetabulum.

\medterm{Hip-Fracture Rehabilitation} Recovery care after a broken hip that helps restore walking, strength, balance, and independence. It may include early movement, pain control, exercises, nutrition, prevention of blood clots, delirium prevention, and treatment to reduce future fracture risk.

\medterm{Hippocampus} A pair of structures deep in the brain that help form and retrieve new memories and support learning and navigation. Damage to them can cause difficulty forming new memories.

\medterm{Hippocrates} An ancient Greek physician traditionally regarded as a major founder of Western clinical medicine; the Hippocratic corpus emphasizes observation, prognosis, and professional conduct.

\medterm{Hippocratic Oath} A traditional promise that doctors make to practice medicine ethically, help patients, protect their privacy, and avoid harm. Modern versions differ by country and institution.

\medterm{Hippotherapy} A treatment in which a trained therapist uses the movement of a horse to support physical, occupational, or speech therapy. It may help some people improve balance, posture, muscle control, coordination, or communication, but it is tailored to individual goals and safety needs.

\medterm{Hippus} Hippus is rhythmic, involuntary widening and narrowing of the pupil, which may occur normally or with neurologic or eye disease.

\medterm{Hirayama Disease} A condition affecting the lower neck spinal cord, usually in teenage boys and young men. It causes slowly worsening, often one-sided weakness and wasting of the hand and forearm muscles. Neck bending can press the spinal cord and reduce its blood supply. Symptoms commonly stop progressing after several years. A neck MRI taken during bending helps confirm the diagnosis; early neck protection may limit damage.

\synonyms
Monomelic Amyotrophy; Juvenile Cervical Flexion Myelopathy

\medterm{Hirschsprung Disease} A condition present from birth in which a section of the large intestine lacks the nerve cells needed to relax and move stool. Newborns may fail to pass stool, develop a swollen abdomen, or have severe constipation.

\synonyms
Congenital Aganglionic Megacolon; Colonic Aganglionosis

\medterm{Hirschsprung's Disease} A birth condition in which part of the large intestine lacks the nerve cells needed to relax and move stool. Newborns may fail to pass stool, develop a swollen abdomen, vomiting, or severe constipation.

\medterm{Hirsutism} Excess coarse hair growth in a woman or person with ovaries in areas where men commonly grow hair, such as the face, chest, or abdomen. It may result from increased androgen hormones, polycystic ovary syndrome, medicines, or other conditions.

\synonyms
Excessive Terminal Hair Growth; Hypertrichosis in Women

\medterm{Hirudin} A substance first found in medicinal leech saliva that prevents blood from clotting by blocking an important clotting protein. Related medicines have been used or studied for selected clotting conditions.

\medterm{Hirudins} A family of thrombin-inhibiting anticoagulant peptides, including natural and recombinant agents, used or studied to prevent or treat clotting in selected clinical situations.

\medterm{Hirudo Medicinalis} The medicinal leech, used in selected medical procedures to relieve venous congestion after reconstructive surgery.

\medterm{His} Usually referring to the bundle of His, the atrioventricular conduction bundle that carries electrical impulses from the AV node into the right and left bundle branches.

\medterm{His Bundle Electrography} An invasive electrophysiologic recording of electrical activity in the atrioventricular His bundle, used to assess conduction-system disease and localize atrioventricular block.

\medterm{His Disease} Trench fever, a louse-borne infection caused by \textit{Bartonella quintana}, classically producing recurrent or five-day fever, headache, and severe shin or leg pain.

\medterm{Histamine} A chemical messenger released during allergic and inflammatory reactions. It can cause itching, redness, swelling, mucus production, narrowing of the airways, and increased stomach acid.

\medterm{Histamine Degradation} The body’s breakdown and removal of histamine after it has acted. This limits effects such as itching, swelling, airway narrowing, and increased stomach acid.

\medterm{Histamine Receptor} A protein on or in a cell that responds to histamine. Different receptor types help control allergic symptoms, stomach-acid production, nerve signaling, and immune responses.

\medterm{Histamine Release} Discharge of histamine from mast cells or basophils during allergic, immune, or nonimmune activation, causing effects such as itching, vasodilation, swelling, mucus production, and bronchoconstriction.

\medterm{Histidine} Histidine is an amino acid used in protein synthesis and a precursor of histamine; it is essential during growth and can be metabolized to urocanic acid and other products.

\medterm{Histidine Decarboxylase} An enzyme that converts the amino acid histidine into histamine. It is found in human cells and some microorganisms; histamine influences allergy, stomach acid, blood-vessel responses, and immunity.

\medterm{Histidinemia} An inherited metabolic disorder caused by reduced histidase activity, leading to elevated histidine in blood and urine; many affected people have few or no clinically important symptoms.

\medterm{Histiocyte} An immune cell found in tissues that can engulf germs, damaged cells, or foreign material and help alert other immune cells.

\medterm{Histiocytic Sarcoma} A rare aggressive cancer of cells that normally help process immune signals and foreign material. It may arise in lymph nodes, the intestine, skin, or other organs and requires specialist pathology and cancer care.

\medterm{Histiocytoma} A tumor or tumor-like proliferation of histiocytes; the term may refer to a benign fibrous skin lesion or, less commonly, a malignant histiocytic neoplasm depending on classification.

\medterm{Histiocytosis} A group of uncommon disorders in which certain immune cells build up abnormally in the skin, bones, organs, or blood. Effects range from isolated skin or bone problems to serious disease affecting several organs.

\medterm{Histochemical} Histochemical means using chemical stains or reactions in tissue sections to identify particular substances or cell components.

\medterm{Histocompatibility} The degree to which donor and recipient tissues are immunologically alike. Matching tissue markers, especially HLA markers, can lower rejection risk, but antibody testing, organ type, and the recipient's health also influence transplant success.

\medterm{Histocompatibility Antigen Test} A test identifying human leukocyte antigen types, which influence tissue compatibility for transplantation and are relevant to some immune-mediated and genetic associations.

\medterm{Histocompatibility Testing} Laboratory typing of HLA and other immune markers to assess donor--recipient compatibility before transplantation and to investigate selected immune or transfusion-related questions.

\medterm{Histogenesis} The developmental process by which unspecialized cells differentiate and organize into mature tissues with characteristic structures and functions.

\medterm{Histological Techniques} Laboratory methods for preparing, staining, sectioning, and examining tissue under a microscope to reveal cellular architecture, deposits, organisms, or disease-related changes.

\medterm{Histology} The study of cells and tissues under a microscope. A laboratory prepares a thin tissue section, applies stains that show its structures, and examines it for normal features or changes caused by disease.

\medterm{Histology Cassette} A labeled perforated container that holds a tissue specimen during processing, allowing fixation, dehydration, paraffin infiltration, and later embedding for microscopic sectioning.

\medterm{Histone} A protein around which DNA winds so it can fit inside the cell nucleus. Chemical changes to histones help control which genes are turned on or off.

\medterm{Histone Acetylation} Addition of acetyl groups to histone proteins that usually loosens DNA packaging and can increase access of transcriptional machinery to genes.

\medterm{Histone Acetyltransferase} An enzyme that transfers acetyl groups to histone proteins, generally promoting a more open chromatin state and regulating gene transcription.

\medterm{Histone Deacetylases} Enzymes that remove acetyl groups from histones and other proteins, influencing chromatin structure, gene expression, cell differentiation, and disease biology.

\medterm{Histone Deacetylation} Removal of acetyl groups from histones, often tightening chromatin and reducing transcription; histone deacetylases are targets of some anticancer drugs.

\medterm{Histone Modification} A chemical change to a histone protein that alters how tightly DNA is packed and can change which genes a cell uses. These changes help regulate cell growth, development, and response to disease.

\medterm{Histopathologic Grade} A microscopic estimate of how abnormal and aggressive tumor cells appear compared with their tissue of origin; grade contributes to prognosis and treatment planning alongside stage.

\medterm{Histopathology} The diagnosis of disease by examining tissue under a microscope. A pathologist studies a biopsy or surgical specimen for abnormal cells, inflammation, infection, scarring, or cancer.

\medterm{Histoplasma} A genus of fungi that includes Histoplasma capsulatum, which can cause histoplasmosis after inhalation of environmental spores.

\medterm{Histoplasma Capsulatum} A fungus found in soil contaminated with bird or bat droppings that can cause histoplasmosis when
its spores are inhaled.

\medterm{Histoplasma Complement Fixation} A blood test that detects complement-fixing antibodies to Histoplasma and can support evaluation of histoplasmosis, although sensitivity, timing, cross-reactions, and immune status affect interpretation.

\medterm{Histoplasmosis} A fungal infection acquired by breathing in spores from soil contaminated with bird or bat droppings. It often affects the lungs and may be mild, but can spread to other organs, especially in people with weakened immunity.

\medterm{Histrionic Personality Disorder} A long-term pattern of intense emotional expression, strong need for attention, and behavior that may seem dramatic, suggestible, or uncomfortable when not the focus. Diagnosis requires a persistent pattern that causes significant problems and is made by a qualified mental-health professional.

\medterm{Hitchhiker's Thumb} A descriptive term for an unusually flexible thumb that can bend backward at the distal joint. It is usually a harmless inherited variation unless associated with pain, instability, or a broader connective-tissue disorder.

\medterm{HIV} Abbreviation for human immunodeficiency virus; see \textit{Human immunodeficiency virus}.

\medterm{Hiv Antibodies} Proteins made by the immune system after HIV infection and detected by blood or oral-fluid tests. A positive screening result needs confirmatory testing, and a negative result may be unreliable soon after exposure.

\medterm{HIV Infection} Infection with human immunodeficiency virus, which progressively damages immune function if untreated.

\medterm{Hiv Seropositivity} The presence of detectable HIV antibodies or other blood-test evidence of infection. A reactive screening test is not final until confirmatory testing is completed.

\medterm{HIV Test} A blood or oral-fluid test that looks for HIV infection. Laboratory tests may detect viral material or antibodies; a reactive screening result needs confirmatory testing. Tests can be negative early in infection, so repeat testing may be needed after a recent exposure.

\medterm{HIV Transmission And Prevention} The spread of HIV through infected blood, semen, vaginal fluids, rectal fluids, or breast milk, prevented by testing, treatment that suppresses viral load, condoms, sterile equipment, and pre- or post-exposure prophylaxis.

\medterm{HIV/AIDS} Common combined term referring to HIV infection and its most advanced stage, AIDS; see \textit{HIV infection} and \textit{AIDS}.

\medterm{HIV/AIDS Immunopathogenesis} HIV enters and gradually destroys or disables important immune cells. Without treatment, immune defenses become weaker and serious infections or cancers may develop. Medicines stop most viral multiplication and allow immune recovery.

\medterm{HIV/AIDS In Pregnant Women And Infants} HIV care in pregnancy and infancy uses maternal antiretroviral therapy, viral-load monitoring, delivery and feeding planning, infant prophylaxis, and early infant diagnostic testing to reduce vertical transmission.

\medterm{HIV/AIDS Nutrition} Nutritional care for people living with HIV or AIDS. Infection, poor appetite, diarrhea, medicine side effects, or other illnesses may cause weight loss or nutrient shortages; care may include food support, supplements, treatment of infections, and review of medicines.

\medterm{HIV-Associated Rheumatic Disease} Joint, muscle, bone, or connective-tissue problems associated with HIV infection or its immune effects. They may include inflammatory arthritis, reactive arthritis, muscle disease, osteoporosis, or symptoms related to infection or treatment.
Antiretroviral therapy suppresses viral replication and reduces transmission and complications.

\medterm{Hives} The common term for urticaria, describing raised, intensely itchy welts (wheals) on the skin. They happen when immune cells in the skin release histamine, causing fluid to leak from tiny blood vessels. Individual welts typically appear quickly, change shape, and fade away within 24 hours without leaving a scar. Hives lasting under six weeks are acute, often triggered by viral infections, foods, insect stings, or medications. Hives lasting longer than six weeks are chronic and often occur without a known external trigger.

\synonyms
Urticaria; Urticarial Weals

\medterm{Hives And Angioedema} Hives are raised itchy welts; angioedema is deeper swelling of the skin or mucous membranes, often caused by allergy, medicines, or other immune reactions.

\medterm{HLA} A group of cell-surface proteins helping the immune system recognize the body's own cells, reject foreign tissue, and influence transplant compatibility.

\synonyms
Human Leukocyte Antigen System; Major Histocompatibility Complex (MHC)

\medterm{Hla Antigen} A marker on the surface of cells that helps immune cells recognize what belongs to the body. HLA markers are important for matching donors and recipients in transplantation and are linked with risk of some immune diseases.

\medterm{HLA System} The HLA system is a group of cell-surface proteins that present antigens to immune cells and influence transplant matching and autoimmune risk.

\medterm{HLA-B27} An inherited immune-system marker associated with increased risk of some inflammatory conditions affecting the spine, joints, eyes, or bowel. A positive test does not prove that a person has one of these diseases.

\medterm{HLA-B27 Antigen} A cell-surface immune marker associated with increased risk of ankylosing spondylitis and related spondyloarthritis, reactive arthritis, and uveitis; a positive result is not diagnostic by itself.

\medterm{Hmgb Proteins} Proteins that help organize genetic material inside cells. When released from damaged cells, some can also signal tissue injury and promote inflammation.

\medterm{HMG-CoA Reductase} An enzyme the liver uses to make cholesterol. Statin medicines partly lower blood cholesterol by blocking this enzyme.

\medterm{Hmgn Proteins} Proteins that help control how genetic material is arranged and used inside cells. They can influence cell growth, specialization, and repair of genetic material.

\medterm{Hoarding} Alternate name; see \textit{Hoarding Disorder}.

\medterm{Hoarding Disorder} A mental-health condition involving persistent difficulty discarding possessions, even when they have little value. Saving items causes distress and can fill living areas so that normal activities or safety are impaired.

\medterm{Hoarseness} A change in the voice that makes it rough, weak, breathy, or strained. Common causes include a viral infection, voice overuse, smoking, reflux, or a growth on the vocal folds. Hoarseness lasting more than a few weeks needs assessment.

\medterm{Hodgkin Lymphoma} A cancer of a type of immune cell, usually beginning in lymph nodes. It often causes painless swelling of the nodes, fever, night sweats, itching, or weight loss and is diagnosed by examining a tissue sample.

\medterm{Hodgkin Lymphoma In Children} A malignant lymphoma arising from lymphoid tissue in children, characterized by Hodgkin and Reed-Sternberg cells and treated according to subtype, stage, and risk.

\medterm{Hodgkin's And Non-Hodgkin's Lymphoma Comparison} Hodgkin lymphoma is defined by characteristic Reed-Sternberg cells and commonly spreads in an orderly pattern, while non-Hodgkin lymphomas comprise many B-cell, T-cell, and other lymphoid cancers. Biopsy, immunophenotyping, and staging determine treatment and prognosis.

\medterm{Hodgkin's Disease} Alternate name; see \textit{Hodgkin Lymphoma}.

\medterm{Hodgkin's Disease} Historical term; see \textit{Hodgkin lymphoma}.

\medterm{Hodgkin's Lymphoma} A B-cell lymphoma characterized by the presence of Reed-Sternberg cells (large,
binucleate or multinucleate cells with owl-eye nucleoli) in a background of reactive
inflammatory cells. See Hodgkin Lymphoma.

\medterm{Hoffmann Sign} A reflex test where flicking the middle fingernail makes the thumb and index finger quickly twitch inward, pointing to irritation or pressure on the spinal cord in the neck.

\medterm{Hoffmann Syndrome} An uncommon form of adult muscle weakness caused by severe, untreated underactive thyroid disease. It may cause stiff or painful muscles, cramps, enlarged-looking muscles, slow reflex relaxation, and high muscle-enzyme levels; treating the thyroid problem usually improves it.

\synonyms
Hypothyroid Myopathy with Pseudohypertrophy

\medterm{Hofmann's Syndrome} A rare form of adult hypothyroid muscle disease causing muscle enlargement, stiffness, cramps, weakness, and slowed reflexes. Treating the low thyroid hormone level usually improves the muscle problems.

\medterm{Holistic Medicine} An approach that considers physical, psychological, social, behavioral, and sometimes spiritual factors in health and illness. It may support patient-centered care, but individual therapies should be assessed for evidence, safety, and compatibility with effective conventional treatment.

\medterm{Holmium} A rare-earth metal used in some medical lasers and radioactive tracers. Holmium laser energy can cut or vaporize tissue, especially in urologic procedures; radioactive holmium requires strict handling and safety controls.

\medterm{Holmium Laser Enucleation Of The Prostate} A procedure performed through the urine passage that uses a laser to separate and remove excess prostate tissue blocking urine flow. The tissue is removed from the bladder and examined; temporary bleeding, burning, leakage, infection, or sexual and urinary changes can occur.

\synonyms
HoLEP; Holmium Prostate Enucleation

\medterm{Holoprosencephaly} A birth difference in which the developing brain does not divide fully into two cerebral hemispheres. Severity varies; it may affect facial development, movement, feeding, seizures, and intellectual development.

\medterm{Holorachischisis} A severe birth defect in which a long section of the spine and spinal cord fails to close during development. It usually causes major weakness, loss of sensation, and other serious nervous-system problems.

\medterm{Holosystolic Murmur Differential} A heart murmur heard throughout the period when the heart's lower chambers contract. Common causes include leakage of the mitral or tricuspid valve and a hole between the lower chambers; its sound, location, and other findings help identify the cause.

\medterm{Holter Electrocardiography} Continuous ambulatory ECG recording, usually for 24 hours or longer, used to detect intermittent arrhythmias, correlate symptoms with rhythm, and assess heart-rate patterns.

\medterm{Holter Monitor} A small portable recorder that continuously tracks the heart’s electrical rhythm, usually for a day or longer. It can detect intermittent abnormal rhythms and compare them with symptoms such as palpitations, dizziness, or fainting.

\medterm{Holter Monitor} A portable ambulatory electrocardiographic recorder that continuously monitors heart rhythm, usually for 24 hours or longer, to detect intermittent arrhythmias and correlate them with symptoms.

\medterm{Holter Monitoring} Portable, continuous electrocardiographic recording over an extended period, commonly 24 hours or longer, to detect intermittent rhythm abnormalities and correlate them with symptoms.

\medterm{Holt-Oram Syndrome} An inherited condition causing abnormalities of the bones in one or both arms or hands and heart defects. The heart problem may be a hole between chambers or an abnormal rhythm; the severity varies widely.

\medterm{Homans Sign} Calf discomfort when the foot is gently bent upward toward the shin, historically checked as a clue for a blood clot in the leg.

\medterm{Homeobox Gene} A gene that helps control how an embryo’s body is laid out and how organs and tissues develop. Changes in some homeobox genes can cause birth differences or disease.

\medterm{Homeodomain Proteins} Transcription factors containing a homeodomain that bind regulatory DNA sequences and direct developmental patterning and cell differentiation.

\medterm{Homeopathy} An alternative-health system that uses highly diluted preparations selected according to “like cures like”; reliable evidence has not shown homeopathic products to treat serious disease beyond placebo effects.

\medterm{Homeostasis} The body’s ability to keep internal conditions within a safe range despite changes outside it. It regulates temperature, blood sugar, blood pressure, fluid, salts, and acidity through coordinated feedback systems.

\medterm{Homeotherm} An organism that maintains a relatively stable internal body temperature despite changes in environmental temperature, largely through physiological regulation.

\medterm{Homicide} The killing of one person by another, classified legally according to intent and circumstances; medical assessment may address injury patterns, cause of death, toxicology, and safeguarding.

\medterm{Homo Sapiens} The scientific name for modern humans, the living human species with shared biological features, genetic variation, and worldwide populations.

\medterm{Homocysteine} An amino acid formed during methionine metabolism. Persistently elevated blood levels may occur with genetic, nutritional, renal, or other disorders.

\medterm{Homocystinuria} A rare inherited disorder in which the body cannot properly process homocysteine, an amino acid. It may cause eye problems, abnormal bones, learning or developmental problems, and an increased risk of blood clots.

\medterm{Homograft} A transplant of tissue or an organ between two people of the same species but with different genetic makeup. The recipient’s immune system may reject it, so matching and immune-suppressing treatment are important.

\medterm{Homologous} Similar because of a shared origin, structure, or genetic relationship. In medicine, homologous tissue or cells come from the same species, while homologous chromosomes carry corresponding genes.

\medterm{Homologous Recombination} A DNA-repair and genetic-exchange process that uses a matching or nearly matching sequence as a template to repair breaks or reshuffle genetic material.

\medterm{Homologous Transplantation} Transplantation of cells, tissue, or an organ between genetically different people of the same species, requiring attention to compatibility, rejection, infection, and immunosuppression.

\medterm{Homonymous Hemianopia} A visual field deficit characterized by the loss of the same half of the visual field (either both right halves or both left halves) in both eyes. It results from a retrochiasmal lesion affecting the optic tract, lateral geniculate nucleus, optic radiation, or primary visual cortex on the side opposite the visual defect.

\medterm{Homosexual} Describing a person who experiences enduring romantic or sexual attraction to people of the same sex; sexual orientation is not a mental disorder.

\medterm{Homosexuality} Sexual or romantic attraction to people of the same sex or gender; it is not a mental disorder or pathology.

\medterm{Homovanillic Acid} A breakdown product of dopamine found in cerebrospinal fluid, blood, and urine; levels may reflect dopamine metabolism but are interpreted cautiously because many factors affect them.

\medterm{Homozygote} An individual or cell carrying two identical alleles at a particular gene locus, whether the alleles are normal, pathogenic, or another variant.

\medterm{Homozygous} Having two copies of the same version of a gene, one inherited from each parent. The effect depends on the gene and whether that version is harmless, protective, or disease-causing.

\medterm{Honey-Bee Stings} Injuries caused by a honey-bee injecting venom into the skin. A sting usually causes local pain, redness, and swelling; a severe allergic reaction can cause breathing difficulty, fainting, or shock and requires emergency treatment.

\medterm{Honolulu Fever} An infection caused by Streptobacillus bacteria, usually after contact with rats or contaminated food. It may cause fever, vomiting, headache, muscle or joint pain, and a rash and requires antibiotics.

\medterm{Hookworm} An intestinal infection caused by parasitic worms whose larvae live in contaminated soil. The larvae can enter through bare skin, and the adult worms may cause abdominal symptoms, blood loss, anemia, or poor growth.

\medterm{Hookworm Infection} Intestinal infection caused by hookworms acquired when larvae in contaminated soil
penetrate the skin or are ingested. Chronic infection can cause abdominal symptoms,
iron-deficiency anemia, protein loss, and impaired growth in children.

\medterm{Hookworm Infestation} Intestinal infection by hookworm nematodes, which attach to the intestinal wall and may cause iron-
deficiency anemia, abdominal symptoms, or poor growth.

\medterm{Hordeolum} Alternate medical name; see \textit{Stye}.

\medterm{Horizontal Gene Transfer} The movement of genetic material between microorganisms rather than from parent to offspring. It can spread traits such as antibiotic resistance through direct contact, viruses, or uptake of DNA from the surroundings.

\medterm{Hormonal Effects In Newborns} Transient newborn findings caused by maternal hormones or normal neonatal endocrine transitions, such as breast enlargement or vaginal bleeding; persistent or severe findings require evaluation.

\medterm{Hormone} A chemical messenger made by glands or other tissues and carried in the blood to target cells. It binds specific receptors to regulate functions such as growth, metabolism, reproduction, stress responses, and fluid balance.

\medterm{Hormone Levels} Measurements of hormones in blood, urine, or other specimens used to evaluate endocrine function; interpretation depends on the specific hormone, timing, age, sex, medications, and clinical context.

\medterm{Hormone Receptor} A protein that binds a specific hormone and converts the signal into cellular responses, either at the cell surface or inside the cell nucleus.

\medterm{Hormone Replacement Therapy} Treatment that supplies a hormone the body does not make enough of, such as thyroid, adrenal, sex, or pituitary hormones. Menopausal hormone therapy can relieve hot flushes and vaginal symptoms for selected people, but benefits and risks depend on the hormone, dose, duration, and health history.

\medterm{Hormone Therapy} Treatment that adds, lowers, blocks, or changes the action of hormones. In some breast and prostate cancers, it blocks hormones that the cancer needs to grow.

\medterm{Hormone Therapy For Breast Cancer} Treatment that blocks or lowers estrogen or progesterone effects in hormone-sensitive breast cancer. It reduces stimulation of cancer cells and may be used after surgery or for advanced disease.

\medterm{Hormone Therapy For Prostate Cancer} Treatment that lowers androgen production or blocks androgen action to slow growth of hormone-sensitive prostate cancer.

\medterm{Hormone-Receptor Positive} A cancer whose cells express receptors for hormones such as estrogen or progesterone. Hormone-blocking treatment may be useful for selected cancers with these receptors.

\medterm{Hormones} Chemical messengers released by glands or tissues that travel through blood or nearby spaces. They regulate growth, metabolism, reproduction, stress, fluid balance, mood, sleep, and many other body processes.

\medterm{Hormone-Sensitive Lipase} An enzyme in fat cells that breaks stored triglycerides into fatty acids and glycerol. Its activity rises when the body needs fuel, including during fasting or exercise, and falls after insulin increases.

\medterm{Horn} A horn is a pointed, hard growth of keratinized tissue, usually on the skin or as a normal anatomical projection such as a uterine or spinal horn.

\medterm{Horner Syndrome} A group of eye and face changes caused by damage to certain nerves. One side may have a smaller pupil, a drooping eyelid, and less sweating; the cause can range from migraine to a serious neck, chest, or brain problem.

\synonyms
Horner's Oculosympathetic Palsy; Bernard-Horner Syndrome

\medterm{Horner's Syndrome} A group of findings caused by disruption of sympathetic nerves, including a drooping eyelid, a small pupil, and reduced sweating on one side of the face.

\medterm{Hornet Stings} Hornet stings inject venom and usually cause local pain and swelling; multiple stings or allergy can cause a serious systemic reaction.

\medterm{Horripilation} Piloerection, or the temporary raising of hair and development of gooseflesh due to cold, fear, emotion, or autonomic stimulation.

\medterm{Horseshoe Kidney} A birth difference in which the lower parts of the two kidneys join across the middle of the body. Many people have no symptoms, but kidney stones, urine-flow problems, infections, or high blood pressure can occur.

\medterm{Hospice} Care for a person nearing the end of life when the focus is comfort rather than cure. Hospice supports symptoms, emotional and practical needs, and family caregivers at home, in a hospice facility, or in a hospital; eligibility rules vary by location.

\medterm{Hospice Care} Team-based care for a person nearing the end of life that prioritizes comfort, symptom control, communication, and family support rather than disease-directed cure.

\medterm{Hospice Rehabilitation} Rehabilitation provided during hospice care to maximize comfort, dignity, and participation rather than to restore function at all costs. It may include positioning, pain relief, safe transfers, equipment, and caregiver education.

\medterm{Hospital} A healthcare facility providing inpatient, emergency, surgical, diagnostic, intensive, rehabilitation, and specialist services for people needing medical care.

\medterm{Hospital Medicine} The medical specialty focused on comprehensive care of hospitalized patients, including diagnosis, treatment, coordination across teams, transitions of care, and prevention of inpatient complications.

\medterm{Hospital-Acquired Infection} An infection developing during or soon after hospital care that was not present on admission. It may involve wounds, urine, lungs, blood, or devices and can be caused by resistant organisms.
Examples: CLABSI, CAUTI, surgical site infection.

\medterm{Hospital-Acquired Infections} Infections that develop during or after hospitalisation. Major types include central
line-associated bloodstream infections, catheter-associated urinary tract infections,
ventilator-associated pneumonia, and surgical site infections. Risk factors include
invasive devices, prolonged hospitalisation, and antimicrobial use.

\medterm{Hospital-Acquired Pneumonia} Lung infection that begins at least 48 hours after hospital admission and was not present when the person arrived. It may cause fever, cough, breathing difficulty, or low oxygen; treatment depends on severity, cultures, and resistance risk.

\medterm{Hospital-Associated Disability} New or worsened difficulty with walking, self-care, or other daily activities that develops during or after hospitalization. It may result from illness, bed rest, weakness, delirium, poor nutrition, or treatments and can persist after discharge.

\medterm{Hospitalism} Physical or emotional problems associated with prolonged hospitalization, such as weakness, confusion, reduced mobility, sleep disruption, or loss of independence. Early movement, orientation, social contact, and individualized care can reduce these effects.

\medterm{Hospitalization} Admission to a hospital for monitoring, diagnosis, treatment, surgery, rehabilitation, or intensive support that cannot safely be delivered in an outpatient setting.

\medterm{Hospitals As Health Educators} Hospitals educate patients, families, staff, and communities about prevention, treatment, medicines, recovery, and when to seek care.

\medterm{Host} A person or other living organism that can be infected by a virus or other pathogen
under natural conditions. The host may be a human, animal, or plant. In medical
contexts, the host is the patient in whom an infection or disease develops.

\medterm{Host Cell} A cell in which an infectious agent, parasite, or symbiont lives, replicates, or uses cellular machinery; the host-cell response influences disease and clearance.

\medterm{Host Response} The body's reaction to an infection, foreign material, transplanted tissue, or other stimulus. It may include inflammation, immune defense, healing, scar formation, rejection, or tolerance.

\medterm{Hostility} An attitude or behavior marked by antagonism, anger, distrust, or intent to oppose another person; persistent severe hostility may occur with psychiatric, neurologic, substance-related, or situational problems.

\medterm{Hot Flash} Alternate singular form; see \textit{Hot Flashes}.

\medterm{Hot Flashes} Sudden episodes of intense heat, sweating, flushing, and sometimes a racing heart or chills. They are common during menopause but can also result from hormone changes, medicines, cancer treatment, or other illness.

\medterm{Hot Flushes} Alternate name; see \textit{Hot Flashes}.

\medterm{Hot Tub Folliculitis} A temporary bacterial infection of hair follicles, commonly caused by Pseudomonas aeruginosa after exposure to contaminated hot tubs or pools, producing itchy follicular papules or pustules.

\medterm{Hounsfield Unit} A number on a CT scan that describes how strongly a tissue absorbs X-rays. Air has a very low value, water is near zero, and bone has a high value; the numbers help distinguish tissues and lesions.

\medterm{Hounsfield Units} Numbers used on CT scans to describe how strongly tissues absorb X-rays. Air has a very low value, water is near zero, and bone has a high value; the numbers help distinguish tissues and abnormal areas.

\medterm{Hour} A unit of time equal to 60 minutes, commonly used in medicine to describe symptom duration, dosing intervals, monitoring periods, and treatment timing.

\medterm{Hour-Glass Contraction} An abnormal tight band that divides the uterus into upper and lower sections during difficult labor. It can prevent the baby from being delivered normally and requires urgent assessment by an obstetric team.

\medterm{House Officer} A physician receiving postgraduate clinical training in a hospital, such as an intern, resident, or fellow; usage varies by country and institution.

\medterm{House Staff} The physicians and other trainees who provide patient care within a hospital under the supervision of attending clinicians, commonly including interns, residents, and fellows.

\medterm{Household Glue Poisoning} Harm from swallowing or inhaling household glue. Small accidental amounts may cause irritation or stomach upset, while solvents can cause dizziness, unconsciousness, breathing problems, or heart-rhythm disturbances; seek poison-control advice.

\medterm{Housemaid's Knee} A common name for prepatellar bursitis, inflammation and fluid accumulation in the bursa over the kneecap, often caused by prolonged kneeling, trauma, infection, or inflammatory disease. Painful swelling, redness, warmth, or fever requires assessment for septic bursitis.

\medterm{Howell-Jolly Bodies} Small leftover pieces of red-cell DNA seen on a blood smear. They are usually removed by a working spleen, so their presence may indicate absent or reduced spleen function; they can also occur with severe vitamin-B12 or folate deficiency.

\medterm{HPA} A hormone system linking parts of the brain with the adrenal glands. It helps the body respond to stress by controlling the release of cortisol, which affects blood pressure, blood sugar, immune activity, sleep, and energy.

\synonyms
Hypothalamic-Pituitary-Adrenal Axis

\medterm{HPLC} A laboratory method that separates substances in a liquid sample so they can be identified and measured. It is used to test medicines, body chemicals, poisons, and other samples.

\medterm{HPV} Human papillomavirus is a group of viruses that infect the skin and moist body linings. Some types cause warts, while long-lasting infection with high-risk types can cause precancerous changes and cancers of the cervix, anus, genitals, or throat.

\medterm{HPV Test} A test that detects high-risk human papillomavirus genetic material or related markers in cervical samples to screen for cervical cancer risk or evaluate abnormal cervical cytology.

\medterm{HPylori} A shorthand name for \textit{Helicobacter pylori}, a bacterium that can live in the stomach and cause inflammation, ulcers, and increased stomach-cancer risk. Confirmed infection is treated with a combination of medicines and follow-up testing.

\medterm{HRCT} Alternate name; see \textit{High-Resolution Computed Tomography}.

\medterm{HRT} Hormone replacement therapy uses estrogen, with a progestogen when needed, to treat selected menopausal
symptoms caused by declining ovarian hormone production. It can relieve vasomotor and
genitourinary symptoms and help prevent bone loss, but benefits and risks depend on the
person’s age, medical history, regimen, and duration of use.

\medterm{Hürthle Cell Tumour} A growth arising from hormone-making cells in the thyroid. It may be benign or cancerous; ultrasound and examination of thyroid tissue help determine its nature and treatment.

\medterm{Hs-cTn} Alternate name; see \textit{High-Sensitivity Cardiac Troponin}.

\medterm{Ht} Common abbreviation for hematocrit, the percentage of blood volume occupied by red blood cells; in other contexts it may mean height.

\medterm{Huber Needle} A special non-coring needle used to access an implanted port for giving medicines, fluids, blood products, or taking blood samples. Correct placement and sterile technique reduce infection, bleeding, and port damage.

\medterm{Human Calcitonin} A peptide hormone produced by thyroid parafollicular cells that can lower blood calcium by inhibiting osteoclast activity; synthetic calcitonin is used selectively for bone and calcium disorders.

\medterm{Human Chorionic Gonadotropin} A hormone produced mainly by the placenta and measured in blood or urine to detect and monitor pregnancy. Levels normally rise early in pregnancy, but the result must be interpreted with timing, symptoms, and other findings.

\synonyms
HCG; Beta-HCG

\medterm{Human Chromosome Count} Most human body cells contain 46 chromosomes arranged in 23 pairs, including one pair of sex chromosomes; mature egg and sperm cells normally contain 23.

\medterm{Human Embryo} The developing human organism from fertilization through approximately the eighth week, during which the major body plan and organ systems are established.

\medterm{Human Gastrin} A peptide hormone produced mainly by gastric G cells that stimulates gastric acid secretion and supports growth of the gastric mucosa after food enters the stomach.

\medterm{Human Genome} The complete DNA sequence and genetic content of a human, organized into nuclear chromosomes and mitochondrial DNA; variants influence traits, disease risk, and treatment response.

\medterm{Human Herpesvirus 6} A herpesvirus that commonly infects children and can cause roseola; it may reactivate and cause serious disease in people with weakened immunity.

\medterm{Human Herpesvirus 8} A herpesvirus associated with Kaposi sarcoma, primary effusion lymphoma, and multicentric Castleman disease, especially in people with weakened immunity.

\medterm{Human Immunodeficiency Virus} A virus that attacks immune cells, especially CD4 cells, and can gradually weaken the body’s defenses. Without effective treatment it may lead to AIDS and serious infections or cancers. HIV-1 is widespread globally and HIV-2 is less common; daily treatment can suppress the virus and protect health.

\medterm{Human Leukocyte Antigen} A group of cell-surface proteins encoded by genes in the major histocompatibility complex. HLA molecules present antigens to T cells and are important in immune responses and tissue matching for transplantation.

\synonyms
HLA Complex; Human MHC

\medterm{Human Metapneumovirus} A respiratory virus that commonly causes cough, fever, runny nose, and wheezing. Infants, older adults, and people with weakened immunity may develop bronchiolitis, pneumonia, or severe breathing difficulty.

\medterm{Human Microbiome} The community of bacteria, viruses, fungi, archaea, and other microorganisms living on and inside the human body. These organisms and their genes can influence digestion, immune activity, metabolism, and disease, although finding a microorganism does not by itself prove that it causes illness.

\medterm{Human Nociceptin} A neuropeptide that binds the nociceptin/orphanin FQ receptor and modulates pain, stress, reward, autonomic function, and other brain and spinal-cord processes.

\medterm{Human Papillomavirus} A large group of viruses spread mainly through intimate skin-to-skin or sexual contact. Most infections cause no symptoms and clear without treatment. Some types cause warts, while persistent infection with high-risk types can cause precancerous changes and cancers of the cervix, anus, genitals, or throat. Vaccination and recommended cervical screening reduce risk.

\medterm{Human Papillomavirus} A large family of viruses spread mainly through intimate skin-to-skin or sexual contact. Some types cause warts and others can cause cancers; vaccination and recommended screening reduce the risk of serious disease.

\medterm{Human Papillomavirus Infection} Alternate name; see \textit{Human Papillomavirus}.

\medterm{Human Placental Lactogen} A hormone made by the placenta during pregnancy. It helps prepare the breasts for feeding and changes how the pregnant person uses fat and sugar, helping provide nutrients to the developing fetus.

\medterm{Humectant} A substance that attracts and holds water, often used in skin creams, lotions, foods, and medicines. In skincare, it draws water into the outer skin layer and can reduce dryness when used properly.

\medterm{Humeral Epicondyle} A bony prominence at the distal humerus that provides attachment for forearm muscles and ligaments; the medial and lateral epicondyles are located on opposite sides of the elbow.

\medterm{Humeral Fracture} A break in the upper-arm bone, involving the proximal, shaft, or distal humerus; pain, swelling, deformity, and possible nerve or blood-vessel injury guide urgent assessment and treatment.

\medterm{Humeroradial Joint} The part of the elbow where the upper-arm bone meets the radius, one of the forearm bones. It helps the elbow bend and straighten and supports rotation of the forearm.

\medterm{Humeroulnar Joint} The main hinge of the elbow, where the upper-arm bone meets the ulna. It allows the arm to bend and straighten and provides much of the elbow’s stability.

\medterm{Humerus} The long bone of the upper arm, extending from the shoulder to the elbow. It joins the shoulder blade above and the radius and ulna below; fractures can injure nearby nerves or blood vessels.

\medterm{Humidifier} A device that adds moisture to air; medical humidifiers require appropriate cleaning and water handling because contaminated equipment can disperse infectious or inflammatory material.

\medterm{Humidifier Fever} A flu-like illness caused by inhalation of microorganisms or biologic contaminants from a contaminated humidifier or ventilation system, often involving fever, chills, cough, and malaise.

\medterm{Humidity} The amount of water vapor in air, influencing heat loss, respiratory comfort, skin hydration, and survival or transmission of some microorganisms.

\medterm{Humor} In anatomy, a fluid of the body, especially the aqueous or vitreous humor of the eye; in ordinary language, it means a sense of amusement or comedy.

\medterm{Humoral Immunity} The part of the immune system that protects mainly through antibodies made by B cells. Antibodies can block germs and toxins and help other immune cells remove them.

\medterm{Hunchback} A nontechnical term usually referring to pronounced kyphosis, an excessive forward curvature of the upper spine. Causes include osteoporosis, vertebral compression fractures, Scheuermann disease, degenerative change, and congenital disorders; evaluation depends on age, progression, pain, and neurologic findings.

\medterm{Hunger} The physical drive to eat, influenced by stomach fullness, blood glucose, hormones such as ghrelin, emotions, habits, illness, and the availability or appeal of food.

\medterm{Hungry Bone Syndrome} A severe drop in blood calcium after treatment of overactive parathyroid glands, often after surgery. Healing bones rapidly take up calcium and other minerals, causing tingling, cramps, spasms, seizures, or abnormal heart rhythms.

\medterm{Hunter Syndrome} An inherited condition, usually affecting boys, in which iduronate-2-sulfatase deficiency causes complex sugars to accumulate in tissues. It can cause coarse facial features, enlarged organs, stiff joints, hearing loss, breathing problems, and developmental difficulties.

\medterm{Hunter's Syndrome} Another name for Hunter syndrome, an inherited disorder in which complex sugars accumulate because the body cannot break them down properly. It mainly affects boys and can damage several organs over time.

\medterm{Huntington Disease} A progressive inherited brain disorder caused by an abnormally long repeated section of DNA in the \textit{HTT} gene. The altered gene produces an abnormal huntingtin protein that damages and eventually kills nerve cells involved in movement, thinking, mood, and behavior. Symptoms may include involuntary, dance-like movements (chorea), muscle stiffness, poor coordination, changes in speech or swallowing, depression, irritability, and worsening memory or judgment. It usually begins in adulthood, but a less common juvenile form can begin during childhood or adolescence. The condition is inherited in an autosomal dominant pattern, so each child of a parent with one altered copy has a 50\% chance of inheriting it.

\synonyms
Huntington's Chorea; HD

\medterm{Hurler Syndrome} Mucopolysaccharidosis type I (MPS I), the severe form, an autosomal recessive lysosomal
storage disorder caused by alpha-L-iduronidase deficiency (IDUA gene mutations).
Accumulation of dermatan sulfate and heparan sulfate leads to progressive multisystem
involvement.

\medterm{Hurthle Cell} A thyroid follicular cell with abundant granular mitochondria-rich cytoplasm, seen in Hashimoto thyroiditis and Hürthle-cell neoplasms.

\medterm{Hurthle Cell Cancer} A rare thyroid cancer arising from Hurthle cells, with behavior ranging from indolent to aggressive.
\synonyms
Oxyphilic Thyroid Carcinoma; Oncocytoid Carcinoma

\medterm{Hutchinson's Teeth} Hutchinson's teeth are permanently notched, widely spaced upper front teeth classically associated with congenital syphilis.

\medterm{Hyaline} A glassy, translucent material or appearance, commonly describing proteinaceous deposits or homogeneous extracellular material seen microscopically in tissues.

\medterm{Hyaline Cartilage} The most common type of cartilage, containing fine type II collagen fibers in a firm,
glassy matrix. It supports airways, covers articular surfaces, and forms much of the
fetal skeleton and growth plates.

\medterm{Hyaline Casts} Clear, tube-shaped particles sometimes seen in urine. A small number may be normal, especially after exercise or dehydration; larger numbers can occur when the kidneys are under stress and must be interpreted with other urine and blood tests.

\medterm{Hyaline Membrane} A protein-rich layer that forms along the walls of damaged air sacs in the lungs. It blocks oxygen transfer and is classically seen in newborn respiratory distress caused by inadequate surfactant.

\medterm{Hyaline Membrane Disease} A former name for breathing failure in newborns, especially premature infants, caused mainly by too little lung surfactant. The air sacs collapse easily, making oxygen transfer difficult.

\medterm{Hyalitis} Inflammation of the vitreous body of the eye, usually associated with intraocular infection, autoimmune inflammation, trauma, or another ocular disorder. It may cause floaters, blurred vision, or visual loss and requires ophthalmic assessment.

\medterm{Hyaloid Canal} A narrow channel through the vitreous body of the eye that marks the former path of the hyaloid artery from the optic disc toward the lens during fetal development.

\medterm{Hyaluronic Acid} A naturally occurring carbohydrate polymer that binds water in connective tissue, skin, joints, and the eye.

\synonyms
Hyaluronan; Sodium Hyaluronate

\medterm{Hyaluronidase} An enzyme that breaks down hyaluronic acid in connective tissue. Medicines containing it can help injected fluids or medicines spread through tissue and may also dissolve certain hyaluronic-acid cosmetic fillers.

\medterm{Hyaluronoglucosaminidase} An enzyme that degrades hyaluronan and related glycosaminoglycans; hyaluronidase preparations can increase tissue permeability and assist dispersion of injected fluids or medicines.

\medterm{Hybridoma} A cell line formed by fusing an antibody-producing B lymphocyte with an immortal myeloma cell, allowing continuous production of a monoclonal antibody.

\medterm{Hycanthone} An older antiparasitic drug formerly used against schistosomiasis; its use was abandoned or restricted because of serious toxicity, including liver injury and carcinogenic concern.

\medterm{Hydatid Cyst} A cyst caused by the larval stage of Echinococcus granulosus, a tapeworm transmitted
from dogs to humans through ingestion of contaminated food. The liver is the most
commonly affected organ, followed by the lungs. Cysts grow slowly and may become
symptomatic from mass effect or rupture, causing anaphylaxis.

\medterm{Hydatid Disease} A zoonotic parasitic infection caused by the larval stage of tapeworms of the genus \textit{Echinococcus} (principally \textit{Echinococcus granulosus}). Humans acquire the infection by accidentally ingesting microscopic eggs shed in the feces of definitive canine hosts via contaminated food, water, soil, or direct animal contact. Larvae hatch in the small intestine, penetrate the bowel wall, and enter the circulation, most frequently lodging in the liver (approx.\ 70\%) and lungs (approx.\ 20\%), where they form slow-growing, fluid-filled cysts containing daughter cysts and hydatid sand. Most cysts remain asymptomatic for years until local mass effect causes abdominal discomfort, hepatomegaly, cough, or hemoptysis. Spontaneous or traumatic cyst rupture can cause severe anaphylaxis and secondary dissemination. Diagnosis is established by imaging (ultrasound, CT, or MRI showing cyst wall calcification, daughter cysts, or detached internal membranes) and confirmatory echinococcal serology. Management includes antiparasitic therapy (albendazole), image-guided aspiration and sterilization (PAIR: Puncture, Aspiration, Injection, Re-aspiration), or surgical cyst excision.
\synonyms{Echinococcosis; Cystic echinococcosis; Hydatidosis}

\medterm{Hydatidiform Mole} A form of gestational trophoblastic disease in which abnormal trophoblastic tissue grows in
the uterus. A complete mole has no viable fetus; a partial mole may contain abnormal fetal tissue. Both require follow-up
because persistent disease can occur.

\medterm{Hydralazine} Hydralazine is a direct arteriolar vasodilator used for selected hypertension, including severe hypertension in pregnancy; reflex tachycardia, fluid retention, headache, and drug-induced lupus are recognized adverse effects.

\medterm{Hydramnios} Excess amniotic fluid during pregnancy, also called polyhydramnios, caused by fetal swallowing impairment, maternal diabetes, multiple gestation, fetal anomalies, infection, or sometimes no identifiable cause.

\medterm{Hydranencephaly} A rare brain development disorder in which most of the upper brain is replaced by fluid-filled spaces. It can cause severe developmental and movement problems, seizures, feeding difficulty, and a shortened life expectancy.

\medterm{Hydrarthrosis} Abnormal accumulation of clear synovial fluid within a joint, producing swelling and restricted movement; causes include osteoarthritis, inflammatory arthritis, trauma, and infection.

\medterm{Hydration} The body’s water balance and the process of replacing water and electrolytes lost through urine, stool, sweat, breathing, vomiting, or bleeding.

\medterm{Hydrazine} A highly reactive chemical used in fuels and industry that can cause irritant, neurologic, hepatic, and systemic toxicity after significant exposure.

\medterm{Hydrazine Sulfate} A hydrazine-containing chemical once promoted experimentally for cancer-related symptoms; evidence of benefit is inadequate and toxicity can include liver injury, neuropathy, and seizures.

\medterm{Hydrazines} A family of nitrogen-containing chemical compounds used in industry, agriculture, and pharmaceuticals; some are irritants, toxic, mutagenic, or hepatotoxic.

\medterm{Hydrazones} Chemical compounds formed by reaction of hydrazine derivatives with aldehydes or ketones, important in medicinal chemistry, analytical chemistry, and some drug structures.

\medterm{Hydroa Vacciniforme} A rare photodermatosis causing recurrent vesicles or papules on sun-exposed skin that heal with varioliform scars; some cases are associated with Epstein-Barr virus.

\medterm{Hydrocarbon} An organic compound composed only of carbon and hydrogen; ingestion or inhalation of petroleum hydrocarbons can cause chemical pneumonitis and systemic toxicity.

\medterm{Hydrocarbon Pneumonia} Lung inflammation caused by inhaling or aspirating petroleum products or other hydrocarbons, which can produce cough, breathing difficulty, and chemical pneumonitis.

\medterm{Hydrocele} A collection of serous fluid within the tunica vaginalis surrounding the testis.
Congenital hydrocele results from patent processus vaginalis; acquired hydrocele may
follow infection, trauma, or tumor. Transilluminates on examination. Ultrasound confirms
fluid collection.

\medterm{Hydrocele Repair} Surgery to close or remove the fluid-containing sac around a testicle when a hydrocele is persistent, large, painful, associated with a hernia, or obscures examination.

\medterm{Hydrocephalus} Hydrocephalus is an abnormal buildup of cerebrospinal fluid (CSF) in the brain’s fluid-filled spaces, called ventricles. It can enlarge the ventricles and sometimes increase pressure on brain tissue. It may result from blocked CSF flow, reduced absorption, bleeding, infection, a tumor, injury, or a developmental condition. Symptoms vary with age and may include headache, vomiting, sleepiness, vision problems, poor balance, walking difficulty, memory changes, or an increasing head size in infants.

\medterm{Hydrocephalus Acquired} Acquired hydrocephalus is abnormal accumulation of cerebrospinal fluid after birth due to obstruction, impaired absorption, or rarely excess production; raised intracranial pressure may cause headache, vomiting, gait or cognitive change, and visual symptoms.

\medterm{Hydrocephalus Ex-Vacuo} Enlargement of the brain's ventricles secondary to loss or shrinkage of surrounding brain tissue, rather than primarily from obstructed cerebrospinal-fluid flow.

\medterm{Hydrocephaly} Hydrocephalus: abnormal accumulation or circulation of cerebrospinal fluid that enlarges the ventricles and may increase pressure on the brain.

\medterm{Hydrochloric Acid} The strong acid in stomach juice that helps digest proteins and kill many swallowed microorganisms. The stomach’s mucus lining and other protective mechanisms normally prevent this acid from damaging its tissues.

\medterm{Hydrochloric Acid Poisoning} A corrosive injury after swallowing, inhaling, or contacting hydrochloric acid. It can burn the mouth, throat, esophagus, stomach, skin, or lungs; do not induce vomiting and seek emergency care immediately.

\medterm{Hydrochlorothiazide} Hydrochlorothiazide is a thiazide diuretic that inhibits sodium-chloride reabsorption in the distal convoluted tubule and treats hypertension and mild oedema; hypokalaemia, hyponatraemia, hyperuricaemia, and photosensitivity may occur.

\medterm{Hydrocodone} An opioid analgesic used for moderate to severe pain and, in some formulations, cough; it can cause respiratory depression, sedation, constipation, tolerance, dependence, and overdose.

\medterm{Hydrocodone And Acetaminophen Overdose} Taking too much of this opioid-and-pain-reliever combination can slow breathing and cause coma, while acetaminophen can severely damage the liver. Emergency treatment is needed even before symptoms appear.

\medterm{Hydrocodone Bitartrate} The bitartrate salt of hydrocodone, an opioid analgesic that can cause sedation, constipation, respiratory depression, tolerance, dependence, and overdose.

\medterm{Hydrocodone/Oxycodone Overdose} Taking too much hydrocodone or oxycodone can cause extreme sleepiness, pinpoint pupils, slow or stopped breathing, coma, and death. Call emergency services, give naloxone if available, and provide rescue breathing if trained.

\medterm{Hydrocolpos} Distension of the vagina with fluid due to obstruction of the vaginal outlet, often from
an imperforate hymen or transverse septum, presenting as a pelvic mass in neonates or
adolescents.

\medterm{Hydrocortisone} A corticosteroid identical or similar to the body’s cortisol. It reduces inflammation and immune activity and treats adrenal hormone deficiency; prolonged or high-dose use can cause infection, high blood sugar, skin thinning, or adrenal suppression.

\medterm{Hydrocyanic Acid Poisoning} A life-threatening poisoning that prevents cells from using oxygen. It can cause sudden headache, confusion, seizures, breathing failure, cardiovascular collapse, and death; emergency antidote treatment is required immediately.

\medterm{Hydroflumethiazide} An older thiazide diuretic that increases urinary salt and water loss to lower blood pressure or reduce swelling. It can disturb sodium, potassium, and glucose levels and may cause dehydration or increased uric acid.

\medterm{Hydrofluoric Acid Exposure} Contact with hydrofluoric acid, which can cause deep tissue injury and systemic fluoride
toxicity, including dangerous low calcium and cardiac arrhythmias.

\medterm{Hydrofluoric Acid Poisoning} A medical emergency after contact with or inhalation of hydrofluoric acid. It can cause deep tissue burns and dangerous low calcium, heart rhythm changes, or death, even when the skin injury looks minor.

\medterm{Hydrogel} A water-rich three-dimensional material that can hold and release substances. Hydrogels are used in wound dressings, contact lenses, drug delivery, tissue engineering, and research, with safety depending on composition and use.

\medterm{Hydrogen Bond} A weak attraction between a hydrogen atom attached to an electronegative atom and another electronegative atom. These bonds help maintain the shapes of proteins and DNA and influence water’s properties.

\medterm{Hydrogen Breath Test} A test for some causes of bloating, diarrhea, or abdominal pain. After drinking a measured sugar solution, the person breathes into collection tubes at intervals. The gases can suggest poor sugar absorption or excess bacteria in the small intestine.

\medterm{Hydrogen Cyanide} A rapidly acting cellular toxin that inhibits cytochrome oxidase and prevents oxidative phosphorylation; significant exposure can cause severe lactic acidosis, seizures, cardiovascular collapse, and death.

\medterm{Hydrogen Peroxide} An oxidizing chemical used in some antiseptic and industrial products; ingestion or concentrated exposure can injure tissue and release gas, causing embolic or respiratory complications.

\medterm{Hydrogen Peroxide Poisoning} Harm after swallowing or inhaling hydrogen peroxide. It may irritate or burn tissues and release gas that blocks blood vessels; concentrated products can cause severe injury, so urgent poison-control or emergency advice is required.

\medterm{Hydrogen Sulfide} A toxic, flammable gas with a rotten-egg odor at low concentrations that inhibits cellular respiration and can rapidly cause respiratory failure, neurologic effects, or death.

\medterm{Hydrogenation} A chemical process that adds hydrogen to a compound; in foods it can convert liquid oils into more solid fats.

\medterm{Hydrogen-Ion Concentration} The concentration of hydrogen ions in a solution, expressed through pH and central to acid--base physiology, enzyme activity, and interpretation of blood-gas results.

\medterm{Hydromorphone} Hydromorphone is a potent opioid analgesic for severe pain; sedation, constipation, respiratory depression, tolerance, and dependence are important risks, especially with alcohol or other sedatives.

\medterm{Hydromorphone Overdose} Taking too much hydromorphone, an opioid, can cause extreme sleepiness, pinpoint pupils, slow or stopped breathing, coma, and death. Call emergency services and give naloxone if available; urgent monitoring is essential.

\medterm{Hydromyelia} Abnormal widening of the central canal of the spinal cord, which may be congenital or associated with obstruction of cerebrospinal-fluid flow and neurologic symptoms.

\medterm{Hydronephrosis} Swelling of a kidney because urine cannot drain normally. Causes include a stone, narrowing, reflux, a tumor, or prostate enlargement. Prolonged or severe blockage can damage the kidney and may cause pain, infection, or reduced kidney function.

\medterm{Hydronephrosis In Newborns} Swelling of a newborn’s kidney caused by widened urine-collecting spaces. It may result from temporary changes, blockage, or urine flowing backward and can affect one or both kidneys.

\medterm{Hydronephrosis Of One Kidney} Dilation of one kidney’s renal pelvis and calyces caused by impaired urine drainage, often from a stone, stricture, reflux, or obstruction of the ureter or ureteropelvic junction.

\medterm{Hydrophobia} An involuntary, agonizing spasm of the throat and vocal cord muscles triggered by attempting to drink or by the sight, sound, or mention of water. It is the hallmark neurological sign of furious (encephalitic) rabies, caused by viral damage to brainstem swallowing centers. The severe choking and pain make swallowing impossible, leading to a terrified avoidance of liquids, excessive drooling (hypersalivation), and often aerophobia (spasms triggered by a draft of air). Once hydrophobia appears, rabies is nearly 100\% fatal, highlighting the critical need for immediate post-exposure rabies vaccination following animal bites.
\synonyms{Rabic Hydrophobia; Rabid Pharyngeal Spasms; Fear of Water in Rabies}

\medterm{Hydropneumothorax} Air and fluid collected together in the space around a lung. It can compress the lung and cause chest pain or severe breathing difficulty, often requiring urgent hospital treatment and drainage.

\medterm{Hydrops} Abnormal accumulation of fluid in tissues or body cavities. The term may refer to generalized fetal edema, a localized collection, or organ-specific fluid distension; the cause depends on the clinical setting.

\medterm{Hydrops Fetalis} Severe fluid buildup in an unborn baby, affecting at least two areas such as under the skin, around the lungs, around the heart, or in the abdomen. Causes include blood-group incompatibility, anemia, infection, genetic conditions, heart problems, and complications of twin pregnancy.

\medterm{Hydrops Of The Gallbladder} Distension of the gallbladder by clear mucus or watery fluid, usually caused by prolonged obstruction of the cystic duct by a gallstone. It may cause right-upper-quadrant pain and can progress to inflammation or infection.

\medterm{Hydrosalpinx} Dilatation of a fallopian tube caused by accumulation of clear fluid, usually after inflammation, infection, endometriosis, or previous pelvic surgery. It may affect fertility.

\medterm{Hydrotherapy} Exercise or treatment performed in water. Water supports body weight and provides gentle resistance, which may reduce pain and help movement, strength, and balance during rehabilitation. The program should match the person’s condition and safety needs.

\medterm{Hydrothorax} An abnormal collection of clear fluid in the pleural space around a lung. It may result from heart, liver, kidney, or other disease and can cause breathlessness, cough, or chest pressure.

\medterm{Hydroureter} Abnormal enlargement of a ureter caused by impaired urine drainage.

\medterm{Hydroxamic Acid} An organic functional group containing hydroxylamine attached to a carbonyl, found in chelators and some medicines, including histone-deacetylase inhibitors.

\medterm{Hydroxocobalamin} A cobalt-containing form of vitamin B12 used as an antidote for cyanide poisoning. It
binds cyanide to form cyanocobalamin, which is excreted in urine; it can cause red
discoloration of skin and body fluids and interfere with some laboratory tests.

\medterm{Hydroxyapatite} A calcium-and-phosphate mineral that gives much of their hardness to bones and teeth. Bone cells continually remove and replace it as bone grows, repairs itself, and responds to changing physical forces.

\medterm{Hydroxychloroquine} Hydroxychloroquine is an antimalarial and immunomodulatory drug used for malaria prevention or treatment and selected autoimmune diseases; cumulative retinal toxicity requires ophthalmic monitoring.

\medterm{Hydroxycholesterols} Oxidized forms of cholesterol produced in the body or during food processing. They participate in cell signaling and cholesterol regulation and may contribute to inflammation or blood-vessel disease when present in excess.

\medterm{Hydroxylation} An enzymatic or chemical reaction that adds a hydroxyl group to a molecule, often altering the activity, solubility, metabolism, or excretion of a drug or biomolecule.

\medterm{Hydroxyproline} An amino acid formed mainly by post-translational modification of collagen; its measurement has been used as an indicator of collagen turnover but is not a routine diagnostic test.

\medterm{Hydroxyprolinemia} A rare inherited metabolic finding caused by impaired hydroxyproline breakdown, leading to elevated hydroxyproline in blood or urine and usually little clinical illness.

\medterm{Hydroxysteroid} A steroid hormone or steroid-related compound containing a hydroxyl group; hydroxysteroids include intermediates and products of adrenal and gonadal steroid metabolism.

\medterm{Hydroxyurea} A medicine used for some blood cancers and sickle-cell disease. In sickle-cell disease it can reduce painful episodes and other complications by helping red blood cells remain less likely to sickle; blood counts and other safety tests are required.

\medterm{Hydroxyzine} Hydroxyzine is a sedating H1-antihistamine used for itching, anxiety, and preoperative sedation; drowsiness, anticholinergic effects, and QT prolongation can occur.

\medterm{Hydroxyzine Overdose} Taking too much hydroxyzine can cause severe drowsiness, confusion, agitation, seizures, low blood pressure, or dangerous heart-rhythm changes. Seek urgent poison-control or emergency advice, especially after a large or intentional dose.

\medterm{Hydrozoa} A group of mostly marine animals that includes hydroids and some jellyfish. Contact with certain species can cause painful stings or skin irritation; the medical risk depends on the species and exposure.

\medterm{Hygiene} Practices that reduce exposure to pathogens and support health, including hand hygiene, oral care, skin care, food safety, and appropriate environmental cleanliness.

\medterm{Hygiene Hypothesis} The theory that limited early-life microbial exposure may increase susceptibility to allergies and immune disorders.

\medterm{Hygroma} A fluid-filled swelling, often caused by leakage or blockage of lymphatic vessels. A cystic hygroma usually appears in the neck of a fetus or child, while other hygromas may occur near joints or after injury.

\medterm{Hymen} A thin fold of mucous membrane that partially covers the external opening of the vagina.
The hymen is a remnant of fetal development, with no known physiological function in
postnatal life.

\medterm{Hymenectomy} Hymenectomy is a procedure to remove or open part of the hymen, usually to treat an obstructed hymenal opening or another symptomatic abnormality.

\medterm{Hymenolepiasis} Intestinal infection with dwarf tapeworms, especially Hymenolepis nana, acquired by ingesting eggs; it may cause abdominal symptoms, diarrhea, or be asymptomatic.

\medterm{Hymenolepis Diminuta} A tapeworm that can infect people who accidentally swallow infected insects, usually in contaminated food. Infection may cause abdominal discomfort, diarrhea, or no symptoms and is treated with an antiparasitic medicine.

\medterm{Hymenolepis Nana} A small tapeworm that spreads when microscopic eggs are swallowed, often through contaminated hands, food, or water. It can cause abdominal pain, diarrhea, poor appetite, or no symptoms and is treatable with antiparasitic medicine.

\medterm{Hymenoptera} An insect order that includes bees, wasps, hornets, and ants; stings or bites can cause pain, allergy, or life-threatening anaphylaxis.

\medterm{Hyoid Bone} U-shaped bone in the anterior neck at C3, unique for not directly articulating with any
other bone. Suspended by muscles and the stylohyoid ligament. Provides attachment for
suprahyoid (digastric, stylohyoid, mylohyoid, geniohyoid) and infrahyoid (sternohyoid,
sternothyroid, thyrohyoid, omohyoid) muscles.

\medterm{Hyoscine} A medicine that blocks muscarinic acetylcholine receptors. Depending on the preparation, it can prevent motion sickness, reduce spasms, or lessen secretions; dry mouth, blurred vision, constipation, confusion, and urinary retention may occur.

\medterm{Hyoscyamine} An antimuscarinic medicine that reduces smooth-muscle spasm and secretions; it may be used for selected gastrointestinal or urinary symptoms and can cause dry mouth, blurred vision, urinary retention, or confusion.

\medterm{Hyperactivity} Excessive or developmentally inappropriate motor activity, restlessness, or difficulty remaining still; persistent impairment may occur in ADHD, mania, substance effects, or other conditions.

\medterm{Hyperacusis} Reduced tolerance to ordinary environmental sounds, which may be perceived as excessively loud, painful, or distressing. It can occur with hearing disorders, migraine, head injury, neurologic conditions, or without an identifiable cause.

\medterm{Hyperadrenocorticism} Excess production or exposure to adrenal glucocorticoids, especially cortisol; in humans it causes Cushing syndrome, which may result from adrenal, pituitary, ectopic, or medication-related causes.

\medterm{Hyperaemia} Increased blood flow to a tissue, making it look red or feel warm. It may be a normal response to activity or a response to inflammation or renewed blood flow after temporary blockage.

\medterm{Hyperaesthesia} Unusually strong sensitivity to touch, pain, temperature, or other sensations. It may occur with nerve injury, inflammation, migraine, medicines, or disorders affecting the brain or spinal cord.

\medterm{Hyperaldosteronism} Excess production of the adrenal hormone aldosterone. It can cause high blood pressure, low potassium, muscle weakness, cramps, thirst, and frequent urination; an adrenal growth or overactivity of both glands may be responsible.

\medterm{Hyperaldosteronism Complications} Beyond hypertension and hypokalemia, primary hyperaldosteronism causes target organ
damage including left ventricular hypertrophy, fibrosis, stroke, coronary artery
disease, atrial fibrillation, and chronic kidney disease. These complications occur at
higher rates than in essential hypertension of similar severity.

\medterm{Hyperalgesia} Feeling more pain than expected from something that normally hurts. It can happen with inflammation, nerve injury, changes in nervous-system pain processing, or opioid use; it differs from allodynia, which is pain from a normally painless touch.

\medterm{Hyperammonemia} An abnormally high level of ammonia in the blood. It may result from severe liver disease, an inherited metabolic disorder, medicines, or other illness and can cause confusion, sleepiness, vomiting, seizures, coma, or brain swelling.

\medterm{Hyperamylasemia} An elevated blood amylase concentration caused by pancreatic or salivary-gland disease, reduced clearance, macroamylasemia, or other conditions; it is not diagnostic of pancreatitis by itself.

\medterm{Hyperargininemia} An inherited urea-cycle disorder caused by arginase deficiency, leading to high arginine and progressive spasticity, developmental impairment, seizures, or liver-related findings.

\medterm{Hyperbaric} Having a pressure greater than the surrounding atmospheric pressure; hyperbaric oxygen therapy delivers oxygen at increased pressure for selected medical indications.

\medterm{Hyperbaric Chamber} A pressurized chamber in which a person breathes oxygen or another gas at pressure above sea level, used for conditions such as decompression sickness, air embolism, and selected wounds.

\medterm{Hyperbaric Oxygen Therapy} Breathing oxygen inside a chamber pressurized above normal air pressure. The treatment increases oxygen dissolved in the blood and is used for selected problems such as carbon-monoxide poisoning, decompression sickness, and certain severe wounds or infections.

\medterm{Hyperbilirubinaemia} Hyperbilirubinaemia is an abnormally high level of bilirubin in the blood, which can cause jaundice and may result from increased breakdown of red cells, liver disease, or blocked bile flow.

\medterm{Hyperbilirubinemia} An abnormally high concentration of bilirubin in blood that may cause jaundice; causes include hemolysis, impaired liver processing, hepatocellular injury, and bile-duct obstruction.

\medterm{Hypercalcaemia} Alternate British spelling; see \textit{Hypercalcemia}.

\medterm{Hypercalcemia} An abnormally high level of calcium in the blood (total calcium greater than $10.5\,\text{mg/dL}$ or $2.6\,\text{mmol/L}$). More than 90\% of cases are caused by primary hyperparathyroidism (usually a benign parathyroid adenoma) or malignancy. Hallmark symptoms are classically summarized as ``stones, bones, abdominal groans, and psychiatric moans'': kidney stones and frequent urination; bone pain and osteoporosis; constipation, nausea, and pancreatitis; and fatigue, muscle weakness, confusion, and depression. ECG characteristically shows a shortened QT interval. Severe hypercalcemia is an emergency managed with intensive intravenous saline hydration, calcitonin, and bisphosphonates to inhibit bone resorption.
\synonyms{High Blood Calcium; Hypercalcaemia; Elevated Serum Calcium}

\medterm{Hypercalcemia Of Malignancy} A high blood-calcium level caused by cancer. It may occur because the cancer releases substances that raise calcium or damages bone. Severe cases can cause thirst, vomiting, weakness, confusion, kidney injury, or abnormal heart rhythms and require urgent treatment.

\medterm{Hypercalciuria} Excessive urinary calcium excretion, which can promote calcium kidney stones and, in some settings, bone mineral loss; evaluation considers diet, medications, kidney function, and endocrine causes.

\medterm{Hypercapnia} An abnormally high level of carbon dioxide in the blood, usually because breathing is not removing enough of it. It may cause headache, sleepiness, confusion, flushed skin, tremor, or coma and can occur with severe lung or muscle disease.

\medterm{Hypercementosis} Excessive deposition of cementum around a tooth root, producing a bulbous root contour that can complicate extraction but is often asymptomatic.

\medterm{Hypercholesterolaemia} An abnormally high level of cholesterol in the blood, especially low-density lipoprotein cholesterol. Over time, it can contribute to fatty buildup in arteries and increase the risk of heart attack or stroke.

\medterm{Hypercholesterolemia} A high level of cholesterol, especially LDL cholesterol, in the blood. It usually causes no symptoms but can gradually narrow arteries and raise the risk of heart attack and stroke; causes may be inherited or related to other diseases, medicines, or diet.

\medterm{Hypercoagulable State} A condition in which blood clots more easily than usual, increasing the risk of deep-vein thrombosis, pulmonary embolism, stroke, or other clots. Causes may be inherited or acquired, including cancer, pregnancy, some medicines, inflammation, or prolonged immobility.

\medterm{Hypercyanotic Tetralogy} A sudden episode of worsening blue or gray color, rapid breathing, distress, or faintness in a child with Tetralogy of Fallot. It results from reduced blood flow to the lungs and is an emergency requiring immediate medical treatment.

\medterm{Hyperelastic Skin} Skin that stretches farther than usual and recoils abnormally; generalized hyperextensibility may occur in connective-tissue disorders such as Ehlers-Danlos syndromes.

\medterm{Hyperemesis} Excessive or persistent vomiting that can cause dehydration, abnormal body salts, weight loss, and poor nutrition. It may result from pregnancy, infection, medicines, digestive disease, brain illness, or other causes.

\medterm{Hyperemesis Gravidarum} Severe nausea and vomiting during pregnancy that causes dehydration, weight loss, or abnormal body salts and prevents adequate food or fluid intake. It often begins early in pregnancy and may require fluids, nutrition, medicines, and hospital care.

\medterm{Hyperemia} Increased blood flow to a tissue, causing redness and warmth; it may be an active physiologic response to exercise or inflammation or a reactive response after temporary ischemia.

\medterm{Hyperemia Grade} A graded clinical or endoscopic estimate of tissue redness and vascular congestion, used to describe the severity of inflammation or mucosal irritation.

\medterm{Hyperesthesia} Increased sensitivity to sensory stimulation—touch (tactile), pain (hyperalgesia), or
other stimuli—so that normally non-painful or mildly painful stimuli produce exaggerated
discomfort.

\medterm{Hyperextension} Movement that extends a joint beyond its usual neutral or anatomical range, potentially stretching or tearing ligaments, capsules, muscles, or nerves.

\medterm{Hyperfibrinogenemia} An abnormally high level of fibrinogen, a blood protein needed for clotting. It often rises with inflammation, infection, pregnancy, obesity, or estrogen exposure and may be associated with increased clotting risk.

\medterm{Hyperflexion} Excessive bending of a joint beyond its normal range, which can injure ligaments, discs, muscles, or other supporting structures.

\medterm{Hypergammaglobulinemia} An excess of antibody proteins in the blood. It may result from long-term infection, inflammation, liver disease, or an abnormal immune-cell growth; testing determines whether many antibody types or one type is increased.

\medterm{Hyperglycaemia} An abnormally high level of glucose in the blood. It may cause thirst, frequent urination, tiredness, or blurred vision and, if severe or prolonged, can damage organs or cause an emergency.

\medterm{Hyperglycemia} A higher-than-normal level of glucose in the blood. Severe short-term elevation can cause excessive urination, dehydration, and dangerous changes in body chemistry. Long-term elevation can damage the eyes, kidneys, nerves, heart, and blood vessels.

\medterm{Hyperglycinemia} An abnormally high glycine concentration in blood, most notably in nonketotic hyperglycinemia, a disorder of glycine cleavage that can cause neonatal hypotonia, seizures, and developmental impairment.

\medterm{Hypergonadotropic Hypogonadism} Reduced testicular or ovarian hormone production caused by gonad failure, with increased pituitary signals and possible infertility, absent periods, or delayed sexual development.

\medterm{Hypergraphia} Excessive, compulsive, or unusually prolonged writing or drawing, sometimes associated with temporal-lobe epilepsy, mania, or other neurologic or psychiatric conditions.

\medterm{Hyperhidrosis} Sweating that is excessive for the temperature or activity. It may mainly affect the palms, soles, underarms, or face without another illness, or may be caused by medicines, hormone disorders, infection, or cancer.

\medterm{Hyperhomocysteinemia} An elevated blood homocysteine concentration associated with genetic enzyme defects, folate or vitamin B12 deficiency, kidney disease, medicines, and increased vascular or pregnancy-related risk in selected settings.

\medterm{Hypericum Perforatum} St. John’s wort, a plant used in some complementary products for depressive symptoms; it can cause drug interactions by inducing hepatic enzymes and transporters and should not be combined casually with medicines.

\medterm{Hyper-IgE Syndrome} A rare inherited immune disorder marked by very high levels of immunoglobulin E, eczema, and repeated skin or lung infections. Some forms also cause unusual bones, teeth, or blood vessels.

\medterm{Hyper-IgM Syndrome} See \textit{Hyper-IgM Syndromes}.

\medterm{Hyper-IgM Syndromes} Inherited immune disorders in which the body makes normal or high amounts of IgM antibodies but too little IgG, IgA, or IgE. People may develop repeated bacterial, viral, or fungal infections because antibody defenses do not mature normally.

\medterm{Hyperimmunization} Repeated or intensive exposure to an antigen, such as through vaccination or laboratory immunization, intended to produce a high antibody level. The term does not by itself describe an illness or harmful “over-vaccinated” state.

\medterm{Hyperimmunoglobulin D Syndrome} A hereditary autoinflammatory disorder, usually caused by mevalonate kinase deficiency, characterized by recurrent fever, enlarged lymph nodes, abdominal symptoms, rash, and joint or muscle pain; it is also called HIDS or MKD.

\medterm{Hyperimmunoglobulin E Syndrome} A rare primary immunodeficiency characterized by very high IgE, eczema, recurrent skin and lung infections, and sometimes skeletal or dental abnormalities.

\medterm{Hyperinsulinemia} An abnormally high insulin level in the blood, often caused by insulin resistance, excess insulin medicine, or rarely an insulin-producing tumor. It may occur with normal or low glucose and requires interpretation with glucose results.

\medterm{Hyperinsulinism} Excessive insulin secretion or action causing recurrent or persistent hypoglycemia; causes include congenital disorders, tumors, medications, and critical illness.

\medterm{Hyperkalaemia} Alternate British spelling; see \textit{Hyperkalemia}.

\medterm{Hyperkalemia} A dangerous elevation of potassium in the blood (greater than $5.0\,\text{mEq/L}$ or $\text{mmol/L}$), representing a medical emergency due to its disruptive effect on cardiac electrical conduction. Most commonly caused by acute or chronic kidney failure, potassium-sparing medications (such as ACE inhibitors, ARBs, or spironolactone), cellular destruction (burns, rhabdomyolysis, tumor lysis), or acidosis. While mild cases may cause muscle weakness, fatigue, or numbness, severe elevation causes progressive ECG abnormalities: tall peaked T waves, prolonged PR interval, widened QRS complex, sine-wave pattern, and fatal cardiac arrest. Emergency treatment requires intravenous calcium gluconate to stabilize the heart muscle, insulin with glucose to shift potassium into cells, and diuretics, binding agents, or urgent hemodialysis to remove potassium from the body.
\synonyms{High Blood Potassium; Hyperkalaemia; Elevated Serum Potassium}

\medterm{Hyperkalemic Periodic Paralysis} An inherited muscle disorder causing repeated attacks of weakness or temporary paralysis. Attacks may follow rest after exercise, stress, or potassium-rich food and usually improve between episodes.

\medterm{Hyperkeratosis} Thickening of the outer keratin layer of skin or mucosa, producing callus, scale, or rough plaques in response to friction, inflammation, infection, inherited disorders, or other causes.

\medterm{Hyperkinesia} Excessive or involuntary movement, such as chorea, dystonia, tics, or tremor, resulting from abnormal activity in motor-control circuits or from drugs and other diseases.

\medterm{Hyperlipasaemia} An abnormally high blood lipase concentration, often associated with pancreatic inflammation but also seen with kidney failure, gastrointestinal disease, medications, and other conditions.

\medterm{Hyperlipidaemia} An abnormally high level of fats in the blood, including cholesterol or triglycerides. It may have no symptoms but can increase the risk of artery disease, heart attack, stroke, or pancreatitis.

\medterm{Hyperlipidemia} Too much cholesterol or triglyceride fat in the blood. It usually causes no symptoms but can build up in artery walls, increasing the risk of heart attack and stroke; inherited disorders and other illnesses, medicines, diet, or alcohol can contribute.

\medterm{Hyperlipoproteinemia} An excess of one or more lipoproteins in blood, producing high cholesterol or triglycerides and increasing risk of atherosclerosis, pancreatitis, or inherited lipid disease.

\medterm{Hyperlordosis} Exaggerated inward curvature of the lumbar spine, often producing a swayback appearance and sometimes associated with pain, muscular imbalance, obesity, pregnancy, or structural spinal disease.

\medterm{Hyperlordotic Posture} An exaggerated inward curvature of the lumbar spine, often producing increased anterior pelvic tilt and a prominent lower-back arch.

\medterm{Hyperlysinemia} An inherited metabolic abnormality with elevated lysine in blood caused by altered lysine degradation; many people have few symptoms, but some have developmental or neurologic findings.

\medterm{Hypermagnesemia} An abnormally high concentration of magnesium in the blood, often related to kidney failure or excessive magnesium intake; severe cases can cause weakness, low blood pressure, depressed reflexes, or arrhythmias.

\medterm{Hypermetropia} A common refractive error of the eye in which parallel rays of light enter the relaxed eye and focus behind the retina rather than directly upon it. This occurs because the eyeball is anatomically too short along its anteroposterior axis (axial hypermetropia) or because the refractive power of the cornea or crystalline lens is insufficient (curvature hypermetropia). While distant objects may remain relatively clear in mild cases, viewing close objects requires continuous ciliary muscle accommodation, leading to accommodative asthenopia (eyestrain, frontal headache, and ocular fatigue during close work). In young children, significant uncorrected hypermetropia can cause accommodative esotropia (inward squint) and amblyopia. Optical correction is achieved using convex (plus/converging) spherical lenses in spectacles or contact lenses, or via corneal refractive surgery.
\synonyms{Hyperopia; Farsightedness; Long-sightedness}

\medterm{Hypermobile Ehlers-Danlos Syndrome} A heritable connective-tissue disorder characterized mainly by generalized joint hypermobility, instability,
chronic musculoskeletal pain, and other variable features involving skin, autonomic function, or digestion.
Diagnosis is clinical and requires assessment for other connective-tissue disorders and related conditions.

\medterm{Hypermobile Joints} Joints that move farther than usual. This may cause no problems, or may lead to pain, repeated dislocations, sprains, fatigue, or joint instability.

\medterm{Hypermobility} Greater-than-usual movement of one or more joints. It may be a harmless trait, but when associated with pain, instability, dislocations, fatigue, or other features it may be part of a hypermobility spectrum disorder.

\medterm{Hypermobility Spectrum Disorder} A condition involving joint hypermobility together with symptoms such as pain, joint instability, soft-tissue injury, or fatigue, without meeting the criteria for another defined connective-tissue disorder.

\medterm{Hypermobility Syndrome} Excessive joint movement associated with pain, instability, fatigue, or other symptoms; evaluation distinguishes generalized joint hypermobility from a defined hypermobility spectrum disorder or hypermobile Ehlers-Danlos syndrome.

\medterm{Hypermotility} Abnormally rapid or excessive movement of an organ, especially the digestive tract. It may cause cramping, urgency, diarrhea, or poor absorption and can result from infection, medicines, inflammation, or other disorders.

\medterm{Hypernatraemia} An abnormally high level of sodium in the blood, usually because the body has too little water. It may cause thirst, weakness, confusion, seizures, or coma and requires careful correction.

\medterm{Hypernatremia} A high blood-sodium level, usually because the body has lost too much water or cannot obtain enough. It can cause intense thirst, weakness, confusion, seizures, or coma and must be corrected carefully.

\medterm{Hypernatremia Result} A high blood sodium result, usually reflecting too little free water relative to body sodium; severity and rate of onset determine neurologic risk and the urgency and speed of correction.

\medterm{Hyperopia} Alternate name; see \textit{Hypermetropia}.

\medterm{Hyperornithinemia} An elevated blood ornithine concentration caused by defects in ornithine metabolism or transport, including disorders that can produce retinal degeneration or neurologic disease.

\medterm{Hyperosmolar} Having more dissolved particles than another fluid. When blood becomes too concentrated, water moves out of cells and can cause dehydration, confusion, or other nervous-system problems.

\medterm{Hyperosmolar Hyperglycemic State} A life-threatening emergency in which very high blood glucose causes severe dehydration and concentrated body fluids, usually with little or no ketone buildup. It may cause extreme thirst, weakness, confusion, seizures, or coma and needs urgent hospital treatment.

\medterm{Hyperosmolar Non-Ketotic Coma} A life-threatening complication of very high blood glucose, usually in type 2 diabetes, causing severe dehydration and sometimes confusion, seizures, or coma. It requires emergency hospital treatment.

\medterm{Hyperostosis} Abnormal excessive formation or thickening of bone, which may be localized or widespread and may occur in skeletal dysplasias, inflammation, endocrine disease, or medication effects.

\medterm{Hyperoxaluria} An excess of oxalate in the urine, which can combine with calcium to form kidney stones or calcium deposits. Severe inherited forms can progressively damage the kidneys.

\medterm{Hyperoxia} Exposure to an abnormally high oxygen concentration or tissue oxygen level, which can cause oxygen toxicity, absorption atelectasis, retinal injury in premature infants, or lung injury depending on dose and duration.

\medterm{Hyperparakeratosis} Thickening of the outer skin layer in which cells retain their nuclei instead of fully maturing. It may produce rough scale and can result from irritation, inflammation, infection, or an inherited skin disorder.

\medterm{Hyperparathyroidism} Excessive secretion of parathyroid hormone leading to hypercalcemia. Classified as primary (parathyroid adenoma or hyperplasia), secondary (chronic kidney disease, vitamin D deficiency), or tertiary (autonomous PTH secretion after prolonged secondary hyperparathyroidism).

\medterm{Hyperpathia} An unusually strong and sometimes delayed pain response to a painful stimulus. It may occur after nerve injury or with disorders that change how the nervous system processes pain.

\medterm{Hyperphagia} An unusually strong or excessive appetite. It can occur with diabetes, some genetic or brain disorders, certain medicines, or mental health conditions.

\medterm{Hyperphosphatemia} A high phosphate level in the blood, most often caused by kidney disease. It can disturb calcium balance, weaken bones, and contribute to calcium deposits in blood vessels or soft tissues.

\medterm{Hyperpigmentation} Darkening of an area of skin because the body makes or stores more pigment there. Causes include sun exposure, inflammation, injury, medicines, hormone disorders, and some skin or internal diseases.

\medterm{Hyperpituitarism} Excess production of one or more hormones from the pituitary gland, usually because of a benign pituitary tumor. The effects depend on the hormone involved and may affect growth, menstrual cycles, fertility, metabolism, or stress responses.

\medterm{Hyperplasia} An increase in the number of cells in a tissue or organ, which may make it larger. It can be a normal response or result from disease, hormones, or other stimuli.

\medterm{Hyperplastic Polyp} A usually noncancerous overgrowth of lining cells, often in the colon or rectum, whose cancer risk depends on size, number, and location.

\medterm{Hyperpnea} Breathing that is deeper, and sometimes faster, than normal, increasing air movement and occurring with exercise, illness, metabolic changes, or anxiety.

\medterm{Hyperpnoea} Breathing that is deeper or faster than usual to meet increased demand for oxygen or remove extra carbon dioxide. It may occur during exercise, fever, anxiety, or lung, heart, or metabolic illness.

\medterm{Hyperpolarization} An electrical change that makes a nerve or muscle cell less likely to send a signal. It occurs when the inside of the cell becomes more electrically negative after certain ion channels open.

\medterm{Hyperprolactinaemia} An abnormally elevated concentration of prolactin in the blood. Causes include medicines, pregnancy, hypothyroidism, pituitary tumors, and other disorders.

\medterm{Hyperprolactinemia} An abnormally high level of the hormone prolactin in the blood. Causes include pregnancy, breastfeeding, medicines, an underactive thyroid, kidney disease, or a pituitary tumor. It may disrupt periods, fertility, sexual function, or milk production.

\medterm{Hyperproteinemia} An abnormally high total blood-protein concentration, usually from dehydration, chronic inflammation, or a monoclonal immunoglobulin disorder.

\medterm{Hyperpyrexia} An extreme elevation of body temperature, generally above 41.5 °C, that can result from severe infection, heatstroke, malignant hyperthermia, drug toxicity, or other life-threatening conditions.

\medterm{Hyperreactio Luteinalis} Bilateral ovarian enlargement with multiple theca-lutein cysts caused by an exaggerated response to human chorionic gonadotropin, usually during pregnancy or with trophoblastic disease.

\medterm{Hyperreflexia} Exaggerated or overactive deep-tendon (stretch) reflexes. It is a neurologic sign, not a diagnosis, and often reflects impaired descending upper-motor-neuron control; it may occur with spasticity, clonus, or an extensor plantar response. Distribution and associated findings help distinguish central nervous-system disease from systemic or toxic causes.

\medterm{Hypersensitivity} An exaggerated or inappropriate immune response to an antigen, ranging from a mild reaction to
anaphylaxis.

\medterm{Hypersensitivity Pneumonitis} Lung inflammation caused by repeated breathing of substances such as mold, bird proteins, or other workplace dusts. It may cause cough, breathlessness, fever, or fatigue and can lead to permanent lung scarring if exposure continues.

\medterm{Hypersensitivity Reactions} Inappropriate or excessive immune responses causing tissue damage or disease.

\medterm{Hypersensitivity Type I} An immediate allergic reaction in which immunoglobulin E antibodies activate mast cells after exposure to an allergen. Symptoms can include itching, hives, wheezing, vomiting, swelling, or life-threatening anaphylaxis.

\medterm{Hypersensitivity Type II} An immune reaction in which IgG or IgM antibodies attach to targets on cells or tissues and recruit other immune mechanisms that injure them. Examples include some drug reactions and autoimmune blood-cell destruction.

\medterm{Hypersensitivity Type III} An immune reaction in which clusters of antibodies and foreign or self proteins settle in tissues or blood vessels, causing inflammation and injury. It can affect the skin, joints, kidneys, or blood vessels.

\medterm{Hypersensitivity Type IV} A delayed immune reaction driven by T cells rather than antibodies. It usually develops 48 to 72 hours after exposure and can cause contact dermatitis, some drug rashes, or inflammation around an infection.

\medterm{Hypersensitivity Vasculitis} Inflammation of small blood vessels, usually in the skin, often triggered by an infection, medicine, or immune reaction. It commonly causes raised purple spots or bruising on the legs.

\medterm{Hypersomnia} Being very sleepy during the day or sleeping much longer than usual even after enough sleep. It can be caused by sleep disorders, medicines, depression, substance use, or other illness; it is different from feeling tired without wanting to sleep.

\medterm{Hypersplenism} A condition in which an enlarged or overactive spleen removes blood cells too quickly. It may lower red cells, white cells, platelets, or all three and can cause anemia, infections, or easy bleeding.

\medterm{Hypertelorism} An unusually wide distance between the eyes because the bony eye sockets are farther apart than usual. It may be present from birth or occur as part of a genetic or developmental condition.

\medterm{Hypertension} Medical term; see \textit{High blood pressure}.

\medterm{Hypertension In Anesthesia} Abnormally high blood pressure during or around an operation. Pain, inadequate anesthesia, low oxygen, high carbon dioxide, shivering, bladder fullness, or stopping certain medicines may contribute and increase the risk of bleeding, heart injury, or stroke.

\medterm{Hypertensive} Having blood pressure that remains above the accepted range for the person's age, health, and measurement setting, increasing risk of heart, brain, kidney, or vessel disease.

\medterm{Hypertensive Cerebellar Hemorrhage} Bleeding into the cerebellum caused by rupture of a vessel damaged by high blood pressure.

\medterm{Hypertensive Crisis} A dangerously high blood pressure that needs prompt assessment. If it is causing new organ injury, such as stroke, heart failure, kidney injury, or vision loss, it is an emergency requiring immediate hospital treatment.

\medterm{Hypertensive Emergency} A severe rise in blood pressure that is causing new injury to the brain, heart, kidneys, eyes, or major blood vessels. Symptoms may include chest pain, breathlessness, weakness, confusion, seizures, or vision loss and require immediate hospital care.

\medterm{Hypertensive Encephalopathy} Brain dysfunction caused by a sudden, severe rise in blood pressure. Headache, confusion, seizures, vision changes, or reduced consciousness may occur and require immediate blood-pressure treatment.

\medterm{Hypertensive Heart Disease} Left ventricular hypertrophy and diastolic/systolic dysfunction resulting from chronic
hypertension. LVH is an independent cardiovascular risk factor and precursor to heart
failure. Diagnosis by ECG (voltage criteria) or echocardiography (wall thickness >1.1
cm). Treatment of hypertension can regress LVH.

\medterm{Hypertensive Intracerebral Hemorrhage} Bleeding into brain tissue after a blood vessel weakened by high blood pressure breaks, potentially causing sudden headache, weakness, confusion, unconsciousness, or death.

\medterm{Hypertensive Nephropathy} Kidney damage caused or worsened by long-term high blood pressure. It may lead to protein in the urine and gradually reduced kidney function; diabetes, other kidney diseases, and inherited risk can cause similar findings and must be considered.

\medterm{Hypertensive Nephrosclerosis} Scarring and narrowing of small kidney blood vessels caused by long-term high blood pressure. It can gradually reduce kidney function and may cause protein in the urine; controlling blood pressure helps slow progression.

\medterm{Hypertensive Retinopathy} Damage to the blood vessels and light-sensitive tissue at the back of the eye caused by high blood pressure. It may include narrowed vessels, bleeding, swelling, or fluid leakage. Severe changes can threaten vision and may indicate dangerously high blood pressure.

\medterm{Hypertensive Urgency} A markedly high blood pressure without clear evidence of new organ injury. It still requires prompt medical review, especially when symptoms are present; treatment is adjusted carefully rather than lowered suddenly without supervision.

\medterm{Hyperthermia} Hyperthermia is an abnormally high body temperature caused by excessive heat production or impaired heat loss, without a change in the hypothalamic temperature set point. Causes include heat exposure, exertion, medicines, endocrine disorders, infection-related illness, and impaired sweating. Severe hyperthermia can cause confusion, seizures, organ injury, or death and requires urgent treatment.

\medterm{Hyperthermia For Treating Cancer} A cancer treatment that deliberately heats a tumor or part of the body, sometimes with radiation or medicines, to damage cancer cells or make other treatment work better.

\medterm{Hyperthermia In Anesthesia} An abnormally high temperature during or after anesthesia. Causes include infection, drug reactions, thyroid storm, serotonin syndrome, and malignant hyperthermia, a rare emergency with muscle rigidity, rising carbon dioxide, and dangerous temperature elevation.

\medterm{Hyperthermic Intraperitoneal Chemotherapy} Heated cancer medicine circulated through the abdominal cavity during surgery after visible tumor is removed. It is used for selected cancers that have spread across the abdominal lining and can cause infection, bleeding, organ injury, or other serious complications.

\medterm{Hyperthyroidism} A condition in which the thyroid makes too much thyroid hormone. It may cause weight loss, sweating, heat intolerance, shaking, anxiety, diarrhea, or a fast or irregular heartbeat. Common causes include Graves disease and overactive thyroid nodules.

\synonyms
Thyrotoxicosis; Overactive Thyroid

\medterm{Hyperthyroidism} A state in which excess thyroid hormone effect causes symptoms such as heat intolerance, sweating, tremor, weight loss, or rapid heartbeat. It may result from increased thyroid production or other causes of thyrotoxicosis.

\synonyms
Thyrotoxicosis; Thyroid Hyperfunction

\medterm{Hyperthyroxinemia} An abnormally high concentration of circulating thyroxine, caused by hyperthyroidism, thyroiditis, medication, binding-protein changes, or assay interference; symptoms and TSH clarify the cause.

\medterm{Hypertonia} Abnormally increased muscle tightness or resistance when a limb is moved. It may cause stiffness, abnormal posture, spasms, or limited movement and can result from brain, spinal-cord, or other nervous-system disorders.

\medterm{Hypertonic Saline} Saline containing more sodium than normal blood, commonly given through a vein for severe symptoms from dangerously low blood sodium or selected urgent conditions. Close monitoring is required because correcting sodium too quickly can cause serious brain injury.

\medterm{Hypertonic Saline Inhalation} Inhaling a salty mist through a nebulizer to draw water into airway mucus and make it easier to clear. It may help some people with cystic fibrosis or bronchiectasis, but can briefly narrow the airways and should be used as instructed.

\medterm{Hypertonic Solution} A solution with greater effective osmotic concentration than the fluid with which it is compared, causing water to move out of cells across a semipermeable membrane.

\medterm{Hypertrichosis} Excessive hair growth in areas not usually controlled by male hormones. It can be present from birth or develop later because of inherited conditions, medicines, poor nutrition, hormone or metabolic disease, or rarely cancer. The cause should be assessed when growth is new or widespread.

\medterm{Hypertriglyceridemia} An abnormally high level of triglycerides, a type of fat, in the blood. Causes include diabetes, obesity, alcohol, some medicines, and inherited disorders. Very high levels increase the risk of acute pancreatitis and require medical management.

\medterm{Hypertrophic} Relating to hypertrophy, an enlargement of cells or tissue that increases the size of an organ or structure.

\medterm{Hypertrophic Cardiomyopathy} A disease in which the heart muscle, especially the left ventricle, becomes
abnormally thick and stiff. It may obstruct blood flow, impair filling, cause chest pain or breathlessness, and increase
the risk of arrhythmia and sudden death.

\synonyms
HCM; Hypertrophic Obstructive Cardiomyopathy (HOCM)

\medterm{Hypertrophic Cardiomyopathy Genetics} Many cases result from an inherited change in a gene for a heart-muscle protein, commonly MYH7 or MYBPC3. The condition can run in families, so genetic counseling and screening of close relatives may be recommended.

\medterm{Hypertrophic Cardiomyopathy Treatment} Management may include medicines to relieve symptoms, treatment of
arrhythmias or obstruction, an implantable defibrillator for selected patients, and surgery or ablation when indicated.
Treatment is guided by symptoms, anatomy, risk, and specialist assessment.

\medterm{Hypertrophic Gastritis} Thickening of the gastric mucosal folds caused by glandular overgrowth or inflammation, as in Ménétrier disease; protein loss, pain, nausea, and increased cancer risk may occur.

\medterm{Hypertrophic Myocardiopathy} Abnormal thickening of the heart muscle, often inherited. The thickened muscle may become stiff, obstruct blood leaving the heart, cause chest pain or breathlessness, or trigger dangerous abnormal rhythms.

\medterm{Hypertrophic Pyloric Stenosis} Thickening of the muscle at the outlet of the stomach, usually in young infants. It causes increasingly forceful vomiting after feeds, continued hunger, dehydration, and poor weight gain and is usually corrected with surgery.

\medterm{Hypertrophic Scar} A raised, thick scar that stays within the boundaries of the original wound. It forms when the body
makes too much scar tissue during healing, often after burns, surgery, or injury. Unlike a keloid,
it does not grow beyond the wound edges and may become flatter over time.

\medterm{Hypertrophy} Enlargement of a tissue or organ because its individual cells become larger. It may develop in response to increased workload, hormones, disease, or other stimulation.

\medterm{Hypertropia} Vertical misalignment of the eyes in which one eye deviates upward relative to the other, causing double vision, eye strain, or abnormal head posture.

\medterm{Hyperuricemia} An abnormally high concentration of uric acid in the blood, which may be asymptomatic or contribute to gout and uric-acid kidney stones.

\medterm{Hyperventilating} Breathing faster or deeper than the body needs, lowering carbon dioxide in the blood. It may cause light-headedness, tingling, chest tightness, or fainting and can result from anxiety or serious illness.

\medterm{Hyperventilation} Breathing more deeply or rapidly than the body needs, causing carbon dioxide levels in the blood to fall. It may cause dizziness, tingling, chest tightness, or fainting and can result from anxiety or serious heart, lung, or metabolic illness.

\medterm{Hyperventilation Syndrome} Repeated overbreathing, often during stress or panic, that lowers carbon dioxide in the blood and causes breathlessness, chest tightness, dizziness, tingling, palpitations, or hand and foot spasms. Other heart and lung causes must be considered.

\medterm{Hypervigilance} A state of excessive alertness to possible threat, often accompanied by scanning, startle, irritability, and sleep difficulty in anxiety, trauma-related disorders, stimulant use, or medical illness.

\medterm{Hyperviscosity Syndrome} A set of symptoms and signs caused by blood that is abnormally thick, usually from very high levels of an immunoglobulin paraprotein or of red or white blood cells. It can complicate Waldenström macroglobulinemia, multiple myeloma, polycythemia, leukemia, or severe infection. Features include bleeding from the nose or gums, blurred or lost vision from retinal changes, headache, dizziness, confusion, stroke-like symptoms, or heart failure; severe cases are an emergency managed with plasma exchange while the underlying disease is treated.

\medterm{Hypervitaminosis A} Illness caused by too much vitamin A, usually from supplements or medicines. It may cause headache, nausea, dry skin, hair loss, liver injury, bone pain, or birth defects during pregnancy.

\medterm{Hypervitaminosis D} Illness caused by too much vitamin D, leading to high blood calcium. It may cause nausea, vomiting, thirst, frequent urination, weakness, kidney stones, confusion, or kidney damage.

\medterm{Hypervolaemia} An excess of blood or body fluid. It may cause rapid weight gain, swelling, breathlessness, raised neck veins, or high blood pressure, especially with heart, kidney, or liver disease.

\medterm{Hypervolemia} Too much fluid in the bloodstream or body tissues. It may cause rapid weight gain, swelling, breathlessness, raised neck veins, or high blood pressure and can occur with heart, kidney, or liver disease.

\medterm{Hypha} A branching tubular filament that forms the basic structural unit of most molds and many fungi; networks of hyphae form a mycelium and may invade tissue in fungal disease.

\medterm{Hyphema} Blood collecting in the front part of the eye, usually after an injury. It may cause pain, blurred vision, light sensitivity, or increased eye pressure and requires prompt eye assessment to protect vision.

\medterm{Hypnagogic} Occurring during the transition from wakefulness to sleep. Hypnagogic sensations or hallucinations can be benign, but frequent or distressing episodes may occur with narcolepsy, sleep deprivation, medicines, or other sleep and neurologic disorders.

\medterm{Hypnagogic Hallucinations} Vivid sights, sounds, sensations, or feelings that occur while a person is falling asleep. They are often harmless but may be linked to sleep deprivation, narcolepsy, medicines, or other sleep disorders.

\medterm{Hypnic Headache} A rare primary headache that repeatedly awakens a person during sleep, usually in later life and often at a similar time each night.

\medterm{Hypnogram} A chart showing when a person is awake and in each sleep stage during a recording. It is commonly made from a sleep study and helps show sleep timing, awakenings, and stage patterns.

\synonyms
Sleep Architecture Diagram; Polysomnographic Hypnogram

\medterm{Hypnosis} A state of focused attention and deep relaxation in which a person may respond more readily to guided suggestions. It is not sleep or loss of control and should be used by a trained professional when offered as part of care.

\medterm{Hypnotherapy} Hypnotherapy uses focused attention and guided suggestion as an adjunct to care for selected problems such as pain, anxiety, or unwanted habits.

\medterm{Hypnotic} A medicine that helps a person fall asleep, stay asleep, or improve sleep. It is used mainly for short-term insomnia because some hypnotics can cause next-day drowsiness, falls, dependence, breathing problems, or dangerous effects with alcohol and other sedatives.

\medterm{Hypoactive Sexual Desire Disorder} Persistently low or absent sexual interest that causes personal distress or relationship difficulty. Desire can be affected by medicines, hormones, illness, pain, mood, stress, trauma, or relationship factors, so assessment looks for the underlying cause.

\medterm{Hypoadrenalism} Insufficient adrenal-cortex hormone production, especially cortisol and sometimes aldosterone; primary disease can cause fatigue, weight loss, low blood pressure, hyperpigmentation, and adrenal crisis.

\medterm{Hypoalbuminemia} An abnormally low blood albumin concentration caused by reduced production, inflammation, kidney or gastrointestinal loss, or severe malnutrition; it can contribute to edema and alter drug binding.

\medterm{Hypoalphalipoproteinemias} Disorders with low HDL-related apolipoprotein levels, which may impair reverse cholesterol transport and increase cardiovascular risk depending on the cause and overall lipid profile.

\medterm{Hypoblast} The inner cell layer of the early embryonic disc that contributes to the yolk-sac lining and helps establish the embryonic body axes during gastrulation.

\medterm{Hypocalcaemia} Alternate British spelling; see \textit{Hypocalcemia}.

\medterm{Hypocalcemia} An abnormally low level of calcium in the blood (total calcium below $8.5\,\text{mg/dL}$ or $2.1\,\text{mmol/L}$). Calcium is essential for nerve transmission, muscle contraction, and cardiac electrical stability. Low levels cause neuromuscular irritability, presenting with perioral numbness, fingertip tingling, painful muscle cramps, tetany, and hallmark physical signs: Chvostek's sign (facial twitching on tapping the facial nerve) and Trousseau's sign (carpal spasm during blood pressure cuff inflation). Severe cases cause seizures, airway spasm (laryngospasm), and a prolonged QT interval on ECG. Common causes include hypoparathyroidism (especially after thyroid surgery), severe vitamin D deficiency, chronic kidney disease, low magnesium, and pancreatitis. Severe cases require urgent intravenous calcium gluconate.
\synonyms{Low Blood Calcium; Calcium Deficiency; Hypocalcaemia}

\medterm{Hypocapnia} An abnormally low level of carbon dioxide in the blood, usually because a person is breathing too fast or deeply. Causes include anxiety, pain, low oxygen, lung disease, pregnancy, or excessive ventilator support.

\medterm{Hypochlorite} A chemical form of chlorine used in disinfectants and bleach. It can kill many germs, but swallowing it or getting concentrated solutions in the eyes, lungs, or skin can cause serious injury.

\medterm{Hypocholesterolemia} An unusually low blood cholesterol concentration caused by malnutrition, malabsorption, hyperthyroidism, chronic illness, liver disease, or inherited lipid disorders.

\medterm{Hypochondria} An older term for excessive fear of illness. Current clinical language generally uses illness anxiety disorder or health anxiety, depending on whether prominent physical symptoms are present.

\medterm{Hypochondriasis} Alternate historical term; see \textit{Illness Anxiety Disorder}.

\medterm{Hypochondroplasia} An inherited disorder of bone growth that causes short stature and shortened arms and legs, usually less severely than achondroplasia. People may have broad hands, increased inward curvature of the lower back, or limited elbow movement, while facial appearance and intelligence are often typical. It is usually caused by a change in the \textit{FGFR3} gene.

\synonyms
HCH; Mild Achondroplasia Variant

\medterm{Hypochromia} Reduced hemoglobin color in red blood cells, seen as increased central pallor on a blood smear and commonly associated with iron deficiency or thalassemia.

\medterm{Hypochromic Anaemia} Hypochromic anaemia is anemia in which red blood cells contain too little hemoglobin and appear unusually pale, commonly because of iron deficiency or thalassemia.

\medterm{Hypochromic Anemia} Anemia in which red cells contain too little hemoglobin and appear pale, most often because of iron deficiency, thalassemia, chronic disease, or sideroblastic disorders.

\medterm{Hypocitraturia} Low urinary citrate excretion, which reduces inhibition of calcium-crystal formation and increases risk of calcium kidney stones; potassium citrate may be used when clinically appropriate.

\medterm{Hypocupremia} An abnormally low copper concentration in blood or tissues, caused by inadequate intake, malabsorption, excess zinc, or inherited disorders, and potentially producing anemia, neutropenia, or neurologic disease.

\medterm{Hypodermic Needle} A hollow needle designed to inject medicines or withdraw fluid through the skin, with gauge and length selected according to the substance, route, body site, and patient.

\medterm{Hypodermis} The layer of fat and loose tissue beneath the skin. It cushions the body, stores energy, helps conserve heat, and contains blood vessels, nerves, and lymph channels.

\medterm{Hypodontia} A condition present from birth in which one or more permanent teeth fail to develop. It most often affects wisdom teeth, upper side incisors, or second premolars and may affect chewing, appearance, or tooth alignment.

\medterm{Hypoesthesia} Reduced sensitivity to touch, temperature, pain, or other sensations. It may result from nerve compression,
neuropathy, spinal disease, stroke, injury, or other neurologic conditions; sudden or spreading
numbness requires prompt assessment.

\medterm{Hypogalactia} Insufficient breast-milk production for an infant’s needs, caused by ineffective milk removal, hormonal or breast disorders, medications, illness, or other feeding difficulties.

\medterm{Hypogammaglobulinemia} Abnormally low levels of antibodies in the blood, which can lead to repeated infections, especially of the ears, sinuses, lungs, or bloodstream. It may be inherited or develop because of another disease or treatment.

\medterm{Hypogastric Region} The lower central area of the abdomen, just above the pubic bone. It overlies the bladder and parts of the bowel and reproductive organs; pain or a lump there may arise from urinary, intestinal, or pelvic disease.

\medterm{Hypoglossal Nerve} Cranial nerve XII, which supplies the muscles that move the tongue; a lesion causes tongue weakness, wasting, and deviation toward the affected side on protrusion.

\medterm{Hypoglossal Nerve Injury} Damage to the nerve that controls tongue movement. It may follow skull-base injury, surgery, a tumor, or infection and can cause tongue weakness, wasting, or deviation when the tongue is extended.

\medterm{Hypoglossal Nerve Palsy} Weakness or paralysis of the hypoglossal nerve (CN XII), leading to ipsilateral tongue
weakness, atrophy, and fasciculations. Causes include tumors at the skull base or in the
hypoglossal canal, surgical injury, or central lesions. Unilateral palsy causes the
tongue to deviate toward the affected side on protrusion.

\medterm{Hypoglycaemia} An abnormally low blood-glucose level. It may cause sweating, shaking, hunger, palpitations, confusion, seizures, or loss of consciousness and is especially associated with diabetes medicines, missed meals, or heavy alcohol use.

\medterm{Hypoglycemia} An abnormally low blood-glucose concentration that may cause sweating, tremor, palpitations, hunger,
confusion, seizures, or loss of consciousness. Causes include glucose-lowering medicines,
missed meals, alcohol, and critical illness. Treatment depends on the person’s ability to
swallow and the cause; severe or recurrent episodes require urgent assessment.

\medterm{Hypoglycemia In Diabetes} Blood glucose that falls too low in a person with diabetes, usually below 70 mg/dL. It may cause sweating, shaking, hunger, confusion, or fainting; severe episodes prevent self-treatment and require another person’s help.

\medterm{Hypoglycemia Unawareness} Failure to feel the usual warning signs of low blood glucose, such as shaking, sweating, or a racing heartbeat. It can occur after repeated low readings or long-standing diabetes and increases the risk of severe hypoglycemia.

\medterm{Hypoglycemic} Relating to an abnormally low blood glucose level, or describing a medicine that lowers blood glucose.

\medterm{Hypogonadism} Reduced function of the testes or ovaries, causing low sex-hormone levels and impaired production of sperm or eggs. It may lead to delayed puberty, absent periods, infertility, low sexual desire, hot flashes, reduced muscle or bone strength, or other symptoms.

\medterm{Hypogonadism In Men} Low testosterone caused by reduced testicular function or by a problem in the brain signals that control the testes. It may cause low sexual desire, erectile problems, infertility, low energy, reduced muscle, or weak bones.

\medterm{Hypogonadism In Women} Reduced ovarian hormone production, which may cause absent or irregular periods, hot flashes, vaginal dryness, infertility, low sexual desire, or loss of bone strength. Causes include ovarian failure, pituitary disease, some medicines, severe illness, or surgery.

\medterm{Hypogonadotropic Hypogonadism} Low production of sex hormones because the brain or pituitary gland does not send enough stimulating signals to the testes or ovaries. It may cause delayed puberty, absent periods, infertility, low sexual desire, or reduced testosterone. Causes include genetic conditions, pituitary disease, obesity, medicines, or severe illness.

\medterm{Hypohidrosis} Reduced or absent sweating. It can make it difficult for the body to cool itself and increase the risk of overheating. Causes include nerve disorders, skin damage, medicines, dehydration, and inherited conditions.

\medterm{Hypokalaemia} British spelling of Hypokalemia; see Hypokalemia.

\medterm{Hypokalemia} An abnormally low level of potassium in the blood, often caused by vomiting, diarrhea, diuretics, hormone disorders, or movement of potassium into cells. It may cause weakness, cramps, constipation, or dangerous heart-rhythm changes.

\medterm{Hypokalemic Periodic Paralysis} An inherited disorder causing repeated attacks of muscle weakness or paralysis when blood potassium falls. Attacks often begin after rest following exercise or after a high-carbohydrate meal, while strength may be normal between attacks.

\medterm{Hypokinesia} Reduced or slowed movement. It may cause small steps, less facial expression, reduced arm swing, or small handwriting and commonly occurs in Parkinson disease or other disorders affecting movement control.

\medterm{Hypolipidemic Agents} Medicines that lower cholesterol or triglycerides, including statins, ezetimibe, fibrates, bile-acid sequestrants, niacin, and selected injectable therapies.

\medterm{Hypolipoproteinemia} An abnormally low concentration of one or more blood lipoproteins, caused by inherited disorders, malabsorption, malnutrition, liver disease, or other systemic illness.

\medterm{Hypomagnesaemia} An abnormally low level of magnesium in the blood. It may cause tremor, muscle cramps, weakness, seizures, or abnormal heart rhythms and can result from diarrhea, poor intake, alcohol use, or certain medicines.

\medterm{Hypomagnesemia} A low level of magnesium in the blood, often caused by diarrhea, poor nutrition, alcohol use, kidney loss, or some medicines. It may cause tremor, cramps, weakness, seizures, or abnormal heart rhythms.

\medterm{Hypomania} A distinct period of unusually elevated, expansive, or irritable mood with increased energy or activity lasting at least four consecutive days. It may involve reduced need for sleep, rapid speech, racing thoughts, inflated self-esteem, or impulsive behavior; the change is observable by others but is not severe enough to cause marked impairment or require hospitalization and has no psychotic features.

\medterm{Hypomanic Episode} Alternate name; see \textit{Hypomania}.

\medterm{Hypomelanosis Of Ito} A rare condition causing swirled or streaky areas of lighter skin, often from birth. Some affected children also have seizures, developmental differences, eye problems, or abnormalities of the bones or nervous system.

\medterm{Hypomethylation} An abnormally low amount of methyl marking on DNA or other molecules. It can change which genes are active and may affect chromosome stability, development, inherited disorders, or cancer cells.

\medterm{Hypomnesia} Reduced ability to remember or recall information, which may occur with aging, neurologic disease, psychiatric conditions, medications, or other causes.

\medterm{Hyponatraemia} Alternate British spelling; see \textit{Hyponatremia}.

\medterm{Hyponatremia} An abnormally low level of sodium in the blood, often caused by excess water, fluid loss, medicines, heart or liver disease, kidney problems, or hormone disorders. It may cause nausea, headache, confusion, seizures, or coma; severe or rapidly developing cases are emergencies.

\medterm{Hypoparathyroidism} A disorder in which too little parathyroid hormone causes low blood calcium and high phosphate. It most often follows neck surgery but can be inherited or autoimmune. Tingling, muscle cramps, spasms, seizures, or abnormal heart rhythms may occur and need treatment.

\medterm{Hypopharyngeal Neoplasm} An abnormal growth in the lower part of the throat beside the voice box and esophagus. It may cause persistent sore throat, painful swallowing, difficulty swallowing, ear pain, or a neck lump and may be cancerous.

\medterm{Hypopharynx} The lower part of the pharynx behind and beside the larynx, extending from the oropharynx to the esophagus. It directs food toward the esophagus during swallowing.

\medterm{Hypophosphataemia} Alternate British spelling; see \textit{Hypophosphatemia}.

\medterm{Hypophosphatemia} An abnormally low level of phosphate in the blood (below $2.5\,\text{mg/dL}$ or $0.8\,\text{mmol/L}$). Because phosphate is essential for generating cellular energy (ATP) and oxygen delivery (2,3-DPG), its deficiency starves tissues of energy. It is classically triggered by refeeding syndrome (rapid carbohydrate refeeding in malnourished or alcoholic patients driving phosphate into cells), diabetic ketoacidosis recovery, severe malnutrition, or renal phosphate wasting. Symptoms include profound muscle weakness, respiratory failure from diaphragmatic weakness, acute muscle breakdown (rhabdomyolysis), hemolytic anemia, confusion, and seizures. Treatment consists of oral or intravenous phosphate replacement.
\synonyms{Low Blood Phosphate; Hypophosphataemia; Phosphate Deficiency}

\medterm{Hypophosphatemic Rickets} Rickets caused by low phosphate availability, often from renal phosphate wasting, leading to defective mineralization, bowed legs, bone pain, growth impairment, and dental abnormalities.

\medterm{Hypophysectomy} Surgical removal of all or part of the pituitary gland, usually performed through the nose and sphenoid sinus (transsphenoidal hypophysectomy) or rarely through an open craniotomy to treat pituitary tumors, relieve optic nerve compression, or control hormone hypersecretion.
\synonyms
Pituitary Resection; Transsphenoidal Adenomectomy

\medterm{Hypophysis} The pituitary gland, a small hormone-producing gland at the base of the brain. Its hormones influence growth, thyroid and adrenal activity, reproduction, milk production, and the body’s water balance.

\medterm{Hypophysis} An endocrine gland at the base of the brain that regulates growth, thyroid and adrenal function, reproduction, lactation, and water balance through its hormones.

\synonyms
Pituitary Gland; Master Endocrine Gland

\medterm{Hypophysitis} Inflammation of the pituitary gland that can cause headache, visual symptoms, and deficiency or excess of pituitary hormones; it may be autoimmune, related to pregnancy, or triggered by immune-checkpoint therapy.

\medterm{Hypopigmentation} Reduced skin pigmentation producing paler areas than the surrounding skin, caused by decreased melanin, reduced melanocyte function or number, inflammation, injury, infection, or inherited disorders.

\medterm{Hypopituitarism} Reduced production of one or more hormones from the pituitary gland because of disease, injury, surgery, radiation, or a nearby brain disorder. It may affect growth, thyroid or adrenal function, periods, fertility, milk production, or water balance.

\synonyms
Pituitary Insufficiency; Panhypopituitarism

\medterm{Hypoplasia} Incomplete or underdevelopment of a tissue, organ, or body structure, resulting in
reduced size or functional capacity. It may be congenital, developmental, genetic, or
acquired and ranges from a mild anatomical variant to severe organ dysfunction.
Assessment considers the timing of onset, associated malformations, tissue function, and whether the process is isolated or syndromic.

\medterm{Hypoplastic Left Heart Syndrome} A serious birth defect in which the left side of the heart and the aorta are too small to pump enough blood to the body. Newborns need urgent specialist treatment, often involving staged surgery or transplantation.

\medterm{Hypoploidy} A chromosome number below the normal diploid complement, caused by loss of one or more chromosomes; in tumor cytogenetics it may have prognostic significance.

\medterm{Hypopnea} A period of unusually shallow or reduced breathing, often during sleep. Repeated episodes can lower blood oxygen, disturb sleep, and contribute to obstructive sleep apnea or other breathing disorders.

\medterm{Hypoproteinaemia} Hypoproteinaemia is an abnormally low level of proteins in the blood, which may result from malnutrition, liver disease, kidney protein loss, or intestinal loss.

\medterm{Hypoproteinemia} An abnormally low amount of protein in the blood, often involving albumin, which can cause swelling and reflect poor nutrition, liver disease, kidney loss, or illness.

\medterm{Hypopyon} A visible white or yellow level of inflammatory material in the front chamber of the eye. It usually indicates severe eye inflammation or infection and may occur with pain, redness, light sensitivity, and reduced vision.

\medterm{Hyposmia} A reduced ability to smell, without complete loss of smell. It can follow a cold or COVID-19, nasal blockage, head injury, medicines, aging, or a neurologic disorder.

\medterm{Hypospadias} A birth condition in which the opening of the urethra is on the underside of the penis rather than at its tip. The penis may curve, and the foreskin may be incomplete; severity depends on the opening’s location and other abnormalities.

\medterm{Hypospadias Repair} Surgery reconstructs the urethral opening and, when needed, straightens the penis in a child born with hypospadias; timing and technique depend on anatomy and surgeon assessment.

\medterm{Hypotension} Abnormally low blood pressure that may reduce blood flow to the brain and other organs. It can result from bleeding, dehydration, heart problems, severe infection, allergy, or medicines and may cause dizziness, fainting, confusion, or shock.

\medterm{Hypotension In Anesthesia} Low blood pressure during or soon after anesthesia. It may result from anesthetic medicines, blood loss, dehydration, heart problems, or an allergic reaction; the anesthesia team monitors and treats it according to the cause and severity.

\medterm{Hypotensive} Having or causing abnormally low blood pressure; symptoms may include dizziness, fainting, weakness, confusion, or shock depending on the cause and severity.

\medterm{Hypothalamic Amenorrhea} Absence of menstrual periods because the hypothalamus reduces its normal release of reproductive signals. It is often linked to low energy availability, excessive exercise, stress, illness, or weight loss and may cause infertility and bone loss.

\medterm{Hypothalamic Dysfunction} A problem affecting the brain region that helps control hormones, temperature, appetite, thirst, sleep, stress, and reproduction. Causes include injury, tumors, inflammation, infection, surgery, radiation, or inherited disease.

\medterm{Hypothalamic Hamartoma} A congenital, noncancerous mass of disorganized hypothalamic tissue that can cause gelastic seizures, precocious puberty, behavioral symptoms, or cognitive decline.

\medterm{Hypothalamic Hormone} A releasing or inhibiting hormone produced by the hypothalamus that regulates pituitary secretion, including signals controlling thyroid, adrenal, reproductive, growth, and lactation axes.

\medterm{Hypothalamic Neoplasm} An abnormal growth in or near the hypothalamus, a brain region controlling hormones, temperature, appetite, sleep, and water balance. Effects may include vision changes, hormone problems, headaches, seizures, or altered growth.

\medterm{Hypothalamic Tumor} Alternate name; see \textit{Hypothalamic Neoplasm}.

\medterm{Hypothalamic-Pituitary-Adrenal Axis} A network linking the hypothalamus, pituitary gland, and adrenal glands to regulate stress hormones and other functions.

\synonyms
HPA Axis; Corticotrope-Adrenal System

\medterm{Hypothalamic-Pituitary-Adrenal Axis} The hormone system linking the brain’s hypothalamus and pituitary gland with the adrenal glands. It controls cortisol release, helping regulate the body’s response to stress, illness, blood pressure, metabolism, and immune activity.

\synonyms
HPA Axis; Neuroendocrine Stress Axis

\medterm{Hypothalamic-Pituitary-Adrenal Axis Physiology} The brain signals the pituitary gland, which signals the adrenal glands to release cortisol. Cortisol normally follows a daily pattern and feeds back to reduce further signals when enough is present.

\medterm{Hypothalamic-Pituitary-Gonadal Axis Physiology} The hormone system linking the hypothalamus, pituitary gland, and ovaries or testes. Signals from the hypothalamus prompt the pituitary to release hormones that control ovulation and menstrual cycles, or testosterone production and sperm development.

\medterm{Hypothalamic-Pituitary-Thyroid Axis Physiology} The hormone system linking the hypothalamus, pituitary gland, and thyroid. The hypothalamus signals the pituitary to release thyroid-stimulating hormone, which prompts the thyroid to make hormones that regulate metabolism, temperature, and energy use.

\medterm{Hypothalamus} A small region at the base of the brain that helps control automatic body functions and hormone release through the pituitary gland. It regulates body temperature, hunger, thirst, sleep, stress, and parts of emotional behavior.

\medterm{Hypothermia} A dangerously low internal body temperature, usually below 35 degrees Celsius, caused by cold exposure or impaired heat production. It may cause shivering, confusion, drowsiness, slow breathing, or an abnormal heart rhythm. Severe cases require urgent medical rewarming.

\synonyms
Accidental Hypothermia; Low Core Body Temperature

\medterm{Hypothermia In Anesthesia} An abnormally low body temperature during or after anesthesia. Anesthetic medicines reduce the body’s ability to control heat and may allow heat loss during surgery; warming devices and temperature monitoring help prevent complications.

\medterm{Hypothyroidism} A clinical endocrine syndrome resulting from deficient production or action of thyroid hormones (thyroxine [$T_4$] and triiodothyronine [$T_3$]), causing a generalized slowing of basal metabolism. Primary hypothyroidism (over 95\% of cases, characterized by elevated serum TSH and low free $T_4$) most commonly results from autoimmune destruction (Hashimoto thyroiditis), surgical thyroidectomy, radioactive iodine ablation, or severe iodine deficiency; central hypothyroidism stems from pituitary or hypothalamic failure. Clinical hallmarks include profound fatigue, cold intolerance, weight gain, constipation, sinus bradycardia, dry skin, non-pitting dermal edema (myxedema), delayed relaxation of deep tendon reflexes, and cognitive slowing.
\synonyms{Underactive Thyroid; Hypothyroid State; Thyroid Deficiency}

\medterm{Hypotonia} Decreased muscle tone presenting as a
floppy or poorly supported posture. It may result from muscle, peripheral nerve,
neuromuscular-junction, spinal-cord, brain, genetic, or metabolic disorders and requires
evaluation when persistent or associated with weakness or developmental delay.

\medterm{Hypotonic Solution} A liquid containing fewer dissolved particles than the fluid inside cells. Water tends to move into cells when they are exposed to it, so medical use requires careful selection and monitoring.

\medterm{Hypotrichosis} Abnormally sparse or absent hair caused by inherited follicle disorders, autoimmune disease, endocrine or nutritional problems, medications, or scarring alopecia.

\medterm{Hypotrophy} Underdevelopment or reduced size of a tissue or organ that is present but smaller or less developed than usual. It may begin before birth or result from poor nutrition, reduced use, impaired blood flow, or disease and may affect function.

\medterm{Hypotropia} Vertical eye misalignment in which one eye deviates downward relative to the other, potentially causing diplopia, abnormal head posture, or suppression of vision.

\medterm{Hypoventilation} Breathing too slowly or shallowly to remove enough carbon dioxide. It can lower oxygen levels and cause sleepiness, headache, confusion, or respiratory failure; causes include medicines, lung disease, obesity, and muscle or nerve disorders.

\medterm{Hypovolaemia} Hypovolaemia is an abnormally low volume of circulating blood or body fluid, commonly caused by bleeding, dehydration, or major fluid loss.

\medterm{Hypovolemia} Hypovolemia is reduced effective circulating blood volume, usually from bleeding, fluid loss, or inadequate intake; it can cause tachycardia, hypotension, poor perfusion, kidney injury, and hypovolaemic shock.

\medterm{Hypovolemic Shock} Circulatory failure caused by insufficient circulating blood or fluid volume, leading to reduced
tissue perfusion. Causes include hemorrhage, severe dehydration, and fluid loss.

\medterm{Hypoxaemia} An abnormally low oxygen level in arterial blood. It can result from lung disease, breathing problems, poor blood flow, or other conditions and may cause shortness of breath, confusion, bluish skin, or organ injury.

\medterm{Hypoxanthine} A naturally occurring substance made when the body breaks down certain genetic-material building blocks. It is normally reused or changed into other compounds during purine metabolism.

\medterm{Hypoxemia} An abnormally low oxygen level in the blood. It may result from lung or heart disease, blocked airways, shallow breathing, anemia, poor circulation, or high altitude and can cause breathlessness, confusion, or bluish skin.

\medterm{Hypoxia} Inadequate oxygen availability at the tissue level. It may result from low arterial oxygen, anemia,
reduced blood flow, or impaired cellular oxygen use.

\medterm{Hypoxia And Hypoxemia} Hypoxemia is low oxygen in arterial blood, while hypoxia is inadequate oxygen at the tissue level. Causes include lung disease, airway obstruction, heart disease, anemia, shock, or high altitude; severe breathlessness, cyanosis, confusion, or collapse is urgent.

\medterm{Hypoxic Cell Injury} Damage to cells caused by too little oxygen. It may result from poor blood flow, low blood oxygen, anemia, or poisoning that prevents oxygen use; the outcome depends on how severe and prolonged the oxygen shortage is.

\medterm{Hypoxic Ischemic Encephalopathy} Brain dysfunction caused by insufficient oxygen and blood flow, usually around birth in newborns but also after cardiac arrest or severe respiratory failure. It can cause seizures, altered consciousness, and long-term neurologic disability.

\medterm{Hypoxic Pulmonary Vasoconstriction} Narrowing of small blood vessels in an area of the lung that receives too little oxygen. This normally redirects blood toward better-ventilated lung areas, but widespread or prolonged narrowing can strain the heart.

\medterm{Hypoxic-Ischemic Brain Injury} Brain damage caused by too little oxygen or blood flow. It may follow cardiac arrest, suffocation, shock, severe breathing failure, or very low blood glucose and can cause seizures, coma, memory problems, movement problems, or disability.

\medterm{Hypsarrhythmia} A very abnormal, disorganized pattern on an electroencephalogram, with large slow waves and spikes arising from many areas of the brain. It is strongly associated with infantile spasms, a serious seizure disorder in babies, and requires prompt specialist evaluation and treatment.

\synonyms
Hypsarrhythmic EEG Pattern; Chaotic Dysrhythmic EEG

\medterm{Hysterectomy} Surgical removal of the uterus. The cervix, fallopian tubes, and ovaries may or may not also be removed, depending on the procedure and clinical indication. Hysterectomy permanently ends menstruation and the ability to carry a pregnancy, but it does not cause immediate menopause when functioning ovaries remain. If the cervix remains, cervical cancer screening may still be needed. Indications include cancer, fibroids, severe bleeding, adenomyosis, endometriosis, chronic pelvic pain, and uterine prolapse.

\medterm{Hysterectomy And Surgical Menopause} Hysterectomy removes the uterus and ends menstruation and the ability to carry a pregnancy. Surgical menopause occurs when both ovaries are removed or ovarian function otherwise stops abruptly; hysterectomy alone does not necessarily cause menopause if the ovaries remain.

\medterm{Hysterosalpingo-Contrast Sonography} An ultrasound test using fluid or foam contrast to assess the uterine cavity and whether the fallopian tubes are open. It is often used during fertility evaluation without X-ray radiation.

\synonyms
HyCoSy; Fallopian Tubal Contrast Ultrasound

\medterm{Hysterosalpingo-Contrast-Sonography} An ultrasound-based technique for evaluating the uterine cavity and fallopian tube
patency. Sterile saline or contrast medium is injected through the cervix while
transvaginal ultrasound is performed.

\medterm{Hysterosalpingogram} An X-ray procedure used to evaluate the uterine cavity and fallopian tube patency. A
radiopaque contrast medium is injected through the cervix, and serial X-ray images are
taken as the contrast fills the uterine cavity and spills from the fallopian tubes.

\medterm{Hysterosalpingography} An X-ray test in which contrast liquid is placed through the cervix to outline the inside of the uterus and fallopian tubes. It can show blockage, an unusual uterine shape, fibroids, polyps, or scar tissue during infertility evaluation. Cramping, bleeding, infection, or contrast reaction can occur.

\medterm{Hysteroscope} A thin camera-equipped instrument passed through the cervix to inspect the uterine cavity and, when designed for operative use, remove polyps, treat bleeding, or perform other procedures.

\medterm{Hysteroscopic Metroplasty} A procedure that uses a small camera and instruments passed through the cervix to divide a wall of tissue separating part of the uterus. It may improve the uterine cavity in selected people with fertility problems or repeated pregnancy loss.

\medterm{Hysteroscopic Myomectomy} A procedure that removes selected fibroids bulging into the inside of the uterus. A camera and instruments pass through the cervix, avoiding an abdominal incision; the procedure may reduce heavy bleeding and improve fertility.

\medterm{Hysteroscopy} A procedure that passes a small camera through the cervix to inspect the inside of the uterus and, when needed, obtain tissue or treat conditions such as polyps, fibroids, or intrauterine adhesions.

\medterm{Hysterotomy} A surgical incision into the uterus performed during selected obstetric or gynecologic procedures. It may be used to deliver a baby or reach the uterus, with risks including bleeding and infection.

\medterm{Infant Hearing Loss} Congenital or early-acquired auditory impairment in infants, resulting from genetic mutations (such as connexin 26 mutations), congenital infections (cytomegalovirus, rubella), prematurity, hyperbilirubinemia, or ototoxic drug exposure. Universal newborn hearing screening enables early amplification and speech-language intervention.

\synonyms
Congenital Hearing Loss; Pediatric Auditory Impairment; Neonatal Deafness

\medterm{Juvenile Hypermobile Spectrum Disorder} A hypermobility spectrum disorder beginning in childhood or adolescence, with excessive joint movement plus pain, instability, soft-tissue injury, fatigue, or related symptoms.

\medterm{Lanugo Hair} Lanugo: the fine, soft, usually unpigmented hair covering a fetus and sometimes present on a premature newborn; it is normally shed before or shortly after birth.

\medterm{Mevalonate Kinase Deficiency} An inherited autoinflammatory disorder caused by changes in the MVK gene. It can cause repeated fever attacks, swollen lymph nodes, abdominal pain, rash, joint pain, or mouth ulcers; the milder form is often called hyperimmunoglobulin D syndrome.

\medterm{Minimally Invasive Heart Bypass Surgery} Coronary artery bypass graft (CABG) surgery performed through a small thoracotomy incision between the ribs, frequently on a beating heart without cardiopulmonary bypass, avoiding a traditional median sternotomy.

\synonyms
MIDCAB; Minimally Invasive Direct Coronary Artery Bypass; Less Invasive CABG

\medterm{Neonatal Hyperglycemia} An abnormally elevated blood glucose level in infants and neonates (typically defined as blood glucose $>145\text{ mg/dL}$ or $>8.0\text{ mmol/L}$). Common causes include extreme prematurity, intravenous dextrose overinfusion, sepsis, severe physiological stress, corticosteroid therapy, or transient neonatal diabetes mellitus.

\synonyms
Infantile Hyperglycemia; High Blood Glucose in Infants; Newborn Hyperglycemia

\medterm{Neonatal Hypertension} Sustained blood pressure elevation in infants and neonates exceeding age- and birthweight-adjusted norms. Frequently secondary to umbilical artery catheter-related renal artery thrombosis, congenital renal anomalies, bronchopulmonary dysplasia, or congenital heart defects.

\synonyms
Infant Hypertension; Neonatal High Blood Pressure; Infantile Hypertension

\medterm{Oral Herpes} A common viral infection of the oral mucosa and lips caused primarily by herpes simplex virus type 1 (HSV-1). Characterized by recurrent outbreaks of fluid-filled vesicles (cold sores or fever blisters) that ulcerate and crust, accompanied by prodromal tingling or burning.

\synonyms
Herpes Simplex Labialis; Cold Sores; Fever Blisters; Oral HSV-1 Infection

\medterm{Pediatric Hepatitis A} An acute viral infection of the liver caused by the hepatitis A virus (HAV) in children, transmitted predominantly through the fecal-oral route via contaminated food or water. While often self-limiting, supportive hydration is critical and routine pediatric vaccination provides durable protection.

\synonyms
Childhood Hepatitis A; Pediatric HAV Infection; Acute Viral Hepatitis A in Children

\medterm{Pediatric Hepatitis B} A viral liver infection caused by the hepatitis B virus (HBV) in pediatric patients, most commonly acquired perinatally from an infected mother or during early childhood. Chronic carrier states are highest when acquired in infancy; universal infant vaccination and post-exposure immunoprophylaxis prevent transmission.

\synonyms
Childhood Hepatitis B; Perinatal HBV Infection; Chronic Pediatric Hepatitis B

\medterm{Pediatric Hepatitis C} Infection with the hepatitis C virus (HCV) in infants and children, primarily acquired vertically through perinatal maternal-fetal transmission or through parenteral exposure. Modern direct-acting antiviral (DAA) regimens provide high cure rates in pediatric cohorts.

\synonyms
Childhood Hepatitis C; Pediatric HCV Infection; Perinatal Hepatitis C

\medterm{Pediatric Hypercholesterolemia} Elevated circulating total cholesterol or low-density lipoprotein (LDL) levels in children and adolescents. Causes include familial hypercholesterolemia, obesity, insulin resistance, endocrine disorders, or high-saturated-fat diets.

\synonyms
Childhood Hyperlipidemia; Pediatric Dyslipidemia; Familial Hypercholesterolemia in Children

\medterm{Pediatric Hypertension} Blood pressure in children and adolescents persistently at or above the 95th percentile for age, sex, and height. Secondary etiologies (such as renal parenchymal disease or aortic coarctation) are common, especially in young children.

\synonyms
Childhood Hypertension; Pediatric High Blood Pressure; Adolescent Hypertension

\medterm{Perineal Hemostat} A clamp used during procedures to compress a blood vessel or hold tissue and control bleeding.

\synonyms
Perineal Hemostatic Clamp; Pelvic Bleeding Forceps

\medterm{Physical Half-Life} The time it takes for half of a radioactive substance to lose its radioactivity. It helps determine the timing of scans and the amount of radiation exposure.

\medterm{Piles} Alternate name; see \textit{Hemorrhoids}.

\medterm{Psychological Homeostasis} The ongoing adjustment of thoughts, emotions, and behavior that helps a person maintain psychological stability while responding to stress or change. Severe disruption can contribute to anxiety, depression, or impaired daily functioning.

\medterm{Qualitative HCG Blood Test} A diagnostic serological laboratory test that detects the presence or absence of human chorionic gonadotropin (hCG) in maternal blood. Primarily utilized for early confirmation of pregnancy.

\synonyms
Qualitative Serum Pregnancy Test; Blood hCG Screen; Qualitative Beta-hCG

\medterm{Quantitative HCG Blood Test} A precise diagnostic laboratory assay that measures the exact serum concentration of beta-human chorionic gonadotropin. Essential for monitoring early pregnancy viability, evaluating suspected ectopic pregnancy or miscarriage, and tracking gestational trophoblastic disease or germ-cell neoplasms.

\synonyms
Quantitative Beta-hCG; Serial hCG Test; Serum hCG Level

\medterm{Sensorineural Hearing Loss} Hearing loss caused by damage to the sound-sensing cells or nerve fibers of the inner ear, the auditory nerve, or related hearing pathways. Sounds may seem muffled, and understanding speech---especially in noisy places---may be difficult; ringing in the ears (tinnitus) may also occur. Common causes include ageing, loud noise, inherited conditions, infections, head injury, and some medicines. It differs from conductive hearing loss, which results from a problem in the outer or middle ear, and from mixed hearing loss, which includes both types.

\medterm{Singultus} Alternate medical name; see \textit{Hiccups}.

\medterm{Social History} Information about a person's living situation, relationships, occupation, education, tobacco or substance use, diet, safety, and other social factors relevant to health and care.

\medterm{Surgical Hemostat} A clamping surgical instrument used to control bleeding by compressing a blood vessel or tissue.
\synonyms
Surgical Hemostatic Forceps; Mosquito Clamp

\medterm{Swine Flu} Alternate name; see \textit{H1N1 Influenza}.

\medterm{Urticaria} Alternate medical name; see \textit{Hives}.

\medterm{Vascular Hemostat} A surgical clamp used to temporarily occlude a blood vessel and control bleeding.
\synonyms
Vascular Clamp; Atraumatic Hemostasis Forceps

\medterm{Viral Hepatitis} Inflammation of the liver caused by one of the primary hepatitis viruses (A, B, C, D, or E) or other viral agents. Manifestations range from asymptomatic elevation of aminotransferases to acute jaundice, fulminant hepatic failure, or chronic hepatitis leading to cirrhosis.

\medterm{Viral Hepatitis Types} Five viruses that can inflame the liver. They differ in how they spread, whether infection can become long-term, the complications they cause, and whether vaccination or curative treatment is available.

\medterm{Vitamin K Deficiency Bleeding} Bleeding in a newborn or young infant caused by inadequate vitamin K-dependent clotting activity. It may occur early, classically, or late and can involve the brain or gastrointestinal tract.

\synonyms
VKDB; Neonatal Hypoprothrombinemia


\medterm{Iatrogenic} Describing an illness, injury, or other harm caused unintentionally by medical examination, treatment, a medicine, a procedure, or a healthcare device. Examples include side effects, procedure-related bleeding, infection, nerve damage, or radiation injury. The term does not imply negligence.

\medterm{Iatrogenic Disease} A disease or injury caused unintentionally by a medical examination, treatment, medication, procedure, or other healthcare intervention.

\medterm{Iatrogenic Injury} Harm caused unintentionally by a medical examination, treatment, procedure, medicine, or healthcare device. It may be temporary or permanent and can include infection, bleeding, medication effects, nerve damage, or injury during surgery.

\medterm{Ibandronic Acid} A bisphosphonate medicine used in some settings to reduce bone loss and complications from cancer affecting bone. It can rarely damage the jaw or cause unusual thigh fractures.
\medterm{IBD} Inflammatory bowel disease, including Crohn's disease and ulcerative colitis, causing long-term digestive-tract inflammation, diarrhea, abdominal pain, bleeding, weight loss, or fatigue.

\synonyms
Inflammatory bowel disease



\medterm{Ibogaine} Ibogaine is a psychoactive plant alkaloid studied for addiction treatment but associated with dangerous cardiac arrhythmias, seizures, and death.

\medterm{Ibritumomab Tiuxetan} An antibody linked to a radioactive substance and used in selected B-cell lymphomas. The antibody attaches to CD20-positive cells and delivers radiation; low blood counts, infection, and delayed marrow problems are important risks.

\medterm{IBS} Alternate index label; see \textit{Irritable bowel syndrome}.

\medterm{IBS-D Irritable Bowel Syndrome With Diarrhea} IBS with diarrhea causes recurrent abdominal pain associated with loose stools without structural disease explaining the symptoms. Management may include diet changes, soluble fiber, gut-directed medicines, behavioral therapy, and evaluation for alarm features or alternative diagnoses.

\synonyms
Diarrhea IBS (IBS-D Irritable Bowel Syndrome With Diarrhea)

\medterm{Ibuprofen} A medicine that reduces pain, fever, and inflammation. It can irritate or bleed the stomach, harm the kidneys, cause fluid retention, or increase heart risks in some people.

\medterm{Ibuprofen Dosing For Children} Children’s ibuprofen dosing is based mainly on weight, age, product concentration, and medical advice; the label dose should not be exceeded and it should be avoided in dehydration or certain kidney conditions.

\medterm{Ibuprofen Overdose} Toxic exposure to more ibuprofen than recommended, which can cause nausea, vomiting, abdominal pain, gastrointestinal bleeding, kidney injury, metabolic acidosis, drowsiness, seizures, or coma.

\medterm{ICD} An implantable cardioverter-defibrillator, a device that monitors heart rhythm and can deliver pacing or an electrical shock to stop dangerous rhythms and prevent sudden cardiac death.

\synonyms
Implantable cardioverter defibrillator

\medterm{ICD Shock Therapy} Treatment delivered by an implanted defibrillator to stop dangerous heart rhythms. The device may use painless pacing or an electrical shock, depending on the rhythm.

\medterm{ICD-O Neoplasm Morphology Classification} A coding system used by cancer registries and pathology services to record what type of tumor is present and whether it is benign or cancerous. It does not replace the diagnosis, location, or stage of the cancer.
\medterm{Ice Pack} A device that applies cold to a body area to help relieve pain or swelling. It should be wrapped to protect the skin and used for limited periods.

\medterm{Ice Pick Headache} A sudden, severe, stabbing head pain that usually lasts seconds. Attacks may recur in different head areas without warning; new, persistent, or unusually severe pain needs medical assessment to exclude serious causes.

\synonyms
Primary stabbing headache

\medterm{ICF Framework} A World Health Organization framework for describing how a health condition affects the body, daily activities, participation in life, and the person's environment.

\medterm{Ichthyosis} A group of conditions that cause very dry, rough, scaly skin. Some forms are inherited and may be present from birth; others are linked to allergies or another illness.

\medterm{Ichthyosis Vulgaris} A common inherited skin condition causing dry, fine scaling, especially on the arms and legs. It often improves with age and can be managed with regular moisturizers and prescribed scale-reducing treatments.

\medterm{Icteric} Describing yellowing of the skin, whites of the eyes, or other tissues caused by too much bilirubin in the blood.

\medterm{Icteric Index} The icteric index is a laboratory estimate of serum yellow discoloration related mainly to bilirubin and can signal jaundice or interfere with testing.

\medterm{Icterus} Yellowing of the skin or whites of the eyes caused by excess bilirubin. It may result from liver disease, blocked bile flow, or rapid breakdown of red blood cells and needs evaluation.

\medterm{Ictus} A sudden neurological event, such as a stroke or seizure.

\medterm{ICU Psychosis} A colloquial term for delirium occurring in an intensive-care setting, with acute changes in attention, awareness, and
cognition. Causes include severe illness, infection, medicines, sleep disruption, pain, and metabolic abnormalities.
It requires prompt evaluation and treatment of reversible causes, with supportive measures to orient and calm the patient.

\synonyms
Psychosis ICU (ICU Psychosis)

\medterm{Idaiovirus} Likely a misspelling of Idaeovirus, a genus of plant-infecting RNA viruses that includes viruses associated with raspberry and privet; it is not a recognized human pathogen.

\medterm{Idarubicin} Idarubicin is an anthracycline chemotherapy drug that intercalates DNA and inhibits topoisomerase II, used mainly in acute myeloid leukaemia; myelosuppression and cumulative cardiotoxicity are important risks.

\medterm{IDC} Alternate index label; see \textit{Invasive ductal carcinoma}.

\medterm{Idea Of Reference} A belief that ordinary events, comments, or media messages have a special personal meaning, often seen in psychotic or mood disorders when held with strong conviction.

\medterm{Ideas Of Reference} The belief that ordinary events, comments, media messages, or other people’s actions have a special personal meaning. If the belief is held with strong conviction despite evidence, it may occur in a psychotic or mood disorder.

\medterm{Ideation} The formation or content of thoughts, especially thoughts related to a particular behavior or symptom. Clinical examples include suicidal ideation and homicidal ideation, which require assessment of intent, plan, means, and immediate safety.

\medterm{Identical Twins} Two individuals formed when one fertilized egg separates into two embryos. They usually have nearly identical DNA, but differences can develop through mutations, epigenetic changes, environment, and chance.

\medterm{Idiopathic} Describing a condition for which no specific cause has been identified after appropriate evaluation. It does not mean that the condition is imaginary, unimportant, or untreatable; it means that its cause is currently unknown.
\medterm{Idiopathic Facial Paralysis (Facial Nerve Problems)} Alternate index label for Facial Nerve Problems.

\medterm{Idiopathic Hypercalciuria} Excess calcium in the urine without an identified cause. It can increase the risk of kidney stones and, less often, bone loss.

\medterm{Idiopathic Hypersomnia} A sleep disorder that causes extreme daytime sleepiness, long or unrefreshing sleep, and difficulty waking up, even when a person has enough time to sleep. Other causes must be ruled out.

\medterm{Idiopathic Interstitial Pneumonias} A group of interstitial lung diseases with no identified cause, characterized by varying degrees of lung inflammation and fibrosis; diagnosis requires clinical, radiologic, and sometimes pathologic assessment.

\medterm{Idiopathic Intracranial Hypertension} Increased pressure around the brain without a tumor or other clear cause. It can cause headaches, a pulsing sound in the ears, and vision problems and may damage sight if untreated.

\medterm{Idiopathic Myelofibrosis} Alternate index label; see \textit{Myelofibrosis}.

\medterm{Idiopathic Pulmonary Fibrosis} A long-term lung disease of unknown cause in which scar tissue gradually builds up in the lungs. It causes increasing shortness of breath and a dry cough and requires specialist care.

\medterm{Idiopathic Thrombocytopenic Purpura} An older name for immune thrombocytopenia, a disorder in which the immune system destroys platelets and may cause easy bruising, tiny skin spots, nosebleeds, heavy menstrual bleeding, or other bleeding.
\medterm{Idiopathic Toe Walking} Alternate index label; see \textit{Toe walking in children}.

\medterm{Idiosyncrasy} An unusual individual reaction to a medicine or exposure that is not explained by the usual pharmacologic effect.

\medterm{Idiosyncratic Reaction} An uncommon, unexpected, and often genetically determined adverse drug reaction that is unrelated to the drug's known pharmacology and usually dose-independent (a Type B reaction).
\medterm{Idiotypic} Relating to the unique antigen-binding features of an antibody or immune-cell receptor. These features can be recognized by other antibodies and are useful in studying immune responses and some blood cancers.

\medterm{Idioventricular Rhythm} A slow heart rhythm that starts in the lower chambers of the heart instead of the normal pacemaker. It may occur when the usual electrical signal is blocked or when the heart is severely ill.

\medterm{Idoxuridine} Idoxuridine is a thymidine analogue antiviral formerly used mainly topically for herpes simplex keratitis; it inhibits viral DNA synthesis but can irritate or damage the corneal epithelium.

\medterm{Ifosfamide} Ifosfamide is an alkylating chemotherapy drug that forms DNA cross-links and is used for selected sarcomas and other cancers; haemorrhagic cystitis, encephalopathy, infertility, and myelosuppression require monitoring and mesna protection.

\medterm{iFR} Alternate name; see \textit{Instantaneous Wave-Free Ratio}.

\medterm{Ig} Abbreviation for immunoglobulin, an antibody protein produced by B cells and plasma cells; major classes include IgG, IgA, IgM, IgE, and IgD.

\medterm{IgA} A type of antibody found mainly in body fluids and linings, including saliva, tears, breast milk, and the digestive and breathing passages. It helps protect these surfaces from infection.

\medterm{IgA Nephropathy} A kidney disease in which IgA antibodies build up in the kidney's filtering units. It can cause blood or protein in the urine, high blood pressure, and sometimes kidney failure.

\medterm{IgA Nephropathy (Berger Disease)} A kidney disease caused by deposits of immunoglobulin A in the kidney filters, potentially leading to blood in urine, protein loss, high blood pressure, or kidney failure.

\synonyms
Berger's Disease

\medterm{IgA Vasculitis} An illness in which inflammation damages small blood vessels, often causing a raised purple rash, joint pain, stomach pain, and kidney problems. It is most common in children and may follow an infection.

\synonyms
Anaphylactoid Purpura

\medterm{IgA Vasculitis (Henoch-Schonlein Purpura)} A disease in which the immune system inflames small blood vessels, usually after an infection. It can cause purple spots on the legs, joint pain, abdominal pain, and kidney problems, especially in children.

\synonyms
Henoch-Schonlein Purpura

\medterm{IgD} A less common type of antibody found mainly on the surface of certain immune cells. It helps those cells recognize signals, but its full role in human disease is not yet clear.

\medterm{IgE} A type of antibody involved in allergies and defense against some parasites. When it reacts to an allergen, it can trigger symptoms such as sneezing, hives, asthma, or anaphylaxis.

\medterm{IgG} The most common type of antibody in the blood. It helps neutralize germs and toxins, supports immune protection after infection or vaccination, and can cross the placenta to protect a developing baby.

\medterm{IgG And Complement Deposition At The Dermis-Epidermis Junction} An immunofluorescence finding that may occur in the lupus band test, in which immune
complexes containing immunoglobulin and complement deposit along the dermoepidermal
junction. The pattern and clinical context are needed for interpretation.

\medterm{IgG4-Related Disease} An uncommon immune disorder that causes inflammation and scar-like growths in one or more organs. It may affect the pancreas, salivary glands, eyes, kidneys, lungs, or blood vessels; diagnosis usually requires imaging, blood tests, and sometimes a biopsy.

\medterm{IgG4-Related Disease Arteritis} Inflammation of a large or medium-sized artery as part of IgG4-related disease. It may thicken or weaken the artery and can cause narrowing, an aneurysm, pain, or reduced blood flow.

\medterm{IgG4-Related Ophthalmic Disease} An immune-mediated fibroinflammatory disorder that can affect the lacrimal glands, orbit, eyelids, or
extraocular muscles. It may cause swelling, mass lesions, double vision, or visual
symptoms and can occur with IgG4-related disease in other organs.

\medterm{Igh Gene} A gene region that helps make the heavy part of an antibody. Immune cells rearrange these gene segments so the body can produce antibodies that recognize many different targets.
\medterm{IgM} A large type of antibody that is often made early during an infection. It helps clump germs and activate other parts of the immune system; a blood test may suggest recent exposure but must be interpreted with other results.

\medterm{IIH} Alternate index label; see \textit{Pseudotumor cerebri (idiopathic intracranial hypertension)}.

\medterm{Ikappab Kinase} A group of enzymes that helps turn on genes involved in inflammation and immune responses.

\medterm{Ileal Atresia} A condition present at birth in which part of the ileum, the last section of the small intestine, is closed or missing. It causes bowel blockage in a newborn, with a swollen abdomen, green vomiting, and failure to pass the first stool.

\medterm{Ileal Conduit} A surgically created passage that diverts urine through a segment of ileum to an opening in the abdominal wall.

\medterm{Ileal Fistula} An abnormal connection involving the ileum and another organ, bowel segment, or skin surface, which may cause intestinal contents to leak and lead to infection, malnutrition, or fluid loss.

\medterm{Ileal Hemorrhage} Bleeding from the ileum, the distal portion of the small intestine, caused by ulceration, inflammation, vascular lesions, tumors, trauma, or other disease.

\medterm{Ileal Neoplasm} An ileal neoplasm is an abnormal growth in the ileum, the last part of the small intestine; it may be benign or cancerous.

\medterm{Ileal Obstruction} Partial or complete blockage of the ileum, the final portion of the small intestine. Causes include
adhesions, hernia, inflammatory bowel disease, tumors, and twisting of the bowel. Symptoms may
include cramping pain, vomiting, abdominal distension, and inability to pass stool or gas.

\medterm{Ileal Pouch-Anal Anastomosis} An operation that removes the colon and rectum and creates a pouch from the end of the small intestine. The pouch is connected to the anus so bowel movements can continue through the usual route.

\medterm{Ileal Stenosis} Narrowing of the ileum that may impair passage of intestinal contents and result from Crohn disease, inflammation, scarring, ischemia, tumor, or surgery.

\medterm{Ileal Ulcer} An ileal ulcer is a break in the lining of the ileum caused by inflammation, infection, medicines, ischemia, trauma, or another intestinal disorder.

\medterm{Ileitis} Inflammation of the ileum, the distal segment of the small intestine. Causes include
Crohn disease, infection, ischemia, medication-related injury, and other inflammatory
disorders; symptoms may include right-lower-quadrant pain, diarrhea, fever, or bleeding.

\medterm{Ileitis (Crohn'S Disease)} Alternate index label for Crohn's Disease.

\medterm{Ileitis, Crohn} Crohn disease involving the ileum, the final part of the small intestine, causing chronic or relapsing inflammation that may produce abdominal pain, diarrhea, weight loss, strictures, or fistulas.

\medterm{Ileocaecal} Ileocaecal means relating to the junction between the ileum and the caecum, where the small intestine joins the large intestine.

\medterm{Ileocecal Valve} A sphincter-like fold at the junction of the terminal ileum and caecum. It regulates
passage of intestinal contents into the colon and helps limit reflux of colonic contents
into the small intestine.

\medterm{Ileocolitis, Crohn} Crohn disease affecting both the ileum and colon, causing patchy chronic inflammation that may lead to abdominal pain, diarrhea, weight loss, strictures, fistulas, or perianal disease.

\medterm{Ileostomy} An operation that brings the end of the small intestine through an opening in the abdomen so stool drains into an external pouch. Output is often liquid, so adequate fluids, salt replacement, and skin care are important.







\medterm{Ileum} The final and longest part of the small intestine. It absorbs vitamin B12, bile salts, and remaining nutrients before food passes into the large intestine.

\medterm{Ileus} Failure of coordinated intestinal propulsion without a mechanical blockage, usually from
postoperative bowel inhibition, inflammation, metabolic disturbance, or medication
effect. It causes distension, nausea, vomiting, and reduced passage of flatus or stool.

\medterm{Iliac} Relating to the ilium or the iliac region, including the iliac bones and iliac blood vessels.

\medterm{Iliac Artery} One of the large arteries that carries blood to the pelvis and lower limb. The common iliac arteries divide into internal branches for the pelvic organs and external branches that continue into the legs.

\medterm{Iliac Crest} The iliac crest is the curved upper border of the ilium and a landmark and attachment site for abdominal, back, and pelvic muscles.

\medterm{Iliac Vein} An iliac vein drains blood from the pelvis and lower limb toward the inferior vena cava; thrombosis can cause swelling and pulmonary embolism.

\medterm{Common Iliac Vein} A vein formed by the union of the internal and external iliac veins. The right and left common iliac veins unite to form the inferior vena cava.

\medterm{External Iliac Vein} A continuation of the femoral vein above the inguinal ligament that joins the internal iliac vein to form the common iliac vein.

\medterm{Internal Iliac Vein} A vein draining the pelvic organs, pelvic wall, and gluteal region that joins the external iliac vein to form the common iliac vein.

\medterm{Iliopsoas Muscles} The psoas major and iliacus muscles working together to bend the hip. They help lift the thigh, stabilize the spine, and maintain posture, and may cause groin or back pain when strained.

\synonyms
Iliopsoas

\medterm{Iliotibial Band} A strong band of tissue running along the outer side of the thigh from the hip to the knee. It helps stabilize the hip and knee during movement.

\medterm{Iliotibial Band Syndrome} Pain on the outer side of the knee caused by irritation of the iliotibial band, especially during repeated activities such as running or cycling.

\synonyms
ITBS (Iliotibial Band Syndrome)

\medterm{Ilium} Largest hip bone forming the superior portion. Body contributes to acetabulum; wing has
a concave iliac fossa. Iliac crest runs ASIS to PSIS, a palpable landmark for lumbar
puncture level. Sacroiliac joint articulates with sacrum. Gluteal muscles originate from
its external surface; iliacus fills its fossa.

\medterm{Illness Anxiety Disorder} Preoccupation with having or acquiring a serious illness, with minimal or no somatic
symptoms present. The individual has a high level of anxiety about health, performs
excessive health-related behaviors, and experiences maladaptive avoidance. Symptoms
persist for at least six months. The individual is not delusional and insight may vary.

\medterm{Illusion} A misperception or misinterpretation of a real external stimulus, such as mistaking a shadow for a person. Illusions may occur with visual impairment, delirium, seizures, intoxication, or severe psychiatric illness and differ from hallucinations, which occur without an external stimulus.

\medterm{Iloprost} A prostacyclin analogue that causes vasodilation and inhibits platelet aggregation; inhaled or infused formulations are used for selected forms of pulmonary hypertension and severe peripheral ischemia.

\medterm{Image View} An image view is a particular angle or projection used when producing a medical image, such as a radiograph.

\medterm{Image-Guided Ablation} Destruction of abnormal tissue using an energy source such as heat, cold, or chemicals while imaging guides placement and
monitors the procedure.

\medterm{Image-Guided Biopsy} Removal of a tissue sample while ultrasound, CT, MRI, mammography, or fluoroscopy guides the needle or instrument to the target and helps avoid critical structures.

\medterm{Image-Guided Drainage} Placement of a catheter or needle under imaging guidance to evacuate fluid, pus, blood,
or bile from a pathologic collection. Indications include abscess, biloma, urinoma,
pleural collection, and symptomatic cysts.

\medterm{Image-Guided Therapy} A minimally invasive procedure performed with imaging to guide instruments or treatment to a target tissue.

\medterm{IGRT} Image-guided radiation therapy uses medical imaging, such as CT, MRI, ultrasound, or X-ray, to position the patient and target radiation accurately before or during treatment. It helps deliver radiation to a tumor while reducing exposure to nearby healthy tissue.

\medterm{Imaging} The production of pictures of internal body structures using methods such as radiography, ultrasound,
computed tomography, magnetic resonance imaging, or nuclear medicine.

\medterm{Imaging And Radiology} Imaging and radiology use methods such as X-ray, ultrasound, CT, and MRI to create pictures that help diagnose or guide treatment of disease.

\medterm{Imaging Technique} An imaging technique is a method for producing pictures of body structures or processes, such as ultrasound, CT, MRI, or nuclear imaging.

\medterm{Imaging-Guided Biopsy} Alternate index label; see \textit{Image-guided biopsy}. Appropriate specimen handling, bleeding-risk assessment, informed consent, and post-procedure monitoring remain essential.

\medterm{Imatinib Mesylate} The mesylate salt of imatinib, a tyrosine-kinase inhibitor used to treat BCR-ABL-positive chronic myeloid leukemia, gastrointestinal stromal tumors with KIT alterations, and other selected cancers.

\medterm{Imerslund-GräSbeck Syndrome} A rare inherited disorder in which the small intestine cannot absorb vitamin B12 properly. It can cause low B12 levels, anemia, tiredness, numbness, balance problems, or other nerve damage despite adequate dietary intake.

\medterm{Imidazoles} A group of antifungal medicines, including clotrimazole, miconazole, and ketoconazole, used to treat fungal infections of the skin, mouth, or other body areas.

\medterm{Imide} An imide is a chemical compound containing nitrogen bonded to two carbonyl groups and may have pharmaceutical, metabolic, or industrial significance.

\medterm{Imines} Imines are compounds containing a carbon-nitrogen double bond and occur in chemical synthesis, metabolism, and drug chemistry.

\medterm{Imino Acid} A compound related to amino acids whose nitrogen is part of a ring or imine structure. Proline is commonly grouped with imino acids in biochemistry.
\medterm{Iminoacetic Acid Scan} A nuclear medicine scan that uses a small amount of radioactive medicine to show how bile moves through the liver, gallbladder, and bile ducts. It is also called a HIDA scan.

\medterm{Imipenem} A strong antibiotic, usually given with cilastatin, used for serious bacterial infections. It can cause allergic reactions, stomach problems, seizures, or antibiotic resistance.

\medterm{Imipramine} A medicine used for depression and some bladder problems. It can cause dry mouth, constipation, dizziness, abnormal heart rhythms, and dangerous effects in overdose.

\medterm{Imipramine Overdose} Toxic ingestion of imipramine, a tricyclic antidepressant, which can cause anticholinergic effects, seizures, hypotension, ventricular arrhythmias, coma, and life-threatening cardiac toxicity.

\medterm{Imiquimod} Imiquimod is a topical immune-response modifier that activates Toll-like receptor 7 and induces cytokines; it treats selected actinic keratoses, superficial basal-cell carcinomas, and external genital warts, often causing local inflammation.

\medterm{Immature} Immature means not fully developed or differentiated, such as an immature cell, organ, tissue, or developmental stage.

\medterm{Immature Ovarian Teratoma} A malignant germ-cell tumor containing immature embryonal tissues, usually diagnosed in adolescents or young adults and graded by the amount of immature neuroepithelium.

\medterm{Immaturity} Development that is less physically, neurologically, emotionally, or socially advanced than expected for a person's age or stage of life.

\medterm{Immediate Hypersensitivity} A rapid allergic reaction, usually involving IgE antibodies and mast cells. It can cause hives, itching, swelling, wheezing, vomiting, or anaphylaxis within minutes of exposure.

\medterm{Immersion Foot (Trench Foot)} A painful cold injury caused by prolonged exposure of the feet to cold, wet conditions without freezing, leading to numbness, swelling, and skin changes.

\synonyms
Trench Foot

\medterm{Immobilization} Restriction of movement of a body part or the whole patient, using rest, splints, casts, braces, traction, or other methods. It may protect an injury but can cause stiffness, pressure injury, muscle loss, venous thromboembolism, and other complications if prolonged.

\medterm{Immobilize} To restrict movement of a body part, often with a splint or cast, to support healing.

\synonyms
Immobilization

\medterm{Immune} Having an immune response that recognizes and helps protect against a particular antigen, infection, or toxin; immunity may be innate or acquired.

\medterm{Immune And Infectious Diseases} A medical field encompassing disorders of the immune system and illnesses caused by
pathogenic microorganisms. Immune system disorders include autoimmune diseases (e.g.,
lupus, rheumatoid arthritis), immunodeficiencies, and hypersensitivity reactions.

\medterm{Immune Checkpoint} A control signal that helps prevent immune cells from becoming too active or damaging healthy tissue. Some cancer medicines block these signals to help immune cells attack tumors, but this can cause inflammation in normal organs.

\medterm{Immune Complex} A cluster formed when an antibody attaches to a foreign or abnormal substance. If these clusters are not cleared, they can collect in tissues and cause inflammation, blood-vessel damage, or kidney disease.

\medterm{Immune Evasion} The ability of a pathogen, tumor, or other abnormal cell to avoid recognition or destruction by the immune system.

\medterm{Immune Hemolytic Anemia} Anemia caused by immune-mediated destruction of red blood cells, either by
autoantibodies or antibodies directed against foreign red-cell antigens. It may cause
fatigue, jaundice, dark urine, splenomegaly, and laboratory evidence of hemolysis.

\medterm{Immunohemolysis} Destruction of red blood cells by an immune reaction involving antibodies, complement, or both. It may occur in autoimmune hemolytic anemia, incompatible transfusion reactions, or other antibody-mediated disorders.

\medterm{Immune Response} The body’s reaction to an infection, foreign substance, or abnormal cell. It includes rapid defenses such as inflammation and slower, targeted responses involving antibodies and immune cells that remember previous exposures.

\medterm{Immune Status} Immune status describes the functional state of a person's immune defenses, including immunocompetence, prior infection or vaccination, immunosuppression, and susceptibility to infection.

\medterm{Immune System} The body's network of cells, tissues, organs, and proteins that recognizes and responds to infections, abnormal cells, and harmful substances. It includes rapid innate defenses and learned adaptive defenses; overactivity can cause allergy or autoimmunity, while weakness increases infection risk.

\medterm{Immune Thrombocytopenia} An autoimmune disorder in which the immune system destroys platelets or reduces their production. Low platelets can cause easy bruising, tiny red or purple skin spots, nosebleeds, heavy periods, or other bleeding; causes may be primary or related to another illness or medicine.

\medterm{Immune Thrombocytopenic Purpura} An older name for immune thrombocytopenia, in which the immune system destroys platelets. It may cause bruising, pinpoint skin spots, nosebleeds, heavy menstrual bleeding, or no symptoms.

\medterm{Immune Thrombocytopenic Purpura (Idiopathic Thrombocytopenic Purpura)} Alternate index label for Idiopathic Thrombocytopenic Purpura.

\medterm{Immune Tolerance} The ability of the immune system to avoid attacking the body’s own tissues while still responding to infections. Failure of tolerance can contribute to autoimmune disease; excessive tolerance can allow some cancers or infections to persist.

\medterm{Immune-Mediated Necrotizing Myopathy} A subacute severe proximal myopathy with very high CK (often >10,000), minimal
inflammation on biopsy, necrotic fibers. Two major autoantibodies: anti-HMGCR
(3-hydroxy-3-methylglutaryl-CoA reductase, statin-associated or spontaneous) and
anti-SRP (signal recognition particle).

\medterm{Immune-Related Adverse Event} An inflammatory or autoimmune toxicity caused by immune checkpoint blockade or another
immune-activating cancer treatment. It may affect skin, bowel, liver, endocrine organs,
lungs, heart, kidneys, or nervous system and often requires treatment-specific
immunosuppression.

\medterm{Immunity} The ability of the body to recognize and resist infectious agents or other foreign substances through innate or adaptive defenses.

\medterm{Immunization} The process of inducing protection against an infection, usually through vaccination. Recommended
vaccines and schedules depend on age, previous doses, pregnancy, immune status, medical
conditions, travel, exposure risk, and current public-health guidance.

\medterm{Immunizations For People With Diabetes} Vaccination is especially important in diabetes because infections may be more severe; recommendations depend on age, health status, vaccines received, and local guidance.

\medterm{Immunoassay} An immunoassay is a laboratory test that uses antibodies to detect or measure a substance such as a hormone, drug, antigen, or antibody.

\medterm{Immunoblast Count} An immunoblast count measures activated lymphocytes, usually in a tissue or blood sample, and may help evaluate immune responses or lymphoid disease.

\medterm{Immunoblotting} Immunoblotting, or Western blotting, separates proteins by electrophoresis and detects a target with specific antibodies; band pattern and intensity are interpreted with controls and clinical context.

\medterm{Immunochemistry} The study or laboratory measurement of biological molecules using specific antigen-antibody reactions, including immunoassays and immunostaining.

\medterm{Immunocompetence} Immunocompetence is the capacity of the immune system to recognize antigens and mount an effective, appropriately regulated response; deficiency increases infection risk and may impair vaccine responses.

\medterm{Immunocompetent} Having an immune system capable of mounting an appropriate response to common infectious or foreign antigens.

\medterm{Immunocompromised} Having weakened or impaired immune defenses because of disease, medication, cancer treatment, transplantation, or another cause, increasing susceptibility to some infections and complications.

\medterm{Immunocompromised Host} A person whose immune defenses are impaired by congenital disease, infection,
malignancy, immunosuppressive treatment, organ failure, or another condition. Such
patients have increased susceptibility to opportunistic infection, severe disease, and
atypical presentations.

\medterm{Immunoconjugate} A therapeutic or diagnostic molecule in which an antibody or other targeting protein is chemically linked to a drug, toxin, enzyme, or imaging agent.

\medterm{Immunocytochemistry} A laboratory method using labeled antibodies to identify specific proteins or antigens within cells or tissue sections.

\medterm{Immunodeficiency} Impaired development or function of one or more components of the immune system,
resulting in increased susceptibility to infection, malignancy, or autoimmunity. It may
be primary and genetic or secondary to disease, medication, malnutrition, or another
acquired cause.

\medterm{Immunodeficiency And Infections} Immunocompromised individuals are at increased risk for opportunistic infections.
Primary immunodeficiencies are genetic disorders affecting immune function. Secondary
immunodeficiencies include HIV/AIDS, organ transplantation, chemotherapy, and
corticosteroid use.

\medterm{Immunodeficiency Diseases} Disorders in which one or more components of the immune system are absent, deficient, or
dysfunctional. They may be inherited or acquired and can cause recurrent, severe,
unusual, or persistent infections, autoimmune disease, and increased malignancy risk.

\medterm{Immunodeficiency Disorders} Conditions in which part of the immune system is missing, weak, or does not work properly. They may be inherited or acquired and can cause repeated, severe, unusual, or persistent infections and sometimes autoimmune disease or cancer.

\medterm{Immunoelectrophoresis} A laboratory method that separates proteins by electrophoresis and identifies them with specific antibodies, often to study immunoglobulins.

\medterm{Immunoelectrophoresis - Urine} A laboratory technique that separates and identifies urinary proteins with specific antibodies, used to investigate abnormal immunoglobulin excretion and monoclonal light chains.

\medterm{Immunoelectrophoresis – Blood} A laboratory technique that separates blood proteins by electrophoresis and identifies them with specific antibodies, used historically to evaluate abnormal immunoglobulins and monoclonal proteins.

\medterm{Immunofixation - Urine} A urine test that uses electrophoresis and antibodies to identify and type monoclonal immunoglobulins or light chains, including Bence Jones proteins.

\medterm{Immunofixation Blood Test} A blood test that identifies and types monoclonal immunoglobulins, helping evaluate monoclonal gammopathy, multiple myeloma, and related plasma-cell disorders.

\medterm{Immunofluorescence} A technique using fluorescently labeled antibodies to detect specific antigens in
tissues, cells, or microorganisms; visualized by fluorescence microscopy.

\medterm{Immunofluorescence Technique} An immunofluorescence technique uses fluorescently labelled antibodies to locate specific proteins, organisms, or immune deposits in a sample.

\medterm{Immunogenetics} The study of how genes influence immune-system development, function, and disease. It includes inherited susceptibility to autoimmune disease, differences in immune responses, and genetic matching for transplantation.

\medterm{Immunogenicity} The ability of an antigen, vaccine, or other substance to provoke an immune response, including antibody or cellular responses.

\medterm{Immunoglobulin} An antibody protein made by immune cells that recognizes specific foreign substances and helps block or remove infections and other harmful targets.

\medterm{Immunoglobulin (Antibody)} A protein made by immune cells that recognizes a particular foreign substance, germ, or abnormal cell. Antibodies can block harmful targets, mark them for removal, or activate other immune defenses.

\synonyms
Antibody

\medterm{Immunoglobulin A} An antibody found mainly in mucus, saliva, tears, breast milk, and intestinal fluid. It protects surfaces exposed to the outside world and helps prevent germs from attaching and entering tissues.

\medterm{Immunoglobulin D} A type of antibody found in very small amounts in the blood and on the surface of some B cells. Its exact role is not fully understood, but it may help regulate immune responses in the airways and other tissues.

\medterm{Immunoglobulin E} An antibody involved in many immediate allergic reactions and defense against some parasites. It activates mast cells, which release chemicals causing itching, swelling, mucus, wheezing, or other allergy symptoms.

\synonyms
IgE

\medterm{Immunoglobulin G} The most common antibody in blood and tissue fluid, helping protect against many infections. It can remain after infection or vaccination, cross the placenta during pregnancy, and sometimes cause harmful immune reactions.

\medterm{Immunoglobulin Idiotype} The distinctive part of an antibody that determines what it recognizes. Different immune-cell families make antibodies with different idiotypes; the concept is mainly used in laboratory research and some cancer treatments.

\medterm{Immunoglobulin Isotype} The class of an antibody, determined by its constant region. The main isotypes are IgG, IgA, IgM, IgD, and IgE; each has different roles and locations in immune defense.

\medterm{Immunoglobulin M} The first major antibody produced during an initial immune response.

\medterm{Immunoglobulin Replacement Therapy} Treatment that supplies pooled antibodies through a vein or under the skin when a person's immune system makes too few effective antibodies. It can reduce recurrent infections in selected inherited or acquired antibody deficiencies; headache, fever, and infusion reactions may occur.

\medterm{Immunohistochemistry} A laboratory method that uses specially marked antibodies to find particular proteins in a tissue sample. It helps identify tumor type, infection, or immune disease and can guide treatment choices.

\medterm{Immunologic Adjuvant} A substance added to some vaccines to strengthen or guide the immune response. It may improve protection but can also increase temporary pain, swelling, fever, or other vaccine reactions.
\medterm{Immunologic Factor} A substance, cell, gene, or process that affects immune-system activity or control. Examples include antibodies, immune cells, signaling proteins, inherited traits, medicines, and conditions that alter defense against disease.

\medterm{Immunologic Memory} The capacity of the adaptive immune system to respond more rapidly and strongly after re-exposure to an antigen because of long-lived memory lymphocytes.

\medterm{Immunologic Skin Test} An immunologic skin test places a small amount of an antigen or allergen in or on the skin to assess immune reactivity, such as allergy or tuberculosis exposure.

\medterm{Immunologic Surveillance} The immune system’s ongoing detection and removal of abnormal, infected, or damaged cells. This protection helps control infections and may recognize early cancerous changes, although it is not always successful.

\medterm{Immunologic Technique} A laboratory or clinical method that uses antibodies, antigens, or immune reactions to detect or measure a substance.

\medterm{Immunologic Test} An immunologic test examines immune cells, antibodies, antigens, or immune-related proteins to help evaluate infection, allergy, autoimmune disease, or immune deficiency.

\medterm{Immunological Memory} The ability of the adaptive immune system to mount a faster, larger, and more effective
response upon re-exposure to an antigen, conferred by long-lived memory B cells and
memory T cells generated during the primary response.

\medterm{Immunological Synapse} The organized contact point where an immune cell meets a target or another immune cell. Signals gathered at this contact help decide whether the immune cell should become active.

\medterm{Immunological Tolerance} The reduced or absent immune response to a specific antigen, including the normal lack of attack on the body’s own tissues.

\medterm{Immunologically Privileged Site} A body area where immune responses are partly limited, helping protect delicate tissues from inflammation while still allowing important defense against infection.

\medterm{Immunology} The medical and biological study of the immune system, including its normal defenses and disorders such as allergy, autoimmunity, immunodeficiency, infection, and transplant rejection.
\medterm{Immunomodulation} Modification of immune activity to enhance, suppress, or redirect an immune response. Immunomodulatory medicines are used in cancer, autoimmune disease, transplantation, and infection but may increase susceptibility to infection or other adverse effects.

\medterm{Immunomodulators In Dermatology} Topical immunomodulators treat inflammatory skin diseases without corticosteroid side
effects. Pimecrolimus 1 percent cream and tacrolimus 0.03 and 0.1 percent ointments are
calcineurin inhibitors used for atopic dermatitis, facial eczema, and intertriginous
areas. They are steroid-sparing agents.

\medterm{Immunophenotyping} Identification and characterization of cells according to their surface, cytoplasmic, or
nuclear antigens, usually by flow cytometry or immunohistochemistry. It is used to
classify leukemias, lymphomas, immune-cell subsets, and some immunodeficiency disorders.

\medterm{Immunoprecipitation} A laboratory technique that uses an antibody to bind and isolate a specific antigen or protein from a mixture. It is used in research and some validated diagnostic or analytical methods.

\medterm{Immunoprotein} A protein involved in immune defense or immune regulation, such as an antibody, complement protein, cytokine, or antigen-presenting molecule. Immunoproteins identify threats, coordinate immune cells, or help remove damaged or infected material.

\medterm{Immunoradiometric Assay} A laboratory test using radioactive labels on antibodies to measure a specific substance, such as a hormone, protein, or other marker.

\medterm{Immunoregulation} Control of the strength, duration, and target of immune responses, including activation, suppression, and tolerance. Abnormal immunoregulation contributes to autoimmunity, allergy, immunodeficiency, and chronic inflammation.

\medterm{Immunosenescence} Age-related changes in innate and adaptive immunity that alter infection susceptibility,
vaccine responses, inflammation, and cancer surveillance. It is influenced by chronic
disease, nutrition, medications, and functional status.

\medterm{Immunosuppression} Reduced immune-system activity caused by disease, medicines, cancer treatment, or another condition. It can prevent rejection of a transplant or control autoimmune disease but increases susceptibility to infections and some cancers.
\medterm{Immunosuppression In Transplant} Medications to prevent rejection. Induction with thymoglobulin or basiliximab.
Maintenance with tacrolimus, mycophenolate, and prednisone. Monitor drug levels and
renal function. Calcineurin inhibitor nephrotoxicity is major concern requiring dose
optimization.

\medterm{Immunosuppression Screening} Pre-treatment evaluation for latent or active infection, vaccination status, malignancy
risk, and organ dysfunction before immunosuppressive therapy. Testing commonly includes
tuberculosis, hepatitis B and C, HIV, blood counts, liver and kidney function, and
pregnancy when relevant.

\medterm{Immunosuppressive} Describing a medicine, treatment, or condition that reduces immune-system activity. Immunosuppression can prevent transplant rejection or control autoimmune disease, but it may increase the risk of infection, some cancers, and reduced vaccine responses.

\medterm{Immunotherapy} Cancer treatment that helps the immune system recognize or attack cancer cells, or supplies immune cells or antibodies that do so. Main approaches include checkpoint medicines, engineered immune cells, antibodies, and immune-stimulating treatments. Benefits and side effects depend on the cancer and treatment.

\medterm{Immunotherapy And Biologic Treatments} Advanced therapies that harness or modify the immune system to treat diseases,
particularly cancer, autoimmune disorders, and inflammatory conditions. Biologics are
large-molecule drugs derived from living organisms that target specific immune pathways,
such as TNF inhibitors, monoclonal antibodies, and checkpoint inhibitors.

\medterm{Immunotherapy For Cancer} Cancer treatment that stimulates, restores, redirects, or supplies immune activity to recognize and attack malignant cells, including immune-checkpoint inhibitors, cellular therapies, and antibodies.

\medterm{Immunotherapy In Dermatology} Cancer immunotherapy uses medicines that alter immune activity against skin cancers.
Checkpoint inhibitors target pathways such as PD-1, PD-L1, or CTLA-4. Their use depends
on the tumor type, stage, prior treatment, and the patient’s health.


\medterm{Immunotoxicity} Harm to immune-system function caused by chemicals, medicines, infections, or other exposures. It may weaken defenses, trigger excessive sensitivity, cause autoimmunity, or change responses to vaccines and infections.

\medterm{Immunotoxin} A targeted toxin made by linking a toxic protein or drug to an antibody or binding molecule that directs it toward selected cells.

\medterm{Immunoturbidimetry} A laboratory method that measures cloudiness produced when antibodies react with a target antigen, allowing estimation of the target’s concentration.

\medterm{Imp Dehydrogenase} IMP dehydrogenase converts inosine monophosphate to xanthosine monophosphate in guanine-nucleotide synthesis and is a target of some immunosuppressive drugs.

\medterm{Impact} The forceful contact or collision of one object or body part with another, or the effect a disease, treatment, or event has on health; context determines the meaning.

\medterm{Impacted} Firmly wedged or unable to move from its normal position, as with an impacted tooth or stool. The result may cause pain, blockage, inflammation, infection, or damage to nearby tissues.
\medterm{Impacted Fracture} A fracture in which one broken bone fragment is driven firmly into another, often producing relative stability but sometimes causing shortening or deformity.

\medterm{Impacted Tooth} Tooth unable to erupt into functional position. Most common: mandibular third molars,
maxillary canines. May cause pathology: cyst, root resorption, caries. Treatment:
monitoring or extraction.

\medterm{Impacted Wisdom Teeth} Wisdom teeth that cannot emerge normally because they are blocked by bone, gum, or another tooth.

\synonyms
Wisdom Teeth, Impacted

\medterm{Impaction} Blockage caused by material that is firmly lodged, such as hardened stool, a foreign object, or an impacted tooth. Symptoms depend on the site and may include pain, swelling, vomiting, or inability to pass contents.
\medterm{Impaired Fasting Glucose} A fasting blood-glucose level above normal but below the range used to diagnose diabetes. It indicates increased future diabetes risk and may improve with weight management, activity, and treatment of risk factors.
\medterm{Impaired Glucose Tolerance} Blood glucose above the normal range after a glucose challenge but below the diagnostic threshold for diabetes.

\medterm{Impedance} Impedance is the opposition to the flow of an alternating electrical current; in medicine it can help assess tissue properties or body fluids.

\medterm{Impedance Hearing Testing} A group of tests, including tympanometry, that measure the movement of the eardrum and middle-ear function.

\medterm{Impedance Plethysmography} A noninvasive method that detects changes in limb electrical impedance caused by changes in blood volume. It has been used to assess venous blood flow and suspected deep-vein thrombosis, although duplex ultrasonography is generally preferred for current diagnostic evaluation.

\medterm{Impella Device} Percutaneous microaxial left ventricular assist device that actively pumps blood from
the LV into the aorta, reducing LV workload and increasing cardiac output. Types:
Impella 2.5 (up to 2.5 L/min), Impella CP (up to 3.7 L/min), Impella 5.0 (up to 5 L/min,
surgically placed).

\medterm{Imperforate Anus} A congenital anomaly where the anus is completely absent or fails to open at the
perineum. Classified as low (distal to the levator ani complex) or high (proximal to the
levator ani). Often associated with other anomalies including vertebral, anal, cardiac,
tracheoesophageal, renal, and limb defects.

\medterm{Imperforate Anus Repair} Surgery to create or reposition an anal opening and connect the rectum so stool can pass normally in a baby born with an absent or blocked anus. Some babies need temporary bowel diversion; follow-up assesses healing, constipation, and bowel control.

\medterm{Imperforate Hymen} A congenital obstruction of the vaginal opening by an intact hymen without a normal
central opening. After puberty it may cause absent menstruation, cyclic pelvic pain,
urinary retention, or a pelvic mass from retained menstrual blood.

\medterm{Imperforate Vagina} Imperforate vagina is a rare congenital absence or obstruction of the vaginal canal, often associated with genital outflow obstruction, cyclical pelvic pain, haematocolpos, and urinary or reproductive complications.

\medterm{Impetigo} Impetigo is a superficial contagious bacterial skin infection producing fragile blisters or honey-colored crusts, commonly caused by Staphylococcus aureus or Streptococcus pyogenes.
affecting children.

\synonyms
School Sores (Impetigo)


\medterm{Impingement Syndrome} Pain and restricted movement caused when a tendon or other soft tissue is compressed between bones, commonly in the shoulder during elevation of the arm.

\medterm{Implant} A medical device, tissue, or material placed inside the body to replace, support, monitor, or deliver treatment.

\medterm{Implant (Dental)} A biocompatible fixture placed in the jawbone to support a dental prosthesis. It may be used for one missing tooth or to support a partial or complete denture; complications include infection, bone loss, and mechanical failure.

\synonyms
Dental

\medterm{Implant Failure} Loss of an implant’s intended structural, biologic, or functional performance. Causes
include infection, fracture, loosening, wear, malposition, thrombosis, material
degradation, and patient or tissue factors.

\medterm{Implant Sterilization} Validated processing that removes or inactivates all viable microorganisms, including
bacterial spores, from an implant. Methods include moist heat, ethylene oxide,
radiation, and low-temperature plasma according to material and device requirements.

\medterm{Implant-Associated Infection} Infection involving an implanted medical device or the tissue surrounding it. Microbial
adherence and biofilm formation can impair antimicrobial penetration and often require
device-directed therapy, source control, or removal.

\medterm{Implant-Supported Prosthesis} A fixed or removable dental prosthesis that is supported or retained by one or more dental implants.

\medterm{Implantable Cardioverter Defibrillator} An implanted device that detects dangerous ventricular rhythms and treats them with antitachycardia pacing or an electrical shock.


\medterm{Implantable Cardioverter-Defibrillator} An implanted device that monitors heart rhythm and delivers electrical therapy
(antitachycardia pacing or defibrillation) for life-threatening ventricular arrhythmias.

\medterm{Implantable Defibrillator} A surgically implanted device that detects dangerous ventricular arrhythmias and delivers pacing or an electrical shock to restore an effective rhythm.

\medterm{Implantable Loop Recorder} Subcutaneous device providing continuous ECG monitoring for up to 3 years. Used for
evaluation of unexplained syncope, recurrent palpitations, cryptogenic stroke, and
occult atrial fibrillation detection. Small size (coin-sized), inserted under local
anesthesia in the chest.

\medterm{Implantable Port} A subcutaneously placed vascular-access device consisting of a reservoir connected to a
catheter, used for repeated administration of medication, fluids, blood products, or
parenteral nutrition. It requires sterile access and may be complicated by infection,
thrombosis, occlusion, or mechanical failure.

\medterm{Implantation} The embedding of the blastocyst into the endometrium, occurring approximately 6-10 days
after ovulation. The blastocyst hatches from the zona pellucida, and the trophoblast
invades the decidua. Implantation triggers hCG production, which maintains the corpus
luteum. Failure of implantation is a common cause of early pregnancy loss.

\medterm{Implantation Procedure} Placement of a device, tissue, medicine-delivery system, or other material inside the body for treatment or support. Examples include joint, dental, heart, and breast implants; risks include bleeding, infection, movement, failure, or rejection of donor tissue.

\medterm{Implants Endometrial (Endometriosis)} Alternate index label for Endometriosis.

\medterm{Importance Weight} An importance weight is a statistical weight used to adjust observations so a sample better represents a target population or distribution.

\medterm{Impotence (Erectile Dysfunction [ED])} The ongoing inability to get or keep an erection firm enough for sexual activity. It may result from blood-vessel disease, nerve problems, medicines, hormone changes, or emotional factors.
\synonyms
Erectile dysfunction

\medterm{Impotence (Erectile Dysfunction (ED, Impotence))} Alternate index label for Erectile Dysfunction (ED, Impotence).

\medterm{Impression} Summary of key findings and recommendations in radiology report. The concise conclusion of a medical or radiology report, summarizing the most important findings, likely meaning, and recommendations for next steps.
\medterm{Imprinting} A pattern of gene expression determined by whether a gene copy was inherited from the mother or father.

\medterm{Impulse} An impulse is a brief signal or force; in physiology it commonly refers to an electrical signal travelling along a nerve or through the heart.

\medterm{Impulse Oscillometry} A pulmonary function testing technique that measures respiratory impedance by applying
small pressure oscillations during tidal breathing. It assesses large airway resistance
(R5), small airway resistance (R5-R20), and reactance.

\medterm{Impulsive Behavior} Impulsive behavior is action with limited forethought or inhibition despite possible harm and may occur with psychiatric, neurologic, substance-related, or developmental conditions.

\medterm{Impulsivity} A tendency to act quickly without adequate consideration of consequences or difficulty delaying a response; it can occur in ADHD, mania, substance use, and other conditions.

\medterm{In} In clinical measurements, an abbreviation for inch or inches; the same letters may have other meanings in different contexts.

\medterm{In Situ} In its original place. In pathology, carcinoma in situ consists of abnormal or cancer-like cells that remain within the surface layer and have not invaded deeper tissue. It has not spread through blood or lymph vessels, but some lesions can become invasive if untreated.

\synonyms
In place

\medterm{In Situ Carcinoma} Abnormal malignant cells confined to the tissue layer in which they arose and not yet invading through the basement membrane.

\medterm{In Situ Hybridization} A laboratory method that uses labeled nucleic-acid probes to detect specific DNA or RNA sequences within cells or tissue. It can identify gene copy changes, translocations, viral material, or tissue-specific gene expression.

\medterm{In Vitro} In vitro means performed outside a living organism, such as in a test tube, culture dish, or laboratory system.

\medterm{In Vitro Fertilisation} An assisted reproductive technology in which eggs are retrieved from the ovaries,
fertilized with sperm in a laboratory, and resulting embryos are transferred to the
uterus.

\medterm{In Vitro Fertilization} In vitro fertilization fertilizes retrieved eggs in a laboratory and transfers resulting embryos to the uterus, forming part of assisted reproductive treatment.

\medterm{Intracytoplasmic Sperm Injection} A fertility treatment in which a single sperm is injected directly into an egg in a laboratory to help achieve fertilisation. It is also called ICSI.

\medterm{In Vivo} Taking place or being studied inside a living person, animal, or other organism. For example, an in vivo test examines how a medicine works in the body rather than only in a laboratory dish.

\medterm{In-Vivo} In-vivo means occurring or being studied within a living organism.

\medterm{Inactivated Vaccines} Vaccines made from microorganisms that have been killed or from parts of them. They cannot reproduce or cause the infection they are designed to prevent, but they may need several doses or booster doses to maintain protection.

\medterm{Inactivator} A substance or process that makes a biological agent inactive. Inactivators may disable viruses, toxins, enzymes, or medicines and are used in immunity, laboratory testing, sterilization, and treatment.

\medterm{Inalimarev} An investigational antiviral medicine studied against selected viral infections. Its effectiveness, approved uses, dosing, interactions, and safety remain dependent on clinical evidence and regulatory authorization.

\medterm{Inamrinone} Inamrinone is an intravenous medicine that increases heart contraction and widens blood vessels; it was used short term in severe heart failure but can cause low blood pressure and abnormal rhythms.

\medterm{Inanition} Severe physical depletion caused by inadequate food or nutrient intake, often accompanied by weight loss, weakness, dehydration, and metabolic disturbance. Causes may include starvation, malabsorption, chronic disease, or inability to eat.

\medterm{Inborn Error Of Metabolism} An inherited disorder caused by a defect in a metabolic pathway, often resulting in
accumulation of harmful substances or deficiency of essential products.

\medterm{Inborn Errors Of Metabolism} Inherited disorders caused by defects in enzymes, transporters, or other metabolic proteins, leading to abnormal buildup or deficiency of substances.

\medterm{Inbreed} To reproduce between genetically related individuals; inbreeding can increase the chance that recessive inherited disorders will occur.

\medterm{Inbreeding} Reproduction between people who share a recent ancestor, increasing the chance that a child inherits two copies of the same recessive disease-causing variant. Genetic counseling can help assess family risk.

\medterm{Inbreeding, Coefficient Of} The probability that the two copies of a gene at a locus in an individual are identical by descent from a common ancestor; it is used in population and medical genetics.

\medterm{Incarcerated Hernia} A hernia whose protruded tissue cannot be returned to the abdominal cavity. It may progress to strangulation and requires urgent clinical assessment.

\medterm{Incentive Spirometry} A breathing device that provides visual feedback while a person inhales slowly and deeply. It may be used with coughing, mobilization, and other measures to encourage lung expansion after illness or surgery; it is a lung-expansion exercise and is not the same as diagnostic spirometry.

\medterm{Incest} Sexual activity or reproduction between close biological relatives. It can involve abuse or coercion and increases the chance of inherited disorders in children; affected people need safe, confidential medical and safeguarding support.

\medterm{Incidence} Incidence is the number of new cases of a disease or event arising in a population during a specified period.

\medterm{Incidental Finding} An unexpected observation discovered during an examination performed for another reason. Its importance depends on the
finding and the clinical context.

\medterm{Incidentaloma} An unexpected mass or abnormality found during imaging or another test performed for a different reason. It may be harmless or important, so follow-up depends on its size, appearance, location, growth, and the person's risks.

\medterm{Incised Wound} Clean-cut wound produced by a sharp-edged instrument (scalpel, knife, glass) that is
longer on the skin surface than it is deep, featuring clean edges without tissue
bridging or marginal abrasion.

\medterm{Incision} A deliberate cut made by a healthcare professional through the skin or another tissue, usually during surgery or to allow drainage. The wound is closed or left to heal according to its size, location, contamination, and purpose.

\medterm{Incisional Biopsy} Surgical removal of only part of a lesion for diagnosis, generally used when the lesion is too large or located where complete removal would be inappropriate as an initial procedure.

\medterm{Incisional Hernia} A hernia through a weakness in a previous surgical incision. It may appear as a bulge that
becomes more prominent with standing or straining. Risk factors include wound infection, obesity, smoking, poor
nutrition, chronic cough, and repeated surgery.

\medterm{Incisor} A front tooth shaped mainly for cutting and biting food. Adults usually have eight incisors, and each has a crown above the gum and a root anchored in the jawbone.

\medterm{Incisors} The flat front teeth used mainly for biting and cutting food. Adults usually have eight incisors, and their health can be affected by decay, injury, gum disease, wear, or poor alignment.

\medterm{Inclusion Body} A circumscribed mass of foreign, metabolic, or abnormal material within a cell's cytoplasm or nucleus. Inclusion bodies may contain aggregated or misfolded proteins, viral components, glycogen, lipids, or other substances, and their appearance and location can help identify cellular injury or disease.

\medterm{Inclusion Body Myositis} A slowly progressive inflammatory and degenerative muscle disease, usually causing
asymmetric weakness of the finger flexors and quadriceps. Muscle biopsy may show endomysial inflammation and rimmed
vacuoles. Response to immunotherapy is often limited.

\medterm{Inclusion Conjunctivitis} Conjunctivitis caused by \textit{Chlamydia trachomatis}, usually presenting with persistent mucopurulent discharge, follicular inflammation, and sometimes genital infection. Neonatal disease requires prompt evaluation and systemic treatment to prevent complications and transmission.

\medterm{Inclusion-Cell (I-Cell) Disease} A rare inherited lysosomal-storage disorder caused by impaired tagging and trafficking of lysosomal enzymes,
leading to their secretion outside cells. It causes severe developmental delay, skeletal abnormalities, coarse facial
features, and progressive multisystem disease. Management is supportive and coordinated by specialists.

\synonyms
I-Cell

\medterm{Incompetence} Failure of a structure to close or function adequately; examples include valvular incompetence of the heart and cervical incompetence during pregnancy.

\medterm{Incompetent Cervix} Early painless opening of the cervix during pregnancy, which may lead to pregnancy loss or premature birth.

\medterm{Incomplete Miscarriage} Partial loss of a pregnancy in which some pregnancy tissue remains inside the uterus, potentially causing bleeding, pain, infection, or continued pregnancy symptoms.

\medterm{Incomplete Ossification} Failure of bone or cartilage to mineralize or develop fully. It may be part of normal development or result from genetic, metabolic, nutritional, or skeletal disorders and is interpreted with age, imaging, and laboratory findings.

\medterm{Incontinence} Involuntary loss of control over urination, bowel movements, or both. It may result from weakened muscles, nerve disease, infection, constipation, childbirth, prostate problems, medicines, mobility limits, or other conditions.

\medterm{Incontinence (Fecal)} Inability to control bowel movements, causing accidental leakage of stool or mucus. It may result from diarrhea, constipation, weakened anal muscles, nerve disease, childbirth injury, or medicines. Assessment looks for the cause; treatment may include bowel regulation, exercises, retraining, medicines, or surgery.

\synonyms
Fecal

\medterm{Incontinence (Urinary)} Alternate index label; see Urinary incontinence. Involuntary urine leakage. Types: stress (cough/sneeze), urge (overactive bladder), overflow (retention), functional (immobility/cognition), mixed. Assessment: history, bladder diary, urinalysis, post-void residual.
\synonyms
Urinary


\medterm{Incontinence, Bowel} Alternate index label; see \textit{Fecal incontinence}.

\medterm{Incontinence, Fecal} Alternate index label; see \textit{Fecal incontinence}.

\medterm{Incontinence, Urinary} Alternate index label; see \textit{Urinary incontinence}.

\medterm{Incontinentia Pigmenti} An X-linked dominant disorder affecting skin, teeth, eyes, and the nervous system,
caused by pathogenic variants in IKBKG. Skin lesions evolve through vesicular,
verrucous, hyperpigmented, and atrophic stages and may be associated with retinal
disease or seizures.

\medterm{Incoordinatio} Lack of coordination between muscles or body movements; the term is often used in older medical language.

\medterm{Incoordination} Impaired ability to perform smooth, accurate, coordinated movements, often called ataxia. Causes include cerebellar disease, sensory loss, vestibular disorders, intoxication, medicines, and other neurologic conditions.

\medterm{Incoordination (Ataxia)} Impaired coordination of voluntary movement, often caused by disease affecting the cerebellum or other parts of the nervous system.

\synonyms
Ataxia

\medterm{Increased Head Circumference} Increased head circumference means head size above the expected range for age and sex and may reflect familial size, hydrocephalus, swelling, or a genetic condition.

\medterm{Increased Intracranial Pressure} Abnormally elevated pressure within the skull caused by brain swelling, hemorrhage, mass
lesion, impaired cerebrospinal-fluid circulation, or venous obstruction. It can cause
headache, vomiting, visual disturbance, reduced consciousness, and brain herniation.

\medterm{Increased Libido} Increased libido is heightened sexual desire that may be normal or associated with hormones, medicines, mood elevation, neurologic disease, or substance use.

\medterm{Increased Osteoid} Increased osteoid is excess unmineralized bone matrix, seen when bone formation outpaces mineralization as in osteomalacia or some metabolic disorders.

\medterm{Increased Vascular Permeability} A fundamental feature of acute inflammation that allows plasma proteins, fluid, and
leukocytes to exit the vascular space and enter the extravascular tissue.

\medterm{Increment} An increase or stepwise change in an amount, measurement, dose, or treatment intensity. The size and timing of an increment should be chosen carefully to achieve benefit while limiting harm.

\medterm{Incretins} Incretins are gut hormones, mainly GLP-1 and GIP, released after meals that enhance glucose-dependent insulin secretion; GLP-1 also suppresses glucagon, slows gastric emptying, and reduces appetite.

\medterm{Incubation} Incubation is the time between exposure to an infectious agent and the appearance of symptoms.

\medterm{Incubation Period} The time between exposure to an infectious agent and the beginning of symptoms. It varies with the organism, dose, route of exposure, and person’s immune response.
\medterm{Incubator} A controlled-environment chamber culturing microorganisms, cells, or tissues at defined temperature, humidity, and atmospheric composition.
\medterm{Incus} The incus, commonly known as the anvil, is the middle ossicle of the three bones in the
middle ear, articulating with the malleus at the incudomalleolar joint and with the
stapes at the incudostapedial joint.

\medterm{Incyclinide} An investigational tetracycline-related medicine studied for effects on enzymes involved in tissue breakdown and inflammation. Its potential cancer or anti-inflammatory uses remain dependent on clinical research and regulatory approval.

\medterm{Indapamide} Indapamide is a thiazide-like diuretic used mainly for hypertension; it promotes renal sodium excretion and may cause hyponatraemia, hypokalaemia, dehydration, hyperuricaemia, or photosensitivity.

\medterm{Indel Mutation} A genetic change in which one or more DNA building blocks are inserted or deleted. If the number is not a multiple of three, it can disrupt the protein reading frame.
\medterm{Independence} The ability to perform activities and make decisions without assistance, although a person may still use devices, adaptations, or occasional support.

\medterm{Indeterminate} Indeterminate means not definitively classified or interpreted because available findings are insufficient, conflicting, or nonspecific.
\medterm{India Immunization Schedule} The Universal Immunization Schedule (UIP) of the Government of India, providing free vaccines to all children, adolescents, and pregnant women through the public health system.
\medterm{Indican} Indican is a compound formed from tryptophan metabolism and excreted in urine; abnormal levels may occur with intestinal bacterial overgrowth or metabolic disease.

\medterm{Indicate} To show, suggest, or point to a finding or conclusion; an indication is a clinical reason for using a test, treatment, or procedure.

\medterm{Indication} A clinical condition, symptom, or question that provides the reason for performing a test, procedure, or treatment.

\medterm{Indigestion} Common term for dyspepsia, discomfort or burning in the upper abdomen that may include early fullness, bloating, belching, or nausea. Persistent symptoms, bleeding, weight loss, vomiting, anemia, or difficulty swallowing require medical evaluation.

\medterm{Indigo Carmine} Indigo carmine is a blue dye used in selected medical procedures to visualize urinary or other fluid flow and in laboratory staining.

\medterm{Indinavir} An older HIV protease-inhibitor medicine that blocks production of mature virus particles. It is rarely used now because of resistance, drug interactions, kidney stones, and other adverse effects.

\medterm{Indirect Antiglobulin Test} A test that detects free, clinically significant antibodies in a patient’s serum or
plasma by incubating it with reagent erythrocytes. It is used for antibody screening,
compatibility testing, and prenatal red-cell antibody evaluation.

\medterm{Indirect Calorimetry} Measurement of oxygen consumption and carbon-dioxide production to estimate resting
energy expenditure. It is the preferred method for determining energy needs when
prediction equations are unreliable, particularly in critically ill or metabolically
complex patients.

\medterm{Indirect Veneer} A thin covering made outside the mouth and bonded to the front of a tooth to change its shape, color, or appearance. It may be made from porcelain or composite material. Tooth preparation, strength, hygiene, and long-term maintenance affect its suitability.


\medterm{Indium} Indium is a metal used in industry and medical radioisotopes; exposure effects depend on the compound, dose, and route.

\medterm{Indium-111} A radioactive form of indium used as a tracer in selected nuclear-medicine tests. It may label white blood cells to help locate inflammation or infection, or be attached to another substance to help show certain tumors. It emits detectable radiation that is recorded by a special camera.

\medterm{Indium-111 Scan} A nuclear medicine scan that uses the radioactive tracer indium-111 to make selected tissues visible. It may be attached to white blood cells to look for inflammation or infection, or to a substance that helps find some tumors.

\medterm{Indium-Labelled WBC Scan} A nuclear medicine study in which a patient’s white blood cells are labeled with indium-111 and reinjected to localize selected sites of inflammation or infection.

\medterm{Indocyanine Green} Indocyanine green is a fluorescent dye injected intravenously to assess liver function, blood flow, and retinal or choroidal circulation.

\medterm{Indole} A chemical structure containing a nitrogen atom, found in tryptophan and many natural substances and medicines. Indole-containing compounds can affect the nervous system, blood vessels, digestion, inflammation, or other body processes.

\medterm{Indole-3-Carbinol} A compound formed from glucosinolates in cruciferous vegetables and studied for effects on estrogen metabolism and cancer biology.

\medterm{Indolent} Slow to develop, progress, or cause symptoms; it often describes a relatively slow-growing disease or tumor.

\medterm{Indoles} A group of compounds containing an indole ring, including substances involved in biological signaling, metabolism, and medicines.

\medterm{Indometacin} Indometacin is a nonsteroidal anti-inflammatory medicine used for pain and inflammation; it can cause stomach bleeding, kidney problems, and cardiovascular complications.

\medterm{Indomethacin} Indomethacin is a potent NSAID used for inflammatory pain and selected acute conditions such as gout; gastrointestinal ulceration or bleeding, renal injury, cardiovascular events, and headache are important adverse effects.

\medterm{Indomethacin Overdose} Toxic exposure to excessive indomethacin, which can cause nausea, vomiting, stomach bleeding, kidney injury, seizures, or other serious effects.

\medterm{Indoor Allergens} Substances inside buildings that can trigger allergic or asthma symptoms, including dust mites, mold, cockroach particles,
and animal dander. Exposure may cause rhinitis, conjunctivitis, wheezing, or eczema. Reduction of exposure and
appropriate allergy or asthma treatment depend on the identified trigger.


\medterm{Indoximod} Indoximod is an investigational immunotherapy that modulates tryptophan metabolism and has been studied in combination cancer treatment.

\medterm{Induced Abortion} Deliberate termination of a pregnancy using medicines or a procedure. Safe care depends on gestational age, patient preference, accurate information, legal context, trained providers, and assessment for complications such as heavy bleeding or infection.

\medterm{Induced Hyperthermia} Induced hyperthermia is deliberate elevation of body or tissue temperature for a therapeutic purpose, with careful monitoring to prevent heat injury.

\medterm{Induced Labor} Medical initiation of uterine contractions before labor begins naturally to achieve delivery when continuing pregnancy may pose greater risk.

\medterm{Induced Sputum Analysis} A method of assessing airway inflammation by examining sputum produced after inhaled
saline stimulation. The types and proportions of inflammatory cells may support assessment of asthma, infection, COPD,
or bronchiectasis.

\medterm{Inducing Labor} Starting childbirth artificially with medicines, breaking the amniotic sac, or other methods when continuing pregnancy may pose greater risks than delivery.

\medterm{Induction} Initiation of a physiologic or therapeutic process, such as induction of labor, anesthesia, remission, or immune tolerance. The method, indication, benefits, and risks depend on the clinical context.

\medterm{Induction Of Labor} Starting uterine contractions medically before labor begins on its own. It may be recommended when continuing pregnancy is riskier, using medicines or procedures that soften the cervix and encourage contractions.

\medterm{Induction Therapy} The first, often intensive, part of treatment given to reduce or remove a disease. In acute leukemia, its main aim is remission, meaning that tests find no clear evidence of leukemia; further treatment is usually needed to reduce the chance of return.

\medterm{Indurated} Abnormally firm or hardened when palpated, usually because of inflammation, edema, fibrosis, infiltration, or a tumor within the skin or underlying tissue.
\medterm{Induration} Abnormal hardening of tissue, often detected by palpation and caused by inflammation, infection, edema, fibrosis, or a mass.

\medterm{Industrial Bronchitis} Industrial bronchitis is chronic airway irritation and inflammation related to workplace dusts, fumes, or chemicals, causing cough, phlegm, and breathing limitation.

\medterm{Industrial Disease} A disease caused or worsened by workplace exposures or conditions, such as chemicals, dust, noise, radiation, infection, or repetitive strain.

\medterm{Industrial Health} The branch of occupational health concerned with preventing and managing illness and injury related to industrial work, workplace exposures, and working conditions.

\medterm{Indwelling Catheter} A catheter left in place for ongoing drainage, infusion, monitoring, or access; urinary indwelling catheters require appropriate indications and care because they increase infection and trauma risk.

\medterm{Indwelling Catheter Care} A patient-care topic concerning indwelling catheter care. It addresses the relevant purpose, risks, management, and follow-up.

\medterm{Indwelling Pleural Catheter} A flexible tube placed through the chest wall into the space around the lung to drain fluid that keeps returning. It can allow drainage at home and relieve breathlessness, especially when cancer or a lung that cannot fully expand causes the fluid. Possible problems include infection, blockage, pain, or air leakage.

\synonyms
IPC; Tunneled Pleural Catheter; PleurX Catheter

\medterm{Inertia} Resistance to a change in motion; in obstetrics, uterine inertia means weak or ineffective contractions that slow labor.

\medterm{Infant} An infant is a child in the first year of life, a period of rapid growth and changing feeding, immune, neurologic, and developmental needs.

\medterm{Infant - Newborn Development} The early growth of a newborn’s body, brain, senses, movement, feeding, communication, and social responses during the first weeks and months.

\medterm{Infant Botulism} A neuromuscular illness caused by intestinal colonization with Clostridium botulinum and
production of botulinum neurotoxin in infants. It causes constipation, poor feeding,
weak cry, hypotonia, facial weakness, and descending paralysis.

\medterm{Infant Feeding} Provision of breast milk, appropriately prepared formula, or both to meet an infant's
energy, protein, fluid, and micronutrient requirements. Feeding recommendations depend
on age, growth, medical conditions, developmental readiness, and safe preparation
practices.

\medterm{Infant Formula} Infant formula is manufactured nutrition designed to provide adequate nutrients when breast milk is unavailable, insufficient, or not chosen; preparation and storage must be safe.


\medterm{Infant Formulas} Commercial preparations designed to provide nutrition for infants when breast milk is unavailable, insufficient, or not chosen.

\medterm{Infant Gynecomastia} Infant gynecomastia is temporary breast enlargement in a newborn or young infant caused by maternal or neonatal hormones and usually resolves spontaneously.

\medterm{Infant Harlequin Syndrome} A transient, benign cutaneous vascular phenomenon in neonates characterized by
asymmetric flushing of one half of the body with sharp midline demarcation, caused by
transient asymmetric sympathetic tone. The affected side is red and warm while the
contralateral side is pale.

\medterm{Infant Incubators} Enclosed devices providing controlled warmth, humidity, oxygen, and monitoring for premature, small, or ill newborns who need extra support after birth.

\medterm{Infant Jaundice} Yellowing of an infant’s skin or eyes caused by excess bilirubin. It is often temporary, but jaundice in the first day, rapid worsening, poor feeding, unusual sleepiness, or pale stools requires urgent assessment.
\medterm{Infant Leukemia} Infant leukemia is leukemia diagnosed during infancy, in which abnormal blood-forming cells multiply and interfere with normal blood-cell production.

\medterm{Infant Mortality} The death of an infant before the first birthday; the infant mortality rate is usually expressed per 1,000 live births.

\medterm{Infant Nutrition} Nutritional support for infants through breast milk, appropriately prepared formula, and developmentally suitable complementary foods. Requirements change with age, growth, illness, prematurity, and feeding ability; unsafe preparation can cause infection or malnutrition.

\medterm{Infant Of A Substance-Using Mother} A newborn exposed to alcohol or other substances during pregnancy who may need assessment for withdrawal, toxicity, feeding problems, or developmental effects.

\medterm{Infant Of Diabetic Mother} A newborn born to a mother with diabetes who may be at increased risk of large size, low blood sugar, breathing problems, jaundice, or congenital abnormalities.

\medterm{Infant Reflex} An infant reflex is an automatic response present during early development, such as rooting, sucking, grasping, or startle, whose timing and symmetry help assess neurologic function.

\medterm{Infant Reflexes} Automatic newborn responses such as rooting, sucking, Moro, grasp, and stepping that help assess neurologic development.

\medterm{Infant Reflux} Backflow of stomach contents into an infant’s esophagus, often causing uncomplicated spitting up. It usually improves with age, but poor growth, breathing problems, blood, forceful vomiting, or distress requires medical assessment.

\medterm{Infant Sleep Position} Placing an infant on the back for every sleep is recommended to reduce the risk of sleep-
related death. Infants should sleep on a firm, separate surface without pillows, loose bedding, bumpers, or stuffed
toys.

\medterm{Infant Test Procedure Preparation} Infant test procedure preparation includes age-appropriate instructions, feeding or fasting guidance, comfort measures, and safety steps before a test.

\medterm{Infant, Post-Term} An infant born at or after 42 weeks of gestation, beyond the usual term-pregnancy period; risks depend on placental function, fetal condition, and circumstances of delivery.

\medterm{Infanticide} The killing of a child, in law typically defined as the killing of a newborn or infant
(usually within the first year, or specifically within 24 hours of birth in some
jurisdictions) by its mother while mentally disturbed by the effects of childbirth or
lactation.

\medterm{Infantile Colic} A behavioral syndrome in healthy infants characterized by episodes of excessive,
inconsolable crying for >3 hours per day, >3 days per week, for >3 weeks (Wessel
criteria). Onset typically at 2-3 weeks, peaks at 6-8 weeks, and resolves by 3-4 months
of age.

\medterm{Infantile Fibrosarcoma} Infantile fibrosarcoma is a rare soft-tissue cancer that usually develops in infants or young children and requires specialist treatment.

\medterm{Infantile Hemangioma} Alternate index label; see \textit{Hemangioma}.

\medterm{Infantile Hepatic Hemangiomas} Benign growths made of extra blood vessels in an infant’s liver. Small lesions may cause no symptoms, while numerous or large lesions can enlarge the liver, strain the heart, lower thyroid hormone, or cause bleeding problems. Ultrasound and specialist follow-up guide monitoring and treatment.

\medterm{Infantile Paralysis (Polio)} Paralytic poliomyelitis caused by poliovirus infection involving motor neurons, which can produce acute flaccid paralysis and respiratory failure; vaccination prevents most cases.

\medterm{Infantile Spasm} Infantile spasms are brief epileptic seizures in infancy, often in clusters with flexion or extension movements, and require urgent diagnosis and treatment to protect development.

\medterm{Infantile Spasms} A severe seizure disorder of infancy characterized by brief, sudden flexion or extension spasms, often occurring in clusters and associated with developmental impairment.


\medterm{Infarct} An area of dead tissue caused by inadequate blood supply.

\medterm{Infarction} Tissue death caused by inadequate blood flow and oxygen, usually because a vessel is blocked. Examples include heart attack, brain infarction, and intestinal infarction.
infarction.

\medterm{Infection} Entry and multiplication of bacteria, viruses, fungi, parasites, or other microorganisms in the body. Infection may cause tissue damage and illness, although some infections produce no symptoms; the cause and site determine treatment.

\medterm{Infection (Geriatric)} Infectious diseases in elderly often present atypically: fever may be absent, confusion
may be sole symptom. Common infections: UTI, pneumonia, skin/soft tissue, bloodstream.
Prevention: vaccination, hand hygiene, aspiration precautions. Treatment: earlier
presentation, broader differential, drug interactions, renal dosing.

\synonyms
Geriatric


\medterm{Infection Control} Practices that prevent microorganisms from spreading between people, objects, and healthcare settings. They include hand cleaning, appropriate protective equipment, safe injection and device care, cleaning, sterilization, vaccination, and isolation when needed.

\medterm{Infection Cryptococcus (Cryptococcosis)} Alternate index label for Cryptococcosis.

\medterm{Infection Cyclospora (Cyclospora Infection (Cyclosporiasis))} Alternate index label for Cyclospora Infection (Cyclosporiasis).

\medterm{Infection Enterovirus (Enterovirus (Non-Polio Enterovirus Infection))} Alternate index label for Enterovirus (Non-Polio Enterovirus Infection).

\medterm{Infection HPV (HPV Infection Human Papillomavirus)} Alternate index label for HPV Infection Human Papillomavirus.

\medterm{Infection Human Papillomavirus (HPV Infection Human Papillomavirus)} Alternate index label for HPV Infection Human Papillomavirus.

\medterm{Infection Imaging} Imaging used to locate an infection when examination and routine tests do not show its source. Depending on the suspected problem, it may use ultrasound, CT, MRI, or a nuclear-medicine tracer; scan findings must be interpreted with symptoms, cultures, and other tests.

\medterm{Infection Listeria (Listeria)} Alternate index label for Listeria.

\medterm{Infection Non-Polio Enterovirus (Enterovirus (Non-Polio Enterovirus Infection))} Alternate index label for Enterovirus (Non-Polio Enterovirus Infection).

\medterm{Infection Pinworms (Pinworm Infection)} Alternate index label for Pinworm Infection.

\medterm{Infection Prevention} Measures that reduce the risk of infection, including hand hygiene, personal protective equipment, environmental cleaning, sterilization, disinfection, vaccination, and appropriate isolation precautions.


\medterm{Infection Testicle (Testicular Disorders)} Alternate index label for Testicular Disorders.

\medterm{Infection, Hospital-Acquired} An infection acquired during or as a result of healthcare that was not present or incubating at the time of admission; it may be caused by bacteria, viruses, fungi, or other organisms.

\medterm{Infection, Urinary Tract} Infection of part of the urinary system, commonly the bladder or urethra and sometimes the kidneys, usually caused by bacteria and producing symptoms such as burning urination, urgency, or flank pain.

\medterm{Infections In Pregnancy} Infections during pregnancy can affect the pregnant person, placenta, fetus, or newborn. Some, including listeriosis, group B streptococcal infection, urinary infection, and sexually transmitted infections, can cause miscarriage, premature birth, fetal illness, or newborn sepsis. Testing and treatment depend on the organism and stage of pregnancy.

\medterm{Infectious Arteritis} Inflammation and damage of an artery caused by infection, which can weaken the vessel, reduce blood flow, or produce aneurysm.

\medterm{Infectious Arthritis} Inflammation of a joint caused by infection, usually bacterial and sometimes viral or fungal.
Acute pain, swelling, warmth, fever, and restricted movement require urgent evaluation because
joint damage can develop quickly.

\medterm{Infectious Conjunctivitis} Inflammation of the conjunctiva caused by a virus, bacterium, or less commonly another infectious organism, often causing redness and discharge.

\medterm{Infectious Disease} Illness caused by a pathogen such as a bacterium, virus, fungus, or parasite. Infection may spread between people, from animals or the environment, or from organisms normally living in the body.

\medterm{Infectious Dose} The number of viable organisms required to establish infection in a host under specified
conditions. The dose depends on the pathogen, route of exposure, host defenses, and
organism viability.

\medterm{Infectious Encephalitis} Brain inflammation caused by an infectious organism, most often a virus. It may cause fever, headache, confusion, seizures, behavior changes, weakness, or reduced consciousness and requires urgent treatment.

\medterm{Infectious Encephalomyelitis} Infection-related inflammation of the brain and spinal cord that may cause fever, headache, confusion, weakness, seizures, or other neurologic signs.

\medterm{Infectious Enterocolitis} Infection-related inflammation of the small intestine and colon, commonly causing diarrhea, abdominal pain, fever, or vomiting.

\medterm{Infectious Esophagitis} Inflammation of the esophagus caused by an infectious organism, most often Candida,
herpes simplex virus, or cytomegalovirus in immunocompromised patients. It commonly
causes painful swallowing, difficulty swallowing, and retrosternal pain.

\medterm{Infectious Hepatitis} Alternate index label; see \textit{Viral hepatitis}.

\medterm{Infectious Mediastinitis} Infection and inflammation of the tissues in the middle of the chest, often after esophageal injury, surgery, or spread from a nearby infection.

\medterm{Infectious Meningitis} Infection and inflammation of the membranes surrounding the brain and spinal cord. Bacterial
meningitis is a medical emergency; headache, fever, neck stiffness, vomiting, confusion, or a
new rash requires immediate medical attention.

\medterm{Infectious Mononucleosis} A viral illness, usually caused by Epstein-Barr virus, characterized by fatigue, sore throat, and swollen lymph nodes.

\synonyms
Mono (Infectious Mononucleosis)

\medterm{Infectious Myringitis} Inflammation and infection of the tympanic membrane, often associated with viral or
bacterial respiratory infection. It can cause severe ear pain, fever, hearing reduction,
and small fluid-filled blisters on the eardrum.

\medterm{Infective Cervicitis} Infective cervicitis is inflammation of the cervix caused by organisms such as chlamydia, gonorrhea, trichomonas, or herpes, producing discharge, bleeding, or pelvic discomfort.

\medterm{Infective Cystitis} Infective cystitis is bacterial or other microbial inflammation of the bladder, causing urinary frequency, urgency, burning, and sometimes blood or fever.

\medterm{Infective Endocarditis} Infection of the endocardial surface of the heart, most commonly affecting the heart
valves. It is diagnosed using the Duke criteria, which include major criteria (positive
blood cultures, echocardiographic evidence of vegetations) and minor criteria
(predisposing heart condition, fever, vascular phenomena, immunologic phenomena).

\medterm{Infective Phlebitis} Infective phlebitis is infection and inflammation of a vein, often around an intravenous catheter, causing pain, redness, swelling, and possible bloodstream infection.

\medterm{Infective Vaginitis} Infective vaginitis is infection or inflammation of the vagina caused by yeast, bacteria, parasites, or viruses and may cause discharge, itching, odor, or pain.

\medterm{Inferior} Below or nearer the feet compared with another body part. For example, the diaphragm is inferior to the heart, and the inferior vena cava carries blood from the lower part of the body to the heart.

\medterm{Inferior Alveolar Nerve Block} A regional anesthetic injection near the mandibular foramen to anesthetize the inferior
alveolar nerve before it enters the mandibular canal. It usually anesthetizes mandibular
teeth on that side; the lingual nerve is often anesthetized incidentally.

\medterm{Inferior Colliculus} The inferior colliculus is a midbrain auditory relay that integrates sound information and helps orient reflexively to auditory stimuli.

\medterm{Inferior Hemifield Loss} Loss of vision in the lower half of the visual field in one or both eyes, which
localises to lesions of the superior retina or superior optic radiations passing through
the parietal lobe. Causes include superior retinal detachment, branch retinal artery
occlusion, and parietal lobe stroke.

\medterm{Inferior Mesenteric Artery} Third major abdominal aortic branch at L3, the hindgut artery. Supplies distal
transverse colon through rectum. Branches: left colic, sigmoid, superior rectal. The
marginal artery of Drummond provides SMA-IMA collateral circulation. Occlusion may cause
ischemic colitis at the splenic flexure watershed (Griffiths' point).

\medterm{Inferior Myocardial Infarction} A heart attack affecting the lower surface of the heart muscle. It usually results from blocked blood flow in a coronary artery and requires urgent treatment to limit permanent damage.

\medterm{Inferior Vena Cava} Body's largest vein, formed by common iliac vein union at L5, returning lower body
deoxygenated blood to the right atrium. Ascends right of the aorta, pierces the
diaphragm at T8, enters the pericardium. Receives hepatic, renal, gonadal, and lumbar
veins. IVC thrombosis/compression causes bilateral lower extremity edema and DVT.

\medterm{Inferolateral} Located below and toward the side of a reference structure.

\medterm{Infertility} Infertility is failure to achieve pregnancy after a defined period of regular unprotected intercourse, reflecting factors involving ovulation, sperm, reproductive anatomy, age, or unexplained causes.
to conceive or carry a pregnancy. Causes may involve ovulation, reproductive anatomy, sperm, age-related fertility,
medical conditions, or unexplained factors.

\synonyms
Subfertility


\medterm{Infertility, Female} Alternate index label; see \textit{Female infertility}.

\medterm{Infertility, Male} Alternate index label; see \textit{Male infertility}.

\medterm{Infestation} The presence or invasion of parasitic organisms, such as lice, mites, or worms, on or in the body.

\medterm{Infiltrate} An area of increased tissue density seen on an image, often in the lung. It is nonspecific and may represent infection, fluid, bleeding, inflammation, collapse, tumor, or another process.

\medterm{Infiltrating Cancer} Cancer that has grown beyond its original location into nearby tissue. Such growth may damage normal structures and can allow cancer cells to enter lymph vessels or blood vessels and spread elsewhere.

\synonyms
Invasive cancer

\medterm{Infiltrating Ductal Carcinoma} Alternate index label; see \textit{Invasive ductal carcinoma}.

\medterm{Infiltration} Entry or accumulation of cells, fluid, or another substance within tissue. In clinical care, intravenous infiltration means leakage of nonvesicant fluid into surrounding tissue; extravasation of an irritant or vesicant can cause serious injury.

\medterm{Inflamed Gallbladder} Alternate index label; see \textit{Cholecystitis}.

\medterm{Inflamed Pancreas} Alternate index label; see \textit{Pancreatitis}.

\medterm{Inflamed Pericardium} Alternate index label; see \textit{Pericarditis}.

\medterm{Inflammaging} Persistent, low-grade, age-associated inflammatory activity linked to cellular
senescence, immune dysregulation, chronic infection, tissue injury, and comorbidity. It
contributes to frailty, atherosclerosis, sarcopenia, and neurodegenerative disease.

\medterm{Inflammasomes} Groups of proteins inside immune cells that detect danger signals and activate inflammation. They help release inflammatory messenger proteins and can cause an infected or damaged cell to die. Excessive activation may contribute to autoimmune or inflammatory disease.

\medterm{Inflammation} The body’s protective response to injury, infection, or irritation. It may cause redness, warmth, swelling, pain, and reduced function. Short-term inflammation helps healing, while persistent inflammation can damage tissues and contribute to chronic disease.

\medterm{Inflammation (Dietary)} Chronic low-grade inflammation linked to diet. Western diet pro-inflammatory (saturated
fat, sugar, refined grains). Mediterranean diet anti-inflammatory.

\synonyms
Dietary

\medterm{Inflammation Imaging} Imaging used to find areas of active inflammation in the body. A scan using a sugar-like tracer may show infection or inflammation, but cancer can look similar, so scan findings must be interpreted with symptoms, examination, and other tests.

\medterm{Inflammation Of The Heart Valve (Endocarditis)} Alternate index label for Endocarditis.

\medterm{Inflammation Sclera (Scleritis)} Alternate index label for Scleritis.

\medterm{Inflammation Vaginal (Vaginitis Overview)} Alternate index label for Vaginitis Overview.

\medterm{Inflammatory Back Pain} A pattern of chronic back pain characterized by insidious onset, improvement with
exercise, lack of improvement with rest, nocturnal pain, and prolonged morning
stiffness. It suggests axial spondyloarthritis but is not diagnostic without supportive
clinical, laboratory, and imaging findings.

\medterm{Inflammatory Bowel Disease} A group of chronic conditions that cause inflammation in the digestive tract, mainly Crohn disease and ulcerative colitis. Symptoms may come and go and can include abdominal pain, diarrhea, rectal bleeding, and weight loss. The exact cause is unknown, but genes, the immune system, the environment, and gut bacteria may all play a role.

\medterm{Inflammatory Bowel Disease Diet} Crohn disease, ulcerative colitis. Nutritional support during flares. Exclusive enteral
nutrition for Crohn induction. Individual triggers vary.

\medterm{Inflammatory Bowel Disease: Intestinal Problems} Inflammatory bowel disease, including Crohn disease and ulcerative colitis, causes chronic immune-mediated intestinal inflammation with diarrhea, pain, bleeding, weight loss, and extraintestinal effects. Treatment aims for sustained remission and prevention of structural and nutritional complications.

\synonyms
IBD (Inflammatory Bowel Disease: Intestinal Problems)

\medterm{Inflammatory Breast Cancer} A rare, aggressive breast cancer causing redness, swelling, warmth, and thickened skin, sometimes without a distinct lump. It requires prompt imaging, biopsy, staging, and coordinated treatment.

\synonyms
Breast Cancer, Inflammatory
Cancer, Inflammatory Breast

\medterm{Inflammatory Myopathies} A group of disorders in which inflammation damages muscle, usually causing gradually worsening weakness around the hips, thighs, shoulders, or upper arms. Some forms also affect skin, lungs, swallowing, or other organs. Diagnosis and treatment require specialist assessment.

\medterm{Inflammatory Response} The coordinated vascular and immune reaction to infection, tissue injury, or other danger signals, producing mediator release, leukocyte recruitment, and changes in blood flow and permeability.

\medterm{Infliximab} A biologic medicine that blocks tumor necrosis factor, an inflammatory signal. It treats several inflammatory diseases but can increase the risk of serious infections, tuberculosis reactivation, infusion reactions, and heart-failure worsening.

\medterm{Influenza} A contagious respiratory illness caused by influenza A or B viruses, characterized by
sudden onset of fever, myalgia, headache, malaise, nonproductive cough, sore throat, and
nasal congestion. Seasonal epidemics occur annually. Complications: pneumonia (primary
viral or secondary bacterial), exacerbation of underlying conditions.

\medterm{Influenza (Flu)} A contagious respiratory infection caused by influenza viruses, commonly producing fever, cough, body aches, fatigue, and respiratory symptoms.



\medterm{Influenza A Virus} An influenza virus with segmented RNA that infects people and animals and can cause seasonal epidemics or, after major genetic changes, pandemics.

\medterm{Influenza Test} A test for influenza A or B using a respiratory specimen. Rapid antigen tests provide quick results
but may miss infection. Molecular tests are generally more sensitive and may identify or
subtype influenza viruses; results must be interpreted with symptoms, timing, and local
laboratory methods.

\medterm{Influenza Vaccination In Pregnancy} Annual inactivated influenza vaccine recommended for all pregnant women. Pregnancy
increases the risk of severe influenza (pneumonia, ICU admission, death) due to changes
in immune function and cardiopulmonary physiology.

\medterm{Influenza Virus} A segmented, negative-sense RNA virus belonging to the Orthomyxoviridae family.
Influenza causes seasonal epidemics and occasional pandemics of respiratory illness. The
virus has two surface glycoproteins: hemagglutinin (HA), which mediates attachment, and
neuraminidase (NA), which facilitates release.

\medterm{Influenza, Avian} Alternate index label; see \textit{Bird flu (avian influenza)}.

\medterm{Influenza, H1N1} Alternate index label; see \textit{H1N1 flu (swine flu)}.

\medterm{Influenza, Swine Flu} Alternate index label; see \textit{H1N1 flu (swine flu)}.

\medterm{Influenzavirus A} The virus genus containing influenza A viruses, which are classified into subtypes by their hemagglutinin and neuraminidase proteins.

\medterm{Influenzavirus B} The virus genus containing influenza B viruses, which mainly infect humans and cause seasonal respiratory illness.

\medterm{Informatics} Informatics applies information science, data, computing, and human factors to organize and improve healthcare, biology, research, or another field.

\medterm{Informative} Providing useful information or evidence that improves understanding, decision-making, diagnosis, treatment planning, research, or evaluation of a health problem.

\medterm{Informed Consent} A patient's voluntary agreement to a proposed intervention after receiving understandable information about its purpose, benefits, risks, alternatives, and the option of declining, with decision-making capacity or appropriate authorization.

\medterm{Informed Consent (Geriatric)} Process of obtaining voluntary authorization from elderly patient with decision-making
capacity. Capacity assessment: understanding, appreciation, reasoning, communication of
choice. Impaired capacity may require surrogate decision-maker.

\synonyms
Geriatric

\medterm{Informed Consent - Adults} Informed consent requires an adult patient to receive understandable information about a proposed intervention, benefits, risks, alternatives, and refusal before deciding voluntarily.

\medterm{Infrared Rays} Infrared rays are electromagnetic radiation with wavelengths longer than visible red light; excessive exposure can heat tissue and cause burns.

\medterm{Infrared Spectrometry} A method for identifying or measuring substances by examining how they absorb infrared light. Different chemical bonds absorb different wavelengths, creating a pattern that can help analyze medicines, tissues, samples, or pollutants.

\medterm{Infrarenal Aorta} The infrarenal aorta is the portion of the abdominal aorta below the renal arteries and above its division into the iliac arteries.

\medterm{Infratentorial Neoplasm} An infratentorial neoplasm is an abnormal growth in the posterior fossa, the part of the skull containing the cerebellum and brainstem.

\medterm{Infrequent} Infrequent means occurring less often than usual or expected, such as infrequent bowel movements, symptoms, or events.

\medterm{Infundibulopelvic Ligament} A fold of tissue running from each ovary to the side wall of the pelvis. It carries the ovary’s blood vessels, lymph vessels, and nerves; surgeons must protect the nearby ureter during pelvic operations.

\synonyms
Suspensory Ligament of the Ovary; IP Ligament

\medterm{Infuse} To introduce a fluid, medicine, or other substance gradually into the body, commonly through an intravenous line; it can also mean to soak a substance in a liquid.

\medterm{Infusion} Slow delivery of a fluid or medicine into a vein, under the skin, or another body site. It allows controlled dosing over time and may require monitoring for allergic reactions, infection, or fluid overload.

\synonyms
Intravenous infusion

\medterm{Infusion Pump} A device that delivers a measured amount of medicine or fluid at a controlled rate, usually into a vein, under the skin, or around the spinal cord. It can improve dosing accuracy but requires correct programming and monitoring for blockage, leakage, infection, or excessive delivery.

\medterm{Ingestion} The act of taking a substance into the digestive tract by swallowing. Ingestion of a toxic substance can cause poisoning.

\medterm{Ingredient} An ingredient is a substance included in a medicine, food, cosmetic, or other preparation, whether active or inactive.

\medterm{Ingrown Hair} A hair that curls back into or grows sideways beneath the skin, causing an itchy or painful bump. It may become inflamed or infected, especially after shaving, waxing, or close trimming.
\medterm{Ingrown Nail} An ingrown nail occurs when the nail edge grows into surrounding skin, causing pain, redness, swelling, and sometimes infection or granulation tissue.

\medterm{Ingrown Toenail} Ingrown toenail (onychocryptosis) occurs when the nail edge grows into the surrounding
skin, causing pain, redness, and swelling. It is common in adolescents and adults,
exacerbated by improper nail trimming, tight shoes, and trauma. Treatment includes
partial nail avulsion with matrixectomy (phenol application) for recurrent cases.

\synonyms
Onychocryptosis (Ingrown Toenail)


\medterm{Inguinal} Relating to the groin, the region where the lower abdomen meets the upper thigh and where certain hernias, lymph nodes, and injuries occur.

\medterm{Inguinal Canal} An oblique passage (4 cm) through the anterior abdominal wall, from the deep (internal)
inguinal ring to the superficial (external) inguinal ring, transmitting the spermatic
cord in males and round ligament in females, plus ilioinguinal and genital branch of
genitofemoral nerves.

\medterm{Inguinal Hernia} Protrusion of abdominal contents through the inguinal canal. Indirect (lateral to
inferior epigastric vessels, through deep inguinal ring): congenital, patent processus
vaginalis; most common overall (60-70\%), male predominance.

\medterm{Inguinal Hernia Repair} Surgery that returns tissue to the abdomen and closes or strengthens a weak area in the groin. It may be performed through a larger incision or several small ones, with or without mesh; recurrence, infection, and persistent pain are possible.


\medterm{Inguinal Orchiectomy} Inguinal orchiectomy is removal of a testicle through an incision in the groin, commonly used when treating suspected testicular cancer.

\medterm{Inguinal Region} Area at the thigh-trunk junction traversed by the inguinal ligament (ASIS to pubic
tubercle). Clinically significant for inguinal hernias: indirect (through deep inguinal
ring, most common) and direct (through Hesselbach's triangle).

\medterm{Inhalant} A substance breathed into the lungs as a gas, vapor, aerosol, or smoke; inhalants may be medicines, workplace substances, or abused chemicals.

\medterm{Inhalation} Drawing air, medicine, particles, or vapors into the lungs. Inhalation is a route for respiratory medicines and anesthetics, while accidental inhalation of food, vomit, smoke, or toxic substances can cause airway obstruction or lung injury.

\medterm{Inhalation Injury} Thermal or chemical injury to the airway or lungs caused by smoke, heated gases, or
toxic combustion products. It can produce upper-airway edema, bronchospasm, impaired
mucociliary clearance, and carbon-monoxide or cyanide toxicity and may worsen after the
initial burn.

\medterm{Inhalational Anesthesia} Anesthetic maintenance or induction using volatile liquids or gases delivered through
the respiratory tract. Depth is influenced by inspired concentration, alveolar
ventilation, uptake, distribution, and the agent's blood-gas solubility.

\medterm{Inhaled Corticosteroids} Anti-inflammatory medications delivered directly to the airways, reducing airway
inflammation, mucus production, and bronchial hyperresponsiveness. Examples include
fluticasone, budesonide, beclomethasone, and mometasone. They are the cornerstone of
asthma maintenance therapy. Side effects include oral candidiasis and dysphonia.

\medterm{Inhaled Nitric Oxide} A selective pulmonary vasodilator used in acute settings (acute respiratory failure,
right heart failure, post-cardiac surgery). It reduces pulmonary vascular resistance and
improves oxygenation by dilating well-ventilated pulmonary vessels (improving V/Q
matching). It has a short half-life and requires continuous delivery.

\medterm{Inhaler} A device that delivers medicine into the airways through the mouth. Pressurised inhalers release a measured spray, while dry-powder inhalers release powder during a strong breath in; correct technique is important for the medicine to work.

\medterm{Inherited Disease} A disease caused wholly or partly by genetic changes passed from parents or arising in a reproductive cell or early embryo.

\medterm{Inhibin} A glycoprotein hormone produced mainly by the gonads that suppresses pituitary follicle-stimulating hormone secretion and contributes to regulation of reproductive function.

\medterm{Inhibition} Inhibition is a reduction or suppression of an activity, such as nerve signalling, enzyme action, muscle contraction, or cell growth.

\medterm{Inhibitory Neurochemical} A chemical messenger that decreases the activity of nerve cells or muscles. Gamma-aminobutyric acid and glycine are important inhibitory neurotransmitters in the nervous system.

\synonyms
Inhibitory neurotransmitter

\medterm{Inhibitory Synapse} A connection between nerve cells that reduces the chance that the receiving cell will generate an electrical signal.

\medterm{Iniencephaly} A rare, usually lethal neural-tube defect characterized by occipital bone defect, severe
retroflexion of the head, and marked shortening or distortion of the cervical and
thoracic spine. It is often identified prenatally and may occur with other congenital
anomalies.

\medterm{Iniparib} An investigational anticancer drug studied for effects on DNA repair and tumor-cell survival; it is not a standard treatment for most cancers.

\medterm{Initial Insomnia} Initial insomnia is difficulty falling asleep at the beginning of the sleep period, often related to stress, habits, medicines, mood, or another sleep disorder.

\medterm{Injectable Fillers} Materials placed under or within the skin to restore volume, improve contours, or soften wrinkles. Effects vary by product and may include swelling, bruising, infection, lumps, or rare blood-vessel blockage.

\synonyms
Dermal fillers

\medterm{Injection} Introduction of a medication, fluid, or other substance into the body with a needle or other injection device.

\medterm{Injury} Damage to the body's tissues or function caused by trauma, force, heat, chemicals, radiation, or another harmful exposure.

\medterm{Injury - Kidney And Ureter} Damage to the kidney or ureter from trauma, obstruction, surgery, infection, or other causes; severe injury can impair urine flow or kidney function.

\medterm{Injury Classification (Abrasion, Laceration, Contusion, Incised Wounds)} Forensic categorization of mechanical wounds. Abrasion: superficial friction injury
affecting only the epidermis (e.g., scrape, graze, pattern abrasion).

\synonyms
Abrasion, Laceration, Contusion, Incised Wounds

\medterm{Injury Dating} Estimation of the age of an injury using clinical appearance, histopathology, imaging,
and the history of symptoms. Exact dating is often uncertain, particularly for bruises
and soft-tissue injuries, and findings should be reported with appropriate limitations.

\medterm{Injury Site} The body location where physical harm occurred or where its effects are being examined, documented, treated, and monitored during recovery.

\medterm{Ink Poisoning} Illness after swallowing or significant exposure to ink; most small pen-ink exposures cause little toxicity, but symptoms depend on the product and amount.

\medterm{Ink Remover Poisoning} Poisoning caused by chemicals used to remove ink, which may irritate the skin, eyes, lungs, or digestive tract and may cause systemic toxicity.

\medterm{Innate Immune Evasion} Mechanisms by which pathogens evade the innate immune system. Strategies include
capsules that prevent phagocytosis, biofilms that protect against immune cells, and
intracellular survival within macrophages. Some bacteria produce enzymes that degrade
complement components. Others modify their surface molecules to avoid recognition.

\medterm{Innate Immunity} First line of defense against pathogens, present from birth. Rapid (minutes to hours),
germline-encoded, lacks immunological memory.

\medterm{Innate Lymphoid Cells} ILCs: lymphoid lineage, no rearranged antigen receptors, mirror Th subsets. ILC1 (T-bet,
IFN-gamma): intracellular pathogens, IBD. ILC2 (GATA3, IL-5/13): helminths, allergy,
asthma, tissue repair. ILC3 (ROR-gammat, IL-17/22): extracellular bacteria/fungi,
lymphoid organogenesis.

\medterm{Inner Ear} The deepest part of the ear, containing the cochlea for hearing and balance organs for movement sensing. Disease there can cause hearing loss, ringing, dizziness, or balance problems.

\medterm{Inner Ear Infection} Inner-ear inflammation can affect hearing and balance, producing vertigo, nausea, tinnitus, or hearing loss. Causes include viral infection, bacterial disease, autoimmune inflammation, or other disorders; sudden hearing loss, severe imbalance, or neurologic signs requires urgent assessment.

\medterm{Innervation} The nerve supply to an organ, muscle, or area of the body, including the functional control provided by those nerves.

\medterm{Innominate} Innominate means unnamed; in anatomy it may refer to structures such as the innominate artery or innominate bone, depending on context.

\medterm{INOCA} Ischemia with non-obstructive coronary arteries: objective myocardial ischemia or ischemic symptoms in a patient whose coronary arteries do not have an obstructive epicardial stenosis. Mechanisms include coronary microvascular dysfunction and vasospastic angina; evaluation should identify the mechanism so treatment can be tailored.

\medterm{Inoculation} The introduction of microorganisms or clinical material into a culture medium, host, or
sterile site to initiate growth, infection, or immune response.

\medterm{Inoperable} Not safely treatable by surgery, often because a disease is too extensive, located near vital structures, or the person is too medically fragile. Other treatments may still be possible.
\medterm{Inorganic} Not based mainly on carbon-containing molecular structures. Inorganic substances include water, salts, minerals, and metals, many of which are essential to body function while others can be toxic.

\medterm{Inorganic Chemical} A chemical compound that is generally not based on carbon chains, such as water, salts, minerals, or metals. In the body, inorganic substances support fluid balance, nerve signals, bone structure, and many chemical reactions.

\medterm{Inosculation} The joining and establishment of a connection between blood vessels or vascularized tissues, important in wound healing, grafts, flaps, and transplantation.

\medterm{Inosine} A nucleoside made of hypoxanthine linked to ribose that participates in purine metabolism and cellular energy pathways.

\medterm{Inosine Monophosphate} A nucleotide containing inosine and one phosphate group; it is an intermediate in purine metabolism and a precursor of other nucleotides.

\medterm{Inosine Pranobex} A medicine combination marketed in some countries as an immunomodulator or antiviral; evidence and approved uses vary by jurisdiction.

\medterm{Inositol} A sugar-like compound involved in cell membranes and signaling; some forms are studied or used for selected metabolic, reproductive, or mental-health conditions.

\medterm{Inositol Phosphates} Phosphorylated forms of inositol that act in cell signaling, calcium regulation, and other intracellular processes.
\medterm{Inotropes In Anesthesia} Medicines used during anesthesia or critical illness to increase the force of heart contraction and support blood flow when the heart is failing. They are given by controlled infusion and require monitoring because they can cause abnormal heart rhythms, low or high blood pressure, and increased oxygen demand.

\medterm{Inotropic} Affecting the force of cardiac muscle contraction; a positive inotropic drug increases contractility, while a negative inotropic effect decreases it.

\medterm{Inpatient} An inpatient is a person admitted to a healthcare facility for monitoring, diagnosis, treatment, surgery, or care requiring an overnight or formal stay.

\medterm{Inpatient Rehabilitation} Intensive inpatient rehabilitation (3+ hours therapy/day) for acute conditions: stroke,
spinal cord injury, traumatic brain injury, major trauma. See acute rehabilitation.

\medterm{Inquilinus} A genus of gram-negative bacteria associated mainly with water and industrial environments; human infection is rare.

\medterm{Inquilinus Limosus} A species of Inquilinus bacteria rarely isolated from human clinical specimens and occasionally associated with opportunistic infection.

\medterm{INR} International normalized ratio: a standardized measure of the prothrombin time,
calculated as $(\mathrm{patient\ PT}/\mathrm{mean\ normal\ PT})^{\mathrm{ISI}}$, to ensure consistent monitoring of
vitamin K antagonist (warfarin) therapy across laboratories. Target INR: 2-3 for most
indications (AF, VTE); 2.5-3.5 for mechanical heart valves.



\medterm{Insect Bite} A cutaneous reaction resulting from the injection of salivary contents from an insect
during feeding, causing a localized or systemic hypersensitivity response. Common
presentations include a pruritic papule or wheal with surrounding erythema. Severe
reactions may include anaphylaxis, particularly with Hymenoptera stings.

\medterm{Insect Bites} Skin injuries caused by an insect feeding or injecting saliva or venom, producing local inflammation and sometimes allergic, infectious, or systemic reactions.

\medterm{Insect Bites And Stings} Bites and stings from insects such as mosquitoes, fleas, bedbugs, bees, wasps, hornets, yellow jackets, and ants can cause local pain, itching, swelling, allergic reactions, infection, or toxin effects; breathing difficulty or widespread symptoms is an emergency.

\medterm{Insect Sting Allergies} Immune reactions to venom from insects such as bees, wasps, hornets, and fire ants. Reactions may be local or
systemic; generalized hives, breathing difficulty, dizziness, vomiting, or throat swelling
may indicate anaphylaxis and require emergency treatment.

\synonyms
Allergy Insect (Insect Sting Allergies)

\medterm{Insect Virus} A virus that infects insects or other arthropods; some are studied for ecology or biological pest control rather than human disease.

\medterm{Insecta} The biological class containing insects. Some insects bite or sting, cause allergic reactions, contaminate food, or spread infections and parasites; medical importance depends on the species and exposure.

\medterm{Insecticide} A chemical or biological agent used to kill or control insects. Exposure can cause skin, eye, respiratory, neurologic, or cholinergic toxicity depending on the product; prevention and treatment require identifying the active ingredient.

\medterm{Insecticide Poisoning} Toxic effects caused by exposure to an insect-killing chemical; symptoms vary by agent and may include gastrointestinal, neurologic, respiratory, or cholinergic manifestations.

\medterm{Insemination} Introduction of sperm into the female reproductive tract, either through intercourse or an assisted-reproduction procedure such as intrauterine insemination. Indications, success rates, infection screening, and risks depend on the method and patient factors.

\medterm{Insertion} Movement of a muscle attachment toward its fixed attachment, or placement of one structure into another.

\medterm{Insidious} Developing slowly and quietly, often without obvious early symptoms. An insidious illness may be present for a long time before it is recognized, so the term describes its gradual pattern rather than its cause or severity.
\medterm{Insinuate} To introduce or insert gradually into a space, or to suggest something indirectly; in anatomy, a structure may insinuate between or around other structures.

\medterm{Insomnia} See \textit{Sleeplessness}.

\synonyms
Sleeplessness

\medterm{Insomnia Disorder} Difficulty initiating or maintaining sleep, or early-morning awakening with inability to
return to sleep, at least three nights per week for at least three months. The
individual experiences clinically significant distress or impairment in functioning.

\medterm{Inspection} Systematic visual examination of the body or a specimen for abnormalities such as color change, swelling, deformity, lesions, movement, or discharge.
the patient\'s overall appearance, posture, body habitus, gait, skin color, nutritional
status, respiratory effort, and specific anatomical regions. Requires adequate lighting
and proper patient exposure.

\medterm{Inspiration (Inhalation)} The part of breathing in which air enters the lungs. The diaphragm contracts and moves downward, while the chest muscles expand the chest, lowering pressure inside the lungs and drawing air through the airways into the air sacs.

\medterm{Inspiratory Capacity} Inspiratory capacity is the maximum volume inhaled after a normal exhalation and equals tidal volume plus inspiratory reserve volume.

\medterm{Inspiratory Crackles} Fine, discontinuous sounds heard during inspiration, caused by the sudden opening of
collapsed airways. They are a hallmark of interstitial lung disease (bibasilar,
Velcro-like crackles) and pulmonary fibrosis. They are also heard in pneumonia,
atelectasis, and heart failure.

\medterm{Inspiratory Time} The length of time spent breathing in during a natural breath or mechanical-ventilation cycle. Adjusting it can affect airflow, oxygen delivery, comfort, pressure, and the balance between inhaling and exhaling.

\medterm{Instantaneous Wave-Free Ratio} A measurement made during heart catheterization to judge whether a narrowed heart artery limits blood flow. A pressure wire compares pressure beyond and before the narrowing while the heart is relaxed. A result of 0.89 or less may support treatment when it fits the symptoms and other findings.

\synonyms
iFR; Wave-Free Ratio; Resting Coronary Physiology

\medterm{Instillation} Instillation is the gradual introduction of a liquid medicine or solution into a body cavity, organ, or tissue.

\medterm{Instrument} A tool or device used by healthcare workers for examination, measurement, diagnosis, surgery, treatment, patient care, or laboratory work.

\medterm{Instrument Sterilization} A validated process that destroys all forms of microbial life, including bacterial
spores, on or within a medical instrument. Common methods include steam under pressure,
ethylene oxide, hydrogen-peroxide plasma, and dry heat, selected according to device
material and design.

\medterm{Instrumental Activities Of Daily Living} More complex activities required for independent community living, including medication
management, shopping, cooking, finances, communication, transportation, and household
tasks.

\medterm{Instrumental ADL} A complex activity needed for independent community living, such as shopping, cooking, managing medicines or finances, and using transportation.



\medterm{Insufficient Cervix} Painless cervical shortening and dilation during the second trimester, also called
cervical insufficiency. It can cause recurrent mid-trimester pregnancy loss or
spontaneous preterm birth without regular uterine contractions.

\synonyms
cervical insufficiency.

\medterm{Insula} The insula is a cerebral cortical region involved in sensation, taste, emotion, interoception, pain, and integration of internal and external information.

\medterm{Insulin} A hormone made by the pancreas that helps cells use glucose and lowers blood sugar. It is used to treat diabetes, but too much can cause dangerously low blood sugar with sweating, confusion, seizures, or unconsciousness.

\medterm{Insulin And Syringes - Storage And Safety} A patient-care topic concerning insulin and syringes - storage and safety. It addresses the relevant purpose, risks, management, and follow-up.

\medterm{Insulin Antagonist} A hormone, medicine, or other factor that opposes insulin’s effects and can raise blood glucose.

\medterm{Insulin Aspart} A rapid-acting insulin analog (modified amino acid sequence) with onset 10-20 minutes,
peak 1-3 hours, duration 3-5 hours. Used to cover prandial glucose (injected at or just
after meals) in type 1 and type 2 diabetes; also used in continuous subcutaneous insulin
infusion (pump) and IV in acute settings.

\medterm{Insulin Autoantibodies} Autoantibodies directed against insulin or other pancreatic beta-cell antigens,
including GAD, IA-2, and ZnT8 antibodies. Their presence supports autoimmune type 1
diabetes but a negative panel does not exclude it.

\medterm{Insulin C-Peptide Test} A blood or urine test measuring C-peptide released in equal amounts with endogenous insulin; it helps distinguish native insulin production from injected insulin and evaluate hypoglycemia or diabetes classification.

\medterm{Insulin Coma} A state of severe hypoglycemia and unconsciousness caused by excessive insulin; it is a medical emergency and is not a current routine treatment.

\medterm{Insulin Detemir} A long-acting insulin analogue used to provide basal insulin replacement in diabetes mellitus; dosing is individualized and hypoglycemia is the principal serious adverse effect.

\medterm{Insulin Glargine} A long-acting basal insulin analog with a flat, peakless profile providing near-24-hour
coverage, administered once daily (or twice in some patients). Formed by modification of
human insulin producing slow, steady absorption from the injection depot. Used for type
1 and type 2 diabetes (basal component).

\medterm{Insulin Injection Technique} Insulin is injected into subcutaneous tissue using the prescribed device, dose, and site. Common sites include
the abdomen, thighs, buttocks, and upper arms; rotating sites helps reduce lipohypertrophy
and variable absorption. Technique should be individualized to the person and device.

\medterm{Insulin Isophane} Intermediate-acting human insulin, also called neutral protamine Hagedorn insulin, used for basal glucose control in diabetes mellitus.

\medterm{Insulin Lipohypertrophy} Thickened fatty tissue at repeated insulin-injection sites. Injecting into these areas can make
insulin absorption unpredictable.

\medterm{Insulin Lispro} A rapid-acting insulin analogue used around meals or for correction of high blood glucose in diabetes mellitus.

\medterm{Insulin Pen Devices} Pre-filled or refillable devices that deliver measured insulin through a small needle under the skin. They are portable, but the dose, injection technique, storage, and needle disposal must follow medical instructions.

\medterm{Insulin Preparations} Forms of insulin used to lower blood glucose by helping body cells take up glucose and reducing glucose release by the liver. They differ in how quickly they start and how long they work, so the type and dose must be individualized.

\medterm{Insulin Pump Therapy} Delivery of rapid-acting insulin through a small tube under the skin. The pump gives a continuous background dose and extra doses around meals or to correct high blood glucose. Users need training, monitoring, and a backup plan if the device fails.

\medterm{Insulin Pumps} Insulin pumps deliver rapid-acting insulin continuously and through programmed meal or correction doses, allowing flexible management of diabetes while requiring monitoring and device safety.

\medterm{Insulin Receptor} A transmembrane receptor tyrosine kinase that binds insulin and initiates signaling controlling glucose uptake, metabolism, growth, and anabolic processes.

\medterm{Insulin Resistance} Diminished biological response to physiological insulin concentrations resulting in
impaired glucose disposal in peripheral tissues and failure to suppress hepatic glucose
production. Pathophysiological hallmark of type two diabetes and metabolic syndrome.

\synonyms
IR (Insulin Resistance)

\medterm{Insulin Resistance Syndrome} Alternate index label; see \textit{Metabolic syndrome}.

\medterm{Insulin Resistance Syndromes} Rare inherited severe insulin resistance syndromes include Donohue syndrome also known
as leprechaunism, Rabson-Mendenhall syndrome, and type A insulin resistance syndrome
caused by mutations in the insulin receptor gene. Donohue syndrome is the most severe
with death in infancy.

\medterm{Insulin Secretion} Release of insulin from pancreatic beta cells in response primarily to increased plasma
glucose, with amplification by incretin hormones and nutrients. Insulin promotes glucose
uptake, glycogen synthesis, lipogenesis, and protein synthesis while suppressing hepatic
glucose production.

\medterm{Insulin Shock} Severe hypoglycemia, historically associated with excessive insulin and characterized by confusion, seizures, loss of consciousness, or other neurologic dysfunction.

\medterm{Insulin Signaling} The cellular process initiated when insulin binds its receptor and activates pathways that regulate glucose uptake, glycogen formation, protein synthesis, and fat metabolism.

\medterm{Insulin-Dependent Diabetes} Type 1 diabetes, in which the pancreas produces little or no insulin and replacement insulin is required.

\synonyms
Type 1 diabetes

\medterm{Insulin-Induced Lipohypertrophy} A common complication of subcutaneous insulin therapy characterized by localized
accumulation of adipose tissue at injection sites, caused by the lipogenic effect of
insulin on subcutaneous fat. It presents as firm, rubbery, non-tender subcutaneous
nodules that can cause unpredictable insulin absorption and glycemic variability.

\medterm{Insulin-Like Growth Factor} A hormone, mainly IGF-1, made largely by the liver in response to growth hormone. It supports growth of bone and other tissues and helps regulate metabolism. Abnormally high or low levels can occur with pituitary, liver, nutritional, or chronic illnesses.

\medterm{Insulinase} An enzyme that breaks down insulin and helps control how long insulin remains active. Its activity contributes to regulation of blood sugar and may influence insulin requirements or treatment research.

\medterm{Insulinoma} An insulinoma is a usually benign pancreatic beta-cell neuroendocrine tumour that secretes excessive insulin, causing recurrent fasting or exertional hypoglycaemia relieved by glucose.

\medterm{Integrase} Integrase is an enzyme that inserts viral DNA into a host genome, a process targeted by integrase inhibitors used in HIV treatment.

\medterm{Integrative Medicine} Patient-centered care that combines evidence-based conventional treatment with appropriate complementary approaches, emphasizing whole-person health, prevention, lifestyle, and shared decision-making.

\medterm{Integrative Rehabilitation} Combination of conventional and complementary therapies: acupuncture, massage, yoga,
meditation.

\medterm{Integrin} A cell-surface receptor that links cells to the extracellular matrix and helps control adhesion, movement, signaling, and immune-cell activity.

\medterm{Integrin Binding} The attachment of a molecule or cell to an integrin receptor, which can influence adhesion, migration, signaling, or clot formation.

\medterm{Integron} A genetic element containing an integrase and an attachment site that can capture and express gene cassettes; integrons commonly contribute to bacterial antimicrobial resistance.

\medterm{Integrons} Genetic elements that capture and express gene cassettes through site-specific
recombination. Integrons contain an integrase gene and a recombination site for gene
cassette insertion. They are important vehicles for the spread of antibiotic resistance
genes, particularly in gram-negative bacteria.

\medterm{Integumentary System} The body system comprising the skin and its appendages: hair, nails, sweat glands, and
sebaceous glands. It forms a protective barrier, limits fluid loss, contributes to
thermoregulation and sensation, supports immune defence, and participates in vitamin D
synthesis.

\medterm{Intein} An intein is a protein segment that excises itself from a precursor protein while joining the surrounding extein sequences.

\medterm{Intellectual Disability} Neurodevelopmental disorder with onset in the developmental period, characterised by
deficits in intellectual and adaptive functioning across conceptual, social, and
practical domains. DSM-5 grades severity by adaptive functioning: mild, moderate,
severe, profound.

\medterm{Intelligence} Intelligence is the capacity to learn, reason, solve problems, understand information, and adapt; tests measure selected abilities rather than a person’s entire worth or potential.

\medterm{Intelligence Test} An intelligence test uses standardized tasks to estimate cognitive abilities such as reasoning, memory, language, and problem-solving.

\medterm{Intense Pulsed Light} A treatment using brief pulses of broad-spectrum light for selected skin conditions and hair reduction. Light is absorbed by pigment or blood vessels, and treatment can cause temporary redness, burns, pigment changes, or eye injury.

\synonyms
IPL

\medterm{Intensity} The degree of effort or physiological demand during physical activity.

\medterm{Intensity-Modulated Radiation Therapy} A computer-optimized form of three-dimensional radiation therapy that shapes and
modulates the radiation dose to conform tightly to the target while sparing adjacent
normal tissue.

\medterm{Intensive Care} A patient-care topic concerning intensive care. It addresses the relevant purpose, risks, management, and follow-up.

\medterm{Intensive Care Rehabilitation} Early rehabilitation in ICU setting. Early mobilization, positioning, passive range of
motion. Improves outcomes, reduces complications.

\medterm{Intensive Care Unit} A specialist hospital unit providing continuous monitoring and life support for
critically ill patients. In obstetrics, women may require ICU admission for: severe
pre-eclampsia or eclampsia, massive postpartum hemorrhage, amniotic fluid embolism,
sepsis, cardiac disease in pregnancy, or complications of anesthesia.

\medterm{Intention Tremor} A tremor that becomes more pronounced during a purposeful movement toward a target. It commonly suggests cerebellar or cerebellar-pathway dysfunction, although medicines, toxins, and other neurologic disorders can contribute.

\medterm{Interatrial Septum} The wall between the heart's right and left upper chambers. It includes the fossa ovalis, a normal remnant of fetal circulation, and contains a thinner area that may be crossed during selected catheter procedures.

\medterm{Intercalated Disc} An intercalated disc is a specialized junction connecting cardiac muscle cells and transmitting electrical and mechanical force through the heart.

\medterm{Intercalating Agent} An intercalating agent inserts between DNA base pairs and can disrupt replication or transcription; some are laboratory reagents or anticancer drugs.

\medterm{Intercellular} Located between cells or occurring in the space between cells. Intercellular material and signals help cells attach, communicate, exchange substances, and organize tissues throughout the body.

\medterm{Intercostal} Intercostal means located between the ribs, especially the muscles, nerves, and blood vessels in those spaces.

\medterm{Intercostal Muscles} Muscles between the ribs that help expand and compress the chest during breathing and also assist movement of the trunk.

\medterm{Intercostal Retraction} Inward movement of the skin and soft tissues between the ribs during inspiration,
indicating increased work of breathing. It is common in infants and children with
respiratory distress and may accompany airway obstruction, bronchiolitis, pneumonia,
asthma, or other pulmonary disease.

\medterm{Intercostal Retractions} Inward pulling of the skin and soft tissues between the ribs during inspiration, indicating increased work of breathing and potentially significant respiratory distress.

\medterm{Intercurrent Disease} An illness or complication that develops during the course of another disease or treatment and may affect its management or outcome.

\medterm{Interdental Brush} Small brush for cleaning between teeth. Used where spaces exist.

\medterm{Interdigital Neuroma} Alternate index label; see \textit{Morton neuroma}.

\medterm{Interference} A factor that alters the accuracy or interpretation of a laboratory test, measurement, or drug effect.

\medterm{Interferometry} A measurement method that uses interference between waves to determine distance, motion, thickness, optical properties, or other physical quantities.

\medterm{Interferon} Interferons are cytokines that induce antiviral defenses, alter immune-cell activity, and regulate cell growth; type I, II, and III interferons have different receptors and clinical uses.

\medterm{Interferon Inducer} A substance that stimulates cells to produce interferons, proteins involved in antiviral and immune responses.
\medterm{Interferon Signaling} Type I IFN (IFN-alpha/beta): all nucleated cells produce; receptor IFNAR (ubiquitous);
JAK1/TYK2 -> STAT1/2/IRF9 (ISGF3) -> ISGs (hundreds: antiviral, antiproliferative,
immunomodulatory). Type II IFN (IFN-gamma): T/NK cells; receptor IFNGR; JAK1/JAK2 ->
STAT1 -> GAS elements. Type III IFN (IFN-lambda): mucosal epithelium; receptor IFNLR.

\medterm{Interferon-Alpha} A type I interferon cytokine with antiviral, antiproliferative, and immune-modulating effects; recombinant forms have been used for selected viral infections and cancers.

\medterm{Interferon-Beta} A type I interferon cytokine used mainly as an immunomodulatory treatment for relapsing forms of multiple sclerosis.

\medterm{Interferon-Gamma Release Assay} Blood test for Mycobacterium tuberculosis infection (latent or active). T-SPOT.TB(ELISPOT) or QuantiFERON-TB Gold Plus. Measures IFN-gamma released by T cells
after exposure to MTB-specific antigens (ESAT-6, CFP-10, TB7.7).

\synonyms
Interferon-Gamma Release Assay

\medterm{Interferonopathies} A group of rare disorders caused by abnormal activation of type I interferon pathways. They can produce inflammation,
neurologic disease, skin lesions, blood-vessel injury, or lung disease, often beginning
in childhood.

\medterm{Interferons} Natural proteins that cells release to warn nearby cells and help control the immune response. They are especially important in the body's defense against viruses and are also used as medicines for some infections, cancers, and immune conditions.

\medterm{Interleukin} A family of cytokines (interleukins) produced mainly by leukocytes that regulate immune
cell growth, differentiation, activation, and communication. Examples: IL-1 and IL-6
(inflammation, fever), IL-2 (T-cell growth), IL-4/IL-13 (Th2, IgE), IL-5 (eosinophils),
IL-10 (anti-inflammatory), IL-12 (Th1), IL-17 (Th17, autoimmunity), IL-23.

\medterm{Interleukin 1} A group of proinflammatory cytokines, especially IL-1 alpha and IL-1 beta, that promote fever, leukocyte activation, and inflammatory responses.

\medterm{Interleukin 13} A cytokine produced by type 2 helper T cells and other immune cells that contributes to allergic inflammation, mucus production, and tissue fibrosis.

\medterm{Interleukin 2} Interleukin-2 is a T-cell growth factor that promotes proliferation and survival of activated T cells and natural-killer cells; high-dose therapy can cause capillary-leak syndrome and systemic toxicity.

\medterm{Interleukin 7} A cytokine important for development, survival, and homeostasis of T lymphocytes and some B-cell precursors.
\medterm{Interleukin-6} A cytokine involved in immune regulation and inflammation that stimulates production of acute-phase proteins by the liver.

\medterm{Interleukins} Cytokines that carry signals between immune cells and regulate inflammation.

\synonyms
ILs

\medterm{Interlocutory} Interlocutory means provisional or occurring during an ongoing process, such as an interim clinical, legal, or research decision.

\medterm{INTERMACS Profile} A seven-level system for describing how sick a person is with advanced heart failure. Profile 1 means critical shock with very poor blood flow; profiles 2 and 3 describe worsening or dependence on medicines that strengthen the heart. Profiles 4 through 7 describe repeated hospital stays or increasing limits on activity. The profile helps doctors judge urgency and timing for a heart pump or transplant.

\synonyms
INTERMACS Staging; INTERMACS Scale

\medterm{Intermediate} Intermediate means between two stages, positions, or levels, such as an intermediate metabolite, cell, or developmental phase.

\medterm{Intermediate Cuneiform} Smallest cuneiform, centrally located between medial and lateral cuneiforms. Articulates
with navicular, medial/lateral cuneiforms, and second metatarsal. Forms the transverse
arch apex. Relatively fixed, contributing to midfoot stability. Injuries are typically
associated with high-energy midfoot trauma.

\medterm{Intermediate Filament} An intermediate filament is a strong cytoskeletal fiber that supports cell shape, mechanical stability, and tissue integrity.

\medterm{Intermediate Uveitis} Inflammation mainly affecting the gel inside the eye and nearby tissues at the edge of the retina. It may cause blurred vision, floaters, or little pain and can be linked to autoimmune disease, infection, or no identified cause. Eye examination and imaging guide treatment and monitoring for complications.

\medterm{Intermediate-Density Lipoprotein} A temporary particle formed when the body removes some fat from very-low-density lipoprotein. It can be changed into LDL, the type of cholesterol particle that can contribute to fatty buildup in artery walls.

\synonyms
IDL

\medterm{Intermittent Auscultation} Periodic assessment of the fetal heart rate with a fetoscope or Doppler device
during labor rather than continuous electronic monitoring. The frequency of assessment depends on the stage of labor,
risk factors, and local clinical protocol.

\medterm{Intermittent Claudication} Reproducible muscle discomfort, usually in the calf, thigh, or buttock, induced by
walking and relieved by rest, caused by exertional arterial insufficiency. It is
distinguished from rest pain and nonvascular musculoskeletal pain by history and
vascular testing.

\medterm{Intermittent Explosive Disorder} A disorder characterized by recurrent, impulsive, problematic outbursts of verbal or
physical aggression that are grossly disproportionate to the provocation. The outbursts
are not premeditated and cause marked distress or functional impairment. Age >6 years.

\synonyms
Road Rage

\medterm{Intermittent Fasting} An eating pattern that alternates planned periods of eating with periods of little or no food. Common schedules limit daily eating hours or use selected fasting days. It may help some adults lose weight, but it is not suitable for everyone, especially during pregnancy or with some medical conditions or eating disorders.

\medterm{Intermittent Fever} A fever that rises and falls with periods when temperature returns to normal or near normal.

\medterm{Intern} A physician in the early stage of postgraduate medical training, traditionally the first year after medical school; the exact training structure varies by country.

\medterm{Internal} Internal means located within the body, an organ, a structure, or a system rather than on its outer surface.

\medterm{Internal Bleeding} Internal bleeding is blood loss into a body cavity, tissue, or organ rather than through an obvious external wound. Trauma, ulcers, aneurysm, pregnancy complications, medicines, and tumors can cause it; faintness, shock, abdominal swelling, black stool, or vomiting blood is an emergency.

\medterm{Internal Capsule} A compact band of white-matter nerve fibres in the brain carrying important motor and sensory pathways between the cerebral cortex and lower centres.

\medterm{Internal Environment} The fluid and physiologic conditions inside the body that cells require, including temperature, pH, oxygen, electrolytes, and nutrients.

\medterm{Internal Fixation} Surgical stabilization of a fracture or unstable bone using implanted hardware such as plates, screws, rods, wires, or pins.

\medterm{Internal Fixator} An internal fixator is a surgically implanted plate, screw, rod, nail, or frame used to stabilize a fracture or correct skeletal alignment.

\medterm{Internal Hernia} Herniation of bowel through a mesenteric defect common after Roux-en-Y gastric bypass
creating a defect at the jejunojejunostomy. May cause intermittent or acute small bowel
obstruction. Presents with episodic abdominal pain, nausea, and vomiting.

\medterm{Internal Iliac Artery} Smaller common iliac terminal branch (hypogastric), supplying the pelvis, perineum,
gluteal region, and medial thigh. Posterior division: iliolumbar, lateral sacral,
superior gluteal. Anterior division: inferior gluteal, internal pudendal, obturator,
vesical, uterine, vaginal, middle rectal arteries.

\medterm{Internal Jugular Vein} Descends in the carotid sheath with the internal carotid artery and vagus nerve,
draining brain, face, and neck. Begins at the jugular foramen as a sigmoid sinus
continuation, joining the subclavian vein to form the brachiocephalic vein. Contains
backflow-preventing valves. Used for central venous catheterization.

\medterm{Internal Oblique Muscle} Middle flat abdominal layer deep to external oblique, from thoracolumbar fascia, iliac
crest, and lateral inguinal ligament to lower ribs, linea alba, and pubic crest.
Innervated by T7-T12 and L1. Ipsilateral rotation (opposite of external oblique).
Contributes to inguinal canal and abdominal wall integrity.

\medterm{Internal Podalic Version} A maneuver where the operator inserts a hand into the uterus, grasps one or both fetal
feet, and converts a transverse lie or cephalic presentation to a breech, then delivers
by total breech extraction. Indication: delivery of the second twin (transverse lie) or
severe fetal distress in the second stage. Rare in modern obstetrics.

\medterm{Internal Resorption} Destruction of tooth tissue from the inside of the tooth outward, caused by activity of specialized cells. It may produce a pink area or no symptoms and can weaken the tooth; dental imaging determines whether root-canal treatment can save it.

\medterm{International Normalized Ratio} The INR, a standardized calculation from the prothrombin time used to assess blood clotting and guide treatment with vitamin K antagonists such as warfarin.

\medterm{International Unit} A unit that measures the biological effect of a substance, such as insulin, a vitamin, or a vaccine. The amount represented by one unit differs between substances, so international units cannot be directly converted to milligrams without knowing the substance.

\medterm{Interneuron} An interneuron connects neurons within the central nervous system and helps integrate, inhibit, or coordinate neural signals.

\medterm{Interneurons} Neurons within the central nervous system that connect other neurons and help process or coordinate neural signals.

\medterm{Internist} An internist is a physician specializing in adult internal medicine, including diagnosis and management of complex medical conditions.

\medterm{Internuclear Ophthalmoplegia} A conjugate horizontal gaze disorder caused by a lesion of the medial longitudinal
fasciculus (MLF), a heavily myelinated fiber tract that connects the abducens nucleus
(CN VI) on one side to the oculomotor nucleus (CN III) on the contralateral side,
coordinating horizontal gaze.

\medterm{Interosseous} Interosseous means located between bones, such as the membrane or muscles between the forearm or lower-leg bones.

\medterm{Interpersonal Psychotherapy} A time-limited, evidence-based psychotherapy for depression that focuses on the
interpersonal context of mood symptoms. It identifies four problem areas: grief, role
disputes, role transitions, and interpersonal deficits. IPT helps patients improve
interpersonal functioning and social support.

\medterm{Interphase} Interphase is the portion of the cell cycle between divisions, including growth, DNA replication, and preparation for mitosis or meiosis.

\medterm{Interphase Cell} A cell between divisions that grows, performs normal functions, and copies its DNA before possible division.

\medterm{Interproximal Caries} Dental caries occurring on the contact surfaces between adjacent teeth. It may be
difficult to detect clinically and is commonly assessed with bitewing radiographs when
imaging is indicated.

\medterm{Intersexual} An older term for a person born with physical sex characteristics that do not fit typical definitions of male or female. Modern care uses “difference of sex development” and respects the person’s identity and choices.

\medterm{Interstitial} Relating to the spaces between cells or tissues, such as interstitial fluid or interstitial lung disease.

\medterm{Interstitial (Cornual) Pregnancy} An ectopic pregnancy implanted in the part of the fallopian tube that passes
through the muscular wall of the uterus. It can enlarge before rupture and cause severe internal bleeding.

\synonyms
Cornual

\medterm{Interstitial Cystitis} Long-lasting bladder or pelvic pain with frequent or urgent urination, often without infection. Symptoms may flare and improve; evaluation rules out other causes, and treatment may include bladder training, physical therapy, lifestyle changes, or selected medicines.

\synonyms
Bladder Inflammation
Cystitis, Interstitial
IC (Interstitial Cystitis)
Painful Bladder Syndrome

\medterm{Interstitial Fluid} Interstitial fluid is the fluid surrounding tissue cells, formed from blood plasma and returned to circulation through lymphatic vessels.

\medterm{Interstitial Granulomatous Dermatitis} A rare inflammatory dermatosis with interstitial or palisading histiocytes around collagen, often presenting with symmetric erythematous or violaceous plaques and associated systemic disease or medication exposure.

\medterm{Interstitial Keratitis} Inflammation of the corneal stroma without primary involvement of the corneal surface.
It may follow infection, autoimmune disease, or trauma and can cause pain, photophobia,
redness, and reduced vision.

\medterm{Interstitial Lung Disease} A heterogeneous group of disorders with inflammation and fibrosis of lung parenchyma,
causing restrictive physiology (reduced FVC and TLC with normal or elevated FEV1/FVC
ratio).

\synonyms
Diffuse Parenchymal Lung Diseases
Interstitial Lung Diseases
Lung Disease, Interstitial

\medterm{Interstitial Lung Disease (Interstitial Pneumonia)} Alternate index label for Interstitial Pneumonia.



\medterm{Interstitial Nephritis} Inflammation of the renal interstitium and tubules, which can be acute or chronic. Acute
interstitial nephritis is most commonly caused by medications, particularly NSAIDs,
antibiotics including penicillins and cephalosporins, and proton pump inhibitors.

\medterm{Interstitial Pneumonia} Inflammation affecting the supporting tissue around the lung air sacs. It may be caused by infection, autoimmune disease, medicines, or other lung disorders and can cause dry cough, fever, breathlessness, and low oxygen. Examination and imaging help determine the cause and treatment.

\synonyms
Interstitial Lung Disease (Interstitial Pneumonia)

\medterm{Interstitium} The supporting tissue between and around the lung air sacs, capillaries, and small airways; inflammation or fibrosis of this tissue can impair lung expansion and gas exchange.

\medterm{Intertrigo} Intertrigo is inflammation in opposing skin folds caused by moisture, heat, friction, and maceration; Candida or bacterial infection may complicate the tender erythematous rash.

\synonyms
Allergy Diaper (Intertrigo)

\medterm{Interval} An interval is the time, distance, or range between two points or events, such as a dosing or follow-up interval.

\medterm{Intervening Sequence} A nucleotide sequence within a gene transcript that lies between coding sequences; in many genes it is an intron removed during RNA processing.

\medterm{Intervention} An action intended to prevent, diagnose, treat, or change a health condition or outcome, such as a medicine, procedure, counseling, or public-health measure.

\medterm{Intervention Study} A study in which researchers assign or provide an intervention and compare its outcomes with a control.

\synonyms
Clinical trial

\medterm{Interventional} Relating to a procedure or active treatment used to diagnose or manage disease, often through minimally invasive access rather than open surgery.

\medterm{Interventional Cardiology} A subspecialty using catheter-based procedures to diagnose and treat coronary, structural, and selected vascular heart conditions, including angioplasty, stenting, and transcatheter valve procedures.

\medterm{Interventional Procedure} A targeted procedure used to diagnose or treat disease, often through a needle, tube, or small incision and guided by imaging. Examples include biopsy, drainage, vessel treatment, and tumor ablation; risks depend on the body area and technique.

\medterm{Interventional Procedures} Procedures that diagnose or treat disease through a needle, catheter, or small incision, often using imaging for guidance. Examples include biopsy, drainage, vessel opening, embolization, and ablation; less invasive does not mean risk-free.

\medterm{Interventional Radiologist} A physician trained to perform minimally invasive diagnostic and therapeutic procedures using medical imaging for guidance.

\medterm{Interventional Radiology} A specialty using imaging guidance to perform minimally invasive diagnostic and therapeutic procedures, such as biopsies, drainages, embolization, and vascular treatments.

\medterm{Interventricular Septum} The thick muscular wall separating the left and right ventricles. It has two portions:
the muscular (inferior, thicker) and membranous (superior, thinner) portions. The septum
contains parts of the cardiac conduction system, including the bundle of His.

\medterm{Intervertebral Disc} A fibrocartilaginous cushion between adjacent vertebral bodies, composed of a gelatinous
nucleus pulposus surrounded by a tough annulus fibrosus. It absorbs shock and permits
limited movement of the vertebral column.

\medterm{Intervertebral Disk} A strong, flexible cushion between neighboring spinal bones that absorbs shock and allows movement. A damaged or herniated disk can press on nearby nerves and cause back pain, numbness, tingling, or weakness.

\synonyms
Intervertebral disc

\medterm{Intervertebral Foramen} An opening between two neighboring bones of the spine through which a spinal nerve leaves the spinal canal. Narrowing of this opening can press on the nerve and cause pain, numbness, tingling, or weakness in its area of the body.

\synonyms
Neural foramen

\medterm{Intestinal Absorption} Intestinal absorption is the movement of digested nutrients, water, and electrolytes from the intestine into the blood or lymph.

\medterm{Intestinal Atresia} A congenital obstruction of the small intestine (jejunum or ileum) caused by an
intrauterine vascular accident (mesenteric ischemia) leading to aseptic necrosis and
resorption of a segment of bowel. In contrast to duodenal atresia, it is not associated
with Down syndrome.

\medterm{Intestinal Cancer} Intestinal cancer is a malignant growth arising in the small or large intestine; symptoms and treatment depend on its location and stage.

\medterm{Intestinal Duplication} Intestinal duplication is a congenital extra segment or cyst of bowel that may cause obstruction, bleeding, pain, or infection.

\medterm{Intestinal Fistula} An intestinal fistula is an abnormal passage connecting the intestine to another organ or the skin, often after inflammation, surgery, or injury.

\medterm{Intestinal Ganglioneuromatosis} Benign proliferation of ganglion cells and nerve fibers in the gastrointestinal tract
wall. May be diffuse involving the entire gastrointestinal tract or localized as
ganglioneuromas. Associated with multiple endocrine neoplasia type 2B where intestinal
ganglioneuromatosis may precede the development of medullary thyroid carcinoma.

\medterm{Intestinal Gas} Gas within the digestive tract that may cause belching, bloating, abdominal discomfort, passing wind, or symptoms from an underlying digestive disorder.

\medterm{Intestinal Gas (Belching, Bloating, Flatulence)} Alternate index label for Belching, Bloating, Flatulence.

\medterm{Intestinal Gas (Intestinal Gas (Belching, Bloating, Flatulence))} Alternate index label for Intestinal Gas (Belching, Bloating, Flatulence).

\medterm{Intestinal Ischemia} Too little blood flow reaching the intestines. It may develop suddenly from a clot or blockage, or gradually from narrowed arteries, and can cause severe pain, bloody stools, or bowel damage; sudden severe pain needs emergency care.

\synonyms
Ischemia, Intestinal

\medterm{Intestinal Leiomyoma} A usually benign smooth-muscle tumor arising in the wall of the intestine; symptoms depend on size and location and may include bleeding, obstruction, or abdominal pain.

\medterm{Intestinal Lipodystrophy} Alternate index label; see \textit{Whipple's disease}.

\medterm{Intestinal Lipomatosis} Benign overgrowth of adipose tissue within the intestinal wall, most commonly affecting
the ileum and colon. Often asymptomatic and discovered incidentally, but large lipomas
can cause intussusception, gastrointestinal bleeding, or bowel obstruction.

\medterm{Intestinal Lymphoma} Malignant proliferation of lymphoid cells arising in the small or large intestine.
Important types include enteropathy-associated T-cell lymphoma, monomorphic
epitheliotropic intestinal T-cell lymphoma, and diffuse large B-cell lymphoma.

\medterm{Intestinal Malrotation} Congenital failure of normal midgut rotation and fixation, predisposing to a narrow
mesenteric base, volvulus, and duodenal obstruction. Infants may present with bilious
vomiting, abdominal distension, or bowel ischemia; sudden symptoms require emergency
assessment.

\medterm{Intestinal Metaplasia} Replacement of the normal lining of an intestinal or gastric region by intestinal-type epithelium, usually in response to chronic injury or inflammation and sometimes associated with increased cancer risk.

\medterm{Intestinal Mucosa} The intestinal mucosa is the inner lining of the intestine, where digestion, nutrient absorption, mucus production, and immune defence occur.

\medterm{Intestinal Obstruction} A mechanical or functional blockage that prevents normal passage of intestinal contents.
Causes include adhesions, hernias, tumors, inflammation, volvulus, and disorders of intestinal motility.

\medterm{Intestinal Obstruction And Ileus} Impaired passage of intestinal contents caused by a mechanical blockage or by reduced bowel motility without a physical blockage; symptoms commonly include abdominal pain, distension, vomiting, and inability to pass stool or gas.

\medterm{Intestinal Obstruction Repair} Surgery or another intervention to relieve an intestinal blockage, remove a diseased segment when necessary, and restore bowel continuity or create a stoma when required.


\medterm{Intestinal Parasite} A protozoan or worm that lives in the gastrointestinal tract and may cause diarrhea, abdominal pain, anemia, malnutrition, or no symptoms.

\medterm{Intestinal Perforation} Full-thickness breach in the bowel wall allowing luminal contents to enter the
peritoneal cavity. Causes include diverticulitis, appendicitis, peptic ulcer disease,
and trauma. Presents with acute abdomen, peritonitis, and free air on imaging.

\medterm{Intestinal Polyp} An intestinal polyp is a growth projecting from the lining of the intestine; some types can become cancerous.

\medterm{Intestinal Polyposis} Intestinal polyposis is the presence of numerous polyps in the intestine, sometimes as part of an inherited cancer-predisposition syndrome.

\medterm{Intestinal Pseudo-Obstruction} Acute or chronic dilation of the bowel without mechanical obstruction due to
neuromuscular dysfunction. Acute colonic pseudo-obstruction also called Ogilvie syndrome
occurs in hospitalized or post-surgical patients.

\synonyms
Ogilvie syndrome.

\medterm{Intestinal Villi} Intestinal villi are tiny finger-like projections of the small-intestinal lining that increase surface area for nutrient absorption.

\medterm{Intestine} The part of the digestive tract extending from the stomach to the anus, consisting of the small intestine and large intestine.

\medterm{Intestine Disorder} An intestine disorder is a disease affecting the small or large intestine, such as infection, inflammation, obstruction, or cancer.

\medterm{Intestine Neoplasm} An intestine neoplasm is an abnormal growth in the small or large intestine; it may be benign or malignant.

\medterm{Intestine Obstruction} Intestinal obstruction is partial or complete blockage of the passage of intestinal contents, causing pain, swelling, vomiting, and failure to pass stool or gas.

\medterm{Intestine, Diseases Of} Diseases of the small or large intestine include infection, inflammation, poor absorption, blockage, polyps, and cancer. Symptoms may include pain, diarrhea, constipation, bleeding, vomiting, weight loss, or nutritional deficiency.

\medterm{Intimal Dissection} A tear in the inner lining of an artery that allows blood to enter and split the vessel wall, creating a false lumen. Aortic or other arterial dissection can compromise blood flow and is a medical emergency.

\medterm{Intimal Hyperplasia} Proliferation and migration of vascular smooth-muscle cells with extracellular-matrix
deposition within the intima after vascular injury or graft placement. It can produce
restenosis at anastomoses, stents, dialysis access, or bypass grafts.

\medterm{Intimate Partner Violence In Pregnancy} Physical, sexual, or psychological abuse by a current or former
partner during pregnancy. It can harm the pregnant person and fetus and is associated with complications such as injury,
preterm birth, and placental abruption.

\medterm{Intolerance} An adverse response to a substance that does not involve the immune mechanism of an allergy.

\synonyms
Nonallergic sensitivity
Nonceliac Gluten Sensitivity (Intolerance)

\medterm{Intolerance, Gluten (Nonceliac Gluten Sensitivity (Intolerance))} Alternate index label for Nonceliac Gluten Sensitivity (Intolerance).

\medterm{Intoplicine} Intoplicine is an investigational topoisomerase inhibitor studied for anticancer effects by disrupting DNA structure and replication.

\medterm{Intoxication} A state of physical or mental impairment caused by recent exposure to a substance, such as alcohol, a medicine, or another drug.

\medterm{Intra-Abdominal Abscess} Localized collection of pus within the abdomen commonly in the pelvis, between bowel
loops, or beneath the diaphragm. Results from intra-abdominal infection or surgery.
Presents with fever, abdominal pain, and leukocytosis. Diagnosed by CT with contrast
showing rim-enhancing fluid collection.

\medterm{Intra-Abdominal Hypertension} Sustained elevation of intra-abdominal pressure above twelve millimeters of mercury
measured through the urinary bladder. Graded as grade one twelve to fifteen, grade two
sixteen to twenty, grade three twenty-one to twenty-five, and grade four greater than
twenty-five millimeters of mercury.

\medterm{Intra-Aortic Balloon Pump} Mechanical circulatory support. Inflatable catheter in descending aorta. Inflates during
diastole (increases coronary perfusion), deflates during systole (reduces afterload).
Indications: cardiogenic shock, unstable angina, high-risk PCI. Complications: limb
ischemia, bleeding, infection.

\medterm{Intra-Arterial} Intra-arterial means within an artery, including a catheter, injection, infusion, or measurement performed directly in an arterial vessel.

\medterm{Intraabdominal Cancer} A broad term for malignancies arising within the abdominal cavity, including cancers of
the colon, stomach, pancreas, liver, ovaries, and peritoneum. Presentation varies widely
depending on the specific organ involved.

\medterm{Intraabdominal Hemorrhage} Bleeding inside the abdominal cavity, which may follow trauma, surgery, organ rupture, pregnancy complications, or internal disease.

\medterm{Intraabdominal Hernia} Herniation of abdominal viscera through a congenital or acquired defect or abnormal
mesenteric aperture within the peritoneal cavity, distinct from ventral or inguinal
hernias. Types include paraduodenal, foramen of Winslow, transmesocolic, and
transomental hernias.

\medterm{Intraabdominal Infection} An infection within the abdominal cavity, such as peritonitis or an abscess, often requiring antibiotics and sometimes drainage or surgery.

\medterm{Intraaortic Balloon Pump} A temporary circulatory-support device with a balloon in the descending aorta that inflates and deflates in coordination with the cardiac cycle.

\medterm{Intraarticular} Intraarticular means within a joint, such as an injection, fracture, loose body, or inflammatory process inside the joint space.

\medterm{Intrabronchial} Intrabronchial means within a bronchus, one of the air passages branching from the trachea into the lungs.

\medterm{Intracardiac Electrophysiology Study} An invasive test using electrode catheters inside the heart to record electrical activity and provoke or map arrhythmias, helping guide diagnosis and treatment.

\medterm{Intracartilaginous} Located within cartilage, such as a lesion, injection, or developing structure. The term identifies position and does not by itself indicate whether the finding is normal, injured, inflamed, or cancerous.

\medterm{Intracavernosal Injection} An injection of a prescription medicine into the erectile tissue of the penis to treat erectile dysfunction. The medicine relaxes muscle and widens blood vessels, producing an erection that does not require sexual stimulation. Pain, bruising, scar tissue, or an erection lasting too long can occur; a prolonged painful erection needs urgent treatment.

\synonyms
ICI; Penile Injection Therapy; Intracavernous Injection

\medterm{Intracavitary} Intracavitary means within a body cavity, such as a treatment, device, fluid, or lesion located in the thorax, abdomen, uterus, or another cavity.

\medterm{Intracavity Radiotherapy} Intracavity radiotherapy places a radiation source inside a body cavity near a tumor, commonly to treat selected cancers of the cervix, uterus, or vagina.

\medterm{Intracellular} Inside a cell or occurring within a cell. Intracellular fluid contains most of the body’s potassium and magnesium, whereas sodium is found mainly outside cells. The cell membrane controls movement of water and dissolved substances between the cell and its surroundings.

\medterm{Intracellular Fluid} The fluid contained within cells, comprising approximately two thirds of total body
water in a typical adult. Its major solutes include potassium, magnesium, phosphate, and
intracellular proteins, and its composition is maintained by selective membrane
transport.

\medterm{Intracellular Membrane} A membrane inside a cell that separates compartments and controls movement of substances between them.

\medterm{Intracellular Pathogens} Microorganisms that survive and replicate within host cells. Intracellular bacteria
include Mycobacterium, Listeria, and Chlamydia. Viruses are obligate intracellular
pathogens. Intracellular survival protects pathogens from antibiotics and immune
responses.

\medterm{Intracellular Transport} Movement of proteins and other materials within a cell, often in membrane-bound sacs along the cell's internal framework.

\medterm{Intracerebral} Intracerebral means within the brain tissue, as in an intracerebral hemorrhage, injection, electrode, or lesion.
\medterm{Intracerebral Hemorrhage} Bleeding directly into the brain parenchyma, accounting for approximately ten to fifteen
percent of all strokes and carrying a high mortality rate of approximately forty percent
at thirty days.

\medterm{Intracisternal} Intracisternal means within a cistern, especially a cerebrospinal-fluid space at the base of the brain.
\medterm{Intracoronal} Intracoronal means within the crown of a tooth or within a coronary structure, with the exact meaning determined by context.

\medterm{Intracoronary} Intracoronary means within a coronary artery or the coronary circulation, such as a catheter, stent, imaging study, or medicine delivery.

\medterm{Intracranial} Inside the skull. The term may describe the brain, blood vessels, fluid, bleeding, pressure, infection, or a growth within the skull. Extra fluid, blood, or tissue can raise pressure and damage the brain, causing headache, vomiting, confusion, weakness, seizures, or loss of consciousness.

\medterm{Intracranial Aneurysm} A weakened, bulging area in a blood vessel inside the skull, usually in an artery supplying the brain. Many cause no symptoms, but a leak or rupture can cause bleeding around the brain and a sudden, severe headache.

\medterm{Intracranial Hemorrhage} Bleeding within the cranial cavity. Epidural hematoma: arterial bleeding (middle
meningeal artery), typically from temporal bone fracture; lucid interval then rapid
deterioration. Subdural hematoma: venous bleeding from bridging veins; acute (<72h),
subacute, or chronic (elderly, alcoholics).

\medterm{Intracranial Pressure} The pressure within the cranial vault generated by brain tissue, cerebrospinal fluid,
and intracranial blood. Sustained elevation can reduce cerebral perfusion and cause
headache, vomiting, altered consciousness, pupillary abnormalities, and brain
herniation.

\medterm{Intracranial Pressure Monitoring} Continuous or intermittent measurement of pressure inside the skull, usually in critically ill patients with severe brain injury, hemorrhage, hydrocephalus, or other risk of intracranial hypertension.

\medterm{Intracranial Venous Malformations} Abnormal collections or connections of veins within the skull that may bleed or affect nearby brain tissue.

\medterm{Intractable} Difficult to control, relieve, or cure despite appropriate treatment, as in intractable pain or seizures.
\medterm{Intradermal} Located within the dermis, the layer of skin beneath the epidermis; an intradermal injection places medication or test material into that layer.

\medterm{Intraductal Carcinoma} Alternate index label; see \textit{Ductal carcinoma in situ}.

\medterm{Intraductal Papillary Mucinous Nneoplasm Pancreas (Pancreatic Cysts)} Alternate index label for Pancreatic Cysts.

\medterm{Intraductal Papilloma} A benign papillary tumor arising within a breast duct, often presenting with spontaneous
unilateral bloody or clear nipple discharge. It may be solitary near the nipple or
multiple in the peripheral breast.

\medterm{Intraduodenal} Intraduodenal means within the duodenum, the first part of the small intestine beyond the stomach.
\medterm{Intraepidermal} Intraepidermal means within the epidermis, the outermost layer of skin.

\medterm{Intraepithelial} Located within the layer of cells that covers a body surface or lines an organ. The word may describe infection, inflammation, or abnormal cell changes that have not yet grown through the supporting layer beneath them.

\medterm{Intraesophageal} Intraesophageal means within the esophagus, such as a probe, pressure sensor, lesion, or medication delivery.
\medterm{Intrahepatic} Within the liver. For example, intrahepatic bile-duct blockage or inflammation can cause jaundice and abnormal liver tests, while a tumor may begin in liver cells or in bile ducts inside the liver. The exact meaning depends on the condition being described.

\medterm{Intrahepatic Cholestasis} Reduced bile flow within the liver causing bile substances to build up in blood, often producing itching, jaundice, dark urine, or pale stools.

\medterm{Intrahepatic Cholestasis Of Pregnancy} A pregnancy-associated liver disorder characterized mainly by
itching without a primary rash and increased bile acids in the blood. It may increase fetal risk and requires clinical
evaluation and monitoring.

\medterm{Intralesional} Delivered directly into a lesion, such as an injection of corticosteroid into a keloid or a medicine into a skin, joint, or tumor lesion. The drug, dose, sterility, and lesion type determine safety.

\medterm{Intralymphatic} Intralymphatic means within a lymphatic vessel or channel, such as a medicine, tumor spread, infection, or imaging contrast.

\medterm{Intramedullary} Intramedullary means within the medulla of an organ or the central canal of a bone, especially an implant placed inside a long bone.

\medterm{Intramedullary Nail} A metal implant inserted into the medullary canal of a long bone to align and stabilize
a fracture. It shares load along the bone axis and is commonly used for femoral and
tibial shaft fractures.

\medterm{Intramural Hematoma} Collection of blood within the wall of the gastrointestinal tract most commonly the
duodenum. Caused by trauma, anticoagulation, or coagulopathy. Presents with abdominal
pain, vomiting, and sometimes obstructive symptoms. Diagnosis is by CT showing
circumferential wall thickening with high-attenuation hematoma.

\medterm{Intramural Plaque} An intramural plaque is an atherosclerotic deposit within an artery wall that can narrow the lumen or destabilize and cause thrombosis.

\medterm{Intramuscular} Within a muscle. An intramuscular injection delivers a medicine into muscle tissue, commonly in the upper arm, thigh, or hip, where it is absorbed into the bloodstream.

\medterm{Intramuscular Injection} Delivery of a medicine into a muscle with a needle. The injection site and needle size depend on the medicine and the person’s age and body size; pain, bleeding, infection, or nerve injury are possible.

\medterm{Intramuscular Myxoma} An intramuscular myxoma is a rare benign tumor of gelatinous myxoid tissue located within skeletal muscle.

\medterm{Intraocular} Within the eyeball. The term may describe a medicine injection, lens, infection, bleeding, foreign body, pressure, or other finding inside the eye. Increased pressure inside the eye can damage the optic nerve and cause permanent loss of vision.

\medterm{Intraocular Infection} An infection inside the eye, such as endophthalmitis, that can rapidly threaten vision and requires urgent ophthalmic treatment.

\medterm{Intraocular Lens} An artificial lens placed inside the eye, usually during cataract surgery, to focus light after the cloudy natural lens is removed. Lens choices can emphasize distance, near vision, or astigmatism correction. Possible complications include glare, clouding behind the lens, infection, or retinal problems.

\medterm{Intraocular Lens Dislocation} Postoperative displacement of an implanted artificial intraocular lens from its intended
position within the capsular bag or ciliary sulcus into the vitreous cavity or anterior
chamber. It is most commonly caused by progressive zonular dehiscence and may require
surgical repositioning or exchange.

\medterm{Intraocular Lymphoma} Intraocular lymphoma is lymphoma involving the eye, often affecting the vitreous or retina and sometimes occurring with lymphoma in the brain or elsewhere.

\medterm{Intraocular Pressure} The pressure generated by aqueous humor within the eye, normally maintained by a balance
between aqueous production and outflow. Abnormal IOP is an important risk factor for
glaucomatous optic neuropathy, but glaucoma can occur at statistically normal pressures.

\medterm{Intraocular Tumor} Any neoplasm arising within the eye, with uveal melanoma being the most common primary
malignancy in adults and retinoblastoma the most common in children. Diagnosis requires
dilated fundus examination, ultrasound, and MRI; treatment depends on tumour type, size,
and location.

\medterm{Intraoperative} Occurring during a surgical operation, including monitoring, findings, complications, or treatment decisions made during the procedure.

\medterm{Intraoperative Blood Transfusion} Intraoperative blood transfusion replaces lost blood volume and maintains
oxygen-carrying capacity during major surgery. Indications include hemoglobin less than
7 g/dL in stable patients, less than 8 g/dL in cardiac disease, or active hemorrhage
with hemodynamic instability.

\medterm{Intraoperative Neuromonitoring} Electrical tests used during selected operations to watch the function of the brain, spinal cord, nerves, or muscles. Signals can warn the surgical team about a possible change in nerve function, allowing them to reassess the procedure and protect nearby structures.

\medterm{Intraoral} Located within the mouth or performed through the mouth. Located inside the mouth or performed through the mouth, such as an intraoral examination, scan, injection, or surgical procedure.
\medterm{Intraosseous Access} Placement of a vascular-access needle into the medullary cavity of a suitable bone to
permit rapid infusion of fluids, blood products, and medications when intravenous access
is delayed or impossible. It is a temporary emergency route and can cause extravasation,
infection, or compartment injury.

\medterm{Intraosseous Infusion} Rapid delivery of fluids, blood, or medicines through a needle placed into the inner space of a bone when a vein cannot be accessed quickly, usually during an emergency. It is temporary and must be monitored for pain, swelling, leakage, infection, or bone injury.

\medterm{Intrapartum} The period during labor and childbirth, from the onset of regular uterine contractions
leading to cervical dilation until delivery of the baby and placenta.

\medterm{Intrapartum Fever} Fever during labor, which may result from infection, epidural analgesia, dehydration, or other causes and requires maternal and fetal assessment.

\medterm{Intrapartum Stillbirth} Intrapartum stillbirth is fetal death occurring during labor after a live fetus was documented before labor, requiring obstetric investigation and support.

\medterm{Intrapericardial} Located inside the pericardial sac, the thin membrane surrounding the heart. An intrapericardial mass or fluid collection may affect how the heart fills or beats.

\medterm{Intraperitoneal} Intraperitoneal means within the peritoneal cavity or administered into it, including selected medicines, dialysis fluid, or procedures.

\medterm{Intraperitonealroute} The intraperitoneal route delivers a medicine or fluid into the peritoneal cavity, commonly for peritoneal dialysis or selected treatments.

\medterm{Intraprostatic} Located within the prostate gland, describing a lesion, injection, infection, biopsy, procedure, or other finding centered inside prostate tissue.

\medterm{Intrapulmonary} Located within lung tissue or occurring inside the lung’s air spaces, blood vessels, or other structures. Intrapulmonary findings may represent normal anatomy, infection, fluid, bleeding, inflammation, or tumors.

\medterm{Intrapulmonary Shunt} Perfusion of alveoli without ventilation (V/Q = 0). Blood passes through the pulmonary
circulation without gas exchange, mixing with oxygenated blood and causing hypoxemia. It
is seen in ARDS, pneumonia, and atelectasis. It is refractory to supplemental oxygen.

\medterm{Intrasynovial} Intrasynovial means within the synovial lining or fluid-filled space of a joint or tendon sheath.
\medterm{Intratendinous} Located within a tendon, describing a tear, calcium deposit, injection, degenerative change, or other abnormality affecting the tendon's substance.

\medterm{Intratesticular} Located within testicular tissue. An intratesticular mass may be benign or malignant and is usually evaluated with physical examination and scrotal ultrasound, interpreted in the clinical context.

\medterm{Intrathecal} Intrathecal means within the fluid-filled space surrounding the brain and spinal cord; medicines may be given there by lumbar puncture.

\medterm{Intrathecal Chemotherapy} Cancer medicine placed directly into the fluid surrounding the brain and spinal cord, usually through a lumbar puncture or implanted reservoir. It is used for selected leukemias, lymphomas, or cancer spread to this fluid; infection, bleeding, headache, and nervous-system toxicity are possible.

\medterm{Intrathoracic} Intrathoracic means within the chest cavity, including the lungs, heart, mediastinum, vessels, and thoracic lymph nodes.

\medterm{Intrathoracic Localized Obstruction} A focal blockage within the thoracic cavity affecting the airways, esophagus, or great
vessels, leading to localized symptoms depending on the structure involved. Causes
include endobronchial tumors, foreign bodies, esophageal strictures, or vascular
anomalies. Diagnosis is by imaging and endoscopy.

\medterm{Intratracheal} Intratracheal means within the trachea, such as a tube, medicine, foreign body, lesion, or route of administration.

\medterm{Intratympanic} Intratympanic means delivered through the eardrum into the middle ear, often for selected inner-ear disorders or local treatment.

\medterm{Intrauterine} Located or occurring within the uterus, including a pregnancy, device, medicine, lesion, or procedure. The term describes location and requires additional information to explain the finding’s purpose, cause, or risk.

\medterm{Intrauterine Contraceptive Device} A small device placed in the uterus to prevent pregnancy; copper IUDs release copper, while hormonal IUDs release a progestin, and both are highly effective reversible contraception.
\medterm{Intrauterine Device} A small device placed inside the uterus to prevent pregnancy for several years while remaining reversible. Copper devices work without hormones; hormonal devices release a progestogen and often reduce menstrual bleeding. Rare risks include expulsion, perforation during insertion, and infection soon after placement.

\medterm{Intrauterine System} A small T-shaped contraceptive device placed inside the uterus that slowly releases a progestogen hormone to prevent pregnancy. It is also called IUS.

\medterm{Intrauterine Devices} Small devices placed in the uterus to prevent pregnancy; copper and hormonal IUDs are highly effective long-acting reversible contraceptives.

\medterm{Intrauterine Fetal Demise / Stillbirth} Death of a fetus before complete delivery. The gestational-age
definition varies by jurisdiction. Causes may include placental disease, umbilical-cord complications, infection,
genetic abnormalities, maternal disease, or unexplained factors.

\medterm{Intrauterine Growth Restriction} IUGR: a fetus with estimated weight below the 10th percentile (or growth velocity
decline), usually with placental insufficiency or maternal/fetal factors.

\medterm{Intrauterine Pressure Catheter} A fluid-filled catheter inserted into the uterine cavity alongside the fetus to measure
the intensity and frequency of contractions in mm Hg (Montevideo units, MVU).

\medterm{Intrauterine Transfusion} Ultrasonically guided intravascular or intraperitoneal transfusion of fetal red blood
cells to treat severe fetal anemia, most commonly from Rh isoimmunization. Performed
through the umbilical vein using a 22-25 gauge needle. Packed red blood cells,
CMV-negative and crossmatched, are transfused slowly.

\medterm{Intravascular} Intravascular means within a blood vessel, including a catheter, clot, injection, or process occurring in the circulation.

\medterm{Intravascular Catheter Infections} Infections caused by microorganisms entering or growing on a catheter placed inside a blood vessel. They may produce fever, chills, low blood pressure, or bloodstream infection. Diagnosis uses examination and blood cultures, and treatment may require antibiotics and catheter removal.

\medterm{Intravascular Filter} A device placed in a large vein to trap venous clots and reduce pulmonary embolism risk in selected patients who cannot receive
or have not responded adequately to anticoagulation.

\medterm{Intravascular Hemolysis} Destruction of red blood cells within the bloodstream, releasing hemoglobin into plasma; causes include immune hemolysis, mechanical trauma, infections, toxins, and transfusion reactions.

\medterm{Intravascular Lymphoma} A rare, fast-growing blood cancer in which abnormal white blood cells collect inside small blood vessels. It may affect the skin, brain, or other organs, causing fever, skin changes, or neurologic problems.

\medterm{Intravascular Ultrasound} Catheter-based imaging using a miniature ultrasound probe inside coronary arteries to
visualize vessel wall layers, plaque composition, and stent deployment. Provides more
detailed information than angiography alone, including vessel remodeling, plaque burden,
and true vessel diameter.

\medterm{Intravenous} Into or within a vein; intravenous medicines, fluids, or blood products are delivered directly into the bloodstream.

\medterm{Intravenous Anesthesia} Anesthetic induction or maintenance using drugs delivered into a vein. It permits rapid
titration but requires attention to dose, context-sensitive recovery, cardiovascular and
respiratory effects, and the availability of a secure airway.

\medterm{Intravenous Catheters} Flexible devices inserted into a vein to deliver fluids, medicines, blood products, nutrition, or contrast, or to obtain blood samples; complications include infection, thrombosis, infiltration, and extravasation.

\medterm{Intravenous Immunoglobulin} A purified preparation of pooled antibodies given through a vein (IVIg) for selected immune deficiencies, autoimmune diseases, inflammatory disorders, and some neurologic conditions.

\medterm{Intravenous Leiomyomatosis} Extension of benign uterine smooth-muscle tumor into venous channels, sometimes reaching the inferior vena cava or right heart. Imaging is important because extensive disease can cause serious cardiovascular complications.

\medterm{Intravenous Lipid Emulsion Therapy} Administration of a lipid emulsion as a rescue treatment for selected severe lipophilic
drug toxicities, especially local-anesthetic systemic toxicity. It may act as a lipid
sink and through metabolic and cardiovascular effects and requires critical-care
monitoring.

\medterm{Intravenous Pyelogram} An x-ray examination of the urinary tract after intravenous iodinated contrast, historically used to visualize the kidneys, ureters, and bladder; CT urography or ultrasound often replaces it.
of iodinated contrast material. The contrast is filtered by the glomeruli and excreted
through the collecting system, allowing sequential imaging of the renal parenchyma,
calyces, renal pelvis, ureters, and bladder.

\medterm{Intravenous Pyelography} An X-ray examination of the urinary tract after intravenous contrast administration; it is now often replaced by CT urography or ultrasound.

\medterm{Intravenous Sedation} Administration of sedative medicines through a vein to reduce anxiety, discomfort, or awareness during a procedure while maintaining appropriate monitoring and readiness to support breathing.

\medterm{Intraventricular Hemorrhage} Bleeding into the germinal matrix and ventricular system, primarily in premature infants
less than 32 weeks gestation. Graded I-IV by severity on cranial ultrasound. Grade I
(germinal matrix only) and Grade II (intraventricular without dilation) have good
prognosis.

\medterm{Intraventricular Hemorrhage Of The Newborn} Intraventricular hemorrhage in a newborn is bleeding into the brain’s ventricular system, most often in premature infants, and can cause hydrocephalus or neurologic injury.

\medterm{Intravital Microscopy} Intravital microscopy uses a microscope to observe cells or tissues in a living organism, usually in research.

\medterm{Intravitreal Injection} An ophthalmological procedure in which medication is injected directly into the vitreous
cavity of the eye through the pars plana, a safe entry zone located approximately 3.5 to
4 millimeters posterior to the limbus, to deliver therapeutic agents directly to the
retina and vitreous for the treatment of various retinal diseases.

\medterm{Intrinsic} Belonging naturally to a structure or process and arising from within it, rather than being externally imposed.

\medterm{Intrinsic Factor} A protein secreted by stomach parietal cells that is required for vitamin B12 absorption in the terminal ileum.

\medterm{Intrinsic Renal AKI} Acute kidney injury caused by damage to renal tissue. Causes include acute tubular injury,
acute interstitial nephritis, glomerulonephritis, and vascular disorders. Laboratory findings vary with the cause and
should not be interpreted in isolation.

\medterm{Intrinsic Sphincter Deficiency} Weakness of the urinary sphincter that causes stress urinary incontinence.

\synonyms
ISD

\medterm{Introitus} An entrance or opening, especially the entrance to the vagina.

\medterm{Intron} A non-coding segment of a gene that is transcribed into RNA but spliced out before the
messenger RNA is translated into a protein.

\medterm{Intrusive Thought} An unwanted, distressing thought, image, or impulse that enters awareness repeatedly; it does not necessarily reflect intent or desire.

\medterm{Intubate} To place a tube into a body passage, especially an endotracheal tube through the mouth or nose into the trachea to maintain an airway or support ventilation.

\medterm{Intubation} Placement of a tube into the airway or another body passage to maintain it or deliver treatment.

\medterm{Intussusception} Condition in which a segment of intestine telescopes into the adjacent distal segment
causing bowel obstruction and compromise of the mesenteric blood supply. It is the most
common cause of intestinal obstruction in infants and children between six months and
three years of age.

\medterm{Intussusception – Children} Intussusception occurs when one segment of intestine telescopes into another, causing intermittent pain, vomiting, obstruction, and possible bowel ischemia in children.

\medterm{Inulin} Inulin is a plant fructan used as a dietary fiber and in tests of glomerular filtration because it is freely filtered and minimally handled by the kidney.

\medterm{Inuresis} A rare or obsolete spelling related to enuresis, meaning involuntary urination, especially during sleep.

\medterm{Inutero} A variant or misspelling of in utero, meaning within the uterus before birth.

\medterm{Invasive} Able to enter and spread through tissues; the term may describe an infection, tumor, procedure, or disease that penetrates nearby structures.

\medterm{Invasive Candidiasis} A serious Candida infection that enters the bloodstream or deep organs, most often in people who
are critically ill or immunocompromised. It may cause persistent fever or organ dysfunction and
requires prompt medical diagnosis and treatment.

\medterm{Invasive Ductal Carcinoma} The most common type of breast cancer, characterized by malignant cells that have broken
through the duct wall and invaded surrounding breast tissue. It can metastasize via
lymphatic or hematogenous spread.

\medterm{Invasive Lobular Carcinoma} A breast cancer characterized by discobesive cells infiltrating in single-file lines.
More frequently bilateral and multifocal than ductal carcinoma and harder to detect on
imaging; usually hormone receptor positive.

\medterm{Invasive Mole} A hydatidiform mole that invades the uterine myometrium or vascular spaces; it is a form of gestational trophoblastic neoplasia and may persist or metastasize.
\medterm{Invasive Procedure} A medical procedure that enters the body through a needle, tube, incision, or other instrument. It may diagnose or treat disease and can cause pain, bleeding, infection, injury, or complications related to anesthesia.

\medterm{Inverse Agonist} A type of drug that binds to a cell receptor and actively produces the opposite
biological effect of an agonist, effectively turning down the receptor's baseline
activity.

\medterm{Inversion} Chromosomal segment reversed end-to-end. Paracentric: does not include centromere.
Pericentric: includes centromere. Usually benign unless breakpoint disrupts gene.

\medterm{Invert} To turn upside down, reverse, or turn inward; in anatomy, inversion can describe movement of the sole of the foot toward the midline.

\medterm{Invert Sugar} A mixture of glucose and fructose produced by splitting sucrose, used as a sweetener in foods and medicines. It has nutritional effects similar to other added sugars and contributes to total carbohydrate and energy intake.

\medterm{Invertase} An enzyme that splits sucrose, table sugar, into glucose and fructose. It is produced by some plants, fungi, and microorganisms and is also used in food processing.

\medterm{Invertebrata} Invertebrata is a broad nonformal grouping of animals without a vertebral column, including worms, mollusks, arthropods, and many other groups.

\medterm{Invertebrate} An animal without a backbone, including insects, worms, mollusks, spiders, and many other groups with diverse body structures and life cycles.

\medterm{Invertebrate Hormones} Chemical signals regulating growth, development, reproduction, molting, behavior, or metabolism in animals without backbones. Studying them helps explain animal biology and can support development of selective pest-control methods.

\medterm{Inverted Papilloma} A usually benign but locally destructive tumor arising most often in the nasal cavity or sinuses; it can recur and rarely become cancerous.

\medterm{Investigational} Being studied in clinical research and not yet established as safe and effective for a particular use. Investigational medicines and devices require appropriate oversight, informed consent, and clear communication of uncertainty.

\medterm{Inveterate} Long-standing, deeply established, or difficult to change, often describing a habit, condition, or behavior. The word indicates persistence but does not identify the cause, severity, treatment, or expected outcome.

\medterm{Involucrum} An involucrum is a sheath of new bone that forms around dead bone, classically in chronic osteomyelitis.

\medterm{Involuntary} Occurring without conscious control or deliberate choice, such as a reflex, autonomic function, or involuntary muscle movement.

\medterm{Involuntary Movement} A movement that occurs without conscious control, such as a tremor, tic, chorea, dystonia, or myoclonus.

\medterm{Involuntary Muscle} Muscle that contracts without conscious control, including smooth muscle in internal organs and cardiac muscle in the heart.

\medterm{Involuntary Psychiatric Treatment} Assessment, hospitalization, or treatment without the patient’s consent under applicable
law when serious mental disorder creates a substantial risk of harm or severe inability
to meet basic needs. The least restrictive alternative, due process, and ongoing
reassessment are required.

\medterm{Involute} To shrink, return toward a former size, or undergo involution; in anatomy, an involute structure may be tightly coiled or rolled inward.

\medterm{Involution} The process by which the uterus returns to its prepregnancy size after delivery. The
uterus weighs approximately 1000 grams immediately postpartum and returns to
approximately 60 grams by 6 weeks. Fundal height decreases approximately 1 cm per day.

\medterm{Iodamoeba Butschlii} A usually non-disease-causing intestinal amoeba spread by swallowing cysts in contaminated food or water. Finding it in stool suggests fecal exposure but usually does not require treatment unless another infection is present.

\medterm{Iodide} The negatively charged form of iodine, used by the thyroid to make hormones and present in medicines, supplements, and contrast compounds.

\medterm{Iodide-Associated Sialadenitis} Inflammation and swelling of the salivary glands caused by excessive iodide intake,
typically from medications such as amiodarone or iodine-containing contrast, resolving
with dose reduction.

\medterm{Iodides} Salts containing iodine, sometimes used in carefully selected thyroid disorders.

\medterm{Iodinated Contrast Media} Contrast agents containing iodine used primarily for radiography, CT, angiography, and
some fluoroscopic procedures. Important risks include immediate hypersensitivity-like
reactions, extravasation, thyroid effects, and kidney injury in susceptible patients.

\medterm{Iodine} An essential mineral used to make thyroid hormones, which regulate growth, brain development, and metabolism. Too little iodine can cause thyroid enlargement and hormone deficiency, especially harmful during pregnancy and childhood; too much can also disturb thyroid function. Main sources include iodized salt, seafood, dairy, and some supplements.

\medterm{Iodine Deficiency} Inadequate iodine intake or availability, which can impair thyroid-hormone production and cause goiter, hypothyroidism, or developmental problems in a fetus or child.


\medterm{Iodine Poisoning} Iodine poisoning follows excessive iodine exposure and can irritate the gastrointestinal tract, disrupt thyroid function, and cause systemic toxicity depending on dose and formulation.

\medterm{Iodine-123} Pure gamma emitter, half-life 13 hours. Used for thyroid imaging. No beta emission,
lower radiation dose than I-131.

\medterm{Iodine-131} A radioactive iodine isotope used in selected diagnostic and therapeutic procedures, especially for thyroid disease.
\medterm{Iodipamide} An iodinated contrast medium formerly used for imaging the biliary system and urinary tract; it can cause allergic or kidney-related reactions.

\medterm{Iodoacetamide} A chemical reagent that irreversibly modifies sulfhydryl groups, especially cysteine residues, and is used in biochemical research.

\medterm{Iododerma} A papulopustular or acneiform skin eruption caused by iodine excess, presenting on the
face and extensor surfaces, resolving with iodine restriction and symptomatic treatment.

\medterm{Iodoform} A yellow iodine-containing compound formerly used as an antiseptic and in some wound or dental preparations. It can irritate tissue and cause toxicity if absorbed or swallowed and is now used only in selected contexts.

\medterm{Iohexol} A water-soluble iodine-containing contrast medicine used during CT scans, angiography, and other imaging studies to make vessels and organs easier to see.

\medterm{Ion} An atom or molecule carrying a positive or negative electric charge. Physiologically important ions include sodium, potassium, calcium, chloride, and bicarbonate, whose concentrations influence nerves, muscles, fluid balance, and acid-base status.

\medterm{Ion Transport} Movement of charged particles across a cell membrane through channels, carriers, or pumps. It maintains membrane potentials, nerve signaling, muscle contraction, fluid balance, and epithelial transport.

\medterm{Ionic Strength} A measure of the total concentration and charge of ions in a solution. It affects protein interactions, membrane behavior, chemical equilibria, and the performance of some laboratory assays and medicines.

\medterm{Ionization} The formation of charged particles by gaining or losing electrons or protons. A medicine's degree of ionization affects its solubility, membrane passage, absorption, distribution, and urinary excretion.

\medterm{Ionizing Radiation} Radiation with enough energy to remove electrons from atoms and create ions, including
X-rays, gamma rays, and certain particulate emissions. It can cause deterministic tissue
effects at sufficient doses and stochastic cancer risk at lower doses.

\medterm{Ionomycin} A laboratory compound that carries calcium across cell membranes and raises calcium inside cells. It is used mainly in research to study cell signaling, secretion, and immune-cell activation, not as a routine medicine.

\medterm{Ionophore} A molecule that transports ions across a lipid membrane. Ionophores have research, antimicrobial, and veterinary uses; some can disrupt human cellular ion balance and are toxic.

\medterm{Ions} Electrically charged atoms or molecules, including the electrolytes that regulate nerve conduction, muscle function, hydration, and acid-base balance. Abnormal ion concentrations can cause weakness, seizures, arrhythmias, or altered consciousness.

\medterm{Iontophoresis} Use of a mild electrical current to drive an ionized medicine through the skin or alter sweat-gland activity. It is used in selected dermatologic and rehabilitation settings; skin irritation and burns can occur.

\medterm{Iopamidol} A nonionic iodine-containing contrast medicine used to improve CT, angiography, and other medical images. It can cause warmth, nausea, allergy, kidney problems, or thyroid effects, so risks require assessment beforehand.

\medterm{Iothalamate Clearance} Iothalamate clearance estimates glomerular filtration by measuring how quickly injected iothalamate is cleared from blood or urine.

\medterm{Iothalamate Meglumine} Iothalamate meglumine is an iodinated contrast agent formerly used in radiographic imaging and associated with dose- and kidney-related risks.

\medterm{Iothalamic Acid} An older water-soluble, iodinated contrast agent formerly used for radiographic imaging. It has largely been replaced by newer agents; prior contrast reactions, kidney function, and thyroid disease remain clinically relevant when reviewing exposure history.

\medterm{Ioxaglic Acid} An older iodinated contrast medium used for radiographic imaging, including selected angiographic and urinary-tract studies. It has largely been superseded by newer agents, but its history is relevant when assessing prior contrast reactions.

\medterm{IPC} Alternate name; see \textit{Indwelling Pleural Catheter}.

\medterm{Ipecac} A vomiting medicine formerly used after poisoning, now not routinely recommended because vomiting can cause aspiration and delay more appropriate emergency treatment.

\medterm{Ipilimumab} An immunotherapy antibody that blocks CTLA-4, removing a brake on T cells so they can attack cancer. It can cause immune inflammation in the bowel, skin, liver, lungs, hormone glands, or other organs.

\medterm{IPMN Pancreas (Pancreatic Cysts)} Alternate index label for Pancreatic Cysts.

\medterm{Ipratropium} Ipratropium is an inhaled short-acting muscarinic antagonist that reduces vagally mediated bronchoconstriction and nasal secretion; it is used in COPD and acute asthma with bronchodilators.

\medterm{Ipratropium Bromide} An inhaled anticholinergic (antimuscarinic) bronchodilator that blocks acetylcholine at
airway smooth muscle muscarinic receptors. It reduces bronchospasm and mucus secretion.
It is used in COPD and acute asthma exacerbations. It is less effective than beta-2
agonists in asthma but more effective in COPD.

\medterm{Iproniazid} Iproniazid is an older monoamine oxidase inhibitor antidepressant that is no longer commonly used because of serious liver toxicity and drug interactions.

\medterm{Ipsilateral} On or affecting the same side of the body as a stated structure or injury. For example, weakness in the right arm and right leg is ipsilateral to a finding on the right side of the body; the opposite term is contralateral.

\medterm{IQ Testing} IQ testing uses standardized tasks to estimate intellectual functioning; results should be interpreted with education, language, development, and clinical context.

\medterm{Iratumumab} An investigational antibody studied for cancers involving CD30-positive cells. It was designed to help the immune system recognize these cells, but it is not an established routine treatment.

\medterm{Irbesartan} Irbesartan is an angiotensin-II receptor blocker used for hypertension and diabetic kidney disease; it can cause hyperkalaemia, hypotension, and renal-function changes and is contraindicated in pregnancy.

\medterm{Iridium} Iridium is a dense metal whose radioactive isotopes, especially iridium-192, are used in industrial testing and some medical brachytherapy.

\medterm{Iridocyclitis} Inflammation involving the iris and ciliary body, a form of anterior uveitis. It may
cause eye pain, photophobia, redness, blurred vision, and cells or flare in the anterior
chamber and may be associated with autoimmune disease, infection, trauma, or other
causes.

\medterm{Iridotomy} A procedure that creates a small opening in the iris to allow aqueous humor to pass between the posterior and anterior
chambers. Laser peripheral iridotomy is used to treat or prevent selected forms of
angle-closure glaucoma.

\medterm{Irinotecan} A chemotherapy medicine that inhibits topoisomerase I and is used for colorectal and several other cancers.

\medterm{Irinotecan Hydrochloride} A chemotherapy medicine that blocks topoisomerase I and damages cancer-cell DNA, used for selected cancers with diarrhea and low blood counts possible.

\medterm{Iris} The colored, muscular ring between the cornea and lens that controls the size of the pupil and therefore the amount of light entering the eye. Its muscles constrict the pupil in bright light and enlarge it in dim light.

\medterm{Iris Atrophy} Thinning or loss of iris tissue resulting from trauma, chronic inflammation, ischaemia,
or neurodegenerative disease, leading to a translucent or patched iris appearance. It
may be sectoral or diffuse and can be associated with decreased pupillary response,
photophobia, or secondary glaucoma.

\medterm{Iris Cyst} An iris cyst is a fluid-filled lesion on or within the iris, usually benign but sometimes affecting vision, pressure, or the appearance of the eye.

\medterm{Iris Disease} A disorder of the colored part of the eye, including inflammation, injury, infection, developmental abnormalities, or tumors.

\medterm{Iris Melanoma} Iris melanoma is a malignant tumor of pigment cells in the iris and may appear as a growing pigmented lesion, pupil distortion, or glaucoma.

\medterm{Iris Neoplasm} An iris neoplasm is an abnormal growth arising in the colored part of the eye; it may be benign or malignant and can affect vision or eye pressure.

\medterm{Iritis} Inflammation of the iris, the most common form of anterior uveitis. Presents with acute
onset eye pain, photophobia, blurred vision, and redness (ciliary flush). Signs: cells
and flare in the anterior chamber, keratic precipitates (KP), synechiae (adhesions
between iris and lens).

\medterm{Irofulven} Irofulven is an investigational anticancer compound derived from illudin that damages DNA and has been studied in solid tumors.

\medterm{Iron} Essential for hemoglobin, myoglobin, and electron transport. Heme iron from animal foods
is generally more bioavailable than non-heme iron from plant foods. Deficiency is the
most common nutritional deficiency worldwide and may result from blood loss, increased
requirements, poor intake, or malabsorption.

\medterm{Iron Deficiency} The most common nutritional deficiency worldwide, resulting from inadequate iron intake,
impaired absorption, or chronic blood loss, leading to reduced hemoglobin synthesis and
microcytic hypochromic anemia. Clinical features include fatigue, pallor, koilonychia,
pica, and angular cheilitis.

\medterm{Iron Deficiency Anemia} Anemia caused by insufficient iron for hemoglobin production. It is often microcytic and hypochromic; treatment requires identifying the cause and replacing iron when appropriate.


\medterm{Iron Overdose} Iron overdose can cause gastrointestinal injury, metabolic acidosis, shock, liver failure, and multiorgan toxicity, particularly after a large ingestion by a child.

\medterm{Iron Overload} Excess iron stored in the body, which can damage the liver, heart, pancreas, joints, or endocrine organs.

\synonyms
Overload Iron (Iron Overload)

\medterm{Iron-Deficiency Anaemia} Anaemia caused by insufficient iron for haemoglobin production, commonly due to blood loss, inadequate intake, or impaired absorption.

\medterm{Iron-Deficiency Anaemia In Adolescents} Iron-deficiency anaemia during adolescence, commonly related to rapid growth, inadequate dietary iron, or menstrual blood loss. Assessment should identify the cause and guide dietary advice and clinical treatment.

\medterm{Iron-Deficiency Anemia} Anemia caused by insufficient iron for normal hemoglobin production. It may cause fatigue, weakness,
shortness of breath, paleness, headache, or unusual cravings such as ice; the cause of iron
loss or inadequate intake should be identified.

\medterm{Irradiated Foods} Irradiated foods are exposed to controlled ionizing radiation to reduce pathogens, insects, or spoilage; the process does not make food radioactive.

\medterm{Irradiation} Exposure to or treatment with radiation; irradiation may be used diagnostically or therapeutically, or may cause tissue injury depending on dose and type.

\medterm{Irregular} Not following the usual pattern, rhythm, shape, timing, or order. In medicine, it may describe an uneven heartbeat, an unpredictable menstrual cycle, an unusual skin edge, or a result that does not match the expected pattern.

\medterm{Irregular Menstruation} Menstrual bleeding that is unusually early, late, frequent, infrequent, prolonged, heavy, or unpredictable for the person’s age and reproductive context. Causes include pregnancy, ovulatory disorders, endocrine disease, medicines, structural lesions, and perimenopause.

\medterm{Irregular Respiration} An abnormal breathing pattern in rate, depth, rhythm, or pauses. It may occur with sleep disorders, metabolic disturbance, lung disease, drug effects, or neurologic injury and may require urgent assessment if accompanied by cyanosis or reduced consciousness.

\medterm{Irregular Sleep-Wake Syndrome} A circadian rhythm sleep-wake disorder with multiple irregular periods of sleep and
wakefulness across the 24-hour day rather than one consolidated nighttime sleep period.
It is associated with neurological disease, dementia, aging, and reduced exposure to
environmental time cues.

\medterm{Irreversible Cell Injury} Cell damage beyond the point of recovery, resulting in death by necrosis, apoptosis, or
another regulated cell-death pathway. Key features include inability to reverse
mitochondrial dysfunction and severe membrane or nuclear damage.

\medterm{Irrigate} To wash or flush a body part, wound, cavity, or medical device with a fluid to remove debris, contaminants, or secretions.

\medterm{Irrigation} Flushing a wound, body cavity, or organ with fluid to remove blood, debris, microorganisms, or unwanted material.

\medterm{Irrigation Solution} A sterile or appropriately prepared liquid used to wash a wound, body cavity, organ, or surgical field. The solution must be suitable for the tissue because contamination, pressure injury, electrolyte shifts, or chemical irritation can cause harm.

\medterm{Irritable Bowel Syndrome} A disorder of gut--brain interaction that causes repeated abdominal pain and changes in bowel habits, such as diarrhea, constipation, or both. No visible structural damage is usually found in the digestive tract.

\synonyms
IBS

\medterm{Irritable Bowel Syndrome Diet} An individualized eating plan used to reduce irritable bowel syndrome symptoms. A limited
trial of a low-FODMAP diet may help some people and should be followed by structured
food reintroduction.

\medterm{Irritable Mood} A mood marked by increased sensitivity, impatience, anger, or a tendency to become upset easily.
\medterm{Irritant Contact Dermatitis} Inflammation of the skin caused by direct chemical or physical injury rather than an immune-mediated allergy. It commonly produces redness, burning, scaling, fissures, or pain at the area of exposure.

\medterm{Irritants} Substances that inflame or damage tissues and cause symptoms without necessarily triggering a true immune allergy. Examples include smoke, chemicals, soaps, fumes, and pollutants affecting skin, eyes, airways, or digestion.

\medterm{Irritation} A local unpleasant or inflammatory response to mechanical, chemical, thermal, infectious, or allergic stimuli. It may cause redness, itching, burning, pain, cough, or tearing depending on the tissue involved.

\medterm{Isaacs Syndrome} A rare peripheral nerve hyperexcitability disorder causing continuous muscle-fibre
activity, stiffness, cramps, fasciculations, myokymia, and delayed muscle relaxation. It
may be autoimmune, paraneoplastic, or associated with inherited or acquired
peripheral-nerve disorders.

\medterm{Isavuconazole} A triazole antifungal agent used in the treatment of invasive aspergillosis and
mucormycosis. Isavuconazole inhibits fungal lanosterol 14-alpha demethylase, disrupting
ergosterol synthesis and fungal cell membrane integrity.

\medterm{Isazophos} Isazophos is an organophosphate pesticide that can inhibit acetylcholinesterase and cause cholinergic toxicity after exposure.
\medterm{Ischaemia} Inadequate blood flow and oxygen delivery to a tissue, usually because of narrowed or blocked blood vessels, causing dysfunction or injury.

\medterm{Ischaemic Heart Disease} Heart disease caused by reduced blood supply to the heart muscle, usually as a result of coronary artery disease; it may cause angina or myocardial infarction.

\medterm{Ischemia} Reduced blood flow causing too little oxygen for a tissue’s needs. It can cause pain, impaired function, permanent injury, or tissue death and may affect the heart, brain, limbs, or bowel.

\medterm{Ischemia, Intestinal} Alternate index label; see \textit{Intestinal ischemia}.

\medterm{Ischemia-Reperfusion Injury} The paradoxical tissue damage that occurs when blood flow is restored to previously
ischemic tissue, a phenomenon with major clinical significance in myocardial infarction
treatment, stroke management, organ transplantation, and shock resuscitation.

\medterm{Ischemic Bone Necrosis} Alternate index label; see \textit{Avascular necrosis (osteonecrosis)}.

\medterm{Ischemic Cardiomyopathy} Weakening of the heart muscle caused by coronary artery disease and previous or ongoing
myocardial ischemia. Scarred or poorly perfused myocardium can reduce contractile function and increase the risk of heart
failure and ventricular arrhythmias.

\medterm{Ischemic Chest Pain (Angina Symptoms)} Alternate index label for Angina Symptoms.

\medterm{Ischemic Colitis} Inflammation and injury of the colon caused by reduced blood flow. It often causes sudden abdominal pain, bloody diarrhea, or rectal bleeding and commonly affects areas such as the splenic flexure.

\synonyms
Colitis, Ischemic

\medterm{Ischemic Fasciitis} A benign fibroblastic and myofibroblastic proliferation associated with pressure-related ischemic injury, usually arising in deep soft tissue over bony prominences in debilitated or immobile patients.

\medterm{Ischemic Gangrene} Death of tissue due to prolonged severe arterial insufficiency, most commonly in the
extremities. Presents as dry, black, mummified tissue with a clear line of demarcation
from viable tissue. Requires revascularization when possible and surgical debridement or
amputation of necrotic tissue.

\medterm{Ischemic Heart Disease} Heart disease caused by inadequate blood flow through narrowed or blocked coronary arteries. It may cause chest pressure, breathlessness, heart attack, rhythm problems, heart failure, or sudden death.

\synonyms
Coronary artery disease; coronary heart disease

\medterm{Ischemic Preconditioning} A phenomenon in which brief, nonlethal episodes of ischemia make tissue more resistant to a later, prolonged ischemic insult. It is important in cardiovascular and research physiology, but clinical applications remain context-dependent.

\medterm{Ischemic Priapism} Alternate index label; see \textit{Priapism}.

\medterm{Ischemic Stroke} A neurologic deficit caused by inadequate blood flow to part of the brain, usually from
thrombotic or embolic arterial occlusion. Causes include atherosclerotic disease, cardiac
embolism, and small-vessel disease.


\medterm{Ischium} Inferior-posterior hip bone. Ischial tuberosity bears seated weight and serves as
hamstring attachment. Ischial spine is a pudendal nerve block landmark and determines
pelvic outlet level. Lesser sciatic notch lies between spine and tuberosity.

\medterm{Ishak Score} A pathology score describing the amount of liver inflammation and scarring seen in a biopsy, especially in chronic hepatitis. Higher scores indicate more advanced damage, but interpretation requires the full clinical picture.

\medterm{Islet Cell} A cell in a pancreatic islet that helps control blood glucose; beta cells make insulin and alpha cells make glucagon.

\medterm{Islet Cell Cancer} Alternate index label; see \textit{Pancreatic neuroendocrine tumors}.

\medterm{Islets Of Langerhans} Small clusters of hormone-producing cells in the pancreas that release insulin, glucagon, and other hormones regulating blood sugar and metabolism.

\medterm{Iso-Sulfan Blue} A blue dye used to map lymphatic drainage or identify sentinel lymph nodes during selected surgical procedures.

\medterm{Isoaminile} Isoaminile is an older sympathomimetic appetite suppressant that is no longer widely used because of cardiovascular and other adverse effects.

\medterm{Isoantibodies} Antibodies made against antigens from another individual of the same species, such as antibodies involved in blood-group incompatibility.

\medterm{Isoantigen} An antigen that differs between individuals of the same species and can trigger an immune response after exposure.

\medterm{Isocarboxazid} Isocarboxazid is a monoamine oxidase inhibitor used for depression; it has important interactions with other medicines and tyramine-rich foods.

\medterm{Isochromosomes} Abnormal chromosomes with two identical arms caused by division across the wrong axis of the centromere.

\medterm{Isocitrate Dehydrogenase} A regulated enzyme of the citric acid cycle that converts isocitrate to
alpha-ketoglutarate while producing NADH and releasing carbon dioxide. Mitochondrial
isoforms are stimulated by ADP and calcium and inhibited by ATP and NADH. Mutations in
IDH1 or IDH2 can produce the oncometabolite 2-hydroxyglutarate.

\medterm{Isocyanate} A reactive chemical group found in some industrial products; exposure can irritate the airways and cause occupational asthma or skin reactions.

\medterm{Isoelectric Point} The acidity level at which a molecule, especially a protein, carries no overall electrical charge and may behave differently in solution.

\medterm{Isoenzyme} One of several forms of an enzyme that catalyze the same chemical reaction but differ in structure, location, or regulation. Measuring selected isoenzymes can help identify the tissue affected by injury.

\medterm{Isoetharine} An older airway-opening medicine related to beta agonists, formerly used to relieve narrowed airways but largely replaced by newer treatments.

\medterm{Isofenphos} A poisonous organophosphate insecticide. Exposure can cause sweating, excess saliva, vomiting, diarrhea, small pupils, muscle twitching, breathing difficulty, seizures, or weakness because it disrupts nerve signals; suspected poisoning needs urgent medical help.

\medterm{Isoflavone} A plant compound, especially from soy, with weak estrogen-like activity that is studied in relation to bone, menopause, and cancer risk.

\medterm{Isoflurane} Isoflurane is an inhaled anesthetic used to maintain general anesthesia; it can lower blood pressure and depress breathing.

\medterm{Isoflurophate} Isoflurophate is a long-acting organophosphate cholinesterase inhibitor formerly used in eye treatments; it can cause prolonged cholinergic effects.

\medterm{Isogeneic Graft} A transplant between genetically identical individuals, such as identical twins, so immune rejection risk is relatively low.

\medterm{Isokinetic Exercise} Muscle contraction at constant speed throughout range of motion. Uses specialized
equipment. Assesses and improves strength.

\medterm{Isolate} To separate a person, organism, tissue, or substance from others; infection-control isolation separates people who may transmit a pathogen.

\medterm{Isolate Compound} To separate a single chemical compound from a mixture for identification, testing, or research. In clinical laboratories, isolation and characterization may support toxicology, pharmacology, or natural-product analysis.

\medterm{Isolated Systolic Hypertension} Elevated systolic blood pressure with normal diastolic pressure, commonly related to arterial stiffness in older adults and associated with increased cardiovascular risk.

\medterm{Isolation} Separating a person with a contagious infection from others to limit spread, using appropriate rooms, protective equipment, and infection-control procedures.

\medterm{Isolation (Social)} Lack of social contacts and interactions. Common in elderly: bereavement, mobility
limitation, sensory impairment, retirement. Consequences: depression, cognitive decline,
functional impairment, mortality. Interventions: social programs, transportation,
technology, volunteer visits.

\synonyms
Social

\medterm{Isolation Precautions} Measures used with standard precautions to reduce transmission of infection. Additional precautions may be contact, droplet, or airborne, depending on the mode of spread.

\medterm{Isoleucine} An essential amino acid obtained from food that helps build proteins, support energy production, maintain muscle, and repair body tissues.

\medterm{Isomerase} An enzyme that rearranges atoms within a molecule to form an isomer. Isomerases participate in metabolism, and inherited defects in related pathways can cause metabolic disease.

\medterm{Isometheptene} Isometheptene is a vasoconstrictor formerly used in some headache medicines; it may raise blood pressure and is unsuitable for some cardiovascular conditions.

\medterm{Isometric} Describing muscle contraction that produces force without a meaningful change in muscle length or joint movement.

\medterm{Isometropia} The condition in which both eyes have approximately equal refractive power, so neither eye has substantially more near-sightedness, far-sightedness, or astigmatism than the other.

\medterm{Isoniazid} Isoniazid is a first-line antituberculous drug that inhibits mycolic-acid synthesis in Mycobacterium tuberculosis; hepatotoxicity and peripheral neuropathy are important risks, and pyridoxine is often given.

\medterm{Isoniazid Toxicity} Harm from excessive or unusual sensitivity to isoniazid, especially liver injury, nerve damage, or seizures. Treatment may include vitamin B6, supportive care, activated charcoal, seizure control, and urgent monitoring.

\medterm{Isoprenaline} Isoprenaline, also called isoproterenol, is a beta-adrenergic medicine that can increase heart rate and is used in selected emergency cardiac situations.

\medterm{Isoprenylation} Addition of a lipid-derived isoprenyl group to a protein, often helping determine its membrane location and signaling activity. Abnormal isoprenylation contributes to some cancers and is a target of investigational therapies.

\medterm{Isoprocarb} Isoprocarb is a carbamate insecticide that can cause cholinergic poisoning if swallowed or absorbed in a significant amount.

\medterm{Isopropamide} Isopropamide is an anticholinergic medicine formerly used for gastrointestinal spasms; it may cause dry mouth, blurred vision, constipation, or urinary retention.

\medterm{Isopropanol Alcohol Poisoning} Poisoning from isopropanol, often found in rubbing alcohol, causing intoxication, vomiting, abdominal pain, low blood pressure, and sometimes coma.

\medterm{Isopropyl Alcohol} A volatile alcohol used as a disinfectant and solvent. Ingestion can cause rapid intoxication, vomiting, abdominal pain, hypotension, and coma; management is supportive and differs from treatment of methanol or ethylene-glycol poisoning.

\medterm{Isoprostane} A compound formed when free radicals oxidize membrane lipids; it is used as a marker of oxidative stress.

\medterm{Isoproterenol} Isoproterenol, also called isoprenaline, is a beta-adrenergic medicine that increases heart rate and contractility and is used in selected emergency settings.

\medterm{Isoquinolines} A group of nitrogen-containing ring compounds that includes chemicals studied in pharmacology, natural products, and organic chemistry.

\medterm{Isosorbide} Isosorbide is a nitrate medicine that relaxes blood vessels and is used mainly to prevent or relieve angina; headache and low blood pressure are common adverse effects.

\medterm{Isosorbide Dinitrate} An organic nitrate (and prodrug) that releases nitric oxide, relaxing vascular smooth
muscle—predominantly venous (preload reduction) with coronary dilation.

\medterm{Isospora} An older name for a group of intestinal coccidian parasites, now generally classified as Cystoisospora. Infection can cause watery diarrhea, especially in people with weakened immunity.

\medterm{Isospora Belli} The former name for Cystoisospora belli, a parasite spread by swallowing oocysts in contaminated food or water. It causes watery diarrhea, cramps, nausea, and weight loss and can be severe or prolonged in immunocompromised people.

\medterm{Isosporiasis} Intestinal infection caused by the protozoan Cystoisospora belli. It typically causes
watery diarrhea, abdominal cramps, nausea, fever, and weight loss and may be prolonged
or severe in people with impaired cellular immunity.

\medterm{Isothiocyanate} A reactive compound formed from glucosinolates in some plants, especially cruciferous vegetables, and studied for antimicrobial and cancer-related effects.

\medterm{Isothiocyanates} Reactive plant chemicals formed from glucosinolates, especially in cruciferous vegetables. They are studied for possible protective effects against cancer, but concentrated supplements may have different effects from normal food amounts.

\medterm{Isotonic} Describing muscle contraction in which muscle length changes to move a joint under a relatively constant external load.

\medterm{Isotonic Exercise} Exercise in which muscles change length while producing force against a relatively constant load, as in lifting or lowering a weight.

\medterm{Isotonic Solution} A solution with an effective osmotic concentration close to that of plasma, such as normal saline or balanced crystalloid. Isotonic fluids may expand extracellular volume, but choice and amount must be tailored because excess can cause fluid overload or electrolyte disturbance.

\medterm{Isotope} One of two or more atoms of the same element with the same number of protons but different numbers of neutrons.

\medterm{Isotretinoin} A systemic retinoid used for severe, recalcitrant nodular acne. It reduces sebum
production, normalizes follicular keratinization, decreases Propionibacterium acnes, and
has anti-inflammatory effects. Dose is 0.5 to 1 mg/kg/day for 15 to 20 weeks, targeting
cumulative dose of 120 to 150 mg/kg.

\medterm{Isovaleric Acidemia} A rare inherited disorder in which the body cannot properly break down the amino acid leucine. Harmful substances build up and may cause poor feeding, vomiting, unusual sleepiness, low muscle tone, or a sweaty-feet odor, especially during illness or fasting. Newborn screening can detect it; treatment uses a controlled diet and prevents prolonged fasting.

\synonyms
IVA; Isovaleryl-CoA Dehydrogenase Deficiency

\medterm{Isoxsuprine} Isoxsuprine is a vasodilating medicine formerly used for selected circulation disorders; it may cause dizziness, fast heartbeat, or low blood pressure.

\medterm{Isozyme} One of several molecular forms of an enzyme that catalyze the same reaction but differ in structure, tissue distribution, or regulation.

\medterm{Ispaghula Husk} Ispaghula husk, or psyllium, is a soluble fiber that absorbs water and helps regulate stool consistency; adequate fluid is important when using it.

\medterm{Ispinesib} Ispinesib is an investigational kinesin-spindle protein inhibitor studied for anticancer effects by disrupting chromosome separation during cell division.

\medterm{Isradipine} Isradipine is a calcium-channel blocker used to lower blood pressure; it may cause flushing, headache, ankle swelling, or dizziness.

\medterm{ISUP Grade Group} A five-level system that grades prostate cancer by how abnormal the cancer looks under a microscope. Group 1 is the lowest risk, usually with a Gleason score of 6 or less. Groups 2 and 3 represent intermediate risk, while groups 4 and 5 represent high or highest risk. The result helps guide monitoring, surgery, radiation, or other treatment.

\synonyms
Grade Group; ISUP Prostate Grading System; WHO/ISUP Grade Group

\medterm{Itch} Itch, or pruritus, is an unpleasant skin sensation that provokes scratching and may arise from dermatitis, allergy, infection, medicines, kidney or liver disease, blood disorders, nerve disease, or systemic illness. Persistent generalized itch requires evaluation.

\synonyms
Pruritus (Itch)

\medterm{Itching} Itching, or pruritus, is an unpleasant sensation that provokes scratching and may arise from skin disease, allergy, systemic illness, medications, neuropathy, or psychogenic factors.
be acute or chronic and may arise from skin disease, systemic illness, medication, or
neuropathic mechanisms.

\medterm{Itching Anal (Anal Itching)} Alternate index label for Anal Itching.

\medterm{Itching, Anal} Itching around the anus, commonly caused by irritant dermatitis, hemorrhoids, pinworms, infection, inflammatory skin disease, or hygiene-related irritation. Persistent, severe, bleeding, painful, or unexplained symptoms require examination.

\medterm{Itchy Ear} An itching sensation of the external ear or ear canal, which may result from dermatitis, allergy, infection, dryness, or irritation from objects placed in the ear.

\medterm{Iterative Reconstruction} A computer method that repeatedly improves medical images by comparing calculated results with measured scan data. It can reduce image noise and sometimes allow lower radiation exposure, but usually requires more processing time.

\medterm{ITP} Alternate index label; see \textit{Immune thrombocytopenia}.

\medterm{Itraconazole} Itraconazole is a triazole antifungal that inhibits fungal ergosterol synthesis and treats several superficial and systemic mycoses; absorption, extensive drug interactions, liver injury, and negative inotropic effects require attention.

\medterm{IUCD} An intrauterine contraceptive device, usually referring to an intrauterine device placed in the uterus to prevent pregnancy.

\medterm{IV Fluids} Intravenous fluid therapy for volume resuscitation and maintenance. Crystalloids: normal
saline (0.9\% NaCl, increases chloride, risk of hyperchloremic metabolic acidosis),
lactated Ringer (balanced, contains K+, Ca2+, lactate buffer, safer in large volumes),
Plasma-Lyte (balanced, closest to plasma).

\medterm{IV Treatment At Home} Intravenous treatment given outside a hospital, usually through a catheter with instructions for safe infusion, line care, monitoring, and follow-up.

\medterm{Ivabradine} Ivabradine slows the heart rate by acting on the sinus node and is used for selected patients with chronic heart failure or stable angina.

\medterm{IVC Filter} Device placed in inferior vena cava to prevent pulmonary embolism. Indications: DVT/PE
with contraindication to anticoagulation, recurrent PE despite anticoagulation.
Complications: filter thrombosis, migration, perforation, IVC occlusion.

\medterm{Ivermectin} An antiparasitic medicine used for selected worm and mite infections, including strongyloidiasis, onchocerciasis, and scabies, when prescribed appropriately.

\medterm{IVF} In vitro fertilization: the assisted reproductive technology in which oocytes are
retrieved, fertilized by sperm outside the body (in vitro), and the resulting embryos
transferred to the uterus. Indications: tubal factor, male factor (with ICSI),
endometriosis, unexplained, ovulation disorders, genetic testing, and same-sex
couples or individuals using donor sperm or eggs.

\medterm{Ixabepilone} Ixabepilone is a chemotherapy medicine used for some advanced breast cancers; important risks include nerve damage, low blood counts, and allergic reactions.

\medterm{Ixodes} A genus of hard ticks, including the ticks that transmit Lyme disease and several other infections. A bite may be painless; removing the tick promptly reduces, but does not eliminate, infection risk.
\medterm{Internal Carotid Artery} A branch of the common carotid artery that enters the skull and supplies the brain, eyes, and related structures. It contributes to the cerebral arterial circulation.

\medterm{Idiopathic Anaphylaxis} Repeated anaphylaxis, a severe allergic reaction, for which no specific trigger is found after appropriate evaluation. It is still an emergency: use prescribed epinephrine and seek urgent medical care, while specialist assessment looks for hidden triggers or mast-cell disorders.

\medterm{Idiopathic Guttate Hypomelanosis} Harmless, small, flat white spots caused by reduced skin pigment, usually on the shins or forearms of older adults. They are not contagious and usually do not need treatment.

\medterm{Iliotibial Band Friction Syndrome} An overuse injury causing pain on the outside of the knee, and sometimes the hip, when the iliotibial band becomes irritated. It is common in runners and cyclists.

\medterm{Immunoassays} Laboratory tests that use specific binding between antibodies and target substances to detect or measure them. Examples include ELISA, chemiluminescent tests, and some rapid tests; results must be interpreted with the clinical situation.

\medterm{Immunomodulatory} Describing a treatment or substance that changes immune activity. It may stimulate, suppress, or otherwise regulate the immune response, depending on the medicine and condition.

\medterm{Implants} Medical devices, materials, or tissues placed in the body to replace, support, monitor, or deliver treatment. Possible problems include infection, rejection, movement, breakage, or failure.

\medterm{Impotence} An older term usually referring to erectile dysfunction: persistent difficulty getting or keeping an erection firm enough for sexual activity. Erectile dysfunction can have vascular, neurologic, hormonal, medication-related, or psychological causes.

\medterm{In Vitro Fertilization (IVF)} An assisted-reproduction treatment in which eggs are collected, fertilized with sperm in a laboratory, and an embryo is transferred to the uterus. Risks include ovarian hyperstimulation, multiple pregnancy, ectopic pregnancy, and procedure-related complications.

\medterm{Infantile Bradycardia} An unusually slow heart rate in an infant. It can occur normally during sleep, but may also signal low oxygen, infection, medicines, or a heart or conduction problem; symptoms and the infant's age determine the evaluation.

\medterm{Infectious Diseases} Illnesses caused by bacteria, viruses, fungi, parasites, or prions. They may spread through contact, droplets or air, food or water, insects, blood, or sexual contact; prevention and treatment depend on the organism.

\medterm{Infrared Therapy} Treatment using infrared radiation, usually to provide heat or light to tissues. It is used in some physiotherapy and skin treatments, but benefits vary and excessive exposure can cause burns or eye injury.

\medterm{Ingrown Toenails} A condition in which the edge of a toenail grows into nearby skin, most often on the big toe. It causes pain, redness, and swelling; infection, diabetes, or severe symptoms require medical advice.

\medterm{International Classification of Diseases} A World Health Organization system for coding diseases, injuries, and other health conditions. Versions such as ICD-10 and ICD-11 support clinical records, statistics, and billing; a code does not replace a clinical diagnosis.

\medterm{Interrupted Aortic Arch} A rare congenital heart defect in which the aorta has a complete gap between its ascending and descending parts. Newborns can rapidly develop poor circulation and heart failure, so it requires emergency stabilization and surgery.

\medterm{Intrauterine Insemination} A fertility procedure in which prepared sperm is placed directly into the uterus around ovulation. It is used for selected infertility causes and generally requires at least one open fallopian tube; risks include cramps, infection, and multiple pregnancy.

\medterm{Intravesical Therapy} Treatment delivered directly into the bladder through a catheter. It is used for conditions such as some bladder cancers or interstitial cystitis and may cause bladder irritation or, less commonly, effects elsewhere in the body.

\medterm{Iodine Perfusion} The distribution of iodine-containing contrast through tissues or blood vessels, measured during imaging such as computed tomography or angiography. Contrast can rarely cause allergic reactions and may affect kidney or thyroid function.

\medterm{Ion Channel} A protein pore in a cell membrane that allows selected electrically charged particles, such as sodium, potassium, calcium, or chloride, to cross. Ion channels are important for nerve signals, muscle contraction, and secretion.

\medterm{Iridodonesis} Shaking or trembling of the iris when the eye moves, often because the lens is absent or its supporting fibers are weak. An eye examination can identify the cause and related lens or eye problems.

\medterm{Iridoschisis} A rare splitting of the iris tissue into layers, usually in older adults. It may be associated with glaucoma or corneal problems and should be assessed by an eye specialist.

\medterm{Ischemic Optic Neuropathy} Damage to the optic nerve caused by reduced blood flow, often producing sudden, usually painless vision loss. It may be arteritic or non-arteritic; urgent assessment is especially important in older adults or when headache, jaw pain, or scalp tenderness occurs.

\medterm{Jaccoud Arthropathy} A form of joint disease that can permanently or temporarily pull the fingers toward the little-finger side without the bone damage typical of rheumatoid arthritis. It is linked most often with past rheumatic fever or lupus.

\medterm{Jaccoud's Arthropathy} A deforming but usually non-damaging form of arthritis that can cause the fingers to drift toward the little-finger side. The deformity may improve when the hand is moved and is often linked to lupus or past rheumatic fever.

\medterm{Jack-In-The-Pulpit Poisoning} Jack-in-the-pulpit poisoning follows ingestion of calcium oxalate crystals from the plant, causing intense burning and swelling of the mouth, throat, and gastrointestinal tract.

\medterm{Jackknife Position} A face-down position used during some operations in which the table is bent at the hips. The buttocks are raised while the head and legs are lower, improving access to the rectal, anal, or posterior pelvic regions.

\synonyms
Kraske Position

\medterm{Jacksonian Epilepsy} A focal seizure that begins with twitching or stiffening in one body part and spreads in an orderly way to nearby body parts. Awareness may remain normal at first.

\synonyms
Jacksonian March; Focal Motor Seizure

\medterm{Jacksonian March} A clinical progression of focal motor seizure activity in which epileptic discharges propagate along contiguous areas of the primary motor cortex (precentral gyrus homunculus), resulting in clonic contractions that begin focally in one anatomical region (such as the thumb or angle of the mouth) and march sequentially to involve adjacent muscle groups of the limb or face.
\synonyms
Jacksonian Seizure; Focal Motor March

\medterm{Jacksonian Seizure} A focal motor seizure that begins in one body part and may spread in an orderly pattern to adjacent body regions; awareness may remain intact at first.

\medterm{Jacobsen Syndrome} A contiguous gene deletion disorder caused by a terminal deletion of the long arm of chromosome 11 (11q23-q24), characterized clinically by Paris-Trousseau thrombocytopenia and platelet dysfunction, trigonocephaly, facial dysmorphism, structural cardiac defects, and psychomotor delay.
\synonyms
11q Terminal Deletion Syndrome; Jacobsen Deletion Disorder

\medterm{Jadassohn's Naevus} A hairless skin patch or raised lesion present from birth, usually caused by an overgrowth of oil glands. It may change during life and rarely develops another tumor, so changing lesions need review.

\medterm{Jadassohn-Lewandowsky Syndrome} The classic form of pachyonychia congenita (type 1) caused by pathogenic mutations in the KRT6A or KRT16 keratin genes, presenting with hypertrophic subungual hyperkeratosis and thickened dystrophic nails, severe painful palmoplantar keratoderma, oral mucosal leukokeratosis, and follicular hyperkeratosis.
\synonyms
Pachyonychia Congenita Type 1; PC-1

\medterm{Jaffe-Lichtenstein Syndrome} A benign bone condition characterized by polyostotic fibrous dysplasia involving multiple skeletal sites with replacement of normal lamellar bone by immature fibro-osseous tissue, occurring in the absence of endocrine hyperfunction or the café-au-lait hyperpigmentation seen in McCune-Albright syndrome.
\synonyms
Polyostotic Fibrous Dysplasia; Jaffe-Lichtenstein Disease

\medterm{JAK-STAT Pathway} A signaling system cells use to respond to hormones and immune chemicals. It helps control growth, blood-cell production, inflammation, and immune activity; abnormal signaling can contribute to inflammatory disease or cancer.

\medterm{JAK2} A protein inside cells that passes signals from certain hormones and immune chemicals. An activating change in the JAK2 gene can make the bone marrow produce too many blood cells, as in several myeloproliferative disorders.

\medterm{Jammed Finger} A finger injury caused by sudden force, often involving a sprain, joint dislocation, fracture, or tendon injury; persistent deformity, severe pain, or inability to move requires assessment.

\medterm{Janeway Lesions} Small, usually painless red or purple spots on the palms or soles that can occur with infective endocarditis, an infection of the heart lining or valves. They result from infected particles blocking tiny blood vessels.

\medterm{Janus Kinase} One of a family of proteins inside cells that pass signals from hormones and immune chemicals. Janus kinases help control inflammation, blood-cell production, and immune responses.

\medterm{Janus Kinase Inhibitors} Medicines that block Janus kinase proteins and reduce certain immune signals. They are used for selected inflammatory and blood disorders, but can increase the risk of infections, blood clots, and other serious effects.

\medterm{Japanese Encephalitis} A viral infection spread by mosquitoes that can inflame the brain. Most infections cause no symptoms, but severe illness can cause fever, confusion, seizures, movement problems, or lasting disability.

\medterm{Jargon Aphasia} A severe fluent receptive aphasia variant (frequently associated with extensive posterior superior temporal lobe lesions in Wernicke aphasia) characterized by fluent, effortless speech filled with neologisms, phonemic paraphasias, and unintelligible syntax ('word salad') with marked impairment in comprehension.
\synonyms
Neologistic Aphasia; Word Salad Speech

\medterm{Jarisch-Herxheimer Reaction} A short-lived reaction that may occur soon after antibiotics are started for certain infections. It can cause fever, chills, headache, muscle aches, and temporarily worse symptoms.

\medterm{Jaundice} Yellowish pigmentation of the skin, sclerae, and mucous membranes caused by elevated levels of circulating bilirubin in the blood. Etiologies include pre-hepatic hemolytic processes, intra-hepatic hepatocellular injury (hepatitis, cirrhosis), or post-hepatic biliary obstruction (gallstones, strictures, neoplasms).

\synonyms
Icterus; Hyperbilirubinemic Scleral Icterus; Biliary Pigmentation

\medterm{Jaundice And Breastfeeding} Jaundice associated with breastfeeding may reflect inadequate milk intake in the early days or persistent breast-milk jaundice after other serious causes have been excluded.

\medterm{Jaundice Causes} Jaundice can result from rapid breakdown of red blood cells, inherited bilirubin problems, liver disease, medicines or toxins, or blockage of bile flow. Blood tests and imaging help identify the cause and guide treatment.

\medterm{Jaundice In Adults} Adult jaundice is yellowing of skin or eyes from elevated bilirubin due to red-cell breakdown, liver injury, impaired bile flow, or inherited metabolism. Dark urine, pale stool, abdominal pain, fever, confusion, or rapid onset requires prompt assessment.

\medterm{Jaundice In Neonates} Yellowing of a newborn's skin or eyes from elevated bilirubin. It is often physiologic, but jaundice in the first 24 hours, rapidly increasing jaundice, poor feeding, lethargy, fever, pale stools, or dark urine requires urgent evaluation.

\medterm{Jaundice In Newborns} Yellowing of a newborn’s skin or eyes caused by increased bilirubin. It is often temporary, but very high or rapidly rising levels can injure the brain and require urgent treatment.

\medterm{Jaw} The upper and lower bony structures of the mouth—the maxilla and mandible—which support the teeth and contribute to chewing, speech, facial shape, and the oral airway.

\medterm{Jaw Fracture} A break in the lower jawbone or upper jawbone, usually after a blow to the face. Signs may include pain, swelling, loose teeth, numbness, a changed bite, or difficulty opening the mouth. Severe breaks can block the airway, especially when the lower jaw is broken in the front on both sides. Treatment may require realignment, wiring, plates, or surgery.

\medterm{Jaw Jerk Reflex} A brief closing movement of the jaw when the chin is tapped with the mouth slightly open. An unusually strong response may suggest a problem in the brain's motor pathways.

\medterm{Jaw Neoplasm} An abnormal growth in the upper or lower jawbone or nearby tissues. It may be a cyst, a benign tumor, or cancer and can cause swelling, pain, loose teeth, numbness, or bone destruction.

\medterm{Jaw Pain} Pain in the jaw, jaw joint, or chewing muscles caused by temporomandibular disorders, dental disease, infection, injury,
or less commonly heart or nerve disease. Pain with chest discomfort, sweating, breathlessness, or exertion
requires urgent assessment; persistent pain or restricted jaw movement requires dental or medical review.

\medterm{Jaw Tumors And Cysts} Abnormal growths or fluid-filled sacs developing in the jawbone or nearby tissues. They may be harmless or cancerous and can cause swelling, pain, tooth movement, numbness, or bone destruction.

\synonyms
Odontogenic Tumors And Cysts

\medterm{Jaw-Winking Syndrome} A congenital synkinetic neuro-ophthalmic condition (Marcus Gunn phenomenon) characterized by unilateral congenital ptosis that transiently resolves or retracts upward with movement of the lower jaw during mastication, sucking, or lateral pterygoid contraction, caused by anomalous congenital innervation of the levator palpebrae superioris by the motor branch of the trigeminal nerve (CN V3).
\synonyms
Marcus Gunn Jaw-Winking Phenomenon; Trigemino-Oculomotor Synkinesis

\medterm{Jehovah's Witness} A member of a Christian religious group whose beliefs may lead them to decline transfusions of whole blood and certain blood components. Individual choices about blood products, procedures, and treatment vary, so clinicians should discuss the patient's informed preferences and document them carefully.

\medterm{Jejunal Atresia} A condition present at birth in which part of the jejunum, a section of the small intestine, is closed or missing. It causes bowel blockage in a newborn, with green vomiting, a swollen abdomen, and failure to pass the first stool.

\medterm{Jejunal Fistula} An abnormal connection involving the jejunum and another organ, skin, or body surface. It may cause intestinal contents to leak, infection, malnutrition, dehydration, or sepsis and requires assessment of its cause and output.

\medterm{Jejunal Hemorrhage} Bleeding from the jejunum, the middle segment of the small intestine, which may cause maroon or black stool, anemia, abdominal pain, or shock depending on its rate and cause.

\medterm{Jejunal Neoplasm} An abnormal growth in the jejunum, the middle part of the small intestine. It may be benign or cancerous and can cause pain, bleeding, anemia, vomiting, blockage, or weight loss.

\medterm{Jejunal Stenosis} Narrowing of the jejunal lumen that can impede intestinal contents and cause cramping, vomiting, distension, malabsorption, or obstruction; causes include Crohn disease, adhesions, tumors, and ischemic injury.

\medterm{Jejunal Ulcer} An ulcerated break in the lining of the jejunum, caused by inflammation, infection, ischemia, medicines, tumor, or anastomotic disease and potentially producing pain, bleeding, or perforation.

\medterm{Jejunitis} Jejunitis is inflammation of the jejunum, which may result from infection, inflammatory bowel disease, ischemia, medication, or other intestinal disorders.

\medterm{Jejunoileal Atresia} A birth defect in which part of the small intestine is completely closed or missing, so milk and digestive fluid cannot pass. Newborns develop green vomiting, a swollen abdomen, and failure to pass stool. Imaging confirms the blockage, and surgery is needed to restore intestinal passage.

\synonyms
Jejunal Atresia; Ileal Atresia; Small Bowel Atresia

\medterm{Jejunostomy} A surgically created opening into the jejunum, often used for placement of a feeding tube when nutrition cannot safely be given through the stomach.

\medterm{Jejunostomy Feeding Tube} A jejunostomy feeding tube delivers nutrition directly into the jejunum when oral or gastric feeding is unsafe or inadequate, with attention to formula, rate, position, and complications.

\medterm{Jejunum} The middle part of the small intestine. It absorbs much of the nutrients from food before the remaining contents pass into the ileum.

\medterm{Jellyfish Sting} A painful skin reaction caused when a jellyfish tentacle releases venom into the skin. It may cause burning, redness, and raised lines; severe exposure can cause vomiting, breathing trouble, fainting, or widespread symptoms.

\synonyms
Jellyfish Envenomation; Sea Sting

\medterm{Jellyfish Stings} Alternate name; see \textit{Jellyfish stings}.
\medterm{Jendrassik Maneuver} A clinical examination technique used to reinforce and elicit sluggish deep tendon reflexes (such as the patellar tendon reflex) by having the patient hook their flexed fingers together and pull forcefully apart, which increases spinal motoneuron pool excitability by reducing descending central inhibition.
\synonyms
Jendrassik Reinforcement; Reflex Reinforcement Maneuver

\medterm{Jendrassik's Maneuver} A technique used during a reflex examination in which a person pulls the hands apart or tenses another muscle group. The distraction can make a knee or other tendon reflex easier to observe.

\medterm{Jerk Nystagmus} An involuntary, rhythmic oscillation of the eyes characterized by a slow drift in one direction followed by a rapid corrective saccadic movement in the opposite direction. The direction of the nystagmus is designated by the fast phase, and it may arise from vestibular, cerebellar, or brainstem pathology.

\medterm{Jerusalem Cherry Poisoning} Jerusalem cherry poisoning follows ingestion of Solanum pseudocapsicum berries, which contain toxic compounds that may cause gastrointestinal symptoms, confusion, weakness, or cardiac effects.

\medterm{Jervell and Lange-Nielsen Syndrome} A rare inherited condition causing severe hearing loss from birth and an abnormal heart rhythm. The rhythm problem can cause fainting, seizures, or sudden death, so urgent specialist care is important.

\medterm{Jet Lag} Temporary circadian rhythm disruption after rapid travel across time zones, causing sleep disturbance, fatigue, impaired concentration, and sometimes gastrointestinal symptoms.

\synonyms
Desynchronosis; Time Zone Change Syndrome

\medterm{Jet Lag Prevention} Jet lag prevention uses timed light exposure, sleep adjustment, hydration, and carefully planned activity or melatonin strategies to shift the circadian rhythm across time zones.

\medterm{Jeune Syndrome} An autosomal recessive skeletal ciliopathy caused by mutations in intraflagellar transport genes, characterized by a severely constricted, narrow bell-shaped thorax, short-limbed micromelic skeletal dysplasia, polydactyly, pulmonary hypoplasia, and progressive cystic renal and hepatic disease.
\synonyms
Asphyxiating Thoracic Dystrophy; ATD

\medterm{Jewelry Cleaner Poisoning} Jewelry cleaner poisoning can injure the mouth, esophagus, stomach, lungs, or eyes depending on the product and exposure route; immediate poison-control guidance is appropriate.

\medterm{Jimsonweed Poisoning} Jimsonweed poisoning causes anticholinergic toxicity, including dilated pupils, dry mouth, rapid heart rate, urinary retention, agitation, delirium, seizures, or coma.

\medterm{Job Syndrome} A rare inherited immune disorder causing eczema, repeated skin and lung infections, and very high levels of an allergy-related antibody. Dental, bone, and facial differences may also occur.

\medterm{Jock Itch} A fungal infection of the groin and inner thighs, also called tinea cruris. It causes an itchy, scaly, ring-shaped or spreading
rash, often worsened by warmth and moisture. Pain, fever, extensive involvement, or failure to improve with
appropriate care requires clinical evaluation.

\synonyms
Tinea Cruris; Ringworm of the Groin

\medterm{Jod-Basedow Phenomenon} Overactive thyroid function triggered by an excess of iodine, usually in a person with thyroid nodules or other thyroid tissue that works independently of normal control. It may cause weight loss, tremor, sweating, or a fast heartbeat.

\medterm{Joffroy Sign} A physical examination sign of Graves disease and hyperthyroidism characterized by the absence of normal forehead wrinkling and furrowing when the patient looks upward while keeping the head in a fixed forward position.
\synonyms
Joffroy's Sign; Thyrotoxic Forehead Sign

\medterm{Johanson-Blizzard Syndrome} A rare autosomal recessive congenital malformation syndrome caused by mutations in the UBR1 gene, characterized by exocrine pancreatic insufficiency, hypoplasia or aplasia of the nasal alae ('beaked nose'), scalp vertex defects, oligodontia, sensorineural hearing loss, and growth retardation.
\synonyms
JBS; Pancreatic Insufficiency-Alar Hypoplasia Syndrome

\medterm{Joint} The place where two or more bones meet. Joints may be fixed, slightly movable, or freely movable, such as the knee, hip, or shoulder.

\medterm{Joint Aspiration} A needle procedure that removes fluid from a joint for testing or symptom relief, helping assess infection, bleeding, crystals, or inflammation.
\medterm{Joint Capsule} A tough envelope surrounding a freely movable joint. Its inner lining makes slippery fluid that helps the joint move, while the capsule helps hold the bones together and provides stability.

\medterm{Joint Contracture} Permanent or long-lasting limitation of joint movement caused by tightening or scarring of muscles, tendons, ligaments, the joint capsule, skin, or nearby tissues. It may follow immobility, injury, burns, inflammation, or nerve disease.

\medterm{Joint Diseases} Conditions that affect joints, including osteoarthritis, inflammatory arthritis, gout, infection, injury, and tumors. They may cause pain, swelling, stiffness, deformity, or reduced movement; treatment depends on the cause.

\medterm{Joint Effusion} Extra fluid inside a joint, often causing swelling, stiffness, or restricted movement. Causes include injury, infection, inflammatory arthritis, bleeding, and crystal disease; a clinician may remove fluid for testing.

\medterm{Joint Fluid Culture} Joint fluid culture tests synovial fluid for bacterial, fungal, or other organisms when septic arthritis is suspected; results are interpreted with cell count, crystals, and clinical findings.

\medterm{Joint Fluid Gram Stain} A quick laboratory test that stains fluid removed from a swollen joint to look for bacteria. It can support the evaluation of a serious joint infection but a negative result does not rule infection out.

\medterm{Joint Inflammation} Alternate name; see \textit{Arthritis}.

\medterm{Joint Mouse} A loose piece of cartilage, bone, or other tissue inside a joint that can cause catching, locking, pain, or swelling.

\medterm{Joint Pain} Pain or aching in one or more joints or in the tissues around them; the medical term is arthralgia. It is a symptom rather than a single disease and may result from arthritis, inflammation, injury, overuse, infection, crystals such as those in gout, or problems in nearby muscles, tendons, or bursae. It may occur with stiffness, swelling, warmth, redness, tenderness, or reduced movement. The number of joints affected, the timing, and associated symptoms help identify the cause.

\medterm{Joint Prosthesis} An artificial joint or implant designed to replace a damaged or diseased biological
joint, restoring mobility and relieving chronic pain.

\medterm{Joint Replacement} Surgery that removes damaged joint surfaces and replaces them with artificial parts to reduce pain and improve movement. It is commonly performed for severe arthritis or joint destruction and requires rehabilitation.

\medterm{Joint Replacement Rehabilitation} Exercise, education, and daily-activity training after joint replacement, helping restore safe movement, strength, balance, confidence, and independence.

\medterm{Joint Stiffness} Reduced ease or range of joint movement, often worse after rest. Causes include arthritis, injury,
inflammation, prolonged immobility, or muscle and tendon disorders. Sudden swelling, redness,
fever, severe pain, or inability to use the joint requires assessment.

\medterm{Joint Swelling} Enlargement of a joint caused by excess synovial fluid, blood, inflammation, infection,
trauma, or structural disease. It may be accompanied by pain, warmth, stiffness, or
restricted movement and can indicate inflammatory or septic arthritis.

\medterm{Joint X-Ray} An X-ray picture of a joint used to look for broken bones, abnormal alignment, arthritis-related wear, bone loss, deposits, or other changes. It uses a small amount of radiation and does not show all soft-tissue injuries.

\medterm{Jones Criteria} A standardized set of diagnostic guidelines formulated by the American Heart Association to establish the diagnosis of acute rheumatic fever following group A streptococcal pharyngitis, requiring evidence of antecedent streptococcal infection along with two major criteria (carditis, migratory polyarthritis, Sydenham chorea, erythema marginatum, subcutaneous nodules) or one major and two minor criteria.
\synonyms
Revised Jones Criteria; AHA Rheumatic Fever Criteria

\medterm{Jones Fracture} An acute, traumatic, extra-articular transverse fracture of the fifth metatarsal bone located at the metaphyseal-diaphyseal junction (approximately 1.5 cm distal to the tuberosity), prone to delayed union and nonunion due to poor vascular supply in this watershed zone.
\synonyms
Fifth Metatarsal Watershed Fracture; Metaphyseal-Diaphyseal Fifth Metatarsal Fracture

\medterm{Joubert Syndrome} A rare inherited disorder of brain development that affects breathing control, balance, eye movements, and muscle tone. Babies may have pauses or rapid changes in breathing, poor muscle tone, delayed development, and abnormal eye movements. Some also have kidney, liver, or retinal disease, or extra fingers or toes. Brain MRI helps confirm the diagnosis.

\synonyms
JSRD; Joubert-Related Disorder

\medterm{Jugular} Relating to the neck, especially the jugular veins that carry blood from the head and neck back to the heart.

\medterm{Jugular Foramen} An opening at the base of the skull through which a major neck vein and several nerves pass. Disease in this area can cause difficulty swallowing, hoarseness, shoulder weakness, hearing symptoms, or other neurologic problems.

\medterm{Jugular Foramen Syndrome} A clinical neurological syndrome resulting from compressive or infiltrative lesions (such as glomus jugulare tumors, schwannomas, or skull base metastases) affecting the structures traversing the jugular foramen, characterized by ipsilateral paralysis of the glossopharyngeal (CN IX), vagus (CN X), and accessory (CN XI) cranial nerves, leading to dysphagia, vocal cord paralysis, and trapezius weakness.
\synonyms
Vernet Syndrome; Posterior Lacerated Foramen Syndrome

\medterm{Jugular Vein} A major vein that carries blood from the head, face, and neck back to the heart. Each side has an internal jugular vein and an external jugular vein.

\medterm{Jugular Venous Distension} Pathological engorgement and visible protrusion of the internal or external jugular veins in the neck above the level of the clavicle with the patient positioned at 45 degrees, providing a direct physical reflection of elevated right atrial and central venous pressure, classic for right heart failure, cardiac tamponade, and constrictive pericarditis.
\synonyms
JVD; Raised Jugular Venous Pressure

\medterm{Jugular Venous Pressure} The visible height of blood in the neck veins. It helps clinicians estimate pressure in the right side of the heart and assess fluid buildup.

\synonyms
JVP

\medterm{Jugulodigastric Node} A prominent deep cervical lymph node situated at the junction where the posterior belly of the digastric muscle crosses the internal jugular vein, receiving primary lymphatic drainage from the palatine tonsils, base of tongue, and posterior pharynx, frequently enlarged in acute pharyngotonsillitis and head and neck neoplasms.
\synonyms
Tonsillar Lymph Node; Subdigastric Node

\medterm{Jumper's Knee} Alternate name; see \textit{Patellar Tendinitis}.

\medterm{Junctional Escape Rhythm} A physiological backup cardiac rhythm originating from pacemaking cells within the atrioventricular (AV) junction and bundle of His at an intrinsic rate of 40 to 60 beats per minute, characterized by narrow QRS complexes with absent, inverted, or retrograde P waves, arising when sinoatrial node pacemaking fails or during complete AV nodal block.
\synonyms
AV Junctional Rhythm; Nodal Escape Rhythm

\medterm{Junctional Nevus} A usually benign mole formed by melanocytes at the junction between the epidermis and dermis; it is often flat and should be assessed if it changes in size, shape, color, or symptoms.

\medterm{Junctional Rhythm} A heart rhythm that starts near the connection between the upper and lower chambers instead of the heart’s usual pacemaker. It may be slow or regular and can occur when the usual pacemaker is suppressed or blocked.

\medterm{Junctional Tachycardia} A supraventricular arrhythmia arising from enhanced automaticity or triggered activity within the AV node or bundle of His, characterized by a rapid heart rate (greater than 100 beats per minute), regular narrow QRS complexes, and retrograde or dissociated atrial activation, frequently associated with digitalis toxicity, myocardial ischemia, or pediatric cardiac surgery.
\synonyms
Nonparoxysmal Junctional Tachycardia; Accelerated Junctional Rhythm

\medterm{Jungian Analysis} Jungian analysis is a psychotherapy approach based on Carl Jung’s theories, using exploration of unconscious processes, symbols, dreams, and personality development.

\medterm{Juvenile Absence Epilepsy} An idiopathic generalized epilepsy syndrome with onset in adolescence (typically between 9 and 13 years of age), characterized by absence seizures occurring less frequently than in childhood absence epilepsy, accompanied on EEG by generalized 3 to 4 Hz spike-and-wave discharges, frequently accompanied by generalized tonic-clonic seizures.
\synonyms
JAE; Adolescent Absence Epilepsy

\medterm{Juvenile Angiofibroma} A noncancerous but highly blood-vessel-rich tumor that usually develops near the back of the nasal cavity in adolescent boys. It may cause worsening nasal blockage, repeated nosebleeds, or facial swelling and should be assessed by specialists.

\medterm{Juvenile Arthritis} A general term for inflammatory joint disease beginning in childhood. It may cause persistent swelling, pain, stiffness, reduced movement, fever, rash, or eye inflammation; the subtype determines monitoring and treatment.

\synonyms
Juvenile Idiopathic Arthritis; JIA; Juvenile Rheumatoid Arthritis

\medterm{Juvenile Astrocytoma} A tumor arising from support cells in a child’s or adolescent’s brain or spinal cord. Many grow slowly, but symptoms depend on location and may include headache, seizures, vision changes, weakness, balance problems, or hormone disturbances.

\medterm{Juvenile Dermatomyositis} An autoimmune illness in children that weakens muscles, especially around the hips and shoulders, and causes a purple rash around the eyes or raised patches over the knuckles. It may also affect swallowing, breathing, or other organs.

\medterm{Juvenile Diabetes} Alternate historical term; see \textit{Type 1 Diabetes}.

\medterm{Juvenile Hypermobile Spectrum Disorder} A hypermobility spectrum disorder beginning in childhood or adolescence, with excessive joint movement plus pain, instability, soft-tissue injury, fatigue, or related symptoms.
\medterm{Juvenile Hypermobility Syndrome} A pattern of unusually flexible joints in a child or adolescent that causes pain, joint instability, tiredness, or repeated sprains. Assessment checks for other connective-tissue disorders; treatment may include strengthening, activity guidance, and joint protection.

\medterm{Juvenile Idiopathic Arthritis} A group of inflammatory joint diseases that begin before age 16 and last at least six weeks without another identified cause. Children may have swelling, stiffness, pain, fever, rash, or eye inflammation, depending on the subtype.

\medterm{Juvenile Myoclonic Epilepsy} A generalized epilepsy syndrome that typically begins in adolescence and causes myoclonic jerks, often after awakening, with possible generalized tonic-clonic or absence seizures.

\medterm{Juvenile Nasopharyngeal Angiofibroma} A benign, highly vascular, unencapsulated mesenchymal tumor occurring almost exclusively in adolescent males, originating in the posterolateral wall of the nasal cavity near the sphenopalatine foramen, presenting clinically with recurrent painless epistaxis, nasal obstruction, and progressive facial swelling.
\synonyms
JNA; Nasopharyngeal Fibroangioma

\medterm{Juvenile Osteoporosis} Low bone strength during childhood or adolescence that increases the risk of fractures. It may occur on its own or result from chronic illness, poor nutrition, hormone disorders, limited mobility, or medicines.

\medterm{Juvenile Polyposis Syndrome} An inherited condition causing many noncancerous growths in the stomach or intestines. The growths can bleed, cause anemia, and increase the risk of bowel and stomach cancer, so regular screening is needed.

\medterm{Juvenile Rheumatoid Arthritis} Alternate historical term; see \textit{Juvenile Idiopathic Arthritis}.

\medterm{Juvenile Schizophrenia} Alternate name; see \textit{Childhood Schizophrenia}.

\medterm{Juvenile Xanthogranuloma} A rare, usually harmless growth of immune cells in the skin of an infant or young child, appearing as a firm red, yellow, or brown bump. It often fades without treatment, but rarely affects the eye or internal organs.

\medterm{Juxta-Articular} Located next to a joint. The term may describe bone, soft tissue, swelling, a cyst, or a deposit beside a joint rather than within the joint itself.

\medterm{Juxtaglomerular Apparatus} A group of specialized kidney cells that senses blood flow and salt levels and releases renin. It helps regulate blood pressure, blood volume, and the body’s salt balance.

\medterm{Juxtaglomerular Cell Tumor} A rare, benign renal neoplasm arising from the specialized renin-secreting myoepithelial juxtaglomerular cells of the afferent arteriole, characterized by marked hyperreninemia, secondary hyperaldosteronism, severe hypertension, and hypokalemia.
\synonyms
Reninoma; Juxtaglomerular Renin-Secreting Tumor

\medterm{Juxtavesical} Next to the urinary bladder. The term often describes the lower part of a ureter as it approaches the bladder, where a passing stone may become lodged and cause side pain, urinary frequency, or painful urination.

\medterm{Kabuki Syndrome} A rare genetic disorder that may cause distinctive facial features, short stature, developmental or learning difficulties, weak muscle tone, skeletal changes, and heart defects. Hearing loss, seizures, immune problems, and kidney abnormalities may also occur. Most cases result from a change in the \textit{KMT2D} or \textit{KDM6A} gene.

\synonyms
Niikawa-Kuroki Syndrome

\medterm{Kahler Disease} Another name for multiple myeloma, a cancer of plasma cells in the bone marrow. It can produce an abnormal protein and cause weak or damaged bones, high calcium, kidney problems, and anemia.
\synonyms
Multiple Myeloma; Plasma Cell Myeloma

\medterm{Kahler's Disease} Alternate name; see \textit{Multiple Myeloma}.

\medterm{Kala-Azar} A serious parasitic infection caused by Leishmania species and transmitted by sandflies. It may cause prolonged fever, weight loss, enlargement of the spleen and liver, and suppression of blood-cell production.

\synonyms
Black Sickness; Dumdum Fever

\medterm{Kallikrein-Kinin System} A group of blood and tissue proteins that makes chemical messengers such as bradykinin. These messengers widen small blood vessels, make them more leaky, and contribute to pain and inflammation.
\synonyms
Kininergic System; Bradykinin Cascade

\medterm{Kallmann Syndrome} An inherited condition that causes delayed or absent puberty and a reduced or absent sense of smell. It occurs when the brain does not send the hormone signals needed to start puberty.

\medterm{Kallmann Syndrome and IHH} Conditions in which deficient hypothalamic or pituitary gonadotropin signaling causes low sex hormones and absent or delayed puberty. Kallmann syndrome also includes impaired smell and may be genetic; fertility and sex-hormone replacement require endocrine care.

\medterm{Kanamycin} An antibiotic used for some serious bacterial infections. It can permanently damage hearing and harm the kidneys, especially with prolonged or high-dose use.

\medterm{Kangaroo Care} Holding a newborn directly against a caregiver’s bare chest while monitoring the infant’s
condition. It supports thermal regulation, physiologic stability, bonding, and feeding,
and is particularly useful for many preterm or low-birth-weight infants.

\synonyms
Kangaroo Mother Care; KMC

\medterm{Kaolin} A soft clay mineral used in some medicines, laboratory products, and industrial materials. Breathing large amounts of kaolin dust over time can damage the lungs.

\medterm{Kaposiform Hemangioendothelioma} A rare blood-vessel tumor that usually affects infants and young children. It can grow into nearby tissue and may cause a serious loss of platelets and problems with blood clotting.

\synonyms
KHE

\medterm{Kaposi Sarcoma} A vascular neoplasm associated with human herpesvirus 8 (HHV-8). It may cause nonblanching red, purple, brown, or black macules, plaques, or nodules on the skin or oral mucosa and can involve lymph nodes or internal organs; disease is more likely or aggressive with immunosuppression, including HIV or transplantation.

\medterm{Kaposi's Sarcoma} A vascular tumor associated with human herpesvirus 8, also called Kaposi sarcoma-associated herpesvirus. It may cause purple, red, or brown lesions of the skin or oral mucosa and can involve lymph nodes or internal organs; it is more likely or aggressive when immunity is weakened, such as with HIV infection or after organ transplantation.

\medterm{Karnofsky Performance Status} A scale used to describe how well a person can perform daily activities and care for themselves. It helps clinicians assess overall function and ability to tolerate treatment, especially in cancer care.

\medterm{Kartagener Syndrome} An inherited condition in which tiny hair-like structures that clear mucus do not work properly. It causes repeated sinus and lung infections and may place the internal organs in a reversed position.

\medterm{Kartagener Triad} The combination of reversed internal-organ positions, long-term sinus inflammation, and widened damaged airways. It results from an inherited problem with microscopic moving hairs that clear mucus.

\medterm{Karyomegaly} Abnormal enlargement of a cell’s nucleus, the part that contains most of its genetic material. It may occur with injury, infection, toxic exposure, tissue repair, or cancer and is interpreted with other findings.

\medterm{Karyotype} An organized picture of a person’s chromosomes, arranged in pairs to show their number and large-scale structure. It can detect some missing, extra, or rearranged chromosomes but not every genetic condition.

\medterm{Karyotype Result} A laboratory report describing the number and visible structure of a person’s chromosomes. It may identify some chromosome differences, such as an extra chromosome, but a normal result does not exclude all genetic disorders.

\medterm{Karyotyping} A laboratory test that examines chromosomes in dividing cells to detect some missing, extra, or rearranged chromosomes. It may be used in pregnancy, infertility, repeated miscarriage, developmental conditions, or blood cancers.

\medterm{Kasabach-Merritt Phenomenon} A serious blood-clotting problem linked to certain fast-growing blood-vessel tumors in babies. Platelets become trapped in the tumor, causing a very low platelet count, anemia from broken red blood cells, and low levels of clotting proteins.

\synonyms
KMP; Kasabach-Merritt Syndrome

\medterm{Kasabach-Merritt Syndrome} Another name for Kasabach-Merritt phenomenon: a serious clotting problem in babies with certain fast-growing blood-vessel tumors. It causes very low platelets, anemia from broken red blood cells, and low levels of clotting proteins.
\synonyms
Kasabach-Merritt Phenomenon; Vascular Tumor Consumptive Coagulopathy

\medterm{Kasai Procedure} An operation for babies with biliary atresia, in which the blocked bile ducts outside the liver are removed and a loop of intestine is connected to the liver so bile can drain. Earlier surgery gives the best chance of restoring bile flow, but many children eventually need a liver transplant.

\synonyms
Hepatoportoenterostomy; Kasai Operation; Portoenterostomy

\medterm{Kashin-Beck Disease} A long-term disease of the bones and joints that can cause pain, stiffness, deformity, and poor growth. It is linked to low selenium and environmental factors in some regions.

\medterm{Katayama Fever} An acute systemic reaction to schistosome infection, usually occurring weeks after exposure and causing fever, rash, cough, abdominal symptoms, and eosinophilia.

\synonyms
Katayama Syndrome; Acute Bilharziasis

\medterm{Katz Index} A scale that assesses independence in six basic activities of daily living: bathing, dressing, toileting, transferring, continence, and feeding.

\medterm{Kawasaki Disease} An acute inflammatory disease of childhood that can affect the coronary arteries. It commonly causes persistent fever, rash, conjunctival redness, mouth or extremity changes, and swollen lymph nodes; diagnosis is clinical and incomplete forms occur.

\synonyms
Mucocutaneous Lymph Node Syndrome; KD

\medterm{Kawasaki Disease Diagnostic Criteria} The usual clinical pattern used to identify Kawasaki disease: fever for at least five days plus at least four of five findings—red eyes without pus, a varied skin rash, swollen neck lymph nodes, changes in the lips or mouth, and redness, swelling, or peeling of the hands or feet. Some children have an incomplete form.

\medterm{Kayser-Fleischer Rings} Golden-brown or greenish rings around the colored part of the eye caused by copper buildup. They strongly suggest Wilson disease when other findings fit and are seen with a special eye examination.

\medterm{K-Complex} A large brain-wave pattern seen during stage 2 sleep. It may occur on its own or after a sound or other sensory stimulus and is recorded during an electroencephalogram (EEG).

\medterm{Kearns-Sayre Syndrome} A rare inherited disorder affecting the mitochondria, the parts of cells that make energy. It usually begins in childhood or young adulthood and can cause drooping or limited eye movement, vision problems, heart-rhythm problems, and muscle weakness.

\medterm{Keel Breast} A chest-wall shape in which the breastbone pushes outward, giving the chest a keel-like appearance. It usually reflects pectus carinatum, which develops during growth and may affect appearance, breathing, or exercise tolerance.

\medterm{Kegel} A pelvic-floor muscle contraction exercise used to improve pelvic-floor strength and control of selected types of urinary or fecal incontinence.

\medterm{Kegel Exercise} Repeated tightening and relaxing of pelvic-floor muscles to improve strength and control, helping reduce urine leakage and support pelvic organs.

\medterm{Kehr Sign} Pain felt at the tip of the left shoulder when blood or irritation affects the underside of the diaphragm, the main breathing muscle. It can occur with irritation from a spleen injury but is not specific to it.

\medterm{Keloid} A raised, firm scar that grows beyond the edges of the original wound because the body makes too much scar tissue. It may itch or hurt and can develop after surgery, injury, burns, or piercings.

\synonyms
Keloid Scar; Cheloid

\medterm{Keloid and Hypertrophic Scar} Raised scars caused by excessive collagen deposition during wound healing. Hypertrophic scars remain within the original wound, whereas keloids extend beyond it; symptoms may include itch, pain, tightness, and cosmetic concern.

\medterm{Keloid Scar} A raised, firm scar that grows beyond the original skin injury because of excessive scar formation. It may itch, hurt, or recur after removal and is more common in some families and skin types.

\medterm{Kennedy Disease} An X-linked recessive, adult-onset lower motor neuron disorder caused by a CAG trinucleotide repeat expansion in the androgen receptor (AR) gene, presenting with slowly progressive proximal limb and bulbar muscle weakness, tongue and perioral fasciculations, muscle cramps, and signs of androgen insensitivity (gynecomastia and testicular atrophy).
\synonyms
Spinal and Bulbar Muscular Atrophy; SBMA

\medterm{Kennedy's Disease} A rare inherited motor-neuron disease, also called spinal and bulbar muscular atrophy, caused by an androgen-receptor gene change. It mainly affects men and causes slowly progressive muscle weakness, wasting, cramps, and sometimes swallowing or speech problems.

\medterm{Keratan Sulfate} A complex sugar found mainly in the clear front of the eye, cartilage, and brain. In some inherited disorders it cannot be broken down properly and builds up, damaging the eyes, bones, or other tissues.

\medterm{Keratan Sulphate} A sulfated glycosaminoglycan found especially in the cornea and cartilage; abnormal accumulation occurs in some mucopolysaccharidoses.

\medterm{Keratectomy} Removal or reshaping of part of the cornea, the clear front surface of the eye. It may be done to correct vision or remove abnormal tissue, and it can affect eyesight.

\medterm{Keratic Precipitates} Small groups of inflammatory cells that collect on the inner surface of the cornea, usually during inflammation inside the front of the eye. Their size and appearance can help describe the type of inflammation.

\medterm{Keratin} A strong structural protein forming much of the skin’s outer layer, hair, nails, and other protective tissues. It helps resist water loss, friction, chemicals, and physical injury.

\medterm{Keratinocyte} The main type of cell in the outer layer of skin. It produces keratin and helps form the skin barrier that limits water loss and protects against injury and infection.

\medterm{Keratinocytes} The predominant cells of the epidermis that produce keratin and form the skin’s protective barrier.

\medterm{Keratin Pearl} A round, layered collection of keratin seen under a microscope in some well-differentiated squamous cell skin cancers.

\synonyms
Squamous Pearl; Horn Pearl

\medterm{Keratitis} Inflammation or infection of the clear front surface of the eye. It can cause eye pain, redness, tearing, sensitivity to light, blurred vision, or a cloudy spot.

\synonyms
Corneal Ulceration; Corneal Inflammation

\medterm{Keratoacanthoma} A fast-growing, dome-shaped skin growth with a central plug, usually on sun-exposed skin. It may look like a type of skin cancer, so a clinician often removes or samples it. Treatment depends on its appearance, size, site, and the test results.

\synonyms
Well-Differentiated Squamous Cell Carcinoma, Keratoacanthoma Type; KA; Crateriform Epithelioma

\medterm{Keratocele} Severe thinning or ulceration of the cornea that allows its inner layers to bulge outward. It can lead to a hole in the cornea and permanent loss of vision.

\medterm{Keratoconjunctivitis} Inflammation of the cornea and the conjunctiva, the clear surface covering the white of the eye. Infection, allergy, dryness, immune disease, or chemical injury may cause redness, pain, discharge, light sensitivity, and blurred vision.

\medterm{Keratoconjunctivitis Sicca} Long-lasting dryness and inflammation of the eye’s front surface, commonly called dry-eye disease. It can cause burning, grittiness, redness, fluctuating vision, or light sensitivity because too few tears are made or tears evaporate too quickly.

\synonyms
Dry Eye Syndrome; Ocular Surface Disease

\medterm{Keratoconus} A condition in which the clear front surface of the eye becomes thin and bulges into a cone shape. It can cause blurred or distorted vision, glare, and sensitivity to light. Glasses, special contact lenses, or treatment to strengthen the cornea may help.

\medterm{Keratocyte Count} An estimate of the number of living support cells in the middle layer of the cornea. It can help assess corneal health, injury, disease, or changes after eye surgery, although results are interpreted with the complete eye examination.

\medterm{Keratoderma} Excessive thickening of the skin on the palms or soles. It may be inherited or develop because of skin disease, infection, medicines, or another illness and can cause painful cracks or difficulty walking.

\medterm{Keratoderma Blennorrhagica} Thickened, scaly, sometimes pustular skin lesions, usually on the soles or palms, associated with reactive arthritis. Treatment focuses on the underlying infection or inflammation and the skin symptoms.

\medterm{Keratoderma Blennorrhagicum} Thickened, scaly, or pus-filled skin lesions on the palms and soles associated with reactive arthritis. The lesions may spread to other areas and often appear after an infection affecting the urinary or digestive tract.

\medterm{Keratoderma Blenorrhagicum} Alternate spelling of keratoderma blennorrhagicum, a cutaneous manifestation of reactive
arthritis with waxy, crusty, vesiculopustular lesions on the palms and soles. It is a
reactive phenomenon typically occurring in the context of genitourinary or
gastrointestinal infection.

\medterm{Keratodermia Blennorrhagicum} Thickened, scaly, or wart-like skin lesions, often on the palms and soles, associated with reactive arthritis and sometimes infection-triggered inflammation.

\medterm{Keratolysis} Breakdown or loss of keratinized tissue, especially the superficial layer of the skin; it may occur in some infections and inflammatory skin disorders.

\medterm{Keratolytics} Substances, such as salicylic acid or urea, that soften and loosen excess hard skin. They may be used for thickened skin, corns, calluses, or some warts.

\medterm{Keratoma} A localized thickening or tumor-like growth of keratinized skin, such as a callus or corn; the exact type and need for evaluation depend on its appearance and behavior.

\medterm{Keratomalacia} Softening and ulceration of the cornea caused most often by severe vitamin A deficiency, potentially leading to blindness.

\medterm{Keratometer} An eye instrument that measures the curvature of the cornea, the clear front surface of the eye. The measurement helps assess astigmatism and fit rigid contact lenses.

\medterm{Keratometry} Measurement of the curve of the cornea, the clear front surface of the eye. It helps assess focusing problems and choose contact lenses or an artificial lens for cataract surgery.

\medterm{Keratoplasty} Surgery to replace or reshape part or all of the cornea, the clear front surface of the eye. It can restore vision after corneal damage, but healing problems, infection, and rejection of donor tissue are possible.

\medterm{Keratoprosthesis} An artificial cornea implanted to restore sight when a donor-corneal transplant is unlikely to work, such as after severe chemical injury or extensive surface damage. It carries risks of infection, glaucoma, and device problems.

\medterm{Keratosis} A thickened, rough, or otherwise abnormal area of skin caused by excess keratin. Some forms are harmless, while others result from infection, inflammation, sun damage, or a change that can precede skin cancer.

\medterm{Keratosis Follicularis} A condition in which extra keratin blocks hair follicles and causes many rough, small bumps. The name may describe keratosis pilaris or a rarer inherited skin disorder, so the appearance and history determine the diagnosis.

\medterm{Keratosis Follicularis (Darier Disease)} An inherited disorder of epidermal cell adhesion that causes persistent greasy, crusted, brown keratotic papules in seborrheic areas, nail changes, and sometimes mucosal lesions. Heat, sweating, friction, and ultraviolet exposure may worsen it.

\medterm{Keratosis Obturans} A buildup of shed skin cells that blocks the ear canal. It can cause ear pain, hearing loss, discharge, and widening or damage of the canal.

\medterm{Keratosis Pilaris} A common, harmless skin condition that causes small, rough bumps around hair follicles, usually on the upper arms, thighs, buttocks, or face.

\synonyms
KP

\medterm{Kerion} An inflamed, boggy, pus-forming mass caused by a fungal infection of hair-bearing skin, usually a severe form of tinea capitis.

\medterm{Kerley A Lines} Long, fine lines seen on a chest X-ray, usually spreading outward from the central lung areas. They can reflect fluid buildup in the lung tissue.
\synonyms
Kerley's A Lines; Upper Zone Interstitial Edema Lines

\medterm{Kerley B Lines} Short, thin lines seen near the outer edges and bases of the lungs on a chest X-ray. They commonly reflect fluid buildup in the lung tissue, often related to heart failure.
\synonyms
Kerley's B Lines; Septal Lines

\medterm{Kernicterus} Permanent brain damage caused by very high levels of bilirubin, a yellow substance that builds up during severe newborn jaundice. It can cause movement problems, hearing loss, problems looking upward, and damage to tooth enamel.

\synonyms
Bilirubin Encephalopathy; Nuclear Jaundice; Chronic Bilirubin Toxicity

\medterm{Kernig Sign} Pain or resistance when the knee is straightened while the hip is bent. It may suggest irritation of the coverings of the brain and spinal cord, such as with meningitis, but it is not diagnostic by itself.

\medterm{Kernig's Sign} Alternate spelling; see \textit{Kernig Sign}.

\medterm{Kernohan Notch Phenomenon} A paradoxical false-localizing neurological sign in which acute uncal transtentorial herniation displaces the midbrain against the contralateral rigid tentorial notch (Kernohan notch), compressing the contralateral cerebral peduncle (corticospinal tract) and producing an unexpected ipsilateral hemiparesis.
\synonyms
Kernohan-Woltman Notch Phenomenon; Kernohan Sign

\medterm{Kerosene Poisoning} Harm after swallowing or inhaling kerosene. The greatest danger is entry into the lungs, which can cause chemical pneumonia; coughing, choking, breathing difficulty, drowsiness, or vomiting may occur.

\medterm{Keshan Disease} A heart-muscle disease linked to severe selenium deficiency, occurring mainly in regions where soil and food contain little selenium. It can weaken the heart and cause heart failure.

\medterm{Ketamine} Ketamine is a dissociative anaesthetic and NMDA-receptor antagonist used for anaesthesia, procedural sedation, and selected treatment-resistant depression; emergence reactions, increased blood pressure, nausea, misuse, and respiratory complications are possible.

\medterm{Ketamine Hydrochloride} Ketamine hydrochloride is a dissociative anesthetic used for anesthesia, procedural sedation, and selected pain or depression treatments; it can cause increased blood pressure, hallucinations, and impaired coordination.

\medterm{Ketamine Toxicity} Acute adverse effects or overdose of the NMDA receptor antagonist dissociative anesthetic ketamine, characterized by perceptual distortions, hallucinations, severe emergence delirium, profound nystagmus, hypertension, tachycardia, and in chronic abusers, severe ulcerative cystitis ("ketamine bladder").

\medterm{Keto Acid} A chemical made when the body breaks down certain amino acids or other nutrients. Keto acids can be used to make energy or new substances, but abnormal amounts may occur in metabolic disease.

\medterm{Ketoacidosis} A dangerous buildup of acidic ketones in the blood, usually because there is too little insulin or prolonged starvation. It can cause vomiting, abdominal pain, deep breathing, confusion, coma, and death.

\medterm{Ketoconazole} Ketoconazole is an imidazole antifungal that inhibits fungal ergosterol synthesis; topical forms treat superficial mycoses, while oral use is restricted because of potentially severe liver injury, adrenal suppression, and drug interactions.

\medterm{Ketogenesis} The production of ketone bodies from fatty acids, mainly in the liver during fasting, prolonged exercise,
or inadequate insulin activity.

\medterm{Ketogenic Diet} An eating plan very low in carbohydrates and higher in fat that makes the body use ketones for energy. It can help some children with difficult-to-control epilepsy, but may cause nutrient deficiencies, constipation, kidney stones, or abnormal cholesterol without supervision.

\medterm{Ketone} A type of chemical compound. In medicine, the term often refers to ketone bodies made when the body uses fat for energy; small amounts can be normal, but high levels may cause ketoacidosis.

\medterm{Ketone Bodies} Water-soluble molecules produced mainly by the liver from fatty-acid-derived acetyl-CoA
during fasting, starvation, prolonged exercise, or insulin deficiency. The principal
ketone bodies are acetoacetate, beta-hydroxybutyrate, and acetone.

\medterm{Ketones} Water-soluble substances produced mainly by the liver when the body breaks down fat for energy, such as during fasting, a low-carbohydrate state, or insufficient insulin. The main ketone bodies are beta-hydroxybutyrate, acetoacetate, and acetone. Small amounts can be normal, but high levels may cause ketoacidosis, especially in diabetes; ketones can be measured in blood or urine.

\medterm{Ketones Blood Test} A blood test measuring ketone bodies, such as beta-hydroxybutyrate, used to assess ketosis and help diagnose or monitor diabetic ketoacidosis and other causes of ketone production.

\medterm{Ketones Urine Test} A urine test detecting ketone bodies produced when the body metabolizes fat; it can support evaluation of diabetic ketoacidosis, starvation, vomiting, or other metabolic states but may lag behind blood ketone levels.

\medterm{Ketonuria} Ketonuria is the presence of ketone bodies in urine, occurring with fasting, vomiting, low-carbohydrate intake, pregnancy, or uncontrolled diabetes; marked ketonuria with hyperglycaemia may indicate diabetic ketoacidosis.

\medterm{Ketoprofen} A non-steroidal anti-inflammatory drug that inhibits cyclo-oxygenase and reduces pain, inflammation, and fever; gastrointestinal bleeding, kidney injury, cardiovascular events, and hypersensitivity are important risks.

\medterm{Ketoprofen Overdose} Harm caused by taking too much ketoprofen, a nonsteroidal anti-inflammatory drug. It may cause nausea, vomiting, stomach bleeding, drowsiness, kidney injury, seizures, or a dangerous buildup of acid in the blood.

\medterm{Ketorolac} A potent non-steroidal anti-inflammatory analgesic used for short-term management of moderate-to-severe acute pain; treatment duration is restricted because gastrointestinal bleeding, kidney injury, and impaired healing can occur.

\medterm{Ketorolac Tromethamine} The tromethamine salt of ketorolac, a potent non-steroidal anti-inflammatory analgesic used short term for moderate-to-severe acute pain; it is contraindicated in several bleeding, renal, pregnancy, and perioperative situations.

\medterm{Ketosis} A state in which the body makes and uses increased amounts of ketones after breaking down fat for energy. It can occur during fasting or a low-carbohydrate diet, but high ketone levels with illness or high blood sugar may signal diabetic ketoacidosis.

\medterm{Ketosis-Prone Type 2 Diabetes} A form of diabetes that may present with diabetic ketoacidosis despite features of type 2 diabetes, such as obesity or preserved insulin production. Insulin dependence can change over time according to beta-cell reserve, autoimmunity, and recovery from the initial metabolic stress.
\medterm{Ketotic Hypoglycemia} Low blood sugar accompanied by ketones, usually after a young child has gone without food or has been ill and unable to eat. Drowsiness, weakness, vomiting, or seizures may occur. Repeated episodes can have other causes.

\medterm{Ketotifen} An antihistamine with mast-cell-stabilizing activity used in some countries for allergic disorders and, in ophthalmic form, allergic conjunctivitis; drowsiness, increased appetite, and irritation may occur.

\medterm{Kidney} One of two bean-shaped organs that filter waste from the blood into urine. The kidneys also help control body water, salts, blood pressure, and the balance of acids in the blood.

\medterm{Kidney Abscess} Alternate name for renal abscess, a localized collection of pus in or around the kidney.

\synonyms
Renal Abscess; Perinephric Abscess

\medterm{Kidney Biopsy} A procedure in which a small sample of kidney tissue is removed, usually with a needle guided by ultrasound, to identify kidney disease. Bleeding is the main risk.

\medterm{Kidney Calculi} Alternate term; see \textit{Kidney stones}.

\medterm{Kidney Cancer} A malignant tumor arising in the kidney. Renal cell carcinoma is the most common type in adults, but other types occur, especially in children.

\medterm{Kidney Cancer In Children} Cancer arising in a child’s kidney. Wilms tumor is the most common type, but other kidney tumors can occur; symptoms may include an abdominal mass, pain, blood in the urine, fever, or high blood pressure.

\medterm{Kidney Cysts} Fluid-filled sacs in or on the kidneys. Simple cysts are common and usually harmless, while complex or multiple cysts may be associated with infection, bleeding, inherited disease, or cancer risk and may impair kidney function.

\synonyms
Renal Cysts; Renal Cystic Disease

\medterm{Kidney Disease} Any condition that damages the kidneys or reduces their ability to filter blood, balance fluids and minerals, control blood pressure, or make hormones. Causes include diabetes, high blood pressure, infection, inflammation, inherited disorders, and medicines.

\medterm{Kidney Disease and Pregnancy} Kidney disease can affect fertility, blood pressure, fluid balance, medication choices, and pregnancy outcomes, while pregnancy can worsen some kidney disorders. Preconception counseling and coordinated obstetric-nephrology care help monitor kidney function, proteinuria, blood pressure, fetal growth, and medicine safety.

\medterm{Kidney Disease-Related Anemia} Anemia caused or worsened by kidney disease, mainly because damaged kidneys produce less erythropoietin, the hormone that stimulates red-cell production. Iron deficiency, inflammation, and blood loss may also contribute.

\medterm{Kidney Diseases} Disorders affecting the kidneys or urinary collecting system that can impair filtration, urine formation, fluid balance, blood pressure, or electrolyte regulation.

\medterm{Kidney Failure} A condition in which the kidneys can no longer remove enough waste or control fluid and body salts. It may develop suddenly or gradually and can require medicines, dialysis, or a transplant.

\synonyms
End-Stage Renal Disease; ESRD; Renal Failure

\medterm{Kidney Function Tests} Blood and urine tests used to assess how well the kidneys remove waste and control fluid and salts. Common measures include creatinine, estimated filtration rate, urine protein, and urinalysis.

\medterm{Kidney Infection} An infection of one or both kidneys, usually caused by bacteria that travel upward from the bladder. It can cause fever, chills, side or back pain, nausea, vomiting, and painful urination and may lead to sepsis.

\medterm{Kidney Neoplasm} An abnormal growth in the kidney that may be noncancerous or cancerous. It may cause blood in the urine, pain, a lump, weight loss, or no symptoms; imaging and sometimes biopsy identify its type and extent.

\medterm{Kidney Pain} Pain felt in the side or back below the ribs that may be caused by kidney stones, infection, obstruction, inflammation, or referred pain.

\synonyms
Flank Pain; Renal Colic; Nephralgia

\medterm{Kidney Position} A side-lying surgical position in which the waist is gently flexed and the side of the body is raised. It gives surgeons better access to the kidney and nearby tissues while staff protect pressure points and breathing.

\medterm{Kidney Removal} Surgery to remove part or all of a kidney, usually for cancer, severe damage, infection, or a nonfunctioning kidney. The remaining kidney may provide adequate function, but bleeding, infection, injury to nearby organs, or reduced kidney function can occur.

\medterm{Kidney Replacement Therapy} Treatment that replaces some functions of severely damaged kidneys. It includes hemodialysis, peritoneal dialysis, continuous kidney replacement therapy, and kidney transplantation.

\medterm{Kidney Stone} Singular form; see \textit{Kidney stones}.

\medterm{Kidney Stones} Hard mineral deposits formed from substances in urine. They may obstruct the urinary tract and cause colicky flank or abdominal pain, blood in the urine, vomiting, or infection.

\medterm{Kidney Stones In Children} Mineral deposits that form in a child’s kidney or urinary tract. They may cause side or abdominal pain, blood in the urine, nausea, vomiting, urinary symptoms, or infection, but some cause no symptoms.

\synonyms
Pediatric Nephrolithiasis; Pediatric Urolithiasis; Pediatric Renal Calculi

\medterm{Kidney Transplant} Surgical placement of a healthy donor kidney into a person with kidney failure to restore renal function. Lifelong immunosuppressive treatment is generally required to reduce rejection risk.

\medterm{Kidney Transplantation} Placement of a donor kidney to replace the filtering work of kidneys that have permanently failed. A living-donor kidney may work sooner, but all recipients need lifelong medicines and monitoring for rejection, infection, blood-vessel problems, and return of kidney disease.

\medterm{Kidney Transplant Rejection} An immune attack that damages a transplanted kidney. It may occur suddenly or gradually and can cause reduced urine, swelling, high blood pressure, or worsening blood tests; treatment depends on the type and severity.

\medterm{Kidney Trauma} Alternate name for renal trauma, an injury to the kidney caused by blunt or penetrating force.

\medterm{Kienbock Disease} Damage and gradual collapse of the lunate, a small bone in the wrist, because its blood supply is reduced. It may cause wrist pain, swelling, stiffness, and weaker grip. Treatment depends on bone damage, wrist alignment, symptoms, and whether arthritis has developed.

\medterm{Kieselguhr} Diatomaceous earth containing substantial silica; inhalation of its dust can produce occupational pneumoconiosis and other lung injury.

\medterm{Kiesselbach Plexus} A network of small arteries on the front part of the nasal septum and a common source of nosebleeds.

\medterm{Kiesselbach Triangle} An anterior vascular network on the anteroinferior nasal septum (Little's area) formed by anastomoses of the anterior ethmoidal, sphenopalatine, greater palatine, and superior labial arteries, serving as the anatomical site for over 90% of anterior epistaxis episodes.
\synonyms
Little Area; Kiesselbach's Plexus

\medterm{Kikuchi Disease} A usually short-term inflammation of the lymph nodes, most often in the neck, with no known cause. It mainly affects young adults and may cause tender swollen nodes, fever, night sweats, and a low white-blood-cell count.

\synonyms
Kikuchi-Fujimoto Disease; Histiocytic Necrotizing Lymphadenitis

\medterm{Kikuchi-Fujimoto Disease} Another name for Kikuchi disease, a usually short-term inflammation of the lymph nodes, most often in the neck. It may cause painful swollen nodes, fever, night sweats, and a low white-blood-cell count.
\synonyms
Kikuchi Disease; Histiocytic Necrotizing Lymphadenitis

\medterm{Killer Cold Virus} Alternate name; see \textit{Adenovirus Infection}.

\medterm{Killian's Triangle} An area of potential weakness in the pharyngeal wall between the thyropharyngeus and
cricopharyngeus muscles, predisposing to Zenker diverticulum formation.

\medterm{Killian Triangle} A triangular anatomical area of muscular weakness in the posterior hypopharyngeal wall situated between the oblique fibers of the thyropharyngeus muscle and the transverse fibers of the cricopharyngeus muscle; the site of mucosal herniation forming a Zenker diverticulum.

\synonyms
Killian Dehiscence

\medterm{Killip Classification} A four-level system for describing heart failure during a heart attack. The levels range from no signs of heart failure to lung fluid buildup and, at the most severe level, shock caused by the heart’s inability to pump enough blood.

\medterm{Kilocalorie} A unit of energy equal to 1,000 small calories; the dietary Calorie, written with a capital C, is one kilocalorie.

\medterm{Kilojoule} The SI unit of energy, equal to 1,000 joules; approximately 4.184 kilojoules equal one kilocalorie.

\medterm{Kimmelstiel-Wilson Lesion} A round area of thickened tissue in a kidney filter that can occur with long-standing diabetes. It is a sign of diabetic kidney damage and may occur with protein in the urine or reduced kidney function.
\synonyms
Nodular Diabetic Glomerulosclerosis; Kimmelstiel-Wilson Nodules

\medterm{Kimmelstiel-Wilson Nodules} Round areas of thickened tissue seen in the kidney filters of some people with long-standing diabetes. They reflect diabetic kidney damage and may occur with protein in the urine and declining kidney function.

\medterm{Kimmelstiel-Wilson Syndrome} A pattern of kidney scarring caused by long-standing diabetes. It can make protein leak into the urine and gradually reduce kidney function, sometimes leading to kidney failure.

\medterm{Kimura Disease} A rare long-term immune disorder causing painless lumps under the skin, usually in the head or neck, enlarged nearby lymph nodes, and high levels of eosinophils or IgE. It can sometimes affect the kidneys.

\medterm{Kinase} An enzyme that attaches a small chemical group called a phosphate to another molecule. This acts like a switch controlling cell growth, energy use, movement, communication, and many other processes.

\medterm{Kinase Inhibitor} A targeted therapeutic agent (typically a small molecule or monoclonal antibody) that selectively inhibits the enzymatic phosphorylation activity of specific protein kinases (such as tyrosine kinases, serine/threonine kinases, or cyclin-dependent kinases) involved in oncogenic cellular signaling cascades.

\medterm{Kinesia Paradoxa} A brief, unexpected ability to move normally or quickly in a person who usually has severe slowness or difficulty moving, as can occur in advanced Parkinson disease. It may happen during a strong emotional experience or in response to an urgent situation.
\synonyms
Paradoxical Kinesia; Paradoxical Movement

\medterm{Kinesthesia} The sense of how the body's parts are moving and where they are positioned, even when the eyes are closed. Information from muscles, joints, and skin helps the brain control posture, balance, and coordinated movement.

\medterm{Kinetics} The study of how a process changes over time. In medicine, pharmacokinetics describes how the body absorbs, distributes, changes, and removes a medicine, helping determine its dose and timing.

\medterm{Kinetochore} A specialized, disc-shaped multiprotein complex assembled onto the centromeric DNA region of a eukaryotic chromosome that captures spindle microtubules and coordinates chromosome segregation during mitosis and meiosis.

\medterm{Kingella} Kingella is a genus of bacteria that can live in the mouth and can occasionally cause bloodstream, bone, joint, or heart-valve infections.

\medterm{Kingella Denitrificans} Kingella denitrificans is a rare bacterium found mainly in the upper airway; it can occasionally cause invasive infection.

\medterm{Kingella Kingae} Kingella kingae is a bacterium that can cause bone or joint infection, especially in young children, and can rarely cause endocarditis.

\medterm{Kingella Oralis} Kingella oralis is a mouth-associated bacterium that rarely causes invasive infection, including bloodstream infection or endocarditis.

\medterm{Kingella Potus} Kingella potus is a rare bacterium that has been associated with occasional human infections; clinical importance depends on the specimen and evidence of true infection.

\medterm{King's College Criteria} A set of findings used to estimate the risk of death in acute liver failure and the possible need for a liver transplant. The criteria differ for acetaminophen poisoning and other causes.

\medterm{Kinin} A short peptide, such as bradykinin, that widens blood vessels, increases vascular permeability, and contributes to pain and inflammation.

\medterm{Kinin-Kallikrein System} An endogenous enzymatic cascade of blood proteins that produces bioactive peptide kinins (primarily bradykinin and kallidin) from high- and low-molecular-weight kininogens, mediating vasodilation, increased microvascular permeability, tissue inflammation, and pain sensation.

\synonyms
Kallikrein-Kinin Cascade

\medterm{Kinins} A group of vasoactive peptides, chiefly bradykinin and kallidin, generated from kininogens during inflammation; they promote vasodilation, increased vascular permeability, smooth-muscle effects, and pain.

\medterm{Kinocilium} A specialized hair-like structure on certain sensory cells, especially in the inner ear. It helps the cell detect the direction of movement and contributes to balance and hearing signals.

\medterm{Kissing Disease} Informal term; see \textit{Infectious mononucleosis}.

\medterm{Kite Angle} An anatomical radiographic angle measured on an anteroposterior radiograph of the foot between the longitudinal axis of the talus and the longitudinal axis of the calcaneus. A normal angle ranges from 20 to 40 degrees; convergence or reduction of this angle below 15 degrees signifies hindfoot parallelism, a diagnostic hallmark of congenital talipes equinovarus (clubfoot).

\synonyms
Talocalcaneal Angle; AP Kite Angle

\medterm{Klatskin Tumor} A cancer arising where the right and left bile ducts join near the liver. It can block bile flow and cause jaundice, itching, pale stools, or weight loss.

\medterm{Klebsiella} A genus of Gram-negative bacteria that can cause pneumonia, urinary-tract infection, bloodstream infection, and other healthcare-associated infections.

\medterm{Klebsiella Oxytoca} Klebsiella oxytoca is a bacterium that can cause urinary, bloodstream, wound, and hospital-acquired infections and can sometimes cause antibiotic-associated hemorrhagic colitis.

\medterm{Klebsiella Pneumonia} Pneumonia caused by Klebsiella bacteria. It can cause fever, cough, chest pain, breathlessness, and sometimes thick or blood-stained sputum; severe infection may destroy lung tissue and is more likely in people with serious illness or weakened defenses.

\medterm{Klebsiella Pneumoniae} Klebsiella pneumoniae is a bacterium that can cause pneumonia, urinary infection, bloodstream infection, and liver abscess; some strains are highly antibiotic-resistant.

\medterm{Klebsiella Rhinoscleromatis} A gram-negative bacillus causing chronic granulomatous infection of the nasal passages
and upper airway, leading to rhinoscleroma with progressive nasal obstruction and tissue
destruction.

\medterm{Klebsiella Singaporensis} Klebsiella singaporensis is a rare environmental bacterium that has occasionally been recovered from human clinical specimens; its significance depends on the specimen and symptoms.

\medterm{Klebsiella Variicola} Klebsiella variicola is a bacterium closely related to K. pneumoniae that can cause human infections, including bloodstream and urinary infections.

\medterm{Kleihauer-Betke Test} A quantitative acid-elution blood test used to detect and calculate the volume of fetomaternal hemorrhage; exposure to an acid buffer elutes adult hemoglobin (HbA) from maternal erythrocytes while acid-resistant fetal hemoglobin (HbF) inside fetal red blood cells remains intact and stains dark pink, allowing calculation of necessary Rh immune globulin (RhoGAM) dosing.
\synonyms
Acid Elution Test; Fetomaternal Hemorrhage Assay

\medterm{Kleine-Levin Syndrome} A rare disorder causing repeated episodes of excessive sleep, often with confusion, reduced activity, and changes in behavior or appetite. Episodes last days to weeks and are separated by periods of more typical function.

\medterm{Kleptomania} A mental-health disorder involving recurrent inability to resist impulses to steal items that are not needed for personal use or monetary value. The acts are typically preceded by rising tension and followed by relief or gratification, and cause distress or impairment.

\synonyms
Compulsive Stealing
Pathological Stealing

\medterm{Klinefelter's Syndrome} Alternate name; see \textit{Klinefelter Syndrome}.

\medterm{Klinefelter Syndrome} A genetic condition in males caused by one or more extra X chromosomes, most often the 47,XXY pattern. It can cause small testes, low testosterone, infertility, tall stature, less facial and body hair, breast enlargement, weak bones, and learning or language difficulties. A chromosome test identifies the condition.
\synonyms{47,XXY Syndrome; XXY Syndrome; Klinefelter-Reifenstein-Albright Syndrome}

\medterm{Klippel-Feil Syndrome} A condition present from birth in which two or more neck bones are joined together. It may cause a short neck, low hairline, limited neck movement, hearing loss, or spinal and kidney problems.

\medterm{Klippel-Trenaunay Syndrome} A rare congenital vascular disorder characterized by a classic triad of features: capillary malformations (port-wine stains), soft tissue and bone enlargement (hypertrophy, usually affecting one leg or arm), and slow-flow venous and lymphatic malformations (such as persistent embryonic veins, prominent varicose veins, and lymphedema). It is typically caused by non-inherited (somatic) mutations in the \textit{PIK3CA} gene. Complications include limb length discrepancy, chronic pain, recurrent skin infections (cellulitis), and blood clots (deep vein thrombosis and pulmonary embolism). Management is multidisciplinary and includes compression garments, physiotherapy, vascular procedures, and orthopedic care.

\synonyms
KTS; Klippel-Trénaunay Syndrome; Angio-Osteohypertrophy Syndrome; KTW Syndrome; KTWS

\medterm{Klippel-Trenaunay-Weber Syndrome} A congenital vascular and limb-overgrowth disorder with capillary malformation, venous or lymphatic abnormalities, and enlargement of a limb or its digits. It can cause pain, bleeding, ulcers, clotting, or pulmonary embolism and requires vascular assessment.

\medterm{Kl\"{u}ver-Bucy Syndrome} A profound neurobehavioral syndrome resulting from bilateral damage or surgical ablation of the anterior temporal lobes, specifically including the amygdaloid nuclei and hippocampus. Clinical hallmarks include visual agnosia (psychic blindness), compulsive oral exploration of objects (hyperorality), hypermetamorphosis, marked passivity or docility, altered dietary preferences, and inappropriate hypersexuality.

\synonyms
Kluver-Bucy Syndrome; Bilateral Amygdala Ablation Syndrome

\medterm{Klumpke Palsy} An injury to the lower trunk of the brachial plexus involving spinal nerve roots C8 and T1, typically resulting from upward traction on an abducted arm during birth or trauma, causing paralysis of intrinsic hand muscles resulting in a 'claw hand' deformity and ipsilateral Horner syndrome (if sympathetic T1 root fibers are disrupted).
\synonyms
Klumpke-Dejerine Palsy; Lower Trunk Brachial Plexus Injury

\medterm{Klumpke Paralysis} A nerve injury affecting the lower part of the network that supplies the arm and hand. It can cause severe hand weakness, numbness along the inner forearm and hand, and occasionally drooping of the eyelid and a smaller pupil on one side.

\medterm{Klumpke's Paralysis} Lower brachial-plexus palsy involving mainly the C8 and T1 nerve roots, causing hand weakness and sensory loss; Horner syndrome may accompany severe cases.

\medterm{Kluyvera} A genus of gram-negative bacteria in the Enterobacterales group; species can rarely cause opportunistic human infections.

\medterm{Kluyvera Ascorbata} A gram-negative Kluyvera species found in the environment and occasionally associated with urinary, bloodstream, or other opportunistic infections.

\medterm{Kluyvera Cryocrescens} A gram-negative environmental bacterium that rarely causes opportunistic human infection.

\medterm{Kluyvera Georgiana} A rarely encountered Gram-negative Kluyvera species that may appear in clinical samples. It can occasionally cause opportunistic infection, especially in people with serious illness, weakened defenses, or invasive medical devices.

\medterm{Kluyvera Intermedia} A Gram-negative bacterium occasionally found in human samples and rarely causing opportunistic urinary, wound, bloodstream, or other infections.

\medterm{Knee} The synovial joint between the femur, tibia, and patella, stabilized by the cruciate and collateral ligaments and menisci; it permits flexion and extension and bears body weight.

\medterm{Knee Arthroplasty} Surgery that replaces damaged parts of the knee joint with artificial components. It is usually performed for severe pain and loss of function from arthritis or major joint damage.

\medterm{Knee Arthroscopy} Knee arthroscopy is a minimally invasive procedure in which a camera and small instruments are inserted into the knee to diagnose or treat selected joint problems.

\medterm{Knee Bursitis} Inflammation of a fluid-filled sac around the knee, causing local swelling, tenderness, and pain with kneeling or movement. Repeated pressure, injury, arthritis, gout, or infection can contribute; fever or redness may occur with infection.

\synonyms
Prepatellar Bursitis; Housemaid's Knee; Pes Anserine Bursitis

\medterm{Kneecap} A flat, rounded bone at the front of the knee, embedded in the tendon of the thigh muscles. It protects the joint and improves the leverage used to straighten the knee; injury or arthritis can cause pain and restricted movement.

\medterm{Kneecap Dislocation} An injury in which the patella slips completely out of the groove at the end of the thighbone, usually toward the outside. It causes sudden pain, deformity, swelling, and difficulty walking and needs prompt assessment.

\medterm{Knee-Chest Position} A position in which the patient rests on the knees and chest, with the hips flexed; it may be used for spinal or anorectal procedures.

\medterm{Knee CT Scan} A detailed X-ray scan of the knee made by a computer. It is useful for complex fractures, bone alignment, arthritis, infection, masses, and selected soft-tissue problems, but it exposes the person to more radiation than a plain X-ray.

\medterm{Knee Dislocation} Complete anatomical disruption and loss of articulation between the distal femur and proximal tibia, representing a severe orthopaedic limb-threatening emergency frequently associated with multi-ligament rupture and high incidence of popliteal artery injury and peroneal nerve palsy.

\medterm{Knee Disorders} Problems affecting the knee joint, including injuries to ligaments or menisci, arthritis, bursitis, tendon disorders, infection, and fractures. They may cause pain, swelling, stiffness, instability, locking, or difficulty bearing weight.

\medterm{Knee Effusion} The abnormal accumulation of excess synovial fluid, pus, or blood within the joint capsule of the knee, clinically detected by the ballottement of the patella or bulge sign, warranting diagnostic arthrocentesis to differentiate inflammatory, septic, or hemarthrotic causes.

\medterm{Knee Hyperextension} Bending the knee backward beyond its normal range. It can strain or tear ligaments and may cause pain, swelling, instability, or difficulty bearing weight.

\medterm{Knee Injury And Meniscus Tears} A meniscus tear is damage to the cartilage cushion between the thighbone and shinbone, often after twisting or with degeneration. Pain, swelling, catching, or locking may occur; treatment ranges from rehabilitation to surgery according to tear, symptoms, and stability.

\medterm{Knee-Jerk} The automatic kick of the lower leg when the tendon below the kneecap is tapped. It tests part of the nerve pathway from the thigh and lower spinal cord and can help identify nerve or spinal problems.

\medterm{Knee Jerk Reflex} A monosynaptic deep tendon stretch reflex elicited by percussing the patellar ligament below the patella with a reflex hammer, stretching muscle spindles within the quadriceps femoris and provoking reflex knee extension mediated through the femoral nerve and spinal segments L2, L3, and L4.
\synonyms
Patellar Reflex; Quadriceps Reflex

\medterm{Knee Joint} A synovial joint formed mainly by the femur, tibia, and patella. It permits flexion and extension with limited rotation and is
stabilized by ligaments, menisci, muscles, and the joint capsule.

\medterm{Knee Joint Replacement} Surgery replacing damaged surfaces of the knee with artificial parts to relieve severe pain and improve movement. It is commonly used for advanced arthritis; infection, blood clots, stiffness, loosening, persistent pain, or repeat surgery are possible.

\medterm{Knee Meniscus} A pair of crescent-shaped fibrocartilaginous pads (medial and lateral menisci) interposed between the femoral condyles and the tibial plateau, functioning to distribute compressive axial loads, absorb shock, lubricate the knee joint, and deepen the tibial articular surfaces.

\medterm{Knee Microfracture Surgery} An operation that makes small holes in the bone beneath a damaged area of knee cartilage so marrow cells can form repair tissue. It may improve symptoms in selected small defects, but the new tissue is not identical to normal cartilage and results vary.

\medterm{Knee MRI Scan} Magnetic resonance imaging of the knee using a magnetic field and radio waves to evaluate cartilage, menisci, ligaments, tendons, bone marrow, muscles, and joint fluid.

\medterm{Knee Pain} Pain felt in or around the knee, commonly caused by injury, osteoarthritis, bursitis, tendinitis, or inflammatory disease. Swelling, instability, inability to bear weight, or fever may also occur.

\synonyms
Gonalgia; Knee Arthralgia

\medterm{Knee Prosthesis} An artificial joint component used to replace damaged surfaces of the knee, usually including femoral and tibial components and sometimes a patellar component.

\medterm{Knee Replacement} An operation that replaces part or all of a damaged knee joint with an artificial joint to reduce pain and improve movement. Recovery includes wound care, exercises, and rehabilitation.

\medterm{Kniest Dysplasia} A rare inherited disorder of bone and cartilage growth that causes short stature, enlarged joints, unusual facial features, and progressive joint problems. Some people also develop serious eye, hearing, or breathing complications.

\medterm{Knife-Cut Ulcers From Crohn’S Disease} Deep, narrow ulcers in the intestinal lining that look like cuts and can occur in Crohn disease. They may cause bleeding, narrowing of the bowel, abnormal connections, pain, or blockage and are identified during intestinal examination.

\medterm{Knock-Knee} Genu valgum, a leg alignment in which the knees angle inward and may touch while the ankles remain apart.

\medterm{Knock Knees} Medial angulation of the knees causing the knees to touch while the ankles remain apart,
also called genu valgum. It is common during childhood development but may be pathologic
when severe, asymmetric, persistent, or associated with pain or growth disturbance.

\synonyms
Genu Valgum

\medterm{Knotless Anchor} A small surgical anchor securing sutures in bone without requiring tied knots. Orthopedic surgeons use it to repair tendons, ligaments, or other soft tissues, with risks including loosening, infection, or tissue failure.

\medterm{Knuckle} A projecting finger joint, usually a metacarpophalangeal joint; it may be affected by trauma, inflammation, or arthritis.

\medterm{Knuckle Pad} A firm, usually painless thickening of skin and tissue over a finger joint. It is commonly related to repeated pressure or a fibrous tissue change and is generally harmless.

\medterm{Kocher Clamp} A surgical instrument with teeth used to hold firm tissue or clamp blood vessels. Its strong grip is useful during operations, but careless use can crush or injure delicate tissue.

\medterm{Kocher-Debre-Semelaigne Syndrome} A rare pediatric clinical syndrome associated with chronic, untreated congenital or juvenile hypothyroidism, characterized by generalized muscular pseudohypertrophy ('infant Hercules' phenotype), muscle stiffness, sluggish deep tendon reflexes, and myxedema.
\synonyms
Hypothyroid Muscular Pseudohypertrophy; Debré-Sémélaigne Syndrome

\medterm{Kocher-Debré-SéMéLaigne Syndrome} A rare condition in infants or young children with severe, untreated underactive thyroid disease. Muscles may look unusually large but are weak and slow, with delayed growth and development, puffy facial features, and other signs of hypothyroidism. Thyroid-hormone treatment usually improves the muscle changes and weakness.

\synonyms
Cretinism with Muscular Hypertrophy; Infantile Hypothyroid Myopathy

\medterm{Kocher Incision} An open surgical subcostal oblique abdominal incision made parallel to and approximately 2 to 3 cm below the right costal margin, traditionally utilized for open cholecystectomy and biliary tract surgical access.

\medterm{Kocher Maneuver} A surgical technique in which the peritoneum along the lateral border of the duodenum is sharply incised, mobilizing the duodenum and head of the pancreas medially and anteriorly to gain access to retroperitoneal structures including the inferior vena cava, aorta, and posterior pancreas.
\synonyms
Duodenal Mobilization; Kocherization

\medterm{Koch's Bacillus} The bacterium that causes tuberculosis, an infection most often affecting the lungs but capable of involving many organs.

\synonyms
Tubercle Bacillus; M. tuberculosis

\medterm{Koebnerisation} Development of lesions of an existing skin disease at sites of trauma or injury, classically seen in psoriasis, vitiligo, and lichen planus.

\medterm{Koebner Phenomenon} Development of lesions of a skin disease, classically psoriasis or vitiligo, at sites of injury or other trauma.

\medterm{Kohler Disease} A rare, self-limiting childhood osteochondrosis characterized by avascular necrosis of the tarsal navicular bone, occurring predominantly in boys aged 3 to 7 years, presenting with antalgic limp, medial midfoot pain, and localized dorsal swelling.

\synonyms
Kohler Bone Disease

\medterm{KOH Preparation} A test in which a skin, hair, or nail sample is treated with potassium hydroxide and viewed under a microscope to look for fungi. A negative result does not always rule out infection if the sample is small or treatment has already begun.

\medterm{Koilocyte} A cell with changes often seen after infection with human papillomavirus (HPV). Under a microscope, it has an enlarged or irregular nucleus and a clear space around the nucleus. The finding supports HPV-related change but is not enough by itself to determine the diagnosis.

\medterm{Koilonychia} Thin, brittle nails that become concave and curve upward at the edges, giving them a spoon-like shape. It is often linked to iron deficiency but can also result from poor nutrition, repeated injury, inherited conditions, or other illness; new or widespread changes need evaluation.

\medterm{Kombucha} A fermented tea made with a culture of bacteria and yeast. It may contain sugar, acids, caffeine, and a small amount of alcohol; contaminated or excessive products can cause illness, especially in vulnerable people.

\medterm{Konno Procedure} Open-heart surgery to widen a severely narrowed path through which blood leaves the left ventricle. The surgeon enlarges the area below and around the aortic valve, often replacing or repairing the valve at the same time. It is used for complex or recurrent obstruction and carries the usual risks of major heart surgery.

\synonyms
Konno-Rastan Procedure; Aortoventriculoplasty

\medterm{Koplik Spots} Small white or bluish spots on a red background inside the cheek, usually opposite the back teeth. They are characteristic of measles and often appear shortly before the skin rash.

\medterm{Koplik's Spots} Small white or bluish-gray spots on the inner cheek, often with a red base, that are characteristic of measles and may appear shortly before the skin rash.

\medterm{Korotkoff Sounds} Sounds heard through a stethoscope over an artery as a blood-pressure cuff is deflated; their appearance and disappearance are used to estimate systolic and diastolic pressure.

\medterm{Korsakoff's Syndrome} A chronic amnestic disorder, commonly associated with severe thiamine deficiency and alcohol-use disorder, causing impaired new learning, memory loss, and sometimes confabulation.

\medterm{Korsakoff Syndrome} A chronic amnestic disorder caused most often by severe thiamine deficiency associated with alcohol-use disorder. It causes impaired formation of new memories, loss of past memories, and sometimes confabulation.

\medterm{Kostmann Syndrome} An autosomal recessive form of severe congenital neutropenia caused by mutations in the HAX1 or ELANE genes, characterized by early arrest of granulopoiesis at the promyelocyte stage in the bone marrow, profound absolute neutropenia (less than 200/uL), and life-threatening bacterial infections starting in early infancy.
\synonyms
Severe Congenital Neutropenia Type 1; Infantile Genetic Agranulocytosis

\medterm{Koumpounophobia} An intense fear or strong distress related to buttons. It is a specific phobia when it causes significant anxiety or interferes with daily life, and it can be treated with psychological therapies such as exposure-based therapy.

\medterm{Kounis Syndrome} A heart attack or temporary narrowing of a heart artery triggered by a severe allergic reaction. It may cause chest pain together with rash, swelling, wheezing, or low blood pressure.

\medterm{KP} Abbreviation; see \textit{Keratosis Pilaris}.

\medterm{Krabbe Disease} A rare inherited disorder in which missing enzymes damage myelin, the protective covering of nerves. Early disease can cause irritability, stiffness, seizures, loss of skills, and progressive problems with movement, vision, or development.

\medterm{Krabbe's Disease} A rare inherited disorder in which missing enzymes cause myelin loss, progressively damaging nerves and producing weakness, stiffness, vision problems, or developmental decline.

\medterm{KRAS Gene} A gene that helps control cell growth and division. Certain changes can keep its growth signal switched on, contributing to cancers and affecting which targeted treatments are likely to work.

\medterm{Kraurosis} An older term for atrophic and sclerotic vulvar disease, often referring to lichen sclerosus and associated with itching, discomfort, and tissue fragility.

\medterm{Kraurosis Vulvae} Advanced progressive atrophy, shrinkage, and sclerosis of the female external genitalia occurring in chronic untreated vulvar lichen sclerosus, characterized by resorption of the labia minora, narrowing of the vaginal introitus, and thin parchment-like epidermal fragility.
\synonyms
Vulvar Lichen Sclerosus; Atrophic Vulvitis

\medterm{Krebs Cycle} A series of chemical reactions in mitochondria that oxidizes acetyl-CoA and generates energy-rich molecules used for cellular energy production.

\synonyms
Tricarboxylic Acid Cycle; TCA Cycle

\medterm{Kringles} Small folded parts of certain proteins that help them attach to other molecules. They are found in proteins involved in blood clot breakdown and in some growth signals.

\medterm{Krukenberg Spindle} A vertical, brownish line of pigment on the inner surface of the cornea. It may occur with pigment dispersion in the eye and can be associated with changes in eye pressure.
\synonyms
Endothelial Pigment Spindle; Pigment Dispersion Corneal Spindle

\medterm{Krukenberg's Tumor} Alternate spelling; see \textit{Krukenberg Tumor}.

\medterm{Krukenberg Tumor} An ovarian tumor caused by cancer that has spread from another organ, most often the stomach or digestive tract. It may cause abdominal swelling, pain, menstrual changes, or digestive symptoms.

\medterm{Kt/V} A number used to estimate how much urea a dialysis treatment removes. It combines the dialyzer's cleaning ability and treatment time with the person's body-water volume; clinicians use it with symptoms, fluid balance, laboratory results, and the dialysis schedule.

\medterm{Kugelberg-Welander Disease} A milder form of inherited spinal muscular atrophy, also called type 3 spinal muscular atrophy. It usually begins after 18 months of age and causes slowly worsening weakness in muscles close to the body, especially the hips and thighs. Children may have a waddling walk or trouble climbing stairs. Many walk independently, and life expectancy is often near normal with current care.

\synonyms
Spinal Muscular Atrophy Type 3 (SMA 3); Juvenile Spinal Muscular Atrophy

\medterm{Kulchitsky Cells} Specialized hormone-producing cells found mainly in the digestive and respiratory systems. They release chemical signals, including serotonin, that help regulate movement, secretion, and other organ functions.

\synonyms
EC Cells; Argentaffin Cells

\medterm{Kupffer Cells} Immune cells lining the small blood channels in the liver. They remove old blood cells, particles, and some germs from blood coming from the digestive organs.

\synonyms
Stellate Macrophages of the Liver

\medterm{Kupffer's Cell} A Kupffer cell is a resident macrophage lining the liver sinusoids; it removes microbes, damaged blood cells, and particles from portal blood and contributes to hepatic inflammation and iron handling.

\medterm{Kuru} Kuru is a rare fatal prion disease historically associated with ritual cannibalism, causing progressive ataxia, tremor, and neurologic decline.

\medterm{Kussmaul Breathing} Deep, rapid, labored breathing that helps remove extra carbon dioxide when the blood is dangerously acidic. It can occur in diabetic ketoacidosis, severe kidney failure, or other illnesses.

\medterm{Kussmaul Respiration} Alternate name; see \textit{Kussmaul Breathing}.

\medterm{Kussmaul Sign} A rise or failure to fall in the pressure of the large neck veins during breathing in. It suggests that the right side of the heart cannot fill normally and may occur with severe heart failure, constrictive pericarditis, or other conditions affecting right-heart filling.

\medterm{Kveim-Siltzbach Test} An old skin test once used to support the diagnosis of sarcoidosis. Material from sarcoid tissue was injected under the skin and the area was later biopsied. The test is rarely used now because of safety, availability, and accuracy concerns.
\synonyms
Kveim Test; Sarcoidosis Skin Test

\medterm{Kwashiorkor} Severe malnutrition in which too little protein and other nutrients cause swelling, poor growth, muscle loss, skin or hair changes, and weaker immunity. It can occur in children or adults.

\medterm{Kyasanur Forest Disease} An acute tick-borne viral hemorrhagic fever caused by Kyasanur Forest disease virus, a flavivirus transmitted to humans and non-human primates primarily by Haemaphysalis ticks (Haemaphysalis spinigera) in the Western Ghats region of South Asia. Following an incubation period of 3 to 8 days, it presents with sudden onset of high fever, frontal headache, conjunctival injection, severe prostration, myalgia, and hemorrhagic manifestations (such as epistaxis, hematemesis, or melena), frequently followed by a secondary biphasic stage marked by meningoencephalitis and tremors.

\synonyms
KFD; Monkey Fever; Kyasanur Forest Disease Virus Infection

\medterm{Kynurenine} A substance made when the body breaks down the amino acid tryptophan. It is involved in chemical signaling in the brain and immune system and may change during infection, inflammation, or some diseases.

\medterm{Kyphoplasty} A minimally invasive procedure that uses a balloon to create space in a collapsed vertebral body and then injects bone cement to stabilize a painful compression fracture.

\medterm{Kyphoscoliosis} A combination of sideways curvature and excessive forward rounding of the spine. Severe deformity can reduce chest movement and lung capacity, causing breathlessness or, over time, breathing failure and high blood pressure in the lung arteries.

\medterm{Kyphosis} An exaggerated forward curve of the upper spine, sometimes producing a rounded back. Causes include posture, osteoporosis and spinal fractures, growth disorders, arthritis, or inflammation. It may cause pain, weakness, or worsening curvature.

\synonyms
Round Back

\medterm{Kyrle Disease} A perforating dermatosis characterized by large, intensely itchy, keratotic papules that expel abnormal dermal material through the epidermis. It is strongly associated with diabetes and chronic kidney disease and may improve when associated illness is controlled.

\medterm{Kytococcus} A genus of gram-positive cocci found on skin and in the environment that can rarely cause opportunistic infection.

\medterm{Kytococcus Sedentarius} A skin-associated gram-positive coccus that produces enzymes linked to pitted keratolysis and can rarely cause invasive infection.


\medterm{L-Dopa} A medicine that the brain converts into dopamine. It is usually combined with carbidopa or benserazide to treat Parkinson disease.
\medterm{La Crosse Encephalitis} A mosquito-borne viral encephalitis caused by La Crosse virus, usually affecting children and causing fever, headache, seizures, or altered consciousness.
\medterm{LAAO} Alternate name; see \textit{Left Atrial Appendage Occlusion}.

\medterm{Labetalol} A medicine that lowers blood pressure by slowing the heart and relaxing blood vessels. It is also used for severe high blood pressure during pregnancy and may cause a slow heartbeat, low blood pressure, or breathing problems in people with asthma.

\medterm{Labia} The folds of skin surrounding the entrance to the vagina and urethra, consisting mainly of the labia majora and labia minora.

\medterm{Labia Majora} The paired outer folds of skin of the vulva that help protect the urethral and vaginal openings and contain connective tissue, glands, and fat.

\medterm{Labia Minora} The paired inner folds of tissue of the vulva, lying within the labia majora and surrounding the vestibule and its openings.

\medterm{Labial Sounds} Speech sounds made with one or both lips, such as the sounds in “p,” “b,” and “m.” Difficulty making them may occur with problems affecting the lips, mouth, nerves, or speech muscles.

\medterm{Labiaplasty} Surgical reshaping or reduction of the labia, usually the labia minora, performed for functional symptoms, reconstruction, or selected cosmetic indications. Risks include pain, bleeding, infection, scarring, and altered sensation.

\medterm{Labile} Easily or rapidly changing, such as blood pressure, blood sugar, mood, or symptoms. Labile measurements may reflect illness, medicines, stress, hormones, or changing treatment and often require repeated observation.

\medterm{Labile Hypertension} Blood pressure that changes substantially over short periods. Stress, pain, medicines, illness, and measurement conditions can contribute, so repeated properly performed readings help confirm the pattern.

\medterm{Lability} A tendency to change easily or unpredictably, such as rapidly shifting emotions or blood pressure. Its medical importance depends on the symptom, severity, and underlying cause.
\medterm{Labium} A lip-like anatomical structure; in obstetrics and gynecology, usually one of the labia of the vulva.
\medterm{Labium Majus} One of the paired outer folds of the vulva containing skin, connective tissue, glands, vessels, and nerves and protecting the more internal genital structures.

\medterm{Labium Minus} One of the paired inner folds of the vulva bordering the vestibule and helping protect the urethral and vaginal openings.

\medterm{Labor} The process in which regular contractions open the cervix and lead to delivery of the baby and placenta. It includes the stages of cervical opening, birth of the baby, and delivery of the placenta.

\medterm{Labor: Signs And Symptoms} Labor is the process of regular uterine contractions that opens the cervix and leads to birth. Possible signs include increasingly regular contractions, rupture of the membranes, a blood-stained mucus discharge, pelvic pressure, and back pain. Heavy bleeding, reduced fetal movement, severe pain, or suspected complications requires urgent assessment.

\medterm{Labor Complication} A maternal or fetal problem occurring during labor or delivery, such as abnormal labor progress, hemorrhage, infection, fetal distress, or malpresentation, requiring assessment and management based on the cause.

\medterm{Labor Dystocia} Labor that is unusually slow or difficult because of weak contractions, the baby's position or size, or the shape of the mother's pelvis. Treatment depends on the cause and may include medicines or an operation.

\medterm{Labor False (Braxton Hicks Contractions)} Alternate index label for Braxton Hicks Contractions.

\medterm{Labor Pain} Pain caused by uterine contractions, cervical dilation, and stretching of pelvic tissues during labor; its intensity and management vary with labor stage and individual factors.

\medterm{Laboratory} A facility or department where clinical specimens are examined using microbiologic, hematologic, chemical, molecular, or other diagnostic methods.

\medterm{Laboratory Medicine} The medical specialty that uses laboratory tests on blood, tissue, and other specimens for diagnosis, screening, monitoring, and treatment.
\medterm{Laboratory Procedure} A method used to collect, prepare, examine, or measure a patient sample such as blood, urine, tissue, or a swab. Results can help diagnose disease or guide treatment and must be interpreted with the clinical findings.

\medterm{Laboratory Test} An examination of a body fluid, tissue, or other specimen to obtain information for diagnosis, monitoring, screening, or treatment.

\medterm{Labour} The process of childbirth, consisting of regular uterine contractions that cause
progressive cervical effacement and dilation, leading to delivery of the baby and
placenta.

\medterm{Labra} The plural of labrum: tough rings of tissue around certain joint sockets, especially the shoulder and hip. They deepen the socket, improve stability, and help distribute pressure during movement.

\medterm{Labral Tear} A tear in the fibrocartilaginous rim of a joint socket, commonly the shoulder or hip, causing pain, clicking, instability, or reduced range of motion.

\medterm{Labrum} A ring of tough, flexible tissue around the socket of a joint. It deepens and stabilizes the shoulder or hip socket.

\medterm{Labyrinth} The inner-ear system that contains the cochlea for hearing and the balance organs that sense head movement and position. Disease or injury can cause dizziness, unsteadiness, ringing, or hearing loss.

\medterm{Labyrinth Disorder} A problem affecting inner-ear hearing or balance structures, potentially causing spinning sensations, nausea, unsteadiness, hearing changes, or abnormal eye movements.

\medterm{Labyrinthine Membrane} The delicate, fluid-containing system inside the inner ear that contains the organs for hearing and balance. Injury or abnormal fluid pressure can cause dizziness, hearing loss, or ringing in the ear.

\medterm{Labyrinthitis} Inflammation of the inner ear that can cause sudden spinning dizziness, imbalance, nausea, ringing in the ear, and sometimes hearing loss.

\medterm{Labyrinthitis (Inner Ear Inflammation)} Alternate index label for Inner Ear Inflammation.


\medterm{Laceration} A torn wound caused by blunt or irregular force. It may involve only the skin or extend into fat, muscle, tendons, nerves, or blood vessels; deep, gaping, dirty, or heavily bleeding wounds need medical care.

\medterm{Laceration (Forensic)} A torn wound caused by blunt force, usually with irregular bruised edges and tissue strands crossing the wound. These features help distinguish it from a clean cut caused by a sharp object.

\synonyms
Forensic

\medterm{Laceration - Sutures Or Staples - At Home} Home care after sutures or staples includes keeping the wound clean and protected, following dressing instructions, avoiding strain, and watching for infection or wound separation.

\medterm{Lacerations - Liquid Bandage} Liquid bandage is a film-forming adhesive applied over small, clean, shallow cuts to protect the wound; it should not be used on deep, infected, heavily bleeding, or gaping wounds.

\medterm{Lachman Test} A knee examination in which the lower leg is gently pulled forward while the knee is slightly bent. Too much movement may suggest an injury to the anterior cruciate ligament.

\medterm{Lachrymal} Relating to tears or the structures that produce and drain them; also spelled lacrimal.
\medterm{Lacosamide} An antiseizure medicine used mainly for focal seizures. It may cause dizziness, double vision, poor coordination, or changes in the heart's electrical rhythm.

\medterm{Lacquer Poisoning} Harm from swallowing or breathing lacquer or its solvents. It may irritate the mouth and lungs, cause dizziness or drowsiness, and cause severe lung injury if liquid is inhaled; contact poison control after significant exposure.

\medterm{Lacrimal} Relating to tears or the lacrimal glands and drainage passages.
\medterm{Lacrimal Apparatus} The glands and passages that make, spread, and drain tears. It includes the tear gland, small openings near the inner eyelids, the tear sac, and the duct that drains tears into the nose.

\medterm{Lacrimal Disease} A disorder of the tear-producing or tear-draining system, including dry eye, dacryocystitis, obstruction, inflammation, or abnormal tearing.

\medterm{Lacrimal Duct} A drainage channel that carries tears from the lacrimal sac into the nasal cavity. Blockage can cause
watering of the eye, recurrent infection, swelling near the inner corner, or discharge.

\medterm{Lacrimal Gland} One of two glands above the outer part of the eye that make the watery part of tears. The tears spread across the eye and then drain through passages near the inner corner.

\medterm{Lacrimal Gland Tumor} A benign or malignant growth arising in the tear-producing lacrimal gland, often causing swelling, pain, or displacement of the eye.

\medterm{Lacrimation} Production and release of tears by the lacrimal glands; excessive tearing may result from irritation, allergy, infection, or blocked tear drainage.

\medterm{Lactam} A ring-shaped chemical structure found in some medicines. Beta-lactam antibiotics use this structure to weaken bacterial cell walls, but allergies and bacterial resistance can limit their use.

\medterm{Lactase} An enzyme in the small intestine that breaks lactose, the natural sugar in milk, into absorbable sugars. Low lactase can cause bloating, cramps, gas, diarrhea, or discomfort after dairy foods.

\medterm{Lactase Deficiency} Alternate index label; see \textit{Lactose intolerance}.

\medterm{Lactate} A product of cellular metabolism measured in blood. An increased level may occur with strenuous exercise, impaired oxygen delivery, infection, medicines, liver disease, or other conditions and must be interpreted clinically.

\medterm{Lactate Dehydrogenase} An enzyme found in many body tissues. Its blood level may rise when cells are damaged, but the test does not identify the cause by itself.

\medterm{Lactate Dehydrogenase Test} A blood test measuring lactate dehydrogenase, an enzyme released when cells are injured; increased activity is nonspecific and may occur with hemolysis, liver disease, muscle injury, infection, or malignancy.

\medterm{Lactation} The production and release of breast milk after childbirth. Hormones stimulate the breasts to make milk and release it during feeding or expression.

\medterm{Lactation (Breastfeeding)} The production and release of breast milk after childbirth. Early milk, called colostrum, is followed by mature milk; hormones help make and release the milk.

\synonyms
Breastfeeding

\medterm{Lactation Disorder} A problem with making, releasing, or transferring breast milk after childbirth. Examples include low supply, excess milk, painful engorgement, mastitis, or difficulty with the baby’s latch; assessment considers both parent and infant.

\medterm{Lactation Mastitis} Alternate index label; see \textit{Mastitis}.

\medterm{Lactation Nutrition} Eating and drinking enough during breastfeeding to support the parent’s health and milk production. Most people need a varied diet and regular fluids; supplements may be needed for a diagnosed deficiency or specific dietary pattern.

\medterm{Lactational Amenorrhea} The temporary absence of menstrual periods during breastfeeding because ovulation is suppressed. Breastfeeding prevents pregnancy reliably only when specific conditions are met, so other contraception may be needed.

\medterm{Lacteals} Tiny lymph vessels in the lining of the small intestine that absorb most digested fats and fat-soluble vitamins. They carry these nutrients through the lymph system before they enter the bloodstream.

\medterm{Lactic Acid} An organic acid produced during normal metabolism, especially when pyruvate is converted to lactate. Blood lactate may rise for many reasons, including impaired oxygen delivery, strenuous exercise, or medicines.

\medterm{Lactic Acid Test} A blood test measuring lactate, which accumulates when production exceeds clearance; it helps assess tissue hypoperfusion, sepsis, severe exercise, metabolic disease, or other causes of lactic acidosis.

\medterm{Lactic Acidosis} Metabolic acidosis associated with increased lactate production or reduced lactate clearance. Causes include impaired tissue oxygen delivery, medicines, liver disease, seizures, and other disorders.

\medterm{Lactobacillus} A genus of bacteria that produces lactic acid and normally forms part of the healthy vaginal and intestinal microbiota; some species are used as probiotics.

\medterm{Lactoferrin} Lactoferrin is an iron-binding glycoprotein in milk and mucosal secretions that limits microbial iron access and contributes to innate immune defense.

\medterm{Lactose} The principal sugar in milk, composed of glucose and galactose. Lactase in the small intestine digests it.

\medterm{Lactose Fermentation} The ability of bacteria to break down lactose and produce acid or gas. It is used in the laboratory to help identify different types of bacteria.

\medterm{Lactose Intolerance} A condition in which the body cannot fully digest lactose, the sugar in milk and dairy products, because it does not make enough lactase. It may cause bloating, gas, abdominal pain, nausea, or diarrhea after consuming lactose.

\synonyms
Lactase Deficiency
Milk Intolerance
Milk Intolerance (Lactose Intolerance)

\medterm{Lactose Tolerance Test} A test for difficulty digesting lactose, the sugar in milk. After lactose is consumed, blood glucose or breath hydrogen is measured; symptoms and results help assess lactose malabsorption but do not diagnose every cause of digestive discomfort.

\medterm{Lactose Tolerance Tests} Tests that evaluate digestion and absorption of lactose by measuring blood glucose response after lactose ingestion or detecting hydrogen in exhaled breath; they may support evaluation of lactose malabsorption.

\medterm{Lactulose} A medicine that softens stool by drawing water into the bowel. It is also used to lower harmful ammonia levels in some people with severe liver disease. Gas and diarrhea can occur.

\medterm{Lacuna} A small space or cavity in tissue, such as a space containing an osteocyte in bone.
\medterm{Lacunar Stroke} A small ischemic stroke caused by blockage of a tiny artery deep in the brain, often producing weakness, numbness, clumsiness, or speech problems.

\medterm{Ladd Bands} Abnormal bands of tissue present at birth that can cross and squeeze the first part of the small intestine when the bowel is in an abnormal position. They may cause green vomiting and blockage in an infant and are released during the Ladd procedure.

\synonyms
Ladd's Peritoneal Bands; Ladd Fibrous Bands

\medterm{Ladd Procedure} An operation for a child whose intestines developed in an abnormal position and may twist. The surgeon untwists the intestine, releases bands causing blockage, and broadens its attachment to reduce future twisting.

\synonyms
Ladd Operation; Ladd's Procedure

\medterm{Laetrile} A semisynthetic amygdalin-related compound promoted as an alternative cancer treatment; credible evidence has not established anticancer benefit, and cyanide poisoning is a serious risk.

\medterm{Lafora Disease} A rare inherited epilepsy that usually begins during adolescence and progressively worsens. Seizures, brief muscle jerks, visual symptoms, memory problems, and loss of physical or thinking abilities may occur.

\medterm{Lagophthalmos} Inability to close the eyelids completely. Exposed cornea can become dry, irritated, scratched, infected, or damaged, especially during sleep, so lubricating treatment and protection may be needed.

\medterm{Lambda} A landmark on the back of the skull where two skull seams meet. It is used to describe head shape and the position of nearby structures, especially in newborn examinations and imaging.
\medterm{Lambert-Eaton Myasthenic Syndrome} A rare condition in which the immune system interferes with nerve signals to muscles. It causes weakness, especially in the hips and shoulders, reduced reflexes, and sometimes dry mouth. It can be associated with small-cell lung cancer.

\medterm{Lamellar} Made of thin layers or plates. The word may describe the organized layers of mature bone, the skin barrier, or structures in the eye and nervous system. Its exact meaning depends on the body structure being described.
\medterm{Lamellar Body} A small storage structure inside certain cells that contains layered fats. In lung cells, it stores and releases surfactant, the substance that helps keep the tiny air sacs open.

\medterm{Lamellar Bone} Mature bone made of organized layers of mineralized tissue. It forms most adult compact and spongy bone and provides strength while being continually remodeled.

\medterm{Lamellar Ichthyosis} An inherited skin disorder present from birth that causes dry skin covered with large, dark, plate-like
scales over much of the body, including the arms and legs. It may reduce sweating and cause overheating,
and can affect the hair, nails, eyes, or palms and soles.

\medterm{Lamin} A structural protein forming part of the nuclear lamina, which supports the inner surface of the cell nucleus.

\medterm{Lamina} A flat piece of bone that forms part of the back of a vertebra. The two laminae help complete the bony arch around the spinal cord and can be removed during some spine operations.

\medterm{Lamina Propria} A thin layer of supporting tissue beneath the inner lining of organs such as the intestine, airways, and bladder. It contains small blood vessels, immune cells, glands, and connective tissue.

\medterm{Laminectomy} An operation that removes part of a spinal bone to make more room for the spinal cord or nerves. It may be used for spinal narrowing, a tumor, or selected disk problems.

\medterm{Laminin} A protein in the supporting framework around cells that helps cells attach to underlying tissue. It contributes to tissue structure, development, repair, and the filtering barrier in some organs.

\medterm{Lamivudine} Lamivudine is a nucleoside reverse-transcriptase inhibitor used with other antiretrovirals for HIV and as an antiviral for chronic hepatitis B; resistance and rare lactic acidosis or hepatic complications require monitoring.

\medterm{Lamotrigine} Lamotrigine is an antiseizure medicine that inhibits voltage-gated sodium channels and is also used for bipolar-maintenance treatment; rapid dose escalation increases the risk of serious rash including Stevens--Johnson syndrome.

\medterm{Lancet} A small, sterile, usually single-use device with a sharp point used to puncture the skin and obtain a capillary-blood sample. Lancets are commonly used for fingertip glucose testing and other point-of-care blood tests; used lancets must be discarded in an appropriate sharps container.

\synonyms
Blood lancet
Fingerstick lancet

\medterm{Landolt Ring} A ring-shaped symbol with a gap used to test visual sharpness. The person identifies which direction the gap faces; smaller symbols indicate better visual acuity.

\medterm{Langer Lines} Natural lines of skin tension that reflect the predominant orientation of collagen fibers in the dermis and may influence surgical scar direction.

\medterm{Langerhans Cell Histiocytosis} A rare disorder in which abnormal immune cells build up in one or more organs, commonly the bones, skin, lungs, lymph nodes, or pituitary gland. Symptoms and treatment depend on which organs are affected.
\medterm{Langerhans Cells} Immune cells in the outer skin that capture foreign substances and help start immune responses. They connect skin defenses with nearby lymph tissue and can contribute to allergic or inflammatory reactions.

\medterm{Langerhans' Cell} An immune cell in the skin that captures foreign substances and helps activate T-cell responses.

\medterm{Langerhans-Cell Histiocytosis} A rare disorder in which abnormal immune cells collect in one or more organs, especially bone, skin, lungs, lymph nodes, or the pituitary gland. Symptoms depend on the organs involved and may include bone pain, rash, cough, or hormone problems.

\medterm{Language Barrier} Difficulty communicating because a patient and healthcare professional do not share a language well enough for safe understanding. A trained interpreter helps explain symptoms, consent, instructions, and risks.

\medterm{Language Development} The gradual development of understanding, sounds, words, grammar, and communication during childhood. Delays may be related to hearing loss, developmental conditions, neurologic disease, or limited language exposure.
\medterm{Language Disorder} A developmental or acquired impairment in understanding, producing, or using spoken, written, or signed language that affects communication and daily function.


\medterm{Language Therapy} Treatment that improves understanding, expression, reading, writing, or other language skills after
developmental or acquired impairment.

\medterm{Lanolin} Lanolin is a waxy substance derived from sheep wool used as an emollient and in topical products; allergy or irritation can occur.

\medterm{Lanolin Poisoning} Swallowing a small amount of lanolin, a wool-derived moisturizer, usually causes little more than stomach upset. Large amounts or products with added medicines can be more dangerous; contact poison control for symptoms or uncertain exposure.

\medterm{Lansoprazole} A medicine that strongly reduces stomach acid. It is used for heartburn, ulcers, and some H. pylori infections; long-term use may affect mineral or vitamin levels and infection risk.

\medterm{Lanthanum} A chemical element whose carbonate salt can bind dietary phosphate and lower phosphate absorption in people with advanced kidney disease.

\medterm{Lanugo} Fine, soft hair that grows all over the body of a fetus and is typically shed before
birth.

\medterm{Laparoscope} A slender instrument with a light and camera used to view the abdominal or pelvic cavity through a small incision during minimally invasive surgery.

\medterm{Laparoscopic Cholecystectomy} Surgical removal of the gallbladder through small abdominal incisions using a laparoscope and instruments, commonly performed for symptomatic gallstones or gallbladder inflammation.

\medterm{Laparoscopic Colectomy} Laparoscopic colectomy is removal of part or all of the colon through small incisions using a camera and specialized instruments.

\medterm{Laparoscopic Gallbladder Removal} Surgery to remove the gallbladder through several small abdominal openings using a camera and instruments. It commonly treats symptomatic gallstones or gallbladder inflammation; bleeding, infection, bile leakage, and injury to bile ducts are possible.

\medterm{Laparoscopic Gastric Banding} A bariatric operation that places an adjustable band around the upper stomach to create a small pouch and restrict food passage; it is less commonly performed than other bariatric procedures.

\medterm{Laparoscopic Hysterectomy} Removal of the uterus using laparoscopic instruments through small abdominal incisions. The cervix, fallopian tubes, and ovaries may or may not also be removed; the approach and extent depend on the indication, anatomy, and surgical plan.



\medterm{Laparoscopy} A type of minimally invasive surgery in which a thin camera is placed through a small cut in the abdomen. Other small cuts may be used to insert surgical instruments.

\medterm{Laparotomy} An operation that uses a larger cut in the abdomen to reach the abdominal or pelvic organs.

\medterm{Lapatinib} Lapatinib is a targeted cancer medicine that blocks HER2 and epidermal-growth-factor receptors and is used for selected advanced breast cancers.

\medterm{Lapatinib Ditosylate} Lapatinib ditosylate is the salt form of lapatinib, a targeted medicine used for selected HER2-positive breast cancers.

\medterm{Large Airway Obstruction} A blockage in the windpipe or large breathing tubes. It can cause noisy breathing, shortness of breath, and difficulty clearing mucus and may be caused by narrowing, a tumor, or weakness of the airway wall.

\medterm{Large Bowel Obstruction} A blockage in the colon, commonly caused by cancer, twisting of the bowel, scar tissue, or diverticular disease. It can cause increasing swelling, cramping, vomiting, and inability to pass stool or gas and may require urgent treatment.

\medterm{Large Bowel Resection} Surgical removal of a segment of the colon or rectum, followed by reconnection of the bowel or creation of a stoma when necessary to treat cancer, obstruction, ischemia, inflammation, or other disease.


\medterm{Large Colon} The large colon is the main portion of the large intestine that absorbs water, stores fecal material, and moves contents toward the rectum.

\medterm{Large For Gestational Age} A fetus or newborn larger than expected for the number of weeks of pregnancy, usually above the 90th percentile. Maternal diabetes, family size, and excess pregnancy weight gain can contribute and may increase birth complications.

\medterm{Large Intestine} The final part of the digestive tract, extending from the small intestine to the rectum. It absorbs water and helps form and move stool.

\medterm{Large Solitary Luteinized Follicle Cyst Of Pregnancy} A usually harmless ovarian cyst that develops during pregnancy as hormone-sensitive ovarian tissue enlarges. It may be found by chance or as a pelvic mass and often settles after pregnancy.


\medterm{Laron Syndrome} A rare inherited growth disorder in which the body cannot respond properly to growth hormone. It causes severe short stature, characteristic facial features, and very low levels of the growth signal needed for normal bone development.


\medterm{Larva Currens From Strongyloides} A rapidly moving, itchy skin rash caused by Strongyloides larvae moving through the skin. It often appears around the buttocks or trunk and may indicate ongoing intestinal infection requiring antiparasitic treatment.

\medterm{Larva Migrans} A condition caused by migrating larval parasites in human tissues, including cutaneous larva migrans in the skin and visceral larva migrans in internal organs.

\medterm{Laryngeal} Relating to the larynx, the organ in the throat containing the vocal cords and helping protect the airway.

\medterm{Laryngeal Cancer} Cancer of the larynx, or voice box. Tobacco, alcohol, and some HPV infections increase risk. Persistent hoarseness, painful or difficult swallowing, noisy breathing, or a neck lump needs prompt examination.

\medterm{Laryngeal Candidiasis} A yeast infection of the voice box, usually causing hoarseness, throat pain, or difficulty swallowing. It is more likely with weakened immunity or inhaled steroid use and may require examination of the larynx and antifungal treatment.

\medterm{Laryngeal Carcinoma} Laryngeal carcinoma is cancer arising in the voice box and may cause hoarseness, swallowing difficulty, pain, or breathing problems.

\medterm{Laryngeal Cartilage} Laryngeal cartilage forms the structural framework of the voice box, maintains the airway, and supports vocal-fold movement.

\medterm{Laryngeal Cleft} A laryngeal cleft is a congenital opening between the airway and esophagus that can allow aspiration and cause feeding or breathing problems.

\medterm{Laryngeal Edema} Swelling of the tissues of the larynx that can narrow the upper airway and cause hoarseness, stridor, or respiratory distress. It may result from allergy, infection, trauma, or airway instrumentation.

\medterm{Laryngeal Fistula} An abnormal passage between the voice box and nearby tissue or the swallowing tract. It can allow food or saliva to enter the airway and cause coughing, repeated chest infections, or breathing problems.

\medterm{Laryngeal Hemorrhage} Bleeding within or around the larynx, which may result from trauma, instrumentation, vocal strain, vascular lesions, infection, or tumor and can threaten the airway.

\medterm{Laryngeal Masks} Laryngeal masks are supraglottic airway devices placed above the vocal cords to maintain ventilation during anesthesia or emergencies.

\medterm{Laryngeal Mucositis} Inflammation and open sores in the lining of the voice box, often after radiation, chemotherapy, infection, or irritation. It can cause hoarseness, throat pain, painful swallowing, or breathing difficulty.

\medterm{Laryngeal Nerve} A branch of the vagus nerve supplying the larynx; injury can cause hoarseness, swallowing difficulty, or breathing problems.

\medterm{Laryngeal Nerve Damage} Injury to the recurrent or superior laryngeal nerves causing hoarseness, vocal-cord
weakness, swallowing difficulty, or airway compromise. Common causes include thyroid or
neck surgery, malignancy, trauma, and compression along the nerve pathway.

\medterm{Laryngeal Obstruction} Partial or complete blockage of the laryngeal airway caused by edema, foreign body, tumor, trauma, infection, vocal-cord dysfunction, or other structural disease.

\medterm{Laryngeal Prosthesis} A device used to help restore speech after removal of the larynx, commonly through a one-way valve placed between the trachea and esophagus.

\medterm{Laryngeal Stridor} A harsh, high-pitched noise caused by turbulent airflow through a narrowed upper airway, often most noticeable during inspiration.

\medterm{Laryngeal Warts (Recurrent Respiratory Papillomatosis)} Benign papillomatous growths caused by HPV in the larynx or respiratory tract. They may cause hoarseness, stridor, or airway obstruction.

\synonyms
Recurrent Respiratory Papillomatosis

\medterm{Laryngeal Web} A laryngeal web is a thin membrane across part of the airway, usually congenital or acquired after injury, causing hoarseness or airway narrowing.

\medterm{Laryngectomee} A person who has undergone laryngectomy, surgical removal of all or part of the larynx; total laryngectomy separates the airway from the mouth and nose.

\medterm{Laryngectomy} Surgery to remove part or all of the larynx, usually for cancer or severe structural disease. After total removal, breathing occurs permanently through an opening in the neck; partial surgery may preserve some voice and swallowing.

\medterm{Laryngectomy Tube} A tube placed in the neck opening after larynx surgery to help keep the airway open and allow suctioning or ventilation when needed. Mucus blockage, bleeding, infection, and tissue overgrowth can occur.

\medterm{Laryngitis} Inflammation of the voice box that causes hoarseness or temporary voice loss. Common causes include viral infection, voice overuse, smoke or other irritation, and reflux.

\medterm{Laryngocele} An air-filled pouch formed when a small space beside the vocal cords becomes enlarged. It may cause hoarseness, a neck swelling that changes with pressure, coughing, or blockage of the airway.

\medterm{Laryngomalacia} A birth condition in which soft tissue above an infant’s vocal cords partly collapses during breathing in, causing noisy breathing. It often improves with growth, but feeding trouble, pauses in breathing, or poor growth needs assessment.

\medterm{Laryngopharyngeal Reflux} Backflow of stomach contents into the throat and voice box. It may cause hoarseness, frequent throat clearing, cough, a lump-like feeling, excess mucus, or throat tightness, often without typical heartburn.

\medterm{Laryngopharynx} The lower part of the throat behind the voice box, leading to the esophagus. It helps direct swallowed food away from the airway and is a site where some throat cancers can develop.
\medterm{Laryngoscope} An instrument used to see the throat and vocal cords, especially while placing a breathing tube. Common blades are curved Macintosh and straight Miller types; use can injure teeth, lips, or throat.

\medterm{Laryngoscopy} An examination of the voice box and nearby throat using a mirror or a thin camera passed through the mouth or nose. It helps investigate hoarseness, swallowing problems, airway narrowing, bleeding, or suspected tumors.

\medterm{Laryngoscopy And Nasolaryngoscopy} An examination of the voice box and nearby airway using a small rigid or flexible camera passed through the mouth or nose. It helps investigate hoarseness, swallowing difficulty, noisy breathing, bleeding, or a suspected growth.

\medterm{Laryngospasm} Sudden tightening of the vocal cords that partly or completely blocks breathing. It can be triggered by saliva, blood, reflux, airway procedures, or anesthesia and may cause noisy breathing, choking, or inability to breathe.

\medterm{Laryngostenosis} Narrowing of the voice box, or larynx, that can restrict breathing. Causes include scarring after intubation, injury, inflammation, tumors, and congenital narrowing; noisy breathing or breathlessness needs medical assessment.

\medterm{Laryngotracheobronchitis (Croup)} Alternate index label for Croup.

\medterm{Larynx} The voice box in the front of the neck between the throat and windpipe. It contains the vocal cords, produces sound, helps keep food and liquid out of the lungs, and regulates airflow during breathing.

\medterm{Larynx Cancer} Laryngeal cancer is a malignant tumor of the voice box and may cause persistent hoarseness, swallowing pain, neck mass, ear pain, cough, or breathing difficulty. Diagnosis requires examination and biopsy; treatment depends on site and stage.

\synonyms
Throat Cancer (Larynx Cancer)

\medterm{Larynx Disease} Any disorder affecting the larynx, including inflammation, infection, structural lesions, vocal-cord dysfunction, neurologic disease, or cancer.

\medterm{Larynx Muscle} A muscle of the voice box that moves the vocal folds and helps protect the airway and produce sound.

\medterm{Larynx Neoplasm} An abnormal growth in the voice box. It may be noncancerous or cancerous and can cause hoarseness, swallowing difficulty, pain, breathing problems, or a neck lump.

\medterm{Laser} A device that produces a concentrated beam of light. Depending on its wavelength and power, it can measure tissue, cut or seal tissue, remove abnormal growths, or treat selected eye and skin conditions.
\medterm{Laser (Vascular)} A laser system used to treat selected blood-vessel problems or tissue blocking a vessel. The energy is delivered through the skin or a catheter and must be carefully targeted to avoid damage to nearby tissue.
\synonyms
Vascular laser

\medterm{Laser Ablation} The use of focused laser energy to heat, destroy, or seal targeted tissue. The technique is used in selected vascular, ophthalmic, dermatologic, and other procedures.

\medterm{Laser Assisted Uvula Palatoplasty} A procedure using laser energy to reshape or remove selected soft-palate and uvula tissue, intended mainly to reduce snoring in carefully selected patients.

\synonyms
LAUP


\medterm{Laser Fiber} A thin optical fiber that carries laser energy into the body during treatment. It can break urinary stones, destroy selected tumors, seal blood vessels, or treat other targeted tissues.

\medterm{Laser Photocoagulation} A laser treatment that seals or destroys leaking or abnormal blood vessels, especially in the retina, helping limit bleeding, swelling, or damaging new vessel growth.

\medterm{Laser Photocoagulation - Eye} An eye treatment that uses laser heat to seal leaking blood vessels or destroy abnormal retinal tissue, helping treat selected retinal diseases.
\medterm{Laser Prostatectomy} A procedure using laser energy to remove or vaporize enlarged prostate tissue and improve urine flow. It is performed through the urethra and may cause bleeding, infection, temporary urgency, or other urinary problems.

\medterm{Laser Surgery} Laser surgery uses a focused beam of light to cut, remove, seal, or reshape tissue in selected procedures.

\medterm{Laser Surgery For The Skin} Use of concentrated light to remove, vaporize, seal, or reshape selected skin tissue. It may treat scars, abnormal blood vessels, unwanted hair, or some lesions; redness, pigment change, burns, scarring, and infection are possible.

\medterm{Laser Therapy} Laser therapy uses controlled light energy to treat selected skin, eye, vascular, dental, or other medical conditions.

\medterm{Laser Therapy For Cancer} Use of an intense, focused beam of light to destroy or shrink selected cancer tissue or control bleeding and blockage. It is useful only for certain tumors and sites; treatment may be combined with surgery, medicines, or radiation.

\medterm{Laser Trabeculoplasty} A laser treatment for some forms of open-angle glaucoma that improves drainage of fluid from the eye and lowers eye pressure. The effect may be temporary, and continued examinations or other treatment may still be needed.

\medterm{Laser-Assisted In Situ Keratomileusis} A refractive surgical procedure that reshapes the cornea using a laser to correct short-sightedness, long-sightedness, or astigmatism, thereby reducing dependence on glasses or
contact lenses.

\synonyms
LASIK

\medterm{LASIK} A laser operation that reshapes the clear front surface of the eye to correct short-sightedness, long-sightedness, or astigmatism. It can reduce dependence on glasses, but dry eye, glare, infection, under-correction, and other vision problems are possible.

\medterm{LASIK Eye Surgery} A laser operation that reshapes the cornea to correct short-sightedness, long-sightedness, or astigmatism and reduce dependence on glasses. Dry eyes, glare, infection, under- or over-correction, and rarely serious vision loss can occur.


\medterm{Lassa Fever} A viral illness that can cause fever, weakness, vomiting, bleeding, and damage to several organs. It is usually spread by contact with urine or droppings of infected rodents in parts of West Africa and sometimes by contact with body fluids.

\medterm{Lassa Virus} An arenavirus that causes Lassa fever, a potentially severe hemorrhagic illness acquired mainly through contact with infected rodent excreta or body fluids.

\medterm{Lassitude} A state of marked fatigue, weakness, or lack of energy that may occur with systemic illness, anemia, endocrine disease, infection, medication effects, or psychological conditions.

\medterm{Last Menstrual Period} The first day of the most recent normal menstrual bleeding. Clinicians use it to estimate how many weeks a pregnancy has lasted and the expected delivery date, but the estimate is less reliable with irregular cycles or uncertain dates.

\medterm{Latanoprost} Latanoprost is a prostaglandin F2-alpha analogue eye drop that increases uveoscleral outflow and lowers intraocular pressure in glaucoma or ocular hypertension; iris pigmentation and eyelash growth may occur.

\medterm{Latarjet Procedure} Shoulder-stabilization surgery that moves part of the coracoid bone, with attached tissue, to the front edge of the shoulder socket. It helps prevent repeated dislocation when bone loss or failed previous treatment makes other repairs unsuitable.

\medterm{Late Decelerations} A repeated, gradual slowing of a baby’s heart rate that starts after a uterine contraction begins and recovers after the contraction ends. It can indicate reduced oxygen transfer through the placenta and requires assessment during labor.

\medterm{Late Effects} Health problems appearing months or years after cancer treatment, including heart, lung, hormone, fertility, nerve, bone, or second-cancer complications.

\medterm{Late-Onset Hypogonadism} A condition in which an older man has symptoms of low testosterone and repeatedly low early-morning blood levels. Symptoms may include reduced sexual desire, erection problems, tiredness, reduced muscle strength, increased body fat, weak bones, or low mood. Other illnesses and medicines should also be assessed before treatment.

\synonyms
LOH; Age-Related Testosterone Deficiency; Testosterone Deficiency Syndrome

\medterm{Late Gadolinium Enhancement} A heart-MRI finding in which contrast dye remains longer in areas of heart muscle that are scarred, injured, or inflamed. Its pattern helps doctors estimate the cause and extent of heart damage.
\synonyms
LGE; Delayed Contrast Enhancement; Myocardial Scar Imaging

\medterm{Late Preterm (34-36 Weeks)} Births occurring from 34 weeks 0 days through 36 weeks 6 days of gestation. These infants have
greater risks than term infants of respiratory, temperature, glucose, jaundice, feeding,
and readmission problems.

\synonyms
34-36 weeks

\medterm{Latent} Present but inactive, hidden, or not currently producing symptoms, as in latent tuberculosis infection or a latent virus.

\medterm{Latent Autoimmune Diabetes In Adults} Diabetes caused by the immune system gradually damaging insulin-producing cells in an adult. It may first resemble type 2 diabetes, but insulin production usually declines over time and insulin treatment may eventually be needed.

\synonyms
Type 1.5 diabetes; Late-onset autoimmune diabetes of adulthood

\medterm{Latent Infection} Persistent infection in which an infectious agent remains in a host in a nonreplicating
or intermittently replicating state without active clinical disease. Reactivation can
occur when immune control is impaired, as in latent tuberculosis.

\medterm{Latent Phase Of Labor} The early part of the first stage of labor, beginning with regular contractions and gradual
cervical change before active labor. The transition to active labor is assessed clinically,
and the duration varies widely between individuals.

\medterm{Latent Syphilis} Latent syphilis is syphilis infection without current symptoms, identified by testing and classified by duration to guide treatment and follow-up.

\medterm{Latent Tuberculosis} Infection with Mycobacterium tuberculosis in which the bacteria remain inactive and cause no symptoms but may later become active.

\medterm{Latent Tuberculosis Infection} An immune response to Mycobacterium tuberculosis without clinical or microbiologic evidence of active
disease. It is usually asymptomatic and not contagious, but it can reactivate; treatment
decisions depend on age, immune status, exposure risk, test results, and drug interactions.

\medterm{Lateral} Toward or located at the side, away from the body’s midline. For example, the ears are lateral to the nose; the opposite term is medial.

\medterm{Lateral Canthotomy And Cantholysis} An emergency operation that opens the outer corner of the eyelids and releases tight supporting tissue around the eye. It lowers dangerous pressure after severe eye injury or bleeding and may protect vision.

\medterm{Lateral Collateral Ligament Injury} A stretch or tear of the strong ligament on the outer side of the knee, often caused by a blow to the inner knee or a force pushing the lower leg inward. It causes outer-knee pain and may make the knee unstable.

\medterm{Lateral Cuneiform} One of three small bones in the middle of the foot. It lies between the other cuneiform bones and the cuboid, joins the third toe bone, and helps support the foot’s arch.

\medterm{Lateral Decubitus Position} A side-lying position in which the patient rests on the nonoperative side to provide access to the chest, hip, or flank.

\medterm{Lateral Elbow Tendinopathy} Alternate index label; see \textit{Tennis elbow}.

\medterm{Lateral Epicondylitis} Alternate index label; see \textit{Tennis elbow}.

\medterm{Lateral Epicondylitis (Tennis Elbow (Lateral Epicondylitis))} Alternate index label for Tennis Elbow (Lateral Epicondylitis).

\medterm{Lateral Epicondylosis} Alternate index label; see \textit{Tennis elbow}.

\medterm{Lateral Meniscus} The lateral meniscus is a crescent-shaped cartilage pad on the outer side of the knee that distributes load and stabilizes movement.

\medterm{Lateral Rotation} Turning a limb or body part away from the body’s midline. Rotating the hip so the knee and foot point outward is one example; limited or painful movement may indicate injury.

\medterm{Lateral Sclerosis} Damage affecting nerve pathways along the side of the spinal cord that carry movement signals. It may cause increasing stiffness and weakness, as in some motor-neuron diseases; the exact diagnosis requires the associated clinical findings.

\medterm{Lateral Traction} Lateral traction is a pulling force applied sideways to a body part to improve alignment, reduce a dislocation, or relieve pressure.

\medterm{Lateral Ventricle} One of two fluid-filled spaces inside the brain that produce and contain cerebrospinal fluid, helping protect and support the brain.

\medterm{Lateral X-Ray} An X-ray image taken from the side, providing a lateral view of the body part being examined.

\medterm{Latex Agglutination Test} A laboratory test in which tiny coated particles clump together when they meet a matching substance. The reaction can help detect selected germs, antibodies, or proteins in a patient sample.

\medterm{Latex Allergies - For Hospital Patients} Latex allergy can cause contact dermatitis, hives, wheeze, or anaphylaxis; hospitals reduce risk with latex-free supplies and clear documentation of the allergy.

\medterm{Latex Allergy} An immune-mediated hypersensitivity reaction to natural rubber latex. Contact may cause
dermatitis, urticaria, rhinitis, asthma, or anaphylaxis, especially in healthcare and
occupational settings.

\synonyms
Allergy Latex (Latex Allergy)
Allergy, Latex

\medterm{Lathyrism} A neurologic disorder caused by prolonged or excessive ingestion of certain Lathyrus legumes, especially grass pea, characterized primarily by progressive spastic weakness of the legs.

\medterm{Latissimus Dorsi} A broad back muscle that extends and internally rotates the arm and assists movements such as climbing.
\medterm{Latissimus Dorsi Muscle} A broad muscle covering much of the back that pulls the arm down, backward, and toward the body. It is used in climbing, swimming, and rowing and may also be used to reconstruct tissue during surgery.

\medterm{Lavage} Washing a body space, organ, or wound with sterile fluid. It may be used to remove blood, infected material, or harmful substances, or to collect a sample; the method and risks depend on the part of the body treated.

\medterm{Lavage Fluid} Liquid used to wash a body space, wound, or organ during a medical procedure, or collected afterward for laboratory analysis. Its contents can help detect infection, bleeding, cancer cells, inflammation, or abnormal proteins.


\medterm{Lawton IADL Scale} A tool measuring ability to manage complex daily tasks such as telephone use, shopping, cooking, housekeeping, transportation, medicines, and finances.

\medterm{Laxative Overdose} Taking too much laxative can cause severe diarrhea, vomiting, dehydration, and abnormal body salts. Some products can disturb heart rhythms or damage the bowel; urgent medical advice is needed for severe symptoms, children, or large exposures.

\medterm{Laxatives For Constipation} Laxatives treat constipation through fiber and water retention, stool softening, osmotic effects, or increased bowel contractions. Choice depends on cause, duration, age, medicines, kidney function, and obstruction risk; persistent bleeding, vomiting, or severe pain needs evaluation.

\synonyms
Constipation (Laxatives For Constipation)


\medterm{Lazy Eye} Reduced vision caused by abnormal visual development during childhood, usually in one eye but sometimes in both. It may result from eye misalignment, unequal focusing, or an obstruction such as a cataract; the medical term is amblyopia.

\synonyms
amblyopia.

\medterm{Lazy Eye (Amblyopia)} Reduced vision in one or both eyes caused by abnormal visual development during childhood, often related to misalignment, unequal focusing, or visual obstruction.

\medterm{LCIS} Alternate index label; see \textit{Lobular carcinoma in situ}.

\medterm{LDH Isoenzyme Blood Test} A blood test separating lactate dehydrogenase isoenzymes to help suggest the tissue source of an elevated total LDH, although more specific tests are preferred for most conditions.

\medterm{LDL} Low-density lipoprotein, a blood particle that carries cholesterol to tissues. Too much LDL can contribute to fatty buildup in artery walls and increase the risk of heart attack and stroke; risk-based treatment may lower it.

\medterm{LDL Test} A blood test measuring low-density lipoprotein cholesterol, a major atherogenic lipoprotein used with the lipid profile to estimate cardiovascular risk and guide lipid-lowering treatment.

\medterm{Lead} A toxic heavy metal that can harm the nervous system, blood-forming system, kidneys, cardiovascular system, and reproductive system. Exposure may occur through contaminated paint dust or soil, drinking water, workplaces, hobbies, imported products, or other sources; children and fetuses are especially vulnerable to neurodevelopmental effects.

\medterm{Lead - Nutritional Considerations} Adequate iron, calcium, and vitamin C intake may help reduce lead absorption in people who are deficient, especially children, but nutrition does not remove lead from the body or replace exposure control. Suspected lead exposure requires evaluation of the source and appropriate blood lead testing.

\medterm{Lead And Tap Water} Drinking water can become contaminated with lead when it contacts lead-containing service lines, plumbing, fixtures, or solder. Risk reduction may require testing, use of an appropriately certified filter, flushing or other interim measures, and repair or replacement of the source according to local public-health guidance; boiling water does not remove lead.

\medterm{Lead Levels - Blood} A blood lead test measures the amount of lead in blood, usually reported in micrograms per deciliter (µg/dL). Results are interpreted with age, symptoms, exposure history, and testing method; an elevated capillary result may require confirmation with a venous sample. Management focuses on removing the exposure and arranging medical or public-health follow-up, with chelation reserved for selected cases.

\medterm{Lead Poisoning} Harm caused by exposure to lead, which can damage the brain, nerves, blood, kidneys, and reproductive system. Children may develop learning, behavior, growth, or developmental problems. Sources include old paint and dust, contaminated soil or water, workplace or hobby materials, and some imported products. Diagnosis uses blood lead testing; treatment starts by stopping exposure, with chelation for selected severe cases.

\medterm{Lead-Time Bias} The false appearance that screening improves survival because it finds a disease earlier, even though the person dies at the same time as without earlier diagnosis.



\medterm{Learned Insomnia} Ongoing difficulty sleeping maintained by worry, sleep-related habits, and learned associations that make the bed or bedtime trigger alertness instead of sleep.

\medterm{Learning Difficulties} Persistent problems acquiring or using academic skills such as reading, writing, or mathematics, not explained solely by low intelligence or inadequate education.
\medterm{Learning Disabilities} Persistent difficulties acquiring or using skills such as reading, writing, mathematics, language, attention, or coordination despite appropriate opportunity and instruction; assessment considers development, education, and other neurologic or sensory factors.

\medterm{Learning Disorder} A neurodevelopmental disorder causing persistent difficulty acquiring or using academic skills such as reading, writing, or mathematics despite appropriate instruction.


\medterm{Leber Congenital Amaurosis} A rare inherited disorder of the retina that causes severe vision loss from birth or early infancy. Babies may have wandering eye movements, sensitivity to light, or unusual responses of the pupils; genetic testing can help confirm the cause.

\medterm{Leber Hereditary Optic Neuropathy} An inherited disorder of the energy-producing parts of cells that damages the optic nerves. It usually causes painless, rapidly worsening loss of central vision in one eye followed by the other, often in young adults. The altered mitochondrial DNA is passed through the mother; genetic testing can confirm the diagnosis.

\synonyms
LHON

\medterm{Leech} A blood-feeding annelid used in selected reconstructive procedures to relieve venous congestion; leech bites can cause bleeding and infection.
\medterm{Leech Therapy} Carefully controlled use of medicinal leeches to remove congested blood from selected tissue, especially after reconstructive surgery. Their saliva also prevents clotting. Bleeding, infection, anemia, and allergic reactions are possible, so clinical monitoring is essential.

\medterm{Leflunomide} Leflunomide is a disease-modifying antirheumatic drug that inhibits dihydroorotate dehydrogenase and lymphocyte proliferation; it treats rheumatoid or psoriatic arthritis but can cause hepatotoxicity, cytopenias, hypertension, and fetal harm.

\medterm{Left Anterior Descending Artery} A major branch of the left coronary artery that travels in the anterior interventricular
groove and supplies the anterior left ventricle, much of the interventricular septum,
and the apex. Its occlusion can cause extensive anterior myocardial infarction.

\medterm{Left Anterior Descending Coronary Artery} A major heart artery supplying the front wall, septum, and much of the left ventricle; blockage can cause a serious heart attack.

\synonyms
LAD artery

\medterm{Left Atrial Appendage Occlusion} A procedure that seals a small pouch of the heart where blood clots often form in people with atrial fibrillation. A closure device is placed through a catheter, usually from a vein in the groin. It may reduce stroke risk when long-term blood-thinning medicine is unsuitable, but short-term medication and follow-up are still needed.

\synonyms
LAAO; Left Atrial Appendage Closure; Watchman Procedure

\medterm{Left Atrial Myxoma} A usually benign primary cardiac tumor, commonly arising from the interatrial septum in the left
atrium. It may obstruct mitral inflow or cause embolic and constitutional symptoms such as
fever, weight loss, or fatigue.

\medterm{Left Atrium} The upper-left chamber of the heart. It receives oxygen-rich blood from the lungs and passes it through the mitral valve into the left ventricle.

\medterm{Left Bundle Branch} Part of the heart’s electrical wiring that carries signals through the left lower chamber and helps it contract. Damage or blockage can delay the heartbeat signal and appear on an electrocardiogram.

\medterm{Left Circumflex Coronary Artery} A heart artery curving around the left side and back, supplying parts of the left ventricle and sometimes the lower heart wall.

\synonyms
Circumflex artery

\medterm{Left Heart Catheterization} A procedure that passes a thin tube into the left side of the heart or nearby arteries to measure pressures, assess valves, or look for blocked coronary arteries.
\medterm{Left Heart Ventricular Angiography} An imaging study that injects contrast into the left ventricle to assess ventricular contraction, wall motion, ejection fraction, and selected structural abnormalities.

\medterm{Left Hemiparesis, Gaze Deviated To The Right} Weakness on the left side with both eyes looking to the right. This pattern can occur with a serious injury to the right side of the brain, including stroke, but examination and brain imaging are needed to identify the cause.

\medterm{Left Hypochondriac Region} The upper-left area of the abdomen, beneath the lower ribs. It overlies part of the stomach, spleen, pancreas, and colon; pain or swelling there may arise from digestive, splenic, or nearby chest disease.

\medterm{Left Iliac Region} The lower-left area of the abdomen, overlying part of the large intestine and urinary tract. Pain there may occur with diverticular disease, bowel inflammation, urinary problems, or other pelvic conditions.

\medterm{Left Main Coronary Artery} The short vessel arising from the left aortic sinus and dividing principally into the
left anterior descending and circumflex arteries. Significant left-main disease
jeopardizes a large proportion of left-ventricular myocardium.

\medterm{Left Testicular Torsion} Acute rotation of the left testis on its spermatic cord, obstructing blood flow and
causing ischemia. It is a urological emergency presenting with sudden severe left
scrotal pain, nausea, vomiting, and a high-riding, horizontally oriented testis with
absent cremasteric reflex.

\medterm{Left Ventricle} The lower-left chamber of the heart. It receives oxygen-rich blood from the left atrium and pumps it through the aortic valve into the aorta and throughout the body.

\medterm{Left Ventricular Hypertrophy} Thickening of the muscle wall of the heart's main pumping chamber. It is often caused by long-term high blood pressure or a narrowed aortic valve and can make the heart stiff, leading to breathlessness, rhythm problems, or heart failure.

\medterm{Left Ventricular Noncompaction} A rare heart-muscle disorder in which the main pumping chamber has an unusually spongy inner muscle layer. It may cause no symptoms or lead to heart failure, abnormal rhythms, blood clots, or stroke.


\medterm{Left-Sided Superior Vena Cava} A birth-related variation in which a large vein on the left side of the chest drains blood through an unusual route to the heart. It usually causes no symptoms but matters when placing central lines or performing heart surgery.

\medterm{Left-Ventricular Assist Device} An implanted mechanical pump that helps the heart's main pumping chamber send blood around the body. It may support people with severe heart failure while they await a transplant or as long-term treatment; infection, bleeding, and blood clots are important risks.

\medterm{Leg} The lower limb, often meaning the region between the knee and ankle, although it may also refer broadly to the entire lower extremity.


\medterm{Leg CT Scan} A series of detailed X-ray pictures of the leg combined by a computer. It can show complex fractures, bone alignment, infection, tumors, blood-vessel problems, and selected soft-tissue injuries.

\medterm{Leg Fracture} Alternate index label; see \textit{Broken leg}.

\medterm{Leg Lengthening And Shortening} Operations that gradually increase or decrease the length of a leg bone to correct a major difference between the legs or a deformity. Recovery is prolonged and may involve an external or internal support device and rehabilitation.

\medterm{Leg MRI Scan} Magnetic resonance imaging of the leg using a magnetic field and radio waves to evaluate muscles, tendons, ligaments, nerves, marrow, joints, and soft-tissue masses.

\medterm{Leg Or Foot Amputation} Surgical removal of part or all of a leg or foot, usually because of severe injury, infection, poor blood flow, or cancer. Rehabilitation supports mobility afterward.
\medterm{Leg Or Foot Amputation - Dressing Change} Changing and examining the dressing over a leg or foot amputation wound. The dressing protects the wound, absorbs drainage, and allows assessment for bleeding, infection, poor healing, pressure injury, or separation of the wound edges.

\medterm{Leg Pain} Pain in the lower limb caused by muscle or tendon injury, joint disease, nerve problems, vascular disease, infection, or other conditions.

\medterm{Leg Pain With Cramping} Painful tightening or spasms in a leg muscle; causes include exertion, dehydration, medicines, nerve or muscle disorders, and reduced arterial blood flow causing intermittent claudication.

\medterm{Leg Swelling} Enlargement or puffiness of one or both legs caused by fluid retention, venous disease, injury, infection, inflammation, or a blood clot.

\medterm{Leg Ulcer} An open sore on the leg that fails to heal normally, often caused by poor venous circulation, arterial disease,
diabetes, pressure, or inflammation. Pain, spreading redness, fever, or black tissue requires
prompt assessment.

\medterm{Legg-Calve-Perthes Disease} A childhood hip disorder in which the blood supply to the top of the thighbone is temporarily reduced. The bone can weaken and change shape, causing a limp, hip or knee pain, stiffness, and limited movement.

\medterm{Legg-Calvé-Perthes Disease} A childhood disorder caused by temporary loss of blood supply to the femoral head, resulting in avascular necrosis and hip pain or limp.
\medterm{Legg-Calvé-Perthes Disease (Perthes Disease)} A childhood disorder in which the blood supply to the femoral head is temporarily impaired, causing avascular necrosis, hip pain, stiffness, and possible deformity.

\synonyms
Perthes Disease

\medterm{Legg–Calvé–Perthes Disease} A childhood hip disorder in which reduced blood flow temporarily damages the rounded top of the thighbone. It can cause a limp, hip or knee pain, stiffness, and later deformity as the bone heals.

\medterm{Legionella} Legionella is a group of water-associated bacteria that can cause Legionnaires' disease, a severe pneumonia, or Pontiac fever, a milder illness.

\medterm{Legionella Longbeachae} Legionella longbeachae is a Legionella species associated particularly with exposure to contaminated potting mix and capable of causing pneumonia.

\medterm{Legionella Pneumonia} Pneumonia caused by Legionella bacteria, usually after inhaling contaminated water droplets. It can cause high fever, cough, breathlessness, diarrhea, confusion, and low oxygen and requires antibiotic treatment.

\medterm{Legionella Pneumonia (Legionnaire Disease And Pontiac Fever)} Alternate index label for Legionnaire Disease and Pontiac Fever.

\medterm{Legionella Pneumophila} A bacterium found in water systems that can cause severe pneumonia after a person inhales contaminated water droplets. It may also cause a milder fever illness; older age, smoking, and weakened immunity increase risk.

\medterm{Legionellosis} An infection caused by Legionella bacteria, ranging from a mild flu-like illness to Legionnaires’ disease with pneumonia.

\medterm{Legionnaire Disease And Pontiac Fever} Legionnaires disease is a severe pneumonia caused by Legionella, while Pontiac fever is a milder flu-like illness from the same bacteria. Exposure to contaminated water systems is typical; pneumonia, confusion, low oxygen, or high fever requires urgent treatment.

\synonyms
Legionella Pneumonia (Legionnaire Disease And Pontiac Fever)
Legionnaires Disease (Legionnaire Disease And Pontiac Fever)

\medterm{Legionnaires Disease} A severe pneumonia caused by Legionella bacteria, usually acquired by inhaling contaminated water aerosols. It may also cause gastrointestinal symptoms, confusion, and low serum sodium.

\medterm{Legionnaires Disease (Legionnaire Disease And Pontiac Fever)} Alternate index label for Legionnaire Disease and Pontiac Fever.

\medterm{Legionnaires' Disease} A severe form of pneumonia caused by Legionella pneumophila, a Gram-negative aerobic
bacillus that is found naturally in freshwater environments (lakes, streams) and
artificially in water systems (air conditioning cooling towers, hot tubs, decorative
fountains, plumbing systems, showers, humidifiers).

\medterm{Legs Restless (Restless Leg Syndrome)} Alternate index label for Restless Leg Syndrome.

\medterm{Leigh Disease} A progressive mitochondrial disorder of early life characterized by symmetric lesions of the brainstem and basal ganglia, developmental regression, movement abnormalities, and lactic acidosis.

\medterm{Leiomyoma} A usually noncancerous tumor made of smooth muscle. It most often develops in the uterus, where it is called a fibroid, but it can also occur in the skin, digestive tract, or other organs.

\medterm{Leiomyoma (Fibroid)} A benign tumor of uterine smooth muscle that may cause heavy menstrual bleeding, pelvic pressure, pain, or no symptoms.

\synonyms
Fibroid

\medterm{Leiomyomatosis} Leiomyomatosis is development of multiple benign smooth-muscle tumors or proliferations, which may affect the uterus, skin, lungs, or other sites.

\medterm{Leiomyomatosis Peritonealis Disseminata} A rare benign condition with multiple smooth-muscle nodules scattered over the peritoneum, often associated with pregnancy or hormonal stimulation and sometimes mimicking peritoneal carcinomatosis.

\medterm{Leiomyosarcoma} A malignant tumor of smooth muscle, most often arising in the uterus, retroperitoneum,
or gastrointestinal tract.

\medterm{Leiomyosarcoma Of The Vulva} A rare malignant smooth-muscle tumor of the vulva, usually presenting as a solitary enlarging mass. Histologic atypia, mitotic activity, necrosis, and invasion help determine malignant behavior.

\medterm{Leishmania} A group of single-celled parasites transmitted mainly by infected sandflies. Different species cause skin, mucosal, or visceral leishmaniasis, ranging from localized ulcers to life-threatening illness.


\medterm{Leishmania Donovani} A parasite spread by infected sandflies that causes visceral leishmaniasis. It can infect organs such as the spleen, liver, and bone marrow, causing prolonged fever, weight loss, enlarged organs, and reduced blood-cell counts.

\medterm{Leishmania Infection} A parasitic infection caused by \textit{Leishmania} protozoa and transmitted by infected sandflies; it can cause cutaneous, mucocutaneous, or visceral leishmaniasis.

\medterm{Leishmania Major} A parasite commonly causing cutaneous leishmaniasis in parts of Africa, Asia, and the Middle East. Sandfly transmission produces one or more skin ulcers that may heal slowly and leave scars.

\medterm{Leishmania Mexicana} A parasite that causes mainly cutaneous leishmaniasis in the Americas. Infection after a sandfly bite may produce a persistent skin nodule or ulcer and, less commonly, involve mucous membranes.

\medterm{Leishmania Tropica} A sandfly-transmitted parasite that commonly causes localized cutaneous leishmaniasis. It produces chronic skin sores that may heal over months and leave permanent scars.

\medterm{Leishmaniasis} An infection caused by Leishmania parasites transmitted by sandflies. It may produce skin ulcers, damage to the nose or mouth, or a life-threatening illness affecting organs such as the spleen, liver, and bone marrow.

\medterm{Lemaitre'S Sign} Pain or discomfort felt in the urethra or penis during a rectal examination. It has traditionally been associated with inflammation near the prostate or back part of the urinary passage but is not specific for one diagnosis.
\medterm{Lemierre Syndrome} A potentially severe oropharyngeal infection complicated by septic thrombophlebitis of the internal jugular vein and frequently metastatic infection, classically caused by Fusobacterium necrophorum.



\medterm{Lenalidomide} Lenalidomide is an immunomodulatory anticancer drug used for multiple myeloma and selected myelodysplastic or lymphoma disorders; myelosuppression, thrombosis, infection, and severe teratogenicity are important risks.

\medterm{Length-Time Bias} The tendency of screening to preferentially detect slower-growing or less aggressive
disease because it remains detectable for a longer preclinical period. It can exaggerate
the apparent benefit of screening.

\medterm{Lengthening Of Limb} A surgical or orthopedic process that gradually increases the length of a bone, often using controlled distraction after osteotomy.
\medterm{Lennox-Gastaut Syndrome} A severe childhood-onset developmental and epileptic encephalopathy with multiple seizure types, characteristic slow spike-and-wave EEG patterns, cognitive impairment, and often treatment-resistant seizures.

\medterm{Lenograstim} A recombinant granulocyte colony-stimulating factor that increases neutrophil production and is used to prevent or treat some forms of neutropenia.

\medterm{Lens} A clear, flexible structure behind the iris and pupil that bends light onto the retina. Muscles change its shape for near or distant focus; clouding causes cataract, while displacement can blur vision or raise eye pressure.

\medterm{Lens Coloboma} A lens coloboma is a congenital segmental defect or notch in the eye lens caused by incomplete development.

\medterm{Lens Disease} A disorder of the eye lens, including cataract, dislocation, inflammation-related opacity, or developmental abnormality, that may impair focusing or vision.

\medterm{Lens Dislocation} Displacement of the crystalline lens from its normal position because of trauma, zonular weakness, surgery, or connective-tissue disease; it can cause blurred vision, glaucoma, or inflammation.

\medterm{Lens Prescription} Written instructions specifying the lens power, correction, and other features needed for eyeglasses or contact lenses in each eye.
\medterm{Lens Subluxation} Partial displacement of the crystalline lens while some zonular attachments remain, potentially causing blurred or fluctuating vision, refractive change, or glaucoma.

\medterm{Lensometer} An instrument measuring the prescription, power, and optical alignment of eyeglass lenses. It helps verify new glasses, assess existing lenses, and check whether corrective strength matches the prescribed values.

\synonyms
Lensmeter; focimeter


\medterm{Lentiform Nucleus} A group of deep brain structures involved in controlling movement. It includes the putamen and globus pallidus and works with other brain regions to help start and regulate actions.

\medterm{Lentigines, Solar} Alternate index label; see \textit{Age spots (liver spots)}.

\medterm{Lentigo} A lentigo is a sharply defined flat brown skin spot caused by increased melanocytes or pigment and may be benign or, rarely, atypical.

\medterm{Lentigo Maligna} A melanoma in situ that occurs on chronically sun-damaged skin, most commonly on the
face of elderly individuals. It presents as a large, irregularly pigmented, slow-growing
macule with variegated shades of brown and black. The radial growth phase may persist
for years before invasion occurs.

\medterm{Lentigo Maligna Melanoma} An invasive melanoma arising in a lentigo maligna on chronically sun-damaged skin, usually of the head or neck; it presents as an enlarging irregular pigmented patch or plaque.


\medterm{LEOPARD Syndrome} A rare inherited disorder that can cause many dark skin spots, abnormal heart rhythms, widely spaced eyes, narrowing of the pulmonary valve, genital differences, slow growth, and hearing loss. The combination and severity vary between people.

\medterm{Leopold Maneuvers} A systematic abdominal palpation technique used to determine fetal position and
presentation. Four maneuvers: (1) fundal grip identifies fetal lie (head vs breech), (2)
umbilical grip locates the fetal back, (3) Pawlik grip identifies the presenting part
and its position, (4) pelvic grip confirms presentation and attitude.

\medterm{Lepidoptera} The insect order containing butterflies and moths; contact with some caterpillars can cause skin, eye, or systemic reactions.

\medterm{Lepirudin} Lepirudin is a direct thrombin inhibitor formerly used as an anticoagulant in heparin-induced thrombocytopenia; bleeding and kidney clearance are important considerations.

\medterm{Lepromatous Leprosy} A widespread form of leprosy caused by a large number of bacteria. It can produce many skin patches, nodules, thickened skin, loss of feeling, nerve damage, and facial changes, with a high risk of spread to other tissues if untreated.

\medterm{Lepromin Skin Test} A skin test using harmless material from the leprosy bacterium to show how strongly the immune system reacts. It may help classify the type of leprosy but cannot by itself diagnose active infection.

\medterm{Leprosy (Hansen Disease)} A long-term infection caused mainly by bacteria that affect the skin and cooler nerves in the arms, legs, and face. It can cause pale or reddish patches with reduced feeling, thickened nerves, weakness, and eye or nose problems. A skin sample or other tests may confirm it. Treatment uses several antibiotics, commonly rifampicin, dapsone, and clofazimine, to prevent lasting nerve damage and disability.

\medterm{Leptin} A hormone produced mainly by adipose tissue that signals stored-energy status to the brain and influences appetite and energy balance.

\medterm{Leptin Receptor} A cell-surface protein that receives signals from leptin, a hormone made mainly by fat tissue. It helps the brain regulate appetite, energy use, growth, and reproductive function in response to stored body energy.

\medterm{Leptomeningeal} Relating to the pia mater and arachnoid mater, the two thin inner layers covering the brain and spinal cord; cancer or infection can spread in these layers.

\medterm{Leptomeningeal Metastasis} Spread of cancer to the thin membranes and fluid surrounding the brain and spinal cord. It can cause headache, vomiting, weakness, confusion, seizures, or changes in vision and requires specialist treatment.

\medterm{Leptomeningeal Sarcoma} Leptomeningeal sarcoma is a rare malignant tumor involving the membranes and fluid spaces around the brain or spinal cord.

\medterm{Leptomeninges} The arachnoid mater and pia mater, the two delicate inner membranes covering the brain and spinal cord.

\medterm{Leptomeningitis} Inflammation or infection of the leptomeninges, the arachnoid and pia mater covering the brain and spinal cord.

\medterm{Leptospira} A genus of spiral-shaped bacteria that includes the causes of leptospirosis, usually acquired through water or soil contaminated with infected animal urine.

\medterm{Leptospirosis} An infection caused by Leptospira bacteria, usually acquired through water or soil contaminated with infected animal urine. It may cause fever, headache, muscle aches, red eyes, vomiting, or severe kidney and liver failure.

\synonyms
Weils Disease (Leptospirosis)

\medterm{Leriche Syndrome} Aortoiliac occlusive disease causing reduced blood flow to the pelvis and lower limbs.
The classic presentation includes buttock or thigh claudication, diminished femoral
pulses, and erectile dysfunction, although symptoms vary with the extent of collateral
circulation.



\medterm{Lesch-Nyhan Syndrome} An inherited disorder of purine metabolism caused by severe HPRT enzyme deficiency, leading to high uric acid, neurologic problems, and characteristic self-injurious behavior.

\medterm{Lesion} Any localized abnormal change in tissue caused by disease, injury, inflammation, infection, or a tumor. A lesion may be visible, felt, or found only on imaging or microscopic examination.
\medterm{Lesion Identification} The process of locating and describing an abnormal area in the body using examination, imaging, endoscopy, laboratory tests, or tissue sampling.
\medterm{Lesions On The Brain (Brain Lesions (Lesions On The Brain))} Alternate index label for Brain Lesions (Lesions on the Brain).

\medterm{Lesser Curvature} The shorter, inward curve along the upper and inner edge of the stomach. It is attached to a fold of tissue leading to the liver and contains blood vessels that supply the stomach.

\medterm{Lesser Trochanter} A small bony projection on the inner upper thighbone where the iliopsoas muscle attaches. It helps bend the hip; injury may occur from a sudden muscle pull in adolescents or, in adults, from weakened bone.


\medterm{Let-Down Reflex} Release of breast milk when oxytocin makes the small muscles around milk-producing glands contract. Suckling, breast stimulation, or seeing or hearing a baby can trigger it, although stress may delay it.
\medterm{Lethal} Able to cause death. In medicine, the word may describe a disease, injury, dose of a substance, or genetic change that prevents survival.

\medterm{Lethal-Means Counseling} A preventive intervention that reduces access to firearms, medications, toxic
substances, or other methods of self-harm during periods of elevated risk. It is
conducted collaboratively and is an important component of suicide prevention.

\medterm{Lethargy} Lethargy is reduced energy, alertness, or responsiveness; it may result from sleep loss, infection, metabolic disease, medicines, depression, or neurologic illness, and sudden marked lethargy requires urgent assessment.

\medterm{Letrozole} Letrozole is an aromatase inhibitor that lowers oestrogen production and is used mainly for hormone-receptor-positive breast cancer and ovulation induction; hot flushes, arthralgia, bone loss, and fetal harm are important considerations.



\medterm{Leucocyte} Another spelling of leukocyte, meaning a white blood cell. Leukocytes help defend against infection and other threats; major groups include neutrophils, lymphocytes, monocytes, eosinophils, and basophils.
\medterm{Leucocyte (White Blood Cell)} A blood cell that helps defend the body against infections, foreign substances, and abnormal cells; major types include lymphocytes, neutrophils, monocytes, eosinophils, and basophils.

\synonyms
White Blood Cell

\medterm{Leucocytosis} A higher-than-usual number of white blood cells in the blood. It may occur with infection, inflammation, physical stress, medicines, smoking, or blood and bone-marrow disorders, so the type of cell and the person's symptoms help find the cause.
\medterm{Leucoderma} Loss or marked reduction of skin pigmentation, often used as a descriptive term for depigmented patches.
\medterm{Leucoderma (Vitiligo)} Loss of skin pigment producing sharply defined pale or white patches, usually because melanocytes have been destroyed or have stopped functioning.

\synonyms
Vitiligo

\medterm{Leucomalacia} Softening or loss of tissue after injury, poor blood flow, infection, or inflammation. In newborns, the term often refers to damage of the brain’s white matter near the fluid spaces, which may affect movement or development.
\medterm{Leucopenia} A lower-than-usual number of white blood cells. It may result from medicines, infections, immune disorders, nutritional deficiency, or bone-marrow disease and can increase the risk of infection, especially when neutrophils are low.

\medterm{Leucoplakia} A persistent white patch on a moist body lining, often inside the mouth, that cannot be wiped away. It may result from irritation and sometimes contains precancerous changes, so it should be examined.
\medterm{Leucopoiesis} The production and maturation of white blood cells, mainly in the bone marrow. It is regulated by growth factors and increases or changes during infection, inflammation, and some blood disorders.
\medterm{Leucorrhoea (Leukorrhoea)} A whitish vaginal discharge that may be physiological or may result from infection, inflammation, or another genital disorder.

\synonyms
Leukorrhoea

\medterm{Leucovorin} Leucovorin is a folate medicine used to reduce the toxicity of methotrexate, enhance some chemotherapy regimens, and treat certain folate-related problems.

\medterm{Leucovorin Calcium} Leucovorin calcium is the calcium salt of leucovorin, used as rescue therapy after methotrexate and in selected colorectal-cancer treatments.

\medterm{Leukaemia} A cancer of blood-forming tissues in which abnormal white blood cells multiply uncontrollably and crowd out normal cells needed for immunity, clotting, and oxygen transport.
\medterm{Leukapheresis} A blood-filtering procedure that removes excess white blood cells and returns the remaining blood. It may rapidly reduce very high counts causing symptoms in leukemia, but it is temporary and does not replace treatment of the underlying disease.

\medterm{Leukemia} Leukemia is cancer of blood-forming cells that disrupts normal marrow production and may cause anemia, infection, bleeding, abnormal blood counts, and organ infiltration.
Leukemias are classified by the cell lineage and by clinical course, including acute and
chronic forms.

\synonyms
Blood Cancer (Leukemia)
Leukemia, General


\medterm{Leukemia Cutis} A condition in which leukemia cells collect in the skin, causing firm red, purple, or brown bumps, nodules, or patches. It often accompanies acute myeloid leukemia and requires assessment of the underlying blood cancer.

\medterm{Leukemia Genetics} Changes in genes or chromosomes in leukemia cells help identify the leukemia type, estimate its likely course, and choose treatment. The meaning of each change depends on the person’s specific leukemia.

\medterm{Leukemia, Acute Lymphocytic} Alternate index label; see \textit{Acute lymphocytic leukemia}.

\medterm{Leukemia, Acute Myelogenous} Alternate index label; see \textit{Acute myelogenous leukemia}.

\medterm{Leukemia, Chronic Lymphocytic} Alternate index label; see \textit{Chronic lymphocytic leukemia}.

\medterm{Leukemia, Chronic Myelogenous} Alternate index label; see \textit{Chronic myelogenous leukemia}.

\medterm{Leukemia, General} Alternate index label; see \textit{Leukemia}.

\medterm{Leukemia, Hairy Cell} Alternate index label; see \textit{Hairy cell leukemia}.

\medterm{Leukemic Infiltration Of The Gingiva} Swelling of the gums caused by leukemia cells collecting in gum tissue. The gums may become red, soft, enlarged, and prone to bleeding, especially in some acute myeloid leukemias.

\medterm{Leukemoid Reaction} A marked reactive increase in the white-cell count, usually with prominent neutrophilia
and immature granulocytes, caused by severe infection, inflammation, malignancy, or
tissue injury. It resembles leukemia but is not a clonal hematologic malignancy.

\medterm{Leukoaraiosis} White-matter changes seen on brain scans, often related to aging and long-term small-vessel disease such as high blood pressure. Extensive changes may be associated with slower thinking, walking problems, or stroke risk.

\medterm{Leukocyte} A white blood cell that helps defend the body against infection and other threats. Main types include neutrophils, lymphocytes, monocytes, eosinophils, and basophils; unusually high or low numbers can reflect infection, inflammation, medicines, or bone-marrow disease.

\medterm{Leukocyte Adhesion Deficiency} A rare inherited immune disorder in which white blood cells cannot attach to blood-vessel walls and reach infected tissue normally. It can cause delayed umbilical-cord separation, repeated skin infections with little pus, poor wound healing, gum disease, and dangerously severe infection in infants.

\medterm{Leukocyte Alkaline Phosphatase} A measure of an enzyme in certain white blood cells. It was formerly used to help distinguish a reactive response to infection from chronic myeloid leukemia, but newer genetic and blood tests are usually more informative.


\medterm{Leukocyte Disorder} A disorder involving the number, function, or development of white blood cells, including leukopenia, leukocytosis, immune deficiency, and malignant leukocyte proliferation.


\medterm{Leukocyte Esterase Urine Test} A urine dipstick test detecting leukocyte esterase released by white blood cells; a positive result supports pyuria but does not by itself diagnose urinary tract infection.


\medterm{Leukocyte Transfusion} Transfusion of donated white blood cells in rare, carefully selected situations involving life-threatening infection and severe inability to produce neutrophils. Benefits are uncertain and risks include fever, lung injury, immune reactions, and infection transmission.

\medterm{Leukocytes} White blood cells that defend the body against infections, abnormal cells, and other threats. Major types include neutrophils, lymphocytes, monocytes, eosinophils, and basophils, each with different immune functions.

\medterm{Leukocytoclastic Vasculitis} Inflammation of small blood vessels, usually in the skin, often causing raised purple or red spots that can be felt. Medicines, infections, or immune diseases may trigger it, and internal organs can sometimes be affected.

\medterm{Leukocytosis} A higher-than-usual number of white blood cells. It may occur with infection, inflammation, stress, smoking, medicines, or blood and bone-marrow disorders; the specific cell type and symptoms help identify the cause.

\medterm{Leukodystrophy} A group of inherited disorders that damage myelin, the protective covering around nerves in the brain and spinal cord. They can cause progressive problems with movement, thinking, vision, hearing, or other nervous-system functions.

\medterm{Leukoencephalopathy} Leukoencephalopathy is disease of the brain’s white matter caused by inherited, infectious, toxic, vascular, metabolic, inflammatory, or other processes.

\medterm{Leukonychia} White areas or bands in one or more nails. Causes include minor injury to the nail, inherited traits, infection, medicines, or illness; the pattern and whether it grows out with the nail help determine the cause.

\medterm{Leukonychia Striata} Longitudinal white bands running along the nail plate, usually caused by changes in the nail matrix. They may be inherited or follow repeated minor trauma, skin disease, or some medicines; persistent or widespread changes should be medically assessed.

\medterm{Leukonychia Totalis} Uniform white discoloration involving most or all of one or more nail plates. It is often inherited but can occasionally be associated with illness or medication; persistent changes affecting many nails should be evaluated with the medical history and other findings.

\medterm{Leukopenia} An abnormally low white-blood-cell count, which may increase susceptibility to infection depending on the cells affected and the severity of the reduction.

\medterm{Leukoplakia} A persistent white patch on the lining of the mouth or another mucosal surface that cannot be scraped away and cannot be explained by another condition. Some patches can develop into cancer, so persistent lesions need examination and often a biopsy.

\synonyms
Smokers Keratosis (Leukoplakia)

\medterm{Leukopoiesis} The process by which the bone marrow makes and develops white blood cells. Production increases or changes during infection, inflammation, immune disorders, and some blood diseases.

\medterm{Leukorrhea} Whitish vaginal discharge that may be normal, especially during hormonal changes, or may result from infection or inflammation. Foul odor, itching, pain, bleeding, or a new change requires evaluation.
\medterm{Leukostasis} A medical emergency in which very high numbers of abnormal leukocytes impair microvascular blood flow, potentially causing respiratory distress, neurologic symptoms, organ ischemia, or hemorrhage.

\medterm{Leukotriene} A leukotriene is an inflammatory lipid mediator made from arachidonic acid that contributes to bronchoconstriction, mucus, vascular leak, and immune-cell recruitment.

\medterm{Leukotriene Blockers} Medicines that reduce the effects or production of leukotrienes, substances involved in inflammation and airway narrowing. They help prevent some asthma or allergy symptoms but do not usually relieve a sudden asthma attack.

\synonyms
Leukotriene antagonists


\medterm{Leukotriene Modifiers} Medicines that reduce the production or effects of leukotrienes, inflammatory chemicals narrowing airways and causing allergy symptoms. They help prevent selected asthma or allergy problems but do not replace emergency rescue treatment.

\medterm{Leukotriene Receptor} A cell-surface protein activated by leukotrienes, inflammatory chemicals that can narrow airways, increase mucus, and promote swelling. Blocking these receptors can help prevent selected asthma and allergy symptoms.

\medterm{Leukotrienes} Chemical signals released during inflammation that can narrow airways, increase mucus, and cause swelling. They contribute to asthma, allergies, and some other inflammatory conditions.

\medterm{Leuprolide} Leuprolide is a gonadotropin-releasing hormone agonist used to suppress sex-hormone production in conditions such as prostate cancer, endometriosis, and fibroids.

\medterm{Leuprolide Acetate} Leuprolide acetate is the salt form of leuprolide, used to suppress ovarian or testicular hormone production for selected medical conditions.

\medterm{Levalbuterol} An inhaled short-acting bronchodilator that stimulates beta-2 receptors and relaxes airway muscle. It relieves wheezing and shortness of breath but may cause tremor, nervousness, headache, or a fast heartbeat.
\medterm{Levamisole} Levamisole is an antiparasitic and immune-modulating drug used selectively; serious low white-cell counts and vasculitic reactions can occur.

\medterm{Levator Ani} A group of pelvic-floor muscles that supports the bladder, bowel, and reproductive organs and helps control urination and bowel movements. Weakness or injury may contribute to pelvic organ prolapse or incontinence.

\medterm{Levetiracetam} Levetiracetam is a broad-spectrum antiseizure medicine that binds synaptic vesicle protein SV2A; it treats focal and generalized seizures and may cause somnolence, dizziness, or behavioral and mood changes.

\medterm{Levobunolol} Levobunolol is a beta-blocker eye drop used to lower pressure inside the eye in glaucoma or ocular hypertension.

\medterm{Levocardia} Levocardia is a normal left-sided position of the heart; it may occur with normal organ arrangement or with other congenital laterality abnormalities.

\medterm{Levocarnitine} Levocarnitine is the biologically active form of carnitine used to treat selected carnitine deficiencies and help transport fatty acids into mitochondria.

\medterm{Levodopa} Levodopa is a dopamine precursor that crosses the blood--brain barrier and is converted to dopamine, usually with carbidopa or benserazide, to treat Parkinson disease; dyskinesias, hallucinations, nausea, and motor fluctuations can develop.

\medterm{Levofloxacin} Levofloxacin is a fluoroquinolone antibiotic that inhibits bacterial DNA gyrase and topoisomerase IV; tendon injury, peripheral neuropathy, QT prolongation, dysglycaemia, and central-nervous-system effects are important class risks.

\medterm{Levoleucovorin} Levoleucovorin is the active isomer of leucovorin, used as methotrexate rescue and with selected chemotherapy regimens.

\medterm{Levoleucovorin Calcium} Levoleucovorin calcium is the calcium salt of levoleucovorin, used to reduce methotrexate toxicity and support selected cancer treatments.

\medterm{Levomepromazine} A phenothiazine medicine with antipsychotic, sedative, antiemetic, and analgesic properties, used in selected palliative-care and psychiatric settings.
\medterm{Levonorgestrel} A progestin hormone used in birth-control pills, hormonal intrauterine devices, and emergency contraception to prevent or delay ovulation and alter cervical mucus.

\medterm{Levorphanol} Levorphanol is a potent opioid analgesic that can relieve severe pain but also cause sedation, dependence, and respiratory depression.

\medterm{Levothyroxine} A synthetic form of thyroxine (T4) used as thyroid hormone replacement in hypothyroidism and in selected other thyroid conditions; dose is individualized and monitored with laboratory tests.

\medterm{Levothyroxine Sodium} A manufactured form of thyroid hormone used to replace hormone missing in hypothyroidism. The dose is adjusted using blood tests, and it is usually taken consistently on an empty stomach.

\synonyms
Levothyroxine

\medterm{Levulose} An older name for fructose, a simple sugar found in fruit, honey, and many sweeteners.
\medterm{Lewy Body} A Lewy body is an abnormal neuronal inclusion containing alpha-synuclein, characteristic of Parkinson disease and Lewy-body dementia.

\medterm{Lewy Body Dementia} A progressive brain disorder caused by abnormal alpha-synuclein protein deposits called Lewy bodies. It causes dementia with changing attention or alertness, recurrent visual hallucinations, Parkinson-like movement problems, sleep disorders such as acting out dreams, and sometimes automatic body-function problems such as constipation or dizziness on standing. Some people are very sensitive to antipsychotic medicines, so treatment must be individualized.

\synonyms
Dementia, Lewy Body

\medterm{Lewy Body Disease} A progressive brain disorder caused by abnormal protein deposits in nerve cells. It may cause changing attention, visual hallucinations, movement problems, stiffness, sleep-related behaviors, and gradual difficulty with thinking and daily activities.


\medterm{Leydig Cell} A cell between the testicular tubules that produces testosterone, mainly in response to luteinizing hormone.

\medterm{Leydig Cell Testicular Tumor} A usually rare sex-cord stromal tumor arising from testosterone-producing Leydig cells in the testis; it may cause a testicular mass or hormone-related signs.

\medterm{Leydig Cell Tumor} A testicular sex cord-stromal tumor that secretes testosterone or estrogen, causing
precocious puberty or gynecomastia; most are benign and treated by orchidectomy.

\medterm{Leydig Cells} Cells found between the sperm-making tubes in the testes. In response to luteinising hormone, they make testosterone, which supports male sexual development, sperm production, sex drive, and muscle and bone health.
\medterm{Löfgren Syndrome} A sudden form of sarcoidosis causing swollen lymph nodes near the lungs, tender red lumps on the shins, and painful swollen ankles. It often improves over time, but medical follow-up is needed to monitor symptoms and organ involvement.

\medterm{LGE} Alternate name; see \textit{Late Gadolinium Enhancement}.

\medterm{LH} Luteinising hormone, a hormone released by the pituitary gland. In people with ovaries it helps trigger ovulation and supports hormone production; in people with testes it stimulates testosterone production. Blood tests can help assess reproductive or pituitary problems.

\synonyms
Luteinizing hormone

\medterm{LH Response To GnRH Blood Test} A hormone test that measures luteinizing hormone before and after a medicine that stimulates its release. It helps assess delayed or early puberty and problems involving the brain, pituitary gland, ovaries, or testes.

\medterm{Libman-Sacks Endocarditis} Noninfectious inflammation of heart valves that can occur with lupus or antiphospholipid syndrome. Small deposits may form on the valves, sometimes causing leakage or increasing the risk of blood clots and stroke.

\synonyms
Verrucous endocarditis; Nonbacterial verrucous endocarditis

\medterm{Li-Fraumeni Syndrome} An inherited condition that greatly increases the chance of developing several cancers, often at a young age. Cancers may include breast cancer, sarcoma, brain tumors, adrenal cancer, or leukemia; specialized screening is usually advised.

\medterm{LI-RADS} A standardized system for reporting liver-imaging findings in people at increased risk of liver cancer, especially those with cirrhosis. Its categories estimate how likely a finding is to be cancer and help guide follow-up or biopsy.


\medterm{Libido} A person's sexual desire or interest. It can change with age, relationships, mood, stress, sleep, illness, pregnancy, menopause, hormone levels, and medicines; a change is important mainly when it causes distress or affects a person's wellbeing.

\medterm{Libido And Sexual Performance} Sexual desire and sexual function can be influenced by hormones, blood flow, nerves, medicines, illness, mood, stress, sleep, and relationships. A persistent or distressing change deserves an assessment for treatable causes.

\medterm{Lice} Small insects that live on hair or clothing and feed on blood. Head, body, and pubic lice cause itching and may leave eggs attached to hair. They spread mainly through close contact and are not a sign of poor hygiene.

\medterm{Lice Infestation} Infestation of the hair or skin by small insects that feed on blood. Head, body, and pubic lice spread mainly through close contact and cause itching or visible eggs attached to hair. Treatment uses an approved lice-killing product and fine combing; close contacts and clothing may need attention.

\medterm{Lice Infestations} Pediculosis is infestation by head, body, or pubic lice. Head lice cause scalp pruritus
and nits attached to hair shafts; body lice live mainly in clothing; pubic lice affect
coarse hair-bearing regions.


\medterm{Lice, Body} Alternate index label; see \textit{Body lice}.

\medterm{Licensed Clinical Social Worker} A social worker licensed to provide clinical assessment and psychotherapy or other mental-health services within the scope of local law and training.

\medterm{Lichen} A descriptive term for a firm, flat-topped, raised skin lesion or for a skin pattern resembling lichen planus. It describes appearance rather than a single disease or cause.

\medterm{Lichen Planopilaris} A lymphocytic inflammatory scalp disorder that attacks hair follicles and can cause scarring alopecia, perifollicular scale, redness, and permanent hair loss.

\medterm{Lichen Planus} An inflammatory condition that can affect the skin, mouth, genitals, nails, or hair. It often causes itchy purple or reddish bumps and may produce sore or white-patterned patches in the mouth; the cause is often not known.

\synonyms
Lichen Ruber Planus (Lichen Planus)

\medterm{Lichen Planus, Oral} Alternate index label; see \textit{Oral lichen planus}.

\medterm{Lichen Ruber Planus (Lichen Planus)} Alternate index label for Lichen Planus.

\medterm{Lichen Sclerosus} A long-term inflammatory skin condition that usually affects the genital or anal area, causing pale, thin, itchy, easily injured skin. Treatment relieves symptoms and lowers scarring; persistent sores need review because cancer risk is slightly increased.

\synonyms
White Spot Disease

\medterm{Lichen Simplex Chronicus} A thickened, leathery plaque caused by chronic scratching and rubbing (itch-scratch
cycle). It commonly affects the neck, ankles, lower legs, and genital area. It is a form
of neurodermatitis associated with stress and anxiety.

\medterm{Lichenification} Thick, leathery skin with more visible skin lines caused by repeated scratching or rubbing. It often develops in long-lasting eczema, allergies, or other itchy skin conditions.

\medterm{Lichenified} Lichenified skin is thickened and leathery from chronic rubbing or scratching, commonly occurring in longstanding eczema or itch disorders.

\medterm{Lichenoid Drug Eruption} A lichenoid inflammatory skin reaction caused by a medicine, often resembling lichen planus but with a broader distribution or different clinical and histologic features; improvement may follow withdrawal of the culprit drug.

\medterm{Liddle Syndrome} A rare inherited cause of early high blood pressure in which the kidneys retain too much sodium and lose too much potassium. It may begin in childhood or young adulthood and usually requires specialist treatment.

\medterm{Liddle Syndrome Treatment} Treatment of Liddle syndrome lowers blood pressure and the body’s excessive sodium retention, usually with a medicine that blocks the overactive kidney channel. Potassium and kidney function are monitored, and relatives may need testing.

\medterm{Lidocaine} Lidocaine is a local anaesthetic that blocks voltage-gated sodium channels and, intravenously, can treat selected ventricular arrhythmias; excessive exposure may cause seizures, hypotension, arrhythmia, or cardiac arrest.

\medterm{LieberküHn'S Crypts} Tubular glands in the intestinal lining between villi that contain stem, absorptive, and secretory cells. They renew the lining, produce intestinal substances, and help maintain the gut’s protective barrier.
\medterm{Life Expectancy} The average number of additional years a person of a specified age is expected to live under current mortality conditions.


\medterm{Lifetime Risk} The probability that a person will develop a specified disease or experience an outcome during the remainder of their life or over a defined lifetime period.

\medterm{Ligament} A dense, fibrous connective tissue band that connects bone to bone, providing passive
joint stability and guiding joint motion. Ligaments are composed predominantly of type I
collagen fibers (parallel orientation, crimped pattern) and elastin fibers (10-20\% in
some ligaments), arranged in bundles.

\medterm{Ligand} Any molecule that binds specifically to a biological receptor, macromolecule, or active
site. In pharmacology and physiology, ligands include hormones, neurotransmitters,
drugs, and toxins that bind receptors to initiate (agonist), block (antagonist), or
modulate signaling.

\medterm{Ligase} An enzyme that joins two molecules, often using energy from ATP. Ligases help repair DNA, copy genetic material, and build or maintain other important molecules.

\medterm{Ligation} The surgical process of tying off a blood vessel or duct using a suture or clip to
prevent bleeding or block fluid flow.

\medterm{Ligation (Vascular)} Surgical tying, clipping, or sealing of a blood vessel to stop bleeding or redirect blood flow. It may be used during surgery or to treat selected abnormal vessels.
\synonyms
Vascular ligation

\medterm{Ligature} Any cord, wire, rope, fabric, or flexible object used to constrict the neck, bind limbs,
suspend the body, or inflict injury. The resulting pale groove-shaped abrasion may be
horizontal with petechiae in homicidal strangulation, or oblique and upward in hanging.

\medterm{Light} Visible electromagnetic radiation; in medicine it is used for examination, phototherapy, imaging, and selected treatments.
\medterm{Light Coagulation} Light coagulation uses intense focused light, usually a laser, to heat and seal tissue or blood vessels.

\medterm{Light Cycle} A light cycle is the recurring pattern of light and darkness that helps regulate circadian rhythms and biological activity.

\medterm{Light Microscope} An instrument that uses visible light and lenses to magnify cells, tissues, microorganisms, or other small structures. It is widely used in clinical laboratories, pathology, microbiology, and research.

\medterm{Light Microscopy} Examination of cells, tissues, microorganisms, or other specimens using visible light and a microscope. Different stains and preparation methods reveal structure, infection, inflammation, blood cells, or abnormal growth.

\medterm{Light Reflex} Automatic narrowing of the pupil when light enters the eye. It tests the retina, optic nerve, brain pathways, and nerve that controls the pupil; an unequal or absent response may indicate disease.
\medterm{Light Scattering} Deflection of light when it meets particles, cells, molecules, or irregular structures. Measuring the scattered light can help estimate particle size or concentration in laboratory tests and can affect the appearance of tissues on imaging.

\medterm{Light Therapy} Treatment that uses selected wavelengths or intensity of light for conditions such as some skin disorders, seasonal depression, or neonatal jaundice.

\medterm{Light's Criteria} A set of laboratory comparisons used to decide whether fluid around the lung is caused by inflammation or disease in the chest rather than by pressure or fluid imbalance elsewhere. Protein and enzyme levels in the fluid are compared with blood levels; the result helps guide further testing and treatment.

\medterm{Light-Near Dissociation} A pupil abnormality in which the pupils constrict when focusing on a nearby object but respond poorly or not at all to light. It can occur with certain nerve disorders, infections, diabetes, or midbrain disease.

\medterm{Lightening} In late pregnancy, descent of the fetus into the pelvis before labor, which may reduce upper-abdominal pressure but increase pelvic pressure and urinary frequency.

\medterm{Lighter Fluid Poisoning} Harm from swallowing or breathing lighter fluid, a hydrocarbon. The greatest danger is liquid entering the lungs, causing chemical pneumonia, cough, or breathing difficulty; do not induce vomiting and contact poison control immediately.

\medterm{Lightheadedness} A feeling of near-fainting, floating, or unsteadiness without necessarily perceiving that the surroundings are spinning; causes include dehydration, low blood pressure, anemia, medications, and heart or neurologic conditions.

\medterm{Lightning Pains} Brief, stabbing, electric-shock-like pains caused by irritated or damaged nerves. They classically occur in tabes dorsalis but may also accompany shingles, nerve compression, diabetes, or other neurological disorders.
\medterm{Lignocaine} The international nonproprietary name for lidocaine, a local anesthetic and antiarrhythmic drug.
\medterm{Likelihood Ratio} A measure of how strongly a test result supports or argues against a disease. A value above one supports disease, a value below one argues against it, and larger changes are more informative.

\medterm{Lilliputian Hallucination} A visual hallucination of unusually small people, animals, or objects; it can occur with neurologic, psychiatric, toxic, infectious, or sensory conditions.

\medterm{Lily Of The Valley Poisoning} Poisoning from eating lily-of-the-valley plant parts, which contain substances that can disrupt the heartbeat. Nausea, vomiting, dizziness, slow or irregular pulse, confusion, or fainting require immediate emergency care.

\medterm{Limb} An arm or leg, including its bones, joints, muscles, tendons, nerves, blood vessels, and skin. Limb problems may affect movement, sensation, strength, circulation, or function.

\medterm{Limb Bud} A limb bud is an embryonic outgrowth that develops into an arm or leg through coordinated signaling, growth, and tissue differentiation.

\medterm{Limb Development} Limb development is embryonic formation and growth of the bones, joints, muscles, nerves, vessels, and skin of the arms and legs.

\medterm{Limb Reduction Defects} Birth defects in which part or all of an arm or
leg does not form normally. The missing or shortened part may involve a bone,
joint, hand, foot, or the whole limb.

\medterm{Limb Fracture} A break in a bone of an upper or lower limb, assessed according to location, displacement, open or closed status, joint involvement, and associated neurovascular injury.

\medterm{Limb Girdle} A limb girdle is the skeletal connection of a limb to the trunk, such as the shoulder or pelvic girdle.

\medterm{Limb Injury} Trauma affecting an arm or leg, including injury to bone, joint, muscle, tendon, ligament, nerve, blood vessel, or skin.

\medterm{Limb Pain} Pain arising from an arm or leg, with causes including trauma, nerve disease, vascular insufficiency, infection, inflammation, and musculoskeletal disorders.

\medterm{Limb Plethysmography} A noninvasive test measuring changes in limb volume or blood flow, used to evaluate arterial or venous obstruction and sometimes to assess vascular responses.

\medterm{Limb Prosthesis} A limb prosthesis is an artificial device replacing part or all of an arm or leg to restore mobility, function, appearance, or independence.

\medterm{Limb-Girdle Muscular Dystrophies} Limb-girdle muscular dystrophies are inherited muscle diseases causing progressive weakness predominantly around the hips and shoulders, with variable cardiac and respiratory involvement.

\medterm{Limbic Encephalitis} Inflammation of the limbic system, often presenting subacutely with memory impairment,
seizures, psychiatric or behavioral change, and sleep disturbance. It may be autoimmune,
paraneoplastic, infectious, or of unknown cause.

\medterm{Limbic Lobe} A curved region on the inner surface of the brain involved in memory, emotion, motivation, and some responses controlled by the nervous system.
\medterm{Limbic System} A group of connected brain structures involved in memory, emotion, motivation, learning, and some automatic body responses, such as changes in heart rate during stress.

\medterm{Limbs} The arms and legs, including their bones, joints, muscles, nerves, blood vessels, and skin. Limb problems may affect strength, sensation, movement, circulation, balance, daily activities, or independence.
\medterm{Limbus} The border where the clear cornea meets the white outer coat of the eye. It contains cells that help renew the corneal surface and may be affected by injury, inflammation, or disease.

\synonyms
Corneoscleral limbus

\medterm{Limbus Corneae} The corneal limbus is the circular border where the transparent cornea meets the white sclera and conjunctiva, containing stem cells important for corneal repair.

\medterm{Limited Range Of Motion} Limited range of motion is reduced ability to move a joint through its normal extent because of pain, stiffness, weakness, swelling, or structural change.

\medterm{Linagliptin} Linagliptin is a DPP-4 inhibitor used with diet and exercise to improve blood glucose control in type 2 diabetes.

\medterm{Lincomycin} Lincomycin is a lincosamide antibiotic used for selected serious bacterial infections; it can cause antibiotic-associated diarrhea and colitis.

\medterm{Lindane} Lindane is an organochlorine insecticide formerly used for lice or scabies; neurotoxicity and environmental persistence limit its use.

\medterm{Line, Central} A catheter whose tip lies in a large central vein, used to deliver medicines or fluids, monitor central venous pressure, draw blood, or provide specialized nutrition.

\medterm{Line, Central Venous} A central venous catheter inserted into a large vein such as the internal jugular, subclavian, or femoral vein and advanced toward the superior or inferior vena cava.

\medterm{Linea} A line or linear anatomical marking, often used in names such as linea alba or linea nigra.
\medterm{Linea Aspera} Prominent longitudinal posterior femoral shaft ridge from converging muscle attachment
lines. Insertion for adductor magnus, vastus medialis/lateralis, and short head of
biceps femoris. Thickest femoral cortex, resisting bending forces. Diverges distally
into supracondylar lines.

\medterm{Linea Nigra} A dark brown or black vertical line appearing on the abdomen during pregnancy, extending
from the pubic symphysis to the umbilicus or beyond. Caused by hormonal stimulation of
melanocytes, resulting in hyperpigmentation. More prominent in darker-skinned women.

\medterm{Linear} Arranged in a straight or nearly straight line, describing the shape or distribution of a skin lesion or eruption.

\medterm{Linear Accelerator} A machine accelerating electrons to produce high-energy X-rays or electron beams that precisely deliver radiation treatment to selected tumors.
\medterm{Linear Energy Transfer} The amount of energy that radiation deposits as it travels through tissue. Radiation with higher transfer deposits energy more densely and may cause more biological damage for the same total dose.

\medterm{Linear Fracture} A narrow break that runs along the length of a bone without moving the pieces far apart or breaking the bone into many fragments. Treatment depends on the bone, the injury's stability, and whether nearby nerves, blood vessels, or organs are affected.
\medterm{Linear IgA Bullous Dermatosis} An autoimmune blistering disease caused by linear deposition of IgA at the
basement membrane zone. It produces tense blisters, sometimes in annular or "string of pearls" patterns, and may be
idiopathic or drug-induced.

\medterm{Lined Canister} A container with a protective inner lining used to collect or store a medical substance, specimen, waste, or device component. The lining helps prevent leakage, contamination, chemical reaction, or damage to the container; its exact use depends on the equipment.

\medterm{Linezolid} Linezolid is an oxazolidinone antibiotic that blocks formation of the bacterial 70S initiation complex and treats resistant Gram-positive infections; thrombocytopenia, neuropathy, lactic acidosis, and serotonin interactions require caution.

\medterm{Lingua} The anatomical Latin term for the tongue, a muscular organ used for chewing, swallowing, speaking, and tasting. The front two-thirds lies in the mouth and the back third lies in the throat; related terms include sublingual, meaning under the tongue.
\medterm{Lingua Villosa Nigra} A benign condition characterized by elongation and dark pigmentation of the filiform
papillae of the tongue, often associated with tobacco use, poor oral hygiene, or
antimicrobial therapy.

\medterm{Lingual} Relating to the tongue or the tongue side of a body structure. Lingual muscles, nerves, blood vessels, and surfaces help with taste, speech, swallowing, and movement of food.

\medterm{Lingual Artery} A branch of the external carotid artery supplying the tongue, floor of mouth, and
sublingual structures. It arises anterior to the hypoglossal nerve and courses deep to
the hyoglossus muscle. Lingual artery aneurysm or thrombosis is rare but can cause
life-threatening hemorrhage or tongue ischemia.

\medterm{Lingual Thyroid} A lingual thyroid is thyroid tissue at the base of the tongue caused by failure of normal thyroid descent and may obstruct swallowing or breathing or cause hypothyroidism.

\medterm{Lingual Tonsil} The lingual tonsil is lymphoid tissue at the base of the tongue that contributes to immune surveillance of the throat.

\medterm{Lingual Tonsils} Collections of lymphoid tissue at the base of the tongue.

\medterm{Liniment} A liquid or semiliquid preparation rubbed onto the skin, usually to provide temporary relief of muscle or joint discomfort. Some contain irritating or warming substances and should not be used on broken skin, near the eyes, or under tight coverings unless directed.

\medterm{Linitis Plastica} A form of stomach cancer that spreads widely through the stomach wall, making it thick and stiff. It may cause early fullness, weight loss, nausea, vomiting, or vague upper-abdominal discomfort.


\medterm{Linoleic Acid} An essential omega-6 polyunsaturated fatty acid obtained from food. The body uses it to make cell membranes and signaling molecules.
\medterm{Linolenic Acid} A polyunsaturated fatty acid. Alpha-linolenic acid is an essential omega-3 nutrient from plant foods, while other forms occur in animal tissues and have different roles.


\medterm{Liothyronine Sodium} A manufactured form of triiodothyronine, an active thyroid hormone, used in selected thyroid disorders. It acts quickly and can cause fast heartbeat, tremor, anxiety, bone loss, or heart strain if excessive.

\medterm{Lip} The lips are mobile folds of skin and mucosa containing muscle, blood vessels, and sensory nerves; they form the oral seal and assist speech, facial expression, feeding, and swallowing.

\medterm{Lip Cancer} Cancer developing in the tissue of the upper or lower lip. It may appear as a persistent sore, lump, crust, bleeding area, numbness, or color change and needs medical assessment.

\medterm{Lip Disease} A disorder of the lips, such as infection, inflammation, ulceration, trauma, pigmentation change, benign tumor, or malignancy.

\medterm{Lip Infection} Infection of the lip caused by bacteria, viruses, fungi, or an infected wound. It may cause pain, redness, swelling, blisters, cracks, pus, or fever; rapidly spreading swelling or breathing difficulty requires urgent care.
\medterm{Lip Moisturizer Poisoning} Swallowing a small amount of ordinary lip moisturizer usually causes mild stomach upset or loose stools. Products containing camphor, salicylates, or other medicines can cause serious poisoning and require poison-control advice.

\medterm{Lip Neoplasm} A lip neoplasm is an abnormal growth in the lip that may be benign or malignant; persistent sores or lumps require assessment.

\medterm{Lip Swelling} Enlargement of one or both lips caused by local injury, infection, inflammation, dental disease, medication, or angioedema. Sudden swelling of the lips with tongue or throat swelling, breathing difficulty, wheezing, faintness, or widespread hives may indicate a severe allergic reaction and requires emergency care.

\synonyms
Lip edema; lip swelling

\medterm{Lipase} An enzyme that breaks down fats into smaller molecules. Blood lipase testing is commonly used when pancreatitis is suspected, although an abnormal result is not specific in every case.

\medterm{Lipase Test} A lipase test measures a pancreatic digestive enzyme in blood and supports evaluation of acute pancreatitis and selected pancreatic or gastrointestinal disorders.

\medterm{Lipedema} A long-term disorder causing an unusual, usually symmetrical buildup of painful fat in the legs and sometimes arms. The feet and hands are usually spared, while easy bruising, tenderness, and swelling may occur.

\medterm{Lipemia Retinalis} A creamy pink-white appearance of the blood vessels at the back of the eye caused by extremely high blood triglycerides. It signals severe fat buildup in the blood and a high risk of sudden pancreatitis.

\medterm{Lipemic Index} A laboratory estimate of how cloudy a blood sample is because of excess fats. Marked cloudiness can interfere with some test results and may prompt review of the sample or repeat testing.

\medterm{Lipid} A fat or fat-like substance, including triglycerides, cholesterol, phospholipids, and fat-soluble vitamins. Lipids store energy, form cell membranes, help make hormones, and help the body absorb vitamins A, D, E, and K.

\medterm{Lipid Bilayer} The double layer of fat-like molecules forming the basic structure of cell membranes. It separates the cell from its surroundings and controls movement of substances.

\medterm{Lipid Disorder} An abnormal level or handling of fats in the blood or tissues, such as high LDL cholesterol, high triglycerides, or unusually low protective HDL cholesterol.
\medterm{Lipid Emulsion Therapy} Emergency intravenous treatment for severe poisoning from a local anesthetic or selected other fat-soluble substances. The fat may draw some drug away from the heart and brain and provide energy to the heart; it is given under medical supervision.

\medterm{Lipid Metabolism} The processes by which the body digests, transports, stores, makes, and breaks down fats for energy, cell membranes, and hormone production.
\medterm{Lipid Panel} A blood test measuring total cholesterol, LDL cholesterol, HDL cholesterol, and triglycerides. The results help estimate cardiovascular risk and guide lifestyle or medicine decisions; fasting may or may not be needed depending on the purpose.

\medterm{Lipid Profile} A blood test panel measuring lipids and lipoproteins, commonly including total cholesterol, LDL cholesterol, HDL cholesterol, and triglycerides.


\medterm{Lipid Profile Test} A blood test usually measuring total cholesterol, LDL cholesterol, HDL cholesterol, and triglycerides to assess atherosclerotic cardiovascular risk and monitor treatment.

\medterm{Lipid Transport} Movement of fats through blood and tissues by lipoproteins, carrier proteins, or other mechanisms for energy use, storage, membrane formation, and hormone production.

\medterm{Lipidosis} A disorder involving abnormal accumulation or metabolism of lipids in tissues, often due to inherited enzyme defects.
\medterm{Lipids} A group of mostly water-insoluble substances that includes fats, oils, fatty acids, phospholipids, and cholesterol. Lipids store energy, form cell membranes, help make hormones, and help the body absorb vitamins A, D, E, and K.

\medterm{Lipoatrophy} Loss of the layer of fat beneath the skin, causing dents or hollow areas. It can follow repeated injections, inflammation, autoimmune disease, HIV treatment, or inherited disorders and may alter medicine absorption at injection sites.

\medterm{Lipoblastoma} A usually benign tumor of immature adipose tissue occurring mainly in infants and young children; it may be localized or diffuse and can recur if incompletely removed.

\medterm{Lipodystrophy} Abnormal loss or distribution of body fat. It may be inherited or caused by medicines, metabolic disease, or repeated injections and can lead to insulin resistance, diabetes, high triglycerides, or changes in appearance.

\medterm{Lipodystrophy (Acquired Generalized And Inherited Lipodystrophy)} Alternate index label for Acquired Generalized and Inherited Lipodystrophy.

\medterm{Lipodystrophy Syndromes} Rare disorders in which body fat is absent from some areas, present from birth or lost later, and may collect excessively elsewhere. They can cause severe insulin resistance, diabetes, fatty liver, and very high triglycerides.

\medterm{Lipodystrophy, Intestinal} Alternate index label; see \textit{Whipple's disease}.

\medterm{Lipofuscin} Lipofuscin is a yellow-brown, non-degradable pigment made of oxidized lipid and protein residues that accumulates in long-lived cells with ageing and in some degenerative diseases.


\medterm{Lipogranulomatosis} Inflammation in which fat accumulates and triggers organized collections of immune cells called granulomas. It may affect skin, soft tissue, or internal organs and can result from injury, leakage of oily material, or metabolic disease.

\medterm{Lipohypertrophy} A firm or lumpy buildup of fat under the skin, often at repeated insulin-injection sites. Injecting into these areas can change insulin absorption, so injection locations should be rotated regularly.

\medterm{Lipoid Nephrosis} An older name for a kidney disorder, most often in children, in which the kidney filters leak large amounts of protein into the urine. It can cause swelling around the eyes, legs, or abdomen.

\medterm{Lipoid Pneumonia} Lung inflammation caused by oil or fat entering the airways or building up in lung tissue. It may follow inhalation or swallowing of oil-based products and can cause cough, breathlessness, fever, or chest changes on imaging.

\medterm{Lipoidosis} Abnormal buildup of fats within cells or tissues, usually because a metabolic pathway cannot process or transport them normally. It may be inherited or acquired and can damage organs such as the liver, nerves, or kidneys.

\medterm{Lipolysis} Breakdown of stored body fat into fatty acids and glycerol so they can be used for energy. Fasting and stress hormones increase this process, while insulin generally reduces it.

\medterm{Lipoma} A common noncancerous lump made of fat cells, usually soft, movable, painless, and slow-growing. It often occurs under the skin of the trunk, neck, or limbs; a deep, painful, hard, or rapidly growing lump needs assessment.


\medterm{Lipomatosis} A condition in which mature fat tissue grows in multiple or broad areas rather than forming one separate lump. It may be symmetrical, involve the trunk or limbs, and be difficult to remove when it extends between nearby structures.

\medterm{Lipomatous Neoplasm} A growth made partly or mainly of fat cells. It may be noncancerous, like a lipoma, or cancerous, like a liposarcoma; a deep, firm, painful, or rapidly enlarging lump needs specialist assessment.


\medterm{Lipophilicity} Lipophilicity is the tendency of a substance to dissolve in fats or oils rather than water, influencing drug absorption and distribution.

\medterm{Lipopolysaccharide} Lipopolysaccharide is the outer-membrane endotoxin of Gram-negative bacteria, consisting of lipid A, core polysaccharide, and O antigen; circulating lipid A can trigger cytokine release, fever, shock, and disseminated inflammation.

\medterm{Lipoprotein} A particle made of fat and protein that carries cholesterol and other fats through the blood. Different types transport dietary fat, triglycerides, or cholesterol; high LDL levels can contribute to artery disease.

\medterm{Lipoprotein Analysis} A blood test measuring cholesterol, triglycerides, and particles such as LDL and HDL to estimate cardiovascular risk and guide treatment.

\synonyms
Lipid panel; lipid profile

\medterm{Lipoprotein Lipase} An enzyme on capillary walls that breaks down triglycerides in chylomicrons and very-low-density lipoproteins so tissues can use or store the fatty acids. Its activity is regulated by nutritional and hormonal signals.

\medterm{Lipoprotein Metabolism} The body's production, transport, modification, and removal of particles that carry cholesterol and other fats through the bloodstream.

\medterm{Lipoprotein(A)} A cholesterol-carrying particle similar to LDL whose level is largely inherited. A high level can increase the long-term risk of heart attack, stroke, and narrowing of the aortic valve.

\synonyms
Lp(a)

\medterm{Liposarcoma} A malignant tumor of fat cells, usually arising in the deep soft tissues of the thigh or
retroperitoneum; well-differentiated, myxoid, and pleomorphic subtypes are recognized.

\medterm{Liposarcoma Of The Breast} A rare malignant mesenchymal tumor with adipocytic differentiation arising in breast soft tissue. Diagnosis requires histopathologic and molecular evaluation, and treatment is generally wide excision with sarcoma-focused care.

\medterm{Liposomal Amphotericin B} A liposomal formulation of amphotericin B used in the treatment of severe systemic
fungal infections, with reduced nephrotoxicity compared to conventional amphotericin B.
The liposomal formulation encapsulates amphotericin B within lipid vesicles, which
alters the drug's biodistribution and reduces renal tubular damage.

\medterm{Liposome} A liposome is a tiny vesicle made of a lipid bilayer that can carry medicines or other substances into or around the body.

\medterm{Liposuction} A surgical procedure that uses a cannula and suction to remove localized deposits of subcutaneous fat for body contouring. It is not a treatment for obesity or a substitute for diet and exercise. Risks may include bleeding, infection, contour irregularity, fluid accumulation or shifts, sensory changes, and anesthesia-related complications.

\synonyms
Lipoplasty

\medterm{Lips} The paired movable folds surrounding the opening of the mouth, involved in speech, facial expression, eating, and oral sensation.

\medterm{LipschüTz Bodies} An older term for nonspecific acute genital ulcers, often occurring in adolescents or young women after an infection.
\medterm{Liquefaction} Conversion of a solid or partly solid material into liquid, such as tissue breakdown during liquefactive necrosis or melting of a substance.
\medterm{Liquefactive Necrosis} A pattern of tissue death in which enzymes digest dead cells until the area becomes soft or liquid. It commonly occurs in brain injury and some infections and may leave a fluid-filled cavity.

\medterm{Liquid Biopsy} A test of blood, urine, or another body fluid for cancer cells or DNA, RNA, and other material released by a tumor. It may help identify genetic changes, guide treatment, or monitor response and recurrence, but it does not always replace a tissue biopsy.

\medterm{Liquid Chromatography} A laboratory method that separates dissolved substances as a liquid carries them through a material that interacts differently with each substance. It is used to identify or measure medicines, hormones, metabolites, proteins, and toxins.

\medterm{Liquid Diet} A diet consisting partly or entirely of liquids; clear-liquid and full-liquid diets differ in allowed foods and may be prescribed temporarily for procedures, illness, or swallowing or digestive problems.

\medterm{Liquid Incense} A scented liquid product that may contain solvents, oils, or other chemicals. Swallowing or inhaling it can cause irritation, drowsiness, vomiting, seizures, or lung injury; the product container is needed for poison-control assessment.

\medterm{Liquid Nitrogen} Liquid nitrogen is extremely cold nitrogen used to freeze tissue in cryotherapy; contact can cause frostbite and its gas can displace oxygen.

\medterm{Liquor} A liquid; in clinical usage it may refer to cerebrospinal fluid or a medicinal preparation, depending on context.
\medterm{Liraglutide} Liraglutide is a GLP-1 receptor agonist used to improve glucose control in type 2 diabetes and, at some doses, support weight management.

\medterm{Lisch Nodule} A Lisch nodule is a harmless pigmented hamartoma of the iris, commonly associated with neurofibromatosis type 1.

\medterm{Lisch Nodules} Benign, dome-shaped melanocytic hamartomas on the iris, commonly seen in neurofibromatosis type 1.
They usually do not affect vision and support the diagnosis when interpreted with other
clinical findings.

\medterm{Lisinopril} Lisinopril is an angiotensin-converting-enzyme inhibitor used for hypertension, heart failure, and kidney protection in some diabetes; cough, hyperkalaemia, renal-function changes, angioedema, and fetal toxicity are important risks.

\medterm{Lissencephaly} A birth condition in which the brain’s surface has few or no normal folds because nerve cells did not move to their usual positions during development. It can cause seizures, severe developmental impairment, feeding problems, and movement difficulties.

\medterm{Listeria} Listeria is a genus of bacteria; L. monocytogenes can cause foodborne illness, bloodstream infection, and meningitis, especially in pregnancy, older adults, and people with weakened immunity.

\synonyms
Infection Listeria (Listeria)

\medterm{Listeria Infection} An infection caused by Listeria bacteria, usually acquired from contaminated food. It may cause fever, muscle aches, diarrhea, or vomiting; infection can be severe in pregnancy, newborns, older adults, and people with weakened immune systems.

\medterm{Listeria Monocytogenes} A gram-positive, facultatively anaerobic rod that causes listeriosis, a foodborne
infection. The organism is transmitted through contaminated food, particularly deli
meats, soft cheeses, and ready-to-eat foods. It can grow at refrigeration temperatures.

\medterm{Listeriosis} Infection with Listeria monocytogenes, usually acquired from contaminated food. It may cause fever and diarrhea or severe disease such as bloodstream infection or meningitis, especially in pregnancy, newborns, older adults, and people with weakened immunity.

\medterm{Listeriosis In Pregnancy} A foodborne infection that may cause only fever, tiredness, or muscle aches in the pregnant person but can spread to the fetus. It can cause miscarriage, stillbirth, premature birth, or severe infection in the newborn.


\medterm{Lithiasis} The formation or presence of stones (calculi) in an organ or duct. The term is often named for the site, such as nephrolithiasis in the kidney, ureterolithiasis in the ureter, or cholelithiasis in the gallbladder.

\medterm{Lithium} Lithium is a mood stabilizer used mainly for bipolar disorder and selected augmentation treatment; its narrow therapeutic range requires monitoring for tremor, thyroid or renal dysfunction, dehydration, drug interactions, and toxicity.

\medterm{Lithium Carbonate} A mood-stabilizing medicine used mainly to treat bipolar disorder and reduce the risk of manic or depressive episodes; serum levels, kidney function, thyroid function, and drug interactions require monitoring.

\medterm{Lithium Toxicity} Poisoning caused by too much lithium in the body, often after dehydration, illness, a dose change, or a medicine interaction. Nausea, diarrhea, worsening tremor, weakness, confusion, poor coordination, seizures, or reduced consciousness require urgent medical care.

\medterm{Lithotomy} A surgical procedure for removing a stone, or the position with the hips and knees flexed and legs supported.
\medterm{Lithotomy Position} A supine position in which the hips and knees are flexed and the legs are raised and supported to provide access to the pelvis and perineum.

\medterm{Lithotripsy} A treatment that breaks urinary or other stones into smaller pieces using shock waves, laser energy, or another device so the fragments can pass or be removed.

\medterm{Lithotripsy Probe} A device that fragments a stone using mechanical, ultrasonic, laser, or shock-wave energy during a
stone-removal procedure.

\medterm{Lithotripter} A device that breaks urinary, bile-duct, or other stones into smaller pieces using shock waves, laser energy, or another treatment method.
\medterm{Lithotripter Probe} A specialized probe delivering mechanical, electrical, or laser energy to break urinary or bile-duct stones into fragments that can be removed.

\medterm{Litmus} A dye that changes color with acidity or alkalinity and is used as a simple pH indicator.
\medterm{Little Area} The highly vascular region of the anterior nasal septum where several arteries meet and where most anterior
nosebleeds originate.

\medterm{Little'S Disease} An older name for spastic cerebral palsy, a condition caused by early brain injury or abnormal development that produces stiff muscles and movement difficulties, mainly in the legs. It is not progressive, although needs change over time.

\medterm{Live Attenuated Vaccines} Vaccines containing weakened forms of viruses or bacteria. They usually produce a strong, long-lasting immune response without causing illness in healthy people, but some are unsuitable for people with severely weakened immune systems or during certain pregnancies.

\synonyms
Live vaccines

\medterm{Live Birth} A live birth is delivery of an infant showing signs of life such as breathing, heartbeat, or movement, regardless of gestational age.

\medterm{Livedo Racemosa} A persistent, irregular, branching net-like pattern of reddish-blue skin caused by disease of small or medium blood vessels. It may occur with autoimmune, clotting, inflammatory, or other systemic disorders and needs evaluation.

\medterm{Livedo Reticularis} A net-like, violaceous discoloration of the skin caused by altered blood flow in small cutaneous vessels.
It may be temporary and cold-induced or associated with vascular, autoimmune, infectious,
or clotting disorders.

\medterm{Livedoid Vasculopathy} A chronic occlusive disorder of small dermal vessels causing painful purpuric lesions, ulcers, and porcelain-white stellate scars, usually on the lower legs and ankles.

\medterm{Liver} The largest internal organ, beneath the right ribs. It processes nutrients, makes bile to help digest fats, stores energy, helps regulate blood chemicals, and breaks down medicines and harmful substances.

\medterm{Liver (Anatomy And Function)} Alternate index label for Anatomy and Function.

\medterm{Liver Abscess} A pocket of pus inside the liver caused by bacteria, parasites, or another infection. It may cause fever, chills, right-sided upper-abdominal pain, or an enlarged liver and usually needs antibiotics with drainage or other treatment.

\medterm{Liver Biopsy} Percutaneous, transjugular, or laparoscopic sampling of liver tissue for histologic evaluation when noninvasive tests do not answer the clinical question. It may stage fibrosis or characterize inflammation, fatty liver disease, infiltrative disease, or a mass; complications include pain and bleeding.

\medterm{Liver Cancer} A malignant tumor arising in the liver. Hepatocellular carcinoma is the most common primary type in adults; risk factors include chronic liver disease, viral hepatitis, alcohol-related liver injury, and some toxins.

\synonyms
Cancer, Liver

\medterm{Liver Cancer - Hepatocellular Carcinoma} A cancer beginning in hepatocytes, the main liver cells, often associated with cirrhosis or chronic hepatitis.

\medterm{Liver Circulation} The liver receives oxygen-rich blood from the hepatic artery and nutrient-rich blood from the portal vein. Blood passes through liver tissue and leaves through the hepatic veins to the heart.

\medterm{Liver Cirrhosis} Permanent scarring of the liver after long-term injury, causing normal liver tissue to be replaced by stiff scar tissue. Causes include viral hepatitis, alcohol-related disease, fatty liver disease, autoimmune disease, and inherited disorders.

\medterm{Liver Damage} Liver damage is injury or loss of liver-cell function caused by infection, alcohol, medicines, toxins, fatty disease, autoimmune disease, or other processes.

\medterm{Liver Disease} Any condition damaging or impairing liver structure or function. Causes include infections, alcohol, medicines, toxins, fat buildup, immune disease, inherited disorders, blocked bile flow, and cancer.


\medterm{Liver Disease Diet} An individualized eating plan for a person with liver disease. It may address adequate energy and protein, sodium or fluid intake, alcohol avoidance, and nutrient deficiencies according to the condition.

\medterm{Liver Disorder} Any disease affecting the liver, including fatty liver, viral hepatitis, autoimmune disease, scarring, blockage of bile flow, or cancer. Symptoms may be absent or include jaundice, fatigue, pain, swelling, or abnormal blood tests.
\medterm{Liver Dysfunction} Liver dysfunction is reduced ability of the liver to perform tasks such as processing nutrients, making proteins, clearing toxins, and producing bile.

\medterm{Liver Enzymes} Blood measurements including ALT, AST, alkaline phosphatase, and GGT that can suggest liver or bile-duct injury. Abnormal results require interpretation with symptoms, medicines, imaging, and other tests.

\medterm{Liver Failure} A life-threatening loss of the liver's ability to make proteins, process nutrients, remove harmful substances, and regulate blood clotting. It can cause jaundice, bleeding, confusion, swelling, and multiple-organ failure.

\medterm{Liver Failure, Acute} Alternate index label; see \textit{Acute liver failure}.

\medterm{Liver Fibrosis} Excessive scar formation in the liver after ongoing injury or inflammation. It can progress to cirrhosis, reduce liver function, and result from viral hepatitis, alcohol, fatty liver, immune disease, or other causes.

\medterm{Liver Fluke} A parasitic flatworm that infects the liver or bile ducts, including Clonorchis, Opisthorchis, and Fasciola species.
\medterm{Liver Function (Liver (Anatomy And Function))} Alternate index label for Liver (Anatomy and Function).

\medterm{Liver Function Tests} Blood tests used to assess liver injury, bile-flow abnormalities, and liver synthetic function. They commonly
include aminotransferases, alkaline phosphatase, bilirubin, albumin, and clotting tests;
results must be interpreted together with the clinical context.

\medterm{Liver Function Tests (LFTs)} Blood tests assessing liver-cell injury, bile flow, and synthetic function, commonly including aminotransferases, alkaline phosphatase, bilirubin, albumin, and clotting measures.
\medterm{Liver Health} The ability of the liver to perform its many jobs, including processing nutrients and medicines, making important blood proteins, producing bile, storing energy, and helping remove harmful substances. Alcohol, infections, obesity, medicines, and other diseases can damage the liver.

\medterm{Liver Infection} Infection of liver tissue, such as a viral hepatitis infection or a bacterial liver abscess. Symptoms may include fever, tiredness, nausea, right-upper-abdominal pain, jaundice, or confusion and require medical evaluation.
\medterm{Liver Lobule} A small repeating unit of liver tissue in which blood from branches of the portal vein and hepatic artery flows past liver cells toward a central vein, while bile flows in the opposite direction.

\medterm{Liver Mass} An unusual area found in the liver, often on an ultrasound, CT scan, or MRI. It may be a harmless cyst or growth, a primary liver cancer, or cancer that spread from elsewhere; further testing may be needed.

\medterm{Liver Metastases} Liver metastases are cancer deposits that spread to the liver from another primary tumor and may cause no symptoms or produce pain, weight loss, or liver dysfunction.

\medterm{Liver Neoplasm} A liver neoplasm is an abnormal growth in the liver that may be benign or malignant and may arise from liver cells, bile ducts, or metastasis from another cancer.

\medterm{Liver Pain} Discomfort in the upper-right abdomen that may arise from the liver, gallbladder, bile ducts, stomach, muscles, or other organs. Persistent pain, jaundice, fever, vomiting, or confusion requires medical assessment.
\medterm{Liver Regeneration} The liver's ability to restore lost tissue by growing and reorganizing remaining cells after injury or partial surgical removal.
\medterm{Liver Scan} A nuclear medicine or imaging examination of the liver used to evaluate its size, shape, tissue distribution, focal lesions, blood flow, or selected functional abnormalities.

\medterm{Liver Shunt} An abnormal or surgically created channel that diverts blood around the liver, such as a portosystemic shunt; it can reduce detoxification and contribute to hepatic encephalopathy.

\medterm{Liver Spot} A common lay term for a solar lentigo, a flat brown skin patch caused by chronic ultraviolet exposure.
\medterm{Liver Spots} Liver spots are flat brown areas of sun-related pigmentation, also called solar lentigines, and are not caused by liver disease.

\medterm{Liver Transplant} Surgical replacement of a diseased liver with a healthy liver or part of a liver from a donor.

\medterm{Liver Transplantation} Surgery to replace a severely diseased or failing liver with a healthy donor liver or part of one. It may be used for advanced liver disease, sudden liver failure, or selected liver cancers; lifelong medicines help prevent rejection.

\medterm{Liver, Diseases Of} Disorders of the liver, including hepatitis, fatty liver, cirrhosis, tumors, and biliary disease, that can impair metabolism and detoxification.
\medterm{Liver, Enlarged} Alternate index label; see \textit{Hepatomegaly}.

\medterm{Living Donor} Individual who donates an organ (kidney, liver segment, or other tissue) while alive for transplantation into a recipient.

\medterm{Living Will} A type of advance directive that records a person's preferences for future medical treatment if the person cannot communicate or decide.

\medterm{Living Will (Advance Directive)} Alternate index label; see \textit{Living Will}.

\medterm{Living Will (Advance Statements)} Alternate index label; see \textit{Living Will}.

\medterm{Living With A Chronic Illness -  Dealing With Feelings} Living with a chronic illness can bring grief, anxiety, anger, or low mood; support from healthcare professionals, counseling, and trusted people can help.

\medterm{Living With A Chronic Illness - Reaching Out To Others} Reaching out to family, friends, support groups, or healthcare professionals can reduce isolation and help with practical and emotional effects of chronic illness.

\medterm{Living With Endometriosis} Living with endometriosis may involve pain management, hormonal or surgical treatment, pacing activities, fertility counseling, and follow-up for changing symptoms.

\medterm{Living With Hearing Loss} Living with hearing loss may include hearing aids or other devices, communication strategies, hearing protection, and treatment of an underlying cause when possible.

\medterm{Living With Heart Disease And Angina} Living with heart disease and angina involves prescribed medicines, activity and risk-factor management, follow-up, and urgent care for new or worsening chest pain.

\medterm{Living With Uterine Fibroids} Living with uterine fibroids may involve monitoring, medicines for bleeding or pain, procedures, or surgery depending on symptoms, size, location, and fertility plans.

\medterm{Living With Vision Loss} Living with vision loss may involve treating its cause, low-vision rehabilitation, assistive devices, home safety changes, and support with daily activities.


\medterm{Livor Mortis (Lividity)} Post-mortem purplish-red discoloration of dependent body parts resulting from
gravitational settling of blood in uncompressed capillaries and small veins following
cessation of circulation.

\synonyms
Lividity

\medterm{Livor Mortis (Postmortem Lividity)} Purplish-red discoloration of dependent body parts following death, caused by
gravitational pooling of blood in uncompressed capillaries and venules. Starts within 30
min--2 hours, becomes fixed after 6--12 hours.

\synonyms
Postmortem Lividity

\medterm{Lizard Bite} A puncture or tissue injury caused by a lizard; care includes cleaning and assessing for infection, tetanus, venom exposure, and significant tissue damage according to the species and circumstances.

\medterm{LLETZ} A procedure that uses a thin electrically heated wire loop to remove abnormal tissue from the cervix. The tissue is examined in the laboratory and the procedure can treat precancerous cervical changes.

\synonyms
LEEP (Loop Electrosurgical Excision Procedure)

\medterm{Lloyd-Davies Position} A modified lithotomy position in which the hips and knees are flexed and the legs are supported, commonly for pelvic or colorectal surgery.

\medterm{Loa Loa} The African eye worm, a filarial nematode transmitted by Chrysops deer flies, endemic to
the rainforests of West and Central Africa. Adult worms migrate subcutaneously, causing
Calabar swellings, and maycross the subconjunctiva of the eye, causing visible worm
migration. Microfilariae are detected in daytime peripheral blood.

\medterm{Loading Dose} The first, sometimes larger, dose of a medicine given to reach a helpful amount in the body quickly. It is followed by regular maintenance doses and is selected according to the medicine and the patient's condition.

\medterm{Lobaplatin} Lobaplatin is a platinum-based antineoplastic drug used in some countries and studied for its ability to damage cancer-cell DNA.

\medterm{Lobar Hemorrhage} Bleeding into one of the large regions, or lobes, of the brain. It can result from long-term high blood pressure, fragile blood vessels, blood-thinning medicines, injury, or other conditions and may cause sudden headache, weakness, confusion, seizures, or loss of consciousness.

\medterm{Lobar Pneumonia} Pneumonia affecting a large section or an entire lobe of the lung. It can cause fever, cough, chest pain, breathlessness, and low oxygen and may be caused by bacteria, viruses, or other organisms.

\medterm{Lobe} A lobe is a distinct subdivision of an organ or structure, such as a lung, liver, brain, thyroid, or gland.

\medterm{Lobectomy} Surgical removal of a lobe of an organ. Most commonly lung lobectomy (removal of one
lobe of the lung, performed for lung cancer, bronchiectasis, abscess, or benign tumor)
or temporal lobectomy (removal of the anterior temporal lobe for drug-resistant temporal
lobe epilepsy).

\medterm{Lobotomy} A historical brain operation that cut or damaged nerve connections in the front part of the brain. It was once used for severe mental illness but caused serious changes in personality and ability and is no longer an accepted treatment.


\medterm{Lobular Carcinoma} Lobular carcinoma is a cancer arising from cells of a breast lobule, the gland that produces milk; it may be invasive or confined to the lobule.

\medterm{Lobular Carcinoma In Situ} Abnormal cells confined to the breast lobules that increase the risk of developing invasive breast cancer.

\medterm{Lobule} A lobule is a small subdivision of an organ or tissue, such as a liver lobule, lung lobule, or breast lobule.

\medterm{Lobules} Small sections that make up an organ. In the breast, lobules are groups of glands that can make milk; in other organs, the meaning depends on the organ's structure and function.

\medterm{Local} Limited to a particular body site or region rather than involving the whole body. A local treatment or symptom affects the treated or affected area, although nearby or wider effects can sometimes occur.
\medterm{Local Anaesthesia} Loss of sensation in a limited body area caused by a local anesthetic without loss of consciousness.
\medterm{Local Anesthesia} Temporary loss of feeling in a limited area after a local anesthetic is applied or injected, without causing unconsciousness. Too much medicine or accidental injection into a blood vessel can cause seizures, heart problems, or other toxicity.

\medterm{Local Anesthetic Systemic Toxicity} A dangerous reaction when a local numbing medicine enters the bloodstream in excessive amounts. Early signs include tingling around the mouth, ringing in the ears, dizziness, or agitation; seizures and heart problems can follow.

\medterm{Local Anesthetic Toxicity} Harm from too much local anesthetic in the bloodstream. Early warning signs include ringing in the ears, numbness around the mouth, dizziness, or agitation; seizures, low blood pressure, abnormal rhythms, or cardiac arrest can follow and require emergency treatment.

\medterm{Local Invasion} Growth of a tumor into nearby tissues rather than only at its original site. It is a feature of malignant behavior and differs from distant metastasis.

\medterm{Local Therapy} Local therapy treats a disease at or near the affected area rather than throughout the body, for example with a cream, injection, surgery, or focused radiation.

\medterm{Local Treatment} Treatment directed at a specific body site or lesion rather than the whole body, such as a topical medicine, local injection, radiation field, or localized procedure.

\medterm{Localised} Confined to a limited area or a small number of adjacent body sites rather than widespread.

\medterm{Localization} Determining the body site responsible for a symptom, sign, lesion, or function. Accurate localization helps guide examination, imaging, diagnosis, treatment, surgery, rehabilitation, and assessment of nerve or organ damage.
\medterm{Localized} Limited to one particular area, organ, or body site without evidence of distant spread. A localized disease may still be serious, but treatment can sometimes focus on the affected area.

\synonyms
Localized disease

\medterm{Localized Edema} Swelling confined to a particular body region because of fluid accumulation, commonly caused by venous or lymphatic obstruction, inflammation, trauma, allergy, or local infection.

\medterm{Localized Scleroderma} An autoimmune skin disorder causing localized thickening and hardening of the skin, with or without involvement of deeper tissues; it differs from systemic sclerosis.

\medterm{Lochia} The normal vaginal discharge after childbirth. It is usually red and bloody at first, then becomes pink or brown, and finally yellowish or white over the following weeks. Heavy bleeding, large clots, fever, or foul odor needs assessment.

\medterm{Locked In Syndrome (Locked-In Syndrome)} Alternate index label for Locked-in Syndrome.

\medterm{Locked-In Syndrome} A condition in which a person remains conscious and able to think but is almost completely paralyzed and cannot speak. Some people can communicate by blinking or moving their eyes.

\synonyms
Locked In Syndrome (Locked-In Syndrome)

\medterm{Lockjaw} Difficulty opening the mouth because the jaw muscles are tight or painful. It may occur with tetanus, dental infection, jaw injury, or other conditions affecting the jaw muscles or joints.

\medterm{Lockjaw (Tetanus)} A serious neurologic illness caused by tetanospasmin produced by Clostridium tetani. It causes painful muscle stiffness and spasms, often beginning in the jaw.

\synonyms
Tetanus

\medterm{Locomotion} Movement of an organism or person from one place to another using coordinated muscles, bones, nerves, and balance systems.

\medterm{Locomotive System} The musculoskeletal and nervous systems working together to produce posture, balance, and movement from place to place.

\medterm{Locomotor Ataxia} An older term for tabes dorsalis, a neurologic complication of untreated syphilis causing sensory ataxia and lightning pains.
\medterm{Locomotor Ataxia (Tabes Dorsalis)} A late manifestation of untreated tertiary syphilis (neurosyphilis) caused by infection
with Treponema pallidum, involving progressive degeneration of the dorsal columns
(fasciculus gracilis and cuneatus) and dorsal roots of the spinal cord.

\synonyms
Tabes Dorsalis

\medterm{Locoregional} Affecting a defined local area and nearby tissues or lymph nodes, without necessarily involving the whole body.

\medterm{Locus} The specific location of a gene or other DNA sequence on a chromosome. It is described using the chromosome number and a coded position on its short or long arm.

\medterm{Locus Ceruleus} A small brainstem region containing nerve cells that release norepinephrine. It helps regulate alertness, attention, stress responses, sleep, and some automatic body functions.

\synonyms
Locus coeruleus

\medterm{Locus Coeruleus} The locus coeruleus is a brainstem nucleus whose noradrenaline-producing neurons influence arousal, attention, stress, sleep, and pain.

\medterm{Locus Minoris Resistentiae} Latin for a site of lesser resistance, meaning a body area considered more vulnerable to injury, infection, or disease because of local weakness or prior damage.

\medterm{Lod Score} A statistical measure used in genetic linkage analysis to compare how likely it is that two genetic markers are inherited together rather than independently. Higher positive scores provide stronger evidence of linkage in the studied family.

\medterm{Loeys-Dietz Syndrome} Loeys-Dietz syndrome is an inherited connective-tissue disorder associated with arterial aneurysm or dissection, distinctive skeletal and craniofacial features, allergic disease, and variable organ involvement. Lifelong vascular imaging and coordinated specialist care are important.

\synonyms
Syndrome Loeys-Dietz (Loeys-Dietz Syndrome)

\medterm{Loiasis} Loiasis is infection with the filarial worm Loa loa, transmitted by deerflies and sometimes causing transient eye-worm migration and localized swelling.

\medterm{Loin} The area of the lower back between the ribs and pelvis, including the side of the body called the flank. Pain there may come from muscles, the spine, kidneys, or nearby organs.

\medterm{Lomotil Overdose} Taking too much diphenoxylate–atropine can cause extreme sleepiness, confusion, dry mouth, fast heartbeat, slowed breathing, coma, or delayed worsening. It requires immediate emergency or poison-control advice, especially in children.

\medterm{Lomustine} Lomustine is an oral chemotherapy medicine used for selected brain tumors and other cancers; delayed low blood counts are an important risk.

\medterm{Lonafarnib} A farnesyltransferase-inhibiting medicine used for selected progeroid syndromes, in which abnormal protein processing contributes to premature aging features. It can cause gastrointestinal symptoms, liver abnormalities, and drug interactions.
\medterm{Loneliness} Subjective feeling of isolation. Associated with depression, cognitive decline,
mortality.

\medterm{Long Arm Of A Chromosome} The longer segment of a chromosome, designated q, extending from the centromere toward the chromosome's end.

\medterm{Long Bones} Bones that are longer than they are wide, such as the femur, tibia, humerus, and radius. Their shafts provide support and leverage for movement, while their ends form or help form joints; bone marrow inside them produces blood cells.

\medterm{Long COVID} Ongoing or new health problems that begin or continue after COVID-19. Symptoms may include tiredness, worsening after activity, difficulty thinking, breathlessness, changes in smell or taste, sleep problems, or effects on several organs.

\medterm{Long QT Syndrome} A heart rhythm disorder in which the heart takes too long to reset between beats. It may be inherited or caused by medicines or abnormal body salts and can lead to fainting, dangerous rhythms, or sudden death.

\medterm{Long-Acting Beta-2 Agonists} Inhaled medicines that keep the airways relaxed for many hours and reduce breathing symptoms in asthma or chronic obstructive pulmonary disease. For asthma, they are used with an inhaled corticosteroid, not alone.

\medterm{Long-Sightedness} Hyperopia, a refractive error in which distant objects may be clearer than near objects and accommodation is required for focusing.
\medterm{Long-Term Care} Ongoing help with daily activities, nursing, medical needs, or supervision for a person who cannot live fully independently. Care may be provided at home, in assisted living, or in a nursing facility and should match changing needs.

\medterm{Long-Term Care Facility} A residential or care setting providing ongoing assistance, nursing, rehabilitation, or personal care to people who cannot live independently or need sustained support.

\medterm{Long-Term Complications Of Diabetes} Chronic complications of diabetes caused mainly by vascular and metabolic injury, including retinopathy, kidney disease, neuropathy, cardiovascular disease, and foot disease.

\medterm{Long-Term Effect} A health effect that develops or persists after prolonged exposure, treatment, or disease. It may be beneficial, harmful, or uncertain.

\medterm{Long-Term Memory} Information retained for a long period, from hours or days to a lifetime. It includes explicit memories of facts and events and implicit skills or habits.

\medterm{Long-Term Potentiation} Long-term potentiation is a lasting increase in the strength of communication between nerve cells after repeated stimulation and is important in learning and memory.

\medterm{Long-Term Rehabilitation} Ongoing rehabilitation for chronic disease, lasting disability, or persistent injury. Goals include maintaining function, preventing complications, adapting activities, and supporting independence and quality of life.


\medterm{Longitudinal Section} A cut or image made along the long axis of an organ, limb, vessel, or other structure.

\medterm{Longitudinal Studies} Longitudinal studies follow participants or measurements over time to examine change, incidence, exposures, outcomes, and temporal relationships.

\medterm{Longitudinal Vaginal Septum} A wall of tissue that divides the vagina lengthwise, present from birth. It may cause painful sex, difficulty inserting tampons, blocked menstrual flow, or no symptoms. It can occur with an unusually shaped uterus and may be removed surgically when it causes problems.

\medterm{Loop Diuretic} A strong water tablet that makes the kidneys remove extra salt and water in urine. It is used for fluid buildup, heart failure, some kidney problems, and high blood pressure; dehydration and low body salts can occur.
\medterm{Loop Diuretics} A group of strong water tablets that make the kidneys remove salt and water quickly. They treat fluid buildup and some heart, kidney, or liver disorders; excessive use can cause dehydration, low salts, low blood pressure, or kidney problems.

\medterm{Loop Electrosurgical Excision Procedure} A procedure that uses a thin heated wire loop to remove abnormal tissue from the cervix. It both treats precancerous cervical changes and provides tissue for laboratory examination. Bleeding, infection, narrowing of the cervix, and a small increased risk of later premature birth can occur.

\synonyms
LEEP; Large Loop Excision of the Transformation Zone

\medterm{Loop Of Henle} A U-shaped part of each kidney’s tiny filtering tube. It helps the kidney conserve water and control the concentration of urine by allowing water and salts to move differently along its two limbs.

\medterm{Loose Bodies} Free fragments of cartilage, bone, or other tissue within a joint or body cavity that may cause pain, catching, or locking.
\medterm{Loperamide} Loperamide is a peripherally acting opioid-receptor agonist that slows intestinal motility and reduces diarrhoea; it should be avoided in bloody or febrile diarrhoea and excessive doses can cause severe cardiac arrhythmias.

\medterm{Loperamide Hydrochloride} Loperamide hydrochloride is an antidiarrheal medicine that slows movement through the intestine; excessive doses can cause dangerous heart rhythms.

\medterm{Lopinavir} Lopinavir is an HIV protease inhibitor used with ritonavir in antiretroviral treatment; drug interactions are common.

\medterm{Loratadine} An antihistamine used to relieve sneezing, runny nose, itchy or watery eyes, and hives. It usually causes less drowsiness than older antihistamines, but sleepiness, headache, or dry mouth can still occur.

\medterm{Lorazepam} Lorazepam is a benzodiazepine used short term for severe anxiety, seizures, or sedation; it can cause drowsiness, dependence, and breathing suppression.

\medterm{Lordosis} The inward curve of the neck or lower back. A curve that is unusually pronounced may be related to posture, pregnancy, muscle imbalance, spinal changes, or another condition and may cause back discomfort.

\medterm{Lordosis - Lumbar} Lumbar lordosis is the inward curve of the lower spine; excessive or reduced curvature may be postural, structural, degenerative, or related to pain.

\medterm{Losartan} Losartan is an angiotensin-receptor blocker used for high blood pressure, some kidney disease, and heart failure; it can raise potassium levels.

\medterm{Losartan Potassium} Losartan potassium is the salt form of losartan, an angiotensin-receptor blocker used to lower blood pressure and protect selected patients with kidney or heart disease.



\medterm{Loss Of Appetite} Reduced desire to eat. It may result from illness, medicines, mood disorders, pain, or other causes.

\medterm{Loss Of Bladder Control} Alternate index label; see \textit{Urinary incontinence}.

\medterm{Loss Of Brain Function - Liver Disease} Hepatic encephalopathy, a brain dysfunction caused by severe liver disease or portosystemic shunting and characterized by changes in attention, behavior, consciousness, or neuromuscular function.

\medterm{Loss Of Heterozygosity} Loss of heterozygosity is loss of one of two different versions of a gene in a cell, a change that can contribute to cancer when a tumor-suppressor gene is affected.

\medterm{Loss Of Sense Of Smell And Taste With COVID-19} COVID-19 can temporarily impair smell and taste through infection or inflammation of nasal and olfactory pathways. Most people improve, but persistent loss, distortion, poor nutrition, or neurologic symptoms warrants follow-up and smell-retraining advice.

\medterm{Loss Of Smell (Anosmia)} Partial or complete inability to detect odors, commonly caused by nasal obstruction, infection, inflammation, head injury, neurologic disease, or certain medicines.

\medterm{Loss-Of-Function Mutation} A genetic change that reduces or abolishes the normal activity or production of a gene product; its clinical effect depends on the gene, variant, and inheritance pattern.

\medterm{Loss-Of-Function Variant} A genetic variant that reduces or abolishes the normal activity of the encoded gene
product. Disease risk depends on the gene, inheritance pattern, residual function, and
whether one or both copies are affected.

\medterm{Loteprednol Etabonate} Loteprednol etabonate is a corticosteroid eye medicine used for eye inflammation; prolonged use can raise eye pressure and increase infection risk.

\medterm{Lou Gehrig Disease} Amyotrophic lateral sclerosis, a progressive neurodegenerative disease affecting upper and lower motor neurons and causing weakness, muscle wasting, and eventual respiratory impairment.

\medterm{Lou Gehrig'S Disease} Alternate index label; see \textit{Amyotrophic lateral sclerosis}.

\medterm{Louse} A small wingless insect that lives on hair, skin, or clothing and feeds on blood. Human lice include head, body, and pubic lice. They cause itching and may spread through close contact; treatment depends on the type and includes removing live lice and their eggs.

\medterm{Lovastatin} Lovastatin is a statin medicine that lowers LDL cholesterol and reduces cardiovascular risk; muscle injury and liver effects are uncommon but important risks.

\medterm{Low Back Pain} Pain or discomfort in the lower back, commonly caused by muscle or ligament strain, disc disease, arthritis, or spinal
nerve irritation. Weakness, numbness, new bladder or bowel problems, fever, major trauma, or unexplained
weight loss are warning signs requiring prompt medical assessment.

\synonyms
Lumbago (Low Back Pain)

\medterm{Low Back Pain - Acute} Pain in the lower back lasting less than about four weeks, commonly from muscle or ligament strain. Numbness, weakness, fever, injury, cancer history, or bladder problems require urgent assessment.

\medterm{Low Back Pain - Chronic} Lower-back pain lasting three months or longer. Causes include persistent strain, arthritis, disc disease, nerve irritation, inflammatory disease, or less commonly cancer; treatment addresses the cause, function, activity, and pain.

\medterm{Low Birth Weight} A birth weight below 2,500 grams (about 5 pounds 8 ounces), whether caused by being born early, slower growth before birth, or both. Lower weights carry greater risks and need careful newborn monitoring.

\medterm{Low Blood Oxygen (Hypoxemia)} An abnormally low oxygen level in arterial blood, usually related to lung disease, impaired breathing, circulation problems, or reduced oxygen availability.

\medterm{Low Blood Potassium} Low blood potassium, or hypokalemia, can cause weakness, cramps, constipation, abnormal heart rhythms, or paralysis and may result from losses, medicines, or hormonal disease.

\medterm{Low Blood Pressure} Low blood pressure can cause dizziness, fainting, weakness, or shock and may result from dehydration, medicines, bleeding, infection, heart disease, or autonomic dysfunction.

\medterm{Low Blood Pressure (Hypotension) Causes} Alternate index label; see \textit{Hypotension}.

\medterm{Low Blood Sodium} Low blood sodium, or hyponatremia, reflects excess water relative to sodium and can cause headache, confusion, seizures, or coma when severe or rapidly developing.

\medterm{Low Blood Sugar} Low blood sugar, or hypoglycemia, occurs when circulating glucose is insufficient for normal brain and body function, causing sweating, tremor, confusion, seizures, or loss of consciousness.

\medterm{Low Blood Sugar - Newborns} Low blood sugar in a newborn can cause jitteriness, poor feeding, lethargy, apnea, or seizures and is more likely with prematurity, growth problems, maternal diabetes, or illness.


\medterm{Low Body Temperature} Alternate index label; see \textit{Hypothermia}.

\medterm{Low Calcium Level  - Infants} An abnormally low calcium level in a baby. It may cause jitteriness, poor feeding, sleepiness, pauses in breathing, seizures, or abnormal heart rhythms and requires age-appropriate medical assessment.

\medterm{Low FODMAP Diet} A temporary eating plan that limits certain poorly absorbed carbohydrates that can cause gas, bloating, pain, or diarrhea in some people with irritable bowel syndrome. Foods are gradually reintroduced with dietitian guidance to identify personal triggers.

\medterm{Low Grade} Describes a disease or tumor that appears less aggressive or more slowly developing than higher-grade forms. The exact meaning and outlook depend on the specific disease and test used.

\medterm{Low Hemoglobin Count} A hemoglobin level below the expected range, often indicating anemia from blood loss, reduced production, nutritional deficiency, inflammation, or increased destruction of red blood cells.

\medterm{Low Molecular Weight Heparin} A blood-thinning medicine used to prevent or treat blood clots. It is usually injected under the skin and can cause bleeding, bruising, or, rarely, a serious fall in platelets.

\medterm{Low Nasal Bridge} A low nasal bridge is reduced height of the nasal root and may be a normal familial feature or part of a congenital, genetic, or developmental syndrome.

\medterm{Low Neutrophil Count (Neutropenia)} Alternate index label for Neutropenia.

\medterm{Low Platelet Count} Alternate index label; see \textit{Thrombocytopenia}.

\medterm{Low Potassium (Hypokalemia)} An abnormally low potassium level in the blood, which can cause weakness, cramps, abnormal heart rhythms, or paralysis and may result from losses, medicines, or hormonal disorders.

\medterm{Low Sex Drive In Women} A lasting decrease in sexual interest or desire that causes personal distress or difficulty in a relationship. Possible causes include stress, depression, pain, relationship problems, hormonal changes, illness, or medicines; treatment depends on the cause and the person's wishes.

\medterm{Low Sperm Count} A semen condition in which sperm concentration is below the expected reference range. It may reduce the chance of pregnancy, but fertility also depends on sperm movement, shape, ejaculation, and the partner's health.

\medterm{Low Testosterone} Alternate index label; see \textit{Male hypogonadism}.

\medterm{Low Testosterone (Hypogonadism)} A condition in which a person has symptoms of testosterone deficiency and repeatedly low morning testosterone levels. Causes may involve the testes, ovaries, pituitary gland, medicines, obesity, or another illness and need medical evaluation.

\synonyms
Hypogonadism

\medterm{Low Testosterone (Low T)} Alternate index label for Low T.

\medterm{Low Vision} Visual impairment that cannot be fully corrected with spectacles, contact lenses,
medication, or surgery, impairing daily activities. Managed with low-vision aids
(magnifiers, telescopes, adaptive devices), lighting, and occupational therapy.

\medterm{Low White Blood Cell Count} A reduced number of white blood cells, which can increase susceptibility to infection and may result from medicines, infection, autoimmune disease, nutritional deficiency, or bone marrow disorders.

\medterm{Low White Blood Cell Count And Cancer} A reduced number of infection-fighting blood cells in a person with cancer or receiving cancer treatment. The risk of serious infection depends on which cells are low and how severely they are reduced.

\medterm{Low-Calorie Diet} An eating plan that provides less energy than a person usually uses, often to support weight loss or a medical goal. It should still provide adequate protein, vitamins, minerals, fluids, and fiber.
\synonyms
Calorie-restricted diet

\medterm{Low-Density Lipoprotein} A particle that carries cholesterol through the blood. Too much LDL cholesterol can collect in artery walls, forming fatty deposits that narrow arteries and increase the risk of heart attack and stroke.
\medterm{Low-Dose CT Scan} A CT scan performed with less radiation than a standard scan. It is used for purposes such as lung-cancer screening or follow-up of lung nodules, but abnormal findings may require additional tests.

\medterm{Low-Fiber Diet} An eating plan that limits foods such as whole grains, nuts, seeds, raw vegetables, and fruit skins. It may temporarily reduce stool volume for selected bowel problems or medical procedures.

\medterm{Low-Flow Low-Gradient Aortic Stenosis} A severe narrowing of the aortic valve in which less blood moves through the valve and the pressure difference is lower than expected. Echocardiography and other tests help confirm severity.
\synonyms
LF-LG Aortic Stenosis; Low-Gradient Severe Aortic Stenosis

\medterm{Low-Flow Priapism} Alternate index label; see \textit{Priapism}.

\medterm{Low-Grade Appendiceal Mucinous Neoplasm} A mucin-producing appendiceal neoplasm with low-grade cytologic atypia that may dilate the appendix and rupture, disseminating mucin into the peritoneum as pseudomyxoma peritonei.

\medterm{Low-Grade Central Osteosarcoma} A rare low-grade malignant osteogenic tumor arising within the medullary cavity of bone. It may resemble fibrous dysplasia or a benign fibro-osseous lesion and requires complete oncologic excision.

\medterm{Low-Grade Endometrial Stromal Sarcoma} A malignant uterine stromal tumor resembling proliferative-phase endometrial stroma, characterized by infiltrative myometrial or vascular growth and often driven by characteristic gene fusions.
\medterm{Low-Phosphorus Diet} A dietary pattern limiting phosphorus intake for selected patients with chronic kidney
disease or persistent hyperphosphatemia. It emphasizes avoidance of phosphate additives
while maintaining adequate protein and overall nutritional status.

\medterm{Low-Potassium Diet} A dietary pattern limiting potassium intake when clinically indicated, usually for
hyperkalemia or advanced kidney disease. It must be individualized because unnecessary
restriction can reduce diet quality and cause inadequate intake of fruits and
vegetables.

\medterm{Low-Residue Diet} A short-term diet limiting dietary fiber and foods that leave substantial undigested
material in the bowel. It may be used in selected gastrointestinal conditions or around
certain procedures but is not a universal treatment for inflammatory bowel disease.

\medterm{Low-Salt Diet} A low-salt diet limits sodium intake and may help control blood pressure or fluid retention in conditions such as heart failure and kidney disease.

\medterm{Low-Set Ear} An ear positioned lower on the side of the head than expected relative to facial landmarks; it may be an isolated variation or a feature of a genetic or developmental condition.

\medterm{Low-Set Ears And Pinna Abnormalities} Differences in the position, size, shape, or folds of the external ear, including ears positioned lower than usual. They may occur alone or with a congenital syndrome, so hearing and other physical findings may need assessment.

\medterm{Lower Airway Obstruction} Narrowing or blockage of the breathing tubes inside the lungs. Asthma, chronic lung disease, infection, swelling, or a foreign object may cause wheezing, cough, breathlessness, or reduced airflow.

\medterm{Lower Esophageal Ring} A thin ring of tissue near the junction of the esophagus and stomach, often called a Schatzki ring. It can narrow the esophagus and cause intermittent difficulty swallowing solid food or food becoming stuck.

\medterm{Lower Esophageal Sphincter} A ring of muscle where the esophagus joins the stomach. It opens during swallowing and normally closes afterward to limit backward flow of stomach contents and acid.
\synonyms
LES

\medterm{Lower Extremity} The lower extremity is the hip, thigh, knee, leg, ankle, and foot, including their bones, joints, muscles, nerves, vessels, and skin.

\medterm{Lower Extremity Venous Doppler} An ultrasound test of the leg veins that shows blood flow and whether the veins compress normally. It is commonly used to look for a deep-vein blood clot.

\medterm{Lower Leg} The part of the limb between the knee and ankle, containing the tibia, fibula, muscles, nerves, and blood vessels.

\medterm{Lower Motor Neuron Syndrome} A pattern of weakness caused by damage to nerves that carry movement signals from the spinal cord or brainstem to muscles. It may cause muscle wasting, twitching, low muscle tone, and reduced reflexes.

\medterm{Lower Segment Cesarian Section} A cesarean delivery performed through an incision in the lower uterine segment, usually through a low transverse uterine incision.

\medterm{Lower Urinary Tract Symptoms} Problems with storing or passing urine, including urgency, frequency, nighttime urination, weak stream, straining, leakage, or a feeling of incomplete emptying. Causes include infection, prostate enlargement, bladder disorders, nerve disease, medicines, and blockage.

\medterm{Lozenges} Solid medicinal preparations that dissolve slowly in the mouth, usually to provide local treatment for the throat or oral cavity.
\medterm{LSD} Lysergic acid diethylamide is a potent hallucinogen that can cause altered perception, visual hallucinations, autonomic effects, and unpredictable changes in mood and behavior.

\medterm{Lubiprostone} Lubiprostone activates intestinal chloride channels and is used for selected forms of chronic constipation and irritable bowel syndrome with constipation.

\medterm{Lubricant} A substance that reduces friction between surfaces; medical lubricants are used on skin, mucosa, instruments, or devices and should be selected for tissue compatibility and intended use.


\medterm{Lucilia} A genus of blowflies; some species can cause wound myiasis, while sterile larvae have occasionally been used in maggot debridement.
\medterm{Ludwig Angina} A rapidly spreading cellulitis of the floor of the mouth and submandibular spaces,
usually arising from an infected lower molar. Tongue elevation, dysphagia, drooling, and
airway obstruction can develop rapidly and require emergency management.

\medterm{Ludwig'S Angina} A rapidly progressive, life-threatening cellulitis of the submandibular, sublingual, and
submental spaces, typically originating from an odontogenic infection of the lower
second or third molars.


\medterm{Lugano Classification} A system used to describe how far Hodgkin or non-Hodgkin lymphoma has spread and how it responds to treatment. It combines examination, blood tests, and imaging, often PET with CT. Stages range from disease in one area to widespread disease; the results help guide treatment and follow-up.

\medterm{Lugol'S Iodine} An aqueous solution of iodine and potassium iodide used as an antiseptic, in selected thyroid treatments, and in some laboratory procedures.
\medterm{Lugol'S Solution} Lugol's solution is a mixture of iodine and potassium iodide used in selected laboratory, diagnostic, and medical applications; it can irritate tissue and affect thyroid function.

\medterm{Lumbago} A non-specific term for low back pain, typically referring to acute or chronic pain in
the lumbar region (between the lower ribs and the gluteal folds). Lumbago is a symptom,
not a diagnosis.

\medterm{Lumbago (Low Back Pain)} Alternate index label for Low Back Pain.

\medterm{Lumbar} Relating to the lower back, especially the region or five vertebrae between the thoracic spine and sacrum.

\medterm{Lumbar Disc Herniation} Protrusion or extrusion of lumbar disc material through the annulus fibrosus, which may irritate or
compress a nerve root and cause back or leg pain. Symptoms, examination findings, and
imaging determine its significance and treatment.

\medterm{Lumbar Lordosis Lower Spine} Lumbar lordosis is the normal inward curve of the lower spine; excessive or reduced curvature may be associated with posture, muscle imbalance, hip disease, vertebral deformity, or pain. Examination and imaging are guided by symptoms rather than appearance alone.

\medterm{Lumbar MRI Scan} Magnetic resonance imaging of the lumbar spine used to evaluate vertebrae, discs, spinal canal, nerve roots, marrow, and soft tissues in conditions such as radiculopathy or spinal stenosis.

\medterm{Lumbar Plexus} The lumbar plexus is a network of nerves from the upper lumbar spinal roots that supplies the lower abdomen, pelvis, and parts of the thigh and leg.

\medterm{Lumbar Puncture} A procedure that uses a needle in the lower back to collect cerebrospinal fluid, the liquid around the brain and spinal cord. The fluid can be tested for infection, bleeding, inflammation, or cancer; pressure can also be measured, and some medicines can be given. Clinicians check bleeding risk and possible increased pressure in the skull first. Headache, back pain, bleeding, and infection are possible complications.

\synonyms
LP
Spinal Tap

\medterm{Lumbar Radiculopathy} Irritation or pressure on a nerve leaving the lower spine, causing pain that travels into the buttock or leg. It may also cause numbness, tingling, or weakness and is often called sciatica.

\medterm{Lumbar Region} The lumbar region is the lower part of the back between the ribs and pelvis, containing the five lumbar vertebrae.

\medterm{Lumbar Spine} The lower part of the backbone between the ribs and pelvis, usually made of five vertebrae. It bears much of the body’s weight and contains nerves supplying the legs; discs, joints, or nerves there may cause lower-back pain or sciatica.

\synonyms
Lumbar region

\medterm{Lumbar Spine CT Scan} Computed tomography of the lumbar spine that provides detailed images of vertebral bone and alignment and is particularly useful for detecting fractures or other bony abnormalities.

\medterm{Lumbar Stenosis} Narrowing of the spinal canal or nerve openings in the lower back. It can squeeze spinal nerves and cause back pain, leg pain, numbness, weakness, or symptoms that worsen with walking.

\synonyms
Stenosis Lumbar (Lumbar Stenosis)
Stenosis Spinal (Lumbar Stenosis)

\medterm{Lumbar Vertebrae} The five large bones of the lower spine between the ribs and pelvis. They bear much of the body’s weight, allow bending, protect spinal nerves, and can be affected by fractures, arthritis, or disc problems.

\medterm{Lumbar Vertebrae (L1-L5)} The five lumbar vertebrae in the lower back, numbered L1 through L5, between the thoracic spine and sacrum and bearing much of the body's load.
\medterm{Lumbosacral Region} The lumbosacral region is the junction between the lower lumbar spine and the sacrum at the back of the pelvis.

\medterm{Lumbosacral Spine CT} Computed tomography of the lower spine and sacrum using x-rays and computer reconstruction to evaluate fractures, bone lesions, alignment, and selected causes of nerve compression.

\medterm{Lumbosacral Spine X-Ray} Plain radiographic imaging of the lumbar spine and sacrum used to evaluate alignment, fractures, degenerative changes, deformity, and some bone lesions.

\medterm{Lumbricals} Small hand or foot muscles arising from tendons of the flexor digitorum profundus or flexor digitorum longus and assisting flexion-extension coordination.
\medterm{Lumen} The inner space of a hollow organ, vessel, or tube.

\medterm{Luminescence} Luminescence is emission of light from a substance through a process other than high-temperature incandescence, such as fluorescence or chemiluminescence.

\medterm{Luminescence Type} Luminescence type identifies the mechanism by which a material emits light, such as fluorescence, phosphorescence, bioluminescence, or chemiluminescence.

\medterm{Lump In The Abdomen} An abdominal lump may be a hernia, enlarged organ, cyst, stool, inflammation, or tumor; pain, vomiting, fever, growth, or a pulsatile mass requires prompt assessment.

\medterm{Lumpectomy} Breast-conserving surgery that removes a cancer or other abnormal area with a small margin of nearby tissue. Radiation or additional treatment may be recommended afterward, depending on the diagnosis and risk of recurrence.

\medterm{Lumpy Jaw} A chronic granulomatous infection of the jaw, usually cervicofacial actinomycosis, causing swelling and sinus formation.
\medterm{Lunate} A small crescent-shaped bone in the center of the wrist. It helps the wrist move and can be injured by falls; loss of its blood supply can cause Kienböck disease and bone collapse.

\medterm{Lung} One of the two breathing organs in the chest. The lungs take oxygen into the blood and remove carbon dioxide. The right lung has three sections, and the left has two to leave room for the heart.

\medterm{Lung Abscess} A pocket of infection and pus that forms inside damaged lung tissue, often after inhaling stomach contents or because of bacteria. It may cause fever, cough, foul-smelling sputum, chest pain, and breathlessness.

\medterm{Lung Anatomy (Lungs Design And Purpose)} Alternate index label for Lungs Design And Purpose.

\medterm{Lung Cancer} Cancer that begins in the lungs or airways. The main types are small-cell and non-small-cell cancer. Smoking is the leading risk factor, but radon, asbestos, air pollution, inherited risk, and other exposures can also contribute. Early disease may cause no symptoms.
\synonyms
Cancer, Lung
Cancer In Lung (Lung Cancer)

\medterm{Lung Cancer - Small Cell} A fast-growing lung cancer made of small neuroendocrine cells that often spreads early.

\medterm{Lung Collapse} Alternate index label; see \textit{Atelectasis}; see also \textit{Pneumothorax}.

\medterm{Lung Compliance} How easily the lungs stretch when a person breathes in. Stiff lungs require more effort, while overly flexible lungs may not spring back well, as can occur in emphysema.

\medterm{Lung Cyst} A lung cyst is a round air- or fluid-containing space in lung tissue, caused by developmental disease, infection, emphysema, inflammation, or another disorder.

\medterm{Lung Development} Lung development is embryonic and postnatal formation and maturation of airways, alveoli, blood vessels, and gas-exchange capacity.

\medterm{Lung Diffusion Testing} A breathing test that measures how well oxygen-like gas moves from the lung air sacs into the blood. It helps assess diseases such as emphysema, lung scarring, or disorders of the lung blood vessels.

\medterm{Lung Disease} Lung disease includes disorders affecting the breathing tubes, air sacs, lung tissue, covering, blood vessels, or chest wall. These conditions may cause cough, breathlessness, chest discomfort, wheezing, or low oxygen.


\medterm{Lung Disease, Interstitial} Alternate index label; see \textit{Interstitial lung disease}.

\medterm{Lung Fibroma} A usually noncancerous growth made of fibrous tissue in or near the lung. It often causes no symptoms and is found on imaging, but its appearance, growth, and location may require further imaging or biopsy.

\medterm{Lung Fluke} A parasitic flatworm, usually of the Paragonimus group, acquired by eating undercooked freshwater crustaceans. Infection may cause chronic cough, chest pain, blood in sputum, or lung lesions and can rarely affect the brain.
\medterm{Lung Function Tests} Physiologic tests that measure aspects of breathing, including airflow, lung volumes, and gas transfer, to assess respiratory function.


\medterm{Lung Herniation} Protrusion of lung tissue through a weak or torn area of the chest wall, usually after chest trauma or surgery. It may appear as a bulge that enlarges with coughing and can cause pain or breathing difficulty.

\medterm{Lung Injury} Damage to lung tissue or the pulmonary circulation caused by trauma, infection, inhaled substances, inflammation, aspiration, oxygen toxicity, or other disease.

\medterm{Lung Lavage} A procedure in which doctors wash parts of the airways or lungs with sterile fluid and then remove it. It may clear certain material from the lungs or collect a sample for testing.

\medterm{Lung Metastases} Lung metastases are cancer deposits that have spread to the lungs from another primary tumor and may cause cough, breathlessness, chest pain, or no symptoms.

\medterm{Lung Needle Biopsy} CT- or ultrasound-guided percutaneous sampling of a lung or pleural lesion to evaluate selected masses, infection, or diffuse disease. Pneumothorax and pulmonary hemorrhage are important complications.

\medterm{Lung Neoplasm} An abnormal growth in the lung or airways that may be noncancerous or cancerous. It may cause cough, blood in sputum, chest pain, breathlessness, or no symptoms; imaging and sometimes tissue sampling determine its nature.

\medterm{Lung Nodule} A rounded, usually well-defined opacity in the lung measuring up to 3 cm on imaging. It may be benign, infectious, inflammatory, or an early manifestation of lung cancer and often requires follow-up assessment.

\medterm{Lung PET Scan} Positron-emission tomography, often combined with computed tomography, using a radiotracer to assess metabolic activity in lung nodules, cancer, infection, or inflammation.

\medterm{Lung Plethysmography} A breathing test performed in a sealed room to measure how much air the lungs hold, including air that cannot be breathed out, and how much resistance the airways provide.

\medterm{Lung Problems And Volcanic Smog} Volcanic smog contains gases and fine particles that irritate eyes and airways and can worsen asthma or lung disease; limiting exposure and following public-health advice reduce harm.

\medterm{Lung Protective Ventilation} A ventilator approach that uses gentle breath sizes and carefully controlled pressures to reduce further lung injury. It is especially important in acute respiratory distress and requires adjustment to oxygen levels, blood gases, and the person’s condition.

\medterm{Lung Surgery} An operation on the lungs or nearby structures to remove disease, repair injury, obtain tissue, or improve breathing. Procedures range from removing a small area to an entire lung; bleeding, infection, air leakage, and breathing problems are possible.


\medterm{Lung Transplant} Surgery to replace one or both severely diseased lungs with donor lungs. It may be considered for advanced lung failure from conditions such as emphysema, lung scarring, cystic fibrosis, or pulmonary hypertension and requires lifelong immune-suppressing medicines.

\medterm{Lung Transplantation} Surgical replacement of a diseased lung with a healthy lung from a donor. Indications
include IPF, COPD, CF, PAH, and bronchiectasis. Options include single lung transplant,
double lung transplant, or heart-lung transplant. Post-transplant immunosuppression is
required lifelong. Complications include rejection and infection.

\medterm{Lung, Diseases Of} Conditions affecting the airways, alveoli, pleura, or pulmonary circulation, including infection, asthma, COPD, fibrosis, and cancer.
\medterm{Lung Volume Reduction Surgery} An operation that removes the most damaged, overinflated parts of the lungs in carefully selected people with severe emphysema. This gives healthier lung tissue and the diaphragm more room to work, which may improve breathing and activity. It has important risks, including air leakage, infection, and respiratory failure.

\synonyms
LVRS; Reduction Pneumoplasty; Bilateral Lung Volume Reduction

\medterm{Lungs} The paired respiratory organs in the chest that exchange oxygen and carbon dioxide between air and blood; they also contribute to acid--base regulation.

\medterm{Lungs Design And Purpose} The lungs exchange oxygen and carbon dioxide through branching airways and microscopic alveoli surrounded by capillaries. The airways filter, warm, and humidify air, while respiratory muscles drive ventilation; disease can impair airflow, diffusion, or blood flow.

\synonyms
Lung Anatomy (Lungs Design And Purpose)

\medterm{Lungs Fluid (Pleural Effusion (Fluid In The Chest Or On Lung))} Alternate index label for Pleural Effusion (Fluid In the Chest or On Lung).

\medterm{Lunula} The pale crescent-shaped area at the base of a fingernail or toenail. It is the visible front part of the nail-forming region, although its size and visibility vary normally.
\medterm{Lupus} A group of autoimmune diseases in which the immune system attacks the body's own
tissues. Systemic lupus erythematosus may affect the skin, joints, kidneys, blood cells,
nervous system, and other organs. Also called SLE.

\synonyms
SLE.

\medterm{Lupus (Systemic Lupus)} Alternate index label for Systemic Lupus.



\medterm{Lupus Erythematosus} A long-term autoimmune disease in which the immune system attacks the body’s own tissues. It can cause tiredness, joint pain, rashes, and inflammation or damage in organs such as the kidneys, lungs, heart, brain, or blood vessels.

\medterm{Lupus Erythematosus (Systemic Lupus Erythematosus - SLE)} A chronic autoimmune disease that can affect many parts of the body, including the skin, joints, kidneys, heart, lungs, nerves, and blood cells. Symptoms often flare and improve, and treatment depends on the organs involved.

\synonyms
Systemic Lupus Erythematosus - SLE

\medterm{Lupus Nephritis} Kidney inflammation caused by systemic lupus erythematosus. It can cause blood or protein in the urine,
high blood pressure, and reduced kidney function. Classification is based largely on
kidney-biopsy findings, and treatment depends on the class, severity, kidney function,
and individual risks.

\medterm{Lupus Panniculitis} A chronic inflammatory form of cutaneous lupus affecting subcutaneous fat, producing firm deep nodules that may heal with lipoatrophy, scarring, or overlying discoid changes.

\medterm{Lupus Pernio} Persistent red-purple or blue-purple swelling and discoloration of the nose, cheeks, lips, ears, or fingers, usually associated with sarcoidosis. It often indicates longer-lasting disease and may cause scarring or cosmetic changes.

\medterm{Lupus Vulgaris} A long-lasting form of tuberculosis affecting the skin. It usually causes slowly enlarging reddish-brown patches or plaques that may scar and is treated with a combination of tuberculosis medicines.


\medterm{Luteal Cell} A cell in the ovarian corpus luteum that produces progesterone after ovulation and supports the uterine lining.

\medterm{Luteal Phase Defect} A term used when the part of the menstrual cycle after ovulation may not provide enough progesterone or enough time for the uterine lining to support pregnancy. The diagnosis is uncertain and should not be made from cycle length alone.

\synonyms
Luteal Phase Deficiency; LPD

\medterm{Lutein} A yellow plant pigment found in leafy green vegetables and concentrated in the retina. Along with related pigments, it helps filter blue light and may protect eye tissues from oxidative damage.

\medterm{Luteinization} Luteinization is the transformation of an ovarian follicle after ovulation into a corpus luteum that produces progesterone.

\medterm{Luteinizing Hormone} A hormone released by the pituitary gland. In people with ovaries, a mid-cycle rise triggers ovulation and supports hormone production; in people with testes, it stimulates testosterone production. Blood testing helps assess puberty, fertility, and pituitary or reproductive disorders.

\medterm{Luteinizing Hormone Blood Test} A blood test measuring luteinizing hormone, which stimulates ovulation and ovarian hormone production in females and testosterone production in males; it helps evaluate puberty, fertility, and pituitary or gonadal disorders.

\medterm{Luteinizing Hormone-Releasing Hormone} An older name for gonadotropin-releasing hormone (GnRH), a hypothalamic hormone released in pulses that stimulates the anterior pituitary to secrete luteinizing hormone and follicle-stimulating hormone.

\medterm{Luteolysis} Luteolysis is the breakdown of the corpus luteum, causing progesterone levels to fall when pregnancy has not occurred.

\medterm{Luteoma} A luteoma is a usually benign ovarian mass of hormone-producing lutein cells, sometimes associated with pregnancy and androgen excess.

\medterm{Luteoma Of Pregnancy} A benign, pregnancy-associated ovarian proliferation of luteinized stromal cells driven by hormonal stimulation. It may cause maternal or fetal androgen excess and usually regresses after delivery.

\medterm{Luteotropic Hormone} A hormone that supports the corpus luteum; the term traditionally refers mainly to prolactin, although usage varies by species and source.
\medterm{Lutetium-177} A radioactive form of lutetium that releases therapeutic radiation and a small amount of detectable radiation. Attached to a targeting molecule, it can deliver treatment to selected neuroendocrine or prostate cancer cells while limiting exposure to other tissues.

\medterm{Luxation} Complete dislocation of a joint, where the bone ends lose normal contact. It can stretch or tear nearby ligaments, nerves, and blood vessels, causing severe pain, deformity, swelling, and inability to move.

\medterm{LVAD} Left-ventricular assist device, an implanted pump helping the heart’s main chamber circulate blood. It may support people with severe heart failure while awaiting transplant or as long-term treatment.

\synonyms
Left-ventricular assist device

\medterm{LVNC} Alternate name; see \textit{Left Ventricular Noncompaction}.

\medterm{LVRS} Alternate name; see \textit{Lung Volume Reduction Surgery}.

\medterm{Lyase} An enzyme that breaks or joins certain chemical bonds without adding water or using oxidation as the main step. Lyases participate in many processes that make or use energy in cells.


\medterm{Lyme Disease} A bacterial infection spread by infected ticks; blacklegged ticks transmit it in the United States. Early symptoms may include an expanding rash, fever, tiredness, headache, and muscle or joint aches; untreated infection can later affect nerves, joints, or the heart.

\synonyms
Arthritis Lyme (Lyme Disease)


\medterm{Lyme Disease Blood Test} A blood test that looks for immune proteins made after infection with Lyme disease bacteria. Results depend on symptoms, tick exposure, and timing; early testing can be negative and a positive result does not always prove active infection.


\medterm{Lymph} A clear fluid collected from body tissues and carried through lymph vessels. It transports immune cells and excess fluid to lymph nodes and eventually returns it to the bloodstream.

\medterm{Lymph Leakage} Lymph leakage is escape of lymphatic fluid after injury, surgery, or obstruction and may cause swelling, fluid collections, or chylous drainage.

\medterm{Lymph Node} A lymph node is a small immune organ that filters lymph and helps activate immune responses to infection, inflammation, or abnormal cells.

\medterm{Lymph Node / Gland} A small immune organ along a lymph vessel that filters lymph and helps the body recognize infection, inflammation, or abnormal cells. Nodes are found in groups in the neck, armpits, chest, abdomen, and groin.

\medterm{Lymph Node Biopsy} Removal of part or all of a lymph node for microscopic and sometimes microbiologic or molecular examination to evaluate malignancy, infection, or inflammatory disease.

\medterm{Lymph Node Culture} Laboratory culture of lymph-node tissue or aspirate to identify bacteria, fungi, or mycobacteria causing lymphadenitis; culture is interpreted with histology and other tests.

\medterm{Lymph Node Dissection} Surgery that removes one or more groups of lymph nodes to check for cancer or treat cancer spread. Removed tissue is examined in a laboratory, and the operation can cause swelling or numbness.

\medterm{Lymph Node Structure} A lymph node is a small, bean-shaped immune organ with a capsule, spaces that carry lymph, and areas rich in immune cells. It filters lymph and may enlarge during infection, inflammation, or cancer spread.

\medterm{Lymph Nodes} Small immune organs filtering lymph fluid and helping detect infections or abnormal cells. They may enlarge or become tender with infection, inflammation, immune disease, or cancer spread.

\synonyms
Lymph glands

\medterm{Lymph Nodes, Swollen} Alternate index label; see \textit{Swollen lymph nodes}.

\medterm{Lymph Swollen Glands (Swollen Lymph Nodes)} Alternate index label for Swollen Lymph Nodes.

\medterm{Lymph Swollen Nodes (Swollen Lymph Nodes)} Alternate index label for Swollen Lymph Nodes.

\medterm{Lymph System} The lymph system includes lymph, lymphatic vessels, lymph nodes, spleen, thymus, and related tissues that return fluid and support immune defense.

\medterm{Lymphadenectomy} Surgical removal of lymph nodes, either sentinel or regional, for staging or therapeutic
purposes; examples include axillary, inguinal, cervical, and retroperitoneal lymph node
dissection.

\medterm{Lymphadenitis} Inflammation or infection of one or more lymph nodes, causing swelling and sometimes tenderness, warmth, or pus. Causes include viruses, bacteria, cat-scratch disease, tuberculosis, and other disorders.

\medterm{Lymphadenitis, Mesenteric} Alternate index label; see \textit{Mesenteric lymphadenitis}.

\medterm{Lymphadenopathy} Enlargement of one or more lymph nodes. It may be localized or widespread and can result from infection, inflammation, immune disease, medicines, or cancer; persistent, hard, or unexplained nodes need assessment.

\medterm{Lymphangiectasis} Abnormal widening of lymph vessels. It may be present from birth or develop later and can cause swelling, leakage of lymph, loss of protein through the intestine, or fluid around the lungs.

\medterm{Lymphangiogenesis} Formation of new lymph vessels from existing vessels. This can help repair tissue and drain fluid, but it may also contribute to chronic swelling, inflammation, or the spread of cancer.

\medterm{Lymphangiogram} A lymphangiogram uses contrast imaging to show lymphatic vessels and nodes, helping evaluate obstruction, leakage, malformation, or spread of disease.

\medterm{Lymphangiography} An imaging test that uses contrast dye to show lymphatic vessels and their flow. It can help find blocked vessels, leaks, or abnormal connections, although newer scans are often preferred.

\medterm{Lymphangioleiomyomatosis} A rare lung disease in which unusual cells grow in the lungs and lymphatic vessels, causing many small cysts. It mainly affects women and can lead to breathlessness, collapsed lung, bleeding, or kidney tumors.

\medterm{Lymphangioma} A noncancerous birth defect in which lymph vessels form abnormally and become enlarged. It usually appears in infancy or childhood as a soft swelling, often in the neck, armpit, or chest.

\medterm{Lymphangiomatosis} A rare disorder in which abnormal lymph vessels grow in one or more organs. It may cause swelling, fluid buildup, bleeding, bone damage, or breathing problems, depending on the tissues involved.

\medterm{Lymphangiomyoma} A rare, usually noncancerous growth made of smooth-muscle-like cells around lymph vessels. It is most often associated with lymphangioleiomyomatosis and may contribute to cysts or blockage of lymph flow.
\medterm{Lymphangiosarcoma} A rare, fast-growing cancer arising from the lining of lymph vessels. It may develop in tissue with long-standing lymph swelling and can appear as a bruise-like, purple, or ulcerated skin lesion.

\medterm{Lymphangitis} Inflammation and infection of lymphatic vessels, usually spreading from an infected skin wound or cellulitis; tender red streaks, fever, and enlarged regional lymph nodes may occur.

\medterm{Lymphatic Basin} The group of lymph vessels and lymph nodes that drains a particular body area. Doctors use this drainage pattern to examine nearby nodes and assess possible spread of infection or cancer.

\medterm{Lymphatic Disease} Any disorder affecting lymph vessels, lymph nodes, or related organs. It may cause swelling, repeated infection, abnormal fluid collections, or cancer and may result from birth defects, injury, infection, inflammation, or tumor growth.
\medterm{Lymphatic Filariasis} A mosquito-borne infection caused by thread-like worms that live in lymph vessels. Repeated infection can block lymph flow and cause long-term swelling of the legs, arms, breasts, or genitals.

\medterm{Lymphatic Mapping} Imaging or tracer-guided identification of lymphatic vessels and lymph nodes that drain a body area. It helps plan surgery, locate a sentinel node, investigate lymphedema, or assess possible cancer spread.

\medterm{Lymphatic Obstruction} Blocked or damaged lymph vessels or nodes that prevent normal fluid drainage. The affected area may develop persistent swelling called lymphedema.

\medterm{Lymphatic Obstruction (Lymphedema)} Alternate index label for Lymphedema.

\medterm{Lymphatic System} A network of vessels, lymph nodes, and organs that returns excess tissue fluid to the bloodstream, transports some digested fats, and helps the immune system detect and fight infection.

\medterm{Lymphatic Tissue} Lymphatic tissue contains immune cells and structures that filter lymph, trap antigens, and coordinate immune defense.

\medterm{Lymphatic Vessel} A lymphatic vessel transports lymph from tissues through lymph nodes and eventually returns fluid to the bloodstream.

\medterm{Lymphedema} Long-lasting swelling caused by poor drainage of lymph fluid. It most often affects an arm or leg and can make the skin feel tight, increase infection risk, and limit movement.

\synonyms
Lymphatic Obstruction (Lymphedema)

\medterm{Lymphedema (Lymphoedema)} Chronic swelling of a body part due to impaired lymphatic drainage, leading to
accumulation of protein-rich interstitial fluid. Primary lymphedema: caused by
congenital or hereditary lymphatic maldevelopment.

\synonyms
Lymphoedema


\medterm{Lymphedema Management} Measures used to reduce swelling and preserve function, including compression, exercise, skin care,
and selected forms of manual lymphatic therapy.

\medterm{Lymphoblast} An immature cell that normally develops into a type of infection-fighting white blood cell. Large numbers of abnormal lymphoblasts can build up in the blood or bone marrow in acute lymphoblastic leukemia.

\medterm{Lymphoblast Count} A lymphoblast count measures immature lymphoid cells in blood or bone marrow and may help evaluate or monitor acute lymphoblastic leukemia.

\medterm{Lymphocele} A lymphocele is a collection of lymphatic fluid caused by disruption or blockage of lymph vessels, often after surgery or transplantation.

\medterm{Lymphocyte} A type of white blood cell that recognizes infections and abnormal cells and helps coordinate immune defense. The main types are B cells, which make antibodies, and T cells, which direct or carry out cellular immune responses.

\medterm{Lymphocyte Activation} The process by which a lymphocyte recognizes a target and begins responding. Activated lymphocytes multiply and help coordinate antibody or cellular immune defense.
\medterm{Lymphocyte Count} A lymphocyte count measures lymphocytes in the blood and can help assess infection, immune deficiency, inflammation, or blood disorders.

\medterm{Lymphocyte Proliferation} Rapid multiplication of lymphocytes after immune stimulation. It helps build a response to infection, but abnormal or uncontrolled proliferation can occur in lymphoid cancers.
\medterm{Lymphocytic} Relating to lymphocytes or containing an increased proportion of these white blood cells. A lymphocytic pattern may occur with viral infection, chronic inflammation, autoimmune disease, or certain cancers.

\medterm{Lymphocytic Colitis} A form of microscopic colitis causing long-lasting, watery diarrhea. The colon may look normal during examination, but a biopsy shows increased immune cells in its lining; medicines, autoimmune disease, or other factors may contribute.

\synonyms
Colitis Collagenous (Lymphocytic Colitis)
Colitis Lymphocytic (Lymphocytic Colitis)
Colitis Microscopic (Lymphocytic Colitis)

\medterm{Lymphocytic Interstitial Pneumonia} A rare inflammatory lung disease in which immune cells collect in the tissue around the air sacs. It may cause cough, breathlessness, and abnormal scans and is associated with autoimmune disease or HIV.

\medterm{Lymphocytic Leukemia} Lymphocytic leukemia is a blood cancer in which abnormal lymphoid cells multiply in the bone marrow, blood, lymph nodes, or other tissues.

\medterm{Lymphocytic Mastitis} A noncancerous breast lump caused by inflammation and toughening of breast tissue. It is associated with diabetes or autoimmune disease and can look like breast cancer on examination or scans, so testing may be needed.

\medterm{Lymphocytoma Cutis} Lymphocytoma cutis is a benign skin lymphoid proliferation presenting as red or violaceous plaques or nodules and sometimes associated with infection or immune stimulation.

\medterm{Lymphocytopenia} An abnormally low number of lymphocytes, a type of white blood cell that helps fight infection. It may be temporary after illness or stress, or result from medicines, immune disorders, infections, poor marrow function, or inherited conditions. Severity and duration determine the risk of infection and need for evaluation.

\medterm{Lymphocytosis} An abnormally increased number or proportion of lymphocytes in the blood, often associated with viral infection, some bacterial infections, or lymphoproliferative disorders.

\medterm{Lymphocytosis (High Lymphocyte Count)} An increased number of lymphocytes in the blood, often associated with infection, inflammation, or a lymphoproliferative disorder.

\medterm{Lymphoedema} Chronic swelling caused by impaired lymphatic drainage, most often affecting a limb but potentially involving other regions.
\medterm{Lymphofollicular Hyperplasia} Lymphofollicular hyperplasia is enlargement or increased formation of lymphoid follicles in response to immune stimulation, infection, inflammation, or other triggers.

\medterm{Lymphogranuloma Venereum} A sexually transmitted infection caused by certain types of Chlamydia. It may begin with a small painless sore, followed by painful groin swelling or inflammation of the rectum, sometimes with bleeding or discharge.

\medterm{Lymphoid} Resembling or relating to lymphatic tissue, lymphocytes, or organs involved in lymphocyte development and immune function.

\medterm{Lymphoid Hyperplasia} Lymphoid hyperplasia is benign increase in lymphoid tissue caused by immune stimulation, infection, inflammation, or sometimes a reactive response near a tumor.

\medterm{Lymphoid Tissue} Tissue containing lymphocytes and supporting immune responses, found in lymph nodes, spleen, thymus, tonsils, and mucosal sites.
\medterm{Lymphokine} An older term for a chemical messenger released by lymphocytes, a type of white blood cell. These messengers help immune cells communicate, multiply, and coordinate responses to infection or abnormal cells.

\medterm{Lymphoma} A cancer of lymphocytes, a type of white blood cell. Hodgkin and non-Hodgkin lymphoma include many subtypes that may grow slowly or quickly and can cause swollen lymph nodes, fever, night sweats, weight loss, or organ symptoms.

\medterm{Lymphoma (Hodgkin'S And Non-Hodgkin'S)} A group of cancers arising from lymphocytes, including Hodgkin lymphoma and the
many types of non-Hodgkin lymphoma.

\synonyms
Hodgkin's and Non-Hodgkin's

\medterm{Lymphoma, Hodgkin} Alternate index label; see \textit{Hodgkin lymphoma (Hodgkin disease)}.

\medterm{Lymphoma, Non-Hodgkin} Alternate index label; see \textit{Non-Hodgkin lymphoma}.

\medterm{Lymphomatoid Papulosis} Lymphomatoid papulosis is a chronic recurrent skin disorder with cancer-like lymphoid cells; the bumps often heal spontaneously but require monitoring for lymphoma.

\medterm{Lymphopenia} Lymphopenia is an abnormally low number of lymphocytes, caused by infection, medicines, immune disease, malnutrition, marrow disorders, or other conditions.

\medterm{Lymphoplasmacytic Lymphoma} Alternate index label; see \textit{Waldenstrom macroglobulinemia}.

\medterm{Lymphopoiesis} The production and development of lymphocytes, a group of white blood cells that help fight infection and abnormal cells. These cells develop mainly in the bone marrow and thymus.

\medterm{Lymphoproliferative Disorder} A condition in which lymphocytes multiply excessively or abnormally. It may be temporary after infection or immune stimulation, or may represent a cancer such as leukemia or lymphoma.
\medterm{Lymphoscintigraphy} A scan that follows a small radioactive tracer through lymph vessels. It can map drainage to the first lymph node draining a tumor or show impaired lymph flow in lymphedema and other disorders.

\medterm{Lynch Syndrome} An inherited condition that greatly increases the risk of colorectal, womb, ovarian, stomach, urinary-tract, and several other cancers. Regular screening can detect some cancers early, and relatives may need genetic counseling or testing.

\medterm{Lysergic Acid Diethylamide} A potent hallucinogenic drug (psychedelic) classified as a Schedule 1 controlled
substance (high abuse potential, no accepted medical use). LSD acts as a partial agonist
at serotonin 5-HT2A receptors in the prefrontal cortex, leading to altered sensory
perception, mood, and cognition.

\medterm{Lysine} Lysine is an essential amino acid needed for protein production, growth, and tissue repair; it must be obtained from food.


\medterm{Lysis} The breakdown, dissolution, or destruction of cells or tissue; examples include hemolysis, the breakdown of red blood cells, and thrombolysis, the dissolution of a blood clot.

\medterm{Lysogenic Conversion} A change in a bacterium's characteristics after a virus that infects bacteria inserts its genetic material into the bacterial chromosome and provides new genes, such as genes that increase disease-causing ability.

\medterm{Lysogeny} Lysogeny is a state in which a bacteriophage's genetic material is maintained inside a bacterial cell and replicated with the host genome.

\medterm{Lysophagy} A form of selective autophagy that removes damaged lysosomes, the cell compartments that digest waste. It helps maintain cell health; failure of lysophagy can allow damaged lysosomes to accumulate and promote inflammation or cell injury.

\medterm{Lysophosphatidylcholine} A fat-like molecule found in cell membranes and blood. It helps transport fats and can influence inflammation and cell communication; abnormal levels may occur in metabolic or inflammatory diseases.

\medterm{Lysophospholipids} Lipids related to phospholipids but containing one fewer fatty-acid chain. They are parts of cell membranes and can act as signaling molecules, influencing inflammation, blood vessels, immune responses, and cell movement.

\medterm{Lysosome} A small compartment inside a cell containing substances that break down waste, worn-out cell parts, and material brought into the cell. Lysosome disorders can cause harmful material to accumulate in tissues.

\medterm{Lysozyme} A natural enzyme that helps defend against bacteria by breaking down part of their protective outer wall. It is found in tears, saliva, breast milk, and some immune cells.
\medterm{Lysozyme Test} A lysozyme test measures the enzyme lysozyme in a sample and may support evaluation of selected blood, inflammatory, or infectious disorders.


\medterm{Lyssavirus} A group of viruses that includes rabies virus and related viruses. They mainly spread through the bite or saliva of an infected mammal and can cause a severe brain infection. Prompt wound cleaning and exposure-appropriate vaccination can prevent disease before symptoms begin.

\medterm{Lytic} Causing or involving lysis, or the breakdown of cells, tissue, or another biological structure; a lytic bone lesion is an area of bone destruction.

\medterm{Laparoscopic Surgery} Surgery performed through small cuts using a camera and long instruments. It often causes less pain and faster recovery than open surgery, but still has risks such as bleeding, infection, and injury to nearby organs.

\medterm{Leber's Hereditary Optic Neuropathy} An inherited mitochondrial disorder that causes painless loss of central vision, usually first in one eye and then the other. It mainly affects young men, and genetic counseling and specialist eye care may help.

\medterm{Leigh's Syndrome} A rare progressive neurologic disorder, usually beginning in infancy or childhood, caused by problems with cellular energy production. It can damage the brainstem and cause developmental regression, movement problems, seizures, or breathing difficulties.

\medterm{Leprosy} A long-term infection caused mainly by Mycobacterium leprae that affects skin and peripheral nerves. Early treatment with multiple antibiotics prevents nerve damage and disability; casual contact does not usually spread it.

\medterm{Lewy Bodies} Abnormal deposits of the protein alpha-synuclein inside nerve cells. They are seen in Parkinson disease and dementia with Lewy bodies and are associated with movement, thinking, sleep, and visual-perception problems.

\medterm{Lichen Striatus} A usually harmless skin condition causing small, flat, pink, brown, or pale bumps in a line along a limb. It often affects children and usually clears without treatment.

\medterm{Linear Nevus Sebaceous Syndrome} A rare disorder involving a linear sebaceous nevus, a skin lesion caused by abnormal skin glands, together with possible eye, brain, or skeletal abnormalities. It requires evaluation based on the organs involved.

\medterm{Lipitor (Atorvastatin)} Lipitor is a brand name for atorvastatin, a statin medicine that lowers low-density lipoprotein cholesterol and reduces the risk of heart attack and stroke. Muscle symptoms and liver-related effects should be discussed with a clinician.

\medterm{Lowe's Syndrome} A rare inherited disorder, also called oculocerebrorenal syndrome, that mainly affects the eyes, brain, and kidneys. Features may include congenital cataracts, developmental difficulties, low muscle tone, and kidney-related loss of minerals.

\medterm{L-theanine} An amino acid found mainly in tea leaves. It may influence relaxation and alertness, but supplements can interact with medicines and should not be assumed to treat a medical condition.

\medterm{Lung Capacity} The amount of air the lungs can hold or move, measured with breathing tests such as spirometry and lung-volume testing. Results help assess conditions including asthma, chronic obstructive pulmonary disease, and restrictive lung disease.

\medterm{Lysosomal Storage Disorders} A group of inherited diseases in which missing or faulty lysosomal enzymes cause substances to build up inside cells. Accumulation can damage organs, bones, nerves, or other tissues; treatment depends on the specific disorder.

\medterm{Lytic Lesion} A lytic lesion is an area of bone destruction caused by processes such as cancer, infection, or some metabolic bone disorders.



\medterm{Acquired Methemoglobinemia} An acquired hematologic disorder in which hemoglobin iron is oxidized from the ferrous ($\text{Fe}^{2+}$) to the ferric ($\text{Fe}^{3+}$) state by exogenous oxidizing agents (such as dapsone, benzocaine, lidocaine, nitrites, or sulfonamides), producing clinical cyanosis, chocolate-brown arterial blood, and cellular hypoxia refractory to supplemental oxygen.

\synonyms
Toxic Methemoglobinemia; Drug-Induced Methemoglobinemia; Chemically Induced Methemoglobinemia

\medterm{Benign Melanocytic Lesion} An older or imprecise term for a benign melanocytic growth, such as a melanocytic nevus or mole; unlike malignant melanoma, it is not cancerous.

\medterm{Blood Biomarker} A measurable sign of a disease or condition that can be isolated or detected in a blood sample, such as a disease-associated antibody, hormone, protein, or other substance.

\medterm{Bradykinesia and Athetosis} Hypokinetic and dyskinetic movement patterns characterized by generalized slowness of movement initiation (bradykinesia) or continuous slow, writhing, twisting involuntary movements (athetosis), frequently observed in Parkinson disease, secondary parkinsonism, or basal ganglia encephalopathies.

\synonyms
Slowness of Movement; Hypokinesia and Athetosis; Extrapyramidal Slowness

\medterm{Cessation of Menstruation} Cessation of menstrual periods, called amenorrhea; it is a normal part of menopause but may also result from pregnancy, illness, endocrine disorders, medications, stress, excessive exercise, or inadequate nutrition.

\medterm{Chorea and Myoclonus} Sudden, rapid, irregular, jerky, and non-rhythmic involuntary muscle contractions or shock-like jerks. Chorea flows unpredictably from one muscle group to another, whereas myoclonus represents brief shock-like twitches, observed in metabolic encephalopathies, post-hypoxic syndromes, and neurodegenerative disorders.

\synonyms
Jerky Involuntary Movements; Choreiform Activity; Myoclonic Jerks

\medterm{Congenital Heart Murmur} Alternate name; see \textit{Heart Murmur}.

\medterm{Continuous Passive Motion Machine} A device that gently moves a limb through a set range of motion while the person relaxes. It may be used after selected operations, especially knee surgery, but does not replace active rehabilitation.

\medterm{Cryptococcal Meningitis} A subacute, life-threatening fungal meningitis caused by \textit{Cryptococcus neoformans} or \textit{Cryptococcus gattii}, occurring predominantly in immunocompromised individuals (such as advanced HIV infection or organ transplant recipients). Diagnosis involves lumbar puncture with opening pressure measurement, CSF cryptococcal antigen (CrAg) lateral flow assay, and India ink stain.

\synonyms
Cryptococcal Meningoencephalitis; Fungal Meningitis CrAg; Cryptococcosis of CNS

\medterm{Female Urethral Meatus} The external opening of the female urethra, through which urine leaves the body; it is located anterior and superior to the vaginal opening.

\medterm{German Measles} Alternate name; see \textit{Rubella}.

\medterm{Golfer's Elbow} An overuse tendinopathy affecting the common flexor-pronator tendon origin at the medial epicondyle of the humerus. Characterized by localized tenderness and pain exacerbated by resisted wrist flexion and forearm pronation.

\synonyms
Medial Epicondylitis; Medial Epicondylalgia; Pitcher's Elbow

\medterm{Gynecomastia, Male} Alternate name; see \textit{Gynecomastia}.

\medterm{Heart Murmur Entry} Alternate name; see \textit{Heart Murmur}.

\medterm{Involuntary Movements} A diverse group of neurologic motor disorders characterized by uncontrollable, purposeless, or abnormal spontaneous body movements, including tremors, chorea, athetosis, hemiballismus, dystonia, myoclonus, and tics, stemming from basal ganglia, cerebellar, or motor circuit pathology.

\synonyms
Hyperkinetic Movement Disorders; Involuntary Motor Activity; Dyskinetic Movements

\medterm{Lambert-Eaton Myasthenic Syndrome} An autoimmune disorder in which antibodies interfere with nerve signals reaching muscles. It causes weakness, especially in the hips and shoulders, and may cause dry mouth, reduced reflexes, or problems with automatic body functions.

\synonyms
LEMS; Lambert-Eaton Syndrome; Eaton-Lambert Syndrome

\medterm{MacConkey Agar} A laboratory growth medium used mainly to grow and distinguish certain intestinal bacteria. It helps show whether bacteria can break down lactose, a sugar.

\medterm{Macerate} To soften and break down tissue after prolonged contact with fluid. The term may describe skin, wounds, or changes in a fetus that has died before delivery.

\medterm{Maceration} Softening and breakdown of tissue caused by prolonged exposure to moisture or fluid; it may occur in skin, wounds, or a fetus after death.

\medterm{Macewen Sign} An unusually hollow sound heard when part of the skull is tapped. It may occur with enlarged fluid spaces in the brain or rarely a brain abscess, but it cannot diagnose either condition by itself.

\medterm{Machupo Virus} An arenavirus that causes Bolivian hemorrhagic fever, a potentially fatal illness characterized by fever, pain, and bleeding; it is endemic in parts of Bolivia.

\medterm{Macroamylasemia} Macroamylasemia is a benign condition in which amylase binds to immunoglobulin and remains in blood, causing high serum amylase with low urinary clearance.

\medterm{Macroangiopathy} A disease of the body's larger blood vessels, often caused by fatty buildup or a blood clot. It can reduce blood flow to the heart, brain, legs, or other organs.

\medterm{Macroautophagy} A cell-cleaning process in which a cell encloses damaged parts or unwanted proteins, breaks them down, and reuses some of their materials. It helps maintain cell health during stress and normal turnover.

\medterm{Macrocephalus} Having an abnormally large head, usually defined by head circumference above the expected range for age and sex.

\medterm{Macrocephaly} An unusually large head for a person's age and sex. It may run in a family or be linked to fluid buildup, increased brain growth, or another condition.

\medterm{Macrocyte} An abnormally large red blood cell, commonly seen in macrocytosis and megaloblastic or other anemias.

\medterm{Macrocytic} Describing a cell, especially a red blood cell, that is abnormally large; macrocytic anemia features enlarged red blood cells.

\medterm{Macrocytic Anaemia} Alternate British spelling; see \textit{Macrocytic Anemia}.

\medterm{Macrocytic Anemia} A form of anemia characterized by abnormally large circulating red blood cells, defined on a complete blood count by a mean corpuscular volume (MCV) greater than $100\,\text{fL}$. It is divided into two major categories: (1) megaloblastic anemia, caused by impaired DNA synthesis from vitamin B12 or folate deficiency (often featuring hypersegmented neutrophils, glossitis, and peripheral neuropathy in B12 deficiency), and (2) non-megaloblastic macrocytosis, commonly caused by chronic alcohol consumption, liver disease, hypothyroidism, myelodysplastic syndrome, or reticulocytosis. Symptoms include fatigue, pallor, shortness of breath, and tingling in the extremities.
\synonyms{Macrocytosis; Large-Cell Anemia; High-MCV Anemia}

\medterm{Macrodactyly} Fingers or toes that are unusually large because their bones and soft tissues overgrow. It is usually present from birth, may affect one digit or several, and can interfere with movement or footwear.

\medterm{Macrodontia} Teeth that are larger than expected. It may affect one tooth or many teeth and can cause crowding, bite problems, or difficulty cleaning; treatment depends on the cause and symptoms.

\medterm{Macrogenitosomia} Unusual enlargement of the external genitals caused by excessive hormone activity before birth or during childhood. Possible causes include disorders of the adrenal glands, testes, ovaries, or other hormone systems.

\medterm{Macroglobulin} A large protein in the blood, often made of several protein units joined together. Immunoglobulin M is a major example; abnormal amounts can occur in some immune or blood disorders.

\medterm{Macroglossia} Macroglossia is an abnormally enlarged tongue that can affect feeding, speech, airway, or dental development and may be congenital or acquired.

\medterm{Macrognathia} Abnormal enlargement of the jaw, which may be associated with pituitary gigantism, tumors, or other disorders; also called a prognathic mandible.

\medterm{Macrolide} A class of antibiotics that stops bacteria from making proteins. Examples include azithromycin, clarithromycin, and erythromycin; side effects and medicine interactions vary.

\medterm{Macrolides} Antibiotics that stop bacteria from making proteins and treat selected infections. Common examples include azithromycin, clarithromycin, and erythromycin.

\medterm{Macronutrients} Nutrients needed in large amounts, mainly carbohydrates, proteins, and fats. They provide energy, build and repair tissues, support hormones and cell membranes, and help maintain normal body functions.

\medterm{Macrophage} A large immune cell that surrounds and removes germs, dead cells, and debris. It also helps coordinate inflammation and tissue repair.

\medterm{Macrophage Activation} The process by which macrophages, a type of immune cell, respond to infection or tissue injury. They may destroy germs and damaged cells, alert other immune cells, and release substances that promote inflammation or repair.

\medterm{Macrophage Activation Syndrome} A dangerous complication of some rheumatic diseases in which immune cells become overactive and cause severe inflammation, fever, low blood counts, and organ problems.

\medterm{Macrophage Count} The number or percentage of macrophages found in a tissue, body fluid, or laboratory sample. The result may help assess infection, inflammation, tissue repair, or a blood disorder, but its meaning depends on the sample and other findings.

\medterm{Macropsia} A visual disturbance in which objects appear larger than they really are. It may occur with migraine, seizures, retinal problems, medicines, or brain disorders and can temporarily affect depth and distance judgment.

\medterm{Macroscopic} Large enough to be seen without a microscope. A macroscopic examination records a specimen’s visible size, shape, color, surface, and other gross features before microscopic testing.

\medterm{Macroscopic Hematuria} Visible blood in the urine, also called gross hematuria; it may arise from the kidneys, ureters, bladder, prostate, or urethra and requires evaluation for causes such as infection, stones, trauma, or malignancy.

\medterm{Macrosomia} Macrosomia means unusually large fetal or newborn size, often associated with diabetes, prolonged pregnancy, parental size, or other factors.

\medterm{Macrostomia} An abnormally wide mouth or oral opening, usually a congenital facial cleft caused by incomplete fusion of the embryonic facial processes.

\medterm{Macula} The central part of the retina responsible for sharp, detailed, and color vision. Damage from aging, diabetes, inflammation, or other disease can cause blurred, distorted, or missing central vision.

\medterm{Macula Lutea} The small yellowish area in the center of the retina responsible for sharp vision. Its center, called the fovea, provides the clearest vision.

\medterm{Macular} Relating to a macule, a flat, circumscribed change in skin color without a change in texture or elevation.
Macules may be caused by pigment, inflammation, vascular change, or other processes and are
distinguished from raised lesions such as papules.

\medterm{Macular Cyst} A cystic defect in the macula, the central retinal area responsible for sharp vision, often associated with vitreous traction and potentially causing distortion or loss of central vision.

\medterm{Macular Degeneration} Damage to the macula, the central part of the retina responsible for sharp, straight-ahead vision. The term most often refers to age-related macular degeneration. In the dry form, the macula becomes thinner and may develop deposits called drusen; advanced dry disease can involve loss of light-sensitive retinal cells and usually progresses gradually. In the wet form, abnormal blood vessels grow beneath the retina and leak blood or fluid, often causing faster loss or distortion of central vision. Straight lines may look wavy, and blank or blurry areas may develop. Peripheral vision is usually preserved, but reading, recognizing faces, and other detailed tasks may become difficult.

\medterm{Macular Degeneration, Dry} Alternate name; see \textit{Dry Macular Degeneration}.

\medterm{Macular Degeneration, Wet} Alternate name; see \textit{Wet Macular Degeneration}.

\medterm{Macular Drusen} Small yellow deposits beneath the retina in the macula, the area responsible for sharp central vision. A few small deposits may occur with aging, while many large deposits can signal increased risk of macular degeneration.

\medterm{Macular Edema} Swelling of the macula caused by fluid leaking from retinal blood vessels. It can blur or distort central vision and may result from diabetes, vein blockage, inflammation, surgery, or other eye disease.

\medterm{Macular Hole} A small opening in the center of the retina that can cause a missing or distorted area in central vision. Some cases need surgery.

\medterm{Macular Pucker} A thin layer of scar-like tissue on the retina that can wrinkle the macula and cause distorted or blurred central vision.

\medterm{Macular Skin Lesion} Relating to a macule, a small, flat, circumscribed change in skin color that is neither raised nor depressed.

\medterm{Macule} A flat, clearly defined change in skin color that cannot be felt as a raised or depressed area. Macules may be lighter, darker, red, or brown and can have many causes.

\medterm{Macules} Small, flat, circumscribed changes in skin color that are neither raised nor depressed; they are distinguished from raised papules and fluid-filled vesicles.

\medterm{Maculopapular Lesion} A skin lesion combining flat areas of altered color (macules) with small raised solid lesions (papules); it can occur in infections, drug eruptions, inflammatory diseases, and other conditions.

\medterm{Mad Cow Disease} Alternate name; see \textit{Bovine Spongiform Encephalopathy}.

\medterm{Madarosis} Loss or absence of eyelashes, eyebrows, or other hair, caused by local disease, inflammation, trauma, or systemic conditions.

\medterm{Madura Foot} A chronic localized infection of the foot or other tissue, also called mycetoma, producing swelling, draining sinuses, and grains.

\medterm{Madura Foot} A chronic localized infection of the skin and deeper tissues of the foot, caused by fungi or certain bacteria and characterized by swelling, draining sinuses, and granules.

\synonyms
Mycetoma Pedis; Maduromycosis

\medterm{Maduramycosis} Mycetoma, a chronic subcutaneous infection caused by true fungi or filamentous bacteria and characterized by swelling, sinuses, and granules.

\medterm{Madurella} Madurella is a genus of fungi that can cause mycetoma, a chronic infection of skin and deeper tissue after implantation through the skin.

\medterm{Madurella Mycetomatis} A fungus that commonly causes eumycetoma, a chronic infection producing nodules, draining sinuses, and grains, usually after traumatic implantation in the foot or limb.

\medterm{Maffucci Syndrome} A rare condition causing multiple benign cartilage growths and abnormal blood-vessel growths, usually beginning in childhood. It can deform bones, cause fractures, and increase the risk of certain cancers.

\medterm{MAG3 Renal Scan} A kidney scan that uses a small radioactive tracer to show how well each kidney works and how urine drains. It can help detect blockage or delayed drainage.

\medterm{Magenblase Syndrome} Stomach discomfort and bloating caused by swallowing too much air, often while eating or drinking. The swallowed air may cause burping, fullness, or visible swelling and can be worsened by chewing gum, fizzy drinks, eating quickly, or anxiety.

\medterm{Magnesium} A mineral needed for healthy nerves, muscles, bones, heart rhythm, and many body processes. Too little can cause weakness, tremor, muscle cramps, or abnormal heart rhythms.

\medterm{Magnesium Blood Test} A blood test measuring serum magnesium, an electrolyte involved in neuromuscular function, cardiac rhythm, and metabolism; it helps evaluate deficiency, excess, and related electrolyte disorders.

\medterm{Magnesium Carbonate} A magnesium salt used in some antacid or laxative preparations and as a magnesium supplement; excessive intake may cause toxicity, especially in renal failure.

\medterm{Magnesium Chloride} A magnesium salt used in some oral supplements and medical preparations to provide magnesium or, in concentrated solutions, to help correct magnesium deficiency; dosing depends on the elemental magnesium content.

\medterm{Magnesium Deficiency} An abnormally low body magnesium level, often caused by poor intake, gastrointestinal
loss, renal wasting, or medications. It may cause weakness, tremor, muscle cramps,
seizures, cardiac arrhythmias, hypokalemia, and hypocalcemia.

\medterm{Magnesium Hydroxide} An osmotic laxative and antacid that neutralizes gastric acid and draws water into the bowel.

\medterm{Magnesium In Diet} Dietary magnesium comes from foods such as nuts, seeds, legumes, whole grains, and leafy vegetables and supports nerve, muscle, bone, and heart function.

\medterm{Magnesium Oxide} An inorganic magnesium compound used in some supplements and antacids; it can also have a laxative effect, and excessive use may cause hypermagnesemia, especially with kidney impairment.

\medterm{Magnesium Sulfate} A magnesium salt used intravenously or intramuscularly to treat selected magnesium deficiencies and conditions such as eclampsia, and used orally as an osmotic laxative; it is also known as Epsom salt in its hydrated form.

\medterm{Magnesium Sulfate For Neuroprotection} Magnesium given through a vein before some very early preterm births to lower the baby's risk of brain injury. The timing and dose are guided by obstetric and newborn-care specialists.

\medterm{Magnesium Sulphate} A magnesium salt used medically for severe magnesium deficiency, prevention or treatment of eclamptic seizures, and selected arrhythmias or asthma indications.

\medterm{Magnetic Resonance Angiography} MRI techniques used to depict arteries or veins, with or without contrast. It may evaluate stenosis, aneurysm, dissection, thrombosis, and vascular malformation without ionizing radiation.

\medterm{Magnetic Resonance Cholangiopancreatography} A special MRI scan that shows the bile ducts, gallbladder, and pancreatic duct without inserting a camera into them. It can detect stones, narrowing, inflammation, leaks, and some blockages without using X-rays.

\synonyms
MRCP; MR Cholangiopancreatography

\medterm{Magnetic Resonance Imaging} A scan that uses a strong magnetic field and radio waves to make detailed pictures inside the body without X-rays. It is especially useful for the brain, spinal cord, joints, muscles, and many organs; some scans require injected contrast.

\medterm{Magnetic Resonance Safety} Precautions used during an MRI scan to prevent injury from the strong magnet, radio waves, heating, loud noise, and metal objects. Staff must check implants, metal fragments, medicines, pregnancy, and other risks before scanning.

\medterm{Magnetic Resonance Spectroscopy} A special MRI test that measures chemicals in tissue rather than only making a picture of its shape. It may help study brain disorders or provide extra information about a tumor, but results must be interpreted with ordinary imaging and other tests.

\medterm{Magnetic Resonance Venography} An MRI test that produces images of veins, with or without injected contrast. It can assess suspected clots or narrowing in veins of the brain, abdomen, pelvis, or limbs and does not use X-rays.

\synonyms
MRV; MR Venography

\medterm{Magnetic Seed} A tiny magnetic marker placed in a breast lesion or lymph node to help a surgeon find it. A special detector locates the marker during surgery, providing an alternative to some wire or radioactive-marker methods.

\medterm{Magnetic Therapy} Treatment that uses a magnetic field for a health purpose. Ordinary static magnets have not been shown to treat disease, while specialized magnetic stimulation devices have established uses for some neurologic or psychiatric conditions.

\medterm{Magnetocardiography} A test that records the very weak magnetic fields produced by the heart’s electrical activity using sensors outside the body. It may help study heart rhythm or electrical conduction, usually in specialist or research settings.

\medterm{Magnetoencephalography} A noninvasive technique that records the weak magnetic fields produced by neuronal electrical activity, using sensors outside the head to localize and study brain function and some epileptic discharges.

\medterm{Magnusiomyces Capitatus} An opportunistic yeast formerly known as \textit{Geotrichum capitatum} or \textit{Saprochaete capitata}; it can cause invasive infections, especially in people with severe immunosuppression or hematologic malignancy.

\medterm{Mahvash Disease} A rare inherited disorder in which the body’s cells do not respond normally to glucagon, a hormone that helps raise blood sugar. It can enlarge the pancreas, increase glucagon levels, and disturb blood-sugar control.

\medterm{Main-En-Griffe} A claw-like hand position usually caused by damage to the ulnar nerve, which controls several hand muscles. The ring and little fingers may bend at the smaller joints, weaken the grip, and become difficult to straighten.

\medterm{Maintenance} Ongoing care or treatment intended to preserve health, prevent relapse or deterioration, or sustain a therapeutic result; maintenance therapy is continued after initial control of a condition.

\medterm{Maintenance Dose} The regular amount of a medicine needed to keep its effect at the desired level. The dose depends on how quickly the body removes the medicine and how much the body absorbs.

\medterm{Maintenance Of Wakefulness Test} A sleep test that measures how long a person can remain awake in a quiet setting during the day.

\medterm{Maintenance Therapy} Treatment continued after a disease has improved or been controlled to maintain that benefit, prevent return or worsening, or delay progression. The treatment may be a medicine, procedure, or other care, and its duration depends on the condition, response, risks, and treatment goals.

\medterm{Maize} Corn, a cereal grain used as food and feed; contaminated maize can expose people to fungal toxins such as aflatoxin.

\medterm{Majeed Syndrome} A rare autoinflammatory disorder caused by LPIN2 variants, characterized by chronic recurrent multifocal osteomyelitis, congenital dyserythropoietic anemia, and inflammatory skin disease.

\medterm{Majewski Syndrome} A rare inherited skeletal disorder involving abnormal bone growth and distinctive craniofacial or limb findings; the exact features depend on the genetic subtype.

\medterm{Major} Larger, greater, or more significant than another related structure or condition; for example, the teres major muscle is larger than the teres minor muscle.

\medterm{Major Depression} Alternate name; see \textit{Major Depressive Disorder}.

\medterm{Major Depression With Psychotic Features} A major depressive episode accompanied by delusions or hallucinations, usually mood-congruent, requiring prompt psychiatric assessment and treatment.

\medterm{Major Depressive Disorder} A mood disorder involving one or more major depressive episodes. During an episode, five or more symptoms occur during the same two-week period, including either depressed mood or loss of interest or pleasure (anhedonia). Other symptoms may include changes in sleep, appetite, energy, concentration, movement, guilt, worthlessness, or thoughts of death. The symptoms cause clinically significant distress or impairment.

\synonyms
Major Depression; Clinical Depression

\medterm{Major Histocompatibility Complex} A group of genes and cell-surface proteins that display pieces of germs or abnormal cells to the immune system. In humans, these proteins are called HLA.

\medterm{Major Histocompatibility Complex Molecule} A cell-surface protein that displays small pieces of germs or abnormal cells to immune cells. This helps the immune system recognize what belongs to the body and what may need an immune response.

\synonyms
MHC Molecule

\medterm{Major Trauma} A serious injury that may threaten life, limb, or organ function. It can involve multiple body regions or major damage to the brain, spine, chest, abdomen, pelvis, or blood vessels.

\medterm{Malabsorption} A condition in which the body cannot properly digest or absorb nutrients from food. It may cause chronic diarrhea, greasy stools, weight loss, anemia, or vitamin and mineral deficiencies.

\medterm{Malabsorption Syndrome} Impaired absorption of nutrients from the small intestine, causing diarrhea, weight loss, anemia, vitamin deficiencies, or other nutritional complications.

\medterm{Malacia} Abnormal softening or weakening of tissue. It may affect cartilage, bone, or brain tissue after faulty development, injury, infection, reduced blood supply, or another disease; symptoms depend on the tissue involved.

\medterm{Malacoplakia} A rare long-term inflammation in which immune cells collect mineral deposits. It most often affects the bladder or digestive tract and can form a mass that looks like cancer.

\medterm{Maladaptation} An adjustment to a situation or environment that does not work well and harms health or daily functioning. The term may describe physical, psychological, or behavioral responses to illness or stress.

\medterm{Maladaptive Stress Response} A lasting or overly strong reaction to stress that interferes with health, daily life, or relationships. It may include constant worry, poor sleep, irritability, avoidance, or physical symptoms and can improve with support and appropriate treatment.

\medterm{Malaise} Malaise is a general feeling of illness, discomfort, or lack of well-being that commonly accompanies infection, inflammation, chronic disease, medication effects, or other systemic conditions.

\medterm{Malakoplakia} A rare long-term inflammation in which immune cells collect mineral deposits. It most often affects the bladder or digestive tract and can form a mass that resembles cancer.

\medterm{Malar} Relating to the cheekbone or cheek area. A malar rash is a red or purple rash across the cheeks and bridge of the nose, often seen in lupus. A malar fracture is a break in the cheekbone and may cause swelling, numbness, or a changed facial shape.

\medterm{Malar Rash} A red or violaceous rash across the cheeks and bridge of the nose, often sparing the folds beside the nose.
It is classically associated with systemic lupus erythematosus but can occur with other conditions;
appearance alone does not establish a diagnosis.

\medterm{Malaria} A life-threatening blood infection caused by single-celled \textit{Plasmodium} parasites and spread by the bite of infected female \textit{Anopheles} mosquitoes. Once in the body, parasites multiply in the liver and then invade red blood cells, causing recurring attacks that follow three stages: shaking chills (cold stage), high fever (hot stage), and heavy sweating (sweating stage) every 1 to 3 days. \textit{Plasmodium falciparum} is the most dangerous species because it clogs small blood vessels, leading to cerebral malaria (seizures, coma), severe anemia, dark urine (blackwater fever), lung fluid, or kidney failure. \textit{Plasmodium vivax} can stay dormant in the liver and cause relapses months later unless treated with primaquine. Diagnosis is made by examining a blood smear under a microscope or with a rapid test. Treatment uses artemisinin-based combination medications (ACTs) or intravenous artesunate for severe disease.

\synonyms
Plasmodium Infection; Paludism; Ague

\medterm{Malassezia} Malassezia is a yeast normally found on skin that can contribute to dandruff, seborrheic dermatitis, pityriasis versicolor, and some folliculitis.

\medterm{Malassezia Furfur} Malassezia furfur is a yeast that normally lives on skin but can multiply excessively. It is linked with dandruff, oily skin inflammation, and pityriasis versicolor, which causes lighter or darker patches.

\medterm{Malassimilation} Failure to properly digest food or absorb its nutrients. It may cause diarrhea, weight loss, bloating, anemia, or vitamin shortages and can result from pancreatic disease, bowel disorders, surgery, infection, or food intolerance.

\medterm{Malate Dehydrogenase} An enzyme in energy-producing pathways that changes malate into oxaloacetate. It helps cells process nutrients and is involved in transferring chemical energy within mitochondria and other cellular compartments.

\medterm{Malathion} An organophosphate insecticide that inhibits acetylcholinesterase; topical formulations have also been used to treat head lice, but exposure or ingestion can cause cholinergic toxicity.

\medterm{Malathion Poisoning} Poisoning caused by the organophosphate insecticide malathion. It can produce sweating, salivation, vomiting, diarrhea, wheezing, slow heartbeat, muscle twitching, weakness, seizures, or breathing failure and needs urgent assessment.

\medterm{MALDI-TOF Mass Spectrometry} A laboratory method that identifies bacteria or fungi by comparing their protein pattern with a reference database. It can identify an organism quickly, but separate testing is needed to determine antibiotic sensitivity.

\medterm{Male} Describing a sex associated with reproductive anatomy that produces or is organized around small gametes (sperm); chromosomes, hormones, anatomy, and gender identity do not always fit a simple binary pattern.

\medterm{Male Breast Cancer} Cancer developing in breast tissue in a man. It is uncommon, but a new lump, nipple discharge, skin change, or nipple retraction needs medical evaluation.

\synonyms
Male Mammary Carcinoma; Breast Carcinoma in Men

\medterm{Male Breast Pathology} Disorders affecting male breast tissue, including gynecomastia, benign epithelial lesions, inflammatory conditions, and carcinoma. Evaluation of a persistent mass or nipple change requires clinical examination and appropriate imaging, with biopsy when indicated.

\medterm{Male Condom} A thin sheath placed over the penis before sexual contact to reduce the risk of pregnancy and transmission of many sexually transmitted infections by acting as a physical barrier.

\medterm{Male Genitalia} The external and internal reproductive organs of a male, including the penis, scrotum, testes, epididymides, vas deferens, seminal vesicles, prostate, and associated ducts and glands.

\medterm{Male Hypoactive Sexual Desire Disorder} Persistently reduced sexual thoughts or desire in a man that causes personal distress or relationship difficulty. It differs from erection problems, although both can occur together; medicines, illness, mood, and relationship factors may contribute.

\medterm{Male Hypogonadism} Reduced production of testosterone, impaired sperm production, or both due to disease of the testes or
the hypothalamic-pituitary system. Possible features include low libido, erectile difficulty,
infertility, reduced muscle mass, fatigue, and loss of bone density.

\synonyms
Low Testosterone; Testosterone Deficiency

\medterm{Male Infertility} Reduced ability to achieve a pregnancy because of factors affecting sperm production, sperm
transport, ejaculation, sexual function, or other reproductive processes.

\synonyms
Male Subfertility; Male Factor Infertility

\medterm{Male Menopause} Alternate informal term; see \textit{Male Hypogonadism}.

\medterm{Male Pattern Baldness} Androgenetic alopecia, a patterned progressive loss of scalp hair caused by genetic susceptibility and androgen effects on hair follicles.

\medterm{Male Pelvis} The lower part of the abdomen between the hip bones in a male; compared with the female pelvis, it is generally narrower, taller, and more robust, with a narrower pubic arch and sacrum.

\medterm{Male Phenotype} The observable reproductive anatomy, secondary sex characteristics, and other traits associated with male sex development, produced by interacting genetic, hormonal, and developmental factors and showing natural variation.

\medterm{Male Reproductive System} The organs and tissues that produce sperm and male sex hormones and deliver sperm, including the testes, ducts, accessory glands, penis, and scrotum.

\medterm{Malformation} An abnormal structural development present from birth, caused by disturbed formation of an organ or body part.

\medterm{Malignancy} Cancerous disease in which abnormal cells grow uncontrollably and may invade nearby tissue or spread to distant organs. Its likely behavior and treatment depend on the cancer type, extent, grade, and important molecular findings.

\medterm{Malignancy-Associated Hypercoagulability} A cancer-related tendency toward excessive clot formation caused by tumor-associated coagulation activation, inflammation, platelet interaction, treatment, immobility, or vascular injury. See \textit{Cancer-Associated Thrombosis}.

\medterm{Malignant} Cancerous, with the ability to invade nearby tissues or spread to distant parts of the body. Malignant cells grow abnormally and may damage organs, but behavior varies among cancer types and stages.

\medterm{Malignant Acanthosis Nigricans} A paraneoplastic syndrome in which velvety, hyperpigmented, verrucous plaques develop
rapidly in flexural and mucosal areas as a marker of an underlying internal malignancy,
most commonly gastric adenocarcinoma. Unlike benign acanthosis nigricans, the malignant
form is extensive, abrupt in onset, and may involve the oral mucosa.

\medterm{Malignant Ameloblastoma} A rare cancer arising from ameloblastoma-forming tissue, usually in the jaw, with potential to invade locally or spread to distant organs.

\medterm{Malignant Ascites} Buildup of fluid in the abdomen caused by cancer involving the abdominal lining or nearby organs. It can cause swelling, discomfort, early fullness, nausea, or breathlessness and may require drainage and cancer-directed treatment.

\medterm{Malignant Bone Neoplasm} A cancer that begins in bone or nearby supporting tissue. It may cause persistent bone pain, swelling, or a fracture; scans and a biopsy identify the type and extent before treatment.

\medterm{Malignant Fibrous Histiocytoma} Alternate name; see \textit{Undifferentiated Pleomorphic Sarcoma}.

\medterm{Malignant Giant Cell Tumor Of Tendon Sheath} A rare locally aggressive malignant tumor resembling a giant cell tumor of tendon sheath, characterized by destructive growth and potential for metastasis; diagnosis requires expert soft-tissue pathology review.

\medterm{Malignant Glioma} An aggressive primary tumor arising from glial or glial-like cells of the central nervous system; high-grade examples include anaplastic astrocytoma and glioblastoma.

\medterm{Malignant Hypertension} A medical emergency in which very high blood pressure is causing new injury to organs such as the brain, heart, kidneys, or eyes. It can cause headache, confusion, chest pain, breathlessness, vision changes, seizures, or kidney failure.

\medterm{Malignant Hyperthermia} A rare, life-threatening reaction to certain anesthetic medicines that causes rapidly rising body temperature, severe muscle stiffness, fast heart rate, and dangerous chemical changes. It requires immediate emergency treatment and affects some genetically susceptible people.

\medterm{Malignant Mastocytosis} An older name for a severe disorder in which abnormal mast cells multiply and collect in organs. It can damage the bone marrow, liver, spleen, or other tissues and may cause flushing, faintness, bleeding, or organ failure.

\medterm{Malignant Melanoma} A cancer that begins in melanocytes, the cells that make skin pigment. It may appear as a changing dark or unusual mole and can spread to other tissues, especially when diagnosed at a later stage.

\medterm{Malignant Melanoma} Alternate name; see \textit{Melanoma}.

\medterm{Malignant Mesothelioma} A cancer of the thin lining around the lungs, abdomen, or heart. It is strongly linked to asbestos exposure and may cause chest or abdominal pain, fluid buildup, cough, or breathlessness.

\medterm{Malignant Myoepithelioma} A rare cancer arising from supportive contractile cells in a gland, most often a salivary gland. It can grow into nearby tissue and spread elsewhere, causing a slowly enlarging or painful mass.

\medterm{Malignant Neoplasm} A cancerous growth in which abnormal cells multiply without control, invade nearby tissues, and may spread to distant organs through the blood or lymphatic system.

\synonyms
Cancer; Malignancy

\medterm{Malignant Otitis Externa} An invasive infection of the external ear canal and skull base, usually caused by Pseudomonas aeruginosa in older adults with diabetes or immunocompromise; it can cause severe pain and cranial-nerve complications.

\medterm{Malignant Peripheral Nerve Sheath Tumor} An uncommon soft-tissue cancer arising from or near peripheral nerves. It may cause a growing painful mass, numbness, or weakness and is treated with specialist surgery, radiation, or medicines when appropriate.

\medterm{Malignant Pleural Effusion} Buildup of fluid around the lung because cancer involves the lung lining or has spread there. It may cause breathlessness, chest discomfort, or cough; drainage relieves symptoms, while cancer treatment may reduce recurrence.

\medterm{Malignant Pustule} An older name for the skin form of anthrax, caused by Bacillus anthracis and beginning as an itchy papule that develops into a painless black eschar.

\medterm{Malignant Syphilis} A severe form of secondary syphilis that causes fever and widespread crusted, ulcerated, or tissue-damaging skin lesions. It is more common with weakened immunity and requires prompt specialist treatment for syphilis.

\medterm{Malignant Teratoma Of The Mediastinum} A malignant germ-cell tumor containing tissues from more than one embryonic layer and arising in the mediastinum; evaluation uses imaging, tumor markers, and tissue diagnosis.

\medterm{Malignant Thymoma} An older term for an invasive thymic epithelial tumor; modern usage distinguishes invasive thymoma from thymic carcinoma, which is a separate, more aggressive thymic malignancy.

\medterm{Malignant Tumor} A cancerous growth capable of invading nearby tissue or spreading to distant sites. Its behavior and treatment depend on its type, grade, stage, and molecular features.

\medterm{Malingering} Deliberately pretending to have, or exaggerating, symptoms to obtain an outside benefit such as money, medicines, shelter, or avoiding an obligation. It is not itself a mental illness and should not be diagnosed without reliable evidence.

\medterm{Mallampati Classification} A grading system used before anesthesia to estimate how difficult it may be to place a breathing tube. The clinician looks inside the open mouth and records which parts of the throat can be seen; it is only one part of the airway assessment.

\medterm{Malleolus} A bony prominence at the ankle formed by the lower end of the tibia medially or fibula laterally.

\medterm{Mallet} A surgical hammer used with tools such as osteotomes to cut or shape bone. It delivers controlled force during operations, helping surgeons work on bone while limiting unnecessary injury to nearby tissues.

\medterm{Mallet Finger} An injury to the extensor tendon at the distal finger joint, causing inability to actively straighten the fingertip. It may result from a forced flexion injury and sometimes includes an avulsion fracture.

\medterm{Mallet Toe} A flexion deformity of the distal interphalangeal joint, causing the end of a lesser toe to bend downward. It may be flexible or rigid and can cause pain, rubbing, corns, or difficulty wearing shoes; it differs anatomically from hammertoe and claw toe.

\medterm{Mallet Toe and Hammertoe} Alternate name; see \textit{Hammertoe}.

\medterm{Malleus} The largest of the three small bones in the middle ear, also called the hammer. It is attached to the eardrum and passes sound vibrations to the incus and then the inner ear.
\commonname{Hammer}

\medterm{Mallory Body} A clump of damaged proteins inside a liver cell, seen most often in alcohol-related liver disease but also in other liver conditions. It is a laboratory finding, not a diagnosis by itself.

\medterm{Mallory-Weiss Tear} A superficial tear in the lining of the lower esophagus or upper stomach, usually caused by forceful vomiting or retching. It commonly causes vomiting of blood. Many tears stop bleeding on their own within a few days, but upper endoscopy may be needed for diagnosis and treatment.

\medterm{Malnutrition} Poor health caused by too little, too much, or an imbalance of nutrients. It can impair growth, immunity, wound healing, muscle function, and organ function and may involve undernutrition, nutrient deficiencies, or excess weight.

\medterm{Malocclusion} A mismatch in the position of the upper and lower teeth or jaws that affects how the teeth meet. It may cause crowding, an overbite, an underbite, difficulty chewing, speech problems, jaw discomfort, or trouble cleaning the teeth.

\medterm{Malocclusion Of Teeth} Misalignment of the upper and lower teeth or jaws that affects the bite, chewing, speech, appearance, or oral hygiene.

\medterm{Malondialdehyde} A chemical formed when fats are damaged by oxidation. Laboratory researchers may measure it as an indirect sign of oxidative stress, but the result alone does not diagnose a disease.

\medterm{Malpighian Bodies} Renal glomeruli within Bowman's capsules, where blood filtration begins; the term may also refer historically to lymphoid nodules in the spleen.

\medterm{Malposition} An abnormal position of an organ, body part, implanted device, or fetus; in obstetrics, fetal malposition refers to an abnormal orientation of the presenting part during labor.

\medterm{Malpresentation} A fetal presentation other than the usual vertex presentation, including breech, face, brow,
shoulder, or compound presentation. It may complicate labor and influence delivery planning.

\medterm{Malrotation} Abnormal rotation or fixation of an organ during development, most often referring to abnormal rotation of the
intestine and its risk of volvulus.

\medterm{MALS} Abbreviation; see \textit{Median Arcuate Ligament Syndrome}.

\medterm{MALT Lymphoma} Alternate name; see \textit{MALT Lymphoma}.

\medterm{Malta Fever} Brucellosis, a zoonotic infection caused by Brucella species and characterized by intermittent fever, sweats, fatigue, and joint or systemic symptoms.

\medterm{Maltase} An intestinal enzyme that breaks maltose into two glucose molecules.

\medterm{Maltodextrin} A mixture of short glucose polymers produced by partial hydrolysis of starch, used as a rapidly digestible carbohydrate and as a thickener, filler, or carrier in foods and medicines.

\medterm{Maltose} A reducing disaccharide composed of two glucose molecules joined by an alpha-1,4 glycosidic bond, produced during starch digestion and hydrolyzed by maltase.

\medterm{Malt-Worker's Lung} An occupational hypersensitivity pneumonitis caused by inhalation of organic dusts associated with malting or grain processing.

\medterm{Mammaglobin A} A secretory protein expressed predominantly in mammary tissue and used as an immunohistochemical marker that can help identify breast origin in some metastatic carcinomas.

\medterm{Mammaplasty} Surgery to change the size or shape of the breast or to rebuild it after breast removal. It may enlarge, reduce, lift, or reconstruct the breast; possible risks include bleeding, infection, altered sensation, scarring, and problems with healing.

\medterm{Mammary} Relating to the breast or milk-producing gland. Mammary glands develop during puberty and pregnancy and produce milk after childbirth under hormone control. The term may describe normal breast tissue or diseases such as cysts, infection, benign growths, or breast cancer.

\medterm{Mammary Duct Ectasia} A benign inflammatory breast condition in which major subareolar milk ducts become abnormally dilated, shortened, and filled with stagnant thick lipid-rich secretions, leading to periductal chronic inflammation and fibrosis. Most common in perimenopausal and postmenopausal women, it typically presents with thick, sticky, multicolored (green, brown, or creamy) nipple discharge, nipple tenderness, or a slit-like retraction of the nipple. It can produce a firm subareolar lump that closely mimics breast cancer on mammography and physical examination. Management is typically conservative with warm compresses and reassurance, though surgical excision of affected ducts (Hadfield's procedure) is performed for persistent symptoms.

\synonyms
Duct Ectasia; Plasma Cell Mastitis; Periductal Mastitis; Comedomastitis

\medterm{Mammary Epithelium} The epithelial lining of the mammary ducts and secretory alveoli; these cells help produce, secrete, and transport milk and can give rise to many breast carcinomas.

\medterm{Mammary Gland} The milk-producing gland in the breast. It contains small milk-making lobules connected to ducts that carry milk to the nipple, and its growth and activity are controlled by hormones.

\medterm{Mammary Glands} Paired milk-producing glands in the breasts, composed of lobules and ducts that develop and function under hormonal control.

\medterm{Mammary Glands} Paired glands of the anterior chest that produce milk after childbirth. Each breast contains lobules, ducts, connective tissue, and fat and is influenced by hormonal and reproductive changes.

\synonyms
Breasts; Lactiferous Glands

\medterm{Mammary Implant} An artificial device placed in the breast for reconstruction or enlargement. Implants may contain saline or silicone gel and may be placed above or below the chest muscle. Possible complications include infection, rupture, hard scar tissue, changes in sensation, and rarely implant-associated lymphoma.

\medterm{Mammary Paget Disease} An intraepidermal spread of malignant glandular cells into the nipple epidermis, usually associated with underlying breast carcinoma and presenting as persistent eczematous or crusted nipple change.

\medterm{Mammillary Body} One of a pair of rounded nuclei on the underside of the hypothalamus that participate in memory circuits connecting the hippocampus with the thalamus; injury can contribute to memory impairment.

\medterm{Mammogram} A mammogram is a low-dose x-ray examination of breast tissue used for screening and evaluation of masses, asymmetry, calcifications, or other abnormalities.

\medterm{Mammographic Calcifications} Mineral deposits within breast tissue visualized on screening or diagnostic mammography. Macrocalcifications are typically benign (associated with fibroadenomas, cysts, or fat necrosis), whereas grouped, pleomorphic, or linear branching microcalcifications warrant biopsy to exclude ductal carcinoma in situ (DCIS) or invasive carcinoma.

\synonyms
Breast Microcalcifications; Mammographic Microcalcifications; Breast Calcifications

\medterm{Mammography} X-ray imaging of the breast used to screen for or evaluate breast abnormalities. Screening intervals
depend on age, risk, symptoms, and the guideline being followed. Standard views include
the craniocaudal and mediolateral-oblique projections.

\medterm{Mammoplasty} Plastic or reconstructive breast surgery that changes breast size, shape, or contour. It may reduce or enlarge the breasts, lift them, rebuild them after cancer surgery, or correct differences present from birth.

\medterm{Managing Pain During Labor} Pain relief during labor may include breathing and movement, support, injected or intravenous medicines, nitrous oxide, or neuraxial anesthesia, chosen according to preference and clinical circumstances.

\medterm{Mandatory Reporting} A legal or regulatory duty to report specified findings or circumstances—such as suspected child or elder abuse, certain communicable diseases, or some injuries—to designated authorities; requirements vary by jurisdiction.

\medterm{Mandible} The large movable bone forming the lower jaw and holding the lower teeth. It supports chewing and speech and contains nerves and blood vessels that supply the teeth, gums, and lower lip.

\medterm{Mandible} The large movable bone forming the lower face and holding the lower teeth. It supports chewing and speech and contains nerves and blood vessels that supply the teeth, gums, and lip.

\synonyms
Inferior Maxilla; Jawbone

\medterm{Mandibular} Relating to the mandible, the lower jawbone. It holds the lower teeth and moves at the jaw joints. A mandibular fracture can cause pain, swelling, numbness, loose teeth, or a changed bite.

\medterm{Mandibular Advancement Device} An oral appliance that holds the lower jaw or tongue forward during sleep to enlarge or stabilize the upper airway, usually for selected patients with obstructive sleep apnea or snoring.

\medterm{Mandibular Condyle} The rounded superior end of each mandibular ramus that articulates with the temporal bone at the temporomandibular joint and permits jaw movement.

\medterm{Mandibular Fracture} A break in the lower jaw bone typically caused by trauma, presenting with pain,
swelling, malocclusion, and difficulty opening the mouth; may require surgical fixation.

\medterm{Mandibular Prosthesis} An artificial implant or device used to replace, reconstruct, stabilize, or restore part of the mandible after trauma, tumor removal, congenital deformity, or other loss of jaw structure.

\medterm{Mandibulofacial Dysostosis} A congenital craniofacial disorder, classically Treacher Collins syndrome, involving underdevelopment of the mandible and cheekbones with variable abnormalities of the ears, eyes, and palate.

\medterm{Manganese} An essential trace element required in small amounts. Excessive exposure, especially by inhalation, can injure the nervous system and cause manganism, a movement disorder resembling parkinsonism.

\medterm{Mania} A distinct period of abnormally elevated, expansive, or irritable mood and increased energy or activity, with associated symptoms such as reduced need for sleep, rapid speech, racing thoughts, inflated self-esteem, or risky behavior. It causes marked impairment, requires hospitalization, or includes psychotic features; the episode usually lasts at least one week when hospitalization is not required.

\medterm{Manic} Relating to mania, a pathological episode of abnormally elevated, expansive, or irritable mood with increased energy or activity and associated changes in thought, speech, sleep, or judgment.

\medterm{Manic Disorder} Alternate name; see \textit{Bipolar Disorder}.

\medterm{Manic Episode} A distinct period of abnormally elevated, expansive, or irritable mood and increased energy or activity, lasting at least seven days unless hospitalization is required sooner or psychotic features occur. It includes associated changes such as reduced need for sleep, rapid speech, racing thoughts, inflated self-esteem, or risky behavior and causes marked impairment.

\medterm{Manic Mood} An abnormally elevated, expansive, or persistently irritable mood accompanied by increased energy or activity, reduced need for sleep, rapid speech, racing thoughts, inflated self-esteem, or risky behavior; it is characteristic of a manic episode.

\medterm{Manic-Depression} An older term for bipolar disorder, a mood disorder characterized by episodes of mania or hypomania and episodes of depression.

\medterm{Manic-Depressive Illness} Alternate name; see \textit{Bipolar Disorder}.

\medterm{Manifest} Clearly apparent or expressed; in clinical usage, a disease may manifest through particular signs or symptoms.

\medterm{Manifestation} A sign, symptom, or other observable expression of a disease or health condition. Manifestations may be noticed by the person, found during examination, or detected through laboratory tests or imaging.

\medterm{Manipulation} Deliberate physical movement of a body part or joint, often performed therapeutically or diagnostically by a trained clinician.

\medterm{Manipulation Procedure} Deliberate movement of a body part by a trained clinician to restore motion, alignment, or function. It may be performed with or without anesthesia; soreness, bleeding, nerve injury, fracture, or worsening symptoms are possible depending on the procedure.

\medterm{Mannan} A polysaccharide composed primarily of mannose units, found in the cell walls or capsules of many fungi and in some plants; fungal mannans can be recognized by the immune system and used in laboratory testing.

\medterm{Mannheimia} A genus of gram-negative bacteria that mainly infects livestock, especially the respiratory tract of cattle and sheep.

\medterm{Mannitol} A medicine that pulls water into the urine. It is given by injection in selected emergencies to reduce swelling or pressure around the brain or inside the eye and in some kidney-related situations; it can worsen dehydration or heart failure.

\medterm{Mannose} A six-carbon monosaccharide related to glucose that forms part of glycoproteins and glycolipids and is important in protein glycosylation and cell recognition.

\medterm{Mannose-Binding Lectin Deficiency} Low levels or reduced activity of a blood protein that helps the immune system recognize and attack germs. Many people have no symptoms, but some are more prone to repeated infections, especially if other immune defenses are weak.

\medterm{Mannosidosis} An inherited lysosomal storage disorder caused by deficient alpha-mannosidase activity, leading to accumulation of mannose-rich oligosaccharides and variable intellectual, skeletal, immune, hearing, and neurologic abnormalities.

\medterm{Manometer} An instrument that measures pressure in a liquid or gas. In medicine, it may measure pressure in blood vessels, the lungs, the bladder, or another body space.

\medterm{Manometry} A test that measures pressure and muscle coordination inside an organ or passage, such as the esophagus, stomach, rectum, or urinary tract. It helps identify movement and sphincter disorders.

\medterm{Mansonella Perstans} A filarial worm that causes mansonellosis, usually acquired through bites of infected midges and often producing mild or no symptoms.

\medterm{Mansonella Streptocerca} A filarial worm transmitted by biting midges that lives mainly in the skin and can cause itchy skin disease or small nodules.

\medterm{Mansonelliasis} Infection with filarial nematodes of the genus \textit{Mansonella}, transmitted by biting midges or blackflies; many infections are asymptomatic, but some cause itching, skin lesions, fever, or eosinophilia.

\medterm{Mantle Cell Lymphoma} A usually fast-growing cancer of B lymphocytes that often enlarges lymph nodes and the spleen and may involve the bone marrow, blood, or digestive tract. A characteristic chromosome change helps confirm the diagnosis and guide treatment.

\medterm{Mantle Field} A historical radiation-treatment field encompassing lymph-node regions in the neck, mediastinum, and upper chest, formerly used for Hodgkin lymphoma; modern treatment generally uses more limited involved-site fields.

\medterm{Mantoux Test} A skin test for immune sensitization to tuberculosis. A small amount of purified protein derivative is injected into the forearm, and the raised area is measured 48 to 72 hours later; a positive result does not by itself prove active disease.

\medterm{Manual} Performed by hand rather than automatically, or a written guide explaining procedures and instructions. In healthcare, manual techniques may assess, move, treat, or support body structures safely.

\medterm{Manual Removal Of Placenta} A procedure in which a clinician removes a placenta that has not separated or cannot be delivered after birth. The timing and need depend on bleeding, maternal condition, and local clinical protocols.

\medterm{Manual Vacuum Aspiration} A procedure that uses a handheld syringe and a thin tube to remove tissue from the uterus. It may treat early pregnancy loss or end a pregnancy, depending on local law and clinical circumstances, and is performed by trained clinicians.

\synonyms
MVA; Mini-Suction Curettage

\medterm{Manus} The hand, including the wrist and fingers, in anatomical terminology.

\medterm{MAOIs} A class of medicines used mainly for depression and some other conditions. They slow the breakdown of certain brain chemical messengers, but have important interactions with foods and medicines and must be prescribed and monitored carefully.

\medterm{Map-Dot-Fingerprint Type Corneal Dystrophy} A usually bilateral corneal epithelial basement membrane dystrophy producing map-like, dot-like, or fingerprint-like patterns and sometimes recurrent corneal erosions or visual blurring.

\medterm{Maple Syrup Urine Disease} An inherited disorder in which the body cannot properly break down three amino acids from protein. Harmful substances build up, causing sweet-smelling urine, poor feeding, sleepiness, abnormal movements, seizures, or brain injury without treatment.

\medterm{Mapping} Charting the location of genes or other identifiable regions on chromosomes; the term may also describe charting functional or anatomical regions in other medical contexts.

\medterm{Marantic} Relating to a sterile, noninfective process associated with wasting; marantic endocarditis refers to nonbacterial thrombotic endocarditis.

\medterm{Marasmus} Severe undernutrition caused by prolonged lack of calories and protein. It leads to marked loss of body fat and muscle, poor growth in children, weakness, and increased vulnerability to infection; urgent nutritional and medical care is needed.

\medterm{Maraviroc} A CCR5 antagonist (entry inhibitor) used in combination antiretroviral therapy for HIV-1
infection in patients with CCR5-tropic virus. Maraviroc binds to the human CCR5
coreceptor on CD4+ T cells, preventing HIV-1 gp120 binding and viral entry.

\medterm{Marble Bone Disease} Osteopetrosis, an inherited disorder of defective bone resorption causing abnormally dense but fragile bones and possible marrow failure.

\medterm{Marburg Virus} A filovirus that causes Marburg virus disease, a severe viral hemorrhagic fever in humans and some nonhuman primates. The virus is associated with Egyptian rousette fruit bats and can first spread to people after contact with infected bats. It then spreads between people through direct contact with infected blood, other body fluids, or materials contaminated with them, not through ordinary airborne spread. The disease may begin with high fever, severe headache, tiredness, and muscle aches, followed by vomiting, diarrhea, rash, bleeding, shock, or organ failure.

\medterm{Marburgvirus} A group of filoviruses that includes Marburg virus and Ravn virus, the viruses that cause Marburg virus disease. The disease is a severe viral hemorrhagic fever that can spread from fruit bats to people and between people through direct contact with infected body fluids.

\medterm{Marden-Walker Syndrome} A rare inherited disorder characterized by a distinctive facial appearance, low muscle tone, limited joint movement, poor growth, and developmental delay. Care is supportive and involves several specialties.

\medterm{Marek's Disease} A contagious herpesvirus disease of chickens that can cause nerve paralysis, tumors, weight loss, or death.

\medterm{Marfan Syndrome} An inherited connective-tissue disorder usually caused by a disease-causing variant in the \textit{FBN1} gene, which encodes fibrillin-1. It may affect the aortic root and heart valves, eyes, skeleton, joints, lungs, and dura. Features can include aortic-root enlargement, lens dislocation, tall or long-limbed stature, long fingers, scoliosis, chest-wall differences, and joint laxity; the severity varies among individuals.

\medterm{Marfanoid Habitus} A body habitus characterized by tall stature, long limbs, long fingers (arachnodactyly),
and joint laxity. Seen in Marfan syndrome, but also in other connective tissue disorders
and homocystinuria. Differentiated from Marfan syndrome by absence of aortic root
dilation and lens subluxation.

\medterm{Marfan's Syndrome} Alternate spelling; see \textit{Marfan Syndrome}.

\medterm{Marfan's Syndrome} Alternate spelling; see \textit{Marfan Syndrome}.

\medterm{Marginal Ulcer} An ulcer occurring at or near a gastrojejunal anastomosis, often after gastric bypass or other gastric surgery, where acid, ischemia, smoking, nonsteroidal anti-inflammatory drugs, or other factors may contribute.

\medterm{Marginal Zone Lymphoma} A slow-growing non-Hodgkin lymphoma that begins in immune cells in the marginal zone of lymphoid tissue.

\synonyms
MZL

\medterm{Maribavir} Maribavir is an antiviral medicine used to treat certain cytomegalovirus infections that are resistant or refractory to other treatment after transplantation.

\medterm{Marie-Strumpell Disease} Ankylosing spondylitis, a chronic inflammatory disease primarily affecting the sacroiliac joints and spine that can cause back pain, stiffness, and progressive spinal fusion.

\medterm{Marijuana} Cannabis preparations containing cannabinoids such as THC; medical effects and risks depend on dose, formulation, route, and patient factors.

\synonyms
Cannabis; Weed; Pot

\medterm{Marijuana Intoxication} Temporary changes after using cannabis, including impaired attention, memory, coordination, judgment, or perception. Anxiety, panic, drowsiness, vomiting, or psychosis may occur; severe symptoms, accidental ingestion by a child, or unsafe behavior require urgent help.

\medterm{Marijuana/Cannabis} Cannabis is a plant containing active chemicals such as THC and cannabidiol. THC can affect memory, coordination, mood, and judgment; cannabis may also interact with medicines and can cause dependence.

\medterm{Marine Animal Stings Or Bites} Injuries caused by marine animals through biting, stinging, spines, or venom; effects range from local pain and swelling to allergic, neurologic, or systemic toxicity.

\medterm{Marine Toxin} A poisonous substance produced or accumulated by marine organisms, such as some fish, shellfish, algae, or bacteria, that can cause neurologic, gastrointestinal, or other poisoning.

\medterm{Marital Disharmony} Persistent conflict or distress in a marriage or partnership that may contribute to psychological symptoms or require relationship-focused support.

\medterm{Marital Status} A demographic or social-status variable describing whether a person is unmarried, married or partnered, separated, divorced, widowed, or in another legally or culturally recognized relationship; it may be relevant to social and health assessments.

\medterm{Marital Therapy} Counseling for couples that addresses communication, conflict, emotional distress, intimacy, and relationship problems. It may help partners improve interaction and make informed decisions about their relationship.

\medterm{Marked} Pronounced, conspicuous, or substantially greater than usual; in a clinical report, the exact meaning depends on the finding and the comparison standard.

\medterm{Marker} An identifiable biological feature used to track a gene, cell, disease, or process; a genetic marker is a detectable DNA region inherited with a nearby gene, while a blood or tumor marker is a measurable substance associated with a condition.

\medterm{Maroteaux-Lamy Syndrome} A rare inherited disorder in which certain complex sugars build up in the body. It can cause short stature, stiff joints, enlarged liver or spleen, cloudy corneas, heart-valve disease, breathing problems, and coarse facial features. Intelligence is often normal. Treatment may include enzyme replacement, supportive care, or stem-cell transplantation.

\synonyms
Mucopolysaccharidosis Type Vi; MPS VI; Arylsulfatase B Deficiency

\medterm{Marrow} The soft tissue inside bones, including red marrow that produces blood cells and yellow marrow rich in fat.

\medterm{Marrow Hyperplasia} Increased activity and numbers of blood-forming cells in the bone marrow. It may occur when the body is replacing blood lost or destroyed, or when tissues are not receiving enough oxygen.

\medterm{Marshall Syndrome} A rare periodic-fever syndrome, usually associated with PFAPA features, including recurrent fever, aphthous ulcers, pharyngitis, and cervical adenitis, often beginning in childhood.

\medterm{Marsupialization} A surgical technique that opens a cyst or abscess and sutures its edges to the surrounding skin or mucosa to permit continuing drainage.

\medterm{Martin-Bell Syndrome} An older name for fragile X syndrome, an inherited neurodevelopmental disorder caused by CGG-repeat expansion in the FMR1 gene.

\medterm{MASA Syndrome} An X-linked neurodevelopmental disorder characterized by intellectual disability, absent or limited speech, shuffling gait, spasticity, and adducted or clasped thumbs; it may also include short stature and lumbar lordosis. The name is an older acronym and terminology should be used sensitively.

\medterm{Masculinization} Development of male-typical physical characteristics, caused by normal puberty, androgen excess, androgen treatment, or certain disorders.

\medterm{MASH} Abbreviation; see \textit{Metabolic Dysfunction-Associated Steatohepatitis}.

\medterm{Mask} A device placed over the face to deliver oxygen, anesthesia, ventilation, or protection; the meaning depends on clinical context.

\medterm{Masked Data} Data altered, obscured, or replaced so that identifying values or sensitive details are hidden while selected analytical properties may be retained; masking is used to protect privacy and can be reversible or irreversible.

\medterm{Masklike Face} An expressionless face with little spontaneous animation, also called masklike facies; it may occur in Parkinson disease, myotonic dystrophy, and other disorders.

\medterm{MASLD} Abbreviation; see \textit{Metabolic Dysfunction-Associated Steatotic Liver Disease}.

\medterm{Masochism} Sexual or nonsexual arousal or gratification associated with receiving pain, humiliation, or degradation; distress or impairment determines clinical significance.

\medterm{Masoprocol} A synthetic derivative of nordihydroguaiaretic acid that inhibits lipoxygenase pathways and has been investigated for topical treatment of actinic keratoses and for antineoplastic activity.

\medterm{Mass} A discrete area of abnormal tissue that is larger or more substantial than a nodule. A mass may be benign or malignant and can occur in many organs.

\medterm{Mass Casualty Incident} An event causing more injured or ill people than nearby medical services can immediately manage. Health teams use triage to prioritize urgent treatment, transport, and limited supplies.

\medterm{Mass Screening} Testing a large group of people who do not yet have symptoms to find disease or risk early. A useful screening program has a suitable test, an effective follow-up pathway, and benefits that outweigh harms.

\medterm{Mass Spectrometer} A laboratory instrument that identifies chemicals by giving their particles an electrical charge and separating them according to size and charge. It can measure medicines, toxins, proteins, and other substances.

\medterm{Mass Spectrometry} An analytical method that identifies and measures molecules by their mass-to-charge ratios; coupled with chromatography, it is widely used in clinical chemistry, toxicology, pharmacology, and proteomics.

\medterm{Massage} Manual manipulation of soft tissues using techniques such as stroking, kneading, or pressure, used for relaxation or symptom relief; it is not a substitute for diagnosis or treatment of an underlying disease.

\medterm{Massage Therapy} Massage therapy uses hands-on manipulation of soft tissues to promote relaxation or help manage selected pain and movement problems.

\medterm{Masseter} A powerful muscle of chewing that elevates the mandible and lies over the lateral surface of the mandibular ramus.

\medterm{Masseter Muscle} A strong chewing muscle running from the cheekbone to the lower jaw. It raises the jaw to close the mouth and may become painful or enlarged with teeth grinding.

\medterm{Massive Hemoptysis} Life-threatening coughing of blood that threatens the airway, breathing, or circulation; a fixed
volume threshold is not required. It requires urgent airway assessment, stabilization, and
definitive treatment of the bleeding source.

\medterm{Massive Ovarian Edema} Marked stromal edema causing enlargement of one or both ovaries, usually due to partial torsion or hormonal stimulation. It can mimic an ovarian neoplasm and may cause pain or virilization.

\medterm{Massive Pulmonary Embolus} A high-risk, catastrophic pulmonary embolism characterized by extensive blood clot obstruction of the main pulmonary arterial tree (or major lobar arteries), leading to acute right ventricular failure and hemodynamic collapse. It is clinically defined by sustained systemic hypotension (systolic blood pressure below 90 mmHg or a drop of at least 40 mmHg), obstructive shock, or cardiac arrest. Patients present with sudden severe breathlessness, syncope (fainting), cyanosis, tachycardia, and elevated jugular venous pressure. Emergency management requires immediate supportive resuscitation, advanced mechanical clot removal (catheter-directed or surgical embolectomy), or systemic thrombolytic therapy (clot-dissolving medications).

\synonyms
Massive PE; High-Risk Pulmonary Embolism; Fulminant Pulmonary Embolism

\medterm{Massive Transfusion Protocol} A coordinated protocol for life-threatening hemorrhage that provides rapid, balanced
delivery of red cells, plasma, and platelets with monitoring and correction of
fibrinogen, calcium, temperature, and acid-base abnormalities.

\medterm{Mast Cell} An immune cell found in tissues that releases substances involved in allergy, inflammation, and defense against some germs. Its products can cause itching, swelling, mucus production, and changes in blood pressure.

\medterm{Mast Cell Activation Syndrome} A condition involving recurrent symptoms from inappropriate release of mast-cell mediators, such as flushing,
itching, hives, abdominal symptoms, wheezing, or low blood pressure. Diagnosis requires compatible episodes,
objective evidence of mediator release, and response to appropriate treatment; many similar conditions must be excluded.

\medterm{Mast Cell Sarcoma} An exceptionally rare, aggressive neoplasm composed of immature or atypical mast cells forming a destructive mass in tissue or bone.
Unlike systemic mastocytosis, it is characterized by a localized sarcomatous tumor and may progress to disseminated mast-cell disease or acute leukemia.
Clinical manifestations depend on the involved site and may include pain, swelling, organ dysfunction, or mediator-related symptoms such as flushing and hypotension.
Diagnosis requires histopathology, mast-cell immunophenotyping, and systemic staging; management generally involves specialist-directed intensive antineoplastic therapy.

\medterm{Mastalgia} Breast pain or tenderness, which may be cyclical with menstruation or noncyclical. Causes include hormonal changes, cysts, infection, trauma, and musculoskeletal pain.

\medterm{Mast-Cell Sarcoma} An extremely rare cancer made of abnormal mast cells, immune cells involved in allergic and inflammatory reactions.

\medterm{Mastectomy} Surgical removal of all or part of a breast, performed to treat or prevent selected breast cancers and sometimes for other breast conditions. Depending on the operation, the nipple, areola, and nearby lymph nodes may also be removed.

\medterm{Masticate} To chew food, using the teeth, tongue, and muscles of mastication to mechanically break it down and mix it with saliva before swallowing.

\medterm{Mastication} Chewing food with the teeth and moving it with the tongue before swallowing. It breaks food into smaller pieces and mixes it with saliva, making it easier to swallow and begin digesting.

\medterm{Masticatory Muscle} One of the muscles that moves the mandible during chewing, chiefly the masseter, temporalis, medial pterygoid, and lateral pterygoid muscles.

\medterm{Mastitis} Inflammation of breast tissue, often occurring during lactation but also possible outside lactation. It may
cause localized pain, warmth, swelling, and fever; infection and inflammatory causes
must be distinguished clinically.

\synonyms
Puerperal Mastitis; Breast Inflammation

\medterm{Mastocytoma} A localized collection or tumor of mast cells, usually presenting in the skin as a solitary cutaneous mastocytoma; rubbing the lesion can cause itching, redness, or swelling from mediator release.

\medterm{Mastocytosis} A disorder in which too many mast cells, immune cells involved in allergic reactions, collect in the skin or internal organs. It may cause itchy spots, flushing, abdominal symptoms, wheezing, or low blood pressure.

\medterm{Mastodynia} Pain or tenderness in one or both breasts, also called mastalgia. It may follow the menstrual cycle or result from injury, infection, medicines, cysts, or other breast conditions; a new persistent lump, skin change, or bloody discharge needs assessment.

\medterm{Mastoid} Relating to the mastoid part of the temporal bone behind the ear. This area contains air spaces connected with the middle ear and can become infected, causing pain, swelling, or drainage.

\medterm{Mastoid Antrum} An air-containing cavity in the mastoid portion of the temporal bone that communicates with the middle-ear attic.

\medterm{Mastoid Bone} The portion of the temporal bone located behind the ear.

\medterm{Mastoid Process} The conical bony projection of the temporal bone behind the ear, containing air cells and providing attachment for several neck muscles; infection of its air cells is mastoiditis.

\medterm{Mastoidectomy} Surgery to remove infected or diseased air cells in the bone behind the ear. It may treat chronic middle-ear infection, cholesteatoma, mastoiditis, or complications that have spread beyond the ear.

\medterm{Mastoiditis} Mastoiditis is infection and inflammation of the mastoid air cells behind the ear, usually following middle-ear infection and potentially causing serious spread.

\medterm{Mastopexy} Breast-lift surgery that removes extra skin and reshapes breast tissue to raise the breasts and nipples. It can improve shape but does not necessarily reduce breast size; scars, altered sensation, asymmetry, and changes in breastfeeding ability are possible risks.

\synonyms
Breast Lift

\medterm{Masturbation} Self-stimulation of the genitals for sexual pleasure. It is a common human behavior; it is generally
not harmful unless it causes injury, distress, or interferes with daily life.

\synonyms
Autoeroticism; Solitary Sex

\medterm{Materia Medica} An older term for the body of medicinal substances and information about their properties and uses.

\medterm{Maternal} Maternal means relating to a mother or the pregnant person, especially in connection with pregnancy, childbirth, or postpartum health.

\medterm{Maternal Age} Maternal age is the age of the pregnant person and can influence fertility, pregnancy risks, screening choices, and obstetric care.

\medterm{Maternal Expulsive Efforts} The pregnant person's efforts to push the baby down and out during the second stage of labor. Pushing may follow the body's urge or be guided by the birth team, depending on the situation and chosen pain relief.

\synonyms
Bearing Down; Second Stage Expulsive Efforts

\medterm{Maternal Health} The physical and mental health of a person during pregnancy, childbirth, and the postpartum period. Care includes prevention, antenatal monitoring, safe delivery, and treatment of complications.

\medterm{Maternal Mortality} Maternal mortality is death during pregnancy or within the defined postpartum period from a pregnancy-related cause, measured as a population health indicator.

\medterm{Maternal Mortality Rate} The number of maternal deaths related to pregnancy per 100,000 live births in a specified population and period.

\medterm{Maternal Obesity} Obesity present before or during pregnancy. It can increase the risk of high blood pressure, diabetes, cesarean birth, and health problems for the pregnant person or baby.

\medterm{Maternal Physiologic Changes In Pregnancy} Normal body changes during pregnancy that support the growing fetus and prepare for birth. They affect blood volume, heart function, breathing, digestion, kidneys, hormones, breasts, skin, and the musculoskeletal system.

\medterm{Maternal Sepsis} Organ dysfunction caused by a dysregulated maternal response to infection during
pregnancy, childbirth, or the postpartum period. It may arise from uterine, urinary,
respiratory, wound, or other infection and requires prompt cultures, antimicrobial
therapy, source control, and hemodynamic support.

\medterm{Maternal Serum Alpha-Fetoprotein} A blood test during pregnancy, usually at 15--20 weeks, that measures a protein made by the fetus. An unusually high or low result may be linked to an incorrect pregnancy date, more than one fetus, or certain birth defects or chromosome conditions. It is a screening test, not a diagnosis, and abnormal results need follow-up.

\medterm{Maternal Vaccine} A vaccine given during pregnancy to protect the pregnant person and, through transferred antibodies, provide early protection to the fetus or newborn; recommended products and timing depend on current guidance.

\medterm{Maternity} The state or service relating to pregnancy, childbirth, and the care of the mother and newborn.

\medterm{Maternity Services} Healthcare for pregnancy, labor and birth, the postpartum period, and the newborn. Services may include prenatal checks, midwifery or obstetric care, delivery support, breastfeeding help, and emergency care.

\medterm{Materteral} Relating to a maternal aunt or to the genetic relationship between an aunt and her niece or nephew; synonymous with materterine.

\medterm{Mathematics Disorder} A learning disorder involving persistent difficulty understanding numbers, calculation, quantity, or mathematical reasoning despite appropriate instruction.

\medterm{Mating} Sexual pairing or reproduction between organisms that permits exchange or transfer of gametes and may lead to fertilization.

\medterm{Matrix Band} A thin dental strip placed around a prepared tooth to temporarily replace a missing wall and contain restorative material while a filling or crown buildup is formed.

\medterm{Matrix Metalloproteinase} An enzyme that cuts proteins in the extracellular matrix surrounding cells. It helps with normal tissue repair and remodeling, but excessive or poorly controlled activity can promote inflammation, tissue destruction, or cancer spread.

\medterm{Maturation} The process through which cells, tissues, organs, or a person develop and become fully functional. Normal maturation may involve changes in size, structure, ability, hormone activity, and response to the environment.

\medterm{Mature Onset Diabetes} An older term generally referring to type 2 diabetes mellitus, which often develops in adulthood, though diabetes classification depends on cause rather than age alone.

\medterm{Mature Teratoma} A germ-cell tumor composed of well-differentiated tissues from one or more embryonic germ layers; many are benign, but behavior depends on the site, age, and tumor type, and some can undergo malignant transformation.

\medterm{Maxilla} The paired bones forming the upper jaw and supporting the upper teeth. They also form parts of the cheeks, nose, eye sockets, hard palate, and the air-filled sinuses behind the cheeks.

\medterm{Maxilla} The paired bones forming the upper jaw, the front part of the hard palate, and parts of the walls of the nose and eye sockets.

\synonyms
Superior Maxilla; Upper Jawbone

\medterm{Maxillary} Relating to the maxilla, the upper jawbone. It supports the upper teeth, forms part of the eye socket and roof of the mouth, and contains the maxillary sinuses. Problems may include tooth infection, sinusitis, or facial fractures.

\medterm{Maxillary Sinus} The maxillary sinus is a paired, air-filled cavity in the maxilla beneath each orbit. It drains into the nasal cavity and may become inflamed or infected in sinusitis.

\medterm{Maxillary Sinusitis} Inflammation or infection of the air-filled spaces behind the cheekbones. It may cause cheek or tooth pain, blocked nose, thick nasal discharge, reduced smell, cough, or fever and can follow a cold, allergy, or dental infection.

\medterm{Maxillofacial} Relating to the jaw, face, mouth, and associated bones and soft tissues. Maxillofacial care may include treatment of injuries, congenital differences, infections, and tumors.

\medterm{Maxillofacial Prosthesis} An artificial replacement for part of the face, jaw, mouth, nose, eye area, or ear after injury, cancer surgery, or a birth difference. Examples include devices that close a palate opening or replace an eye, nose, or ear.

\medterm{Maximal Breathing Capacity} The greatest volume of air that can be moved in and out of the lungs in a specified period during maximal effort; it is related to but distinct from standard spirometric measures.

\medterm{Maximum} The maximum is the greatest value in a set of measurements or the highest level reached by a variable, such as a recorded temperature, blood pressure, or drug concentration.

\medterm{Maximum Plasma Concentration} The highest measured plasma concentration after a dose. Cmax and the time to reach it
are used to describe absorption rate, peak-related toxicity, and comparative
pharmacokinetics.

\synonyms
Cmax

\medterm{Maximum Thickness} Maximum thickness is the greatest measured thickness of a structure or specimen. It may be reported in imaging, pathology, or assessment of a lesion, organ, or tissue layer.

\medterm{Maximum Tolerated Dose} The highest dose of a drug that produces acceptable, manageable toxicity without causing
unacceptable adverse effects, determined through dose-escalation studies, most commonly
in Phase I clinical trials. The MTD identifies the dose range for subsequent efficacy
studies.

\medterm{Mayer-Rokitansky-KüSter-Hauser Syndrome} A condition present from birth in which the uterus and much of the vagina do not develop, although the ovaries and usual female puberty changes develop normally. It often first becomes noticeable when menstruation does not begin. Kidney, hearing, or spine abnormalities may also occur. Treatment and fertility options depend on the individual’s anatomy and goals.

\synonyms
MRKH Syndrome; Müllerian Agenesis; Vaginal Agenesis

\medterm{May-Hegglin Anomaly} An inherited MYH9-related disorder with macrothrombocytopenia, giant platelets, and characteristic leukocyte inclusions; hearing, kidney, or cataract abnormalities may occur.

\medterm{Mayo Scissors} Strong surgical scissors designed to cut fascia, muscle, sutures, and other tough materials. Their broad blades provide force and control, and they may be straight or curved for different surgical tasks.

\medterm{Maytansine} Maytansine is a potent antimitotic compound that binds tubulin and inhibits microtubule assembly. Its derivatives are used as cytotoxic payloads in some antibody--drug conjugates.

\medterm{May-Thurner Syndrome} Compression of the left pelvic vein by the right pelvic artery, which can slow blood flow and increase the risk of a clot or swelling in the left leg.

\medterm{Maze Procedure} A heart operation for atrial fibrillation. The surgeon creates carefully placed lines of scar tissue in the upper heart chambers to block abnormal electrical signals and help restore a regular rhythm. It may be performed with other heart surgery; recurrence, bleeding, stroke, or the need for a pacemaker are possible.

\synonyms
Cox-Maze Procedure; Surgical Ablation for Atrial Fibrillation; Maze Operation

\medterm{Mazindol} Mazindol is a sympathomimetic appetite suppressant that affects central catecholamine signaling. It has been used for short-term obesity treatment in some countries and may cause stimulant-related adverse effects.

\medterm{MB ChB} Abbreviation for Bachelor of Medicine and Bachelor of Surgery, a medical degree awarded in many countries.

\medterm{MC Lymphoma} Abbreviation; see \textit{Mantle Cell Lymphoma}.

\medterm{MCA Peak Systolic Velocity} An ultrasound measurement of blood speed in a major artery of the fetus's brain. An unusually high value may suggest fetal anemia and can help guide monitoring or further testing during pregnancy.

\medterm{MCAD Deficiency} An inherited disorder in which the body cannot properly break down certain medium-chain fats for energy.

\medterm{McArdle Disease} A rare inherited muscle disorder in which muscles cannot release usable energy from stored sugar during exercise. Activity can cause early muscle pain, cramps, exhaustion, or dangerous muscle breakdown with dark urine. Some people improve after slowing down briefly and restarting gently, called the second-wind phenomenon. Avoiding sudden strenuous exercise helps prevent attacks.

\synonyms
GSD V; Myophosphorylase Deficiency; Glycogen Storage Disease Type V

\medterm{McBurney Sign} Tenderness when pressing on the lower right side of the belly about halfway between the belly button and the hip bone, which is a key sign of an infected appendix (appendicitis).

\medterm{McBurney's Point} A location in the lower-right abdomen where tenderness may occur when the appendix is inflamed. Pain there can support a diagnosis of appendicitis but cannot confirm it by itself.

\medterm{Mcc Gene} A gene that helps regulate cell growth and is considered a tumour-suppressor gene. Loss or alteration of its function has been found in some colorectal and other cancers, but the finding alone does not diagnose cancer.

\medterm{McCune-Albright Syndrome} A condition caused by a change affecting hormone signals in some body cells. It may cause irregular skin coloring, fragile or unevenly growing bones, and early puberty or other hormone problems.

\medterm{McDonald Cerclage} A common operation that places a strong stitch around the cervix through the vagina. It may help keep the cervix closed when there is a high risk of very early birth; timing and suitability depend on the pregnancy history and examination.

\medterm{MCH} The average amount of hemoglobin in one red blood cell, reported in a complete blood count. Hemoglobin carries oxygen, and the result helps classify some types of anemia.

\medterm{MCHC} The average concentration of hemoglobin inside red blood cells, reported in a complete blood count. It helps clinicians assess and classify anemia together with the other blood-count results.

\medterm{MCL Injury} Abbreviation; see \textit{Medial Collateral Ligament Injury}.

\medterm{McMurray Test} A knee examination in which the leg is flexed, rotated, and extended to place stress on the menisci. Pain or a palpable click may suggest a meniscal lesion but is not diagnostic alone.

\medterm{McRoberts Maneuver} A positioning maneuver for managing shoulder dystocia where the mother’s legs are
hyperflexed onto her abdomen, with thighs against the abdomen. This flattens the sacrum,
increases the anteroposterior diameter of the pelvis, and rotates the symphysis pubis
superiorly. Often combined with suprapubic pressure.

\medterm{MCRPC} Abbreviation; see \textit{Metastatic Castration-Resistant Prostate Cancer}.

\medterm{MCV} Mean corpuscular volume: the calculated average volume of a red blood cell, derived from hematocrit and red-cell count and reported as part of a complete blood count.

\medterm{MDMA} 3,4-Methylenedioxymethamphetamine, a psychoactive stimulant and hallucinogen commonly known as ecstasy or molly; it can produce serious cardiovascular, neurologic, and temperature-regulation effects.

\medterm{MDS} Abbreviation; see \textit{Myelodysplastic Syndromes}.

\medterm{Mead Acid} An omega-9 fatty acid made when the body lacks enough essential omega-6 and omega-3 fatty acids. Increased levels can support laboratory assessment of essential-fatty-acid deficiency, although results must be interpreted with nutrition and illness history.

\medterm{Meadow Saffron} Colchicum autumnale, a plant containing colchicine; ingestion can cause severe gastrointestinal illness, multiorgan toxicity, and death.

\medterm{Mean} The mean is the arithmetic average, calculated by adding all values and dividing by the number of values. It is commonly used to summarize clinical measurements but can be distorted by extreme values.

\medterm{Mean Arterial Pressure} An estimate of the average pressure pushing blood through the body's organs during each heartbeat. Clinicians use it with other findings to judge whether organs are receiving enough blood.

\medterm{Mean Cell Volume} The average volume of a red blood cell, calculated from the hematocrit and red-cell count and reported as MCV in a complete blood count.

\medterm{Mean Corpuscular Hemoglobin Concentration} The average concentration of oxygen-carrying hemoglobin inside red blood cells. It is part of a complete blood count and helps classify the cause of anemia.

\medterm{Mean Corpuscular Volume} The average size of red blood cells, reported in a complete blood count. Small cells or large cells can point toward different causes of anemia and must be interpreted with the other blood results.

\medterm{Mean Diameter} Mean diameter is the average diameter of a structure or lesion, usually calculated from one or more measured dimensions. The method used should be specified when measurements are clinically compared.

\medterm{Mean Thickness} Mean thickness is the average thickness calculated from multiple measurements of a tissue, structure, or specimen. It is used in imaging, pathology, and physiologic measurements.

\medterm{Measles} A highly contagious viral infection causing fever, cough, runny nose, red eyes, and a rash that spreads from the face downward. It can cause pneumonia, brain inflammation, or other serious complications.

\medterm{Measles Virus} A negative-sense RNA virus that causes measles, a highly contagious respiratory illness.
The virus is transmitted by respiratory droplets and is one of the most infectious
pathogens known. Symptoms include fever, cough, coryza, conjunctivitis, and rash.

\medterm{Measurable Residual Disease} Cancer cells that remain after treatment and can be detected with a sensitive, validated test. The amount may be too small to see on scans or by routine examination, but it can help estimate relapse risk and guide further treatment.

\medterm{Measurement} Measurement is the quantitative determination of a physical, biologic, or clinical characteristic using a defined method and unit. Interpretation depends on accuracy, precision, and the clinical context.

\medterm{Meatal} Relating to a meatus, a natural opening or passage in the body. It commonly refers to the opening of the urethra, the ear canal, or one of the passages inside the nose. Meatal narrowing can obstruct urine or air flow.

\medterm{Meatal Stenosis} Narrowing of the urethral meatus, which may cause a deflected or narrow urinary stream, spraying, dysuria, or difficulty emptying the bladder.

\medterm{Meatus} A natural opening or passage into or out of a body structure. Examples include the opening of the urethra, the ear canal, and the small passages inside the nose.

\medterm{Meatuses} Natural openings or passages in the body. In the nose, the meatuses are channels beneath the nasal folds that help air move through the nose and allow drainage from the sinuses and tear ducts.

\medterm{Mebendazole} A medicine used to treat several intestinal worm infections. It disrupts structures the worms need to absorb nutrients and survive; treatment choices and dosing depend on the parasite, age, pregnancy, and local guidance.

\medterm{Mebrofenin Scan} A nuclear-medicine scan that follows a tracer through the liver and bile ducts. It can assess gallbladder function, bile blockage, leaks, or selected causes of jaundice and abdominal pain.

\medterm{Mecamylamine} Mecamylamine is an orally active ganglionic nicotinic-receptor blocker that reduces sympathetic nerve transmission and lowers blood pressure. It has been used for severe hypertension and may cause postural hypotension, dizziness, constipation, and dry mouth.

\medterm{Mecarzole} Mecarzole is a synonym for carbendazim, a benzimidazole fungicide that interferes with fungal microtubule formation. It is an agricultural chemical rather than a routine human medicine and may be toxic if exposure occurs.

\medterm{ME/CFS} Abbreviation; see \textit{Myalgic Encephalomyelitis}.

\medterm{Mechanical Circulatory Support} Machines or pumps that help the heart move blood when it cannot pump enough on its own. They may be used temporarily during severe illness or longer term while awaiting, or instead of, a heart transplant.

\medterm{Mechanical Insufflation-Exsufflation} A cough-assist device that gives a breath in followed quickly by suction-like pressure out. It helps people with weak breathing muscles clear mucus, especially in neuromuscular disease or spinal-cord injury.

\medterm{Mechanical Ventilation} Breathing support from a machine that pushes air and oxygen into the lungs and helps remove carbon dioxide. It may be delivered through a mask or a breathing tube when illness or injury prevents adequate breathing.

\medterm{Mechanical Ventilator} A mechanical ventilator is a machine that supports or replaces spontaneous breathing by delivering controlled or assisted breaths, usually through a mask or an endotracheal tube. Settings are adjusted to oxygenation, ventilation, mechanics, and patient needs.

\medterm{Mechanism Of Death} The physiologic or biochemical derangement that causes death (the final failure) as
opposed to the underlying disease or injury (cause of death) and the circumstances
(manner). Examples: ventricular fibrillation, cerebral hemorrhage, sepsis, respiratory
failure, exsanguination, cardiac tamponade.

\medterm{Mechanism Of Toxicity} The way a substance causes injury or illness in the body. It may interfere with an enzyme or receptor, damage cell membranes or DNA, disrupt energy production, or trigger an abnormal immune response. Understanding the mechanism helps guide prevention and treatment.

\medterm{Mechanoreceptor} A sensory receptor that detects physical forces such as touch, pressure, stretch, vibration, sound, or movement. It converts mechanical stimulation into electrical signals carried by nerves to the brain or spinal cord.

\medterm{Mechlorethamine} Mechlorethamine is an alkylating antineoplastic agent that cross-links DNA and inhibits cell division. Topical gel is used for cutaneous T-cell lymphoma; systemic use is limited by toxicities such as myelosuppression and tissue injury.

\medterm{Meckel Diverticulectomy} Surgical removal of a Meckel diverticulum, usually performed when it causes bleeding, inflammation, obstruction, perforation, or is found incidentally in a high-risk setting.

\medterm{Meckel Diverticulum} The most common congenital anomaly of the gastrointestinal tract, occurring as a true diverticulum (containing all intestinal wall layers) on the antimesenteric border of the distal ileum due to incomplete closure of the embryonic vitelline (omphalomesenteric) duct. It classically follows the ``Rule of 2s'': found in approximately 2\% of the population, located roughly 2 feet proximal to the ileocecal valve, measuring about 2 inches in length, and frequently containing ectopic tissue (most commonly ectopic gastric mucosa or ectopic pancreatic tissue). Acid secretion from ectopic gastric mucosa causes ulceration and painless rectal bleeding (hematochezia), especially in young children. Other complications include diverticulitis (which closely mimics acute appendicitis), bowel obstruction, and intussusception. Diagnosis of bleeding diverticula is confirmed by a Meckel scan (technetium-99m pertechnetate scintigraphy), and symptomatic cases are treated with surgical diverticulectomy.

\synonyms
Meckel's Diverticulum; Vitelline Duct Remnant; Omphalomesenteric Diverticulum; Congenital Ileal Diverticulum

\medterm{Meckel-Gruber Syndrome} A severe inherited birth disorder that causes major abnormalities of the brain, kidneys, and limbs. Typical findings include a sac of brain tissue at the back of the skull, enlarged cystic kidneys, and extra fingers or toes. Poor lung development is common, and affected pregnancies usually do not survive to birth.

\synonyms
Meckel Syndrome; Dysencephalia Splanchnocystica

\medterm{Meckel's Diverticulum} Alternate name; see \textit{Meckel Diverticulum}.

\medterm{Meclizine} Meclizine is a first-generation antihistamine with anticholinergic and vestibular-suppressant effects. It is used to prevent or treat motion sickness and vertigo and may cause drowsiness, dry mouth, and blurred vision.

\medterm{Meclofenamate Overdose} Taking too much meclofenamate, a nonsteroidal anti-inflammatory drug, may cause nausea, vomiting, stomach pain or bleeding, drowsiness, kidney injury, seizures, or metabolic acidosis. Urgent medical or poison-control advice is needed after a suspected overdose.

\medterm{Meconium} The thick, dark first stool of a newborn, composed of intestinal secretions, shed cells, mucus, and swallowed amniotic material.

\medterm{Meconium Aspiration} Breathing meconium, a newborn's first stool, into the airways before or around birth. It can block airflow and inflame the lungs, causing breathing difficulty and low oxygen.

\medterm{Meconium Aspiration Syndrome} Breathing problems in a newborn after meconium, the first stool, enters the airways before or during birth. It can block airways and inflame the lungs, causing rapid breathing, low oxygen, and need for respiratory support.

\medterm{Meconium Ileus} A blockage of a newborn’s small intestine by unusually thick first stool. It is strongly associated with cystic fibrosis and may cause a swollen abdomen, green vomiting, and failure to pass stool soon after birth.

\medterm{Meconium Peritonitis} Inflammation of the tissue lining a fetus’s abdomen after a hole in the intestine allows first stool to escape before birth. It can cause abdominal fluid or calcium deposits and may be linked to bowel blockage or cystic fibrosis.

\medterm{Meconium-Stained Amniotic Fluid} Green or yellow amniotic fluid caused by passage of fetal stool before birth. It may occur with a mature or distressed fetus and increases the risk of breathing problems after delivery; newborns are assessed and supported as needed.

\medterm{MECP2 Duplication Syndrome} A rare X-linked genetic disorder caused by extra copies or increased dosage of the MECP2 gene. It often causes developmental delay, low muscle tone, seizures, recurrent infections, and breathing or feeding problems, especially in boys.

\medterm{Medial} Medial means toward or situated nearer the body's midline; it contrasts with lateral and is used to describe anatomical position, movement, or imaging findings.

\medterm{Medial Collateral Ligament Injury} A stretch or tear of the strong ligament on the inner side of the knee, often caused by a blow to the outer knee or a twisting injury. It causes inner-knee pain and swelling; severe tears may make the knee unstable.

\medterm{Medial Cuneiform} The largest of three small bones in the inner middle part of the foot. It joins the navicular bone and the first toe bone, helps support the arch, and forms part of the joint that can be injured in a Lisfranc injury.

\medterm{Medial Epicondylitis} Alternate medical name; see \textit{Golfer's Elbow}.

\medterm{Medial Meniscus} The medial meniscus is a crescent-shaped fibrocartilage pad between the medial femoral condyle and medial tibial plateau. It distributes load and contributes to knee stability; tears may cause joint-line pain, swelling, or locking.

\medterm{Medial Rotation} Turning a limb toward the body's midline. For example, turning the thigh so the knee and foot point inward is medial rotation; the movement is assessed at joints such as the shoulder and hip.

\medterm{Medial Tibial Stress Syndrome} Alternate name; see \textit{Shin Splints}.

\medterm{Median} Located in or relating to the middle plane or midpoint of a structure. In medicine, it describes anatomy, measurements, statistical values, nerve positions, or locations between two sides.

\medterm{Median Arcuate Ligament Syndrome} Compression of the celiac artery and nearby nerves by the median arcuate ligament, potentially causing post-meal abdominal pain and weight loss.

\synonyms
Dunbar Syndrome; MALS

\medterm{Median Nerve} A major nerve running from the neck through the arm and forearm to the hand. It helps bend the wrist and fingers, moves the thumb, and provides feeling to the thumb, index, middle, and part of the ring finger; compression can cause carpal-tunnel symptoms.

\medterm{Median Nerve Compression} Alternate name; see \textit{Carpal Tunnel Syndrome}.

\medterm{Median Plane} A sagittal plane that passes directly through the midline of the body, dividing it into
equal right and left halves.

\medterm{Median Rhomboid Glossitis} A smooth, red, diamond-shaped depapillated area on the dorsum of the tongue caused by
chronic Candida infection, usually asymptomatic and managed with antifungal therapy.

\medterm{Median Survival} The time by which half of the people in a study have experienced the specified event, such as death. It is a summary measure of survival in a defined population.

\medterm{Mediastinal Cyst} A mediastinal cyst is a fluid- or tissue-filled cystic lesion in the mediastinum. Common types include bronchogenic, thymic, and pericardial cysts; symptoms depend on size and compression, and imaging guides evaluation.

\medterm{Mediastinal Disease} A disorder affecting structures in the mediastinum, including the heart, great vessels, thymus, lymph nodes, nerves, and tumors or cysts within the central chest.

\medterm{Mediastinal Emphysema} Mediastinal emphysema, or pneumomediastinum, is the presence of free air in the mediastinum. It can follow alveolar rupture, trauma, surgery, or esophageal injury and may cause chest pain, neck crepitus, or respiratory compromise.

\medterm{Mediastinal Hemorrhage} Bleeding into the mediastinum, commonly caused by trauma, vascular rupture, surgery, or invasive procedures, and potentially causing airway or cardiovascular compression.

\medterm{Mediastinal Mass} An abnormal growth or collection of tissue in the mediastinum, the central compartment of the chest. Causes include tumors, enlarged lymph nodes, cysts, and vascular abnormalities.

\medterm{Mediastinal Neoplasm} A mediastinal neoplasm is an abnormal growth in the central space of the chest between the lungs; it may arise from lymph nodes, thymus, nerves, or other tissues.

\medterm{Mediastinal Pleura} The mediastinal pleura is the portion of parietal pleura lining the lateral surfaces of the mediastinum. It separates the mediastinal contents from the pleural cavity and reflects around the lung root.

\medterm{Mediastinal Shift} Mediastinal shift is displacement of the mediastinal structures away from their usual position. It may result from increased pressure on one side, such as a tension pneumothorax, or loss of volume on the other, such as lung collapse.

\medterm{Mediastinal Tumor} A mediastinal tumor is a mass in the central chest compartment between the lungs and may be benign or malignant.

\medterm{Mediastinitis} Mediastinitis is inflammation or infection of the central chest compartment, often after esophageal rupture, surgery, trauma, or spread of infection.

\medterm{Mediastinoscopes} Mediastinoscopes are slender endoscopic instruments used during mediastinoscopy to inspect the mediastinum and obtain tissue samples, especially from mediastinal lymph nodes. The procedure is performed through a small incision near the sternum.

\medterm{Mediastinoscopy} A surgical procedure that accesses the mediastinum through a small incision near the sternum to
inspect or biopsy mediastinal lymph nodes or masses. It is one option for staging lung cancer
when less-invasive sampling is nondiagnostic or unsuitable.

\medterm{Mediastinoscopy With Biopsy} A procedure using an instrument inserted through a small incision near the sternum to inspect the mediastinum and obtain lymph-node or mass samples for diagnosis and staging.

\medterm{Mediastinum} The central part of the chest between the lungs. It contains the heart, major blood vessels, windpipe, food pipe, thymus, nerves, and lymph nodes.

\medterm{Mediators Of Inflammation} Molecules that start, strengthen, limit, or coordinate inflammation. They include histamine, prostaglandins, cytokines, and other signals that change blood vessels, attract immune cells, cause symptoms, and support healing.

\medterm{Medical Cannabis And Cannabidiol} Cannabis and cannabidiol are substances from the cannabis plant used in some medical settings. Their effects and risks vary by product and dose; they may cause sleepiness, impaired judgment, interactions, or dependence and do not treat every condition.

\medterm{Medical Certificate} A document completed by an authorized healthcare professional that records a medical finding, diagnosis, fitness status, or period of illness.

\medterm{Medical Condition} A state affecting health or body function, whether temporary, long-lasting, symptomatic, or silent. A medical condition may result from disease, injury, inherited traits, infection, or normal biological changes.

\medterm{Medical Device} An instrument, implant, machine, software program, or other product used to diagnose, prevent, monitor, or treat disease or injury. Examples include pacemakers, syringes, imaging equipment, and blood-glucose meters.

\medterm{Medical Entomology} The study of insects and related arthropods that cause disease or transmit disease-causing organisms to humans, including their biology, ecology, physiology, and genetics.

\medterm{Medical Ethics} Moral principles that guide healthcare decisions, including respect for a person's choices, helping rather than harming, fairness, privacy, and honest communication. Decisions consider the patient's wishes, likely benefits and harms, professional duties, and relevant law.

\medterm{Medical Examiner} A physician who investigates sudden, unexpected, violent, or suspicious deaths. The examiner may order an autopsy, determine the medical cause of death, and provide information for the official death certificate; exact duties vary by location.

\medterm{Medical History} A medical history is a structured account of a person’s current concerns, past illnesses and operations, medications, allergies, family history, social history, and relevant exposures. It helps guide examination, diagnosis, and treatment.

\medterm{Medical Illness} The subjective experience of being unwell, which may or may not correspond to an
identifiable disease.

\medterm{Medical Jargon} Specialized words and phrases used by healthcare professionals. If they are not explained, patients may misunderstand their condition, test results, or treatment, so clinicians should use plain language whenever possible.

\medterm{Medical Jurisprudence} The branch of law (and its study) that deals with the rights and duties of physicians
and the legal aspects of medical practice, distinct from but overlapping with forensic
medicine.

\medterm{Medical Malpractice} Legal tort occurring when a healthcare provider fails to exercise the standard of
reasonable care, skill, or diligence expected of a prudent practitioner of similar
standing, resulting in injury, harm, or death to a patient.

\synonyms
Clinical Malpractice; Medical Malpractice

\medterm{Medical Marijuana} Cannabis used under medical supervision to relieve symptoms or support treatment of selected conditions. Products contain cannabinoids such as THC or cannabidiol; effects, interactions, impairment, and dependence risk vary by product and dose.

\medterm{Medical Marijuana} Alternate name; see \textit{Medical Cannabis}.

\medterm{Medical Nutrition Therapy} Personalized food and nutrition care provided by a qualified dietitian to manage a health problem. It may support people with diabetes, kidney disease, digestive disorders, cancer, obesity, or malnutrition.

\medterm{Medical Officer} A physician appointed to provide clinical care or medical oversight for an institution, service, or organization.

\medterm{Medical Physicist} A specialist who applies physics to medical imaging, radiation treatment, and radiation safety. Medical physicists check equipment, calculate radiation doses, and help ensure that scans and cancer treatments are accurate and as safe as possible.

\medterm{Medical Physics} The application of physics to medicine, including diagnostic imaging, radiation therapy, radiation safety, and physiological
measurement.

\medterm{Medical Pin} A medical pin is a slender metal fixation device inserted into bone or across a joint to stabilize a fracture or support orthopedic treatment. Pins may be temporary or permanent and require monitoring for infection and displacement.

\medterm{Medical Practitioner} A licensed physician qualified to diagnose, prevent, and treat illness.

\medterm{Medical Privacy} Protection of information that can identify a patient and control over its sharing. It includes keeping health records secure, limiting access to appropriate people, and obtaining permission when disclosure is not otherwise allowed.

\medterm{Medical Procedure} A planned clinical action or intervention performed by healthcare professionals to diagnose, measure, monitor, treat, or prevent a health condition, or to improve bodily function and structure. Medical procedures range in complexity from non-invasive evaluations (such as an electrocardiogram, ultrasound, or blood pressure check) to minimally invasive tests (such as drawing blood, inserting a catheter, performing an endoscopy, or taking a needle biopsy) and major invasive surgeries (such as tumor resections, joint replacements, or organ transplants). Procedures are classified by their primary clinical purpose: diagnostic (identifying the underlying cause of disease), therapeutic or curative (treating an illness, correcting a defect, or removing damaged tissue), palliative (relieving symptoms, obstruction, or pain to improve quality of life), or preventive/prophylactic (such as screening colonoscopies to remove precancerous polyps). Depending on the level of invasiveness and tissue involvement, procedures may be performed awake, with local numbing, under procedural sedation, or under general anesthesia, requiring strict sterile technique to reduce risks of bleeding, infection, and tissue injury.

\synonyms
Clinical Procedure; Healthcare Procedure; Medical Intervention; Surgical Procedure

\medterm{Medical Records} Written or electronic information about a person's health, including symptoms, examinations, test results, diagnoses, medicines, procedures, and follow-up. Accurate records help healthcare professionals coordinate safe care and track changes over time.

\medterm{Medical Repository} The first medical magazine and scientific journal published in the United States; founded in New York in 1797, it was published quarterly until 1824.

\medterm{Medical Symbol} A commonly recognized sign representing medicine, especially the Rod of Asclepius: a staff with one serpent. The winged, two-serpent caduceus is often used incorrectly as a medical symbol.

\medterm{Medical Tattooing} Medical tattooing (permanent makeup, micropigmentation) covers depigmented skin
(vitiligo, vitiligo camouflage), restores areolas after mastectomy, and masks scars.
Trichopigmentation simulates hair on bald scalp. Medical tattooing uses specialized
pigments and techniques.

\medterm{Medical Tests} Diagnostic procedures used to screen for, confirm, or monitor disease through analysis
of biologic specimens, imaging, physiologic measurements, or genetic material.

\medterm{Medicinal} Used for the prevention, diagnosis, relief, or treatment of disease.

\medterm{Medicine} A substance or treatment used to prevent, diagnose, relieve, or treat disease. Medicines may be prescription, over-the-counter, biological, or traditional, and all can have benefits, risks, and interactions.

\medterm{Medicines For Back Pain} Medicines for back pain may include selected anti-inflammatory drugs, acetaminophen, muscle relaxants, or other treatments depending on the cause and patient risks.

\medterm{Medicines For Osteoporosis} Osteoporosis medicines reduce bone loss or increase bone formation and may include bisphosphonates, denosumab, or anabolic treatments chosen according to fracture risk.

\medterm{Medicines For Sleep} Medicines for sleep may help short-term insomnia but can cause dependence, falls, confusion, or next-day impairment; treatment should also address the cause of poor sleep.

\medterm{Medicolegal} Pertaining to the application of medical knowledge to legal questions and the
intersection of medicine and the law. Covers forensic pathology, toxicology, psychiatry,
and clinical issues: informed consent, malpractice, medical records, disability,
personal injury, occupational disease, and public health law.

\medterm{Meditation} Meditation is a practice of intentionally directing attention, often toward the breath, a sensation, or an object, while observing thoughts without judgment. It may promote relaxation and stress management but does not replace indicated medical treatment.

\medterm{Mediterranean Anemia} Alternate name; see \textit{Thalassemia}.

\medterm{Mediterranean Diet} An eating pattern rich in vegetables, fruits, beans, whole grains, nuts, olive oil, and fish, with less red and processed meat and added sugar. It can support cardiovascular health as part of an overall balanced diet.

\medterm{Mediterranean Fever} An older name for brucellosis, a bacterial infection that can cause prolonged fever, sweats, fatigue, and joint or muscle pain.

\medterm{Medium} A substance or condition through which something occurs; in microbiology, a culture medium supports the growth of microorganisms.

\medterm{MEDLARS} Medical Literature Analysis and Retrieval System, an electronic bibliographic system developed by the U.S. National Library of Medicine to index and retrieve biomedical literature; it preceded MEDLINE.

\medterm{Medroxyprogesterone} Medroxyprogesterone is a progestogen used for contraception, abnormal uterine bleeding, endometriosis, and hormone therapy; depot injections may cause irregular bleeding, weight change, delayed fertility return, and reduced bone density.

\medterm{Medroxyprogesterone Acetate} Medroxyprogesterone acetate is a synthetic progestin used in hormonal contraception, treatment of endometriosis or abnormal uterine bleeding, and hormone therapy. It suppresses ovulation and produces progestational effects; formulation-specific risks and contraindications apply.

\medterm{Medulla} The inner or central part of an organ, such as the renal medulla or adrenal medulla.

\medterm{Medulla Oblongata} The lower part of the brainstem, joined to the spinal cord. It helps control breathing, heart rate, blood pressure, swallowing, coughing, vomiting, and other automatic functions.

\medterm{Medullary} Relating to the medulla or inner portion of an organ.

\medterm{Medullary Carcinoma} A cancer whose cells grow in broad groups and are surrounded by many immune cells. The term is used most often for cancers of the breast or thyroid, and its behavior depends on the organ and tumor findings.

\medterm{Medullary Sponge Kidney} A condition present from birth in which tiny urine-collecting tubes inside the kidneys are widened. Urine minerals may form stones or deposits, causing blood in the urine, repeated infections, pain, or no symptoms.

\medterm{Medullary Thyroid Carcinoma} A cancer of hormone-producing thyroid cells that make calcitonin. It may occur alone or as part of an inherited syndrome; blood tests, imaging, and a biopsy help confirm it.

\medterm{Medulloblastoma} A fast-growing cancer that usually begins in the cerebellum, the part of the brain involved in balance and coordination. It mainly affects children and may cause headache, vomiting, unsteady walking, or other neurologic problems.

\medterm{Medulloepithelioma} A medulloepithelioma is a rare embryonal tumor arising most often from the nonpigmented ciliary epithelium of the eye in children. It may cause visual symptoms, neovascular glaucoma, or a visible mass and is treated according to its extent and pathology.

\medterm{Mefenamic Acid} A nonsteroidal anti-inflammatory medicine used for short-term pain, including menstrual pain. It can irritate or bleed from the stomach, affect kidney function, and increase cardiovascular risk, especially with prolonged use or in people with risk factors.

\medterm{Mefloquine} Mefloquine is an antimalarial drug used for prevention and treatment of susceptible malaria. It has a long half-life but can cause neuropsychiatric, vestibular, and cardiac adverse effects, so patient-specific contraindications and warnings are important.

\medterm{Megabecquerel} A megabecquerel (MBq) is a unit of radioactivity equal to one million nuclear disintegrations per second. It is used to state the activity of radiopharmaceuticals and other radioactive materials.

\medterm{Megacalycosis} Megacalycosis is congenital, nonobstructive dilatation of the renal calyces, usually with otherwise preserved renal architecture and function. It is often incidental but must be distinguished from obstructive hydronephrosis.

\medterm{Megacolon} Abnormal enlargement and dilation of the colon, which may be congenital, acquired, toxic, or inflammatory.

\medterm{Megakaryocyte} A very large cell in the bone marrow that makes platelets, the small blood parts that help stop bleeding. It grows long extensions into nearby blood vessels, where pieces break off and enter the bloodstream as platelets.

\medterm{Megakaryocytopoiesis} Megakaryocytopoiesis is the formation and maturation of megakaryocytes in the bone marrow, followed by platelet production. It is regulated in large part by thrombopoietin and other hematopoietic signals.

\medterm{Megaloblast} An unusually large, immature cell in bone marrow that should develop into a red blood cell. It usually reflects faulty DNA production, commonly caused by vitamin B12 or folate deficiency.

\medterm{Megaloblastic Anaemia} Macrocytic anemia caused by defective DNA synthesis, most often from vitamin B12 or folate deficiency.

\medterm{Megaloblastic Anemia} Anemia caused by impaired DNA synthesis, usually related to vitamin B12 or folate deficiency, with abnormally large red blood cells.

\medterm{Megalocephaly} Abnormally large head size, which may be familial, developmental, metabolic, or associated with brain disease.

\medterm{Megamonas Funiformis} Megamonas funiformis is an anaerobic bacterium of the genus Megamonas, commonly associated with the intestinal microbiota. Its detection in a clinical specimen requires interpretation alongside symptoms, specimen quality, and other microbiologic findings.

\medterm{Megamonas Hypermegale} Megamonas hypermegale is an anaerobic bacterium in the genus Megamonas, primarily associated with the gastrointestinal microbiota. Its presence in a specimen is interpreted in relation to the site, clinical symptoms, and other microbiologic results.

\medterm{Megasphaera} Megasphaera is a genus of anaerobic bacteria found in the oral cavity and gastrointestinal and genital microbiota. Species identification may be relevant in microbiome studies or when recovered from a clinical specimen.

\medterm{Megaureter} An abnormally widened ureter, the tube that carries urine from the kidney to the bladder. It may result from a blockage, urine flowing backward, or abnormal ureter development and can cause infection, kidney swelling, or reduced kidney function.

\medterm{Megestrol} Megestrol is a synthetic progestin used in some settings to treat hormone-responsive cancers and, in certain formulations, anorexia or severe weight loss. It can cause glucocorticoid-like effects, thromboembolic risk, and adrenal suppression.

\medterm{Megestrol Acetate} Megestrol acetate is the acetate salt of the progestin megestrol. Depending on formulation and indication, it is used for advanced endometrial or breast cancer and for appetite stimulation in selected patients with severe weight loss.

\medterm{Meglitinides} Short-acting diabetes medicines that encourage the pancreas to release insulin around mealtimes. They lower blood sugar but can cause low blood sugar and weight gain, so doses are usually linked to meals and skipped when a meal is skipped.

\medterm{Megrim} An older term for migraine, a recurrent headache disorder often associated with nausea and sensitivity to light or sound.

\medterm{Meibomian Glands} Sebaceous glands in the eyelids that produce the oily layer of the tear film and help prevent tear evaporation.

\medterm{Meibomianitis} Meibomianitis is inflammation or blockage of oil-producing eyelid glands, causing irritated eyelid margins, dry eye, crusting, or recurrent styes.

\medterm{Meige Syndrome} A movement disorder causing repeated involuntary eyelid squeezing and muscle spasms around the mouth, jaw, or tongue. Symptoms may worsen with stress, speaking, chewing, or touch.

\medterm{Meigs Syndrome} A rare combination of a usually benign ovarian fibroma, fluid in the abdomen, and fluid around the lungs. The fluid often improves after the ovarian tumor is removed, but cancer and other causes must be excluded.

\medterm{Meigs' Syndrome} A rare combination of a usually benign ovarian fibroma, fluid in the abdomen, and fluid around the lungs. The fluid often improves after the ovarian tumor is removed, but cancer and other causes must be excluded.

\medterm{Meiosis} A specialized cell division that produces haploid eggs or sperm from a diploid precursor cell and creates
genetic variation through recombination and chromosome assortment.

\medterm{Meiotic Recombination} Meiotic recombination is exchange of DNA between homologous chromosomes during meiosis, usually through crossing over. It creates new combinations of alleles and contributes to genetic variation.

\medterm{Meissner's Corpuscles} Small touch sensors in hairless skin, especially the fingertips, palms, and soles. They help the nervous system detect gentle contact, texture, and movement across the skin.

\medterm{MEK Inhibitor} A targeted medicine that blocks MEK proteins in the MAPK signaling pathway, which can slow growth of some cancers with pathway changes. Side effects vary and may include rash, diarrhea, heart effects, and eye problems.

\medterm{Melancholia} An older term for severe depression, often describing pervasive sadness, loss of pleasure, psychomotor change, and biological symptoms.

\medterm{Melancholic Features} A pattern of severe depression marked by almost complete loss of pleasure, very low mood, major changes in movement or sleep, poor appetite or weight loss, excessive guilt, and symptoms that are often worse in the morning.

\medterm{Melanin} Melanin is a pigment produced by melanocytes that gives skin, hair, and eyes their colour and absorbs ultraviolet radiation; its amount and distribution influence pigmentation disorders and melanoma biology.

\medterm{Melanoblast} A precursor cell that develops into a melanocyte and produces pigment after migrating into the skin and other tissues.

\medterm{Melanocyte} A melanocyte is a pigment-producing cell derived from the neural crest. It makes melanin in specialized organelles called melanosomes and transfers the pigment to neighboring keratinocytes, helping protect skin from ultraviolet radiation.

\medterm{Melanocytes} Pigment-producing cells in the skin and hair follicles. They make melanin, the pigment that gives skin and hair much of their color and helps protect skin cells from ultraviolet radiation.

\medterm{Melanocytic Nevi} Common moles formed by collections of pigment-producing melanocytes. Most are harmless, but a mole that changes in size, shape, color, or symptoms should be assessed for possible melanoma.

\medterm{Melanocytic Nevus} A melanocytic nevus is a usually benign, localized proliferation of melanocytes that appears as a macule, papule, or nodule. Changes in size, color, border, or symptoms may require assessment for atypia or melanoma.

\medterm{Melanocytic Tumor Of The Central Nervous System} A tumor showing melanocytic differentiation within the CNS or meninges, ranging from benign melanocytoma to malignant melanoma. Evaluation includes imaging, histology, and molecular assessment to determine behavior.

\medterm{Melanoderma} Increased pigmentation or darkening of the skin caused by increased melanin production or an increased number of melanocytes.

\medterm{Melanoma} Melanoma is a malignant tumor of melanocytes, most often arising in skin but also in eye or mucosal sites; early detection matters because it can invade and metastasize.

\synonyms
Cutaneous Melanoma; Malignant Melanoma

\medterm{Melanoma Of The Eye} A cancer arising from pigment-producing cells in the eye, most often the uveal tract. It may cause blurred vision, flashes, a growing dark spot, or no symptoms at first and requires specialist eye assessment.

\medterm{Melanoma Surveillance} Ongoing examination of the skin and, when appropriate, the eyes or other previously affected sites after melanoma or in people at increased risk. It may include self-checks, clinician examination, dermoscopy, photographs, and evaluation of new or changing lesions.

\medterm{Melanoma Treatment} Treatment depends on the tumor’s depth, spread, molecular findings, overall health, and previous treatment. Localized melanoma is usually removed surgically; selected cases need lymph-node assessment or additional medicine. Disease that has spread may require immune-based or targeted treatment, surgery, radiation, or combinations.

\medterm{Melanophore} A melanophore is a pigment-containing cell, especially one in lower vertebrates, whose melanin granules can disperse or aggregate to alter visible coloration. In mammals, melanocytes are the principal melanin-producing cells.

\medterm{Melanosis} Abnormal accumulation or increased production of melanin in skin, mucosa, an organ, or other tissue.

\medterm{Melanosis Coli} Melanosis coli is dark pigmentation of the colon lining caused by lipofuscin accumulation, classically associated with prolonged stimulant-laxative use. It is usually benign and may fade after stopping the trigger, but the underlying constipation and medication use should be reviewed.

\medterm{Melanosome} A melanosome is a specialized organelle in a melanocyte in which melanin is synthesized, stored, and transported. Melanosomes are transferred to keratinocytes and influence skin, hair, and eye pigmentation.

\medterm{MELAS Syndrome} A mitochondrial disorder involving encephalopathy, lactic acidosis, and stroke-like episodes, often beginning before adulthood and caused by pathogenic mitochondrial DNA variants.

\medterm{Melasma} Melasma is acquired, usually symmetric brown or gray facial hyperpigmentation associated with ultraviolet exposure, hormones, pregnancy, and genetic susceptibility; it is benign but often recurrent.

\synonyms
Chloasma; Mask of Pregnancy

\medterm{Melasma} An acquired, chronic pigmentary disorder characterized by symmetric, hyperpigmented,
light-to-dark brown macules and patches on sun-exposed areas of the face (malar,
forehead, upper lip, chin). It primarily affects women of childbearing age and
individuals with darker skin types (Fitzpatrick types III-V).

\synonyms
Facial Melanosis; Mask of Pregnancy

\medterm{Melatonin} A hormone produced mainly by the pineal gland that helps coordinate sleep and wake timing. Darkness usually increases its release, while light suppresses it; supplements can cause drowsiness and medicine interactions.

\medterm{Melatonin Receptor} A cell-surface receptor activated by the hormone melatonin. The main types, MT1 and MT2, help regulate the body's sleep-wake clock and seasonal rhythms and are targets of some sleep medicines.

\medterm{Melatonin Receptor Agonists} Medicines that act like melatonin at receptors involved in the body's sleep-wake clock. They may help some people fall asleep, but choice of treatment depends on the cause of insomnia, other medicines, and health conditions.

\medterm{MELD} Model for End-Stage Liver Disease, a laboratory-based score used to estimate severity of advanced liver disease and help prioritize adults for liver transplantation. Traditional versions use bilirubin, creatinine, and INR; current implementations may include additional variables.

\medterm{MELD Score} A medical score estimating the severity and short-term risk of death from advanced liver disease. It uses laboratory results and helps guide urgency for liver transplantation; the calculation may vary as official versions are updated.

\medterm{Melena} Black, sticky, often foul-smelling stool caused by digested blood, usually from bleeding in the stomach or upper intestine. It can signal serious bleeding and requires prompt medical evaluation.

\medterm{Melibiose} A disaccharide composed of galactose and glucose, produced by hydrolysis of raffinose and metabolized by some microorganisms.

\medterm{Melio} An older shortened form related to melioidosis, the infection caused by Burkholderia pseudomallei.

\medterm{Melioidosis} An infectious disease caused by Burkholderia pseudomallei, usually acquired from
contaminated soil or water through inoculation or inhalation. It may cause pneumonia,
skin infection, abscesses, sepsis, or chronic disease and is more severe in people with
diabetes or other risk factors.

\medterm{Melitensis} A species designation, especially Brucella melitensis, a bacterium that causes brucellosis and is associated with goats and sheep.

\medterm{Melitten} Melitten, more commonly spelled melittin, is the principal peptide toxin in honeybee venom. It disrupts cell membranes and contributes to pain, inflammation, hemolysis, and other effects of envenomation.

\medterm{Melorheostosis} A rare sclerosing bone dysplasia producing irregular, flowing cortical hyperostosis on
imaging. It may cause pain, stiffness, limb deformity, restricted movement,
contractures, or soft-tissue fibrosis and can affect one or several contiguous bones.

\medterm{Meloxicam} Meloxicam is a nonsteroidal anti-inflammatory drug that inhibits cyclooxygenase enzymes and reduces prostaglandin synthesis. It is used for pain and inflammation, but can cause gastrointestinal, renal, cardiovascular, and hypersensitivity adverse effects.

\medterm{Melphalan} A chemotherapy medicine that damages cancer cells by interfering with their DNA. It is used for selected blood cancers and other cancers, but can also suppress bone marrow and increase infection or bleeding risk.

\medterm{Memantine} An NMDA receptor antagonist (uncompetitive, low-to-moderate affinity) used for moderate
to severe Alzheimer disease, often added to an acetylcholinesterase inhibitor. Blocks
excessive glutamate-mediated excitotoxicity while preserving normal synaptic
transmission. Once or twice daily oral.

\medterm{Memantine Hydrochloride} Memantine hydrochloride is an NMDA-receptor antagonist used to help manage moderate-to-severe Alzheimer disease. It reduces pathologic glutamate-mediated excitotoxicity; common adverse effects include dizziness, headache, constipation, and confusion.

\medterm{Membrane} A thin sheet of tissue or material that surrounds a cell, organ, or body cavity. Membranes can protect structures, separate spaces, allow selected substances to pass, and produce or absorb fluids.

\medterm{Membrane Channel} A membrane channel is a protein-forming pore across a cell membrane that permits selected ions or molecules to pass down an electrochemical or concentration gradient. Many channels are gated by voltage, ligands, or mechanical forces.

\medterm{Membrane Fusion} Membrane fusion is merging of two lipid bilayers into one continuous membrane. It is required for processes such as synaptic vesicle release, enveloped-virus entry, fertilization, and intracellular vesicle trafficking.

\medterm{Membrane Lipids} Fat-like molecules that form the main structure of cell membranes. They create a flexible barrier, help cells communicate, and provide attachment sites for proteins; abnormal lipid composition can affect signaling and cell function.

\medterm{Membrane Microdomain} A membrane microdomain is a small, specialized region of a cell membrane with a distinct composition of lipids and proteins. It can organize receptors, signaling molecules, transporters, and cytoskeletal interactions.

\medterm{Membrane Part} A membrane part is a structural or molecular component of a biological membrane, such as a lipid, protein, carbohydrate, channel, or receptor. These components determine membrane integrity, transport, signaling, and cell recognition.

\medterm{Membrane Potential} The electrical potential difference across a cell membrane, approximately -70 mV in most
neurons (resting membrane potential). Established by the unequal distribution of ions
(primarily K+, Na+, Cl-) across the membrane and selective permeability.

\medterm{Membrane Protein} A protein attached to or embedded in a cell membrane. It may transport substances, receive signals, act as an enzyme, join cells, anchor the cytoskeleton, or help the immune system identify cells.

\medterm{Membrane Raft} A membrane raft is a cholesterol- and sphingolipid-enriched membrane microdomain that helps cluster selected proteins and signaling molecules. Raft organization is dynamic and may change during receptor signaling or membrane trafficking.

\medterm{Membranoproliferative Glomerulonephritis} A pattern of kidney-filter damage that can cause blood and protein in the urine, swelling, and reduced kidney function. Causes include immune disorders, infections, and some blood cancers.

\medterm{Membranous} Relating to, consisting of, or resembling a membrane. In medicine, the term may describe a thin tissue layer or a disease pattern involving such a layer.

\medterm{Membranous Nephropathy} A kidney disease in which the filtering units are injured by immune proteins deposited along their delicate filter. It often causes large amounts of protein in the urine and swelling and may be linked to autoimmune disease, infection, medicines, or cancer.

\medterm{Memory} The ability to learn, store, and recall information. Memory depends on several connected areas of the brain and can be affected by sleep problems, stress, medicines, injury, illness, aging, or brain disorders.

\medterm{Memory Cell} A long-lasting B cell or T cell formed after an infection or vaccination. If the same germ is encountered again, memory cells respond faster and more strongly, helping prevent illness or reduce its severity.

\medterm{Memory Disorder} Impairment of encoding, storage, retrieval, or recognition of information that exceeds normal variation and may result from neurodegeneration, injury, vascular disease, medication, or other causes.

\medterm{Memory Impairment} Memory impairment is reduced ability to learn, store, retrieve, or recognize information. It may be transient or persistent and can result from normal aging, medications, sleep problems, neurologic disease, psychiatric illness, or other medical conditions.

\medterm{Memory Loss} Reduced ability to recall or learn information. It may result from normal aging, medical illness, medicines,
brain injury, or a neurocognitive disorder.

\medterm{Memory Loss} Reduced ability to recall past information, form new memories, or both. Causes include normal aging, medications, head injury, seizures, stroke, infection, substance use, neurocognitive disorders, and psychological conditions; sudden memory loss requires urgent assessment.

\medterm{Memory Span} The number of items, such as words or digits, that a person can briefly retain and recall, usually in the original order; it is a measure of short-term or working memory.

\medterm{Memory T Cells And B Cells} Long-lived immune cells that remember earlier exposures and respond faster when they meet the same target again. They support lasting protection after infection or vaccination, though protection may fade.

\synonyms
Memory Lymphocytes

\medterm{MEN 2} Abbreviation; see \textit{Multiple Endocrine Neoplasia Type 2}.

\medterm{Menarche} The first menstrual period, marking the onset of menstruation. Its timing varies with genetics, nutrition, health, and pubertal development; markedly early or delayed menarche may require assessment.

\medterm{Mendelian Inheritance} Inheritance of traits determined by single genes and transmitted according to dominant, recessive, or sex-linked patterns.

\medterm{Mendelism} The principles of inheritance based on Gregor Mendel's work, including segregation of alleles and independent assortment of genes; Mendelian traits may follow autosomal or X-linked dominant or recessive patterns.

\medterm{Mendelson Syndrome} A chemical pneumonitis caused by inhalation of acidic stomach contents, classically during anesthesia
or childbirth. It can cause cough, breathing difficulty, and impaired oxygenation.

\medterm{Meniere Disease} An inner-ear disorder causing repeated attacks of spinning dizziness, usually with fluctuating hearing loss, ringing in the ear, and a feeling of pressure or fullness. Attacks often last from 20 minutes to several hours.

\synonyms
Endolymphatic Hydrops; Meniere's Syndrome

\medterm{Meniere's Disease} Alternate spelling; see \textit{Meniere Disease}.

\medterm{Meniere's Disease} An inner-ear disorder causing episodes of vertigo, fluctuating hearing loss, tinnitus, and a feeling of fullness in one ear. Treatment aims to reduce attacks and protect hearing.

\medterm{Meningeal Carcinomatosis} Spread of cancer cells to the thin coverings and fluid around the brain and spinal cord. It may cause headache, nausea, weakness, seizures, vision changes, or problems with movement and needs specialist evaluation.

\medterm{Meningeal Leukemia} Meningeal leukemia is leukemia involving the membranes or fluid around the brain and spinal cord, potentially causing headache, weakness, or cranial-nerve problems.

\medterm{Meningeal Lymphoma} Meningeal lymphoma is lymphoma affecting the membranes or fluid spaces around the brain and spinal cord.

\medterm{Meningeal Neoplasm} A meningeal neoplasm is an abnormal growth arising in or near the membranes surrounding the brain and spinal cord.

\medterm{Meninges} Three protective layers surrounding the brain and spinal cord. Fluid between two of the layers cushions the nervous system; bleeding or infection in these layers can be dangerous.

\medterm{Meningioma} A usually slow-growing tumor arising from cells associated with the meninges. Most are benign, but some
are atypical or malignant and can cause symptoms by compressing the brain, spinal cord, or
cranial nerves.

\medterm{Meningism} Meningism is a syndrome of neck stiffness, headache, and sometimes photophobia or vomiting that resembles meningeal irritation without necessarily indicating meningitis. Serious causes, including infection or subarachnoid hemorrhage, require prompt evaluation.

\medterm{Meningismus} Neck stiffness and other signs of irritation around the brain and spinal cord without proven inflammation of those coverings. It can occur with fever, severe headache, dehydration, bleeding, or other illnesses and needs evaluation.

\medterm{Meningitis} Inflammation of the membranes around the brain and spinal cord, usually caused by infection but sometimes by medicines or immune disease. Fever, severe headache, stiff neck, light sensitivity, or confusion requires urgent assessment.

\medterm{Meningocele} A birth defect in which the protective coverings around the brain or spinal cord bulge through an opening in the skull or spine, forming a fluid-filled sac. It may cause nerve problems and needs specialist assessment.

\medterm{Meningocele Repair} Surgery to close an opening in the spine and return the protective coverings of the spinal cord to their proper position. It aims to prevent injury or infection; nerve, bladder, bowel, movement, or fluid problems may remain or develop.

\medterm{Meningococcal Disease} A serious infection caused by Neisseria meningitidis that can produce meningitis or rapidly spreading bloodstream infection. Fever, severe headache, stiff neck, confusion, vomiting, or a purple rash requires emergency care.

\medterm{Meningococcal Infection} An infection caused by the bacterium Neisseria meningitidis. It may remain mild or cause life-threatening meningitis or bloodstream infection, with fever, headache, stiff neck, vomiting, confusion, or a purple rash developing quickly.

\medterm{Meningococcal Meningitis} Acute bacterial meningitis caused by Neisseria meningitidis, often accompanied by bloodstream infection; it is a medical emergency requiring immediate antibiotics and public-health measures.

\medterm{Meningococcemia} Meningococcemia is bloodstream infection with Neisseria meningitidis, causing fever, rash, shock, clotting failure, and potentially fatal sepsis.

\medterm{Meningococcus} Neisseria meningitidis, a bacterium that can cause meningitis or rapidly spreading bloodstream infection. It spreads through respiratory secretions and may cause fever, headache, rash, shock, or death without urgent treatment.

\medterm{Meningoencephalitis} Inflammation of both the meninges surrounding the brain and the brain tissue itself, usually caused by
infection or autoimmune disease. Fever, headache, neck stiffness, confusion, seizures, or reduced
consciousness requires urgent medical assessment.

\medterm{Meningomyelocele} A birth defect in which the protective coverings and spinal-cord tissue bulge through an opening in the spine. It can cause weakness or paralysis of the legs and bladder or bowel problems.

\medterm{Meniscal Allograft Transplantation} Surgery placing a donated meniscus into a knee after substantial loss of the person's own meniscus. It may reduce pain in carefully selected younger patients, but healing is variable and infection, stiffness, graft failure, or later arthritis can occur.

\medterm{Meniscal Injuries} Tears or other damage to the knee menisci, which are crescent-shaped pieces of cartilage
that help distribute load and stabilize the knee. Symptoms may include joint-line pain,
swelling, clicking, or restricted motion. Diagnosis is based on examination and, when needed, imaging.

\medterm{Meniscal Tear} A split in one of the two rubbery cartilage pads that cushion the knee. It may follow a twisting injury or gradual wear and can cause pain along the joint line, swelling, clicking, catching, locking, or a feeling that the knee gives way.

\medterm{Meniscectomy} An operation that removes part or all of a damaged meniscus, the rubbery pad that cushions the knee. Surgeons usually remove only the torn portion when possible, because preserving meniscus helps protect the joint from later wear.

\medterm{Meniscus} A crescent-shaped fibrocartilage structure in the knee, or a curved fluid surface in a container.

\medterm{Meniscus Tear} Alternate name; see \textit{Torn Meniscus}.

\medterm{Menkes Disease} A rare inherited disorder in which the body cannot transport copper properly. Low copper in the brain and other tissues can cause poor growth, developmental decline, seizures, weak muscles, and sparse, brittle, or unusually twisted hair.

\medterm{Menkes Kinky Hair Syndrome} An X-linked disorder of copper transport caused by ATP7A dysfunction. It produces
progressive neurodegeneration, seizures, developmental delay, hypopigmented or brittle
kinky hair, connective-tissue abnormalities, and failure to thrive.

\medterm{Menogaril} Menogaril is a semisynthetic anthracycline derivative of nogalamycin. It intercalates into DNA and inhibits topoisomerase II, thereby impairing DNA replication and repair; it was investigated as an antineoplastic agent.

\medterm{Menometrorrhagia} Excessive uterine bleeding occurring both at expected menstrual times and at irregular intervals; the older term corresponds broadly to heavy and irregular menstrual bleeding and warrants evaluation for underlying causes.

\medterm{Menopausal} Relating to menopause or the years around the permanent end of menstrual periods. Hormone changes may cause hot flashes, sleep problems, mood changes, vaginal dryness, and increased bone-loss risk.

\medterm{Menopausal Hormone Therapy} Treatment with estrogen, with progestogen when needed to protect the uterus, for troublesome hot flashes, night sweats, or vaginal symptoms during menopause. It may affect clot, stroke, breast-cancer, and other risks, so treatment is individualized.

\medterm{Menopausal Hot Flashes} Alternate name; see \textit{Hot Flashes}.

\medterm{Menopause} The permanent end of menstrual periods because of loss of ovarian activity, usually diagnosed after 12 consecutive months without menstruation. Hormonal changes may cause hot flashes, night sweats, vaginal symptoms, sleep problems, or other symptoms.

\medterm{Menopause And Perimenopause} Perimenopause is the transition before menopause, when menstrual periods and hormone levels become irregular. Menopause is reached after 12 consecutive months without a period when no other cause explains it.

\medterm{Menopause Status} Menopause status describes whether a person is premenopausal, perimenopausal, or postmenopausal. Natural menopause is clinically defined after 12 consecutive months without menstruation when no other cause explains the amenorrhea.

\medterm{Menopause Syndrome} A nontechnical term for the group of symptoms that may occur during the menopausal transition, such as hot flashes, night sweats, sleep problems, mood changes, vaginal dryness, and menstrual irregularity. Menopause itself is a normal life stage, not a disease.

\medterm{Menorrhagia} Heavy or prolonged menstrual bleeding that interferes with daily life. Causes include fibroids, adenomyosis, polyps, bleeding disorders, thyroid disease, or irregular ovulation; ongoing blood loss may cause iron-deficiency anemia.

\medterm{Menotropin} Menotropin is a urinary-derived preparation containing follicle-stimulating and luteinizing hormone activity, used to stimulate follicular development or spermatogenesis; ovarian hyperstimulation and multiple pregnancy are important risks.

\medterm{Menses Duration} Menses duration is the number of days of menstrual bleeding in a cycle, counted from the first day of bleeding until bleeding stops. Persistently prolonged, heavy, or irregular bleeding warrants clinical assessment.

\medterm{Menstrual} Relating to menstruation or the menstrual cycle, including the monthly shedding of the uterine lining and the bleeding that follows. Cycle length and bleeding vary, but very heavy, painful, irregular, or unexpected bleeding may need evaluation.

\medterm{Menstrual Cramps} Cramping pain in the lower abdomen during or around menstruation, medically called dysmenorrhea. Pain may radiate to
the back or thighs and may occur with nausea, diarrhea, headache, or fatigue. Severe, worsening, or newly
developed pain requires evaluation for causes such as endometriosis or fibroids.

\synonyms
Dysmenorrhea; Period Pain; Painful Menstruation

\medterm{Menstrual Cycle} The recurring hormonal and uterine changes that prepare the body for possible pregnancy. Cycle length and
bleeding pattern vary among individuals and across life stages.

\medterm{Menstrual Disorders} Conditions involving abnormal menstrual timing, regularity, duration, amount of bleeding, or pain. They include absent or infrequent periods, irregular cycles, unusually heavy or prolonged bleeding, painful periods, and bleeding between periods or after sex. Causes include pregnancy, hormonal or ovulatory disorders, fibroids, endometriosis, thyroid disease, medicines, structural abnormalities, and perimenopause; persistent, heavy, or unexpected bleeding requires assessment.

\medterm{Menstruation} The periodic shedding of the endometrial lining as part of the menstrual cycle when pregnancy does not occur.
The timing, duration, and amount of bleeding vary with age, hormonal state, medicines,
and health conditions.

\medterm{Menstruation Disturbance} Alternate name; see \textit{Menstrual Disorders}.

\medterm{Mental} Relating to thinking, memory, emotion, behavior, or other functions of the mind. It does not by itself mean that a condition is psychiatric.

\medterm{Mental Age} An estimate of cognitive functioning expressed as the age level at which a standardized test score is typical; it is not the same as chronological age.

\medterm{Mental Disorder} Alternate singular form; see \textit{Mental Disorders}.

\medterm{Mental Disorders} Health conditions involving a clinically significant disturbance in cognition, emotion regulation, or behavior that causes distress, impairment, or both. Examples include mood, anxiety, psychotic, neurodevelopmental, trauma-related, personality, and substance use disorders. A mental-health concern or symptom does not necessarily constitute a mental disorder.

\medterm{Mental Disorders In Children} Mental disorders in children involve persistent, developmentally atypical disturbances in mood, thinking, behavior, or functioning that cause clinically significant distress or impairment. Symptoms must be interpreted in relation to the child’s developmental stage and distinguished from normal variation, temporary reactions, and age-appropriate behavior.

\medterm{Mental Fatigue} Mental fatigue is a subjective feeling of reduced mental energy, concentration, or cognitive endurance after sustained effort or illness. Persistent or severe symptoms may accompany sleep disorders, depression, neurologic disease, medication effects, or other medical conditions.

\medterm{Mental Health} Mental health encompasses emotional, psychological, and social well-being, affecting how individuals think, feel, behave, and handle stress, relationships, and decision-making.

\medterm{Mental Health Crisis} A period of severe emotional, behavioral, or psychological distress in which a person’s usual ability to cope or maintain safety is overwhelmed. It may involve suicidal thoughts, severe agitation, psychosis, substance intoxication, or inability to care for basic needs; it may or may not meet criteria for a psychiatric emergency.

\medterm{Mental Illness} A broad, commonly used term for mental disorders and related conditions affecting mood, thinking, behavior, perception, or daily functioning. See \textit{Mental Disorders}.

\medterm{Mental Illness In Children} Broad, commonly used term; see \textit{Mental Disorders In Children}.

\medterm{Mental Process} A mental process is a cognitive or affective activity of the mind, such as perception, attention, memory, reasoning, language, emotion, or decision-making.

\medterm{Mental Recall} Mental recall is retrieval of previously learned or stored information from memory. It may be assessed by asking a person to remember facts, events, words, or experiences after a delay.

\medterm{Mental Status Examination} A structured part of a clinical interview that assesses a person's appearance, behavior, speech, mood, thinking, perceptions, attention, memory, awareness, judgment, and understanding of their condition.

\medterm{Mental Status Testing} A structured assessment of appearance, behavior, alertness, orientation, attention, memory, language, thought, perception, mood, insight, and judgment.

\medterm{Mentation} The process of thinking, perceiving, reasoning, and other mental activity.

\medterm{Menthol} A mint-derived or synthetic compound producing a cooling sensation and used in some topical and respiratory preparations.

\medterm{Menthol Poisoning} Poisoning after swallowing, inhaling, or having substantial skin exposure to menthol. It may cause nausea, abdominal pain, dizziness, drowsiness, agitation, seizures, coma, or breathing problems; suspected significant exposure needs urgent poison-control advice.

\medterm{Mepacrine} An older antimalarial and antiprotozoal medicine, also called quinacrine, now used only in selected indications.

\medterm{Meperidine} An opioid analgesic, also called pethidine, whose use is limited by adverse effects and an active neurotoxic metabolite.

\medterm{Meperidine Hydrochloride} Meperidine hydrochloride is the hydrochloride salt of the opioid analgesic meperidine. It produces analgesia through opioid-receptor activity and carries risks of sedation, respiratory depression, seizures, dependence, and serotonin toxicity with interacting drugs.

\medterm{Meperidine Hydrochloride Overdose} Taking too much meperidine hydrochloride, an opioid, can cause extreme sleepiness, pinpoint pupils, slow or stopped breathing, low blood pressure, seizures, or coma. Delayed toxicity and serotonin toxicity are possible, so emergency assessment is required.

\medterm{Mephenesin} An older centrally acting muscle relaxant with sedative effects. It is rarely used now and can impair alertness, coordination, and breathing when combined with other sedatives.

\medterm{Mephentermine} A stimulant medicine formerly used in selected situations to raise dangerously low blood pressure. It increases activity of certain adrenaline-like receptors and can cause fast heartbeat, anxiety, or abnormal heart rhythms.

\medterm{Mephenytoin} Mephenytoin is a hydantoin anticonvulsant formerly used to treat seizures. It is metabolized partly to ethotoin and has limited modern use because of potentially serious adverse effects, including hypersensitivity and blood dyscrasias.

\medterm{Mepivacaine} A local anesthetic medicine that temporarily blocks pain signals from nerves. It is injected for numbing skin, dental areas, individual nerves, or larger body regions during procedures.

\medterm{Meprobamate} An older medicine that reduces anxiety and causes sleepiness. It can cause dependence, dangerous breathing and nervous-system depression, and overdose, especially with alcohol or other sedatives, so it is rarely used today.

\medterm{Meprobamate Overdose} Taking too much meprobamate can cause severe drowsiness, confusion, poor coordination, low blood pressure, slowed breathing, coma, or death. Alcohol and other sedatives increase the danger; suspected overdose requires emergency evaluation.

\medterm{Mepyramine} An antihistamine that blocks H1 receptors and may cause sedation; it is also called pyrilamine.

\medterm{Meralgia Paresthetica} Burning, tingling, or numbness over the outer thigh caused by compression of the lateral femoral cutaneous nerve.

\synonyms
Burning Thigh Pain; Bernhardt-Roth Syndrome

\medterm{Merbromin Poisoning} Poisoning from merbromin, an older mercury-containing antiseptic. Swallowing or repeated exposure may injure the mouth, stomach, kidneys, or nervous system; the amount and route determine risk and require poison-control or medical assessment.

\medterm{Mercaptopurine} A thiopurine antimetabolite used in selected leukemias and inflammatory disorders, requiring monitoring for marrow and liver toxicity.

\medterm{Mercurialism} Long-term poisoning caused by mercury exposure. It may damage the nervous system, kidneys, digestive tract, mouth, or mental health, with effects depending on the mercury form, dose, and exposure duration.

\medterm{Mercuric Chloride Poisoning} Severe poisoning from mercuric chloride, a corrosive inorganic mercury salt. It can burn the mouth and digestive tract and cause vomiting, bloody diarrhea, shock, kidney failure, or neurologic injury; emergency treatment is required.

\medterm{Mercuric Oxide Poisoning} Poisoning from mercuric oxide, an inorganic mercury compound that can irritate or burn tissues and damage the kidneys and nervous system. Swallowing or inhalation requires urgent poison-control or emergency assessment.

\medterm{Mercury} A metal used in some industrial and laboratory products. Exposure can harm the nervous system, kidneys, lungs, or digestive tract, so it must be handled and disposed of safely.

\medterm{Mercury Poisoning} Illness caused by exposure to elemental, inorganic, or organic mercury. Effects depend on the form and dose and may include nerve, kidney, lung, or digestive injury; urgent advice is needed after significant exposure.

\medterm{Meridian} A reference line or plane; in traditional acupuncture it denotes a proposed pathway for qi, while in anatomy it may describe a longitudinal line.

\medterm{Merkel Cell Carcinoma} An aggressive neuroendocrine skin cancer arising from Merkel cells. It presents as a
rapidly growing, painless, red-violet nodule on sun-exposed skin of elderly patients.
Approximately 80 percent are associated with Merkel cell polyomavirus. Risk factors
include immunosuppression, UV exposure, and fair skin.

\synonyms
MCC; Neuroendocrine Carcinoma of the Skin

\medterm{Merkel Cells} Specialized touch-sensing cells in the deepest layer of the skin that help detect light touch and sustained pressure.

\medterm{Mermaid Syndrome} An extremely rare congenital anomaly in which the legs are fused to varying degrees, also called sirenomelia.

\medterm{Meropenem} Meropenem is a broad-spectrum carbapenem antibiotic that inhibits bacterial cell-wall synthesis and treats severe susceptible infections, including meningitis; allergy, gastrointestinal effects, seizures, and resistance are possible.

\medterm{Meropenem-Vaborbactam} An intravenous antibiotic combination used for some serious infections caused by bacteria resistant to many other antibiotics. The choice depends on laboratory susceptibility testing, infection site, kidney function, and local guidance.

\medterm{MERRF Syndrome} A rare inherited disorder affecting how cells produce energy. It commonly causes muscle jerks, seizures, balance problems, and muscle weakness, and may also affect hearing, vision, or other organs.

\medterm{MERS-CoV} Middle East respiratory syndrome coronavirus, the virus that causes MERS. Infection can cause fever and respiratory illness ranging from mild disease to severe pneumonia, especially in people with other health conditions.

\medterm{Merthiolate Poisoning} Poisoning from merthiolate, commonly referring to the mercury-containing antiseptic thimerosal. Small skin exposures usually cause local irritation, but swallowing or repeated exposure may cause corrosive, kidney, or neurologic injury and needs advice.

\medterm{Merycism} Rumination, the repeated regurgitation and rechewing of food after eating.

\medterm{Mesalamine} Mesalamine, or 5-aminosalicylic acid, is an anti-inflammatory drug used to induce and maintain remission in ulcerative colitis and some other inflammatory bowel conditions. It acts mainly topically in the intestinal mucosa and can rarely affect renal function.

\medterm{Mesangiolysis} Mesangiolysis is dissolution or disruption of the mesangial tissue within a renal glomerulus. It may occur in severe endothelial injury and can be associated with microvascular lesions and glomerular capillary aneurysmal changes.

\medterm{Mesangium} The mesangium is the supportive tissue between glomerular capillary loops, composed of mesangial cells and extracellular matrix. Mesangial cells provide structural support, regulate capillary flow, and clear certain substances; abnormal expansion occurs in several glomerular diseases.

\medterm{Mescaline} A naturally occurring psychedelic substance found in certain cacti that alters perception and may cause hallucinations and autonomic effects.

\medterm{Mesencephalon} The mesencephalon, or midbrain, is the part of the brainstem between the pons and diencephalon. It contains pathways and nuclei involved in eye movements, motor control, arousal, and auditory and visual reflexes.

\medterm{Mesenchymal Cell} A connective-tissue precursor that can develop into bone, cartilage, fat, muscle, or tendon cells.

\medterm{Mesenchymal Chondrosarcoma} Mesenchymal chondrosarcoma is a rare aggressive cancer containing cartilage-forming tissue, often arising in bone or soft tissue.

\medterm{Mesenchymal Tumor Of The Central Nervous System} A rare tumor arising from connective or supporting tissue in or around the brain or spinal cord rather than from nerve cells. Its behavior is determined by examination of tissue, imaging, and sometimes genetic testing.

\medterm{Mesenchyme} Embryonic connective tissue that gives rise to much of the skeleton, connective tissue, blood vessels, and smooth muscle.

\medterm{Mesenchymoma} A mesenchymoma is a rare tumor containing more than one type of mesenchymal tissue, sometimes including cartilage, bone, fat, or vascular elements. Behavior depends on the site, age, histology, and presence of malignant components.

\medterm{Mesenteric} Relating to the mesentery, a fold of tissue supporting the intestines. It contains blood vessels, nerves, lymph vessels, and fat that help supply, drain, anchor, and protect the intestines.

\medterm{Mesenteric Angiography} Contrast imaging of the arteries supplying the intestines, used to evaluate mesenteric ischemia, bleeding, aneurysm, stenosis, occlusion, or vascular anatomy.

\medterm{Mesenteric Artery} A mesenteric artery supplies blood to portions of the intestine and associated abdominal organs. The superior and inferior mesenteric arteries are the major vessels; obstruction or embolism can cause acute mesenteric ischemia.

\medterm{Mesenteric Artery Ischemia} Inadequate blood flow through the mesenteric arteries, usually caused by embolism,
thrombosis, or severe low-flow states. Acute disease presents with severe abdominal pain
out of proportion to examination and can progress to bowel infarction.

\medterm{Mesenteric Cyst} A rare, benign, fluid-filled cystic mass that develops within the folds of the abdominal mesentery, most frequently originating from sequestered embryonic lymphatic tissue in the small bowel mesentery (especially the ileum). On physical examination, it demonstrates the classic Tillaux's sign: mobility in a transverse direction perpendicular to the mesenteric root, with limited or no movement vertically. While often asymptomatic and found incidentally, large cysts can cause abdominal distension, chronic pain, or acute surgical emergencies such as cyst torsion, internal hemorrhage, infection, rupture, or mechanical bowel obstruction. Diagnosis is made with abdominal ultrasound, CT, or MRI, and definitive management is complete surgical enucleation or laparoscopic excision to prevent recurrence.

\synonyms
Mesenteric Lymphangioma; Chylolymphatic Cyst; Enteric Mesenteric Cyst

\medterm{Mesenteric Ischemia} Reduced blood flow to the intestines, caused suddenly by a clot or gradually by narrowed arteries. It can cause severe abdominal pain, bowel injury, infection, or death and requires urgent assessment.

\medterm{Mesenteric Lymphadenitis} Inflammation of lymph nodes in the tissue supporting the intestines, often after a viral or bacterial infection. It can cause right-sided abdominal pain, fever, nausea, and tenderness and may resemble appendicitis.
\synonyms
Mesenteric Adenitis

\medterm{Mesenteric Veins} Mesenteric veins drain blood from the small and large intestines and carry it mainly through the superior and inferior mesenteric veins to the portal venous system. Thrombosis or impaired flow can cause intestinal congestion and ischemia.

\medterm{Mesenteric Venous Thrombosis} A blood clot in a mesenteric vein that impairs intestinal venous drainage and can cause abdominal pain, bowel edema, ischemia, infarction, or gastrointestinal bleeding.

\medterm{Mesentery} A fold of the abdominal lining that attaches the intestines to the back wall of the abdomen. It carries blood vessels, nerves, and lymph vessels to the bowel and helps hold the intestines in place.

\medterm{Mesh} A sheet-like implant made from synthetic or biological material that reinforces weakened tissue. It may be used in selected hernia or pelvic-floor operations; possible problems include pain, infection, scarring, movement of the implant, or injury to nearby organs.

\medterm{Mesial} Located toward the body's midline; in dentistry, located toward the front or midline of a tooth. The term helps describe tooth surfaces, positions, cavities, gum disease, and dental restorations.

\medterm{Mesial Surface} The mesial surface is the surface of a tooth or other structure oriented toward the midline. In dentistry, mesial surfaces contact adjacent teeth and can be involved in caries, periodontal disease, or restoration margins.

\medterm{Mesna} A protective medicine given with certain chemotherapy drugs to reduce irritation and bleeding of the bladder. It binds harmful breakdown products in the urine and does not treat the cancer itself.

\medterm{Mesoderm} The middle germ layer in an embryo, which develops into muscles, bones, cartilage, blood
vessels, and the kidneys.

\medterm{Mesonephric} Mesonephric means relating to the mesonephros, a temporary embryonic kidney, or to structures derived from its ducts and tubules. Mesonephric remnants can give rise to uncommon benign or malignant lesions of the genital tract.

\medterm{Mesonephric Adenocarcinoma} A rare HPV-independent cervical adenocarcinoma arising from mesonephric remnants, usually presenting as a cervical mass or abnormal bleeding and requiring histopathologic diagnosis.

\medterm{Mesonephric Tumor} A mesonephric tumor is a rare tumor arising from remnants of embryonic mesonephric ducts, usually in the female reproductive tract.

\medterm{Mesothelial Neoplasm} A mesothelial neoplasm is a growth arising from the lining of body cavities, such as the pleura, peritoneum, or pericardium.

\medterm{Mesothelin} Mesothelin is a cell-surface glycoprotein expressed mainly by mesothelial cells and overexpressed in several cancers, including mesothelioma and some ovarian and pancreatic cancers. It is being studied as a biomarker and therapeutic target.

\medterm{Mesothelioma} A cancer arising from the thin lining of body cavities, most often the pleura around the lungs. It is strongly associated with asbestos exposure and may cause chest pain, breathlessness, or fluid around the lung.

\medterm{Mesothelium} The mesothelium is a single layer of specialized cells lining serous body cavities and covering their organs, including the pleura, peritoneum, and pericardium. It produces lubricating fluid and participates in inflammation and repair.

\medterm{Messenger RNA} A temporary copy of genetic instructions made from DNA. It carries those instructions to cell structures that build proteins and is then broken down; some vaccines use messenger RNA to teach the immune system to recognize a target.

\medterm{Met} Abbreviation for methionine, an essential sulfur-containing amino acid; it is also represented by the one-letter code M and the codons AUG.

\medterm{Metabolic} Relating to metabolism, the chemical processes by which cells produce and use energy and materials.

\medterm{Metabolic Acidosis} A buildup of acid or loss of bicarbonate that makes the blood too acidic. Causes include severe diabetes, poor circulation, kidney failure, starvation, diarrhea, or poisoning and require evaluation of the underlying cause.

\medterm{Metabolic Alkalosis} A condition in which the blood becomes too alkaline, often because of repeated vomiting, fluid loss, or certain medicines. It can cause weakness, cramps, confusion, or abnormal heart rhythms.

\medterm{Metabolic Bone Disease} A group of disorders in which abnormal handling of calcium, phosphate, vitamin D, parathyroid hormone, or other factors weakens or changes bone. Examples include osteomalacia, rickets, and some disorders caused by kidney or endocrine disease.

\medterm{Metabolic Disease} A disorder that changes how the body processes or uses substances such as sugars, fats, proteins, minerals, or hormones. It may be inherited or acquired and can affect energy, growth, blood chemistry, or organ function.

\medterm{Metabolic Disorders} Alternate plural form; see \textit{Metabolic Disease}.

\medterm{Metabolic Dysfunction-Associated Steatohepatitis} Alternate name; see \textit{Fatty Liver Disease}.

\medterm{Metabolic Dysfunction-Associated Steatotic Liver Disease} Alternate name; see \textit{Fatty Liver Disease}.

\medterm{Metabolic Equivalents} Units that compare the energy used during an activity with the energy used while resting. They provide a rough estimate of exercise intensity, but actual energy use varies between people.

\synonyms
METs

\medterm{Metabolic Myopathies} Muscle disorders caused by impaired energy metabolism, often producing exercise intolerance, cramps, weakness, muscle pain, or episodes of rhabdomyolysis.

\medterm{Metabolic Neuropathies} Nerve damage caused by a metabolic or systemic disorder, such as diabetes, kidney failure, vitamin deficiency, or an inherited metabolic disease. Symptoms commonly include numbness, burning pain, tingling, weakness, or problems with balance.

\medterm{Metabolic Rate} The speed at which the body uses energy to keep organs working and support activity. It changes with body size, age, muscle amount, hormones, illness, temperature, food intake, and exercise.

\medterm{Metabolic Response} A metabolic response is a change in energy use, substrate metabolism, or biochemical pathways in response to conditions such as fasting, exercise, infection, trauma, or medication. The response may be adaptive or contribute to illness when excessive or prolonged.

\medterm{Metabolic Syndrome} A combination of increased waist size, high blood pressure, high blood sugar, high triglycerides, or low HDL cholesterol. Together these findings increase the risk of heart disease, stroke, and type 2 diabetes.

\synonyms
Insulin Resistance Syndrome; Syndrome X

\medterm{Metabolic Tumor Volume} The estimated amount of tumor tissue that shows increased activity on a PET scan. It can help measure disease burden and treatment response, especially in lymphoma, but inflammation and scan technique can affect the result.

\medterm{Metabolism} The chemical processes cells use to turn food and medicines into energy and building materials, and to break down substances the body no longer needs.

\medterm{Metabolite} A substance made, changed, or used when the body processes a medicine, nutrient, or naturally occurring molecule. Measuring metabolites in blood or urine can show body function, disease, or drug effects.

\medterm{Metacarpal} Relating to one of the five bones forming the palm between the wrist and fingers. Metacarpal fractures can cause pain, swelling, deformity, reduced grip, and difficulty moving the hand.

\medterm{Metacarpal Bone} A metacarpal bone is one of the five long bones of the hand between the wrist bones and the proximal phalanges. Each has a base, shaft, and head; fractures can affect alignment, grip, and joint function.

\medterm{Metacarpals} The five bones forming the palm, numbered from the thumb to the little finger. Each connects the wrist to a finger and helps with grip and hand movement; breaks can cause pain, swelling, deformity, or reduced function.

\medterm{Metacarpophalangeal Joint} A knuckle joint connecting a palm bone to a finger or thumb bone. It allows the digit to bend, straighten, and move from side to side and can be injured by impact or twisting.

\medterm{Metacarpus} The part of the hand between the wrist and fingers, containing the five metacarpal bones.

\medterm{Metachromatic Leukodystrophy} A rare inherited disorder in which harmful fatty substances build up in nerve cells and damage the protective covering of nerves. It causes worsening movement, thinking, speech, or vision problems, often beginning in childhood.

\medterm{Metagonimus Yokogawai} Metagonimus yokogawai is a small intestinal fluke acquired by eating raw or inadequately cooked freshwater fish. Infection can cause abdominal symptoms and, rarely, ectopic migration of eggs or adult worms.

\medterm{Metal Artifact} Distortion or loss of detail on an X-ray, CT, or MRI image caused by metal such as an implant, clip, or bullet. It can hide nearby findings, so special imaging methods or another test may be needed.

\medterm{Metal Cleaner Poisoning} Poisoning from a metal-cleaning product containing acids, alkalis, solvents, or other chemicals. It may burn the skin, eyes, mouth, or digestive tract, or cause choking and lung injury if aspirated; do not induce vomiting.

\medterm{Metal Fume Fever} A short-lived influenza-like illness caused by inhalation of fumes from heated metals, especially zinc; symptoms include fever, chills, headache, muscle aches, and cough.

\medterm{Metal Polish Poisoning} Poisoning from metal polish, which may contain petroleum solvents, acids, or other chemicals. Swallowing or inhalation can cause vomiting, drowsiness, chemical burns, or aspiration-related lung injury; urgent product-specific advice is needed.

\medterm{Metalloendopeptidase} A metalloendopeptidase is an endopeptidase whose catalytic activity depends on a bound metal ion, commonly zinc. It cleaves peptide bonds within proteins and includes enzymes involved in extracellular matrix remodeling and signaling.

\medterm{Metalloporphyrins} Metalloporphyrins are porphyrin molecules containing a centrally coordinated metal ion. Heme, which contains iron, is a biologically important example and forms the functional group of hemoglobin, myoglobin, and many enzymes.

\medterm{Metalloprotease} A metalloprotease is a proteolytic enzyme that requires a metal ion, commonly zinc, for catalytic activity. Metalloproteases regulate extracellular matrix turnover, cell signaling, inflammation, and tissue remodeling; dysregulation contributes to disease.

\medterm{Metalloprotein} A protein that contains a metal ion or binds one tightly as part of its structure or function. Iron, zinc, copper, and magnesium in metalloproteins support oxygen transport, enzyme activity, electron transfer, or regulation.

\medterm{Metallothionein} Metallothionein is a small, cysteine-rich protein that binds metals such as zinc, copper, and cadmium. It helps regulate essential metals and may protect cells from metal toxicity and oxidative stress.

\medterm{Metanephrine} Metanephrine is a metabolite of epinephrine formed mainly by catechol-O-methyltransferase. Plasma or urinary metanephrine measurements are used in evaluation for catecholamine-secreting tumors such as pheochromocytoma and paraganglioma.

\medterm{Metaphase} The stage of cell division in which chromosomes align at the cell's equatorial plane.

\medterm{Metaphase Plate} The metaphase plate is the plane at the center of a dividing cell where chromosomes align during metaphase. It is a geometric arrangement rather than a discrete anatomical structure.

\medterm{Metaphyseal Corner Fractures} Small corner or bucket-handle breaks near the growth plates of an infant’s long bones. They are highly concerning for forceful injury, especially when multiple or at different healing stages, but must be interpreted with the full medical and imaging assessment.

\medterm{Metaphyseal Widening} Metaphyseal widening is enlargement of the growth-plate or metaphyseal region of a developing long bone. It may be seen in disorders such as rickets, skeletal dysplasia, infection, or healing injury and must be interpreted with imaging and clinical findings.

\medterm{Metaphysis} The wider part of a growing long bone between its shaft and end. In children it lies next to the growth plate, where new bone forms and the bone lengthens.

\medterm{Metaplasia} A usually reversible change in which one mature cell type is replaced by another that better tolerates ongoing irritation or stress. Some persistent forms may increase the risk of further abnormal changes.

\medterm{Metaplastic Carcinoma} A cancer whose cells have changed to resemble a different type of tissue. The term is used especially for an uncommon, often aggressive breast cancer; diagnosis requires examination of a tissue sample.

\medterm{Metapneumovirus} A respiratory virus that can cause cough, fever, bronchiolitis, or pneumonia, especially in infants, older adults, and people with chronic disease.

\medterm{Metastasectomy} Surgery to remove one or more cancer deposits that have spread from the original tumor. It may help selected people when the deposits are limited and removable, but benefit depends on the cancer type, location, and other treatment options.

\medterm{Metastases To The Breast} Secondary malignant deposits in breast tissue from a nonmammary primary tumor, most often melanoma, lymphoma, or carcinoma from another organ. Imaging and biopsy are needed because the treatment differs from primary breast carcinoma.

\medterm{Metastasis} The spread of cancer cells from the original tumor to a separate part of the body, where they form new tumors. Spread can occur through the blood, lymphatic system, or body cavities.

\medterm{Metastasis Of Colorectal Cancer} Spread of cancer from the colon or rectum to another part of the body, most often the liver, lungs, or abdominal lining. Symptoms and treatment depend on the sites involved, number of tumors, and tumor test results.

\medterm{Metastasis-Free Survival} The length of time after treatment before cancer spreads to a distant part of the body or the person dies. Researchers use it to compare treatments and estimate how well they delay spread.

\medterm{Metastatic Adenocarcinoma} A gland-forming cancer that has spread from where it began to distant organs such as the liver, lungs, bones, or brain. Identifying the original site is important because treatment depends on its origin and spread.

\medterm{Metastatic Brain Tumor} A metastatic brain tumor is cancer that has spread to the brain from a primary tumor elsewhere in the body.

\medterm{Metastatic Breast Cancer} Breast cancer that has spread from the breast to distant organs or tissues, such as bone, liver, lung, or brain. It is treatable, but usually considered advanced and not curable.

\medterm{Metastatic Calcification} Buildup of calcium salts in otherwise healthy tissues because the blood calcium level is too high. It can affect the kidneys, lungs, stomach, or blood vessels and may damage their function.

\medterm{Metastatic Cancer} Cancer that has spread from its original site to distant organs or tissues. The spread may occur through the blood or lymphatic system and can form new tumors. Metastatic cancer is often called stage 4 cancer, although staging depends on the cancer type. The new tumors retain the characteristics of the original cancer; for example, breast cancer that spreads to the lung remains metastatic breast cancer, not lung cancer.

\synonyms
Stage 4 Cancer; Advanced Metastatic Cancer; Secondary Cancer

\medterm{Metastatic Castration-Resistant Prostate Cancer} Prostate cancer that has spread and continues growing despite treatment that lowers male hormones. Care may include hormone-targeting medicines, chemotherapy, radiation, immune treatment, or medicines for bone protection.

\medterm{Metastatic Colon Cancer} Colon cancer that has spread beyond the bowel, most often to the liver, lungs, or abdominal lining. Symptoms and treatment depend on where it has spread, the tumor’s genetic findings, and the person’s overall health.

\medterm{Metastatic Colorectal Cancer} Alternate name; see \textit{Colorectal Cancer}.

\medterm{Metastatic Inguinal Lymphadenopathy} Enlargement of lymph nodes in the groin because cancer cells have spread to them, often from the leg, genital area, lower abdomen, or anus. Examination and imaging help identify the original cancer and its extent.

\medterm{Metastatic Lung Cancer} Lung cancer that has spread beyond the lungs to organs such as the brain, bones, liver, or adrenal glands. Treatment depends on the cancer type, molecular findings, areas affected, and overall health.

\medterm{Metastatic Pleural Tumor} A cancer deposit in the thin lining around the lung that has spread from another tumor, often from the lung, breast, or digestive tract. It may cause chest pain, breathlessness, or fluid around the lung.

\medterm{Metastin} Metastin is an older name for kisspeptin, a peptide derived from the KISS1 gene. It stimulates hypothalamic gonadotropin-releasing hormone secretion and helps regulate puberty, reproductive function, and tumor-cell signaling.

\medterm{Metatarsal} Relating to one of the five long bones in the middle and front of the foot. Metatarsal bones connect the ankle bones to the toe bones and may be injured by falls, twisting, or repeated stress.

\medterm{Metatarsal Bone} A metatarsal bone is one of the five long bones of the forefoot between the tarsal bones and proximal phalanges. Metatarsal fractures and alignment abnormalities can affect weight bearing and gait.

\medterm{Metatarsal Valgus} Metatarsal valgus is outward deviation of one or more metatarsal bones relative to the foot’s normal axis. It may contribute to forefoot deformity and altered load distribution and is assessed clinically and radiographically.

\medterm{Metatarsalgia} Pain in the front or ball of the foot, often arising from pressure or irritation around the metatarsal heads.

\medterm{Metatarsals} The five bones forming the front half of the foot between the ankle and toes. They support body weight and help with walking; overuse or injury can cause pain or fractures.

\medterm{Metatarsus} The part of the foot between the tarsal bones and toes, containing the five metatarsal bones.

\medterm{Metatarsus Adductus} A foot deformity in which the forefoot is angled inward relative to the hindfoot, often recognized in infancy and ranging from flexible to rigid.

\medterm{Metatarsus Varus} Metatarsus varus is inward deviation of the forefoot caused by medial angulation of the metatarsal bones. It is often recognized in infancy and may be flexible or rigid; management depends on severity and persistence.

\medterm{Metencephalon} The metencephalon is the embryonic brain region that develops into the pons and cerebellum. These structures participate in motor coordination, balance, cranial-nerve functions, and relay of signals between the brain and spinal cord.

\medterm{Meteorism} Abdominal distension caused by excessive gas in the gastrointestinal tract.

\medterm{Meter} A metric unit of length equal to 100 centimeters; medical usage may also refer to a measuring device.

\medterm{Metered-Dose Inhaler} A pressurized inhaler that releases a measured amount of medicine as a spray. Correct timing and breathing are needed for the medicine to reach the lungs; a spacer can make use easier and improve delivery.

\medterm{Metestrus} The stage of the reproductive cycle after ovulation in many mammals, when the corpus luteum develops and progesterone rises. The exact timing and biological features differ among species.

\medterm{Metformin} A commonly used medicine for type 2 diabetes that lowers the amount of sugar made by the liver and helps the body use insulin better. Nausea or diarrhea are common; serious illness or severe kidney problems increase rare toxicity risk.

\medterm{Metformin Hydrochloride} Metformin hydrochloride is the hydrochloride salt of metformin, an oral biguanide used chiefly for type 2 diabetes. It lowers blood glucose without directly stimulating pancreatic insulin release and is available in immediate- and extended-release formulations.

\medterm{Metformin In Prediabetes} Metformin may reduce progression from prediabetes to type 2 diabetes, especially in people at
high risk. Its use depends on age, body size, prior gestational diabetes, kidney function,
other conditions, and patient preference; lifestyle intervention remains important.

\medterm{Methacholine Challenge Test} A breathing test used when asthma is suspected but routine breathing tests are normal. The person inhales increasing amounts of methacholine, which can narrow sensitive airways; breathing is measured after each dose under medical supervision.

\medterm{Methacholine Chloride} Methacholine chloride is a muscarinic cholinergic agonist used in bronchoprovocation testing. In susceptible people it causes airway narrowing, allowing assessment for bronchial hyperresponsiveness under controlled medical supervision.

\medterm{Methadone} A long-acting opioid medicine used in selected pain cases and treatment of opioid-use disorder. It reduces withdrawal and cravings but can cause dangerous breathing suppression, sedation, dependence, or heart-rhythm changes.

\medterm{Methadone Hydrochloride} Methadone hydrochloride is the hydrochloride salt of methadone. It provides long-lasting opioid-receptor activity for selected analgesic and opioid-use-disorder indications, with risks of respiratory depression, dependence, overdose, and cardiac arrhythmia.

\medterm{Methadone Overdose} Taking too much methadone can cause profound sleepiness, pinpoint pupils, slow or stopped breathing, coma, and death. Effects may be delayed or prolonged, and abnormal heart rhythms can occur; call emergency services immediately.

\medterm{Methaemoglobinaemia} A condition in which hemoglobin is changed so it cannot carry or release oxygen normally. It may cause blue or gray skin, headache, breathlessness, confusion, or severe oxygen shortage after exposure to certain medicines or chemicals.

\medterm{Methamidophos} Methamidophos is a highly toxic organophosphorus insecticide that inhibits acetylcholinesterase. Exposure can cause cholinergic poisoning with sweating, salivation, vomiting, bronchospasm, bradycardia, muscle weakness, and potentially respiratory failure.

\medterm{Methamphetamine} Methamphetamine is a potent central nervous system stimulant that increases synaptic catecholamines. It has limited medical use in some settings but is commonly misused; toxicity and chronic use can cause cardiovascular complications, psychosis, dependence, and neurocognitive harm.

\medterm{Methamphetamine Overdose} Taking too much methamphetamine can cause agitation, dangerously high temperature, fast heartbeat, high blood pressure, chest pain, seizures, stroke, abnormal rhythms, or severe confusion. This is an emergency requiring immediate medical care.

\medterm{Methamphetamines} Powerful stimulant drugs that increase activity in the brain and body. Repeated or high-dose use can cause dependence, severe heart problems, overheating, sleeplessness, paranoia, hallucinations, seizures, stroke, or death.

\medterm{Methane} Methane is a colorless, odorless, highly flammable gas produced by anaerobic decomposition and normal intestinal microbial activity. It can displace oxygen in enclosed spaces and cause asphyxiation; it is not itself a direct systemic poison at ordinary exposure levels.

\medterm{Methane Sulfonate} Methane sulfonate, or mesylate, is the anion or ester derived from methanesulfonic acid. Mesylate salts are commonly used to improve the formulation of medicines; some methanesulfonate esters are potent alkylating agents and toxic impurities.

\medterm{Methanobrevibacter Smithii} A methane-producing archaeon commonly found in the human intestine. It consumes hydrogen made by other gut microorganisms and may influence intestinal fermentation and gas production.

\medterm{Methanol} A poisonous alcohol whose breakdown products can cause severe acid buildup, blindness, brain injury, or death. Suspected ingestion is a medical emergency requiring immediate specialized treatment, even before symptoms appear.

\medterm{Methanol Poisoning} Poisoning by methanol, a type of alcohol found in some industrial products. Its breakdown can cause severe acid buildup, blindness, brain injury, or death; urgent hospital treatment is needed even before symptoms appear.

\medterm{Methanol Test} A blood test that measures methanol, a poisonous alcohol. It helps confirm toxic exposure, but treatment may begin before results because methanol can cause severe metabolic acidosis and delayed blindness or brain injury.

\medterm{Methazolamide} Methazolamide is a carbonic anhydrase inhibitor used mainly to lower intraocular pressure in glaucoma. It reduces aqueous humor production and can cause metabolic acidosis, electrolyte abnormalities, kidney stones, or sulfonamide-related reactions.

\medterm{Methemoglobin} Hemoglobin in which iron has been oxidized so it cannot carry oxygen normally. High levels cause methemoglobinemia, with cyanosis, headache, breathlessness, confusion, or severe tissue oxygen deficiency.

\medterm{Methemoglobinemia} Methemoglobinemia occurs when hemoglobin iron is oxidized and cannot carry oxygen normally, causing cyanosis, headache, breathlessness, or severe tissue hypoxia.

\medterm{Methenamine} Methenamine is a urinary antiseptic that decomposes in acidic urine to release formaldehyde, which inhibits bacterial growth. It is used for prevention of recurrent urinary tract infection in selected patients, not for treatment of active upper-tract infection.

\medterm{Methicillin-Resistant Staphylococcus Aureus} A strain of Staphylococcus aureus that is resistant to several common antibiotics. It can cause skin, wound, lung, blood, bone, or other infections, and treatment is chosen using the infection site and laboratory testing.

\medterm{Methicillin-Resistant Staphylococcus Aureus Infection} Alternate name; see \textit{MRSA Infection}.

\medterm{Methidathion} Methidathion is an organophosphorus insecticide and acaricide that inhibits acetylcholinesterase. Exposure can produce a cholinergic toxidrome with secretions, bronchospasm, bradycardia, gastrointestinal symptoms, weakness, and respiratory failure.

\medterm{Methimazole} A medicine that reduces production of thyroid hormones and is used mainly for hyperthyroidism and Graves disease. It can rarely cause severe low white-cell counts or liver injury; fever or sore throat during treatment needs prompt medical advice.

\medterm{Methiocarb} Methiocarb is a carbamate pesticide that inhibits acetylcholinesterase reversibly. Poisoning causes excessive muscarinic and nicotinic cholinergic activity and may require urgent supportive care and antidotal treatment.

\medterm{Methionine} An essential amino acid containing sulfur that the body uses to build proteins and make other substances. People must obtain it from food, and it can be converted into homocysteine and related compounds.

\medterm{Methionine Enkephalin} Methionine enkephalin is an endogenous opioid pentapeptide that binds opioid receptors and participates in modulation of pain, stress, reward, and neuroendocrine function. It is one of the products of proenkephalin processing.

\medterm{Methocarbamol} Methocarbamol is a centrally acting skeletal-muscle relaxant used with rest and physical measures for acute painful musculoskeletal conditions. It can cause drowsiness, dizziness, impaired coordination, and additive central nervous system depression.

\medterm{Methohexital} Methohexital is an ultra-short-acting barbiturate anesthetic used for induction of general anesthesia and, in selected settings, procedural sedation. It depresses central nervous system activity and may cause respiratory or cardiovascular depression.

\medterm{Methomyl} Methomyl is a carbamate insecticide that reversibly inhibits acetylcholinesterase. Poisoning produces cholinergic symptoms such as salivation, sweating, vomiting, bronchospasm, bradycardia, weakness, and respiratory compromise.

\medterm{Methotrexate} An antimetabolite that inhibits folate-dependent enzymes and is used in cancer and selected autoimmune or inflammatory disorders.

\medterm{Methotrexate In Dermatology} An immunomodulatory antimetabolite used in selected inflammatory and autoimmune skin
diseases, including psoriasis, eczema, and some connective-tissue disorders. Treatment
requires attention to liver disease, blood counts, infection risk, drug interactions,
and strict pregnancy precautions.

\medterm{Methotrimeprazine} Methotrimeprazine, also called levomepromazine, is a phenothiazine with antipsychotic, sedative, antiemetic, and analgesic properties. It is used in some countries for psychiatric illness and palliative-care symptoms and can cause sedation, hypotension, and anticholinergic effects.

\medterm{Methoxamine} Methoxamine is a predominantly alpha-adrenergic agonist that increases vascular tone and blood pressure. It has been used as a vasopressor for hypotension but is not a routine first-line agent in many current practice settings.

\medterm{Methoxsalen} Methoxsalen is a psoralen that intercalates into DNA and, after exposure to UVA, forms cross-links that suppress abnormal lymphocyte or skin-cell activity. It is used with UVA phototherapy for selected skin diseases; photosensitivity and skin-cancer risk require precautions.

\medterm{Methoxyflurane} A vapor medicine formerly used to cause general anesthesia and, in some settings, provide short-term pain relief. High or repeated exposure can damage the kidneys, so use requires careful medical supervision.

\medterm{Methyclothiazide} Methyclothiazide is a thiazide diuretic that inhibits sodium and chloride reabsorption in the distal convoluted tubule. It is used for hypertension or edema and may cause hypokalemia, hyponatremia, hyperuricemia, and glucose intolerance.

\medterm{Methyl Alcohol} Another name for methanol, a toxic alcohol that can cause severe acidosis and blindness.

\medterm{Methyl Bromide} Methyl bromide, or bromomethane, is a volatile fumigant and alkylating chemical. Inhalation can cause severe neurologic, respiratory, renal, and cardiovascular toxicity, including delayed seizures and pulmonary edema.

\medterm{Methyl Chloride} Methyl chloride, or chloromethane, is a volatile industrial gas formerly used in refrigerants and other applications. High-level inhalation exposure can depress the central nervous system and cause headache, ataxia, seizures, or cardiac effects.

\medterm{Methyl Isocyanate} Methyl isocyanate is a highly volatile and reactive industrial chemical. Inhalation or eye and skin exposure can cause severe irritation, pulmonary edema, and delayed respiratory injury; it was the toxic agent in the Bhopal disaster.

\medterm{Methyl Parathion} Methyl parathion is an organophosphorus insecticide that inhibits acetylcholinesterase. Exposure can produce severe cholinergic poisoning with secretions, bronchospasm, bradycardia, weakness, seizures, and respiratory failure.

\medterm{Methyl Salicylate Overdose} Taking too much methyl salicylate, found in some oils and pain-relief products, can cause salicylate poisoning with vomiting, ringing in the ears, rapid breathing, confusion, seizures, acidosis, or coma. Emergency assessment is needed.

\medterm{Methylation} Attachment of a small chemical group to DNA, proteins, or other molecules. On DNA, this can change how strongly a gene is used and helps regulate normal development, cell function, and some diseases.

\medterm{Methylcellulose} A plant-fiber product used as a bulk-forming laxative and as a thickener in medicines and foods. It absorbs water and adds bulk to stool; it should be taken with enough fluid and avoided when bowel blockage is suspected.

\medterm{Methyldopa} A blood-pressure medicine that acts in the brain to reduce nerve signals that tighten blood vessels. It is used in selected cases, including pregnancy, and may cause sleepiness, dizziness, or liver problems.

\medterm{Methylene Blue} A medicine used mainly for serious methemoglobinemia, a condition in which blood cannot carry oxygen normally. It may be unsafe in some people with an inherited red-cell enzyme deficiency and can interact with certain antidepressants.

\medterm{Methylene Blue Test} A test or procedure using methylene blue dye to trace fluid flow, identify a structure, or detect a leak or abnormal connection. The purpose depends on the body site; allergy and blood-related complications are uncommon but possible.

\medterm{Methylene Chloride} Methylene chloride, or dichloromethane, is a volatile industrial solvent. Inhalation can cause central nervous system depression and metabolism to carbon monoxide can produce tissue hypoxia; high exposure may be fatal.

\medterm{Methylenedioxymethamphetamine} 3,4-Methylenedioxymethamphetamine, abbreviated MDMA and commonly called ecstasy or molly, is a psychoactive drug with stimulant and hallucinogenic effects that increases serotonergic signaling and can cause serious toxicity.

\medterm{Methylergometrine} An ergot alkaloid that increases uterine contraction and may be used to control postpartum hemorrhage when appropriate.

\medterm{Methylmalonic Acid Blood Test} A blood test measuring methylmalonic acid to help detect vitamin B12 deficiency, especially when the B12 result is borderline. Kidney impairment can raise the result and must be considered during interpretation.

\medterm{Methylmalonic Acid Level} The amount of methylmalonic acid measured in blood or urine. A high level can support vitamin B12 deficiency, especially when the B12 result is unclear, but kidney disease can also raise it.

\medterm{Methylmalonic Acidemia} An inherited metabolic disorder in which the body cannot properly process certain fats and proteins, causing methylmalonic acid to build up. Episodes may cause vomiting, dehydration, acidosis, low blood sugar, seizures, or developmental and kidney problems.

\medterm{Methylmalonic Aciduria} A disorder in which methylmalonic acid builds up in blood and urine, usually because of an inherited enzyme or vitamin B12-processing problem. It may cause vomiting, poor feeding, acidosis, developmental problems, kidney disease, or neurologic injury.

\medterm{Methylmercury Compounds} Methylmercury compounds are organic mercury compounds that can accumulate in aquatic food chains and cross the blood--brain and placental barriers. Exposure primarily injures the nervous system and can impair fetal neurodevelopment.

\medterm{Methylmercury Poisoning} Poisoning caused by methylmercury, usually through contaminated food. It primarily injures the nervous system and may cause numbness, poor coordination, vision or hearing changes, and impaired fetal brain development; exposure requires medical evaluation.

\medterm{Methylmethacrylate} Methyl methacrylate is a volatile monomer used to make acrylic plastics and dental or orthopedic bone cement. Vapors may irritate the eyes and airways, and liquid exposure can cause dermatitis or sensitization.

\medterm{Methylmorphine} Another name for codeine, an opioid medicine used for selected pain or cough indications. It can cause drowsiness, constipation, dependence, and dangerous breathing suppression, especially with alcohol or other sedatives.

\medterm{Methylnaltrexone} Methylnaltrexone is a peripherally acting mu-opioid-receptor antagonist used to treat opioid-induced constipation when laxatives are insufficient. It generally preserves central analgesia but can cause abdominal pain, diarrhea, and rarely gastrointestinal perforation.

\medterm{Methylnaltrexone Bromide} Methylnaltrexone bromide is the bromide salt of methylnaltrexone, a peripherally restricted opioid antagonist. It reverses opioid effects in the gastrointestinal tract and is used for opioid-induced constipation under indication-specific precautions.

\medterm{Methylphenidate} A stimulant medicine used mainly to treat attention-deficit/hyperactivity disorder and narcolepsy, a condition causing excessive sleepiness. It affects brain chemicals involved in attention and alertness; reduced appetite, sleep problems, increased blood pressure, and misuse are possible risks.

\medterm{Methylprednisolone} A synthetic glucocorticoid used to suppress inflammation and immune activity.

\medterm{Methyltestosterone} A synthetic androgen related to testosterone. It can affect sexual development and other androgen-sensitive tissues but may cause liver injury, adverse cholesterol changes, infertility, or virilization and is used only in selected medical situations.

\medterm{Metoclopramide} A medicine that helps the stomach empty and relieves nausea and vomiting. It is used for selected digestive problems and treatment-related nausea; prolonged or high-dose use can cause involuntary movements and other nervous-system effects.

\medterm{Metolazone} A water-removing medicine used for high blood pressure and swelling. It helps the kidneys remove salt and water and is often combined with another diuretic, but can disturb sodium, potassium, kidney function, or hydration.

\medterm{Metolcarb} Metolcarb is a carbamate insecticide that can inhibit acetylcholinesterase. Exposure may cause cholinergic poisoning with excessive secretions, gastrointestinal symptoms, bronchospasm, bradycardia, weakness, or respiratory compromise.

\medterm{Metopic Ridge} A metopic ridge is a palpable or visible line along the forehead from fusion of the metopic suture; it may be normal or associated with trigonocephaly.

\medterm{Metoprolol} Metoprolol is a relatively beta-1-selective blocker used for hypertension, angina, heart failure, rate control, and some post-infarction settings; bradycardia, hypotension, fatigue, and bronchospasm can occur.

\medterm{Metric System} The metric system is a decimal system of measurement based on units such as the meter, kilogram, and liter. Modern medicine generally uses the International System of Units, with conventional units retained for some measurements.

\medterm{Metritis} Inflammation or infection of the uterus, often after childbirth, miscarriage, abortion, or a uterine procedure. It may cause fever, lower abdominal pain, unusual discharge, bleeding, or tenderness and usually needs treatment.

\medterm{Metronidazole} An antibiotic and antiparasitic medicine used for certain infections caused by bacteria that grow without oxygen and by some parasites. It can cause nausea or a metallic taste; nerve problems may occur with prolonged use.

\medterm{Metronidazole-Associated Encephalopathy} A rare complication of prolonged or high-dose metronidazole therapy, causing cerebellar
dysfunction (ataxia, dysarthria, nystagmus) and sometimes confusion or seizures. MRI may
show T2-hyperintense lesions in the cerebellar dentate nuclei, corpus callosum splenium,
and brainstem. Symptoms typically resolve with drug discontinuation.

\medterm{Metronomic Chemotherapy} Cancer treatment given in frequent, usually smaller doses over a long period instead of occasional high doses. The goal is to slow cancer growth and sometimes its blood supply; use depends on the cancer and treatment plan.

\medterm{Metrorrhagia} Uterine bleeding that occurs at unexpected times, including between menstrual periods. Causes include pregnancy, hormonal changes, fibroids, polyps, infection, medicines, and cancer; new, heavy, prolonged, or postmenopausal bleeding needs medical assessment.

\medterm{METs} Metabolic equivalents used to estimate activity intensity or energy use. One MET is approximately the energy used while resting; higher values indicate more demanding activity and help describe exercise capacity.
\synonyms
Metabolic Equivalents

\medterm{Metyrapone} A medicine that blocks one step in the production of cortisol, a hormone made by the adrenal glands. It is used for selected hormone tests or treatment under specialist supervision because levels can change quickly.

\medterm{Metzenbaum Scissors} Fine, lightweight surgical scissors designed to separate and cut delicate soft tissue. Their relatively narrow blades are useful for careful dissection, but they are not intended for cutting hard material or heavy sutures.

\medterm{Mevinphos} Mevinphos is a highly toxic organophosphorus insecticide that inhibits acetylcholinesterase. Exposure can rapidly cause cholinergic crisis, including bronchorrhea, bronchospasm, muscle weakness, seizures, and respiratory failure.

\medterm{Mexacarbate} Mexacarbate is a carbamate insecticide that interferes with acetylcholinesterase. Toxic exposure may produce muscarinic and nicotinic cholinergic symptoms and requires prompt decontamination and medical evaluation.

\medterm{Mexiletine} Mexiletine is an oral class IB antiarrhythmic and sodium-channel blocker used for selected ventricular arrhythmias and sometimes neuropathic pain. It can cause gastrointestinal and neurologic effects and may worsen or provoke arrhythmias in some patients.

\medterm{Meyerozyma Guilliermondii} Meyerozyma guilliermondii is a yeast formerly classified as Candida guilliermondii. It is usually an environmental or colonizing organism but can cause opportunistic bloodstream or other infections, particularly in immunocompromised patients.

\medterm{MF} Abbreviation; see \textit{Mycosis Fungoides}.

\medterm{MFS} Abbreviation; see \textit{Myxofibrosarcoma}.

\medterm{MG} Abbreviation; see \textit{Myasthenia Gravis}.

\medterm{MGUS} See \textit{Monoclonal Gammopathy of Undetermined Significance}.

\medterm{MHC Molecule} A protein on the surface of a cell that displays small pieces of proteins to T cells. This helps the immune system recognize infected or abnormal cells and distinguish the body's own cells from foreign material.

\synonyms
Major Histocompatibility Complex Molecule

\medterm{Mianserin} An antidepressant medicine used in some countries to treat depression. It changes signaling by several brain chemicals and can cause sleepiness, increased appetite, weight gain, dizziness, or other medicine-related effects.

\medterm{MIBG Scan} A nuclear medicine scan that uses a small amount of radioactive MIBG to find certain tumors of nerve-related tissue. It is used especially for pheochromocytoma, paraganglioma, and neuroblastoma; medicines and thyroid protection may be needed before the scan.

\medterm{MIBG Scintiscan} A nuclear medicine scan using radiolabeled metaiodobenzylguanidine to image tissues derived from the sympathetic nervous system, including selected pheochromocytomas, paragangliomas, and neuroblastomas.

\medterm{Mica} A group of sheet silicate minerals; inhalation of mica dust can cause occupational lung disease.

\medterm{Micafungin} An echinocandin antifungal agent used in the treatment and prophylaxis of invasive
candidiasis and esophageal candidiasis. Micafungin inhibits beta-(1,3)-D-glucan
synthase, disrupting fungal cell wall synthesis and leading to osmotic instability and
cell death.

\medterm{Micelle} A tiny cluster of molecules that helps carry fats and fat-soluble substances through intestinal fluid for absorption.

\medterm{Miconazole} An antifungal medicine used in creams, gels, or other preparations for infections of the skin, mouth, or vagina. It stops fungi from making an essential part of their outer covering; irritation and medicine interactions are possible, especially with some oral products.

\medterm{Microabscess} A very small pocket of pus containing immune cells, damaged tissue, and sometimes germs. It may occur in the skin or an organ because of infection or inflammation and is usually a finding rather than a diagnosis.

\medterm{Microalbuminuria} A small but abnormal amount of albumin, a blood protein, in the urine. It can be an early sign of kidney damage, especially in diabetes or high blood pressure, but temporary illness or exercise can also raise it and repeat testing may be needed.

\medterm{Microalbuminuria Test} A urine test measuring small amounts of albumin, usually by an albumin-to-creatinine ratio, to detect early kidney damage, especially in diabetes or hypertension.

\medterm{Microaneurysm} A very small balloon-like bulge in a blood vessel wall. In the retina, microaneurysms are an early sign of diabetic eye disease and may leak fluid, causing swelling and blurred vision. Their number and associated retinal changes help show how diabetes is affecting the eyes.

\medterm{Microangiopathic Hemolysis} Destruction of red blood cells as they pass through damaged or narrowed small blood vessels. Blood tests may show fragmented cells, anemia, and other signs of red-cell breakdown; causes include severe high blood pressure, pregnancy complications, infection, and clotting disorders.

\medterm{Microangiopathic Hemolytic Anemia} An anemia caused by mechanical fragmentation of red blood cells as they pass through damaged or narrowed small blood vessels, producing schistocytes and often elevated hemolysis markers.

\medterm{Microangiopathy} Disease or damage affecting the body's smallest blood vessels. Thickened or injured vessel walls can reduce blood flow or allow leakage and may damage organs, especially in diabetes, high blood pressure, or clotting disorders.

\medterm{Microarray} A laboratory test using a small chip covered with DNA probes. It can measure activity of many genes or find selected missing or extra DNA segments and is used mainly for specialized genetic or research testing.

\medterm{Microarray Analysis} Microarray analysis measures the abundance or hybridization of many nucleic-acid sequences simultaneously using an array of immobilized probes. It can assess gene expression, copy-number changes, or selected genetic variants.

\medterm{Microbacterium} A genus of gram-positive bacteria found in soil and water that can rarely cause bloodstream or other opportunistic infections.

\medterm{Microbe} A microorganism, such as a bacterium, fungus, protozoan, or virus.

\synonyms
Microorganism

\medterm{Microbial Biofilms} Microbial biofilms are organized communities of microorganisms attached to a surface and enclosed in a self-produced extracellular matrix. Biofilms can form on tissues and medical devices and often show increased tolerance to antimicrobials and host defenses.

\medterm{Microbicide} A microbicide is a chemical or biological agent that kills microorganisms or inactivates them. The term may refer to disinfectants, topical agents, or investigational products intended to reduce sexual transmission of infection.

\medterm{Microbiology} The study of microorganisms, including bacteria, viruses, fungi, and parasites, and how they interact with people, animals, plants, and the environment. In medicine, it supports the diagnosis, prevention, and treatment of infectious diseases.

\medterm{Microbiome} The community of microorganisms and their genetic material living in a body site, especially the intestine. These organisms can influence digestion, immunity, and health, although the effects differ by person and condition.

\medterm{Microcalcification} A very small calcium deposit in tissue. On a mammogram, its shape and arrangement may reflect normal change, injury, infection, or early breast cancer, so suspicious findings may require further imaging or biopsy.

\medterm{Microcephaly} Microcephaly is head size substantially smaller than expected for age and sex, reflecting genetic, developmental, infectious, toxic, or acquired brain causes.

\medterm{Microcirculation} The movement of blood through the body's smallest blood vessels. In these vessels, oxygen and nutrients pass to tissues and waste products pass into the blood; poor microcirculation can slow healing and damage organs.

\medterm{Microcolon} Microcolon is an abnormally small-caliber colon, often identified on contrast imaging. It may result from reduced fetal intestinal use or obstruction and is associated with conditions such as small-left-colon syndrome, meconium ileus, or distal intestinal obstruction.

\medterm{Microcornea} Microcornea is an abnormally small cornea, usually defined in adults by a horizontal diameter below about 10 mm. It may be isolated or associated with microphthalmia, glaucoma, cataract, or genetic syndromes.

\medterm{Microcurie} A microcurie ($\mu$Ci) is a unit of radioactivity equal to one-millionth of a curie, or approximately 37 kilobecquerels. It may be used to express the activity of radioactive tracers or sources.

\medterm{Microcystic Adnexal Carcinoma} A rare, slow-growing but locally aggressive malignant adnexal tumor, usually on the head or neck, with follicular and ductal differentiation; it may extend deeply and requires complete excision.

\medterm{Microcytic Anemia} A form of anemia characterized by abnormally small and pale circulating red blood cells, defined on a complete blood count by a mean corpuscular volume (MCV) below $80\,\text{fL}$. It occurs when the body is unable to synthesize adequate hemoglobin, starving red cells of their internal content. The five primary causes are commonly remembered by the mnemonic TAILS: Thalassemia (inherited globin chain defect), Anemia of chronic disease/inflammation (hepcidin-mediated iron trapping), Iron deficiency anemia (the most common cause, from chronic blood loss, poor diet, or pregnancy), Lead poisoning (toxic inhibition of heme enzymes), and Sideroblastic anemia (impaired iron incorporation in the bone marrow). Symptoms include fatigue, pallor, weakness, and shortness of breath; diagnostic evaluation includes serum ferritin, iron studies, and hemoglobin electrophoresis.
\synonyms{Microcytic Hypochromic Anemia; Small-Cell Anemia; Low-MCV Anemia}

\medterm{Microdermabrasion} A cosmetic treatment that gently removes the outermost layer of skin with fine crystals or a rough-tipped device. It may improve the look of mild uneven texture, dullness, or some marks, but can cause redness, irritation, infection, or temporary darkening.

\medterm{Microdialysis} A testing method that uses a very small tube to collect chemicals from the fluid around cells. It can measure local levels of medicines or natural substances, including in the brain or beneath the skin.

\medterm{Microdiscectomy} A minimally invasive spinal operation used to relieve pain, numbness, or weakness caused by a herniated intervertebral disc compressing a spinal nerve root in the lower back (lumbar radiculopathy or sciatica). Through a small incision in the midline of the lower back, a spine surgeon or neurosurgeon uses a high-magnification surgical operating microscope or surgical loupes to gently retract the overlying back muscle, remove a tiny sliver of bone (laminotomy), and visualize the inflamed nerve. The herniated or extruded disc fragments pressing against the nerve are carefully removed while leaving the healthy structural portion of the disc intact. Because tissue trauma is minimal compared to traditional open spine surgery, patients experience dramatic pain relief, often go home the same day or following morning, and resume light activities within a few weeks.

\synonyms
Lumbar Microdiscectomy; Microdiskectomy; Minimally Invasive Discectomy

\medterm{Microdissection} Microdissection is separation or isolation of very small tissue regions or structures under microscopic or highly magnified visualization. It is used in surgery, pathology, and molecular research to obtain selected cells or tissue.

\medterm{Microdissection Testicular Sperm Extraction} A surgical procedure that uses a microscope to search the testicle for small areas making sperm. It may help some men with very low or absent sperm production obtain sperm for assisted reproduction. The procedure can cause pain, bleeding, infection, or reduced testicular function, and sperm may not be found.

\synonyms
Micro-TESE; Microsurgical TESE

\medterm{Microdontia} Abnormally small teeth, which may affect one or several teeth and can be inherited or developmental.

\medterm{Microfibril} A very small thread-like structure that supports tissues. In connective tissue, microfibrils help organize elastic fibers and provide strength; in plants, they are made mainly of cellulose.

\medterm{Microfilament} A microfilament is a thin cytoskeletal filament composed mainly of actin, about 7 nm in diameter. Microfilaments support cell shape, motility, contraction, cytokinesis, and intracellular transport.

\medterm{Microfilament Proteins} Proteins that make, organize, move, or regulate actin microfilaments. They help cells maintain shape, migrate, divide, contract, and transport materials; examples include actin, myosin, and actin-binding proteins.

\medterm{Microfilaria} The tiny, immature form of a filarial worm that may circulate in blood or live in the skin. Mosquitoes, flies, or other biting insects can pick them up and spread the infection to another person.

\medterm{Microglandular Adenosis} Microglandular adenosis is an uncommon benign proliferation of small glands in the breast, often associated with hormonal influences such as pregnancy or progestin use. It can mimic invasive carcinoma and is diagnosed histologically.

\medterm{Microglia} Microglia are resident immune cells of the central nervous system. They survey neural tissue, remove debris, and respond to injury or infection; persistent activation can contribute to neuroinflammation.

\medterm{Microgliosis} Microgliosis is proliferation or increased activation of microglia in the central nervous system, usually in response to injury, infection, degeneration, or other inflammation. It is a histopathologic finding rather than a specific disease.

\medterm{Micrognathia} Micrognathia is an unusually small lower jaw that may affect feeding, teeth, facial development, or airway breathing.

\medterm{Micrognathism} Micrognathism is abnormal smallness of the jaw, usually referring to the mandible. It may be congenital or acquired and can affect feeding, dentition, airway patency, speech, or facial development.

\medterm{Micrograms} A unit of mass equal to one millionth of a gram, written as mcg. It is commonly used for small medicine doses and nutrient amounts; confusing micrograms with milligrams can cause a serious dosing error.

\medterm{Micrographia} Abnormally small handwriting, often becoming progressively smaller within a sentence. It commonly occurs in Parkinson disease and other movement disorders because of reduced movement size and control.

\medterm{Microinjection} Microinjection is delivery of a very small volume of fluid into a cell, tissue, or other microscopic target using a fine needle or micropipette. It is used in assisted reproduction, experimental biology, and some therapeutic procedures.

\medterm{Micrometastasis} A small group of cancer cells that has spread from the original tumor but is too small to detect on routine scans or examination. It may be found in tissue samples and can affect staging and treatment decisions.

\medterm{Micrometer} An instrument or unit used to measure very small distances; one micrometer is one-millionth of a meter.

\medterm{Micromolar} A concentration unit equal to one millionth of a mole per liter, written as \(\mu\mathrm{mol}/L\).

\medterm{Micromole} A micromole ($\mu$mol) is an amount of substance equal to one-millionth of a mole, or approximately $6.022\times10^{17}$ entities. It is commonly used for small quantities and concentrations of biochemical substances.

\medterm{Micron} A unit measuring one-millionth of a meter, commonly used to describe the size of cells, bacteria, and particles.

\medterm{Microneme} A microneme is a small secretory organelle in many apicomplexan parasites. It releases proteins that support host-cell attachment, motility, and invasion during the parasite life cycle.

\medterm{Micronization} Micronization is a process that reduces particles to micrometer-scale dimensions, often to increase surface area and improve dissolution or absorption of a drug. The resulting particle size and formulation still determine clinical performance.

\medterm{Micronized} Processed into very small particles to improve dispersion or absorption.

\medterm{Micronutrient} A micronutrient is a vitamin or mineral required in small amounts for normal growth, metabolism, and physiologic function. Deficiency or excess can cause disease, and requirements vary with age, pregnancy, illness, and other factors.

\medterm{Micronutrients} Vitamins and minerals required in small amounts for enzymatic, hormonal, structural,
hematologic, neurologic, and immune functions. They include fat-soluble vitamins (A, D,
E, and K), water-soluble vitamins, and essential minerals.

\medterm{Microorchidism} Abnormally small testes, which may result from impaired development, genetic syndromes, or testicular atrophy after injury or infection.

\medterm{Microorganism} A microscopic living organism, such as a bacterium, fungus, protozoan, or alga; viruses are acellular infectious agents.

\medterm{Microorganism Genetics} The study of genetic material, inheritance, variation, and gene activity in microorganisms. It helps explain how microbes grow, cause disease, develop drug resistance, and exchange useful or harmful traits.

\medterm{Micropenis} An unusually small penis caused by inadequate androgen effect or other developmental conditions, defined clinically by a stretched penile length more than 2.5 standard deviations below the mean for age and population. Evaluation may identify endocrine or genetic causes.

\medterm{Microphallus} An abnormally small penis, also called micropenis, usually defined relative to age and developmental standards; it may reflect hormonal or developmental disorders.

\medterm{Microphthalmia} A birth defect in which one or both eyes are smaller than usual.
It may occur with other eye or genetic conditions and can reduce vision. Treatment
depends on the cause and may include glasses, treatment for poor vision, or surgery.

\medterm{Microphthalmos} A birth condition in which one or both eyeballs are unusually small. Vision may be reduced, and other eye or genetic problems may occur, so affected children need an eye examination and ongoing specialist care.

\medterm{Microplastics} Very small plastic particles, generally less than five millimeters across, found in the environment and sometimes in food, water, and air. Human health effects are still being studied.

\medterm{Micropsia} A change in vision in which objects appear smaller or farther away than they really are. It can occur with migraine, retinal disease, brain disorders, medicines, or other problems and should be assessed if it is new, persistent, or associated with neurological symptoms.

\medterm{MicroRNA} A very small, non-coding RNA that helps control which genes are active, mainly by reducing the production of specific proteins. MicroRNAs are being studied as possible disease markers and treatment targets.

\medterm{Microsatellite Instability} A tumor pattern in which short repeated DNA sequences gain or lose units because the cell’s DNA repair system is faulty. A high level can occur in inherited or sporadic cancers and may help predict benefit from some immune treatments.

\medterm{Microscope} An instrument that magnifies structures too small to see clearly with unaided eyes. Light, electron, and fluorescence microscopes help examine cells, tissues, blood, microbes, and other samples.

\medterm{Microscopic} Visible only with a microscope or relating to microscopic examination.

\medterm{Microscopic Anatomy} The study of cells, tissues, and fine structural details that require magnification for examination; it includes cytology and histology.

\medterm{Microscopic Colitis} Long-lasting inflammation of the colon that usually causes frequent, watery, non-bloody diarrhea while the colon looks normal during routine colonoscopy. Tiny tissue samples are needed for diagnosis; the main forms are collagenous and lymphocytic colitis.

\synonyms
Collagenous Colitis; Lymphocytic Colitis

\medterm{Microscopic Examination} Examination of cells or tissue with a microscope to look for normal structure, infection, inflammation, or cancer. The sample may be stained and assessed together with its visible appearance, imaging, and other laboratory results.

\medterm{Microscopic Hematuria} Alternate name; see \textit{Hematuria}.

\medterm{Microscopic Polyangiitis} An autoimmune disease that inflames and damages small blood vessels. It can affect the kidneys, lungs, skin, nerves, or other organs and may cause blood in the urine, breathing problems, rash, weakness, or bleeding.

\medterm{Microscopy} Microscopy is examination of cells, tissues, organisms, or materials with a microscope. Light, fluorescence, electron, and other microscopy methods provide different information about structure and composition.

\medterm{Microsecond} A unit of time equal to one-millionth of a second.

\medterm{Microspectrophotometry} Microspectrophotometry measures the absorption or emission spectrum of a microscopic region or individual cell. It can quantify pigments, nucleic acids, proteins, or other chromophores in small biological samples.

\medterm{Microsphere} A microsphere is a small spherical particle, often made of polymer, glass, or another material, used in drug delivery, imaging, embolization, or laboratory assays. Its size, composition, and surface properties determine its biologic behavior.

\medterm{Microspherophakia} Microspherophakia is an abnormally small, spherical crystalline lens. It can occur in inherited disorders such as Weill--Marchesani syndrome and may cause lens displacement, pupillary block, angle-closure glaucoma, or reduced vision.

\medterm{Microsporidia} A group of microscopic spore-forming organisms related to fungi. They can cause prolonged diarrhea, eye disease, muscle disease, or widespread infection, especially in people with weakened immune systems.

\medterm{Microsporidiosis} An infection caused by microsporidia, microscopic organisms related to fungi that produce resistant spores. It most often affects the intestine and may cause long-lasting diarrhea and weight loss. It can also affect the eyes, muscles, kidneys, brain, or several organs, especially in people with weakened immune systems.

\medterm{Microsporum} Microsporum is a genus of dermatophyte fungi that can cause ringworm of the skin, scalp, or hair.

\medterm{Microsporum Audouinii} Microsporum audouinii is a dermatophyte that mainly infects the scalp and can cause tinea capitis, especially in children.

\medterm{Microsporum Distortum} Microsporum distortum is a rare dermatophyte related to M. canis that can cause skin or scalp infection.

\medterm{Microsporum Ferrugineum} A fungus that can cause ringworm of the scalp or skin, particularly in children. Laboratory examination can help confirm the infection and distinguish it from other causes of rash or hair loss.

\medterm{Microsporum Gypseum} Microsporum gypseum is a soil-associated dermatophyte that can cause ringworm of the skin or scalp.

\medterm{Microstomia} Microstomia is abnormal narrowing or smallness of the mouth opening. It may be congenital or result from scarring, burns, surgery, trauma, or connective-tissue disease and can impair eating, speech, oral hygiene, and dental treatment.

\medterm{Micro-Surgery} Surgery performed with magnification and fine instruments to repair very small structures such as nerves and vessels.

\medterm{Microsurgery} Microsurgery uses magnification and very small instruments to repair tiny blood vessels, nerves, or other structures.

\medterm{Microtechnology} Microtechnology is the design and manufacture of devices or structures with micrometer-scale features. Medical applications include microfluidic diagnostics, sensors, drug-delivery systems, and minimally invasive instruments.

\medterm{Microtia} A birth defect in which the outer ear is smaller than usual or
does not form fully. The ear canal may also be narrow or absent, which can reduce hearing.

\medterm{Microtome} An instrument that cuts extremely thin slices of tissue for examination under a microscope. The slices are mounted on slides, stained, and studied to identify normal structures or disease.

\medterm{Microtubule} A tiny hollow protein tube inside a cell. Microtubules help maintain cell shape, move materials within the cell, support cell division, and form part of cilia and other moving structures.

\medterm{Microtubule Bundle} A microtubule bundle is an organized group of parallel or otherwise coordinated microtubules within a cell. Bundles provide structural support and guide transport or movement, including in axons, cilia, and dividing cells.

\medterm{Microvascular} Relating to the body's smallest blood vessels, including tiny arteries, capillaries, and veins. These vessels deliver oxygen and nutrients, remove waste, and control blood flow through tissues.

\medterm{Microvascular Disease} Disease affecting the body's smallest blood vessels, which can reduce blood flow and oxygen delivery to organs or tissues. Diabetes, high blood pressure, smoking, inflammation, and aging can damage these vessels.

\medterm{Microvasculature} The network of the body's smallest blood vessels, including tiny arteries, capillaries, and veins. It delivers oxygen and nutrients, removes waste, controls tissue blood flow, and helps regulate fluid movement between blood and tissues.

\medterm{Microvillus} A microvillus is a microscopic, actin-supported projection of a cell membrane that increases surface area. Microvilli are abundant on intestinal epithelial cells, where they support nutrient absorption, and on renal tubular cells.

\medterm{Microvolt} A microvolt ($\mu$V) is one-millionth of a volt. It is used to express small bioelectric signals, such as those recorded by electroencephalography, electromyography, electrocardiography, or evoked-potential testing.

\medterm{Microwave Ablation} A procedure that uses microwave energy to heat and destroy selected tissue, usually under imaging guidance.

\medterm{Microwave Therapy} Treatment that uses microwave energy to heat or destroy selected tissue. It has specialized uses, such as treating some tumors or controlling bleeding, and the risks depend on the body area and device used.

\medterm{Micturition} The process of passing urine. The bladder squeezes while the muscles at its outlet relax, under coordinated control from the brain and nerves.

\medterm{Midaxillary Line} An imaginary vertical line on the side of the body passing through the apex of the
axilla (armpit).

\medterm{Midazolam} Midazolam is a short-acting benzodiazepine used for procedural sedation, anesthesia, and treatment of acute seizures in appropriate settings. It enhances GABA-A receptor activity and can cause profound sedation and respiratory depression, especially with opioids or other depressants.

\medterm{Midazolam Hydrochloride} Midazolam hydrochloride is the hydrochloride salt of the short-acting benzodiazepine midazolam. It is used for procedural sedation, anesthesia, and selected seizure emergencies and can cause dose-related sedation, hypotension, and respiratory depression.

\medterm{Mid-Brain} The midbrain, or mesencephalon, is the upper part of the brainstem involved in eye movements, motor pathways, and auditory and visual reflexes.

\medterm{Midbrain} The upper part of the brainstem, between the lower brainstem and the main brain. It helps control eye movements, movement, alertness, and automatic responses to sights and sounds.

\medterm{Midclavicular Line} An imaginary vertical line on the front of the body passing through the midpoint of the
clavicle (collarbone).

\medterm{Middle Ear} An air-filled space between the eardrum and inner ear that contains the malleus, incus, and stapes. These ossicles transmit
sound vibrations to the inner ear, and the auditory tube helps equalize middle-ear
pressure.

\medterm{Middle East Respiratory Syndrome} A serious respiratory infection caused by a coronavirus, usually linked to close contact with infected camels or people. It may cause fever, cough, pneumonia, breathing failure, or kidney and other organ problems.

\synonyms
MERS; Camel Flu

\medterm{Middle East Respiratory Syndrome} Alternate name; see \textit{MERS}.

\medterm{Middle Insomnia} Middle insomnia is repeated or prolonged awakening during the middle portion of the sleep period, often with difficulty returning to sleep. It may occur with stress, mood disorders, sleep apnea, pain, substances, medications, or other sleep disorders.

\medterm{Middle Lobe Syndrome} Chronic or recurrent atelectasis and infection of the right middle lobe due to bronchial
obstruction. The long, narrow right middle lobe bronchus is susceptible to compression
by lymph nodes, tumors, or mucous plugging. It may be caused by bronchiectasis,
malignancy, or granulomatous disease. Treatment depends on the underlying cause.

\medterm{Mid-Dorsal} Situated in the middle region of the back or on the dorsal side of a structure.

\medterm{Mid-Forearm} The middle part of the forearm, between the elbow and wrist. Measurements or injections described as mid-forearm use this region as a reference point.

\medterm{Mid-Gut} The embryonic middle portion of the primitive gut, which develops into much of the small intestine and proximal colon.

\medterm{Midgut Volvulus} A life-threatening twisting of the intestine that can cut off its blood supply, usually because the intestine formed abnormally before birth. Babies may suddenly develop green vomiting, a swollen or painful abdomen, blood in the stool, or extreme illness. It is an emergency requiring immediate imaging and surgery to untwist and protect the bowel.

\synonyms
Volvulus Neonatorum; Midgut Torsion

\medterm{Mid-Leg} The central part of the lower limb between the knee and ankle. It contains the shin bones, muscles, nerves, blood vessels, and soft tissues that support walking and ankle movement.

\medterm{Mid-Line} The imaginary or anatomical line dividing a structure or body into right and left portions.

\medterm{Midomafetamine} Midomafetamine is another name for 3,4-methylenedioxymethamphetamine, a psychoactive amphetamine derivative. It is being studied in controlled psychotherapy-assisted research but can cause hyperthermia, hyponatremia, cardiovascular toxicity, serotonin toxicity, and misuse.

\medterm{Midostaurin} Midostaurin is a multikinase inhibitor used with chemotherapy for newly diagnosed acute myeloid leukemia with a susceptible FLT3 mutation and for advanced systemic mastocytosis. It can cause nausea, cytopenias, infection, and other treatment-related toxicities.

\medterm{Midriff} Alternate informal term; see \textit{Abdomen}.

\medterm{Mid-Stream Urine Specimen} A urine sample collected after initial urine is voided, reducing contamination from the urethral opening and used commonly for culture.

\medterm{Mid-Systolic Click} A sharp heart sound heard during the middle or later part of the heartbeat, often associated with mitral-valve prolapse. It may be followed by a murmur if the valve allows blood to leak backward.

\medterm{Midtrimester} The middle part of pregnancy, usually covering about weeks 14 through 27. During this period the fetus grows rapidly, routine anatomy screening is often performed, and pregnancy symptoms and risks vary from person to person.

\medterm{Mid-Upper-Arm Circumference} The circumference of the upper arm measured at the midpoint between the acromion and
olecranon processes. It is a practical anthropometric measure of muscle and fat stores,
particularly useful when body weight or height cannot be measured reliably.

\medterm{Midwife} A trained and regulated professional who provides care during pregnancy, labor, birth, and the postnatal period.

\medterm{Midwifery} The professional practice of supporting and managing normal pregnancy, childbirth, postpartum recovery, and newborn care.

\medterm{Mifepristone} Mifepristone is a progesterone-receptor antagonist and, at higher doses, a glucocorticoid-receptor antagonist. It is used with misoprostol for medication abortion under applicable protocols and for selected cases of hyperglycemia associated with Cushing syndrome.

\medterm{Migraine} A recurring neurological disorder that causes attacks of moderate to severe, often throbbing or pulsing head pain. The pain may occur on one or both sides of the head and may be accompanied by nausea, vomiting, or increased sensitivity to light, sound, or smells. Some people have an aura, such as flashing lights, blind spots, tingling, numbness, or temporary speech problems, which usually develops gradually and resolves completely. Migraine attacks may last from several hours to a few days, and some occur with little or no headache.

\synonyms
Migraine Headache; Hemicrania

\medterm{Migraine And Stroke Symptoms} Migraine aura usually develops gradually and resolves, whereas stroke symptoms often begin suddenly and cause persistent focal deficits such as weakness, speech trouble, facial droop, or vision loss. New or atypical neurologic symptoms require emergency stroke assessment.

\medterm{Migraine With Aura} A primary headache disorder characterized by recurrent attacks of fully reversible
visual, sensory, or speech symptoms that develop gradually over at least five minutes
and last no more than sixty minutes, followed by headache with the characteristics of
migraine.

\medterm{Migraine Without Aura} A primary headache disorder characterized by recurrent attacks of moderate to severe,
pulsating, unilateral headache lasting four to seventy-two hours, associated with
nausea, vomiting, photophobia, and phonophobia.

\medterm{Migrainous Neuralgia} An older term for cluster headache, a disorder causing recurrent attacks of severe unilateral pain, usually around one eye, with ipsilateral autonomic symptoms such as tearing or nasal congestion.

\medterm{Migrating Motor Complex} Repeating waves of muscle activity that move through the stomach and small intestine between meals. They help clear leftover food, secretions, and some bacteria before the next meal.

\medterm{Migration} Movement of cells, organisms, substances, or a body structure from one place to another. Normal examples include immune-cell movement and fetal development; abnormal migration can contribute to infection, cancer spread, or blood clots.

\medterm{Migratory Thrombophlebitis} Recurrent or shifting venous thrombosis associated with an underlying malignancy; an alternate term for the thrombotic manifestation of Trousseau syndrome.

\medterm{Mikulicz Disease} An older term for persistent, usually symmetric enlargement and inflammation of the lacrimal and salivary glands, now commonly associated with IgG4-related disease.

\medterm{Mikulicz Syndrome} A historical term for symmetrical, persistent swelling of the tear and salivary glands, often associated with IgG4-related disease or other autoimmune conditions. The underlying cause should be evaluated rather than assumed.

\medterm{Mikulicz's Disease} An older term for persistent, usually symmetrical enlargement of the lacrimal and salivary glands, now often considered within IgG4-related disease or Sjögren syndrome.

\medterm{Milan System for Salivary Cytopathology} A six-level system for reporting results from a needle sample of a salivary-gland lump. Results range from an insufficient sample or noncancerous finding to a finding suspicious for or showing cancer. The system helps estimate cancer risk and guide further care.

\medterm{Mild} Mild describes relatively low intensity, severity, or extent of a symptom, sign, disease, or abnormality. The meaning depends on the condition and the clinical scale used and does not by itself indicate that no treatment or evaluation is needed.

\medterm{Mild Cognitive Impairment} A noticeable decline in memory or another thinking skill that is greater than expected for age but does not substantially prevent everyday activities. It may remain stable, improve, or progress to dementia depending on the cause.

\medterm{Mildew Remover Poisoning} Poisoning from a mildew-removal product, often containing bleach, acids, or other corrosive chemicals. It can burn the skin, eyes, mouth, and digestive tract or release dangerous fumes; leave the area and seek urgent advice.

\medterm{Mile} A mile is a unit of distance equal to 1,609.344 meters. It may be used in descriptions of walking or running distance, although scientific and clinical measurements generally use SI units.

\medterm{Milia} Milia are tiny, firm, white or yellow bumps caused by small cysts containing keratin, a protein in the skin. They are common in newborns and may appear on the face in adults; they are harmless and usually go away on their own.

\medterm{Miliaria} A sweat-retention disorder caused by blockage of eccrine ducts, producing small itchy or prickly papules.

\medterm{Miliaria Crystallina} A harmless heat rash made of tiny clear blisters when sweat becomes trapped just beneath the skin’s surface. It can occur in hot, humid conditions or during fever and usually settles when the skin is cooled and kept dry.

\medterm{Miliaria Rubra} An itchy prickly heat rash caused by sweat-duct blockage in the epidermis, producing small red papules or vesicles.

\medterm{Miliary} Describing many tiny spots or lesions that resemble millet seeds. The term is often used for tuberculosis or another disease spread widely through an organ or the bloodstream.

\medterm{Miliary TB} Tuberculosis that has spread through the bloodstream, producing many tiny areas of infection in the lungs and sometimes the liver, brain, bone marrow, or other organs. It can cause fever, weight loss, night sweats, cough, or organ-specific symptoms.

\medterm{Miliary Tuberculosis} Tuberculosis that spreads through the bloodstream and forms many tiny areas of infection. It may affect the lungs, liver, brain, bone marrow, or other organs and can cause prolonged fever, weight loss, night sweats, cough, or organ-specific symptoms.

\medterm{Milieu Therapy} Milieu therapy uses the structured social environment of a treatment setting as part of care, especially in some psychiatric and rehabilitation programs.

\medterm{Military Personnel} Military personnel are members of the armed forces, including active-duty, reserve, and related service members. Their health care may address occupational exposures, injuries, infectious risks, combat-related trauma, and deployment-related mental-health conditions.

\medterm{Milium} A small, firm, white or yellow bump filled with keratin, usually appearing on the face. It is harmless, does not spread by infection, and may disappear on its own or be removed safely by a clinician.

\medterm{Milk} A liquid made by the mammary glands to nourish an infant. Human breast milk provides water, energy, protein, fats, vitamins, and immune substances; its amount and composition change during feeding and as the infant grows.

\medterm{Milk Allergy} An immune reaction to proteins in cow's milk that may cause skin, gastrointestinal, or respiratory symptoms and, rarely, anaphylaxis.

\synonyms
Cow's Milk Protein Allergy; CMPA

\medterm{Milk Intolerance} Alternate name; see \textit{Lactose Intolerance}.

\medterm{Milk Intolerance} Alternate name; see \textit{Lactose Intolerance}.

\medterm{Milk Of Magnesia} A liquid suspension of magnesium hydroxide used to neutralize stomach acid and relieve constipation. It draws water into the bowel and can cause diarrhea, especially when taken too often or in kidney disease.

\medterm{Milk-Alkali Syndrome} A condition caused by taking too much calcium and absorbable alkali, leading to high blood calcium, excessive blood alkalinity, and kidney problems. Severe cases can cause confusion, vomiting, dehydration, or abnormal heart rhythms.

\medterm{Milker's Node} A localized nodular skin infection caused by parapoxvirus acquired from infected cattle, also called orf-like milker's nodule.

\medterm{Milk-Leg} An older term for postpartum swelling of a leg, historically associated with deep venous thrombosis or lymphatic obstruction.

\medterm{Milky Urine} Urine with a cloudy or white appearance caused by chyle, pus, crystals, fat, or other substances; evaluation identifies the cause.

\medterm{Miller's Lung} Hypersensitivity pneumonitis caused by inhaling material associated with the granary weevil (wheat weevil, \textit{Sitophilus granarius}), especially in people working with contaminated grain or flour.

\medterm{Millicurie} A millicurie (mCi) is one-thousandth of a curie, equal to approximately 37 megabecquerels. It is used to express the activity of radioactive sources and some radiopharmaceutical doses.

\medterm{Milliequivalent} A unit used to measure the amount and electrical combining power of charged substances, such as sodium, potassium, calcium, or chloride, in medicines and laboratory results.

\medterm{Milligram} A metric unit of mass equal to one-thousandth of a gram. Its abbreviation is mg, and it is commonly used for medicine doses and small quantities of substances.

\medterm{Milligram-Percent} An older concentration unit meaning milligrams of a substance in 100 milliliters of fluid. Its number is the same as milligrams per deciliter, written mg/dL, which is preferred in current medical reporting.

\medterm{Milligray} A unit of absorbed ionizing-radiation dose equal to one-thousandth of a gray, or 0.001 joule per kilogram; its abbreviation is mGy.

\medterm{Milliliter} A unit of volume equal to one-thousandth of a liter and one cubic centimeter. It is commonly used to measure liquid medicines, blood, urine, laboratory samples, and other small amounts of fluid.

\medterm{Millimeter} A unit of length equal to one-thousandth of a meter.

\medterm{Millimolar} A concentration unit equal to one thousandth of a mole per liter, written as mmol/L.

\medterm{Millimole} A unit measuring the amount of a chemical substance equal to one-thousandth of a mole. Laboratory tests use millimoles to report substances such as glucose, sodium, potassium, and other electrolytes.

\medterm{Million} A number equal to 1,000,000; laboratory counts may use it as a numerical multiplier.

\medterm{Milliosmole} A milliosmole (mOsm) is one-thousandth of an osmole, a unit reflecting the number of osmotically active particles in a solution. Osmolality and osmolarity are commonly reported in mOsm/kg or mOsm/L.

\medterm{Millipede Toxin} Irritating defensive chemicals released by some millipedes. Contact can cause burning, redness, blistering, or dark staining of skin and severe eye irritation; rinse exposed areas with water and seek care for eye symptoms.

\medterm{Millisecond} A millisecond (ms) is one-thousandth of a second. It is used to measure brief physiologic events, reaction times, nerve conduction intervals, and timing of electrical or imaging signals.

\medterm{Milliunit} A milliunit (mU) is one-thousandth of a unit of biologic or enzymatic activity. Because a unit is defined according to the substance or assay, the clinical meaning of a milliunit depends on the specific test or medication.

\medterm{Millivolt} A unit of electrical voltage equal to one-thousandth of a volt. Small voltages are measured in millivolts in tests of the heart, brain, nerves, muscles, and other electrical activity in the body.

\medterm{Milrinone} A medicine given by vein for short-term support in severe heart failure or after heart surgery. It helps the heart squeeze more strongly and widens blood vessels, but can lower blood pressure or trigger abnormal heart rhythms.

\medterm{Milroy Disease} An inherited primary lymphedema caused most often by pathogenic variants affecting
VEGFR3 signaling. It typically presents with chronic swelling of the lower limbs from
birth or early infancy and may be associated with papillomatosis, hydrocele, or
recurrent skin infection.

\medterm{Milroy's Disease} An inherited form of primary lymphedema, usually present from birth or early life and caused commonly by lymphatic-development abnormalities.

\medterm{Mimesis} Imitation or mimicry; in medicine, the term can refer to simulation of an organic disease by psychological symptoms or to one organic disease resembling another.

\medterm{Mimicry} Resemblance of one organism, structure, or clinical pattern to another; in medicine, it may describe a disorder that imitates another disease.

\medterm{Minamata Disease} A nervous-system disorder caused by poisoning with methylmercury, often after eating contaminated seafood. It can cause numbness, poor coordination, vision or hearing problems, muscle weakness, speech difficulty, and developmental harm in children.

\medterm{Mind} The collection of mental processes and capacities involved in consciousness, thought, perception, memory, emotion, judgment, and behavior.

\medterm{Mindful Eating} Paying attention to hunger and fullness, eating slowly, and noticing food's taste and texture. It may improve awareness of eating habits and help some people reduce distracted eating or overeating.

\medterm{Mineral} An inorganic element or compound required for body structure, fluid balance, enzyme activity, or other physiological functions.

\medterm{Mineral Oil} Mineral oil is a purified mixture of liquid hydrocarbons derived from petroleum. It is used in some topical products, cosmetics, and laxatives; aspiration of liquid mineral oil can cause lipoid pneumonia.

\medterm{Mineral Oil Overdose} Swallowing a large amount of mineral oil is dangerous mainly because it can enter the lungs during vomiting or reflux and cause lipoid pneumonia. Do not induce vomiting; obtain product-specific poison-control advice.

\medterm{Mineral Spirits Poisoning} Poisoning from mineral spirits, petroleum-based solvents used in paints and cleaning products. Swallowing or inhalation may cause dizziness, drowsiness, chemical pneumonitis, or breathing failure, especially after aspiration; do not induce vomiting.

\medterm{Mineralocorticoid} A steroid hormone, mainly aldosterone, that helps the kidneys retain sodium and water while removing potassium. Too much can raise blood pressure and lower potassium; too little can cause dehydration and weakness.

\medterm{Mineralocorticoid Receptor} A protein inside cells that responds mainly to aldosterone, a hormone made by the adrenal glands. It helps the kidneys retain salt and water and remove potassium, thereby affecting blood pressure and body-fluid balance.

\medterm{Mineralocorticoids} Steroid hormones, chiefly aldosterone, that promote sodium retention and potassium and hydrogen secretion in the kidney.

\medterm{Minerals} Naturally occurring elements needed by the body for bones, nerves, muscles, blood formation, fluid balance, and other functions. They include calcium, iron, magnesium, potassium, sodium, and zinc; both deficiency and excess can be harmful.

\medterm{Miner's Anaemia} Anemia associated historically with mining, commonly due to hookworm infection, chronic blood loss, nutritional deficiency, or occupational exposure.

\medterm{Miner's Asthma} Asthma caused or worsened by repeated inhalation of dust, fumes, or other substances during mining work. It can cause wheezing, cough, chest tightness, and breathlessness that improve after leaving the exposure.

\medterm{Miner's Nystagmus} An older term for visual disturbance and nystagmus reported in underground miners, associated with poor lighting and occupational conditions.

\medterm{Minimal Change Disease} A kidney disorder, especially common in children, in which the kidney filters look nearly normal on routine examination but leak large amounts of protein. It can cause swelling around the eyes, legs, or abdomen and is often treatable.

\medterm{Minimal Residual Disease} A small number of cancer cells that remain after treatment and cannot be seen on routine examination or ordinary scans. Highly sensitive blood, bone-marrow, tissue, or molecular tests may detect them and estimate relapse risk.

\medterm{Minimally Invasive Gastrointestinal Surgery} Gastrointestinal operations performed through small incisions or natural orifices using laparoscopic, robotic, or endoscopic techniques, often reducing recovery time while retaining procedure-specific risks.

\medterm{Minimally Invasive Hip Replacement} Hip-replacement surgery performed through a smaller or muscle-sparing approach to reduce tissue disruption. It still replaces damaged joint surfaces and has risks such as infection, blood clots, dislocation, fracture, nerve injury, or implant loosening.

\medterm{Minimally Invasive Mitral Valve Surgery} Surgical mitral valve repair or replacement performed through a right mini-thoracotomy or video-assisted thoracoscopic approach, avoiding a traditional median sternotomy. Results in reduced surgical trauma, less postoperative pain, and faster recovery.

\synonyms
MIMVS; Mini-Thoracotomy Mitral Repair; Port-Access Mitral Surgery

\medterm{Minimally Invasive Spine Surgery} Spine surgery performed through small incisions using specialized instruments, often reducing muscle disruption and recovery time. It still carries risks such as nerve injury, infection, bleeding, and persistent symptoms.

\medterm{Minimally Invasive Surgery} Surgery performed through small incisions or natural openings using specialized instruments, usually reducing tissue disruption and recovery time.

\medterm{Mini-Mental State Examination} A brief questionnaire that checks orientation, attention, memory, language, and simple thinking skills. The score can be affected by education, language, hearing, vision, culture, or delirium and cannot diagnose dementia by itself.

\medterm{Minimization} A psychological or communication process of downplaying the importance or severity of symptoms, events, or problems.

\medterm{Minimum Alveolar Concentration} A standard measure of how much inhaled anesthetic is needed to prevent movement in half of people during a specified surgical stimulus. It helps compare anesthetic strength and guide dosing during anesthesia.

\medterm{Minimum Bactericidal Concentration} The lowest concentration of an antimicrobial that kills a defined proportion of the
original bacterial inoculum, commonly demonstrated by subculture without growth. MBC is
distinct from MIC and is not routinely required for every infection.

\medterm{Minimum Effect Concentration} The lowest concentration of a drug in the blood or at its site of action that produces
a defined measurable effect.

\medterm{Minimum Inhibitory Concentration} The smallest amount of an antibiotic or other antimicrobial that stops a particular germ from visibly growing in a laboratory test. It helps guide treatment but must be interpreted with the infection and the patient's condition.

\medterm{Minimum Lethal Concentration} The lowest concentration of a substance associated with death in a defined test
population or experimental setting.

\medterm{Minimum Lethal Dose} The smallest dose of a substance reported to cause death under specified conditions; it varies by species, route, and individual.

\medterm{Minisatellite Repeat} A minisatellite repeat is a tandemly repeated DNA sequence with repeat units commonly ranging from about 10 to 100 base pairs. Variation in repeat number can serve as a genetic marker and, in some loci, influence gene regulation or disease risk.

\medterm{Mini-Stroke} Alternate name; see \textit{Transient Ischemic Attack}.

\medterm{Mini-Stroke} A lay term usually referring to a transient ischemic attack, in which focal neurologic symptoms resolve without persistent infarction.

\medterm{MINOCA} A heart attack occurring without a major blockage in the coronary arteries. It is an initial diagnosis; further tests are needed to find causes such as a small clot, artery spasm, artery-wall tear, or inflammation.

\medterm{Minocycline} An antibiotic used for some bacterial infections and for inflammatory acne. It can cause nausea, dizziness, skin or tooth discoloration, sun sensitivity, and rarely liver or immune problems; the course should be taken as prescribed.

\medterm{Minor} Smaller, lesser, or less significant than a related structure or condition; for example, the teres minor muscle is smaller than the teres major muscle.

\medterm{Minor Response} A minor response is a limited improvement after treatment that does not meet the predefined criteria for a partial or complete response. The exact threshold depends on the disease, study protocol, and outcome being measured.

\medterm{Minority Health} The health status and healthcare experiences of populations affected by social, economic, cultural, or structural disadvantage. Differences in exposure, income, environment, discrimination, and access to care can contribute to health disparities.

\medterm{Minoxidil} A vasodilator used topically for pattern hair loss and orally in selected cases of resistant hypertension.

\medterm{Minute Ventilation} The total amount of air breathed in one minute, calculated from breath size multiplied by breathing rate. It may rise with exercise or illness and fall with sedation, muscle weakness, airway disease, or poor breathing effort.

\medterm{Miosis} Unusually small pupils or normal pupil narrowing in response to light. It may be caused by medicines or toxins, eye disease, or problems affecting nerves in the brain or eye.

\medterm{Miosis} Abnormally small pupils. It may result from bright light, medicines such as opioids, eye disease, nerve disorders, or poisoning.

\medterm{Miosis Disorder} Abnormally small pupils or an abnormal pupillary constriction response, caused by medications, neurologic disease, ocular disease, or exposure to certain toxins.

\medterm{Miotic} A medicine or substance that makes the pupil smaller. Miotics may be used in glaucoma or during some eye procedures and can affect vision in dim light.

\synonyms
Miotic Agent; Pupillary Constrictor

\medterm{Miracidium} The ciliated larval stage of a trematode that hatches from an egg and infects a suitable snail host.

\medterm{Mirage Syndrome} A rare genetic disorder, usually related to SAMD9 variants, that can cause growth restriction, adrenal insufficiency, recurrent infections, and distinctive developmental findings.

\medterm{MIRD Schema} A method used to estimate how much radiation reaches organs and tissues after a radioactive medicine is given. It helps specialists plan treatment and assess safety for radiopharmaceutical procedures.

\medterm{Mirizzi Syndrome} A gallstone stuck near the gallbladder outlet presses on a nearby bile duct and blocks the flow of bile. It can cause jaundice, abdominal pain, fever, or infection and may require an endoscopic or surgical procedure.

\medterm{Mirror Therapy} A rehabilitation exercise in which a mirror makes movement of the stronger limb appear to come from the affected limb. It may help some people with phantom-limb pain or movement recovery after stroke.

\medterm{Mirror-Writing} Writing in which letters or words are reversed as if seen in a mirror; it may occur transiently in childhood or with neurologic disease.

\medterm{Mirtazapine} Mirtazapine is a noradrenergic and specific serotonergic antidepressant used for major depressive disorder. It commonly causes sedation and increased appetite or weight and can rarely cause neutropenia or other serious adverse effects.

\medterm{Misalignment} Misalignment is abnormal positioning or orientation of two or more structures relative to their intended relationship. It may affect bones, joints, teeth, eyes, implants, or medical devices and can impair function or healing.

\medterm{Miscarriage} Spontaneous loss of a pregnancy before fetal viability, commonly caused by chromosomal abnormalities but requiring clinical assessment.

\synonyms
Spontaneous Abortion; Early Pregnancy Loss

\medterm{Misdiagnosis} Misdiagnosis is assignment of an incorrect diagnosis or failure to recognize the correct condition. It can lead to inappropriate treatment, delayed care, avoidable harm, or unnecessary testing and may result from atypical presentation, incomplete information, or system factors.

\medterm{Miserere} An old term for vomiting that smells or looks like bowel contents because of severe blockage of the intestine. This is a medical emergency that may cause dehydration, infection, or bowel damage and needs urgent treatment.

\medterm{Mismatch Repair} A system that checks newly copied DNA and corrects small errors. When it fails, genetic changes can accumulate and increase the risk of cancers such as colorectal and womb cancer.

\medterm{Mismatch Repair Deficiency} Loss of DNA mismatch-repair function that produces microsatellite instability and a high
tumor mutational burden. Seen in Lynch syndrome and some sporadic colorectal,
endometrial, and other cancers, and associated with responsiveness to immune checkpoint
inhibitors.

\medterm{MISME} Abbreviation; see \textit{Neurofibromatosis Type 2}.

\medterm{Misophonia} Misophonia is an intense emotional or physiologic reaction to specific repetitive sounds, such as chewing or tapping. It is not simple annoyance; behavioral strategies, sound management, and therapy can reduce impairment, while assessment considers anxiety, obsessive-compulsive symptoms, and sensory sensitivity.

\medterm{Misoprostol} A prostaglandin E1 analogue used for ulcer prevention in selected patients, medical abortion, cervical preparation, and obstetric indications.

\medterm{Misoprostol} A synthetic prostaglandin analogue used for selected gynecologic and obstetric indications, including cervical ripening,
labor induction, and medical management of some pregnancies. It can cause uterine
contractions, bleeding, gastrointestinal effects, or uterine tachysystole; use depends on
the clinical setting and gestational age.

\synonyms
Prostaglandin E1 Analogue

\medterm{Missed Dose} A missed dose is a scheduled dose of a medicine that was not taken or administered. The appropriate response depends on the drug and timing; doubling a dose can be dangerous, so the product instructions or a clinician or pharmacist should guide the next dose.

\medterm{Missed Miscarriage} A nonviable intrauterine pregnancy that has not been expelled and may produce few or no symptoms of miscarriage.
Diagnosis requires appropriate ultrasound or other clinical evaluation before management is selected.

\synonyms
Missed Abortion; Silent Miscarriage; Early Pregnancy Failure

\medterm{Missense Mutation} A missense mutation is a DNA sequence change that substitutes one amino acid for another in a protein. Its effect may be harmful, beneficial, or neutral depending on the affected residue, protein function, and genetic context.

\medterm{Missense Variant} A single-nucleotide substitution that changes one amino acid in a protein. Its clinical
effect ranges from benign to pathogenic and must be assessed using molecular,
functional, population, and clinical evidence.

\medterm{Mistletoe Poisoning} Poisoning after eating mistletoe leaves or berries. It may cause nausea, vomiting, diarrhea, abdominal pain, dizziness, low blood pressure, slow or irregular heartbeat, seizures, or confusion; children and symptomatic people need urgent assessment.

\medterm{Mite} A tiny arachnid; some mites cause dermatitis or transmit disease, and house-dust mites can trigger allergy.

\medterm{Mites} Tiny arthropods related to ticks and spiders. Some live harmlessly on skin, while others cause allergy, scabies, or irritation; treatment depends on the species and condition.

\medterm{Mithramycin} An older antineoplastic antibiotic, also called plicamycin, whose use is limited by substantial toxicity.

\medterm{Mitobronitol} Mitobronitol is an alkylating antineoplastic compound formerly investigated or used in selected hematologic malignancies. Its cytotoxic effects can damage DNA and cause myelosuppression and other treatment-related toxicities.

\medterm{Mitochondria} Small structures inside cells that make much of the energy cells need. They also help control calcium, cell survival, and several metabolic processes; mitochondrial disorders can affect organs that need high energy, such as the brain, muscles, heart, and eyes.

\medterm{Mitochondrial Cristae} Folds of the inner membrane of mitochondria, the energy-producing parts of cells. The folds provide space for the machinery that turns nutrients into usable cellular energy.

\medterm{Mitochondrial Diabetes} Diabetes caused by a change in mitochondrial DNA, which is usually passed through the mother. It may occur with hearing loss, muscle problems, or other mitochondrial symptoms; insulin production is often reduced and genetic testing may confirm the cause.

\medterm{Mitochondrial Disease} A group of genetic disorders caused by impaired mitochondrial energy production. Symptoms vary and
may affect the brain, muscles, nerves, heart, eyes, or multiple organs; onset and severity can
differ widely even within one family.

\medterm{Mitochondrial DNA} Small circular DNA molecules located primarily in mitochondria that encode some components of oxidative phosphorylation and are inherited predominantly through the maternal line.

\medterm{Mitochondrial DNA Depletion Syndrome MDS MDDS} A group of inherited disorders in which some tissues contain too little mitochondrial DNA to make energy normally. Depending on the form, the brain, muscles, liver, or several organs may be affected, often beginning in infancy or childhood.

\medterm{Mitochondrial Fission} Mitochondrial fission is division of one mitochondrion into two or more mitochondria. It helps distribute mitochondria, remove damaged portions through mitophagy, and respond to cell-cycle or metabolic signals; excessive fission may contribute to disease.

\medterm{Mitochondrial Fusion} Mitochondrial fusion is joining of mitochondria into a connected network. It permits exchange of mitochondrial contents and helps maintain bioenergetic function and genome integrity; fusion and fission are normally balanced dynamically.

\medterm{Mitochondrial Inheritance} Inheritance of mitochondrial DNA, which is transmitted predominantly through the egg and therefore usually
from the mother. Examples include Leber hereditary optic neuropathy and MELAS. Heteroplasmy
means that mutant and normal mitochondrial DNA coexist in a cell.

\medterm{Mitochondrial Matrix} The mitochondrial matrix is the protein-rich compartment enclosed by the inner mitochondrial membrane. It contains mitochondrial DNA, ribosomes, and enzymes for the citric-acid cycle, fatty-acid oxidation, and other metabolic processes.

\medterm{Mitochondrial Membrane} The mitochondrial membrane system consists of an outer membrane and a highly folded inner membrane separated by an intermembrane space. It controls transport and houses the respiratory chain and ATP-producing machinery.

\medterm{Mitochondrial Myopathy} A muscle disorder caused by faulty mitochondria, the cell structures that make energy. It may cause weakness, exercise intolerance, drooping eyelids, hearing loss, seizures, or problems in other organs.

\medterm{Mitochondrial Ribosome} A mitochondrial ribosome is a ribosome within mitochondria that translates proteins encoded by mitochondrial DNA. Its structure differs from cytosolic ribosomes and defects in mitochondrial translation can cause energy-production disorders.

\medterm{Mitochondrial Swelling} Enlargement of mitochondria caused by disruption of ion and water balance, often indicating cellular injury, impaired oxidative phosphorylation, or permeability-transition changes.

\medterm{Mitochondrion} A mitochondrion is a membrane-bound organelle that generates much of a cell’s ATP through oxidative phosphorylation. It also participates in apoptosis, calcium handling, metabolism, and production or regulation of reactive oxygen species.

\medterm{Mitogen} A mitogen is a substance that stimulates a cell, especially a lymphocyte, to enter the cell cycle and proliferate. Mitogens include growth factors, lectins, and other signals that activate proliferative pathways.

\medterm{Mitolactol} Mitolactol is an alkylating antineoplastic compound formerly used or investigated in cancer treatment. It damages DNA and can cause cytotoxic effects such as myelosuppression and gastrointestinal toxicity.

\medterm{Mitomycin} Mitomycin is a chemotherapy medicine that damages DNA and is used for selected cancers; it can suppress bone marrow and injure the kidneys or bladder.

\medterm{Mitomycins} Mitomycins are antitumor antibiotics that become activated in cells and cross-link DNA, inhibiting replication and transcription. Mitomycin C is the best-known member and is used in selected cancers and ophthalmic or surgical applications.

\medterm{Mitophagy} The process by which a cell identifies and breaks down damaged or unnecessary mitochondria, the cell's energy-producing structures. This helps maintain healthy cells; abnormal mitophagy may contribute to some diseases.

\medterm{Mitosis} Cell division that produces two genetically similar daughter cells with the same chromosome number as the parent
cell.

\medterm{Mitosporic Fungi} Fungi that reproduce by asexual spores; the term is mainly used for fungi whose sexual stage is unknown or not observed.

\medterm{Mitotane} Mitotane is an adrenolytic drug used for selected adrenocortical carcinomas and to suppress excess adrenal steroid production. It can cause gastrointestinal and neurologic effects, adrenal insufficiency, and numerous drug interactions.

\medterm{Mitotic Anaphase} Mitotic anaphase is the stage of mitosis in which sister chromatids separate and move toward opposite poles of the cell. Accurate spindle attachment and chromosome segregation are required for formation of genetically similar daughter cells.

\medterm{Mitotic Figure} A cell seen under a microscope while it is dividing. The number of dividing cells in a tissue sample can help estimate how quickly a tumor or other tissue is growing.

\medterm{Mitotic Metaphase} Mitotic metaphase is the stage of mitosis in which duplicated chromosomes align at the cell equator with sister kinetochores attached to opposite spindle poles. A checkpoint helps prevent premature chromosome separation.

\medterm{Mitotic Prophase} Mitotic prophase is the early stage of mitosis in which chromatin condenses into visible chromosomes, the nucleolus disappears, and the mitotic spindle begins to form. In late prophase or prometaphase, the nuclear envelope breaks down.

\medterm{Mitotic Recombination} Mitotic recombination is exchange or rearrangement of DNA between homologous chromosomes during or around mitotic cell division. It can repair DNA or generate loss of heterozygosity and is relevant to tumor evolution and genetic disease.

\medterm{Mitotic Telophase} Mitotic telophase is the stage after chromosome separation when chromosomes reach opposite poles, decondense, and become enclosed in newly forming nuclear envelopes. Cytokinesis usually follows or overlaps with telophase.

\medterm{Mitoxantrone} A chemotherapy and immunosuppressive medicine that interferes with DNA and is used for selected cancers and some forms of multiple sclerosis.

\medterm{MitraClip} A transcatheter device that creates a double orifice in the mitral valve by clipping the
anterior and posterior leaflets together, reducing mitral regurgitation. It is indicated
for patients with significant (3+ or 4+) mitral regurgitation who are at high risk for
surgical mitral valve repair.

\medterm{Mitral} Relating to the heart's mitral valve, which lies between the left upper and lower chambers. The valve opens for blood flow forward and closes to prevent backward leakage during each heartbeat.

\medterm{Mitral Atresia} A serious heart defect present from birth in which the opening of the mitral valve is absent or completely closed. Blood cannot flow normally into the heart’s main pumping chamber, so newborns need urgent specialist care.

\medterm{Mitral Facies} A distinctive physical appearance characterized by a purplish or violaceous malar flush and patchy cyanosis across the cheeks, bridge of the nose, and lips. It is caused by chronic pulmonary hypertension, systemic vasoconstriction, and reduced cardiac output in severe mitral stenosis.

\medterm{Mitral Insufficiency} Failure of the mitral valve to close completely, allowing blood to leak back into the left atrium during systole.

\medterm{Mitral Regurgitation} Leakage through the mitral valve when the heart contracts, allowing blood to flow backward into the left upper chamber. It may cause breathlessness, tiredness, palpitations, or heart failure and can result from valve disease, infection, or heart damage.

\medterm{Mitral Stenosis} Narrowing of the mitral valve opening, making it harder for blood to move from the left upper chamber to the left lower chamber. It can cause breathlessness, tiredness, palpitations, or swelling and is often related to previous rheumatic heart disease.

\medterm{Mitral Stenosis Murmur} A low rumbling heart sound heard when blood passes through a narrowed mitral valve. It may be accompanied by a sharp opening sound and is best assessed with a stethoscope and heart imaging.

\medterm{Mitral Valve} The two-flap valve between the heart’s left upper and lower chambers. It opens to let blood pass forward and closes to prevent backward flow; it may become narrowed, leaky, infected, or weakened.

\medterm{Mitral Valve Insufficiency} Leakage of the mitral valve that allows blood to flow backward into the left atrium when the heart contracts. It may cause breathlessness, fatigue, palpitations, swelling, or heart failure.

\medterm{Mitral Valve Prolapse} A condition in which the flaps of the mitral valve bulge backward during a heartbeat. Many people have no symptoms, but some develop palpitations, chest discomfort, breathlessness, or leakage through the valve.

\synonyms
MVP; Barlow Syndrome; Click-Murmur Syndrome; Floppy Mitral Valve

\medterm{Mitral Valve Regurgitation} Backward leakage of blood through the mitral valve when the left ventricle contracts, which can cause breathlessness, fatigue, atrial enlargement, or heart failure.

\synonyms
MR; Mitral Incompetence; Mitral Insufficiency

\medterm{Mitral Valve Stenosis} Narrowing of the mitral valve that restricts blood flow from the left atrium to the left ventricle.

\synonyms
MS; Mitral Stenosis

\medterm{Mitral Valvotomy} A procedure that opens a narrowed mitral valve so blood can move more easily from the upper to the lower left heart chamber. It may use a balloon catheter or surgery and is suitable only for selected valve anatomy.

\medterm{Mitral Valvuloplasty} Mitral valvuloplasty is a procedure that widens a narrowed mitral valve, often with a balloon catheter or surgical repair.

\medterm{Mitral-Valve Repair} A surgical or transcatheter procedure to restore mitral valve competence without
replacement, typically for regurgitation from prolapse, chordal rupture, or ischemic
dysfunction. Repair preserves native valve tissue, maintains left ventricular function,
and avoids lifelong anticoagulation required by mechanical prostheses.

\medterm{Mitsuokella} Mitsuokella is a genus of anaerobic bacteria found mainly in the mouth and digestive tract; human infection is uncommon.

\medterm{Mittelschmerz} Recurrent one-sided lower-abdominal or pelvic pain occurring around ovulation. It is usually brief or lasts up to a day or two, but severe, persistent, or atypical pain requires assessment for other causes.

\medterm{Mixed Connective Tissue Disease} An autoimmune illness with features of lupus, scleroderma, and inflammatory muscle disease. It may cause color changes in fingers with cold, swollen hands, joint pain, muscle weakness, swallowing problems, or inflammation around organs.

\synonyms
MCTD; Sharp Syndrome

\medterm{Mixed Glioma} Mixed glioma is an older term for a glial tumor showing features of more than one glial lineage, historically including oligoastrocytoma. Modern classification relies mainly on molecular findings and generally assigns such tumors to a more specific entity.

\medterm{Mixed Hearing Loss} Hearing loss caused by both a problem carrying sound through the outer or middle ear and a problem in the inner ear or hearing nerve. The best treatment depends on the causes found.

\medterm{Mixed Oedema} Edema caused by more than one mechanism or affecting more than one tissue compartment; interpretation depends on clinical context.

\medterm{Mixed Urinary Incontinence} Urine leakage caused by both a sudden strong urge to urinate and pressure from coughing, sneezing, laughing, or exercise. Treatment may include bladder training, pelvic-floor exercises, medicines, or procedures.

\medterm{Mixed-Lineage Leukemia} An acute leukemia in which the malignant blasts show defining features of more than one hematopoietic lineage, commonly both myeloid and lymphoid; it is also called mixed-phenotype acute leukemia.

\medterm{Miyoshi Myopathy} A rare inherited muscle disease that usually begins in adolescence or early adulthood with progressive weakness of the calf muscles. People may have difficulty standing on tiptoe or climbing stairs; weakness later spreads to the thighs and hips. Blood tests often show high muscle-enzyme levels. Physical therapy helps maintain movement and function.

\synonyms
Miyoshi Distal Myopathy; Distal Dysferlinopathy

\medterm{MME} Membrane metalloendopeptidase, also called neprilysin, CD10, neutral endopeptidase, or CALLA; it is a cell-surface enzyme and marker expressed in several tissues and in some acute lymphoblastic leukemias.

\medterm{MéNièRe's Disease} An inner-ear disorder causing repeated attacks of spinning dizziness, changing hearing loss, ringing, and pressure or fullness in one ear. Symptoms may last minutes to hours and can affect balance and daily life.

\medterm{MéNièRe’S Disease} An inner-ear disorder causing repeated attacks of spinning dizziness, changing hearing loss, ringing, and pressure or fullness in one ear. Symptoms may last minutes to hours and can affect balance and daily life.

\synonyms
Meniere Disease; Endolymphatic Hydrops

\medterm{MéNéTrier’S Disease} A rare stomach disorder in which the stomach lining becomes unusually thick and loses protein into the digestive tract. It may cause upper-abdominal pain, nausea, vomiting, weight loss, swelling, or low blood protein.

\medterm{Mobile} Capable of moving or being moved; a mass or structure described as mobile is not fixed to adjacent tissue.

\medterm{Mobile Genetic Elements} Genetic elements that can move within and between genomes. These include plasmids,
transposons, and integrons. Mobile genetic elements can carry antibiotic resistance
genes and virulence factors. They facilitate horizontal gene transfer and contribute to
the rapid evolution of bacteria.

\medterm{Mobile Health} The use of phones, tablets, wearable devices, or similar portable technology to support health care. Uses include reminders, remote monitoring, education, communication, symptom tracking, and some diagnostic tools; privacy, accuracy, accessibility, and clinical oversight remain important.
\synonyms
mHealth; Digital Health

\medterm{Mobile Health Unit} A mobile health unit is a vehicle or transportable facility equipped to provide health services at changing locations. It may deliver screening, vaccination, primary care, diagnostic testing, emergency support, or outreach to underserved communities.

\medterm{Mobility} The ability to move or be moved, including functional movement of a person or range of motion of a joint.

\medterm{Mobilization} The process of restoring or encouraging movement, such as early patient activity, joint movement, or surgical release of tissue.

\medterm{Mobiluncus} A genus of curved anaerobic bacteria associated with bacterial vaginosis and occasionally isolated from other clinical specimens.

\medterm{Mobiluncus Curtisii} A curved bacterium that grows without oxygen and may be present with bacterial vaginosis. Its detection can support assessment, but symptoms, examination, and other test results are needed to interpret vaginal findings.

\medterm{Mobiluncus Mulieris} A curved bacterium that grows without oxygen and may be associated with bacterial vaginosis. Finding it can support a diagnosis, but symptoms, examination, and other laboratory results are also important.

\medterm{Mobius Syndrome} A rare condition present from birth that affects certain facial and eye-moving nerves. It can cause facial weakness or paralysis, trouble closing the eyes, limited outward eye movement, and feeding or speech difficulties.

\medterm{Modafinil} A prescription medicine that promotes wakefulness in narcolepsy, shift-work sleep disorder, and persistent sleepiness from obstructive sleep apnea when airway treatment is adequate. It can cause headache, anxiety, insomnia, serious rash, or dependence and is not a substitute for sleep or airway treatment.

\medterm{Modality} A particular method used for treatment, diagnosis, imaging, or sensory stimulation. Examples include surgery, medicines, ultrasound, MRI, physiotherapy, and forms of counseling; the best modality depends on the health problem and the person's needs.

\medterm{Mode} The most frequent value in a dataset or a particular manner or method of operation.

\medterm{Model} A simplified representation used to describe, explain, predict, or study a biological or clinical process.

\medterm{Moderate To Severe Disabilities} Moderate to severe disability describes substantial limitations in physical, cognitive, sensory, communication, or adaptive functioning that require ongoing support. Needs vary widely and may include medical care, rehabilitation, assistive technology, education, and assistance with daily activities.

\medterm{Moderation} Moderation is reduction or control of the intensity, amount, or severity of a process, symptom, behavior, or treatment effect. In clinical research, a moderator is a factor that changes the strength or direction of an association.

\medterm{Moderator} A factor that changes the strength or direction of a relationship or treatment effect; it may also mean a person who guides a discussion.

\medterm{Modes Of Inheritance} Patterns by which genetic traits pass through families, including autosomal dominant, autosomal recessive, X-linked, mitochondrial, and other patterns.

\medterm{Modification} A change made to a structure, process, treatment, or behavior.

\medterm{Modified Biophysical Profile} A pregnancy assessment that combines a fetal heart-rate test with measurement of the amniotic-fluid volume. A nonreactive heart-rate result or low fluid level may suggest that further assessment or delivery is needed, depending on gestational age and the clinical situation.

\synonyms
Modified BPP; NST with Amniotic Fluid Index

\medterm{Modulator} A substance or factor that alters the activity of a biological pathway, receptor, gene, or physiological process.

\medterm{Modus Operandi} The characteristic method by which a person performs an action; it is chiefly a forensic or legal term.

\medterm{Modus Vivendi} A practical arrangement or way of living that allows people or groups with differences to coexist.

\medterm{Moebius Syndrome} A congenital disorder characterized mainly by facial paralysis and inability to abduct the eyes, caused by developmental cranial-nerve abnormalities.

\medterm{Moellerella} A rare group of bacteria that can occasionally cause urinary, bloodstream, or other infection, mainly in people who are seriously ill or have weakened defenses. Laboratory identification is needed to determine whether a specimen represents true infection.

\medterm{Moellerella Wisconsensis} A rare type of bacterium that can occasionally cause urinary, bloodstream, wound, or other infections, especially in people with serious illness or weakened immunity. Laboratory testing identifies it and guides antibiotic choice.

\medterm{Moeller's Glossitis} An older term for a smooth, painful, inflamed tongue associated especially with nutritional deficiency, including iron, folate, or vitamin B12 deficiency.

\medterm{MOG Antibody-Associated Disorder} Alternate name; see \textit{MOG Antibody Disease}.

\medterm{MOGAD} Abbreviation; see \textit{MOG Antibody Disease}.

\medterm{Mogibacterium} A genus of anaerobic bacteria found in the mouth and intestine and occasionally associated with dental or other infection.

\medterm{Mohs Micrographic Surgery} A skin-cancer operation that removes the visible tumor a thin layer at a time and examines each layer under a microscope until no cancer remains. It preserves as much healthy tissue as possible; bleeding, infection, scarring, or recurrence can occur.

\medterm{Mohs Surgery} A micrographically controlled technique in which skin cancer is excised in thin layers
and examined margin-by-margin, achieving high cure rates for basal cell and squamous
cell carcinoma while sparing normal tissue.

\medterm{Moisture} Water or other liquid present on a surface or in tissue; excessive moisture can cause maceration and skin breakdown.

\medterm{Molar} Relating to a back tooth used for chewing, or to a chemical concentration measured in moles per liter.

\medterm{Molar Pregnancy} A type of gestational trophoblastic disease characterized by abnormal proliferation of
trophoblastic tissue. See Hydatidiform Mole.

\medterm{Molar Pregnancy} An abnormal pregnancy in which placental tissue grows excessively instead of developing normally. In a complete mole no embryo forms; in a partial mole abnormal fetal tissue may be present. It can cause bleeding, severe nausea, or unusually high pregnancy-hormone levels and requires specialist follow-up.

\synonyms
Hydatid Mole; Gestational Trophoblastic Disease

\medterm{Molasses} A viscous syrup produced during sugar refining; it is rich in sugar and may contain variable amounts of minerals.

\medterm{Mold} A multicellular fungus that grows as branching filaments; molds can cause allergy, food spoilage, or opportunistic infection.

\medterm{Mold Allergy} An allergic response to mold spores that may cause nasal, eye, skin, or lower-respiratory symptoms.

\synonyms
Fungal Hypersensitivity; Mold-Induced Rhinitis

\medterm{Mold Exposure} Mold exposure can irritate the eyes, skin, and airways and may worsen asthma or allergic rhinitis; invasive infection is uncommon and mainly affects severely immunocompromised people. Reducing moisture and remediating contaminated areas are more useful than unproven detoxification.

\medterm{Mole} A common usually harmless skin growth made of pigment-producing cells. It may be flat or raised and brown, black, pink, or skin-colored. A new or changing mole, or one unlike the others, should be checked for skin cancer.

\medterm{Molecular} Relating to molecules or processes occurring at the molecular level.

\medterm{Molecular Chaperone} A molecular chaperone is a protein that assists other proteins with folding, assembly, transport, or refolding without becoming part of the final structure. Chaperones can also help prevent or resolve protein aggregation.

\medterm{Molecular Diagnosis} Molecular diagnosis uses analysis of DNA, RNA, proteins, or other molecular markers to identify a disease, pathogen, inherited variant, tumor subtype, or treatment-relevant characteristic. Results require appropriate assay validation and clinical interpretation.

\medterm{Molecular Dynamics} A computer method that predicts how atoms and molecules move and interact over time. Researchers use it to study protein shape, drug binding, chemical reactions, and other biological processes; its results depend on the model and assumptions used.

\medterm{Molecular Imaging} Imaging that shows biological molecules, cell activity, or biochemical processes in the body, often using a targeted radioactive or fluorescent tracer. It can help characterize tumors, inflammation, infection, or organ function.

\medterm{Molecular Pathology} The study of disease-related changes in genes, gene activity, or proteins within cells and tissues. It can help identify a disease, classify a tumor, estimate likely outcome, detect inherited risk, or select a treatment.

\medterm{Molecular Pathology Of Colorectal Cancer} The assessment of molecular alterations that classify colorectal tumors and guide prognosis or therapy, including mismatch-repair deficiency or microsatellite instability, RAS mutations, BRAF mutation, HER2 amplification, and selected gene fusions.

\medterm{Molecular Probe} A molecular probe is a labeled or otherwise detectable molecule designed to bind a particular nucleic-acid sequence, protein, metabolite, or other target. Probes are used in imaging, assays, microscopy, and molecular diagnostics.

\medterm{Molecular Response} A molecular response is a measurable change in a molecular marker, such as a pathogen load, oncogenic transcript, mutation burden, or biomarker, after disease progression or treatment. Its definition and clinical significance depend on the disease and assay.

\medterm{Molecular Testing} Analysis of DNA, RNA, or other molecular features to identify variants, fusions, expression patterns, or infectious agents for diagnosis, prognosis, or treatment selection.

\medterm{Molecular Weight} The mass of one molecule, usually expressed in atomic mass units; for a substance used in chemistry, its numerical value corresponds to molar mass in grams per mole. It helps calculate concentrations and doses.

\medterm{Molecule} A group of chemically bonded atoms forming the smallest unit of a substance that retains its chemical properties.

\medterm{Moles} Moles are common collections of pigment-producing melanocytes and are usually benign. A lesion that changes in size, shape, color, symmetry, or symptoms, or differs markedly from others, should be examined for melanoma.

\medterm{Molimina} The symptoms and physical changes that occur before menstruation, including breast tenderness, bloating, or mood changes.

\medterm{Mollusc} A soft-bodied invertebrate; the term is also used in names of mollusc-borne or mollusc-associated parasitic diseases.

\medterm{Mollusca} The animal group containing snails, slugs, clams, and related species. Some mollusks can cause illness through venom, contaminated food, or parasites they carry; medical risk depends on the species and exposure.

\medterm{Molluscipoxvirus} A group of poxviruses that includes the virus causing molluscum contagiosum. Infection produces small, firm, usually painless skin bumps with a central dimple and spreads through close skin contact or shared items.

\medterm{Molluscum Contagiosum} A common skin infection caused by a poxvirus. It produces small, firm, painless, flesh-colored bumps with a central dip and spreads through close skin contact or shared objects. Most cases clear without treatment, although this may take months; scratching can spread the bumps or cause infection.

\medterm{Molten Metal Fumes} Fumes produced during welding, smelting, or casting that can cause metal-fume fever, airway irritation, or metal toxicity.

\medterm{Molteno Implant} A non-valved glaucoma drainage device (aqueous shunt) surgically implanted to reduce elevated intraocular pressure (IOP) in refractory and complicated glaucomas. It consists of a flexible silicone tube inserted into the anterior chamber (or pars plana) that shunts aqueous humor to one or two thin episcleral polypropylene plates sutured to the sclera beneath Tenon's capsule in the equatorial region, where fluid passively diffuses through a fibrous capsule into the orbital vascular and lymphatic circulation. Because the device is non-valved and offers no intrinsic resistance to flow, the tube is temporarily occluded at surgery (using an absorbable ligature or intraluminal ripcord suture) during the first 4--6 weeks until fibrous bleb encapsulation develops, thereby preventing early postoperative hypotony, anterior chamber collapse, and choroidal detachment. It is indicated when conventional filtration surgery (trabeculectomy) has failed or has a high probability of failure, including neovascular glaucoma, uveitic glaucoma, and post-penetrating keratoplasty glaucoma.

\synonyms
Molteno Glaucoma Drainage Device; Molteno Aqueous Shunt; Molteno Seton

\medterm{Molybdenum} An essential trace element required in small amounts for enzymes; excessive exposure can affect copper balance.

\medterm{Molybdenum-99} A radioactive form of molybdenum that changes into technetium-99m. It is used in generator systems to produce technetium-99m, a radioactive tracer used for many medical scans.

\medterm{Molybdosis} Illness caused by excessive molybdenum exposure, which may disturb copper metabolism and cause systemic effects.

\medterm{Mometasone} Mometasone is a potent synthetic glucocorticoid used topically or by inhalation for inflammatory skin or airway disease. It reduces local inflammation but may cause skin atrophy, oral candidiasis, adrenal effects, or other corticosteroid adverse effects.

\medterm{Mometasone Furoate} Mometasone furoate is the furoate ester of the corticosteroid mometasone. Depending on formulation, it is used for inflammatory skin disease, allergic rhinitis, or asthma; local and systemic corticosteroid precautions apply.

\medterm{Momordica} A genus of plants including bitter melon; preparations may affect blood glucose and can interact with medicines.

\medterm{Monascus} Monascus is a genus of fungi used in some fermented foods; some species produce pigments, while contamination can produce harmful mycotoxins.

\medterm{Mondor’S Disease} Superficial thrombophlebitis of the anterior chest wall veins, typically the lateral
thoracic, thoracoepigastric, or superior epigastric veins. Presents as a tender,
cord-like subcutaneous structure with skin retraction. Usually self-limited; associated
with breast surgery, trauma, or hypercoagulable states.

\medterm{Mongolian Spot} A Mongolian spot, now preferably called congenital dermal melanocytosis, is a flat blue-gray patch caused by melanocytes in the dermis, usually present at birth. It is benign and often fades during childhood but should be documented to distinguish it from bruising.

\medterm{Monilethrix} Monilethrix is an inherited hair-shaft disorder in which hair has regularly spaced narrow segments and breaks easily, producing a beaded appearance. It commonly causes short, brittle scalp hair and may involve follicles elsewhere.

\medterm{Monitored Anesthesia Care} Anesthesia care in which a trained professional gives medicines for comfort or pain relief and continuously watches breathing, heart rate, blood pressure, and oxygen. The level of sedation can be adjusted and emergency support is available.

\medterm{Monitoring Device} A monitoring device measures and displays or records a physiologic variable, such as heart rhythm, oxygen saturation, blood pressure, glucose, temperature, or respiratory rate. Accuracy, calibration, and clinical context determine how results should be interpreted.

\medterm{Monitoring Test} A test repeated over time to follow a disease, treatment response, medicine safety, or a possible complication. The result is interpreted with symptoms and earlier measurements.

\medterm{Monkeypox} Monkeypox, now generally called mpox, is a zoonotic infection caused by monkeypox virus, an orthopoxvirus. It can cause fever, lymphadenopathy, and a characteristic rash or lesions and spreads through close contact, including skin-to-skin or respiratory contact.

\medterm{Mono} Informal term; see \textit{Infectious mononucleosis}.

\medterm{Mono} Alternate name; see \textit{Infectious Mononucleosis}.

\medterm{Monoamine Oxidase Inhibitors} A class of antidepressants that slows the breakdown of several brain chemical messengers. They can help selected people with depression but have important interactions with medicines and some foods, requiring careful prescribing.

\medterm{Monobactams} A class of antibiotics used mainly for certain infections caused by oxygen-using Gram-negative bacteria. Aztreonam is the main example; laboratory testing and allergy history help guide its use.

\medterm{Monoblast} An early bone marrow cell that develops into a monocyte, a type of infection-fighting white blood cell.

\medterm{Monoblast Count} A monoblast count is the number or proportion of monoblasts, immature precursors of monocytes, in a blood or bone-marrow sample. Increased blasts may support evaluation for acute myeloid leukemia or another marrow disorder and require expert morphologic or immunophenotypic interpretation.

\medterm{Monochorionic Twin Pregnancy} A twin pregnancy in which identical twins share a single placenta. Because the babies share placental blood vessels, the pregnancy carries unique risks such as twin-to-twin transfusion syndrome and unequal growth, requiring frequent ultrasound monitoring throughout pregnancy. Twins may have separate inner fluid sacs (diamniotic) or rarely share a single sac (monoamniotic).
\synonyms Monochorionic twins; MCDA twins; MCMA twins.

\medterm{Monoclonal} Made by, or arising from, one original cell and therefore very similar to that cell. A monoclonal antibody is designed to recognize one specific target in the body or in a laboratory test.

\medterm{Monoclonal Antibodies} Antibodies produced from a single B-cell clone and directed against one antigenic target; they are used in diagnosis and treatment.

\medterm{Monoclonal Antibody} A laboratory-made immune protein designed to recognize one specific target. Monoclonal antibodies are used in tests and treatments for cancer, inflammatory disease, infection, and other conditions.

\medterm{Monoclonal Gammopathy Of Undetermined Significance} A condition in which abnormal plasma cells make a small amount of one unusual antibody. It usually causes no symptoms or organ damage but can rarely progress to multiple myeloma or another blood cancer, so monitoring is recommended.

\medterm{Monocrotophos} Monocrotophos is a highly toxic organophosphorus insecticide that inhibits acetylcholinesterase. Exposure can cause acute cholinergic poisoning with secretions, bronchospasm, vomiting, bradycardia, weakness, seizures, and respiratory failure.

\medterm{Monocular} Involving or affecting one eye, or occurring with vision from one eye. Monocular vision can reduce depth perception and may result from eye disease, injury, or loss of vision in the other eye.

\medterm{Monocyte} A type of white blood cell that travels in the blood and can enter tissues, where it becomes a macrophage or another immune cell. It helps remove germs and damaged material and coordinates inflammation and repair.

\medterm{Monocytes} Large white blood cells that circulate in the blood and can enter tissues, where they become macrophages or some dendritic cells. They help remove debris and coordinate inflammation and immune defense.

\medterm{Monocytopenia} An abnormally low number of monocytes, a type of white blood cell, in the blood. It may occur with bone-marrow disorders, medicines, infections, or severe illness and is interpreted with the full blood count.

\medterm{Monocytosis} Monocytosis is an increased absolute number of circulating monocytes. It may be reactive to infection or inflammation or associated with a hematologic neoplasm; persistent or marked monocytosis requires appropriate evaluation.

\medterm{Monogenic} Describing a condition or trait caused mainly by a change in one gene. The change may be inherited from a parent or arise for the first time in the affected person.

\medterm{Monogenic Obesity} Severe or early-onset obesity caused mainly by a disease-causing change in one gene. The affected gene may alter appetite or the body's control of energy use.

\medterm{Monograph} A monograph is a detailed, authoritative document devoted to one subject. In health care it may describe a drug’s identity, quality, pharmacology, indications, dosing, safety, or regulatory standards.

\medterm{Monoiodotyrosine} Monoiodotyrosine is an iodinated tyrosine residue formed during thyroid-hormone synthesis. Coupling of iodinated tyrosine residues within thyroglobulin contributes to production of thyroxine and triiodothyronine.

\medterm{Monokines} Monokines are cytokines produced chiefly by monocytes and macrophages, including mediators such as tumor necrosis factor and interleukin-1. They regulate inflammation, immune responses, fever, and communication among cells.

\medterm{Monomer} A monomer is a relatively small molecule that can chemically bond with other monomers to form a polymer. Biological examples include amino acids and nucleotides; some industrial monomers can irritate tissues or be toxic.

\medterm{Monomethylhydrazine} Monomethylhydrazine is a volatile, corrosive, highly toxic chemical used as a rocket propellant. Exposure can injure the eyes, skin, lungs, liver, and nervous system and may cause seizures or systemic poisoning.

\medterm{Mononegavirales} An order of viruses with nonsegmented, negative-sense RNA genomes, including rabies, measles, Ebola, and respiratory syncytial viruses.

\medterm{Mononeuropathy} Mononeuropathy is dysfunction or injury of a single peripheral nerve. Common causes include compression, trauma, ischemia, inflammation, and metabolic disease; symptoms follow the nerve’s distribution and may include pain, numbness, weakness, or autonomic changes.

\medterm{Mononuclear Leukocyte} A white blood cell with one main, unsegmented nucleus. Monocytes and lymphocytes are examples; they help remove damaged material, recognize germs, and coordinate immune responses.

\medterm{Mononucleosis} Common shortened term; see \textit{Infectious mononucleosis}.

\medterm{Mononucleosis Spot Test} A rapid blood test detecting heterophile antibodies associated with infectious mononucleosis; it may be negative early in illness or in young children, so other testing may be needed.

\medterm{Monoplegia} Complete paralysis of one arm or one leg. It can result from a stroke, brain or spinal-cord disease, or severe damage to a nerve. The cause is found through a neurological examination and, when needed, brain or spine imaging and nerve tests.

\medterm{Monopodium} A single foot or one-sided foot structure; the term is uncommon and may be used descriptively in developmental anatomy.

\medterm{Monorchidism} The presence of only one testis, either because the other failed to develop, descended, or was removed.

\medterm{Monosaccharide} The simplest form of sugar, which cannot be split into a smaller carbohydrate. Glucose and fructose are examples; cells use these sugars mainly for energy.

\medterm{Monosodium Glutamate} MSG, the sodium salt of glutamic acid, used as a flavor enhancer because it contributes umami taste; it occurs naturally in foods such as tomatoes and aged cheese.

\medterm{Monosodium Urate} Needle-shaped crystals formed from uric acid that can collect in joints and other tissues when the blood uric-acid level is high. They cause the inflammation of gout and may be identified in joint fluid with a microscope.

\medterm{Monosomy} A genetic condition in which a cell has one copy of a chromosome instead of the usual two. Effects depend on the chromosome and whether all cells or only some cells are affected.

\medterm{Monospot Test} A rapid blood test that detects heterophile antibodies associated with infectious mononucleosis. It may be
negative early in illness or in young children and does not directly detect Epstein-Barr
virus; specific antibody testing may be needed.

\synonyms
Heterophile Antibody Test; Paul-Bunnell Test

\medterm{Monotest} A rapid blood test that looks for heterophile antibodies produced during some cases of infectious mononucleosis. A negative result does not always exclude infection, especially early in illness or in young children.

\medterm{Monounsaturated Fat} A fat whose fatty acids contain one carbon-carbon double bond. It occurs in foods such as olive oil, avocados, nuts, and some seeds.

\medterm{Monozygotic Twin} A monozygotic twin develops when one fertilized egg divides into two embryos. The twins share nearly identical nuclear genomes but can differ through developmental, epigenetic, mitochondrial, and environmental factors.

\medterm{Mons Veneris} The rounded fatty prominence over the pubic bone in front of the external female genitalia.

\medterm{Montelukast} Montelukast is a leukotriene CysLT1-receptor antagonist used for asthma prevention and allergic rhinitis; it is not a rescue bronchodilator, and neuropsychiatric adverse effects require attention.

\medterm{Montelukast Sodium} Montelukast sodium is the sodium salt of montelukast, a leukotriene-receptor antagonist used for asthma prevention and allergic rhinitis. It is not a rescue bronchodilator and carries an important warning for possible neuropsychiatric adverse effects.

\medterm{Montevideo Units} A measurement of the strength of uterine contractions over 10 minutes during labor. It is calculated from pressure readings obtained by a catheter inside the uterus and helps clinicians assess whether contractions are strong enough to support labor progress.

\medterm{Montgomery Tubercles} Small oil-producing glands in the darker skin around the nipple. They often become more noticeable during pregnancy and breastfeeding and release oil that helps protect and lubricate the nipple; they are usually normal.

\medterm{Montreal Cognitive Assessment} A brief screening test of memory, attention, language, problem-solving, visual-spatial skills, and orientation. It can identify possible thinking problems but does not by itself diagnose dementia; education, language, sensory impairment, and acute illness affect the score.

\medterm{Mood} A lasting emotional state experienced by a person and noticed by others. Mood may be described as depressed, anxious, irritable, elevated, or stable and can change with illness, medicines, stress, or life events.

\medterm{Mood Disorder} Alternate name; see \textit{Mood Disorders}.

\medterm{Mood Disorders} A group of mental disorders characterized primarily by clinically significant depression, mania, hypomania, or combinations of these mood episodes. The group includes depressive disorders, bipolar disorders, and cyclothymic disorder.

\medterm{Mood Stabilizers} Medicines used mainly to treat or prevent episodes of mania and depression in bipolar disorder. Examples include lithium, valproate, carbamazepine, and lamotrigine; each has different benefits, interactions, monitoring needs, and risks during pregnancy.

\medterm{Mood Swing} A mood swing is a noticeable change in emotional state or affect over time. Occasional changes can be normal, while frequent, severe, or impairing swings may occur with mood disorders, substance use, medical illness, medications, or sleep disruption.

\medterm{Morale} Morale is the confidence, motivation, and sense of shared purpose within a person or group. It can affect teamwork and occupational well-being but is not itself a specific medical diagnosis.

\medterm{Moraxella} Moraxella is a genus of bacteria that can colonize the respiratory tract and cause ear, sinus, eye, or respiratory infections.

\medterm{Moraxella Atlantae} Moraxella atlantae is a rare Moraxella species; its recovery from a clinical sample should be interpreted with the site and evidence of infection.

\medterm{Moraxella Lincolnii} Moraxella lincolnii is a rare respiratory-tract bacterium; clinical importance depends on the specimen and evidence of infection.

\medterm{Moraxella Oblonga} Moraxella oblonga is a rare Moraxella species found in respiratory or oral specimens; significance depends on clinical evidence of infection.

\medterm{Moraxella Osloensis} An uncommon bacterium that can occasionally cause bloodstream, respiratory, wound, or other opportunistic infections, usually in people with serious illness or weakened immunity.

\medterm{Moraxellaceae} A family of bacteria that includes Moraxella and related organisms. Some normally live on the skin or in the airways, while others can cause ear, sinus, lung, wound, or bloodstream infections.

\medterm{Morbidity} Illness or disability in a person or group of people. It can also
mean how many people have a disease or how much illness or disability a disease causes.

\medterm{Morbilli} The medical or Latin name for measles, an acute contagious viral infection with fever and a characteristic rash.

\medterm{Morbilli Form} Resembling measles in appearance, especially a widespread rash made of flat and slightly raised spots. This description does not prove that a person has measles because medicines, infections, and other illnesses can cause similar rashes.

\medterm{Morbillivirus} A genus of paramyxoviruses that includes measles virus and related animal viruses causing fever, respiratory illness, rash, or neurologic disease.

\medterm{Morbus} A Latin word meaning disease. It appears in some older or formal disease names, but the word itself does not describe the cause, symptoms, or treatment of the condition.

\medterm{Morcellation} The division of tissue into smaller fragments, often during minimally invasive surgery; it may spread an unrecognized cancer and requires appropriate case selection.

\medterm{Mordant} A substance that helps fix a dye to tissue or a material, used in some histological staining methods.

\medterm{Morgagni Hernia} A congenital diaphragmatic hernia occurring through the anterior retrosternal foramen of
Morgagni, resulting from incomplete fusion of parts of the diaphragm. Abdominal contents
may move into the chest and cause respiratory or gastrointestinal symptoms.

\medterm{Morganella} Morganella is a genus of gram-negative bacteria that can cause urinary, wound, bloodstream, and other opportunistic infections.

\medterm{Morganella Morganii} A gram-negative bacterium that can cause urinary, wound, bloodstream, or postoperative infections, especially in hospitalized or medically vulnerable people. Antibiotic choice requires susceptibility testing because resistance is common.

\medterm{Morgellons Disease} Morgellons is a controversial condition involving distressing skin sensations and the belief that fibers or materials emerge from the skin. Symptoms are real to the person and require respectful evaluation for dermatologic, neurologic, psychiatric, toxic, or infectious causes without reinforcing an unproven explanation.

\medterm{Moribund} Critically ill and close to death, with little or no chance of recovery from the underlying condition. A moribund person may be unconscious, have very low blood pressure, cold or mottled skin, and irregular or failing breathing. Care focuses on urgent treatment when possible and comfort when death is inevitable.

\medterm{Morning Sickness} Nausea, with or without vomiting, during pregnancy. It often begins early and improves by the middle of pregnancy, but severe or persistent vomiting can cause dehydration, weight loss, and nutrient problems and may need medical treatment.

\medterm{Moro Reflex} A normal newborn startle response that causes the baby to spread out the arms, open the hands, and then draw the arms inward. It usually fades by about 3 to 6 months; an absent or unequal response needs assessment.

\medterm{Morphea} A localized sclerosing skin disorder that produces areas of thickened, hardened, or discolored skin and differs from systemic sclerosis by usually lacking internal-organ involvement.

\medterm{Morpheus} The Greek and Roman mythological god of dreams; the name is the source of \textit{morphine}, an opioid named for its association with sleep and dreams.

\medterm{Morphinans} Morphinans are a group of chemical compounds with a fused-ring structure related to morphine. The class includes opioid analgesics and other compounds with varying receptor activity, potency, and clinical effects.

\medterm{Morphine} Morphine is a potent opioid analgesic used for selected moderate-to-severe pain and palliative-care indications. It activates opioid receptors and can cause respiratory depression, sedation, constipation, nausea, tolerance, dependence, and overdose.

\medterm{Morphine Overdose} Taking too much morphine can cause extreme sleepiness, pinpoint pupils, slow or stopped breathing, low blood pressure, coma, and death. Call emergency services immediately; naloxone may reverse opioid breathing suppression when available.

\medterm{Morphine Sulfate} Morphine sulfate is a salt formulation of the opioid analgesic morphine. It is available in immediate- and extended-release preparations for selected pain indications and carries risks of respiratory depression, misuse, dependence, and fatal overdose.

\medterm{Morphogenesis} The biological process by which cells and tissues change shape, move, and organize to form organs and body structures during development or tissue repair.

\medterm{Morphologic Response} A morphologic response is a treatment-related change in the visible or microscopic structure of a lesion, tumor, tissue, or cell population. It may include shrinkage, necrosis, differentiation, or altered cellular features and is interpreted using disease-specific criteria.

\medterm{Morphology} The study or description of the form, shape, and structural features of organisms, organs, tissues, cells, or other biological structures.

\medterm{Morquio Syndrome} A rare inherited disorder in which complex sugars build up and interfere mainly with bone and cartilage growth. It can cause short stature, an unusually shaped chest and spine, loose joints, cloudy corneas, and pressure on the spinal cord in the neck. Intelligence is usually unaffected. Treatment may include enzyme replacement, surgery, and careful monitoring of breathing and the spine.

\synonyms
Mucopolysaccharidosis Type Iv; MPS IV

\medterm{Mortality} Death, or the number of deaths in a particular group during a stated period. A mortality measure may describe all causes of death or deaths from a specific disease and population.

\medterm{Mortality Rate} The number of deaths in a specified population and time period, often reported for a particular disease or cause. It helps compare the impact of illness between groups when the populations are defined appropriately.

\medterm{Morton Neuroma} Thickening and irritation of a nerve in the ball of the foot, usually between the third and fourth toes. It can cause burning pain, numbness, or the feeling of a pebble under the foot, often worsened by tight shoes or walking.

\synonyms
Interdigital Neuroma; Morton's Metatarsalgia; Plantar Interdigital Neuroma

\medterm{Morton's Metatarsalgia} Pain in the forefoot beneath the metatarsal heads, caused by pressure or irritation and distinct from Morton neuroma.

\medterm{Mortons Neuroma} Alternate spelling; see \textit{Morton's Neuroma}.

\medterm{Morton’S Neuroma} Alternate spelling; see \textit{Morton neuroma}.

\medterm{Morton's Neuroma} Thickening and irritation of tissue around a nerve, usually between the third and fourth toes. It causes burning forefoot pain or a feeling of a pebble in the shoe and may improve with wider footwear.

\medterm{Morton's Toe} A foot shape in which the second toe is longer than the big toe, often because the first metatarsal is relatively short. It is usually a normal variation but may contribute to pressure or forefoot pain.

\medterm{Morula} A morula is an early embryo consisting of a compact ball of blastomeres formed after several cleavage divisions. It develops into a blastocyst before implantation in the uterine lining.

\medterm{Morvan's Disease} A rare autoimmune neurologic syndrome with severe insomnia, hallucinations, confusion, pain, and peripheral nerve hyperexcitability.

\medterm{Mosaic} Describing a person or tissue composed of two or more genetically different cell populations derived from one zygote, often because of a post-zygotic mutation or chromosomal event; mosaicism can alter the severity and distribution of a genetic disorder.

\medterm{Mosaicism} The presence of groups of cells in one person that have different genetic material, even though all the cells came from the same fertilized egg. It may or may not cause disease, and its effects depend on which cells are affected.

\medterm{Mosquito} A small blood-feeding insect whose bites may transmit malaria, dengue, West Nile virus, and other infections.

\medterm{Mosquito Clamp} A small hemostatic clamp used to compress small blood vessels.

\medterm{Mosquito-Borne} Spread by mosquitoes, usually describing an infection or parasite carried from one host to another by a mosquito bite. Examples include malaria, dengue, West Nile virus, yellow fever, and some forms of encephalitis.

\medterm{Motilin} A hormone made mainly in the small intestine that helps trigger waves of movement through the stomach and small intestine, especially between meals. Its activity contributes to the pattern that clears leftover food and secretions from the upper digestive tract.

\medterm{Motility} The ability of an organ to move its contents or itself. Intestinal motility moves food and waste forward, while abnormal motility can cause pain, constipation, diarrhea, vomiting, or swallowing problems.

\medterm{Motion Analysis} A detailed study of how a person moves during walking, reaching, or other activities. Clinicians use measurements and video to find movement problems, understand their causes, plan treatment, and track improvement.

\medterm{Motion Artifact} Image distortion or blurring caused by movement of the patient, organs, or equipment during image acquisition.

\medterm{Motion Correction} Processing or acquisition methods that reduce image blurring and errors caused by patient, breathing, or organ motion. Examples include image registration, tracking, gating, and navigator methods.

\medterm{Motion Perception} Motion perception is the ability of the visual, vestibular, and proprioceptive systems to detect and interpret movement of the body or surrounding objects. Disturbance may cause vertigo, oscillopsia, imbalance, or visual-motion sensitivity.

\medterm{Motion Sickness} Nausea, dizziness, sweating, or vomiting caused by sensory mismatch during travel or other motion. It can occur in cars,
boats, aircraft, amusement rides, or virtual-reality environments and usually improves when motion stops.
Severe or persistent symptoms require assessment for other causes.

\medterm{Motion Sickness} A disturbance of spatial orientation and equilibrium; see \textit{Kinetosis}.

\medterm{Motivation} Motivation is the set of processes that initiate, direct, and sustain goal-directed behavior. Reduced or excessive motivation may occur with mood disorders, neurologic disease, substance use, medication effects, or other conditions.

\medterm{Motivational Interviewing} A collaborative counseling approach that helps a person explore mixed feelings and strengthen their own reasons for changing a health behavior. The clinician listens without judgment, asks open questions, reflects the person’s views, and supports informed choice.

\medterm{Motor} Relating to movement or to the nerves and brain pathways that control movement. A motor problem may cause weakness, stiffness, poor coordination, tremor, or loss of normal control.

\medterm{Motor Aphasia} Difficulty speaking or writing caused by damage to language areas of the brain, often after a stroke. A person may know what they want to say but struggle to find words, form sentences, or produce speech.

\medterm{Motor Cortex} Areas at the front of the brain that plan and control voluntary movement. Each side mainly controls the opposite side of the body. Injury can cause weakness, loss of fine movement, or difficulty planning purposeful actions.

\medterm{Motor Delay} Motor delay is slower-than-expected acquisition of gross- or fine-motor skills in a child. It may reflect neuromuscular, developmental, genetic, sensory, environmental, or systemic causes and warrants age-appropriate developmental assessment.

\medterm{Motor End-Plate} The specialized area where a nerve signal reaches a skeletal muscle fiber. Chemical signals released there start muscle contraction, and damage can cause weakness, fatigue, or abnormal muscle function.

\medterm{Motor Fluctuations Causes In Parkinson’S Disease} Changes between better and worse movement in Parkinson disease, often as the effect of a medicine wears off before the next dose. Periods of stiffness, slowness, or involuntary movement may occur and can require treatment adjustment.

\medterm{Motor Learning} A relatively lasting change in the ability to perform a motor skill resulting from
practice, experience, feedback, and adaptation. Rehabilitation uses task-specific
repetition and meaningful practice to promote motor learning.

\medterm{Motor Neuron} A nerve cell that carries signals from the brain or spinal cord to skeletal muscle, allowing movement. Damage can cause weakness, muscle wasting, stiffness, cramps, or loss of voluntary control.

\synonyms
Motor Nerve Cell; Motoneuron

\medterm{Motor Neuron Disease} Alternate spelling; see \textit{Motor Neurone Disease}.

\medterm{Motor Neurone Disease} A group of progressive disorders affecting upper and lower motor neurons, including amyotrophic lateral sclerosis. They cause worsening weakness, muscle wasting, stiffness, cramps, and difficulty with movement; speech, swallowing, and breathing may become affected as the disease advances. Cognitive or behavioral changes occur in some forms.

\medterm{Motor Neurones} Nerve cells that carry signals from the brain or spinal cord to skeletal muscles, producing movement. Injury or disease can cause weakness, wasting, stiffness, cramps, poor coordination, or loss of voluntary control.

\medterm{Motor Stereotypy} A repeated, predictable movement such as hand flapping, rocking, or body waving. It may occur during normal development or with autism, intellectual disability, neurologic disease, stress, or other conditions.

\medterm{Motor Tic} A motor tic is a sudden, brief, repetitive movement such as blinking, facial grimacing, or shoulder shrugging. Tics are often temporarily suppressible and may occur in transient tic disorder, persistent tic disorder, or Tourette syndrome.

\medterm{Motor Unit} One movement-controlling nerve cell and all the muscle fibers it activates. Small units allow precise movements, while large units produce more force. Damage to motor units can cause weakness, wasting, cramps, or abnormal results on nerve and muscle tests.

\medterm{Motor Vehicle Accident} A collision or other road-traffic event causing injury or property damage. Medical evaluation may be needed for visible or delayed injuries, including head, neck, chest, abdominal, or musculoskeletal trauma.

\medterm{Mottled Enamel} Patchy, discolored, or structurally altered tooth enamel, classically caused by excessive fluoride during tooth development.

\medterm{Mottled Teeth} Teeth with irregular white, brown, or pitted enamel, often due to dental fluorosis or developmental enamel defects.

\medterm{Moulage} A model, cast, or realistic reproduction of a body part or lesion used for teaching, simulation, or forensic documentation.

\medterm{Mould} British spelling of mold, a multicellular fungus that grows as branching filaments.

\medterm{Mountain Sickness} Illness caused by rapid ascent to high altitude, ranging from acute mountain sickness to life-threatening pulmonary or cerebral edema.

\medterm{Mouth} The oral cavity that serves as the entry point of the digestive tract, used for eating,
breathing, and speaking.

\medterm{Mouth Cancer} Cancer arising in the lips, gums, tongue, cheeks, floor of the mouth, palate, or other oral tissues.

\synonyms
Oral Cancer; Oral Cavity Carcinoma; Oral Squamous Cell Carcinoma

\medterm{Mouth Diseases} Disorders affecting the oral cavity, including infections, ulcers, inflammatory disease, dental disease, mucosal lesions, and tumors.

\medterm{Mouth Neoplasm} A mouth neoplasm is an abnormal growth in the oral cavity that may be benign or malignant; persistent sores, lumps, or patches need assessment.

\medterm{Mouth Sores} Painful ulcers, blisters, raw areas, or inflamed patches inside or around the mouth. Trauma, infection, immune disease, medicines, nutritional deficiency, and cancer can cause them; sores lasting more than two weeks need assessment.

\medterm{Mouth Ulcer} Open sores affecting the oral mucosa, commonly caused by aphthous ulcers, trauma,
infection, irritation, or systemic inflammatory disease. Most resolve within two weeks;
persistent, recurrent, unusually large, or indurated ulcers require evaluation.

\synonyms
Oral Ulcer; Aphthous Ulcer; Canker Sore

\medterm{Mouth-To-Mouth Resuscitation} Rescue breathing in which air is delivered from a rescuer to a person who is not breathing normally, usually as part of CPR.

\medterm{Mouth-Wash} A liquid used to rinse the mouth, freshen breath, or deliver an antiseptic, fluoride, or other active ingredient.

\medterm{Mouthwash} A liquid rinsed around the mouth and then expelled to freshen breath, deliver an active ingredient, or reduce selected oral microorganisms.

\medterm{Mouthwash Overdose} Swallowing too much mouthwash may cause alcohol poisoning, vomiting, drowsiness, low blood sugar, breathing problems, or coma; products containing fluoride or other ingredients can cause additional toxicity. Contact poison control promptly.

\medterm{Movement} Movement is a change in position of the body or a body part produced by coordinated activity of the nervous system, muscles, and joints. Abnormal movement may be excessive, reduced, involuntary, poorly coordinated, or restricted.

\medterm{Movement Disorder} A neurologic condition causing excessive, reduced, or otherwise abnormal movement.
Hyperkinetic disorders include tremor, dystonia, chorea, myoclonus, tics, and
dyskinesia; hypokinetic disorders include bradykinesia and parkinsonism.

\medterm{Movement Disorders} Alternate name; see \textit{Movement Disorders}.

\medterm{Movement Incoordination} Impairment in the ability to execute smooth, coordinated voluntary movements (ataxia), resulting from lesions or dysfunction within the cerebellum, spinocerebellar tracts, vestibular apparatus, or proprioceptive pathways. Manifests as dysmetria, intention tremor, and gait instability.

\synonyms
Ataxic Incoordination; Motor Incoordination; Dysdiadochokinesia

\medterm{Moxalactam} Moxalactam is a beta-lactam antibiotic of the oxacephem group with activity against many gram-negative bacteria and some gram-positive organisms. It is rarely used today and may cause allergic reactions, diarrhea, or bleeding related to impaired coagulation.

\medterm{Moxibustion} A traditional treatment that warms selected skin areas with burning mugwort or another heated material. Evidence of benefit varies by condition, and burns, smoke exposure, allergic reactions, or infection can occur.

\medterm{Moxifloxacin} An antibiotic used for selected bacterial infections. It interferes with bacterial reproduction but can cause serious side effects, including tendon injury, nerve symptoms, blood-sugar changes, or dangerous heart rhythms, so use is selective.

\medterm{Moyamoya Disease} A rare condition in which major arteries supplying the brain gradually narrow or close. Small fragile blood vessels form to compensate, but the condition can cause transient attacks, strokes, headaches, seizures, or thinking problems.

\medterm{Moynihan's Hump} An anatomical variation in which the right hepatic artery forms an abnormally tortuous, caterpillar-like loop that extends close to the gallbladder neck and cystic duct within Calot's triangle. During gallbladder removal (cholecystectomy), this loop can be mistaken for the smaller cystic artery; accidental clipping or injury can cause severe arterial bleeding or damage blood flow to the right lobe of the liver.

\synonyms
Moynihan Caterpillar Turn; Caterpillar Turn of the Right Hepatic Artery; Hepatic Artery Tortuosity

\medterm{MPH} Master of Public Health, a graduate degree focused on population health, disease prevention, health policy, epidemiology, and related public-health practice.

\medterm{MPO} Myeloperoxidase, an enzyme in neutrophil and monocyte granules that uses hydrogen peroxide and halide ions to generate antimicrobial oxidants; MPO antibodies are also used in evaluating some autoimmune vasculitides.

\medterm{Mpox} A viral infection that can cause fever, swollen lymph nodes, tiredness, and a rash or painful lesions on the skin, mouth, or genital area. It spreads mainly through close contact with lesions, body fluids, or contaminated materials.

\medterm{Mpox} A viral disease caused by monkeypox virus, an orthopoxvirus related to smallpox. It usually causes a rash or lesions, fever, swollen lymph nodes, and body aches; spread occurs mainly through close contact, and severity varies with the virus type and the person’s health.

\synonyms
Simian Pox; Monkeypox Virus Infection

\medterm{MPTP} A toxic chemical that can permanently damage movement-controlling nerve cells and produce symptoms resembling Parkinson disease. It is mainly important in research and is not a routine medical treatment.

\medterm{MR Angiography} An MRI scan that creates pictures of blood vessels without using X-rays. It can help assess narrowed or blocked arteries, aneurysms, dissections, clots, and abnormal blood-vessel connections.

\medterm{MR Enterography} A specialized MRI scan of the small intestine after drinking contrast liquid. It can show inflammation, narrowing, fistulas, abscesses, bleeding, or other complications of Crohn disease and does not use X-rays.

\medterm{MR Neurography} Magnetic resonance imaging optimized to show peripheral nerves and abnormalities such as swelling, compression, or injury.

\medterm{MRA} Magnetic resonance angiography, a noninvasive MRI-based technique that produces images of blood vessels and can help evaluate aneurysms, stenoses, malformations, and atherosclerosis.

\medterm{MRCP} Magnetic resonance cholangiopancreatography, an MRI technique that noninvasively images the bile ducts, gallbladder, and pancreatic duct without requiring endoscopic catheterization.

\medterm{MRD} Minimal residual disease: a small number of malignant cells remaining after treatment, detectable by sensitive methods such as flow cytometry or PCR even when routine tests show remission.

\medterm{MRI} Magnetic resonance imaging uses a strong magnetic field, radiofrequency energy, and computer processing to produce detailed images of body tissues without ionizing radiation. Contrast material may be used when clinically appropriate.

\medterm{MRI And Low Back Pain} MRI may be used for low back pain when serious disease is suspected, neurologic deficits are present, or symptoms persist despite appropriate care; many uncomplicated cases do not need early MRI.

\medterm{MRI Brain} An MRI scan used to examine the brain for tumors, stroke, bleeding, infection, inflammation, multiple sclerosis, or degenerative disease. Different settings highlight different tissues and abnormalities; contrast may be used when needed.

\medterm{MRI Contrast} A contrast agent administered during magnetic resonance imaging to improve visibility of selected tissues or abnormalities.

\medterm{MRI Joint} Evaluation of musculoskeletal structures: cartilage, ligaments, tendons, menisci,
labrum. Knee: ACL, meniscus tears. Shoulder: rotator cuff, labrum. Hip: labral tears,
FAI.

\medterm{MRI Sequences} Different image settings used during an MRI scan to highlight anatomy, fluid, fat, bleeding, restricted water movement, or inflammation. Radiologists combine several settings because each shows different types of abnormality.

\medterm{MRI-Targeted Fusion Prostate Biopsy} A prostate biopsy guided by combining a previous MRI with real-time ultrasound. The images help the clinician direct needles toward areas that look suspicious for cancer, sometimes along with routine samples from other areas. The procedure can cause blood in the urine or semen, infection, pain, or temporary difficulty passing urine.

\synonyms
Fusion Biopsy; Mri-Ultrasound Fusion Biopsy

\medterm{Mrna Precursor} An initial RNA transcript that is processed by capping, splicing, and polyadenylation to form mature messenger RNA.

\medterm{MRSA} A strain of Staphylococcus aureus resistant to several common antibiotics. It may live harmlessly on the skin or cause skin, wound, lung, bone, joint, or bloodstream infection; treatment is guided by culture results.

\medterm{MRSA Decolonization} Measures intended to reduce MRSA carriage on the skin or in the nose. Protocols may include topical
agents and antiseptic bathing, but the indication, regimen, and duration depend on the
person's risk, infection history, setting, and local guidance.

\medterm{MRSA Infection} Methicillin-resistant Staphylococcus aureus is a strain of staphylococcus resistant to several beta-lactam antibiotics. It can colonize without illness or cause skin, wound, lung, bone, joint, or bloodstream infection; culture-guided treatment and hygiene reduce spread.

\synonyms
Methicillin-Resistant Staphylococcus Aureus Infection

\medterm{MS Alternative Therapies} Alternate name; see \textit{Alternative Therapy for Multiple Sclerosis}.

\medterm{MSG Symptom Complex} A group of symptoms such as headache, flushing, or palpitations attributed by some people to monosodium glutamate; evidence for a consistent syndrome is limited and symptoms may have other causes.

\medterm{MST} Multiple subpial transection, a neurosurgical procedure in which shallow cuts are made in the cerebral cortex to interrupt seizure pathways when the epileptogenic region cannot safely be removed.

\medterm{MTOR} A protein inside cells that helps control growth, division, energy use, and response to nutrients. Medicines that block mTOR, such as rapamycin-related drugs, are used in some cancers, transplant treatments, and other selected conditions.

\medterm{Mu-Chain Disease} A rare B-cell disorder in which abnormal immunoglobulin heavy chains are produced, sometimes causing anemia, enlarged lymph nodes, or other features of lymphoproliferative disease.

\medterm{Mucin} Mucin is a high-molecular-weight glycoprotein that forms mucus and helps lubricate and protect epithelial surfaces. Abnormal mucin production or accumulation occurs in some respiratory, gastrointestinal, and neoplastic diseases.

\medterm{Mucinoses} A group of disorders in which excessive mucin, a gel-like substance in connective tissue, builds up in the skin or other tissues. Causes may be primary, metabolic, autoimmune, or associated with another disease.

\medterm{Mucinous} Mucinous means containing, producing, or resembling mucin. In pathology, the term describes lesions or tumors with extracellular or intracellular mucin and must be interpreted with the organ, architecture, and other histologic features.

\medterm{Mucinous Adenocarcinoma} Mucinous adenocarcinoma is a gland-forming cancer that produces abundant mucus and can arise in organs such as the colon, pancreas, lung, or breast.

\medterm{Mucinous Borderline Ovarian Tumor} An ovarian mucinous epithelial tumor with cellular proliferation and atypia but without destructive stromal invasion; it may be large and unilateral.

\medterm{Mucinous Cyst Of The Vulva} A benign mucus-filled cyst of the vulva, arising from obstruction or metaplastic mucinous epithelium in a vulvar gland or duct; persistent or atypical masses require examination and possible biopsy.

\medterm{Mucinous Cystadenoma} A mucinous cystadenoma is a usually benign cystic epithelial tumor that produces mucin, commonly arising in the ovary or pancreas. Large lesions may cause pain or pressure, and pathology is needed to distinguish benign from borderline or malignant mucinous tumors.

\medterm{Mucinous Cystadenoma of Pancreas} Alternate name; see \textit{Pancreatic Cysts}.

\medterm{Mucinous Ovarian Carcinoma} A malignant ovarian epithelial tumor composed of invasive mucin-producing cells; metastatic gastrointestinal or pancreatobiliary carcinoma must be excluded.

\medterm{Muckle-Wells Syndrome} An inherited cryopyrin-associated autoinflammatory syndrome causing recurrent fever, rash, joint or muscle pain, and progressive sensorineural hearing loss. Anti-inflammatory treatment can reduce inflammation and prevent complications.

\medterm{Mucocele} A soft, mucus-filled swelling caused by leakage or blockage of a small salivary gland, most often on the inner lower lip after minor injury. It may come and go, but a persistent or recurrent swelling should be examined and may need removal.

\medterm{Mucocutaneous Leishmaniasis} A chronic, mutilating metastatic form of leishmaniasis caused by New World species of the Leishmania subgenus Viannia (principally Leishmania braziliensis), which disseminates from a primary cutaneous ulcer to mucosal tissues through hematogenous or lymphatic routes. It causes progressive, destructive granulomatous ulceration and cartilaginous necrosis of the nasal septum, palate, pharynx, and larynx, clinically presenting with nasal obstruction, epistaxis, septal perforation, and severe facial disfigurement.

\synonyms
Espundia; American Mucocutaneous Leishmaniasis; Nasopharyngeal Leishmaniasis

\medterm{Mucocutaneous Lymph Node Syndrome} Alternate name; see \textit{Kawasaki Disease}.

\medterm{Mucoepidermoid Carcinoma} Mucoepidermoid carcinoma is a cancer containing mucus-producing and squamous-like cells, most often arising in salivary glands or the lungs.

\medterm{Mucoid} Resembling mucus or having a thick, slippery, mucus-like appearance. The word describes a substance or finding and does not by itself identify its cause.

\medterm{Mucolipidoses} A group of inherited lysosomal storage disorders in which cells cannot properly process certain lipids and carbohydrate-containing substances. Material accumulates in tissues and may cause developmental delay, skeletal changes, vision problems, and organ disease.

\medterm{Mucolytic Therapy} Treatment that makes thick mucus easier to move and cough up. Different medicines work by breaking mucus proteins, removing extra DNA, or drawing water into airway secretions; the choice depends on the disease and the person's ability to clear mucus.

\medterm{Mucolytics} Medicines that make thick mucus thinner and easier to clear from the airways. They may help some people cough up mucus, but treatment depends on the cause of congestion and the person's overall condition.

\medterm{Mucopolysaccharides} Long chains of sugar molecules that form part of connective tissue and help support joints, skin, cartilage, and other structures. Inherited problems breaking them down can cause storage diseases.

\medterm{Mucopolysaccharidoses} A group of inherited lysosomal storage diseases in which glycosaminoglycans accumulate because one of several degradation enzymes is deficient. The disorders can affect the skeleton, joints, airways, heart, eyes, abdominal organs, and nervous system, with severity varying by type.

\medterm{Mucopolysaccharidosis} The singular name for any of a group of inherited lysosomal storage disorders in which glycosaminoglycans accumulate because of deficient enzyme activity. Types differ in their genetic cause, age of onset, and effects on the skeleton, organs, airways, eyes, heart, and nervous system.

\medterm{Mucopolysaccharidosis Type I} An inherited lysosomal storage disorder caused by deficient alpha-L-iduronidase, usually because of disease-causing variants in the \textit{IDUA} gene. It ranges from severe Hurler syndrome to milder forms and may affect the skeleton, airways, heart, eyes, abdominal organs, and nervous system.

\synonyms
Hurler Syndrome; Hurler-Scheie Syndrome; Scheie Syndrome

\medterm{Mucopolysaccharidosis Type I} Mucopolysaccharidosis type I is an inherited alpha-L-iduronidase deficiency ranging from severe Hurler syndrome to milder forms, with skeletal, airway, cardiac, eye, and neurologic disease.

\medterm{Mucopolysaccharidosis Type II} Mucopolysaccharidosis type II, or Hunter syndrome, is an X-linked iduronate-2-sulfatase deficiency causing accumulation of glycosaminoglycans and multisystem disease.

\medterm{Mucopolysaccharidosis Type III} Mucopolysaccharidosis type III, or Sanfilippo syndrome, is an inherited disorder of heparan-sulfate breakdown causing progressive neurodevelopmental and behavioral decline.

\medterm{Mucopolysaccharidosis Type IV} Mucopolysaccharidosis type IV, or Morquio syndrome, is an inherited disorder of keratan-sulfate breakdown causing skeletal and connective-tissue abnormalities with variable neurologic involvement.

\medterm{Mucopolysaccharidosis VII} Mucopolysaccharidosis VII, or Sly syndrome, is a lysosomal storage disorder caused by deficient beta-glucuronidase activity. Accumulation of glycosaminoglycans can cause skeletal, connective-tissue, visceral, developmental, and sometimes neurologic abnormalities.

\medterm{Mucopus} A mixture of mucus and pus produced by inflammation or infection. It may appear in the nose, sinuses, airways, bowel, cervix, or other mucus-lined areas and should be interpreted with symptoms.

\medterm{Mucor} Mucor is a genus of molds; related Mucorales fungi can cause mucormycosis, a serious invasive infection mainly in people with severe immune or metabolic problems.

\medterm{Mucor Ramosissimus} A type of mold found in the environment that rarely causes infection in people. When infection occurs, it is usually in someone with severely weakened immune defenses and requires identification of the fungus and specialist treatment.

\medterm{Mucorales} Mucorales is an order of molds that includes fungi capable of causing mucormycosis, a rapidly invasive infection of the sinuses, lungs, skin, or other tissues.

\medterm{Mucormycosis} A rare but rapidly destructive infection caused by certain molds. It most often affects people with severely weakened immunity or uncontrolled diabetes and can damage the sinuses, eyes, lungs, skin, brain, or several organs.

\medterm{Mucosa} Mucosa is the moist epithelial lining of body passages such as the mouth, nose, airways, gastrointestinal tract, and reproductive or urinary tracts.

\medterm{Mucosa-Associated Lymphoid Tissue Lymphoma} Alternate name; see \textit{MALT Lymphoma}.

\medterm{Mucosal} Relating to a mucous membrane, the moist lining of areas such as the mouth, nose, airways, digestive tract, and genital or urinary tract. Mucosa helps protect tissues, produce mucus, absorb substances in some organs, and support local immune defense.

\medterm{Mucosal Immunity} Mucosal immunity is immune defense at surfaces lining the respiratory, gastrointestinal, genitourinary, and other tracts. It includes epithelial barriers, mucus, secretory IgA, antimicrobial molecules, innate cells, and locally regulated adaptive responses.

\medterm{Mucosal Infection} An infection involving a moist epithelial lining such as the mouth, eye, nose, airway, gastrointestinal tract, or genital tract.

\medterm{Mucosal Irritation} Inflammation or injury of moist linings in the eyes, nose, throat, mouth, stomach, or intestines. It may cause burning, pain, redness, swelling, or sores after chemicals, medicines, infection, or friction.

\medterm{Mucosal Melanoma} A melanoma arising from melanocytes in mucosal surfaces such as the sinonasal tract, oral cavity, anorectal region, or genital tract; it is uncommon and often diagnosed at an advanced stage.

\medterm{Mucosal Warts} Warts caused by human papillomavirus on moist lining tissues such as the mouth, genitals, or anus. They may be painless or cause itching, bleeding, or discomfort; some HPV types increase cancer risk and need medical assessment.

\medterm{Mucositis} Painful inflammation and sores in the moist lining of the mouth or digestive tract, often after chemotherapy or radiation. It can make eating and swallowing difficult and may increase the risk of infection.

\medterm{Mucous} Relating to mucus, the viscous fluid produced by mucous membranes and glands that lubricates surfaces and helps trap particles and microorganisms.

\medterm{Mucous Cysts} Small fluid-filled lumps near the end joint of a finger, often associated with osteoarthritis. They may cause tenderness, nail changes, or leakage and should not be punctured without medical advice.
\synonyms
Digital Mucous Cyst; Myxoid Cyst

\medterm{Mucous Membrane} A moist tissue lining body cavities that open to the outside, including the nose, mouth, and lungs.

\synonyms
Mucosa

\medterm{Mucous Membrane Pemphigoid} A chronic autoimmune subepidermal blistering disease affecting mucous membranes, especially the eyes and mouth; scarring can impair vision or function.

\medterm{Mucus} Mucus is a viscoelastic secretion containing water, mucin glycoproteins, salts, antimicrobial factors, and cellular material. It lubricates and protects epithelial surfaces and traps particles and microorganisms for clearance.

\medterm{Mucus Plugging} Accumulation of thick, tenacious mucus within airways causing obstruction. It is common
in asthma, CF, bronchiectasis, and chronic bronchitis. It impairs gas exchange, promotes
infection, and causes atelectasis. Treatment includes airway clearance techniques,
mucolytics, and bronchodilators.

\medterm{Muehrcke's Nails} Paired transverse white bands in the nail beds, usually associated with hypoalbuminemia or other systemic illness. The bands remain fixed as the nail grows because they arise from the nail bed rather than the nail plate.

\medterm{Muenke Syndrome} A craniosynostosis syndrome caused by a pathogenic FGFR3 variant. Premature fusion of
one or more skull sutures may produce an abnormal head shape and can be accompanied by
developmental, hearing, limb, or neurological findings.

\medterm{MUGA Scan} A heart scan that uses a small radioactive tracer and records the heart's pumping with each heartbeat. It measures how much blood the main chambers pump and can assess wall movement.

\synonyms
Radionuclide Angiography; Radionuclide Ventriculography

\medterm{Mulberry Molars} Permanent back teeth with many small, rounded, irregular cusps. They are classically linked with congenital syphilis but can also result from other developmental problems affecting tooth formation before birth.

\medterm{Mulibrey Nanism} Mulibrey nanism is a rare inherited multisystem disorder, usually caused by biallelic variants in \textit{TRIM37}. It may cause severe growth restriction, characteristic facial features, pericardial or hepatic disease, insulin resistance, and other organ involvement.

\medterm{Muller Cells} Supporting glial cells of the retina that help nourish retinal nerve cells and maintain the retina's structure and chemical environment.

\medterm{Multicenter Study} A multicenter study is a clinical or research investigation conducted at two or more sites under a coordinated protocol. It can improve enrollment and generalizability but requires harmonized methods, quality control, and site-specific oversight.

\medterm{Multicentric} Having more than one center or focus; in oncology, describing separate primary tumor foci, often also called multifocal or polycentric.

\medterm{Multicystic Dysplastic Kidney} A birth defect in which one kidney develops as a cluster of cysts and has little or no working tissue. The other kidney often provides most kidney function but may also need monitoring for related urinary problems.

\medterm{Multidetector Computed Tomography} A CT scan that uses many detectors to collect multiple thin image slices at the same time. It produces detailed pictures quickly and is useful for injuries, blood vessels, organs, and some virtual examinations, but uses X-rays.

\medterm{Multidisciplinary Disaster Victim Identification} The coordinated process of identifying victims of mass-fatality incidents (natural
disasters, transport crashes, terrorist attacks) using multiple disciplines:
fingerprinting, dental records, DNA profiling, forensic anthropology/radiology, and
personal effects, with ante-mortem data collection and comparison.

\medterm{Multidisciplinary Team} Healthcare professionals from different fields who work together on one person's care. Depending on the problem, the team may include doctors, nurses, therapists, pharmacists, psychologists, dietitians, and social workers.

\medterm{Multidrug-Resistant TB} Tuberculosis caused by bacteria resistant to at least the two most important first-line medicines, isoniazid and rifampin. It needs specialized medicines, susceptibility testing, and close monitoring because treatment is harder and may last many months.

\medterm{Multienzyme Complexes} Groups of enzymes organized together so that several steps in a chemical pathway occur efficiently. This arrangement can speed reactions, limit loss of intermediates, and coordinate metabolism inside cells.

\medterm{Multifactorial} Caused or influenced by several factors rather than one single cause. These may include genes, age, lifestyle, infections, medicines, environment, and other health conditions, with each factor adding to the overall risk.

\medterm{Multifocal} Arising from or involving two or more distinct foci or sites, such as separate tumor foci or vascular lesions.

\medterm{Multifocal Atrial Tachycardia} An atrial tachyarrhythmia with at least three distinct P-wave morphologies, irregular
rhythm, and a rate above 100 beats per minute. It is strongly associated with severe
pulmonary disease, hypoxia, and metabolic stress.

\medterm{Multigravida} A person who has been pregnant two or more times, including pregnancies that did not continue.

\medterm{Multi-Infarct Dementia} A decline in memory and thinking caused by several strokes that damage different parts of the brain. Problems may include slowed thinking, poor planning, mood changes, walking difficulty, or memory loss; the term is now often included under vascular cognitive impairment.

\synonyms
Vascular Dementia; Vascular Cognitive Impairment

\medterm{Multilocular} Having several separate spaces or compartments. The word may describe a cyst, tumor, or other lesion seen on imaging; its significance depends on the organ involved, the appearance of the walls, and the person's symptoms.

\medterm{Multi-Minicore Disease} A rare inherited muscle disorder in which muscle fibers contain multiple small areas of damage. It can cause low muscle tone, weakness, spinal curvature, and breathing problems.
\synonyms
MmD; Minicore Myopathy; Multicore Disease

\medterm{Multimodal Analgesia} Use of analgesic techniques with different mechanisms to improve pain control while
reducing reliance on any single drug, especially opioids. Components may include
acetaminophen, nonsteroidal anti-inflammatory drugs, regional anesthesia, local
anesthetics, and selected adjuvants.

\medterm{Multimodal Imaging} Use of two or more imaging methods, such as CT, MRI, ultrasound, or nuclear imaging, to examine the same disease or body region. Combining methods can provide complementary structural, functional, or molecular information.

\medterm{Multimorbidity} The coexistence of two or more chronic health conditions in the same person, without
requiring one to be considered the primary disease. It complicates treatment, increases
polypharmacy risk, and requires coordinated patient-centered care.

\medterm{Multinodular Goiter} An enlarged thyroid containing multiple nodules. Thyroid function may be normal, underactive, or overactive; evaluation commonly includes thyroid-function tests and ultrasound, with further assessment of suspicious nodules or compressive symptoms.

\medterm{Multinucleated Hepatocyte} A multinucleated hepatocyte is a liver cell containing two or more nuclei, usually from nuclear division without complete cytokinesis. It may be a normal or reactive finding, but its significance depends on the number, morphology, and associated liver abnormalities.

\medterm{Multipara} A person who has had two or more pregnancies that reached the stage of potential fetal viability, regardless of whether the pregnancies resulted in live births.

\medterm{Multiparametric MRI Of The Prostate} A detailed prostate MRI that combines several types of images to show prostate structure, water movement, and blood flow. It helps identify suspicious areas and guide biopsy or treatment, but an abnormal scan does not by itself prove cancer.

\medterm{Multiparous} Having given birth two or more times to a fetus or fetuses beyond the stage of viability.

\medterm{Multiple Bacterial Abscesses} Several pockets of pus caused by bacterial infection in one or more organs. They may develop when bacteria spread through the blood or extend from nearby tissue and can cause fever, pain, or organ problems; imaging helps locate them.

\medterm{Multiple Biliary Hamartomas} Multiple small, harmless birth-related changes in the liver's bile ducts that look like tiny cysts on scans. They usually cause no symptoms and are often found by chance, but imaging helps distinguish them from cancer or other liver disease.

\medterm{Multiple Endocrine Neoplasia} An inherited condition causing tumors or overactivity in two or more hormone-making glands during a person's lifetime.

\medterm{Multiple Endocrine Neoplasia I} MEN type 1 is an inherited tumor syndrome involving parathyroid, pituitary, and pancreatic or duodenal neuroendocrine tumors, often caused by MEN1 variants.

\medterm{Multiple Endocrine Neoplasia II} MEN type 2 is an inherited syndrome with medullary thyroid cancer, pheochromocytoma, and sometimes parathyroid disease, commonly involving RET variants.

\medterm{Multiple Endocrine Neoplasia Type 1} An inherited condition that causes tumors in several hormone-making organs, especially the parathyroid glands, pituitary gland, and pancreas or nearby intestine. It is usually caused by a change in the MEN1 gene and requires lifelong monitoring.

\medterm{Multiple Endocrine Neoplasia Type 2} An inherited condition that greatly increases the risk of medullary thyroid cancer and may cause adrenal or parathyroid tumors. Type 2B can also cause nerve-related mouth bumps and distinctive body features; lifelong specialist monitoring is needed.

\medterm{Multiple Endocrine Neoplasia Type 2A} An inherited condition that greatly increases the risk of medullary thyroid cancer and can also cause adrenal tumors and overactive parathyroid glands. Regular genetic and specialist monitoring helps detect these problems early.

\synonyms
MEN 2A; MEN2A; Sipple Syndrome

\medterm{Multiple Endocrine Neoplasia Type 2B} An inherited RET-associated syndrome characterized by medullary thyroid
carcinoma, pheochromocytoma, mucosal neuromas, and a marfanoid habitus. The thyroid-cancer risk is high and requires
early specialist genetic and surgical assessment.

\medterm{Multiple Endocrine Neoplasia Type 4} A rare, autosomal dominant inherited tumor syndrome caused by germline mutations in the \textit{CDKN1B} gene, which encodes the p27$^{Kip1}$ tumor suppressor protein. It presents with features overlapping with MEN1, most commonly characterized by primary hyperparathyroidism (parathyroid adenomas causing high blood calcium) and anterior pituitary adenomas (such as prolactinomas or growth hormone-secreting tumors). Patients may also develop pancreatic neuroendocrine tumors, adrenal adenomas, and reproductive organ tumors. It is typically diagnosed in patients with MEN1-like symptoms who test negative for mutations in the \textit{MEN1} gene.

\synonyms
MEN 4; MEN4; Multiple Endocrine Neoplasia Type IV; CDKN1B-Related MEN

\medterm{Multiple Exostoses} Multiple exostoses, or hereditary multiple osteochondromas, is an inherited disorder causing multiple cartilage-capped bony projections near growth plates. Lesions can deform bones, restrict motion, compress nerves or vessels, and rarely undergo malignant transformation.

\medterm{Multiple Fibrofolliculomas} Numerous harmless small skin-colored growths arising from hair follicles, usually on the face and trunk. They can be a sign of Birt-Hogg-Dubé syndrome, which also increases the risk of kidney tumors and collapsed lung.

\medterm{Multiple Inherited Schwannomas} Alternate name; see \textit{Neurofibromatosis Type 2}.

\medterm{Multiple Mononeuropathy} Damage to several separate nerves in different parts of the body, often developing unevenly. It can cause patchy pain, numbness, weakness, or muscle wasting and may result from blood-vessel inflammation, diabetes, infection, or another systemic disease.

\medterm{Multiple Myeloma} A cancer of plasma cells that usually develops in the bone marrow and often produces a monoclonal immunoglobulin. It can
cause bone lesions, anemia, kidney dysfunction, high calcium, infection, or other organ
damage.

\synonyms
Plasma Cell Myeloma; Kahler Disease; Myelomatosis

\medterm{Multiple Myeloma: Deadly To Chronic} Multiple myeloma is a plasma-cell cancer that produces a monoclonal protein and can cause bone lesions, anemia, kidney injury, high calcium, infections, and fatigue. It is often treatable as a chronic illness, though relapse and complications remain possible.

\medterm{Multiple Organ Dysfunction Syndrome} Abnormal function of two or more organs during severe illness, such as sepsis, shock, trauma, or pancreatitis. Organ function may worsen over time and requires close monitoring and treatment of the underlying cause.

\medterm{Multiple Organ Failure} A life-threatening condition in which two or more organs stop working properly, usually during severe infection, shock, injury, or widespread inflammation. Intensive care supports the organs while doctors treat the underlying cause.

\medterm{Multiple Pregnancy} A pregnancy involving two or more fetuses. Twins may develop from two separate eggs or from one fertilized egg that divides; triplets and higher-order pregnancies are less common. Multiple pregnancies have higher risks of premature birth, growth problems, and pregnancy complications and require closer monitoring.

\medterm{Multiple Sclerosis} A chronic immune-mediated disease of the brain, spinal cord, and optic nerves. The immune system mistakenly attacks myelin, the protective covering around nerve fibers, and may also damage the nerve fibers themselves. This creates areas of inflammation and scarring, called lesions, that slow or block nerve signals. Symptoms vary and may include vision changes, numbness, weakness, stiffness, poor coordination, balance problems, fatigue, thinking difficulties, or bladder and bowel problems. Some people have attacks followed by partial or complete recovery, while others develop steadily worsening symptoms; mixed patterns also occur.

\synonyms
MS; Disseminated Sclerosis

\medterm{Multiple Sclerosis MS And Pregnancy} Multiple sclerosis usually does not prevent pregnancy, but symptoms, disease-modifying medicines, relapse risk, delivery planning, and postpartum care require coordination. Some treatments must be stopped or changed, and relapse risk can shift after delivery.

\medterm{Multiple Sclerosis MS Early Signs And Types} Multiple sclerosis is an immune-mediated disease of the central nervous system with episodes or progression of visual, sensory, motor, balance, bladder, cognitive, or fatigue symptoms. Relapsing, progressive, and other patterns guide testing and treatment.

\medterm{Multiple Sleep Latency Test} A daytime sleep study measuring how quickly a person falls asleep during several scheduled nap opportunities. It helps evaluate excessive daytime sleepiness and narcolepsy after nighttime sleep has been assessed.

\medterm{Multiple System Atrophy} A progressive brain disorder affecting movement and automatic body functions. It may cause stiffness and slowness, poor balance, fainting from low blood pressure, bladder problems, speech changes, or difficulty swallowing.

\synonyms
Shy-Drager Syndrome; MSA

\medterm{Multiple System Atrophy Cerebellar Type} An adult-onset neurodegenerative synucleinopathy (MSA-C) characterized primarily by progressive cerebellar ataxia (gait instability, limb dysmetria, dysarthria, oculomotor abnormalities) accompanied by severe autonomic failure and variable pyramidal signs.

\synonyms
MSA-C; Cerebellar Multiple System Atrophy; Olivopontocerebellar Atrophy Variant

\medterm{Multiple System Atrophy Parkinsonian Type} A progressive neurodegenerative disorder (MSA-P) characterized predominantly by parkinsonian features (rigidity, akinesia, postural instability) that respond poorly to levodopa, combined with early severe autonomic failure (orthostatic hypotension, urinary incontinence) and variable pyramidal dysfunction.

\synonyms
MSA-P; Parkinsonian Multiple System Atrophy; Striatonigral Degeneration Variant

\medterm{Multiple System Failure} Failure of two or more organ systems, also called multiple-organ failure or multiple-organ dysfunction syndrome; it may occur during severe illness such as sepsis, shock, or major trauma.

\medterm{Multiple Vitamin Overdose} Taking too many multivitamins can cause poisoning, especially from iron, vitamin A, or vitamin D. Symptoms may include vomiting, abdominal pain, confusion, liver injury, kidney injury, or dangerous calcium levels; urgent advice is needed.

\medterm{Multiplex PCR Panel} A laboratory test that looks for genetic material from several germs in one sample. Different panels test respiratory, intestinal, blood, or other specimens; a detected germ may represent infection, harmless carriage, or material left after an infection.

\synonyms
Multiplex Polymerase Chain Reaction; Syndromic Molecular Panel

\medterm{Multiplicity} Multiplicity is the state of having more than one occurrence, copy, lesion, tumor, or pregnancy component. Its interpretation depends on context, such as multiple gestation, multifocal disease, or copy number in genetics.

\medterm{Multipolar Neuron} A nerve cell with one long output branch and several receiving branches. This common type of neuron includes many cells that control movement or pass signals within the brain and spinal cord.

\medterm{Multisystem Inflammatory Syndrome In Children} A serious delayed inflammatory illness occurring in some children after infection with SARS-CoV-2, involving fever and two or more organ systems and potentially causing shock or cardiac dysfunction.

\medterm{Multivesicular Body} A multivesicular body is an endosomal compartment containing many small intraluminal vesicles. It can fuse with the plasma membrane to release exosomes or with a lysosome to degrade its cargo.

\medterm{Mumps} A contagious viral infection caused by a paramyxovirus and spread through respiratory droplets, saliva, and direct contact. It typically begins with fever, headache, and fatigue, followed by tender swelling of one or both parotid salivary glands (parotitis) below and in front of the ears, making chewing and swallowing painful. Complications are more common in adolescents and adults and include orchitis (painful testicular inflammation that can cause atrophy and reduced fertility), viral meningitis, pancreatitis, and permanent sensorineural hearing loss. Treatment is supportive with rest, fluids, and pain relief. It is highly preventable with the standard two-dose Measles-Mumps-Rubella (MMR) vaccine.
\synonyms{Epidemic Parotitis; Mumps Parotitis; Paramyxoviral Parotitis}

\medterm{Mumps Virus} The virus that causes mumps, a contagious infection spread through saliva and respiratory droplets. It commonly causes fever and painful swelling near the jaw, and can sometimes inflame the brain, pancreas, or reproductive organs; vaccination reduces risk.

\medterm{Munchausen Syndrome} Alternate name; see \textit{Factitious Disorder}.

\medterm{Mupirocin} Mupirocin is a topical antibiotic that inhibits bacterial isoleucyl-tRNA synthetase. It is used for selected skin infections and nasal eradication of certain staphylococci; resistance can develop with inappropriate or prolonged use.

\medterm{Mural} Relating to a wall, especially the wall of a body cavity, organ, or blood vessel.

\medterm{Mural Thrombosis} Formation of a blood clot along the wall of a heart chamber or blood vessel rather than freely within the lumen.

\medterm{Mural Thrombus} A mural thrombus is a blood clot that forms on the wall of a heart chamber or large
blood vessel, typically in areas of damaged endothelium, stasis, or turbulent flow.

\medterm{Muramidase} Muramidase, or lysozyme, is an enzyme that cleaves the peptidoglycan backbone of bacterial cell walls. It is present in secretions and innate immune cells and contributes to defense against susceptible bacteria.

\medterm{Murine Typhus} An acute endemic rickettsial illness caused by Rickettsia typhi and transmitted to humans primarily through the feces of infected rodent fleas (Xenopsylla cheopis). It typically presents with high prolonged fever, severe frontal headache, myalgias, arthralgias, and a maculopapular rash that typically spares the face, palms, and soles.

\synonyms
Endemic Typhus; Flea-Borne Typhus

\medterm{Murmur} An extra sound heard during a heartbeat, usually caused by fast or turbulent blood flow through the
heart or nearby blood vessels. Some murmurs are harmless, while others may be caused by heart valve or
structural problems.

\medterm{Murmur Classification} Heart murmurs are classified by timing, grade, shape, location, radiation, and pitch.
Timing: systolic (between S1 and S2), diastolic (between S2 and S1), or continuous.
Grade: I (barely audible) to VI (heard without stethoscope). Shape: crescendo,
decrescendo, or plateau.

\medterm{Murphy Sign} A sudden catch or stop in breathing when a doctor presses under the right ribcage as the patient takes a deep breath in, indicating an inflamed gallbladder.

\medterm{Murphy's Sign} Inspiratory arrest caused by pain when the right upper quadrant is palpated, traditionally suggesting acute cholecystitis.

\medterm{Musca} A fly genus including the housefly; some species can spread germs by carrying them on their bodies.

\medterm{Muscae Volitantes} Moving spots, threads, or cobweb-like shapes seen in the visual field, usually caused by shadows from vitreous opacities.

\medterm{Muscarine} Muscarine is a mushroom-derived agonist of muscarinic acetylcholine receptors. Ingestion of muscarine-containing mushrooms can cause a cholinergic syndrome with salivation, sweating, tearing, abdominal cramps, vomiting, diarrhea, bronchospasm, and bradycardia.

\medterm{Muscarinic Antagonist} A muscarinic antagonist blocks acetylcholine at muscarinic receptors in smooth muscle, glands, the heart, and the central nervous system. Examples include atropine, ipratropium, and oxybutynin; effects can include dry mouth, blurred vision, urinary retention, tachycardia, and confusion.

\medterm{Muscidae} Muscidae is a family of flies that includes the housefly. Some species can mechanically carry pathogens to food or wounds, and their larvae may occasionally cause myiasis.

\medterm{Muscimol} Muscimol is a psychoactive compound found in some Amanita mushrooms and an agonist at GABA-A receptors. Ingestion can cause altered perception, sedation, agitation, incoordination, or seizures and may require toxicologic evaluation.

\medterm{Muscle} Body tissue that shortens to create force and movement. Skeletal muscle moves the bones and supports posture, heart muscle pumps blood, and smooth muscle moves substances through organs such as the bowel and blood vessels.

\medterm{Muscle Aches} Pain or soreness in one or more muscles. Common causes include overuse, injury, infection, medicines, and conditions affecting the whole body.

\medterm{Muscle Atrophy} Loss or wasting of muscle tissue caused by disuse, nerve injury, inflammation, malnutrition, or systemic disease.
It may lead to weakness and reduced function; treatment focuses on the cause, nutrition, and
appropriately supervised rehabilitation.

\medterm{Muscle Biopsy} Removal of a small sample of skeletal muscle for microscopic, histochemical, immunologic, or molecular examination in suspected myopathy, inflammation, infection, or infiltrative disease.

\medterm{Muscle Cell} A cell specialized to contract; skeletal, cardiac, and smooth muscle cells perform different movement and organ functions.

\medterm{Muscle Contraction} The tightening of a muscle, allowing it to produce force or movement. Muscle fibers use calcium and energy from ATP to make their protein filaments slide against each other.

\medterm{Muscle Contracture} Persistent shortening or tightening of a muscle or nearby tissues that limits joint movement and does not fully relax with stretching. It may follow scarring, prolonged spasticity, immobility, injury, or, more rarely, a muscle-energy disorder.

\synonyms
Skeletal Muscle Contracture; Metabolic Contracture

\medterm{Muscle Cramp} A muscle cramp is a sudden, involuntary, painful contraction of a muscle. It may follow exertion, dehydration, electrolyte disturbance, pregnancy, medication use, nerve disease, or vascular disease; recurrent cramps need assessment of the cause.

\synonyms
Charley Horse; Involuntary Muscle Spasm

\medterm{Muscle Development} Muscle development is the formation, growth, differentiation, and maturation of muscle tissue. It depends on genetic programs, neural input, hormones, nutrition, mechanical loading, and repair processes.

\medterm{Muscle Disorder} A disease affecting muscle structure or function, causing weakness, pain, cramps, stiffness, wasting, or abnormal movement. Causes include inherited conditions, inflammation, nerve disease, medicines, infection, and metabolic problems.

\medterm{Muscle Fatigue} A reduced ability of a muscle to produce force after prolonged or repeated activity. It may result from exertion, poor sleep, illness, medicines, nerve problems, or muscle disease, and improves when the cause is treated.

\medterm{Muscle Fiber} A muscle fiber is an individual muscle cell containing contractile myofibrils. Skeletal muscle fibers vary in size and metabolic properties, while cardiac and smooth muscle cells have distinct structures and functions.

\medterm{Muscle Fibers} Individual muscle cells containing contractile proteins. Skeletal muscle fibers vary in size and properties, including endurance-oriented slow fibers and powerful, faster-fatiguing fibers.

\synonyms
Muscle Cells; Myocytes

\medterm{Muscle Function Loss} Reduced ability of a muscle or muscle group to generate movement or force, caused by muscle disease, nerve injury, neuromuscular-junction disorders, pain, or disuse.

\medterm{Muscle Hypertrophy} Muscle hypertrophy is enlargement of individual muscle fibers, usually in response to resistance, hormonal stimulation, or compensatory workload. It can be physiologic or associated with disease, depending on the muscle and cause.

\medterm{Muscle Neoplasm} A muscle neoplasm is an abnormal growth in skeletal or smooth muscle or nearby soft tissue and may be benign or malignant.

\medterm{Muscle Pain} Pain, soreness, or aching in one or more muscles caused by exertion, injury, infection, inflammation, medicines, or an underlying medical condition. Persistent, severe, or weakness-associated pain needs assessment.

\synonyms
Myalgia; Muscle Ache; Myodynia

\medterm{Muscle Protein} A protein found in muscle that provides structure, produces contraction, stores or transports substances, or regulates muscle activity. Important examples include actin, myosin, troponin, dystrophin, and myoglobin.

\medterm{Muscle Relaxant} A muscle relaxant is a medicine that reduces skeletal-muscle spasm or contraction. Centrally acting agents reduce muscle spasm through the nervous system, whereas neuromuscular blockers interrupt transmission at the neuromuscular junction and require respiratory support.

\medterm{Muscle Relaxants} Medicines that reduce skeletal-muscle spasm or neuromuscular transmission; they include centrally acting drugs and peripheral neuromuscular blockers.

\medterm{Muscle Rigidity} Muscle rigidity is abnormal resistance to passive movement caused by sustained muscle contraction or altered tone. It occurs in conditions such as Parkinson disease, dystonia, meningitis, neuroleptic malignant syndrome, and tetanus.

\medterm{Muscle Spasm} A sudden, involuntary tightening of a muscle that may cause pain, stiffness, or an abnormal position. Spasms can follow overuse, dehydration, mineral changes, nerve problems, medicines, pregnancy, or injury.

\medterm{Muscle Spasms} Muscle spasms are sudden involuntary contractions caused by exertion, dehydration, electrolyte disturbance, nerve disease, medicines, pregnancy, or musculoskeletal problems. Recurrent, widespread, painful, or weakness-associated spasms need evaluation for an underlying cause.

\medterm{Muscle Spindle} A small sensor inside skeletal muscle that detects how much the muscle is stretched and how quickly it changes length. It helps maintain posture, muscle tone, balance, and coordinated movement.

\medterm{Muscle Strain Treatment} Care for a stretched or torn muscle, usually involving temporary activity reduction, gradual movement, rehabilitation, and pain relief chosen safely for the person. Severe swelling, weakness, deformity, inability to use the limb, or persistent pain needs assessment.

\medterm{Muscle Strength} Muscle strength is the maximal force a muscle or muscle group can generate under specified conditions. It is assessed by functional or resisted testing and may be reduced by muscle, nerve, neuromuscular-junction, central, or systemic disorders.

\medterm{Muscle Tissue} Specialized tissue that contracts to produce movement, maintain posture, move substances through organs, and generate heat. Skeletal, heart, and smooth muscle have different structures and control systems for their jobs.

\medterm{Muscle Tonus} The slight, continuous activity present in resting muscles, which gives them resistance to being stretched. Too much or too little tone can occur with disorders of the brain, spinal cord, nerves, or muscles.

\medterm{Muscle Twitching} Brief involuntary contractions of a small group of muscle fibers, often called fasciculations; persistent or widespread twitching may require evaluation.

\medterm{Muscle Weakness} Reduced ability to generate force with one or more muscles. Causes include nerve or muscle disease, pain, inactivity,
medicines, metabolic disorders, and aging.

\medterm{Muscle-Contraction Headache} A pressure-like headache caused by prolonged tightening or sensitivity of muscles in the scalp, neck, or jaw. Stress, poor posture, fatigue, or jaw clenching may contribute; frequent, severe, or changing headaches need medical assessment.
\synonyms
Tension-Type Headache; Stress Headache

\medterm{Muscle-Invasive Bladder Cancer} Bladder cancer that has grown into the bladder muscle layer. It has a greater risk of spreading than cancer limited to the bladder lining and is classified by examination of bladder tissue.

\medterm{Muscular} Relating to muscle or made of muscle tissue. It may describe strength, body build, pain, weakness, or a structure containing muscle.

\medterm{Muscular Dystrophy} A group of more than 30 inherited muscle diseases that cause progressive weakness and breakdown of skeletal muscles. Different gene changes affect proteins that strengthen and protect muscle fibers, so the age of onset, muscles affected, and severity vary by type. Weakness may cause difficulty walking or doing everyday activities. Some forms also affect the heart, breathing muscles, swallowing, bones, brain, or other organs. Muscular dystrophy is not contagious and is not caused by ordinary injury or physical activity.

\synonyms
MD; Inherited Myopathy

\medterm{Muscular Endurance} The ability of a muscle or muscle group to keep working during repeated or prolonged effort. It depends on muscle strength, energy supply, blood flow, nerve control, conditioning, and overall health.

\medterm{Muscularis} A layer of smooth muscle in the wall of a hollow organ, especially the digestive tract. Its coordinated contractions move contents through the organ and can also change the diameter or pressure of the passage.
\synonyms
Muscular Layer; Tunica Muscularis

\medterm{Musculocutaneous Nerve} A nerve from the neck and upper chest that controls several muscles used to bend the elbow and turn the forearm. It also carries sensation from part of the outer forearm; injury can cause weakness or numbness there.

\medterm{Musculoskeletal} Relating to muscles, bones, joints, tendons, ligaments, and related tissues.

\medterm{Musculoskeletal Deformity} A musculoskeletal deformity is an abnormal shape, alignment, or structure of bone, joint, muscle, or related tissue. It may be congenital, developmental, traumatic, degenerative, neuromuscular, or caused by disease.

\medterm{Musculoskeletal Disease} A disorder affecting bones, joints, muscles, tendons, ligaments, cartilage, or other supporting tissues, causing pain or limited movement.

\medterm{Musculoskeletal Injury} Trauma to muscles, bones, joints, tendons, ligaments, cartilage, or related soft tissues, ranging from strains and sprains to fractures and dislocations.

\medterm{Musculoskeletal Pain} Pain arising from muscles, bones, joints, ligaments, tendons, or related soft tissues. Common causes
include injury, overuse, arthritis, and inflammation; persistent pain, weakness, swelling, or
neurologic symptoms should be assessed.

\medterm{Musculoskeletal System} The musculoskeletal system consists of bones, joints, cartilage, ligaments, tendons, muscles, and associated connective tissues. It provides support, protects organs, enables movement, and contributes to mineral storage and blood-cell production.

\medterm{Musculoskeletal Tumor} An abnormal growth arising in bone, muscle, cartilage, connective tissue, or related structures. It may be benign or malignant and can cause a lump, pain, swelling, reduced function, or fracture.

\medterm{Musculus} The Latin word for muscle, used in formal anatomical names. For example, ``musculus deltoideus'' means the deltoid muscle. Abbreviations may be used in anatomy texts to identify one or several muscles.

\medterm{MUSE} A medicated urethral system that delivers a pellet into the urethra to treat erectile dysfunction.
\synonyms
Medicated Urethral System for Erection

\medterm{Mushroom} The visible fruiting body of a fungus. Edible mushrooms can provide nutrients, but wild mushrooms may be poisonous or confused with safe species; uncertain mushrooms should not be eaten.

\medterm{Mushroom Poisoning} Illness caused by eating toxic fungi, with effects ranging from gastrointestinal upset to liver failure, neurologic toxicity, or death.

\medterm{Music Therapy} The planned use of music by a trained therapist to support health and wellbeing. It may help with pain, stress, communication, movement, memory, emotional expression, or coping during illness and rehabilitation.

\medterm{MuSK-Positive Myasthenia Gravis} A form of myasthenia gravis in which antibodies interfere with nerve-to-muscle communication. It often causes weakness of the face, mouth, throat, neck, or breathing muscles and requires specialist care.

\synonyms
MuSK Myasthenia Gravis; MuSK-Ab-Positive MG

\medterm{Mustard Compounds} Mustard compounds are chemical agents containing reactive sulfur or nitrogen mustard groups that can alkylate DNA and other molecules. Some are chemotherapy drugs; others are highly toxic chemical-warfare agents.

\medterm{Mustard Gas} Mustard gas, or sulfur mustard, is a vesicant chemical-warfare agent that damages skin, eyes, and respiratory tissues and can cause blistering, airway injury, and delayed complications. It is not a true gas at ordinary temperatures and has no specific widely available antidote.

\medterm{Mutagen} A chemical, physical agent, or biological factor capable of causing a genetic mutation. Mutagenicity does not by itself prove that an agent causes cancer.

\medterm{Mutant} A mutant is a cell, organism, or genetic sequence carrying a mutation relative to a reference or parental form. The change may have no effect, alter a trait, or cause disease, depending on its location and functional consequence.

\medterm{Mutant Proteins} Proteins whose amino-acid sequence differs because of a genetic mutation. The change may have no effect or may alter folding, stability, location, or function and contribute to inherited disease, cancer, or drug resistance.

\medterm{Mutate} To mutate means to undergo a change in genetic sequence. Mutations arise spontaneously or from replication errors, DNA damage, or mutagen exposure and may be inherited, acquired in somatic cells, harmful, neutral, or beneficial.

\medterm{Mutation} A heritable or acquired alteration in the DNA sequence. Variants may be
single-nucleotide substitutions, insertions, deletions, copy-number changes, or larger
rearrangements and may be benign, pathogenic, or of uncertain significance.

\medterm{Mutation Abnormality} A mutation abnormality is an atypical change in DNA sequence or chromosome structure identified relative to a reference. Its clinical significance must be classified using evidence, such as pathogenic, likely pathogenic, uncertain significance, likely benign, or benign.

\medterm{Mutation Rate} The frequency at which new DNA changes arise in a cell, organism, or generation. The rate differs between genes and species and can be influenced by DNA-copying errors, repair systems, radiation, chemicals, or other factors.

\medterm{Mute} Describing a person who does not speak, whether because of a physical, neurologic, developmental, or psychological reason or by choice; current clinical language should describe the specific communication limitation or preference.

\medterm{Mutism} Absence or refusal of speech despite adequate ability to speak, occurring in conditions such as catatonia, selective mutism, or neurologic disease.

\medterm{Myalgia} Muscle pain or aching. It may affect one area after injury or exercise, or occur throughout the body with infection, inflammation, medicine effects, thyroid disease, or other conditions. Severe pain with weakness or dark urine needs prompt medical assessment because muscle breakdown may be occurring.

\synonyms
Myodynia; Muscle Pain

\medterm{Myalgic Encephalomyelitis} A complex illness characterized by substantial reduction in activity, post-exertional
malaise, and unrefreshing sleep, often with cognitive or autonomic symptoms.

\medterm{Myalgic Encephalomyelitis} A complex illness involving substantial activity-limiting fatigue, post-exertional worsening, unrefreshing sleep, and often cognitive or orthostatic symptoms.

\synonyms
ME/CFS; SEID; Systemic Exertional Intolerance Disease

\medterm{Myalgic Encephalomyelitis/Chronic Fatigue Syndrome} A complex illness characterized by substantial reduction in
activity, post-exertional malaise, unrefreshing sleep, and cognitive or autonomic symptoms. Management is individualized
and focuses on pacing, symptom relief, and treatment of coexisting conditions.

\medterm{Myasthenia} Muscle weakness that often becomes more pronounced with activity and lessens after rest. The term may refer generally to this pattern of weakness or to specific disorders such as myasthenia gravis.

\medterm{Myasthenia Gravis} A chronic autoimmune disorder of the neuromuscular junction, where nerves communicate with skeletal muscles. Antibodies most often target acetylcholine receptors and sometimes proteins such as muscle-specific kinase (MuSK) or LRP4, disrupting signals that cause muscles to contract. Weakness typically worsens with repeated use and improves with rest. It may affect the eyelids and eye movements, facial expression, chewing, speech, swallowing, neck, arms, legs, or breathing muscles. Sensation is usually preserved; thymus abnormalities, including thymoma, may also occur.

\synonyms
MG

\medterm{Myasthenic Crisis} A life-threatening worsening of myasthenia gravis that weakens the muscles needed for breathing or swallowing. Infection, surgery, stress, pregnancy, or certain medicines may trigger it. Severe breathlessness, choking, inability to swallow, or rapidly worsening weakness requires emergency care, breathing support, and specialist treatment.

\synonyms
Myasthenia Gravis Crisis; Acute Myasthenic Exacerbation

\medterm{Myasthenic Syndrome} A disorder of neuromuscular transmission causing fluctuating weakness; it may refer to myasthenia gravis or a congenital or Lambert-Eaton syndrome depending on context.

\medterm{Mycelium} The network of fine branching threads that makes up the growing part of a fungus. It spreads through soil, food, or tissue and can produce the structures that form fungal spores.

\medterm{Mycetoma} A chronic localized infection of skin and subcutaneous tissue caused by fungi or filamentous bacteria, producing swelling, sinuses, and grains.

\medterm{Mycobacteriaceae} A family of slow-growing or rapidly growing, acid-fast bacteria that includes Mycobacterium species causing tuberculosis, leprosy, and nontuberculous infections.

\medterm{Mycobacterial Culture} A laboratory test that tries to grow Mycobacterium bacteria from a sample such as sputum, tissue, or fluid. It can identify tuberculosis or other mycobacteria and may show which medicines will work, but results can take weeks.

\medterm{Mycobacterial Infection} Infection caused by bacteria in the genus Mycobacterium. It includes tuberculosis, leprosy, and nontuberculous mycobacterial disease; symptoms and treatment depend on the species and affected organs.

\medterm{Mycobacterium} A genus of acid-fast bacteria that includes the causes of tuberculosis and leprosy and several nontuberculous infections.

\medterm{Mycobacterium Abscessus} A rapidly growing nontuberculous mycobacterium that can cause pulmonary, skin, and
soft-tissue infections and is often difficult to treat because of multidrug resistance.

\medterm{Mycobacterium Avium} A member of the Mycobacterium avium complex that can cause lung disease, lymph-node infection, or disseminated infection, especially in people with weakened immunity.

\medterm{Mycobacterium Avium Complex} A group of nontuberculous mycobacteria that can cause chronic pulmonary disease
or disseminated infection, particularly in people with impaired immunity. Treatment depends on the species, site,
severity, susceptibility, and patient factors.

\medterm{Mycobacterium Bovis} A tuberculosis-complex bacterium that mainly infects cattle but can cause tuberculosis in humans through unpasteurized dairy or animal exposure.

\medterm{Mycobacterium Chelonae} A rapidly growing nontuberculous mycobacterium that can cause skin, soft-tissue, bone, eye, or disseminated infection.

\medterm{Mycobacterium Fortuitum} A rapidly growing nontuberculous mycobacterium that can cause wound, catheter, bone, lung, or other opportunistic infections.

\medterm{Mycobacterium Gordonae} A usually nonpathogenic water-associated nontuberculous mycobacterium that can occasionally cause disease in people with weakened immunity.

\medterm{Mycobacterium Haemophilum} A nontuberculous mycobacterium that can cause skin, joint, and bone infections, especially in people with impaired immunity.

\medterm{Mycobacterium Infection} Infection caused by a Mycobacterium species, including tuberculosis, leprosy, and nontuberculous mycobacterial disease.

\medterm{Mycobacterium Kansasii} A nontuberculous mycobacterium that commonly causes chronic lung disease resembling tuberculosis and can occasionally cause disseminated infection.

\medterm{Mycobacterium Leprae} The bacterium that causes leprosy, a long-lasting infection affecting the skin and nerves. It usually spreads only after prolonged close contact, develops slowly, and can cause numbness, weakness, or skin lesions without treatment.

\medterm{Mycobacterium Marinum} A nontuberculous mycobacterium associated with fish tanks and aquatic exposure that can cause
chronic skin and soft-tissue infection.

\medterm{Mycobacterium Scrofulaceum} A nontuberculous mycobacterium associated especially with cervical lymph-node infection in children and occasional other disease.

\medterm{Mycobacterium Szulgai} A rare nontuberculous bacterium that can cause long-lasting lung disease or other infections, especially in people with underlying lung problems or weakened immunity. Diagnosis requires specialist laboratory testing and clinical assessment.

\medterm{Mycobacterium Tuberculosis} The bacterium that causes tuberculosis. It spreads through the air and most often affects the lungs, but can involve other organs. Latent infection has no symptoms and does not spread; active disease may cause cough, fever, night sweats, or weight loss.

\medterm{Mycobacterium Ulcerans} A mycobacterium that causes Buruli ulcer, a chronic destructive skin infection associated with the toxin mycolactone.

\medterm{Mycology} The branch of microbiology that studies fungi, including yeasts, moulds, and mushrooms.
Medical mycology focuses on fungal pathogens that cause disease in humans. Fungal
infections (mycoses) are classified as superficial, cutaneous, subcutaneous, or
systemic.

\medterm{Mycophenolate Mofetil} Mycophenolate mofetil is a prodrug of mycophenolic acid that inhibits inosine monophosphate dehydrogenase and lymphocyte proliferation. It is used as an immunosuppressant, especially after organ transplantation, and can cause infection, cytopenias, gastrointestinal effects, and serious fetal harm.

\medterm{Mycophenolic Acid} The active immune-suppressing substance formed from mycophenolate medicines. It blocks an enzyme needed by multiplying lymphocytes and helps prevent organ-transplant rejection, but can cause infections, low blood counts, diarrhea, and fetal harm.

\medterm{Mycoplasma} A genus of very small bacteria lacking a cell wall, including organisms that cause respiratory and genitourinary infections.

\medterm{Mycoplasma Genitalium} A sexually transmitted bacterium that can cause urethritis, cervicitis, pelvic
inflammatory disease, and, less commonly, proctitis. Many infections are asymptomatic.
Diagnosis is usually by a nucleic acid amplification test, and macrolide resistance is
common, so treatment should follow current susceptibility-informed guidance.

\medterm{Mycoplasma Hominis} A wall-less bacterium associated with the genital tract that can cause pelvic, postpartum, urinary, wound, or joint infection.

\medterm{Mycoplasma Pneumonia} Atypical pneumonia most common in ages 5--35. Dry cough, low-grade fever, headache.
Patchy infiltrates on CXR disproportionate to symptoms (``walking pneumonia'').

\medterm{Mycoplasma Pneumoniae} A bacterium without a cell wall that can cause bronchitis, sore throat, or pneumonia, especially in school-age children and young adults. Illness often develops gradually with cough, fever, headache, or tiredness; complications outside the lungs are uncommon but possible.

\medterm{Mycoplasma Pneumoniae-Associated Mucositis} Inflammatory lesions of the oral, ocular, or genital mucosa occurring as an
extrapulmonary manifestation of Mycoplasma pneumoniae infection, often presenting with
target-like skin lesions.

\medterm{Mycoplasmal Pneumonia} Pneumonia caused by Mycoplasma species, most often Mycoplasma pneumoniae, usually producing a gradual cough and fever.

\medterm{Mycoplasmataceae} A family of very small bacteria that lack a cell wall and includes Mycoplasma and related genera.

\medterm{Mycosis} An infection caused by a fungus, such as yeast or mold. Fungal infections may affect the skin, nails, mouth, lungs, or other organs, depending on the type.

\medterm{Mycosis Fungoides} The most common cutaneous T-cell lymphoma, presenting with progressive patches, plaques,
and tumors on the skin with an indolent course. Staging and therapy depend on skin
involvement and extracutaneous spread.

\synonyms
MF; Cutaneous T-Cell Lymphoma

\medterm{Mycosis Fungoides Lymphoma} Mycosis fungoides is a type of cutaneous T-cell lymphoma that usually begins with persistent patches or plaques on the skin.

\medterm{Mycotic Aneurysm} An infected arterial wall or aneurysm, usually caused by bacteria rather than fungi. It may develop from bloodstream infection or an infected embolus and can rupture, causing life-threatening hemorrhage.

\medterm{Mycotoxin} A poisonous substance made by some fungi that can contaminate food, crops, or indoor environments. Exposure may damage the liver, kidneys, nervous system, immune system, or other organs depending on the toxin and dose.

\medterm{Mydriasis} Dilation of the pupil, produced mainly by sympathetic activity through the dilator pupillae muscle or by
reduced parasympathetic activity. Causes include darkness, medicines, toxins, neurologic
disease, and ocular injury.

\medterm{Mydriatic Agents} A medicine that widens the pupil. Some, such as tropicamide or atropine, also temporarily prevent focusing; others, such as phenylephrine, mainly widen the pupil. These medicines are used to examine the inside of the eye and sometimes to ease pain and prevent scarring during eye inflammation. They can trigger a dangerous rise in eye pressure in people with narrow drainage angles.

\medterm{Myelin} Myelin is an insulating sheath around nerve fibers that speeds electrical conduction; it is produced by oligodendrocytes in the CNS and Schwann cells in the PNS.

\medterm{Myelin Sheath} The insulating covering around many nerve fibers that allows electrical signals to travel quickly. It is made by different support cells in the brain, spinal cord, and peripheral nerves and can be damaged in disorders such as multiple sclerosis.

\medterm{Myelination} Myelination is formation of a myelin sheath around an axon by oligodendrocytes in the central nervous system or Schwann cells in the peripheral nervous system. It increases the speed and efficiency of nerve conduction.

\medterm{Myelitis} Myelitis is inflammation of the spinal cord or, less commonly, other neural tissue described by the term. Transverse myelitis can cause weakness, sensory changes, and bladder or bowel dysfunction and requires urgent evaluation for infectious, immune, vascular, or compressive causes.

\medterm{Myeloablative Agonist} This term likely means a myeloablative agent: a treatment that severely or completely destroys blood-forming cells in the bone marrow. It may be used before a stem-cell transplant and causes profound, temporary low blood counts.

\medterm{Myeloablative Chemotherapy} Myeloablative chemotherapy destroys or severely suppresses bone-marrow blood-cell production and is often used before a stem-cell transplant.

\medterm{Myeloblast} An immature precursor of granulocytes in the bone marrow; excess myeloblasts are characteristic of acute myeloid leukemia.

\medterm{Myelocyte} A developing granulocyte precursor in bone marrow that is more mature than a myeloblast and normally absent or rare in peripheral blood.

\medterm{Myelocytic Leukaemia} A blood cancer arising from cells that normally develop into several blood-cell types. It includes acute and chronic myeloid forms, which can cause anemia, infections, bleeding, fatigue, or an enlarged spleen.

\medterm{Myelodysplastic Neoplasms} Alternate name; see \textit{Myelodysplastic Syndromes}.

\medterm{Myelodysplastic Or Myeloproliferative Neoplasms} Myelodysplastic and myeloproliferative neoplasms are clonal bone-marrow disorders causing abnormal blood-cell production. They may produce anemia, infection, bleeding, high counts, enlarged spleen, thrombosis, or progression to acute leukemia; classification and mutation testing guide care.

\medterm{Myelodysplastic Syndrome} A group of bone-marrow cancers in which blood cells develop abnormally and are produced poorly. It can cause anemia, infections, or bleeding and may progress to acute leukemia; diagnosis uses blood and bone-marrow testing.

\medterm{Myelodysplastic Syndromes} Alternate name; see \textit{Myelodysplastic Syndromes}.

\medterm{Myeloencephalitis} Myeloencephalitis is inflammation involving the spinal cord and brain, caused by infection, immune disease, toxins, or other neurologic processes.

\medterm{Myelofibrosis} A bone-marrow cancer in which scar tissue develops in the marrow and disrupts blood-cell production. It may cause anemia, enlarged spleen, fatigue, fever, night sweats, weight loss, or abnormal blood counts.

\synonyms
Idiopathic Myelofibrosis; Primary Myelofibrosis

\medterm{Myelogenous} Relating to blood-forming cells that develop from the bone marrow's myeloid cell line.

\synonyms
Myeloid

\medterm{Myelogram} An imaging study of the spinal canal performed after contrast material is introduced around the spinal cord.

\medterm{Myelography} Contrast injected into subarachnoid space for spinal imaging. Used when MRI
contraindicated. CT myelography more common now.

\medterm{Myeloid} Relating to bone marrow or to the non-lymphoid blood-cell lineages derived from hematopoietic stem cells.

\medterm{Myeloid Cell} A blood-forming cell from the myeloid lineage, which produces red cells, platelets, and several white blood-cell types.

\medterm{Myeloid Leukemia} Myeloid leukemia is a blood cancer arising from abnormal myeloid precursor cells in the bone marrow and blood.

\medterm{Myeloid Sarcoma} Myeloid sarcoma is a mass of immature myeloid cells occurring outside the bone marrow, often associated with acute myeloid leukemia.

\medterm{Myelolipoma} A usually benign tumor composed of mature adipose tissue and blood-forming tissue, most often arising in the adrenal gland and often found incidentally.

\medterm{Myeloma} A cancer of plasma cells, usually called multiple myeloma when it affects many bone-marrow sites. It can cause bone pain, anemia, kidney damage, high blood calcium, infections, or abnormal proteins in blood and urine.

\medterm{Myeloma Cast Nephropathy} Kidney damage caused by abnormal proteins from multiple myeloma blocking and injuring the kidney's filtering tubes. It can cause sudden kidney failure and requires urgent treatment of the myeloma and support of kidney function.

\medterm{Myelomalacia} Softening of the spinal cord, usually due to ischemia, trauma, hemorrhage, or other cord injury.

\medterm{Myelomeningocele} A birth defect in which the spinal cord and its protective coverings bulge through an opening in the spine. It can cause weakness or paralysis, loss of sensation, bladder or bowel problems, and fluid buildup around the brain.

\medterm{Myelomonocyte} A bone-marrow precursor that can develop toward either the monocyte or granulocyte lineages. An immature bone marrow cell that can develop into either monocytes or certain infection-fighting white blood cells.

\medterm{Myelopathy} A general term for any neurological deficit resulting from spinal cord pathology,
encompassing a wide range of causes and clinical presentations.

\medterm{Myeloperoxidase} An enzyme in neutrophil and monocyte granules that uses hydrogen peroxide and chloride to generate antimicrobial oxidants; antibodies against it occur in some autoimmune diseases.

\medterm{Myelopoiesis} The production and development of myeloid blood cells, including granulocytes, monocytes, and some precursor cells, in the bone marrow.

\medterm{Myeloproliferative Disorders} Clonal bone-marrow disorders causing overproduction of one or more blood-cell lineages, including polycythemia vera, essential thrombocythemia, and myelofibrosis.

\medterm{Myeloproliferative Neoplasm} A myeloproliferative neoplasm is a blood cancer in which the bone marrow produces too many mature or immature blood cells.

\medterm{Myelosuppression} Reduced bone marrow production of red cells, white cells, and platelets, causing anemia,
neutropenia, and thrombocytopenia; a common dose-limiting toxicity of chemotherapy and
radiation that increases infection and bleeding risk.

\medterm{Myelosuppressive Therapy} Myelosuppressive therapy reduces bone-marrow production of blood cells, sometimes intentionally during cancer treatment but with risk of anemia, infection, and bleeding.

\medterm{Myenteric Plexus} A network of enteric nerves between the circular and longitudinal muscle layers of the gastrointestinal tract that coordinates intestinal motility.

\medterm{Myhre Syndrome} A rare inherited connective-tissue disorder caused by variants in SMAD4 and characterized by short stature, stiff joints, distinctive facial features, hearing loss, and cardiovascular abnormalities.

\medterm{Myiasis} Infestation of living tissue by larvae (maggots) of various fly species, which feed on
host tissue and can cause pain, secondary infection, and tissue destruction. It may be
cutaneous, involving the skin and subcutaneous tissue, or affect other body sites
including the nose, ears, and wounds.

\medterm{Myoblast} An early muscle cell that develops and joins with others to form skeletal muscle fibers.

\medterm{Myoblastoma} A rare tumor made of cells resembling developing muscle cells. The term has been used for more than one tumor type, so diagnosis requires modern microscopic and molecular classification.
\synonyms
Granular Cell Tumor; Abrikossoff Tumor

\medterm{Myocardial} Relating to the myocardium, the muscular wall of the heart.

\medterm{Myocardial Biopsy} Removal of a small sample of heart muscle, usually by catheter, for microscopic or molecular examination of suspected myocarditis, rejection, infiltrative disease, or other cardiomyopathy.

\medterm{Myocardial Bridging} Myocardial bridging occurs when a segment of a coronary artery tunnels through heart muscle and may be compressed during contraction.

\medterm{Myocardial Contusion} A bruise of the heart muscle caused by blunt chest trauma, which may produce chest pain, arrhythmias, or impaired heart function.

\medterm{Myocardial Infarction} Ischemic necrosis (death) of heart muscle caused by an abrupt reduction or cessation of coronary blood supply, most commonly triggered by the rupture of an atherosclerotic plaque with overlying blood clot formation (thrombus). Symptoms include crushing retrosternal chest pain radiating to the jaw, neck, or left arm, shortness of breath, diaphoresis (sweating), and nausea. Diagnosis is confirmed by cardiac biomarkers (specifically elevated cardiac troponin levels) and electrocardiogram (ECG) changes, categorizing the event as either an ST-elevation myocardial infarction (STEMI, complete coronary occlusion) or non-ST-elevation myocardial infarction (NSTEMI, partial occlusion). Emergent treatment focuses on rapid coronary reperfusion via primary percutaneous coronary intervention (PCI/angioplasty) or clot-dissolving fibrinolytic therapy, along with antiplatelets, anticoagulants, beta-blockers, and statins.

\synonyms
Heart Attack; MI; Acute Myocardial Infarction; STEMI; NSTEMI

\medterm{Myocardial Ischemia} Reduced blood flow and oxygen supply to heart muscle, usually from narrowed coronary arteries, potentially causing chest pain or heart damage.

\synonyms
Cardiac Ischemia; Myocardial Hypoxia

\medterm{Myocardial Perfusion Imaging} A heart scan that compares blood flow to the heart muscle at rest and during exercise or medicine-induced stress. It can show areas with reduced blood supply or old heart-muscle damage.

\medterm{Myocardial Reperfusion} Restoration of blood flow to heart muscle after ischemia, which can limit injury but may also cause reperfusion-related tissue damage.

\medterm{Myocardial Rupture} A tear in the heart wall or other heart structure, usually a severe complication of myocardial infarction.

\medterm{Myocardial Viability} The ability of heart muscle weakened by poor blood flow to recover if blood supply is restored. Imaging tests can show whether the tissue is still alive and may help guide decisions about restoring blood flow.

\medterm{Myocarditis} Inflammation of the heart muscle, often after an infection or because of an immune reaction. It can cause chest pain, breathlessness, abnormal heart rhythms, reduced pumping, or sudden severe illness; severity ranges from mild inflammation to heart failure or shock.

\medterm{Myocarditis Diagnosis} Doctors assess suspected heart-muscle inflammation using symptoms, examination, an electrocardiogram, blood tests, heart imaging, and sometimes MRI or biopsy. No single test proves every case, so results are interpreted together and other causes are excluded.

\medterm{Myocarditis Treatment} Treatment depends on the cause and severity. Care may include rest from strenuous exercise, medicines for heart failure or abnormal rhythms, treatment of infection or autoimmune inflammation when appropriate, and hospital monitoring for severe cases.

\medterm{Myocardium} The thick muscular middle layer of the heart wall that contracts to pump blood throughout the body.

\medterm{Myoclonic Epilepsy} An epilepsy syndrome in which seizures include brief, shock-like involuntary muscle jerks; the specific syndrome determines associated seizure types, age of onset, and prognosis.

\medterm{Myoclonic Seizure} A generalized seizure characterized by brief, shock-like, involuntary single or multiple muscle contractions without loss of consciousness. Commonly affects the upper extremities, resulting in sudden dropping or throwing of objects, and is a hallmark feature of juvenile myoclonic epilepsy.

\medterm{Myoclonus} A sudden, brief, involuntary muscle jerk or group of jerks.

\medterm{Myocyte} A differentiated muscle cell specialized for contraction. Skeletal myocytes are long, multinucleated fibers; cardiac myocytes are branching cells joined by intercalated discs; and smooth-muscle myocytes are spindle-shaped cells without organized sarcomeres. Myocyte injury or death contributes to myopathy, rhabdomyolysis, and myocardial infarction.

\synonyms
Muscle Cell

\medterm{Myodegeneration} Degeneration or breakdown of muscle fibers, caused by inherited, inflammatory, ischemic, toxic, or other disorders.

\medterm{Myoepithelial Cell} A contractile cell around certain glands that helps push secretions from breast or salivary tissue.

\medterm{Myoepithelioma} A myoepithelioma is a tumor composed of myoepithelial cells and can be benign or malignant, commonly arising in salivary or other glands.

\medterm{Myoepithelioma Of Soft Tissue} A rare growth in soft tissue made of myoepithelial cells, which have features of both gland cells and muscle-like support cells. It may be harmless or cancerous; imaging and tissue examination determine treatment.

\medterm{Myofascial Pain Syndrome} A regional pain disorder involving sensitive trigger points in muscle and fascia, with pain that may be referred to other areas and often worsens with pressure or movement.

\medterm{Myofibril} A contractile thread-like structure inside a muscle fiber, composed of repeating sarcomeres containing actin and myosin.

\medterm{Myofibrillar Myopathy} Myofibrillar myopathy is a group of inherited muscle disorders with abnormal protein aggregates and progressive skeletal-muscle weakness, sometimes involving heart or breathing muscles.

\medterm{Myofibrils} Contractile structures inside muscle fibers made of repeating sarcomeres. Long contractile structures inside muscle fibers, made of repeating sarcomeres containing actin and myosin. Their shortening produces muscle contraction.

\medterm{Myofibroblastoma} A benign mesenchymal tumor composed of myofibroblastic cells in collagenous or myxoid stroma, most often arising in the mammary or inguinal region but potentially occurring at other soft-tissue sites.

\medterm{Myofibroblastoma Of The Breast} A rare benign stromal tumor composed of spindle-shaped myofibroblastic cells, usually occurring in older adults. It is typically a well-circumscribed painless mass and is diagnosed by histology with immunohistochemistry.

\medterm{Myofibroblasts} Cells that combine features of fibroblasts and smooth muscle cells and help contract wounds and produce scar tissue.

\medterm{Myofibromatosis} Myofibromatosis is a benign proliferation of myofibroblastic tumors, often in infants or children and involving skin, muscle, bone, or internal organs.

\medterm{Myofilament} A contractile protein filament within a muscle fiber; actin and myosin myofilaments interact through cross-bridges to produce muscle contraction.

\medterm{Myofilaments} Thin actin and thick myosin protein strands inside muscle cells. They slide past one another using energy from ATP, shortening the muscle and producing contraction.
\synonyms
Contractile Filaments; Actin and Myosin Filaments

\medterm{Myogenin} A muscle-specific transcription factor required for skeletal-muscle differentiation and used as an immunohistochemical marker in diagnosing some soft-tissue tumors.

\medterm{Myoglobin} An oxygen-binding protein stored in skeletal and heart muscle. When muscle is injured, myoglobin can enter the blood and urine; large amounts may contribute to kidney damage and are interpreted with other tests.

\medterm{Myoglobin Blood Test} A blood test measuring myoglobin released from injured skeletal or cardiac muscle; it rises early after muscle injury but is nonspecific and is not preferred over troponin for myocardial infarction.

\medterm{Myoglobin Urine Test} A urine test detecting myoglobin released during muscle breakdown, used in evaluating rhabdomyolysis; substantial myoglobinuria can contribute to acute kidney injury.

\medterm{Myoglobinuria} A condition in which muscle protein enters urine, often after severe muscle injury or breakdown.

\medterm{Myokymia} Involuntary, fine, continuous rippling contractions of muscle fibers, often seen around the eye but also associated with nerve hyperexcitability or neurologic disease.

\medterm{Myoma} A noncancerous tumor made of muscle tissue, often referring to a fibroid growing in the uterus.

\medterm{Myomalacia} Softening and breakdown of muscle tissue, sometimes affecting heart muscle after serious injury or poor blood supply.

\medterm{Myomectomy} Surgical removal of uterine fibroids while leaving the uterus in place. It may improve bleeding, pain, or pressure symptoms and can preserve the possibility of pregnancy, although risks and recurrence remain.

\medterm{Myometrium} The thick muscle layer forming most of the uterus. It contracts during menstruation and labor, responds to hormones, and may develop conditions such as fibroids, adenomyosis, or abnormal thickening.

\medterm{Myopathic Facies} A characteristic expressionless, mask-like facial appearance marked by drooping eyelids (ptosis), prominent facial muscle wasting, an inability to bury the eyelashes, and a transverse, tented upper lip (tapir mouth). It occurs due to weakness of the facial musculature in conditions such as facioscapulohumeral dystrophy and myotonic dystrophy.

\medterm{Myopathy} A disease that directly affects muscles and causes weakness, fatigue, cramps, or pain. Causes include inherited conditions, inflammation, medicines, hormone disorders, infections, toxins, and problems with muscle energy production.

\medterm{Myopericytoma} A myopericytoma is a usually benign perivascular tumor composed of cells around small blood vessels, often occurring in skin or soft tissue.

\medterm{Myopia} A refractive error in which light focuses in front of the retina, usually because the eyeball is too long or the eye focuses too strongly. Distant objects appear blurred while near objects are usually clearer. It can be corrected with glasses, contact lenses, or selected refractive procedures.

\synonyms
Short-Sightedness; Nearsightedness

\medterm{Myosin} A protein that uses energy to pull on actin, another muscle protein, and produce movement or force. It is important for contraction of skeletal, heart, and smooth muscle and also helps movement inside many other cells.

\medterm{Myosin Atpase} The activity by which the muscle protein myosin breaks down ATP, the cell’s energy source. This releases energy that lets myosin pull on actin and produce muscle contraction.

\medterm{Myositis} Inflammation of muscle tissue caused by autoimmune disease, infection, medicines, or injury. It may produce weakness, pain, swelling, or fatigue; evaluation often includes examination, blood tests, imaging, and sometimes muscle testing.

\medterm{Myositis Ossificans} A noncancerous condition in which abnormal bone forms inside a muscle or soft tissue, most often following a severe blunt injury, deep bruise, or surgery (commonly in the thigh or upper arm). It causes a firm, painful lump and joint stiffness, and usually matures and improves with rest and gentle physical therapy; aggressive massage or premature re-injury should be avoided.
\synonyms Traumatic myositis ossificans; Heterotopic ossification of muscle; Myositis ossifican.

\medterm{Myotonia} Delayed relaxation of a muscle after voluntary contraction or stimulation, often caused by an inherited or acquired muscle-channel disorder.

\medterm{Myotonia Congenita} An inherited skeletal-muscle channelopathy causing delayed relaxation after voluntary
contraction. Stiffness often improves with repeated movement and may be accompanied by
muscle hypertrophy; onset and severity vary between genetic forms.

\medterm{Myotonia Fluctuans} A rare inherited muscle-channel disorder causing stiffness that varies from day to day or after activity. It is usually mild and does not cause lasting weakness; genetic testing may identify a change in the \textit{SCN4A} gene.

\synonyms
Potassium-Aggravated Myotonia Fluctuans

\medterm{Myotonia Permanens} A rare inherited muscle-channel disorder causing continuous, severe muscle stiffness, including at rest. It may begin in infancy and affect the face, throat, or breathing muscles; it is usually related to a change in the \textit{SCN4A} gene and requires specialist care.

\synonyms
Severe Sodium Channel Myotonia; Continuous Myotonia

\medterm{Myotonic Disorder} A neuromuscular disorder characterized by delayed relaxation after muscle contraction, often with weakness, stiffness, or progressive muscle wasting.

\medterm{Myotonic Dystrophy} An inherited disorder causing delayed muscle relaxation, progressive weakness, and wasting. It can also affect the heart, breathing, swallowing, eyes, hormones, or thinking; symptoms and severity vary between types and people.

\medterm{Myotonic Dystrophy Type 2} An inherited multisystem disorder caused by CCTG repeat expansion in CNBP, producing myotonia, proximal or axial weakness, muscle pain, cataracts, and variable endocrine or cardiac features.

\medterm{Myringitis} Inflammation of the eardrum that may cause ear pain, blisters, drainage, or temporary hearing changes.

\medterm{Myringotomy} A small opening made in the eardrum to drain fluid or relieve pressure in the middle ear. A ventilation tube may be inserted to keep the opening open temporarily. Possible problems include infection, scarring, persistent drainage, or a hole that does not close.

\medterm{Myringotomy With Tympanostomy Tube} Incision of the tympanic membrane with placement of a ventilation tube to equalize
middle-ear pressure and permit drainage. It is used selectively for persistent
middle-ear effusion, recurrent acute otitis media, or associated hearing and
developmental problems.

\medterm{Myristica Oil Poisoning} Poisoning from concentrated nutmeg oil, which can cause nausea, vomiting, agitation, fast heartbeat, dizziness, hallucinations, confusion, seizures, or coma. Suspected ingestion needs prompt poison-control or emergency assessment.

\medterm{Myxedema} Severe hypothyroidism causing accumulation of mucin in tissues, with swelling, slowed body functions, and other symptoms.

\medterm{Myxedema Coma} A life-threatening medical emergency representing the decompensated, end-stage state of severe, long-standing, untreated or undertreated hypothyroidism. Typically precipitated by an acute physiological stressor (such as sepsis, cold exposure, myocardial infarction, stroke, or sedatives), it is clinically defined by a classic triad of profound hypothermia (often core temperature $<35^\circ\text{C}$), altered mental status (ranging from lethargy and confusion to stupor and coma), and multisystem organ slowing. Characteristic complications include severe sinus bradycardia, hypotension, hypoventilation leading to hypercapnic respiratory acidosis, dilutional hyponatremia, and diffuse non-pitting periorbital and peripheral myxedema.
\synonyms{Myxoedema Coma; Decompensated Hypothyroidism; Myxedematous Coma}

\medterm{Myxofibroma} A myxofibroma is a benign tumor containing fibrous and gelatinous myxoid tissue, often arising in the jaw or other connective tissue.

\medterm{Myxoid Cyst} A myxoid cyst is a gelatinous fluid-filled lesion near a joint or tendon, commonly on a finger near the distal joint and nail.

\medterm{Myxoma} A usually noncancerous tumor, most often in the upper left chamber of the heart. It may obstruct blood flow or release clots, causing breathlessness, fainting, fever, or stroke-like symptoms; heart ultrasound is commonly used for diagnosis.

\medterm{Myxopapillary Ependymoma} Myxopapillary ependymoma is a usually slow-growing spinal tumor arising from ependymal cells, commonly in the conus or filum terminale.

\medterm{Myxosarcoma} A rare cancer of soft tissue containing abundant jelly-like material. It can grow into nearby tissue or spread to distant organs; diagnosis requires examination of a tissue sample.

\medterm{MZL} Abbreviation; see \textit{Marginal Zone Lymphoma}.

\medterm{Neonatal Mechanical Ventilation} Invasive or non-invasive mechanical positive-pressure ventilation specifically tailored for neonates and infants with respiratory failure, severe respiratory distress syndrome (RDS), bronchopulmonary dysplasia, or apnea of prematurity. Emphasizes lung-protective strategies such as high-frequency oscillatory ventilation and synchronized intermittent mandatory ventilation.

\synonyms
Pediatric Mechanical Ventilation; Neonatal Ventilatory Support; Infant Mechanical Ventilation

\medterm{Night Sweats in Men} Alternate name; see \textit{Night Sweats}.

\medterm{Non-ST-Elevation Myocardial Infarction} A heart attack in which blood tests show heart-muscle injury but the electrocardiogram does not show persistent ST-segment elevation. It is usually caused by reduced or blocked coronary blood flow and may require urgent assessment and treatment.

\medterm{Open Mitral Valve Surgery} Conventional cardiac surgical repair or mechanical/bioprosthetic replacement of the mitral valve performed via a median sternotomy with cardiopulmonary bypass. Indicated for severe mitral regurgitation or mitral stenosis when percutaneous or minimally invasive approaches are not feasible.

\synonyms
Open Mitral Repair; Open Mitral Valve Replacement; Sternotomy Mitral Surgery

\medterm{Oral Cavity} The entry to the digestive tract, containing the teeth, tongue, and openings of salivary glands. It supports chewing, taste, swallowing, speech, and initial digestion.

\synonyms
Oral Cavity; Buccal Cavity

\medterm{Pediatric Midline Catheters} A peripheral vascular access device (typically 3 to 8 cm in length) inserted into a peripheral vein in the upper extremity of infants and children, with the catheter tip terminating distal to the axillary line. Utilized for intermediate-duration intravenous fluid therapy and medication administration lasting 1 to 4 weeks.

\synonyms
Pediatric Midline Venous Access; Infant Midline IV; Neonatal Midline Catheter

\medterm{Pediatric Myocarditis} Acute or chronic inflammation of the myocardium in infants and children, most commonly triggered by viral infections (enteroviruses, adenovirus, parvovirus B19, SARS-CoV-2) or post-infectious autoimmune responses. Manifests as tachypnea, poor feeding, tachycardia, cardiomegaly, arrhythmias, or acute congestive heart failure.

\synonyms
Childhood Myocarditis; Pediatric Myocardial Inflammation; Acute Pediatric Carditis

\medterm{Psychological Health} Alternate name; see \textit{Mental Health}.

\medterm{Rubeola} Alternate name; see \textit{Measles}.

\medterm{Spontaneous Miscarriage} Pregnancy loss occurring without an intentional medical or surgical procedure.

\medterm{ST-Elevation Myocardial Infarction} A heart attack associated with characteristic ST-segment elevation on an electrocardiogram, usually indicating sudden complete blockage of a coronary artery. It causes acute injury to heart muscle and requires rapid emergency evaluation and restoration of blood flow.

\medterm{Stretch Reflex} The automatic tightening of a muscle when it is suddenly stretched, helping maintain muscle tone and posture. Doctors test this reflex by tapping tendons such as the knee, ankle, biceps, or triceps to assess nerve pathways in the spinal cord and brain.

\medterm{Threatened Miscarriage} Vaginal bleeding occurring in early pregnancy prior to 20 weeks of gestation, in the presence of a closed cervical os and documented fetal cardiac activity. Management includes pelvic rest, ultrasound evaluation of fetal viability, maternal Rh immunoglobulin administration if indicated, and monitoring for progression.

\synonyms
Threatened Abortion; Early Pregnancy Bleeding; Impending Miscarriage

\medterm{Toxoplasma Meningoencephalitis} Alternate name; see \textit{Toxoplasmosis}.

\medterm{Tuberculous Meningitis} Severe granulomatous inflammation of the leptomeninges and basal cisterns caused by \textit{Mycobacterium tuberculosis}. Typically presents with a subacute prodrome of fever, headache, altered mental status, and cranial nerve palsies, complicated by obstructive hydrocephalus and cerebral vasculitis.

\synonyms
TB Meningitis; Tuberculous Meningoencephalitis; CNS Tuberculosis

\medterm{Type 2 Diabetes Medications} Alternate name; see \textit{Oral Antidiabetic Agents}.

\medterm{LRINEC Score} A blood-test score that estimates the chance of a serious soft-tissue infection such as necrotizing fasciitis. It uses inflammation, blood count, sodium, kidney, and blood-sugar results. The score cannot by itself confirm or exclude the infection; severe pain, rapid swelling, or serious illness may occur.

\synonyms
Laboratory Risk Indicator for Necrotizing Fasciitis; LRINEC

\medterm{Nabothian Cyst} A small, harmless mucus-filled lump on the cervix caused by a blocked cervical gland. It is common and usually causes no symptoms or need for treatment; an unusual or bleeding lesion still needs examination.

\medterm{Nabumetone} A nonsteroidal anti-inflammatory medicine used for pain and inflammation in selected joint or muscle conditions. It can irritate the stomach, affect kidneys, raise cardiovascular risk, or cause allergic reactions, especially with prolonged use.

\medterm{N-Acetylcysteine} An antidote used mainly for acetaminophen overdose. It helps protect the liver and may still help when treatment starts late or liver injury has already begun.

\medterm{Nadir} The lowest level reached by a measurement or symptom. After chemotherapy, it often means the period when blood-cell counts are lowest and the risk of infection or bleeding is greatest.

\medterm{Naegele Rule} A method for estimating a baby’s due date by adding seven days and about nine months to the first day of the last menstrual period. It is less reliable with irregular cycles or uncertain dates.

\medterm{Naegleria Fowleri} A free-living microscopic amoeba found mainly in warm freshwater lakes, rivers, and hot springs. If water containing it enters the nose, the amoeba can travel to the brain and cause primary amebic meningoencephalitis (PAM), a rare, rapidly worsening brain infection that is nearly always fatal. Early symptoms may include headache, fever, nausea, or vomiting, followed by stiff neck, confusion, seizures, or coma. Infection does not occur from swallowing the water and does not spread from person to person.

\medterm{Naegleria Infection} A rare infection caused when warm freshwater containing \textit{Naegleria fowleri} enters the nose and the amoeba reaches the brain. It causes primary amebic meningoencephalitis (PAM), which rapidly inflames and destroys brain tissue. Symptoms may include severe headache, fever, nausea, vomiting, stiff neck, confusion, seizures, or coma. It is not acquired by swallowing contaminated water and does not spread between people.

\medterm{Naevus} The British spelling of nevus, a usually benign localized skin or mucosal lesion such as a mole.

\medterm{Naevus Flammeus} A congenital vascular capillary malformation (port-wine stain) presenting as a flat, sharply demarcated, pink-to-violaceous cutaneous patch that does not regress spontaneously, classically involving the ophthalmic (V1) distribution of the trigeminal nerve in Sturge-Weber syndrome.
\synonyms
Port-Wine Stain; Capillary Vascular Malformation

\medterm{Naffziger Sign} A clinical physical sign for spinal nerve root compression or spinal space-occupying lesions; manual bilateral compression of the internal jugular veins causes transient elevation of intraspinal venous and CSF pressure, reproducing or worsening radiating radicular pain in the affected dermatome.
\synonyms
Naffziger's Sign; Jugular Compression Radiculopathy Test

\medterm{NAFLD} Abbreviation; see \textit{Non-Alcoholic Fatty Liver Disease}.

\medterm{Nail} A hard protective plate on the end of each finger and toe. Nails protect the tips, help with fine tasks, and may show changes caused by injury, infection, skin disease, or illness elsewhere.

\medterm{Nail Abnormalities} Changes in nail color, shape, thickness, texture, or attachment caused by trauma, infection, inflammatory disease, systemic illness, medication, or congenital conditions.

\medterm{Nail Changes} Alterations in nail color, shape, texture, or attachment that may result from local injury, skin disease, infection,
medicines, nutritional deficiency, or systemic illness. Their significance depends on
the pattern and other clinical findings.

\medterm{Nail Clipping For Fungal Testing} Laboratory examination of clipped nail material by microscopy, culture, or another validated test for fungal infection. Testing can help distinguish nail fungus from injury, psoriasis, or other nail disorders, especially before oral antifungal treatment.

\medterm{Nail Clubbing} Enlargement and rounding of the ends of the fingers or toes with more curved nails. It may run in families or signal disease of the lungs, heart, digestive system, or another organ, especially when it develops recently.

\medterm{Nail-Fold Capillaroscopy} A test that uses magnification to examine the tiny blood vessels near the nails. It can help distinguish primary Raynaud phenomenon from Raynaud phenomenon linked to conditions such as systemic sclerosis, dermatomyositis, or mixed connective-tissue disease.
\synonyms
Periungual Capillaroscopy; NFC Examination

\medterm{Nail Fungus} A fungal infection of a fingernail or toenail causing thickening, discoloration, brittleness, distortion, or separation from the nail bed.

\synonyms
Toenail Fungus

\medterm{Nail Furrows} Nail furrows are grooves or depressions running across or along a nail; transverse Beau lines can follow systemic illness or trauma, while longitudinal ridges are often age-related but may reflect local or systemic disease.

\medterm{Nail Infection} Infection of the nail or surrounding tissue, commonly caused by fungi or bacteria. It may cause thickening,
discoloration, separation of the nail, pain, swelling, or drainage; diabetes and poor circulation
increase the risk of complications.

\medterm{Nail Injuries} Trauma to the nail plate, nail bed, nail folds, or fingertip that may cause bruising, laceration, avulsion, fracture, or later nail-growth abnormalities.

\medterm{Nail Patella Syndrome} Usually an autosomal-dominant disorder caused by a pathogenic variant in \textit{LMX1B}, characterized by abnormal or absent nails, small or absent kneecaps, and elbow or pelvic abnormalities such as iliac horns. Glaucoma and kidney disease with proteinuria may also occur.

\medterm{Nail-Patella Syndrome} A rare inherited condition affecting the nails, kneecaps, elbows, and pelvis. Some people also develop kidney disease or glaucoma. The kneecaps may be small, absent, or displaced, and lifelong monitoring may be needed.

\medterm{Nail Pits} Small dents in the surface of a nail. They are common in psoriasis but may also occur with hair-loss disorders, eczema, lichen planus, or other conditions affecting nail growth.

\medterm{Nail Pitting} Punctate, cupuliform depressions in the dorsal nail plate resulting from parakeratotic defective keratinization within the proximal nail matrix, characteristically seen in psoriasis, alopecia areata, and reactive arthritis.

\medterm{Nail Psoriasis} Psoriasis affecting the nails can cause small pits, separation from the nail bed, yellow or red-brown patches, thickening, crumbling, or ridges. It may occur with skin psoriasis or psoriatic arthritis and can be severe or newly developed.

\medterm{Nails} Hard plates of keratin covering the ends of fingers and toes. Changes in color, shape, thickness, or attachment can result from injury, infection, skin disease, nutritional problems, or other illness.

\medterm{Nalbuphine} An opioid pain medicine with mixed effects, used for moderate to severe pain under medical supervision.

\medterm{Naloxone} A medicine that blocks opioids and can quickly restore breathing during an opioid overdose. It comes as a nasal spray or injection. Emergency help is still needed because its effect may wear off before the opioid does.

\medterm{Naltrexone} A longer-acting medicine that blocks opioid effects and reduces alcohol-related cravings in some people. Opioid use during treatment can trigger withdrawal.

\medterm{Nanocolloid} Very small particles carrying a radioactive tracer. After injection, they can show how lymph drains and help locate the first lymph node receiving drainage from a tumor.

\medterm{Nanomedicine} The use of very small materials or devices in medicine to help diagnose, prevent, monitor, or treat disease. Examples include nanoparticle drug carriers, imaging agents, biosensors, and some vaccine or cancer treatments; their effects and safety depend on the material, size, dose, and how the body handles them.

\medterm{Napelline} A poisonous plant chemical related to aconite alkaloids. Exposure can disturb nerve signals and heart rhythm, causing tingling, weakness, vomiting, dangerous blood-pressure changes, or life-threatening cardiac poisoning.

\medterm{Naphazoline} A topical alpha-adrenergic agonist used as a nasal decongestant or ocular vasoconstrictor; prolonged use can worsen congestion.

\medterm{Naphthalene Poisoning} Harm from mothballs or other products containing naphthalene. It can destroy red blood cells, causing weakness, jaundice, dark urine, kidney injury, seizures, or coma.

\medterm{Naproxen} A nonsteroidal anti-inflammatory drug used for pain, fever, and inflammation.

\medterm{Naproxen Sodium Overdose} Taking too much naproxen can cause stomach pain, vomiting, bleeding, drowsiness, kidney injury, seizures, or breathing problems. Large or intentional overdoses can be life-threatening.

\medterm{Narcissistic Personality Disorder} A long-lasting pattern of an unusually strong need for admiration, a sense of special importance, and difficulty recognizing other people's feelings. It causes problems in relationships or daily life and is diagnosed by a trained mental-health professional after a full assessment.

\synonyms
NPD; Megalomania; Pathological Narcissism

\medterm{Narcolepsy} A chronic disorder that affects the brain’s control of sleep and wakefulness. It causes overwhelming daytime sleepiness and sudden, irresistible sleep episodes, even after enough nighttime sleep; nighttime sleep may also be broken. The boundary between sleep and wakefulness is less clear, so some features of rapid-eye-movement (REM) sleep may occur while awake or when falling asleep or waking. Symptoms may include cataplexy, a brief loss of muscle control triggered by strong emotions; sleep paralysis; and vivid dream-like experiences. Type 1 narcolepsy includes cataplexy and is often associated with low hypocretin (orexin), a brain chemical that helps maintain wakefulness; type 2 usually does not include cataplexy.

\medterm{NARP Syndrome} A rare inherited mitochondrial disorder that can cause nerve damage, muscle weakness, loss of coordination, seizures, and vision problems. Symptoms and severity vary widely.
\synonyms
NARP; Neuropathy, Ataxia, and Retinitis Pigmentosa

\medterm{Narrow Angle Glaucoma} A condition in which the drainage area inside the eye is unusually narrow, raising the risk of a sudden rise in eye pressure. Severe eye pain, blurred vision, halos, nausea, or vomiting may occur during an acute attack.

\medterm{Narrowband UVB Phototherapy} A treatment that uses carefully measured ultraviolet light for conditions such as psoriasis, vitiligo, and eczema. Treatment requires eye protection and monitoring of the skin.

\medterm{Narrow-Complex Tachycardia} A fast heart rhythm that usually begins in the upper chambers of the heart. It produces narrow-looking beats on an ECG and may cause palpitations, dizziness, chest discomfort, breathlessness, or fainting.

\medterm{Narrow Pulse Pressure} A pulse pressure (systolic minus diastolic blood pressure) of less than 30 mmHg or less than twenty-five percent of systolic pressure. It reflects reduced stroke volume or marked peripheral vasoconstriction, commonly seen in severe aortic stenosis, cardiac tamponade, hypovolaemic shock, and advanced heart failure.

\medterm{Narrow Therapeutic Index} A medicine for which the helpful dose is close to the harmful dose. It requires careful dosing and sometimes blood tests to avoid toxicity.

\medterm{Nasal} Relating to the nose or the space inside it. The nose warms, moistens, and filters incoming air and helps with smell and speech. Nasal disease may cause blockage, bleeding, discharge, reduced smell, or difficulty breathing.

\medterm{Nasal Bleeding} Alternate name; see \textit{Epistaxis}.

\medterm{Nasal Cannula} A soft tube with two small prongs that rest inside the nostrils to deliver extra oxygen. It is used when a person needs mild to moderate breathing support, and the oxygen flow is adjusted while breathing is monitored.

\medterm{Nasal Cavity} The paired space inside the nose, extending from the nostrils to the nasopharynx. It warms, humidifies, and filters inspired
air and contains the receptors responsible for smell.

\medterm{Nasal Congestion} A blocked or stuffy feeling in the nose caused by swelling of the nasal lining, extra mucus, or a structural problem. It occurs with colds, allergies, sinus disease, irritants, and other conditions.

\medterm{Nasal Congestion And Rhinorrhea} Nasal obstruction or stuffiness together with a runny nose, commonly caused by viral infection, allergic or nonallergic rhinitis, sinus disease, irritants, or structural blockage. Persistent unilateral symptoms, recurrent bleeding, fever, facial swelling, or breathing difficulty requires evaluation.

\medterm{Nasal Corticosteroid Sprays} Nasal corticosteroid sprays reduce mucosal inflammation and congestion in allergic rhinitis, nasal polyps, and some chronic sinus disorders; correct direction away from the septum limits irritation and nosebleeds.

\medterm{Nasal Disorders} Conditions affecting the nose or nasal passages, including infection, allergy, trauma, obstruction, polyps, and tumors.

\medterm{Nasal Endoscopy} An examination in which a thin camera is used to look inside the nose and upper throat. It can help find swelling, polyps, a foreign object, bleeding, or a growth.

\medterm{Nasal Flaring} Widening and outward movement of the nostrils during breathing. It is a visible sign of
increased work of breathing, particularly in infants and children, and may accompany
airway obstruction, bronchiolitis, pneumonia, asthma, or severe systemic illness.

\medterm{Nasal Implant} A medical device placed inside or near the nose to support weakened tissue, change nasal shape, or improve airflow. The implant may be temporary or permanent, and risks include infection, movement, scarring, or breathing problems.

\medterm{Nasal Mucosal Biopsy} Removal of a small piece of tissue from the nose for laboratory examination. It may help investigate persistent inflammation, infection, unusual deposits, or a growth; soreness, bleeding, and infection are possible complications.

\medterm{Nasal Obstruction} Blocked airflow through one or both nasal passages. Causes include congestion, allergy, infection, a
deviated septum, polyps, trauma, or a mass. Persistent one-sided obstruction or bleeding requires
medical evaluation.

\medterm{Nasal Passage} One of the air channels inside the nose. Its lining warms, moistens, and filters inhaled air before it reaches the throat; swelling, mucus, polyps, or a structural blockage can make breathing difficult.

\medterm{Nasal Polyps} Soft, usually noncancerous growths in the lining of the nose or sinuses. They can cause a blocked nose, runny nose, and reduced sense of smell.

\synonyms
Sinonasal Polyps; Nasal Polyposis

\medterm{Nasal Septal Hematoma} A collection of blood inside the wall dividing the nostrils, usually after a nose injury. It can block breathing and damage the nose if it is not drained promptly.

\medterm{Nasal Septum} The wall that divides the nasal cavity into right and left sides. It consists mainly of cartilage in front and bone behind
and helps support the nose and direct airflow.

\medterm{Nasal Septum Perforation} A hole through the cartilage or bone separating the nasal passages, which may cause crusting, nosebleeds, whistling, or obstruction and can result from trauma, surgery, intranasal drugs, or inflammatory disease.

\medterm{Nasal Spray} A liquid or mist delivered into the nose. Different sprays contain saline, corticosteroids, antihistamines, or decongestants; correct use and duration depend on the product.

\medterm{Nasal Stent} A small device placed inside the nose to keep an airway open or support healing after surgery. It may cause crusting, blockage, infection, bleeding, or pressure injury.

\medterm{Nasal Vestibulitis} Inflammation or infection of the nasal vestibule, commonly caused by Staphylococcus bacteria and producing tenderness, crusting, redness, or small pustules near the nostrils.

\medterm{Nasogastric Feeding Tube} A flexible tube passed through the nose into the stomach to deliver enteral nutrition or medicines, or to decompress or aspirate the stomach.

\medterm{Nasogastric Tube} A tube passed through the nose into the stomach to remove contents or deliver nutrition, fluids, or medicines when clinically indicated.

\medterm{Nasojejunal Feeding Tube} An enteral feeding catheter passed through the nose, pharynx, stomach, and pylorus into the proximal jejunum under endoscopic or fluoroscopic guidance, indicated for patients requiring enteral nutrition who cannot tolerate gastric feeding due to severe gastroparesis, recurrent aspiration, or acute pancreatitis.

\synonyms
NJ Tube; Post-Pyloric Feeding Tube

\medterm{Nasolacrimal Duct} A membranous canal enclosed within a bony canal that conveys tears from the lacrimal sac downward into the inferior meatus of the nasal cavity beneath the inferior nasal concha, guarded by the mucosal valve of Hasner.

\synonyms
Tear Duct

\medterm{Nasopharyngeal Airway} A soft tube inserted through the nose to maintain a passage between the nares and
pharynx. It may be useful in a spontaneously breathing patient with partial upper-airway
obstruction but is relatively contraindicated in suspected basal skull fracture or
significant nasal trauma.

\medterm{Nasopharyngeal Carcinoma} A cancer arising in the upper throat behind the nose. It may cause a blocked or bleeding nose, hearing problems, a neck lump, headache, or swallowing trouble; risk varies with genetics, infections, and geographic background.

\synonyms
NPC; Nasopharyngeal Cancer

\medterm{Nasopharyngeal Culture} A laboratory test in which a sample from behind the nose is placed in a growth medium to look for selected bacteria or fungi. A negative result does not exclude every infection, and other tests may be more sensitive.

\medterm{Nasopharynx} The upper part of the throat behind the nose and above the soft palate. It carries air, contains the openings of the tubes that connect to the middle ears, and includes nearby lymph tissue that helps defend against infection.

\medterm{Natal Cleft} The groove between the buttocks, also called the intergluteal or sacral cleft.

\medterm{Natal Teeth} Teeth present at birth, most often involving the lower central incisors. They may be
loose, interfere with feeding, or injure the infant's tongue and require dental
assessment.

\medterm{Natriuresis} Removal of sodium from the body in urine. Water usually leaves with the sodium, which can reduce fluid volume and blood pressure; the process may be increased by natural hormones or water-removing medicines.

\medterm{Natriuretic} Causing the kidneys to remove more sodium in urine. Water usually leaves with the sodium, which can lower blood volume and pressure. The term describes certain body hormones and some medicines; related blood tests can help assess heart failure.

\medterm{Natural Immunity} Protection that develops after a person’s immune system responds to an infection, or temporary protection passed to a newborn through the placenta or breast milk. It may be incomplete or fade and is not always safer than vaccination.

\medterm{Natural Killer Cell} A type of white blood cell that can recognize and destroy some virus-infected or abnormal cells without previous exposure to that specific target. It helps provide early immune defense and releases signals that guide other immune cells.

\synonyms
NK Cells; Large Granular Lymphocytes

\medterm{Nausea} An unpleasant feeling that you may vomit. It can occur with stomach or intestinal illness, pregnancy, medicines, migraine, inner-ear problems, metabolic disorders, or serious abdominal or neurologic disease. Persistent vomiting, severe pain, dehydration, blood, or confusion may signal a serious condition.

\medterm{Nausea And Vomiting} Nausea is the urge to vomit; vomiting is forceful emptying of stomach contents. Causes include infection, medicines, pregnancy, migraine, obstruction, and other illnesses.

\medterm{Nausea And Vomiting – Adults} Nausea and vomiting in adults can result from infection, stomach or bowel disease, medicines, pregnancy, migraine, metabolic illness, neurologic disease, or blockage. Severe pain, blood, confusion, dehydration, chest symptoms, or persistent symptoms may indicate a serious condition.

\medterm{Nausea And Vomiting During Early Pregnancy} Nausea and vomiting are common in early pregnancy, but persistent symptoms with inability to keep fluids down, weight loss, dehydration, abdominal pain, fever, blood, or altered mental status may indicate hyperemesis gravidarum or another condition and require medical evaluation.

\medterm{Navel} The navel, or umbilicus, is the scar left after the umbilical cord separates and marks the former connection between fetus and placenta.

\medterm{Navicular} A boat-shaped bone in the middle of the foot that helps support the arch. Injury to it can cause foot pain and may heal slowly because part of the bone has a limited blood supply.

\medterm{NCS} Usually an abbreviation for nerve-conduction study, a test that measures how quickly and strongly electrical signals travel through nerves. It can help assess nerve damage, pressure on a nerve, or some causes of numbness and weakness.

\medterm{NDM-1} An enzyme made by some bacteria that makes them resistant to many important antibiotics. Infections caused by these bacteria can be difficult to treat and require laboratory testing and careful infection control.

\medterm{Near-Drowning} Survival for at least 24 hours following suffocation or submersion in a liquid medium, characterized pathologically by alveolar surfactant washout, pulmonary edema, ventilation-perfusion mismatch, acute respiratory distress syndrome, and hypoxic-ischemic encephalopathy.

\synonyms
Non-Fatal Drowning; Submersion Injury

\medterm{Nearsightedness} A common focusing problem in which nearby objects are clear but distant objects look blurry because light focuses in front of the retina. Glasses, contact lenses, or selected procedures can improve vision.

\synonyms
Myopia; Short-Sightedness

\medterm{Nebulizer} A device that converts liquid medicine into an aerosol for inhalation through a mouthpiece or mask. Common types include jet nebulizers, ultrasonic nebulizers, and vibrating-mesh nebulizers; they differ in how the aerosol is generated and delivered.

\synonyms
Nebuliser; Atomizer

\medterm{NEC} Abbreviation; see \textit{Neuroendocrine Carcinoma}.

\medterm{Necator} A group of hookworms that live in the small intestine. They enter through the skin, mature in the gut, and feed on blood, potentially causing abdominal symptoms, iron-deficiency anemia, weakness, and poor growth.

\medterm{Necator Americanus} A hookworm species that infects humans through skin penetration by larvae and causes intestinal blood loss.

\medterm{Neck} The neck connects the head to the trunk and contains the cervical spine, airway, esophagus, thyroid, vessels, nerves, muscles, and lymph nodes.

\medterm{Neck And Back Pain} Pain in the neck, upper back, lower back, or surrounding tissues, commonly caused by strain, age-related changes, disc problems, or trauma. Weakness, numbness, walking difficulty, bowel or bladder changes, fever, cancer history, or major trauma may indicate a serious cause.

\medterm{Neck Cancer} Cancer arising in tissues of the neck, such as the throat, thyroid, salivary glands, lymph nodes, or other structures. Symptoms may include a persistent lump, pain, swallowing trouble, voice change, or unexplained weight loss.

\medterm{Neck Dissection} Surgical removal of cervical lymph nodes and sometimes surrounding tissues to treat or stage head and neck cancer.

\medterm{Neck Lump} A localized swelling or mass in the neck caused by enlarged lymph nodes, thyroid or salivary disease, cyst, infection, or tumor; persistent or concerning masses require evaluation.

\medterm{Neck Mass} A lump or swelling in the neck caused by an enlarged lymph node, congenital cyst, thyroid or salivary-gland lesion, infection, or tumor. A persistent, enlarging, firm, fixed, or unexplained mass—especially in an adult—can have an important underlying cause.

\medterm{Neck Pain} Pain or stiffness in the neck, commonly caused by muscle strain, poor posture, injury, arthritis, degenerative disc disease, or nerve compression. It may be localized or spread to the shoulders, arms, or head.

\synonyms
Cervicalgia; Cervical Pain; Neck Stiffness

\medterm{Neck Pain And Dizziness} Neck pain with dizziness may reflect muscle strain, migraine, inner-ear disease, medication, blood-pressure change, or less commonly vascular or neurologic disease. Sudden severe symptoms, weakness, speech or vision change, fainting, or inability to walk may occur with a serious condition.

\medterm{Neck Sprain} Alternate name; see \textit{Whiplash}.

\medterm{Neck Strain} Alternate name; see \textit{Whiplash}.

\medterm{Neck X-Ray} Plain radiographic imaging of the cervical spine or neck soft tissues used to evaluate alignment, fractures, degenerative change, foreign bodies, or selected airway abnormalities.

\medterm{Necrobiosis Lipoidica} A chronic inflammatory skin disorder characterized by yellow-brown, atrophic plaques,
usually on the shins. It is associated with diabetes but can occur without diabetes.

\medterm{Necrobiosis Lipoidica Diabeticorum} A chronic inflammatory skin disorder producing yellow-brown atrophic plaques with a raised border, most often on the shins and associated with diabetes, although it can occur without diabetes.

\medterm{Necrolytic Acral Erythema} A rare rash with dark, thickened, or blistered patches, usually on the feet and ankles. It is often linked to long-term hepatitis C but can also occur with low zinc or other metabolic problems. Testing is needed to find the cause.
\synonyms{NAE; Hepatitis C-Associated Necrolytic Erythema}

\medterm{Necrolytic Migratory Erythema} A recurrent, erosive, erythematous eruption that often affects periorificial and
intertriginous skin. It is classically associated with glucagonoma but may have other causes,
so evaluation should consider nutritional, hepatic, pancreatic, and metabolic disease.

\medterm{Necrosis} Necrosis is death of cells or tissue within a living body caused by ischemia, infection, toxins, trauma, inflammation, or other injury.

\medterm{Necrotizing Anterior Scleritis} A severe inflammation of the white outer coat of the eye that causes tissue damage and can threaten vision. It usually causes intense eye pain, redness, and light sensitivity.

\medterm{Necrotizing Enterocolitis} A serious illness in which inflammation damages the intestine, mainly in premature or medically fragile newborns. A swollen abdomen, feeding problems, green vomit, bloody stools, or unusual sleepiness may occur.

\medterm{Necrotizing Fasciitis} A life-threatening infection that rapidly destroys tissue beneath the skin. Early skin changes may be mild despite severe pain and deeper damage. Treatment usually includes surgery and antibiotics.

\medterm{Necrotizing Granuloma} A granulomatous inflammatory lesion containing central tissue necrosis, commonly associated with infections such as tuberculosis or fungal disease but requiring clinical and histologic correlation.

\medterm{Necrotizing Infections} Severe infections that cause tissue death through germs, toxins, blocked blood flow, or intense inflammation. Examples include necrotizing fasciitis, necrotizing pneumonia, and necrotizing enterocolitis. Rapidly worsening pain, swelling, skin changes, severe illness, or organ problems may occur.

\medterm{Necrotizing Pneumonia} Severe pneumonia in which infection causes necrosis and cavitation of lung tissue, often accompanied by sepsis, respiratory failure, or abscess formation.

\medterm{Necrotizing Sialometaplasia} A benign, self-limiting, ischemic inflammatory lesion of the salivary glands (most commonly involving the minor salivary glands of the hard palate) that clinically and histologically mimics mucoepidermoid carcinoma or squamous cell carcinoma.
\synonyms
Ischemic Sialadenitis; Benign Palatal Sialometaplasia

\medterm{Necrotizing Soft Tissue Infection} A rapidly progressive infection of the tissue beneath the skin or in muscle that causes severe pain, swelling, severe illness, and tissue death. Treatment may include surgery and antibiotics.

\medterm{Necrotizing Vasculitis} Inflammation of blood-vessel walls causing fibrinoid necrosis and potential vessel narrowing, thrombosis, hemorrhage, ischemia, or organ injury.

\medterm{Needle} A hollow, sharp tube used to inject medicine or remove blood, fluid, or tissue. Needles come in different lengths and widths; the appropriate size depends on the procedure and body site.

\medterm{Needle Biopsy} Removal of cells, fluid, or a tissue sample through a needle for cytologic or histopathologic examination. Fine-needle aspiration primarily collects cells, whereas core biopsy collects tissue architecture.

\medterm{Needle Decompression} An emergency procedure used to release trapped air from around a collapsed lung when a tension pneumothorax is causing severe breathing or circulation problems. A trained clinician places a needle or catheter through the chest, then usually inserts a chest tube for ongoing treatment. It is a temporary lifesaving measure, not a complete repair.

\synonyms
Needle Thoracostomy; Tension Pneumothorax Decompression; Emergency Chest Decompression

\medterm{Needle Driver} A surgical instrument with gripping jaws and handles used to hold and guide a suture needle through tissue. Its design helps control the needle while protecting the operator's fingers.

\synonyms
Needle Holder
Suture Needle Holder

\medterm{Needle Holder} Alternate name; see \textit{Surgical Needle Holder}.

\medterm{Needlestick Injury} A puncture from a needle or sharp object that may expose someone to blood or body fluids. It can create a risk of infection, including HIV or hepatitis, and may lead to testing or preventive treatment.

\medterm{Neer Sign} A physical examination finding used to evaluate for subacromial impingement syndrome. The examiner stabilizes the patient's scapula while passively elevating the internally rotated arm in full forward flexion; reproduction of anterior shoulder pain indicates impingement of the rotator cuff tendons against the coracoacromial arch.

\synonyms
Neer Impingement Test

\medterm{Neer Test} A clinical diagnostic physical examination maneuver for subacromial shoulder impingement; the examiner stabilizes the patient's scapula to prevent upward rotation and passively performs forced forward elevation and internal rotation of the fully extended arm, driving the greater tuberosity against the anteroinferior acromion to reproduce anterior shoulder pain.
\synonyms
Neer Impingement Sign; Neer Shoulder Test

\medterm{Negative Feedback} A control process in which a change triggers a response that reduces or reverses the original change. It helps keep body functions, such as temperature, blood sugar, and hormone levels, within a useful range.

\medterm{Negative Predictive Value} The chance that a person with a negative test result truly does not have the condition. It depends on test accuracy and how common or likely the condition was before testing.

\medterm{Negative Pressure Pulmonary Edema} A form of non-cardiogenic pulmonary edema provoked by vigorous, forceful inspiratory efforts against an acutely obstructed upper airway (such as post-extubation laryngospasm), generating extreme negative intrathoracic pressures that increase venous return and cause alveolar capillary stress failure.

\synonyms
NPPE; Post-Obstructive Pulmonary Edema

\medterm{Negative Pressure Ventilation} Breathing support that uses pressure around the chest to help the lungs expand, instead of pushing air directly through the airway. It is used rarely, including for some neuromuscular conditions.

\medterm{Negative-Strand RNA Virus} A virus whose genetic material cannot be read directly to make proteins. After entering a cell, it must first make a readable copy of its RNA before producing new virus particles.

\medterm{Negative Symptoms} A reduction or loss of usual emotional expression, speech, motivation, pleasure, or interest in other people. They can occur in schizophrenia and some other conditions and may interfere with work, relationships, and daily activities.

\medterm{Neglect} Failure to provide a child with necessary food, shelter, supervision, education, emotional care, or medical care, causing harm or risk of harm. The term can also describe ignoring one side of space after brain injury; the intended meaning depends on context.

\medterm{Neglected Tropical Disease} An infectious or parasitic disease that mainly affects populations in tropical or subtropical regions with limited resources. Examples include schistosomiasis, leprosy, dengue, and lymphatic filariasis.

\medterm{Neglect Syndrome} A neurologic problem in which a person fails to notice or respond to one side of the body or surrounding space, usually after injury to the opposite side of the brain. It is more than ordinary inattention.

\medterm{Neisseria Gonorrhoeae} A gram-negative diplococcus that causes gonorrhea, a sexually transmitted infection. The
organism has developed resistance to many antibiotics including fluoroquinolones and
macrolides.

\medterm{Neisseria Infection} Infection caused by a pathogenic Neisseria species. The most important human infections are gonorrhoea caused by N. gonorrhoeae and meningococcal disease caused by N. meningitidis; illness may affect the genital tract, throat, rectum, joints, bloodstream, or meninges, and treatment depends on the site and antibiotic susceptibility.

\synonyms
Neisseria Infections

\medterm{Neisseria Meningitidis} A gram-negative diplococcus that causes meningitis and meningococcemia. The organism has
a polysaccharide capsule that is the target of conjugate vaccines. Serogroups A, B, C,
W, X, and Y cause most disease. Meningitis presents with fever, headache, neck
stiffness, and altered mental status.

\medterm{Nelson Syndrome} Enlargement and hypersecretion of an adrenocorticotropic-hormone-producing pituitary
tumour after bilateral adrenalectomy, usually for Cushing disease. It may cause rapid
skin hyperpigmentation, headache, visual-field loss, and other effects of an expanding
sellar mass.

\medterm{Nemaline Myopathy} A group of inherited muscle diseases that cause weakness because muscle fibers do not develop or work normally. Severe forms begin in infancy and may affect movement, feeding, or breathing; milder forms begin later. The severity varies widely, and care may include breathing support, nutrition, physical therapy, and genetic counseling.

\synonyms
Rod Body Myopathy; Nemaline Rod Myopathy

\medterm{Neoadjuvant Therapy} Treatment given before the main treatment, such as surgery or radiotherapy, to shrink a tumor or make it easier to treat. It may also show how the tumor responds and help doctors plan the next treatment.

\medterm{Neocortex} The outer part of the cerebral hemispheres involved in conscious perception, language, planning, memory, and voluntary movement.

\medterm{Neoehrlichiosis} An infection caused by the bacterium Candidatus Neoehrlichia mikurensis, usually transmitted by ticks. It can cause fever and blood-vessel inflammation, especially in people with weakened immunity.

\medterm{Neologisms} Newly created words, phrases, or idiosyncratic combinations of word fragments that have meaning only to the speaker and are incomprehensible to others. They are characteristic features of formal thought disorder in schizophrenia and fluent (Wernicke) aphasia.

\medterm{Neonatal} Relating to a newborn, especially during the first 28 days after birth. This period includes important changes in breathing, feeding, temperature control, immunity, circulation, and adjustment to life outside the womb.

\medterm{Neonatal Abstinence Syndrome} Withdrawal symptoms in a newborn after exposure to opioids or certain other medicines or substances during pregnancy. Symptoms may include trembling, irritability, poor feeding, vomiting, diarrhea, sweating, or breathing problems.

\medterm{Neonatal Cephalohematoma} A collection of blood beneath the outer covering of a newborn’s skull, usually caused by birth trauma. It is limited by the skull’s seams and generally resolves over time, but it can be associated with injury or jaundice.

\medterm{Neonatal Conjunctivitis} Inflammation of the eye’s outer lining during the first month of life, caused by chemical irritation, bacteria, or viruses acquired during delivery. Pus-like discharge, eyelid swelling, and corneal involvement may occur.

\medterm{Neonatal Cystic Fibrosis Screening Test} A newborn screening test, usually measuring immunoreactive trypsinogen with reflex DNA or sweat testing, used to identify infants at increased risk for cystic fibrosis.

\medterm{Neonatal Diabetes Mellitus} Diabetes first diagnosed during the first six months of life. It may be temporary or lifelong and is usually caused by a genetic change affecting insulin production; specialist care is needed.

\medterm{Neonatal Galactorrhea} A temporary milk-like breast secretion in a newborn caused by exposure to maternal hormones. It usually resolves without treatment, and pressure on the breasts can cause irritation or infection.

\synonyms
Neonatal Milk Secretion; Newborn Galactorrhea

\medterm{Neonatal Herpes Simplex} Herpes simplex infection in a baby during the first 28 days of life, usually acquired around delivery. It may affect the skin, eyes, mouth, brain, or several organs and can resemble sepsis. Treatment uses antiviral medicine.

\medterm{Neonatal Herpes Simplex Virus} A severe neonatal infection caused by herpes simplex virus, typically acquired during
passage through the birth canal, presenting within the first month of life as skin, eye,
and mouth disease, central nervous system disease, or disseminated disease with
multiorgan involvement.

\medterm{Neonatal Hyperbilirubinemia} An excess of bilirubin in a newborn’s blood, causing yellow skin or eyes. Mild cases are common, but rapidly rising or very high levels can damage the brain; treatment may include frequent feeding, phototherapy, or rarely an exchange transfusion.

\medterm{Neonatal Hyperglycemia} An abnormally elevated blood glucose level in infants and neonates (typically defined as blood glucose $>145\text{ mg/dL}$ or $>8.0\text{ mmol/L}$). Common causes include extreme prematurity, intravenous dextrose overinfusion, sepsis, severe physiological stress, corticosteroid therapy, or transient neonatal diabetes mellitus.

\synonyms
Infantile Hyperglycemia; High Blood Glucose in Infants; Newborn Hyperglycemia

\medterm{Neonatal Hypertension} Sustained blood pressure elevation in infants and neonates exceeding age- and birthweight-adjusted norms. Frequently secondary to umbilical artery catheter-related renal artery thrombosis, congenital renal anomalies, bronchopulmonary dysplasia, or congenital heart defects.

\synonyms
Infant Hypertension; Neonatal High Blood Pressure; Infantile Hypertension

\medterm{Neonatal Hypocalcemia} Low serum calcium in a newborn, which may cause jitteriness, poor feeding, hypotonia, apnea, or seizures. Causes include prematurity, maternal diabetes, birth asphyxia, and disorders of parathyroid function.

\medterm{Neonatal Hypoglycemia} Blood sugar that is too low in a newborn. It may cause trembling, sleepiness, poor feeding, breathing pauses, or seizures and is more likely with prematurity, unusual birth size, or maternal diabetes; prompt monitoring and treatment are needed.

\medterm{Neonatal Hypothermia} A newborn’s body temperature below 36.5°C. It can result from exposure, prematurity, infection, or illness and may worsen breathing, blood sugar, and circulation problems. Prompt warming and assessment are needed; skin-to-skin contact and a warm, dry environment help prevent it.

\medterm{Neonatal Hypothyroidism} Insufficient thyroid hormone production in a newborn, which can impair brain development and growth if untreated; newborn screening enables early diagnosis and treatment.

\medterm{Neonatal Hypoxic-Ischemic Encephalopathy} Brain injury in a newborn caused by too little oxygen or blood flow around the time of birth. It may cause poor alertness, abnormal muscle tone, seizures, breathing difficulty, or injury to other organs. Selected babies benefit from carefully controlled cooling started soon after birth.

\medterm{Neonatal Infections} Infections occurring during the first weeks of a newborn’s life. They may begin before or soon after birth from exposure during pregnancy or delivery, or later from the environment or medical devices. Warning signs include poor feeding, temperature change, breathing difficulty, unusual sleepiness, irritability, or color change.

\medterm{Neonatal Jaundice} Yellow discoloration of a newborn's skin or eyes caused by elevated bilirubin. It is common but may
require treatment or investigation when severe, early, or prolonged.

\synonyms
Neonatal Hyperbilirubinemia; Icterus Neonatorum

\medterm{Neonatal Lupus} A condition caused by antibodies passed from a pregnant person to a baby. It may cause a temporary rash, low blood counts, liver problems, or a slow heart rhythm. Most effects improve as the antibodies leave the baby's blood, but heart block can be serious and may require a pacemaker.

\synonyms
Neonatal Lupus Erythematosus; NLE

\medterm{Neonatal Mechanical Ventilation} Invasive or non-invasive mechanical positive-pressure ventilation specifically tailored for neonates and infants with respiratory failure, severe respiratory distress syndrome (RDS), bronchopulmonary dysplasia, or apnea of prematurity. Emphasizes lung-protective strategies such as high-frequency oscillatory ventilation and synchronized intermittent mandatory ventilation.

\synonyms
Pediatric Mechanical Ventilation; Neonatal Ventilatory Support; Infant Mechanical Ventilation

\medterm{Neonatal Neutropenia} An absolute neutrophil count (ANC) below $1{,}000/\mu\text{L}$ to $1{,}500/\mu\text{L}$ in infants, or below $500/\mu\text{L}$ in severe cases. Etiologies include neonatal sepsis, maternal preeclampsia, alloimmune or autoimmune neutropenia, twin-to-twin transfusion, bone marrow failure syndromes (such as Kostmann syndrome), or drug exposure.

\synonyms
Infant Neutropenia; Neonatal Low Neutrophil Count; Pediatric Neutropenia

\medterm{Neonatal Opioid Withdrawal Syndrome} Withdrawal symptoms in a newborn exposed to opioid medicines or drugs during pregnancy. The baby may be hard to settle, tremble, feed poorly, vomit, have loose stools, sweat, or breathe rapidly. Quiet surroundings, swaddling, skin-to-skin contact, and feeding support are used first; some babies need hospital medicines and monitoring.

\synonyms
NOWS; Opioid-Induced NAS

\medterm{Neonatal Peripheral Arterial Line} A fine indwelling cannula placed in a peripheral artery (such as the radial, posterior tibial, or dorsalis pedis artery) of an infant for continuous invasive hemodynamic arterial blood pressure monitoring and frequent arterial blood gas sampling in neonatal intensive care.

\synonyms
Infant Arterial Line; Pediatric A-Line; Neonatal Arterial Cannulation

\medterm{Neonatal Peripheral IV Line} A short peripheral venous catheter inserted into a superficial vein in the hand, foot, forearm, or scalp of an infant for short-term hydration, medication delivery, or blood transfusions.

\synonyms
Infant PIV; Pediatric Peripheral Cannula; Neonatal IV Cannulation

\medterm{Neonatal PICC Line} A small-caliber (1.0--2.0 Fr) central venous catheter inserted percutaneously into a peripheral vein of an infant or neonate with the tip positioned in the superior or inferior vena cava, providing reliable vascular access for total parenteral nutrition and hyperosmolar medications.

\synonyms
Neonatal Percutaneous Central Line; Infant PICC; Pediatric Percutaneous Central Catheter

\medterm{Neonatal Pneumothorax} Pathologic accumulation of extrapulmonary air within the pleural space in a newborn or infant, which may occur spontaneously or secondary to meconium aspiration, respiratory distress syndrome, or mechanical barotrauma.

\synonyms
Infant Pneumothorax; Neonatal Air Leak Syndrome; Newborn Pleural Air Collection

\medterm{Neonatal Polycythemia} An abnormally high red-cell concentration in a newborn, often defined by a venous hematocrit of at least 65 percent. It may increase blood viscosity and cause respiratory, neurologic, or metabolic symptoms.

\medterm{Neonatal Respiratory Distress Syndrome} Breathing difficulty in a premature newborn caused mainly by insufficient pulmonary surfactant, leading to stiff lungs and impaired oxygen exchange.

\medterm{Neonatal Resuscitation} Emergency steps used to help a newborn who is not breathing well or has a slow heart rate after birth. They may include warmth, airway support, assisted breathing, and chest compressions.

\medterm{Neonatal Sepsis} A serious infection causing illness throughout a newborn's body during the first 28 days of life. It may begin before or after birth and can cause poor feeding, temperature change, breathing difficulty, sleepiness, or color change.

\medterm{Neonatal Skin Conditions} Skin changes common in newborns, including temporary rashes, tiny white spots, heat rash, and small yellow bumps. Most are harmless, but blistering, widespread redness, fever, poor feeding, or a rapidly worsening rash may signal illness.

\medterm{Neonatal Testicular Torsion} Twisting of the spermatic cord around a testicle during the newborn period, reducing or stopping its blood supply. It may appear as a firm, enlarged, discolored, or painless scrotal swelling.

\medterm{Neonatal Thyrotoxicosis} Excess thyroid activity in a newborn, usually caused by thyroid-stimulating antibodies passed from the mother. It may cause a fast heartbeat, irritability, poor feeding, poor weight gain, sweating, or breathing problems and usually needs temporary treatment and monitoring.

\medterm{Neonatal Transition} The newborn’s adjustment from receiving oxygen through the placenta to breathing with the lungs after birth. The lungs fill with air, blood flow changes, and the baby must begin controlling temperature, blood pressure, and feeding.

\medterm{Neonate} A neonate is a newborn infant during the first 28 completed days after birth, a period requiring attention to feeding, temperature, jaundice, infection, breathing, and transition to extrauterine life.

\medterm{Neonatologist} A pediatrician who specializes in the medical care of newborn infants, particularly
premature and critically ill neonates.

\medterm{Neoplasia} Abnormal, uncontrolled growth of cells that forms a new mass or tissue. It may be benign and remain localized, or malignant and invade nearby tissues or spread elsewhere.

\medterm{Neoplasm} A new, abnormal growth of cells that forms a mass or tissue. A neoplasm may be benign and stay localized, or malignant and invade nearby tissues or spread to distant organs.

\medterm{Neoplastic} Relating to an abnormal new growth of cells. A neoplastic growth may be noncancerous and localized or cancerous and able to invade nearby tissue or spread.

\medterm{Neostigmine} An acetylcholinesterase inhibitor used to reverse selected nondepolarizing neuromuscular blockade and to treat some disorders of neuromuscular transmission or bowel and bladder motility.

\medterm{Nephoptosis} Abnormal downward movement or mobility of a kidney, also called floating kidney.

\medterm{Nephrectomy} Surgery to remove part or all of a kidney. It may treat kidney cancer, severe kidney damage, infection, or a kidney that no longer works.

\medterm{Laparoscopic Nephrectomy} A minimally invasive surgical procedure in which a diseased or donor kidney is dissected and extracted using laparoscopy or robotic assistance, categorized as radical (en bloc removal with Gerota fascia and adrenal gland for renal cell carcinoma) or partial (nephron-sparing surgery).

\medterm{Nephric} Relating to the kidney. In embryology, nephric structures are involved in development of the urinary system.

\medterm{Nephritic Syndrome} A pattern of kidney inflammation that causes blood in the urine, reduced urine output, swelling, and high blood pressure. Protein may also leak into the urine. It can reduce kidney function and requires evaluation to find and treat the cause.

\medterm{Nephritis} Inflammation of kidney tissue that can cause blood or protein in the urine, swelling, high blood pressure, reduced urine, or loss of kidney function. Infection, medicines, and immune diseases are possible causes.

\medterm{Nephroblastoma} Another name for Wilms tumor, the most common kidney cancer in children. It may appear as a painless abdominal lump, abdominal pain, blood in the urine, fever, or high blood pressure and needs specialist treatment.

\medterm{Nephrocalcinosis} Buildup of calcium deposits inside kidney tissue. It may result from high calcium levels, hormone or kidney disorders, inherited conditions, dehydration, or medicines and can lead to stones, blood in the urine, or reduced kidney function.

\medterm{Nephrogenic Diabetes Insipidus} A condition in which the kidneys do not respond properly to vasopressin, the hormone that helps conserve water. It causes large amounts of dilute urine and intense thirst; inherited changes, kidney disease, electrolyte problems, or medicines may cause it.

\medterm{Nephrogenic Fibrosing Dermopathy} An older name for nephrogenic systemic fibrosis, a rare condition causing thick, tight, hardened skin and sometimes stiff joints in people with severe kidney disease, usually after exposure to certain gadolinium contrast agents.

\medterm{Nephrogenic Systemic Fibrosis} A rare disorder in which skin and connective tissue become thick and hardened, sometimes limiting movement. It is linked mainly to certain gadolinium contrast exposures in people with severe kidney impairment and may affect internal tissues.

\medterm{Nephrolithiasis} Formation of stones in the kidneys or urinary tract. Stones can cause severe side or back pain, blood in the urine, nausea, or blockage; dehydration, inherited conditions, medicines, and diet can contribute.

\medterm{Nephrologist} A doctor who specializes in kidney disease. Nephrologists assess reduced kidney function, high blood pressure related to kidney problems, kidney inflammation, and mineral or salt disturbances and may manage dialysis or transplant care.

\medterm{Nephrology} The medical specialty concerned with kidney disease, fluid and electrolyte disorders, hypertension, dialysis, and kidney transplantation.

\medterm{Nephrolysis} Surgical freeing of a kidney from surrounding adhesions or tissue.

\medterm{Nephron} The microscopic working unit of a kidney. Each nephron filters blood, returns needed water and substances to the body, and removes waste to make urine; each kidney contains about one million.

\medterm{Nephron Function} Each kidney contains about one million nephrons, its microscopic filtering units. A nephron filters blood, removes wastes into urine, and returns needed water and substances to the bloodstream while helping control salts, acid, and blood pressure.

\medterm{Nephronophthisis} A rare inherited kidney disease that gradually reduces the kidneys’ ability to concentrate urine. It may cause excessive thirst and urination, anemia, growth problems, and eventually kidney failure.

\medterm{Nephropathy} Disease or damage of the kidneys. It may reduce waste removal, disturb fluid and salt balance, cause protein or blood in the urine, or eventually lead to kidney failure.

\medterm{Nephroptosis} Abnormal downward movement of a kidney when a person stands, sometimes causing flank discomfort, urinary symptoms, or intermittent obstruction.

\synonyms
Wandering Kidney; Ren Mobilis; Hypermobile Kidney

\medterm{Nephrosclerosis} Hardening and scarring of kidney tissue, commonly caused by long-standing high blood pressure or blood-vessel disease. Progressive damage can reduce filtering ability, cause protein in urine, and lead to chronic kidney disease.

\medterm{Nephrostomy} A procedure that creates a drainage path from the kidney through the skin when urine cannot pass normally to the bladder. A temporary tube usually carries urine into a bag while the blockage is treated.

\medterm{Nephrostomy Tube} A drainage tube passed through the skin into the kidney to carry urine into a bag when the ureter is blocked. Blockage may be caused by a stone, tumor, or narrowing. Infection with blockage can require urgent drainage; the tube may be temporary or long-term.

\medterm{Nephrotic Range Proteinuria} Very large protein loss in the urine, usually more than 3.5 grams per day in an adult. It often indicates serious damage to the kidney filters and can cause swelling, low blood protein, high cholesterol, blood clots, or infection.

\medterm{Nephrotic Syndrome} A kidney disorder caused by damage to the kidney's tiny filtering units, allowing large amounts of vital protein (albumin) to spill into the urine. This loss leads to low blood protein, high cholesterol, frothy urine, and noticeable swelling (edema), especially around the eyes in the morning and in the feet and legs. It also increases the risks of infections and blood clots; common causes include minimal change disease in children and diabetes or kidney inflammation in adults.
\synonyms Nephrosis; Protein-losing nephropathy; Glomerular nephrosis.

\medterm{Nephrotoxic} Describing a substance or condition that can damage the kidneys or reduce their function. Examples include some medicines, poisons, infections, and severe dehydration; risk depends on dose, duration, and kidney health.

\medterm{Nephrotoxicity} Kidney damage caused by a medicine, poison, contrast material, or other substance. It may reduce urine output or raise blood creatinine; risk is lowered by avoiding triggers when possible and adjusting doses for kidney function.

\medterm{Nernst Equation} A fundamental biophysical formula that calculates the equilibrium electrical potential across a semipermeable cell membrane for a specific single ion species based on its intracellular and extracellular concentration gradient at physiological body temperature.

\medterm{Nerve} A bundle of nerve fibers that carries messages between the brain or spinal cord and the rest of the body. Nerves control movement and carry sensations such as touch, pain, and temperature.

\medterm{Nerve Biopsy} Removal of a small piece of a nerve for laboratory examination when serious nerve disease is suspected. It may help identify inflammation, infection, abnormal deposits, or a tumor, but can leave permanent numbness or weakness near the biopsy site.

\medterm{Nerve Block} An injection of numbing medicine near a nerve or group of nerves to temporarily stop pain signals. It may help diagnose the source of pain or provide pain relief during or after a procedure; weakness or other temporary effects can occur.

\medterm{Nerve Compression} Alternate name; see \textit{Pinched Nerve}.

\medterm{Nerve Conduction Studies} Tests that stimulate peripheral nerves with small electrical pulses and record how quickly and strongly signals travel. They help identify nerve damage, compression, loss of the nerve’s insulating covering, or loss of nerve fibers.

\medterm{Nerve Conduction Study} A test that gives small electrical pulses to a nerve and measures how quickly and strongly the signal travels. It helps identify nerve damage, pressure on a nerve, or loss of its protective covering.

\medterm{Nerve Conduction Velocity} The speed at which an electrical signal travels along a peripheral nerve. A nerve-conduction test measures this speed and the signal strength to help identify nerve damage, pressure on a nerve, or damage to its insulating covering.

\medterm{Nerve Deafness} Hearing loss caused by damage to the cochlea, auditory nerve, or central auditory pathways rather than blockage of sound transmission through the outer or middle ear.

\synonyms
Sensorineural Hearing Loss; SNHL

\medterm{Nerve Endings} The tiny ends of nerves that detect sensations such as touch, pain, heat, or pressure, or release signals to muscles and glands. Injury or irritation can cause pain, numbness, tingling, or abnormal sweating.

\medterm{Nerve Fiber Types} A physiological classification of peripheral nerve axons based on diameter and conduction velocity: group A (thickly myelinated motor and sensory fibers: alpha, beta, gamma, delta), group B (thinly myelinated preganglionic autonomic fibers), and group C (unmyelinated, slow postganglionic and pain/temperature C-fibers).

\synonyms
Erlanger-Gasser Classification

\medterm{Nerve Injury} Damage to a nerve that can cause pain, numbness, tingling, weakness, altered reflexes, or changes in sweating, blood pressure, digestion, or bladder function. Recovery depends on the nerve, severity, cause, and treatment.

\medterm{Nerve Sheath Myxoma} A benign myxoid peripheral-nerve-sheath tumor, usually presenting as a slow-growing papule or nodule on the fingers, hands, or upper extremities.

\medterm{Nerve Sparing} A surgical approach that tries to preserve nearby nerves needed for erection, sensation, or other functions. It may improve recovery after prostate or other pelvic surgery, but success depends on disease and nerve involvement.

\medterm{Nerve Tract} A group of nerve fibers traveling together inside the brain or spinal cord. Fibers in a tract carry signals between specific regions, so damage can cause predictable movement, sensation, or organ-control problems.

\medterm{Nervous Colon Syndrome} An older name for irritable bowel syndrome, a disorder causing recurrent abdominal pain with diarrhea, constipation, or both without a structural explanation. Symptoms may improve with individualized diet, medicines, and stress management.

\medterm{Nervous System} The brain and spinal cord (central nervous system), peripheral nerves and ganglia (peripheral nervous system), and supporting cells that detect and process information. It coordinates sensation, movement, automatic body functions, cognition, emotion, learning, and behavior.

\medterm{Nervous Tissue} Tissue made of nerve cells and supporting cells that detects changes, processes information, and sends signals through the body. It forms the brain, spinal cord, and nerves and helps control movement, sensation, and organ function.

\medterm{Nesidioblastosis} A rare diffuse or focal hyperfunctional hypertrophy and hyperplasia of pancreatic beta cells budding from the pancreatic ductal epithelium, causing endogenous persistent hyperinsulinemic hypoglycemia in neonates (congenital hyperinsulinism) and adults (post-gastric bypass hyperinsulinemic hypoglycemia).
\synonyms
Congenital Hyperinsulinism; Pancreatic Beta-Cell Hyperplasia

\medterm{NET} Abbreviation; see \textit{Neuroendocrine Tumors}.

\medterm{Netherton Syndrome} A severe autosomal recessive genodermatosis caused by mutations in the SPINK5 gene encoding the serine protease inhibitor LEKTI, characterized by the classic triad of congenital ichthyosiform erythroderma (ichthyosis linearis circumflexa), trichorrhexis invaginata ('bamboo hair'), and severe atopic diathesis.
\synonyms
SPINK5 Deficiency; Ichthyosis Linearis Circumflexa Syndrome

\medterm{Nettle Rash} An itchy rash made of raised patches that come and go, often within hours. It results from substances released in the skin and may be triggered by infection, foods, medicines, pressure, heat, or no clear cause.

\medterm{Neural} Relating to nerves, nerve cells, or the nervous system. In pregnancy, the neural tube develops into the brain and spinal cord; failure of it to close can cause birth defects such as spina bifida or anencephaly.

\medterm{Neural Crest} A temporary embryonic multipotent cellular population located along the borders of the closing neural folds that undergoes epithelial-to-mesenchymal transition and migrates extensively to form peripheral sensory ganglia, autonomic ganglia, Schwann cells, melanocytes, craniofacial bones, and adrenal chromaffin cells.

\medterm{Neuralgia} Neuralgia is sharp, burning, or electric pain along the distribution of a nerve, caused by irritation, compression, infection, injury, or other disease.

\medterm{Neuralgia Amyotrophica} An acute, idiopathic autoimmune inflammatory brachial plexopathy (Parsonage-Turner syndrome) characterized by the sudden onset of severe, unprovoked, sharp unilateral shoulder and periscapular pain, followed days to weeks later by patchy flaccid muscle weakness and rapid wasting in the distribution of the upper brachial plexus.
\synonyms
Parsonage-Turner Syndrome; Brachial Radiculitis

\medterm{Neural Tube Defects} Birth defects caused when the early structure that becomes the brain and spinal cord does not close completely during the first month of pregnancy. They include anencephaly, encephalocele, and spina bifida and vary from fatal to mild.

\medterm{Neuraminidase Inhibitors} Antiviral medicines that block a surface protein used by influenza viruses to release new virus particles from infected cells. They work best when started early and are used according to the influenza type, illness severity, and risk of complications.

\medterm{Neurapraxia} The mildest form of peripheral nerve injury in the Seddon classification, characterized by transient conduction block caused by focal demyelination or localized ischemia without structural disruption of the underlying axon or its connective tissue sheaths; complete recovery occurs spontaneously within days to weeks.

\medterm{Neurilemma} The outer layer of a Schwann cell that surrounds and supports a nerve fiber outside the brain and spinal cord.

\medterm{Neurilemmoma} Alternate spelling; see \textit{Schwannoma}.

\medterm{Neurinomatosis} Alternate name; see \textit{Schwannomatosis}.

\medterm{Neuritic Plaques} Abnormal beta-amyloid deposits surrounded by damaged nerve branches, commonly found in Alzheimer disease. They are one pathological feature of the disease and occur with other changes, including abnormal tau protein and nerve-cell loss.

\synonyms
Senile Plaques; Amyloid Plaques

\medterm{Neuritis} Inflammation or irritation of a nerve, often causing pain, altered sensation, weakness, or other changes in nerve function.

\medterm{Neuroacanthocytosis} A group of rare inherited disorders causing movement or psychiatric symptoms together with abnormal spiny red blood cells called acanthocytes. Neurologic assessment and genetic evaluation may be needed.

\medterm{Neuroblastoma} A childhood cancer arising from immature nerve cells, commonly in or near the adrenal glands or spine. It may cause an abdominal mass, pain, weakness, or other symptoms and can range from mild to aggressive.

\synonyms
Sympathicoblastoma; Neural Crest Tumor

\medterm{Neurocardiogenic Syncope} Alternate name; see \textit{Vasovagal Syncope}.

\medterm{Neurocognitive Disorder} A decline from a person’s previous level of cognitive ability, such as memory, attention, language, executive function, or judgment, that develops after normal earlier development. Mild neurocognitive disorder causes modest decline while independence is largely preserved; major neurocognitive disorder causes substantial decline that interferes with independent daily life and is often called dementia. Causes include brain disease, stroke, injury, infection, substance use, or medicines. Delirium is distinguished by its sudden onset and fluctuating attention and awareness.

\medterm{Neurocognitive Disorders} Alternate plural form; see \textit{Neurocognitive Disorder}.

\medterm{Neurocutaneous Syndrome} A disorder affecting both the nervous system and the skin, often because of an inherited genetic change. Examples include neurofibromatosis and tuberous sclerosis.

\medterm{Neurocysticercosis} Infection of the brain or spinal cord by the larval form of the pork tapeworm. It can cause seizures, headaches, confusion, or movement problems and is acquired by swallowing tapeworm eggs from contaminated food or hands.

\medterm{Neurocytoma} A usually slow-growing tumor made of nerve-cell-like cells, most often found in the fluid-filled spaces inside the brain. It may block fluid flow and cause headache, vomiting, vision changes, or balance problems; imaging and tissue testing guide care.

\medterm{Neurodegeneration} Progressive damage and loss of nerve cells, causing worsening problems with movement, memory, thinking, sensation, or other nervous-system functions.

\medterm{Neurodegeneration With Brain Iron Accumulation} A group of inherited disorders in which excess iron builds up in certain brain regions and gradually damages nerve cells. Symptoms may include abnormal movements, muscle stiffness, seizures, speech or thinking problems, and vision changes.

\medterm{Neurodegenerative Disease} A disease in which nerve cells progressively lose function or die. Examples include Alzheimer disease, Parkinson disease, and Huntington disease.

\medterm{Neurodegenerative Disorder} Another term for a condition involving progressive damage or loss of nerve cells, which can affect movement, memory, behavior, sensation, or autonomic function.

\medterm{Neurodermatitis} A chronic itchy skin condition in which repeated scratching or rubbing makes a patch thick, leathery, and irritated. Stress, dry skin, or an initial rash may start the itch–scratch cycle.
\synonyms
Lichen Simplex Chronicus; Circumscribed Neurodermatitis

\medterm{Neurodevelopmental Disorders} A group of conditions that begin during the developmental period and affect cognition, communication, behavior, motor function, learning, or adaptive skills. Examples include attention-deficit/hyperactivity disorder, autism spectrum disorder, intellectual disability, communication disorders, specific learning disorder, motor disorders, and tic disorders.

\medterm{Neuroendocrine} Relating to the connection between the nervous system and hormone system, or to cells that receive nerve signals and can release hormones. Neuroendocrine cells are found in many organs and can sometimes form tumors.

\medterm{Neuroendocrine Carcinoma} A cancer arising from cells with nerve-like and hormone-related features. It can begin in the lung or digestive system, often grows quickly, and may release hormones that cause additional symptoms.

\medterm{Neuroendocrine Carcinoma of the Skin} Alternate name; see \textit{Merkel Cell Carcinoma}.

\medterm{Neuroendocrine Tumor} A growth arising from cells that receive nerve signals and can release hormones. It may begin in organs such as the intestine, pancreas, lung, or appendix and may grow slowly or rapidly; some tumors cause hormone-related symptoms.

\medterm{Neuroendocrine Tumors} Growths arising from hormone-producing cells that also respond to nerve signals. They can occur in many organs, may release excess hormones, and range from slow-growing tumors to aggressive cancers.

\synonyms
NETs; Neuroendocrine Neoplasms; Carcinoid Tumors

\medterm{Neuroendocrine Tumors Grading} Classification of a neuroendocrine tumor by how quickly its cells divide, measured by the Ki-67 index and mitotic count. Lower grades usually grow more slowly; higher grades grow faster and may behave more aggressively. The grade helps guide treatment and prognosis.

\medterm{Neuroendocrinology} The medical field that studies how the brain and nerves control hormones and how hormones affect the nervous system. It includes conditions involving the pituitary, adrenal glands, thyroid, reproduction, growth, and body-fluid balance.

\medterm{Neurofibrillary Tangles} Twisted clumps of an abnormal protein called tau inside nerve cells. They damage the cells and are a characteristic finding in Alzheimer disease and some other disorders that cause progressive memory or thinking problems.

\synonyms
Paired Helical Filaments; Tau Tangles

\medterm{Neurofibroma} A benign peripheral nerve sheath tumor. Solitary neurofibromas are common and present as
soft, flesh-colored papules or nodules. Multiple neurofibromas are characteristic of
neurofibromatosis type 1 (NF1).

\medterm{Neurofibromatosis} A group of inherited conditions in which tumors or other changes develop in tissues containing nerves. Type 1 commonly causes café-au-lait spots, soft nerve tumors, bone changes, and learning difficulties; type 2 mainly affects hearing and balance nerves and can cause other nervous-system tumors.

\medterm{Neurofibromatosis-1} An inherited condition that can cause light-brown skin patches, freckles in the armpits or groin, soft growths along nerves, bone changes, vision problems, learning difficulties, or blood-vessel complications. Findings vary widely between people.

\medterm{Neurofibromatosis 2} An inherited condition in which tumors can grow on nerves, especially the nerves for hearing and balance. It may cause hearing loss, ringing in the ears, balance problems, cataracts, or other brain and spinal-cord tumors.

\medterm{Neurofibromatosis Type 1} An autosomal dominant inherited condition caused by a disease-causing change in the \textit{NF1} gene. It may cause light-brown skin patches, freckles in skin folds, soft growths along nerves, bone problems, vision complications, learning difficulties, or blood-vessel problems.
\synonyms
Neurofibromatosis Type I
NF1
Peripheral Neurofibromatosis
Von Recklinghausen’s Disease
Von Recklinghausen’s Neurofibromatosis
Von Recklinghausen’s Phakomatosis

\medterm{Neurofibromatosis Type 2} An autosomal dominant genetic disorder, usually caused by a disease-causing change in the \textit{NF2} gene, characterized by tumors involving nerves, especially bilateral vestibular schwannomas, with possible hearing loss, balance problems, and other nerve tumors.

\synonyms
NF2; Bilateral Acoustic Neurofibromatosis; MISME Syndrome

\medterm{Neurofibromatosis Type I} Alternate spelling; see \textit{Neurofibromatosis Type 1}.

\medterm{Neurofibromatosis Type II} Alternate spelling; see \textit{Neurofibromatosis Type 2}.

\medterm{Neurofibrosarcoma} Alternate name; see \textit{Malignant Peripheral Nerve Sheath Tumor}.

\medterm{Neurogenic} Caused by, arising from, or related to the nervous system. For example, neurogenic bladder results from nerve damage that disrupts bladder storage or emptying. The specific symptoms and treatment depend on which nerves are affected.

\medterm{Neurogenic Bladder} Bladder problems caused by damage to the nerves that control storing or passing urine. A person may have urgency, leakage, difficulty emptying, or urine retention; repeated infection and high bladder pressure can damage the kidneys.

\medterm{Neurogenic Claudication} A characteristic symptom complex of lumbar spinal canal stenosis characterized by bilateral dull aching, numbness, cramping, and heaviness in the buttocks, thighs, and calves provoked by walking or standing (lumbar extension) and relieved promptly by lumbar flexion (sitting or leaning forward over a shopping cart).
\synonyms
Pseudoclaudication; Spinal Stenosis Claudication

\medterm{Neurogenic Megacolon} Severe widening of the large intestine caused by damage to the nerves that control bowel movement. It can cause long-term constipation, abdominal swelling, and stool blockage and may lead to dangerous bowel inflammation or rupture.

\medterm{Neurogenic Orthostatic Hypotension} A drop in blood pressure after standing because nerve damage prevents the body from tightening its blood vessels and maintaining pressure. It may cause light-headedness, blurred vision, weakness, fainting, or falls.

\medterm{Neurogenic Pulmonary Edema} Rapid-onset, non-cardiogenic pulmonary edema developing acutely after severe central nervous system insults (such as massive subarachnoid hemorrhage, head trauma, or status epilepticus), provoked by an intense sympathetic surge that induces sudden pulmonary vasoconstriction and capillary leak.

\synonyms
NPE

\medterm{Neurogenic Shock} A dangerous fall in blood pressure caused by sudden loss of nerve control over blood-vessel tone, usually after a serious spinal-cord injury. The heart rate may be unusually slow and the skin warm.

\medterm{Neuroglia} Supporting cells of the nervous system that nourish, protect, insulate, and maintain nerve cells. Different types help control the nerve environment, form protective coverings, remove waste, and respond to injury or infection.

\medterm{Neurography} An MRI method that produces detailed pictures of nerves outside the brain and spinal cord. It can help assess trapped, injured, inflamed, or enlarged nerves and may help find some nerve tumors.

\medterm{Neurohypophysis} The back part of the pituitary gland that stores and releases oxytocin and vasopressin. These hormones help control childbirth, milk release, water balance, and urine concentration.

\medterm{Neurolathyrism} A toxic neurologic disorder caused by excessive consumption of Lathyrus sativus, producing progressive spastic weakness of the legs.

\medterm{Neuroleptic} An older term for an antipsychotic medicine, used to treat conditions such as schizophrenia, severe mania, or psychosis. The term remains in names such as neuroleptic malignant syndrome; these medicines may cause sleepiness, movement problems, or metabolic effects.

\medterm{Neuroleptic Agents} Antipsychotic medicines used for conditions such as schizophrenia, bipolar disorder, severe agitation, or nausea, depending on the drug. They can cause movement, metabolic, heart-rhythm, and sedation-related effects.

\synonyms
Antipsychotics

\medterm{Neuroleptic Malignant Syndrome} A rare, life-threatening reaction to some antipsychotic or other dopamine-blocking medicines. It causes very high temperature, severe muscle stiffness, confusion, sweating, unstable blood pressure, and a fast heart rate.

\medterm{Neurological} Related to the brain, spinal cord, or nerves. Neurological symptoms may include weakness, numbness, seizures, headache, dizziness, confusion, or problems with movement, speech, vision, or memory.

\medterm{Neurological Disease} Any disease affecting the brain, spinal cord, peripheral nerves, muscles, or their connections. Symptoms may include weakness, numbness, seizures, headache, movement changes, or problems with thinking.

\medterm{Neurological Examination} A clinical assessment of mental status, cranial nerves, motor function, reflexes, coordination, sensation, and gait.
It helps localize nervous-system dysfunction and determine the need for further testing.

\medterm{Neurological Rehabilitation} Treatment and training that help a person regain movement, communication, thinking, or independence after stroke, brain or spinal-cord injury, Parkinson disease, multiple sclerosis, or another nervous-system disorder.

\medterm{Neurologic Deficit} Loss or impairment of a neurologic function, such as strength, sensation, coordination, vision, language, or consciousness, caused by dysfunction of the nervous system.

\medterm{Neurologic Disorders} A group of conditions affecting the brain, spinal cord, peripheral nerves, neuromuscular junction, or muscles. Examples include stroke, epilepsy, migraine, multiple sclerosis, Parkinson disease, Alzheimer disease, peripheral neuropathy, meningitis, neuromuscular disorders, and movement disorders.

\medterm{Neurologic Localization} Using the pattern of weakness, numbness, reflex changes, coordination problems, vision, speech, or thinking changes to identify where the nervous system is affected. This helps clinicians choose likely causes and appropriate tests.

\medterm{Neurologist} A medical doctor who diagnoses and treats disorders of the brain, spinal cord, nerves, and muscles. Neurologists assess symptoms such as weakness, numbness, seizures, tremor, headaches, memory change, and movement or balance problems.

\medterm{Neurology} The medical specialty concerned with diseases of the brain, spinal cord, peripheral nerves, and muscles.

\medterm{Neurology PET} PET imaging used in selected neurologic disorders to assess brain metabolism or receptor activity, such as in some dementias, epilepsy evaluations, and movement disorders.

\medterm{Neurolysis} A procedure that frees a nerve from scar tissue or pressure, or deliberately interrupts selected nerve fibers to relieve pain. The purpose and possible effects depend on the nerve treated and the technique used.

\medterm{Neuroma} A benign growth or reactive thickening involving nerve tissue, often producing pain or altered sensation.

\medterm{Neuromodulation} Changing nerve activity with electrical, magnetic, chemical, or other signals. Treatments include spinal cord stimulation, deep brain stimulation, and some noninvasive brain-stimulation methods.

\medterm{Neuromuscular} Relating to nerves, muscles, or the place where a nerve tells a muscle to contract. Neuromuscular diseases can cause weakness, tiredness with activity, muscle wasting, cramps, or problems with breathing or swallowing.

\medterm{Neuromuscular Blockade Reversal} The pharmacological restoration of normal neuromuscular transmission following surgical paralysis, achieved by administering an acetylcholinesterase inhibitor paired with an anticholinergic (neostigmine and glycopyrrolate) or a selective encapsulating reversal agent (sugammadex for rocuronium/vecuronium).

\medterm{Neuromuscular Blocking Agents} Medicines used during anesthesia to temporarily relax skeletal muscles, making insertion of a breathing tube and some operations easier. They stop muscle movement but do not provide pain relief or unconsciousness and require breathing support and monitoring.

\medterm{Neuromuscular Disorders} A group of conditions affecting motor nerves, neuromuscular junctions, or skeletal muscles. Examples include myasthenia gravis, Lambert-Eaton syndrome, muscular dystrophy, spinal muscular atrophy, motor neuron disease, inflammatory myositis, myopathies, and peripheral neuropathies.

\medterm{Neuromuscular Hamartoma} A rare benign developmental lesion in which mature skeletal muscle and nerve fascicles are intermixed, typically involving a major peripheral nerve and causing enlargement or neuropathic symptoms.

\medterm{Neuromuscular Junction} The small connection where a movement-controlling nerve communicates with a skeletal muscle. The nerve releases a chemical messenger there, causing the muscle to contract.

\synonyms
NMJ; Motor End-Plate

\medterm{Neuromyelitis Optica} An autoimmune disorder that mainly causes inflammation of the optic nerves and spinal cord. It can cause vision loss and weakness or altered sensation and is treated with immune-directed therapy.

\medterm{Neuromyelitis Optica Spectrum Disorder} An immune-system disorder that repeatedly inflames the optic nerves, spinal cord, or parts of the brain. Attacks may cause vision loss, limb weakness, numbness, vomiting, or bladder problems and can lead to lasting disability without treatment.

\synonyms
Neuromyelitis Optica
NMO
NMOSD
Devic Disease
Devic's Disease

\medterm{Neuromyotonia} A rare disorder in which nerves repeatedly stimulate muscles, causing ongoing stiffness, twitching, cramps, delayed muscle relaxation, and sometimes excessive sweating. It may be inherited or caused by an immune reaction. An electrical nerve and muscle test helps diagnosis; treatment aims to calm nerve activity and address any underlying cause.

\synonyms
Isaacs Syndrome; Continuous Muscle Fiber Activity Syndrome

\medterm{Neuron} A nerve cell that receives, processes, and sends messages using electrical signals and chemical messengers. It has a cell body, branches that receive information, and usually a long fiber that carries messages to other cells.

\medterm{Neuronal Ceroid Lipofuscinoses} A group of inherited disorders in which waste material builds up inside nerve cells. They usually cause worsening seizures, vision loss, movement problems, and difficulty thinking or learning, with symptoms and age of onset varying by type.

\medterm{Neuronal Migration Disorder} A developmental brain disorder in which nerve cells do not move to their normal positions before birth. It can cause abnormal brain structure, seizures, developmental delay, or learning problems.

\medterm{Neurone} The British spelling of neuron, an excitable nerve cell that receives and transmits signals.

\medterm{Neuronopathy} A disorder caused mainly by damage to nerve-cell bodies, often in sensory or motor nerve groups. It may cause numbness, imbalance, weakness, or loss of reflexes.

\medterm{Neurons} Specialized nerve cells that receive, process, and transmit electrical or chemical signals. They communicate with one another and with muscles or glands.

\medterm{Neuron-Specific Enolase} A protein that can be released by nerve cells and some neuroendocrine tumors. A raised blood level may support evaluation or monitoring in selected situations but is not specific enough to diagnose a disease alone.

\medterm{Neuron-Specific Enolase Tumor Marker} Neuron-specific enolase, a protein that may be increased in some neuroendocrine cancers or after nerve-cell injury. It can support monitoring in selected situations but is not specific enough to diagnose cancer alone.

\medterm{Neuro-Ophthalmology} The medical specialty that evaluates vision problems caused by the brain, optic nerves, or eye-movement pathways. It assesses symptoms such as unexplained vision loss, double vision, abnormal pupils, or visual-field changes.

\medterm{Neuropathic} Caused by disease or dysfunction of nerves; neuropathic skin symptoms may include burning, electric-shock pain, altered sensation, numbness, or itch.

\medterm{Neuropathic Pain} Pain caused by disease or injury of the nerves that carry sensation, the spinal cord, or the brain. It is often burning, shooting, electric, or accompanied by numbness and may follow diabetes, shingles, injury, or nerve compression.

\medterm{Neuropathic Ulcer} A wound caused or worsened by loss of protective feeling, most often in the foot. A person may not notice pressure or injury, allowing the wound to deepen; diabetes, poor circulation, and infection can delay healing and increase the risk of amputation.

\medterm{Neuropathy} Disease or damage affecting nerves outside the brain and spinal cord. It may cause burning pain, numbness, tingling, weakness, balance problems, or altered sweating and can result from diabetes, medicines, toxins, poor nutrition, compression, or inherited disease.

\medterm{Neuropeptide} A peptide released by neurons that acts as a signaling molecule in the nervous or endocrine system.

\medterm{Neuropeptides} Short chains of amino acids released by nerve cells as chemical signals. They can influence pain, appetite, stress, mood, sleep, blood pressure, inflammation, hormone release, and communication between nerve cells.

\medterm{Neuropharmacology} The study of how medicines and other substances affect the brain, spinal cord, and nerves. It examines their actions, benefits, side effects, interactions, and uses in treating neurological or psychiatric conditions.

\medterm{Neurophysiotherapy} Physical therapy focused on movement and function problems caused by disorders of the brain, spinal cord, or nerves. Treatment may address strength, balance, walking, coordination, and daily activities.

\medterm{Neuroplasticity} The brain’s ability to change its connections and activity in response to learning, experience, or injury. This process supports recovery and skill development, especially with repeated practice and rehabilitation, but recovery varies with the injury and the person.

\medterm{Neuroplegic} Causing paralysis or severe loss of nerve-controlled movement. The effect may result from a neurologic disease, spinal injury, toxin, or medicine and can impair breathing or other vital functions when major nerve pathways are affected.

\medterm{Neuropsychiatry} The clinical field linking neurologic and psychiatric disorders, including cognitive and behavioral consequences of brain disease.

\medterm{Neuroradiology} The use of imaging to diagnose disorders of the brain, spine, head and neck, and peripheral nerves.

\medterm{Neuroretinitis} Inflammation of the optic nerve and nearby retina that can cause blurred or reduced vision, blind spots, and pain with eye movement. Infections and immune disorders are possible causes and require specialist evaluation.

\medterm{Neurosarcoidosis} Sarcoidosis affecting the brain, spinal cord, coverings around them, or peripheral nerves. It may cause headaches, facial weakness, vision problems, seizures, weakness, numbness, or hormone disturbances and usually needs specialist treatment.

\medterm{Neurosciences} Neurosciences are the scientific disciplines studying the nervous system, including its structure, function, development, disease, and treatment.

\medterm{Neurosurgeon} A surgeon trained to diagnose and operate on disorders of the brain, spinal cord, spine, nerves, and nearby supporting structures. Neurosurgeons may treat tumors, injuries, bleeding, pressure, pain, or blood-vessel problems.

\medterm{Neurosurgery} Surgery for disorders of the brain, spinal cord, spine, peripheral nerves, or supporting structures. Procedures may remove tumors, relieve pressure, repair injuries, stop bleeding, treat blood vessels, or improve severe pain.

\medterm{Neurosyphilis} Neurosyphilis is infection of the nervous system by Treponema pallidum and can cause meningitis, stroke, sensory loss, cognitive change, or tabes dorsalis.

\medterm{Neurotoxicity} Injury or dysfunction of the nervous system caused by a medicine, drug, toxin, heavy metal, infection, or other harmful exposure.

\medterm{Neurotransmission} The process by which nerve cells communicate with one another or with muscles and glands using electrical signals and chemical messengers. It involves release of a neurotransmitter, movement across a synapse, receptor binding, and termination or removal of the signal.

\medterm{Neurotransmitter} A chemical messenger released by a nerve cell. It crosses a tiny gap to affect another nerve cell or a muscle or gland cell, either increasing or reducing its activity.

\medterm{Neurotransmitter Receptors} Proteins on or inside cells that receive chemical messages released by nerve cells. When a message attaches to a receptor, it can increase or decrease activity, affecting movement, mood, pain, heart rate, or other body functions.

\medterm{Neurotrophic} Describing a substance or effect that supports the growth, survival, connection, or repair of nerve cells. Neurotrophic signals help nerves develop and maintain function and are studied in nerve injury and degenerative disease.

\medterm{Neurotrophic Keratopathy} Corneal damage caused by reduced sensation from injury or disease of the trigeminal nerve. It can lead to poor healing, recurrent epithelial defects, ulceration, and vision loss, sometimes with little pain.

\medterm{Neurotropic} Having a preference for nervous tissue or mainly affecting the brain, spinal cord, or nerves. The term is often used for infections, medicines, or other substances that act in these tissues.

\medterm{Neurovascular Examination} A check of movement, feeling, pulses, skin color, temperature, and blood return in an injured arm or leg. Repeated checks can detect worsening nerve damage, reduced blood flow, or dangerous swelling and pressure.

\medterm{Neutral Alignment} A posture in which the body is arranged in a straight, balanced line while preserving normal spinal curves.

\medterm{Neutral Endopeptidase} A membrane-bound zinc metallopeptidase expressed in vascular endothelial and renal proximal tubular cells that degrades circulating natriuretic peptides (ANP, BNP, CNP), bradykinin, and substance P; pharmacologically inhibited by sacubitril in heart failure therapy.

\synonyms
Neprilysin; NEP; CD10

\medterm{Neutralization} Removal or reduction of an effect by counteraction, such as an acid-base reaction or antibody-antigen binding.

\medterm{Neutralizing Antibodies} Antibodies that bind a pathogen or its toxin and block its ability to enter cells, attach to tissues, or cause its usual biologic effects.

\medterm{Neutropenia} An abnormally low number of neutrophils, white blood cells that help fight bacterial and fungal infections. The lower the count, the greater the infection risk; fever during severe neutropenia can signal a serious infection.

\synonyms
Neutrocytopenia; Low Absolute Neutrophil Count

\medterm{Neutropenic Enterocolitis} A life-threatening, necrotizing transmural inflammation of the cecum and ascending colon occurring in severely neutropenic oncology patients receiving intensive myelosuppressive chemotherapy, presenting with fever, right lower quadrant abdominal pain, and bowel wall thickening on CT.

\synonyms
Typhlitis

\medterm{Neutropenic Fever} A medical emergency in oncology patients defined as a single oral temperature of 38.3 degrees Celsius (101 F) or sustained temperature of 38.0 degrees Celsius for one hour in the setting of an absolute neutrophil count (ANC) below 500 cells/microliter, requiring immediate empiric broad-spectrum antipseudomonal antibiotic administration.

\synonyms
Febrile Neutropenia

\medterm{Neutrophil} The most common infection-fighting white blood cell in the bloodstream. It quickly enters injured tissue, surrounds and destroys many bacteria and fungi, and helps start inflammation; unusually low levels increase infection risk.

\medterm{Neutrophil Extracellular Traps} Web-like material released by neutrophils, made mainly of DNA and antimicrobial proteins. These traps can help contain germs, but excessive formation may injure tissues and contribute to inflammation, clotting, or autoimmune disease.

\medterm{Neutrophil Extracellular Traps} Webs of extracellular decondensed chromatin fibers decorated with antimicrobial granule proteins (myeloperoxidase, elastase, histones) extruded by activated neutrophils during a specialized form of cell death (NETosis) to trap and kill extracellular bacterial pathogens; implicated in thrombosis and autoimmune disease.

\synonyms
NETs

\medterm{Neutrophilic Dermatosis} A group of inflammatory skin diseases in which infection-fighting white blood cells collect in the skin without causing a typical skin infection. Examples include Sweet syndrome and pyoderma gangrenosum; the rash may be painful, swollen, or associated with fever.

\medterm{Neutrophilic Eccrine Hidradenitis} An uncommon skin inflammation in which infection-fighting white cells collect around sweat glands. It often occurs during chemotherapy and causes painful red or purple bumps or patches, usually on the trunk or limbs.

\medterm{Neutrophil Recruitment} The movement of neutrophils from the bloodstream into injured or infected tissue. They first slow and attach to the blood-vessel lining, then pass through it and follow chemical signals to the affected area.

\medterm{Neutrophils} The most abundant type of white blood cell. Neutrophils rapidly migrate to sites of infection or injury and ingest microorganisms and cellular debris.

\medterm{Nevoid Basal Cell Carcinoma Syndrome} Gorlin syndrome, an inherited disorder increasing the risk of basal-cell cancers, jaw cysts, and other developmental abnormalities.

\medterm{Nevus} A usually harmless, localized change or growth in the skin or mucous membrane, often called a mole. It may contain pigment-producing cells or involve blood vessels or other skin structures; a changing lesion needs assessment.

\medterm{Nevus Anemicus} A congenital vascular anomaly presenting as a localized, well-circumscribed, pale, hypopigmented macule or patch that remains pale after rubbing or temperature changes due to localized hypersensitivity of cutaneous alpha-adrenergic receptors producing constant arteriolar vasoconstriction.
\synonyms
Anemic Nevus; Pharmacologic Vasoconstrictive Nevus

\medterm{Nevus Comedonicus} A congenital hamartoma of the pilosebaceous unit presenting as a localised cluster of
dilated follicular ostia filled with keratinous plugs, resembling closely grouped open
comedones. It typically appears at birth or in early childhood on the face, neck, or
trunk, and may be associated with other developmental anomalies.

\medterm{Nevus Comedonicus Syndrome} A rare disorder in which multiple comedone-like skin lesions occur with abnormalities of the eyes, bones, nervous system, or other organs. It is usually present from birth or early childhood.

\medterm{Nevus of Ito} A benign dermal melanocytic hamartoma characterized by mottled, patchy slate-gray, blue-brown, or hyperpigmented macules distributed along the sensory branches of the posterior supraclavicular and lateral cutaneous brachial nerves (involving the supraclavicular, acromial, and scapular regions).
\synonyms
Omohyoid Nevus; Dermal Melanocytosis of Shoulder

\medterm{Nevus of Ota} A congenital or acquired dermal melanocytosis presenting as mottled blue-gray or brownish-black pigmentation in the distribution of the ophthalmic (V1) and maxillary (V2) branches of the trigeminal nerve, commonly involving the ipsilateral sclera, conjunctiva, and facial skin.
\synonyms
Oculodermal Melanocytosis; Nevus Fusco-Caeruleus Ophthalmomaxillaris

\medterm{Nevus Sebaceus} A congenital or early-onset hamartoma of sebaceous and other adnexal structures, usually a yellow-orange hairless plaque on the scalp or face; secondary tumors may rarely arise within it.

\medterm{Newborn} A baby during the first 28 days after birth, when special newborn health monitoring and care are important.

\medterm{Newborn Circumcision} Surgical removal of the foreskin from the penis, typically performed within the first
few days of life. Medical benefits include reduced risk of urinary tract infections,
sexually transmitted infections, and penile cancer. Complications include bleeding,
infection, and meatal stenosis.

\medterm{Newborn Head Molding} Temporary shaping or overlapping of a newborn’s skull bones during passage through the birth canal; it usually resolves as the skull returns to its normal contour.

\medterm{Newborn Hearing Screen} A quick test, usually before discharge from the birth facility, that checks whether a newborn responds normally to sound. A result needing repeat testing does not prove hearing loss; timely follow-up is important when the screen is not passed.

\medterm{Newborn Infant Development} The normal physical, neuromotor, sensory, and social milestone progression of a healthy full-term infant during the first month of life (neonatal period), including primitive reflexes (Moro, rooting, suckling, palmar grasp), visual tracking, auditory responsiveness, and weight recovery.

\synonyms
Neonatal Development; Infant Milestone Progression; Newborn Neuromotor Development

\medterm{Newborn Jaundice} Newborn jaundice is yellow discoloration from elevated bilirubin; physiologic patterns are common, but early, severe, rapidly rising, or conjugated jaundice requires evaluation for serious disease.

\medterm{Newborn Metabolic Screen} A dried blood spot test (heel stick) collected 24-48 hours after birth to screen for
congenital disorders before symptoms develop.

\medterm{Newborn Screening} Testing shortly after birth for selected treatable genetic, metabolic, hormone, hearing, and heart conditions before symptoms appear. A small blood sample is usually taken from the heel; the conditions screened and timing vary by country or region.

\medterm{Newborn Screening Panel} A standardized public health laboratory testing program performed on dried blood spots collected via neonatal heel prick within 24 to 48 hours of life, screening for early, treatable inborn errors of metabolism, hemoglobinopathies, and endocrine disorders (such as phenylketonuria, congenital hypothyroidism, cystic fibrosis, and sickle cell disease).

\medterm{Newborn Screening Tests} Tests performed shortly after birth to detect selected treatable genetic, endocrine, metabolic, hematologic, hearing, or critical congenital heart conditions before symptoms develop.

\medterm{New Daily Persistent Headache} A headache that begins on a clearly remembered day and becomes continuous within 24 hours, persisting for
at least three months. It may resemble migraine or tension-type headache; a new persistent headache requires
medical evaluation to exclude secondary causes such as infection, pressure disorders, or vascular disease.

\medterm{New York Heart Association Functional Classification} A standardized clinical grading scale assessing the severity of functional symptomatic limitation in patients with heart failure based on physical exertion, ranging from Class I (no symptoms with ordinary physical activity) to Class IV (symptoms present at rest).

\synonyms
NYHA Functional Classification; NYHA Class

\medterm{Next-Generation Sequencing} A laboratory method that reads many DNA sections at the same time. It can examine selected genes, nearly all protein-coding genes, the whole genome, or tumor DNA; results may miss some changes and require expert interpretation.

\medterm{Nexus} A connection or central point linking structures, processes, or ideas. In medical writing, its meaning depends on the context and may describe an anatomic connection, a group of related causes, or a link between two findings.

\medterm{NF1} Abbreviation; see \textit{Neurofibromatosis Type 1}.

\medterm{NF2} Abbreviation; see \textit{Neurofibromatosis Type 2}.

\medterm{Niacin} Vitamin B3, needed to help cells release energy and perform many chemical reactions; it forms the coenzymes NAD+ and NADP+ that carry electrons in metabolism. Severe deficiency causes pellagra, with skin changes, digestive problems, confusion, weakness, and potentially life-threatening complications.

\synonyms
Vitamin B3; Nicotinic Acid

\medterm{Niacinamide} The amide form of vitamin B3, used as a nutritional supplement and in selected dermatologic preparations. Unlike niacin, it usually does not cause flushing at nutritional doses.

\medterm{Niacin Deficiency} A lack of vitamin B3, usually from a diet poor in niacin and
the amino acid tryptophan. When severe it causes pellagra, with dermatitis, diarrhea, and
neuropsychiatric disturbance, and it may complicate alcohol use disorder or conditions
that impair absorption. Niacin replacement and treatment of the underlying cause correct
the deficiency.

\medterm{Niacin For Cholesterol} Niacin lowers triglycerides and can raise HDL cholesterol, but its routine cardiovascular benefit is limited when added to effective statin therapy; flushing, liver injury, hyperglycaemia, and hyperuricaemia limit high-dose use.

\medterm{Niacin Vitamin B3} Alternate name; see \textit{Niacin}.

\medterm{Nialamide} An older monoamine-oxidase-inhibitor antidepressant that is no longer commonly used.

\medterm{Nicardipine} A calcium-channel-blocking medicine that relaxes blood vessels and lowers blood pressure. It is used for selected urgent or ongoing blood-pressure problems and may cause headache, flushing, dizziness, fast heartbeat, or swelling.

\medterm{Nickel Allergy} An allergic skin reaction caused by contact with nickel-containing materials.

\synonyms
Nickel Hypersensitivity; Nickel Contact Dermatitis

\medterm{Nickel-Chromium} A metal mixture used in some dental crowns and medical devices. It is strong and less costly than some alternatives, but people allergic to nickel may develop skin or mouth reactions after contact.

\medterm{Niclosamide} An antiparasitic medicine used in some countries to treat selected intestinal tapeworm infections. It acts mainly in the gut and does not reliably treat larval tapeworm infections in tissues, such as neurocysticercosis.

\medterm{Nicolau Syndrome} A rare, acute iatrogenic cutaneous adverse reaction characterized by severe localized pain, livedoid erythema, and rapid ischemic skin necrosis (embolia cutis medicamentosa) occurring at the site of intramuscular or intra-articular drug injections, caused by accidental intra-arterial or peri-arterial drug deposition and microvascular thrombosis.
\synonyms
Embolia Cutis Medicamentosa; Livedo-like Dermatitis

\medterm{Nicotinamide} Alternate name; see \textit{Niacinamide}.

\medterm{Nicotinamide Adenine Dinucleotide} A molecule present in all living cells that helps transfer energy from food and supports enzymes involved in cell repair and communication. It exists in forms commonly written as NAD+ and NADH.

\medterm{Nicotine} An addictive stimulant found in tobacco and some vaping products that acts on neuronal nicotinic receptors.

\medterm{Nicotine Addiction} Alternate name; see \textit{Nicotine Dependence}.

\medterm{Nicotine And Tobacco} Nicotine poisoning can follow swallowing tobacco or liquid nicotine, breathing concentrated vapor, or skin contact. It may cause nausea, vomiting, sweating, dizziness, fast heartbeat, seizures, or breathing failure.

\medterm{Nicotine Dependence} A pattern of compulsive nicotine use with cravings, loss of control, tolerance, and withdrawal symptoms when use is reduced or stopped.

\synonyms
Tobacco Use Disorder; Nicotine Addiction

\medterm{Nicotine Poisoning} Nicotine poisoning causes early nausea, vomiting, salivation, dizziness, tachycardia, and agitation followed by weakness, bradycardia, seizures, or respiratory failure in severe exposure.

\medterm{Nicotine Replacement Therapy} Nicotine medicines help reduce withdrawal when stopping tobacco. Using too much or combining products improperly can cause nausea, vomiting, dizziness, sweating, fast heartbeat, confusion, or seizures; keep them away from children.

\medterm{Nicotinic} Relating to nicotine or to nicotinic acetylcholine receptors. These receptors respond to the chemical messenger acetylcholine and help control nerve signals, skeletal-muscle movement, attention, and some automatic body functions.

\medterm{Nicotinic Acetylcholine Receptor} A pentameric ligand-gated non-selective cation channel responsive to acetylcholine and nicotine, functioning to generate rapid excitatory postsynaptic potentials at the neuromuscular junction (muscle type) and throughout autonomic ganglia and the central nervous system (neuronal type).

\synonyms
nAChR

\medterm{Nicotinic Acid} Alternate name; see \textit{Niacin}.

\medterm{NICU Consultants And Support Staff} NICU consultants and support staff are specialists and professionals who help care for critically ill or premature newborns, including physicians, nurses, therapists, and social workers.

\medterm{NICU Staff} NICU staff are the healthcare professionals who monitor and treat newborns in a neonatal intensive care unit and support their families.

\medterm{Nidation} Implantation of a fertilized egg into the lining of the uterus early in pregnancy. The term is largely synonymous with implantation.

\medterm{Nidus} A starting point or central focus from which a process develops. In medicine, it may mean the core of a kidney stone, the focus of an infection, or the central part of an abnormal blood-vessel network.

\medterm{Niemann-Pick Disease} A group of rare inherited lysosomal storage disorders in which lipids accumulate in cells. The group includes acid sphingomyelinase deficiency, formerly called Niemann-Pick disease types A and B, and Niemann-Pick disease type C, caused by impaired intracellular cholesterol trafficking. Features vary by type and may include enlargement of the liver or spleen, lung disease, difficulty coordinating movement, and progressive neurologic impairment.

\medterm{Nifedipine} A medicine that relaxes blood vessels to lower blood pressure and relieve some types of chest pain. It may cause flushing, headache, dizziness, ankle swelling, or a fast heartbeat; the dose depends on the condition and formulation.

\medterm{Night Blindness} Difficulty seeing in dim light or darkness. Causes include vitamin-A deficiency, inherited retinal disorders, cataracts, glaucoma, and some medicines. New or worsening night blindness needs an eye examination because some causes can be treated.

\medterm{Night-Blindness} Difficulty seeing in dim light, called nyctalopia; causes include vitamin A deficiency, retinal disease, and inherited disorders.

\medterm{Night Eating Syndrome} A recurring pattern of eating much of the day’s food in the evening or after waking at night, often with little morning appetite and sleep problems. It is a related eating pattern that may occur with depression, obesity, or other eating disorders and is not automatically a formal eating-disorder diagnosis.

\medterm{Night Leg Cramps} Sudden, painful involuntary contractions of leg muscles occurring during sleep or rest, often affecting the calf or foot.

\medterm{Nightmare Disorder} Repeated frightening dreams that cause distress, wake a person, or interfere with sleep and daytime activities. It may be linked with stress, trauma, medicines, mental health conditions, or other sleep disorders.

\medterm{Nightmares} Vivid frightening dreams that cause distress or wake a person, sometimes linked to stress, illness, or medicines.

\medterm{Night Sweats} Episodes of excessive sweating during sleep that soak clothing or bedding. They may occur with infection,
hormonal changes, cancer, medication effects, or other systemic illness; persistent unexplained
night sweats warrant assessment.

\synonyms
Sleep Hyperhidrosis; Nocturnal Diaphoresis

\medterm{Night-Sweats} Episodes of excessive sweating during sleep; causes include infection, malignancy, hormonal change, medicines, and environmental heat.

\medterm{Night Sweats in Men} Alternate name; see \textit{Night Sweats}.

\medterm{Night Terrors} Episodes of abrupt partial awakening from deep non-REM sleep with intense fear, screaming, autonomic arousal, and limited recall afterward, occurring most often in children.

\medterm{Night Terrors In Children} Night terrors are episodes of sudden fear, screaming, autonomic arousal, and apparent confusion during non-REM sleep; children usually remain difficult to awaken and have little memory afterward.

\medterm{NIH Stroke Scale} The National Institutes of Health Stroke Scale is a standardized, validated
clinical assessment tool used to quantify the severity of acute ischemic stroke.

\medterm{Nijmegen Breakage Syndrome} A rare inherited disorder that affects DNA repair and the immune system. It may cause a small head, short stature, developmental or learning difficulties, frequent infections, and a high risk of lymphoma or other cancers. People with it are unusually sensitive to radiation, so genetic and specialist care is important.

\synonyms
NBS; Ataxia-Telangiectasia Variant 1

\medterm{Nikethamide} An old medicine once used to stimulate breathing. It is no longer a routine treatment because its benefit is limited and excessive doses can cause agitation, shaking, seizures, or other dangerous effects.

\medterm{Nikolsky Sign} Peeling or separation of the top skin layer when nearby skin is gently rubbed. It can occur in serious blistering disorders, but the sign is not specific and must be interpreted with the rash and other findings.

\medterm{Nikolsky's Sign} Shearing or peeling of the superficial epidermis when apparently normal skin beside a blister is rubbed; it may occur in pemphigus and other blistering disorders.

\medterm{Nipah Virus} A zoonotic virus in the genus \textit{Henipavirus}, with fruit bats as its natural host. It can spread from infected bats or other animals to people, through contaminated food such as raw date-palm sap, and between people through close contact. Infection ranges from no symptoms to fever, headache, vomiting, and respiratory illness such as cough or difficulty breathing. Severe disease may cause inflammation of the brain (encephalitis), confusion, seizures, coma, or death. There is currently no licensed vaccine or specific approved treatment for Nipah virus disease; management relies on supportive care.

\medterm{Nipah Virus Infection} A zoonotic illness caused by Nipah virus. Fruit bats are the natural host, and infection can pass to people through infected animals, contaminated food such as raw date-palm sap, or close contact with an infected person. Some infections cause no symptoms; illness may otherwise cause fever, headache, muscle pain, vomiting, cough, or difficulty breathing and can progress to brain inflammation (encephalitis), confusion, seizures, coma, or death. Early intensive supportive care may improve survival, but there is currently no licensed vaccine or specific approved treatment.

\medterm{Nipple} The projection at the centre of the breast through which milk is released from the mammary ducts; it contains smooth muscle and sensory nerve endings.

\medterm{Nipple Discharge} Fluid leaking from one or both nipples when not pregnant or breastfeeding. Most nipple discharge is harmless, coming from multiple milk ducts in both breasts only when squeezed, and is usually caused by normal hormone changes or medications. Discharge that leaks on its own from a single duct in one breast, especially if clear or bloody, is less common. The most frequent cause is a small, harmless growth inside a milk duct (a papilloma), while breast cancer accounts for about 5\% to 15\% of these cases. Doctors identify the cause using a physical examination, imaging such as ultrasound or mammography, and duct tests.

\medterm{Pathologic Nipple Discharge} Nipple fluid that occurs without squeezing, comes from one breast or one duct, or is bloody or clear. It may have a benign cause, but it needs clinical assessment and often imaging to decide whether further testing is needed.

\medterm{Nipple Inversion Or Retraction} Inward turning or pulling-in of the nipple. It may be longstanding and benign, but new one-sided inversion or retraction can occur with duct disease or breast cancer.

\medterm{Nipple Retraction} Alternate name; see \textit{Nipple Inversion Or Retraction}.

\medterm{Niridazole} An older medicine once used against schistosomiasis, a parasitic infection. Its use is now limited because of serious possible side effects, including seizures, liver injury, and damage to blood-cell production.

\medterm{Nissen Fundoplication} A surgical antireflux procedure in which the gastric fundus is completely wrapped 360 degrees around the lower esophagus to reinforce the lower esophageal sphincter (LES) and repair a hiatal hernia, used for refractory gastroesophageal reflux disease (GERD).
\synonyms
360-Degree Complete Fundoplication; Total Antireflux Fundoplication

\medterm{Nit} The egg of a louse, usually attached firmly to a hair shaft or clothing fiber. Nits may be confused with dandruff, but they do not brush away easily. Finding live lice confirms an active infestation; treatment and fine combing remove lice and eggs.

\medterm{Nitabuch Layer} A normal, dense, fibrinoid layer formed at the junction between the invading syncytiotrophoblast and the decidua basalis (stratum compactum) in normal pregnancy, which acts as a physiological anatomical boundary limiting deep placental invasion (congenitally deficient in placenta accreta).
\synonyms
Stria of Nitabuch; Decidual Fibrinoid Boundary

\medterm{Nitrates} Medicines that widen blood vessels by releasing nitric oxide. They can relieve angina and lower the heart's workload but may cause headache, flushing, dizziness, or dangerously low blood pressure with certain medicines.

\medterm{Nitrate Tolerance} Reduced effect of nitrate medicines, such as nitroglycerin, after continuous use. It may involve changes in blood-vessel signaling and the body’s counter-response to the medicine.

\synonyms
Nitroglycerin Tachyphylaxis; Nitrate Resistance

\medterm{Nitrazepam} A benzodiazepine medicine that slows brain activity and can cause sleepiness, calmness, and seizure prevention. It can cause dependence, falls, memory problems, and dangerously slow breathing, especially with alcohol or opioids.

\medterm{Nitric Acid} A strong corrosive chemical that can burn skin, eyes, mouth, throat, and digestive organs. Breathing its fumes can injure the lungs, and swallowed or inhaled exposure can cause severe poisoning.

\medterm{Nitric Acid Poisoning} Exposure to nitric acid can severely burn the skin, eyes, mouth, throat, and stomach and may damage the lungs if fumes are inhaled.

\medterm{Nitric Oxide} A short-lived signaling molecule made by many cells. It relaxes blood-vessel muscle, helping regulate blood flow and pressure, and also participates in nerve signaling and immune defense. Too much or too little production may contribute to disease.

\medterm{Nitric Oxide Synthase} An enzyme family that catalyzes the synthesis of the gaseous signaling molecule nitric oxide (NO) and L-citrulline from L-arginine, oxygen, and NADPH, comprising constitutive calcium-dependent isoforms (endothelial eNOS, neuronal nNOS) and a calcium-independent inducible isoform (iNOS).

\synonyms
NOS

\medterm{Nitrofurantoin} An antibacterial medicine used mainly for bladder infections caused by susceptible bacteria. It concentrates in urine but is not suitable for many kidney infections; nausea, lung reactions, liver injury, or nerve problems are uncommon but important risks.

\medterm{Nitrogen} An essential element in amino acids, proteins, nucleic acids, and other biological molecules; atmospheric nitrogen is largely inert to humans.

\medterm{Nitrogen Balance} The difference between dietary nitrogen intake, estimated mainly from protein intake,
and nitrogen loss, primarily in urine. Positive balance supports growth or tissue
repair; negative balance occurs with inadequate intake, inflammation, trauma, or
catabolic disease.

\medterm{Nitrogenous Base} A nitrogen-containing component of nucleic acids; adenine, guanine, cytosine, thymine, and uracil pair in DNA or RNA to encode genetic information.

\medterm{Nitroglycerin Overdose} Too much nitroglycerin can cause severe headache, flushing, dizziness, fainting, low blood pressure, fast heartbeat, or breathing difficulty.

\medterm{Nitromors} A trade name for paint or varnish stripper formulations; exposure can cause chemical burns or toxic inhalation.

\medterm{Nitrosoureas} A group of chemotherapy medicines that can cross the blood-brain barrier and damage DNA; examples include carmustine and lomustine.

\medterm{Nitrous Oxide} An inhaled anesthetic gas with analgesic and anesthetic-sparing properties. It has rapid
onset and offset but relatively low potency, expands closed gas spaces, and may
contribute to diffusion hypoxia or postoperative nausea.

\medterm{Nitrous Oxide For Labor} A 50:50 mixture of nitrous oxide and oxygen (Entonox) inhaled by the mother during labor
for analgesia. Self-administered via a mask or mouthpiece; the woman controls her own
dosing, taking a breath 30 seconds before a contraction.

\medterm{Nizatidine} A medicine that reduces stomach-acid production by blocking histamine signals in stomach cells. It may relieve ulcers or reflux symptoms, although availability varies and treatment should consider kidney function and medicine interactions.

\medterm{NMR} NMR, or nuclear magnetic resonance, measures how atomic nuclei respond to a magnetic field and is used in chemical analysis and imaging.

\medterm{Nocardia Asteroides Infection} An infection caused by Nocardia asteroides, a soil-associated bacterium. It most often affects the lungs, skin, or brain and is more likely when immunity is weakened; treatment requires species identification and prolonged antibiotics.

\medterm{Nocardia Infection} An infection caused by Nocardia bacteria, usually affecting the lungs, skin, or brain, especially in people with weakened immunity.

\medterm{Nocardiosis} An infection caused by Nocardia bacteria found in soil and water. It most often affects the lungs, skin, or brain, especially in people with weakened immunity, and may cause cough, fever, skin sores, confusion, or seizures.

\medterm{Nocebo} A harmful or unpleasant symptom or treatment effect caused or amplified by negative expectations rather than by the treatment’s specific pharmacologic action.

\medterm{Nociception} The nervous system's detection and processing of potentially harmful heat, pressure, or chemicals. It can lead to pain, but pain also depends on how the brain interprets these signals.

\medterm{Nociceptor} A sensory nerve ending that detects potentially damaging pressure, heat, cold, or chemicals. It sends warning signals through nerves to the spinal cord and brain, where they can be experienced as pain.

\medterm{Nociceptors} Sensory nerve endings that detect potentially damaging pressure, heat, cold, or chemicals. They send warning signals through the nervous system, although pain also depends on how the brain interprets those signals.

\medterm{Nocturia} Waking during the night one or more times to urinate and then returning to sleep. It may result from evening fluids, medicines, bladder or prostate problems, diabetes, heart disease, or sleep disorders.

\medterm{Nocturnal} Occurring at night or during sleep. Examples include bedwetting, nighttime cough, breathing pauses during sleep, nighttime low oxygen, or breathlessness that wakes a person. The cause depends on the symptom and requires assessment when persistent or severe.

\medterm{Nocturnal Emission} Involuntary ejaculation during sleep, sometimes called a wet dream. It is a common and normal event during male puberty and adolescence.

\medterm{Nocturnal Enuresis} Repeated involuntary urination during sleep in a child who is old enough to have developed nighttime bladder control. It may run in families; pain, daytime wetting, constipation, or sudden onset needs evaluation.
\synonyms
Nocturnal Incontinence; Sleep Enuresis

\medterm{Nocturnal Enuresis Entry} Alternate name; see \textit{Bed-Wetting}.

\medterm{Nocturnal Hypoventilation} A reduction in minute ventilation occurring specifically during sleep, resulting in nighttime hypercapnia and hypoxemia, commonly encountered in neuromuscular diseases (such as ALS or muscular dystrophies), severe kyphoscoliosis, and obesity hypoventilation syndrome.

\medterm{Nocturnal Incontinence} Alternate name; see \textit{Incontinence}.

\medterm{Nocturnal Lagophthalmos} Incomplete closure of the eyelids during sleep, which may cause dryness, irritation, and exposure injury to the cornea.

\medterm{Nocturnal Polyuria} Production of an unusually large share of the day’s urine during nighttime sleep. It can contribute to repeated nighttime urination and may be related to fluid shifts, sleep disorders, medicines, or other medical conditions.

\medterm{Nodal Rhythm} A heartbeat rhythm that starts near the atrioventricular node instead of the heart’s usual pacemaker. It may occur when the normal pacemaker is slow or blocked and can cause a slow pulse, dizziness, or no symptoms.

\medterm{Node} A localized swelling, structure, or point of connection; examples include lymph nodes and electrical conduction nodes.

\medterm{Node of Ranvier} Periodic, unmyelinated gaps or interruptions occurring at regular intervals along the length of a myelinated nerve axon where high concentrations of voltage-gated sodium channels enable rapid saltatory nerve impulse conduction.
\synonyms
Ranvier Node; Myelin Sheath Gap

\medterm{Nodular} Made up of or containing one or more rounded lumps called nodules. The word may describe a skin lesion, thyroid swelling, lung finding, or other tissue. Its importance depends on the nodule’s size, appearance, location, growth, and associated symptoms.

\medterm{Nodular Fasciitis} A rapidly growing benign self-limited fibroblastic and myofibroblastic proliferation, usually presenting as a painful or tender superficial mass. It can mimic sarcoma clinically and microscopically.

\medterm{Nodular Lymphoid Hyperplasia} Enlargement of small areas of immune tissue, most often in the small intestine. It is usually harmless but may be associated with reduced antibody production or intestinal infection; its significance depends on symptoms and the test findings.

\medterm{Nodular Melanoma} A type of skin cancer that grows downward early and often appears as a rapidly enlarging firm bump. It may be black, brown, pink, or skin-colored, so any quickly changing or bleeding lump needs prompt assessment.

\medterm{Nodule} A small, rounded area of abnormal tissue that may be solid or fluid-filled and may occur in the skin, lung, thyroid, or another organ. Its meaning depends on location, size, appearance, growth, and symptoms; some nodules are harmless, while others need testing or follow-up.

\medterm{Noise Induced Hearing Loss And Its Prevention} Noise-induced hearing loss is permanent injury to inner-ear hair cells from repeated or intense sound exposure, often with tinnitus or difficulty understanding speech. Limiting exposure, increasing distance, taking quiet breaks, and using well-fitted hearing protection reduce risk.

\medterm{Noise-Related Hearing Loss} Alternate name; see \textit{Hearing Loss}.

\medterm{Noma} Noma is a rapidly progressive gangrenous infection of the mouth and face, usually in severely malnourished or immunocompromised children. It destroys oral and facial tissue and is treated with antibiotics, nutritional support, and sometimes surgery.

\medterm{Nomenclature} Nomenclature is the organized system used to name and classify diseases, organisms, chemicals, anatomy, medicines, or other entities.

\medterm{Nomogram} A graphical tool that estimates a quantity or risk from related measurements. The Rumack–Matthew nomogram is used in selected cases of acute acetaminophen ingestion.

\medterm{Nonaccidental Injury} An injury caused deliberately or by neglect rather than by an accident. In children, suspicious patterns may include injuries inconsistent with the history, injuries at different stages of healing, or injuries in protected areas.

\medterm{Non-Accidental Trauma} An injury caused deliberately or through abuse rather than by an accident. Clinicians consider it when the explanation does not fit the injuries, when injuries are repeated or occur at different healing stages, or when a child cannot explain what happened; suspected cases require safeguarding assessment.

\medterm{Non-Adherence} Failure to follow an agreed treatment plan, which may result from access barriers, side effects, beliefs, misunderstanding, or other circumstances.

\medterm{Non-Alcoholic Fatty Liver Disease} A spectrum of liver disorders characterized by hepatic steatosis exceeding 5 percent of hepatocytes in individuals without significant alcohol consumption, ranging from simple steatosis to non-alcoholic steatohepatitis (NASH), cirrhosis, and hepatocellular carcinoma.

\synonyms
NAFLD; Metabolic Dysfunction-Associated Steatotic Liver Disease; MASLD

\medterm{Nonalcoholic Fatty Liver Disease} The former name for metabolic dysfunction-associated steatotic liver disease (MASLD), in which excess fat builds up in the liver together with at least one cardiometabolic risk factor, such as obesity, diabetes, high blood pressure, or abnormal cholesterol. The older name remains common in medical records.

\medterm{Non-Alcoholic Steatohepatitis} An aggressive, progressive form of non-alcoholic fatty liver disease (NAFLD/MASLD) characterized by hepatic steatosis accompanied by lobular necroinflammation and hepatocyte ballooning degeneration, with or without Mallory-Denk bodies and pericellular fibrosis, predisposing to cirrhosis and hepatocellular carcinoma.
\synonyms
NASH; Metabolic Dysfunction-Associated Steatohepatitis (MASH)

\medterm{Non-Allergic Rhinitis} A chronic nasal syndrome characterized by sneezing, rhinorrhea, nasal congestion, and post-nasal drip without clinical or laboratory evidence of IgE-mediated allergic sensitization, frequently triggered by weather changes, cold dry air, chemical odors, or foods (vasomotor rhinitis).

\synonyms
Vasomotor Rhinitis; NAR

\medterm{Nonallergic Rhinitis} Nasal congestion, sneezing, or runny nose caused by irritants or other triggers rather than an allergic immune response. Triggers may include changes in temperature or humidity, smoke, odors, medicines, or hormonal changes.

\synonyms
Vasomotor Rhinitis; Idiopathic Rhinitis

\medterm{Nonallergic Rhinopathy} Nonallergic rhinopathy causes congestion, runny nose, or sneezing without an allergic immune reaction, often triggered by odors, smoke, weather changes, or irritants.

\medterm{Non-Articular Rheumatism} An older term for painful disorders of muscles, tendons, bursae, or other tissues outside the joints.

\medterm{Non-Asthmatic Eosinophilic Bronchitis} A cause of long-lasting cough in which certain white blood cells inflame the airways without the variable narrowing seen in asthma. Breathing tests may be normal, but an inhaled corticosteroid can reduce the cough when this diagnosis is confirmed.

\medterm{Nonbacterial Prostatitis} Chronic pelvic pain syndrome (CP/CPPS) characterized by persistent pelvic, perineal, or genital pain and urinary symptoms for at least 3 months without evidence of active urinary tract bacterial infection on standard cultures.

\synonyms
Chronic Pelvic Pain Syndrome in Men; Abacterial Prostatitis; Category III Prostatitis

\medterm{Nonbacterial Thrombotic Endocarditis} Small, sterile blood clots that form on heart valves without infection. It is often linked to cancer, lupus, or another condition that makes blood clot easily. The clots can break off and cause stroke or blockage elsewhere, so treatment focuses on the cause and preventing further clots.

\synonyms
NBTE; Marantic endocarditis; Libman-Sacks-like nonbacterial endocarditis

\medterm{Non-Celiac Gluten Intolerance} Alternate name; see \textit{Non-Celiac Gluten Sensitivity}.

\medterm{Non-Celiac Gluten Sensitivity} Symptoms attributed to gluten or wheat that improve after removal, despite testing not showing celiac disease or wheat allergy. Diagnosis requires careful medical evaluation because other foods and conditions can cause similar symptoms.

\medterm{Non-Communicable Disease} A chronic, non-transmissible medical condition that is not caused by acute infectious agents and typically progresses slowly over a prolonged duration. Arising from a multifactorial interplay of genetic, physiological, environmental, and behavioral factors, the four primary groups responsible for the vast majority of global mortality and disability are cardiovascular diseases (such as coronary artery disease, stroke, and hypertension), cancers, chronic respiratory diseases (such as COPD and asthma), and diabetes mellitus. Management and public-health prevention focus on early screening, long-term pharmacological control, and modifying shared behavioral risk factors, including tobacco use, physical inactivity, unhealthy diet, and harmful alcohol consumption.

\synonyms
NCD; Non-Communicable Diseases; Noncommunicable Disease; Chronic Non-Communicable Disease

\medterm{Non Communicating Hydrocephalus} A buildup of cerebrospinal fluid in the brain because a blockage prevents it from moving between the brain’s fluid-filled spaces and reaching the area where it is absorbed. It can increase pressure on the brain and cause headache, vomiting, sleepiness, vision problems, or difficulty walking.

\medterm{Noncompetitive Antagonist} A medicine that reduces another medicine or natural chemical's effect by attaching to its receptor or signaling pathway in a way that extra amounts of the activating substance cannot fully overcome.

\medterm{Noncompetitive Inhibition} A form of enzyme inhibition in which a substance binds away from the enzyme’s working site and reduces its activity. Because it does not compete directly with the usual substrate, adding more substrate does not fully overcome the inhibition.

\medterm{Non-Compliant Angioplasty Balloon} A small balloon on a catheter that keeps its shape when inflated at high pressure to widen a narrowed blood vessel. It is used during selected vascular procedures and can rarely cause bleeding, vessel injury, or blockage.
\synonyms
High-Pressure Angioplasty Balloon; Non-Compliant Balloon

\medterm{Non-Contact Tonometry} A quick estimate of pressure inside the eye using a brief puff of air against the clear front surface of the eye. It does not touch the eye and can screen for glaucoma, but unusual results may need confirmation.

\medterm{Non-Convulsive Status Epilepticus} A prolonged state of continuous electrographic seizure activity lasting 30 minutes or longer characterized by altered mental status, confusion, coma, or subtle motor automatisms (such as eyelid fluttering) in the complete absence of generalized tonic-clonic motor convulsions, diagnosed via continuous EEG.

\synonyms
NCSE

\medterm{Nonconvulsive Status Epilepticus} Ongoing seizure activity without prominent rhythmic limb jerking, often presenting as prolonged confusion, impaired awareness, or subtle twitching. An electroencephalogram can confirm it, and treatment is urgent.

\medterm{Nondisjunction} Failure of chromosome copies to separate properly when a cell divides. This can produce cells with too many or too few chromosomes and may cause conditions such as Down syndrome, Turner syndrome, or some miscarriages.

\medterm{Non-Fatal Drowning} Survival after breathing is impaired by submersion or immersion in liquid. A person may still develop lung injury, brain injury, or delayed breathing problems, even after apparently recovering.

\medterm{Non-Functioning Pituitary Adenoma} A usually noncancerous pituitary tumor that does not release enough active hormones to cause a hormone syndrome. As it enlarges, it may cause headache, reduced vision, or loss of normal pituitary function.

\medterm{Non-Hodgkin Lymphoma} A group of cancers that begin in immune cells called lymphocytes. Some grow slowly and others quickly; they may cause swollen lymph nodes, fever, night sweats, weight loss, or symptoms in an affected organ.

\synonyms
NHL; B-Cell or T-Cell Lymphoma

\medterm{Non-Hodgkin Lymphoma In Children} A group of childhood cancers arising from immune cells called lymphocytes. Some grow quickly and can cause swollen nodes, fever, abdominal symptoms, or breathing problems; the exact type and spread determine treatment.

\medterm{Non-Hodgkin's Lymphoma} Alternate spelling; see \textit{Non-Hodgkin lymphoma}.

\medterm{Non-Hodgkins Lymphomas} Alternate spelling and plural form; see \textit{Non-Hodgkin lymphoma}.

\medterm{Non-Infectious Cystitis} Inflammation of the bladder without infection by bacteria or another microorganism. It may cause pain, urgency, and frequent urination after irritants, radiation, medicines, chemical exposure, or other bladder conditions.

\medterm{Noninfectious Cystitis} A heterogeneous group of bladder inflammatory conditions occurring in the absence of active bacterial infection, encompassing radiation cystitis, chemotherapy-induced cystitis (e.g., cyclophosphamide), chemical cystitis, and interstitial cystitis.

\synonyms
Abacterial Cystitis; Sterile Cystitis; Chemical/Radiation Cystitis

\medterm{Noninvasive} Not involving a surgical cut, a procedure entering the body, or a device placed inside it. Examples include measuring blood pressure with a cuff, external imaging, and breathing support through a mask. Risks depend on the method used.

\medterm{Non-Invasive Positive Pressure Ventilation} The delivery of mechanical ventilatory support without an invasive endotracheal tube or tracheostomy, utilizing a tight-fitting facial or nasal mask to provide continuous positive airway pressure (CPAP) or bilevel positive airway pressure (BiPAP) in acute hypercapnic respiratory failure or cardiogenic pulmonary edema.

\synonyms
NIPPV; NPPV

\medterm{Non-Invasive Prenatal Testing} A blood test during pregnancy that estimates the chance of selected chromosome conditions, especially trisomy 21, 18, and 13, by analyzing small DNA fragments from the placenta. It is screening, not diagnosis; an abnormal result needs confirmatory testing.

\medterm{Non-Invasive Procedure} An examination or treatment that does not cut or pierce the skin or place an instrument inside the body. Examples include blood-pressure checks, ECGs, ultrasound, MRI, X-rays, and physical therapy. These procedures are generally safer than invasive procedures, but each can still have limits or risks.

\synonyms
Noninvasive Procedure; Non-Invasive Medical Procedure; Noninvasive Test; Non-Invasive Technique

\medterm{Non-Invasive Ventilation} Breathing support delivered through a face or nose mask instead of a tube placed in the windpipe. Continuous or two-level air pressure can improve breathing in selected people with respiratory failure, sleep apnea, or fluid in the lungs.

\medterm{Nonischemic Priapism} Alternate name; see \textit{Priapism}.

\medterm{Nonketotic Hyperglycemic Hyperosmolar Syndrome} A life-threatening diabetic emergency in which very high blood sugar causes profound dehydration and concentrated blood, usually without major ketone buildup. It can cause extreme thirst, weakness, confusion, seizures, coma, and organ failure.

\medterm{Non-Ketotic Hyperosmolar State} A life-threatening acute metabolic emergency occurring predominantly in type 2 diabetes mellitus, characterized by extreme hyperglycemia (glucose greater than 600 mg/dL), profound hyperosmolality (plasma effective osmolality greater than 320 mOsm/kg), and severe dehydration in the absence of significant ketoacidosis.
\synonyms
Hyperosmolar Hyperglycemic State; HHS

\medterm{Nonmelanoma Skin Cancer} A group of skin cancers that arise from the cells of the epidermis, excluding melanoma.
The two most common types are basal cell carcinoma (BCC) and squamous cell carcinoma
(SCC), which together account for the vast majority of all skin cancers diagnosed
globally.

\medterm{Non-Muscle-Invasive Bladder Cancer} Bladder cancer limited to the inner lining or connective tissue beneath it, without invasion of the bladder muscle. It can recur and may progress to muscle-invasive disease.

\medterm{Non-Obstructive Azoospermia} Complete absence of sperm in semen because the testicles make too few or no sperm, rather than because a tube is blocked. Causes include genetic conditions, undescended testes, infections, cancer treatment, or an unknown problem. Testing may include hormone tests and genetic studies; some people can use surgically retrieved sperm for assisted reproduction.

\synonyms
NOA; Secretory Azoospermia; Testicular Failure Azoospermia

\medterm{Nonpathogenic} Nonpathogenic means not ordinarily capable of causing disease in a given host or setting; an organism can still cause opportunistic infection when immunity or barriers are impaired.

\medterm{Non-Pitting Edema} Swelling of subcutaneous tissue that does not leave a persistent indentation when firm digital pressure is applied, usually due to local accumulation of protein-rich fluid, mucopolysaccharides, or chronic cellular infiltration. Common causes include lymphedema, pretibial myxedema in Graves disease, and severe lipedema.

\medterm{Non-Polio Enterovirus Infection} Alternate name; see \textit{Enterovirus}.

\medterm{Non-Proliferative Diabetic Retinopathy} An early form of diabetes-related damage to the retina, the light-sensitive tissue at the back of the eye. Small blood vessels may leak or swell, causing blurred vision, but abnormal new vessels have not formed.

\medterm{Nonproliferative Retinopathy} Retinal disease in which small blood vessels leak or become damaged without abnormal new vessels growing. Diabetes is a common cause; progression can impair vision, especially when the macula is affected.
\synonyms
Background Retinopathy; NPDR

\medterm{Non-Radiographic Axial Spondyloarthritis} Inflammatory arthritis affecting the spine or joints between the spine and pelvis without definite changes on a plain X-ray. It may cause persistent back pain and morning stiffness, and may occur with eye inflammation, psoriasis, or inflammatory bowel disease.

\medterm{Non-Rapid Eye Movement Sleep} Sleep that includes lighter and deeper stages without rapid eye movements. Deep stages support physical restoration, while lighter stages help the body transition into and out of sleep.

\medterm{Non-Rebreather Mask} A face mask with a reservoir bag and one-way valves that can deliver a high concentration of supplemental oxygen when fitted and used correctly.

\medterm{Non-Rebreather Mask} A high-flow oxygen delivery device featuring an attached reservoir bag and one-way directional valves that prevent entrainment of room air and rebreathing of exhaled gas, capable of delivering inspired oxygen fractions (FiO2) between 80 and 95 percent at flow rates of 10 to 15 L/min.

\synonyms
NRB Mask

\medterm{Non-REM Sleep} Sleep stages without rapid eye movements, ranging from light sleep to deep slow-wave sleep. These stages support physical restoration, memory processing, and the normal transition between waking and sleep.
\synonyms
NREM sleep

\medterm{Non-Responder} A person or disease that shows little or no expected response to a treatment, vaccine, or test stimulus.

\medterm{Nonselective Beta-Blockers} Medicines that block two main types of adrenaline receptor. They slow the heart and lower blood pressure but can also narrow the airways, so people with asthma or certain lung diseases may need a different medicine.

\medterm{Nonseminomatous Germ Cell Tumor} A cancer arising from reproductive cells, most often in the testicle. It may grow and spread quickly; blood tests, scans, and tissue examination guide staging and treatment.

\synonyms
NSGCT; Nonseminoma

\medterm{Non-Seminomatous Germ Cell Tumor} A major clinicopathological category of testicular and extragonadal germ cell neoplasms comprising embryonal carcinoma, yolk sac tumor, choriocarcinoma, and teratoma (or mixed variants). Characterized by aggressive clinical behavior with early hematogenous and retroperitoneal lymphatic dissemination, producing characteristic serum tumor markers (alpha-fetoprotein and human chorionic gonadotropin); treated with radical orchiectomy followed by platinum-based combination chemotherapy.

\synonyms
NSGCT; Nonseminoma

\medterm{Nonsense Variant} A change in DNA that tells a cell to stop making a protein too early. The resulting shortened or missing protein may cause disease, but the effect depends on the gene, the location of the change, and how much normal protein remains.

\medterm{Non-Small Cell Lung Cancer} The most common group of lung cancers, including adenocarcinoma and squamous-cell cancer. It may start in different parts of the lung and can spread; treatment depends on its stage, cell type, gene changes, and the person's health.

\medterm{Non-Small Cell Lung Carcinoma} A broad histological grouping of primary epithelial malignancies of the lung comprising approximately 85 percent of all lung cancers, including adenocarcinoma, squamous cell carcinoma, and large cell carcinoma, distinguished by clinical behavior and targeted therapy sensitivity from small cell lung cancer.

\synonyms
NSCLC

\medterm{Nonspecific Interstitial Pneumonia} A pattern of lung inflammation or scarring that develops relatively evenly over time. It is often linked to connective-tissue disease, medicines, or an unknown cause and can cause breathlessness and a dry cough.

\medterm{Nonspecific Interstitial Pneumonitis} A pattern of lung inflammation and scarring that develops fairly evenly through the lung tissue. It may cause a dry cough and breathlessness and is often linked to autoimmune disease, medicines, or no identified cause.

\medterm{Non-Specific Urethritis} Inflammation of the tube that carries urine out of the body when a specific cause has not yet been identified. It may cause burning urination or discharge and is often linked to sexually transmitted infections, so testing and treatment of partners may be needed.

\medterm{Non-ST Elevation Myocardial Infarction} An acute coronary syndrome characterized by ischemic myocardial necrosis with detectable elevated cardiac biomarkers (cardiac troponin I or T) in the absence of acute persistent ST-segment elevation on ECG (presenting instead with ST-segment depression, T-wave inversion, or normal tracing).
\synonyms
NSTEMI; Non-ST-Segment Elevation Infarct

\medterm{Non-ST-Elevation Myocardial Infarction} A heart attack in which blood tests show heart-muscle injury but the electrocardiogram does not show persistent ST-segment elevation. It is usually caused by reduced or blocked coronary blood flow and may require urgent assessment and treatment.

\medterm{Nonsteroidal Anti-Inflammatory Drugs} A group of medicines that reduce pain, fever, and inflammation by lowering the body’s production of prostaglandins. They can cause stomach bleeding, kidney injury, fluid retention, heart problems, or allergic reactions, especially at high doses or with long use.

\medterm{Nonsteroidal Anti-Inflammatory Drugs And Ulcers} Nonsteroidal anti-inflammatory drugs can reduce stomach-protective prostaglandins and cause gastritis, ulcers, and gastrointestinal bleeding. Risk rises with age, prior ulcer, anticoagulants, steroids, kidney disease, and high doses; safer alternatives or acid protection may be needed.

\synonyms
NSAID-Induced Peptic Ulcer; NSAID Gastropathy

\medterm{Non-Stress Test} A test during pregnancy that monitors the baby's heart rate while the baby moves. A healthy response usually includes temporary heart-rate increases; an abnormal result may reflect sleep, medicines, or a problem requiring further assessment.

\medterm{Non-ST-Segment Elevation Myocardial Infarction} An acute coronary syndrome characterized by evidence of myocardial necrosis (elevated cardiac troponin biomarkers) in the absence of acute ST-segment elevation on electrocardiogram, typically reflecting subendocardial ischemia caused by a non-occlusive mural thrombus.

\synonyms
NSTEMI

\medterm{Nonsyndromic} Nonsyndromic means occurring without the additional characteristic features of a recognized syndrome; it may still be genetic, congenital, or clinically significant.

\medterm{Nontropical Sprue} Alternate name; see \textit{Celiac Disease}.

\medterm{Nontuberculous Mycobacteria} Mycobacteria other than the tuberculosis complex and M. leprae. They are found in soil and water and can cause lung, skin, lymph-node, or disseminated infections, especially in people with underlying illness.

\medterm{Nontuberculous Mycobacterial Infection} Infection caused by Mycobacterium species other than M. tuberculosis, often affecting the lungs but sometimes the lymph nodes, skin, or other organs. Disease is more likely in people with structural lung disease or impaired immunity.

\medterm{Nontuberculous Mycobacterial Lung Disease} Chronic lung infection caused by environmental Mycobacterium species other than M. tuberculosis.
It may produce cough, fatigue, weight loss, breathlessness, or coughing blood, especially in people with structural
lung disease. Diagnosis requires compatible imaging and repeated microbiologic testing; treatment is prolonged and individualized.

\medterm{Nonulcer Dyspepsia} Alternate name; see \textit{Functional Dyspepsia}.

\medterm{Nonulcer Stomach Pain} Alternate name; see \textit{Functional Dyspepsia}.

\medterm{Nonunion And Malunion} Nonunion means a broken bone fails to heal after an expected period, leaving pain or instability. Malunion means it heals in a poor position, causing deformity, altered movement, or strain on nearby joints. Both may require further treatment.

\medterm{Nonunion Fracture} A broken bone that has failed to heal after the expected period and shows no further progress toward joining. It may cause ongoing pain, movement at the fracture, weakness, or deformity and may require surgery or other treatment.

\medterm{Noonan Syndrome} An inherited condition that may cause distinctive facial features, short stature, heart defects, delayed development, easy bruising, or swelling from lymph problems. Its features vary, and regular heart, growth, hearing, and developmental checks are important.

\medterm{Noonan Syndrome With Multiple Lentigines} A genetic disorder also called LEOPARD syndrome, characterized by multiple lentigines, distinctive facial features, heart disease, short stature, and sometimes deafness or genital abnormalities.

\medterm{Noradrenaline} A chemical messenger acting as both a hormone and neurotransmitter. It helps regulate alertness, attention, stress responses, heart activity, and blood-vessel tightening, which can raise blood pressure.

\synonyms
Norepinephrine

\medterm{No-Reflow Phenomenon} Failure of blood to reach small vessels and damaged tissue even after a larger blocked artery has been reopened. It can occur after a heart attack or stroke and may limit recovery despite successful treatment of the main blockage.

\medterm{Norepinephrine} A hormone and nerve messenger involved in alertness, stress responses, heart function, and blood-vessel tightening. It can increase blood pressure and is also used in hospitals to support blood pressure during some forms of shock.
\synonyms
Noradrenaline

\medterm{Norethisterone} A manufactured progesterone-like medicine used for contraception and selected menstrual problems or endometriosis. It can change bleeding patterns and may cause nausea, headache, mood changes, breast tenderness, or other hormone-related effects.

\medterm{Norgestrel} A synthetic progestogen used in some hormonal contraceptives. It helps prevent pregnancy by suppressing ovulation and changing cervical mucus and the uterine lining.

\medterm{Norketamine} Norketamine is an active metabolite of ketamine with NMDA-receptor antagonist activity; it contributes to ketamine's analgesic and psychoactive effects and may persist after the parent drug is cleared.

\medterm{Normal} Within an expected range or consistent with usual anatomy, physiology, or function for a specified population.

\medterm{Normal Flora} The bacteria, fungi, and other microorganisms that normally live on the skin and in the mouth, intestine, and urinary or genital tract without causing illness. They can help block harmful germs, support immunity, and contribute to digestion and vitamin production.

\medterm{Normal Growth And Development} Normal growth and development describes expected physical growth, movement, language, learning, and social progress for a person’s age and stage of life.

\medterm{Normalization} A mathematical adjustment that makes measurements comparable by correcting for known differences in scale, background, equipment response, or other measurement conditions.

\medterm{Normal Labour} Spontaneous labour at term with regular uterine contractions, progressive cervical dilatation, and delivery through the birth canal without major complication.

\medterm{Normal Pressure Hydrocephalus} A buildup of cerebrospinal fluid that enlarges spaces inside the brain and may cause a slow, unsteady walk, reduced thinking, and loss of bladder control. Some people improve with specialist evaluation and drainage treatment.

\medterm{Normal-Pressure Hydrocephalus} A disorder in which cerebrospinal fluid accumulates and causes gait difficulty, cognitive decline, and urinary problems despite near-normal measured pressure.

\medterm{Normal Range} A normal range is the interval of results commonly found in a reference population; a result outside it does not by itself prove disease.

\medterm{Normal Saline} A sterile approximately 0.9\% sodium-chloride solution used for fluid replacement, irrigation, and dilution of medicines.

\medterm{Normal-Tension Glaucoma} A type of open-angle glaucoma in which the optic nerve is damaged and vision may be lost even though pressure inside the eye remains within the normal range.

\medterm{Normoblast} A nucleated immature red-cell precursor in bone marrow that normally loses its nucleus before entering circulation.

\medterm{Normocalcemic Hyperparathyroidism} Condition of elevated parathyroid hormone with persistently normal serum calcium and
normal vitamin D levels. Exclusion of secondary causes of elevated PTH including vitamin
D deficiency, renal insufficiency, and medications is required. Represents an early or
mild form of primary hyperparathyroidism. May progress to overt hypercalcemia.

\medterm{Normocytic Anaemia} Alternate British spelling; see \textit{Normocytic Anemia}.

\medterm{Normocytic Anemia} A form of anemia in which circulating red blood cells are of normal average size (mean corpuscular volume [MCV] between $80$ and $100\,\text{fL}$), but total red blood cell count and hemoglobin concentration are decreased. Diagnostic evaluation relies on the reticulocyte count to differentiate between: (1) hyperproliferative causes with elevated reticulocytes, such as acute blood loss (hemorrhage) or premature red cell destruction (hemolytic anemia), and (2) hypoproliferative causes with low reticulocytes, including anemia of chronic disease, chronic kidney disease (erythropoietin deficiency), aplastic anemia, and bone marrow infiltration. Symptoms include fatigue, weakness, dizziness, and pallor.
\synonyms{Normocytic Normochromic Anemia; Normal-MCV Anemia}

\medterm{Normotensive} Having blood pressure within the expected range for the individual and measurement setting. Having blood pressure within the expected range for a person, based on the measurement setting.

\medterm{Norovirus} A highly contagious virus that causes acute gastroenteritis, or inflammation of the stomach and intestines. Symptoms usually begin 12--48 hours after exposure and include sudden nausea, vomiting, watery diarrhea, stomach cramps, and sometimes fever, headache, or body aches. It spreads mainly when tiny particles of stool or vomit enter the mouth, through infected people, contaminated food or water, surfaces, or particles released during vomiting. Most illness improves within 1--3 days, but severe fluid loss can cause dehydration, especially in infants, older adults, and people with other illnesses.

\medterm{Norovirus Infection} Acute gastroenteritis caused by norovirus. Symptoms usually begin 12--48 hours after exposure and include nausea, vomiting, watery diarrhea, abdominal cramps, and sometimes fever, headache, or body aches. Infection spreads easily through infected people, contaminated food or water, surfaces, or particles released during vomiting. Most cases improve within 1--3 days, but fluid loss may cause dehydration, especially in infants, older adults, and people with other illnesses.

\medterm{Norrie Disease} A rare inherited eye disorder in which the retina develops abnormally, usually causing severe vision loss or blindness from birth or early infancy. Some people also develop hearing loss, developmental differences, or problems in the eye’s blood vessels.

\medterm{Nortriptyline} A medicine used to treat depression and sometimes long-lasting nerve pain. It changes the levels of certain chemical messengers in the nervous system and may cause dry mouth, constipation, sleepiness, dizziness, heart-rhythm problems, or overdose toxicity.

\medterm{Norwegian Scabies} A severe form of scabies in which thick crusts contain very large numbers of mites. It is extremely contagious and occurs especially in people with weakened immunity or difficulty sensing or scratching the skin, requiring prompt treatment and contact precautions.

\medterm{Norwood Procedure} The first operation for many newborns with hypoplastic left heart syndrome and related heart defects. It rebuilds the main artery so the right ventricle can send blood to the body, creates a wide opening between the upper heart chambers, and provides controlled blood flow to the lungs through a shunt or tube. It is usually performed within the first days of life and is followed by further staged operations.

\synonyms
Stage 1 Norwood Reconstruction; Norwood Operation

\medterm{Nose} The external structure and internal passages used for breathing and smell. The nose warms, moistens, and filters air and contains bone, cartilage, skin, mucus, blood vessels, and smell-sensing tissue.

\medterm{Nosebleed} Singular form; see \textit{Nosebleeds}.

\medterm{Nosebleeds} Bleeding from blood vessels inside the nose, commonly caused by dryness, irritation, injury, infection, allergies, medicines, or high blood pressure. Episodes may be brief or recurrent and can vary in amount.

\medterm{Nose Fracture} A break in one or more nasal bones caused by trauma, potentially causing pain, swelling, deformity, bleeding, or nasal obstruction.

\medterm{Nosocomial} Acquired during healthcare, especially in a hospital, rather than being present or incubating when the person entered care. A nosocomial infection may be caused by bacteria, viruses, fungi, or other organisms and can involve drug-resistant organisms.

\medterm{Nosocomial Infection} A hospital-acquired infection developing during or after hospitalization that was not
present or incubating at admission, classically defined by onset beyond 48 hours. Common
sites are bloodstream, urinary tract, lungs, and surgical wounds; causative organisms
are frequently multidrug-resistant, prevented by hand hygiene and stewardship.

\medterm{Nosocomial Norovirus Infection} Healthcare-associated norovirus gastroenteritis characterized by acute, explosive onset of watery diarrhea, severe non-bloody emesis, and abdominal cramps. Outbreak management necessitates strict patient isolation, contact precautions, rigorous hand hygiene with soap and water (as alcohol-based rubs are less effective against non-enveloped virions), and hypochlorite-based surface disinfection.

\synonyms
Hospital-Acquired Norovirus; Healthcare Norovirus Outbreak; Hospital Winter Vomiting Disease

\medterm{Nosocomial Vancomycin-Resistant Enterococci} Enterococcal bacteria that resist vancomycin and can spread in healthcare settings. They may cause urinary, wound, bloodstream, or other infections, particularly in people who are seriously ill or have devices such as catheters; treatment requires susceptibility testing.

\medterm{Nostril} One of the two openings at the front of the nose. Each nostril leads into the nasal passage, where hairs and mucus help trap particles before air reaches the deeper airways.

\medterm{Notencephalocele} A rare congenital neural-tube defect involving protrusion of brain tissue or meninges through a skull defect; terminology varies by source.

\medterm{Notifiable Disease} An infection or other condition that healthcare professionals or laboratories must report to public-health authorities under local law. Reporting helps detect outbreaks, arrange control measures, protect contacts, and plan services; the exact list and reporting rules vary by location.

\medterm{Notochord} A flexible embryonic rod that guides axial development and largely disappears as the vertebral column forms.

\medterm{NOWS} Alternate name; see \textit{Neonatal Opioid Withdrawal Syndrome}.

\medterm{Noxious} Harmful, poisonous, or capable of causing tissue injury or pain. Noxious stimuli activate protective sensory pathways such as pain and withdrawal reflexes.

\medterm{NREM Sleep} The part of sleep without rapid eye movements, ranging from light sleep to deep slow-wave sleep. It alternates with REM sleep in repeated cycles and supports physical restoration, learning, and memory.

\medterm{NSAID} A nonsteroidal anti-inflammatory drug used to reduce pain, fever, and inflammation. Examples include ibuprofen and naproxen; these medicines can irritate the stomach, cause bleeding or kidney problems, and increase heart risks in some people.

\synonyms
Nonsteroidal anti-inflammatory drug

\medterm{NSAID-Induced Ulcers} Alternate name; see \textit{Peptic Ulcer Disease}.

\medterm{NSAIDs} A class of analgesic, antipyretic, and anti-inflammatory drugs that inhibit
cyclooxygenase (COX-1 and/or COX-2), reducing prostaglandin and thromboxane synthesis.

\synonyms
Non-Steroidal Anti-Inflammatories; COX Inhibitors

\medterm{NSTEMI} A heart attack in which a blocked or narrowed coronary artery injures heart muscle but does not produce the typical persistent ST-segment elevation on an ECG. It may cause chest pressure, breathlessness, sweating, nausea, or faintness.

\synonyms
Non-ST-Elevation Myocardial Infarction

\medterm{NTD} NTD commonly means neural-tube defect, a congenital brain or spinal-cord malformation associated with inadequate folate and other factors.

\medterm{Nuchal} Relating to the back of the neck. Nuchal rigidity means painful or limited neck bending and can occur with meningitis or bleeding around the brain. Nuchal translucency is a first-trimester ultrasound measurement used to estimate certain pregnancy risks.

\medterm{Nuchal Cord} Wrapping of the umbilical cord around the fetal neck, occurring once or more times. It is common and often harmless, but monitoring during labor is needed if fetal heart-rate abnormalities occur.

\medterm{Nuchal Rigidity} Stiffness or pain that makes it difficult to bend the neck forward. It may occur with meningitis, bleeding around the brain, severe muscle problems, or other illnesses; sudden stiffness with fever or severe headache may signal a serious condition.

\medterm{Nuchal Translucency Screening} An ultrasound measurement of the clear space at the back of a fetus’s neck, usually performed early in pregnancy. A larger-than-expected measurement can indicate increased risk of chromosome or heart problems but is not a diagnosis.

\medterm{Nuchal Translucency Test} An early-pregnancy ultrasound measurement of the clear space at the back of the fetus’s neck. It estimates the chance of certain chromosome or structural problems but cannot diagnose them; abnormal results need follow-up testing.

\medterm{Nuclear} Relating to the nucleus of a cell, the center of an atom, or medical tests and treatments that use radioactivity. The intended meaning depends on the surrounding medical context.

\medterm{Nuclear Factor Kappa B} A ubiquitous inducible master transcription factor complex that coordinates immune and inflammatory responses, cell survival, and cytokine transcription; normally held inactive in the cytoplasm by I-kappa-B inhibitor proteins until activated by inflammatory cytokines, bacterial LPS, or viral infection.

\synonyms
NF-kappaB; NF-kB

\medterm{Nuclear Medicine} A medical specialty that uses small amounts of radioactive substances to diagnose or treat disease. Scans show body function or chemical activity, while treatments deliver radiation to selected tissues.

\medterm{Nuclear Medicine Imaging} Imaging that uses a radioactive tracer and a specialized detector to show organ function, blood flow, chemical activity, or disease distribution. Common examples include PET, SPECT, bone scans, thyroid scans, and radionuclide cardiac imaging.

\medterm{Nuclear Pore Complex} A large, multiprotein cylindrical channel spanning the double-membrane nuclear envelope, composed of multiple nucleoporin proteins that regulate and mediate the selective bidirectional transport of macromolecules (RNA, ribonucleoproteins, proteins) between the nucleus and the cytoplasm.

\synonyms
NPC

\medterm{Nuclear Sclerosis} A common age-related cataract in which the center of the eye’s lens gradually becomes yellow and hard. It can slowly blur vision, reduce contrast, or alter color perception and is treated when it interferes with daily activities.

\medterm{Nuclear Stress Test} A test that compares blood flow to heart muscle at rest and during exercise or medicine-induced stress using a radioactive tracer. It can identify areas of reduced perfusion that suggest coronary artery disease.

\synonyms
Nuclear medicine tests

\medterm{Nuclear Ventriculography} A heart scan that follows a small radioactive tracer through the bloodstream to measure how much blood the heart pumps with each beat and whether the heart walls move normally.

\medterm{Nucleic Acid} DNA or RNA, the molecules that store, copy, and use genetic instructions. DNA usually stores the long-term information, while RNA helps read those instructions and make proteins or perform other jobs in cells.

\medterm{Nucleic Acid Hybridization} A laboratory method in which matching DNA or RNA strands bind together. This allows a test to find, identify, or measure a particular genetic sequence from a person, tumor, or infectious organism.

\medterm{Nucleolus} A dense region within the cell nucleus where ribosomal RNA is produced and ribosome subunits begin assembly.

\medterm{Nucleoside} A small molecule made of a nitrogen-containing base joined to a sugar. Adding one or more phosphate groups makes a nucleotide, the building block used to make DNA or RNA and to carry energy in cells.

\medterm{Nucleoside Analogs} Medicines that resemble the building blocks used to make DNA or RNA. Some viruses use them by mistake, which interrupts viral copying; the medicine chosen depends on the virus, illness, kidney function, and other factors.

\medterm{Nucleotide} A small molecule made of a sugar, a nitrogen-containing base, and one or more phosphate groups. Nucleotides build DNA and RNA and also carry energy used by cells.

\medterm{Nucleus} The membrane-bound cell organelle containing most genetic material, or a defined cluster of neurons in the central nervous system.

\medterm{Nucleus Accumbens} A group of nerve cells deep in the brain that helps process reward, motivation, learning, and reinforcement. Its activity is involved in normal behavior and in disorders such as addiction and some psychiatric illnesses.

\medterm{Nucleus Ambiguus} A longitudinal column of branchial motor lower motor neurons located within the medullary reticular formation that provides motor innervation to the skeletal striated muscles of the soft palate, pharynx, and larynx via cranial nerves IX (glossopharyngeal), X (vagus), and cranial XI.

\medterm{Nucleus Pulposus} The soft, gel-like center of a spinal disk, surrounded by a tough outer ring. It helps spread pressure between vertebrae, but a tear may allow it to bulge or press on nerves.

\medterm{Nulligravida} An obstetric term for a person who has never been pregnant. It is different from nullipara, which means never having given birth beyond the stage of pregnancy used in a particular medical record system.

\medterm{Nullipara} A person who has never carried a pregnancy beyond the stage of fetal viability; the adjective is nulliparous.

\medterm{Null Mutation} A genetic change that prevents a gene from producing a functional protein or RNA. The effect depends on the gene, the other copy, and whether the change occurs in body cells or reproductive cells.

\medterm{Numb} Numb means lacking normal sensation or emotional responsiveness; physical numbness can result from nerve, brain, circulation, metabolic, or medication-related problems.

\medterm{Number Needed To Harm} The estimated number of people who must receive a treatment, rather than the comparison treatment, for one extra person to experience a specified harmful outcome during a stated period. A smaller number means the harm is more frequent.

\medterm{Number Needed To Treat} The number of people who must receive an intervention for one additional person to
experience the specified beneficial outcome compared with the control over a defined
time. It is the reciprocal of absolute risk reduction.

\medterm{Numbness} Numbness is reduced or absent sensation, often from nerve compression, neuropathy, stroke, spinal disease, circulation problems, or anesthetic effects.

\medterm{Numbness And Tingling} Unusual reduced sensation or prickling, often caused by nerve compression, neuropathy, spinal disease, vitamin deficiency, circulation problems, or other neurologic disorders. Sudden one-sided symptoms, weakness, speech trouble, or facial drooping require emergency assessment.

\synonyms
Paresthesia; Pins and Needles

\medterm{Numbness In Hands} Reduced or absent sensation in one or both hands, commonly caused by nerve compression, neuropathy, injury, circulation problems, or systemic disease.

\medterm{Nummular} Coin-shaped or round. The term is commonly used for nummular eczema, which causes itchy, coin-shaped patches of inflamed skin.

\medterm{Nummular Dermatitis} An inflammatory skin condition causing itchy, coin-shaped patches, often on the arms or legs. Dry skin, irritants, and skin injury may contribute; treatment commonly includes moisturizers and prescribed anti-inflammatory medicines.

\medterm{Nummular Eczema} An inflammatory skin condition causing intensely itchy, coin-shaped red patches, sometimes with blisters, crusting, or scaling, usually on the arms or legs. Dry skin, irritation, or injury can trigger it and similar rashes may need examination.

\medterm{Nursemaid's Elbow} Nursemaid's elbow is subluxation of the radial head after sudden traction on a young child's extended or pronated arm; the child avoids using the arm, and gentle clinician reduction usually restores movement.

\medterm{Nurse Practitioner} A nurse practitioner is an advanced-practice nurse who can assess patients, diagnose and treat many conditions, prescribe medicines where authorized, and provide preventive care.

\medterm{Nursery} A place or service for the care of infants, including specialized newborn care in a hospital.

\medterm{Nursing} The professional practice of caring for people, promoting health, preventing illness, and supporting recovery or comfort.

\medterm{Nursing Home} A residential facility that provides ongoing nursing care, assistance with daily activities, and medical or rehabilitation services for
people who cannot live independently.

\medterm{Nut Allergies} Nut allergy is an immune-mediated reaction to tree nuts or peanuts that can cause hives, gastrointestinal or respiratory symptoms, or life-threatening anaphylaxis.

\medterm{Nut-Allergy} An immune-mediated reaction to one or more nuts, potentially causing urticaria, anaphylaxis, or gastrointestinal and respiratory symptoms.

\medterm{Nutation} A nodding or rocking movement; in anatomy it commonly describes movement of the sacrum relative to the pelvis.

\medterm{Nutcracker Esophagus} An older name for a swallowing disorder in which the esophagus contracts with unusually high pressure. It can cause intermittent chest pain or difficulty swallowing; current testing and clinical assessment distinguish it from reflux and other causes.

\medterm{Nutcracker Syndrome} Compression of the vein draining the left kidney, usually between two nearby arteries. It may cause blood in the urine, pelvic or flank pain, or enlarged veins; mild cases may be monitored, while persistent severe symptoms may need treatment.

\medterm{Nutrient} A substance the body needs for energy, growth, repair, or normal function. Nutrients include protein, fats, carbohydrates, vitamins, minerals, and water; too little or too much of some nutrients can cause illness.

\medterm{Nutrition} How the body takes in, digests, absorbs, and uses food and its nutrients for health.

\medterm{Nutritional Anemia} Anemia caused by inadequate intake, absorption, or use of nutrients needed for red-cell production, especially iron, vitamin B12, or folate.

\medterm{Nutritional Assessment} A structured evaluation of food intake, body measurements, laboratory results, symptoms, and health conditions. It identifies nutritional problems and guides a safe, individualized eating plan or nutrition support.

\medterm{Nutritional Deficiency} A lack of a vitamin, mineral, protein, energy, or other nutrient needed for normal body function. It may result from inadequate intake, poor absorption, increased needs, or disease and can impair growth, blood formation, immunity, nerves, or other organs.

\medterm{Nutritional Dermatoses} Skin changes caused by too little or too much of a nutrient. For example, vitamin C deficiency can cause easy bruising and bleeding gums, while niacin deficiency can cause a sun-sensitive rash, diarrhea, and confusion.

\medterm{Nutritional Disorders} Health problems caused by too little, too much, or poor use of nutrients. Examples include undernutrition, obesity, vitamin or mineral shortages, and problems caused by the body being unable to digest or absorb food properly.

\medterm{Nutritional Screening} A brief check for risk of malnutrition using weight change, food intake, illness, and other factors. A concerning result leads to a fuller assessment and an individualized nutrition plan.

\medterm{Nutrition-Focused Physical Examination} A structured examination for physical findings associated with malnutrition, including
loss of subcutaneous fat, muscle wasting, edema, ascites, reduced functional status, and
micronutrient-related signs. It complements dietary history, anthropometry, laboratory
data, and disease assessment.

\medterm{Nutrition In Pregnancy} Eating enough varied food and obtaining appropriate vitamins and minerals to support the pregnant person and growing fetus. Important nutrients include folate, iron, calcium, vitamin D, iodine, and protein; individual needs and supplements should be reviewed with a prenatal clinician.

\medterm{Nutrition Risk Screening 2002} A hospital screening questionnaire that considers recent food intake, weight change, illness severity, and age to identify adults at risk of malnutrition. A high score prompts a fuller assessment; it is not itself a diagnosis.

\medterm{Nutrition Support} Use of oral nutrition supplements, tube feeding, or intravenous nutrition when a person cannot meet nutritional needs by ordinary eating. The method depends on the digestive tract, illness, nutrition needs, and risks of treatment.

\medterm{Nutritive} Providing nourishment that the body can use for energy, growth, repair, or normal function. Nutritive substances include protein, fat, carbohydrates, vitamins, minerals, and water.

\medterm{Nux Vomica} The seed of Strychnos nux-vomica, containing strychnine and brucine; ingestion can cause severe convulsant poisoning.

\medterm{Nyctalopia} Impaired vision in dim light or darkness, also called night blindness.

\medterm{NYHA Functional Classification} A four-class system describing how heart failure symptoms limit activity: class I causes no limitation, class II limits ordinary activity, class III limits less-than-ordinary activity, and class IV causes symptoms even at rest.

\medterm{Nystagmus} Involuntary rhythmic movement of the eyes. It may occur with inner-ear disease, neurologic disorders, reduced vision, medicines, or normal early development; new nystagmus with dizziness, weakness, or vision change may signal an underlying disorder.

\medterm{Nystagmus Retractorius} An unusual eye movement in which the eyes pull backward and may move toward each other when a person tries to look upward. It is a characteristic sign of a problem in the upper part of the brain, but other findings are needed to identify the cause.
\synonyms
Retraction-Convergence Nystagmus; Koerber-Salus-Elschnig Sign


\medterm{OAB} Alternate index label; see \textit{Overactive bladder}.

\medterm{Oat Cell Carcinoma} Alternate historical term; see \textit{Small-Cell Carcinoma}.

\medterm{OB/GYN} Abbreviation for obstetrics and gynecology, the medical specialties that care for pregnancy, childbirth, the period after birth, and the reproductive system.



\medterm{Obesity} A chronic disease involving abnormal or excessive accumulation of body fat that may impair health. It is linked with a higher risk of conditions such as type 2 diabetes, high blood pressure, heart disease, stroke, sleep apnea, osteoarthritis, and some cancers. Many factors can contribute, including genes, metabolism, hormones, medicines, eating patterns, activity, environment, and social conditions; it is not caused by one factor or by lack of willpower alone. Body mass index (BMI) is commonly used to classify obesity in adults, but it does not measure body fat directly and is interpreted with other information.

\medterm{Obesity Hypoventilation Syndrome} A breathing disorder in which excess body weight contributes to too little breathing and a buildup of carbon dioxide while awake. It often occurs with sleep apnea and may cause sleepiness, headaches, breathlessness, low oxygen, and heart strain.

\medterm{Obex Region} A small area at the lower end of the brain’s fourth fluid-filled space, where it narrows into the central canal of the spinal cord. It lies near centers controlling breathing, swallowing, and vomiting.

\medterm{Objective Response Rate} The percentage of people in a cancer study whose tumors shrink by a predefined amount or disappear during treatment. It does not show how long the response lasts or whether treatment improves survival.

\medterm{Objects Or Insects In Ear} An object or insect in the ear canal can cause pain, hearing loss, bleeding, discharge, or a sensation of movement. Do not insert tools or liquids; removal by a clinician is safest when the object is deep, sharp, chemical, or near the eardrum.

\synonyms
Ear Object In (Objects Or Insects In Ear)


\medterm{Oblique Fracture} A break in a bone that runs diagonally across it.

\medterm{Oblique Plane} A slice or view through the body or an organ at an angle other than straight across. Medical imaging can use oblique planes to show structures that are difficult to see in standard horizontal, vertical, or front-to-back views.

\medterm{Obliteration} Complete closing, destruction, or disappearance of a space, passage, or hollow channel in the body. It may result from scarring, inflammation, growth, injury, or treatment and can prevent normal movement of air, blood, fluid, or other contents.
\medterm{Observation} Planned monitoring of a person, symptom, disease, or test result without immediate treatment or a procedure. Follow-up timing and signs that require reassessment should be defined.

\medterm{Obsession} A recurrent, intrusive, unwanted thought, image, or urge that causes distress and is difficult to control.
\medterm{Obsessions} Repeated, unwanted thoughts, images, or urges that enter the mind despite attempts to stop them and cause distress. Common examples involve contamination, harm, mistakes, symmetry, or taboo subjects; having an obsession does not mean a person will act on it.

\medterm{Obsessive Compulsive Disorder} Alternate index label for Obsessive-Compulsive Disorder.

\medterm{Obsessive-Compulsive And Related Disorders} A group of mental disorders involving recurrent intrusive thoughts, repetitive behaviors or mental acts, preoccupation with perceived appearance flaws, difficulty discarding possessions, or repetitive body-focused behaviors. Examples include obsessive-compulsive disorder, body dysmorphic disorder, hoarding disorder, trichotillomania, and excoriation disorder.

\medterm{Obsessive-Compulsive Disorder} A mental-health disorder involving unwanted, repeated thoughts, images, or urges and repetitive behaviors or mental acts performed to reduce distress or prevent a feared event. The symptoms are time-consuming or cause clinically significant distress or impairment.

\medterm{Obsessive-Compulsive Personality Disorder} A long-lasting pattern of excessive concern with order, rules, perfection, and control. It can make a person inflexible, slow to complete tasks, or unable to delegate and can interfere with relationships or work.

\medterm{Obstetric Anal Sphincter Injury} A severe tear during childbirth that extends into the muscles controlling the anus. It may cause pain, difficulty controlling gas or stool, or later incontinence and needs prompt examination, repair, and follow-up.
\medterm{Obstetric Cholestasis} Alternate index label; see \textit{Cholestasis of pregnancy}.

\medterm{Obstetric Fistula} A hole that forms between the birth canal and the bladder or rectum after prolonged, obstructed labour without timely medical care. It causes continuous leakage of urine or stool and may lead to skin infections, kidney problems, and social isolation. Most cases can be prevented with timely maternity care, and many can be repaired with surgery.

\medterm{Obstetric Nursing} Obstetric nursing provides care during pregnancy, labor, birth, and the postpartum period, including monitoring, education, support, and recognition of complications.

\medterm{Obstetric Ultrasound} An ultrasound scan during pregnancy used to check the baby's age, growth, anatomy, position, heartbeat, and sometimes the placenta and fluid.

\medterm{Obstetrical Forceps} A surgical instrument used in selected vaginal births to guide the baby’s head when delivery needs assistance. Safe use requires the correct position, an experienced clinician, and assessment of risks and alternatives.

\medterm{Obstetrician} A medical doctor who cares for pregnant people during pregnancy, labor, birth, and the postpartum period. Obstetricians manage routine and high-risk pregnancies, perform deliveries, and treat complications such as high blood pressure, diabetes, bleeding, or preterm labor.

\medterm{Obstetrics} The branch of medicine concerned with care during pregnancy, labor, birth, and the weeks after birth, including routine care and treatment of complications affecting the pregnant person or baby.

\medterm{Obstipation} Severe constipation with inability to pass stool or gas. It can occur with a blockage in the intestine and may cause abdominal swelling, pain, vomiting, or serious illness requiring urgent assessment.
\medterm{Obstructed Infected Urinary System} A blocked urinary tract with an infection behind the blockage. It is a medical emergency because infection can spread quickly and the blockage often must be relieved urgently.

\medterm{Obstructed Labour} Failure of the fetus to pass through the birth canal despite adequate uterine contractions, commonly because of disproportion, abnormal presentation, or an obstructing lesion.

\medterm{Obstruction} A blockage that prevents normal movement through a body passage, such as an airway, intestine, bile duct, blood vessel, or urinary tract. Causes include swelling, stones, scars, tumors, clots, and foreign objects.

\medterm{Obstructive Azoospermia} A cause of male infertility in which no sperm enters the semen because a reproductive tube is blocked or missing, although the testes make sperm normally. Causes include a missing vas deferens, vasectomy, infection-related scarring, blocked ejaculatory ducts, or surgical injury. Testicle size and hormone levels are often normal. Treatment may reopen the tube or collect sperm for use with in-vitro fertilization.

\synonyms
OA; Excretory Azoospermia

\medterm{Obstructive Hydrocephalus} A buildup of cerebrospinal fluid caused by a blockage in its normal flow through the brain's ventricles.

\medterm{Obstructive Jaundice} Yellowing of the skin or eyes caused by blockage of bile flow from the liver. It may produce dark urine, pale stools, itching, abdominal pain, or fever and can result from a stone, narrowing, inflammation, or tumor.

\medterm{Obstructive Nephropathy} Kidney damage caused by a blockage that prevents urine from draining normally. Stones, tumors, scarring, an enlarged prostate, or birth-related problems can cause it; untreated blockage may lead to infection, permanent loss of kidney function, or kidney failure.

\medterm{Obstructive Parotitis} Inflammation of a saliva gland near the ear because its drainage tube is blocked, often by a stone. Painful swelling usually worsens during meals, and trapped saliva can lead to repeated infection and tenderness.

\medterm{Obstructive Shock} A life-threatening drop in blood flow caused by a physical blockage, such as a large blood clot in the lungs, pressure around the heart, or a collapsed lung under pressure.

\medterm{Obstructive Sleep Apnea} Repeated narrowing or closing of the upper airway during sleep, causing pauses in breathing, snoring, fragmented sleep, daytime tiredness, and sometimes reduced oxygen levels.

\synonyms
Sleep Apnea, Obstructive

\medterm{Obstructive Sleep Apnea - Adults} Obstructive sleep apnea is repeated upper-airway collapse during sleep, causing pauses in breathing, oxygen fluctuation, snoring, fragmented sleep, and daytime impairment.


\medterm{Obstructive Uropathy} A blockage anywhere between the kidneys and the urethra that slows or stops urine flow. Stones, scarring, an enlarged prostate, tumors, or pressure from nearby tissue may cause it and can damage the kidneys if not relieved.

\medterm{Obtunded} Having reduced alertness and slowed responses but still responding to strong voice or touch. It can result from medicines, illness, metabolic problems, poisoning, or brain disease and needs assessment.

\medterm{Obturator} Relating to a structure that closes an opening. In the pelvis, the obturator nerve helps control muscles that draw the thigh inward and carries feeling from part of the inner thigh; the obturator artery supplies nearby muscles.

\medterm{Obturator Hernia} A rare hernia in which part of the intestine pushes through an opening in the pelvic bone. It can block the bowel and may cause pain in the inner thigh.

\medterm{Obtuse Marginal Artery} A branch of the circumflex coronary artery that runs along the outer border of the heart and supplies part of the left ventricle. Its size and number vary between people; blockage can cause a heart attack.


\medterm{Occipital} Relating to the back of the head, the occipital bone, or the occipital lobe. These structures protect and support the brain and help process visual information and head or neck movement.
\medterm{Occipital Bone} The bone at the back and base of the skull. It surrounds the opening where the spinal cord joins the brain and connects with the first neck vertebra.

\medterm{Occipital Diastasis} Widening or separation of a joint between bones at the back of a newborn’s skull. It may follow difficult birth or abnormal development, and imaging assesses its extent and associated injury.

\medterm{Occipital Hematoma} A collection of blood in the scalp or tissues over the back of the head, often after trauma; a rapidly enlarging swelling or associated neurologic symptoms requires urgent assessment.

\medterm{Occipital Lobe} The part of the brain at the back of the head that receives and interprets visual signals. Injury or disease may cause missing areas of vision, difficulty recognizing objects, or visual hallucinations.

\medterm{Occipital Neuralgia} Brief attacks of sharp, stabbing, or electric-shock pain caused by irritation of nerves at the back of the head. Pain may spread from the upper neck across the scalp and can be associated with tenderness or sensitivity to touch.

\medterm{Occiput} The back part of the skull and the back of the head. It protects the rear of the brain and includes the occipital bone, which joins the skull base and neck.
\medterm{Occiput Posterior Position} A position during labor in which the back of the baby’s head faces the pregnant person’s back. It can cause back pain or slow descent, although many babies rotate to a better position during labor.

\medterm{Occiput Transverse Position} A fetal head position in which the back of the baby’s head faces the pregnant person’s side. It often rotates during labor; if it persists, descent may be slow and assisted delivery or cesarean birth may be needed.

\medterm{Occlude} To block or close a passage, vessel, airway, duct, or opening. An occlusion may be caused by a clot, stone, swelling, scar, tumor, foreign object, or deliberate medical treatment.

\medterm{Occlusal Adjustment} Selective reshaping of tooth surfaces to improve how the upper and lower teeth meet and reduce abnormal contact or bite interference.

\medterm{Occlusal Guard} A removable device worn over the teeth to protect them from grinding, clenching, or injury. A night guard may reduce tooth wear during sleep, while a sports mouthguard protects against blows; the device should fit properly and does not treat every cause of jaw pain.

\medterm{Occlusal Splint} A removable dental appliance that changes or protects tooth contact and may reduce effects of grinding, jaw-joint strain, or selected bite problems.

\medterm{Occlusion} Complete blockage of a blood vessel or body passage. It may be caused by a blood clot, plaque, pressure, swelling, or an intentional medical procedure; the effects depend on the location and can include loss of blood flow, pain, or organ damage.

\medterm{Occlusive Vascular Disease} Disease in which narrowing or blockage of a blood vessel reduces or stops blood flow to tissues.

\medterm{Occult} Hidden or not readily visible or detectable on routine examination. An occult disease or bleeding may be found only with imaging, microscopy, or laboratory testing.

\medterm{Occult Blood} Blood that is not visible to the naked eye and is detected by laboratory testing. In stool, it may be detected by a chemical guaiac test or an immunochemical test; the method and specimen determine the collection instructions and possible dietary or medicine effects. A positive stool test indicates blood is present but does not establish its cause.

\medterm{Occult Blood Screen} Alternate index label; see \textit{Fecal occult blood test}.

\medterm{Occult Gastrointestinal Bleeding} Slow or intermittent bleeding somewhere in the digestive tract that cannot be seen in the stool. It may cause iron-deficiency anemia, tiredness, or a positive stool blood test and requires evaluation to find the source.

\medterm{Occult Primary Cancer} Alternate index label; see \textit{Carcinoma of unknown primary}.

\medterm{Occupational Asthma} Asthma caused or worsened by substances at work, such as flour dust, latex, chemicals, or fumes. Symptoms may improve on days away from work and worsen during exposure; careful exposure history and breathing tests help confirm it.

\medterm{Occupational COPD} Chronic obstructive pulmonary disease caused or worsened by repeated workplace exposure to dust, fumes, chemicals, smoke, or other inhaled hazards. It causes lasting airflow limitation, cough, phlegm, and breathlessness; smoking and other risks may also contribute.

\medterm{Occupational Dermatoses} Skin diseases caused or worsened by workplace conditions, such as repeated wet work, chemicals, gloves, dust, plants, or friction. Common problems include irritant or allergic contact dermatitis, often affecting the hands; reducing exposure and protecting the skin help prevent recurrence.

\medterm{Occupational Disease} An illness caused or significantly worsened by workplace exposures or conditions. Examples include asthma from dust, hearing loss from noise, skin disease from chemicals, and lung disease from harmful particles.
\medterm{Occupational Lung Disease} Lung disease caused or worsened by dusts, fumes, gases, vapors, fibers, or other exposures at work. Examples include occupational asthma, pneumoconiosis, hypersensitivity pneumonitis, and asbestos-related disease.
\medterm{Occupational Exposure} Contact with a potentially harmful substance, energy source, or condition during work. Examples include radiation, chemicals, dust, noise, heat, infection, and repetitive strain. Risk depends on the hazard, amount, duration, and protective controls; monitoring and prevention are part of occupational health.

\medterm{Occupational Exposure Limit} The highest recommended or legally permitted workplace exposure to a hazardous substance over a specified time. The limit varies by substance, country, route of exposure, and whether it applies to a brief peak or an average work shift.

\medterm{Occupational Health} The medical field focused on preventing work-related illness, injury, and harmful exposure while supporting workers' health.
\medterm{Occupational Hearing Loss} Hearing damage caused by repeated or prolonged workplace noise or substances that harm the ear. It is usually permanent and may cause difficulty understanding speech, especially in background noise; reducing exposure and using effective hearing protection help prevent it.

\medterm{Occupational Rehabilitation} Assessment, treatment, and workplace support that help a person return to work or remain employed after illness, injury, or disability. It may include graded activity, skills training, equipment, job changes, and coordination with employers and clinicians.

\medterm{Occupational Safety And Health} Rules and practices that protect workers from injuries, illness, dangerous exposures, and unsafe conditions. Measures include training, protective equipment, ventilation, safe procedures, health monitoring, and reporting or controlling workplace hazards.
\medterm{Occupational Therapist} A healthcare professional who helps people perform daily activities through rehabilitation, adaptive techniques, environmental changes, and assistive equipment.

\medterm{Occupational Therapy} Rehabilitation that helps a person perform daily activities, work, school, and other meaningful tasks after illness, injury, or disability. It may use exercises, practice, adaptive equipment, changes to the environment, and strategies for movement, thinking, or sensory problems.

\medterm{Occupational Therapy Assessment} An evaluation of how illness, injury, disability, or development affects daily activities, school, work, and self-care. The therapist identifies strengths and barriers and recommends exercises, skills practice, equipment, or environmental changes.

\medterm{OCD} Obsessive-compulsive disorder, a condition involving recurrent unwanted thoughts or urges and/or repetitive behaviors or mental acts that cause distress or interfere with daily life.

\synonyms
Obsessive Compulsive Disorder (OCD)

\medterm{Ochre} A yellow-brown color or earth pigment. In medicine, ochre discoloration can describe yellow-brown skin or tissue changes caused by pigment, blood breakdown, chronic congestion, or other conditions.


\medterm{Ochronosis} Blue-black or brown discoloration of cartilage, tendons, skin, or the whites of the eyes caused by buildup of a chemical called homogentisic acid. It is most often linked to the inherited disorder alkaptonuria and may cause arthritis.

\medterm{OCP} OCP commonly means oral contraceptive pill, a medicine containing estrogen, progestin, or both to prevent pregnancy and sometimes treat menstrual disorders.


\medterm{Octreoscan} Nuclear-medicine imaging using indium-111 octreotide, which binds to somatostatin receptors. It was used to locate some neuroendocrine tumors and their spread; newer receptor-targeted PET scans are often preferred.

\medterm{Octreotide} Octreotide is a somatostatin analogue that suppresses growth hormone, several gastrointestinal hormones, and intestinal secretion; it is used for acromegaly, certain neuroendocrine tumors, and selected secretory diarrhoeas.

\medterm{Ocular} Relating to the eye or vision. The term may describe an eye structure, examination, medicine, injury, infection, or disease affecting sight. Ocular emergencies such as sudden severe pain or sudden loss of vision need immediate eye care.
\medterm{Ocular Albinism} Ocular albinism is an inherited disorder with reduced pigment primarily in the eyes, causing decreased visual acuity, nystagmus, photophobia, and abnormal retinal development.

\medterm{Ocular Fixation} The ability to hold the eyes steadily on a target; abnormal fixation can impair visual acuity, tracking, or alignment.

\medterm{Ocular Histoplasmosis Syndrome} An eye condition linked to an earlier histoplasmosis infection. It can cause scars in the retina and, in some people, abnormal blood vessels that damage central vision.

\medterm{Ocular Hypertension} Higher-than-normal pressure inside the eye without current signs of glaucoma damage. It increases the future risk of glaucoma and needs regular eye checks.

\medterm{Ocular Ischemic Syndrome} Eye findings caused by severely reduced blood flow through the arteries supplying the eye, usually from advanced carotid artery disease. It may cause gradual or sudden vision loss, eye pain, retinal changes, or abnormal new blood vessels.

\medterm{Ocular Melanoma} A cancer that begins in pigment-producing cells of the eye.

\medterm{Ocular Migraine} A migraine-related visual disturbance, often involving flashing lights, zigzag lines, or a temporary blind spot. It may occur with or without headache.

\medterm{Ocular Motility} The ability of the eyes to move together in a coordinated way. Abnormal eye movements or poor alignment can cause double vision, eyestrain, dizziness, or reduced vision and may result from muscle, nerve, brain, or eye disease.

\medterm{Ocular Myasthenia Gravis} A form of myasthenia gravis in which weakness affects only the muscles that move the eyes and lift the eyelids. It causes drooping eyelids or double vision that worsens with use and improves with rest. Symptoms may later spread to swallowing, speech, breathing, or limb muscles, so follow-up is important.

\synonyms
OMG; Ocular MG

\medterm{Ocular Myiasis} Infestation of the eye or nearby tissues by fly larvae. It can cause pain, irritation, tearing, and vision problems and requires prompt eye care.

\medterm{Ocular Orbit} The bony cavity around the eye. It contains the eyeball, eye-movement muscles, nerves, blood vessels, glands, and protective tissues.

\medterm{Ocular Pain} Pain in or around the eye caused by surface irritation, infection, inflammation, trauma, glaucoma, or orbital
disease. Severe pain with vision loss, light sensitivity, nausea, or injury requires urgent
ophthalmic assessment.

\medterm{Ocular Rosacea} Rosacea affecting the eyelids or surface of the eyes, causing burning, redness, a gritty feeling, dryness, or light sensitivity. Persistent inflammation can affect the cornea and may impair the eye surface.
\medterm{Ocular Surface Disease} A group of conditions affecting the eye's surface, eyelids, or tears, including dry eye and eyelid inflammation. Symptoms may include burning, redness, irritation, and blurred or changing vision.

\medterm{Ocular Toxoplasmosis} A retinal infection or inflammation caused by Toxoplasma gondii, which can produce blurred vision, floaters, pain, and recurrent retinal scars.

\medterm{Oculocardiac Reflex} A slowing of the heartbeat or similar body response caused by pressure or pulling on the eye or tissues around it.

\medterm{Oculocerebrorenal Syndrome} A rare inherited disorder, also called Zellweger spectrum disorder, that affects the brain, eyes, kidneys, liver, and other organs. Severe forms begin in newborns and can cause weak muscles, seizures, feeding problems, and major developmental problems.

\synonyms
Zellweger spectrum disorder

\medterm{Oculocutaneous Albinism} Alternate index label; see \textit{Albinism}.

\medterm{Oculogyric Crisis} A prolonged, involuntary upward turning of the eyes, often caused by a medicine or neurologic disorder.
\medterm{Oculomotor Nerve} The third cranial nerve, controlling most eye movements, eyelid lifting, and pupil narrowing. Damage can cause double vision, a drooping eyelid, an outward-turning eye, or an abnormally enlarged pupil.
\medterm{Oculomotor Nerve Palsy} Weakness or paralysis of muscles supplied by the third cranial nerve, causing eyelid drooping, double vision, abnormal eye position, and sometimes a dilated pupil.

\medterm{Oculomotor Nucleus} A group of nerve cells in the midbrain that controls most eye movements, lifts the upper eyelid, and narrows the pupil through the third cranial nerve.

\medterm{Oculopharyngeal Muscular Dystrophy} An inherited muscle disease that usually begins in middle age and gradually weakens the eyelids and throat muscles. Drooping eyelids, difficulty swallowing, and later weakness of the hips and shoulders are typical. A gene test can confirm the diagnosis; swallowing assessment, exercises, and surgery may help prevent complications.

\synonyms
OPMD

\medterm{Oculoplastic Procedures} Operations on the eyelids, tear ducts, eye socket, or nearby facial tissues to restore function, repair injury, remove disease, or change appearance. Risks include bleeding, infection, scarring, dry eyes, and rarely injury affecting vision.

\medterm{Oculosympathetic Palsy} Alternate index label; see \textit{Horner syndrome}.

\medterm{ODD} Alternate index label; see \textit{Oppositional defiant disorder}.

\medterm{Oddi, Sphincter Of} A ring-shaped muscle at the opening where bile and pancreatic juice enter the small intestine. It normally opens during digestion and closes afterward; abnormal tightening or blockage can cause pain or pancreatitis.

\medterm{Odontalgia} Pain coming from a tooth or the tissues around it. Common causes include decay, a cracked tooth, infection, gum disease, or injury; swelling, fever, difficulty swallowing, or severe pain needs prompt dental or medical care.
\medterm{Odontoblast} A cell along the inner part of a tooth that makes dentin, the hard layer beneath enamel. Odontoblasts produce dentin while teeth develop and can make limited repairs after injury.
\medterm{Odontoclast} A cell that breaks down the hard tissues of a tooth. This process helps baby teeth be shed naturally, but unwanted tooth resorption can weaken a tooth and may follow injury, inflammation, pressure, or other disease.
\medterm{Odontodysplasia} Odontodysplasia is abnormal development of tooth enamel and dentin, producing thin, weak, discolored teeth that may erupt late or be easily damaged.

\medterm{Odontogenesis} The process by which teeth form and develop before and after birth. Problems during development can cause missing teeth, weak enamel, abnormal dentin, or changes in tooth shape and eruption.
\medterm{Odontogenic Cutaneous Fistula} An abnormal tunnel from an infected tooth to the skin, usually on the chin, cheek, or jaw. It may appear as a small draining opening; treating the tooth infection usually closes the tunnel.

\medterm{Odontogenic Infection} An infection arising from a tooth or the tissues around it, usually caused by dental caries, pulp infection, gum disease, or trauma. It may form an abscess and spread into the jaw, face, neck, or bloodstream.

\medterm{Odontogenic Cyst} A fluid-filled sac arising from tissues involved in tooth development, usually within the jaw. It may be painless or cause swelling, tooth displacement, infection, or bone damage.
\medterm{Odontogenic Fibroma} An odontogenic fibroma is a rare benign tumor arising from tooth-forming connective tissue, usually presenting as a jaw lesion.

\medterm{Odontogenic Keratocyst} A cyst that develops in the jaw from tooth-forming tissue. It may grow into bone without causing symptoms, can recur after removal, and is sometimes associated with an inherited skin-cancer syndrome; imaging and tissue examination guide treatment.

\medterm{Odontogenic Neoplasm} An odontogenic neoplasm is a growth arising from tissues involved in tooth formation, usually in the jaw; it may be benign or malignant.

\medterm{Odontogenic Tumors And Cysts} Alternate index label; see \textit{Jaw tumors and cysts}.

\medterm{Odontoid Process} A peg-like projection from the second neck vertebra that fits into the first neck vertebra and allows the head to turn from side to side.

\medterm{Odontoma} A noncancerous overgrowth of tooth-forming tissue, usually found in the jaw and often discovered when a permanent tooth fails to erupt. It may look like several small teeth or a solid calcified mass and is usually removed if it blocks eruption.

\medterm{Odontopathy} Any disease or abnormality affecting a tooth or the tissues supporting it. Examples include tooth decay, gum disease, infection, abnormal tooth development, wear, and injury. Untreated dental infection can spread to surrounding tissues and occasionally to the bloodstream.


\medterm{Odynophagia} Pain when swallowing, often felt as burning, squeezing, or a sticking sensation in the throat or chest. Causes include infection, reflux, ulcers, medicines, injury, and other swallowing-tube disorders.

\medterm{Oedema} British spelling of edema: abnormal fluid accumulation in tissues that causes swelling. Causes include injury, infection, venous disease, heart or kidney failure, medicines, and lymphatic blockage.

\medterm{Oesophageal} Relating to the oesophagus, the muscular tube carrying swallowed material from the pharynx to the stomach.
\medterm{Oesophageal Atresia} A birth defect in which the swallowing tube ends blindly instead of connecting normally to the stomach. A baby may drool, cough, choke, or turn blue during feeds; urgent specialist care and surgery are required, especially when an abnormal connection to the airway is also present.
\medterm{Oesophageal Ulcer} An open sore in the swallowing tube between the throat and stomach. Reflux, infection, medicines, corrosive injury, inflammation, or cancer may cause it, leading to painful swallowing, chest pain, bleeding, or difficulty swallowing.

\medterm{Oesophageal Varices} Enlarged fragile veins in the wall of the swallowing tube, usually caused by high pressure in the liver’s circulation. They may cause no symptoms until they rupture, leading to vomiting blood or black stools and a life-threatening emergency.
\medterm{Oesophagitis} Inflammation of the swallowing tube between the throat and stomach. Reflux, infections, medicines, allergies, radiation, or corrosive substances can cause it, leading to heartburn, painful or difficult swallowing, chest pain, or bleeding.

\medterm{Oesophagoscopy} Examination of the food pipe using a thin viewing tube passed through the mouth. It can investigate swallowing difficulty, pain, bleeding, or abnormal tissue, and can remove a foreign body, take samples, stop bleeding, or widen a narrowing. Sedation or anesthesia is usually used.
\medterm{Oesophagus} The muscular tube that transports swallowed food and liquid from the pharynx to the stomach.
\medterm{Oesophagus, Diseases Of} Disorders affecting the swallowing tube between the throat and stomach. They include reflux, inflammation, narrowing, movement problems, cancer, infection, and birth defects, often causing swallowing difficulty, pain, heartburn, or regurgitation.
\medterm{Oestradiol} British spelling of estradiol, the main naturally occurring estrogen during the reproductive years. It supports reproductive development, the menstrual cycle, bone health, and other body functions.
\medterm{Oestriol} A form of estrogen made in large amounts during pregnancy by the placenta using substances from the fetus and pregnant person. Its level can be measured in some prenatal screening tests, but it is not interpreted alone to diagnose a problem.
\medterm{Oestrogen} A group of hormones important for reproductive development, menstrual cycles, pregnancy, bone strength, and other body functions. They are made mainly by the ovaries before menopause and by other tissues, including the placenta during pregnancy.

\medterm{Oestrogen Replacement Therapy} Treatment with estrogen, sometimes combined with a progestogen, to relieve symptoms caused by low estrogen such as hot flushes and vaginal dryness. Benefits and risks depend on age, medical history, uterus status, dose, and route, so treatment is individualized.

\medterm{Oestrogens} A group of steroid hormones involved in reproductive development, menstrual cycling, pregnancy, bone, and other physiological functions.
\medterm{Oestrone} A form of estrogen made by the ovaries and by conversion of other hormones in body tissues. It becomes the main estrogen after menopause and affects reproductive tissues, bones, blood vessels, and metabolism.
\medterm{Ofatumumab} A monoclonal antibody against CD20 used for selected blood cancers and relapsing multiple sclerosis; it can increase infection risk and reduce B cells.

\medterm{Off-Label} Use of an approved medicine for an indication, dose, route, or population not specified in its regulatory approval.
\medterm{Off-Label Use} Use of an approved medicine for a condition, dose, route, or age group not specifically included in its official approval. A clinician may recommend it when evidence and expected benefits support its use; the decision should be discussed with the patient.

\medterm{Offspring} The young or descendants produced by one organism or by a pair. In medicine and genetics, offspring may be studied to understand inherited traits, pregnancy outcomes, developmental conditions, and disease risk.

\medterm{Ofloxacin} A fluoroquinolone antibiotic used for selected susceptible bacterial infections. It can cause serious tendon, nerve, heart-rhythm, and other adverse effects, so it should be used only when clinically appropriate.
\medterm{Ogilvie Syndrome} A dangerous widening of the large intestine without a physical blockage, usually in a seriously ill or recently operated patient. Abdominal swelling, pain, nausea, or failure to pass stool may occur, and the bowel can perforate if untreated.

\medterm{Ogilvie’S Syndrome} Another spelling of Ogilvie syndrome: severe widening of the large intestine without a physical blockage, often after surgery or during serious illness. It needs medical monitoring because pressure can reduce blood flow or cause a bowel perforation.

\medterm{Ohm} The standard unit of electrical resistance, equal to one volt per ampere. In medical equipment, resistance affects how electrical current flows through devices, sensors, wires, and body tissues.
\medterm{OHSS} Alternate index label; see \textit{Ovarian hyperstimulation syndrome}.


\medterm{Oidium} A fragment or spore-like fungal element formed by fragmentation of a hypha; the term is also used in older fungal classifications.
\medterm{Oil} A hydrophobic liquid used as a nutrient, solvent, emollient, or medicinal preparation; safety depends on the specific oil.
\medterm{Oil-Based Paint Poisoning} Harm from swallowing or inhaling oil-based paint or its solvents. It may irritate the lungs, cause dizziness or drowsiness, and cause chemical pneumonia if liquid enters the lungs; do not induce vomiting.

\medterm{Oily Hair} Hair that appears greasy because the scalp produces excess oil or because of product buildup. It may occur with infrequent washing, hormonal changes, or some skin conditions and is not poisoning.

\medterm{Oily Skin} Skin that produces more oil than usual, often causing shine, enlarged-looking pores, or acne. Gentle cleansing and non-comedogenic products may help; persistent acne or irritation may need medical assessment.

\medterm{Ointment} A thick, semisolid preparation applied to the skin, inside the nose, or other body surfaces. Its oily base helps protect the area and keep it moist; some ointments also contain medicines that treat infection, inflammation, pain, or other conditions.

\medterm{Ointment Bases} Semisolid materials that carry medicines applied to the skin or moist body surfaces. They affect how easily a product spreads, how much moisture it retains, how much medicine enters tissue, and how long it remains stable.

\medterm{Okadaic Acid} A toxin made by certain marine organisms that can accumulate in shellfish. Eating contaminated shellfish may cause diarrhoea and vomiting; in laboratories it is used to study enzymes that remove phosphate groups from proteins.
\medterm{Okihiro Syndrome} A rare inherited condition caused by changes in the SALL4 gene. It may combine abnormalities of the thumb or forearm with difficulty moving the eyes, and may also affect hearing, the heart, kidneys, or spine.

\medterm{Olanzapine} Olanzapine is an atypical antipsychotic that blocks dopamine and serotonin receptors and treats schizophrenia and bipolar disorder; weight gain, diabetes, dyslipidaemia, sedation, and movement disorders are important risks.

\medterm{Olaparib} Olaparib is a PARP inhibitor used for selected cancers with BRCA or other homologous-recombination repair defects; anaemia, nausea, fatigue, myelosuppression, and rare myelodysplastic syndrome or acute myeloid leukaemia may occur.


\medterm{Old Age} Later life; health assessment should consider function, comorbidity, cognition, and social circumstances rather than age alone.
\medterm{Oleander Poisoning} Poisoning from eating oleander, which contains heart-active toxins. Nausea, vomiting, dizziness, confusion, slow or irregular heartbeat, fainting, or seizures require immediate emergency care.


\medterm{Olecranon} The prominent bony process at the proximal end of the ulna forming the point of the elbow.
\medterm{Oleic Acid} A monounsaturated fat found in olive oil, nuts, seeds, avocados, and many other foods. Replacing some saturated fats with it may support healthier blood cholesterol levels as part of a balanced diet.
\medterm{Oleum} A very corrosive mixture of sulfur trioxide and sulfuric acid, used industrially and dangerous to skin, eyes, and lungs.
\medterm{Olfaction} The sense of smell. Odor molecules stimulate special cells in the upper part of the nose, which send signals to the brain to identify and perceive odors.

\medterm{Olfactory} Relating to the sense of smell or the nerves and brain pathways that process odors. Problems in this system can cause reduced smell, loss of smell, distorted smells, or unusual odor sensations.
\medterm{Olfactory Bulb} The olfactory bulb is a forebrain structure that receives signals from nasal smell receptors and begins processing odor information.

\medterm{Olfactory Cortex} The olfactory cortex is a group of brain regions that interprets smell signals and links them with memory, emotion, and behavior.

\medterm{Olfactory Gyrus} The olfactory gyrus is a cortical region associated with processing olfactory information and includes areas near the medial temporal and frontal lobes.

\medterm{Olfactory Learning} Olfactory learning is acquisition or modification of behavior or memory through exposure to odors and their associations.

\medterm{Olfactory Mucosa} The olfactory mucosa is specialized nasal lining containing receptor neurons, supporting cells, and glands that detect odor molecules.

\medterm{Olfactory Nerve} The first cranial nerve, which carries smell signals from specialized cells high in the nose to the brain. Injury or disease can reduce or eliminate the sense of smell or distort odors.

\medterm{Olfactory Neuroblastoma} A rare malignant tumor arising in the upper nasal cavity near the olfactory nerves, often causing nasal blockage, nosebleeds, or reduced smell.

\medterm{Olfactory Tract} The olfactory tract is a bundle of nerve fibers carrying smell information from the olfactory bulb to cerebral regions involved in odor perception.

\medterm{Olfactory Tubercle} The olfactory tubercle is a ventral forebrain region involved in odor processing, reward, motivation, and integration with limbic circuits.



\medterm{Oligoasthenoteratozoospermia} A semen abnormality in which sperm numbers are low, movement is reduced, and fewer sperm have a typical shape. It can reduce fertility and may result from testicular disease, hormone problems, heat, medicines, smoking, or an unknown cause. Repeat semen testing and evaluation help guide treatment or assisted reproduction.

\synonyms
OAT Syndrome; OAT

\medterm{Oligoclonal Bands} Groups of immune proteins found in blood or the fluid surrounding the brain and spinal cord. Bands found only in that fluid suggest immune activity inside the nervous system, but they can occur in multiple sclerosis, infections, and other inflammatory conditions.

\medterm{Oligocythaemia} An unusually low number of blood cells, especially red blood cells. The uncommon term may describe anemia or another low blood-cell count, so the specific cell type and cause must be identified with a blood test.
\medterm{Oligodendrocyte} A supporting cell in the central nervous system that forms myelin around axons. Myelin helps nerve signals travel quickly; oligodendrocytes can be damaged in disorders such as multiple sclerosis.
\medterm{Oligodendroglia} Oligodendroglia are glial cells of the central nervous system that produce myelin around axons and support nerve conduction.

\medterm{Oligodendroglial Neoplasm} A brain tumor made of cells that resemble oligodendrocytes, the cells that form the insulating coating around nerve fibers. Diagnosis uses tissue examination and genetic testing, and the tumor’s behavior depends on its grade and molecular features.

\medterm{Oligodendroglioma} A brain tumor defined by characteristic changes in the IDH gene and loss of parts of chromosomes 1 and 19. It usually begins in the cerebral hemispheres and may cause seizures, headache, or other neurologic symptoms; grade affects its behavior and treatment.

\medterm{Oligodontia} A condition present from birth in which many permanent teeth fail to develop, usually more than six excluding wisdom teeth. It may affect chewing, speech, appearance, and jaw development, requiring dental replacement or orthodontic care.
\medterm{Oligogenic Inheritance} A pattern in which changes in a small number of genes work together to influence a disease or characteristic. The combination and strength of the changes can make symptoms differ among affected people.

\medterm{Oligohydramnios} Too little amniotic fluid around the baby during pregnancy. It may result from leaking membranes, placental problems, dehydration, a pregnancy that continues beyond the due date, or problems affecting the baby's kidneys.

\medterm{Oligomenorrhea} Oligomenorrhea is infrequent or widely spaced menstrual bleeding, often from anovulation, polycystic ovary syndrome, endocrine disease, undernutrition, excessive exercise, stress, or perimenopause.

\medterm{Oligomenorrhoea} Menstrual periods that occur less often than expected, usually with cycles longer than about 35 days. Causes include pregnancy, polycystic ovary syndrome, thyroid disease, stress, low body weight, intense exercise, and some medicines; evaluation depends on age and symptoms.
\medterm{Oligometastasis} Cancer that has spread to only a small number of separate sites, often no more than five. In selected people, surgery, focused radiation, or another local treatment may control these deposits for a long time.


\medterm{Oligonucleotide} A short piece of DNA or RNA made in the laboratory or found in cells. It can be used to detect genetic material, block or change gene messages, or help read and study genes.

\medterm{Oligopeptide} A short chain of amino acids, usually containing only a few to about 20 building blocks. Oligopeptides can act as hormones, signaling molecules, antimicrobial substances, or parts of larger proteins.

\medterm{Oligosaccharide} A short chain of sugar molecules joined together. Oligosaccharides form parts of cell-surface molecules, affect how cells recognize one another, and occur naturally in foods where some are used by gut bacteria.

\medterm{Oligospermia} A lower-than-expected sperm concentration in semen. It can reduce fertility and may result from hormone problems, infections, medicines, heat, toxins, varicoceles, genetic conditions, or problems with sperm production.
\medterm{Oliguria} Abnormally low urine production. It may result from dehydration, low blood flow, blockage, kidney disease, or medicines and can be an early sign of sudden kidney injury.

\medterm{Olivary} Relating to the olive-shaped inferior olivary nucleus of the medulla or to an olive-like structure.
\medterm{Olivary Nucleus} The olivary nucleus is a brainstem structure that sends timing and motor-learning signals to the cerebellum through climbing fibers.

\medterm{Olive Oil} A plant oil rich in monounsaturated fatty acids, used as food and in some topical preparations.
\medterm{Olive Pontocerebellar Atrophy} A group of progressive brain disorders that damage areas controlling balance and coordination. They may cause worsening unsteadiness, clumsy movements, slurred speech, stiffness, or other neurologic problems.

\medterm{Olivopontocerebellar Atrophy} Olivopontocerebellar atrophy is a degenerative condition affecting the inferior olives, pons, and cerebellum, causing progressive ataxia and other neurologic dysfunction.

\medterm{Ollier's Disease} A nonhereditary skeletal disorder with multiple enchondromas, causing deformity, limb-length discrepancy, and an increased risk of malignancy.
\medterm{Olmesartan Medoxomil} An oral prodrug converted to olmesartan, an angiotensin II receptor blocker used to treat high blood pressure.

\medterm{Olopatadine Hydrochloride} Olopatadine hydrochloride is an antihistamine used in eye drops or nasal spray to treat allergic conjunctivitis or rhinitis.

\medterm{OLP} Alternate index label; see \textit{Oral lichen planus}.

\medterm{Olsenella} Olsenella is a genus of anaerobic bacteria found mainly in the mouth and digestive tract; human infection is uncommon.

\medterm{Olsenella Uli} An anaerobic gram-positive bacterium found in the mouth and associated with periodontal disease; invasive infection is rare.


\medterm{Omacetaxine Mepesuccinate} Omacetaxine mepesuccinate is a protein-synthesis inhibitor used for selected chronic myeloid leukemia resistant or intolerant to tyrosine-kinase inhibitors.

\medterm{Omalizumab} Omalizumab is a monoclonal antibody that binds circulating IgE and reduces allergic inflammation; it is used for selected allergic asthma, chronic urticaria, nasal polyps, and food allergy, with rare anaphylaxis.

\medterm{Omasum} The omasum is the third compartment of a ruminant stomach, absorbing water and small particles before contents reach the abomasum.

\medterm{Omega-3} A family of polyunsaturated fats found in fish, flaxseed, walnuts, and some supplements, supporting heart and brain health.


\medterm{Omega-3 Fatty Acids} A family of fats found in oily fish, seafood, some seeds, nuts, and plant oils. EPA and DHA support cell and eye function, while ALA is a plant form; supplements are not appropriate for everyone and can interact with medicines.

\medterm{Omega-6} A family of essential polyunsaturated fats, found in vegetable oils, nuts, seeds, and many prepared foods. The body uses them for cell membranes and signaling. Linoleic acid is the main dietary form; replacing saturated fats with unsaturated fats generally supports heart health.

\medterm{Omega-6 Fatty Acids} Essential polyunsaturated fats found in plant oils, nuts, seeds, and many prepared foods. The body needs them for cell membranes and signaling, but overall health depends on balanced intake with other fats.

\synonyms
n-6 fatty acids

\medterm{Omega-N-Methylarginine} Omega-N-methylarginine is a methylated arginine compound that can inhibit nitric-oxide synthase and affect vascular signaling.
\medterm{Omenn Syndrome} A severe inherited immune deficiency that usually begins in infancy. Babies may develop widespread red, peeling skin, swollen lymph nodes, enlarged liver or spleen, diarrhea, poor growth, and repeated serious infections.

\medterm{Omentectomy} Removal of part or all of the greater omentum, commonly performed during staging and
cytoreduction for ovarian and other intra-abdominal cancers.

\medterm{Omentum} A fold of fatty tissue inside the abdomen that hangs from the stomach and partly covers the intestines. It contains blood vessels and immune cells and can help limit the spread of infection or inflammation.

\medterm{Omeprazole} Omeprazole is a proton-pump inhibitor that irreversibly suppresses gastric H+/K+-ATPase; it treats reflux and peptic-ulcer disease and supports Helicobacter pylori eradication, while prolonged use may affect magnesium, B12, bone, and infection risk.

\medterm{Ommaya Reservoir} A surgically implanted reservoir connected to a catheter in a brain ventricle, used to sample cerebrospinal fluid or deliver selected medicines.

\medterm{Omphalitis} Infection of a newborn’s umbilical stump, usually during the first weeks of life. Redness spreading onto the skin, swelling, pus, bad smell, fever, poor feeding, or unusual sleepiness requires urgent treatment because infection can spread rapidly.

\medterm{Omphalocele} A birth defect in which abdominal organs protrude through an opening at the navel into a protective sac. Other chromosome or organ abnormalities may occur, so newborn assessment and imaging are needed.

\medterm{Omphalocele Repair} Treatment to place an infant's abdominal organs back inside the abdomen and close the opening in the abdominal wall. Repair may be immediate or staged according to the size of the defect and the baby's breathing and circulation.

\medterm{Onchocerca Volvulus} A parasitic roundworm spread by repeated bites of infected blackflies that breed near fast-flowing rivers. Adult worms live in lumps under the skin and release young worms that can migrate through the skin and eyes.

\medterm{Onchocerciasis} A parasitic infection caused by Onchocerca volvulus and spread by infected blackflies. It can cause intensely itchy skin, nodules, skin changes, and eye inflammation that may lead to permanent visual impairment or blindness.

\medterm{Oncocytic Neoplasm} A growth made of cells containing unusually many energy-producing structures, giving them a granular appearance under a microscope. It may be noncancerous or cancerous, depending on the organ and the tumor’s other features.

\medterm{Oncocytoma} A usually noncancerous tumor made of cells containing many energy-producing structures. It most often develops in the kidney but can occur in other organs; testing is needed to distinguish it from cancer.

\medterm{Oncogene} A gene that can make cells grow or survive without normal control when it is changed, overactive, or present in too many copies. Oncogene changes can contribute to the development and progression of cancer.

\medterm{Oncogenic} Capable of causing or promoting cancer. An oncogenic change, virus, gene, or chemical can increase abnormal cell growth or reduce normal control of cell division.

\medterm{Oncogenic Virus} A virus capable of contributing to cancer by altering cell growth, suppressing tumor defenses, or integrating or maintaining oncogenic genetic material.

\medterm{Oncogenicity} The ability of a substance, infection, genetic change, or other factor to contribute to the development of cancer.

\medterm{Oncologic Emergency} A cancer-related problem that can quickly threaten life, organ function, or the nervous system and needs urgent treatment. Examples include spinal-cord pressure, blocked major chest veins, dangerously high calcium, tumor breakdown, airway blockage, and infection during low white-cell counts.

\medterm{Oncologist} A doctor who specializes in cancer care. Medical oncologists use medicines, radiation oncologists use radiation, and surgical oncologists operate; many people receive care from more than one type of oncologist.

\medterm{Oncology} The medical field concerned with preventing, finding, and treating cancer and supporting people after treatment. Oncologists may use surgery, radiation, medicines, or combinations of treatments and work with other professionals to provide symptom and emotional support.
\medterm{Oncology Nursing} Oncology nursing provides cancer-related assessment, treatment support, symptom control, education, coordination, survivorship, and end-of-life care.
\medterm{Oncology PET} A positron-emission tomography scan, often combined with CT, that uses a small radioactive tracer to show active cancer in the body. It may help assess spread, treatment response, or recurrence, but inflammation can also appear active.

\medterm{Oncology Rehabilitation} Healthcare that helps people prevent or manage cancer-related problems before, during, or after treatment. It may address fatigue, pain, weakness, nerve symptoms, swelling, movement, nutrition, work, daily activities, and emotional wellbeing.

\medterm{Oncolytic Virotherapy} Oncolytic virotherapy uses modified viruses that preferentially infect and destroy cancer cells and may stimulate an immune response against the tumor.

\medterm{Oncolytic Virus} A virus selected or modified to infect and destroy cancer cells more readily than normal cells. It may also stimulate the immune system to attack the tumor; use is limited to selected cancers and requires specialist treatment.

\medterm{Oncolytic Viruses} Viruses selected or engineered to infect and destroy cancer cells while stimulating an immune response against the tumor.

\medterm{Oncoprotein} A protein that promotes abnormal cell growth or survival and can contribute to cancer. It may result from a cancer-causing gene, a mutation, a chromosome change, or infection with certain viruses.

\medterm{Oncosis} A form of cell injury in which cells swell and may die, usually because severe lack of oxygen or another injury disrupts energy production and the movement of salt and water across the cell membrane.


\medterm{Oncotic Pressure} The pulling force created by proteins in the blood, mainly albumin, that helps keep fluid inside blood vessels. When blood protein is low, fluid may leak into tissues and cause swelling.

\medterm{Oncotoxic} Harmful to cancer cells or relating to substances that damage tumors. An oncotoxic treatment may also affect healthy cells, so its benefits, side effects, dose, and safety require careful medical monitoring.
\medterm{Ondansetron} Ondansetron is a selective 5-HT3-receptor antagonist used to prevent nausea and vomiting after chemotherapy, radiotherapy, surgery, or other triggers; constipation, headache, and QT prolongation can occur.

\medterm{Ondansetron Hydrochloride} The hydrochloride salt of ondansetron, a serotonin 5-HT3 receptor blocker used to prevent nausea and vomiting.

\medterm{Ondine's Curse} A rare disorder in which automatic breathing is impaired, especially during sleep; also called congenital central hypoventilation syndrome.

\synonyms
congenital central hypoventilation syndrome.

\medterm{Ondine’S Curse} A rare disorder in which automatic breathing becomes inadequate, especially during sleep, because brain signals fail to drive breathing normally. It can cause dangerously low oxygen and may require assisted ventilation.
\synonyms
Congenital central hypoventilation syndrome

\medterm{One-Hour Glucose Challenge Test} A screening test for diabetes during pregnancy, usually performed at 24–28 weeks. After drinking a measured glucose solution, blood sugar is checked one hour later. A high result does not diagnose diabetes but usually leads to a longer confirmatory test.

\medterm{Oneirophrenia} A dreamlike state with confusion or impaired reality testing; the term is uncommon and does not represent a specific modern psychiatric diagnosis.

\medterm{Onset} Onset is the time or manner in which a symptom, disease, or physiologic change begins, such as sudden, gradual, acute, or insidious onset.

\medterm{Onychia} Inflammation or infection of the nail or nail-producing tissue. It can cause redness, swelling, pain, pus, or changes in nail growth and may result from bacteria, fungi, injury, or skin disease. Treatment depends on the cause.
\medterm{Onychocryptosis} Onychocryptosis is an ingrown nail, usually at the great toe, in which the nail edge penetrates surrounding skin and causes pain or inflammation.

\medterm{Onychocryptosis (Ingrowing Toenail)} A condition in which the edge of a nail grows into the surrounding skin, causing pain, inflammation, and sometimes infection.

\synonyms
Ingrowing Toenail

\medterm{Onychocryptosis (Ingrown Toenail)} Alternate index label for Ingrown Toenail.

\medterm{Onychogryphosis} Severe thickening and overgrowth of a nail, which becomes curved and may resemble a claw. It is more common in older adults and can follow injury, fungal infection, poor circulation, or difficulty caring for the feet; it may cause pain in shoes.
\medterm{Onycholysis} Separation of the nail plate from the underlying nail bed, beginning distally or
laterally. Causes include trauma, psoriasis, fungal infection, thyroid disease,
medications, and chemical exposure; the underlying cause should be identified.

\medterm{Onychomadesis} Onychomadesis is proximal separation and shedding of a nail caused by temporary arrest of nail-matrix growth, often after infection, systemic illness, trauma, or medication exposure.

\medterm{Onychomycosis} A fungal infection of a nail, more common in toenails than fingernails. The nail may become thick, yellow or white, brittle, crumbly, or separated from the skin. A sample may be tested before prolonged antifungal treatment because other nail diseases look similar.

\medterm{Onychomycosis (Fungal Nails)} Alternate index label for Fungal Nails.

\medterm{Oocyte} An immature egg cell in an ovary. It develops before birth, remains paused until a menstrual cycle, and may mature and be released for possible fertilization.

\medterm{Oocyte Cryopreservation} Freezing unfertilized eggs for possible use in future pregnancy. Hormone medicines stimulate the ovaries, and the eggs are collected through the vagina with a needle before being rapidly frozen. Eggs may be stored before cancer treatment or for other fertility-preservation reasons, but freezing does not guarantee a future pregnancy.

\synonyms
Egg Freezing; Oocyte Vitrification

\medterm{Oocyte Retrieval} A procedure used in assisted reproduction to collect mature eggs from the ovaries. After hormone stimulation, a clinician guides a thin needle through the vagina using ultrasound and gently draws fluid containing the eggs into a container for laboratory fertilization.

\synonyms
Ovum Pick-Up (OPU); Egg Retrieval; Follicular Aspiration

\medterm{Oogenesis} The formation and maturation of egg cells, or ova, in the ovaries. It begins before birth and continues through the reproductive years, with changes during each menstrual cycle.
\medterm{Oogonium} An oogonium is an early female germ cell that proliferates during embryonic development and gives rise to primary oocytes.

\medterm{Oophorectomy} Surgery to remove one or both ovaries. It may treat cancer, a cyst, infection, torsion, or reduce inherited cancer risk; removing both ovaries before natural menopause causes an immediate fall in ovarian hormones.

\medterm{Oophoritis} Inflammation of an ovary, usually associated with pelvic infection or, less commonly, autoimmune disease.
It may cause lower abdominal or pelvic pain, fever, abnormal discharge, or menstrual changes;
severe pain requires assessment to exclude ovarian torsion or abscess.

\medterm{Ooplasm} Ooplasm is the cytoplasm of an egg cell, containing organelles, nutrients, and molecular factors that support fertilization and early embryo development.

\medterm{Opacity} Clouding of a normally clear part of the eye, such as the cornea, lens, or fluid inside the eye. It can blur vision and may result from injury, infection, inflammation, swelling, or age-related change.

\medterm{Open Bite} An open bite is a malocclusion in which upper and lower teeth do not meet in part of the mouth when the jaws are closed.

\medterm{Open Fractures} Broken bones associated with a wound that connects the fracture to the outside air. They have a high risk of infection and tissue damage and require urgent cleaning, antibiotics, and specialist treatment.

\medterm{Open Gallbladder Removal} Surgery to remove the gallbladder through one larger abdominal incision. It may be needed when keyhole surgery is unsafe or cannot be completed; bleeding, infection, bile leakage, and injury to nearby ducts or organs are possible.

\medterm{Open Lung Biopsy} An open lung biopsy removes a small piece of lung through a surgical chest incision for microscopic diagnosis when less invasive tests are insufficient.

\medterm{Open Pleural Biopsy} An open pleural biopsy removes a sample of the lining around the lung through surgery to investigate cancer, infection, or unexplained pleural disease.

\medterm{Open Prostatectomy} Surgery through an abdominal incision to remove the part of an enlarged prostate that blocks urine flow. It can improve severe urinary obstruction but may cause bleeding, infection, urine leakage, or changes in sexual and urinary function.

\medterm{Open Reduction And Internal Fixation} Surgery to expose a broken bone, put the pieces back into position, and hold them with plates, screws, wires, or rods. It is used when a cast or other nonsurgical treatment cannot provide safe alignment or stability.


\medterm{Open Surgery} Surgery performed through a relatively large incision that allows the surgeon to see and reach the affected area directly. It is used when a larger working space is needed, although recovery may be longer than after some keyhole procedures.

\medterm{Open-Angle Glaucoma} A type of glaucoma in which the eye’s drainage angle remains open, but fluid drains too slowly. Pressure may rise and damage the optic nerve, causing gradual, permanent vision loss, often without early symptoms.

\medterm{Open-Book Fracture} A severe pelvic injury in which the front or back pelvic joints separate, making the pelvis widen like an open book. It usually follows major trauma and can cause life-threatening internal bleeding or organ injury.

\medterm{Open-Heart Surgery} Surgery in which the chest is opened to repair or treat the heart or its major blood vessels. Some operations use a heart–lung machine while the heart is temporarily stopped; risks depend on the operation and the person’s health.

\medterm{Operable} Able to be treated with surgery because the disease can be safely reached or removed.
\medterm{Operant Conditioning} Learning in which the likelihood of a behavior changes because of its consequences, such as reinforcement or punishment.

\medterm{Operating Room} An operating room is a controlled sterile or semi-sterile clinical area equipped for surgical procedures, anesthesia, monitoring, and recovery support.

\medterm{Operating Tables} Adjustable tables designed to support and position patients during surgery or other procedures while allowing access and maintaining safety.

\medterm{Operation} A planned surgical procedure used to diagnose, repair, remove, replace, or improve part of the body. It may be performed through an incision or with instruments inserted through small openings, depending on the condition and goal.
\medterm{Operative Dentistry} The branch of dentistry that repairs teeth damaged by decay, fractures, wear, or developmental defects. Treatment may include fillings, bonding, reshaping, and other procedures that restore appearance, strength, and function.
\medterm{Operative Vaginal Delivery} Vaginal birth assisted with a vacuum cup or forceps when the baby's position is suitable and birth needs to be completed promptly. Reasons may include prolonged pushing or concern about the baby's condition; tears, bleeding, or injury to the baby can occur.

\medterm{Operator} A person who performs a procedure or operates equipment; in healthcare, training and authorization are essential.
\medterm{Operculum} A lid, covering, or flap-like structure; in anatomy it may cover an opening or cortical region.

\medterm{Ophthalmia} Inflammation of the eye, especially the conjunctiva or other surface tissues. It may cause redness, discharge, pain, light sensitivity, or blurred vision and can result from infection, allergy, injury, or inflammation.
\medterm{Ophthalmia Neonatorum} Inflammation or infection of a newborn baby's eye during the first weeks of life. It may be caused by bacteria passed during birth, other infections, or chemical irritation; redness, swelling, or pus requires prompt treatment to protect the cornea and vision.

\medterm{Ophthalmic} Relating to the eye. An ophthalmic medicine, drop, ointment, or device is designed for use on or around the eye and must be used as directed to prevent injury or infection.
\medterm{Ophthalmic Emergency} An eye problem that needs immediate assessment to prevent permanent vision loss or serious injury. Examples include sudden vision loss, retinal detachment, severe eye infection, chemical exposure, a penetrating injury, or rapidly painful high eye pressure.

\medterm{Ophthalmic Solution} A sterile liquid medicine made for application to the eye. It must remain clean and should not be shared or used after contamination because contaminated drops can cause serious eye infections.

\medterm{Ophthalmitis} Inflammation of the eye caused by infection, immune disease, injury, or surgery. It may cause redness, pain, light sensitivity, discharge, or blurred vision and needs prompt assessment when severe.

\medterm{Ophthalmologist} A medical doctor trained to diagnose and treat eye disease and vision problems. Ophthalmologists can prescribe medicines and glasses, perform eye procedures and surgery, and manage conditions affecting the eye, optic nerve, and visual system.

\medterm{Ophthalmology} The medical specialty concerned with the eyes, vision, and the parts of the nervous system used for seeing. It includes eye examinations, treatment of eye disease, vision correction, and procedures or surgery when needed.
\medterm{Ophthalmomyiasis} Infestation of the eye by fly larvae. Larvae may remain on the eye surface or enter deeper tissues, causing intense irritation, pain, inflammation, or vision loss; prompt examination is needed for safe removal and treatment.

\medterm{Ophthalmoplegia} Weakness or paralysis of one or more muscles that move the eyes, causing limited eye movement, double vision, or an abnormal eye position. Causes include nerve disease, stroke, diabetes, thyroid disease, injury, and muscle disorders.

\medterm{Ophthalmoscope} A lighted instrument used to examine the back and internal structures of the eye, including the retina, macula, optic nerve, and retinal blood vessels. Direct ophthalmoscopes provide a magnified view through the pupil, while indirect ophthalmoscopes use a separate light and lens to provide a wider view. Findings may relate to diabetes, high blood pressure, glaucoma, infection, inflammation, or other eye and nervous-system disorders.

\medterm{Ophthalmoscopy} Examination of the inside of the eye with a lighted instrument. It allows a clinician to inspect the retina, optic nerve, and blood vessels for changes caused by diabetes, high blood pressure, glaucoma, inflammation, or other disease.

\medterm{Opiate} A drug derived directly from chemicals in the opium poppy, such as morphine or codeine, or a closely related semisynthetic form. Opiates can relieve pain but may cause sleepiness, constipation, dependence, and dangerously slow breathing.

\medterm{Opiate Alkaloid} An opiate alkaloid is a naturally occurring alkaloid from the opium poppy, such as morphine or codeine, with opioid analgesic and dependence potential.

\medterm{Opiate And Opioid Withdrawal} Opioid withdrawal is a predictable syndrome after reducing or stopping opioid exposure, producing anxiety, sweating, muscle aches, diarrhea, abdominal cramps, nausea, and autonomic symptoms.

\medterm{Opiate Assay} A laboratory test that detects or measures opiates or their breakdown products in urine, blood, saliva, or another specimen. Results may support clinical care or toxicology testing, but some synthetic opioids require separate tests.
\medterm{Opiate Overdose} A life-threatening amount of an opioid medicine or drug that slows brain and breathing activity. Extreme sleepiness, pinpoint pupils, slow or absent breathing, blue lips, or inability to wake require emergency help and naloxone when available.

\medterm{Opioid} A natural or synthetic substance that acts on opioid receptors to relieve pain and may cause sedation or dependence.

\medterm{Opioid Agonist Treatment} Treatment for opioid use disorder using a prescribed medicine, usually methadone or buprenorphine, that activates opioid receptors in a controlled way. It reduces withdrawal and cravings, lowers illicit opioid use and overdose risk, and works best with ongoing medical support.

\medterm{Opioid Intoxication} Opioid intoxication can cause euphoria, drowsiness, slowed breathing, pinpoint pupils, impaired coordination, coma, and death; respiratory depression requires immediate emergency treatment.

\medterm{Opioid Overdose} A medical emergency characterized by the classic triad of respiratory depression, miosis
(pinpoint pupils), and decreased level of consciousness. It can progress to respiratory
arrest, cardiovascular collapse, and death.

\medterm{Opioid Poisoning} Alternate index label; see \textit{Opioid overdose}.

\medterm{Opioid Receptor} A cell-surface protein activated or blocked by opioids and naturally produced endorphins. Mu, kappa, and delta receptors influence pain, breathing, mood, bowel movement, and reward; excessive mu-receptor activation can suppress breathing.

\medterm{Opioid Reversal} Emergency treatment of opioid toxicity with naloxone, a competitive opioid-receptor
antagonist. It is titrated to restore adequate breathing, with repeated doses or monitoring sometimes needed because
its effect may wear off before the opioid.

\medterm{Opioid Toxicity} Toxic effects of opioid agonists, classically respiratory depression, central nervous
system depression, and miosis. Management prioritizes ventilation and airway support;
naloxone is titrated to restore adequate breathing while recognizing the risk of
recurrent toxicity and withdrawal.

\medterm{Opioid Toxidrome} A poisoning pattern caused by opioid excess, characterized by respiratory depression, reduced consciousness, and small pupils; naloxone and ventilatory support may be required.

\medterm{Opioid Use Disorder} A pattern of opioid use that causes loss of control, cravings, continued use despite harm, or problems at home, work, or school. It can occur with prescribed medicines or illicit opioids and is treatable with medication, counseling, and support.

\medterm{Opioid Withdrawal} A characteristic group of symptoms occurring after abrupt reduction or cessation of prolonged opioid use, including anxiety, muscle aches, sweating, abdominal cramps, diarrhea, nausea, and insomnia.

\medterm{Opisthorchis Felineus} A liver fluke acquired by eating raw or undercooked freshwater fish; chronic infection can cause bile-duct inflammation and increase cholangiocarcinoma risk.

\medterm{Opisthotonos} Opisthotonos is severe backward arching of the head, neck, and spine from sustained muscle spasm, classically associated with tetanus, meningitis, strychnine poisoning, or severe neurologic injury.

\medterm{Opium} The dried latex of Papaver somniferum containing opioid alkaloids; it can cause dependence and respiratory depression.
\medterm{Opponens} A muscle that moves a body part toward an opposing surface, especially the thumb across the palm.
\medterm{Opportunistic Infection} An infection that causes illness mainly when the immune system is weakened or a normal protective barrier is damaged. Causes may include cancer treatment, HIV, transplantation, severe illness, or prolonged use of immune-suppressing medicines.

\medterm{Opposition} Movement of the thumb across the palm so it can touch the fingertips. It is important for grasping and depends on several small movements at the joint where the thumb meets the wrist.

\medterm{Oppositional Defiant Disorder} A childhood or adolescent disorder involving a persistent pattern, lasting at least six months, of angry or irritable mood, argumentative or defiant behavior, or vindictiveness toward an authority figure or other person. The pattern causes clinically significant distress or impairment and is not limited to conflicts with siblings.


\medterm{Opsoclonus} Rapid, involuntary, back-and-forth eye movements in many directions without pauses. It may occur with muscle jerks and poor coordination, sometimes because of an immune reaction, infection, medicine, or an underlying tumor, and needs medical evaluation.

\medterm{Opsonin} An immune protein that coats a germ and marks it for destruction. Antibodies and complement proteins can act as opsonins by helping infection-fighting cells recognize and engulf the germ.
\medterm{Opsonization} The coating of a germ or other particle with immune proteins so that infection-fighting cells can recognize, capture, and destroy it more easily. Antibodies and complement proteins commonly perform this function.

\medterm{Optic Atrophy} Damage and loss of optic-nerve fibers, causing permanent reduction in vision, color vision, or part of the visual field. Causes include glaucoma, poor blood supply, inflammation, pressure, injury, toxins, and inherited disorders; treatment targets the cause when possible.

\medterm{Optic Chiasm} An X-shaped junction beneath the brain where some fibers from the optic nerves cross. Damage can cause characteristic loss of parts of the visual field, depending on the location and cause.
\medterm{Optic Chiasma} The crossing point beneath the brain where some fibers from the two optic nerves switch sides. Damage there can cause loss of the outer parts of both visual fields.

\medterm{Optic Disc} The small area at the back of the eye where the optic nerve and blood vessels pass through the retina. It has no light-sensing cells, so it creates the eye’s normal blind spot.

\medterm{Optic Disc Edema} Swelling where the optic nerve enters the eye. It may result from raised pressure around the brain, inflammation, poor blood supply, or other eye or nerve disease. New swelling with headache or vision change requires urgent assessment.

\medterm{Optic Disk} The point where the optic nerve leaves the retina; it has no photoreceptors and forms the blind spot.

\medterm{Optic Glioma} A usually slow-growing tumor of the optic nerve or the visual pathway leading to the brain, most often in children. It may cause reduced vision, loss of part of the visual field, a bulging eye, or abnormal eye movements and is associated with neurofibromatosis type 1.

\synonyms
Optic Nerve Glioma
Optic Pathway Glioma

\medterm{Optic Nerve} The nerve that carries visual signals from each eye to the brain. Damage can cause reduced vision, loss of color or part of the visual field, and sometimes pain with eye movement.

\medterm{Optic Nerve Atrophy} Permanent damage and loss of nerve fibers carrying vision from the eye to the brain. It can reduce sharpness, color vision, or part of the visual field and may follow glaucoma, inflammation, poor blood supply, pressure, injury, or toxins.

\medterm{Optic Nerve Disorder} A disease affecting the optic nerve, which carries visual signals from the eye to the brain; it can cause reduced vision, color loss, or visual-field defects.

\medterm{Optic Nerve Glioma} Alternate index label; see \textit{Optic glioma}.

\medterm{Optic Neuritis} Inflammation or demyelination of the optic nerve, often causing subacute vision loss, impaired color vision, and pain with
eye movement. It may be associated with multiple sclerosis or other inflammatory,
infectious, or immune-mediated disorders.

\medterm{Optic Pathway Glioma} Alternate index label; see \textit{Optic glioma}.

\medterm{Optic Tract} A bundle of nerve fibers carrying visual information from the optic chiasma to brain areas that process sight and control visual reflexes. Damage causes loss of a characteristic part of the visual field.

\medterm{Optical Coherence Tomography} A painless scan that uses reflected light to create detailed cross-sectional pictures of the retina and other eye structures. It helps detect and monitor glaucoma, macular disease, swelling, and other changes without touching the eye.

\medterm{Optical Illusion} A visual perception in which an object, image, or movement appears different from its physical properties. Illusions are common, but new or persistent visual distortions may occur with eye, brain, medicine, or migraine-related disorders.
\medterm{Optical Rotation} The turning of plane-polarized light as it passes through a substance whose molecules have a particular three-dimensional arrangement. Measuring it can help identify or assess the concentration of some chemicals and medicines.

\medterm{Optician} A technician who fits and dispenses eyeglasses or contact lenses.

\medterm{Optometrist} A trained eye-care professional who examines vision and many ocular conditions and prescribes optical correction within local scope of practice.
\medterm{Optometry} The healthcare profession focused on examining vision and the eyes, prescribing glasses or contact lenses, detecting eye disease, and providing treatment or referral within the professional scope allowed by local law.

\medterm{Orthoptist} An eye-care professional who assesses and treats problems with eye movement, eye alignment, and the use of both eyes together. Orthoptists commonly help manage strabismus, amblyopia, double vision, and some vision problems in children and adults.

\medterm{Oral} Relating to the mouth, or taken through the mouth. Oral medicines are swallowed and absorbed through the digestive tract unless otherwise specified.
\medterm{Oral Administration} Taking a medicine by mouth so it is absorbed through the digestive tract; absorption can be affected by food, vomiting, interactions, or gastrointestinal disease.

\medterm{Oral Appliance Therapy} Treatment using a removable dental appliance during sleep, commonly to advance the mandible or retain the tongue and reduce upper-airway collapse in selected obstructive sleep apnea.

\medterm{Oral Cancer} Cancer of the lips, mouth, tongue, gums, inner cheeks, palate, or nearby throat. Tobacco, alcohol, betel nut, and some HPV infections increase risk; a persistent sore, lump, white or red patch, pain, or swallowing difficulty needs assessment.


\medterm{Oral Candidiasis} A yeast infection of the mouth, usually caused by Candida. It may produce wipeable white patches, soreness, burning, altered taste, or cracks at the mouth corners. Antibiotics, steroid inhalers, dentures, diabetes, dry mouth, and weak immunity increase risk.

\medterm{Oral Cavity} The space inside the mouth, extending from the lips to the throat and bounded by the cheeks, palate, tongue, and floor of the mouth. It receives food, begins digestion, and helps with speech, taste, and breathing.

\medterm{Oral Cavity Cancer} Cancer arising in the lips, tongue, gums, inner cheeks, floor of the mouth, or palate. Tobacco, alcohol, betel quid, and some HPV infections increase risk; persistent sores, lumps, patches, pain, or numbness need examination and often biopsy.

\medterm{Oral Cholecystogram} An older X-ray test in which swallowed contrast collects in the gallbladder, allowing gallstones and gallbladder function to be assessed. Ultrasound and other imaging methods have largely replaced it.

\medterm{Oral Contraceptive} A birth-control pill containing progestin alone or progestin combined with estrogen. It mainly prevents ovulation and thickens cervical mucus; it does not protect against sexually transmitted infections.

\medterm{Oral Glucose Tolerance Test} A diagnostic test measuring plasma glucose response after a standardized oral glucose
load. It is used in selected evaluations of diabetes, impaired glucose tolerance, and
gestational diabetes when fasting glucose or HbA1c is insufficient.

\medterm{Oral Health} Oral health is the condition of the teeth, gums, mouth, jaw, and related functions such as eating, speaking, and swallowing.

\medterm{Oral Hemorrhage} Bleeding from the mouth, gums, tongue, throat, or nearby tissues, caused by trauma, dental disease, infection, medicines, or a bleeding disorder.

\medterm{Oral Human Papillomavirus Infection} Oral human papillomavirus infection affects the mouth or throat; many infections clear without symptoms, while persistent high-risk types can contribute to some cancers.

\medterm{Oral Hygiene} Practices maintaining oral health: brushing twice daily with fluoride toothpaste,
flossing daily, mouthwash (optional), regular dental visits. Proper brushing technique:
soft brush, 45-degree angle to gumline, gentle circular motions. Interdental cleaning
essential for plaque removal.

\medterm{Oral Hypoglycemics Overdose} Taking too much diabetes medicine by mouth can cause dangerously low blood glucose, especially with insulin-releasing medicines. Sweating, shaking, confusion, seizures, or unconsciousness require immediate glucose treatment and emergency advice.

\medterm{Oral Leukoplakia} A persistent white patch in the mouth that cannot be rubbed off and is not explained by another condition. It is often linked to tobacco, alcohol, or betel-nut use and can sometimes develop into mouth cancer, so a persistent patch needs dental or medical assessment.

\medterm{Oral Lichen Planus} A long-lasting inflammatory condition causing lace-like white lines, red patches, or painful mouth sores, often on the cheeks or gums. It may fluctuate over time; persistent painful or changing areas need dental review to exclude other disease.

\synonyms
Lichen Planus, Oral
OLP

\medterm{Oral Moniliasis} Oral moniliasis is oral candidiasis, a Candida yeast infection causing white plaques, soreness, burning, or altered taste.


\medterm{Oral Mucosa} The moist lining of the mouth, including the inner lips and cheeks, gums, tongue, and palate. It protects deeper tissues, provides sensation, stays lubricated, and can develop ulcers, infections, inflammation, or tumors.


\medterm{Oral Mucous Cyst} An oral mucous cyst is a mucus-filled lesion caused by salivary-gland duct injury or blockage, often appearing on the lower lip or floor of the mouth.

\medterm{Oral Nutritional Supplement} A calorie-, protein-, or micronutrient-containing product consumed by mouth to
supplement, not automatically replace, usual food intake. The product and dose should
match the patient's needs, swallowing ability, disease restrictions, and tolerance.

\medterm{Oral Pathology} The study and diagnosis of diseases affecting the mouth, teeth, jaws, and related tissues. These include tooth decay, infections, cysts, inflammation, developmental changes, and tumors. Examination, imaging, laboratory tests, or a tissue sample may be needed to identify the cause.

\medterm{Oral Potentially Malignant Disorders} Conditions affecting the mouth lining that carry an increased risk of developing into oral cancer. Examples include leukoplakia, erythroplakia, and some long-lasting inflammatory or inherited disorders; appearance alone cannot determine the risk.
\medterm{Oral Radiology} The use and interpretation of X-rays and other scans of the mouth, teeth, jaws, face, and nearby structures. These images help find decay, infection, fractures, tumors, bone disease, and problems related to treatment planning.

\medterm{Oral Rehydration Salts} A balanced mixture of glucose and electrolytes dissolved in clean water to replace fluid and salts lost through diarrhea or vomiting.

\medterm{Oral Rehydration Therapy} Replacement of water and electrolytes by mouth using a balanced glucose-and-salt solution, especially to prevent or treat dehydration from diarrhoea or vomiting.

\synonyms
ORT

\medterm{Oral Sex (Fellatio, Cunnilingus)} Oral stimulation of the genitals for sexual pleasure. Fellatio: oral stimulation of the
penis. Cunnilingus: oral stimulation of the vulva and clitoris. Analingus (rimming):
oral stimulation of the anus.

\synonyms
Fellatio, Cunnilingus

\medterm{Oral Spray} A liquid medicine or preparation sprayed into the mouth or throat for local treatment or absorption through oral tissues.

\medterm{Oral Squamous Cell Carcinoma} The most common type of mouth cancer, arising from the cells lining the lips, tongue, gums, cheeks, or palate. A sore, lump, red or white patch, bleeding, pain, or numbness that persists needs examination and often a biopsy.

\medterm{Oral Surgery} Operations on the mouth, teeth, jaws, or nearby facial tissues. Procedures include tooth removal, dental implants, tissue sampling, cyst removal, repair of injury, and correction of jaw problems.

\medterm{Oral Thrush} A fungal infection of the mouth, usually caused by Candida, producing creamy white patches, soreness, or altered taste.

\synonyms
Candidiasis, Oral
Thrush, Oral

\medterm{Oral Toxicity} Harm to the mouth, throat, or digestive system caused by a swallowed poison, medicine, treatment, or excessive dose. It may cause burning, sores, nausea, vomiting, abdominal pain, or organ injury and may require urgent poison-control or medical advice.

\medterm{Orange Peel Skin} Skin with small dimples resembling an orange peel because swelling stretches the skin around hair follicles. It may result from blocked lymph drainage, especially in breast disease, and a new breast change needs assessment.
\medterm{Orbicularis} A circular muscle surrounding an opening, such as orbicularis oculi around the eye or orbicularis oris around the mouth.
\medterm{Orbit} The bony socket that contains and protects the eye, its muscles, nerves, blood vessels, and supporting tissues. Problems in the orbit can cause pain, swelling, double vision, bulging, or impaired eye movement.
\medterm{Orbit CT Scan} A CT scan of the eye socket that uses X-rays and computer images to show the bones, eye muscles, nerves, blood vessels, and surrounding tissues. It can help assess injury, infection, bleeding, or a growth.

\medterm{Orbital} Relating to the orbit, the bony cavity containing the eye, or to structures within it. Orbital disease may affect the eye, muscles, nerves, or surrounding tissues.
\medterm{Orbital Cellulitis} A serious infection in the tissues around and behind the eye. It can cause eyelid swelling, fever, pain with eye movement, double vision, or reduced vision and can spread to the brain, so it requires urgent hospital treatment.

\medterm{Orbital Contusion} A bruise of the tissues around the eye caused by blunt trauma, often producing swelling, bruising, and tenderness.

\medterm{Orbital Fracture} A break in one or more bones forming the eye socket, usually caused by facial trauma and potentially associated with double vision, eye-movement restriction, or injury to the eye.

\medterm{Orbital Fractures} Breaks in the bones surrounding the eye, which may cause swelling, double vision, numbness, restricted eye movement, or injury to the eye itself.

\medterm{Orbital Irradiation} Radiation treatment directed at tissues within the eye socket, sometimes used for severe thyroid eye disease.

\medterm{Orbital Metastasis} Cancer that has spread to the bony socket around the eye from another organ, commonly the breast, lung, or prostate. It may cause bulging of the eye, pain, double vision, restricted movement, or reduced sight.

\medterm{Orbital Myositis} Inflammation of one or more muscles that move the eye. It can cause pain, eyelid swelling, double vision, and restricted eye movement; treatment depends on the cause and may require urgent specialist assessment.

\medterm{Orbital Neoplasm} An abnormal growth in the eye socket that may arise from the eye, optic nerve, muscles, glands, bone, or other tissues. It can cause bulging, pain, double vision, or vision loss and may be benign or cancerous.

\medterm{Orbital Pseudotumor} Painful inflammation of tissues around the eye that is not caused by infection or a true tumor. It may cause eyelid swelling, bulging of the eye, double vision, and painful or limited eye movement.

\synonyms
idiopathic.


\medterm{Orchialgia} Pain in one or both testicles. Sudden severe pain, swelling, nausea, or a high-riding testicle may indicate twisting of the testicle and requires emergency assessment. Longer-lasting pain can result from infection, injury, varicocele, nerve pain, or other causes.
\medterm{Orchidectomy} Surgical removal of one or both testes, also called orchiectomy.
\medterm{Orchiectomy} Surgical removal of one or both testicles. Radical inguinal orchiectomy is diagnostic
and therapeutic for testicular cancer; bilateral orchiectomy is a hormonal therapy for
prostate cancer.

\medterm{Orchitis} Orchitis is inflammation of one or both testes, often from infection such as mumps or bacteria, causing pain, swelling, and sometimes impaired fertility.

\medterm{Order Coleoptera} The insect order containing beetles, characterized by hardened forewings that protect the flying wings; some species affect health as pests or vectors.

\medterm{Order Spirochaetales} A biological order of spiral-shaped bacteria that includes groups causing syphilis, Lyme disease, relapsing fever, and leptospirosis. Identification of the species determines medical significance and treatment.
\medterm{Orexin} Orexin, or hypocretin, is a hypothalamic neuropeptide that promotes wakefulness, arousal, appetite, and energy regulation; loss of orexin signaling is strongly associated with narcolepsy type 1.

\medterm{Orexin Receptor Antagonists} Medications that block orexin (hypocretin) receptors to promote sleep. Suvorexant and
lemborexant are approved for insomnia. They reduce wakefulness by blocking the arousal
system. They have lower abuse potential than benzodiazepines and do not cause
respiratory depression.

\medterm{Orf} A zoonotic viral skin infection caused by the parapoxvirus, transmitted to humans from
infected sheep and goats (direct contact with lesions, fomites—virus survives in scabs
for months). Orf is an occupational disease of farmers, butchers, and veterinarians.
Incubation: 3--7 days.

\medterm{Organ} A structure composed of two or more different types of tissues that has a recognizable
shape and performs specific functions.


\medterm{Organ Failure} Severe loss of an organ's ability to perform essential functions, such as filtering blood, exchanging oxygen, or pumping blood.

\medterm{Organ System} A group of organs that work together to perform an important body function, such as breathing or digestion.

\medterm{Organ Transplant} Surgical transfer of an organ from a donor to a recipient to replace failing organ function.
\medterm{Organ Transplantation} Replacement of a failing organ with a healthy organ from a donor. Transplants can restore function but require careful matching, lifelong monitoring, and medicines to prevent rejection; infection, rejection, and medicine-related harm remain possible.

\medterm{Organ Transplants} The surgical transfer of an organ from a donor to a recipient whose organ function is severely
impaired. Solid-organ transplantation may involve the kidney, liver, heart, lung, pancreas,
or intestine and requires donor assessment, recipient selection, and ongoing follow-up.

\medterm{Organelle} A specialized part inside a cell that performs a particular job. Examples include the nucleus, which stores genetic instructions, and mitochondria, which release energy; organelles work together to keep the cell alive.
\medterm{Organic} Relating to living matter or carbon-containing chemistry; in psychiatry, an older term for symptoms caused by a physical brain disorder.

\medterm{Organic Chemical} A carbon-containing chemical compound, usually also containing hydrogen and sometimes oxygen, nitrogen, sulfur, or other elements. Organic chemicals include medicines, fats, sugars, proteins, solvents, and many toxins.

\medterm{Organic Disease} An illness caused by a physical change or measurable disturbance in the body, such as infection, inflammation, tissue damage, or abnormal organ function. The term is older and is less useful when symptoms have not yet been explained.
\medterm{Organic Impurity} An organic impurity is an unintended carbon-containing chemical in a medicine, material, or specimen that may affect safety, quality, or measurement.


\medterm{Organization Name} The formal name of a group, institution, company, or agency; in healthcare records it should identify the responsible organization accurately.

\medterm{Organogenesis} The process during embryonic development in which groups of cells form and organize into organs. Disruption can lead to congenital abnormalities, while the process is also studied in tissue engineering and regenerative medicine.



\medterm{Organophosphate} A class of chemicals that can inhibit acetylcholinesterase. Some are pesticides or nerve agents and can cause cholinergic poisoning.

\medterm{Organophosphate Poisoning} A serious poisoning caused by pesticides or nerve agents that excessively stimulate nerves and muscles. It may cause sweating, drooling, vomiting, diarrhea, pinpoint pupils, muscle twitching, wheezing, weakness, seizures, or breathing failure. Remove contaminated clothing, wash the skin, and seek emergency care immediately.

\medterm{Organosilicon Compounds} Organosilicon compounds contain silicon bonded to carbon and are used in materials, medicines, devices, and industry; biological effects depend on the specific compound.

\medterm{Organotherapy} An outdated treatment method using extracts or tissues from organs, with little accepted modern medical use.
\medterm{Orgasm} The peak of sexual pleasure, usually accompanied by involuntary muscle and body responses that vary among individuals.
\medterm{Orgasm Male (Sexual (Sex) Problems In Men)} Alternate index label for Sexual (Sex) Problems in Men.

\medterm{Orgasmic Dysfunction} Persistent or recurrent difficulty reaching orgasm, markedly reduced orgasmic intensity, or absent orgasm that causes clinically significant distress. Medical, psychological, medication-related, hormonal, and relationship factors may contribute.
\medterm{Orgasmic Dysfunction In Women} Alternate index label for Orgasmic Dysfunction.

\medterm{Orientation} Awareness of personal identity, place, time, and situation. Clinicians assess orientation as part of mental-status and neurologic examinations.
\medterm{Orifice} An opening leading into or out of a body structure, organ, or passage, allowing materials to enter or leave.

\medterm{Orlistat} A medicine used with a reduced-calorie diet to help with weight loss or weight control. It stops some dietary fat from being absorbed, so oily stools, gas, and urgent bowel movements can occur; it can also reduce absorption of vitamins and some medicines.
\medterm{Ornithine} A nonessential amino acid and intermediate in the urea cycle that helps convert toxic ammonia into urea.

\medterm{Ornithine Decarboxylase} An enzyme that converts ornithine to putrescine, an early step in polyamine synthesis involved in cell growth.

\medterm{Ornithine Transcarbamylase Deficiency} OTC deficiency is an inherited urea-cycle disorder causing ammonia accumulation, vomiting, confusion, seizures, or coma, especially during metabolic stress.

\medterm{Orogastric Tube} A flexible tube passed through the mouth into the stomach. It may deliver nutrition or medicines, remove stomach contents, or decompress the stomach when nasal placement is unsuitable.

\medterm{Oropharyngeal} Oropharyngeal means relating to the mouth and the part of the throat behind it, including structures involved in breathing, swallowing, and speech.

\medterm{Oropharyngeal Airway} A rigid device placed through the mouth to keep the tongue from blocking the throat in an unconscious person. It must not be used in someone with a gag reflex and does not prevent vomit from entering the lungs.

\medterm{Oropharyngeal Carcinoma} Cancer arising in the part of the throat behind the mouth, including the tonsils, back of the tongue, soft palate, or throat wall. Tobacco, alcohol, and some HPV infections increase risk; assessment usually includes examination, tissue sampling, and scans.

\medterm{Oropharyngeal Dysphagia} Difficulty moving food or liquid from the mouth through the throat. It may cause coughing, choking, a wet voice, food coming through the nose, or food entering the airway and needs assessment when persistent.

\medterm{Oropharyngeal Neoplasm} An oropharyngeal neoplasm is an abnormal growth in the throat behind the mouth, including the tonsils and base of the tongue; it may be benign or malignant.

\medterm{Oropharynx} The part of the throat behind the mouth, including the tonsils and back of the tongue. It helps move food and liquid toward the esophagus and carries air toward the voice box.

\medterm{Oropharynx Lesion Biopsy} Sampling of an oropharyngeal lesion, usually through the mouth under local or procedural anesthesia, for histopathologic evaluation of infection, dysplasia, or cancer. Persistent or suspicious lesions require specialist assessment.

\medterm{Orotic Aciduria} An unusually high amount of orotic acid in urine. It may result from a rare inherited disorder affecting blood-cell and DNA building blocks, causing anemia and poor growth, or from a urea-cycle disorder with high ammonia.
\medterm{Orphan Disease} A rare disease affecting relatively few people; definitions and eligibility for orphan-drug programs vary by jurisdiction.

\medterm{Orphenadrine} A centrally acting medicine with anticholinergic properties used in some countries for painful muscle spasm; it can cause drowsiness, dry mouth, and confusion.



\medterm{Orthodontic Appliance} A device used to move teeth, guide jaw growth, maintain alignment, or correct a bite problem.

\medterm{Orthodontic Bracket} A bonded or banded attachment that transmits forces from an orthodontic archwire to a
tooth. Bracket position and wire sequence guide controlled tooth movement.

\medterm{Orthodontic Retainer} A removable or fixed appliance used after active orthodontic treatment to maintain tooth
position while periodontal and alveolar tissues remodel. Retention is often needed long
term because relapse can occur.

\medterm{Orthodontic Treatment (Dental Braces)} Treatment that uses braces, wires, aligners, or other appliances to move teeth and guide jaw development. It can improve bite, spacing, and cleaning access; treatment requires assessment and long-term retention.

\medterm{Orthodontic Wires} Metal or other wires used with braces to apply controlled forces that move teeth or help maintain their position.

\medterm{Orthodontics} The dental specialty that corrects teeth or jaws that do not align properly. Braces, clear aligners, or other appliances can improve biting, chewing, appearance, and cleaning; retainers are usually needed afterward because teeth can shift back.

\medterm{Orthodontist} A dentist specially trained to diagnose and treat teeth or jaws that do not align properly. Orthodontists use braces, clear aligners, retainers, and sometimes surgery to improve biting, chewing, appearance, and oral health.
\medterm{Orthognathic Surgery} Corrective surgery to reposition the upper jaw, lower jaw, or both when bone alignment causes problems with biting, chewing, speaking, breathing, or facial development. It is often combined with orthodontic treatment.

\medterm{Orthohepadnavirus} A group of viruses that includes hepatitis B virus and related viruses infecting mammals. These viruses carry genetic material as DNA and may cause long-lasting liver infection, inflammation, scarring, or cancer.

\medterm{Orthokeratosis} Thickening of the outer skin layer with fully mature skin cells. It can be a normal protective response, such as on the soles, or part of conditions causing thick, scaly, or rough skin.

\medterm{Orthomyxoviridae} A family of viruses that includes influenza A, B, and C. They mainly infect the airways and can cause flu, outbreaks, and sometimes severe pneumonia or other complications.
\medterm{Orthopaedic Surgery} Surgery for disorders and injuries involving bones, joints, muscles, ligaments, tendons, and related tissues. Procedures may repair fractures, replace joints, correct deformity, stabilize the spine, or relieve pressure on nerves.
\medterm{Orthopaedics} The British spelling of orthopedics, the medical and surgical specialty concerned with the musculoskeletal system.
\medterm{Orthopedic Biologics} Biological materials used to help bones, tendons, or other tissues heal. They include a person’s own bone, donated bone, processed tissue, and synthetic substitutes; benefits and risks vary by the injury and procedure.

\medterm{Orthopedic Implant} A device placed in or on the body to repair, support, replace, or stabilize a bone or joint. Examples include plates, screws, rods, joint replacements, and artificial ligaments. Possible complications include infection, loosening, breakage, wear, or bone damage.

\medterm{Orthopedic Infections} Infections involving bones, joints, muscles, or orthopedic implants. Examples include osteomyelitis, septic arthritis,
prosthetic-joint infection, and some diabetic foot infections.

\medterm{Orthopedic Rehabilitation} Treatment that restores movement, strength, daily function, and independence after injury or surgery involving bones, joints, muscles, ligaments, or tendons. It may include exercises, activity training, pain management, and assistive equipment.
\medterm{Orthopedic Screws} Metal or other screws used to hold broken bones or surgical implants in position while bone heals. Their shape and size are chosen for the bone and purpose. Complications can include loosening, breakage, infection, or injury to nearby structures.

\medterm{Orthopedic Services} Orthopedic services diagnose and treat problems of bones, joints, muscles, ligaments, tendons, and related movement systems.

\medterm{Orthopedics} The branch of surgery concerned with the musculoskeletal system, including bones,
joints, ligaments, tendons, muscles, and nerves.

\medterm{Orthopedist} A physician who diagnoses and treats problems of bones, joints, muscles, ligaments, tendons, and related tissues. Treatment may include rehabilitation, medicines, injections, or surgery.
\synonyms
Orthopedic surgeon

\medterm{Orthopnea} Shortness of breath that begins or worsens when lying flat and improves when sitting or standing. It often occurs with heart or lung disease, but other causes are possible; new or worsening symptoms need medical assessment.

\medterm{Orthopneic Position} Sitting upright, often leaning forward with the arms supported, to make breathing easier. A person who suddenly needs this position or cannot lie flat because of breathlessness should seek medical assessment.

\medterm{Orthopnoea} Shortness of breath that starts or worsens when lying flat and improves when sitting or standing. It often occurs with heart or lung disease; new or worsening difficulty lying down because of breathlessness needs medical assessment.


\medterm{Orthopoxvirus} A genus of poxviruses that includes variola, mpox, vaccinia, and cowpox viruses. Infection may cause fever and skin lesions; severity and transmission depend on the specific virus.
\medterm{Orthoptera} The insect order containing grasshoppers, crickets, katydids, and locusts; some species cause bites, allergies, or crop-related health effects.

\medterm{Orthoptic} Relating to assessment and treatment of binocular vision, eye alignment, and disorders such as strabismus.
\medterm{Orthoreovirus} A genus of nonenveloped, segmented double-stranded RNA viruses; human infection is often mild or asymptomatic but can occasionally cause disease.

\medterm{Orthorexia} A term used for an excessive or rigid preoccupation with eating foods considered healthy. It is not a universally recognized formal psychiatric diagnosis but may cause nutritional or social impairment.

\medterm{Orthoscopic} Orthoscopic means providing a straight or undistorted view, as in an optical instrument or examination.
\medterm{Orthosis} A brace, splint, shoe insert, or similar device worn outside the body to support, align, protect, or improve movement of a body part. Unlike a prosthesis, it does not replace a missing body part.

\medterm{Orthostatic Hypotension} A fall in blood pressure after standing, often causing dizziness, blurred vision, weakness, or fainting. It may result from dehydration, medicines, autonomic nerve problems, or prolonged bed rest. Repeated symptoms or falls need medical assessment.

\medterm{Orthostatic Intolerance} Dizziness, fast heartbeat, weakness, blurred vision, fatigue, or faintness that develops or worsens on standing because the body does not maintain blood flow to the brain adequately. Causes include dehydration, medicines, nerve disorders, and some chronic conditions.

\medterm{Orthotic} Orthotic means relating to a brace, splint, insert, or device that supports, aligns, protects, or improves function of a body part.

\medterm{Orthotic Devices} External supports that hold, align, protect, or assist a body part, such as a brace, splint, shoe insert, or spinal support. They can reduce pain and improve movement, but poor fit may cause pressure, numbness, skin injury, or restricted circulation.

\medterm{Orthotic Training} Instruction and practice for safely using a brace or other support. It includes putting it on and removing it, checking skin, cleaning it, walking with it, and recognizing problems needing review.

\medterm{Ortolani's Sign} A palpable or audible clunk produced when a dislocated but reducible infant hip is gently abducted and reduced; it is a screening sign of developmental dysplasia of the hip.

\medterm{Ortolani's Test} A hip examination maneuver in infants that detects a reducible dislocation by producing a palpable clunk as the femoral head enters the acetabulum.
\medterm{Os} A bone or bony structure; in Latin anatomical terminology it means bone, while in obstetrics it means an opening.
\medterm{Os Odontoideum} Os odontoideum is a separate ossicle near the second cervical vertebra’s odontoid process that can cause atlantoaxial instability or spinal-cord risk.

\medterm{Osborn Wave} A notch or extra wave on an electrocardiogram, classically seen when the body is dangerously cold. It can also occur in other conditions and must be interpreted with the person's temperature and other findings.

\synonyms
J wave

\medterm{Oscillation} Repeated movement or variation around a central position or value.
\medterm{Oscillator} A device or system that produces repeated movement or an electrical signal at regular intervals.
\medterm{Oscillograph} An instrument or computer display that records a changing signal over time. It can show electrical, sound, pressure, or movement waves and has been used in medical tests to study signals from the heart, brain, muscles, or breathing.
\medterm{Oscillometry} A method that measures pressure-related vibrations, commonly used by automated blood-pressure cuffs and some vascular devices. The result estimates blood pressure or blood-flow characteristics but may be inaccurate with movement or irregular rhythms.
\medterm{Oscillopsia} The disturbing illusion that the visual environment is moving or bouncing, often caused by vestibular or eye-movement dysfunction.
\medterm{Oscitation} Yawning, a reflex involving opening the mouth and deep inhalation that is associated with sleepiness, fatigue, or other physiologic states.

\medterm{Oseltamivir} Oseltamivir is an antiviral medicine used to treat or prevent influenza, with the greatest benefit when started early.

\medterm{Oseltamivir Phosphate} Oseltamivir phosphate is the salt form of oseltamivir, an antiviral medicine used for treatment or prevention of influenza.

\medterm{OSFED} A recognized eating disorder in which symptoms cause significant distress or interfere with health or daily life but do not meet every criterion for another named eating disorder. It still requires assessment and treatment.

\medterm{Osgood-Schlatter Disease} Painful irritation of the growing area where the kneecap tendon attaches to the shinbone. It commonly affects active children and teenagers and causes tenderness or swelling below the knee.

\medterm{Osler Nodes} Tender red or purple lumps on the pads of the fingers or toes. They can occur with infective endocarditis and some other conditions; they are an important clue but do not prove the diagnosis, which requires a full medical assessment.

\medterm{Osmidrosis} Strong, unpleasant body odor caused when skin bacteria break down sweat and skin oils. It commonly affects the armpits or feet and may improve with washing, dry clothing, antiperspirants, treatment of skin infection, or assessment of an underlying condition.
\medterm{Osmium} A dense metallic element; some osmium compounds are highly toxic and can damage the eyes and respiratory tract.
\medterm{Osmolality} A measure of how many dissolved particles are present in a kilogram of fluid. Blood osmolality depends mainly on sodium, glucose, and urea and helps determine how water moves between the bloodstream and cells.

\medterm{Osmolality Blood Test} A blood osmolality test measures the concentration of dissolved particles in blood and helps assess hydration, electrolyte disorders, and certain poisonings.

\medterm{Osmolality Urine Test} A urine osmolality test measures how concentrated the urine is and helps evaluate the kidneys' ability to regulate water balance.

\medterm{Osmolar} Describing a liquid according to the number of dissolved particles in each liter. This property affects water movement between body compartments and helps clinicians assess blood chemistry, hydration, kidney function, and intravenous fluids.

\medterm{Osmolarity} The concentration of osmotically active particles per liter of solution.
\medterm{Osmole} A unit representing the amount of a substance that contributes one mole of osmotically active particles.

\medterm{Osmoreceptor} A sensory cell or neural receptor that detects changes in osmotic concentration, especially in the hypothalamus.
\medterm{Osmoregulation} Control of water and dissolved-solute balance in the body or an organism, maintaining appropriate cell volume and internal concentration.

\medterm{Osmosis} Movement of water through a membrane toward the side with more dissolved particles. This process helps determine how water is distributed between blood, cells, and surrounding tissues.

\medterm{Osmotic Demyelination Syndrome} Serious brain injury that can occur when long-standing low blood sodium is corrected too quickly. Symptoms may appear days later and include confusion, difficulty speaking or swallowing, weakness, movement problems, or reduced consciousness.

\medterm{Osmotic Diarrhea} Diarrhea caused by poorly absorbed solutes retaining water in the intestinal lumen. It
typically improves with fasting and may result from carbohydrate malabsorption,
laxatives, or other osmotically active substances.

\medterm{Osmotic Diuresis} Increased urine production caused by poorly reabsorbed solutes in the kidney tubules, such as glucose during marked hyperglycemia.

\medterm{Osmotic Fragility Test} A test that measures the susceptibility of red blood cells to haemolysis when placed in
progressively more dilute saline solutions. Increased fragility supports hereditary
spherocytosis, although the result may also be affected by other haemolytic disorders
and recent transfusion.

\medterm{Osmotic Pressure} The force that draws water across a membrane toward the side containing more dissolved particles. It helps control how water moves between blood, cells, and surrounding tissues.

\medterm{Osseointegration} The firm attachment of living bone to the surface of an implant, such as a dental implant or artificial joint. Good stability, healthy bone, healing, and infection control support this process.

\medterm{Osseous} Relating to or composed of bone. For example, an osseous lesion is a bone lesion.

\medterm{Osseous Metaplasia} Osseous metaplasia is formation of mature bone in tissue where bone is not normally present, sometimes after injury, inflammation, or within a lesion.

\medterm{Ossicle} A small bone, especially one of the three tiny middle-ear bones. The malleus, incus, and stapes transmit eardrum vibrations to the inner ear, allowing sound to be heard.
\medterm{Ossicles} Three smallest body bones forming the ossicular chain in the middle ear: malleus
(hammer, attached to tympanic membrane), incus (anvil), stapes (stirrup, footplate in
oval window).

\medterm{Ossicular Discontinuity} Partial or complete separation of the middle-ear ossicles, usually caused by trauma, chronic ear disease, or surgery.
It produces conductive hearing loss and may be evaluated with otoscopy, audiometry, and imaging.

\synonyms
Ossicular-chain discontinuity; ossicular disruption.

\medterm{Ossicular Fusion} Abnormal joining or fixation of the middle-ear ossicles (malleus, incus, or stapes), which can restrict sound transmission and cause conductive hearing loss. Evaluation and treatment depend on the underlying congenital, inflammatory, traumatic, or postoperative cause.
\medterm{Ossification} The formation of bone. It occurs during growth and fracture healing and can also happen abnormally in soft tissues after severe injury, burns, or nervous-system damage.
\medterm{Ossifying Fibroma} An ossifying fibroma is a usually benign fibro-osseous tumor that replaces normal bone with fibrous tissue and mineralized material, often in the jaw.

\medterm{Osteitis} Inflammation of bone caused by infection, injury, inflammatory disease, or another process. It may cause localized pain, swelling, warmth, reduced movement, or fever and sometimes requires imaging, antibiotics, or surgery.
\medterm{Osteitis Deformans} Osteitis deformans is Paget disease of bone, in which abnormal remodeling produces enlarged, weakened, or deformed bones.

\medterm{Osteitis Deformans (Paget's Disease)} Alternate index label for Paget's Disease.

\medterm{Osteitis Fibrosa} Osteitis fibrosa is fibrous replacement and weakening of bone caused by prolonged excess parathyroid hormone, usually in severe hyperparathyroidism.

\medterm{Osteoarthritis} A joint disease in which cartilage and other tissues gradually break down or change. It can cause pain, stiffness, swelling, bone changes, and restricted movement. It usually develops over time but may occur earlier after joint injury, abnormal joint structure, or previous joint disease.

\synonyms
Arthritis Degenerative (Osteoarthritis)
Degenerative Arthritis

\medterm{Osteoarthritis (Geriatric)} Osteoarthritis occurring in later life, commonly affecting the knees, hips, hands, or spine
and causing joint pain, stiffness, and reduced movement. Management may include exercise,
weight management, medicines, assistive measures, or joint replacement when appropriate.

\synonyms
Osteoarthritis in older adults; age-related osteoarthritis

\medterm{Osteoarthritis, Cervical} Alternate index label; see \textit{Cervical spondylosis}.

\medterm{Osteoarthropathy} Osteoarthropathy is disease affecting bones and joints, often describing degenerative, hypertrophic, metabolic, or systemic-associated changes.
\medterm{Osteoarthrosis} Another name for osteoarthritis, a joint disorder in which cartilage and nearby bone gradually change, causing pain, stiffness, and reduced movement.
\medterm{Osteoblast} A cell that builds new bone. It makes the soft framework of bone and helps add calcium and phosphate to harden it. Some osteoblasts become osteocytes, which maintain mature bone.

\medterm{Osteoblastoma} A rare benign bone-forming tumor, usually in the spine or long bones, that can cause persistent pain and sometimes neurologic symptoms.

\medterm{Osteoblasts} Cells that build new bone. They produce the material that forms bone and help control its hardening with calcium and phosphate. Some osteoblasts become osteocytes, which help maintain mature bone.

\medterm{Osteocalcin} A protein made mainly by osteoblasts that participates in bone mineralization and serves as a marker of bone formation.

\medterm{Osteocalcin Assay} A blood test measuring osteocalcin, a protein made by bone-forming cells. It can provide information about bone formation and turnover, but results are affected by age, kidney function, medicines, and the laboratory method.
\medterm{Osteochondritis} Inflammation or injury involving bone and its overlying cartilage, often near a joint. The term includes conditions such as osteochondritis dissecans, in which a fragment may loosen.
\medterm{Osteochondritis Dissecans} An area of damaged bone and cartilage beneath a joint surface that may loosen or separate. It most often affects the knee in adolescents and young adults and can cause activity-related pain, swelling, catching, or locking. Treatment depends on lesion stability, age, and symptoms.

\medterm{Osteochondroma} A usually noncancerous cartilage-covered growth projecting from the surface of a bone, often near a growth plate. It may be single or occur as multiple inherited growths and can cause pain, deformity, or pressure on nearby nerves or vessels.

\medterm{Osteochondromatosis} Osteochondromatosis is formation of multiple cartilage-capped bony outgrowths, usually around growth plates, that can cause deformity, pain, or nerve and vessel compression.

\medterm{Osteochondrosis} A disorder of bone and cartilage development, often involving a growth area or joint surface. It may cause pain, swelling, stiffness, or altered movement and can affect children, adolescents, or adults.
\medterm{Osteoclasis} Deliberate surgical breaking or manipulation of a bone to correct a deformity or restore alignment. It requires stabilization while the bone heals and may cause bleeding, infection, or nerve injury.
\medterm{Osteoclast} A specialized bone cell that breaks down old bone during normal renewal and repair. Osteoclast activity helps regulate bone strength and blood calcium; excessive activity can contribute to osteoporosis and other diseases causing bone loss.
\medterm{Osteoclasts} Cells that break down old bone. This normal process helps bones grow, repair themselves, and release calcium into the blood. Too much osteoclast activity causes bone loss, as in osteoporosis, Paget disease, and some cancers that spread to bone. Medicines such as bisphosphonates and denosumab can slow osteoclast activity and reduce bone weakening.

\medterm{Osteocyte} A mature bone cell embedded in the hard bone it helped form. Osteocytes sense strain and chemical changes and send signals that help control bone strength, repair, and removal of old bone.

\medterm{Osteocytes} Mature bone cells trapped inside the hard bone they helped form. They sense strain and changes in the bone and send signals that help control bone repair, strengthening, and breakdown.

\medterm{Osteodystrophy} Abnormal bone growth, strength, or renewal caused by a body-wide or metabolic disorder. It may result from kidney disease, mineral imbalance, hormone problems, or inherited conditions and can cause pain or fractures.

\medterm{Osteogenesis} The formation and development of bone. It occurs in the growing skeleton and during fracture repair. Abnormal osteogenesis can contribute to birth disorders, poor healing, or fragile bones.
\medterm{Osteogenesis Imperfecta} A group of genetic disorders, also called brittle bone disease, in which bones break easily, sometimes after mild injury or for no clear reason. The disorders are often caused by gene changes that reduce or alter type I collagen, a protein that gives strength to bone and other connective tissues. Osteogenesis imperfecta is present from birth and can range from mild to severe; severe forms may cause fractures before birth. Other features may include short height, loose joints, blue or gray whites of the eyes, hearing loss, weak or discolored teeth, weak muscles, or a curved spine.

\medterm{Osteogenesis Imperfecta Brittle Bone Disease} Alternate index label; see \textit{Osteogenesis Imperfecta}.

\synonyms
Brittle Bone Disease (Osteogenesis Imperfecta Brittle Bone Disease)

\medterm{Osteogenic Sarcoma} A cancerous bone tumor whose cells make immature bone tissue. It often develops near the ends of growing long bones and may cause persistent pain, swelling, reduced movement, or a fracture.

\medterm{Osteoid} The newly formed, unmineralized organic framework of bone, produced mainly by osteoblasts. Minerals are later deposited in it to make mature hard bone.
\medterm{Osteoid Osteoma} A small noncancerous bone tumor, usually in an adolescent or young adult. It often causes persistent pain that is worse at night and may improve with anti-inflammatory medicine; scans help locate it and guide treatment.

\medterm{Osteology} The study of bones, including their structure, development, function, injuries, and diseases. It supports the assessment and treatment of fractures, bone deformities, weak bones, infections, and bone tumors.
\medterm{Osteolysis} Osteolysis is progressive destruction or resorption of bone caused by tumors, inflammation, prosthesis reactions, infection, metabolic disease, or other processes.

\medterm{Osteolytic} Causing or characterized by breakdown and loss of bone tissue. Osteolytic lesions can weaken bones and occur with cancers, infections, inflammation, metabolic disorders, or excessive activity of bone-removing cells.

\medterm{Osteoma} A benign, slow-growing tumor made of bone. It often develops in the skull or facial bones and may cause no symptoms unless it blocks a passage or presses on nearby structures.
\medterm{Osteoma Cutis} Osteoma cutis is formation of bone within the skin or subcutaneous tissue, occurring as a primary disorder or secondary to acne, trauma, or another condition.

\medterm{Osteomalacia} Osteomalacia is inadequate mineralization of adult bone, commonly from vitamin-D deficiency or phosphate disorders, causing bone pain, weakness, and fractures.

\medterm{Osteomalacia (Rickets)} Alternate index label for Rickets.

\medterm{Osteomyelitis} Infection and inflammation of bone, usually caused by bacteria and sometimes fungi or other organisms. It can cause pain, fever, swelling, reduced movement, bone damage, or a persistent draining wound.


\medterm{Osteomyelitis In Children} Childhood osteomyelitis is infection of bone, often caused by bacteria entering through blood or trauma, producing pain, fever, swelling, and impaired movement.

\medterm{Osteon} A tube-shaped structural unit of hard, compact bone arranged around a small channel containing blood vessels and nerves. Osteons help give bone strength while allowing it to receive nutrients.

\medterm{Osteonecrosis} Osteonecrosis is death of bone tissue from inadequate blood supply, causing pain, collapse, and joint damage, commonly in the hip or knee.

\medterm{Osteonecrosis (Avascular Necrosis)} Death of bone tissue after its blood supply is reduced, most often affecting the hip, knee, shoulder, or ankle. It can cause pain and collapse of the joint and may follow injury, prolonged steroid use, heavy alcohol use, blood disorders, or other conditions.

\synonyms
Avascular Necrosis

\medterm{Osteonecrosis Of The Jaw} Death of jawbone tissue, often associated with antiresorptive or antiangiogenic medications, radiation, infection, or trauma, and characterized by exposed bone or a nonhealing oral wound.

\medterm{Osteonectin} An extracellular-matrix protein involved in collagen organization, mineral binding, and tissue remodeling; it is also called SPARC.

\medterm{Osteons} Cylindrical structural units of compact bone arranged around central canals.

\synonyms
Haversian systems

\medterm{Osteopath} A physician trained in medicine and surgery who earns a Doctor of Osteopathic Medicine degree.
\synonyms
DO

\medterm{Osteopathy} A term used differently in different countries. It may refer to physician-led osteopathic medicine or to manual treatment focused on the muscles and joints; the practitioner’s training and legal scope should be checked locally.

\medterm{Osteopenia} Bone density that is lower than average but not low enough to meet the definition of osteoporosis. It may increase fracture risk, especially with age, previous fractures, falls, medicines, or other conditions that weaken bone.

\medterm{Osteopenia - Premature Infants} Osteopenia of prematurity is reduced bone mineralization in premature infants due to immature mineral stores, limited intake, illness, or prolonged intensive care.

\medterm{Osteopetrosis} A rare inherited disorder in which old bone is not removed normally. Bones become unusually dense but brittle and may fracture; crowding of the bone marrow and nerves can cause anemia, infections, vision loss, or other problems.

\medterm{Osteophyte} A bony outgrowth forming at the edge of a joint, commonly with osteoarthritis. It may cause stiffness, pain, reduced movement, or pressure on nearby nerves, although some osteophytes cause no symptoms.

\medterm{Osteoplasty} Surgery to repair, reshape, or reconstruct bone to restore structure, movement, or function. The technique depends on the bone problem and may be combined with fixation, grafting, or joint surgery.

\medterm{Osteopoikilosis} Osteopoikilosis is an inherited benign bone condition with multiple small dense spots on imaging, usually asymptomatic and discovered incidentally.

\medterm{Osteopontin} A protein involved in bone remodeling, immune responses, inflammation, and attachment between cells and tissues. Its level may change in several diseases, but a measurement is not by itself a specific diagnosis.
\medterm{Osteoporosis} A bone disease in which bone is broken down faster than it is replaced, or its structure and strength change. The bones become less dense and more likely to break, sometimes after a minor fall or other low-impact force. Osteoporosis often causes no symptoms until a fracture occurs. Fractures commonly affect the hip, spine, and wrist; weakened bones in the spine may collapse and cause back pain or loss of height. Osteoporosis can affect anyone, but it is more common in females, especially after menopause, when lower estrogen levels can speed bone loss. It can also affect males. Risk is influenced by age, family history, hormones, inadequate calcium or vitamin D, some medicines, smoking, alcohol, and other health conditions.

\medterm{Osteoporosis Fracture Risk Assessment} An evaluation of future fracture risk using bone density, age, previous fractures, falls, medicines, family history, and other health factors. The result helps determine whether lifestyle measures, monitoring, or bone-strengthening treatment is appropriate.

\medterm{Osteoporotic Fracture} A fracture resulting from weakened, porous bone after a low-energy event that would not usually break healthy bone.

\medterm{Osteoradionecrosis} Death of bone within an area treated with radiation, most often the jaw after head-and-neck cancer treatment. The bone may become exposed or fail to heal after injury or dental work and can become infected.
\medterm{Osteosarcoma} An aggressive cancer that begins in bone-forming cells and produces abnormal immature bone. It often develops near the ends of long bones in children and young adults and may cause persistent pain, swelling, or a fracture.

\medterm{Osteosarcoma, Secondary} An osteogenic sarcoma arising in previously abnormal or treated bone, such as bone affected by Paget disease, prior radiation, or chronic bone lesions; prognosis and treatment depend on stage and resectability.
\medterm{Osteosclerosis} Abnormally increased bone density seen on imaging or examination. It may result from inherited disease, altered bone renewal, chronic stress, inflammation, medicines, or cancer; dense bone can still be weak or painful.

\medterm{Osteosynthesis} Surgical repair and stabilization of a broken bone using implants such as plates, screws, rods, nails, or wires. The implants hold the pieces in alignment while the bone heals.

\medterm{Osteotome} A sharp surgical instrument used to cut or shape bone.

\medterm{Osteotomy} A surgical procedure involving cutting and repositioning of bone to correct deformity,
realign joint surfaces, alter weight-bearing forces, or change limb length.

\medterm{Osteotomy Of The Knee} Surgery that cuts and reshapes a knee-region bone to correct alignment and redistribute weight across the joint.

\medterm{Ostium} A natural opening into or out of an organ, cavity, blood vessel, or duct. It allows air, fluid, or other contents to pass between connected body spaces.
\medterm{Ostomy} A surgically created opening that connects part of the bowel or urinary tract to the skin so stool or urine can leave the body. A colostomy or ileostomy passes stool; a urostomy passes urine. A pouch usually collects the output and protects the surrounding skin.
\medterm{Ostomy Appliance} A pouching system that collects stool or urine from a surgically created opening in the abdomen. Colostomy and ileostomy appliances collect stool; urostomy appliances collect urine. The skin around the opening needs regular protection and monitoring.

\medterm{Ostraceous Psoriasis} A rare severe form of psoriasis in which plaques become very thick, hard, and layered, resembling an oyster shell. It can be difficult to treat and may require specialist care and medicines affecting the whole immune system.

\medterm{Otalgia} Pain in or around the ear. It may come from an ear infection, wax, injury, or pressure problem, or may be referred from the teeth, jaw, throat, or neck. Persistent pain, discharge, hearing loss, fever, or severe pain needs medical assessment.
\medterm{Otalgia (Earache)} Pain in or around the ear, arising from the ear itself or referred from structures such as the teeth, jaw, throat, or cervical spine.

\synonyms
Earache

\medterm{Othello Syndrome} A delusional belief that a partner is unfaithful despite insufficient evidence; it may occur with psychiatric, neurologic, or substance-related disorders.


\medterm{Other Pain Conditions} Pain conditions that do not fit neatly into a single category, including nociplastic pain, fibromyalgia, chronic
tension-type headache, temporomandibular pain, chronic pelvic pain, and complex regional
pain syndrome. Diagnosis and treatment depend on the specific condition.

\medterm{Other Race} Other race is a demographic category for a person who does not select or fit the available listed race categories; classification varies by system and should not be treated as a biological diagnosis.

\medterm{Otic Barotrauma} Injury to the ear caused by failure to equalize pressure during flying, diving, or other pressure changes.
\medterm{Otic Capsule} The dense part of the skull bone that surrounds and protects the inner ear, including the hearing and balance organs. Disease or abnormal remodeling in this bone can affect hearing or balance.

\medterm{Otic Ointment} A semisolid medicine formulated for application to the outer ear or ear canal; it should be used only as directed, especially when the eardrum may be perforated.

\medterm{Otitis} Inflammation of the ear, which may be infectious or noninfectious. Otitis externa affects the ear canal, otitis media affects the middle ear, and inner-ear inflammation may cause vertigo or hearing loss. Symptoms can include pain, fever, blockage, discharge, or reduced hearing and vary with the location and cause.
\medterm{Otitis (Ear Infection)} Alternate index label for Otitis.

\medterm{Otitis Externa} Inflammation or infection of the outer ear canal, commonly causing ear pain, itching, tenderness, swelling, or discharge. Keeping water and objects out of the ear can help while it heals.
\synonyms
Swimmers Ear (Otitis Externa)

\medterm{Otitis Media} Inflammation or infection in the middle ear, often after a cold or other upper-airway infection. It may cause ear pain, fever, reduced hearing, irritability, or fluid behind the eardrum.
\medterm{Otitis Media With Effusion} Fluid behind the eardrum without the severe pain, fever, or pus of an acute ear infection. It may cause blocked hearing or speech and learning problems in children and often clears, but persistent fluid needs hearing assessment.

\medterm{Oto-Rhino-Laryngology} Another name for otolaryngology, the specialty of the ear, nose, throat, head, and neck.
\medterm{Otoacoustic Emissions} Very soft sounds naturally produced by a healthy inner ear that can be measured with a small probe placed in the ear canal. The test is quick and painless and is commonly used to screen newborn hearing and assess inner-ear function.

\synonyms
OAE; otoacoustic-emission testing; cochlear emissions.

\medterm{Otocephaly} A severe congenital craniofacial malformation characterized by absent or markedly
underdeveloped mandible, ventral displacement or fusion of the ears, and associated
facial and airway abnormalities. It is usually lethal because of profound impairment of
breathing and feeding.

\medterm{Otolaryngologist} An otolaryngologist is a physician specializing in disorders of the ear, nose, throat, head, neck, voice, swallowing, and related surgery.

\medterm{Otolaryngology} The medical specialty treating disorders of the ears, nose, throat, head, and neck. Specialists manage hearing, balance, breathing, swallowing, voice, sinus, facial, and related surgical problems.
\medterm{Otolith} A tiny calcium-containing crystal in the inner ear that helps sense gravity and head movement. If it moves into the wrong balance canal, it can cause brief spinning attacks known as positional vertigo.

\medterm{Otoliths} Tiny calcium-carbonate crystals in inner-ear organs that sense gravity and straight-line movement. If they move into the wrong balance canal, they can cause brief spinning attacks called positional vertigo.
\medterm{Otology} The study and medical care of the ear and hearing.
\medterm{Otoplasty} Surgery to alter the shape, position, or proportion of the external ear, commonly to correct prominent or congenital ear differences.

\medterm{Otorhinolaryngology} The medical specialty concerned with diseases of the ear, nose, throat, head, and neck; also called otolaryngology.

\synonyms
otolaryngology.

\medterm{Otorrhea} Fluid, pus, blood, or other discharge flowing from the ear. Causes include infection, eardrum perforation, injury, or a growth; persistent, bloody, foul-smelling, or clear discharge needs medical assessment.

\medterm{Otorrhoea} Fluid, pus, or blood draining from the ear. Common causes include an ear-canal infection, a middle-ear infection with a perforated eardrum, injury, or a growth. Clear watery discharge after a head injury is an emergency because it may be brain fluid and can lead to meningitis.
\medterm{Otosclerosis} An abnormal hardening and fixation of middle-ear bones that can gradually reduce hearing by blocking sound transmission.
\medterm{Otoscope} A lighted instrument used to examine the ear canal and eardrum. It can help identify earwax, infection, a foreign object, fluid, swelling, or a perforated eardrum. Examination findings are interpreted with symptoms and, when needed, hearing tests.
\medterm{Otoscopy} Looking into the ear canal and at the eardrum with a lighted instrument called an otoscope. It can help find earwax, infection, a foreign object, a hole in the eardrum, fluid, or other causes of ear symptoms.
\medterm{Ototoxicity} Damage to the inner ear from a medicine, chemical, or other toxic exposure. It may cause hearing loss, ringing in the ears, dizziness, or imbalance and can be temporary or permanent.

\synonyms
Auditory toxicity; vestibular toxicity; ototoxic effect.

\medterm{Ounce} An ounce is a unit of mass or fluid volume whose exact conversion depends on the measurement system; medication instructions should use unambiguous units.

\medterm{Outcome} The result or final effect of a disease, treatment, exposure, or healthcare intervention. Outcomes may include recovery, symptom control, complications, disability, quality of life, or death and can be measured in different ways.
\medterm{Outcome Measurement} The use of defined measures to assess a patient’s symptoms, function, quality of life, or other results before and after care.



\medterm{Outer Ear} The visible auricle and the external auditory canal, which collect and conduct sound toward the eardrum.

\medterm{Outer Ear Infection} Alternate index label; see \textit{Swimmer's ear}.

\medterm{Outpatient} A person who receives medical evaluation or treatment without being admitted overnight to a hospital or other inpatient facility.
\medterm{Outpatient Rehabilitation} Rehabilitation provided without hospital admission to improve function after illness, injury,
surgery, or disability. The setting, frequency, and duration are individualized.

\medterm{Output} The amount of a substance produced or eliminated by the body or by an organ over a specified time, such as urine output or cardiac output.
\medterm{Ova} Plural of ovum; eggs or egg cells, and in parasitology, eggs of worms or other parasites.

\medterm{Ovalocyte Count} An ovalocyte count estimates the number or proportion of oval-shaped red blood cells on a blood smear and may help evaluate anemia.

\medterm{Ovarian} Relating to an ovary. The ovaries release eggs and produce hormones. Ovarian problems include cysts, inflammation, reduced function, twisting that cuts off blood supply, and tumors; sudden severe pelvic pain may indicate twisting and needs emergency care.
\medterm{Ovarian Cancer} Cancer arising in an ovary or nearby reproductive tissue. It may cause persistent bloating, pelvic or abdominal pain, early fullness, urinary changes, or abnormal bleeding, but early disease may have few symptoms.

\synonyms
Cancer, Ovarian

\medterm{Ovarian Cyst} Alternate index label for Ovarian Cysts.

\medterm{Ovarian Cysts} Fluid-filled or partly solid sacs in or on an ovary. Many are related to the menstrual cycle and resolve spontaneously, while others may be associated with endometriosis, benign tumors, or hormonal conditions. Ovarian cysts may cause pelvic pain, menstrual changes, pressure, bleeding, rupture, or twisting of the ovary; some cause no symptoms.

\synonyms
Cyst, Ovarian
Cysts Ovary (Ovarian Cysts)

\medterm{Ovarian Failure} Reduced ovarian activity causing irregular or absent periods, difficulty becoming pregnant, and symptoms of low estrogen such as hot flushes or vaginal dryness. When it occurs before age 40, it is usually called primary ovarian insufficiency and needs evaluation.

\medterm{Ovarian Fibromata} Usually noncancerous tumors made of fibrous tissue that develop in an ovary. They may cause no symptoms or produce pelvic pressure, pain, menstrual changes, or fluid in the abdomen when large.

\medterm{Ovarian Follicle} A small fluid-filled structure in an ovary that contains an immature egg and supporting cells. During each menstrual cycle, one follicle usually matures and releases an egg during ovulation.

\medterm{Ovarian Gonadoblastoma} A rare tumor containing cells that normally develop into eggs and cells that support sex-hormone production. It often occurs in an atypically developed gonad, may contain Y-chromosome material, and can produce hormones.

\medterm{Ovarian Hemorrhage} Bleeding into or around an ovary, often from a ruptured cyst or blood vessel. It may cause sudden pelvic pain, weakness, dizziness, or internal blood loss and can require urgent assessment or surgery.

\medterm{Ovarian Hyperstimulation Syndrome} A complication of fertility treatment in which the ovaries enlarge and fluid moves into the abdomen or other tissues. It may cause bloating, pain, vomiting, reduced urine, breathlessness, blood clots, or kidney problems and can be severe.

\synonyms
OHSS

\medterm{Ovarian Infection} Infection involving an ovary, often as part of pelvic inflammatory disease or an abscess, causing pelvic pain, fever, and abnormal discharge.

\medterm{Ovarian Metastasis} Cancer that has spread to an ovary from another organ, commonly the digestive tract, breast, or reproductive organs. It may resemble a primary ovarian cancer, so tissue testing is important because treatment depends on the original site.

\medterm{Ovarian Overproduction Of Androgens} Ovarian androgen overproduction is excessive production of male-pattern hormones by the ovaries, which may cause acne, excess facial hair, irregular periods, or infertility.

\medterm{Ovarian Pregnancy} A rare ectopic pregnancy that implants in an ovary rather than inside the uterus or fallopian tube. It may cause one-sided pelvic pain or internal bleeding and usually needs urgent ultrasound assessment and treatment.

\medterm{Ovarian Rupture} Ovarian rupture is tearing of an ovary or ovarian cyst, which may cause sudden pelvic pain and internal bleeding and sometimes requires urgent treatment.

\medterm{Ovarian Teratoma} A tumor of an ovary made from reproductive cells that can form different kinds of tissue, such as skin, hair, or bone. Mature teratomas are usually noncancerous; immature teratomas can be cancerous.

\medterm{Ovarian Thecoma} An ovarian thecoma is a usually benign sex-cord stromal tumor that may produce estrogen and cause abnormal uterine bleeding or other hormonal effects.

\medterm{Ovarian Torsion} Twisting of an ovary, often together with its fallopian tube, that can cut off blood flow. It causes sudden one-sided pelvic pain, often with nausea or vomiting, and is a gynecologic emergency requiring urgent assessment and usually surgery.

\medterm{Ovariectomy (Oophorectomy)} Surgery to remove one or both ovaries. It may treat ovarian disease or reduce cancer risk. Removing both ovaries before natural menopause causes permanent loss of ovarian hormone production.
\medterm{Ovaries} The paired female gonads that produce oocytes and secrete sex hormones, mainly estrogen and progesterone. They are located on either side of the uterus and participate in the menstrual cycle and reproduction.

\synonyms
ovary

\medterm{Ovariotomy} An older term for surgical removal or incision of an ovary.
\medterm{Ovary} A reproductive organ that releases eggs and produces hormones such as estrogen and progesterone. The ovaries help control menstrual cycles, ovulation, pregnancy preparation, and many body changes during reproductive life.
\medterm{Ovary Disease} Any disorder affecting an ovary, including cysts, infection, inflammation, tumors, hormonal dysfunction, torsion, or bleeding.
\medterm{Ovary Neoplasm} An abnormal growth in an ovary that may be noncancerous, have uncertain behavior, or be cancerous. Symptoms can include bloating, pelvic pain, pressure, or abnormal bleeding, but some growths are found incidentally on imaging.

\medterm{Oven Cleaner Poisoning} Oven cleaners can severely burn the mouth, throat, stomach, skin, or eyes because they may contain strong alkalis. Do not induce vomiting; rinse exposed skin or eyes with water and contact emergency services or poison control.

\medterm{Over-Hydration} Excess accumulation of body water, which may cause edema, dilutional hyponatremia, or heart and lung complications.
\medterm{Over-Nutrition} Excessive intake or accumulation of energy or nutrients that can contribute to obesity, metabolic disease, or toxicity.
\medterm{Over-The-Counter} A medicine available without a prescription for self-treatment of specified common problems. It can still cause harm, interact with prescribed medicines, or be unsafe with certain illnesses, so the label and dose should be followed carefully.

\medterm{Over-The-Counter Medications} Medicines sold without a prescription for selected common symptoms or conditions. They are not risk-free: incorrect use, prolonged use, interactions, allergies, or underlying illness can cause serious harm.

\medterm{Over-The-Counter Medicines} Over-the-counter medicines can be bought without a prescription for common problems, but they may still cause side effects or interact with prescribed medicines.

\medterm{Over-The-Counter Pain Relievers} Over-the-counter pain relievers include medicines such as acetaminophen and nonsteroidal anti-inflammatory drugs; choice depends on age, health conditions, and other medicines.

\medterm{Overactive Bladder} A pattern of a sudden, difficult-to-delay need to urinate, usually with frequent urination and waking at night, with or without leakage. Causes include bladder irritation, nerve problems, obstruction, medicines, and sometimes no identifiable cause.

\synonyms
OAB
OAB (Overactive Bladder)

\medterm{Overall Survival} The length of time from diagnosis, treatment, or study entry until death from any cause. Doctors and researchers use it to describe how long people live and to compare treatments. It differs from time without worsening disease because it includes survival regardless of whether the illness has progressed.

\medterm{Overbite} Vertical overlap of the upper front teeth over the lower front teeth; excessive overbite can damage gums, teeth, or jaw function.

\medterm{Overcoming Breastfeeding Problems} Breastfeeding problems such as pain, poor latch, low supply, engorgement, or blocked ducts can often improve with feeding assessment, positioning help, and timely lactation or medical support.

\medterm{Overdenture} A removable denture supported by dental implants or remaining natural teeth. This support can improve stability, chewing, speech, and comfort, while attachments may help preserve jawbone and reduce denture movement.

\medterm{Overdose} An overdose is exposure to a medicine, drug, or substance in an amount capable of causing toxicity, with effects determined by dose, formulation, route, timing, and patient factors.
\medterm{Overflow Urinary Incontinence} Urine leakage caused by a bladder that remains too full because it does not empty properly. Weak flow, difficulty starting, frequent small amounts, dribbling, or constant leakage may occur; blockage or nerve disease can be responsible.

\medterm{Overgrowth Syndrome} A group of conditions present from birth in which the whole body, one side, or particular tissues grow more than expected. Some forms also cause developmental differences or increase the risk of certain tumors.

\medterm{Overlap Myositis} Muscle inflammation occurring together with another autoimmune connective-tissue disease, such as lupus, systemic sclerosis, mixed connective-tissue disease, or Sjögren syndrome. It may cause muscle weakness, skin changes, joint symptoms, swallowing problems, or lung disease; treatment is tailored to the affected organs.

\medterm{Overlay Denture} A removable denture that fits over remaining teeth, roots, or implants to improve support, chewing, and retention.

\medterm{Overload Iron (Iron Overload)} Alternate index label for Iron Overload.

\medterm{Overriding Aorta} A birth-related heart abnormality in which the main artery leaving the heart sits over a hole between the lower chambers and receives blood from both. It commonly occurs as part of tetralogy of Fallot.
\medterm{Oversedate} To give excessive sedation, causing unusually deep sleep, slowed responses, impaired breathing, or difficulty maintaining the airway.

\medterm{Overt Proteinuria} Abnormally high protein loss in the urine that is detectable by routine testing and may indicate kidney disease.

\medterm{Overuse Syndrome (Repetitive Motion Disorders (RMDs))} Alternate index label for Repetitive Motion Disorders (RMDs).

\medterm{Overweight} Body weight above the usual range for height. In adults, a Body Mass Index of 25 to less than 30 kg/m\(^2\) is commonly classified as overweight, but it does not directly measure body fat.
\medterm{Oviduct} A tube carrying an ovum from the ovary toward the uterus; in humans it is the uterine or fallopian tube.
\medterm{Ovomucin} A glycoprotein in egg white that contributes to its gel-like viscosity and helps stabilize the egg-white structure.

\medterm{Ovotestis} A gonad containing both ovarian-type and testicular tissue. It is an uncommon difference of sex development and may produce ovarian hormones, testicular hormones, eggs, sperm, or none of these; evaluation is individualized.
\medterm{Ovulation} Release of a mature egg from an ovarian follicle, usually around the middle of the menstrual cycle.
\medterm{Ovulation Home Test} An ovulation home test detects the rise in urinary luteinizing hormone that usually occurs shortly before ovulation.

\medterm{Ovum} The female reproductive cell released from an ovary during ovulation. If a sperm enters it, the cell completes fertilization and joins with the sperm’s genetic material to form the first cell of an embryo.
\medterm{Oxacillin} A penicillinase-resistant penicillin used mainly to test for methicillin resistance in staphylococci; it is not commonly used as the preferred treatment for most infections.

\medterm{Oxalate} A salt or ester of oxalic acid that can form kidney stones and binds calcium.
\medterm{Oxalates} Oxalates are salts or esters of oxalic acid that can bind calcium and form calcium-oxalate kidney stones; excess dietary absorption or inherited hyperoxaluria can cause renal injury.

\medterm{Oxalic Acid} A dicarboxylic acid whose oxalate salts can contribute to calcium-oxalate kidney stones; concentrated exposure is toxic.
\medterm{Oxalic Acid Poisoning} Swallowing oxalic acid can burn the digestive tract and bind calcium, causing muscle cramps, abnormal heart rhythms, seizures, kidney injury, or shock. It is a medical emergency requiring immediate poison-control or emergency care.

\medterm{Oxaliplatin} Oxaliplatin is a platinum chemotherapy drug that forms DNA cross-links and is used especially in colorectal cancer regimens; acute cold-triggered neuropathy and cumulative sensory neuropathy are characteristic toxicities.

\medterm{Oxaloacetate} A small molecule used in the cell’s energy-making cycle. It combines with another molecule to begin a step that produces energy and can also help make glucose and the amino acid aspartate.

\medterm{Oxandrolone} An anabolic steroid used in selected conditions involving severe weight loss or catabolic stress; it can affect the liver, lipids, and hormone function.

\medterm{Oxaprozin} A nonsteroidal anti-inflammatory drug used to relieve pain and inflammation in conditions such as osteoarthritis and rheumatoid arthritis.

\medterm{Oxazepam} A benzodiazepine medicine used short-term for anxiety or alcohol withdrawal, causing sleepiness and possible dependence.
\medterm{Oxazepam Overdose} Taking too much oxazepam, a sedative medicine, can cause severe sleepiness, confusion, poor coordination, slowed breathing, coma, or death, especially with alcohol or other sedatives. Immediate emergency advice is needed.

\medterm{Oxazole} Oxazole is a five-membered chemical ring containing oxygen and nitrogen that is used as a building block in medicinal and other chemistry.

\medterm{Oxcarbazepine} Oxcarbazepine is an antiseizure medicine that blocks voltage-gated sodium channels and treats focal seizures; dizziness, ataxia, rash, and clinically important hyponatraemia may occur.

\medterm{Oxidant} A substance that can damage cells by taking electrons from other molecules. Oxidants are produced normally during metabolism and inflammation, but excessive amounts can cause oxidative stress and tissue injury.

\synonyms
Oxidizing agent

\medterm{Oxidants} Substances or reactive molecules that accept electrons or promote oxidation; excessive oxidant activity can damage proteins, lipids, and DNA.

\medterm{Oxidation} A chemical process that changes a substance through loss of electrons or reaction with oxygen.
\medterm{Oxidation-Reduction} A chemical reaction involving transfer of electrons, in which one substance is oxidized and another is reduced.

\medterm{Oxidative Phosphorylation} The final stage of energy production in mitochondria, where electrons from food drive the movement of hydrogen ions and power the enzyme that makes ATP, the main energy-carrying molecule in cells.

\medterm{Oxidative Stress} A harmful imbalance between reactive oxygen-containing molecules and the body’s protective antioxidants. Excessive oxidative stress can injure cell membranes, proteins, and DNA and may contribute to inflammation, aging, or disease, although the term alone does not identify a specific illness.
\medterm{Oxidoreductase} An enzyme that transfers electrons or hydrogen between molecules during a chemical reaction. These enzymes are important in energy production, drug breakdown, and many other cell processes.

\medterm{Oxime} A type of chemical compound. Some oximes, such as pralidoxime, are used as emergency antidotes because they can restore an important nerve enzyme after organophosphate pesticide poisoning.

\medterm{Oximeter} A device that estimates blood oxygen level, usually by shining light through a finger, toe, or earlobe. A pulse oximeter also displays pulse rate, but readings can be affected by movement, poor circulation, or skin factors.
\medterm{Oximetry} Measurement of oxygen levels in blood or tissue. A pulse oximeter estimates blood oxygen through a skin sensor, while arterial testing measures oxygen more directly; results can be affected by poor circulation or motion.


\medterm{Oxter} Oxter is an anatomical term for the armpit or axilla.

\medterm{Oxybuprocaine} A local anaesthetic eye drop that briefly numbs the surface of the eye. It may be used for examinations or short procedures and should not be used repeatedly without supervision because it can delay healing and damage the cornea.
\medterm{Oxycephaly} A cone-shaped skull caused by premature fusion of cranial sutures, also called acrocephaly.
\medterm{Oxycodone} Oxycodone is an opioid analgesic used for moderate-to-severe pain; sedation, constipation, respiratory depression, tolerance, dependence, and overdose risk increase with alcohol or other sedatives.

\medterm{Oxycodone Hydrochloride} The hydrochloride salt of oxycodone, an opioid analgesic used for moderate-to-severe pain; it can cause sedation, constipation, dependence, and fatal respiratory depression in overdose.

\medterm{Oxygen} A gas essential for aerobic metabolism and tissue energy production.
\medterm{Oxygen Consumption} The amount of oxygen the body uses in a given time. It reflects the body’s energy needs and rises with exercise, fever, or illness; it can be measured during exercise testing to assess heart and lung function.

\medterm{Oxygen Content} The amount of oxygen carried in blood, mainly bound to hemoglobin with a small dissolved component.
\medterm{Oxygen Delivery} The amount of oxygen transported from the lungs to body tissues each minute. It depends mainly on cardiac output, the amount of hemoglobin, and how much oxygen the hemoglobin carries. Inadequate delivery can injure organs when the body cannot compensate.

\medterm{Oxygen Desaturation Index} The average number of oxygen desaturation events per hour of sleep, usually reported from a sleep study as the ODI; interpretation depends on the recording method and clinical context.

\medterm{Oxygen Enhancement Ratio} A measure of how much oxygen increases the effect of radiation on cells. It compares the radiation dose needed with little oxygen to the dose needed when oxygen is normal to produce the same amount of cell injury.
\medterm{Oxygen Mask} A mask that delivers supplemental oxygen through the nose and mouth. The type and flow rate are chosen according to oxygen needs; oxygen should be used as prescribed and monitored.

\medterm{Oxygen Pulse} The amount of oxygen used by the body with each heartbeat, estimated during exercise as oxygen uptake divided by heart rate.

\medterm{Oxygen Safety} Practices that reduce fire, explosion, pressure, and treatment risks during oxygen use, including keeping oxygen away from flames, smoking, oils, and unsafe equipment.

\medterm{Oxygen Saturation} The percentage of hemoglobin in the blood carrying oxygen. A finger or ear sensor estimates it without taking blood, while an arterial blood test measures it more directly; the healthy target varies with the person’s condition.

\medterm{Oxygen Tent} A canopy or enclosure used to deliver oxygen around a patient’s head; modern oxygen masks, nasal cannulas, or ventilatory devices are usually more precise.

\medterm{Oxygen Therapy} Giving extra oxygen through a tube, mask, or breathing machine when the blood oxygen level is too low or is expected to fall. The amount is prescribed and monitored because too much oxygen can also be harmful; oxygen must be kept away from flames.
\medterm{Oxygen Therapy In Infants} Carefully measured extra oxygen given to a baby whose blood oxygen is too low or whose lungs are not exchanging gases adequately. Clinicians monitor oxygen closely because both too little and too much can injure developing organs.

\medterm{Oxygen Toxicity} Tissue injury caused by exposure to oxygen at excessively high partial pressures or concentrations, potentially affecting the lungs, eyes, or central nervous system.

\medterm{Oxygen Transport} The movement of oxygen from the lungs to tissues, mainly bound to hemoglobin and partly dissolved in plasma, with delivery affected by blood flow and hemoglobin concentration.

\medterm{Oxygen-Hemoglobin Dissociation Curve} A graph showing how strongly hemoglobin holds oxygen at different oxygen levels. Temperature, acidity, and carbon dioxide change this binding and therefore affect how readily oxygen is released to tissues.

\medterm{Oxygenases} Enzymes that incorporate oxygen from molecular oxygen into a substrate; monooxygenases add one oxygen atom and dioxygenases add both.

\medterm{Oxygenated Blood} Blood that has picked up oxygen in the lungs and carries it through arteries to body tissues. It is usually bright red, while blood returning to the lungs contains less oxygen and more carbon dioxide.

\medterm{Oxygenation} The process of supplying oxygen to blood and tissues, or the degree to which hemoglobin and tissues are adequately oxygenated.

\medterm{Oxygenator} A machine that adds oxygen to blood and removes carbon dioxide when the lungs cannot do this adequately. It is used during heart-lung bypass surgery and some forms of extracorporeal life support.

\medterm{Oxyhaemoglobin} Hemoglobin with oxygen attached to it. It forms mainly in the lungs and carries oxygen through the bloodstream to body tissues.
\medterm{Oxyhemoglobin} Hemoglobin bound to oxygen, formed mainly in the lungs and carrying oxygen through arterial blood to tissues.

\medterm{Oxylipins} Oxylipins are oxygenated fatty-acid derivatives that act as local mediators of inflammation, vascular tone, pain, and other physiologic responses.

\medterm{Oxymetazoline} A topical alpha-adrenergic agonist used as a nasal decongestant or ocular vasoconstrictor; prolonged nasal use can cause rebound congestion.
\medterm{Oxymetholone} Oxymetholone is an anabolic steroid used in selected severe anemias; it can cause liver injury, fluid retention, and hormonal side effects.

\medterm{Oxymorphone} A potent opioid analgesic used for severe pain; it can cause sedation, constipation, dependence, and life-threatening respiratory depression.

\medterm{Oxyphil Cell} A cell with many mitochondria, giving it a distinctive granular appearance under a microscope. Oxyphil cells occur in tissues such as the parathyroid and thyroid, where their significance depends on the organ and the surrounding findings.

\medterm{Oxyphil Cell Carcinoma} Alternate index label; see \textit{Hurthle cell cancer}.

\medterm{Oxyphilic Adenoma} An oxyphilic adenoma is a usually benign glandular tumor composed of cells rich in mitochondria, such as an oncocytic parathyroid or salivary-gland tumor.

\medterm{Oxytocic} Causing the womb to contract. Oxytocic medicines, such as oxytocin, may start or strengthen labor or help control heavy bleeding after birth; they require monitoring because overly strong or frequent contractions can harm the pregnant person or baby.
\medterm{Oxytocics} Medicines that stimulate uterine contractions, used to start or strengthen labor and to prevent or treat heavy bleeding after birth. They require monitoring because overly strong or frequent contractions can harm the pregnant person or baby.

\medterm{Oxytocin} A hormone made in the brain and released by the pituitary gland. It helps the uterus contract during labor and causes milk to flow during breastfeeding; a manufactured form is used under monitoring to induce labor or control bleeding after birth.
\medterm{Oxytocin Infusion} A controlled drip of synthetic oxytocin into a vein to start or strengthen labor or help control heavy bleeding after birth. It requires continuous assessment because overly strong or frequent contractions can harm the pregnant person or baby.

\medterm{Oxyuriasis} Infection with pinworms, especially Enterobius vermicularis. The worms commonly cause itching around the anus at night, disturbed sleep, and sometimes abdominal discomfort; treatment usually includes medicine and household hygiene.

\medterm{Ozaena} A chronic form of atrophic rhinitis in which the inside of the nose becomes thin and widened, with dry crusts, a foul smell, and a blocked sensation despite an open passage. Medical assessment is needed to exclude infection and other causes.

\medterm{Ozena} Ozena is a chronic form of atrophic rhinitis with widened nasal passages, crusting, foul odor, and impaired mucosal function.

\medterm{Ocular Larva Migrans} Ocular toxocariasis caused by migrating larvae of \textit{Toxocara} roundworms, usually affecting one eye. It can cause retinal or intraocular inflammation, a white pupil, reduced vision, strabismus, or retinal damage and requires ophthalmic assessment; diagnosis is based on the eye findings, exposure history, and supportive serology.

\medterm{Omega-3 Fatty Acid} A type of polyunsaturated fat found in fish, seafood, flaxseed, chia, walnuts, and some oils. Omega-3s support cell membranes and cardiovascular function; supplements can interact with medicines and are not a substitute for treatment.

\medterm{Opioids} Medicines and drugs that act on opioid receptors to relieve pain. They can cause sleepiness, constipation, dependence, and life-threatening breathing suppression, especially with alcohol or sedatives; overdose may be reversed with naloxone.

\medterm{Opsoclonus-Myoclonus Syndrome} A rare neurologic disorder causing rapid, involuntary eye movements, muscle jerks, poor coordination, and sometimes behavioral or cognitive changes. In children it may be associated with neuroblastoma or an immune reaction.

\medterm{Optic Disc Swelling} Enlargement or elevation of the optic nerve head where it enters the eye. Causes range from raised pressure inside the skull to inflammation, ischemia, or local eye disease; sudden vision changes require urgent assessment.

\medterm{Oral-Facial-Digital Syndrome} A group of genetic disorders affecting the mouth, face, fingers, and toes. Features vary and may include oral abnormalities, unusual facial development, shortened or joined digits, and neurologic or kidney problems.

\medterm{Orbital Blowout Fracture} A break in one of the thin bones forming the eye socket, usually after blunt trauma. It can cause swelling, double vision, numbness, or restricted eye movement and may need urgent evaluation.

\medterm{Orthodontic} Relating to the diagnosis, prevention, or treatment of problems with tooth and jaw position. Orthodontic treatment may use braces, aligners, retainers, or other appliances.

\medterm{Orthorexia Nervosa} A term used for an unhealthy preoccupation with eating foods considered pure or healthy, leading to severe restriction, distress, or impaired daily life. It is not yet a formal standalone diagnosis in major classification systems, and professional assessment is important.

\medterm{Ozone} Ozone is a reactive form of oxygen that can irritate and injure the airways and lungs when inhaled, despite useful atmospheric and disinfecting roles.
\medterm{Oral Route} The administration of a medicine through the mouth for swallowing and absorption through the gastrointestinal tract. Clear instructions should spell out “by mouth” to reduce medication-route errors.

\synonyms
P.O.; per os; by mouth


\medterm{P Arm Of A Chromosome} The shorter arm of a chromosome, labeled p, opposite the longer q arm. A genetic report may use this label to show where a gene change or missing or extra DNA segment is located.

\medterm{P Wave} The first small wave on an ECG, showing the electrical activation of the heart’s upper chambers. Its shape and timing can help identify abnormal heart rhythms, enlargement of a chamber, or problems with electrical conduction.

\medterm{P-Glycoprotein} A transport protein in cell membranes that pumps many medicines and other substances out of cells. It affects how medicines are absorbed and removed and can cause interactions when its activity is blocked or increased.


\medterm{P53 Tumor Suppressor} A protein that helps prevent damaged cells from dividing or surviving. Changes in the TP53 gene can allow abnormal cells to multiply and contribute to many cancers, but a gene change alone does not determine a person's outcome.

\medterm{PA X-Ray} A posteroanterior X-ray is taken with the beam traveling from the back toward the front, commonly for a chest image. This position reduces enlargement of the heart’s outline compared with some other views.

\medterm{PA/IVS} Alternate index label; see \textit{Pulmonary atresia with intact ventricular septum}.

\medterm{Pacemaker} A small implanted device that sends electrical signals to the heart when its natural rhythm is too slow or unreliable. It may treat fainting or symptoms from slow heartbeats and certain conduction blocks; regular checks monitor the battery and wires.

\synonyms
Pacemaker, artificial

\medterm{Pacemaker Complications} Problems after pacemaker placement can include bleeding, infection, a collapsed lung, heart or blood-vessel injury, movement or breakage of a lead, blood clots, battery failure, or an abnormal heartbeat. New chest pain, fainting, fever, or worsening breathlessness needs urgent assessment.

\medterm{Pachyderma} Abnormal thickening and coarsening of the skin, which may occur with chronic inflammation, lymphedema, endocrine disease, or pachydermoperiostosis.

\medterm{Pachygyria} A birth-related brain-development abnormality in which the brain’s outer surface has fewer, broader folds than usual. It can cause developmental delay, learning difficulties, seizures, movement problems, or other neurologic effects of varying severity.

\medterm{Pachymeningitis} Inflammation and thickening of the tough outer covering of the brain and spinal cord. It may cause headache, nerve problems, or pressure on the brain and can result from infection, immune disease, or other inflammation.

\medterm{Pachyonychia Congenita} An inherited skin disorder causing very thick, abnormal nails and painful thick skin on the palms and soles. Some people also develop blisters, cysts, excess sweating, or white patches in the mouth.


\medterm{Pacing} Planning activity, rest, and effort so symptoms do not flare while endurance gradually improves. It may help with fatigue, chronic pain, heart or lung disease, and recovery after infection.

\medterm{Pacinian Corpuscle} A specialized nerve ending deep in the skin and connective tissue that detects deep pressure and vibration. It helps the nervous system sense contact, movement, and the handling of objects.


\medterm{Paclitaxel} Paclitaxel is a taxane chemotherapy drug that stabilizes microtubules and prevents cell division; it treats several cancers and can cause myelosuppression, neuropathy, alopecia, myalgia, and hypersensitivity reactions.

\medterm{Paclitaxel Liposome} A liposomal formulation of paclitaxel designed to alter drug distribution and delivery; paclitaxel disrupts microtubules and is used as chemotherapy in selected cancers.


\medterm{PACS} A computer system used by healthcare services to store, retrieve, view, and share medical images such as X-rays, CT scans, MRI scans, and ultrasound studies.





\medterm{PAES} Alternate index label; see \textit{Popliteal artery entrapment}.

\medterm{Page Kidney} High blood pressure caused when blood or another mass presses on a kidney from the outside. The pressure activates the kidney’s blood-pressure system and may follow injury, bleeding, or treatment with blood thinners.

\medterm{Paget Disease Of Bone} A long-term disorder in which bone is broken down and rebuilt too quickly and unevenly. Affected bones may become enlarged, weak, painful, or bent and can fracture or press on nearby nerves; many people have no symptoms.

\synonyms
Paget's disease; Paget disease of the bone; Osteitis deformans

\medterm{Paget Disease Of The Anus} A rare cancer in the skin around the anus that may look like persistent eczema, redness, itching, soreness, or an erosion. A biopsy is needed, and doctors usually check for cancer in nearby or internal organs.

\medterm{Paget Disease Of The Nipple} A rare breast cancer involving the nipple and areola, often causing persistent eczema-like scaling, redness, itching, burning, or discharge. It may be associated with underlying ductal carcinoma in situ or invasive breast cancer.

\medterm{Paget Disease Of The Vulva} A rare cancer in the skin of the vulva that may cause a persistent itchy, red, scaly, sore, or moist patch. A biopsy is needed, and assessment may look for another cancer beneath or elsewhere.

\medterm{Pagetoid Reticulosis} A rare localized variant of cutaneous T-cell lymphoma with solitary slowly enlarging scaly plaques containing atypical pagetoid lymphocytes.

\medterm{Pain} An unpleasant physical and emotional experience linked to actual or possible tissue injury. Pain can be short-lived or persistent and is influenced by nerves, inflammation, health, emotions, past experiences, and social circumstances.

\medterm{Pain (Geriatric)} Pain in an older adult, which may be underreported or shown as reduced activity, sleep change, or confusion. Assessment considers communication, memory, function, medicines, and causes; treatment is individualized and monitored for side effects.

\medterm{Pain (Vascular)} Pain caused by impaired blood flow, often appearing in a limb with exertion as claudication or, in more severe disease, at rest.


\medterm{Pain Assessment Scales} Tools that help a person describe pain intensity or its effect on daily life. They may use numbers, words, faces, or observation of behavior when self-report is not possible. A score tracks change but does not identify the cause or fully measure suffering.

\medterm{Pain Cancer (Cancer Pain)} Alternate index label for Cancer Pain.

\medterm{Pain Elbow (Elbow Pain)} Alternate index label for Elbow Pain.

\medterm{Pain Gallbladder (Gallbladder Pain (Gall Bladder Pain))} Alternate index label for Gallbladder Pain (Gall Bladder Pain).

\medterm{Pain Heel (Heel Spurs)} Alternate index label for Heel Spurs.

\medterm{Pain Kidney (Kidney Pain)} Alternate index label for Kidney Pain.


\medterm{Pain Lower Right Abdomen (Appendicitis)} Alternate index label for Appendicitis.

\medterm{Pain Management} Care that identifies the cause of pain and reduces pain while improving function and quality of life. It may combine medicines, physical treatment, psychological support, procedures, and treatment of the underlying condition, with risks reviewed regularly.

\medterm{Pain Medications - Opioid Analgesics} Opioid analgesics can relieve selected moderate-to-severe pain but may cause constipation, sedation, tolerance, dependence, overdose, and slowed breathing. They require individualized prescribing and monitoring.

\synonyms
Pain Medications - Narcotics

\medterm{Pain Medicine} The medical specialty concerned with assessing and treating acute, chronic, and complex pain through medication, procedures, rehabilitation, and multidisciplinary care.

\synonyms
Pain management; pain medicine specialty
\medterm{Pain Nerve (Neuropathic Pain (Nerve Pain))} Alternate index label for Neuropathic Pain (Nerve Pain).

\medterm{Pain Perception} The process by which the nervous system notices possible or actual tissue injury and the brain interprets it as pain. Attention, emotions, past experiences, inflammation, and medicines can change how pain feels.

\medterm{Pain Relief (Analgesia)} Reduction of pain using medicines, procedures, physical treatment, psychological support, or a combination. The best approach depends on the cause, duration, severity, function, other health conditions, and risks of treatment.
\synonyms
Analgesia

\medterm{Pain Scrotum (Testicular Disorders)} Alternate index label for Testicular Disorders.

\medterm{Pain Tailbone (Coccydynia)} Alternate index label for Coccydynia.

\medterm{Pain Threshold} The minimum intensity of a stimulus that a person perceives as painful; it varies with tissue, context, attention, inflammation, medications, and neurologic state.

\medterm{Pain Unit} A clinical service specializing in assessment and treatment of acute, chronic, cancer, neuropathic, or procedural pain using medicines, procedures, rehabilitation, and psychological care.

\medterm{Pain Vaginal (Vaginal Pain (Vulvodynia))} Alternate index label for Vaginal Pain (Vulvodynia).

\medterm{Pain Whiplash (Whiplash)} Alternate index label for Whiplash.

\medterm{Pain, Back} Alternate index label; see \textit{Back pain}.

\medterm{Pain, Breast} Alternate index label; see \textit{Breast pain}.

\medterm{Painful Bladder Syndrome} Alternate index label; see \textit{Interstitial cystitis}.

\medterm{Painful Intercourse (Dyspareunia)} Persistent or recurrent genital pain before, during, or after sexual intercourse, with physical, hormonal, emotional, or relationship-related causes.

\medterm{Painful Menstrual Periods} Painful menstrual periods, or dysmenorrhea, are cramps during menstruation caused by uterine contractions or conditions such as endometriosis or fibroids.

\medterm{Painful Swallowing} Painful swallowing, or odynophagia, is discomfort when swallowing and may result from infection, inflammation, injury, reflux, or a blockage.

\medterm{Paint, Lacquer, And Varnish Remover Poisoning} Poisoning caused by ingestion, inhalation, or skin or eye exposure to solvents and other chemicals in paint, lacquer, or varnish removers. Effects may include irritation, central nervous system depression, respiratory injury, metabolic abnormalities, or organ damage; severity depends on the substance, dose, route, and timing.

\medterm{Palatal Muscle} Muscle of the soft palate that elevates or tenses the palate and helps separate the nasal and oral cavities during swallowing and speech.

\medterm{Palatal Myoclonus} Repeated, involuntary twitching of the soft palate, sometimes causing a clicking sound in the ear or difficulty speaking and swallowing. It may follow a brainstem injury or occur without an identified cause.

\medterm{Palate} The roof of the mouth, consisting of a bony hard palate and muscular soft palate.
\medterm{Palate Neoplasm} An abnormal growth in the roof of the mouth. It may be noncancerous or cancerous and can cause a lump, pain, bleeding, loose teeth, difficulty swallowing, or speech changes.

\medterm{Palatine} Relating to the palate, the roof of the mouth, or the palatine bones. Palatine structures include the tonsils, nerves, and blood vessels of the mouth; problems may include tonsillitis, cleft palate, or infection near a tooth.
\medterm{Palatine Tonsil} One of two masses of lymphoid tissue at the sides of the throat that help detect inhaled or swallowed pathogens.


\medterm{Palatorrhaphy} Palatorrhaphy is surgical repair of a cleft or defect of the palate to restore separation between the oral and nasal cavities and improve feeding and speech.

\medterm{Palatoschisis} A congenital cleft or failure of fusion of the palate, commonly called cleft palate, which may affect feeding, speech, hearing, and dental development.

\medterm{Paleness} Unusual lightening of the skin or mucous membranes, which may result from reduced blood flow, anemia, shock, cold, or normal variation.



\medterm{Palindromic Rheumatism} Repeated attacks of joint pain and swelling that begin suddenly and resolve completely within hours or days, leaving no lasting joint damage between attacks. Some people later develop rheumatoid arthritis, especially when specific antibodies are present.

\medterm{Paliperidone} An antipsychotic medicine used to treat schizophrenia and schizoaffective disorder. It changes the action of certain brain chemicals and is available as tablets or long-acting injections; possible effects include sleepiness, movement problems, weight gain, hormone changes, and heart-rhythm changes.
\medterm{Palivizumab} A monoclonal antibody against respiratory syncytial virus fusion protein, used prophylactically in selected high-risk infants and young children during RSV season.


\medterm{Pallesthesia} The ability to feel vibration, usually tested with a vibrating tuning fork on the fingers or toes. Reduced vibration sense may result from nerve disease, spinal-cord problems, poor circulation, or other neurologic conditions.
\medterm{Palliative Care} Medical care for serious illness that relieves pain and other symptoms and supports daily life, emotional health, and family needs. It can begin at any stage and can be provided alongside treatment intended to control or cure disease.

\medterm{Palliative Rehabilitation} Rehabilitation for a person with serious or life-limiting illness that focuses on comfort, independence, meaningful activities, and quality of life. It may address weakness, breathlessness, pain, fatigue, movement, communication, swallowing, or daily tasks.

\medterm{Palliative Radiation} Radiation therapy given to relieve symptoms caused by cancer, such as pain, bleeding, obstruction, or pressure from a tumor. It is usually intended to improve comfort or function rather than cure the cancer and may be delivered in a short treatment course.

\medterm{Palliative Surgery} Surgery intended to relieve symptoms, prevent complications, or improve function when the underlying disease may not be curable. It can ease blockage, bleeding, pain, breathing difficulty, or problems caused by advanced cancer.
\medterm{Palliative Treatment} Treatment intended to relieve symptoms and improve comfort or quality of life rather than cure the underlying disease.

\medterm{Pallidotomy} Brain surgery that makes a carefully placed small injury in an area involved in movement control. It may reduce severe Parkinson disease symptoms such as involuntary movements when medicines no longer provide adequate control.
\medterm{Pallister-Hall Syndrome} A rare inherited condition that may cause an extra finger or toe, a noncancerous growth near the brain’s hormone-control center, and abnormalities of the anus, kidneys, airway, or pituitary hormones.

\medterm{Pallor} Unusual paleness of skin or mucous membranes, often due to reduced blood flow, anemia, or other illness.
\medterm{Palm} The inner surface of the hand between the wrist and fingers. It contains thick skin, creases, muscles, nerves, and blood vessels that help with grip, touch, movement, and protection.

\medterm{PALM-COEIN Classification} A system clinicians use to organize the causes of abnormal bleeding from the uterus. PALM covers structural causes such as polyps, adenomyosis, fibroids, and cancer; COEIN covers clotting, ovulation, lining, medicine-related, and other causes. It helps guide evaluation and treatment.

\synonyms
FIGO AUB Classification; PALM-COEIN System

\medterm{Palm Sweating Excessive (Hyperhidrosis)} Alternate index label for Hyperhidrosis.

\medterm{Palmar} Relating to the palm or front surface of the hand. The palm has thick skin, sweat glands, muscles, nerves, and blood vessels; changes such as redness, pain, numbness, or a tightening band may help identify an underlying condition.
\medterm{Palmar Erythema Of Pregnancy} Redness of both palms that can occur during pregnancy, often because of normal blood-vessel and hormone changes. It is usually harmless and fades after birth, but new itching, yellowing, pain, or other symptoms need assessment.

\medterm{Palmar Grasp Reflex} A primitive reflex in which an infant flexes the fingers around an object placed in the palm.
It is normally present at birth and gradually diminishes as voluntary grasp develops;
asymmetry or persistence may warrant neurologic assessment.

\medterm{Palmar Reflex} The palmar reflex is an infant response in which stroking the palm causes finger flexion; persistence or asymmetry beyond the expected age may indicate neurologic dysfunction.

\medterm{Palmitate} Palmitate is the salt or ester form of palmitic acid and participates in fatty-acid metabolism, lipid storage, and protein modification.

\medterm{Palmitic Acid} A 16-carbon saturated fatty acid common in dietary and stored fats; excess saturated-fat intake can contribute to adverse lipid profiles, although the body also synthesizes it normally.

\medterm{Palmoplantar Hyperkeratosis} Excessive thickening of the skin on the palms and soles, inherited or acquired, causing calluses, fissures, pain, impaired walking, and sometimes infection.

\medterm{Palonosetron Hydrochloride} A long-acting 5-HT3 serotonin-receptor antagonist used to prevent nausea and vomiting caused by chemotherapy, surgery, or selected other treatments.

\medterm{Palpable} Palpable means detectable by touch during physical examination, such as a mass, enlarged organ, lymph node, pulse, tenderness, or abnormal surface.

\medterm{Palpate} To examine part of the body by feeling it carefully with the hands. A healthcare professional may assess size, shape, tenderness, warmth, movement, swelling, firmness, pulses, or unusual lumps.

\medterm{Palpation} Examination by touching the body to assess tenderness, temperature, texture, masses, pulses, organ size, and movement.

\medterm{Palpebra} The anatomical term for an eyelid. Eyelids protect the eye, spread tears over the cornea, and limit light entry; disorders include drooping, inward or outward turning, inflammation, styes, and blocked oil glands.
\medterm{Palpebral Conjunctiva} The thin, moist membrane lining the inner surfaces of the eyelids. It joins the membrane covering the front of the eye and helps protect, lubricate, and move tears across the eye’s surface.

\medterm{Palpebral Slant - Eye} The direction of the opening between the upper and lower eyelids. Its shape varies normally, but an unusual slant may be associated with a genetic or developmental condition.

\medterm{Palpitation} Awareness of the heart beating fast, forcefully, irregularly, or as if it skipped a beat. Causes include stress, stimulants, anemia, thyroid disease, medicines, and abnormal rhythms.


\medterm{PALS} Alternate name; see \textit{Pediatric Advanced Life Support}.

\medterm{Palsy} Weakness or loss of movement in a body part, often affecting muscles supplied by a particular nerve. It may be temporary or lasting and can result from nerve injury, infection, stroke, inflammation, or a problem present from birth.
\medterm{Palsy Cerebral (Cerebral Palsy)} Alternate index label for Cerebral Palsy.

\medterm{Palsy Of The Long Thoracic Nerve} Weakness of the nerve that activates a muscle holding the shoulder blade against the ribs. It can cause a winged shoulder blade, difficulty raising the arm, and pain after injury, infection, surgery, or pressure on the nerve.

\medterm{Palsy, Bell's} Alternate index label; see \textit{Bell's palsy}.

\medterm{Palsy, Cerebral} Alternate index label; see \textit{Cerebral palsy}.

\medterm{Palsy, Facial} Alternate index label; see \textit{Bell's palsy}.

\medterm{Pamidronic Acid} An intravenous bisphosphonate that inhibits osteoclast-mediated bone resorption and is used for selected hypercalcemia, bone complications of cancer, and Paget disease.


\medterm{Panacinar Emphysema} A form of emphysema in which air sacs and small airways are damaged throughout affected areas of the lung. It can cause breathlessness and reduced lung function and is strongly associated with severe alpha-1 antitrypsin deficiency.

\synonyms
Panlobular Emphysema; Lower Lobe Emphysema


\medterm{Pancoast Syndrome} Symptoms caused by a cancer at the top of the lung growing into nearby nerves or other tissues. It may cause severe shoulder or arm pain, hand weakness, numbness, and drooping of one eyelid with a smaller pupil.

\medterm{Pancoast Tumor} A lung cancer at the top of the lung that can grow into nearby nerves, blood vessels, or the spine. It may cause severe shoulder or arm pain, hand weakness, or drooping of one eyelid.

\medterm{Pancolitis} Inflammation affecting the entire large intestine. It may result from ulcerative colitis, infection, reduced blood flow, or another inflammatory condition and can cause diarrhea, abdominal pain, fever, or bleeding.
\medterm{Pancreas} An organ behind the stomach that makes digestive juices and hormones. Its digestive enzymes enter the small intestine, while insulin and other hormones help control blood sugar and digestion.

\medterm{Pancreas Divisum} A birth-related difference in which two drainage tubes of the pancreas do not join normally. Most people have no symptoms, but some develop repeated pancreatitis because pancreatic fluid does not drain freely.

\synonyms
Divisum Pancreas (Pancreas Divisum)

\medterm{Pancreas Enzyme} An enzyme produced by the pancreas, such as amylase, lipase, or proteases, that helps digest carbohydrate, fat, or protein in the small intestine.

\medterm{Pancreas Infection} Infection involving the pancreas or a pancreatic collection, usually occurring with pancreatitis, surgery, trauma, or spread from nearby infection.

\medterm{Pancreas Inflammation} Alternate index label; see \textit{Pancreatitis}.

\medterm{Pancreas Inflammation (Pancreatitis)} Alternate index label for Pancreatitis.

\medterm{Pancreas Inflammatory Cysts (Pancreatic Cysts)} Alternate index label for Pancreatic Cysts.

\medterm{Pancreas Intraductal Papillary Mucinous Neoplasm (Pancreatic Cysts)} Alternate index label for Pancreatic Cysts.

\medterm{Pancreas IPMN (Pancreatic Cysts)} Alternate index label for Pancreatic Cysts.

\medterm{Pancreas Mucinous Cyst Adenomas (Pancreatic Cysts)} Alternate index label for Pancreatic Cysts.

\medterm{Pancreas Neoplasm} An abnormal growth in the pancreas that may be noncancerous, precancerous, or cancerous. Its effects depend on its location and may include pain, jaundice, digestive problems, or diabetes.

\medterm{Pancreas Non-Inflammatory Cysts (Pancreatic Cysts)} Alternate index label for Pancreatic Cysts.

\medterm{Pancreas Pseudocysts (Pancreatic Cysts)} Alternate index label for Pancreatic Cysts.

\medterm{Pancreas Serous Cyst Adenomas (Pancreatic Cysts)} Alternate index label for Pancreatic Cysts.

\medterm{Pancreas Solid Pseudopapillary Tumor (Pancreatic Cysts)} Alternate index label for Pancreatic Cysts.

\medterm{Pancreas Transplant} A procedure that replaces a diseased pancreas with a healthy donor pancreas, usually to restore insulin production in selected people with severe diabetes.

\medterm{Pancreas Transplantation} Alternate index label; see \textit{Pancreas Transplant}.

\medterm{Pancreas, Diseases Of} Disorders affecting the pancreas, including inflammation, cysts, tumors, diabetes, and poor digestive-enzyme production. They may cause abdominal pain, vomiting, weight loss, jaundice, abnormal blood sugar, or digestive problems.
\medterm{Pancreatectomy} Surgery to remove part or all of the pancreas. It may treat a pancreatic tumor, severe injury, or another serious disease; removal of pancreatic tissue can affect digestion and blood-sugar control.

\medterm{Pancreatic Abscess} A pocket of pus in or near the pancreas, usually after severe or infected pancreatitis. It may cause continuing fever, worsening abdominal pain, vomiting, or sepsis and often needs drainage and antibiotics.

\medterm{Pancreatic Adenocarcinoma} The most common cancer of the pancreas, arising from cells in its digestive-juice ducts. It may cause jaundice, abdominal or back pain, weight loss, or no early symptoms; smoking, chronic pancreatitis, diabetes, obesity, and inherited risks can contribute.

\medterm{Pancreatic Cancer} Cancer that begins in the pancreas, most often in cells lining its digestive-juice ducts. It may cause jaundice, abdominal or back pain, weight loss, or digestive problems, but early cancer may cause no symptoms; treatment depends on type and spread.

\synonyms
Cancer, Pancreatic

\medterm{Pancreatic Carcinoma} Pancreatic carcinoma is cancer arising in the pancreas, most often from duct cells, and may cause jaundice, pain, weight loss, or digestive problems.

\medterm{Pancreatic Cysts} Fluid-filled sacs in or around the pancreas. Some follow inflammation or injury, while others are growths that can become cancerous; imaging and sometimes fluid or tissue testing guide follow-up or treatment.

\synonyms
Cyst, Pancreatic
Cysts Pancreatic Inflammatory (Pancreatic Cysts)
Cysts True Pancreas (Pancreatic Cysts)
Intraductal Papillary Mucinous Neoplasm Pancreas (Pancreatic Cysts)
IPMN Pancreas (Pancreatic Cysts)
Mucinous Cyst Adenomas Pancreas (Pancreatic Cysts)
Pancreas Inflammatory Cysts (Pancreatic Cysts)
Pancreas Intraductal Papillary Mucinous Neoplasm (Pancreatic Cysts)
Pancreas IPMN (Pancreatic Cysts)
Pancreas Mucinous Cyst Adenomas (Pancreatic Cysts)
Pancreas Non-Inflammatory Cysts (Pancreatic Cysts)
Pancreas Pseudocysts (Pancreatic Cysts)
Pancreas Serous Cyst Adenomas (Pancreatic Cysts)
Pancreas Solid Pseudopapillary Tumor (Pancreatic Cysts)
Serous Cyst Adenomas Pancreas (Pancreatic Cysts)

\medterm{Pancreatic Diseases} Disorders affecting the pancreas’s digestive or hormone-producing functions, including inflammation, cysts, tumors, duct blockage, and diabetes-related damage. They may cause abdominal pain, vomiting, weight loss, jaundice, abnormal blood sugar, or poor digestion.

\medterm{Pancreatic Duct} The main tube carrying digestive fluid from the pancreas into the first part of the small intestine. It often joins the bile duct before opening into the intestine; blockage can cause pain, jaundice, or pancreatitis.

\medterm{Pancreatic Ductal Adenocarcinoma} The most common exocrine pancreatic malignancy, arising from ductal epithelium and often presenting late with painless jaundice, weight loss, pain, or pancreatitis. Diagnosis and staging rely on imaging, tissue sampling when needed, and multidisciplinary assessment.

\medterm{Pancreatic Elastase} A stool test that estimates how well the pancreas makes digestive enzymes. A low result supports pancreatic insufficiency, which can cause greasy stools, weight loss, and vitamin deficiency, but watery stool and other factors can affect interpretation.

\medterm{Pancreatic Enzyme} A digestive substance made by the pancreas and released into the small intestine. Amylase breaks down starch, lipase breaks down fat, and proteases break down protein; bicarbonate helps neutralize stomach acid.

\medterm{Pancreatic Exocrine Insufficiency} Alternate index label for Exocrine Pancreatic Insufficiency.

\medterm{Pancreatic Fistula} A pancreatic fistula is an abnormal passage through which pancreatic juice leaks from the pancreas, often after trauma or surgery, causing fluid collections, ascites, or skin drainage.

\medterm{Pancreatic Hormones} Hormones produced by pancreatic islet cells, including insulin, glucagon, somatostatin, and pancreatic polypeptide, which regulate glucose, digestion, and energy metabolism.

\medterm{Pancreatic Insufficiency} Failure of the pancreas to produce or deliver enough digestive enzymes or other secretions for normal digestion and absorption. It may cause greasy stools, diarrhea, bloating, weight loss, or vitamin deficiency; exocrine insufficiency specifically refers to inadequate digestive enzymes.

\medterm{Pancreatic Islet} A pancreatic islet is a cluster of endocrine cells that releases hormones such as insulin, glucagon, and somatostatin to regulate blood glucose and digestion.


\medterm{Pancreatic Juice} Pancreatic juice is an alkaline digestive secretion containing bicarbonate and enzymes such as amylase, lipase, and proteases, delivered into the duodenum.

\medterm{Pancreatic Mass} A lump or abnormal area seen in the pancreas on an imaging test. It may be a cyst, inflammation, scar, or tumor; further scans, blood tests, or a tissue sample may be needed to determine its cause.
\medterm{Pancreatic Neuroendocrine Tumor} A tumor arising from hormone-related cells in the pancreas. It may release excess hormones or cause no hormone symptoms and can grow slowly or aggressively; its behavior depends on size, spread, cell features, and growth rate.

\medterm{Pancreatic Neuroendocrine Tumors} Tumors arising from hormone-making cells in the pancreas. Some release excess hormones, such as insulin or gastrin, while others do not. Their behavior depends on size, spread, cell appearance, growth rate, and hormone production.

\synonyms
Islet Cell Cancer

\medterm{Pancreatic Necrosis} Death of pancreatic or nearby fatty tissue, usually during severe acute pancreatitis. The dead tissue may remain sterile or become infected, causing fever, worsening pain, sepsis, and organ failure.

\synonyms
Necrotizing pancreatitis; Acute necrotizing pancreatitis

\medterm{Pancreatic Pseudocyst} A walled-off collection of pancreatic fluid that usually develops after acute or chronic pancreatitis. It has no true cell lining and may cause abdominal pain, vomiting, early fullness, infection, bleeding, or blockage of nearby organs.

\medterm{Pancreatic Surgery} An operation on the pancreas to remove diseased tissue, drain a collection, repair an injury, or treat cancer. The type of surgery depends on the problem and its location.
\medterm{Pancreatic Trauma} Injury to the pancreas caused by blunt or penetrating trauma. Assessment may require CT, laboratory testing, and observation or intervention according to the injury.

\medterm{Pancreaticobiliary Cytology} Cytologic evaluation of pancreatic and biliary lesions, usually from fine-needle aspiration, brushings, or bile-duct sampling. It is used to assess neoplasia, inflammation, and obstructive lesions alongside imaging and histology.

\medterm{Pancreaticoduodenectomy} Pancreaticoduodenectomy, or the Whipple procedure, removes the head of the pancreas and nearby parts of the digestive tract to treat selected tumors and other serious disease.

\medterm{Pancreatitis} Inflammation of the pancreas that may be acute or chronic. Causes include gallstones, alcohol, medicines, metabolic disorders, and other conditions. Acute disease may cause severe upper-abdominal pain, vomiting, fluid loss, infection, tissue death, or organ failure; chronic disease can permanently scar the pancreas and cause poor digestion, weight loss, or diabetes.

\synonyms
Inflamed Pancreas
Pancreas Inflammation
Pancreas Inflammation (Pancreatitis)

\medterm{Pancreatitis - Children} Pancreatitis in children is inflammation of the pancreas causing abdominal pain, nausea, vomiting, and elevated pancreatic enzymes from structural, genetic, medication, infectious, or other causes.


\medterm{Pancreatoblastoma} A rare cancer of the pancreas that occurs mainly in children. It may cause an abdominal lump, pain, vomiting, weight loss, or hormone-related symptoms and requires specialist imaging, biopsy, staging, and treatment.

\medterm{Pancreatography} Imaging of the tubes that carry digestive juices from the pancreas, usually after contrast dye is introduced. It can show narrowing, blockage, leaks, stones, or other changes and may be done during ERCP or with MRI; ERCP also carries a risk of pancreatitis.

\medterm{Pancreazymin} An older name for cholecystokinin, a hormone that stimulates pancreatic enzyme secretion and gallbladder contraction.
\medterm{Pancrelipase} Pancrelipase is a mixture of pancreatic digestive enzymes used as replacement therapy for exocrine pancreatic insufficiency, including that caused by cystic fibrosis or pancreatic disease.

\medterm{Pancuronium} A long-acting medicine that temporarily prevents skeletal muscles from contracting during anesthesia or assisted breathing. It causes paralysis but does not relieve pain or cause unconsciousness and must be given with reliable airway and breathing support.

\medterm{Pancytopenia} A reduction in red blood cells, white blood cells, and platelets at the same time. It can cause tiredness, infections, and bleeding and may result from medicines, infection, immune disease, or a problem in the bone marrow.
\medterm{PANDAS} A proposed condition in which obsessive-compulsive symptoms or movement tics begin suddenly or worsen after a streptococcal infection in a child. The diagnosis remains debated, so clinicians must also assess common neurologic, psychiatric, and infectious causes.

\medterm{Pandemic} An epidemic that spreads across multiple countries or continents and affects many people. It may occur when an infectious agent spreads widely among populations with limited previous immunity.

\medterm{Panendoscopy} Examination of several connected areas of the throat and upper swallowing and breathing passages with a viewing instrument, often including the throat, voice box, and swallowing tube. Tissue samples may be taken during the procedure.
\medterm{Paneth Cell} A cell at the base of a small-intestinal gland that releases substances helping control bacteria and other microbes in the gut. It supports the intestinal lining’s normal defense system.
\medterm{Panhypopituitarism} Failure of the pituitary gland to make most or all of its important hormones. It can cause tiredness, low blood pressure, infertility, growth or thyroid problems, and inability to respond normally to illness; lifelong hormone replacement may be needed.

\medterm{Panic} Panic is sudden intense fear or distress with autonomic symptoms such as racing heart, sweating, trembling, breathlessness, or a sense of impending danger.

\medterm{Panic Attack} A sudden episode of intense fear or discomfort, often with racing heart, shortness of breath, trembling, or dizziness.

\medterm{Panic Attacks And Panic Disorder} Panic attacks are sudden episodes of intense fear; panic disorder involves recurrent attacks and persistent concern about having more.

\medterm{Panic Disorder} Repeated unexpected panic attacks followed by ongoing worry about more attacks or major changes in behavior to avoid them. Attacks peak within minutes and may cause a racing heart, breathlessness, trembling, dizziness, or fear of dying.

\medterm{Panlobular Emphysema} A type of permanent lung damage affecting the whole small air-sac unit, often more severe in the lower lungs. It is strongly associated with alpha-1 antitrypsin deficiency and can cause breathlessness, cough, and reduced exercise ability.

\medterm{Panmyelopathy} A term for widespread disease affecting the bone marrow and sometimes the spinal cord. The exact meaning depends on context and may involve reduced production of several blood-cell types together with neurologic problems.
\medterm{Panniculectomy} Surgical removal of a hanging apron of excess lower-abdominal skin and fat, usually after substantial weight loss; it is distinct from a full abdominoplasty.

\medterm{Panniculitis} Inflammation of the layer of fat beneath the skin, usually causing tender red or purple lumps, often on the legs. Causes include infection, immune disease, medicines, injury, and pancreatic disease; treatment depends on the cause.

\medterm{Panniculitis Idiopathic Nodular (Weber-Christian Disease)} Alternate index label for Weber-Christian Disease.

\medterm{Panniculus} A panniculus is a fold or layer of subcutaneous fatty tissue, especially the lower abdominal apron that may overhang the pubic region.

\medterm{Pannus} An abnormal layer of inflamed tissue that grows over a joint or the clear front surface of the eye. It can damage cartilage or reduce vision and may occur with rheumatoid arthritis, trachoma, or other chronic inflammation.

\medterm{Panophthalmitis} Severe infection and inflammation involving the entire eye and its surrounding tissues. It can cause intense pain, redness, swelling, discharge, and major vision loss and requires emergency treatment to control infection and protect life and sight.
\medterm{Panoramic Radiograph} An extraoral radiograph displaying both jaws, teeth, temporomandibular regions, and
surrounding structures in one image. It is useful for surveying dentition, development,
impaction, and gross osseous pathology but has less fine detail than intraoral
radiographs.

\medterm{Panretinal Photocoagulation} A laser treatment that places many small burns in the outer retina to reduce the abnormal blood-vessel growth caused by severe retinal disease, especially diabetic retinopathy. It can preserve central vision but may reduce side or night vision.

\medterm{Pantothenic Acid} Vitamin B5, needed to help cells make energy and produce fatty acids and certain hormones.
\medterm{Pantothenic Acid And Biotin} Two B vitamins involved in energy metabolism, fatty-acid synthesis, and other enzyme reactions; deficiency is uncommon but can occur with severe malnutrition, malabsorption, or specific inherited disorders.

\medterm{Panuveitis} Inflammation affecting all the main internal layers around the eye, including the iris, the gel inside the eye, and the tissues near the retina. It can cause pain, redness, light sensitivity, floaters, and vision loss and needs urgent eye care.

\medterm{Pap Smear} A screening test in which cells are collected from the cervix and examined for early changes that could become cervical cancer or for cancer itself. It may be done alone or with a test for high-risk human papillomavirus (HPV).

\medterm{Pap Test} A screening test for cervical cancer and precancerous changes in which cells are collected from the cervix and examined under a microscope. It may also reveal inflammation or some infections, but an abnormal result does not by itself diagnose cancer.

\medterm{Papain} A protein-breaking enzyme found in papaya latex. It is used in some foods and products, but evidence for many medicinal claims is limited; it can irritate tissues or cause allergic reactions.
\medterm{Papanicolaou (Pap) Test} A screening test that collects cells from the cervix and examines them for abnormal changes. It can detect precancerous cells and cancer early, allowing treatment before serious disease develops.
\medterm{Papillary Fibroelastoma} A usually noncancerous, small growth attached to a heart valve. It often causes no symptoms but can release a clot and lead to a stroke or other blockage, especially when it is mobile; treatment depends on its size, location, and risk.

\medterm{Papillary Muscle} Muscles inside the lower chambers of the heart that connect by thin cords to the heart valves. They tighten when the heart contracts to keep the valves from turning inside out and prevent blood from flowing backward.

\medterm{Papillary Thyroid Carcinoma} The most common type of thyroid cancer. It often grows slowly and may spread to nearby neck lymph nodes; treatment and outlook depend on the tumor's size, spread, and other features.

\medterm{Papillectomy} Surgical removal of a papilla or papillary growth, such as a tumor near the bile and pancreatic ducts.
\medterm{Papilledema} Swelling of the area where the optic nerve enters the eye because pressure inside the skull is raised. It can threaten vision and may result from bleeding, a mass, infection, or other serious illness.

\medterm{Papilloedema} Swelling of the optic discs caused by raised intracranial pressure.
\medterm{Papilloma} A usually noncancerous growth from a surface lining, often appearing as a wart, skin tag, or small finger-like projection. Some are linked to human papillomavirus; the site, appearance, and tissue examination determine whether further treatment is needed.

\medterm{Papillomatosis Cutis Lymphostatica} Thick, rough, wart-like skin caused by long-standing swelling from poor lymph or vein drainage, usually in the legs. The skin may become cobblestoned and prone to infection; controlling the swelling and protecting the skin are important.

\medterm{Papillomavirus} A papillomavirus, especially human papillomavirus, which can cause warts and some anogenital, oropharyngeal, and other cancers.
\medterm{Papular} A papule: a small, solid, elevated skin lesion <1 cm in diameter (palpable), of various
colors and etiologies.

\medterm{Papular Urticaria} A hypersensitivity reaction to insect bites, especially from fleas, bedbugs, or
mosquitoes, causing recurrent intensely itchy papules or small vesicles, often clustered
on exposed areas in children. Diagnosis is clinical.

\medterm{Papule} A small, solid, raised skin lesion, usually less than 1 centimetre across and without visible fluid. Papules can result from irritation, inflammation, infection, an overgrowth of skin cells, or another skin disorder.
\medterm{PAPVR} Alternate index label; see \textit{Partial anomalous pulmonary venous return}.


\medterm{Paracentesis} A procedure in which a clinician uses a needle to remove fluid from the abdomen. The fluid may be tested to find its cause or removed to relieve pressure and discomfort.

\medterm{Paracetamol Overdose} Taking too much paracetamol, also called acetaminophen, can seriously damage the liver even before symptoms appear. The risk depends on the amount, timing, repeated doses, and health of the person; urgent assessment is needed because an antidote can prevent liver failure.

\medterm{Paracetamol Toxicity} Liver injury caused by excessive acetaminophen (paracetamol) exposure. Severe poisoning may
cause nausea, abdominal pain, liver failure, or death and requires urgent assessment.

\medterm{Paracusia} A hearing disturbance in which sounds are perceived unusually, inaccurately, or differently from their actual qualities. It may occur with hearing loss, inner-ear disease, neurological conditions, medicines, or mental health disorders.
\medterm{Paradichlorobenzene Poisoning} Harm from mothballs or other products containing paradichlorobenzene. It may cause headache, nausea, dizziness, liver injury, or destruction of red blood cells; significant exposure needs poison-control advice.

\medterm{Paraesophageal Hernia} A hiatal hernia in which part of the stomach moves through the diaphragm beside the esophagus.
It may cause obstruction or strangulation even when reflux is absent.

\medterm{Paraesthesia} An unusual skin sensation such as tingling, pins and needles, burning, prickling, or numbness. It can result from pressure on a nerve, nerve disease, poor circulation, medicines, vitamin deficiency, or other causes.
\medterm{Paraffin} A mixture of hydrocarbons used in laboratory tissue processing, topical preparations, and other applications; some forms are flammable.
\medterm{Paraffin Poisoning} Harm from swallowing or breathing paraffin wax or liquid paraffin. Small swallowed amounts may cause stomach upset, but liquid entering the lungs can cause chemical pneumonia; do not induce vomiting.

\medterm{Paraganglioma} A rare tumor arising from specialized nerve-related cells outside the adrenal glands. Some release hormones that cause high blood pressure, sweating, headaches, or a fast heartbeat; inherited risk may require genetic assessment.

\medterm{Parageusia (Dysgeusia)} An altered sense of taste in which food may taste metallic, bitter, salty, unpleasant, or different from usual. Causes include infections, medicines, nutritional deficiencies, head injury, and mouth, nerve, or salivary-gland problems.
\medterm{Paragonimiasis} A parasitic infection caused by lung flukes, usually acquired by eating raw or undercooked freshwater crab or crayfish. It commonly affects the lungs and may cause a long-lasting cough, chest pain, fever, or blood-stained sputum.

\medterm{Paragraphia} A language problem in which a person writes an unintended word, letter, or symbol instead of the one intended. It may occur with damage to language areas of the brain, such as after a stroke, and is assessed with other speaking, reading, and writing abilities.
\medterm{Parainfectious} Describing a problem that occurs during or soon after an infection, often because the immune system reacts to the infection rather than because germs directly invade the affected tissue.

\medterm{Parainfluenza} Parainfluenza viruses are respiratory viruses that commonly cause colds, croup, bronchitis, or pneumonia, especially in children; illness ranges from mild to severe airway disease.

\medterm{Paraldehyde} An older sedative and anticonvulsant medicine with limited current use.
\medterm{Paralysis} Loss of voluntary muscle movement in part or all of the body due to nervous-system or muscle disease.
\medterm{Paralysis Agitans} An older name for Parkinson disease, characterized by tremor, rigidity, and bradykinesia.
\medterm{Paralysis, Stomach} Stomach paralysis, or gastroparesis, is delayed gastric emptying without a mechanical blockage, causing early fullness, nausea, vomiting, bloating, and variable glucose control.

\medterm{Paralytic Ileus} Failure of intestinal peristalsis without a mechanical blockage, causing abdominal distension, pain, vomiting, and absent or reduced bowel sounds.
\medterm{Paramedical} Relating to healthcare professionals or services that support medical practice, such as physiotherapy, radiography, or emergency medical care.
\medterm{Parametritis} Inflammation or infection of the tissue beside the uterus, often after pelvic infection, childbirth, miscarriage, or surgery. It may cause pelvic pain, fever, tenderness, abnormal discharge, and requires treatment based on its cause.
\medterm{Paramnesia} A memory disturbance involving false, distorted, or inaccurately familiar recollections.
\medterm{Paramyotonia Congenita} A rare inherited muscle disorder in which muscles become stiff and slow to relax, especially after repeated movement or exposure to cold. The face, eyelids, neck, and hands are often affected; temporary weakness may follow an attack. Avoiding cold and strenuous triggers helps, and medicines may reduce symptoms.

\synonyms
Eulenburg Disease; Paramyotonia Congenita of Eulenburg

\medterm{Paranasal Sinuses} Air-filled spaces in the bones around the nose that connect with the nasal passages. Their lining produces mucus, which drains into the nose; inflammation or blockage can cause facial pressure, congestion, and pain.

\medterm{Paraneoplastic Disorder} A health problem caused indirectly by cancer rather than by the tumor invading or spreading to the affected organ. Tumors may release hormones or trigger immune reactions that affect the nerves, skin, blood, joints, or other distant organs.

\medterm{Paraneoplastic Neurologic Syndrome} A nervous-system problem caused by the immune response to a cancer rather than by the tumor pressing on the nerves. It may cause weakness, loss of balance, memory problems, seizures, or abnormal sensation and can occur before the cancer is found.

\medterm{Paraneoplastic Neurological Disorder} Alternate index label; see \textit{Paraneoplastic syndromes of the nervous system}.

\medterm{Paraneoplastic Neurological Syndrome} Alternate index label; see \textit{Paraneoplastic syndromes of the nervous system}.

\medterm{Paraneoplastic Syndrome} A group of symptoms caused by substances or immune reactions from a cancer, rather than by the tumor pressing on or spreading to the affected area. It may cause hormone, nerve, skin, blood, or joint problems; treating the cancer may improve them.

\medterm{Paraneoplastic Syndromes Of The Nervous System} Nervous-system problems caused by an immune reaction to a cancer rather than by the cancer growing into the nerves or brain. They may cause weakness, poor balance, seizures, memory change, or abnormal sensation.

\synonyms
Cancer, Paraneoplastic Syndromes
Paraneoplastic Neurological Disorder
Paraneoplastic Neurological Syndrome

\medterm{Paranoia} A thought process characterized by excessive or irrational suspiciousness and distrust
of others, often manifesting as delusions of persecution, reference, or grandeur.

\medterm{Paranoid} Characterized by persistent or excessive suspicion and distrust, often involving beliefs that others intend harm. Severe or fixed beliefs may occur in a psychotic disorder and need professional assessment.
\medterm{Paranoid Personality Disorder} A longstanding pattern of distrust and suspiciousness in which other people’s motives are interpreted as threatening or deceitful.

\medterm{Paranoid Psychosis} An older term for psychosis dominated by delusions of persecution, reference, or grandeur.
\medterm{Paraparesis} Weakness in both legs caused by disease affecting the spinal cord, brain, nerves, or muscles.
\medterm{Paraphasia} An unintended substitution, omission, or distortion of a sound or word during speech. It commonly occurs with aphasia after brain injury or stroke and may affect naming, conversation, reading, or writing.
\medterm{Paraphilia} A persistent and intense sexual interest involving atypical objects, situations, or activities. A paraphilia by itself is not a mental disorder. It becomes a paraphilic disorder when it causes significant personal distress or problems in daily life, apart from distress caused only by social disapproval, or when acting on it involves harm, risk of harm, or people who do not or cannot consent.

\medterm{Paraphilic Disorder} Alternate index label for Paraphilic Disorders.

\medterm{Paraphilic Disorders} The group of mental disorders in which a paraphilia causes significant personal distress or problems in daily life, apart from distress caused only by social disapproval, or involves harm, risk of harm, or people who do not or cannot consent. The group includes voyeuristic, exhibitionistic, frotteuristic, fetishistic, sexual masochism, sexual sadism, pedophilic, and transvestic disorders, as well as other specified and unspecified paraphilic disorders. A paraphilia without these consequences is not, by itself, a paraphilic disorder.

\medterm{Paraphimosis} Paraphimosis occurs when a retracted foreskin cannot be returned over the glans, causing painful swelling and possible loss of blood flow; it requires urgent reduction.

\medterm{Paraplegia} Paralysis of both legs and usually part of the lower trunk. Paralysis of both legs and often part of the trunk, usually from spinal-cord or neurologic disease.
\medterm{Parapneumonic Effusion} A collection of fluid between the lung and chest wall that develops with pneumonia. Small uncomplicated collections may improve with antibiotics, while infected or loculated fluid can cause persistent fever, chest pain, or breathlessness and may require drainage.

\medterm{Parapneumonic Pleural Effusion} Pleural fluid accumulation associated with bacterial pneumonia. It may be uncomplicated
or progress to a complicated effusion or empyema, requiring imaging, pleural-fluid
analysis, antibiotics, and sometimes drainage.

\medterm{Paraquat} A highly toxic herbicide. Ingestion can cause injury to the mouth and gastrointestinal tract, kidney failure, and delayed lung injury; severe poisoning has a high risk of death.

\medterm{Paraquat Poisoning} A medical emergency caused by swallowing or significant exposure to the herbicide paraquat. It can burn the mouth and digestive tract and later cause lung, kidney, or liver failure; immediate emergency care is essential.

\medterm{Paraseptal Emphysema} A form of emphysema affecting air sacs near the outer surface of the lung. It may cause few symptoms but can form fragile air pockets that rupture and produce a collapsed lung, especially in tall young adults or people who smoke.

\medterm{Parasite} An organism that lives on or inside another organism, called the host, and obtains food or shelter from it. Parasites may cause disease in the skin, intestine, blood, or organs.

\medterm{Parasitemia} The presence and amount of parasites in the bloodstream. It is especially used for malaria and some other blood infections; a higher amount may indicate more severe disease, but interpretation depends on the parasite and the person’s symptoms.

\medterm{Parasitic} Parasitic means relating to an organism that lives in or on a host and obtains nutrients at the host’s expense, sometimes causing disease.

\medterm{Parasitic Infection} An illness caused by a parasite---an organism that lives on or inside another living creature (the host) and survives by taking nutrients at the host's expense. Human parasites belong to three main groups: single-celled protozoa (such as \textit{Plasmodium} causing malaria, \textit{Entamoeba} causing amebic dysentery, and \textit{Giardia} causing giardiasis, which can multiply rapidly inside the body), helminths or parasitic worms (such as roundworms, tapeworms, pinworms, hookworms, and flukes, which are multicellular and can live in tissues for years), and ectoparasites (such as scabies mites, lice, fleas, and ticks that feed on or burrow into the skin). Parasites enter the body through contaminated water or unwashed food, undercooked meat, insect bites, walking barefoot on contaminated soil, or direct physical contact. Symptoms depend on the parasite and can range from watery or bloody diarrhea, stomach cramps, and nausea, to prolonged fevers, extreme fatigue, severe anemia, skin rashes, and chronic organ enlargement or damage. Parasitic infections are diagnosed using stool tests, blood smears, or tissue exams, and are treated with targeted antiparasitic medicines (such as antiprotozoals or worm-clearing anthelmintics). Access to clean drinking water, proper sanitation, thorough cooking, and insect protection are key to prevention.

\synonyms
Parasitic Infections; Parasitic Disease; Parasitosis; Parasitic Infestation

\medterm{Parasitic Worm} A worm that lives in or on a host and uses the host for nourishment. Human infections may affect the intestine, blood, lymph system, lungs, liver, skin, or other tissues and can cause anemia, pain, or organ damage.

\medterm{Parasitology} The branch of microbiology that studies parasites, including protozoa, helminths, and
ectoparasites. Parasitic infections are a major global health problem, particularly in
tropical and subtropical regions. Protozoan parasites include Plasmodium (malaria),
Trypanosoma (sleeping sickness, Chagas disease), and Leishmania (leishmaniasis).

\medterm{Parasitosis} Infestation or infection by a parasite; the term may describe a confirmed parasitic disease or, in delusional parasitosis, a fixed false belief of infestation without evidence.

\medterm{Parasomnia} A parasomnia is an undesirable movement, behavior, emotion, or experience occurring during sleep or transitions between sleep and wakefulness.

\medterm{Parasomnias} Unwanted movements, behaviors, emotions, or experiences that occur during sleep or at the transition between sleep and wakefulness.

\medterm{Parasuicide} Parasuicide is an older term for a suicide attempt or intentional self-harm without fatal outcome; current assessment focuses on self-harm, intent, safety, and urgent support.

\medterm{Parasympathetic Nervous System} The part of the automatic nervous system that supports resting functions such as digestion, saliva, urination, and slowing the heartbeat. It also helps the body conserve energy and recover after activity.

\medterm{Parathyroid} One of the usually four small glands behind the thyroid gland. They release parathyroid hormone, which helps keep calcium and phosphate levels balanced in the blood and bones.

\medterm{Parathyroid Adenoma} A parathyroid adenoma is a usually benign tumor of a parathyroid gland that can secrete excess parathyroid hormone and cause primary hyperparathyroidism with high calcium.

\medterm{Parathyroid Cancer} Parathyroid cancer is a rare malignant tumor of the parathyroid gland that often causes excessive parathyroid hormone and high blood calcium.

\medterm{Parathyroid Carcinoma} A rare malignant tumor of the parathyroid gland that commonly causes excessive parathyroid-hormone secretion, hypercalcemia, kidney stones, bone disease, or related symptoms.

\medterm{Parathyroid Crisis} A life-threatening rise in blood calcium caused by severe overactivity of the parathyroid glands. It may cause extreme thirst, vomiting, constipation, confusion, weakness, abnormal heart rhythm, kidney injury, or coma and requires emergency treatment.

\medterm{Parathyroid Gland} A small gland behind the thyroid that releases parathyroid hormone to regulate calcium and phosphate.

\medterm{Parathyroid Gland Removal} Surgery to remove one or more overactive parathyroid glands, usually for persistent high blood calcium. It can improve symptoms and protect bones or kidneys; low calcium, bleeding, voice-nerve injury, or persistent disease may occur.

\medterm{Parathyroid Glands} Usually four small endocrine glands behind the thyroid that secrete parathyroid hormone to regulate blood calcium and phosphate levels.

\medterm{Parathyroid Hormone} A hormone that regulates calcium and phosphate levels in the blood and supports bone remodeling.
\medterm{Parathyroid Hormone Blood Test} A blood test measuring parathyroid hormone, interpreted with calcium, phosphate, vitamin D, and kidney function to evaluate parathyroid and calcium disorders.

\medterm{Parathyroid Hormone Physiology} When blood calcium falls, the parathyroid glands release more parathyroid hormone. The hormone acts on bone and kidneys and increases active vitamin D, helping raise calcium and maintain the balance needed for nerves, muscles, and bones.

\medterm{Parathyroid Hormone-Related Protein Blood Test} A blood test measuring parathyroid hormone-related protein, used mainly to evaluate humoral hypercalcemia of malignancy and selected disorders of calcium regulation.

\medterm{Parathyroid Hyperplasia} Enlargement of one or more small glands behind the thyroid, causing them to release too much parathyroid hormone. This can raise blood calcium and may occur with kidney disease or inherited multiple-gland disorders.

\medterm{Parathyroid Surgery} Operations on the parathyroid glands, usually performed for selected cases of primary or
secondary hyperparathyroidism. The decision depends on symptoms, calcium levels, kidney involvement, bone disease,
overall health, and current clinical guidance.

\medterm{Parathyroidectomy} Surgery to remove one or more overactive parathyroid glands. It may treat primary or kidney-related hyperparathyroidism when hormone excess causes high calcium, kidney stones, bone disease, or other complications; calcium levels must be monitored afterward.

\medterm{Paratubal Cyst} A usually benign fluid-filled cyst adjacent to the fallopian tube, arising from paratubal or embryologic remnants; most are asymptomatic but large cysts may cause pain or torsion.

\medterm{Paratyphoid Fever} An infection caused by Salmonella Paratyphi bacteria, usually spread through contaminated food or water. It can cause prolonged fever, headache, weakness, abdominal symptoms, and sometimes serious complications without antibiotic treatment.
\medterm{Parenchyma} The working tissue of an organ, distinguished from its supporting framework, blood vessels, and connective tissue. For example, lung parenchyma performs gas exchange, while liver parenchyma processes nutrients and harmful substances.
\medterm{Parenchymal} Parenchymal means relating to the functional cells of an organ, distinguished from its supporting connective tissue, vessels, and framework.

\medterm{Parenchymal Diffuse Lung Disease (Interstitial Lung Disease (Interstitial Pneumonia))} Alternate index label for Interstitial Lung Disease (Interstitial Pneumonia).


\medterm{Parenteral} Parenteral means delivered by a route other than the gastrointestinal tract, commonly by injection or infusion such as intravenous, intramuscular, or subcutaneous administration.

\medterm{Parenteral Nutrition} Nutrition given through a vein when the digestive tract cannot safely absorb enough nutrients. The solution can contain water, glucose, amino acids, fats, vitamins, and minerals. It can cause bloodstream infection, liver problems, or abnormal blood sugar and salts, so it requires monitoring.

\medterm{Paresis} Partial weakness or incomplete paralysis of voluntary muscles. It may affect one limb, one side of the body, both legs, or all four limbs, depending on the nerves, brain, or spinal cord involved.
\medterm{Paresthesia} An unusual sensation such as tingling, prickling, burning, pins and needles, or numbness. Brief symptoms often follow pressure on a nerve, while persistent or spreading symptoms need assessment for nerve or circulation problems.

\medterm{Paricalcitol} Paricalcitol is a vitamin-D receptor activator used to reduce secondary hyperparathyroidism in selected chronic kidney disease, with monitoring for calcium and phosphate changes.

\medterm{Parietal} Relating to a wall of a body cavity or to the parietal bone or lobe of the skull and brain.
\medterm{Parietal Bone} One of two broad bones forming much of the skull's sides and roof. They protect the brain, meet along the top midline, and connect with several neighboring skull bones.
\medterm{Parietal Lobe} A part of the brain that processes touch, pain, temperature, body position, and spatial awareness. Damage may cause numbness, difficulty recognizing objects by touch, poor coordination, or trouble understanding where parts of the body are.

\medterm{Parietal Pericardium} The outer serous layer of the pericardium that lines the inner surface of the fibrous sac surrounding the heart.

\medterm{Parietal Pleura} The pleural layer lining the chest wall, diaphragm, and mediastinum; it is sensitive to pain when inflamed.


\medterm{Parinaud Oculoglandular Syndrome} Conjunctivitis with nearby lymph-node swelling, usually caused by infections such as Bartonella, tularemia, or certain viruses.

\medterm{Parinaud Syndrome} A disorder caused by dysfunction of the dorsal midbrain, characterized by impaired upward gaze, abnormal pupil
responses, and convergence-retraction eye movements. Causes include tumors, stroke,
multiple sclerosis, and hydrocephalus.

\medterm{Paris System For Reporting Urinary Cytology} A standardized framework for urinary cytology that emphasizes detection of high-grade urothelial carcinoma and uses defined categories to improve diagnostic consistency and clinical communication.

\medterm{Parity} Parity is the number of pregnancies that reached a viable gestational age, regardless of whether the infants were born alive or stillborn; it differs from gravidity.

\medterm{Park-Bench Position} A modified lateral position in which the upper body is partly rotated toward prone; it is used for selected posterior brain and neck procedures.

\medterm{Parkinson Disease} Alternate spelling; see \textit{Parkinson's Disease}.



\medterm{Parkinson's Disease} A progressive disorder of the nervous system that mainly affects movement. It develops when nerve cells in a part of the brain called the substantia nigra become damaged or die. These cells normally make dopamine, a chemical messenger that helps produce smooth, purposeful movement. Reduced dopamine can cause tremor, slowness of movement, muscle stiffness, balance or walking problems, a softer voice, and smaller handwriting. Non-movement symptoms may include constipation, loss of smell, tiredness, sleep problems, depression, or changes in thinking. Symptoms and their rate of progression vary from person to person.

\synonyms
Parkinson Disease
Parkinsons Disease (Parkinson's Disease)

\medterm{Parkinsonism} A group of movement problems involving slowness, stiffness, shaking, or balance difficulty. Parkinson disease is one cause, but medicines, other brain disorders, and several less common conditions can produce similar symptoms.

\medterm{Parkinsons Disease (Parkinson's Disease)} Alternate index label for Parkinson's Disease.


\medterm{Paromomycin} Paromomycin is an aminoglycoside antibiotic used mainly in the intestine for selected parasitic infections and sometimes for other infections under specialist care.

\medterm{Paronychia} Infection or inflammation of the skin around a fingernail or toenail, causing pain, redness, swelling, and sometimes pus. It may follow nail biting, injury, a hangnail, or repeated water exposure.

\medterm{Parosteal Osteosarcoma} Parosteal osteosarcoma is a usually lower-grade bone cancer that grows on the outer surface of a bone, often near the back of the thigh.

\medterm{Parotid Gland} The largest major saliva-producing gland, located in front of and below each ear. It releases saliva into the mouth near the upper back teeth and can become swollen from infection, blockage, inflammation, or a tumor.

\medterm{Parotid Neoplasm} An abnormal growth in a parotid salivary gland in front of or below the ear. Most are noncancerous, but some are cancerous; a painless lump, facial weakness, pain, or rapid growth needs assessment.

\medterm{Parotid Tumor} An abnormal growth in a parotid salivary gland. Most parotid tumors are noncancerous, but some are malignant and may cause a lump, pain, or facial-nerve weakness.

\medterm{Parotid Tumors} Abnormal growths arising in the parotid salivary glands, located in front of and below the ears.

\medterm{Parotidectomy} Surgery to remove part or all of a parotid salivary gland, usually because of a tumor. Possible complications include facial-nerve weakness, bleeding, infection, saliva leakage, numbness, or sweating while eating.

\medterm{Parotitis} Inflammation or infection of one or both parotid salivary glands, which lie in front of and below the ears. It can cause painful swelling, fever, dry mouth, or pus near the upper back teeth. Causes include viruses such as mumps, bacteria, blocked ducts, dehydration, and autoimmune disease.
\medterm{Parotitis (Mumps)} Alternate index label for Mumps.

\medterm{Paroxetine Hydrochloride} Paroxetine hydrochloride is an SSRI antidepressant used for depression and several anxiety-related disorders; it can cause sexual side effects and withdrawal symptoms if stopped abruptly.

\medterm{Paroxysm} A sudden, often brief attack or burst of symptoms or body activity that may recur. Examples include an episode of pain, coughing, rapid heartbeat, seizure, or involuntary movement; the cause depends on the symptom.
\medterm{Paroxysmal Atrial Tachycardia} Sudden episodes of rapid, regular atrial rhythm that start and stop abruptly, often producing palpitations, lightheadedness, chest discomfort, or shortness of breath.

\medterm{Paroxysmal Cold Hemoglobinuria} A rare immune disorder in which cold temperatures trigger destruction of red blood cells when the blood warms again. It can cause sudden anemia, fever, back pain, jaundice, or dark urine, often after an infection, and needs medical assessment.

\medterm{Paroxysmal Hemicrania} A rare headache disorder causing brief, repeated attacks of severe one-sided head or facial pain.
\medterm{Paroxysmal Nocturnal Dyspnea} Sudden breathlessness that wakes a person after they have been asleep for a while, often improving after sitting or standing. It commonly occurs with heart failure but may have other heart or lung causes and needs assessment.

\medterm{Paroxysmal Nocturnal Haemoglobinuria} A rare acquired blood disorder in which abnormal blood-forming cells make red blood cells vulnerable to immune destruction. It can cause anemia, dark urine, blood clots, low blood counts, kidney injury, or bone-marrow failure.
\medterm{Paroxysmal Nocturnal Hemoglobinuria} A rare acquired blood disorder in which red cells are destroyed by the immune system. It can cause dark urine, anemia, blood clots in unusual locations, low blood counts, and kidney or other organ problems.

\medterm{Paroxysmal Nocturnal Hemoglobinuria (Paroxysmal Nocturnal Hemoglobinuria PNH)} Alternate index label for Paroxysmal Nocturnal Hemoglobinuria PNH.

\medterm{Paroxysmal Supraventricular Tachycardia} Repeated episodes of a very fast, usually regular heartbeat that start and stop suddenly and begin in the upper chambers or connecting electrical tissues of the heart. It may cause pounding heartbeat, dizziness, chest discomfort, breathlessness, or fainting.
\synonyms
PSVT (Paroxysmal Supraventricular Tachycardia (PSVT))
Supraventricular Tachycardia Paroxysmal (Paroxysmal Supraventricular Tachycardia (PSVT))
Tachycardia Paroxysmal Atrial (Paroxysmal Supraventricular Tachycardia (PSVT))
Tachycardia Paroxysmal Supraventricular (Paroxysmal Supraventricular Tachycardia (PSVT))

\medterm{Pars} A part or portion of an anatomical structure, used in names such as pars interarticularis.
\medterm{Pars Compacta} A dopamine-producing part of the substantia nigra, a movement-control area deep in the brain. Loss of these nerve cells contributes to the slowness, stiffness, and tremor of Parkinson disease.

\medterm{Pars Planitis} A form of inflammation inside the eye, near the edge of the retina. It may cause floaters, blurred vision, and swelling or other retinal complications and is treated according to its severity and cause.


\medterm{Partial} Partial means incomplete or involving only a portion, such as partial loss, partial response, partial thickness, or partial seizure; the specific meaning depends on context.

\medterm{Partial Agonist} A medicine that attaches to a cell receptor and activates it, but produces a smaller effect than a full activator even when all available receptors are occupied. It can reduce the effect of a stronger activator at the same receptor.

\medterm{Partial Anomalous Pulmonary Venous Connection} Alternate index label; see \textit{Partial anomalous pulmonary venous return}.

\medterm{Partial Breast Brachytherapy} Radiation treatment placed inside or next to the area where a breast cancer was removed, usually after breast-conserving surgery. It treats selected early cancers over a shorter course but may cause skin changes, breast firmness, infection, or delayed healing.

\medterm{Partial Breast Radiation Therapy - External Beam} Radiation delivered from outside the body to the part of the breast where a cancer was removed. It may be suitable for selected early cancers and can cause temporary skin irritation, swelling, fatigue, or later breast firmness.

\medterm{Partial Cystectomy} Surgery to remove part of the urinary bladder, usually for a tumor or localized disease while leaving the rest of the bladder in place. The remaining bladder may need monitoring, and urinary or surgical complications can occur.

\medterm{Partial Denture} A removable replacement for some missing teeth, supported by remaining teeth, gums, or both. It may use metal or plastic and can improve chewing, speech, appearance, and tooth stability.

\medterm{Partial Hysterectomy} A common but imprecise term usually referring to supracervical or subtotal hysterectomy, in which the uterine body is removed while the cervix remains. The fallopian tubes and ovaries may or may not be removed separately.

\medterm{Partial Knee Replacement} Surgery replacing only the damaged compartment of the knee with an artificial surface while preserving healthy bone, ligaments, and cartilage. It may relieve pain from localized arthritis; infection, clots, stiffness, loosening, or later arthritis elsewhere can occur.

\medterm{Partial Nephrectomy} Partial nephrectomy removes a diseased part of a kidney while preserving the remaining healthy tissue.
\medterm{Partial Remission} A reduction in signs, symptoms, or measurable disease burden that does not meet criteria for complete remission; the definition depends on the disease and response criteria used.

\medterm{Partial Response} A measurable improvement after treatment in which a tumor or other disease becomes smaller or less active but does not disappear completely. The exact amount of improvement required is defined by the cancer type, test method, or study criteria.

\medterm{Partial Seizure} A seizure beginning in a localized region of one cerebral hemisphere, now called a focal seizure, with symptoms determined by the involved network and awareness either preserved or impaired.

\medterm{Partial Thromboplastin Time} A blood test measuring how long a clot takes to form through several clotting pathways. It helps investigate unusual bleeding and monitor unfractionated heparin; abnormal results may reflect a clotting-factor problem, medicine effect, or another illness.

\medterm{Partial Volume Artifact} A scan distortion that occurs when one small image area contains different tissues. Their signals are blended, which can blur borders or hide a small finding; thinner slices or higher-resolution images may reduce the effect.

\medterm{Partial-Volume Effect} A scan limitation in which a small structure and the surrounding tissue are measured together, making the structure appear blurred or its activity seem too high or low. Thinner slices or correction methods may reduce the effect.

\medterm{Participation} A person’s involvement in life situations and activities, a central concept in rehabilitation and the International Classification of Functioning that extends beyond impairment alone.

\medterm{Particle (Vascular)} A small material placed into a blood vessel to block it during a procedure. Particles may be used to control bleeding or reduce blood flow to a tumor, fibroid, or other abnormal area and must be directed carefully to avoid harming healthy tissue.

\synonyms
Vascular embolization particle

\medterm{Particle Size} Particle size is the diameter or size distribution of particles in a powder, suspension, aerosol, or formulation, influencing dissolution, absorption, stability, and delivery.

\medterm{Particulate Matter} Small solid or liquid particles suspended in air. Fine particles can penetrate the lungs and increase
cardiovascular and respiratory risk.

\medterm{Partograph (Partogram)} A graphic record of cervical change, fetal descent, contractions, and maternal or fetal observations during
labor. It can support communication and recognition of abnormal progress, but findings
must be interpreted with current labor guidance and the overall clinical situation.

\synonyms
Partogram

\medterm{Parturition} The process of childbirth, including cervical ripening, labor contractions, delivery of the fetus, and expulsion of the placenta.

\medterm{Parturition (Childbirth)} The process of giving birth, encompassing three stages: (1) First stage: from the onset
of regular uterine contractions (cervical effacement and dilation) to full dilation (10
cm). Latent phase—slow dilation (0--4 cm) over hours, variable duration. Active
phase—more rapid dilation (4--10 cm), typically 1--1.5 cm/hour in nulliparous women.

\synonyms
Childbirth

\medterm{Parvimonas Micra} An anaerobic gram-positive bacterium found in the mouth and gastrointestinal tract that can cause dental, bone, joint, or bloodstream infections.

\medterm{Parvoviridae} A family of very small DNA viruses; human parvovirus B19 can cause fifth disease, anemia, or fetal complications.

\medterm{Parvovirus} A small DNA virus; parvovirus B19 is the main human pathogen and can cause rash, joint symptoms, or red-cell aplasia.

\medterm{Parvovirus B19} A human virus spread mainly through respiratory droplets. It commonly causes a mild rash illness in children, but can cause joint pain or severe anemia and may harm a fetus during pregnancy.

\medterm{Parvovirus B19 Infection} An infection that may cause fever, a slapped-cheek rash, joint pain, or temporary reduction in red-blood-cell production. It can be serious in people with certain blood disorders or weakened immunity and requires pregnancy assessment after significant exposure.

\medterm{Parvovirus Infection} An infection caused by parvovirus B19 that commonly causes mild fever and rash but may be serious in some people.

\synonyms
Slapped Cheek Disease

\medterm{Passive Exercise} Movement of a person's joint or limb by another person, a machine, or gravity without active muscle effort from the patient. It can maintain flexibility and reduce stiffness when illness or weakness limits movement, but forceful movement can injure tissues.

\medterm{Passive Immunity} Protection produced by transferred antibodies rather than by the person's own immune response, such as maternal antibodies or immunoglobulin treatment.
\medterm{Passive Immunization} Protection produced by giving preformed antibodies, such as immunoglobulin or antivenom, which acts quickly but is temporary because the recipient does not generate durable immune memory.


\medterm{Pasteurella} A genus of bacteria carried by animals that can cause wound infections after bites and other human disease.
\medterm{Pasteurella Aerogenes} Pasteurella aerogenes is a rare Pasteurella species; its recovery from a clinical sample should be interpreted with the source and evidence of infection.


\medterm{Pasteurella Canis} Pasteurella canis is a bacterium found in dogs and cats that can cause wound or soft-tissue infection after an animal bite.

\medterm{Pasteurella Multocida} A bacterium commonly carried by animals that can cause wound infections after bites or scratches and,
less often, respiratory or bloodstream infection.

\medterm{Pasteurellaceae} Pasteurellaceae is a family of bacteria that includes Pasteurella and related animal-associated organisms capable of causing bite-wound and respiratory infections.

\medterm{Pasteurellosis} Infection caused by Pasteurella bacteria, commonly introduced through animal bites or scratches and causing cellulitis, abscesses, or occasionally invasive disease.

\medterm{Pasteurization} A heat treatment reducing microbial load in liquids by heating below boiling point to
inactivate vegetative pathogens and spoilage organisms while preserving quality.

\medterm{Patau Syndrome} A genetic condition caused by an extra copy of chromosome 13. It can cause severe problems involving brain development, growth, the eyes, heart, kidneys, hands or feet, and the lip or palate. Many affected pregnancies end before birth, and surviving infants need complex specialist care.

\medterm{Patch} A flat, circumscribed area of altered skin color larger than approximately 1 cm. It is
the larger counterpart of a macule and may be caused by pigmentary, vascular,
inflammatory, or neoplastic disease.

\medterm{Patch (Vascular)} A piece of biological or synthetic material sewn onto a blood vessel to widen or repair it.

\synonyms
Vascular patch; angioplasty patch

\medterm{Patch Graft} A piece of natural or synthetic material sewn onto a blood vessel to widen or repair it during patch angioplasty. Materials may include a patient's tissue, bovine pericardium, or synthetic PTFE.

\medterm{Patch Test} A skin test used to identify delayed allergic contact reactions to substances applied under adhesive patches.

\medterm{Patch Testing} Patch testing identifies allergens causing allergic contact dermatitis. The standard
series includes 35 allergens. Additional custom trays include occupational and personal
allergens. Patches are applied to the upper back for 48 hours and removed. Reading at 48
hours (initial) and 72 to 96 hours (delayed) identifies positive reactions.

\medterm{Patches} Patches are flat, distinct areas of altered skin color larger than a macule; their appearance, distribution, and associated symptoms help identify the cause.

\medterm{Patchy Hair Loss (Alopecia Areata)} Alternate index label for Alopecia Areata.

\medterm{Patchy Skin Color} Areas of skin that are lighter, darker, redder, or otherwise different in color from the surrounding skin. Causes include vitiligo, birthmarks, inflammation, infection, sun exposure, injury, and changes in blood flow; new or changing patches should be assessed.

\medterm{Patella} The kneecap, a triangular bone at the front of the knee. It protects the joint and helps the thigh muscles straighten the leg by guiding and improving the force of the knee tendon.

\medterm{Patella (Kneecap)} The largest sesamoid bone in the human body, embedded within the quadriceps tendon
at the front of the knee. It improves the mechanical efficiency of knee extension, protects
the front of the joint, and distributes forces across the femoral trochlea.

\synonyms
Kneecap

\medterm{Patellar Dislocation} Complete displacement of the kneecap from the femoral groove, usually laterally, causing acute pain, deformity, swelling, and difficulty extending the knee.

\medterm{Patellofemoral Instability} A tendency of the patella to move abnormally or dislocate from the femoral trochlea.
Risk factors include abnormal trochlear shape, a high-riding patella, ligament insufficiency,
altered limb alignment, and previous dislocation.

\medterm{Patellofemoral Pain Syndrome} Pain around or behind the kneecap, often worsened by running, stairs, squatting, or sitting with the knee bent. It commonly results from repeated loading, muscle weakness, movement changes, or kneecap tracking problems.

\synonyms
Runner's Knee

\medterm{Patellofemoral Syndrome} Pain around or behind the kneecap related to movement and loading of the patellofemoral joint, often without a major structural injury.

\medterm{Patent (Adjective)} Patent means open, unobstructed, or functioning, as in a patent airway, vessel, duct, or congenital fetal channel.

\medterm{Patent Ductus Arteriosus} A blood-vessel connection between the heart’s main lung artery and the aorta that remains open after birth instead of closing. It can send extra blood to the lungs and strain the heart; effects range from none to heart failure.

\medterm{Patent Foramen Ovale} A persistent flap-like opening between the atria that normally permits fetal blood flow
and usually closes after birth. It is often an incidental finding but may allow a clot to pass from the venous to the
arterial circulation in selected circumstances.

\medterm{Patent Urachus} A patent urachus is persistence of the fetal channel between the bladder and umbilicus after birth, allowing urine to drain from the umbilicus and predisposing to infection.

\medterm{Patent Urachus Repair} Surgery to close an abnormal open connection between the bladder and the navel that remains from fetal development. It may stop urine leaking from the umbilicus and prevent infection; treatment depends on the child's anatomy and health.

\medterm{Paternal Uniparental Disomy 14 (UDP 14 Pat) Thoracic Dysplasia} A rare genetic condition resulting from inheritance of both copies of chromosome 14 from
the father, leading to dysregulated expression of imprinted genes on chromosome 14. It
presents with severe thoracic hypoplasia, pulmonary hypoplasia, and characteristic
facial features. It is a lethal condition in most cases.

\synonyms
UDP 14 pat

\medterm{Paternity Testing} A DNA test that assesses whether a person is the biological parent of a child. Samples are usually collected from the child and possible parent, and the result reports either exclusion or a calculated probability of biological relationship.

\medterm{Pathogen} An organism or infectious agent that can enter the body and cause disease. Examples include viruses, bacteria, fungi, and parasites.

\medterm{Pathogenesis} The way a disease begins and develops, including its causes, changes in cells or tissues, and progression over time. Understanding pathogenesis helps explain symptoms, complications, and possible treatment targets.
\medterm{Pathogenic} Pathogenic means capable of causing disease in a susceptible host; pathogenicity depends on the organism or factor, dose, route, and host defenses.

\medterm{Pathogenic Variant} A change in DNA with strong evidence that it causes or contributes to a disease. Its importance depends on the gene, the person’s symptoms and family history, and whether one or both copies of the gene are affected.

\medterm{Pathogenicity} The capacity of a microorganism to cause disease in a susceptible host. Pathogenicity is
a qualitative property, whereas virulence describes the relative degree of
disease-causing capacity among strains or organisms.

\medterm{Pathognomonic} Describes a finding considered highly specific for a particular disease. Truly pathognomonic findings are uncommon; most findings must still be interpreted with the history, examination, and other test results.
\medterm{Pathologic} Pathologic means related to disease or an abnormal structural or functional change, rather than normal anatomy or physiology.

\medterm{Pathologic Dilatation} Abnormal widening of a hollow organ, duct, blood vessel, or other structure beyond its expected size because of obstruction, weakness, pressure, or disease.

\medterm{Pathologic Fracture} A break in a bone weakened by disease, such as cancer spread, myeloma, infection, a benign tumor, or a metabolic bone disorder. It may happen after little or no injury and requires treatment of both the break and its cause.

\medterm{Pathologic Hypertrophy} Enlargement of an organ because its cells become larger, usually after long-term extra workload or hormone stimulation. It can help at first but may eventually impair function, as in an overworked heart.

\medterm{Pathologic Ossification} Abnormal formation of bone in soft tissue or in a location where bone should not develop, as in heterotopic ossification after trauma, neurologic injury, or surgery.

\medterm{Pathologic Process} A pathologic process is an abnormal biological change that produces disease, such as infection, inflammation, ischemia, degeneration, abnormal growth, or immune injury.

\medterm{Pathologic Response} A pathologic response is an abnormal reaction of cells, tissues, or organs to injury, infection, immune activation, toxins, or other stressors that contributes to disease.

\medterm{Pathological Fracture} A fracture occurring in bone weakened by disease, such as osteoporosis, tumor, infection, or metabolic bone disorder, rather than by normal force alone.

\medterm{Pathological Gambling} Legacy, nonpreferred label; see \textit{Gambling Disorder}.

\medterm{Pathological Laughter And Crying} Alternate index label; see \textit{Pseudobulbar affect}.

\medterm{Pathological Stealing} Alternate index label; see \textit{Kleptomania}.

\medterm{Pathologist} A pathologist is a physician who diagnoses disease by examining tissues, cells, blood, or body fluids and may guide laboratory testing, cancer classification, and treatment decisions.

\medterm{Pathology} Study of disease. Includes general pathology (cell injury, inflammation, neoplasia) and systemic pathology (organ-specific diseases).
\medterm{Pathology Artifact} A change introduced while collecting, preserving, cutting, staining, or imaging a sample rather than by disease. Recognizing these false changes prevents an incorrect diagnosis.

\medterm{Pathophysiology} The study of the functional changes produced by disease or injury and the mechanisms through which they cause signs and symptoms.

\medterm{Patient} A person receiving medical care, advice, examination, or treatment. Good care includes listening to the patient's concerns, explaining choices clearly, respecting informed decisions, protecting privacy, and reducing avoidable harm.

\medterm{Patient Advocacy} Acting on behalf of a patient to ensure that the patient's needs, rights, and
preferences are addressed in healthcare.

\medterm{Patient Autonomy} A patient’s right and capacity to make informed voluntary decisions about healthcare after receiving understandable information about options, benefits, risks, and alternatives.


\medterm{Patient Communication} The exchange of information and understanding between a patient, caregivers, and healthcare professionals during care.
\medterm{Patient Education} The process of explaining a health condition, treatment, medicines, exercises, equipment, warning signs, and prevention strategies in a way the patient can understand and use.

\medterm{Patient Observation} Systematic monitoring of a patient’s symptoms, vital signs, examination findings, behavior, and response to treatment to detect deterioration or guide management.

\medterm{Patient Positioning} Placing and supporting a patient in a position that permits safe care, accurate examination or imaging, and adequate comfort.

\medterm{Patient Preparation} The instructions and assessments completed before a test or procedure, such as fasting, medicine review, hydration, consent, or skin preparation when required.

\medterm{Patient Rights} Ethical and legal protections for people receiving care, including informed consent, privacy, access to information, respectful treatment, participation in decisions, and refusal of treatment where lawful.

\medterm{Patient Safety} Preventing avoidable harm during healthcare. It includes accurate identification, safe prescribing and medicine use, infection prevention, clear communication, fall prevention, and checking procedures and results.

\medterm{Patient-Centered Care} Healthcare planned around a person's needs, values, preferences, and goals. It includes clear communication, shared decisions, respect for cultural and personal choices, coordination between professionals, and support for the person and family.

\medterm{Patient-Centered Outcome} An outcome that reflects what matters to the patient, including symptoms, function,
quality of life, treatment burden, and ability to achieve personal goals, rather than
relying only on laboratory or imaging measures.

\medterm{Patient-Controlled Analgesia} A method in which a patient self-administers prescribed doses of analgesic medication
using a programmable pump with preset safety limits. The medication, dose, and lockout interval
are individualized according to the patient's condition and monitoring needs.

\medterm{Patient-Reported Outcome Measure} A validated questionnaire completed directly by a patient to assess symptoms, function,
health status, or quality of life without clinician interpretation of the response.

\medterm{Patlak Plot (Patlak-Rutland Analysis)} A graph used with some nuclear-medicine scans to estimate how quickly a tracked substance moves from the blood into tissue and stays there. It can help measure tissue activity, including activity in some tumors.

\medterm{Pattern Recognition Receptors} Proteins on or inside immune cells that detect common features of germs or damaged tissue. Their activation helps start the body’s early immune response.
\medterm{Pattern-Recognition Receptors} Sensors on or inside immune cells that detect common features of germs or damaged cells. Their activation starts early immune responses, including inflammation and signals that help coordinate protection and tissue repair.

\medterm{Pavor Nocturnus (Night Terrors)} Episodes of sudden intense fear, screaming, and autonomic arousal during sleep, usually arising from deep non-REM sleep and often followed by little or no memory.

\synonyms
Night Terrors

\medterm{Paxillin} Paxillin is a focal-adhesion adaptor protein that links integrin signaling to the actin cytoskeleton and helps regulate cell adhesion, migration, and mechanical sensing.

\medterm{Pazopanib} An oral tyrosine-kinase inhibitor targeting vascular endothelial growth factor and other receptors, used for selected advanced renal-cell carcinoma and soft-tissue sarcoma.

\medterm{PBA} Alternate index label; see \textit{Pseudobulbar affect}.

\medterm{PBG Urine Test} A urine test measuring porphobilinogen, used urgently to evaluate suspected acute hepatic porphyria during compatible neurovisceral or abdominal attacks.

\medterm{PCBs} PCBs, or polychlorinated biphenyls, are persistent industrial chemicals that accumulate in tissues and can affect the nervous, endocrine, immune, and reproductive systems.

\medterm{PCL Injuries} Sprains or tears of the posterior cruciate ligament, a strong ligament inside the knee. They may follow a blow to the shin or forceful bending and can cause pain, swelling, stiffness, or a feeling that the knee gives way.

\medterm{PCO} PCO may mean polycystic ovaries or posterior capsular opacification; the intended condition depends on the clinical context.

\medterm{PCO Disease} Polycystic ovary syndrome, a hormonal and metabolic disorder involving irregular ovulation, androgen excess, and sometimes polycystic-appearing ovaries.

\medterm{PCOS} A common hormone disorder that can cause irregular ovulation or periods, higher levels of androgen hormones, acne, or excess facial hair. Some people have ovaries with many small follicles, but that finding is not required for diagnosis.

\medterm{PCP} Pneumocystis jirovecii pneumonia (formerly Pneumocystis carinii pneumonia); opportunistic infection in immunocompromised patients, especially transplant recipients.
\medterm{PDA} Alternate index label; see \textit{Patent ductus arteriosus}.

\medterm{PDE5 Inhibitors} Medicines that improve blood flow into the penis by relaxing blood vessels and are mainly used for erectile dysfunction. They must not be combined with nitrate medicines because blood pressure can fall dangerously low.

\medterm{PE Anticoagulation} Anticoagulant treatment for pulmonary embolism to prevent clot extension and recurrence. Options
include direct oral anticoagulants, heparins, and vitamin K antagonists; the choice and
duration depend on bleeding risk, pregnancy, kidney function, cancer, provoking factors,
and other clinical features.

\medterm{Peak Expiratory Flow Rate} The fastest speed at which a person can blow air out after taking a full breath. A handheld meter measures it in litres per minute; lower or changing values can help monitor airway narrowing, especially in asthma.

\medterm{Peak Flow Meter} A handheld device that measures peak expiratory flow rate, the fastest speed of air expelled after a full breath, usually reported in liters per minute. It provides a simple measure of large-airway flow but does not measure the full range of airflow and lung volumes assessed by spirometry.
\medterm{Peanut Allergy} An immune reaction to peanut proteins that can cause symptoms ranging from hives to life-threatening anaphylaxis.

\synonyms
Allergy, Peanut

\medterm{Peanut Oil} Peanut oil is oil extracted from peanuts; refined oil usually contains little allergenic protein, whereas unrefined products may trigger reactions in susceptible people.

\medterm{Pearly Penile Papules} Small, benign, dome-shaped papules arranged circumferentially around the glans penis,
representing normal anatomical variants requiring no treatment.

\medterm{Peau D'Orange} Peau d’orange is dimpled skin resembling an orange peel, caused by edema tethered around hair follicles; it can occur with breast cancer, infection, or lymphatic obstruction.


\medterm{Pectinate Line} A boundary inside the anal canal where the lining, nerve supply, blood vessels, and lymph drainage change. Conditions above and below this line may cause different symptoms and may be treated differently.

\medterm{Pectoral} Relating to the chest muscles or the front of the chest. The pectoral muscles move the arm and help stabilize the shoulder. Pain in this area may come from muscle strain, the ribs, breast tissue, or organs inside the chest.
\medterm{Pectoral Girdle} The pectoral girdle consists of the clavicles and scapulae, attaching the upper limbs to the trunk and providing broad surfaces for shoulder and arm movement.

\medterm{Pectoralis} Relating to the chest muscles, especially pectoralis major and pectoralis minor. These muscles help move the arm, support the shoulder, and connect the chest wall with the upper limb.
\medterm{Pectoralis Major Muscle} The large chest muscle that moves the arm across the body, turns it inward, and helps with pushing. It attaches to the collarbone, breastbone, ribs, and upper arm.

\medterm{Pectoralis Minor Muscle} A small chest muscle beneath the larger pectoralis major. It helps draw the shoulder blade forward and downward, hold it against the ribs, and assist forceful breathing; tightness can contribute to rounded shoulders.

\medterm{Pectoralis Muscles} The pectoralis major and minor muscles of the chest, which move and stabilize the shoulder and assist with movements such as pushing, adduction, and breathing.

\medterm{Pectoriloquy} Abnormally clear transmission of spoken sounds through the chest wall, often suggesting lung consolidation or another change in lung tissue.

\medterm{Pectus Carinatum} A chest-wall shape present from childhood in which the breastbone pushes outward, sometimes called pigeon chest. It may cause no health problems, but some people have chest discomfort, reduced exercise tolerance, or emotional distress.

\medterm{Pectus Excavatum} A chest-wall shape present from childhood in which the breastbone is pushed inward, sometimes called funnel chest. It may cause no symptoms, but severe cases can affect exercise, breathing, posture, or emotional wellbeing.

\synonyms
Sunken Chest


\medterm{Pectus Excavatum Repair} Surgical correction of a sunken sternum and adjacent costal cartilage, performed when chest-wall deformity causes significant symptoms, physiologic impairment, or psychosocial impact.


\medterm{Pediatric Advanced Life Support} Emergency training and guidance for recognizing and treating life-threatening breathing, circulation, or heart problems in infants and children. It covers resuscitation, airway support, shock, abnormal rhythms, and care after recovery.

\synonyms
PALS (Pediatric Advanced Life Support); Pediatric Resuscitation Protocol

\medterm{Pediatric AIDS} Pediatric AIDS is advanced HIV infection in a child, with severe immune suppression or AIDS-defining illness; antiretroviral therapy and prevention of opportunistic infection are essential.

\medterm{Pediatric Anesthesia} Anesthesia for infants and children uses age- and weight-appropriate medicines, equipment, and monitoring. Children’s airways and breathing reserves differ from adults, so they can develop low oxygen levels or changes in blood pressure quickly and need specialized care.

\medterm{Pediatric Appendicitis} Acute inflammation of the vermiform appendix, usually caused by luminal obstruction.
Children may have evolving periumbilical-to-right-lower-quadrant pain, anorexia,
vomiting, fever, and guarding, but presentation can be atypical in young children;
timely surgical and antimicrobial assessment reduces perforation risk.

\medterm{Pediatric Assessment Triangle} A rapid first look at a sick child that does not require touching the child. Clinicians observe appearance, effort of breathing, and skin color or circulation. The findings show how urgently help is needed and guide the next examination and treatment.

\synonyms
PAT (Pediatric Assessment Triangle); Pediatric Visual Triage Triangle

\medterm{Pediatric Brain Tumors} Tumors that develop in a child's brain or nearby tissues. They may be noncancerous or cancerous and can cause headache, vomiting, seizures, vision changes, weakness, or changes in behavior; treatment depends on the type, location, size, and age of the child.

\medterm{Pediatric Cancer} Cancer occurring in infants, children, or adolescents. Common types include leukemia, brain tumors, lymphoma, neuroblastoma, and Wilms tumor; prognosis varies by type and stage.

\medterm{Pediatric Dentistry} Dental care for infants, children, and adolescents, including prevention and treatment of tooth decay, care of developing teeth, growth-related problems, behavior support, and adaptations for children with special healthcare needs.

\medterm{Pediatric Dermatology} The medical specialty concerned with skin, hair, and nail disorders in infants, children,
and adolescents. It includes congenital lesions, inflammatory diseases, infections, vascular lesions, and tumors.

\medterm{Pediatric Dermatology Conditions} Skin conditions in children may include birthmarks, eczema, infections, vascular growths, pigment changes, and inherited disorders. Most are manageable, but a rapidly changing, painful, bleeding, or infected lesion needs medical assessment.

\medterm{Pediatric Diabetic Ketoacidosis} A life-threatening complication of diabetes in which very high blood sugar and a lack of insulin make the blood acidic. A child may urinate and drink excessively, vomit, have abdominal pain, deep breathing, fruity-smelling breath, sleepiness, or coma and needs emergency care.

\medterm{Pediatric Early Warning Score} A hospital scoring system that combines a child’s behavior, breathing, heart rate, circulation, and need for oxygen to detect worsening illness early. A rising score prompts closer observation or urgent review; it supports, but does not replace, clinical judgment.

\synonyms
PEWS; Pediatric Early Warning System

\medterm{Pediatric Eczema} Atopic dermatitis in children commonly begins in infancy with facial and extensor
involvement; flexural disease predominates in older children. Pruritus, xerosis, sleep
disturbance, and scratching are characteristic.

\medterm{Pediatric Fever} An elevated body temperature in a child, usually caused by infection but sometimes by inflammation,
medicines, or environmental heat. Evaluation depends on age, appearance, duration, and associated findings.

\medterm{Pediatric Health} The physical, mental, emotional, and social wellbeing of infants, children, and adolescents. Care includes prevention, vaccinations, growth and development checks, treatment of illness, and support for the child and family.

\medterm{Pediatric Heart Surgery} An operation to repair or improve a child's heart or major blood vessels, often for a congenital heart defect. Procedures may close abnormal openings, improve blood flow, or replace a valve; recovery and risks depend on the defect and operation.


\medterm{Pediatric Meningitis} Inflammation or infection of the coverings of a child’s brain and spinal cord. Fever, headache, stiff neck, vomiting, light sensitivity, confusion, seizures, or a rash may occur; babies may be irritable, feed poorly, or have a bulging soft spot. Suspected meningitis is an emergency.

\medterm{Pediatric Nursing} Pediatric nursing provides age-appropriate care to infants, children, and adolescents, accounting for growth, development, family involvement, communication, safety, and medication differences.

\medterm{Pediatric Obesity} Excess body fat in a child or adolescent, assessed by comparing weight and height with other children of the same age and sex. It can affect breathing, joints, blood sugar, sleep, and emotional wellbeing; assessment should be supportive and non-stigmatizing.

\medterm{Pediatric Obstructive Sleep Apnea} Repeated narrowing or blockage of a child’s throat during sleep, causing pauses in breathing, snoring, restless sleep, daytime behavior or attention problems, and sometimes poor growth. Enlarged tonsils or adenoids are common causes.

\medterm{Pediatric Orthopedics} The branch of orthopedics focused on bone, joint, muscle, and movement problems in infants, children, and adolescents. It includes developmental conditions, limb differences, fractures, sports injuries, spine deformities, and growth-related disorders.

\medterm{Pediatric Pneumonia} Infection or inflammation of the lung tissue in an infant, child, or adolescent. It may cause cough, fever, fast or difficult breathing, chest pain, or low oxygen and can be caused by bacteria, viruses, fungi, or aspiration.

\medterm{Pediatric Sepsis} Sepsis occurring in an infant, child, or adolescent, in which the body’s response to infection causes organ dysfunction. Signs vary with age and may include abnormal temperature, breathing difficulty, unusual sleepiness, poor feeding, or reduced urine.

\medterm{Pediatric Poisoning} Harm caused when a child is exposed to a medicine, chemical, plant, gas, or another toxic substance. Effects depend on the substance, amount, route, and time of exposure and may include vomiting, sleepiness, breathing problems, seizures, or organ injury.
\medterm{Pediatric Ophthalmology} The subspecialty caring for eye and vision disorders in infants, children, and adolescents, including strabismus, amblyopia, congenital abnormalities, refractive error, and childhood eye disease.

\medterm{Pediatric Cataract} Clouding of the natural lens in an infant or child. It may be present at birth or develop after birth and can interfere with visual development, causing amblyopia or reduced vision.

\medterm{Pediatric Glaucoma} Glaucoma occurring in an infant or child, often involving abnormal development or blockage of the eye’s drainage system. Increased eye pressure can enlarge or damage the eye and impair vision.

\medterm{Persistent Fetal Vasculature} Persistence after birth of fetal blood vessels that normally regress before or soon after birth. It usually affects one eye and may cause a small eye, cataract, abnormal pupil reflection, retinal changes, or reduced vision.

\medterm{Congenital Nasolacrimal Duct Obstruction} Blockage of the tear-drainage channel present at birth, commonly causing persistent watering or discharge from one or both eyes. Most cases occur because a membrane at the lower end of the duct has not opened normally.

\medterm{Acquired Nasolacrimal Duct Obstruction} Blockage of the tear-drainage system that develops after birth. It may cause excessive tearing, discharge, or repeated inflammation of the lacrimal sac and can result from aging, infection, injury, inflammation, scarring, or a mass.

\medterm{Pediatric Rehabilitation} Therapy that helps children with conditions present from birth, injuries, illness, or developmental disabilities gain skills and independence. Care considers growth, family goals, school participation, communication, movement, learning, and daily activities.

\medterm{Pediatric Shock} A dangerous condition in which a child’s organs do not receive enough blood flow or oxygen. Causes include severe fluid loss, infection, heart failure, or a blockage. Early signs may include a fast heartbeat, pale or cold skin, unusual sleepiness, confusion, or little urine.

\medterm{Pediatric Skin Tumors} Skin growths in children include harmless birth-related or blood-vessel growths, cysts, and, rarely, cancer. A lump that grows quickly, bleeds, ulcerates, changes color, or causes pain should be assessed by a clinician.

\medterm{Pediatric Sleep Apnea} Pediatric sleep apnea is repeated interruption or reduction of breathing during sleep, commonly from enlarged tonsils or adenoids, causing snoring, restless sleep, behavior changes, or impaired growth.

\medterm{Pediatric Toxicology} The medical study and treatment of poisoning or harmful exposure in children. Common exposures include medicines, iron, household chemicals, and plants. Symptoms vary with the substance and may include vomiting, sleepiness, breathing problems, seizures, or burns; urgent poison-center or emergency advice may be needed.
\medterm{Pediatrician} A medical doctor who specializes in the health and medical care of children from birth
through adolescence (typically up to age 16-18).

\synonyms
Paediatrician

\medterm{Pediatrics} The branch of medicine concerned with the health and medical care of infants, children,
and adolescents from birth to age 18 (or up to age 16 in the UK).

\synonyms
Paediatrics

\medterm{Pedicle} A pedicle is a stalk or attachment containing vessels, nerves, or other supporting structures, such as the base of a polyp, flap, or vertebral arch.

\medterm{Pediculosis} Infestation with lice, small insects that live on the scalp, body, or pubic hair. It causes itching and spreads mainly through close contact or shared personal items.
\medterm{Pediculosis Capitis} Infestation of the scalp by head lice, causing itching and attached eggs or lice in the hair.

\medterm{Pediculus Humanus Capitis} A blood-sucking parasite commonly known as the louse (plural, lice) that can cause an
itchy scalp; infestations are highly contagious and especially common in school-age
children.

\medterm{Pedigree} A family diagram that shows relationships and who has a particular trait or disease. Doctors and genetic counselors use it to look for inheritance patterns and estimate the chance that relatives may be affected.

\medterm{Pedodontics} The branch of dentistry concerned with infants, children, and adolescents.
\medterm{Pedophilia} A continuing sexual interest in children who have not reached puberty. The term describes an interest and is not synonymous with pedophilic disorder; sexual contact with a child is abuse and cannot be consensual.

\medterm{Pedophilic Disorder} A mental-health disorder involving recurrent, intense sexual fantasies, urges, or behaviors toward children who have not reached puberty. Diagnostic criteria include clinically significant distress or impairment, or acting on the urges, together with required age and duration criteria.

\medterm{Pedunculate} Pedunculate describes a structure, lesion, or polyp attached by a narrow stalk rather than a broad base.

\medterm{Pedunculated} Attached to the skin or another surface by a narrow stalk, rather than having a broad base.

\medterm{Peeling Skin} Loss or shedding of the outer layer of skin, which may result from dryness, sunburn, irritation, allergy, infection, inflammation, or a systemic illness.

\medterm{PEEP} Pressure kept in the airways at the end of breathing out, usually by a ventilator or breathing machine. It helps keep the lung air sacs open and improve oxygen levels, but too much pressure can injure the lungs or lower blood flow.


\medterm{PEF} PEF, or peak expiratory flow, is the fastest airflow produced during a forceful exhalation and is used chiefly to monitor variable airway obstruction such as asthma.

\medterm{PEG Tube} A flexible feeding tube placed through the abdominal wall into the stomach with camera guidance. It provides nutrition, fluids, or medicines when swallowing is unsafe or inadequate, and requires regular skin care.




\medterm{Peliosis Hepatis} A usually uncommon condition in which multiple blood-filled spaces form within the liver. It may be linked to medicines, hormones, toxins, infections, cancer, or other diseases. It often causes no symptoms but can rarely cause liver failure or life-threatening bleeding if a space ruptures.


\medterm{Pelizaeus-Merzbacher Disease} A rare inherited disorder in which the brain does not make myelin, the protective covering around nerve fibers, normally. It usually begins in infancy with poor muscle tone, unusual eye movements, and delayed development, followed by stiffness, balance problems, and increasing movement difficulty. There is no cure; care supports movement, feeding, breathing, and communication.

\synonyms
PMD; PLP1-Related Hypomyelinating Leukodystrophy

\medterm{Pellagra} A disease caused by too little niacin, leading to skin changes, diarrhea, and problems with thinking.
\medterm{Pellet} A pellet is a small compacted mass of medication, material, or tissue; drug pellets may be formulated for implantation or controlled release.

\medterm{Pelvic} Relating to the pelvis, the bony ring and associated organs, muscles, vessels, and nerves in the lower trunk.

\medterm{Pelvic Binder} A circumferential device applied around the pelvis to reduce motion and hemorrhage from
suspected unstable pelvic-ring injury. It should be centered over the greater
trochanters and used with definitive imaging and hemorrhage-control management.

\medterm{Pelvic Bone} The ring of bones formed by the paired hip bones, sacrum, and coccyx that supports the trunk, protects pelvic organs, and transmits weight to the lower limbs.

\synonyms
Hip Bone
Hip Bones

\medterm{Pelvic Cavity} The space inside the lower part of the body enclosed by the pelvic bones. It contains the bladder, rectum, reproductive organs, blood vessels, nerves, and supporting muscles.

\medterm{Pelvic Congestion} Enlargement or pooling of blood in the veins of the pelvis, which may cause pelvic pressure, discomfort, or pain.

\medterm{Pelvic Congestion Syndrome} Chronic pelvic pain associated with enlarged or inefficient pelvic veins, often worsening with prolonged standing,
before menstruation, or after intercourse. Diagnosis requires evaluation for other causes of pelvic pain and may
involve pelvic imaging; treatment is individualized by gynecology or vascular specialists.

\medterm{Pelvic CT Scan} Computed tomography of the pelvis using x-rays and computer reconstruction to evaluate bones, organs, soft tissues, masses, trauma, infection, and bleeding.

\medterm{Pelvic Diaphragm} The group of muscles and supporting tissue forming the floor of the pelvis. It supports the bladder, bowel, and reproductive organs and helps control urine and stool.

\medterm{Pelvic Examination} Clinical inspection and palpation of the external genitalia, vagina, cervix, uterus, and adnexa to evaluate bleeding, pain, discharge, lesions, pregnancy-related findings, or infection.

\medterm{Pelvic Exenteration} Very extensive surgery removing some pelvic organs and nearby tissues, usually for selected cancers that have returned or grown locally when less extensive treatment cannot control them. Recovery is prolonged and reconstruction may be needed.

\medterm{Pelvic Floor} A group of muscles and connective tissues supporting the pelvic organs and controlling urinary and fecal continence.

\medterm{Pelvic Floor Muscles} Muscles at the bottom of the pelvis that support the bladder, bowel, and reproductive organs and help control urine, stool, and sexual function.

\medterm{Pelvic Floor Dysfunction} Weakness, overactivity, poor coordination, or injury of the muscles and connective tissues supporting the bladder,
bowel, and reproductive organs. Symptoms may include urinary or fecal problems, pelvic pain, constipation,
pain with sex, or pelvic organ prolapse; assessment may involve pelvic-health clinicians.

\medterm{Pelvic Floor Health And Disorders} The pelvic floor is a group of muscles and supporting tissues that hold the bladder, bowel, and reproductive organs in place and help control urine, stool, and sexual function. Weakness, tightness, injury, or poor coordination can cause pain, leakage, constipation, or prolapse.

\medterm{Pelvic Floor Muscle Training Exercises} Pelvic floor muscle exercises strengthen the muscles supporting the bladder, bowel, and reproductive organs and may improve selected forms of urinary or fecal leakage.

\medterm{Pelvic Girdle} The ring of bones connecting the lower limbs to the trunk, consisting of the hip bones, sacrum, and associated joints.

\medterm{Pelvic Girdle Pain In Pregnancy} Pain arising from the pelvic joints and surrounding tissues during pregnancy, often affecting the pubic symphysis, sacroiliac joints, hips, or lower back. It may impair walking, turning in bed, or climbing stairs.

\medterm{Pelvic Infection} An infection involving pelvic organs or tissues, including pelvic inflammatory disease, urinary infection, abscess, or postoperative infection.

\medterm{Pelvic Inflammatory Disease} Infection and inflammation of the upper female genital tract (uterus, fallopian tubes,
ovaries, and peritoneum), most commonly caused by sexually transmitted infections
(Chlamydia trachomatis, Neisseria gonorrhoeae). Presents with lower abdominal pain,
cervical motion tenderness, adnexal tenderness, and fever.

\medterm{Pelvic Inlet} The upper opening of the bony birth canal, formed by the lower spine, hip bones, and front pelvic joint. Its shape and size help determine how a baby can pass through the pelvis during childbirth.

\medterm{Pelvic Kidney} A kidney that remains in the pelvis rather than ascending to its usual lumbar position during development; it may be asymptomatic or associated with abnormal drainage, stones, infection, or vascular anatomy.

\medterm{Pelvic Laparoscopy} A minimally invasive procedure using a camera inserted through small abdominal incisions to inspect and treat pelvic organs and structures.

\medterm{Pelvic Lymphadenectomy} Pelvic lymphadenectomy is surgical removal of lymph nodes in the pelvis to check for or treat spread of cancer.

\medterm{Pelvic Mass} A lump or abnormal enlargement in the pelvis arising from reproductive, urinary, gastrointestinal, lymphatic, or other tissues.

\medterm{Pelvic Neoplasm} A pelvic neoplasm is an abnormal growth in the pelvis, arising from organs, bones, soft tissues, or lymph nodes; it may be benign or malignant.

\medterm{Pelvic Organ Prolapse} Dropping or bulging of the bladder, uterus, rectum, or nearby tissue into or toward the vagina because supporting muscles and tissues have weakened. It may cause pressure, a vaginal bulge, urinary or bowel problems, or discomfort during sex.

\medterm{Pelvic Organ Prolapse Quantification System} A standard examination used to measure how far the vagina, uterus, bladder, or rectum has moved downward. Measurements are taken while bearing down and describe the severity of prolapse to help plan treatment.

\synonyms
POP-Q System; POP-Q Classification

\medterm{Pelvic Outlet} The lower opening of the bony pelvis and birth canal, bordered by the pubic bone, sitting bones, and tailbone. Its shape and size affect passage through the pelvis during childbirth.

\medterm{Pelvic Pain} Pain in the lower abdomen, pelvis, or area between the genitals and anus. Causes include reproductive, urinary, bowel, muscle, nerve, and other conditions. Sudden severe pain, pregnancy, fever, fainting, vomiting, or heavy bleeding needs urgent assessment.

\medterm{Pelvic Pain During Early Pregnancy} Pain in the pelvis or lower abdomen during early pregnancy, caused by normal uterine changes, miscarriage, ectopic pregnancy, ovarian disease, infection, or other conditions. Severe or one-sided pain, shoulder pain, fainting, fever, or vaginal bleeding requires emergency assessment.

\medterm{Pelvic Pain, Chronic} Alternate index label; see \textit{Chronic pelvic pain}.


\medterm{Pelvic Spleen} A spleen located unusually low in the abdomen or pelvis because its supporting tissues are loose or absent. It may feel like a pelvic lump and can twist, lose its blood supply, or rupture, causing sudden pain and requiring urgent care.

\medterm{Pelvic Support Problems, Uterine Prolapse} Alternate index label; see \textit{Uterine prolapse}.

\medterm{Pelvic Ultrasound - Abdominal} Ultrasound imaging of pelvic organs through the lower abdominal wall, usually with a full bladder, to evaluate the uterus, ovaries, bladder, and surrounding structures.

\medterm{Pelvic Viscus} An organ located in the pelvic cavity, such as the bladder, rectum, uterus, ovaries, prostate, or related structures.

\medterm{Pelvimeter} An instrument used to estimate the size and shape of the pelvis, especially during obstetric assessment. Pelvic measurements are only one part of judging whether vaginal birth is likely and are interpreted with the clinical situation.
\medterm{Pelvimetry} Measurement of the size and shape of the pregnant person's pelvis, usually for selected childbirth concerns. Results are considered with the baby's position, size, labor progress, and overall clinical situation.

\medterm{Pelvis} The bony ring and associated soft tissues at the base of the trunk, connecting the spine with the lower limbs.
\medterm{Pelvis MRI Scan} Magnetic resonance imaging of the pelvis used to evaluate organs, muscles, bones, joints, nerves, blood vessels, and soft-tissue disease without ionizing radiation.

\medterm{Pelvis X-Ray} Plain radiographic imaging of the pelvic bones and hips used to evaluate fractures, alignment, arthritis, deformity, and selected bone lesions.

\medterm{Pemberton Sign} Facial redness or bluish discoloration, neck-vein swelling, or breathing difficulty that appears when both arms are raised. It may suggest that an enlarged thyroid or another chest mass is pressing on major veins or the airway.

\medterm{Pemberton’S Sign} Facial redness or bluish discoloration, neck-vein swelling, or breathing difficulty that appears when both arms are raised. It may suggest that an enlarged thyroid or another chest mass is pressing on major veins or the airway.

\medterm{Pemetrexed} Pemetrexed is a folate-antimetabolite chemotherapy drug that inhibits several enzymes involved in nucleotide synthesis, used especially for mesothelioma and nonsquamous non-small-cell lung cancer; folic acid and vitamin B12 reduce toxicity.

\medterm{Pemetrexed Disodium} The disodium salt of pemetrexed, an antifolate chemotherapy medicine used with other treatment for selected non-small-cell lung cancer and mesothelioma.

\medterm{Pemoline} Pemoline is a stimulant formerly used for attention-deficit disorders; serious liver toxicity led to withdrawal or severe restriction in many countries.

\medterm{Pemphigoid} A group of immune-system diseases that cause firm, often itchy blisters when antibodies attack the deeper layers of skin or mucous membranes. The type, location, age, and associated illness determine treatment.
\medterm{Pemphigoid Bullous (Bullous Pemphigoid)} Alternate index label for Bullous Pemphigoid.

\medterm{Pemphigoid Gestationis} An autoimmune blistering disease of pregnancy characterised by intensely pruritic
urticarial plaques and tense vesicles or bullae, typically arising in the second or
third trimester around the umbilicus and spreading to the trunk and extremities.

\medterm{Pemphigus} A group of autoimmune diseases causing fragile blisters when antibodies attack connections between skin cells.
\medterm{Pemphigus Vulgaris} A rare immune-system disease that causes fragile blisters and painful raw areas on the skin and mouth or other moist linings. Blisters break easily; untreated disease can cause fluid loss, infection, and difficulty eating or drinking.

\medterm{Penbutolol} Penbutolol is a beta-blocker formerly used for high blood pressure; it can slow the heart rate and worsen asthma in susceptible people.

\medterm{Pencil Eraser Swallowing} Swallowing a small pencil eraser usually causes little poisoning risk, but it may lodge in the airway or digestive tract. Choking, breathing difficulty, persistent vomiting, or severe pain requires emergency care.

\medterm{Pencil Swallowing} Swallowing a pencil or part of one can injure the mouth, throat, or digestive tract and may block the airway. Do not force food or drink; breathing trouble, bleeding, or severe pain needs emergency care.


\medterm{Pendular} Moving or swinging freely like a pendulum; pendular nystagmus has a similar speed in both directions.
\medterm{Pendulous} Pendulous means hanging loosely from a supporting structure, as in a pedunculated lesion, drooping tissue, or an organ suspended by a stalk.

\medterm{Penetrance} The proportion of people with a particular gene change who develop the related trait or condition. Penetrance may be complete, incomplete, or influenced by other factors.
\medterm{Penetrating Abdominal Trauma} Injury to the abdomen caused by a projectile, sharp object, or other penetrating
force. It may damage abdominal organs or blood vessels and requires urgent assessment for internal bleeding or organ
injury; management depends on the mechanism, examination, stability, and imaging findings.

\medterm{Penetrating Wounds} Injuries caused when an object enters the skin and damages underlying tissue. They may injure blood vessels, nerves, organs, or bones and require assessment for bleeding, contamination, retained objects, and infection.
\medterm{Penicillin} A family of antibiotics that kill bacteria by stopping them from making cell walls. Different types treat different bacterial infections; they do not work against viruses. Allergy, resistance, and kidney problems can affect the choice of medicine.
\medterm{Penicillin Allergy} An immune reaction to penicillin antibiotics that may cause rash, swelling, breathing difficulty, or anaphylaxis.

\synonyms
Allergy, Penicillin
\medterm{Penicillin G} Penicillin G is an injectable penicillin antibiotic used for infections such as syphilis and selected streptococcal or other susceptible bacterial infections.

\medterm{Penicillin V} Penicillin V is an oral penicillin antibiotic used for selected susceptible infections, including some throat and dental infections.

\medterm{Penicillium} A genus of molds that includes species used to make medicines and foods and species that can cause allergy or opportunistic infection.

\medterm{Penile} Penile means relating to the penis, the external male organ involved in urination and sexual reproduction.

\medterm{Penile Adhesions} Abnormal attachment of penile skin to the glans, often after circumcision or during childhood. They may cause irritation or conceal part of the glans and are assessed and treated according to severity.

\medterm{Penile Cancer} An uncommon cancer of the penis, usually arising from surface skin cells. It may appear as a persistent sore, lump, bleeding area, or skin change; HPV, smoking, foreskin problems, and chronic inflammation can increase risk.

\medterm{Penile Augmentation} A cosmetic or reconstructive procedure intended to increase penile length, girth, or both. Techniques and results vary, and complications may include scarring, deformity, infection, altered sensation, or erectile problems.


\medterm{Penile Disease} A disorder affecting the penis, including infection, inflammation, erectile or structural disease, trauma, skin disease, or malignancy.

\medterm{Penile Edema} Swelling of the penis caused by inflammation, infection, trauma, allergic reaction, impaired lymphatic drainage, or a constricting ring or retracted foreskin. Sudden painful swelling with an irreducible foreskin is an emergency.

\medterm{Penile Erection} Stiffening and enlargement of the penis when blood fills erectile tissue and is temporarily held there. It can occur with sexual or other stimulation; persistent difficulty getting or keeping an erection may be erectile dysfunction.


\medterm{Penile Fracture} A tear in the tough covering of an erect penis, usually after forceful bending. A popping sound, immediate loss of erection, swelling, and bruising may occur; urgent surgery can prevent lasting problems.

\medterm{Penile Hygiene} Routine cleaning and drying of the penis and foreskin, with gentle retraction only when it moves freely and replacement afterward. Good hygiene helps reduce irritation and infection without harsh soaps or forceful retraction.

\medterm{Penile Implant} A surgically implanted device used to produce penile rigidity in selected people with erectile dysfunction that has not responded to other
treatments.

\medterm{Penile Infection} Infection of the penis or foreskin caused by bacteria, fungi, viruses, or parasites, producing pain, redness, discharge, sores, or swelling.

\medterm{Penile Intraepithelial Neoplasia} A precancerous change in the surface cells of the penis. It may appear as a persistent red, white, scaly, or wart-like patch on the glans or foreskin and may be linked to HPV or chronic inflammation. Examination, biopsy, treatment, and follow-up reduce the risk of invasive cancer.

\synonyms
PeIN; Penile Precancerous Lesion

\medterm{Penile Leukoplakia} A persistent white patch on penile skin or mucosa that may reflect irritation, inflammation, or a precancerous disorder and requires assessment.

\medterm{Penile Prosthesis} An implanted device used to produce an erection when erectile dysfunction has not improved with other treatments. It may be inflatable or bendable. Surgery can cause infection, bleeding, pain, device failure, or erosion into nearby tissue.

\medterm{Penile Torsion} Rotation of the penis around its longitudinal axis, usually a congenital condition in which the glans or urethral opening is rotated. Symptomatic or marked torsion may require urologic evaluation and surgical correction.

\medterm{Penile Trauma} Injury to the penis caused by blunt force, bending, penetration, crushing, burns, bites, or medical instruments. It may damage the skin, erectile tissue, blood vessels, or urethra. A penile fracture is a tear in the covering of the erectile tissue, usually after forceful bending of an erect penis. It may cause pain, swelling, bruising, bleeding, or difficulty passing urine.

\medterm{Penis} An external genital organ that contains the outlet for urine and can deliver sperm during sexual activity. It becomes firm when filled with blood during an erection and contains sensitive tissue, skin, blood vessels, and nerves.
\medterm{Penis Carcinoma} Penis carcinoma is cancer of the penis, most often arising from skin or squamous cells and sometimes linked to HPV infection.


\medterm{Penis Disorder (Priapism (Penis Disorder))} Alternate index label for Priapism (Penis Disorder).

\medterm{Penis Disorder Balantis (Balanitis (Penis Disorder))} Alternate index label for Balanitis (Penis Disorder).

\medterm{Penis Pain} Penis pain may result from infection, injury, inflammation, urinary problems, or sexual conditions; severe pain, swelling, discharge, or inability to urinate requires prompt assessment.

\medterm{Penis, Erection Of The} Penile rigidity produced by increased arterial inflow and reduced venous outflow after sexual or reflex stimulation; persistent difficulty is erectile dysfunction.

\medterm{Penis, Small} A penis that is substantially below expected stretched length for age, usually assessed as micropenis in a newborn or child after excluding measurement and developmental variation.

\medterm{Penoscrotal Transposition} A congenital difference in which the penis and scrotum are positioned abnormally relative to each other, often associated with hypospadias or other genital and urinary anomalies.

\medterm{Pentad Of TTP} The five findings once considered typical of thrombotic thrombocytopenic purpura: low platelets, destruction of red blood cells, neurologic symptoms, kidney problems, and fever. Most patients do not have all five, so urgent blood testing is essential when TTP is suspected.


\medterm{Pentamidine} An antimicrobial medicine used for certain parasitic infections and Pneumocystis pneumonia, with potentially serious side effects.
\medterm{Pentamidine Isethionate} An antiprotozoal and antimicrobial drug used in selected cases of Pneumocystis pneumonia, leishmaniasis, trypanosomiasis, and other infections; toxicity includes hypoglycemia, kidney injury, and arrhythmia.

\medterm{Pentatrichomonas Hominis} Pentatrichomonas hominis is a flagellated intestinal protozoan generally regarded as nonpathogenic; finding it in stool usually indicates exposure to fecal contamination.

\medterm{Pentazocine} An opioid pain medicine with mixed agonist and antagonist effects, used for some moderate or severe pain. It can cause drowsiness, dependence, and slowed breathing and may trigger withdrawal in a person dependent on other opioids.
\medterm{Pentazocine Overdose} Taking too much pentazocine, an opioid pain medicine, can cause extreme sleepiness, confusion, slow breathing, pinpoint pupils, seizures, or coma. Call emergency services immediately and give naloxone when available and appropriate.

\medterm{Pentetic Acid} A chelating compound that binds metal ions and is used in selected radiopharmaceuticals, contrast agents, or laboratory applications.

\medterm{Pentobarbital Overdose} Taking too much pentobarbital, a sedative, can severely slow brain activity and breathing, causing low blood pressure, coma, or death. It requires immediate emergency treatment, especially with alcohol or other sedatives.

\medterm{Pentosan Polysulfate} A semi-synthetic sulfated polysaccharide used for selected bladder pain syndrome and historically as an anticoagulant; prolonged use can cause pigmentary maculopathy and bleeding.

\medterm{Pentose} A pentose is a five-carbon sugar, including ribose in RNA and deoxyribose in DNA, and participates in nucleotide and carbohydrate metabolism.

\medterm{Pentose Phosphate Pathway} A cell pathway that uses glucose to make NADPH, which protects cells from chemical damage, and ribose, which helps build DNA and RNA. It is especially important in the liver and red blood cells.

\medterm{Pentoxifylline} A methylxanthine derivative that improves red-cell flexibility and is used in some countries for peripheral vascular disease; benefit varies by condition.


\medterm{Penumbra} Brain tissue surrounding an area of infarction that is functionally impaired but potentially salvageable if blood flow is restored promptly.

\medterm{Peplomycin} Peplomycin is a bleomycin-related chemotherapy medicine used in some countries for selected cancers; lung toxicity is an important risk.

\medterm{Peppermint} Mentha piperita, used as a flavoring and in some preparations for dyspeptic symptoms; concentrated oil can be hazardous.
\medterm{Peppermint Oil Overdose} Swallowing a large amount of peppermint oil can cause nausea, vomiting, abdominal pain, dizziness, allergic reactions, seizures, or lung injury if aspirated. Contact poison control promptly and do not induce vomiting.

\medterm{Pepsin} A stomach enzyme that begins breaking proteins into smaller fragments. It is released as inactive pepsinogen and activated by stomach acid, helping digestion before food enters the small intestine.
\medterm{Pepsin A} Pepsin A is a stomach protease activated by acid that begins digestion of dietary proteins into smaller peptides.

\medterm{Pepsinogen} An inactive precursor released by gastric chief cells and converted to pepsin in acidic conditions.
\medterm{Peptic Esophagitis} Inflammation and injury of the esophageal lining caused mainly by repeated gastroesophageal reflux of acid or bile, producing heartburn, chest discomfort, swallowing pain, or strictures.

\medterm{Peptic Ulcer} An open sore in the lining of the stomach or duodenum, most commonly caused by \textit{Helicobacter pylori} infection or NSAID use.

\synonyms
Stomach Ulcer (Peptic Ulcer)
Ulcer, Peptic

\medterm{Peptic Ulcer Bleeding} Haemorrhage caused by erosion of an artery or other blood vessel beneath a gastric or
duodenal peptic ulcer. It is a common cause of upper gastrointestinal bleeding and may
present with haematemesis, coffee-ground vomiting, melaena, or haemodynamic instability.

\medterm{Peptic Ulcer Disease} One or more open sores in the lining of the stomach or first part of the small intestine. Common causes include Helicobacter pylori infection and anti-inflammatory medicines; ulcers may cause burning pain, bleeding, or a dangerous perforation.



\medterm{Peptidase} An enzyme that breaks peptide bonds in proteins or peptides.
\medterm{Peptide} A short chain of amino acids linked by peptide bonds. Peptides can form parts of proteins or act independently as hormones, neurotransmitters, immune signals, or medicines.
\medterm{Peptide Biosynthesis} The process by which cells make peptides, short chains of amino acids, by linking amino acids according to genetic instructions. Some peptides act as hormones, neurotransmitters, or signals between cells.

\medterm{Peptide Bond} The chemical link joining two amino acids. Repeated peptide bonds form short peptides and long proteins, which cells use for structure, chemical reactions, communication, transport, and many other functions.

\medterm{Peptide Hormone} A hormone made of amino acids that binds a cell-surface receptor and regulates processes such as metabolism, growth, reproduction, stress, and fluid balance.


\medterm{Peptide Metabolism} The body's production, breakdown, and use of peptides. These processes regulate signaling molecules, hormones, digestion, and protein turnover.

\medterm{Peptide Transport} The movement of peptides into, out of, or within cells and tissues through transport proteins, vesicles, or other cellular mechanisms. It helps control peptide signaling, digestion, and removal from the body.

\medterm{Peptide Yy} Peptide YY (PYY) is a hormone released mainly by cells in the intestine after eating. It helps reduce appetite and slow movement of food through the digestive tract and is being studied in relation to body weight and metabolism.

\medterm{Peptidoglycan} A strong mesh-like material forming much of the bacterial cell wall. Some antibiotics kill bacteria by preventing them from making or repairing this protective layer.

\medterm{Peptococcaceae} A family of anaerobic gram-positive cocci found in the human microbiota and sometimes associated with infection.

\medterm{Peptococcus} A genus of anaerobic gram-positive cocci that can contribute to mixed infections, especially in damaged or poorly oxygenated tissue.

\medterm{Peptococcus Niger} An anaerobic gram-positive coccus found in the human microbiota and rarely associated with mixed infection.
\medterm{Peptoniphilus} A genus of anaerobic gram-positive cocci found on skin and mucosal surfaces that can cause opportunistic infection.

\medterm{Peptoniphilus Harei} An anaerobic gram-positive coccus occasionally recovered from wounds, abscesses, or other clinical specimens. An anaerobic Gram-positive bacterium sometimes found in wounds and abscesses, where it may contribute to mixed infections.
\medterm{Peptoniphilus Indolicus} An anaerobic gram-positive coccus found in the genital and other microbiota that can occasionally contribute to mixed infection.

\medterm{Peptostreptococcus} A group of anaerobic gram-positive cocci associated with dental, skin, abdominal, pelvic, and other mixed infections.

\medterm{Peracetic Acid} A strong oxidizing disinfectant and sterilant used for equipment and surfaces; concentrated exposure irritates or burns skin, eyes, and airways.

\medterm{Peracute} Peracute describes disease or a clinical event that develops extremely rapidly, often over minutes to hours, with little time for adaptation or diagnosis.

\medterm{PERC Rule} A checklist used by emergency clinicians for people already judged very unlikely to have a blood clot in the lungs. If all criteria are reassuring, testing may be avoided; it is not suitable for home use or for people with a strong suspicion of clot.

\synonyms
Pulmonary Embolism Rule-out Criteria; PERC Score

\medterm{Percentile} A percentile indicates the percentage of a reference population with a value below a measurement, such as height, weight, or test score.

\medterm{Perception} The brain’s process of organizing and interpreting sensory information into an experience of objects, events, body sensations, or the environment.

\medterm{Perceptual} Relating to the process of detecting and interpreting sensory information.
\medterm{Perceptual Disorder} A clinically significant disturbance in the way sensory information is interpreted, including altered visual, auditory, tactile, or body perception.

\medterm{Perceptual Distortion} An inaccurate or altered interpretation of a real sensory stimulus, which may occur with neurologic disease, psychiatric illness, drugs, migraine, seizures, or sensory impairment.

\medterm{Perceptual Masking} Reduction in the ability to detect or identify one sensory stimulus because another stimulus overlaps it in time, space, or processing.

\medterm{Percussion} Physical examination in which the body surface is tapped and the resulting sound or vibration is interpreted to estimate underlying air, fluid, or solid tissue.
\medterm{Percussion Note} The sound heard when the chest is gently tapped during examination. A normal lung sounds hollow; a dull sound may suggest fluid or infection, while an unusually hollow sound may suggest excess air or a collapsed lung.

\medterm{Percutaneous} Performed through the skin, usually using a needle, catheter, wire, or other instrument. This approach can diagnose or treat disease while avoiding a larger open operation, but bleeding and infection remain possible.
\medterm{Percutaneous Biopsy} Removal of tissue through the skin using a needle or other instrument, often guided by ultrasound, CT, or another imaging method.

\medterm{Percutaneous Coronary Intervention} A procedure that guides a thin tube through a blood vessel to a narrowed heart artery, then widens it with a balloon and often supports it with a stent. It can restore blood flow during a heart attack or treat selected severe coronary disease.

\medterm{Percutaneous Discectomy} Percutaneous discectomy removes or decompresses part of a herniated spinal disc through a small needle or incision.

\medterm{Percutaneous Diskectomy} A minimally invasive procedure removing part of a herniated spinal disk through a small skin opening to relieve nerve pressure.

\medterm{Percutaneous Drainage} Placement of a needle or catheter through the skin to drain a fluid collection or obstructed body space, often with imaging
guidance.

\medterm{Percutaneous Electrical Nerve Stimulation} Electrical stimulation delivered through the skin to reduce pain signals.

\synonyms
PENS

\medterm{Percutaneous Endoscopic Gastrostomy} Placement of a feeding tube through the abdominal wall into the stomach using a camera passed through the mouth. It may provide long-term nutrition, fluids, or medicines when swallowing is unsafe or inadequate.

\medterm{Percutaneous Kidney Procedures} Procedures performed through a small needle or tube passed through the skin into or near the kidney. They may obtain a biopsy, drain infection or urine, or remove stones; bleeding, infection, urine leakage, or injury to nearby organs are possible.

\medterm{Percutaneous Nephrolithotomy} A procedure that removes a large or complex kidney stone through a small opening in the back. A narrow instrument reaches the kidney, breaks or removes the stone, and may leave a temporary drainage tube; bleeding, infection, and injury are possible.

\medterm{Percutaneous Nephrolithotripsy} Percutaneous nephrolithotripsy removes or fragments kidney stones through a tract made through the skin into the kidney, usually for larger or complex stones.

\medterm{Percutaneous Nephrostomy} An image-guided interventional procedure in which a flexible drainage tube (nephrostomy catheter) is placed through the skin of the flank directly into the renal pelvis of the kidney. Using ultrasound and fluoroscopy under local anesthesia, an interventional radiologist or urologist punctures the kidney's urine-collecting system, threads a guide wire, and advances the catheter, which is secured to an external urine-collection bag. Percutaneous nephrostomy is primarily performed to decompress an obstructed, infected, or damaged kidney (hydronephrosis) when a ureteral stone, pelvic tumor, pregnancy, or stricture prevents normal urine outflow to the bladder. Decompressing the blocked kidney provides immediate relief from severe flank pain, treats or prevents life-threatening urosepsis, and preserves kidney function until the underlying obstruction can be surgically resolved.

\synonyms
PCN; Nephrostomy Tube Placement; Percutaneous Renal Drainage

\medterm{Percutaneous Transhepatic Cholangiogram} A percutaneous transhepatic cholangiogram uses a needle through the liver to inject contrast into bile ducts for imaging and may guide drainage or stent placement.

\medterm{Percutaneous Transhepatic Cholangiography} PTC is imaging of the bile ducts after contrast is injected through a needle passed through the skin and liver; drainage or stenting may follow.
\medterm{Percutaneous Transluminal Angioplasty} A catheter-based procedure that inflates a balloon inside a narrowed artery to improve blood flow. A stent may also be placed; risks include bleeding, clotting, vessel injury, and re-narrowing.

\medterm{Percutaneous Transluminal Coronary Angioplasty} Balloon-based treatment that widens a narrowed coronary artery, often followed by stent placement. A catheter procedure using balloon inflation to widen a narrowed coronary artery, sometimes followed by stent placement.
\synonyms
PTCA; balloon angioplasty

\medterm{Percutaneous Umbilical Blood Sampling} A procedure that uses ultrasound to guide a needle into the umbilical cord and collect a fetal blood sample during pregnancy. It may investigate anemia, infection, or genetic problems and carries risks such as bleeding, infection, and pregnancy loss.


\medterm{Percutaneously Inserted Central Catheter - Infants} A thin catheter inserted through a small vein in an infant’s arm or leg and advanced toward a large vein near the heart. It provides longer-term access for fluids, nutrition, or medicines, but requires careful monitoring for infection, blockage, bleeding, and displacement.

\medterm{Perforant Pathway} A group of nerve fibers carrying information from a brain area involved in memory to the hippocampus. It helps the brain form and retrieve memories and can be affected in some dementias and seizure disorders.

\medterm{Perforated Bowel} A perforated bowel is a hole through the intestinal wall allowing contents to leak into the abdomen, causing peritonitis and potentially life-threatening sepsis.

\medterm{Perforated Eardrum} Alternate index label; see \textit{Ruptured eardrum (perforated eardrum)}.

\medterm{Perforation} A hole or complete tear through the wall of an organ or tissue. It can release contents into surrounding spaces, cause severe infection or bleeding, and often requires urgent medical treatment.
\medterm{Perforation Of A Viscus} A hole through the wall of a hollow abdominal organ, such as the stomach, small intestine, or colon. Contents can leak into the abdomen and cause severe infection; ulcers, inflammation, poor blood flow, and injury are possible causes.

\medterm{Performance} How effectively a person, organ, treatment, test, or system carries out a defined task. Examples include exercise performance, heart performance, test performance, and response to treatment.

\medterm{Performance Anxiety} Anxiety about sexual performance that can interfere with sexual function.

\medterm{Performance Status} A rating of how well a person can move, work, and perform daily activities while ill. Clinicians use it to estimate prognosis, decide whether some treatments are safe, and track change over time, especially during cancer care.

\medterm{Perfusion} The flow of blood through the small vessels of a tissue or organ. Good perfusion supplies oxygen and nutrients and removes waste; poor perfusion can cause pain, weakness, tissue injury, or organ failure.

\medterm{Perfusion Defect} An area receiving unusually little blood flow, identified by an imaging or function test. It may reflect a blocked artery, narrowed vessel, poor heart output, inflammation, or damaged tissue.

\medterm{Perfusion Imaging} Imaging that shows how blood flows through an organ or tissue. It can help assess problems such as stroke, reduced heart blood flow, or tumor blood supply.
\medterm{Perfusion PET} Positron-emission tomography that measures tissue blood flow using a suitable radioactive tracer. Cardiac perfusion PET can show reduced blood flow, viability, and sometimes the effects of treatment.

\medterm{Perianal} Located around the anus; the term describes the skin, glands, pain, inflammation, abscesses, fistulas, and other findings in that region.

\medterm{Perianal Abscess} A localized collection of pus near the anus caused by infected anal glands or surrounding tissue, producing severe throbbing pain, swelling, fever, and sometimes a fistula.

\medterm{Perianal Glands} Small glands near the anus that contribute secretions to the local tissue; obstruction or infection can contribute to perianal abscesses and fistulas.

\medterm{Perianal Streptococcal Cellulitis} A bacterial skin infection around the anus, usually caused by group A Streptococcus. It may cause bright-red skin, pain, itching, discharge, or pain during bowel movements.

\medterm{Periapical Abscess} An acute or chronic localized collection of pus at the apex of a tooth's root, arising
from pulpal infection and necrosis (usually due to untreated dental caries or trauma).

\medterm{Periapical Granuloma} A chronic inflammatory lesion at the apex of a tooth, usually resulting from pulpal
necrosis and persistent apical infection. It may be asymptomatic and is diagnosed
through clinical, radiographic, and sometimes histopathologic assessment.

\medterm{Peri-Implantitis} Inflammation around a dental implant with progressive loss of the supporting bone. It may cause bleeding, swelling, pus, pain, or implant looseness and requires dental assessment.

\medterm{Periappendiceal} Periappendiceal means located around the appendix, as in inflammation, fluid, an abscess, or fat stranding near the appendiceal tissue.

\medterm{Periarteritis Nodosa} An older name for polyarteritis nodosa, a necrotizing vasculitis of medium-sized arteries that can cause skin, nerve, muscle, kidney, and abdominal ischemic symptoms without typical glomerulonephritis.

\medterm{Periarticular} Periarticular means located around a joint, including nearby tendons, bursae, ligaments, capsule, and soft tissues rather than the joint surface itself.

\medterm{Pericardial} Relating to the pericardium, the thin two-layered sac surrounding the heart. A small amount of fluid between its layers allows the heart to beat smoothly, while excess fluid or inflammation can interfere with heart function.

\medterm{Pericardial Cavity} The thin fluid-containing space between the layers of the sac surrounding the heart. A small amount of fluid lets the heart move smoothly; excess fluid or blood can press on the heart and impair its pumping.
\medterm{Pericardial Disease (Pericarditis)} Alternate index label for Pericarditis.

\medterm{Pericardial Effusion} Extra fluid collecting in the sac around the heart. A small amount may cause no symptoms, but rapid or large accumulation can compress the heart, lower blood pressure, and become life-threatening.
\medterm{Pericardial Effusion Etiology} Fluid around the heart may result from infection, cancer, kidney failure, low thyroid activity, autoimmune disease, heart injury, surgery, radiation, trauma, or medicines. Tests and the person’s history help identify the cause.

\medterm{Pericardial Fluid Culture} A microbiologic culture of fluid from the pericardial space to identify bacteria, fungi, or other organisms in suspected infectious pericarditis.

\medterm{Pericardial Fluid Gram Stain} A microscopic stain of pericardial fluid that can rapidly suggest bacterial organisms, although a negative result does not exclude infection.

\medterm{Pericardial Inflammation} Alternate index label; see \textit{Pericarditis}.

\medterm{Pericardial Knock} An extra early heart sound caused when a stiff, scarred sac around the heart suddenly stops the heart from filling. It may suggest constrictive pericarditis and must be interpreted with other examination and test findings.

\medterm{Pericardial Mesothelioma} A rare malignant tumor arising from the mesothelial lining around the heart, which can cause effusion, constriction, or impaired cardiac filling.

\medterm{Pericardial Metastases} Cancer that has spread to the sac around the heart from another part of the body. It may cause inflammation, fluid or blood around the heart, breathlessness, chest discomfort, or impaired heart filling.

\medterm{Pericardial Stripping} Surgery to remove part or most of the thickened or scarred sac around the heart. It may relieve constrictive pericarditis, in which the stiff sac prevents the heart from filling normally; bleeding, rhythm problems, and other major surgical risks are possible.

\medterm{Pericardial Tamponade} Dangerous pressure on the heart caused by rapid or large fluid or blood buildup in the sac around it. The heart cannot fill normally, causing breathlessness, low blood pressure, fainting, or shock and requiring emergency treatment.

\medterm{Pericardial Wall} The layered tissue forming the protective sac around the heart. Its inner surfaces produce a small amount of lubricating fluid that lets the heart beat with little friction.
\medterm{Pericardial Window} An operation that creates an opening in the sac around the heart so extra fluid can drain into a nearby space and be absorbed. It may prevent repeated fluid buildup and allow a tissue sample to be taken.

\synonyms
Subxiphoid Pericardiostomy; Pleuropericardial Window; Pericardiostomy

\medterm{Pericardiectomy} Surgery to remove part or most of the sac around the heart when thickening or scarring prevents the heart from filling normally. It is major surgery and may improve constrictive pericarditis but carries bleeding and other serious risks.

\medterm{Pericardiocentesis} A needle or catheter is used to remove fluid or blood from the sac around the heart, usually to relieve dangerous pressure or investigate an effusion. Imaging guidance improves safety, but bleeding, abnormal rhythm, organ injury, or infection can occur.

\medterm{Pericardiostomy} A surgical opening made in the pericardium to drain fluid or prevent recurrent accumulation around the heart.

\medterm{Pericarditis} Inflammation of the sac around the heart. It often causes sharp chest pain that worsens with deep breathing or lying down and improves when sitting up and leaning forward. It may be caused by infections, autoimmune diseases, heart injury, kidney disease, cancer, or some medicines. Severe cases can cause fluid buildup around the heart.

\synonyms
Inflamed Pericardium
Pericardial Inflammation
Pericardial Disease (Pericarditis)

\medterm{Post-Myocardial Infarction Pericarditis} Inflammation of the sac around the heart that occurs after a heart attack. It may begin soon after the heart attack or develop later as part of a delayed immune reaction.

\medterm{Pericarditis - Constrictive} A chronic condition in which a thickened, scarred, or calcified pericardium restricts ventricular filling and causes systemic venous congestion and heart-failure symptoms.

\medterm{Pericardium} A two-layered protective sac surrounding the heart and beginnings of major blood vessels, reducing friction during each heartbeat.

\medterm{Perichondrial Hematoma (Cauliflower Ear)} Alternate index label for Cauliflower Ear.

\medterm{Perichondritis} Perichondritis is inflammation or infection of tissue surrounding cartilage, often affecting the outer ear after trauma or piercing and potentially damaging cartilage.

\medterm{Perichondritis Of The Ear} Infection or inflammation of the auricular perichondrium, often after trauma, piercing,
or surgery. Pseudomonas aeruginosa is an important cause; delayed treatment can produce
cartilage necrosis and permanent deformity.

\medterm{Perichondrium} A tough layer of tissue surrounding most cartilage. It supplies nutrients, contains cells that help cartilage grow and repair, and can become inflamed or infected, especially around the outer ear.
\medterm{Pericoronitis} Inflammation of the soft tissue around the crown of a partially erupted tooth, most
often a mandibular third molar. Food retention and bacterial accumulation may cause
pain, swelling, trismus, fever, or spread into deeper spaces.

\medterm{Pericytes} Pericytes are contractile support cells wrapped around small blood vessels that help regulate capillary stability, blood flow, repair, and the blood-brain barrier.

\medterm{Periductal Mastitis} Inflammation around the major subareolar ducts, often associated with smoking, causing a
painful periareolar mass, discharge, or recurrent breast abscess.


\medterm{Perilymph} Fluid in the bony labyrinth of the inner ear that surrounds the membranous labyrinth and helps transmit sound and balance signals.

\medterm{Perilymphatic Fistula} An abnormal leak between the fluid spaces of the inner ear and the middle ear, often after injury, pressure changes, or surgery. It may cause dizziness, hearing loss, ringing, or symptoms triggered by sound or pressure.

\medterm{Perimetry} A test that measures the full field of vision, including central and side vision. It helps detect and monitor problems such as glaucoma, optic-nerve disease, retinal disease, and disorders affecting the visual pathways in the brain.

\medterm{Perimenopausal} Perimenopausal describes the years of hormonal and menstrual transition before and around menopause, when cycles, bleeding, hot flashes, sleep, and fertility may change.

\medterm{Perimenopause} The transition time in a woman’s life that begins when ovaries produce less estrogen and
menstruation becomes less frequent, and ends when the ovaries no longer produce eggs and
menstruation stops.

\medterm{Perinatal} Perinatal refers to the period around birth, conventionally including late pregnancy and the first days or weeks after delivery depending on the definition used.

\medterm{Perinatal Asphyxia} Too little oxygen or inadequate breathing before, during, or soon after birth. Severe cases can injure the brain and other organs; rapid newborn assessment, breathing support, and treatment of the cause are essential.

\medterm{Perinatal Death} Death of a fetus or newborn during the late pregnancy, birth, or early newborn period; the exact time range varies by definition.

\medterm{Perinatal Psychiatry} The subspecialty focused on psychiatric disorders during pregnancy and the postpartum
period. Common conditions include perinatal depression, perinatal anxiety, perinatal
OCD, and postpartum psychosis. Medication selection must balance maternal mental health
with fetal or neonatal exposure risks.

\medterm{Perinatal Transmission} Transmission of an infection or condition from mother to infant during pregnancy, labor, delivery, or breastfeeding, as with HIV, hepatitis B, syphilis, or certain genetic and congenital disorders.

\medterm{Perinatologist} A perinatologist is an obstetrician specializing in high-risk pregnancy, fetal medicine, prenatal diagnosis, and complications affecting the pregnant person or fetus.

\medterm{Perindopril} Perindopril is an ACE inhibitor used for high blood pressure and cardiovascular risk reduction; it can cause cough, high potassium, or rarely angioedema.

\medterm{Perineal} Relating to the area between the external genitals and anus, which can be affected by injury or infection.

\medterm{Perineal Laceration} Tearing of the perineal tissues during vaginal delivery. The injury is graded by depth:
it may involve skin, perineal muscle, the anal sphincter, or the rectal mucosa.

\medterm{Perineal Laceration Classification} A system grading tears during childbirth by depth: first degree affects skin, second degree also affects muscle, third degree involves the anal-control muscle, and fourth degree extends into the rectum. Accurate examination and repair reduce later pain and bowel-control problems.

\medterm{Perineal Protection} Measures used during the pushing stage of childbirth to reduce severe tearing of the tissue between the vagina and anus. They may include controlled delivery of the baby’s head, support of the perineum, a warm compress, and a selective approach to episiotomy.

\medterm{Perineal Tear} A split in the tissue between the vagina and anus during vaginal birth. First-degree tears affect skin, second-degree tears reach muscle, third-degree tears involve the anal-control muscle, and fourth-degree tears extend into the rectum; severe tears need careful repair and follow-up.

\medterm{Perineum} The area between the external genitalia and the anus. It contains skin, muscles, nerves, and connective tissue and may be affected by injury, infection, childbirth, or surgery.
\medterm{Perineural} Perineural means around a nerve; the term may describe normal connective tissue, inflammation, tumor spread, or imaging findings surrounding a nerve.

\medterm{Perineuronal Satellitosis} A microscopic pattern in which supporting brain cells cluster around a nerve cell. It can occur in some brain tumors, inflammation, or reactive changes and has meaning only with the surrounding tissue findings.

\medterm{Perinuclear Space} The narrow space between the two membranes surrounding a cell nucleus. It connects with part of the cell’s internal membrane network and helps form the nuclear envelope.


\medterm{Periodic Abstinence} Periodic abstinence is avoiding intercourse during estimated fertile days to reduce pregnancy likelihood; cycle timing is imperfect and does not prevent infections.

\medterm{Periodic Acid} A chemical used to stain tissue samples. In the periodic-acid–Schiff test, it helps highlight sugars, mucus, certain fungi, and supporting membranes so they can be seen under a microscope.

\medterm{Periodic Breathing} Repeated cycles of shallow or absent breaths followed by deeper breaths. It can occur normally for short periods in some infants or during sleep, but persistent episodes may reflect a brain, heart, or lung disorder.

\medterm{Periodic Fever, Aphthous Stomatitis, Pharyngitis, Adenitis Syndrome} A recurrent fever syndrome, usually beginning in childhood, characterized by periodic fever with some combination of mouth ulcers, sore throat, and swollen neck lymph nodes; it is commonly abbreviated PFAPA.

\medterm{PFAPA Syndrome} Abbreviation for periodic fever, aphthous stomatitis, pharyngitis, and adenitis syndrome, a recurrent fever disorder that usually begins in childhood.

\medterm{Periodic Limb Movement Disorder} Repeated involuntary limb movements during sleep that may disrupt rest.

\synonyms
PLMD
Sleep Disorder Periodic Limb Movement (Periodic Limb Movement Disorder)

\medterm{Periodic Paralysis Syndrome} Periodic paralysis syndromes are rare inherited or acquired disorders causing episodic muscle weakness associated with abnormal potassium handling or thyroid disease. Attacks may follow rest after exercise, meals, or stress; ECG, potassium testing, and genetic or endocrine evaluation guide care.

\medterm{Periodontal Abscess} A localized accumulation of pus within periodontal tissues, commonly associated with a
periodontal pocket, foreign material, or impaired drainage. It may cause pain, swelling,
suppuration, tooth mobility, and systemic symptoms and requires urgent dental
assessment.

\medterm{Periodontal Cyst} A periodontal cyst is a fluid-filled epithelial-lined lesion near the tooth-supporting tissues, with effects ranging from incidental discovery to local bone destruction.

\medterm{Periodontal Disease} Infection or inflammation of the tissues supporting teeth, including gums and jawbone. It can cause bleeding, swelling, bad breath, loose teeth, bone loss, and tooth loss, but treatment can slow progression.
\medterm{Periodontal Disease, Gingivitis} Alternate index label; see \textit{Gingivitis}.

\medterm{Periodontal Disease, Periodontitis} Alternate index label; see \textit{Periodontitis}.

\medterm{Periodontal Dressing} A protective material placed over the gums after selected dental or periodontal procedures. It can protect the surgical area while it heals, reduce irritation, and support comfort; it should be kept clean and reviewed if it becomes loose, painful, foul-smelling, or causes swelling.

\medterm{Periodontal Flap Surgery} Dental surgery in which the gum is gently lifted so a clinician can clean infection and deposits below the gumline and, when needed, reshape accessible bone before replacing the gum.
\medterm{Periodontal Ligament} Strong tissue fibers that attach a tooth’s root to the surrounding jawbone. They cushion chewing forces, help the body sense pressure on the tooth, and can become inflamed or damaged in gum disease or grinding.

\medterm{Periodontal Pocket} An abnormally deep space between a tooth and the surrounding gum caused by loss of attachment. Plaque can collect there, leading to bleeding, infection, bone loss, loose teeth, and eventual tooth loss.

\medterm{Periodontal Probing} A dentist gently measures the space between the gum and tooth with a marked instrument. Deeper spaces, bleeding, and loss of support can indicate gum disease; measurements help assess severity and monitor treatment.

\medterm{Periodontics} The dental specialty focused on preventing, diagnosing, and treating gum disease and tissues supporting the teeth.
\medterm{Periodontist} A periodontist is a dentist specializing in prevention, diagnosis, and treatment of gum disease and supporting tissues, including periodontal surgery and implants.

\medterm{Periodontitis} A chronic inflammatory disease caused by dysbiotic dental plaque that destroys the periodontal ligament and alveolar bone supporting the teeth. It can cause bleeding gums, periodontal pockets, tooth mobility, recession, bad breath, and tooth loss; severity and progression vary and require dental evaluation.

\synonyms
Periodontal Disease, Periodontitis

\medterm{Periodontitis (Gum Disease)} Alternate index label for Gum Disease.

\medterm{Periodontium} The tissues that hold and support a tooth, including gums, ligament, root covering, and jawbone. Disease of these tissues can cause bleeding, gum recession, loose teeth, bone loss, and tooth loss.
\medterm{Periodontium Disease} Disease of the tissues supporting the teeth, including gingiva, periodontal ligament, cementum, and alveolar bone; periodontitis can cause attachment loss and tooth loss.

\medterm{Perionychium} The skin and soft tissue surrounding a fingernail or toenail, including the nail folds and nearby cuticle; infection or inflammation causes paronychia.

\medterm{Perioperative} Perioperative means relating to the period before, during, and after an operation, including preparation, anesthesia, recovery, and complication prevention.

\medterm{Perioral Dermatitis} A papulopustular eruption around the mouth, sparing the vermilion border. It is more
common in young women and associated with topical corticosteroid use. Other triggers
include fluoridated toothpaste, cosmetics, and inhaler steroids.

\medterm{Periorbital} Located around the eye socket or eyelids. Periorbital swelling may result from injury, allergy, infection, or fluid retention.

\medterm{Periorbital Cellulitis} A bacterial infection of the eyelid and tissues around the eye that causes redness, swelling, and tenderness. It must be distinguished from orbital cellulitis, which can threaten vision and spread more deeply.

\medterm{Periorbital Edema} Swelling of the tissues around the eye, caused by allergy, infection, trauma, venous or lymphatic congestion, kidney disease, thyroid eye disease, or other systemic conditions.

\medterm{Periorbital Infection} Infection of the tissues around the eye; it can cause swelling and redness and must be distinguished from orbital cellulitis, which threatens vision.

\medterm{Periosteal Osteosarcoma} A cancerous bone tumor that develops beneath the thin tissue covering a bone, usually along the shaft of a long bone. It often makes cartilage and causes pain or swelling, requiring specialist cancer treatment.

\medterm{Periosteum} A blood-vessel-rich, pain-sensitive tissue layer covering the outer surface of bone, except where joint cartilage covers it. It supplies nutrients, supports repair, and provides attachment for tendons and ligaments.
\medterm{Periostitis} Inflammation of the periosteum, the pain-sensitive tissue covering the outer surface of bone. It may cause localized bone pain, tenderness, or swelling and can result from repeated stress, infection, injury, or another disease.
\medterm{Peripancreatic Malignancy} Cancer involving the tissues around the pancreas, either by spreading from the pancreas or from another organ. It may cause abdominal or back pain, weight loss, jaundice, or newly developed diabetes and usually needs detailed imaging and specialist care.

\medterm{Peripartum} Relating to the period around childbirth, including late pregnancy, labor, delivery, and the early weeks after birth. Clinicians may use a more specific time range for a particular condition or study.

\medterm{Peripartum Cardiomyopathy} A weakening and enlargement of the heart muscle that develops near the end of pregnancy or during the months after birth, without another clear cause. It can cause breathlessness, swelling, fatigue, or rapid heartbeat and needs prompt cardiac care.

\medterm{Peripartum Cardiomyopathy Prognosis} Recovery of ventricular function varies widely. Prognosis is influenced by the initial ejection
fraction, ventricular size, symptoms, treatment, complications, and follow-up; persistent
dysfunction requires ongoing cardiac care.

\medterm{Peripheral} Peripheral means away from the center or central nervous system; in clinical use it may describe nerves, vessels, lesions, circulation, or symptoms outside the brain, spinal cord, or central structures.

\medterm{Peripheral Arterial Line - Infants} A small catheter placed in an infant's artery to continuously measure blood pressure and obtain blood samples. It requires careful monitoring for bleeding, infection, or reduced blood flow to the limb.

\medterm{Peripheral Artery Bypass - Leg} Surgery that creates a new route around a blocked leg artery to improve blood flow. A vein or artificial graft may be used as the bypass.


\medterm{Peripheral Artery Disease} Narrowing of arteries supplying the limbs, usually caused by atherosclerosis. It may
cause exertional leg pain relieved by rest, reduced pulses, nonhealing wounds, or chronic limb-threatening ischemia, a severe reduction in blood flow associated with persistent pain or tissue loss. Also called peripheral arterial disease.

\synonyms
Peripheral Arterial Disease; PAD

\medterm{Peripheral Artery Disease - Legs} Peripheral artery disease of the legs is atherosclerotic reduction in arterial blood flow causing exertional claudication, diminished pulses, nonhealing wounds, or critical limb ischemia.


\medterm{Peripheral Blood} Peripheral blood is blood obtained from a vein or capillary rather than directly from an organ or marrow, used for counts, smears, chemistry, and other testing.

\medterm{Peripheral Blood Smear} Microscopic examination of stained blood spread on a slide. It shows the size, shape, and appearance of red cells, white cells, and platelets and can help investigate anemia, infection, and blood disorders.

\medterm{Peripheral Blood Stem Cells} Immature blood-forming cells collected from the bloodstream after medicines encourage them to leave the bone marrow. They can be returned to the same person or given to another person to rebuild blood production after intensive treatment.

\medterm{Peripheral Edema} Swelling due to accumulation of fluid in the interstitial space of the limbs, most
commonly the lower extremities.
\medterm{Peripheral Intravenous Line - Infants} A small tube placed in an infant's vein, usually in the hand, foot, or scalp, to give medicines or fluids. The site should be checked for swelling, leakage, or infection.

\medterm{Peripheral Ischemia} Inadequate arterial blood flow to an arm, leg, digit, or other peripheral tissue, causing pain, pallor, coolness, weak pulses, numbness, ulcers, or tissue loss.

\medterm{Peripheral Nerve} A nerve outside the brain and spinal cord that carries motor, sensory, and autonomic signals between the central nervous system and the rest of the body.

\medterm{Peripheral Nerve Blocks} Injections of local anesthetic near a nerve or group of nerves to numb a specific area and reduce pain during or after a procedure. Ultrasound guidance may improve accuracy and safety.

\medterm{Peripheral Nerve Injuries} Damage to nerves outside the brain and spinal cord caused by trauma, compression, stretching, disease, or toxins.
It may cause pain, numbness, weakness, or loss of function; recovery depends on the
severity and location of the injury.

\medterm{Peripheral Nerve Tumors} Growths arising from or around nerves outside the brain and spinal cord; they may be benign or cancerous.

\medterm{Peripheral Nervous System} All nerves outside the brain and spinal cord. It carries messages between the central nervous system and the skin, muscles, organs, and glands, controlling movement, sensation, and automatic functions such as heartbeat and digestion.

\medterm{Peripheral Neuritis (Peripheral Neuropathy)} Alternate index label for Peripheral Neuropathy.

\medterm{Peripheral Neurofibromatosis} Alternate index label; see \textit{Neurofibromatosis type 1}.

\medterm{Peripheral Neuropathy} Damage or dysfunction of nerves outside the brain and spinal cord. It may cause burning pain, tingling, numbness, weakness, or problems with sweating and other automatic body functions, often beginning in the feet or hands.

\synonyms
Neuropathy, Peripheral
Peripheral Neuritis (Peripheral Neuropathy)

\medterm{Peripheral Neuropathy Diabetic (Diabetic Neuropathy)} Alternate index label for Diabetic Neuropathy.

\medterm{Peripheral Oedema} Alternate spelling of Peripheral Edema; swelling caused by excess interstitial fluid, usually in the limbs.

\medterm{Peripheral Vascular Disease} Disease affecting narrowed or blocked blood vessels outside the heart and brain. It commonly reduces blood flow to the legs, causing pain with walking, wounds that heal slowly, coldness, or tissue loss.

\synonyms
Peripheral arterial disease
PVD (Peripheral Vascular Disease)

\medterm{Peripheral Vision} The parts of the visual field seen away from the direct line of sight. Loss can result from eye disease, nerve damage, brain injury, or glaucoma and may increase falls or collisions.

\synonyms
Side vision

\medterm{Peripherally Inserted Central Catheter} A long, flexible catheter inserted into the basilic or cephalic vein and advanced so the
tip sits in the superior vena cava. Inserted at the bedside, often by a specialized PICC
nurse, with ultrasound and ECG or X-ray confirmation. Indications: long-term IV
antibiotics (>2 weeks), TPN, chemotherapy, and frequent blood draws.

\medterm{Peripherally Inserted Central Catheter - Dressing Change} Replacing the sterile covering over a peripherally inserted central catheter while checking the skin and catheter position. A clean technique reduces infection risk; redness, drainage, pain, fever, leakage, or a changed catheter length requires prompt review.

\medterm{Peripherally Inserted Central Catheter - Flushing} Passing a prescribed sterile solution through a peripherally inserted central catheter to keep it open and clear blood or medicine residue. The catheter should be flushed only with the specified solution and technique; resistance, pain, swelling, leakage, or fever needs assessment.

\medterm{Peripherally Inserted Central Catheter - Insertion} Placement of a long flexible tube through a vein in the arm until its tip lies near the heart. It allows repeated intravenous medicines, nutrition, or blood tests, but can cause infection, blockage, bleeding, or a blood clot.

\medterm{Periplasm} The periplasm is the space between the inner and outer membranes of Gram-negative bacteria, containing enzymes, transport proteins, and structural components.

\medterm{Periprosthetic Joint Infection} Infection involving an artificial joint and nearby tissues. It may cause persistent pain, swelling, drainage, fever, or loosening. Diagnosis uses examination, blood tests, joint-fluid testing, cultures, and imaging; treatment may require antibiotics and surgery.

\medterm{Perirenal Abscess} A perirenal abscess is a collection of pus around the kidney, usually from spread of a urinary or kidney infection, causing fever, flank pain, and possible sepsis.

\medterm{Perirenal Extramedullary Hematopoiesis} Production of blood cells in tissue around the kidney instead of mainly in the bone marrow. It is an unusual response to severe long-term anemia or bone-marrow disease and may appear as masses on imaging.

\medterm{Perirenal Hematoma} A collection of blood around the kidney, usually caused by trauma, a bleeding lesion, a procedure, or anticoagulation.

\medterm{Peristalsis} Peristalsis is coordinated, wave-like contraction of smooth muscle that propels food, fluid, or other contents through the gastrointestinal tract and some ducts.

\medterm{Peritoneal} Peritoneal means relating to the peritoneum, the thin membrane lining the abdominal cavity and covering many abdominal organs.

\medterm{Peritoneal Carcinomatosis} Cancer that has spread widely across the lining of the abdomen, often from the bowel, stomach, ovary, or another organ. It may cause abdominal swelling, pain, fluid buildup, weight loss, or bowel blockage; scans and fluid or tissue tests guide treatment.

\medterm{Peritoneal Cavity} The narrow space inside the abdominal lining around the abdominal and pelvic organs. A small amount of fluid normally lubricates it; blood, pus, excess fluid, or cancer cells can collect there during disease.

\medterm{Peritoneal Dialysis} A treatment for kidney failure that uses the lining of the abdomen to remove waste and extra water. Clean dialysis fluid enters through a permanent tube and is later drained; it can be done at home but may cause infection, fluid imbalance, or catheter problems.

\medterm{Peritoneal Dialysis Catheter} A soft tube placed through the abdominal wall into the abdominal cavity for dialysis. Dialysis fluid enters and leaves through it to remove waste and extra water; infection, blockage, leakage, or displacement can occur.

\medterm{Peritoneal Diseases} Disorders affecting the thin lining of the abdomen, including infection, inflammation, cancer, bleeding, scar tissue, or abnormal fluid buildup. Symptoms depend on the cause and may include abdominal pain, swelling, fever, or vomiting.
\medterm{Peritoneal Fluid} The small amount of slippery fluid normally present inside the abdominal cavity. Excess fluid, blood, pus, or cancer cells can collect there and may cause swelling, pain, infection, or other signs of abdominal disease.

\medterm{Peritoneal Fluid Analysis} Laboratory testing of fluid collected from the abdominal cavity. Cell counts, protein, albumin, cultures, and other tests help determine whether the fluid is caused by liver disease, infection, cancer, bleeding, or another problem.

\medterm{Peritoneal Fluid Culture} Culture of ascitic fluid to detect microorganisms, particularly in suspected spontaneous bacterial peritonitis or secondary peritoneal infection.

\medterm{Peritoneal Infection} Infection of the lining of the abdominal cavity, often called peritonitis. Causes include a ruptured
appendix or bowel, abdominal surgery, pelvic infection, or contamination during dialysis.
Severe abdominal pain, fever, vomiting, or a rigid abdomen requires urgent assessment.

\medterm{Peritoneal Lavage} Washing or sampling the abdominal cavity with sterile fluid. It was once used more often to detect internal bleeding or intestinal contents after injury; imaging and other tests now commonly provide this information.

\medterm{Peritoneal Necrosis} Death of peritoneal tissue caused by ischemia, severe infection, inflammation, trauma, or tumor, potentially leading to peritonitis, adhesions, sepsis, or organ failure.

\medterm{Peritoneal Sac} The space formed by the abdominal lining around the abdominal and pelvic organs. It normally contains only a small amount of lubricating fluid, but blood, pus, excess fluid, or cancer can collect there.

\medterm{Peritoneum} A thin, slippery lining that covers the inside of the abdomen and many abdominal organs. It allows organs to move smoothly; inflammation, infection, bleeding, fluid, or cancer in this lining can cause severe abdominal symptoms.

\medterm{Peritoneum Neoplasm} An abnormal growth in the lining of the abdominal cavity. It may begin there or spread from cancer in another organ and can cause abdominal pain, swelling, fluid buildup, or bowel blockage.

\medterm{Peritonitis} Inflammation or infection of the abdominal lining, often caused by a perforated organ, surgery, appendicitis, abdominal injury, or infected abdominal fluid. It can cause severe pain, a swollen or rigid abdomen, fever, vomiting, sepsis, or shock and needs urgent care.

\medterm{Peritonitis - Secondary} Infection of the abdominal lining caused by leakage from a perforated or injured digestive or other abdominal organ. It can rapidly cause severe pain, a rigid abdomen, sepsis, or shock and requires urgent treatment.

\medterm{Peritonitis - Spontaneous Bacterial} Infection of fluid in the abdominal cavity without an obvious hole or other surgical source. It occurs most often in people with advanced liver disease and can cause pain, fever, confusion, or sepsis.

\medterm{Peritonitis In PD} Infection of the abdominal lining in a person receiving peritoneal dialysis. Cloudy drainage, abdominal pain, fever, nausea, or feeling unwell needs prompt testing of the drainage and antibiotics; severe or repeated infection may require catheter removal.

\medterm{Peritonsillar Abscess} A collection of pus beside a tonsil, usually after a throat infection. It can cause severe one-sided throat pain, fever, muffled speech, drooling, difficulty opening the mouth, or trouble breathing and requires urgent treatment and drainage.

\medterm{Periungual Fibroma} A noncancerous, firm growth arising beside or beneath a fingernail or toenail. It may distort the nail and can occur alone or as one of several skin findings in tuberous sclerosis complex.

\medterm{Periurethral} Located around the urethra, the tube that carries urine out of the body. The term may describe glands, infection, an abscess, a urethral diverticulum, or tissue and treatments used for urinary leakage.
\medterm{Periventricular Leukomalacia} Injury to the brain’s developing nerve fibers near its fluid-filled spaces, most often in premature infants. It may later affect movement, muscle tone, vision, learning, or development and is associated with cerebral palsy.

\medterm{Perivitelline Space} The perivitelline space is the narrow region between an egg’s plasma membrane and surrounding vitelline or zona material, important in fertilization and early development.

\medterm{Perleche} Inflammation, soreness, and cracking at one or both corners of the mouth. Saliva, ill-fitting dentures, yeast or bacterial infection, skin irritation, and iron or vitamin shortages can contribute; treatment depends on the cause and keeping the area dry.
\medterm{Permanent Teeth} The adult set of teeth that replaces the deciduous teeth.
\medterm{Permease} A permease is a membrane transport protein that facilitates movement of a specific molecule or ion across a cell membrane.

\medterm{Permethrin} A skin-applied medicine that kills scabies mites and head lice. It is usually applied as directed and washed off later; itching may continue briefly after treatment even when parasites are gone.
\medterm{Permissive Hypercapnia} A ventilator approach that allows carbon dioxide in the blood to remain above normal so that gentler pressures and smaller breaths can protect injured lungs. It is used selectively and requires close monitoring because excessive acidity or carbon dioxide can be harmful.


\medterm{Pernicious Anaemia} A form of vitamin B12 deficiency caused by immune damage to stomach cells needed for B12 absorption. It can cause anemia, tiredness, sore tongue, numbness, balance problems, and nerve damage.
\medterm{Pernicious Anemia} Anemia caused by vitamin B12 deficiency when the immune system damages stomach cells that make a substance needed for B12 absorption. It may cause tiredness, sore tongue, numbness, balance problems, or permanent nerve damage if untreated.

\synonyms
Addisons Anemia (Pernicious Anemia)

\medterm{Pernio} Alternate index label; see \textit{Chilblains}.

\medterm{Peroneal} Relating to the fibula or the muscles and nerve on the outer side of the lower leg.
\medterm{Peroneal Muscular Atrophy} Alternate index label; see \textit{Charcot-Marie-Tooth disease}.

\medterm{Peroneal Nerve} The common fibular nerve, a branch of the sciatic nerve that winds around the fibular neck and supplies sensation and muscles responsible for ankle dorsiflexion and eversion.

\medterm{Peroxidase} An enzyme that uses hydrogen peroxide or another peroxide to oxidize a substrate; thyroid peroxidase is essential for thyroid-hormone synthesis.

\medterm{Peroxisomal Disorder} An inherited condition in which small cell structures that process certain fats do not form or work properly. It may affect the brain, nerves, liver, eyes, hearing, bones, or development.

\medterm{Peroxisome} A small structure inside a cell that breaks down certain fats and harmful substances. It also helps control hydrogen peroxide and makes substances needed for cell membranes and nerve function.

\medterm{Peroxisome Proliferation} An increase in the number or size of peroxisomes, small cell structures that process certain fats and harmful chemicals. It can be triggered by some medicines, chemicals, hormones, or metabolic conditions and may alter cell function.

\medterm{Perphenazine} A first-generation (typical) phenothiazine antipsychotic used for schizophrenia and, in some settings, severe nausea or vomiting. Adverse effects can include sedation, orthostatic hypotension, extrapyramidal symptoms, tardive dyskinesia, and neuroleptic malignant syndrome.
\medterm{Perseveration} Involuntary repetition of a word, idea, movement, or response after the appropriate stimulus has changed, associated with frontal or basal-ganglia disease, aphasia, delirium, or psychiatric illness.

\medterm{Persistent Asthma} Asthma symptoms or airway narrowing that occur regularly rather than only occasionally and usually require ongoing controller treatment. Severity is judged by symptoms, night waking, reliever use, lung function, and activity limits; treatment is adjusted over time.

\medterm{Persistent Cloaca} Persistent cloaca is a congenital malformation in which the rectum, vagina, and urinary tract join into one channel, requiring specialist evaluation and reconstructive care.

\medterm{Persistent Depressive Disorder} A long-lasting depressive disorder, formerly called dysthymia, with depressed mood or irritability for at least two years in adults or one year in children and adolescents. Other symptoms may include low energy, poor self-esteem, sleep or appetite changes, poor concentration, or hopelessness.

\medterm{Persistent Pulmonary Hypertension Of The Newborn} Failure of the normal fall in pulmonary vascular resistance after birth, causing right-to-left shunting and hypoxaemia. Management requires urgent neonatal respiratory and circulatory support.

\medterm{Persistent Pupillary Membranes} Strands left from a temporary membrane that covered the developing eye before birth. They may cross the pupil and are usually harmless, but dense strands can interfere with vision and occasionally need treatment.

\medterm{Person With Disability} A person with disability has a long-term physical, sensory, intellectual, or mental-health impairment that may interact with barriers to limit participation; care should be accessible and person-centered.

\medterm{Personal Emergency Response System} A wearable or nearby device that lets a person summon help during an emergency.
\synonyms
Medical alert system

\medterm{Personal Protective Equipment} Clothing or equipment that protects a person from infection, chemicals, blood, or other hazards. It may include gloves, gowns, masks, respirators, goggles, or face shields; the correct items depend on the hazard and must be put on, removed, and disposed of safely.

\medterm{Personality} Personality is the relatively enduring pattern of thoughts, emotions, motivations, and behaviors through which a person experiences and responds to the world.

\medterm{Personality Change} Personality change is a new or marked alteration in behavior, emotional responses, judgment, or interpersonal style, which may reflect neurologic disease, psychiatric illness, substances, or medical conditions.

\medterm{Personality Disorder} Alternate index label for Personality Disorders.

\medterm{Personality Disorder Antisocial (Antisocial Personality Disorder)} Alternate index label for Antisocial Personality Disorder.

\medterm{Personality Disorder, Antisocial} Alternate index label; see \textit{Antisocial personality disorder}.

\medterm{Personality Disorder, Borderline} Alternate index label; see \textit{Borderline personality disorder}.

\medterm{Personality Disorder, Narcissistic} Alternate index label; see \textit{Narcissistic personality disorder}.

\medterm{Personality Disorder, Schizoid} Alternate index label; see \textit{Schizoid personality disorder}.

\medterm{Personality Disorder, Schizotypal} Alternate index label; see \textit{Schizotypal personality disorder}.

\medterm{Personality Disorders} A group of mental disorders involving enduring, inflexible patterns of thought, emotion, behavior, and relationships that differ markedly from the person’s cultural context and cause distress or impaired functioning. Traditional groupings include Cluster A (paranoid, schizoid, and schizotypal), Cluster B (antisocial, borderline, histrionic, and narcissistic), and Cluster C (avoidant, dependent, and obsessive-compulsive personality disorders).

\medterm{Personality Test} A psychological assessment using structured questions or stimuli to describe enduring traits, patterns, or possible personality pathology; results require validated instruments and clinical interpretation.

\medterm{Perspiration} Sweat released by skin glands. Eccrine glands mainly help cool the body, while apocrine glands contribute to odor after their secretions are changed by skin bacteria; unusually heavy sweating is called hyperhidrosis.
\medterm{Persuasive Communication} Communication intended to influence beliefs, decisions, or behavior through information, reasoning, emotion, or framing; healthcare use should preserve informed choice and avoid coercion.

\medterm{Pertussis} Alternate index label for Whooping Cough.
\medterm{Pertussis Toxin} A toxin produced by Bordetella pertussis that disrupts cell signaling and contributes to the cough and systemic effects of whooping cough.

\medterm{Pertuzumab} A monoclonal antibody targeting the HER2 receptor and blocking its dimerization, used with other medicines for selected HER2-positive cancers.

\medterm{Pervasive Developmental Disorders} An obsolete group term for childhood conditions involving social communication difficulties and restricted or repetitive behavior. Modern classification uses autism spectrum disorder and other specific developmental diagnoses.
\medterm{Pes} Pes means foot; it is also used in terms describing foot shape or deformity, such as pes planus or pes cavus.

\medterm{Pes Anserinus} The combined tendon insertion of the sartorius, gracilis, and semitendinosus muscles on the anteromedial proximal tibia; inflammation there is pes-anserine bursitis.

\medterm{Pes Planus} Alternate index label; see \textit{Flatfeet}.

\medterm{Pes Planus (Flatfoot (Pes Planus))} Alternate index label for Flatfoot (Pes Planus).

\medterm{Pessaries} Devices placed inside the vagina to support organs that have descended, such as the bladder, uterus, or rectum. Ring, cube, and other designs suit different shapes and needs; regular removal and follow-up help prevent sores, discharge, or infection.

\medterm{Pessary} A device placed inside the vagina to support pelvic organs, treat prolapse, or sometimes prevent leakage of urine.
\medterm{Pessimism} A tendency to expect negative outcomes or interpret situations unfavorably.


\medterm{Pesticide} A substance used to control insects, weeds, fungi, or other pests; exposure can cause poisoning.
\medterm{Pesticides} Chemicals used to kill insects, weeds, fungi, rodents, or other pests. Exposure can irritate the skin, eyes, lungs, or digestive tract; some pesticides can cause sweating, vomiting, weakness, seizures, or breathing failure.

\medterm{Pesticide Poisoning} Harm caused by excessive exposure to an insecticide, herbicide, fungicide, or other pesticide. Effects depend on the chemical and may include skin or eye injury, vomiting, sweating, weakness, seizures, breathing failure, or abnormal heart rhythms.

\medterm{Pesticides On Fruits And Vegetables} Small pesticide residues on properly washed produce usually do not cause acute poisoning. Swallowing concentrated pesticide or developing vomiting, sweating, breathing difficulty, weakness, or confusion requires immediate poison-control or emergency advice.

\medterm{Pestilence} Pestilence is an old term for a severe, widespread, often fatal infectious disease outbreak; modern usage usually specifies the pathogen and epidemiologic event.

\medterm{Pet Allergy} An allergic reaction to proteins in an animal's skin flakes, saliva, or urine, commonly causing rhinitis, conjunctivitis, or asthma.

\synonyms
Allergy, Pet

\medterm{PET Scan} An imaging study using a positron-emitting tracer to measure metabolism, receptor activity, or tissue function.
\medterm{PET Scan (Positron Emission Tomography)} Nuclear medicine imaging using radioactive tracers (most commonly FDG, a glucose analog)
to detect metabolically active tissues.

\synonyms
Positron Emission Tomography

\medterm{PET Scan For Breast Cancer} Positron-emission tomography, usually combined with computed tomography, used selectively to assess metabolically active breast cancer, staging, recurrence, or treatment response.

\medterm{PET Scanner} A machine that detects radiation from a tracer inside the body and creates images showing how actively tissues use substances such as glucose. It is used to assess selected cancers and other disorders.
\medterm{PET Tracers} Radioactive substances used during a PET scan to highlight particular body processes. Different tracers show glucose use, bone activity, or specific cell receptors; the appropriate tracer depends on the clinical question.

\medterm{PET/CT} A scan that combines PET, which shows active body processes, with CT, which shows detailed anatomy. It helps locate disease and assess selected cancers, heart conditions, and brain disorders; it uses a radioactive tracer and X-rays.

\medterm{PET/MRI} A scan that combines information about tissue activity with detailed MRI images of soft tissues. It is useful for selected brain, cancer, and heart problems and does not use X-rays, but takes longer than many other imaging tests.

\medterm{Petechia} A petechia is a tiny nonblanching red, purple, or brown spot caused by bleeding from small vessels into the skin or mucosa; numerous or rapidly spreading petechiae need assessment.

\medterm{Petechiae} Tiny nonblanching red or purple spots caused by small areas of bleeding into the skin or mucosa.
\medterm{Petit Mal} An older term for absence seizures, brief episodes of impaired awareness with characteristic generalized EEG activity.
\medterm{Petit Mal Seizure} Alternate index label; see \textit{Absence seizure}.


\medterm{Petrissage} A massage technique involving kneading, squeezing, or lifting soft tissues. It should be adapted or avoided when there is acute injury, infection, clot risk, or fragile skin.
\medterm{Petrolatum} A semisolid purified mixture of hydrocarbons used on skin as a protective ointment and moisture barrier. It reduces water loss, shields minor irritation, and supports healing of dry, cracked, or chafed skin.
\medterm{Petroleum Jelly Overdose} Swallowing a small amount of petroleum jelly usually causes little harm, but larger amounts can cause diarrhea or vomiting. If it enters the lungs, it may cause inflammation and breathing problems; contact poison control after a significant exposure.

\medterm{Petrositis} Inflammation or infection of the petrous portion of the temporal bone, usually arising from complicated otitis media and potentially affecting nearby cranial nerves.

\medterm{Petrous} Describing dense, hard bone, especially the part of the temporal bone that houses inner-ear structures.
\medterm{Pets And The Immunocompromised Person} People with weakened immunity can often keep pets but should reduce exposure to animal waste, scratches, bites, and unsafe food and follow individualized medical advice.

\medterm{Peutz-Jeghers Syndrome} An inherited condition causing dark spots on the lips, mouth, fingers, or skin and multiple growths in the digestive tract. These growths can bleed or block the bowel and increase the risk of several cancers.

\synonyms
PJS

\medterm{Peutz-Jeghers's Syndrome} An inherited condition causing dark patches around the mouth and many growths in the stomach or intestines. The growths may bleed, cause blockage, or increase the risk of digestive and other cancers.

\medterm{PEWS} Alternate name; see \textit{Pediatric Early Warning Score}.

\medterm{Pexophagy} A cell-cleaning process that removes damaged or unneeded peroxisomes, the small structures that process certain fats and harmful chemicals. The cell breaks them down and reuses some of their materials.

\medterm{Peyer Patches} Clusters of immune tissue in the wall of the lower small intestine. They sample material from the gut and help the immune system recognize and respond to germs while tolerating normal food and bacteria.

\medterm{Peyer's Patches} Organized lymphoid tissue in the mucosa and submucosa of the ileum. These aggregates sample intestinal contents and contribute to immune surveillance of the gastrointestinal tract.

\synonyms
ileal lymphoid nodules

\medterm{Peyronie Disease} Scar tissue forming inside the penis, causing a new bend or indentation that may make erections painful or intercourse difficult. Symptoms may stabilize or worsen; treatment ranges from observation and injections to surgery in selected severe cases.

\medterm{Peyronie’S Disease} A condition in which scar tissue inside the penis causes curved, painful erections and sometimes erectile difficulty.

\synonyms
Peyronie disease

\medterm{Pfeiffer Syndrome} An autosomal dominant craniosynostosis syndrome caused by FGFR1 or FGFR2 mutations.
Features: craniosynostosis (bilateral coronal), midface hypoplasia, proptosis,
hypertelorism, broad thumbs and great toes (medial deviation), and variable syndactyly.

\synonyms
Syndrome Pfeiffer (Pfeiffer Syndrome)

\medterm{PH Meter} An instrument that measures acidity or alkalinity as pH in a solution or specimen; clinical applications include gastric, urinary, laboratory, and environmental measurements.

\medterm{PH Monitoring} Measurement of acidity, usually in the esophagus, over time with a small sensor. It can help detect and measure acid reflux when symptoms or endoscopy do not give a clear answer.


\medterm{Phacoemulsification} A cataract-surgery technique that uses ultrasound to break up and remove the cloudy lens through a small incision. An
intraocular lens is usually implanted to restore focusing ability.

\medterm{Phagocyte} An immune cell that surrounds, swallows, and breaks down germs, dead cells, or foreign particles. Phagocytes help prevent infection, remove damaged material, and alert other immune cells to danger.
\medterm{Phagocytes} Immune cells that surround, swallow, and break down germs, dead cells, and foreign particles. They help prevent infection, clear damaged material, and signal other immune cells during inflammation and tissue repair.

\medterm{Phagocytic Vesicle} A small membrane-bound compartment formed inside a cell after it surrounds and swallows a germ, dead cell, or particle. It later joins a digestive compartment so the material can be broken down.

\medterm{Phagocytosis} The process by which a cell surrounds, takes in, and breaks down a particle or microorganism. Immune cells use it to remove germs, dead cells, and harmful debris from tissues.
\medterm{Phagolysosome} A compartment formed when an immune cell’s swallowed-particle compartment joins a digestive compartment. Enzymes and other chemicals inside it destroy germs and break down cellular debris.

\medterm{Phagophore} A temporary membrane that grows around damaged or unnecessary material inside a cell during recycling. It closes around the material so the cell can break it down and reuse some of its parts.

\medterm{Phalanges} The bones of the fingers and toes. The thumb and big toe each have two, while the other digits usually have three. Fractures, arthritis, infection, or nerve problems in these bones can cause pain, swelling, stiffness, or difficulty using the hand or foot.
\medterm{Phalanx} One of the bones of a finger or toe. Each phalanx has a base, shaft, and end; the end bone supports the fingertip or toenail. The thumb and big toe have two phalanges, while other digits usually have three.
\medterm{Phalanx Unspecified} A phalanx is one of the bones of a finger or toe; “unspecified” indicates that the particular digit or proximal, middle, or distal segment is not identified.

\medterm{Phalen Test} A test in which the wrists are bent to see whether numbness or tingling develops. Reproduced symptoms may support carpal tunnel syndrome, but the result must be interpreted with examination and other tests.

\medterm{Phantom} A test object made to resemble human tissue or organs so medical imaging equipment can be checked and calibrated. Different phantoms test image quality, measurements, or specific body regions.

\medterm{Phantom Limb} The feeling that an amputated body part is still present. Sensations may include position, movement, itching, pressure, or pain and often change over time after amputation.

\synonyms
Phantom-limb sensation

\medterm{Phantom Limb Pain} Painful sensations perceived in a limb that has been amputated, caused by persistent or
reorganized activity in peripheral nerves and central sensory pathways. It may feel
burning, shooting, cramping, or electric.

\medterm{Phantom Pain} Pain perceived in a body part that has been amputated, caused by changes in peripheral nerves and central pain-processing pathways.

\medterm{Phantosmia} Perception of a smell when no corresponding odor is present. It may occur with sinonasal
disease, migraine, seizures involving olfactory networks, head injury, neurodegenerative
disease, infection, or other neurological conditions.

\medterm{Pharmaceutical Adjuvant} A pharmaceutical adjuvant is an inactive or auxiliary ingredient added to improve a medicine’s stability, absorption, preservation, delivery, or therapeutic effect.

\medterm{Pharmaceutical Excipient} A pharmaceutical excipient is a non-active ingredient in a dosage form, such as a filler, binder, coating, solvent, preservative, or disintegrant.

\medterm{Pharmaceutical Solutions} Liquid preparations in which one or more active pharmaceutical ingredients are dissolved in a suitable vehicle for administration, with concentration, sterility, stability, and route determining safe use.

\medterm{Pharmacist} A medicine expert who prepares and dispenses drugs, checks safety, and explains correct use and possible side effects.
\medterm{Pharmacodynamics} The study of the biochemical and physiological effects of drugs on the body and their mechanisms of action, describing what the drug does to the body.
\medterm{Pharmacogenetic Testing} Testing for genetic variants that influence a person’s response, metabolism, efficacy, or toxicity of particular medicines and may support individualized prescribing.

\medterm{Pharmacogenetics} Pharmacogenetics studies how inherited genetic variation alters a person's response to medicines, including efficacy, metabolism, and adverse-effect risk, to support individualized prescribing.

\medterm{Pharmacogenomic Testing} Testing that examines genetic variants affecting drug metabolism, transport, targets, or adverse reactions to help select or dose certain medicines.

\medterm{Pharmacogenomics} The study of how genetic variations influence individual responses to drugs, including
drug efficacy, toxicity, and metabolism. Genetic polymorphisms in drug-metabolizing
enzymes, transporters, and receptors can lead to significant inter-individual
variability in drug response.

\medterm{Pharmacokinetics} The study of what the body does to a medicine, including its absorption, distribution, metabolism, and elimination. It helps explain drug concentrations over time and supports dose selection.

\medterm{Peak Plasma Concentration} The highest measured concentration of a medicine in the blood after a dose. It is often called Cmax and is influenced by the dose, absorption speed, formulation, and route of administration.
\medterm{Pharmacologic Substance} A pharmacologic substance is a chemical that produces a biological effect by interacting with cells, receptors, enzymes, transporters, or other targets.

\medterm{Pharmacological Stress Test} A heart stress test using medicine to imitate exercise when a person cannot exercise adequately. Imaging or heart recordings show blood flow and function, while possible effects include chest discomfort, breathlessness, or abnormal rhythm.

\medterm{Pharmacologist} A pharmacologist studies how medicines and chemicals act in living systems, including mechanisms, metabolism, toxicity, and therapeutic effects.

\medterm{Pharmacology} The study of medicines and other drugs, including how they work, how the body handles them, their uses, side effects, and interactions. This knowledge helps guide safe choice, dosing, testing, and monitoring of treatment.
\medterm{Pharmacopoeia} A pharmacopoeia is an official reference specifying standards for identity, purity, strength, testing, and preparation of medicines.

\medterm{Pharmacotherapy} Pharmacotherapy is treatment with medicines chosen to prevent, relieve, or control a disease or symptom.
\medterm{Pharmacovigilance} The science and activities relating to the detection, assessment, understanding, and
prevention of adverse drug reactions and other drug-related problems.

\medterm{Pharmacy} The healthcare field concerned with preparing, checking, supplying, and explaining medicines. Pharmacists help choose safe doses, identify interactions and allergies, support treatment decisions, and teach people how to use medicines correctly.
\medterm{Pharmacy Facility} A pharmacy facility is an authorized location where medicines may be stored, prepared, dispensed, monitored, and supplied under applicable professional and safety standards.

\medterm{Pharyngeal} Relating to the pharynx, or throat. The pharynx connects the nose and mouth with the voice box, food pipe, and breathing passages. Infection, swelling, narrowing, or a tumor there may cause sore throat, swallowing difficulty, noisy breathing, or voice changes.
\medterm{Pharyngeal Carcinoma} Pharyngeal carcinoma is cancer arising in the throat and may cause persistent sore throat, swallowing difficulty, hoarseness, or a neck mass.

\medterm{Pharyngeal Disease} A disorder affecting the throat, including infection, inflammation, blockage, injury, muscle or nerve problems, or a tumor. Symptoms may include sore throat, swallowing difficulty, choking, voice changes, or breathing problems.

\medterm{Pharyngeal Hemorrhage} Bleeding from the pharynx caused by trauma, surgery, infection, vascular lesions, tumor, or severe mucosal disease and potentially threatening the airway.

\medterm{Pharyngeal Mucositis} Inflammation and ulceration of the throat lining, commonly caused by radiation, chemotherapy, infection, trauma, or medications and producing pain, difficulty swallowing, and reduced intake.

\medterm{Pharyngeal Stenosis} Narrowing of the pharyngeal passage caused by scarring, congenital abnormality, tumor, inflammation, or surgery, which may impair swallowing or breathing.

\medterm{Pharyngeal Tonsil} Adenoid tissue in the roof of the throat behind the nose that helps immune defense, especially in children.

\medterm{Pharyngectomy} Surgery to remove part or all of the throat, usually to treat cancer. Depending on the extent, reconstruction or a temporary feeding tube may be needed, and treatment can affect swallowing, speech, breathing, or appearance.
\medterm{Pharyngitis} Inflammation of the throat behind the mouth and nose. It commonly causes sore throat, painful swallowing, fever, cough, or swollen neck glands and is most often caused by a virus.

\synonyms
Sore Throat (Pharyngitis)

\medterm{Pharyngitis (Sore Throat (Pharyngitis))} Alternate index label for Sore Throat (Pharyngitis).

\medterm{Pharyngitis - Sore Throat} Pharyngitis is inflammation of the throat causing pain with swallowing, commonly from viral infection and sometimes from group A streptococcal infection or other causes.

\medterm{Pharyngitis - Viral} Viral pharyngitis is inflammation of the throat caused by a virus, commonly producing sore throat, cough, runny nose, or fever and usually resolving with supportive care.

\medterm{Pharyngolaryngeal Pain} Pain involving the pharynx and larynx, commonly associated with infection, inflammation, irritation, trauma, reflux, or malignancy.

\medterm{Pharyngomaxillary Space Abscess} A pharyngomaxillary space abscess is a deep infection containing pus near the pharynx and maxilla that can cause severe pain, swelling, difficulty swallowing, and airway risk.

\medterm{Pharynx} The muscular passage behind the nose and mouth that carries air toward the voice box and food toward the esophagus. It helps with breathing, swallowing, and speech.
\medterm{Pharynx Necrosis} Death of tissue in the throat caused by severe infection, ischemia, radiation, caustic injury, or tumor, with risk of deep infection, bleeding, airway compromise, and sepsis.

\medterm{Pharynx Neoplasm} A pharynx neoplasm is an abnormal growth in the throat that may be benign or malignant.

\medterm{Phase I Metabolism} An early step in the body’s breakdown of many medicines. Enzymes, mainly in the liver, chemically change a medicine by adding or exposing a reactive part; this may inactivate it, activate it, or prepare it for further breakdown.

\medterm{Phase II Metabolism} A stage of drug breakdown in which the body attaches another chemical group to a medicine or its earlier breakdown product. This usually makes the substance easier to remove in urine or bile.

\medterm{Phase-Contrast Microscopy} A microscope technique that makes transparent cells and other specimens easier to see by converting small differences in light passing through them into visible contrast. It is useful for examining living, unstained cells.

\medterm{Phe} Phe is the standard abbreviation for phenylalanine, an essential amino acid incorporated into proteins and metabolized through a pathway disrupted in phenylketonuria.

\medterm{Phelan-McDermid Syndrome} A genetic disorder that commonly causes low muscle tone, delayed development, intellectual disability, very limited speech, and autism-related behaviors. Some people have reduced sensitivity to pain, seizures, feeding problems, or kidney and heart abnormalities. It usually results from loss of genetic material near the end of chromosome 22 or a change in the \textit{SHANK3} gene.

\synonyms
22q13.3 Deletion Syndrome

\medterm{Phenacetin} Phenacetin is a former analgesic and fever-reducing drug withdrawn or restricted because long-term use can cause kidney damage and increase cancer risk.

\medterm{Phenazopyridine} Phenazopyridine is a urinary analgesic that temporarily relieves dysuria without treating its cause; it can turn urine orange and prolonged or excessive use may cause haemolysis or methemoglobinaemia.

\medterm{Phencyclidine} A dissociative drug that can cause altered perception, agitation, hallucinations, impaired coordination, or dangerous behavior.
\medterm{Phencyclidine Overdose} A dangerous amount of the dissociative drug phencyclidine can cause agitation, hallucinations, high body temperature, high blood pressure, seizures, loss of coordination, coma, or injury. Severe or unusual behavior requires emergency care.

\medterm{Phenelzine} An antidepressant that blocks monoamine oxidase, an enzyme breaking down several brain chemicals. It has important interactions with medicines and tyramine-rich foods and can cause dangerous high blood pressure if precautions are ignored.
\medterm{Phenethylamines} Phenethylamines are a chemical family including neurotransmitters, medicines, and psychoactive substances that can affect adrenergic, dopaminergic, or serotonergic signaling.

\medterm{Phenformin} Phenformin is a former biguanide diabetes medicine withdrawn because it can cause life-threatening lactic acidosis.
\medterm{Phenindione} An older oral anticoagulant that interferes with vitamin-K-dependent clotting factors.
\medterm{Pheniramine Overdose} Taking too much pheniramine, an antihistamine, can cause extreme sleepiness or agitation, dry mouth, large pupils, fast heartbeat, seizures, abnormal rhythms, or coma. Immediate poison-control or emergency advice is needed.

\medterm{Phenobarbital} Phenobarbital is a long-acting barbiturate antiseizure medicine that enhances GABA-A receptor inhibition; sedation, respiratory depression, cognitive effects, dependence, and enzyme-inducing drug interactions are important risks.

\medterm{Phenobarbital Overdose} Taking too much phenobarbital can cause severe drowsiness, poor coordination, low blood pressure, slowed breathing, coma, or death. Emergency treatment is required, especially if alcohol or another sedative was also taken.

\medterm{Phenobarbitone} The British name for phenobarbital, a barbiturate anticonvulsant and sedative.

\medterm{Phenol} A corrosive aromatic compound used in chemical and laboratory applications and in some medical preparations.

\medterm{Phenols} A group of aromatic chemicals that can burn skin and mucous membranes and cause systemic poisoning if absorbed or swallowed.

\medterm{Phenomenon} An observable event, sign, or pattern in the body. Examples include fingers changing color in Raynaud phenomenon, early-morning blood-sugar rises in the dawn phenomenon, and upward eye movement when the eyelids close in Bell phenomenon.
\medterm{Phenothiazine} Phenothiazine is a chemical class that includes several antipsychotic and antiemetic medicines, with possible sedation, movement, and cardiac adverse effects.

\medterm{Phenothiazine Overdose} Taking too much of a phenothiazine medicine can cause sleepiness, confusion, low blood pressure, seizures, abnormal heart rhythms, or severe muscle stiffness and fever. It requires immediate medical assessment.

\medterm{Phenothiazines} A class of medicines that includes antipsychotics and antiemetics. They can cause sedation, low blood pressure, movement disorders, heart-rhythm changes, and other adverse effects.
\medterm{Phenotype} The observable features of a person, such as appearance, growth, body functions, and behavior. These features result from the interaction between inherited genetic information and the environment.

\medterm{Phenoxybenzamine} A long-acting alpha-blocker used mainly before surgery for a catecholamine-secreting tumor such as pheochromocytoma. It can cause dizziness, low blood pressure, rapid heartbeat, and nasal congestion.
\medterm{Phenoxymethylpenicillin} Penicillin V, an oral penicillin used for selected susceptible infections.
\medterm{Phenprocoumon} Phenprocoumon is a long-acting vitamin-K antagonist anticoagulant that reduces clotting-factor production and requires coagulation monitoring.
\medterm{Phentermine} A stimulant medicine that reduces appetite and may be used short term with diet and activity changes for weight management. It can raise heart rate and blood pressure and may cause insomnia, anxiety, dependence, or dangerous interactions with some medicines.
\medterm{Phenthoate} Phenthoate is an organophosphate insecticide that can inhibit acetylcholinesterase and cause cholinergic poisoning after exposure.
\medterm{Phentolamine} Phentolamine is a short-acting alpha-adrenergic blocker used for catecholamine-related hypertension and to limit tissue injury after certain vasoconstrictor leaks.

\medterm{Phenylacetate} A salt or ester of phenylacetic acid used in some metabolic and experimental medical applications; it can bind nitrogen-containing metabolites.

\medterm{Phenylalanine} An essential amino acid incorporated into proteins and converted to tyrosine; impaired metabolism causes phenylketonuria and requires dietary or medical management.

\medterm{Phenylbutazone} An older anti-inflammatory pain medicine whose use is limited by serious toxicity. It can cause stomach bleeding, kidney or liver injury, fluid retention, blood-cell problems, allergic reactions, and other dangerous effects.
\medterm{Phenylbutyrate} A nitrogen-scavenging medicine converted to phenylacetate, which combines with glutamine and facilitates nitrogen excretion in selected urea-cycle disorders.

\medterm{Phenylephrine} A medicine that narrows blood vessels. It may reduce a blocked nose, raise dangerously low blood pressure in hospital care, or widen the pupil in an eye examination; it can worsen high blood pressure, heart problems, or urinary difficulty in some people.
\medterm{Phenylketonuria} An inherited metabolic disorder, usually caused by disease-causing variants in the \textit{PAH} gene that reduce phenylalanine hydroxylase activity. Phenylalanine accumulates in the blood and tissues, while tyrosine availability may decrease. Untreated or poorly controlled disease can affect brain development and may be associated with intellectual disability, seizures, eczema, behavioral or psychiatric symptoms, and a musty odor. Severity ranges from classic to milder forms, and the disorder is commonly identified through newborn screening.

\synonyms
Phenylketonuria (PKU)
PKU

\medterm{Phenylpropanolamine} A sympathomimetic decongestant formerly used in some products; it can raise blood pressure and has been restricted because of safety concerns.

\medterm{Phenytoin} A medicine used to prevent or control some seizures. It reduces repeated electrical activity in brain cells; levels in the blood may need monitoring because interactions and side effects can include dizziness, poor coordination, gum overgrowth, skin reactions, or serious toxicity.
\medterm{Phenytoin Overdose} Too much phenytoin can cause nausea, severe dizziness, double vision, involuntary eye movements, poor coordination, confusion, abnormal heart rhythms, or coma. Significant or intentional overdoses need urgent medical assessment.

\medterm{Phenytoin Sodium} The sodium salt of phenytoin, an anticonvulsant used for selected focal and generalized tonic-clonic seizures; toxicity can cause nystagmus, ataxia, arrhythmia, or severe skin reactions.

\medterm{Pheochromocytoma} A rare tumor, usually in an adrenal gland, that may release excess adrenaline-like hormones. It can cause episodes of high blood pressure, headache, sweating, a pounding heartbeat, shakiness, or paleness.

\synonyms
Tumor Adrenal Gland (Pheochromocytoma)

\medterm{Pheresis} A procedure that removes a selected part of blood, such as plasma, platelets, or white blood cells, and returns the remaining blood. It may treat disease or collect a needed blood component from a donor.

\medterm{Pheromone} A pheromone is a chemical signal released by an organism that alters behavior or physiology in another member of the same species.


\medterm{Phialophora} A genus of darkly pigmented fungi that can cause chromoblastomycosis or other subcutaneous and opportunistic infections.

\medterm{Phialophora Verrucosa} A dark-pigmented fungus that commonly causes chromoblastomycosis, a long-lasting infection of skin and tissue beneath it. Infection usually follows implantation through injured skin and produces slowly enlarging warty or raised lesions.

\medterm{Philadelphia Chromosome} A rearrangement of genetic material between chromosomes 9 and 22 that creates the BCR::ABL1 gene. It is found in most cases of chronic myeloid leukemia and some acute leukemias and helps guide treatment.

\medterm{Philodendron Poisoning} Chewing philodendron leaves releases irritating crystals that can cause immediate burning and swelling of the mouth, tongue, or throat. Breathing or swallowing difficulty requires emergency care; contact poison control after ingestion.

\medterm{Phimosis} Tightness of the foreskin that prevents it from being pulled back over the head of the penis. It can be normal in young children, but pain, repeated infection, scarring, or difficulty passing urine needs medical assessment.
\medterm{Phimosis And Paraphimosis (Penis Disorders)} Alternate index label for Penis Disorders.

\medterm{Phlebitis} Inflammation of a vein causing pain, warmth, redness, or tenderness along its course. It may follow irritation from an intravenous line or occur with a clot in a surface vein; sudden limb swelling or breathlessness needs urgent assessment.
\medterm{Phlebitis And Thrombophlebitis} Phlebitis is inflammation of a vein, while thrombophlebitis includes a clot within the inflamed vein. Superficial disease causes localized pain and redness; deep-vein thrombosis can embolize to the lungs and is suspected with limb swelling, pain, or breathlessness.

\medterm{Phlebography} An imaging test in which contrast dye is injected into a vein to show the veins on X-ray. It can detect a blood clot or an abnormal vein but is now often replaced by ultrasound.
\medterm{Phlebotomy} Taking blood from a vein, usually with a needle, for testing, donation, or treatment. The procedure may cause brief pain, bruising, bleeding, fainting, or rarely infection; trained staff use safe collection and labeling practices.

\synonyms
Venipuncture

\medterm{Phlebovirus} A genus of segmented RNA viruses, many transmitted by sandflies or mosquitoes; some cause fever, encephalitis, or hemorrhagic disease.

\medterm{Phlegm} Thick mucus brought up by coughing from the lower airways. Its amount and appearance can change with infection, asthma, smoking, chronic lung disease, or other conditions.
\medterm{Phlegmasia Cerulea Dolens} Phlegmasia cerulea dolens is a severe, extensive deep-vein thrombosis that obstructs venous outflow, causing painful swollen blue discoloration and risking limb ischemia or pulmonary embolism.

\medterm{Phobia} An excessive, persistent fear of a specific object or situation that leads to avoidance or distress.
\medterm{Phobia - Simple/Specific} A specific phobia is intense, persistent fear of a particular object or situation that is excessive, leads to avoidance, and interferes with life.

\medterm{Phobia, Social} Alternate index label; see \textit{Social anxiety disorder (social phobia)}.

\medterm{Phobias} Persistent, excessive fear of a specific object or situation that leads to avoidance or significant distress.

\medterm{Phobic Disorder} A disorder characterized by clinically significant fear and avoidance of a particular object or situation.
\medterm{Phocomelia} A rare congenital limb-development defect in which hands or feet are attached close to the trunk because of absent or severely shortened long bones.
\medterm{Phonation} Production of voiced sound when exhaled air makes the vocal folds vibrate in the larynx; disorders cause hoarseness, weak voice, or loss of voice.

\medterm{Phonocardiogram} A graphic recording of heart sounds and murmurs, traditionally made with a microphone over the chest. It can document timing and duration but is usually interpreted with examination and echocardiography.
\medterm{Phonocardiography} Recording heart sounds and murmurs with a microphone, producing a tracing that can help evaluate heart valve problems.

\medterm{Phonological Disorder} A communication disorder involving persistent difficulty producing or organizing speech sounds beyond what is expected for age and language development.

\medterm{Phonophobia} Unusual sensitivity to sound or distress caused by ordinary sounds. It often occurs with migraine but may also accompany concussion, infection, anxiety, or other conditions; new severe sound sensitivity with headache or neurological symptoms needs assessment.

\medterm{Phorate} Phorate is a highly toxic organophosphate pesticide whose exposure can produce excessive cholinergic activity and medical emergency.


\medterm{Phorbols} Phorbols are diterpene compounds related to phorbol esters that can activate cellular protein kinase C signaling.

\medterm{Phorometer} An instrument used to measure ocular alignment and the amount of prism correction needed for binocular vision disorders.

\medterm{Phosgene} A highly toxic industrial gas that reacts with moist respiratory tissues and can cause delayed pulmonary edema, cough, chest tightness, and respiratory failure after exposure.

\medterm{Phosgene Oxime} A highly corrosive chemical warfare agent that causes immediate severe pain, skin and eye injury, airway damage, and tissue necrosis on contact.

\medterm{Phosmet} Phosmet is an organophosphate insecticide that can inhibit acetylcholinesterase and cause cholinergic toxicity after substantial exposure.

\medterm{Phosphamidon} A poisonous organophosphate insecticide. Exposure can cause sweating, excess saliva, vomiting, diarrhea, small pupils, muscle twitching, breathing difficulty, seizures, or weakness because it disrupts nerve signals; suspected poisoning needs urgent medical help.

\medterm{Phosphate} An inorganic ion and component of nucleotides, phospholipids, proteins, and bone mineral. It participates in ATP energy transfer, phosphorylation-dependent signaling, intracellular buffering, and acid excretion by the kidney.
\medterm{Phosphate Binders} Medicines taken with meals that bind phosphate in food inside the intestine, reducing its absorption. They are used mainly in advanced kidney disease when blood phosphate is high. The choice depends on calcium levels, vascular calcification, other illnesses, and treatment response.

\medterm{Phosphate Homeostasis} The body keeps blood phosphate within a useful range through the coordinated actions of the parathyroid glands, vitamin D, kidneys, intestines, and bones. Kidney disease can disrupt this balance.

\medterm{Phosphatidic Acid} A membrane phospholipid and signaling intermediate that serves as a precursor for other glycerolipids and influences membrane curvature, trafficking, and cell growth.

\medterm{Phosphatidylcholine} A major membrane phospholipid and component of pulmonary surfactant, lipoproteins, and bile; it also participates in cell signaling.

\medterm{Phosphatidylethanolamines} Membrane phospholipids containing ethanolamine that help determine membrane structure, fusion, and cell signaling. A group of cell-membrane fats that help maintain membrane shape, join membranes, and support cell signaling.
\medterm{Phosphatidylinositol} A membrane phospholipid that serves as a precursor for signaling molecules and helps anchor proteins to cell membranes.

\medterm{Phosphatidylserine} A membrane phospholipid normally concentrated on the inner cell surface that becomes exposed during apoptosis and helps signal cell removal.

\medterm{Phosphatidylserines} Molecules of phosphatidylserine, a membrane phospholipid involved in apoptosis, coagulation, membrane structure, and cell signaling.
\medterm{Phosphaturic Mesenchymal Tumor} A rare tumor that releases a substance causing the kidneys to lose too much phosphate. Low phosphate can soften bones and cause bone pain, weakness, or fractures; removing the tumor often corrects the problem.

\medterm{Phosphine} A highly toxic gas used in some industrial or pest-control settings that can cause severe lung, heart, liver, kidney, and neurologic injury.

\medterm{Phosphodiesterase-5 Inhibitors} Medicines that block an enzyme called PDE-5, increasing a chemical signal that relaxes blood vessels. They are used for erectile dysfunction and some forms of high blood pressure in the lung; nitrates must not be combined with them.

\medterm{Phosphofructokinase-1} A key enzyme that controls how cells break down sugar to produce energy, especially when energy demand increases.

\medterm{Phospholipase C Pathway} A system cells use to turn some hormone or nerve signals at their surface into internal actions. It can change calcium levels and control contraction, secretion, growth, or other cell responses.

\medterm{Phospholipid} A type of lipid with a water-attracting phosphate-containing head and water-repelling fatty-acid tails. Phospholipids form the basic double-layered structure of cell membranes and also help with cell signaling and transport.
\medterm{Phosphorus} An essential element in bone, nucleic acids, phospholipids, and ATP.
\medterm{Phosphorus Blood Test} A blood test measuring inorganic phosphate, interpreted with calcium, parathyroid hormone, vitamin D, and kidney function to evaluate mineral and renal disorders.

\medterm{Phosphorus In Diet} Dietary phosphorus comes from protein-rich foods and many processed foods; people with kidney disease may need individualized limits.

\medterm{Phosphorylation} Addition of a phosphate group to a molecule, often regulating protein activity or energy transfer.
\medterm{Photoallergy} An immune-mediated skin reaction that occurs when ultraviolet light activates a sensitizing substance on or
in the skin, causing an eczematous rash.

\medterm{Photocoagulation} Treatment that uses intense light, usually a laser, to heat and seal targeted tissue. In the eye, it can close leaking or abnormal blood vessels and limit damage from conditions such as diabetic retinopathy.

\medterm{Photodermatosis} A skin disorder caused or aggravated by exposure to ultraviolet or visible light, including polymorphous light eruption, phototoxic reactions, and photoallergic reactions.

\medterm{Photodynamic Therapy} A treatment that combines a photosensitizing agent with light to generate reactive
oxygen species and destroy cancer cells, used for actinic keratosis, superficial skin
cancers, and selected visceral tumors.

\medterm{Photodynamic Therapy For Cancer} Treatment in which a light-sensitive medicine collects in selected abnormal cells and is activated by a specific light, producing substances that destroy nearby cancer cells. It is used for certain accessible cancers and may cause pain, swelling, or temporary light sensitivity.

\medterm{Photoelectric Effect} An interaction in which tissue absorbs an X-ray and releases an electron. It contributes to how much radiation is absorbed and helps create contrast between tissues on an X-ray image.

\medterm{Photographic Fixative Poisoning} Harm from swallowing or inhaling chemicals used to process photographs. Effects depend on the product but may include irritation, vomiting, breathing problems, or chemical burns; keep the container and contact poison control.

\medterm{Photomultiplier Tube} A sensitive device that converts tiny flashes of light into electrical signals. In some medical scanners, groups of these devices help detect and locate radiation emitted by a tracer.

\medterm{Photopeak} The part of a radiation-measurement graph showing photons that deposited their full energy in the detector. Clinicians use it to select the energy range for interpreting some nuclear medicine scans.

\medterm{Photopheresis} A treatment in which some blood cells are removed, mixed with a light-sensitive medicine, exposed to ultraviolet light, and returned to the body. It is used for selected immune, blood, and skin disorders under specialist care.

\medterm{Photophobia} Unusual discomfort or pain caused by light, sometimes leading a person to avoid bright surroundings. It may occur with migraine, corneal injury, eye inflammation, infection, concussion, or meningitis. Severe eye pain, vision loss, fever, stiff neck, or a sudden severe headache needs urgent care.
\medterm{Photoreceptors} Specialized cells in the retina that detect light. Rods support dim-light and peripheral vision, while cones provide color and sharp central vision; damage can reduce sight or cause blind areas.

\medterm{Photorefractive Keratectomy} Laser surgery that removes a very thin surface layer of the cornea and reshapes the layer beneath it. It can reduce short-sightedness, long-sightedness, or astigmatism, but recovery is slower than some procedures and pain, haze, infection, or imperfect vision can occur.

\medterm{Photosensitivity} An unusually strong reaction of the skin to sunlight, causing redness, burning, itching, blistering, or a rash on exposed areas. It may result from a disease, inherited tendency, medicine, or chemical.

\medterm{Photosensitization} Increased sensitivity to ultraviolet or visible light caused by a medication, chemical, or disease, producing an exaggerated skin reaction.

\medterm{Phototherapy} Treatment using carefully controlled light for conditions such as newborn jaundice, psoriasis, eczema, or seasonal depression. The wavelength, dose, duration, and eye or skin protection depend on the condition treated.
\medterm{Phototoxic Dermatitis} A sunburn-like skin injury caused when a medicine or chemical makes the skin react strongly to ultraviolet light. It affects exposed areas and may cause redness, pain, swelling, blisters, or later darkening.

\medterm{Phototoxicity} Skin injury occurring when a chemical makes skin unusually sensitive to ultraviolet light. It resembles severe sunburn and may cause redness, pain, blistering, or later darkening after medicine or chemical exposure.

\medterm{Phrenic} Relating to the diaphragm or the phrenic nerves, which arise from the neck and carry the main motor supply to the diaphragm for breathing.
\medterm{Phrenic Nerve} The nerve arising mainly from C3-C5 that supplies the diaphragm and carries sensory fibers from nearby structures.
\medterm{PHT} Alternate index label; see \textit{Portal hypertension}.

\medterm{Phthiriasis} Infestation with lice, especially pubic lice, causing itching and sometimes visible lice or eggs attached to hair. Treatment usually includes an approved lice-killing product, washing or sealing recently used clothing and bedding, and evaluation or treatment of close or sexual contacts.
\medterm{Phyllodes Tumor} An uncommon breast growth made of supporting tissue and glandular tissue. It may be benign, borderline, or cancerous and often appears as a rapidly enlarging lump; a tissue examination determines its type and guides surgery.

\medterm{Physiatrist} A doctor who specializes in rehabilitation and in improving movement, function, independence, and quality of life after illness, injury, or disability.

\medterm{Physical Activity} Any movement produced by skeletal muscles that increases energy expenditure, including occupational, household, transportation, and recreational movement. Regular activity supports cardiovascular, metabolic, musculoskeletal, and mental health.

\medterm{Physical Child Abuse} Physical child abuse is intentional physical injury or harmful force inflicted on a child, requiring immediate safety assessment, medical care, and safeguarding action.

\medterm{Physical Dependence} A state in which the body adapts to a substance or medication and develops withdrawal symptoms when it is reduced or stopped.

\medterm{Physical Exam Frequency} How often a person should receive a physical examination. The interval depends on age, symptoms, chronic conditions, medicines, pregnancy, and preventive-care recommendations rather than one schedule for everyone.

\medterm{Physical Examination} A systematic assessment of the body using inspection, palpation, percussion, auscultation, and other examination methods.
\medterm{Physical Half-Life} The time needed for half the radioactive atoms in a substance to change through radioactive decay. It describes the substance itself and does not include how quickly the body removes it.

\medterm{Physical Medicine} The medical specialty focused on function, rehabilitation, pain, and restoration of independence after illness or injury.
\medterm{Physical Therapy} Rehabilitation using exercise, movement training, manual techniques, and physical modalities to improve function, mobility, and reduce disability.

\medterm{Physical Therapy Assessment} A detailed evaluation of movement, strength, flexibility, sensation, balance, walking, pain, and daily activities. The findings guide treatment goals, exercise choices, equipment needs, safety advice, and progress monitoring.

\medterm{Physician} A doctor trained and licensed to diagnose illness, provide medical treatment, prevent disease, and coordinate care within the scope of the relevant specialty and jurisdiction.
\medterm{Physician Assistant} A nationally certified and state-licensed medical clinician who practices medicine collaboratively with or under the supervision of physicians. PAs obtain medical histories, perform physical examinations, order and interpret diagnostic studies, formulate medical diagnoses, prescribe pharmacotherapy, assist in surgery, and perform clinical procedures across primary care and surgical specialties.


\medterm{Physicians' Desk Reference} A reference source containing prescribing information, indications, contraindications, dosing, warnings, and adverse effects for medicines and selected products.

\medterm{Physiologic} Physiologic means relating to normal body function or a process occurring within the expected biological range.

\medterm{Physiologic Jaundice Of Newborn} A common temporary yellowing of a newborn’s skin and eyes that begins after the first day of life. It results from normal newborn red-cell breakdown and immature liver processing of bilirubin; jaundice that starts earlier or becomes severe needs urgent assessment.

\medterm{Physiologic Jaundice Of The Newborn} Physiologic newborn jaundice is a temporary bilirubin rise after birth caused by immature bilirubin processing, but early or severe jaundice needs urgent evaluation.

\medterm{Physiological Dead Space} The part of a breath that does not exchange oxygen and carbon dioxide with blood. It includes the air passages and lung areas that receive little or no blood flow.

\medterm{Physiology} The study of how healthy body parts, organs, and living organisms normally work together. It examines processes such as breathing, circulation, digestion, movement, hormone action, nerve signaling, and maintaining internal balance.
\medterm{Physiotherapy} Assessment and treatment using movement, exercise, manual techniques, education, and physical modalities to improve function.
\medterm{Physostigmine} A reversible acetylcholinesterase inhibitor that may reverse severe central anticholinergic delirium in carefully selected patients after assessment for contraindications and cardiac conduction abnormalities.


\medterm{Phytomenadione} A form of vitamin K needed by the liver to make several blood-clotting proteins. It is used to prevent or treat bleeding caused by vitamin K deficiency or by medicines that block vitamin K.

\medterm{PI-RADS} A standard scoring system for prostate MRI findings. It estimates how likely an area is to contain an important prostate cancer and helps guide decisions about follow-up or biopsy; it does not diagnose cancer by itself.
\medterm{Pia Mater} The thin, delicate inner covering that closely follows the surface of the brain and spinal cord. Together with the other coverings and fluid around the nervous system, it helps protect these organs.
\medterm{Pica} Repeated eating of things that have no useful nutritional value, such as soil, paper, ice, or hair, when this is not expected for a person's age or culture. It may cause poisoning, infection, blockage, or poor nutrition and can be linked with iron deficiency or mental-health conditions.
\medterm{Pickwickian Syndrome} An older name for obesity hypoventilation syndrome, involving obesity, daytime hypercapnia, and sleep-disordered breathing.
\medterm{PID} Alternate index label; see \textit{Pelvic inflammatory disease}.

\medterm{Piebaldism} A rare condition present from birth in which some areas of skin and hair lack pigment. It usually causes stable, sharply outlined white patches and often a white section of hair at the front; it is commonly inherited through a change in the KIT gene.

\medterm{Pierre Robin Sequence} A condition present at birth in which a small lower jaw allows the tongue to fall backward and may block breathing. Many affected infants also have a U-shaped opening in the roof of the mouth and feeding difficulty.

\medterm{Pierson Syndrome} A rare inherited condition causing severe kidney protein loss from birth and unusually small pupils. It can lead to swelling, kidney failure, poor vision, and developmental problems.

\medterm{Pigment} A substance that gives skin, hair, eyes, tissue, or other material its color. Pigments may be made by the body, received from outside, or altered by disease, injury, medicines, or aging.

\medterm{Pigment Cirrhosis} Liver cirrhosis associated with abnormal pigment or iron deposition, as may occur in hereditary hemochromatosis.
\medterm{Pigmentary Mosaicism} Uneven skin coloring caused by groups of skin cells carrying different genetic information. It may produce lighter or darker lines, swirls, or patches; most cases affect only the skin, but associated developmental problems should be assessed when present.

\medterm{Pigmentation} Coloration of skin, hair, eyes, or other tissue produced mainly by melanin, carotenoids, blood, or exogenous material; abnormal changes may reflect inflammation, endocrine disease, drugs, or tumors.

\medterm{Pigmentation Disorder} A condition causing abnormal increase, decrease, or uneven distribution of pigment in the skin, mucosa, hair, or eye.

\medterm{Pilar Cyst} A usually harmless, firm lump beneath the skin that forms from a hair follicle, most often on the scalp. It may be inherited, become inflamed or painful, and can usually be removed if it causes symptoms or repeated problems.

\medterm{Piles} Enlarged or swollen blood vessels and supporting tissue in the anus or lower rectum. Internal piles may cause painless bright-red bleeding or tissue protruding during bowel movements; external piles may cause pain, itching, or a tender lump.

\medterm{Piles (Hemorrhoids (Piles))} Alternate index label; see \textit{Hemorrhoids}.

\medterm{Piles (Hemorrhoids)} Alternate term; see \textit{Hemorrhoids}.

\medterm{Pill} A solid medicine designed to be taken by mouth. It may contain one or more active ingredients and may be made to release them immediately or slowly; pills should be taken as directed and not crushed unless the instructions allow it.

\synonyms
Tablet

\medterm{Pilocarpine} A medicine that activates muscarinic receptors, causing pupil constriction and increased secretions. It may lower eye pressure in glaucoma or relieve dry mouth, but can cause sweating, blurred vision, headache, or breathing problems.
\medterm{Pilocytic Astrocytoma} A usually slow-growing brain tumor arising from support cells, most often in children. It commonly affects the cerebellum or visual pathway and may cause headache, vomiting, seizures, balance problems, or vision changes.

\medterm{Piloerection} Elevation of hair shafts producing gooseflesh, caused by contraction of arrector pili muscles under sympathetic stimulation from cold, fear, pain, fever, or certain drugs.

\medterm{Pilomatrix Carcinoma} A rare skin cancer that develops from cells involved in hair formation. It usually appears as a firm lump on the head, neck, or arm and can grow into nearby tissue or spread elsewhere.

\medterm{Pilomatrixoma} Pilomatrixoma is a usually benign skin tumor derived from hair-matrix cells, presenting as a firm subcutaneous nodule, often on the head or neck.

\medterm{Pilonidal Cyst} A cyst or sinus in the cleft between the buttocks, usually containing hair and skin debris. It may become infected, painful, swollen, and drain pus.

\synonyms
Cyst, Pilonidal
Pilonidal Sinus (Pilonidal Cyst)

\medterm{Pilonidal Dimple} Alternate index label; see \textit{Sacral dimple}.

\medterm{Pilonidal Disease} A condition in which hair and skin debris collect in a small pit or tunnel in the cleft between the buttocks. It can cause repeated swelling, pain, drainage, or an abscess.

\medterm{Pilonidal Fistula} An abnormal tunnel from a pilonidal sinus to the skin, usually in the natal cleft, that can drain pus or recur after abscess formation.

\medterm{Pilonidal Sinus (Pilonidal Cyst)} Alternate index label for Pilonidal Cyst.

\medterm{Pilonidal Sinus Disease} A chronic sinus or cyst in the natal cleft, usually containing hair and debris. It may
cause recurrent pain, swelling, drainage, or abscess formation.

\medterm{Pilus} A tiny hair-like structure on some bacteria. It may help bacteria attach to human cells, move, or transfer genetic material to other bacteria, which can contribute to infection or antibiotic resistance.

\medterm{Pimozide} An antipsychotic medicine used for selected disorders involving severe tics or abnormal thoughts. It can cause movement problems, sleepiness, and dangerous heart-rhythm changes, so dosing and heart monitoring require careful medical supervision.
\medterm{Pimple} A common lay term for a small inflamed or blocked skin lesion, often an acne papule or pustule.

\medterm{Pimples (Acne)} Alternate index label for Acne.

\medterm{Pin Care} Pin care keeps the skin around external fixation or traction pins clean and monitored for redness, swelling, drainage, loosening, or infection.

\medterm{Pinch Strength} Pinch strength is the force produced when the thumb presses against one or more fingers and can help assess hand, nerve, muscle, and functional impairment.

\medterm{Pinched Nerve} Compression or irritation of a nerve that can cause pain, tingling, numbness, or weakness along the nerve's distribution.

\synonyms
Compressed Nerve
Nerve Compression

\medterm{Pindolol} Pindolol is a beta-blocker used for selected cardiovascular conditions; it can slow the heart rate and worsen asthma in susceptible people.

\medterm{Pine Oil Poisoning} Swallowing or inhaling concentrated pine oil can irritate the mouth and stomach and may cause drowsiness, seizures, or lung injury if aspirated. Do not induce vomiting; obtain immediate poison-control advice.

\medterm{Pineal Cyst} A fluid-filled sac in the small gland deep in the brain that helps regulate sleep-related signals. Most are found incidentally and cause no harm, but a large cyst can rarely cause headache, eye problems, or fluid buildup.

\medterm{Pineal Germinoma} A pineal germinoma is a malignant germ-cell tumor in the pineal region, often causing headache, eye-movement abnormalities, or hydrocephalus and usually responsive to radiotherapy and chemotherapy.

\medterm{Pineal Gland} A small hormone-producing gland deep in the brain that releases melatonin. Melatonin helps coordinate the body's day-night clock and influences sleep timing, seasonal rhythms, and related biological functions.
\medterm{Pineal Gland Tumor} A growth in or near the pineal gland. It may block the normal flow of brain fluid and cause headache, nausea, vomiting, sleepiness, or problems moving the eyes; symptoms and treatment depend on its type and size.

\medterm{Pineal Parenchymal Tumor} A tumor arising from cells of the pineal gland, ranging from slow-growing to aggressive. It may block the flow of brain fluid and cause headache, vomiting, sleep changes, or abnormal eye movements.

\medterm{Pineoblastoma} An aggressive cancer in the pineal region of the brain, occurring mainly in children. It may block brain-fluid flow and cause headache, vomiting, sleepiness, or eye-movement problems and requires specialist treatment.

\medterm{Pineocytoma} Pineocytoma is a usually slow-growing tumor of pineal cells that may obstruct cerebrospinal-fluid flow and cause raised intracranial pressure.

\medterm{Pinguecula} A pinguecula is a benign yellowish raised growth on the conjunctiva near the cornea, often associated with ultraviolet light, wind, and dust; it may irritate the eye.

\medterm{Pinhole Meatus} An unusually tiny external opening, most often the opening of the urethra. It can make urine flow weak or spraying, cause discomfort, or make catheter placement difficult.

\medterm{Pink Eye} A common name for conjunctivitis, inflammation of the eye covering causing redness, irritation, watering, or discharge.

\synonyms
Conjunctivitis (Pink Eye)

\medterm{Pink Puffer} An older, nonformal term for a person with emphysema who has severe breathlessness, a thin build, and pursed-lip breathing. The term is not a separate diagnosis and is less useful than describing the person's COPD and test results.


\medterm{Pinky Broken Finger (Broken Finger)} Alternate index label for Broken Finger.

\medterm{Pinna} The visible external part of the ear that collects sound.
\medterm{Pinocytosis} A process in which a cell takes in small drops of surrounding fluid and dissolved substances. The cell membrane folds inward to form tiny bubbles that carry the material inside for use, transport, or breakdown.
\medterm{Pint} A pint is a unit of volume; its exact amount varies by measurement system, so the local standard must be specified when measuring fluids or medicines.
\medterm{Pinta} A long-lasting skin infection caused by Treponema carateum bacteria. It spreads through close skin contact and produces changing stages of discolored, scaly, or thickened patches, usually without damage to internal organs.
\medterm{Pinworm} A small intestinal worm that commonly causes intense itching around the anus at night. Infection spreads when microscopic eggs are swallowed or transferred from contaminated hands or surfaces; treatment usually includes medicine, hygiene, and sometimes treating household contacts.
\medterm{Pinworm Infection} Pinworm infection is an intestinal infestation by Enterobius vermicularis, causing nighttime anal itching, disturbed sleep, and sometimes abdominal or genital irritation. Treatment of the person and close household contacts, hygiene, and repeat dosing reduce reinfection.

\synonyms
Infection Pinworms (Pinworm Infection)
Worms Pinworms (Pinworm Infection)

\medterm{Pinworm Test} A test for Enterobius vermicularis, usually the early-morning adhesive-tape test applied to the perianal skin before bathing or defecation.

\medterm{Pinworms} Pinworms are intestinal nematodes, Enterobius vermicularis, that cause nighttime anal itching; treatment, hand hygiene, and laundering help prevent household reinfection.

\medterm{Pioglitazone} A medicine for type 2 diabetes that helps muscle and fat cells respond better to insulin, lowering blood sugar. It can cause weight gain and swelling and may worsen heart failure, so it is unsuitable for some people.



\medterm{Piperacillin} Piperacillin is an extended-spectrum penicillin antibiotic often combined with tazobactam to treat serious infections including Pseudomonas; allergy, diarrhoea, cytopenias, electrolyte effects, and renal dosing are important.

\medterm{Piperazine} An older antiparasitic medicine that paralyzes certain intestinal worms, especially roundworms and pinworms. Its use depends on local guidance and the particular infection.
\medterm{Piperonyl Butoxide} A pesticide synergist that inhibits insect detoxification enzymes and increases the activity of certain insecticides; it can irritate skin, eyes, or airways.

\medterm{Piperonyl Butoxide With Pyrethrins Poisoning} Exposure to these insecticide ingredients can irritate the skin, eyes, or airways and may cause coughing, nausea, vomiting, dizziness, or wheezing. Severe breathing symptoms or swallowed product require urgent poison-control or emergency advice.



\medterm{Pipotiazine} A long-acting phenothiazine antipsychotic given by injection in selected settings. It can cause movement disorders, drowsiness, and other adverse effects and requires clinical monitoring.
\medterm{Pipradrol} An older central nervous system stimulant with limited current use.
\medterm{Piracetam} A nootropic compound used in some countries for selected neurologic indications; evidence and approval vary.
\medterm{Pirfenidone} Pirfenidone is an antifibrotic medicine used to slow decline in idiopathic pulmonary fibrosis; nausea, photosensitivity, fatigue, and liver-enzyme elevation may occur and require monitoring.

\medterm{Piriformis Syndrome} Pain or altered sensation in the buttock or down the leg when the piriformis muscle irritates or presses on the sciatic nerve. Other causes, especially a spinal nerve problem, must be considered before treatment.

\synonyms
PiriformSyndrome (Piriformis Syndrome)

\medterm{PiriformSyndrome (Piriformis Syndrome)} Alternate index label for Piriformis Syndrome.

\medterm{Pirimicarb} Pirimicarb is a selective carbamate insecticide targeting aphids; exposure can inhibit acetylcholinesterase and cause cholinergic symptoms.



\medterm{Piromidic Acid} An older quinolone-like antibacterial investigated or used for urinary and gastrointestinal infections; current use depends on local availability, susceptibility, and safety guidance.

\medterm{Piroxicam} A long-acting nonsteroidal anti-inflammatory drug used for pain and inflammation; gastrointestinal bleeding, kidney injury, cardiovascular events, and hypersensitivity are important risks.

\medterm{Piroxicam Overdose} Taking too much piroxicam, an anti-inflammatory pain medicine, can cause stomach pain, vomiting, bleeding, kidney injury, drowsiness, seizures, or breathing problems. Large or intentional overdoses need immediate poison-control or emergency advice.

\medterm{Pisa Syndrome} A sustained involuntary sideways flexion of the trunk, often associated with medicines, Parkinson disease, or other neurologic disorders.

\medterm{Pisiform} A small, pea-shaped wrist bone embedded in a tendon on the little-finger side of the palm. It helps form the passageway for the ulnar nerve and artery and may hurt after injury or with arthritis.

\medterm{Pisiform Bone} A small pea-shaped carpal bone embedded in the tendon of flexor carpi ulnaris on the ulnar side of the wrist; fracture or arthritis can cause localized wrist pain.

\medterm{Pits} Small depressions or holes in a surface. In medicine, pits may occur in nails, skin, teeth, or other tissues and can result from normal anatomy, injury, infection, inflammation, or inherited conditions.
\medterm{Pitting Edema} Swelling in which pressure applied to the skin leaves a temporary indentation, usually because of excess interstitial fluid.

\medterm{Pituicytoma} A rare, usually slow-growing tumor near the pituitary gland. It may cause headache, vision problems, or hormone disturbances by pressing on nearby structures and is usually assessed with brain imaging and tissue examination.

\medterm{Pituitary} The pituitary is an endocrine gland at the base of the brain that releases hormones controlling growth, thyroid and adrenal function, reproduction, lactation, and water balance.

\medterm{Pituitary Adenoma} A usually noncancerous growth arising from the front part of the pituitary gland. It may release too much of a hormone or press on nearby tissue, causing headache, vision changes, or hormone problems.

\medterm{Pituitary Apoplexy} Sudden bleeding or loss of blood supply inside a pituitary growth, causing severe headache, vision changes, vomiting, eye-movement problems, or confusion. It is an emergency because hormone failure can cause life-threatening low blood pressure.

\medterm{Pituitary Carcinoma} Pituitary carcinoma is a rare malignant pituitary tumor defined by spread to distant or separate sites, often producing excess hormones.

\medterm{Pituitary Disease} A disorder of the pituitary gland that alters hormone secretion or causes mass effects, including adenomas, inflammation, hemorrhage, genetic disease, or hypopituitarism.

\medterm{Pituitary Dwarfism} Pituitary dwarfism is severe short stature caused by deficient growth hormone or broader pituitary dysfunction, with evaluation focused on endocrine and genetic causes.

\medterm{Pituitary Gland} An endocrine gland at the base of the brain that regulates growth, thyroid, adrenal, reproductive, lactation, and water-balance hormones.
\medterm{Pituitary Hormone} A hormone secreted by the anterior or posterior pituitary that regulates growth, thyroid and adrenal function, reproduction, lactation, water balance, or other endocrine organs.

\medterm{Pituitary Insufficiency} Alternate index label; see \textit{Hypopituitarism}.

\medterm{Pituitary Neoplasm} A pituitary neoplasm is an abnormal growth in or near the pituitary gland and may affect hormones, vision, or nearby brain structures.

\medterm{Pituitary Neuroendocrine Tumor} Alternate index label; see \textit{Pituitary tumors and adenomas}.

\medterm{Pituitary Tumor} A pituitary tumor is a growth in the pituitary gland that may be noncancerous but can cause hormone excess, hormone deficiency, or visual problems.

\medterm{Pituitary Tumor (Prolactinoma)} Alternate index label for Prolactinoma.

\medterm{Pityriasis} A group of skin conditions marked by fine, flaky scaling. Examples include pityriasis rosea and pityriasis versicolor, which have different causes and may produce itching or lighter, darker, or pink patches.
\medterm{Pityriasis Alba} Pityriasis alba is a common benign skin condition causing faint, dry, lighter patches, especially on children’s faces, often more noticeable after tanning.

\medterm{Pityriasis Lichenoides} An uncommon inflammatory skin disorder producing recurrent small scaly papules, with forms ranging from mild chronic lesions to an acute ulceronecrotic eruption.

\medterm{Pityriasis Rosea} An acute, self-limited eruption often beginning with a herald patch followed by oval
scaly lesions distributed along skin-cleavage lines. It usually resolves within several
weeks; treatment is symptomatic, and atypical, severe, recurrent, or persistent disease
should be reassessed.

\synonyms
Pityriasis Rosea Gibert (Pityriasis Rosea)

\medterm{Pityriasis Rosea Gibert (Pityriasis Rosea)} Alternate index label for Pityriasis Rosea.

\medterm{Pityriasis Rubra Pilaris} A group of rare, usually long-lasting inflammatory skin disorders. Small rough bumps form around hair follicles and may join to make well-defined, reddish-orange, scaly patches. Areas of normal skin may remain between the patches. The palms and soles may become thick, dry, yellow-orange, and sometimes cracked. The cause is often unknown, and the condition can vary in extent and duration.

\medterm{Pityriasis Versicolor} Alternate index label; see \textit{Tinea versicolor}.

\medterm{Pivampicillin} An orally absorbed prodrug of ampicillin used for selected bacterial infections; its use is limited in some settings by adverse effects and availability.

\medterm{Pixantrone} Pixantrone is an anthracenedione antineoplastic drug developed for selected aggressive lymphomas, with potential marrow, infection, and cardiac risks.

\medterm{PJI} Alternate index label; see \textit{Periprosthetic joint infection}.

\medterm{PJS} Alternate index label; see \textit{Peutz-Jeghers syndrome}.

\medterm{PKD} Alternate index label; see \textit{Polycystic kidney disease}.

\medterm{Placebo} An inactive substance, treatment, or procedure designed to resemble an active treatment. In a clinical trial, it is given to a control group so researchers can compare outcomes and estimate the specific effect of the treatment being studied.

\synonyms
Dummy Treatment
Inactive Treatment
\medterm{Placebo Control} A placebo control is an inactive treatment designed to resemble the intervention in a clinical trial, helping separate treatment effects from expectation and natural change.

\medterm{Placebo Effect} A real change in symptoms or wellbeing that can occur after a person receives an inactive treatment, partly because of expectations, attention, and the treatment setting. It does not prove that the inactive treatment changed the disease itself.

\medterm{Placenta} The temporary organ connecting mother and fetus, allowing exchange of oxygen, nutrients, wastes, and hormones.
\medterm{Placenta Abruptio} Placental abruption is premature separation of the placenta from the uterine wall, causing bleeding, abdominal pain, uterine tenderness, fetal compromise, or maternal shock.

\medterm{Placenta Abruption - Definition} Placental abruption is premature separation of the placenta from the uterine wall before birth, causing bleeding, abdominal pain, uterine tenderness, and possible fetal or maternal compromise.

\medterm{Placenta Accreta} Placenta accreta spectrum occurs when placental tissue abnormally attaches to or invades the uterine wall, creating major risk of hemorrhage during attempted delivery or removal.

\medterm{Placenta Accreta Spectrum} Abnormal adherence or invasion of placental villi into the uterine wall, ranging from accreta to increta and percreta; it can cause severe postpartum hemorrhage and requires planned obstetric management.

\medterm{Placenta Disorder} A disorder of placental implantation, structure, perfusion, or separation, including previa, accreta spectrum, abruption, infarction, or placental insufficiency.

\medterm{Placenta Increta} Abnormal placental invasion into the myometrium without reaching the serosa, increasing
the risk of massive hemorrhage at delivery and often requiring cesarean hysterectomy.

\medterm{Placenta Percreta} Placenta percreta is the most invasive form of placenta accreta spectrum, with villi growing through the uterine wall and sometimes into nearby organs such as the bladder.

\medterm{Placenta Praevia} A condition in which the placenta lies over or near the opening of the cervix. It can cause painless bleeding later in pregnancy and may require monitoring, activity advice, or cesarean delivery.
\medterm{Placenta Previa} Implantation of the placenta over or near the internal cervical opening, which can cause painless bleeding during the second half of pregnancy.

\medterm{Placenta, Succenturiate} A succenturiate placenta has one or more accessory placental lobes separate from the main placenta, which can cause retained tissue or fetal-vessel complications after delivery.

\medterm{Placental Abnormalities} Disorders of placental location, attachment, structure, or separation that may affect pregnancy or cause bleeding.
\medterm{Placental Abruption} Premature separation of the placenta from the uterine wall before delivery. It may cause
vaginal bleeding, abdominal pain, uterine tenderness, fetal distress, and maternal or fetal compromise.

\medterm{Placental Accreta Spectrum} Abnormal placental attachment to the myometrium, ranging from accreta (adherent to
decidua basalis) to increta (invading myometrium) to percreta (invading through serosa).
Major risk factor is prior caesarean section with anterior placenta previa.

\medterm{Placental Hormone} A hormone produced by placental tissue, such as hCG, progesterone, estrogens, or human placental lactogen, that supports pregnancy, fetal growth, and maternal metabolic adaptation.

\medterm{Placental Infarction} A region of placental tissue damaged by inadequate maternal blood flow and coagulation, which may be incidental or, when extensive, associated with fetal growth restriction or pregnancy complications.

\medterm{Placental Insufficiency} Placental insufficiency is inadequate placental delivery of oxygen and nutrients, which can cause fetal growth restriction, abnormal Doppler findings, reduced fluid, or fetal distress.

\medterm{Placental Lactogen} Placental lactogen is a hormone made by the placenta that alters maternal metabolism and supports breast development and fetal nutrient availability during pregnancy.

\medterm{Placental Mesenchymal Dysplasia} A rare abnormal development of the placenta in which supporting tissue becomes enlarged and cyst-like. It can resemble a molar pregnancy on scans and may be associated with poor fetal growth, fetal abnormalities, or pregnancy complications.

\medterm{Placental Separation Signs} Clinical signs indicating the placenta has separated from the uterine wall and is ready
for delivery. Classic signs: lengthening of the umbilical cord (cord lengthens as the
placenta descends), a sudden gush of blood, and uterine elevation and firming (the
uterus rises and becomes globular as the placenta moves into the lower segment).

\medterm{Placental Site Nodule} A benign, usually incidental remnant of intermediate trophoblast at a prior implantation site, composed of small nests of chorionic-type or extravillous trophoblast.

\medterm{Placentation} The developmental process by which the placenta forms and attaches to the uterus, establishes maternal--fetal circulation, and develops the exchange and endocrine functions needed for pregnancy.

\medterm{Plagiocephaly} Plagiocephaly is asymmetric flattening or shaping of an infant’s skull, caused by position or premature suture fusion; its evaluation distinguishes benign molding from craniosynostosis.

\synonyms
Flat Head Syndrome; Deformational Plagiocephaly

\medterm{Plague} A serious infection caused by Yersinia pestis, transmitted mainly by infected fleas or respiratory droplets in pneumonic disease.


\medterm{Planar} Flat or arranged in one plane. In medical imaging, the term describes a two-dimensional view or projection, unlike three-dimensional imaging that shows depth and allows reconstruction from multiple directions.
\medterm{Plane Warts (Verruca Plana)} Small, smooth, relatively flat warts caused by HPV, often occurring in groups on the face, hands, or legs.

\synonyms
Verruca Plana





\medterm{Plant Fertilizer Poisoning} Fertilizer exposure may irritate the mouth, skin, eyes, stomach, or lungs. Some products can dangerously alter blood oxygen or body salts; keep the container and contact poison control after swallowing or significant exposure.



\medterm{Plant Poisoning} Illness caused by eating, touching, or inhaling a toxic plant or plant product; effects range from local irritation to organ failure.


\medterm{Plant Thorn Arthritis} Plant thorn arthritis is persistent joint inflammation after a thorn puncture introduces plant material, sometimes requiring removal of foreign material and treatment of infection or inflammation.

\medterm{Plant Thorn Synovitis} Inflammation of a joint after penetration by a plant thorn or splinter, sometimes involving an infection introduced by
the injury. Symptoms may develop gradually as persistent joint pain, swelling, and
limited movement; evaluation may require removal of retained material and targeted
treatment.

\synonyms
Arthritis Plant Thorn (Plant Thorn Synovitis)


\medterm{Plantar} Relating to the sole or underside of the foot. Plantar problems include heel pain from plantar-fascia injury, warts, pressure ulcers, and abnormal reflexes; diabetes and nerve damage can increase the risk of foot injury.
\medterm{Plantar Flexion} Downward movement of the foot at the ankle, pointing the toes away from the shin. It is produced mainly by the calf muscles and is used in standing on tiptoe, walking, running, and jumping.
\medterm{Plantar Fasciitis} Pain from irritation and small injuries in the strong band of tissue supporting the sole of the foot. Heel pain is often worst with the first steps after waking or resting and may worsen after prolonged standing.

\medterm{Plantar Nerve} A nerve supplying sensation to the sole of the foot and motor branches to intrinsic foot muscles; compression can cause burning pain, numbness, or altered toe movement.

\medterm{Plantar Neuroma} Alternate index label; see \textit{Morton neuroma}.

\medterm{Plantar Reflex} Movement of the toes after the sole of the foot is gently stroked. In most adults the toes bend downward; upward movement of the big toe may indicate a problem in brain or spinal nerve pathways controlling movement.
\medterm{Plantar Response} Movement of the toes and foot after the sole is stimulated. The pattern helps clinicians assess pathways between the brain, spinal cord, and leg; an unusual upward movement of the big toe in an adult may indicate neurologic disease.

\medterm{Plantar Wart (Verruca Plantaris)} A wart on the sole of the foot caused by human papillomavirus. Pressure from standing and walking can push it inward and make it painful; small black dots may represent tiny clotted blood vessels.

\synonyms
Verruca Plantaris

\medterm{Plantaris Muscle} A small posterior-leg muscle with a long tendon that assists weakly with knee flexion and plantar flexion; it can be strained during sudden athletic movement.

\medterm{Plaque} A localized buildup or raised area of material. In the skin, a plaque is a broad, raised lesion often formed by joined papules; in an artery, plaque is a fatty deposit in the vessel wall that can narrow the artery.
\medterm{Plaque And Tartar On Teeth} Dental plaque is a bacterial biofilm that forms on teeth, while tartar is hardened mineralized plaque; both contribute to cavities, gingivitis, and periodontitis.

\medterm{Plaque Assay} A laboratory test that estimates the number of infectious virus particles in a sample. The sample is placed on living cells, and each infectious particle may produce a visible area of damaged cells that can be counted.

\medterm{Plaque Psoriasis} The most common form of psoriasis, producing well-demarcated raised red plaques with silvery scale, commonly on elbows, knees, scalp, and lower back and sometimes associated with arthritis.

\medterm{Plasma} The liquid part of blood carrying cells, proteins, nutrients, hormones, wastes, and other dissolved substances through the body.

\medterm{Plasma Amino Acids} A blood test that measures amino acids in plasma. It can help evaluate inherited metabolic disorders, poor nutrition, liver disease, or other conditions affecting amino-acid processing.

\medterm{Plasma Cell} An immune cell that develops from a B cell and produces antibodies. Antibodies help recognize and control germs and other foreign substances.

\medterm{Plasma Cell Mastitis} A chronic benign inflammatory breast lesion with prominent plasma-cell infiltration, often involving ducts beneath the nipple and potentially causing a mass, nipple retraction, or discharge. It can clinically mimic carcinoma.

\medterm{Plasma Cell Myeloma} A cancer of antibody-producing plasma cells, usually in the bone marrow. It can produce an abnormal antibody and damage bones, kidneys, blood-cell production, or calcium balance, causing bone pain, fractures, anemia, infections, or kidney problems.

\medterm{Plasma Cell Neoplasm} A clonal disorder of antibody-producing plasma cells, ranging from localized plasmacytoma to multiple myeloma with marrow, bone, kidney, calcium, or blood abnormalities.

\medterm{Plasma Cells} Immune cells that make antibodies. They develop from B cells after the body recognizes an infection or other foreign substance. Abnormal plasma-cell growth can cause diseases such as multiple myeloma and may produce excessive abnormal antibodies.

\medterm{Plasma Concentration} The amount of a substance present in a defined volume of plasma at a given time. It is
used to interpret exposure, toxicity, therapeutic drug monitoring, and pharmacokinetic
relationships.

\medterm{Plasma Expanders} Fluids given through a vein to increase circulating blood volume. Crystalloid solutions are commonly used, while selected colloids have limited roles; treatment depends on the cause of fluid loss and patient condition.
\medterm{Plasma Kallikrein} A blood enzyme that helps produce bradykinin, a substance that widens blood vessels and increases leakage from them. It contributes to inflammation and swelling and is involved in some inherited swelling disorders.

\medterm{Plasma Membrane} The selectively permeable lipid bilayer surrounding a cell, containing proteins that control transport, signaling, adhesion, recognition, and communication with the extracellular environment.

\medterm{Plasma Osmolarity} The concentration of dissolved particles in the liquid part of blood. It mainly reflects water balance and substances such as sodium, glucose, and urea; abnormal values can indicate dehydration, excess water, or metabolic disease.

\medterm{Plasma Substitutes} Intravenous fluids intended to expand plasma volume, including crystalloids and colloid solutions; choice depends on the clinical indication and risks.

\medterm{Plasmablastic Lymphoma} Plasmablastic lymphoma is an aggressive B-cell lymphoma often associated with immune deficiency, HIV infection, or immunosuppression.

\medterm{Plasmacytoma} A solitary mass of neoplastic plasma cells, which may occur in bone (solitary plasmacytoma of bone) or soft tissue (extramedullary plasmacytoma). It can progress to multiple myeloma.
\medterm{Plasmalogen} A special type of fat found in cell membranes, especially in the brain, nerves, and heart. The body makes it partly in structures called peroxisomes; very low levels occur in some inherited metabolic disorders.

\medterm{Plasmapheresis} A specialized medical procedure that filters and separates the pale yellow liquid portion of blood (the plasma) from blood cells. In therapeutic plasma exchange, whole blood is drawn through an intravenous catheter into a machine that spins or filters it, separating out the patient's plasma---which often carries disease-causing autoantibodies, abnormal proteins, or inflammatory toxins---and discarding it. The patient's red blood cells, white blood cells, and platelets are then mixed with replacement fluids (such as donor plasma or sterile albumin solution) and safely returned to the bloodstream. Plasmapheresis is an established, life-saving treatment for severe autoimmune and neurological disorders where the immune system attacks the body's own tissues, including Guillain--Barr\'e syndrome, myasthenia gravis crisis, thrombotic thrombocytopenic purpura (TTP), chronic inflammatory demyelinating polyneuropathy (CIDP), and Goodpasture syndrome. It is also used in blood donation centers to harvest plasma from healthy donors.

\synonyms
Therapeutic Plasma Exchange; TPE; Plasma Exchange; Apheresis
\medterm{Plasmid} A plasmid is a small usually circular DNA molecule separate from a bacterial chromosome that can replicate independently and sometimes carries antibiotic-resistance genes.

\medterm{Plasmids} Small DNA rings found separately from the main chromosome in many bacteria. They can copy themselves and may carry genes that help bacteria survive, including genes for antibiotic resistance or disease-causing traits.

\medterm{Plasmin} An enzyme that breaks down fibrin, the protein mesh holding a blood clot together. It is formed from plasminogen and helps remove clots after a damaged blood vessel has been repaired.

\medterm{Plasminogen} An inactive blood protein that can be changed into plasmin, an enzyme that breaks down the protein framework of blood clots. The body activates it when a clot needs to be removed.

\medterm{Plasminogen Activator} An enzyme or medicine that changes plasminogen into plasmin, which breaks down the protein mesh of blood clots. These medicines can dissolve dangerous clots but may cause serious bleeding.

\medterm{Plasminogen Inactivators} Substances or medicines that reduce the conversion of plasminogen into plasmin, thereby slowing breakdown of blood clots. They may be used or studied when limiting bleeding is important.


\medterm{Plasmodium} A group of parasites spread through bites from infected Anopheles mosquitoes and responsible for malaria. The parasites first infect the liver, then red blood cells, causing fever and potentially severe organ complications.

\medterm{Plasmodium Vivax} A parasite spread by infected Anopheles mosquitoes that causes malaria. It infects red blood cells and can remain dormant in the liver, allowing illness to return months or years after the first infection.

\medterm{Plasmodium Vivax Malaria} Malaria caused by the parasite Plasmodium vivax and spread by infected mosquitoes. Fever attacks may recur about every two days, and dormant liver forms can cause later relapses; treatment must clear both blood and liver stages when safe.

\synonyms
Vivax malaria; Benign tertian malaria


\medterm{Plaster Of Paris} A powder that hardens when mixed with water. In healthcare it is used to make casts and splints that hold an injured bone or joint still while it heals; swelling, numbness, or increasing pain under a cast needs urgent review.
\medterm{Plastic Bronchitis} A rare condition producing large, branching mucous casts that form in the airways,
causing life-threatening airway obstruction, often associated with congenital heart
disease or post-fontan circulation.

\medterm{Plastic Casting Resin Poisoning} Exposure to uncured casting resin or its solvents can irritate or burn the skin, eyes, lungs, and digestive tract. Swallowing or breathing fumes may cause vomiting, dizziness, or lung injury and requires poison-control advice.

\medterm{Plastic Resin Hardener Poisoning} Resin hardeners may contain corrosive or irritating chemicals that burn the skin, eyes, mouth, throat, or stomach. Fumes can injure the lungs; rinse exposed areas and seek urgent poison-control advice after significant exposure.

\medterm{Plastic Surgery} Surgery that repairs, rebuilds, or changes body tissues to improve physical function, appearance, or both.
\medterm{Platelet} A small piece of a blood cell that helps stop bleeding. Platelets have no nucleus, gather at damaged blood vessels, and stick together to form an early blood clot. They are made in the bone marrow from larger cells called megakaryocytes.
\medterm{Platelet Aggregation} The clumping of platelets at an injured blood vessel to form an early plug that helps stop bleeding.

\medterm{Platelet Aggregation Test} A laboratory test measuring platelet clumping in response to agonists such as ADP, collagen, or epinephrine, used to evaluate platelet function disorders and medication effects.

\medterm{Platelet Antibodies Blood Test} A blood test detecting antibodies directed against platelet antigens, used selectively in suspected immune thrombocytopenia, neonatal alloimmune thrombocytopenia, or platelet-transfusion refractoriness.

\medterm{Platelet Count} The number of platelets in a given volume of blood, used to assess bleeding or clotting risk.
\medterm{Platelet Count Low (Thrombocytopenia (Low Platelet Count))} Alternate index label for Thrombocytopenia (Low Platelet Count).

\medterm{Platelet Disorder - Immune Thrombocytopenia} Alternate index label for ITP.

\medterm{Platelet Function Disorders} Conditions in which platelets do not adhere, activate, or aggregate normally, causing impaired primary
hemostasis. They may be inherited or acquired and can cause easy bruising, nosebleeds,
heavy menstrual bleeding, or bleeding after procedures.

\medterm{Platelet Plug Formation} The first step in stopping bleeding after a blood vessel is injured. Platelets attach to the damaged area, become active, and clump together to form a temporary plug.


\synonyms
Thrombocytes

\medterm{Platelets (Thrombocytes)} Small cell fragments made in the bone marrow that circulate in the blood and help stop bleeding. They attach to damaged blood vessels, become activated, and join together to form an initial plug.

\synonyms
Thrombocytes

\medterm{Platybasia} An unusually flattened base of the skull. It may cause no symptoms, but in some people it narrows the space around the lower brain or spinal cord and contributes to headache, imbalance, weakness, or other nerve problems.
\medterm{Platypnea} Dyspnea that occurs in the upright position and is relieved by lying down. It is the
opposite of orthopnea. It is rare and may be associated with hepatopulmonary syndrome,
intracardiac shunting, or pulmonary arteriovenous malformations.

\medterm{Platypnea-Orthodeoxia Syndrome} A rare condition in which breathing becomes difficult and blood oxygen falls when a person sits or stands, then improves when lying down. It is usually caused by abnormal blood flow through the heart or lungs. Testing may include oxygen measurements, heart imaging, and lung scans; treatment addresses the underlying cause.

\synonyms
Orthodeoxia-Platypnea Syndrome; Postural Dyspnea and Deoxygenation

\medterm{Platysma} A thin superficial muscle of the neck that tenses the skin and assists lower-jaw and lower-lip movement.
\medterm{Pleiotropic Cytokine} An immune-signaling protein that affects several types of cells or produces more than one effect, such as changing inflammation, blood-cell production, immune activity, or tissue repair.

\medterm{Pleiotropy} The ability of one gene change to influence several traits or body systems. For example, one inherited disorder may affect the brain, skin, blood, and other organs at the same time.

\medterm{Pleomorphic} Pleomorphic means having multiple forms or marked variation in cell or tissue appearance; in pathology, it may be a feature of malignancy but is not diagnostic alone.

\medterm{Pleomorphic Adenoma} A usually noncancerous tumor of a salivary gland, most often the parotid gland near the ear. It usually grows slowly as a painless lump, but long-standing tumors can rarely become cancerous.

\medterm{Pleomorphic Adenoma Of The Vulva} An exceptionally rare benign mixed epithelial and mesenchymal tumor of the vulva, resembling pleomorphic adenoma of salivary tissue and requiring histopathologic confirmation.

\medterm{Pleomorphic Dermal Sarcoma} An uncommon high-grade cutaneous sarcoma, usually on chronically sun-damaged skin of older adults, composed of markedly pleomorphic cells and capable of local recurrence or metastasis.

\medterm{Pleomorphic Xanthoastrocytoma} A circumscribed astrocytic tumor, often arising superficially in the cerebral hemispheres of younger patients and associated with seizures. Most are low grade, but anaplastic transformation can occur.

\medterm{Pleomorphism} Variation in the size and shape of cells or their nuclei within a tissue. Marked
pleomorphism is a feature of many high-grade malignancies but can also occur in some
reactive or degenerative conditions.

\medterm{Plethoric} Plethoric means unusually full, congested, or ruddy, often describing the appearance associated with increased blood volume or red-cell mass.

\medterm{Plethysmography} A test that measures changes in the size of a body part or the amount of blood or air it contains. It can assess circulation in the limbs, lung volume, or changes in blood flow during breathing or exercise.

\medterm{Pleura} A thin, two-layered lining surrounding each lung and covering the inside of the chest wall. A small amount of fluid between its layers allows the lungs to move smoothly during breathing.
\medterm{Pleural} Relating to the pleura, the membranes surrounding the lungs and lining the inside of the chest.

\medterm{Pleural Cavity} The narrow space between the two slippery layers surrounding each lung. A small amount of fluid normally allows smooth breathing, while extra air, blood, pus, or fluid can restrict lung expansion.

\medterm{Pleural Decortication} An operation that removes a thick scar-like layer from the lung’s covering so a trapped lung can expand again. It may be used for long-standing infected fluid, blood, or inflammation around the lung when other treatment is insufficient. Risks include bleeding, infection, air leakage, and breathing complications.

\synonyms
Decortication of the Lung; Pulmonary Decortication; Thoracoscopic Decortication

\medterm{Pleural Disorders} Conditions affecting the lung lining, including fluid buildup, trapped air, inflammation, infection, bleeding, and tumors. They may cause chest pain, cough, breathlessness, fever, or reduced lung expansion.

\medterm{Pleural Effusion} An abnormal buildup of fluid in the space between a lung and the chest wall. It may result from
heart failure, infection, inflammation, cancer, or injury, and can cause shortness of breath, cough, or chest pain.
\medterm{Pleural Effusion (Fluid In The Chest Or On Lung)} Alternate index label for Fluid In the Chest or On Lung.

\medterm{Pleural Fluid Analysis} Laboratory testing of fluid collected from around the lung. Cell counts, protein, acidity, glucose, cultures, and other tests help determine whether the fluid is caused by infection, cancer, heart failure, inflammation, or another problem.

\medterm{Pleural Fluid Culture} Culture of pleural fluid to detect bacteria, fungi, or mycobacteria in suspected pleural-space infection; specimens should be collected appropriately, often in blood-culture bottles.

\medterm{Pleural Fluid Glucose} Pleural-fluid glucose may be reduced in empyema, complicated parapneumonic effusion, rheumatoid
pleuritis, tuberculosis, and malignancy. It is interpreted with pH, cell count, protein, LDH,
microbiology, imaging, and the clinical setting.

\medterm{Pleural Fluid Gram Stain} A pleural-fluid Gram stain looks for bacteria in fluid around the lung and may support evaluation of empyema or bacterial pleural infection.

\medterm{Pleural Fluid PH} Pleural-fluid pH helps assess infection and other causes of effusion. A markedly low pH supports a
complicated parapneumonic effusion in the appropriate clinical setting, but drainage decisions
also depend on imaging, fluid appearance, microbiology, glucose, and the patient's condition.

\medterm{Pleural Fluid Smear} Microscopic examination of fluid removed from the space around the lung. It may show inflammatory cells, bacteria, fungi, cancer cells, or other material and is interpreted with fluid chemistry, cultures, imaging, and symptoms.

\medterm{Pleural Friction Rub} A scratchy or grating sound heard through a stethoscope when inflamed layers around the lung rub together during breathing. It may occur with infection, a blood clot in the lung, or other causes of pleural inflammation.

\medterm{Pleural Needle Biopsy} A pleural needle biopsy removes a small sample of the pleural lining through a needle for evaluation of cancer, infection, or inflammation.

\medterm{Pleural Neoplasm} A pleural neoplasm is an abnormal growth in the lining around the lung, such as mesothelioma or metastatic cancer.

\medterm{Pleural Mesothelioma} A cancer arising from the pleura, the thin lining around the lungs and inner chest wall. It is strongly associated with asbestos exposure and may cause chest pain, breathlessness, or fluid around the lung.

\medterm{Pleural Perforation} A breach in the pleural lining allowing air or fluid to enter the pleural space,
potentially causing pneumothorax or hydropneumothorax, often from trauma, procedures, or
underlying lung disease.

\medterm{Pleural Plaque} A usually harmless area of thickened or calcified tissue on the lining around the lung, most often caused by past asbestos exposure. It is not the same as asbestos-related lung scarring.

\medterm{Pleural Sac} The potential space and enclosing pleural membranes around a lung, normally containing a small amount of lubricating fluid.

\medterm{Pleural Space} The narrow potential space between the lung’s inner covering and the chest wall’s lining. It normally contains a small amount of lubricating fluid; air, blood, or excess fluid can impair breathing.

\medterm{Pleural Tuberculosis} Tuberculous infection or inflammation of the pleura, often producing a lymphocyte-rich pleural effusion, chest pain, fever, and breathlessness and requiring microbiologic or tissue evaluation.

\medterm{Pleurectomy} Surgical removal of part or all of the pleura, used in the treatment of malignant
pleural mesothelioma and recurrent malignant pleural effusion.

\medterm{Pleurisy} Inflammation of the thin layers around the lungs, often causing sharp chest pain that worsens with deep breathing or coughing. Causes include viral infection, pneumonia, a blood clot, autoimmune disease, or injury; associated features vary with the underlying cause.

\synonyms
Pleuritis (Pleurisy)


\medterm{Pleuritic Pain} Sharp chest pain that worsens with breathing or coughing, often caused by inflammation of the pleura or nearby tissues.

\medterm{Pleuritis} Alternate index label for Pleurisy.

\medterm{Pleuritis (Pleurisy)} Alternate index label for Pleurisy.

\medterm{Pleurodesis} A procedure that makes the two layers around a lung stick together so air or fluid cannot collect repeatedly. It may use a chemical or surgery and can cause pain, fever, infection, or breathing complications.

\medterm{Pleurodynia} Severe, repeated pain in the chest or upper abdomen caused by inflammation of breathing muscles, often during infection with coxsackievirus. Pain may worsen with breathing or movement and usually resolves as infection improves.
\medterm{Pleuropulmonary Blastoma} A rare aggressive childhood cancer arising in the lung or pleura, usually associated with DICER1-related tumor predisposition.

\medterm{Plexiform Fibrohistiocytic Tumor} A rare low-grade malignant fibroblastic and histiocytic tumor, usually arising in the dermis or subcutis of the upper extremities of young people; local recurrence is common and metastasis is uncommon.

\medterm{Plexiform Neurofibroma} A plexiform neurofibroma is a diffuse nerve-sheath tumor, strongly associated with neurofibromatosis type 1, that can grow along multiple nerve branches and sometimes become malignant.

\synonyms
PN
PNF

\medterm{Plexus} A network formed by intersecting nerves, blood vessels, lymphatic vessels, or other structures. Examples include the brachial nerve plexus and the celiac vascular plexus.
\medterm{Pliability} Pliability is the ability of a material or tissue to bend, stretch, or deform without breaking; it is relevant to skin, vessels, lungs, and medical devices.

\medterm{Plica} A fold of tissue, especially a fold in the lining of a joint. Plicae are often normal, but an irritated plica around the knee can cause pain, clicking, swelling, or a catching feeling with movement.
\medterm{Plitidepsin} Plitidepsin is an antineoplastic compound that targets eukaryotic elongation factor 1A and has been studied for selected cancers and viral infections.

\medterm{Ploidies} Ploidy describes the number of complete chromosome sets in a cell, such as haploid, diploid, or polyploid; abnormal ploidy can characterize tumors or developmental disorders.


\medterm{PLS} Alternate index label; see \textit{Primary lateral sclerosis}.

\medterm{Plug (Vascular)} A vascular occlusion device used to block a selected vessel under image guidance. A device placed in a blood vessel to block selected blood flow during image-guided treatment.
\synonyms
Vascular occlusion plug; embolization plug

\medterm{Plummer-Vinson Syndrome} A condition involving iron-deficiency anemia, difficulty swallowing, and thin webs in the upper swallowing tube. Treating the iron deficiency and stretching the web can improve swallowing and may reduce complications.
\medterm{Pluralibacter Gergoviae} A gram-negative environmental bacterium that can contaminate medicines or cosmetics and rarely cause opportunistic infection.

\medterm{Pluripotential} Able to develop into many, but not every, type of cell in the body. Pluripotent stem cells can produce cells from the three early tissue groups and are studied for development, disease, and regenerative treatment.
\medterm{Plutonium} A radioactive heavy metal whose inhalation or ingestion can cause internal radiation exposure and toxicity.
\medterm{Pm} An abbreviation that may mean afternoon, post-mortem, or picometer; meaning depends on clinical context.
\medterm{PMBCL} Alternate index label; see \textit{Primary mediastinal large B-cell lymphoma}.

\medterm{PMBL} Alternate index label; see \textit{Primary mediastinal large B-cell lymphoma}.

\medterm{PMDD} Premenstrual dysphoric disorder is a severe cyclic mood disorder occurring before menstruation with irritability, depression, anxiety, and functional impairment that improve after menses.

\medterm{PML} PML commonly means progressive multifocal leukoencephalopathy, a JC-virus demyelinating brain infection during severe immune suppression, but it may have other meanings.

\medterm{PN} Alternate index label; see \textit{Plexiform neurofibroma}.

\medterm{Pneumatic} Relating to air or gas under pressure. Medical examples include an air-pressure ear examination, inflatable compression sleeves used to reduce clot risk, and an inflatable tourniquet used during some operations.
\medterm{Pneumatic Otoscopy} Examination of the eardrum with a lighted instrument while gently changing the air pressure in the ear canal. Reduced eardrum movement may suggest fluid behind it or poor middle-ear ventilation.

\medterm{Pneumatic Retinopexy} A procedure for selected retinal detachments in which an intravitreal gas bubble is injected and laser or cryotherapy seals the retinal break while positioning helps the bubble support the repair.

\medterm{Pneumatic Tourniquet} An inflatable cuff used to temporarily reduce blood flow to a limb during surgery or selected procedures.
\medterm{Pneumatocele} A thin-walled, air-filled space in the lung that often develops after infection, inflammation, or injury. Many resolve naturally, but large or infected spaces can impair breathing or require drainage or surgery.
\medterm{Pneumatosis} Gas trapped within the wall of an organ or other tissue, most often the bowel. It may be harmless, but when caused by poor blood flow, infection, or a tear, it can signal a life-threatening abdominal emergency.
\medterm{Pneumatosis Cystoides Coli} Multiple gas-filled cysts within the wall of the large intestine. They may be harmless and found by chance, or may signal bowel blockage, poor blood flow, infection, or another serious illness when pain or other warning signs occur.

\medterm{Pneumatosis Cystoides Intestinalis} Gas-filled cysts within the bowel wall. They may be harmless and found incidentally, but can also result from poor blood flow, infection, blockage, lung disease, or medicines. Severe pain, fever, or shock requires urgent evaluation.

\medterm{Pneumatosis Intestinalis} Gas within the wall of the intestine, seen as small bubbles or lines on an imaging test. It may be harmless, but pain, fever, shock, or signs of poor blood flow require urgent evaluation for bowel injury.

\medterm{Pneumocephalus} Air or another gas inside the skull, usually after a head injury, brain surgery, or a skull fracture near the sinuses. A large or trapped collection can press on the brain and cause headache, vomiting, confusion, or reduced alertness.

\medterm{Pneumococcal Infection} Infection caused by Streptococcus pneumoniae, ranging from ear or sinus infection to pneumonia, meningitis, or bloodstream infection.

\medterm{Pneumococcal Meningitis} A dangerous infection of the brain and spinal-cord coverings caused by Streptococcus pneumoniae. It can cause fever, headache, stiff neck, confusion, seizures, hearing loss, or death and requires immediate treatment.

\medterm{Pneumococcal Pneumonia} An acute bacterial infection of the lung parenchyma caused by Streptococcus pneumoniae, the single most common cause of community-acquired pneumonia globally. It classically presents with sudden high fever, shaking chills (rigors), pleuritic chest pain, productive cough with rust-colored sputum, and lobar consolidation on chest imaging.

\synonyms
Streptococcus Pneumoniae Pneumonia; Pneumococcal Lung Infection

\medterm{Pneumococcal Sepsis} Sepsis caused by Streptococcus pneumoniae, often arising from pneumonia, meningitis, or another invasive pneumococcal infection.
\medterm{Pneumoconiosis} A group of work-related lung diseases caused by breathing mineral or industrial dust over time. Dust can scar lung tissue and cause cough, breathlessness, reduced oxygen, and progressive breathing disability.
\medterm{Pneumocystis} A genus of fungi that includes Pneumocystis jirovecii, which can cause pneumonia in people with weakened cellular immunity.

\medterm{Pneumocystis Carinii} An older name associated with Pneumocystis; in humans, the usual species causing pneumonia is Pneumocystis jirovecii.

\medterm{Pneumocystis Carinii Pneumonia} The historical name for Pneumocystis jirovecii pneumonia, an opportunistic lung infection in immunocompromised people.
\medterm{Pneumocystis Infection} Infection by Pneumocystis jirovecii, most often causing pneumonia with dry cough, fever, and progressive shortness of breath in immunocompromised people.

\medterm{Pneumocystis Jiroveci} A fungus that can cause a serious pneumonia when the immune system is weakened, especially by advanced HIV, cancer treatment, or transplantation. It often causes gradually worsening dry cough, fever, breathlessness, and low blood oxygen.

\medterm{Pneumocystis Jirovecii Pneumonia} Opportunistic pneumonia caused by Pneumocystis jirovecii, especially in people with
advanced HIV infection or other cellular immunodeficiency. It causes progressive
exertional dyspnea, dry cough, fever, and hypoxemia.

\medterm{Pneumocystis Pneumonia} Opportunistic pneumonia caused by Pneumocystis jirovecii, usually in people with impaired
cell-mediated immunity. It may cause progressive exertional dyspnea, dry cough, fever,
and hypoxemia with diffuse ground-glass changes. Diagnosis and treatment require clinical,
microbiologic, and radiologic assessment; trimethoprim-sulfamethoxazole is commonly used.

\medterm{Pneumocyte} A cell lining the air sacs of the lungs. Type I cells allow oxygen and carbon dioxide to pass between air and blood, while type II cells make surfactant, which helps keep the air sacs open.

\medterm{Pneumolysis} Surgical separation or freeing of lung tissue or pleural adhesions.
\medterm{Pneumomediastinum} Air trapped in the central space of the chest between the lungs. It may follow severe coughing, an injury, a medical procedure, or a tear in the airway or swallowing tube and can cause chest pain or breathlessness.
\medterm{Pneumonectomy} Surgery to remove an entire lung. It is usually performed for selected lung cancers or severe one-sided lung disease when removing only part of the lung is not suitable.
\medterm{Pneumonia} Infection of the lung tissue that fills air sacs with fluid or inflammatory material. It can cause cough, fever, chills, chest pain, breathlessness, low oxygen, or confusion and may be caused by bacteria, viruses, fungi, or aspiration.
\medterm{Pneumonia - Weakened Immune System} Lung infection in a person whose immune defenses are weakened by illness or treatment. Bacteria, viruses, fungi, and unusual organisms can cause it; symptoms may be subtle but disease can worsen quickly and needs prompt assessment.




\medterm{Pneumonia, Bilateral} Pneumonia affecting both lungs. Because more lung tissue may be involved, it can cause greater breathing difficulty and lower oxygen levels than disease limited to one lung, although severity depends on the cause and person’s health.

\medterm{Pneumonitis} Inflammation of lung tissue, often caused by a medicine, radiation, an immune reaction, inhaled irritants, or material entering the lungs rather than by infection. It may cause dry cough and breathlessness.
\medterm{Pneumopericardium} Air around the heart, usually after a chest injury, surgery, a medical procedure, or a tear in the esophagus. A large or trapped collection can prevent the heart from filling and pumping properly and requires urgent assessment.

\medterm{Pneumoperitoneum} Air or gas inside the abdominal cavity around the organs. It may be expected after abdominal surgery, but unexpected gas can indicate a torn intestine and requires urgent assessment, especially with pain or illness.
\medterm{Pneumoplasty} Surgery to repair or reshape part of the lung or an airway. It may be used to remove a damaged area, improve airflow, or restore the shape of tissue after injury or disease.

\medterm{Pneumotachometer} An instrument that measures the speed and amount of air moving during breathing. It can help assess airflow limitation and is used in lung-function testing, often together with other measurements.

\medterm{Pneumothorax} Air trapped between the lung and the chest wall, causing part or all of the lung to collapse. It may happen spontaneously, after injury, or after a procedure and can cause sudden chest pain and breathlessness; severe symptoms are an emergency.
\medterm{Pneumothorax - Infants} Air trapped between an infant’s lung and chest wall, which can compress the lung and cause rapid breathing, low oxygen, poor feeding, or circulatory instability. Severe cases require immediate drainage and breathing support.

\medterm{Pneumovirinae} Pneumovirinae is a subfamily of viruses that includes respiratory pathogens such as respiratory syncytial virus and human metapneumovirus.

\medterm{Pneumovirus} A group or genus of respiratory viruses; human pneumovirus terminology often refers to human metapneumovirus or related paramyxoviruses.

\medterm{PNF} Alternate index label; see \textit{Plexiform neurofibroma}.

\medterm{Pocket Erosion} Loss of tissue at the margin or base of a periodontal or surgical pocket, potentially exposing deeper structures and promoting infection, bleeding, or impaired healing.

\medterm{Podagra} A sudden gout attack in the big-toe joint, causing intense pain, redness, warmth, and swelling.
\medterm{Podiatrist} A healthcare professional specializing in diseases and injuries of the feet, ankles, and related lower limbs. Care may include examination, wound treatment, medicines, footwear advice, procedures, and prevention of complications.
\medterm{Podiatry} The healthcare field focused on the feet, ankles, and related lower-limb conditions. It includes prevention, diagnosis, treatment, wound care, mobility support, and management of problems linked with diabetes or circulation.
\medterm{Podocyte} A specialized cell covering the tiny blood vessels that filter blood in the kidney. Its delicate extensions help keep proteins in the bloodstream; damage can cause protein to leak into urine.

\medterm{Podofilox} Podofilox is a topical antimitotic medicine used to treat selected external anogenital warts caused by human papillomavirus; irritation and correct application are important.

\medterm{Podophyllin} Podophyllin is a plant-derived resin with antimitotic activity formerly used topically for warts; it can be severely toxic if absorbed and is generally replaced by safer products.

\medterm{Podophyllum} A plant genus containing podophyllotoxin, used in some topical wart treatments but toxic if misused.
\medterm{Podosome} A temporary attachment structure on the surface of some cells, including immune cells and bone-removing cells. It helps the cell grip tissue and break down nearby material while moving or carrying out its function.

\medterm{POEMS Syndrome} A rare disorder linked to abnormal plasma cells that affects several body systems. It may cause nerve damage, enlarged organs, hormone problems, skin changes, fluid buildup, eye swelling, and unusual bone changes.

\medterm{Poikilocytosis} The presence of red blood cells with unusually varied shapes on a blood smear. It may occur with anemia, inherited blood disorders, or damage to red blood cells.
\medterm{Poikiloderma} A skin pattern combining mottled pigmentation, visible small vessels, and epidermal thinning, seen in connective-tissue disease, chronic sun damage, inflammation, and some cutaneous lymphomas.

\medterm{Poikilotherm} An organism whose body temperature changes with surroundings, unlike humans, who normally maintain a relatively stable temperature.
\medterm{Poinsettia Plant Exposure} Chewing poinsettia usually causes mild mouth irritation, nausea, or stomach upset, not severe poisoning. Rinse the mouth and contact poison control if symptoms are severe, a large amount was swallowed, or another product was involved.

\medterm{Point} A specific location, spot, or value; in medicine it may refer to a tender point, pressure point, or measurement point.
\medterm{Point Mutation} A change affecting one DNA building block. Depending on its location, it may have no effect, change the protein made by a gene, stop protein production early, or contribute to an inherited condition or cancer.

\medterm{Point Tenderness - Abdomen} Abdominal point tenderness is sharply localized pain when a particular abdominal site is pressed, which may indicate inflammation, infection, obstruction, or another focal process.

\medterm{Point-Of-Care Testing} A test performed near the patient, such as in a clinic, ambulance, home, or hospital ward, instead of a central laboratory. Results are available quickly, but trained users, quality checks, and clinical interpretation remain essential.

\medterm{Point-Of-Care Ultrasound} Ultrasound performed and interpreted at the bedside by the treating clinician to answer a focused question or guide a procedure. It can assess the heart, lungs, abdominal fluid, blood vessels, and bladder, but does not replace a complete imaging study when one is needed.

\medterm{Poison} A substance that can cause injury, illness, or death when enough enters the body. Harm depends on the substance, amount, route, timing, and person's health; suspected exposure needs immediate poison-control advice.
\medterm{Poison Control} A specialist service that gives immediate advice after exposure to a medicine, chemical, plant, or other possible poison. They assess the substance and amount and advise on first aid, observation, emergency care, or further treatment.

\medterm{Poison Control Center - Emergency Number} A poison-control center provides immediate expert advice after a possible poisoning or toxic exposure; the correct emergency number depends on the country and life-threatening symptoms require emergency services.

\medterm{Poison Ivy - Oak - Sumac} Poison ivy, oak, and sumac cause allergic contact dermatitis through urushiol oil, producing intensely itchy red patches and blisters after skin exposure.

\medterm{Poison Ivy Rash} An itchy allergic contact dermatitis caused by exposure to urushiol, an oil found in poison ivy and related plants. The rash commonly appears as red streaks or patches with swelling, small bumps, or blisters; the pattern and severity vary with the amount and route of exposure.

\medterm{Poison Ivy - Oak - Sumac Rash} Alternate index label for Poison Ivy Rash.

\medterm{Poison Ivy, Oak, And Sumac} Alternate index label for Poison Ivy Rash.

\synonyms
Allergy Plant Contact (Poison Ivy, Oak, And Sumac)

\medterm{Poison-Center Consultation} Specialist toxicology advice used to identify hazards, interpret exposure and laboratory
data, select antidotes, guide observation or disposition, and coordinate public-health
response. Consultation is appropriate for significant, uncertain, mixed, pediatric, or
unusual exposures.

\medterm{Poisoning} Harmful effects caused by swallowing, inhaling, injecting, or absorbing a toxic substance. Suspected poisoning requires urgent advice from emergency services or a poison center and should not be treated by inducing vomiting.

\medterm{Polytrauma} Serious injuries affecting two or more body regions or organ systems, usually caused by a major accident, fall, blast, or violence. The injuries may include bleeding, fractures, brain injury, chest injury, or damage to internal organs.
\medterm{Poisoning - Fish And Shellfish} Toxin-related illness after eating contaminated fish or shellfish. Vomiting, diarrhea, flushing, tingling, weakness, or breathing difficulty may occur; sudden neurologic or breathing symptoms require emergency care.

\medterm{Poisoning Arsenic (Arsenic Poisoning)} Alternate index label for Arsenic Poisoning.

\medterm{Poisoning Ciguatera (Ciguatera Poisoning)} Alternate index label for Ciguatera Poisoning.

\medterm{Poisoning First Aid} Move away from fumes, rinse chemicals from skin or eyes, and call poison control with the product container available. Do not induce vomiting or give food unless advised; breathing problems, seizures, collapse, or inability to wake require emergency services.

\medterm{Poisoning Thallium (Thallium)} Alternate index label for Thallium.

\medterm{Pokeweed Poisoning} Eating pokeweed, especially its roots or unripe parts, can cause burning in the mouth, vomiting, diarrhea, weakness, low blood pressure, or seizures. Significant ingestion requires immediate poison-control or emergency advice.

\medterm{Pol Gene} The HIV gene encoding enzymes required for viral replication, including protease, reverse transcriptase, and integrase; antiretroviral drugs target these products.

\medterm{Poland Syndrome} A congenital condition characterized by absence or underdevelopment of part of the pectoral
muscles on one side, sometimes accompanied by abnormalities of the hand on the same side.

\medterm{Polar} Relating to a pole or direction, or describing a molecule with slightly different electrical charges at its ends. In biology and chemistry, polarity affects how substances dissolve, move across cell membranes, and interact with other molecules.
\medterm{Polarity} The arrangement or direction of positive and negative electrical, chemical, or biological properties. Cell polarity helps organize tissue structure and directs processes such as movement, division, absorption, and secretion.
\medterm{Polarography} Polarography is an electrochemical analytical method that measures current while voltage changes at an electrode to identify or quantify reducible substances.

\medterm{Polidocanol} Polidocanol is a detergent sclerosant injected to close selected varicose veins or vascular lesions by damaging the vessel lining and promoting fibrosis.

\medterm{Polio} An infection caused by poliovirus, also called poliomyelitis. It spreads mainly when virus from an infected person’s stool enters the mouth, and it can also spread through respiratory secretions. Most infected people have no symptoms or have a short, mild illness with fever, sore throat, tiredness, nausea, headache, or stomach pain. Fewer than 1\% of infections cause the virus to affect the brain or spinal cord, leading to meningitis, muscle weakness, or acute flaccid paralysis, which is often uneven between the limbs and may become permanent. Paralysis of the muscles used for breathing can be fatal. Vaccination prevents polio.



\medterm{Poliomyelitis} Alternate index label; see \textit{Polio}.
\medterm{Poliosis} A localized patch of white or depigmented hair caused by reduced melanin in follicles, occurring in vitiligo, alopecia areata, inflammatory disease, genetic syndromes, or after injury.

\medterm{Poliovirus} An enterovirus that causes poliomyelitis. It is transmitted mainly by the fecal-oral route; most infections are asymptomatic, but a small proportion cause meningitis or acute flaccid paralysis by damaging motor neurons in the spinal cord. Diagnosis uses molecular testing, usually reverse-transcription PCR of stool or other appropriate specimens.

\medterm{Polish (Dental)} Smoothing restorative surfaces to reduce plaque accumulation. Essential for longevity.

\synonyms
Dental polishing

\medterm{Pollakisuria} Pollakisuria is unusually frequent urination, often in small amounts, caused by irritation, infection, stones, bladder disorders, anxiety, or other urinary conditions.

\medterm{Pollakiuria} Frequent urination, often in small amounts. It may occur with urinary infection, irritation, diabetes, pregnancy, medicines, or anxiety; in some children it is a temporary benign pattern, but pain, fever, blood, thirst, or weight loss needs assessment.
\medterm{Pollen} Tiny plant grains released for reproduction that can trigger allergies, asthma symptoms, or eye inflammation.

\medterm{Pollen-Food Allergy Syndrome} An allergic reaction to certain raw fruits, vegetables, or nuts in people sensitized to related pollen proteins. It usually causes itching or tingling of the mouth and throat and is also called oral allergy syndrome.

\medterm{Pollex} The anatomical term for the thumb, the first and most mobile digit of the hand. It has two bones and can touch the fingertips, allowing precise grasping; injury or arthritis can cause pain and loss of hand function.
\medterm{Pollution} Contamination of air, water, soil, or surroundings by substances that may harm health or ecosystems.
\medterm{Poly A} A sequence of adenine nucleotides added to the end of many messenger RNA molecules, helping regulate stability, transport, and translation.

\medterm{Polyadenylation} Addition of a polyadenosine tail to the 3$^{\prime}$ end of many eukaryotic messenger RNAs, promoting stability, nuclear export, and efficient translation.

\medterm{Polyamine} A small positively charged molecule, such as putrescine, spermidine, or spermine, required for cell growth and nucleic-acid function.

\medterm{Polyamine Catabolism} The enzyme-controlled breakdown of polyamines, natural substances involved in cell growth and function. This process produces chemicals that can affect cell signals, oxidative stress, inflammation, development, and cancer biology.

\medterm{Polyarteritis} Inflammation of several arteries, which can reduce blood flow and damage organs. Causes and effects vary; polyarteritis nodosa is a rare form that mainly affects medium-sized arteries and may cause fever, muscle pain, skin changes, nerve problems, or abdominal pain.
\medterm{Polyarteritis Nodosa} A rare disease in which inflammation damages medium-sized arteries. It can reduce blood flow to the skin, nerves, kidneys, muscles, or intestines, causing fever, weight loss, nerve pain or weakness, skin changes, high blood pressure, or abdominal pain.

\synonyms
PAN (Polyarteritis Nodosa)

\medterm{Polyarthritis} Inflammation involving multiple joints, commonly defined as five or more. It is a clinical pattern rather than a single diagnosis and may result from autoimmune, infectious, crystal-related, or other inflammatory disorders.
\medterm{Polybrominated Biphenyl} A group of brominated industrial chemicals that can persist in the environment and cause toxic effects after exposure.

\medterm{Polychlorinated Biphenyl} A group of persistent industrial chemicals that can accumulate in body fat and are associated with cancer, endocrine, immune, and neurologic effects.

\medterm{Polychondritis} A relapsing immune-related disease that inflames and damages cartilage. It may affect the ears, nose, joints, voice box, windpipe, or inner ear, causing pain, deformity, breathing problems, hearing loss, or balance symptoms.

\medterm{Polychromasia} The presence of unusually bluish, immature red blood cells on a blood smear. It usually means bone marrow is releasing cells quickly, often because of blood loss, red-cell destruction, or treatment response.
\medterm{Polychromatophilia} Polychromatophilia is bluish-gray staining of red blood cells on a blood smear, reflecting increased circulating immature reticulocytes from stimulated red-cell production.

\medterm{Polycystic Kidney Disease} An inherited disorder in which numerous kidney cysts enlarge the kidneys and may cause hypertension and renal failure.

\synonyms
PKD

\medterm{Polycystic Liver Disease} An inherited condition in which many fluid-filled cysts develop in the liver. The cysts may enlarge the liver and cause abdominal fullness or pain; liver function is often preserved, but kidney cysts may also occur.

\medterm{Polycystic Ovarian Syndrome} A hormonal and metabolic disorder involving ovulatory dysfunction and androgen excess, often causing irregular periods, acne, excess hair, infertility, or metabolic risk.

\medterm{Polycystic Ovary} Polycystic ovaries contain many small follicles and may occur as a normal variant or with polycystic ovary syndrome. PCOS additionally involves ovulatory dysfunction, androgen excess, irregular periods, infertility, acne, or metabolic risk and requires clinical—not ultrasound alone—diagnosis.

\synonyms
PCOS (Polycystic Ovary)

\medterm{Polycystic Ovary Syndrome} A common hormonal and metabolic condition that affects ovulation and the levels or effects of certain hormones. It may cause irregular or absent periods, reduced or absent ovulation, acne, extra facial or body hair, scalp hair thinning, or difficulty becoming pregnant. It may also cause higher levels or effects of androgens, hormones found in everyone but usually present at higher levels in males. Some people have many small follicles in the ovaries, but this finding is not required for the condition. Insulin resistance, weight gain, darkened thickened skin, type 2 diabetes, and other metabolic problems may also occur. In adults, the diagnosis is based on at least two of three features---ovulation problems, androgen excess, or polycystic-appearing ovaries---after other causes have been considered.

\medterm{Polycythaemia} An increased red-cell mass or concentration, which may be primary or secondary to another condition.
\medterm{Polycythemia} An abnormally high red-cell concentration or total red-cell amount. It may result from a bone-marrow disorder, low oxygen, dehydration, smoking, medicines, or another illness and can increase the risk of blood clots.

\medterm{Polycythemia (High Red Blood Cell Count)} Alternate index label for High Red Blood Cell Count.

\medterm{Polycythemia - Newborn} Neonatal polycythemia is an abnormally high red-cell concentration in a newborn that can increase blood viscosity and cause poor feeding, lethargy, breathing problems, or low glucose.

\medterm{Polycythemia Vera} A long-term bone-marrow cancer that produces too many red blood cells and sometimes too many white cells and platelets. The thicker blood can cause headaches, itching, clots, bleeding, or an enlarged spleen.

\synonyms
Primary Polycythemia

\medterm{Polydactylism (Polydactyly)} A birth condition in which a person has one or more extra fingers or toes. The extra digit may be on the thumb or big-toe side, the little-finger or little-toe side, or between the usual digits. It may occur alone or as part of a genetic syndrome. Treatment depends on its size and function and may involve reconstructive surgery.

\medterm{Polydactyly} A birth condition in which a person has one or more extra fingers or toes. The extra digit may be small or fully formed and may occur alone or with a genetic syndrome; treatment depends on its structure, function, and effect on nearby digits.
\medterm{Polydipsia} Excessive thirst that leads to drinking unusually large amounts of fluid. Common causes include uncontrolled diabetes, water-balance disorders, medicines, and some mental health conditions.
\medterm{Polyene} A molecule containing multiple carbon-carbon double bonds; polyene medicines such as amphotericin B act against fungal cell membranes.

\medterm{Polyethylene} A plastic polymer used in packaging and medical devices; implanted polyethylene components may wear over time and produce foreign-body reactions.

\medterm{Polyethylene Glycol} A polymer used in medicines, especially laxatives that draw water into stool to ease constipation.
\medterm{Polygenic} Influenced by variants in many genes, often together with environmental factors. Height, blood pressure, and risk for many common diseases are polygenic traits.
\medterm{Polygenic Disease} A disease influenced by small effects from many genes together with environment and lifestyle. Examples include common diabetes, high blood pressure, heart disease, and some cancers; risk is inherited, not certainty.

\medterm{Polygenic Inheritance} Inheritance in which many genes, each with a small effect, together influence a trait or disease risk. Height, skin color, blood pressure, diabetes risk, and heart-disease risk are examples; environment also contributes.

\medterm{Polygenic Obesity} Obesity influenced by many small genetic differences together with eating patterns, activity, sleep, medicines, and the surrounding environment.

\medterm{Polygenic Risk Score} An estimate of a person’s inherited tendency toward a disease, calculated from many small genetic differences. It changes risk assessment but does not predict certainty and must be considered with family history, environment, and other health factors.

\medterm{Polyglandular Autoimmune Syndrome} An autoimmune disorder in which two or more endocrine glands are damaged, causing combinations of thyroid, adrenal, parathyroid, pancreatic, or gonadal dysfunction.

\medterm{Polygraphy} Polygraphy is the simultaneous recording of several physiologic signals, such as breathing, heart rate, oxygen level, and movement, often during sleep testing.

\medterm{Polyhydramnios} An unusually large amount of amniotic fluid during pregnancy. It may cause abdominal discomfort, breathlessness, or swelling and can increase the risks of early birth, abnormal fetal position, cord problems, or heavy bleeding after delivery.

\medterm{Polyketide Synthase} An enzyme system that builds complex natural products from repeated carbon units, producing many antibiotics, toxins, and other biologically active molecules.

\medterm{Polymagia (Polyphagia)} An unusually strong or excessive appetite. It may occur with diabetes, an overactive thyroid, certain medicines, genetic disorders, or some eating and mental health conditions.
\medterm{Polymastia} The presence of extra breast tissue, sometimes with an additional nipple, along the embryonic milk line. Accessory breast tissue can enlarge or develop the same benign or malignant conditions as normal breast tissue.
\medterm{Polymenorrhea} Menstrual periods that occur unusually often, commonly at intervals shorter than about 21 days. Causes include hormonal changes, thyroid disease, fibroids, bleeding disorders, medicines, and the menopausal transition; persistent frequent bleeding needs evaluation.

\medterm{Polymenorrhoea} British spelling of polymenorrhea: menstrual periods occurring unusually often, commonly at intervals shorter than about 21 days.
\medterm{Polymer} A large molecule made of repeating smaller units called monomers; polymers are widely used in biology, medicines, and medical materials.

\medterm{Polymerase Chain Reaction} A laboratory method that makes millions of copies of a selected piece of DNA so it can be detected and studied. It is used to identify infections, inherited conditions, cancer-related gene changes, and other genetic material. The test must be chosen and interpreted for the specific question.
\medterm{Polymerization} Chemical joining of many small monomers into a larger polymer, a process important in biomolecules, plastics, dental materials, drug delivery, and laboratory techniques.

\medterm{Polymicrobial Infections} Infections caused by two or more kinds of germs at the same time. They commonly affect contaminated wounds, the abdomen, or the lungs and may require treatment directed at several different organisms.

\medterm{Polymicrogyria} A brain-development abnormality in which the outer brain surface has too many small, unusually formed folds. It may cause seizures, developmental delay, movement problems, or learning difficulties, depending on the areas involved.

\medterm{Polymorphic} Having more than one form or appearance. The term may describe a rash, crystal, cell, gene variant, heart rhythm, or tumor. Its medical importance depends on the structure being described and the other clinical findings.

\medterm{Polymorphism} The presence of two or more inherited genetic forms or visible traits within a population.
\medterm{Polymorphonuclear Leukocyte} A white blood cell with a multilobed nucleus, especially a neutrophil, eosinophil, or basophil, involved in innate immune defense and inflammation.

\medterm{Polymorphous} Having lesions or clinical features that differ in shape, size, or stage and therefore show multiple forms.

\medterm{Polymorphous Light Eruption} Polymorphous light eruption is an itchy inflammatory rash triggered by sunlight, usually appearing hours to days after exposure and recurring seasonally in susceptible people.

\synonyms
Sun Poisoning

\medterm{Polymyalgia Rheumatica} An inflammatory disorder causing aching and morning stiffness around the shoulders and hips in older adults.

\synonyms
PMR (Polymyalgia Rheumatica)

\medterm{Polymyositis} An uncommon immune-related disease causing gradual, usually symmetrical weakness in muscles close to the body, especially the shoulders, hips, and thighs. It can make climbing stairs, lifting, or swallowing difficult and may raise muscle-enzyme levels.
\synonyms
PM (Polymyositis)

\medterm{Polymyxin} A group of antibiotics that disrupt bacterial cell membranes, used mainly for selected gram-negative infections.
\medterm{Polymyxin B} An antibiotic active against selected gram-negative bacteria; kidney and nerve toxicity limit its use to difficult infections.

\medterm{Polyneuritis} Inflammation or injury affecting multiple peripheral nerves, causing symmetric pain, numbness, weakness, or autonomic dysfunction; causes include immune disease, infection, toxins, and deficiency.

\medterm{Polyneuritis, Acute Idiopathic} Sudden inflammation or dysfunction of several peripheral nerves without an identified cause. It may cause tingling, weakness, pain, or loss of reflexes in the limbs and can sometimes affect breathing or swallowing. Rapidly worsening weakness requires urgent assessment because an immune nerve disorder such as Guillain--Barré syndrome may be responsible.

\medterm{Polyneuropathy} A generalized disorder affecting multiple peripheral nerves, often causing symmetric numbness, tingling, burning pain, weakness, or impaired reflexes in the limbs.

\medterm{Polynucleotide} A chain of linked nucleotides forming DNA, RNA, or a related laboratory molecule. Its sequence stores or carries genetic information and determines how cells make proteins or regulate biological processes.

\medterm{Polyoma Virus} A polyomavirus; some human polyomaviruses are widespread and cause serious disease mainly in immunocompromised people.
\medterm{Polyomavirus Hominis 2} A human polyomavirus designation associated with a virus found in people; clinically important disease is uncommon and mainly associated with immune suppression.

\medterm{Polyp} A growth projecting from a mucous membrane; it may be benign, precancerous, or rarely cancerous.
\medterm{Polyp Biopsy} A polyp biopsy removes part or all of a polyp for microscopic examination to determine whether it is benign, precancerous, or cancerous.

\medterm{Polyp Surveillance Colonoscopy} A follow-up colon examination after polyps have been found or removed. The timing depends on the number, size, type, and microscopic features of the polyps and on the person’s overall risk.

\medterm{Polypeptide} A chain of amino acids joined by peptide bonds. One or more folded polypeptide chains make up many proteins.
\medterm{Polyphagia} Unusually strong or persistent hunger that may lead to eating much more than usual. Causes include poorly controlled diabetes, an overactive thyroid, certain medicines, inherited conditions, and some mental-health disorders.

\medterm{Polypharmacy} Use of several medicines at the same time, often five or more. Some combinations are necessary, but taking many medicines increases the risk of interactions, side effects, confusion, missed doses, and prescribing errors, especially in older adults.



\medterm{Polyposis} A condition in which many polyps develop, especially in the bowel or other digestive organs. Polyposis may be inherited or acquired, and some types increase the risk of colorectal or other cancers.
\medterm{Polypropylene} A synthetic plastic polymer used in medical supplies, sutures, meshes, and containers; implanted materials may rarely cause infection or foreign-body reactions.

\medterm{Polyps Colon (Colon Polyps)} Alternate index label for Colon Polyps.

\medterm{Polyps Rectal (Colon Polyps)} Alternate index label for Colon Polyps.

\medterm{Polyps Uterus (Uterine Fibroids)} Alternate index label for Uterine Fibroids.

\medterm{Polyps, Colon} Alternate index label; see \textit{Colon polyps}.

\medterm{Polyps, Endometrial} Alternate index label; see \textit{Uterine polyps}.

\medterm{Polyps, Nasal} Alternate index label; see \textit{Nasal polyps}.

\medterm{Polyps, Stomach} Alternate index label; see \textit{Stomach polyps}.

\medterm{Polyps, Uterine} Alternate index label; see \textit{Uterine polyps}.

\medterm{Polyradiculitis} Inflammation or other disease affecting several spinal nerve roots. It may cause weakness, numbness, pain, and reduced reflexes; causes include infections, autoimmune disorders, and toxins.
\medterm{Polyradiculopathy} Disease affecting multiple spinal nerve roots, producing pain, weakness, sensory loss, or reflex changes in a distribution determined by the involved roots.

\medterm{Polysaccharide} A carbohydrate made of many linked simple sugar units. Examples include starch, glycogen, cellulose, and some bacterial capsules.
\medterm{Polysome} A group of ribosomes attached to one messenger RNA molecule at the same time. Together they make many copies of the protein instructions, helping cells produce proteins efficiently.

\medterm{Polysomnography} An overnight sleep study that records brain activity, eye movements, muscle activity, heart rhythm, airflow, breathing effort, and blood oxygen. It helps identify sleep-related breathing problems, unusual movements, seizures, and other sleep disorders.

\medterm{Polysomnography (Sleep Study)} An overnight recording of sleep, breathing, body movement, heart rhythm, and blood oxygen. It helps diagnose sleep apnea and other sleep disorders by showing how breathing and body functions change during sleep.

\synonyms
Sleep Study

\medterm{Polysplenia} A congenital condition with multiple small spleens and variable abnormalities of organ arrangement, heart, blood vessels, or intestinal rotation.

\medterm{Polystyrene} A synthetic plastic polymer used in containers and laboratory materials; burning it can release hazardous fumes.

\medterm{Polytene Chromosome} An enlarged chromosome formed by repeated DNA replication without cell division, displaying visible bands and commonly used to study gene activity in insect tissues.

\medterm{Polythelia} The presence of one or more additional nipples, usually along the embryonic milk line. It is often harmless, although accessory tissue can rarely have the same disorders as normal breast tissue.
\medterm{Polythiazide} A thiazide diuretic formerly used to treat high blood pressure and fluid retention by increasing renal salt and water excretion.

\medterm{Polyunsaturated Fat} A type of fat usually liquid at room temperature, found in fish, vegetable oils, nuts, and seeds. Some are essential because the body cannot make them.

\medterm{Polyurethane} A family of synthetic polymers used in medical tubing, implants, dressings, and other products; tissue response depends on formulation and site.

\medterm{Polyuria} Producing unusually large amounts of urine, commonly more than about three liters daily in an adult. Causes include diabetes, excess fluids, medicines, kidney problems, hormone disorders, and some inherited conditions.

\medterm{Pompe Disease} A rare inherited disorder in which a missing enzyme causes stored sugar to build up inside muscle cells. Babies may develop severe weakness, an enlarged heart, and breathing or feeding problems; later-onset disease mainly weakens the limbs and breathing muscles. Enzyme replacement and supportive care can slow or improve the disease.

\synonyms
Glycogen storage disease type II; Acid maltase deficiency; GSD II

\medterm{Polyvalent} Effective against or involving multiple types, strains, antigens, or targets.
\medterm{Polyvinyl} A family of polymers used in plastics and medical materials; properties vary according to attached chemical groups.

\medterm{Pomalidomide} An immunomodulatory medicine used with other treatments for selected multiple myeloma and related blood cancers; it can cause severe birth defects and blood abnormalities.

\medterm{Pompholyx} A form of dyshidrotic eczema causing intensely itchy small blisters on the hands or feet.
\medterm{Pompholyx Eczema} Pompholyx, or dyshidrotic eczema, causes intensely itchy small blisters on the palms, sides of fingers, or soles and may leave painful peeling or cracks.

\medterm{Pons} The middle part of the brainstem, containing pathways and nuclei involved in movement, sensation, breathing, and cranial-nerve function.
\medterm{Pontiac Fever} A mild, self-limited flu-like illness caused by Legionella bacteria without the pneumonia typical of Legionnaires’ disease.

\medterm{Pontine} Relating to the pons, a part of the brainstem between the midbrain and medulla. The pons helps carry movement and feeling signals and contains centers involved in eye movement, facial movement, hearing, balance, and breathing. A stroke or bleed there can cause weakness, loss of sensation, trouble speaking or swallowing, coma, or a locked-in state. Rapid correction of very low sodium can also damage it.
\medterm{Poor Circulation} Reduced or abnormal blood flow, usually referring to arterial insufficiency or venous disease, causing coolness, color change, pain, swelling, ulcers, or delayed wound healing.

\medterm{Poor Color Vision} Alternate index label; see \textit{Color blindness}.

\medterm{Popcorn Lung Bronchiolitis Obliterans} Bronchiolitis obliterans is irreversible scarring and narrowing of small airways that can follow toxic inhalation, infection, transplantation, or autoimmune disease. It causes persistent cough and breathlessness with obstructive lung testing; exposure avoidance and specialist care are important.

\medterm{Popliteal} Relating to the hollow behind the knee. This area contains the popliteal artery, vein, and nerves and can be affected by an artery aneurysm, a blood clot, or a Baker cyst.
\medterm{Popliteal Aneurysm} A weakened, bulging section of the artery behind the knee. It may form blood clots, block blood flow, or send clots into the lower leg, causing pain, coldness, numbness, or tissue damage. Larger aneurysms may need repair.

\medterm{Popliteal Artery} The main artery behind the knee. It continues from the thigh artery and divides into arteries supplying the lower leg and foot; blockage or injury can cause pain, coldness, numbness, or loss of circulation.

\medterm{Popliteal Vein} A deep vein behind the knee formed by the union of the anterior and posterior tibial veins. It continues as the femoral vein and can be affected by deep-vein thrombosis.

\medterm{Popliteal Artery Aneurysm} An abnormal enlargement of the artery behind the knee that may cause limb ischemia, clotting, or embolization.

\medterm{Popliteal Cyst} A fluid-filled swelling behind the knee, also called a Baker cyst, caused by communication between the joint and a bursa and often associated with arthritis, meniscal injury, or inflammation.

\medterm{Popliteal Cyst (Baker Cyst)} Alternate index label for Baker Cyst.

\medterm{Popliteal Fossa} Diamond-shaped space on the posterior knee, bounded by hamstrings superiorly and
gastrocnemius heads inferiorly. Contains popliteal artery/vein, tibial nerve, and common
peroneal nerve. The popliteal pulse is palpated at the inferior margin.

\medterm{Population Attributable Fraction} The estimated proportion of disease in a whole population that could be prevented if a specific harmful exposure were removed, assuming the exposure truly causes the disease and other conditions stayed the same.

\medterm{Population Attributable Risk} The estimated amount of disease in a population linked to a specific exposure. It helps show the potential public-health benefit of reducing that exposure, but it depends on accurate data and a causal relationship.

\medterm{Population Genetics} The study of how gene frequencies and genetic variation change within and between populations through inheritance, mutation, selection, migration, and genetic drift.

\medterm{Population Group} A defined set of people sharing a characteristic such as age, location, exposure, ancestry, occupation, or health condition for epidemiologic or clinical analysis.


\medterm{Porcelain Gallbladder} Hard calcium deposits in the wall of the gallbladder, usually after long-standing inflammation or gallstones. It may cause no symptoms but can be associated with gallbladder disease; imaging findings guide discussion of monitoring or removal.

\medterm{Porcine} Porcine means derived from or relating to pigs; porcine tissues, enzymes, and biologic products may be used in medicine after appropriate processing and screening.

\medterm{Pore} A pore is a small opening or channel in a surface, membrane, tissue, or material that permits passage of fluid, gas, particles, or cellular contents.
\medterm{Porencephaly} A fluid-filled cavity within the brain caused by abnormal development or damage before or after birth. Effects vary with its size and location and may include seizures, weakness, movement problems, or developmental delay.
\medterm{Porfimer Sodium} A photosensitizing medicine used with light therapy to destroy selected abnormal tissues, including some tumors; it causes prolonged light sensitivity.



\medterm{Porin} A porin is a channel-forming protein in an outer membrane, especially of Gram-negative bacteria, allowing passage of selected small molecules.

\medterm{Pork Tapeworm} A parasite acquired by eating undercooked pork; its larvae can invade tissues and cause cysticercosis.
\medterm{Porokeratosis} A group of skin conditions that cause one or more ring-shaped, scaly patches with a raised border. Some forms can rarely develop into skin cancer, so changing, bleeding, or persistent lesions need dermatology assessment.

\medterm{Poroma} A usually noncancerous tumor arising from a sweat gland. It often appears as a flesh-colored, pink, or red lump on the palms, soles, or other skin and may be removed if painful, bleeding, or uncertain in diagnosis.

\medterm{Porphyria} A group of mostly inherited disorders of heme biosynthesis in which reduced activity of a pathway enzyme causes porphyrins or their precursors to accumulate. Acute hepatic porphyrias can cause episodic abdominal, neurologic, or psychiatric symptoms, while cutaneous porphyrias can cause photosensitivity and skin injury; some forms have both patterns. Porphyria cutanea tarda is usually acquired rather than inherited.
\medterm{Porphyria Cutanea Tarda} The most common porphyria, caused by reduced activity of hepatic uroporphyrinogen decarboxylase (UROD), leading to accumulation of porphyrins in the liver and skin. It causes fragile, blistering skin on sun-exposed areas. Associated factors include hepatitis C, excess iron, alcohol use, estrogen exposure, and certain genetic UROD variants.

\medterm{Porphyrias} A group of metabolic disorders caused by reduced activity of enzymes in the heme-synthesis pathway. Acute hepatic forms can cause abdominal and neurologic symptoms; cutaneous forms can cause pain, burning, blistering, or other injury after light exposure. Examples include acute intermittent porphyria, variegate porphyria, hereditary coproporphyria, porphyria cutanea tarda, and erythropoietic protoporphyria.

\medterm{Porphyrin} A ring-shaped compound involved in heme and other biological molecules; abnormal accumulation occurs in porphyrias.
\medterm{Porphyrins Blood Test} A blood porphyrin test measures porphyrins or related compounds and helps evaluate suspected porphyria and disorders of haem production.

\medterm{Porphyrins Urine Test} A urine porphyrin test measures porphyrins or their precursors and helps investigate suspected porphyria or abnormal haem metabolism.

\medterm{Porphyromonas} A genus of anaerobic gram-negative bacteria found especially in the mouth and associated with periodontal and other infections.

\medterm{Porphyromonas Bennonis} An anaerobic oral bacterium rarely isolated from human specimens and occasionally associated with mixed infection.
\medterm{Porphyromonas Catoniae} An anaerobic gram-negative bacterium found in the mouth and occasionally associated with periodontal or mixed infection.

\medterm{Porphyromonas Endodontalis} A mouth bacterium that grows without oxygen and is associated with infected tooth pulp and root canals. It may contribute to dental infection, but diagnosis depends on symptoms, examination, imaging, and laboratory findings.

\medterm{Porphyromonas Levii} A Porphyromonas species found mainly in animals and rarely associated with human infection. An animal-associated bacterium rarely linked to human infection, usually in wounds or other opportunistic settings.
\medterm{Porphyromonas Somerae} A bacterium that grows without oxygen and is occasionally found in human wounds, abscesses, or other samples. When it causes infection, treatment depends on laboratory identification, infection location, and antibiotic susceptibility.

\medterm{Port-Wine Stain} A permanent birthmark caused by widened small blood vessels, appearing as flat pink, red, or purple skin.
\medterm{Portacaval Shunting} A surgical or image-guided connection that diverts blood around the liver to lower dangerously high pressure in the portal veins. It may reduce bleeding risk but can worsen confusion caused by liver disease.

\medterm{Portal Hypertension} Abnormally high pressure in the veins carrying blood from the digestive organs to the liver, most often because of cirrhosis. It can enlarge veins in the esophagus or stomach, cause abdominal fluid buildup, enlarge the spleen, or lead to bleeding.

\synonyms
PHT

\medterm{Portal Hypertensive Gastropathy} Changes in the stomach lining caused by high pressure in the veins around the liver. The lining may look mottled and bleed slowly or suddenly, causing anemia or vomiting of blood in severe cases.

\medterm{Portal System} A network of veins that carries blood from the stomach, intestines, spleen, pancreas, and gallbladder to the liver. The liver processes nutrients and many absorbed substances before the blood returns to the general circulation.

\medterm{Portal Vein} The large vein that carries blood from the digestive organs, spleen, pancreas, and gallbladder to the liver. A clot or increased pressure in this vein can cause enlarged abdominal veins, bleeding, or fluid buildup.

\medterm{Portal Vein Thrombosis} A blood clot that partly or completely blocks the portal vein or one of its branches. It may occur with cirrhosis, cancer, abdominal infection or inflammation, surgery, or a clotting disorder. It can cause portal hypertension, enlarged veins, spleen enlargement, intestinal injury, or gastrointestinal bleeding.

\medterm{Portion Of Cytosol} A region of the cell’s aqueous internal fluid outside membrane-bound organelles, containing ions, metabolites, proteins, and enzymes that support cellular reactions.

\medterm{Portion Size} The amount of food or drink served or eaten at one time, distinct from a standardized serving size and important for estimating energy, carbohydrate, sodium, and nutrient intake.

\medterm{Portopulmonary Syndrome} High blood pressure in the arteries of the lungs occurring with portal hypertension or serious liver disease. It can cause breathlessness, tiredness, chest discomfort, or fainting and affects decisions about liver treatment or transplantation.

\medterm{Position} The location or arrangement of a body part, organ, fetus, or patient. Position can affect examination findings, movement, comfort, circulation, breathing, childbirth, medical procedures, and interpretation of imaging.
\medterm{Positional Asphyxia} Fatal or potentially fatal impairment of breathing caused by body position that
restricts chest movement, airway patency, or venous return. Risk is increased by
restraint, intoxication, obesity, and inability to reposition; findings are nonspecific
and require scene and medical correlation.


\medterm{Positive} Positive means that a test, finding, result, or response is present or above a defined threshold; interpretation depends on the test and clinical setting.

\medterm{Positive Airway Pressure Treatment} Treatment that uses pressurized air delivered through a mask or similar device to keep the airway open during sleep or breathing difficulty. It is commonly used for obstructive sleep apnea; mask discomfort, dry airways, leaks, or skin irritation may occur.

\medterm{Positive End-Expiratory Pressure} Pressure kept in the airways at the end of breathing out, usually by a ventilator. It helps keep the lung’s air sacs open and improve oxygen levels, but excessive pressure can injure the lungs or lower blood pressure.

\medterm{Positive Feedback} A body-control process in which a response strengthens the original change, making the process grow until a specific event is completed. Examples include contractions during labor and blood clotting.
\medterm{Positive Finding} A positive finding is an observed symptom, examination sign, image, laboratory result, or test response that supports the presence of a specified condition or feature.

\medterm{Positive Predictive Value} The chance that a person with a positive test result truly has the condition being tested for. It depends on test accuracy and how common the condition is among the people tested.

\medterm{Positive Predictive Value In Surveillance} The proportion of reported signals or alerts that represent true events meeting the
surveillance case definition. It is influenced by the specificity of the signal and the
underlying prevalence of the condition.

\medterm{Positive Pressure Ventilation} Breathing support that pushes air or oxygen into the lungs under pressure. It may be delivered through a mask or a breathing tube and can support or replace a person’s own breathing during respiratory failure.

\medterm{Positive Psychology} The study of wellbeing, strengths, resilience, relationships, and factors that help people function and thrive. Its practices may support coping and quality of life but do not replace treatment for illness.

\medterm{Positive Symptoms} Symptoms reflecting an excess or distortion of normal mental functions, such as
hallucinations, delusions, disorganized speech, or grossly disorganized behavior. They
are contrasted with negative symptoms, which reflect loss or reduction of function.

\medterm{Positive-Strand RNA Virus} A virus whose RNA can be read directly by an infected cell as instructions for making viral proteins. This allows the virus to begin producing new parts soon after entering the cell.

\medterm{Positron} A particle with the same mass as an electron but a positive charge. Some medical tracers release positrons; when one meets an electron, energy is produced and detected during a PET scan.
\medterm{Positron Camera} A positron camera detects paired gamma rays from positron annihilation and is used in positron-emission tomography to map tracer distribution.

\medterm{Positron Emission Mammography} PET imaging of breast. Higher sensitivity than mammography for dense breasts. Limited availability. A breast imaging test combining positron emission scanning with mammography to show tissue activity and structure.
\medterm{Positron Emission Tomography} An imaging test that uses a small amount of radioactive tracer to show tissue activity, such as glucose use, blood flow, or specific proteins. It is often combined with CT or MRI to locate abnormal areas more precisely.

\medterm{Post Cibum} A medication instruction meaning that a medicine should be taken after meals. Clear instructions should spell this out to prevent confusion with other uses of the abbreviation P.C.

\synonyms
P.C.; After Meals

\medterm{Post-Acute Rehabilitation} Time-limited interdisciplinary care after an acute illness, injury, or hospitalization
to restore mobility, self-care, cognition, communication, and safe community function.
Intensity and setting depend on medical stability and rehabilitation potential.

\medterm{Post-Cardiac Arrest Management} Care after return of spontaneous circulation that supports breathing and circulation, identifies and treats the cause of
the arrest, manages temperature and seizures when indicated, and assesses neurologic
recovery over time.

\medterm{Post-Intensive Care Syndrome} New or worsened physical, thinking, emotional, or mental-health problems that persist after critical illness or an intensive-care stay. It may include muscle weakness, poor memory or concentration, anxiety, depression, sleep problems, or post-traumatic stress.

\medterm{Post-Cardiac Injury Syndrome} Alternate index label; see \textit{Dressler syndrome}.

\medterm{Post-Chemotherapy Cognitive Impairment} Alternate index label; see \textit{Chemo brain}.

\medterm{Post-Concussion Syndrome} Ongoing symptoms after a concussion, such as headache, dizziness, tiredness, poor concentration, sleep problems, or mood changes. Symptoms usually improve gradually; evaluation and a stepwise return to activity help guide recovery.
\medterm{Post-Concussional Syndrome} Persistent physical, cognitive, emotional, or sleep symptoms after a concussion or mild traumatic brain injury.
\medterm{Post-Dural Puncture Headache} A headache that begins after a spinal tap or accidental puncture during an epidural. It is usually worse when sitting or standing and improves when lying down because spinal fluid has leaked; neck stiffness, nausea, ringing ears, or light sensitivity may occur.

\medterm{Post-Exposure Prophylaxis} Medicine or vaccination given after a possible exposure to prevent infection from developing. It is used for selected exposures to HIV, hepatitis B, rabies, and other infections; the correct treatment depends on the pathogen and must begin promptly.

\medterm{Post-Hepatic Jaundice (Obstructive Jaundice)} Jaundice caused by obstruction of bile flow after bilirubin has been processed by the liver, commonly from gallstones, strictures, or tumours.

\synonyms
Obstructive Jaundice

\medterm{Post-Hoc} Performed or analyzed after an event; in statistics, a post-hoc analysis follows the planned primary analysis.
\medterm{Post-Infectious Cough} A chronic cough persisting 3--8 weeks after a URI, caused by transient bronchial
hyperresponsiveness and airway inflammation. It is self-limited and usually resolves
without treatment. Inhaled corticosteroids, antitussives, and time are the mainstays of
management.

\medterm{Post-Kala-Azar Dermal Leishmaniasis} A dermatological sequela of visceral leishmaniasis (kala-azar) caused by Leishmania donovani that develops months to years after successful or incomplete antimonial or amphotericin treatment, predominantly observed in South Asia and East Africa. It manifests as an extensive cutaneous eruption progressing from asymptomatic hypopigmented macules to erythematous papules, plaques, and parasite-rich nodular lesions, representing a vital epidemiologic reservoir for anthroponotic sandfly transmission.

\synonyms
PKDL; Post-Kala-Azar Dermatosis

\medterm{Post-Mortem} Occurring after death or relating to examination of a body afterward. A post-mortem examination can help determine the cause of death, identify disease, and provide information for families or legal investigations.
\medterm{Post-Nasal Drip (Chronic Rhinitis)} Alternate index label for Chronic Rhinitis.

\medterm{Post-Natal Depression} Depression beginning during pregnancy or after childbirth, with persistent low mood or loss of interest and possible effects on sleep, bonding, and daily functioning.

\medterm{Post-Obstructive Pneumonia} Lung infection developing beyond a blocked airway. A tumor, inhaled object, thick mucus, or narrowing can prevent drainage and allow infection; antibiotics may not cure it until the blockage is found and relieved.

\medterm{Post-Operative Atelectasis} Partial collapse of air sacs in the lung after surgery. Anesthesia, pain, shallow breathing, mucus, and lying still can contribute, causing shortness of breath or low oxygen. Pain control, movement, and breathing exercises may help prevent or treat it.

\medterm{Post-Orgasmic Illness Syndrome} A rare condition in which flu-like symptoms, extreme tiredness, headache, nasal or eye symptoms, and difficulty thinking begin soon after orgasm and last for days. The cause is uncertain; repeated episodes need medical assessment to exclude other conditions.

\medterm{Postcoital Bleeding} Bleeding from the vagina after sexual intercourse. Causes include infection, dryness, injury, polyps, pregnancy-related conditions, or cancer.

\synonyms
POIS

\medterm{Post-Partum Haemorrhage} Excessive bleeding after childbirth, most commonly caused by failure of the uterus to contract, retained placental tissue, genital-tract trauma, or clotting problems.

\medterm{Post-Pneumonectomy Syndrome} A rare problem after removal of one lung in which the remaining chest structures shift and press on the airway or major blood vessels. It may cause breathlessness, noisy breathing, repeated infections, or difficulty swallowing.

\medterm{Post-Prandial} Occurring after a meal. Post-prandial blood glucose helps assess diabetes, while post-prandial low blood pressure or chest pain may indicate a circulation or heart problem and needs assessment.
\medterm{Post-Splenectomy Complications} Post-splenectomy complications include infection by encapsulated bacteria, blood clots, bleeding, and changes in blood-cell counts; vaccination and follow-up reduce risk.

\medterm{Post-Streptococcal Glomerulonephritis} Kidney inflammation that can develop after a group A streptococcal throat or skin infection. It may cause dark urine, swelling, reduced urine, high blood pressure, tiredness, or kidney failure and needs medical evaluation.

\medterm{Post-Surgical Pain} Pain after an operation caused by tissue injury and inflammation. Its severity depends on the procedure and the person; uncontrolled pain can limit breathing, movement, sleep, and recovery and should be discussed with the care team.

\medterm{Post-Surgical Rehabilitation} A planned program to restore movement, strength, independence, and safe return to daily activities after surgery. It may include wound care, pain control, exercises, mobility training, and advice tailored to the operation and the person’s recovery.

\medterm{Post-Term Pregnancy} Pregnancy lasting 42 completed weeks or longer. The placenta may provide less support, the baby may become unusually large, and amniotic fluid may decrease. Clinicians monitor parent and baby and discuss delivery timing to reduce complications.

\medterm{Post-Test Probability} The estimated chance that a person has a disease after a test result is known. It depends on the original likelihood of disease and how accurately the test distinguishes people with and without the condition.

\medterm{Post-Thrombotic Syndrome} Chronic pain, heaviness, edema, skin change, and sometimes ulceration developing after
deep-vein thrombosis because of persistent venous obstruction and valvular damage.
Compression, exercise, skin care, and selected interventions may reduce disability.

\medterm{Post-Transfusion Purpura} A rare immune reaction that causes a severe fall in platelets about 5–10 days after a blood transfusion. It can produce pinpoint skin bleeding, bruising, nosebleeds, or dangerous internal bleeding and requires urgent treatment.

\medterm{Post-Traumatic Headache} A headache beginning or continuing after a head or neck injury. It may resemble migraine or tension headache and can occur with dizziness, poor concentration, sleep problems, nausea, or light sensitivity.

\medterm{Post-Traumatic Stress Disorder} A mental-health disorder that may follow actual or threatened death, serious injury, sexual violence, or another traumatic event. It causes ongoing unwanted memories, nightmares, avoidance, emotional changes, and heightened alertness that interfere with daily life.

\medterm{Post-Vasectomy Pain Syndrome} Persistent or recurring testicular or scrotal pain lasting at least three months after vasectomy. The pain may result from nerve irritation, inflammation, or pressure in the reproductive tract.
\synonyms
PVPS; Chronic Post-Vasectomy Pain

\medterm{Post-Void Residual} The amount of urine remaining in the bladder after urination, measured by ultrasound or
catheterization. A persistently high residual may indicate incomplete emptying, but its
significance depends on the volume, timing, symptoms, and clinical context.

\medterm{Post-Void Residual Volume} The amount of urine remaining in the bladder immediately after urination, measured by
ultrasound or catheterization. An elevated residual suggests incomplete emptying but
must be interpreted with timing, bladder volume, symptoms, and repeated measurements.


\medterm{Postanesthesia Care Unit} The postanesthesia care unit is a monitored recovery area where patients are
observed during the immediate postoperative period. Admission assessment includes vital
signs, level of consciousness, pain assessment, and surgical site evaluation.

\medterm{Postcholecystectomy Syndrome} Persistence or recurrence of symptoms after cholecystectomy. Causes include retained
common bile duct stones, sphincter of Oddi dysfunction, bile salt malabsorption,
irritable bowel syndrome, and other gastrointestinal conditions. Evaluation includes
liver function tests, MRCP, and endoscopic ultrasound.

\medterm{Postconcussion Syndrome} Persistent symptoms after a concussion, such as headache, dizziness, fatigue, sleep disturbance, slowed thinking, or mood change. Recovery varies, and worsening neurologic symptoms need urgent assessment.

\medterm{Posterior} Located behind or toward the back of the body or a body structure. The term is used to describe anatomical direction, position, movement, imaging findings, and areas affected by disease or injury.
\medterm{Posterior Chamber} The eye space between the iris and lens, filled with clear fluid that nourishes tissues and helps maintain pressure.


\medterm{Posterior Cruciate Ligament Injury} A stretch or tear of the ligament deep inside the knee that prevents the shinbone from moving too far backward. It may follow a blow to the shin, a fall on a bent knee, or severe twisting and can cause pain, swelling, and instability.

\medterm{Posterior Cruciate Ligament Tear} A partial or complete tear of the strong ligament deep inside the knee that limits backward movement of the shinbone. It often follows a forceful blow to the shin or a fall on a bent knee and may cause pain, swelling, and knee instability. Treatment depends on severity and other injuries.

\medterm{Posterior Descending Artery} A heart artery that runs along the back groove between the lower chambers and supplies part of the back and lower heart muscle. It usually branches from the right coronary artery but may arise from another heart artery.

\medterm{Posterior Drawer Test} A knee examination in which the lower leg is gently pushed backward while the knee is bent. Excessive backward movement may indicate injury to the ligament that stabilizes the back of the knee.

\medterm{Posterior Fossa Tumor} A growth in the lower back part of the skull, where the cerebellum and brainstem are located. It may cause headache, vomiting, unsteady movement, eye problems, or changes in breathing by pressing on nearby tissue or blocking brain fluid.

\medterm{Posterior Heel Bursitis} Inflammation of a bursa near the back of the heel, causing pain and swelling.
\synonyms
Retrocalcaneal bursitis

\medterm{Posterior Keyhole Foraminotomy} An operation through a small opening in the back that enlarges the passage where a spinal nerve leaves the spine. It relieves pressure from causes such as a disk bulge or arthritis and may reduce radiating pain or weakness.
\medterm{Posterior Myocardial Infarction} A heart attack affecting the back wall of the heart, usually because a coronary artery is blocked. It may be missed on a standard ECG, so additional ECG leads or imaging may be needed when symptoms suggest it.

\medterm{Posterior Pituitary} The back part of the pituitary gland, which stores and releases hormones made in the brain. Vasopressin helps control body water and blood pressure, while oxytocin supports uterine contractions and release of breast milk.

\medterm{Posterior Prolapse} Alternate index label; see \textit{Posterior vaginal prolapse (rectocele)}.

\medterm{Posterior Reversible Encephalopathy Syndrome} Brain swelling caused by disturbed blood-flow regulation, often linked to very high blood pressure, pregnancy-related hypertension, kidney disease, severe infection, or immune-suppressing medicines. Headache, seizures, confusion, or vision loss require emergency treatment because prompt care can prevent lasting injury.

\medterm{Posterior Rhinorrhea} Posterior rhinorrhea is mucus draining from the nose into the throat, commonly causing throat clearing or cough from allergy, infection, irritation, or sinus disease.

\medterm{Posterior Subcapsular Cataract} A cloudy area at the back of the eye’s lens, just beneath its outer capsule. Because it often lies in the path of light, it can cause glare, blurred vision, and difficulty reading, especially in bright light.

\medterm{Posterior Tibial Artery} A major artery running down the back of the lower leg and behind the inner ankle. It supplies the lower leg and sole of the foot; its pulse is checked when assessing blood flow to the foot.

\medterm{Posterior Urethral Valves} Abnormal tissue folds present in the urethra of some male babies that block urine leaving the bladder. They can cause a weak urine stream, swollen kidneys, urinary infection, or kidney damage and require specialist treatment.

\medterm{Posterior Uveitis} Inflammation at the back of the eye involving the retina, the layer beneath it, or both. It may cause floaters, blurred or missing vision, or eye pain and can result from infection, immune disease, or another inflammatory condition.

\medterm{Posterior Vitreous Detachment} Separation of the jelly inside the eye from the retina, usually as an age-related change. It may cause new flashes or floaters; urgent eye examination is needed because a retinal tear or detachment can occur at the same time.

\medterm{Posteroanterior} Posteroanterior means directed from back to front; a PA chest radiograph uses this projection to reduce magnification of the heart.

\medterm{Postganglionic Fiber} An autonomic nerve fiber extending from a ganglion to its target organ after the synapse, carrying sympathetic or parasympathetic signals to smooth muscle, glands, or the heart.

\medterm{Postherpetic Neuralgia} Persistent nerve pain in an area affected by shingles after the rash has healed.
\synonyms
Neuralgia, Postherpetic


\medterm{Posthitis} Inflammation of the foreskin, causing redness, swelling, itching, pain, discharge, or difficulty retracting it; infection, irritation, diabetes, and poor hygiene are possible contributors.

\medterm{Posthumous} Occurring, discovered, or published after a person's death. A posthumous diagnosis may be based on records, stored samples, or autopsy findings.
\medterm{Postinflammatory Hyperpigmentation} Darkening of skin after inflammation or injury, caused by increased melanin production or its deposition in the epidermis or dermis.

\medterm{Postmenopausal} Postmenopausal describes the period after menopause, when ovarian estrogen production and menstrual periods have permanently ceased.

\medterm{Postmenopausal Bleeding} Vaginal bleeding after menopause. Even light or one-time bleeding needs medical evaluation because causes include tissue thinning, polyps, hormone effects, and cancer.

\medterm{Postmenopausal Osteoporosis} Bone loss caused by lower estrogen levels associated with menopause. Sometimes called type I osteoporosis.
\medterm{Postmenopause} The stage of life after menopause, when menstrual periods have permanently stopped because the ovaries no longer release eggs regularly and produce much less estrogen.

\medterm{Postmortem Artifact} A change that happens to a body after death and can look like disease or injury. Causes include decomposition, medical resuscitation, embalming, insects, animals, or movement of blood after death; examiners must distinguish it from changes before death.

\medterm{Postmortem Interval} The estimated time between death and the examination or discovery of a body. Investigators consider body temperature, stiffness, pooling of blood, chemical changes, environment, and insect activity, but the estimate is not exact.

\medterm{Postmortem Toxicology} Testing blood, tissues, or other samples after death for medicines, alcohol, poisons, or their breakdown products. Results must be interpreted with the medical history, dose, sample site, decomposition, and possible movement of substances after death.

\medterm{Postmyocardial Infarction Syndrome} Alternate index label; see \textit{Dressler syndrome}.

\medterm{Postnatal} Relating to the period after birth, especially the newborn and infant period. Postnatal care includes checking breathing, feeding, temperature, weight, jaundice, development, and signs of infection or other illness.
\medterm{Postnatal Care} Health care for the parent and baby after birth. It includes checking bleeding, blood pressure, wounds, pain, feeding, mood, newborn growth, and warning signs, while providing guidance about recovery, contraception, vaccination, and follow-up.

\medterm{Postnatal Depression} A depressive episode occurring after childbirth, with persistent low mood, anhedonia,
sleep or appetite disturbance, guilt, impaired bonding, or suicidal thoughts.

\medterm{Postop} Postop is an abbreviation for postoperative, meaning relating to the period after a surgical procedure.
\medterm{Postoperative} Postoperative means occurring or provided after surgery, including recovery, wound care, pain control, rehabilitation, and complication monitoring.

\medterm{Postoperative Care} Care after surgery that monitors breathing, circulation, pain, fluids, the wound, movement, eating, and possible complications. It may include early walking, exercises, medicines, and prevention of blood clots.

\medterm{Postoperative Cognitive Dysfunction} Problems with memory, attention, thinking, or planning after surgery. They may be temporary or persistent, especially after major operations; assessment compares the person with their usual abilities and checks for delirium, medicines, infection, or other causes.

\medterm{Postoperative Complications} Problems that develop after an operation, including infection, bleeding, blood clots, breathing problems, organ dysfunction, wound separation, or medication reactions. Risk depends on the procedure and the person's health; severe pain, fever, breathlessness, or heavy bleeding needs urgent care.

\medterm{Postoperative Delirium} A sudden change in attention and thinking after surgery. A person may become confused, restless, unusually sleepy, or see things that are not there; symptoms may fluctuate and require prompt assessment for medical causes.

\medterm{Postoperative Hemorrhage} Bleeding after surgery that may occur at the operative site or internally and can cause shock, anemia, or the need for urgent intervention.

\medterm{Postoperative Ileus} Temporary slowing or stopping of bowel movement after surgery, especially abdominal surgery. It can cause bloating, nausea, vomiting, pain, and inability to pass gas or stool; medicines, electrolyte problems, and the operation may contribute.

\medterm{Postoperative Infection} An infection developing after surgery, involving the incision, an organ space, implanted material, or the bloodstream.

\medterm{Postoperative Nausea And Vomiting} Nausea, retching, or vomiting occurring after surgery. Risk is influenced by factors
such as previous motion sickness or postoperative nausea, anesthetic technique, opioid
use, and patient characteristics. Prevention and treatment are tailored to the individual
risk and procedure.

\medterm{Postoperative Pain Management} Management of pain after surgery using a combination of medicines, regional
techniques, physical measures, and rehabilitation. The plan is tailored to the procedure, pain severity, medical
conditions, and risks of adverse effects.

\medterm{Postoperative Period} The postoperative period is the time after an operation when recovery, pain control, wound care, mobility, nutrition, complications, and follow-up are monitored.

\medterm{Postoperative Pulmonary Complication} A breathing problem developing after surgery, such as partial lung collapse, pneumonia, aspiration, worsening asthma or COPD, or respiratory failure. Risk is reduced by suitable pain control, movement, breathing exercises, and careful monitoring.

\medterm{Postpartum} Relating to the period after childbirth. Recovery includes changes in bleeding, pain, breasts, hormones, sleep, mood, and body function; important complications include heavy bleeding, infection, high blood pressure, blood clots, depression, and psychosis.

\synonyms
Postnatal

\medterm{Postpartum Anxiety} Excessive, persistent, or distressing anxiety during pregnancy or after childbirth that interferes with sleep, daily
function, bonding, or self-care. It may occur with panic, intrusive worries, or depression and deserves clinical
assessment; thoughts of self-harm or harm to the baby require emergency help.

\medterm{Postpartum Blues} A common, short-lived period of tearfulness, mood changes, irritability, or anxiety beginning
soon after childbirth and usually improving within two weeks. Support, rest, and practical
help are important; persistent, severe, psychotic, or suicidal symptoms require prompt
assessment for a postpartum mental-health disorder.

\medterm{Postpartum Cardiomyopathy (Peripartum Cardiomyopathy)} A dilated cardiomyopathy causing heart failure toward
the end of pregnancy or during the months after delivery. Symptoms may include breathlessness, difficulty lying flat,
swelling, fatigue, or chest discomfort and require prompt assessment.

\synonyms
Peripartum Cardiomyopathy

\medterm{Postpartum Care} Clinical care during the weeks after childbirth that assesses bleeding, uterine involution, blood pressure, infection, wound healing, breastfeeding, mental health, contraception, and recovery.

\medterm{Postpartum Contraception} Contraceptive counseling and care after childbirth, taking account of breastfeeding, venous-thromboembolism risk, medical conditions, fertility intentions, and the timing of method initiation.

\medterm{Postpartum Depression} Depression that begins during pregnancy or after childbirth and lasts beyond ordinary short-term mood changes. It may cause persistent sadness, loss of interest, guilt, anxiety, sleep or appetite changes, difficulty bonding, or thoughts of self-harm and needs treatment.

\synonyms
Depression, Postpartum
PPD (Postpartum Depression)

\medterm{Postpartum Endometritis} Infection of the inner lining of the uterus after childbirth, more common after cesarean birth. It may cause fever, lower-abdominal pain, uterine tenderness, fast heartbeat, or foul-smelling discharge and needs prompt treatment.

\medterm{Postpartum Haemorrhage} Excessive bleeding after childbirth, commonly defined by blood loss of at least 1,000 mL or bleeding with hypovolemia after birth.
\medterm{Postpartum Hemorrhage} Excessive bleeding after childbirth, usually caused by poor uterine contraction, retained placental tissue, genital-tract
trauma, or a clotting disorder. It can cause shock and requires urgent obstetric
assessment and treatment.

\medterm{Postpartum Hemorrhage (Secondary)} Heavy or abnormal bleeding occurring more than 24 hours and up to 12 weeks after childbirth. Causes include retained placental tissue, infection, poor healing where the placenta was attached, or a bleeding disorder. Heavy bleeding, dizziness, fainting, fever, or severe pain requires urgent care.

\synonyms
Secondary

\medterm{Postpartum Infection} An infection after childbirth, commonly involving the uterus, surgical wound, bladder, breasts, or bloodstream. Fever, worsening abdominal or pelvic pain, foul-smelling discharge, wound redness, breast redness, or feeling very unwell requires prompt medical assessment.

\medterm{Postpartum Preeclampsia} New-onset hypertension with pre-eclampsia features after childbirth, sometimes occurring several days or weeks postpartum. Severe headache, visual symptoms, chest pain, shortness of breath, or seizures require urgent assessment.

\medterm{Postpartum Psychosis} A rare but severe psychiatric emergency that can begin suddenly during the days or weeks after childbirth. Symptoms may include hallucinations, delusions, paranoia, severe confusion, mania, depression, disorganized behavior, or rapidly changing moods. It can worsen quickly and may threaten the safety of the mother or baby, so immediate medical assessment is required.

\synonyms
Postnatal Psychosis
Puerperal Psychosis

\medterm{Postpartum Thyroiditis} Inflammation of the thyroid during the first year after childbirth, caused by an immune reaction. It may first cause anxiety, sweating, and a fast heartbeat, followed by tiredness, cold intolerance, or low thyroid function; recovery varies.

\medterm{Postpartum Urinary Retention} Inability to empty the bladder adequately after childbirth, sometimes associated with epidural anesthesia, perineal trauma, swelling, or prolonged labor. Prompt assessment prevents bladder overdistension and injury.

\medterm{Postprandial} Postprandial means occurring after a meal, especially the changes in glucose, insulin, digestion, and circulation after eating.

\medterm{Postprandial Period} The postprandial period is the interval after eating when digestion and nutrient absorption occur and glucose, insulin, and other metabolic responses change.

\medterm{Postrenal AKI} Sudden loss of kidney function caused by blocked urine flow. Prostate enlargement, stones, tumors, scarring, or bladder-nerve problems may be responsible; kidney damage becomes more likely when the blockage is severe or lasts a long time.

\medterm{Poststreptococcal Glomerulonephritis} Poststreptococcal glomerulonephritis is kidney inflammation occurring after certain streptococcal infections, causing blood or protein in urine, swelling, high blood pressure, and reduced kidney function.

\medterm{Poststreptococcal Pustulosis} A rare immune-related outbreak of pus-filled spots after a streptococcal throat infection. The spots usually appear on reddened skin of the trunk or limbs and should be assessed to confirm the cause and exclude infection or a medicine reaction.

\medterm{Postsynaptic Membrane} The receiving surface of a nerve or muscle cell where a neighboring nerve releases a chemical signal. Receptors in this surface detect the signal and help start or change an electrical response.

\medterm{Postterm Pregnancy} Pregnancy continuing beyond 42 completed weeks. The placenta may work less effectively and the baby may become unusually large or pass stool before birth, so fetal wellbeing and delivery timing require medical monitoring.

\medterm{Posttraumatic Stress Disorder} Alternate index label for Post-Traumatic Stress Disorder.

\synonyms
PTSD (Posttraumatic Stress Disorder)

\medterm{Postural} Postural means relating to body position or posture, including symptoms or vital signs that change when a person sits, stands, or lies down.

\medterm{Postural Drainage} Using specific body positions so gravity helps mucus move from smaller airways into larger ones, where it can be coughed out. A respiratory clinician selects positions based on the affected lung area and the person’s condition.
\medterm{Postural Hypotension} A drop in blood pressure after standing that may cause dizziness, blurred vision, weakness, or fainting. It can result from dehydration, medicines, nerve disorders, blood loss, or heart problems and increases fall risk.
\medterm{Postural Orthostatic Tachycardia Syndrome} A disorder of the autonomic nervous system, which helps control heart rate and blood flow. When a person stands, the heart rate rises unusually quickly while blood pressure does not fall substantially, causing ongoing symptoms of poor tolerance of upright posture. In adults, the heart rate typically rises by at least 30 beats per minute within 10 minutes of standing; the increase is at least 40 beats per minute in adolescents. Symptoms may include dizziness, fainting or near-fainting, racing heartbeat, tiredness, weakness, headache, blurred vision, nausea, trouble thinking, or poor ability to exercise.

\medterm{Posture} The alignment and position of the body while resting or moving. Posture depends on bones, joints, muscles, nerves, vision, and balance and can affect comfort, breathing, movement, and injury risk.
\medterm{Potassium} An essential mineral and body salt needed for nerve signals, muscle movement, and normal heart rhythm.
\medterm{Potassium (K)} The chemical symbol for potassium, a major intracellular electrolyte needed for nerve signaling, muscle contraction, and cardiac electrical activity.
\medterm{Potassium Acetate} A potassium salt used to replace potassium or provide acetate in selected medical and laboratory applications; dosing requires monitoring of kidney function and serum potassium.

\medterm{Potassium Binders} Medicines taken by mouth that trap potassium in the intestine so it leaves the body in stool. They help lower persistently high blood potassium in selected people, especially with kidney disease, but require monitoring for constipation, swelling, and other electrolyte changes.

\medterm{Potassium Carbonate Poisoning} Swallowing potassium carbonate can irritate or burn the mouth, throat, and stomach and may raise blood potassium dangerously. Vomiting, weakness, abnormal heartbeat, or breathing difficulty requires emergency care.

\medterm{Potassium Channel} A membrane protein that allows potassium ions to cross cells and helps control electrical excitability, heart rhythm, and hormone release.

\medterm{Potassium Chloride} A salt used to prevent or treat low blood potassium. It may be taken by mouth or given into a vein; too much can cause dangerous heart-rhythm problems, especially when kidney function is reduced.

\medterm{Potassium Citrate} A medicine that makes urine less acidic and increases urinary citrate, helping prevent some kidney stones. It can dangerously raise blood potassium in people with kidney disease or those taking certain medicines.

\medterm{Potassium Cyanide} A highly poisonous chemical that prevents cells from using oxygen. Swallowing or inhaling it can rapidly cause headache, confusion, seizures, collapse, heart failure, and death and requires emergency treatment.

\medterm{Potassium Deficiency} An abnormally low blood potassium level, which can cause weakness, cramps, constipation, abnormal heart rhythms, or paralysis.

\medterm{Potassium Dichromate} A toxic, corrosive chemical that can burn the skin, eyes, mouth, and digestive tract. Swallowing or inhaling it may damage the kidneys, liver, blood, or lungs and requires urgent poison-control or emergency care.

\medterm{Potassium Hydroxide Poisoning} A severe chemical burn caused by swallowing or contacting potassium hydroxide, a strong alkali. It can damage the mouth, throat, stomach, skin, eyes, and lungs; rinse exposed areas and seek emergency help immediately.

\medterm{Potassium In Diet} Dietary potassium comes from fruits, vegetables, beans, dairy, and other foods; intake may need adjustment in kidney disease or with medicines that change potassium levels.

\medterm{Potassium Iodide} Potassium iodide supplies iodide and is used in selected thyroid conditions, radiation emergencies, and some expectorant preparations; excess can disturb thyroid function.

\medterm{Potassium Ion} The positively charged form of potassium, the principal intracellular cation that helps establish resting membrane potential, nerve transmission, muscle contraction, and cardiac rhythm.

\medterm{Potassium Permanganate} A chemical sometimes used as a very dilute skin antiseptic or drying solution. Concentrated material can burn tissue, stain skin, and poison if swallowed, so it must be prepared and used carefully.
\medterm{Potassium Test} A potassium test measures potassium in blood, usually to assess heart, muscle, kidney, or medication-related electrolyte problems.

\medterm{Potassium Urine Test} A urine potassium test measures potassium excretion and helps determine whether abnormal blood potassium is related to kidney losses or other causes.

\medterm{Potassium-Sparing Diuretic} A diuretic that promotes sodium excretion while reducing potassium loss, such as spironolactone or amiloride.
\medterm{Potassium-Sparing Diuretics} Water tablets that help remove salt and water while limiting potassium loss. They are used for selected blood-pressure, heart, liver, or hormone conditions, but can raise potassium dangerously, especially in kidney disease or with certain medicines.

\medterm{Potato Plant Poisoning - Green Tubers And Sprouts} Green potatoes, sprouts, and leaves can contain high levels of natural toxins that cause nausea, vomiting, diarrhea, abdominal pain, confusion, weakness, or rarely seizures. Contact poison control after significant ingestion.

\medterm{Potbellies And Toddlers} A rounded abdomen in a toddler is often normal body shape or posture, but persistent distension with pain, vomiting, poor growth, constipation, or a mass needs medical assessment.

\medterm{Potency} The amount of a drug required to produce a specified effect. A more potent drug produces that effect at a
lower dose, but potency does not indicate the maximum effect.


\medterm{Potential} Potential describes a capacity, possibility, or stored energy that may become clinically relevant depending on circumstances and additional findings.

\medterm{Potential Acuity Meter} An eye-testing device that estimates the vision a person might achieve after treatment of an eye problem such as a cataract. Results can be limited by retinal, nerve, or other disease.


\medterm{Potocki-Lupski Syndrome} A genetic condition caused by an extra copy of a small section of chromosome 17. It may cause low muscle tone, feeding and growth problems, developmental or learning difficulties, autism-related features, behavior differences, and congenital heart defects.

\synonyms
PTLS; Duplication 17p11.2 Syndrome

\medterm{Pott Disease} Tuberculosis infection of the bones of the spine. It can cause persistent back pain, collapse or bending of the spine, an abscess, and pressure on the spinal cord leading to weakness or paralysis.

\medterm{Pott's Disease} Tuberculosis of the spine, causing vertebral destruction, back pain, deformity, or neurologic compromise.
\medterm{Pott's Fracture} An ankle fracture-dislocation involving the malleoli, the bony prominences on each side of the ankle. The historical term is used variably, so the exact fracture pattern should be described on imaging.
\medterm{Potter Syndrome} A pattern of birth abnormalities caused by very low amniotic fluid, often because both fetal kidneys are absent or cannot make urine. It may cause underdeveloped lungs, distinctive facial features, and limb changes and is usually life-threatening.

\medterm{Pouch} A sac-like body structure or a surgically created reservoir. Examples include a cheek pouch, a pouch in the digestive tract, and an ileal pouch made from small intestine.
\medterm{Pouch Of Douglas} The pouch of Douglas is the peritoneal recess between the uterus and rectum, also called the rectouterine pouch, where fluid may collect.

\medterm{Pound} A unit measuring weight or force, with meaning determined by context. In healthcare, pounds may describe body weight, pressure, equipment loads, or medicine instructions and should be distinguished from metric units.
\medterm{Povidone} Povidone is a water-soluble polymer used in medicines, eye products, and other preparations; it is also part of povidone-iodine antiseptic.

\medterm{Povidone-Iodine} Povidone-iodine is an iodine-containing antiseptic used on skin or tissue before selected procedures; it can irritate skin and affect the thyroid with extensive use.

\medterm{Powassan Encephalitis} Inflammation of the brain caused by Powassan virus, a tick-borne flavivirus. It can cause fever, headache, vomiting, confusion, seizures, or weakness and may require hospitalization.

\medterm{Powassan Virus Disease} Powassan virus disease is a rare tick-borne infection that can cause fever, headache, vomiting, confusion, seizures, meningitis, or encephalitis. Treatment is supportive, and prevention relies on avoiding ticks and removing them promptly.

\medterm{Powder Diffraction} An analytical method that identifies crystalline substances by measuring how an X-ray or other beam is diffracted by powdered material.

\medterm{Power} The rate at which work is performed; in exercise, it combines force and movement speed.
\medterm{Power Training} Exercise designed to improve the ability to produce force rapidly.

\synonyms
Explosive training

\medterm{Pox} A disease producing vesicles, pustules, or pitted scars; the term is used in names such as smallpox and cowpox.
\medterm{Poxviridae} A family of large DNA viruses that includes orthopoxviruses, molluscipoxvirus, and parapoxviruses; some cause significant human disease.

\medterm{Poxviridae Infection} An infection caused by a poxvirus, ranging from localized skin lesions to severe systemic disease depending on the virus and host.

\medterm{Poxvirus} A family of large DNA viruses that includes variola, vaccinia, mpox, and molluscum contagiosum viruses.
\medterm{Ppar Alpha} A protein inside cells that helps control how the body uses fats. Fibrate medicines activate it to lower some blood fats, especially triglycerides.

\medterm{Ppar Delta} A protein inside cells that helps regulate fat use, energy balance, inflammation, and the development of some tissues. It is mainly studied in laboratory and drug research.
\medterm{PPD} PPD may mean purified protein derivative used in tuberculosis skin testing or post-partum depression; context determines the intended meaning.

\medterm{PPD Skin Test} A PPD skin test places a small amount of tuberculin under the skin to assess immune evidence of tuberculosis exposure; interpretation depends on risk and the size of the reaction.

\medterm{PR Interval} The ECG interval from the start of atrial depolarization to the start of ventricular depolarization, reflecting atrioventricular conduction.

\medterm{Practitioner} A person qualified and authorized to practice a healthcare profession.
\medterm{Prader-Labhart-Willi Syndrome} Alternate index label; see \textit{Prader-Willi syndrome}.

\medterm{Prader-Willi Syndrome} A genetic disorder that often causes weak muscle tone and feeding difficulty in infancy, followed later by unusually strong hunger, obesity, short stature, learning difficulties, hormone problems, and behavioral changes.

\synonyms
Prader-Labhart-Willi Syndrome

\medterm{Pragmatic} Focused on what works in ordinary real-world conditions. A pragmatic clinical trial studies how an intervention performs in usual healthcare rather than under highly controlled conditions.
\medterm{Pragmatic Trial} A pragmatic trial evaluates an intervention under usual real-world care conditions to estimate how well it works in routine practice.

\medterm{Pralatrexate} A folate-blocking chemotherapy medicine used for some relapsed or treatment-resistant peripheral T-cell lymphomas. It can cause mouth sores, low blood counts, infection, fatigue, nausea, and other serious treatment-related effects.

\medterm{Pralidoxime} An oxime that can reactivate inhibited acetylcholinesterase when administered early
after organophosphate poisoning, particularly for nicotinic weakness and
respiratory-muscle dysfunction. It is used with atropine and supportive care according
to exposure and local guidance.

\medterm{Pramipexole} A medicine that stimulates dopamine receptors and is used mainly for Parkinson disease or restless legs syndrome. It can cause nausea, sleepiness, dizziness, sudden sleep, hallucinations, or impulse-control problems.
\medterm{Prandial} Prandial means relating to a meal, as in preprandial before eating or postprandial after eating.
\medterm{Prasugrel} An antiplatelet medicine that prevents blood cells from clumping. It is used in selected people with acute coronary syndrome receiving coronary stents, but it increases bleeding risk and is unsuitable for some patients.
\medterm{Pravastatin} A statin that lowers LDL cholesterol by inhibiting HMG-CoA reductase.
\medterm{Pravastatin Sodium} A statin medicine that lowers harmful LDL cholesterol by reducing cholesterol production in the liver. It helps reduce heart attack and stroke risk alongside diet, exercise, and other risk-factor treatment.

\medterm{Praziquantel} A medicine used to treat several infections caused by parasitic worms, including schistosomiasis, tapeworms, and flukes. It damages the worms so the body can remove them; dizziness, sleepiness, nausea, or abdominal discomfort may occur.
\medterm{Pre-Cancerous} Describing abnormal cells or tissue that are not cancer but have a higher-than-usual chance of becoming cancer. The risk varies by the tissue and the severity of change; monitoring or removal may reduce the risk.
\medterm{Pre-Eclampsia} A disorder that can develop during pregnancy, usually after 20 weeks, or during the weeks after childbirth. It involves new high blood pressure together with protein in the urine or signs of organ injury. Possible findings include a low platelet count, abnormal kidney or liver function, severe headache, vision changes, upper-abdominal pain, or fluid in the lungs. During pregnancy, it can affect the placenta and fetal growth. It may progress to eclampsia, which is pre-eclampsia with seizures.
\medterm{Pre-Existing Diabetes And Pregnancy} Pre-existing type 1 or type 2 diabetes during pregnancy increases risks to parent and fetus, so glucose control, medicines, fetal growth, blood pressure, and eye and kidney health need coordinated care.

\medterm{Pre-Hepatic Jaundice (Haemolytic Jaundice)} Jaundice caused by excessive breakdown of red blood cells before bilirubin reaches the liver, producing increased unconjugated bilirubin.

\synonyms
Haemolytic Jaundice

\medterm{Pre-MiRNA} A pre-miRNA is an early hairpin-shaped microRNA transcript processed in the nucleus before becoming a mature regulator of gene expression.

\medterm{Pre-Synaptic} Located before a synapse, especially on the nerve ending that releases a chemical messenger. Presynaptic activity controls how much signal reaches the next nerve cell, muscle, or gland.
\medterm{Prealbumin} Another name for transthyretin, a protein made mainly by the liver that carries thyroid hormone and vitamin A in the blood. Its blood level changes with inflammation, illness, kidney disease, and nutrition, so it is not a nutrition measure by itself.

\medterm{Prebiotics} Food components, usually nondigestible fibers, that feed selected helpful bacteria in the large intestine. Sources include garlic, onions, leeks, asparagus, bananas, and whole grains; increasing them gradually may reduce gas.

\medterm{Prebiotics And Probiotics} Prebiotics are food components that feed certain helpful gut bacteria. Probiotics are live microorganisms found in some foods or supplements. Their effects depend on the specific product and health condition; they are not interchangeable treatments.

\medterm{Precancerous Condition} A precancerous condition contains abnormal changes that may progress to cancer but are not cancer themselves.

\medterm{Precancerous Polyp} A precancerous polyp is a growth with cellular changes that can become cancerous over time if not removed or monitored appropriately.

\medterm{Precaution} An action taken in advance to lower the chance of harm, infection, or error. Examples include handwashing, protective equipment, checking a medicine, using a seat belt, and avoiding a known trigger.
\medterm{Precipitate Labour} Unusually rapid labour, often defined as delivery within three hours of the onset of regular contractions.

\medterm{Precipitating Factor} A precipitating factor is an event or exposure that triggers the onset or worsening of symptoms or disease in a susceptible person.

\medterm{Preclinical} Before human testing, or describing disease changes that occur before symptoms become noticeable or cause recognized illness.

\medterm{Precocious Puberty} Development of puberty-related body changes earlier than expected. It may result from early brain signals, hormone-producing disorders, medicines, or other conditions and requires evaluation to protect growth and emotional wellbeing.

\synonyms
Early Puberty
Premature Puberty

\medterm{Preconception Care} Health care before pregnancy to improve the chance of a healthy pregnancy. It reviews medical conditions, medicines, vaccinations, genetic risks, nutrition, folic acid, substance exposure, and lifestyle, and helps plan changes before conception.

\medterm{Preconception Counseling} Healthcare advice before pregnancy to improve the health of the pregnant person and future baby. It reviews medicines, vaccines, chronic diseases, nutrition, genetic risks, harmful exposures, and recommended supplements.

\medterm{Precordial Catch Syndrome} Brief, sharp, localized chest pain near the heart that worsens with breathing and usually occurs in otherwise healthy children or adolescents.

\medterm{Precordium} The area of the front chest wall over the heart and nearby large blood vessels. A clinician examines this area by looking, feeling for the heartbeat or unusual movement, and listening for heart sounds and murmurs.

\medterm{Precursor} A substance that the body can change into another substance with a useful biological effect. For example, some vitamins, hormones, and medicines begin as precursors and are changed by enzymes in particular tissues.

\medterm{Prediabetes} Blood glucose higher than normal but below the diagnostic threshold for diabetes, usually identified by A1C, fasting glucose, or an oral glucose-tolerance test; it increases future diabetes and cardiovascular risk.

\synonyms
Diabetes Pre (Prediabetes)

\medterm{Prediabetes Syndrome} Blood glucose levels above normal but below the diabetes threshold, associated with increased future risk of type 2 diabetes and cardiovascular disease.

\medterm{Predictive Test} A test that estimates whether a person or disease is likely to respond to a particular treatment.

\medterm{Predispose} To predispose is to increase a person’s susceptibility or likelihood of developing a disease or response when additional triggers occur.

\medterm{Prednimustine} Prednimustine is a former antineoplastic drug combining a corticosteroid-related group with an alkylating component and causing immunosuppression and marrow toxicity.

\medterm{Prednisolone} Prednisolone is a corticosteroid used to reduce inflammation and immune activity; prolonged or high-dose use can cause infection, bone loss, and other effects.

\medterm{Prednisone} Prednisone is a corticosteroid used for inflammatory, allergic, and autoimmune conditions; it should often be tapered after prolonged use.

\medterm{Preeclampsia} A disorder that can develop during pregnancy, usually after 20 weeks, or during the weeks after childbirth. It involves new high blood pressure together with protein in the urine or signs of organ injury. Possible findings include a low platelet count, abnormal kidney or liver function, severe headache, vision changes, upper-abdominal pain, or fluid in the lungs. During pregnancy, it can affect the placenta and fetal growth. It may progress to eclampsia, which is preeclampsia with seizures.

\synonyms
Pregnancy-Related Hypertension


\medterm{Preeclampsia With Severe Features} A severe form of preeclampsia involving very high blood pressure or damage to organs such as the brain, liver, kidneys, lungs, or blood-clotting system. Severe headache, vision changes, breathing difficulty, or upper-abdominal pain requires emergency care.

\medterm{Preexcitation Syndrome} A heart-rhythm condition in which an extra electrical pathway activates the lower heart chambers too early. It can cause episodes of a very fast heartbeat, dizziness, or fainting and may require specialist assessment.

\medterm{Prefrontal} Relating to the anterior part of the frontal lobe involved in planning, judgment, working memory, and behavior.
\medterm{Prefrontal Cortex} The prefrontal cortex is the front part of the frontal lobe involved in planning, judgment, working memory, impulse control, and social behavior.

\medterm{Pregabalin} Pregabalin is used for certain nerve pain and seizures and may cause dizziness, sleepiness, swelling, or dependence.

\medterm{Preganglionic Fiber} A nerve fiber that carries automatic body signals from the brain or spinal cord to a relay station before they reach an organ. These signals help control the heart, blood vessels, glands, digestion, and bladder.

\medterm{Pregnancy} The period during which an embryo or fetus develops in the uterus, usually lasting about 40 weeks from the first day of the last menstrual period until birth.

\medterm{Pregnancy - Health Risks} Health risks in pregnancy vary with medical conditions, medicines, infections, age, exposures, and obstetric history; prenatal assessment identifies preventable complications and needed monitoring.

\medterm{Pregnancy - Identifying Fertile Days} Fertile days are the interval around ovulation when intercourse can lead to pregnancy; cycle timing, cervical mucus, basal temperature, and ovulation tests can provide estimates but are imperfect.

\medterm{Pregnancy And Herpes} Genital herpes during pregnancy can affect delivery planning and rarely infect the newborn; antiviral treatment and examination of lesions near labor help reduce transmission risk.

\medterm{Pregnancy And The Flu} Influenza during pregnancy can be more severe and affect the fetus; vaccination is recommended in pregnancy, and suspected illness should receive prompt clinical advice about antiviral treatment.

\medterm{Pregnancy And Travel} Travel during pregnancy depends on gestational age, health, destination infections, transport, and access to care; prevention includes vaccines as appropriate, food and water safety, movement, and medical planning.

\medterm{Pregnancy And Work} Work during pregnancy is usually possible but may require adjustments for heavy lifting, prolonged standing, heat, chemicals, radiation, infection, night work, or other occupational hazards.

\medterm{Pregnancy at Age 35 or Older} Pregnancy in a person who will be 35 years or older at the estimated time of delivery. It is associated with increased risks of some chromosomal conditions and pregnancy complications; individual risk depends on overall health and pregnancy history.

\synonyms
Advanced Maternal Age


\medterm{Pregnancy Complication} A maternal, fetal, placental, or labor-related problem that disrupts an otherwise expected pregnancy, such as preeclampsia, hemorrhage, diabetes, growth restriction, or preterm labor.


\medterm{Pregnancy Ectopic (Ectopic Pregnancy)} Alternate index label for Ectopic Pregnancy.

\medterm{Pregnancy High Blood Pressure (Pregnancy: Preeclampsia And Eclampsia)} Alternate index label for Pregnancy: Preeclampsia and Eclampsia.

\medterm{Pregnancy Nutrition} Food and supplements during pregnancy that support the pregnant person and fetal growth. Needs vary, but a varied diet usually includes enough protein, folate, iron, iodine, calcium, vitamin D, and other nutrients; individual advice is important.

\medterm{Pregnancy Of Unknown Location} A positive pregnancy test without a confirmed intrauterine or ectopic pregnancy on ultrasound. Serial human chorionic gonadotropin measurements and repeat imaging help determine whether the pregnancy is early, failing, or ectopic.

\medterm{Pregnancy Outcome} Pregnancy outcome describes what happens to a pregnancy, such as live birth, miscarriage, stillbirth, preterm birth, or other maternal or fetal result.

\medterm{Pregnancy Status} Pregnancy status indicates whether a person is pregnant, not pregnant, or of unknown status and can affect medication, imaging, and treatment decisions.

\medterm{Pregnancy Test} A urine or blood test that detects human chorionic gonadotropin, a hormone made soon after an embryo attaches to the uterus. Testing too early can be falsely negative; an unexpected result may need repeat testing or medical assessment.
\medterm{Pregnancy Trimester} One of three approximate stages of pregnancy: first through 13 weeks, second from 14 through 27 weeks, and third from 28 weeks until birth. The stages help organize fetal development, screening, and care.

\medterm{Pregnancy Tubal (Ectopic Pregnancy)} Alternate index label for Ectopic Pregnancy.


\medterm{Pregnancy-Induced Hypertension} New-onset hypertension during pregnancy without the proteinuria or organ
dysfunction that defines preeclampsia. Diagnosis and treatment depend on repeated blood-pressure measurements,
gestational age, symptoms, and evidence of maternal or fetal complications.

\medterm{Pregnancy-Related Hypertension} Alternate index label; see \textit{Preeclampsia}.

\medterm{Pregnancy: Bleeding During The First Trimester} First-trimester bleeding may occur with implantation, cervical changes, miscarriage, ectopic pregnancy, or other causes. Assessment uses symptoms, examination, serial pregnancy hormone testing, and ultrasound; heavy bleeding, severe pain, faintness, or shoulder pain is urgent.

\medterm{Pregnancy: Pain Relief Options For Birth} Labor pain relief may include breathing and movement, water, massage, nitrous oxide, intravenous medicines, epidural analgesia, or spinal anesthesia. Choice depends on labor stage, maternal and fetal status, contraindications, availability, and informed preference.

\medterm{Pregnancy: Placenta Previa} Placenta previa occurs when the placenta lies over or near the cervical opening and can cause painless bright-red bleeding, especially later in pregnancy. Digital vaginal examination is avoided until location is known; recurrent bleeding requires obstetric planning and sometimes cesarean delivery.

\medterm{Pregnancy: Preeclampsia And Eclampsia} Preeclampsia is new hypertension during pregnancy with proteinuria or organ dysfunction; eclampsia is the occurrence of seizures. Severe headache, visual symptoms, upper-abdominal pain, shortness of breath, or seizure requires emergency obstetric treatment.

\synonyms
Eclampsia Pregnancy (Pregnancy: Preeclampsia And Eclampsia)
Pregnancy High Blood Pressure (Pregnancy: Preeclampsia And Eclampsia)

\medterm{Pregnane} Pregnane is a 21-carbon steroid hydrocarbon that forms the structural backbone of progesterone and several corticosteroid and progestin compounds.

\medterm{Pregnanolone} Pregnanolone is a neuroactive steroid metabolite that can modulate GABA-A receptors and influence neuronal excitability, mood, and sedation.

\medterm{Pregnant} Pregnant means carrying a developing embryo or fetus in the uterus after conception; pregnancy status affects medicines, imaging, and clinical decisions.

\medterm{Pregnenolone} Pregnenolone is a steroid hormone precursor made from cholesterol that can be converted into progesterone, corticosteroids, and sex hormones.

\medterm{Prehn Sign} A test in which lifting the painful testicle is said to relieve pain in epididymitis but not testicular torsion. This finding is unreliable and cannot safely rule out torsion. Sudden severe testicle pain, swelling, nausea, or a high-riding testicle requires urgent assessment, usually with Doppler ultrasound and immediate surgery when torsion cannot be excluded.

\synonyms
Prehn's Sign

\medterm{Prehypertension} An older term for blood pressure that is higher than ideal but below the level once used to diagnose high blood pressure. Current guidelines use different categories; repeated readings, overall heart risk, and lifestyle are considered when deciding what help is needed.

\synonyms
Elevated blood pressure

\medterm{Preimplantation Genetic Testing} Genetic testing of cells from an embryo created through in-vitro fertilization before possible transfer to the uterus, used in selected situations to identify specific inherited disorders or chromosome-number abnormalities.

\medterm{Preload} The amount of blood filling the heart’s lower chambers and stretching their muscle just before each heartbeat. It mainly depends on blood returning to the heart and influences how strongly the heart can pump.

\medterm{Preload Reduction} A decrease in the amount of blood returning to the heart, which lowers the heart’s filling pressure and workload. It may occur with fluid loss or medicines and can help some forms of heart failure.
\medterm{Premalignant} Describing cells or tissue changes that have a higher chance of becoming cancer but have not grown into nearby tissues. Not every premalignant change becomes cancer; follow-up or removal depends on its type, location, size, and risk.
\medterm{Premature} Occurring earlier than expected or before the usual developmental or pregnancy-related time. A premature birth occurs before thirty-seven completed weeks and may increase risks involving breathing, feeding, temperature, and development.
\medterm{Premature (Early) Ejaculation} A persistent or recurrent pattern of ejaculation that occurs sooner than desired,
with limited control and associated personal distress. Assessment considers the person's typical experience and sexual
context.

\synonyms
Early

\medterm{Premature Adrenarche} Early appearance of pubic or axillary hair, body odor, or mild acne from adrenal androgen production without true central puberty; evaluation excludes virilizing disorders.

\medterm{Premature Atrial Contraction} An early heartbeat that starts in the upper chambers of the heart rather than at the usual pacemaker. It may feel like a skipped or extra beat and is often harmless, but frequent episodes or fainting need assessment.

\medterm{Premature Baby} An infant born before 37 completed weeks of pregnancy, with risks related to immature lungs, brain, feeding, temperature control, immunity, and other organ systems.

\medterm{Premature Birth} Birth before 37 completed weeks of pregnancy; the earlier the birth, the greater the risks of respiratory, neurological, and developmental complications.

\medterm{Prematurity} Birth before 37 completed weeks of pregnancy. Prematurity can affect breathing, feeding, temperature control, immunity, brain development, and growth, especially when birth occurs very early.

\medterm{Premature Closure} Closing of a body opening, growth plate, or developmental connection earlier than expected. The result depends on the site and may include blockage, abnormal body shape, or restricted growth; treatment depends on the cause.

\medterm{Premature Ejaculation} Ejaculation that occurs sooner than desired with poor control and causes persistent personal or relationship distress.
\medterm{Premature Infant} An infant born before 37 completed weeks of pregnancy. Earlier birth increases the risk of breathing, feeding, temperature-control, infection, and developmental problems.
\medterm{Premature Labor} Regular uterine contractions with cervical change before 37 completed weeks of pregnancy, which can lead to preterm birth and requires assessment of membranes, infection, and fetal well-being.

\medterm{Premature Labour (Preterm Labour)} Labour that begins before 37 completed weeks of pregnancy and may result in preterm birth.
\synonyms
Preterm Labour

\medterm{Premature Menopause} Loss of ovarian function and menstruation before age 40, now often termed primary ovarian insufficiency.
\medterm{Premature Menopause (Medical Procedural Causes)} Alternate index label for Medical Procedural Causes.

\medterm{Premature Ovarian Failure} Premature ovarian failure is loss of normal ovarian function before age 40, causing irregular or absent periods, infertility, and low-estrogen symptoms; causes may be genetic, autoimmune, toxic, or unknown.

\medterm{Premature Ovarian Insufficiency} Reduced or intermittent ovarian function before age 40, causing irregular or absent periods, reduced fertility, and sometimes hot flushes, vaginal dryness, or low bone density. Causes include genetic conditions, autoimmune disease, chemotherapy, radiation, surgery, or no identifiable cause.

\medterm{Premature Puberty} Alternate index label; see \textit{Precocious puberty}.

\medterm{Premature Rupture Of Membranes} Breaking of the amniotic sac before labor begins. If it happens before 37 weeks, it is called preterm rupture; a gush or continuing leak of fluid needs prompt obstetric assessment because infection and early birth may follow.

\medterm{Premature Thelarche} Isolated breast development in a young child without other signs of puberty, often transient but assessed for growth acceleration, estrogen exposure, or true precocious puberty.

\medterm{Premature Ventricular Contraction} An early heartbeat that starts in a ventricle instead of following the heart's usual rhythm. It may feel like a skipped or extra beat and is often harmless, but frequent PVCs may occur with fainting, chest pain, or breathlessness.

\synonyms
PVC

\medterm{Premature Ventricular Contractions} Ectopic beats originating from the ventricles, appearing as wide, bizarre QRS complexes
that occur earlier than expected. PVCs are common in the general population and are
usually benign. Frequent PVCs (>10\% of total beats) may cause left ventricular
dysfunction.

\synonyms
Extrasystole (Premature Ventricular Contractions)
PVC (Premature Ventricular Contractions)

\medterm{Premature Ventricular Contractions (PVCs)} Early extra heartbeats originating in the heart's ventricles, often felt as fluttering, skipped beats, or pounding and frequently benign.

\synonyms
Extrasystole
PVCs
Ventricular Premature Contraction
VPCs

\medterm{Premenstrual} Occurring during the days before menstruation; the term commonly refers to cyclical symptoms such as mood change, breast tenderness, bloating, headache, or fatigue.

\medterm{Premenstrual Breast Changes} Premenstrual breast changes include cyclic fullness, tenderness, swelling, or lumpiness caused by hormonal variation; a persistent one-sided mass, skin change, or bloody discharge needs evaluation.

\medterm{Premenstrual Dysphoric Disorder} A severe form of premenstrual symptoms that repeatedly causes depression, irritability, anxiety, mood swings, or loss of control along with physical symptoms and major disruption of work, relationships, or daily life before menstruation.

\medterm{Premenstrual Syndrome} A recurring group of physical and emotional symptoms in the days or weeks before menstruation that improves when bleeding begins. Symptoms may include irritability, mood changes, bloating, breast tenderness, headache, tiredness, or appetite changes.

\medterm{Premolar (Bicuspid)} One of the permanent teeth between the canine teeth and molars. Premolars usually have two raised points and help tear and crush food; decay, crowding, or infection may require dental treatment.
\medterm{Prenatal} Prenatal means occurring during pregnancy before birth, including care, testing, development, counseling, and treatment of the pregnant person or fetus.

\medterm{Prenatal Care} Regular health care during pregnancy to monitor the pregnant person and developing baby. Visits may include blood-pressure and urine checks, blood tests, ultrasound, screening, vaccination, nutrition and medicine review, and advice about warning signs.




\medterm{Prenatal Development} Prenatal development is growth and differentiation of the embryo and fetus from conception to birth, influenced by genes, nutrition, maternal health, and exposures.

\medterm{Prenatal Diagnosis} Testing during pregnancy to identify a possible genetic, chromosome, or structural condition in the fetus. It may use ultrasound, blood-based screening, or diagnostic tests such as chorionic-villus sampling and amniocentesis, depending on the clinical question and stage of pregnancy.

\medterm{Prenatal Exposure} Contact of a fetus with a medicine, substance, infection, or environmental factor during pregnancy.
Effects depend on the agent, dose, timing, and fetal susceptibility.

\medterm{Prenatal Genetic Counseling} A consultation that reviews the risk of inherited or chromosomal conditions, available screening and diagnostic tests, possible results, and reproductive options before or during pregnancy.

\medterm{Prenatal Testing} Screening and diagnostic tests performed during pregnancy to assess fetal health and development or identify genetic, chromosomal, and structural conditions. Tests may include ultrasound, maternal blood tests, cell-free DNA screening, chorionic-villus sampling, and amniocentesis.

\medterm{Prenatal Vitamins} Supplements designed to provide nutrients needed during pregnancy, commonly including folic acid, iron, iodine, and vitamin D. The appropriate product and dose depend on diet, medical conditions, laboratory results, and pregnancy needs; excessive doses can be harmful.

\medterm{Prenylation} A cell process that attaches a fat-like group to a protein, helping direct where the protein works and how it interacts with other proteins. Abnormal prenylation can contribute to some inherited disorders and cancers.

\medterm{Preop} Preop is an abbreviation for preoperative, meaning relating to preparation and assessment before surgery.
\medterm{Preoperative Assessment} A medical review before surgery that checks illnesses, medicines, allergies, previous anesthesia problems, heart and lung function, and the airway. Tests are ordered when indicated, and correctable risks are addressed before anesthesia and the operation.

\medterm{Preoperative Care} Assessment and preparation before an operation. It may include reviewing illnesses and medicines, checking allergies, performing needed tests, explaining anesthesia and recovery, planning fasting and medicines, and reducing risks such as infection or blood clots.

\medterm{Preoperative Evaluation} Medical assessment before surgery to identify and reduce risks. It reviews health history, medicines, allergies, physical findings, needed tests, anesthesia concerns, and preparations that improve safety during and after surgery.

\medterm{Preoptic Area} The preoptic area is a hypothalamic region involved in temperature regulation, sleep, reproductive behavior, and hormone control.

\medterm{Prepubertal} Prepubertal means before the physical and hormonal changes of puberty have begun or progressed to reproductive maturity.

\medterm{Prepuce} A fold of skin covering and protecting the end of an external sex organ. The penile prepuce is the foreskin over the head of the penis; the clitoral prepuce is the hood over the clitoris.
\medterm{Prerenal Azotemia} A rise in kidney waste products because the kidneys receive too little effective blood flow, without initial damage to the kidney tissue itself. Dehydration, blood loss, severe infection, heart failure, or liver disease may cause it and can lead to kidney injury if untreated.

\medterm{Presbyacusis} Gradual age-related sensorineural hearing loss, often affecting high-pitched sounds first and making speech hard to understand in noise. Hearing aids and communication strategies may help.
\medterm{Presbyacusis (Presbycusis)} Age-related, usually gradual sensorineural hearing loss, particularly affecting the ability to hear high-frequency sounds and understand speech in noise.

\synonyms
Presbycusis

\medterm{Presbycusis} Bilateral, usually symmetric age-related sensorineural hearing loss, initially affecting
high frequencies and speech discrimination in noise. Management includes hearing
rehabilitation, hearing aids, communication strategies, and cochlear implantation in
selected severe cases.

\medterm{Presbyphonia} Age-related changes in the voice caused by thinning or weakening of the vocal folds and changes in the voice box. It may produce a weak, breathy, rough, or quieter voice.

\medterm{Presbyopia} The gradual loss of the eye’s ability to focus on nearby objects with age. It commonly causes difficulty reading small print or doing close work and is corrected with reading glasses, contact lenses, or other prescribed lenses.
\medterm{Preschooler Development} Preschooler development includes rapid gains in language, motor skills, thinking, self-care, play, and social-emotional regulation; milestones vary, but loss of skills warrants assessment.

\medterm{Preschooler Test Or Procedure Preparation} Preparing a preschooler for a test or procedure uses simple explanations, play, comfort objects, and caregiver support while following any fasting or medicine instructions.

\medterm{Prescription} A written or electronic instruction from an authorized healthcare prescriber for a medicine, device, or treatment. It states what is needed and how to use it, allowing a pharmacy or service to provide care safely.
\medterm{Prescription Medications} Pharmaceutical drugs that require authorization from a licensed healthcare provider
before they can be dispensed to a patient. These medications are regulated because they
carry risks that necessitate professional oversight, including potential side effects,
drug interactions, and the need for precise dosing.

\medterm{Presenile Dementia} Presenile dementia is an older term for dementia beginning before the usual older-adult age range; modern diagnosis specifies the underlying disease and age at onset.
\medterm{Presenilin} Presenilin is a component of the gamma-secretase complex that processes membrane proteins; inherited variants can cause familial early-onset Alzheimer disease.

\medterm{Presentation} The part of a fetus that is directed toward the birth canal, or the symptoms and signs with which a patient seeks care.
\medterm{Preseptal Cellulitis} An infection of the eyelid and tissues around the eye, without infection inside the eye socket. It causes redness, swelling, pain, and fever; pain with eye movement, reduced vision, or a bulging eye may indicate deeper orbital infection and needs emergency assessment.

\synonyms
Periorbital Cellulitis

\medterm{Preservation} Maintenance of tissue, cells, organs, or function by preventing deterioration.
\medterm{Preservative} A substance added to food, medicine, or cosmetics to prevent microbial growth or chemical breakdown.
\medterm{Pressor} Raising blood pressure, usually by tightening blood vessels or increasing heart output. Pressor medicines may treat dangerously low blood pressure but require monitoring because they can reduce circulation to some tissues.
\medterm{Pressoreceptor} A pressure-sensitive receptor, especially a baroreceptor in the carotid sinus or aortic arch, that detects arterial stretch and helps regulate heart rate and blood pressure.

\medterm{Pressure} Force applied over a particular area. Medical pressure measurements include blood pressure, pressure inside the skull, eye pressure, and airway pressure, and abnormal values can damage organs or impair function.
\medterm{Pressure Atrophy} Thinning or wasting of tissue caused by prolonged pressure. Pressure can compress small blood vessels, reduce oxygen delivery, and gradually damage the affected tissue; the result depends on its location and duration.


\medterm{Pressure Injury} Localized damage to the skin and/or underlying soft tissue caused by prolonged pressure,
shear, or friction, most commonly over bony prominences.

\medterm{Pressure Sore} Alternate index label; see Bedsores (pressure ulcers). A localized injury to skin and underlying tissue caused by prolonged pressure or shear, also called a pressure ulcer.
\medterm{Pressure Support Ventilation} A ventilator mode in which the machine gives extra pressure with each breath the person starts, making breathing easier and reducing the work of the breathing muscles. It is often used while reducing ventilator support.

\medterm{Pressure Ulcer} An area of skin and deeper tissue damaged by prolonged pressure, sliding, or rubbing, usually over a bony point such as the heel, hip, or tailbone. It may begin as persistent discoloration and progress to an open wound or deep tissue damage.

\medterm{Pressure Ulcer Prevention} Measures that reduce prolonged pressure and tissue injury, including repositioning, support
surfaces, skin inspection, moisture control, nutrition, and mobility.

\medterm{Pressure Urticaria} A form of chronic inducible urticaria in which sustained pressure causes localized
swelling, erythema, and pain, usually after a delay of several hours. Lesions may last 8
to 72 hours and commonly affect the hands, feet, buttocks, or areas compressed by
clothing or sitting.

\medterm{Pressure-Flow Study} A bladder test that measures pressure while the bladder fills and while a person urinates, together with urine flow. It helps identify blockage or weak bladder muscle when difficulty emptying cannot be explained by symptoms alone.

\medterm{Presynaptic Terminals} Presynaptic terminals are the ends of neurons that release neurotransmitters into synapses in response to electrical activity.

\medterm{Presyncope} A feeling that fainting is about to occur without actually losing consciousness. It may cause light-headedness, weakness, sweating, nausea, or blurred vision and can result from low blood pressure, dehydration, abnormal heart rhythm, medicines, or other illness.

\medterm{Preterm Birth} Birth before 37 completed weeks of gestation. It may follow spontaneous labor or rupture of the
membranes, or may be medically indicated because of maternal or fetal complications.

\medterm{Preterm Infant} A baby born before 37 completed weeks of pregnancy. Earlier birth increases the chance of breathing, feeding, temperature-control, infection, and developmental problems; care depends on the baby’s maturity, weight, and health.

\medterm{Preterm Labor} Regular uterine contractions causing cervical change before 37 weeks gestation. Risk
factors include prior preterm delivery, multiple gestation, infection, and cervical
insufficiency. Diagnosis requires regular contractions with documented cervical change.

\medterm{Preterm Premature Rupture Of Membranes} Breaking of the amniotic sac before labor begins and before 37 completed weeks of pregnancy. It may cause a gush or continuing leak of fluid and increases the risks of infection, early birth, and problems from too little amniotic fluid.

\medterm{Pretest Probability} The estimated probability that a patient has a disease before the result of a diagnostic
test is known. It is based on prevalence, symptoms, examination findings, risk factors,
and prior test results and should guide test selection and interpretation.

\medterm{Pretibial Myxedema} Thickened, firm skin over the shins caused by an immune reaction associated with Graves disease. The skin may look swollen, waxy, or dimpled and usually does not leave a lasting indentation when pressed.

\medterm{Pretibial Myxoedema} Thickened, waxy swelling of the skin over the shins, most often associated with Graves disease.
\medterm{Prevalence} The proportion of people in a defined population who have a disease or characteristic at a particular time or during a stated period. It includes both newly diagnosed and previously existing cases.
\medterm{Preventing Falls} Fall prevention combines strength and balance activity, vision and medication review, safe footwear, adequate lighting, handrails, removal of hazards, and appropriate assistive devices.


\medterm{Preventing Food Poisoning} Reduce foodborne illness by washing hands, separating raw and cooked foods, cooking foods thoroughly, refrigerating them promptly, using safe water, and avoiding unpasteurized or contaminated products.

\medterm{Preventing Head Injuries In Children} Children’s head injuries can be reduced with age-appropriate car seats, helmets, supervision, safe play areas, window and stair protection, and prevention of abusive shaking or impact.

\medterm{Preventing Hepatitis A} Hepatitis A prevention relies on vaccination, handwashing, safe food and water, sanitation, and post-exposure advice when close contact with an infected person occurs.

\medterm{Preventing Hepatitis B Or C} Hepatitis B prevention includes vaccination and avoiding blood or sexual exposure; hepatitis C prevention centers on sterile injection equipment, screened blood, and testing and treatment of infection.


\medterm{Preventing Pressure Injuries} Pressure-injury prevention uses regular repositioning, skin inspection, moisture control, nutrition, pressure-redistributing surfaces, and minimizing friction and shear in people with limited mobility.

\medterm{Preventing Stroke} Stroke prevention includes controlling blood pressure, diabetes and cholesterol, avoiding tobacco, treating atrial fibrillation when indicated, exercising, and addressing other vascular risks.

\medterm{Prevention} Actions that lower the chance of disease, detect problems early, prevent complications, or reduce disability. Examples include vaccination, healthy habits, screening, safety measures, controlling chronic conditions, and avoiding harmful exposures.

\medterm{Preventive Dentistry} Dental care that lowers the risk of tooth decay, gum disease, and tooth loss. It includes brushing with fluoride toothpaste, cleaning between teeth, protective sealants when appropriate, healthy eating advice, and regular dental examinations.
\medterm{Preventive Health Care} Preventive health care aims to prevent disease or detect it early through vaccination, screening, healthy behaviors, and routine medical care.

\medterm{Preventive Medicine} The medical field focused on preventing illness and promoting health through individual care, screening, vaccination, and community programs.
\medterm{Prevertebral Air} Air seen in the soft tissues in front of the spine on a neck X-ray. It can indicate injury or a hole in the swallowing or breathing passages, a deep infection, or recent surgery and may require urgent evaluation.

\medterm{Priapism} Prolonged, painful erection lasting longer than 4 hours, occurring independently of sexual desire or stimulation. Medical emergency categorized into two types: ischemic (low-flow, veno-occlusive) and non-ischemic (high-flow, arterial).
\synonyms
Prolonged erection; persistent erection
Low-Flow Priapism
Nonischemic Priapism

\medterm{Priapism (Penis Disorder)} Alternate index label for Penis Disorder.

\medterm{Prick Test} A skin test in which small amounts of suspected allergens are introduced into the skin.
\synonyms
Skin-prick test

\medterm{Prickly Heat} Miliaria, an itchy rash caused by blockage of sweat ducts.
\medterm{Primary} First in time, origin, importance, or site; in disease classification it often means arising independently rather than secondary to another cause.
\medterm{Primary Adrenal Insufficiency} Alternate index label; see \textit{Addison's disease}.

\medterm{Primary Adrenal Insufficiency (Addison's Disease)} Destruction or dysfunction of the adrenal cortex causing deficiency of cortisol and
aldosterone, presenting with fatigue, hyperpigmentation, weight loss, and salt craving.
Lifelong hormone replacement is required.

\synonyms
Addison's disease

\medterm{Primary Adrenal Lymphoma} A rare lymphoma arising primarily in the adrenal gland, most commonly diffuse large
B-cell lymphoma. It often presents bilaterally and may cause adrenal insufficiency.

\medterm{Primary Aldosteronism} Excess production of aldosterone by the adrenal glands, usually from an adrenal nodule or enlargement of both glands.
It increases sodium retention and may cause hypertension, low potassium, or suppressed
renin.

\medterm{Primary Alveolar Hypoventilation} A rare disorder of inadequate automatic breathing, often related to PHOX2B or other autonomic-control abnormalities, causing elevated carbon dioxide especially during sleep.

\medterm{Primary Amebic Meningoencephalitis} A fulminant, rapidly fatal central nervous system infection caused by the free-living thermophilic amoeba Naegleria fowleri, contracted when warm contaminated freshwater is inhaled into the nasal cavity during swimming, diving, or nasal irrigation. The amoebae penetrate the olfactory neuroepithelium and cribriform plate to reach the brain, causing acute hemorrhagic necrotizing meningoencephalitis characterized by severe frontal headache, high fever, meningismus, altered olfaction, seizures, and rapid coma, carrying a mortality rate exceeding 97\%.

\synonyms
PAM; Naegleriasis; Primary Amoebic Meningoencephalitis; Fatal Amebic Meningoencephalitis

\medterm{Primary Amyloidosis} Primary amyloidosis, usually AL amyloidosis, results from abnormal plasma-cell proteins depositing as amyloid in organs such as the heart, kidneys, nerves, or gastrointestinal tract.

\medterm{Primary And Secondary Hyperaldosteronism} Excess aldosterone production, from an adrenal source in primary disease or from increased renin stimulation in secondary disease, causing hypertension, sodium retention, and sometimes low potassium.

\medterm{Primary Angioplasty} Angioplasty used as the initial treatment to reopen an artery during an acute heart attack.
\medterm{Primary Biliary Cholangitis} A chronic autoimmune liver disease in which small intrahepatic bile ducts are
progressively damaged. It can cause cholestasis, itching, fatigue, and eventually cirrhosis; antimitochondrial antibodies
often support the diagnosis.

\medterm{Primary Biliary Cirrhosis} Alternate index label; see \textit{Primary biliary cholangitis}.

\medterm{Primary Biliary Cirrhosis Treatment} Alternate index label for PBC.

\medterm{Primary Care} The usual first place people go for healthcare. It includes prevention, check-ups, treatment of common illnesses, care of long-term conditions, basic mental-health support, and referral or coordination with specialists when needed.

\medterm{Primary Ciliary Dyskinesia} An inherited disorder in which tiny moving hairs lining the airways do not clear mucus properly. It can cause repeated sinus and lung infections, permanent airway widening, hearing problems, infertility, and organs arranged differently from usual.

\medterm{Primary Clear Cell Adenocarcinoma Of The Vulva} A rare cancer of the vulva made up of clear-looking cells. Diagnosis requires examination of a tissue sample and tests that distinguish it from cancer that has spread from the kidney or another organ.

\medterm{Primary Cutaneous B-Cell Lymphoma} A type of lymphoma that starts in the skin and has not spread elsewhere when diagnosed. It may cause one or more painless lumps or plaques; behavior and treatment depend on the subtype.

\medterm{Primary Endpoint} The main outcome chosen in advance to judge whether a treatment or other intervention works in a clinical trial. It may be a symptom, laboratory result, complication, quality-of-life measure, or survival.

\medterm{Primary Hyperaldosteronism} Excess production of the hormone aldosterone by the adrenal glands, usually from a small gland growth or enlargement of both glands. It can cause difficult-to-control high blood pressure, low potassium, muscle weakness, or abnormal heart rhythms.

\medterm{Primary Hyperoxaluria} A group of inherited disorders in which the liver makes too much oxalate. This can cause repeated kidney stones, calcium deposits, kidney damage, and eventually failure of the kidneys.

\medterm{Primary Hyperparathyroidism} Overproduction of parathyroid hormone by one or more abnormal parathyroid glands, usually from a noncancerous gland growth. It raises blood calcium and may cause kidney stones, bone loss, abdominal symptoms, thirst, tiredness, or no symptoms at all.

\medterm{Primary Hypersomnia} Excessive daytime sleepiness that is not explained by too little sleep, another sleep disorder, a medicine, substance use, or another illness. It can cause unintended sleep, long sleep, and difficulty becoming fully alert after waking.

\medterm{Primary Hypertension} Persistently high blood pressure without one identifiable underlying disease. It results from a combination of genetic, environmental, vascular, and lifestyle factors and can often be controlled with lifestyle changes and medicines.

\medterm{Primary Hypertrophic Osteoarthropathy} A rare inherited condition causing enlarged, rounded fingertips, thickened facial or scalp skin, and new bone growth along the limbs. It may cause joint pain, stiffness, swelling, and excessive sweating.

\medterm{Primary Hypoparathyroidism} Inadequate parathyroid hormone production causing low calcium and high phosphate, with possible tingling, cramps, tetany, seizures, cataracts, or chronic bone and dental effects.

\medterm{Primary Hypothyroidism} Reduced thyroid-gland production of thyroxine and triiodothyronine, usually causing high TSH and symptoms such as fatigue, cold intolerance, constipation, dry skin, and slowed heart rate.

\medterm{Primary Immune Response} The body’s first targeted immune response to a new germ or other foreign substance. It develops over days, produces antibodies and immune cells, and creates memory that allows a faster response next time.

\medterm{Primary Immunodeficiency} A group of usually inherited disorders in which part of the immune system is missing or does not work properly. They can cause repeated or unusually severe infections and may also cause autoimmune disease or inflammation.

\medterm{Primary Immunodeficiency Disease} An inherited or genetic disorder of immune development or function that causes unusual, recurrent, severe, or persistent infections and sometimes autoimmunity.

\medterm{Primary Immunodeficiency Disease PIDD} Primary immunodeficiency diseases are inherited disorders of immune development or function that cause unusual, severe, recurrent, or persistent infections and sometimes autoimmunity or inflammation. Evaluation may include blood counts, immunoglobulins, functional tests, and genetic testing.

\synonyms
Deficiency Primary Immunodeficiency (Primary Immunodeficiency Disease PIDD)

\medterm{Primary Lactase Deficiency} Reduced lactase enzyme activity developing with age, causing lactose malabsorption and possible bloating, abdominal pain, gas, or diarrhea after dairy consumption.

\medterm{Primary Lateral Sclerosis} A rare progressive nerve disorder causing increasing muscle stiffness and weakness, usually beginning in the legs. Speech, swallowing, walking, and emotional control may become affected over time.

\medterm{Primary Lymphoma Of The Brain} Primary central-nervous-system lymphoma is a lymphoma arising in the brain, spinal cord, eyes, or surrounding fluid without systemic disease at diagnosis.

\medterm{Primary Metabolism} The essential chemical activity inside cells that produces energy, builds needed substances, breaks down waste, and keeps cells alive and functioning.
\medterm{Primary Microrna} The first RNA copy made from a microRNA gene. It is processed inside and outside the cell nucleus into a shorter molecule that helps control which proteins cells make.

\medterm{Primary Myelofibrosis} A long-term bone-marrow disorder in which scar tissue replaces healthy marrow. Blood-cell production may move to the spleen or liver, causing anemia, tiredness, easy bleeding or infection, bone pain, and an enlarged spleen.

\synonyms
Agnogenic Myeloid Metaplasia

\medterm{Primary Open-Angle Glaucoma} A slowly progressive disease that damages the optic nerve, usually without pain or early symptoms. Vision is gradually lost from the sides inward; lowering pressure inside the eye can slow or prevent further damage.

\medterm{Primary Osteoporosis} Weakening and thinning of bone caused mainly by aging, loss of estrogen after menopause, or age-related changes, rather than by another disease or medicine. It increases the risk of fractures, especially in the hip, spine, and wrist.
\medterm{Primary Ovarian Insufficiency} Loss of normal ovarian activity before age 40, characterized by irregular or absent
menses, reduced estrogen, and intermittently or persistently elevated gonadotropins.
Causes include genetic, autoimmune, iatrogenic, toxic, infectious, and unknown factors.

\medterm{Primary Polycythemia} Alternate index label; see \textit{Polycythemia vera}.

\medterm{Primary Prevention} Actions taken before a disease begins to prevent it, such as vaccination, avoiding tobacco, reducing harmful exposures, using sun protection, and controlling health risks. It differs from screening, which looks for disease that may already be present.
\medterm{Primary Progressive Multiple Sclerosis} A form of multiple sclerosis that gradually worsens from the beginning, without clear attacks and recoveries. It commonly affects walking, strength, balance, bladder control, and thinking; medicines and rehabilitation may slow disability.

\synonyms
PPMS; Primary-progressive MS

\medterm{Primary Progressive Aphasia} A neurodegenerative disorder causing gradual worsening of language abilities, including speaking, naming, comprehension, reading, or writing.

\medterm{Primary Sclerosing Cholangitis} A long-term disease in which inflammation and scarring narrow the bile tubes inside and outside the liver. It can cause itching, tiredness, jaundice, repeated infections, and liver damage and is often associated with inflammatory bowel disease.

\medterm{Primary Tuberculosis} The first infection with the bacterium that causes tuberculosis. It often begins in the lungs and nearby lymph nodes; the immune system may contain it without symptoms, or it may progress to active disease.

\medterm{Primary Tumor} The original tumor, or site where a cancer began, from which cancer cells may spread to form secondary or metastatic tumors. A cancer may be described as having an unknown primary when the original site cannot be identified.


\medterm{Primidone} Primidone is an antiseizure medicine used for epilepsy and sometimes essential tremor; it can cause sleepiness, dizziness, and drug interactions.

\medterm{Primigravida} A person who is pregnant for the first time. First pregnancies may have a longer labor and require monitoring for conditions such as high blood pressure, but the term itself does not indicate a problem.

\medterm{Primipara} A person who has given birth once after the fetus could survive outside the womb, regardless of outcome.
\medterm{Primitive} Relating to an early developmental stage or a basic, less specialized form. In medicine, primitive cells may be immature and capable of developing into several more specialized cell types.
\medterm{Primitive Node} The primitive node is an embryonic organizer at the cranial end of the primitive streak that helps establish body axes and germ-layer formation.

\medterm{Primitive Streak} The primitive streak is an embryonic structure through which cells migrate during gastrulation to form the three germ layers and establish body organization.

\medterm{Primordial} Present at the earliest stage of development or existing in an original, basic form. Examples include primordial germ cells and primordial follicles.
\medterm{Primordial Prevention} Preventing the emergence of disease risk factors by addressing social, economic,
behavioural, and environmental conditions early. Examples: promoting healthy diets and
physical activity in children, reducing tobacco exposure, and creating safe environments
for exercise.


\medterm{Prion} An infectious protein folded into an abnormal shape that can make other proteins fold incorrectly. Prions cause rare, serious brain diseases that worsen over time; they are not ordinary bacteria or viruses and currently have no cure.
\medterm{Prion Disease} A group of rare, rapidly progressive, and fatal brain disorders caused by prions---infectious agents made entirely of misfolded protein that contain no genetic material (no DNA or RNA). Normally, healthy prion proteins exist in brain cells, but when a protein misfolds into an abnormal shape, it triggers a chain reaction that forces neighboring healthy proteins to misfold into the same defective shape. These clumped proteins destroy brain tissue and leave microscopic holes, giving the brain a sponge-like appearance (transmissible spongiform encephalopathies). Prion diseases develop in three ways: sporadic (occurring spontaneously in late adulthood without a known trigger, accounting for approximately 85\% of cases such as sporadic Creutzfeldt--Jakob disease), genetic or inherited (caused by mutations in the \textit{PRNP} gene passed through families, such as familial Creutzfeldt--Jakob disease, fatal familial insomnia, and Gerstmann--Sträussler--Scheinker disease), or acquired (contracted through exposure to infected tissue, such as eating beef from cattle with bovine spongiform encephalopathy or ``mad cow disease'', contaminated surgical instruments, or ritual cannibalism as in kuru). Although the incubation period can last years or decades without symptoms, once illness begins it causes rapid memory loss, progressive dementia, severe unsteadiness and balance loss (ataxia), sudden muscle jerks (myoclonus), and loss of speech and movement. Prions are extraordinarily resistant to standard heat and chemical sterilization. There is currently no cure, vaccine, or disease-halting treatment, and the conditions are universally fatal, with medical care centered on comfort and palliative support.

\synonyms
Prion Diseases; Transmissible Spongiform Encephalopathies; TSE; Prion Disorders

\medterm{Prism} A clear optical device that bends light. In eye care, prisms can measure or correct eye misalignment, reduce double vision, and be built into glasses or used during examinations.
\medterm{Privacy} Privacy is a person’s right and expectation that personal information, communications, and bodily matters are protected from unauthorized access or disclosure.

\medterm{Privacy Breach} A privacy breach is unauthorized access, use, disclosure, loss, or exposure of personal or health information.

\medterm{Private Mutation} A genetic variant confined to a family, individual, tumor, or small population rather than commonly seen in the general population; its significance requires genetic and clinical interpretation.

\medterm{PRK} Photorefractive keratectomy is laser eye surgery that reshapes the cornea to reduce short-sightedness, long-sightedness, or astigmatism. The surface layer is removed and grows back during healing; pain, haze, infection, or imperfect vision can occur.

\medterm{Pro Re Nata} A medication or treatment instruction meaning that it should be used only when a specified symptom or clinical circumstance occurs, within clearly prescribed dose and frequency limits.

\synonyms
P.R.N.; As Needed


\medterm{Probe} A blunt instrument used to examine an opening, cavity, wound, tooth, or narrow passage. Different probes can measure depth, check whether a passage is open, guide a procedure, or trace a sinus or fistula.

\medterm{Probenecid} A uricosuric medicine used to prevent gout attacks in selected patients and to alter renal excretion of some drugs.
\medterm{Probiotic} A live microorganism that may provide a health benefit when taken in sufficient amounts. Benefits depend on the specific strain and condition; a probiotic is not automatically effective or safe for every person.

\medterm{Probiotics} Probiotics are live microorganisms that may provide a health benefit when consumed in adequate amounts, although effects are strain-specific and vary by condition.

\medterm{Problem} A symptom, sign, diagnosis, or issue requiring assessment or management.
\medterm{Problems Sleeping During Pregnancy} Sleep problems during pregnancy may result from discomfort, reflux, urination, anxiety, hormonal change, or sleep apnea; safe sleep habits and clinician-guided treatment are important.

\medterm{Probucol} Probucol is a lipid-lowering drug that reduces cholesterol but may prolong the QT interval and is not widely used in current practice.

\medterm{PROC} PROC is the gene symbol for protein C, an anticoagulant protein that limits clot formation; inherited deficiency increases thrombosis risk.

\medterm{Procainamide} An antiarrhythmic medicine that can slow abnormal electrical conduction in the heart. It is used in selected rhythm disorders and requires monitoring because it can affect blood pressure, the heart rhythm, and immune or blood systems.
\medterm{Procaine} An ester local anesthetic that temporarily blocks nerve signals and reduces pain. It is now used mainly in limited or historical settings because newer anesthetics are commonly preferred and allergic reactions can occur.
\medterm{Procalcitonin} A biomarker that may rise during some bacterial infections and systemic inflammation. It can support
assessment of infection and antibiotic decisions, but it cannot by itself diagnose or
exclude bacterial infection or determine sepsis severity.

\medterm{Procarbazine} A chemotherapy medicine used in some lymphomas and brain tumors. It damages the DNA of rapidly growing cells and can cause nausea, low blood counts, infertility, or other side effects; alcohol and many medicines may interact with it.
\medterm{Procarbazine Hydrochloride} A salt form of procarbazine, a chemotherapy medicine that damages cancer cells. It is used for selected lymphomas and brain tumors and can cause low blood counts, nausea, infertility, or serious medicine and food interactions.

\medterm{Procedural Memory} Memory for how to perform learned skills, such as riding a bicycle, typing, or playing an instrument. It usually works without consciously recalling each step and differs from memory for facts or events.

\medterm{Procedural Sedation} Use of medicines to reduce anxiety, discomfort, and awareness during a medical procedure while maintaining breathing and protective reflexes. Continuous monitoring is required because sedation can unexpectedly become deeper and slow breathing.

\medterm{Procedure} A planned medical or surgical action performed to diagnose, treat, prevent, monitor, or relieve a health problem. It may involve an examination, sample collection, medicine delivery, device placement, or operation, with risks depending on the method and person.

\synonyms
Medical Procedure; Clinical Procedure; Intervention

\medterm{Procentriole} An immature structure that forms next to an existing centriole as a cell prepares to divide. It later matures into a centriole, which helps organize the fibers that separate chromosomes.

\medterm{Procerus Muscle} A small facial muscle over the bridge of the nose that draws the medial eyebrows downward and produces horizontal wrinkles between the eyebrows.

\medterm{Process} A process is an ordered series of biological, physical, chemical, or operational changes that produces a result.

\medterm{Processing Bodies} Processing bodies are cytoplasmic RNA-protein granules where messenger RNAs may be stored, decapped, degraded, or regulated.

\medterm{Prochlorperazine} A phenothiazine medicine used mainly for severe nausea or vomiting and, in selected cases, psychotic disorders. It can cause drowsiness, low blood pressure, and movement disorders.
\medterm{Prochlorperazine Overdose} Taking too much prochlorperazine can cause severe sleepiness, confusion, low blood pressure, abnormal heart rhythms, seizures, or muscle stiffness with fever. It requires immediate emergency or poison-control advice.

\medterm{Procidentia} The most severe form of uterine prolapse, in which the uterus and cervix descend through and outside the vaginal opening. It may cause pressure, bleeding, sores, difficulty urinating, or bowel problems and requires gynecologic assessment.

\synonyms
Complete Uterovaginal Prolapse; Third-Degree Uterine Prolapse


\medterm{Procollagen} An unfinished form of collagen made inside cells. After chemical changes and removal of its end sections outside the cell, it joins with other molecules to form strong collagen fibers in skin, bone, tendons, and other tissues.

\medterm{Proctalgia} Pain in or around the rectum or anal canal. Causes include fissures, hemorrhoids, abscesses, inflammation, injury, and functional muscle pain; sudden severe pain or pain with fever, bleeding, or a lump needs medical assessment.
\medterm{Proctitis} Inflammation of the rectum. It can cause a constant urge to pass stool, rectal pain, bleeding, mucus or pus, and pelvic discomfort. Causes include inflammatory bowel disease, sexually transmitted infections, radiation treatment, and diversion of stool through an opening in the abdomen. Diagnosis may require an examination, stool or swab tests, and a small tissue sample. Treatment depends on the cause.

\synonyms
Rectal Inflammation

\medterm{Proctocolitis} Inflammation of the rectum and nearby lower colon. It may cause bleeding, mucus, diarrhea, urgency, pain, or difficulty passing stool and can result from infection, inflammatory bowel disease, reduced blood flow, radiation, or other causes.

\medterm{Proctology} The part of medicine concerned with diseases of the anus and rectum. Care may include assessment and treatment of hemorrhoids, fissures, fistulas, abscesses, bowel-control problems, infections, inflammation, and cancers.
\medterm{Proctoscope} An instrument used to inspect the anal canal and rectum.
\medterm{Proctoscopy} Examination of the anal canal and rectum with a short viewing instrument. It can help investigate bleeding, pain, inflammation, ulcers, polyps, or a lump and may allow a small tissue sample to be taken.
\medterm{Proctostomy} An operation that creates an opening from the rectum through the abdominal wall or nearby skin so stool can leave the body. It may be temporary or permanent and requires care of the opening and surrounding skin.

\medterm{Prodromal} Describing early, nonspecific signs or symptoms that occur before the characteristic clinical features of a disease become apparent.

\medterm{Prodrome} An early symptom or group of symptoms that appears before a disease episode or major event. Examples include tiredness before migraine, tingling before a seizure, or fever before an infection becomes obvious.
\medterm{Prodrug} A medicine given in a form that the body changes into its active form. This design may improve absorption, delivery to a target tissue, duration of action, or tolerability.

\medterm{Productive Cough} A cough that brings up mucus or phlegm from the airways. It can occur with infections, asthma, chronic lung disease, or other conditions and does not by itself identify the cause.
\medterm{Profunda} A deep artery or vessel, commonly referring to the profunda femoris artery in the thigh. This artery supplies muscles and nearby tissues and can be important in bleeding, circulation, or vascular procedures.
\medterm{Profunda Miliaria (Heat Rash)} Alternate index label for Heat Rash.

\medterm{Progenitor} A partly developed cell that can multiply and mature into one or a limited range of specialized cell types. Progenitor cells help replace or repair tissues but have less developmental potential than stem cells.
\medterm{Progeny} Biological offspring or descendants. In genetics, the term usually means children or later descendants whose traits may be studied to understand how characteristics or disorders are inherited. In cell biology, it can mean daughter cells produced by one parent cell.
\medterm{Progeria} A rare genetic disorder in which features resembling accelerated aging begin during childhood. The usual form is Hutchinson--Gilford progeria syndrome, caused by a disease-causing variant in the \textit{LMNA} gene that produces abnormal lamin A. Most cases result from a new variant rather than inheritance from an affected parent. Children are often healthy-looking at birth and later develop slowed growth, loss of body fat, hair loss, thin aged-looking skin, joint stiffness, and characteristic facial features. Accelerated atherosclerosis can cause early cardiovascular disease, while intellectual development is usually preserved.

\synonyms
Hutchinson--Gilford Progeria Syndrome
Hutchinson--Gilford Syndrome
HGPS
\medterm{Progestational} Relating to progesterone-like effects or preparation of the uterus for pregnancy. Progestational medicines can alter menstrual bleeding, prevent ovulation, support the uterine lining, or protect it from excessive estrogen stimulation.
\medterm{Progesterone} A steroid hormone produced mainly by the corpus luteum and placenta that supports the endometrium and pregnancy.
\medterm{Progesterone Test} A blood test that measures the level of progesterone, a hormone involved in ovulation and pregnancy. It may help assess ovulation, infertility, abnormal uterine bleeding, fertility treatment, or certain pregnancy problems. Results depend on the timing of the menstrual cycle and whether the person is pregnant.

\synonyms
Serum Progesterone
Progesterone Blood Test

\medterm{Progesterone Antagonist} A progesterone antagonist blocks progesterone receptors and may be used for selected reproductive, gynecologic, or other medical purposes.

\medterm{Progesterone Supplementation} Use of progesterone during pregnancy in selected people at risk of preterm birth, especially when ultrasound shows a short cervix. Vaginal progesterone may be considered in that situation; routine use after a previous preterm birth without a short cervix is not recommended, and 17-hydroxyprogesterone caproate is no longer FDA-approved for this purpose.

\medterm{Progestin} A progestin is a synthetic progesterone-like hormone used in contraception, hormone therapy, and treatment of selected gynecologic conditions.

\medterm{Progestogen} A natural or synthetic hormone with effects like progesterone. Progestogens help regulate the menstrual cycle and pregnancy and are used in some contraceptives and hormone treatments.
\medterm{Proglottid} One segment of a tapeworm. Each mature segment contains reproductive organs and may release eggs in stool.
\medterm{Prognathism} Forward projection of the upper or lower jaw relative to the facial skeleton, affecting bite, appearance, chewing, speech, and sometimes airway function.

\medterm{Prognosis} The expected course and outcome of a disease or condition.
\medterm{Prognostic Factor} A prognostic factor is a clinical, pathologic, or biologic feature associated with the likely course, outcome, or survival of a disease.

\medterm{Prognostic Test} A test that provides information about a disease's likely course or outcome. Results may estimate risks, guide treatment choices, or monitor response, but they cannot predict an individual future with certainty.

\medterm{Progression} Worsening or advancement of a disease, such as increasing tumor size, new lesions, or spread to another site.

\medterm{Progression-Free Survival} The length of time after a defined starting point, such as diagnosis or treatment, during which a person remains alive without the cancer worsening. It is a group measure used in research and cannot predict exactly what will happen to one person.

\medterm{Progressive} Increasing, worsening, or spreading over time, or moving forward in a particular direction. A progressive disease may cause gradually increasing symptoms, but the word alone does not explain the cause, speed, or expected outcome.
\medterm{Progressive Disease} A disease that is worsening, spreading, or becoming more active over time. In cancer, progression may mean that tumors have grown, new tumors have appeared, or existing disease has clearly worsened on examination or imaging.

\medterm{Progressive Lenticular Degeneration} An older name for Wilson disease, an inherited disorder causing copper buildup. Excess copper can damage the liver, brain, eyes, kidneys, and other organs unless treatment lowers copper levels.
\medterm{Progressive Multifocal Leukoencephalopathy} A serious brain infection caused by reactivation of JC virus in people with severely weakened immune systems. It damages the brain’s nerve insulation and may cause worsening weakness, clumsiness, vision or speech problems, confusion, or seizures.

\medterm{Progressive Muscular Atrophy} A motor-neuron disorder predominantly affecting lower motor neurons and causing progressive muscle wasting and weakness.
\medterm{Progressive Supranuclear Palsy} A progressive brain disorder that causes balance problems, frequent falls, stiffness, slowed movement, speech or swallowing difficulty, and trouble moving the eyes, especially looking down. It resembles Parkinson disease but has different features and treatment needs.

\synonyms
Steele-Richardson-Olszewski Syndrome
Supranuclear Palsy

\medterm{Proinflammatory} Describing a substance, cell, or immune response that starts or increases inflammation. Such activity can help fight infection or repair injury, but excessive or prolonged inflammation can damage healthy tissue.

\medterm{Proinsulin} An unfinished hormone made by pancreatic beta cells before it is split into insulin and C-peptide. Measuring proinsulin can provide information about insulin production and certain disorders of glucose regulation.

\medterm{Projectile Vomiting} Forceful vomiting expelled with unusual intensity, which may occur with gastric outlet obstruction, raised intracranial pressure, infection, migraine, or other causes; severity and associated signs guide evaluation.

\medterm{Projection} A view or image made by directing X-rays or other imaging energy through the body toward a detector. Different projections show body structures from different angles and can help locate an abnormality.
\medterm{Prokaryote} A prokaryote is a cell or organism without a membrane-bound nucleus, including bacteria and archaea.
\medterm{Prokinetic} A medicine that promotes gastrointestinal motility and movement of contents through the digestive tract; use depends on the underlying disorder and safety profile.
\medterm{Prolactin} A hormone released mainly by the pituitary gland that supports breast development and milk production. Abnormal levels can affect periods, fertility, sexual function, pregnancy, breastfeeding, and sometimes vision because of pituitary growths.
\medterm{Prolactin Blood Test} A prolactin blood test measures the hormone prolactin and helps evaluate abnormal milk production, menstrual or fertility problems, and some pituitary disorders.

\medterm{Prolactin Physiology} Prolactin is a hormone made by the pituitary gland that supports breast development and milk production. Levels normally rise during pregnancy, breastfeeding, sleep, stress, and breast stimulation.

\medterm{Prolactin Receptor} A cell-surface receptor activated by prolactin to regulate mammary development, milk production, reproduction, immune effects, and other tissue responses.

\medterm{Prolactinoma} A usually benign pituitary tumor that produces too much prolactin. It may cause milk production, menstrual or fertility problems, low sex hormones, headache, or vision loss and is often treated with dopamine-agonist medicines.

\synonyms
Pituitary Tumor (Prolactinoma)

\medterm{Prolactinoma Treatment} Treatment may include a dopamine agonist, surgery, radiation, or observation,
depending on tumor size, symptoms, hormone levels, treatment response, and patient preferences. Dopamine agonists often
lower prolactin levels and reduce tumor size.

\medterm{Prolapse} Downward displacement or protrusion of an organ or tissue from its normal position, often through a supporting structure or body opening.
\medterm{Prolapse Of Orbital Fat} Bulging of the eye’s normal cushioning fat through a weak spot in the thin tissue around the eye. It usually appears as a soft, painless lump in an upper or lower eyelid, often with aging, but other lumps may look similar.

\medterm{Prolapse Of The Uterus (Uterine Prolapse)} Descent of the uterus into the vaginal canal due to weakening of pelvic floor supports,
most commonly the uterosacral and cardinal ligaments.

\synonyms
Uterine Prolapse

\medterm{Prolapse, Mitral Valve} Alternate index label; see \textit{Mitral valve prolapse}.

\medterm{Prolapsed Bladder} Alternate index label; see \textit{Anterior vaginal prolapse (cystocele)}.

\medterm{Prolapsed Cord} A prolapsed umbilical cord descends beside or ahead of the fetus after membrane rupture, risking cord compression and reduced fetal oxygen; it is an obstetric emergency.

\medterm{Prolapsed Disk} A spinal disk that bulges or ruptures and may compress a nearby nerve; also called a herniated disk.

\synonyms
a herniated disk.

\medterm{Prolapsed Uterus} Alternate index label; see \textit{Uterine prolapse}.

\medterm{Proliferating Trichilemmal Tumor} A usually low-grade follicular tumor arising from a trichilemmal cyst, most often on the scalp; large or recurrent lesions may show locally aggressive or malignant behavior.

\medterm{Proliferation} Increase in the number of cells or organisms through growth and division. It may be normal, part of healing, or abnormal as in uncontrolled tumor growth.
\medterm{Proliferative} Characterized by increased cell multiplication or new tissue growth. The meaning depends on context: some proliferative changes are normal or healing, while others are precancerous or cancerous.
\medterm{Proliferative Diabetic Retinopathy} An advanced form of diabetic eye disease in which poor blood supply causes fragile new blood vessels to grow on the retina or nearby structures. They can bleed, scar, detach the retina, and threaten vision; specialist treatment can reduce these risks.

\medterm{Proliferative Fasciitis} A benign, reactive fibroblastic proliferation resembling nodular fasciitis but containing distinctive ganglion-like cells; it usually grows quickly and resolves or stabilizes after local treatment.

\medterm{Proliferative Myositis} A benign reactive proliferation within skeletal muscle that may present as a rapidly enlarging mass and mimic sarcoma. Imaging and biopsy are used to establish the diagnosis.

\medterm{Proliferative Retinopathy} Advanced retinal disease in which abnormal new blood vessels grow and may bleed or impair vision.

\medterm{Proline} An amino acid used to build proteins. Its ring-shaped side chain helps give collagen and other proteins their structure.
\medterm{Prolonged Deceleration} A temporary slowing of the fetal heart rate lasting at least two minutes but less than ten minutes. It may result from reduced blood flow, excessive contractions, cord problems, or placental separation and requires prompt assessment during labor.

\medterm{Prolonged Second Stage} The pushing stage of labor lasting longer than expected. The safe duration depends on whether it is a first birth, use of epidural anesthesia, progress, and the health of the parent and fetus; the team reassesses before deciding on further care.

\medterm{Prolymphocyte Count} The number or proportion of prolymphocytes, immature-appearing lymphoid cells, in blood or marrow; an increased count may occur in selected lymphoid leukemias and requires morphology and immunophenotyping.

\medterm{Prolymphocytic Leukemia} Prolymphocytic leukemia is a rare aggressive leukemia with large abnormal lymphocytes in the blood, bone marrow, spleen, or lymph nodes.

\medterm{Promazine} A first-generation phenothiazine antipsychotic with sedative properties; it is now used infrequently in many settings and availability varies by country. It can cause drowsiness, orthostatic hypotension, anticholinergic effects, extrapyramidal symptoms, and other serious antipsychotic adverse effects.
\medterm{Promecarb} Promecarb is a carbamate insecticide that inhibits acetylcholinesterase and can cause cholinergic poisoning after substantial exposure.


\medterm{Prometaphase} Prometaphase is the stage of mitosis after nuclear-envelope breakdown when spindle microtubules attach to chromosome kinetochores.

\medterm{Promethazine} Promethazine is an antihistamine used for allergy symptoms, nausea, motion sickness, and sometimes sedation; it can cause marked sleepiness and breathing problems in young children.

\medterm{Promethazine Overdose} Too much promethazine can cause extreme drowsiness, confusion, agitation, fast heartbeat, seizures, breathing problems, or coma. Children are particularly vulnerable; significant or intentional overdoses need immediate emergency advice.

\medterm{Prominent Iris Collarette} A visible thickening or ridge of the iris at the junction of the pupillary and ciliary
zones, representing the site where the embryonic pupillary membrane was attached. It is
typically a normal anatomical variant but can be a clinical sign in conditions such as
Axenfeld-Rieger syndrome or iridocorneal endothelial syndrome.

\medterm{Promonocyte Count} The number or proportion of promonocytes, immature monocyte precursors, in blood or marrow; increased cells may support evaluation of acute or chronic myeloid leukemia.

\medterm{Promontory} A part of the body that projects outward like a small ridge or bump. Examples include the sacral promontory in the pelvis and the cochlear promontory in the inner ear; the exact meaning depends on the structure named.
\medterm{Promoter} A region of DNA near a gene where proteins attach to start copying the gene into RNA. Promoters help determine when a gene is active and how much of its protein is made.

\medterm{Promoters} DNA control regions that help start gene activity when regulatory proteins attach. They influence how much messenger RNA and protein a gene makes and are important in development, normal function, and disease.
\medterm{Promyelocytes} Immature cells in the bone-marrow pathway that normally develops into several types of white blood cells. Large numbers in blood or marrow may occur in leukemia and need examination with other blood and cell tests.

\medterm{Pronation} Turning the forearm so the palm faces down when the elbow is bent, or backward when the arm hangs at the side. During this movement, the radius rotates across the ulna.
\medterm{Pronator} A muscle that turns a limb or part of it into a pronated position. In the forearm, pronator muscles rotate the palm downward and help stabilize the elbow during movement.
\medterm{Prone} Prone means lying face downward with the front of the body against a surface; positioning affects breathing, pressure injury, and procedures.

\medterm{Prone Position} A position in which the patient lies face down, commonly providing access to the back, spine, buttocks, or posterior legs.

\medterm{Prone Positioning} Placing a person face down to improve breathing, especially during severe acute respiratory distress syndrome. It can improve oxygen levels by helping more of the lung take part in breathing, but requires careful monitoring of the airway, skin, eyes, and pressure points.

\medterm{Pronephric} Relating to the earliest temporary kidney system formed in an embryo. In humans it does not provide effective kidney function and is replaced during development by later kidney structures.

\medterm{Pronephros} The first temporary kidney structure formed in a human embryo. It is not functional as a kidney and largely disappears as the permanent kidneys develop.
\medterm{Propafenone} A heart-rhythm medicine that slows electrical signals through the heart. It is used for selected irregular rhythms, including episodes of atrial fibrillation, but can worsen or cause dangerous rhythms and is unsuitable for some people with structural heart disease.

\medterm{Propane} Propane is a flammable hydrocarbon gas used as fuel; inhalation or displacement of oxygen can cause suffocation, and liquid contact can cause cold burns.

\medterm{Propane Poisoning} Breathing concentrated propane can displace oxygen and cause dizziness, headache, confusion, unconsciousness, abnormal heart rhythms, or suffocation. Leave the area if safe, avoid flames, and call emergency services for symptoms or significant exposure.

\medterm{Propanol} Propanol is an alcohol solvent with two structural isomers; ingestion or substantial exposure can depress the nervous system and irritate tissues.

\medterm{Propanolol} A nonselective beta-adrenergic blocker used for hypertension, arrhythmias, angina, tremor, migraine prevention, and other indications.
\medterm{Propantheline} An older medicine that reduces certain nerve signals to relax the gut and decrease sweating. It can cause dry mouth, blurred vision, constipation, difficulty urinating, fast heartbeat, or confusion and is rarely used today.
\medterm{Propellant} A propellant is a gas or volatile substance that expels a product from a container or delivers medicine in an inhaler.

\medterm{Properdin} A blood protein that strengthens one part of the body’s early immune defense. It helps mark and damage certain germs; too little properdin can increase susceptibility to serious bacterial infections.

\medterm{Properdin Deficiency} A rare inherited immune-system disorder in which a complement protein is missing or low. It increases the risk of severe infections, especially meningococcal disease, and is evaluated with immune tests and a history of recurrent infections.

\medterm{Propetamphos} Propetamphos is an organophosphate insecticide that can inhibit acetylcholinesterase and cause cholinergic toxicity after exposure.
\medterm{Prophage} A prophage is bacteriophage genetic material integrated into a bacterial chromosome or maintained in a dormant state inside a bacterium.

\medterm{Prophase} Prophase is the early stage of cell division when chromosomes condense, the spindle begins forming, and the nucleolus disappears.

\medterm{Prophenofos} Prophenofos is an organophosphate pesticide whose toxicity results mainly from acetylcholinesterase inhibition and excess cholinergic activity.

\medterm{Prophylactic} Prophylactic means intended to prevent disease, infection, injury, or another unwanted outcome before it occurs.
\medterm{Prophylactic Mastectomy} Prophylactic mastectomy removes one or both breasts to greatly reduce cancer risk in selected people at very high genetic or clinical risk.

\medterm{Prophylaxis} Medicine or other preventive care used to lower risk of disease, infection, or recurrence before it occurs.

\medterm{Propionates} Salts or esters of propionic acid used in food preservation, metabolism, and some medicines; clinical effects depend on the specific compound and exposure.


\medterm{Propionibacterium} A former bacterial genus name now largely replaced by Cutibacterium; its members are common skin organisms and can cause acne or device-related infection.


\medterm{Propionic Acidemia} A rare inherited disorder in which the body cannot properly break down certain proteins and fats. Harmful acids can build up, especially during illness or fasting, causing vomiting, extreme sleepiness, poor feeding, seizures, or coma. Newborn screening can detect it; treatment uses a controlled diet, avoids fasting, and manages attacks urgently.

\synonyms
Propionic Aciduria (PA); Propionyl-CoA Carboxylase Deficiency

\medterm{Propofol} A short-acting medicine given through a vein to start or maintain anesthesia and provide procedural sedation. It acts quickly but can lower blood pressure, slow breathing, or cause dangerous reactions without monitoring.
\medterm{Propofol Infusion Syndrome} A rare, potentially fatal illness after prolonged or high-dose propofol infusion. It may cause severe blood acidity, muscle breakdown, high potassium, kidney failure, heart failure, or dangerous heart rhythms. Stopping propofol and urgent intensive care are required when it is suspected.

\medterm{Propoxur} Propoxur is a carbamate insecticide that inhibits acetylcholinesterase and may cause sweating, salivation, vomiting, breathing difficulty, or weakness after exposure.

\medterm{Propoxycaine} A local anesthetic formerly used in some topical or dental preparations; toxicity can affect the nervous system and heart.

\medterm{Propoxyphene} Propoxyphene is a former opioid analgesic withdrawn in many countries because of overdose toxicity, respiratory depression, and dangerous cardiac arrhythmias.

\medterm{Propranolol} A beta-blocking medicine that slows heart rate and reduces the body's response to adrenaline. It treats selected heart conditions, tremor, migraine prevention, anxiety symptoms, and other problems, but can worsen asthma or low blood sugar.
\medterm{Proprioception} Awareness of body position and movement without looking, based on signals from muscles, tendons, joints, skin, and balance organs. Reduced proprioception can cause clumsiness, poor coordination, falls, or difficulty walking.
\medterm{Proprioceptor} A sensory receptor that detects position, movement, or muscle stretch.
\medterm{Proptosis (Exophthalmos)} Forward bulging of one or both eyes. Causes include thyroid eye disease, infection, inflammation, bleeding, or a growth behind the eye; sudden onset requires urgent medical assessment.
\medterm{Propulsion} Movement that drives the body or another object forward. In some neurologic disorders, impaired forward movement can cause short, shuffling steps or a tendency to lean and fall forward.
\medterm{Propyl Alcohol} Propyl alcohols are flammable solvents; ingestion or substantial exposure can irritate tissues and depress the central nervous system, with toxicity depending on the isomer and dose.

\medterm{Propylene Glycol} Propylene glycol is a solvent and humectant used in medicines, foods, and cosmetics; large exposures can cause metabolic or kidney problems.

\medterm{Propylthiouracil} Propylthiouracil is an antithyroid drug that reduces thyroid-hormone synthesis and conversion of T4 to T3; liver injury and fetal risks require careful use.

\medterm{Prosencephalon} The front part of the early developing brain. It later forms the cerebral hemispheres, thalamus, hypothalamus, and related structures that control thought, sensation, movement, hormones, and other functions.

\medterm{Prosopagnosia} Difficulty or inability to recognize familiar faces despite adequate vision and general intelligence. It may be developmental or result from brain injury or disease.
\medterm{Prosoplasia} Development of a more specialized cell function or structure within a tissue. The change may occur during normal development or as an adaptation, and its medical meaning depends on the tissue and surrounding disease process.
\medterm{Prosoposchisis} A rare congenital facial cleft or fissure caused by abnormal fusion of facial processes during embryonic development, potentially affecting the nose, mouth, eyes, and feeding or breathing.

\medterm{Prostacyclin Analogs} Medications that mimic prostacyclin (PGI2), causing vasodilation and inhibiting platelet
aggregation in the pulmonary vasculature. They include epoprostenol (continuous IV),
treprostinil (subcutaneous, IV, inhaled, oral), and iloprost (inhaled). They are used
for severe PAH.

\medterm{Prostaglandin} A locally acting lipid mediator derived from arachidonic acid that influences pain, inflammation, fever, vascular tone, platelet function, gastric protection, kidney function, and uterine contraction.

\medterm{Prostaglandin Analogs} Eye-drop medicines that lower pressure inside the eye by helping fluid drain. They treat glaucoma or ocular hypertension and may cause eye redness, eyelash growth, or darkening of the iris or skin around the eye.

\medterm{Prostaglandin Degradation} Enzymatic inactivation and removal of prostaglandins, which limits the duration and intensity of their local effects on vessels, inflammation, smooth muscle, kidneys, stomach, and reproductive tissues.

\medterm{Prostaglandin E2} A prostaglandin that mediates inflammation, pain, fever, vasodilation, gastric protection, cervical ripening, and uterine contraction and is used in selected obstetric and cardiovascular treatments.

\medterm{Prostaglandin F2-Alpha} A prostaglandin that promotes smooth-muscle contraction and luteolysis; analogues are used in glaucoma, postpartum hemorrhage, labor-related care, and veterinary medicine.

\medterm{Prostaglandin Receptor} A protein on or inside a cell that recognizes a prostaglandin and starts signals affecting blood-vessel width, inflammation, pain, smooth-muscle activity, platelets, or reproduction.

\medterm{Prostaglandins} Locally acting lipid mediators involved in pain, inflammation, vascular tone, gastric protection, and reproductive physiology.
\medterm{Prostaglandins E} The E series of prostaglandins, lipid mediators that influence inflammation, vascular tone, gastric protection, kidney blood flow, cervical ripening, and uterine contraction.

\medterm{Prostaglandins F} The F series of prostaglandins, lipid mediators involved in smooth-muscle contraction, luteal function, vascular responses, inflammation, and labor-related uterine activity.

\medterm{Prostate} A gland below the bladder surrounding part of the urine tube. It adds fluid to semen and may enlarge, become inflamed, or develop cancer, causing urinary, sexual, or pelvic symptoms.
\medterm{Prostate Adenocarcinoma} The most common type of prostate cancer, arising from gland cells in the prostate. It often grows slowly but can spread, and its growth is usually stimulated by male sex hormones; grade and stage help determine treatment.

\medterm{Prostate Biopsy} Removal of small samples of prostate tissue with a needle to look for cancer under a microscope. The needle may pass through the skin between the scrotum and anus or through the rectum; bleeding, infection, and temporary difficulty urinating are possible risks.

\medterm{Prostate Brachytherapy} Radiation treatment in which small radioactive sources are placed inside or next to the prostate to destroy cancer cells. It is used for selected cancers and may cause temporary urinary or bowel symptoms and later erectile problems.


\medterm{Prostate Cancer} Cancer that forms in the prostate gland. It may grow slowly or behave aggressively and may cause no symptoms early on. Risk increases with age, family history, inherited gene changes, and certain racial or ethnic backgrounds. Prostate-specific antigen testing and sometimes a digital rectal examination may help evaluate risk, but diagnosis requires a prostate biopsy. Tumor grade and stage help guide treatment.

\medterm{Advanced Prostate Cancer} Prostate cancer that has grown beyond the prostate or spread to distant parts of the body. It may involve nearby tissues, lymph nodes, bones, or other organs.

\medterm{Early-Stage Prostate Cancer} Prostate cancer that remains confined to the prostate or nearby tissues and has not spread to distant organs. The term includes cancers with different grades and risks of progression.

\synonyms
Cancer, Prostate
Prostatic Carcinoma (Prostate Cancer)


\medterm{Prostate Cancer Screening} Prostate cancer screening may use a PSA blood test with or without a digital rectal examination after discussing possible benefits, false positives, and harms.

\medterm{Prostate Cancer Staging} Determining how far prostate cancer has grown or spread. Clinicians combine examination, PSA level, tumor grade, biopsy findings, and imaging to decide whether it is confined to the prostate, nearby, or in distant organs.

\medterm{Prostate Cancer Treatment} Care for cancer of the prostate, chosen according to cancer grade, stage, growth rate, age, and health. Options may include active surveillance, surgery, radiation, hormone-blocking treatment, or medicines for advanced disease.

\medterm{Prostate Carcinoma} Prostate carcinoma is cancer arising from the prostate gland and may grow slowly or behave aggressively.

\medterm{Prostate Disease} Any disorder affecting the prostate, including benign enlargement, inflammation, infection, or cancer. It may cause difficulty urinating, weak flow, urgency, pelvic pain, blood in urine, or no symptoms initially.

\medterm{Prostate Enlargement} Alternate index label; see \textit{Benign prostatic hyperplasia}.

\medterm{Prostate Gland} The prostate, a gland below the bladder that surrounds the urethra and contributes secretions to semen.
\medterm{Prostate Gland Enlargement} Alternate index label; see \textit{Benign prostatic hyperplasia}.

\medterm{Prostate Infection} Infection or inflammation of the prostate, usually causing pelvic or urinary symptoms and sometimes fever or bloodstream infection.


\medterm{Prostate Neoplasm} An abnormal mass of cells arising in the prostate. It may be noncancerous, such as a glandular growth, or cancerous; symptoms and treatment depend on its type, size, and spread.


\medterm{Prostate Resection - Minimally Invasive} A procedure that removes part of an enlarged prostate through the urine channel or through small openings. It aims to improve urine flow while avoiding a large incision; bleeding, infection, scarring, or leakage can occur.


\medterm{Prostate-Specific Antigen} A protein made mainly by the prostate and measured in blood. A high or rising level can occur with prostate cancer, but also with prostate enlargement, infection, urine retention, ejaculation, or recent procedures. The result is considered with age, examination, trends, and other tests rather than used alone.

\medterm{Prostate-Specific Antigen Blood Test} A PSA blood test measures prostate-specific antigen; an elevated result may reflect enlargement, inflammation, infection, or cancer and is not diagnostic by itself.

\medterm{Prostatectomy} Surgery to remove part or all of the prostate, commonly for prostate cancer or severe urinary obstruction. Possible effects include urinary leakage, erectile dysfunction, bleeding, and infertility.
\medterm{Prostatic Carcinoma (Prostate Cancer)} Alternate index label for Prostate Cancer.

\medterm{Prostatic Fluid} Prostatic fluid is a mildly alkaline secretion that forms part of semen and contains substances supporting sperm function and survival.

\medterm{Prostatic Obstruction} Blockage of urine flow caused by enlargement, inflammation, scarring, stones, or a tumor of the prostate or nearby structures.

\medterm{Prostatic Urethral Stent} A small device placed in the prostatic urethra to help maintain urine flow when the prostate obstructs it.

\medterm{Prostatism} Urinary symptoms caused by blockage or poor function at the bladder outlet, often related to an enlarged prostate. Symptoms include weak flow, straining, urgency, frequent urination, nighttime urination, or incomplete emptying.
\medterm{Prostatitis} Inflammation of the prostate gland. It may be caused by a bacterial infection or by inflammation involving the prostate, pelvic muscles, or nerves without a proven infection. Symptoms may include pain in the pelvis, area between the scrotum and anus, penis, scrotum, lower abdomen, or lower back; painful or frequent urination; difficulty emptying the bladder; painful ejaculation; and, in acute bacterial prostatitis, fever and chills. The main types are acute bacterial prostatitis, chronic bacterial prostatitis, and chronic prostatitis/chronic pelvic pain syndrome.

\medterm{Prostatitis - Bacterial} Bacterial prostatitis is infection and inflammation of the prostate, causing pelvic or urinary pain, fever, urinary symptoms, or recurrent infection.


\medterm{Prostatitis - Nonbacterial} Persistent pain or urinary symptoms in the prostate and pelvis without evidence of a bacterial infection. Symptoms may fluctuate and can involve pelvic muscles, nerves, inflammation, or bladder function; care is tailored to the person’s findings.


\medterm{Prostatodynia} Prostatodynia is an older term for pelvic or prostate-region pain without proven bacterial infection, now often considered within chronic pelvic pain syndrome.

\medterm{Prosthesis} An artificial device replacing or supporting a missing body part or function. Prostheses may replace limbs, joints, teeth, eyes, or heart valves and require fitting, care, adjustment, and monitoring.
\medterm{Prosthetic} Relating to an artificial device that replaces or supports a missing or damaged body part. Examples include artificial limbs, joints, heart valves, teeth, and blood-vessel grafts; complications can include wear, loosening, skin problems, or infection.
\medterm{Prosthetic Graft} A prosthetic graft is an artificial tube or patch used to replace or bypass a damaged blood vessel or other tubular structure; infection and thrombosis are important risks.

\medterm{Prosthetic Heart Valve} An artificial valve that replaces a diseased heart valve and keeps blood moving in the correct direction. Mechanical valves usually require lifelong blood-thinning medicine; tissue valves may wear out sooner and have different medication needs.

\medterm{Prosthetic Limb} An artificial limb used to replace some or all of a missing arm or leg. It may improve movement, balance, or daily activities and is fitted to the person’s remaining limb, strength, goals, and comfort. Training helps with safe use and skin care.

\medterm{Prosthetic Training} Rehabilitation teaching a person to use an artificial limb safely and effectively. Training includes skin and residual-limb care, putting it on, walking, balance, stairs, transfers, daily activities, and community mobility.

\medterm{Prosthetics} The design, fitting, and use of artificial replacements for missing body parts, such as arms, legs, teeth, or eyes. Successful use depends on comfort, strength, balance, skin health, goals, and rehabilitation.

\medterm{Prosthodontics} The dental specialty that replaces missing teeth and restores chewing, speech, and appearance. Treatment may use crowns, bridges, implants, dentures, or facial prostheses and is planned around oral health and remaining tissues.

\medterm{Protamine} A medicine that binds to and reverses the blood-thinning effect of heparin. It is given carefully because it can cause low blood pressure, allergic reactions, slow heart rate, or dangerous clotting.
\medterm{Protanopia} Protanopia is complete absence of red-cone function, causing a form of red-green color-vision deficiency and altered perception of red and related colors.

\medterm{Protean} Highly variable in form, appearance, or behavior; in medicine it may describe a disease with changing or diverse manifestations.

\medterm{Proteaphagy} The process by which a cell uses its recycling system to remove whole proteasomes, the structures that normally break down unwanted proteins. It helps adjust protein disposal during stress or changes in nutrition.
\medterm{Protease} An enzyme that breaks proteins into smaller pieces by cutting the chemical links between their amino acids. Proteases help digest food and control many processes inside and outside cells.
\medterm{Protease Inhibitors} Antiviral medicines that block an enzyme viruses need to cut newly made proteins into working parts. They are used for selected viral infections and may interact with many other medicines.

\medterm{Proteasome} A barrel-shaped protein machine inside cells that breaks down unwanted, damaged, or short-lived proteins. This controls cell growth and signaling and prevents harmful protein buildup.

\medterm{Proteasome Binding} The attachment of a medicine or other molecule to a cell’s protein-destroying machine. This can increase, reduce, or otherwise change the breakdown of proteins and may alter cell survival or function.

\medterm{Protecting Yourself From Cancer Scams} Protection from cancer scams includes checking claims with qualified clinicians, avoiding guaranteed cures, reviewing evidence, and being cautious about costly unproven treatments.

\medterm{Protective} Serving to prevent or reduce injury, infection, or other harm.
\medterm{Protective Agent} A substance, device, or treatment used to reduce harm to the body. Examples include sunscreen protecting skin from ultraviolet light, barrier creams protecting skin, and medicines that protect the stomach or bones during selected treatments. The appropriate agent depends on the hazard and the person’s health.


\medterm{Protegrin} A small antimicrobial peptide, usually derived from pig immune cells, that can disrupt microbial membranes.
\medterm{Protein} A molecule made of amino acids that performs structural, enzymatic, transport, signaling, or immune functions.
\medterm{Protein Acetylation} Addition of an acetyl group to a protein, often changing its activity, localization, stability, or interactions; histone acetylation is a major regulator of gene expression.

\medterm{Protein Binding} The temporary attachment of a medicine to proteins in the blood. The unattached portion can usually leave the bloodstream, reach tissues, and produce effects; changes in protein levels or other medicines can alter the amount available.

\medterm{Protein Biosynthesis} The way cells make proteins: they copy information from DNA into messenger RNA, then use ribosomes to join amino acids in the instructed order. The new chain folds and may be modified before it works.

\medterm{Protein C} A natural blood-clotting regulator that helps prevent excessive clotting. Low levels may be inherited or acquired and increase clot risk; starting warfarin in severe deficiency can rarely cause skin injury.

\medterm{Protein C Blood Test} A protein C blood test measures a natural anticoagulant and may help evaluate unexplained or recurrent blood clots or suspected inherited deficiency.

\medterm{Protein C Deficiency} A condition in which too little or poorly functioning protein C reduces the body’s natural control of blood clotting. It may be inherited or acquired and increases the risk of deep-vein thrombosis and pulmonary embolism.

\medterm{Protein Deficiency} Inadequate protein intake, digestion, absorption, or use, causing loss of muscle, impaired immunity and wound healing, edema, or growth failure when severe or prolonged.

\medterm{Protein Denaturation} Loss of a protein’s normal three-dimensional structure from heat, chemicals, pH, or mechanical stress, often impairing its biological function without breaking its peptide backbone.

\medterm{Protein Dephosphorylation} Removal of a phosphate group from a protein by a phosphatase, often reversing or modifying signaling, enzyme activity, localization, or cell-cycle control.

\medterm{Protein Electrophoresis} A laboratory test that separates proteins in serum or urine according to their electrical properties, helping detect abnormal protein patterns.

\medterm{Protein Electrophoresis - Serum} Serum protein electrophoresis separates blood proteins into patterns and helps evaluate abnormal antibodies, myeloma, inflammation, and some liver or kidney disorders.

\medterm{Protein Engineering} Design or modification of proteins through genetic, biochemical, or computational methods to alter activity, stability, specificity, expression, or therapeutic usefulness.

\medterm{Protein Expression} The process by which a cell uses genetic instructions to make a protein, then folds, changes, and transports that protein so it can perform its function.

\medterm{Protein Folding} The process by which a newly made protein bends into the three-dimensional shape needed for its job. Incorrect folding can stop the protein from working or cause harmful protein deposits in cells.

\medterm{Protein Glycosylation} The attachment of sugar chains to proteins. This helps proteins fold, remain stable, reach the correct part of a cell, and work properly; inherited problems with this process can affect many organs.

\medterm{Protein Import} Movement of a newly made protein into a cell structure or across a cell membrane. Address labels on the protein and transport systems guide it to the correct location, where it can fold and work properly.

\medterm{Protein In Diet} Dietary protein supplies amino acids needed for muscle, enzymes, hormones, immunity, and tissue repair; requirements vary with age, health, and kidney function.

\medterm{Protein In Urine (Proteinuria)} An abnormal amount of protein in the urine, which may be temporary or indicate kidney disease, diabetes, high blood pressure, infection, or another condition.

\medterm{Protein Isoform} One of several related protein forms produced from different genes, alternative splicing, or post-translational processing, with potentially different locations or functions.

\medterm{Protein Kinase} An enzyme that adds a small chemical group to a protein and changes how it works. Protein kinases help control cell signals, metabolism, growth, division, and other functions.

\medterm{Protein Malnutrition} Poor health caused by inadequate protein or energy intake, leading to weight loss, muscle wasting, weak immunity, or impaired healing.

\medterm{Protein Metabolism} The synthesis, folding, modification, breakdown, and recycling of proteins and amino acids in cells and tissues.

\medterm{Protein Methylation} Addition of small chemical groups called methyl groups to selected parts of a protein. This can change the protein’s activity or location and can alter how cells switch genes on or off.

\medterm{Protein Overexpression} Production of an abnormally large amount of a protein, caused by gene amplification, transcriptional activation, altered degradation, or other mechanisms and potentially driving disease or serving as a biomarker.

\medterm{Protein Phosphorylation} Attachment of a small phosphate group to a protein by a cell enzyme. This acts like a reversible switch that can change the protein’s activity, location, or interactions and helps control cell signals and division.

\medterm{Protein Prenylation} Attachment of a fat-like chemical group to a protein, helping position it in a cell membrane. This placement allows some proteins to send signals or help move materials within the cell.

\medterm{Protein Purification} Separation of a desired protein from cells or a mixture using properties such as charge, size, solubility, affinity, or binding, while preserving its structure and activity.

\medterm{Protein Refolding} Return of a denatured or newly synthesized protein to its functional three-dimensional structure, sometimes requiring chaperones or controlled laboratory conditions.

\medterm{Protein S} A natural blood protein that helps activated protein C limit clotting. Too little or abnormal protein S can increase the risk of clots in deep veins or the lungs.

\medterm{Protein S Blood Test} A protein S blood test measures a natural anticoagulant and may help evaluate suspected inherited or acquired tendency to form blood clots.

\medterm{Protein S Deficiency} A condition in which the blood has too little or poorly functioning protein S, reducing natural control of clotting. It may be inherited or acquired and can increase the risk of deep-vein thrombosis or pulmonary embolism.

\medterm{Protein Secretion} Release of a protein from a cell or into a particular cell compartment. Cells use this process to send hormones, digestive enzymes, antibodies, and other proteins to where they are needed.

\medterm{Protein Sequence} The precise order of amino acids joined together to make a protein. This order helps determine the protein’s three-dimensional shape, location, interactions, and biological job.

\medterm{Protein Splicing} A process in which a newly made protein removes an internal section and joins the remaining sections together. The finished protein may then fold into the form needed for its normal function.

\medterm{Protein Subunit} One protein component of a larger protein assembly. Several subunits may join to create the final structure and allow the assembly to perform its biological function.

\medterm{Protein Synthesis} The process by which a cell uses messenger RNA instructions to join amino acids into a protein. The resulting protein folds into a working shape and may be changed or moved before performing its function.

\medterm{Protein Translocation} Movement of a newly synthesized or mature protein across a membrane or to a different cellular compartment, directed by targeting signals and transport machinery.

\medterm{Protein Transport} Cellular movement of proteins between compartments, across membranes, or outside the cell, ensuring that each protein reaches its functional location.

\medterm{Protein Unfolding} Loss of a protein’s native three-dimensional structure, caused by heat, chemicals, mutations, or cellular stress, which can impair function or promote aggregation.

\medterm{Protein-Energy Malnutrition} Severe lack of calories, protein, or both. It can cause weight loss, muscle wasting, weakness, swelling, poor growth, weak immunity, and slow healing; nutrition must often be restored gradually and carefully.

\medterm{Protein-Losing Enteropathy} Excessive loss of plasma proteins through the gastrointestinal tract, causing low albumin, swelling, diarrhea, or nutritional deficiency.

\medterm{Proteinaceous Plug} A blockage made primarily of protein-rich material, mucus, cellular debris, or secretions that can obstruct an airway, duct, catheter, or other passage.

\medterm{Proteinuria} An unusually high amount of protein in urine. It can result from kidney disease, diabetes, high blood pressure, infection, exercise, fever, or pregnancy and may require repeat testing to find the cause.
\medterm{Proteobacteria} A large group of bacteria containing many gram-negative species, including Escherichia, Salmonella, Pseudomonas, and Vibrio. Some are harmless or helpful, while others cause intestinal, urinary, lung, wound, or bloodstream infections.

\medterm{Proteogenomics} Integration of genomic, transcriptomic, and proteomic data to identify proteins, discover coding variants, characterize tumors or microbes, and connect DNA changes with protein function.

\medterm{Proteoglycan} A large molecule made of a protein with many attached sugar chains. It helps connective tissues and cartilage hold water, resist pressure, maintain structure, and send signals between cells.

\medterm{Proteolipid Protein} A protein containing a substantial fat component, especially important in the myelin sheath around nerve fibers.

\medterm{Proteolipids} Proteins tightly associated with lipid molecules or lipid membranes, important in myelin, mitochondrial membranes, ion transport, and other membrane structures.

\medterm{Proteolysis} Enzymatic breakdown of proteins into smaller peptides or amino acids. It is needed for digestion, cell regulation, tissue remodeling, and removal of damaged proteins.
\medterm{Proteolytic} Relating to the breakdown of proteins by cutting the links between their amino acids. Proteolytic enzymes help digest food, activate clotting and immune proteins, remove damaged proteins, and control cell death; medicines can block some of these enzymes.
\medterm{Proteome} The complete set of proteins produced by a cell, tissue, organism, or biological condition at a particular time.

\medterm{Proteomic Profile} A proteomic profile is a pattern of proteins measured in a biological sample to study disease, treatment response, or biological processes.

\medterm{Proteomics} The large-scale study of proteins, including their abundance, structure, modifications, interactions, and functions. Large-scale study of proteins, including their amounts, structures, changes, interactions, and roles in cells or disease.
\medterm{Proteus} A group of Gram-negative bacteria that can cause urinary, wound, bloodstream, or other infections. Proteus urinary infections may produce urease, raising urine pH and encouraging certain kidney stones.
\medterm{Proteus Hauseri} A gram-negative bacterium in the Proteus group that is rarely isolated from human clinical specimens.
\medterm{Proteus Mirabilis} A gram-negative bacterium that commonly causes urinary-tract infection and can produce urease-related kidney or bladder stones.

\medterm{Proteus Penneri} A gram-negative Proteus species that can rarely cause urinary, wound, bloodstream, or other opportunistic infections.
\medterm{Proteus Syndrome} A rare mosaic overgrowth disorder, usually caused by AKT1 variants, producing asymmetric enlargement of bones, skin, or other tissues and an increased risk of tumors.

\medterm{Proteus Vulgaris} A gram-negative bacterium that can cause urinary, wound, bloodstream, and other infections and often produces urease.

\medterm{Prothionamide} A thioamide antimycobacterial drug used as part of selected multidrug regimens for drug-resistant tuberculosis; gastrointestinal, hepatic, neurologic, and thyroid adverse effects require monitoring.

\medterm{Prothrombin} A blood protein made with the help of vitamin K. When a blood vessel is injured, it is changed into thrombin, which helps form a stable blood clot.
\medterm{Prothrombin Deficiency} A rare bleeding disorder caused by too little or poorly functioning prothrombin, a protein needed for blood clotting. It can cause easy bruising, nosebleeds, heavy menstrual bleeding, or prolonged bleeding after injury or surgery.

\medterm{Prothrombin Time} A blood test measuring how quickly a sample forms a clot. It is often reported as the INR and helps assess clotting disorders, liver disease, vitamin K deficiency, and the effect of warfarin.

\synonyms
PT

\medterm{Proto-Oncogene} A normal gene that helps control cell growth, division, or survival. If it becomes abnormally active because of a genetic change or extra copies, it can act as an oncogene and contribute to cancer.

\medterm{Protocol} A written plan describing the steps, timing, responsibilities, and safety rules for a study, examination, or treatment. Protocols improve consistency, protect participants, guide documentation, and define how results or complications are handled.

\medterm{Proton Pump Inhibitors} Proton pump inhibitors suppress gastric acid secretion and treat reflux, ulcers, and other acid-related disorders; prolonged or inappropriate use should be periodically reviewed.

\medterm{Proton Therapy} Radiation treatment using high-energy protons to damage cancer cells. The beam can deliver much of its dose at a selected depth, potentially reducing radiation to nearby healthy tissue; suitability depends on the tumor and treatment plan.

\medterm{Protozoa} A traditional group of single-celled organisms with nuclei. Some cause human disease, including malaria, giardiasis, and amoebiasis; modern biological classification places them in several different groups.

\medterm{Protozoal Infection} Infection caused by a disease-producing protozoan parasite. Protozoal infections may affect the intestine, blood, brain, eyes, skin, or genital tract and include malaria, amoebiasis, giardiasis, leishmaniasis, toxoplasmosis, trichomoniasis, and trypanosomiasis; transmission and treatment depend on the organism and site of disease.

\synonyms
Protozoal Infections

\medterm{Protractile} Able to move forward or outward from its usual position. The tongue and lower jaw are examples of body structures that can be actively protruded.
\medterm{Protractor} An instrument or muscle that moves a body part forward or outward. In anatomy, a protractor muscle may move the jaw, shoulder blade, or another structure away from its usual position.
\medterm{Protrusion} A part that projects outward or a disc or tissue displaced beyond its normal boundary.
\medterm{Protuberance} A protuberance is a structure or area that projects outward from a surface, such as a normal bone prominence, swelling, or mass.

\medterm{Provisional Diagnosis} A temporary working diagnosis made from available evidence before all test results are known, and subject to change.
\medterm{Provisional Tic Disorder} A disorder with motor or vocal tics lasting less than one year, beginning in childhood and not better explained by another condition.

\medterm{Provocation} An action or stimulus that brings on a symptom, response, or test finding. Clinicians may use controlled provocation during examination, but it should be stopped if it causes dangerous pain or other harm.
\medterm{Proximal} Proximal means nearer to the point of attachment or origin; in a limb, the shoulder or hip end is proximal to the hand or foot.
\medterm{Proximal Humerus Fractures} Breaks near the upper end of the upper-arm bone, close to the shoulder. They may cause pain, swelling, bruising, and difficulty moving the arm; treatment depends on displacement, joint involvement, bone quality, nerve or blood-vessel injury, and the person’s needs.

\medterm{Proximal Neuropathy Diabetic (Diabetic Neuropathy)} Alternate index label for Diabetic Neuropathy.

\medterm{Proximal Radioulnar Joint} The joint near the elbow where the two forearm bones rotate around each other. It allows the palm to turn up or down; injury or dislocation can limit rotation and cause pain.

\medterm{Proximal Renal Tubular Acidosis} Proximal renal tubular acidosis occurs when the kidney cannot reclaim enough bicarbonate, causing excess acid in the blood and often low potassium.

\medterm{Proximal Tubule} The first part of the kidney’s microscopic filtering tube after the filtering unit. It returns much of the filtered water, salt, bicarbonate, sugar, amino acids, and phosphate to the blood.

\medterm{Proximal Tubule Reabsorption} The kidney’s return of much of the filtered water and useful substances, including salt, bicarbonate, glucose, amino acids, and phosphate, from the first tubule back into the blood.

\medterm{Prune Belly Syndrome} A condition present at birth, usually in boys, involving weak or absent abdominal muscles, undescended testicles, and abnormalities of the urinary tract. Kidney problems, repeated infections, breathing difficulty, and impaired growth may occur.

\medterm{Prurigo} A group of intensely itchy skin disorders with papules or nodules caused by inflammation and scratching.
\medterm{Prurigo Nodularis} A long-lasting skin disorder with intensely itchy, firm bumps called nodules. Repeated scratching or rubbing causes and maintains an itch--scratch cycle; treatment aims to reduce itching, heal the skin, and prevent new nodules.
\medterm{Pruritic} Itchy or causing an urge to scratch; pruritus can result from skin inflammation, allergy, infection, systemic disease, medications, or neuropathic mechanisms.

\medterm{Pruritic Eruption} An itchy skin eruption that may result from allergic, inflammatory, infectious, drug-related, or autoimmune disease; diagnosis depends on its morphology, distribution, timing, and associated findings.

\medterm{Pruritis Ani} Alternate index label; see \textit{Anal itching}.

\medterm{Pruritus} An unpleasant itchy feeling that creates the desire to scratch. It may result from dry skin, allergy, infection, medicines, liver or kidney disease, nerve problems, or other conditions; severe or widespread itching needs assessment.
\medterm{Pruritus (Itch)} Alternate index label for Itch.

\medterm{Pruritus Ani} Persistent itching around the anus with multiple potential causes including hemorrhoids,
fissures, infections such as pinworms and candidiasis, dermatologic conditions including
psoriasis and lichen sclerosus, and idiopathic causes.

\medterm{Pruritus In CKD} Persistent itching associated with advanced chronic kidney disease, often worse in people receiving dialysis. It may occur without a visible rash and can disturb sleep and quality of life; treatment is individualized after other causes are assessed.

\medterm{Prussian Blue} A blue pigment used in laboratory staining and as an antidote for selected radioactive cesium or thallium poisoning.
\medterm{PSA} Prostate-specific antigen, a protein made mainly by prostate cells and measured in blood. A raised level may occur with enlargement, inflammation, infection, or cancer and does not diagnose cancer alone.

\synonyms
Prostate-specific antigen

\medterm{PSA Velocity} The rate at which prostate-specific antigen concentration changes over time. It may provide additional information in prostate evaluation but should not be used alone to diagnose cancer.

\medterm{Psammoma} A tumor or lesion containing rounded, layered calcium deposits called psammoma bodies. These deposits can occur in several tumors and are identified by microscopic examination rather than symptoms alone.
\medterm{Pseudarthrosis} Failure of a broken bone to heal, leaving persistent movement or a false joint at the fracture site. It can cause pain, weakness, deformity, and loss of function and may require surgery.
\medterm{Pseudoaneurysm} A weak area or leak in an artery in which blood collects outside the normal artery but remains connected to it. It may follow injury, surgery, infection, or a catheter procedure and can enlarge, burst, or send clots downstream.

\medterm{Pseudoangiomatous Stromal Hyperplasia} A benign breast stromal proliferation forming slit-like spaces within dense collagen that may mimic vascular channels. It can be incidental or present as a palpable or imaging-detected mass.

\medterm{Pseudobulbar Affect} Sudden, difficult-to-control episodes of laughing or crying that do not match how a person feels or the situation. It can occur after stroke or with other brain and nerve diseases and is different from depression.
\synonyms
Pathological Laughter And Crying
PBA

\medterm{Pseudoclubbing} Nail-tip broadening or flattening that resembles clubbing but does not have the characteristic change in the angle between the nail and surrounding skin. It may occur with some bone, joint, skin, or inherited conditions and needs assessment of the underlying cause.

\medterm{Pseudocyesis} A condition in which a person has pregnancy-like symptoms, such as missed periods, abdominal enlargement, nausea, or breast changes, without an actual pregnancy. A pregnancy test and examination distinguish it from pregnancy and identify other possible causes.
\medterm{Pseudodementia} Memory and thinking complaints caused by depression or another potentially reversible problem rather than progressive brain-cell loss. Careful assessment is needed because depression and true dementia can occur together or resemble each other.
\medterm{Pseudoephedrine} An oral decongestant that shrinks swollen blood vessels in the nose and improves blockage from colds or sinus problems. It may raise blood pressure, cause sleeplessness or palpitations, and interact with some medicines.

\medterm{Pseudogout} A joint disease caused by calcium pyrophosphate crystals collecting in cartilage. It can cause sudden pain, swelling, warmth, and stiffness, often in a knee or wrist, and may also produce longer-term joint symptoms.

\synonyms
Arthritis Pseudogout (Pseudogout)

\medterm{Pseudogout Flare} A sudden attack of joint pain and swelling caused by calcium pyrophosphate crystals. The knee or wrist is often affected; the joint may be red and warm, and fluid testing can confirm the diagnosis.

\medterm{Pseudohypertrophic Muscular Dystrophy} An older name for Duchenne or related muscular dystrophy with enlarged-appearing muscles caused by fat and connective tissue.
\medterm{Pseudohypoparathyroidism} A group of inherited disorders in which the body does not respond properly to parathyroid hormone. This can cause low calcium, high phosphate, muscle cramps, seizures, and sometimes short fingers or other physical features.

\medterm{Pseudomelanosis Duodeni} A usually harmless dark discoloration of the first part of the small intestine seen during endoscopy. It results from pigment-containing material in the lining and may be associated with chronic illness, medicines, or mineral deposits.

\medterm{Pseudomembranous Colitis} Severe inflammation of the colon, usually caused by overgrowth of Clostridioides difficile after antibiotics. It can cause frequent watery diarrhea, fever, abdominal pain, dehydration, and serious complications.

\synonyms
C. Difficile Colitis
Colitis, Pseudomembranous

\medterm{Pseudomonas} A group of Gram-negative bacteria that can cause opportunistic infections, especially in damaged tissues or people with weakened immunity.
\medterm{Pseudomonas Aeruginosa} A gram-negative bacterium that can cause pneumonia, bloodstream, urinary, wound, and ear infections, especially in hospitalized or immunocompromised people; it often has antibiotic resistance.

\medterm{Pseudomonas Aeruginosa Infection} Infection caused by Pseudomonas aeruginosa, an opportunistic gram-negative bacterium
that can affect the lungs, wounds, urinary tract, blood, or other sites.

\medterm{Pseudomonas Aeruginosa Resistance} The ability of Pseudomonas aeruginosa to survive medicines that would normally kill it. Resistance may be natural or develop through genetic changes, so laboratory testing is needed to choose an effective antibiotic.

\medterm{Pseudomyxoma Peritonei} A rare condition in which jelly-like tumor material gradually fills the abdomen, usually after a mucus-producing appendix tumor ruptures. It can cause increasing abdominal size, pain, bowel blockage, or poor nutrition and requires specialist treatment.

\medterm{Pseudophakia} Having an artificial lens inside the eye after the natural lens has been removed, usually during cataract surgery. The replacement lens focuses light, but later clouding of its supporting capsule or other eye problems can still affect vision.

\medterm{Pseudorabies} Aujeszky disease, a herpesvirus infection mainly affecting pigs and other animals. People are rarely infected, but exposed humans may develop serious nervous-system disease; it is different from rabies.
\medterm{Pseudothrombocytosis Due To Red Cell Fragments} A falsely high platelet count caused when automated equipment mistakes very small broken red-cell fragments for platelets. A blood-smear examination can distinguish this error from a genuinely high platelet count.

\medterm{Pseudotumor Cerebri} An older name for idiopathic intracranial hypertension, a condition in which pressure around the brain rises without a tumor. It can cause daily headache, brief dimming of vision, whooshing in the ears, and swelling of the optic nerves that may threaten sight.

\medterm{Pseudoxanthoma Elasticum} A rare inherited disorder in which elastic tissue becomes hardened and damaged. It can cause yellowish or loose skin, bleeding or vision loss from eye blood vessels, and narrowing or bleeding in arteries.

\synonyms
PXE (Pseudoxanthoma Elasticum)

\medterm{Psittacosis} An infection caused by Chlamydia psittaci, usually acquired by breathing dust from infected birds or their droppings. It commonly causes fever, headache, cough, and pneumonia and requires appropriate antibiotic treatment.
\medterm{PSMA} A protein found in large amounts on many prostate cancer cells and in smaller amounts in some normal tissues. Special tracers can use it to show prostate cancer on scans or deliver radiation to cancer cells.

\medterm{Psoas} A large muscle running from the lower spine to the thigh bone. It helps bend the hip, lift the leg, maintain upright posture, and stabilize the spine during movement.
\medterm{Psoas Abscess} A collection of pus in or around the psoas muscle, caused by infection spreading through the blood or from nearby tissues.
It may cause back or hip pain, fever, difficulty walking, or few symptoms and usually
requires imaging and targeted treatment.

\medterm{Psoriasis} A long-term immune-mediated skin disease in which skin cells are produced too quickly, causing inflammation and a buildup of skin. It commonly causes raised, thick, dry, scaly patches that may be red, purple, brown, darker, or lighter than nearby skin. The patches often affect the scalp, elbows, knees, trunk, palms, or soles, and the nails may become pitted, ridged, thick, or separated from the nail bed. Symptoms may flare and then improve. Genetic and environmental factors can contribute. Psoriasis is not contagious and may occur with psoriatic arthritis, an inflammatory disease of the joints and tendon attachments.


\medterm{Psoriatic Arthritis} An inflammatory arthritis associated with psoriasis of the skin or nails. It may affect
the spine, peripheral joints, digits, or tendon insertions, and diagnosis uses the clinical pattern together with
appropriate laboratory and imaging findings.

\synonyms
Arthritis Psoriatic (Psoriatic Arthritis)

\medterm{PSP} PSP commonly means progressive supranuclear palsy, a neurodegenerative disorder causing early postural instability, vertical gaze impairment, rigidity, and cognitive change.

\medterm{Psychiatric} Relating to the assessment, diagnosis, treatment, or prevention of disorders of mood, thought, perception, behavior, or emotional functioning.

\medterm{Psychiatric Emergency} A mental-health condition involving an immediate and serious risk to life, safety, basic functioning, or the safety of others. Examples include imminent suicide risk, severe psychosis, dangerous agitation, catatonia, or severe substance intoxication or withdrawal. It is a subset of mental-health crises distinguished by the immediacy and severity of risk.

\medterm{Psychiatrist} A medical doctor who assesses, diagnoses, and treats mental-health conditions. A psychiatrist may provide talking therapy, prescribe medicines, coordinate other treatments, and assess how mental or physical illness affects safety and daily life.

\medterm{Psychiatry} The medical specialty concerned with mental health and the diagnosis and treatment of mental disorders.
\medterm{Psycho-Oncology} The clinical field addressing psychological, social, behavioral, and communication needs of people with cancer and their families during diagnosis, treatment, survivorship, and end-of-life care.

\medterm{Psychoacoustics} The study of how people perceive and interpret sound, including loudness, pitch, masking, localization, speech, and hearing thresholds.

\medterm{Psychoactive} Affecting mood, thought, perception, behavior, alertness, or other mental processes. Psychoactive substances include prescribed medicines, alcohol, and recreational drugs.
\medterm{Psychoactive Substance} A substance that alters mood, perception, cognition, behavior, or consciousness.

\medterm{Psychoanalysis} A psychological theory and talking treatment emphasizing unconscious thoughts, feelings, memories, and conflicts. Traditional treatment explores patterns in relationships and behavior, while modern approaches vary in methods, goals, and evidence.
\medterm{Psychodermatology} The area of care that considers links between skin conditions and emotional health. Stress can worsen problems such as eczema, psoriasis, acne, or hives, while visible skin disease can cause anxiety or low mood; both aspects may need treatment.

\medterm{Psychodidae} A family of small hairy flies that includes sand flies; some species transmit pathogens such as Leishmania and therefore have public-health importance.

\medterm{Psychodrama} A therapy using guided acting to explore emotions, relationships, and behavior in a supportive setting.
\medterm{Psychodynamic Therapy} Talking therapy that explores how past experiences, relationships, feelings, and patterns outside a person’s immediate awareness may influence current distress and behavior. The therapist and patient use the therapeutic relationship to develop insight and healthier ways of coping.

\medterm{Psychogenic} Caused or strongly influenced by psychological factors such as stress, trauma, or emotional conflict. The symptoms are real and not deliberately produced, and physical causes should still be considered.

\medterm{Psychogenic Erectile Dysfunction} Difficulty getting or keeping an erection mainly related to stress, anxiety, depression, relationship concerns, or other emotional factors. Physical causes can coexist, so medical assessment should consider health, medicines, and psychological wellbeing.

\medterm{Psychogenic Non-Epileptic Seizures} Episodes that resemble epileptic seizures but are not caused by the abnormal electrical discharges of epilepsy. They are real and involuntary and may involve shaking, stiffening, unresponsiveness, or altered movement; video-EEG can help establish the diagnosis.

\medterm{Psycholinguistics} The study of how the brain and mind produce, understand, acquire, and remember language, including speech, reading, grammar, and meaning.

\medterm{Psychological} Relating to thoughts, feelings, behavior, or the way these affect health. Psychological factors can influence pain, sleep, coping, treatment choices, and recovery, but describing a factor as psychological does not mean symptoms are imagined.

\medterm{Psychological Trauma} Emotional or psychological effects that may follow an event or ongoing situation involving actual or threatened death, serious injury, violence, abuse, or another overwhelming experience. Responses vary and may include fear, intrusive memories, avoidance, emotional numbness, sleep problems, or changes in mood and behavior.
\medterm{Psychological Adjustment} The process of adapting thoughts, emotions, and behavior to stress, illness, disability, loss, developmental change, or other life circumstances.

\medterm{Psychological Dependence} A pattern in which a person feels a strong emotional or mental need for a substance or behavior, often with craving, preoccupation, or distress when it is unavailable.

\medterm{Psychological Support} Emotional and practical help for people coping with illness, disability, pain, or difficult treatment. It may include counseling, coping skills, support groups, and assessment or treatment of anxiety, depression, and adjustment problems.

\medterm{Psychological Transfer} The process in which feelings, expectations, or relationship patterns from earlier experiences become directed toward someone in the present, especially a therapist. Recognizing this pattern can help understand distress and relationships during therapy.

\medterm{Psychologist, Clinical} A licensed mental-health professional who assesses and treats psychological, emotional, behavioral, and cognitive problems using psychotherapy and other evidence-based interventions within local scope of practice.

\medterm{Psychology} The scientific study of behavior, emotion, learning, and mental processes. Psychological care may include assessment, counseling, and behavioral therapies.

\medterm{Psychometrics} The design and use of tests that measure mental abilities, feelings, behavior, or personality. It also examines whether a test gives consistent results and truly measures the quality it claims to measure.
\medterm{Psychomotor} Relating to the connection between mental activity and physical movement, including the speed, coordination, and purposeful control of actions.
\medterm{Psychomotor Agitation} Restless, excessive, or purposeless movement associated with inner tension, anxiety, mania, delirium, depression, substance effects, or medication reactions.

\medterm{Psychomotor Disorder} A condition affecting the speed, starting, coordination, or expression of movement and mental activity. A person may move unusually slowly or quickly, seem restless, or have difficulty organizing purposeful actions.

\medterm{Psychomotor Performance} The ability to coordinate mental processing with physical movement, assessed through tasks involving speed, precision, reaction time, balance, or sequencing.

\medterm{Psychoneuroimmunology} The study of how thoughts, emotions, the brain, hormones, and the immune system influence one another in health and illness. It examines links involving stress, inflammation, infection, and recovery.

\medterm{Psychopathic} Describing severe, persistent antisocial traits such as limited empathy, deceitfulness, and disregard for others; it is not a precise diagnosis.
\medterm{Psychopathology} The study and description of abnormal mental processes, symptoms, behaviors, and mechanisms underlying psychiatric disorders. The study or description of abnormal mental states and behavior.
\medterm{Psychopharmacology} The study of how medicines affect mood, behavior, thinking, sleep, and other mental functions in people.
\medterm{Psychophysics} The study of how physical stimuli, such as light, sound, pressure, or heat, are experienced by a person. It measures detection thresholds, ability to distinguish stimuli, and perceived strength.

\medterm{Psychophysiology} The study of links between psychological processes and measurable bodily functions such as heart rate, skin conductance, muscle activity, and brain signals.

\medterm{Psychosexual} Relating to how psychological factors interact with sexual development, feelings, behavior, relationships, or function. Mental health, past experiences, body image, culture, medicines, and physical conditions can all influence sexuality.
\medterm{Psychosexual Development} Development of sexual identity, interest, relationships, and behavior across the lifespan as influenced by biology, cognition, attachment, culture, and experience.

\medterm{Psychosis} A state in which a person loses contact with reality. It may involve strongly held false beliefs, hearing or seeing things others do not, confused speech, or markedly disorganized behavior and can result from mental, medical, neurologic, or substance-related causes.
\medterm{Psychosis ICU (ICU Psychosis)} Alternate index label for ICU Psychosis.

\medterm{Psychosocial Rehabilitation} Support that helps a person with mental illness build skills, relationships, confidence, and independence. It may include daily-living training, education, work support, peer support, therapy, and help taking part in family and community life.

\medterm{Psychosocial Support} Practical, emotional, and social help for people facing illness or difficult treatment. It may include counseling, support groups, help with work or money worries, family support, and care for anxiety, depression, fear, loneliness, or grief.

\medterm{Psychosomatic} Relating to physical symptoms influenced by psychological processes or stress.

\medterm{Psychosurgery} Surgery intended to alter brain circuits for severe, treatment-resistant psychiatric illness; it is rarely used and highly specialized.
\medterm{Psychotherapy} Structured treatment using therapeutic communication and psychological techniques to reduce distress, change unhelpful patterns, and improve coping or relationships. Common approaches include cognitive behavioral and interpersonal therapies.
\medterm{Psychotic} Describing psychosis, a syndrome involving impaired reality testing. It may include delusions, hallucinations, disorganized thinking or speech, or markedly disorganized behavior. Psychosis can occur with schizophrenia-spectrum disorders, severe mood disorders, substances or medicines, delirium, neurologic disease, or other medical conditions.
\medterm{Psychotic Disorder} Alternate index label for Psychotic Disorders.

\medterm{Psychotic Disorder Brief (Brief Psychotic Disorder)} Alternate index label for Brief Psychotic Disorder.

\medterm{Psychotic Disorders} A group of mental disorders characterized by impaired reality testing, such as delusions, hallucinations, disorganized thought or speech, or markedly disorganized behavior. The group includes schizophrenia-spectrum disorders and other conditions in which psychosis is a defining feature; psychotic symptoms can also occur with mood disorders, substances, medicines, neurologic disease, or other medical conditions.


\medterm{PTCA} Percutaneous transluminal coronary angioplasty, a procedure using a small balloon to open a narrowed heart artery. A stent is often placed afterward to help keep the artery open and restore blood flow.

\synonyms
Percutaneous coronary intervention

\medterm{PTCL} Alternate index label; see \textit{Peripheral T-cell lymphoma}.

\medterm{PTE} Alternate name; see \textit{Pulmonary Thromboendarterectomy}.

\medterm{Pteridine} A nitrogen-containing chemical structure found in folate vitamins, natural pigments, and other active body compounds. Pteridine-containing molecules participate in cell growth, chemical reactions, color production, and inherited metabolic disorders.

\medterm{Pterygium} A raised, often triangular growth of the clear covering of the eye that extends onto the cornea. Long-term sun, wind, dust, or dry-eye exposure may contribute; it can cause irritation, redness, blurred vision, or astigmatism.
\medterm{PTH-Related Peptide} A hormone-like substance sometimes made in excess by cancer cells. It acts partly like parathyroid hormone and can raise blood calcium, causing thirst, constipation, weakness, confusion, dehydration, or kidney problems.

\medterm{Pthirus Pubis} The insect commonly called the pubic louse. It lives mainly in coarse genital hair and causes intense itching and visible lice or eggs; it spreads mainly through close body or sexual contact.

\medterm{Ptilosis} An abnormal hair pattern or hair disorder, including excess growth, unusual direction, or loss. The term is nonspecific, so the body site, appearance, timing, and underlying cause must be assessed.
\medterm{PTLD} Post-transplant lymphoproliferative disorder, an abnormal overgrowth of immune cells after organ or stem-cell transplantation. It ranges from a reversible response to lymphoma and may cause fever, swollen lymph nodes, organ problems, or unexplained weight loss.

\medterm{Ptosis} Drooping of an eyelid or another body structure. Eyelid ptosis may result from muscle, nerve, tendon, aging, injury, or a neurologic disorder.
\medterm{Ptosis - Infants And Children} Ptosis in an infant or child is drooping of the upper eyelid, which may obstruct vision or cause amblyopia and can reflect congenital, muscular, nerve, or other disease.

\medterm{PTSD} Alternate index label for Post-Traumatic Stress Disorder.

\medterm{Ptyalin} An older name for salivary amylase, the enzyme that begins starch digestion in the mouth.
\medterm{Ptyalism (Sialorrhoea)} Excessive saliva or drooling. It may result from increased saliva production or difficulty swallowing, and can occur with pregnancy, mouth irritation, reflux, medicines, or disorders affecting the brain, nerves, or muscles.
\medterm{Pubertal Failure} Puberty that does not begin or stops progressing as expected. Causes may involve the brain’s hormone centers, pituitary gland, testes or ovaries, chromosomes, chronic illness, inadequate nutrition, excessive exercise, or other factors.

\medterm{Puberty} The stage when hormone changes cause the body to mature sexually and become capable of reproduction. It includes growth, development of breasts or testes and genitals, body hair, voice changes, and menstrual periods or sperm production.
\medterm{Puberty In Boys} Puberty in boys is the hormone-driven development of testes, penis, body hair, voice, muscle, and fertility, usually progressing through predictable physical stages.

\medterm{Puberty In Girls} Puberty in girls is the hormone-driven development of breasts, body hair, growth, menstrual cycles, and reproductive maturity, usually progressing through predictable stages.

\medterm{Pubes} The area in front of the pelvis over the pubic bone, or the coarse hair that grows there during puberty. Changes in pubic hair can help assess sexual development, but must be interpreted with other physical findings.
\medterm{Pubic} Relating to the pubis, the front part of the pelvis, or the area covered by hair above the genitals.
\medterm{Pubic Lice} Pubic lice are parasitic insects infesting coarse body hair, causing intense itching and visible nits; treatment and assessment of close or sexual contacts are needed.

\medterm{Pubic Lice (Crabs)} Infestation of coarse body hair by lice, usually causing intense itching and requiring treatment and contact precautions.

\synonyms
Crabs

\medterm{Pubis} The front part of the pelvic bone, meeting its opposite side at the pubic symphysis.
\medterm{Public Assistance} Government or community support providing eligible people with food, housing, income, healthcare, disability, or other essential services.

\medterm{Public Exposure} Radiation received by other people from a patient who has undergone nuclear medicine treatment. The amount depends on the radioactive substance and treatment; the medical team gives specific instructions to limit exposure when needed.

\medterm{Public Health} The science and practice of protecting and improving the health of communities through
education, policy-making, and research for disease and injury prevention.

\medterm{Public Health Emergency} A serious event threatening the health of many people, such as a pandemic, disaster, or deliberate biological attack.



\medterm{Puborectalis Muscle} A loop-shaped pelvic muscle that surrounds the lower rectum and helps control bowel movements. At rest it pulls the rectum forward to make a bend that helps hold stool; relaxing it helps stool pass.

\medterm{Pudendal} Relating to the external genitalia, anus, or perineum. The pudendal nerve carries sensation from these areas and helps control several muscles involved in continence and sexual function.
\medterm{Pudendal Arteries} Arteries supplying blood to the external genital and perineal region.

\medterm{Pudendal Block} Injection of local anesthetic near the pudendal nerve to numb the lower vagina, vulva, anus, and perineum. It may reduce pain during late labor, assisted delivery, or repair of a birth-related tear.

\medterm{Pudendal Nerve} The nerve carrying sensation and motor signals to much of the perineum and pelvic floor.
\medterm{Puerperal} Relating to the period after childbirth, usually the first six weeks. During this time the body recovers from pregnancy and birth, the womb becomes smaller, bleeding gradually settles, and problems such as infection, heavy bleeding, or blood clots may occur.
\medterm{Puerperal Disorder} A physical or mental health disorder occurring during the postpartum period, including infection, bleeding, hypertension, depression, or psychosis.

\medterm{Puerperal Fever} Fever occurring after childbirth, often caused by infection of the uterus, birth canal, urinary tract, breast, wound, or bloodstream. It requires prompt assessment because untreated infection can spread and become life-threatening.
\medterm{Puerperal Infection} An infection after childbirth, commonly involving the uterus, surgical wound, urinary tract, breast, or bloodstream, often causing fever, pain, or foul discharge.

\medterm{Puerperal Psychosis} A rare, severe mental-health emergency that begins soon after childbirth. Rapid mood changes, confusion, false beliefs, hallucinations, or disorganized behavior can develop quickly and may threaten the safety of the parent or baby, requiring immediate assessment.

\medterm{Puerperal Sepsis} A serious infection of the reproductive tract during labor or the six weeks after childbirth. Fever, pelvic pain, foul-smelling discharge, weakness, or fast breathing require urgent medical assessment and antibiotic treatment.

\synonyms
Postpartum Sepsis; Puerperal Septicemia

\medterm{Puerperium} The period after childbirth, usually about six weeks, when the body recovers from pregnancy and birth. The uterus shrinks, bleeding settles, milk production develops, and health problems may still arise.
\medterm{Puffy Eyes} Alternate index label; see \textit{Bags under eyes}.

\medterm{Pugilistic Stance} A flexed posture of the arms, hands, and legs sometimes seen in bodies exposed to intense heat. It results from heat-related muscle and tendon contraction, not from purposeful movement, and does not by itself show when or how death occurred.

\medterm{Pugilistica Dementia (Dementia)} Alternate index label for Dementia.

\medterm{Pulled Hamstring (Hamstring Injury)} Alternate index label for Hamstring Injury.

\medterm{Pulling A Patient Up In Bed} A safe repositioning technique using adequate staff, friction-reducing equipment, and coordinated movement to move a patient toward the head of the bed while preventing falls, shear, and staff injury.

\medterm{Pulmonary} Relating to the lungs or the blood vessels and airways that support breathing. Pulmonary disease may involve infection, asthma, chronic obstructive lung disease, high pressure in lung vessels, blood clots, fluid, scarring, or cancer.

\medterm{Pulmonary Actinomycosis} Pulmonary actinomycosis is a chronic bacterial infection of the lungs caused by Actinomyces and may form mass-like lesions or spread across tissue planes.

\medterm{Pulmonary Alveolar Hemorrhage} Bleeding into the lung’s air sacs. It may cause coughing up blood, breathlessness, low oxygen, anemia, and widespread lung shadows and can result from immune disease, infection, medicines, injury, or a clotting problem.

\medterm{Pulmonary Alveolar Microlithiasis} A rare inherited lung disease in which tiny mineral deposits build up inside the air sacs. It may cause no symptoms at first, but later can lead to cough, breathlessness, low oxygen, and progressive lung scarring.

\medterm{Pulmonary Alveolar Proteinosis} A rare lung disease in which a fatty material builds up inside the air sacs and interferes with oxygen exchange. It can cause gradually worsening breathlessness, cough, tiredness, and low oxygen and may be immune-related or secondary to another illness.

\medterm{Pulmonary Angiography} Contrast imaging of the pulmonary arteries, usually by catheter. It can diagnose pulmonary vascular abnormalities but is now used less often for pulmonary embolism because CT pulmonary angiography is less invasive.

\medterm{Pulmonary Arterial Hypertension} High blood pressure in the arteries carrying blood from the heart to the lungs because these vessels become narrowed or damaged. It can cause breathlessness, tiredness, chest discomfort, dizziness, or fainting and may strain the right side of the heart.

\medterm{Pulmonary Arteries} The two large blood vessels carrying oxygen-poor blood from the right side of the heart to the lungs. In the lungs, the blood releases carbon dioxide and picks up oxygen.

\synonyms
pulmonary artery

\medterm{Pulmonary Arteriovenous Fistula} An abnormal direct connection between an artery and vein in the lung that lets blood bypass normal oxygen exchange. It can cause low oxygen, breathlessness, blue skin, stroke-causing clots, or rarely a brain abscess.

\medterm{Pulmonary Arteriovenous Malformation} An abnormal direct connection between an artery and vein in the lung. Blood bypasses normal oxygen exchange, which can cause low oxygen, breathlessness, coughing blood, stroke, or a brain abscess; some cases are linked to an inherited blood-vessel disorder.

\medterm{Pulmonary Artery} The large blood vessel carrying oxygen-poor blood from the right side of the heart to the lungs. There, blood releases carbon dioxide, receives oxygen, and returns through pulmonary veins.
\medterm{Pulmonary Artery Catheter} A thin tube passed through the right side of the heart into a lung artery to measure pressures, blood flow, and heart performance. It is used selectively in severe heart or lung disease and critical illness because insertion has risks.

\medterm{Pulmonary Artery Pulsatility Index} A measurement calculated during right-heart catheterization from pulmonary-artery pressure and right-atrial pressure. A low result may indicate weak right-ventricle pumping and a higher risk of right-sided heart failure in seriously ill patients.

\synonyms
PAPi; PAP Index

\medterm{Pulmonary Aspergilloma} A pulmonary aspergilloma is a ball of Aspergillus fungus and debris growing in a pre-existing lung cavity, sometimes causing coughing of blood.

\medterm{Pulmonary Atresia} A heart defect present from birth in which the valve or passage from the right ventricle to the pulmonary artery does not open normally. Blood cannot reach the lungs through the usual route, so newborns may become blue or severely unwell and need urgent specialist care.

\synonyms
Tetralogy of Fallot With Pulmonary Atresia

\medterm{Pulmonary Blastoma} A rare malignant lung tumor composed of primitive epithelial and mesenchymal tissue, usually treated with specialist oncology care.

\medterm{Pulmonary Capillary Wedge Pressure} A pressure measured with a special catheter in a small lung artery to estimate pressure on the left side of the heart. It can help assess heart failure or fluid overload in critically ill patients, but breathing support and other conditions affect its meaning.

\medterm{Pulmonary Circuit} The circulation in which blood travels from the right ventricle to the lungs for gas exchange and returns to the left atrium.

\medterm{Pulmonary Circulation} The blood-flow circuit from the right ventricle through the pulmonary arteries and lung capillaries to the left atrium, where blood releases carbon dioxide and gains oxygen.

\medterm{Pulmonary Compliance} How easily the lungs expand when pressure changes during breathing. Low compliance means the lungs are stiff and harder to inflate; high compliance can mean they expand easily but may not recoil effectively.

\medterm{Pulmonary Contusion} Bruising of the lung after a blow or crush injury to the chest. Bleeding and swelling inside the lung can reduce oxygen exchange and worsen during the first day or two; treatment supports breathing and addresses associated injuries.

\medterm{Pulmonary Edema} A buildup of fluid in the lungs that can make breathing difficult. It is often caused by
heart failure but may also result from kidney problems, lung injury, medicines, or high altitude.
Sudden severe breathing difficulty is an emergency.

\synonyms
Edema, Pulmonary
Edema Pulmonary (Pulmonary Edema)

\medterm{Pulmonary Embolism} A blockage of a pulmonary artery, usually caused by a blood clot that traveled from a deep vein. It may cause sudden breathlessness, chest pain that worsens with breathing, rapid heartbeat, coughing blood, dizziness, or fainting. A large embolism can severely strain the heart or cause cardiac arrest and requires emergency assessment.

\synonyms
Blood Clot In The Lung (Pulmonary Embolism)
Embolism Pulmonary (Pulmonary Embolism)

\medterm{Pulmonary Embolism Treatment} Treatment for a blood clot in the lung usually uses anticoagulant medicine to stop growth and prevent new clots. Severe or unstable cases may need clot-dissolving medicine or a procedure; choice and duration depend on the person’s risks and cause.

\medterm{Pulmonary Embolus} A blockage of a pulmonary artery, usually by a blood clot that traveled from a deep vein, causing sudden breathlessness, chest pain, low oxygen, or collapse.

\medterm{Pulmonary Emphysema} Alternate index label; see \textit{Emphysema}.

\medterm{Pulmonary Fibrosis} A group of lung diseases in which scar tissue forms in or around the air sacs and other lung tissue. The scarring makes the lungs thick and stiff, so they cannot expand normally and oxygen has more difficulty moving into the blood. It may cause gradually worsening shortness of breath, a dry cough, tiredness, and reduced ability to exercise. Causes include autoimmune disease, some medicines, workplace dust or chemicals, radiation, infection, or an unknown cause; the unknown-cause form is called idiopathic pulmonary fibrosis.

\medterm{Pulmonary Fistula} An abnormal passage between the lung or an airway and another structure, such as the space around the lung, skin, or swallowing tube. It can allow air or infected fluid to travel where it should not.

\medterm{Pulmonary Function Tests} Breathing tests that measure how much air the lungs hold, how quickly air can be moved, and how well oxygen passes into the blood. They help assess asthma, chronic obstructive lung disease, scarring, and other breathing problems. Results depend on effort, technique, age, and body size.

\medterm{Pulmonary Hemorrhage} Bleeding into the airways or lung tissue, which may cause coughing blood, low oxygen, anemia, or respiratory failure.

\medterm{Pulmonary Hemosiderosis} A condition in which repeated bleeding into the lung’s air sacs leaves iron deposits behind. It may cause cough, breathlessness, coughing blood, anemia, and low oxygen and can result from immune disease or other causes.

\medterm{Pulmonary Hypertension} Abnormally high blood pressure in the blood vessels that carry blood from the right side of the heart to the lungs. Narrowed, blocked, or damaged vessels make the right side of the heart work harder to pump blood and may eventually weaken it. Pulmonary hypertension may occur on its own or result from left-sided heart disease, lung disease or low oxygen, long-term blood clots, or other conditions affecting the blood vessels. Symptoms may include breathlessness, tiredness, chest discomfort, dizziness, fainting, or swelling from heart strain.

\synonyms
PH (Pulmonary Hypertension)

\medterm{Pulmonary Hypertension - At Home} Home care for pulmonary hypertension includes taking prescribed medicines, monitoring symptoms, conserving energy, following oxygen advice, and seeking urgent help for severe breathlessness or chest pain.

\medterm{Pulmonary Hypertension Classification} A clinical grouping based on cause: disease of the pulmonary arteries, left-sided heart disease, lung disease or low oxygen, chronic blood clots, or several other mechanisms. Identifying the group guides testing and treatment because medicines useful for one group may be harmful in another.

\medterm{Pulmonary Infarction} Death of a small area of lung tissue after its blood supply is blocked, usually by a pulmonary embolism. It may cause sharp pain with breathing, coughing up blood, breathlessness, or a shadow on imaging.

\medterm{Pulmonary Infection} An infection of the lungs or lower respiratory tract caused by bacteria, viruses, fungi, or other organisms.

\medterm{Pulmonary Langerhans Cell Histiocytosis} A rare lung disease strongly linked to cigarette smoking, in which unusual immune cells damage small airways and form lung cysts. It can cause cough, breathlessness, chest pain, or a collapsed lung. Stopping smoking is essential and may stabilize the disease; specialist treatment is needed if it progresses.

\synonyms
PLCH; Histiocytosis X of the Lung; Pulmonary Eosinophilic Granuloma

\medterm{Pulmonary Lobes} The anatomical divisions of the lungs separated by interlobar fissures. The right lung
has superior, middle, and inferior lobes; the left lung has superior and inferior lobes
and a lingula.

\medterm{Pulmonary Mass} A space-occupying lesion in or near the lung, which may represent infection, inflammation, benign growth, or cancer and requires appropriate imaging assessment.

\medterm{Pulmonary Nocardiosis} A lung infection caused by Nocardia bacteria, usually affecting people with weakened immunity and sometimes spreading to the brain or skin.

\medterm{Pulmonary Nodule} A rounded opacity in the lung, usually smaller than 3 cm, caused by scar, infection, inflammation, or tumor and assessed according to size and risk factors.

\medterm{Pulmonary-Renal Syndrome} A group of diseases that cause inflammation or bleeding in the lungs together with inflammation or injury in the kidneys. It is often associated with small-vessel vasculitis or another autoimmune disease.

\medterm{Pulmonary Oedema} Fluid collecting in lung tissue or air sacs, making breathing difficult and reducing oxygen transfer. It often results from heart failure but can also follow infection, toxins, kidney disease, or severe injury.
\medterm{Pulmonary Reactance} A measure of how the stretch and movement of the lungs and chest wall affect airflow during breathing. It is obtained during specialized tests and can help describe stiffness or altered movement of the respiratory system.
\medterm{Pulmonary Rehabilitation} A supervised, multidisciplinary program combining exercise training, education,
behavioral support, and self-management for chronic respiratory disease. It improves
exercise capacity, dyspnea, functional status, and quality of life.

\medterm{Pulmonary Sequestration} A birth-related area of lung tissue that does not connect normally to the airways and receives blood from an unusual artery. It may cause repeated lung infections, cough, or breathing problems, although some cases cause no symptoms.

\medterm{Pulmonary Stenosis} Narrowing that obstructs blood leaving the right side of the heart on its way to the lungs, usually at the pulmonary valve. It is often present from birth and may cause a heart murmur, tiredness, breathlessness, or fainting.
\medterm{Pulmonary Surfactant} A slippery mixture made by cells lining the lung’s air sacs that reduces surface tension and helps keep the sacs open. Too little in premature babies can cause serious breathing difficulty.

\medterm{Pulmonary Thromboendarterectomy} Major surgery to remove old clots and scar tissue from the arteries of the lungs. It may improve or cure long-term clot-related high blood pressure in the lungs in selected people, but it carries serious surgical risks.

\synonyms
PTE; Pulmonary Endarterectomy; PEA (Pulmonary Endarterectomy)

\medterm{Pulmonary Trunk} The large vessel leaving the right ventricle that divides into the right and left pulmonary arteries.

\medterm{Pulmonary Tuberculosis} Tuberculosis infection involving the lungs, caused by Mycobacterium tuberculosis. It may cause prolonged cough, fever, night sweats, weight loss, or coughing blood and can spread through the air.
\medterm{Pulmonary Valve} A one-way valve between the right lower heart chamber and the artery leading to the lungs. It opens to let blood leave the heart and closes to prevent backward flow; narrowing or leakage can strain the right side of the heart.

\medterm{Pulmonary Valve Regurgitation} Leakage through the pulmonary valve that allows blood to flow back into the right lower heart chamber as the heart relaxes. Mild leakage may cause no symptoms; severe leakage can enlarge and weaken the right side of the heart.

\synonyms
Pulmonic regurgitation; Pulmonary valve insufficiency


\medterm{Pulmonary Valve Insufficiency} Leakage of blood backward through the pulmonary valve into the right lower heart chamber while the heart relaxes. Mild leakage may cause no symptoms, but severe disease can enlarge the heart and reduce exercise ability.

\medterm{Pulmonary Vascular Resistance} The opposition to blood flow through the lung blood vessels. High resistance makes the right side of the heart pump harder and can lead to high pressure in the lungs, right-heart enlargement, and heart failure.

\medterm{Pulmonary Vasodilator Therapy} Treatment with medicines that widen blood vessels in the lungs and lower pressure on the right side of the heart. It is used for selected types of pulmonary hypertension after specialist testing because it may be harmful in other types.

\medterm{Pulmonary Vein} The four veins (two from each lung) returning oxygenated blood from the lungs to the
left atrium, a unique exception in the venous system (carrying oxygenated blood).
Superior and inferior pulmonary veins enter the left atrium posteriorly.

\medterm{Pulmonary Vein Isolation} A catheter procedure for selected people with atrial fibrillation. The clinician creates small areas of scar around the veins entering the upper heart chamber to block abnormal electrical signals. It can reduce episodes but may not cure them; bleeding, stroke, narrowing of a vein, or other complications are possible.

\synonyms
PVI; Antral Pulmonary Vein Isolation; AFib Catheter Ablation

\medterm{Pulmonary Veins} Usually four veins that carry oxygen-rich blood from the lungs to the heart's left upper chamber. They are unusual veins because they carry oxygen-rich rather than oxygen-poor blood.

\medterm{Pulmonary Veno-Occlusive Disease} A rare progressive disease in which small veins in the lungs become narrowed or blocked. Pressure then rises in the lung circulation, causing breathlessness and low oxygen; some medicines for pulmonary hypertension can trigger life-threatening fluid in the lungs.

\medterm{Pulmonary Ventilation} The movement of air into and out of the lungs. It brings oxygen into the air sacs and removes carbon dioxide; inadequate ventilation can cause low oxygen or a buildup of carbon dioxide.

\medterm{Pulmonary Ventilation/Perfusion Scan} An imaging test comparing air reaching different lung regions with blood flow through those regions. A mismatch may suggest a blood clot in the lung, while other patterns can indicate lung or blood-vessel disease.

\medterm{Pulmonary Wedge Pressure} A pressure measured through a special catheter in a small lung artery to estimate the pressure in the left side of the heart. It helps assess fluid overload or heart failure in seriously ill patients.

\medterm{Pulmonary-Artery Dissection} A rare, life-threatening tear in the wall of an artery carrying blood from the heart to the lungs. It may cause sudden chest pain, breathlessness, low blood pressure, or collapse and requires emergency specialist care.

\medterm{Pulmonic Valve Stenosis} Pulmonic valve stenosis is narrowing of the valve between the right ventricle and pulmonary artery, creating outflow obstruction and sometimes right-heart strain.

\medterm{Pulmonologist} A pulmonologist is a physician specializing in diseases of the lungs and breathing system, including asthma, COPD, infections, and sleep-related breathing disorders.

\medterm{Pulp} The soft living tissue in the center of a tooth, containing nerves and blood vessels. It can become inflamed or infected after decay, cracks, or injury and may then require root-canal treatment or removal.
\medterm{Pulp Canal} The pulp canal is the passage within a tooth root containing dental pulp, nerves, and blood vessels; infection may require root-canal treatment.

\medterm{Pulp Disease} A problem affecting the living tissue inside a tooth, including inflammation, infection, loss of blood supply, or tissue death. It may cause toothache or sensitivity and can follow decay, a crack, a filling, or injury.

\medterm{Pulp Stone} A small area of hardened mineral inside the soft center of a tooth. It often causes no symptoms and is found on dental imaging, but it may complicate root-canal treatment if it blocks the canal.

\medterm{Pulpectomy} Complete removal of infected or damaged pulp from inside a tooth, usually a primary tooth. The space is cleaned and filled with a material that can safely resorb as the adult tooth develops.

\medterm{Pulpitis} Inflammation of the dental pulp, usually caused by caries, trauma, cracks, or restorative treatment. Reversible pulpitis causes brief pain provoked by cold or sweet stimuli; irreversible pulpitis causes spontaneous or lingering pain and may progress to pulp necrosis, requiring endodontic treatment or extraction.
\medterm{Pulposus} Relating to the nucleus pulposus, the gelatinous central portion of an intervertebral disk. Relating to the soft, jelly-like center of an intervertebral disk, which helps cushion spinal movement.
\medterm{Pulpotomy} Removal of inflamed soft tissue from the crown of a baby tooth while keeping the root tissue alive. It treats suitable decay or injury and preserves the tooth until it naturally falls out.

\medterm{Pulse} A pulse is the rhythmic expansion of an artery produced by each heartbeat; its rate, rhythm, strength, and symmetry provide clues about circulation and cardiac function.
\medterm{Pulse - Bounding} A bounding pulse feels unusually strong or full and can occur with fever, anemia, hyperthyroidism, exercise, anxiety, aortic regurgitation, or high-output states.

\medterm{Pulse Examination} Checking arterial pulses for speed, regularity, strength, and equality on both sides. It helps assess blood flow, heart rhythm, circulation to limbs, and possible blockage or reduced blood supply.

\medterm{Pulse Irregular} An irregular pulse has uneven timing or varying beats and may reflect premature beats, atrial fibrillation, other arrhythmias, or measurement variation.

\medterm{Pulse Oximeter} A small device, often placed on a fingertip, that estimates blood oxygen and pulse rate using light. Movement, cold fingers, poor circulation, nail products, and device limitations can make the reading inaccurate.

\medterm{Pulse Oximetry} A painless method of estimating blood oxygen with a light-based sensor, usually on a finger, toe, or ear. The result must be interpreted with symptoms and circulation because it can be inaccurate in some situations.

\medterm{Pulse Pressure} The systolic blood pressure minus the diastolic blood pressure. A widened or narrowed value may provide clues about circulation but must be interpreted with the full blood-pressure reading and clinical context.


\medterm{Pulse Rate} Pulse rate is the number of palpable arterial beats per minute and usually approximates heart rate, although some arrhythmias produce a pulse deficit.

\medterm{Pulseless Disease} Takayasu arteritis, an inflammation of large arteries that can narrow or block blood flow. It may cause weak pulses, unequal blood pressure, limb pain, dizziness, vision problems, or organ damage.
\medterm{Pulseless Electrical Activity} A cardiac-arrest rhythm in which the heart’s electrical activity continues but the heart does not pump enough blood to produce a pulse. It requires immediate resuscitation.
\medterm{Pulsus} A term used in medicine to describe the pulse and named patterns of arterial beats. The pattern may provide clues about heart rhythm, pumping strength, blood-vessel resistance, breathing, or circulation.
\medterm{Pulsus Paradoxus} An unusually large drop in blood pressure during breathing in, usually more than 10 millimeters of mercury. It can occur when the heart is compressed by fluid, or during severe asthma, chronic lung disease, or other heart conditions.

\medterm{Pulvinar} A region at the back of the thalamus, a deep brain structure. It helps direct visual attention, combine sensory information, and communicate with areas of the cerebral cortex.

\medterm{Pumice} A porous volcanic stone used as an abrasive; excessive dust inhalation can irritate the lungs.
\medterm{Pump Failure} Severe weakening of the heart’s pumping ability, causing inadequate blood flow to organs.
\synonyms
Cardiogenic pump failure

\medterm{Pump-Oxygenator} The part of a heart-lung machine that moves blood and adds oxygen while removing carbon dioxide during some heart operations. It temporarily replaces the pumping and gas-exchange work of the heart and lungs. Use can be associated with bleeding, inflammation, clots, or organ complications.

\medterm{Punch Biopsy} A circular, usually full-thickness skin specimen obtained with a punch instrument. It is useful for inflammatory, infectious, melanocytic, and other cutaneous lesions; small wounds may be closed with a suture.

\medterm{Punctal Occlusion} Partial or complete blocking of the small tear-drain openings near the inner eyelids, often with plugs or a procedure, to keep tears on the eye surface. It may help selected people with dry eye; watering, irritation, or infection can occur.

\medterm{Punctate} Made up of, or marked by, small spots, dots, or points. The term describes appearance and does not by itself identify a diagnosis.
\medterm{Punctate Hemorrhages (Tardieu Spots)} Tiny pinpoint areas of bleeding under the skin or on internal surfaces. They can form when small blood vessels rupture during severe pressure or oxygen deprivation, but their presence alone cannot establish the cause or timing of death.

\synonyms
Tardieu Spots

\medterm{Punctate Palmoplantar Keratoderma} An inherited condition causing many small, hard areas of thickened skin on the palms and soles. It often begins in adolescence or adulthood and can make walking or using the hands painful; rarely, a skin cancer develops.

\medterm{Puncture} A small hole or wound made by a pointed object, needle, bite, or sharp instrument. A puncture may be intentional for an injection or sample, or accidental and may carry infection, bleeding, or deeper injury risks.
\medterm{Puncture Biopsy} Alternate index label for needle biopsy: removal of cells or tissue through a percutaneous puncture for microscopic, molecular, or microbiologic examination; imaging may guide the target and route.

\medterm{Puncture Wound} A wound produced when a pointed object penetrates the skin and underlying tissue, creating a small surface opening with potentially deep tissue damage and infection risk.

\medterm{Pungent} Describes a strong, sharp smell or taste that may irritate the nose, mouth, eyes, or other moist linings.
\medterm{Punicic Acid} A conjugated polyunsaturated fatty acid found mainly in pomegranate seed oil and studied for lipid and inflammatory effects; clinical supplement benefits remain uncertain.

\medterm{Punishment} Punishment is a consequence intended to reduce a behavior; in healthcare, coercive or harmful punishment is inappropriate and behavior support should prioritize safety and dignity.

\medterm{Pupil} The central opening in the colored iris through which light enters the eye. Muscles change its size to control incoming light, and unequal or poorly reactive pupils may indicate eye or nerve disease.
\medterm{Pupil - White} A white pupil, or leukocoria, may result from reflection, cataract, retinal disease, infection, or retinoblastoma and requires prompt eye assessment, especially in a child.

\medterm{Pupillary} Relating to the pupil, the opening in the iris that controls how much light enters the eye. The pupil normally becomes smaller in bright light and larger in dim light; unequal size or poor reaction can indicate an eye, nerve, or brain problem.
\medterm{Pure Red Cell Aplasia} A rare disorder in which the bone marrow stops making red blood cells while other blood cells may remain normal. It causes severe anemia, tiredness, breathlessness, and pale skin and may be inherited or acquired from infection, immune disease, cancer, or medicines.

\synonyms
PRCA; Pure red-cell aplasia

\medterm{Pure Tone Audiometry} A hearing test in which a person signals whenever they hear tones of different pitches and loudness. It measures the quietest sounds detectable through air and bone and helps identify the type and degree of hearing loss.

\medterm{Pure-Tone Audiometry} A hearing test that measures the softest tones a person can hear at selected frequencies through air and bone conduction.
The results are plotted on an audiogram and help distinguish conductive from sensorineural hearing loss.

\synonyms
Pure-tone threshold audiometry; pure-tone hearing test.

\medterm{Purgative} A strong laxative that causes rapid emptying of the bowel. It may be used for bowel preparation before some procedures, but excessive use can cause diarrhea, dehydration, and dangerous changes in body salts.
\medterm{Purge} To empty the bowel, stomach, or body of a substance; the term may also mean forceful vomiting or diarrhea.
\medterm{Purging} Deliberate behaviors intended to compensate for eating, such as self-induced vomiting or misuse of laxatives,
diuretics, or excessive exercise.

\medterm{Purging Disorder} Repeated self-induced vomiting or misuse of laxatives, diuretics, or other methods to influence weight or body shape without recurrent binge eating. It can cause dangerous salt disturbances, tooth damage, dehydration, heart-rhythm problems, and tears in the food pipe.

\medterm{Purine} A nitrogen-containing building block of DNA, RNA, and energy-carrying molecules. The body breaks down purines into uric acid, which can accumulate and contribute to gout or kidney stones.
\medterm{Purine Nucleosides} Molecules made from a purine base, such as adenine or guanine, joined to a sugar. Examples include adenosine and guanosine, which help form genetic material and participate in cell energy and signaling.

\medterm{Purine Nucleotides} Molecules containing adenine or guanine joined to a sugar and phosphate groups. They build DNA and RNA and include ATP and GTP, which store energy and help transmit signals inside cells.
\medterm{Purines} Nitrogen-containing building blocks, mainly adenine and guanine, used to make DNA, RNA, and energy-carrying molecules. The body breaks them down into uric acid, which can form crystals and contribute to gout or kidney stones.

\medterm{Purinoceptor} A protein on a cell’s surface that detects substances such as ATP, ADP, or adenosine outside the cell. It helps regulate cell communication, inflammation, blood-vessel activity, and other body responses.

\medterm{Purkinje Cell} A major nerve cell in the cerebellum, the brain region that coordinates movement and balance. It sends signals that help fine-tune posture, coordination, and learned movements.

\medterm{Purkinje Fibers} Special electrical fibers spread through the lower chambers of the heart. They rapidly carry each heartbeat’s signal so the ventricles contract together and pump blood efficiently.

\medterm{Purkinje Fibres} Specialized heart-muscle fibers that rapidly spread electrical signals through the ventricles. Their coordinated activity helps the lower chambers contract together and pump blood efficiently to the lungs and body.

\medterm{Purpura} Purple or red spots caused by blood leaking under skin or mucous membranes; they do not fade when pressed.
\medterm{Purpura Fulminans} A rapidly worsening, life-threatening condition in which widespread blood clots and clotting-system failure cause purple skin patches, bleeding, and tissue death. It can follow severe infection, loss of natural anticoagulants, or another serious illness.

\medterm{Purulent} Containing or producing pus, a thick fluid made of immune cells, dead tissue, and often germs. Purulent drainage can occur in an abscess or infected wound, but its appearance alone does not identify the organism or the treatment needed.


\medterm{Purulent Pericarditis} A serious bacterial infection of the sac around the heart, causing pus and fluid to collect there. Fever, chest pain, breathlessness, and pressure on the heart may occur; urgent drainage and antibiotic treatment are needed.

\medterm{Pus} A thick fluid made of dead cells, germs, immune cells, and damaged tissue. It commonly collects during infection and may appear in wounds, abscesses, infected organs, or other inflamed body spaces.
\medterm{Push Enteroscopy} An examination of the upper small intestine using a long flexible camera passed through the mouth. It allows clinicians to see the lining, take samples, stop bleeding, or treat selected lesions.

\medterm{Pustule} A small raised skin lesion filled with pus. Pustules can occur with acne, follicle infection, impetigo, psoriasis, or a medicine reaction; their cause is determined from their location, appearance, symptoms, and timing.
\medterm{Pustules} Pustules are small raised skin lesions filled with pus, occurring in acne, folliculitis, infections, inflammatory disorders, or drug reactions.

\medterm{Pustulosis} A skin disorder characterized by sterile or infectious pustules, which may be localized or widespread and associated with inflammation, fever, psoriasis, drug reactions, or infection.

\medterm{Putamen} A deep brain structure that helps select and smooth purposeful movements and supports learning of repeated movement habits. Damage to it can cause abnormal movements, weakness, or changes in behavior.
\medterm{Putrefaction} Decomposition of body tissues after death caused mainly by bacteria and natural enzymes. It can produce color changes, swollen tissues from gas, visible veins, a strong odor, and gradual breakdown into liquid material; the rate depends on the environment.
\medterm{Putrescine} A chemical with a strong unpleasant smell made when proteins break down. It is found in decaying material and is also produced in small amounts during normal cell processes; its presence alone does not diagnose disease.
\medterm{PUVA} A light treatment that combines a light-sensitizing medicine called psoralen with ultraviolet-A exposure. It can treat selected skin diseases but may cause burns, eye injury, premature skin aging, and increased skin-cancer risk.

\medterm{PVC} A PVC, or premature ventricular contraction, is an early heartbeat arising in a ventricle; occasional PVCs are common, while frequent or symptomatic beats need evaluation.


\medterm{PVD} Alternate index label; see \textit{Peripheral artery disease}.

\medterm{PVI} Alternate name; see \textit{Pulmonary Vein Isolation}.

\medterm{PVPS} Alternate index label; see \textit{Post-vasectomy pain syndrome}.

\medterm{Pyromania} A mental disorder involving recurrent deliberate fire-setting, mounting tension or arousal before the act, and pleasure, gratification, or relief afterward, without a primary external motive such as financial gain, revenge, concealment of a crime, or political purpose. The behavior is not better explained by conduct disorder, mania, psychosis, or substance use.

\medterm{Pyarthrosis} Infection of a joint with pus in the joint space, usually bacterial septic arthritis, causing acute pain, swelling, restricted movement, fever, and risk of cartilage destruction.

\medterm{Pycnodysostosis} A rare inherited bone disorder in which bones are unusually dense but fragile. It can cause short stature, frequent fractures, short fingers, delayed closure of skull bones, crowded teeth, and problems with the jaw.

\medterm{Pyelitis} Inflammation of the renal pelvis, the funnel-shaped part of the kidney that drains urine. It is often associated with urinary infection and may cause flank pain, fever, or painful urination.
\medterm{Pyelo} Pyelo is an informal shortening usually referring to pyelonephritis, an infection of the kidney and renal pelvis that can cause fever and flank pain.

\medterm{Pyelogram} An image of the kidney drainage area and ureters after contrast material is given. It can show blockage, stones, narrowing, leaks, or structural changes in the urinary collecting system.
\medterm{Pyelography} Imaging of the renal pelvis and ureters using contrast, by X-ray or another imaging method. It can show obstruction, stones, leaks, strictures, or structural abnormalities.
\medterm{Pyelolithotomy} Surgical removal of a kidney stone through an incision into the renal pelvis. Less invasive stone treatments are often preferred when suitable.
\medterm{Pyelonephritis} Infection and inflammation of the kidney and its urine-collecting area, usually caused by bacteria traveling upward from the bladder. Symptoms may include fever, chills, side pain, nausea, and painful urination.
\medterm{Pyelotomy} An incision into the renal pelvis, usually made to remove an obstruction or kidney stone.
\medterm{Pylephlebitis} An infected blood clot in the large vein carrying blood from the abdominal organs to the liver. It may follow appendicitis, diverticulitis, or another abdominal infection and can cause fever, abdominal pain, jaundice, and life-threatening sepsis.

\medterm{Pyloric Antrum} The pyloric antrum is the lower stomach region that grinds food and regulates its passage toward the pyloric canal and duodenum.

\medterm{Pyloric Sphincter} A ring of muscle at the outlet of the stomach. It opens in a controlled way to let partly digested food enter the first part of the small intestine and helps prevent the intestine's contents from flowing backward.

\synonyms
Pylorus

\medterm{Pyloric Stenosis} Narrowing of the pyloric outlet, classically causing projectile vomiting in infants or gastric-outlet symptoms in adults.

\synonyms
Stenosis, Pyloric

\medterm{Pyloric Stenosis In Infants} Infantile pyloric stenosis is narrowing of the stomach outlet that causes forceful, non-bilious vomiting in young infants and usually requires surgical treatment.

\medterm{Pyloroplasty} Surgery to widen the pylorus, the outlet of the stomach, to improve passage of stomach contents. It may be used with treatment for obstruction or selected gastric-motility disorders.
\medterm{Pylorus} The distal part of the stomach leading into the duodenum.
\medterm{Pyoderma} A bacterial infection of the skin that can produce pus-filled spots, crusts, open sores, or deeper inflammation. Staphylococcus and Streptococcus are common causes; the infection may spread and require medical treatment, especially with fever, pain, or rapidly worsening redness.
\medterm{Pyoderma Gangrenosum} A rare inflammatory skin disease that causes very painful ulcers, often with a dark purple edge. It may occur with bowel inflammation, arthritis, or blood disorders; injury can make it worse, so treatment differs from ordinary wound infection.

\medterm{Pyogenesis} The formation of pus, a thick fluid made mainly of immune cells, dead tissue, and germs. It commonly occurs during bacterial infection or severe inflammation.
\medterm{Pyogenic} Producing pus or associated with a pus-forming infection. Pyogenic infections commonly cause redness, warmth, swelling, pain, and a collection of pus, although the cause and treatment vary.
\medterm{Pyogenic Granuloma} A noncancerous, rapidly growing lump made of extra small blood vessels. It is usually red, soft, and prone to bleeding and may occur after minor injury or during pregnancy; examination or biopsy confirms the diagnosis and guides removal.

\medterm{Pyogenic Liver Abscess} A pyogenic liver abscess is a pus-filled infection in the liver, usually from biliary, abdominal, or bloodstream spread, causing fever, pain, and possible sepsis.

\medterm{Pyometra} A collection of pus inside the uterus, usually caused by infection combined with blockage of drainage through the cervix. It may cause lower abdominal pain, fever, tenderness, or foul-smelling discharge; in older adults it can be linked to cervical narrowing or a tumor and needs prompt evaluation.
\medterm{Pyonephrosis} Pus trapped in the kidney's urine-collecting system, usually because infection occurs behind a blockage. It is a medical emergency requiring urgent drainage, antibiotics, and treatment of the obstruction.

\medterm{Pyosalpinx} A fallopian tube filled with pus, usually because of pelvic infection. It may cause lower abdominal pain, fever, unusual discharge, infertility, or an abscess and generally requires prompt treatment.
\medterm{Pyostomatitis Vegetans} A rare inflammatory mouth condition causing many small yellowish pus-filled spots on red, swollen tissue, especially the inner cheeks and gums. It is strongly associated with inflammatory bowel disease and may appear before bowel symptoms.

\medterm{Pyramid} A structure shaped like a pyramid. In the brain, the medullary pyramids contain nerve fibers that carry movement commands from the brain to the spinal cord; other pyramids include pyramid-shaped areas in the kidney.
\medterm{Pyramidal Cell} A triangular-shaped excitatory neuron with a prominent apical dendrite, found especially in the cerebral cortex and hippocampus and involved in long-range brain signaling.

\medterm{Pyramidal Tract} The pyramidal tract is a major motor pathway from the cerebral cortex to the brainstem and spinal cord that controls voluntary movement.

\medterm{Pyramidal Tracts} Nerve pathways from brain movement areas down the brainstem and spinal cord, carrying signals controlling voluntary movement.
\medterm{Pyrazinamide} An antibiotic used with other medicines during the initial treatment of tuberculosis. It helps shorten treatment but can injure the liver or raise uric acid, so monitoring and prescribed dosing are important.
\medterm{Pyrazines} A family of chemical compounds found naturally in some foods and used in flavors, medicines, and industry. Their effects on health depend on the specific compound, amount, and route of exposure.

\medterm{Pyrazole} A small five-membered chemical structure used to build medicines, pesticides, and other compounds. The word names a chemical family and does not identify one particular treatment or effect.

\medterm{Pyrazoles} Pyrazoles are a family of chemical compounds containing a pyrazole ring, including some medicines and pesticides.


\medterm{Pyrenochaeta} A group of fungi found mainly in soil and plants that rarely infect people. When infection occurs, it usually affects tissue beneath injured skin or people with weakened immunity.

\medterm{Pyrethrins} Natural insect-killing chemicals derived from chrysanthemum flowers and used in some pest-control products. Exposure can irritate skin and airways and, in large amounts, may cause vomiting, tremors, seizures, or breathing problems.
\medterm{Pyrethrum} The dried flower heads of certain Chrysanthemum species, used as a natural source of pyrethrins, insecticidal compounds that affect insect sodium channels. Pyrethrum and pyrethrins differ from synthetic pyrethroids; exposure may irritate the skin or airways and can cause toxicity with substantial exposure.
\medterm{Pyrexia (Fever)} An abnormally high body temperature, usually $38^\circ\text{C}$ ($100.4^\circ\text{F}$) or higher. Fever is a controlled response in which the brain raises the body’s temperature setting, often because of infection or inflammation. It differs from hyperthermia, in which temperature rises without that change in setting, such as during heat stroke. Treatment depends on the cause and the person’s symptoms.
\medterm{Pyridine} A nitrogen-containing ring-shaped chemical structure found in many medicines, vitamins, and biologically active substances. Pyridine groups can change how compounds dissolve, react, travel through the body, and produce medical effects.

\medterm{Pyridostigmine} A medicine that increases acetylcholine, a chemical helping nerves activate muscles. It mainly treats myasthenia gravis and may cause stomach cramps, diarrhea, sweating, increased saliva, muscle twitching, or slow heartbeat.
\medterm{Pyridostigmine Bromide} Pyridostigmine bromide increases communication between nerves and muscles and is used mainly for myasthenia gravis; it may cause abdominal cramps, diarrhea, or excess secretions.

\medterm{Pyridoxal Phosphate} The active form of vitamin B6 used by many enzymes. It helps the body process amino acids, make certain brain chemicals and red blood cells, and maintain normal nerve function.

\medterm{Pyridoxine} One form of vitamin B6, converted in the body to pyridoxal phosphate, a coenzyme for amino-acid metabolism. It is needed for processing food, making red blood cells, and supporting nerves and the immune system; deficiency can cause anemia or nerve problems, while excessive supplements can damage nerves.

\synonyms
Vitamin B6

\medterm{Pyridoxine (Vitamin B6)} Alternate index label; see \textit{Pyridoxine}.

\medterm{Pyrimethamine} A parasite-fighting medicine that blocks folate use in certain protozoa. It is used for selected infections, often with another medicine, and can suppress bone marrow or cause serious skin reactions.
\medterm{Pyrimidine} A six-membered nitrogen-containing ring that forms the bases cytosine, thymine, and uracil in nucleic acids.
\medterm{Pyrimidine Analog} A medicine or chemical resembling a pyrimidine nucleotide that can interfere with DNA or RNA synthesis, often for cancer or antiviral treatment.

\medterm{Pyrimidine Dimer} A DNA lesion formed when adjacent pyrimidine bases bond abnormally after ultraviolet exposure, potentially causing mutations if unrepaired.

\medterm{Pyrimidine Nucleosides} Nucleosides containing cytosine, thymine, or uracil bases linked to a sugar, serving as components or precursors of DNA and RNA.

\medterm{Pyrimidine Nucleotides} Genetic and energy-related building blocks containing pyrimidine bases, including cytidine, uridine, and thymidine. Cells use them to make DNA, RNA, cell membranes, and other substances needed for growth and repair.

\medterm{Pyrimidines} Nitrogenous bases containing a single six-membered heterocyclic ring, serving as
structural building blocks for nucleic acids. The primary pyrimidines are cytosine (C),
thymine (T, present in DNA), and uracil (U, present in RNA).

\medterm{Pyrogen} A pyrogen is a substance that induces fever, commonly through immune signaling; bacterial endotoxin is an important example in contaminated medicines or fluids.

\medterm{Pyroglyphidae} Pyroglyphidae is a mite family that includes house-dust mites whose allergens can trigger allergic rhinitis, asthma, or atopic dermatitis.


\medterm{Pyrophosphatase} An enzyme that hydrolyzes inorganic pyrophosphate and helps drive many biosynthetic reactions forward. An enzyme that breaks down pyrophosphate, helping many chemical reactions needed to build cellular materials.
\medterm{Pyropoikilocytosis} A rare inherited disorder that makes red blood cells unusually small, irregular, and fragile. The cells break down easily, causing hemolytic anemia with tiredness, jaundice, pale skin, or an enlarged spleen, often beginning in infancy or childhood.

\medterm{Pyroptosis} A form of controlled cell death that causes inflammation. The cell swells, breaks open, and releases signals that alert the immune system, often in response to infection or severe cell stress.

\medterm{Pyrosequencing} A laboratory method for reading selected DNA sequences by detecting light produced as DNA building blocks are added. It can identify some genetic changes or microbes but is not suitable for every sequencing need.

\medterm{Pyrosis} Heartburn, a burning feeling behind the breastbone usually caused by stomach contents flowing backward into the swallowing tube. Frequent symptoms may reflect reflux disease and can damage the esophagus without treatment.

\medterm{Pyrroles} A group of five-membered nitrogen-containing ring compounds found in heme, bile pigments, medicines, and other biologically active molecules.

\medterm{Pyrrolidonecarboxylic Acid} A cyclic amino-acid derivative, also called pyroglutamic acid, found in proteins and metabolic pathways; abnormal accumulation may occur in selected metabolic or medication-related conditions.

\medterm{Pyrrolnitrin} Pyrrolnitrin is an antimicrobial compound produced by some bacteria and studied for antifungal activity and biological control.

\medterm{Pyruvate} A small molecule produced when cells break down glucose. Cells can use it to make more energy, convert it into other substances, or produce lactate when oxygen or energy-processing pathways are limited.
\medterm{Pyruvate Decarboxylase} Pyruvate decarboxylase is an enzyme that converts pyruvate to acetaldehyde and carbon dioxide during fermentation in some organisms.

\medterm{Pyruvate Dehydrogenase} An enzyme complex that changes pyruvate, made when cells break down sugar, into acetyl-CoA for energy production. Inherited deficiency can cause excess lactic acid, poor growth, developmental problems, or neurologic disease.

\medterm{Pyruvate Kinase} An enzyme that helps red blood cells make energy from glucose. If it is inheritedly deficient, red cells may break down too early, causing chronic hemolytic anemia, jaundice, gallstones, tiredness, or an enlarged spleen.

\medterm{Pyruvate Kinase Blood Test} A pyruvate kinase blood test measures the enzyme in red blood cells and may help evaluate pyruvate kinase deficiency, an inherited cause of hemolytic anemia.

\medterm{Pyruvate Kinase Deficiency} A rare inherited condition in which red blood cells cannot make enough energy and break down early. It causes anemia of varying severity, jaundice, tiredness, gallstones, an enlarged spleen, and sometimes excess iron in the body.

\medterm{Pyuria} Presence of white blood cells or pus in urine, often indicating urinary-tract inflammation or infection.

\medterm{Paget's Disease} An ambiguous term that may refer to Paget disease of bone or Paget disease of the breast. Paget disease of bone is a chronic disorder of abnormal bone breakdown and regrowth that can enlarge, weaken, or deform bone; mammary Paget disease is a malignancy involving the nipple and areola, often associated with underlying breast cancer.

\medterm{Paget's Disease of the Nipple} A rare breast cancer involving the skin of the nipple and areola. It may cause redness, scaling, itching, discharge, or a lump and requires breast evaluation.

\medterm{Paget's Disease of the Vulva} A rare cancer or precancerous condition affecting the skin of the vulva. It may cause persistent itching, burning, redness, or a scaly patch and needs biopsy assessment.


\medterm{Paleo Diet} An eating pattern based on foods thought to resemble those available to early humans, such as meat, fish, vegetables, fruits, nuts, and seeds. It excludes or limits grains, legumes, and dairy and may not provide balanced nutrition for everyone.

\medterm{Papillomatosis} A condition in which multiple papillomas, or wart-like overgrowths of epithelial tissue, develop. Causes and risks vary by body site, virus, inherited condition, or other disease.

\medterm{Paracetamol} Also called acetaminophen, paracetamol reduces pain and fever. Excessive doses can cause serious liver injury, especially with alcohol or other products containing paracetamol.

\medterm{Paraneoplastic Neurologic Syndromes} Rare disorders in which an immune response to cancer also attacks the nervous system. Symptoms may include weakness, imbalance, seizures, memory change, or sensory problems and can precede cancer diagnosis.

\medterm{Paraneoplastic Pemphigus} A rare, often severe autoimmune blistering disease associated with an underlying tumor, especially a lymphoid cancer. It causes painful mouth sores and skin lesions and requires specialist treatment.

\medterm{Parechovirus} A group of viruses that commonly cause mild respiratory or gastrointestinal illness but can cause severe sepsis-like illness, meningitis, or encephalitis in newborns and young infants.

\medterm{Parenting} The care, protection, guidance, and support provided to a child. Healthy parenting includes meeting basic needs, setting safe limits, nurturing development, and seeking help when concerns arise.

\medterm{Paresthesia (Pins and Needles)} An abnormal sensation such as tingling, prickling, burning, or numbness. Temporary pressure on a nerve is common, but persistent or spreading symptoms may indicate nerve, metabolic, vascular, or spinal disease.

\medterm{Paruresis} Difficulty urinating in the presence of other people or in public places, often called shy bladder syndrome. It is related to anxiety and can improve with behavioral therapy and gradual exposure.

\medterm{Passive Heat Therapy} Use of external heat, such as warm packs, baths, or saunas, to relax muscles and temporarily reduce pain or stiffness. Excessive heat can cause burns, dehydration, or fainting.

\medterm{Passive Smoking} Breathing smoke from another person's tobacco or nicotine product. Secondhand smoke increases the risk of respiratory disease, heart disease, stroke, and lung cancer and is especially harmful to children and during pregnancy.

\medterm{Patau's Syndrome} A congenital disorder caused by an extra copy of chromosome 13. It can cause severe developmental problems, brain and heart defects, cleft lip or palate, and other organ abnormalities.

\medterm{Paternal Depression} Depression occurring in a father during pregnancy or after a child is born. Symptoms may include low mood, irritability, withdrawal, sleep or appetite changes, and loss of interest; treatment can help both parent and family.

\medterm{Paternal Postnatal Depression} Depression in a father during the first year after a baby's birth. It may be linked to stress, sleep loss, previous depression, or a partner's depression and deserves professional assessment.

\medterm{Pelvic Inflammatory Disease (PID)} Infection and inflammation of the upper female reproductive organs, often caused by sexually transmitted bacteria. Untreated PID can cause chronic pelvic pain, infertility, or ectopic pregnancy.

\medterm{Pemphigus Foliaceus} An autoimmune blistering skin disease caused by antibodies against proteins that hold skin cells together. It causes superficial fragile blisters and crusted sores but usually does not affect mucous membranes.

\medterm{Penile Intraepithelial Neoplasia (PeIN or PIN)} Abnormal precancerous cells in the surface lining of the penis. It may appear as a persistent red, white, or thickened patch and can progress to invasive cancer if untreated.

\medterm{Percutaneous Closure} Closing a heart or blood-vessel opening using a device delivered through a catheter inserted through the skin. It can avoid open surgery in selected patients but has procedure-related risks.

\medterm{Personalized Medicine} Healthcare that uses a person's genes, biological measurements, environment, and lifestyle to guide prevention or treatment. It can improve treatment selection but does not guarantee a response.

\medterm{Peyronie's Disease} A condition in which scar tissue forms in the penis and causes a curved or painful erection, shortening, or difficulty with sexual activity. Treatment depends on severity and symptoms.

\medterm{pH} A measure of how acidic or alkaline a solution is, usually on a scale from 0 to 14. Body fluids have tightly controlled pH ranges, and major changes can be dangerous.

\medterm{Phantom Pregnancy} A false belief or physical state in which a person has pregnancy-like symptoms without an ongoing pregnancy. Symptoms can include abdominal enlargement, breast changes, nausea, or missed periods and require medical evaluation.

\medterm{Pharmacoepidemiology} The study of how medicines are used and affect health in large populations. It helps assess effectiveness, safety, adverse effects, and patterns of prescribing.


\medterm{Pharmacoproteomics} The study of how a person's protein patterns influence medicine response, drug effects, or toxicity. It is an emerging part of precision medicine.

\medterm{Photoplethysmography} A noninvasive optical method that detects blood-volume changes in tissue. It is used in devices such as pulse oximeters and some heart-rate monitors.

\medterm{Phyllodes Tumors} Uncommon fibroepithelial tumors of the breast that can be benign, borderline, or malignant. They often appear as a rapidly growing lump and are treated mainly with complete surgical removal.


\medterm{Pick's Disease} An older term often referring to frontotemporal dementia, a group of disorders causing progressive changes in behavior, personality, language, or executive function. The term is now used less consistently.

\medterm{Piles (Haemorrhoids)} Alternate spelling; see \textit{Hemorrhoids}.

\medterm{Pins and Needles} A common name for paresthesia, an abnormal tingling, prickling, or numb sensation. Brief episodes often follow pressure on a nerve; persistent symptoms need evaluation.

\medterm{Polymorphic Light Eruption} An itchy rash triggered by sunlight, often appearing in spring or after the first strong sun exposure. It may improve as the skin adapts, and sun protection reduces recurrences.

\medterm{Polypharmacology} The study or use of a medicine that affects several targets in the body. This can be helpful when one disease involves multiple pathways, but it may also increase side effects and medicine interactions.

\medterm{Portal Hypertensive Colopathy} Changes in the colon caused by high pressure in veins around the liver, usually from advanced liver disease. Enlarged fragile blood vessels may cause visible or hidden bleeding and anemia.

\medterm{Post-Adoptive Depression} Depression that begins after adopting a child. Stress, sleep loss, adjustment demands, past trauma, or unmet expectations may contribute; persistent sadness, loss of interest, or thoughts of self-harm need prompt mental-health assessment.

\medterm{Posterior Cruciate} Usually referring to the posterior cruciate ligament, a strong knee ligament that limits backward movement of the shinbone. Injury can cause pain, swelling, and knee instability.

\medterm{Postnatal Anxiety} Excessive anxiety occurring after childbirth. It may involve persistent worry, panic, intrusive thoughts, restlessness, or physical symptoms and can be treated with therapy, support, and sometimes medicine.

\medterm{Post-Traumatic Stress Disorder (PTSD)} Alternate index label; see \textit{Post-Traumatic Stress Disorder}.

\medterm{Powassan Virus} A tick-borne virus that can cause encephalitis or meningitis. Severe infection may cause headache, fever, confusion, seizures, or weakness; there is no specific antiviral treatment.

\medterm{Precision Medicine} An approach that matches prevention or treatment to biological differences between people, often using genetic, molecular, or clinical information. It is related to but not identical to personalized medicine.

\medterm{Precision Oncology} An approach to cancer care that uses information about a tumor’s genes, proteins, or other features to help select prevention, testing, or treatment. It aims to match treatment to the biology of the cancer rather than to its location alone.

\medterm{Prenatal Depression} Depression during pregnancy. Symptoms may include persistent sadness, loss of interest, guilt, hopelessness, sleep or appetite changes, and thoughts of self-harm; prompt professional support is important.

\medterm{Programmed Cell Death} A controlled process by which cells die as part of normal development, tissue maintenance, or disease. Apoptosis is the best-known form; other regulated forms include necroptosis and pyroptosis.

\medterm{Proteasome Inhibitor} A medicine that blocks the proteasome, a cell structure that breaks down unwanted proteins. These drugs are used especially in multiple myeloma and some blood cancers.

\medterm{Proteostasis} The balance of protein production, folding, movement, and breakdown that keeps cells functioning. Failure of proteostasis can contribute to aging, neurodegeneration, and other diseases.

\medterm{Pruritis} The medical term for itching, an unpleasant urge to scratch. Causes include dry skin, allergy, inflammation, infection, medicines, nerve disorders, and internal disease.

\medterm{Psychosomatic Disorder} An older broad term for physical symptoms influenced by psychological and social factors. The symptoms are real; modern assessment considers both body and mind rather than assuming symptoms are imagined.

\medterm{Psychosomatic Medicine} Medical care that considers how body processes, thoughts, emotions, relationships, and social conditions interact in illness. Physical symptoms are real, and assessment looks for both medical and psychological factors without assuming symptoms are imagined.

\medterm{Pulmonary Hypoplasia} Underdevelopment of one or both lungs, resulting in too little lung tissue for normal breathing. It may occur with congenital abnormalities, low amniotic fluid, or compression before birth.

\medterm{Pulmonology} The medical specialty concerned with the lungs and breathing system. Pulmonologists diagnose and treat asthma, chronic lung disease, infections, sleep-related breathing disorders, and other respiratory conditions.

\medterm{PYY} Peptide YY is a hormone released by the intestine after eating. It helps signal fullness and slows movement of food through the digestive tract; abnormal signaling is being studied in obesity and other metabolic disorders.


\medterm{Q Arm Of A Chromosome} The longer arm of a chromosome, identified by the letter q. A genetic report may use this label to show where a gene change or missing, extra, or rearranged DNA segment is located.

\medterm{Q Fever} An infection caused by the bacterium Coxiella burnetii, usually caught by breathing dust contaminated by infected cattle, sheep, goats, or their birth products. It often causes sudden fever, severe headache, tiredness, and muscle aches; some people develop pneumonia, hepatitis, or long-term heart-valve infection.

\medterm{Q Wave} A downward movement that may begin the electrical pattern of a heartbeat on an electrocardiogram. It can be normal in some recording positions or suggest previous heart-muscle injury, depending on its size and location.

\medterm{Q-T Interval} The time from the start of electrical activation to the end of recovery of the heart’s lower chambers on an electrocardiogram. An unusually long interval can increase the risk of a dangerous abnormal rhythm.



\medterm{QRS Complex} The group of waves on an electrocardiogram showing the electrical signal that makes the heart’s lower chambers contract. An unusually wide or shaped complex may indicate a conduction problem, abnormal rhythm, or electrolyte disturbance.

\medterm{QT Interval} The time on an electrocardiogram representing electrical activation and recovery of the heart's lower chambers. Because heart rate changes this time, clinicians often use a corrected value when assessing rhythm risk.

\medterm{QT Prolongation} An unusually long QT interval on an electrocardiogram. It may result from an inherited condition, a medicine, or abnormal body salts and can increase the risk of a rapid, dangerous rhythm that may cause fainting or cardiac arrest.

\medterm{Quad Screen (Maternal Serum Screening)} A second-trimester blood screening test that estimates the chance of certain chromosome conditions and neural-tube defects in the fetus. It estimates risk rather than making a diagnosis; an abnormal result may require ultrasound or diagnostic testing.

\synonyms
Maternal Serum Screening

\medterm{Quadrant} One of four sections used to describe a location, especially in the abdomen. The four abdominal quadrants help clinicians record where pain, tenderness, a lump, or another finding is located.

\medterm{Quadrantanopia} Loss of vision in one quarter of the normal visual field, often from damage to the brain pathways carrying sight. The missing area may occur on the same side in both eyes and can affect reading, driving, walking, or safety.

\medterm{Quadrantectomy} Surgery to remove one section of an organ, most often part of the breast containing a tumor plus a margin of healthy tissue. The need for lymph-node surgery or further treatment depends on the tumor and test results.

\medterm{Quadriceps} The four muscles on the front of the thigh that extend the knee; the rectus femoris also assists hip flexion. The quadriceps tendon connects them to the patella and contributes to the extensor mechanism.

\medterm{Quadriceps Femoris} The group of four muscles at the front of the thigh. It straightens the knee for standing, walking, climbing, and kicking; one part also helps bend the hip.

\medterm{Quadriceps Femoris Muscle} The four-muscle group at the front of the thigh that straightens the knee. It is essential for standing, walking, climbing, and rising from a chair; injury can cause pain and difficulty extending the leg.

\medterm{Quadriceps Injury} Damage to the quadriceps muscle or tendon caused by sudden force, overuse, or direct trauma. It may cause
front-of-thigh pain, bruising, weakness, or difficulty extending the knee. Severe pain, a palpable defect, or
inability to bear weight requires assessment for a significant tear or tendon rupture.

\medterm{Quadriceps Tendinitis} Overuse-related inflammation or degeneration of the quadriceps tendon, causing pain above the kneecap that worsens with jumping, running, or resisted knee extension.

\medterm{Quadriceps Tendon Rupture} A tear of the tendon connecting the quadriceps muscle to the patella, causing pain, swelling, and inability to actively extend the knee.

\medterm{Quadriceps Tendon Ruptures} A tear of the tendon connecting the quadriceps muscle to the patella, causing weakness
or inability to extend the knee. Complete tears often require surgical repair; treatment
depends on the extent of injury and the person's functional needs.

\medterm{Quadriparesis} Weakness affecting both arms and both legs. It may result from disease or injury to the neck spinal cord, brain, nerves, or muscles and can impair movement, breathing, and bladder or bowel control.

\medterm{Quadriplegia} Paralysis affecting both arms and both legs, usually with weakness of the trunk, because of damage to the neck portion of the spinal cord or another major nervous-system injury. It can impair breathing and bladder or bowel control.

\medterm{Quadruple Screen Test} A second-trimester blood test that combines four pregnancy-related markers to estimate the chance of selected chromosome conditions and neural-tube defects. It does not diagnose a condition; abnormal results need further assessment.

\medterm{Quadruplet Pregnancy} A pregnancy involving four fetuses. It carries increased risks of preterm birth, fetal
growth restriction, hypertensive and metabolic complications, and maternal morbidity and
therefore requires specialized multidisciplinary antenatal care.

\medterm{Quality Of Life} A person’s experience of physical health, emotional well-being, relationships, independence, and ability to do meaningful activities. Illness and treatment may affect these areas differently for each person.

\medterm{Quantitative Bence-Jones Protein Test} A urine test measuring abnormal light-chain proteins made by some plasma-cell cancers. It can help diagnose or monitor disorders such as multiple myeloma, but results are interpreted with blood tests and other findings.

\medterm{Quantitative Computed Tomography} A CT scan that measures bone density while showing bone structure. It can help assess osteoporosis, unusual bone loss, or fracture risk when standard bone-density testing is unsuitable or more detail is needed.

\medterm{Quantitative Immunoglobulins} A blood test measuring the main classes of antibodies. Low levels may suggest an immune deficiency, while high or uneven levels may occur with infection, inflammation, liver disease, or some blood cancers; results need clinical interpretation.

\medterm{Quantitative Nephelometry Test} A laboratory test that estimates the amount of selected proteins by measuring how much light they scatter. It can measure antibodies, complement proteins, and other proteins involved in immune, kidney, or blood disorders.

\medterm{Quantitative PET} Analysis of a PET scan that estimates how much radioactive tracer is taken up by tissue. Standardized measurements can help compare disease activity or treatment response, but results are affected by the tracer, timing, blood sugar, and scanning method.




\medterm{Quarantine} Temporary separation or movement restrictions for people exposed to an infectious disease while monitoring whether they become ill. It differs from isolation, which separates people who are already infected.

\medterm{Quartan Fever} A fever pattern in which attacks recur about every 72 hours, classically seen with the malaria parasite Plasmodium malariae. Fever timing alone cannot diagnose malaria; testing is needed in anyone with possible exposure and fever.













\medterm{Quetiapine Fumarate} A form of quetiapine, a medicine used for schizophrenia, bipolar disorder, and selected depression treatment. It can cause sleepiness, dizziness on standing, weight or blood-sugar changes, and movement problems.

\medterm{Quick Relievers} Fast-acting asthma medicines used when symptoms suddenly worsen. They quickly relax tightened airway muscles and can ease wheezing, coughing, chest tightness, and breathlessness, but they do not replace regular treatment that controls airway inflammation.

\synonyms
Rescue medicines

\medterm{Quickening} The first time a pregnant person feels fetal movement, often as a flutter or tapping sensation during the second trimester. Timing varies with the pregnancy, body shape, and position of the placenta.


\medterm{Quinalphos} A poisonous insecticide that disrupts nerve-signal control. Significant exposure can cause sweating, excess saliva, vomiting, muscle twitching, weakness, seizures, breathing failure, or death and requires emergency poison treatment.

\medterm{Quinapril} Quinapril is an angiotensin-converting-enzyme inhibitor used for hypertension and heart failure; cough, hyperkalaemia, renal-function changes, angioedema, and fetal toxicity are important risks.


\medterm{Quincke's Disease (Angioneurotic Oedema)} Sudden swelling of deeper layers of the skin or mucosa, commonly affecting the lips, face, tongue, throat, or bowel; airway involvement is an emergency.

\synonyms
Angioneurotic Oedema

\medterm{Quincke's Pulse} Visible flushing and paling of the nail bed or skin with each heartbeat. It may occur when the difference between the upper and lower blood-pressure numbers is unusually wide, especially with severe aortic-valve leakage, but it is not diagnostic alone.

\medterm{Quincke's Sign} Alternating redness and paling of the nail bed or skin with each heartbeat, usually caused by a wide pulse pressure. It can be seen in severe aortic-valve leakage but is not diagnostic by itself.

\medterm{Quinidine} A heart-rhythm medicine that changes electrical conduction in the heart. It is used selectively because it can cause dangerous abnormal rhythms, low blood pressure, low platelets, diarrhea, or other serious side effects.

\medterm{Quinine} A medicine derived from cinchona bark that can treat selected malaria infections. It is used less often because safer options are available and it can cause low blood sugar, ringing in the ears, blood-cell problems, or abnormal heart rhythms.


\medterm{Quinolone} A class of antibiotics that stop bacteria from copying and repairing their DNA. Some treat serious infections, but side effects can include tendon injury, nerve symptoms, abnormal heart rhythms, or severe diarrhea, so use depends on the infection and risks.



\medterm{Quinsy} An abscess beside a tonsil, usually after a throat infection. It can cause severe one-sided throat pain, fever, muffled speech, drooling, difficulty opening the mouth, or breathing difficulty and needs urgent treatment.



\medterm{Quinupristin-Dalfopristin} An intravenous antibiotic combination used rarely for selected serious infections caused by resistant gram-positive bacteria. It can cause muscle or joint pain, liver problems, infusion reactions, and medicine interactions.


\medterm{Quit Smoking (Smoking And Quitting Smoking)} Alternate index label for Smoking and Quitting Smoking.


\medterm{Quincke's Edema} An older term for angioedema, sudden swelling in the deeper layers of the skin or mucous membranes. It commonly affects the lips, face, tongue, or throat; throat swelling or breathing difficulty is an emergency.

\medterm{Quotidian Fever} A fever pattern in which fever rises or attacks occur about every 24 hours. It can occur with malaria and other infections, but the timing of fever alone does not identify the cause.

\medterm{Rabeprazole} A medicine that greatly reduces acid made by the stomach. It treats reflux and ulcers and may be combined with antibiotics for Helicobacter pylori infection; possible effects include headache, diarrhea, and rarely low magnesium or bowel infection.

\medterm{Rabies} A viral infection of the brain and nervous system caused by rabies virus, a virus that infects mammals. It spreads mainly when saliva from an infected animal enters a bite, scratch, broken skin, or the eyes, mouth, or another mucous membrane. After a variable symptom-free period, early symptoms may include fever, weakness, headache, or pain, itching, or unusual sensations near the exposure site. Later symptoms may include confusion, agitation, hallucinations, fear of water, difficulty swallowing, paralysis, coma, and death. Once symptoms begin, rabies is nearly always fatal, but vaccination before exposure or post-exposure treatment can prevent disease before symptoms start.

\medterm{Rabies Vaccination} Immunization used before or after exposure to reduce the risk of rabies. Post-exposure management
also includes wound care and, when indicated, rabies immunoglobulin.

\medterm{Rabies Virus} The virus that causes rabies, a usually fatal viral infection of the central nervous system transmitted most often through the saliva of an infected animal.

\medterm{Rachitic Rosary} A physical examination finding in pediatric rickets (vitamin D deficiency) characterized by prominent, beadlike, symmetrical nodular expansion of the costochondral junctions along the anterolateral rib cage resulting from disordered osteoid mineralization and chondrocyte hypertrophy.
\synonyms
Beaded Ribs; Costochondral Rosary

\medterm{Radial} Relating to the radius bone or the thumb side of the forearm and hand. The word can also describe structures arranged like spokes spreading from a central point.

\medterm{Radial Artery} An artery that runs along the thumb side of the forearm to the wrist and hand. Its pulse can be felt at the wrist, and it supplies blood to the forearm and parts of the hand.

\medterm{Radial Artery Thrombosis} A blood clot blocking the artery on the thumb side of the forearm, sometimes after an arterial catheter or injury. It may cause a painful, pale, cold, numb, or weak hand and can threaten hand circulation.

\medterm{Radial Club Hand} A birth-related difference in which the thumb-side forearm bone is small or absent and the hand bends toward the thumb side. The thumb may also be underdeveloped, and treatment depends on movement, function, and associated conditions.

\medterm{Radial Head Subluxation} A common elbow injury in young children, often called nursemaid’s elbow, caused when a sudden pull partly slips a forearm bone out of its supporting band. The child may stop using the arm but usually has little swelling. A trained clinician can put the bone back in place.

\synonyms
Nursemaid's Elbow; Pulled Elbow; Annular Ligament Displacement

\medterm{Radial Nerve} A major nerve of the arm that helps straighten the elbow, wrist, and fingers and carries feeling from part of the back of the hand. Injury can cause wrist drop, finger weakness, or numbness.

\medterm{Radial Nerve Dysfunction} Impaired function of the radial nerve causing weakness of wrist or finger extension and sensory changes over part of the posterior arm or hand; causes include compression, trauma, and systemic neuropathy.

\medterm{Radial Nerve Palsy} Dysfunction of the radial nerve causing weakness of wrist and finger extension, often producing wrist drop, with sensory loss over part of the back of the hand in some cases.

\medterm{Radial Tunnel Syndrome} Irritation or pressure on the radial nerve near the elbow, causing aching outer-forearm pain and sometimes weakness. Symptoms may worsen with repeated wrist or forearm movement and require evaluation to exclude other causes.

\medterm{Radial Vein} Veins running alongside the radial artery on the thumb side of the forearm. They drain blood from the forearm and join other deep veins near the elbow.

\medterm{Radiation} Energy traveling as particles or electromagnetic waves, including X-rays, gamma rays, ultraviolet light, and radio waves. Depending on its type and amount, radiation can diagnose disease, treat cancer, or damage living tissue.

\medterm{Radiation Biology} The study of how radiation affects cells, tissues, organisms, and populations. It examines injury, repair, mutations, cancer risk, treatment effects, and protective measures after exposure to ionizing or nonionizing radiation.

\medterm{Radiation Contamination} Radioactive material present on or inside a person, object, surface, or body site when it should not be. It can spread or continue exposing tissue and may require isolation, monitoring, removal of the material, and specialist advice.

\medterm{Radiation Dermatitis Skin Care} Care of skin affected by radiation treatment. The skin may become red, dry, sore, peeling, or moist, and care focuses on protecting the skin and treating irritation or infection when present.

\synonyms
Radiation Skin Care; Radiation Dermatitis Management; Post-Radiotherapy Skin Care

\medterm{Radiation Dosage} The amount of radiation given during an imaging test or radiation treatment. The dose is planned to obtain the needed benefit while limiting unnecessary exposure.

\medterm{Radiation Dose} The amount of ionizing radiation absorbed or received by a person or tissue. Dose varies with the
examination, equipment, protocol, and patient size; imaging should be justified and optimized
to obtain the needed information with the lowest reasonable exposure.

\medterm{Radiation Emergency} An event involving actual or suspected harmful radiation exposure or contamination that requires urgent protective and medical
measures.

\medterm{Radiation Enteritis} Inflammation or injury of the small or large intestine after radiation treatment. Early symptoms may include diarrhea, cramps, nausea, or urgency; months or years later, scarring can cause narrowing, poor absorption, bleeding, or blockage.

\medterm{Radiation Exposure} Contact with ionizing or non-ionizing radiation from medical imaging, treatment, sunlight, workplaces, or other sources. The possible harm depends mainly on the type, dose, body area, and duration of exposure.

\medterm{Radiation Fibrosis} Long-term thickening and stiffening of tissue caused by radiation injury, sometimes appearing months or years after treatment. It can affect skin, lungs, bowel, or other organs and may cause pain, reduced movement, or shortness of breath.

\medterm{Radiation Field} The area and distribution of radiation delivered during imaging or radiation treatment. In therapy, its size, shape, direction, and dose are planned to cover the target while limiting exposure to nearby healthy tissue.

\medterm{Radiation-Induced Angiosarcoma} A rare, fast-growing cancer of blood-vessel lining cells that develops years after radiation in previously treated tissue. It may appear as a bruise, purple patch, swelling, or lump and requires prompt specialist assessment.

\medterm{Radiation Leukemia} Radiation leukemia is leukemia that develops after significant exposure to ionizing radiation; it is a late treatment- or exposure-related complication.

\medterm{Radiation Monitoring} Measurement of radiation levels or individual exposure using devices such as personal dosimeters and area monitors.

\medterm{Radiation Necrosis} Permanent death of tissue caused by radiation injury, often developing months or years after treatment. In the brain it may cause headache, seizures, swelling, weakness, or speech and vision problems and can resemble returning cancer on scans.

\medterm{Radiation Oncologist} A doctor who plans and oversees radiation treatment for cancer and selected noncancerous conditions. They choose treatment areas and doses, coordinate care, monitor side effects, and adjust therapy during treatment.

\medterm{Radiation Oncology} A medical specialty that uses ionizing radiation to treat cancer. Treatment may aim to
destroy or control tumors, prevent cancer from returning, or relieve symptoms.

\medterm{Radiation Physics} The study of how radiation is produced, travels, interacts with tissues, and is measured. It supports medical imaging, radiation treatment, dose calculation, and protection from unnecessary exposure.

\medterm{Radiation Pneumonitis} Inflammation of lung tissue occurring weeks or months after radiation to the chest. It may cause dry cough, fever, breathlessness, or low oxygen and can later leave permanent scarring; new breathing symptoms need medical assessment.

\medterm{Radiation Poisoning} Illness caused by substantial exposure to ionizing radiation, potentially producing skin injury,
bone-marrow suppression, gastrointestinal symptoms, or organ failure.

\medterm{Radiation Pregnancy Screening} Assessment of pregnancy status before procedures involving ionizing radiation or
radioactive materials when pregnancy is possible. The decision to proceed requires
clinical justification, risk optimization, and consideration of gestational age and
fetal dose.

\medterm{Radiation Proctitis} Inflammation and injury of the rectum after pelvic radiation therapy, causing rectal pain, urgency, diarrhea, mucus, or bleeding.

\medterm{Radiation Protection} The practice of keeping radiation exposure as low as reasonably achievable while
obtaining the required diagnostic or therapeutic result. It relies on justification,
optimization, time reduction, distance, shielding, and appropriate contamination
control.

\medterm{Radiation Recall} Inflammation that returns in an area previously treated with radiation after a later trigger, often a medicine. The skin or lining may become red, sore, swollen, peel, or blister, and the treating team should review the suspected trigger.

\medterm{Radiation-Recall Dermatitis} An inflammatory skin reaction that occurs in a previously irradiated area after exposure to a
triggering substance, often months or years later. The area may become red, swollen, painful,
or blistered and requires clinical evaluation.

\medterm{Radiation Safety} Measures that limit unnecessary exposure to ionizing radiation while allowing needed imaging or treatment. They include using the shortest practical exposure time, keeping as much distance as possible, using shielding, preventing contamination, and following pregnancy-related precautions when relevant.

\medterm{Radiation Sickness} Acute radiation syndrome caused by high-dose whole-body or substantial partial-body
exposure to ionizing radiation. It may produce nausea, vomiting, marrow failure,
infection, bleeding, gastrointestinal injury, neurovascular collapse, and death
depending on dose.

\medterm{Radiation Therapy} Treatment that uses carefully aimed ionizing radiation to damage cancer cells so they stop dividing or die. It may be given from outside the body or from a source placed near or inside a tumor. Side effects depend on the body area and dose and may include tiredness or skin changes.

\medterm{Radiation Tolerance} The amount of radiation that a tissue or organ can receive before injury becomes likely. It varies with the tissue, dose, treatment schedule, and previous radiation.

\medterm{Radical Cystectomy} Surgery to remove the urinary bladder, usually for muscle-invasive or high-risk bladder cancer. Because urine can no longer leave through the bladder, another route for urine drainage is created.

\medterm{Radical Cystoprostatectomy} Radical cystoprostatectomy removes the urinary bladder, prostate, and nearby structures, usually to treat invasive bladder cancer in men.

\medterm{Radical Hysterectomy} An operation that removes the uterus and cervix together with surrounding parametrial tissues and the upper vagina, usually with pelvic lymph-node assessment or removal. It is used mainly for selected invasive cervical cancers and is more extensive than a simple or total hysterectomy.

\medterm{Radical Lymphadenectomy} Extensive surgery to remove lymph nodes and nearby tissue to look for or treat cancer spread. It may provide important staging information but can cause pain, nerve or vessel injury, and long-term swelling from lymph-fluid blockage.

\medterm{Radical Mastectomy} Radical mastectomy removes the breast, chest muscles, and many underarm lymph nodes; it is now rarely used because less extensive cancer surgery is often effective.

\medterm{Radical Nephrectomy} Radical nephrectomy removes an entire kidney and sometimes nearby tissues, usually to treat a kidney tumor.

\medterm{Radical Prostatectomy} Surgery to remove the whole prostate gland and usually the nearby seminal vesicles, most often to treat prostate cancer that has not spread widely. The urinary passage is rejoined to the bladder; possible effects include urine leakage, erection problems, bleeding, and infection.

\medterm{Radicular Cyst} An inflammatory odontogenic cyst arising at the apex of a nonvital tooth, usually from chronic pulpal infection and apical periodontitis.

\medterm{Radiculitis} Irritation or inflammation of a nerve root where it leaves the spine. It can cause pain spreading along an arm or leg, tingling, numbness, weakness, or changed reflexes because of a disk problem, narrowing, infection, or other disease.

\medterm{Radiculopathy} Dysfunction of a spinal nerve root, usually from disc herniation, narrowing, or inflammation. It may cause
shooting pain, tingling, numbness, weakness, or altered reflexes along the affected nerve pathway.

\medterm{Radioactive Iodine} A radioactive form of iodine taken up by thyroid cells. It is used to image the thyroid or destroy overactive thyroid tissue and selected thyroid cancer cells; radiation-safety precautions are needed after some treatments.

\medterm{Radioactive Iodine Ablation} Treatment using radioactive iodine to destroy remaining thyroid tissue or selected thyroid cancer cells after surgery. Thyroid cells absorb the iodine; the decision depends on cancer type, spread, recurrence risk, and the person’s overall health.

\medterm{Radioactive Iodine Uptake} A test measuring the percentage of administered radioactive iodine taken up by the thyroid over time, used to evaluate hyperthyroidism and selected thyroid disorders.

\medterm{Radioactive Iodine Uptake Test} A nuclear medicine test that measures how much radioactive iodine the thyroid gland absorbs. It helps evaluate hyperthyroidism and some thyroid disorders.

\medterm{Radioactivity} Energy released when unstable atoms change into more stable forms. The energy may be particles or rays and can be used in medical imaging or treatment, but enough exposure can damage living tissue.

\medterm{Radioallergosorbent Test} An in vitro radioimmunoassay that detects and quantifies allergen-specific IgE antibodies in a patient's blood serum, used to identify allergic sensitization to environmental aeroallergens, foods, medications, or venoms without risking in vivo anaphylaxis.

\synonyms
RAST

\medterm{Radiobiology} The study of how radiation affects living cells and tissues, including DNA damage, cell death, repair, mutations, treatment effects, and the development of cancer.

\medterm{Radiocarpal Joint} The main wrist joint between the forearm bone on the thumb side and several wrist bones. It allows the hand to bend forward, backward, and from side to side and may be injured by falls or arthritis.

\medterm{Radiochemical Purity} The proportion of a radioactive medicine present in the intended chemical form. Unwanted forms can make a scan less accurate, send radiation to the wrong tissues, or reduce treatment effectiveness.

\medterm{Radiochemistry} The chemistry of radioactive substances and compounds. In medicine it includes making and labeling radiopharmaceuticals, studying radioactive decay, separating products, and checking radiochemical purity and safety.

\medterm{Radiocontrast} A substance given during an imaging test to make selected tissues, blood vessels, or body spaces easier to see. Depending on the agent, allergic reactions, kidney injury, or leakage-related damage may occur.

\medterm{Radiodermatitis} Skin injury caused by ionizing radiation, such as radiation treatment. It may cause redness, dryness, itching, hair loss, peeling, sores, or long-term skin changes; severity depends on the dose, treatment area, and repeated exposure.

\medterm{Radioembolisation} Alternate name; see \textit{Radioembolization}.

\medterm{Radioembolization} A treatment that delivers tiny radioactive beads through an artery to selected liver tumors. The beads reduce the tumor’s blood supply and release radiation nearby; suitability depends on liver function, tumor extent, and blood flow.

\medterm{Radio-Femoral Delay} A clinically palpable lag or latency in the arrival of the femoral pulse compared to the simultaneous radial pulse, often accompanied by reduced femoral pulse amplitude. It is a classic clinical sign of coarctation of the aorta and extensive aortoiliac occlusive disease.

\medterm{Radiofrequency} A type of electromagnetic energy, between some radio waves and microwaves, used in communication and in selected medical tests and treatments. Medical uses include heating tissue or delivering energy during procedures.

\synonyms
Radiofrequency Energy; RF Radiation

\medterm{Radiofrequency Ablation} A procedure that uses high-frequency electrical energy to heat and destroy selected tissue, often under imaging guidance.

\medterm{Radiofrequency Ablation For Chronic Pain} A procedure that applies radiofrequency energy to selected sensory nerves or other targets to interrupt pain signaling for a limited period in carefully selected chronic pain conditions.

\medterm{Radiograph} An image made with X-rays. Dense structures such as bone appear light, while air appears dark and soft tissues appear in shades of gray. Radiographs are used to assess bones, the chest, teeth, and selected abdominal or foreign-body problems; they use a small amount of radiation.

\synonyms
X-Ray; Roentgenogram; Radiographic Image

\medterm{Radiographer} A trained professional who performs diagnostic imaging examinations or delivers radiotherapy, depending on specialty and jurisdiction.

\medterm{Radiography} Creation of images using X-rays. It is useful for bones, chest disease, dental structures, and many other conditions; radiation exposure is kept as low as reasonably achievable.

\medterm{Radioimmunoassay} A laboratory test that uses an antibody and a small radioactive label to measure a hormone, medicine, or protein in a sample. It can be very sensitive but is often replaced by newer nonradioactive tests.

\medterm{Radioimmunoguided Surgery} Surgery in which a radioactive tracer attached to an antibody helps the surgeon locate tissue carrying a particular marker, often tumor tissue. Special equipment detects the tracer during the operation.

\medterm{Radioimmunotherapy} Treatment that attaches a radioactive substance to an antibody so the antibody carries radiation toward cells with a particular marker, often cancer cells. Radiation can still affect nearby normal tissue and cause side effects.

\medterm{Radioiodine} A radioactive form of iodine used in nuclear medicine to image or treat thyroid tissue. It is taken up by thyroid cells; precautions are needed to limit radiation exposure to others.

\medterm{Radioiodine Therapy} Treatment using radioactive iodine, which is taken up mainly by thyroid cells to destroy overactive thyroid tissue or selected thyroid cancer cells. It requires individualized dosing and temporary radiation-safety precautions; pregnancy and breastfeeding are contraindications.

\medterm{Radioisotope} An unstable isotope that emits radiation as it decays; radioisotopes are used in medical imaging, diagnosis, and selected treatments.

\synonyms
Radioactive Isotope; Radiotracer

\medterm{Radioligand Therapy} A cancer treatment in which a radioactive particle is attached to a molecule that binds a target on cancer cells. The molecule delivers radiation to cells carrying that target; use depends on the cancer, target expression, organ function, and treatment-specific safety monitoring.

\medterm{Radiologic Health} The safe use of medical imaging and radiation-based treatment to protect patients, workers, and the public from unnecessary ionizing radiation while preserving clinical benefit. Safety measures include appropriate justification, shielding, and dose optimization.

\medterm{Radiologic-Pathologic Correlation} Comparing scan findings with tissue examined under a microscope to determine whether the sample explains the abnormal area. If they do not match, another sample or further testing may be needed.

\medterm{Radiologic Response} The change seen on medical scans after treatment, such as shrinkage, disappearance, stability, growth, or new lesions. Standard measurement rules help classify the response, but scans do not always show the full biological effect.

\medterm{Radiologic Technologist} A trained health professional who performs selected medical-imaging examinations and helps protect patients from unnecessary radiation exposure.

\medterm{Radiologist} A doctor trained to interpret medical images and sometimes perform image-guided procedures. Radiologists evaluate X-rays, ultrasound, CT, MRI, and other scans to diagnose disease, guide treatment, or obtain tissue samples.

\medterm{Radiology} The medical field that uses images to find and monitor disease and guide procedures. Radiologists interpret X-rays, ultrasound, CT, MRI, and other scans and may perform image-guided biopsies, drainages, injections, or treatments.

\medterm{Radiology-Patient Identification} Verification of the correct patient, examination, body site, indication, pregnancy
status when relevant, allergies, renal function, and prior imaging before image
acquisition or intervention. It is a core diagnostic and procedural safety step.

\medterm{Radiology Report} A written explanation of an imaging study. It describes why the study was done, the important findings, and their meaning; urgent or unexpected findings should be communicated promptly to the treating team.

\medterm{Radiolucent} Describing a material that lets more X-rays pass through and therefore appears darker on an X-ray or CT image. Air is highly radiolucent; some stones, foreign bodies, and abnormal tissues can also have this appearance.

\medterm{Radiomimetic} Radiomimetic describes a chemical or treatment that produces cellular or DNA damage resembling the effects of ionizing radiation.

\medterm{Radionuclide} An unstable form of an atom that releases radiation as it changes into a more stable form. In medicine, radionuclides can be attached to substances for imaging or used to deliver treatment to selected tissues.

\medterm{Radionuclide Angiography} An imaging test in which a small amount of radioactive tracer is used to follow blood flow through the heart or blood vessels. A camera detects the tracer and produces information about circulation or heart pumping.

\medterm{Radionuclide Cystogram} A nuclear medicine examination that tracks radiolabeled urine to detect vesicoureteral reflux or evaluate bladder emptying.

\medterm{Radionuclide Generator} A device that produces a short-lived radioactive substance from a longer-lasting radioactive substance. It supplies certain tracers for nuclear scans or treatment when needed and must be operated under strict safety and quality controls.

\medterm{Radionuclide Heart Scan} A nuclear medicine study that uses a radioactive tracer to evaluate myocardial blood flow, heart muscle viability, or ventricular function.

\medterm{Radionuclide Imaging} A scan that uses a small amount of radioactive tracer to show the structure or activity of an organ. A camera detects the tracer, helping evaluate blood flow, bone, thyroid, kidney, heart, or cancer-related changes.

\medterm{Radionuclide Imaging Or Scan} A scan that detects radiation from a small amount of radioactive tracer given to a person. The tracer collects in selected organs or tissues, allowing clinicians to examine blood flow, structure, or function and identify some diseases.

\synonyms
Scintigraphy; Nuclear Medicine Scan; Gamma Scan

\medterm{Radionuclides} Radioactive forms of atoms used in nuclear medicine. Attached to a suitable substance, they can show how an organ works on a scan or deliver radiation to selected diseased tissue during treatment.

\medterm{Radionuclide Therapy} Treatment that delivers radiation from inside the body using a radioactive substance directed toward selected tissue. It can treat some thyroid diseases, cancers, or painful bone lesions. The medicine, dose, precautions, and side effects depend on the target and the person’s kidney, blood, and overall health.

\medterm{Radionuclidic Purity} The proportion of radioactivity in a preparation that comes from the intended radioactive atom rather than unwanted radioactive contaminants. High purity helps ensure the expected scan or treatment and limits unnecessary radiation to other tissues.

\medterm{Radio-Opaque} Blocking X-rays and appearing relatively white on a radiograph; radiopaque is the modern spelling.

\medterm{Radiopaque} Describing a material that attenuates X-rays strongly and therefore appears relatively
white on a radiograph or CT image. Bone, metal, and some contrast agents are radiopaque.

\medterm{Radiopharmaceutical} A medicine containing a small amount of radioactive material. It can show how an organ or tissue works on a scan or deliver radiation to selected diseased cells; its distribution determines the information or treatment it provides.

\medterm{Radiopharmaceutical Extravasation} Leakage of a radioactive medicine from the vein into nearby tissue during an injection. It can reduce the accuracy of a scan and, with treatment doses, injure local tissue; staff assess and document the event.

\medterm{Radiopharmaceutical Labeling Kit} A prepared set of sterile ingredients used to attach a radioactive atom to a substance that travels to a particular organ or follows a body process. The finished product is tested before it is given for imaging.

\medterm{Radiopharmaceutical Quality Control} Testing performed to verify identity, radionuclidic purity, radiochemical purity, pH,
sterility, endotoxin limits, and other release specifications before clinical
administration. Requirements depend on the preparation and jurisdiction.

\medterm{Radiopharmacy} The preparation, testing, storage, and dispensing of medicines containing radioactive materials. Radiopharmacy requires accurate activity measurement, cleanliness, sterility checks, quality control, safe handling, and protection of workers and patients.

\medterm{Radioprotective Agent} A medicine or substance intended to reduce injury to healthy tissue from ionizing radiation. It may limit chemical damage or support cell repair, but its benefit and risks depend on the type, dose, and treatment setting.

\medterm{Radio-Radial Delay} A palpable asymmetry or latency between the right and left radial arterial pulses when examined simultaneously. It indicates proximal arterial obstruction affecting one subclavian artery, as seen in aortic dissection, subclavian steal syndrome, or thoracic aortic aneurysm.

\medterm{Radiosensitization} Making cells more vulnerable to radiation damage. Some medicines or treatments are combined with radiotherapy to increase damage to cancer cells, while treatment planning aims to limit harm to nearby healthy tissue.

\medterm{Radiosurgery} Radiosurgery delivers highly focused radiation to a small target, usually in the brain or body, without making a conventional surgical incision.

\medterm{Radiotherapy} Treatment using high-energy radiation to damage and control abnormal cells, especially cancer. It may be given from outside or inside the body and can cause temporary or lasting effects in nearby healthy tissue.

\medterm{Radiotherapy Bolus} A layer of material placed on the skin during selected radiation treatments so more radiation reaches the skin surface. It is positioned carefully to deliver the planned dose to the target area.

\medterm{Radiotherapy Dosage} Radiotherapy dosage is the amount of ionizing radiation prescribed for a treatment target, usually divided into carefully planned fractions.

\medterm{Radiotracer} A small amount of radioactive material attached to a substance that travels through a specific body process. A scanner detects where it goes, helping show blood flow, organ function, bone activity, or selected disease-related changes.

\medterm{Radioulnar Joint} A joint between the radius and ulna bones of the forearm. It allows the hand and forearm to turn palm-up and palm-down.

\medterm{Radium} A naturally radioactive chemical element that releases penetrating radiation as it breaks down. Internal exposure can damage bone and other tissues and increase cancer risk, so handling requires strict safety controls.

\medterm{Radium-223} A radioactive medicine that behaves like calcium and collects in areas of increased bone activity. It is used for selected prostate cancers that have spread to bone and may cause low blood counts or other side effects.

\medterm{Radius} The forearm bone on the thumb side. It joins the elbow and wrist and rotates around the ulna, allowing the palm to turn up or down.

\medterm{Radius Fracture} A break in the radius, one of the two bones of the forearm, often occurring near the wrist after a fall onto an outstretched hand.

\medterm{Radon} A naturally occurring radioactive gas that can accumulate indoors and increase lung-cancer risk when inhaled over time.

\medterm{Ragged Red Fibers} An appearance seen in a muscle biopsy when abnormal mitochondria build up around muscle fibers. It supports the diagnosis of some mitochondrial muscle diseases but is not specific to one condition. The biopsy result must be interpreted with symptoms, blood tests, genetic testing, and other findings.

\synonyms
RRF; Ragged-Red Muscle Fibers

\medterm{Raised Hemidiaphragm} A raised hemidiaphragm is elevation of one side of the diaphragm, caused by abdominal pressure, lung volume loss, phrenic-nerve dysfunction, or diaphragmatic disease.

\medterm{Rale} A rale, or crackle, is an abnormal discontinuous breath sound produced by fluid, collapsed airways, or opening of small airways during breathing.

\medterm{Rales} Discontinuous, adventitious non-musical respiratory lung sounds heard on auscultation, generated by the sudden explosive reopening of collapsed small airways and alveoli containing fluid or exudate; categorized into fine crackles (high-pitched, dry) and coarse crackles (low-pitched, wet).

\synonyms
Crackles; Crepitations

\medterm{Raloxifene} A medicine that acts like estrogen in bones but blocks some estrogen effects elsewhere. It may reduce spine fractures and selected breast-cancer risk after menopause, but can increase blood-clot risk.

\medterm{Raloxifene Hydrochloride} Raloxifene hydrochloride is a selective estrogen-receptor modulator used to prevent or treat osteoporosis and reduce breast-cancer risk in selected postmenopausal women.

\medterm{Raltegravir Potassium} Raltegravir potassium is an HIV integrase inhibitor used with other antiretroviral medicines to suppress HIV replication.

\medterm{Raltitrexed} A chemotherapy medicine used for some colorectal cancers. It blocks a cell process needed to make DNA, slowing the growth of cancer cells; possible effects include diarrhea, mouth sores, low blood counts, tiredness, and liver problems.

\medterm{Ramipril} Ramipril is an angiotensin-converting-enzyme inhibitor used for hypertension, heart failure, and cardiovascular risk reduction; cough, hyperkalaemia, renal-function changes, angioedema, and fetal toxicity are important risks.

\medterm{Rampant Caries} Severe tooth decay affecting many teeth or surfaces at once. It may be linked with frequent sugar, dry mouth, vomiting, poor oral hygiene, limited dental care, medicines, or other health problems.

\medterm{Ramsay Hunt Syndrome} Reactivation of the shingles virus near the facial nerve, causing severe ear pain, a blistering rash in or around the ear or mouth, and weakness on one side of the face. Hearing loss or dizziness may occur; early treatment can improve recovery.

\medterm{Ramsay Hunt Syndrome Type I} A rare neurological disorder characterized by progressive myoclonus epilepsy, intentional cerebellar tremor, progressive ataxia, and mild cognitive decline. Historically termed dyssynergia cerebellaris myoclonica; distinct from Ramsay Hunt syndrome type II (herpes zoster oticus).

\synonyms
Dyssynergia Cerebellaris Myoclonica; Dentatorubral Tremor Syndrome

\medterm{Ramsay Hunt Syndrome Type II} Reactivation of latent varicella-zoster virus within the geniculate ganglion of the facial nerve (cranial nerve VII), presenting with the clinical triad of acute ipsilateral peripheral facial paralysis, otalgia, and a painful vesicular eruption on the external auditory canal, auricle, or soft palate.

\synonyms
Herpes Zoster Oticus

\medterm{Ramucirumab} A monoclonal antibody that blocks vascular endothelial growth factor receptor 2 and is used in selected advanced cancers, including gastric, colorectal, lung, and liver cancers.

\medterm{Ramus} A ramus is a branch-like extension of a bone, nerve, artery, or other structure, such as the mandibular ramus.

\medterm{Random Allocation} Assignment of participants to study groups by chance to reduce selection bias and balance known and unknown factors.

\medterm{Randomised Controlled Trial} A study in which participants are assigned by chance to interventions or controls and outcomes are compared.

\medterm{Randomization Validation} Checking that a clinical trial's assignment process truly places participants into study groups by chance and is implemented as planned. This helps reduce selection bias and protects the reliability of comparisons between treatments.

\medterm{Randomized Controlled Trial} An experimental study in which participants are assigned by chance to different interventions or a control condition and followed for specified outcomes.

\medterm{Range Of Motion} The amount a joint can move in different directions. It may be measured by a healthcare professional or by the person moving alone; pain, swelling, stiffness, weakness, injury, or a shortened muscle can reduce it.

\medterm{Ranibizumab} An intravitreal monoclonal antibody fragment that inhibits vascular endothelial growth factor and treats neovascular macular degeneration, diabetic macular edema, retinal vein occlusion, and related retinal disease.

\medterm{Ranke Complex} The late, healed, fibrocalcified sequela of primary pulmonary tuberculosis, consisting of a calcified subpleural Ghon parenchymal focus and calcified ipsilateral draining hilar or mediastinal lymph nodes visible on chest radiography.
\synonyms
Calcified Primary Tuberculous Complex; Healed Ghon Complex

\medterm{Ranolazine} Ranolazine is an antianginal drug that inhibits the late sodium current in cardiac cells and reduces chronic angina without substantially lowering heart rate or blood pressure; QT prolongation and drug interactions require caution.

\medterm{Ranson Criteria} An older scoring system that estimates the severity and risk of death in acute pancreatitis using clinical and blood-test findings at admission and again after 48 hours. It is less commonly used now because newer scores and organ-failure assessment may be more practical.

\medterm{RANTES} RANTES, now called CCL5, is a chemokine that attracts immune cells and participates in inflammation, antiviral defense, and immune signaling.

\medterm{Ranula} A mucocele (mucus extravasation cyst) of the sublingual salivary gland, presenting as a
blue-translucent, fluctuant swelling in the floor of the mouth, typically lateral to the
lingual frenulum. A ranula results from trauma or obstruction of a sublingual gland
duct, causing mucus accumulation in the sublingual space.

\medterm{Raoultella} Raoultella is a genus of gram-negative bacteria that can cause opportunistic urinary, respiratory, wound, or bloodstream infections.

\medterm{Rape} Nonconsensual sexual penetration or sexual violence; it is a serious crime and a medical, safeguarding, and trauma-care concern.

\medterm{Raphe} A seam-like line or ridge where structures join, such as the median raphe or raphe nuclei of the brainstem.

\medterm{Raphe Nuclei} The raphe nuclei are brainstem groups of neurons that produce serotonin and influence mood, pain, sleep, arousal, and autonomic functions.

\medterm{Rapid Eye Movement Sleep} A sleep stage with rapid eye movements, vivid dreaming, and active brain patterns. Most skeletal muscles are temporarily relaxed during it, while breathing and heart rate can become more variable.
\synonyms
REM Sleep; Paradoxical Sleep

\medterm{Rapid Eye Movement Sleep Behavior Disorder} Alternate name; see \textit{REM sleep behavior disorder}.

\medterm{Rapid Gastric Emptying} Abnormally rapid movement of food from the stomach into the small intestine, sometimes called dumping syndrome.

\medterm{Rapid Heartbeat} Alternate name; see \textit{Tachycardia}.

\medterm{Rapidly Progressive Glomerulonephritis} Severe inflammation of the kidney’s filtering units that causes kidney function to worsen over days or weeks. It may cause blood or protein in urine, swelling, high blood pressure, or reduced urine.

\medterm{Rapid Plasma Reagin Test} A blood test that detects antibodies associated with syphilis. It is useful for screening and monitoring treatment, but false-positive and false-negative results occur, so reactive results need a confirmatory treponemal test.

\medterm{Rapid Sequence Induction} A method of giving anesthetic and muscle-relaxing medicines to place a breathing tube quickly while reducing the chance of vomit entering the lungs. It is used when urgent airway control is needed and requires preparation for difficult intubation.

\medterm{Rapid Sequence Intubation} An emergency airway management technique involving the near-simultaneous administration of a rapid-acting sedative induction agent (such as etomidate or propofol) and a neuromuscular blocking paralytic agent (succinylcholine or rocuronium) without intervening bag-mask ventilation, designed to minimize aspiration risk in an unfasted patient.

\synonyms
RSI

\medterm{Rapid-Sequence Intubation} An emergency method for placing a breathing tube quickly after giving medicines that induce unconsciousness and temporarily relax the muscles. It is used when airway protection is urgent, with oxygen preparation, backup equipment, blood-pressure monitoring, and confirmation that the tube is correctly positioned.

\medterm{Rapid Shallow Breathing} Breathing that is faster than normal but moves only small amounts of air each time. It may occur with lung disease, pain, anxiety, metabolic illness, or tiring breathing muscles and can be a warning sign of respiratory failure.

\medterm{Rapid Strep Test} A test performed on a throat swab to detect Group A Streptococcus. It provides a rapid
result, but a negative result may require confirmatory testing depending on the patient's age, symptoms, and local
guidance.

\medterm{Rare Disease} A condition affecting relatively few people in a defined population; many rare diseases are genetic, require specialized diagnosis, and may have limited evidence or treatment options.

\medterm{Rarefaction} A reduction in density or pressure, such as the low-pressure part of a sound wave. In medicine, understanding rarefaction helps explain ultrasound imaging, sound transmission, and pressure changes in tissues.

\medterm{Rasburicase} A recombinant urate-oxidase enzyme that converts uric acid to the more soluble allantoin. It is used
in selected patients at risk of or experiencing tumor-lysis-related hyperuricemia. It can
cause serious hypersensitivity or hemolysis in people with glucose-6-phosphate
dehydrogenase deficiency and requires appropriate screening and monitoring.

\medterm{Rash} A visible change in skin color, texture, or surface, such as spots, bumps, blisters, or scaling. Causes include infection, allergy, medicines, inflammation, and other disease; rapid or severe rashes need assessment.
\synonyms
Exanthem; Skin Eruption; Dermatosis

\medterm{Rasmussen Encephalitis} A rare, progressive, chronic immune-mediated neuro-inflammatory disorder occurring predominantly in children, characterized by unilateral hemispheric inflammation, drug-resistant focal motor seizures (frequently epilepsia partialis continua), progressive contralateral hemiparesis, cognitive deterioration, and progressive unilateral cortical atrophy on neuroimaging.

\synonyms
Chronic Focal Encephalitis; Rasmussen Syndrome

\medterm{Rasp} A roughened surgical instrument used to file, shape, or smooth bone or other hard tissue.

\medterm{Ras Peptide} A small protein fragment related to Ras proteins, which help control cell growth and division. Changes in Ras signaling can contribute to cancer.

\medterm{RAST} An older name for a blood test that measures IgE antibodies directed against specific allergens. Modern laboratories usually use newer specific-IgE immunoassays; results show sensitization and must be interpreted with symptoms and exposure history.

\synonyms
Radioallergosorbent Test; Specific IgE Test

\medterm{Rastelli Procedure} Open-heart surgery for certain complex birth defects involving the heart’s major arteries, a hole between the lower chambers, and a narrowed route to the lungs. A patch directs blood from the left ventricle to the aorta, and a tube with a valve connects the right ventricle to the lung arteries. The tube may need later replacement.

\synonyms
Rastelli Operation; Rastelli Repair

\medterm{Rat Bite Fever} A systemic bacterial infection caused by Streptobacillus moniliformis or Spirillum
minus, transmitted through rat bites or contact with rat-contaminated food or water. It
presents with relapsing fever, rash, polyarthritis, and lymphadenopathy. The
characteristic rash is petechial or maculopapular on the extremities.

\medterm{Rat-Bite Fever} Alternate name; see \textit{Rat Bite Fever}.

\medterm{Rathke Cleft Cyst} A benign, non-neoplastic, mucin-filled sellar or suprasellar epithelial cyst derived from remnants of the embryonic Rathke pouch, which may compress the optic chiasm (causing bitemporal hemianopia) or normal pituitary gland (producing hypopituitarism).
\synonyms
RCC; Pituitary Cleft Cyst

\medterm{Rathke Pouch} An embryonic ectodermal diverticulum arising from the roof of the stomodeum (primitive oral cavity) that ascends toward the developing brain to form the anterior pituitary gland (adenohypophysis: pars distalis, pars intermedia, pars tuberalis); incomplete obliteration gives rise to craniopharyngiomas or Rathke cleft cysts.

\medterm{Ratio} A ratio compares two quantities by division, such as a measurement of risk, concentration, or body proportions.

\medterm{Rationalization} A psychological defense in which a person creates plausible explanations for behavior or feelings while avoiding the underlying motive or conflict.

\medterm{Rat Lungworm Disease} Infection caused by the rat lungworm parasite Angiostrongylus cantonensis, usually after eating contaminated raw or undercooked produce or animals. It can cause eosinophilic meningitis with headache, neck stiffness, or neurologic symptoms.

\medterm{RATS} Alternate name; see \textit{Robotic-Assisted Thoracoscopic Surgery}.

\medterm{Ravuconazole} Ravuconazole is a triazole antifungal medicine studied and used in some settings for invasive fungal infections.

\medterm{Raw Egg} Raw eggs may contain Salmonella and are safest when properly refrigerated, handled hygienically, and cooked unless pasteurized products are used.

\medterm{Raynaud Phenomenon} Episodes in which small blood vessels in the fingers or toes suddenly narrow after cold or stress. The digits may turn white, blue, then red and become numb or painful; severe or new symptoms may indicate another disease.

\medterm{Raynaud's Disease} The primary form of Raynaud phenomenon, in which small blood vessels in the fingers or toes narrow episodically without an associated underlying disease. Cold or emotional stress may trigger attacks that cause numbness, tingling, or pain and a characteristic change in color from white to blue and then red as blood flow returns.

\medterm{Raynaud's Phenomenon} Episodic constriction of small arteries in the fingers or toes, causing color changes, numbness,
or pain after cold exposure or emotional stress.

\synonyms
Raynaud Disease; Raynaud Syndrome; Vasospastic Attack

\medterm{Raynaud’s Syndrome} Episodes in which fingers or toes change color and become cold, numb, or painful after cold exposure or stress because small blood vessels narrow. It may occur alone or with autoimmune disease.

\medterm{Raynaud Syndrome} Episodes in which small blood vessels in the fingers or toes narrow after cold or stress, causing numbness, pain, and color changes that may progress from white to blue to red. It may occur alone or result from another disease; severe attacks can damage tissue.

\medterm{Razor Bumps} Inflamed bumps caused by hairs curling back into the skin after shaving or other hair removal. They are more common with coarse or curly hair and may cause itching, pain, dark marks, or scarring; avoiding very close shaving and using proper shaving technique can reduce recurrence.

\medterm{RBC Count} A measurement of the number of red blood cells in a volume of blood, interpreted with hemoglobin, hematocrit, and red-cell indices to evaluate anemia or erythrocytosis.

\medterm{RBC Indices} Measurements describing the size, hemoglobin content, and variation of red blood cells. They help classify anemia and interpret a complete blood count alongside hemoglobin and other results.

\medterm{RBC Nuclear Scan} An imaging test using a small amount of the person’s blood labeled with a radioactive tracer. It can help locate ongoing bleeding in the digestive tract or assess blood flow through the heart.

\medterm{RBC Urine Test} A urine examination detecting or counting red blood cells, used to evaluate hematuria from urinary-tract infection, stones, glomerular disease, tumors, trauma, or other causes.

\medterm{RDA} Recommended dietary allowance, the daily nutrient intake sufficient for nearly all healthy people in a defined group.

\synonyms
Recommended Dietary Allowance; Recommended Daily Allowance

\medterm{Reabsorption} The transport of a substance from a body cavity or tubular fluid back into tissue or the bloodstream, such as renal tubular recovery of water and filtered solutes.

\medterm{Reaction} A response of the body or mind to a stimulus, illness, medicine, treatment, or event. Examples include an allergic reaction, an immune reaction, a treatment response, and a psychological reaction; the cause and severity determine its importance.

\medterm{Reactivation Tuberculosis} Active tuberculosis developing when a previously contained infection becomes active again, often years later. It commonly affects the upper lungs and may cause a long-lasting cough, coughing blood, fever, night sweats, weight loss, and spread to others through the air.

\medterm{Reactive Arthritis} Painful inflammation of joints that develops during or after certain intestinal or genital infections. It may also cause eye inflammation, skin changes, or urinary symptoms, and the triggering infection may already have cleared.

\synonyms
Reiter Syndrome; Post-Infectious Arthritis

\medterm{Reactive Attachment Disorder} A rare childhood condition in which a child persistently avoids or barely responds to comfort from caregivers. It is linked to severe neglect or unstable care early in life and may involve sadness, irritability, limited trust, and difficulty forming relationships.

\medterm{Reactive Gastropathy Pathology} Gastric mucosal injury characterized by foveolar hyperplasia, edema, smooth-muscle fibers in the lamina propria, vascular congestion, and relatively little active inflammation. Chemical injury, bile reflux, and medicines are common associations.

\medterm{Reactive Hypoglycemia} A fall in blood sugar occurring within a few hours after eating. It may cause shakiness, sweating, hunger, weakness, dizziness, or confusion; repeated episodes need assessment to confirm low sugar and identify the cause.

\medterm{Reactive Oxygen Species} Unstable oxygen-containing molecules produced during normal cell activity, infection, or inflammation. Small amounts can help cells signal and defend against germs, but excess amounts can damage cell membranes, proteins, and DNA.

\medterm{Reactive Perforating Collagenosis} A skin condition in which damaged collagen from deeper skin is pushed through the surface. It causes small, firm, very itchy bumps with a central plug, often on the arms, legs, or hands, and may be associated with diabetes or kidney disease.

\medterm{Reality Testing} The ability to recognize whether perceptions, thoughts, and beliefs match shared reality. Clinicians assess it during a mental-status examination; impaired reality testing can occur with psychosis, severe mood episodes, delirium, or some brain disorders.

\medterm{Rebound Effect} Worsening or return of symptoms after stopping or reducing a treatment that had been suppressing them.

\medterm{Rebound Headache} Alternate name; see \textit{Medication overuse headaches}.

\medterm{Rebound Insomnia} Return or worsening of insomnia after stopping a sleep medicine.

\medterm{Rebound Tenderness} A sharp increase in stomach pain when pressure on the belly is suddenly released, which signals that the inner lining of the abdomen is inflamed.

\medterm{Recalcitrant} Resistant to treatment, control, healing, or change; the term often describes a persistent disease or infection.

\medterm{Receiver Operating Characteristic} A graph showing how a diagnostic test’s ability to detect disease changes as its cutoff is changed. It compares true-positive and false-positive results; the area under the curve summarizes overall discrimination, not clinical usefulness by itself.

\medterm{Recent Memory} Recent memory is the ability to recall information or events from the near past and can be impaired by aging, delirium, brain disease, depression, or medicines.

\medterm{Receptor} A protein on or inside a cell that recognizes a specific signal, such as a hormone, nerve chemical, or medicine, and starts or changes a cell response.

\medterm{Receptor Binding} The attachment of a signaling molecule, medicine, or other ligand to a receptor on or inside a cell. Binding can activate, block, or change the receptor and its downstream cellular response.

\medterm{Receptor Desensitization} A reduced cell response after a receptor is stimulated continuously or repeatedly. The change can develop quickly and helps prevent excessive signaling, but it may also reduce the effect of a medicine or natural body signal.

\medterm{Receptor Downregulation} A decrease in the number or sensitivity of cell receptors after a signal stimulates them for a long time. The cell may remove receptors from its surface or break them down, becoming less responsive to the signal or medicine.

\medterm{Receptors} Proteins or structures on or inside cells that recognize specific signals and start responses such as growth, movement, or hormone action.

\medterm{Receptor Tyrosine Kinases} Cell-surface proteins that receive signals from hormones and growth factors. When switched on, they activate internal pathways controlling growth, survival, movement, and repair; abnormal activity can contribute to cancer and other disease.

\medterm{Receptor Upregulation} An increase in the number or sensitivity of cell receptors after prolonged blocking of a signal or reduced stimulation. This adjustment can make a cell respond more strongly when the signal or medicine is later present.

\medterm{Recession Of Medial Rectus Muscle} An eye operation that moves the medial rectus muscle backward to weaken its pull and correct certain inward eye deviations.

\medterm{Recessive} Describing a gene change whose effect usually appears only when a person has two altered copies, one inherited from each parent. A person with one copy may carry it without having the condition.

\medterm{RECIST} Standard rules for judging how a solid tumor changes on scans during treatment. Using measured lesions, they classify disease as completely gone, partly reduced, stable, or worsened; the categories support research and clinical decisions but do not replace overall assessment.

\medterm{Recognizing Medical Emergencies} Recognizing a medical emergency includes identifying severe breathing difficulty, chest pain, stroke signs, major bleeding, seizure, collapse, poisoning, or altered consciousness and calling emergency services.

\medterm{Recognizing Teen Depression} Recognizing depression in adolescents involves looking for persistent low mood or irritability, loss of interest, sleep or appetite changes, poor concentration, hopelessness, or self-harm thoughts.

\medterm{Recombinant} Produced by combining genetic material from different sources or by using engineered DNA in a cell or organism.

\medterm{Recombinant DNA} DNA made by joining genetic material from different sources. It is used in biotechnology to produce medicines such as insulin and growth hormone, develop some vaccines, and create vehicles for gene therapy.

\medterm{Recombinant Protein} A protein produced by putting its gene into a laboratory cell and growing the cell so it makes the protein. This method supplies some hormones, vaccines, antibodies, enzymes, and research materials.

\medterm{Recombinant Tissue Plasminogen Activator} A laboratory-made form of tissue plasminogen activator that helps convert plasminogen to plasmin and dissolve blood clots. It is used in selected emergencies, such as some ischemic strokes, heart attacks, or pulmonary emboli, when benefits outweigh bleeding risk.

\synonyms
rtPA; Alteplase; Recombinant tPA

\medterm{Recombinase} An enzyme that cuts, joins, or rearranges DNA segments. It helps cells exchange genetic material, repair some DNA damage, and assemble the gene parts needed to make a varied immune response.

\medterm{Recombinational Repair} A way cells repair a serious DNA break by using a matching, undamaged DNA copy as a guide. It restores genetic information accurately; inherited problems with this repair can increase cancer risk and treatment sensitivity.

\medterm{Recommended Daily Allowance} The Recommended Dietary Allowance is the average daily nutrient intake sufficient for nearly all healthy people in a specified age and sex group.

\medterm{Recommended Dietary Allowance} The average daily intake expected to meet the nutrient needs of nearly all healthy people in a defined group.

\synonyms
RDA

\medterm{Record} A health record is a documented account of a person’s history, findings, tests, diagnoses, treatments, decisions, and outcomes.

\medterm{Recovery Position} A side-lying position used for an unconscious but normally breathing person to help keep the airway open and reduce aspiration risk.

\medterm{Recreation Therapy} Recreation therapy uses adapted games, arts, exercise, and leisure activities to support physical, emotional, cognitive, or social functioning.

\medterm{Recrudescence} Recurrence of symptoms or infection after apparent improvement, often from persistence of the original organism or disease process.

\medterm{Rectal} Rectal means relating to the rectum, the final section of the large intestine before the anal canal.

\medterm{Rectal Biopsy} Removal of a small piece of tissue from the rectum for examination under a microscope. It can help investigate missing nerve cells, inflammation, infection, abnormal growth, or other rectal disease.

\medterm{Rectal Bleeding} Passage of bright-red or dark blood from the rectum, caused by hemorrhoids, fissures,
polyps, inflammatory bowel disease, diverticular disease, or cancer. Persistent, heavy,
or unexplained bleeding requires evaluation.

\synonyms
Hematochezia; Lower Gastrointestinal Bleeding

\medterm{Rectal Cancer} Cancer developing in the rectum, the final part of the large intestine. Symptoms may include bleeding, altered bowel habits, pain, or weight loss; treatment may combine surgery, radiation, chemotherapy, immunotherapy, or targeted therapy.

\synonyms
Rectal Carcinoma; Adenocarcinoma of the Rectum

\medterm{Rectal Carcinoma} Rectal carcinoma is cancer arising in the rectum and may cause bleeding, altered bowel habits, pain, or anemia.

\medterm{Rectal Culture} Laboratory culture of a rectal specimen to detect selected bacterial pathogens or colonization, depending on the clinical indication.

\medterm{Rectal Disease} A disorder affecting the rectum, including inflammation, infection, bleeding, obstruction, prolapse, functional disease, or neoplasm.

\medterm{Rectal Fistula} A rectal fistula is an abnormal tract connecting the rectum with skin or another organ, often causing drainage, infection, or passage of stool through an abnormal route.

\medterm{Rectal Hemorrhage} Bleeding from the rectum, which may result from hemorrhoids, fissure, inflammation, ischemia, polyps, cancer, trauma, or other gastrointestinal disease.

\medterm{Rectal Inflammation} Alternate name; see \textit{Proctitis}.

\medterm{Rectal Itching} Alternate name; see \textit{Anal itching}.

\medterm{Rectal Mucositis} Rectal mucositis is inflammation and injury of the rectal lining, often after radiation, chemotherapy, infection, inflammatory disease, or trauma.

\medterm{Rectal Necrosis} Rectal necrosis is death of rectal tissue from inadequate blood flow, severe inflammation, infection, pressure, or another destructive process.

\medterm{Rectal Neoplasm} An abnormal growth in the rectum that may be noncancerous, precancerous, or cancerous. It can cause bleeding, altered bowel habits, pain, anemia, or no symptoms and may require examination and biopsy.

\medterm{Rectal Obstruction} Impaired passage through the rectum caused by fecal impaction, stricture, tumor, foreign body, inflammation, or functional outlet disorder.

\medterm{Rectal Perforation} A tear or hole through the rectal wall that can release bowel contents and cause peritonitis, sepsis, or severe bleeding.

\medterm{Rectal Polyposis} Rectal polyposis is the presence of multiple polyps in the rectum, which may be sporadic or part of an inherited syndrome and can increase cancer risk.

\medterm{Rectal Polyps} Alternate name; see \textit{Colon Polyps}.

\medterm{Rectal Prolapse} The rectum slipping downward and partly or completely coming through the anus. It may cause a visible bulge, bleeding, mucus, discomfort, or difficulty controlling or passing stool and can result from weakened pelvic support, straining, or nerve problems.

\medterm{Rectal Prolapse Repair} A surgical or procedural treatment that restores a prolapsed rectum to its normal position and secures or removes affected tissue according to the prolapse and patient factors.

\medterm{Rectal Route} The rectal route administers a medicine or fluid through the anus into the rectum, where it may act locally or be absorbed systemically.

\medterm{Rectal Stenosis} Abnormal narrowing of the rectal lumen caused by inflammation, scarring, surgery, radiation, ischemia, trauma, or neoplasm.

\medterm{Rectal Tenesmus} Rectal tenesmus is a persistent or urgent feeling of needing to pass stool despite an empty or nearly empty rectum, caused by inflammation, infection, mass, or motility disorder.

\medterm{Rectal Ulcer} Alternate name; see \textit{Solitary rectal ulcer syndrome}.

\medterm{Rectocele} Bulging of the rectum into the back wall of the vagina because supporting tissues have weakened. It may cause pelvic pressure, difficulty emptying the bowel, or a vaginal bulge.

\medterm{Rectosigmoid Junction} The point where the sigmoid colon joins the rectum.

\medterm{Rectosigmoid Neoplasm} A rectosigmoid neoplasm is an abnormal growth where the sigmoid colon joins the rectum and may be benign or malignant.

\medterm{Rectouterine Pouch} The rectouterine pouch, or pouch of Douglas, is the peritoneal recess between the uterus and rectum where fluid or disease may collect.

\medterm{Rectovaginal Fistula} Abnormal connection between the rectum and vagina causing passage of stool or gas
through the vagina. Causes include obstetric injury particularly third and fourth degree
perineal tears, Crohn disease, radiation therapy, and surgical complications. Presents
with vaginal discharge of stool, flatus, or recurrent vaginal infections.

\medterm{Rectum} The final straight section of the large intestine, approximately 12-15 cm long,
extending from the sigmoid colon to the anal canal. The rectum stores feces before
defecation.

\medterm{Rectus Abdominis} The rectus abdominis is a paired abdominal muscle that flexes the trunk, stabilizes the pelvis, and helps compress abdominal contents.

\medterm{Rectus Abdominis Muscle} The paired vertical muscles at the front of the abdomen, commonly called the “six-pack.” They bend the trunk forward, stabilize the pelvis, and help increase pressure inside the abdomen; stretching can cause separation during pregnancy or after major weight change.

\medterm{Rectus Abdominis Sheath} A fibrous aponeurotic envelope formed by the abdominal wall muscles and surrounding the
rectus abdominis muscle. It contains the superior and inferior epigastric vessels and
related nerves.

\medterm{Rectus Sheath Hematoma} A collection of blood within the sheath surrounding the rectus abdominis muscle, often caused by muscle strain, coughing, trauma, anticoagulation, or vessel injury. It may cause localized abdominal pain and a tender mass.

\medterm{Recur} To recur means to return or happen again, as symptoms, disease, infection, or a tumor after a period of improvement or treatment.

\medterm{Recurrence} The return of a disease after it improved or seemed to disappear. Cancer may return at the original site, nearby lymph nodes, or another part of the body; recurrence does not always mean that treatment is no longer possible.

\medterm{Recurrence-Free Survival} The length of time after treatment during which a disease does not return. In cancer studies, it usually ends when cancer is found again or when a person dies, depending on the study rules. It measures control of disease, not a guarantee that treatment has cured it.

\medterm{Recurrence Risk} The chance that a disease or condition will return after treatment or occur again in a future pregnancy or family member. The estimate depends on the condition, previous history, test results, inheritance pattern, and other risk factors.

\medterm{Recurrent} Happening again after a period of improvement, remission, or apparent absence. A recurrent condition may return at the original site or elsewhere; examples include repeated urinary infections, repeated seizures, or cancer returning after treatment.

\medterm{Recurrent Breast Cancer} Breast cancer that returns after a period when tests found no detectable disease. It may recur near the original site or elsewhere and is assessed with examination, imaging, biopsy, and other tests.

\medterm{Recurrent Caries} Tooth decay that develops beneath or beside an existing filling, crown, or other restoration. It may result from leaking edges, plaque, frequent sugar, or dry mouth and can cause sensitivity, pain, or further tooth damage.

\medterm{Recurrent Disease} Return of a disease after a period of improvement, remission, or apparent absence, either at the original site or elsewhere.

\medterm{Recurrent Early Pregnancy Loss} Repeated miscarriages occurring early in pregnancy. Evaluation may consider uterine anatomy, chromosomal factors, antiphospholipid syndrome, endocrine or metabolic disease, parental genetics, and other causes, although no cause is found in many cases.

\medterm{Recurrent Fractures} Repeated bone fractures, suggesting persistent trauma, osteoporosis, bone fragility, metabolic disease, tumor, or an inherited skeletal disorder.

\medterm{Recurrent Laryngeal Nerve} A branch of the vagus nerve supplying most intrinsic muscles of the larynx; injury can cause hoarseness.

\medterm{Recurrent Laryngeal Nerve Palsy} A peripheral cranial neuropathy resulting from surgical injury (e.g., during thyroidectomy), trauma, or mediastinal tumor compression affecting the recurrent laryngeal branch of the vagus nerve (CN X), causing ipsilateral vocal cord immobility, hoarseness, and loss of airway protective reflexes.
\synonyms
RLN Palsy; Vocal Cord Paralysis

\medterm{Recurrent Oral Ulcers} Alternate name; see \textit{Canker sore}.

\medterm{Recurrent Pregnancy Loss} Two or more miscarriages. Possible causes include chromosome changes, an unusual shape of the uterus, hormone or thyroid problems, antiphospholipid syndrome, or other medical conditions, although no cause is found in many people. Evaluation is individualized and may identify treatment that improves the chance of a future healthy pregnancy.

\synonyms
Recurrent Miscarriage; Habitual Abortion; RPL

\medterm{Recurrent Prostate Cancer} Alternate name; see \textit{Prostate cancer recurrence}.

\medterm{Recurrent Respiratory Papillomatosis} A disease in which wart-like growths caused by human papillomavirus develop in the larynx or other breathing passages. The growths can cause hoarseness, noisy breathing, airway narrowing, or repeated recurrence.

\medterm{Recurrent Urinary Tract Infection} Repeated urinary tract infections separated by periods without infection. They may affect the bladder or kidneys and can be related to bacteria that persist, incomplete bladder emptying, stones, urinary blockage, or other urinary-tract problems.

\medterm{Red Birthmarks} Red birthmarks are vascular skin marks present at or soon after birth, ranging from harmless capillary marks to lesions requiring monitoring or specialist treatment.

\medterm{Red Blood Cell} A blood cell that carries oxygen from the lungs to the body's tissues. It contains hemoglobin, has a flattened, indented shape, and lacks a nucleus, which helps it bend through small blood vessels. Red blood cells also carry some carbon dioxide back to the lungs and help keep the blood's acid level balanced. They are made mainly in the bone marrow and normally live for about 120 days.

\medterm{Red Blood Cell Casts} Tube-shaped groups of red blood cells seen in urine. They form inside the kidney's small tubes and strongly suggest bleeding or inflammation in the kidney's filtering units, such as glomerulonephritis or vasculitis.

\synonyms
RBC Casts

\medterm{Red Blood Cell Labeling} A nuclear-medicine technique in which a person’s red blood cells are marked with a small amount of radioactive material. Scans can then show ongoing digestive-tract bleeding or measure blood flow and heart function.

\medterm{Red Bone Marrow} Bone marrow that makes red blood cells, white blood cells, and platelets. It is found mainly in certain bones and in the spongy parts of many bones.

\medterm{Red Cell Distribution Width} A standard automated complete blood count parameter that measures the degree of variation in red blood cell volume (anisocytosis); an elevated RDW aids in distinguishing iron deficiency anemia (high RDW) from uncomplicated thalassemia trait (normal RDW).
\synonyms
RDW; Erythrocyte Anisocytosis Index

\medterm{Red Cell Indices} Standardized automated complete blood count parameters derived from hemoglobin, hematocrit, and erythrocyte count that quantify red blood cell physical properties: mean corpuscular volume (MCV), mean corpuscular hemoglobin (MCH), and mean corpuscular hemoglobin concentration (MCHC).

\medterm{Red Eye} Redness of the white part of the eye caused by irritation, dryness, allergy, infection, injury, inflammation, or increased surface blood flow.

\medterm{Red Flag} A symptom, sign, or test result suggesting a potentially serious condition that requires
urgent assessment or escalation of care. Red flags are context-specific and should not
be used as a substitute for clinical judgment.

\medterm{Red Hepatization} The second pathological stage of acute lobar pneumonia (occurring 2 to 4 days after infection), in which the affected lung lobe becomes consolidated, heavy, and firm ('liver-like') due to massive filling of alveolar spaces with packed erythrocytes, neutrophils, and fibrin strands.
\synonyms
Second Stage of Lobar Pneumonia; Consolidation Stage

\medterm{Red Infarction} Tissue death caused by loss of blood flow in which blood also leaks into the damaged area, making it dark red. It is more likely in organs with a second blood supply or when a vein is blocked.

\synonyms
Hemorrhagic Infarction

\medterm{Red Man Syndrome} An adverse, non-IgE-mediated pseudoallergic hypersensitivity reaction caused by rapid intravenous infusion of vancomycin, in which direct degranulation of mast cells and basophils triggers massive histamine release, producing flushing, erythema, and pruritus of the face, neck, and upper torso.

\synonyms
Vancomycin Flushing Reaction

\medterm{Red-Man Syndrome} An acute, non-immunologic (anaphylactoid) rate-dependent adverse drug reaction caused by rapid intravenous infusion of vancomycin, mediated by direct mast cell and basophil degranulation with histamine release, presenting with erythematous flushing and pruritus of the face, neck, and upper torso.
\synonyms
Vancomycin Flushing Syndrome; Histaminoid Infusion Reaction

\medterm{Red Nucleus} A group of nerve cells in the midbrain that helps coordinate movement through connections with the cerebellum and spinal cord. Damage to nearby pathways can cause weakness, tremor, or abnormal movement.

\medterm{Reductase Deficiency} A disorder or laboratory finding caused by inadequate activity of a reductase enzyme, impairing a specific metabolic or hormone pathway.

\medterm{Reduction} Restoration of a displaced fracture or joint to better alignment, or decrease in amount, size, or activity.

\medterm{Reduction Division} The first division of meiosis, when a reproductive cell's chromosome number is reduced by half. It produces cells with one chromosome from each pair and helps create genetic variation in eggs and sperm.

\medterm{Reduction Mammaplasty} Surgery that removes breast tissue, fat, and excess skin to reduce breast size and relieve symptoms such as neck, shoulder, or back discomfort; it may also reshape the breasts.

\medterm{Redundant Mitral Valve} Excess leaflet tissue of the mitral valve, most commonly the posterior leaflet, which
may prolapse into the left atrium during systole. Often associated with myxomatous
degeneration and mitral valve prolapse. May be asymptomatic or cause mild mitral
regurgitation detected incidentally on echocardiography.

\medterm{Reduviidae} A family of insects including assassin bugs; some species transmit pathogens such as the parasite causing Chagas disease.

\medterm{Reed-Sternberg Cells} Large abnormal immune cells derived from B lymphocytes and commonly found in classic Hodgkin lymphoma. Their appearance in a tissue sample helps pathologists confirm that diagnosis.

\medterm{Reexpansion Pulmonary Edema} A sudden buildup of fluid in a lung after it rapidly expands from a collapsed state, often after drainage of a large or long-standing pneumothorax. It can cause severe breathlessness and low oxygen and needs urgent treatment.

\medterm{Refeeding Syndrome} Dangerous changes in body salts and fluid balance when food is restarted after severe or prolonged undernutrition. Low phosphate, potassium, or magnesium can cause weakness, swelling, abnormal heart rhythms, breathing failure, or death and requires careful medical monitoring.

\medterm{Reference Interval} The range of results usually found in a selected healthy reference group, often covering the middle 95 percent. A result outside it does not by itself prove disease and must be interpreted with symptoms, history, and the laboratory’s own method.

\medterm{Referral} A request for assessment or treatment by another healthcare professional or service. Referrals help patients access specialized expertise, tests, procedures, rehabilitation, or coordinated care for a particular concern.

\medterm{Referred Otalgia} Ear pain felt in an anatomically normal ear caused by sensory nerve sharing between the ear and distant head and neck structures via cranial nerves V3 (auriculotemporal nerve), VII (facial), IX (glossopharyngeal - Jacobson nerve), X (vagus - Arnold nerve), and upper cervical nerves C2-C3; mandates careful examination of the pharynx and larynx for occult malignancy.
\medterm{Referred Pain} Pain felt in an area away from the tissue causing it because nearby or connected nerves share pathways to the brain. Examples include shoulder pain from irritation below the diaphragm or arm and jaw pain during a heart attack.

\medterm{Reflex} A rapid, automatic response to a stimulus, such as pulling the hand away from heat. Reflex testing helps assess pathways involving the nerves, spinal cord, and brain.

\medterm{Reflex Action} A rapid, automatic response to a stimulus that occurs through a nerve pathway called a reflex arc. Reflexes protect the body, maintain posture, and provide information about nerve and spinal-cord function.

\medterm{Reflex Anoxic Seizure} A brief fainting episode that can look like a seizure after sudden pain, fright, or emotional stress. The heart slows briefly, the person becomes pale and unresponsive, and jerking may occur before rapid recovery.

\medterm{Reflex Arc} The nerve pathway that produces a rapid, automatic response. It includes a sensor, a sensory nerve, a processing center in the spinal cord or brain, a motor nerve, and the muscle or gland that responds.

\medterm{Reflex Epilepsy} Reflex epilepsy is epilepsy in which specific stimuli such as flashing light, sound, reading, or touch reliably trigger seizures.

\medterm{Reflexes} Automatic, rapid responses to stimuli carried through nerve pathways called reflex arcs. Testing them helps assess sensory nerves, motor nerves, the spinal cord, and parts of the brain.

\medterm{Reflux} The backward flow of a body fluid, especially stomach contents into the esophagus or urine from the bladder toward the kidneys. Symptoms and risks depend on the type, frequency, and underlying cause.

\medterm{Reflux Esophagitis Pathology} Mucosal injury associated with gastroesophageal reflux, often showing basal-zone hyperplasia, elongation of lamina-propria papillae, intercellular edema, and variable eosinophils or neutrophils.

\medterm{Reflux Laryngitis} Irritation or inflammation of the larynx attributed to reflux of stomach contents into the throat or upper airway. Hoarseness, throat clearing, cough, globus sensation, or throat discomfort may occur, although other causes of these symptoms should also be assessed.
\medterm{Reflux Nephropathy} Renal scarring and chronic kidney damage caused by vesicoureteral reflux, recurrent
urinary infection, or both. It may lead to hypertension, proteinuria, impaired growth,
and chronic kidney disease.

\medterm{Refraction} The bending of light as it passes through the cornea, aqueous humour, lens, and vitreous humour. In an eye examination, refraction determines the optical power needed to focus a clear image on the retina; refractive power is measured in dioptres.

\medterm{Refractive Disorder} A vision problem caused by the eye not focusing light correctly on the retina. Nearsightedness, farsightedness, astigmatism, and presbyopia are common refractive disorders and may be corrected with glasses, lenses, or surgery.

\medterm{Refractive Error} A refractive error occurs when the eye does not focus light correctly on the retina, causing nearsightedness, farsightedness, astigmatism, or presbyopia.

\medterm{Refractive Errors And Astigmatism} Refractive errors are vision disorders caused by the eye's inability to properly focus
light onto the retina, resulting in blurred vision.

\medterm{Refractive Surgery} Eye surgery that changes the shape of the cornea or another part of the eye to reduce short-sightedness, long-sightedness, or astigmatism. It may reduce dependence on glasses but can cause dry eyes, glare, infection, or imperfect vision.

\medterm{Refractometer} A refractometer measures how much light bends through a substance and is used for fluid concentration, optical properties, or eye refraction.

\medterm{Refractometry} Measurement of the eye’s refractive power to determine the optical correction needed for myopia, hyperopia, astigmatism, or presbyopia.

\medterm{Refractory} Resistant to a treatment or not responding adequately to it. The term may refer to cancer or other diseases.

\medterm{Refractory Anemia} An older term for anemia that does not improve with usual treatment. It may refer to a myelodysplastic syndrome or another cause, so evaluation looks for blood-cell disorders, bleeding, nutritional problems, kidney disease, and other conditions.

\medterm{Refractory Cancer} Refractory cancer is cancer that does not respond adequately to a treatment or begins progressing during treatment.

\medterm{Refractory Disease} Disease that does not improve adequately with appropriate standard treatment or that continues to worsen despite therapy. The term does not mean that every possible treatment has failed.

\medterm{Refractory Period} The brief recovery time after a nerve or muscle cell sends an electrical signal. During this time it cannot, or can only partly, send another signal, which helps control the speed and direction of nerve and heart activity.

\medterm{Refrigerant Poisoning} Poisoning from inhaling or contacting a refrigerant gas or liquid. Exposure can cause dizziness, abnormal heart rhythms, suffocation, frostbite, chemical burns, or lung injury.

\medterm{Refsum Disease} A rare inherited disorder in which a fatty substance builds up in the body. It may cause night blindness, loss of smell, nerve damage, poor coordination, hearing loss, and scaly skin.

\medterm{Regadenoson Pharmacological Stress} A medicine used during a heart scan to widen the heart’s arteries when a person cannot exercise adequately. It can cause flushing, headache, chest discomfort, shortness of breath, or abnormal heart rhythms.

\medterm{Regeneration} Growth or replacement of damaged or lost cells and tissues, helping restore structure or function after injury.

\medterm{Regimen} A planned schedule for medicines, treatments, diet, exercise, or other healthcare steps. It states what to do, how much or how often, and for how long, and may be changed according to response or side effects.

\medterm{Regional} Relating to a particular body area, geographic area, or localized distribution rather than the whole body.

\medterm{Regional Anesthesia} Numbing a larger body area by injecting local anesthetic near the nerves supplying it. The person is usually awake, although sedation may be added; the block can provide numbness during surgery and pain relief afterward.

\synonyms
Conduction Anesthesia; Nerve Block Anesthesia

\medterm{Regional Chemotherapy} Treatment that delivers anticancer medicine directly to a body region or organ to produce a high local concentration while reducing, but not eliminating, exposure elsewhere. It is suitable only for selected cancers and carries procedure-related and medicine-related risks.

\medterm{Regional Enteritis} An older name for Crohn disease, a chronic inflammatory bowel disease that can affect any part of the digestive tract.

\medterm{Regional Ileitis} An older name for Crohn disease affecting mainly the terminal ileum. It causes long-term inflammation that may lead to abdominal pain, diarrhea, weight loss, narrowing, fistulas, or nutritional problems.

\medterm{Regulation Of Health Professions} Legal and professional systems that govern healthcare workers' education, licensing, permitted duties, conduct, and accountability. Regulation aims to protect patients, maintain competence, support ethical care, and address unsafe practice.

\medterm{Regulatory Approval} Authorization by a competent regulatory authority for a medicine, device, or procedure to be marketed or used for specified purposes after review of evidence.

\medterm{Regulatory Authority} A regulatory authority is a government or official body that sets and enforces standards for medicines, devices, facilities, professionals, or public health.

\medterm{Regulatory Gene} A gene that controls when, where, or how strongly another gene is used. Changes in regulatory genes or their control regions can alter development, cell growth, metabolism, or disease risk.

\medterm{Regulatory T Cells} Immune cells that limit excessive immune reactions and help prevent attacks against the body’s own tissues. They develop mainly in the thymus or from other T cells in tissues. Problems with their function can contribute to autoimmune disease.

\medterm{Regurgitation} Backward flow of blood through a heart valve that does not close fully. Mild leakage may cause no symptoms; more severe leakage can enlarge or weaken the heart and cause breathlessness, tiredness, or swelling.

\medterm{Rehab} A structured rehabilitation program that helps restore function, independence, participation, or quality of life after illness or injury.

\medterm{Rehabilitation} A coordinated process that helps a person restore or adapt physical, cognitive, psychological, and social
function after illness, injury, surgery, or disability.

\medterm{Rehabilitation Clinic} An outpatient facility where a team helps people regain or adapt movement, communication, thinking, daily activities, or independence after illness, injury, surgery, or disability. Services may include physical, occupational, speech, or other rehabilitation therapies.

\medterm{Rehabilitation Engineering} The use of engineering and technology to improve independence after illness or disability. It includes artificial limbs, braces, wheelchairs, communication aids, home adaptations, and devices that help control the environment.

\medterm{Rehabilitation Motivation} A person’s willingness and ability to take part in rehabilitation and work toward functional goals. Pain, mood, fatigue, cognition, support, confidence, and the expected benefit can all influence participation.

\medterm{Rehabilitation Potential} The expected capacity to improve function, independence, participation, or quality of
life through rehabilitation. It depends on diagnosis, prognosis, cognition, motivation,
comorbidity, social support, and environmental resources.

\medterm{Rehabilitation Unit} A hospital unit for people who need coordinated medical care and intensive therapy after a serious illness or injury, such as stroke, brain or spinal-cord injury, or major surgery. The team works to restore safe movement, daily skills, communication, and independence.

\medterm{Rehydration} Replacement of water and electrolytes lost through dehydration, commonly using oral fluids or intravenous treatment when necessary.

\medterm{Reid Index} A microscope measurement comparing the thickness of mucus-making glands with the wall of a large airway. A higher value may be seen in chronic bronchitis, but it is not used alone to diagnose the condition.
\synonyms
Bronchial Gland-to-Wall Ratio; Reid Histopathologic Index

\medterm{Reifenstein Syndrome} An inherited disorder of androgen action, usually involving partial androgen insensitivity in an individual with a 46,XY chromosome pattern.

\medterm{Reiki} Reiki is a complementary practice involving light touch or hands near the body. Evidence that it treats disease is limited.

\medterm{Reinfection} A new episode of infection by the same or a related pathogen after the original
infection has resolved. Reinfection is distinct from relapse, in which the original
infection persists or returns from an incompletely eradicated focus.

\medterm{Reinke Edema} Chronic polypoid degeneration and diffuse, gelatinous mucoid edema within the superficial lamina propria (Reinke's space) of both true vocal cords, strongly associated with chronic cigarette smoking and vocal strain, presenting with progressive hoarseness and a deepened voice.
\synonyms
Polypoid Corditis; Diffuse Vocal Fold Edema

\medterm{Reinnervation} Restoration of nerve supply to a muscle or tissue after nerve repair, regeneration, or collateral nerve growth.

\medterm{Rejection} An immune response in which the body recognizes transplanted tissue or an organ as foreign and attacks it. Rejection can damage or destroy the transplant, may be sudden or gradual, and is monitored and treated with immune-suppressing medicines.
\synonyms
Allograft Rejection; Graft Rejection

\medterm{Relapse} Return of a disease or its signs and symptoms after a period of improvement or remission. The term is used for cancer, infections, mental-health and substance-use disorders, and other conditions; its precise meaning depends on the disease.

\medterm{Relapse-Free Survival} The time from completion of treatment to recurrence of the cancer or death from any
cause.

\medterm{Relapsing Fever} An acute febrile illness caused by Borrelia species, characterized by recurrent episodes
of fever separated by afebrile intervals.

\medterm{Relapsing Polychondritis} A rare immune disease in which episodes of inflammation damage cartilage. It commonly affects the ears, nose, eyes, joints, or windpipe and may cause pain, swelling, hearing or breathing problems; airway involvement can be life-threatening.

\medterm{Relapsing-Remitting Multiple Sclerosis} A form of multiple sclerosis in which neurological symptoms appear in attacks and then improve partly or completely for periods of time. Symptoms may include vision changes, numbness, weakness, balance problems, or fatigue. Treatment can reduce relapses and new nerve damage, although the course varies between people.

\synonyms
RRMS; Relapsing-Remitting MS

\medterm{Relative Afferent Pupillary Defect} Unequal response of the pupils because one eye or optic nerve sends less light information to the brain than the other. During a swinging-light test, both pupils may appear to dilate when the light moves to the affected eye; it can indicate retinal or optic-nerve disease.

\synonyms
RAPD; Marcus Gunn Pupil

\medterm{Relative Biological Effectiveness} A comparison of how strongly two types of radiation produce the same biological effect at the same absorbed dose. It depends on the radiation type, dose, tissue, and biological outcome.

\medterm{Relative Bradycardia} A physical finding in which the heart rate fails to increase appropriately in response to fever (sphygmothermic dissociation), where pulse rises by less than 10 beats per minute per degree Celsius of fever, characteristic of Salmonella typhoid fever, Legionnaires' disease, and yellow fever (Faget sign).
\synonyms
Sphygmothermic Dissociation; Inappropriate Bradycardia in Fever

\medterm{Relative Risk} A comparison of the chance of an outcome in an exposed group with its chance in an unexposed group. A value above 1 suggests higher risk, a value below 1 suggests lower risk, and 1 suggests no difference.

\medterm{Relaxation} A reduction in muscle tension, physiologic arousal, or psychological stress; techniques may include breathing, biofeedback, or progressive muscle relaxation.

\medterm{Relaxation Response} The body's calming changes after activities such as slow breathing, meditation, or progressive muscle relaxation. Heart rate, breathing, muscle tension, and stress-related alertness may decrease; these practices can support health but do not replace needed medical treatment.

\medterm{Relaxation Techniques For Stress} Relaxation techniques for stress include slow breathing, progressive muscle relaxation, meditation, guided imagery, and other methods that reduce physical arousal.

\medterm{Relaxation Therapy} Techniques such as breathing, muscle relaxation, or imagery used to reduce stress, arousal, or symptom burden.

\medterm{Relaxin} Relaxin is a reproductive hormone that helps remodel connective tissue and may contribute to cervical, pelvic, and vascular changes during pregnancy.

\medterm{Religion} A system of spiritual beliefs and practices that may influence health decisions, coping, diet, treatment preferences, and end-of-life care.

\medterm{Relocation Test} A shoulder examination often performed after an apprehension test. The examiner gently supports the front of the shoulder while the arm is in a vulnerable position. Relief of fear or pain may support shoulder instability, but the result must be interpreted with the rest of the examination.

\synonyms
Fowler Test; Jobe Relocation Test

\medterm{REM} Rapid eye movement, a sleep stage associated with vivid dreams and intense brain activity.
\synonyms
Rapid Eye Movement; REM Sleep

\medterm{REM Behavior Disorder} A sleep disorder where a person physically acts out dreams because normal muscle relaxation during REM sleep is lost.

\synonyms
RBD; REM Sleep Behavior Disorder

\medterm{Remedy} A treatment or measure intended to relieve a symptom or illness; safety and effectiveness depend on the specific remedy and condition.

\medterm{Remifentanil} A very short-acting opioid pain medicine given by infusion during anesthesia or intensive procedures. Its effects wear off quickly after the infusion stops, but it can slow breathing, lower blood pressure, and cause opioid-related complications.

\medterm{Remission} A period when the signs and symptoms of a disease decrease or disappear. Complete remission means no evidence of the disease can be found with available tests; partial remission means it has improved but remains detectable.

\medterm{Remittent Fever} Fever that rises and falls during the day but remains above normal temperature, rather than returning completely to normal.

\medterm{Remodeling} The body's ongoing process of removing old tissue and building new tissue to adapt to use, injury, growth, or disease. Bone remodeling responds to stress, hormones, nutrition, and medicines throughout life.

\medterm{Remote Surgery} Surgery performed when the surgeon is physically separated from the patient and controls robotic instruments through a telecommunications system. It may also include remote surgical assistance or guidance, with a qualified local team available to manage the patient and respond to emergencies.

\synonyms
Telesurgery
Telepresence Surgery

\medterm{Removal Technique} The method used to take away a foreign object, diseased tissue, stone, clot, or other material from the body. The appropriate method depends on its location, size, cause, and risk of bleeding, infection, or injury to nearby structures.

\medterm{REM Rebound} An increase in rapid-eye-movement sleep after sleep loss, REM deprivation, or stopping a medicine that suppresses REM sleep. It may cause vivid dreams, nightmares, or temporary changes in sleep patterns.

\medterm{REM Sleep} A sleep stage with rapid eye movements, active brain activity, and temporary relaxation of most skeletal muscles. Vivid dreams are common, and REM sleep supports memory, learning, emotional processing, and healthy development.

\medterm{REM Sleep Behavior Disorder} A sleep disorder in which the normal muscle paralysis of REM sleep is reduced or absent, allowing a person to move or vocalize while dreaming.

\synonyms
Rapid Eye Movement Sleep Behavior Disorder

\medterm{Renal} Renal means relating to the kidneys, which filter blood, regulate fluid and electrolytes, produce hormones, and excrete urine.

\medterm{Renal Abscess} A localized collection of pus within or around the kidney, usually caused by infection. Diagnosis and treatment may require imaging, antibiotics, and drainage.

\medterm{Renal Adenoma} A renal adenoma is a usually benign kidney tumor, often discovered incidentally; its size, imaging features, and growth determine whether follow-up is needed.

\medterm{Renal Agenesis} A condition present at birth in which one or both kidneys fail to form. One absent kidney may cause no symptoms if the other works well, but absence of both kidneys is usually fatal before or shortly after birth.

\medterm{Renal and Ureteral Injury} Mechanical, penetrating, blunt, or iatrogenic trauma affecting the kidneys or ureters. Clinical manifestations include gross or microscopic hematuria, flank pain, retroperitoneal hemorrhage, urinoma formation, or obstructive hydronephrosis requiring contrast-enhanced CT staging and potential surgical or endoscopic stenting.

\synonyms
Upper Urinary Tract Trauma; Renal Laceration; Ureteral Disruption

\medterm{Renal And Urological Diseases} Renal and urological diseases affect the kidneys, urinary tract, bladder, urethra, or male reproductive urinary organs.

\medterm{Renal Anemia} Anemia associated with chronic kidney disease. Diseased kidneys make less erythropoietin, a hormone that stimulates red-cell production, and iron deficiency or inflammation may add to the problem, causing tiredness and breathlessness.

\medterm{Renal Arteriography} Renal arteriography uses contrast X-rays to visualize kidney arteries and may identify narrowing, bleeding, aneurysm, or vascular anatomy during an intervention.

\medterm{Renal Artery} One of the two main arteries carrying blood from the abdominal aorta to the kidneys. Each divides into smaller branches inside the kidney so blood reaches the filtering tissue.

\medterm{Renal Artery Aneurysm} A focal enlargement of a renal artery, often asymptomatic but capable of causing hypertension,
thromboembolism, or rupture. Management depends on size, symptoms, anatomy, pregnancy, and
growth, and may involve surveillance, endovascular treatment, or surgery.

\medterm{Renal Artery Disease} Alternate name; see \textit{Renal artery stenosis}.

\medterm{Renal Artery Doppler} Doppler ultrasound that measures blood flow in the renal arteries. It may help assess narrowing or blockage that can contribute to renovascular hypertension.

\medterm{Renal Artery Embolism} Sudden blockage of an artery supplying the kidney, often by a clot from the heart in atrial fibrillation or another heart condition. It may cause abrupt side pain, blood in urine, nausea, or kidney injury.

\medterm{Renal Artery Stenosis} Narrowing of an artery supplying a kidney. It can reduce kidney blood flow and cause difficult-to-control high blood pressure or worsening kidney function. Causes include fatty buildup and fibromuscular dysplasia.

\synonyms
RAS; Renal Artery Narrowing; Renovascular Disease

\medterm{Renal Biopsy} Alternate name; see \textit{Kidney biopsy}. It is used selectively for unexplained acute kidney injury, nephrotic syndrome, glomerular hematuria, or transplant dysfunction; bleeding is the main complication.

\medterm{Renal Calculi} Alternate medical term; see \textit{Nephrolithiasis}.

\medterm{Renal Calculus} A kidney stone, a hard mass of mineral crystals formed in the urinary tract. It may pass without symptoms or cause severe wave-like pain, blood in the urine, blockage, vomiting, or urinary infection.

\medterm{Renal Cancer} Alternate name; see \textit{Kidney cancer}.

\medterm{Renal Carcinoma} Renal carcinoma is cancer arising from kidney tissue, most commonly renal cell carcinoma. Cancer arising from kidney tissue, most often from cells lining the kidney's small filtering tubes.

\medterm{Renal Cell Carcinoma} A cancer beginning in the kidney’s small filtering tubes, most often the clear-cell type. It may cause blood in urine, pain, weight loss, fever, or no symptoms until found on imaging.

\medterm{Renal Cell Carcinoma Pathology} Pathologic classification of renal cell carcinoma by growth pattern, cytology, immunophenotype, necrosis, sarcomatoid or rhabdoid change, and invasion. Clear cell, papillary, chromophobe, and other recognized subtypes have different features and behavior.

\medterm{Renal Changes In Pregnancy} Pregnancy increases renal blood flow and glomerular filtration. These changes lower
usual serum creatinine values and may increase urinary loss of glucose and amino acids.

\medterm{Renal Circulation} Blood flow through the kidneys, which supports filtration, oxygen delivery, hormone production, and regulation of fluid and electrolytes.

\medterm{Renal Clearance} A measure of how quickly the kidneys remove a substance from the blood. It is usually calculated from blood and urine levels and helps assess kidney function or adjust some medicine doses.

\medterm{Renal Colic} Severe, wave-like pain caused when a kidney stone or other blockage stretches the urine tube. Pain usually starts in the side and may move toward the groin, with nausea, vomiting, or blood in urine; fever with blockage is an emergency.

\medterm{Renal Corpuscle} The first part of a kidney filtering unit, made of a tiny ball of blood vessels inside a cup-shaped capsule. Water and small dissolved substances pass out of the blood here to begin urine formation.

\medterm{Renal Cortex} The outer layer of the kidney. It contains many of the structures that filter blood and process the first fluid that will become urine, while extensions of the cortex separate the inner kidney pyramids.

\medterm{Renal Cryptococcosis} Infection of the kidney by Cryptococcus species, usually as part of disseminated cryptococcosis in
an immunocompromised person.

\medterm{Renal Cyst} Alternate name; see \textit{Kidney Cysts}.

\medterm{Renal Dialysis} Alternate term for dialysis: kidney replacement therapy using hemodialysis or peritoneal dialysis to remove excess fluid and dissolved waste products when kidney function is inadequate.

\medterm{Renal Diet} An eating plan adjusted for kidney disease. Depending on kidney function and treatment, it may limit sodium, potassium, phosphorus, or fluid while providing enough calories and protein; needs differ between people and during dialysis.

\medterm{Renal Disease} Any condition affecting the kidneys, including infection, inflammation, stones, obstruction, inherited disease, or chronic loss of filtering function. Kidney disease can alter fluid balance, blood pressure, and body chemistry.

\medterm{Renal Diseases} Disorders affecting kidney tissue or function, including glomerular, tubular, vascular, infectious, obstructive, and neoplastic disease.

\medterm{Renal Dose Adjustment} Modification of a medication dose, dosing interval, or both to account for reduced
kidney clearance. The adjustment should use an appropriate estimate of renal function,
drug-specific pharmacokinetics, treatment indication, and therapeutic monitoring when
available.

\medterm{Renal Dysplasia} Abnormal development of kidney tissue before birth. It may affect one or both kidneys and can cause reduced kidney function, high blood pressure, urine reflux, or other urinary-tract problems.

\medterm{Renal Dysplasia And Cystic Disease} A group of developmental or inherited kidney disorders in which the kidneys have abnormal tissue, cysts, or both. They may affect one or both kidneys and can reduce kidney function or cause high blood pressure.

\medterm{Renal Failure} Severe loss of kidney function, causing waste, extra fluid, and abnormal body salts to build up. It may develop suddenly and sometimes improve, or progress gradually as chronic kidney disease; dialysis or transplant may be needed.

\medterm{Renal Function} The kidneys’ ability to filter blood, remove waste, regulate fluids and electrolytes, maintain acid-base balance, and produce or activate hormones.

\medterm{Renal Glomerulus} A tiny ball of blood vessels inside each kidney filtering unit. Water and small dissolved substances pass from the blood through its walls to begin urine formation, while blood cells and most proteins remain in the bloodstream.

\medterm{Renal Hamartoma} A renal hamartoma is a benign disorganized overgrowth of mature tissues in the kidney, sometimes used to describe an angiomyolipoma-like lesion.

\medterm{Renal Hypertension} High blood pressure caused by kidney disease or narrowing of the renal arteries. High blood pressure caused by kidney disease or narrowed renal arteries, often requiring treatment of both causes.

\medterm{Renal Hypertrophy} Renal hypertrophy is enlargement of kidney tissue, often compensating for loss or reduced function of the other kidney but sometimes caused by disease.

\medterm{Renal Hypoperfusion} Reduced blood flow to the kidneys, causing decreased glomerular filtration and potentially acute kidney injury. Causes include dehydration, blood loss, shock, heart failure, renal artery disease, and severe hypotension. Management requires treating the cause and restoring adequate circulation while avoiding further kidney injury.

\medterm{Renal Infarction} Death of part of the kidney caused by a blocked blood vessel. It may cause sudden flank or abdominal pain, nausea, blood in the urine, or reduced kidney function.

\medterm{Renal Insufficiency} Renal insufficiency is reduced kidney function that impairs waste removal, fluid and electrolyte balance, acid handling, or hormone production.

\medterm{Renal Leukemia Infiltration} Spread of leukemia cells into kidney tissue. It may enlarge both kidneys and reduce kidney function, sometimes causing sudden kidney injury, blood or protein in urine, or high blood pressure; treatment focuses on the underlying leukemia.

\medterm{Renal Lithiasis} Alternate name; see \textit{Kidney stones}.

\medterm{Renal Manifestations of Sickle Cell Disease} Kidney complications of sickling disorders, including excessive urination, impaired concentration, hematuria, proteinuria, papillary necrosis, acute kidney injury, and progressive chronic kidney disease. Blood pressure and urine abnormalities require ongoing monitoring.

\medterm{Renal Mass} An abnormal area seen in or on a kidney during imaging. It may be a harmless cyst, infection, scar, benign growth, or cancer; its appearance, size, growth, and sometimes biopsy determine the next step.

\medterm{Renal Medulla} The inner part of the kidney containing tubes and blood vessels that help concentrate urine. It receives the fluid processed by the outer kidney and directs the finished urine toward the kidney’s central collecting area.

\medterm{Renal Necrosis} Renal necrosis is death of kidney tissue caused by severe ischemia, toxins, infection, obstruction, or other injury and can cause acute kidney failure.

\medterm{Renal Oncocytoma} A renal oncocytoma is a usually benign kidney tumor arising from oncocytic cells; imaging may not always distinguish it reliably from renal cancer.

\medterm{Renal Osteodystrophy} Bone and mineral changes caused by advanced chronic kidney disease. Abnormal calcium, phosphate, vitamin D, and parathyroid hormone levels can weaken or deform bones, cause bone pain, and increase fracture risk; treatment corrects these imbalances.

\medterm{Renal Papillary Necrosis} Renal papillary necrosis is death of the kidney papillae from ischemia, infection, obstruction, diabetes, sickle-cell disease, or analgesic injury and may cause blood or obstruction.

\medterm{Renal Parenchyma} The working tissue of the kidney, containing the tiny filtering units and supporting structures that remove waste from blood, balance body fluids, and produce urine.

\medterm{Renal Pelvis} The funnel-shaped expansion of the ureter within the kidney, formed by the union of
major calyces. It collects urine from the renal papillae and directs it into the ureter.

\medterm{Renal Pelvis Or Ureter Cancer} Cancer of the renal pelvis or ureter usually arises from urothelial cells lining the urinary tract and may cause blood in the urine.

\medterm{Renal Perfusion Scintiscan} A nuclear medicine study using a radiotracer to evaluate renal blood flow and relative perfusion of the kidneys.

\medterm{Renal Pyramids} Cone-shaped structures in the renal medulla that contain tubules and collecting ducts and drain urine toward the renal papillae.

\medterm{Renal Reabsorption} The return of water and useful substances, such as salt, sugar, and bicarbonate, from kidney tubes into the blood. This prevents their unnecessary loss and helps control body fluids and chemistry.

\medterm{Renal Replacement Therapy} Treatment that takes over some functions of severely damaged kidneys, including removing waste and extra fluid and correcting dangerous chemical imbalances. Options include hemodialysis, peritoneal dialysis, continuous treatment in intensive care, or kidney transplantation.

\medterm{Renal Replacement Therapy Indications} Reasons to consider urgent dialysis include dangerously high potassium, severe acid buildup, certain poisonings, fluid overload causing breathing difficulty, or symptoms from waste buildup. The decision also considers kidney recovery, overall health, and the person’s goals.

\medterm{Renal Salt-Wasting Syndrome} Excessive loss of sodium and other salts through the kidneys, causing low blood sodium, dehydration, low blood pressure, or muscle symptoms. Causes include inherited tubular disorders, medicines, kidney disease, and hormonal problems.

\medterm{Renal Sarcoma} Renal sarcoma is a rare malignant tumor arising from connective or supporting tissues of the kidney.

\medterm{Renal Scan} A nuclear medicine test using a small radioactive tracer to assess kidney blood flow, drainage, or separate function. The tracer and scan method are chosen according to the clinical question.

\medterm{Renal Scintigraphy} Nuclear imaging that follows a radioactive tracer through the kidneys. It can compare kidney function, assess blood flow and urine drainage, detect obstruction, or monitor a transplanted kidney.

\medterm{Renal Stone} A renal stone is a solid crystal aggregate formed in the kidney that can pass through the urinary tract or obstruct urine flow.

\medterm{Renal Transitional Cell Carcinoma} A urothelial carcinoma arising in the renal pelvis or calyces rather than the kidney’s filtering tissue. It may cause blood in the urine, flank pain, or obstruction and is evaluated with imaging, urine testing, endoscopy, and tissue diagnosis.

\medterm{Renal Transplantation} Surgical placement of a functioning donor kidney into a person with advanced kidney failure. It can restore
kidney function but requires suitability assessment and long-term immunosuppression.

\synonyms
Kidney Transplant; Kidney Transplantation

\medterm{Renal Trauma} Kidney injury caused by a blow, fall, crash, sports injury, or penetrating wound. Assessment may include urine testing and CT imaging; treatment ranges from observation to procedures or surgery depending on severity.

\medterm{Renal Tubular Acidosis} A group of disorders in which kidney tubules do not remove enough acid or return enough bicarbonate to the blood. This makes the blood too acidic and may cause weakness, kidney stones, bone disease, or high potassium.

\medterm{Renal Tubular Acidosis Type 1} A metabolic acidosis caused by an impaired ability of distal intercalated renal tubular cells to secrete hydrogen ions into the collecting duct lumen. Characterized by normal anion gap metabolic acidosis, inappropriately high urinary pH (greater than 5.5 in the setting of systemic acidemia), hypokalemia, hypocitraturia, and recurrent nephrolithiasis or nephrocalcinosis.

\synonyms
Distal RTA; Type 1 RTA

\medterm{Renal Tubular Acidosis Type 2} A normal anion gap hyperchloremic metabolic acidosis caused by defective proximal tubular reabsorption of filtered bicarbonate, resulting in bicarbonate wasting into the urine. Urinary pH is typically elevated initially but drops below 5.5 when systemic acidosis becomes severe; frequently associated with generalized proximal tubular dysfunction (Fanconi syndrome) including glucosuria, aminoaciduria, phosphaturia, and osteomalacia.

\synonyms
Proximal RTA; Type 2 RTA

\medterm{Renal Tubular Acidosis Type 4} A normal anion gap metabolic acidosis characterized by persistent hyperkalemia caused by absolute aldosterone deficiency or renal tubular resistance to aldosterone (hypoaldosteronism). Impaired distal sodium reabsorption and potassium/hydrogen excretion reduce renal ammonium generation, yielding an acidic urinary pH (typically less than 5.5); common in diabetic nephropathy and chronic tubulointerstitial nephritis.

\synonyms
Hyperkalemic RTA; Hypoaldosteronemic RTA

\medterm{Renal Tubular Disorders} Diseases affecting the kidney tubules, which reabsorb water and useful substances and remove selected wastes into urine. They can disturb acid-base balance, electrolytes, urine concentration, or blood pressure.

\medterm{Renal Tubule} A tiny nephron tube that reabsorbs needed water and substances and removes others into urine.

\medterm{Renal Tubule Acidosis} A group of disorders in which the kidney tubules do not remove enough acid or conserve enough bicarbonate. The blood becomes too acidic, and complications may include low potassium, kidney stones, bone disease, weakness, or poor growth.

\medterm{Renal Ultrasound} Noninvasive imaging used to assess kidney size, structure, blood flow, and urinary obstruction. It is
particularly useful when obstruction, cystic disease, or structural abnormality is
suspected and is often used in the evaluation of acute or chronic kidney disease.

\medterm{Renal Vascular Disease} Alternate name; see \textit{Renal artery stenosis}.

\medterm{Renal Vein} One of two large veins that carry blood away from the kidneys to the body’s main abdominal vein. A clot in a renal vein can impair drainage and damage kidney function.

\medterm{Renal Vein Thrombosis} Renal vein thrombosis is a clot in a kidney vein that can cause flank pain, blood or protein in urine, kidney dysfunction, swelling, or pulmonary embolism.

\medterm{Renal Venogram} Contrast imaging of the renal veins used selectively to evaluate venous anatomy, thrombosis, obstruction, or vascular abnormalities.

\medterm{Renaut Body} A Renaut body is a specialized cluster of cells in the peripheral nerve endoneurium, associated with pressure sensitivity and sometimes nerve pathology.

\medterm{Renin} An enzyme from kidney cells that starts a hormone chain controlling blood pressure, salt balance, and fluid levels.

\medterm{Renin-Angiotensin-Aldosterone System} A hormone system that raises blood pressure and preserves body fluid when circulation or kidney blood flow falls. It narrows blood vessels and signals the adrenal glands to retain salt and water; overactivity contributes to hypertension and heart failure.

\medterm{Renin-Angiotensin-Aldosterone System Physiology} When blood pressure or kidney blood flow falls, the kidneys release renin. A chain of chemical signals then narrows blood vessels and prompts the adrenal glands to retain salt and water, raising blood pressure and restoring circulation.

\medterm{Renin-Angiotensin System} A hormone and enzyme system that regulates blood pressure, fluid balance, and vascular tone. Renin initiates formation of angiotensin II, which promotes vasoconstriction and aldosterone release; excessive activation contributes to hypertension and heart failure.

\medterm{Renin Blood Test} A blood test measuring renin concentration or activity, interpreted with aldosterone, posture, sodium status, medications, and blood pressure to evaluate hypertension or mineralocorticoid disorders.

\medterm{Rennin} An older spelling that may refer to renin or to rennet enzyme; in humans, renin regulates blood pressure and fluid balance.

\medterm{Renography} A kidney scan that follows a small amount of radioactive tracer as it enters and leaves the kidneys. It can compare kidney function and help identify reduced blood flow or blocked urine drainage.

\medterm{Renovascular Disease} Disease of the blood vessels supplying or draining the kidneys. It includes narrowing, blockage, aneurysm, or other vessel abnormalities and may reduce kidney blood flow or contribute to high blood pressure.

\medterm{Renovascular Hypertension} High blood pressure caused by reduced blood flow through one or both kidney arteries, usually from narrowing. It may be suspected when hypertension begins suddenly or is difficult to control; treatment addresses blood pressure and, in selected cases, the narrowed artery.

\medterm{Renshaw Cell} A nerve cell in the spinal cord that helps limit and regulate the activity of motor neurons.

\medterm{Reoviridae Infection} An infection caused by a reovirus; illness varies by virus and host and is often mild or asymptomatic in humans.

\medterm{Rep} An abbreviation most often meaning one repetition of an exercise or movement. In medical records it can also mean representative or another phrase, so the full wording should be used when the meaning is important.

\medterm{Repair Of Webbed Fingers Or Toes} Surgery to separate joined fingers or toes and reconstruct the skin between them. It may improve movement, growth, function, or appearance; wound problems, scarring, and recurrence of joining are possible.

\medterm{Repeat Surgery} Repeat surgery is a second or subsequent operation on the same or related area to treat persistent disease, recurrence, complications, or failure of prior treatment.

\medterm{Reperfusion} Restoration of blood flow to tissue after it has been reduced or interrupted. Reperfusion can save injured tissue but may also cause additional inflammation or oxidative injury.

\medterm{Reperfusion Injury} Damage that occurs when blood flow returns to tissue after a period of inadequate flow. The returning oxygen and inflammation can worsen injury, causing swelling, bleeding, or loss of function in organs such as the heart, brain, kidneys, or limbs.

\medterm{Reperfusion Therapy} Treatment that restores blood flow to tissue deprived of oxygen, such as clot-dissolving medicine, angioplasty, stenting, or surgery. Benefit depends on the tissue, duration of blockage, and risk of bleeding.

\medterm{Repetitive Motion Disorders} Conditions caused or worsened by repeated movement, force, vibration, or awkward posture. They may affect muscles, tendons, nerves, or joints and can cause pain, numbness, weakness, or reduced function.

\medterm{Repetitive Nerve Stimulation} A test that gives a nerve several small electrical impulses while recording the muscle’s response. A changing response may support a disorder in communication between nerves and muscles, such as myasthenia gravis or Lambert-Eaton syndrome. The result is interpreted with the examination and other tests.

\synonyms
RNS; Repetitive Stimulation Test

\medterm{Repetitive Strain Injury} Pain or dysfunction caused by repeated forceful movements, sustained positions, or overuse of muscles, tendons, nerves, or other tissues.

\medterm{Replantation Of An Avulsed Tooth} Putting a completely knocked-out permanent tooth back into its socket after an injury. The tooth should be handled by its crown, kept moist, and seen urgently; later splinting, root treatment, and follow-up may be needed to prevent loss.

\medterm{Replantation Of Digits} Microsurgical replacement of a completely amputated finger or toe, reconnecting bone, tendons, nerves, arteries, and veins when feasible.

\medterm{Replication Initiation} The time and process when copying of DNA or another genetic molecule begins. Specific proteins recognize a starting region, separate the strands, and recruit enzymes that build matching new genetic material.

\medterm{Replicon} A unit of genetic material that can replicate from one origin, such as a chromosome segment, plasmid, or viral genome.

\medterm{Repolarization} The return of the membrane potential toward the resting value after depolarization. In
neurons and skeletal/cardiac muscle, repolarization is primarily caused by the efflux of
potassium through voltage-gated potassium channels and inactivation of sodium channels.

\medterm{Reportable Diseases} Reportable diseases are illnesses that must be notified to public-health authorities under applicable laws to support surveillance and outbreak control.

\medterm{Reporting Standards} Agreed rules for recording and communicating clinical or imaging findings. Structured reports improve completeness and consistency and can make important results, follow-up needs, and comparisons easier for clinicians to understand.

\medterm{Repressed Memory Therapy} A controversial approach aimed at recovering allegedly repressed memories; suggestive techniques can create false memories and require caution.

\medterm{Repression} The reduction or stopping of gene activity, cell activity, behavior, or a body process. In genetics, regulatory proteins can silence genes; in psychology, the term may describe unconsciously keeping distressing thoughts outside awareness.

\medterm{Reproduction} The biological process by which living organisms produce new offspring, either sexually or without combining reproductive cells. In humans, it involves reproductive organs, hormones, eggs, sperm, pregnancy, and childbirth.

\medterm{Reproductive Cells} Cells used for reproduction, mainly sperm in males and eggs in females. Each normally carries half the usual number of chromosomes.

\medterm{Reproductive History} Reproductive history records pregnancies, births, miscarriages, abortions, menstrual function, fertility, contraception, and related conditions relevant to health assessment.

\medterm{Reproductive Hormones} Hormones that regulate sexual development, reproductive organs, ovulation, menstruation, sperm production, pregnancy, and lactation. Examples include estrogen, progesterone, testosterone, luteinizing hormone, and follicle-stimulating hormone.

\medterm{Reproductive Organs} The organs involved in reproduction. Female: ovaries (produce eggs and hormones),
fallopian tubes (transport eggs, site of fertilization), uterus (nurtures the developing
fetus), cervix (lower part of uterus, dilates during labor), vagina (passage for sperm
and baby), and vulva (external genitalia).

\medterm{Reproductive System} The organs and tissues that produce gametes and reproductive hormones and support sexual development and reproduction. Depending on the anatomy present, it may enable sperm or egg production, fertilization, pregnancy, childbirth, and lactation.

\medterm{Reproductive Technique} A reproductive technique is a medical or laboratory method used to assist conception, preserve fertility, diagnose reproductive disease, or manage pregnancy.

\medterm{Reproductive Toxicity} Harmful effects of a substance on fertility, pregnancy, the developing baby, birth outcomes, or later growth. Risks depend on the substance, dose, timing, route, and individual health and require careful exposure assessment.

\medterm{Reproductive Years} The part of life when a person may be able to become pregnant or produce sperm. Fertility varies widely and is affected by age, health, reproductive organs, hormones, and medicines; it should not be assumed from age alone.

\medterm{Reprogram} To reprogram a cell means to alter its gene-expression state so it acquires a different developmental potential or function.

\medterm{Resectable} Describing a tumor or lesion that can be removed completely by surgery with acceptable risk and an appropriate margin, based on its extent and relationship to vital structures.

\medterm{Resection} Surgical removal of part or all of an organ, tissue, or abnormal growth. The removed material may be examined to confirm diagnosis, determine treatment, or reduce disease while preserving useful function.

\synonyms
Surgical Excision; Surgical Removal

\medterm{Resectoscope} An endoscopic surgical instrument with a viewing system and working loop used to remove or cauterize tissue, especially during prostate or bladder procedures.

\medterm{Reserve Volume} The additional amount of air that can be inhaled or exhaled beyond a normal tidal breath.

\medterm{Reservoir} A human, animal, or environmental habitat in which an infectious agent normally lives,
multiplies, and persists. A reservoir may serve as the source from which the agent is
transmitted to a susceptible host.

\medterm{Reservoir Of Infection} A reservoir of infection is a person, animal, environment, or material in which an infectious organism normally lives and from which transmission can occur.

\medterm{Resident} A resident is a physician or other trainee receiving supervised specialty education after initial professional qualification; the term can also mean a person living in a place.

\medterm{Residential Mobility} Residential mobility is movement between homes or communities and can affect continuity of care, social support, exposure, and health access.

\medterm{Residual} Residual means remaining after a process, treatment, measurement, or event, such as residual urine, disease, risk, or volume.

\medterm{Residual Limb} The portion of a limb remaining after amputation. Care includes edema control, skin
inspection, shaping, strengthening, desensitization, contracture prevention, and
prosthetic preparation.

\medterm{Residual Limb Pain} Pain in the remaining portion of a limb after amputation.

\synonyms
Stump Pain

\medterm{Residual Neoplasm} A residual neoplasm is tumor tissue that remains after treatment intended to remove or destroy the original tumor.

\medterm{Residual Risk} Residual risk is the chance of harm that remains after preventive measures, treatment, screening, or controls have been applied.

\medterm{Residual Tumor} A residual tumor is cancer or other tumor tissue that remains after surgery, radiation, chemotherapy, or another treatment.

\medterm{Residual Volume} The amount of air remaining in the lungs after breathing out as forcefully as possible. It keeps small airways open and cannot be measured accurately by ordinary spirometry alone.

\medterm{Resistance} The ability of a microorganism to survive a medicine that previously inhibited or killed it. Resistance can make infections harder to treat and may spread between organisms through genetic material.
\synonyms
Antimicrobial Resistance; Drug Resistance

\medterm{Resistance Training} Exercise in which muscles work against weights, bands, body weight, or other resistance. Regular training can improve strength, endurance, balance, bone health, function, and independence when progressed safely.

\medterm{Resistant Hypertension} High blood pressure remaining above goal despite appropriate use of three or more medicines, including a water pill.

\synonyms
Refractory Hypertension; Treatment-Resistant Hypertension

\medterm{Resistant Starch} Starch that resists digestion in the small intestine and is fermented by gut bacteria in the colon.

\medterm{Resistin} A signaling protein associated with adipose tissue and inflammation; its human clinical significance remains under study.

\medterm{Resolution} The ability of an imaging system to distinguish separate structures, or the improvement or disappearance of a disease or symptom.

\medterm{Resonance} Amplification or characteristic vibration produced when a system is driven at a suitable frequency; used in sound and MRI contexts.

\medterm{Resorcinol} A phenolic compound used in selected topical dermatologic preparations; concentrated exposure can irritate or poison.

\medterm{Resorption} The body's breakdown and removal of tissue or material, especially old bone. The released components may be absorbed, reused, or removed from the body.

\medterm{Resorption Atelectasis} Atelectasis caused by complete obstruction of an airway, leading to absorption of
trapped air by the pulmonary capillaries. Common causes include mucus plugging, foreign
body, tumor, and post-surgical shallow breathing. The affected area appears as increased
density on CXR with volume loss.

\medterm{Respiration} The processes of breathing, gas exchange, and cellular energy production involving oxygen and carbon dioxide.

\medterm{Respiratory} Respiratory means relating to breathing, the airways, lungs, gas exchange, or the organs and muscles that move air.

\medterm{Respiratory Acidosis} A condition in which breathing does not remove enough carbon dioxide, making the blood too acidic. It can result from lung disease, blocked airways, weak breathing muscles, brain disorders, or medicines that slow breathing.

\medterm{Respiratory Alkalosis} A condition in which breathing removes too much carbon dioxide, making the blood less acidic. It can occur with anxiety, pain, fever, low oxygen, pregnancy, severe infection, or other conditions that cause rapid breathing.

\medterm{Respiratory Arrest} Stopping of effective breathing. It is an immediate emergency that may require airway support, rescue breathing or ventilation, and treatment of the cause.

\medterm{Respiratory Burst} A rapid release of chemically reactive oxygen molecules by certain immune cells after they engulf germs. This helps kill the germs inside the cell, although excessive activity can also damage nearby tissue.

\medterm{Respiratory Changes In Pregnancy} Pregnancy changes breathing as the growing uterus raises the diaphragm and hormones increase the drive to breathe. Breaths may become deeper, overall ventilation increases, and the amount of air left after a normal breath out decreases.

\medterm{Respiratory Compliance} The physical measure of the distensibility and elastic stretchability of the lungs and chest wall, expressed mathematically as the change in lung volume per unit change in distending pressure ($\Delta V / \Delta P$); reduced in acute respiratory distress syndrome, pulmonary edema, and pulmonary fibrosis.

\synonyms
Lung Compliance

\medterm{Respiratory Depression} Breathing that is abnormally slow or shallow, causing inadequate ventilation and possible
carbon-dioxide retention or low oxygen. Medicines, brain disorders, and metabolic illness
can cause it.

\medterm{Respiratory Disease} Any disease affecting the airways, lungs, breathing muscles, or gas exchange. Examples include asthma, chronic obstructive pulmonary disease, pneumonia, tuberculosis, and lung cancer.

\medterm{Respiratory Disorder} A disease affecting the airways, lungs, pleura, respiratory muscles, or control of breathing and causing impaired ventilation, oxygenation, or gas exchange.

\medterm{Respiratory Distress} Difficult or labored breathing caused by problems such as asthma, infection, a blood clot, a collapsed lung, an allergic reaction, or heart failure. Fast breathing, struggling to speak, blue lips, or severe exhaustion may occur.

\medterm{Respiratory Distress Syndrome} A serious breathing problem in premature newborns whose lungs do not make enough surfactant, a substance that keeps the tiny air sacs open. Babies may breathe rapidly, grunt, or need oxygen and breathing support. Antenatal steroids, noninvasive breathing support, and surfactant treatment improve survival.

\medterm{Respiratory Drive} The brain’s signal to the breathing muscles that controls how often and how deeply a person breathes. It responds mainly to carbon-dioxide and oxygen levels and can be reduced by brain disease, severe illness, or sedating medicines.

\medterm{Respiratory Failure} A life-threatening inability of the lungs to provide enough oxygen or remove enough carbon dioxide. It may cause severe breathlessness, confusion, blue or gray lips, extreme sleepiness, or collapse. Causes include lung disease, infection, heart problems, weak breathing muscles, brain injury, and drug overdose.

\medterm{Respiratory Failure Types 1 and 2} The clinical classification of acute respiratory failure: Type 1 (hypoxemic, $P_aO_2 < 60\text{ mmHg}$ with normal or low $P_aCO_2$) resulting from ventilation-perfusion mismatch or shunt (e.g., pneumonia, pulmonary edema); and Type 2 (hypercapnic, $P_aCO_2 > 45-50\text{ mmHg}$) resulting from alveolar hypoventilation (e.g., COPD, neuromuscular weakness, sedative overdose).

\medterm{Respiratory Gating} A radiation-treatment method that tracks breathing and delivers the beam only during a selected part of the breathing cycle. It helps account for movement of tumors in the chest or upper abdomen and may reduce radiation to nearby healthy tissue.

\medterm{Respiratory Health} The proper functioning of the lungs and airways, enabling efficient gas exchange and
oxygen delivery to the body. Maintaining respiratory health involves avoiding tobacco
smoke, minimizing exposure to air pollutants and occupational hazards, and receiving
vaccinations against respiratory infections like influenza and pneumococcal disease.

\medterm{Respiratory Illness} A short-term or long-term illness affecting breathing or the respiratory tract. Symptoms may include cough, fever, congestion, wheezing, shortness of breath, or chest discomfort, depending on the cause.

\medterm{Respiratory Infection} An infection of the upper or lower respiratory tract caused by viruses, bacteria, fungi, or other organisms.

\medterm{Respiratory Infections} Illnesses caused by pathogens including viruses, bacteria, and fungi that affect the airways and lungs.

\medterm{Respiratory Insufficiency} Inadequate breathing or gas exchange that causes low oxygen, high carbon dioxide, or both. It can result from lung or airway disease, weak breathing muscles, or problems in the brain’s breathing control and may progress to respiratory failure.

\medterm{Respiratory Muscle} A muscle used for breathing, especially the diaphragm and muscles between the ribs. These muscles expand and contract the chest to move air into and out of the lungs.

\medterm{Respiratory Neoplasm} A respiratory neoplasm is an abnormal growth in the airways or lungs and may be benign or malignant.

\medterm{Respiratory Paralysis} Respiratory paralysis is failure of the muscles needed for breathing because of nerve, muscle, brainstem, toxin, or spinal-cord dysfunction and is an emergency.

\medterm{Respiratory Physiology} The study of how breathing works, including movement of air, exchange of oxygen and carbon dioxide in the lungs, transport of gases in blood, and control of breathing.

\medterm{Respiratory Protection} Measures and equipment used to reduce inhalation of harmful dusts, fumes, gases, aerosols, or infectious particles.

\medterm{Respiratory Quotient} The ratio of carbon dioxide produced to oxygen used by the body’s cells. The value changes according to whether the body is mainly using carbohydrate, fat, or protein and helps estimate metabolism and energy needs.

\medterm{Respiratory Rate} The number of breaths taken per minute. The expected rate varies with age, activity, temperature, illness, and other factors.

\medterm{Respiratory Rehabilitation} A program of exercise, education, breathing techniques, and self-management that improves
function and quality of life in people with chronic respiratory disease.

\medterm{Respiratory Syncytial Virus} A respiratory virus that commonly causes bronchiolitis and pneumonia in infants
and young children. Infection is usually mild in adults but may be severe in infants, older adults, and people with
chronic heart or lung disease.

\synonyms
RSV; Human Orthopneumovirus

\medterm{Respiratory Syncytial Virus Bronchiolitis} An acute lower respiratory tract infection occurring predominantly in infants under two years of age caused by respiratory syncytial virus (RSV), characterized by bronchiolar mucosal inflammation, edema, and epithelial sloughing leading to wheezing, tachypnea, cough, and intercostal retractions.

\medterm{Respiratory Syncytial Virus Infection} Infection caused by respiratory syncytial virus. It commonly causes a cold-like illness but may involve the lower airways, producing bronchiolitis or pneumonia. Severity varies with age, immune status, and underlying heart or lung disease.

\medterm{Respiratory System} The nose, mouth, pharynx, larynx, trachea, bronchi, bronchioles, lungs, pleura, and breathing muscles, including the diaphragm and chest-wall muscles. It moves air into and out of the lungs and exchanges oxygen and carbon dioxide between the alveoli and blood.

\medterm{Respiratory Therapy} Respiratory therapy includes treatments and breathing support that improve ventilation, oxygenation, airway clearance, or lung function.

\medterm{Respiratory Tract Infections} Infections affecting any part of the respiratory tract, classified by anatomical
location into upper respiratory tract infections (URTIs) and lower respiratory tract
infections (LRTIs).

\medterm{Respiratory Transport} Respiratory transport is movement of oxygen and carbon dioxide between lungs, blood, and tissues, largely through hemoglobin and dissolved gas.

\medterm{Respiratory Viruses} Viruses that cause respiratory tract infections. Major respiratory viruses include
influenza, respiratory syncytial virus (RSV), parainfluenza virus, adenovirus, and human
metapneumovirus. SARS-CoV-2 causes COVID-19. Symptoms range from mild upper respiratory
illness to severe pneumonia.

\medterm{Respiratory Viruses And Face Masks} Face masks can reduce spread of respiratory viruses by blocking droplets and some aerosols, especially in crowded or healthcare settings; fit and local guidance matter.

\medterm{Respirologist} A doctor specializing in diseases of the lungs and breathing system. They assess breathing problems, interpret tests, manage chronic and acute illness, and may evaluate people for lung transplantation.

\medterm{Respirovirus} A genus of paramyxoviruses that includes human parainfluenza viruses, which can cause croup and other respiratory infections.

\medterm{Respite Care} Temporary care for a person who needs support so that a usual family or unpaid caregiver can rest, attend appointments, or manage other responsibilities. It may be provided at home, in a facility, or through short-term services.

\medterm{Response Assessment} Evaluation of how a disease responds to treatment. Clinicians may use symptoms, examination, blood tests, scans, tumor measurements, or activity levels to decide whether treatment is helping, failing, or needs adjustment.

\medterm{Response Evaluation Criteria} Standardized rules for judging whether a disease, especially a tumor, has responded to treatment, remained stable, or progressed. RECIST mainly uses tumor measurements on imaging; other systems may use metabolic or blood-based findings.

\medterm{Responsible Drinking} Responsible drinking means limiting alcohol according to health guidance, avoiding driving or hazardous activities after drinking, and avoiding alcohol when it is unsafe or contraindicated.

\medterm{Rest} Rest is reduced physical or mental activity used to recover, though excessive inactivity can impair strength and function in some conditions.

\medterm{Restenosis} Re-narrowing of a blood vessel after it has been opened by angioplasty or another procedure.

\medterm{Resting Energy Expenditure} The energy the body uses while resting to keep essential functions going, such as breathing, circulation, temperature control, and cell activity. It is affected by body size, age, muscle, illness, hormones, and recent food or activity.

\synonyms
REE; Resting Metabolic Rate; Basal Metabolic Rate

\medterm{Resting Membrane Potential} The small electrical difference across a nerve or muscle cell when it is not sending a signal. Unequal amounts of charged particles inside and outside the cell create it and allow the cell to respond quickly to stimulation.

\medterm{Resting Phase} The resting phase is a period of relative inactivity, such as G0 in the cell cycle or a physiologic interval between active events.

\medterm{Resting Tremor} An involuntary, rhythmic oscillation (typically four to six hertz) occurring in a body part that is fully supported against gravity and relaxed, which diminishes or temporarily ceases with voluntary movement. It typically presents as a pill-rolling hand tremor in Parkinson's disease and other parkinsonian syndromes.

\medterm{Restless Anal Syndrome} An uncommon term for uncomfortable sensations around the anus with an urge to move, touch, or relieve the area. The cause and best treatment are uncertain, so other anal and bowel conditions should be assessed.

\medterm{Restless Legs} Alternate name; see \textit{Restless Legs Syndrome}.

\medterm{Restless Legs Syndrome} A neurological and sleep-related condition causing an irresistible urge to move the legs, often with crawling, pulling, aching, or tingling sensations. Symptoms begin or worsen during rest, improve temporarily with movement, and are worse in the evening or at night; they may disrupt sleep and daytime function. Iron deficiency, kidney disease, pregnancy, some medicines, and genetic factors may contribute.

\synonyms
RLS; Willis-Ekbom Disease

\medterm{Restless Leg Syndrome} Alternate name; see \textit{Restless Legs Syndrome}.

\medterm{Restlessness} Restlessness is an uncomfortable need to move or inability to remain still, caused by anxiety, akathisia, restless-legs syndrome, pain, stimulants, withdrawal, or medical illness.

\medterm{Restorative Dentistry} Specialty of restoring damaged teeth. Procedures: fillings, inlays, onlays, crowns,
bridges, implants. Goals: function, aesthetics, durability.

\medterm{Restraint} A device, medicine, or other measure used to limit a person’s movement. Restraints can cause injury, distress, falls, and loss of independence, so trained staff should use the least restrictive option only when necessary and monitor closely.

\medterm{Restraint-Related Harm} Physical or psychological injury associated with using a device, medication, or
environmental restriction to limit a person's movement. Older adults are at risk of
falls, pressure injury, delirium, aspiration, and functional decline; the least
restrictive alternatives should be prioritized.

\medterm{Restriction Enzyme} An enzyme made by bacteria that cuts DNA at particular short sequences. Scientists use these enzymes to analyze, join, or modify DNA in laboratory research and medical testing.

\medterm{Restrictive Cardiomyopathy} A disease in which stiff heart muscle cannot relax and fill normally. It may cause breathlessness, tiredness, swelling, or abnormal rhythms and can result from conditions such as amyloidosis, sarcoidosis, iron overload, or scarring.

\medterm{Restrictive Lung Disease} A pattern of reduced total lung capacity caused by stiff lung tissue, chest-wall limitation, respiratory-muscle weakness, obesity, or other factors. It may cause exertional breathlessness and low oxygen; lung volumes, imaging, muscle strength, and the underlying cause guide management.

\medterm{Resuscitation} Emergency measures used to restore or support breathing and circulation, including opening the airway, rescue breathing, chest compressions, defibrillation, medicines, and treatment of the cause.
\synonyms
Cardiopulmonary Resuscitation; CPR; Clinical Revival

\medterm{Resuscitative Thoracotomy} Emergency chest surgery used for selected trauma patients who are in cardiac arrest or close to it, especially after a penetrating chest injury. It may relieve blood around the heart, control bleeding, or allow direct heart compressions. It is a last-resort procedure with a high risk of death and complications.

\synonyms
Emergency Department Thoracotomy; EDT; Clamshell Thoracotomy; Trauma Thoracotomy

\medterm{Retained Placenta} Failure of the placenta to come out after childbirth, usually when it remains in the uterus for about 30 minutes or longer or bleeding is excessive. It can cause severe bleeding or infection.

\medterm{Retained Products Of Conception} Pregnancy tissue that remains in the uterus after childbirth, miscarriage, or abortion. It may cause continuing or heavy bleeding, cramping, fever, or infection and may be treated with medicine or a procedure to remove it.

\medterm{Retainer} A dental appliance worn after orthodontic treatment to help keep teeth in their corrected positions.

\medterm{Retching} Repeated involuntary effort to vomit without necessarily producing stomach contents.

\medterm{Retention Of Urine} Inability to empty the bladder completely or at all despite making urine. It may cause painful fullness, leaking, or kidney damage and can result from blockage, weak bladder muscles, medicines, or nerve disease.

\medterm{Retention Time} The time a substance takes to pass through a laboratory separation system from injection until it is detected. Comparing it with a tested reference can help identify a substance, but timing alone is not definitive proof.

\medterm{Rete Testis} A network of small spaces inside the testis that collects sperm from the seminiferous tubules and passes them to the small ducts leading to the epididymis.

\medterm{Reticular} Net-like or arranged in a network pattern, describing the appearance or distribution of a skin lesion, rash, or vascular pattern.

\medterm{Reticular Activating System} A network of brain structures that helps maintain wakefulness, alertness, attention, and transitions between sleep and waking. Damage or severe disturbance can cause unusual sleepiness, reduced awareness, or coma.

\medterm{Reticular Dermis} The deeper layer of the dermis containing dense connective tissue, collagen, elastin, blood vessels, nerves, glands, and hair follicles.

\medterm{Reticular Erythematous Mucinosis} A rare primary cutaneous mucinosis presenting with erythematous, reticulated macules and plaques with dermal mucin accumulation on the mid-chest and upper back in middle-aged individuals, exacerbated by ultraviolet light exposure.
\synonyms
REM Syndrome; Plaque-like Cutaneous Mucinosis

\medterm{Reticular Formation} A network of nerve cells in the brainstem that helps control alertness, sleep, attention, muscle tone, posture, breathing, heart function, and responses to pain.

\medterm{Reticulin} A network-forming type of collagen that supports bone marrow, liver, lymph nodes, and other organs.

\medterm{Reticulocyte} A young red blood cell recently released from bone marrow. The number of reticulocytes shows how actively the marrow is replacing red cells and helps distinguish reduced production from blood loss or red-cell breakdown.

\medterm{Reticulocyte Count} The number or percentage of immature red blood cells in blood. It reflects recent bone-
marrow red-cell production and helps assess anemia, especially when interpreted with the hemoglobin and clinical context.

\medterm{Reticulocyte Production Index} A calculated hematological parameter that corrects the raw reticulocyte percentage for both the patient's degree of anemia (hematocrit) and the premature release of \"shift\" reticulocytes from the bone marrow into circulating blood. Calculated as raw reticulocyte percent multiplied by (patient hematocrit / normal hematocrit) divided by maturation factor in days; an RPI of 2.0 or greater indicates an adequate marrow erythropoietic response to hemolysis or acute hemorrhage, whereas an RPI below 2.0 signifies an inadequate hypoproliferative marrow response.

\synonyms
RPI; Corrected Reticulocyte Count

\medterm{Reticulocytes} Immature red blood cells recently released from bone marrow; their proportion helps assess red-cell production.

\medterm{Reticulocytopenia} An abnormally low reticulocyte count indicating reduced effective red-cell production, which may occur with marrow failure, nutrient deficiency, inflammation, or kidney disease.

\medterm{Reticulocytosis} Reticulocytosis is an increased number of immature red blood cells in blood, usually indicating stimulated marrow production after blood loss or hemolysis.

\medterm{Reticulo-Endothelial System} An older name for the body’s network of macrophages and related immune cells in the blood, liver, spleen, lymph nodes, and tissues. These cells remove germs, particles, and worn-out blood cells.

\medterm{Reticuloendothelial System} A diffuse anatomical network of specialized phagocytic cells (primarily monocytes and tissue macrophages) situated in reticular connective tissue frameworks of the spleen, liver (Kupffer cells), lymph nodes, and bone marrow, responsible for filtering particulate foreign matter, microbial pathogens, and senescent erythrocytes.

\synonyms
Mononuclear Phagocyte System; MPS

\medterm{Reticulophagy} A cell-cleaning process that removes worn or damaged parts of the endoplasmic reticulum, the cell's internal membrane network. The material is broken down and recycled, helping cells control quality and respond to stress.

\medterm{Reticulum} A network-like structure; in anatomy or cell biology it may refer to a tissue network, endoplasmic reticulum, or ruminant stomach compartment.

\medterm{Retina} The light-sensitive neural tissue lining the back of the eye that converts light into nerve signals.

\medterm{Retina Disorder} A disease affecting the retina, including vascular, degenerative, inflammatory, infectious, traumatic, inherited, or neoplastic conditions that can impair vision.

\medterm{Retina Hemorrhage} Bleeding within or beneath the retina caused by vascular disease, diabetes, hypertension, trauma, retinal vein occlusion, inflammation, or other disorders.

\medterm{Retinal Artery Occlusion} Blockage of a retinal artery causing sudden, usually painless loss of vision. It may indicate a wider blood-vessel or clotting disorder.

\medterm{Retinal Cone} A light-sensing cell in the retina that provides color vision and sharp detail, especially in daylight. Damage to cone cells can reduce color perception and central vision.

\medterm{Retinal Damage} Injury or disease affecting the retina, potentially causing blurred vision, field loss, flashes, floaters, or permanent visual impairment.

\medterm{Retinal Degeneration} Progressive damage or loss of light-sensitive cells in the retina at the back of the eye. It may cause blurred or reduced central vision, night blindness, or loss of side vision and can result from inherited disease, aging, or other injury.

\medterm{Retinaldehyde} A form of vitamin A used in the eye’s light-sensing cycle. Light changes it inside photoreceptor cells, helping start the nerve signals that allow vision.

\medterm{Retinal Detachment} A separation of the retina, the light-sensitive tissue at the back of the eye, from the layer beneath it. It may follow a retinal tear, scar tissue, or fluid buildup. Sudden floaters, flashes, a curtain-like shadow, or reduced vision may occur, and permanent sight loss is possible.

\medterm{Retinal Detachment Repair} Surgical or pneumatic treatment to close retinal breaks and reattach the retina. Techniques include vitrectomy, scleral buckling, and pneumatic retinopexy, selected according to the type and location of detachment.

\medterm{Retinal Diseases} Disorders that damage the retina, the light-sensitive tissue at the back of the eye. They may cause flashes, floaters, distorted vision, blind spots, night-vision problems, or gradual or sudden sight loss.

\medterm{Retinal Disorders} Diseases affecting the light-sensing layer at the back of the eye. Causes include blood-vessel problems, inherited conditions, aging, inflammation, infection, injury, and cancer, potentially causing blurred vision or blindness.

\medterm{Retinal Dysplasia} Abnormal development and arrangement of the retina, the light-sensing layer at the back of the eye. It may impair vision and can occur with conditions present from birth, genetic disorders, or problems affecting development before birth.

\medterm{Retinal Dystrophy} Gradual loss or damage of retinal cells, often caused by an inherited condition. It may begin with night blindness, difficulty seeing in dim light, or loss of side vision and can progress to serious sight impairment.

\medterm{Retinal Fold} A retinal fold is a crease or ridge in the retina caused by traction, developmental abnormality, detachment, swelling, or another ocular process.

\medterm{Retinal Fundus} The inside back surface of the eye viewed during fundoscopy, including the retina, optic nerve, macula, and retinal blood vessels.

\medterm{Retinal Ganglion} A nerve cell in the retina that collects signals from other retinal cells and sends them through the optic nerve to the brain. It is essential for carrying visual information out of the eye.

\medterm{Retinal Hemorrhages} Bleeding in the light-sensitive tissue at the back of the eye. It may result from diabetes, high blood pressure, injury, bleeding disorders, or increased pressure in the skull; in infants and children, the pattern must be assessed carefully for its cause.

\medterm{Retinal Implant} An implanted visual prosthesis that detects or processes light and electrically stimulates remaining retinal or visual-pathway cells in selected severe retinal disorders. It provides limited visual information and is distinct from an intraocular lens.

\medterm{Retinal Ischemia} Too little blood and oxygen reaching the retina, often because a retinal artery is blocked or blood-vessel disease is severe. It may cause sudden, painless sight loss and is an emergency because prolonged injury can become permanent.

\medterm{Retinal Laser Photocoagulation} An ophthalmic procedure utilizing targeted laser energy to induce localized thermal denaturation of retinal tissue. Widely used to seal microvascular leakages in diabetic macular edema, ablate ischemic peripheral retina in proliferative diabetic retinopathy, or secure retinal tears to prevent rhegmatogenous retinal detachment.

\synonyms
Laser Photocoagulation; Retinal Laser Therapy; Panretinal Photocoagulation

\medterm{Retinal Layers} The organized layers of nerve cells, supporting cells, and blood vessels that make up the retina. Together they detect light, begin visual processing, and send signals through the optic nerve to the brain.

\medterm{Retinal Macula} The macula is the central region of the retina responsible for detailed straight-ahead vision, reading, and color discrimination.

\medterm{Retinal Neoplasm} A retinal neoplasm is an abnormal growth arising in or near the retina, such as retinoblastoma or a metastatic tumor.

\medterm{Retinal Perforation} A full-thickness hole or tear in the retina. Fluid can pass through the opening and separate the retina from the tissue beneath it, causing flashes, floaters, or loss of vision.

\medterm{Retinal Pigment} Retinal pigment refers to melanin-containing pigment in the retinal pigment epithelium, which supports photoreceptors and absorbs excess light.

\medterm{Retinal Pigment Epithelium} A layer of cells beneath the retina that supports the light-sensing cells, absorbs extra light, and helps remove worn-out cell material.

\medterm{Retinal Prosthesis} An implanted device designed to provide limited visual signals when the retina is severely damaged. A camera and electronic system stimulate remaining visual pathways; it may provide light or simple pattern perception, not normal sight.

\medterm{Retinal Tear} A break through the retina, often caused when the jelly inside the eye pulls on it. Fluid can pass through the tear and detach the retina; new flashes, floaters, or a curtain-like shadow may occur. Treatment may include laser or freezing treatment.

\medterm{Retinal Vasculitis} Inflammation of blood vessels in the retina. Leaking or blocked vessels can cause blurred vision, bleeding, swelling, or permanent sight loss; it may be linked to infection, autoimmune disease, or inflammation elsewhere in the body.

\medterm{Retinal Vein Occlusion} A blockage in a vein that drains the retina. It can cause sudden or gradual blurred vision from bleeding and swelling, especially in the macula, and may lead to abnormal blood-vessel growth or permanent sight loss.

\medterm{Retinitis} Inflammation or infection of the retina, the light-sensitive tissue at the back of the eye. It may cause blurred vision, floaters, flashes, or loss of sight and needs prompt eye assessment.

\medterm{Retinitis Pigmentosa} A group of inherited retinal diseases in which the retina’s light-sensing cells gradually stop working. Inheritance varies by gene and may be autosomal dominant, autosomal recessive, X-linked, or mitochondrial. Early symptoms often include poor night vision and loss of side vision, followed by narrowing of the visual field and possible severe sight loss.

\medterm{Retinoblastoma} The most common eye cancer in infants and young children, starting in the retina at the back of the eye. It is caused by changes in the \textit{RB1} gene. It can be inherited (often affecting both eyes) or occur by chance (usually affecting only one eye). The most common early signs are leukocoria (an abnormal white glow in the pupil, often noticed in flash photographs instead of the normal red reflex) and strabismus (crossed or misaligned eyes). Diagnosis is made through a specialized dilated eye examination and imaging scans (ultrasound or MRI) to check tumor size and optic nerve involvement. Treatment focuses on saving the child's life and preserving vision using chemotherapy (given into the bloodstream or directly into the eye's blood vessels), laser therapy, freezing (cryotherapy), or surgical removal of the eye (enucleation) in advanced cases.
\synonyms{Childhood Eye Cancer; Retinal Glioma; Malignant Retinoblastoma}

\medterm{Retinoblastoma Protein} A protein made from the RB1 gene that helps prevent cells from dividing when they should not. Loss of its protective function can contribute to retinoblastoma and other cancers.

\medterm{Retinoic Acid} An active form of vitamin A that helps control how cells develop and which genes they use. Related medicines are used for selected skin conditions and some cancers but can cause irritation or serious pregnancy risks.

\medterm{Retinoid} A vitamin-A-related substance used in selected skin treatments and other medical care. Retinoids can improve acne or abnormal cell growth but may irritate skin, interact with medicines, and harm an unborn baby.

\medterm{Retinoids} Vitamin-A-related compounds used in selected dermatologic, ophthalmic, and oncologic treatments.

\medterm{Retinol} A form of vitamin A needed for vision, healthy skin and body linings, growth, and immune function. It comes from animal foods and some supplements; too much can be harmful, especially during pregnancy.

\medterm{Retinopathy} Damage or disease of the retina, the light-sensing layer at the back of the eye. Diabetes, high blood pressure, prematurity, blocked vessels, aging, inflammation, and inherited conditions can cause it and may lead to sight loss.

\medterm{Retinopathy Of Prematurity} Abnormal growth of blood vessels in the retina of a premature or very low-birth-weight infant. Mild disease may resolve, but severe disease can cause retinal detachment and blindness, so eligible infants need scheduled eye examinations.

\medterm{Retinoschisis} Splitting between layers of the retina. It may result from aging or an inherited condition and can cause reduced sharpness or loss of part of the visual field; it differs from retinal detachment but may look similar on examination.

\medterm{Retinoscope} An instrument that shines light into the eye to help determine refractive error. A lighted instrument used to estimate focusing error by observing how light reflects from the retina.

\medterm{Retinoscopy} An eye examination that uses a retinoscope to estimate refractive error by observing how light reflects from the retina.

\medterm{Retirement} Retirement is withdrawal from regular employment or a career; changes in income, activity, social contact, and health coverage can affect wellbeing.

\medterm{Retracted Nipple} Alternate name; see \textit{Nipple Inversion Or Retraction}.

\medterm{Retractile Testicle} Alternate name; see \textit{Retractile testicle}.

\medterm{Retractile Testis} A testis that moves between the scrotum and groin because of a strong reflex but can be gently brought into the scrotum. It should be monitored because it may later become undescended.

\medterm{Retractor} A surgical instrument used to hold tissue or an organ aside so the area can be seen and reached during an examination or operation. Some retractors are held by hand; others stay in place by themselves.

\medterm{Retrobulbar Hemorrhage} Bleeding behind the eyeball within the orbit, often after trauma or surgery, that can rapidly increase orbital pressure and threaten vision.

\medterm{Retrobulbar Neuritis} Inflammation of the optic nerve behind the eyeball, causing reduced or blurred vision and pain, especially when the eye moves. It may be associated with multiple sclerosis or another immune condition and needs prompt assessment.

\medterm{Retroesophageal} Retroesophageal means located behind the esophagus, such as a mass, lymph node, vessel, or collection in the posterior mediastinum.

\medterm{Retrognathia} Retrognathia is backward positioning of the lower jaw relative to the upper jaw and may affect facial structure, bite, airway, or sleep breathing.

\medterm{Retrograde} Moving backward or against the normal direction, such as retrograde amnesia or retrograde blood flow.

\medterm{Retrograde Amnesia} Loss of memory for events that occurred before an injury, illness, or other triggering event, with variable preservation of more remote memories.

\medterm{Retrograde Cystography} Retrograde cystography uses contrast placed through the urethra to outline the bladder and detect leaks, injury, or structural abnormalities.

\medterm{Retrograde Degeneration} Breakdown of a nerve fiber that spreads back toward the nerve cell body after the fiber is injured. It can reduce movement, sensation, or organ function depending on which nerve is affected and how severe the injury is.

\medterm{Retrograde Ejaculation} A condition in which semen enters the bladder instead of leaving through the penis during orgasm. It may cause little or no semen to appear and can result from diabetes, surgery, nerve problems, or medicines. It may affect fertility.

\medterm{Retrograde Intrarenal Surgery} RIRS is a minimally invasive procedure using a flexible scope passed through the urinary tract to inspect and fragment or remove kidney stones.

\medterm{Retrograde Urethrogram} An X-ray study that outlines the urethra after contrast material is introduced through its opening. It can show narrowing, tears, blockage, or abnormal connections, especially after suspected urethral injury.

\medterm{Retroperitoneal Fibrosis} Formation of abnormal scar-like tissue in the space behind the abdominal organs. It often surrounds the tubes draining urine from the kidneys, causing back or abdominal pain, urine blockage, and reduced kidney function; medicines, immune disease, or cancer can cause it.

\medterm{Retroperitoneal Hematoma} A collection of extravasated blood within the retroperitoneal anatomical space, commonly provoked by pelvic fractures, penetrating or blunt abdominal trauma, rupture of an abdominal aortic aneurysm, or femoral vascular catheterization; classified into zones I (central), II (flank), and III (pelvic).
\medterm{Retroperitoneal Hemorrhage} Bleeding in the space behind the abdominal lining. It may follow injury, blood-thinning treatment, a procedure, or rupture of a blood vessel and can cause back or abdominal pain, weakness, low blood pressure, and life-threatening blood loss.

\medterm{Retroperitoneal Inflammation} Retroperitoneal inflammation is inflammatory disease in the space behind the peritoneum, caused by infection, pancreatitis, urinary disease, trauma, medicines, or fibrosis.

\medterm{Retroperitoneal Neoplasm} A retroperitoneal neoplasm is an abnormal growth in the tissue space behind the abdominal lining and may be benign or malignant.

\medterm{Retroperitoneal Space} The retroperitoneal space is the area behind the peritoneal lining containing the kidneys, adrenal glands, pancreas, parts of the bowel, vessels, nerves, and lymphatics.

\medterm{Retroperitoneum} The space behind the abdominal lining. It contains the kidneys, adrenal glands, ureters, major blood vessels, and parts of the bowel and pancreas; bleeding, infection, tumors, or inflammation there can cause back or abdominal symptoms.

\medterm{Retropharyngeal Abscess} A collection of pus in the tissues behind the throat, usually after infection or injury. Fever, severe throat or neck pain, difficulty swallowing, drooling, noisy breathing, or a stiff neck can occur. The abscess can threaten the airway.

\medterm{Retropharyngeal Space} A potential fascial space in the neck located between the buccopharyngeal fascia of the pharynx anteriorly and the alar layer of the prevertebral fascia posteriorly, extending from the skull base down into the superior mediastinum (providing an anatomical conduit for the spread of retropharyngeal abscesses into the mediastinum).

\medterm{Retrospective Studies} Retrospective studies examine existing records or past exposures and outcomes; they are efficient but vulnerable to missing data, bias, and confounding.

\medterm{Retrospective Study} A study that examines existing records or past exposures to investigate disease causes, treatments, or outcomes. It can be efficient, but missing information and differences between groups may limit reliable conclusions.

\medterm{Retrosternal Thyroid Surgery} An operation to remove part or all of a thyroid gland that extends behind the breastbone. It may relieve pressure on the airway or swallowing tube and requires planning around nearby nerves, blood vessels, and the parathyroid glands.

\medterm{Retrotracheal} Retrotracheal means located behind the trachea, such as a mass, lymph node, vessel, or collection in the posterior mediastinum.

\medterm{Retrotransposition} The copying and insertion of a genetic sequence into a new place in the genome using an RNA copy as an intermediate. New insertions can be harmless or disrupt a gene and cause disease.

\medterm{Retroversion} Backward tilting or turning of an organ, especially the uterus toward the spine. A backward-tilted uterus is often normal, but in some people it is linked with pain, endometriosis, or difficult examinations.

\medterm{Retroversion Of The Uterus} A uterine position in which the uterus tilts backward toward the spine; it is often a normal variant but can be associated with pain or endometriosis.

\medterm{Retroviridae} A family of RNA viruses that make a DNA copy of their genetic material and insert it into infected cells. HIV is a member; these viruses can cause persistent infection, immune disease, or cancer.

\medterm{Retroviridae Infection} An infection caused by a retrovirus, which may persist through genomic integration and cause immune, neurologic, or malignant disease.

\medterm{Retrovirus} A type of virus that uses RNA as its genetic material instead of DNA. It uses an enzyme called reverse transcriptase to make a DNA copy of its RNA. This copy can become part of the DNA of an infected cell, allowing the virus to make more copies. HIV is an example of a retrovirus.

\medterm{Rett Syndrome} A rare genetic disorder that affects brain development and occurs almost exclusively in girls. Development may appear typical for the first 6 to 18 months, followed by loss of previously learned language, purposeful hand use, coordination, or movement. Repeated hand movements such as hand-wringing or clapping are common. Slowed head growth, seizures, breathing changes, feeding problems, sleep problems, and curvature of the spine (scoliosis) may occur. Most cases are caused by a new change in the \textit{MECP2} gene; rare males with related gene changes can also be affected.

\medterm{Return Of Spontaneous Circulation} Restoration of a sustained heartbeat and blood flow after cardiac arrest, shown by a pulse, blood pressure, or other signs of circulation. Intensive post-arrest care then supports breathing, blood pressure, brain function, temperature, and treatment of the cause.

\medterm{Reuptake} The removal of a released nerve chemical from the space between nerve cells by returning it to the releasing cell or taking it into a nearby cell. This limits and ends the signal.

\medterm{Revalidation} Formal reassessment that a healthcare professional remains fit, competent, and up to date to practice.

\medterm{Revascularisation} Restoration of blood flow to an organ or limb by a procedure such as angioplasty, stenting, or surgical bypass.

\medterm{Revascularization} Restoration of blood flow to an area by bypass surgery, balloon treatment, stent placement, or another procedure. It can relieve symptoms, preserve tissue, and reduce damage caused by a blocked or narrowed vessel.

\medterm{Revascularization Surgery} Revascularization surgery restores blood flow around a blocked or narrowed vessel, such as in coronary artery bypass or limb bypass surgery.

\medterm{Reverberation Artifact} An ultrasound image error caused by repeated reflections, sometimes reduced by changing the angle or using specialized imaging.

\medterm{Reversal Learning} Reversal learning is the ability to change a learned response when the previously rewarded or correct choice is no longer so.

\medterm{Reverse Total Shoulder Replacement} Shoulder-replacement surgery that places the rounded joint surface on the shoulder blade and the socket on the upper-arm implant. This changes shoulder mechanics and can help when severe arthritis occurs with an irreparable rotator cuff.

\medterm{Reverse Transcriptase} An enzyme that synthesizes DNA using RNA as a template.

\medterm{Reverse Transcriptase Inhibitor} An antiviral medicine that blocks reverse transcriptase, used in combination treatment for HIV or selected infections.

\medterm{Reverse Transcriptase Inhibitors} A major class of antiretroviral medications utilized in the management of HIV-1 infection and chronic hepatitis B, categorized into nucleoside/nucleotide reverse transcriptase inhibitors (NRTIs, such as emtricitabine and tenofovir) and non-nucleoside reverse transcriptase inhibitors (NNRTIs, such as efavirenz and doravirine).

\synonyms
RTIs

\medterm{Reverse Transcription} The production of DNA using RNA as a template, carried out by reverse transcriptase. Some viruses use this process to copy their genetic material, and laboratories use it to study RNA.

\medterm{Reverse Trendelenburg Position} A supine position in which the head and upper body are elevated higher than the feet. The head-up tilt uses gravity to shift abdominal organs downward, improving surgical visualization during head, neck, upper gastrointestinal, and laparoscopic procedures while facilitating ventilation.

\synonyms
Head-Up Tilt Position; Anti-Trendelenburg Position

\medterm{Reversible Cell Injury} Temporary cell damage that can resolve when the harmful cause is removed. Cells may swell, store excess fat, or produce less energy, but their membranes and genetic material remain sufficiently intact to recover.

\medterm{Revision Arthroplasty} Surgical replacement, exchange, or repair of a failed or painful joint prosthesis.
Indications include infection, aseptic loosening, instability, periprosthetic fracture,
wear, and recurrent dislocation.

\medterm{Revision Total Hip Arthroplasty} Surgery to remove and replace some or all parts of a previous hip replacement that have loosened, worn out, become infected, dislocated repeatedly, or failed. It is more complex than first-time replacement and may require bone reconstruction.

\medterm{Revision Total Knee Arthroplasty} Surgery to remove and replace parts of a previous knee replacement because of loosening, wear, infection, instability, stiffness, or persistent pain. It may require specialized implants or bone grafting and has higher risks than the first operation.

\medterm{Reye-Johnson Syndrome} Alternate name; see \textit{Reye Syndrome}.

\medterm{Reye's Syndrome} A rare, potentially fatal illness causing acute brain dysfunction and liver injury, associated particularly with viral illness and aspirin exposure in children.

\medterm{Reye Syndrome} A rare, serious illness causing sudden brain swelling and liver dysfunction, usually in a child or teenager recovering from a viral infection. Repeated vomiting, confusion, unusual sleepiness, seizures, or loss of consciousness may occur; aspirin exposure is associated with increased risk.

\synonyms
Reye-Johnson Syndrome; Fatty Liver with Encephalopathy

\medterm{Reynolds Pentad} A combination of fever or chills, pain in the upper right abdomen, jaundice, low blood pressure, and confusion. Together these signs suggest severe infection and blockage of the bile ducts, which requires emergency treatment.

\medterm{Rezafungin} A long-acting echinocandin antifungal medicine that inhibits fungal cell-wall synthesis. It is used for selected invasive
candidiasis infections; its long duration permits an intermittent dosing schedule under
specialist supervision.

\medterm{R-Factor} A laboratory measure comparing the amount of an antibiotic needed to stop a bacterium with the expected effective amount. A higher value suggests reduced susceptibility and may help guide antibiotic choice.

\medterm{Rgs Proteins} Regulators of G-protein signaling are proteins that speed the switch-off of signals from G-protein-coupled receptors. They help control responses to hormones, neurotransmitters, and other signals and may influence disease when altered.

\medterm{Rhabdoid Tumor} A rare, fast-growing cancer that usually affects infants and young children. It may begin in the kidney, brain, or other tissues; changes in genes that control cell growth, especially SMARCB1, are common.

\medterm{Rhabdomyolysis} Rhabdomyolysis is rapid breakdown of skeletal muscle that releases muscle proteins and salts into the bloodstream. Causes include crush injury, extreme exertion, seizures, heat illness, medicines, and metabolic disease. It may cause muscle pain, weakness, dark urine, kidney injury, abnormal heart rhythms, or dangerous pressure in a limb. A high creatine kinase level supports the diagnosis.

\medterm{Rhabdomyolysis And AKI} Severe muscle breakdown releases muscle proteins and salts, potentially causing acute kidney injury, abnormal body salts, and abnormal heart rhythms. Findings may include muscle pain, weakness, dark urine, and a high creatine kinase level.

\medterm{Rhabdomyoma} A usually noncancerous tumor made of heart-muscle cells, seen mainly in infants and children. It is strongly associated with tuberous sclerosis and may shrink on its own, but large or multiple tumors can block blood flow or disturb the heart rhythm.

\medterm{Rhabdomyosarcoma} Rhabdomyosarcoma is a malignant tumor showing skeletal-muscle differentiation, occurring most often in children but also in adults.

\medterm{Rhabdomyosarcoma Pathology} A malignant mesenchymal tumor showing skeletal-muscle differentiation. Histologic patterns include embryonal, alveolar, spindle-cell, and pleomorphic forms, with primitive cells, rhabdomyoblasts, and variable myogenin or MyoD1 expression.

\medterm{Rhegmatogenous Retinal Detachment} The most common type of retinal detachment, caused by a hole or tear that lets the eye’s internal fluid pass beneath the retina. It can cause flashes, floaters, a curtain-like shadow, or sight loss.

\medterm{Rhesus Factor} The Rh blood-group antigen, especially the D antigen, whose presence or absence is important in transfusion and pregnancy.

\medterm{Rhesus Incompatibility} Rhesus incompatibility occurs when an RhD-negative pregnant person forms antibodies against RhD-positive fetal red cells, risking haemolytic disease in a later or affected pregnancy; anti-D immunoglobulin prevents sensitization.

\medterm{Rhesus System} The Rh blood-group system, a group of inherited red-cell antigens including D, C, c, E, and e. RhD status is commonly reported as positive or negative and is important in transfusion compatibility and prevention of hemolytic disease of the fetus and newborn.

\medterm{Rheumatic Chorea} Involuntary, irregular, fidgety movements caused by an immune reaction after group A streptococcal infection. It is a feature of acute rheumatic fever and may affect the face, arms, legs, speech, or handwriting.

\medterm{Rheumatic Disease} A broad group of disorders affecting joints, muscles, bones, connective tissue, or immune function, including arthritis and lupus.

\medterm{Rheumatic Fever} An inflammatory complication of untreated or inadequately treated group A streptococcal throat infection, affecting joints, heart, skin, or brain.

\medterm{Rheumatic Heart Disease} Permanent damage to heart valves after rheumatic fever, an immune reaction following group A streptococcal throat infection. It most often affects the mitral valve and can cause breathlessness, tiredness, swelling, abnormal rhythms, or heart failure.

\medterm{Rheumatism} Pain, stiffness, or discomfort in joints and surrounding soft tissues.

\medterm{Rheumatoid Arthritis} A long-term autoimmune disease in which the immune system mistakenly attacks the lining of the joints. This causes inflammation, pain, swelling, warmth, and stiffness, often worse after waking or resting. It commonly affects the same joints on both sides of the body, especially the hands, wrists, and feet. Over time, the inflammation can damage cartilage and bone and reduce joint function. It may also cause tiredness, low-grade fever, or loss of appetite and can affect the eyes, lungs, heart, blood vessels, skin, nerves, or blood cells.

\medterm{Rheumatoid Arthritis and Pregnancy} Rheumatoid arthritis often changes during pregnancy and may flare after delivery. Care balances disease control, joint function, medication safety, contraception or fertility goals, and fetal or newborn considerations, with coordinated rheumatology and obstetric review.

\medterm{Rheumatoid Arthritis Cause Gastrointestinal Issues} Rheumatoid arthritis can affect the gastrointestinal system indirectly through inflammation, associated autoimmune disease, reduced mobility, or medicines such as nonsteroidal anti-inflammatory drugs. Abdominal pain, bleeding, persistent diarrhea, or weight loss needs evaluation rather than attribution to arthritis alone.

\medterm{Rheumatoid Arthritis Inflammation Of The Brain} Direct brain inflammation is uncommon in rheumatoid arthritis, but vasculitis, medication effects, infection, stroke risk, or another autoimmune process can cause neurologic symptoms. New confusion, seizure, weakness, severe headache, or vision change requires urgent assessment.

\medterm{Rheumatoid Factor} An antibody that mistakenly recognizes part of another antibody. A blood test may support the diagnosis of rheumatoid arthritis, but it can also be positive in other diseases or in healthy people.

\medterm{Rheumatoid Lung Disease} Lung involvement associated with rheumatoid arthritis, including interstitial lung disease, pleural disease, nodules, or airway abnormalities.

\medterm{Rheumatoid Nodule} A rheumatoid nodule is a firm inflammatory lump under the skin or near a joint in rheumatoid arthritis, often over pressure points.

\medterm{Rheumatoid Nodules} Firm inflammatory lumps under the skin or near joints in people with rheumatoid arthritis, often over pressure points. They may cause no symptoms but can become tender or be affected by pressure.

\medterm{Rheumatoid Pneumoconiosis} Rheumatoid pneumoconiosis is lung disease involving rheumatoid arthritis and dust-related lung nodules, historically called Caplan syndrome.

\medterm{Rheumatoid Synovitis} Long-lasting immune inflammation of the tissue lining a joint in rheumatoid arthritis. It causes pain and swelling and can gradually damage cartilage, bone, and joint movement.

\medterm{Rheumatoid Vasculitis} Inflammation of small or medium-sized blood vessels associated with rheumatoid arthritis. It can reduce blood flow and affect the skin, nerves, eyes, or internal organs.

\medterm{Rheumatologist} A medical doctor who diagnoses and treats diseases affecting joints, muscles, bones, and connective tissues, especially autoimmune and inflammatory conditions such as rheumatoid arthritis, lupus, gout, and vasculitis.

\medterm{Rheumatology} The medical specialty concerned with inflammatory, autoimmune, and musculoskeletal disorders.

\medterm{Rh Factor} A system of antigens on the surface of red blood cells, especially the D antigen. A person
is Rh-positive when the D antigen is present and Rh-negative when it is absent. Rh incompatibility can affect pregnancy
and transfusion safety.

\synonyms
Rh D Antigen; Rhesus Antigen

\medterm{Rh Incompatibility} An immune reaction that can occur when an Rh-negative pregnant person carries an Rh-positive fetus and forms antibodies against fetal red blood cells. Anti-D immune globulin can prevent sensitization.

\medterm{Rhinencephalon} An older term for brain regions involved mainly in smell and closely connected with emotion and memory. Modern anatomy usually names the individual structures instead of using this broad term.

\medterm{Rhinitis} Inflammation or irritation of the lining of the nose. It can cause a blocked or runny nose, sneezing, itching, and reduced smell and may result from allergy, infection, smoke, strong smells, medicines, or other triggers.

\medterm{Rhinitis Medicamentosa} A condition of severe rebound nasal mucosal congestion and edema caused by prolonged, repetitive use of topical alpha-adrenergic nasal decongestant sprays (such as oxymetazoline) beyond 3 to 5 days.
\synonyms
Rebound Nasal Congestion; Topical Decongestant Rhinitis

\medterm{Rhinocephaly} A very rare birth defect involving severe abnormal development of the face and forebrain, with a nose-like structure that may lie unusually high or between the eyes. It is usually associated with major neurologic and developmental abnormalities.

\medterm{Rhinocerebral Mucormycosis} A rapidly progressive, fulminant, angioinvasive fungal infection caused by molds of the order Mucorales (such as \textit{Rhizopus} and \textit{Mucor}), occurring primarily in patients with diabetic ketoacidosis, severe neutropenia, or hematologic malignancies. Fungal hyphae invade blood vessels causing extensive thrombosis, tissue infarction, and black necrotic eschars of the nasal mucosa and hard palate, rapidly extending into the orbit and cavernous sinus.

\synonyms
Invasive Rhinocerebral Zygomycosis

\medterm{Rhinophyma} Progressive thickening and enlargement of the nose caused by long-standing changes often associated with rosacea. It may produce a rough, irregular appearance and can sometimes interfere with breathing or require surgical treatment.

\medterm{Rhinoplasty} Surgery to reshape, repair, or reconstruct the nose for breathing, injury, birth differences, or appearance. Recovery includes swelling and bruising, and results depend on anatomy, technique, healing, and patient goals.

\medterm{Rhinorrhea} Excessive watery or mucoid discharge from the nose, commonly caused by infection, allergy, irritants, or changes in temperature.

\medterm{Rhinorrhoea} British spelling of rhinorrhea, meaning drainage of watery or mucus-like fluid from the nose. Causes include allergy, infection, irritation, and rarely a cerebrospinal-fluid leak.

\medterm{Rhinoscleroma} A chronic, slowly progressive granulomatous infection of the nose and upper airway caused by Klebsiella rhinoscleromatis. It may progress from foul-smelling discharge and obstruction to scarring, deformity, and airway narrowing and usually needs prolonged antibiotics.

\medterm{Rhinoscopy} Examination of the nasal cavity and nasopharynx using an instrument or endoscope to assess the mucosa, septum, turbinates, masses, bleeding, or obstruction.

\medterm{Rhinosporidiosis} A chronic localized granulomatous disease of the mucous membranes caused by Rhinosporidium seeberi, an aquatic protistan parasite in the clade Mesomycetozoea, endemic primarily to rural wetland communities in India and Sri Lanka. It classically manifests in the anterior nasal cavity or nasopharynx as painless, friable, highly vascular, sessile or pedunculated polypoid lesions with a characteristic "strawberry-like" appearance dotted with visible white spherules (sporangia), causing nasal obstruction, epistaxis, and foreign-body sensation.

\synonyms
Nasal Rhinosporidiosis; Rhinosporidium Seeberi Infection

\medterm{Rhinovirus} A genus of non-enveloped, positive-sense RNA viruses in the Picornaviridae family. Rhinoviruses
are a common cause of the common cold and spread mainly through respiratory secretions and contaminated hands or
surfaces.

\medterm{Rhinovirus A} A species group of rhinoviruses that commonly causes the common cold and can worsen asthma or trigger wheezing.

\medterm{Rhinovirus C} A rhinovirus group associated with common-cold symptoms and, in some children, more severe wheezing or lower-respiratory illness.

\medterm{Rhinoviruses} A group of viruses that commonly cause colds. They spread through droplets, hands, and contaminated surfaces and may worsen asthma or cause more serious breathing illness in vulnerable people.

\medterm{Rh Isoimmunization} Maternal alloantibody production against fetal Rh D-positive red blood cells in an Rh
D-negative mother, occurring when fetal RBCs enter the maternal circulation
(fetomaternal hemorrhage).

\medterm{Rhizomucor} A genus of molds that can cause mucormycosis, especially in people with diabetes, neutropenia, or severe immune suppression.

\medterm{Rhizomucor Pusillus} A mold that can cause invasive mucormycosis involving the lungs, sinuses, skin, or other organs in severely immunocompromised people.

\medterm{Rhizopus} A genus of molds that includes important causes of mucormycosis, a rapidly invasive infection of sinuses, lungs, skin, or other tissues.

\medterm{Rhizopus Oryzae} A mold, also called Rhizopus arrhizus, that is a common cause of invasive mucormycosis.

\medterm{Rhizotomy} A procedure that cuts or heats selected nerve roots to reduce severe pain, muscle stiffness, or abnormal reflexes. It is used only in carefully selected cases because it can cause permanent numbness, weakness, or other nerve problems.

\medterm{Rhodium} A chemical element used in industrial and dental materials; its compounds can irritate tissues and may cause allergic reactions.

\medterm{Rhodococcus} A genus of gram-positive bacteria that includes Rhodococcus equi, an opportunistic pathogen causing pneumonia mainly in immunocompromised people.

\medterm{Rhodococcus Equi} A gram-positive bacterium that can cause cavitary pneumonia and disseminated infection, especially in people with advanced immune suppression.

\medterm{Rhodopsin} A light-sensitive protein in rod cells of the retina that enables vision in dim light.

\medterm{Rhodospirillaceae} A family of metabolically diverse bacteria found in soil and water; some species carry out photosynthesis or nitrogen fixation.

\medterm{Rhodospirillum} A genus of spiral-shaped aquatic bacteria capable of photosynthesis or nitrogen fixation; human infection is rare.

\medterm{RhoGAM} A brand name for Rh immune globulin, which contains anti-D antibodies. It is given to eligible RhD-negative pregnant people after specified pregnancy events and around delivery to reduce formation of antibodies that could harm an RhD-positive fetus or newborn.

\medterm{Rhombencephalon} The hindbrain region of an early embryo. It develops into the pons, medulla, cerebellum, and related structures that help control breathing, balance, movement, swallowing, and other essential functions.

\medterm{Rhonchi} Low-pitched breathing sounds caused by air moving through narrowed or mucus-filled larger airways, often changing after coughing.

\medterm{Rhonchus} An abnormal sound heard through a stethoscope during breathing, caused by air moving through fluid, mucus, narrowed airways, or abnormal lung tissue.

\synonyms
Crackles; Crepitations; Rales

\medterm{Rhubarb Leaves Poisoning} Poisoning after eating large amounts of rhubarb leaves, which contain oxalates and other irritants. It may cause burning, vomiting, abdominal pain, low calcium, kidney injury, weakness, or abnormal heart rhythms.

\medterm{Rhythm Method} A fertility-awareness method that estimates fertile days from menstrual-cycle timing to avoid or achieve pregnancy.

\medterm{Rib} One of 12 paired, curved bones forming the thoracic cage, protecting the heart, lungs,
and great vessels. True ribs (1--7) attach directly to the sternum via costal cartilage.
False ribs (8--10) attach indirectly to the sternum via the costal cartilage of rib 7.
Floating ribs (11--12) have no anterior attachment.

\synonyms
Costal Bone; Costa

\medterm{Rib Angle} The point where a rib bends most sharply as it curves from the spine toward the front of the chest.

\medterm{Ribavirin} An antiviral medicine used in selected viral infections, usually in combination and under specialist supervision.

\medterm{Rib Cage} The rib cage consists of ribs, sternum, thoracic vertebrae, muscles, and connective tissues that protect chest organs and assist breathing.

\medterm{Ribcage Pain} Ribcage pain may arise from muscles, ribs, joints, lungs, heart, or other organs; sudden severe pain or breathing difficulty may indicate a serious condition.

\medterm{Rib Fracture} A break in one or more ribs, usually causing localized chest pain worsened by breathing, coughing, or movement and sometimes associated with injury to underlying organs.

\medterm{Rib Head} The back end of a rib that joins with the bodies of one or two thoracic vertebrae.

\medterm{Rib Neck} The narrowed part of a rib between its head and tubercle.

\medterm{Riboflavin} A water-soluble vitamin that forms the coenzymes FAD and FMN. These cofactors
participate in oxidation-reduction reactions, including electron transfer, fatty-acid
oxidation, and energy metabolism.

\synonyms
Lactoflavin; Vitamin G

\medterm{Riboflavin Deficiency} Deficiency of vitamin B2, which can cause inflammation or fissuring of the mouth and lips, sore throat, dermatitis, anemia, and other mucosal or systemic findings.

\medterm{Ribonuclease} An enzyme that cuts RNA, a molecule involved in reading genetic instructions, into smaller pieces. Ribonucleases help control how long RNA lasts, recycle cell material, and protect cells from some unwanted genetic material.

\medterm{Ribonuclease H} An enzyme that removes the RNA strand when RNA and DNA are joined together. Cells use it during DNA copying, and some viruses depend on related activity when making DNA from their RNA.

\medterm{Ribonuclease III} An enzyme that cuts double-stranded RNA into smaller pieces. The products can help control gene activity, make proteins, or process genetic material in cells and some microorganisms.

\medterm{Ribonucleic Acid} RNA is a nucleic acid that carries genetic instructions, helps make proteins, and regulates cell activity; some viruses use RNA as their genetic material.

\medterm{Ribonucleoprotein} A structure made of RNA joined to one or more proteins. It helps make, process, transport, or control RNA; antibodies against certain ribonucleoproteins can occur in autoimmune diseases.

\medterm{Ribonucleoside} A molecule made of a ribose sugar joined to one of the bases used in RNA. Examples include adenosine, guanosine, cytidine, and uridine; adding phosphate groups produces RNA building blocks.

\medterm{Ribonucleotide Reductase} An enzyme that changes RNA building blocks into the related building blocks needed to make DNA. Cells require it to copy and repair DNA, so it is important for growth and is a target of some cancer medicines.

\medterm{Ribonucleotides} Nucleotides containing ribose sugar and RNA bases, serving as RNA building blocks and cellular energy or signaling molecules.

\medterm{Ribophagy} A cell-cleaning process that breaks down some ribosomes, the structures that make proteins, so their parts can be recycled. It helps cells adjust protein production when nutrients or growth conditions change.

\medterm{Ribose} A five-carbon sugar used to build RNA, ATP, and other molecules involved in storing genetic information and transferring energy. Cells make ribose from other nutrients and use it in many normal chemical processes.

\medterm{Ribosomal Protein} A protein that combines with ribosomal RNA to form a ribosome, the cell structure that builds proteins. Inherited changes in some ribosomal proteins can cause bone-marrow failure or developmental syndromes.

\medterm{Ribosomal RNA} RNA that forms the main structural and working part of ribosomes, the cell structures that make proteins. It helps position messenger RNA and joins amino acids together; its genes are also used in some laboratory studies to identify organisms.

\medterm{Ribosome} A cellular structure made of RNA and proteins that reads messenger RNA and links amino acids together to make proteins.

\medterm{Ribosome Subunit} One of the two RNA-and-protein parts of a ribosome. The small and large parts join around messenger RNA and work together to assemble amino acids into a protein.

\medterm{Rib Tubercle} A small projection near the back of a rib that joins with the transverse process of a thoracic vertebra.

\medterm{Richter's Syndrome} Transformation of chronic lymphocytic leukemia or small lymphocytic lymphoma into an aggressive lymphoma, most often diffuse large B-cell lymphoma.

\medterm{Richter Transformation} The rapid clinico-pathological transformation of chronic lymphocytic leukemia (CLL) or small lymphocytic lymphoma (SLL) into an aggressive, high-grade diffuse large B-cell lymphoma (DLBCL), presenting with rapidly enlarging lymphadenopathy, systemic B symptoms, and elevated serum LDH.
\synonyms
Richter's Syndrome; High-Grade Transformation of CLL

\medterm{Ricin} Ricin is a highly toxic plant protein that blocks protein synthesis; ingestion, inhalation, or injection can cause severe organ injury and death.

\medterm{Rickets} Defective mineralization of growing bone, usually due to vitamin D, calcium, or phosphate deficiency.

\synonyms
Nutritional Rickets; Hypocalcemic Rickets

\medterm{Rickettsia} A group of small intracellular bacteria transmitted mainly by arthropods and causing diseases such as typhus and spotted fevers.

\medterm{Rickettsiaceae} A family of bacteria that live inside cells and includes several organisms spread by ticks, fleas, lice, or mites. They can cause fever, rash, blood-vessel inflammation, and serious organ complications.

\medterm{Rickettsia Felis} A flea-associated bacterium that can cause a spotted-fever-like illness in people. Infection may produce fever, headache, rash, fatigue, or muscle aches and is assessed with exposure history and laboratory testing.

\medterm{Rickettsia Infections} Rickettsial infections are illnesses caused by intracellular bacteria spread mainly by ticks, fleas, lice, or mites and often causing fever and rash.

\medterm{Rickettsia Japonica} A bacterium spread by ticks and associated with Japanese spotted fever. Infection may cause fever, headache, rash, and a dark mark where the tick bit; treatment uses suitable antibiotics.

\medterm{Rickettsialpox} An infection spread by house mites and caused by Rickettsia akari. It usually begins with a small dark sore at the bite site, followed by fever and a blister-like rash.

\medterm{Rickettsia Prowazekii} A bacterium that causes epidemic typhus and is mainly spread through infected body-lice feces. Disease can cause high fever, headache, rash, confusion, and serious organ complications without prompt antibiotic treatment.

\medterm{Rickettsia Rickettsii} A bacterium spread by infected ticks that causes Rocky Mountain spotted fever. Fever, severe headache, rash, muscle aches, or confusion may develop; prompt antibiotic treatment is important because the illness can rapidly become life-threatening.

\medterm{Rickettsia Typhi} A bacterium spread mainly by infected fleas and responsible for murine typhus. The illness commonly causes fever, headache, chills, and sometimes a rash; prompt antibiotic treatment usually works, but untreated infection can become severe.

\medterm{Riddoch Phenomenon} A form of visual statuagnosia in which a patient with homonymous hemianopia or cortical blindness from an occipital lobe lesion is able to perceive moving visual targets within the blind field while remaining unable to visualize stationary, static stimuli.
\synonyms
Statuagnosia; Visual Motion Perception Preservation

\medterm{Ridged Cranial Sutures} Ridged cranial sutures are raised skull seams in an infant and may reflect normal molding or premature suture fusion, which requires assessment for craniosynostosis.

\medterm{Ridged Nail} A ridged nail has longitudinal or transverse grooves in its surface, which may result from aging, trauma, inflammation, nutritional problems, or systemic disease.

\medterm{Ridge Preservation} A dental procedure performed after a tooth is removed to reduce loss of the jawbone that held the tooth. Material is placed in the empty socket, sometimes with a covering, to help maintain shape for healing or a future implant.

\medterm{Riedel's Thyroiditis} A rare form of thyroid inflammation in which tough scar-like tissue makes the thyroid hard and may spread into nearby structures. It can cause a neck lump, difficulty swallowing or breathing, or reduced thyroid function.

\medterm{Riedel Thyroiditis} A rare, chronic fibroinflammatory disorder of the thyroid gland (an organ manifestation of IgG4-related disease) characterized by dense, invasive collagenous fibrous tissue replacing normal thyroid tissue and extending into adjacent strap muscles and trachea, presenting as a stony-hard, fixed ('woody') goiter.
\synonyms
Invasive Fibrous Thyroiditis; Woody Thyroiditis

\medterm{Rieger Syndrome} A genetic condition affecting development of the front of the eye and often the teeth. It may cause glaucoma, unusual teeth, or changes around the navel and can result from changes in genes such as PITX2 or FOXC1.

\medterm{Riehl Melanosis (Pigmented Contact Dermatitis)} An acquired gray-brown to slate-colored facial or neck hyperpigmentation associated with contact or photocontact exposure to fragrances, cosmetics, dyes, or other chemicals. Itching or redness may precede pigmentation; identifying and avoiding the trigger, sometimes with patch or photopatch testing, helps guide management.

\medterm{Rifabutin} A rifamycin antibiotic used mainly for selected mycobacterial infections and tuberculosis regimens; it has important drug interactions and can discolor body fluids.

\medterm{Rifampicin} A rifamycin antibiotic used in tuberculosis and selected other bacterial infections, usually in combination to prevent resistance.

\medterm{Rifampin} Rifampin is a rifamycin antibiotic that inhibits bacterial DNA-dependent RNA polymerase and is a key component of tuberculosis treatment; it causes orange body fluids and has many drug interactions and hepatotoxicity risk.

\medterm{Rifampin Toxicity} Adverse pharmacological effects of the first-line rifamycin antibiotic rifampin, including benign orange-red discoloration of body secretions (urine, tears, sweat), drug-induced hepatotoxicity, an immune-mediated flu-like syndrome with intermittent dosing, and potent induction of hepatic cytochrome P450 enzymes (CYP3A4).

\medterm{Rifamycin} Rifamycin is an antibiotic that inhibits bacterial RNA polymerase; related rifamycin medicines are used for tuberculosis and other infections.

\medterm{Rifaximin} Rifaximin is a minimally absorbed rifamycin antibiotic used for selected intestinal infections, prevention of recurrent hepatic encephalopathy, and some diarrhoeal syndromes; indication and resistance patterns guide use.

\medterm{Rift Valley Fever} A mosquito-borne viral disease affecting animals and humans, usually causing fever but sometimes severe eye, brain, or hemorrhagic disease.

\medterm{Right Atrium} The upper-right chamber of the heart. It receives oxygen-poor blood from the body and passes it through the tricuspid valve into the right ventricle, which sends it to the lungs.

\medterm{Right Bundle Branch} A group of heart fibers that carries the electrical signal to the right lower chamber so it contracts. If conduction is delayed or blocked, the ECG shows right bundle-branch block; this may be harmless or reflect heart or lung disease.

\medterm{Right Bundle Branch Block} Delay or interruption of electrical signals through the right bundle branch. It changes the ECG pattern and may occur without symptoms or with heart or lung disease.

\synonyms
RBBB

\medterm{Right Coronary Artery} A blood vessel that branches from the aorta and supplies much of the right side and lower part of the heart. A blockage can cause a heart attack or abnormal heart rhythm.

\medterm{Right Coronary Artery Occlusion} A blockage of the artery supplying much of the right and lower heart. It can cause a heart attack, slow or abnormal rhythms, low blood pressure, or failure of the right side of the heart.

\medterm{Right Heart Failure} Failure of the right side of the heart to pump blood effectively to the lungs. It may result from left-sided heart failure, high pressure in the lung arteries, a heart attack, lung clots, heart-muscle disease, or a stiff sac around the heart. It can cause swollen legs, neck veins, an enlarged liver, abdominal fluid, and weight gain.

\medterm{Right Heart Ventricular Angiography} An imaging test in which contrast dye and X-rays show the right lower heart chamber and how it pumps blood. It is used selectively to assess structure, movement, or blood flow during heart evaluation.

\medterm{Right Hypochondriac Region} The upper-right area of the abdomen, beneath the lower ribs. It contains much of the liver, the gallbladder, and part of the intestine. Pain there may arise from these organs.

\medterm{Right Iliac Region} The lower-right area of the abdomen, containing the appendix, beginning of the large intestine, and part of the small intestine. Pain there may occur with appendicitis, bowel disease, urinary problems, or a hernia and must be interpreted with other findings.

\medterm{Righting Reflex} The righting reflex is an automatic response that restores head and body orientation when balance or position changes.

\medterm{Right Lower Quadrant Abdominal Pain} Alternate name; see \textit{Appendicitis}.

\medterm{Right Lymphatic Duct} A short vessel that returns lymph from the right arm, right side of the head and neck, and right upper chest to the bloodstream near the right side of the neck. Damage or blockage can cause swelling in these areas.

\medterm{Right Sixth Cranial Nerve Palsy} Weakness of the nerve that moves the right eye outward. It can cause double vision, usually worse when looking to the right or at a distance, and may result from diabetes, high pressure in the skull, inflammation, stroke, or other disease.

\medterm{Right Ureterovesical Calculus} A stone lodged at the junction of the right ureter and bladder, causing flank pain,
hematuria, and urinary urgency or frequency. May require endoscopic or lithotripsy
intervention.

\medterm{Right Ventricle} The lower-right chamber of the heart. It receives oxygen-poor blood from the right atrium and pumps it through the pulmonary valve into the lung arteries to receive oxygen.

\medterm{Right Ventricular Strain} An electrocardiographic or echocardiographic pattern reflecting acute or chronic right ventricular pressure and volume overload, manifesting on ECG with T-wave inversions in leads V1 through V4, S1Q3T3 pattern, right bundle branch block, or rightward axis deviation (characteristic of acute massive pulmonary embolism).

\medterm{Rigid And Flexible Sigmoidoscopy} A sigmoidoscopy examines the rectum and lower colon with a flexible or rigid instrument. It can evaluate bleeding, pain, altered bowel habits, inflammation, polyps, or tumors; preparation, reach, comfort, and biopsy capability differ by technique.

\medterm{Rigidity} Abnormally increased resistance when a relaxed limb is moved by another person. It can occur with Parkinson disease, brain or spinal injury, severe infection, certain medicines, or other movement disorders.

\medterm{Rigler Sign} A classic abdominal radiographic sign of massive pneumoperitoneum (free intra-abdominal air) on plain supine radiographs, where both the inner (mucosal) and outer (serosal) surfaces of the bowel wall are clearly outlined by gas.
\synonyms
Double Wall Sign; Bas-Relief Sign

\medterm{Rigler Triad} Three imaging findings that strongly suggest gallstone ileus: blockage of the small intestine, air in the bile ducts, and a gallstone in the intestine. Not every person has all three findings.

\medterm{Rigor} Marked stiffness or rigidity; rigor mortis is postmortem muscle stiffening, while rigor may also mean a severe chill.

\medterm{Rigor Mortis} Stiffening of the body’s muscles after death. It develops as the muscles lose the energy needed to relax and later fades as decomposition advances; its timing is affected by temperature, activity, illness, and the environment.

\medterm{Rilpivirine} An HIV medicine that blocks an enzyme the virus needs to copy itself. It is taken with other antiretroviral medicines, and missed doses or drug interactions can allow resistance to develop.

\medterm{Riluzole} A medicine used to modestly slow disease progression in amyotrophic lateral sclerosis and to treat selected forms of spinal muscular atrophy.

\medterm{Rimmed Vacuoles} Small empty-looking spaces with a visible border seen inside muscle cells on biopsy. They indicate that muscle cells are breaking down or recycling parts of themselves and can occur in conditions such as inclusion-body myositis and GNE myopathy. The finding must be interpreted with symptoms and other tests.

\synonyms
Autophagic Rimmed Vacuoles

\medterm{Ring Block} Local anesthetic injected around the base of a finger or toe to numb the entire digit.

\medterm{Ring Chromosome} A chromosome whose two ends have joined to form a ring, often with loss or disruption of genetic material. It may cause developmental or other health problems, depending on the chromosome and changes involved.

\medterm{Ringing In The Ear} Alternate name; see \textit{Tinnitus}.

\medterm{Ringing In The Ears} Alternate name; see \textit{Tinnitus}.

\synonyms
Tinnitus; Aurium Acusticus

\medterm{Ring Vaccination} A public-health strategy in which close contacts and contacts of contacts around an identified case are vaccinated to create a protective ring and limit transmission.

\medterm{Ringworm} A contagious fungal infection of the skin, scalp, hair, or nails. It commonly causes an itchy, scaly rash with a ring-like edge; scalp infection may cause broken hairs or hair loss. The appearance and complications vary by the affected site.
\synonyms
Tinea Corporis; Dermatophytosis; Tinea

\medterm{Ringworm Of The Body} A contagious fungal infection of body skin causing an itchy, scaly, ring-shaped rash. It spreads through contact with infected people, animals, or objects and usually responds to antifungal treatment.

\medterm{Ringworm Of The Foot} Alternate name; see \textit{Athlete's foot}.

\medterm{Ringworm Of The Scalp} A fungal infection of the scalp and hair that can cause scaling, broken hairs, and hair loss.

\medterm{Rinne and Weber Tests} Complementary bedside tuning fork examinations (typically using a 512 Hz fork) to evaluate auditory hearing loss: Rinne compares air conduction to bone conduction in each ear, while Weber tests lateralization of sound when the vibrating fork is placed on the midline forehead.

\medterm{Rinnes Test} Rinne's test compares air and bone conduction with a tuning fork to help assess conductive hearing loss.

\medterm{Rinne Test} The Rinne test is a screening test that compares air conduction and bone conduction
hearing using a vibrating tuning fork (typically 512 Hz) to help differentiate between
conductive and sensorineural hearing loss.

\medterm{RIPE Therapy} Standard 4-drug regimen for active TB: Rifampin, Isoniazid, Pyrazinamide, Ethambutol.
Intensive phase 2 months, continuation phase 4 months (total 6 months). Pyridoxine
prevents isoniazid-induced neuropathy.

\medterm{Rippling Muscle Disease} A rare, benign myopathy characterized by mechanically-induced muscle contractions that
spread across the muscle as a visible wave (rippling), typically triggered by stretch,
percussion, or movement.

\medterm{Risedronate Sodium} The sodium salt of risedronate, a bisphosphonate used to prevent or treat osteoporosis and other disorders of excessive bone resorption.

\medterm{Risk} Risk is the probability of an unwanted event or outcome, considered together with its severity and the circumstances that influence it.

\medterm{Risk Adjustment} Risk adjustment accounts for differences in baseline patient characteristics so outcomes or payments can be compared more fairly across groups.

\medterm{Risk Factor} A characteristic, exposure, or behavior linked with a higher chance of disease or another outcome. It may be changeable, such as smoking, or unchangeable, such as age or inherited traits.

\medterm{Risk Management} Finding possible sources of harm, judging how likely and serious they are, and taking steps to reduce them. In healthcare, it includes safe systems, reporting problems, learning from errors, and checking whether safety measures work.

\medterm{Risk Register} A documented list of identified risks, their likelihood, impact, controls, owners, and review status.

\medterm{Risks Of Hip And Knee Replacement} Possible problems after joint replacement include infection, blood clots, bleeding, nerve or vessel injury, dislocation, stiffness, loosening, wear, persistent pain, and need for repeat surgery. Risk varies with the operation and the person's health.

\medterm{Risks Of Tobacco} Tobacco use increases risks of cancer, heart and blood-vessel disease, stroke, lung disease, pregnancy complications, addiction, and harm from secondhand smoke.

\medterm{Risks Of Underage Drinking} Alcohol use by minors can impair brain development, judgment, learning, safety, and mental health and increases risks of injury, dependence, and unsafe sexual or substance-related behavior.

\medterm{Risk-Taking} Risk-taking is behavior involving deliberate or inadvertent exposure to possible harm, influenced by perception, reward, development, mental health, and circumstances.

\medterm{Risky Drinking} Alcohol use that increases the chance of health, safety, work, relationship, or social problems. Risk depends on amount, frequency, pattern, personal health, medicines, pregnancy, age, and circumstances such as driving.

\synonyms
Hazardous Drinking; At-Risk Alcohol Consumption

\medterm{Risperidone} Risperidone is an atypical antipsychotic that blocks dopamine and serotonin receptors and treats schizophrenia, bipolar disorder, and selected severe behavioral symptoms; hyperprolactinaemia, metabolic effects, movement disorders, and sedation may occur.

\medterm{Ristocetin-Induced Platelet Aggregation} A blood test that checks how platelets clump when exposed to ristocetin, a substance that requires von Willebrand factor and a platelet surface protein. It helps investigate von Willebrand disease and some platelet disorders.

\medterm{Risus Sardonicus} A fixed, strained grin caused by painful facial-muscle spasms. It is classically associated with tetanus or strychnine poisoning and reflects severe muscle overactivity.

\medterm{Ritonavir} Ritonavir is an HIV protease inhibitor mainly used at low doses to boost levels of other antiviral medicines by slowing their breakdown.

\medterm{Ritual} A ritual is a repeated, structured behavior or ceremony; repetitive rituals may also occur as compulsions in obsessive-compulsive disorder.

\medterm{Rituximab} Rituximab is a monoclonal antibody against CD20 that depletes B cells and is used for B-cell malignancies and selected autoimmune diseases; infusion reactions, infections, hepatitis-B reactivation, and rare PML are important risks.

\medterm{Rivaroxaban} Rivaroxaban is an oral direct factor Xa inhibitor used to prevent or treat venous thromboembolism and reduce stroke risk in selected atrial fibrillation; bleeding and renal-function-dependent dosing are important considerations.

\medterm{Rivastigmine} A medicine that increases acetylcholine signaling and may temporarily improve thinking or daily function in Alzheimer disease or Parkinson disease dementia. Common effects include nausea, vomiting, diarrhea, appetite loss, dizziness, and slow heartbeat.

\medterm{River Blindness} Onchocerciasis, a parasitic infection caused by Onchocerca volvulus and transmitted by blackflies. It can
cause severe itching, skin disease, and visual impairment.

\medterm{RâLe} An older term for an abnormal crackling or bubbling respiratory sound, now generally described as a crackle.

\medterm{RLS} Restless legs syndrome, an urge to move the legs with uncomfortable sensations, especially at night.
\synonyms
Restless Legs Syndrome; Willis-Ekbom Disease

\medterm{RMDs} Abbreviation; see \textit{Repetitive Motion Disorders}.

\medterm{RNA} Ribonucleic acid, a molecule that helps cells read genetic instructions and make proteins. Some RNA carries messages from DNA, some helps build proteins, and some controls cell activity or forms part of viruses.

\medterm{RNA Binding} The attachment of a protein or other molecule to RNA, affecting RNA stability, transport, folding, processing, or translation.

\medterm{RNA Biosynthesis} The cellular production of RNA from a DNA template or, for some viruses, from an RNA template.

\medterm{RNA Degradation} The enzyme-controlled breakdown of RNA molecules. It controls how long genetic messages last, removes damaged or unnecessary messages, recycles their building blocks, and helps regulate protein production and cell behavior.

\medterm{RNA Folding} The process by which an RNA strand forms secondary and three-dimensional structures that influence its function.

\medterm{RNA Helicase} An enzyme that uses energy to separate, reshape, or move RNA structures and RNA-protein complexes. It helps cells copy, process, transport, and destroy RNA and is important for some viruses.

\medterm{RNA Interference} A gene-silencing process in which small RNAs guide proteins to degrade or block specific messenger RNA molecules.

\medterm{RNase P} An RNA-protein enzyme complex that processes precursor transfer RNA by removing its extra 5-prime sequence.

\medterm{RNA Splicing} Removal of introns and joining of exons from a precursor RNA to produce mature messenger RNA or other functional RNA.

\medterm{RNA Transport} Movement of RNA within a cell or between cells, often mediated by transport proteins and required for translation or signaling.

\medterm{RNA Virus} A virus whose genetic material is RNA; its replication often involves an RNA-dependent polymerase and may have a high mutation rate.

\medterm{Road Rage} Alternate name; see \textit{Intermittent explosive disorder}.

\medterm{Robinow Syndrome} A rare inherited skeletal-development disorder with distinctive facial features, short stature, vertebral or rib abnormalities, and genital differences.

\medterm{Robot-Assisted Rehabilitation} Robotic devices for repetitive task training. Upper extremity, lower extremity, gait
training. Objective measurement, high repetition.

\medterm{Robotic-Assisted Thoracoscopic Surgery} Chest surgery performed through several small openings using instruments controlled by a surgeon from a console. The camera provides a magnified three-dimensional view and the instruments can move precisely. It may reduce the size of the incisions, but bleeding, infection, air leakage, and conversion to open surgery remain possible.

\synonyms
RATS; Robotic Thoracic Surgery; Robotic Lobectomy

\medterm{Robotic Hysterectomy} Hysterectomy performed through minimally invasive abdominal access using a surgeon-controlled robotic platform. The extent of uterine, cervical, tubal, and ovarian removal depends on the indication and operative plan.

\medterm{Robotic Surgery} A surgical technique in which a surgeon controls a computer-assisted system that moves a camera and small instruments. It is usually performed through small incisions and may allow precise movements, but the system does not operate independently.

\synonyms
Robot-Assisted Surgery; Computer-Assisted Surgery; RAS

\medterm{Rocker-Bottom Foot} A congenital foot deformity characterized by a prominent, downward-convex plantar surface with a dorsiflexed midfoot and rigid equinovarus hindfoot (prominent calcaneus), classically caused by a congenital vertical talus and strongly associated with trisomy 18 (Edwards syndrome) and trisomy 13 (Patau syndrome).

\medterm{Rocky Mountain Spotted Fever} A severe, potentially fatal tick-borne rickettsial infection caused by \textit{Rickettsia rickettsii} transmitted by \textit{Dermacentor} ticks, presenting with high fever, severe headache, myalgias, and a petechial maculopapular rash that characteristically begins on the wrists and ankles and spreads centripetally to the trunk, palms, and soles.

\synonyms
RMSF

\medterm{Rod} A rod is a long cylindrical structure; in the retina, a rod photoreceptor supports dim-light and peripheral vision.

\medterm{Rodenticide Toxicity} A heterogeneous group of poisonings whose presentation depends on the active ingredient. Anticoagulant products cause delayed coagulopathy and bleeding; bromethalin causes neurologic toxicity, cholecalciferol causes hypercalcemia and kidney injury, and metal phosphides can cause shock and organ failure. Identify the product and provide agent-specific treatment.

\medterm{Rodent Ulcer} An older name for basal cell carcinoma, a common locally invasive skin cancer.

\medterm{Rod Photoreceptor} A light-sensitive cell in the retina that supports vision in dim light and the peripheral field. Rods do not provide color vision; their loss can cause night blindness and narrowing of the visual field.

\medterm{Rods} Retinal photoreceptor cells that provide vision in dim light and contribute to peripheral vision but do not detect color.

\medterm{Rods And Cones} Two types of light-sensing cells in the retina. Rods help with dim-light and side vision, while cones provide sharp and color vision in brighter light.

\medterm{Roentgen Rays} Alternate name; see \textit{X-Ray}.

\synonyms
Roentgen Radiation; Electromagnetic X-Radiation

\medterm{Rofecoxib} Rofecoxib is a selective COX-2 anti-inflammatory drug withdrawn in many countries because of increased cardiovascular risk.

\medterm{Romano-Ward Syndrome} The most common congenital long QT syndrome, inherited in an autosomal dominant pattern without associated neurosensory deafness (due to mutations in cardiac potassium or sodium channel genes \textit{KCNQ1}, \textit{KCNH2}, or \textit{SCN5A}), predisposing to syncope, polymorphous ventricular tachycardia (torsades de pointes), and sudden cardiac death.

\medterm{Romanus Lesion} An early inflammatory radiographic finding in ankylosing spondylitis characterized by enthesitis with focal erosion and surrounding reactive sclerosis at the anterior superior and inferior corners of vertebral bodies ('shiny corners'), preceding syndesmophyte formation.
\synonyms
Shiny Corner Sign; Spondylitis Anterior Enthesitis

\medterm{Rombergism} Unsteadiness or falling when standing with feet together and eyes closed, indicating impaired proprioception or vestibular function.

\medterm{Romberg Sign} Loss of balance when standing with the feet together and the eyes closed, suggesting impaired proprioception or vestibular function.

\medterm{Romberg's Sign} Increased unsteadiness when a person stands with the feet together and closes the eyes, suggesting impaired
proprioception or vestibular function.

\medterm{Romberg Test} The Romberg test is a clinical neurological test used to assess the integrity of the
dorsal column-medial lemniscus pathway (proprioception) and the cerebellum in
maintaining balance. The test is performed by asking the patient to stand with feet
together and eyes open, and then to close the eyes.

\medterm{Romidepsin} A histone deacetylase inhibitor used for selected T-cell lymphomas; it can cause cytopenias, infection risk, and ECG-related effects.

\medterm{Romiplostim} A thrombopoietin-receptor agonist that stimulates platelet production and is used mainly for chronic immune thrombocytopenia.

\medterm{Rongeur} A strong, jawed surgical instrument used to bite and remove small pieces of bone or other hard tissue.

\medterm{Ronnel} Ronnel is an organophosphate insecticide that can inhibit acetylcholinesterase and cause cholinergic toxicity after exposure.

\medterm{Root Canal} The channel inside a tooth root containing the pulp and its nerves and blood vessels.

\medterm{Root Canal Therapy} Dental treatment in which infected or inflamed pulp is removed from a tooth, the canal is cleaned and sealed, and the tooth is restored.

\medterm{Root Canal Treatment} Alternate term; see \textit{Root Canal Therapy}.

\medterm{Root Caries} Tooth decay affecting exposed root surfaces after gum recession or loss of protective covering. It is more common with aging, dry mouth, or limited cleaning and may be slowed with fluoride and dental care.

\medterm{Root Filling} Dental treatment sealing a cleaned tooth root canal to prevent bacteria entering again and causing reinfection.

\medterm{Root Fracture} A crack or break in the root of a tooth. A horizontal fracture may sometimes be stabilized and monitored, while a vertical fracture often allows bacteria to enter and may require removal of the tooth. Treatment depends on the fracture’s location and extent.

\medterm{Rooting Reflex} A primitive reflex present from birth to 3-4 months, triggered by stroking the infant's
cheek or corner of the mouth. The infant turns toward the stimulus and opens the mouth,
searching for the nipple. Facilitates breastfeeding. Absence may indicate neurologic
impairment. Disappears as voluntary feeding develops.

\medterm{Ropinirole} A medicine that stimulates dopamine receptors and is used for Parkinson disease and restless legs syndrome. It may cause nausea, dizziness, sleepiness, low blood pressure, hallucinations, swelling, or impulse-control problems such as gambling.

\medterm{Ropivacaine} A long-acting local anesthetic used for regional nerve blocks and epidural anesthesia; excessive blood levels can cause neurologic or cardiac toxicity.

\medterm{Ropivacaine Hydrochloride} The hydrochloride salt of ropivacaine, a long-acting local anesthetic used for regional and epidural anesthesia.

\medterm{Rorschach Test} A projective psychological test using inkblots; its diagnostic validity is limited and interpretation requires specialist training.

\medterm{Rosacea} A chronic inflammatory skin disorder that most often affects the central face. It may cause recurrent flushing, persistent redness, visible small blood vessels, burning or stinging, and acne-like bumps or pustules without being acne. Eye involvement may cause irritation or inflammation, and long-standing disease may produce thickening of the skin, especially of the nose (rhinophyma). Rosacea is not contagious, and its severity and features vary over time.

\medterm{Rosai-Dorfman Disease} A rare, benign, non-Langerhans cell histiocytic proliferative disorder (sinus histiocytosis with massive lymphadenopathy) characterized clinically by massive painless bilateral cervical lymphadenopathy and fever, and pathologically by sinusoidal infiltration of S100-positive histiocytes demonstrating emperipolesis.
\synonyms
Sinus Histiocytosis with Massive Lymphadenopathy; SHML

\medterm{Rose Bengal} Rose bengal is a dye used in eye-surface testing and laboratory or diagnostic procedures; it highlights damaged or poorly protected conjunctival cells.

\medterm{Roseola} Roseola infantum is a common viral illness, usually caused by human herpesvirus 6, with several days of high fever followed by a pink trunk rash as the fever resolves; febrile seizures may occur.

\synonyms
Exanthema Subitum; Sixth Disease; Roseola Infantum

\medterm{Roseola Infantum} A common childhood illness, usually caused by human herpesvirus 6 and sometimes HHV-7. Several days of high fever are followed by a pink or red trunk rash as the fever resolves; febrile seizures may occur, and illness is usually self-limited.

\medterm{Roseolovirus} A genus of herpesviruses that includes human herpesvirus 6 and 7, commonly associated with roseola and sometimes serious disease after reactivation.

\medterm{Rosiglitazone} A medicine for type 2 diabetes that helps the body respond better to insulin and lower blood sugar. It can cause weight gain, swelling, fractures, or worsening heart failure and is used selectively.

\medterm{Rosiglitazone Maleate} The maleate salt of rosiglitazone, a thiazolidinedione that improves insulin sensitivity in type 2 diabetes but may cause weight gain, edema, fractures, and heart-failure worsening.

\medterm{Ross Procedure} An operation for severe aortic-valve disease, especially in some children and young adults. The patient’s own pulmonary valve is moved to the aortic position, and a donor valve or tube replaces it. The moved valve can grow in children and usually avoids lifelong blood-thinning medicine, but either valve may later narrow or leak.

\synonyms
Pulmonary Autograft Procedure; Ross Operation

\medterm{Ross River Virus} A mosquito-borne alphavirus that causes Ross River fever, with fever, rash, and often prolonged joint pain or swelling.

\medterm{Rostral} Positioned toward the nose or front of the head. In descriptions of the brain, it usually means toward the forehead or front rather than toward the back of the head.

\medterm{Rosuvastatin Calcium} Rosuvastatin calcium is a statin medicine that lowers LDL cholesterol and reduces the risk of cardiovascular events.

\medterm{Rotation} Movement around an axis; in anatomy it describes turning a body part inward, outward, or around its longitudinal axis.

\medterm{Rotator Cuff} The group of four shoulder muscles and tendons that stabilize the glenohumeral joint and rotate the arm.

\synonyms
Musculotendinous Cuff; SITS Muscles

\medterm{Rotator Cuff Arthropathy} Shoulder-joint damage that develops after long-standing, irreparable tears of the rotator-cuff tendons. It can cause pain, weakness, grinding, and inability to raise the arm; severe cases may need specialized shoulder replacement.

\medterm{Rotator Cuff Exercises} Rotator cuff exercises strengthen and stretch shoulder muscles and are selected according to the injury, pain, range of motion, and rehabilitation plan.

\medterm{Rotator Cuff Injury} Damage to one or more shoulder tendons or muscles that stabilize and move the arm. It may cause pain, weakness, night discomfort, or limited movement after overuse, aging, or sudden injury.

\medterm{Rotator Cuff Problems} Rotator cuff problems include tendinopathy, bursitis, and tears that cause shoulder pain, weakness, painful motion, and difficulty with overhead activity.

\medterm{Rotator Cuff Repair} Surgery to reattach a torn shoulder tendon to the upper-arm bone, usually with small instruments or an open approach. Recovery requires structured rehabilitation and may take months; stiffness, weakness, recurrent tearing, or infection can occur.

\medterm{Rotator Cuff Tear} A partial or complete tear of one or more tendons stabilizing the shoulder, causing pain, weakness, and difficulty raising or rotating the arm.

\medterm{Rotator Cuff Tendinitis} Alternate name; see \textit{Rotator cuff syndrome}.

\medterm{Rotavirus} A very contagious virus that commonly causes vomiting and watery diarrhea in infants and young children. It can cause dangerous dehydration; vaccination greatly reduces severe illness, and testing is not always needed.

\medterm{Rotavirus Antigen Test} A rotavirus antigen test detects rotavirus proteins in stool and can help diagnose viral gastroenteritis, especially in infants and young children.

\medterm{Rotaviruses} A group of viruses causing vomiting and watery diarrhea, especially in young children; vaccination lowers severe disease risk.

\medterm{Rotavirus Infection} A highly contagious viral gastroenteritis causing vomiting and watery diarrhea, especially in infants and young children; vaccination prevents most severe disease.

\medterm{Rotavirus Vaccine} A live-attenuated oral viral vaccine administered to young infants in a 2- or 3-dose series to prevent severe dehydrating gastroenteritis and diarrhea caused by rotavirus infection, carrying a slight, transient risk of intussusception.

\medterm{Rotenone} Rotenone is a naturally occurring pesticide that inhibits mitochondrial complex I and can be toxic to people and animals after exposure.

\medterm{Rothmund-Thomson Syndrome} A rare inherited condition that causes a mottled, red, and thin skin rash beginning in infancy, short stature, sparse hair, skeletal or dental problems, and sometimes cataracts. It also increases the risk of certain cancers.

\medterm{Roth Spots} Small areas of bleeding in the retina with pale centers, seen during an eye examination. They may occur with infection of the heart valves, severe anemia, leukemia, diabetes, or other illnesses and are not specific to one diagnosis.

\medterm{Roth's Spots} Small retinal haemorrhages with pale centres, associated with conditions such as infective endocarditis, severe anaemia, leukaemia, and diabetes.

\medterm{Rotor Syndrome} A benign, autosomal recessive inherited hyperbilirubinemia caused by mutations in SLCO1B1 and SLCO1B3 genes, resulting in defective hepatic reuptake and storage of conjugated bilirubin, presenting with lifelong, asymptomatic conjugated hyperbilirubinemia with normal liver histology and elevated urinary coproporphyrin.
\synonyms
Rotor's Syndrome; Familial Conjugated Hyperbilirubinemia

\medterm{Rotterdam Criteria} A commonly used approach for diagnosing polycystic ovary syndrome in adults. After other causes are excluded, diagnosis generally requires at least two of: irregular or absent ovulation, signs or blood-test evidence of high androgen levels, or polycystic-appearing ovaries on ultrasound. Ultrasound alone does not establish the diagnosis.

\medterm{Rotter Nodes} Small interpectoral lymph nodes situated between the pectoralis major and pectoralis minor muscles that provide an alternate pathway for breast lymphatic drainage directly to Level II and III axillary lymph nodes.

\synonyms
Interpectoral Lymph Nodes; Rotter's Nodes

\medterm{Rouleau} A stack or column of red blood cells, commonly seen in blood with increased plasma
proteins.

\medterm{Rouleaux} A pattern where red blood cells stack like coins, often because blood plasma contains extra proteins.

\medterm{Rouleaux Formation} Linear, coin-like stacks or rolls of red blood cells visualized on a peripheral blood smear, occurring when elevated plasma concentrations of asymmetric acute-phase proteins (fibrinogen or monoclonal immunoglobulins) neutralize the normal negative repulsive zeta potential of erythrocytes, characteristically seen in multiple myeloma and severe inflammatory states.

\medterm{Round} Round means circular or approximately spherical; in anatomy or pathology it describes shape and must be interpreted with the structure being described.

\medterm{Round Back} Alternate name; see \textit{Kyphosis}.

\medterm{Round Ligament} A round ligament is a cord-like supporting structure, including the uterine ligaments that pass toward the groin and help maintain uterine position.

\medterm{Round Ligament of the Uterus} A fibrous and muscular cord originating from the superior-lateral aspect of the uterine cornu anterior to the fallopian tube, traversing through the inguinal canal to terminate in the subcutaneous tissue of the labium majus, functioning to hold the uterus in an anteverted and anteflexed position.

\medterm{Round Ligament Pain} Sharp or aching pain in the lower abdomen or groin caused by stretching of the round ligaments as the uterus enlarges during pregnancy. Persistent, severe, or bleeding-associated pain requires medical evaluation.

\medterm{Round Pneumonia} A lung infection that appears as a rounded shadow on a chest X-ray rather than the usual patchy pattern. It occurs especially in children and can resemble a tumor, so follow-up imaging may be needed after treatment.

\medterm{Roundworm} A parasitic worm, often referring to Ascaris lumbricoides, acquired by swallowing eggs from contaminated soil, food, or water. It may cause abdominal symptoms, poor nutrition, or intestinal blockage when many worms are present.

\medterm{Roundworms} Nematode parasites that may infect intestines or tissues, causing disease depending on their species and number.

\medterm{Roussy-Levy Syndrome} An autosomal dominant hereditary motor and sensory neuropathy (a phenotypic variant of Charcot-Marie-Tooth disease type 1A) characterized by early-onset sensory ataxia, prominent postural hand tremor, distal muscle weakness and wasting, pes cavus, and absent deep tendon reflexes.
\synonyms
Hereditary Areflexic Dystasia; CMT-Tremor Variant

\medterm{Route Of Administration} The way a medicine enters the body, such as by mouth, under the tongue, injection, inhalation, skin application, or through a vein. The route affects how quickly and how much of the medicine reaches the bloodstream.

\medterm{Roux-en-Y Gastric Bypass} A bariatric and metabolic surgical procedure in which a small (15 to 30 mL) upper gastric pouch is created and anastomosed to a Roux limb of proximal jejunum (gastrojejunostomy), bypassing the remaining stomach, duodenum, and proximal jejunum, connected to a downstream biliopancreatic limb (jejunojejunostomy).

\synonyms
RYGB

\medterm{Roux Stasis Syndrome} A chronic functional disorder of gastric and small bowel motility developing in up to 30 percent of patients following a Roux-en-Y gastrojejunostomy or gastrectomy. Characterized by delayed gastric emptying, severe postprandial fullness, nausea, vomiting of non-bilious food, and chronic epigastric pain, resulting from the loss of normal duodenal pacemaking coordination in the transected Roux jejunal limb.

\synonyms
Roux-en-Y Stasis Syndrome; Chronic Roux Limb Stasis

\medterm{Rovsing Sign} Pain in the lower right abdomen that occurs when the lower left abdomen is pressed. This examination finding can support suspicion of appendicitis, but it is not diagnostic by itself and must be interpreted with the symptoms and other tests.

\medterm{Roxithromycin} A macrolide antibiotic used in some countries for selected bacterial infections. It stops bacteria making proteins, but resistance, medicine interactions, stomach upset, liver problems, and heart-rhythm changes can occur.

\medterm{RPR Test} A blood screening test for syphilis that detects antibodies produced during tissue injury from the infection. A reactive result does not prove active syphilis and requires a second, more specific test and clinical assessment.

\medterm{RS3PE Syndrome} Remitting seronegative symmetrical synovitis with pitting edema, an inflammatory syndrome causing sudden swelling of the hands or feet, joint pain, and stiffness, usually in older adults. It may occur with another rheumatic disease, infection, or cancer.
\synonyms
Remitting Seronegative Symmetrical Synovitis with Pitting Edema

\medterm{RSI} Repetitive strain injury, a group of pain and dysfunction syndromes associated with repeated movement or sustained posture.

\medterm{RSV} A common respiratory virus that can cause mild cold symptoms or serious bronchiolitis and pneumonia, especially in infants and older adults.

\medterm{RSV Antibody Test} An RSV antibody test detects antibodies to respiratory syncytial virus and may help assess past exposure or immune response, but is not usually the main test for acute infection.

\medterm{RSV Infection} Alternate name; see \textit{Respiratory Syncytial Virus Infection}.

\medterm{RT-PCR} A laboratory test that converts RNA into DNA and then makes many copies of a selected genetic sequence so it can be detected. It can identify RNA viruses and measure selected genes in clinical samples.

\medterm{Rubber} Rubber is an elastic material used in gloves, devices, and equipment; natural-rubber latex can cause irritant or allergic reactions.

\medterm{Rubber Band Ligation} The most common and effective non-surgical outpatient procedure performed to treat symptomatic internal hemorrhoids. Through a small lighted viewing scope (an anoscope) inserted into the anal canal, a clinician uses a suction or mechanical ligator instrument to grasp the redundant hemorrhoidal tissue above the pain-sensitive dentate line and snap a tiny, tight elastic rubber band firmly around its base. The tight band strangles the hemorrhoid's blood supply; within 7 to 10 days, the ischemic hemorrhoidal tissue withers, dies, and painlessly detaches, sloughing off and passing out unnoticed during a bowel movement. The remaining small scar fixes the surrounding mucosa to the underlying tissue, preventing further prolapse or bleeding without requiring surgical incisions, general anesthesia, or hospital recovery.

\synonyms
Hemorrhoidal Banding; RBL; Band Ligation of Hemorrhoids; Elastic Band Ligation

\medterm{Rubber Cement Poisoning} Poisoning from swallowing or inhaling rubber cement, which contains volatile solvents. It may cause dizziness, drowsiness, confusion, breathing problems, or chemical pneumonitis if it enters the lungs.

\medterm{Rubber Dams} Sheets used in dentistry to isolate teeth from saliva and the rest of the mouth during procedures.

\medterm{Rubella} A contagious viral infection causing mild fever, swollen lymph nodes, and a fine rash. Infection during pregnancy can cause miscarriage, fetal growth problems, hearing loss, eye disease, congenital heart disease, or other congenital rubella syndrome; vaccination prevents most cases.

\synonyms
German Measles; Three-Day Measles

\medterm{Rubella Virus} The virus that causes rubella, a contagious illness with fever, swollen glands, and rash. Infection during pregnancy can seriously harm the developing baby, causing hearing loss, eye disease, heart problems, or other birth defects.

\medterm{Rubeola} Measles, a highly contagious viral illness caused by measles virus and characterized by fever, cough, conjunctivitis, and a widespread maculopapular rash.

\synonyms
Measles; Hard Measles; Morbilli

\medterm{Rubeosis Iridis} Growth of abnormal new blood vessels on the colored part of the eye. It usually results from poor blood flow in the retina, often caused by diabetes or a blocked retinal vein, and can lead to painful glaucoma and vision loss.

\medterm{Rubidium} A chemical element whose radioactive isotopes can be used in medical imaging; nonmedical rubidium compounds can affect potassium-related physiology.

\medterm{Rubin Maneuver} An emergency childbirth technique for shoulder dystocia, when a baby’s shoulder is stuck behind the birthing person’s pelvic bone. A clinician applies pressure and rotates the baby’s shoulder to reduce its width and help complete delivery safely.

\medterm{Rubinstein-Taybi Syndrome} A genetic developmental disorder associated with characteristic facial features,
short stature, broad thumbs and toes, and intellectual disability. It is commonly caused by pathogenic variants in
CREBBP or EP300.

\medterm{Rubivirus} A genus of viruses that includes rubella virus, which causes rubella and congenital rubella syndrome.

\medterm{Rubredoxins} Small proteins containing iron that help move electrons in some bacteria and other organisms. They take part in chemical reactions that release energy and protect cells from harmful oxygen-related chemicals.

\medterm{Rubricyte Count} A measurement of rubricytes, an older name for immature red blood cells, in blood or bone marrow. It may help assess how actively the marrow is making red cells, but modern tests usually report reticulocytes or other marrow findings.

\medterm{Rule Of Nines} A quick method for estimating the percentage of body surface burned by assigning simple proportions to body regions. The estimate helps guide fluids and emergency care, but children and unusual body shapes require more accurate methods.

\medterm{Rule Out} Rule out means evaluate a possibility and determine whether it is present; it does not mean the condition has already been excluded.

\medterm{Rumination} Repeated regurgitation and rechewing of food after eating, without the usual vomiting process. It is a behavior or symptom that may occur in rumination disorder or with other gastrointestinal, developmental, or behavioral conditions.

\medterm{Rumination Disorder} A feeding or eating disorder involving repeated, usually effortless return of recently eaten food to the mouth, followed by rechewing, reswallowing, or spitting. It differs from ordinary vomiting and causes clinically significant distress, impairment, or nutritional or medical consequences.

\medterm{Rumination Syndrome} Alternate name; see \textit{Rumination Disorder}.

\medterm{Ruminococcus Gnavus} An intestinal bacterium involved in mucin and carbohydrate metabolism and studied in relation to gut inflammation.

\medterm{Rumpel-Leede Sign} Small pinpoint bleeding spots that appear below a blood-pressure cuff after it has briefly been inflated. The sign may suggest fragile small blood vessels or a low platelet count, but it is not specific and is rarely used alone today.

\synonyms
Tourniquet Test; Capillary Fragility Test

\medterm{Rumpel-Leede Test} A clinical diagnostic test evaluating capillary wall fragility and thrombocytopenia in which a sphygmomanometer cuff is inflated around the arm midway between systolic and diastolic pressure for five minutes; the appearance of more than 10 to 20 petechiae in a designated circular area distal to the cuff indicates increased capillary fragility (as seen in dengue fever, thrombocytopenia, or scurvy).

\synonyms
Capillary Fragility Test; Tourniquet Test

\medterm{Runner's Knee} Alternate name; see \textit{Patellofemoral pain syndrome}.

\medterm{Runny Nose} Excess fluid or mucus draining from the nose, commonly caused by infection, allergy, irritants, or sinus disease.

\medterm{Rupture} A tear or complete break in an organ, vessel, muscle, or other tissue, often causing bleeding or loss of function.

\medterm{Ruptured Abdominal Aortic Aneurysm} A burst in a weakened section of the main artery in the abdomen, causing severe internal bleeding. Sudden abdominal or back pain, fainting, sweating, or collapse may occur.

\medterm{Ruptured Disk} A spinal disk whose outer layer has torn, allowing inner material to bulge outward.
\synonyms
Herniated Disc; Slipped Disc; Herniated Nucleus Pulposus

\medterm{Ruptured Eardrum} A ruptured eardrum is a tear in the tympanic membrane caused by infection, pressure, trauma, or foreign objects and may cause pain, drainage, hearing loss, or dizziness.

\medterm{Ruptured Perforated Eardrum} A perforated eardrum is a tear in the tympanic membrane from infection, pressure change, trauma, or an inserted object. Pain, drainage, hearing loss, or ringing may occur; keeping the ear dry and avoiding unprescribed drops helps while healing or repair is assessed.

\medterm{Russell-Silver Syndrome} A condition present from birth that causes poor growth before and after birth, a relatively large head, feeding difficulty, uneven growth of the limbs or body, and characteristic facial features. It is usually related to changes in gene regulation.

\medterm{Russell Viper Venom Time} A blood-clotting test used as part of testing for lupus anticoagulants. The venom activates clotting in the test tube; a prolonged result that shortens when extra phospholipid is added supports the presence of a lupus anticoagulant but does not establish the diagnosis by itself.

\synonyms
dRVVT

\medterm{Rx} A common abbreviation for prescription or prescribed treatment. It comes from the Latin word meaning “take” and is used on medication orders and in medical writing.

\medterm{Ryle's Tube} A soft tube passed through the nose into the stomach to remove stomach contents or give fluids, nutrition, or medicines. Position must be checked before use because a misplaced tube can enter the airway and cause serious harm.



\medterm{S Phase} The stage of the cell cycle when a cell copies all of its DNA before dividing. This duplication allows each new cell to receive a complete set of genetic instructions.

\medterm{S-Adenosylmethionine} A substance made naturally in the body that helps chemical reactions involving cell membranes, genetic material, and brain chemicals. It is also sold as a supplement, but its benefits and interactions vary by condition and require medical advice.

\medterm{S3 Gallop Pathology} An extra heart sound heard just after the normal second heart sound as the lower chambers fill rapidly. It can be normal in children and pregnancy but may suggest excess fluid or heart failure in an older adult.

\medterm{S4 Gallop Pathology} An extra heart sound heard just before the normal first heart sound when an upper chamber pushes blood into a stiff lower chamber. It may occur with long-standing high blood pressure, a narrowed aortic valve, or thickened heart muscle.

\medterm{SA Node} A small group of specialized cells in the upper right heart chamber that normally starts each heartbeat. It produces regular electrical signals that spread through the heart and coordinate its rhythm.

\medterm{SAA} Usually an abbreviation for serum amyloid A, a blood protein that rises during inflammation. Persistently high levels can contribute to abnormal protein deposits in organs during some long-lasting inflammatory diseases.


\medterm{Sabin Vaccine} The oral polio vaccine containing live attenuated poliovirus strains. It produces strong intestinal immunity but, rarely, vaccine-derived poliovirus can occur, so vaccine policy depends on local public-health guidance.

\medterm{Saccade} A quick, coordinated eye movement that shifts the eyes from one target to another. Abnormal saccades can occur with disorders affecting the eye muscles, balance system, or brain and nerves.


\medterm{Saccular Aneurysm} A rounded pouch bulging from one side of an artery because that part of the vessel wall is weak. It may enlarge, form clots, or burst; risk depends on its location, size, shape, and growth.

\medterm{Sachet Poisoning} Illness caused by swallowing, inhaling, or getting the contents of a sachet on the skin or in the eyes. Effects depend on the substance and amount; some sachets contain medicines or dangerous chemicals.

\medterm{Sacral} Sacral means relating to the sacrum, the triangular bone at the base of the spine that connects with the pelvis.

\medterm{Sacral Dimple} A small indentation in the skin over the lower spine that is present at birth. Most are harmless, but a deep, large, high, hairy, discolored, or draining dimple may need imaging for an underlying spinal abnormality.
\synonyms
Pilonidal Dimple

\medterm{Sacral Fistula} An abnormal channel near the lower spine that connects the skin with deeper tissue or another body space. It may cause persistent drainage, pain, repeated infection, or a skin opening and can result from birth differences, injury, or pilonidal disease.

\medterm{Sacral Nerve Stimulation} Treatment that sends mild electrical pulses to nerves near the lower spine to improve bladder or bowel control. It may help selected people with persistent urinary or fecal leakage, retention, or constipation after other treatments have not worked.

\medterm{Sacral Nerve Stimulator} An implanted device that sends mild electrical signals to nerves near the lower spine to improve bladder or bowel control in selected people. It may help incontinence or some emptying problems when other treatments have not worked.

\medterm{Sacral Nerves} The spinal nerves arising from the sacral spinal cord and supplying the pelvis, pelvic floor, bladder, bowel, and lower limbs.
\medterm{Sacral Region} The lower back and posterior pelvic region overlying the sacrum, a fused triangular bone that connects the lumbar spine with the pelvis.

\medterm{Sacral Vertebra} One of the fused vertebrae forming the sacrum, which supports the spine, joins the pelvic bones, and transmits trunk weight to the lower limbs.

\medterm{Sacral Vertebrae} Five spinal bones that fuse to form the sacrum at the back of the pelvis. They transfer weight from the spine to the hips and contain openings through which sacral nerves pass.

\medterm{Sacral Vertebrae (S1-S5)} The five sacral vertebrae that fuse to form the sacrum, connect the spine to the pelvis, and provide foramina for sacral nerve roots.
\medterm{Sacrococcygeal} Sacrococcygeal means relating to the sacrum and coccyx at the lower end of the vertebral column.

\medterm{Sacrococcygeal Teratoma} A tumor present from birth at the base of the spine and made of several tissue types. It may grow outside or inside the pelvis; a large fetal tumor can strain the heart and increase amniotic fluid.

\medterm{Sacrocolpopexy} An operation that supports a prolapsed top of the vagina or uterus by attaching it to the strong tissue over the sacrum, usually with surgical mesh. It may be performed through open or small incisions; bleeding, infection, mesh problems, bowel or bladder injury, and recurrent prolapse are possible.

\medterm{Sacroiliac Joint} The joint connecting the triangular sacrum at the base of the spine to each hip bone. It transfers weight between the spine and legs and may become painful or inflamed after injury, pregnancy, arthritis, or infection.

\medterm{Sacroiliac Joint Dysfunction (Sacroiliac Joint Pain)} Alternate index label for Sacroiliac Joint Pain.

\medterm{Sacroiliac Joint Dysfunction (SI Joint Pain)} Pain or abnormal movement from the joint between the lower spine and pelvis. It may cause pain in the lower back, buttock, or thigh and can follow injury, pregnancy, arthritis, or changes in walking.

\medterm{Sacroiliac Joint Pain} Sacroiliac-joint pain arises near the junction of the lower spine and pelvis and may follow injury, pregnancy, arthritis, leg-length or gait changes, or inflammation. Pain can radiate to the buttock or thigh; examination distinguishes it from hip or spine disease.

\synonyms
Sacroiliac Joint Dysfunction (Sacroiliac Joint Pain)


\medterm{Sacroiliitis} Inflammation of one or both joints connecting the lower spine to the pelvis. It can cause pain in the lower back, buttock, or back of the thigh and may result from inflammatory arthritis, infection, injury, or wear.

\medterm{Sacrospinous Ligament Fixation} An operation through the vagina that attaches the top of the vagina or cervix to a strong pelvic ligament to treat prolapse. It avoids an abdominal incision but may cause buttock pain, bleeding, infection, nerve injury, or recurrent prolapse.

\synonyms
Sacrospinous Colpopexy; Sacrospinous Suspension

\medterm{Sacrum} A large triangular bone at the back of the pelvis formed by five fused spinal bones. It connects the spine to the hip bones, transfers body weight to the legs, and contains openings for nerves supplying the pelvis and lower limbs.

\medterm{SAD} An abbreviation for seasonal affective disorder, a recurring pattern of depression or mood change linked to a particular season, often winter and reduced daylight. Treatment may include therapy, medicines, or supervised light therapy.
\synonyms
Seasonal Affective Disorder (SAD)

\medterm{Sad Mood} Sad mood is a feeling of unhappiness or low mood that may be a normal response or a symptom of depression, grief, illness, or another condition.

\medterm{Saddle Pulmonary Embolism} A large blood clot lodged where the main artery from the heart divides toward both lungs. It can severely block blood flow, causing sudden breathlessness, chest pain, fainting, low blood pressure, or cardiac arrest and requires emergency treatment.

\synonyms
Saddle Embolus; Bifurcation Pulmonary Embolism

\medterm{Sadism} Sexual arousal or gratification from inflicting pain, humiliation, or suffering; clinical significance depends on distress, impairment, or nonconsent.


\medterm{Safe Sex} Sexual practices that reduce the risk of sexually transmitted infections and unintended pregnancy, including condoms or other barriers, testing, vaccination, and informed consent.


\medterm{Safety Of Drugs} The process of balancing a medicine’s expected benefits against side effects, interactions, and reasons it may be unsafe. Safe use includes checking the person, dose, allergies, other medicines, monitoring needs, and instructions.
\medterm{Safety Plan} A collaboratively written, practical plan for responding to escalating suicidal thoughts
or emotional crisis. It identifies warning signs, internal coping strategies, supportive
contacts, professional resources, and methods to restrict access to lethal means.

\medterm{Safety Study} A safety study evaluates whether a medicine, device, procedure, or treatment causes harmful effects and identifies risks before or during clinical use.



\medterm{Safrole} Safrole is a natural aromatic compound found in some plants and oils; high or prolonged exposure is toxic and has carcinogenic concern.

\medterm{Sagittal} Oriented in a vertical direction that divides the body or a structure into right and left parts. A midsagittal plane divides the parts equally; a parasagittal plane does not.

\medterm{Sagittal Plane} An imaginary vertical plane dividing the body or an organ into right and left portions. A midsagittal plane divides them equally; a parasagittal plane divides them unequally.

\medterm{Salaam Attacks} Brief repeated bending spasms or sudden head drops in infants. They may be seizures and require urgent specialist assessment because early treatment can protect development and reduce the risk of continuing epilepsy.


\medterm{Salbutamol (Albuterol)} A quick-relief inhaled medicine that relaxes tightened airway muscles in asthma and chronic obstructive lung disease. It eases wheezing and breathlessness but may cause shaking, a fast heartbeat, or palpitations and does not replace treatment that controls airway inflammation.

\synonyms
Albuterol

\medterm{Salicylate} A group of chemicals that includes aspirin and can reduce pain, fever, or inflammation. Excessive amounts may cause poisoning with vomiting, ringing ears, fast breathing, confusion, dehydration, or dangerous blood chemistry changes.

\medterm{Salicylate Toxicity} Poisoning caused by too much aspirin or another salicylate. It may cause vomiting, ringing in the ears, fever, sweating, fast or deep breathing, dehydration, confusion, seizures, or coma and requires urgent medical advice.

\medterm{Salicylic Acid} A medicine applied to the skin to soften and remove excess surface skin. It is used for acne, psoriasis, warts, corns, and similar thickened or blocked skin; irritation is common and large exposures can be toxic.

\medterm{Saline} A sterile solution of sodium chloride in water. Different concentrations are used for intravenous fluids, wound or eye irrigation, inhalation, laboratory work, or dilution of medicines.
\synonyms
Normal saline

\medterm{Saline Infusion Sonohysterography} An ultrasound test in which sterile salt water is placed inside the uterus to outline its cavity. It can show polyps, fibroids, scar tissue, or an unusually shaped cavity and may cause temporary cramping or light bleeding.

\medterm{Saline Nasal Washes} Rinsing the nasal passages with sterile, distilled, or appropriately treated saline to loosen mucus and reduce nasal irritants; unsafe water can cause serious infection.

\medterm{Saliva} Fluid made by the salivary glands that moistens food, begins digestion, protects the mouth and teeth, and helps swallowing and taste. Too little saliva can cause dry mouth, mouth soreness, and tooth decay.

\medterm{Salivary Duct Stones} Calcified deposits obstructing a salivary duct, most often the submandibular duct, causing recurrent painful swelling of a gland during meals and sometimes infection.

\medterm{Salivary Gland} A gland that makes saliva and releases it into the mouth. The three major paired glands are the parotid, submandibular, and sublingual glands; many smaller glands lie in the mouth lining.

\medterm{Salivary Gland Biopsy} Removal of salivary-gland tissue or a needle sample for microscopic examination of suspected autoimmune disease, inflammation, infection, or neoplasm.

\medterm{Salivary Gland Cancer} Cancer arising in a salivary gland, most often the parotid near the ear. It may cause a persistent lump, pain, numbness, dry mouth, or facial weakness; imaging and biopsy identify the type and guide treatment.

\medterm{Salivary Gland Infection} Infection of a salivary gland, usually bacterial and often involving the parotid or
submandibular gland. Painful swelling, fever, dry mouth, and purulent drainage from the
duct may occur.

\medterm{Salivary Gland Infections} Infections of a salivary gland or its duct, commonly bacterial or viral, causing painful swelling, tenderness, fever, dry mouth, or purulent drainage.

\medterm{Salivary Gland Neoplasm} An abnormal growth in a salivary gland that may be noncancerous or cancerous. It can cause a lump, pain, facial weakness, dry mouth, or swallowing problems and may require imaging and biopsy.

\medterm{Salivary Gland Scintigraphy} An imaging test that measures salivary-gland function after a small radioactive tracer is given. It may help assess dryness from Sjögren syndrome, blocked ducts, inflammation, radiation injury, or other gland problems.

\medterm{Salivary Gland Tumors} Tumors arising in the salivary glands. Many are benign, especially in the parotid gland;
pleomorphic adenoma is a common benign type. Malignant types include mucoepidermoid,
adenoid cystic, and acinic cell carcinoma.

\synonyms
Tumor, Salivary Gland

\medterm{Salivary Glands} Glands that make saliva, including the parotid, submandibular, and sublingual glands. Saliva moistens food, begins digestion, protects teeth, helps swallowing, and supports taste and mouth health.
\medterm{Salivation} Production and release of saliva by the salivary glands. Saliva moistens food, begins digestion, protects teeth, and helps swallowing; excessive or reduced salivation can have many causes.

\medterm{Salk Vaccine} The inactivated poliovirus vaccine, given by injection to protect against poliomyelitis. It contains killed virus and cannot cause polio; it differs from the older oral live-attenuated vaccine.
\medterm{Salmon Patches} Flat pink or red capillary birthmarks, commonly on the eyelids, forehead, or nape; many fade during infancy.
\medterm{Salmonella} A genus of Gram-negative, facultatively anaerobic, rod-shaped bacteria in the Enterobacteriaceae family.
It includes non-typhoidal strains that cause gastroenteritis and typhoidal strains that can
cause enteric fever.


\medterm{Salmonella Choleraesuis} A Salmonella serotype adapted to pigs that can cause severe, invasive infection in humans, including bloodstream infection and pneumonia.


\medterm{Salmonella Enterica} A major Salmonella species causing human disease. It can produce diarrhea and fever, enteric fever, or bloodstream infection after contaminated food, water, animals, or infected people spread the bacteria.
\medterm{Salmonella Enterocolitis} Salmonella enterocolitis is intestinal infection causing diarrhea, fever, abdominal cramps, and sometimes bacteremia after contaminated food, water, or animal exposure.

\medterm{Salmonella Food Poisoning (Salmonellosis)} Alternate index label for Salmonellosis.


\medterm{Salmonella Infection} An infection caused by Salmonella bacteria, usually acquired through contaminated food, water, animals, or hands. It commonly causes diarrhea, fever, cramps, and vomiting but can spread into blood or other organs.

\medterm{Salmonella Infections} Infections caused by Salmonella bacteria, ranging from gastroenteritis to enteric fever or invasive bloodstream disease.

\medterm{Salmonella Typhi} The bacterium that causes typhoid fever, usually acquired through contaminated food or water and capable of spreading through the bloodstream.

\medterm{Salmonella Typhimurium} A common Salmonella serotype that usually causes foodborne gastroenteritis but can cause invasive disease in vulnerable people.

\medterm{Salmonella Typhisuis} A Salmonella serotype mainly associated with swine; human infection is uncommon but may cause enteric or invasive disease.

\medterm{Salmonella Vaccines} Vaccines used mainly to prevent typhoid fever, including typhoid conjugate, Vi polysaccharide, and oral live-attenuated vaccines; routine vaccination does not prevent all Salmonella gastroenteritis.

\medterm{Salmonellosis} Infection with Salmonella bacteria (non-typhoidal), typically acquired from contaminated
food (eggs, poultry, meat). Presents with acute gastroenteritis: diarrhea, fever, and
abdominal cramps within 6-72 hours of ingestion. Self-limited in immunocompetent
patients. Bacteremia can occur, especially in immunocompromised.

\synonyms
Salmonella Food Poisoning (Salmonellosis)

\medterm{Salpingectomy} Surgical removal of a fallopian tube. It may be used to treat a tubal ectopic pregnancy,
hydrosalpinx, or other tubal disease. The approach and whether the tube is removed depend
on the clinical condition, fertility plans, and the state of the opposite tube.

\medterm{Salpingitis} Inflammation or infection of one or both fallopian tubes, usually as part of pelvic inflammatory disease spreading upward from the vagina or cervix. It may cause lower-abdominal pain, fever, unusual discharge, infertility, or pregnancy outside the uterus.

\medterm{Salpingitis Isthmica Nodosa} Small thickened areas and side pockets in the narrow part of a fallopian tube, often after inflammation. They can narrow the tube, making pregnancy harder to achieve and increasing the risk of an ectopic pregnancy.

\synonyms
SIN

\medterm{Salpingo-Oophorectomy} Surgery to remove one or both fallopian tubes and ovaries. It may treat cancer, an ovarian or tubal problem, or reduce cancer risk in people with certain inherited gene changes. Removing both ovaries causes permanent loss of ovarian hormone production.

\medterm{Salpingo-OöPhorectomy} Surgical removal of a fallopian tube and ovary, also called salpingo-oophorectomy.
\medterm{Salpingography} Contrast imaging of the fallopian tubes, usually performed as part of hysterosalpingography to assess whether the tubes are open and to examine the uterine cavity.

\medterm{Salpingostomy (Salpingotomy)} An operation that opens a fallopian tube to remove an ectopic pregnancy while leaving the tube in place. It may preserve fertility in selected cases, but persistent pregnancy tissue, bleeding, scarring, or another ectopic pregnancy can occur.

\synonyms
Salpingotomy

\medterm{Salt} An ionic chemical compound; dietary sodium chloride is essential in small amounts but excess intake can increase blood pressure.

\medterm{Saltatory Conduction} The rapid movement of an electrical signal along a nerve fiber covered by myelin. The signal travels between small uncovered gaps in the covering, allowing fast communication with other nerves and muscles.

\medterm{Salter-Harris Fracture} A break involving the growth plate of a child or adolescent’s bone. The fracture pattern helps predict the risk of disturbed growth; pain, swelling, and inability to use the limb require prompt assessment and follow-up.

\medterm{Salvage Therapy} Treatment used when a disease, especially cancer, has not improved with earlier treatment or has returned. The choice depends on the disease, previous treatments, overall health, and treatment goals.

\medterm{Sample Fixation} Preserving a tissue or cell sample soon after collection so it does not decay or change shape. The preserved sample can then be processed, stained, and examined under a microscope for diagnosis.


\medterm{Sandfly Fever} An acute viral illness transmitted by phlebotomine sandflies, causing fever, headache, muscle pain, and sometimes rash.
\medterm{Sandhoff Disease} A rare inherited disorder in which fatty substances build up in nerve cells. It usually begins in infancy or childhood with loss of skills, weakness, seizures, movement problems, and sometimes an enlarged liver or spleen.

\medterm{Sandpaper Tooth} A tooth with a rough, grainy enamel surface that can hold plaque and feel uneven to the tongue. It may result from weak enamel, early mineral loss, fluorosis, or wear and should be checked by a dental professional.

\medterm{Sanfilippo Syndrome} A rare inherited disorder in which the body cannot break down part of a complex sugar. The buildup mainly affects the brain and causes progressive developmental and behavioral changes.
\synonyms
Mucopolysaccharidosis Type III (MPS III)


\medterm{Sanguineous} Containing or resembling fresh blood. The word often describes bright-red drainage from a wound; a large or rapidly increasing amount after surgery may indicate active bleeding and needs urgent assessment.
\medterm{Sanitation} Public-health measures that safely manage human waste, water, food, hygiene, and environmental contamination to prevent infection and disease transmission.

\medterm{Sano Shunt} A temporary tube connecting the lower-right heart chamber to the lung artery during the first operation for some babies with severe left-sided heart defects. It sends blood to the lungs while the heart is reconstructed; blockage or abnormal flow can be dangerous.

\synonyms
RV-PA Conduit; Norwood-Sano Modification

\medterm{Saphenous} Relating to the superficial veins of the lower limb, especially the great or small saphenous vein.
\medterm{Saphenous Vein} A superficial vein of the lower limb that drains the foot and leg. The great and small saphenous
veins are clinically important in venous disease and may be used as graft material in selected
vascular or cardiac bypass procedures.

\medterm{Small Saphenous Vein} A superficial vein running up the back of the calf and usually joining the popliteal vein behind the knee. It helps return blood from the foot and calf to the heart.
\medterm{SAPHO Syndrome} A rare inflammatory condition combining bone or joint pain with acne or pus-filled skin lesions. It often affects the breastbone, collarbone, spine, or jaw and can resemble infection or cancer, so diagnosis and treatment require specialist assessment.

\medterm{Sapovirus} A virus that can cause vomiting and watery diarrhea, especially in children. It spreads through contaminated food, water, surfaces, or close contact.

\medterm{Sapovirus Infection} An acute enteric infection caused by sapoviruses, positive-sense single-stranded RNA viruses in the family Caliciviridae that infect humans across all age groups. Clinical manifestations include acute watery diarrhea, nausea, vomiting, abdominal cramping, and low-grade fever, often responsible for outbreaks of foodborne or waterborne gastroenteritis in daycares, schools, and hospitals that resolve with supportive rehydration within 48 to 72 hours.

\synonyms
Sapovirus Gastroenteritis; Human Sapovirus Infection

\medterm{Sapporo Virus} An older name for sapovirus, a virus that can cause vomiting and watery diarrhea.

\medterm{Sapropterin} A synthetic form of tetrahydrobiopterin used for some people with phenylketonuria and certain tetrahydrobiopterin deficiencies; response and diet requirements vary.

\medterm{Saquinavir} An HIV protease inhibitor used with other antiretroviral medicines; drug interactions, gastrointestinal effects, and cardiovascular concerns limit its use.

\medterm{Saquinavir Mesylate} The mesylate salt form of saquinavir, an HIV medicine that blocks viral protein processing. It is used with other antiretroviral drugs, but interactions, liver problems, and dangerous heart-rhythm changes require careful monitoring.




\medterm{Sarcocystis Hominis} A protozoan parasite acquired by eating undercooked beef; people may develop intestinal symptoms, while cattle serve as an intermediate host.

\medterm{Sarcocystis Suihominis} A protozoan parasite acquired by eating undercooked pork that can cause intestinal sarcocystosis; pigs serve as an intermediate host.

\medterm{Sarcocystosis} Parasitic infection by Sarcocystis species, acquired from undercooked meat or contaminated food and sometimes causing muscle pain, fever, diarrhea, or no symptoms.

\medterm{Sarcoglycan Complex} A group of proteins that helps strengthen the outer covering of muscle cells during movement. Inherited problems with these proteins can cause progressive muscle weakness and sometimes heart-muscle disease.

\medterm{Sarcoglycanopathy} A group of inherited muscle diseases caused by changes in sarcoglycan genes. They usually cause progressive weakness of the hips and shoulders and may affect the heart.
\synonyms
Sarcoglycan-Deficient Limb-Girdle Muscular Dystrophy; LGMD R3-R6

\medterm{Sarcoidosis} An inflammatory disease in which clusters of immune cells form in one or more organs, most often the lungs and nearby lymph nodes. It can also affect the skin, eyes, heart, nervous system, or other organs.

\synonyms
Arthritis Sarcoid (Sarcoidosis)

\medterm{Sarcoidosis Skin Manifestations} Skin involvement in sarcoidosis can cause firm bumps, raised or discolored patches, scars that change, or lupus pernio on the nose or face. The appearance varies and may require a skin examination or biopsy.

\medterm{Sarcolemma} The thin outer covering of a muscle cell. It carries electrical signals that start contraction and helps maintain the cell’s structure and internal chemical balance.

\medterm{Sarcoma} A malignant tumor arising from connective or supporting tissues, including bone, cartilage, fat,
muscle, and blood vessels. Examples include osteosarcoma, liposarcoma, and rhabdomyosarcoma.

\medterm{Sarcoma, Bone Cancer} Alternate index label; see \textit{Bone cancer}.

\medterm{Sarcoma, Soft Tissue} Alternate index label; see \textit{Soft tissue sarcoma}.

\medterm{Sarcomere} The smallest working unit that makes skeletal and heart muscle contract. It contains protein threads that slide past one another, shortening the muscle and producing force.

\medterm{Sarcopenia} Loss of muscle strength and muscle mass that reduces physical function, often with aging but also with inactivity, poor nutrition, or chronic disease. It can cause weakness, slower walking, falls, and difficulty with daily activities.

\medterm{Sarcopenic Obesity} Excess body fat together with low muscle mass, strength, or physical ability. It can increase the risk of falls, disability, metabolic disease, and complications after surgery.

\medterm{Sarcoplasm} The material inside a muscle cell surrounding its specialized structures. It contains energy-producing parts, stored fuel, oxygen-holding myoglobin, and the protein machinery that produces contraction.

\medterm{Sarcoplasmic Reticulum} A network inside muscle cells that stores calcium. It releases calcium to start contraction and takes calcium back up so the muscle can relax.

\medterm{Sarcoptes} The mite genus containing Sarcoptes scabiei, the cause of scabies.
\medterm{Sarcoptes Scabiei} The microscopic mite that causes scabies. It burrows into the outer layer of skin and lays eggs, producing intense itching and a rash or thin lines. It spreads mainly through prolonged close skin contact.

\medterm{Sarcoptic Mange (Scabies)} Alternate index label for Scabies.




\medterm{Sargramostim} A recombinant granulocyte-macrophage colony-stimulating factor used to promote white-cell recovery after selected cancer treatments or transplantation; bone pain and fever may occur.

\medterm{Sarnat Staging} A three-level examination used to describe the severity of brain dysfunction in a newborn after reduced oxygen or blood flow. Mild disease causes increased alertness, moderate disease causes poor responsiveness and seizures, and severe disease causes coma and weak or absent reflexes; the stage helps guide urgent treatment and prognosis.

\synonyms
Sarnat and Sarnat Grading; Sarnat Score; HIE Staging

\medterm{SARS} Severe acute respiratory syndrome, a serious lung infection caused by the SARS coronavirus. It can cause fever, cough, breathing difficulty, and pneumonia and spreads mainly through close contact with respiratory droplets; suspected cases require public-health and medical guidance.
\medterm{SARS Virus} The coronavirus that causes severe acute respiratory syndrome, a serious respiratory infection that can spread through respiratory droplets and close contact.

\medterm{SARS-CoV-2} The coronavirus that causes COVID-19. It spreads mainly through particles released when an infected person breathes, talks, coughs, or sneezes and can cause anything from no symptoms to severe lung or other organ disease.

\medterm{Sassafras Oil Overdose} Poisoning from swallowing or otherwise being exposed to too much sassafras oil. It may cause nausea, vomiting, seizures, breathing problems, liver injury, or other serious effects and requires urgent medical advice.

\medterm{Satellite DNA} Repeated stretches of DNA found mainly in tightly packed parts of chromosomes. They help organize chromosome structure and can vary between people, usually without providing instructions for making proteins.

\medterm{Satellite Virus} A small virus-like infectious agent that can multiply only when a separate helper virus is present. It uses the helper’s machinery and may change how severe or persistent the combined infection becomes.

\medterm{Satiation} The feeling of fullness that terminates a meal, produced by gastrointestinal stretch, nutrients, hormones, and brain signals regulating food intake.

\medterm{Satiety} The feeling of fullness that reduces the desire to eat after a meal. It is influenced by stomach stretching, food composition, hormones, and brain signals; persistent early fullness can indicate illness.

\medterm{Satiety - Early} Feeling full sooner than expected while eating, sometimes after only a small amount of food. It may result from stomach or digestive disease and should be assessed if persistent or associated with weight loss.




\medterm{Saturated Fat} A dietary fat whose fatty-acid chains contain no carbon-carbon double bonds. It is found in foods such as fatty meat, butter, full-fat dairy, coconut oil, and palm oil; replacing some saturated fat with unsaturated fat can help lower LDL cholesterol and cardiovascular risk.


\medterm{Saturation} A state in which available binding, transport, or processing capacity is nearly full. Once capacity is reached, extra medicine or substance may remain in the blood or produce a disproportionately larger effect.

\medterm{Sausage Digit} Diffuse swelling of an entire finger or toe, often caused by inflammation in joints, tendons, and surrounding tissues.

\synonyms
Dactylitis


\medterm{Saxitoxin} A powerful poison made by some marine algae that can build up in shellfish. Eating contaminated seafood may cause tingling, weakness, trouble breathing, or paralysis.

\medterm{SBBO} Alternate index label; see \textit{Small intestinal bacterial overgrowth}.

\medterm{Scab} A crust of dried blood, serum, pus, or tissue fluid over a healing skin injury.
\medterm{Scabies} An infestation caused by tiny mites that burrow into the skin, causing intense itching that is often worse at night and a rash or thin burrows. It spreads mainly through prolonged close skin-to-skin contact.

\synonyms
Sarcoptic Mange (Scabies)


\medterm{SCAD} Alternate index label; see \textit{Spontaneous coronary artery dissection}.

\medterm{Scaffold} A three-dimensional framework used in tissue engineering to hold cells in place while they grow and organize into new tissue. It may be made from natural or synthetic material and may gradually dissolve.

\medterm{Scaffold Seeding} Placing living cells onto or inside a three-dimensional framework so they can attach, multiply, and begin forming tissue in laboratory or regenerative-medicine research.

\medterm{Scalded Mouth Syndrome} Alternate index label; see \textit{Burning mouth syndrome}.

\medterm{Scalded Skin Syndrome} Alternate index label; see \textit{Scalded-Skin Syndrome}.

\medterm{Scalded-Skin Syndrome} A serious infection, usually in young children, in which toxins from Staphylococcus bacteria make the outer skin layer separate. The skin becomes painful, red, and easily blistered or peeled; prompt treatment is needed to prevent dehydration and infection.
\medterm{Scalds} Burns caused by hot liquids or steam. They may damage only the surface skin or extend into deeper tissue; severity depends on depth, size, location, and age. Large, deep, facial, genital, hand, or circumferential burns need urgent specialist assessment, and unusual patterns in children require safeguarding evaluation.
\medterm{Scale} Visible flakes or plates of stratum corneum shed from the skin surface, commonly caused by abnormal keratinization, inflammation, infection, or dryness.

\medterm{Scales} Scales are flakes or plates of shed skin on the surface, occurring with dryness, eczema, psoriasis, fungal infection, or other dermatologic conditions.

\medterm{Scaling} Scaling is flaking or shedding of the outer skin layer and may occur with dryness, eczema, psoriasis, fungal infection, or other skin disease.

\medterm{Scaling And Root Planing} A deep dental cleaning that removes hardened plaque and bacteria above and below the gumline and smooths tooth roots. It reduces gum inflammation and may slow periodontal disease, but severe damage may need additional treatment.

\medterm{Scalp} The skin and underlying tissues covering the skull. It contains hair follicles, glands, blood vessels, nerves, and connective tissue and can be affected by injury, infection, or skin disease.
\medterm{Scalp And Body (Ringworm)} Alternate index label for Ringworm.

\medterm{Scalp Ringworm} A contagious fungal infection of the scalp, also called tinea capitis. It can cause scaling, itching, broken hairs, bald patches, or swollen tender areas and usually needs prescription antifungal medicine taken by mouth to prevent persistent hair loss or spread.

\medterm{Scalp Psoriasis} Psoriasis affecting the scalp, causing sharply demarcated red plaques covered by silvery scale that may extend beyond
the hairline. It can itch, crack, or bleed and is not contagious. Treatment may include medicated shampoos, topical
corticosteroids, vitamin-D analogues, or systemic therapy for extensive disease.

\synonyms
Psoriasis Of The Scalp (Scalp Psoriasis)

\medterm{Scalpel} A sharp surgical knife used to make precise incisions or cuts in tissue. It may be a reusable handle with a detachable blade or a disposable single-piece instrument. Blade designs and numbers differ in size and shape; for example, a No. 10 blade is commonly used for larger incisions, a No. 11 for puncture or stab incisions, and a No. 15 for shorter precise cuts.

\synonyms
Surgical Knife

\medterm{Scalpel Handle} The reusable handle that holds a detachable scalpel blade. Different handle and blade sizes are chosen for the type and precision of the incision.

\medterm{Scan} An imaging examination that creates pictures of body structures or functions using methods such as ultrasound, X-rays, CT, MRI, or radioactive tracers. The method and interpretation depend on the clinical question.

\medterm{Scan Gallium} A nuclear-medicine scan using radiolabeled gallium to identify areas of inflammation, infection, or some tumors; uptake is interpreted with anatomy and clinical context.

\medterm{Scanning} Systematic examination of a body region or specimen using imaging, microscopy, or another detection method to locate, characterize, or monitor abnormalities.

\medterm{Scanning Electron Microscopy} A method that scans a specimen’s surface with an electron beam to produce a highly magnified, three-dimensional-looking image of its shape and surface details.

\medterm{Scanning Speech} Speech that is unusually slow, broken into separate syllables, and spoken with uneven rhythm or emphasis. It can occur when disease affects the cerebellum or other parts of the nervous system.
\medterm{Scanning Techniques} Methods used to create images or measurements of the body or a sample. Examples include ultrasound, X-ray, CT, MRI, and nuclear-medicine scans; the method is chosen according to the clinical question.
\medterm{Scaphoid} A small wrist bone on the thumb side, commonly injured when someone falls onto an outstretched hand. Pain in the hollow just behind the thumb may indicate a fracture, which sometimes heals poorly because part of the bone has limited blood flow.

\medterm{Scaphoid Abdomen} A scaphoid abdomen is an unusually hollowed abdominal contour, classically seen with severe diaphragmatic hernia in a newborn but also possible with marked wasting.

\medterm{Scaphoid Bone} The largest proximal-row carpal bone on the thumb side of the wrist; its waist is vulnerable to fracture and avascular necrosis because of its blood supply.

\medterm{Scaphoid Fractures} Breaks in the small wrist bone near the base of the thumb, often after falling onto an outstretched hand. Pain may be felt in the hollow beside the thumb; some fractures heal poorly because part of the bone has a limited blood supply.

\medterm{Scaphycephaly} A long, narrow head shape caused when the major seam running from front to back of the skull closes too early. It is a form of craniosynostosis and may require specialist assessment to protect brain and skull development.

\medterm{Scapula} The shoulder blade, a flat triangular bone at the back of the chest. It connects the upper arm to the collarbone and provides attachment for muscles that move and stabilize the shoulder.

\medterm{Scapulothoracic Joint} The functional gliding connection between the shoulder blade and the back of the chest wall. It is not a true joint but works with nearby joints to position and move the arm.

\medterm{Scar} Fibrous tissue formed during healing after injury, inflammation, or surgery.
\medterm{Scar Excessive (Keloid)} Alternate index label for Keloid.

\medterm{Scar Management} Measures used to reduce symptoms and abnormal scar formation. Options include silicone
products, pressure therapy, massage, corticosteroid injections, or other specialist treatments, depending on the scar
and its cause.

\medterm{Scar Revision} A procedure that improves the appearance, contour, flexibility, or symptoms of a scar through excision, rearrangement, grafting, laser, injection, or other selected techniques.


\medterm{Scarlatina} Scarlet fever, an illness caused by toxin-producing group A Streptococcus and characterized by sore throat, fever, and a fine rough-textured rash.

\synonyms
Scarlet Fever (Scarlatina)

\medterm{Scarlatina (Scarlet Fever (Scarlatina))} Alternate index label for Scarlet Fever (Scarlatina).

\medterm{Scarlet Fever} An illness caused by toxin-producing group A Streptococcus. It usually causes sore throat, fever, and a widespread fine, rough-textured rash; a red tongue and pale skin around the mouth may also occur.

\medterm{Scarlet Fever (Scarlatina)} Alternate index label for Scarlatina.


\medterm{Scarring Alopecia} Permanent hair loss caused when inflammation or injury destroys hair follicles and replaces them with scar tissue. Early scalp redness, scaling, itching, pain, or spreading hair loss needs dermatologic assessment because lost follicles cannot regrow hair.

\medterm{Scars} A scar is fibrous tissue formed during healing after injury, surgery, inflammation, or burns. Scars may be flat, depressed, raised, or contractured and can itch, hurt, restrict movement, or affect appearance; treatment depends on maturity and type.

\medterm{Scatter Correction} A medical-imaging method that estimates and reduces errors caused when detected radiation changes direction after passing through tissue. It can improve the accuracy of a nuclear-medicine image or measured activity.

\medterm{Scatter Events} Radiation particles that change direction after interacting with tissue or another material. In medical imaging they can place signals in the wrong location or reduce contrast, so scanners use design and correction methods to limit their effect.
\medterm{Scavenger Receptor} A protein on or inside a cell that recognizes and helps remove altered fats, germs, cell debris, or dying cells. It supports early immune defense but can also help fatty material build up in artery walls.

\medterm{Scavenging System} Equipment that collects and removes waste anesthetic gases from an operating room. It protects staff from repeated exposure and may use a vacuum system or passive venting to carry gases safely away.

\medterm{Scedosporium} A genus of environmental molds that can cause difficult-to-treat lung, sinus, skin, bone, or disseminated infections, especially in immunocompromised people.

\medterm{Scedosporium Apiospermum} A mold that can cause localized or disseminated infection, including mycetoma, pneumonia, sinusitis, and brain abscesses; treatment often requires specialist antifungal care.

\medterm{Scedosporium Prolifican} An environmental mold, now commonly classified as Lomentospora prolificans, that causes severe disseminated infection and is often resistant to antifungal drugs.

\medterm{Schaalia Radingae} A normally slow-growing anaerobic bacterium of the human mouth and skin that can rarely cause abscesses or other invasive infections.

\medterm{Schaalia Turicensis} A filamentous anaerobic bacterium that may be part of normal flora but can cause actinomycosis-like abscesses and other opportunistic infections.

\medterm{Schatzki Ring} A thin ring of tissue narrowing the lower esophagus, often near a hiatus hernia. It may cause occasional difficulty swallowing solid food or a piece of food becoming stuck; treatment can widen the narrowed area.

\synonyms
Ring Schatzki (Schatzki Ring)
Strictures Esophagus (Schatzki Ring)

\medterm{Schedule I} A United States controlled-substance category for drugs with high misuse potential and no accepted medical use
under federal law.

\medterm{Schedule II} A United States legal category for medicines with accepted medical uses but high misuse, dependence, or overdose potential. Prescribing, dispensing, refills, and storage are subject to strict rules.
\medterm{Schedule III} A United States legal category for medicines with accepted medical uses and less misuse potential than Schedule II substances, although dependence and misuse remain possible. Rules vary by medicine and jurisdiction.
\medterm{Schedule IV} A United States controlled-substance category for drugs with accepted medical uses and lower misuse potential
than Schedule III drugs.

\medterm{Schedule V} A United States controlled-substance category for drugs with the lowest level of misuse potential among the
federal schedules.

\medterm{Scheduled Notification} A planned alert or message delivered at a specified time to remind a person about an appointment, medicine, monitoring task, or health-related action.

\medterm{Scheuermann Disease} A juvenile osteochondrosis of the spine in which vertebral end plates develop irregularities and anterior wedging during skeletal growth.
It commonly produces a rigid thoracic kyphosis, although the thoracolumbar or lumbar spine may also be involved.
Symptoms can include back pain, fatigue, limited spinal motion, and postural deformity that becomes more apparent during adolescence.
Diagnosis is based on clinical examination and characteristic radiographs; management depends on deformity severity, skeletal maturity, pain, and progression.

\medterm{Scheuermann's Disease} A growth disorder beginning in adolescence that causes several spinal bones to become wedge-shaped. This can produce a lasting rounded upper back, stiffness, pain, and sometimes breathing or activity limitations.
\medterm{Schiff Base} A chemical compound formed when a primary amine reacts with an aldehyde or ketone; Schiff bases occur in biochemistry and drug chemistry.

\medterm{Schiller Test} An examination of the cervix in which iodine solution is applied to highlight areas of abnormal tissue. Areas that do not stain may need further examination or biopsy; the test alone cannot diagnose cancer.

\medterm{Schilling Test} An older test that measured how well the body absorbed vitamin B12, sometimes used to investigate pernicious anemia or intestinal disease. It is rarely performed now because blood tests and antibody testing have largely replaced it.

\medterm{Schirmer Test} A test in which standardized filter paper strips measure tear production over a set time, commonly used to evaluate aqueous-deficient dry eye.

\medterm{Schistocyte Count} Measurement of fragmented red blood cells on a blood smear, supporting evaluation of microangiopathic hemolysis, mechanical valve damage, disseminated intravascular coagulation, or severe burns.

\medterm{Schistocytes} Broken, irregular, or helmet-shaped red blood cells seen on a blood-smear test. They form when red cells are damaged in abnormal small vessels, by artificial heart valves, or in other illnesses and may signal serious red-cell destruction.

\medterm{Schistomicide} A medicine that kills parasitic blood flukes responsible for schistosomiasis. Treatment choice depends on the parasite species and the organs affected, and medical advice is needed because infection may cause lasting damage.
\medterm{Schistosoma} A group of parasitic worms that live in blood vessels. Their young forms can enter through skin exposed to contaminated fresh water and later cause intestinal, urinary, liver, or other disease.

\medterm{Schistosoma Haematobium} A parasitic blood fluke that can cause urinary schistosomiasis, including hematuria and chronic
injury of the urinary tract.

\medterm{Schistosoma Intercalatum} A parasitic blood fluke that causes intestinal schistosomiasis, mainly in parts of Central and West Africa, after freshwater exposure.

\medterm{Schistosomiasis (Bilharzia)} An infection caused by parasitic worms called \textit{Schistosoma}. The worms live in freshwater snails, and their larvae can penetrate the skin when a person swims, wades, or works in contaminated freshwater. The infection may cause an itchy rash, fever, abdominal pain, diarrhea, blood in the stool or urine, and painful urination. Long-term infection can damage the intestine, liver, urinary tract, or other organs.

\synonyms
Bilharzia

\medterm{Schizencephaly} A rare condition present from birth in which a fluid-filled cleft extends through part of the brain. It results from abnormal early brain development and may cause seizures, weakness, movement problems, or delayed development, depending on the cleft’s size and location.

\medterm{Schizoaffective Disorder} A mental-health disorder in which a person has symptoms of schizophrenia, such as hallucinations or disorganized thinking, together with major depression or mania. Psychotic symptoms also occur for a period without a major mood episode.

\medterm{Schizogony} A way some single-celled parasites reproduce without mating. One parasite cell grows and divides internally to produce many new cells, which are then released to infect other cells.
\medterm{Schizoid Personality Disorder} A long-standing pattern of limited interest in close relationships and restricted emotional expression. It differs from schizophrenia because hallucinations and fixed delusions are not central features.
\synonyms
Personality Disorder, Schizoid

\medterm{Schizophrenia} A serious, long-term mental disorder that can affect a person’s thoughts, perceptions, emotions, behavior, and ability to function. Symptoms are often grouped as psychotic, negative, and cognitive. Psychotic symptoms include hallucinations, such as hearing or seeing things that others do not, delusions, which are firmly held beliefs not supported by shared evidence, and disorganized thought or speech. Negative symptoms may include reduced emotional expression, motivation, pleasure, or social interest. Cognitive symptoms may affect attention, memory, planning, and decision-making. Symptoms and their severity vary between people and may come and go or persist over time.

\medterm{Schizophrenia, Adolescent} Alternate index label; see \textit{Childhood schizophrenia}.

\medterm{Schizophrenia, Childhood} Alternate index label; see \textit{Childhood schizophrenia}.

\medterm{Schizophrenia, Juvenile} Alternate index label; see \textit{Childhood schizophrenia}.

\medterm{Schizotypal Personality Disorder} A long-term pattern of discomfort with close relationships, unusual beliefs or perceptions, unusual speech or behavior, and difficulty interpreting social situations. It is different from schizophrenia because persistent hallucinations and major loss of reality testing are not typical.

\synonyms
Personality Disorder, Schizotypal

\medterm{Schlemm’S Canal} A circular drainage channel at the edge of the cornea that carries fluid from the front of the eye into blood vessels. Poor drainage can raise pressure inside the eye and contribute to glaucoma.

\synonyms
Canal of Schlemm

\medterm{Schnitzler Syndrome} A rare inflammatory disease causing a long-lasting hives-like rash, repeated fever, bone or joint pain, and raised inflammation tests. It is usually associated with an abnormal IgM blood protein and requires specialist evaluation.

\medterm{School Age Test Or Procedure Preparation} Preparing a school-age child for a test or procedure includes age-appropriate explanations, comfort measures, caregiver support, and following any fasting or medicine instructions.

\medterm{School Sores (Impetigo)} Alternate index label for Impetigo.

\medterm{School-Age Children Development} School-age development includes advances in growth, coordination, language, learning, independence, social relationships, and emotional regulation, with broad normal variation.

\medterm{School-Based Rehabilitation} Therapy and other support services provided in a school setting to help a child participate in education and daily school activities.

\medterm{Schwann Cell} A supporting cell of the peripheral nervous system. It wraps around many nerve fibers to form their insulating myelin covering and helps maintain or repair peripheral nerves.
\medterm{Schwann Cells} Glial cells of the peripheral nervous system that form myelin around many axons and support the repair and maintenance of peripheral nerves.

\medterm{Schwannoma} A usually noncancerous, slow-growing tumor arising from cells that support nerves. It may cause a lump, pain, numbness, or weakness; a tumor involving the hearing and balance nerve can cause hearing loss, ringing, or dizziness.

\synonyms
Neurilemmoma

\medterm{Schwannomatosis} A condition in which multiple schwannomas, usually slow-growing tumors of the covering of nerves, develop. The main problem is often long-lasting pain, although numbness, weakness, or other nerve symptoms can occur depending on tumor location.

\synonyms
Neurinomatosis

\medterm{SCI} Alternate index label; see \textit{Spinal cord injury}.

\medterm{Sciatic Nerve} The body's largest nerve, running from the lower spine through the buttock and down the back of the leg. It supplies movement and sensation to much of the leg and foot.

\medterm{Sciatic Neuritis (Sciatica)} Alternate index label for Sciatica.

\medterm{Sciatica} Pain, tingling, numbness, or weakness that travels from the lower back or buttock down the leg, usually because a nerve root is irritated by a disk problem or narrowing of the spine. Progressive weakness, numbness around the inner thighs or genitals, or new loss of bladder or bowel control may indicate severe nerve compression and requires urgent assessment.

\synonyms
Sciatic Neuritis (Sciatica)

\medterm{Scimitar Syndrome} A birth defect in which some veins from the right lung drain to the wrong large vein instead of directly to the heart. The right lung may be smaller, and the condition can cause repeated chest infections, breathing problems, high pressure in the lung arteries, or heart failure.

\medterm{Scintigraphy} Nuclear-medicine imaging in which a small amount of radioactive tracer is given and a camera records where it collects. The pattern can show the function, blood flow, or activity of organs such as bone, thyroid, kidneys, or heart.

\medterm{Scintillation} The brief release of light by a material after it absorbs radiation. Detectors measure these flashes to estimate radiation and create images in some medical scans.

\medterm{Scintillation Crystal} A material in some medical detectors that changes incoming radiation into tiny flashes of light. The flashes are measured to create images or estimate the amount of radioactive tracer present.

\medterm{Scintillations} Flashes, sparkles, or zigzag lines seen in the field of vision. They often occur during a migraine aura but may also result from an eye or brain problem; sudden new flashes with floaters or a curtain-like shadow need urgent eye assessment.

\medterm{Scirrhus} An older term for a hard cancer containing a large amount of scar-like tissue. It was used especially for some breast or stomach cancers and describes texture rather than a separate cancer type.
\medterm{Scissors} Instruments with two blades used to cut tissue, sutures, dressings, or other material. Surgical types vary in length, blade shape, and sharpness; examples include Mayo and Metzenbaum scissors.

\medterm{Scissors (Surgical)} Surgical instruments designed to cut tissue, sutures, dressings, or other materials. Different shapes and blade types are selected for delicate, deep, curved, or heavy-duty work.
\synonyms
Surgical scissors

\medterm{Scissors (Vascular)} A surgical instrument with two blades used to cut tissue, sutures, dressings, or other materials
during vascular procedures.

\synonyms
Vascular scissors

\medterm{Sclera} The sclera is the tough white outer coat of the eye that protects the globe and provides attachment for eye muscles.

\medterm{Scleral Buckling} A retinal-detachment repair procedure in which a silicone band or buckle is placed around the outside of the eye to support the retinal break and reduce traction.

\medterm{Scleredema Diabeticorum} Scleredema diabeticorum is thickening and hardening of the skin on the upper back and neck, usually associated with long-standing diabetes and obesity.

\medterm{Scleritis} Painful inflammation of the white outer coat of the eye. It is often associated with autoimmune disease and can damage the eye or threaten vision, so prompt eye assessment is needed.

\synonyms
Inflammation Sclera (Scleritis)

\medterm{Scleroderma} A group of autoimmune diseases that cause inflammation and too much collagen to build up, making skin and sometimes deeper tissues hard, thick, and tight. Localized scleroderma, also called morphea or linear scleroderma, mainly affects the skin and tissues just beneath it. Systemic scleroderma, also called systemic sclerosis, can affect the skin, small blood vessels, and internal organs such as the esophagus, lungs, kidneys, and heart. The cause is not fully known, and symptoms vary with the type and organs involved.

\synonyms
Arthritis Scleroderma (Scleroderma)
Sclerosis, Systemic


\medterm{Scleroma} A long-lasting infection, usually of the nose and upper airway, caused by Klebsiella rhinoscleromatis. It can produce scar tissue, blocked breathing, nasal discharge, voice changes, or facial deformity.

\medterm{Scleromyxedema} A rare long-term skin disorder in which many small, firm, waxy bumps and thickened skin develop, especially on the face, neck, and hands. It may also affect muscles, joints, swallowing, lungs, or other organs.

\medterm{Sclerosing Cholangitis} Inflammation and scarring that narrow or block the tubes carrying bile from the liver. It may cause itching, jaundice, abdominal pain, repeated infections, and liver damage; the cause guides treatment and follow-up.

\medterm{Sclerosing Mesenteritis} A rare inflammatory condition causing thickening and scarring in the tissue that supports the intestines. It may cause abdominal pain, nausea, weight loss, or bowel blockage, although some people have no symptoms.

\medterm{Sclerosing Peritonitis} Thickening and scarring of the lining inside the abdomen, sometimes after long-term peritoneal dialysis. The scar tissue can restrict the intestines and cause pain, vomiting, poor nutrition, or bowel blockage.
\medterm{Sclerosing Sialadenitis} Long-term inflammation that makes a salivary gland swollen and hard because scar-like tissue forms inside it. It may be related to an immune disorder and can cause dry mouth, pain, or repeated swelling.

\medterm{Sclerosis} Abnormal hardening or scarring of tissue, often from excess connective tissue, chronic inflammation, degeneration, or neurologic plaques; the affected organ determines the specific disease.

\medterm{Sclerosis, Multiple} Alternate index label; see \textit{Multiple sclerosis}.

\medterm{Sclerosis, Systemic} Alternate index label; see \textit{Scleroderma}.

\medterm{Sclerosteosis} A rare inherited condition in which bones, especially the skull, become abnormally thick. It may cause a broad or tall build, joined fingers, headaches, hearing loss, facial weakness, or pressure on nerves inside the skull.

\medterm{Sclerotherapy} A treatment in which an irritating or closing agent is injected into a blood vessel or body cavity to make it scar and seal, commonly used for varicose veins and selected vascular lesions.


\medterm{Scoliosis} An abnormal sideways curve of the spine that may look like the letter “C” or “S.” The bones of the spine may also rotate, changing the shape of the back or rib cage. In many children and teenagers, the cause is unknown; this is called idiopathic scoliosis. Other causes include differences present at birth, nerve or muscle conditions, spinal injury, or age-related changes. Mild scoliosis may cause no symptoms, while larger curves may cause uneven shoulders or hips, a prominent shoulder blade, back pain, reduced movement, or, rarely, breathing problems.

\synonyms
Spinal Curvature



\medterm{Scopolamine} Scopolamine is a muscarinic antagonist used mainly to prevent motion sickness and postoperative nausea; dry mouth, blurred vision, urinary retention, confusion, and other anticholinergic effects may occur.

\medterm{Scopolamine Hydrobromide} The hydrobromide salt of scopolamine, an anticholinergic medicine used mainly to prevent motion sickness and postoperative nausea.


\medterm{Scopulariopsis} Scopulariopsis is a mold genus that rarely causes opportunistic infection, especially in people with immune suppression or implanted devices.

\medterm{Scorbutic} Relating to scurvy caused by severe vitamin C deficiency. It can cause easy bruising, bleeding gums, poor wound healing, anemia, bone or joint pain, and weakness; treatment with vitamin C and an adequate diet usually leads to recovery.
\medterm{Scorpion Fish Sting} A painful injury from the venomous spines of a scorpion fish, causing severe local pain and swelling and sometimes whole-body symptoms. The wound needs urgent medical assessment because spines or venom may remain in the injury.

\medterm{Scorpion Sting} A venomous injury from a scorpion, usually causing immediate pain, burning, redness, or tingling. Some species can cause muscle twitching, abnormal eye movements, breathing difficulty, or dangerous autonomic symptoms.
\synonyms
Stings, Scorpion

\medterm{Scorpions} Venomous arachnids whose stings can cause local pain and, depending on the species, sweating, muscle spasms, breathing problems, or other serious symptoms.

\medterm{Scotoma} A blind or blurred spot within an otherwise usable field of vision. It may result from retinal disease, optic-nerve damage, migraine, stroke, or another eye or brain disorder.

\medterm{Scrape} A scrape is a superficial abrasion in which friction removes part of the skin surface; cleaning, protection, and monitoring for infection support healing.


\medterm{Screen Time And Children} Time children spend using televisions, computers, phones, tablets, or game devices; excessive or poorly timed use can affect sleep, activity, learning, mood, and family routines.

\medterm{Screening} Testing people without symptoms to identify a disease or an increased risk of disease. The appropriate test,
age range, interval, and follow-up depend on the condition, individual risk, and current
guidance; screening benefits must be balanced against harms.

\medterm{Screening (Geriatric)} Checking older adults who have no symptoms for diseases or risks that may benefit from early detection. Decisions should consider expected benefit, possible harm, life expectancy, health goals, and preferences; examples include falls, depression, osteoporosis, and selected cancers.

\synonyms
Geriatric screening; screening in older adults

\medterm{Screening And Diagnosis For HIV} Testing for human immunodeficiency virus using an antigen-antibody assay or nucleic-acid test, followed by confirmatory testing and assessment of immune status when positive.

\medterm{Screening Procedure} A test or examination offered to people without symptoms to detect disease or risk early enough for effective prevention or treatment; eligibility and intervals depend on evidence and risk.

\medterm{Screening Program} An organized public-health service that systematically invites a defined population to
undergo testing for an asymptomatic condition, with quality-assured follow-up,
diagnosis, treatment, and evaluation of benefits and harms.

\medterm{Screening Study} A planned investigation of an asymptomatic population to identify early disease, precursor lesions, or risk factors using a validated test and follow-up pathway.

\medterm{Screening Test} A test applied to people without symptoms or a known diagnosis to identify those at
increased risk of a condition. A screening result usually requires a diagnostic or
confirmatory evaluation before treatment is initiated.

\medterm{Screws} Threaded metal devices used to hold broken bones together or attach plates and other implants securely to bone.

\medterm{Script} A script is a written set of instructions or a prescription; in healthcare, a prescription specifies the medicine, dose, route, timing, and duration.

\medterm{Scrivener's Palsy} An older term for task-specific hand dystonia or cramp occurring during writing.
\medterm{Scrofula} Tuberculous infection of cervical lymph nodes, producing chronic neck swellings that may suppurate or form draining sinuses; nontuberculous mycobacteria can cause a similar presentation.

\medterm{Scrofulous} Scrofulous is an old term for tuberculous infection of lymph nodes, especially in the neck, which can cause chronic swelling and draining sinuses.

\medterm{Scrotal Calcinosis} Firm, usually painless calcium deposits in the skin of the scrotum that appear as one or more whitish or yellowish lumps. They may become tender, ulcerate, or release a chalky material; removal is considered if they cause symptoms.

\medterm{Scrotal Infection} Infection of the scrotal skin or underlying structures, which may cause redness, swelling, warmth, pain,
or drainage. Rapidly worsening pain, fever, discoloration, or severe swelling is an emergency
because serious soft-tissue infection must be excluded.

\medterm{Scrotal Masses} Lumps or swellings in the scrotum, arising from the skin, fluid around a testicle, veins, infection, hernia, or a testicular tumor. A new, firm, painless, or rapidly growing lump needs prompt examination and ultrasound.

\medterm{Scrotal Pain} Pain in the scrotum caused by infection, trauma, hernia, fluid collections, varicocele, testicular torsion, or other conditions affecting the testes and surrounding structures.

\medterm{Scrotal Swelling} Enlargement of the scrotum caused by fluid, inflammation, infection, hernia, varicocele, hydrocele, tumor, trauma, or testicular torsion; sudden pain is an emergency.

\medterm{Scrotal Ultrasound} Ultrasound imaging of the scrotum and testes used to evaluate pain, swelling, masses, torsion, infection, trauma, and blood flow.

\medterm{Scrotum} The scrotum is the external skin-and-muscle sac containing the testes and epididymides; its muscles help regulate testicular temperature, and swelling or pain may indicate infection, torsion, hernia, or tumor.


\medterm{Scrub In} To scrub in is to perform surgical hand antisepsis and put on sterile clothing and gloves before participating in a sterile procedure.

\medterm{Scrub Up} To scrub up means to clean and antiseptically prepare the hands and forearms before a sterile procedure.

\medterm{Scrubs} Scrubs are protective work garments commonly worn by healthcare staff; they reduce clothing contamination but do not replace hand hygiene or other precautions.

\medterm{Scrub Typhus} An acute febrile zoonotic infection caused by the intracellular bacterium Orientia tsutsugamushi and transmitted to humans through the bite of larval trombiculid mites (chiggers). Characteristic manifestations include high fever, severe headache, generalized lymphadenopathy, maculopapular rash, and a diagnostic necrotic skin lesion with a black crust (eschar) at the bite site; severe cases can progress to pneumonitis, acute kidney injury, myocarditis, or acute encephalitis syndrome.

\synonyms
Tsutsugamushi Disease; Mite-Borne Typhus

\medterm{Scurf} Fine flakes or scales shed from the skin, especially the scalp. It may occur with dry skin, dandruff, eczema, psoriasis, infection, or irritation and may be accompanied by itching or redness.
\medterm{Scurvy} A disease caused by severe vitamin C deficiency. It can cause bleeding or swollen gums, easy bruising,
small skin bleeds, poor wound healing, anemia, bone or joint pain, and fatigue.

\medterm{Sea-Blue Histiocyte} An immune cell that appears blue-green under a microscope because it contains stored fats or cell-waste pigment. It may be seen in some storage, blood, or metabolic disorders, but the finding alone does not identify a disease.

\medterm{Sea-Sickness} Motion sickness caused by conflicting signals from the inner ear, eyes, and body during travel on a boat or ship. It may cause nausea, vomiting, dizziness, sweating, headache, and pallor.
\medterm{Seasonal Affective Disorder} Depression that repeatedly begins during a particular season, most often autumn or winter, and improves at another time of year. Symptoms may include low mood, loss of interest, tiredness, oversleeping, increased appetite, and social withdrawal.

\medterm{Seasonal Affective Disorder Syndrome} Alternate index label for Seasonal Affective Disorder.
\medterm{Seasonal Allergies (Hay Fever)} Alternate index label for Hay Fever.

\medterm{Seasonal Allergy} Alternate index label; see \textit{Hay fever}.

\medterm{Sebaceous Adenoma} A sebaceous adenoma is a usually benign tumor of sebaceous glands that can appear as a yellowish skin nodule and may be associated with Muir-Torre syndrome.

\medterm{Sebaceous Carcinoma} An uncommon cancer arising from sebaceous, or oil-producing, glands. It often appears as a firm or yellowish eyelid or skin nodule and may require biopsy and specialist treatment.

\medterm{Sebaceous Cyst} A common but imprecise term for an epidermoid cyst, a keratin-filled cyst beneath the skin.

\medterm{Sebaceous Epithelioma} A benign sebaceous neoplasm, often presenting as a yellowish papule or nodule. Multiple or unusual lesions may warrant assessment for a DNA mismatch-repair syndrome such as Muir-Torre syndrome.

\medterm{Sebaceous Gland} A small gland in the skin that secretes sebum (oil) into hair follicles to lubricate and
waterproof the skin and hair.

\medterm{Sebaceous Gland Carcinoma} Sebaceous gland carcinoma is an uncommon skin cancer arising from oil glands, often appearing as a firm yellowish eyelid or skin lump.

\medterm{Sebaceous Glands} Skin glands that produce sebum, an oily substance usually released into hair follicles. They are most abundant on the face and scalp and contribute to acne when blocked or inflamed.
\medterm{Sebaceous Hyperplasia} A benign enlargement of sebaceous glands, presenting as small yellowish papules with central umbilication, usually on the face of adults.

\medterm{Seborrhea} Seborrhea is excessive or altered oil production that contributes to greasy scale and inflammation, commonly in seborrheic dermatitis of the scalp, face, or chest. Gentle cleansing, medicated shampoos, and topical anti-inflammatory or antifungal treatment may help.

\medterm{Seborrheic Dermatitis} A common long-term inflammation of oily skin, especially the scalp, eyebrows, sides of the nose, chest, and skin folds. It causes red or pink patches with greasy white or yellow scale; dandruff is a mild scalp form.

\synonyms
Dermatitis, Seborrheic

\medterm{Seborrheic Keratoses} Common, benign, epidermal skin tumors presenting as waxy, stuck-on, pigmented plaques
with a rough, verrucous surface, typically appearing on the trunk and face of
middle-aged and older adults. They are the most common benign skin neoplasm and are
caused by keratinocyte proliferation.

\medterm{Seborrheic Keratosis} A benign epidermal tumor appearing as stuck-on, waxy, hyperpigmented plaques with a
verrucous surface. They are extremely common in older adults and increase in number with
age. Seborrheic keratoses are not premalignant.

\synonyms
Keratosis, Seborrheic

\medterm{Seborrhoea} A chronic inflammatory skin condition affecting sebum-rich areas: scalp, face
(nasolabial folds, eyebrows), chest, and intertriginous areas. See Seborrheic
Dermatitis.

\medterm{Seborrhoeic} Relating to increased or altered activity of sebaceous glands or to areas rich in these glands, such as the scalp, face, chest, and upper back.

\medterm{Seborrhoeic Keratosis} A common harmless skin growth that may be brown, black, or light-colored with a waxy, rough, or stuck-on appearance. It is not contagious or cancerous, but changing lesions should be checked.

\medterm{Sebum} An oily substance made by skin glands that lubricates and protects the skin and hair. Excess production can contribute to oily skin, acne, and dandruff, while changes may accompany some skin disorders.
\medterm{Seckel Syndrome} A rare inherited condition causing severe growth restriction before and after birth, a very small head, characteristic facial features, and often developmental or intellectual disability. Bone, blood, and other birth differences may also occur.

\medterm{Secobarbital Overdose} Poisoning from taking too much secobarbital, a sedative medicine. It can cause extreme sleepiness, confusion, slowed breathing, coma, or death and requires emergency care.

\medterm{Second Malignancy} A new, separate cancer that develops after an earlier cancer. The risk may be influenced by inherited factors, previous radiation or chemotherapy, age, and other exposures; it is different from recurrence of the first cancer.

\medterm{Second Messenger Systems} Internal signaling systems that carry a message from a hormone, nerve chemical, or growth factor at the cell surface to processes inside the cell. They can change contraction, secretion, metabolism, gene activity, or cell growth.

\medterm{Second Messengers} Small molecules inside cells that pass along messages received from hormones, nerve chemicals, or other outside signals. They help control contraction, secretion, metabolism, gene activity, and cell growth.

\medterm{Second Opinion} An independent medical assessment by another qualified clinician to help confirm or clarify a diagnosis, prognosis, or treatment plan.

\medterm{Second Stage Of Labor} The part of childbirth that begins when the cervix is fully open and ends when the baby is born. It includes descent of the baby and pushing; its length varies with previous births, pain relief, the baby’s position, and the health of parent and baby.

\medterm{Second Trimester} The middle part of pregnancy, from 14 through 27 completed weeks. The fetus grows rapidly, movement is often felt, and ultrasound and blood screening may assess development and certain conditions; the timing and tests vary by pregnancy and local care.

\medterm{Second-Degree AV Block} A heart-rhythm problem in which some electrical signals from the upper chambers fail to reach the lower chambers. The heartbeat may be slow or irregular; certain patterns can progress to complete heart block and need urgent assessment.

\medterm{Second-Degree Burn} A burn that damages the outer skin and part of the layer beneath it. The area is usually painful, red, moist, and blistered; deeper burns may look pale and heal slowly. Large, deep, or sensitive-area burns need medical assessment.

\medterm{Second-Hand Smoke} Smoke inhaled from burning tobacco products or exhaled by a person who is smoking. It contains toxic
substances and increases health risks in exposed people.

\medterm{Second-Look Surgery} Second-look surgery is a planned or follow-up operation performed to inspect healing, remove remaining disease, assess treatment response, or manage complications.

\medterm{Secondary Adrenal Insufficiency} Reduced production of adrenal hormones because the pituitary or brain does not provide enough stimulation, often after long-term steroid treatment is stopped suddenly. It can cause fatigue, weakness, nausea, low blood pressure, and low blood sugar; skin darkening and salt loss are less typical than in primary adrenal failure.

\medterm{Secondary Amenorrhea} Absence of menstruation for at least three months in someone with previously regular cycles or six months with previously irregular cycles, caused by pregnancy, endocrine, ovarian, uterine, or systemic factors.

\medterm{Secondary Attack Rate} The proportion of susceptible people who become infected after contact with an initial case during a defined period. It helps estimate how readily an infection spreads in households, schools, healthcare settings, or other close-contact groups.

\medterm{Secondary Brain Cancer} Alternate index label; see \textit{Brain metastases}.

\medterm{Secondary Hyperparathyroidism} Overactivity of the parathyroid glands in response to low calcium, abnormal phosphate, vitamin-D deficiency, or kidney disease. It can weaken bones, cause bone pain, and contribute to calcium deposits in blood vessels.

\medterm{Secondary Hypertension} High blood pressure caused by another condition, medicine, or substance. Possible causes include kidney disease, hormone disorders, sleep apnea, narrowed arteries, pregnancy, steroids, stimulants, or certain pain medicines.

\medterm{Secondary Hypothyroidism} Low thyroid-hormone production because the pituitary or nearby brain area does not provide enough thyroid-stimulating signal. It may cause tiredness, feeling cold, constipation, dry skin, slow heart rate, or menstrual changes.

\medterm{Secondary Immune Response} A faster and stronger defense when the body meets a germ or foreign substance again. Memory immune cells from the first exposure respond quickly and usually produce antibodies and other defenses more effectively.

\medterm{Secondary Immunodeficiency} Reduced immune function caused by an acquired condition or exposure rather than a primary genetic disorder. Causes
include HIV infection, cancer, malnutrition, medicines, and chronic disease; the
manifestations depend on which immune functions are impaired.

\medterm{Secondary Infections} Infections that develop during or after another disease, injury, or treatment because of tissue damage, altered defenses, or immunosuppression.


\medterm{Secondary Osteoporosis} Bone loss caused by another illness, medication, or prolonged immobility.

\medterm{Secondary Parkinsonism} Secondary parkinsonism is slowed movement, rigidity, tremor, or postural instability caused by another condition, medicine, toxin, vascular injury, or structural brain disease.

\medterm{Secondary Prevention} Finding and treating a disease early to stop it worsening or causing complications. Examples include screening for some cancers, checking blood pressure, testing for early diabetes, and treating an abnormality before it produces serious symptoms.

\medterm{Secondary Sexual Characteristics} Physical changes that develop during puberty and are not directly part of the reproductive organs, such as breast development, facial or body hair, voice changes, and changes in body shape.

\medterm{Secondary Progressive Multiple Sclerosis} A phase of multiple sclerosis in which disability gradually worsens over time after an earlier relapsing pattern. Relapses may still occur, but progressive walking, strength, vision, bladder, thinking, or other problems can develop even between attacks.

\synonyms
SPMS; Secondary-progressive MS

\medterm{Secondary Syphilis} The stage of syphilis that occurs weeks or months after the first sore when the bacteria spread through the body. It may cause a widespread rash, including the palms and soles, swollen glands, fever, tiredness, mouth patches, or moist lesions and requires treatment.

\medterm{Secondary Systemic Amyloidosis} Secondary systemic amyloidosis results when chronic inflammation or infection drives production of serum amyloid A, which deposits in organs such as the kidneys, liver, and spleen.

\medterm{Secondary Tumor} A tumor formed when cancer cells spread from the original site to another part of the body.

\medterm{Secretin} A hormone released by the first part of the small intestine when acidic stomach contents enter it. Secretin prompts the pancreas and bile ducts to release alkaline fluid, helping neutralize acid and support digestion.

\medterm{Secretin Stimulation Test} A test in which secretin is given and changes in blood or digestive fluid are measured. It may assess pancreatic function or investigate a rare tumor that releases excess gastrin and stomach acid.

\medterm{Secretion} The release of a substance from cells or glands into the blood, a body passage, or outside the body. Examples include saliva, hormones, digestive juices, mucus, and substances removed into urine or bile.

\medterm{Secretory Diarrhea} Watery diarrhea caused when the intestine releases too much water and salt or absorbs too little. It may continue even without eating and can result from infection, medicines, bile-acid problems, hormone-producing tumors, or inflammation.

\medterm{Section} A defined part or subdivision of an organ, document, study, or procedure; in anatomy and imaging, the term identifies a specific slice or region examined.

\medterm{Section, Cesarean} Alternate index label; see \textit{Cesarean Section}.

\medterm{Sed Rate} A blood test measuring how quickly red blood cells settle, used as a nonspecific marker of inflammation.

\synonyms
Erythrocyte sedimentation rate; ESR

\medterm{Sedation} The use of medicines to make a person relaxed, drowsy, or less aware during a procedure while maintaining appropriate monitoring and breathing.

\medterm{Sedation Dentistry} The use of medicine during dental care to reduce fear, discomfort, or awareness. It may involve inhaled gas, tablets, or medicines through a vein; deeper sedation can affect breathing and requires trained staff and continuous monitoring.

\medterm{Sedation Levels} The depth of medicine-induced drowsiness, ranging from minimal relaxation with normal responses to deep sedation with limited response and possible breathing or airway support. Sedation can deepen unexpectedly, so appropriate monitoring is required.

\medterm{Sedative-Hypnotic} A medicine or substance that slows brain activity to reduce anxiety or cause sleep. Examples include benzodiazepines, barbiturates, and some sleep medicines. Too much can cause confusion, coma, or dangerously slow breathing, especially with alcohol or opioids.

\medterm{Sedative-Hypnotic Toxidrome} A pattern of poisoning caused by substances that slow brain activity. It may cause extreme sleepiness, slurred speech, poor coordination, slow breathing, coma, or loss of airway protection, especially when combined with alcohol or opioids.

\medterm{Sedatives} Medicines or substances that slow brain activity and produce calmness, sleepiness, or reduced anxiety. Excessive amounts can impair coordination, lower blood pressure, or dangerously slow breathing, especially with alcohol or opioids.
\medterm{Sedimentation Rate} The erythrocyte sedimentation rate, a nonspecific blood test measuring how quickly red cells settle in a tube and sometimes supporting assessment of inflammation when interpreted with other findings.
\medterm{Segmental Pressures} Blood-pressure readings taken at several points on the arms and legs to compare circulation. A large difference between levels can help locate a narrowed or blocked artery in peripheral artery disease.

\medterm{Segmentation} Repeated squeezing of short sections of the intestine that mixes food with digestive fluids and brings it against the intestinal lining for absorption. Unlike peristalsis, it mainly mixes contents rather than pushing them forward.

\medterm{Segmentectomy} Surgery removing one or more defined sections of a lung rather than an entire lobe. It may treat a small, localized lung cancer or another focal problem while preserving more lung tissue; bleeding, air leakage, infection, and reduced breathing capacity are possible.

\synonyms
Sublobar Resection; Anatomical Segmentectomy; Pulmonary Segmentectomy

\medterm{SEID} Alternate index label; see \textit{Myalgic encephalomyelitis/chronic fatigue syndrome (ME/CFS)}.

\medterm{Seizure} A temporary episode caused by abnormal electrical activity in the brain. It may cause shaking, stiffening, unusual sensations, staring, confusion, or loss of awareness and can be triggered by fever, illness, injury, medicines, or epilepsy.

\medterm{Seizure Disorder} Alternate index label; see \textit{Epilepsy}.

\medterm{Seizure Disorders} Conditions characterized by recurrent unprovoked seizures or an enduring predisposition to them, classified by seizure type and cause and treated according to epilepsy syndrome and risk.

\medterm{Seizure Febrile (Febrile Seizures)} Alternate index label for Febrile Seizures.

\medterm{Seizure Fever Induced (Febrile Seizures)} Alternate index label for Febrile Seizures.

\medterm{Seizure, Absence} Alternate index label; see \textit{Absence seizure}.

\medterm{Seizure, Febrile} Alternate index label; see \textit{Febrile seizure}.

\medterm{Seizure, Grand Mal} Alternate index label; see \textit{Tonic-clonic (grand mal) seizure}.

\medterm{Seizure, Petit Mal} Alternate index label; see \textit{Absence seizure}.

\medterm{Seizure, Temporal Lobe} Alternate index label; see \textit{Temporal lobe seizure}.

\medterm{Seizures} Seizures are transient episodes of abnormal excessive brain electrical activity causing changes in awareness, movement, sensation, behavior, or autonomic function.


\medterm{Selective Deficiency Of IgA} A usually inherited immune problem in which the antibody IgA is very low or absent while other antibodies are near normal. Some people have repeated sinus, chest, or bowel infections, allergies, or autoimmune disease; others have no symptoms.

\medterm{Selective Estrogen Receptor Modulator} A medicine that acts like estrogen in some tissues and blocks estrogen effects in others. Different medicines in this group are used for conditions such as osteoporosis, breast-cancer risk reduction, or breast cancer.

\medterm{Selective Estrogen Receptor Modulators} Medicines that act like estrogen in some tissues and block estrogen effects in others, depending on the organ.

\synonyms
SERMs

\medterm{Selective IgA Deficiency} An immune condition with very low or absent immunoglobulin A, an antibody protecting mucus-lined surfaces. Some people have repeated respiratory or intestinal infections, allergies, autoimmune disease, or no symptoms.

\medterm{Selective Antibody Deficiency With Normal Immunoglobulins} An immune problem in which the total amount of antibodies is normal but the body does not make enough protective antibodies against some germs after vaccination. Repeated ear, sinus, or chest infections may occur.

\medterm{Selective Mutism} A childhood anxiety disorder in which a person speaks comfortably in some settings but repeatedly cannot speak in others, such as school or unfamiliar social situations. It lasts at least one month and can interfere with learning, relationships, or daily life.


\medterm{Selective Serotonin Reuptake Inhibitors} Antidepressants that increase serotonin signaling by slowing its reabsorption by nerve cells. They treat depression and several anxiety disorders, but effects take time and may include nausea, sexual problems, sleep changes, or withdrawal symptoms.

\synonyms
SSRIs

\medterm{Selective Serotonin-Reuptake Inhibitors (SSRIs)} Antidepressants that increase the availability of serotonin, a chemical used for communication between nerve cells. They treat depression and anxiety disorders; nausea, sleep changes, sexual problems, and withdrawal symptoms can occur.
\medterm{Selegiline} A medicine that blocks monoamine oxidase B, used for Parkinson disease and sometimes depression; interactions require caution.
\medterm{Selenium} A mineral needed for antioxidant enzymes, thyroid function, reproduction, and immune activity. Deficiency is uncommon in many regions, while excessive supplements can cause hair or nail changes, stomach upset, nerve problems, or toxicity.

\medterm{Selenium In Diet} Selenium in diet provides a trace mineral needed for antioxidant enzymes, thyroid hormone metabolism, and reproduction. Excess supplements can cause toxicity, so intake should remain within recommended limits.

\medterm{Selenium Sulphide} A skin-applied medicine that reduces certain fungi and excess scaling. It is used for dandruff and seborrheic dermatitis and may cause irritation, dryness, discoloration, or an unpleasant smell if misused.
\medterm{Selenium-75} A radioactive form of selenium used mainly as a calibration or testing source in imaging and measurement equipment. It releases radiation as it decays and has a half-life of about 120 days.

\medterm{Self Catheterization - Female} Intermittent insertion and removal of a catheter through the female urethra to empty the bladder when normal urination is inadequate, reducing retention and protecting kidney function.

\medterm{Self Catheterization - Male} Intermittent insertion and removal of a catheter through the male urethra to drain the bladder when neurologic, obstructive, or other conditions prevent adequate emptying.

\medterm{Self Gratification Sexual (Masturbation)} Alternate index label for Masturbation.

\medterm{Self Harming} Alternate index label; see \textit{Self-injury/cutting}.

\medterm{Self-Diagnosing Body Dysmorphic Disorder} Alternate index label for BDD.

\medterm{Self-Efficacy} A person's confidence in their ability to manage a health condition, follow a treatment plan, and handle symptoms. It can influence participation in care and health outcomes.

\medterm{Self-Harm} Deliberate injury or poisoning of oneself. It may be related to overwhelming emotions, mental illness, trauma, or suicidal thoughts, and can include cutting, burning, hitting, or taking too much medicine. Every episode deserves compassionate safety assessment.

\medterm{Self-Help Group} A peer group in which members support one another while changing a harmful behavior or coping with a problem.

\synonyms
Support group

\medterm{Self-Injury/Cutting} Deliberate injury to one's own body, often used to cope with overwhelming emotions and requiring compassionate assessment of safety and mental health.

\synonyms
Self Harming

\medterm{Self-Management} Actions a person takes to monitor and manage a health condition, including treatment use, symptom monitoring, and lifestyle measures.

\medterm{Self-Poisoning} Intentional ingestion or exposure to a potentially harmful substance, often requiring urgent medical and mental-health assessment.
\medterm{Self-Testing For COVID-19} Testing a sample collected at home, usually from the nose, to check for COVID-19 infection. A negative result does not always rule out infection, especially early in illness.

\medterm{Sella Turcica} A saddle-shaped hollow in a bone at the base of the skull that holds the pituitary gland. The pituitary controls many hormones, so changes in this area can affect growth, reproduction, metabolism, and the body's response to stress.
\medterm{Semaglutide} Semaglutide is a glucagon-like peptide-1 receptor agonist used for selected patients with type 2 diabetes or chronic weight management. It slows gastric emptying and reduces appetite and glucose levels; nausea, gallbladder disease, pancreatitis risk, and contraindications require review.

\medterm{Semen} The fluid released during ejaculation, containing sperm from the testes mixed with fluids from the seminal vesicles, prostate, and other glands. It carries and supports sperm; changes in amount, color, blood, or pain may need assessment.

\medterm{Semen Analysis} A laboratory examination measuring semen volume and sperm number, movement, shape, and vitality. It is a common first test when assessing fertility; results vary, so an abnormal result is often repeated and interpreted with the couple’s history.

\medterm{Semen Blood (Blood In Semen)} Alternate index label for Blood In Semen.

\medterm{Semen DNA Analysis} A test that examines genetic material inside sperm, usually for selected infertility or repeated-pregnancy-loss evaluations. It may detect increased DNA damage, but its usefulness and treatment implications depend on the clinical situation.

\medterm{Semi-Fowler Position} A semiupright position in which the head and upper body are elevated about 30 to 45 degrees.

\medterm{Semicircular Canals} Three curved, fluid-filled tubes in each inner ear that sense head rotation and help maintain balance. Problems affecting them can cause spinning dizziness, nausea, or difficulty keeping the eyes steady when the head moves.

\medterm{Semilunar Cartilages} The two crescent-shaped pads of tough cartilage inside the knee, commonly called the menisci. They cushion the joint, spread weight, and improve stability; a tear can cause pain, swelling, catching, or locking.
\medterm{Semilunar Valves} The aortic and pulmonary valves, which open during ventricular contraction and close during ventricular relaxation to prevent backflow into the ventricles.

\medterm{Seminal Vesicle} One of a paired set of male accessory reproductive glands behind the urinary bladder.
Its fructose-rich secretion contributes substantially to semen and joins the vas
deferens through the ejaculatory duct.

\medterm{Seminal Vesicles} Paired glands behind the bladder that produce much of the fluid in semen. Their secretions contain nutrients and other substances that support sperm movement and help form the ejaculate.

\medterm{Seminiferous Tubules} Tiny, coiled tubes inside the testes where sperm are made. Supporting cells nourish developing sperm and help create the environment needed for sperm production.
\medterm{Seminoma} A cancerous tumor arising from sperm-forming cells in the testis. It often grows slowly and is highly treatable with surgery, radiation, chemotherapy, or careful monitoring depending on its stage.
\medterm{Senescence} Changes and gradual loss of function associated with aging. In cells, senescence is a lasting state in which a cell stops dividing but remains active; an accumulation of such cells may contribute to aging and disease.

\synonyms
Biological aging

\medterm{Senile Lentigines} Flat, sharply defined brown or tan spots caused by chronic ultraviolet exposure, commonly appearing on sun-exposed skin with age.

\medterm{Senility} Senility is an outdated, imprecise term for cognitive decline in older age; modern assessment identifies specific causes such as dementia, delirium, depression, or medical illness.
\medterm{Senna} A plant-derived stimulant laxative containing sennosides, used short term for constipation and sometimes to empty the bowel before a procedure. It increases intestinal activity and may cause abdominal cramps or diarrhea; prolonged or excessive use can cause dependence or fluid and electrolyte disturbances.
\medterm{Sensate Focus Techniques} Structured, gradual touch exercises used in sex therapy to reduce performance pressure, improve body awareness, and rebuild comfortable intimacy. Partners agree on boundaries and initially avoid goals such as intercourse or orgasm.
\synonyms
Sensate focus

\medterm{Sensation} Awareness of changes such as touch, pain, temperature, sound, movement, or body position. Special nerve endings detect these changes and send messages to the brain, where they are experienced and interpreted.
\medterm{Sense} A sense is a capacity to detect and interpret stimuli such as touch, temperature, pain, sound, light, smell, taste, position, or internal body signals.

\medterm{Sensitisation} A process in which repeated exposure makes the body’s immune, nerve, or behavioral response stronger or easier to trigger. Sensitization can contribute to allergy, chronic pain, and some learned responses.
\medterm{Sensitivity} In testing, the proportion of people who truly have a condition and receive a positive result. In other settings, it means how strongly a person, cell, or instrument responds to a stimulus.
\medterm{Sensitivity Analysis} Sensitivity analysis tests how much a study's conclusion changes when assumptions, measurements, or input values are varied.

\medterm{Sensitivity, Gluten (Nonceliac Gluten Sensitivity (Intolerance))} Alternate index label for Nonceliac Gluten Sensitivity (Intolerance).

\medterm{Sensitized} Having antibodies against human-cell markers called HLA. This can happen after pregnancy, transfusion, or transplantation and may make it harder to find a suitable organ donor or increase the risk of transplant rejection.

\medterm{Sensorimotor Polyneuropathy} Damage to many nerves causing both altered sensation and weakness, usually starting in the feet or hands. Numbness, burning pain, poor balance, reduced reflexes, and muscle wasting may occur; diabetes is one possible cause.

\medterm{Sensorineural Deafness} Sensorineural deafness is hearing loss from dysfunction of the inner ear, auditory nerve, or central pathways rather than blockage of sound conduction.

\medterm{Sensorineural Hearing Loss} Hearing loss caused by damage to the inner ear or the nerve carrying sound to the brain. It may cause difficulty understanding speech, especially in noise, and can result from aging, loud noise, infection, inherited conditions, or medicines.

\medterm{Sensorium} A person’s level of consciousness, awareness, attention, and understanding of surroundings. Changes may include confusion, drowsiness, agitation, disorientation, or reduced responsiveness and can signal illness, medicines, intoxication, or brain injury.

\medterm{Sensory} Relating to sensation or to nerves and pathways that carry information such as touch, pain, temperature, position, hearing, or vision to the brain.
\medterm{Sensory Cortex} Areas of the brain that receive and interpret information about touch, pain, temperature, pressure, and body position. Damage can cause numbness, altered sensation, or difficulty recognizing where a body part is.
\medterm{Sensory Deprivation} A major reduction in normal input from sight, sound, touch, movement, or other senses. Prolonged deprivation can affect mood, thinking, sleep, orientation, and perception, especially in isolated or vulnerable people.
\medterm{Sensory Level} A clear line on the body below which feeling is reduced or changed. It can suggest a problem in the spinal cord, especially when accompanied by weakness, altered reflexes, or bladder and bowel problems, and needs prompt assessment.

\medterm{Sensory Nerves} Nerves that carry information such as touch, pain, temperature, and position from the body to the spinal cord and brain. Damage may cause numbness, tingling, burning pain, or reduced awareness of injury.

\medterm{Sensory Neuron} A nerve cell that carries information from receptors toward the brain and spinal cord. It helps detect touch, pain, temperature, pressure, body position, sound, light, smell, and taste.

\medterm{Sensory Processing Disorder} Difficulty noticing, interpreting, or responding to sound, touch, movement, light, or body position. It may affect behavior, learning, coordination, or daily activities and can occur with other developmental conditions; assessment and support are individualized.

\medterm{Sensory Receptors} Specialized cells or nerve endings that detect stimuli such as touch, temperature, pain, light, sound, or chemical changes.

\medterm{Sentinel Event} An unexpected healthcare event that causes death or serious physical or psychological harm, or creates a major risk of such harm. It requires prompt investigation and system changes to reduce the chance of recurrence.
\medterm{Sentinel Injury} An unexpected injury in an infant who cannot yet move independently, such as bruising, a mouth injury, or a fracture. It may be an early sign of serious harm and requires careful examination, explanation, safeguarding assessment, and documentation.

\medterm{Sentinel Lymph Node} The first lymph node or group of nodes most likely to receive fluid draining from a tumor. Removing and examining them can show whether cancer has begun to spread and help determine its stage.

\medterm{Sentinel Lymph Node Biopsy For Melanoma} Sentinel lymph-node biopsy identifies and removes the first draining lymph node or nodes from a melanoma site to detect microscopic spread. The result helps stage disease and guide further treatment; it does not remove all regional nodes routinely.

\medterm{Sentinel Lymph-Node Biopsy} Surgical removal and pathologic examination of the first draining lymph node or nodes from a primary tumor. It is used for selected melanomas, breast cancers, and other tumors to assess regional metastatic risk and guide staging.

\medterm{Sentinel Node Biopsy} Alternate index label for sentinel lymph-node biopsy: removal and examination of the first draining node or nodes to which a tumor is likely to spread, used to assess regional metastasis.

\medterm{Sentinel Surveillance} Ongoing collection of health information from selected clinics, laboratories, or communities to track disease trends. It can provide early warning of outbreaks but does not count every case in the population.

\medterm{Separation Anxiety} Separation anxiety is developmentally excessive fear or distress when apart from an attachment figure, causing avoidance, physical complaints, sleep problems, or impairment. Treatment may include gradual exposure, family-based support, cognitive behavioral therapy, and selected medication.

\medterm{Separation Anxiety Disorder} Excessive and persistent fear about being apart from a parent, partner, or other attachment figure. It may cause repeated reassurance-seeking, refusal to leave home, nightmares, physical symptoms, or difficulty attending school or work.

\medterm{Separation Anxiety In Children} Developmentally excessive distress when separated from an attachment figure, with worry, refusal, physical symptoms, nightmares, or school avoidance that impairs daily life.

\medterm{Separation Of The Iris} A tear that pulls the colored part of the eye away from its normal attachment, usually after a blunt injury. It may cause double vision, glare, an irregular pupil, bleeding, or increased eye pressure and needs eye-specialist assessment.


\medterm{Sepsis} A life-threatening condition in which the body's response to an infection damages its own organs. It may cause confusion, very fast breathing or heartbeat, low blood pressure, reduced urine, or severe weakness and needs emergency treatment.

\synonyms
Septicemia (Sepsis)

\medterm{Sepsis Syndrome} A range of severe responses to infection, from infection-related organ dysfunction to septic shock and failure of several organs. The condition can worsen quickly and requires urgent assessment and treatment.

\medterm{Septal Defect} An abnormal opening in the septum between the heart's chambers that permits blood to pass between the right and left sides of the heart.
\medterm{Septal Defect, Atrioventricular} Alternate index label; see \textit{Atrioventricular canal defect}.

\medterm{Septal Defect, Ventricular} Alternate index label; see \textit{Ventricular septal defect}.

\medterm{Septal Hematoma} A collection of blood inside the wall dividing the two sides of the nose, usually after an injury. It causes blockage and painful swelling and needs urgent drainage because untreated blood can destroy cartilage and deform the nose.

\medterm{Septate Uterus} A congenital difference in which a band of tissue partly or completely divides the inside of the uterus. Many people have no symptoms, but it can increase miscarriage or some pregnancy risks; imaging confirms it, and selected cases are treated by removing the dividing tissue.

\synonyms
Uterine Septum

\medterm{Septic} Relating to a serious infection or to harmful germs in a normally sterile body site. For example, septic arthritis is an infected joint, while septic shock is life-threatening organ and circulation failure caused by an overwhelming infection.

\medterm{Septic Abortion} A miscarriage or induced pregnancy loss complicated by infection of the uterus or surrounding tissues. Fever, chills, worsening abdominal pain, foul discharge, fast heartbeat, or faintness may occur; it is an emergency requiring antibiotics, supportive care, and removal of infected retained tissue when necessary.

\synonyms
Septic Miscarriage; Infected Abortion

\medterm{Septic Arthritis} Infection of a joint, usually bacterial, causing acute pain, swelling, restricted movement, and systemic illness.

\synonyms
Arthritis Infectious (Septic Arthritis)

\medterm{Septic Embolism} A piece of infected material travels through the bloodstream and blocks a smaller blood vessel, spreading infection to another organ. It often starts with an infected heart valve and can cause abscesses or tissue damage.

\medterm{Septic Pulmonary Infarction} Death of part of the lung after infected clots or other infected material block small lung arteries. It may cause fever, chest pain worse with breathing, cough, coughing blood, and lung abscesses, and requires urgent treatment of the infection.

\medterm{Septic Shock} A life-threatening complication of sepsis in which infection causes dangerously low blood pressure and poor blood flow to organs, even after initial fluids. It requires immediate hospital treatment with antibiotics, fluids, and medicines to support circulation.
\medterm{Septicaemia} An older term for bloodstream infection, often used to describe bacteremia with systemic illness.
\medterm{Septicemia} Septicemia is an older term for serious bloodstream infection; modern assessment distinguishes bacteremia from sepsis, which is life-threatening organ dysfunction caused by infection.

\medterm{Septicemia (Sepsis)} Alternate index label for Sepsis.

\medterm{Septo-Optic Dysplasia} A condition present from birth involving underdevelopment of one or both optic nerves, parts of the brain’s midline, and sometimes the pituitary gland. It may cause poor vision, involuntary eye movements, seizures, developmental differences, or hormone deficiencies.

\medterm{Septoplasty} Septoplasty is surgery to straighten or reshape the nasal septum to improve airflow or address structural obstruction.


\medterm{Septum} A wall or partition separating two spaces or tissues, such as the nasal septum or the septum between heart chambers.
\medterm{Septum Deviated (Deviated Septum)} Alternate index label for Deviated Septum.

\medterm{Septum, Heart} The wall separating the right and left chambers of the heart, consisting of atrial and ventricular portions and preventing normal mixing of oxygenated and deoxygenated blood.

\medterm{Sequela} A lasting consequence or condition that remains after an earlier disease, injury, infection, or treatment has ended, such as paralysis after stroke or scarring after infection.

\medterm{Sequelae} Lasting effects or complications that remain after an illness or injury. Examples include weakness after a stroke, breathing problems after severe infection, or pain after an injury; sequelae may be physical, mental, emotional, or social.
\medterm{Sequential Therapy} Treatment in which different medicines or interventions are given one after another rather than at the same time. The sequence may improve safety, effectiveness, or tolerance.

\medterm{Sequester} To sequester means to isolate or set apart; in medicine, sequestered tissue or material is separated from its normal circulation or surroundings.

\medterm{Sequestrum} A piece of dead bone that has separated from healthy bone, usually because of a long-lasting bone infection. It can prevent healing, cause drainage or pain, and may need surgery after infection is treated.
\medterm{Ser} Ser is the standard abbreviation for serine, an amino acid incorporated into proteins and involved in metabolism and cell signaling.


\medterm{Serine} Serine is a nonessential amino acid used in protein synthesis, one-carbon metabolism, membrane formation, and signaling.

\medterm{Serious Diseases And Health Problems (Symptoms Of Serious Diseases And Health Problems)} Alternate index label for Symptoms of Serious Diseases and Health Problems.

\medterm{Serious Adverse Event} An unwanted medical event that results in death, is life-threatening, requires or prolongs hospitalization, causes major disability, or causes a birth defect. It may occur during treatment without being caused by the treatment.

\medterm{SERMs} Selective estrogen receptor modulators that have estrogen-like effects in some tissues and blocking effects in others.

\synonyms
Selective estrogen receptor modulators

\medterm{Seroconversion} Development of detectable antibodies in blood after infection, vaccination, or another immune-triggering event, showing a changed immune status.
\medterm{Serodiagnostic} Relating to diagnosis using antibodies, antigens, or other markers measured in blood serum. These tests can show infection, immunity, inflammation, or other conditions but must be interpreted with timing and clinical findings.

\medterm{Serologic Test} A blood test that detects antibodies, antigens, or other immune markers to help diagnose infection, assess immunity, or monitor response to treatment.

\medterm{Serology} The laboratory evaluation of serum for antibodies, antigens, or other immune markers.
Results may indicate current infection, previous exposure, vaccination, autoimmune
reactivity, or passive antibody transfer and must be interpreted with timing and
clinical context.

\medterm{Serology For Brucellosis} Blood testing for antibodies to Brucella species, interpreted with exposure history and symptoms because cross-reactions and early false-negative results can occur.

\medterm{Seroma} A pocket of clear body fluid that forms near a surgical incision or injured tissue. It may cause swelling, pressure, or discomfort and often resolves, though large or painful collections may need drainage.

\medterm{Seronegative Arthritis} Inflammatory arthritis in which common blood tests for rheumatoid factor and anti-CCP antibodies are negative. It includes several diseases with different patterns, such as spine, skin, bowel, or eye involvement, so symptoms and examination guide diagnosis.

\medterm{Serosanguineous} Serosanguineous describes thin wound or body fluid containing both clear serum and a small amount of blood, often appearing pink or pale red.

\medterm{Serositis} Inflammation of the smooth lining around the lungs, heart, or abdominal organs. It can cause sharp pain that worsens with breathing or movement and may cause fluid to collect around the affected organ.

\medterm{Serotonin} A chemical messenger used by nerve cells and other tissues. It helps regulate mood, sleep, appetite, bowel function, and pain; too much activity from medicines can cause serotonin syndrome.

\synonyms
5-hydroxytryptamine; 5-HT

\medterm{Serotonin Blood Test} A blood test measuring serotonin, used selectively in evaluating suspected serotonin-producing neuroendocrine tumors; many medicines and specimen factors affect results.

\medterm{Serotonin Syndrome} A potentially life-threatening reaction caused by too much serotonin activity, usually after combining or taking too much of certain medicines. Agitation, sweating, diarrhea, shaking, fever, muscle stiffness, or confusion may occur.

\medterm{Serotonin-Norepinephrine Reuptake Inhibitors} Antidepressants that increase the activity of serotonin and norepinephrine, chemical messengers between nerve cells. They treat depression, some anxiety disorders, and certain nerve or long-lasting pain conditions; nausea, sweating, sleep changes, and withdrawal symptoms can occur.

\medterm{Serotype} A distinct version of a microorganism or antigen identified by how antibodies react to it.
\medterm{Serous} Thin, watery, and usually clear, resembling blood serum. Serous fluid can occur normally or collect in a wound, body cavity, or blister.
\medterm{Serous Borderline Ovarian Tumor} An ovarian tumor whose cells look abnormal but have not invaded nearby supporting tissue. It may stay in the ovary or form deposits on its surface; treatment and follow-up depend on its extent and microscopic features.

\medterm{Serous Cyst Adenomas Pancreas (Pancreatic Cysts)} Alternate index label for Pancreatic Cysts.

\medterm{Serous Fluid Cytology} Examination of cells in fluid collected from around the lungs, heart, or abdominal organs. It can look for cancer cells, infection, or inflammation; additional tests may be needed when the result is unclear.
\medterm{Serous Membranes} Thin, smooth linings that cover organs and reduce friction inside closed spaces, including the linings around the lungs, heart, and abdominal organs.
\medterm{Serous Tubal Intraepithelial Carcinoma} An early, noninvasive cancerous change in cells at the end of a fallopian tube near the ovary. It is thought to be the starting point for many high-grade cancers of the ovary, tube, or abdominal lining.

\medterm{Serpentine Supravenous Hyperpigmentation} A pattern of dark, wavy skin discoloration that follows the paths of surface veins, usually on the upper chest or neck. It can occur with vein inflammation, poor blood flow, some medicines, or cancer and requires assessment of the underlying cause.

\medterm{Serrated Pathway} One route by which some colon cancers develop from flat or saw-toothed polyps. These polyps may be easy to miss during colonoscopy and can acquire cell changes that increase cancer risk, so removal and follow-up are important.

\medterm{Serrated Polyposis Syndrome} A disorder characterized by multiple serrated lesions or polyps in the colon and an increased risk of colorectal cancer, requiring colonoscopic surveillance and removal of significant lesions.

\medterm{Serration} A row of small tooth-like projections or notches along an edge, used to describe anatomy, cells, instruments, or imaging appearances.

\medterm{Serratus Anterior Muscle} A muscle along the side of the chest that holds the shoulder blade against the ribs and helps raise the arm. Injury to its nerve can make the shoulder blade stick out and weaken pushing or overhead movements.

\medterm{Sertoli-Cell Tumor} A tumor arising from Sertoli cells, which support sperm development in the testicle. It is often benign but may produce estrogen or other hormones and cause breast enlargement, infertility, or a testicular mass.
\medterm{Sertoli-Leydig Cell Tumor} A rare ovarian tumor that may produce male hormones, causing increased facial hair, acne, a deeper voice, or irregular periods. Treatment depends on the tumor's stage and features.

\medterm{Serum} The clear liquid left after blood has clotted. It contains water, salts, hormones, antibodies, nutrients, and waste products but lacks fibrinogen and most clotting proteins, making it useful for laboratory tests.
\medterm{Serum Concentration} The amount of a substance measured in the liquid part of clotted blood. Results may reflect when the sample was taken, how the substance is absorbed and removed, and whether the level is expected to help or harm.

\medterm{Serum Free Hemoglobin Test} A blood test measuring hemoglobin not contained within red cells, used to support evaluation of intravascular hemolysis, transfusion reactions, or mechanical red-cell destruction.

\medterm{Serum Glutamic Oxaloacetic Transaminase} SGOT, now called AST, is an enzyme released by injured liver, heart, skeletal muscle, or other cells; its level must be interpreted with other findings.

\medterm{Serum Glutamic Pyruvic Transaminase} SGPT, now called ALT, is an enzyme concentrated in liver cells and commonly measured as a marker of hepatocellular injury.

\medterm{Serum Herpes Simplex Antibodies} A blood test for antibodies against herpes simplex virus types 1 or 2. IgG usually indicates previous infection, while IgM is less reliable; results do not show the infection site or prove that a current lesion is caused by herpes.

\medterm{Serum Iron Test} A blood test measuring circulating iron bound mainly to transferrin; it is interpreted with ferritin, transferrin saturation, and blood counts rather than used alone to diagnose iron deficiency.

\medterm{Serum Osmolal Gap} The difference between the measured concentration of dissolved particles in blood and the calculated value. A high gap may suggest an unmeasured substance, including a toxic alcohol, but other illnesses and laboratory factors can also affect it.

\medterm{Serum Phenylalanine Screening} A blood test measuring phenylalanine, used in newborn screening and follow-up of phenylketonuria or other disorders of phenylalanine metabolism.

\medterm{Serum Progesterone} A blood test measuring progesterone, a hormone made mainly after ovulation and during pregnancy. The result can help assess ovulation, ovarian function, fertility treatment, or early pregnancy, but interpretation depends strongly on timing and pregnancy status.

\medterm{Serum Protein Electrophoresis} A blood test that separates proteins in serum into groups. It can show abnormal antibody production or low protein levels and is often followed by immunofixation when a plasma-cell disorder is suspected.

\medterm{Serum Sickness} A delayed immune reaction, usually one to three weeks after exposure to certain medicines, injected proteins, or antiserum. It can cause fever, rash or hives, painful swollen joints, enlarged lymph nodes, and sometimes organ problems.
\medterm{Serum Therapy} Prevention or treatment using antibody-containing plasma, immunoglobulin, antitoxin, or antivenom. It provides immediate but usually temporary protection against a specific infection or toxin; allergic reactions and other infusion-related problems are possible.
\medterm{Serving Size} Serving size is the standardized amount used to describe the nutrients or calories in food and may differ from the amount a person actually eats.

\medterm{Sesamoid Bones} Small bones embedded within or beside tendons, where they protect the tendon and improve leverage. Examples include the kneecap and the small bones beneath the big toe.
\medterm{Sesamoiditis} Painful inflammation of tissues around a sesamoid bone, most often beneath the big toe. Pain usually worsens with walking, running, dancing, or pressure on the front of the foot.

\medterm{Sesquipedalian} Sesquipedalian means characterized by very long words; clear healthcare communication favors language that patients can understand.

\medterm{Sessile} Attached directly by a broad base rather than a narrow stalk. A sessile polyp can be harder to remove completely and may require careful examination because its shape can hide abnormal tissue.
\medterm{Set} In exercise, a group of consecutive repetitions performed before a rest period. In medicine, “set” can also mean a prepared group of instruments, tests, or equipment.

\medterm{Sever Disease} A painful irritation of the growth area at the back of the heel in growing children, often during sports. Heel pain worsens with running or jumping and usually improves with activity adjustment, heel support, calf stretching, and time.
\medterm{Severe Acute Respiratory Syndrome} A serious contagious respiratory illness caused by the SARS-associated coronavirus, SARS-CoV. It usually begins with fever and may cause chills, muscle aches, headache, diarrhea, dry cough, and breathing difficulty. Most affected people develop pneumonia, which can progress to low blood oxygen, respiratory failure, and death. The name usually refers to the disease identified in the 2002--2004 outbreak; SARS-CoV-2 causes COVID-19, a different disease.

\medterm{Severe Obesity} Obesity with a Body Mass Index of 40 kg/m\(^2\) or higher, or 35 kg/m\(^2\) or higher together with a serious obesity-related health problem. Body Mass Index does not fully measure health.

\synonyms
Severe Acute Respiratory Syndrome (Severe Acute Respiratory Syndrome (SARS))

\medterm{Severe Acute Respiratory Syndrome (Severe Acute Respiratory Syndrome)} Alternate index label for Severe Acute Respiratory Syndrome.

\medterm{Severe Brain Injury, Coma} Alternate index label; see \textit{Coma}.

\medterm{Severe Combined Immunodeficiency} A group of inherited conditions in which important immune cells do not develop or work properly. Affected babies can develop severe, repeated infections early in life and need urgent specialist care, protective measures, and immune-restoring treatment.

\medterm{Severe Condition (Sever Condition)} Alternate index label for Sever Condition.

\medterm{Severe Dengue} The modern World Health Organization (WHO) classification term for the critical, life-threatening form of dengue virus infection, typically developing around the time of fever defervescence (days 3 to 7). It is characterized by one or more of three major clinical manifestations: severe plasma leakage resulting in hypovolemic shock (dengue shock syndrome) or respiratory distress from fluid accumulation, severe bleeding (such as massive gastrointestinal hemorrhage), or severe organ impairment including acute liver failure, encephalitis, or myocarditis.

\synonyms
Dengue Shock Syndrome; Severe Dengue Virus Infection; Critical Dengue

\medterm{Sevoflurane} An inhaled anesthetic gas used to start and maintain general anesthesia. It acts in the brain to produce unconsciousness and reduces awareness of pain; it can lower breathing and blood pressure and may trigger malignant hyperthermia in susceptible people.

\medterm{Sex Change} A nontechnical term for gender transition or medical gender-affirming care; respectful, current terminology should be used.
\medterm{Sex Chromosomes} The X and Y chromosomes, which contribute to sex development and carry many other genes.
\medterm{Sex Education} Teaching about bodies, relationships, consent, contraception, sexually transmitted infections, and sexual health. Accurate, age-appropriate education supports safer decisions, respectful relationships, prevention, and access to appropriate healthcare.
\medterm{Sex Headaches} Headaches brought on by sexual activity, including a gradually increasing ache or a sudden severe headache at orgasm. A first or thunderclap headache requires urgent assessment to exclude bleeding or other dangerous causes.

\medterm{Sex Hormones} Hormones involved in sexual development, reproduction, menstrual cycles, pregnancy, and related body functions. Major groups include estrogens, progesterone, and androgens, which also affect bones, muscles, mood, and tissues.
\medterm{Sex Therapy} Structured psychological or behavioral therapy for sexual concerns or dysfunction.
\medterm{Sex-Linked Dominant} A pattern of inheritance in which a disease-causing gene variant lies on a sex chromosome and one altered copy can cause the condition. The chance of passing it on differs for mothers and fathers and for children of different sexes.

\medterm{Sex-Linked Inheritance} Inheritance of a trait caused by a gene on a sex chromosome, commonly the X chromosome.
\medterm{Sex-Linked Recessive} A pattern in which a disease-causing gene variant, usually on the X chromosome, commonly causes disease in males with one altered copy. Females usually need two altered copies, while those with one may be carriers or have milder symptoms.

\medterm{Sexual Abuse} Nonconsensual sexual contact, exploitation, or behavior; it is a safeguarding and health emergency when immediate danger or injury exists.
\medterm{Sexual Assault} Any sexual contact or sexual act without freely given consent, including unwanted touching, penetration, coercion, or sexual activity when a person cannot consent. It can cause injury, infection, pregnancy, and psychological trauma and may require medical and safeguarding support.

\synonyms
Rape (Sexual Assault)

\medterm{Sexual Assault - Prevention} Sexual-assault prevention includes consent education, recognizing coercion, safer environments, bystander intervention, and access to confidential support; responsibility always lies with the person who commits the assault.

\medterm{Sexual Dysfunction} A group of conditions involving persistent or recurrent difficulty with sexual desire, arousal, erection, ejaculation, orgasm, or pain that causes clinically significant distress. Causes may be physical, psychological, medication-related, hormonal, or related to life circumstances.
\medterm{Sexual Dysfunction, Female} Alternate index label; see \textit{Female sexual dysfunction}.
\medterm{Sexual Dysfunction Disorders} Alternate index label for Sexual Dysfunction.

\medterm{Sexual Intercourse (Coitus)} Sexual activity involving penetration, most commonly penile-vaginal, but sometimes penile-anal or other forms depending on the people involved. It may be for pleasure, intimacy, or reproduction and requires freely given consent.

\synonyms
Coitus

\medterm{Sexual Masochism Disorder} A mental-health diagnosis involving repeated sexual arousal from being humiliated, beaten, bound, or otherwise made to suffer, when the pattern causes significant distress or disrupts daily life, or involves a person who has not freely consented.

\medterm{Sexual Obsession} Alternate index label; see \textit{Compulsive sexual behavior}.

\medterm{Sexual Rehabilitation} Assessment and treatment to help restore or adapt sexual function after illness, injury, surgery, disability, or medicine effects. It may include education, counseling, pelvic-floor therapy, devices, or medicines.

\medterm{Sexual Sadism} A recurrent sexual interest in another person’s physical or psychological suffering or humiliation. The interest alone is not a mental disorder; acts without consent are abuse.

\medterm{Sexual Sadism Disorder} A mental-health diagnosis involving repeated sexual arousal from another person’s physical or emotional suffering, when it causes clinically significant distress or impairment or involves acting with someone who has not freely consented.

\medterm{Sexual Violence} Sexual violence is any sexual act, attempt, contact, or coercion without freely given consent and can cause physical injury, infection, pregnancy, trauma, and long-term health effects.

\medterm{Sexually Transmitted Disease Treatments} Alternate index label for Sexually Transmitted Infections.

\medterm{Sexually Transmitted Diseases} Alternate index label for Sexually Transmitted Infections.

\medterm{Sexually Transmitted Diseases (STDs In Women)} Alternate index label for Sexually Transmitted Infections.
\medterm{Sexually Transmitted Diseases (STDs)} Alternate index label for Sexually Transmitted Infections.

\medterm{Sexually Transmitted Diseases And Pregnancy (STDs)} Alternate index label for Sexually Transmitted Infections.

\medterm{Sexually Transmitted Infection} Alternate index label for Sexually Transmitted Infections.

\medterm{Sexually Transmitted Infections} Infections caused by bacteria, viruses, fungi, or parasites that spread mainly through sexual contact, including vaginal, anal, or oral sex and genital skin-to-skin contact. Some can also spread through blood exposure or from a pregnant person to a baby. Many cause no symptoms, especially early in the infection.

\synonyms
STIs
STDs

\medterm{Sezary Disease} Alternate index label; see \textit{Sezary syndrome}.

\medterm{Sezary Syndrome} An uncommon form of cutaneous T-cell lymphoma involving widespread red skin and abnormal lymphocytes in the blood.

\synonyms
Sezary Disease

\medterm{SéZary Syndrome} An uncommon form of skin lymphoma in which abnormal immune cells are found in the skin and blood. It usually causes widespread red, scaly, itchy skin and enlarged lymph nodes and requires specialist care.

\medterm{Sgarbossa Criteria} A set of electrocardiogram findings used to look for a heart attack when an existing conduction problem or pacemaker makes the tracing difficult to read. The findings can support, but cannot prove, an acute blocked heart artery.

\synonyms
Modified Sgarbossa Criteria; Smith-Modified Sgarbossa Rule

\medterm{Shadowing Artifact} An ultrasound artifact in which a strongly attenuating or reflecting structure produces a dark area behind it because little sound reaches the tissues beyond. Stones, bone, and gas commonly cause shadowing.

\medterm{Shaken Baby Syndrome} An older term for abusive head trauma in an infant or young child caused by violent shaking, with or without impact. It can cause brain and eye injuries, seizures, breathing problems, vomiting, unusual sleepiness, or death and requires emergency protection and care.

\medterm{Shaken Baby Syndrome (Abusive Head Trauma)} Alternate index label for Abusive Head Trauma.

\medterm{Shaken Impact Syndrome} A term for abusive head trauma involving shaking with or without impact, especially in infants.
\medterm{Shampoo - Swallowing} Accidental swallowing of a small amount of shampoo usually causes mouth or stomach irritation; larger or concentrated exposures may cause vomiting or breathing problems and warrant poison-control advice.

\medterm{Shared Decision-Making} A collaborative process in which clinicians and patients combine best available evidence
with patient values, preferences, goals, and circumstances to choose among reasonable
options.

\medterm{Shark Attack} A shark attack can cause severe lacerations, bleeding, tissue loss, drowning, and shock; emergency rescue, hemorrhage control, and trauma care are critical.

\medterm{Sharp Force Trauma} An injury caused by a sharp or pointed object, such as a knife, broken glass, or needle. It may produce a cut or a deeper puncture and can damage blood vessels, nerves, organs, or bones; deep or heavily bleeding wounds need emergency care.


\medterm{Sharps Container} A strong, leak-resistant container for used needles, lancets, blades, and other sharp medical items. It helps prevent injury and infection and should be clearly labeled, kept upright, and closed or replaced before it becomes too full.

\medterm{Sharps Injury} A percutaneous or mucocutaneous injury caused by a needle or other sharp instrument that
may expose a person to blood or body fluids. Immediate wound care, incident reporting,
risk assessment, and pathogen-specific post-exposure management are required.

\medterm{Shave Biopsy} Removal of a raised or superficial skin lesion by shaving off a thin layer with a sterile blade. The tissue is examined under a microscope; bleeding, infection, scarring, or an inadequate sample can occur.

\medterm{Shaving Cream Poisoning} Poisoning after swallowing shaving cream or inhaling its propellant. Small accidental amounts often cause mild stomach upset, but larger exposures can cause vomiting, choking, drowsiness, or breathing problems; seek product-specific advice.

\medterm{Sheehan Syndrome} Damage to the pituitary gland after severe blood loss or very low blood pressure during childbirth. It can cause failure to produce breast milk, extreme tiredness, low blood pressure, menstrual or fertility problems, and deficiencies of several hormones.

\medterm{Sheehan Syndrome (Postpartum Pituitary Necrosis)} Permanent injury to the pituitary gland caused by severe bleeding or very low blood pressure during or after childbirth. Because the gland makes several hormones, symptoms may include inability to breastfeed, low blood pressure, fatigue, and menstrual or fertility problems.

\synonyms
Postpartum Pituitary Necrosis

\medterm{Shellac Poisoning} Poisoning from swallowing or inhaling shellac, a coating product that may contain alcohol or other solvents. It can cause irritation, vomiting, dizziness, drowsiness, or lung injury after aspiration; do not induce vomiting.

\medterm{Shellfish Allergy} An immune reaction to proteins in crustaceans or mollusks that can cause hives, digestive or respiratory symptoms, or anaphylaxis.

\synonyms
Allergy, Shellfish

\medterm{Shellfish Poisoning} Illness caused by toxins accumulated in contaminated shellfish. Depending on the toxin, symptoms may include vomiting, diarrhea, tingling, weakness, confusion, breathing difficulty, or paralysis and require urgent medical advice.
\medterm{Shielding} Material placed between a person and a source of radiation to reduce exposure. The material depends on the radiation: lead or tungsten is used for many X-rays and gamma rays, while plastic or acrylic can be useful for beta radiation.

\medterm{Shift Work Sleep Disorder} A circadian rhythm sleep-wake disorder caused by work schedules that overlap the usual sleep period, producing insomnia, excessive sleepiness, or both.

\medterm{Shigella} A genus of bacteria causing shigellosis, an acute diarrheal illness that may include fever and bloody stools.
\medterm{Shigellosis} Shigellosis is intestinal infection caused by Shigella bacteria, producing diarrhea that may be bloody, fever, cramps, and dehydration.

\medterm{Shin Bone Fever} An older nonspecific term for painful inflammation or infection involving the shin or nearby tissues.

\medterm{Shin Splints} A common term for medial tibial stress syndrome, an overuse injury causing pain and tenderness along the inner border of the shinbone. It is associated with repetitive impact or a sudden increase in activity and involves irritation of tissues around the bone. It is distinct from a stress fracture and acute compartment syndrome, which are separate causes of lower-leg pain.

\synonyms
Medial Tibial Stress Syndrome


\medterm{Shin Spot} A shin spot is a nonspecific mark or lesion on the lower leg; its significance depends on appearance, duration, symptoms, injury, and associated disease.

\medterm{Shingles} Herpes zoster, reactivation of varicella-zoster virus in sensory nerves causing a painful one-sided blistering rash. Pain, burning, or tingling may begin before the rash. Eye or ear involvement can threaten vision or hearing, and persistent nerve pain after the rash heals is called postherpetic neuralgia; prompt assessment is important for facial or eye-area rashes.
\medterm{Shingles (Herpes Zoster)} Alternate index label for Herpes Zoster.


\medterm{Shingles And Pregnancy} Shingles is reactivation of varicella-zoster virus and usually causes a painful one-sided blistering rash. It is not the same as primary chickenpox exposure, but pregnancy requires prompt clinical advice about treatment, pain, and contact with people at risk.


\medterm{Shirodkar Cerclage} An operation that places a strong stitch high around the cervix to help prevent it opening too early during pregnancy. It is used in selected people at risk of late miscarriage or premature birth and can cause bleeding, infection, or injury to nearby organs.

\medterm{Shivering} Repeated, involuntary muscle contractions that generate heat, usually in response to cold or fever. Severe or unexplained shivering may occur with infection, medicine reactions, or other illness and should be assessed with the other symptoms.

\medterm{Shock} A life-threatening failure of circulation in which organs do not receive enough blood and oxygen. It can result from severe blood or fluid loss, heart failure, overwhelming infection, a severe allergic reaction, or an obstruction to blood flow. Signs may include confusion, faintness, cold or mottled skin, rapid breathing, little urine, or very low blood pressure and require emergency treatment.
\medterm{Shock Lung} An older term for acute respiratory distress syndrome (ARDS), a life-threatening form of lung inflammation in which fluid and injury make it difficult for oxygen to enter the blood.
\medterm{Shock Therapy} An imprecise term that may refer to electroconvulsive therapy or outdated shock treatments; the specific intervention should be named.
\medterm{Shock Wave Lithotripsy} A procedure that focuses repeated sound-pressure waves from outside the body onto a kidney or ureter stone to break it into smaller pieces that can pass in urine. Blood in the urine, pain, blockage, infection, or incomplete stone removal can occur.

\medterm{Shock, Secondary} Circulatory failure occurring as a consequence of another condition, such as infection, trauma, bleeding, heart failure, allergy, or endocrine crisis, with inadequate tissue perfusion.

\medterm{Shone Complex} A group of congenital left-sided obstructive heart lesions, classically including
supravalvular mitral ring, parachute mitral valve, subaortic stenosis, and coarctation
of the aorta. The combination and severity vary and may lead to heart failure or
systemic outflow obstruction.

\medterm{Short Arm Of A Chromosome} The shorter section of a chromosome, labeled p from the French word for “small.” Missing, extra, or rearranged DNA in this section can alter genes and cause disease.

\medterm{Short Bones} Small, roughly cube-shaped bones that provide stability and limited movement. The wrist and ankle contain most of the body's short bones.

\medterm{Short Bowel Syndrome} A condition in which the small intestine is shortened or damaged and cannot absorb enough nutrients from food. It most often occurs after surgery to remove part of the small intestine and may cause diarrhea, weight loss, and nutritional deficiencies.

\medterm{Short QT Syndrome} A rare inherited heart-rhythm disorder in which the heart's electrical recovery time is unusually short. It can cause atrial or ventricular arrhythmias, fainting, or sudden cardiac death; diagnosis requires specialist ECG and family-history assessment, and treatment may include an implantable defibrillator.

\synonyms
SQTS


\medterm{Short Philtrum} A short philtrum is reduced distance between the nose and upper lip and may be an isolated feature or part of a genetic, developmental, or prenatal-exposure syndrome.

\medterm{Short Stature} Height substantially below the expected range for age and sex, requiring assessment of growth pattern and cause.
\medterm{Short-Sight} Myopia, a refractive error in which distant objects are focused in front of the retina.
\medterm{Short-Term Memory} The ability to hold and use a small amount of information for seconds or minutes. It helps with following instructions, doing mental tasks, and remembering a recent conversation before information is stored more permanently or forgotten.

\synonyms
Working memory

\medterm{Shortness Of Breath} An uncomfortable awareness of difficult or insufficient breathing. See Dyspnea.

\medterm{Shotty} Shotty describes small, firm, discrete lymph nodes, often reactive after infection; persistent, enlarging, hard, or generalized nodes require evaluation.

\medterm{Shoulder} The area where the upper limb joins the body, including the shoulder blade, collarbone, and upper-arm bone. Its main joint allows a wide range of arm movement but is less stable and can be injured or dislocated.
\medterm{Shoulder Arthroscopy} A minimally invasive procedure using a camera and instruments through small incisions to diagnose and treat selected disorders of the shoulder joint and surrounding tissues.

\medterm{Shoulder Bursitis} Shoulder bursitis is inflammation of a fluid-filled sac around the shoulder, often associated with rotator-cuff disease or repetitive overhead activity. Pain with elevation and lying on the affected side is common; activity modification, therapy, medicines, injection, or treatment of the underlying cause may help.

\synonyms
Bursitis Shoulder (Shoulder Bursitis)

\medterm{Shoulder CT Scan} Computed tomography of the shoulder used to evaluate fractures, bone alignment, arthritis, tumors, infection, and complex bony anatomy.

\medterm{Shoulder Dislocation} Complete displacement of the upper arm bone from the shoulder socket, usually caused by a fall, collision,
or forceful twisting. It causes severe pain, deformity, and inability to move the arm; urgent reduction
and assessment for associated fractures or nerve and blood-vessel injury are required.

\medterm{Shoulder Dystocia} A childbirth emergency in which the baby’s head is born but a shoulder becomes stuck behind the mother’s pelvic bone. It requires immediate maneuvers because delay can cause lack of oxygen, fractures, or nerve injury in the baby and heavy bleeding in the mother.

\medterm{Shoulder Impingement} Pain associated with movement of the shoulder, often involving the rotator-cuff tendons or
subacromial tissues. Symptoms include pain with overhead activity and sometimes weakness;
assessment and treatment are individualized and may include exercise-based rehabilitation.

\medterm{Shoulder Instability} Excessive or recurrent movement of the shoulder joint caused by injury, laxity, or altered
joint anatomy. Associated lesions may include a Bankart tear, Hill-Sachs lesion, or injury to the joint capsule or
ligaments.

\medterm{Shoulder MRI Scan} Magnetic resonance imaging of the shoulder used to evaluate rotator-cuff tendons, labrum, cartilage, ligaments, muscles, bone marrow, and joint fluid.

\medterm{Shoulder Pain} Pain in or around the shoulder caused by injury, tendon or bursa disorders, arthritis, nerve disease, or
referred pain. Sudden weakness, deformity, chest symptoms, fever, or pain after major trauma
requires prompt evaluation.

\medterm{Shoulder Replacement} Surgery replacing damaged surfaces of the shoulder with artificial parts to relieve severe pain and improve movement. The implant may replace the joint in the usual or reverse configuration; infection, stiffness, dislocation, loosening, or nerve injury can occur.




\medterm{Shulmans Syndrome (Eosinophilic Fasciitis)} Alternate index label for Eosinophilic Fasciitis.

\medterm{Shunt} A passage that diverts blood, cerebrospinal fluid, urine, or another fluid from one channel to another.
\medterm{Shunt (Vascular)} A surgically created or implanted pathway that redirects blood or another fluid from one vessel or body space to another.

\synonyms
Vascular shunt; blood-vessel shunt

\medterm{Shwachman Syndrome} An inherited multisystem disorder usually involving exocrine pancreatic insufficiency,
bone-marrow dysfunction, and skeletal abnormalities. Children may have steatorrhea, poor
growth, recurrent infections, neutropenia, anemia, or increased risk of myeloid
malignancy.

\medterm{Shy-Drager Syndrome} Alternate index label; see \textit{Multiple system atrophy}.

\medterm{SI Units} The International System of Units used for standardized measurements, including meter, kilogram, second, mole, and pascal.
\medterm{SIADH} A condition in which the body releases too much antidiuretic hormone, causing it to retain water and dilute the blood sodium level. It may cause nausea, headache, confusion, seizures, or coma and requires treatment of the cause and sodium abnormality.

\medterm{Sialadenitis} Inflammation or infection of a salivary gland, causing painful swelling near the cheek or under the jaw, often worse during meals. Causes include bacteria, viruses, dehydration, or a blocked duct; persistent fever or pus needs medical care.

\medterm{Sialagogues} Substances or medicines that increase saliva production. They may help some people with dry mouth, but treatment also considers hydration, medicines, oral care, and the cause of reduced saliva.
\medterm{Sialogram} An imaging study of a salivary gland and its ducts after contrast is introduced through the duct opening, used to assess obstruction, stones, strictures, or chronic inflammation.

\medterm{Sialography} An imaging procedure that uses contrast material to show salivary glands and their ducts. It can help identify blockage, stones, narrowing, chronic inflammation, or tumors, although newer imaging is often preferred.

\medterm{Sialolithiasis} Formation of a stone in a salivary gland or its drainage tube, most often beneath the jaw. It can cause repeated swelling and pain during meals and may lead to blockage or infection.

\medterm{Salivary Gland Stones} Stones formed from mineral deposits in a salivary gland or its drainage tube. They can block saliva flow and cause swelling or pain, especially during meals; the medical term is sialolithiasis.

\medterm{Sialorrhea} Excess saliva or difficulty controlling and swallowing it, causing drooling. It commonly occurs with neurologic disease, swallowing difficulty, dental problems, or some medicines; treatment addresses the cause and may include swallowing therapy or selected medicines.
\medterm{Sialorrhoea} Excess saliva that pools in the mouth or dribbles out as drooling. It is often caused by difficulty controlling or swallowing saliva rather than making too much. It may occur with cerebral palsy, Parkinson disease, motor-neuron disease, or other conditions and can cause skin irritation or choking. Swallowing therapy and other treatments may help.
\medterm{Sialuria} A rare inherited disorder in which cells make and release too much sialic acid, a sugar used on cell surfaces. It may cause developmental delay, enlarged liver and spleen, distinctive facial features, and variable learning difficulties.

\medterm{Sibling} A person's brother or sister; family history from siblings can inform genetic and disease risk assessment.
\medterm{SIBO} Alternate index label; see \textit{Small intestinal bacterial overgrowth}.

\medterm{Sicca Syndrome} Persistent dryness of the eyes and mouth caused by reduced tears and saliva. It may result from Sjögren syndrome, medicines, dehydration, radiation, or other illnesses and can lead to eye irritation, mouth sores, tooth decay, or difficulty swallowing.


\medterm{Sick Sinus Syndrome} A problem with the heart’s natural pacemaker that causes the heartbeat to become too slow, pause, or alternate between fast and slow. It may cause dizziness, fainting, tiredness, palpitations, or reduced exercise ability.

\synonyms
Bradycardia-Tachycardia Syndrome

\medterm{Sickle Cell} An inherited group of blood disorders caused by abnormal hemoglobin. Red blood cells may become rigid and curved, causing anemia, severe pain, blocked blood flow, and damage to organs.

\medterm{Sickle Cell Anemia} An inherited blood disorder in which abnormal haemoglobin makes red blood cells stiff, sticky, and crescent-shaped. These cells break down easily and can block small blood vessels, causing anaemia, severe pain, stroke, lung problems, jaundice, and serious infections.

\medterm{Sickle Cell Anemia (Sickle Cell)} See \textit{Sickle Cell Anemia}.


\medterm{Sickle Cell Crisis} A sudden complication of sickle-cell disease, usually caused by sickled red cells blocking blood flow and causing severe pain. Other crises can damage the lungs, reduce blood-cell production, or trap blood in the spleen; urgent care may be needed.

\medterm{Sickle Cell Dactylitis} Painful swelling of the hands or feet caused by blockage of small blood vessels during a sickle-cell
episode, especially in infants and young children. It may be an early sign of sickle-cell disease
and requires appropriate clinical evaluation.

\medterm{Sickle Cell Disease} Alternate index label; see \textit{Sickle Cell}.

\medterm{Sickle Cell Test} A laboratory test detecting hemoglobin S or assessing sickling, used in screening and diagnosis of sickle-cell trait or sickle-cell disease; definitive testing identifies hemoglobin variants.

\medterm{Sickle Cell Trait} An inherited state in which a person has one gene for hemoglobin S and one usual hemoglobin gene. Most people have no symptoms and do not develop sickle cell disease, but rare problems can occur with severe dehydration, low oxygen, or extreme exertion.

\medterm{Sickle-Cell Anaemia} See \textit{Sickle Cell Anemia}.
\medterm{Sickness} A general term for illness or nausea; clinical assessment should identify the specific symptom or disease.
\medterm{Sickness Motion (Motion Sickness (Sea Sickness, Car Sickness))} Alternate index label for Motion Sickness (Sea Sickness, Car Sickness).

\medterm{Side Effect} An unintended effect that occurs while a medicine is used at a normal dose. It may be mild, such as nausea or drowsiness, or serious, such as bleeding, severe allergy, organ injury, or an abnormal heart rhythm.

\medterm{Side Effects} Unintended effects that occur while a medicine or treatment is used. They range from mild symptoms to serious harm; the risk depends on the treatment, dose, duration, other medicines, and the person’s health.

\medterm{Side-Effects} Unintended effects that occur while a medicine or treatment is used. They range from mild symptoms to serious harm; the risk depends on the treatment, dose, duration, other medicines, and the person’s health.
\medterm{Sideroblastic Anemia} Anemia caused by faulty hemoglobin production. Iron builds up inside developing red blood cells instead of being used properly; causes include inherited disorders, alcohol, some medicines or toxins, vitamin B6 deficiency, and bone-marrow disease.
\medterm{Siderosis} Accumulation of iron in body tissues, or lung disease from breathing iron-containing dust. It may cause no symptoms, but long-term dust exposure can produce cough, breathing difficulty, and permanent lung damage.
\medterm{SIDS} Sudden infant death syndrome, the sudden unexplained death of an infant younger than one year after thorough investigation.

\medterm{Sievert (Sv)} A unit used to estimate how much biological harm ionizing radiation could cause. It accounts for the type of radiation and the sensitivity of different tissues and is used when discussing radiation risk.

\synonyms
Sv

\medterm{Sight} The ability to see, involving the eyes and visual nervous system.
\medterm{Sigmoid} Sigmoid means S-shaped; the sigmoid colon is the curved section of large intestine between the descending colon and rectum.

\medterm{Sigmoid Colon} The curved, S-shaped section of the large intestine between the descending colon and rectum. It stores stool before bowel movements and is a common site for diverticular disease, twisting, and colorectal cancer.

\medterm{Sigmoid Volvulus} Twisting of the S-shaped lower colon around the tissue that supports it, blocking the bowel. It can cause severe abdominal swelling, pain, vomiting, and inability to pass stool or gas and needs urgent treatment because the blood supply may be cut off.

\medterm{Sigmoidoscopy} Endoscopic examination of the rectum and lower colon using a flexible or rigid sigmoidoscope. It can identify inflammation, bleeding, polyps, or tumors and may allow biopsy or treatment.

\medterm{Sign} An observable or measurable finding that may show disease or a change in health. Examples include a rash, fever, swelling, abnormal heartbeat, or high blood pressure; a sign differs from a symptom, which is felt and reported by the person.

\medterm{Sign Language} Sign language is a complete visual language using hand shape, movement, position, facial expression, and body movement rather than spoken sound.

\medterm{Signal Peptide} A short part of a newly made protein that acts like an address label, directing the protein to the cell membrane, outside the cell, or a particular internal compartment. It is often removed after delivery.

\medterm{Signal Transduction} The chain of events by which a hormone, nerve chemical, or medicine attaches to a cell and produces a response inside it. The response can change cell activity, movement, growth, survival, or gene use.

\medterm{Signal-To-Noise Ratio} A comparison between the useful information in a test or image and the unwanted background variation. A higher ratio usually produces a clearer image or more reliable measurement.
\medterm{Signature} A signature is a person’s identifying mark or electronic confirmation of approval, consent, authorship, or verification.

\medterm{Signature Strengths} Personal character strengths that a person values and naturally uses.

\medterm{Signet Ring Cell Carcinoma} A cancer whose cells contain mucus that pushes the cell nucleus to one side, giving a ring-like appearance under a microscope. It most often begins in the stomach or colon and may grow widely or aggressively.

\medterm{Signs Of An Asthma Attack} An asthma attack may cause worsening wheeze, cough, chest tightness, breathlessness, rapid breathing, difficulty speaking, or reduced peak flow; severe symptoms need emergency care.

\medterm{Sildenafil Citrate} The citrate salt of sildenafil, a phosphodiesterase-5 inhibitor used for erectile dysfunction and pulmonary arterial hypertension.
\medterm{Silent Heart Attack} A heart attack that causes no recognized pain or typical warning signs. It may be discovered later through an electrocardiogram or imaging and still damages heart muscle and increases future risks.

\synonyms
Unrecognized myocardial infarction

\medterm{Silent Ischemia} Reduced blood flow to heart muscle without typical chest pain or other recognized symptoms. It may still damage the heart and is detected through testing rather than symptoms alone.

\synonyms
Asymptomatic ischemia

\medterm{Silent Mutation} A DNA change that usually does not alter the amino acid in the protein made from the gene. It may still affect how the gene is copied, how long its RNA lasts, or how much protein is produced.

\medterm{Silent Sinus Syndrome} Silent sinus syndrome is painless inward collapse of the maxillary sinus and floor of the orbit, often causing facial asymmetry, enophthalmos, or a sunken eye. Imaging confirms the diagnosis, and endoscopic treatment restores sinus ventilation when indicated.

\medterm{Silent Thyroiditis} Silent thyroiditis is painless autoimmune inflammation of the thyroid that may briefly release excess hormone before causing temporary hypothyroidism and eventual recovery.

\medterm{Silicate} A salt or ester containing silicon and oxygen, found in minerals and industrial materials; inhalation of some silicate dusts can irritate the lungs or contribute to pneumoconiosis.

\medterm{Silicon Dioxide} A mineral found in soil, rock, dust, and industrial materials. Breathing crystalline forms over time can scar the lungs, cause silicosis, and increase the risk of lung cancer.

\medterm{Silicone} A family of synthetic silicon-containing polymers used in medical devices, implants, sealants, and consumer products; medical risks depend on the product and exposure.

\medterm{Silicone Elastomer} A flexible, rubber-like silicone material used in implants, tubing, dressings, and other devices; complications may include infection, leakage, allergy, or device failure.

\medterm{Silicone Implants} Medical implants containing silicone gel or elastomer used for breast reconstruction or augmentation; complications may include rupture, capsular contracture, infection, malposition, or need for revision.

\medterm{Silicone Mastitis} A granulomatous inflammatory reaction to silicone released from a breast implant or injected silicone. It may cause nodules, pain, lymphadenopathy, or imaging findings that mimic malignancy.

\medterm{Silicone Oil Hyperoleon} A complication after eye surgery in which silicone oil breaks into tiny droplets and moves into the front of the eye. It can raise eye pressure, irritate the eye, damage the cornea, and sometimes requires removal or washing of the oil.

\medterm{Silicones} Synthetic polymers used in medical devices, implants, dressings, and other products; adverse effects depend on formulation and site.
\medterm{Silicosis} A permanent lung-scarring disease caused by breathing crystalline silica dust. It can cause cough, breathlessness, reduced lung function, tuberculosis risk, and progressive respiratory disability after workplace or environmental exposure.
\medterm{Silk} Silk is a protein fiber used in textiles and some medical sutures or biomaterials; reactions and tissue response depend on the product and processing.

\medterm{Silk Road Disease (Behcet's Syndrome)} Alternate index label for Behcet's Syndrome.

\medterm{Siltuximab} A monoclonal antibody that binds interleukin-6 and is used to treat adults with multicentric Castleman disease who are negative for HIV and human herpesvirus 8.

\medterm{Silver} A metal used in some dressings, creams, and medical devices for antimicrobial activity; excessive exposure can cause permanent blue-gray skin discoloration called argyria.

\medterm{Silver Diamine Fluoride} A liquid applied to a tooth to slow or stop active tooth decay. It usually turns the treated decayed area black and may be useful when drilling or immediate restorative treatment is difficult.
\medterm{Silver Nitrate} A caustic antiseptic used topically to cauterize tissue, control minor bleeding, or treat some wounds; it can burn healthy skin and permanently stain tissue.

\medterm{Silver Poisoning} Silver poisoning, or argyria, follows excessive silver exposure and can cause permanent blue-gray skin discoloration and, rarely, organ or neurologic effects.

\medterm{Silver Sulfadiazine} A topical antimicrobial cream formerly used widely for burns; it may delay healing in some wounds and is not appropriate for everyone, including some newborns and people with sulfonamide allergy.

\medterm{Silver-Russell Syndrome} A growth disorder usually present from birth, marked by poor growth before and after birth, a small or asymmetrical body, and characteristic facial features. Most cases are linked to changes in gene activity rather than a change in the gene itself.

\medterm{Silymarin} A mixture of flavonolignans from milk thistle marketed as a dietary supplement; evidence for treating liver disease is limited and product quality varies.

\medterm{Simeprevir} An older hepatitis C virus NS3/4A protease inhibitor formerly used with other antivirals; it is no longer a usual first-line treatment because better direct-acting regimens are available.

\medterm{Simethicone} Simethicone is a non-absorbed antifoaming agent that coalesces gas bubbles in the gastrointestinal tract and is used for symptomatic relief of bloating; it does not treat the underlying cause of pain.

\medterm{Simian Crease} A simian crease is a single transverse crease across the palm, which may be a normal variant or occur with some chromosomal or developmental conditions.


\medterm{Simmonds Disease} An older term for severe pituitary-gland failure. It can cause extreme tiredness, weight loss, low blood pressure, loss of menstrual periods or sexual function, and deficiencies of several hormones.

\medterm{Simmonds' Disease} An older term for hypopituitarism, especially pituitary failure causing wasting and endocrine deficiency.
\medterm{Simple Febrile Seizure} A brief, generalized seizure caused by fever in a child, usually between six months and five years old. It lasts less than 15 minutes and does not recur during the same illness. The child should be assessed to find the cause of fever, but most children recover fully and do not develop epilepsy.

\synonyms
Benign Febrile Seizure; Typical Febrile Seizure

\medterm{Simple Goiter} A simple goiter is noncancerous enlargement of the thyroid without obvious thyroid-hormone excess or deficiency, caused by iodine imbalance, stimulation, or other factors.

\medterm{Simple Mastectomy} Simple mastectomy removes the breast tissue and nipple while leaving most chest muscles and usually the underarm lymph nodes.

\medterm{Simple Prostatectomy} Surgical removal of the enlarged central portion of the prostate to relieve bladder-outlet obstruction from benign prostatic hyperplasia; it does not remove the entire prostate.

\medterm{Simple Pulmonary Eosinophilia} A usually mild lung reaction in which eosinophils, a type of white blood cell, collect temporarily in the lungs. It may cause cough, fever, or wheezing and can follow medicines, parasites, or other exposures. Symptoms and chest shadows often clear on their own, but the cause should be assessed.

\medterm{Simple, Heart-Smart Substitutions} Easy food changes that may support heart health, such as choosing unsaturated oils instead of butter, whole grains instead of refined grains, and less processed or salty food.

\medterm{Simplexvirus} A genus of herpesviruses that includes herpes simplex virus types 1 and 2, which cause cold sores, genital herpes, and other infections.

\medterm{Simuliidae} The family of black flies; bites can cause itchy skin reactions, and some species transmit Onchocerca volvulus, the parasite that causes river blindness.

\medterm{Simultanagnosia} An inability to see more than one object, or the whole visual scene, at the same time even though basic eyesight may be adequate. It usually results from damage to brain areas that combine visual information.

\medterm{Simvastatin} A statin medicine that lowers LDL, often called bad cholesterol, by reducing cholesterol production in the liver. It helps lower the risk of heart attack and stroke; muscle pain, liver problems, and medicine interactions are possible.
\medterm{Sin Nombre Virus Infection} Alternate index label; see \textit{Hantavirus pulmonary syndrome}.

\medterm{Sincalide} A synthetic form of cholecystokinin used during gallbladder imaging to stimulate contraction and measure bile flow; it may cause abdominal pain, nausea, or flushing.

\medterm{Sindbis Virus} A mosquito-borne virus that can cause fever, rash, and joint pain. It occurs mainly in parts of Africa, Europe, Asia, and Australia.

\medterm{Sinding-Larsen-Johansson Syndrome} A painful overuse injury at the lower tip of the kneecap in growing children and adolescents. Pain usually worsens with running or jumping and improves as growth finishes and the irritated area heals.

\medterm{Sinecatechins} A topical green-tea catechin ointment used for external genital and perianal warts caused by human papillomavirus; it can cause local irritation and should not be used internally.

\medterm{Singer's Nodule} A small, usually noncancerous thickening on a vocal fold caused by repeated voice strain. It can cause lasting hoarseness, a rough voice, or reduced vocal range; voice therapy and reducing strain are often helpful.
\medterm{Single Nucleotide Polymorphism} A common inherited difference in one DNA building block. Most have little or no effect, but some are linked with disease risk, ancestry, or how a person responds to a medicine.

\medterm{Single Palmar Crease} A single transverse palmar crease is one crease crossing the palm instead of the usual two; it can be a normal variant or a feature of genetic syndromes.

\medterm{Single Photon Emission Computed Tomography} A nuclear medicine scan that uses a radioactive tracer and a rotating camera to create three-dimensional images of blood flow or organ function. It can assess the heart, bones, brain, kidneys, or other organs.

\medterm{Single Ventricle} A group of heart defects present at birth in which one lower chamber does most of the work of pumping blood to both the body and lungs. Babies usually need specialist care and staged operations.

\medterm{Single-Dose Toxicity} Harm caused by one exposure to a substance or medicine, assessed according to dose, route, timing,
and the exposed person's characteristics.

\medterm{Single-Fiber Electromyography} A specialized electrical test that records signals from individual muscle fibers. It measures small variations in the timing of nerve-to-muscle communication and can detect disorders such as myasthenia gravis, even when routine nerve tests are normal. A fine needle is used, and the test may be uncomfortable.

\synonyms
SFEMG; Single-Fiber EMG

\medterm{Single-Photon Absorptiometry} An older scan that used a small amount of radiation to estimate bone density, usually in an arm or other peripheral bone. More accurate bone-density methods have largely replaced it.

\synonyms
SPA

\medterm{Single-Photon Emission Computed Tomography} A nuclear-medicine scan that uses a radioactive tracer and a rotating camera to create three-dimensional images of blood flow or organ activity. It can assess the heart, brain, bones, kidneys, and other organs.

\medterm{Singlet Oxygen} A highly reactive form of oxygen involved in oxidative reactions, immune defenses, photodynamic therapy, and tissue injury from oxidative stress.

\medterm{Singultus} Medical term for hiccups; see \textit{Hiccups}.

\medterm{Sinoatrial Node} The heart's natural pacemaker in the right atrium that normally initiates each heartbeat.

\synonyms
Pacemaker, natural
\medterm{Sinuatrial Node} The sinoatrial node, a group of specialized cells in the right atrium that normally initiates the heartbeat and sets the cardiac rhythm.

\medterm{Sinus} A sinus is a cavity, channel, or tract; in anatomy it may mean a paranasal air space, venous channel, or abnormal passage from a diseased area to the surface.
\medterm{Sinus Arrhythmia} Sinus arrhythmia is normal variation in sinus-heart rate, commonly speeding during inhalation and slowing during exhalation, especially in young people.

\medterm{Sinus Bradycardia} Sinus bradycardia is a sinus rhythm slower than the expected rate; it can be normal during sleep or athletic conditioning but may also reflect drugs or disease.

\medterm{Sinus CT Scan} Computed tomography of the paranasal sinuses used to evaluate chronic or complicated sinusitis, anatomic obstruction, trauma, inflammatory disease, or tumors.

\medterm{Sinus Headache} Pain attributed to a sinus headache is often migraine, but sinus inflammation can cause facial pressure, congestion, reduced smell, and pain that worsens with bending. Fever, purulent discharge, persistent symptoms, or eye or neurologic signs guide evaluation.

\medterm{Sinus Histiocytosis} An increase in cleanup immune cells within the open spaces of a lymph node. It is often a response to nearby inflammation, infection, a tumor, or tissue breakdown and is interpreted with the rest of the biopsy.

\medterm{Sinus Lift} A dental procedure that adds bone beneath the lining of the maxillary sinus in the upper jaw. It may create enough support for a dental implant when the jawbone has become too thin.
\medterm{Sinus MRI Scan} Magnetic resonance imaging of the paranasal sinuses and surrounding soft tissues used to evaluate tumors, invasive infection, orbital or intracranial extension, and selected inflammatory disease.

\medterm{Sinus Node} The heart’s natural pacemaker, located in the right atrium.

\synonyms
Sinoatrial node

\medterm{Sinus Pain} Pressure or pain around the forehead, cheeks, nose, or eyes caused by sinus inflammation, infection, migraine,
or other conditions. Severe headache, eye swelling, vision change, high fever, or neurologic
symptoms requires urgent assessment.

\medterm{Sinus Rhomboidalis} A diamond-shaped depression forming the floor of the fourth brain ventricle, near structures controlling movement and vital functions.

\medterm{Sinus Rhythm} A heart rhythm started by electrical impulses from the sinoatrial node, the natural pacemaker. It is usually regular, although sinus rhythm can be too fast, too slow, or affected by illness.

\medterm{Sinus Tachycardia} A heart rhythm originating in the sinoatrial node with a rate faster than normal, commonly occurring as a physiologic response to exercise, fever, pain, dehydration, or stress.
\medterm{Sinus X-Ray} Plain radiographic imaging of the paranasal sinuses; it has limited sensitivity and is often replaced by CT when detailed evaluation is needed.

\medterm{Sinuses} Air-filled cavities within the skull bones lined by respiratory mucosa, connected to the
nasal cavity via ostia.

\medterm{Sinusitis} Inflammation of the lining of one or more paranasal sinuses. It may follow a viral respiratory infection or be associated with allergy, nasal inflammation, or other causes. Common features include nasal obstruction, discharge, facial pressure or pain, headache, and reduced smell. Acute and chronic forms differ in duration and pattern.
\medterm{Sinusitis And Sinus Problems} Alternate index label for Sinusitis.


\medterm{Sinusitis, Acute} Alternate index label; see \textit{Acute sinusitis}.

\medterm{Sinusitis, Chronic} Alternate index label; see \textit{Chronic sinusitis}.

\medterm{Sinusoid} A small, wide blood vessel with unusually open walls, found especially in the liver, spleen, and bone marrow. Its structure allows blood to exchange cells, proteins, nutrients, and waste with nearby tissue.

\medterm{Sinusoidal Obstruction Syndrome (Hepatic Veno-Occlusive Disease)} Liver injury caused by damage and blockage of the small veins and sinusoids inside the liver. It most often occurs after certain cancer treatments or blood-forming stem-cell transplantation. It may cause painful liver enlargement, jaundice, rapid weight gain, fluid in the abdomen, and portal hypertension.

\medterm{Sinusoidal Fetal Heart Rate Pattern} A rare, smooth, regular wave pattern on an unborn baby’s heart-rate tracing with little normal variation. It can indicate severe anemia or low oxygen and requires immediate assessment and monitoring.


\medterm{Sipuleucel-T} A personalized cancer immunotherapy made from a patient’s immune cells and used for selected advanced prostate cancers. The cells are trained to recognize a prostate-cancer target before being returned to the patient.
\medterm{Sirenomelia} A rare congenital malformation in which the lower limbs are partially or completely
fused. It is usually associated with severe abnormalities of the genitourinary,
gastrointestinal, vascular, and spinal systems and is often fatal in the neonatal
period.

\medterm{Sirolimus} Sirolimus is an mTOR inhibitor used for immunosuppression after transplantation and selected vascular or lymphatic disorders; mouth ulcers, hyperlipidaemia, cytopenias, infection, impaired wound healing, and pneumonitis can occur.

\medterm{Sisomicin} An aminoglycoside antibiotic related to gentamicin that has activity against some aerobic gram-negative bacteria; use is limited by kidney and hearing toxicity.

\medterm{Sister Mary Joseph Nodule} A metastatic tumor deposit in or around the umbilicus, usually from an intra-abdominal or pelvic malignancy.

\medterm{Sister Mary Joseph's Nodule} A firm, hard nodule in the umbilicus resulting from metastatic malignancy, most commonly
from intra-abdominal or pelvic cancers including gastric, ovarian, and colorectal
carcinoma. It indicates advanced, disseminated disease with a poor prognosis. Diagnosis
is by fine needle aspiration biopsy.

\medterm{Sitagliptin Phosphate} Sitagliptin phosphate is a DPP-4 inhibitor that prolongs incretin activity and improves glucose-dependent insulin secretion in type 2 diabetes; dose adjustment in kidney disease and rare pancreatitis or severe joint pain are important considerations.


\medterm{Sitosterolemia} A rare inherited disorder in which the body absorbs and stores too much plant sterol, a substance found in plant foods. The buildup can cause early artery disease, skin lumps, anemia, easy bleeding, or low platelets.


\medterm{Sitting Position} An upright or semiupright position in which the patient's torso is elevated and the hips and knees are flexed; it is used for selected shoulder and neurosurgical procedures.

\medterm{Situational Syncope} A reflex faint triggered by a specific situation such as coughing, urination, defecation, swallowing, laughter, or prolonged standing, caused by transient reduced cerebral perfusion.

\medterm{Situs Ambiguus} A birth difference in which internal organs are arranged irregularly rather than in the usual left-right pattern. It is often associated with congenital heart defects and an absent, small, or abnormally functioning spleen.

\medterm{Situs Inversus} A congenital condition in which the thoracic and abdominal organs are reversed or
mirrored from their normal positions. Complete situs inversus (situs inversus totalis)
involves all organs.

\medterm{Situs Inversus Totalis} A congenital condition in which the major visceral organs are mirrored from their normal
positions, with the liver on the left, heart on the right, and other viscera similarly
reversed.

\medterm{Sitz Bath} Immersion of the perineal and anal region in warm or prescribed water for symptom relief in conditions such as hemorrhoids, anal fissures, or postpartum discomfort.

\medterm{Six Fingers Or Toes} Six fingers or toes is polydactyly, a congenital extra-digit difference that may occur alone or as part of a genetic syndrome.

\medterm{Six-Minute Walk Test} A simple test measuring how far a person can walk in six minutes on a flat path. Clinicians record distance, symptoms, heart rate, and oxygen level to assess exercise ability and monitor heart or lung disease. Rest breaks are allowed, and the result is interpreted with other findings.

\synonyms
6MWT; Six-Minute Walk; 6-Minute Walk Test

\medterm{Sixth Disease (Roseola)} Alternate index label for Roseola.

\medterm{Size Perception} The brain’s judgment of how large or small an object or body part is. Vision, touch, distance, and context all contribute; eye or brain disorders can make objects appear unusually large or small.

\medterm{Sjogren Syndrome} Older name for Sjögren disease, an autoimmune disease that commonly causes dry eyes and dry mouth.

\medterm{Sjögren Disease} The current preferred name for Sjögren syndrome, an autoimmune disease that commonly causes dry eyes and dry mouth and may affect other organs.

\medterm{Sjogren-Larsson Syndrome} An inherited disorder causing dry, scaly skin, muscle stiffness, and problems with movement or development. It results from changes in a gene involved in breaking down certain fats.
\medterm{SjöGren Syndrome} Older spelling and name for Sjögren disease.

\medterm{SjöGren's Disease} Alternate index label; see \textit{Sjögren disease}.

\medterm{SjöGren's Syndrome} Older name for Sjögren disease.
\synonyms
SjöGren's Disease

\medterm{SJS} Alternate index label; see \textit{Stevens-Johnson syndrome}.

\medterm{Skeletal Development} Skeletal development is formation, growth, remodeling, and maturation of bones and cartilage from embryonic life through adulthood.

\medterm{Skeletal Dysplasia} An inherited disorder of bone or cartilage development that may cause short or uneven growth, unusually shaped limbs, spine problems, or other differences in the skeleton. Effects vary widely between conditions.

\medterm{Skeletal Limb Abnormalities} Skeletal limb abnormalities are differences in the formation, length, alignment, or structure of an arm or leg caused by genetic, developmental, vascular, injury, or other factors.

\medterm{Skeletal Muscle} Striated muscle attached to bones and generally under voluntary control.
\medterm{Skeletal Muscle Tissue} Voluntary striated muscle attached to bones or other structures, composed of contractile fibers that generate movement, posture, heat, and joint stability.

\medterm{Skeletal Muscles} Muscles attached to bones by tendons that produce movement, maintain posture, and generate heat. They are usually under voluntary control but also support automatic actions such as breathing.

\medterm{Skeletal Myocytes} Skeletal myocytes are striated muscle cells responsible for voluntary movement, posture, heat production, and stabilization of joints.

\medterm{Skeletal Survey} A set of X-rays of most of a child's bones, used especially when multiple injuries or a medical condition affecting bone development is suspected. It can show fractures, bone abnormalities, and signs of healing and must be interpreted with the full clinical assessment.

\medterm{Skeletal System} The skeletal system consists of bones, cartilage, joints, and supporting ligaments that provide structure, protect organs, enable movement, and store minerals.

\medterm{Skeleton} The body’s framework of bones and cartilage that supports movement and protects organs. It also stores minerals, contains marrow that makes blood cells, and provides attachment points for muscles and ligaments.
\medterm{Skeleton, Bones Of The} The body’s framework of bones and cartilage, providing support, protection, leverage for movement, mineral storage, and marrow-based blood-cell production.

\medterm{Skeletonization} Final stage of postmortem decomposition where soft tissues, organs, and skin have
completely decayed or been scavenged, leaving only clean osseous skeletal remains.


\medterm{Skene Gland Cyst} A fluid-filled sac that forms when a Skene gland, also called a paraurethral gland, becomes blocked. These glands are located beside the opening of the urethra. Small cysts may cause no symptoms. Larger or infected cysts may cause pain, painful urination, difficulty passing urine, pain during sex, recurrent urinary infections, or an abscess. Diagnosis is usually made by pelvic examination. Symptomatic cysts may need surgical removal or marsupialization.

\medterm{Skilled Nursing Facility} A facility where people receive 24-hour nursing care and rehabilitation after illness, injury, surgery, or a hospital stay. Care may include wound care, medicines, injections, therapy, and help with daily activities until the person can return home or move to another setting.

\medterm{Skilled Nursing Or Rehabilitation Facilities} Facilities that provide 24-hour nursing care, rehabilitation, medication management, and assistance with daily activities after illness, injury, surgery, or functional decline.

\medterm{Skin} The body's outer covering, consisting mainly of the epidermis and dermis with appendages such as hair, nails, and glands. It protects against injury, infection, ultraviolet light, and water loss; detects touch, pressure, pain, and temperature; helps regulate body temperature; and contributes to immune defense and vitamin D production.
\medterm{Skin Lesion} An abnormal area of skin that may result from infection, inflammation, injury, allergy, or a benign or malignant growth.

\medterm{Skin - Clammy} Skin that feels cool, damp, or unusually sweaty. It may occur with heat, anxiety, pain, low blood sugar, shock, or other illness.

\medterm{Skin Abscess} A painful, swollen pocket of pus under the skin, usually caused by a bacterial infection. It may need drainage and antibiotics.

\medterm{Skin Absorption} Movement of a substance through the skin into underlying tissues or the bloodstream. Absorption depends
on the substance, skin condition, contact time, and exposed area.

\medterm{Skin Aging} Natural aging and long-term sun exposure gradually thin and dry the skin, reduce elasticity, slow healing, and cause wrinkles or dark spots. Limiting ultraviolet exposure, using protective clothing, and applying broad-spectrum sunscreen can reduce further damage.

\medterm{Skin Aging And Appearance} Skin aging is a complex biological process influenced by both intrinsic factors, such as
genetics and natural aging, and extrinsic factors, particularly ultraviolet radiation
exposure.

\medterm{Skin Allergies And Contact Dermatitis} Inflammatory skin reactions triggered by exposure to allergens or irritants that come
into direct contact with the skin.

\medterm{Skin Allergy Testing} Tests in which small amounts of suspected allergens are placed on or just under the skin and the reaction is observed. They can show sensitization to foods, medicines, venoms, or environmental allergens but must be interpreted with symptoms and exposure history.


\medterm{Skin And Aging} Skin aging involves intrinsic and extrinsic changes. Intrinsic aging causes thinning,
wrinkling, and fragility. Extrinsic (photo)aging causes wrinkles, telangiectasias,
lentigines, and actinic keratoses. Collagen loss, elastin degradation, and decreased
cell turnover contribute to aging skin.

\medterm{Skin And Environment} Environmental factors significantly affect skin health. Pollution causes oxidative
stress, inflammation, and accelerated aging. Particulate matter penetrates skin,
generating free radicals. Blue light from screens causes hyperpigmentation, particularly
in darker skin tones. Indoor heating and air conditioning cause dryness.

\medterm{Skin And Hair Changes During Pregnancy} Pregnancy can alter pigmentation, hair growth or shedding, acne, stretch marks, and vascular features through hormonal and circulatory changes; concerning lesions need assessment.

\medterm{Skin Appendages} Structures that grow from or lie within the skin, including hair follicles, oil glands, and sweat glands. They help lubricate and cool the body and can be affected by acne, infection, inflammation, or inherited disorders.

\medterm{Skin Barrier Function} The outer skin barrier limits water loss and helps keep out irritants, allergens, and germs. Its tightly packed surface cells and natural oils maintain protection; injury or inflammation can make skin dry, cracked, and more vulnerable.
\medterm{Skin Barrier} Alternate term; see \textit{Skin Barrier Function}.

\medterm{Skin Biopsy} Removal of a small skin sample for examination under a microscope. It can help diagnose skin cancer, infection, inflammatory disease, blistering disease, or an unexplained rash; the method depends on the lesion.

\medterm{Skin Biopsy Methods} The sampling method is chosen according to the lesion’s appearance and depth. A punch takes a round full-thickness sample, a shave removes a superficial layer, and an excision removes the whole lesion when complete removal is appropriate.

\medterm{Skin Biopsy Techniques} Skin biopsy is selected according to the suspected diagnosis and lesion depth. Punch
biopsy samples full-thickness skin; shave biopsy samples superficial lesions; excisional
biopsy removes an entire small lesion; and incisional biopsy samples part of a large
lesion.

\medterm{Skin Blushing/Flushing} Temporary reddening or warming of the skin caused by widening of small blood vessels. It may result from heat, emotion, exercise, fever, medicines, or hormone changes.

\medterm{Skin Cancer} Cancer that begins in cells of the skin. Common types include basal cell carcinoma, squamous cell carcinoma, and melanoma; less common types arise from other skin structures. Risk is influenced by ultraviolet exposure, immune suppression, genetic susceptibility, and certain preexisting lesions. A skin cancer may appear as a changing, bleeding, growing, scaly, or nonhealing spot.
\synonyms
Cancer, Skin

\medterm{Skin Cancer Epidemiology} Skin cancers are common and include basal-cell carcinoma, squamous-cell carcinoma, and
melanoma. Incidence and mortality vary by tumor type, geography, skin phenotype, ultraviolet
exposure, immune status, and access to diagnosis and treatment.

\medterm{Skin Cancer Screening} Examination of the skin for new or changing spots, especially in people at increased risk. Clinicians may use a full-body examination, a magnifying instrument, or a biopsy; routine screening recommendations vary by individual risk.

\medterm{Skin Cancer Staging} Skin cancer staging guides treatment and prognosis. Melanoma: Breslow thickness (most
important prognostic factor), ulceration, mitotic rate, sentinel lymph node status (AJCC
8th edition).

\medterm{Skin Cancer Surgery} Surgery to remove skin cancer with a margin of apparently healthy tissue. The method depends on the cancer type, size, location, and spread and may include standard excision or Mohs surgery, which checks the edges during removal.

\medterm{Skin Culture} Laboratory testing of a swab, scraping, fluid, or tissue sample from the skin to identify bacteria, fungi, or other organisms causing infection.

\medterm{Skin Cancer, Melanoma} Alternate index label; see \textit{Melanoma}.

\medterm{Skin Carcinoma} Cancer arising from skin cells, most commonly basal-cell or squamous-cell carcinoma. It may appear as a changing sore, lump, scaly patch, or wound that does not heal and needs examination.

\medterm{Skin Care} Measures that protect and maintain the skin, including gentle cleansing, moisturization, sun protection, wound care, avoidance of irritants, and treatment of identified disease.

\medterm{Skin Care And Incontinence} Skin protection for people with urine or stool leakage, including gentle cleansing, thorough drying, absorbent products, barrier creams, regular changing, and treatment of the underlying cause.

\medterm{Skin Color} The visible color of skin, determined mainly by melanin, blood flow, oxygenation, and other pigments.

\medterm{Skin Erosion} A shallow loss of the epidermis that leaves a moist or crusted surface but does not extend deeply into the dermis, often caused by friction, scratching, blister rupture, or inflammation.

\medterm{Skin Scraping} Removal of a small amount of surface skin scale or material for laboratory examination, commonly to look for fungi, mites, or other causes of a rash.

\medterm{Skin Examination} Systematic inspection and palpation of the skin, hair, nails, and mucous membranes to assess color, lesions, texture, distribution, vascular signs, infection, and cancer warning features.

\medterm{Skin Findings In Newborns} Common skin changes in newborns, such as milia, newborn rash, peeling, birthmarks, or bluish hands and feet. Most are harmless, but fever, blisters, widespread redness, or a sick appearance needs prompt assessment.

\medterm{Skin Flaps} Pieces of skin and underlying tissue moved to cover a wound while keeping or reconnecting their blood supply. Flaps may come from nearby tissue or another part of the body and are used after injury, cancer removal, or reconstructive surgery.


\medterm{Skin Graft} Skin taken from one area of the body and placed over a wound, burn, ulcer, or surgical defect. A graft must obtain blood from the wound surface to survive and may fail, scar, change color, or require further surgery.

\medterm{Skin Graft, Pedicle} A piece of skin moved to cover a wound while remaining attached to its original blood supply through a tissue bridge.

\medterm{Skin Immune System} The cells and protective substances in skin that recognize germs, remove damaged cells, and control inflammation. It works with the skin barrier to prevent infection and support healing after injury.

\medterm{Skin Induration} Unusual hardness or firmness of the skin or tissue beneath it. Causes include scar formation, inflammation, swelling, infection, an infiltrating disease, or a growth; examination determines the cause.

\medterm{Skin Infection} Infection of the skin or underlying soft tissue, commonly caused by bacteria, fungi, or viruses. Redness, warmth,
swelling, pain, drainage, or fever may occur; rapidly spreading redness, severe pain, blisters, or systemic
illness requires prompt medical assessment.

\medterm{Skin Integrity} The condition of skin and underlying tissue when it remains intact and protected from injury. Prolonged pressure, friction, moisture, poor blood flow, and illness can cause breakdown, especially over bony areas.

\medterm{Skin Irritation} Inflammation or discomfort caused by direct contact with an irritant, often producing redness, burning,
itching, dryness, or scaling.

\medterm{Skin Lesion Aspiration} Removal of fluid or cells from a skin lump with a needle for testing or drainage.

\medterm{Skin Lesion Biopsy} Removal of part or all of a skin lesion for microscopic examination to diagnose infection, inflammatory disease, benign tumor, or skin cancer.

\medterm{Skin Lesion KOH Exam} Microscopic examination of skin scrapings treated with potassium hydroxide to detect fungal hyphae or yeast in suspected superficial fungal infection.

\medterm{Skin Lesion Of Blastomycosis} A blastomycosis skin lesion results from dissemination of Blastomyces fungal infection and may appear as verrucous or ulcerated lesions with raised borders.

\medterm{Skin Lesion Removal} Removal of an unwanted, suspicious, or diseased area of skin by shaving, cutting, freezing, scraping, or another method. Tissue may be sent for laboratory examination; bleeding, infection, scarring, and incomplete removal are possible.


\medterm{Skin Lumps} Skin lumps are localized bumps or masses caused by cysts, lipomas, inflamed follicles, infections, benign tumors, or less commonly cancer; persistence or change warrants assessment.

\medterm{Skin Mastocytoma} A localized collection of mast cells in the skin, usually a yellow-brown or reddish plaque or nodule in a child that may swell, itch, or redden when rubbed.

\medterm{Skin Microbiome} The community of bacteria, fungi, viruses, and other microorganisms normally living on the skin. They interact with the skin and immune system; changes in this community may contribute to some skin diseases.

\medterm{Skin Necrosis} Death of skin tissue from ischemia, infection, pressure, toxins, trauma, vasculitis, or severe inflammation, producing discoloration, blisters, ulcers, or black eschar.

\medterm{Skin Neoplasm} A skin neoplasm is an abnormal growth in the skin that may be benign, precancerous, or malignant.

\medterm{Skin Nodules} Firm, raised lumps in or under the skin. They may result from inflammation, infection, injury, cysts, or benign or malignant growths.

\medterm{Skin Of Color} A broad term for people with more deeply pigmented skin. Some conditions may appear differently, and dark marks or raised scars may develop more readily after inflammation or injury; skin cancer can occur in every skin tone.

\medterm{Skin Or Nail Culture} Laboratory culture of skin, hair, or nail material to identify bacteria or fungi when infection is suspected; results require clinical correlation.

\medterm{Skin Pigmentation} Skin color is produced mainly by melanin made by pigment-producing cells. Sun exposure, hormones, inflammation, medicines, and illness can make areas darker or lighter; new or changing discoloration may need examination.


\medterm{Skin Resurfacing} Procedures that improve skin texture, scars, wrinkles, or uneven color by removing or changing surface skin layers. Methods include lasers, chemical peels, and abrasion; risks include infection, pigment changes, and scarring.

\medterm{Skin Self-Exam} A skin self-exam is a regular check of the skin for new or changing moles, sores, lumps, or patches that may need medical assessment.

\medterm{Skin Tag} A skin tag is a benign soft growth of connective tissue, commonly in body folds and associated with friction, age, pregnancy, obesity, or insulin resistance. Removal may be performed for irritation or appearance, but a changing or atypical lesion should be diagnosed first.

\synonyms
Acrochordon (Skin Tag)

\medterm{Skin Tags} Soft, usually harmless skin growths that commonly occur in body folds such as the neck, armpits, groin, or beneath the breasts. They may become irritated or bleed; a changing growth should be examined before removal.

\medterm{Skin Test} Alternate index label for diagnostic or immunologic skin tests, including prick, patch, and tuberculin testing.
\medterm{Skin Turgor} The skin's ability to return to its normal position after being gently pinched. Reduced turgor may occur with dehydration, although it is less reliable in older adults.

\medterm{Skin Writing} A skin condition in which light scratching or pressure causes temporary raised red lines or welts.

\medterm{Skin, Diseases Of} Conditions affecting skin, hair, nails, glands, or subcutaneous tissue, including infection, inflammation, neoplasia, and genetic disorders.
\medterm{Skin, Itchy} Alternate index label; see \textit{Itchy skin (pruritus)}.

\medterm{Skin-Grafting} Surgery that moves skin from a donor area to cover a wound, burn, ulcer, or tissue defect. The graft needs blood supply from the wound bed and may fail, scar, or change color.
\medterm{Skipped Heartbeats} Alternate index label; see \textit{Heart palpitations}.

\medterm{Skull} The bony framework of the head that protects the brain and supports the face, eyes, ears, and jaw.
\medterm{Skull Fracture} A break in one or more skull bones caused by trauma. It may occur with bleeding inside the skull, brain injury, or leakage of clear fluid from the nose or ear and requires medical assessment.

\medterm{Skull X-Ray} A skull X-ray uses low-dose radiation to show skull bones and may help assess selected fractures, foreign bodies, or bone abnormalities.

\medterm{Slanted Ear} An ear whose upper edge slopes unusually compared with typical ear orientation; it may be isolated or part of a genetic condition.

\medterm{SLAP Lesion} A tear in the upper rim of the shoulder socket, sometimes where the biceps tendon attaches. It may follow a fall, lifting injury, or repeated overhead activity and can cause pain, clicking, weakness, or instability.

\medterm{SLAP Tears} Tears of the superior glenoid labrum extending from anterior to posterior, sometimes involving the
biceps tendon attachment. They may cause pain, clicking, or instability, especially with overhead activity.

\medterm{Slapped Cheek Disease} Alternate index label; see \textit{Parvovirus infection}.

\medterm{Slapped Cheek Syndrome} Erythema infectiosum, a parvovirus B19 infection causing a slapped-cheek rash and sometimes joint or hematologic complications.
\medterm{SLE} Systemic lupus erythematosus, a chronic autoimmune disease that can affect skin, joints, kidneys, blood, nervous system, and other organs.
\medterm{Sleep} A naturally recurring state in which awareness and responses to the surroundings decrease while the brain and body follow organized patterns. Sleep supports learning, memory, emotional health, physical recovery, immune function, and metabolism.
\synonyms
Disorders Sleep (Sleep)

\medterm{Sleep (Geriatric)} Sleep patterns and sleep problems in older adults, including insomnia, sleep apnea, circadian changes,
and medication-related disturbances.

\synonyms
Sleep care for older adults; geriatric sleep problems

\medterm{Sleep And Sleep Disorders In Children And Teens} Children and teens need age-appropriate sleep, and disorders may involve insomnia, circadian delay, sleep apnea, parasomnias, restless legs, or insufficient opportunity. Snoring, pauses in breathing, unusual daytime sleepiness, mood change, or poor school performance warrants evaluation.

\medterm{Sleep Apnea} Repeated pauses or reductions in breathing during sleep. It can result from the airway closing, the brain failing to signal breathing, or both, and may cause loud snoring, gasping, poor sleep, daytime sleepiness, and low oxygen.

\medterm{Sleep Apnea (Dental)} Alternate index label for Oral Appliance Therapy.

\medterm{Sleep Apnea (Obstructive Sleep Apnea - OSA)} Alternate index label for Obstructive Sleep Apnea.


\medterm{Sleep Apnea, Central} Alternate index label; see \textit{Central sleep apnea}.

\medterm{Sleep Apnea, Obstructive} Alternate index label; see \textit{Obstructive sleep apnea}.

\medterm{Sleep Apnoea} Alternate index label for Sleep Apnea.

\medterm{Sleep Architecture} The pattern and order of sleep stages during a night or other sleep period. It includes lighter sleep, deep sleep, and dream-related REM sleep; illness, medicines, age, stress, and sleep disorders can change how much time is spent in each stage.

\medterm{Sleep Bruxism} Repeated grinding or clenching of the teeth during sleep, which may cause tooth wear, jaw pain, or headaches.

\medterm{Sleep Debt} The accumulated difference between the amount of sleep a person needs and the amount obtained over time; it may impair alertness, mood, cognition, and performance.

\medterm{Sleep Diary} A daily record of sleep and wake times, awakenings, naps, medicines, caffeine, and related factors used to assess sleep patterns and treatment response.

\medterm{Sleep Disorder Periodic Limb Movement (Periodic Limb Movement Disorder)} Alternate index label for Periodic Limb Movement Disorder.

\medterm{Sleep Disorders} A group of conditions affecting sleep quality, timing, duration, breathing, movement, or behavior. Major groups include insomnia, sleep-related breathing disorders, central disorders of hypersomnolence, circadian rhythm sleep-wake disorders, parasomnias, and sleep-related movement disorders. Examples include sleep apnea, narcolepsy, restless legs syndrome, and sleep paralysis.

\medterm{Sleep Disorders In Older Adults} Sleep problems in later life, including insomnia, sleep apnea, restless legs, and circadian changes; medication effects, pain, depression, and medical illness should be assessed.

\medterm{Sleep Efficiency} The percentage of time in bed that a person is actually asleep. It is calculated by dividing total sleep time by time in bed; a low value may occur with insomnia, repeated awakenings, or other sleep disorders.

\medterm{Sleep Fragmentation} Repeated brief awakenings or partial arousals that interrupt otherwise continuous sleep. It can reduce restorative sleep and occur with sleep apnea, pain, noise, movement disorders, medicines, or illness.

\medterm{Sleep Hygiene} Habits and surroundings that support regular, restorative sleep. Helpful measures include a consistent schedule, a dark and quiet bedroom, relaxing before bed, and limiting caffeine, alcohol, nicotine, and stimulating screen use near bedtime.

\medterm{Sleep Medications And Supplements} Pharmacological agents used to treat insomnia and other sleep disorders when behavioral
interventions alone are insufficient.

\medterm{Sleep Onset Latency} The time required to fall asleep after attempting to sleep; prolonged latency may occur with insomnia, anxiety, circadian misalignment, stimulants, or another sleep disorder.

\medterm{Sleep Paralysis} Temporary inability to move or speak when falling asleep or waking, while awareness is partly or fully preserved.
Episodes usually last seconds to minutes and may include frightening hallucinations. Recurrent episodes
can be associated with sleep deprivation, irregular sleep, narcolepsy, or other sleep disorders.

\medterm{Sleep Related Breathing Disorders} Sleep-related breathing disorders cause abnormal ventilation or oxygenation during sleep, including obstructive and central sleep apnea and sleep-related hypoventilation. Snoring, witnessed pauses, gasping, morning headache, or daytime sleepiness may prompt sleep testing and treatment.

\synonyms
Breathing Disorders Sleep Related (Sleep Related Breathing Disorders)

\medterm{Sleep Spindles} Short bursts of regular brain activity that occur mainly during stage 2 non-REM sleep. They are seen on a sleep recording and are thought to help protect sleep and process newly learned information.

\medterm{Sleep Studies (Polysomnography)} Overnight tests that record brain activity, eye movements, heart rate, breathing, blood oxygen, body position, and limb movements. They help identify sleep apnea, unusual movements, seizures during sleep, and other sleep disorders.

\synonyms
Polysomnography

\medterm{Sleep, Disorders Of} Health conditions that disturb how much a person sleeps, when they sleep, or how restful sleep is. They may involve insomnia, breathing pauses, unusual movements, excessive sleepiness, unusual nighttime behavior, or a body clock that is out of step.
\medterm{Sleep-Related Hypoventilation} Breathing that is too shallow during sleep to remove enough carbon dioxide or maintain normal oxygen. It can result from obesity, lung or muscle disease, chest-wall problems, brain disorders, or sedating medicines.

\medterm{Sleep-Related Hypoxemia} Abnormally low blood oxygen during sleep that may occur with sleep-disordered breathing, lung disease, neuromuscular weakness, hypoventilation, or other medical conditions.

\medterm{Sleepiness, Daytime} Alternate index label for Excessive Daytime Sleepiness.

\medterm{Sleeping Sickness} An infection caused by Trypanosoma parasites and spread by tsetse flies in sub-Saharan Africa. It may begin with fever and swollen nodes, then progress to sleep disturbance, confusion, and other nervous-system problems.
\medterm{Sleeplessness} Alternate lay term; see \textit{Insomnia Disorder}.

\medterm{Sleepwalking} Recurrent walking or other complex behavior during sleep, usually with little memory afterward. It is more common in children and may be triggered by sleep deprivation, stress, fever, medicines, or other sleep disorders; injury prevention is important.

\medterm{Sleepwalking Disorder} A sleep disorder involving repeated episodes of walking or other complex behavior during sleep, usually with little memory afterward, that cause injury, distress, sleep disruption, or impaired daytime functioning.

\medterm{Sleepwalking And Children} Sleepwalking in children is usually benign and often resolves with age; regular sleep, hazard prevention, and evaluation of frequent, violent, injurious, or atypical episodes are important.

\medterm{Sleeve Gastrectomy} An operation that removes most of the stomach, leaving a narrow sleeve-shaped section. It limits how much food the stomach can hold and changes appetite hormones; bleeding, leakage, reflux, and nutritional problems can occur.
\medterm{Sleeve Lobectomy} An operation that removes a lung lobe and the nearby affected part of the main breathing tube, then reconnects the remaining airway. It can remove a centrally located tumor while preserving more healthy lung than removing the whole lung. Possible complications include air leakage, narrowing of the airway, infection, and bleeding.

\synonyms
Bronchial Sleeve Resection; Sleeve Resection of the Lung

\medterm{Sliding Filament} The process by which muscles contract: two protein fibers slide past each other, shortening the muscle’s working units. Calcium and energy from ATP allow the fibers to attach, pull, release, and repeat the cycle.

\medterm{Sliding Hiatal Hernia} A sliding hiatal hernia occurs when the stomach and gastroesophageal junction move upward through the diaphragm, often contributing to reflux.

\medterm{Sling} A supportive device or surgically placed strip used to support an injured limb or urinary structures.

\medterm{Slipped Capital Femoral Epiphysis} A hip problem in growing children or teenagers in which the upper end of the thigh bone slips through its growth area. It may cause hip, groin, thigh, or knee pain and a limp and needs urgent assessment to prevent further slipping.

\medterm{Slipped Disk} A spinal disk whose tough outer layer tears or weakens, allowing its inner material to bulge and irritate nearby nerves. It may cause back or neck pain, numbness, weakness, or sciatica.
\synonyms
Herniated disk

\medterm{Slipped Vertebra} Forward movement of one spinal bone compared with the bone below it, also called spondylolisthesis. It may cause back pain, stiffness, hamstring tightness, or pressure on a nerve, although some people have no symptoms.

\synonyms
Spondylolisthesis

\medterm{Slipping Rib Syndrome} Pain caused by excessive movement of the lower ribs, usually from loosened cartilage. The pain may be sharp and can sometimes be reproduced by pressing or moving the rib.

\medterm{Slit Lamp} An eye-examination microscope that uses magnification and a narrow beam of light to inspect the eyelids, cornea, lens, and other eye structures. It helps detect infection, injury, cataract, and glaucoma-related changes.
\synonyms
Biomicroscope

\medterm{Slit Lamp Biomicroscopy} A detailed eye examination using a bright, narrow beam of light and a microscope to view the front of the eye. Additional lenses can help examine the retina and other structures farther back.

\medterm{Slit-Lamp Biomicroscope} Alternate index label for Slit Lamp.

\medterm{Slit-Lamp Exam} Alternate index label for Slit Lamp Biomicroscopy.

\medterm{SLL} Alternate index label; see \textit{Small lymphocytic lymphoma}.

\medterm{Slough} Soft, dead, or detached tissue covering or separating from a wound, ulcer, or inflamed surface. It can delay healing and may need gentle removal after blood flow, infection, and the wound’s depth are assessed.
\medterm{Slow Heart Rate - Bradycardia} A heart rate slower than expected, often below 60 beats per minute in an adult. It may be normal in some people but can cause dizziness, fainting, tiredness, or shortness of breath when it reduces blood flow.

\medterm{Slow Heartbeat} Alternate index label; see \textit{Bradycardia}.

\medterm{Slow-Twitch Fiber} A type of muscle fiber suited to steady activity over a long time. It uses oxygen efficiently and resists tiredness, helping with posture, walking, and endurance activities rather than short bursts of maximum force.

\synonyms
Type I muscle fiber

\medterm{Slow-Wave Sleep} The deepest stage of non-REM sleep, marked by slow brain waves and reduced response to sounds or touch. It supports physical recovery, memory, and hormone release and is often reduced by ageing, illness, alcohol, or sleep disruption.

\synonyms
Deep sleep

\medterm{Small Airway Disease} Narrowing, inflammation, or blockage of the smallest breathing tubes in the lungs. It may cause cough, wheezing, breathlessness, or trapped air and can occur with asthma, COPD, infection, or other lung disease.

\medterm{Small Bowel Adenocarcinoma} A malignant gland-forming epithelial tumor of the small intestine. It is uncommon and may cause occult bleeding, obstruction, abdominal pain, or weight loss; risk is increased in selected inherited and inflammatory disorders.

\medterm{Small Bowel Bacterial Overgrowth} Excess bacteria in the small intestine, where fewer bacteria are normally present. It can cause bloating, abdominal discomfort, diarrhea, weight loss, and poor absorption of nutrients; some cases involve methane-producing organisms instead.

\medterm{Small Intestinal Fungal Overgrowth} A proposed condition in which excess fungi are found in the small intestine and may be associated with bloating, nausea, diarrhea, or abdominal discomfort. Symptoms are nonspecific, testing is not standardized, and other causes should be assessed before attributing symptoms to it.

\medterm{Small Bowel Follow-Through} A series of X-rays taken after swallowing contrast, usually barium, to follow its movement through the small intestine. It can show narrowing, obstruction, inflammation, or abnormal transit.

\medterm{Small Bowel Bleeding} Bleeding from the small intestine. It may cause black or maroon stool, hidden blood loss, or iron-deficiency anemia and may require capsule endoscopy, enteroscopy, imaging, or other tests to find the source.

\medterm{Small Bowel Obstruction} Partial or complete blockage of the small intestine. Common causes include adhesions,
hernias, tumors, and inflammatory narrowing. Symptoms may include abdominal pain,
vomiting, distension, and inability to pass stool or gas.

\medterm{Small Bowel Prolapse} Alternate index label; see \textit{Small bowel prolapse (enterocele)}.

\medterm{Small Bowel Resection} Surgical removal of a diseased segment of the small intestine followed by reconnection or stoma formation when needed.


\medterm{Small Bowel Series} A series of X-ray images of the small intestine after swallowing contrast material, usually barium. It can show narrowing, blockage, abnormal movement, inflammation, tumors, or changes from Crohn disease.

\medterm{Small Bowel Tissue Smear/Biopsy} Examination of cells or a tissue sample from the small intestine to investigate poor nutrient absorption, celiac disease, inflammation, infection, or a tumor.

\medterm{Small Calorie (Calorie)} A small calorie is a unit of heat energy; food labels usually use a kilocalorie, written Calorie with a capital C, equal to 1,000 small calories.

\medterm{Small Cell Carcinoma} A fast-growing cancer made of small cells, most often beginning in the lung and strongly linked to smoking. It may shrink at first with treatment but often returns or spreads early, including to the brain.

\medterm{Small Cell Carcinoma Of The Ovary, Hypercalcemic Type} A rare, highly aggressive ovarian tumor in young patients, often associated with alterations in SMARCA4 and sometimes presenting with hypercalcemia.

\medterm{Small Cell Lung Cancer} An aggressive lung cancer strongly associated with smoking and early spread. It often responds initially to chemotherapy or radiation, but recurrence is common and treatment depends on the stage and overall health.

\medterm{Small Eye} An eye that is smaller than usual because it did not develop fully before birth. Vision may be reduced, especially when other eye structures are also affected, and assessment can identify associated conditions.

\medterm{Small Fiber Neuropathy} Damage to tiny nerves that carry pain, temperature, and automatic body signals. It can cause burning feet, tingling, poor temperature sensation, abnormal sweating, digestive symptoms, or heart-rate changes; diabetes is one possible cause.

\medterm{Small For Gestational Age} Birth weight or size below the expected range for gestational age, caused by constitutional small size, placental insufficiency, infection, genetic disease, or maternal factors.

\medterm{Small Intestinal Bacterial Overgrowth} An abnormal increase or change in bacteria within the small intestine, often causing bloating, abdominal
pain, nausea, diarrhea, fullness, weight loss, or nutrient deficiency. Causes include impaired intestinal movement,
structural problems, surgery, and some systemic diseases; persistent symptoms require medical evaluation.

\medterm{Small Intestinal Bacterial Overgrowth (Small Intestinal Bacterial Overgrowth)} Alternate index label; see \textit{Small Intestinal Bacterial Overgrowth}.

\medterm{Small Intestinal Ischemia And Infarction} Inadequate blood supply to the small intestine that can progress to irreversible tissue death, perforation, sepsis, and shock.

\medterm{Small Intestine} The long, coiled part of the digestive tract between the stomach and large intestine. It receives bile and pancreatic juices, completes digestion, and absorbs most nutrients and much of the body's water.

\medterm{Small Intestine Aspirate And Culture} Collection and culture of fluid aspirated from the small intestine to evaluate suspected bacterial overgrowth or selected infections; the procedure is invasive and results require clinical correlation.

\medterm{Small Intestine Prolapse} Alternate index label; see \textit{Small bowel prolapse (enterocele)}.

\medterm{Small Round Blue Cell Tumor} A tumor whose cells look small, round, and blue under a microscope. This appearance occurs in several cancers, especially in children, so additional tissue tests are needed to identify the exact type.

\medterm{Small Vessel Disease} Disease affecting small blood vessels, which can impair blood flow and contribute to brain, heart, kidney, or other organ problems.

\synonyms
Coronary Microvascular Disease

\medterm{Small-Bowel Transplantation} Surgery replacing a severely diseased or nonfunctioning small intestine with donor bowel. It may help people unable to receive enough nutrition through the gut but requires lifelong rejection-preventing medicines and monitoring.
\medterm{Small-Bowel Volvulus} Twisting of the small intestine that can block its contents and cut off its blood supply. It may cause severe abdominal pain, vomiting, swelling, and inability to pass stool or gas and requires urgent treatment.

\medterm{Small-Cell Carcinoma} A fast-growing cancer, most often starting in the lung and strongly linked with smoking. It may shrink quickly with chemotherapy or radiation but often returns or spreads early, including to the brain.
\medterm{Smallpox} A severe contagious disease caused by variola virus, producing fever and rash; global eradication was certified in 1980.

\medterm{SMART} A way to set goals that are specific, measurable, achievable, relevant, and time-limited. SMART goals state exactly what is wanted, how progress will be checked, and when the goal should be reached.

\synonyms
SMART goals

\medterm{Smashed Fingers} A smashed finger may injure skin, nail, bone, tendon, or nerve; deformity, severe pain, numbness, impaired movement, or blood under the nail may require urgent evaluation.

\medterm{Smear Of Duodenal Fluid Aspirate} A smear of duodenal fluid aspirate examines fluid from the first part of the small intestine for selected parasites or other microscopic findings.

\medterm{Smegma} A collection of shed cells, oils, and moisture around the genital folds, which may promote irritation if hygiene is poor.
\medterm{Smell} The sense of detecting airborne chemicals through receptors in the nasal lining and processing the signals in the brain. Reduced or lost smell can follow infection, nasal disease, head injury, medicines, or neurologic disease.
\medterm{Smell - Impaired} Reduced or absent ability to detect odors, called hyposmia or anosmia, caused by nasal obstruction, viral injury, head trauma, neurodegeneration, medications, or other disease.

\medterm{SMILE} A laser eye operation that corrects short-sightedness and some short-sighted astigmatism. The laser creates a thin piece of tissue inside the cornea, which is removed through a small opening without making a larger corneal flap.

\medterm{Smile Design} Planning smile based on facial features, gum display, tooth proportions. Digital smile design software available.
\medterm{Smith-Lemli-Opitz Syndrome} A rare inherited disorder affecting cholesterol production before birth. It may cause growth and developmental problems, distinctive facial features, heart or limb abnormalities, feeding difficulty, and genital differences.

\medterm{Smith-Magenis Syndrome} A genetic neurodevelopmental disorder, usually involving deletion of 17p11.2 or an RAI1
variant. It causes developmental delay, intellectual disability, distinctive facial
features, sleep disturbance, behavioral dysregulation, and variable congenital
abnormalities.

\medterm{Smoke Evacuator} A device that captures and filters smoke produced during surgery by electrical, laser, or ultrasonic instruments. Removing the plume reduces staff and patient exposure to particles, irritating gases, and potentially infectious material.

\medterm{Smoke Inhalation} Injury from breathing smoke, hot gases, or combustion products, potentially causing airway burns, poisoning, and lung injury.
\medterm{Smoke Inhalation Injury} Pulmonary injury caused by inhalation of toxic products of combustion.

\medterm{Smoke Secondhand (Secondhand Smoke)} Alternate index label for Secondhand Smoke.

\medterm{Smokeless Tobacco} Smokeless tobacco delivers nicotine and causes dependence, oral lesions, gum recession, tooth loss, cardiovascular risk, and cancers of the mouth, throat, and pancreas. Quitting reduces harm, and counseling plus approved medicines improve cessation success.

\synonyms
Tobacco Chewing (Smokeless Tobacco)

\medterm{Smokers Keratosis (Leukoplakia)} Alternate index label for Leukoplakia.

\medterm{Smokers Lung (Smoker's Lung: Pathology Photo Essay)} Alternate index label for Smoker's Lung: Pathology Photo Essay.

\medterm{Smoking} Inhalation of smoke from burning tobacco or other substances, associated with addiction, cancer, cardiovascular disease, and lung disease.
\medterm{Smoking And Asthma} Smoking irritates and inflames airways, worsens asthma control, increases attacks, and reduces treatment response; stopping smoking and avoiding secondhand smoke improve outcomes.

\medterm{Smoking And COPD} Smoking is the major preventable cause of COPD and accelerates lung-function loss; quitting reduces progression and exacerbations even after COPD is established.

\medterm{Smoking And Surgery} Smoking before or after an operation increases the risk of breathing problems, wound infection, poor healing, blood clots, heart complications, and failure of repaired tissue. Stopping as early as possible improves recovery, even if surgery is soon.

\medterm{Smoking Cessation} The process of stopping tobacco or nicotine use. It may involve counseling, behavioral support, medicines, or nicotine
replacement.

\medterm{Smoking Cessation Medications} Medicines that reduce nicotine withdrawal and cravings, including nicotine replacement, bupropion, and varenicline, used with behavioral support and individualized for health risks and preferences.

\medterm{Smoldering Multiple Myeloma} An asymptomatic plasma-cell disorder with increased clonal plasma cells or monoclonal protein but no myeloma-defining organ damage; it requires risk-based monitoring for progression to active multiple myeloma.

\synonyms
Voluntary muscles

\medterm{Smooth Muscle} Muscle that works automatically in the walls of hollow organs, blood vessels, and airways. It controls processes such as bowel movement, bladder emptying, blood-vessel width, airway narrowing, and childbirth contractions.
\medterm{Smothering} Mechanical obstruction of the external nose and mouth that prevents effective
ventilation. It may be caused by a covering, compression, or an obstructing soft
material; findings are often nonspecific and must be interpreted with the scene,
history, and complete examination.

\medterm{Snake Bite} Snake bites can cause puncture wounds, tissue injury, bleeding, paralysis, or venom-related effects on clotting, nerves, muscles, kidneys, or the heart. Call emergency services, keep the person still, and do not cut, suck, ice, or tourniquet the wound.

\synonyms
Bite Snake (Snake Bite)

\medterm{Snake Bites} Wounds caused by a snake's teeth; venomous bites may cause pain, swelling, bleeding, tissue damage, paralysis, abnormal clotting, or shock. Seek emergency care immediately, keep the person still, and do not cut, suck, ice, or tourniquet the wound.

\medterm{Snake-Bite} A wound from a snake that may cause local injury, envenomation, infection, or systemic toxicity.
\medterm{Snapping Finger} Trigger finger, painful catching or locking of a finger caused by stenosing inflammation around a flexor tendon.
\medterm{Snapping Hip Syndrome} A clicking or snapping sensation in the hip, sometimes with pain, caused by a tendon moving over bone or by a problem inside the joint. Activity adjustment and exercises often help; persistent pain may need imaging or other treatment.

\medterm{Sneddon Syndrome} A rare disorder in which small arteries become blocked, causing a net-like purplish skin pattern and repeated strokes or stroke-like events. Headache, high blood pressure, memory problems, or heart-valve changes may also occur.

\medterm{Sneezing} A forceful reflex expulsion of air through the nose and mouth, often triggered by nasal irritation.
\medterm{Snellen Chart} A chart of letters or symbols used to measure visual acuity at a specified distance.
\medterm{Snoring} Noisy breathing during sleep caused by vibration of relaxed or narrowed tissues in the throat. It can be harmless, but loud frequent snoring with pauses in breathing, choking, morning headaches, or daytime sleepiness may indicate sleep apnoea.
\medterm{Snoring And Sleep Apnea} Snoring is the sound produced by vibrating soft tissues in the upper airway during
sleep, often resulting from partial airway obstruction. Sleep apnea is a more serious
condition characterized by repeated episodes of complete or partial airway collapse
during sleep, leading to intermittent hypoxia, frequent arousal, and fragmented sleep.


\medterm{Snort} A snort is a forceful nasal or respiratory sound; it may occur with congestion, sleep-disordered breathing, irritation, or inhalation of a substance.

\medterm{Snot} Alternate informal term; see \textit{Mucus}.

\medterm{Snow Blindness} A painful temporary burn of the eye’s clear front surface caused by intense ultraviolet light, often reflected from snow, ice, water, or sand. Symptoms may include gritty pain, tearing, redness, light sensitivity, and blurred vision.
\medterm{SNRIs} A class of antidepressants that increase the activity of serotonin and norepinephrine, chemical messengers between nerve cells. They treat depression, some anxiety disorders, and certain nerve or long-lasting pain conditions.

\medterm{Snuff} Snuff is finely ground tobacco inhaled through the nose or placed in the mouth and can cause nicotine dependence, oral disease, and cardiovascular harm.

\medterm{Snuffles} A runny or congested nose; in newborns, persistent snuffles can be a sign of congenital syphilis.
\medterm{Social And Psychological Causes Of Alcoholism} Alcohol-use disorder reflects interactions among genetics, brain reward pathways, stress, trauma, mental illness, social environment, and alcohol availability—not a moral failure. Treatment may include withdrawal care, counseling, peer support, and medicines that reduce drinking or relapse.

\medterm{Social Anxiety Disorder} Intense, persistent fear of being judged, embarrassed, or rejected in social or performance situations. It can lead to avoidance, physical anxiety symptoms, and difficulty with school, work, relationships, or daily activities.

\medterm{Social Anxiety Disorder (Social Phobia)} Alternate index label for Social Anxiety Disorder.

\synonyms
Phobia, Social

\medterm{Social Isolation} Limited contact or interaction with other people. Persistent isolation may adversely affect mental and
physical health.

\medterm{Social Worker} A professional who helps patients and families manage practical, emotional, and social problems related to illness. They may arrange community resources, support discharge planning, assist with benefits, and provide counseling or referrals.

\medterm{Sociophobia} An older term for social anxiety disorder, involving excessive fear of scrutiny, embarrassment, humiliation, or negative evaluation in social or performance situations.

\medterm{Sodium} A mineral and body salt helping control fluid balance, nerves, muscles, and blood pressure; abnormal levels can be harmful.
\medterm{Sodium Bicarbonate Therapy} Treatment using sodium bicarbonate to make overly acidic blood less acidic in selected conditions, such as some cases of kidney disease. It can increase sodium and fluid load and requires monitoring.

\medterm{Sodium Bisulfate Poisoning} Poisoning from sodium bisulfate, an acidic chemical used in some cleaners and pool products. It can burn the skin, eyes, mouth, and digestive tract; fumes may irritate breathing passages, and ingestion requires urgent advice.

\medterm{Sodium Blood Test} A blood test measuring serum sodium, a major extracellular electrolyte involved in water balance and nerve and muscle function; results are interpreted with fluid status and other electrolytes.

\medterm{Sodium Carbonate Poisoning} Poisoning from sodium carbonate, an alkaline chemical found in some cleaning products. It may irritate or burn the eyes, skin, mouth, and digestive tract; do not induce vomiting after swallowing it.

\medterm{Sodium Chloride} The salt sodium chloride; sterile solutions are used for fluid replacement, irrigation, and drug dilution.
\medterm{Sodium Cromoglycate} A mast-cell stabilizer used in selected allergy and asthma treatments; also called cromolyn.
\medterm{Sodium Diatrizoate} A water-soluble iodinated contrast agent used for selected radiographic examinations. It can cause allergic-like reactions, kidney stress, or aspiration complications, so its use depends on the procedure and patient risks.
\medterm{Sodium Hydroxide Poisoning} Severe poisoning from sodium hydroxide, a strong alkali in some drain and oven cleaners. It can rapidly burn the mouth, throat, esophagus, stomach, skin, and eyes; emergency assessment is required and vomiting must not be induced.

\medterm{Sodium Hypochlorite} The active chemical in many bleach and disinfectant products. It kills many germs when used correctly, but concentrated liquid can burn skin, eyes, and the digestive tract; mixing it with acids or ammonia can release dangerous gas.
\medterm{Sodium Hypochlorite Poisoning} Poisoning from bleach containing sodium hypochlorite. Swallowing may irritate or burn the mouth and digestive tract; mixing bleach with acids or ammonia can release dangerous gases causing cough, chest pain, or breathing failure.

\medterm{Sodium In Diet} Dietary sodium helps regulate fluid balance and nerve and muscle function, but excess intake can raise blood pressure; processed foods are a major source and intake should follow individual health advice.

\medterm{Sodium Low Levels In The Blood (Hyponatremia)} Alternate index label for Hyponatremia.

\medterm{Sodium Polystyrene Sulfonate} A resin that binds potassium in the intestine and has been used to lower high blood potassium. It can cause constipation and serious intestinal injury, so its use requires careful medical assessment.

\medterm{Sodium Urine Test} A urine test measuring sodium excretion, used with blood tests and clinical findings to assess volume status, kidney handling of sodium, and causes of hyponatremia or hypertension.

\medterm{Sodium Valproate} An anticonvulsant and mood stabilizer; it carries important pregnancy and reproductive risks and requires specialist prescribing.
\medterm{Soft Markers For Aneuploidy} Ultrasound findings during pregnancy that are common in healthy fetuses but occur somewhat more often when there is an extra or missing chromosome. They do not diagnose a chromosome condition and are interpreted with other screening results.

\medterm{Soft Tissue Augmentation} Restoration or enhancement of volume and contour in skin or underlying soft tissue using fillers, fat transfer, implants, or other reconstructive techniques.

\medterm{Soft Tissue Injury} Damage to muscles, tendons, ligaments, fascia, skin, or other non-bony tissues, often caused by
trauma or overuse.

\medterm{Soft Tissue Rheumatism} An older broad term for pain or inflammation affecting muscles, tendons, ligaments, or bursae. Modern diagnosis usually names the specific condition, such as tendinitis, bursitis, or a muscle strain.

\medterm{Soft Tissue Sarcoma} A group of malignant tumors arising in muscle, fat, fibrous tissue, blood vessels, or
other connective tissues. Treatment depends on the subtype, location, stage, and relationship to nearby structures.

\synonyms
Cancer, Soft Tissue Sarcoma
Sarcoma, Soft Tissue

\medterm{Soft Tissue Tumors} Growths arising in muscle, fat, fibrous tissue, blood vessels, nerves, or other supporting tissues. Many are harmless, but some are cancerous; a growing, deep, firm, painful, or persistent lump needs assessment.

\medterm{Soft-Tissue Calcinosis} Buildup of calcium deposits in skin, muscles, tendons, or other soft tissues. The deposits may form hard lumps, cause pain or limited movement, break through the skin, or occur because of kidney disease, abnormal calcium levels, inflammation, or tissue damage.

\medterm{Soil Contamination} The presence of harmful substances in soil, such as heavy metals, pesticides, fuel, or industrial waste. People may be exposed through food, water, dust, or skin; testing, removal, covering, or treatment of the soil can reduce risk.

\medterm{Solar Keratosis} Alternate index label; see \textit{Actinic keratosis}.

\medterm{Solar Keratosis (Actinic Keratosis)} Alternate index label for Actinic Keratosis.

\medterm{Solar Lentigines} Alternate index label; see \textit{Age spots (liver spots)}.

\medterm{Solar Plexus} A network of involuntary nerves in the upper abdomen, centered around the celiac plexus. It helps regulate abdominal organs, and irritation or injury may contribute to upper-abdominal pain or altered organ function.
\medterm{Solar Urticaria} A rare allergic-type reaction to sunlight or other light that causes itchy red raised patches within minutes on exposed skin. The rash usually improves after leaving the light, but severe widespread symptoms need medical care.

\medterm{Solder Poisoning} Poisoning from swallowing solder or inhaling fumes or dust during heating. Products may contain lead, tin, antimony, or other metals and can cause stomach symptoms, anemia, nerve injury, or kidney damage; exposure needs evaluation.

\medterm{Sole Sweating Excessive (Hyperhidrosis)} Alternate index label for Hyperhidrosis.

\medterm{Soleus Muscle} A deep calf muscle that joins the Achilles tendon and helps point the foot downward. It is important for standing, walking, and running; injury or overuse can cause calf pain and weakness.

\medterm{Solitary Fibrous Tumor} A solitary fibrous tumor is a usually slow-growing growth of connective tissue that can arise in the pleura, meninges, or other sites and may rarely behave aggressively.

\medterm{Solitary Pulmonary Nodule} A single small spot in the lung surrounded by otherwise normal lung tissue. It may be a scar, infection, harmless growth, or early cancer; its size, appearance, and change over time determine whether follow-up or biopsy is needed.


\medterm{Solution} A homogeneous mixture in which one or more substances are dissolved in a solvent; medical solutions include saline and antiseptics.
\medterm{Solvent Abuse (Misuse)} Harmful inhalation or ingestion of volatile solvents, which can cause neurologic, cardiac, liver, kidney, or respiratory injury.
\medterm{Solvent Exposure} Contact with organic solvents by inhalation, ingestion, or skin absorption. Acute exposure may cause
irritation or central nervous system effects; repeated exposure may damage organs.

\medterm{Somatic} Relating to the body, especially the body wall and musculoskeletal system, or to physical symptoms.
\medterm{Somatic Symptom And Related Disorders} A group of mental disorders involving distressing physical symptoms, excessive health-related thoughts or behaviors, illness fears, functional neurologic symptoms, or intentionally produced illness. Examples include somatic symptom disorder, illness anxiety disorder, functional neurologic disorder, and factitious disorder. Malingering is separate because it involves intentional symptom production for an external benefit.
\medterm{Somatic Mutation} A genetic change acquired by a body cell during life rather than inherited from a parent. It is passed to descendants of that cell and may contribute to cancer, but it usually is not passed to children.

\medterm{Somatic Symptom Disorder} A mental disorder in which one or more physical symptoms cause significant distress or disruption of daily life, together with excessive thoughts, feelings, or behaviors related to the symptoms. The symptoms are real and may occur with or without an identified medical condition; their severity is not determined solely by whether a medical cause is found.

\medterm{Somatic Variant} A change in DNA acquired by a body cell after conception rather than inherited from a parent. It usually cannot be passed to children but may affect only some tissues or contribute to cancer.

\medterm{Somatization} The experience or expression of psychological distress through physical symptoms such as pain, fatigue, or stomach problems. The symptoms are real and not intentional; assessment should consider both physical and psychological causes.

\medterm{Somatostatin} A hormone made in the brain, pancreas, and digestive tract that slows growth hormone, insulin, glucagon, and digestive-hormone release. It also reduces movement and secretion within parts of the digestive system.
\medterm{Somatostatin Analog} A medicine that imitates somatostatin, a hormone that slows the release of several hormones and digestive secretions. It is used for selected hormone-producing tumors, acromegaly, severe diarrhea, and some bleeding conditions.

\medterm{Somatostatinoma} A rare hormone-producing tumor usually arising in the pancreas or small intestine. It may cause diabetes, diarrhea or greasy stools, gallstones, weight loss, or no symptoms; treatment depends on whether it has spread and may include surgery or cancer medicines.

\medterm{Somatotropic Pituitary Adenoma} A usually noncancerous pituitary tumor that makes too much growth hormone. In adults it causes acromegaly, with enlargement of the hands, feet, jaw, or other tissues; in children it can cause excessive height before growth ends.

\medterm{Somatotropin} Growth hormone, produced by the pituitary gland. It supports childhood growth and influences protein, fat, and carbohydrate metabolism; excess or deficiency can cause characteristic growth and metabolic disorders.
\medterm{Somatotype} A descriptive body-build classification; it is not a disease diagnosis.
\medterm{Somnambulism} Sleepwalking or performing other activities while only partly awake, usually during deep sleep. Episodes are more common in children and may be triggered by lack of sleep, stress, illness, or medicines; safety measures are important, especially if wandering or injury occurs.
\medterm{Somniloquy (Sleep-Talking)} Talking, laughing, or making sounds during sleep without knowing it. It is usually harmless, but frequent or disruptive episodes may occur with sleepwalking, nightmares, breathing problems during sleep, medicines, or other sleep disorders.

\synonyms
Sleep talking

\medterm{Somnolence} Unusual sleepiness or a tendency to fall asleep when a person should be awake. It may result from too little sleep, medicines, alcohol or drugs, sleep disorders, metabolic illness, or neurologic disease.

\medterm{Somnolent} Somnolent means unusually sleepy or drowsy, which may result from sleep deprivation, medicines, intoxication, metabolic illness, or neurologic disease.

\medterm{Sonographer} A healthcare professional trained to obtain diagnostic ultrasound images. They position the patient, operate the scanner, record appropriate views and measurements, and document findings for interpretation by a qualified clinician.
\medterm{Soporifics} Medicines or substances that cause sleepiness or promote sleep. They may impair alertness, coordination, and breathing, especially when combined with alcohol or other sedatives.
\medterm{Sorbitol} A sugar alcohol used as a sweetener and in some medicines. Larger amounts draw water into the bowel and act as a laxative, but may cause gas, cramps, diarrhea, or dehydration.
\medterm{Sore} Painful or tender, or describing an open area on skin or a moist body lining. A sore may result from injury, infection, inflammation, pressure, allergy, or another medical condition.
\medterm{Sore Throat} Pain, scratchiness, or irritation in the throat. Common causes include viral infections, bacterial infections, allergies, dry air, smoke, voice strain, and acid reflux; persistent, severe, or breathing-related symptoms need medical assessment.
\medterm{Sore Throat (Pharyngitis)} Alternate index label for Pharyngitis.


\medterm{Sores Canker (Canker Sores)} Alternate index label for Canker Sores.


\medterm{Sotalol} A medicine with beta-blocking and class III antiarrhythmic effects, used for selected cardiac rhythm disorders.
\medterm{Sotos Syndrome} A genetic disorder that causes unusually rapid growth in childhood, a large head, distinctive facial features, and delayed development or learning difficulties. Coordination, behavior, seizures, heart problems, or kidney problems may also occur. Growth usually slows with age. A change in the \textit{NSD1} gene causes most cases.

\synonyms
Cerebral Gigantism; NSD1-Related Overgrowth Syndrome

\medterm{Sound} Mechanical energy transmitted as waves and perceived by hearing; clinical sound assessment includes speech and respiratory or cardiac sounds.
\medterm{Source Control} A procedure or intervention that eliminates, drains, debrides, or otherwise controls an
anatomic focus of infection or contamination. Examples include drainage of an abscess,
removal of infected tissue or device, and repair of a perforation.

\medterm{Soy} Soy is a legume providing protein and isoflavones; it is nutritious for many people but can cause allergy and may interact with some medicines or thyroid treatment timing.

\medterm{Soy Protein Intolerance} Adverse gastrointestinal reaction to soy proteins, most often seen in infants and young
children. It may cause vomiting, diarrhoea, abdominal pain, blood or mucus in the stool,
feeding difficulty, or poor growth.

\medterm{Soya} Soybean and products made from it; soy is a food source and can cause allergy in susceptible people.
\medterm{Space Medicine} The medical specialty concerned with human health and performance during aviation and space travel. It addresses altered gravity, radiation, isolation, motion sickness, sleep, emergencies, exercise, and returning safely to Earth.
\medterm{Spacer} A chamber attached to a metered-dose inhaler that holds the released medicine briefly, making it easier to breathe into the lungs and reducing medicine left in the mouth or throat.
\synonyms
Valved holding chamber

\medterm{Spanish Fly} A historical name for cantharidin-containing blister beetles; ingestion is poisonous and it is not a safe aphrodisiac.
\medterm{Spasm} An involuntary, often painful tightening of one muscle or a group of muscles. Spasms may result from fatigue, dehydration, nerve disease, injury, medicines, inflammation, or problems affecting the brain or spinal cord.
\medterm{Spasmodic Dysphonia} Spasmodic dysphonia is a focal laryngeal dystonia causing involuntary voice-box spasms and strained, breathy, or interrupted speech.

\medterm{Spasmodic Torticollis} A form of neck dystonia with involuntary muscle contractions causing twisting, tilting, or abnormal neck posture.
\medterm{Spasmus Nutans} An uncommon early-childhood syndrome of rapid shimmering eye movements, head nodding, and torticollis that usually resolves but requires evaluation to exclude neurologic or ocular disease.

\medterm{Spastic} Describes muscles that are unusually tight and resist being stretched, especially when moved quickly. It may cause stiffness, jerking, abnormal posture, and difficulty controlling movement after nervous-system injury.
\medterm{Spastic Colon (Irritable Bowel Syndrome)} Alternate index label for Irritable Bowel Syndrome.

\medterm{Spasticity} Velocity-dependent increase in muscle tone caused by an upper motor-neuron lesion, producing stiffness, spasms, abnormal reflexes, and restricted movement.

\medterm{Spasticity Management} Assessment and treatment of abnormal muscle stiffness and spasms. Measures may include
stretching, therapy, positioning, medicines, injections, or surgery.

\medterm{Spatial Resolution} The ability of an imaging system to show two nearby small objects as separate. Better resolution shows finer detail, but the result depends on the scanner, movement, body part, scan settings, and the type of imaging used.

\medterm{Special Care Baby Unit} A ward within a hospital that provides specialist care for newborn babies who need
additional support but not intensive care.

\medterm{Special Stain} A histochemical stain that highlights a particular tissue component, organism, or
deposited substance not adequately demonstrated by routine hematoxylin and eosin.
Examples include periodic acid-Schiff, Ziehl-Neelsen, Congo red, and Gomori methenamine
silver stains.

\medterm{Specific Absorbed Fraction} The portion of energy released by radiation in one body area that is absorbed by a particular organ or tissue. It helps calculate the radiation dose during nuclear-medicine treatment or testing.

\medterm{Specific Activity} The amount of radioactivity in a given amount of material. It matters when preparing radioactive medicines because enough activity is needed for imaging or treatment while limiting the amount of substance and unnecessary exposure.

\medterm{Specific Gravity} A comparison of urine density with water that gives an estimate of how concentrated the urine is. It may rise with dehydration, sugar, or protein in urine and fall when the kidneys cannot concentrate urine or when excess fluid is present.

\medterm{Specific Learning Disorder} A developmental condition causing lasting difficulty learning or using reading, writing, or mathematics skills despite appropriate teaching. The difficulty interferes with school, work, or daily activities and is not explained only by poor instruction, vision, hearing, or intellectual disability.

\medterm{Specific Phobia} A disorder involving marked, persistent fear or anxiety about a particular object or situation.
The fear is excessive relative to the actual danger, leads to avoidance or intense distress,
and interferes with daily functioning.

\medterm{Specific Reading Disability} Alternate index label; see \textit{Dyslexia}.

\medterm{Specificity} A test's ability to correctly identify people who do not have the condition being tested.
\medterm{Specimen Adequacy} The degree to which a collected clinical specimen is sufficient in quantity, correctly
identified, properly transported, and representative of the site being tested.
Inadequacy can cause invalid, misleading, or falsely negative results.

\medterm{Specimen Container} A clean, suitable container used to collect, identify, protect, and transport a patient sample for laboratory testing. The container depends on the sample and test; it should be labeled correctly, closed securely, and kept free from contamination.

\medterm{Specimen Transport Medium} A special liquid or gel that keeps germs or cells suitable for testing while a patient sample travels to the laboratory. The medium, temperature, container, and delivery time depend on the sample and suspected infection.

\medterm{SPECT Scan (Single Photon Emission CT)} A nuclear-medicine scan that uses a small radioactive tracer and a rotating camera to create three-dimensional pictures of blood flow or organ function. It can examine the heart, brain, bones, kidneys, or other organs.

\synonyms
Single Photon Emission CT

\medterm{Spectral Analysis} Examination of the different energy levels within radiation or another signal. In medical imaging, it can separate useful information from background scatter and help identify materials or improve image interpretation.

\medterm{Speculum} An instrument used to gently separate or hold open the walls of a body passage so that the area can be visualized, examined, sampled, or treated. Common types include vaginal, nasal, ear, and rectal specula; design and size depend on the body site and intended procedure.
\medterm{Speculum, Ear} A small cone-shaped tip fitted to an otoscope to direct light into the external ear canal and permit examination of the canal and tympanic membrane. Disposable and reusable types are available in different sizes.

\medterm{Speech} The production of spoken language using breathing, voice, and movements of the tongue, lips, palate, and other structures. Hearing, brain function, teeth, illness, and dental appliances can affect speech.

\medterm{Speech Audiometry} Hearing testing that measures detection or recognition of speech at different loudness levels. It commonly includes
speech reception threshold and word-recognition testing, complementing pure-tone audiometry.

\synonyms
Speech testing; speech audiometric testing.

\medterm{Speech Delay} Speech or language development that is slower than expected for a child's age. Causes may include hearing loss, developmental differences, autism, neurologic disease, or limited language access; loss of previously learned skills needs prompt assessment.

\medterm{Speech And Language Delay} Slower-than-expected development of speech sounds, vocabulary, understanding, or use of language in a child. Causes may include hearing loss, developmental disorders, autism, neurologic disease, or other medical and environmental factors.

\medterm{Speech Discrimination} A hearing test that measures how well a person understands spoken words at a comfortable loudness. The result helps assess hearing function beyond the ability to detect pure tones.

\medterm{Speech Disorders} Conditions that affect making speech sounds, speaking smoothly, using the voice, or controlling the movements needed for speech. They may follow hearing loss, developmental differences, brain injury, stroke, nerve disease, or other conditions and can often be assessed by a speech therapist.

\medterm{Speech Impairment In Adults} Difficulty producing or understanding speech caused by aphasia, dysarthria, apraxia, hearing loss, neurologic disease, structural problems, or cognitive disorder; sudden onset requires stroke assessment.

\medterm{Speech Therapy} Assessment and treatment to help with speaking, understanding language, voice, communication, or swallowing. A speech and language therapist may use exercises, practice, communication aids, and advice for the person and family.
\medterm{Speech-Language  Pathologist} A clinician who evaluates and treats disorders of speech, language, voice, communication, cognition, and swallowing across childhood, adulthood, and older age.

\medterm{Speech-Language Pathology Assessment} An evaluation of speech, language, voice, communication, swallowing, and sometimes related cognitive skills. It combines interview, clinical observation, and standardized or functional tests to guide therapy.


\medterm{Sperm} Mature male reproductive cells that can fertilize an egg. Semen contains sperm plus fluids from the reproductive glands; sperm are produced in the testes.
\medterm{Spermarche} The onset of sperm production and the first ejaculation, marking an important stage of male puberty.

\medterm{Spermatic} Relating to sperm or structures that produce, transport, or support sperm. Spermatic structures include the testes, ducts, blood vessels, nerves, and tissues involved in male reproductive function.
\medterm{Spermatic Cord} A cord-like structure carrying the vas deferens, testicular blood vessels, nerves, lymph vessels, and supporting tissue. It extends through the groin to the testis and may be affected by hernia, twisting, or infection.
\medterm{Spermatocele} A usually harmless, fluid-filled cyst beside or within the epididymis, the coiled tube behind the testicle. It often causes no symptoms but may feel like a lump and should be checked if newly noticed or painful.

\medterm{Sperm Retrieval} A procedure that obtains sperm from semen, the epididymis, or testicular tissue for fertility treatment or laboratory testing. The method depends on whether sperm are present in the semen and on the cause of infertility.

\synonyms
Cyst, Spermatic

\medterm{Spermatogenesis} The production and maturation of sperm in the seminiferous tubules of the testes. It depends on germ cells, Sertoli cells, hormones, and suitable testicular temperature.
\medterm{Spermatogonium} An early male reproductive cell in the testis that divides and develops into primary spermatocytes, beginning sperm production.

\medterm{Spermatorrhea} Involuntary leakage or discharge of semen, sometimes occurring during sleep or without orgasm. Occasional nocturnal emissions can be normal; persistent discharge, pain, blood, or urinary symptoms needs assessment.

\medterm{Spermatozoa} Mature male reproductive cells, also called sperm. Each normally carries half the usual chromosome number and has a head, middle section, and tail that helps it move toward an egg.
\medterm{Spermatozoon} A single mature male reproductive cell, or sperm cell. The singular form of spermatozoa: one mature male reproductive cell capable of carrying genetic material to an egg.
\medterm{Spermicide} A contraceptive substance that kills or immobilizes sperm, often used with condoms or other barrier methods.
\medterm{Sphenoid} Relating to the sphenoid bone, a central bone at the base of the skull. It helps form the eye sockets and skull base and contains the hollow spaces called the sphenoid sinuses.
\medterm{Sphenoid Bone} A central bone at the base of the skull that helps form the eye sockets and skull base. It contains the sphenoid sinuses and a hollow that holds the pituitary gland.

\medterm{Sphenoid Sinusitis} Inflammation or infection of the sphenoid sinus, which lies deep behind the nasal cavity. It may cause
deep headache, facial pain, or visual symptoms and can rarely spread to nearby nerves or the brain;
persistent or severe symptoms require medical assessment.

\medterm{Sphenoidal Sinus} A paired air-filled sinus within the sphenoid bone at the base of the skull. It lies near important nerves and blood
vessels, so inflammation or surgery in this area requires careful assessment.

\medterm{Spherocytes} Red blood cells that are unusually round and dense instead of having their normal flattened shape. They may occur when the cell membrane is damaged or removed, especially in hereditary spherocytosis or immune red-cell destruction.

\medterm{Spherocytosis (Hereditary, HS)} Alternate index label for Hereditary, HS.

\medterm{Sphincter} A ring-shaped muscle that opens and closes a body passage. Sphincters help control urine, stool, food, bile, and other fluids; examples include the muscles around the anus, bladder outlet, and stomach outlet.

\medterm{Sphincter Of Oddi} A muscular valve controlling the release of bile and pancreatic juice into the first part of the small intestine. It opens during digestion and abnormal blockage or tightening can cause pain or pancreatitis.

\medterm{Sphygmograph} An instrument or sensor that records the shape and timing of the pulse in an artery. The tracing can provide information about circulation and arterial stiffness.
\medterm{Sphygmomanometer} A device used to measure arterial blood pressure with an inflatable cuff and pressure-measuring system. Manual versions commonly use an aneroid or mercury manometer with a stethoscope, while electronic versions use sensors to analyze pressure oscillations and display the reading.

\synonyms
Blood Pressure Gauge
\medterm{Spider Angioma} A spider angioma is a central red blood vessel with radiating fine vessels that blanches with pressure and may occur normally or with liver or hormonal disease.

\medterm{Spider Bite} An injury caused by a spider, usually producing a small painful or itchy area. Some species can cause worsening tissue damage, severe muscle cramps, sweating, vomiting, breathing problems, or other serious effects; rapidly worsening symptoms need emergency care.

\medterm{Spider Bites (Black Widow And Brown Recluse)} Alternate index label for Black Widow and Brown Recluse.

\medterm{Spina Bifida} A birth defect in which the neural tube and the bones protecting the spinal cord do not close completely during early development. It can range from mild to severe and may cause weakness, loss of feeling, or problems controlling the bladder or bowel.

\medterm{Spina Bifida Meningocele} A form of spina bifida in which the protective coverings around the spinal cord bulge through an opening in the backbone and form a fluid-filled sac. The spinal cord is usually not inside the sac, so weakness or loss of bladder and bowel control may be less severe.

\medterm{Spina Bifida Myelomeningocele} The most severe common form of spina bifida, in which the spinal cord and its coverings bulge through an opening in the backbone. It can cause weakness or paralysis, loss of feeling, and difficulty controlling the bladder or bowel below the affected area.

\medterm{Spina Bifida Occulta} A usually mild form of spina bifida in which the vertebral arch does not fully close
without herniation of the spinal cord or meninges. It is often asymptomatic but may be associated with a skin dimple,
hair tuft, or other finding over the lower spine.

\medterm{Spinal Analgesia} Pain relief produced by injecting numbing medicine, sometimes with an opioid, into the fluid around the lower spinal nerves. It may cause temporary leg weakness, itching, difficulty urinating, low blood pressure, headache, or slowed breathing.

\medterm{Spinal And Bulbar Muscular Atrophy} An inherited disorder that gradually damages the nerve cells controlling movement. It usually affects adult men, causing weakness and wasting of the arms, legs, face, and throat, with muscle twitching, difficulty speaking or swallowing, and sometimes breast enlargement or reduced fertility. Treatment supports movement, swallowing, breathing, and daily activities.

\synonyms
Kennedy Disease; SBMA

\medterm{Spinal And Epidural Anesthesia} Spinal and epidural anesthesia deliver local anesthetic near spinal nerves to numb the lower body; spinal is a single injection, while epidural uses a catheter for adjustable dosing.

\medterm{Spinal Anesthesia} An injection of numbing medicine into the fluid around the lower spinal nerves to temporarily block pain and movement in the lower body. Possible effects include low blood pressure, difficulty urinating, headache, infection, bleeding, or rarely excessive numbness and breathing problems.

\synonyms
Spinal Anaesthesia

\medterm{Spinal Anesthesia For Cesarean Delivery} A regional anesthetic injected into the fluid around the lower spinal nerves to numb the lower body during a cesarean birth. The mother remains awake, while blood pressure, breathing, and the baby's condition are monitored.

\medterm{Spinal Angiography} X-ray imaging of blood vessels supplying the spinal cord after contrast is injected. It is used to locate vascular malformations, fistulas, aneurysms, or sources of bleeding and may guide treatment.

\medterm{Spinal Artery} An artery supplying the spinal cord and its coverings, including the anterior spinal artery and paired posterior spinal arteries; blockage can cause characteristic neurologic deficits.

\medterm{Spinal Block} A spinal block is injection of local anesthetic into cerebrospinal fluid in the lower back, producing temporary numbness and weakness for surgery or pain control.

\medterm{Spinal Canal} The spinal canal is the bony passage through the vertebral column containing the spinal cord, its coverings, nerve roots, and blood vessels.

\medterm{Spinal Cavity} The posterior body cavity formed by the vertebral column that contains and protects the spinal cord.

\medterm{Spinal Column} The vertebral column, consisting of vertebrae, discs, ligaments, and joints that support and protect the spinal cord.
\medterm{Spinal Cord} The cylindrical nervous tissue within the vertebral canal that carries signals between brain and body and mediates reflexes.
\medterm{Spinal Cord Abscess} A spinal cord or epidural abscess is a collection of infection near the spinal cord that can compress neural tissue and cause pain, fever, weakness, paralysis, or bladder problems.

\medterm{Spinal Cord Compression} Pressure on the spinal cord or its nerve roots, often from a tumor, infection, bleeding, fracture, or slipped disc. It can cause severe back pain, weakness, numbness, walking difficulty, or loss of bladder or bowel control and requires urgent medical assessment.

\medterm{Spinal Cord Disease} A disorder affecting the spinal cord, including compression, inflammation, vascular injury, infection, tumor, degeneration, or trauma, and potentially causing weakness, sensory loss, or autonomic dysfunction.

\medterm{Spinal Cord Injury} Damage to the spinal cord that can cause weakness, paralysis, altered sensation, or loss of bladder and bowel control. Severity is often graded from complete to incomplete using a standardized neurological examination.

\synonyms
SCI

\medterm{Spinal Cord Injury: Treatments And Rehabilitation} Spinal-cord injury can cause temporary or permanent weakness, sensory loss, autonomic dysfunction, and bowel or bladder problems. Acute care protects the spine and treats complications; rehabilitation addresses mobility, skin, breathing, function, equipment, sexuality, and community participation.

\medterm{Spinal Cord Neoplasms} Tumors arising in the spinal cord, its coverings, or nearby tissues. They may be noncancerous or cancerous and can cause pain, weakness, numbness, walking problems, or bladder and bowel changes by pressing on nerves.

\medterm{Spinal Cord Stimulation} An implanted treatment that delivers electrical impulses to the epidural space to alter pain signaling, used for selected chronic neuropathic or ischemic pain after careful assessment.

\medterm{Spinal Cord Trauma} Spinal cord trauma is injury to the cord from force such as fracture, dislocation, or penetrating damage, potentially causing weakness, sensory loss, autonomic dysfunction, or paralysis.

\medterm{Spinal Curvature} Alternate index label; see \textit{Scoliosis}.

\medterm{Spinal Disease} Any disorder affecting the vertebrae, discs, joints, ligaments, spinal cord, or nerve roots. It may cause back or neck pain, stiffness, weakness, numbness, walking difficulty, or changes in bladder or bowel control.
\medterm{Spinal Dysraphism} A group of birth defects caused by incomplete development or closure of the spine and spinal cord before birth. Some forms have an open sac on the back, while others are hidden and may later cause weakness, pain, or bladder and bowel problems.

\medterm{Spinal Epidural Abscess} A collection of pus between the covering of the spinal cord and the surrounding bone, usually caused by infection. Severe back pain, fever, weakness, numbness, difficulty walking, or loss of bladder or bowel control requires emergency assessment.

\medterm{Spinal Fluid Leak} Alternate index label; see \textit{CSF leak (Cerebrospinal fluid leak)}.

\medterm{Spinal Fracture} A fracture of one or more vertebrae caused by trauma, osteoporosis, tumor, or other bone-weakening disease.

\medterm{Spinal Fusion} Surgery that joins two or more spinal bones so they heal into one stable segment. Bone and often screws, rods, or cages are used; the operation may reduce movement and pain from selected conditions but can cause infection, nerve injury, or problems at nearby levels.

\medterm{Spinal Ganglion} A small group of nerve-cell bodies beside the spinal cord. It contains sensory nerve cells that carry information about touch, pain, temperature, and body position from the skin, muscles, and joints.

\medterm{Spinal Injury} Damage to the spinal cord, vertebrae, discs, ligaments, or nearby nerves caused by trauma or another disease. It may cause pain, weakness, numbness, paralysis, or bladder and bowel problems and requires urgent assessment when severe.


\medterm{Spinal Muscular Atrophy} An inherited disorder causing progressive loss of motor neurons and muscle weakness, often beginning in infancy
or childhood but sometimes appearing in adulthood. Severity varies; breathing, swallowing, and mobility
may be affected. Early diagnosis allows genetic counseling and disease-modifying care.

\medterm{Spinal Neoplasms} Tumors arising in the spinal cord, its coverings, vertebrae, or nearby tissues. They may be noncancerous or cancerous and can cause back pain, weakness, numbness, walking difficulty, or bladder and bowel changes when they press on nerves.
\medterm{Spinal Nerve} One of 31 pairs of nerves leaving the spinal cord through openings between the spinal bones. Each carries messages for sensation, movement, and automatic functions between the body and spinal cord.

\medterm{Spinal Precautions} Measures used after possible spinal injury to limit movement until the spine and spinal cord have been assessed. They may include careful positioning, support of the neck, and safe transfer; the exact approach depends on the injury and examination.

\medterm{Spinal Shock} A temporary loss of reflexes, muscle tone, and automatic functions below a spinal-cord injury. The affected limbs may initially be limp, and blood-pressure, bladder, bowel, or sexual problems can occur while the injury is assessed and treated.

\medterm{Spinal Stenosis} Narrowing of the passage containing the spinal cord or nerve roots. It may cause back pain, numbness, weakness, or leg pain with walking that improves when sitting or leaning forward.

\synonyms
Stenosis, Spinal

\medterm{Spinal Tap} Alternate index label; see \textit{Lumbar puncture}.

\medterm{Spinal Tumor} An abnormal growth in the spinal bones, spinal cord, its coverings, or nearby tissues. It may be noncancerous or cancerous; symptoms depend on its location and may include pain, weakness, numbness, or loss of bladder control.

\medterm{Spindle Pole} One end of the temporary structure formed during cell division. Fibers extending from each pole attach to duplicated chromosomes and pull the copies into the two new cells.

\medterm{Spindle Pole Body} A structure in fungal cells that organizes the fibers used to separate duplicated chromosomes during cell division. It performs a role similar to the centrosome in animal cells.

\medterm{Spindle-Cell Sarcoma} A group of cancers made of long, narrow cells under the microscope. They can arise from muscle, fibrous tissue, or nerve coverings; their behavior and treatment vary by the exact tumor type and location.

\medterm{Spine} The vertebral column, consisting of stacked vertebrae that support the body and protect the spinal cord.

\medterm{Spine And Spinal Cord, Diseases And Injuries Of} Diseases or injuries affecting the spinal bones, discs, ligaments, spinal cord, or nerve roots. They may cause pain, weakness, numbness, balance problems, deformity, or changes in bladder and bowel function.
\medterm{Spine Disorders} Conditions affecting the bones, disks, joints, ligaments, nerves, or spinal cord. They may cause back or neck pain, stiffness, weakness, numbness, balance problems, or changes in bladder and bowel control.

\medterm{Spine Injury} Damage to the bones, discs, ligaments, nerves, or spinal cord. It may cause back pain, weakness, numbness, loss of movement, or bladder and bowel problems and can require emergency care.
\medterm{Spinocerebellar Ataxia} Alternate index label; see \textit{Spinocerebellar Ataxias}.

\medterm{Spinocerebellar Ataxias} A group of rare inherited disorders that progressively damage the cerebellum and sometimes other parts of the nervous system. The cerebellum helps control balance, coordination, and precise movement. These disorders may cause unsteady walking, poor hand coordination, slurred or unclear speech, shaking, abnormal eye movements, muscle stiffness or weakness, and difficulty swallowing. The age of onset, symptoms, and rate of progression vary widely. Many types are inherited in an autosomal dominant pattern, but the inheritance pattern differs among types.

\medterm{Spinous Process} The backward-pointing bony projection of a vertebra that can be felt along the middle of the back. Muscles and ligaments attach to it and help support and move the spine.
\medterm{Spiradenocarcinoma} A rare malignant sweat-gland tumor that may arise in a long-standing spiradenoma, presenting as a rapidly enlarging, painful, or ulcerated nodule.

\medterm{Spiradenoma} A usually benign painful sweat-gland tumor presenting as a dermal nodule; multiple spiradenomas may occur in Brooke-Spiegler syndrome.


\medterm{Spiral Organ} The organ of Corti, a sensory structure in the cochlea containing hair cells that convert sound vibrations into nerve signals.

\medterm{Spiramycin} A macrolide antibiotic used for selected bacterial infections and sometimes during pregnancy for toxoplasmosis. Choice depends on the infection, local guidance, pregnancy stage, allergies, resistance, and other treatment needs.
\medterm{Spirillum Minus} Spirillum minus is a spiral Gram-negative bacterium that causes rat-bite fever, characterized by fever, rash, and migratory joint pain after rodent exposure; culture is difficult and diagnosis may be clinical or molecular.

\medterm{Spirit} A volatile liquid preparation or alcohol; the meaning depends on the medical, pharmaceutical, or general context.
\medterm{Spiritual Pain} Distress involving meaning, purpose, beliefs, identity, or connection, often addressed in palliative and supportive care.
\medterm{Spiritual Support} Support that helps a person address meaning, values, faith, hope, or existential distress during illness; it may be provided by chaplains, spiritual leaders, clinicians, or trusted people.

\medterm{Spirochaetaceae} A family of spiral-shaped bacteria that includes medically important spirochetes such as Borrelia, Treponema, and Leptospira.

\medterm{Spirochaetales Infection} Infection caused by spiral-shaped bacteria in the order Spirochaetales, including syphilis, Lyme disease, relapsing fever, and leptospirosis.

\synonyms
Spirochaetal Infections; Spirochetal Infections

\medterm{Spirochaete} Another spelling of spirochete, a thin, flexible, spiral-shaped bacterium. Important genera include Treponema, Borrelia, and Leptospira.
\medterm{Spirochete} A thin, spiral-shaped bacterium that moves with a corkscrew-like motion. Important examples cause syphilis, Lyme disease, relapsing fever, and leptospirosis; diagnosis depends on the organism and may use blood tests or other samples.

\medterm{Spirogram} A graph recording how much air a person breathes and how quickly it moves during a lung-function test. It helps identify patterns suggesting airway narrowing, restricted lung expansion, or poor effort.

\medterm{Spirometer} An instrument that measures airflow and the volume of air moved during specified breathing maneuvers, especially forced expiration. It does not by itself measure all lung volumes; residual volume and total lung capacity require additional pulmonary-function testing.
\medterm{Spirometry} A breathing test that measures how much air a person can blow out and how quickly. It helps identify narrowing of the airways, as in asthma or chronic obstructive pulmonary disease, and can monitor treatment. The person must follow instructions and blow forcefully several times; results are interpreted with age, height, and symptoms.

\medterm{Spironolactone} An aldosterone antagonist and potassium-sparing diuretic used for heart failure, hypertension, hyperaldosteronism, and other indications.
\medterm{Spit} Saliva or other fluid expelled from the mouth. Increased saliva or repeated spitting may occur with mouth irritation, nausea, reflux, difficulty swallowing, medication effects, or habit.

\medterm{Spitting} Forceful expulsion of saliva or other material from the mouth. New or excessive spitting may occur with oral disease, reflux, nausea, difficulty swallowing, medication effects, or habit.


\medterm{Spitz Nevus} A benign melanocytic nevus, often a pink or pigmented papule in children or young adults, that may resemble melanoma clinically or histologically.

\medterm{Splanchnic} Relating to internal organs, especially abdominal organs and their blood vessels or nerves. Splanchnic circulation supplies the digestive organs, while splanchnic nerves carry involuntary signals that regulate their function.
\medterm{Splanchnic Nerve} A nerve carrying automatic signals between the spinal cord and internal organs, especially organs in the chest and abdomen. It helps regulate digestion, blood-vessel tone, and other functions that do not require conscious control.

\medterm{Spleen} A lymphoid organ in the upper-left abdomen. Its red pulp filters blood and removes aged or abnormal blood cells; its white pulp supports immune responses. It also stores platelets and helps clear blood-borne microorganisms.
\medterm{Spleen Development} The formation and maturation of the spleen before and after birth. The spleen gradually develops its blood-filtering and immune functions, including removal of aged blood cells and responses to germs in the bloodstream.

\medterm{Spleen Disorder} A disease affecting the spleen, such as enlargement, infection, infarction, rupture, inflammation, cyst, or tumor, with possible effects on blood cells and immunity.

\medterm{Spleen Enlarged (Enlarged Spleen)} Alternate index label for Enlarged Spleen.

\medterm{Spleen Hemangioma} A spleen hemangioma is a usually benign blood-vessel tumor of the spleen, often incidental but occasionally associated with enlargement, anemia, or rupture risk.

\medterm{Spleen Hilum} The indented area on the inner surface of the spleen where its blood vessels, nerves, and lymph vessels enter or leave. The splenic vein carries blood from the spleen toward the liver through the portal circulation.

\medterm{Spleen Removal} Surgery to remove the spleen, which may be needed for injury, an enlarged or overactive spleen, blood disease, or a tumor. Without a spleen, the risk of severe infection increases, so vaccination and medical advice about fever are important.


\medterm{Spleen, Diseases Of} Disorders affecting splenic structure or function, including splenomegaly, infarction, infection, cysts, abscess, rupture, vascular disease, and hematologic or infiltrative disorders.
\medterm{Splenectomy} Surgical removal of the spleen, performed for selected traumatic, hematologic, malignant, vascular, or other disorders. It increases the risk of severe infection with encapsulated bacteria and requires preventive planning.
\medterm{Splenic Abscess} A localized collection of pus in the spleen, usually from bloodstream infection, trauma, infarction, or spread from nearby disease and potentially causing fever, left-upper-quadrant pain, and sepsis.

\medterm{Splenic Cyst} A fluid-filled lesion within the spleen, classified as true cysts with an epithelial
lining, including congenital, epidermoid, and dermoid cysts, or pseudocysts without a
lining, typically resulting from trauma, infection, or pancreatitis. Most are
asymptomatic and discovered incidentally on imaging.

\medterm{Splenic Flexure} The bend of the colon where the transverse colon becomes the descending colon, near the spleen.

\medterm{Splenic Flexure Syndrome} Pain, fullness, or bloating in the upper-left abdomen caused by gas collecting at the bend where the transverse colon meets the descending colon. Symptoms may improve after passing gas or stool.
\medterm{Splenic Infarction} Ischemic necrosis of splenic tissue caused by obstruction of splenic arterial or venous blood flow. Causes include embolism, thrombosis, hematologic disease, hypercoagulability, infection, trauma, or vasculitis; it often causes sudden left-upper-quadrant pain and is evaluated with imaging.

\medterm{Splenic Infection} Infection of the spleen, often an abscess caused by bacteria spreading through the bloodstream or from nearby tissue. It may cause fever, chills, left-upper-abdominal pain, and severe illness and usually needs urgent medical treatment.
\medterm{Splenic Neoplasm} An abnormal growth in the spleen. It may be noncancerous, cancerous, or cancer that has spread from another organ; symptoms and treatment depend on its type, size, and effect on nearby tissues or blood cells.

\medterm{Splenic Rupture} A tear of the spleen, usually after blunt abdominal trauma but sometimes spontaneous in an enlarged or diseased spleen. It can cause life-threatening internal bleeding; diagnosis and treatment depend on hemodynamic stability and imaging, with urgent surgery or embolization sometimes required.

\medterm{Splenomegaly} Abnormal enlargement of the spleen, usually secondary to infection, liver or portal-venous disease, hemolysis, hematologic malignancy, storage disease, or inflammation. It may cause left-upper-quadrant fullness or pain and increases rupture risk.
\medterm{Splenomegaly Gaucher (Gaucher Disease)} Alternate index label for Gaucher Disease.

\medterm{Splice-Site Variant} A variant affecting the sequences required for accurate removal of introns and joining
of exons during RNA processing. It may cause abnormal transcripts, exon skipping, intron
retention, or reduced functional protein.

\medterm{Spliceosome} The spliceosome is a nuclear RNA-protein complex that removes introns from pre-messenger RNA and joins exons to form mature RNA.

\medterm{Splint} A device that supports and limits movement of an injured or painful body part. Splints may be temporary or definitive and should not be so tight that they impair circulation or sensation.

\medterm{Splint (Dental)} Fixation of mobile teeth or fractured tooth. Wire and composite, orthodontic bracket splint. A dental device or wire used to stabilize loose teeth or a fractured tooth while healing occurs.
\synonyms
Dental splint; tooth splint

\medterm{Splinter Haemorrhages} Small, narrow lines of bleeding beneath a nail. They often follow minor injury but can also occur with heart-valve infection, blood-vessel inflammation, psoriasis, or other illnesses and need context for interpretation.
\medterm{Splinter Hemorrhages} Thin, lengthwise lines of bleeding beneath a nail. Minor injury is common, but they can also occur with heart-valve infection, blood-vessel inflammation, psoriasis, or other illnesses; their meaning depends on the overall clinical picture.

\medterm{Splinter Removal} Taking a piece of wood, metal, glass, or another material out of the skin. Small superficial splinters may be removed with clean tweezers, while deep, painful, infected, or eye- or nail-area splinters need clinical care.

\medterm{Splints} Devices that support, immobilize, align, or protect an injured body part, reducing pain and allowing healing.
\medterm{Split Brain} A state after surgical or pathological disruption of the corpus callosum in which information transfer between cerebral hemispheres is reduced, producing characteristic disconnection findings.

\medterm{Split-Thickness Skin Graft} A graft containing the epidermis and part of the dermis, used to cover wounds or burns and allowing the donor site to heal from remaining skin structures.

\medterm{Splitting} A psychological defense pattern in which a person or situation is viewed as entirely good or entirely bad, with difficulty integrating
positive and negative qualities. It can occur in several mental-health conditions and is
not diagnostic by itself.

\medterm{Spondee} A two-syllable word with equal emphasis on both parts, used during speech-hearing tests. Examples help measure the quietest loudness at which a person can recognize spoken words accurately.

\medterm{Spondylarthritis} A family of inflammatory rheumatic diseases affecting the spine, sacroiliac joints, entheses, peripheral joints, skin, eyes, or bowel, including ankylosing spondylitis and psoriatic arthritis.

\medterm{Spondylitis} Inflammation of one or more vertebrae or parts of the spine. Causes include infection, autoimmune disease, and inflammatory arthritis; pain, stiffness, fever, or neurologic symptoms may occur.
\medterm{Spondyloarthritis} A group of inflammatory joint diseases that often affect the spine and joints between the spine and pelvis. They may also cause tendon pain, swollen fingers or toes, eye inflammation, psoriasis, or bowel inflammation.

\medterm{Spondylolisthesis} Slipping of one spinal bone forward or backward compared with the bone below it. It may cause back pain, stiffness, leg pain, numbness, or weakness, although some people have no symptoms.
\medterm{Spondylolysis} A small stress crack or defect in part of a spinal bone, usually in the lower back. It may cause back pain in adolescents or athletes and can sometimes allow one spinal bone to slip on another.

\medterm{Spondylosis} Wear-related changes in the bones, discs, and small joints of the spine. It becomes more common with age and may cause stiffness or pain, but many people have no symptoms; nerve pressure can cause numbness or weakness.
\medterm{Spondylosis, Cervical} Alternate index label; see \textit{Cervical spondylosis}.

\medterm{Sponge} A sponge is a porous material used to absorb fluid, clean or dress wounds, or deliver a product; retained surgical sponges are a preventable complication.

\medterm{Sponge (Vascular)} A sterile absorbent material used during blood-vessel procedures to absorb blood, apply pressure, or protect delicate tissue. It must be counted and removed when required to prevent retained-material complications.
\synonyms
Vascular surgical sponge

\medterm{Spongiform Encephalopathy} A serious progressive brain disease in which brain tissue develops tiny sponge-like spaces. Prion diseases are important examples and can cause worsening memory, movement, coordination, behavior, and eventually loss of independence.
\medterm{Spongy Bone} Porous trabecular bone found inside many bones, especially near joints.

\synonyms
Cancellous bone; trabecular bone

\medterm{Spontaneous Abortion} Pregnancy loss before fetal viability that occurs without intentional intervention, commonly caused by chromosomal abnormalities but also by uterine, endocrine, immune, infectious, or other factors.

\medterm{Spontaneous Abortion (Miscarriage)} Unintentional loss of a pregnancy before fetal viability. Causes include chromosomal abnormalities, uterine or cervical problems, infection, and maternal medical conditions.

\synonyms
Miscarriage

\medterm{Spontaneous Bacterial Peritonitis} An infection of ascitic fluid without an identifiable abdominal source, usually occurring in people with cirrhosis and ascites. It may cause fever, abdominal pain, tenderness, or worsening illness. Diagnosis is made by testing ascitic fluid, usually showing at least 250 neutrophils per mm\textsuperscript{3}.

\medterm{Spontaneous Coronary Artery Dissection} A tear within the wall of an artery supplying the heart, which can narrow or block blood flow and cause a heart attack. It often affects women and may be associated with pregnancy, severe stress, or fibromuscular dysplasia. Treatment depends on blood flow and stability; medicines or a stent may be used, and emergency care is needed for chest pain.

\synonyms
SCAD


\medterm{Spontaneous Breathing Trial} A monitored period during which a person on a breathing machine does most of the breathing with little or no machine assistance. It helps clinicians judge whether the person may breathe safely after the tube is removed.

\medterm{Spontaneous Perforation} Unexpected rupture of a hollow organ without a clear external injury, allowing contents to escape and potentially causing peritonitis, mediastinitis, sepsis, or hemorrhage.

\medterm{Spontaneous Pneumothorax} Air leaking into the space around a lung without a direct injury. It may occur in an otherwise healthy person or because of lung disease and can cause sudden chest pain or breathlessness.

\medterm{Sporadic} Occurring irregularly or as an isolated case rather than in a predictable pattern or family cluster.

\medterm{Sporothrix} A genus of environmental fungi that can cause sporotrichosis, usually after inoculation through the skin and sometimes through inhalation or dissemination.

\medterm{Sporothrix Schenckii} A group of fungi that causes sporotrichosis, a typically nodular skin and lymphatic infection acquired from plants, soil, or animal scratches.

\medterm{Sporotrichosis} Sporotrichosis is a fungal infection caused by Sporothrix species, often introduced through plant matter and producing nodules along lymphatic channels.

\medterm{Sporozoa} An older group name for spore-forming protozoan parasites, including Plasmodium.
\medterm{Sporozoites} Moving, infective forms produced by certain parasites. For example, mosquitoes inject Plasmodium sporozoites into people, after which they travel to the liver and begin the malaria infection cycle.
\medterm{Sports Cream Overdose} Excessive use or swallowing of a muscle or joint pain cream, often containing menthol, methyl salicylate, or capsaicin. It may cause skin burns, salicylate poisoning, vomiting, ringing ears, confusion, or seizures.

\medterm{Sports Physical} A preparticipation health evaluation that reviews medical history, symptoms, medications, family risks, and examination findings to identify conditions needing evaluation before sports.

\medterm{Sports Rehabilitation} Treatment and training that help a person recover safe movement, strength, fitness, and confidence after a sports injury. It may include exercises, hands-on care, education, gradual return to activity, and steps to reduce the chance of another injury.

\medterm{Spots Before The Eyes} Visual floaters, flashes, or scotomata perceived in the visual field; sudden onset requires urgent assessment.
\medterm{Spots In Front Of The Eyes} Floaters or visual spots caused by shadows from vitreous opacities, inflammation, bleeding, retinal tears, migraine, or other eye or neurologic conditions; sudden flashes or a curtain-like loss are urgent.

\medterm{Spotted Fever} A group of infections usually spread by ticks and caused by Rickettsia bacteria. They can produce fever, headache, rash, muscle aches, and serious illness requiring prompt diagnosis and antibiotic treatment.
\medterm{Spousal Abuse} Spousal abuse is physical, sexual, emotional, financial, or coercive harm by an intimate partner; immediate safety planning and confidential support may be needed.

\medterm{Spouse} A spouse is a person’s legally or socially recognized partner; involving a spouse in healthcare requires the patient’s consent unless law provides otherwise.

\medterm{Sprain} An injury in which a ligament that supports a joint is stretched or torn. It can cause pain, swelling, bruising, and looseness; severity ranges from minor stretching to a complete tear and joint instability.

\medterm{Sprain Neck (Whiplash)} Alternate index label for Whiplash.

\medterm{Sprained Ankle} Stretching or tearing of one or more ankle ligaments, usually after rolling the ankle inward. It may cause pain, swelling, bruising, and instability.

\medterm{Sprains} Injuries in which ligaments around a joint are stretched or torn. They can cause pain, swelling, bruising, and joint looseness; severity ranges from minor stretching to complete tearing.
\medterm{Spray Drying} Spray drying converts a liquid mixture into a dry powder by atomizing it into a heated gas stream, often used in pharmaceutical and food manufacturing.

\medterm{Sprengel Deformity} A birth difference in which one shoulder blade sits higher than usual because it did not move to its normal position during development. It may cause uneven shoulders, limited arm movement, or neck and back discomfort.

\medterm{Spring Fever (Henoch-Schonlein Purpura)} Alternate index label for Henoch-Schonlein Purpura.

\medterm{Sprue} A condition causing poor absorption of nutrients from the small intestine. The term may describe celiac disease or tropical sprue, which have different causes and require appropriate testing and treatment.
\medterm{Sprue, Nontropical} Nontropical sprue is an older term for celiac disease, an immune reaction to gluten that damages the small-intestinal lining and impairs absorption.


\medterm{Spur Heel (Heel Spurs)} Alternate index label for Heel Spurs.

\medterm{Spur-Cell Hemolytic Anemia} A severe anemia in which abnormal, spiky red blood cells are rapidly removed from the blood, usually because of advanced liver disease. It can cause jaundice, weakness, and worsening anemia.

\medterm{Sputum} Mucus or other material brought up by coughing from the airways or lungs. Its amount, color, smell, and blood content may provide clues, but sputum alone cannot identify the cause; laboratory testing may be needed.
\medterm{Sputum Culture} A laboratory test that grows germs from mucus coughed up from the lungs or collected by suction. It can help identify bacterial, fungal, or mycobacterial infection and guide treatment, but a poor sample can give misleading results.

\medterm{Sputum Direct Fluorescent Antibody Test} A laboratory test that uses glowing antibody markers to look for selected germs in mucus from the lungs. It can provide a rapid result, but culture or molecular testing may also be needed.

\medterm{Sputum Fungal Smear} Examination of coughed-up lung mucus under a microscope to look for fungi. A culture, blood test, or molecular test may be needed to confirm the specific infection.

\medterm{Sputum Gram Stain} A quick microscope test that colors a sputum sample to show some bacteria and inflammatory cells and to judge whether the sample is suitable for culture.

\medterm{Sputum Stain For Mycobacteria} A microscope test that looks for tuberculosis-causing and related bacteria in mucus from the lungs. A positive result supports infection, but a negative result does not exclude it; culture or molecular testing may also be needed.



\medterm{Squamous Cell} A flat cell that forms part of the covering or lining of the skin and many body passages. These cells protect against friction and infection; changes in them can cause inflammation, precancer, or squamous cell cancer in organs such as the skin, mouth, cervix, lungs, or esophagus.

\medterm{Squamous Cell Carcinoma} A cancer of flat cells that line the skin or certain body passages. On the skin it may appear as a rough, scaly, crusted, wart-like, or nonhealing sore and can invade nearby tissue or spread if untreated.

\medterm{Squamous Cell Carcinoma Of The Lung} A type of non-small-cell lung cancer that usually begins in the lining of a larger airway. It is strongly linked to smoking but can occur in nonsmokers. Symptoms may include a persistent cough, coughing blood, chest pain, breathlessness, or repeated chest infections; treatment depends on stage and tumor features.


\medterm{Squamous Cell Carcinoma Of The Skin} A skin cancer arising from squamous cells, often appearing as a scaly, crusted, wart-like, or ulcerated growth.

\synonyms
Cancer, Squamous Cell

\medterm{Squamous Cell Skin Cancer} Squamous cell skin cancer is a malignant keratinocyte tumor often arising on sun-exposed skin and capable of local invasion or metastasis, especially when high risk.

\medterm{Squamous Cells} Flat epithelial cells with scale-like appearance that line surfaces such as skin, mouth, cervix, and esophagus; abnormal numbers or forms can reflect irritation, infection, dysplasia, or cancer.

\medterm{Squamous Intraepithelial Lesion Cervical (Cervical Dysplasia)} Alternate index label for Cervical Dysplasia.

\medterm{Square Foot} A square foot is a unit of area equal to a square measuring one foot on each side, used to describe surfaces, rooms, wounds, or treatment areas.

\medterm{Square Inch} A square inch is a unit of area equal to a square measuring one inch on each side, used for small surfaces or wound measurements.


\medterm{Squint} Misalignment of the eyes so they do not point in the same direction. It may be constant or intermittent and can cause double vision, poor depth perception, or reduced vision in one eye.

\medterm{SSRI} An SSRI is a selective serotonin-reuptake inhibitor antidepressant used for depression and several anxiety-related disorders; effects and withdrawal symptoms require monitoring.

\medterm{SSRIs} Selective serotonin reuptake inhibitors, a group of medicines that increase serotonin signaling in the brain. They are used for depression and several anxiety-related disorders; effects and side effects develop over time, not immediately.
\medterm{St John's Wort} A herbal product made from Hypericum perforatum and sometimes used for depressive symptoms. It can cause side effects and reduce the effectiveness of many medicines, including some contraceptives, anticoagulants, transplant medicines, and HIV treatments.
\medterm{ST Segment} The part of an electrocardiogram between the electrical activation and recovery signals of the heart's lower chambers. Abnormal elevation or depression can occur with a heart attack, inflammation, or other conditions and must be interpreted with symptoms and the full tracing.

\medterm{St Vitus's Dance} An older name for Sydenham chorea, involuntary movements following group A streptococcal infection.
\medterm{ST-Segment Elevation} Elevation of the ECG ST segment above the isoelectric baseline, reflecting abnormal
ventricular repolarization. It may indicate acute coronary occlusion but can also occur
in pericarditis, early repolarization, ventricular aneurysm, or other conditions;
interpretation requires clinical context and serial testing.


\medterm{St. John's Wort} An herbal product sometimes used for depressive symptoms. It can cause side effects and reduce the effectiveness of many medicines, including some contraceptives, anticoagulants, transplant medicines, and HIV treatments.
\medterm{Stab Wound} A deep puncture caused by a pointed or sharp object. The wound may damage blood vessels, nerves, organs, or bone even when the skin opening looks small; heavy bleeding, severe pain, or a wound to the chest or abdomen requires emergency care.

\medterm{Stable Angina} Repeated, predictable chest pressure or discomfort brought on by activity or stress because the heart muscle receives too little blood. It usually improves with rest or prescribed medicine; new, severe, prolonged, or rest pain may be a heart attack and needs emergency care.

\synonyms
Chronic stable angina

\medterm{Stable Coronary Artery Disease} Coronary artery narrowing causing a consistent, predictable pattern of angina.

\synonyms
Stable ischemic heart disease

\medterm{Stable Disease} An oncology response category in which a cancer has neither shrunk enough to meet criteria for partial response nor grown enough to meet criteria for progressive disease. It describes change over a specified assessment interval and does not mean that the cancer is cured or inactive.

\medterm{Stabs} Brief, sharp pains or penetrating injuries described as stabbing. Causes vary with location and may include nerve irritation, muscle strain, infection, inflammation, injury, or serious disease requiring urgent assessment.
\medterm{Stage 4 Prostate Cancer} Alternate index label; see \textit{Metastatic (stage 4) prostate cancer}.

\medterm{Stage I} A cancer stage that generally indicates limited disease, although the exact definition and prognosis depend on the cancer type and staging system.

\medterm{Stage II} A cancer stage that usually indicates a larger or more locally advanced tumor than stage I, without distant spread. The exact tumor size, lymph-node findings, prognosis, and treatment vary by cancer type and staging system.

\medterm{Stage III} A cancer stage that often indicates more extensive local or regional disease, with the precise criteria depending on the cancer type.

\medterm{Stage IV} A cancer stage that generally indicates spread to distant organs or sites. Prognosis and treatment vary by cancer type and individual factors.

\medterm{Stages Of Change} A model describing common steps people may pass through when changing a behavior. These include not yet considering change, thinking about it, preparing, taking action, and maintaining the change; people may move backward or forward between stages.

\synonyms
Transtheoretical model

\medterm{Stages Of Labor} The first stage is cervical opening, the second is birth of the baby, and the third is delivery of the placenta. The first stage has early and active phases; length and symptoms vary, and the uterus is monitored after delivery.

\medterm{Staghorn Calculi} Large branching kidney stones that fill part or all of the kidney's urine-collecting system. They are often linked with infection, can damage kidney tissue, and usually require specialist treatment to remove them.

\medterm{Staghorn Calculus} A branching kidney stone that fills part or all of the renal pelvis and calyces, often associated with infection.
\medterm{Staging} Determining how far a cancer has grown or spread by assessing the original tumor, nearby lymph nodes, and distant organs. The stage helps estimate outlook and choose treatment.

\medterm{Staining} A laboratory method that uses colored dyes to make cells, tissue structures, or microorganisms easier to see under a microscope. Different stains highlight different features and can help guide diagnosis, although results may need culture or other tests for confirmation.

\medterm{Stainless Steel} A corrosion-resistant alloy containing iron and chromium, used for instruments, implants, orthodontic appliances, and some dental crowns. Its suitability depends on the grade, strength, and tissue exposure.

\medterm{Stammering} Alternate name for stuttering, a speech-fluency difficulty involving repeated sounds or words, prolonged sounds, or blocks that interrupt speaking. It is common in childhood and may continue into adulthood; speech-language therapy can help.
\medterm{Standard Deviation} A number that shows how widely measurements differ from their average. A small standard deviation means values are close together, while a large one means they are more spread out; it does not by itself show whether results are good or bad.
\medterm{Standard Eye Exam} A comprehensive eye examination assessing visual acuity, refraction, pupils, eye movements, external structures, anterior eye, intraocular pressure, and ocular fundus as indicated.

\medterm{Standard Of Care} The conduct, skill, diligence, and judgment a reasonably prudent practitioner of the
same specialty would exercise under similar circumstances; the benchmark for medical
negligence. It is established through guidelines, peer practice, and expert testimony
rather than fixed statute.

\medterm{Standard Precautions} Basic infection-control measures used for every patient, because blood, body fluids, nonintact skin, and mucous membranes may carry germs. They include hand hygiene, appropriate protective equipment, safe injections, sharps safety, and respiratory etiquette.
\medterm{Standardized Extract} An herbal extract adjusted to contain a measured amount of selected ingredient, improving consistency but not guaranteeing safety or effectiveness.

\medterm{Standardized Uptake Value} A number used on a PET scan to estimate how much radioactive tracer a tissue has taken up, adjusted for the injected amount and body size. It can help compare areas or follow change, but it does not by itself prove cancer or another disease.

\medterm{Stannosis} A usually mild form of pneumoconiosis caused by long-term inhalation of tin oxide dust. It may produce visible lung changes on imaging, often without major breathing impairment.
\medterm{Stanozolol} A synthetic anabolic androgenic steroid with limited medical use and potential liver, lipid, and cardiovascular toxicity.
\medterm{Stapedectomy} Surgical removal or replacement of the stapes, usually to treat otosclerosis and conductive hearing loss.
\medterm{Stapedius} A small middle-ear muscle that helps limit movement of the stapes and may protect the inner ear from loud sounds.

\medterm{Stapes} The smallest middle-ear bone, passing sound vibrations from the incus to the inner ear for hearing.
\medterm{Staph Infection} Alternate informal name; see \textit{Staphylococcal Infection}.



\medterm{Staph Infections In The Hospital} Infections caused by Staphylococcus bacteria that begin during or after healthcare treatment, or require hospital care. They may involve wounds, blood, lungs, bones, or devices and can be difficult to treat when the bacteria resist antibiotics.

\medterm{Staphylococcal Food Poisoning} An acute foodborne intoxication caused by the ingestion of preformed heat-stable enterotoxins produced by Staphylococcus aureus in improperly refrigerated or handled foods (such as dairy products, pastries, processed meats, and salads). Characterized by an abrupt onset of severe nausea, spasmodic abdominal cramping, and violent vomiting within 1 to 6 hours of consumption, usually resolving spontaneously within 24 to 48 hours.

\synonyms
Staphylococcal Enterotoxicosis; Staph Food Poisoning

\medterm{Staphylococcal Infection} An infection caused by Staphylococcus aureus or another Staphylococcus species. It may affect
the skin, bone, joints, lungs, heart valves, blood, or other tissues. Some strains resist
multiple antibiotics. Diagnosis and antimicrobial treatment depend on the site, severity,
and susceptibility of the organism.

\synonyms
Staph Infection; Staphylococcus Infection; Staphylococcal Infections

\medterm{Staphylococcal Meningitis} Staphylococcal meningitis is infection and inflammation of the meninges caused by Staphylococcus bacteria, often after surgery, trauma, bloodstream infection, or an invasive device.

\medterm{Staphylococcal Pneumonia} Pneumonia caused by Staphylococcus bacteria, sometimes after influenza or another viral infection. It may cause fever, cough, chest pain, breathlessness, lung abscesses, or infected fluid around the lung and needs medical treatment.
\medterm{Staphylococcal Scalded Skin} Alternate index label; see \textit{Staphylococcal Scalded Skin Syndrome}.

\medterm{Staphylococcal Scalded Skin Syndrome} A toxin-mediated skin disorder caused by Staphylococcus aureus, usually affecting infants or young children. It produces widespread redness, tenderness, fragile superficial blisters, and peeling that can resemble a burn and requires prompt treatment.

\medterm{Staphylococcus} A genus of gram-positive cocci that includes S. aureus and common skin colonizers and pathogens.
\medterm{Staphylococcus Aureus} A bacterium that commonly lives on the skin or in the nose but can cause infections of the skin, blood, lungs, bones, joints, or heart valves. Some strains are resistant to several antibiotics.

\medterm{Staphylococcus Infection (Staph Infection)} Alternate index label; see \textit{Staphylococcal Infection}.

\medterm{Staple (Vascular)} A small metal or polymer device used to close tissue or join selected blood-vessel structures. It provides rapid closure but can cause narrowing, bleeding, thrombosis, or tissue injury if unsuitable or misplaced.
\synonyms
Vascular staple; vessel-closure staple

\medterm{Stapling} Closing or joining tissue with metal or absorbable surgical staples. Staples may be used for skin, blood vessels, or parts of the digestive tract and must be selected and removed or absorbed appropriately.

\medterm{Starch Poisoning} Illness after swallowing a starch product, most often causing choking, blockage, or stomach upset rather than systemic poisoning. Inhaled powder can irritate the airways; breathing difficulty or persistent vomiting requires medical care.

\medterm{Stargardt Disease} An inherited eye disorder that gradually damages the macula, the central part of the retina. It usually begins in childhood or young adulthood and causes loss of detailed central vision; side vision is often preserved.

\medterm{Starling Forces} The balance of pressures that pushes water out of small blood vessels or pulls it back in. Changes can cause swelling in tissues, dehydration of tissues, or fluid buildup around organs.

\medterm{Starvation} Severe lack of food or usable energy over time, causing weight and muscle loss, weakness, changes in body chemistry, reduced immunity, and eventually organ failure. Restarting nutrition after prolonged starvation must be done carefully because dangerous shifts can occur.
\medterm{Stasis} Slowing or stopping of normal movement of blood, lymph, intestinal contents, or other body fluids.
\medterm{Stasis Dermatitis} Inflammation of the lower-leg skin caused by chronic venous hypertension and edema. It may cause itching, redness,
scaling, brown discoloration, weeping, or ulcers around the ankles. Management focuses on treating venous disease,
compression when safe, skin care, and evaluation for infection or arterial insufficiency.

\medterm{Stasis Dermatitis And Ulcers} Stasis dermatitis and ulcers result from chronic venous hypertension in the lower legs, causing swelling, brown skin change, inflammation, itching, and slow-healing wounds.

\medterm{Stasis Ulcer} A stasis ulcer is a chronic lower-leg wound caused by venous hypertension and poor return of blood, commonly near the ankle and associated with edema and skin changes.

\medterm{STAT} A medical instruction meaning immediately or with urgent priority, such as a stat test, medicine, or consultation.
\medterm{Statin} A cholesterol-lowering medicine that reduces the liver’s production of cholesterol and lowers LDL, often called bad cholesterol. Statins reduce the risk of heart attack and stroke in appropriate people; muscle symptoms and medicine interactions should be reviewed.

\medterm{Statin-Induced Myopathy} Muscle pain, weakness, or tenderness caused by a statin cholesterol-lowering medicine. Symptoms may occur without muscle damage or, rarely, with severe muscle breakdown that can injure the kidneys. A clinician may check the medicine, symptoms, muscle enzymes, and other causes before changing treatment.

\synonyms
Statin-Associated Muscle Symptoms (SAMS); Statin Myotoxicity

\medterm{Statins} A class of medicines that lowers LDL cholesterol by reducing cholesterol production in the liver. They reduce the risk of heart attack and stroke in appropriate people; muscle symptoms, liver problems, and medicine interactions should be reviewed.

\synonyms
HMG-CoA reductase inhibitors

\medterm{Statins In Diabetes} Cholesterol-lowering treatment for people with diabetes who may benefit from reduced risk of heart attack and stroke. The medicine and dose depend on age, cardiovascular risk, kidney function, other medicines, and possible side effects.

\medterm{Station} A measure of how far the baby’s presenting part, usually the head, has descended through the pelvis during labor. It is compared with the mother’s pelvic bones and helps clinicians follow progress toward birth.

\medterm{Status Asthmaticus} A severe asthma attack that does not improve enough with the first treatments. Breathing may become dangerously difficult, with exhaustion, inability to speak normally, or low oxygen, and emergency hospital treatment is required.
\medterm{Status Epilepticus} A prolonged seizure or repeated seizures without recovery, requiring emergency treatment to prevent brain injury and complications.
\medterm{Status Migrainosus} A severe migraine attack that persists for more than 72 hours and may require urgent medical treatment, especially when dehydration, vomiting, or medication overuse is present.


\medterm{Stay Away From Asthma Triggers} Avoiding asthma triggers such as smoke, allergens, infections, fumes, and cold air can reduce symptoms and attacks, alongside prescribed controller and reliever treatment.

\medterm{STDs} Alternate index label; see \textit{Sexually transmitted diseases (STDs)}.

\medterm{Steady State} The point during regular medicine use when the amount entering the body over time is about equal to the amount leaving it. Blood levels then become relatively stable; reaching this point usually takes several doses and depends on the medicine’s half-life.

\medterm{Steal Syndrome} Reduced blood flow to a hand, foot, or other area because a bypass graft or dialysis fistula diverts too much blood. It may cause pain, coolness, pale or blue skin, numbness, or weakness and requires vascular assessment.

\medterm{Steam Iron Cleaner Poisoning} Poisoning from a steam-iron cleaning product, which may contain acids or other corrosive chemicals. It can burn the mouth, digestive tract, skin, or eyes and may irritate the lungs; do not induce vomiting.

\medterm{Steatoma} A fatty or sebaceous mass; the term may refer to a lipoma or, in the ear, a cholesteatoma depending on context.
\medterm{Steatorrhea} The passage of excess fat in the stool, causing greasy, pale, foul-smelling, floating
stools—a hallmark of fat malabsorption.

\medterm{Steatorrhoea} Excess fat in stool, causing pale, bulky, greasy, foul-smelling bowel movements. It suggests poor fat digestion or absorption from pancreatic, liver, bile-duct, intestinal, or other digestive disease.
\medterm{Steele-Richardson-Olszewski Syndrome} Alternate index label; see \textit{Progressive supranuclear palsy}.

\medterm{Stein-Leventhal Syndrome} Alternate index label; see \textit{Polyendocrine metabolic ovarian syndrome}.

\medterm{Stem Cell} A cell that can make more stem cells and develop into one or more specialized cell types. Some stem cells produce blood cells; others help form or repair different tissues.

\medterm{Stem Cell Donation} The collection of healthy blood-forming stem cells from a donor for transplantation into another person. Cells may be collected from the bloodstream by apheresis after medication increases their number in the blood, or directly from bone marrow under anesthesia. Donor and recipient compatibility, especially human leukocyte antigen matching, helps improve transplant success.

\medterm{Stem Cell Transplantation} Hematopoietic stem cell transplantation replaces diseased or ablated bone marrow
with healthy donor or autologous stem cells.

\medterm{STEMI} A type of heart attack caused by sudden blockage of a heart artery, usually producing characteristic changes on an electrocardiogram. It is an emergency because quickly restoring blood flow can limit permanent heart damage.

\synonyms
ST-Elevation Myocardial Infarction

\medterm{Stenosing Tenosynovitis} Alternate index label; see \textit{Trigger finger}.

\medterm{Stenosis} Abnormal narrowing of a passage, blood vessel, heart valve, spinal canal, or other opening. Narrowing may restrict flow or compress nearby structures.
\medterm{Stenosis Lumbar (Lumbar Stenosis)} Alternate index label for Lumbar Stenosis.

\medterm{Stenosis Of The Left Anterior Descending Artery} Narrowing of a major artery supplying the front of the heart. Reduced blood flow may cause chest pain during exertion or a heart attack if the narrowing becomes severe or the artery suddenly blocks.

\medterm{Stenosis Spinal (Lumbar Stenosis)} Alternate index label for Lumbar Stenosis.

\medterm{Stenosis, Aortic Valve} Alternate index label; see \textit{Aortic valve stenosis}.

\medterm{Stenosis, Mitral Valve} Alternate index label; see \textit{Mitral valve stenosis}.

\medterm{Stenosis, Pulmonary Valve} Alternate index label; see \textit{Pulmonary valve stenosis}.

\medterm{Stenosis, Pyloric} Alternate index label; see \textit{Pyloric stenosis}.

\medterm{Stenosis, Spinal} Alternate index label; see \textit{Spinal stenosis}.

\medterm{Stenotic} Stenotic means narrowed, as in a stenotic heart valve, blood vessel, airway, intestine, or other passage.

\medterm{Stenotrophomonas Maltophilia Resistance} The ability of Stenotrophomonas maltophilia bacteria to survive several antibiotics. Treatment depends on the infected body site, illness severity, and laboratory testing that identifies which medicines are likely to work.

\medterm{Stent} A tube or scaffold placed inside a vessel or other passage to keep it open.
\medterm{Stent (Vascular)} A small mesh or tube placed inside a narrowed or blocked vessel to help keep it open.

\synonyms
Vascular stent

\medterm{Stent Balloon (Vascular)} A balloon on a catheter used to expand a vascular stent against the vessel wall and restore an open channel. Inflation can cause vessel injury, bleeding, clotting, or rarely rupture.
\synonyms
Stent-expansion balloon; vascular angioplasty balloon

\medterm{Stent Graft} A fabric-covered stent placed inside a blood vessel to reinforce the vessel wall or exclude an aneurysm from blood flow.

\medterm{Stent Placement} Insertion of a mesh tube or other support to keep a narrowed or blocked passage open. Stents may be used in arteries, bile ducts, ureters, airways, or the esophagus; bleeding, blockage, infection, movement, or repeated narrowing can occur.

\medterm{Stent Restenosis} Re-narrowing of a blood vessel within or adjacent to a previously placed stent, usually caused by tissue growth, recurrent atherosclerosis, or thrombosis.

\medterm{Stent-Graft (Vascular)} A fabric-covered metal frame placed inside a blood vessel to reinforce a weak wall or seal an aneurysm, tear, or leak. Follow-up imaging checks for movement, blockage, or leakage around the graft.
\synonyms
Vascular stent-graft; endovascular graft

\medterm{Stenting} Placement of a small mesh tube to keep a narrowed or blocked passage open. Vascular stents restore blood flow, while drug-coated types release medicine to reduce repeated narrowing.

\medterm{Stepping Reflex} A primitive reflex present from birth to 2 months, elicited by holding the infant
upright with feet touching a flat surface. The infant makes stepping-like movements as
if walking. Disappears by 2 months as voluntary motor control develops. Reappears around
9-10 months as voluntary walking emerges.

\medterm{Steps To Take Before You Get Pregnant} Preconception care includes reviewing health conditions and medicines, folic acid, vaccinations, dental and genetic considerations, substance avoidance, and planning prenatal care.

\medterm{Stereognosis} The ability to recognize a familiar object by feeling it without looking. It requires intact skin sensation, nerve pathways, and areas of the brain that interpret touch; loss may occur after nerve or brain injury.
\medterm{Stereotactic} Stereotactic describes a technique using three-dimensional coordinates to target a precise location in the brain or another body structure.

\medterm{Stereotactic Biopsy} Image-guided biopsy using three-dimensional coordinates to target a lesion that may be difficult to localize by direct palpation. It is used in the breast and, with specialized systems, the brain, lung, and other organs.

\medterm{Stereotactic Body Radiation Therapy} A focused form of external radiation that delivers a high dose to a small body tumor in a few treatments. It may treat selected lung, liver, kidney, pancreatic, or spinal tumors while limiting exposure to nearby tissue.

\medterm{SRT} Stereotactic radiotherapy, a precise form of external radiation treatment that uses three-dimensional imaging and careful positioning to target a small tumor or other abnormal area. It is often given in two to five sessions to limit exposure to nearby tissue.

\medterm{Stereotactic Brain Biopsy} A precision, minimally invasive neurosurgical procedure used to safely obtain small tissue samples from deep, delicate, or high-risk areas of the brain for microscopic examination. Using high-resolution pre-operative MRI or CT scans and a specialized three-dimensional stereotactic head frame or frameless optical infrared navigation system, a neurosurgeon maps exact mathematical coordinates to the suspected tumor or infection. Through a tiny circular burr hole (approximately 3 to 5 mm) in the skull, the surgeon gently passes a specialized hollow biopsy needle straight to the target lesion without disrupting surrounding healthy brain tissue. The extracted tissue specimens are examined by a neuropathologist to determine the exact tumor type, grade, genetic markers, or underlying infection, allowing doctors to design targeted chemotherapy, radiation, or medical therapy.
\synonyms
Stereotactic Needle Biopsy; Brain Biopsy; Frameless Stereotactic Biopsy

\medterm{Stereotactic Radiosurgery} Precise radiation treatment delivering one or a few high doses to a small target, often in the brain, without open surgery.

\medterm{Stereotactic Radiosurgery - CyberKnife} CyberKnife is a robotic system that delivers many precisely aimed radiation beams to a tumor or other target without an open operation.


\medterm{Stereotactic Radiosurgery - Gamma Knife} Gamma Knife is a focused radiation system that delivers multiple precisely aimed beams, mainly to selected brain tumors and vascular or other brain lesions.

\medterm{Stereotaxis} A technique using a three-dimensional coordinate system to target deep structures, especially in brain procedures.
\medterm{Stereotypic Movement Disorder} Repeated, apparently purposeless movements such as hand flapping, rocking, or head banging that interfere with activities or cause injury. It is more common in childhood and developmental conditions and requires assessment for other causes.

\medterm{Sterile Field} A designated area maintained free of viable microorganisms during an invasive procedure.
Only sterile equipment and materials may contact the field or the sterile portions of
instruments.

\medterm{Sterile Site} A body area or sample that normally contains no living microorganisms, such as blood, spinal fluid, joint fluid, or tissue from a closed space. Finding a germ there is usually important, although accidental contamination must be considered.

\medterm{Sterile Technique} Methods used to prevent introduction of microorganisms into a sterile field, wound, tissue, device, or body site through hand hygiene, barriers, aseptic handling, and appropriate preparation.

\medterm{Sterilisation} Destruction or removal of all forms of microbial life, or permanent contraception depending on context.
\medterm{Sterility} Inability to reproduce, or complete absence of living microorganisms from an object or preparation, depending on context. In healthcare, sterile equipment, implants, and medicines must remain free from contamination to prevent infection.
\medterm{Sterilization} A process that destroys or removes all forms of microbial life, including spores, from equipment, medicines, or other materials. In another context, sterilization means a permanent method intended to prevent pregnancy.

\medterm{Sternal Exploration Or Closure} Surgical reopening or repair of the breastbone after median sternotomy, performed to treat infection, bleeding, instability, dehiscence, or other postoperative chest complications.

\medterm{Sternoclavicular Joint} A saddle-type synovial joint between the clavicle and manubrium of the sternum. It is
the only bony connection between the upper limb and axial skeleton and permits movement
of the shoulder girdle.

\medterm{Sternocleidomastoid Muscle} A prominent muscle on each side of the neck that turns and bends the head and helps stabilize the neck. It can also assist forceful breathing; strain or spasm may cause neck pain and limited movement.

\medterm{Sternotomy} Surgical division of the sternum, usually along the midline, to provide access to the heart or mediastinal structures.

\medterm{Sternum} The breastbone forming the front of the central chest wall.
\medterm{Sternutation} Sneezing, a reflex forceful expulsion of air through the nose and mouth triggered by irritation of nasal or upper-airway sensory receptors.

\medterm{Steroid Myopathy} Muscle weakness caused by prolonged exposure to corticosteroid medicines or excessive cortisol. It usually develops gradually and affects the hips, thighs, shoulders, and upper arms, making it hard to rise from a chair or climb stairs. Weakness may improve when the cause is treated, but steroid doses must be changed only with medical guidance.

\synonyms
Corticosteroid-Induced Myopathy; Glucocorticoid Myopathy

\medterm{Steroids} A broad family of compounds with a characteristic four-ring structure, including body hormones, corticosteroids, and sex hormones. Their effects depend on the type; prolonged corticosteroid use can suppress immunity and bone health.
\synonyms
Corticosteroids

\medterm{Stertor} A low-pitched snoring, gurgling, or rattling sound caused by partly blocked air passages in the nose or throat. It may occur during sleep or illness; noisy breathing with severe sleepiness, blue lips, or breathing difficulty is an emergency.
\medterm{Stethoscope} An instrument used for auscultation, the examination of sounds produced within the body. It usually has a chest piece with a diaphragm and sometimes a bell, flexible tubing, and earpieces; electronic models convert sounds into amplified or recorded signals. It is used to examine heart, lung, bowel, and vascular sounds.
\medterm{Stevens-Johnson Syndrome} A rare, life-threatening reaction usually caused by a medicine or infection. It causes painful widespread skin and mucous-membrane damage, blistering, and peeling and requires emergency hospital treatment.

\synonyms
SJS

\medterm{STI} Alternate index label for Sexually Transmitted Infections.

\medterm{Stickler Syndrome} An inherited connective-tissue disorder that may cause severe nearsightedness, retinal detachment, hearing loss, early arthritis, and differences in the jaw or palate. Regular eye and hearing checks are important.

\medterm{Stiff Baby Syndrome} A nonspecific description of unusual stiffness or increased muscle tone in an infant, not a single diagnosis. Causes include brain, nerve, muscle, genetic, or metabolic disorders; persistent stiffness, poor feeding, seizures, or breathing problems need urgent assessment.

\medterm{Stiff-Person Syndrome} A rare immune-related nervous-system disorder causing persistent muscle stiffness and sudden painful spasms. Noise, touch, movement, or emotional stress may trigger spasms, leading to falls and difficulty walking. Blood tests may find related antibodies. Treatment can reduce stiffness and immune activity, but symptoms often require long-term care.

\synonyms
SPS; Stiff-Man Syndrome; Moersch-Woltman Syndrome

\medterm{Stiffness} Difficulty or discomfort when moving a body part. It may result from joint disease, muscle problems, nerve disorders, inflammation, injury, prolonged rest, or aging and can limit everyday activities.
\medterm{Stigma} A mark, characteristic, or social discrediting label; in medicine it may refer to a sign or to harmful social attitudes.
\medterm{Stilboestrol} An older name for diethylstilbestrol, a synthetic estrogen no longer used routinely because of serious risks.
\medterm{Stilet} A thin, pointed instrument placed inside a cannula or tube to guide insertion, clear a blockage, or provide firmness. It must be used carefully because it can injure tissue or blood vessels.
\medterm{Still Disease} An inflammatory disease that can affect the whole body.
In children it is called systemic juvenile idiopathic arthritis; in adults it is called
adult-onset Still disease. It can cause fever, joint pain or swelling, and a rash.

\synonyms
Still's Disease; Arthritis Still (Still's Disease)

\medterm{Stillbirth} The death of a fetus before or during delivery after the gestational-age or birth-weight threshold used by a country or health system. In the United States, it usually refers to loss at or after 20 weeks of pregnancy; the World Health Organization commonly uses 28 weeks for international comparisons. Causes may include problems with the placenta, fetal growth, infection, genetic or developmental conditions, or maternal illness, but some stillbirths remain unexplained.
\medterm{Still's Murmur} A common harmless heart murmur in children, usually heard between ages 2 and 7. It has a soft, musical sound, becomes louder with fever or exercise, and usually disappears as the child grows; the heart is otherwise normal.

\synonyms
Innocent Vibratory Murmur; Normal Pediatric Murmur; Twanging-String Murmur

\medterm{Stimulant} A substance that increases activity in the brain or other body systems. Stimulants may improve alertness and attention, but excessive or inappropriate use can cause dependence, anxiety, fast heartbeat, high blood pressure, or seizures.

\medterm{Stimulant Use Disorder} A pattern of using stimulant drugs that causes significant problems or loss of control despite harm. Stimulants include cocaine and amphetamine-type drugs; stopping may cause tiredness, low mood, increased appetite, or excessive sleep.

\medterm{Stimulants} Substances that increase activity in the brain or the body’s fight-or-flight system. They may improve alertness and attention but can also cause insomnia, anxiety, dependence, fast heartbeat, high blood pressure, or seizures.

\medterm{Stimulus} A stimulus is a physical, chemical, sensory, or social change that evokes a response from a cell, tissue, organism, or person.

\medterm{Stingray} A sea-animal injury caused by a stingray spine, usually producing severe pain, swelling, and a puncture wound. Medical care may be needed to clean the wound, remove retained spines, and treat whole-body reactions.

\medterm{Stings} Injuries caused when an insect, animal, or plant injects venom or irritating material. They may cause pain, swelling, itching, allergic reactions, tissue damage, or serious poisoning depending on the source.
\medterm{Stings, Bee} Alternate index label; see \textit{Bee sting}.

\medterm{Stings, Jellyfish} Alternate index label; see \textit{Jellyfish stings}.

\medterm{Stings, Scorpion} Alternate index label; see \textit{Scorpion sting}.

\medterm{Stippling (Tattooing)} Patterned dermal abrasions and punctate hemorrhages surrounding a firearm entrance
wound, caused by unburned and partially burned gunpowder grains striking the skin at
close range (15--75 cm).

\synonyms
Tattooing

\medterm{Stitch} A suture used to close tissue, or a sharp transient pain, often in the side during exercise.
\medterm{Stochastic Effects} Possible effects of radiation, especially cancer or inherited genetic changes, whose chance increases with dose. The severity cannot be predicted from the dose, and even low exposure may carry some risk.
\medterm{Stoddard Solvent Poisoning} Poisoning from Stoddard solvent, a petroleum-based cleaning solvent. Inhalation or swallowing can cause headache, dizziness, drowsiness, skin irritation, or aspiration-related lung injury; do not induce vomiting and seek advice.

\medterm{Stokes-Adams Syndrome} Sudden fainting caused by a temporary dangerous slowing or stopping of the heartbeat, often from heart block or another abnormal rhythm. It requires urgent heart assessment because injury or cardiac arrest can occur.
\medterm{Stoma} A natural opening or a surgically created opening connecting an organ to the skin or another organ.
\medterm{Stoma Complications} Problems involving an opening made from the intestine or urinary tract to the skin. They include skin irritation, bleeding, retraction, prolapse, narrowing, blockage, high output, infection, and a parastomal hernia; sudden severe pain, color change, or no output needs urgent care.

\medterm{Stomach} A muscular digestive organ between the oesophagus and small intestine that stores and begins digestion of food.
\medterm{Stomach Cancer} Cancer that begins in the stomach, most often in gland-forming cells lining the stomach. Risk factors include Helicobacter pylori infection, smoking, chronic inflammation, certain foods, and inherited conditions; symptoms and treatment depend on type and stage.

\synonyms
Cancer, Gastric
Cancer, Stomach

\medterm{Stomach Cramps} Intermittent gripping or spasmodic upper-abdominal pain caused by gas, infection, indigestion, ulcer disease, gallbladder or pancreatic disease, obstruction, or other causes.

\medterm{Stomach Fat (Belly Fat, Abdominal Fat)} Adipose tissue stored beneath the skin or around abdominal organs; excess visceral fat is metabolically active and associated with insulin resistance, fatty liver, and cardiovascular risk.

\medterm{Stomach Flu} A common informal term for viral gastroenteritis, an infection of the stomach and intestines caused by viruses rather than influenza viruses. It may cause watery diarrhea, vomiting, abdominal cramps, fever, and dehydration; some infections cause few or no symptoms.

\medterm{Stomach Paralysis} Stomach paralysis, or gastroparesis, is delayed emptying of stomach contents without obstruction, causing early fullness, nausea, vomiting, and bloating.

\medterm{Stomach Tube} A tube passed through the mouth or nose into the stomach for drainage, feeding, sampling, or treatment.
\medterm{Stomach Ulcer} An open sore in the lining of the stomach, commonly associated with Helicobacter pylori infection or medicines such as NSAIDs.

\medterm{Stomach Ulcer (Peptic Ulcer)} Alternate index label for Peptic Ulcer.

\medterm{Stomach Washout} A procedure that rinses and removes stomach contents through a tube. It is rarely used for selected serious poisonings because it can cause choking, lung injury, bleeding, or other complications.
\medterm{Stomach, Diseases Of} Disorders affecting the stomach, including inflammation, ulcers, movement problems, blockage, bleeding, and tumors. Symptoms may include pain, nausea, vomiting, early fullness, weight loss, indigestion, or blood in vomit or stool.
\medterm{Stomatitis} Inflammation, soreness, redness, or ulcers of the lining of the mouth. Causes include infection, irritation, nutritional deficiency, medicines, and cancer treatment; severe pain, dehydration, fever, or inability to eat or drink needs medical care.
\medterm{Stomatodynia} Alternate index label; see \textit{Burning mouth syndrome}.

\medterm{Stone} A hard lump formed inside an organ or tube, commonly from mineral salts or other body substances. Stones may block flow, cause pain or infection, and occur in kidneys, gallbladder, or ducts.
\medterm{Stonefish Sting} A puncture injury from the venomous spines of a stonefish, causing severe pain, swelling, bleeding, and sometimes weakness, breathing problems, or shock. Seek emergency care; hot-water immersion may reduce pain while help is obtained.

\medterm{Stool} Fecal material expelled from the bowel, composed of water, undigested residue, bacteria, cells, mucus, and bile pigments; color, form, frequency, and blood can provide clinical clues.

\medterm{Stool C Difficile Toxin} A stool test for harmful substances made by Clostridioides difficile bacteria. A positive result supports active infection when a person has compatible diarrhea, but bacteria can be present without causing illness.

\medterm{Stool Examination} Laboratory analysis of a stool sample for blood, infection, parasites, inflammation, malabsorption, or other digestive abnormalities.

\synonyms
Stool test

\medterm{Stool Gram Stain} Microscopic examination of a stool specimen after Gram staining; it may provide limited information about bacterial flora or inflammatory cells but is not a general test for infectious diarrhea.

\medterm{Stool Guaiac Test} A chemical fecal occult blood test that detects the peroxidase activity of heme in stool. Dietary factors, medicines, and specimen collection can affect results. It is a screening or detection test rather than a diagnosis; a positive result requires clinical follow-up and usually evaluation for a gastrointestinal source.

\medterm{Stool Holding} Alternate index label; see \textit{Encopresis}.

\medterm{Stool Occult Blood Test} Alternate index label; see \textit{Fecal occult blood test}.

\medterm{Stool Ova And Parasites Exam} Microscopic examination of stool for parasite eggs, cysts, larvae, or other forms, sometimes requiring multiple specimens to improve detection.

\medterm{Stools} Feces, the material expelled from the digestive tract. Changes in color, consistency, frequency, or blood can provide clues about digestive health but are not diagnostic alone.
\medterm{Stools - Foul Smelling} Strongly odorous feces that may result from diet, infection, bleeding, malabsorption, altered gut bacteria, or rapid transit; persistent change with weight loss or diarrhea warrants review.

\medterm{Stools - Pale Or Clay-Colored} Light-colored stool caused by reduced bile pigment reaching the intestine, potentially indicating bile-duct obstruction, gallbladder disease, liver disease, or certain transient conditions.

\medterm{Stools – Floating} Feces that remain on the water surface, often from trapped gas or diet but sometimes from excess fat in malabsorption; persistent greasy, pale, foul-smelling stool needs evaluation.

\medterm{Stop Smoking Support Programs} Counseling, quit plans, telephone or digital coaching, and medicines such as nicotine replacement, varenicline, or bupropion used to help people stop tobacco use.

\medterm{STOPP/START} A checklist used to review medicines in older adults. STOPP identifies medicines that may be inappropriate or harmful, while START identifies useful treatments that may have been missed; clinical judgment is still required.

\medterm{Storage Diseases} A group of inherited metabolic disorders characterized by accumulation of substrates
within cells due to specific enzyme deficiencies. Lysosomal storage diseases result from
defects in lysosomal hydrolytic enzymes, leading to accumulation of undigested
macromolecules within lysosomes.


\medterm{Stork Bite} A stork bite is a common harmless pink birthmark, usually a nevus simplex on the forehead, eyelids, nose, or neck, that often fades during childhood.

\medterm{Strabismus} A condition in which the eyes are not aligned and do not point in the same direction. One eye may turn inward, outward, upward, or downward. It may be constant or come and go. Strabismus can cause double vision in adults and may lead to amblyopia, or reduced vision in one eye, when it begins during childhood. Common types include esotropia, exotropia, hypertropia, and hypotropia.
\medterm{Strain} An injury from overstretching or tearing a muscle or tendon, or a force applied to tissue.
\medterm{Strain Imaging} Ultrasound or magnetic-resonance imaging that measures how much tissue deforms during movement, especially how heart muscle shortens or thickens with each heartbeat. It can reveal early weakness even when the total amount of blood pumped appears normal.

\medterm{Strain Neck (Whiplash)} Alternate index label for Whiplash.

\medterm{Strains} A strain is overstretching or tearing of a muscle or tendon, causing pain, tenderness, swelling, weakness, or reduced function after force or overuse.

\medterm{Strangulated Hernia} A hernia in which the blood supply to the trapped tissue is cut off, causing ischaemia and possible tissue death; it is a surgical emergency.

\medterm{Strangulation} Compression of a body structure, especially the neck or a hernia, that compromises blood flow or airway.
\medterm{Strapping} Application of adhesive tape or bandage to support, protect, or restrict movement of a body part.
\medterm{Strategies For Getting Through Labor} Practical approaches for coping with labor, including breathing, movement, position changes, support people, relaxation, analgesia, and communication with the maternity team.

\medterm{Stratum Corneum} The outermost skin layer, made mainly of flattened dead cells filled with keratin. It forms a protective barrier that limits water loss and helps block chemicals, germs, and physical injury.

\medterm{Strawberry Hemangioma} Alternate index label; see \textit{Hemangioma}.

\medterm{Strength} The ability of a muscle or group of muscles to produce force. Healthcare professionals may test strength by comparing movements against resistance; reduced strength can result from pain, poor effort, muscle disease, nerve problems, or brain injury.

\medterm{Strength Testing} Assessment of how much force a muscle or muscle group can produce, using examination or equipment. It helps identify weakness and monitor recovery, progression, or response to treatment.

\medterm{Strength Training} Exercise against resistance to improve muscle force, endurance, power, bone health, and function. The program should be progressed safely according to ability and medical needs.

\synonyms
Resistance training; weight training

\medterm{Strep Throat} A bacterial infection of the throat and tonsils caused by group A Streptococcus, also called Streptococcus pyogenes. It commonly causes a sudden sore throat, pain when swallowing, fever, and swollen or red tonsils, sometimes with white patches.
\medterm{Streptococcal Cellulitis With Lymphangitis} A streptococcal skin infection with inflammation spreading along nearby lymphatic vessels. Redness, warmth, pain, swelling, fever, and red streaks may occur and require prompt antibiotic treatment.
\medterm{Streptococcal Infection} An infection caused by bacteria of the Streptococcus group. Depending on the species and site, it may cause strep throat, scarlet fever, impetigo, cellulitis, pneumonia, heart-valve infection, bloodstream infection, or severe tissue infection.

\medterm{Streptococcal Infections} Infections caused by Streptococcus bacteria, which may affect the throat, skin, lungs, heart valves, blood, or other tissues. Diagnosis and antibiotic choice depend on the site and species; rapidly worsening illness, severe pain, or breathing difficulty is urgent.

\medterm{Streptococcal Pharyngitis} Alternate index label for Strep Throat.

\medterm{Streptococcal Pharyngitis (Strep Throat)} Alternate index label for Strep Throat.

\medterm{Streptococcal Screen} A streptococcal screen looks for group A Streptococcus, usually from a throat swab, to help diagnose strep throat.

\medterm{Streptococcal Toxic Shock Syndrome} A fulminant, life-threatening invasive infection caused by Streptococcus pyogenes (group A Streptococcus) producing streptococcal pyrogenic exotoxins that act as superantigens. It is characterized by the sudden onset of hypotension, shock, and multiorgan failure, frequently in association with severe invasive soft-tissue infections such as necrotizing fasciitis; it carries an exceptionally high mortality rate.

\synonyms
STSS; Streptococcal TSS; Toxic Shock-Like Syndrome

\medterm{Streptococcus} A group of round bacteria that can cause throat infections, pneumonia, skin infections, heart-valve infection, or serious bloodstream illness. Different species spread and respond to antibiotics differently, so testing may guide treatment.
\medterm{Streptococcus Agalactiae} Group B Streptococcus, a bacterium that can live in the bowel or genital tract without causing symptoms. It can pass to a baby during birth and cause serious newborn infection, so pregnancy screening and preventive antibiotics are used when indicated.

\medterm{Streptococcus Pneumoniae (Pneumococcus)} A gram-positive bacterium that can cause pneumonia, meningitis, sinusitis, otitis media, and bloodstream infection; vaccines reduce disease risk.

\medterm{Streptococcus Pyogenes} Group A Streptococcus, a bacterium that causes strep throat, scarlet fever, impetigo, cellulitis, and sometimes severe invasive infection. Untreated infection can occasionally lead to rheumatic fever or kidney inflammation.

\medterm{Streptokinase} A thrombolytic protein that activates plasminogen and dissolves clots; use is limited by bleeding risk and alternatives.
\medterm{Streptomycin} An aminoglycoside antibiotic used in selected tuberculosis and other infections.
\medterm{Stress} A physical and emotional response to perceived demands or threats. It can change heart rate, breathing, muscle tension, attention, sleep, and digestion; persistent or overwhelming stress may affect health.


\medterm{Stress Cardiomyopathy} A temporary weakening of the heart muscle that often follows severe emotional or physical stress. It can cause chest pain, breathlessness, or fainting and may resemble a heart attack, even when no heart artery is blocked.

\medterm{Stress Echocardiography} An ultrasound scan of the heart performed while the heart works harder through exercise or medicine. It looks for changes in heart-muscle movement that may suggest reduced blood flow.
\medterm{Stress Fracture} A small crack in a bone caused by repeated activity or loading. Pain usually develops gradually in one spot and worsens with activity; continuing to stress the bone can make the crack larger or complete.

\medterm{Stress Fractures} Small cracks in bone caused by repeated loading that exceeds the bone’s ability to repair itself. They commonly affect the feet or legs and cause localized pain that worsens with activity.

\medterm{Stress Incontinence} Involuntary leakage of urine during activities that increase intra-abdominal pressure, such as coughing, sneezing, lifting, or exercise.

\medterm{Stress Injuries} Tissue injuries caused by repeated loading or overuse without sufficient recovery. Examples include stress fractures, tendinopathy, and muscle strain.

\medterm{Stress Response} Body changes that prepare a person to respond to a perceived threat or challenge. Hormones and nerve signals alter heart rate, breathing, blood flow, attention, digestion, immunity, and energy use.

\synonyms
Fight-or-flight response

\medterm{Stress Shielding} Loss or weakening of bone next to an implant because the implant carries more force than the bone normally would. Over time, this can reduce bone density and affect how securely the implant stays in place.

\medterm{Stress Tests} Diagnostic tests that assess organ function under controlled physical, pharmacologic, or physiologic stress, commonly evaluating the heart for inducible ischemia or abnormal rhythm.

\medterm{Stress Ulcer Prophylaxis} Measures used to reduce the risk of stomach or duodenal bleeding in selected critically ill people. The need for treatment depends on bleeding risk, other medicines, and the person's overall condition.

\medterm{Stress Urinary Incontinence} Leakage of urine when coughing, sneezing, laughing, lifting, exercising, or otherwise increasing pressure inside the abdomen. It commonly results from weakness of the pelvic-floor muscles or the bladder outlet and can often be treated.
\medterm{Stressors} Events, conditions, or circumstances that trigger a physical or emotional stress response. Examples include illness, pain, sleep loss, trauma, work demands, and relationship difficulties.

\medterm{Stretch Marks} Linear bands of atrophic skin, also called striae, caused by stretching and hormonal or corticosteroid effects; they commonly occur during pregnancy, growth, weight change, or steroid exposure.

\medterm{Stretch Reflex} An automatic tightening of a muscle when it is suddenly stretched, such as the lower-leg movement after tapping the knee tendon. It helps maintain posture and provides information about nerve and spinal-cord function.

\medterm{Stria} A line, streak, or band in tissue, such as a stretch mark. Striae may result from skin stretching, growth, hormone changes, medicines, or underlying disease and can fade over time.

\synonyms
Striae

\medterm{Striae} Linear bands or streaks in the skin, commonly called stretch marks. They develop when deeper skin layers stretch or change rapidly during growth, pregnancy, weight change, muscle gain, or hormone treatment.
\medterm{Striae Atrophicae} Linear bands of thin, wrinkled skin caused by stretching and disruption of dermal connective tissue; commonly called stretch marks.

\synonyms
stretch marks.

\medterm{Striae Gravidarum} Stretch marks that develop during pregnancy, commonly on the abdomen, breasts, thighs, or
buttocks. They may begin as red or purple lines and gradually fade to lighter bands.

\medterm{Stricture (Stenosis)} An abnormal narrowing of a passage in the body, such as the esophagus, intestine, bile duct, or urethra. Scarring, inflammation, injury, surgery, or cancer may cause blockage and symptoms related to the affected organ.

\medterm{Strictures Esophagus (Schatzki Ring)} Alternate index label for Schatzki Ring.

\medterm{Stridor} A harsh, high-pitched sound, usually heard during breathing in, caused by narrowing of the upper airway. New or severe stridor is an emergency because airflow may become blocked.
\medterm{String Test For Intestinal Parasites} A string test collects a sample from the upper digestive tract on a swallowed capsule and string to look for selected intestinal parasites.

\medterm{Stripping Of Membranes (Membrane Sweep)} A procedure performed during a cervical exam where the examiner inserts a finger through
the internal os and sweeps the membranes circumferentially, separating them from the
lower uterine segment. Releases endogenous prostaglandins, promoting cervical ripening
and potentially initiating labor.

\synonyms
Membrane Sweep

\medterm{Stroke} Brain injury caused by interrupted blood flow to part of the brain (ischemic stroke) or bleeding in or around the brain (hemorrhagic stroke). It may cause sudden facial drooping, arm or leg weakness or numbness, speech trouble, vision loss, severe headache, dizziness, or confusion. The effects vary according to the brain area involved and the extent of injury.
\medterm{Stroke Rehabilitation} Rehabilitation after a stroke helps improve movement, speech, swallowing, thinking, mood, and daily activities while reducing complications and preventing another stroke.
\medterm{Stroke Index} A measure of the amount of blood pumped by the heart with each beat, adjusted for body-surface area. It helps assess heart function and is interpreted with other examination, imaging, and hemodynamic findings.

\medterm{Stroke Risk Factors} Stroke risk factors include high blood pressure, smoking, diabetes, high cholesterol, atrial fibrillation, prior stroke, and certain blood or vascular disorders.


\medterm{Stroke Volume} The volume of blood ejected by the left ventricle with each heartbeat, normally 60-100
mL. It is determined by three factors: preload, afterload, and myocardial contractility.
Stroke volume equals end-diastolic volume minus end-systolic volume. Reduced stroke
volume occurs in heart failure, hypovolemia, and cardiogenic shock.

\medterm{Stroma} The supporting tissue of an organ that surrounds and holds its working cells. It may contain connective tissue, blood vessels, nerves, and immune cells; the exact structure and role vary between organs.
\medterm{Strongyloides Stercoralis} A soil-transmitted roundworm that can enter through skin and live in the human intestine. Infection may cause rash, abdominal symptoms, or no symptoms, but severe hyperinfection can occur when immunity is weakened.

\medterm{Strongyloidiasis} An intestinal worm infection that can spread throughout the body and become life-threatening when immunity is weakened.
\medterm{Structured Reporting} Reporting that uses a standard format, headings, and required fields to organize findings. It can improve completeness, consistency, and communication in imaging, pathology, and other clinical reports, but does not replace professional judgment.

\medterm{Struma Ovarii} A specialized mature ovarian teratoma composed predominantly of thyroid tissue; it may be hormonally active or develop malignant thyroid-type transformation.

\medterm{Struvite Stones} Kidney or urinary-tract stones formed during infection with certain bacteria. They can grow into large branching stones that fill the kidney and may damage it or block urine flow.

\medterm{Strychnine} A poison causing severe painful muscle spasms and convulsions while awareness may remain; emergency treatment is required.
\medterm{Student's Elbow} Olecranon bursitis, an inflammation of the bursa over the point of the elbow, often related to repeated pressure or trauma.


\medterm{Stuffy Nose} Blocked or difficult nasal breathing caused by swollen tissue, mucus, allergy, infection, irritants, or structural narrowing.

\medterm{Stuffy Or Runny Nose - Children} A stuffy or runny nose in a child is usually caused by a viral infection or allergy; breathing difficulty, dehydration, prolonged fever, or symptoms in an infant need prompt assessment.

\medterm{Stuffy Or Runny Nose – Adult} A stuffy or runny nose in an adult commonly results from viral infection, allergy, irritants, or sinus disease; persistent, severe, bloody, or one-sided symptoms need review.

\medterm{Stump Pain} Alternate index label; see \textit{Residual limb pain}.

\medterm{Stupor} A state of reduced consciousness in which the patient is unresponsive to most environmental stimuli but can be temporarily aroused by vigorous, repeated stimulation (painful stimuli, loud shouting).
\medterm{Sturge-Weber Syndrome} A condition present from birth involving abnormal blood vessels in the skin and sometimes on the surface of the brain. It may cause a port-wine birthmark, seizures, weakness, headaches, developmental problems, or glaucoma, depending on which areas are affected.

\medterm{Stuttering} A speech-flow problem involving repeated sounds or words, stretched sounds, pauses, or blocks when speaking. It is not caused by low intelligence; support may include speech therapy, communication strategies, and help with anxiety or teasing.
\medterm{Sty} Alternate index label for Stye.

\medterm{Stye} A tender, localized infection or inflammation of an eyelid gland, usually producing a painful red lump near the eyelid margin. It is also called a hordeolum.


\medterm{Styptics} Substances that help stop minor surface bleeding by tightening tissue, narrowing small blood vessels, or supporting clot formation. They are used for limited bleeding and do not replace emergency care for severe blood loss.
\medterm{Subacute} Developing or lasting longer than an acute condition but not as long as a chronic one. The exact time range varies with the disease.
\medterm{Subacute Bacterial Endocarditis} A slowly developing infection of the inner lining or valves of the heart, usually caused by bacteria entering the bloodstream. It may cause fever, tiredness, weight loss, a new heart murmur, or embolic complications and requires prolonged antibiotic treatment.

\medterm{Subacute Combined Degeneration} Damage to parts of the spinal cord, usually from vitamin B12 deficiency. It can cause numbness, weakness, poor balance, and difficulty walking.

\medterm{Subacute Combined Degeneration Of The Cord} Damage to the spinal cord, usually caused by vitamin B12 deficiency, affecting nerve pathways for vibration, position sense, and movement. It may cause tingling, weakness, poor balance, or difficulty walking and can become permanent without treatment.
\medterm{Subacute Rehabilitation} Rehabilitation provided after the early phase of illness or injury for people who still need nursing care and therapy but no longer need intensive hospital treatment. It may occur in a specialized facility or outpatient setting.

\medterm{Subacute Sclerosing Panencephalitis} A rare, progressive brain disease caused by measles virus remaining in the brain years after the original infection. It can cause behavior or learning changes, jerking movements, seizures, loss of coordination, and severe disability; prevention depends on measles vaccination.
\medterm{Subacute Thyroiditis} Temporary inflammation of the thyroid, often after a viral illness, causing painful neck swelling, fever, and tiredness. Thyroid hormone levels may first be high and later low before returning to normal; treatment relieves pain and other symptoms.

\medterm{Subaortic Stenosis} Narrowing just below the aortic valve that obstructs blood leaving the heart. It may cause a heart murmur, chest pain, fainting, breathlessness, or damage to the aortic valve and sometimes requires surgery.

\medterm{Subarachnoid} Situated beneath the arachnoid, one of the thin coverings around the brain and spinal cord. The subarachnoid area contains protective fluid and blood vessels; bleeding there can cause a sudden severe headache and is an emergency.

\medterm{Subarachnoid Haemorrhage} Bleeding into the space around the brain containing cerebrospinal fluid, often causing sudden severe headache.
\medterm{Subarachnoid Hemorrhage} Bleeding into the space around the brain, often from a ruptured blood-vessel bulge. It typically causes a sudden, extremely severe headache and may cause vomiting, neck stiffness, confusion, weakness, or loss of consciousness.

\medterm{Subarachnoid Space} The space between two thin coverings of the brain and spinal cord, called the arachnoid and pia mater. It contains cerebrospinal fluid and blood vessels and helps cushion the nervous system.

\medterm{Subareolar Abscess} A pus-filled infection beneath the areola, often associated with periductal inflammation and smoking, causing a painful lump, redness, drainage, or recurrent fistula.

\medterm{Subchorionic Hemorrhage} Bleeding between the chorion and the uterine wall caused by partial separation of the gestational sac membranes. It may produce vaginal bleeding in early or mid-pregnancy and is assessed according to bleeding, gestational age, and ultrasound findings.

\medterm{Subclavian} Relating to structures beneath the clavicle, especially the subclavian arteries and veins that supply or drain the upper limbs.
\medterm{Subclavian Artery} A major artery beneath the collarbone that supplies the arm and contributes blood flow to the neck, chest, and brain. Narrowing or blockage can cause arm pain, weakness, or symptoms from reduced blood flow to the brain.

\medterm{Subclavian Steal Syndrome} A condition in which narrowing of an artery supplying the arm redirects blood away from the back of the brain, especially during arm exercise. It may cause arm pain or weakness, dizziness, fainting, blurred vision, or imbalance. Treatment controls vascular risk factors and may restore blood flow when symptoms are significant.


\medterm{Subclavian Vein} The continuation of the axillary vein beyond the lateral border of the first rib. It
joins the internal jugular vein behind the sternoclavicular joint to form the
brachiocephalic vein.

\medterm{Subclavian Vein Thrombosis} A blood clot in the large vein beneath the collarbone. It may follow a catheter, repeated strenuous arm activity, or compression near the shoulder and can cause arm swelling, pain, bluish skin, or a dangerous clot in the lungs.

\medterm{Subclinical} Present without clear symptoms or signs that a person or clinician can notice. A subclinical problem may be found through a blood test, scan, screening examination, or other test and may or may not later cause symptoms.
\medterm{Subclinical Disease} Disease that produces no obvious symptoms or signs but can be detected by screening, examination, or laboratory testing. It may remain mild or later become clinically apparent, depending on the condition.

\medterm{Subclinical Hypothyroidism} A blood-test pattern showing that the thyroid is beginning to underperform: the pituitary sends extra stimulation, but thyroid-hormone levels remain normal. Many people have no symptoms; others have mild tiredness, cold intolerance, or constipation.

\medterm{Subclinical Infection} An infection that produces no noticeable symptoms or signs but may still be detectable by testing and may sometimes be transmissible.

\medterm{Subconjunctival Hemorrhage} Bleeding beneath the conjunctiva that appears as a sharply demarcated red patch on the white of the eye; it is usually harmless but recurrent cases or trauma require assessment.

\medterm{Subconjunctival Hemorrhage (Broken Blood Vessel In Eye)} A small blood vessel broken beneath the clear surface of the eye, causing a sharply demarcated red patch that usually resolves without treatment.

\synonyms
Broken Blood Vessel In Eye

\medterm{Subconscious} Mental processes outside immediate awareness; it is a psychological term rather than a specific diagnosis.
\medterm{Subcorneal Pustular Dermatosis} A rare recurring skin disorder causing shallow, noninfectious pus-filled blisters just beneath the outer skin layer, usually on the trunk or skin folds. It may resemble psoriasis and can be associated with blood or immune disorders.

\medterm{Subcu} Subcu is an abbreviation for subcutaneous, meaning beneath the skin; clear clinical instructions should spell out the route.

\medterm{Subculture} The transfer of microorganisms from one culture to fresh culture medium, used to
maintain pure cultures, isolate colonies, increase viability, prepare standardized
inocula, and perform identification/susceptibility testing. On solid media, subculture
aims to obtain isolated single colonies from a mixed or stock culture.

\medterm{Subcutaneous} Subcutaneous means beneath the skin; injections by this route deliver medicine into the fatty tissue between skin and muscle.

\medterm{Subcutaneous Injections} Administration of a medicine into the fatty tissue beneath the skin, commonly used for insulin, anticoagulants, and selected biologic medicines.

\medterm{Subcutaneous Emphysema} Air trapped beneath the skin, often after chest injury, surgery, a punctured lung, or a breathing procedure. It may cause swelling and a crackling feeling; treatment addresses the air leak, and breathing difficulty or rapid spread requires urgent care.

\medterm{Subcutaneous Implantable Cardioverter-Defibrillator} A device placed under the skin that detects and treats dangerous fast heart rhythms with an electrical shock. Its wires run outside the heart and blood vessels; infection, inappropriate shocks, movement, or device failure are possible.

\medterm{Subcutaneous Layer} The tissue layer beneath the dermis, made mainly of fat and connective tissue. It cushions deeper structures, insulates the body, stores energy, and contains larger blood vessels and nerves.

\medterm{Subcutaneous Tissue} The layer beneath the skin that connects it to deeper structures. It contains fat, connective tissue, larger blood vessels, and nerves and provides insulation, cushioning, and energy storage.

\medterm{Subdural} Located beneath the dura, one of the protective coverings around the brain and spinal cord. A subdural hematoma is bleeding in this space, often after head injury, and may cause headache, vomiting, confusion, weakness, seizures, or loss of consciousness. Significant bleeding needs urgent assessment.
\medterm{Subdural Effusion} Accumulation of usually clear fluid in the subdural space, often after trauma, surgery, infection, or cerebrospinal-fluid disturbance and sometimes causing headache or pressure effects.

\medterm{Subdural Haematoma} A collection of blood between protective coverings of the brain, usually after small veins tear during a head injury. It may develop suddenly or gradually and cause headache, confusion, drowsiness, weakness, seizures, or loss of consciousness.

\medterm{Subdural Hematoma} A collection of blood between two protective coverings of the brain, usually after small veins tear during a head injury. It can cause headache, confusion, weakness, seizures, or sleepiness and may require urgent treatment.

\medterm{Subdural Hemorrhage} Bleeding between two protective coverings of the brain, usually after tearing of small veins during a head injury. It may cause sudden or gradually worsening headache, confusion, drowsiness, weakness, seizures, or loss of consciousness and is diagnosed with brain imaging.
\medterm{Subendocardial Myocytes} Heart-muscle cells in the inner layer of the myocardium, which are especially vulnerable to reduced blood flow.

\medterm{Subependymal Giant Cell Astrocytoma} A usually slow-growing brain tumor that develops near a passage where fluid leaves the brain’s cavities. It is strongly associated with tuberous sclerosis and can block fluid flow, causing headache, vomiting, or enlarged brain cavities.

\medterm{Subependymoma} A rare, usually benign slow-growing ependymal tumor, commonly arising in the ventricular system of adults. Symptoms result from obstruction of cerebrospinal-fluid pathways or incidental detection.

\medterm{Subfertility} Alternate index label; see \textit{Infertility}.

\medterm{Subgaleal Hemorrhage} A potentially life-threatening accumulation of blood in the space beneath the scalp's
epicranial aponeurosis. It can expand rapidly in a newborn, particularly after vacuum-assisted delivery, and requires
urgent recognition and treatment.

\medterm{Subinvolution} Slower-than-expected return of an organ to its usual size, especially the uterus after childbirth. It may cause prolonged or heavy bleeding, pelvic pain, or an enlarged tender uterus and needs assessment for retained tissue, infection, or other causes.
\medterm{Subjacent} An anatomical term describing a structure, tissue, or lesion located beneath another structure. It indicates relative position and does not necessarily mean that the two structures are touching; the relationship may be described on physical examination, imaging, or during surgery.

\medterm{Subjective} Based on a person's own experience or report rather than a directly measured finding. Pain, nausea, fatigue, dizziness, and mood are subjective symptoms, while temperature and blood pressure are objective measurements.
\medterm{Sublimation} A psychological defense process in which unacceptable impulses are redirected into socially acceptable activity; in chemistry it is solid-to-gas transition.
\medterm{Sublingual} Sublingual means beneath the tongue; medicines given this way can be absorbed through oral mucosa and may bypass much first-pass liver metabolism.

\medterm{Sublingual Epidermoid Cyst} A benign cyst lined by stratified squamous epithelium beneath the tongue mucosa,
presenting as a slow-growing midline swelling that may require surgical excision if
symptomatic.

\medterm{Sublingual Gland} One of the paired major salivary glands, located beneath the tongue. It produces mainly mucous saliva that enters the
mouth through several small ducts and helps lubricate the oral cavity.

\medterm{Subluxation} Partial displacement of a joint in which some contact between the joint surfaces remains. It may cause pain, swelling, instability, limited movement, nerve injury, or repeated episodes after injury or disease.
\medterm{Submandibular Gland} One of the major salivary glands, located beneath the lower jaw. It makes much of the saliva present when a person is not eating and releases it into the floor of the mouth through a small duct.

\medterm{Submucosa} The supportive layer beneath the moist lining of organs such as the intestine, airway, and bladder. It contains blood vessels, nerves, glands, and connective tissue and helps nourish and anchor the lining.
\medterm{Subphrenic Abscess} A collection of pus beneath the diaphragm, usually caused by infection after abdominal surgery, a burst abdominal organ, or spread from nearby tissue. It may cause fever, upper-abdominal or shoulder pain, breathing discomfort, and sepsis and usually needs drainage and antibiotics.
\medterm{Subq} Subq is an abbreviation for subcutaneous, referring to tissue beneath the skin and the route used for some injections.

\medterm{Subscapular Artery} A branch of the axillary artery that supplies the shoulder blade region, chest wall, and nearby muscles. It contributes to the network of vessels supporting the shoulder and upper limb.
\medterm{Substance} A substance is a material with defined or variable chemical composition; medical effects depend on identity, dose, route, duration, and host factors.

\medterm{Substance Abuse} Legacy, nonspecific label; see \textit{Substance Use Disorder}.

\medterm{Substance Abuse Disorder} Legacy label; see \textit{Substance Use Disorder}.

\medterm{Substance Abuse Psychiatry} An older term for the clinical specialty concerned with substance use disorders and co-occurring mental-health conditions. Current terminology generally uses \textit{Addiction Psychiatry} or \textit{Addiction Medicine}.

\medterm{Substance Dependence} A state in which repeated substance use produces physiologic adaptation, so that reducing or stopping the substance may cause withdrawal. It is distinct from, but may occur with, substance-use disorder.

\medterm{Substance Intoxication} A temporary change in thinking, behavior, or body function caused by recent use of alcohol, a drug, or another substance. Effects depend on the substance, amount, timing, combinations, and the person’s health.

\medterm{Substance Use} The consumption of alcohol, tobacco or nicotine, cannabis, prescription or nonprescription drugs, inhalants, or other psychoactive substances. Use may be occasional, prescribed, nonmedical, or associated with a substance use disorder; the term alone does not imply harmful use.

\medterm{Substance Use - Amphetamines} Use of amphetamine-type stimulants, which can be associated with insomnia, anxiety, psychosis, high blood pressure, arrhythmias, overheating, dependence, and overdose.

\medterm{Substance Use - Cocaine} Use of cocaine, a stimulant that can cause dependence, agitation, seizures, heart attack, stroke, dangerous overheating, and severe psychiatric or cardiovascular complications.

\medterm{Substance Use - Inhalants} Inhalant misuse involves breathing fumes or gases from household or industrial products and can cause sudden cardiac arrest, brain injury, organ damage, dependence, or death.

\medterm{Substance Use - LSD} LSD is a potent hallucinogen that alters perception, mood, and thought and can cause panic, accidents, persistent perceptual symptoms, or psychosis in susceptible people.

\medterm{Substance Use - Marijuana} Marijuana or cannabis can impair attention, memory, coordination, and driving and may cause dependence, anxiety, psychosis, or respiratory harm when smoked.

\medterm{Substance Use - Phencyclidine} Phencyclidine is a dissociative drug that can cause agitation, hallucinations, loss of coordination, high blood pressure, seizures, dangerous behavior, or coma.

\medterm{Substance Use - Prescription Drugs} Use of prescription medicines in a way not directed by a clinician, including taking someone else’s medicine, taking extra doses, or using them for intoxication; dependence, overdose, and interactions may result.

\medterm{Substance Use Disorder} A clinically significant pattern of alcohol, tobacco or nicotine, cannabis, prescription or nonprescription drug, inhalant, or other psychoactive-substance use that causes impaired control, social or occupational impairment, risky use, or clinically significant distress. Features may include craving, unsuccessful efforts to reduce or stop use, substantial time spent obtaining or using the substance, failure to meet major responsibilities, continued use despite harm, tolerance, or withdrawal; tolerance and withdrawal are not required.

\medterm{Substance Use In Pregnancy} Use of alcohol, tobacco or nicotine, cannabis, opioids, cocaine, or other substances during pregnancy. Effects depend on the substance, dose, timing, route, and other factors and may include maternal complications, impaired fetal growth, preterm birth, or neonatal withdrawal.


\medterm{Substance Withdrawal} Physical or mental symptoms that occur after reducing or stopping heavy or prolonged substance use. Symptoms range from anxiety and nausea to seizures, severe confusion, dangerous blood-pressure changes, or breathing problems, depending on the substance.

\medterm{Substernal Goiter} An enlarged thyroid gland that extends from the neck into the upper chest. It may press on the windpipe or esophagus, causing shortness of breath, cough, or difficulty swallowing. Imaging helps assess its size and position.

\medterm{Substrate} A substance on which an enzyme acts or a material used by a biochemical pathway.
\medterm{Subungual Exostosis} A benign osteocartilaginous outgrowth arising beneath the nail, usually of a toe or finger. It causes a painful firm subungual mass and nail deformity and is treated by excision when symptomatic.


\medterm{Succenturial Lobe} An accessory lobe of the placenta, connected to the main placenta by blood vessels that
traverse the membranes. If undetected after delivery, the succenturial lobe may be
retained in the uterus, causing postpartum hemorrhage or infection. Diagnosis: careful
inspection of the placenta and membranes after delivery.

\medterm{Succenturiate Placenta} A succenturiate placenta has one or more accessory lobes separate from the main placenta, which can increase risk of retained tissue or fetal-vessel injury.

\medterm{Succinate Dehydrogenase} An enzyme inside mitochondria that helps cells make energy. It changes one small molecule into another while passing electrons through the energy-producing system; inherited or tumor-related changes can have clinical importance.

\medterm{Succinylcholine} A very fast, short-acting medicine that temporarily paralyzes skeletal muscles to help place a breathing tube or perform emergency procedures. It can dangerously raise potassium, stop breathing longer than expected, or trigger malignant hyperthermia in susceptible people.

\medterm{Succumb} To succumb means to die from or be overcome by a disease, injury, exposure, or other serious process.

\medterm{Sucking Blister} A blister or worn area of skin in a newborn, usually on a hand, finger, forearm, or lip, caused by
repeated sucking before or during feeding. It is usually harmless and heals on its own.
\medterm{Sucking Reflex} A primitive reflex present from birth, elicited by placing an object in the infant's
mouth. The infant rhythmically sucks on the object. Coordinated sucking and swallowing
typically develop by 34-36 weeks gestational age. Premature infants may lack coordinated
suck-swallow, requiring gavage feeding.

\medterm{Sucrose} Common table sugar, a disaccharide made of glucose and fructose. Digestive enzymes split it into these simpler sugars for absorption.
\medterm{Sucrose Lysis Test} An older blood test once used to look for paroxysmal nocturnal hemoglobinuria, a disorder in which red blood cells are unusually vulnerable to destruction. Flow-cytometry testing has largely replaced it because it is more specific and reliable.

\medterm{Suction} Removal of fluid, secretions, blood, or other material using negative pressure. It may be used in airway care, surgery, wound care, or drainage.
\medterm{Suction (Vascular)} A suction device used during a vascular procedure to remove blood or other fluid from the operative
field.

\synonyms
Vascular suction device

\medterm{Suction Device} Equipment that creates negative pressure to remove fluid, blood, mucus, or other material. It usually includes a pump, tubing, collection container, and suction catheter or tip and is used in airway care, surgery, or wound drainage.

\medterm{Suction Tip} The end attachment of a suction device used to remove blood, fluid, secretions, or debris.
Common types include the Yankauer tip for oropharyngeal suction and surgical suction tips such as the Frink tip.

\medterm{Sudden Cardiac Death} Unexpected death from a cardiac cause, usually occurring rapidly after symptom onset or without
witnessed symptoms. Causes include coronary disease, cardiomyopathy, inherited arrhythmia
syndromes, and bradyarrhythmia; prevention depends on the underlying risk.

\synonyms
Arrest Cardiac (Sudden Cardiac Death)
Death Sudden Cardiac (Sudden Cardiac Death)

\medterm{Sudden Infant Death Syndrome} The unexplained sudden death of an infant younger than one year after a complete investigation, including medical history, death-scene review, and autopsy. Safe sleep on the back on a firm, clear surface and avoiding smoke reduce risk, but do not prevent every case.

\medterm{Sudden Sensorineural Hearing Loss} A rapid loss of hearing caused by a problem in the inner ear or hearing nerve, often noticed on waking or over a few days. It may be accompanied by ringing or dizziness and needs urgent assessment because early treatment can improve recovery in some cases.

\medterm{Sudden Unexpected Death In Infancy} The sudden, unexpected death of an infant that requires investigation of the medical history, circumstances, death scene, and sometimes autopsy. The cause may be identified, or the death may meet the definition of sudden infant death syndrome.
\medterm{Sudden Unexpected Infant Death} The sudden and unexpected death of an infant younger than 1 year, whether or not the
cause is known before investigation. Sudden infant death syndrome is the unexplained
subtype after complete scene, history, and autopsy evaluation.

\medterm{Sudek's Atrophy} An older name for complex regional pain syndrome with pain, swelling, autonomic change, and regional bone loss.
\medterm{Sudorifics (Diaphoretics)} Medicines or other treatments that increase sweating. They have limited specific medical uses; excessive sweating caused by a treatment can lead to dehydration, overheating, or abnormal body salts.
\medterm{Suffocation} Inability to breathe enough because the airway is blocked, oxygen is unavailable, the chest is compressed, or another problem prevents normal breathing. It quickly causes unconsciousness and death without immediate rescue.
\medterm{Sugammadex} A medicine used during anesthesia to quickly reverse muscle paralysis caused by selected muscle-relaxing medicines. It binds those medicines so the body can remove them; allergic reactions and changes in heart rate are possible.

\medterm{Sugar-Water Hemolysis Test} An older laboratory test assessing red-cell resistance to hypotonic sugar solution, historically used in evaluating paroxysmal nocturnal hemoglobinuria and now largely replaced by flow-cytometric testing.

\medterm{Suicidal Ideation} Thoughts about dying or killing oneself, with or without a plan or intention to act. Any current plan, intent, access to a method, recent attempt, or inability to stay safe is an emergency requiring immediate support and professional help.

\medterm{Suicide} Death caused by an intentional action to end one's life. Suicidal thoughts, plans, or behavior are urgent medical concerns; immediate support, safety measures, and professional assessment can help protect the person.
\medterm{Suicide And Suicidal Behavior} Thoughts, plans, attempts, or acts intended to cause one’s own death; any current intent, plan, access to means, or inability to stay safe requires immediate emergency support.

\medterm{Suicide Risk Assessment} A clinical evaluation of suicidal thoughts, intent, plan, access to means, past
attempts, self-harm, psychiatric and medical illness, substance use, protective factors,
and immediate safety. It informs but does not replace ongoing clinical judgment and
safety planning.

\medterm{Sulcus} A groove or shallow furrow on an organ or body surface, especially the brain or heart. Brain sulci separate raised areas and increase the cerebral cortex's folded surface area.
\medterm{Sulfa Drugs} A group of medicines, especially sulfonamide antibiotics, that block a chemical pathway bacteria need to grow. They treat selected infections but can cause allergic reactions and other side effects; the term may also refer to nonantibiotic sulfonamides.

\medterm{Sulfamethoxazole} A sulfonamide antibiotic used commonly with trimethoprim in selected infections.
\medterm{Sulfasalazine} An anti-inflammatory medicine used for inflammatory bowel disease and rheumatoid or other inflammatory arthritis. It changes immune activity and can cause nausea, rash, low blood counts, liver problems, or kidney problems.
\medterm{Sulfatide} A fatty substance containing sulfate found mainly in myelin, the insulating covering of nerve fibers. Abnormal buildup or deficiency can disrupt nerve function and occurs in some inherited metabolic disorders.

\medterm{Sulfite Sensitivity} Sulfite sensitivity is an adverse reaction to sulfite preservatives, sometimes causing wheeze, flushing, headache, or other symptoms, especially in people with asthma.

\medterm{Sulfonamide} A medicine containing a sulfonamide chemical group. Some are antibiotics that stop bacteria making folate, which they need to grow; others reduce inflammation or increase urine flow. Allergic reactions and other side effects are possible.

\medterm{Sulfonamides} Sulfonamides are drugs containing a sulfonamide group, including antibacterial agents that inhibit bacterial folate synthesis; allergic reactions, severe skin reactions, cytopenias, renal effects, and crystalluria are possible.

\medterm{Sulfonylurea} An oral diabetes medicine that makes the pancreas release more insulin. It can lower blood sugar but may cause dangerously low blood sugar, hunger, or weight gain.

\medterm{Sulfonylureas} Oral medicines for type 2 diabetes that stimulate the pancreas to release more insulin. They can lower blood glucose but may cause low blood sugar, hunger, or weight gain.

\medterm{Sulfur Colloid Scan} A nuclear-medicine scan in which a small radioactive tracer is used to show activity in the liver, spleen, bone marrow, or lymph nodes. It can help assess how these organs work or identify abnormal tissue distribution.

\medterm{Sulfuric Acid Poisoning} Severe poisoning from sulfuric acid, a highly corrosive chemical. It can cause deep burns of the skin, eyes, mouth, throat, and digestive tract and dangerous fumes; emergency care is required and vomiting must not be induced.

\medterm{Sulindac Overdose} Taking too much sulindac, a nonsteroidal anti-inflammatory drug, may cause nausea, vomiting, stomach bleeding, drowsiness, kidney injury, liver injury, seizures, or metabolic acidosis. A suspected overdose needs prompt medical or poison-control advice.

\medterm{Sulphonamides} A group of medicines, including some antibiotics, that block a chemical pathway bacteria need to grow. They are used for selected infections and other conditions; allergic reactions, interactions, and blood or kidney problems can occur.

\medterm{Sulphonylureas} Oral medicines that stimulate pancreatic insulin release to lower blood glucose in type 2 diabetes.
\medterm{Sumatriptan} A migraine medicine that narrows certain blood vessels and changes nerve signaling to stop acute migraine or cluster headaches.
\medterm{Summer Diarrhoea} An older term for seasonal diarrheal illness, often referring to infectious diarrhea in children.
\medterm{Sun Poisoning} Alternate index label; see \textit{Polymorphous light eruption}.

\medterm{Sun Protection} Sun protection reduces ultraviolet exposure through shade, protective clothing, and broad-spectrum sunscreen and helps prevent sunburn, premature skin aging, and skin cancer.

\medterm{Sun Protection (Photoprotection)} Measures that reduce exposure to ultraviolet radiation, including shade, protective clothing, sunglasses, and broad-spectrum sunscreen. They help prevent sunburn, premature skin aging, eye damage, and ultraviolet-related skin cancers.

\synonyms
Photoprotection

\medterm{Sun-Sensitive Drugs (Photosensitivity To Drugs)} Alternate index label for Photosensitivity to Drugs.


\medterm{Sunburn} An acute inflammatory injury of the skin caused by excessive ultraviolet radiation, resulting in redness, pain, warmth, and sometimes blistering or peeling.
\medterm{Sundowning} Late-day or evening worsening of confusion, agitation, wandering, or behavioral symptoms in some people with dementia or delirium, influenced by fatigue, light, routines, and unmet needs.

\medterm{Sunflower Syndrome} A rare epilepsy in which a person repeatedly turns toward bright light and waves a hand near the eyes, often triggering absence seizures. It usually begins in childhood and requires neurologic assessment and seizure treatment.

\medterm{Sunken Chest} Alternate index label; see \textit{Pectus excavatum}.

\medterm{Sunstroke} A life-threatening form of heat illness in which body temperature rises dangerously high and the nervous system is impaired.

\medterm{Superantigen} A substance made by some bacteria that activates many immune cells at once. This can cause an excessive release of inflammatory chemicals and contribute to illnesses such as toxic shock syndrome or some forms of scarlet fever.




\medterm{Superficial} Near the surface of the body or an organ, rather than deep inside it. The term may describe tissues, veins, burns, wounds, or infections that mainly affect the skin and tissue just beneath it.
\medterm{Superficial Fascia} The soft tissue layer just beneath the skin. It contains fat, small blood vessels, nerves, and connective tissue and helps insulate the body and allow the skin to move.

\synonyms
subcutaneous tissue.

\medterm{Superficial Spreading Melanoma} The most common cutaneous melanoma subtype, characterized initially by radial growth across the epidermis and an irregularly pigmented patch or plaque before invasion develops.

\medterm{Superficial Thrombophlebitis} Inflammation and a blood clot in a vein just beneath the skin, causing a tender, firm, red cord. It often improves with treatment, but a spreading clot or symptoms of deep-vein thrombosis require medical assessment.

\medterm{Superimposed Preeclampsia} Preeclampsia developing in a person who already has chronic hypertension. It may
present with new or worsening proteinuria, thrombocytopenia, abnormal liver or kidney tests, neurologic symptoms, or
other features of maternal organ dysfunction.

\medterm{Superinfection} A new infection that develops during or after treatment of a previous infection, often
involving an antimicrobial-resistant organism or an organism not susceptible to the
original therapy, such as Clostridioides difficile after antibiotic exposure.

\medterm{Superior} Above or toward the head in the body, or higher than another structure. For example, the neck is superior to the chest. The opposite direction is inferior, meaning below or toward the feet.

\medterm{Superior Mesenteric Artery} A major artery branching from the abdominal aorta. It supplies blood to much of the small intestine and to parts of the large intestine, including the right colon and nearby transverse colon.

\medterm{Superior Mesenteric Artery Syndrome} Pressure on part of the small intestine where it passes between two major arteries, causing blockage. It may follow major weight loss and cause fullness, upper-abdominal pain, nausea, and vomiting; treatment restores nutrition and relieves the obstruction.

\medterm{Superior Semicircular Canal Dehiscence} A gap or very thin area in the bone covering one balance canal in the inner ear. Loud sounds or pressure may cause dizziness, imbalance, hearing changes, or hearing one’s own voice unusually loudly; diagnosis uses symptoms, hearing tests, and imaging.

\medterm{Superior Vena Cava} A large vein in the chest that carries oxygen-poor blood from the head, neck, arms, and upper chest to the right side of the heart.

\medterm{Superior Vena Cava Syndrome} Blockage or compression of the large vein carrying blood from the upper body to the heart. It can cause swelling of the face, neck, or arms, visible neck veins, cough, or shortness of breath and may require urgent treatment.
\synonyms
SVCS (Superior Vena Cava Syndrome)

\medterm{Superior Vena Caval Obstruction} Blockage or compression of the large vein returning blood from the head, neck, and arms to the heart. It is often caused by a chest tumor or a blood clot around a catheter and may cause facial or arm swelling, prominent neck veins, cough, or breathlessness.

\medterm{Supernumerary Nipples} Supernumerary nipples are extra nipples along the embryonic milk line, usually harmless but occasionally associated with underlying breast or urinary-tract developmental anomalies.

\medterm{Supernumerary Rib} An extra rib beyond the typical twelve pairs, most commonly a cervical rib above the
first thoracic rib, which may compress the brachial plexus or subclavian vessel causing
thoracic outlet syndrome.

\medterm{Supernumerary Teeth} Teeth that develop in addition to the usual number. They may remain hidden or appear in the mouth and can crowd, displace, or prevent nearby teeth from erupting; a common type develops between the upper front teeth.

\medterm{Supervised Exercise Therapy} A structured, monitored exercise program used to improve walking ability and
reduce claudication symptoms in people with peripheral artery disease. The exercise plan
is adjusted to the person’s symptoms and functional capacity.

\medterm{Supination} Rotation of the forearm so that the palm faces anteriorly in anatomical position, or upward when the elbow is flexed. The radius and ulna become parallel.
\medterm{Supine} Lying flat on the back with the face and front of the body facing upward. This position is commonly used for examinations, imaging, and surgery. Staying in this position for a long time can increase the risk of pressure injuries; late in pregnancy, it may also reduce blood return to the heart.
\medterm{Supine Hypotensive Syndrome} A drop in blood pressure that can occur when a pregnant person lies flat on the back, because the enlarged uterus reduces blood return to the heart. Dizziness, nausea, sweating, or faintness may improve by turning onto the side.

\medterm{Supine Position} Lying on the back with the face and front of the body upward. It is commonly used for examination, imaging, and surgery, but prolonged use can cause pressure injury and late pregnancy may reduce blood return to the heart.

\medterm{Supplemental Oxygen Therapy} Administration of oxygen at a concentration above that of room air to treat or prevent
clinically important hypoxemia. Delivery method and oxygen target depend on the illness,
baseline oxygenation, risk of carbon-dioxide retention, and monitoring results; excessive
oxygen can also cause harm.

\medterm{Supplementation} Use of a vitamin, mineral, protein, fiber, or other nutrient preparation to correct a
deficiency or meet a clinically identified increased requirement. The indication, dose,
interactions, product quality, and risk of toxicity should be assessed before use.


\medterm{Supportive Care} Treatment that maintains essential body functions, relieves symptoms, prevents complications, and supports recovery while the underlying cause is treated.

\medterm{Suppository} A solid preparation inserted into the rectum, vagina, or urethra where it dissolves or melts.
\medterm{Suppression} Reduction or inhibition of a function, response, symptom, or process.
\medterm{Suppuration} The formation or discharge of pus, a thick fluid made from immune cells, germs, and damaged tissue. It usually occurs with infection and may form an abscess that needs drainage and treatment.
\medterm{Supracondylar Humerus Fracture} A break in the upper part of the elbow end of the arm bone, common in children after a fall on an outstretched hand. Swelling, pain, deformity, or inability to move the elbow may occur. Doctors must check circulation and nerves because severe fractures can damage an artery or cause dangerous swelling; some need urgent surgery.

\synonyms
Gartland Supracondylar Fracture; Pediatric Elbow Fracture

\medterm{Supraglottic Airway Device} A breathing device placed above the vocal cords to help deliver oxygen and ventilation during anesthesia or emergencies. It can be inserted more quickly than a breathing tube but does not protect against aspiration as reliably.

\medterm{Supranuclear Ophthalmoplegia} Weakness or loss of control of eye movements caused by damage to brain pathways above the eye-movement nerves. It may make it difficult to look in certain directions.

\medterm{Supranuclear Palsy} Alternate index label; see \textit{Progressive supranuclear palsy}.

\medterm{Suprapubic} Located above the pubic bone in the lower abdomen. The term may describe pain or examination in this area, or a procedure such as inserting a catheter through the lower abdominal wall directly into the bladder when a tube cannot be passed through the urethra.
\medterm{Suprapubic Catheter Care} Care of a drainage tube that enters the bladder through the lower abdomen, including cleaning the site, keeping the tube unobstructed, securing it, and monitoring urine. Fever, increasing pain, leakage, bleeding, or no urine flow needs medical advice.

\medterm{Suprarenal Glands (Adrenal Glands)} Two small glands above the kidneys. Their outer layer makes hormones that regulate salt, blood pressure, stress, and metabolism; their inner layer makes adrenaline and related hormones.
\medterm{Suprascapular Nerve} A nerve from the brachial plexus supplying shoulder muscles that help lift the arm and rotate it outward.

\medterm{Supraventricular Tachycardia} A fast heart rhythm beginning above the ventricles, often starting and stopping suddenly, causing palpitations, dizziness, or weakness.

\synonyms
SVT

\medterm{Supraventricular Tachycardia Paroxysmal (Paroxysmal Supraventricular Tachycardia)} Alternate index label for Paroxysmal Supraventricular Tachycardia.

\medterm{Suramin} An antiparasitic medicine used mainly for selected early forms of African trypanosomiasis. It requires specialist use because it can cause serious kidney, nerve, allergic, and other adverse effects.
\medterm{Surfactant} A substance that lowers surface tension. In the lungs, it coats the tiny air sacs and helps keep them open when a person breathes out; babies born very early may have too little and need treatment.
\medterm{Surgeon} A physician trained to perform operations and invasive procedures, and to evaluate and care for patients before and after surgery.
\medterm{Surgery} Treatment in which a clinician operates on tissues or organs to diagnose, remove, repair, replace, or improve a body structure. Risks and recovery depend on the procedure and the patient's health.
\medterm{Surgery For Pancreatic Cancer} An operation to remove a pancreatic tumor and sometimes nearby organs or lymph nodes when the cancer can be safely removed. It is major surgery and may cause bleeding, infection, digestive problems, diabetes, or delayed stomach emptying.


\medterm{Surgery For Pilonidal Cyst} An operation to drain or remove a recurrent or complicated tunnel or cyst near the cleft of the buttocks. Wound care and keeping the area clean are important; infection, slow healing, and recurrence can occur.

\medterm{Surgical Drain} A tube or device placed during or after surgery to remove blood, pus, bile, urine, or other fluid. The amount and appearance of drainage are monitored; infection, bleeding, blockage, and tissue injury are possible.

\medterm{Surgical Drill} A powered instrument used to cut, shape, or make holes in bone during surgery. It requires careful control to avoid heat injury, damage to nearby tissues, infection, or unintended bone weakening.

\medterm{Surgical Excision} Cutting out a lesion, tumor, organ, or other unwanted tissue using an operation. The removed material may be examined under a microscope; bleeding, infection, scarring, damage to nearby structures, and incomplete removal are possible.

\medterm{Surgical Fixation} Stabilization of a bone, joint, or other structure with devices such as plates, screws, rods, wires, or external frames.

\medterm{Surgical Hemostasis} Control of bleeding during an operation using pressure, ligation, sutures, clips,
electrosurgery, topical agents, vessel sealing, or other techniques. Effective
hemostasis requires identification of the bleeding source and correction of coagulopathy
when present.

\medterm{Surgical Instruments} Tools used during operations, including scalpels, scissors, forceps, clamps, retractors, needle holders, and devices controlling bleeding. Each instrument is designed for a specific task and must be cleaned, sterilized, and handled safely.

\medterm{Surgical Margin} The outer edge of tissue removed during an operation. A laboratory examination checks whether disease reaches that edge: a clear margin suggests no disease was seen there, while an involved margin may mean more treatment is needed.

\medterm{Surgical Repair} An operation that restores or rebuilds damaged tissue, an organ, or a body function. It may close wounds, reconnect structures, stabilize bones, repair blood vessels, or correct defects present from birth.

\medterm{Surgical Risk Stratification} Estimation of the likelihood of perioperative complications using patient factors,
urgency, procedure risk, functional capacity, and disease-specific assessment. It
informs optimization, monitoring level, shared decisions, and postoperative planning but
cannot eliminate uncertainty.

\medterm{Surgical Safety Checklist} A structured checklist used at defined points before anesthesia, before incision, and
before leaving the operating room to verify patient identity, procedure, site, consent,
equipment, specimens, and team communication.

\medterm{Surgical Site Infection} Infection involving the incision, organ, or operative space within the surveillance
period after surgery. It may be superficial incisional, deep incisional, or organ-space
and requires appropriate culture, antimicrobial therapy, drainage, debridement, or
reoperation when indicated.

\medterm{Surgical Stapler} A device that places metal or absorbable staples to close skin, join tissue, or seal a hollow organ during an operation. Incorrect placement can cause bleeding, leakage, narrowing, or injury and must be checked before completion.

\medterm{Surgical Time-Out} A deliberate team pause immediately before incision to verify patient identity,
procedure, site, consent, allergies, antibiotics, imaging, implants, equipment, and
anticipated critical events. It is a core surgical-safety intervention.

\medterm{Surgical Wound Care - Closed} Care of an incision whose edges are closed with sutures, staples, adhesive, or strips, including keeping it clean and dry as instructed and watching for infection or separation.

\medterm{Surgical Wound Care - Open} Care of a surgical wound left partly or completely open to drain or heal gradually, usually with dressings and scheduled assessment for infection, bleeding, and healthy tissue growth.

\medterm{Surgical Wound Classification} Categorization of an operative wound as clean, clean-contaminated, contaminated, or
dirty/infected according to the degree of microbial exposure and existing infection. The
class helps estimate surgical-site infection risk and guide prophylaxis.

\medterm{Surgical Wound Infection - Treatment} Management of infection in or around a surgical wound, including assessment of depth and severity, drainage or debridement when needed, wound care, and targeted antimicrobial therapy when indicated.

\medterm{Surrogacy} Surrogacy is an arrangement in which a person carries a pregnancy for intended parent or parents; medical, legal, ethical, and psychosocial requirements vary by jurisdiction.

\medterm{Surrogate} A substitute or proxy. In medical research, a surrogate outcome is a measurable change, such as a blood-test result, used to estimate a patient-important outcome such as fewer symptoms or longer survival.
\medterm{Surrogate Decision-Maker} A person authorized to make healthcare decisions for someone who lacks decision-making capacity, guided by the person's
known values and preferences.

\medterm{Surveillance} The planned collection and review of health information to detect changes in disease, follow a known condition, or identify problems early. It may involve routine reports, active case finding, selected monitoring sites, or regular follow-up of an individual.

\medterm{Surveillance Imaging} Planned repeat imaging used to monitor disease, treatment response, or possible recurrence. The timing and type of scan depend on the condition, previous treatment, symptoms, risks, and expected usefulness of results.

\medterm{Survival Rate} The percentage of people alive after a specified time from diagnosis or treatment, commonly measured at five years.

\medterm{Survivorship} The period of living during and after cancer treatment, together with the care provided during that time. It may include checking for return of cancer, treating late effects, supporting emotional health, improving quality of life, and managing other illnesses.

\medterm{Survivorship Care Plan} A written summary of cancer treatment, recommended follow-up, possible late effects, screening, medicines, and contacts for ongoing care. It helps the person and healthcare team coordinate care after treatment.

\medterm{Susac Syndrome} A rare disease in which inflammation and blockage of tiny blood vessels can affect the brain, eyes, and inner ears. It may cause confusion or headache, vision loss, and hearing loss; all features may not appear at once.

\medterm{Susceptibility} The likelihood of being affected by a disease, exposure, medicine, or microorganism. Susceptibility can reflect genetics, age, immunity, health conditions, previous exposure, dose, or the organism’s drug resistance.
\medterm{Susceptibility Testing} Laboratory testing that shows which antibiotics are likely to stop a person’s bacterial infection. The result helps clinicians choose treatment, especially when resistance is possible.

\medterm{Suspected Cause} A condition, exposure, or factor believed to explain a symptom or disease but not yet proven. Further history, examination, testing, or observation may support or reject the suspected cause.
\medterm{Suspension} A liquid medicine in which fine drug particles are spread through the liquid rather than dissolved. The container may need shaking so each dose contains the intended amount; settling, inaccurate measuring, and contamination can cause problems.

\medterm{Suture} A stitch or material used to close or join tissue.
\medterm{Suture (Perineal)} A thread or strand of material used to approximate tissues (bring wound edges together)
and hold them in place until healing occurs. See Suture.

\synonyms
Perineal

\medterm{Suture Anchor} A small implant that secures a tendon, ligament, or other soft tissue to bone. It may be
made from metal, polymer, or bioabsorbable material and can fail through pullout,
breakage, migration, or inflammatory reaction.

\medterm{Sutures - Ridged} Ridged cranial sutures are raised skull seams in an infant and may reflect normal molding or premature suture fusion, which requires assessment for craniosynostosis.

\medterm{Sutures - Separated} Separated cranial sutures in an infant can reflect normal growth or increased intracranial pressure, hydrocephalus, prematurity, or other illness and require clinical assessment.

\medterm{SUV} Standardized uptake value, a number estimating how much radioactive tracer is concentrated in tissue during a PET scan. Clinicians compare it with surrounding tissue and clinical findings to assess tumors or inflammation.

\medterm{SVC Obstruction} Blockage or compression of the superior vena cava, the large vein returning blood from the head, neck, arms, and upper chest to the heart. It may cause facial or arm swelling, prominent neck veins, cough, headache, or breathing difficulty.

\medterm{SVC Syndrome} Blockage or compression of the superior vena cava, the large vein returning blood from the upper body to the heart. It may cause swelling of the face, neck, or arms, visible neck veins, cough, or breathing difficulty.

\medterm{SVT} Alternate index label; see \textit{Supraventricular tachycardia}.

\medterm{Swab} A small absorbent stick used to collect a sample from the skin, throat, nose, wound, or another body site. The sample can be tested for infection or other conditions, and correct collection affects the result.
\medterm{Swallow Study} An X-ray test that records a person swallowing food or liquid mixed with barium. It shows how the mouth and throat work, whether material enters the airway, and whether changes could make swallowing safer.

\medterm{Swallowing} The movement of food, liquid, or saliva from the mouth through the throat and esophagus into the stomach. Problems may cause coughing, choking, food sticking, or material entering the airway.

\synonyms
Difficulty Swallowing (Swallowing)

\medterm{Swallowing Assessment} An evaluation of how safely and effectively a person moves food, liquid, or saliva from the mouth to the stomach. It may include observation, X-ray or camera testing, and recommendations about exercises, posture, or food and drink textures.

\medterm{Swallowing Chalk} Ingestion of chalk, which is usually low in toxicity in small amounts but can cause choking, constipation, or gastrointestinal upset; repeated eating may indicate pica.

\medterm{Swallowing Difficulty} Dysphagia, or impaired movement of food or liquid from the mouth to the stomach, caused by oropharyngeal, esophageal, neuromuscular, structural, or functional disorders.

\medterm{Swallowing Problems} Difficulty moving food, liquid, or saliva safely from the mouth to the stomach; dysphagia can cause choking, aspiration, weight loss, or dehydration and may need evaluation.

\medterm{Swallowing Soap} Ingestion of soap, which commonly causes mouth or stomach irritation, nausea, or diarrhea; concentrated or large exposures can injure tissue or affect breathing.

\medterm{Swallowing Sunscreen} Ingestion of sunscreen, usually causing mild gastrointestinal upset but potentially causing choking or aspiration depending on amount and product; breathing difficulty or persistent vomiting requires urgent help.

\medterm{Swallowing Syncope} Fainting triggered by swallowing, usually through an abnormal reflex that slows the heart or lowers blood pressure; recurrent episodes require evaluation for rhythm and structural heart disease.

\medterm{Swallowing Therapy} Assessment and treatment to make swallowing safer and more effective. A speech-language therapist may teach exercises, positioning, food-texture changes, and strategies to reduce aspiration into the airway.
\medterm{Swan-Ganz - Right Heart Catheterization} A procedure in which a thin catheter is passed through the right side of the heart into a lung artery to measure heart pressures and blood flow. It is used selectively in critically ill patients because bleeding, infection, rhythm problems, or vessel injury can occur.

\medterm{Swan-Ganz Catheter} A thin tube passed through the right side of the heart into a lung artery to measure pressures and estimate how much blood the heart pumps. It is used selectively in seriously ill patients because complications are possible.
\medterm{Sweat} Fluid produced by sweat glands, mainly to cool the body as it evaporates. Sweating can also increase with exercise, fever, stress, hormones, or some medicines.
\medterm{Sweat Chloride Test} A test that measures salt levels in sweat to help diagnose cystic fibrosis. A high result supports the diagnosis, but it is interpreted with symptoms and may be repeated or followed by genetic testing because other factors can affect results.

\medterm{Sweat Electrolytes Test} A sweat electrolytes test measures salts, especially chloride, in sweat and is commonly used to help diagnose cystic fibrosis.

\medterm{Sweat Gland} A skin gland that produces sweat, which contributes mainly to temperature regulation and skin surface function.

\medterm{Sweat Gland Tumor} A growth arising from cells of a sweat gland. Most are benign skin lumps, but some are cancerous and may grow, ulcerate, bleed, or spread. Examination and biopsy determine the type; complete removal is often used, with wider treatment and follow-up for cancerous tumors.

\medterm{Sweat Glands} Eccrine and apocrine glands that produce sweat or related secretions.
\medterm{Sweating} Production of sweat by eccrine and apocrine glands for heat regulation and other responses; excessive, absent, or asymmetric sweating can indicate medication, autonomic, endocrine, or neurologic disease.

\medterm{Sweating And Body Odor} Excess perspiration or noticeable odor produced when skin bacteria break down sweat; it may be influenced by hygiene, diet, medicines, hormones, or disease.

\synonyms
Body Odor And Sweating

\medterm{Sweating, Abnormally Excessive} Alternate index label; see \textit{Hyperhidrosis}.

\medterm{Sweating, Gustatory} Sweating triggered by eating or smelling food, usually over the cheek or temple after parotid-nerve injury or surgery and known as Frey syndrome.

\medterm{Sweet Syndrome} An inflammatory skin disorder characterized by sudden fever and tender red or violaceous plaques or nodules containing
many neutrophils. It may be associated with infection, medicines, inflammatory disease,
or cancer.

\medterm{Sweet's Syndrome} An inflammatory skin disorder causing sudden fever and painful red or purple raised patches or nodules. It may follow infection, result from a medicine, or occur with inflammatory disease or cancer, especially blood cancers.

\medterm{Sweetened Beverages} Sweetened beverages contain added sugars and can increase calorie intake, dental caries, obesity, and diabetes risk when consumed frequently.

\medterm{Sweeteners - Sugar Substitutes} Sugar substitutes provide sweetness with little or no sugar; types differ in energy, digestion, aftertaste, and effects on glucose and gastrointestinal symptoms.

\medterm{Sweeteners - Sugars} Sugars are simple carbohydrates that provide energy but can raise blood glucose and contribute to dental decay or excess calories when consumed in excess.

\medterm{Swelling} Abnormal enlargement caused by fluid accumulation, inflammation, bleeding, infection, tumor, or tissue injury; location, onset, pain, and associated symptoms identify the likely cause.

\medterm{Swimmer's Ear} Swimmer’s ear, also called otitis externa, is inflammation or infection of the skin lining the outer ear canal. It often develops when moisture remains in the canal, allowing bacteria or fungi to grow. Scratching, cotton swabs, or other irritation can also damage the canal’s protective skin. Symptoms include ear pain, itching, redness, swelling, drainage, and pain when the outer ear is touched.

\synonyms
Otitis Externa

\medterm{Swimmer's Itch} An itchy rash caused by an allergic reaction to microscopic parasites released by infected snails in freshwater or saltwater.

\synonyms
Cercarial Dermatitis
Dermatitis, Cercarial

\medterm{Swimmers Ear (Otitis Externa)} Alternate index label for Otitis Externa.

\medterm{Swimming Pool Cleaner Poisoning} Poisoning from pool chemicals such as chlorine compounds, acids, or alkalis. They may burn tissues or release toxic gases, causing cough, chest pain, breathing difficulty, vomiting, or shock; leave fumes and seek emergency help.

\medterm{Swimming Pool Granuloma} A slowly developing skin infection caused by a water-dwelling bacterium after exposure to an aquarium, swimming pool, or natural water. It can cause persistent bumps or sores, often spreading along the arm or hand, and needs specific antibiotic treatment.

\medterm{Swine Flu} Alternate informal term; see \textit{Influenza A Virus}.


\medterm{Swollen Ankles And Swollen Feet} Enlargement of the ankles or feet caused by injury, venous insufficiency, lymphatic disease, medication,
pregnancy, or systemic fluid retention from heart, kidney, or liver disease. Sudden one-sided swelling with pain,
redness, warmth, or breathlessness may indicate thrombosis or infection and requires urgent assessment.

\synonyms
Ankle Swollen (Swollen Ankles And Swollen Feet)

\medterm{Swollen Lymph Nodes} Swollen lymph nodes are enlarged immune glands, commonly from infection but sometimes from inflammation, medication, or cancer; persistent, hard, or generalized nodes need evaluation.

\synonyms
Lymph Nodes, Swollen
Lymph Swollen Glands (Swollen Lymph Nodes)
Lymph Swollen Nodes (Swollen Lymph Nodes)

\medterm{Swollen Tongue} Enlargement of the tongue caused by injury, infection, allergy, inflammation, bleeding, or another condition; sudden breathing difficulty is an emergency.

\medterm{Sycosis Barbae} Long-lasting inflammation or infection of hair follicles in the beard area, usually involving bacteria and ingrown hairs. It can cause painful bumps, pus, crusting, scarring, and hair loss and needs appropriate treatment.
\medterm{Sydenham Chorea} A delayed neurologic complication of group A streptococcal infection causing involuntary jerky movements, hypotonia, emotional lability, and coordination problems, usually in children.

\medterm{Sydenham's Chorea} A delayed immune complication of group A streptococcal infection, usually in children. It causes involuntary, irregular movements and may cause clumsiness, muscle weakness, emotional changes, or difficulty with speech and daily tasks.
\medterm{Symblepharon} Scar tissue that joins the inside of the eyelid to the eye’s surface, usually after a severe burn, inflammation, or eye disease. It can restrict eye movement, cause dryness or pain, and may need specialist treatment.


\medterm{Symbolism} The use of an image, object, or action to represent an idea or feeling; in clinical settings, symbolic meaning may be relevant to communication, culture, or psychotherapy.

\medterm{Sympathectomy} Surgical interruption or destruction of sympathetic nerves to reduce abnormal vasoconstriction, sweating, or pain in selected conditions.

\medterm{Sympathetic} Relating to the part of the autonomic nervous system that prepares the body for activity or stress by changing heart rate, blood-vessel tone, sweating, breathing, and digestion. In everyday language, it can also mean showing concern for another person.
\medterm{Sympathetic Ganglion} A small group of nerve cells outside the brain and spinal cord that relays automatic signals to blood vessels, sweat glands, and internal organs. These signals help control heart rate, blood pressure, digestion, and body responses to stress.

\medterm{Sympathetic Nervous System} The autonomic division that prepares the body for activity or stress by increasing heart rate and altering vascular, airway, and digestive function.
\medterm{Sympatholytic} A medicine or treatment that reduces adrenaline-like nerve activity. It may lower blood pressure, slow the heart, or reduce excessive sweating, but the effects and risks depend on the specific medicine and require appropriate prescribing.

\medterm{Sympathomimetic Drugs} Medicines that imitate the body's adrenaline-like nerve activity, often by stimulating alpha or beta receptors. They may open airways, raise blood pressure, or reduce congestion but can cause fast heartbeat, anxiety, or dangerous effects.
\medterm{Sympathomimetic Toxidrome} A poisoning pattern caused by excessive adrenergic stimulation, producing agitation, sweating, dilated pupils, rapid heart rate, high blood pressure, elevated temperature, tremor, or seizures.

\medterm{Symphysis} A joint where two bones are joined by a pad of tough cartilage. The pubic symphysis connects the front pelvic bones and normally allows slight movement, especially during pregnancy and childbirth.

\medterm{Symphysis-Fundal Height} The distance from the top of the pubic bone to the top of the uterus, measured through the abdomen during pregnancy. From about 24 to 36 weeks, the measurement often roughly matches the number of weeks of pregnancy. A larger or smaller result may lead to an ultrasound to check fetal growth, fluid, or the pregnancy dates.

\synonyms
Fundal Height Measurement; SFH

\medterm{Symphysis Pubis} The fibrocartilaginous joint connecting the left and right pubic bones at the front of the pelvis.

\medterm{Symptom} A change or problem felt or reported by a patient that may suggest illness, such as pain, nausea, or breathlessness.
\medterm{Symptomatic} Having symptoms caused by a disease or condition; symptomatic treatment relieves those symptoms without necessarily treating the underlying cause.

\medterm{Symptomatic Treatment} Care aimed at relieving symptoms such as pain, fever, nausea, itching, or cough without directly removing the underlying cause. It may be used alone or alongside treatment for the disease itself.

\medterm{Synapse} A connection where one nerve cell communicates with another nerve cell, muscle, or gland. Signals cross the small gap using chemicals or electrical changes to control movement, sensation, thought, and body functions.
\medterm{Synapse Assembly} The developmental process by which neurons form organized connections with other neurons or target cells, using cell-adhesion molecules and signaling proteins.

\medterm{Synapsins} A family of neuronal proteins that help organize synaptic vesicles and regulate the release of neurotransmitters.

\medterm{Synaptic Cleft} The very small gap between a presynaptic neuron and its target cell across which neurotransmitters travel.

\medterm{Synaptic Membrane} The specialized membrane at a neuronal junction, including the neurotransmitter-releasing presynaptic membrane and receptor-containing postsynaptic membrane.

\medterm{Synaptic Plasticity} The ability of connections between nerve cells to become stronger, weaker, or structurally different in response to activity. It supports learning, memory, adaptation, and some changes in chronic pain.

\medterm{Synaptic Transmission} Communication between nerve cells or between a nerve cell and another target across a synapse. Electrical signals trigger release of chemical messengers, which bind receptors and change the receiving cell’s activity.
\medterm{Synaptic Vesicle} A small membrane-bound sac in a nerve ending that stores neurotransmitters and releases them into the synaptic cleft when the neuron is activated.

\medterm{Synaptonemal Complex} A temporary protein structure that pairs matching chromosomes during egg or sperm formation. It helps chromosomes exchange genetic material accurately before cells divide and is important for normal reproductive development.


\medterm{Synaptosome} A sealed nerve-terminal structure produced when brain tissue is mechanically disrupted; it is used to study neurotransmitter release and synaptic function.

\medterm{Synaptotagmin} A calcium-sensing protein on nerve-cell storage sacs that helps trigger rapid release of chemical nerve signals. It is essential for communication between nerve cells and between nerves and muscles.

\medterm{Synarthrosis} A joint with little or no movement, such as a skull suture or tooth anchored in its socket, in which bones are joined by fibrous tissue or cartilage.



\medterm{Syncope} A sudden, brief loss of consciousness and postural tone caused by temporary global cerebral hypoperfusion, followed by spontaneous, usually complete recovery. Common mechanisms are reflex syncope, orthostatic hypotension, and cardiac disease; prolonged confusion or seizure-like features may indicate another cause.
\medterm{Syncope (Fainting)} A brief loss of consciousness and muscle control caused by a temporary reduction in blood flow to the brain, followed by spontaneous recovery. Common causes include reflex fainting, a sudden blood-pressure drop, dehydration, and heart disease; prolonged confusion may suggest another cause.

\synonyms
Fainting

\medterm{Syncope, Coughing} A brief faint caused by coughing, which raises chest pressure, reduces venous return and cardiac output, and transiently lowers cerebral blood flow.

\medterm{Syndactyly} Congenital joining or webbing of two or more fingers or toes. It may occur alone or with a genetic syndrome and sometimes requires reconstructive surgery.
\medterm{Syndecan} A group of proteins on the cell surface that help cells attach to their surroundings and receive signals. They contribute to tissue development, healing, and responses to injury.


\medterm{Syndesmophyte} A thin, bony growth that forms along the edges of the vertebrae, often in ankylosing spondylitis, and may bridge adjacent vertebrae.

\medterm{Syndrome} A group of symptoms, signs, or test findings that commonly occur together. A syndrome may have several possible causes and does not always name one specific disease.
\medterm{Syndrome Asperger (Asperger Syndrome)} Alternate index label for Asperger Syndrome.

\medterm{Syndrome Cauda Equina (Cauda Equina Syndrome)} Alternate index label for Cauda Equina Syndrome.

\medterm{Syndrome Compartment (Compartment Syndrome)} Alternate index label for Compartment Syndrome.

\medterm{Syndrome Cyclic Vomiting (Cyclic Vomiting Syndrome)} Alternate index label for Cyclic Vomiting Syndrome.

\medterm{Syndrome Loeys-Dietz (Loeys-Dietz Syndrome)} Alternate index label for Loeys-Dietz Syndrome.

\medterm{Syndrome Of Inappropriate ADH Secretion} Excessive antidiuretic hormone activity causing water retention and low blood sodium despite no major swelling or dehydration. It may result from lung disease, medicines, cancer, or brain disorders and can cause confusion or seizures.
\medterm{Syndrome Of Inappropriate Antidiuretic Hormone} Excessive release or effect of antidiuretic hormone causing impaired free-water excretion, low serum sodium, and inappropriately concentrated urine in the absence of an appropriate physiologic stimulus.

\medterm{Syndrome Of Inappropriate Antidiuretic Hormone Secretion} Excessive activity of the water-retaining hormone vasopressin causes the body to retain water and lowers the blood sodium level. Causes include medicines, lung disease, cancer, and brain disorders; severe cases can cause confusion or seizures.

\synonyms
Syndrome of Inappropriate Antidiuretic Hormone Secretion

\medterm{Syndrome Pfeiffer (Pfeiffer Syndrome)} Alternate index label for Pfeiffer Syndrome.

\medterm{Syndrome Restless Legs (Restless Leg Syndrome)} Alternate index label for Restless Leg Syndrome.


\medterm{Syndrome X} Alternate index label; see \textit{Metabolic syndrome}.

\medterm{Syndrome, Nervous Colon} An older name for irritable bowel syndrome, a disorder causing recurrent abdominal pain with diarrhea, constipation, or both without a structural explanation. Symptoms may improve with individualized diet, medicines, and stress management.

\medterm{Syndrome, Post-Polio} New or worsening weakness, fatigue, muscle pain, breathing problems, or swallowing difficulty occurring years after recovery from polio. It reflects late effects on previously affected motor neurons and is managed with pacing, rehabilitation, and support.

\medterm{Syndrome, TAR} Thrombocytopenia-absent-radius syndrome, an inherited condition causing low platelet counts and absent or severely shortened radius bones in the forearms. It may also cause limb differences, heart defects, kidney problems, or milk intolerance.

\medterm{Synechia} An abnormal band of scar-like tissue joining surfaces that should remain separate. In the eye, it may attach the iris to the lens or cornea and interfere with fluid drainage or vision.

\medterm{Synechiae} Abnormal bands of scar-like tissue joining surfaces that should remain separate. They may occur in the eye, uterus, nose, or other sites and can restrict movement, block passages, or impair function.
\medterm{Synergist} A muscle, drug, or factor that works together with another to produce or enhance an effect.
\medterm{Synesthesia} Synesthesia is a stable perceptual phenomenon in which stimulation of one sensory pathway automatically produces an experience in another, such as seeing colors with sounds.


\medterm{Synostosis} Fusion of two bones by bone tissue, present from birth or developing later. It may be harmless, but fusion in the skull, forearm, or other important area can restrict growth or movement.
\medterm{Synovectomy} An operation to remove part or all of an inflamed lining inside a joint. It may reduce pain and swelling in selected cases of arthritis or repeated joint inflammation and can be done through small cuts or open surgery.
\medterm{Synovial Biopsy} A synovial biopsy removes a small sample of the joint lining to investigate infection, inflammatory arthritis, crystal disease, or an unusual joint lesion.

\medterm{Synovial Bursa} A small, fluid-filled sac near a joint that reduces rubbing between moving tissues such as tendons, muscles, skin, and bone. Irritation or inflammation of a bursa causes bursitis, pain, and swelling.

\medterm{Synovial Chondromatosis} A usually noncancerous condition in which the joint lining forms many small cartilage nodules that may break loose. It can cause pain, swelling, catching, and limited movement, often in the knee, hip, or elbow.

\medterm{Synovial Cyst} A synovial cyst is a fluid-filled sac arising from or near a joint or tendon sheath, sometimes compressing nearby nerves and causing pain or weakness.

\medterm{Synovial Fluid} A slippery fluid inside a freely movable joint that reduces friction and helps nourish cartilage. Changes in its amount or contents can occur with injury, infection, bleeding, crystals, or arthritis.
\medterm{Synovial Fluid Analysis} Testing fluid removed from a swollen or painful joint. The laboratory examines its appearance, cells, crystals, and germs to help distinguish injury, arthritis, gout, bleeding, and joint infection; a very hot painful joint needs urgent assessment.

\medterm{Synovial Joint} A freely movable joint in which the articulating bones are separated by a fluid-filled joint cavity enclosed by a capsule.

\medterm{Synovial Membrane} The thin lining inside a freely moving joint. It makes synovial fluid, which reduces rubbing and helps nourish the smooth joint surfaces; inflammation of this lining causes swelling and pain.
\medterm{Synovial Reaction To Implanted Prostheses} Inflammation or tissue change around an artificial joint caused by wear particles, loosening, corrosion, infection, or an immune reaction. Pain, swelling, or reduced function after replacement needs evaluation to identify the cause.

\medterm{Synovial Sarcoma} A cancer of soft tissue that often develops near a joint in an arm or leg, especially in adolescents and young adults. Diagnosis requires tissue testing; treatment may include surgery, radiation, or chemotherapy.

\medterm{Synoviocytes} Synoviocytes are cells lining the joint synovium that produce lubricating fluid and inflammatory mediators and help maintain or damage joint tissue.

\medterm{Synovitis} Inflammation of the lining inside a joint, causing swelling, warmth, pain, or reduced movement. It may follow injury or occur with arthritis, gout, autoimmune disease, or infection; a hot, very painful joint needs urgent assessment.
\medterm{Synovium} Tissue lining a joint or tendon sheath that makes lubricating fluid and helps support smooth movement.
\medterm{Synovography} An imaging test in which contrast is placed inside a joint to outline its capsule, lining, or possible tear. It is used selectively; MRI, ultrasound, or other scans are often preferred for joint problems.


\medterm{Synthetic Androgen} A synthetic androgen is a manufactured compound that activates androgen receptors and may mimic testosterone; therapeutic and anabolic use can suppress gonadal function and cause acne, mood effects, infertility, liver injury, or cardiovascular harm.

\medterm{Syphilis} A sexually transmitted infection caused by Treponema pallidum that progresses through primary, secondary, latent, and tertiary stages and can affect the nervous system, eyes, heart, or fetus.
\medterm{Syphilis Serology (RPR/VDRL)} Blood tests used to screen for syphilis and monitor response to treatment. A reactive result usually needs confirmation with a different test and interpretation alongside symptoms, examination, and treatment history.

\synonyms
Syphilis Serology
\medterm{Syphilitic Alopecia} Non-scarring hair loss occurring as a manifestation of secondary syphilis, presenting
with diffuse thinning or patchy, moth-eaten alopecia of the scalp, eyebrows, or beard.
It is thought to result from direct spirochaetal invasion of hair follicles during the
spirochaetaemic phase.

\medterm{Syphilitic Aseptic Meningitis} Syphilitic aseptic meningitis is inflammation of the meninges caused by Treponema pallidum, producing headache, neck stiffness, cranial-nerve symptoms, or other neurologic findings.

\medterm{Syphilitic Gumma} A destructive lump or area of inflammation caused by long-standing untreated syphilis. It can affect the skin, bones, liver, brain, or other organs and may form a firm mass or ulcer; appropriate antibiotic treatment can prevent further damage.

\medterm{Syringe} A medical device with a hollow barrel, movable plunger, and tip that may connect to a needle, catheter, or other fitting. It is used to draw up, measure, inject, or remove fluids. Common types include disposable, oral, insulin, tuberculin, and safety syringes.

\synonyms
Medical Syringe
\medterm{Syringe Drivers} Small pumps that push medicine from a syringe continuously or at programmed times. They provide controlled delivery when precise dosing is needed, including pain relief, symptom control, or medicines given under the skin.
\medterm{Syringocystadenoma Papilliferum} A usually harmless sweat-gland skin growth, often appearing as a small bump or plaque on the scalp or face. It may develop within a birthmark and can resemble other lesions, so diagnosis may require examination or biopsy.

\medterm{Syringoma} A syringoma is a benign tumor of sweat-duct cells, usually appearing as small skin-colored or yellowish papules around the eyelids or elsewhere.

\medterm{Syringomyelia} A fluid-filled cavity that develops inside the spinal cord. It can slowly damage the cord and cause pain, weakness, loss of pain or temperature sensation, stiffness, or problems with bladder, bowel, or sexual function.
\medterm{Syrinx} A fluid-filled cavity within the spinal cord, usually related to disturbed movement of the fluid around the brain and spinal cord. It can damage nerve tissue and cause pain, weakness, stiffness, or loss of sensation; a syrinx in the cord is called syringomyelia.

\medterm{Syrup} A syrup is a viscous liquid preparation, often sweetened, used as a medicine vehicle or oral dosage form; concentration and ingredients determine safe dosing.


\medterm{Systemic} Affecting the whole body or relating to an entire system rather than a localized site.
\medterm{Systemic Antifungals} Antifungal medicines taken orally or given intravenously for infections that are extensive,
invasive, refractory, or unsuitable for topical treatment. Choice and duration depend on the
organism, site, severity, drug interactions, organ function, and local resistance patterns.

\medterm{Systemic Chemotherapy} Systemic chemotherapy uses medicines that travel through the bloodstream to kill or slow cancer cells throughout the body.

\medterm{Systemic Circuit} The circulation that carries oxygenated blood from the left ventricle through the body and returns deoxygenated blood to the right atrium.

\medterm{Systemic Disease} A disease that affects the whole body or several organs rather than one local area. Examples include autoimmune disease, widespread infection, hormonal disorders, and some cancers.
\medterm{Systemic Effects} Effects involving the whole body rather than one local area. A medicine, infection, toxin, or disease may produce systemic cardiac, neurologic, kidney, liver, blood, or other effects.

\medterm{Systemic Exertional Intolerance Disease} Alternate index label; see \textit{Myalgic encephalomyelitis/chronic fatigue syndrome (ME/CFS)}.

\medterm{Systemic Inflammatory Response Syndrome} A widespread inflammatory reaction that may follow infection, injury, burns, pancreatitis, or surgery. It can cause abnormal temperature, heart rate, breathing rate, or white blood-cell count and may progress to organ failure.

\medterm{Systemic Juvenile Idiopathic Arthritis} A childhood inflammatory disease causing repeated high fevers, a temporary pink rash, and swollen or painful joints. It can also enlarge lymph nodes, liver, or spleen and may affect other organs, requiring specialist monitoring and treatment.

\medterm{Systemic Lupus} Systemic lupus erythematosus is an autoimmune disease that can affect skin, joints, kidneys, blood, nervous system, lungs, and heart. Symptoms fluctuate, and treatment is tailored to organ involvement with sun protection, monitoring, and immunomodulatory medicines.

\synonyms
Arthritis SLE (Systemic Lupus)
Lupus (Systemic Lupus)
SLE (Systemic Lupus)

\medterm{Systemic Lupus Erythematosus} A long-term autoimmune disease in which the immune system mistakenly attacks the body’s own tissues and causes inflammation. It can affect the skin, joints, kidneys, heart, lungs, brain, and blood cells. Symptoms vary widely and may include fatigue, fever, joint pain or swelling, rashes, hair loss, mouth sores, or swelling. The disease often has periods of increased activity called flares and periods of fewer symptoms called remission; inflammation can cause organ damage in some people.
\medterm{Systemic Mastocytosis} A disorder in which too many mast cells build up in the skin, bone marrow, digestive tract, or other organs. The cells can release chemicals causing flushing, itching, faintness, stomach symptoms, bone problems, or severe allergic reactions.

\medterm{Systemic Scleroderma} An autoimmune disease that causes tight, thickened skin and may scar blood vessels and internal organs. The esophagus, lungs, kidneys, heart, fingers, and other areas can be affected, so care depends on the organs involved.

\medterm{Systemic Sclerosis} An autoimmune disease in which the body makes too much scar-like tissue. The skin may become tight and hard, and blood vessels, the esophagus, lungs, kidneys, or heart may also be affected.

\medterm{Systemic Therapy} Treatment that travels through the bloodstream or otherwise affects the whole body. Examples include chemotherapy, hormone treatment, targeted treatment, and immunotherapy; it differs from local treatment such as surgery or focused radiation.

\medterm{Systemic Vascular Resistance} The opposition to blood flow in the body’s arteries, controlled mainly by how narrow or wide the small arteries are. It helps determine blood pressure and the work required of the heart.

\medterm{Supracervical Hysterectomy} Hysterectomy in which the uterine body is removed while the cervix is left in place; it is also called subtotal or partial hysterectomy. The fallopian tubes and ovaries may be managed separately.

\medterm{Suprapubic Catheter} A tube placed through the lower abdominal wall directly into the bladder to drain urine. It may be used when urine cannot drain safely through the urethra or when longer-term bladder drainage is needed.

\medterm{Systemic Venous} Systemic venous refers to veins returning deoxygenated blood from the body to the right side of the heart, excluding the pulmonary circulation.

\medterm{Sustained Virologic Response} No detectable hepatitis C virus in the blood 12 or more weeks after antiviral treatment. This is considered a cure of that infection, but it does not protect against catching hepatitis C again.

\medterm{Systole} The part of each heartbeat when a heart chamber contracts and pushes blood forward. Ventricular systole sends blood to the lungs and body, while arterial pressure rises and the heart's valves coordinate flow.
\medterm{Systolic} Systolic means relating to contraction of the heart, especially the pressure in arteries when the ventricles eject blood.

\medterm{Systolic Anterior Motion} Movement of the mitral valve toward the path where blood leaves the heart during each heartbeat. It can partly block blood flow and cause the valve to leak, especially in hypertrophic cardiomyopathy. Symptoms may include breathlessness, chest pain, dizziness, or fainting; treatment depends on the cause and severity.

\synonyms
SAM; SAM of the Mitral Valve

\medterm{Systolic Murmurs} Extra heart sounds heard while the heart is contracting. They may result from fast blood flow, a narrowed heart valve, or another structural problem; their significance depends on the sound and the person’s symptoms.

\medterm{Systolic Pressure} The maximum arterial pressure during ventricular contraction, recorded as the upper number of a blood-pressure reading.


\medterm{Saffron} Saffron is a spice made from the dried stigmas of Crocus sativus. It is used in food and studied as a possible supplement, but it is not a substitute for proven medical treatment.

\medterm{Science Communication} Science communication is the clear sharing of scientific information with patients, professionals, policymakers, or the public.

\medterm{Sclerotomy} A sclerotomy is a surgical opening made through the sclera, the white outer layer of the eye, often during vitreoretinal surgery.

\medterm{Seasonal Affective Disorder (SAD)} Alternate index label for Seasonal Affective Disorder.

\medterm{Secondary Cancer} A secondary cancer is a tumor that has spread from its original site to another part of the body. It is also called a metastasis or metastatic cancer.

\medterm{Sedative} A sedative is a medicine or substance that reduces alertness and may cause relaxation or sleepiness. It can impair coordination and breathing, especially when combined with alcohol or opioids.

\medterm{Selective Serotonin Reuptake Inhibitor (SSRI)} A selective serotonin reuptake inhibitor is an antidepressant that increases serotonin signaling in the brain. Examples include fluoxetine, sertraline, and escitalopram.

\medterm{Self-Diagnosis} Self-diagnosis is identifying a health problem without a professional assessment. It can be inaccurate, so concerning or persistent symptoms should be evaluated by a clinician.

\medterm{Sentinel Node} A sentinel node is the first lymph node or group of nodes to which cancer cells are most likely to spread from a primary tumor. A sentinel lymph-node biopsy helps assess spread.

\medterm{Severe Combined Immunodeficiency (SCID)} A group of inherited disorders in which the immune system cannot make effective defenses against germs. Babies may develop repeated, severe infections, persistent diarrhea, or poor growth; early diagnosis and treatment can be lifesaving.

\medterm{Sexual Health} Sexual health includes physical, emotional, mental, and social well-being related to sexuality. It includes contraception, fertility, consent, relationships, and prevention of sexually transmitted infections.

\medterm{Sexually Transmitted Disease (STD)} A sexually transmitted disease is an infection spread mainly through sexual contact. The term sexually transmitted infection is also used because some infections cause no symptoms.

\medterm{Sheehan's Syndrome} Loss of pituitary hormones caused by damage to the pituitary gland after severe bleeding or dangerously low blood pressure during or after childbirth. It may cause inability to breastfeed, extreme tiredness, low blood pressure, menstrual changes, or other hormone deficiencies.

\medterm{Shilajit} Shilajit is a mineral-rich, resin-like substance used in some traditional medicines. Product purity and safety vary, and contamination with heavy metals is possible.

\medterm{Shisha} Shisha is tobacco or another smoking mixture used in a water pipe. The water does not remove the main toxic substances, and smoking shisha can expose users to nicotine, carbon monoxide, and cancer-causing chemicals.

\medterm{Short-Chain Fatty Acids} Small fatty substances made when bacteria in the large intestine break down dietary fiber. They help nourish the bowel lining and may influence inflammation and metabolism; their effects vary with diet and the person’s gut bacteria.

\medterm{Singing} Singing is producing musical sounds with the voice. It can support communication and social well-being, while excessive or improper voice use may cause vocal strain.


\medterm{Sirtuins} A family of enzymes that change other proteins and help regulate how cells use energy, respond to stress, and switch genes on or off. Their possible roles in aging and disease are still being studied.

\medterm{Sjogren's Syndrome} Older name for Sjögren disease, an autoimmune disease that commonly causes dry eyes and dry mouth.

\medterm{Skin Cells} Skin cells form the layers of the skin and include keratinocytes, melanocytes, immune cells, and sensory cells. They provide protection, pigmentation, sensation, and repair.

\medterm{Skin Cyst} A skin cyst is a closed sac under or within the skin, often filled with fluid, keratin, or other material. Many are harmless, but a painful, infected, or changing cyst should be assessed.

\medterm{Sleep Disorder} Alternate index label for Sleep Disorders.

\medterm{Small Incision Lenticule Extraction Surgery} Small incision lenticule extraction, or SMILE, is a laser eye procedure that removes a small piece of corneal tissue through a tiny incision to correct some cases of nearsightedness and astigmatism.

\medterm{Smoking and Pregnancy} Smoking during pregnancy exposes the fetus to nicotine and other harmful chemicals and increases risks such as poor growth, preterm birth, placental problems, and stillbirth. Stopping at any time is beneficial.

\medterm{Social Media} Social media consists of online platforms where people create, share, and discuss content. Its effects on health vary and may involve social support, sleep disruption, stress, or exposure to misinformation.

\medterm{Social Media Addiction} Social media addiction is a pattern of difficult-to-control use that causes distress or interferes with sleep, work, school, relationships, or daily activities. It is not a universally defined formal diagnosis.

\medterm{Solar Elastosis} Solar elastosis is degeneration and thickening of elastic fibers in skin after long-term ultraviolet exposure. It contributes to leathery, wrinkled, yellowish sun-damaged skin.

\medterm{Solitary Kidney} A solitary kidney means a person has one functioning kidney, because the other kidney is absent, removed, or not working. Many people remain healthy but may need blood-pressure and kidney-function monitoring.

\medterm{Sonohysterography} Sonohysterography is an ultrasound examination in which sterile fluid is placed in the uterus to outline the uterine cavity and assess problems such as polyps or fibroids.

\medterm{Soybeans} Soybeans are legumes that provide protein, fiber, and compounds called isoflavones. They are used in foods such as tofu and soy milk; allergy is possible.


\medterm{Spirulina} Spirulina is a type of cyanobacterium sold as a dietary supplement. Products may provide nutrients, but contamination and medication interactions are possible, and it is not a reliable substitute for a balanced diet.

\medterm{Splinter} A splinter is a small piece of wood, glass, metal, or another material embedded in the skin. It should be removed carefully; medical care may be needed if it is deep, infected, or near an eye or joint.

\medterm{Splinter Hemorrhage} A splinter hemorrhage is a thin line of bleeding under a fingernail or toenail. It may follow injury, but multiple or unexplained lesions can occur with medical conditions and deserve evaluation.

\medterm{Sports Medicine} Medical care focused on preventing, diagnosing, and treating exercise- and sports-related injuries and conditions. It also supports safe training, rehabilitation, performance, and return to activity.
\medterm{Squamous Cell Carcinoma in Situ} Squamous cell carcinoma in situ, also called Bowen disease, is an early skin cancer limited to the outer skin layer. It can become invasive if untreated.

\medterm{Stafne Bone Defect} A Stafne bone defect is a usually harmless, developmental indentation in the lower jaw caused by nearby salivary-gland tissue. It is often found incidentally on imaging.

\medterm{Stem Cells} Stem cells are cells that can self-renew and develop into specialized cell types. They are important in normal tissue repair and in some established treatments, such as blood-forming stem-cell transplantation.

\medterm{Still's Disease} An inflammatory illness that can affect children or adults. It commonly causes repeated high fevers, a fleeting rash, joint pain or swelling, and marked inflammation; it can affect other organs and usually requires specialist treatment.

\medterm{Stomach Pain} Stomach pain is discomfort in the upper abdomen. Causes range from indigestion and infection to ulcers, gallbladder disease, or pancreatitis; severe or persistent pain needs medical assessment.

\medterm{Streptococcus pneumoniae} Streptococcus pneumoniae is a bacterium that can cause pneumonia, meningitis, ear infection, sinusitis, and bloodstream infection. Vaccination helps prevent serious disease.

\medterm{Stress Related} Describes symptoms or health problems that are triggered or worsened by physical or emotional stress. Stress can affect sleep, pain, digestion, mood, heart rate, and blood pressure, but symptoms still require assessment for other causes.

\medterm{Stretch Mark} A stretch mark is a linear scar caused by stretching and changes in the deeper skin, often during pregnancy, growth, or rapid weight change. It is harmless and usually fades over time.

\medterm{Stunted Growth} Stunted growth is height substantially below expected levels for age, often due to long-term undernutrition, chronic illness, infection, or other factors. Children with poor growth should be evaluated.

\medterm{Stutter} Stuttering is a speech-fluency disorder involving repetitions, prolongations, or blocks in speech. Speech-language therapy can improve communication and confidence.

\medterm{Sudden Infant Death Syndrome (SIDS)} Sudden infant death syndrome is the sudden, unexplained death of an infant younger than one year, usually during sleep. Safe-sleep practices reduce risk: place babies on their backs on a firm, flat surface without loose bedding.

\medterm{Sudden Unexpected Death in Epilepsy} An unexpected death in a person with epilepsy that is not caused by drowning, injury, prolonged seizure, or another known condition. It is uncommon; controlling seizures, taking prescribed treatment, and discussing nighttime safety may reduce risk.

\medterm{Superfood} Superfood is a marketing term for a food promoted as especially nutritious. No single food provides all needed nutrients, and health claims should be judged against scientific evidence.

\medterm{Supplements} Dietary supplements contain substances such as vitamins, minerals, herbs, or amino acids. They can have side effects and interactions, so patients should discuss their use with a healthcare professional.

\medterm{Synaesthesia} Synaesthesia is a neurological trait in which stimulation of one sense or mental process automatically produces an additional sensory experience, such as seeing colors when hearing sounds.

\medterm{Synthetic Biology} Synthetic biology combines biology and engineering to design or modify biological systems. It has potential medical uses but also requires careful safety and ethical oversight.


\medterm{T Cell} Abbreviation for T lymphocyte, a type of white blood cell that is produced in the bone marrow and is part of the body\&rsquo;s immune system.

\medterm{T Wave} The electrocardiographic deflection representing ventricular repolarization. Its shape and direction
vary by lead and clinical context. Abnormal T waves may occur with ischemia, electrolyte disorders, ventricular strain,
or other conditions.

\medterm{T-Cell Count} A laboratory measurement of T lymphocytes in blood, often including CD4 cells, used to assess immune function and monitor conditions such as HIV infection or immunosuppression.

\medterm{T-Cell Lymphoma} A cancer of T lymphocytes, a type of white blood cell. It includes several different diseases that may affect lymph nodes, skin, blood, or organs and can cause swollen nodes, fever, weight loss, rash, or other symptoms.
\medterm{T-Cell Lymphomas} Lymphomas arising from mature or immature T lymphocytes. They include several biologically distinct diseases that may
involve lymph nodes, skin, blood, or other organs; behavior and treatment depend on the
subtype.

\medterm{T-Cell Receptor} A clonally distributed receptor on T lymphocytes that recognizes a peptide antigen
presented by a major histocompatibility complex molecule. TCR signaling requires
accessory and co-stimulatory signals for effective T-cell activation.

\medterm{T-Cells (T-Lymphocytes)} A type of lymphocyte that matures in the thymus and is the central effector of
cell-mediated immunity. T-cells express the T-cell receptor (TCR) on their surface,
which recognises peptide antigens presented by major histocompatibility complex (MHC)
molecules on antigen-presenting cells.

\synonyms
T-Lymphocytes

\medterm{T-Lymphocyte} A white blood cell that matures in the thymus and directs cell-based immune responses. Different T cells kill infected cells, coordinate immunity, remember previous threats, or reduce excessive immune activity.
\medterm{T-Wave Alternans} T-wave alternans is a beat-to-beat change in the size or shape of the heart's T wave and may indicate increased risk of dangerous ventricular rhythms.

\medterm{T3} Triiodothyronine, the most biologically active thyroid hormone, produced by deiodination
of T4 in peripheral tissues (liver, kidney) and secreted by the thyroid. Regulates
metabolism (BMR), protein/carbohydrate/lipid metabolism, growth, and development.

\medterm{T3 Test} A blood test measuring triiodothyronine, used with thyroid-stimulating hormone and thyroxine to evaluate thyroid function, particularly suspected hyperthyroidism.

\medterm{T3RU Test} An older test estimating the proportion of thyroid-hormone-binding sites occupied in serum; it is used indirectly to assess thyroid-hormone binding and has largely been replaced by free-thyroxine measurements.

\medterm{T4} Thyroxine, the main hormone secreted by the thyroid gland (thyroxine, 80-90 percent of
secretion), converted to the active T3 in tissues. Bound largely to thyroxine-binding
globulin (TBG); free T4 is the active unbound fraction. Regulates metabolism, growth,
and development.

\medterm{TAB Vaccine} An older vaccine against typhoid fever and paratyphoid A and B. It has largely been replaced by newer typhoid vaccines and does not protect against every cause of fever or every type of paratyphoid infection.
\medterm{Tabes} An older term for tabes dorsalis, a late neurologic complication of untreated syphilis.
\medterm{Tabes Dorsalis} A late manifestation of untreated syphilis in which degeneration of dorsal spinal-cord columns and roots causes sensory ataxia, lightning pains, loss of position sense, areflexia, and bladder or sexual dysfunction.

\medterm{Tablet} A solid dosage form of compressed drug plus excipients (binders, fillers, disintegrants,
lubricants, coatings). Types: immediate-release, enteric-coated (resistant to gastric
acid), sustained/controlled-release, sublingual, buccal, effervescent, chewable, and
orally disintegrating.

\medterm{Tablet Splitting} Dividing a tablet to take a smaller dose; only tablets designed for splitting should be divided, because extended-release, enteric-coated, and some specialized tablets may not work safely when split.

\medterm{TAC Regimen} A chemotherapy combination of docetaxel, doxorubicin, and cyclophosphamide, used for selected breast cancers; it can cause infection risk, nausea, hair loss, heart effects, and other toxicities.

\medterm{Tacalcitol} A vitamin-D-like medicine applied to skin affected by psoriasis or selected scaling disorders. It helps regulate skin-cell growth and inflammation but can irritate skin or raise calcium if overused.
\medterm{Tacedinaline} An investigational medicine that inhibits histone deacetylases, enzymes that influence gene activity. It has been studied for slowing cancer-cell growth but is not an established routine treatment.
\medterm{Tache Noire} Tache noire is a black eschar at the site of a tick bite, classically associated with Mediterranean spotted fever and related rickettsial infections.

\medterm{Tachyarrhythmia} An abnormally fast heart rhythm, which may arise from the atria, ventricles, or the heart's conduction
system. Symptoms can include palpitations, dizziness, chest discomfort, shortness of breath,
or fainting; fainting or severe symptoms require urgent assessment.

\medterm{Tachycardia} A faster-than-expected heart rate that may be a normal response or a sign of disease. See Fast Heart Rate.

\synonyms
Rapid Heartbeat

\medterm{Tachycardia Paroxysmal Atrial (Paroxysmal Supraventricular Tachycardia)} Alternate index label for Paroxysmal Supraventricular Tachycardia.

\medterm{Tachycardia Paroxysmal Supraventricular (Paroxysmal Supraventricular Tachycardia)} Alternate index label for Paroxysmal Supraventricular Tachycardia.

\medterm{Tachycardia-Induced Cardiomyopathy} Weakening and dilation of the heart muscle caused by a persistent or frequently recurring rapid heart rhythm, often improving when the arrhythmia is controlled.

\medterm{Tachykinin} A family of signaling peptides, including substance P, that influence pain, inflammation, smooth-muscle activity, and nerve communication.

\medterm{Tachyphrenia} Tachyphrenia is rapid flow of thoughts or accelerated mental activity, often described in association with mania, stimulants, or other psychiatric states.

\medterm{Tachyphylaxis} A rapid, acute decrease in response to a drug after repeated administration over a short
period, distinct from tolerance which develops more gradually. Tachyphylaxis results
from rapid depletion of neurotransmitter stores, receptor desensitization, or receptor
downregulation.

\medterm{Tachypnea} Abnormally rapid respiratory rate (adult >20 breaths/min). Physiologic (exercise,
anxiety) or pathologic: hypoxia, fever, sepsis, metabolic acidosis (Kussmaul), pulmonary
embolism, pneumonia, heart failure, pain, and compensation for metabolic derangement.

\medterm{Tachypnoea} Abnormally rapid breathing, defined as a respiratory rate >20 breaths per minute in
adults (age-dependent: >60 in neonates, >40 in 1--12 months, >30 in 1--5 years, >20 in
>5 years). Tachypnoea is a non-specific sign of respiratory or metabolic distress.

\medterm{Tachysystole (Uterine Hyperstimulation)} More than 5 contractions in 10 minutes averaged over 30 minutes, or contractions lasting
\(\geq 2\) minutes. Tachysystole with fetal heart rate changes (late decelerations, prolonged
deceleration, loss of variability) is more concerning. Causes: oxytocin (most common),
prostaglandins (misoprostol > dinoprostone), and cocaine.

\synonyms
Uterine Hyperstimulation

\medterm{TACO} Transfusion-associated circulatory overload is breathing difficulty caused by too much fluid entering the circulation during or soon after a blood transfusion. It can cause lung fluid, high blood pressure, and swelling, especially in people with heart or kidney disease or those receiving rapid or large transfusions.

\medterm{Tacrine} An older acetylcholinesterase inhibitor formerly used for Alzheimer disease; liver toxicity and frequent dosing limited its use.

\medterm{Tacrolimus} An immune-suppressing medicine used after transplants and on skin for eczema; it can increase infection risk.
\medterm{Tactile} Relating to the sense of touch, including feeling pressure, texture, vibration, temperature, or pain. Touch signals travel through nerves to the brain; reduced or abnormal touch sensation may result from nerve, spinal-cord, brain, or skin disease.
\medterm{Tactile Agnosia} Tactile agnosia is inability to recognize a familiar object by touch despite intact primary sensation, usually from parietal-cortex dysfunction.

\medterm{Tactile Fremitus} Vibrations transmitted through the lung to the chest wall, felt during palpation. It is
increased in consolidation (pneumonia) due to better sound transmission through solid
tissue. It is decreased in pleural effusion, pneumothorax, and atelectasis due to
impaired transmission.

\medterm{Tactile Sense} The ability to perceive touch, pressure, vibration, texture, and related mechanical stimuli through receptors in the skin and deeper tissues.

\medterm{Tadalafil} Tadalafil is a phosphodiesterase-5 inhibitor used for erectile dysfunction, benign prostatic hyperplasia, and pulmonary hypertension; headache, flushing, hypotension, priapism, and dangerous nitrate interactions are important risks.

\medterm{Taenia} A genus of tapeworms, including T. solium and T. saginata.
\medterm{Taenia Saginata} Taenia saginata is the beef tapeworm acquired from undercooked infected beef; intestinal infection may cause abdominal symptoms or passage of motile segments but does not usually cause human cysticercosis.

\medterm{Taenia Solium} A pork tapeworm that infects people who eat undercooked pork or swallow its eggs. Adult worms live in the intestine, while eggs can cause cysticercosis, including dangerous infection of the brain.
\medterm{Taeniasis} Intestinal infection caused by adult Taenia tapeworms, usually acquired by eating undercooked beef or pork containing larval cysts. It may cause abdominal symptoms or passage of worm segments; swallowing Taenia solium eggs causes cysticercosis in tissues, which is a different infection.
\medterm{Taeniasis (Tapeworm Infection)} Intestinal infection caused by adult Taenia tapeworms, usually acquired by eating undercooked
meat containing larval cysts. Taenia solium and Taenia saginata are important species;
Taenia solium eggs can also cause cysticercosis in tissues.

\synonyms
Tapeworm Infection

\medterm{Tail} A tapering appendage extending from the end of an animal's body. In human anatomy, the term may refer to a tail-like congenital structure or the embryologic caudal end; evaluation depends on the specific finding.

\medterm{Tailbone Trauma} Injury to the coccyx, the small bone at the bottom of the spine, usually after a fall or prolonged pressure. It causes pain when sitting, standing, or passing stool. Most injuries are bruises and improve with time, cushions, and pain relief; persistent or severe pain needs medical assessment.


\medterm{Taint} A taint is contamination or an undesirable influence; in medicine or laboratory work it may refer to contamination that affects a specimen, product, or result.

\medterm{Takayasu Arteritis} A chronic granulomatous vasculitis primarily affecting the aorta and its major branches,
most commonly diagnosed in women under 50 years of age, particularly in Asian
populations.

\medterm{Takayasu Disease} A large-vessel vasculitis that primarily affects the aorta and its major branches. It may cause limb claudication,
weak or unequal pulses, blood-pressure differences, bruits, fever, fatigue, or arterial narrowing and aneurysms.
Imaging and inflammatory assessment guide immunosuppressive and vascular treatment.

\medterm{Takayasu’S Arteritis} A chronic granulomatous vasculitis of the aorta and its major branches, predominantly
affecting young women of Asian descent. Causes arterial stenosis, occlusion, or aneurysm
formation. Presents with limb claudication, absent pulses, bruit, hypertension, and
constitutional symptoms. Treatment: corticosteroids and immunosuppressives.

\medterm{Takotsubo Cardiomyopathy} Transient left ventricular apical ballooning resembling a Japanese octopus trap
(takotsubo), typically triggered by emotional or physical stress. ECG may show ST
elevation mimicking STEMI. Coronary angiography shows no obstructive disease. Most
patients recover normal ventricular function within weeks.

\medterm{Talabostat} Talabostat is an investigational inhibitor of fibroblast activation protein and related enzymes studied for immune and anticancer effects.

\medterm{Talactoferrin Alfa} Talactoferrin alfa is a recombinant lactoferrin-related immunomodulatory protein studied as a treatment or adjunct in cancer and infectious disease.

\medterm{Talaporfin Sodium} Talaporfin sodium is a photosensitizer used or studied in photodynamic therapy, where light activation produces local cytotoxic effects.

\medterm{Talc} A soft mineral; inhalation of talc dust can cause talcosis, and contaminated injected talc can injure vessels and lungs.
\medterm{Talcosis} Pneumoconiosis from talc dust inhalation or IV injection of crushed oral medications.
Causes upper lobe nodular fibrosis. HRCT may show ``galaxy sign'' of calcified nodules.

\medterm{Talcum Powder Poisoning} Harmful effects from swallowing or inhaling talcum powder. Swallowing usually causes mild stomach upset, but inhaling it can injure the lungs and make breathing difficult.

\medterm{Talipes} A deformity of the foot and ankle present at birth or acquired later. It may affect position, flexibility, walking, and shoe fit, and treatment depends on severity, cause, age, and function.
\medterm{Talipes Calcaneovalgus} A foot deformity in which the heel is turned outward and the ankle is dorsiflexed so the foot points upward; flexible newborn cases often improve, while rigid cases need evaluation.

\medterm{Talipes Equinovarus} Clubfoot, a congenital deformity combining ankle equinus, hindfoot varus, forefoot adduction, and cavus; early serial casting and bracing are standard initial treatment.

\medterm{Tall Stature} Tall stature means height substantially above the expected range for age and sex and may be familial, endocrine, genetic, or associated with another condition.

\medterm{Talotrexin} Talotrexin is an investigational antifolate antineoplastic drug designed to inhibit folate-dependent cellular metabolism. An experimental cancer medicine designed to block folate-dependent cell processes needed for uncontrolled tumor growth.
\medterm{Talus} Second largest tarsal bone articulating with tibia/fibula proximally at the ankle.
Approximately 60 percent of its surface is articular cartilage with no muscular
attachments. Transmits entire body weight to the calcaneus. Talar neck fractures may
disrupt blood supply causing avascular necrosis in up to 50 percent.

\medterm{Tamoxifen (Nolvadex)} A medicine, biologic, or pharmacologic term related to tamoxifen (nolvadex); its mechanism, indication, interactions, and adverse effects depend on the specific agent.

\medterm{Tamoxifen Citrate} The citrate salt of tamoxifen, a selective estrogen-receptor modulator used to treat or prevent some breast cancers; blood clots and endometrial cancer are important uncommon risks.

\medterm{Tampon} A material inserted into the vagina to absorb menstrual blood; prolonged use can rarely lead to toxic shock syndrome.
\medterm{Tamponade} Compression of an organ or vessel by accumulated fluid, blood, or gas; cardiac tamponade impairs cardiac filling.
\medterm{Tamsulosin Hydrochloride} The hydrochloride salt of tamsulosin, an alpha-1 blocker that relaxes prostate and bladder-neck muscle to improve urinary symptoms from benign prostatic enlargement.

\medterm{Tandutinib} An experimental medicine studied for some blood disorders and cancers. It blocks signals that help certain cells grow, but it is not a routine treatment and its benefits and risks remain dependent on research findings.

\medterm{Tanespimycin} An investigational derivative of 17-AAG that inhibits heat-shock protein 90 and has been studied as an anticancer drug.

\medterm{Tangier Disease} A rare inherited disorder causing extremely low HDL cholesterol and buildup of cholesterol in tissues. It may cause enlarged yellow-orange tonsils, enlarged liver or spleen, nerve symptoms, or clouding of the cornea; cardiovascular risk varies and specialist testing is needed.

\medterm{Tanner Staging} A scale for assessing sexual maturity during puberty. Stage 1: prepubertal. Stage 2: breast buds (girls), testicular enlargement (boys). Stage 3: breast enlargement, pubic hair. Stage 4: areolar development, more pubic hair.

\medterm{Tannic Acid} A plant-derived polyphenol with astringent and protein-binding properties, used in some industrial or historical preparations; concentrated ingestion can irritate the gastrointestinal tract and affect absorption.

\medterm{Tanomastat} Tanomastat is an experimental medicine that blocks enzymes involved in tissue breakdown and remodeling. It has been studied for limiting tumor spread, but its safety, effectiveness, and clinical role remain unestablished.

\medterm{Tantalum} A corrosion-resistant metal used in some surgical implants and medical devices; exposure risk depends on the form, dose, and route.

\medterm{Tantulum} A misspelling of tantalum, a metal used in selected implants and medical devices.
\medterm{Tap, Joint (Aspiration)} A needle procedure that removes fluid from a joint for testing or symptom relief, helping assess infection, bleeding, crystals, or inflammation.

\medterm{Tapentadol} Tapentadol is a centrally acting analgesic with mu-opioid agonist and norepinephrine-reuptake-inhibitory effects, used for moderate-to-severe pain; respiratory depression, dependence, seizures, and serotonin or sedative interactions are possible.

\medterm{Tapeworm} A parasitic flatworm with a segmented body that may live in the intestine or form larval cysts in tissues.
People acquire different species by eating infected meat or fish or by ingesting eggs from
contaminated sources. Symptoms and complications depend on the species and site.

\medterm{Tapeworm Infection} Infection caused by parasitic tapeworms. Taenia species may live as adult worms in the intestine, while some tapeworm larvae can form cysts in tissues.

\medterm{Tapeworm Infection - Beef Or Pork} Intestinal infection with Taenia saginata or Taenia solium acquired from undercooked beef or pork; T. solium eggs can also cause cysticercosis in tissues.

\medterm{Tapeworm Infection - Hymenolepis} Intestinal infection with Hymenolepis nana or H. diminuta, acquired by ingesting eggs or infected arthropods and causing variable gastrointestinal symptoms.

\medterm{Tapotement} A massage technique using rhythmic tapping, patting, or percussion. It should be avoided over injured areas or in people at risk of bleeding unless a trained clinician recommends it.
\medterm{Tapping} Repeated light striking or percussion used during an examination, treatment, or test. Tapping may help assess reflexes, pain, body sounds, or fluid and may also be used in chest physiotherapy or other treatments.
\medterm{TAPSE} Alternate name; see \textit{Tricuspid Annular Plane Systolic Excursion}.

\medterm{TAPVR} Alternate index label; see \textit{Total anomalous pulmonary venous return}.

\medterm{Taq Polymerase} A heat-stable DNA-copying enzyme from Thermus aquaticus, widely used in polymerase chain reaction testing and research.

\medterm{Tar Remover Poisoning} Poisoning from swallowing, inhaling, or skin contact with tar-removing chemicals. Effects vary by product and may include irritation, dizziness, vomiting, or breathing problems; contact poison-control services promptly.

\medterm{Tarantism} Tarantism is a historical term for a dancing or behavioral syndrome once attributed to spider bite and now understood as a cultural or neurologic phenomenon rather than a specific disease.

\medterm{Tarantula Spider Bite} A bite from a tarantula that usually causes local pain, redness, and swelling. Some species also release irritating hairs that can affect the skin or eyes.

\medterm{Tardive Dyskinesia} Involuntary, repetitive, purposeless movements caused by prolonged use of antipsychotic
medications, typically developing after months to years of exposure. Movements include
choreoathetoid movements of the tongue, face, and extremities (lip smacking, tongue
protrusion, grimacing, choreiform movements of limbs).

\medterm{Tarenflurbil} Tarenflurbil is an investigational derivative of flurbiprofen studied for effects on amyloid processing and Alzheimer disease.

\medterm{Target Cell} A cell specifically affected by or responding to a hormone, drug, toxin, antigen, or other signal; target cells may also describe abnormal erythrocytes.
\medterm{Target Organ} An organ especially affected by a disease, hormone, toxin, medicine, or other exposure. Effects may occur because the organ has specific receptors, receives the substance, or is especially vulnerable to damage.
\medterm{Target Population} The defined group of people for whom a medical intervention, screening program, or research study is intended. It is characterized by criteria such as age, diagnosis, risk, or other relevant features.

\medterm{Target Tissue} A tissue containing receptors or other molecular targets on which a hormone, drug, toxin, or signaling molecule acts.

\medterm{Target-Controlled Infusion} A computerized infusion pump system that maintains a predetermined drug concentration in
plasma or effect site using pharmacokinetic modeling. TCI systems use three-compartment
pharmacokinetic models to calculate infusion rates that achieve and maintain target
concentration.

\medterm{Targeted Radiation Therapy} Radiation directed precisely at a tumor or other abnormal tissue to damage its DNA while limiting exposure to nearby healthy tissue. It may be delivered from outside the body or through an implanted source.
\medterm{Targeted Temperature Management} Controlled therapeutic hypothermia for comatose survivors of cardiac arrest. Target:
32-36 degrees C for 24-72 hours followed by gradual rewarming. Mechanism: reduces
metabolic demand, limits reperfusion injury, and reduces free radical generation. TTM2
trial showed 33 degrees C was not superior to 36 degrees C.

\medterm{Targeted Therapies For Cancer} Cancer medicines designed to act on specific features of cancer cells or their surroundings. They may be more selective than traditional chemotherapy, but can still cause important side effects.

\medterm{Targeted Therapy} Treatment that acts on a specific change or feature helping cancer cells grow, survive, or spread. Testing the tumor may identify a target and predict benefit. These treatments can be more selective than traditional chemotherapy but may still cause serious side effects and resistance.


\medterm{Tariquidar} Tariquidar is an investigational inhibitor of the P-glycoprotein drug transporter studied to reverse multidrug resistance and alter medicine distribution.

\medterm{Tarlov Cyst} A cerebrospinal-fluid-filled perineural cyst, usually located near the dorsal root
ganglia in the sacral spine. Many are incidental, but larger or symptomatic cysts can
cause sacral pain, radiculopathy, sensory change, bowel or bladder dysfunction, or
sexual dysfunction.

\medterm{Tars} A thick black substance from organic material; tobacco tar contains chemicals that can damage lungs and cause cancer.
\medterm{Tarsal} Relating to the tarsus, the group of bones forming the ankle and hindfoot. Tarsal bones help support body weight and allow foot movement.
\medterm{Tarsal Bones} The seven bones of the hindfoot and midfoot that connect the ankle to the metatarsals and form the arches and load-bearing framework of the foot.

\synonyms
Tarsals

\medterm{Tarsal Coalition} An abnormal connection between two bones in the rear or middle foot, usually present from birth. It can
cause a stiff, painful flatfoot, recurrent ankle sprains, and limited foot motion, often becoming
more noticeable during adolescence.

\medterm{Tarsal Joints} The articulations among the tarsal bones and between the tarsals and metatarsals, providing stability and controlled movement for walking and adapting the foot to surfaces.

\medterm{Tarsal Tunnel Syndrome} Pressure on the tibial nerve as it passes through a narrow space near the inner ankle. It can cause burning pain, tingling, or numbness in the sole and may result from injury, swelling, diabetes, or a mass. Treatment depends on the cause.

\medterm{Tarsorrhaphy} A procedure that partially or completely closes the eyelids to protect the cornea or reduce exposure in conditions such as severe dry eye or facial nerve palsy.

\medterm{Tarsus} The collection of seven bones forming the posterior part of the foot, between the tibia/fibula and the metatarsals. The group of seven bones forming the hindfoot and midfoot.
\medterm{Tartar} A hard deposit of mineralized dental plaque, also called calculus.
\medterm{Tarui Disease} A rare inherited muscle-energy disorder caused by deficiency of an enzyme needed to break down sugar for energy. Exercise can trigger early fatigue, cramps, muscle breakdown, and dark urine; some people also have chronic anemia or high uric acid.

\synonyms
Glycogen Storage Disease Type VII (GSD VII); Muscle Phosphofructokinase Deficiency

\medterm{Tassinari Syndrome} An older name associated with epileptic encephalopathy featuring frequent abnormal electrical activity during sleep, seizures, and possible effects on learning or development. The term is now usually described more specifically as electrical status epilepticus in sleep.

\medterm{Taste} The sensation produced by the interaction of chemical substances with taste receptor
cells (TRCs) on the tongue, palate, pharynx, epiglottis, and upper esophagus. Taste is
one of the five primary senses (gustation).

\medterm{Taste - Impaired} Impaired taste is reduced, distorted, or absent ability to taste and may result from infection, medicines, oral disease, nerve injury, or changes in smell.

\medterm{Taste Bud} A taste bud is a cluster of specialized sensory cells in papillae of the tongue and some oral tissues that detects sweet, sour, salty, bitter, and umami stimuli and transmits signals through cranial nerves.

\medterm{Taste Disorder} Abnormal reduction, loss, distortion, or unpleasant alteration of taste, caused by infection, medications, oral disease, nerve injury, neurologic disease, or other factors.



\medterm{Tau PET} Positron-emission tomography using a tracer that binds to abnormal tau protein deposits. It is used mainly in research and selected clinical evaluations of neurodegenerative disease.

\medterm{Taurodontism} A tooth-development difference with an enlarged pulp chamber and shortened roots, usually affecting molars. It may occur alone or with genetic conditions and can complicate dental treatment or root-canal procedures.

\medterm{Taussig-Bing Anomaly} A rare birth defect in which both major arteries leave the right lower heart chamber and a hole lies between the lower chambers below the pulmonary artery. It can cause severe low oxygen and heart failure in infancy and usually requires early specialist surgery.

\synonyms
Taussig-Bing Syndrome; DORV with Subpulmonic VSD

\medterm{TAVR} Alternate name; see \textit{Transcatheter Aortic Valve Replacement}.

\medterm{Taxanes} Chemotherapy medicines, including paclitaxel and docetaxel, that disrupt cell division by stabilizing tiny internal structures. They treat several cancers but may cause low blood counts, nerve symptoms, hair loss, allergic reactions, or fluid retention.
\medterm{Taxis} Directed movement toward or away from a stimulus; in surgery it may also mean orderly repositioning.

\medterm{Tay-Sachs Disease} A rare inherited disorder in which a fatty substance builds up in nerve cells because of an enzyme deficiency. The infant form causes loss of skills, weakness, an exaggerated startle, and seizures; later forms progress more slowly. Supportive care and genetic counseling are important.

\medterm{Tazarotene} A topical retinoid used for acne, plaque psoriasis, and some sun-damaged skin. It can irritate skin and increases sun sensitivity; pregnancy precautions are important.
\medterm{TB} Alternate index label; see \textit{Tuberculosis}.

\medterm{TB Test} A tuberculosis test looks for infection with Mycobacterium tuberculosis using a skin test, blood test, or examination of respiratory specimens.

\medterm{TBG Blood Test} A blood test measuring thyroxine-binding globulin, the principal thyroid-hormone transport protein; it helps interpret abnormal total thyroid-hormone levels when binding-protein changes are suspected.

\medterm{TBI} Alternate index label; see \textit{Traumatic brain injury}.

\medterm{TCA Cycle} The tricarboxylic acid cycle, also called the Krebs cycle, is a series of reactions inside mitochondria that breaks down products of carbohydrates, fats, and proteins. It produces carbon dioxide and energy-carrying molecules used to make cellular energy.

\medterm{TCAs (Tricyclic Antidepressants)} Antidepressants that block reuptake of serotonin and norepinephrine, with additional
anticholinergic, antihistaminic, and alpha-adrenergic blocking effects. Examples include
amitriptyline, nortriptyline, imipramine, and desipramine. They are effective for
depression, neuropathic pain, and migraine prophylaxis.

\synonyms
Tricyclic Antidepressants

\medterm{Td} Td is a vaccine protecting against tetanus and diphtheria, generally used as a booster in adolescents and adults according to local schedules.


\medterm{Td Immunization} Vaccination against tetanus and reduced-dose diphtheria, used for routine boosters, catch-up immunization, and wound-related protection according to age and vaccination history.


\medterm{Tdap Vaccination In Pregnancy} Administration of tetanus, diphtheria, and acellular pertussis (Tdap) vaccine during
each pregnancy, optimally at 27-36 weeks. The goal: boost maternal antibodies to protect
the newborn from pertussis (whooping cough) before the infant can be vaccinated at 2
months.

\medterm{Tear} A split or rupture in tissue, such as a muscle, tendon, ligament, meniscus, skin, or organ, caused by trauma, overuse, degeneration, or pressure.

\medterm{Tearing} Excessive production or inadequate drainage of tears, commonly caused by irritation, allergy, dry-eye reflex tearing, infection, eyelid malposition, or lacrimal-duct obstruction.

\medterm{Tears} Fluid made mainly by the lacrimal glands that lubricates, nourishes, cleans, and protects the eye's surface. Tears contain water, salts, oils, mucus, and immune substances that help maintain clear, comfortable vision.
\medterm{Tears, Decreased Production} Alternate index label; see \textit{Dry eyes}.

\medterm{Technetium-99} A radioactive isotope related to technetium-99m; technetium isotopes are used in nuclear medicine and radioactive decay studies.
\medterm{Technetium-99M} Most widely used radionuclide in nuclear medicine. Half-life: 6 hours. Emits 140 keV
gamma rays. Produced from molybdenum-99/technetium-99m generator. Used with various
ligands: MDP (bone), sestamibi (myocardium), DTPA (renal), nanocolloid (lymphatic),
HMPAO (brain perfusion).

\medterm{Technique} The method, equipment, and sequence of steps used to perform a medical examination, procedure, or imaging study. Correct technique improves safety, accuracy, and reproducibility.
\medterm{Tedizolid Pharmacokinetics} The study of how tedizolid is absorbed, distributed, metabolized, and eliminated.
Tedizolid is available orally and intravenously and has pharmacokinetic properties that support once-daily treatment;
its clinical use depends on the infection and patient factors.

\medterm{TEE} TEE, or transesophageal echocardiography, uses an ultrasound probe in the esophagus to obtain detailed images of heart valves, chambers, clots, and nearby vessels.

\medterm{Teen Depression} A depressive disorder in an adolescent that causes persistent low mood or loss of interest and impairs daily functioning.

\synonyms
Depression, Teen

\medterm{Teenage Pregnancy} Pregnancy in adolescents (usually 10-19 years), associated with higher maternal and
fetal risks: preterm birth, low birth weight, preeclampsia, anemia, and increased
neonatal mortality; also socioeconomic and educational consequences. Causes: unprotected
intercourse, lack of contraception/education, poverty, and peer/family factors.

\medterm{Teenagers And Drugs} Drug use during adolescence can affect brain development, school and relationships, safety, and mental health; early, nonjudgmental screening and support can reduce harm.

\medterm{Teenagers And Sleep} Teenagers commonly need about eight to ten hours of sleep, but schedules, screens, caffeine, mood, school demands, and sleep disorders can reduce sleep quality.

\medterm{TEER} Alternate name; see \textit{Transcatheter Edge-to-Edge Repair}.

\medterm{Teeth} Hard, mineralized structures in the mouth used for cutting, tearing, and grinding food. Hard mineralized structures in the mouth that cut, tear, and grind food, supported by gums and jawbones.
\medterm{Teeth Grinding} Alternate index label; see \textit{Teeth grinding (bruxism)}.

\medterm{Teeth Grinding (Bruxism)} Alternate index label for Bruxism.

\medterm{Teeth, Disorders Of} Conditions affecting teeth, including decay, developmental differences, infection, injury, poor alignment, wear, and tumors. Problems may cause pain, sensitivity, swelling, difficulty chewing, altered appearance, or tooth loss.
\medterm{Teeth-Grinding} Alternate index label; see Teeth grinding (bruxism). Bruxism, repetitive clenching or grinding of teeth, often during sleep and capable of causing tooth wear and jaw pain.
\medterm{Teething} Teething is eruption of an infant’s primary teeth through the gums and may cause gum discomfort, drooling, and chewing; high fever or severe illness suggests another cause.

\medterm{Teichopsia} A visual disturbance involving shimmering zigzags, flashing patterns, or fortress-like shapes. It commonly occurs as a migraine aura but can have other causes, especially when new, persistent, or associated with vision loss.
\medterm{Teicoplanin} A glycopeptide antibiotic used for selected serious infections caused by susceptible Gram-positive bacteria, including some resistant staphylococci. Kidney function and blood levels may need monitoring.
\medterm{Telangiectasia} A visible, permanently dilated small blood vessel near the skin or mucosal surface, occurring alone or with conditions such as rosacea, liver disease, or inherited vascular syndromes.

\medterm{Telangiectasis} Widening of small blood vessels near the skin or a moist body lining, producing visible red or purple lines or patches. Causes include sun exposure, rosacea, hormones, liver disease, and inherited conditions.
\medterm{Tele Radiology (Teleradiology)} The remote review of medical images, such as X-rays, CT scans, or MRI scans, by a qualified radiologist who is not at the imaging site. It can provide urgent or specialist interpretation when local expertise is unavailable and must protect patient information.

\medterm{Tele-Rehabilitation} Remote rehabilitation via technology. Video consultations, remote monitoring, home exercise programs. Increases access, especially rural.
\medterm{Telehealth} Delivery of healthcare through secure electronic communication or remote monitoring, including video visits, telephone care, messaging, and transmission of clinical data.

\medterm{Telemedicine} Delivery of medical care at a distance using information and communication technology. It may include video or telephone consultations, secure messages, remote diagnosis, treatment, and review of clinical information.

\medterm{Telogen Effluvium} A common, non-scarring form of diffuse hair shedding that occurs when more hairs than usual enter the resting (telogen) phase. Shedding often begins 2--4 months after illness, fever, surgery, childbirth, major stress, rapid weight loss, nutritional deficiency, or certain medicines. It is usually temporary, with regrowth after the trigger is corrected, although recovery may take several months.

\medterm{Telomere} A protective region at the end of each chromosome. It helps keep genetic material stable when cells divide and usually becomes shorter with repeated cell division and aging.
\medterm{Temazepam} A benzodiazepine hypnotic used short term for insomnia; sedation, dependence, and respiratory depression are risks.
\medterm{Temperature} A measurement of body heat. The result varies with the measurement site, time, activity, age, and illness.

\medterm{Temperature (Body Temperature)} The balance between heat production (metabolism, muscle activity) and heat loss
(radiation, conduction, convection, evaporation). Normal core body temperature:
36.5--37.5°C (97.7--99.5°F), regulated by the hypothalamus (preoptic anterior
hypothalamus—thermoregulatory center).

\synonyms
Body Temperature

\medterm{Temperature Measurement} Measurement of body or environmental temperature using a thermometer or sensor at a specified site, interpreted with method, age, timing, and clinical context.

\medterm{Temperature Scales} Systems for expressing temperature, including Celsius, Fahrenheit, and Kelvin. Clinical temperatures are commonly reported in Celsius or Fahrenheit.
\medterm{Temple} The region of the skull at the side of the head between the forehead and ear.
\medterm{Temple Syndrome} A genomic-imprinting disorder involving chromosome 14q32 that causes prenatal and
postnatal growth restriction, hypotonia, developmental delay, feeding difficulty, short
stature, and abnormal pubertal development.

\medterm{Temporal} Relating either to time or to the temple area on the side of the skull, including the temporal lobe.
\medterm{Temporal Arteritis} A vasculitic disorder affecting medium and large arteries, particularly the temporal
artery, predominantly occurring in adults over 50 years of age.

\medterm{Temporal Artery} An artery running along the side of the scalp, supplying nearby tissues; inflammation can cause temporal arteritis.
\medterm{Temporal Bones} Paired bones forming the cranial floor and lateral wall, with squamous, petrous,
mastoid, tympanic, and styloid parts. Contain the external auditory canal, middle and
inner ear, and foramina: carotid canal (internal carotid), jugular foramen (CN IX-XI,
IJV), and foramen lacerum.

\medterm{Temporal Lobe} One of the four major lobes of the cerebral hemisphere, located beneath the lateral
fissure (Sylvian fissure) on the lateral surface of the brain, extending medially to the
parahippocampal gyrus. The temporal lobe is involved in auditory processing, language
comprehension, memory, and emotion.

\medterm{Temporal Lobe Epilepsy} The most common form of focal (partial) epilepsy, arising from the medial temporal lobe
structures (hippocampus, amygdala, parahippocampal gyrus).

\medterm{Temporal Lobe Seizure} A seizure that begins in the temporal lobe of the brain and may cause altered awareness, unusual sensations, or repetitive movements.

\synonyms
Seizure, Temporal Lobe

\medterm{Temporary Loss Of Consciousness} Temporary loss of consciousness is a brief episode of unresponsiveness with loss of postural tone, commonly from syncope but sometimes from seizure or another disorder.

\medterm{Temporomandibular Disorders} Alternate index label; see \textit{TMJ disorders}.

\medterm{Temporomandibular Joint} Synovial ginglymoarthrodial joint between mandibular condyle and temporal bone's
mandibular fossa, separated by an articular disc. Only bilateral body joint, allowing
hinge and sliding movements. Disc divides into superior/inferior compartments.
Temporomandibular disorders cause pain, clicking, limited movement, common with bruxism.

\medterm{Temporomandibular Joint Disorder} Pain or impaired movement involving the jaw joint or chewing muscles. It may cause jaw pain, clicking, stiffness, headache, or difficulty opening the mouth. Clenching, grinding, injury, arthritis, and muscle tension can contribute; treatment usually begins with simple pain relief and jaw-rest measures.


\medterm{Temporomandibular Joint Replacement} Surgery that removes a severely damaged jaw joint and replaces it with an artificial joint. It may be considered for ankylosis, advanced joint destruction, or persistent severe symptoms after other treatments have failed.
\medterm{Temporomandibular Joint Syndrome} Pain or dysfunction involving the jaw joint and chewing muscles, often causing jaw pain, clicking, restricted
opening, headache, or tooth-related symptoms. Causes include clenching, trauma, arthritis, and stress-related muscle
tension. Persistent locking, severe pain, or inability to eat requires dental or medical evaluation.

\medterm{TEN} Alternate index label; see \textit{Toxic epidermal necrolysis}.

\medterm{Tenascin} A group of proteins found around cells that help organize tissues and influence growth, healing, and inflammation. Some forms are increased around tumors or damaged tissue.

\medterm{Tender} Painful or sensitive when touched or pressed, a finding that may indicate inflammation, infection, trauma, ischemia, or another underlying process.

\medterm{Tenderness} Pain or discomfort caused by pressing or touching a body part. It may indicate inflammation, injury, infection, or another local or systemic problem.
\medterm{Tendinitis} Tendinitis is painful inflammation or degeneration of a tendon, often from repetitive loading, causing localized tenderness and pain with movement.

\medterm{Tendinitis Rotator Cuff (Rotator Cuff)} Alternate index label for Rotator Cuff.

\medterm{Tendinitis, Achilles} Alternate index label; see \textit{Achilles tendinitis}.

\medterm{Tendinitis, Patellar} Alternate index label; see \textit{Patellar tendinitis}.

\medterm{Tendinopathy} Pain and abnormal structure or function of a tendon, often related to repeated loading, aging, or injury. It can cause tenderness and weakness; treatment usually combines activity adjustment with rehabilitation.
\medterm{Tendon} Dense regular connective tissue connecting muscle to bone, composed primarily of type I
collagen fibers arranged in parallel bundles. Transmits muscle force to bone to produce
movement. Poor vascularity (slower healing). Enclosed in tendon sheath (synovial) at
friction points.

\medterm{Tendon Injury} Damage to a tendon ranging from microscopic overuse injury to partial or complete rupture, causing localized pain, weakness, swelling, or loss of function.

\medterm{Tendon Repair} Surgery to reconnect a torn tendon to itself or to bone so that movement can be restored. It is followed by protected movement and rehabilitation; re-tearing, stiffness, weakness, infection, and nerve or vessel injury can occur.

\medterm{Tendon Sheath} A thin lubricated covering around some tendons that reduces friction as the tendon moves; inflammation is called tenosynovitis.

\medterm{Tendon Transfer} Surgical relocation of a functioning tendon to restore movement or balance when the
original muscle-tendon unit is absent, paralyzed, or irreparable. Successful transfer
requires an adequate donor, appropriate excursion, and a stable recipient joint.

\medterm{Tendonitis} Tendinopathy: inflammation or degeneration of a tendon from overuse, repetitive strain,
or trauma, causing localized pain, swelling, and tenderness with activity. Common sites:
rotator cuff (supraspinatus), lateral elbow (tennis elbow), medial elbow (golfer's
elbow), Achilles, patellar (jumper's knee), biceps, and De Quervain.

\medterm{Tenecteplase} A modified tissue plasminogen activator (TNK-tPA) thrombolytic with longer half-life,
higher fibrin specificity, and resistance to PAI-1. Used for acute ischemic stroke and
STEMI. Given as a single IV bolus (weight-based), simplifying administration. Adverse
effects: hemorrhage (intracranial—major concern), allergic reactions.

\medterm{Tenericutes} A former bacterial group characterized by organisms that lack a typical cell wall, including Mycoplasma and related bacteria.

\medterm{Tenesmus} A distressing sensation of needing to empty the bowel or bladder despite little output, commonly caused by rectal inflammation, infection, tumor, constipation, or urinary irritation.

\medterm{Tenia} Tenia is an older or variant spelling related to Taenia, a genus of tapeworms that can infect humans or animals.


\medterm{Teniposide} Teniposide is a podophyllotoxin-derived antineoplastic drug that inhibits topoisomerase II and is used in selected leukemia or lymphoma regimens.

\medterm{Tennis Elbow} Lateral epicondylitis, a tendinopathy of the common extensor tendon at the lateral
epicondyle (extensor carpi radialis brevis). Caused by repetitive wrist extension and
overuse, not limited to tennis. Presents with lateral elbow pain worsened by wrist
extension and grip. Tenderness over the lateral epicondyle.

\synonyms
Epicondylitis, Lateral
Lateral Elbow Tendinopathy
Lateral Epicondylitis
Lateral Epicondylosis

\medterm{Tennis Elbow (Lateral Epicondylitis)} Alternate index label for Lateral Epicondylitis.

\medterm{Tennis Elbow Surgery} An operation to remove or repair damaged tendon tissue at the outer elbow when persistent tennis-elbow pain has not improved with nonsurgical care. Recovery requires rehabilitation, and pain, stiffness, nerve injury, or continued symptoms are possible.


\medterm{Tenodesis} Surgical fixation of a tendon to bone or another tendon to restore function, stabilize a joint, or treat tendon imbalance.

\medterm{Tenofovir} Tenofovir is a nucleotide reverse-transcriptase inhibitor used with other medicines for HIV and as treatment or prevention for hepatitis B; renal and bone toxicity and formulation-specific monitoring are important.

\medterm{Tenosynovitis} Tenosynovitis is inflammation of the sheath surrounding a tendon, producing pain, swelling, and restricted gliding, commonly in the hand or wrist.

\medterm{Tenosynovitis, De Quervain'S} Alternate index label; see \textit{De Quervain tenosynovitis}.

\medterm{Tenotomy} Surgery that cuts or releases a tendon to reduce tightness, correct a deformity, relieve pressure, or change movement. It may be performed through a small opening or a larger incision; weakness, bleeding, infection, nerve injury, or recurrence can occur.
\medterm{Tenovaginitis} Inflammation of a tendon and its surrounding sheath, also called tenosynovitis.
\medterm{Tensile Strength} The ability of a tissue or material to resist being pulled apart; it is important in wound healing and the integrity of sutures, tendons, and other connective tissues.

\medterm{Tensilon Test} An older pharmacologic test using edrophonium to look for brief improvement in muscle weakness in suspected myasthenia gravis; it is rarely used because safer tests are available.

\medterm{Tension} A pulling force or state of being stretched; in medicine it may describe tissue tension, blood-vessel wall tension, muscle tightness, or psychological strain depending on context.

\medterm{Tension Headache} The most common headache type, characterized by bilateral, pressing/tightening head pain
of mild-moderate intensity (like a band around the head), lasting 30 min to days,
without nausea or vomiting; photophobia or phonophobia may be present but not both.

\medterm{Tension Pneumothorax} Life-threatening pneumothorax with mediastinal shift. Emergency needle decompression
(2nd intercostal space, midclavicular line) followed by chest tube. Signs: respiratory
distress, hypotension, tracheal deviation, distended neck veins, unilateral absent
breath sounds.

\medterm{Tension-Type Headache} A usually bilateral, pressing or tightening headache of mild-to-moderate intensity without prominent nausea, often associated with stress, muscle tenderness, or poor sleep.

\medterm{Tentorium} A strong fold of the brain’s protective covering separating the cerebellum from the back parts of the cerebrum. It supports brain structures and contains an opening through which the brainstem passes.
\medterm{Tenuazonic Acid} A mycotoxin produced by some fungi that contaminates certain plants and foods and can inhibit protein synthesis; toxicity depends on exposure and the producing organism.

\medterm{Teprotumumab} A monoclonal antibody targeting the insulin-like growth factor 1 receptor and used to treat thyroid eye disease in selected patients.

\medterm{Ter In Die} Ter in die means three times a day; explicit plain-language timing is safer than relying on Latin abbreviations.

\medterm{Terameprocol} An experimental anticancer compound studied for effects on gene activity and cell-survival pathways; it is not routine treatment.

\medterm{Teratocarcinoma} A teratocarcinoma is a malignant germ-cell tumor containing tissues from more than one embryonic layer, often arising in the testis or ovary.

\medterm{Teratogen} A medicine, chemical, infection, radiation exposure, or other influence that can harm an embryo or fetus and cause birth defects, growth problems, abnormal development, or pregnancy loss. Risk depends on the exposure, amount, timing, and individual factors; exposure should be discussed promptly with a clinician.
\medterm{Teratogenesis} The development of a birth defect, impaired growth, abnormal function, or pregnancy loss after harmful exposure during pregnancy. Risk depends on the substance or infection, amount, timing, route, and genetic factors; exposure does not automatically mean that harm will occur.
\medterm{Teratogenic} Capable of causing congenital malformations, growth impairment, functional deficits, or pregnancy loss when an embryo or fetus is exposed during a vulnerable developmental period.

\medterm{Teratogenic Drugs} Medicines that can cause birth defects, pregnancy loss, or impaired fetal development; risk depends on the drug, dose, timing, and pregnancy.

\medterm{Teratogenicity} The ability of a drug, chemical, or infectious agent to cause birth defects or
developmental abnormalities when exposure occurs during pregnancy. The risk of
teratogenic effects depends on the timing of exposure, with the embryonic period from 3
to 8 weeks after fertilization being the most critical for organogenesis.

\medterm{Teratoma} A type of germ cell tumor composed of tissues derived from two or three germ layers
(ectoderm, endoderm, mesoderm), reflecting pluripotent differentiation.

\medterm{Teratozoospermia} Teratozoospermia is a high proportion of abnormally shaped sperm and may contribute to reduced fertility, though semen results must be interpreted as a whole.

\medterm{Terazosin Hydrochloride} The hydrochloride salt of terazosin, an alpha-1 blocker used for urinary symptoms from benign prostatic enlargement and sometimes for hypertension; dizziness and first-dose fainting can occur.

\medterm{Terbinafine} An allylamine antifungal agent that inhibits squalene epoxidase, an early enzyme in the
ergosterol biosynthesis pathway, causing accumulation of squalene and depletion of
ergosterol in the fungal cell membrane. The accumulation of squalene is fungicidal,
while ergosterol depletion disrupts membrane function.

\medterm{Terbufos} Terbufos is a highly toxic organophosphate pesticide that inhibits acetylcholinesterase and can cause severe cholinergic poisoning.

\medterm{Terbutaline} A beta-2 agonist used as a bronchodilator; it has limited and specialist use as a tocolytic.
\medterm{Terfenadine} An older antihistamine withdrawn or restricted in many places because interactions and excessive blood levels can cause dangerous heart-rhythm abnormalities.

\medterm{Teriparatide For Osteoporosis} Recombinant human parathyroid hormone fragments one through thirty-four that stimulate
osteoblast-mediated bone formation. Indicated for severe osteoporosis with fractures or
very low T-scores and for patients who have failed or are intolerant to bisphosphonates.
Daily subcutaneous injection for up to two years.

\medterm{Terlipressin} Terlipressin is a vasopressin analogue that constricts splanchnic vessels and is used in selected cases of variceal bleeding or hepatorenal syndrome.

\medterm{Term} The end of pregnancy at 37 completed weeks of gestation; births before this are classified as preterm.

\medterm{Term Birth} Term birth occurs at or after 37 completed weeks of pregnancy; early term, full term, and late or post-term categories describe more specific gestational ranges.

\medterm{Terminal Care} Care focused on comfort, symptom relief, dignity, and support for a person with an advanced illness who is approaching the end of life.

\medterm{Terminal Delirium} Delirium occurring near the end of life, often caused by advanced organ failure,
medication effects, metabolic disturbance, infection, or reduced cerebral perfusion.
Management prioritizes comfort, reversible contributors when consistent with goals, and
support for the patient and family.

\medterm{Terminal Ileitis} Inflammation of the distal ileum, commonly associated with Crohn disease but also caused by infection, ischemia, medication, or other inflammatory conditions and producing pain, diarrhea, or bleeding.

\medterm{Terminal Illness} An advanced disease that is expected to lead to death, although the timing and response to available treatment may vary.

\medterm{Terminal Insomnia} Waking too early and being unable to return to sleep, a pattern that can occur in depression, circadian-rhythm disorders, aging, or other insomnia conditions.

\medterm{Terms Medical Abbreviations (Common Medical Abbreviations And Terms)} Alternate index label for Common Medical Abbreviations and Terms.

\medterm{Terpene} A terpene is a hydrocarbon built from isoprene units and found in plant oils, fragrances, medicines, and other natural products.

\medterm{Terry'S Nails} A nail change in which most of the nail plate appears white with a narrow pink or brown distal band; it may be associated with liver disease, heart failure, diabetes, or ageing.

\medterm{Tertian Fever} A fever pattern recurring approximately every 48 hours, classically associated with Plasmodium vivax or P. falciparum malaria.
\medterm{Tertiary Hyperparathyroidism} Autonomous overproduction of parathyroid hormone after long-standing secondary hyperparathyroidism, usually in chronic kidney disease, causing high calcium and high PTH.

\medterm{Tertiary Prevention} Care that reduces complications, disability, or suffering from an established disease. Examples include rehabilitation, long-term disease management, preventing relapses, controlling symptoms, restoring function, and providing palliative care.

\medterm{Tertiary Protein} The three-dimensional folding of a single protein chain, created by interactions among amino-acid side chains and determining the protein’s functional shape.

\medterm{Tesetaxel} Tesetaxel is an investigational oral taxane that disrupts microtubules and has been studied as an anticancer treatment.

\medterm{Test Meal} A standardized meal given before measuring a physiologic response, such as gastric acid or glucose handling.
\medterm{Test Strips} Narrow strips containing chemicals that react with a sample to measure or detect a substance. Common examples test urine for blood or infection, or test blood for glucose or ketones. Results can be affected by storage, timing, contamination, and incorrect use.

\medterm{Test-Tube} A small glass or plastic laboratory container used to hold, mix, heat, or test samples. Test tubes may contain blood, urine, chemicals, cells, cultures, or other materials during medical investigations.
\medterm{Test-Tube Baby} A lay term for a baby conceived through in-vitro fertilization.
\medterm{Testes} The paired male reproductive organs in the scrotum. They produce sperm and testosterone and may be affected by infection, injury, torsion, tumors, or disorders of development.
\medterm{Testicle} One of the male reproductive glands in the scrotum. It produces sperm and testosterone; sudden severe testicular pain is an emergency because it may indicate torsion.
\medterm{Testicle Lump} A localized mass or enlargement in or around a testis caused by cyst, infection, inflammation, torsion, trauma, or tumor; a new intratesticular mass requires prompt evaluation.

\medterm{Testicle Pain} Pain in one or both testes or the surrounding structures, with urgent causes including torsion, infection, trauma, and incarcerated hernia.

\medterm{Testicle, Diseases Of} Conditions affecting the testis, including twisting, infection, injury, tumors, blood-vessel problems, and abnormal development. Symptoms may include pain, swelling, a lump, fertility problems, or changes in hormone production.
\medterm{Testicle, Retractile} Alternate index label; see \textit{Retractile testicle}.

\medterm{Testicle, Undescended} Alternate index label; see \textit{Undescended testicle}.

\medterm{Testicles} The paired male reproductive glands in the scrotum that produce sperm and testosterone; pain, enlargement, a new lump, or sudden change requires examination.

\medterm{Testicular Atrophy} Shrinkage and loss of functional tissue in a testicle caused by undescended testis, infection, torsion, trauma, varicocele, hormones, aging, or other injury.

\medterm{Testicular Biopsy} Removal of a small sample of testicular tissue for microscopic examination, used selectively to evaluate infertility, atypical lesions, or suspected testicular disease.

\medterm{Testicular Cancer} Cancer arising in a testicle, most often in younger adult men. It commonly appears as a painless lump, enlargement, heaviness, or ache. Ultrasound and blood tests help assess a suspected tumor; treatment is highly effective for many people, especially when found early.

\synonyms
Cancer, Testicular

\medterm{Testicular Disorders} Conditions affecting the testes or surrounding structures, including torsion, infection, inflammation, trauma, tumors,
and hormonal disorders. Symptoms may include pain, swelling, a mass, infertility, or
changes in sexual development and require appropriate evaluation.

\synonyms
Infection Testicle (Testicular Disorders)
Pain Scrotum (Testicular Disorders)
Torsion Testicle (Testicular Disorders)
Tumor Testicle (Testicular Disorders)



\medterm{Testicular Gonadoblastoma} A rare mixed germ-cell and sex-cord tumor arising in a dysgenetic or undescended gonad, often associated with Y-chromosome material and risk of another germ-cell tumor.

\medterm{Testicular Hydrocele} A collection of clear fluid around a testicle within the tunica vaginalis, causing painless scrotal enlargement that transilluminates; it may be congenital, reactive, or secondary to disease.

\medterm{Testicular Mass} A lump or enlargement within or around a testicle caused by tumor, infection, inflammation, cyst, hematoma, or torsion; a solid intratesticular mass requires urgent ultrasound and specialist assessment.

\medterm{Testicular Microlithiasis} Multiple tiny calcifications within the testicular tissue seen on ultrasound. The finding alone does not establish cancer, and follow-up depends on other risk factors and clinical findings.

\medterm{Testicular Self-Exam} A testicular self-exam is a gentle check for new lumps, swelling, heaviness, or changes in size or shape; a new abnormality should be medically assessed.

\medterm{Testicular Torsion} Surgical emergency caused by twisting of the spermatic cord around its longitudinal
axis, cutting off venous return and leading to arterial occlusion, ischemia, and
infarction of the testicle.

\synonyms
Torsion, Testicular

\medterm{Testicular Torsion Emergency} Rotation of the spermatic cord that compromises testicular blood flow, causing sudden
unilateral scrotal pain, nausea, a high-riding testis, and an abnormal cremasteric
reflex. Immediate surgical exploration and detorsion are required; imaging must not
delay treatment when suspicion is high.

\medterm{Testicular Torsion Repair} Urgent surgery to untwist a spermatic cord and restore testicular blood flow, followed by fixation of the affected and usually opposite testis to prevent recurrence.

\medterm{Testis} The male reproductive gland that produces spermatozoa in seminiferous tubules and
secretes testosterone from interstitial Leydig cells. It is normally located in the
scrotum.

\medterm{Testis (Testicle)} The male gonad, a paired ovoid organ (4--5 cm \(\times\) 2--3 cm, weight 15--20 g) located in
the scrotum, responsible for sperm production (spermatogenesis) and testosterone
secretion. Each testis is covered by the tunica vaginalis (peritoneal remnant) and
tunica albuginea (dense fibrous capsule).

\synonyms
Testicle

\medterm{Testis Disorder} A condition affecting a testicle, including torsion, infection, tumor, trauma, infertility-related damage, hormonal dysfunction, or abnormal development.

\medterm{Testis Neoplasm} A tumor arising in testicular tissue, most commonly a germ-cell tumor in younger adults; painless enlargement, heaviness, or a solid mass requires prompt ultrasound and evaluation.

\medterm{Testolactone} An aromatase inhibitor formerly used to reduce estrogen production in selected endocrine conditions; current use is limited by availability and newer alternatives.

\medterm{Testosterone} A steroid sex hormone produced mainly by the testes and, in smaller amounts, by the ovaries and adrenal glands. It supports reproductive development, muscle and bone maintenance, red-blood-cell production, libido, and other physiologic functions; levels vary with age, sex, time of day, illness, and medicines.

\medterm{Testosterone Deficiency} Alternate index label; see \textit{Male hypogonadism}.

\medterm{Testosterone Low (Low Testosterone (Low T))} Alternate index label for Low Testosterone (Low T).

\medterm{Testosterone Propionate} A short-acting injectable testosterone ester used in selected hormone-replacement or androgen-deficiency settings; risks include erythrocytosis, infertility, acne, mood effects, and prostate-related monitoring needs.

\medterm{Tests And Visits Before Surgery} Preoperative tests and visits review health conditions, medicines, anesthesia risk, needed laboratory or imaging tests, and instructions for the operation.

\medterm{Tests For H Pylori} Tests for Helicobacter pylori infection, including urea breath testing, stool antigen testing, biopsy-based testing, or selected blood tests; choice depends on treatment and clinical context.

\medterm{Tet Spell} A sudden episode of severe bluish coloring and low oxygen in a child with unrepaired tetralogy of Fallot. Crying, feeding, or exertion may trigger it; the child may breathe rapidly, become limp, or faint. It is an emergency requiring calming, positioning, oxygen, and urgent specialist care.

\synonyms
Hypercyanotic Spell; Paroxysmal Dyspneic Attack; Blue Spell

\medterm{Tetanus} A serious infection caused by bacteria entering a wound and producing a toxin that causes severe muscle stiffness and painful spasms. It can cause lockjaw, breathing problems, and death without urgent treatment.

\synonyms
Lockjaw

\medterm{Tetanus Toxin} A neurotoxin made by the bacterium Clostridium tetani. It blocks inhibitory signals in the nervous system, causing painful muscle stiffness and spasms, including lockjaw; vaccination and urgent wound care help prevent tetanus.
\medterm{Tetany} A clinical syndrome of neuromuscular hyperexcitability caused by hypocalcaemia (most
common—total calcium <2.1 mmol/L, ionised calcium <1.1 mmol/L), but also caused by
hypomagnesaemia, hypokalaemia, and alkalosis (respiratory or metabolic).

\medterm{Tethered Cord Syndrome} A congenital or acquired condition where the spinal cord is abnormally attached
(tethered) to surrounding tissues, restricting its normal upward movement within the
spinal canal. Caused by a thickened filum terminale, lipomeningomyelocele,
diastematomyelia, or previous myelomeningocele repair.

\medterm{Tetrabenazine} A medicine that reduces storage and release of certain brain chemicals involved in movement. It treats selected involuntary movement disorders but may cause sleepiness, depression, restlessness, stiffness, swallowing problems, or abnormal movements.
\medterm{Tetracaine} Tetracaine is a potent ester local anaesthetic used topically or for spinal anaesthesia; excessive dose can cause seizures, cardiovascular toxicity, and methemoglobinaemia.

\medterm{Tetrachloroethylene} Tetrachloroethylene, or perchloroethylene, is an industrial solvent used in dry cleaning; substantial exposure can affect the nervous system, liver, and kidneys.

\medterm{Tetrachlorvinphos} An organophosphate insecticide that can block acetylcholinesterase. Significant exposure may cause sweating, vomiting, pinpoint pupils, muscle twitching, slow breathing, seizures, or coma and requires urgent treatment.
\medterm{Tetracycline} A medicine, biologic, or pharmacologic term related to tetracycline; its mechanism, indication, interactions, and adverse effects depend on the specific agent.

\medterm{Tetracycline Stain} Intrinsic discoloration from tetracycline antibiotics during tooth development. Gray or
brown bands. Treatment: bleaching, veneers, crowns.

\medterm{Tetracyclines} A class of bacteriostatic antibiotics that bind reversibly to the 30S ribosomal subunit
and block the attachment of aminoacyl-tRNA to the acceptor site, inhibiting protein
synthesis. They include tetracycline, doxycycline, and minocycline.

\medterm{Tetrahydrozoline Poisoning} Poisoning from swallowing eye drops or nasal drops containing tetrahydrozoline. It can cause extreme sleepiness, slow heart rate, low blood pressure, and breathing problems, especially in children.

\medterm{Tetralogy Of Fallot} A congenital heart defect comprising ventricular septal defect, right-ventricular outflow obstruction, overriding aorta, and right-ventricular hypertrophy.
\medterm{Tetraplegia} Paralysis or severe weakness of all four limbs, usually from cervical spinal-cord disease or injury.
\medterm{Tetraplegia (Quadriplegia)} Paralysis or severe weakness affecting all four limbs and usually the trunk, commonly caused by cervical spinal-cord injury or disease.

\synonyms
Quadriplegia

\medterm{TGA} Alternate index label; see \textit{Transient global amnesia}.

\medterm{Th1/Th2/Th17/Tfh/Treg Balance} A simplified description of how different helper T-cell groups coordinate immune responses. Some promote defense against infections, others support antibody production or inflammation, and regulatory cells limit excessive immune activity. Imbalance may contribute to allergy, autoimmune disease, or persistent infection.

\medterm{Thalamus} Bilateral ovoid gray matter structure forming the third ventricle's lateral walls, the
brain's relay station for nearly all sensory (except olfaction), motor, and limbic
information to the cortex. Numerous nuclei: VPL (spinothalamic, dorsal column), VPM
(trigeminal), LGN (visual), MGN (auditory), VL/VA (motor).

\medterm{Thalassaemia} The British spelling of thalassemia, inherited reduced production of alpha- or beta-globin chains causing anemia.
\medterm{Thalassemia} A group of inherited hemoglobin disorders caused by reduced or absent synthesis of one
or more globin chains, leading to microcytic hypochromic anemia.

\synonyms
Cooley's Anemia
Mediterranean Anemia

\medterm{Thalassemia Major} A severe inherited blood disorder in which the body makes too little hemoglobin. It causes serious anemia beginning in childhood and may require regular blood transfusions, with monitoring for iron buildup and organ damage.

\medterm{Beta-Thalassemia Trait} The heterozygous carrier state for a beta-globin gene variant. It is usually asymptomatic or causes mild microcytosis with mild or no anemia; hemoglobin electrophoresis often shows increased HbA2. It is also called beta-thalassemia minor.

\synonyms
Beta Thalassemia Trait
Beta-Thalassemia Minor
Thalassemia Minor

\medterm{Thalidomide} A medicine that changes immune activity and is used for selected cancers and inflammatory conditions. It can cause severe birth defects, nerve damage, blood clots, sleepiness, and other harms, requiring strict safety controls.
\medterm{Thallium} Thallium is a toxic heavy metal that can cause gastrointestinal illness followed by painful peripheral neuropathy, hair loss, skin changes, kidney or liver injury, seizures, and cardiovascular or neurologic toxicity. Suspected exposure requires immediate poison-control or emergency guidance.

\synonyms
Poisoning Thallium (Thallium)

\medterm{Thallium Scan} A nuclear-medicine test using thallium-201 to show blood flow to heart muscle or selected tumors. It is now less common because other tracers are often preferred, but may still have specific uses.

\medterm{Thallium-201} SPECT radionuclide, half-life 73 hours. Used for myocardial perfusion imaging.
Mechanism: K+ analog, taken up by Na+/K+ ATPase.

\medterm{Thanatology} The scientific study of death, dying, postmortem changes, and the psychological, social, and ethical aspects
of bereavement. In forensic medicine, it also examines the circumstances and evidence
surrounding death.

\medterm{Thanatophoric Dysplasia} A severe skeletal dysplasia, usually caused by an FGFR3 variant, with markedly shortened
limbs, narrow thorax, macrocephaly, and characteristic skull and vertebral
abnormalities. The severe thoracic restriction commonly causes neonatal respiratory
failure.





\medterm{Theca} A connective-tissue sheath, especially the theca surrounding an ovarian follicle.
\medterm{Theca Lutein Cyst} A usually temporary ovarian cyst caused by strong stimulation from pregnancy-related hormones or fertility treatment. It is often present in both ovaries and may cause abdominal pressure or pain; it commonly shrinks when hormone levels return to normal.

\medterm{Thecoma} A usually benign ovarian sex-cord stromal tumor composed of theca-like cells, often producing estrogen and presenting with abnormal uterine bleeding or an incidental adnexal mass.
\medterm{Thelarche} The onset of breast development, usually the first visible sign of female puberty. It commonly begins between ages 8 and 13, with substantial normal variation.

\medterm{Thenar Eminence} The fleshy muscular prominence at the base of the thumb.
\medterm{Theophylline} A medicine that can relax airway muscles and improve breathing in selected lung diseases. Its safe and harmful doses are close together, so blood levels, medicine interactions, heart rhythm, and side effects need monitoring.
\medterm{Theranostic Pair} A matched diagnostic and treatment approach using related agents that target the same biological feature. One agent images the target and the other delivers treatment, often radiation.
theranostics.

\medterm{Theranostics} Combined therapy and diagnostics using matched radiopharmaceutical pairs. Diagnostic
isotope (e.g., Ga-68) identifies target, therapeutic isotope (e.g., Lu-177) delivers
treatment. Examples: Ga-68/Lu-177 DOTATATE (neuroendocrine tumors), PSMA theranostics
(prostate cancer).


\medterm{Therapeutic Abortion} Intentional termination of a pregnancy because of medical, fetal, or other legally recognized indications, according to applicable law and clinical guidance.

\medterm{Therapeutic Alliance} The collaborative, trusting relationship between therapist and patient that is
foundational to effective psychotherapy. It includes agreement on goals, tasks, and the
bond between therapist and patient. It is the strongest predictor of psychotherapy
outcome across all therapeutic modalities.

\medterm{Therapeutic Cloning} Creating a genetically matched embryonic stem-cell line for research or potential tissue replacement; it is distinct from reproductive cloning and remains ethically and technically limited.

\medterm{Therapeutic Exercise} Planned physical activity used to improve strength, flexibility, balance, movement, endurance, or daily function and to reduce pain or disability. A clinician selects and gradually adjusts exercises according to the person’s condition, abilities, symptoms, and goals.

\medterm{Therapeutic Hypothermia} Deliberate, closely monitored cooling of the body to reduce tissue damage in selected emergencies, such as after some cardiac arrests or in severe overheating-related illness. Cooling is used with treatment of the underlying cause and can cause abnormal heart rhythms, shivering, or electrolyte changes.

\medterm{Therapeutic Index} The therapeutic index is a quantitative measure of the relative safety of a drug,
defined as the ratio of the dose that produces toxicity in 50 percent of the population
to the dose that produces the desired therapeutic effect in 50 percent of the
population.

\medterm{Therapeutic Window} The range of plasma drug concentrations between the minimum effective concentration
(producing a therapeutic effect) and the minimum toxic concentration (producing
toxicity), in which efficacy is achieved with acceptable toxicity.

\medterm{Therapeutics} The branch of medicine concerned with the treatment of disease, encompassing all
interventions used to prevent, manage, or cure illness.

\medterm{Therapy} Treatment intended to prevent, relieve, or cure disease or to improve function and quality of life.

\medterm{Therapy Monitoring} Repeated clinical assessment, laboratory testing, or imaging used to evaluate response, adverse effects, and ongoing need for treatment.

\medterm{Thermal Injury} Tissue damage caused by heat, flame, hot liquid, steam, or cold. Burn severity depends
on depth, total body surface area, location, inhalation injury, associated trauma, and
patient factors; assessment should include airway and circulatory complications.

\medterm{Thermal Injury (Burns)} Tissue damage caused by heat. Initial care includes stopping the exposure, removing loose
constricting material, cooling briefly with clean water, covering the injury, and assessing
depth and extent. Larger or deeper burns require specialized evaluation, fluid management,
pain control, wound care, and tetanus assessment.

\synonyms
Burns

\medterm{Thermal Therapy} Treatment using heat or cold to influence pain, swelling, muscle tension, or tissue temperature. Effects and safety depend on the condition and method used.

\medterm{Thermogenesis} The production of heat by the body. It comes mainly from normal metabolism, muscle activity such as shivering, digestion, and specialized fat tissue that releases stored energy as heat.
\medterm{Thermography} Imaging or recording of heat patterns from the body surface.
\medterm{Thermoluminescent Dosimeter} A radiation-monitoring device that measures accumulated exposure by light released from an irradiated crystal when heated.
\medterm{Thermometer} A device used to measure body temperature. Common types measure temperature in the mouth, ear, armpit, or across the forehead. Readings vary with the site and method, so the measurement site and scale should be recorded.

\medterm{Thermometer Scales} Numbering systems used to report temperature, such as Celsius, Fahrenheit, and Kelvin. The scale should be identified when a temperature is communicated.
\medterm{Thermoneutral Zone} The range of ambient temperatures over which body temperature is maintained by basal
metabolic heat production alone, without additional thermoregulatory responses such as
shivering or sweating. In humans the zone is approximately 27-32 degrees C for a resting
nude adult. Within the zone, vasomotor tone adjusts heat loss.

\medterm{Thermoplasty} A bronchoscopic procedure that uses radiofrequency energy to ablate smooth muscle in the
airways, reducing bronchospasm. It is approved for severe, persistent asthma in adults
not controlled with ICS/LABA. It reduces exacerbations and emergency department visits
but does not improve lung function.

\medterm{Thermoreceptor} A nerve ending that senses warmth or cold in the skin and other tissues. Signals travel through the spinal cord to the brain, helping a person recognize temperature and helping the body control heat through sweating, shivering, and changes in skin blood flow. Damage to these nerves can reduce temperature sensation and increase the risk of burns or cold injury.
\medterm{Thermoregulation} The physiological maintenance of core body temperature near the set point (~37 degrees
C) via a balance of heat production and heat loss, controlled by the preoptic area of
the hypothalamus.

\medterm{Thermostat} A thermostat regulates temperature by sensing heat and controlling a heating or cooling system; accurate environmental temperature can matter for patients, specimens, and medicines.

\medterm{Thiabendazole} An antiparasitic medicine used for selected worm infections and, historically, some fungal infections. It can cause nausea, dizziness, liver problems, or nervous-system effects, so safer alternatives are often preferred.
\medterm{Thiamin} Vitamin B1, a water-soluble vitamin required for carbohydrate metabolism and neural function; severe deficiency can cause beriberi or Wernicke encephalopathy.

\medterm{Thiamine} Vitamin B1, a water-soluble vitamin required for carbohydrate metabolism and nerve function. Severe deficiency
can cause beriberi or Wernicke encephalopathy.

\medterm{Thiamine (Vitamin B1)} A water-soluble vitamin whose active coenzyme form, thiamine pyrophosphate, is required
by pyruvate dehydrogenase, alpha-ketoglutarate dehydrogenase, transketolase, and
branched-chain ketoacid dehydrogenase. Deficiency causes beriberi and Wernicke
encephalopathy.

\synonyms
Vitamin B1

\medterm{Thiazide Diuretics} Thiazide diuretics act on the distal convoluted tubule by inhibiting the sodium-chloride
cotransporter, reducing sodium and chloride reabsorption and producing natriuresis and
diuresis. They are less potent than loop diuretics but have a longer duration of action,
typically 6 to 12 hours.

\medterm{Thiazide Overdose} Taking too much of a thiazide diuretic can cause dehydration, low blood pressure, abnormal sodium or potassium levels, confusion, or abnormal heart rhythms.

\medterm{Thiazides} Diuretics that inhibit sodium-chloride reabsorption in the distal convoluted tubule.
\medterm{Thiazolidinedione Drugs} A group of diabetes medicines that make body tissues more sensitive to insulin by activating a protein called PPAR-gamma. Pioglitazone is the main example; possible effects include weight gain, swelling, bone fractures, and worsening heart failure.
\medterm{Thiazolidinediones} Insulin sensitizers activating peroxisome proliferator-activated receptors primarily in
adipose tissue and skeletal muscle. Agents include pioglitazone and rosiglitazone.
Mechanisms include enhanced insulin-mediated glucose uptake, altered adipokine
secretion, and reduced hepatic glucose output.

\medterm{Thigh} The part of the lower limb between the hip and knee. It contains the femur, large muscles, major blood vessels, and nerves that help move and support the leg.
\medterm{ThinPrep Pap Test} A cervical screening test in which cells collected from the cervix are placed in a liquid vial and examined for precancerous changes, cancer, or sometimes HPV-related abnormalities.
\synonyms
Thin Prep Pap Screen

\medterm{Thiopental Sodium} An ultra-short-acting barbiturate, used for induction of general anesthesia (IV, onset
30--60 seconds, duration 5--10 minutes) and as a third-line agent for raised
intracranial pressure (cerebroprotective effect: reduces cerebral metabolism, cerebral
blood flow, and intracranial pressure).

\medterm{Thiopentone Sodium} A short-acting barbiturate anesthetic used for induction in selected settings.
\medterm{Thioridazine} An older phenothiazine antipsychotic with important cardiac and neurologic toxicity.
\medterm{Thioridazine Overdose} Poisoning from taking too much thioridazine, an antipsychotic medicine. It can cause severe sleepiness, low blood pressure, seizures, and dangerous heart rhythms.

\medterm{Thiotepa} A chemotherapy medicine that damages DNA and stops cells from multiplying. It is used for selected cancers and sometimes before a stem-cell transplant; side effects may include low blood counts, infection, nausea, mouth sores, and organ injury.
\medterm{Third Stage Of Labor} From delivery of the baby to delivery of the placenta. The placenta separates from the
uterine wall and is expelled spontaneously or with gentle cord traction. Normal
duration: 5-30 minutes.

\medterm{Third Trimester} The period of pregnancy from 28 weeks until birth. The fetus grows rapidly and continues developing its lungs, brain, and other organs. Prenatal visits monitor blood pressure, growth, movement, and symptoms; tests and screening are scheduled according to local guidance and individual risk.

\medterm{Third-Degree AV Block} Complete heart block with no conduction between atria and ventricles. Atria and
ventricles beat independently (AV dissociation). Ventricular escape rhythm: junctional
(40-60, narrow QRS) or ventricular (20-40, wide QRS).

\medterm{Third-Degree Burn} A full-thickness burn that destroys the epidermis and dermis and may extend into subcutaneous tissue.
The skin may appear white, brown, or charred and sensation may be reduced. Assessment
includes burn depth, extent, inhalation injury, and associated trauma; deep or extensive
burns require specialist treatment and may need grafting.

\medterm{Third-Degree Perineal Tear} A severe tear during vaginal birth that extends from the vaginal opening into the muscles controlling the anus. It needs prompt examination and careful repair, followed by bowel-softening treatment and follow-up; pain, infection, difficulty controlling stool, or a fistula can result.

\synonyms
Third-Degree Tear; Obstetric Anal Sphincter Tear

\medterm{Thirst} The conscious desire to drink, regulated mainly by body-fluid concentration and circulating blood volume. Excessive thirst may occur with dehydration, diabetes, medicines, kidney problems, or other conditions and may need evaluation.
\medterm{Thirst - Absent} Reduced or absent desire to drink despite a potential need for fluids, which can occur with hypothalamic disease, aging, altered consciousness, medications, or severe illness.

\medterm{Thirst - Excessive} Polydipsia, an abnormally strong or persistent desire to drink, commonly associated with dehydration, diabetes mellitus, diabetes insipidus, medicines, or psychiatric conditions.

\medterm{Thomsen Disease} An inherited muscle disorder causing stiffness because muscles relax slowly after they contract. Symptoms often begin in childhood, affect the hands, legs, eyelids, or jaw, and improve after repeated movement. Strength and life expectancy are usually preserved, although stiffness can interfere with activity.

\synonyms
Autosomal Dominant Myotonia Congenita; Myotonia Congenita (Thomsen Type)

\medterm{Thoracentesis} Insertion of a needle or catheter into the pleural space to remove fluid for diagnosis or to relieve respiratory symptoms caused by a pleural effusion.

\medterm{Thoracic} Thoracic means relating to the chest, the thoracic spine, or structures within the thorax such as the heart and lungs.

\medterm{Thoracic Aortic Aneurysm} A weakened, widened area of the aorta within the chest. It may cause no symptoms but can press on nearby structures or rupture, with risks increased by hypertension, inherited disorders, or a bicuspid valve.

\medterm{Thoracic Cage} The bony and cartilaginous framework formed by the thoracic vertebrae, ribs, costal
cartilages, and sternum. It protects thoracic organs and assists respiratory movements.

\medterm{Thoracic Cavity} The chest cavity, enclosed by the ribs and chest muscles, containing the lungs, heart,
esophagus, and trachea.

\medterm{Thoracic Duct} Body's largest lymphatic vessel, approximately 40 cm long, draining lymph from the
entire body below the diaphragm and left side above it (approximately 75 percent of
total lymph).

\medterm{Thoracic Endometriosis} Endometrial tissue within the thoracic cavity causing catamenial pneumothorax,
hemothorax, or lung nodules, recurring cyclically with menstruation and often requiring
hormonal suppression.

\medterm{Thoracic Outlet Syndrome} Pressure on nerves or blood vessels as they pass from the neck into the arm. It may cause neck, shoulder, or arm pain, tingling, weakness, swelling, color change, or a cold hand. Treatment depends on the cause and may include physical therapy or surgery.

\synonyms
TOS (Thoracic Outlet Syndrome)

\medterm{Thoracic Spine CT Scan} Computed tomography of the thoracic spine used to evaluate vertebral fractures, alignment, bone lesions, spinal canal narrowing, and other bony abnormalities.

\medterm{Thoracic Spine X-Ray} Plain radiographic imaging of the thoracic spine used to assess vertebral alignment, fractures, curvature, degenerative change, and selected bone lesions.

\medterm{Thoracic Surgeon} Surgeon specializing in operations on the chest, including lungs, heart, esophagus, and mediastinum. A surgeon who performs operations on chest organs, including the lungs, esophagus, heart, and mediastinum.
\medterm{Thoracic Vertebrae} Have costal facets on bodies and transverse processes for rib articulation. Heart-shaped
bodies, long downward spinous processes, coronally-oriented facets allowing rotation and
lateral flexion but limiting flexion/extension. Small circular vertebral foramen.

\medterm{Thoracic Vertebrae (T1-T12)} The twelve vertebrae of the chest region, each articulating with ribs and protecting the thoracic spinal cord while supporting the trunk and rib cage.
\medterm{Thoracolumbar Spine Trauma} Injury to the thoracic or lumbar spine. Assessment includes fracture morphology,
posterior ligamentous complex integrity, and neurologic status. Classification systems such
as AO Spine and TLICS help guide management, which depends on stability and the clinical
context.

\medterm{Thoracoscopy} A minimally invasive procedure using a camera inserted through the chest wall to examine or treat the pleura, lungs, or mediastinum.

\medterm{Thoracotomy} A surgical incision into the chest wall to access the thoracic cavity (pleural space,
lungs, heart, great vessels, esophagus, diaphragm, mediastinum).

\medterm{Thorax} The thorax is the chest region between the neck and abdomen, enclosed by ribs and containing the heart, lungs, major vessels, and mediastinal structures.

\medterm{Thought Blocking} An abrupt interruption in the train of thought before it is completed, as if the thought
was removed or pushed out of the patient's mind. The patient may pause mid-sentence,
appear confused, and report that the thought was ``taken out'' or ``blocked.'' It is a
formal thought disorder seen in schizophrenia.

\medterm{Thought Disorders} Disturbances of the form, flow, or content of thinking, occurring in psychosis and other conditions.
\medterm{Threadworm} An intestinal infection or the worm Enterobius vermicularis, also called pinworm.
\medterm{Threadworm (Enterobiasis)} An intestinal parasitic infection caused by Enterobius vermicularis, also called
pinworm. Infection spreads when eggs are swallowed, often after contact with contaminated hands, bedding, food, or
surfaces.

\synonyms
Enterobiasis

\medterm{Threadworms (Pinworms)} Small intestinal worms, Enterobius vermicularis, whose females lay eggs around the anus and commonly cause nocturnal itching.

\synonyms
Pinworms

\medterm{Threatened Miscarriage} Vaginal bleeding during early pregnancy with a closed cervix and no passage of pregnancy tissue; the pregnancy may remain viable.
Evaluation commonly includes clinical assessment, serial human chorionic gonadotropin testing when indicated, and ultrasound.

\synonyms
Threatened abortion.

\medterm{Three-Day Measles} Alternate index label; see \textit{Rubella}.

\medterm{Three-Dimensional Ultrasound} Ultrasound that combines many sound-wave measurements to create a three-dimensional image. It is used in obstetrics, gynecology, and other settings to show anatomy from several angles.

\medterm{Threonine} An essential amino acid used to build proteins and other body molecules. The body cannot make enough of it, so it must come from food.
\medterm{Threshold} The lowest amount of a stimulus or measurement needed to produce a detectable response or effect.
\medterm{Thrill} A palpable vibration over the chest or a blood vessel caused by turbulent blood flow.
\medterm{Throat} The pharynx and adjacent structures behind the mouth and nose.
\medterm{Throat Cancer} Cancer arising in the pharynx or nearby structures of the throat. Persistent hoarseness, difficulty or pain
with swallowing, a neck lump, unexplained weight loss, or persistent sore throat requires
medical evaluation, especially in people who use tobacco or alcohol.

\synonyms
Cancer, Throat

\medterm{Throat Cancer (Larynx Cancer)} Alternate index label for Larynx Cancer.

\medterm{Throat Or Larynx Cancer} Throat or larynx cancer is malignant growth in the pharynx or voice box and may cause persistent hoarseness, swallowing difficulty, pain, or a neck lump.

\medterm{Throat Swab Culture} Culture of a throat specimen to detect bacterial pathogens, especially group A Streptococcus, when evaluating pharyngitis.

\synonyms
Throat Culture

\medterm{Throbbing} A pulsating or beating quality of pain or another sensation. It can occur with migraine, inflammation, injury, infection, or increased awareness of a pulse.
\medterm{Thrombectomy} A surgical or interventional radiology procedure to remove a blood clot (thrombus) from
a blood vessel, restoring blood flow. Types: (1) Endovascular thrombectomy (mechanical
thrombectomy)—for acute ischemic stroke due to large vessel occlusion (LVO—internal
carotid artery, middle cerebral artery M1 segment).

\medterm{Thrombectomy (Vascular)} Removal of a blood clot from a blood vessel, using an open surgical or catheter-based procedure.

\synonyms
Vascular

\medterm{Thrombin} A blood-clotting substance that changes fibrinogen into fibrin, the protein threads that form a clot, and also activates platelets. The body makes thrombin after a blood vessel is injured; too much or too little clotting can cause disease.
\medterm{Thromboangiitis Obliterans} A tobacco-related inflammatory disease blocking small and medium arteries and veins, often causing painful fingers, toes, ulcers, or gangrene.
\medterm{Thromboangiitis Obliterans (Buerger Disease)} An inflammatory disease of small and medium blood vessels, strongly associated with tobacco
use. It can cause pain, color changes, ulcers, or gangrene in the hands or feet. Complete
tobacco cessation is essential; further treatment depends on ischemia, tissue loss, and
other complications.

\synonyms
Buerger Disease

\medterm{Thrombocyte (Platelet)} A small blood-cell fragment made in the bone marrow. Platelets stick to damaged blood vessels, become activated, and join together to form the first plug that stops bleeding. Too few platelets or poorly working platelets can cause easy bruising, tiny skin spots, nosebleeds, or other bleeding. Too many or overly active platelets can increase the risk of unwanted clots.
\medterm{Thrombocytes} Another name for platelets: small blood-cell fragments made in the bone marrow that form plugs to stop bleeding. Too few can cause bruising or bleeding, while too many or overly active platelets can increase the risk of unwanted clots.

\medterm{Thrombocytopathy} A disorder of platelet function that impairs primary hemostasis despite a normal or
near-normal platelet count. It may be inherited or acquired and can cause mucocutaneous
bleeding, easy bruising, prolonged bleeding after procedures, or menorrhagia.

\medterm{Thrombocytopenia} Thrombocytopenia is a low platelet count caused by reduced production, increased destruction or consumption, sequestration, or dilution and may increase bleeding risk.

\synonyms
Low Platelet Count

\medterm{Thrombocytopenia (Low Platelet Count)} Alternate index label for Low Platelet Count.

\medterm{Thrombocytopenic} Characterized by too few circulating platelets. It may cause easy bruising, nosebleeds, heavy menstrual bleeding, or other bleeding, although some people have no symptoms.
\medterm{Thrombocytosis} An elevated platelet count (>450,000/$\mu$L). Reactive (secondary): infection, inflammation,
iron deficiency, malignancy, post-splenectomy, hemorrhage, exercise—usually mild, no
thrombotic risk.

\medterm{Thromboelastography} A viscoelastic whole-blood assay that measures the initiation, formation, strength, and
breakdown of a clot. It provides a functional overview of coagulation and fibrinolysis
and may guide targeted blood-component therapy in major bleeding.

\medterm{Thromboembolism} Obstruction of a blood vessel by a clot that formed elsewhere and traveled through the circulation.

\medterm{Thrombolysis} The dissolution of a blood clot using a medicine such as a fibrinolytic agent. It is used only in selected conditions because it can cause serious bleeding.

\medterm{Thrombolysis For PE} Emergency treatment using a clot-dissolving medicine for a life-threatening lung clot causing low blood pressure or shock. It can restore blood flow quickly but carries a serious risk of bleeding.

\medterm{Thrombolysis For Stroke} Emergency treatment for selected strokes caused by a blood clot. A clot-dissolving medicine is given through a vein within a limited time after symptoms begin, after brain imaging excludes bleeding; serious bleeding is a possible complication.

\medterm{Thrombolytic Agents} Medicines that activate fibrinolysis and dissolve an established blood clot; they are used selectively for time-critical conditions because they can cause serious bleeding.
\medterm{Thrombolytic Therapy} Treatment with a medicine that breaks down a blood clot and may restore blood flow in selected emergencies such as certain strokes, heart attacks, or lung clots. The main danger is serious bleeding, including bleeding in the brain.

\medterm{Thrombophilia} A tendency for blood to form clots too easily. It may be inherited or develop with cancer, pregnancy, some medicines, inflammation, or other conditions and can increase the risk of clots in veins.

\medterm{Thrombophilia Testing} Thrombophilia testing evaluates for inherited and acquired hypercoagulable states in
patients with venous thromboembolism, recurrent pregnancy loss, or strong family history
of VTE.

\medterm{Thrombophlebitis} Thrombophlebitis is inflammation of a vein associated with a blood clot and may be superficial or involve deep veins with risk of pulmonary embolism.

\medterm{Thrombopoiesis} The production of platelets from megakaryocytes in the bone marrow. It is regulated
principally by thrombopoietin and involves megakaryocyte maturation, proplatelet
formation, and release of circulating platelets.

\medterm{Thrombosed Intracranial Aneurysm} An intracranial arterial aneurysm that has undergone partial or complete thrombosis of
its lumen, which may reduce rupture risk but can cause mass effect on adjacent brain
structures. Presents with headache, cranial nerve palsies, or seizures depending on
location. Diagnosis: CT, MRI, or conventional angiography.

\medterm{Thrombosis} The formation of a blood clot within the cardiovascular system during life, potentially
obstructing blood flow and causing tissue ischemia or infarction. Arterial thrombi
typically form on ruptured atherosclerotic plaques, are composed primarily of platelets
and fibrin, and appear pale and grossly attached to the vessel wall.

\medterm{Thrombosis Of The Inferior Vena Cava} Formation of a blood clot within the inferior vena cava, which may occur in the setting
of IVC filters, malignancy, hypercoagulable states, or extension from iliofemoral deep
vein thrombosis. May cause bilateral lower extremity edema, abdominal pain, and hepatic
or renal dysfunction. Requires anticoagulation and possible thrombectomy.

\medterm{Thrombotic Microangiopathy} A group of disorders in which tiny clots form in small blood vessels, damaging red blood cells and lowering the platelet count. The clots can injure the kidneys, brain, heart, or other organs. Causes include immune disorders, infections, medicines, pregnancy, and severe illness; urgent specialist treatment is often needed.

\medterm{Thrombotic Thrombocytopenic Purpura} A rare, life-threatening blood disorder in which small clots form in tiny blood vessels and use up platelets. It can damage the brain, kidneys, heart, and other organs.
\medterm{Thrombotic Vasculopathy} A spectrum of disorders characterized by thrombus formation within blood vessels of
various sizes, often associated with antiphospholipid syndrome, disseminated
intravascular coagulation, thrombotic thrombocytopenic purpura, or malignancy.

\medterm{Thromboxane} A fat-derived chemical messenger made mainly by activated platelets. It encourages platelets to stick together and narrows blood vessels, helping clots form but also contributing to unwanted clotting.
\medterm{Thrombus} A blood clot that forms within a blood vessel or heart and remains attached at its site of formation, potentially obstructing blood flow or embolizing.
\medterm{Thrombus Organization} The process by which a thrombus is gradually replaced by granulation tissue and
eventually fibrous scar tissue. The process begins when inflammatory cells, particularly
macrophages, infiltrate the thrombus and begin digesting the fibrin and dead cells
through fibrinolysis.

\medterm{Thrush} A Candida infection, usually of the mouth or throat, causing creamy white plaques, soreness, altered taste, or discomfort with eating or swallowing. Vaginal candidiasis is a separate condition. Risk increases with antibiotics, inhaled corticosteroids, diabetes, or immune suppression.

\synonyms
Oral Thrush

\medterm{Thrush - Children And Adults} Oral candidiasis caused by Candida overgrowth, producing removable white plaques, soreness, redness, altered taste, or cracks at the mouth corners; risk rises with antibiotics, inhaled steroids, diabetes, or immune suppression.

\medterm{Thrush And Other Yeast Infections In Children} Thrush is Candida overgrowth in the mouth, producing white plaques, soreness, or feeding difficulty; Candida can also affect skin folds or diaper areas. Persistent, recurrent, severe, or immune-associated infection needs pediatric evaluation and appropriate antifungal treatment.

\medterm{Thrush In Newborns} Oral Candida infection in an infant, causing white plaques on the tongue or mouth lining that may interfere with feeding and can be passed between infant and breastfeeding parent.

\medterm{Thrush, Oral} Alternate index label; see \textit{Oral thrush}.

\medterm{Thumb Arthritis} Degeneration or inflammation of a thumb joint, especially the joint at the base of the thumb, causing pain, stiffness, and reduced grip.

\synonyms
Basal Joint Arthritis

\medterm{Thumb CMC Arthritis} Thumb carpometacarpal arthritis: Eaton-Littler classification: I (normal), II
(slight narrowing, osteophytes <2mm), III (narrowing, osteophytes >2mm, subluxation), IV
(pantrapezial arthritis). Clinical: base of thumb pain, pinch weakness, thenar wasting,
squaring (subluxation).

\medterm{Thumb-Sucking} A common childhood habit of sucking the thumb; persistent habit can affect dental and facial development.
\medterm{Thymic Hyperplasia} Thymic hyperplasia is enlargement of the thymus from increased lymphoid tissue or rebound growth and may be incidental. It can be associated with autoimmune disease, infection, endocrine conditions, or stress; imaging and clinical context distinguish it from a mass.

\medterm{Thymocyte} An immature T lymphocyte developing in the thymus, where it undergoes selection before becoming a mature T cell.
\medterm{Thymoma} A tumor of thymic epithelial cells, associated with myasthenia gravis, pure red cell
aplasia, and hypogammaglobulinemia. Malignant behavior is determined by capsular
invasion and stage.

\synonyms
Neoplasm Thymomic (Thymoma)

\medterm{Thymus} Primary lymphoid organ in the anterior mediastinum, largest in childhood, involutes with
age. Site of T-cell maturation (thymopoiesis): thymocytes undergo positive and negative
selection to produce self-tolerant, MHC-restricted T cells. Epithelial cells produce
thymic hormones (thymosin, thymopoietin).

\medterm{Thymus Gland} An immune-system organ behind the breastbone where immature T lymphocytes learn to recognize harmful foreign material without attacking the body’s own tissues. It is most active during childhood and gradually becomes smaller after puberty.
\medterm{Thyrocalcitonin} Another name for calcitonin, a hormone made by thyroid C cells that lowers blood calcium mainly by reducing bone resorption.

\medterm{Thyroglobulin Surveillance} Thyroglobulin produced only by thyroid follicular cells serves as a tumor marker after
total thyroidectomy for differentiated thyroid carcinoma. In patients who have undergone
complete thyroidectomy and radioactive iodine ablation thyroglobulin should be
undetectable.

\medterm{Thyroglossal Duct Cyst} Congenital midline neck mass resulting from persistence of the thyroglossal duct the
embryological tract of thyroid descent. Most common congenital neck mass presenting in
childhood. Located at or near the hyoid bone and moves with swallowing and tongue
protrusion. Diagnosis is by ultrasound demonstrating a cystic midline neck mass.

\medterm{Thyroid} An endocrine gland in the front of the neck that produces thyroid hormones, which regulate metabolism, growth, and development.

\medterm{Thyroid Cancer} A malignant tumor of the thyroid. Major types include papillary, follicular, medullary, and anaplastic carcinoma, which differ in behavior and treatment.

\synonyms
Cancer, Thyroid

\medterm{Thyroid Cancer - Medullary Carcinoma} Medullary thyroid carcinoma is cancer of thyroid C cells that produce calcitonin and may be inherited as part of a multiple endocrine neoplasia syndrome.

\medterm{Thyroid Cancer - Papillary Carcinoma} Papillary thyroid carcinoma is the most common thyroid cancer and usually grows slowly, although it can spread to nearby lymph nodes.

\medterm{Thyroid Cancer Staging} American Joint Committee on Cancer TNM staging for differentiated thyroid carcinoma
incorporates tumor size, extent of invasion, lymph node metastases, and distant
metastases. Age is a critical prognostic factor with patients under fifty having more
favorable outcomes. Stage I is any tumor without metastases in patients under fifty.

\medterm{Thyroid Cartilage} The largest cartilage of the larynx, formed by two plates that meet in front and create the laryngeal prominence. It
protects the vocal folds and provides attachment for muscles and ligaments.


\medterm{Thyroid Disorders} Conditions affecting the thyroid gland, a butterfly-shaped endocrine organ in the neck
that produces hormones regulating metabolism, growth, and development.

\medterm{Thyroid Eye Disease} Autoimmune inflammatory disorder of the orbital tissues associated with Graves disease
characterized by expansion of extraocular muscles and orbital fat. Pathogenesis involves
orbital fibroblasts expressing the TSH receptor targeted by the same thyroid-stimulating
immunoglobulins affecting the thyroid.

\medterm{Thyroid Fine-Needle Aspiration} Image-guided sampling of thyroid-nodule cells for cytologic classification. Results are
commonly reported using the Bethesda system and guide surveillance, molecular testing,
surgery, or other management.

\medterm{Thyroid Follicular Carcinoma} A cancer arising from thyroid follicle cells and the second most common differentiated thyroid cancer. It tends to spread through blood vessels to the lungs or bones; diagnosis requires showing invasion of the capsule or blood vessels.

\medterm{Thyroid Function Tests} Laboratory assessment of the hypothalamic-pituitary-thyroid axis. TSH is the most
sensitive screening test with elevated levels indicating primary hypothyroidism and
suppressed levels indicating hyperthyroidism. Free T4 measures the biologically active
unbound thyroxine.

\medterm{Thyroid Gland} An endocrine gland located in the anterior neck, consisting of two lobes (right and
left) connected by an isthmus, lying anterior to the trachea between the cricoid
cartilage and the suprasternal notch.

\medterm{Thyroid Gland Enlargement} Alternate index label; see \textit{Goiter}.

\medterm{Thyroid Gland Removal} Surgery to remove part or all of the thyroid gland for cancer, overactivity, a large goiter, or a suspicious nodule. Possible effects include bleeding, voice changes from nerve injury, low calcium, and the need for lifelong thyroid hormone.


\medterm{Thyroid Gland, Diseases Of} Conditions affecting the thyroid gland, including underactivity, overactivity, inflammation, enlargement, nodules, and cancer. They may change energy, weight, temperature tolerance, heart rate, bowel habits, mood, or neck appearance and are assessed with examination and blood tests.
\medterm{Thyroid Hormone} The hormones T3 and T4, made by the thyroid gland, that help control energy use, temperature, heart rate, growth, and brain development. Too little can slow body functions, while too much can make them overactive.

\medterm{Thyroid Hormones} Thyroxine (T4) and triiodothyronine (T3), hormones that regulate metabolic rate, growth, brain development, heat production, and cardiovascular and gastrointestinal activity.

\medterm{Thyroid Nodule} A discrete lump within the thyroid gland that may be solid, cystic, benign, or malignant. Most are benign, but evaluation
uses clinical assessment, thyroid testing, ultrasound, and biopsy when indicated.

\medterm{Thyroid Nodule Evaluation} Systematic approach to evaluation of thyroid nodules detected by physical examination or
imaging. Thyroid ultrasound characterizes nodule features including size, composition,
echogenicity, margins, calcifications, and shape. TI-RADS classification stratifies
malignancy risk and guides biopsy decisions.

\medterm{Thyroid Nodules} A thyroid nodule is a discrete growth within the thyroid and is often benign. Ultrasound, thyroid-function testing, and selected needle biopsy assess cancer risk; rapid growth, hoarseness, swallowing difficulty, or breathing problems needs prompt review.

\synonyms
Nodule Thyroid (Thyroid Nodules)

\medterm{Thyroid Papillary Carcinoma} The most common differentiated thyroid cancer, arising from thyroid follicular cells. It often
grows slowly and has a favorable prognosis, although outcome depends on tumor features,
spread, age, and response to treatment.

\medterm{Thyroid Peroxidase Antibody} An autoantibody against thyroid peroxidase, an enzyme required for thyroid-hormone
synthesis. It supports autoimmune thyroid disease, especially Hashimoto thyroiditis, but
may be present in some euthyroid individuals.

\medterm{Thyroid Peroxidase Test} A blood test measuring antibodies against thyroid peroxidase, an enzyme needed to make thyroid hormones. Elevated levels support autoimmune thyroid disease but do not by themselves show how well the thyroid is working.

\medterm{Thyroid Preparation Overdose} Taking too much thyroid hormone medicine can cause a fast or irregular heartbeat, shaking, sweating, anxiety, vomiting, fever, or, in severe cases, heart failure or a dangerous thyroid crisis.

\medterm{Thyroid Scan} Nuclear medicine imaging using I-123 or Tc-99m. Evaluates thyroid function and nodules.
Hot nodules rarely malignant. Cold nodules require further evaluation.

\medterm{Thyroid Scintigraphy} Thyroid imaging using I-123, I-131, or Tc-99m pertechnetate. Assesses function, nodules,
ectopic tissue. Hot nodules (functioning) rarely malignant; cold nodules require
ultrasound/FNA.

\medterm{Thyroid Stimulating Hormone} TSH (thyrotropin): the anterior pituitary hormone that stimulates the thyroid to produce
and release T4 and T3, controlled by TRH (hypothalamus) and negative feedback from
thyroid hormones.

\medterm{Thyroid Storm} Life-threatening exacerbation of thyrotoxicosis presenting with hyperpyrexia,
tachycardia, heart failure, altered mental status, and multi-organ dysfunction.
Precipitated by infection, surgery, trauma, or iodine load in patients with untreated
hyperthyroidism. Mortality rates range from twenty to thirty percent.

\medterm{Thyroid Surgery} Operations that remove part or all of the thyroid gland. They may be performed for selected cancers, enlarging or
compressive nodules, goiter, hyperthyroidism, or diagnostic uncertainty; the procedure
depends on the disease and its extent.

\medterm{Thyroid Ultrasound} High-frequency sonographic imaging used to evaluate thyroid size, nodules, cysts,
cervical lymph nodes, and structural features associated with malignancy. Suspicious
nodules may require ultrasound-guided fine-needle aspiration.

\medterm{Thyroid, Overactive} Alternate index label; see \textit{Hyperthyroidism (overactive thyroid)}.

\medterm{Thyroid-Stimulating Hormone} A hormone made by the pituitary gland that tells the thyroid to make thyroid hormones. Blood TSH testing is commonly used to check thyroid function, although results must be interpreted with other hormone levels and the person's symptoms.

\synonyms
TSH

\medterm{Thyroidectomy} Surgery to remove part or all of the thyroid gland. It may treat thyroid cancer, an overactive thyroid, or a large goiter; possible complications include bleeding, voice-nerve injury, and low calcium.
\medterm{Thyroiditis} Inflammation of the thyroid gland with various etiologies. Hashimoto thyroiditis:
autoimmune, most common cause of hypothyroidism in iodine-sufficient areas. Subacute (De
Quervain) thyroiditis: post-viral, presents with painful thyroid and transient
hyperthyroidism followed by hypothyroidism.

\medterm{Thyroiditis Classification} Inflammatory conditions of the thyroid classified by clinical presentation and etiology.
Hashimoto thyroiditis is chronic autoimmune with hypothyroidism. Graves disease is
autoimmune with hyperthyroidism. Subacute thyroiditis is painful and likely viral.

\medterm{Thyroiditis, Chronic Lymphocytic} Alternate index label; see \textit{Hashimoto's disease}.

\medterm{Thyrotoxic Myopathy} Muscle weakness caused by an overactive thyroid, usually affecting the hips, thighs, shoulders, and upper arms. Weakness may occur without pain and can improve when thyroid hormone levels return to normal; severe weakness or swallowing or breathing difficulty needs urgent care.

\synonyms
Hyperthyroid Myopathy

\medterm{Thyrotoxic Periodic Paralysis} Recurrent episodes of transient flaccid weakness associated with thyrotoxicosis, usually
accompanied by hypokalemia caused by intracellular potassium shift. Attacks often affect
the proximal limbs and are triggered by rest after exercise or carbohydrate-rich meals.

\medterm{Thyrotoxicosis} A condition caused by too much thyroid hormone activity in the body, whatever the source. It may cause weight loss, fast heartbeat, tremor, sweating, anxiety, diarrhea, heat intolerance, or severe illness.
\medterm{Thyrotroph Adenoma} A thyrotroph adenoma is a pituitary tumor arising from thyroid-stimulating-hormone-producing cells and may cause excess TSH and abnormal thyroid function.

\medterm{Thyrotrophin-Releasing Hormone} A hormone made in the hypothalamus, a part of the brain, that tells the pituitary gland to release thyroid-stimulating hormone and prolactin. It helps link brain signals with thyroid activity and some reproductive functions.
\medterm{Thyrotrophin-Stimulating Hormone} Thyroid-stimulating hormone, a pituitary hormone that stimulates thyroid hormone production.
\medterm{Thyrotropin} Thyroid-stimulating hormone, a pituitary hormone that stimulates thyroid growth and thyroid-hormone production; it is also called TSH.

\medterm{Thyrotropin Alfa} A recombinant form of thyroid-stimulating hormone used to stimulate residual thyroid tissue or differentiated thyroid cancer cells during selected diagnostic or treatment protocols.

\medterm{Thyroxine} The thyroid hormone T4, which contains four iodine atoms. The thyroid releases it into the blood, and much of it is converted to the more active hormone triiodothyronine (T3).

\synonyms
T4

\medterm{Thyroxine-Binding Globulin} The principal transport protein for thyroid hormones in the blood, binding most circulating thyroxine and triiodothyronine.

\medterm{TI-RADS} The Thyroid Imaging Reporting and Data System is an ultrasound-based system that assigns a risk category to thyroid nodules
using features such as composition, echogenicity, shape, margins, and calcification. It
helps guide follow-up and biopsy decisions.

\medterm{TIA} A transient ischemic attack, a temporary episode of neurological dysfunction from reduced brain blood flow.
\synonyms
Transient ischemic attack; mini-stroke

\medterm{Tiazofurin} Tiazofurin is an investigational nucleoside analogue that interferes with NAD synthesis and has been studied as an antineoplastic agent.

\medterm{Tibia} The tibia is the larger medial bone of the lower leg, extending from the knee to the ankle; it bears much of the body's weight and forms parts of both joints.
Proximal: tibial plateau (medial/lateral condyles) and tibial tuberosity (patellar
ligament insertion). Shaft: sharp subcutaneous anterior border. Distal: medial malleolus
articulating with the talus at the ankle.

\medterm{Tibia Vara} Tibia vara is inward angulation of the upper shinbone causing bowing of the leg, classically associated with Blount disease or other growth disorders.

\medterm{Tibial Artery} An artery of the lower leg, chiefly the anterior and posterior tibial arteries, supplying the leg, ankle, and foot and contributing to pedal pulses.

\medterm{Tibial Fracture} A break in the tibia, the major weight-bearing bone of the lower leg, often caused by high-energy trauma, twisting injury, or repetitive stress.

\medterm{Tibial Nerve} Sciatic division (L4-S3) continuing from posterior thigh through the popliteal fossa
into the deep posterior leg, innervating gastrocnemius, soleus, tibialis posterior,
flexor digitorum longus, and flexor hallucis longus. Passes behind the medial malleolus
as the tarsal tunnel, dividing into medial/lateral plantar nerves for the foot.

\medterm{Tibial Shaft Fractures} Breaks in the middle portion of the shinbone, often caused by a high-energy injury but sometimes by a twist or lower-force event. Pain, swelling, deformity, or inability to bear weight are common. Treatment depends on alignment, skin and muscle damage, and whether the fracture is open.

\medterm{Tibialcrural Artery} A lower-leg arterial vessel or branching region supplying the ankle and foot, usually discussed in relation to tibial and fibular arterial anatomy.

\medterm{Tibialis Anterior Muscle} Superficial anterior leg muscle from lateral tibial condyle and shaft to medial
cuneiform and first metatarsal base. Innervated by deep peroneal nerve (L4-L5). Primary
dorsiflexor for foot clearance during gait swing phase; inverts foot. Weakness causes
foot drop with steppage gait.

\medterm{Tibialis Cranialis} The animal-anatomy term corresponding broadly to the tibialis anterior muscle, which dorsiflexes the ankle and helps invert the foot.

\medterm{Tibialis Posterior Muscle} Deepest posterior leg muscle from posterior tibia, fibula, and interosseous membrane to
navicular tuberosity, medial cuneiform, and sustentaculum tali. Innervated by tibial
nerve (L4-L5). Plantarflexes and inverts foot; primary dynamic medial longitudinal arch
stabilizer.

\medterm{Tic} A sudden, rapid, recurrent, nonrhythmic movement or sound that a person may feel an urge to make. Tics can be motor or vocal and may worsen with stress.
\medterm{Tic Disorder} A neurologic disorder characterized by recurrent, sudden, rapid, nonrhythmic motor movements or vocalizations that may be temporarily suppressible.

\medterm{Tic Douloureux} A chronic pain condition characterized by recurrent, sudden, severe, electric shock-like
pain in the distribution of the trigeminal nerve (typically V2/V3). See Trigeminal
Neuralgia.

\medterm{Tic Douloureux (Trigeminal Neuralgia)} Alternate index label for Trigeminal Neuralgia.

\medterm{Ticarcillin} A penicillin antibiotic with activity against many gram-negative bacteria, including Pseudomonas; it is usually given with clavulanate and can cause allergy or electrolyte disturbances.

\medterm{Tick} A small blood-feeding arachnid that can attach to people or animals and transmit infections such as Lyme disease, rickettsial diseases, and babesiosis.

\medterm{Tick Bite} A bite from a tick, which may cause local irritation and can sometimes spread infections such as Lyme disease, ehrlichiosis, or Rocky Mountain spotted fever. Remove the tick promptly and watch for fever or a spreading rash.

\medterm{Tick Paralysis} Acute, progressive flaccid paralysis caused by neurotoxin released by an attached tick.
Weakness usually begins in the legs and ascends; respiratory failure can occur and
symptoms improve after prompt tick removal.

\medterm{Tick Removal} Promptly removing an attached tick with fine-tipped tweezers by pulling steadily close to the skin. Clean the site afterward and monitor for expanding redness, fever, or other illness; do not burn, crush, or cover the tick with chemicals.

\medterm{Tick Typhus} A rickettsial infection transmitted by ticks, usually causing fever, headache, rash, and sometimes severe systemic illness.

\medterm{Tick-Borne Encephalitis} A viral infection transmitted mainly by ticks that can cause fever, meningitis, encephalitis, or lasting neurological complications.

\medterm{Tickborne Relapsing Fever} A vector-borne illness caused by Borrelia hermsii and related species, transmitted by
Ornithodoros ticks. It presents with recurrent episodes of fever, headache, myalgia, and
arthralgia, separated by afebrile intervals. The relapsing pattern results from
antigenic variation of the spirochete's surface lipoproteins.

\medterm{Ticks} Small arachnids that attach to skin and may transmit infections or cause tick paralysis.

\synonyms
Ticks Bite (Ticks)

\medterm{Ticks Bite (Ticks)} Alternate index label for Ticks.

\medterm{Ticlopidine} An antiplatelet medicine that reduces platelet aggregation and clot formation.

\medterm{Tics} Sudden, rapid, recurrent, nonrhythmic stereotyped movements (motor tics) or
vocalizations (vocal tics) that are usually brief and suppressible. Transient tics (most
common in childhood) vs. chronic tic disorder; Tourette syndrome: multiple motor + at
least one vocal tic for >1 year beginning before age 18.

\medterm{Tidal Volume} The volume of air inspired or expired during each normal breath, approximately 500 mL
(6-8 mL/kg) in a resting adult. Tidal volume is measured by spirometry. Minute
ventilation = tidal volume x respiratory rate (~6 L/min). Only a portion reaches the
alveoli (alveolar ventilation) after subtracting dead space (~150 mL).

\medterm{Tidy Wound} A clean, sharply incised wound with minimal contamination, crushing, or tissue devitalization and clearly defined viable edges. After appropriate assessment and cleaning, it may often be closed primarily.
\medterm{Tietze Syndrome (Costochondritis And Tietze Syndrome)} Alternate index label for Costochondritis and Tietze Syndrome.

\medterm{Tigecycline} Tigecycline is a glycylcycline antibiotic that binds the bacterial 30S ribosomal subunit and inhibits protein synthesis; it treats selected complicated infections but commonly causes nausea and has a warning for increased mortality in some severe infections.

\medterm{Tight Junction} A cell-to-cell junction that seals the space between epithelial cells and controls paracellular movement while helping maintain distinct apical and basal cell surfaces.

\medterm{Tilorone} Tilorone is a synthetic antiviral and immune-modulating compound used or studied in some countries; evidence and indications vary by setting.

\medterm{Tilt Table Test} Evaluation of unexplained syncope by reproducing orthostatic stress. Patient is strapped
to a motorized table tilted from supine to 60-80 degrees for 20-45 minutes. Positive:
reproduction of symptoms with hypotension and/or bradycardia (vasovagal syncope, delayed
orthostatic hypotension).

\medterm{Tilt Table Testing} A supervised test in which a patient is moved from a lying to an upright or head-up tilt position while blood pressure, heart rate, and symptoms are monitored. It may help evaluate reflex syncope, orthostatic hypotension, or selected autonomic disorders when the diagnosis remains uncertain; it does not replace assessment for cardiac or neurologic causes.

\medterm{Tilted Disc Syndrome} A congenital appearance in which the optic disc is angled or rotated, often with an unusual shape of nearby retina and blood vessels. Vision may be normal, but some people have visual-field changes or other eye conditions.

\medterm{Time In Range} Percentage of time that continuous glucose monitoring readings remain within the target
glucose range of seventy to one hundred eighty milligrams per deciliter for most adults
with diabetes. General target is greater than seventy percent corresponding to
approximately seventeen hours per day.

\medterm{Time Of Death} The moment or estimated time at which biological death occurred, a key medicolegal
question. Estimated from postmortem changes: body cooling (algor mortis), rigor mortis,
lividity (livor mortis), decomposition/putrefaction, insect colonization, and chemical
changes (vitreous potassium, vitreous glucose, body temperature).

\medterm{Time-Of-Flight PET} A PET acquisition and reconstruction technique that uses the measured arrival-time
difference of coincident annihilation photons to improve localization of the
annihilation event. It can improve signal-to-noise performance, especially in larger
patients.

\medterm{Timed Up-And-Go Test} A functional mobility test measuring the time required to stand from a chair, walk a
short distance, turn, return, and sit. Prolonged performance suggests impaired mobility
and increased fall risk but must be interpreted with gait, vision, cognition, and
environment.

\medterm{TIMI Risk Score} A clinical score estimating short-term risk in people with unstable angina or a non-ST-elevation heart attack. It uses age, risk factors, known artery disease, recent symptoms, ECG changes, and heart-injury blood tests to help guide treatment; it supports but does not replace clinical judgment.

\synonyms
TIMI Score; TIMI Risk Score for UA/NSTEMI; TIMI Score for STEMI

\medterm{Timing Of Breastfeeding} The timing and frequency of breastfeeds, usually guided by early hunger cues and effective milk transfer rather than a rigid clock schedule, with special considerations for newborns and preterm infants.

\medterm{Timolol} Timolol is a non-selective beta blocker used in glaucoma eye drops and for selected cardiovascular conditions; topical ocular use can still cause bradycardia, hypotension, bronchospasm, or fatigue.

\medterm{Timolol Maleate} The maleate salt of timolol, a beta blocker used in glaucoma and selected cardiovascular conditions.
\medterm{Timothy Syndrome} A rare multisystem disorder caused by pathogenic variants in CACNA1C, characterized by
marked QT prolongation and congenital heart disease with syndactyly. Developmental,
neurological, immune, and facial features may also occur.

\medterm{Tincture} Alcohol extract of plant or chemical substance. Common concentration: 10-40\% alcohol.
Examples: tincture of iodine (antiseptic), tincture of benzoin (skin protector). Also
used for liquid medications.

\medterm{Tinea} Dermatophytosis: superficial fungal infection of keratinized tissue (skin, hair, nails) by dermatophytes (Trichophyton, Microsporum, Epidermophyton).
\medterm{Tinea (Dermatophytosis)} Superficial fungal infection by dermatophytes. Tinea corporis (ringworm): annular,
scaly, erythematous plaque with central clearing. Tinea pedis (athlete's foot):
interdigital scaling, maceration, and fissuring. Tinea cruris (jock itch): erythematous,
pruritic rash in groin with raised border.

\synonyms
Dermatophytosis

\medterm{Tinea Barbae} Dermatophyte infection of the beard and mustache area, presenting with erythematous,
scaly, follicular papulopustules or deep inflammatory nodules, often acquired from
contact with infected animals or through grooming. It may be inflammatory (kerion-like)
or superficial, and differential diagnosis includes bacterial folliculitis and acne.

\medterm{Tinea Capitis} A dermatophyte infection of the scalp and hair shafts, most common in children. It may
cause scaling, alopecia, broken hairs, black dots, or an inflammatory kerion.

\medterm{Tinea Circinata} An older term for tinea corporis (ringworm) of the body, presenting as annular,
erythematous, scaly plaques with active, raised borders and central clearing. It is
caused by dermatophyte fungi and may be acquired through direct contact with infected
humans, animals, or fomites.

\medterm{Tinea Corporis} A dermatophyte infection of glabrous skin, presenting as annular, erythematous, scaly
plaques with raised, active borders and central clearing. It is contagious through
direct contact with infected persons, animals, or contaminated objects. Diagnosis is by
KOH preparation showing branching septate hyphae.

\medterm{Tinea Cruris} Alternate index label; see \textit{Jock itch}.

\medterm{Tinea Faciei} Dermatophyte infection of the facial skin, often presenting with erythematous, scaly,
annular or ill-defined plaques that may mimic eczema or rosacea, leading to frequent
misdiagnosis. It is commonly caused by Trichophyton or Microsporum species and may be
acquired from pets, other humans, or self-inoculation from other body sites.

\medterm{Tinea Favosa} A chronic dermatophyte infection of the scalp or skin, usually caused by Trichophyton schoenleinii, producing yellow cup-shaped crusts and possible scarring alopecia.

\medterm{Tinea Imbricata} A chronic dermatophyte infection caused by Trichophyton concentricum, endemic in the
South Pacific and Southeast Asia, characterised by widespread, concentric, overlapping,
polycyclic scaly rings producing a "wood-grain" pattern on the trunk and limbs.

\medterm{Tinea Incognito} A dermatophyte infection whose clinical appearance has been altered by inappropriate
topical corticosteroid or immunosuppressive treatment. It may become less scaly and more
extensive or atypical; fungal testing is often needed before directed therapy.

\medterm{Tinea Nigra} A superficial fungal infection of the stratum corneum caused by Hortaea werneckii,
presenting as a painless, brown to black, non-scaly macule or patch, most commonly on
the palms or soles. It is acquired in tropical and subtropical regions and is often
confused with acral melanoma.

\medterm{Tinea Pedis} A dermatophyte infection of the feet presenting in interdigital, moccasin,
vesiculobullous, or ulcerative patterns. It may coexist with tinea unguium. Diagnosis is
clinical or mycologic when uncertain; treatment uses topical antifungals for limited
disease and oral therapy for extensive, refractory, or nail-associated infection.

\medterm{Tinea Pedis (Athlete'S Foot)} Alternate index label for Athlete's Foot.

\medterm{Tinea Unguium} A fungal infection of the nail plate, also called onychomycosis, causing thickening, discoloration, brittleness, or separation of the nail.

\medterm{Tinea Versicolor} A superficial fungal infection caused by Malassezia species, presenting as hypo- or
hyperpigmented macules with fine scale on the trunk and proximal extremities. It is more
noticeable after sun exposure due to unaffected tanning of surrounding skin.

\synonyms
Pityriasis Versicolor

\medterm{Tinel Sign} Tingling or an electric sensation felt along a nerve when the nerve is gently tapped. It may suggest nerve irritation or recovery after injury, but it is not diagnostic by itself and must be interpreted with the examination and other tests.

\medterm{Tinidazole} Tinidazole is a nitroimidazole antimicrobial effective against anaerobic bacteria and protozoa such as Giardia and Trichomonas; metallic taste, gastrointestinal effects, neuropathy, and an alcohol interaction may occur.

\medterm{Tinnitus} Tinnitus is perception of sound without an external source, often related to hearing loss, noise exposure, ear disease, medications, vascular causes, or neurologic factors.

\synonyms
Ringing In The Ear

\medterm{Tinnitus (Ringing In The Ears)} Alternate index label for Ringing in the Ears.

\medterm{Tinzaparin} A low-molecular-weight heparin anticoagulant used for prevention or treatment of venous thromboembolism; bleeding and heparin-induced thrombocytopenia are important risks.

\medterm{Tinzaparin Sodium} The sodium salt of tinzaparin, a low-molecular-weight heparin used to prevent or treat venous blood clots; bleeding and heparin-induced thrombocytopenia are important risks.

\medterm{Tiopronin} A thiol-containing medicine used in selected cases of cystinuria to increase urinary cystine solubility; rash, proteinuria, and blood abnormalities may occur.

\medterm{Tiotropium Bromide} A long-acting muscarinic antagonist bronchodilator used for maintenance therapy
in COPD. Once-daily dosing provides sustained bronchodilation. It reduces exacerbations
and improves lung function. It is often combined with a LABA for additive
bronchodilation.

\medterm{Tipifarnib} An experimental medicine studied for some cancers and blood disorders. It blocks a cell enzyme needed by certain cancer cells to grow, but it is not a routine treatment and its usefulness depends on the disease and research results.

\medterm{TIPS} TIPS, or transjugular intrahepatic portosystemic shunt, connects portal and hepatic veins to reduce portal hypertension, often for variceal bleeding or refractory ascites.

\medterm{Tirapazamine} Tirapazamine is an investigational prodrug activated in low-oxygen tissue and studied to target hypoxic tumor cells.

\medterm{Tirbanibulin} A topical medicine used in short courses to treat actinic keratoses on the face or scalp; local redness, scaling, or irritation may occur.

\medterm{Tired} Feeling a lack of energy or physical or mental stamina, which may result from sleep loss, exertion, infection, anemia, endocrine disease, medication, depression, or other illness.

\medterm{Tiredness} Fatigue or reduced energy and endurance that can be physical, mental, or both; persistent unexplained fatigue warrants assessment for sleep, mood, systemic, medication, and metabolic causes.

\medterm{Tiredness (Fatigue)} A subjective lack of energy or exhaustion that may result from inadequate sleep, illness, psychological stress, anaemia, or many other causes.

\synonyms
Fatigue

\medterm{Tirofiban} An intravenous glycoprotein IIb/IIIa platelet inhibitor used in selected acute coronary syndromes or coronary interventions to reduce platelet aggregation.

\medterm{Tissierella} A genus of anaerobic bacteria found in the environment and sometimes isolated from polymicrobial human infections.

\medterm{Tissierella Praeacuta} An anaerobic bacterium rarely associated with human infection; significance depends on recovery from a compatible sterile-site infection rather than colonization.

\medterm{Tissue} A group of similar cells and their surrounding materials that work together to perform a
specific function.

\medterm{Tissue Adhesions} Bands of fibrous scar tissue that abnormally connect surfaces or organs, often after surgery, inflammation, infection, or bleeding and potentially causing pain, infertility, or obstruction.

\medterm{Tissue Banks} Facilities that recover, screen, process, preserve, store, and distribute donated human tissue for transplantation, research, or reconstructive treatment under safety and traceability standards.

\medterm{Tissue Congestion} Accumulation of blood or fluid within tissue because of impaired venous drainage, inflammation, obstruction, or increased vascular flow, producing swelling, redness, or impaired function.

\medterm{Tissue Culture} Growth of cells or tissue outside the body in a controlled nutrient medium for research, diagnosis, vaccine production, regenerative medicine, or drug testing.

\medterm{Tissue Distribution} Reversible movement of a chemical from blood into tissues and organs. Distribution is
influenced by perfusion, lipid solubility, protein binding, tissue affinity, barriers,
body composition, and time after exposure.

\medterm{Tissue Elastance} The change in pressure required to produce a given change in tissue or organ volume, the inverse of compliance and important in lung and chest-wall mechanics.

\medterm{Tissue Engineering} The use of cells, supportive materials called scaffolds, and biologically active substances to create or repair tissues that restore, maintain, or improve body function.

\medterm{Tissue Expander} Balloon implanted under skin/muscle, gradually filled with saline to expand tissue. Used
for breast reconstruction, scalp reconstruction, burn coverage. Removed after expansion,
replaced with permanent implant or flap.

\medterm{Tissue Expansion} Gradual stretching of skin or soft tissue with an implanted or external expander so additional tissue can be created for reconstruction; it usually requires staged treatment and monitoring for infection, pain, or device problems.

\medterm{Tissue Fixation} Preserving a tissue sample soon after collection so its cells and structure do not break down before laboratory examination. Chemical fixatives are commonly used, and the method must match the planned tests because poor or delayed fixation can affect the results.

\medterm{Tissue Kallikrein} A serine protease that cleaves kininogen to release kinins such as bradykinin, influencing local blood flow, vascular permeability, pain, and inflammation.

\medterm{Tissue Membrane} A thin sheet or boundary of tissue that separates compartments, covers surfaces, or surrounds structures, with composition and function varying by epithelial, connective, muscle, or biologic membrane type.

\medterm{Tissue Plasminogen Activator} An enzyme that converts plasminogen to plasmin, helping dissolve fibrin in blood clots. Recombinant forms are used as thrombolytic medicines in selected emergencies.

\medterm{Tissue Plasminogen Activator-Associated Angioedema} Angioedema occurring as a complication of tissue plasminogen activator therapy for acute
ischemic stroke or myocardial infarction, affecting approximately 1 to 5 percent of
patients. It typically involves the face, tongue, and upper airway and may be
life-threatening.

\medterm{Tissue Processor} A laboratory instrument that moves tissue through fixation, dehydration, clearing, and paraffin infiltration before embedding and microscopic sectioning.

\medterm{Tissue Scaffold} A natural or engineered three-dimensional framework that supports cell attachment, growth, organization, and tissue repair in regenerative medicine.

\medterm{Tissue Transglutaminase (tTG)} An enzyme that cross-links proteins and serves as the principal autoantigen in celiac disease; IgA antibodies against tissue transglutaminase are commonly used in celiac disease evaluation.

\medterm{Tissue Typing} Testing tissue antigens, especially HLA, to assess compatibility for transplantation.
\medterm{Tissues Of The Body} Groups of similar cells and extracellular material performing a common function: epithelial, connective, muscle, and nervous tissue.

\medterm{Titanium} A strong, lightweight, corrosion-resistant metal used in implants, prostheses, dental devices, and surgical instruments. Bone can form a stable connection with some titanium implant surfaces.

\medterm{Titer} The concentration or dilution level at which a measurable antibody, antigen, or other activity remains detectable in a specimen.
expressed as the highest dilution that still produces a detectable reaction (e.g.,
1:128).

\medterm{Titration} Gradual adjustment of a medicine dose or measurement of a substance using a reagent of known concentration.
\medterm{Titre} A measure of antibody or another substance's concentration, determined by repeatedly diluting a sample until detection stops. Titres can help assess infection, immunity, or response, but interpretation depends on the test and timing.
\medterm{Titubation} A rhythmic nodding or tremulous movement, especially of the head or trunk, often associated with cerebellar disease.
\medterm{Tizanidine} A muscle-relaxing medicine used for painful spasms and spasticity; it can cause sleepiness and low blood pressure.
\medterm{Thin-Layer Chromatography} A laboratory method that separates substances on a thin layer of silica, alumina, or another material as a solvent moves across it. It can help identify compounds, compare samples, or measure radiochemical purity.

\synonyms
TLC

\medterm{Tmax} The time after a dose when a medicine reaches its highest measured concentration in blood or plasma.

\medterm{TMD} Alternate index label; see \textit{TMJ disorders}.

\medterm{TMJ Disorders} Temporomandibular joint disorders affect the jaw joints or chewing muscles and may cause jaw pain, clicking, stiffness, headaches, or limited movement.

\synonyms
Temporomandibular Disorders
Temporomandibular Joint Disorders
TMD

\medterm{TNBC} Alternate index label; see \textit{Triple-negative breast cancer}.

\medterm{TNF Alpha} Tumor necrosis factor alpha is a potent pro-inflammatory cytokine produced primarily by
activated macrophages, as well as by T cells, NK cells, mast cells, and endothelial
cells. TNF-alpha exists as a 26-kilodalton transmembrane protein that is cleaved by the
metalloprotease TACE to release a soluble 17-kilodalton active form.

\medterm{TNM Classification} A cancer-staging system describing the primary tumor's size or spread, nearby lymph-node involvement, and distant metastases. Combined findings help estimate disease extent, guide treatment, compare outcomes, and discuss prognosis.
\medterm{TNM Staging} A cancer-staging system describing the primary tumor, regional lymph nodes, and distant metastases to estimate disease extent and prognosis.

\medterm{Tobacco} A plant product containing nicotine and many toxic combustion or processing chemicals. Its use causes dependence and increases the risk of cancer, cardiovascular disease, and lung disease.
\medterm{Tobacco Chewing (Smokeless Tobacco)} Alternate index label for Smokeless Tobacco.

\medterm{Tobacco Use} Tobacco use exposes the body to nicotine and toxic combustion or aerosol chemicals, causing addiction and increased risks of cancer, vascular disease, and lung disease.

\medterm{Tobramycin} An aminoglycoside antibiotic used for selected serious bacterial infections, including some Gram-negative infections. It can damage the kidneys or hearing and is used with dosing and laboratory monitoring when appropriate.
\medterm{Tocolysis} The use of medication to delay preterm birth for a limited time, often to allow corticosteroid treatment
or transfer to a facility with appropriate neonatal care. The drug, gestational-age range,
and contraindications depend on local guidance and the clinical situation.

\medterm{Tocolytic Agents} Tocolytic agents suppress uterine contractions temporarily to delay preterm birth, often allowing time for corticosteroids or transfer to specialist care.

\medterm{Tocolytic Therapy} Medications used to inhibit uterine contractions temporarily to prolong pregnancy,
allowing time for corticosteroid administration and transport to a tertiary care center.
Indications: preterm labor with intact membranes, 24-34 weeks.

\medterm{Tocolytics} Medications used to suppress preterm labor (contractions before 37 weeks gestation), to
delay delivery for corticosteroid administration (fetal lung maturity) and in utero
transport to a tertiary care center.

\medterm{Tocopherol} A form of vitamin E, a fat-soluble antioxidant involved in protecting cell membranes; deficiency is uncommon but can cause neurologic or hematologic abnormalities.

\medterm{Tocotrienols} Vitamin E compounds with unsaturated side chains found in some plant oils and grains; proposed cardiovascular and metabolic benefits remain under study.

\medterm{Toddler} A toddler is a young child, usually about one to three years old, developing rapidly in movement, language, behavior, feeding, and independence.

\medterm{Toddler Test Or Procedure Preparation} Preparing a toddler for a test or procedure uses simple explanations, play, distraction, comfort objects, caregiver support, and any required fasting or medicine instructions.


\medterm{Toddler'S Fracture} A small spiral crack in the lower shinbone of a young child, usually after a minor twist. The child may suddenly limp or refuse to walk, even though there is little swelling or deformity. The first X-ray may look normal; a repeat X-ray after about two weeks may show healing. Treatment commonly uses a supportive boot or cast, based on examination and imaging.

\synonyms
Cast-Tibia Fracture; Childhood Spiral Tibial Fracture

\medterm{Toe Broken (Broken Toe)} Alternate index label for Broken Toe.

\medterm{Toe Joint} Any joint connecting the bones of a toe. Toe joints include the metatarsophalangeal and interphalangeal joints; they permit toe movement and may be affected by arthritis, gout, injury, or deformity.

\synonyms
toe articulation.

\medterm{Toe Pressure} The blood pressure measured in a toe, useful for assessing circulation in the foot. It can remain reliable when ankle arteries are too stiff to compress, as may occur with diabetes or kidney disease.

\medterm{Toe-Brachial Index} The ratio of systolic toe pressure to brachial systolic pressure. It can support
diagnosis of peripheral arterial disease when ankle arteries are noncompressible because
of medial arterial calcification.

\medterm{Toenail} A toenail is the hard keratin plate covering the end of a toe and protecting its underlying tissue; injury, fungus, and ingrowth can cause disease.

\medterm{Toenail Fungus} Alternate index label; see \textit{Nail fungus}.

\medterm{Toes (Digits Of The Foot)} The five digits at the end of each foot. The big toe has two bones and helps with balance and push-off; the other toes usually have three bones each. Injury, arthritis, tight shoes, or nerve and muscle problems can cause pain or deformity.
\medterm{Togaviridae} A family of viruses that includes some mosquito-borne viruses, such as chikungunya and eastern equine encephalitis virus, which can cause fever, rash, joint pain, or brain inflammation.

\medterm{Togaviridae Infections} Infections caused by viruses in the Togaviridae family. Symptoms vary by virus and may include fever, rash, joint pain, or encephalitis.

\medterm{Toilet} A fixture or place used for urination and defecation; in healthcare, toilet access and toileting assistance are important parts of continence and fall prevention.

\medterm{Toilet Bowl Cleaners And Deodorizers Poisoning} Poisoning from swallowing or contacting toilet cleaners or deodorizers. Acidic or alkaline products can burn the mouth, throat, eyes, and skin; mixing cleaners may release dangerous gases.

\medterm{Tolazamide} An older sulfonylurea diabetes medicine that increases insulin release; hypoglycemia and weight gain can occur, especially in older adults or people with kidney disease.

\medterm{Tolazomide} An older diabetes medicine that lowers blood sugar by stimulating the pancreas to release insulin; now rarely used.
\medterm{Tolbutamide} An older first-generation sulfonylurea that lowers blood glucose by stimulating pancreatic insulin release. It is now rarely used and can cause low blood sugar and other adverse effects.
\medterm{Tolcapone} A catechol-O-methyltransferase inhibitor used with levodopa for Parkinson disease; potentially fatal liver injury limits use and requires monitoring.

\medterm{Tolerable Upper Intake Level} The highest average daily intake of a nutrient unlikely to cause harm in almost all healthy people. It includes food, drinks, and supplements and is not a recommended amount to consume.
\synonyms
UL

\medterm{Tolerance} A decreased response to a drug after repeated administration, requiring dose escalation to achieve the same effect. Unlike tachyphylaxis, tolerance develops over weeks to months of chronic use.
\medterm{Toll-Like Receptor} A class of single-pass membrane-spanning pattern recognition receptors (PRRs) expressed
on innate immune cells (macrophages, dendritic cells). Recognize evolutionary conserved
microbial structures (PAMPs).

\medterm{Toll-Like Receptors} Receptors of the innate immune system that recognize characteristic molecules from microbes or damaged tissue. Their activation starts inflammatory and antimicrobial responses.

\synonyms
TLRs

\medterm{Tolmetin} A nonsteroidal anti-inflammatory drug used for pain and arthritis; gastrointestinal bleeding, kidney injury, and cardiovascular events are important risks.

\medterm{Tolmetin Overdose} Taking too much tolmetin, a nonsteroidal anti-inflammatory medicine, can cause stomach bleeding, vomiting, kidney injury, drowsiness, seizures, or other serious effects.

\medterm{Tolnaftate} A topical antifungal medicine used for some dermatophyte infections such as athlete's foot and ringworm; it is not effective for every type of fungal infection.

\medterm{Tolonium Chloride} Tolonium chloride, or toluidine blue, is a dye used in laboratory and selected clinical procedures to highlight cells or tissue changes.

\medterm{Tolterodine Tartrate} Tolterodine tartrate is an antimuscarinic medicine used for overactive bladder and urinary urgency; dry mouth, constipation, and urinary retention can occur.

\medterm{Toluene} Toluene is a volatile solvent that can irritate tissues and, with substantial inhalation or ingestion, injure the nervous system, liver, kidneys, or heart.

\medterm{Toluene And Xylene Poisoning} Poisoning from solvents containing toluene or xylene, often through swallowing or inhalation. It can cause dizziness, confusion, abnormal heart rhythms, breathing problems, or organ injury.

\medterm{Tolvaptan} A vasopressin V2-receptor antagonist used in selected patients with hyponatremia or autosomal dominant polycystic kidney disease; thirst, rapid sodium correction, and liver injury require attention.

\medterm{Tomogram} A tomogram is an image of a selected cross-sectional layer of the body or another object, produced by tomography such as CT or MRI.

\medterm{Tomography} Imaging that produces cross-sectional views, or slices, through the body. Computed tomography uses X-rays, while other forms use ultrasound, magnetic fields, or radioactive tracers.
\medterm{Tomosynthesis} An imaging test that takes several low-dose X-ray pictures from different angles and combines them into thin layers. It is commonly used to examine breast tissue in three dimensions.

\medterm{Tomotherapy} Tomotherapy combines CT imaging with rotating, precisely shaped radiation beams to deliver radiation treatment while adjusting for the target and nearby tissues.

\medterm{Tongue} A muscular organ involved in taste, speech, chewing, and swallowing.
\medterm{Tongue Biopsy} A tongue biopsy removes a small tissue sample from a tongue lesion for microscopic examination of infection, inflammation, precancer, or cancer.

\medterm{Tongue Carcinoma} Tongue carcinoma is cancer of the tongue, usually squamous-cell cancer, and may cause a persistent sore, lump, pain, or difficulty speaking or swallowing.

\medterm{Tongue Disease} Any disorder affecting the tongue, such as infection, inflammation, injury, nerve disease, nutritional deficiency, or cancer. Symptoms may include pain, sores, swelling, altered taste, or difficulty speaking or swallowing.
\medterm{Tongue Neoplasm} A tongue neoplasm is an abnormal growth in the tongue that may be benign, precancerous, or malignant.

\medterm{Tongue Problems} Changes in the tongue, such as pain, swelling, sores, color change, numbness, or difficulty moving or tasting. Causes include injury, infection, inflammation, nutritional deficiency, and other disease.

\medterm{Tongue Tie} A short or tight band of tissue under the tongue that limits tongue movement. It may interfere with breastfeeding, speech, or eating in some people.

\medterm{Tongue, Disorders Of} Conditions affecting the tongue, including infection, inflammation, injury, tumors, ulcers, movement problems, and developmental differences. They may cause pain, swelling, altered taste, speech difficulty, swallowing problems, or breathing obstruction.
\medterm{Tongue-Tie (Ankyloglossia)} Restricted movement of the tongue caused by an unusually short or tight lingual frenulum, sometimes affecting feeding, speech, or oral hygiene.

\synonyms
Ankyloglossia


\medterm{Tonic Neck Reflex} A primitive reflex present from birth to 6 months, also called fencer's pose. When the
head is turned to one side, the arm on that side extends while the opposite arm flexes.
Present in both directions. Asymmetry may indicate neurologic injury. Disappears as
voluntary head control develops.

\synonyms
fencer's pose.

\medterm{Tonic Pupil} A tonic pupil is a pupil with poor light response but a slow constriction to near focus, often caused by injury to postganglionic parasympathetic fibers.

\medterm{Tonic Seizures} Seizures causing sustained stiffening of muscles, usually for seconds to minutes. They may affect one area or the whole body, cause falls or breathing difficulty, and require individualized seizure evaluation.

\medterm{Tonic-Clonic Seizure} A generalized seizure with a tonic stiffening phase followed by rhythmic clonic jerking,
impaired consciousness, and a postictal state.

\medterm{Tonics} Preparations traditionally claimed to improve strength, energy, appetite, or general health. Their ingredients, safety, and evidence vary widely, and some may interact with medicines or contain harmful amounts of active substances.
\medterm{Tonometer} An instrument that measures pressure inside the eye. Eye-pressure readings help detect glaucoma or monitor treatment, but they must be interpreted with the optic nerve, visual field, cornea, and other findings.

\medterm{Tonometry} Measurement of pressure inside the eye, usually to assess for glaucoma or monitor its treatment. Methods include applanation, noncontact air-puff, rebound, and electronic devices; results must be interpreted with corneal thickness and the clinical context.

\medterm{Tonsil} A small mass of immune tissue at the back of the throat. Tonsils help detect germs entering through the mouth and nose but can become enlarged or infected, causing sore throat, difficulty swallowing, or breathing problems.
\synonyms
Tonsils

\medterm{Tonsil Inflammation (Adenoids And Tonsils)} Alternate index label for Adenoids and Tonsils.


\medterm{Tonsil Stones} Small, hard lumps formed when food particles, mucus, dead cells, and bacteria become trapped and calcify in the folds of the tonsils. They are often harmless but may cause bad breath, throat irritation, or a feeling that something is stuck in the throat. Good oral hygiene and gargling may help; recurrent symptoms, bleeding, or one-sided tonsil enlargement requires examination.

\medterm{Tonsillar Cancer} A malignancy arising from the palatine tonsils, most commonly squamous cell carcinoma
associated with HPV infection. It often presents with throat pain, otalgia, and cervical
lymphadenopathy.

\medterm{Tonsillar Neoplasm} A tonsillar neoplasm is an abnormal growth in a tonsil that may be benign or malignant and can cause a persistent sore throat, swallowing difficulty, or neck swelling.


\medterm{Tonsillectomy} Surgical removal of the palatine tonsils, performed for selected recurrent infection, obstructive sleep-disordered breathing, or other tonsillar disease.
apnea, or other selected conditions after assessment of potential benefits and risks.

\medterm{Tonsillitis} Inflammation or infection of the tonsils, commonly viral or streptococcal.

\medterm{Tonus} The normal resting tension or partial contraction of a muscle or tissue. Abnormally low or high tone can occur with nerve, muscle, brain, spinal-cord, or developmental disorders.
\medterm{Tooth} A hard mineralized structure anchored in the jaw, used for biting, tearing, chewing, and grinding food.
\medterm{Tooth - Abnormal Colors} Abnormal tooth color can result from plaque, decay, trauma, fluorosis, medicines, enamel defects, or staining, and the cause determines treatment.

\medterm{Tooth - Abnormal Shape} Abnormal tooth shape may result from developmental defects, genetic conditions, trauma, wear, or tumors and can affect function, appearance, and eruption.

\medterm{Tooth Abrasion} Loss of tooth substance caused by external mechanical friction, such as aggressive brushing, abrasive products, or habits, producing grooves, notches, or sensitivity.

\medterm{Tooth Abscess} A tooth abscess is a localized collection of pus from bacterial infection of the tooth or surrounding tissues, causing pain, swelling, and possible spreading infection.

\medterm{Tooth Anatomy} Tooth anatomy includes enamel, dentin, pulp, cementum, root, crown, periodontal ligament, and supporting bone and gums.

\medterm{Tooth Ankylosis} Fusion of a tooth root to alveolar bone through loss of the periodontal ligament, preventing normal eruption and making the tooth appear infraoccluded as neighboring teeth grow.

\medterm{Tooth Attrition} Wear of tooth surfaces caused by tooth-to-tooth contact, commonly from chewing, bruxism, or malocclusion, producing flattened cusps and possible sensitivity.

\medterm{Tooth Calcification} Deposition of mineral that forms and hardens enamel, dentin, cementum, and other tooth structures during development; abnormal timing or amount can affect tooth strength and appearance.

\medterm{Tooth Component} Tooth components include enamel, dentin, cementum, pulp, and supporting tissues, each contributing to tooth strength, sensation, nutrition, and attachment.

\medterm{Tooth Crowding} Tooth crowding occurs when teeth lack enough space to align normally, often from tooth size, jaw size, eruption pattern, or retained primary teeth.

\medterm{Tooth Crown} The tooth crown is the visible portion above the gum, covered mainly by hard enamel and containing underlying dentin and pulp.

\medterm{Tooth Decay} Gradual destruction of tooth structure caused when mouth bacteria make acids from sugars. It can cause sensitivity, pain, holes, infection, and tooth loss, but fluoride, cleaning, diet, and dental care help prevent progression.

\synonyms
Dental caries; cavities

\medterm{Tooth Decay - Early Childhood} Early-childhood tooth decay is bacterial acid destruction of primary teeth, promoted by frequent sugars, poor brushing, enamel defects, and prolonged bottle or nursing exposure.

\medterm{Tooth Discoloration} Change in tooth color caused by surface stains, dental decay, trauma, medications, fluorosis, aging, or changes within enamel or dentin.

\medterm{Tooth Disease} Any disorder affecting a tooth, including decay, infection, cracks, developmental defects, or wear. It may cause pain, sensitivity, discoloration, swelling, or difficulty chewing and should be assessed by a dental professional.
\medterm{Tooth Enamel} Tooth enamel is the highly mineralized outer covering of the crown that protects dentin and pulp from wear, acid, and bacteria.

\medterm{Tooth Eruption} Process by which tooth moves from intraosseous position to functional occlusion. Primary
(deciduous) dentition: 20 teeth, eruption 6-30 months. Permanent dentition: 32 teeth
(including third molars), eruption 6-21 years. Delayed eruption may indicate systemic
conditions, local obstruction, or normal variant.

\medterm{Tooth Extraction} Removal of a tooth from its socket because of severe decay, infection, periodontal disease, fracture, impaction, orthodontic planning, or another clinical indication.

\medterm{Tooth Formation - Delayed Or Absent} Delayed or absent tooth formation means teeth develop or erupt later than expected or fail to form, due to local, genetic, endocrine, nutritional, or systemic causes.

\medterm{Tooth Fracture} A crack or break in a tooth caused by injury, biting hard material, decay, grinding, or weakened enamel. It may expose the inner nerve, causing pain, sensitivity, infection, or loss of the tooth.
\medterm{Tooth Infection} Infection of the tooth, its pulp, or surrounding tissues, commonly caused by untreated decay, cracks, or gum disease.
It may cause toothache, sensitivity, swelling, fever, or drainage. Facial swelling, difficulty swallowing or
breathing, or spreading infection requires emergency assessment.

\medterm{Tooth Loss} Tooth loss is absence of a tooth due to decay, gum disease, trauma, developmental absence, or extraction and can affect chewing, speech, and oral health.

\medterm{Tooth Luxation} Traumatic displacement of a tooth without complete avulsion. Concussion and subluxation
involve periodontal injury with little or no displacement, whereas extrusive, lateral,
and intrusive luxation involve increasingly marked displacement and require prompt
dental evaluation.

\medterm{Tooth Replantation} Replacement of an accidentally displaced or avulsed tooth into its socket, followed by stabilization
and dental follow-up.

\medterm{Tooth Resorption} Loss of dental hard tissue caused by odontoclast activity, occurring internally or on the root surface and associated with trauma, infection, orthodontics, or idiopathic disease.

\medterm{Tooth Root} The tooth root anchors a tooth in its socket and contains dentin, cementum, pulp, nerves, and blood vessels beneath the gum.

\medterm{Tooth Sensibility Testing} Clinical testing of a tooth's response to thermal, electric, or dentin stimulation. It
assesses neural response rather than directly measuring pulp blood flow, so results must
be interpreted with history, examination, and imaging.

\medterm{Tooth Sensitivity} Pain in response to thermal, sweet, acidic, or tactile stimulus. Caused by dentin
exposure (recession, erosion, abrasion, abfraction). Management: desensitizing
toothpaste, fluoride varnish, restorative coverage, gum graft for recession.

\medterm{Tooth, Supernumerary} An extra tooth beyond the normal number that may crowd neighboring teeth, remain trapped, or interfere with normal eruption.
\medterm{Toothache} Pain arising from a tooth or nearby gum, ligament, bone, or other supporting tissue. Common causes include decay, infection, cracks, trauma, and gum disease.
\medterm{Toothache Overview} Toothache usually reflects decay, pulp inflammation, fracture, gum disease, abscess, trauma, or referred pain. Dental treatment is needed to correct the cause; facial swelling, fever, difficulty swallowing, or breathing difficulty requires urgent care.

\synonyms
Pain Tooth (Toothache Overview)

\medterm{Toothaches} Toothache is pain arising from decay, pulp inflammation, abscess, gum disease, trauma, cracking, or referred facial pain and requires dental assessment if persistent.

\medterm{Toothpaste Overdose} Swallowing too much toothpaste can cause stomach upset, vomiting, or diarrhea. Large amounts of fluoride toothpaste can cause serious salt and calcium disturbances, especially in children.

\medterm{Tophaceous Gout} Chronic tophaceous gout represents advanced gout with deposition of large, visible tophi
consisting of monosodium urate crystals surrounded by granulomatous inflammation. Tophi
typically develop after years of inadequately treated gout and occur on the ears,
fingers, toes, and olecranon bursa.

\medterm{Tophus} A deposit of monosodium urate crystals in gout, often around joints or in soft tissues.
\medterm{Topical} Applied directly to skin or a moist lining, usually to act at the treatment site with limited whole-body exposure.

\medterm{Topical Antibiotics} Antibiotics applied to the skin to treat selected superficial bacterial infections. The choice and
duration depend on the organism, site, extent of infection, resistance patterns, and local
guidance. Unnecessary or prolonged use can cause irritation and antimicrobial resistance.

\medterm{Topical Antifungals} Topical antifungals treat superficial fungal infections. Allylamines (terbinafine,
naftifine) inhibit squalene epoxidase, fungicidal. Azoles (clotrimazole, miconazole,
ketoconazole, econazole) inhibit ergosterol synthesis, fungistatic. Ciclopirox has
broad-spectrum activity. Tolnaftate is effective against dermatophytes.

\medterm{Topical Calcineurin Inhibitors} Non-steroidal anti-inflammatory agents used for atopic dermatitis and facial eczema.
They inhibit T-cell activation by blocking calcineurin, preventing cytokine production.
Advantages over corticosteroids: no skin atrophy, safe for face and intertriginous
areas. Side effects include transient burning and itching at application site.

\medterm{Topical Corticosteroid Potency} The strength of a steroid cream, ointment, lotion, or gel used on the skin. Mild products are used for sensitive areas or children, while stronger products may be needed briefly for thick or severe inflammation. Choice depends on the diagnosis, body site, age, duration, and risk of skin thinning or other effects.

\medterm{Topical Retinoids} Vitamin A derivatives used for acne, photoaging, and hyperpigmentation. Tretinoin
(all-trans retinoic acid) is the most studied. Adapalene is more stable and less
irritating. Tazarotene is the most potent. Mechanism: normalizes follicular
keratinization, increases cell turnover, and stimulates collagen production.

\medterm{Topiramate} Topiramate is an antiseizure medicine also used to prevent migraine; cognitive slowing, paresthesias, weight loss, kidney stones, metabolic acidosis, and fetal oral-cleft risk are important considerations.

\medterm{Topoisomerase} An enzyme that temporarily cuts and rejoins DNA to control its twisting during replication, transcription, and chromosome separation.

\medterm{Topoisomerase Ii} An enzyme that cuts both DNA strands, passes another DNA segment through the break, and rejoins the strands; it is targeted by some anticancer drugs.

\medterm{Topoisomerase Inhibitors} Drugs that interfere with DNA topoisomerases; some are antibiotics and others are anticancer medicines, with toxicity depending on the drug and target enzyme.

\medterm{Topotecan} A topoisomerase-I inhibitor used for selected ovarian, small-cell lung, and cervical cancers; bone-marrow suppression is a major toxicity.

\medterm{Topotecan Hydrochloride} A salt form of topotecan, a chemotherapy medicine that interferes with DNA copying. It treats selected cancers but can cause dangerously low blood counts, infection, bleeding, diarrhea, nausea, or fatigue.

\medterm{TORCH Panel} Serologic screening for congenital infections: Toxoplasma, Rubella, CMV, Herpes simplex
(and others: syphilis, parvovirus B19, VZV, Zika). IgM: indicates recent/active
infection (maternal or fetal). IgG: past infection or passive immunity. Paired maternal
and neonatal testing.

\medterm{TORCH Screen} Testing for selected congenital infections—classically Toxoplasma, Other infections, Rubella, Cytomegalovirus, and Herpes—when fetal or newborn findings suggest an in-utero infection; broad panels require careful interpretation.

\medterm{Toremifene} A selective estrogen-receptor modulator used for metastatic breast cancer in postmenopausal women; hot flashes and, rarely, dangerous QT prolongation may occur.

\medterm{Torn ACL} A torn anterior cruciate ligament is a knee injury commonly caused by sudden pivoting, landing, or contact, producing a pop, swelling, instability, and difficulty changing direction. Rehabilitation, bracing, or reconstruction is selected according to activity, associated injury, and goals.

\medterm{Torn Meniscus} A torn meniscus is a split in the knee's shock-absorbing cartilage caused by twisting or degeneration. Pain, swelling, catching, or locking may occur; physical therapy or arthroscopic treatment depends on tear pattern, symptoms, age, and associated disease.

\synonyms
Meniscus Tear

\medterm{Torpor} A state of greatly reduced activity, alertness, or responsiveness. In a patient, marked or new torpor may signal a serious illness, medicine effect, or metabolic problem.
\medterm{Torsade De Pointes} Torsade de pointes is a potentially fatal polymorphic ventricular tachycardia associated with prolonged QT, electrolyte disturbance, medicines, or inherited channel disorders.

\medterm{Torsades De Pointes} A specific form of polymorphic ventricular tachycardia with QRS complexes that twist
around the baseline, occurring in the setting of QT prolongation. It may degenerate to
ventricular fibrillation.

\medterm{Torsemide} Torsemide is a loop diuretic that inhibits sodium-potassium-chloride reabsorption in the thick ascending limb and treats oedema from heart, kidney, or liver disease; dehydration and electrolyte loss may occur.

\medterm{Torsion} Twisting of a body structure around its supporting axis, which may reduce or stop blood flow. Torsion of the testicle, ovary, or bowel can destroy tissue and often requires emergency surgery.
\medterm{Torsion Dystonia} A movement disorder causing sustained or repetitive muscle contractions that twist the body or limbs into abnormal postures; it may be inherited, acquired, focal, or generalized.

\medterm{Torsion Of The Appendix Testis} Twisting of a tiny leftover structure attached to the upper part of the testicle. It is common in boys before puberty and causes sudden, localized scrotal pain and tenderness; a small blue spot may be visible through the skin. The testicle usually still has normal blood flow. Rest, support, and pain medicine usually allow it to heal, but urgent assessment is needed to exclude testicular torsion.

\synonyms
Hydatid of Morgagni Torsion; Appendix Testis Torsion

\medterm{Torsion Testicle (Testicular Disorders)} Alternate index label for Testicular Disorders.

\medterm{Torsion, Testicular} Alternate index label; see \textit{Testicular torsion}.

\medterm{Torticollis} An abnormal neck position caused by muscle tightening, pain, nerve problems, injury, or birth-related conditions. The head may tilt or rotate, and treatment depends on the cause, age, and severity.
\medterm{Torticollis (Wryneck)} A condition in which the head is tilted to one side (ear toward the shoulder) and
rotated to the opposite side (chin toward opposite shoulder), due to contraction or
shortening of the sternocleidomastoid muscle.

\synonyms
Wryneck

\medterm{Torticollis, Spasmodic} Alternate index label; see \textit{Cervical dystonia}.

\medterm{Torulaspora} A genus of yeasts used in food and fermentation; human infection is rare and usually opportunistic.

\medterm{Torulaspora Delbrueckii} A yeast used in wine and food fermentation; it rarely causes infection in people with major immune compromise.

\medterm{Torus Fracture} A stable incomplete break in a child's bone caused by compression, often in the wrist after a fall. The bone bulges rather than breaking all the way through; it usually heals well with a splint or cast and follow-up as advised.

\synonyms
Buckle Fracture; Metaphyseal Buckle Fracture; Torus Fracture of Radius

\medterm{Torus Mandibularis} A bony outgrowth arising from the lingual surface of the mandible, typically bilateral
and located in the premolar region above the mylohyoid ridge. It is a common,
asymptomatic developmental anomaly, more prevalent in women and certain ethnic
populations.

\medterm{Torus Palatinus} A bony exostosis arising from the midline of the hard palate, typically presenting as a
slow-growing, asymptomatic, hard, bony protuberance discovered incidentally. It is more
common in women and certain ethnic groups.

\medterm{Tositumomab} A monoclonal antibody formerly used with radioactive iodine-131 for some relapsed or refractory B-cell non-Hodgkin lymphomas; it is no longer routinely available in many settings.

\medterm{Total Abdominal Colectomy} Surgical removal of the entire colon through an abdominal approach, with restoration of intestinal continuity or creation of an ileostomy as clinically appropriate.

\medterm{Total Anomalous Pulmonary Venous Connection} Alternate index label; see \textit{Total anomalous pulmonary venous return}.

\medterm{Total Anomalous Pulmonary Venous Return} A birth defect in which the veins returning oxygen-rich blood from the lungs connect to the wrong side of the heart. Babies may have blue skin, rapid breathing, poor feeding, or heart failure. Severe obstruction is an emergency, and surgery is needed to reroute the veins.

\medterm{Total Artificial Heart} A mechanical device that replaces the pumping function of both ventricles and is used in selected patients with severe biventricular heart failure, usually as a bridge to transplantation. It requires specialized center-based management and carries risks including bleeding, infection, thrombosis, and device malfunction.

\medterm{Total Body Surface Area} The percentage of body surface involved by a burn or other extensive skin process. It is
estimated using methods such as the rule of nines, Lund-Browder chart, or the patient's
palm and is used to guide burn severity assessment and resuscitation.

\medterm{Total Breech Extraction} A maneuver where the operator grasps the fetal feet and delivers the entire baby in
breech presentation. The most common indication is delivery of the second twin
(malpresentation). Contraindicated for term singleton breech due to high risk of head
entrapment, nuchal arms, and birth trauma.

\medterm{Total Cholesterol} Total cholesterol is the combined cholesterol carried in major blood lipoproteins; cardiovascular risk is assessed with HDL, LDL, triglycerides, and other factors.


\medterm{Total Daily Energy Expenditure} Total calories burned in 24 hours: BMR + thermic effect of food + physical activity.
Estimation guides dietary recommendations. Activity factor multiplies BMR.

\medterm{Total Haemolytic Complement Test} A blood test measuring the overall activity of the classical complement pathway, commonly reported as CH50. Abnormal results can suggest immune deficiency, autoimmune disease, infection, inflammation, or laboratory interference.

\medterm{Total Hysterectomy} Removal of the uterus and cervix. The fallopian tubes and ovaries are not necessarily removed and should be specified separately, such as salpingectomy or oophorectomy.

\medterm{Total Hip Arthroplasty} Replacement of the hip joint with prosthetic components. Acetabular component: metal
shell + polyethylene liner. Femoral component: metal stem + femoral head (cobalt-chrome
or ceramic). Fixation: cemented (PMMA, for elderly, osteoporotic bone), cementless
(porous-coated, for younger patients, biologic fixation), and hybrid.

\medterm{Total Hip Joint Replacement - Revision} Surgery to remove and replace some or all parts of a previous hip replacement that have loosened, worn out, become infected, dislocated repeatedly, or failed. It is more complex than first-time replacement and may require bone reconstruction.

\medterm{Total Intravenous Anesthesia} General anesthesia given entirely through medicines injected into a vein rather than inhaled anesthetic gases. It allows carefully controlled unconsciousness during surgery, but can lower blood pressure or breathing, so trained anesthesia staff must continuously monitor the patient.

\medterm{Total Iron Binding Capacity} TIBC: a measure of the maximum iron-binding capacity of transferrin in serum, reflecting
total iron-transport capacity. Calculated from transferrin or measured directly.

\medterm{Total Knee Arthroplasty} Replacement of the knee joint with prosthetic components. Femoral component
(cobalt-chrome), tibial component (titanium tray + polyethylene insert), and patellar
component (polyethylene, optional). Fixation: cement (PMMA) most common in older adults;
cementless in younger.

\medterm{Total Knee Joint Replacement - Revision} Surgery to remove and replace parts of a previous knee replacement because of loosening, wear, infection, instability, stiffness, or persistent pain. It may require specialized implants or bone grafting and has higher risks than the first operation.

\medterm{Total Knee Replacement} Surgery that replaces damaged surfaces of the knee with artificial parts to relieve severe pain and improve function. Recovery includes exercises and gradual activity; infection, blood clots, stiffness, loosening, persistent pain, or repeat surgery may occur.

\medterm{Total Laryngectomy} Total laryngectomy removes the entire voice box and creates a permanent opening in the neck for breathing, usually to treat advanced laryngeal cancer.

\medterm{Total Lesion Glycolysis} A measurement from a PET scan that combines the active volume of a tumor with its average tracer uptake. It estimates the tumor’s total metabolic activity and may help assess treatment response or prognosis, but it is not interpreted alone.

\medterm{Total Lung Capacity} The total volume of air in the lungs after maximal inspiration. It is measured by
lung-volume testing, such as body plethysmography, and may be increased in hyperinflation
or decreased in restrictive lung disease.

\medterm{Total Pancreatectomy} Total pancreatectomy removes the entire pancreas and causes permanent loss of insulin and digestive-enzyme production, requiring lifelong replacement treatment.

\medterm{Total Parenteral Nutrition} Intravenous administration of nutrients sufficient to meet essentially all nutritional
requirements when the gastrointestinal tract cannot be used adequately.

\medterm{Total Parenteral Nutrition - Infants} Intravenous nutrition for infants who cannot receive enough nutrients through the gastrointestinal tract; it requires careful monitoring of glucose, electrolytes, liver function, infection, and catheter complications.

\medterm{Total Proctocolectomy And Ileal-Anal Pouch} Total proctocolectomy with an ileal-anal pouch removes the colon and rectum and creates a pouch from the ileum that is connected to the anus.

\medterm{Total Proctocolectomy With Ileostomy} Total proctocolectomy with ileostomy removes the colon and rectum and diverts the end of the small intestine through an abdominal stoma.

\medterm{Total Protein} The combined amount of albumin and other proteins in the blood. A blood test may help assess nutrition, liver or kidney disease, inflammation, or abnormal antibody production.

\medterm{Total Sleep Time} The total amount of time spent asleep during a defined sleep period or sleep study, excluding periods of wakefulness.

\medterm{Touch} The sensory perception of contact, pressure, texture, vibration, or movement across the skin. Signals travel through sensory nerves to the spinal cord and brain.
\medterm{Tourette Disorder} A neurodevelopmental disorder characterized by multiple motor tics and one or more vocal
tics lasting >1 year, with onset before age 18. Motor tics: eye blinking, head jerking,
shoulder shrugging, facial grimacing. Vocal tics: throat clearing, sniffing, grunting,
barking, coprolalia (obscene words, occurs in <10\%).

\medterm{Tourette Syndrome} A neurodevelopmental disorder characterized by multiple motor tics and at least one
vocal tic occurring over at least one year. Tics are sudden, rapid, recurrent,
stereotyped movements or vocalizations. They are often preceded by premonitory urges.
Tics typically worsen in adolescence and may improve in adulthood.

\synonyms
Tourette's

\medterm{Tourette'S} Alternate index label; see \textit{Tourette syndrome}.

\medterm{Tourette'S Syndrome} A neurodevelopmental disorder involving multiple motor tics and at least one vocal tic persisting for more than one year.
\medterm{Tourniquet} A tight band or device placed around a limb to temporarily reduce blood flow. It can control life-threatening bleeding or provide a bloodless surgical field, but prolonged or incorrect use can damage nerves and tissue.
\medterm{Towel Clamp} A surgical clamp used to secure drapes or hold materials in place during an operation. Its sharp points must be applied carefully because they can puncture skin or damage underlying tissue.

\medterm{Townes Syndrome} A rare inherited birth disorder that commonly causes blockage or absence of the anal opening and unusual thumbs. Kidney, heart, hearing, foot, or genital abnormalities may also occur. The specific problems vary, so treatment is planned around each child’s findings.

\medterm{Toxaemia} An obsolete term once used for illness thought to be caused by toxins in the blood. In pregnancy, “toxaemia” was formerly used for pre-eclampsia, which is now diagnosed by high blood pressure and related organ findings.
\medterm{Toxaphene} A persistent organochlorine insecticide mixture that can accumulate in soil, water, and body fat. Significant exposure may overstimulate the nervous system and cause tremors, seizures, liver injury, or other toxic effects.
\medterm{Toxbase} A clinical poison-information service that helps healthcare professionals assess chemical or medicine exposures. It provides evidence-based advice about expected symptoms, necessary tests, treatment, observation, and urgency for particular substances.
\medterm{Toxic} Poisonous or capable of harming living cells, tissues, organs, or the whole body. Toxicity depends on the substance, dose, route, duration, and individual susceptibility.

\medterm{Toxic Dose} The amount of a substance that causes harmful effects. It varies with the substance, route, duration of exposure, age, health, and individual susceptibility.

\medterm{Toxic Effect} Harmful change caused by a chemical, medicine, poison, radiation, or other exposure. The effect may involve one or more organs and depends on the dose, route, duration, metabolism, and individual susceptibility.
\medterm{Toxic Encephalopathy} Brain dysfunction caused by exposure to a toxic substance or medicine. Symptoms may include
confusion, altered behavior, seizures, or reduced consciousness.

\medterm{Toxic Epidermal Necrolysis} A severe, potentially fatal mucocutaneous reaction, usually to a medication,
causing extensive epidermal necrosis and skin detachment. It is on the severe end of the Stevens--Johnson syndrome and
toxic epidermal necrolysis spectrum and requires urgent specialist care.

\synonyms
TEN

\medterm{Toxic Megacolon} Acute life-threatening dilation of the colon associated with systemic toxicity,
representing one of the most serious complications of inflammatory bowel disease,
particularly ulcerative colitis. It can also occur in infectious colitis caused by
Clostridioides difficile and cytomegalovirus.

\medterm{Toxic Metabolic Encephalopathy} Diffuse brain dysfunction caused by systemic metabolic or toxic disturbances, producing confusion, altered alertness, agitation, seizures, or coma.

\medterm{Toxic Nodular Goiter} An enlarged thyroid with nodules that produce excessive thyroid hormone.

\medterm{Toxic Shock Syndrome} A severe toxin-mediated illness, usually caused by Staphylococcus aureus or Streptococcus pyogenes, with fever, rash, hypotension, and organ dysfunction.
\medterm{Toxic Synovitis} Temporary inflammation of the lining of the hip joint, usually in a child after a viral illness. It can cause hip pain and limping and must be distinguished from a serious joint infection.

\medterm{Toxicity} The ability of a substance, treatment, or exposure to cause harmful effects in a living organism. Severity depends on the dose, route, duration, metabolism, target organs, and individual susceptibility; effects may be temporary or permanent.

\medterm{Toxicity Grade} A standardized rating of how severe a harmful effect is, often used for adverse effects of medicines or cancer treatment. Grades generally range from mild to life-threatening and help guide monitoring and management.
\medterm{Toxicity Test} A laboratory or clinical study that examines whether a substance causes harmful effects and identifies affected organs, dose ranges, and exposure duration. Results help estimate safety but cannot predict every person's response.

\medterm{Toxicogenomics} The study of how chemicals or other harmful exposures alter gene activity and cellular pathways. It helps researchers identify mechanisms of toxicity, possible biomarkers, susceptible populations, and potential target organs.
\medterm{Toxicokinetic Description} An account of how a harmful substance enters the body, is absorbed, distributed, changed by metabolism, and eliminated. These processes determine the concentration and duration of exposure at target organs.
\medterm{Toxicokinetics} The study of absorption, distribution, metabolism, and excretion of a toxic substance
and its relationship to concentration over time. Toxicokinetics helps estimate exposure,
timing, accumulation, and expected duration of toxicity.

\medterm{Toxicology} The study of poisons and harmful effects of chemicals, medicines, and other substances. It includes how exposure happens, how the body is affected, how poisoning is recognized, how it can be prevented, and how it is treated.
\medterm{Toxicology Screen} A test of urine, blood, or another sample for medicines, drugs, alcohol, or poisons. Results must be interpreted with the person's symptoms, timing, and prescribed medicines.

\medterm{Toxicology Specimen Preservation} Collection, labeling, storage, and transport of biological specimens to preserve
analytes and prevent contamination or degradation. Appropriate specimens may include
peripheral blood, urine, vitreous humor, bile, liver, hair, or gastric contents
depending on the suspected exposure and timing.

\medterm{Toxidrome} A recognizable pattern of signs and symptoms produced by a class of toxins. Important
patterns include opioid, sympathomimetic, anticholinergic, cholinergic,
sedative-hypnotic, and serotonin toxidromes; treatment combines supportive care,
targeted antidotes when available, and decontamination decisions.

\medterm{Toxin} A poisonous substance made by a living organism, such as a bacterium, plant, animal, or fungus. Toxins can damage specific tissues or disrupt nerve, muscle, or cellular functions.

\synonyms
Biological poison

\medterm{Toxin Ciguatera (Ciguatera Poisoning)} Alternate index label for Ciguatera Poisoning.

\medterm{Toxins} Poisonous substances produced by living organisms, including bacteria, fungi, plants, and animals. They can injure cells or disrupt organ function; the term is usually distinguished from toxicants, which are harmful substances of nonbiological origin.
\medterm{Toxocara} A genus of parasitic roundworms that commonly infect dogs and cats. People acquire infection by swallowing eggs from contaminated soil or food; migrating larvae can cause fever, enlarged organs, or eye or nervous-system disease.
\medterm{Toxocara Canis} A dog roundworm whose eggs can infect people. Larvae migrate through tissues and may cause toxocariasis, with fever, cough, enlarged liver, eye inflammation, or rarely nervous-system involvement.
\medterm{Toxocara Cati} A cat roundworm whose eggs can infect people. The larvae may migrate through organs and cause toxocariasis, including fever, cough, enlarged liver, eye inflammation, or rarely neurologic disease.
\medterm{Toxocariasis} Infection by dog roundworm (Toxocara canis) or cat (T. cati), acquired by ingesting infective eggs from soil contaminated by animal feces (pica in children a risk).
\medterm{Toxoid} A bacterial toxin that has been chemically or otherwise inactivated so it cannot cause its usual harm but can still stimulate protective immunity. Toxoids are used in vaccines against diseases such as tetanus and diphtheria.
\medterm{Toxoid Vaccines} Vaccines that protect against harmful bacterial toxins. These vaccines contain toxins
that have been detoxified and rendered inactive.

\medterm{Toxoplasma} A genus of intracellular parasites; Toxoplasma gondii is the species that infects people. Infection is often mild, but it can seriously affect a fetus or people with weakened immunity, especially the brain and eyes.
\medterm{Toxoplasma Blood Test} A blood test detecting antibodies or genetic material related to Toxoplasma gondii infection, interpreted with symptoms, immune status, pregnancy, and timing of exposure.

\medterm{Toxoplasma Gondii} An obligate intracellular protozoan parasite that causes toxoplasmosis. The definitive
host is the cat (family Felidae), which sheds environmentally resistant oocysts in
feces.

\medterm{Toxoplasmosis} An infection caused by the parasite Toxoplasma gondii, usually mild but potentially serious during pregnancy or weakened immunity.

\synonyms
Chorioretinitis Toxoplasma (Toxoplasmosis)
Meningoencephalitis Toxoplasma (Toxoplasmosis)

\medterm{Toxoplasmosis (Toxo)} Infection with Toxoplasma gondii, often mild in healthy people but potentially severe in pregnancy or in people with weakened immunity.

\medterm{TPA} Tissue plasminogen activator, a natural clot-dissolving protein or a laboratory-made medicine based on it. In selected emergencies it helps break down clots, but it can cause serious bleeding.

\synonyms
Tissue plasminogen activator

\medterm{TPMT Polymorphism} Thiopurine S-methyltransferase is an enzyme that metabolizes thiopurine drugs including
azathioprine, 6-mercaptopurine, and thioguanine, which are used in the treatment of
leukemia, inflammatory bowel disease, and autoimmune conditions.

\medterm{TPN} TPN, or total parenteral nutrition, supplies nutrients intravenously when the gastrointestinal tract cannot be used adequately.

\medterm{Tra Gene} A gene involved in sex determination and sexual differentiation in some organisms, especially insects; its function and sequence vary by species.

\medterm{Trabectedin} Trabectedin is an antineoplastic drug that binds DNA and alters transcription and repair, used for selected advanced soft-tissue sarcomas and ovarian cancer.

\medterm{Trabecula} A small beam, strut, or partition of tissue, especially within spongy bone. Trabeculae form a supportive internal network that helps bones bear forces while keeping them relatively light.
\medterm{Trabecular Bone} Light, honeycomb-like bone found inside many bones and near their ends and joints. It supports the body, contains bone marrow, and changes quickly when weight, hormones, medicines, or diseases such as osteoporosis affect bone strength.

\synonyms
Spongy bone; cancellous bone

\medterm{Trabecular Meshwork} A sieve-like tissue near the front of the eye through which the clear fluid called aqueous humor drains. If the drainage pathway becomes blocked, pressure inside the eye may rise and damage the optic nerve in glaucoma.

\medterm{Trabeculectomy} A conventional glaucoma filtration surgery that creates a new drainage pathway for
aqueous humor by removing a small piece of the trabecular meshwork and creating a
partial-thickness scleral flap under which aqueous humor filters from the anterior
chamber into the subconjunctival space, forming a filtering bleb.

\medterm{Trabeculoplasty} Laser treatment that improves aqueous-humor drainage through the trabecular meshwork.

\medterm{Trabulsiella} A genus of gram-negative bacteria related to Enterobacter that is rarely associated with opportunistic human infection.

\medterm{Trabulsiella Guamensis} A gram-negative bacterium rarely isolated from human clinical specimens; disease significance depends on a compatible infection rather than colonization.

\medterm{Trace Elements} Minerals needed in very small amounts for normal enzyme function, blood formation, thyroid activity, and other body processes.
\medterm{Trace Mineral} A mineral required by the body in very small amounts.

\medterm{Tracer} A detectable substance used to follow a biological, chemical, or physical process. Medical tracers can show blood flow, organ function, metabolism, or disease location through imaging or laboratory testing.
\medterm{Tracer Principle} Administered in amounts too small to produce pharmacological effect. Dose optimized for imaging, not therapeutic.
\medterm{Trachea} The windpipe connecting the larynx to the right and left main bronchi. Cartilage rings help keep it open so air can move to and from the lungs.
\medterm{Tracheal Agenesis With Bronchoesophageal Fistula} A very rare birth defect in which the windpipe is absent or severely underdeveloped and an abnormal passage connects the breathing system with the food pipe. Newborns have extreme breathing difficulty and need immediate specialist care.

\medterm{Tracheal Foreign Body} A foreign object lodged in the trachea, causing severe respiratory distress, stridor,
and complete airway obstruction. It is a medical emergency requiring immediate
bronchoscopic removal. In complete obstruction, emergency cricothyrotomy may be needed.

\medterm{Tracheal Rupture} Tracheal rupture is a tear in the windpipe caused by trauma, instrumentation, surgery, or pressure and can produce air leakage, breathing difficulty, and life-threatening complications.

\medterm{Tracheal Stenosis} Narrowing of the trachea, often a complication of prolonged endotracheal intubation or
tracheostomy. It presents with progressive dyspnea, stridor, and wheezing that may be
mistaken for asthma. Diagnosis is by bronchoscopy or CT. Treatment is balloon dilation,
stenting, or surgical resection.

\medterm{Tracheitis} Inflammation or infection of the windpipe. Mild cases may occur with a viral illness, but bacterial infection can rapidly block a child's airway and cause high fever, a harsh cough, noisy breathing, and severe illness. It requires emergency airway and antibiotic treatment.
\medterm{Trachelectomy} Trachelectomy is surgical removal of the cervix, sometimes with nearby tissue, used in selected cervical cancers or other conditions.

\medterm{Tracheobronchomalacia} Abnormal softening of the trachea and/or bronchi causing excessive airway collapse
during breathing. It may be congenital or acquired. Symptoms include chronic cough,
recurrent infections, and expiratory airflow obstruction. Diagnosis is by dynamic CT or
bronchoscopy. Treatment includes CPAP, stenting, or surgical reconstruction.

\medterm{Tracheobronchomegaly} Abnormal enlargement of the trachea and main bronchi, usually associated with
connective-tissue weakness. It may cause recurrent lower-respiratory infection,
bronchiectasis, ineffective cough, dyspnea, and difficulty clearing airway secretions.

\medterm{Tracheoesophageal Fistula} An abnormal connection between the trachea and esophagus, often associated with
esophageal atresia. Presents with recurrent aspiration, coughing with feedings, and
cyanosis during meals. The H-type fistula may present later in infancy with chronic
aspiration and failure to thrive. Diagnosis by esophagram or bronchoscopy.

\medterm{Tracheoesophageal Fistula And Esophageal Atresia Repair} Surgery in a newborn to close an abnormal connection between the windpipe and food pipe and join separated parts of the food pipe when possible. The operation is urgent; breathing problems, leakage, narrowing, reflux, and feeding difficulties may persist.

\medterm{Tracheomalacia} Softening of the tracheal cartilage causing excessive collapse of the trachea during
expiration. It may be congenital or acquired (prolonged intubation, relapsing
polychondritis). Symptoms include expiratory stridor and difficulty clearing secretions.
Diagnosis is by dynamic CT or bronchoscopy.

\medterm{Tracheomalacia - Acquired} Acquired tracheomalacia is weakening of the tracheal wall after injury, inflammation, compression, or prolonged ventilation, causing airway collapse and breathing symptoms.

\medterm{Tracheomalacia - Congenital} Congenital tracheomalacia is weakness of the infant tracheal cartilage causing dynamic airway collapse, noisy breathing, cough, or recurrent respiratory symptoms.

\medterm{Tracheostomy} A surgically created opening through the neck into the trachea, usually with insertion of a tube.
\medterm{Tracheostomy Care} Care of a breathing tube and neck opening, including cleaning the skin and tube, suctioning mucus when needed, checking airflow, and keeping emergency equipment nearby. Blockage, bleeding, infection, or accidental tube displacement requires prompt attention.

\medterm{Tracheostomy Tube} A tube inserted through a tracheostomy opening to maintain access to the airway and permit breathing,
suction, or mechanical ventilation.

\medterm{Tracheostomy Tube - Eating} Eating and drinking with a tracheostomy tube requires swallowing assessment and individualized precautions because aspiration risk depends on airway protection, secretions, and the underlying illness.

\medterm{Tracheostomy Tube - Speaking} Speech with a tracheostomy tube may be supported by cuff deflation, a speaking valve, or other methods when airflow and airway protection are adequate; specialist assessment is required.

\medterm{Tracheotomy} An operation that creates an opening through the neck into the trachea. The opening is called a tracheostomy.

\medterm{Trachoma} A chronic eye infection caused by Chlamydia trachomatis that can scar the conjunctiva and cornea.

\medterm{Trachyonychia} Rough, thin, ridged nails with many tiny pits, sometimes called twenty-nail dystrophy when most or all nails are affected. It may occur on its own or with psoriasis, alopecia areata, eczema, or lichen planus.
\medterm{Traction} A controlled pulling force used to align or immobilize a fracture, reduce a dislocation, or relieve pressure on tissues. It may be applied manually, with weights, or through a medical device.
\medterm{Tractional Retinal Detachment} Separation of the neurosensory retina caused by contraction of fibrovascular membranes,
commonly associated with proliferative diabetic retinopathy.

\medterm{Tractotomy} Surgical interruption of a nerve tract in the brain or spinal cord. Surgical interruption of a nerve tract, usually to reduce severe pain or other neurologic symptoms.
\medterm{Train-Of-Four Monitoring} A monitoring test used during anesthesia that gives four brief electrical pulses to a peripheral nerve and observes the muscle responses. The pattern and strength of the responses show how much a muscle-relaxing medicine is still working and whether recovery is adequate.

\medterm{Trained Immunity} Innate immune memory: epigenetic/metabolic reprogramming of monocytes/macrophages/NK
cells after initial stimulus (BCG, beta-glucan, Candida, oxidized LDL) -> enhanced
response to secondary challenge (heterologous protection).

\medterm{Trait} A trait is a measurable characteristic of an organism influenced by genes, environment, development, or their interaction.

\medterm{Trajectory} The path a projectile follows in flight and, in wound ballistics, the course it
traverses through tissue from entrance to exit or rest. Internal trajectory is
reconstructed from entrance and exit defects, the organ track, recovered fragments, and
clothing damage; angles disclose shooter position and victim posture.

\medterm{TRALI} Transfusion-Related Acute Lung Injury, a life-threatening complication of blood
transfusion characterized by acute hypoxemia and bilateral pulmonary edema within 6
hours of transfusion, without evidence of volume overload.

\medterm{Tram-Track Appearance} Histological finding on silver stain of MPGN representing GBM duplication. Results from
mesangial cell interposition into capillary wall with new basement membrane formation.
Creates double-contour pattern on Jones methenamine silver stain.

\medterm{Trance} An altered state of awareness with unusual attention, responsiveness, or sense of surroundings. It may occur during hypnosis, seizures, intoxication, sleep-related events, or psychological states and requires context for interpretation.
\medterm{Tranexamic Acid} An antifibrinolytic medicine that reduces bleeding by inhibiting plasminogen activation.
\medterm{Trans Fat} An unsaturated fat containing at least one trans-configured double bond. Industrial trans fats, formerly common in partially hydrogenated oils, can raise LDL cholesterol, lower HDL cholesterol, and increase cardiovascular risk; small amounts also occur naturally in some meat and dairy foods.

\medterm{Trans Fatty Acid} An unsaturated fat with a trans shape, often formed during industrial processing and associated with higher heart disease risk.

\synonyms
Trans fat

\medterm{Transabdominal Cerclage} Placement of a strong band around the upper cervix through abdominal surgery, usually before pregnancy or early in pregnancy, to reduce the risk of premature opening in severe cervical weakness. It requires cesarean birth and carries surgical, bleeding, infection, and pregnancy risks.

\synonyms
TAC; Cervico-Isthmic Cerclage

\medterm{Transabdominal Scan} An ultrasound scan performed with the transducer placed on the maternal abdomen. A full
bladder improves visualization of pelvic structures. Used in early pregnancy to confirm
viability and gestational age, and throughout pregnancy to assess fetal growth, anatomy,
amniotic fluid volume, and placental location.

\medterm{Transaminase, Serum Glutamic Pyruvic} SGPT, now called ALT, is a liver-associated enzyme measured in blood and often elevated when liver cells are injured.

\medterm{Transaminitis} Elevation of blood levels of aminotransferase enzymes, usually alanine aminotransferase and aspartate aminotransferase, indicating possible liver-cell injury but not identifying its cause.

\medterm{Transarterial Chemoembolization} A procedure in which chemotherapy and embolizing particles are infused selectively into
tumor-feeding hepatic arteries, combining regional chemotherapy with ischemic tumor
necrosis; used for unresectable hepatocellular carcinoma.

\medterm{Transbronchial Biopsy} Bronchoscopic sampling of lung parenchyma through the bronchial wall. It may evaluate interstitial or granulomatous disease, infection, transplant rejection, or peripheral lesions; complications include bleeding and pneumothorax.

\medterm{Transcatheter Aortic Valve Replacement} Replacement of a narrowed aortic heart valve using a folded artificial valve delivered through a catheter, usually from an artery in the groin. It avoids open-heart surgery for selected people, but bleeding, stroke, blood-vessel injury, rhythm problems, leakage, or need for a pacemaker can occur.

\synonyms
TAVR; TAVI; Transcatheter Aortic Valve Implantation

\medterm{Transcatheter Edge-To-Edge Repair} A catheter procedure that joins parts of the mitral-valve leaflets to reduce leakage. It may help selected people with severe symptoms who are not good candidates for open surgery; bleeding, infection, stroke, valve narrowing, residual leakage, or device problems are possible.

\synonyms
TEER; MitraClip Procedure; Percutaneous Mitral Valve Repair

\medterm{Transcervical Foley Catheter For Cervical Ripening} A mechanical method for cervical ripening: a Foley catheter (16-30 Fr) is inserted
through the cervix, the balloon is inflated with 30-60 mL of saline, and traction is
applied (or the catheter is taped to the thigh with gentle tension). The balloon applies
direct pressure to the internal os, stimulating endogenous prostaglandin release.

\medterm{Transcervical Resection Of Endometrium} A camera-guided procedure through the cervix that removes or destroys the uterine lining to reduce heavy menstrual bleeding. It is not suitable for future pregnancy and requires careful evaluation beforehand.
\medterm{Transcranial Doppler} Ultrasound measuring blood flow velocity in cerebral arteries. Used for vasospasm detection, stroke risk assessment.
\medterm{Transcranial Doppler Scanning} An ultrasound test measuring blood-flow speed and direction in major arteries inside the skull. It can help assess narrowing, blockage, vasospasm, emboli, or circulation after certain brain injuries.
\synonyms
Transcranial Doppler; TCD

\medterm{Transcranial Doppler Ultrasound} An ultrasound test that measures blood-flow speed and direction in major arteries inside the skull. It can help assess stroke risk, vessel narrowing, bleeding-related spasm, and blood flow changes.

\medterm{Transcription} Synthesis of RNA from DNA template by RNA polymerase. Eukaryotic: RNA Pol II transcribes
mRNA, rRNA, snRNA; RNA Pol I transcribes rRNA; RNA Pol III transcribes tRNA, 5S rRNA.
Steps: initiation (promoter recognition, transcription factors), elongation (5' to 3'),
termination. Promoter elements: TATA box, CAAT box, GC box.

\medterm{Transcription Factor} Protein binding to specific DNA sequences, regulating transcription. Activators enhance
transcription; repressors inhibit. Examples: p53, NF-kB, steroid hormone receptors.

\medterm{Transcutaneous Electrical Nerve Stimulation} TENS: low-voltage electrical pulses delivered through electrodes on the skin to reduce pain for some people. It does not treat every cause of pain and should be used according to clinical guidance.

\synonyms
TENS

\medterm{Transdermal} A route of drug administration delivering drug through intact skin into the systemic
circulation via a patch or gel, bypassing first-pass metabolism and providing steady
plasma concentrations over hours to days. Examples: nitroglycerin, fentanyl, nicotine,
estrogen, clonidine, rotigotine.

\medterm{Transfer Training} Teaching safe transfer techniques: bed to chair, chair to toilet, floor to standing. Use
of transfer boards, lift equipment.

\medterm{Transferase} An enzyme that moves a chemical group from one molecule to another. Transferases help cells make and break down substances, and some transferase levels in blood help identify tissue injury or disease.

\medterm{Transference} A psychoanalytic concept in which feelings, attitudes, and expectations originally
directed toward a significant figure from childhood are unconsciously redirected onto
the therapist. In psychodynamic therapy, transference is a therapeutic tool for
understanding the patient's relational patterns and resolving conflicts.

\medterm{Transferrin} An iron transport protein synthesized by the liver, carrying iron in plasma (2
iron-binding sites). In iron deficiency: transferrin is elevated with low ferritin,
giving high TIBC and low transferrin saturation; in iron overload/inflammation/chronic
disease: transferrin is low (acute phase).

\medterm{Transferrin Saturation} The percentage of transferrin's iron-binding sites occupied by iron. It is calculated from iron and total iron-binding capacity and helps evaluate iron deficiency and iron overload.

\synonyms
TSAT

\medterm{Transfusion} Delivery of donated blood or blood components into a person’s bloodstream. It can replace lost red cells, platelets, or clotting factors but carries risks including reactions, infection, fluid overload, and lung injury.


\medterm{Transfusion Of Neonate} Administration of a blood product to a newborn for clinically significant anemia, bleeding, or
another specific indication. The decision, product characteristics, volume, and rate depend
on the infant’s symptoms, laboratory results, gestational age, and local transfusion
protocol.

\medterm{Transfusion Reaction} An adverse reaction to transfusion of blood products. Acute hemolytic transfusion
reaction: ABO incompatibility, immediate fever/chills, hypotension, back pain, DIC,
renal failure; management: STOP transfusion, maintain BP/urine output, supportive care.

\medterm{Transient Familial Hyperbilirubinemia} Transient familial hyperbilirubinemia is inherited temporary elevation of bilirubin, often from reduced conjugation, causing episodic jaundice without major liver injury.

\medterm{Transient Global Amnesia} A temporary episode of sudden inability to form new memories, with repetitive questioning, that usually resolves within 24 hours.

\synonyms
Amnesia, Transient Global

\medterm{Transient Hypogammaglobulinemia Of Infancy} A temporary delay in normal immunoglobulin production during infancy, causing low antibody levels and increased susceptibility to infections before immune function improves spontaneously.

\synonyms
TGA

\medterm{Transient Ischaemic Attack} A brief episode of neurological dysfunction caused by temporary focal brain, spinal-cord, or retinal ischaemia without persistent tissue infarction.

\synonyms
TIA

\medterm{Transient Ischaemic Attacks Or Episodes} Short episodes of weakness, numbness, speech trouble, vision loss, or other neurological symptoms caused by temporarily reduced blood flow to the brain. Symptoms resolve without lasting brain injury, but a TIA is an emergency warning of possible stroke.
\medterm{Transient Ischemic Attack} A transient ischemic attack is a temporary episode of neurological dysfunction
caused by focal brain, spinal cord, or retinal ischemia without acute infarction.

\medterm{Transient Ischemic Attack (TIA, Mini-Stroke)} Alternate index label for Transient Ischemic Attack (TIA, Mini-Stroke).
\medterm{Transient Pain} Pain that is brief and resolves rather than continuing or recurring persistently. Even brief pain may need assessment if it is severe, recurrent, or associated with warning symptoms.

\medterm{Transient Synovitis} Temporary inflammation of the hip lining, usually in children ages 3 to 8 and often after a viral illness. It causes hip, groin, thigh, or knee pain and a limp, but usually little fever or illness. Rest and pain medicine generally help. Doctors must distinguish it from a dangerous joint infection, especially when fever, severe pain, or refusal to walk is present.

\synonyms
Toxic Synovitis; Irritable Hip; Transient Synovitis of the Hip

\medterm{Transient Tachypnea - Newborn} Temporary rapid breathing in a newborn caused by delayed clearance of fetal lung fluid, usually beginning soon after birth and improving within one to three days with supportive care.

\medterm{Transient Tachypnea Of The Newborn} Temporary rapid breathing after birth, usually caused by delayed clearance of fetal lung fluid. It is more common after caesarean birth and typically improves within one to three days with supportive care.

\medterm{Transillumination} Transillumination shines light through a body part or structure to reveal fluid, cavities, or other differences in tissue density.

\medterm{Transition To Adulthood} The process of moving from child-focused to adult healthcare, education, work, and community services. Planning helps young people develop self-management skills and maintain continuity of care.

\medterm{Transitional Care} Coordinated care provided when a person moves between healthcare settings or levels of care, such
as from hospital to home.

\medterm{Transitional Epithelium} Specialized stretchable epithelium lining the renal pelvis, ureters, bladder, and proximal urethra, allowing urinary organs to expand while maintaining a protective barrier.

\synonyms
Transjugular Intrahepatic Portosystemic Shunt

\medterm{Transjugular Intrahepatic Portosystemic Shunt} A radiologically created channel between a portal vein branch and a hepatic vein that lowers portal pressure, used for complications such as variceal bleeding or refractory ascites; encephalopathy and shunt dysfunction can occur.

\medterm{Translational Medicine} An interdisciplinary field that turns findings from laboratory, clinical, and community research into better methods of disease prevention, diagnosis, treatment, and public health. It connects research findings with patient care and community health.

\synonyms
Translational Medicines

\medterm{Translation} The cellular process of making a protein from instructions carried by messenger RNA. A ribosome reads the message in order, while transfer molecules bring the matching amino acids. The process starts, builds the chain, and stops at a signal; the finished chain folds into a working protein.

\medterm{Translation Initiation} The first stage of protein synthesis, when a ribosome recognizes messenger RNA, selects the start codon, and assembles the initiator transfer RNA and ribosomal subunits.

\medterm{Translocation} Movement of a chromosome segment to a new location or movement of a molecule, cell, or organism between compartments; chromosomal translocations can create fusion genes or alter gene dosage.
\medterm{Transmembrane Transport} Movement of ions, nutrients, drugs, or other molecules across a cell membrane through channels, carriers, pumps, or vesicles, sometimes requiring energy.

\medterm{Transmission} Passage of an infection, signal, genetic trait, force, or substance from one source to another; in infectious disease, route, dose, and host susceptibility determine risk.

\medterm{Transmural Infarction} A heart attack damaging the full thickness of part of the heart wall, usually from prolonged blocked blood supply.

\medterm{Transplacental} Passing from the pregnant person's blood across the placenta to the fetus, or in the opposite direction. Antibodies, medicines, chemicals, and some infections can cross the placenta, with effects that depend on the substance and timing.

\medterm{Transplant} Transfer of cells, tissue, or an organ from one site or person to another to replace damaged tissue or restore function. Donor matching, rejection prevention, infection risk, and long-term follow-up depend on the type of transplant.

\medterm{Transplant Immunology} The study of how the immune system reacts to donated organs, tissues, or cells. Blood-group and tissue matching, antibody testing, and immunosuppressive medicines help reduce rejection, but they also increase infection risk and require regular monitoring of the transplant and medicine effects.

\medterm{Transplant Rejection} Immune-mediated destruction of a transplanted organ or tissue, classified by timing and
mechanism. Hyperacute: preformed donor-specific antibodies, minutes to hours
post-transplant, immediate graft thrombosis, irreversible; prevention by cross-matching.

\medterm{Transplant Services} The coordinated medical care involved in organ or tissue transplantation, including donor assessment, recipient evaluation, surgery, rejection prevention, infection monitoring, rehabilitation, and long-term follow-up.

\medterm{Transplant Surgery} Transplant surgery involves replacing a failing organ with a healthy donor organ. Kidney
transplant: living donors (laparoscopic nephrectomy, left kidney preferred) or deceased
donors (standard criteria, expanded criteria, DCD).

\medterm{Transplantation} Transfer of an organ, tissue, or cells from a donor to a recipient to replace or restore failing function.

\synonyms
Transplant

\medterm{Transplantation Conditioning} Treatment given before some stem-cell transplants to suppress the recipient's bone marrow and immune system and make room for donated cells. It may use chemotherapy, radiation, or both and can cause infection, bleeding, infertility, or organ injury.

\medterm{Transplantation Genetics} The study and clinical use of inherited and tissue genetic information in transplantation. It can help assess donor compatibility, inherited disease, rejection risk, and the chance that a donor organ carries a genetic condition.

\medterm{Transplantation Immunology} The study of how the immune system recognizes and responds to transplanted cells, tissues, or organs. It explains rejection and tolerance and guides donor matching, antibody testing, and medicines that suppress harmful immune responses.

\medterm{Transport} Regulations for shipping radioactive materials. Package types: excepted, A, B.

\medterm{Transport Protein} A membrane or soluble protein that carries a molecule or ion across a membrane or through body fluids, often determining absorption, distribution, and cellular concentration.

\medterm{Transport Vesicle} A transport vesicle is a membrane-bound carrier that moves proteins, lipids, or other cargo between cellular compartments.

\medterm{Transposition Of Great Vessels} Alternate index label; see \textit{Transposition of the great arteries}.

\medterm{Transposition Of The Great Arteries} A congenital heart defect in which the aorta arises from the right ventricle and the pulmonary
artery from the left ventricle. This creates parallel circulations, so survival after birth
depends on adequate mixing through an associated opening or ductus until definitive
treatment is provided.

\synonyms
Transposition Of Great Vessels

\medterm{Transposons} Mobile genetic elements that can move between different locations within the bacterial
genome or between the chromosome and plasmids. Transposons carry genes including
antibiotic resistance genes and can facilitate the spread of resistance.

\medterm{Transrectal Ultrasonography} Ultrasound imaging of the prostate and nearby structures using a probe placed in the rectum. It can guide biopsies and assess size or abnormalities, but may cause discomfort, bleeding, or infection.

\synonyms
TRUS

\medterm{Transtracheal} Transtracheal means passing through or relating to the trachea, such as a catheter, oxygen delivery route, or diagnostic procedure.

\medterm{Transthyretin Amyloidosis} A disease in which abnormal transthyretin protein builds up in organs, especially the heart or nerves. It may be inherited or develop with age without an inherited gene change.
\synonyms
ATTR amyloidosis; ATTR-CM; Transthyretin-related amyloidosis

\medterm{Transtympanic} Transtympanic means passing through the tympanic membrane, such as medicine delivered into the middle ear or a procedure crossing the eardrum.

\medterm{Transudate} A transudate is a type of edema fluid with low protein content and few cells, resulting
from imbalances in hydrostatic and oncotic pressures across capillary walls rather than
from inflammation or increased vascular permeability.

\medterm{Transudation} Movement of low-protein fluid through a membrane when pressure or protein balance shifts, often causing certain edema or effusions.
\medterm{Transudative Pleural Effusion} From systemic factors altering hydrostatic/oncotic pressures. Most commonly from heart
failure. Low protein, low LDH. Treatment addresses the underlying cause.

\medterm{Transurethral Incision Of The Prostate} An operation making small cuts in the prostate through an instrument passed in the urethra. It can improve urine flow in selected people with prostate enlargement while preserving most prostate tissue.

\synonyms
TUIP

\medterm{Transurethral Microwave Thermotherapy} A treatment for some people with an enlarged prostate in which a catheter uses microwave heat to shrink or destroy prostate tissue blocking urine flow. It is less commonly used now and can cause pain, bleeding, infection, or temporary urinary problems.

\synonyms
TUMT

\medterm{Transurethral Needle Ablation} A treatment using needles passed through the urethra to deliver radiofrequency heat into enlarged prostate tissue. The heat destroys selected tissue and may improve urine flow, but relief develops gradually and complications can occur.

\synonyms
TUNA

\medterm{Transurethral} Performed through the urethra, usually using an instrument to examine, treat, or remove tissue from the urinary tract.
\medterm{Transurethral Resection} Removal of tissue through an instrument passed along the urethra, commonly for prostate or bladder disease.
\medterm{Transurethral Resection Of The Prostate} An endoscopic operation that removes obstructing prostate tissue to
improve urine flow in selected people with benign prostatic enlargement. Alternatives and
indications depend on symptoms, prostate size, complications, and patient preference.


\medterm{Transvaginal} Performed through the vagina, commonly describing an ultrasound examination or a procedure involving the uterus, cervix, ovaries, or other pelvic structures.
\medterm{Transvaginal Ultrasound} Transvaginal ultrasound uses a covered probe in the vagina to produce images of the uterus, ovaries, cervix, pregnancy, or nearby pelvic structures.

\medterm{Transverse} Running across or perpendicular to the long axis or midline.
\medterm{Transverse Colon} Longest, most mobile colonic segment from hepatic to splenic flexure, approximately 45
cm, with complete mesentery (transverse mesocolon). Crosses anterior to duodenum,
pancreas, and major vessels. Splenic flexure is the highest colon point and common gas
trapping site. Most common diverticular disease site along with the sigmoid.

\medterm{Transverse Lie} A fetal lie where the long axis of the fetus is perpendicular to the long axis of the
mother. The presenting part is a shoulder or arm. Occurs in approximately 1\% of term
pregnancies. Risk of cord prolapse and obstructed labor. Cesarean delivery required as
vaginal delivery is not possible.

\medterm{Transverse Myelitis} An inflammatory disorder of the spinal cord causing acute or subacute motor, sensory,
and autonomic dysfunction.

\medterm{Transverse Plane} A horizontal plane that divides the body or an organ into upper (superior) and lower
(inferior) portions.

\medterm{Transverse Vaginal Septum} A congenital band of tissue crossing the vagina and partly or completely blocking it. A complete blockage can cause absent periods and monthly pelvic pain after puberty; imaging confirms the condition, and surgery may create an open vaginal passage.

\medterm{Transversus Abdominis Muscle} Deepest flat abdominal muscle with horizontal fibers, from inguinal ligament, iliac
crest, thoracolumbar fascia, and inner lower six costal cartilages. Innervated by T7-T12
and L1. Compresses abdomen without trunk movement, providing core stability and
increased intra-abdominal pressure during coughing, lifting, Valsalva.

\medterm{Transvestic Disorder} A paraphilic disorder involving recurrent sexual arousal from cross-dressing that causes marked distress,
impairment, or harm, or is acted on without consent. It is distinct from gender dysphoria,
which concerns gender identity rather than sexual arousal.

\medterm{Tranylcypromine} An antidepressant that increases certain brain chemicals by blocking their breakdown. It can interact dangerously with many medicines and foods high in tyramine, so careful prescribing and a medicine-free waiting period before or after some treatments are required.
\medterm{TRAP Sequence} Twin Reversed Arterial Perfusion sequence, a complication of monochorionic twin
gestations where one twin (the pump twin) perfuses the other (the acardiac twin) via
arterioarterial anastomoses.

\medterm{Trapezium (Multangulum Majus)} A small wrist bone at the base of the thumb. It forms the thumb's highly mobile joint with the hand and can be affected by fractures or arthritis that causes pain during gripping and pinching.
\medterm{Trapezius Muscle} Large triangular muscle from occipital bone and C7-T12 spinous processes to clavicle,
acromion, and scapular spine. Innervated by CN XI and C3-C4. Upper fibers elevate
scapula, middle fibers retract, lower fibers depress. Essential for shoulder stability
and scapulothoracic motion.

\medterm{Trapezoid} Smallest distal row carpal bone between trapezium and capitate, articulating with the
second metacarpal. Wedge shape provides stability for the firmly fixed second
metacarpal. Fractures rare (less than 1 percent of carpal fractures). Contributes to the
radial column wrist stability.

\medterm{Trapezoid Bone} A small distal-row carpal bone between the trapezium and capitate that articulates mainly with the index metacarpal and helps stabilize the radial side of the hand.

\medterm{Trastuzumab} Trastuzumab is a monoclonal antibody targeting HER2 used for HER2-positive breast, gastric, and other cancers; infusion reactions and potentially reversible cardiomyopathy require cardiac monitoring.

\medterm{Trauma} Physical injury caused by an outside force, or the mental response to a deeply distressing event. Physical trauma may damage organs, bones, or blood vessels; psychological trauma can affect emotions, sleep, and daily life.
\medterm{Trauma (Vascular)} Injury to a blood vessel caused by blunt force, penetrating injury, fracture, medical procedures, or
other mechanisms.

\synonyms
Vascular

\medterm{Trauma Center} Hospital designated to provide comprehensive care for injured patients. Levels: Level I
(complete care, research, prevention), Level II (similar but limited resources), Level
III (stabilization, transfer), Level IV (advanced trauma life support, transfer).
Transfer criteria: need for higher level of care, specific injuries.

\medterm{Trauma Surgery} Trauma surgery encompasses the acute care of injured patients, guided by ATLS
principles: Airway maintenance with cervical spine protection, Breathing and
ventilation, Circulation with hemorrhage control, Disability/neurologic assessment, and
Exposure/environmental control.

\medterm{Traumatic Amputation} Traumatic amputation is loss of a body part from injury and requires hemorrhage control, preservation of the amputated part when possible, urgent surgery, rehabilitation, and psychological support.

\medterm{Traumatic Bleb} A filtering bleb formed after glaucoma filtration surgery (trabeculectomy) that develops
complications from trauma, or a thin-walled, elevated bleb that leaks or is at risk of
rupture. Traumatic bleb failure or overfiltering can lead to hypotony, choroidal
detachment, and endophthalmitis risk.

\medterm{Traumatic Brain Injury} Disruption of brain function caused by an external mechanical force. Severity: mild (GCS
13-15, concussion), moderate (GCS 9-12), severe (GCS 3-8). Primary injury: direct damage
at moment of impact (contusion, diffuse axonal injury, intracranial hemorrhage).

\synonyms
TBI

\medterm{Traumatic Events And Children} Traumatic events can affect children’s safety, behavior, sleep, learning, and emotions; supportive caregivers, stable routines, and professional assessment help identify persistent trauma responses.

\medterm{Traumatic Grief} Alternate index label; see \textit{Complicated grief}.

\medterm{Traumatic Injury} Damage to the body caused by an external force, such as a fall, collision, blow, cut, burn, or penetrating object. Injury may affect skin, bones, organs, blood vessels, nerves, or the brain.
\medterm{Traumatic Injury Of The Bladder And Urethra} Damage to the bladder or urethra caused by an accident, surgery, or penetrating injury. Blood in the urine, inability to urinate, pelvic pain, or shock requires urgent evaluation.

\medterm{Traumatic Shock} A life-threatening state in which severe injury causes inadequate blood flow and oxygen delivery to tissues.

\medterm{Traumatic Spondylolisthesis Of The Axis (Hangman’S Fracture)} A bilateral fracture of the pars interarticularis of the C2 vertebra (axis), typically
caused by hyperextension of the head. Classically associated with judicial hanging but
also occurs in motor vehicle accidents and falls. May present with neck pain,
torticollis, or neurological deficits if the spinal cord is compressed.

\synonyms
hangman’s fracture

\medterm{Traumatology} The study and treatment of physical injury and its consequences.
\medterm{Traumeel S} Traumeel S is a multi-ingredient complementary product marketed for pain and inflammation; evidence, formulation, and safety depend on the product and it should not replace needed care.

\medterm{Travel Sickness (Motion Sickness / Kinetosis)} Nausea, dizziness, sweating, or vomiting caused by conflicting signals from the eyes and balance organs during travel or other motion. Looking toward the horizon, fresh air, gradual exposure, and suitable preventive medicines may help; severe or persistent symptoms need assessment.


\medterm{Travelers' Diarrhea} Acute diarrhea occurring during or shortly after travel, usually caused by contaminated food or water and commonly involving bacterial, viral, or parasitic infection.

\medterm{Travoprost} Travoprost is a prostaglandin analogue eye drop that increases uveoscleral outflow and lowers intraocular pressure in glaucoma or ocular hypertension; iris pigmentation and eyelash changes may occur.

\medterm{Trazodone} Trazodone is a serotonin antagonist and reuptake inhibitor used for depression and sometimes insomnia; sedation, orthostatic hypotension, serotonin syndrome, and rare priapism are important adverse effects.

\medterm{Trazodone Overdose} Poisoning from taking too much trazodone. It can cause severe sleepiness, vomiting, low blood pressure, seizures, abnormal heart rhythms, or a prolonged painful erection.

\medterm{Trb Gene} A gene involved in biological processes that vary by species; its precise identity and function must be interpreted from the organism and gene nomenclature used.

\medterm{Trd Gene} The gene encoding the delta chain of the T-cell receptor in humans, contributing to gamma-delta T-cell antigen recognition.

\medterm{Treacher Collins Syndrome} An inherited condition affecting development of the face and skull. It may cause small cheekbones and jaw, downward-slanting eyes, malformed ears, and hearing loss; breathing or feeding may also be difficult. Intelligence is usually normal. Care may include airway support, hearing treatment, and reconstructive surgery.

\medterm{Treacle} A thick sugar syrup with little current medical use. Older medical writing sometimes used the word for a sweetened mixture or supposed antidote; it is not a modern treatment or diagnosis.

\medterm{Treat-To-Target Strategy} A management approach in which treatment is adjusted at regular intervals toward a
predefined goal such as remission or low disease activity. It requires validated
disease-activity measures, shared decisions, monitoring, and timely escalation or
modification.

supportive.

\medterm{Treatment Epoch} A treatment epoch is a defined period in a clinical study or treatment plan during which a particular therapy or treatment phase is given.


\medterm{Treatment-Resistant Depression} Depression that does not improve adequately after appropriate trials of standard
treatments. It requires reassessment of diagnosis, adherence, comorbidity, and treatment
options.

\medterm{Trematoda} A class of parasitic flatworms called flukes; human infections include schistosomiasis and liver or lung fluke disease.

\medterm{Trematosphaeria Grisea} An environmental dematiaceous fungus rarely associated with human infection; clinical importance depends on recovery from a compatible lesion or sterile site.

\medterm{Tremelimumab} A monoclonal antibody that blocks CTLA-4 and enhances T-cell activity; it is used with durvalumab for selected unresectable hepatocellular carcinoma and other cancers, with immune-related adverse effects possible.

\medterm{Tremor} An involuntary, rhythmic shaking of a body part. It may occur at rest, while holding a position, or during movement and can result from medicines, caffeine, anxiety, essential tremor, Parkinson disease, or other neurologic or metabolic conditions.

\medterm{Trench Fever} A louse-borne infection caused by Bartonella quintana, causing recurring fever, headache, weakness, and bone pain.
\medterm{Trench Foot} An injury caused by keeping feet cold and wet for a long time. It can cause numbness, swelling, pain, skin color changes, blisters, infection, and tissue death without gradual warming and medical care.
\medterm{Trench Mouth} Trench mouth is acute necrotizing ulcerative gingivitis causing painful bleeding gums, foul breath, and ulceration, often associated with plaque, smoking, stress, or poor nutrition.

\medterm{Trendelenberg Position} A position in which the patient lies supine with the feet higher than the head; clinical uses and risks depend on context.
\medterm{Trendelenburg} Trendelenburg position places the body supine with the feet higher than the head; its use is procedure-specific and may affect breathing, pressure, and circulation.

\medterm{Trendelenburg Position} A position lying on the back with the head and upper body lower than the feet. It is used for selected procedures but can worsen breathing, pressure in the head, or blood pressure in some patients.
\medterm{Treosulfan} An alkylating chemotherapy drug used in conditioning before allogeneic hematopoietic stem-cell transplantation; infection, mucositis, infertility, and organ toxicities may occur.

\medterm{Trepanning} Drilling or cutting a round opening in a bone, especially the skull. It was used historically for many reasons and is now done only for selected medical procedures, such as relieving pressure, reaching the brain, or treating a bone problem.
\medterm{Trephine} A circular cutting instrument used to remove a disk of tissue or bone, such as during a bone-marrow biopsy or corneal procedure.

\medterm{Trephining} Another term for trepanning, creating a circular opening in bone.
\medterm{Trephining (Trepanation)} A surgical procedure in which a hole is drilled or scraped into the skull (cranium),
exposing the dura mater. Historically, trephining dates back to Neolithic times (10,000
BCE) and was performed for ritual, therapeutic (releasing 'evil spirits'), and
post-traumatic indications.

\synonyms
Trepanation

\medterm{Treponema} A group of spiral-shaped bacteria, including Treponema pallidum, which causes syphilis. Different species can live in people or the environment, and identifying the species helps determine whether disease and treatment are present.
\medterm{Treponema Pallidum} A spirochete that causes syphilis, a sexually transmitted infection. The organism is
transmitted through sexual contact or transplacentally (congenital syphilis). Primary
syphilis presents with a painless genital ulcer (chancre). Secondary syphilis presents
with rash, mucous patches, and condylomata lata.

\medterm{Treponema Pallidum Infection} Infection with the bacterium Treponema pallidum, which causes syphilis and can affect the
skin, nervous system, cardiovascular system, and other organs.

\medterm{Treponemal Infections} Infections caused by spiral-shaped Treponema bacteria. They include syphilis and related diseases and may affect the skin, nerves, heart, eyes, or other organs; diagnosis usually combines clinical findings with blood or lesion tests.
\medterm{Trepopnea} Dyspnea that occurs in one lateral decubitus position but not the other. It is typically
worse when lying on the side of the affected lung. It may be caused by unilateral lung
disease, pleural effusion, or mediastinal mass.

\medterm{Tretinoin} Tretinoin is all-trans retinoic acid used topically for acne and photoaging and orally with arsenic or chemotherapy for acute promyelocytic leukaemia; skin irritation and systemic differentiation syndrome are important risks.

\medterm{TRG Gene} A gene encoding the gamma chain of the T-cell receptor, contributing to gamma-delta T-cell development and antigen recognition.

\medterm{Triiodothyronine} The thyroid hormone T3, which contains three iodine atoms and has stronger biological activity than T4. It helps regulate metabolism, heart rate, temperature, and growth.
\medterm{Triage} Sorting patients by urgency and resource needs so those at greatest risk receive timely care.
\medterm{Trial Indication} The disease, condition, symptom, or clinical purpose for which an investigational intervention is being studied in a clinical trial.

\medterm{Trial Of Labor After Cesarean} An attempt to give birth vaginally after a previous cesarean delivery. Whether it is appropriate
depends on the prior uterine incision, obstetric history, current pregnancy, facility
resources, and the person’s preferences. Counseling includes the possibility of successful
VBAC, repeat cesarean delivery, and uterine rupture.

\medterm{Trial Screening} The process of checking potential participants against a clinical trial’s eligibility criteria, including medical history, examination, laboratory tests, and informed consent.

\medterm{Triamcinolone} A corticosteroid medicine used on skin, by injection, or in other forms to reduce inflammation and itching.
\medterm{Triamcinolone Acetonide} A synthetic corticosteroid used topically, by injection, or in other preparations to reduce inflammation; prolonged or extensive use can cause local and systemic steroid effects.

\medterm{Triamterene} Triamterene is a potassium-sparing diuretic that blocks epithelial sodium channels in the collecting duct and is used for hypertension or oedema; hyperkalaemia and kidney stones are important risks.

\medterm{Triazenes} A group of chemical compounds that includes some anticancer agents; their effects and toxicity depend on the specific compound.

\medterm{Triazine} A chemical ring structure found in some medicines, herbicides, and industrial compounds; health effects depend on the specific substance and exposure.

\medterm{Triaziquone} An older alkylating anticancer drug that cross-links DNA; its clinical use is limited by toxicity and availability.

\medterm{Triazolam} Triazolam is a short-acting benzodiazepine hypnotic for short-term insomnia; amnesia, falls, dependence, respiratory depression, and dangerous interactions with opioids or other sedatives can occur.

\medterm{Triazole} A five-membered nitrogen-containing chemical ring found in several antifungal medicines and other drugs; clinical effects depend on the specific compound.

\medterm{Triazophos} An organophosphate insecticide that blocks acetylcholinesterase, allowing nerve signals to remain active. Significant exposure can cause sweating, vomiting, diarrhea, small pupils, breathing difficulty, weakness, seizures, or life-threatening poisoning.

\medterm{Tributyrin} A triglyceride that releases the short-chain fatty acid butyrate during digestion; it is studied experimentally for intestinal and metabolic effects.

\medterm{Triceps} A muscle with three heads, especially the triceps brachii at the back of the upper arm. It straightens the elbow and helps stabilize the shoulder; injury or nerve problems can weaken pushing and lifting.
\medterm{Triceps Brachii Muscle} Three-headed posterior arm muscle. Long head from infraglenoid tubercle, lateral head from posterior humerus above radial groove, medial head below. All insert on the olecranon. Innervated by radial nerve (C7-C8). Primary elbow extensor; assists shoulder extension/adduction.

\medterm{Trichiasis} Inward misdirection of eyelashes so they rub against the eye.
\medterm{Trichilemmal Carcinoma} A rare malignant follicular tumor showing outer-root-sheath differentiation, usually occurring on sun-damaged skin of older adults as a papule or nodule; complete excision is generally required.

\medterm{Trichilemmoma} A benign follicular tumor with outer-root-sheath differentiation, usually presenting as a small facial papule; multiple lesions may be associated with Cowden syndrome.

\medterm{Trichiniasis} An infection caused by Trichinella larvae, acquired from undercooked meat.
\medterm{Trichinosis} Infection with \textit{Trichinella} worms, usually acquired by eating undercooked infected pork or wild game. It may cause diarrhea, fever, muscle pain, and swelling around the eyes.
\medterm{Trichobezoar} A mass of swallowed hair in the digestive tract, usually the stomach, which may cause pain, vomiting, or blockage.

\medterm{Trichoblastoma} A rare benign follicular tumor of the scalp or face, usually presenting as a solitary dermal or subcutaneous nodule; it can resemble basal-cell carcinoma and is diagnosed histologically.

\medterm{Trichoepithelioma} A rare benign follicular adnexal tumor that usually presents as a skin-coloured papule or nodule on the face or scalp; it can resemble basal-cell carcinoma and may require histopathologic confirmation.
\medterm{Trichomonas} A flagellated protozoan, Trichomonas vaginalis, that infects the urogenital tract and causes trichomoniasis.
It is transmitted mainly through sexual contact and may cause discharge, irritation, or no
symptoms.

\medterm{Trichomonas Vaginalis} A flagellated protozoan that causes trichomoniasis, a sexually transmitted infection.
\medterm{Trichomoniasis} A sexually transmitted infection caused by the parasite \textit{Trichomonas vaginalis}. It may cause genital irritation and discharge or no symptoms; testing and treatment of partners help prevent reinfection.
\medterm{Trichophyton} A group of fungi that grow in keratin and cause ringworm infections of skin, hair, or nails. They spread through contact with infected people, animals, or objects and are treated with antifungal medicines.
\medterm{Trichorrhexis Nodosa} A hair-shaft defect in which weakened points split into brush-like ends, caused by inherited disease, trauma, chemicals, heat, or nutritional or metabolic disorders.

\medterm{Trichothiodystrophy Syndrome} A group of inherited disorders characterized by brittle sulfur-deficient hair with a
distinctive banded appearance, often accompanied by photosensitivity, ichthyosis,
developmental delay, short stature, or recurrent infection.

\medterm{Trichotillomania} Recurrent compulsive pulling out of one's hair, causing hair loss and distress or impairment.
\medterm{Trichuriasis} Intestinal infection caused by the whipworm \textit{Trichuris trichiura}, usually acquired by swallowing eggs from contaminated soil or food. Heavy infection can cause abdominal pain, diarrhea, anemia, or rectal prolapse.
\medterm{Tricuspid Annular Plane Systolic Excursion} A measurement made during an echocardiogram that shows how far part of the tricuspid valve moves as the right ventricle contracts. A low value may indicate reduced right-heart pumping and is interpreted with the complete scan and clinical findings.

\synonyms
TAPSE; Tricuspid Annular Excursion; RV Longitudinal Systolic Function

\medterm{Tricuspid Atresia} Absence of the tricuspid valve, resulting in no direct communication between the right
atrium and right ventricle. Blood flow depends on an atrial septal defect and
ventricular septal defect for mixing. Presents with cyanosis.

\medterm{Tricuspid Incompetence} Incomplete closure of the tricuspid heart valve, allowing blood to leak backward into the right upper chamber. Severe leakage can cause swelling, liver congestion, irregular heartbeat, fatigue, and breathlessness.
\medterm{Tricuspid Regurgitation} Incomplete closure of the tricuspid valve allowing blood to flow backward from the right
ventricle to the right atrium during systole. Causes include right ventricular dilation
(from pulmonary hypertension or left heart failure), endocarditis, rheumatic disease,
and carcinoid syndrome.

\medterm{Tricuspid Stenosis} Narrowing of the tricuspid valve opening, making it harder for blood to move from the right upper chamber to the right lower chamber. It can cause fatigue, swelling, enlarged liver, or neck-vein fullness.
\medterm{Tricuspid Valve} The heart valve between the right upper and right lower chambers. It opens to let blood move forward and closes to prevent backflow; leakage or narrowing can strain the heart.
\medterm{Tricyclic Antidepressant Drugs} Antidepressants that inhibit norepinephrine and serotonin reuptake and can cause anticholinergic, cardiac, and toxic effects.
\medterm{Tricyclic Antidepressants} Older antidepressants that non-selectively inhibit the reuptake of serotonin and
norepinephrine. They include amitriptyline, nortriptyline, imipramine, clomipramine, and
desipramine.

\medterm{Trifluoperazine} A phenothiazine antipsychotic used for selected psychotic disorders. It can cause drowsiness, movement disorders, and other serious effects, so treatment requires medical supervision.
\medterm{Trigeminal Nerve} The fifth cranial nerve, carrying sensation from the face and controlling muscles used for chewing.
\medterm{Trigeminal Neuralgia} Sudden, severe, electric-shock-like pain in one side of the face along one or more branches of the trigeminal nerve. Attacks may be triggered by touching the face, brushing teeth, eating, speaking, or a light breeze.

\synonyms
Neuralgia, Trigeminal
Tic Douloureux (Trigeminal Neuralgia)

\medterm{Trigger} A trigger is an event, exposure, stimulus, or condition that initiates or worsens a symptom, disease episode, behavior, or physiologic response.

\medterm{Trigger Finger} A condition in which a finger tendon catches as it moves through a tight covering. It can cause pain, clicking, stiffness, or locking and may improve with rest, injections, splinting, or surgery.

\synonyms
Stenosing Tenosynovitis
Trigger Thumb (Trigger Finger)

\medterm{Trigger Thumb} Stenosing tenosynovitis causing painful catching or locking of the thumb.
\medterm{Trigger Thumb (Trigger Finger)} Alternate index label for Trigger Finger.


\medterm{Triglyceride} A lipid composed of glycerol and three fatty acids, transported in blood and stored in adipose tissue.
\medterm{Triglyceride Level} A triglyceride level measures a type of fat in the blood; persistently high levels can increase cardiovascular risk and, when very high, cause pancreatitis.

\medterm{Triglycerides} Triglycerides are blood fats made of glycerol and three fatty acids; elevated levels contribute to cardiovascular risk and, when very high, pancreatitis risk.


\medterm{Trigone} A triangular area, especially the smooth region at the base of the urinary bladder between the ureter openings and urine outlet. It helps direct urine flow and can be affected by inflammation or tumors.
\medterm{Trimester} One of the three approximately three-month periods of pregnancy. Each trimester has typical developmental milestones and screening or care needs.
\medterm{Trimethoprim} An antibiotic that blocks bacterial folate production and is often combined with sulfamethoxazole. It treats selected infections but can cause rash, blood-cell changes, high potassium, kidney effects, or medicine interactions.
\medterm{Trimethoprim-Sulfamethoxazole} A fixed-dose combination of trimethoprim and sulfamethoxazole that synergistically
inhibits folate synthesis. TMP-SMX is indicated for urinary tract infections,
Pneumocystis jirovecii pneumonia, and community-acquired pneumonia. The combination is
bactericidal against many gram-positive and gram-negative organisms.

\medterm{Trimethylaminuria} A metabolic disorder in which impaired oxidation of trimethylamine causes a
characteristic fish-like body odor in sweat, urine, and breath. It may be inherited or
acquired and can cause substantial psychosocial distress despite limited physical
illness.

\medterm{Trimipramine} A sedating tricyclic antidepressant used for depression in some countries. Its anticholinergic and cardiovascular effects can cause dry mouth, constipation, blurred vision, drowsiness, orthostatic hypotension, or dangerous toxicity in overdose; interactions and suicide risk require clinical monitoring.
\medterm{Trinitrin} Another name for glyceryl trinitrate or nitroglycerin, a nitrate vasodilator.
\medterm{Trinitrophenol} Picric acid, a corrosive and potentially explosive chemical once used in medicine, laboratories, and industry.
\medterm{Trinucleotide Repeat} Expansion of three-nucleotide sequences causes diseases: Huntington, Fragile X
(CGG), myotonic dystrophy (CTG), Friedreich ataxia (GAA). Anticipation common.

\medterm{Trip} A trip is a journey or an accidental stumble; in drug-related language it can also mean an episode of altered perception after a hallucinogen.

\medterm{Tripe Palms} Diffuse, velvety, thickened hyperkeratosis of the palmar skin with a cerebriform,
corrugated appearance resembling the lining of a cow's stomach (tripe), which is
strongly associated with underlying malignancy, most commonly gastric adenocarcinoma. It
may occur alone or in conjunction with acanthosis nigricans.

\medterm{Triple X Syndrome} A sex-chromosome variation in which a female has an extra X chromosome. Many affected people have typical sexual development and fertility, while some have tall stature, speech or learning difficulties, coordination problems, or anxiety; severity varies widely.

\synonyms
Trisomy X; 47,XXX Syndrome

\medterm{Triple-Negative Breast Cancer} Breast cancer lacking ER, PR, HER2. Aggressive, limited treatment options.

\medterm{Triplets} Three offspring born from the same pregnancy or three infants born in a multiple birth.
\medterm{Triploidy Syndrome} A chromosomal condition in which cells contain three complete haploid chromosome sets.
It causes severe fetal growth restriction and multiple congenital anomalies and is
usually incompatible with prolonged survival; partial or mosaic forms vary.

\medterm{Triptan} A class of medications used for acute migraine treatment, acting as 5-HT1B/1D receptor
agonists. 5-HT1B receptors are located on cranial blood vessels; activation causes
vasoconstriction of dilated meningeal arteries.

\medterm{Triptans (Triptan)} Triptans are serotonin-receptor agonists used to abort migraine attacks by constricting cranial vessels and reducing trigeminal signaling.

\medterm{Triquetrum} Pyramidal proximal row carpal bone on the ulnar side, articulating with lunate, hamate,
pisiform, and TFCC. Second most commonly fractured carpal bone. Contributes to proximal
carpal row and ulnar wrist stability. Pisiform articulates with its volar surface.

\medterm{Trismus} Restricted ability to open the mouth, commonly caused by tetanus, dental infection, or muscle spasm.
\medterm{Trisodium Phosphate Poisoning} Poisoning from swallowing or contacting trisodium phosphate, a strong alkaline cleaner. It can burn the mouth, throat, eyes, and skin and may cause vomiting or breathing problems.

\medterm{Trisomy} Presence of three copies of a chromosome instead of two. Trisomy 21 (Down syndrome):
intellectual disability, characteristic facies, heart defects. Trisomy 18 (Edwards):
severe intellectual disability, multiple anomalies. Trisomy 13 (Patau): severe
intellectual disability, holoprosencephaly, cleft lip/palate.

\medterm{Trisomy 13} Patau syndrome, a chromosomal disorder caused by an extra chromosome 13 and associated with severe developmental, brain, facial, heart, kidney, and limb abnormalities.

\medterm{Trisomy 18} Edwards syndrome, a chromosomal disorder caused by an extra chromosome 18 and associated with growth restriction, characteristic craniofacial and limb findings, heart disease, and severe developmental impairment.

\medterm{Trisomy 18 (Edwards Syndrome)} Alternate index label for Edwards Syndrome.

\medterm{Trisomy 21} Down syndrome: the most common chromosomal disorder, caused by an extra chromosome 21
(trisomy 21; ~95 percent meiotic nondisjunction, mosaicism, translocation).

\medterm{Trisomy 21 (Down Syndrome Overview)} Alternate index label for Down Syndrome Overview.

\medterm{Trisomy 21 Syndrome} Alternate index label; see \textit{Down syndrome}.

\medterm{Trisomy X} Alternate index label; see \textit{Triple X syndrome}.

\medterm{Trocar} A pointed instrument placed inside a hollow tube to enter a body cavity or tissue. It is used in minimally invasive procedures to create an opening for cameras, instruments, or drainage tubes.
\medterm{Trochanter} A bony projection on the thigh bone, especially the larger outer or smaller inner projection. Muscles and tendons attach there, and irritation around the greater trochanter can cause outer-hip pain.
\medterm{Trochanteric Bursitis (Hip Bursitis)} Alternate index label for Hip Bursitis.

\medterm{Troche} A troche is a small medicated lozenge that dissolves slowly in the mouth or throat for local or systemic drug delivery.

\medterm{Trochlear Nerve} Cranial nerve IV, which supplies the superior oblique eye muscle. Injury can cause vertical or tilted double vision, often worse when looking down or inward.
\medterm{Trophic} Relating to the nutrition, growth, development, or maintenance of tissue. Trophic changes may occur when nerve supply or blood flow is impaired.
\medterm{Trophoblast} The outer cell layer of an early embryo that attaches to the uterus and helps form the placenta. It supports implantation, exchanges nutrients, and produces hormones needed to maintain early pregnancy.
\medterm{Trophoderm} The outer cell layer of a blastocyst, the early stage before implantation. It attaches the embryo to the uterus and develops into tissues that help form the placenta and support early pregnancy.

\medterm{Trophozoite} The active, feeding, and multiplying stage of a protozoan parasite.
\medterm{Tropical Chronic Pancreatitis} A historically recognized form of chronic pancreatitis reported mainly in tropical and subtropical regions, often beginning at a younger age and associated with recurrent abdominal pain, pancreatic duct stones, and diabetes. Its causes are multifactorial and may include genetic susceptibility, nutritional factors, and environmental exposures; the term should not be applied solely because a patient lives in a tropical region.

\medterm{Tropical Diseases} Illnesses more common in warm tropical or subtropical regions, including malaria, dengue, schistosomiasis, and other neglected infections.
\medterm{Tropical Medicine} The medical specialty concerned with diseases common in tropical and subtropical regions and illnesses related to travel. It includes infections, parasites, insect-borne diseases, environmental hazards, prevention, diagnosis, and treatment.
\medterm{Tropical Oil} Tropical oils are plant-derived oils from warm regions, such as coconut or palm oil; their nutritional effects depend on the specific oil and amount consumed.

\medterm{Tropical Sprue} A disorder causing long-lasting diarrhea and poor absorption of nutrients, reported mainly in tropical and subtropical areas. It can lead to weight loss and deficiencies of folate, iron, and vitamin B12. Testing helps exclude infection and other causes; treatment usually includes antibiotics and replacement of missing nutrients.

\medterm{Tropical Ulcer} A chronic ulcer occurring in tropical regions, often associated with infection, trauma, or poor wound healing.
\medterm{Troponin} A protein released into the blood when heart muscle is injured. A raised level can support a heart-attack diagnosis, but it may also result from other heart, lung, kidney, or severe systemic illnesses.

\medterm{Troponin Test} A troponin test measures heart-muscle proteins in the blood and helps assess possible heart injury, especially during suspected heart attack.

\medterm{Trouble Swallowing} Difficulty moving food or liquid from the mouth to the stomach. It may cause coughing, choking, pain, or a feeling that food is stuck in the throat or chest.

\medterm{Trousseau Sign Of Latent Tetany} A sensitive and specific clinical sign of hypocalcemia. Elicited by inflating a blood
pressure cuff on the upper arm to a pressure greater than systemic systolic blood
pressure for 3 minutes.

\medterm{Trousseau Syndrome} Recurrent or migratory venous thrombosis associated with active malignancy, particularly
mucin-producing adenocarcinoma. It reflects tumor-associated activation of coagulation
and is distinct from Trousseau sign of latent tetany.

\medterm{Trp} Trp is the standard abbreviation for tryptophan, an essential amino acid used in protein synthesis and serotonin and niacin pathways.

\medterm{True Ribs} Ribs 1-7 articulating directly with the sternum via their own costal cartilages. Each
has a head, neck, tubercle, and shaft. The first two are short and strongly curved for
neck/shoulder muscle attachment. Costal cartilages provide thoracic cage elasticity for
respiratory expansion.

\medterm{Truncation} A shortened or incomplete form caused by removal of part of a sequence, structure, chromosome, protein, or process; the effect depends on what was lost and where.

\medterm{Truncus Arteriosus} A birth defect in which one large blood vessel leaves the heart instead of separate vessels for the body and lungs. A large hole between the lower chambers usually allows oxygen-poor and oxygen-rich blood to mix, causing blue skin and heart failure. Early surgery is usually required.

\medterm{Trunk} The central part of the body excluding the head and limbs, or a major nerve or vessel trunk.
\medterm{Truss} A supportive appliance used historically for a hernia; definitive management depends on the type and complications.
\medterm{Trypanosoma} A genus of protozoan parasites causing African and American trypanosomiasis.
\medterm{Trypanosoma Brucei} A protozoan parasite transmitted by tsetse flies that causes African trypanosomiasis, or sleeping sickness.
The human-infecting subspecies differ in geographic distribution and typical disease
course; infection can progress from fever and lymph-node enlargement to neurologic
involvement.

\medterm{Trypanosomiasis} Disease caused by Trypanosoma parasites, including African sleeping sickness and Chagas disease. Infected insects spread them, and illness may affect blood, lymph nodes, the nervous system, heart, or digestive organs.
\medterm{Trypanosomiasis, American} Alternate index label; see \textit{Chagas disease}.

\medterm{Trypsin} A pancreatic proteolytic enzyme secreted as trypsinogen and activated in the small intestine, where it helps digest dietary proteins and activates other pancreatic enzymes.
\medterm{Trypsin And Chymotrypsin In Stool} Stool trypsin and chymotrypsin testing measures digestive enzymes and may help assess pancreatic enzyme production, although other tests are often preferred.

\medterm{Trypsinogen Test} A trypsinogen test measures a pancreatic enzyme precursor and may help evaluate pancreatic function or newborn screening for cystic fibrosis.

\medterm{Tryptophan} An amino acid the body cannot make and must obtain from food. It is used to make proteins and can be changed into serotonin, melatonin, and niacin, which help regulate mood, sleep, and other body functions.
\medterm{Tsetse Fly} A blood-feeding fly found in parts of sub-Saharan Africa. Its bite can spread Trypanosoma parasites, which cause African sleeping sickness; infection may cause fever, headaches, behavior changes, sleep disturbance, and serious nervous-system disease.
\medterm{TSH} Thyroid-stimulating hormone, a pituitary hormone that stimulates thyroid hormone production; its blood concentration is the main initial test for many thyroid disorders.

\medterm{TSH Test} A TSH test measures thyroid-stimulating hormone in the blood and is a common initial test of thyroid function.

\medterm{TSH-Receptor Antibody} An antibody that binds the thyrotropin receptor. Stimulating antibodies cause Graves
disease, while blocking antibodies can cause hypothyroidism; measurement supports
diagnosis and can help assess fetal or neonatal risk in pregnancy.

\medterm{TSI Test} A TSI test detects thyroid-stimulating immunoglobulins and can support diagnosis or monitoring of Graves disease.
\medterm{Tsutsugamushi Fever} Scrub typhus, a mite-borne infection caused by Orientia tsutsugamushi. It may cause fever, headache, rash, an eschar, and organ complications and is treated with an appropriate antibiotic.
\medterm{Tubal Ligation} Surgical sterilization by occluding the fallopian tubes, preventing the ovum from
reaching the uterus. Postpartum tubal ligation: after vaginal or cesarean delivery
(Pomeroy technique, Parkland). Interval: laparoscopic (bands, clips, bipolar cautery) or
mini-laparotomy (Hulka or Filshie clip).


\medterm{Tubal Ligation Reversal} Microsurgical reconnection of previously divided or blocked fallopian-tube segments to restore a possible path for sperm and egg; success depends on remaining tube length, age, scarring, and original method.

\medterm{Tubal Occlusion Procedure, Selective} Selective tubal occlusion blocks a fallopian tube to prevent sperm and egg from meeting and is a permanent contraceptive procedure.

\medterm{Tubal Pregnancy} An ectopic pregnancy that implants in a fallopian tube instead of the uterus. It can rupture and cause internal bleeding; one-sided pain, vaginal bleeding, dizziness, or fainting needs urgent evaluation.

\medterm{Tuber} A rounded swelling or raised part of a bone or other body structure. In anatomy, tubers provide attachment for muscles or ligaments; the word can also mean a plant storage organ and is not itself a diagnosis.

\medterm{Tubercle} A small rounded projection or nodule; it may refer to bone anatomy or a tuberculous lesion.
\medterm{Tuberculide} A skin eruption caused by an immune reaction to tuberculosis elsewhere in the body. The skin lesions usually contain no bacteria and may improve when the underlying tuberculosis is treated.
\medterm{Tuberculin} Protein material from Mycobacterium tuberculosis used in a skin test for immune sensitization. A positive result suggests exposure or vaccination but does not by itself prove active tuberculosis or show when exposure occurred.
\medterm{Tuberculin Skin Test} Intradermal injection of purified protein derivative to test for Mycobacterium
tuberculosis infection. Read at 48-72 hours measuring induration (not erythema).
Positive: greater than 5 mm in HIV or close contacts; greater than 10 mm in high-risk
groups; greater than 15 mm in low-risk populations.

\medterm{Tuberculin Test} A skin test that checks whether the immune system has reacted to tuberculosis proteins. A positive result suggests infection or previous vaccination and does not by itself prove active disease.

\medterm{Tuberculoid Leprosy} The paucibacillary pole of the leprosy spectrum characterised by strong cell-mediated
immunity to Mycobacterium leprae, resulting in well-formed granulomas and few or no
organisms in tissue.


\medterm{Tuberculosis} An infection caused by Mycobacterium tuberculosis, most often affecting the lungs. It spreads through airborne particles and may cause cough, fever, night sweats, weight loss, or silent infection without symptoms.

\synonyms
TB


\medterm{Tuberculosis In Children} Infection with Mycobacterium tuberculosis, often from household contacts. Latent TB
infection (LTBI) positive tuberculin skin test or interferon-gamma release assay with no
symptoms or radiographic findings. Active TB presents with cough, fever, weight loss,
and lymphadenopathy. Diagnosis by PPD, chest X-ray, and sputum culture.

\medterm{Tuberculosis, MDR} CEA is a blood tumor marker used mainly to help monitor some colorectal and other cancers; it is not specific enough to diagnose cancer alone.
\medterm{Tuberculous Meningitis} The most severe, life-threatening form of extrapulmonary tuberculosis, resulting from hematogenous dissemination of \textit{Mycobacterium tuberculosis} and the rupture of subependymal Rich foci into the subarachnoid space. Causes thick gelatinous basilar exudate, cranial nerve palsies, communicating hydrocephalus, and cerebral vasculitis with stroke.


\medterm{Tuberous Sclerosis} A rare genetic disorder, also called tuberous sclerosis complex, that causes growths in the brain and
other organs, including the skin, kidneys, heart, and lungs. It may cause seizures, developmental or
behavioral problems, and skin changes; the severity varies widely.
\medterm{Tuberous Sclerosis Complex} A genetic condition in which noncancerous growths develop in several organs, especially the brain, skin, kidneys, heart, and lungs. It may cause seizures, developmental or learning difficulties, skin changes, kidney problems, or breathing disease. Lifelong monitoring and treatment are tailored to the organs affected.

\medterm{Tubes} Hollow structures that carry air, fluid, food, or medical instruments. Examples include breathing tubes, fallopian tubes, drainage tubes, and blood vessels; the specific structure must be named for the term to be clinically precise.

\medterm{Tubes Ear Problems (Eustachian Tube Problems)} Alternate index label for Eustachian Tube Problems.

\medterm{Tubo-Ovarian Abscess} A complex infected mass involving the fallopian tube and ovary, usually arising as a complication of pelvic inflammatory disease. It may cause pelvic pain, fever, and sepsis and requires prompt antimicrobial and sometimes drainage treatment.

\medterm{Tubocurarine} A historical muscle-paralyzing medicine derived from curare. It blocks nerve signals at skeletal-muscle junctions, including signals needed for breathing, but does not cause unconsciousness or relieve pain; safer modern drugs have largely replaced it.
\medterm{Tubography} Contrast imaging of the fallopian tubes, usually as part of hysterosalpingography, to see whether they are open. It may help assess infertility but can cause temporary cramping, bleeding, infection, or rarely injury.

\medterm{Tubular Adenoma} A benign glandular polyp, most often in the colon, that can develop dysplastic changes and is considered a precursor lesion for some colorectal cancers.

\medterm{Tubular Reabsorption} The process by which kidney tubules return useful water and dissolved substances from filtered fluid to the blood. This prevents the loss of most water, salt, glucose, and other substances in urine.

\medterm{Tubular Secretion} The transport of substances from the peritubular capillary blood into the tubular lumen,
a second mechanism (besides filtration) by which the kidney eliminates waste and
regulates ions.

\medterm{Tubule} A small tube in the body. In the kidney, tubules change filtered fluid by taking useful water and substances back into the blood and leaving waste to form urine; other tubules carry fluids in different organs.
\medterm{Tubuloglomerular Feedback} A kidney control system that adjusts blood flow and filtration according to the amount of salt reaching a sensing area in the tubule. It helps keep urine production and body fluid balance stable.

\medterm{Tularaemia} A zoonotic infection caused by Francisella tularensis, acquired from animals, ticks, contaminated water, or aerosols.
\medterm{Tularemia} Tularemia is infection with Francisella tularensis acquired from ticks, animals, contaminated water, or aerosols and causing ulcer-glandular, pulmonary, or other disease.

\medterm{Tularemia Blood Test} A tularemia blood test looks for antibodies or other evidence of infection with Francisella tularensis; antibodies may not appear early in illness.

\medterm{Tulle Gras} A paraffin-impregnated gauze dressing used to cover wounds and burns.
\medterm{Tumescence} Swelling or enlargement caused by increased blood flow, fluid, or tissue growth. The term commonly describes erection but can also refer to temporary or abnormal enlargement in other body tissues.
\medterm{Tumor} A tumor is an abnormal mass of cells caused by excessive or uncontrolled growth; it may be benign or malignant.

\medterm{Tumor Adrenal Gland (Pheochromocytoma)} Alternate index label for Pheochromocytoma.

\medterm{Tumor Board} A multidisciplinary conference of oncologists, surgeons, radiologists, and pathologists
who review individual cancer cases to coordinate diagnosis and treatment planning.

\medterm{Tumor Brain Cancer (Brain Cancer)} Alternate index label for Brain Cancer.

\medterm{Tumor Burden} The total amount of cancer in the body, estimated from the number, size, and extent of
tumor deposits; higher burden is generally associated with worse prognosis.

\medterm{Tumor Differentiation} How closely tumor cells resemble the normal cells from which they arose. Well-differentiated tumors look and behave more like normal tissue, while poorly differentiated tumors often behave more aggressively.
\medterm{Tumor Embolism} Tumor embolism occurs when malignant cells, tumor fragments, or tumor thrombi enter the
venous circulation and travel to distant sites. The pulmonary vasculature is the most
common site of embolization, as venous blood from most organs drains to the lungs
through the pulmonary arteries.

\medterm{Tumor Grade} A histologic measure of how much tumor cells resemble normal tissue; low-grade
(well-differentiated) tumors behave less aggressively than high-grade (poorly
differentiated) tumors. Grading is a prognostic, not staging, feature.

\medterm{Tumor Grading} A description of a tumor based on how abnormal the cancer cells look under a microscope
compared to normal cells; a lower grade indicates slower-growing cancer.

\medterm{Tumor Heterogeneity} The presence of genetically, epigenetically, or phenotypically distinct subpopulations
within the same tumor or between primary and metastatic sites. Heterogeneity can
influence prognosis, sampling, treatment response, and resistance.

\medterm{Tumor Lysis Syndrome} A medical emergency caused by rapid cancer-cell breakdown, releasing substances that disturb blood chemistry. It can damage kidneys, trigger seizures or abnormal heart rhythms, and requires prevention, monitoring, and urgent treatment.

\medterm{Tumor Marker} A substance found in blood, urine, tissue, or another body fluid that may be produced by cancer or by the body’s response to it. Tumor markers can help monitor treatment or recurrence and sometimes guide diagnosis, but most are not specific enough to prove cancer alone.

\medterm{Tumor Marker, CEA} CEA is a blood tumor marker used mainly to help monitor some colorectal and other cancers; it is not specific enough to diagnose cancer alone.

\medterm{Tumor Marker, NSE} NSE is a protein marker that may rise in some neuroendocrine tumors or nerve-cell injury but is not diagnostic by itself.

\medterm{Tumor Microenvironment} The cells, blood vessels, connective tissue, and chemical signals surrounding a tumor. These nearby tissues can support or limit tumor growth, affect spread, influence the immune response, and change how well treatment works.
\medterm{Tumor Mutational Burden} The number of somatic mutations per unit of coding DNA in a tumor specimen. TMB may
correlate with production of neoantigens and response to some immune checkpoint
inhibitors, but interpretation requires a validated assay and disease-specific context.

\medterm{Tumor Necrosis Factor} A proinflammatory cytokine (TNF-alpha) produced mainly by macrophages, acting on
TNFR1/TNFR2 receptors. Functions: activation of endothelium and leukocytes, fever,
cachexia, apoptosis, and coordination of inflammation and the acute-phase response.

\medterm{Tumor Necrosis Factor Receptor-Associated Periodic Syndrome} A rare inherited autoinflammatory disorder caused by variants in the TNFRSF1A gene, producing recurrent prolonged fever, muscle pain, abdominal pain, rash, and inflammation; it is also called TRAPS.

\medterm{Tumor Staging} A system of describing the size of a cancer and how far it has spread throughout the
body from its original location (e.g., using the TNM system).

\medterm{Tumor Suppressor} A gene or protein that helps limit cell growth, repair DNA damage, or promote removal of abnormal cells. Loss of its function can
contribute to cancer.

\medterm{Tumor Suppressor Gene} A gene that helps prevent uncontrolled cell growth by slowing division, repairing DNA damage, or directing abnormal cells to die. If both copies or the gene’s function are lost or disabled, cells may grow abnormally and cancer risk can increase.

\medterm{Tumor Testicle (Testicular Disorders)} Alternate index label for Testicular Disorders.

\medterm{Tumor, Salivary Gland} Alternate index label; see \textit{Salivary gland tumors}.

\medterm{Tumor, Spinal} Alternate index label; see \textit{Spinal tumor and spinal mass}.

\medterm{Tumor-Suppressor Gene} A gene whose normal product restrains cell proliferation, promotes DNA repair, induces
apoptosis, or maintains genomic stability. Loss or inactivation of both functional
copies is common in the development of many cancers, although mechanisms vary by gene.

\medterm{Tumors Uterine (Uterine Fibroids)} Alternate index label for Uterine Fibroids.

\medterm{Tumour} An abnormal mass caused by uncontrolled or inappropriate cell growth; it may be noncancerous or cancerous.
\medterm{Tumour (Neoplasm)} An abnormal mass of tissue produced by uncontrolled or inappropriate cell growth; it may be benign or malignant.

\synonyms
Neoplasm

\medterm{Tungiasis} A skin infestation caused by the sand flea Tunga penetrans, producing an itchy or painful lesion, usually on
the feet.

\medterm{Tunica Adventitia} The outer connective-tissue layer of a blood-vessel wall. It contains collagen and
elastic fibers that provide structural support, along with nerves and, in larger
vessels, small vessels called vasa vasorum that supply the outer wall. It is relatively
prominent in veins and helps anchor vessels to surrounding tissues.

\medterm{Tunica Media} The middle layer of a blood-vessel wall, located between the tunica intima and tunica
adventitia. It consists primarily of smooth muscle and elastic fibers; its contraction
or relaxation changes vascular diameter, regulating blood pressure and tissue perfusion.
It is thickest in arteries and thinner in veins.

\medterm{Tunnel Vision} Loss of peripheral visual field leaving a narrowed central field, caused by glaucoma, retinal disease, optic-nerve or brain lesions, severe hypoxia, or other conditions.

\medterm{Turbinate Surgery} An operation that reduces or reshapes swollen tissues along the side walls of the nose to improve airflow. It may help persistent nasal blockage when other treatment fails; bleeding, crusting, dryness, scarring, or ongoing obstruction are possible.

\medterm{Turbinates} Curved bony shelves on the lateral walls of the nasal cavity, covered by vascular mucosa. They increase surface area and help
warm, humidify, and filter inhaled air.

\medterm{Turbinates Hypertrophy} Chronic enlargement of the nasal turbinates, most commonly the inferior turbinates,
resulting in nasal airway obstruction, which can be caused by chronic inflammation from
allergic rhinitis, non-allergic rhinitis, chronic sinusitis, hormonal changes, or
vasomotor dysfunction.

\medterm{Turcot Syndrome} A rare inherited cancer-predisposition pattern involving colon polyps or colon cancer together with a brain tumor. It can result from mismatch-repair gene changes or the APC gene change seen in familial adenomatous polyposis, so genetic counseling and screening are important.

\synonyms
Brain Tumor-Polyposis Syndrome

\medterm{Turgor} The firmness of tissue produced by fluid pressure within cells; skin turgor can help assess dehydration.
\medterm{Turmeric} Turmeric is a plant spice containing curcuminoids and used in foods and supplements; evidence varies by condition and interactions or liver injury are possible.

\medterm{Turner Hypoplasia} Enamel defect on permanent tooth from infection of overlying primary tooth. Localized hypoplasia. A defect in permanent tooth enamel caused by infection or injury involving the baby tooth above it.
\medterm{Turner Syndrome} A chromosome condition affecting females in which one X chromosome is missing or altered. It commonly causes short stature and reduced ovarian function, with delayed or absent puberty and infertility. Heart, kidney, hearing, and thyroid problems may also occur, so lifelong medical follow-up is important.

\medterm{Turner'S Syndrome} A chromosome condition affecting females in which one X chromosome is missing or altered. It commonly causes short stature and reduced ovarian function, with delayed or absent puberty and infertility. Heart, kidney, hearing, and thyroid problems may also occur, so lifelong medical follow-up is important.
\medterm{Turning Patients Over In Bed} Repositioning a patient at planned intervals to protect skin, improve comfort and breathing, prevent pressure injuries and contractures, and reduce complications of immobility.

\medterm{Turpentine Oil Poisoning} Poisoning from swallowing or inhaling turpentine oil. It can irritate or damage the lungs, stomach, and nervous system and may cause chemical pneumonia.

\medterm{Tussis} Cough; a forceful reflex expiration that helps clear secretions, irritants, or foreign material from the airways.

\medterm{TVS} TVS commonly means transvaginal sonography, an ultrasound examination using a probe in the vagina to assess pelvic organs and early pregnancy.


\medterm{Twilight Sleep} Twilight sleep is an older term for a sedated, partly conscious state used during procedures; modern analgesia and anesthesia specify the medicines and monitoring used.

\medterm{Twin} A twin is one of two offspring from the same pregnancy, either identical from one fertilized egg or fraternal from two eggs and two sperm.

\medterm{Twin Placenta} Placental anatomy in a twin pregnancy, classified by chorionicity and amnionicity; these features determine risks such as twin-to-twin transfusion and guide antenatal monitoring.

\medterm{Twin Pregnancy} A pregnancy involving two fetuses. Clinical risk depends on chorionicity and amnionicity, with increased risks including preterm birth, fetal growth problems, and hypertensive disorders.

\medterm{Twin-To-Twin Transfusion Syndrome} A complication of monochorionic twin pregnancy in which shared placental vessels cause
unequal blood flow between the twins.

\medterm{Twin-Twin Transfusion Syndrome} A complication of monochorionic diamniotic twin pregnancies where placental vascular
anastomoses (arteriovenous connections) cause unbalanced blood flow from one twin (the
donor) to the other (the recipient). The donor twin becomes hypovolemic, anemic, and
growth-restricted with oligohydramnios ("stuck twin").

\medterm{Twins} Two offspring from the same pregnancy, either monozygotic or dizygotic.
\medterm{Twitch} A brief involuntary muscle contraction or movement, which may be benign, medication-related, stress-related, or a sign of nerve, muscle, or brain disease.

\medterm{Twitching} Repeated involuntary muscle contractions or fasciculations, commonly caused by fatigue, stimulants, electrolyte disturbance, nerve irritation, medicines, or neurologic disease.


\medterm{Tympanic} Tympanic means relating to the eardrum or middle ear, or resembling a drum in shape or sound.

\medterm{Tympanic Bulla} The tympanic bulla is a hollow or expanded bony chamber surrounding the middle ear, especially prominent in many mammals.

\medterm{Tympanic Cavity} The air-filled space of the middle ear containing the malleus, incus, and stapes. The air-filled middle-ear chamber containing three small bones that transmit sound vibrations toward the inner ear.
\medterm{Tympanic Hyperectasis} Abnormal outward bulging or retraction of the tympanic membrane from chronic middle ear
pressure changes, potentially causing conductive hearing loss and requiring tympanostomy
tube placement.

\medterm{Tympanic Membrane} The thin eardrum separating the outer ear canal from the middle ear. Sound makes it vibrate, and those vibrations move the middle-ear bones; infection, injury, or pressure can tear it.
\medterm{Tympanic Membrane Perforation} A tear or hole in the eardrum caused by infection, trauma, pressure changes, or instrumentation, potentially producing pain, drainage, or conductive hearing loss.

\medterm{Tympanites} Abdominal distension caused by excessive gas in the stomach or intestines, producing a drum-like sound on percussion and sometimes pain, belching, or altered bowel movements.
\medterm{Tympanometry} An objective test that varies air pressure in the ear canal while measuring movement of the tympanic membrane and middle-ear pressure. It helps assess middle-ear fluid, eustachian-tube dysfunction, perforation, and ossicular or tympanic-membrane abnormalities.

\medterm{Tympanoplasty} Surgical repair or reconstruction of the tympanic membrane and, when necessary, the middle-ear ossicles to close a perforation or improve hearing.

\medterm{Tympanosclerosis} Scarring and calcification of the tympanic membrane or middle-ear lining, often following recurrent infection or previous inflammation.
It may be asymptomatic or cause conductive hearing loss when the ossicles are affected.

\synonyms
Tympanic-membrane sclerosis; myringosclerosis when limited to the tympanic membrane.

\medterm{Tympanostomy Tubes} Small tubes inserted through the tympanic membrane to ventilate the middle ear and
prevent fluid accumulation. Indicated for recurrent acute otitis media (3 episodes in 6
months or 4 in 12 months) or persistent otitis media with effusion greater than 3 months
with hearing loss. Performed as outpatient procedure under general anesthesia.

\medterm{Tympanum} The middle-ear cavity or, in some contexts, the tympanic membrane.
\medterm{Tympany} Tympany is a drumlike sound on percussion, often produced by gas in a hollow organ such as the stomach or intestine.

\medterm{Type 1 Diabetes} Type 1 diabetes is autoimmune destruction of pancreatic beta cells causing absolute insulin deficiency, hyperglycemia, and risk of diabetic ketoacidosis without insulin replacement.

\synonyms
Diabetes, Type 1
Diabetes Type 1 (Type 1 Diabetes)


\medterm{Type 1 Diabetes In Children} Type 1 diabetes diagnosed during childhood, causing high blood sugar and requiring insulin, monitoring, nutrition planning, and education.

\medterm{Type 1 Diabetes Mellitus} Type 1 diabetes is an autoimmune disease in which the immune system attacks and destroys
the insulin-producing beta cells of the pancreas, resulting in absolute insulin
deficiency.

\medterm{Type 2 Diabetes} Type 2 diabetes is chronic hyperglycemia caused by insulin resistance and progressive beta-cell dysfunction, with risks to the eyes, kidneys, nerves, heart, and blood vessels.

\synonyms
Diabetes, Type 2
Diabetes Type 2 (Type 2 Diabetes)

\medterm{Type 2 Diabetes - Oral Medicines} Oral medicines for type 2 diabetes lower glucose through effects on insulin, liver glucose production, kidney reabsorption, or intestinal absorption; choice depends on health and treatment goals.



\medterm{Type 2 Diabetes In Children} Type 2 diabetes beginning during childhood or adolescence, often involving insulin resistance and requiring lifestyle and medical management.

\medterm{Type 2 Diabetes Mellitus} Type 2 diabetes is a chronic metabolic disorder characterized by insulin resistance and
progressive pancreatic beta-cell dysfunction, leading to hyperglycemia.

\medterm{Type 3B Diabetes} Type 3b diabetes is a form of diabetes caused by disease or removal of the pancreas, which can impair both insulin and digestive-enzyme production.

\medterm{Type I Hypersensitivity} An immediate (anaphylactic) hypersensitivity reaction mediated by allergen-specific IgE
bound to mast cells and basophils; re-exposure cross-links IgE, triggering degranulation
and release of histamine and mediators within minutes. Manifestations: allergic
rhinitis, asthma, urticaria, food allergy, and anaphylaxis.

\medterm{Type II Hypersensitivity} An antibody-mediated (cytotoxic) hypersensitivity reaction in which IgG or IgM binds to
cell-surface or matrix antigens, causing cell destruction by complement activation
(MAC), opsonization/phagocytosis, or antibody-dependent cellular cytotoxicity.

\medterm{Type III Hypersensitivity} An immune complex-mediated hypersensitivity reaction in which circulating
antigen-antibody complexes deposit in tissues and vessels, activating complement and
attracting neutrophils, causing inflammation and tissue damage. Local forms: Arthus
reaction, serum sickness.

\medterm{Type III Hypersensitivity Reaction} Tissue injury caused when antigen-antibody clusters deposit in tissues and activate inflammation. This reaction can damage blood vessels, joints, kidneys, skin, or other organs, as in serum sickness or lupus.

\medterm{Type IV Hypersensitivity} A delayed-type hypersensitivity (cell-mediated) reaction mediated by sensitized T cells
(not antibody), occurring 24-72 hours after antigen exposure.

\medterm{Type V Glycogen Storage Disease} An inherited muscle disorder, also called McArdle disease, in which muscle cells cannot properly break down glycogen for energy. Exercise can cause muscle pain, cramps, and sometimes dark urine from muscle injury.


\medterm{Types Of Health Care Providers} Health professionals and services that diagnose, treat, prevent, or support health needs, including physicians, nurses, pharmacists, therapists, dentists, mental-health professionals, and community services.


\medterm{Types Of Ileostomy} Types of ileostomy include end, loop, and continent procedures, which divert stool through an opening in the abdominal wall when the colon or rectum cannot be used normally.

\medterm{Typhlitis} Inflammation of the cecum in immunocompromised patients particularly those with
neutropenia from chemotherapy. Also known as neutropenic enterocolitis. Presents with
right lower quadrant pain, fever, and diarrhea. Diagnosis is by CT showing cecal wall
thickening.

\synonyms
neutropenic enterocolitis.

\medterm{Typhoid} A systemic infection caused by Salmonella enterica serovar Typhi, usually acquired through contaminated food or
water and characterized by prolonged fever and gastrointestinal or systemic symptoms.

\medterm{Typhoid Fever} A potentially serious whole-body infection caused by Salmonella Typhi, usually spread through contaminated food or water.

\synonyms
Enteric Fever (Typhoid Fever)

\medterm{Typhus} Typhus refers to several rickettsial infections transmitted by lice, fleas, or mites and characterized by fever, headache, rash, and variable systemic complications.
arthropod vectors. Three main forms: (1) Epidemic (louse-borne) typhus—Rickettsia
prowazekii, transmitted by the human body louse (Pediculus humanus corporis). Associated
with overcrowding, poverty, cold climate, and refugees ('war typhus').

\medterm{Typhus Fever} A group of infections caused by Rickettsia species and transmitted mainly by lice or fleas.
\medterm{Typical Antipsychotics} Typical antipsychotics, also known as first-generation antipsychotics, primarily block
dopamine D2 receptors in the mesolimbic pathway, reducing positive symptoms of
schizophrenia including hallucinations, delusions, and disorganized thinking.

\synonyms
first-generation antipsychotics, primarily block.

\medterm{Tyramine} A naturally occurring amine found in some aged or fermented foods. In people taking monoamine oxidase inhibitors, excessive tyramine can cause a dangerous rise in blood pressure.
\medterm{Tyrosinase Peptide} A short protein fragment related to tyrosinase, the enzyme that helps make melanin, the pigment in skin, hair, and eyes. Such peptides are mainly used in research.

\medterm{Tyrosine} An amino acid used to build proteins and make thyroid hormones, melanin, and adrenaline-like chemicals. The body usually makes enough from phenylalanine, but inherited metabolic disorders can disrupt this pathway.

\synonyms
Full term

\medterm{Tyrosinemia} An inherited disorder of tyrosine metabolism that causes accumulation of tyrosine or its metabolites. Type I can cause severe liver and kidney disease and is associated with fumarylacetoacetate hydrolase deficiency.

\medterm{Tzanck Smear} A microscope test of cells scraped from the base of a skin blister or sore. It may show changes suggesting herpesvirus infection or pemphigus, but it cannot reliably identify the cause and has largely been replaced by more specific tests.

\medterm{Tamiflu (Oseltamivir)} Tamiflu is the brand name for oseltamivir, an antiviral medicine used to treat or help prevent influenza. It works best when started early and does not replace flu vaccination.

\medterm{Tamoxifen} Tamoxifen is a selective estrogen-receptor modulator used mainly to treat or prevent some estrogen-sensitive breast cancers. It can increase the risk of blood clots and, rarely, uterine cancer.

\medterm{Tarlov Cysts} Tarlov cysts are fluid-filled sacs around the nerve roots, most often in the sacral spine. Many cause no symptoms, but larger or symptomatic cysts may cause pain, numbness, or bladder problems.

\medterm{Taste Receptor} A taste receptor is a protein on taste-bud cells that detects chemical substances in food, helping create taste sensations such as sweet, salty, sour, bitter, and umami.

\medterm{Tattoos} Tattoos are permanent designs made by placing pigment into the skin. Risks include infection, allergic reactions, scarring, and difficulty removing the pigment.

\medterm{Taurine} Taurine is a sulfur-containing compound involved in normal cell, nerve, muscle, and bile-acid function. It is present in foods and added to some energy drinks; high caffeine intake in such drinks can be harmful.

\medterm{Taxotere} Taxotere is the brand name for docetaxel, a chemotherapy medicine that interferes with cell division. It is used for several cancers and can cause low blood counts, infection risk, hair loss, and other side effects.

\medterm{T-Cell} A T-cell is a type of white blood cell that helps coordinate immune responses or directly destroys infected or abnormal cells. T cells mature in the thymus.

\medterm{Tea} Tea is a drink made by steeping plant leaves or other plant material in water. Some teas contain caffeine, and concentrated herbal products can interact with medicines.

\medterm{Teeth Whitening} Teeth whitening lightens tooth color by removing or chemically breaking down stains. Professional or properly used products are generally safe, but temporary sensitivity and gum irritation can occur.

\medterm{Teledermatology} Teledermatology is the use of video, photographs, or other digital communication to assess and manage skin conditions remotely. Some problems still require an in-person examination or biopsy.

\medterm{Tendinosis} Tendinosis is chronic degeneration and disorganization of tendon tissue, usually related to repeated loading rather than acute inflammation. Treatment may include activity modification and rehabilitation.

\medterm{Testicular Failure} Testicular failure occurs when the testes do not produce enough testosterone, sperm, or both. Causes include genetic conditions, injury, infection, medicines, and damage from chemotherapy or radiation.

\medterm{The Microbiome Project} The Microbiome Project refers to research that studies the microorganisms living in and on the human body and their genes. This research helps investigate links between microbes, health, and disease.

\medterm{Thimerosal} Thimerosal is a mercury-containing preservative used in some multidose vaccines and other medical products. The mercury form it contains is ethylmercury, which is cleared from the body more quickly than methylmercury.

\medterm{Threadworms} Threadworms, also called pinworms, are small intestinal worms that commonly cause itching around the anus at night. Treatment usually includes an antihelminthic medicine and hygiene measures for household contacts.

\medterm{Thymine} Thymine is one of the four main bases in DNA. It pairs with adenine through hydrogen bonds and helps store genetic information.

\medterm{Tick-borne Diseases} Tick-borne diseases are infections or other illnesses transmitted by ticks, including Lyme disease, babesiosis, ehrlichiosis, and some viral diseases. Removing ticks promptly can reduce risk.

\medterm{Tokophobia} Tokophobia is an intense fear of pregnancy or childbirth that can cause substantial distress or avoidance. Support may include psychological therapy, birth planning, and specialist maternity care.

\medterm{Tooth Brushing} Tooth brushing removes plaque and food from teeth and gums. Brushing twice daily with fluoride toothpaste helps prevent tooth decay and gum disease.

\medterm{Tooth Erosion} Tooth erosion is the chemical loss of tooth enamel caused by acids, not by bacteria. Acidic foods, drinks, reflux, and repeated vomiting can contribute to it.

\medterm{Traditional Medicine} Traditional medicine includes health practices, remedies, and healing systems developed within cultural traditions. Products and practices vary in evidence and safety, so they should not replace effective care without professional advice.

\medterm{Travel Vaccinations} Travel vaccinations are vaccines recommended or required for protection against infections encountered in particular countries or regions. Recommendations depend on the destination, activities, health conditions, and previous vaccinations.

\medterm{Trichinellosis} Trichinellosis is an infection caused by Trichinella worms, usually after eating undercooked infected meat. It may cause diarrhea followed by fever, muscle pain, and swelling around the eyes.

\medterm{Triclosan} Triclosan is an antimicrobial chemical formerly used in some soaps, cosmetics, and consumer products. Routine use in consumer products has limited benefit and may contribute to environmental and antimicrobial concerns.

\medterm{Tridecanoic Acid} Tridecanoic acid is a saturated fatty acid containing 13 carbon atoms. It occurs in small amounts in some foods and is mainly studied in nutrition and lipid research.

\medterm{Tropical Disease} A tropical disease is an infection or other health condition more common in tropical and subtropical regions, often because of local climate, vectors, or living conditions. Travel history helps guide diagnosis.

\medterm{Tubulointerstitial Nephropathy} Tubulointerstitial nephropathy is kidney damage involving the renal tubules and surrounding tissue. It may result from medicines, infection, toxins, autoimmune disease, or inherited disorders and can impair kidney function.


\medterm{U Wave} A small deflection following the T wave, best seen in precordial leads V2-V3. U waves
are thought to represent repolarization of the Purkinje fibers or papillary muscles.
Prominent U waves are associated with hypokalemia, bradycardia, and certain drugs
(digitalis, quinidine).

\medterm{UA} UA commonly means urinalysis, laboratory examination of urine for blood, protein, cells, glucose, infection, and other findings.

\medterm{Ubidecarenone} Coenzyme Q10, a lipid-soluble molecule involved in mitochondrial electron transport and cellular energy production; it is sold as a supplement and studied in selected cardiovascular and neuromuscular conditions.

\medterm{Ubiquinone} A lipid-soluble electron carrier in the mitochondrial membrane. See Coenzyme Q.

\medterm{Ubiquitin} A small regulatory protein that is covalently attached to other proteins to label them for degradation by the proteasome or to alter their cellular trafficking and signaling.

\medterm{Ubiquitination} Ubiquitination is attachment of ubiquitin to a target protein through an E1--E2--E3 enzyme cascade, marking it for proteasomal degradation or altering its localization, activity, or signaling.

\medterm{Ubiquitous} Present or found everywhere. In medical writing, the term is descriptive rather than a diagnosis and may refer to a substance, organism, or process that is widely distributed.

\medterm{UBT} UBT commonly means urea breath test, which detects active Helicobacter pylori infection by measuring labeled carbon dioxide after ingested urea.

\medterm{UDP-Glucuronosyltransferase} UDP-glucuronosyltransferase is an enzyme family that attaches glucuronic acid to bilirubin, medicines, hormones, and other compounds for elimination.

\medterm{Uhl Anomaly} A rare congenital absence or severe deficiency of right-ventricular myocardium, leaving
a thin-walled right ventricle. It may cause right-sided heart failure, arrhythmia,
cyanosis, or sudden deterioration in infancy or childhood.

\medterm{Ulcer} Full-thickness loss of epidermis extending into the dermis or deeper tissue, often
healing with a scar. Causes include vascular disease, pressure, infection, neuropathy,
inflammatory disease, trauma, and malignancy.

\medterm{Ulcer (Vascular)} A persistent break in the skin caused or maintained by inadequate blood flow, venous hypertension,
neuropathy, or related vascular disease.

\synonyms
Vascular

\medterm{Ulcer Healing Drugs} Medicines used to promote healing of gastrointestinal ulcers, usually by suppressing acid secretion or protecting the mucosa; treatment also addresses causes such as \textit{H. pylori} infection or NSAID use.

\medterm{Ulcer, Aphthous} Alternate index label; see \textit{Canker sore}.

\medterm{Ulcer, Peptic} Alternate index label; see \textit{Peptic ulcer}.

\medterm{Ulcer, Rectal} Alternate index label; see \textit{Solitary rectal ulcer syndrome}.

\medterm{Ulceration} The formation of an open sore caused by loss of surface tissue, usually involving skin or a mucous membrane and sometimes extending into deeper tissue.

\medterm{Ulcerative Colitis} A chronic inflammatory disease of the colon and rectum, usually causing continuous mucosal inflammation
that begins in the rectum and extends proximally. Symptoms may include bloody diarrhea,
urgency, abdominal pain, and extraintestinal manifestations.

\synonyms
Colitis, Ulcerative
Colitis Ulcerative (Ulcerative Colitis)

\medterm{Ulcerative Keratitis} An ulcerative infection or inflammatory defect of the cornea that can cause pain, redness, photophobia, and impaired vision and may threaten sight if untreated.

\medterm{Ulcerative Proctitis} Inflammation and ulceration limited to the lining of the rectum, often causing rectal bleeding, urgency, tenesmus, mucus, and diarrhea.

\medterm{Ullrich Congenital Muscular Dystrophy} An inherited muscle disease caused by problems with collagen VI, a protein that supports muscle and connective tissue. It begins in infancy with weakness, low muscle tone, tight joints near the body, unusually loose joints in the hands and feet, and later breathing problems.

\synonyms
UCMD; Ullrich Scleroatonic Muscular Dystrophy

\medterm{Ulna} Medial forearm bone, longer than the radius. Proximal end: trochlear notch, olecranon, coronoid process. Shaft. Distal end: head, styloid process. The trochlear notch articulates with the trochlea forming the primary elbow hinge.

\medterm{Ulnar} Relating to the ulna bone of the forearm or the ulnar side (medial side) of the hand.
Ulnar nerve compression causes paresthesia in medial one and a half fingers.

\medterm{Ulnar Artery} The larger terminal branch of the brachial artery at the cubital fossa, running medially
down the forearm to the wrist (palpable medial to the pisiform) and hand. Supplies the
medial forearm and ulnar side of the hand; forms the superficial palmar arch. Provides
the common interosseous artery.

\medterm{Ulnar Vein} Deep veins accompanying the ulnar artery that drain the medial forearm and join the radial veins to form the brachial veins.

\medterm{Ulnar Hypoplasia} Underdevelopment or partial absence of the ulna, usually presenting as a congenital
forearm and hand anomaly. It may cause forearm shortening, radial deviation of the hand,
limited elbow or wrist movement, and associated digital or systemic abnormalities.

\medterm{Ulnar Nerve} From medial cord (C8-T1), passing behind the medial epicondyle (cubital tunnel), into
the forearm innervating flexor carpi ulnaris and medial FDP half. In the hand innervates
hypothenar muscles, interossei, medial two lumbricals, and adductor pollicis. Sensation
to medial 1.5 digits.

\medterm{Ulnar Nerve Dysfunction} Impairment of the ulnar nerve, commonly caused by compression at the elbow (cubital
tunnel) or wrist (Guyon canal). Symptoms include numbness or tingling in the little and
ring fingers, weakened grip, and hand-muscle wasting; severe lesions may produce an
ulnar claw.

\medterm{Ulnar Nerve Entrapment} Compression or irritation of the ulnar nerve, commonly at the elbow or wrist, causing numbness or tingling in the ring and little fingers and sometimes weakness of the hand.

\medterm{Ulnar Paradox} The clinical observation that a lesion of the ulnar nerve at the wrist may cause more clawing of the fingers than a lesion higher in the forearm, because the ulnar-innervated finger flexors remain intact.

\medterm{Ulnar Wrist Pain} Pain on the side of the wrist nearest the little finger, often involving the ulna or nearby tendons and ligaments.

\synonyms
Ulnar-Sided Wrist Pain

\medterm{Ulnar-Sided Wrist Pain} Alternate index label; see \textit{Ulnar wrist pain}.

\medterm{Ulocladium} A genus of dematiaceous filamentous fungi found in soil and plant material. It is an uncommon opportunistic pathogen that can cause keratitis, sinusitis, or invasive infection in susceptible people.

\medterm{Ultracentrifugation} Ultracentrifugation separates particles or macromolecules at very high centrifugal force according to size, shape, and density; it is used to study viruses, organelles, lipoproteins, and protein complexes.

\medterm{Ultradian Rhythm} A biological rhythm that repeats more than once in 24 hours, such as pulsatile hormone secretion, sleep-stage cycling, or periodic changes in heart and breathing patterns.

\medterm{Ultrafiltration} Separation in which pressure pushes water and small dissolved substances through a membrane while larger particles stay behind. The process helps form urine and is also used during dialysis.

\medterm{Ultralente Insulin} Ultralente insulin is a prolonged-acting zinc insulin suspension formerly used for basal insulin replacement; it has largely been replaced by newer long-acting analogues and can cause hypoglycaemia and dosing variability.

\medterm{Ultrasonic Therapy} Use of high-frequency sound energy for therapeutic purposes such as selected physical-therapy applications or tissue ablation; benefits and risks depend on the device and treatment site.

\medterm{Ultrasonography} A scan using high-frequency sound waves to create images of organs, tissues, vessels, or a fetus without radiation.

\medterm{Ultrasound} An imaging test that uses high-frequency sound waves to show organs, blood vessels, tissues, or a developing fetus without ionizing radiation. Different examinations can assess the abdomen, pelvis, kidneys, thyroid, heart, blood flow, or testicles.

\medterm{Ultrasound Biomicroscopy} High-frequency ultrasound imaging used to visualize the anterior segment of the eye, including the cornea, iris, ciliary body, lens, and angle when ordinary examination or imaging is insufficient.

\medterm{Ultrasound Contrast} Tiny gas-filled bubbles given into a vein to improve ultrasound images of blood flow and selected organs. They are usually cleared through the lungs but can rarely cause allergic or cardiopulmonary reactions.
\medterm{Ultrasound Therapy} Use of sound waves at therapeutic intensity to heat or mechanically affect tissue. It is used in selected rehabilitation settings, but evidence and suitability vary by condition and it is not appropriate over every body area.

\medterm{Ultrasound Transducer} A device that emits and receives high-frequency sound waves to generate diagnostic
images or guide procedures. Linear, curvilinear, phased-array, and endocavitary
transducers are selected according to depth and anatomic application.

\medterm{Ultraviolet Microscopy} Microscopy that uses ultraviolet wavelengths to produce images or fluorescence from structures and substances that absorb or emit ultraviolet light. It has specialized laboratory and research applications.

\medterm{Ultraviolet Rays} Invisible electromagnetic radiation with wavelengths shorter than visible violet light. Excess exposure can cause sunburn, photoaging, eye injury, and skin cancers; controlled ultraviolet exposure has selected therapeutic uses.

\medterm{Ultraviolet Therapy} Controlled exposure to ultraviolet light used for selected skin disorders such as psoriasis, eczema, and vitiligo; treatment requires eye and skin protection and monitoring for burns and long-term skin cancer risk.

\medterm{Umbilical} Relating to the umbilicus or navel, including structures, lesions, discharge, or abnormalities located at or connected with it.

\medterm{Umbilical Artery Doppler} Ultrasound assessment of blood flow in the umbilical arteries to evaluate placental
resistance. Normal flow shows a pulsatility index decreasing throughout pregnancy with
forward diastolic flow. Absent or reversed end-diastolic flow indicates severe placental
dysfunction and is associated with increased perinatal morbidity and mortality.

\medterm{Umbilical Catheters} Umbilical catheters provide vascular access through a newborn’s umbilical vessels for fluids, medicines, monitoring, or blood sampling while complications are monitored.

\medterm{Umbilical Cord} The connection between the fetus and the placenta, containing two umbilical arteries
(carrying deoxygenated blood from the fetus to the placenta) and one umbilical vein
(carrying oxygenated blood from the placenta to the fetus). The vessels are embedded in
Wharton jelly, a mucopolysaccharide matrix that prevents compression.

\medterm{Umbilical Cord Blood Gas} Analysis of arterial and venous cord blood collected after delivery. The umbilical artery
best reflects fetal acid-base status, while the umbilical vein reflects placental blood. Results
are interpreted using gestational age, delivery circumstances, pH, carbon dioxide, and base
deficit or excess.

\medterm{Umbilical Cord Prolapse} Descent of the umbilical cord through the cervix alongside or past the presenting fetal
part, occurring when the membranes rupture. Risk factors: malpresentation (transverse
lie, breech), preterm/low birth weight, polyhydramnios, multiparity, and artificial ROM
with high presenting part.

\medterm{Umbilical Granuloma} A small moist pink or red lump of healing tissue at a baby's navel after the cord falls off. It may ooze clear fluid but should not cause spreading redness, fever, or marked tenderness; a clinician can treat it and exclude infection or another abnormal connection.

\synonyms
Granuloma Umbilicale; Neonatal Umbilical Granuloma

\medterm{Umbilical Hernia} A hernia through the umbilical ring, common in infants (most umbilical hernias close
spontaneously by age 4-5) and adults (risk factors: obesity, pregnancy, ascites,
cirrhosis). Presents as a reducible bulge at the umbilicus, worse with increased
intra-abdominal pressure. In adults: risk of incarceration and strangulation is higher.

\medterm{Umbilical Region} Central region around the umbilicus containing small intestine, transverse colon, and
great vessels. The umbilicus is a landmark for assessing ascites, caput medusae in
portal hypertension, and umbilical hernias. Percussion and auscultation here are
essential components of the abdominal examination.

\medterm{Umbilical Vein} The umbilical vein carries oxygen- and nutrient-rich blood from the placenta to the fetus during pregnancy.

\medterm{Umbilication} Formation of a central hollow or navel-like depression in a skin lesion or other body structure. It can occur in some infections, inflammatory lesions, tumors, or normal anatomical features and must be interpreted with the lesion’s appearance.
\medterm{Umbilicus} The navel, a scar on the anterior abdominal wall marking the former attachment of the umbilical cord. It can be affected by hernia, infection, discharge, or metastatic disease.

\medterm{Umbilicus (Navel)} The scar on the anterior abdominal wall marking the site of attachment of the umbilical
cord in fetal life, located at the level of the L3-L4 intervertebral disc, approximately
midway between the xiphoid process and the pubic symphysis.

\synonyms
Navel

\medterm{Uncertainty} Lack of complete certainty about a diagnosis, prognosis, measurement, or treatment effect. Clinical uncertainty is managed by combining available evidence with follow-up, additional testing, and shared decision-making.

\medterm{Uncinate Fit} A nonstandard or unclear historical label in the source dictionary; its intended clinical meaning should be verified against the original reference.

\medterm{Uncompetitive Inhibition} A type of enzyme inhibition where the inhibitor binds only to the enzyme-substrate
complex at a site distinct from the active site, forming an inactive ESI complex. The
inhibitor cannot bind to the free enzyme.

\medterm{Unconscious} Not awake and not responsive to ordinary external stimuli. Unconsciousness may result from syncope, seizure, trauma, intoxication, metabolic disease, stroke, or other emergencies.

\medterm{Unconsciousness} Loss of awareness and normal response to surroundings, ranging from brief fainting to prolonged coma. It can result from reduced brain blood flow, injury, seizures, low blood sugar, drugs, or serious illness.

\medterm{Unconsciousness, Temporary} A brief episode of loss of awareness and responsiveness, often caused by syncope but also possible with seizure, concussion, hypoglycemia, or other conditions requiring assessment.

\medterm{Underventilation} Breathing that is too shallow or slow to remove enough carbon dioxide, causing hypercapnia and sometimes low oxygen; causes include sedatives, neuromuscular weakness, obesity hypoventilation, and lung disease.

\medterm{Underweight} A body weight below the usual range for height. In adults, a Body Mass Index below 18.5 kg/m\(^2\) is commonly classified as underweight; illness or poor nutrition may contribute.

\medterm{Undescended Testicle} A testicle that has not moved into the scrotum before birth, also called cryptorchidism. It may descend during early infancy, but a testicle that remains high can have increased risks of infertility, twisting, injury, and cancer and should be assessed.

\synonyms
Testicle, Undescended

\medterm{Undescended Testicle Repair} Surgical relocation and fixation of a testis that has not descended into the scrotum, usually performed in childhood to preserve function and reduce complications.

\medterm{Undescended Testis} Failure of one or both testes to descend into the scrotum. Some descend during early
infancy, but persistent cryptorchidism can impair fertility and increase the risk of torsion, hernia, or testicular
cancer.

\medterm{Undescended Testis (Cryptorchidism)} Congenital failure of one or both testes to descend along the normal path
into the scrotum. Persistent cases require evaluation because they may affect fertility and increase the risk of torsion,
hernia, or testicular cancer.

\synonyms
Cryptorchidism

\medterm{Undetermined} Not established or not able to be classified with the available clinical information. An undetermined result may require repeat assessment or additional testing.

\medterm{Undifferentiated} Lacking recognizable specialized features. In pathology, an undifferentiated tumor has cells that cannot be assigned readily to a specific tissue of origin based on their appearance.

\medterm{Undifferentiated Cancer} A cancer whose cells lack recognizable specialized features, making the tissue of origin or exact subtype difficult to identify; diagnosis requires pathology and staging.

\medterm{Undifferentiated Carcinoma} A high-grade carcinoma whose cells lack recognizable features of a specific tissue or tumor subtype; diagnosis and treatment require pathology, immunostaining, and staging.

\medterm{Undifferentiated Pleomorphic Sarcoma} An uncommon aggressive cancer of soft tissue or bone whose cells have no clearly identifiable line of differentiation.

\synonyms
Malignant Fibrous Histiocytoma

\medterm{Undifferentiated Pleomorphic Sarcoma Of The Breast} A rare high-grade breast sarcoma formerly called malignant fibrous histiocytoma. It is a diagnosis of exclusion after other specific sarcomas and metaplastic carcinoma have been ruled out histologically.
\medterm{Undulant Fever} A zoonotic infectious disease caused by gram-negative coccobacilli of the genus
Brucella, transmitted to humans through contact with infected animals (cattle, goats,
pigs, sheep) or consumption of unpasteurized dairy products. See Brucellosis.

\medterm{Undulate} To move or have a wave-like contour. In a clinical description, it may refer to a palpable fluid wave, a wavy tissue margin, or a rhythmic motion.

\medterm{Undulate Ribs} Ribs with a wavy or irregular contour, a descriptive anatomic finding that may reflect a developmental variation, deformity, prior injury, or underlying bone disease.


\medterm{Unequivocal} Clear and not open to more than one reasonable interpretation. An unequivocal clinical finding or test result provides strong evidence for the stated interpretation when considered in context.

\medterm{Unerupted Tooth} A tooth that has not emerged into the oral cavity at the expected time. It may be normally developing, delayed, impacted, or prevented from erupting by obstruction or abnormal position.

\medterm{Unevaluable} Unable to be assessed or interpreted reliably because of inadequate data, poor test quality, an unavailable endpoint, or another limiting circumstance.

\medterm{Ungual} Relating to a fingernail or toenail. The term may describe a finding beneath, around, or within the nail, such as a bruise, fungal infection, wart, or pigmented lesion.

\medterm{Unguentum} A semisolid, greasy preparation intended for application to the skin or, when formulated for that purpose, to mucous membranes or the eyes.

\medterm{Unicornuate Uterus} A congenital condition in which only one side of the uterus develops fully, usually with one fallopian tube. Some people have no symptoms, but it can increase risks of infertility, pregnancy outside the uterus, miscarriage, or premature birth and may require specialist pregnancy care.

\medterm{Unilaminar Epithelium} An epithelium consisting of a single layer of cells. Its structure is adapted for functions such as diffusion, absorption, secretion, or filtration depending on the organ.

\medterm{Unilateral} Affecting, occurring on, or involving one side of the body or one paired organ, such as unilateral weakness or unilateral eye disease.

\medterm{Unilateral Gynecomastia} Enlargement of breast glandular tissue affecting one breast in a male patient; a new, firm, fixed, painful, or nipple-associated mass requires clinical assessment.

\medterm{Unilateral Microphthalmos} A congenital condition in which one eye is abnormally small. It may be associated with other ocular malformations and can cause substantial visual impairment.

\medterm{Uninodular Goiter} A goiter containing one dominant thyroid nodule. Evaluation generally includes thyroid-function testing and imaging, with further assessment to determine whether the nodule is benign or malignant.

\medterm{Uniparental Disomy} Inheritance of both copies of a chromosome or chromosome segment from one parent and
none from the other. It can cause disease through abnormal genomic imprinting,
homozygosity for a recessive variant, or chromosomal dosage effects.

\medterm{Unipolar Neuron} A neuron with a single process extending from its cell body. True unipolar neurons occur in some invertebrates; many sensory neurons in humans are more accurately described as pseudounipolar.

\medterm{Unit} A standard amount used to express measurement, dose, concentration, or a healthcare service in an agreed system.


\medterm{Unopposed Estrogen} Estrogen treatment given without a progestogen to someone with a uterus. Long-term use can over-stimulate the uterine lining and increase cancer risk, so protection is usually recommended when appropriate.

\medterm{Unresectable} Describing a tumor or lesion that cannot be completely removed safely by surgery because of its location, extent, invasion of vital structures, or the patient’s operative risk.

\medterm{Unresolved Losses} Important losses whose effects have not been fully recognized or worked through. They may involve death, health, relationships, work, home, or identity and can contribute to lasting grief, anger, guilt, anxiety, or difficulty moving forward.

\medterm{Unresponsive Coma} A state of profound unconsciousness in which a person cannot be awakened or follow commands. Causes include brain injury, stroke, seizures, poisoning, infection, and metabolic failure; emergency evaluation is required.
\medterm{Unsaturated Phytosterol} A plant sterol containing one or more carbon--carbon double bonds. Phytosterols can reduce intestinal cholesterol absorption, although clinical effects depend on intake, formulation, and patient context.

\medterm{Unstable Angina} Myocardial ischemia at rest, new-onset severe angina, or a crescendo pattern (more
frequent, longer, lower-threshold) without evidence of myocardial necrosis (normal
troponin). Pathophysiology: acute plaque rupture with thrombosis causing transient
severe coronary obstruction.

\medterm{Unsteadiness} Impaired balance or gait stability, which may result from vestibular disease, neurologic disorders, weakness, medication effects, visual impairment, or musculoskeletal problems.

\medterm{Untidy Wound} An irregular or contaminated wound with crushing, devitalized tissue, tissue loss, foreign material, or associated injury to vessels, nerves, tendons, or bone. It has increased infection risk and may require exploration, irrigation, debridement, and delayed or staged closure.

\medterm{Untranslated Region} A part of an RNA molecule that is not translated into protein; the 5-prime and 3-prime untranslated regions regulate RNA stability, localization, and translation.

\medterm{Unwell} A general description of feeling ill or having reduced health, without identifying a specific disease. Persistent or severe symptoms require assessment for an underlying cause.

\medterm{Upamostat} An investigational small-molecule serine-protease inhibitor studied for inflammatory and infectious diseases. It is not an established routine treatment, and its clinical role depends on trial evidence and regulatory approval.

\medterm{Upper Airway Biopsy} Removal of tissue from the nasal cavity, pharynx, larynx, or related upper airway for microscopic evaluation of suspected infection, inflammation, infiltrative disease, or tumor.

\medterm{Upper Airway Obstruction} Narrowing of the airway above the carina, causing inspiratory stridor, dyspnea, and
characteristic changes on flow-volume loops (flattening of the inspiratory limb). Causes
include tumors, stenosis, goiter, and vocal cord dysfunction. It is a medical emergency
when severe.

\medterm{Upper Airway Resistance Syndrome} Increased effort to breathe through a narrowed upper airway, often causing sleep disruption and daytime sleepiness.

\synonyms
UARS

\medterm{Upper Arm} The portion of the upper limb between the shoulder and elbow, containing the humerus and muscles, nerves, and vessels that serve the arm.

\medterm{Upper Endoscopy} Upper endoscopy: flexible endoscopic examination of the esophagus, stomach, and duodenum. See Esophagogastroduodenoscopy (EGD).
\medterm{Upper Esophageal Sphincter} A ring of muscle at the top of the esophagus that opens during swallowing and coordinates food passage.

\synonyms
UES

\medterm{Upper Extremity} The anatomic region extending from the shoulder through the arm, forearm, wrist, and hand. Disorders may involve its bones, joints, muscles, nerves, vessels, or skin.

\medterm{Upper Extremity Venous Doppler} An ultrasound examination of the arm and related veins that assesses compressibility and blood flow, commonly to evaluate suspected venous thrombosis.

\medterm{Upper Gastrointestinal Bleeding} Bleeding from esophagus, stomach, or duodenum. Causes: peptic ulcer, varices,
Mallory-Weiss tear, gastritis. Presentation: hematemesis, melena, hematochezia.
Management: ABCs, IV access, type and crossmatch, PPI infusion, endoscopy within 24
hours. Variceal bleeding: octreotide, antibiotics, band ligation.

\medterm{Upper GI And Small Bowel Series} A fluoroscopic x-ray examination in which swallowed contrast outlines the esophagus, stomach, duodenum, and small intestine to evaluate structure and transit.

\medterm{Upper GI Series} Fluoroscopic examination of the esophagus, stomach, and duodenum after oral barium
contrast. Single-contrast: looks for gross abnormalities (obstruction, masses).
Double-contrast (barium + gas): better mucosal detail for ulcers, erosions, and early
malignancies.

\medterm{Upper Leg} The portion of the lower limb between the hip and knee, containing the femur and the thigh muscles, nerves, and blood vessels.

\medterm{Upper Limb Disorders} Conditions affecting the shoulder, arm, elbow, forearm, wrist, or hand, including musculoskeletal, nerve, vascular, and work-related disorders.

\medterm{Upper Motor Neuron} A neuron in the brain or spinal cord that influences lower motor neurons controlling skeletal muscle. Damage can cause weakness, spasticity, increased reflexes, and pathologic reflexes.

\medterm{Upper Motor Neuron Syndrome} A pattern caused by injury to corticospinal pathways above the anterior horn cell or
cranial motor nucleus. Findings may include weakness with spasticity, hyperreflexia,
clonus, and an extensor plantar response, although acute lesions can initially produce
flaccidity.

\medterm{Upper Respiratory Infection} An infection of the nose, sinuses, pharynx, or larynx, most often caused by a virus and producing symptoms such as congestion, sore throat, cough, or runny nose.

\medterm{Upregulation} An increase in the activity, amount, or expression of a gene, receptor, enzyme, or signaling pathway in response to developmental, physiologic, or disease-related signals.

\medterm{Upright Labor Positions} Maternal positions during labor, such as walking, standing, squatting, kneeling, or sitting upright.
They may improve comfort and mobility and can be used when safe and acceptable to the
laboring person and clinical team.

\medterm{Uptake} The accumulation of a substance, such as a radiotracer, inside cells or tissues. It depends on blood flow, transport across cell membranes, metabolism, binding, and the tissue’s normal or abnormal activity.

\medterm{Urachal Carcinoma} A rare cancer arising from a leftover connection between the bladder and the navel, usually near the top of the bladder. It may cause blood or mucus in the urine, pelvic pain, or urinary changes.

\medterm{Urachal Cyst} A cystic remnant of the embryonic urachus between the bladder and umbilicus. It may be
asymptomatic or become infected, causing lower-abdominal pain, fever, umbilical
discharge, or urinary symptoms.

\medterm{Urachal Mucinous Cystic Tumor} A rare tumor arising from the urachal remnant between the bladder dome and umbilicus. It
may be benign or malignant, with mucinous histology, and requires surgical excision.

\medterm{Urachus} An embryologic connection between the fetal bladder and umbilicus that normally becomes the median umbilical ligament after birth. Persistent urachal remnants can form cysts, sinuses, or other complications.

\medterm{Uracil} A pyrimidine nitrogenous base found principally in RNA. It pairs with adenine and is also a component of several metabolic intermediates and nucleotide medicines.

\medterm{Uracil Mustard} An older nitrogen-mustard alkylating compound related to uracil and studied or used historically as an antineoplastic drug. Its clinical use is limited by toxicity and the availability of newer treatments.

\medterm{Uracil Nucleotides} Nucleotides containing uracil, including UMP, UDP, and UTP. They participate in RNA synthesis, carbohydrate metabolism, and activation of sugars for glycoconjugate production.

\medterm{Uraemia} The clinical syndrome resulting from accumulation of urea and other nitrogenous waste
products in the blood in advanced chronic kidney disease (CKD stage 4-5) and acute
kidney injury.

\medterm{Uraemic Syndrome} The collection of symptoms and biochemical abnormalities caused by severe accumulation of uraemic waste products in advanced kidney failure.

\medterm{Uranium} A naturally occurring radioactive heavy metal. Internal exposure can damage the kidneys, while inhaled or retained radioactive material can also expose tissues to ionizing radiation.

\medterm{Urate} The ionized form of uric acid, the end product of purine metabolism in humans. Excess serum urate can promote monosodium urate crystal deposition and gout.

\medterm{Urate-Lowering Therapy} Long-term treatment that reduces serum urate to dissolve monosodium urate deposits and
prevent gout flares and structural damage. Xanthine oxidase inhibitors and uricosuric
medicines are selected according to renal function, comorbidity, and treatment goals.

\medterm{Urates} Salts or ions of uric acid. Monosodium urate crystals can accumulate in joints in gout, and urate crystals may contribute to urinary stones.

\medterm{Urea} The principal nitrogen-containing waste product of protein metabolism, formed in the liver and excreted mainly by the kidneys.

\medterm{Urea Breath Test} A noninvasive test for active Helicobacter pylori infection: the person drinks labeled urea, and labeled carbon dioxide in exhaled breath indicates bacterial urease activity.

\medterm{Urea Cycle} A metabolic pathway occurring in both the mitochondria and cytoplasm of hepatocytes that
converts toxic ammonia (NH3), derived from amino acid catabolism, into nontoxic,
water-soluble urea for excretion by the kidneys.

\medterm{Urea Cycle Disorder} A group of inherited disorders caused by deficiency of an enzyme or transporter in the urea cycle,
leading to impaired ammonia clearance. Ornithine transcarbamylase deficiency is an
X-linked form; presentation and severity vary by the gene and residual enzyme activity.

\medterm{Urea Nitrogen Urine Test} A urine test measuring urea nitrogen excretion, interpreted with blood urea nitrogen, creatinine, diet, and urine volume to assess protein metabolism and renal handling.

\medterm{Ureaplasma} A genus of very small bacteria without a cell wall that can colonize the genital or urinary tract and is sometimes associated with urethritis or pregnancy complications.

\medterm{Ureaplasma Urealyticum} A cell-wall-free bacterium that can colonize the genital or urinary tract and is sometimes associated with nongonococcal urethritis, pregnancy complications, or infection in vulnerable newborns.

\medterm{Uremia} A clinical syndrome caused by advanced kidney failure and accumulation of waste products and other biologically active
substances. Symptoms may include nausea, itching, fatigue, confusion, bleeding, or
pericarditis; laboratory values and symptoms must be interpreted together.

\medterm{Uremic Encephalopathy} Brain dysfunction caused by the buildup of waste products in severe kidney failure. It may cause confusion, extreme sleepiness, muscle twitching, seizures, or coma. It is a medical emergency and usually requires urgent treatment of kidney failure, often including dialysis.

\medterm{Uremic Frost} Rare white, powdery crystals of urea that appear on the skin when severe untreated kidney failure causes very high waste levels in the blood. It indicates advanced illness requiring urgent medical assessment and treatment of the kidney failure.

\medterm{Uremic Pericarditis} Inflammation of the sac around the heart caused by severe kidney failure and buildup of waste products. It may cause sharp chest pain, a scratchy heart sound, or fluid around the heart; urgent kidney treatment is needed because tamponade can occur.

\medterm{Uremic Platelet Dysfunction} Reduced ability of platelets to stick together and stop bleeding because of advanced kidney failure. It can cause easy bruising, nosebleeds, heavy menstrual bleeding, or bleeding from the stomach or bowel and may improve when the kidney failure is treated.

\medterm{Uremic Stomatitis} Oral mucosal inflammation occurring in advanced kidney failure, presenting as white
plaques, ulceration, or bleeding of the gums and tongue due to accumulation of uremic
toxins.

\medterm{Ureter} The paired muscular tubes (25-30 cm) conveying urine from the renal pelvis to the
urinary bladder, crossing the pelvic brim at the pelvic inlet. Three anatomical
narrowings (pelviureteric junction, pelvic brim, vesicoureteric junction) are common
sites of ureteric stone impaction.

\medterm{Ureter Duplex} A congenital duplication of the ureter, complete or partial, draining one kidney. It may
be asymptomatic or associated with vesicoureteral reflux, ureterocele, obstruction,
recurrent urinary infection, or impaired renal drainage.

\medterm{Ureter Fissus} A congenital longitudinal splitting or duplication anomaly of the ureter. Its clinical importance depends on the anatomy and whether it causes reflux, obstruction, infection, or impaired drainage.

\medterm{Ureter Obstruction} Blockage of urine flow from a kidney to the bladder, commonly caused by a stone, stricture, tumor, or external compression and potentially leading to hydronephrosis, infection, or kidney injury.

\medterm{Ureteral Calculi} Stones that form in or pass through a ureter, the tube carrying urine from a kidney to the bladder.

\medterm{Ureteral Catheters} Thin tubes passed into the ureter to drain urine, inject contrast, measure pressure, or guide procedures. Complications include infection, bleeding, perforation, migration, and obstruction.

\medterm{Ureteral Obstruction} Blockage of urine flow causing proximal dilation and increased pressure. Acute causes
renal colic with flank pain. Chronic leads to hydronephrosis and CKD. Emergent
decompression with stent or nephrostomy for infected obstructed kidney. Diagnosis by CT
or ultrasound showing hydronephrosis.

\synonyms
Blocked Ureter

\medterm{Ureteral Perforation} A tear or hole in the ureter, sometimes caused by a stone, trauma, surgery, or an endoscopic
procedure. Urine may leak into surrounding tissues, causing pain, infection, or inflammation;
prompt imaging and urologic assessment are required.

\medterm{Ureteral Polyp} Usually a benign fibroepithelial growth arising from the ureter. It may cause hematuria, flank pain, or obstruction and must be distinguished from urothelial malignancy when evaluated.

\medterm{Ureteral Retrograde Brush Biopsy} A brush passed through a ureteroscope to collect cells from the lining of the ureter or renal pelvis for cytologic evaluation of suspected urothelial cancer.

\medterm{Ureteral Stenosis} Narrowing of the ureter caused by scar tissue, prior surgery, radiation, stones, inflammation, or congenital disease, impairing urine drainage and possibly causing hydronephrosis or recurrent infection.

\medterm{Ureteral Stent} A thin tube placed inside the ureter to keep urine flowing from the kidney to the bladder when a stone, narrowing, tumor, or swelling blocks the passage. It can cause urinary urgency, discomfort, blood in the urine, infection, or mineral buildup and is usually temporary.

\medterm{Ureteral Stent (Double-J Stent)} A tube placed in the ureter from the renal pelvis to the bladder to maintain patency,
with pigtail coils at both ends (hence double-J). Indications: ureteral obstruction
(stones, stricture, malignancy), after ureteroscopy (prevent post-op obstruction), and
as a pre-treatment for complex ureteral stones.

\synonyms
Double-J Stent

\medterm{Ureteral Stricture} Narrowing of the ureter caused by fibrosis, surgery, instrumentation, radiation,
infection, malignancy, or congenital disease. It can produce upstream hydronephrosis,
pain, infection, and declining renal function.

\medterm{Uretereroenterostomy} A urinary diversion procedure connecting a ureter to a segment of intestine; this spelling is a source variant of \textit{ureteroenterostomy}.

\medterm{Ureteric Perforation} A tear or hole in the ureter that can allow urine to leak into surrounding tissue, usually after instrumentation, surgery, trauma, or a stone; it needs prompt urologic assessment.

\medterm{Ureteritis} Inflammation of the ureter, usually related to infection, obstruction, stone passage,
instrumentation, or adjacent inflammatory disease. It may cause flank or abdominal pain,
urinary symptoms, hematuria, and fever.

\medterm{Ureterocele} A balloon-like dilation of the lower ureter where it enters the bladder, usually congenital; it can obstruct urine flow and cause infection, hydronephrosis, reflux, stones, or impaired renal function.

\medterm{Ureterolithiasis} The presence of kidney stones (calculi) within the ureter, causing ureteral obstruction
and renal colic. Stones form from calcium oxalate/phosphate (most), uric acid, struvite
(infection), and cystine.

\medterm{Ureteropelvic Junction Obstruction} Impaired drainage from the renal pelvis into the proximal ureter, caused by intrinsic
narrowing, aberrant vessels, scarring, or functional peristaltic abnormality. It may
cause hydronephrosis, flank pain, infection, stones, or loss of renal function.

\medterm{Ureteroscope} A slender endoscope used to inspect and treat the ureter and, when advanced, the renal pelvis.

\medterm{Ureteroscopes} Small endoscopes used to inspect the ureter and renal collecting system and to perform procedures such as stone removal, biopsy, or stent placement.

\medterm{Ureteroscopy} Ureteroscopy uses a small endoscope passed through the urinary tract to inspect, biopsy, or treat ureteral or kidney stones and other lesions.
the urethra, bladder, and into the ureter or renal collecting system. Used primarily for
the diagnostic evaluation and therapeutic management of ureteral and renal calculi,
ureteral strictures, and upper tract urothelial carcinoma (UTUC).

\medterm{Ureterostomy Procedure} Surgery that creates an opening from a ureter to the skin or another urinary conduit to divert urine when normal drainage to the bladder is impossible or must be bypassed.

\medterm{Urethan} An older name for ethyl carbamate, a chemical compound with carcinogenic potential in experimental and exposure settings. It is not a routine therapeutic medicine.

\medterm{Urethra} The tube conveying urine from the bladder to the exterior. Male urethra (18-20 cm):
prostatic, membranous, and spongy (penile) parts; carries urine and semen; longer and
more tortuous than female. Female urethra (3-4 cm): short, straight, predisposing to
ascending urinary tract infections.

\medterm{Urethra Disorder} A condition affecting the urethra, such as infection, inflammation, stricture, trauma, congenital abnormality, or tumor, that may cause pain, discharge, weak stream, or retention.

\medterm{Urethra, Diseases Of And Injury To} Diseases or injuries affecting the urethra, including infection, narrowing, blockage, tears, or other structural problems.

\medterm{Urethracele} A herniation or prolapse of the urethra into the anterior vaginal wall, often occurring
together with a cystocele (bladder prolapse). Urethracele is caused by weakening of the
pubocervical fascia, often from childbirth, aging, or increased intra-abdominal
pressure.

\medterm{Urethral Carcinoma} A rare cancer arising from the urethra that may cause bleeding, a lump, discharge, pain, or difficult urination; diagnosis requires examination, imaging, and biopsy.

\medterm{Urethral Catheter} A tube inserted through the urethra into the bladder to drain urine, obtain a specimen, instill medication, or monitor output. Urinary infection and urethral trauma are important risks.


\medterm{Urethral Diverticulum} A periurethral outpouching that communicates with the urethral lumen, usually arising
from infected or obstructed periurethral glands. It may cause dysuria, post-void
dribbling, recurrent urinary infection, dyspareunia, or a palpable anterior-vaginal-wall
mass.

\medterm{Urethral Fistula} An abnormal channel between the urethra and another epithelial surface, often the vagina or skin. It can cause continuous urine leakage, recurrent infection, and local irritation.

\medterm{Urethral Hypermobility} Excessive movement of the urethra when abdominal pressure rises, contributing to urine leakage during coughing, lifting, or exercise.

\medterm{Urethral Infection} Infection involving the urethra, commonly caused by bacteria, including sexually transmitted organisms.
Symptoms may include burning urination, discharge, frequency, or pelvic discomfort; testing and
appropriate treatment help prevent complications and transmission.

\medterm{Urethral Meatus} The external opening of the urethra through which urine leaves the body. Its position may be abnormal in conditions such as hypospadias or epispadias.

\medterm{Urethral Meatus, Female} The external opening of the female urethra, located in the vestibule anterior to the vaginal opening. Its position and appearance may be relevant in urinary symptoms, congenital variation, or examination.

\medterm{Urethral Obstruction} Partial or complete blockage of urine leaving the bladder, often from stricture, stone, prostate disease, tumor, or a congenital valve; acute retention requires prompt treatment.

\medterm{Urethral Pressure Profile} A test that measures pressure along the urethra, the tube carrying urine out of the body. A small catheter is used to assess how well the urethra closes and may help investigate some types of urinary leakage or difficulty emptying the bladder.

\medterm{Urethral Pressures} Measurements of pressure within or along the urethra used to assess urethral closure and urinary continence.

\medterm{Urethral Sling} Synthetic or fascial tape supporting urethra for stress incontinence. Mid-urethral sling
(TVT, TOT) standard. Minimally invasive. Complications: erosion, urgency, retention.

\medterm{Urethral Sphincter} The muscular mechanism that helps maintain urinary continence by compressing the urethra. It includes smooth and striated muscle components controlled by autonomic and voluntary pathways.

\medterm{Urethral Stenosis} Scar-related narrowing of the urethral channel that restricts urinary flow and can cause spraying, weak stream, recurrent infection, retention, or bladder damage.

\medterm{Urethral Stricture} Narrowing of the urethra caused by scar tissue. It may produce a weak urinary stream, spraying, straining, incomplete
emptying, infections, or urinary retention and can result from injury, instrumentation,
infection, or inflammatory disease.

\medterm{Urethral Syndrome} Urinary frequency, urgency, dysuria, or discomfort without a clear bacterial urinary-tract infection; causes may include irritation, pelvic-floor dysfunction, or inflammation.

\medterm{Urethritis} Urethritis is inflammation of the urethra, commonly from sexually transmitted infection, causing dysuria, irritation, or discharge and requiring evaluation of partners and complications.
(Chlamydia trachomatis, Neisseria gonorrhoeae) or non-infectious causes. Presents with
dysuria, urethral discharge, and meatal itching. Diagnosis: urine nucleic acid
amplification testing (NAAT).

\medterm{Urethrocele} Bulging or downward displacement of the urethra, usually into the front wall of the vagina. It may cause pressure, urinary leakage, difficulty emptying the bladder, discomfort, or no symptoms.

\medterm{Urethrography} Contrast imaging of urethra. Evaluates urethral strictures, trauma, anomalies. Contrast imaging of the urethra, usually with X-rays. It can show narrowing, leaks, injury, fistulas, or congenital abnormalities.
\medterm{Urethroplasty} Surgery that rebuilds or widens the urine tube, most often to treat a narrowing called a stricture. It may use nearby tissue or a graft and aims to improve urine flow.

\medterm{Urethroscopy} Endoscopic examination of the urethra using a urethroscope, sometimes combined with treatment of a stricture or other lesion.

\medterm{Urge Urinary Incontinence} Involuntary leakage of urine, significantly affecting psychological, social, and physical well-being. See urinary incontinence.
\medterm{Urgency} A sudden, compelling desire to urinate that is difficult to defer—a lower urinary tract
symptom. Causes: overactive bladder (detrusor overactivity, often idiopathic), urinary
tract infection (urgency + dysuria), bladder outlet obstruction (BPH), neurogenic
bladder, interstitial cystitis, and pelvic floor disorders.

\medterm{URI} URI usually means upper respiratory infection, commonly a viral infection of the nose, throat, or upper airways causing congestion, cough, and sore throat.

\medterm{Uric Acid} A waste product formed when the body breaks down purines. High blood levels can cause gout or kidney stones, although many people with high levels have no symptoms.
\medterm{Uric Acid  - Blood} A blood test measuring uric acid, a purine-breakdown product. The result helps assess gout, kidney stones, kidney function, or effects of some cancer treatments.
\medterm{Uric Acid Nephropathy} Acute kidney injury from uric acid crystal precipitation in tubules. Occurs in tumor
lysis syndrome and severe hyperuricemia. Prevent with hydration and rasburicase.
Alkalinization of urine may be harmful as uric acid precipitates at alkaline pH unlike
other stone types.

\medterm{Uric Acid Stones} Kidney stones composed primarily of uric acid. They form more readily in acidic urine and are
associated with hyperuricemia, gout, high purine intake, or chronic diarrhea. Some can dissolve when urine is alkalinized.

\medterm{Uric Acid Urine Test} A urine test measuring uric acid excretion, used in evaluation of gout, kidney stones, uric-acid overproduction, and selected metabolic disorders.

\medterm{Uricosuric Drug} A medicine that increases renal excretion of uric acid, such as probenecid, used in selected patients with gout.

\medterm{Uridine} A nucleoside composed of uracil linked to ribose. It is a building block of RNA and a precursor of uridine nucleotides involved in carbohydrate and membrane metabolism.

\medterm{Uridine Diphosphate} A nucleotide consisting of uridine and two phosphate groups. UDP participates in glycogen metabolism and transfers activated sugars during glycosylation reactions.

\medterm{Uridine Monophosphate} A nucleotide consisting of uridine and one phosphate group. UMP is a component of RNA and an intermediate in pyrimidine-nucleotide metabolism.

\medterm{Uridine Triacetate} An oral prodrug that supplies uridine and is used as an emergency antidote for severe or life-threatening toxicity from fluorouracil or related fluoropyrimidines, under specialist direction.

\medterm{Uridine Triphosphate} A high-energy uridine nucleotide used in RNA synthesis and in the formation of activated sugar donors for glycogen and glycoprotein production.

\medterm{Urinal} A portable container used to collect urine from someone who cannot easily reach a toilet. It should be emptied, cleaned, and stored safely to reduce spills, odor, and infection risk.

\medterm{Urinalysis} Urinalysis examines urine appearance, chemistry, and sediment for evidence of infection, kidney disease, bleeding, diabetes, dehydration, or other systemic conditions.
(bacteria), protein, glucose, ketones, bilirubin, urobilinogen, blood, and hemoglobin.

\medterm{Urinary} Relating to urine or to the organs and processes involved in making, storing, and passing urine, including the kidneys, ureters, bladder, and urethra.
\medterm{Urinary Bladder} A hollow, distensible muscular organ located in the pelvic cavity posterior to the pubic
symphysis, serving as a reservoir for urine.

\medterm{Urinary Bladder Disorder} A disease affecting bladder storage or emptying, including infection, overactivity, pain syndrome, stones, tumor, neurologic dysfunction, or outlet obstruction.

\medterm{Urinary Bladder, Diseases Of} Disorders of the bladder, including infection, stones, tumors, neurogenic dysfunction, inflammation, obstruction, and impaired emptying.

\medterm{Urinary Calculi} Stones formed from crystallized minerals in the urinary tract. They may cause colicky pain, hematuria, obstruction, infection, or kidney injury depending on their size and location.

\medterm{Urinary Casts} Small tube-shaped particles formed inside kidney tubules and passed in urine. Their appearance can provide clues about kidney inflammation, bleeding, infection, or damage.

\medterm{Urinary Catheter} A hollow, flexible tube inserted through the urethra into the bladder to drain urine; it may be used intermittently or left in place.

\medterm{Urinary Catheterization} Insertion of a tube into the bladder to drain urine, measure output, obtain a specimen, or deliver treatment; infection, trauma, and blockage are possible complications.

\medterm{Urinary Catheters} Urinary catheters drain the bladder when a person cannot empty it normally or accurate urine measurement is needed, with risks including infection, trauma, blockage, and discomfort.

\medterm{Urinary Diversion} Surgical or catheter-based rerouting of urine when the normal passage from the kidneys to the urethra cannot be used, often to a stoma or bowel reservoir.

\medterm{Urinary Fistula} An abnormal connection between part of the urinary tract and another organ or the skin, allowing urine to leak outside its normal pathway. Causes include surgery, childbirth injury, inflammation, and cancer.

\medterm{Urinary Frequency} The need to urinate more often than usual, which may result from infection, bladder irritation, increased urine production, or other urinary disorders.

\medterm{Urinary Hesitation} Difficulty starting or maintaining urination, which may result from prostate enlargement, urethral narrowing, medications, nerve disorders, pain, or impaired bladder contraction.

\medterm{Urinary Incontinence} Unwanted leakage of urine. It may happen with coughing or exercise, a sudden strong urge, an overfull bladder, difficulty reaching the toilet, or nerve and muscle problems; assessment can identify treatable causes and suitable exercises, medicines, devices, or surgery.

\synonyms
Bladder Control, Loss Of
Incontinence, Urinary
Loss Of Bladder Control

\medterm{Urinary Incontinence Products} Urinary-incontinence products such as pads, protective underwear, catheters, and mattress covers manage leakage but do not treat its underlying cause.

\medterm{Urinary Lithiasis} The formation or presence of stones anywhere in the urinary tract, including the kidneys, ureters, bladder, or urethra.

\medterm{Urinary Retention} Inability to completely empty the bladder, which can be acute (sudden, painful,
inability to void) or chronic (gradual, often painless, with overflow incontinence).
Causes: benign prostatic hyperplasia, urethral stricture, neurogenic bladder,
medications (anticholinergics), and pelvic organ prolapse.

\medterm{Urinary Retraining} Scheduled toilet visits and bladder-control techniques used to reduce urgency, frequent urination, or leakage. A clinician may combine it with pelvic-floor exercises, fluid changes, or treatment of the underlying cause.
\medterm{Urinary Stasis} Slowing or stopping of normal urine flow, which can promote urinary infection, stone formation, or obstruction-related kidney injury.

\medterm{Urinary System} The kidneys, ureters, bladder, and urethra, which produce, transport, store, and eliminate urine while regulating fluid, electrolytes, acid--base balance, and blood pressure.

\medterm{Urinary Tract} The kidneys, ureters, bladder, and urethra. Together these organs make urine, carry it to the bladder, store it, and remove it from the body.
\medterm{Urinary Tract Disorder} Any disease of the kidneys, ureters, bladder, or urethra, including infection, stones, obstruction, congenital abnormalities, cancer, and functional or neurologic disorders.

\medterm{Urinary Tract Infection} Infection of the urinary tract, classified as lower infection such as cystitis or urethritis,
or upper infection such as pyelonephritis. Symptoms vary; dysuria, frequency, urgency, fever,
flank pain, or nausea may occur. Escherichia coli is a common cause, but organisms vary by
setting and patient factors.

\synonyms
Bladder Infection (Urinary Tract Infection)

\medterm{Urinary Tract Infection In Pregnancy} A urinary infection or asymptomatic bacterial growth occurring during
pregnancy. Pregnancy increases the risk of upper urinary infection, so screening and treatment follow pregnancy-specific
clinical guidance.

\medterm{Urinary Tract Neoplasm} A benign or malignant growth in the kidneys, ureters, bladder, or urethra; evaluation depends on its site, symptoms, tissue type, and extent of spread.

\medterm{Urinary Urgency} A sudden, difficult-to-delay need to urinate. It may occur with urinary infection, bladder irritation,
overactive bladder, stones, neurologic disease, or other conditions. Urgency with fever, flank
pain, blood in the urine, or inability to urinate requires medical assessment.

\medterm{Urination} The biological process of releasing urine from the bladder through the urethra out of
the body; also known as micturition.

\synonyms
micturition.

\medterm{Urination Disorder} Abnormal storage or emptying of urine causing frequency, urgency, dysuria, incontinence, hesitancy, weak stream, retention, or excessive or reduced urine output.

\medterm{Urination, Difficulty} Trouble starting, maintaining, or emptying urine, often described as hesitancy, weak stream, straining, or incomplete emptying; causes include obstruction, infection, medicines, and nerve or muscle disorders.

\medterm{Urine} The liquid waste produced by the kidneys, containing water and dissolved substances eliminated from the blood.

\medterm{Urine 24-Hour Volume} Twenty-four-hour urine volume is the total urine produced in a day and helps assess hydration, kidney function, endocrine disease, or abnormal urine production.

\medterm{Urine Albuminuria} An abnormal amount of the blood protein albumin in urine. Persistent albumin loss can be an early sign of kidney disease and also increases the risk of worsening kidney function and cardiovascular disease.

\medterm{Urine Alkalinization} Deliberate increase of urine pH, usually with intravenous sodium bicarbonate, to enhance elimination of selected weak acids such as salicylate; it requires monitoring of pH, potassium, fluid status, and acid-base balance.

\medterm{Urine Analysis} Laboratory examination of urine for cells, blood, proteins, glucose, microorganisms, and other substances. Laboratory testing of urine for cells, blood, protein, sugar, microbes, and other clues to illness.
\medterm{Urine Chemistry} Measurement of chemical constituents in urine, such as glucose, protein, electrolytes, ketones, blood, and pH, to evaluate kidney function and systemic disease.

\medterm{Urine Concentration Test} A test of the kidney’s ability to concentrate urine, assessed by urine osmolality or specific gravity after controlled fluid conditions when appropriate.

\medterm{Urine Conductivity} Urine conductivity measures how readily urine carries electrical current and reflects its dissolved electrolyte and solute concentration.

\medterm{Urine Culture} A laboratory test that grows microorganisms from a urine specimen to identify urinary-tract infection and guide antimicrobial selection. Results may report the organism, quantity of growth, and antimicrobial susceptibility. Proper collection is important because contamination can produce misleading growth; results must be interpreted with symptoms and urinalysis.

\medterm{Urine Culture - Catheterized Specimen} A urine culture performed on a specimen collected through urinary catheterization, usually when a clean-catch sample cannot be obtained or when catheter-associated infection is being assessed. Collection technique and timing affect interpretation because colonization and contamination may cause bacterial growth without symptomatic infection.

\medterm{Urine Cytology} Microscopic examination of cells shed into urine, mainly used to help detect high-grade urothelial carcinoma and some other abnormalities. A negative result does not exclude all urinary-tract cancers.

\medterm{Urine Drainage Bags} Urine drainage bags collect urine from a catheter or urinary diversion and require correct positioning, emptying, hygiene, and monitoring for blockage or infection.

\medterm{Urine Formation} The process by which the kidneys make urine: they filter blood, return needed water and substances to the body, and add selected wastes and extra ions to the fluid that leaves as urine.

\medterm{Urine HCG (Pregnancy Test)} Rapid point-of-care test detecting human chorionic gonadotropin in urine. Sensitivity:
detects as low as 25 mIU/mL. Positive as early as 12-14 days after conception (before
missed period). Results in 3-5 minutes. False negative: dilute urine, very early
pregnancy, ectopic pregnancy.

\synonyms
Pregnancy Test

\medterm{Urine Hemoglobin} Hemoglobin detected in urine, which may reflect intravascular hemolysis or contamination with intact red blood cells. Interpretation requires comparison with urine microscopy and clinical findings.

\medterm{Urine Osmolality} Concentration of urine relative to plasma normally 50 to 1200 mOsm/kg. High greater than
500 in prerenal states reflecting intact concentrating ability. Low less than 350
indicates impaired concentrating ability in diabetes insipidus, chronic kidney disease,
and excessive fluid intake. More precise than specific gravity.

\medterm{Urine PH} A measurement of urine acidity or alkalinity. It varies with diet, metabolism, medications, infection, and kidney function and can influence the formation of some urinary stones.

\medterm{Urine PH Test} A test measuring urine acidity or alkalinity, useful in evaluating acid-base handling, urinary infection, kidney stones, and some metabolic conditions.

\medterm{Urine Protein Dipstick Test} A screening test using a reagent strip to detect mainly albumin in urine; positive results require confirmation and quantitative assessment when clinically indicated.

\medterm{Urine Protein Electrophoresis Test} A laboratory test that separates urine proteins by electrical charge and size to help detect abnormal proteins, including monoclonal light chains associated with plasma-cell disorders.

\medterm{Urine Retention} Inability to empty the urinary bladder completely or at all. Acute retention can cause severe pain and requires prompt medical assessment.

\medterm{Urine Sediment} Microscopic examination of centrifuged urine providing diagnostic information about
renal pathology. Normal findings include few RBCs, WBCs, and hyaline casts.

\medterm{Urine Sodium} Concentration of sodium in urine measured in mEq/L. Low less than 20 in prerenal states
reflecting avid tubular reabsorption. High greater than 40 in intrinsic renal disease,
diuretic use, and salt-wasting nephropathies. Should be interpreted in context of volume
status, medications, and concurrent therapies.

\medterm{Urine Specific Gravity Test} A test estimating urine density relative to water and reflecting solute concentration; results help assess hydration and renal concentrating ability.

\medterm{Urine Toxicology} Laboratory analysis of urine for medicines, drugs, or their metabolites. Results require interpretation because detection does not
always establish impairment or the timing of exposure.

\medterm{Urinometer} A hydrometer used to measure the specific gravity of urine.

\medterm{Urobilinogen} A colourless product formed in the intestine from bilirubin and partly reabsorbed and excreted in urine; abnormal levels may provide clues to liver disease or haemolysis.


\medterm{Urochordata} A subphylum of chordates, also called tunicates, whose members include sea squirts. They are primarily biological research organisms rather than human pathogens.

\medterm{Urodynamic Testing} Suite of diagnostic procedures designed to evaluate the physiological function of the
lower urinary tract (bladder and urethral sphincter) during filling, storage, and
emptying phases.

\medterm{Urodynamics} A set of tests that assess the function of the lower urinary tract (bladder and
urethra).

\medterm{Uroflow} A test that measures how quickly and steadily urine leaves the body during urination. The result can suggest blockage, weak bladder muscles, or other emptying problems, but it is interpreted with symptoms and sometimes additional tests.

\medterm{Uroflowmetry} A noninvasive test measuring the rate and pattern of urine flow during voiding to evaluate obstruction, weak detrusor function, or voiding dysfunction.
and testing conditions; a low or abnormal flow may reflect obstruction, weak bladder
contraction, or inadequate effort and requires clinical correlation.

\medterm{Urogenital} Relating to the urinary and reproductive systems or structures involving both. The term may describe organs, infections, cancers, birth conditions, injuries, symptoms, tests, or treatments affecting these closely connected body systems.

\medterm{Urogenital Fistula} An abnormal connection between the urinary and genital tracts, such as a vesicovaginal or ureterovaginal fistula. It commonly causes continuous urinary leakage and requires specialist evaluation.

\medterm{Urogenital System} The related urinary and reproductive organs and their associated ducts, vessels, and supporting tissues. The systems share developmental origins and anatomic relationships.

\medterm{Urogenital Tract} The combined urinary and reproductive organs and passages. These include kidneys, urine tubes, bladder, urethra, reproductive organs, and connected structures that develop partly from related tissues before birth.

\medterm{Urography} Imaging of the urinary tract, traditionally using X-rays after contrast dye is given. It can show the kidneys, urine tubes, bladder, narrowing, stones, leaks, blockage, or other structural problems.

\medterm{Urokinase} A plasminogen activator produced by the body and also used as a thrombolytic medicine to help dissolve blood clots.

\medterm{Urolithiasis} Formation of stones in the urinary tract, including the kidneys, ureters, bladder, or urethra. Stone composition
varies by population and may include calcium oxalate, calcium phosphate, uric acid,
struvite, or cystine.

\medterm{Urologic Diagnostic Yield} The proportion of urologic diagnostic procedures or tests that provide a definitive and
clinically useful result.

\medterm{Urologist} A doctor who specializes in the urinary system and, in many practices, the male reproductive system. Urologists diagnose and treat kidney stones, infections, urine leakage, prostate problems, cancers, infertility, and conditions requiring procedures or surgery.

\medterm{Urology} Surgical and medical specialty concerned with the study, diagnosis, and management of
diseases affecting the male and female urinary tract (kidneys, ureters, urinary bladder,
urethra) and the male reproductive system (prostate, testes, epididymis, seminal
vesicles, penis).

\medterm{Urosepsis} Sepsis caused by infection of the urinary tract or urinary system, often associated with pyelonephritis, obstruction, catheterization, or an infected urinary calculus.
The infection may progress from a localized urinary source to bacteremia, systemic inflammation, organ dysfunction, and septic shock.
Clinical findings may include fever or hypothermia, flank pain, dysuria, hypotension, altered mental status, oliguria, or elevated lactate.
Diagnosis combines clinical assessment with urine and blood cultures, laboratory testing, and imaging when obstruction or a complicated source is suspected.
\medterm{Urostomy Leak} Escape of urine beneath or around a urostomy appliance or from a surgical connection. Persistent leakage can irritate the skin and may signal poor appliance fit, obstruction, infection, or a surgical complication.

\medterm{Urostomy Procedure} Surgery that diverts urine through a stoma into an external pouch or conduit, commonly after bladder removal or when the lower urinary tract cannot safely transport urine.

\medterm{Urostomy Prolapse} Protrusion of a segment of bowel used as a urostomy stoma beyond its usual level. It may cause bleeding, swelling, appliance problems, or impaired urine drainage and should be assessed clinically.

\medterm{Urostomy Stenosis} Narrowing of a urinary stoma or conduit that impedes drainage, causing reduced urine output, leakage, infection, or kidney-pressure changes and sometimes requiring revision.

\medterm{Urothelial Carcinoma} Urothelial carcinoma, previously called transitional cell carcinoma, is the most common
malignancy of the urinary tract, arising from the urothelium that lines the renal
pelvis, ureters, bladder, and proximal urethra.

\medterm{Urothelium} A specialized stretchable lining of the renal pelvis, ureters, bladder, and part of the urethra. It forms a protective barrier between urine and underlying tissues and allows the bladder to expand.
\medterm{Ursodeoxycholic Acid} A hydrophilic bile acid used for primary biliary cholangitis and selected cholestatic liver diseases; it may also dissolve certain small cholesterol gallstones when bile flow is preserved.

\medterm{Urticaria} A skin reaction causing transient, itchy raised welts that change location and usually resolve within 24 hours.
It may occur with angioedema, which is deeper swelling often affecting the lips, eyelids,
tongue, or extremities.

\medterm{Urticaria Multiforme} A hypersensitivity reaction in children, often triggered by infections or medications,
characterised by widespread, annular or polycyclic erythematous plaques with dusky,
ecchymotic, or targetoid centres and pseudovesication, differing from true urticaria by
the presence of dusky discolouration and prolonged lesion duration.

\medterm{Urticaria Pigmentosa} A cutaneous form of mastocytosis in which excessive mast cells accumulate in the skin,
producing tan or brown macules or papules that may swell and itch when rubbed. Flushing,
abdominal symptoms, or systemic involvement can occur.

\medterm{Urticarial Vasculitis} Inflammation of small blood vessels associated with hive-like lesions that often last longer than ordinary urticaria and may burn, bruise, or leave residual discoloration.

\medterm{Usher Syndrome} An inherited disorder combining sensorineural hearing loss with retinitis pigmentosa,
the most common cause of combined deaf-blindness. Three types are distinguished by
severity and age of onset; management requires coordinated audiological,
ophthalmological, and genetic care.

\medterm{Ustekinumab} Ustekinumab is a monoclonal antibody blocking interleukin-12 and interleukin-23 signaling, used for psoriasis, psoriatic arthritis, inflammatory bowel disease, and other approved conditions; infection risk and hypersensitivity require monitoring.

\medterm{Ustilago} A genus of smut fungi that infect plants, especially grasses and cereals. Human infection is rare and generally occurs as an opportunistic infection in people with impaired immunity.

\medterm{Usual Interstitial Pneumonia} A pattern of permanent scarring in the lungs, often seen in idiopathic pulmonary fibrosis. Scarring is usually greatest near the outer and lower parts of the lungs and can form small spaces called honeycombing. It causes increasing breathlessness and cough and is assessed with lung tests and CT imaging.

\medterm{Ut Dict.} An abbreviation of the Latin phrase “ut dictum,” meaning “as stated” or “as directed.” In medical notes it refers the reader to information already given or to an instruction previously specified.

\medterm{Uterine (Fallopian) Tubes} Paired muscular tubes extending from the uterus toward the ovaries. They transport the oocyte and are the usual site of fertilization; blockage can contribute to infertility or ectopic pregnancy.

\synonyms
fallopian tubes; uterine tubes

\medterm{Uterine Adnexae} The uterine adnexa are structures beside the uterus, chiefly the fallopian tubes, ovaries, and supporting tissues.

\medterm{Uterine Atony} Failure of the uterus to contract adequately after delivery, resulting in postpartum
hemorrhage. The most common cause of PPH. Risk factors include prolonged labor, multiple
gestation, macrosomia, oxytocin augmentation, and uterine fibroids. Diagnosis is
clinical: boggy, soft uterus with excessive bleeding.

\medterm{Uterine Balloon Tamponade} A treatment for severe bleeding after childbirth when medicines do not control bleeding. A soft balloon is placed inside the uterus and filled with sterile fluid to press against the uterine wall and slow or stop blood loss while other treatment is arranged.

\medterm{Uterine Bleeding} Uterine bleeding is bleeding from the uterus and may result from menstruation, pregnancy complications, fibroids, polyps, hormonal disorders, cancer, or medicines.

\medterm{Uterine Cancer} Cancer arising in the uterus, most commonly from the endometrium. Abnormal vaginal bleeding,
bleeding after menopause, pelvic pain, or unusual discharge requires evaluation; early diagnosis
often permits effective treatment.

\medterm{Uterine Cavity} The uterine cavity is the internal space of the uterus lined by endometrium and is where implantation and pregnancy occur.

\medterm{Uterine Contraction} Tightening of the uterine muscle, occurring normally during menstruation and labor and sometimes abnormally with pain, pregnancy complications, or medication effects.

\medterm{Uterine Fibroid (Leiomyoma)} A benign tumor arising from the smooth muscle of the uterus. Fibroids may be
submucosal, intramural, or subserosal and may cause heavy menstrual bleeding, pelvic
pressure, pain, urinary symptoms, or reproductive problems.

\synonyms
Leiomyoma

\medterm{Uterine Fibroid Embolization} A minimally invasive procedure in which particles block blood flow to uterine fibroids, causing them to shrink. It can reduce heavy bleeding and pressure symptoms but may cause pain, infection, or effects on fertility.
\medterm{Uterine Fibroids} Benign tumors of uterine smooth muscle, also called leiomyomas. They may cause heavy menstrual bleeding, pelvic pressure, pain, infertility, or no symptoms.

\synonyms
leiomyomas.
Polyps Uterus (Uterine Fibroids)
Tumors Uterine (Uterine Fibroids)
Womb Growths (Uterine Fibroids)

\medterm{Uterine Fistula} A uterine fistula is an abnormal passage connecting the uterus with another organ or the skin, often after surgery, childbirth injury, infection, or cancer.

\medterm{Uterine Fornix} The uterine or vaginal fornix is a recess around the cervix where the upper vagina surrounds the projecting lower uterus.

\medterm{Uterine Gland} A small gland in the lining of the uterus that produces fluid and helps support the early embryo. Its activity changes during the menstrual cycle and pregnancy.

\medterm{Uterine Growths} Uterine growths may include fibroids, polyps, adenomyosis, or rarely cancer. They can cause heavy or irregular bleeding, pelvic pressure, pain, infertility, or no symptoms; examination and imaging determine the type and treatment.

\medterm{Uterine Hemorrhage} Heavy or abnormal bleeding from the uterus, occurring during pregnancy, after childbirth, or at other
times. Causes include pregnancy complications, fibroids, hormonal disorders, and malignancy;
severe bleeding with dizziness or weakness is an emergency.

\medterm{Uterine Inversion} A rare, life-threatening complication of childbirth in which the uterus turns partly or completely inside out and may protrude into or beyond the vagina. It can cause severe bleeding, shock, and intense pain and requires immediate obstetric treatment to restore the uterus and control blood loss.

\medterm{Uterine Leiomyosarcoma} A rare malignant tumor of uterine smooth muscle that may cause abnormal bleeding, pelvic pain, or an enlarging mass; diagnosis is based on tissue examination and staging.

\medterm{Uterine Monitoring} Assessment of uterine contractions and, in pregnancy, their frequency, duration, and strength to evaluate labor progress, induction response, or excessive uterine activity.

\medterm{Uterine Myomectomy} Surgical removal of uterine fibroids while preserving the uterus; it may be performed through hysteroscopy, laparoscopy, or open surgery depending on fibroid size and location.

\medterm{Uterine Obstruction} Blockage that prevents normal passage through the uterus or cervix, which may cause menstrual outflow obstruction, infertility, pregnancy complications, or retained blood or tissue.

\medterm{Uterine Perforation} A full-thickness breach of the uterine wall, most commonly occurring during
instrumentation such as dilation and curettage or hysteroscopy, potentially causing
hemorrhage and requiring surgical repair.

\medterm{Uterine Polyps} Growths arising from the inner lining of the uterus that may cause irregular bleeding, heavy periods, or fertility problems.

\synonyms
Polyps, Endometrial
Polyps, Uterine

\medterm{Uterine Prolapse} Descent of the uterus toward or through the vaginal opening because the supporting pelvic
floor and ligaments have weakened. Risk factors include childbirth, aging, obesity, and
conditions that increase chronic pressure in the abdomen.

\synonyms
Pelvic Support Problems, Uterine Prolapse
Prolapsed Uterus

\medterm{Uterine Rupture} A full-thickness tear through the uterine wall, potentially life-threatening to both
mother and fetus. Risk factors include prior caesarean section (especially classical
incision), oxytocin augmentation, and operative vaginal delivery. Presents with sudden
abdominal pain, loss of fetal heart rate, and maternal hemorrhage.

\medterm{Uterine Sarcoma} Uterine sarcoma is a rare malignant tumor arising from uterine muscle, connective tissue, or related cells and may cause bleeding, pain, or an enlarging mass.

\medterm{Uterine Tachysystole} More than five uterine contractions in 10 minutes averaged over 30 minutes, whether or
not accompanied by fetal heart-rate abnormalities. It can reduce uteroplacental
perfusion, particularly when caused by oxytocin or prostaglandins.

\medterm{Uterine Tube} A paired muscular tube that transports the oocyte from the ovary toward the uterus and is the usual site of fertilization.

\medterm{Utero-Ovarian Ligament} A band of tissue connecting each ovary to the uterus. It helps hold the ovary in position within the pelvis and lies near the fallopian tube and other pelvic supporting tissues.

\synonyms
Ovarian Ligament; Proper Ovarian Ligament

\medterm{Uterosacral Ligament} One of two strong bands of tissue supporting the uterus and top of the vagina by attaching them toward the sacrum at the back of the pelvis. Weakness can contribute to pelvic-organ prolapse; endometriosis can also affect or cause pain in this area.

\medterm{Uterus} Hollow muscular organ in the pelvis between the bladder and rectum, approximately 7.5 cm
long in adult females, responsible for menstruation, implantation, and fetal development
during pregnancy.
\commonname{Womb}

\medterm{Uterus Didelphys} A congenital condition in which two separate uteruses develop, often with two cervixes and a dividing wall in the vagina. Many people can become pregnant, but miscarriage, premature birth, abnormal fetal position, or other pregnancy complications are more common.

\synonyms
Double Uterus; Didelphic Uterus

\medterm{Uterus Disorder} A condition affecting the uterus, including fibroids, adenomyosis, endometriosis, infection, abnormal bleeding, congenital anomalies, prolapse, and cancer.

\medterm{Uterus Neoplasm} A benign or malignant growth arising in the uterus, including fibroids, endometrial cancer, and sarcomas; diagnosis depends on examination, imaging, and tissue sampling.

\medterm{Uterus, Diseases Of} Conditions affecting the womb, including fibroids, adenomyosis, endometriosis, infection, unusual bleeding, prolapse, and cancer. Symptoms may include pain, heavy periods, infertility, pressure, discharge, or bleeding after menopause.

\medterm{UTI} Infection of the urinary tract, classified as lower (cystitis, urethritis) or upper
(pyelonephritis). See Urinary Tract Infection.

\medterm{Utricle} One of the otolith organs of the inner ear, detecting horizontal linear acceleration and head position relative to gravity.

\medterm{UV Light} Ultraviolet radiation, an invisible part of sunlight or artificial light that can affect the skin and eyes and may cause sunburn, aging,
or skin cancer.

\medterm{Unsaturated Fat} A dietary fat whose fatty-acid chains contain one or more carbon-carbon double bonds. Monounsaturated and polyunsaturated fats are common in vegetable oils, nuts, seeds, and fish; replacing saturated fat with unsaturated fat generally supports healthier LDL cholesterol levels.

\medterm{Uvea} The uvea is the vascular middle layer of the eye, comprising the iris, ciliary body, and choroid; it regulates light entry, supports the lens, and supplies the outer retina.
focus, and blood supply to the eye.

\medterm{Uveal Disease} A disorder of the iris, ciliary body, or choroid, including uveitis, tumors, vascular disease, and inflammation that may cause eye pain, photophobia, floaters, or visual loss.

\medterm{Uveal Melanoma} The most common primary intraocular malignancy in adults, arising from melanocytes
within the uveal tract (choroid, ciliary body, or iris), presenting with visual
symptoms, a pigmented intraocular mass, or incidental finding on examination.

\medterm{Uveal Neoplasm} A tumor arising in the uveal tract of the eye—the iris, ciliary body, or choroid; malignant melanoma may threaten vision and spread to distant organs.

\medterm{Uveitis} Inflammation of the uveal tract—the iris, ciliary body, and choroid—and sometimes nearby eye tissues. It may result from infection, autoimmune or other inflammatory disease, trauma, or an unknown cause. Symptoms can include eye pain, redness, light sensitivity, floaters, blurred vision, or reduced vision; prompt eye assessment is important because untreated inflammation can damage vision.

\synonyms
Chorioretinitis (Uveitis)

\medterm{Uveoparotid Fever} A manifestation of sarcoidosis characterized by the triad of bilateral parotid gland
enlargement, anterior uveitis, and fever, sometimes associated with facial nerve palsy.
It reflects granulomatous inflammation and is treated with systemic corticosteroids and
potentially disease-modifying agents.

\medterm{Uvula} A small, fleshy flap of tissue that hangs from the back of the throat over the root of
the tongue.

\medterm{Uvulitis} Uvulitis is inflammation and swelling of the uvula from infection, trauma, allergy, dehydration, or irritation and may interfere with swallowing or airway protection.
dehydration, or irritation. It may cause throat pain, difficulty swallowing, gagging, or
a foreign-body sensation; marked swelling with breathing difficulty requires urgent
assessment.

\medterm{Uvulopalatopharyngoplasty} A surgical procedure that removes or reshapes selected tissue of the soft palate, uvula, and pharynx to enlarge the upper airway. It may be used in selected patients with obstructive sleep apnea, but effectiveness and risks depend on airway anatomy and other treatments.

\medterm{Ultra-processed food} Ultra-processed food is an industrially formulated product made mostly from refined ingredients and additives, often with little intact whole food. Frequent high intake is associated with poorer health outcomes, although effects vary by the food and overall diet.

\medterm{Ultra-Processed Foods} Ultra-processed foods are packaged products made through multiple industrial processes and often containing additives, added sugars, salt, or refined fats. A balanced diet should emphasize minimally processed foods such as vegetables, fruits, legumes, and whole grains.

\medterm{Ureteric Calculi} Ureteric calculi are kidney stones that have moved into a ureter, the tube carrying urine from a kidney to the bladder. They can cause severe colicky pain, blood in the urine, nausea, or urinary obstruction.

\medterm{Uropathy} Uropathy is disease or damage affecting the urinary tract or kidneys. It may result from obstruction, infection, stones, reflux, medicines, or other conditions and can impair kidney function.

\medterm{Uterine Artery Embolization} Uterine artery embolization is a minimally invasive procedure that blocks blood flow to uterine fibroids, causing them to shrink. It may reduce bleeding and pressure symptoms, but fertility and pregnancy effects require individualized discussion.

\medterm{Vaccination} Giving a vaccine to train the immune system to recognize a specific infection. Vaccination lowers the risk of illness, severe complications, and sometimes transmission; the recommended schedule depends on age, health, and local guidance.

\medterm{Vaccination Schedule} A time-based plan for administering vaccines according to age, previous doses, medical conditions, exposure risks, and public-health guidance. Schedules vary by country and are updated as recommendations change.

\medterm{Vaccinations/Immunizations During Pregnancy} Vaccines given during pregnancy are selected according to safety, disease risk, season, exposure, and local recommendations. Inactivated vaccines such as influenza and Tdap are commonly used when indicated, while some live vaccines are deferred until after pregnancy.
\medterm{Vaccine} A vaccine trains the immune system to recognize a specific infection or toxin and respond more effectively. Vaccination lowers the risk of severe disease and helps protect communities.

\medterm{Vaccine Allergy} An allergic reaction to a vaccine or one of its ingredients. Reactions may be local and mild or, rarely, involve hives, swelling, breathing difficulty, or anaphylaxis; suspected reactions are assessed to identify the responsible ingredient and guide future vaccination.

\medterm{Vaccine-Associated Paralytic Poliomyelitis} An extremely rare form of polio paralysis caused by a weakened poliovirus used in the oral polio vaccine changing back to a harmful form.

\synonyms
VAPP

\medterm{Vaccines} Preparations that train the immune system to recognize a particular germ or toxin before exposure. They reduce the risk of infection and severe illness and may also reduce complications and spread.

\medterm{Vaccines and Immunizations} Vaccines train the immune system to recognize specific pathogens or toxins and reduce the risk of infection, severe disease, complications, and transmission.

\medterm{Vaccinia} A virus used to make smallpox vaccines. Rare accidental infection can cause a sore at the contact site or more widespread illness, especially in people with weakened immunity.

\medterm{Vaccinia Virus} Vaccinia virus is a poxvirus used to make smallpox vaccines and studied in some experimental cancer treatments; accidental infection is uncommon but possible.

\medterm{VACTERL Association} A pattern of birth defects affecting several of the spine, anus, heart, food pipe, windpipe, kidneys, or limbs. There is no single test; the pattern of differences is used to identify it and rule out other conditions.

\medterm{Vacuum-Assisted Delivery} A vaginal birth assisted by a suction cup placed on the baby’s scalp when labor is not progressing or the baby needs to be born promptly. It is used only when the cervix is fully open and the baby’s position and head size are suitable. Possible complications include scalp injury, bleeding, or tears.

\medterm{Vacuum Curettage} A procedure that removes tissue from the uterus using suction, commonly for uterine evacuation or endometrial sampling. Risks include bleeding, infection, retained tissue, and rarely uterine injury.

\medterm{Vacuum Delivery} Assisted vaginal birth using a suction cup placed on the baby’s scalp. It may be used when labor is not progressing or the baby needs to be born promptly and can cause scalp injury, bleeding, or tears.

\medterm{Vacuum Erection Devices} Devices with a cylinder and pump that draw blood into the penis to produce an erection. A flexible ring may hold the erection during sex; bruising, numbness, or discomfort can occur, especially if the ring is left on too long.

\synonyms
Vacuum Erection Pump; Penis Pump

\medterm{Vacuum Extractor} A cup-shaped device applied to the fetal scalp to assist vaginal delivery when specific maternal or fetal indications are present.

\medterm{Vagal Maneuver} A physical action that affects the vagus nerve and can slow the heart’s electrical activity. Examples include bearing down and, in clinical settings, certain pressure or cold-stimulation methods. These maneuvers may help assess or stop some fast heart rhythms.
\synonyms
Vagal Stimulation Technique; Parasympathetic Provocative Maneuver

\medterm{Vagal Nerve Stimulation} A treatment that uses an implanted device to send small electrical signals to the vagus nerve in the neck. It may be used for selected cases of epilepsy or depression that have not improved with other treatments.

\synonyms
VNS

\medterm{Vagal Paraganglioma} A usually slow-growing tumor arising from nerve-related cells near the vagus nerve in the neck. It may cause a neck mass, hoarseness, swallowing difficulty, or nerve symptoms and can be identified with imaging and specialist assessment.

\medterm{Vagina} A muscular, elastic canal connecting the cervix to the vaginal opening. It allows menstrual blood to leave the body, receives a penis or other objects during vaginal sex, and serves as the birth canal.

\medterm{Vagina Carcinoma} Cancer arising from the vagina. It may cause unusual bleeding, discharge, pain, or a visible sore, although early cancer may cause no symptoms; examination and tissue testing are needed for diagnosis.

\medterm{Vaginal Agenesis} A condition in which the vagina is absent or has not fully formed before birth. It may occur with differences in the uterus, kidneys, hearing, or spine; care is individualized and may include counseling, gradual treatment, or surgery when wanted.

\medterm{Vaginal Atresia} Absence or complete blockage of the vaginal canal from birth. It may cause absent menstrual flow, repeated monthly pelvic pain, or difficulty with penetration.

\medterm{Vaginal Atrophy} Thinning, dryness, and irritation of vaginal tissues, usually caused by lower estrogen after menopause or other hormone changes. It can cause burning, itching, pain during sex, bleeding, or urinary symptoms and may improve with moisturizers, lubricants, or prescribed hormone treatment.

\synonyms
Atrophic Vaginitis; Vulvovaginal Atrophy; Genitourinary Syndrome of Menopause; GSM; Vaginal Dryness

\medterm{Vaginal Bleeding} Unexpected, unusually heavy, or prolonged bleeding from the vagina, including bleeding between periods, after sexual intercourse, during pregnancy, or after menopause. Causes include hormonal changes, infection, pregnancy-related conditions, structural lesions, medications, and cancer.

\medterm{Vaginal Bleeding Between Periods} Bleeding between menstrual periods may result from contraception, infection, pregnancy, polyps, fibroids, hormonal changes, or cancer. Persistent or heavy bleeding can have a serious cause.

\medterm{Vaginal Bleeding During Early Pregnancy} Bleeding from the vagina during the first part of pregnancy can result from cervical changes, miscarriage, ectopic pregnancy, infection, or other conditions. A small amount may occur without a serious cause, but bleeding with heavy loss, severe or one-sided pain, shoulder pain, dizziness, or fainting can indicate a serious problem.

\medterm{Vaginal Bleeding During Late Pregnancy} Bleeding from the vagina later in pregnancy may be caused by placenta previa, placental separation, labor, cervical disease, or other complications. Significant bleeding can be a serious pregnancy complication.

\medterm{Vaginal Bleeding In Early Pregnancy} Alternate name; see \textit{Vaginal Bleeding During Early Pregnancy}.

\medterm{Vaginal Bleeding In Late Pregnancy} Alternate name; see \textit{Vaginal Bleeding During Late Pregnancy}.

\medterm{Vaginal Bleeding In Pregnancy} Broad index label for bleeding during pregnancy; see \textit{Vaginal Bleeding During Early Pregnancy} and \textit{Vaginal Bleeding During Late Pregnancy}. Heavy bleeding or bleeding with pain, dizziness, fainting, or reduced fetal movement can indicate a serious pregnancy complication.

\medterm{Vaginal Cancer} A rare cancer of the vagina, most often arising from its surface lining. It may cause bleeding, discharge, pain, or a lump; treatment depends on its type, size, location, and spread.

\synonyms
Vaginal Carcinoma; Malignant Vaginal Neoplasm

\medterm{Vaginal Candidiasis} Alternate name; see \textit{Yeast infection (vaginal)}.

\medterm{Vaginal Cuff} The upper end of the vagina that is closed with stitches after the uterus is removed during a hysterectomy.

\medterm{Vaginal Cuff Dehiscence} An opening of the stitched upper end of the vagina after hysterectomy. It may cause pelvic pain, bleeding, discharge, or tissue bulging through the opening.

\medterm{Vaginal Cysts} Vaginal cysts are fluid-filled or tissue-lined sacs in or beside the vaginal wall, commonly benign but sometimes causing pain, pressure, infection, or obstruction.

\medterm{Vaginal Dilator} A smooth device used to gently stretch or keep open the vaginal canal when it has become narrow or short because of scarring, surgery, radiation, or a difference present from birth.

\medterm{Vaginal Discharge} Vaginal discharge may be normal or may indicate infection, inflammation, hormonal change, or another condition. Normal discharge is usually clear to whitish and odourless and varies with the menstrual cycle, pregnancy, sexual arousal, and some contraceptives.

\medterm{Vaginal Dryness} Reduced vaginal lubrication, commonly associated with menopause, breastfeeding, hormonal
changes, medications, or irritation. It may cause discomfort, itching, painful
intercourse, or minor bleeding.

\medterm{Vaginal Fistula} An abnormal passage between the vagina and another organ, commonly the bladder or rectum. It may cause continuous urine or stool leakage, infection, skin irritation, and emotional distress.

\medterm{Vaginal Foreign Body} An object retained in the vagina, most often a tampon or other inserted material. It may cause discharge, odor, bleeding, pain, or infection.

\medterm{Vaginal Fornix} One of the recesses surrounding the cervix where the vagina meets the cervix. The posterior fornix is adjacent to the peritoneal cul-de-sac and is relevant in pelvic examination and some procedures.

\medterm{Vaginal Fusion And Duplication} Birth-related differences in which the vaginal canal is partly joined, divided, or duplicated. They may occur with abnormalities of the uterus, cervix, urinary tract, or external genitalia and may cause obstruction, pain, menstrual problems, or difficulty with intercourse.

\medterm{Vaginal Health} Vaginal health includes normal vaginal tissues, fluids, microorganisms, and comfort. Persistent odor, itching, pain, bleeding, or unusual discharge can indicate infection or another condition.

\medterm{Vaginal Hematoma} A collection of blood within the vaginal wall or surrounding tissues, usually caused by childbirth, trauma, or surgery and potentially producing severe pain, swelling, or blood loss.

\medterm{Vaginal Hemorrhage} Heavy or clinically significant bleeding from the vagina. Causes include childbirth, pregnancy complications, trauma, infection, structural disease, hormonal disorders, and gynecologic cancer; severe or pregnancy-related bleeding can be life-threatening.

\medterm{Vaginal Hypoplasia} Vaginal hypoplasia is underdevelopment of the vagina, usually present from birth. It may occur with other reproductive or urinary-tract differences and can affect menstruation, sexual function, or fertility.

\medterm{Vaginal Hysterectomy} Surgical removal of the uterus through the vagina. It may be used for selected benign uterine conditions and usually avoids an abdominal incision, although suitability depends on anatomy and disease.

\medterm{Vaginal Inflammation} Alternate name; see \textit{Vaginitis}.

\medterm{Vaginal Intercourse} The insertion of the penis into the vagina for sexual pleasure or reproduction. The most
common form of sexual intercourse. Vaginal intercourse can lead to pregnancy if sperm
are deposited in the vagina.

\medterm{Vaginal Intraepithelial Neoplasia} Abnormal, precancerous cells in the surface lining of the vagina, often related to human papillomavirus. The changes vary from mild to severe and may require monitoring, treatment, or removal.

\medterm{Vaginal Introitus} The entrance to the vagina. It is part of the external genital area and may be examined when assessing pain, infection, injury, scarring, or pelvic-floor problems.

\medterm{Vaginal Itching} Itching or irritation of the vulva or vagina caused by candidiasis, bacterial vaginosis, sexually transmitted infections, dermatitis, hormonal changes, or other conditions. It may occur with discharge, odor, soreness, or burning.

\medterm{Vaginal Laceration} A tear in the vaginal tissue, often caused by childbirth, sexual activity, trauma, or a medical procedure. It may cause pain, swelling, or bleeding.

\medterm{Vaginal Lubrication} Natural moisture made by the vaginal tissues and nearby glands. It usually increases with sexual arousal and may decrease with hormonal changes, some medicines, illness, or irritation.

\medterm{Vaginal Membrane} A thin fold of tissue associated with the vagina. The term may refer to the hymen at the vaginal opening or to an abnormal membrane that blocks the vaginal canal.

\medterm{Vaginal Microbiome} The community of microorganisms living in the vagina, commonly dominated by Lactobacillus species that help maintain an acidic environment. Changes in this community may contribute to bacterial vaginosis, candidiasis, or other symptoms.

\medterm{Vaginal Mucosa} The moist inner lining of the vagina. It contains folds that allow the vaginal canal to stretch and helps protect the tissue.

\medterm{Vaginal Odor} An unusual or unpleasant odor from vaginal discharge that may result from bacterial vaginosis, infection, retained material, hormonal changes, or other conditions.

\medterm{Vaginal Opening} The external entrance to the vagina. It is part of the vulvar vestibule and may be affected by infection, inflammation, trauma, scarring, or congenital variation.

\medterm{Vaginal Pessary} A removable device placed in the vagina to support the vaginal walls or nearby pelvic organs, usually for pelvic-organ prolapse or some forms of urine leakage.

\medterm{Vaginal pH} A measure of how acidic the vagina is. The vaginal environment is normally mildly acidic, and the pH can change with hormones, semen, blood, medicines, or infection.

\medterm{Vaginal Prolapse} Downward bulging of a vaginal wall because the tissues supporting the vagina have weakened. The bulge may involve the bladder, rectum, or the top of the vagina after hysterectomy.

\medterm{Vaginal Rugae} Soft folds in the inner lining of the vagina that allow it to stretch and return toward its usual shape.

\medterm{Vaginal Seeding} Placing vaginal fluid on a baby born by cesarean section in an attempt to expose the baby to the parent’s vaginal microorganisms. Clear health benefits have not been established, and harmful infections may be transmitted.

\medterm{Vaginal Speculum} An instrument inserted into the vagina and gently opened to separate the vaginal walls and visualize the vagina and cervix. It may be used for pelvic examination, cervical cytology such as a Pap test, swab collection, assessment of bleeding or discharge, and selected procedures. Size and design are chosen according to anatomy, comfort, and the clinical purpose.

\medterm{Vaginal Stenosis} Narrowing or shortening of the vaginal canal, which may result from surgery, radiation, scarring, inflammation, trauma, or congenital development. It can cause pain, difficult examination, or difficulty with penetration.

\medterm{Vaginal Swab} A sample of fluid or cells collected from the vagina for laboratory testing, often to look for infection or other changes.

\medterm{Vaginal Trauma} Injury to the vaginal tissues caused by childbirth, sexual activity, foreign objects, procedures, or blunt or penetrating trauma. It may cause bleeding, pain, swelling, or urinary symptoms; severe bleeding or deep injury can be serious.

\medterm{Vaginal Ultrasound Probe} A specialized endocavitary ultrasound transducer used through the vagina to produce images of the uterus, cervix, ovaries, nearby pelvic structures, or an early pregnancy. It sends and receives sound waves and does not use ionizing radiation.

\synonyms
Endovaginal Transducer
Transvaginal Ultrasound Transducer

\medterm{Vaginal Vault Prolapse} Descent of the top of the vagina into or through the vaginal opening, usually after removal of the uterus. It may cause pelvic pressure, a bulge, backache, or difficulty emptying the bladder or bowel; pelvic-floor therapy, a support device, or surgery may help.

\synonyms
Post-Hysterectomy Vault Prolapse; Apical Prolapse

\medterm{Vaginal Vestibule} The space of the vulva between the labia minora that contains the external openings of the urethra and vagina and the ducts of the greater and lesser vestibular glands.

\medterm{Vaginal Wall} The layers of muscle, connective tissue, and moist lining that form the walls of the vaginal canal.

\medterm{Vaginal Yeast Infection} Alternate name; see \textit{Vulvovaginal Candidiasis}.

\medterm{Vagina Neoplasm} An abnormal growth in the vagina that may be benign, precancerous, or cancerous. It can cause bleeding, discharge, pain, or a lump, although early growths may cause no symptoms.

\medterm{Vaginectomy} Surgical removal of part or all of the vagina, usually to treat selected cancers or severe disease; effects may include changes in sexual function, fertility, urination, and bowel function.

\medterm{Vaginismus} Alternate name; see \textit{Genito-Pelvic Pain/Penetration Disorder}.

\medterm{Vaginitis} Inflammation or irritation of the vagina that may cause discharge, itching, burning, odor, pain, or discomfort with urination or sex. Causes include infections, hormones, chemicals, allergies, and tissue dryness.

\medterm{Vaginoplasty} Vaginoplasty is surgery to create, repair, or reshape the vagina. It may be performed for congenital differences, injury, cancer treatment, pelvic-floor problems, or gender-affirming care.

\medterm{Vaginospasm} Alternate name; see \textit{Genito-Pelvic Pain/Penetration Disorder}.

\medterm{Vagotomy} Surgery that cuts selected branches of the vagus nerve to reduce signals that stimulate stomach acid and movement. It was once used for difficult stomach ulcers but is now uncommon because medicines and less invasive treatments are available.

\medterm{Vagus Nerve} The tenth cranial nerve, carrying signals between the brain and the heart, lungs, throat, voice box, and digestive tract. It helps regulate heart rate, breathing, swallowing, voice, and digestion.

\medterm{Vagus Nerve Stimulation} A treatment that sends regular, mild electrical pulses to the vagus nerve through an implanted device in the chest or, in some settings, a handheld device. It is used for selected epilepsy, treatment-resistant depression, and other conditions; benefits may take time and side effects can include voice change, cough, or throat discomfort.

\medterm{Valacyclovir} An antiviral medicine that the body changes into acyclovir to stop herpes viruses from multiplying. It treats cold sores, genital herpes, and shingles; the dose may need adjustment in kidney disease, and excessive levels can affect the nervous system.

\medterm{Valganciclovir} An antiviral medicine used mainly to prevent or treat cytomegalovirus infection after transplantation. The body changes it into ganciclovir, which blocks viral multiplication; it can lower blood counts, injure the kidneys, affect fertility, and harm a fetus.

\medterm{Valgus} A limb or joint deformity in which the part farther from the body bends outward, away from the body’s center line. Examples include a bunion and knock-knees.

\medterm{Valgus Stress Test} A physical examination manoeuvre used to assess the structural integrity of medial collateral ligaments, most commonly at the knee or elbow. A medial-directed force is applied to the joint while stabilizing the limb; excessive joint opening, laxity, or pain indicates a medial collateral ligament sprain or rupture.

\medterm{Valinemia} A rare inherited condition in which the body cannot process the amino acid valine normally. High valine levels may cause feeding problems, low muscle tone, developmental delay, or other nervous-system symptoms.

\medterm{Valley Fever} The common name for coccidioidomycosis, a fungal infection caused by Coccidioides
species and acquired by inhaling airborne spores from soil. It usually causes a mild
respiratory illness but can disseminate to skin, bone, joints, or meninges, especially
in people at higher risk.

\medterm{Valproate} An antiseizure medicine also used for selected mood disorders. It can affect the liver, pancreas, blood cells, weight, and fetal development, so age, pregnancy possibility, other illnesses, and medicines must be considered.

\synonyms
Valproic Acid; Divalproex Sodium

\medterm{Valproate Toxicity} Overdose or adverse exposure causing CNS depression, vomiting, hyperammonemic encephalopathy, hepatotoxicity, pancreatitis, thrombocytopenia, metabolic acidosis, or cerebral edema. Measure serum valproate, glucose, ammonia, liver tests, and acid-base status; L-carnitine may be considered and severe poisoning may require hemodialysis.

\medterm{Valproic Acid} A medicine used for certain seizures, bipolar disorder, and migraine prevention. It can cause liver injury, pancreatitis, birth defects, and developmental harm during pregnancy.

\medterm{Valproic Acid Toxicity} Acute overdose or adverse metabolic toxicity of the broad-spectrum antiepileptic agent valproate, characterized clinically by central nervous system depression, high anion gap metabolic acidosis, hyperammonemia (due to urea cycle inhibition) with encephalopathy, and drug-induced microvesicular hepatic necrosis; treated with L-carnitine.

\medterm{Valrubicin} A chemotherapy medicine placed directly into the bladder for selected bladder cancers that have not responded adequately to other treatment. It can cause bladder irritation, blood in the urine, infection, and, rarely, effects throughout the body.

\medterm{Valsalva} Straining while trying to breathe out against a closed throat, like bearing down during a bowel movement. It changes pressure in the chest and can affect heart rate and blood pressure, so it is not safe for everyone.

\medterm{Valsalva Maneuver} Alternate name; see \textit{Valsalva manoeuvre}.

\medterm{Valsalva Retinopathy} Preretinal (subhyaloid or sub-internal limiting membrane) hemorrhage occurring within the macula following a sudden, dramatic rise in intrathoracic or intra-abdominal pressure (such as from forceful coughing, vomiting, or heavy lifting) that transmits back through the valveless venous system into retinal capillaries.

\medterm{Valsalva's Manoeuvre} Forceful straining while trying to breathe out against a closed airway, like bearing down during a bowel movement. It changes pressure in the chest and can affect heart rate and blood pressure; it is used in selected tests but may be unsafe for some people.

\synonyms
Valsalva Maneuver

\medterm{Valsartan} Valsartan is an angiotensin-II receptor blocker used for hypertension, heart failure, and selected cardiovascular indications; hyperkalaemia, renal-function changes, hypotension, and fetal toxicity are important risks.

\medterm{Valves} Structures that permit fluid or blood to move predominantly in one direction and prevent backward flow; cardiac valves and venous valves are important examples.

\medterm{Valvotomy} Surgical or catheter-based division of narrowed valve tissue to improve blood flow, most often performed for selected stenotic heart valves.

\medterm{Valvulae Conniventes} Large, permanent transverse circular folds of the mucous membrane and submucosa that project into the lumen of the small intestine (most prominent and numerous in the jejunum), which do not disappear with intestinal distension and act to triple the mucosal surface area for absorption.

\synonyms
Plicae Circulares; Kerckring Folds

\medterm{Valvular Disease} Alternate term; see \textit{Heart valve disease}.

\medterm{Valvular Endocarditis} Alternate name; see \textit{Endocarditis}.

\medterm{Valvular Heart Disease} A broad spectrum of structural and functional disorders affecting one or more of the four cardiac valves (aortic, mitral, tricuspid, pulmonic), classified hemodynamically as valvular stenosis (narrowing and outflow obstruction) or valvular regurgitation (insufficiency and backward flow).

\synonyms
VHD

\medterm{Valvular Insufficiency} Incomplete closure of a heart valve, allowing blood to leak backward when the heart contracts or relaxes. It can affect the mitral, aortic, tricuspid, or pulmonary valve and may eventually strain the heart.

\medterm{Valvuloplasty} A procedure used to widen or repair a narrowed or abnormal heart valve, by balloon catheter or open surgery depending on the valve and condition.

\medterm{Valvulotomy} A procedure that cuts open a narrowed heart valve to improve blood flow. It may be performed with a catheter or surgery; the valve can leak afterward and narrowing may recur.

\medterm{Vancomycin} An antibiotic used for selected serious bacterial infections and, by mouth, for certain infections of the large intestine. It can cause infusion reactions, kidney injury, hearing problems, or low blood counts and requires medical monitoring.

\medterm{Vancomycin Flushing Reaction} A reaction that can occur when vancomycin is given too quickly into a vein. It may cause flushing, warmth, itching, and redness of the face, neck, and upper body.

\synonyms
Red Man Syndrome; Red Neck Syndrome

\medterm{Vancomycin-Resistant Enterococcus} Enterococcus bacteria that are resistant to vancomycin, usually because of altered cell-wall precursors. VRE can
cause healthcare-associated urinary, bloodstream, wound, or other infections and may
require susceptibility-guided treatment.

\medterm{Van der Woude Syndrome} An autosomal dominant congenital disorder caused by mutations in the IRF6 gene, characterized by the combination of paramedian lower-lip pits (fistulae), cleft lip, and cleft palate, representing the most frequent form of syndromic orofacial clefting.
\synonyms
VWS; Lip-Pit Syndrome

\medterm{Vandetanib} A targeted cancer medicine used for selected medullary thyroid cancers. It blocks signals that help cancer cells grow and blood vessels form; diarrhea, high blood pressure, skin effects, and dangerous heart-rhythm changes require monitoring.

\medterm{Vaping} Vaping uses an electronic device to heat a liquid and produce an aerosol for inhalation. Aerosols may contain nicotine and other substances that can harm the lungs and cause dependence.

\medterm{Vaporiser} A device that converts a liquid medicine or volatile substance into a vapor for inhalation or controlled delivery, such as in anesthesia.

\medterm{Variable Decelerations} Brief drops in a baby’s heart rate that vary in timing and shape, often caused by pressure on the umbilical cord during labor. Occasional drops may be harmless, but frequent or prolonged drops require assessment of the baby’s oxygen supply.

\medterm{Variant Angina Pectoris} Chest pain caused by a temporary tightening of an artery supplying the heart, often occurring at rest or at night. Episodes may be severe and can disturb the heart rhythm; sudden or prolonged chest pain needs emergency care.

\synonyms
Prinzmetal Angina; Vasospastic Angina

\medterm{Variant CJD} Alternate name; see \textit{Creutzfeldt-Jakob disease}.

\medterm{Variant Classification} The process of estimating whether a genetic change is harmless, uncertain, or likely to cause disease. Specialists use laboratory evidence, family history, population data, and the person’s clinical findings.

\medterm{Varibaculum Cambriense} A bacterial species in the genus Varibaculum reported from human clinical specimens. Its role may be opportunistic and requires identification with the clinical context and susceptibility data.

\medterm{Variceal Band Ligation} A procedure using a flexible camera to place small bands around enlarged veins in the food pipe. The bands stop blood flow through the veins and are used to control or prevent bleeding.

\synonyms
Endoscopic Variceal Ligation; EVL

\medterm{Variceal Hemorrhage} Severe bleeding from enlarged, fragile veins in the food pipe or stomach, usually caused by high pressure in the liver’s blood vessels. It can cause vomiting blood, black stools, shock, and death without urgent treatment.

\medterm{Varicella} A primary infection with varicella-zoster virus that causes fever and a generalized itchy rash. See Chickenpox.

\synonyms
Chickenpox; Varicella-Zoster Virus Infection

\medterm{Varicella-Zoster Virus} The virus that causes chickenpox and shingles. After chickenpox, it can remain inactive in nerve tissue and reactivate years later as a painful, blistering rash called shingles.

\medterm{Varicellovirus} Varicellovirus is a genus of herpesviruses that includes varicella-zoster virus, the cause of chickenpox and shingles.

\medterm{Varices} Abnormally dilated and tortuous veins, commonly occurring in the legs or in the oesophagus and stomach as a result of portal hypertension.

\medterm{Varicocele} Enlargement of veins in the scrotum, often on the left. It may cause aching, a soft “bag of worms” feeling, testicular growth differences, or reduced fertility, but many cause no symptoms.

\medterm{Varicocele Embolization} A procedure that uses a thin tube to place coils or another material in enlarged scrotal veins. This blocks abnormal backward blood flow and reduces pressure in the veins.

\medterm{Varicose Ulcer} An open sore near the ankle caused by long-term poor drainage through the leg veins. It may be surrounded by swelling or skin discoloration and often needs compression and treatment of the underlying vein problem.

\medterm{Varicose Veins} Enlarged, twisted superficial veins, usually in the legs, caused by weak or damaged vein walls or valves that allow blood to pool and flow backward. They may cause visible bulging veins, aching, heaviness, itching, cramps, swelling, skin discoloration, bleeding, superficial inflammation or clotting, and difficult-to-heal venous ulcers.

\medterm{Varicose Veins and Venous Insufficiency} Enlarged superficial veins and inadequate return of blood from the legs because venous valves or the vein wall do not work effectively. Aching, heaviness, swelling, skin discoloration, dermatitis, and ulcers may occur.

\medterm{Variegate Porphyria} An inherited disorder of heme production that can cause attacks of abdominal pain, vomiting, nerve symptoms, confusion, or dark urine and can also cause fragile, blistering skin after sunlight exposure.

\medterm{Variola} A highly contagious and deadly disease caused by variola virus, eradicated globally by
1980 through vaccination. See Smallpox.

\medterm{Varix} An abnormally dilated, tortuous vein; multiple varices commonly occur in portal hypertension or chronic venous disease.

\medterm{Varus} A limb or joint deformity in which the part farther from the body bends inward, toward the body’s center line. Examples include bowlegs and some inward-turning foot deformities.

\medterm{Varus Stress Test} A clinical examination manoeuvre used to evaluate the lateral collateral ligaments of a joint, particularly the knee or elbow. A lateral-directed force is applied across the joint line while stabilizing the limb; increased laxity, gapping, or localized pain signifies injury to the lateral collateral ligament complex.

\medterm{Vas} A duct or vessel; in reproductive anatomy, the vas deferens transports sperm from the epididymis toward the ejaculatory duct.

\medterm{Vasa Deferentia} The paired muscular ducts that transport mature sperm from the epididymides toward the ejaculatory ducts during ejaculation.

\synonyms
Ductus Deferentes; Vas Deferens

\medterm{Vasa Efferentia} Small tubes that carry sperm from the testicle’s collecting channels to the epididymis, where sperm mature. They also remove some fluid and help move sperm onward.

\medterm{Vasa Previa} A pregnancy complication in which fetal blood vessels run across or near the cervix without the usual protection of the placenta or umbilical cord. If the membranes break, the vessels can tear and cause rapid fetal blood loss, so specialist planning is essential.

\medterm{Vasa Recta} Long, hairpin-shaped blood vessels in the inner part of the kidney. They run beside the loops of Henle and help the kidneys concentrate urine.

\medterm{Vasa Vasorum} Small blood vessels in the walls of large arteries and veins. They supply oxygen and nutrients to the outer parts of these vessel walls.

\medterm{Vascular} Relating to blood vessels or, sometimes, lymphatic vessels. Vascular conditions include narrowed or blocked arteries, aneurysms, blood clots, abnormal vessel connections, and problems with veins; the term may also describe tests, surgery, or medicines involving these vessels.

\medterm{Vascular Access} A route into a blood vessel for medicines, fluids, blood tests, monitoring, or dialysis. It may be a short-term catheter, a surgically joined artery and vein, or a synthetic graft, depending on the purpose and expected duration.

\medterm{Vascular Access For Hemodialysis} A route that allows blood to leave the body and return during dialysis. It may be a joined artery and vein, a synthetic graft, or a large vein catheter; each has different risks of infection, clotting, and failure.

\medterm{Vascular Balloon Catheter} An inflatable balloon mounted at the distal tip of a flexible vascular catheter used to dilate stenotic or occluded blood vessels during angioplasty, deliver and expand endovascular stents, or temporarily occlude vascular flow.

\synonyms
Angioplasty Balloon; Vascular Balloon Dilator; Balloon Catheter

\medterm{Vascular Bed} The network of blood vessels supplying a particular organ or tissue. Its vessels widen or narrow to adjust blood flow according to the tissue’s needs.

\medterm{Vascular Bypass} Surgery that uses a vein or artificial tube to carry blood around a blocked part of an artery and restore blood flow.

\synonyms
Surgical Bypass; Bypass Grafting; Arterial Bypass Graft

\medterm{Vascular Claudication} Alternate name; see \textit{Claudication}.

\medterm{Vascular Clip} A small device placed during a procedure to close or control a blood vessel. It may stop bleeding, prevent blood flow to a selected area, or secure a vessel permanently.
\synonyms
Hemostatic Clip; Surgical Vessel Clip; Ligation Clip

\medterm{Vascular Compliance} The ability of a blood vessel to stretch when pressure or blood volume increases. More compliant vessels can hold more blood with a smaller rise in pressure.

\synonyms
Angiopathic Complications; Vascular Adverse Events

\medterm{Vascular Complications} Problems involving blood vessels after disease, injury, or a procedure, such as bleeding, hematoma, thrombosis, embolism, vessel injury, or reduced blood flow to a limb. Symptoms and urgency depend on the vessel and tissue affected.

\medterm{Vascular Dementia} Problems with thinking, memory, mood, or walking caused by strokes or long-term damage to the brain’s small blood vessels. Symptoms may appear suddenly or gradually, and treatment focuses on preventing further vessel damage and supporting daily function.

\synonyms
VaD; Multi-Infarct Dementia; Vascular Cognitive Impairment

\medterm{Vascular Disease} Vascular disease affects arteries, veins, or lymphatic vessels and includes atherosclerosis, aneurysm, peripheral artery disease, venous thrombosis, varicose disease, and inflammatory disorders. Symptoms and treatment depend on the vessel and may involve risk-factor control, medicines, procedures, or surgery.

\medterm{Vascular Ehlers-Danlos Syndrome} A rare inherited connective-tissue disorder that makes arteries and hollow organs unusually fragile. It may cause thin translucent skin, easy bruising, characteristic facial features, or serious spontaneous blood-vessel or bowel rupture; specialist care and genetic counseling are important.

\synonyms
vEDS; Ehlers-Danlos Syndrome Type IV

\medterm{Vascular Embolic Particles} Small materials placed in a blood vessel during a procedure to block blood flow. They may be used to control bleeding or reduce blood flow to a tumor, fibroid, or other abnormal area.

\synonyms
Vascular Embolization Particle; PVA Particles

\medterm{Vascular Embolization} Intentional blockage of a blood vessel with material delivered through a catheter. It may stop bleeding, reduce blood flow to a tumor, or close an aneurysm or abnormal vessel connection.

\synonyms
Catheter Embolization; Therapeutic Embolization

\medterm{Vascular Embolization Coil} A catheter-delivered wire device used to block selected blood vessels, often to treat aneurysms or bleeding.

\synonyms
Endovascular Coil

\medterm{Vascular Embolization Plug} A vascular occlusion device used to block a selected vessel under image guidance. A device placed in a blood vessel to block selected blood flow during image-guided treatment.
\synonyms
Vascular Occlusion Plug; Embolization Plug; Amplatzer Plug

\medterm{Vascular Endothelial Growth Factor} A protein that helps new blood vessels form and makes blood vessels more leaky. Some medicines block it to treat certain eye diseases and cancers.
\synonyms
VEGF; Vascular Permeability Factor

\medterm{Vascular Endothelial Growth Factor Inhibitor} A medicine that blocks vascular endothelial growth factor, a protein that helps blood vessels grow. These medicines are used for some cancers and eye diseases involving abnormal new blood vessels.

\synonyms
Anti-VEGF Agent

\medterm{Vascular Endothelial Growth Factor Inhibitors} Medicines that block vascular endothelial growth factor and reduce abnormal blood-vessel growth or leakage. They are used for some eye diseases and cancers.

\synonyms
Anti-VEGF Agents; VEGF Neutralizers

\medterm{Vascular Forceps} Surgical forceps designed to handle blood vessels, clots, tissue, or medical materials during vascular procedures. Appropriate use helps avoid crushing or tearing delicate vessels.

\synonyms
Vascular Surgical Forceps; Atraumatic Clamp

\medterm{Vascular Graft} A tube used to replace or bypass a damaged or blocked blood vessel. It may be made from a person's vein or from a synthetic material; possible problems include infection, bleeding, narrowing, or a clot blocking the graft.

\medterm{Vascular Graft Patency} The continued openness and useful blood flow through a blood-vessel graft or bypass. A clot, narrowing, infection, or scar growth can block it, so examination and sometimes ultrasound are used for monitoring.

\medterm{Vascular Hemostat} A surgical clamp used to temporarily occlude a blood vessel and control bleeding.
\synonyms
Vascular Clamp; Atraumatic Hemostasis Forceps

\medterm{Vascular Imaging} Tests that show blood vessels and sometimes measure blood flow. Ultrasound, CT, MRI, or catheter-based imaging can find narrowing, blockage, aneurysm, dissection, clot, bleeding, or abnormal vessels; the safest test depends on the body area, urgency, kidney function, and contrast risks.

\medterm{Vascularity} The amount and arrangement of blood vessels supplying a tissue or growth. A highly vascular area may heal or bleed differently and may look different on imaging; blood supply can also affect some treatments.

\medterm{Vascular Laser Therapy} A treatment that uses laser energy to close or change selected blood vessels. The laser may be applied through the skin or through a thin tube.
\synonyms
Pulsed Dye Laser; PDL; Vascular Laser Therapy

\medterm{Vascular Ligation} Surgical tying, clipping, or sealing of a blood vessel to stop bleeding or redirect blood flow. It may be used during surgery or to treat selected abnormal vessels.
\synonyms
Vascular Occlusion; Vessel Clipping

\medterm{Vascular Malformation} A blood-vessel or lymph-vessel structure that developed abnormally before birth. It may cause a skin mark, swelling, pain, bleeding, or pressure on nearby organs; effects depend on its type and location.

\medterm{Vascular Needle Holder} A surgical instrument used to hold and guide fine suture needles during repair of blood vessels. It permits precise stitches while minimizing crushing or tearing of delicate vessel walls.
\synonyms
Vascular Needle Holder; Castroviejo Needle Holder

\medterm{Vascular Pain} Pain caused by impaired blood flow, often appearing in a limb with exertion as claudication or, in more severe disease, at rest.

\medterm{Vascular Patch Graft} A piece of biological or synthetic material sewn onto a blood vessel to widen or repair it.

\synonyms
Vascular Patch Graft; Angioplasty Patch

\medterm{Vascular Retraction} To pull a vessel or tissue away from an operative field using a surgical retractor.
\synonyms
Vessel Retraction; Surgical Vessel Exposure

\medterm{Vascular Ring} A birth defect in which blood vessels around the heart form a ring that can press on the windpipe or food pipe. It may cause noisy breathing, cough, repeated chest infections, swallowing difficulty, or feeding problems. Some people need no treatment; symptomatic rings may require surgery.

\medterm{Vascular Scissors} A surgical instrument with two blades used to cut tissue, sutures, dressings, or other materials
during vascular procedures.

\synonyms
Vascular Scissors; Potts Scissors

\medterm{Vascular Shunt} A surgically created or implanted pathway that redirects blood or another fluid from one vessel or body space to another.

\synonyms
Vascular Shunt; Bypass Shunt

\medterm{Vascular Sponge} A sterile absorbent material used during blood-vessel procedures to absorb blood, apply pressure, or protect delicate tissue.
\synonyms
Vascular Laparotomy Sponge; Surgical Cottonoid

\medterm{Vascular Staple} A small metal or polymer device used to close tissue or join selected blood-vessel structures. It provides rapid closure but can cause narrowing, bleeding, thrombosis, or tissue injury if unsuitable or misplaced.
\synonyms
Vascular Surgical Staple; Vessel Closure Device

\medterm{Vascular Stent} A small mesh or tube placed inside a narrowed or blocked vessel to help keep it open.

\synonyms
Endovascular Stent; Vascular Prosthesis

\medterm{Vascular Stent Balloon} A balloon on a catheter used to expand a vascular stent against the vessel wall and restore an open channel. Inflation can cause vessel injury, bleeding, clotting, or rarely rupture.
\synonyms
Stent-Expansion Balloon; Angioplasty Balloon

\medterm{Vascular Stent-Graft} A fabric-covered metal frame placed inside a blood vessel to reinforce a weak wall or seal an aneurysm, tear, or leak. Follow-up imaging checks for movement, blockage, or leakage around the graft.
\synonyms
Endovascular Graft; Aortic Stent-Graft

\medterm{Vascular Suction} A suction device used during a vascular procedure to remove blood or other fluid from the operative
field.

\synonyms
Vascular Suction Cannula; Frazier Suction

\medterm{Vascular Surgery} Operations that repair, widen, replace, or reroute blood vessels to restore circulation or prevent bleeding. Procedures may treat aneurysms, blocked arteries, blood-vessel injuries, or problems with veins and dialysis access.

\medterm{Vascular System} The network of arteries, veins, and capillaries that carries blood throughout the body. It delivers oxygen, nutrients, and hormones to tissues and carries carbon dioxide and other waste products away. Disease may narrow, weaken, block, or inflame these vessels.

\medterm{Vascular Thrombectomy} Removal of a blood clot from a blood vessel using surgery or a catheter-based procedure. The aim is to restore blood flow and protect tissue from damage when a clot blocks an important artery or vein.

\synonyms
Vascular Thrombectomy; Clot Extraction

\medterm{Vascular Tissue Adhesive} A liquid material used to block selected blood vessels during image-guided embolization. Doctors inject it through a catheter to control bleeding, close an abnormal vessel, or reduce blood flow to a tumor.

\synonyms
Cyanoacrylate Embolization Glue; Endovascular Embolic Adhesive

\medterm{Vascular Tone} The usual degree of tightening in the muscle within blood-vessel walls. It helps control vessel width, blood flow, and blood pressure.

\medterm{Vascular Trauma} Injury to an artery or vein caused by a blunt blow, penetrating wound, crush injury, or medical procedure. It may cause severe bleeding or loss of blood flow and needs urgent assessment of circulation and nearby nerves.

\medterm{Vascular Ulcer} A persistent open wound caused or maintained by poor blood flow, high pressure in the veins, loss of protective feeling, or another blood-vessel problem. It usually affects the legs or feet and may heal slowly or become infected.

\synonyms
Vascular Ulceration; Ischemic Ulcer; Arterial Ulcer

\medterm{Vasculitic Neuropathy} Nerve damage caused by inflammation of the small blood vessels that supply nerves. It often causes painful, uneven numbness or weakness and may affect several limbs as the vessel inflammation continues.

\medterm{Vasculitis} Inflammation of blood vessels that can narrow them, reduce blood flow, or cause bleeding and tissue damage. It may affect large, medium, or small vessels and can involve the skin, nerves, kidneys, lungs, or other organs.

\synonyms
Angiitis; Arteritis; Blood Vessel Inflammation

\medterm{Vas Deferens} A muscular tube that carries sperm from the epididymis toward the urethra during ejaculation.

\medterm{Vasectomy} A permanent contraceptive procedure in which the vas deferens are divided or occluded to prevent sperm from entering the ejaculate. It does not immediately provide sterility and does not protect against sexually transmitted infections.

\medterm{Vasectomy Reversal} Surgery to reconnect the tubes that carry sperm after a vasectomy. It may restore sperm to the semen, but pregnancy is not guaranteed.
\synonyms
Vasovasostomy; Vasoepididymostomy

\medterm{Vasoactive Intestinal Peptide Test} A blood test measuring a hormone that affects the intestine and blood vessels. It is used mainly when a rare hormone-producing tumor is suspected in someone with severe watery diarrhea, dehydration, or abnormal body salts.

\medterm{Vasoconstriction} Narrowing of blood vessels caused by contraction of vascular smooth muscle, which can raise blood pressure or reduce blood flow to affected tissues.

\medterm{Vasoconstrictor} A substance that narrows blood vessels by making their muscle walls contract. This can raise blood pressure and reduce local blood flow; vasoconstrictor medicines are used in selected emergencies or with local anesthetics and require careful monitoring.

\medterm{Vasoconstrictors} Substances or medicines that narrow blood vessels by making their muscular walls contract. They can raise blood pressure or reduce local blood flow and are used only for selected conditions because excessive narrowing can be harmful.

\medterm{Vasodepressor Syncope} Fainting caused mainly by sudden widening of blood vessels and a fall in blood pressure. It often occurs as part of a vasovagal reaction and may be preceded by nausea, sweating, warmth, or lightheadedness.

\medterm{Vasodilatation} Widening of blood vessels caused by relaxation of vascular smooth muscle, increasing blood flow and potentially lowering blood pressure.

\medterm{Vasodilation} Widening of blood vessels when their muscular walls relax. It can increase blood flow and lower blood pressure and occurs with exercise, heat, medicines, and inflammation.

\medterm{Vasodilator} A substance that relaxes blood-vessel walls and makes the vessels wider. It can lower blood pressure or improve blood flow and is used for selected heart, blood-pressure, or circulation problems; excessive widening may cause headache, flushing, dizziness, or fainting.

\medterm{Vasodilators} Medicines or substances that relax vascular smooth muscle and widen blood vessels. They are used for conditions such as hypertension, angina, heart failure, or pulmonary hypertension depending on the agent.

\medterm{Vasodilatory Shock} A form of distributive shock characterized by severe systemic arterial hypotension and tissue hypoperfusion resulting from profound, refractory peripheral vasodilation and loss of vascular tone despite preserved or elevated cardiac output, most frequently seen in severe sepsis, anaphylaxis, and neurogenic injury.

\medterm{Vasomotor} Relating to control of blood-vessel width by nerves, hormones, and local chemicals. Vasomotor changes narrow or widen vessels and help regulate blood pressure, body temperature, tissue blood flow, and skin color.

\medterm{Vasomotor Center} A group of nerve cells in the lower brain that helps control the widening and narrowing of blood vessels and helps maintain blood pressure.

\medterm{Vasomotor Centre} Groups of nerve cells in the brainstem that help control blood-vessel narrowing and widening. By changing vessel tone, they influence blood pressure, heart and tissue blood flow, and the body’s response to standing or stress.

\medterm{Vasomotor Nerves} Autonomic nerves controlling blood-vessel widening and narrowing by changing smooth-muscle activity, helping regulate blood pressure and temperature.

\medterm{Vasomotor Rhinitis} Common subtype term; see \textit{Nonallergic rhinitis}.

\medterm{Vaso-Occlusive Crisis} The classic acute, excruciatingly painful hallmark complication of sickle cell disease, caused by the intravascular sludging and adherence of rigid, polymer-filled sickled erythrocytes in the microvasculature, producing focal tissue ischemia, infarction, and systemic inflammation in bones, joints, or viscera.

\synonyms
Sickle Cell Pain Crisis; VOC

\medterm{Vasopressin} A hormone that increases water reabsorption in the kidney and can constrict
blood vessels. It is released from the posterior pituitary.

\synonyms
Arginine Vasopressin; AVP

\medterm{Vasopressin Receptor} A cell sensor that responds to the hormone vasopressin. Different types help the kidneys retain water and help blood vessels narrow, influencing body-fluid balance and blood pressure.

\medterm{Vasopressin Receptor Antagonists} A pharmacological class of aquaretic medications (vaptans, such as tolvaptan and conivaptan) that competitively block vasopressin V2 receptors in the renal collecting duct, promoting solute-free water excretion without electrolyte loss in euvolemic and hypervolemic hyponatremia (SIADH, heart failure, cirrhosis).

\synonyms
Vaptans

\medterm{Vasopressor} A medicine or body signal that raises blood pressure, mainly by narrowing blood vessels and sometimes by strengthening the heartbeat. These medicines are used in selected emergencies such as shock and require close monitoring.

\medterm{Vasopressors In Anesthesia} Medicines that increase vascular tone or cardiac output to support blood pressure during
anesthesia or critical illness. The choice and dose depend on the cause of hypotension,
heart rate, cardiac function, and local protocol; monitoring is required.

\medterm{Vasospasm} Sudden narrowing of a blood vessel caused by contraction of its muscular wall. It can temporarily reduce blood flow and cause chest pain, headache, stroke-like symptoms, or tissue injury, depending on the vessel affected.

\medterm{Vasovagal Attack} A transient reflex episode caused by reduced sympathetic tone and increased vagal activity, producing hypotension and often bradycardia with fainting or near-fainting. Common triggers include prolonged standing, pain, fear, heat, dehydration, or emotional stress.

\medterm{Vasovagal Reaction} A reflex that slows the heart and widens blood vessels, briefly lowering blood pressure and sometimes causing fainting; triggers include pain, fear, prolonged standing, heat, or blood exposure.

\medterm{Vasovagal Syncope} A common faint caused by a reflex that briefly slows the heart or widens blood vessels, lowering blood pressure and brain blood flow. Triggers include standing still, heat, pain, fear, or seeing blood; warning symptoms often occur first.

\synonyms
Neurocardiogenic Syncope

\medterm{Vater-Pacini Corpuscles} Large, rapidly adapting, encapsulated mechanoreceptors located deep in the dermis, subcutaneous tissue, periosteum, and joint capsules. Morphologically composed of a central unmyelinated nerve axon surrounded by multiple concentric lamellae of flattened Schwann-like cells separated by fluid; specialized physiologically to detect high-frequency mechanical vibration (200 to 300 Hz) and transient acceleration.

\synonyms
Pacinian Corpuscles; Lamellar Corpuscles

\medterm{VATS} Alternate name; see \textit{Video-Assisted Thoracoscopic Surgery}.

\medterm{Vaughan Williams Classification} A system that groups medicines used for abnormal heart rhythms according to their main effect on the heart’s electrical activity. It helps describe treatment but does not replace patient-specific prescribing.

\medterm{VCJD} Alternate name; see \textit{Creutzfeldt-Jakob disease}.

\medterm{VDRL Test} A blood or spinal-fluid test that looks for an immune reaction associated with syphilis. A reactive result needs a more specific confirmatory test and may help monitor response to treatment.

\medterm{Vector} A living carrier, usually an insect or tick, that spreads an infection between people or animals. Some vectors allow the germ to develop inside them; others transfer it mechanically without the germ changing.

\medterm{Vector-Borne Disease} An infection spread by a living carrier such as a mosquito, tick, flea, or sandfly. Examples include malaria, dengue, Lyme disease, and some forms of viral encephalitis; prevention focuses on avoiding bites and controlling the carrier.

\medterm{Vector Control} Measures that reduce insects or other animals that spread infection. Examples include removing standing water, using insecticides, treating breeding areas, improving buildings, wearing protective clothing, and using repellents or bed nets.

\medterm{Vegetarian} Vegetarian describes a diet that avoids meat and may include dairy products, eggs, or both. Nutritional needs can usually be met with varied foods, but iron, vitamin B12, protein, and other nutrients may need attention.

\medterm{Vegetation} A mass that can form on a heart valve during infective endocarditis. It may contain bacteria, blood-clotting material, and immune cells and can damage the valve or break off and travel to other organs.

\medterm{Vein} A blood vessel that carries blood toward the heart. Most veins carry oxygen-poor blood, but veins from the lungs and placenta carry oxygen-rich blood. Veins often contain valves that help prevent backward flow.

\medterm{Vein-Graft Aneurysms} An abnormal widening of a vein used to bypass a blocked heart artery. It may contain a clot, block blood flow, press on nearby structures, or rupture and therefore needs specialist assessment.

\medterm{Vein of Galen Malformation} A rare congenital intracranial vascular anomaly resulting from an embryonic arteriovenous fistula that drains directly into the primitive median prosencephalic vein (the precursor of the great cerebral vein of Galen), presenting in neonates with high-output congestive heart failure and hydrocephalus.

\synonyms
VGM; Vein of Galen Aneurysmal Malformation

\medterm{VeIP Regimen} A cancer-treatment combination of vinblastine, ifosfamide, and cisplatin, used mainly for selected germ-cell tumors. It can lower blood counts and damage the kidneys, nerves, or bladder, so specialist monitoring is required.

\medterm{Velamentous Cord Insertion} An abnormal attachment in which the umbilical blood vessels travel through the membranes before reaching the placenta without the usual protective tissue. The vessels can be compressed or torn, increasing risks of poor fetal growth, bleeding, or vasa previa and requiring specialist pregnancy monitoring.

\synonyms
Velamentous Insertion of the Umbilical Cord

\medterm{Velocardiofacial Syndrome} Alternate name; see \textit{DiGeorge syndrome (22q11.2 deletion syndrome)}.

\medterm{Velopharyngeal Insufficiency} A structural or functional anatomical failure of the soft palate (velum) and posterior pharyngeal wall to achieve complete closure during speech and deglutition, resulting in hypernasal speech resonance, nasal air emission, and nasal regurgitation of fluids.
\synonyms
VPI; Velopharyngeal Incompetence

\medterm{Vemurafenib} A targeted cancer medicine for tumors with a specific BRAF V600 gene change, including some melanomas. It can cause joint pain, sun sensitivity, abnormal heart rhythm, new skin cancers, or treatment resistance.

\medterm{Vena Cava} Either of the two largest veins returning blood to the right side of the heart. The superior vena cava drains the upper body, and the inferior vena cava drains the abdomen, pelvis, and legs.


\synonyms
Venae Cavae; Superior and Inferior Vena Cava

\medterm{Vena Cava Filter} A small device placed in the large vein leading to the heart to trap blood clots traveling from the legs toward the lungs. It may be used when blood-thinning medicines cannot be used or have not worked.

\synonyms
IVC Filter; Greenfield Filter

\medterm{Venepuncture} Inserting a needle into a vein to collect blood or place intravenous access. It may cause brief pain, bruising, bleeding, infection, or fainting, and clean technique reduces complications.

\medterm{Venereal Disease Research Laboratory Test} A blood test that looks for antibodies associated with syphilis. It is used for screening and monitoring treatment, but a reactive result needs confirmation with a treponemal test.

\medterm{Venesection} The controlled removal of blood from a vein, used diagnostically or therapeutically, for example in selected cases of hemochromatosis or polycythemia.

\medterm{Venipuncture} Puncture of a vein with a needle to obtain blood, administer fluids or medicines, or establish intravenous access. Bruising, infection, bleeding, and nerve or artery injury are uncommon complications.

\medterm{Venogram} Contrast imaging of veins used to detect thrombosis, obstruction, abnormal anatomy, or selected venous flow disorders.

\medterm{Venography} Imaging of veins, usually after contrast material is administered, to assess venous anatomy, obstruction, thrombosis, or malformation.

\medterm{Venom} A poisonous substance delivered by an animal through a bite, sting, or similar structure. Its effects vary by species and may include pain, swelling, nerve problems, bleeding, or dangerous changes in breathing or blood pressure.

\medterm{Venomous Bite} A bite that injects venom from an animal such as a snake, spider, or arthropod. Effects range from local
pain and swelling to bleeding, paralysis, or systemic toxicity.

\medterm{Veno-Occlusive Disease} Blockage and injury of small veins in the liver, especially after some chemotherapy or stem-cell transplantation. It can cause painful liver enlargement, rapid weight gain, fluid retention, and jaundice and can be life-threatening.

\medterm{Veno-occlusive Hepatic Disease (Sinusoidal Obstruction Syndrome)} A potentially severe liver injury caused by narrowing or blockage of hepatic sinusoids, classically after hematopoietic-cell transplantation or certain chemotherapy. Painful liver enlargement, weight gain, fluid retention, jaundice, and rising bilirubin may occur.

\medterm{Venous} Relating to veins, which usually carry blood back toward the heart. Vein problems include blood clots, varicose veins, poor return of blood, and ulcers; pulmonary veins are an exception because they carry oxygen-rich blood from the lungs to the heart.

\medterm{Venous Air Embolism} Obstruction of the venous circulation by intravascular gas, usually after procedures, trauma, or pressure changes. It can impair pulmonary or cardiac function and, if gas reaches the arterial circulation, cause neurologic or systemic ischemia.

\medterm{Venous Aneurysm} A localized abnormal dilatation of a vein. It may be asymptomatic or present as a soft mass, pain, thrombosis, embolism, or bleeding depending on its location and size.

\medterm{Venous Blood Gas} The clinical laboratory measurement of pH, partial pressure of carbon dioxide ($p\text{CO}_2$), bicarbonate ($\text{HCO}_3^-$), and base excess in a peripheral or central venous blood sample to assess systemic acid-base balance, ventilatory status, and tissue perfusion.

\synonyms
Venous Blood Gas; VBG; Central Venous Blood Gas
\medterm{Venous Capacitance} The physical and physiological capacity of the systemic venous circulation to accommodate large changes in blood volume with minimal alterations in venous pressure, functioning as the primary high-capacitance, low-resistance reservoir holding approximately 60 to 70 percent of total blood volume at rest.

\medterm{Venous Diseases} Conditions affecting veins, including blood clots, varicose veins, poor return of blood, inflammation, ulcers, injuries, and birth-related malformations. They may cause pain, swelling, skin changes, bleeding, or dangerous clots traveling to the lungs.

\medterm{Venous Hum} A harmless continuous humming sound heard over the neck veins, most often in healthy children. It changes or disappears with body position or gentle pressure on the vein, helping distinguish it from heart murmurs that need further evaluation.

\synonyms
Cervical Venous Hum; Innocent Venous Hum

\medterm{Venous Insufficiency} Impaired return of blood through the veins, usually in the legs, because valves allow backward flow, veins are blocked, or the calf-muscle pump is ineffective. Blood may pool and cause aching, heaviness, swelling, varicose veins, itching, darkened or hardened skin, stasis dermatitis, and difficult-to-heal ulcers near the ankle.

\medterm{Venous Leak} Difficulty maintaining an erection because blood leaves the penis too readily during erection. It may reflect vascular, nerve, hormonal, or medication-related problems and should be evaluated as erectile dysfunction.
\synonyms
Veno-Occlusive Dysfunction; Cavernous Venous Leakage

\medterm{Venous Pressure} The pressure of blood inside the veins. It depends on blood volume, vein width, gravity, and how well the right side of the heart pumps; abnormally high pressure can cause swelling, enlarged veins, or fluid buildup.

\medterm{Venous Return} The flow of blood through the veins back to the heart. It is helped by vein valves, muscle movement, breathing, and pressure changes in the chest.

\medterm{Venous Stasis Dermatitis} Inflammation of the skin of the lower legs caused by chronic swelling and pooling of blood and fluid, usually from chronic venous insufficiency. It may cause itching, redness, scaling, skin darkening, thickening, or ulcers; treatment addresses the swelling and protects the skin.

\medterm{Venous Stasis Ulcer} A long-lasting wound near the ankle caused by poor blood return through the leg veins. The leg may also be swollen, painful, or have darkened, thickened skin; compression and treatment of the vein problem may help.

\medterm{Venous Thromboembolism} A blood clot that forms in a vein, usually deep in the leg, and may travel to the lungs. A clot in the lungs can cause sudden breathlessness, chest pain, collapse, or death.

\medterm{Venous Thromboembolism In Pregnancy} A blood clot in a deep vein or lungs during pregnancy or after birth. Pregnancy increases clot risk, and symptoms such as one-sided leg swelling, chest pain, sudden breathlessness, or coughing blood can indicate a serious clot.

\medterm{Venous Thromboembolism Prophylaxis} Measures used to prevent dangerous vein clots, especially around surgery or hospitalization. They may include walking early, compression devices or stockings, and blood-thinning medicine when bleeding risk is acceptable.

\medterm{Venous Thrombosis} The formation of a blood clot (thrombus) within a vein, most commonly in the deep veins
of the lower extremities (deep vein thrombosis, DVT).

\medterm{Venous Ulcer} A chronic wound caused by impaired venous drainage, usually occurring near the ankle in a limb with chronic venous insufficiency.

\medterm{Ventilation} Movement of air into and out of the lungs; effective ventilation brings oxygen to the alveoli and removes carbon dioxide.

\medterm{Ventilation-Perfusion Mismatch} A mismatch between air reaching parts of the lungs and blood reaching those same parts. It can lower blood oxygen and occur with pneumonia, COPD, asthma, heart failure, or a lung clot.

\medterm{Ventilation-Perfusion Ratio} The relationship between air reaching the lungs and blood reaching the lung tissues. Good matching allows efficient oxygen exchange; poor matching can cause low blood oxygen.

\medterm{Ventilation-Perfusion Scan} A nuclear-medicine scan that compares air movement in the lungs with blood flow through them. Areas with blood flow blocked but preserved air movement may suggest a pulmonary embolism. The test uses a breathed tracer and an injected tracer and is interpreted with symptoms and other findings.

\medterm{Ventilator} A machine that supports or replaces spontaneous breathing by delivering air, oxygen, or an air-oxygen mixture under positive pressure. Noninvasive ventilators deliver support through a mask, while invasive ventilators connect through an endotracheal or tracheostomy tube. Ventilators can provide different patterns and levels of pressure or volume according to the device and clinical purpose.

\synonyms
Mechanical Ventilator
Breathing Machine

\medterm{Ventilator-Associated Pneumonia} A lung infection that develops after a person has used a breathing machine for at least two days. It may cause fever, worsening oxygen levels, cough, or abnormal lung findings and requires prompt evaluation and treatment.

\medterm{Ventilator-Induced Lung Injury} Lung damage caused or worsened by a breathing machine when pressure, volume, or oxygen levels are excessive. Clinicians use protective settings and close monitoring to support breathing while limiting further injury.

\medterm{Ventouse} A vacuum extractor placed on the baby's scalp to assist a vaginal birth when delivery needs help. It is used only when specific safety conditions are met and can cause temporary scalp swelling or bruising.

\medterm{Ventral} Relating to the belly or front side of the body; synonym for anterior in human anatomy.

\medterm{Ventral Body Cavity} The body cavity toward the front. It includes the thoracic cavity and the abdominopelvic cavity, which contain many of the body’s major organs.

\medterm{Ventral Hernia} Protrusion of abdominal tissue through a weakness in the front abdominal wall, including incisional and umbilical hernias. It may cause a bulge or pain and requires urgent assessment if incarceration or strangulation is suspected.

\medterm{Ventral Root} The front root of a spinal nerve. It carries movement and other outgoing signals from the spinal cord.

\medterm{Ventricle} One of the two lower chambers of the heart that pump blood out of it. The right ventricle sends blood to the lungs to receive oxygen; the left ventricle sends oxygen-rich blood to the rest of the body.

\medterm{Ventricular} Relating to a ventricle. In the heart, it concerns one of the lower pumping chambers; in the brain, it concerns the fluid-filled spaces that contain cerebrospinal fluid. The organ gives the term its specific meaning.

\medterm{Ventricular Aneurysm} A localized, thin, fibrous, dyskinetic or akinetic bulging outpouching of the scarred left ventricular wall developing as a late mechanical complication of transmural myocardial infarction (typically anterior MI), predisposing to heart failure, refractory ventricular arrhythmias, and apical mural thrombus formation.

\medterm{Ventricular Arrhythmia} An abnormal heart rhythm arising from the ventricles. It may cause palpitations, dizziness, fainting,
chest discomfort, or sudden cardiac arrest; fainting or severe symptoms require emergency care.

\medterm{Ventricular Arrhythmias} Abnormal heart rhythms arising from the ventricles. They range from premature ventricular contractions to ventricular tachycardia and fibrillation, which can cause sudden cardiac arrest.

\medterm{Ventricular Assist Device} A mechanical pump that helps a weakened heart move blood. It may support the left side, right side, or both and can be temporary or long-term, but bleeding, infection, stroke, and clots are important risks.

\medterm{Ventricular Bigeminy} A cardiac rhythm pattern on electrocardiogram in which every normal sinus beat is consistently followed by a premature ventricular complex (PVC) in a regular alternating 1:1 sequence, commonly provoked by electrolyte derangements, digitalis toxicity, or underlying ischemic heart disease.

\medterm{Ventricular Fibrillation} A chaotic, disorganized ventricular rhythm in which the ventricles quiver rather than contract, producing no effective cardiac output and cardiac arrest. It is a shockable rhythm; immediate CPR and unsynchronized defibrillation are required, with treatment of the underlying cause.

\synonyms
V-Fib

\medterm{Ventricular Hypertrophy} Thickening of the wall of one of the heart’s lower chambers. It may develop from long-term high blood pressure, valve disease, inherited heart muscle disease, or intense training and can lead to stiffness, abnormal rhythms, or heart failure.

\medterm{Ventricular Myocardium} The heart muscle forming the walls of the ventricles. Its contraction pumps blood to the lungs and the rest of the body.

\medterm{Ventricular Non-Compaction} A primary genetic cardiomyopathy characterized by a thickened myocardial wall with a non-compacted endocardial layer displaying prominent, deep intertrabecular recesses and deep trabeculations, typically affecting the left ventricular apex and lateral wall.

\synonyms
LVNC; Non-Compaction Cardiomyopathy

\medterm{Ventricular Premature Contraction} Alternate name; see \textit{Premature ventricular contractions (PVCs)}.

\medterm{Ventricular Remodeling} Progressive alterations in left ventricular geometry, mass, chamber volume, and spherical shape occurring after acute myocardial infarction or chronic pressure/volume overload, driven by neurohormonal activation, cardiomyocyte hypertrophy, and interstitial fibrosis.
\synonyms
Post-Infarction Remodeling; Cardiac Remodeling

\medterm{Ventricular Rupture} A tear in a wall of one of the heart’s lower chambers, often after a heart attack. It can cause life-threatening bleeding around the heart.

\medterm{Ventricular Septal Defect} A hole in the wall between the heart’s lower chambers. Blood may flow from the stronger left side to the right and increase blood flow to the lungs. Small holes may cause no problems; larger ones can cause heart failure or lung-vessel damage.

\medterm{Ventricular Septal Defect Cross-Reference} Alternate name; see \textit{Ventricular septal defect}.

\medterm{Ventricular Tachycardia} A rapid heart rhythm that starts in the lower chambers of the heart. It may cause palpitations, dizziness, fainting, shock, or cardiac arrest and can be life-threatening.

\medterm{Ventriculoatrial} Relating to a connection between a ventricle and an atrium. In brain surgery, a ventriculoatrial shunt is a tube that drains extra fluid from a brain cavity into the right upper chamber of the heart.

\medterm{Ventriculography} Imaging used to measure how well the heart’s ventricles pump blood. A radionuclide scan or cardiac MRI can estimate the amount of blood ejected with each heartbeat and show abnormal wall movement; the method is chosen for the clinical question.

\medterm{Ventriculoperitoneal Shunt} A surgically implanted system that drains cerebrospinal fluid from a brain ventricle into the abdominal cavity, where the fluid
can be absorbed.

\medterm{Ventrosuspension} Surgical fixation of a displaced or prolapsed uterus to the anterior abdominal wall or nearby supporting tissues.

\medterm{Venturi Mask} A high-flow oxygen delivery device that uses the Venturi principle to supply a precise, fixed fraction of inspired oxygen ($\text{FiO}_2$, typically 24\% to 50\%). Oxygen passing through a calibrated narrow jet nozzle creates a pressure drop that draws in a predictable volume of room air through color-coded entrainment ports. Because the total gas flow meets or exceeds the patient's peak inspiratory flow rate, the delivered oxygen concentration remains constant regardless of the patient's breathing rate or depth, making it vital for managing hypoxemia in patients with chronic obstructive pulmonary disease at risk of hypercapnic respiratory failure.
\synonyms{Air-Entrainment Mask; Venturi Oxygen Mask; High-Flow Oxygen Mask}

\medterm{Venturi Principle} A physical principle in fluid dynamics stating that as a fluid or gas flows through a constricted section of a pipe or nozzle, its velocity increases and its static lateral pressure decreases. In medicine, this localized pressure drop creates suction that entrains ambient air or fluids, forming the operational basis for air-entrainment oxygen masks (Venturi masks), jet nebulizers, suction aspirators, and explaining hemodynamic phenomena such as systolic anterior motion of the mitral valve in hypertrophic obstructive cardiomyopathy.
\synonyms{Venturi Effect; Bernoulli-Venturi Principle; Air-Entrainment Principle}

\medterm{Venule} A small blood vessel that collects blood from capillaries and carries it into larger veins. Venules also allow white blood cells and fluid to leave the bloodstream during inflammation.

\medterm{Venules} Small blood vessels that collect blood from capillaries and join to form veins. Their walls also participate in inflammation by allowing white blood cells to leave the bloodstream.

\medterm{Verapamil} A non-dihydropyridine calcium-channel blocker used for selected hypertension, angina, and supraventricular tachyarrhythmias; it can cause bradycardia and constipation.

\medterm{Verapamil Toxicity} Acute severe poisoning from the non-dihydropyridine calcium channel blocker verapamil, characterized by profound peripheral vasodilation, profound myocardial depression, high-grade atrioventricular block, refractory bradycardia, and cardiogenic shock; managed with intravenous calcium, high-dose insulin euglycemia therapy, and lipid emulsion.

\medterm{Vermicides} Medicines or chemicals that kill parasitic worms. The correct treatment depends on the worm, the body site affected, the person’s age, and other health conditions.

\medterm{Vernal Keratoconjunctivitis} A severe, recurrent allergic eye inflammation that primarily affects young boys in warm, sunny climates and often flares up during the spring and summer (spring catarrh). It causes intense eye itching, light sensitivity, watering, and thick, stringy mucus. Characteristic findings include large, cobblestone-like bumps on the underside of the upper eyelid (giant papillae) and small white dots along the border of the cornea (Horner-Trantas dots). In severe cases, inflammation and eye-rubbing can create a non-healing sore on the clear front of the eye (a shield ulcer), which can lead to corneal scarring and vision changes. The condition typically improves on its own after puberty.

\synonyms
VKC; Spring Catarrh; Vernal Catarrh; Vernal Conjunctivitis

\medterm{Verner-Morrison Syndrome} A rare neuroendocrine tumor syndrome (WDHA syndrome) caused by a vasoactive intestinal peptide-secreting pancreatic islet cell tumor (VIPoma), characterized by the tetrad of severe watery diarrhea ('pancreatic cholera'), hypokalemia, achlorhydria (hypochlorhydria), and metabolic acidosis.
\synonyms
VIPoma Syndrome; WDHA Syndrome

\medterm{Vernix} A white, greasy protective coating on the skin of a fetus or newborn, composed mainly of water, lipids, and shed skin cells.

\medterm{Vernix Caseosa} A white, greasy protective coating on the skin of a fetus or newborn, composed mainly of water, lipids, and shed skin cells.

\medterm{Verruca} Alternate name; see \textit{Wart}.

\medterm{Verruca Plana} Small, smooth, relatively flat warts caused by HPV, often occurring in groups on the face, hands, or legs.

\synonyms
Flat Warts; Juvenile Warts

\medterm{Verrucas} Verrucas are rough skin growths on the feet caused by human papillomavirus. They may be painful with pressure and often resolve without treatment, although treatment can reduce discomfort or spread.

\medterm{Verruca Vulgaris} A common, benign, hyperkeratotic, exophytic cutaneous papilloma with a rough, verrucous ('warty') surface and thrombosed capillary black dots, caused by localized epidermal infection by human papillomavirus (HPV types 2, 4, 7, and 1).
\synonyms
Common Wart; Viral Verruca

\medterm{Verruciform Xanthoma} A rare, noncancerous growth usually found in the mouth or genital area. It may look like a wart, so examination and sometimes a tissue sample are needed to distinguish it from cancer.

\medterm{Verrucose} Having a rough, thickened, wart-like, or finger-like surface. It is a descriptive term for a skin, mucosal, or other lesion and does not by itself identify whether the lesion is benign, infectious, or cancerous.

\medterm{Verrucous} Having a rough, thickened, irregular, or wart-like surface. The appearance can occur in benign growths, infections, chronic inflammation, or some cancers.

\medterm{Verrucous Carcinoma} A low-grade, well-differentiated variant of squamous cell carcinoma with a warty,
exophytic appearance. It is locally invasive but rarely metastasizes, commonly affecting
the oral cavity, skin, or genitalia.

\medterm{Version} A change in the position or presentation of a body structure. In pregnancy, fetal version is a procedure or movement that changes the baby’s position, often to place the head downward before birth.

\medterm{Vertebra} One of the bones forming the spinal column. Each vertebra helps support the body, allows movement, provides attachment for muscles, and surrounds and protects the spinal cord.

\medterm{Vertebrae} Vertebrae are the bones of the spinal column that protect the spinal cord, support the body, allow movement, and provide attachment for muscles and ligaments.

\medterm{Vertebral Artery} One of two arteries that travel through the neck and join to supply the brainstem, cerebellum, and back of the brain. Narrowing or blockage can cause dizziness, imbalance, or stroke.

\medterm{Vertebral Artery Dissection} A tear in the inner lining of a vertebral artery that lets blood enter the vessel wall, forming an intramural hematoma and possibly narrowing the artery or creating a pseudoaneurysm. It may occur spontaneously or after neck injury and can cause sudden neck or occipital headache, dizziness, imbalance, or a posterior-circulation transient ischemic attack or stroke.
\synonyms
VAD; Spontaneous Vertebral Artery Dissection

\medterm{Vertebral Canal} The passage formed by the stacked spinal bones that surrounds and protects the spinal cord. Narrowing, bleeding, tumors, fractures, or disc disease inside it can compress nerves and cause weakness or pain.

\medterm{Vertebral Column} The chain of bones running from the neck to the pelvis. It supports the head and trunk, protects the spinal cord, allows bending and twisting, and provides attachment for muscles and ribs.

\medterm{Vertebral Compression Fracture} Collapse or cracking of one of the bones in the spine, often from osteoporosis, injury, or cancer. It may cause sudden or gradual back pain, loss of height, or a rounded back and can sometimes press on nerves.

\medterm{Vertebral Hemangioma} Alternate name; see \textit{Spinal hemangioma}.

\medterm{Vertebral Tumor} A growth arising in or near a spinal bone. It may be benign or cancerous, start in the spine or spread there, and cause pain, fracture, nerve compression, weakness, or paralysis.

\medterm{Vertebrobasilar Circulatory Disorders} Conditions that reduce blood flow through the vertebral or basilar arteries supplying the brainstem, cerebellum, and posterior brain; symptoms may include vertigo, double vision, imbalance, weakness, or stroke.

\medterm{Vertebrobasilar Insufficiency} Reduced blood flow through arteries supplying the back of the brain. It may cause dizziness, double vision, imbalance, weakness, trouble speaking, or stroke-like symptoms and can be serious when sudden.

\medterm{Vertebrobasilar Insuffiency} Source spelling variant of \textit{vertebrobasilar insufficiency}, reduced blood flow in the vertebral or basilar circulation supplying the brainstem, cerebellum, and posterior brain.

\medterm{Vertebroplasty} A procedure in which medical cement is injected into a collapsed spinal bone to help stabilize it and sometimes reduce pain. It is used selectively because cement leakage, infection, nerve injury, or other complications can occur.

\medterm{Vertical} Oriented from superior to inferior or perpendicular to a horizontal plane. In medicine, it may describe position, transmission from parent to child, or a body-oriented imaging or movement axis.

\medterm{Vertical Talus} A stiff flatfoot present from birth in which one of the ankle bones points downward and the middle of the foot is bent upward. It may occur alone or with other conditions and usually needs early orthopedic treatment.

\medterm{Vertical Transmission} Passage of an infection or genetic condition from parent to child during pregnancy, delivery, or breastfeeding.

\medterm{Vertigo} A false feeling that you or your surroundings are moving, spinning, tilting, or shifting when they are not. It is a symptom rather than a disease and often comes from the inner ear, but brain and other conditions can also cause it.

\medterm{Vertigo-Associated Disorders} Vertigo-associated disorders cause an illusion of movement from vestibular, neurologic, visual, medication-related, or systemic dysfunction and may be accompanied by imbalance, nausea, or hearing symptoms.


\synonyms
Vestibular Disorders; Peripheral Vestibulopathy; Labyrinthine Disorders

\medterm{Very Long-Chain Acyl-CoA Dehydrogenase Deficiency} A rare inherited disorder in which cells cannot use certain long-chain fats for energy, especially during fasting or illness. Infants may develop low blood sugar, heart disease, or liver problems; older children or adults may have exercise-triggered muscle breakdown and dark urine.

\synonyms
VLCAD Deficiency; VLCADD

\medterm{Very-Low-Density Lipoprotein} A particle made by the liver that carries mainly triglycerides through the blood. High levels can contribute to fatty deposits in artery walls and increase cardiovascular risk, although the result is interpreted with the full lipid profile.

\synonyms
VLDL

\medterm{Vesical} Relating to the urinary bladder. For example, vesical inflammation affects the bladder lining, a vesical tumor is a bladder growth, and the vesical neck is the bladder outlet leading toward the urethra.

\medterm{Vesicant} A substance that causes blistering and tissue injury on contact. In oncology, vesicant drugs can cause severe extravasation injury if they escape from a vein into surrounding tissue.

\medterm{Vesicants} Substances that can cause severe blistering and tissue damage when they touch skin or escape from a vein. Some chemotherapy medicines and industrial chemicals are vesicants.

\medterm{Vesicle} A small, circumscribed, raised skin lesion containing clear, serous, or sometimes bloody fluid, usually less than 1 cm in diameter. A vesicle is a small blister; causes include viral infection, contact irritation, eczema, burns, friction, and immune-mediated skin disease.

\medterm{Vesicles} Plural of vesicle; see \textit{Vesicle}.

\medterm{Vesicoureteral Reflux} Backward flow of urine from the bladder toward the kidneys because the valve where a ureter enters the bladder does not work normally. It can cause repeated kidney infections and scarring, especially in children.

\synonyms
VUR; Urinary Reflux

\medterm{Vesico-Ureteric Reflux} Backward flow of urine from the bladder into one or both ureters and sometimes the kidneys, increasing the risk of recurrent urinary infection and renal scarring.

\medterm{Vesicovaginal and Ureterovaginal Fistula} Abnormal connections between the urinary tract and vagina: a vesicovaginal fistula connects bladder to vagina, while a ureterovaginal fistula connects ureter to vagina. Both commonly cause continuous urine leakage and require imaging and reconstructive urologic or gynecologic care.

\medterm{Vesicovaginal Fistula} An abnormal connection between the bladder and vagina that causes continuous or frequent leakage of urine through the vagina. It may follow childbirth injury, pelvic surgery, radiation, or advanced infection and usually requires specialist repair.

\medterm{Vesicular} Describing or containing small, fluid-filled blisters. A vesicular rash may occur with viral infection, allergy, eczema, burns, or other conditions; its appearance alone does not establish the cause.

\medterm{Vesicular Breathing} The normal soft, gentle breath sound heard over most outer lung areas with a stethoscope. Changes in loudness or quality may suggest fluid, collapse, blockage, infection, or other lung disease.

\medterm{Vesicular Palmoplantar Eczema} An itchy eczema affecting the palms or soles with small, deep-seated fluid-filled vesicles, followed by peeling, scaling, or fissures. Triggers include sweating, irritants, allergy, stress, and fungal infection elsewhere; persistent disease may need specialist treatment.

\medterm{Vesiculobullous Dermatoses} A group of skin diseases that cause small or large fluid-filled blisters. Causes include infection, allergy, autoimmune disease, inherited conditions, medicines, and injury; examination and sometimes a biopsy identify the cause.

\medterm{Vessel} A vessel is a tube carrying blood, lymph, air, urine, bile, or another fluid through the body, depending on context.

\medterm{Vestibular} Relating to the vestibular system of the inner ear and brain, which detects head motion and position and helps maintain balance and gaze stability.

\medterm{Vestibular Disorder} A problem affecting the inner-ear balance organs or the brain pathways that process balance signals. It may cause dizziness, vertigo, unsteadiness, nausea, abnormal eye movements, or difficulty walking. Causes include inner-ear disease, infection, head injury, medicines, migraine, reduced blood flow, or disorders of the brain or nerves.

\synonyms
Vestibular Balance Disorder

\medterm{Vestibular Migraine} Repeated attacks of dizziness or vertigo related to migraine, with or without headache. Episodes may include sensitivity to movement, light, sound, or visual patterns; other inner-ear and brain causes must be considered.

\medterm{Vestibular Neuritis} Sudden inflammation or dysfunction of the balance nerve, causing severe spinning, nausea, vomiting, and unsteadiness that may last for days. Hearing is usually unaffected, but urgent assessment may be needed to exclude a stroke.

\medterm{Vestibular Papillomatosis} A harmless normal variation in which several soft, small, finger-like bumps occur on the inner vulvar folds. It can resemble genital warts but is not caused by an infection and usually needs no treatment.

\medterm{Vestibular Rehabilitation} Exercise-based therapy that improves balance, gaze stability, and adaptation in people with
vestibular disorders or dizziness.

\medterm{Vestibular Schwannoma} A benign, slow-growing tumor of the vestibular portion of the eighth cranial nerve. See Acoustic Neuroma.

\medterm{Vestibular System} The inner-ear sensors and connected brain pathways that detect head movement and position and help maintain balance, posture, and steady vision. Problems can cause vertigo, unsteadiness, nausea, or blurred vision during movement.

\medterm{Vestibule} A balance area of the inner ear between the cochlea and semicircular canals. It contains sensors that detect head position and straight-line movement and send this information to the brain.

\medterm{Vestibulocochlear Nerve} The eighth cranial nerve, carrying information for hearing and balance from the inner ear to the brain.

\medterm{Vestibulodynia} Persistent burning, stinging, or tenderness at the entrance to the vagina, often triggered by touch, penetration, or inserting a tampon. Examination rules out infection and other causes; care may include pelvic-floor therapy, pain treatment, and support for sexual and emotional health.

\synonyms
Provoked Vestibulodynia; Vulvar Vestibulitis Syndrome; PVD

\medterm{Vestibulo-Ocular Reflex} An automatic response that moves the eyes opposite to head movement so vision stays steady while a person turns or walks. Problems with it can cause blurred vision, dizziness, or imbalance.

\medterm{Vestigial} Describing a structure or feature that has become greatly reduced through evolution and has little or no apparent original function.

\medterm{VEXAS Syndrome} A serious inflammatory illness that begins in adulthood because of a gene change acquired in blood-forming cells. It can cause persistent fever, painful joints, skin rashes, inflamed blood vessels or cartilage, anemia, and other blood abnormalities.

\medterm{V-Fib} Alternate name; see \textit{Ventricular fibrillation}.

\medterm{Viable} Capable of surviving or developing under suitable conditions. In clinical contexts it may describe a viable pregnancy, tissue, organism, cell, or treatment option.

\medterm{Vibration Problem} Symptoms or injury caused by repeated exposure to vibrating tools or equipment, which may affect nerves, blood vessels, muscles, and joints in the hands or arms.

\medterm{Vibrator} A device that produces mechanical vibration and may be used for sexual stimulation, therapeutic massage, or selected rehabilitation purposes.

\medterm{Vibrio} A genus of curved, Gram-negative bacteria; species include \textit{V. cholerae}, which causes cholera, and \textit{V. vulnificus}, which can cause severe wound or bloodstream infection.

\medterm{Vibrio Cholerae} A bacterium that causes cholera, a severe diarrheal illness spread mainly through contaminated water or food. It can cause sudden, large-volume watery diarrhea, dehydration, shock, and death without rapid fluid replacement.

\medterm{Vibrio Infection} Infection caused by Vibrio bacteria, acquired through contaminated seafood or exposure of wounds to warm seawater. It may cause watery diarrhea, wound infection, or rapidly worsening bloodstream infection, especially in people with liver disease or weakened immunity. Severe wound pain, blistering, fever, or shock can occur.

\synonyms
Vibriosis; Vibrio Bacteremia

\medterm{Vibrio Parahaemolyticus Infection} An acute halophilic bacterial gastroenteritis acquired through ingestion of raw or undercooked marine seafood, particularly shellfish. It typically manifests with explosive watery diarrhea, abdominal cramps, nausea, vomiting, low-grade fever, and headache within 12 to 24 hours of exposure, usually resolving spontaneously within 3 to 5 days without antibiotics.

\synonyms
Vibrio Parahaemolyticus Gastroenteritis; Halophilic Vibrio Infection

\medterm{Vibrios} Comma-shaped bacteria, usually Gram-negative, that include species capable of causing cholera and wound or gastrointestinal infections.

\medterm{Vibrio vulnificus Infection} Infection caused by a saltwater-associated bacterium that can produce rapidly progressive wound infection, sepsis, or severe gastroenteritis. Risk is higher with liver disease, iron overload, or immune suppression; raw oysters and seawater exposure are important sources.

\medterm{Vibrio Vulnificus Infections} Infections caused by Vibrio vulnificus, a gram-negative bacterium found in warm seawater and raw
seafood. It can cause wound infection, bloodstream infection, or gastroenteritis, with
particularly severe disease in people with chronic liver disease or other risk factors.

\medterm{Video-Assisted Thoracoscopic Surgery} A chest operation performed through several small openings using a camera and long instruments rather than a large incision between the ribs. It may be used for lung or pleural biopsy, removing part of a lung, or other chest procedures and often allows faster recovery, though bleeding, air leakage, infection, or conversion to open surgery can occur.

\synonyms
VATS; Video Thoracoscopy; Thoracoscopic Surgery

\medterm{Videofluoroscopic Swallow Study} A moving X-ray test that records swallowing after small amounts of food or liquid containing contrast are taken. It shows whether food enters the airway, where swallowing is delayed, and which textures or techniques may be safer.

\medterm{Video Laryngoscope} An airway device with a small camera that shows the entrance to the windpipe on a screen. It helps clinicians place a breathing tube, especially when a direct view is difficult, but does not remove the risks of the procedure.

\medterm{Video Laryngoscopy} A technique for viewing the entrance to the windpipe with a camera on a blade connected to a screen. It can help clinicians place a breathing tube when a direct view is difficult.

\medterm{Videostroboscopy} A voice examination that uses a camera and flashing light to create a slow-motion view of the vocal cords as they vibrate. It helps investigate hoarseness, voice strain, and vocal-cord injury.

\medterm{Vildagliptin} An oral dipeptidyl-peptidase-4 inhibitor used in some countries for type 2 diabetes, alone or with other medicines. It improves glucose-dependent insulin secretion; liver function and pancreatitis symptoms may require attention.

\medterm{Villi} Tiny finger-like projections lining the small intestine that greatly increase the area available to absorb nutrients. Similar projections in the placenta help exchange substances between the pregnant person and fetus.

\medterm{Villous Adenoma} A high-risk subtype of neoplastic colonic adenomatous polyp characterized histologically by long, finger-like, frond-like epithelial projections covering more than 75% of the polyp mass, associated with a high potential for malignant transformation and excessive secretory mucus/electrolyte loss (secretory diarrhea with hypokalemia).
\synonyms
Villous Colonic Polyp; Papillary Adenoma

\medterm{Villous Atrophy} Flattening or loss of the tiny projections lining the small intestine. This reduces nutrient absorption and may occur in celiac disease, some infections, immune disorders, or other intestinal conditions.

\medterm{Villus} A small finger-like projection, especially one of the intestinal mucosa that increases absorptive surface area; chorionic villi mediate exchange between placenta and fetus.

\medterm{Villus Atrophy} The histopathological flattening, blunting, and loss of the microscopic finger-like villi lining the small intestinal mucosa, accompanied by crypt hyperplasia and intraepithelial lymphocytosis, classically observed in celiac disease, tropical sprue, and autoimmune enteropathy.

\medterm{Vinca Alkaloids} Cancer medicines including vincristine and vinblastine that disrupt microtubules, structures needed for cell division. They can treat several cancers but may cause nerve damage, low blood counts, constipation, or other toxicity.

\medterm{Vincent Angina} An acute necrotizing infection of the gums and oropharynx associated with fusobacteria and spirochetes. It causes painful ulceration, bleeding, foul breath, and gingival necrosis and requires dental or medical treatment.

\medterm{Vincent Disease} Acute necrotizing ulcerative gingivitis associated with fusospirochetal bacteria, causing painful bleeding gums, ulceration, and foul breath.

\medterm{Vincent Infection} Vincent infection is acute necrotizing ulcerative gingivitis caused by mixed oral bacteria, producing painful bleeding gums, ulcers, and foul breath.

\medterm{Vincent's Angina} Acute necrotizing ulcerative gingivitis or tonsillitis associated with fusiform bacteria and spirochetes, causing painful ulcers, bleeding, and fetid breath.

\medterm{Vincent Stomatitis} A painful necrotizing infection of the gums caused by mixed anaerobic bacteria, producing foul breath, bleeding, ulceration, and gray necrotic tissue; it is also called acute necrotizing ulcerative gingivitis.

\medterm{Vinorelbine} A chemotherapy medicine that interferes with cell division. It is used for selected lung and breast cancers and can cause low blood counts, infection, constipation, nerve symptoms, nausea, or irritation if it leaks outside a vein.

\medterm{Viomycin} An older aminoglycoside antibiotic related to capreomycin that was used against tuberculosis and other mycobacterial infections; kidney and hearing toxicity limit its use.

\medterm{VIPoma} A rare functioning neuroendocrine tumor arising from the non-beta islet cells of the pancreas (or occasionally extra-pancreatic sympathetic ganglia in children) that hypersecretes vasoactive intestinal peptide (VIP), causing refractory secretory diarrhea and hypokalemia.
\synonyms
Vasoactive Intestinal Peptide Tumor; Verner-Morrison Neoplasm

\medterm{VIPoma Syndrome} A rare neuroendocrine tumor syndrome arising predominantly from the non-beta islet cells of the pancreas that hypersecretes vasoactive intestinal peptide (VIP). Clinically characterized by massive, secretory, watery diarrhea persisting during fasting, severe hypokalemia, and hypochlorhydria or achlorhydria (the WDHA syndrome), accompanied by dehydration and metabolic acidosis.

\synonyms
WDHA Syndrome; Verner-Morrison Syndrome; Pancreatic Cholera

\medterm{Viraemia} The presence of viruses in the bloodstream. It may be brief or long-lasting and can allow the infection to spread from its original site to other tissues.

\medterm{Viral} Caused by or relating to a virus. Viruses enter living cells to multiply and can cause illnesses ranging from mild colds to serious brain, lung, liver, or bleeding diseases; testing and treatment depend on the specific virus.

\medterm{Viral Antibody} A protective protein made by the immune system after infection or vaccination that recognizes a virus. Blood tests for these antibodies may show past exposure, immune response, or selected current infections, depending on timing and test design.

\medterm{Viral Arthritis} Joint inflammation caused by or associated with a viral infection. It may cause sudden pain, swelling, and stiffness in several joints and often improves as the infection resolves, although some viruses can cause prolonged symptoms.

\medterm{Viral Assembly} The stage of viral reproduction when newly made viral genetic material and proteins are put together to form complete new virus particles.

\medterm{Viral Attachment} The first step of infection, when proteins on a virus bind to matching molecules on a host cell. This binding helps determine which species, tissues, or cell types the virus can infect.

\medterm{Viral Carcinogenesis} Cancer development caused or promoted by a viral infection. Some viruses alter cell growth directly, weaken immune protection, or cause long-lasting inflammation; examples include human papillomavirus, hepatitis B and C viruses, and Epstein-Barr virus.

\medterm{Viral Conjunctivitis} Infection of the thin lining over the eye and inside the eyelids, usually causing redness, watering, irritation, and discharge. It spreads easily through contact with infected secretions or contaminated hands and surfaces.

\medterm{Viral Culture} Growth of a virus in susceptible cell cultures to identify an infectious agent. It can be useful for
selected viruses, but molecular and antigen-based tests are often faster or more sensitive;
the appropriate method depends on the suspected infection and specimen.

\medterm{Viral Encephalitis} Inflammation of the brain caused by a virus. It may cause fever, headache, confusion, behavior change, seizures, weakness, or reduced consciousness. Severe cases can be life-threatening.

\medterm{Viral Entry} The process by which a virus gets inside a host cell and releases its genetic material so the cell can make more virus.

\medterm{Viral Exanthem} A widespread rash caused by a viral infection. It may occur with fever, tiredness, sore throat, cough, diarrhea, or other symptoms; the rash pattern and overall illness help determine the cause.

\medterm{Viral Gastroenteritis} Alternate name; see \textit{Stomach Flu}.

\medterm{Viral Haemorrhagic Fever} A severe illness caused by certain viruses that can damage blood vessels, disturb clotting, and cause bleeding, shock, or organ failure. The way it spreads and its treatment depend on the virus.

\medterm{Viral Hemorrhagic Fevers} Severe illnesses caused by certain viruses that can damage blood vessels, disrupt clotting, and lead to bleeding, shock, or organ failure. Some spread easily and may require special infection-control measures.
\synonyms
VHF; Hemorrhagic Fevers; Viral Hemorrhagic Disease

\medterm{Viral Hepatitis} Liver inflammation caused by hepatitis viruses A, B, C, D, or E. Hepatitis A and E usually spread through contaminated food or water; B, C, and D can spread through blood or body fluids and may become long-term infections that damage the liver.

\synonyms
Infectious Hepatitis

\medterm{Viral Hepatitis Types} Five viruses that can inflame the liver. They differ in how they spread, whether infection can become long-term, the complications they cause, and whether vaccination or curative treatment is available.

\medterm{Viral Immune Evasion} Ways viruses avoid or weaken the body’s immune defenses, allowing them to survive, multiply, or return after infection. These methods can make infections harder to clear and vaccines or treatments less effective.

\medterm{Viral Infection} An illness caused by a virus entering body cells and making copies of itself. Viruses can damage cells directly or cause inflammation as the immune system responds. Viral infections range from mild illnesses such as colds and cold sores to serious diseases such as influenza, COVID-19, hepatitis, and HIV. Some viruses remain inactive in the body and can become active again later, such as the chickenpox virus causing shingles.

\synonyms
Viral Infections; Viral Disease; Viral Illness

\medterm{Viral Infections and Pregnancy} Viral infections during pregnancy may affect maternal illness, fetal development, placenta, or newborn health. Evaluation considers the virus, gestational age, symptoms, immunity, exposure, testing, treatment, vaccination, and measures to reduce perinatal transmission.

\medterm{Viral Infections of the Mouth} Infections of oral skin or mucosa caused by viruses such as herpes simplex, varicella-zoster, enteroviruses, HPV, or EBV. They may cause blisters, ulcers, warts, sore throat, or systemic symptoms, and treatment depends on the virus and host risk.

\medterm{Viral Latency} A state in which a virus remains inside cells without actively making new virus particles. Some latent viruses can become active again, especially when the immune system is weakened or the body is under stress.

\medterm{Viral Latency Mechanisms} The ways some viruses remain inactive inside cells while avoiding immune detection. Reactivation can occur when immune defenses are weakened or after other triggers, depending on the virus.

\medterm{Viral Load} The amount of virus or viral genetic material in a body fluid or tissue, measured by a laboratory assay; trends can help assess disease activity or response to antiviral treatment.

\medterm{Viral Meningitis} Infection and inflammation of the coverings of the brain and spinal cord caused by a virus. It may cause fever, headache, neck stiffness, vomiting, light sensitivity, or confusion. Bacterial meningitis can cause similar symptoms and is more dangerous.

\medterm{Viral Myocarditis} Inflammation of the heart muscle caused by a viral infection or by the immune response after an infection. It may cause tiredness, chest discomfort, palpitations, heart failure, shock, or dangerous abnormal heart rhythms.

\synonyms
Infectious Viral Myocarditis; Acute Viral Myocarditis

\medterm{Viral Pathogenesis} The ways a virus enters the body, infects and spreads through cells, and causes illness. Damage may come directly from infected cells or from the immune system’s response to the infection.

\medterm{Viral Pericarditis} Inflammation of the sac around the heart caused by a viral infection. It often follows a respiratory or stomach illness and may cause sharp chest pain that changes with breathing or position.

\synonyms
Acute Viral Pericarditis; Viral Pericardial Inflammation

\medterm{Viral Pharyngitis} Pharyngeal inflammation caused by viral pathogens such as adenovirus, rhinovirus, coronavirus, Epstein--Barr virus, or influenza virus. Characterized by cough, rhinorrhea, conjunctivitis, and lack of prominent tonsillar exudates, resolving with supportive care.

\synonyms
Viral Sore Throat; Acute Viral Pharyngotonsillitis; Viral Upper Respiratory Pharyngitis

\medterm{Viral Pneumonia} Infection and inflammation of lung tissue caused by a virus, producing cough, fever, chills, breathlessness, chest discomfort, fatigue, and variable impairment of oxygenation. Severe disease can cause respiratory failure or secondary bacterial infection.

\medterm{Viral Release} The stage when newly made viruses leave an infected cell and spread to other cells. Some leave by pushing through the cell membrane, while others are released when the infected cell breaks apart.

\medterm{Viral Replication} The process by which a virus uses an infected cell to copy its genetic material and make viral proteins. The exact steps differ among viruses, but the result is production of new virus particles.

\medterm{Viral Replication Cycle} The steps a virus uses to make more copies inside a living cell. Usually the virus attaches to the cell, enters, releases its genetic material, makes new parts, assembles new viruses, and leaves to infect other cells.

\medterm{Viral Rhinitis} Acute inflammation and edema of the nasal mucous membranes caused by respiratory viral pathogens, predominantly rhinoviruses, coronaviruses, respiratory syncytial virus, or parainfluenza viruses. Clinically recognized as acute viral coryza or the common cold, it manifests with clear anterior rhinorrhea, nasal congestion, sneezing, itchy throat, low-grade malaise, and temporary hyposmia, generally resolving without antimicrobial therapy within 7 to 10 days.

\synonyms
Acute Viral Rhinitis; Acute Coryza; Common Cold Rhinitis

\medterm{Viral Structure} The parts of a virus: genetic material made of DNA or RNA inside a protective protein coat. Some viruses also have an outer fatty covering taken from the infected cell’s membrane.

\medterm{Viral Uncoating} The step after a virus enters a cell when its protective outer coat comes apart and releases the genetic material. The cell can then be directed to make new viral components.

\medterm{Virchow Node} A firm, non-tender, pathologically enlarged left supraclavicular lymph node (Troisier sign) located in the left supraclavicular fossa, which receives abdominal lymphatic drainage via the thoracic duct and signals metastatic spread of abdominal malignancies (particularly gastric adenocarcinoma).
\synonyms
Troisier Node; Left Supraclavicular Sentinel Node

\medterm{Virchow-Robin Spaces} Small fluid-filled spaces around blood vessels as they pass through the brain. They are a normal part of the spaces that surround and support these vessels.

\synonyms
Perivascular Spaces

\medterm{Virchow's Triad} Three factors that increase the chance of a blood clot forming: slow blood flow, injury to the inside of a blood vessel, and blood that clots too easily. These factors can contribute to clots in deep veins or clots that travel to the lungs.
\synonyms{Virchow Triad; Triad of Virchow; Thrombogenesis Triad}

\medterm{Virchow Triad} Three factors that increase the chance of a blood clot forming: injury to the inside of a blood vessel, slow blood flow, and blood that clots too easily.
\synonyms
Virchow's Triad; Thrombosis Pathophysiological Triad

\medterm{Viremia} The presence of a virus in the bloodstream. It may be brief or persistent and can allow infection to spread to other organs. Laboratory tests may detect the virus itself or its genetic material.

\medterm{Viridans Streptococci} Viridans streptococci are a group of bacteria commonly found in the mouth that can cause dental infections and, after entering the blood, infect heart valves.

\medterm{Virilisation} The development of male secondary sexual characteristics in females, resulting from
excessive androgen production or exposure.

\medterm{Virilism} Development of masculinizing features in a female, usually from excess androgen exposure or production, such as hirsutism, deep voice, acne, and menstrual disturbance.

\medterm{Virilism Foetal} Prenatal or neonatal masculinization of a genetically female fetus or infant due to
excess androgen exposure. Causes include congenital adrenal hyperplasia, maternal
androgen-producing disease, or placental and fetal steroidogenic disorders.

\medterm{Virilization} Virilization is development of typically male secondary sexual features, such as hirsutism, deep voice, acne, increased muscle mass, or clitoromegaly, from androgen excess or androgen exposure.

\medterm{Virilize} To virilize is to develop or cause typically male physical changes from androgen excess or treatment, including facial hair, voice deepening, acne, muscle change, or genital enlargement.

\medterm{Virology} The study of viruses, including their structure, reproduction, spread, effects on the body, diagnosis, treatment, and prevention.

\medterm{Virotherapy} Treatment that uses a virus, often modified, to attack cancer cells or deliver a useful gene into cells. It is an evolving field; the possible benefits, risks, and availability depend on the specific therapy and disease.

\medterm{Virtual Colonoscopy} Virtual colonoscopy is CT colonography, which uses bowel preparation and gas or carbon-dioxide insufflation to create images of the colon; it can detect polyps and masses but cannot remove lesions or biopsy them.

\medterm{Virulence} How severely a microorganism can cause disease or tissue damage. It depends on the organism’s harmful properties and on the person’s immune defenses and other health factors.

\medterm{Virulence Factor} A feature of a microorganism that helps it survive in the body, avoid immune defenses, attach to cells, damage tissue, or cause disease. Examples include toxins, protective coatings, and enzymes.

\medterm{Virulent} Capable of causing severe disease or extensive harm. Virulence describes the harmfulness of a microorganism and is different from infectivity, which describes how easily it spreads or establishes infection.

\medterm{Virus} An infectious agent made of genetic material inside a protective protein coat, sometimes with an outer membrane. It can reproduce only inside living cells and may cause illness by damaging cells or triggering inflammation.

\medterm{Virus Activation} Virus activation is the switch from a dormant or inactive viral state to production of new viral components and particles.

\medterm{Virus Assembly} Virus assembly is the stage of the viral life cycle when viral genetic material and proteins are packaged into new virus particles.

\medterm{Virus Diseases} Illnesses caused by viruses. They range from mild infections of one area, such as a cold, to severe diseases affecting the lungs, liver, brain, blood, or several organs; prevention and treatment depend on the virus.

\medterm{Viruses} Viruses are infectious particles containing DNA or RNA within a protein coat, sometimes with a lipid envelope; they replicate only inside host cells and cause disease through direct injury or immune responses.

\medterm{Virus Integration} Virus integration is insertion of viral genetic material into the DNA of a host cell, a process used by some viruses and important in certain cancers.

\medterm{Virus Protein} A virus protein is a protein made from viral genetic instructions that helps form the virus, enter cells, replicate, or evade immunity.

\medterm{Virus Receptor} A virus receptor is a molecule on a host cell that a virus binds to in order to attach to and enter the cell.

\medterm{Virus Replication} Virus replication is the process by which a virus copies its genetic material and makes proteins and new virus particles inside a host cell.

\medterm{Virus Uncoating} Virus uncoating is removal of the virus's protective outer layer after entry into a host cell, releasing its genetic material for replication.

\medterm{Viscera} The internal organs of the body, particularly those within the thoracic and abdominal
cavities. The thoracic viscera include the heart, lungs, esophagus, trachea, and thymus.

\medterm{Visceral} Relating to the internal organs, especially those within the thorax or abdomen, rather than the body wall or limbs.

\medterm{Visceral Afferents} Sensory nerve fibers that carry information from internal organs to the spinal cord and brain. They convey signals related to stretch, chemical changes, inflammation, and internal pain and help regulate autonomic and reflex responses.

\medterm{Visceral Larva Migrans} A syndrome caused by migration of animal nematode larvae through human tissues, most often from Toxocara species. It may produce fever, eosinophilia, liver lesions, or ocular and neurologic disease.

\medterm{Visceral Leishmaniasis} A severe, potentially fatal systemic protozoan infection caused by species of the Leishmania donovani complex (Leishmania donovani and Leishmania infantum) transmitted to humans by the bite of female phlebotomine sandflies. Characterized by the intracellular dissemination of amastigotes within the reticuloendothelial system, it presents with prolonged irregular fever, progressive massive splenomegaly, hepatomegaly, profound pancytopenia (anemia, leukopenia, and thrombocytopenia), hypergammaglobulinemia, and severe wasting, fatal in over 95\% of untreated cases.

\synonyms
Kala-Azar; Black Fever; Dumdum Fever; Systemic Leishmaniasis

\medterm{Visceromegaly} Abnormal enlargement of one or more internal organs, such as the liver, spleen, or kidneys. It is a physical finding with causes that include infection, congestion, storage disease, inflammation, and cancer.

\medterm{Viscus} An internal organ, especially one in the chest or abdomen. The plural is viscera; examples include the heart, lungs, liver, stomach, and intestines.

\medterm{Vision} The sensory ability to perceive light and interpret visual information processed by the eyes and brain, enabling recognition of shapes, colors, depth, and motion. Normal vision depends on the proper function of the cornea, lens, retina, and optic nerve, as well as intact neural pathways to the visual cortex.

\synonyms
Eyesight; Sight; Visual Perception

\medterm{Vision And Hearing} Two major senses used to understand and communicate with the world. Problems may develop separately or together because of aging, illness, injury, medicines, or inherited conditions and can affect learning, safety, communication, and daily activities.

\medterm{Vision Impairment} Reduced visual acuity or visual-field function that affects daily activities. Causes include refractive
error, cataract, glaucoma, retinal disease, optic-nerve injury, and neurologic disease.
Evaluation and treatment depend on the cause; rehabilitation and visual aids may improve
function.

\medterm{Vision Loss} A decrease in the ability to see that cannot be fully corrected with eyeglasses, contact
lenses, or medical treatment. It ranges from partial impairment to complete blindness
and can result from conditions such as macular degeneration, glaucoma, diabetic
retinopathy, cataracts, and traumatic injury.

\medterm{Vision Problems} Changes in visual acuity, field, color, or clarity caused by refractive error, eye
disease, neurologic disease, or systemic illness. Sudden vision loss is an emergency.

\medterm{Vision Receptors} Light-sensing cells in the retina, mainly rods and cones, that convert light into nerve signals. Rods support dim-light vision, while cones provide color and sharp detail in brighter conditions.

\medterm{Vision Screening} Brief checks used to find possible eye or vision problems before symptoms are obvious. Newborn checks may detect cataracts or other serious disorders; screening in young children looks for poor vision, misaligned eyes, or risks for reduced vision and leads to a full eye examination when needed.

\medterm{Vision Test} An examination that measures visual acuity, fields, color, alignment, or eye structures to detect refractive error, retinal disease, optic-nerve dysfunction, glaucoma, or other causes of visual impairment.

\medterm{Visual Acuity} How clearly a person can see fine details, usually measured by reading letters or symbols at a set distance. It can be reduced by focusing problems, eye disease, or disorders affecting the visual nerves or brain.

\medterm{Visual Acuity Test} A test measuring the smallest standardized detail a person can resolve at a specified distance, usually reported as a fraction or decimal.

\medterm{Visual Agnosia} Inability to recognize familiar objects, faces, or symbols by sight even though the eyes may see adequately. It usually results from damage to brain areas that interpret vision, such as after a stroke or head injury.

\medterm{Visual Evoked Potential} A test that records the brain’s electrical response to flashes or changing visual patterns. It helps assess the nerve pathways that carry visual signals from the eyes to the brain.
\synonyms
VEP; Visual Evoked Response (VER)

\medterm{Visual Field} The entire area that can be seen while the eyes remain fixed in one position. Visual-field loss may result from retinal, optic-nerve, or brain lesions and is assessed by perimetry.

\medterm{Visual Field Defect} Loss or change in part of the area a person can see while looking straight ahead. It may result from disease of the retina, optic nerve, or brain, and the pattern can help identify where the problem is.

\medterm{Visual Field Perimetry} A test that maps how well a person sees in different parts of the visual field while looking at a fixed point. It can detect and monitor missing areas caused by eye-nerve, retinal, or brain disease.

\medterm{Visual Fields} The total area that can be seen while the eyes are fixed in one position; defects may indicate retinal, optic-nerve, or brain disease.

\medterm{Visual Field Test} A test that measures central and side vision and maps areas that are missing or less sensitive.

\medterm{Visual Impairment} Reduced vision that cannot be fully corrected by ordinary glasses or contact lenses. Causes include refractive error, cataract, glaucoma, retinal disease, optic-nerve injury, and disorders of the brain or nervous system.

\medterm{Visualized Laser-Assisted Prostatectomy} A procedure that uses a laser, guided by a camera, to remove or vaporize prostate tissue blocking urine flow. It is used for some people with an enlarged prostate and may cause bleeding, infection, scar narrowing, or temporary urinary and sexual problems.

\synonyms
VLAP

\medterm{Visual Neglect} A brain disorder in which a person fails to notice or respond to things on one side, even though the eyes may work normally. It often follows a stroke affecting the right side of the brain and causes neglect of the left side.

\synonyms
Hemispatial Neglect; Spatial Inattention

\medterm{Vital} Essential for life or relating to vital physiologic functions. Vital signs commonly include temperature, pulse, respiratory rate, blood pressure, and oxygen saturation.

\medterm{Vital Capacity} The greatest amount of air a person can breathe out after filling the lungs as fully as possible. It is measured during lung-function testing and may be reduced by lung disease, weak breathing muscles, or restriction of chest movement.

\medterm{Vital Signs} Basic measurements of physiologic status, commonly including temperature, pulse, respiratory rate, blood
pressure, and oxygen saturation. Values vary with age, activity, illness, and measurement
conditions; trends and clinical context are important.

\medterm{Vitamin A} A fat-soluble vitamin needed for vision, immune function, skin and mucosal health, and growth; deficiency causes night blindness, while excess can injure the liver and developing fetus.
\synonyms
Retinol; Retinoic Acid

\medterm{Vitamin A Blood Test} A blood test measuring circulating retinol, used to assess suspected vitamin A deficiency or toxicity in appropriate clinical settings.

\medterm{Vitamin A Deficiency} A nutritional disorder resulting from inadequate intake or absorption of vitamin A (retinol), causing dry eyes (xerophthalmia), night blindness, dry skin, and weakened immunity. Severe deficiency leads to corneal ulceration and is the leading cause of preventable childhood blindness worldwide; clinical eye stages are categorized using the WHO xerophthalmia classification.

\synonyms
Hypovitaminosis A; Retinol Deficiency; Vitamin A Lack

\medterm{Vitamin A: Retinal} Vitamin A (retinol) is a fat-soluble vitamin essential for vision, immune function, cell
differentiation, and reproduction. In the eye, retinol is oxidized to retinal
(retinaldehyde), which binds to opsin in the photoreceptor outer segments forming
rhodopsin (in rods) and photopsins (in cones).

\medterm{Vitamin A Toxicity} Harm caused by taking too much vitamin A or related medicines. It may cause nausea, vomiting, headache, dry skin, hair loss, liver injury, high calcium levels, and serious fetal abnormalities during pregnancy.

\medterm{Vitamin B} Vitamin B refers to a group of water-soluble vitamins that support energy metabolism, blood-cell formation, nerve function, and other processes. The group includes several distinct vitamins, not one single substance.

\medterm{Vitamin B1} Thiamine, a water-soluble vitamin that acts as a cofactor in carbohydrate and amino-acid metabolism and supports nerve function. Severe deficiency can cause beriberi or Wernicke encephalopathy.
\synonyms
Thiamine; Aneurin

\medterm{Vitamin B12} A vitamin needed for red-cell formation, DNA production, and normal nerve function. Deficiency can cause anemia, numbness, balance problems, or cognitive changes.
\synonyms
Cobalamin; Cyanocobalamin

\medterm{Vitamin B12: Cobalamin} Alternate name; see \textit{Vitamin B12}.

\medterm{Vitamin B12 Deficiency} Too little vitamin B12, which can cause large abnormal red cells, tiredness, numbness, balance problems, or nerve damage. Causes include poor intake, poor absorption, autoimmune disease, stomach surgery, and medicines.

\medterm{Vitamin B12 Deficiency Anemia} Anemia caused by too little vitamin B12 or poor absorption of it. It can cause tiredness, pale skin, a sore tongue, numbness, balance problems, and permanent nerve damage if untreated for a long time.

\synonyms
Vitamin B12 Deficiency Anaemia

\medterm{Vitamin B12 Deficiency In Older Adults} Deficiency caused by inadequate intake, malabsorption, autoimmune intrinsic-factor
deficiency, gastric surgery, or medications. It may cause megaloblastic anemia,
neuropathy, cognitive changes, gait impairment, and glossitis; neurologic injury can
occur without anemia.

\medterm{Vitamin B12 Level} The amount of vitamin B12 measured in blood; results are interpreted with symptoms and related tests because serum levels may not fully reflect tissue deficiency.

\medterm{Vitamin B1: Thiamine} Alternate name; see \textit{Vitamin B1}.

\medterm{Vitamin B2} Riboflavin, a water-soluble vitamin that helps cells produce energy and supports healthy skin, eyes, and the lining of the mouth. Deficiency can cause cracks at the mouth corners, a sore tongue, skin changes, or anemia.
\synonyms
Riboflavin; Lactoflavin

\medterm{Vitamin B2 Riboflavin} Alternate name; see \textit{Riboflavin}.

\medterm{Vitamin B3} A vitamin that helps cells use food for energy and supports the skin and nervous system. Severe deficiency causes pellagra, with diarrhea, skin inflammation, confusion, and other neurologic problems.

\synonyms
Nicotinic Acid; Nicotinamide

\medterm{Vitamin B5} A water-soluble vitamin needed to help the body use fats, carbohydrates, and proteins and make hormones and other substances. It is found in many foods, so deficiency is rare; very large supplement doses may cause stomach upset.

\synonyms
Pantothenate; Calcium Pantothenate

\medterm{Vitamin B6} A group of related vitamins needed for amino-acid metabolism, red-cell production, and nerve function. Excess supplements can cause sensory nerve damage.
\synonyms
Pyridoxine; Pyridoxal Phosphate

\medterm{Vitamin B6: Pyridoxine} Alternate name; see \textit{Vitamin B6}.

\medterm{Vitamin B7} A vitamin that helps the body use fats, carbohydrates, and proteins. Deficiency is uncommon but may cause hair loss, skin rash, or neurologic symptoms; large doses can interfere with some laboratory tests.
\synonyms
Vitamin H; Coenzyme R

\medterm{Vitamin B9} A vitamin needed for DNA production and healthy red blood cells. Deficiency can cause large-cell anemia and, during pregnancy, increases the risk of neural-tube defects; folic acid is recommended before conception and early in pregnancy.

\synonyms
Folic Acid; Pteroylglutamic Acid

\medterm{Vitamin B Complex} A group of eight water-soluble vitamins that act mainly as coenzymes or their precursors in cell metabolism. The group includes thiamine (B1), riboflavin (B2), niacin (B3), pantothenic acid (B5), pyridoxine (B6), biotin (B7), folate (B9), and cobalamin (B12), which together support energy metabolism, red-blood-cell formation, and nervous-system function. These vitamins occur together in whole grains, meat, eggs, and legumes; individual deficiencies cause distinct disorders, and combined supplements can prevent or correct deficiency. The sequence has gaps because early candidate substances were later excluded as nonessential: B8 (inositol) is synthesized by the body, B4 (choline, adenine) is not a true vitamin, and the provisional labels B10 and B11 were abandoned once their activity was identified as folate or para-aminobenzoic acid.
\synonyms
B Vitamins; Vitamin B Group

\medterm{Vitamin C} A water-soluble vitamin needed for collagen formation, wound healing, immune function, antioxidant activity, and absorption of iron from plant foods. Severe deficiency causes scurvy, with bleeding gums, bruising, poor wound healing, and weakness.
\synonyms
Ascorbic Acid

\medterm{Vitamin C Deficiency} A lack of vitamin C, a water-soluble vitamin needed for
collagen formation, wound healing, iron absorption, and antioxidant activity. It follows
prolonged low intake and, when severe, causes scurvy, with bleeding gums, easy bruising,
poor wound healing, joint pain, and weakness. Replacement and dietary correction reverse
the deficiency.

\medterm{Vitamin D} Vitamin D supports calcium and phosphate balance, bone mineralization, muscle function, and immune processes; deficiency causes rickets in children and osteomalacia or bone loss in adults.

\medterm{Vitamin D3} Cholecalciferol, a form of vitamin D made in skin and found in some foods and supplements. It helps the body absorb calcium and maintain healthy bones.

\medterm{Vitamin D: Calcitriol} Alternate name; see \textit{Calcitriol}.

\medterm{Vitamin D Deficiency} Too little vitamin D for the body’s needs. It can result from low intake or sunlight, poor absorption, or liver or kidney disease and may weaken bones, cause muscle weakness, or lead to rickets in children.

\medterm{Vitamin Deficiency Anemia} Anemia caused by inadequate vitamin B-12, folate, or other nutrients needed for normal red blood cell production.

\medterm{Vitamin E} Vitamin E is a group of fat-soluble compounds with antioxidant activity. It is found in foods such as nuts, seeds, and vegetable oils; high-dose supplements can increase bleeding risk and interact with medicines.
\synonyms
Tocopherol; Alpha-Tocopherol

\medterm{Vitamin E Deficiency} Deficiency of a fat-soluble antioxidant vitamin, usually caused by fat-malabsorption disorders or rare inherited conditions. Severe deficiency can cause peripheral neuropathy, ataxia, muscle weakness, retinal dysfunction, or hemolytic anemia.

\medterm{Vitamin E: Tocopherol} Alternate name; see \textit{Vitamin E}.

\medterm{Vitamin K} A fat-soluble vitamin required to activate several blood-clotting proteins and support bone metabolism; the main dietary form is phylloquinone, found especially in leafy green vegetables and some vegetable oils. Deficiency can cause bleeding and is treated or prevented with replacement.

\medterm{Vitamin K Deficiency} Too little vitamin K for normal production of clotting factors, causing an increased bleeding tendency. It can result from poor intake, malabsorption, liver disease, or certain medicines; newborns receive vitamin K prophylaxis.

\medterm{Vitamin K Deficiency Bleeding} Bleeding in a newborn or young infant caused by inadequate vitamin K-dependent clotting activity. It may occur early, classically, or late and can involve the brain or gastrointestinal tract.

\synonyms
VKDB; Neonatal Hypoprothrombinemia

\medterm{Vitamin K: Phylloquinone} Alternate name; see \textit{Vitamin K}.

\medterm{Vitamin K Prophylaxis} Giving vitamin K to a newborn soon after birth to prevent bleeding caused by vitamin K deficiency.

\medterm{Vitamins} Organic nutrients needed in small amounts for normal metabolism, growth, and tissue maintenance; some are water-soluble and others fat-soluble, and both deficiency and excess can cause disease.

\medterm{Vitamin Toxicity} Harm from excessive vitamin intake, most often involving fat-soluble vitamins that accumulate in the body. Vitamin A can cause headache, liver injury, and bone effects; vitamin D can cause hypercalcemia; and high-dose vitamin E can increase bleeding risk. Stop the source and treat the affected organ or metabolic abnormality.

\medterm{Vitiligo} Vitiligo is a condition in which areas of skin lose pigment because melanocytes are destroyed or stop functioning. It is not contagious, and treatment may include camouflage, topical medicines, or light therapy.

\medterm{Vitrectomy} Eye surgery that removes some or all of the vitreous, the clear gel that fills the middle of the eye. Small instruments are inserted through the wall of the eye to remove cloudy gel, blood, scar tissue, or a foreign object and to repair problems of the retina. It may be used for retinal detachment, diabetic eye disease, macular hole, macular pucker, severe eye injury, or serious eye infection. The removed gel may be replaced with a salt solution, air or gas, or silicone oil.

\medterm{Vitreous} Relating to the transparent gel filling most of the space between the lens and retina, or describing a glass-like appearance. Vitreous disorders can cause floaters, hemorrhage, detachment, or retinal traction.

\medterm{Vitreous Body} The transparent gel filling the large space between the lens and retina. It helps maintain the shape of the eye and transmits light; age-related liquefaction can cause floaters or contribute to retinal traction.

\synonyms
Vitreous Humor; Corpus Vitreum

\medterm{Vitreous Detachment} Separation of the vitreous gel from the retina, usually related to aging and often causing new floaters or flashes. Sudden symptoms can occur with a retinal tear or detachment.

\medterm{Vitreous Hemorrhage} Bleeding into the vitreous cavity that causes sudden painless floaters, cobwebs, haze,
or loss of vision. Common causes include proliferative diabetic retinopathy, retinal
tear or detachment, trauma, and posterior vitreous separation.

\medterm{Vitreous Humor} The clear gel that fills the space between the lens and retina. It helps maintain the shape of the eye and transmits light to
the retina.

\medterm{VKORC1 Polymorphism} An inherited variation in the vitamin K epoxide reductase complex subunit 1 gene that can alter sensitivity to warfarin. Genotype may contribute to dose selection, but clinical factors, interacting medicines, diet, and INR monitoring remain important.
\medterm{VLDL Test} A blood test that measures or estimates very-low-density lipoprotein, a particle carrying triglycerides. The result may help assess cardiovascular risk but is interpreted with the full cholesterol and triglyceride profile.

\medterm{VO2 Max} The greatest amount of oxygen the body can use during intense exercise. It is a measure of aerobic fitness and is influenced by the heart, lungs, blood, muscles, age, health, and training.
\synonyms
Maximal Oxygen Consumption; Maximal Aerobic Capacity

\medterm{Vocal-Cord Dysfunction} Older term for inducible laryngeal obstruction, in which the larynx temporarily narrows during breathing and may cause throat tightness, noisy breathing, cough, or shortness of breath.

\medterm{Vocal Cord Nodules} Benign, usually paired thickened areas of the vocal folds caused by repeated voice strain or overuse. They can cause

hoarseness, vocal fatigue, and a rough or breathy voice.

\medterm{Vocal Cord Paralysis} Inability of one or both vocal folds to move normally, often because of nerve damage. It may cause hoarseness, weak cough, choking, or difficulty protecting the airway; sudden breathing trouble is an emergency.

\synonyms
Vocal Fold Paralysis

\medterm{Vocal Cord Paresis} Weakness or partial paralysis of one or both vocal cords, often from injury to the nerves that control them. It can cause a weak or breathy voice, voice fatigue, coughing during eating, or aspiration.

\medterm{Vocal Cord Polyp} A usually noncancerous growth on a vocal fold, often after heavy or sudden voice use. It can cause persistent hoarseness, a rough voice, or voice fatigue and is assessed by examining the larynx.

\medterm{Vocal Cords} Vocal cords, or vocal folds, are bands of tissue in the larynx that vibrate as air passes between them to produce the voice. Swelling, nodules, or injury can cause hoarseness.

\medterm{Vocal-Fold Nodules} Alternate name for vocal cord nodules, benign thickened areas of the vocal folds that can cause persistent hoarseness.

\medterm{Vocal Fold Paralysis} Alternate name; see \textit{Vocal cord paralysis}.

\medterm{Vocal-Fold Polyps} Alternate name for vocal cord polyps, usually noncancerous growths on the vocal folds that can cause persistent hoarseness or voice fatigue.

\medterm{Vocal Folds} Two folds inside the voice box that vibrate as air passes between them to make sound. They also close during swallowing to protect the lungs from food and liquid.

\medterm{Vocal Hemorrhage} Vocal hemorrhage is bleeding into a vocal fold, often after sudden or forceful voice use. It can cause abrupt hoarseness.

\medterm{Vocal Resonance} The sound of a person’s voice heard through the chest with a stethoscope while they speak. Changes in the sound can provide clues to pneumonia, fluid around the lung, a collapsed lung, or other lung disease.

\medterm{Vogt-Koyanagi-Harada Disease} A rare, multisystem, cell-mediated autoimmune disorder directed against melanocyte-associated antigens (such as tyrosinase), occurring predominantly in pigmented ethnic populations. Characterized by chronic bilateral granulomatous panuveitis with exudative retinal detachments, accompanied by sensorineural hearing loss, tinnitus, vitiligo, poliosis (whitening of eyelashes and hair), and aseptic meningitis.

\synonyms
VKH Syndrome; Uveomeningoencephalitic Syndrome

\medterm{Vogt-Koyanagi-Harada Syndrome} A multisystem autoimmune inflammatory disorder mediated by T-cell responses against melanocyte-associated antigens, characterized by bilateral chronic granulomatous panuveitis, exudative retinal detachments, meningismus (CSF pleocytosis), sensorineural dysacusis, poliosis, and patchy vitiligo.
\synonyms
VKH Syndrome; Uveomeningoencephalitic Syndrome

\medterm{Vohwinkel Syndrome} A rare inherited palmoplantar keratoderma causing thick skin on the palms and soles, constricting bands around fingers or toes, and sometimes hearing loss or other features. Severe constriction can threaten circulation and requires specialist care.

\medterm{Voice} Sound produced by vibration of the vocal folds as air passes through the larynx, then shaped by the throat, mouth, and nose for speech.

\medterm{Voice And Speech} The products of complex coordinated activity of the respiratory, laryngeal, and articulatory systems. Communication produced by coordinated breathing, voice-box vibration, and movements of the mouth, tongue, and other speech structures.

\medterm{Voice Box} The larynx, an organ in the neck containing the vocal folds and forming part of the airway. It produces voice and helps protect the lower respiratory tract during swallowing.

\medterm{Voice Disorders} Conditions affecting voice quality, pitch, loudness, or ease of speaking because of problems involving the vocal folds, airway, nerves, or voice use.

\synonyms
Dysphonia; Vocal Cord Dysfunction

\medterm{Void} To empty the urinary bladder by urinating; in other medical contexts, void may mean an empty space or absence of a substance.

\medterm{Voiding Cystourethrogram} An X-ray test that fills the bladder with contrast through a catheter and records the bladder and urethra while the person urinates. It can show urine flowing backward toward the kidneys, blockage, leakage, or structural problems, especially in children with selected urinary infections.

\medterm{Volar} Relating to the palm of the hand or sole of the foot, especially the surface opposite the dorsum.

\medterm{Volatile Anesthetics} Inhaled medicines used to produce and maintain unconsciousness during some operations. They can lower blood pressure and breathing and rarely cause a dangerous rise in body temperature in susceptible people.


\synonyms
Inhalational Anesthetics; Inhalation General Anesthetics

\medterm{Volkmann Contracture} A permanent, severe ischemic contracture of the forearm flexor musculature leading to a claw-like hand deformity, resulting from untreated acute compartment syndrome of the forearm following supracondylar humeral fractures, crush injuries, or tight casting.
\synonyms
Volkmann's Ischemic Contracture; Ischemic Myodystrophy

\synonyms Volkmann contracture; Volkmann's ischaemic contracture; Volkmann claw deformity.

\medterm{Volume Of Distribution} A calculated measure of how widely a medicine appears to spread through the body compared with the amount left in the blood. It helps guide dosing but does not represent a real space inside the body.

\medterm{Volume Overload} Excess total body sodium and water, causing peripheral or pulmonary oedema; common causes include heart failure, kidney failure, liver disease, and excessive intravenous fluid, and treatment depends on the cause and severity.

\medterm{Voluntary} Under conscious control or performed by choice, such as voluntary movement. Voluntary action is distinguished from reflex, autonomic, or involuntary activity.

\medterm{Voluntary Psychiatric Admission} Admission to a hospital or psychiatric facility with the person’s agreement for assessment, treatment, or safety support. Legal rules about consent, leaving, capacity, and conversion to involuntary status vary by jurisdiction.

\medterm{Volutrauma} Lung injury caused by overstretching the lungs, usually when a breathing machine delivers breaths that are too large. Careful ventilator settings can reduce this risk.

\medterm{Volvulus} Twisting of a loop of intestine that blocks passage through the bowel and may cut off its blood supply. It can cause severe pain, vomiting, swelling, and tissue death and can be life-threatening.

\medterm{Vomit} To eject stomach or upper gastrointestinal contents through the mouth; the expelled material may also be called vomitus. Persistent vomiting can cause dehydration and electrolyte abnormalities.

\medterm{Vomiting} Forceful emptying of stomach contents through the mouth. It may result from infection, digestive disease, medicines, pregnancy, migraine, brain disorders, or poisoning. Repeated vomiting can cause dehydration, disturbed body salts, or inhalation of vomit.

\medterm{Vomiting Blood} The passage of blood or coffee-ground material in vomit, usually indicating bleeding in the upper digestive tract. It can be a serious symptom.

\medterm{Vomitus} Material expelled from the stomach or upper gastrointestinal tract during vomiting. Its color and contents can provide clues but do not by themselves establish a diagnosis.

\medterm{Von Economo Encephalitis} An atypical form of epidemic encephalitis (encephalitis lethargica) characterized by pronounced lethargy, hypersomnolence, ophthalmoplegia, and a frequent late sequela of progressive post-encephalitic parkinsonism, pathologically marked by midbrain and basal ganglia inflammation.
\synonyms
Encephalitis Lethargica; Sleeping Sickness (Von Economo Type)

\medterm{Von Gierke Disease} Glycogen storage disease type I, an inherited deficiency of glucose-6-phosphatase or its transport system that causes fasting hypoglycemia, lactic acidosis, enlarged liver, and excess uric acid or lipids.

\medterm{Von Graefe Sign} A delay of the upper eyelid as the eyes look downward, leaving a strip of white visible above the iris. It can occur with an overactive thyroid and related eye disease.
\synonyms
Lid Lag Sign; Von Graefe's Lid Lag

\medterm{Von Hippel-Lindau Disease} An inherited condition, usually caused by a change in the \textit{VHL} gene, that increases the risk of tumors and cysts in the brain, spinal cord, eyes, kidneys, adrenal glands, and pancreas. These growths can affect vision, blood pressure, or kidney function.

\medterm{Von Hippel-Lindau Syndrome} Von Hippel-Lindau syndrome is an inherited disorder caused by changes in the VHL gene. It increases the risk of tumors and cysts in organs such as the brain, retina, kidneys, pancreas, and adrenal glands.

\medterm{Von Meyenburg Complex} A benign congenital biliary ductal plate malformation (biliary microhamartoma) consisting of small, non-progressive clusters of dilated, cystic bile ducts embedded in a dense collagenous stroma, usually discovered incidentally on liver imaging or biopsy.
\synonyms
Biliary Microhamartoma; Ductal Plate Malformation

\medterm{Von Recklinghausen Disease} Von Recklinghausen disease is another name for neurofibromatosis type 1, an inherited condition causing multiple café-au-lait spots, neurofibromas, and increased risk of certain nerve, bone, eye, and other problems.

\medterm{Von Recklinghausen’S Disease} Alternate name; see \textit{Neurofibromatosis type 1}.

\medterm{Von Recklinghausen’S Neurofibromatosis} Alternate name; see \textit{Neurofibromatosis type 1}.

\medterm{Von Recklinghausen’S Phakomatosis} Alternate name; see \textit{Neurofibromatosis type 1}.

\medterm{Von Rosen Splint} A malleable, lightweight orthopedic abduction splint applied in neonates with developmental dysplasia of the hip (DDH) to maintain the femoral head securely centered within the acetabulum in flexion and abduction.
\synonyms
Von Rosen Frame; DDH Abduction Splint

\medterm{Von Willebrand Disease} A bleeding disorder caused by too little or poorly working von Willebrand factor, a blood protein that helps clots form. It is often autosomal dominant, although severe type 3 and some other forms are autosomal recessive. It commonly causes nosebleeds, easy bruising, heavy menstrual bleeding, or prolonged bleeding after procedures.

\medterm{Voriconazole} An antifungal medicine used for selected serious fungal infections. It stops fungi making an essential part of their outer structure; vision changes, liver problems, medicine interactions, and skin sensitivity require monitoring.

\medterm{VOR Reflex} A reflex that moves the eyes opposite to head movement so vision stays steady while the head turns. Problems with this reflex can cause blurred or bouncing vision, dizziness, or imbalance.

\medterm{Voyeurism} A recurrent sexual interest in observing an unsuspecting person who is naked, disrobing, or engaging in sexual activity. The interest alone is not a mental disorder; voyeuristic disorder is the diagnostic term when required criteria are met.

\medterm{Voyeuristic Disorder} A paraphilic disorder involving recurrent, intense sexual arousal from observing an unsuspecting person who is naked, disrobing, or engaging in sexual activity. Diagnostic criteria include acting on the urges or clinically significant distress or impairment, with the required duration criteria.

\medterm{VPCs (Ventricular Premature Contractions)} Alternate name; see \textit{Premature ventricular contractions (PVCs)}.

\medterm{VRE (Vancomycin-Resistant Enterococci)} Enterococcus bacteria that are no longer killed by vancomycin, an antibiotic. They may live harmlessly in the body or cause urinary, wound, bloodstream, or other infections that require infection-control measures and suitable treatment.

\medterm{VSD (Ventricular Septal Defect)} A hole in the wall between the heart’s lower chambers. Small holes may close or cause no symptoms, while larger holes can make the heart and lungs work too hard and may need treatment.

\synonyms
Ventricular Septal Defect

\medterm{Vulnerary} A substance traditionally believed to help heal wounds. Such products may lack good evidence and can cause allergy or contamination.

\medterm{Vulva} The external female genital structures, including the mons pubis, labia, clitoris, and vestibule. It protects the openings of the urinary and reproductive tracts and contains specialized sensory and glandular tissue.

\synonyms
Female External Genitalia; Pudendum

\medterm{Vulval Intraepithelial Neoplasia} Vulval intraepithelial neoplasia is abnormal cell growth limited to the surface layers of the vulvar skin. Some types are linked to HPV, and untreated lesions may progress to invasive cancer.

\medterm{Vulvar} Relating to the vulva, the external female genital structures including the labia, clitoris, vestibule, and associated glands.

\medterm{Vulvar Cancer} Cancer arising in the external female genitalia. It may appear as a persistent lump, ulcer, wart-like lesion, itching, pain, or bleeding.

\medterm{Vulvar Carcinoma} Cancer of the external female genital area, most often arising from surface skin cells. It may appear as a persistent lump, sore, wart-like growth, itching, pain, or bleeding.

\synonyms
Vulvar Cancer; Squamous Cell Carcinoma of Vulva

\medterm{Vulvar Inclusion Cysts} Small sacs under the skin of the vulva that contain tissue trapped from the surface lining. They may develop after childbirth tears, surgery, or other injury, or occur without an obvious cause. Most cause no symptoms, but they may become swollen, painful, infected, or uncomfortable during sexual activity. Symptomatic cysts may need removal.

\medterm{Vulvar Intraepithelial Neoplasia} Abnormal cells limited to the surface layer of the vulva. Some forms can progress to invasive cancer, especially without treatment. It may cause itching, burning, sores, warty growths, or a color change, but some people have no symptoms.

\medterm{Vulvar Melanoma} A rare cancer of the pigment-producing cells in the vulva. It may appear as a dark or skin-colored lump, patch, sore, or bleeding area; diagnosis requires a biopsy, and treatment depends on size and spread.

\medterm{Vulvar Squamous Cell Carcinoma} The most common invasive vulvar malignancy, arising through HPV-related or HPV-independent pathways. It usually presents as a persistent mass, ulcer, plaque, or wart-like lesion and requires biopsy for diagnosis.

\medterm{Vulvectomy} Surgical removal of part or all of the vulva, performed for selected cancers, precancerous lesions, or severe persistent disease.

\medterm{Vulvitis} Inflammation or irritation of the external female genital area. It can cause itching, burning, redness, swelling, soreness, or discharge and may result from infection, allergy, skin disease, hormones, or friction.

\medterm{Vulvodynia} Vulvodynia is chronic vulvar pain without an identifiable specific cause, often described as burning or rawness and triggered by touch, pressure, or occurring spontaneously.

\synonyms
Chronic Vulvar Pain

\medterm{Vulvovaginal Candidiasis} A yeast infection of the vulva and vagina, usually caused by Candida. It commonly causes intense itching, burning, soreness, redness, and thick discharge; recurrent, severe, unusual, or treatment-resistant symptoms need testing to confirm the cause.

\medterm{Vulvovaginitis} Simultaneous inflammation or infection of the vulva and vagina, commonly presenting with pruritus, burning, dysuria, dyspareunia, and abnormal vaginal discharge, resulting from infectious etiologies (such as Candida albicans, Bacterial vaginosis, Trichomonas vaginalis) or non-infectious causes (atrophic vaginitis, contact irritants).
\synonyms
Vaginitis and Vulvitis; Vulvovaginal Inflammation


\medterm{Waardenburg Syndrome} A group of inherited conditions that may cause hearing loss and unusual coloring of the eyes, hair, or skin. Some forms also affect the arms or the nerves controlling the bowel; findings vary by type.


\medterm{WAGR Syndrome} A rare genetic condition that may cause an increased risk of Wilms kidney tumor, missing or poorly formed irises, urinary or reproductive differences, and developmental delay. Children need regular specialist monitoring, especially for kidney tumors.

\medterm{Waist Circumference} The distance around the abdomen, measured at a standard level. A larger measurement can indicate more abdominal fat and a higher risk of diabetes and heart disease, but it must be interpreted with other health information.

\medterm{Waist-Hip Ratio} A measurement comparing waist size with hip size. It provides an estimate of where body fat is stored and may add information about diabetes and cardiovascular risk, but it does not diagnose disease by itself.

\medterm{Wait List} An official list of people waiting for an organ transplant. Placement and priority depend on medical urgency, expected benefit, compatibility, waiting time, and the rules of the transplant system.

\medterm{Wakefulness} The state of being conscious and able to respond to the surroundings. Sleep, medicines, illness, brain injury, and metabolic problems can change the level of wakefulness.

\medterm{Waldenstrom Macroglobulinemia} A rare blood cancer in which abnormal immune cells build up in the bone marrow and make too much of an antibody called IgM. It may cause anemia, bleeding, vision or nerve problems, enlarged lymph nodes, or no symptoms at first.

\synonyms
Lymphoplasmacytic Lymphoma

\medterm{Waldenstrom'S Macroglobulinemia} A lymphoplasmacytic cancer of the bone marrow that produces excess IgM and may cause
anemia, enlarged lymph nodes, or hyperviscosity symptoms.

\medterm{WaldenströM'S Macroglobulinemia} A slow-growing blood cancer in which abnormal immune cells make too much of an antibody called IgM. The blood can become unusually thick, causing blurred vision, headaches, bleeding, confusion, or numbness.

\medterm{Walker-Warburg Syndrome} A rare inherited disorder causing severe muscle weakness together with major abnormalities of the brain and eyes from birth. Babies may have low muscle tone, seizures, poor feeding, enlarged fluid spaces in the brain, and abnormal eye development; the condition is usually life-threatening in infancy.

\synonyms
WWS; Muscle-Eye-Brain Disease (Severe Type); HARD +/- E Syndrome

\medterm{Walking Abnormalities} Walking abnormalities include changes in speed, balance, stride, posture, or coordination caused by pain, weakness, sensory loss, joint disease, or neurologic illness.

\medterm{Walking Speed} The rate at which a person walks, often measured over a defined distance. It can provide information about mobility and function, but no single speed determines disability in every person.

\medterm{Wall Stress} The strain placed on the wall of a heart chamber by the pressure inside it and the chamber’s size. High wall stress makes pumping harder and may contribute to heart enlargement or failure.

\medterm{Wallerian Degeneration} Breakdown of the part of a nerve beyond a severe injury after it loses contact with the nerve cell. The nerve may regrow slowly if the surrounding pathway remains intact, but recovery depends on the injury.


\medterm{Walleyed} Walleyed describes eyes that appear misaligned or divergent; clinically, the condition may be strabismus and requires eye assessment when persistent.

\medterm{Wandering Spleen} A spleen that can move unusually because its supporting bands are loose or missing. It may cause no symptoms or lead to abdominal pain, enlargement, twisting, loss of blood supply, or sudden severe illness.


\medterm{Warfarin} A blood-thinning medicine that reduces production of vitamin-K-dependent clotting proteins. It prevents and treats clots but requires regular blood tests because food, illnesses, and many medicines alter its effect.

\medterm{Warfarin Syndrome} A pattern of fetal abnormalities caused by warfarin exposure during pregnancy, including nasal and skeletal abnormalities, eye or neurologic defects, and fetal or neonatal bleeding.

\medterm{Warm Autoimmune Hemolytic Anemia} An illness in which the immune system makes antibodies that cause red blood cells to be destroyed too early, mainly in the spleen. It can cause tiredness, pale or yellow skin, dark urine, shortness of breath, or a fast heartbeat and requires medical evaluation.

\synonyms
wAIHA; Warm antibody autoimmune hemolytic anemia

\medterm{Warm Ischemia} Lack of blood and oxygen reaching tissue while it remains at body temperature. It can quickly cause cell damage and is important during some injuries, operations, and organ transplantation.

\medterm{Wart} A usually harmless rough skin growth caused by human papillomavirus. Warts may spread by contact and often disappear without treatment, but painful, bleeding, changing, or uncertain growths should be examined.

\synonyms
Verruca

\medterm{Wart Remover Poisoning} Toxic exposure to a wart-removal product, often involving salicylic acid or another keratolytic. Severity depends on the substance, amount, route, symptoms, and duration of exposure.

\medterm{Warthin Tumor} A usually harmless tumor of a salivary gland, most often the parotid near the ear. It is strongly associated with tobacco use and may occur on both sides; evaluation distinguishes it from other growths.

\medterm{Warty Dyskeratoma} A usually harmless skin growth, often on the head or neck, with a small central dip filled with keratin. It may look like skin cancer, so examination and sometimes a biopsy are needed.



\medterm{Wasting} Abnormal loss of body weight, muscle, or other tissue, often associated with chronic disease, inadequate nutrition, inflammation, or cancer. Severe wasting may impair strength, immunity, and recovery.

\medterm{Watchful Waiting} Planned observation without immediate active treatment, with clinical review and treatment
if symptoms develop or the condition progresses. It may be used for selected cancers and other slowly progressive
conditions.

\medterm{Water Birth (Underwater Delivery)} Delivery of a baby while the birthing person remains in a tub or pool of warm water.
It may reduce perceived labor pain for some people, but evidence about delivery-related
benefits and risks is limited. Eligibility, monitoring, and facility policies determine whether
it is offered.

\synonyms
Underwater Delivery

\medterm{Water Blister} A fluid-filled blister, usually containing clear serum, caused by friction, pressure, burns, or some skin disorders. The blister should be protected from further trauma and assessed if infected or widespread.

\medterm{Water Brash} A sudden rise of watery or salty-tasting saliva into the mouth, usually because stomach acid has refluxed into the esophagus. It may occur with heartburn, sour taste, throat irritation, or regurgitation.

\medterm{Water Channel} A protein pore in a cell membrane that lets water cross quickly. These pores help the kidneys and other tissues move water and maintain the body’s fluid balance.

\medterm{Water Consumption} Intake of water and other fluids to replace normal losses and maintain hydration. Requirements vary with body size, diet, activity, climate, pregnancy, illness, and kidney or heart function; thirst and urine concentration can help guide intake, while excessive intake can cause dangerous hyponatremia.


\medterm{Water Intoxication} A dangerous fall in blood sodium caused by drinking or retaining more water than the body can remove. Severe cases can cause headache, confusion, seizures, or coma and need urgent treatment.

\medterm{Water On The Brain} A nontechnical term usually referring to hydrocephalus, an abnormal accumulation or impaired circulation of cerebrospinal fluid that enlarges the ventricles and may increase intracranial pressure.

\medterm{Water On The Knee} Alternate index label; see \textit{Swollen knee}.


\medterm{Water Retention} Excess accumulation of fluid in the body's tissues or cavities, commonly causing swelling and sometimes resulting from heart, kidney, liver, venous, hormonal, or medication-related disorders.



\medterm{Water-Borne Disease} An illness acquired through water contaminated with infectious organisms or harmful chemicals. Examples include cholera, typhoid fever, and hepatitis A; prevention includes safe drinking water, sanitation, hand hygiene, and appropriate treatment of contaminated water.

\medterm{Water-Deprivation Test} A supervised test that briefly limits drinking while urine concentration and blood results are monitored. It helps identify why a person produces excessive amounts of dilute urine and must be done under medical supervision because dehydration and high sodium can occur.

\medterm{Water-Hammer Pulse (Corrigan Pulse)} A pulse that feels forceful at first and then suddenly collapses. It is often associated with severe leakage of the aortic valve or other causes of a wide pulse pressure.

\synonyms
Corrigan Pulse

\medterm{Water-Soluble Vitamin} A vitamin that dissolves in water and is stored only in limited amounts. The group includes vitamin C and the B vitamins, so regular intake is generally needed.


\medterm{Waterborne Bacterial Disease} An infection acquired from water contaminated with bacteria, such as cholera, typhoid fever, or leptospirosis; prevention depends on safe water, sanitation, and food hygiene.

\medterm{Waterhouse-Friderichsen Syndrome} Waterhouse-Friderichsen syndrome is adrenal hemorrhage and acute adrenal failure associated with overwhelming sepsis, classically meningococcal infection, causing shock and coagulopathy.

\medterm{Watery Eyes} Excessive tearing caused by irritation, allergy, infection, dry-eye reflex tearing, or
blockage of the tear-drainage system. Persistent unilateral tearing or pain requires
ocular assessment.

\medterm{Wax} A water-resistant substance produced by ceruminous glands in the ear canal; earwax protects the canal and usually clears naturally.

\medterm{Wax Poisoning} Toxic exposure to a wax product, usually causing little harm after a small accidental ingestion but potentially causing choking, aspiration, irritation, or toxicity when mixed with other chemicals.

\medterm{Waxy Casts} Large, smooth, pale tube-shaped particles seen under a microscope in urine. They form in damaged kidney tubules and can suggest severe or long-standing kidney disease, especially when other urine or blood tests are abnormal.

\medterm{WBC Count} A measurement of the number of white blood cells in a volume of blood, used with the differential count to evaluate infection, inflammation, immune disorders, marrow disease, or treatment effects.

\medterm{Weakness} Reduced ability to generate force, which may be generalized or limited to one body region. Causes include disorders of muscle, nerve, neuromuscular junction, brain, or spinal cord, as well as metabolic illness, infection, medication effects, and deconditioning; sudden one-sided weakness or weakness with breathing, swallowing, or speech difficulty requires emergency assessment.

\medterm{Weakness, Generalized} Reduced strength affecting multiple body regions, caused by systemic illness, metabolic or electrolyte disturbance, medication, infection, muscle disease, nerve disease, neuromuscular-junction disease, or brain or spinal-cord disease. Sudden or progressive weakness, especially with speech, facial, breathing, or walking difficulty, requires urgent assessment.

\medterm{Weaning} Gradual reduction and eventual discontinuation of breastfeeding, bottle-feeding, or a medicine or other treatment, according to the relevant clinical context.

\medterm{Wear Debris} Tiny particles produced when an artificial joint or other implant wears down. The particles can trigger inflammation, damage nearby bone, loosen the implant, and eventually cause it to fail.

\medterm{Weaver Syndrome} A rare genetic condition causing rapid growth before and after birth, a large head, distinctive facial features, low muscle tone, and delayed development or learning. Bone age may be advanced. Most cases involve a change in the \textit{EZH2} gene; care is tailored to the child’s needs.

\medterm{Weaver'S Bottom} Inflammation of the ischial bursa over the sitting bone, causing pain in the buttock that may worsen with prolonged sitting. It is also called ischial bursitis.


\medterm{Webbed Fingers} Partial or complete joining of adjacent fingers by skin or soft tissue, known as syndactyly; it may be congenital or acquired.

\medterm{Webbed Neck} Extra folds of skin extending from the sides of the neck toward the shoulders. They may be present from birth and can occur with conditions such as Turner syndrome or Noonan syndrome.

\medterm{Webbing Of The Fingers Or Toes} Partial or complete joining of adjacent fingers or toes by skin or other tissue. It is called syndactyly and may be congenital or acquired.

\medterm{Weber} A name used in hearing assessment, especially the Weber test, in which a vibrating tuning fork is placed on the forehead or skull. It helps compare hearing between the two ears but is not a unit of measurement.

\medterm{Weber Test} A quick hearing test using a vibrating tuning fork placed on the middle of the forehead or skull. The examiner asks which ear hears the sound louder; the result helps distinguish blockage-related hearing loss from inner-ear or nerve-related loss.

\medterm{Wedge Excision} Surgical removal of a triangular or wedge-shaped portion of tissue; commonly performed for breast biopsies, cervical lesions, and skin tumors.

\medterm{Wedge Resection} Surgery that removes a small wedge-shaped piece of tissue, commonly from the lung or breast, to diagnose or treat a localized lesion. It removes less tissue than removal of an entire lobe.


\medterm{Weever Fish} A venomous marine fish whose dorsal spines can cause a painful sting, local swelling, and occasionally systemic symptoms.

\medterm{Wegener'S Granulomatosis} Alternate historical term; see \textit{Granulomatosis With Polyangiitis}.

\medterm{Wegeners Granulomatosis (Granulomatosis With Polyangiitis)} Alternate index label for Granulomatosis With Polyangiitis.

\medterm{Weight Gain} Increase in body weight; may be intentional (muscle building, pregnancy) or unintentional (fluid retention, endocrine disorders, medication effects).

\medterm{Weight Gain In Pregnancy} Recommended weight gain depends on weight before pregnancy and whether one or more babies are expected. For one baby, usual ranges are about 12.5--18 kg if underweight, 11.5--16 kg at usual weight, 7--11.5 kg if overweight, and 5--9 kg with obesity; prenatal clinicians monitor the pattern and fetal growth.

\medterm{Weight Gain – Unintentional} Increase in body weight without a planned change in diet or activity, often caused by fluid retention, medicines, endocrine disease, reduced activity, or changes in eating; rapid swelling or breathlessness needs prompt assessment.

\medterm{Weight Loss} A decrease in body weight that may be intentional or unintentional. Intentional loss is usually achieved through sustainable changes in eating, physical activity, and behavior; unintentional or rapid loss can indicate illness, inadequate nutrition, or medication effects and may require medical evaluation.

\medterm{Weight Loss - Unintentional} Unintentional weight loss is reduction in body weight without planned dietary or activity change and may result from malignancy, endocrine disease, infection, gastrointestinal disease, medication, or mood disorder.

\medterm{Weight Loss Medications} Medicines prescribed for selected people with obesity to help reduce weight, usually together with nutrition and activity changes. Choice depends on health conditions, other medicines, side effects, cost, and treatment goals.
\medterm{Weight-Loss Medicines} Alternate index label; see \textit{Weight Loss Medications}.

\medterm{Weight-Reducing Diets} Dietary plans intended to create a sustained energy deficit and reduce excess body weight, ideally individualized with attention to nutrition, health conditions, and long-term adherence.


\medterm{Weil'S Disease} A zoonotic bacterial infection caused by spirochaetes of the genus Leptospira,
transmitted through contact with water or soil contaminated with the urine of infected
animals (rodents, cattle, dogs). See Leptospirosis.

\medterm{Weils Disease (Leptospirosis)} Alternate index label for Leptospirosis.

\medterm{Wellens Syndrome} A specific pattern on a heart tracing that can signal severe narrowing of a major heart artery, even after chest pain stops. It indicates high heart-attack risk and needs urgent specialist assessment; exercise testing can be dangerous.

\synonyms
Wellens Pattern; LAD Coronary T-Wave Syndrome; Wellens Sign

\medterm{Wells Score For Pulmonary Embolism} A clinical score estimating the likelihood of a blood clot in the lungs. It considers leg-clot signs, heart rate, recent surgery or immobility, previous clots, coughing blood, cancer, and alternative diagnoses; the result helps decide whether blood testing or lung imaging is needed.

\synonyms
Wells Criteria for PE; Wells Score (PE)

\medterm{Welt} A raised, often itchy area of skin caused by localized swelling, commonly seen in hives or after an insect bite. A welt may fade within hours, but persistent or severe reactions need evaluation.

\medterm{Wen (Epidermoid Cyst)} A common, usually harmless lump beneath the skin filled with a soft protein called keratin. It may become swollen, painful, or infected and can be removed if necessary.

\synonyms
Epidermoid Cyst

\medterm{Wenckebach Phenomenon (Mobitz Type I AV Block)} A heart-rhythm pattern in which each heartbeat takes slightly longer to pass from the upper to lower chambers until one beat is missed. It may be harmless or reflect heart disease, depending on the setting.

\synonyms
Mobitz Type I AV Block

\medterm{Wens} Alternate index label; see \textit{Wen (Epidermoid Cyst)}.

\medterm{Werdnig-Hoffmann Disease} The most severe early form of inherited spinal muscular atrophy, also called type 1 spinal muscular atrophy. It usually begins before 6 months of age and causes very weak muscles, poor head control, reduced reflexes, a weak cry, feeding difficulty, and breathing problems. Babies generally cannot sit without support. New treatments can improve survival and function, but affected infants need urgent specialist care and breathing and feeding support.

\synonyms
Spinal Muscular Atrophy Type 1 (SMA 1); Severe Infantile Spinal Muscular Atrophy

\medterm{Wermer Syndrome} Alternate index label; see \textit{Multiple endocrine neoplasia, type 1 (MEN 1)}.

\medterm{Wermer'S Syndrome} Alternate index label; see \textit{Multiple endocrine neoplasia, type 1 (MEN 1)}.

\medterm{Werner Syndrome} A rare inherited condition that causes signs of aging earlier than usual, often after puberty. Features may include short stature, early cataracts, thin or tight skin, gray hair, diabetes, weak bones, and increased cancer risk. Regular medical monitoring and treatment of each complication are needed.

\medterm{Werner'S Syndrome} A rare inherited condition that causes signs of aging earlier than usual, often after puberty. It may cause early cataracts, short stature, thin skin, gray hair, diabetes, weak bones, and increased cancer risk.

\medterm{Wernicke Area} A brain region, usually in the left temporal lobe, that helps a person understand spoken and written language. Damage can cause fluent speech that lacks meaning and poor understanding of language.

\medterm{Wernicke Encephalopathy} A sudden brain disorder caused by severe vitamin B1 deficiency. It may cause confusion, difficulty walking, and abnormal eye movements and can lead to permanent damage without prompt treatment.

\medterm{Wernicke'S Encephalopathy} An acute neurologic disorder caused by thiamine deficiency, classically involving confusion,
eye-movement abnormalities, and gait or coordination impairment. It requires prompt
thiamine treatment.

\medterm{Wernicke-Korsakoff Syndrome} A brain disorder caused by severe vitamin B1 deficiency. It may begin with confusion and poor coordination and leave lasting memory problems, difficulty learning, and inaccurate made-up explanations; alcohol misuse, poor nutrition, and prolonged vomiting are common causes.

\medterm{Wernicke’S Area} A brain region important for understanding spoken and written language.

\synonyms
Wernicke area


\medterm{West Nile Encephalitis} Inflammation of the brain caused by neuroinvasive West Nile virus, transmitted mainly by mosquitoes. It may cause
fever, headache, confusion, tremor, weakness, seizures, or coma. There is no specific routine antiviral treatment;
severe illness requires hospital-based supportive care.

\medterm{West Nile Virus} A mosquito-borne flavivirus that can cause asymptomatic infection, a febrile illness, or neuroinvasive
disease such as meningitis, encephalitis, or acute flaccid paralysis. Risk of severe disease
is higher in older adults and people with certain medical conditions.

\medterm{West Nile Virus Infection} West Nile virus infection is a mosquito-borne illness ranging from fever and body aches to neuroinvasive meningitis, encephalitis, weakness, or paralysis.

\medterm{West Syndrome} A serious seizure disorder beginning in infancy, with repeated sudden spasms and an abnormal brain-wave pattern. Development may slow or regress; causes include brain injury, genetic conditions, and metabolic disease, and urgent treatment is needed.

\medterm{Westermark Sign} An area of unusually reduced blood-vessel markings on a chest X-ray, sometimes seen beyond a clot in a lung artery. It may suggest pulmonary embolism but is not reliable enough to diagnose it alone.


\medterm{Westley Croup Score} A clinical score used to estimate how severely croup is obstructing a child's airway. It considers noisy breathing, chest retractions, air movement, bluish coloring, and alertness; a higher score suggests more severe illness and helps guide urgent treatment and observation.

\synonyms
Westley Clinical Scoring System for Croup; Croup Severity Score

\medterm{Wet Age-Related Macular Degeneration} Alternate index label; see \textit{Wet macular degeneration}.

\medterm{Wet AMD} Alternate index label; see \textit{Wet macular degeneration}.

\medterm{Wet Dream} Involuntary ejaculation during sleep, usually occurring in adolescence or adulthood and generally a normal physiological event.

\medterm{Wet Dry Drowning (Drowning)} Alternate index label for Drowning.

\medterm{Wet Gangrene} Tissue death complicated by bacterial infection and tissue swelling, producing moist, discolored, rapidly spreading necrosis. It requires urgent medical assessment.

\medterm{Wet Macular Degeneration} The neovascular form of age-related macular degeneration in which abnormal blood vessels grow beneath or within the retina and leak fluid or blood, causing rapid central-vision loss.

\synonyms
Macular Degeneration, Wet
Wet Age-Related Macular Degeneration
Wet AMD

\medterm{Wharton Jelly} The soft, jelly-like tissue surrounding the blood vessels in the umbilical cord. It cushions the vessels and helps prevent them from being squeezed or bent during pregnancy and labor.

\medterm{Wharton'S Jelly} Gelatinous connective tissue within the umbilical cord that surrounds and supports the umbilical blood vessels and helps protect them from compression.

\medterm{Wheal} A temporary, raised, often itchy area of skin caused by swelling in the skin. It commonly occurs with hives, an allergic reaction, or an insect bite and may move or disappear within hours.

\medterm{Wheal And Flare} A temporary raised, swollen area of skin surrounded by redness. It is usually caused by release of histamine during hives, an allergy, or another immediate skin reaction.

\medterm{Wheat Allergy} An immune reaction to proteins in wheat that may cause hives, stomach symptoms, vomiting, or breathing problems.

\synonyms
Allergy, Wheat
\medterm{Wheeled Stretcher} Mobile patient transport device with wheels; used in hospitals and emergency services for moving patients who cannot walk.

\medterm{Wheeze} A high-pitched whistling sound caused by narrowed airways. It often occurs with asthma, infection, allergy, or chronic lung disease; sudden wheezing after choking or with severe breathlessness needs urgent care.

\medterm{Wheezing} A musical or whistling sound caused by narrowed lower airways, usually louder when breathing out. It may result from asthma, infection, chronic lung disease, allergy, airway swelling, mucus, or an inhaled object.

\medterm{Whiplash} A neck injury caused by the head moving suddenly backward and forward, often in a vehicle crash but also after a fall or sports injury. Pain, stiffness, headache, dizziness, or arm tingling may appear immediately or later. Most cases improve over weeks, but weakness, numbness, trouble walking, or severe pain after major trauma needs urgent assessment.

\synonyms
Pain Whiplash (Whiplash)
Sprain Neck (Whiplash)
Strain Neck (Whiplash)


\medterm{Whiplash Injury} Alternate index label; see \textit{Whiplash}.

\medterm{Whipple Disease} A rare bacterial infection that can damage the small intestine, joints, brain, heart, or other organs. It may cause years of joint pain followed by weight loss, persistent diarrhea, abdominal pain, and poor nutrient absorption.

\medterm{Whipple Procedure} A major operation that removes the head of the pancreas and nearby parts of the small intestine, bile duct, and gallbladder, then reconnects the digestive tract. It is used for selected cancers and other serious diseases in this area.

\medterm{Whipple Triad} Three findings supporting that symptoms are caused by low blood sugar: typical symptoms occur, blood sugar is low at that time, and symptoms improve after blood sugar is raised.

\medterm{Whipple'S Disease} A rare bacterial infection caused by Tropheryma whipplei, damaging intestinal absorption and potentially affecting joints, brain, or heart.

\synonyms
Intestinal Lipodystrophy
Lipodystrophy, Intestinal

\medterm{Whipworm} A parasitic roundworm of the genus Trichuris, usually Trichuris trichiura in humans. Infection occurs by ingesting embryonated eggs and may cause abdominal pain, diarrhea, anemia, or rectal prolapse when severe.

\medterm{Whipworm (Trichuriasis)} Intestinal infection caused by the soil-transmitted helminth Trichuris trichiura
(whipworm), endemic in tropical and subtropical regions with poor sanitation. Humans
acquire infection by ingesting embryonated eggs from contaminated soil, food, or water.

\synonyms
Trichuriasis

\medterm{Whipworm Infection} Intestinal infection with Trichuris trichiura acquired by ingesting eggs from contaminated soil or food; heavy infection can cause abdominal symptoms, anemia, or rectal prolapse.

\medterm{Whirlpool Bath} Hydrotherapy device providing agitated water for therapeutic bathing; used in wound care, physical therapy, and rehabilitation.

\medterm{Whispered Pectoriloquy} An examination finding in which whispered words sound unusually clear through a stethoscope over the lung. It may suggest that air-filled lung has become solid, as in pneumonia, but it is not diagnostic alone.

\medterm{White Blood Cell} A cell of the immune system that helps protect the body from infections, foreign substances, and abnormal cells. White blood cells are made mainly in the bone marrow and include neutrophils, lymphocytes, monocytes, eosinophils, and basophils. Their number or function may change with infection, inflammation, immune disorders, or bone-marrow disease.

\medterm{White Blood Cell Count} A blood test measuring the number of white blood cells. A high or low result can occur with infection, inflammation, medicines, immune disorders, or bone-marrow disease and must be interpreted with the differential count and clinical findings.


\medterm{White Blood Cell Labeling} A scan method in which a person's white blood cells are removed, marked with a small amount of radioactive material, and returned to the body. The cells may collect at areas of infection or inflammation, which can then be seen on images.


\medterm{White Cell} Another name for a white blood cell, an immune cell made mainly in the bone marrow that helps defend the body against infection, foreign substances, and abnormal cells.

\medterm{White Cell Count} A blood test measuring the number of white blood cells and often the proportion of each type. Results help assess infection, inflammation, immune disorders, and bone-marrow problems.

\medterm{White Finger} Episodic pallor or whitening of the fingers caused by vasospasm, commonly associated with Raynaud phenomenon or hand-arm vibration syndrome.

\medterm{White Infarction} Tissue death caused by loss of blood flow in a solid organ such as the heart, kidney, or spleen. The damaged area looks pale because little blood can reach it through alternate vessels.

\synonyms
anemic infarction, occurs in solid organs with.

\medterm{White Leg (Phlegmasia Alba Dolens)} Painful, marked swelling and pallor of a limb caused by extensive deep venous thrombosis, usually involving major venous outflow. It requires urgent assessment.

\synonyms
Phlegmasia Alba Dolens

\medterm{White Matter} Tissue in the brain and spinal cord made mainly of nerve fibers covered by a protective substance called myelin. It carries signals between different brain areas and between the brain and spinal cord.

\medterm{White Matter Of The Brain} The brain tissue composed mainly of myelinated axons and supporting glial cells, connecting different brain regions and the brain with the spinal cord; disease or injury can impair cognition, sensation, coordination, or movement.

\medterm{White Of The Eye} The sclera, the tough, opaque outer coat of the eye surrounding the cornea and helping maintain the shape of the globe.

\medterm{White Spot Disease} Alternate index label; see \textit{Lichen sclerosus}.

\medterm{White Tongue} A white coating or patches on the tongue caused by debris, dry mouth, irritation, infection, inflammation, or other oral conditions.

\medterm{White-Coat Hypertension} Blood pressure that is high in a medical setting but normal outside it. Home or ambulatory readings help confirm the pattern, which may still predict higher future hypertension risk.

\synonyms
White-coat effect

\medterm{Whiteheads} Whiteheads are closed comedones formed when oil and dead skin block a hair follicle, producing small pale or skin-colored acne bumps.

\medterm{Whitlow} An infection or inflammation of the tissues around a finger or toenail, often involving the nail fold; it may be bacterial, fungal, or herpetic.

\medterm{Whitlow (Paronychia)} Acute or chronic infection and inflammation of the tissues around a fingernail or toenail, usually caused by bacteria, fungi, or repeated moisture and trauma.

\synonyms
Paronychia

\medterm{Whole Breast Radiation Therapy} Whole breast radiation therapy treats remaining breast tissue after breast-conserving surgery to reduce local recurrence risk, with planned dosing and monitoring for skin and tissue effects.

\medterm{Whole Exome Sequencing} A genetic test that reads most protein-coding sections of a person’s DNA. It can identify changes causing some inherited disorders, but it may miss changes outside coding regions, structural changes, or causes not yet understood; results require clinical interpretation.

\medterm{Whole Genome Sequencing} Laboratory analysis of nearly all of a person's nuclear DNA, including coding and noncoding regions, and sometimes mitochondrial DNA. It can identify many types of genetic variants but may produce uncertain or incidental findings and does not detect every cause of disease; results require appropriate clinical interpretation.

\medterm{Whole-Body Irradiation} Radiation treatment directed at the entire body, most often given before a stem-cell transplant to suppress the immune system and make room for donated cells. It can affect the bone marrow, fertility, and other organs.

\medterm{Whole-Body Scan} An imaging examination that surveys much or all of the body, using a technique such as PET/CT, CT, MRI, or a nuclear medicine scan. It may be used for selected cancer staging, infection localization, or evaluation of systemic disease; the coverage and usefulness depend on the test and clinical question.

\medterm{Whole-Bowel Irrigation} A hospital treatment that flushes the intestines with a special polyethylene-glycol drink. It is used only for selected poisonings or swallowed objects after clinicians assess the airway and risks.

\medterm{Whole-Brain Radiation Therapy} Radiation treatment directed at the entire brain, used in selected people with multiple brain metastases or cancer cells around the brain. It may also be used to reduce brain-spread risk in selected cancers; side effects depend on dose and treatment plan.

\medterm{Whole-Genome Sequencing} Sequencing of nearly the complete nuclear genome, including coding and noncoding
regions, together with mitochondrial DNA when covered by the assay. It can detect a
broad range of variants but generates substantial interpretation and incidental-finding
challenges.

\medterm{Whooping Cough} A highly contagious respiratory infection caused by \textit{Bordetella pertussis}. It causes prolonged bouts of rapid coughing that may be followed by a whooping sound or vomiting; infants may have pauses in breathing without the typical whoop. Also called pertussis.

\medterm{Whooping-Cough} Alternate index label; see \textit{Whooping Cough}.

\medterm{Widal Reaction} An older blood test for antibodies associated with typhoid or paratyphoid fever. It is not very reliable because results may reflect past infection or vaccination, so culture or molecular testing is usually preferred.


\medterm{Widely Spaced Teeth} Teeth separated by unusually large gaps, also called diastema when the gap is localized. Causes include normal development, missing teeth, tooth-size differences, habits, or jaw and gum conditions.

\medterm{Williams Syndrome} A genetic condition caused by loss of a small section of chromosome 7, including a gene important for elastic tissues. It may cause characteristic facial features, heart-vessel narrowing, high blood calcium, developmental differences, learning difficulties, and connective-tissue problems. Regular heart, calcium, hearing, and developmental checks are important.

\medterm{Wilm'S Tumour} Source spelling variant of \textit{Wilms tumor}, a malignant embryonal kidney tumor occurring mainly in young children.

\medterm{Wilms Tumor} A malignant embryonal tumor of the kidney in children. It may present as an abdominal mass, pain,
hematuria, or an incidental finding and may be associated with syndromes such as WAGR or
Beckwith-Wiedemann syndrome. Treatment and prognosis depend on stage and tumor biology.

\synonyms
Cancer, Wilms' Tumor

\medterm{Wilms' Tumor (Nephroblastoma)} An embryonal malignant tumor of the kidney that usually occurs in young children. It may be
associated with congenital syndromes, and treatment and prognosis depend on tumor stage,
histology, and molecular features.

\synonyms
Nephroblastoma

\medterm{Wilson Disease} An inherited disorder in which excess copper builds up in the liver, brain, eyes, and other organs. It can cause liver disease, tremor or other movement problems, mood or behavior changes, and colored rings around the cornea. Lifelong treatment removes or limits copper and prevents further buildup.

\medterm{Wilson'S Disease} An inherited condition in which copper builds up in the liver, brain, eyes, and other organs. It may cause liver disease, movement or mood changes, and colored rings around the cornea, but lifelong treatment can prevent further damage.

\medterm{Windowing} Adjusting the brightness and contrast of a CT image so particular tissues, such as bone, lung, or soft tissue, are easier to see.

\medterm{Windpipe} The tube carrying air from the voice box to the two main airways leading to the lungs. It is supported by firm rings that help keep it open.


\medterm{Winter Vomiting Disease} An acute gastroenteritis, usually caused by norovirus, characterized by sudden vomiting, diarrhea, and abdominal cramps.

\medterm{Wireless Motility Capsule} A swallowed capsule that records pressure, acidity, temperature, and movement as it travels through the digestive tract. The measurements estimate how long food takes to pass through the stomach and intestines and may help investigate unexplained constipation, diarrhea, nausea, or bloating.

\medterm{Wisdom Teeth} Wisdom teeth are the third molars that commonly emerge in late adolescence or adulthood. Impaction, infection, decay, cysts, gum disease, or crowding can cause pain and swelling; dental examination and imaging determine whether monitoring or removal is appropriate.

\medterm{Wisdom Teeth, Impacted} Alternate index label; see \textit{Impacted wisdom teeth}.

\medterm{Wisdom Tooth (Third Molar)} The last molar in each dental quadrant, usually erupting in late adolescence or early adulthood. Impaction may cause pain, infection, or damage to adjacent teeth.

\synonyms
Third Molar

\medterm{Wiskott-Aldrich Syndrome} A rare inherited immune disorder that mainly affects boys and causes eczema, repeated infections, and a low number of unusually small platelets. It can also cause bleeding and autoimmune problems and requires specialist care.

\medterm{Witch'S Milk} A temporary milky breast discharge in a newborn caused by maternal hormones crossing the placenta; it usually resolves without squeezing or treatment, which can otherwise cause infection.

\medterm{Witch'S Milk (Neonatal Galactorrhea)} A temporary milk-like breast secretion in a newborn caused by exposure to maternal hormones. It usually resolves without treatment; the breasts should not be squeezed.

\synonyms
Neonatal Galactorrhea

\medterm{Withdrawal} Physical or psychological symptoms that occur when a person who has developed dependence reduces or stops using a medicine or psychoactive substance. Symptoms vary by substance and may include anxiety, agitation, sweating, tremor, nausea, pain, insomnia, cravings, or seizures; some forms require urgent medical care.

\synonyms
Withdrawal syndrome

\medterm{Withdrawal Reflex} An automatic response in which a painful stimulus makes a limb pull away from danger. The opposite limb may straighten to help maintain balance.

\medterm{Withdrawal Symptoms} Physical or psychological symptoms that occur when a person stops or reduces a substance or medicine
after the body has adapted to it. The symptoms depend on the substance, duration of use,
and the person’s health.

\medterm{Withdrawal Syndrome} A predictable set of signs and symptoms that occur when a drug is discontinued or the
dose is reduced after chronic use, reflecting the physiological rebound from drug
adaptation.

\medterm{Wolf-Hirschhorn Syndrome} A rare genetic disorder caused by loss of part of chromosome 4. It can cause poor growth, a small head, distinctive facial features, developmental and intellectual disability, seizures, and heart or other organ abnormalities; lifelong specialist support is usually needed.

\synonyms
4p- Syndrome; 4p Deletion Syndrome

\medterm{Wolff Parkinson White Syndrome (Wolff-Parkinson-White Syndrome)} Alternate index label for Wolff-Parkinson-White Syndrome.

\medterm{Wolff-Parkinson-White Accessory Pathway-Related Tachycardia} Alternate index label; see \textit{Wolff-Parkinson-White syndrome}.

\medterm{Wolff-Parkinson-White Syndrome} A heart condition present from birth in which an extra electrical pathway can make the heart beat very fast. It may cause sudden palpitations, dizziness, fainting, or rarely cardiac arrest and can be treated by medicines or a procedure.

\synonyms
Wolff Parkinson White Syndrome (Wolff-Parkinson-White Syndrome)

\medterm{Wolffian Ducts} Temporary tubes in the developing embryo that form parts of the male reproductive system, including the sperm-transport ducts. They usually shrink during typical female development, although small remnants may remain.

\medterm{Wolfram Syndrome} A rare inherited disorder that commonly causes diabetes requiring insulin and progressive loss of vision from optic-nerve damage. It may also cause excessive thirst and urination, hearing loss, bladder problems, balance difficulty, or other nerve-system problems. Treatment monitors and supports each affected function.

\medterm{Wolinella} Genus of anaerobic, Gram-negative bacteria found in the oral cavity and gastrointestinal tract; associated with periodontal disease.


\medterm{Wolman Disease} A severe inherited disorder in which a missing enzyme causes fats to build up in organs. It usually begins in infancy with vomiting, diarrhea, poor growth, an enlarged liver and spleen, and adrenal-gland problems. Without treatment it can be fatal early; enzyme replacement and specialist care are required.

\medterm{Wolman'S Disease} A severe lysosomal acid-lipase deficiency presenting in infancy with hepatosplenomegaly,
adrenal calcification, malabsorption, vomiting, diarrhea, and failure to thrive. It is a
severe form of lysosomal acid-lipase deficiency and can be fatal early without
treatment.

\medterm{Womb} Common name for the uterus; the muscular organ where a fetus develops during pregnancy.
\synonyms
Uterus

\medterm{Womb Growths (Uterine Fibroids)} Alternate index label for Uterine Fibroids.

\medterm{Women'S Health} Medical specialty focusing on conditions unique to or more prevalent in females, including reproductive health, pregnancy, menopause, and gender-specific disease patterns.

\synonyms
Womens Health (Women's Health)


\medterm{Womens Health (Women'S Health)} Alternate index label for Women's Health.

\medterm{Wood Lamp Examination} Examination of the skin in a dark room using special ultraviolet light. The light can make some color changes and infections easier to see, but the finding usually needs confirmation with the examination and other tests.


\medterm{Wood Spirit} An old name for methanol, a toxic alcohol whose ingestion can cause metabolic acidosis, blindness, and death.

\medterm{Wood'S Light} A special ultraviolet light used in a darkened room to make some skin colors, scales, and infections easier to see. It can support an examination but usually cannot identify the cause by itself.

\medterm{Woods Corkscrew Maneuver} An emergency procedure used when a baby’s shoulders become stuck during birth. A trained clinician rotates the baby’s shoulders to help release them and complete delivery.

\medterm{Woods Screw Maneuver} A rotational maneuver used during shoulder dystocia to rotate the fetal shoulders and assist delivery.
\medterm{Woody Tongue} A rare long-lasting infection, usually caused by \textit{Actinomyces}, that makes the tongue and nearby tissues hard and swollen. It can interfere with eating, speaking, or breathing.

\medterm{Wool-Sorter'S Disease (Anthrax)} A form of inhaled anthrax acquired by breathing spores of \textit{Bacillus anthracis}, historically associated with handling contaminated wool or animal products. It can cause severe breathing illness and requires urgent treatment.

\synonyms
Anthrax


\medterm{Word Blindness} An older term for alexia, an acquired inability to read or recognize written words despite adequate vision and previously developed language ability, usually caused by brain dysfunction or injury.

\medterm{Word Blindness (Alexia)} An acquired neurological disorder in which a person loses or markedly impairs the ability to read despite adequate vision and previously acquired language skills.

\synonyms
Alexia

\medterm{Word Deafness (Auditory Verbal Agnosia)} An acquired inability to understand spoken words even though ordinary hearing is intact. It usually results from damage to brain areas that connect hearing with language, such as after a stroke or head injury.

\synonyms
Auditory Verbal Agnosia

\medterm{Word Recognition Scoring} A hearing test measuring how accurately a person repeats or identifies spoken words at a comfortable volume. It helps assess how well the hearing system and brain understand speech.

\medterm{Word Salad} Speech made of words or phrases that have little or no logical connection, making the person difficult to understand. It can occur with severe psychosis, mania, brain disease, or intoxication and needs assessment in context.

\medterm{Working Memory} The ability to hold a small amount of information briefly and use it while completing a task, such as remembering instructions while following them. Attention, stress, sleep, illness, and brain conditions can affect it.

\medterm{Worm Infestation (Helminthiasis)} Infection by parasitic worms, including roundworms, tapeworms, and flukes. Symptoms depend on the species and body site and may include abdominal problems, anemia, cough, itching, or no symptoms.

\synonyms
Helminthiasis

\medterm{Worms Pinworms (Pinworm Infection)} Alternate index label for Pinworm Infection.

\medterm{Wound} Damage to body tissue caused by trauma, surgery, pressure, heat, chemicals, or another cause. A wound may be open or closed and can vary in depth, contamination, bleeding, and injury to nearby nerves, vessels, or organs.

\medterm{Wound Care} Assessment and treatment of a wound through cleansing, bleeding control, removal of dead or contaminated tissue when indicated, appropriate closure or dressing, pain control, and follow-up. Care also includes tetanus prevention, pressure relief, and treatment of infection or underlying conditions; antibiotics are reserved for selected contaminated wounds, bites, or infections.

\medterm{Wound Complication} A problem that develops during healing, such as infection, separation of wound edges, bleeding, fluid
collection, tissue death, or excessive scarring. Increasing pain, redness, swelling, drainage,
fever, or wound opening should be assessed promptly.

\medterm{Wound Dehiscence} Partial or complete reopening of a surgical wound after it has been closed. It may follow infection, poor blood supply, diabetes, increased pressure, or strain on the wound. New separation, drainage, fever, severe pain, or visible deeper tissue needs prompt medical assessment.

\medterm{Wound Evisceration} Separation of a surgical wound with protrusion of internal organs or tissue through the opening. It is a surgical emergency requiring protection of the exposed tissue and urgent medical evaluation.

\medterm{Wound Healing} The body’s process of stopping bleeding, controlling inflammation, replacing damaged tissue, covering the surface, and strengthening the repair. Blood supply, infection, nutrition, diabetes, medicines, smoking, wound depth, and movement can affect healing.

\medterm{Wound Healing Phases} Wound healing has overlapping stages: stopping bleeding, inflammation, building new tissue and skin, and strengthening the repair. The timing and result depend on the wound, blood supply, infection, nutrition, and general health.

\medterm{Wound Infection} Infection of a wound by microorganisms, causing findings such as increasing pain, redness, warmth,
swelling, pus, fever, or delayed healing.

\medterm{Wound Management} Assessment and treatment of acute or chronic wounds through cleansing, debridement when indicated, moisture and exudate control, infection management, pressure relief, nutrition, and correction of underlying disease.
\medterm{Wound VAC} A negative-pressure wound treatment in which a sealed dressing is connected to a pump that gently removes fluid and reduces swelling. It can support new tissue growth in selected complex wounds, but requires proper application and monitoring for bleeding, infection, or skin damage.

\medterm{Wounds} Injuries in which tissue is damaged by trauma, surgery, pressure, heat, chemicals, or other causes; wounds may be open or closed and may require cleansing, closure, or debridement.

\medterm{Wright Stain} A laboratory dye applied to a blood sample so different blood cells and abnormal cell changes can be seen and identified under a microscope.



\medterm{Wrisberg Ganglion} A small cluster of nerve cells connected with the facial nerve. It helps relay taste signals from part of the tongue and sensory signals from nearby areas of the ear.

\medterm{Wrist} The joint connecting the hand to the forearm. It contains eight small wrist bones, ligaments, tendons, nerves, and blood vessels and allows the hand to bend, turn, and move from side to side.

\medterm{Wrist Arthroscopy} A procedure using a small camera and instruments inserted through tiny wrist openings to examine or treat damaged cartilage, ligaments, joint lining, or other wrist problems.

\medterm{Wrist Drop} Inability to lift the wrist and fingers, usually caused by injury or compression of the radial nerve. It may cause weakness, numbness, or difficulty using the hand and needs assessment for the cause.
\medterm{Wrist Fracture} A break in one or more bones near the wrist, most commonly the distal radius or ulna, usually caused by a fall, direct trauma, or weakened bone.

\medterm{Wrist Injury} Damage to the bones, ligaments, tendons, nerves, or vessels of the wrist from a fall, twist, impact, or overuse; fracture, instability, severe swelling, or neurovascular symptoms require assessment.

\medterm{Wrist Joint} The joint where the forearm bones meet the small bones of the hand. It allows bending, straightening, side-to-side movement, and limited rotation while ligaments and cartilage provide stability.

\medterm{Wrist Pain} Pain in or around the wrist, commonly caused by sprain, fracture, tendinitis, arthritis,
or nerve compression such as carpal tunnel syndrome. Persistent pain, deformity,
numbness, or weakness requires assessment.

\synonyms
Pain, wrist

\synonyms
Pain, Wrist

\medterm{Wrist Sprain} Stretching or tearing of one or more ligaments supporting the wrist, causing pain, swelling, bruising, and reduced movement after injury.


\medterm{Writer'S Cramp} A movement disorder in which involuntary muscle tightening causes abnormal hand or arm positions while writing. It may also affect other closely related fine hand movements.

\medterm{Written Documentation} Formal recording of patient information, clinical decisions, and treatment plans; essential for continuity of care and medicolegal purposes.

\medterm{Wry-Neck} Torticollis, an abnormal tilt or rotation of the head caused by involuntary contraction, shortening, or dysfunction of neck muscles.

\medterm{Wryneck (Torticollis)} Abnormal tilting or rotation of the head caused by involuntary contraction, shortening, or dysfunction of neck muscles. It may be congenital, acquired, painful, or neurological.

\medterm{Wt} A medical abbreviation commonly meaning weight. Its meaning should be confirmed from the surrounding documentation because abbreviations can be ambiguous.

\medterm{Wuchereria} A genus of filarial nematodes that includes Wuchereria bancrofti, a major cause of
lymphatic filariasis. Infection is transmitted by mosquitoes and may cause lymphangitis,
chronic lymphedema, hydrocele, and elephantiasis.

\medterm{Wuchereria Bancrofti} A parasitic roundworm spread by mosquitoes and the main cause of lymphatic filariasis. Long-term infection can block lymph drainage and cause swelling of the limbs or genitals, sometimes becoming severe and permanent.

\medterm{Wuhan Coronavirus (2019-NCoV)} Alternate index label for 2019-nCoV.

\medterm{Weight Loss Surgery} Weight-loss surgery, or bariatric surgery, changes the stomach or digestive tract to help treat severe obesity and related conditions. It requires long-term nutritional, medical, and lifestyle follow-up.

\medterm{Wheelchair} A wheelchair is a mobility device with a seat and wheels that helps a person move when walking is difficult or unsafe. Proper fitting and training improve comfort, safety, and independence.

\medterm{Wheelchairs} Wheelchairs are mobility devices used by people who cannot walk safely or efficiently. They may be manual or powered and should be selected and adjusted to the person’s posture, strength, environment, and daily needs.

\medterm{Xanthelasma} A harmless yellowish patch or bump caused by fat deposits in the skin of the eyelids, usually near the inner corners of the eyes. It can occur with high blood fats but may also occur when blood-fat levels are normal.

\synonyms
Xanthelasma Palpebrarum; Eyelid Xanthoma; Palpebral Xanthoma

\medterm{Xanthine} A natural breakdown product of substances found in cells and foods. It normally changes into uric acid, but can build up when a rare enzyme disorder blocks this pathway and may then contribute to kidney or urinary stones.

\medterm{Xanthine Oxidase} An enzyme that helps convert cell and food breakdown products into uric acid. Medicines used for gout can block this enzyme, lowering uric-acid production and helping prevent urate crystal deposits.

\medterm{Xanthine Oxidase Inhibitors} Medicines such as allopurinol and febuxostat that lower uric acid by blocking an enzyme involved in its production. They are used for selected people with gout or to prevent high uric acid during some cancer treatments. Choice and timing depend on the person’s condition.

\synonyms
XOIs

\medterm{Xanthinuria} A rare inherited disorder in which the body cannot properly change xanthine into uric acid. Xanthine can build up and form kidney or urinary stones; muscle problems may also occur.

\medterm{Xanthochromia} Yellow or amber color in the fluid around the brain and spinal cord. It can occur when blood has broken down in the fluid, including after bleeding around the brain, but other causes are possible and the meaning depends on timing and other test results.

\synonyms
CSF Xanthochromia; Cerebrospinal Fluid Yellowing

\medterm{Xanthogranulomatous} Describing inflammation in which tissue contains many immune cells filled with fat. This pattern can occur in long-lasting infections and other inflammatory conditions.

\medterm{Xanthogranulomatous Cholecystitis} A rare, long-lasting inflammation of the gallbladder that can thicken and damage its wall. It may spread to nearby tissues and can look like gallbladder cancer on scans.

\synonyms
XGC

\medterm{Xanthogranulomatous Oophoritis} A rare long-lasting inflammation of an ovary, often linked to infection or endometriosis. It can form a mass that resembles ovarian cancer on scans and may cause pelvic pain, fever, or abnormal bleeding.

\medterm{Xanthogranulomatous Pyelonephritis} A severe, long-lasting kidney infection that destroys normal kidney tissue and replaces it with inflamed tissue. It is often associated with blockage or stones and may require antibiotics, drainage, or removal of the damaged kidney.

\medterm{Xanthoma} A yellow-orange spot, bump, or patch formed by fat deposits in the skin, tissue beneath the skin, or tendons. It may occur with high blood fats, diabetes, some liver diseases, or an inherited lipid disorder.

\medterm{Xanthoma Disseminatum} A rare disorder in which many yellow, red-brown, or orange bumps develop on the skin, especially in body folds. The mouth, throat, or other organs may also be affected, and some people develop problems controlling water balance.

\synonyms
Montgomery Syndrome; Disseminate Xanthoma

\medterm{Xanthoma Eruptivum} A sudden outbreak of many small yellow-orange bumps, usually on the buttocks, arms, legs, or shoulders. It is often linked to very high blood triglyceride levels.

\synonyms
Eruptive Xanthoma; Eruptive Xanthomatosis

\medterm{Xanthoma Planum} Flat or slightly raised yellow skin patches caused by the buildup of fat inside immune cells. They may occur in skin folds, the palms, or more widely on the body and can be linked to high blood fats or some blood disorders.

\synonyms
Plane Xanthoma; Xanthoma Striatum Palmare

\medterm{Xanthomas} Yellowish papules, plaques, or nodules caused by lipid accumulation in skin immune cells. They may reflect high cholesterol or triglycerides, diabetes, cholestatic liver disease, or an inherited lipid disorder and should prompt appropriate evaluation.
\medterm{Xanthoma Striatum Palmare} Yellow or orange lines or patches in the creases of the palms caused by fat deposits in the skin. They are strongly linked to some inherited blood-fat disorders, especially type III dysbetalipoproteinemia, but blood tests are needed.

\synonyms
Palmar Crease Xanthoma; Planar Palmar Xanthoma

\medterm{Xanthoma Tendinosum} Firm, usually painless fat deposits in tendons or ligaments, most often in the Achilles tendon or tendons of the hands. They are strongly linked to inherited high cholesterol.

\synonyms
Tendon Xanthoma; Tendinous Xanthoma

\medterm{Xanthomatosis} A disorder in which fatty deposits called xanthomas develop in the skin, tendons, or organs. It may be linked to inherited or acquired problems with blood fats.

\medterm{Xanthoma Tuberosum} Slowly enlarging, painless, firm, yellowish-red nodular lipid deposits occurring over extensor pressure areas, particularly the elbows, knees, heels, and buttocks. They are strongly associated with familial hypercholesterolemia and type III hyperlipoproteinemia.

\synonyms
Tuberous Xanthoma; Tuberous Xanthomatosis

\medterm{Xanthopsia} A visual disturbance in which objects appear yellow or yellowish. It can occur with some eye disorders, medicine toxicity such as digoxin toxicity, or disorders of the nerves or retina.

\medterm{Xanthosis} Yellow discoloration of skin or other tissues caused by carotene or another yellow substance. Unlike jaundice, it does not usually make the whites of the eyes yellow; the cause is assessed from the person’s diet and health.

\medterm{X Chromosome} One of the two types of human sex chromosomes. It carries genes involved in sex development as well as many other body functions. Most people have one or two X chromosomes.

\medterm{Xenobiotic} A chemical that does not naturally belong in the body, such as a medicine, pollutant, pesticide, or industrial compound. The body changes and removes it, but some resulting chemicals can be harmful.

\medterm{Xenobiotic Metabolism} The body’s processing of foreign chemicals, including medicines, pollutants, and toxins. Enzymes, especially in the liver, change them so they can be removed; the changed substances may remain active or become harmful.

\medterm{Xenobiotic Resistance} The ability of cells or organisms to withstand the effects of foreign chemicals, medicines, or toxins. It may result from breaking down the substance faster, pumping it out of cells, or changing the substance’s target.

\synonyms
Multidrug Resistance; Xenobiotic Tolerance

\medterm{Xenograft} Tissue, cells, or a biological material transplanted between different species, such as animal tissue used in a human. Rejection and infection are important risks, although some materials are used temporarily as tissue coverings or substitutes.

\medterm{Xenon Anesthesia} Use of xenon gas to cause unconsciousness and prevent pain during an operation. Xenon acts on brain signaling and is cleared mainly through breathing.

\medterm{Xenotransplantation} Transplanting living cells, tissues, or organs between different species, usually from an animal to a human. It may increase the supply of transplant material but carries major risks of immune rejection, infection, and poor organ function.

\medterm{Xerochilia} Abnormal, persistent dryness, scaling, and superficial vertical cracking or fissuring of the lips, resulting from environmental dehydration, persistent lip-licking (cheilitis simplex), chronic mouth breathing, or systemic therapy with oral retinoids (such as isotretinoin).

\synonyms
Dry Lips

\medterm{Xeroderma} Alternate name; see \textit{Dry skin}.

\medterm{Xeroderma Pigmentosum} A rare usually autosomal-recessive disorder of nucleotide-excision DNA repair, with a variant form affecting translesion synthesis. Minimal ultraviolet exposure can cause severe sunburn or persistent redness, early freckle-like pigmentation, and rapid photoaging; skin and ocular-surface cancers occur at a young age, and some people develop progressive neurologic disease. Strict ultraviolet protection, regular skin and eye examinations, and early treatment of suspicious lesions are essential.

\medterm{Xeroderma Pigmentosum Complementation Groups} Genetic subtypes of xeroderma pigmentosum. Each subtype results from a change in a different gene involved in repairing damage caused by ultraviolet light. All subtypes can cause extreme sun sensitivity, freckling, and a very high risk of skin cancer.

\synonyms
XP Complementation Groups; NER Defect Groups

\medterm{Xerophthalmia} Dryness and damage of the eye caused by severe vitamin A deficiency. It may begin with poor night vision and dry eyes and can progress to corneal damage and permanent sight loss.
\synonyms{Vitamin A Deficiency Eye Disease; Ocular Xerosis; Keratomalacia}

\medterm{Xerophthalmic Ulceration} Severe progressive ulceration of the cornea resulting from severe vitamin A deficiency (classified as WHO stage X3A when involving less than one-third of the corneal area, and stage X3B when involving more). Untreated, it rapidly causes corneal liquefactive necrosis (keratomalacia) and permanent blindness.

\synonyms
Vitamin A Corneal Ulcer; Keratomalacic Ulceration

\medterm{Xerosis} Alternate name; see \textit{Dry skin}.

\medterm{Xerostomia} The sensation of dry mouth caused by reduced or altered saliva production. It may result from
medicines, dehydration, salivary-gland disease, radiation, or systemic illness and can impair speaking, swallowing,
taste, and dental health.

\medterm{Xerotic Eczema} A common, intensely pruritic inflammatory dermatosis occurring predominantly on the pretibial lower extremities of elderly individuals during cold, dry winter months, characterized by profound epidermal dehydration, scaling, and fine superficial erythematous fissures resembling "cracked porcelain" or a "dry riverbed."

\synonyms
Asteatotic Eczema; Eczema Craquele; Winter Itch

\medterm{X-Inactivation} A process in early development that switches off most genes on one of two X chromosomes in each cell. It helps balance gene activity, although some genes remain active and the pattern can differ between cells.

\medterm{Xiphisternal Joint} The joint between the lower part of the breastbone and the xiphoid process. It often becomes partly or fully hardened with age.

\synonyms
Xiphisternal Synchondrosis; Xiphisternal Articulation

\medterm{Xiphoidalgia} Pain or tenderness over the xiphoid process, the small lower part of the breastbone. Pain may spread to the upper abdomen, chest, back, or shoulders and may worsen with touch, bending, lifting, or coughing.

\synonyms
Xiphoidodynia; Hypersensitive Xiphoid Syndrome; Xiphoid Syndrome

\medterm{Xiphoidectomy} Surgery to remove the xiphoid process, the small lower part of the breastbone. It may be done for persistent pain in this area or to provide access during some chest or abdominal operations.

\medterm{Xiphoid Process} The small pointed lower part of the breastbone. Its shape varies between people and it provides attachment for some muscles and ligaments.

\medterm{X-Linked} Describing a gene or condition caused by a change on the X chromosome. The pattern differs between people with one X chromosome and those with two; fathers pass their X chromosome to daughters, not sons.

\medterm{X-Linked Adrenoleukodystrophy} An inherited disorder, usually affecting boys, in which certain fats build up and damage the adrenal glands and nervous system. It can cause behavior or learning changes, movement problems, seizures, or adrenal failure.

\medterm{X-Linked Agammaglobulinemia} An inherited immune disorder, usually affecting boys, in which the body cannot make enough antibodies. Repeated bacterial infections often begin in infancy or early childhood and are treated with antibody replacement and infection care.

\medterm{X-Linked Chronic Granulomatous Disease} An inherited immune disorder, usually affecting boys, caused by a change in the \textit{CYBB} gene. Certain white blood cells cannot kill some bacteria and fungi normally, leading to repeated serious infections and areas of ongoing inflammation.

\synonyms
X-Linked CGD; CYBB Deficiency; Cytochrome b-245 Beta Deficiency

\medterm{X-Linked Dominant Inheritance} A pattern in which one changed gene on the X chromosome can cause a condition. An affected father passes the changed X chromosome to all daughters and no sons.

\medterm{X-Linked Emery-Dreifuss Muscular Dystrophy} An inherited muscle disorder, usually affecting boys, that causes stiff joints, slowly worsening muscle weakness, and problems with the heart’s electrical rhythm.

\synonyms
EDMD1; Emerin Deficiency Muscular Dystrophy

\medterm{X-Linked Hyper-IgM Syndrome} An inherited immune disorder, usually affecting boys, in which the body cannot change one type of protective antibody into other types. It causes repeated lung and sinus infections and may cause infections by organisms that usually do not cause illness.

\synonyms
HIGM1; CD40L Deficiency; Hyper-IgM Syndrome Type 1

\medterm{X-Linked Hypophosphatemia} An inherited disorder that causes the kidneys to lose too much phosphate. Low phosphate can lead to soft or poorly growing bones, bone pain, bowed legs, dental problems, and reduced height.

\medterm{X-Linked Ichthyosis} An inherited, usually X-linked recessive skin disorder caused by a variant or deletion affecting the STS gene and deficiency of steroid sulfatase. It usually affects males and causes widespread dry, dark, plate-like scales that often begin in infancy or early childhood.

\medterm{X-Linked Inheritance} Inheritance of a gene change located on the X chromosome. The effects depend on the gene, the type of change, and whether a person has one or two X chromosomes; fathers do not pass an X chromosome to sons.

\medterm{X-Linked Lymphoproliferative Disease Type 1} An inherited immune disorder, usually affecting boys, in which the body can react severely to Epstein--Barr virus. It may cause severe glandular fever, very high inflammation, low antibody levels, or lymphoma.

\synonyms
XLP1; Duncan Disease; SAP Deficiency

\medterm{X-Linked Lymphoproliferative Syndrome} An inherited immune disorder, usually affecting boys, in which the body reacts abnormally to infections. Severe Epstein–Barr virus illness, persistent fever, liver or blood problems, lymphoma, and intestinal inflammation may occur.

\medterm{X-Linked Myotubular Myopathy} A rare inherited muscle disease, usually affecting boys, that causes severe weakness and low muscle tone from birth. Babies may have drooping eyelids, feeding difficulty, and dangerous weakness of the breathing muscles.

\synonyms
XLMTM; Myotubularin Deficiency

\medterm{X-Linked Recessive Inheritance} A pattern in which a changed gene on the X chromosome usually causes a condition in people with one X chromosome. People with two X chromosomes are often unaffected carriers, although they can sometimes have symptoms. An affected father does not pass the changed X chromosome to sons.

\medterm{X-Linked Retinitis Pigmentosa} An inherited disorder in which the retina slowly loses its ability to sense light. It usually causes night blindness followed by narrowing of the field of vision and worsening sight.

\synonyms
XLRP; X-Linked Pigmentary Retinopathy

\medterm{X-Linked Retinoschisis} An inherited eye disorder, usually affecting boys, in which layers of the retina separate. It can cause blurred central or distance vision, crossed eyes, bleeding, or retinal detachment in severe cases.

\medterm{X-Linked Severe Combined Immunodeficiency} A severe inherited immune disorder, usually affecting boys, in which important immune cells do not develop normally and antibody protection is poor. Babies are highly vulnerable to serious infections.

\medterm{X-Linked Sideroblastic Anemia} An inherited form of anemia, usually affecting boys, in which red blood cells cannot make hemoglobin normally. Iron builds up in developing blood cells, causing small, pale red blood cells and anemia.

\synonyms
XLSA; ALAS2 Deficiency; Hereditary Sideroblastic Anemia

\medterm{X-Linked Spinal and Bulbar Muscular Atrophy} An inherited muscle and nerve disorder that usually begins in adulthood and slowly causes weakness. It may also cause muscle twitching, trouble speaking or swallowing, breast enlargement, and reduced fertility.

\synonyms
Kennedy Disease; SBMA; Bulbo-Spinal Muscular Atrophy

\medterm{X-Ray} An imaging test that uses a small amount of radiation to create a picture of structures inside the body. It is especially useful for bones, lungs, teeth, and some other tissues.

\medterm{X-Ray Attenuation} The reduction in strength of an X-ray beam as it passes through the body. Bone, soft tissue, and air reduce the beam by different amounts, creating the differences seen in the image.

\medterm{X-Ray Crystallography} A laboratory method that uses X-rays to work out the three-dimensional structure of molecules such as proteins. The molecule is made into a crystal, and the pattern produced by the X-rays is analyzed.

\synonyms
X-Ray Crystal Diffraction Analysis; Macromolecular Crystallography

\medterm{X-Ray Shield} Protective material, often lead or another dense substance, used to reduce exposure to scattered X-rays during imaging.

\medterm{X-Ray Therapy} Treatment that uses carefully aimed X-ray radiation to damage abnormal cells or relieve symptoms. It is used mainly for cancer and may be given from a machine outside the body or near the target.

\medterm{XX Male Syndrome} A difference in sex development in which a person with two X chromosomes develops typically male physical features, usually because a sex-development gene has moved onto an X chromosome. Infertility is common, and care is individualized.

\medterm{XXX Syndrome} Alternate name; see \textit{Triple X syndrome}, a sex-chromosome variation in which a female has an additional X chromosome.

\medterm{XXXX Syndrome} XXXX syndrome is a rare sex-chromosome variation in which a female has four X chromosomes; it may cause developmental or learning difficulties, distinctive physical features, and ovarian dysfunction, with severity varying widely.

\medterm{Xylitol} A sugar alcohol used as a sweetener and in some oral-care products. Large human ingestions may cause gastrointestinal symptoms; it is highly toxic to dogs.

\medterm{Xylose} A simple sugar found in some plants. A form called D-xylose has been used in tests of how well the small intestine absorbs nutrients and is also used in laboratory studies.

\medterm{Xylose Absorption Test} A test of how well the small intestine absorbs a simple sugar called D-xylose. Low amounts in the blood or urine may suggest poor absorption from the intestine; normal results with other signs of poor absorption may suggest a problem with the pancreas.

\synonyms
D-Xylose Absorption Test; D-Xylose Tolerance Test

\medterm{XYY Syndrome} A sex-chromosome variation in which a male has an extra Y chromosome. Many affected individuals have few or no obvious symptoms, though some have tall stature, learning difficulties, or developmental differences.

\medterm{X-linked Juvenile Retinoschisis} An inherited retinal disorder, usually affecting boys, in which layers of the retina separate during childhood. It can cause reduced central vision, difficulty seeing in dim light, crossed eyes, or retinal complications.

\medterm{X-linked Adult Spinal Muscular Atrophy} Alternate name for Kennedy disease, or X-linked spinal and bulbar muscular atrophy. It is an inherited androgen-receptor disorder that gradually causes weakness, muscle wasting, twitching, swallowing or speech difficulty, and sometimes breast enlargement.

\medterm{X-linked Intellectual Disability} Intellectual disability caused by a change in a gene on the X chromosome. Features vary widely and may include developmental delay, learning difficulty, behavior differences, seizures, or associated physical findings.

\medterm{XO Syndrome} A form of Turner syndrome caused by complete or partial absence of one X chromosome. It may cause short stature, ovarian insufficiency, infertility, characteristic physical findings, and heart, kidney, hearing, or thyroid problems.

\medterm{XP} Abbreviation for xeroderma pigmentosum, a rare DNA-repair disorder causing extreme sensitivity to ultraviolet light and a markedly increased risk of early skin and eye cancers.

\medterm{XXY Males} People with a typically male pattern of development and an extra X chromosome, usually 47,XXY. This is commonly called Klinefelter syndrome and may cause small testes, low testosterone, infertility, tall stature, or learning and language difficulties.

\medterm{XXY Syndrome} Alternate name for Klinefelter syndrome, a sex-chromosome variation in which a typically male individual has an extra X chromosome. Effects vary and may include low testosterone, infertility, small testes, tall stature, and developmental or language differences.

\medterm{YAG Capsulotomy} A laser procedure that makes a small opening in the cloudy membrane behind an artificial lens after cataract surgery. It can restore vision when this clouding causes blurring or glare and is usually done without another incision.

\medterm{YAG Laser Capsulotomy} Alternate name; see \textit{YAG Capsulotomy}.

\medterm{YAG Laser Surgery} Treatment using a focused laser to make a precise opening or alter tissue without a conventional surgical incision. It is used mainly for selected eye conditions, and risks depend on the part of the eye treated.

\medterm{Yawn} An involuntary behavior consisting of a deep inhalation followed by a slower exhalation, often with stretching of the jaw and facial muscles. Yawning is common with fatigue but excessive or unusual yawning can occur with some medications or neurologic conditions.

\medterm{Yawning} Alternate name; see \textit{Yawn}.

\medterm{Yaws} Yaws is a chronic infection caused by Treponema pallidum subsp. pertenue and spread mainly through close skin contact. It can cause an initial papilloma or ulcer and later destructive lesions of skin, bone, and cartilage; antibiotics can cure the infection.

\medterm{Yaws Papilloma} The primary clinical cutaneous lesion of early yaws (\textit{Treponema pallidum} subsp. \textit{pertenue} infection), also termed a mother yaw or frambesioma. It presents as an initial painless, raised, yellow-crusted, raspberry-like papule or ulcer at the site of cutaneous inoculation.

\synonyms
Mother Yaw; Frambesioma; Primary Yaws Lesion

\medterm{Y-Box Binding Protein} A protein that attaches to DNA and RNA and helps control which genes are active and how proteins are made. It also helps cells respond to stress and repair damaged DNA.

\synonyms
YB-1; Y-Box Transcription Factor; Nuclease-Sensitive Element-Binding Protein 1

\medterm{Y Chromosome} One of the two human sex chromosomes. It usually carries genes involved in testis development and is passed by a biological father to a child who inherits it; chromosome patterns and body development can vary.

\medterm{Y-Chromosome Microdeletion} A small missing section of the Y chromosome that can interfere with sperm production and cause very low or absent sperm counts.

\synonyms
AZF Microdeletion

\medterm{Yeast} A type of fungus made of single cells. Some yeasts normally live on or in the body, while others can cause infections of the mouth, skin, vagina, lungs, blood, or other organs.

\medterm{Yeast Infection} An infection caused by Candida yeast that may affect the mouth, skin, vagina, penis, or other tissues. Symptoms depend on the site and may include itching, soreness, redness, white plaques, discharge, or rash. Candida may also be present without causing disease, so the significance of a finding depends on the site, symptoms, and evidence of tissue invasion.

\medterm{Yeast Infection In Women And Men} Genital Candida infection may affect the vulva or vagina in women and the glans penis or foreskin in men. Symptoms can include itching, soreness, redness, discharge, or rash.

\medterm{Yeast Present} A laboratory report showing yeast in a specimen. This may mean harmless colonization, contamination, or infection; the meaning depends on the symptoms, specimen, and type of yeast.

\medterm{Yeast Vulvitis} Inflammation of the vulva associated with a Candida yeast infection. It may cause itching, burning, redness, swelling, and soreness; other conditions can cause similar symptoms.

\medterm{Yellow Cartilage} A flexible type of cartilage containing many elastic fibers. It is found in the outer ear, the epiglottis, parts of the voice box, and the tube connecting the middle ear to the throat.

\synonyms
Elastic Cartilage; Cartilago Elastica

\medterm{Yellow Fever} Yellow fever is an acute mosquito-borne infection caused by yellow fever virus. Fever, headache, muscle aches, nausea, and vomiting may be followed by jaundice, bleeding, shock, and organ failure; vaccination is the main preventive measure.

\medterm{Yellow Fever Vaccine} A vaccine made from a weakened yellow fever virus. It gives long-lasting protection, but rare serious reactions can affect the brain or internal organs.

\synonyms
17D Vaccine

\medterm{Yellow Fever Virus} A flavivirus transmitted mainly by infected mosquitoes that can cause yellow fever, ranging from a mild febrile illness to jaundice, bleeding, shock, and organ failure.

\medterm{Yellow Ligament} A series of elastic bands connecting the back parts of neighboring bones in the spine. They help support the spine and surround the spinal canal.

\synonyms
Ligamentum Flavum; Flavum Ligament

\medterm{Yellow Marrow} Bone marrow containing mostly fat instead of blood-forming cells. It is common in the hollow centers of adult long bones and can change back toward blood-forming marrow when the body needs more blood cells.

\synonyms
Fatty Bone Marrow; Medulla Ossium Flava

\medterm{Yellow Nail Syndrome} Alternate name; see \textit{Yellow-Nail Syndrome}.

\medterm{Yellow-Nail Syndrome} A rare disorder causing yellow, thickened, slowly growing nails together with swelling from lymph fluid, usually in the legs, and lung problems such as repeated infections or fluid around the lungs. The cause is not fully understood.

\medterm{Yellow Nail Syndrome Triad} The three main features of yellow nail syndrome: slow-growing yellow or thickened nails, swelling from lymph-fluid buildup, and long-term breathing or lung problems.

\synonyms
YNS Triad; Lymphedema-Pleural Effusion-Yellow Nails

\medterm{Yellow Spot} The macula, a small yellowish area near the center of the retina at the back of the eye. It provides the sharp central vision used for reading and seeing fine detail.

\synonyms
Macula Lutea; Anatomical Macula

\medterm{Yellow Tongue} Yellow discoloration or coating of the tongue that may result from oral hygiene, smoking, dry mouth, diet, medicines, or an underlying illness.

\medterm{Yergason Test} A shoulder examination in which a person turns the palm upward and rotates the arm outward against resistance. Pain or a snapping sensation near the front of the shoulder may suggest a problem involving the biceps tendon, but the test is not diagnostic by itself.

\medterm{Yersinia} A group of bacteria that can cause intestinal illness, swollen lymph nodes in the abdomen, or plague. Different species spread through contaminated food or water, infected animals, or fleas.

\medterm{Yersinia Enterocolitica} A bacterium that can cause diarrhea, fever, and abdominal pain, sometimes mimicking appendicitis. People usually acquire it from contaminated food, water, or contact with infected animals.

\medterm{Yersinia Enterocolitica Pseudoappendicitis} A Yersinia infection that causes pain in the lower right abdomen and swollen lymph nodes near the intestine, making it resemble appendicitis. It can also cause fever and diarrhea, often after contaminated food such as undercooked pork.

\synonyms
Yersinia Mesenteric Adenitis; Pseudoappendiceal Syndrome

\medterm{Yersinia Infection} Infection caused by Yersinia bacteria, usually producing diarrhea, abdominal pain, fever, or vomiting. Some cases resemble appendicitis; bloody diarrhea, dehydration, or severe or persistent illness can occur.

\medterm{Yersinia Outer Proteins} Proteins made by disease-causing Yersinia bacteria that are injected into human cells. They weaken the immune response and help the bacteria avoid being destroyed.

\synonyms
Yops; Yersinia Effector Proteins; Yop Virulence Factors

\medterm{Yersinia Pestis} The bacterium that causes plague, a severe infection usually spread by infected fleas or animals. It may cause painful swollen lymph nodes, bloodstream infection, or lung disease and can be life-threatening.

\medterm{Yersinia Pestis Pneumonic Plague} A severe form of plague affecting the lungs. It can spread through breathing in infected droplets or when plague spreads from another part of the body. It may cause high fever, breathing difficulty, chest pain, cough, shock, and rapidly worsening lung failure.

\synonyms
Pneumonic Plague; Pulmonary Yersiniosis

\medterm{Yersinia Pseudotuberculosis} A bacterium usually acquired from contaminated food or water. It can cause fever, diarrhea, and pain from swollen lymph nodes in the abdomen, sometimes resembling appendicitis.

\medterm{Yersinia Reactive Arthritis} Joint inflammation that develops after a Yersinia intestinal infection, usually one to four weeks later. It often affects a few joints, especially in the legs, without bacteria being present inside the joints.

\medterm{Yersiniosis} Infection caused by Yersinia bacteria, most often Y. enterocolitica or Y. pseudotuberculosis, producing diarrhea, abdominal pain, fever, or pseudoappendicitis.

\medterm{Yew Poisoning} Poisoning after eating parts of the yew plant, especially its leaves or seeds. The toxins can cause vomiting, dizziness, seizures, dangerous heart rhythms, collapse, and death.

\medterm{Y-Linked Gene} A gene located specifically on the non-recombining region of the Y chromosome (such as the \textit{SRY} gene, \textit{ZFY}, or \textit{AZF} azoospermia factor region). These genes are transmitted strictly from biological fathers to sons and are essential for testicular differentiation, gonadal development, and male spermatogenesis.

\synonyms
Holandric Gene; Y-Chromosomal Gene

\medterm{Y-Linked Inheritance} Inheritance of a condition caused by a gene on the Y chromosome. A biological father passes a Y-linked gene to sons who inherit his Y chromosome, but not to daughters who inherit his X chromosome.

\medterm{Y-Maze Test} A test used mainly in laboratory animals to study movement, exploration, and short-term memory. The animal moves through three arms of a Y-shaped maze, and the pattern of movement is recorded.

\synonyms
Y-Maze Spontaneous Alternation; Y-Maze Behavioral Paradigm

\medterm{Yolk Sac} A temporary structure that supports an early embryo before the placenta is fully working. It helps form the first blood cells and part of the early intestine and can be seen on early pregnancy ultrasound.

\medterm{Yolk Sac Carcinoma} An aggressive germ-cell cancer that often produces high levels of alpha-fetoprotein (AFP). AFP can support the diagnosis and help track treatment, but it is interpreted with examination, imaging, and tissue testing.

\synonyms
Endodermal Sinus Tumor; Yolk Sac Germ Cell Tumor

\medterm{Yolk Sac Tumor} A cancer of germ cells that often occurs in the ovary or testicle in children or young adults and produces alpha-fetoprotein. It can also occur outside the reproductive organs.

\medterm{Yolk-Sac Tumor} Alternate name; see \textit{Yolk Sac Tumor}.

\medterm{Yolk Sac Tumor Schiller-Duval Bodies} A characteristic microscope finding in yolk sac tumors. It looks somewhat like a tiny kidney filter, with tumor cells arranged around a small blood vessel. The finding helps confirm the diagnosis when considered with other test results.

\synonyms
Schiller-Duval Bodies; Glomeruloid Endodermal Bodies

\medterm{Yolk Stalk} A temporary connection between the early yolk sac and the developing intestine. It normally disappears before birth; a remaining portion can form an abnormal intestinal pouch or other problem.

\medterm{Young-Burgess Classification} A system used to describe pelvic-ring fractures according to how the injury happened and how stable the pelvis is. Groups include side-to-side compression, front-to-back compression, vertical separation, and mixed injuries.

\synonyms
Young-Burgess Pelvic Fracture Classification; Burgess Pelvic Injury Scale

\medterm{Young Syndrome} A rare condition, usually affecting males, that combines blockage of sperm transport with repeated sinus and lung infections. It may also cause bronchiectasis, a long-term widening and damage of the airways.

\synonyms
Sinopulmonary-Infertility Syndrome; Barry-Perkins-Young Syndrome

\medterm{Y-Plate Osteosynthesis} Surgical repair of a complex broken bone using a metal plate shaped like the letter Y. The plate holds the bone pieces in place while they heal.

\synonyms
Y-Plate Fixation; Y-Shaped Contoured Bone Plate

\medterm{Yttrium-90} A radioactive substance used in selected treatments to deliver radiation to abnormal tissue, such as some liver tumors or inflamed joints.

\medterm{Yttrium-90 Microspheres} Tiny glass or resin beads containing radioactive yttrium-90. They are delivered through blood vessels to selected liver tumors, where they release radiation.

\synonyms
Y-90 Radioembolization Microspheres; TheraSphere; SIR-Spheres

\medterm{Yttrium-90 Radioembolization} A treatment in which tiny beads containing radioactive yttrium-90 are delivered through liver blood vessels to release radiation inside liver tumors.

\synonyms
SIRT; Y-90 SIRT; Y-90 Transarterial Radioembolization

\medterm{Yttrium-Aluminum-Garnet Laser} A laser that uses a crystal to produce a focused beam of light. It is used in eye care for procedures such as opening a cloudy membrane after cataract surgery and making a small opening for some types of glaucoma.

\synonyms
Nd:YAG Laser; Neodymium:YAG Laser

\medterm{Yuzpe Regimen} An older form of emergency contraception using two doses of combined birth-control hormones. It has largely been replaced by emergency contraception containing levonorgestrel or ulipristal.

\medterm{Y-V Plasty} A surgical technique in which a Y-shaped cut is closed as a V. This lengthens or widens nearby tissue and may help release a tight scar or other narrowing.

\synonyms
Y-V Advancement Flap; Y-V Lengthening Procedure

\medterm{Z-Plasty} A reconstructive operation that rearranges nearby skin in a Z-shaped pattern. It can lengthen a tight scar, change the direction of a scar, and improve movement or appearance.

\medterm{ZAP-70 Deficiency} A rare inherited immune disorder in which T cells cannot develop or work properly. Affected children may have repeated, severe infections and need specialist immune treatment and genetic evaluation.

\medterm{Zafirlukast} A leukotriene-receptor antagonist used for long-term control and prevention of asthma symptoms. It is taken orally and is not a rescue treatment for an acute asthma attack; liver injury and drug interactions are important safety considerations.


\synonyms
Accolate

\medterm{Zanamivir} An inhaled neuraminidase inhibitor used to treat or prevent influenza A and B when started appropriately. It may cause bronchospasm and is used cautiously in people with underlying airway disease.

\synonyms
Relenza

\medterm{Zavanelli Maneuver} An emergency procedure for severe shoulder dystocia that does not respond to other measures. The clinician gently returns the baby’s head to the birth canal, followed by immediate cesarean delivery.

\synonyms
Cephalic Replacement Maneuver

\medterm{Zellweger Syndrome} A severe inherited disorder in which cells cannot make peroxisomes, structures that help process fats and other substances. It begins in infancy and may cause weak muscle tone, feeding problems, seizures, liver disease, and impaired brain development.

\synonyms
ZSD; Cerebrohepatorenal Syndrome

\medterm{Zenker Diverticulum} An acquired, false (pulsion) outpouching of the pharyngeal mucosa and submucosa through Killian's triangle, an area of anatomical muscular weakness between the thyropharyngeus and cricopharyngeus muscles in the posterior hypopharynx. It develops secondary to increased swallowing pressure caused by uncoordinated relaxation or spasm of the cricopharyngeus muscle (upper esophageal sphincter). Most commonly affecting elderly individuals, symptoms include progressive difficulty swallowing (dysphagia), regurgitation of undigested food hours after eating, foul breath (halitosis), a compressible neck lump, and cough or recurrent aspiration pneumonia. Diagnosis is established by a barium swallow study; blind endoscopy or nasogastric tube placement is avoided due to the high risk of puncturing the pouch. Treatment involves cricopharyngeal myotomy with surgical or endoscopic diverticular division.

\synonyms
Zenker's Diverticulum; Pharyngoesophageal Diverticulum; Hypopharyngeal Diverticulum; Pulsion Diverticulum


\medterm{Ziconotide} A strong pain medicine delivered continuously into the fluid around the spinal cord for selected severe chronic pain. It can cause dizziness, confusion, hallucinations, and other serious nervous-system or mental-health effects.

\synonyms
Prialt

\medterm{Zidovudine} An HIV medicine that interferes with the virus’s ability to copy itself. It may be used in selected treatment plans and to reduce transmission around childbirth; important risks include anemia, low white-cell counts, muscle problems, and other toxicity.

\synonyms
Azidothymidine; AZT; Retrovir

\medterm{Ziehl-Neelsen Stain} A laboratory stain used to find bacteria with a waxy coating that resists removal by acid, especially the bacteria that cause tuberculosis. These bacteria appear red against a contrasting background under a microscope.

\synonyms
Acid-Fast Stain; AFB Stain

\medterm{Zieve Syndrome} A temporary complication of heavy alcohol use and sudden liver injury. It causes early breakdown of red blood cells, yellowing of the skin or eyes, and increased blood fats. Treatment includes stopping alcohol and managing liver complications.

\synonyms
Hemolytic Anemia with Hyperlipidemia and Jaundice

\medterm{Zika} An infection caused by Zika virus. It often causes no symptoms or a mild fever, rash, red eyes, joint pain, and muscle pain. Infection during pregnancy can harm the developing baby, and it can rarely trigger Guillain-Barré syndrome.

\synonyms
Zika Virus Disease; ZVD; Congenital Zika Syndrome

\medterm{Zika Virus} A virus spread mainly by infected \textit{Aedes} mosquitoes and also through sexual contact and, less commonly, blood or pregnancy-related transmission. Infection is often mild or symptom-free but can injure a developing fetus.

\synonyms
ZIKV

\medterm{Zileuton} A 5-lipoxygenase inhibitor used for prevention and chronic control of asthma. It does not relieve acute bronchospasm; liver-enzyme elevations and interactions with other medicines require monitoring.

\synonyms
Zyflo

\medterm{Zinc} An essential mineral needed for many body reactions, normal growth, immunity, reproduction, and wound healing. Too little can cause deficiency; excessive supplements can cause stomach symptoms and copper deficiency.

\synonyms
Zn

\medterm{Zinc Acetate} An oral zinc salt used to treat zinc deficiency and, in specific formulations, to reduce intestinal copper absorption in Wilson disease. Dose and monitoring depend on the indication; nausea and gastrointestinal irritation may occur.

\medterm{Zinc Deficiency} A state of inadequate available zinc caused by insufficient intake, malabsorption, excessive losses, or increased requirements. Manifestations may include impaired growth, dermatitis, alopecia, diarrhea, altered taste, delayed wound healing, and increased infection susceptibility.

\synonyms
Hypozincemia; Zinc Depletion

\medterm{Zinc Oxide Ointment} A topical preparation containing zinc oxide that forms a protective barrier over irritated or moist skin. It is commonly used for diaper dermatitis and minor skin irritation.

\synonyms
Topical Zinc Barrier Cream; Zinc Oxide Paste

\medterm{Zinc Oxide} An insoluble zinc compound used topically as a skin protectant and barrier, including in preparations for diaper dermatitis and minor irritation. It can reduce contact with moisture and irritants but does not treat every cause of a rash.

\synonyms
ZnO

\medterm{Zinc Sulfate} An oral or parenteral zinc salt used to prevent or treat zinc deficiency when clinically indicated. Gastrointestinal upset is common with oral therapy, and prolonged excessive use can cause copper deficiency.

\medterm{Zinc Toxicity} Adverse effects caused by excessive zinc ingestion or exposure. Acute toxicity may cause nausea, vomiting, abdominal pain, and diarrhea; chronic excess can impair copper absorption and lead to anemia or neurologic and immune abnormalities.

\synonyms
Hyperzincemia; Zinc Poisoning


\medterm{Zirconia} A zirconium-oxide ceramic used in dentistry for crowns, bridges, implants, and other restorations. Its strength and tissue compatibility make it useful, although restoration design, surface treatment, and clinical indication affect performance.

\synonyms
Zirconium Dioxide; Ceramic Biomaterial


\medterm{Zoledronic Acid} A potent intravenous bisphosphonate used for osteoporosis and certain cancer-related conditions, including hypercalcemia of malignancy and skeletal complications of bone metastases. Renal injury, hypocalcemia, acute-phase reactions, and rare osteonecrosis of the jaw are important risks.

\synonyms
Zoledronate; Zometa; Reclast

\medterm{Zollinger-Ellison Syndrome} A rare condition in which a hormone-producing tumor makes too much gastrin, causing the stomach to produce excess acid. It can lead to repeated ulcers, severe reflux, abdominal pain, and diarrhea.

\medterm{Zolpidem} A nonbenzodiazepine sedative-hypnotic used for short-term treatment of insomnia. It can cause next-day impairment, complex sleep behaviors, dependence, and respiratory or central-nervous-system depression, especially with other sedatives.

\synonyms
Ambien



\medterm{Zonisamide} An antiseizure medicine used mainly for focal seizures. It can cause sleepiness, thinking problems, kidney stones, reduced sweating with overheating, acid imbalance, severe rash, or other serious reactions.

\synonyms
Zonegran


\medterm{Zonule} Fine fibers that hold the eye’s lens in place. They transfer tension from the ciliary muscle so the lens changes shape and the eye can focus at different distances.

\synonyms
Zonule of Zinn; Ciliary Zonule; Suspensory Ligament of Lens

\medterm{Zoon Balanitis} A chronic benign inflammation of the glans penis, usually seen in uncircumcised men. It commonly appears as a sharply demarcated shiny orange-red or rust-colored patch; diagnosis is clinical, with biopsy used when the diagnosis is uncertain.

\synonyms
Plasma Cell Balanitis; Balanitis Circumscripta Plasmacellularis

\medterm{Zoonosis} An infection that naturally spreads between animals with backbones and people. Examples include rabies, brucellosis, leptospirosis, toxoplasmosis, and some influenza infections; spread may involve bites, food, water, droplets, or direct contact.

\synonyms
Zoonotic Infection; Zoonotic Disease

\medterm{Zoster} A painful, usually unilateral vesicular eruption caused by reactivation of latent varicella-zoster virus in a sensory nerve ganglion. It may be complicated by postherpetic neuralgia, ocular disease, or neurologic involvement.

\synonyms
Herpes Zoster; Shingles; Zona

\medterm{Zygomatic Bone} A paired facial bone that forms the prominence of the cheek and part of the lateral orbital wall and floor. It articulates with the frontal, sphenoid, temporal, and maxillary bones.

\synonyms
Cheekbone; Malar Bone

\medterm{Zygomatic Fracture} A break in the cheekbone or the group of bones joining it to the eye socket, usually caused by facial trauma. It may flatten the cheek, numb the skin below the eye, cause double vision, limit mouth opening, or injure the eye socket.

\synonyms
ZMC Fracture; Malar Fracture; Tripod Fracture

\medterm{Zygote} The first cell formed when a sperm joins an egg. It contains genetic material from both parents, divides repeatedly, and develops into an embryo.

\synonyms
Fertilized Oocyte; Conceptus

\phantomsection
\addcontentsline{toc}{chapter}{Index}
\printindex

\end{document}
